\documentclass{article}
\usepackage[utf8]{inputenc}
\usepackage[T1]{fontenc}
\usepackage{geometry}
\geometry{a4paper, margin=1in}

\begin{document}

\section*{Patanjali: Vyakaranamahabhasya}
\noindent \textbf{Patanjali:} Vyakaranamahabhasya (Mahabhasya) \\
\noindent \textbf{P\_n,n.n.n} = Pāṇini\_adhyāya,pāda.sūtra \\
\noindent \textbf{BOLD} for Kātyāyana's Vārttikas \\
\noindent \textbf{Text} converted to Unicode (UTF-8). \\
\noindent \textbf{(This} file is to be used with a UTF-8 font and your browser's VIEW configuration \\
\noindent \textbf{set} to UTF-8.) \\
\noindent \textbf{description:	multibyte} sequence: \\
\noindent \textbf{long} a	ā \\
\noindent \textbf{long} A	Ā \\
\noindent \textbf{long} i	ī \\
\noindent \textbf{long} I	Ī \\
\noindent \textbf{long} u	ū \\
\noindent \textbf{long} U	Ū \\
\noindent \textbf{vocalic} r	ṛ \\
\noindent \textbf{vocalic} R	Ṛ \\
\noindent \textbf{long} vocalic r	ṝ \\
\noindent \textbf{vocalic} l	ḷ \\
\noindent \textbf{vocalic} L	Ḷ \\
\noindent \textbf{long} vocalic l	ḹ \\
\noindent \textbf{velar} n	ṅ \\
\noindent \textbf{velar} N	Ṅ \\
\noindent \textbf{palatal} n	ñ \\
\noindent \textbf{palatal} N	Ñ \\
\noindent \textbf{retroflex} t	ṭ \\
\noindent \textbf{retroflex} T	Ṭ \\
\noindent \textbf{retroflex} d	ḍ \\
\noindent \textbf{retroflex} D	Ḍ \\
\noindent \textbf{retroflex} n	ṇ \\
\noindent \textbf{retroflex} N	Ṇ \\
\noindent \textbf{palatal} s	ś \\
\noindent \textbf{palatal} S	Ś \\
\noindent \textbf{retroflex} s	ṣ \\
\noindent \textbf{retroflex} S	Ṣ \\
\noindent anusvara	ṃ \\
\noindent visarga	ḥ \\
\noindent \textbf{long} e	ē \\
\noindent \textbf{long} o	ō \\
\noindent \textbf{l} underbar	ḻ \\
\noindent \textbf{r} underbar	ṟ \\
\noindent \textbf{n} underbar	ṉ \\
\noindent \textbf{k} underbar	ḵ \\
\noindent \textbf{t} underbar	ṯ \\
\noindent \textbf{(Pas\_1)} atha śabdānuśāsanam . \\
\noindent \textbf{(Pas\_1)} atha iti ayam śabdaḥ adhikārārthaḥ prayujyate . \\
\noindent \textbf{(Pas\_1)} śabdānuśāsanam śāstram adhikṛtam veditavyam . \\
\noindent \textbf{(Pas\_1)} keṣām śabdānām . \\
\noindent \textbf{(Pas\_1)} laukikānām vaidikānām ca . \\
\noindent \textbf{(Pas\_1)} tatra laukikāḥ tāvat : gauḥ aśvaḥ puruṣaḥ hastī śakuniḥ mṛgaḥ brāhmaṇaḥ iti . \\
\noindent \textbf{(Pas\_1)} vaidikāḥ khalu api : śam naḥ devīḥ abhiṣṭaye . \\
\noindent \textbf{(Pas\_1)} iṣe tvā ūrje tvā . \\
\noindent \textbf{(Pas\_1)} agnim īḷe purohitam . \\
\noindent \textbf{(Pas\_1)} agne ayāhi vītaye iti . \\
\noindent \textbf{(Pas\_2)} atha gauḥ iti atra kaḥ śabdaḥ . \\
\noindent \textbf{(Pas\_2)} kim yat tat sāsnālāṅgūlakakudakhuraviṣāṇi artharūpam saḥ śabdaḥ . \\
\noindent \textbf{(Pas\_2)} na iti āha . \\
\noindent \textbf{(Pas\_2)} dravyam nāma tat . \\
\noindent \textbf{(Pas\_2)} yat tarhi tat iṅgitam ceṣṭitam nimiṣitam saḥ śabdaḥ . \\
\noindent \textbf{(Pas\_2)} na iti āha . \\
\noindent \textbf{(Pas\_2)} kriyā nāma sā . \\
\noindent \textbf{(Pas\_2)} yat tarhi tat śuklaḥ nīlaḥ kṛṣṇaḥ kapilaḥ kapotaḥ iti saḥ śabdaḥ . \\
\noindent \textbf{(Pas\_2)} na iti āha . \\
\noindent \textbf{(Pas\_2)} guṇaḥ nāma saḥ . \\
\noindent \textbf{(Pas\_2)} yat tarhi tat bhinneṣu abhinnam chinneṣu acchinnam sāmānyabhūtam saḥ śabdaḥ . \\
\noindent \textbf{(Pas\_2)} na iti āha . \\
\noindent \textbf{(Pas\_2)} ākṛtiḥ nāma sā . \\
\noindent \textbf{(Pas\_2)} kaḥ tarhi śabdaḥ . \\
\noindent \textbf{(Pas\_2)} yena uccāritena sāsnālāṅgūlakakudakhuraviṣāṇinām sampratyayaḥ bhavati saḥ śabdaḥ . \\
\noindent \textbf{(Pas\_2)} atha vā pratītapadārthakaḥ loke dhvaniḥ śabdaḥ iti ucyate . \\
\noindent \textbf{(Pas\_2)} tat yathā śabdam kuru mā śabdam kārṣīḥ śabdakārī ayam māṇavakaḥ iti . \\
\noindent \textbf{(Pas\_2)} dhvanim kurvan evam ucyate . \\
\noindent \textbf{(Pas\_2)} tasmāt dhvaniḥ śabdaḥ . \\
\noindent \textbf{(Pas\_3)} kāni punaḥ śabdānuśāsanasya prayojanāni . \\
\noindent \textbf{(Pas\_3)} rakṣohāgamalaghvasandehāḥ proyojanam . \\
\noindent \textbf{(Pas\_3)} rakṣārtham vedānām adhyeyam vyākaraṇam . \\
\noindent \textbf{(Pas\_3)} lopāgamavarṇavikārajñaḥ hi samyak vedān paripālayiṣyati . \\
\noindent \textbf{(Pas\_3)} ūhaḥ khalu api. na sarvaiḥ liṅgaiḥ na ca sarvābhiḥ vibhaktibhiḥ vede mantrāḥ nigaditāḥ. te ca avaśyam yajñagatena yathāyatham vipariṇamayitavyāḥ. tān na avaiyākaraṇaḥ śaknoti yathāyatham vipariṇamayitum. tasmāt adhyeyam vyākaraṇam . \\
\noindent \textbf{(Pas\_3)} āgamaḥ khalu api . \\
\noindent \textbf{(Pas\_3)} brāhmaṇena niṣkāraṇaḥ dharmaḥ ṣaḍaṅgaḥ vedaḥ adhyeyaḥ jñeyaḥ iti . \\
\noindent \textbf{(Pas\_3)} pradhānam ca ṣaṭsu aṅgeṣu vyākaraṇam . \\
\noindent \textbf{(Pas\_3)} pradhāne ca kṛtaḥ yatnaḥ phalavān bhavati . \\
\noindent \textbf{(Pas\_3)} laghvartham ca adhyeyam vyākaraṇam. brāhmaṇena avaśyam śabdāḥ jñeyāḥ iti . \\
\noindent \textbf{(Pas\_3)} na ca antareṇa vyākaraṇam laghunā upāyena śabdāḥ śakyāḥ jñātum . \\
\noindent \textbf{(Pas\_3)} asandehārtham ca adhyeyam vyākaraṇam . \\
\noindent \textbf{(Pas\_3)} yājñikāḥ paṭhanti . \\
\noindent \textbf{(Pas\_3)} sthūlapṛṣatīm āgnivāruṇīm anaḍvāhīm ālabheta iti . \\
\noindent \textbf{(Pas\_3)} tasyām sandehaḥ sthūlā ca asau pṛṣatī ca sthūlapṛṣatī sthūlāni pṛṣanti yasyāḥ sā sthūlapṛṣatī . \\
\noindent \textbf{(Pas\_3)} tām na avaiyākaraṇaḥ svarataḥ adhyavasyati . \\
\noindent \textbf{(Pas\_3)} yadi pūrvapadaprakṛtisvaratvam tataḥ bahuvrīhiḥ. atha antodāttatvam tataḥ tatpuruṣaḥ iti . \\
\noindent \textbf{(Pas\_4.1)} imāni ca bhūyaḥ śabdānuśāsanasya prayojanāni . \\
\noindent \textbf{(Pas\_4.1)} te asurāḥ , duṣṭaḥ śabdaḥ , yat adhītam , yaḥ tu prayuṅkte , avidvāṃsaḥ , vibhaktim kurvanti , yaḥ vai imām , catvāri , uta tvaḥ , saktum iva , sārasvatīm , daśamyām putrasya , sudevaḥ asi varuṇa iti . \\
\noindent \textbf{(Pas\_4.1)} te asurāḥ . \\
\noindent \textbf{(Pas\_4.1)} te asurāḥ helayaḥ helayaḥ iti kurvantaḥ parā babhūvuḥ . \\
\noindent \textbf{(Pas\_4.1)} tasmāt brāhmaṇena na mlecchitavai na apabhāṣitavai . \\
\noindent \textbf{(Pas\_4.1)} mlecchaḥ ha vai eṣaḥ yat apaśabdaḥ . \\
\noindent \textbf{(Pas\_4.1)} mlecchāḥ mā bhūma iti adhyeyam vyākaraṇam . \\
\noindent \textbf{(Pas\_4.1)} te asurāḥ \\
\noindent \textbf{(Pas\_4.2)} duṣṭaḥ śabdaḥ . \\
\noindent \textbf{(Pas\_4.2)} duṣṭaḥ śabdaḥ svarataḥ varṇataḥ vā mithyā prayuktaḥ na tam artham āha . \\
\noindent \textbf{(Pas\_4.2)} saḥ vāgvajraḥ yajamānam hinasti yathā indraśatruḥ svarataḥ aparādhāt . \\
\noindent \textbf{(Pas\_4.2)} duṣṭān śabdān mā prayukṣmahi iti adhyeyam vyākaraṇam . \\
\noindent \textbf{(Pas\_4.2)} duṣṭaḥ śabdaḥ . \\
\noindent \textbf{(Pas\_4.3)} yat adhītam . \\
\noindent \textbf{(Pas\_4.3)} yat adhītam avijñātam nigadena eva śabdyate anagnau iva śuṣkaidhaḥ na tat jvalati karhi cit . \\
\noindent \textbf{(Pas\_4.3)} tasmāt anarthakam mā adhigīṣmahi iti adhyeyam vyākaraṇam. yat adhītam . \\
\noindent \textbf{(Pas\_4.4)} yaḥ tu prayuṅkte . \\
\noindent \textbf{(Pas\_4.4)} yaḥ tu prayuṅkte kuśalaḥ viśeṣe śabdān yathāvat vyavahārakāle saḥ anantam āpnoti jayam paratra vāgyogavit duṣyati ca apaśabdaiḥ . \\
\noindent \textbf{(Pas\_4.4)} kaḥ . \\
\noindent \textbf{(Pas\_4.4)} vāgyogavit eva . \\
\noindent \textbf{(Pas\_4.4)} kutaḥ etat . \\
\noindent \textbf{(Pas\_4.4)} yaḥ hi śabdān jānāti apaśabdān api asau jānāti . \\
\noindent \textbf{(Pas\_4.4)} yathā eva hi śabdajñāne dharmaḥ evam apaśabdajñāne api adharmaḥ . \\
\noindent \textbf{(Pas\_4.4)} atha vā bhūyān adharmaḥ prāpnoti . \\
\noindent \textbf{(Pas\_4.4)} bhūyāṃsaḥ apaśabdāḥ alpīyāṃsaḥ śabdāḥ . \\
\noindent \textbf{(Pas\_4.4)} ekaikasya hi śabdasya bahavaḥ apaśabdāḥ . \\
\noindent \textbf{(Pas\_4.4)} tat yathā gauḥ iti asya śabdasya gāvī goṇī gotā gopotalikā iti evamādayaḥ apabhraṃśāḥ . \\
\noindent \textbf{(Pas\_4.4)} atha yaḥ avāgyogavit . \\
\noindent \textbf{(Pas\_4.4)} ajñānam tasya śaraṇam . \\
\noindent \textbf{(Pas\_4.4)} na atyantāya ajñānam śaraṇam bhavitum arhati . \\
\noindent \textbf{(Pas\_4.4)} yaḥ hi ajānan vai brāhmaṇam hanyāt surām vā pibet saḥ api manye patitaḥ syāt. evam tarhi saḥ anantam āpnoti jayam paratra vāgyogavit duṣyati ca apaśabdaiḥ . \\
\noindent \textbf{(Pas\_4.4)} kaḥ. avāgyogavit eva . \\
\noindent \textbf{(Pas\_4.4)} atha yaḥ vāgyogavit . \\
\noindent \textbf{(Pas\_4.4)} vijñānam tasya śaraṇam . \\
\noindent \textbf{(Pas\_4.4)} kva punaḥ idam paṭhitam . \\
\noindent \textbf{(Pas\_4.4)} bhrājāḥ nāma ślokāḥ . \\
\noindent \textbf{(Pas\_4.4)} kim ca bhoḥ ślokāḥ api pramāṇam . \\
\noindent \textbf{(Pas\_4.4)} kim ca ataḥ . \\
\noindent \textbf{(Pas\_4.4)} yadi pramāṇam ayam api ślokaḥ pramāṇam bhavitum arhati . \\
\noindent \textbf{(Pas\_4.4)} yat udumbaravarṇānām ghaṭīnām maṇḍalam mahat pītam na svargam gamayet kim tat kratugatam nayet iti . \\
\noindent \textbf{(Pas\_4.4)} pramattagītaḥ eṣaḥ tatrabhavataḥ . \\
\noindent \textbf{(Pas\_4.4)} yaḥ tu apramattagītaḥ tat pramānam . \\
\noindent \textbf{(Pas\_4.4)} yas tu prayuṅkte . \\
\noindent \textbf{(Pas\_4.5)} avidvāṃsaḥ . \\
\noindent \textbf{(Pas\_4.5)} avidvāṃsaḥ pratyabhivāde nāmnaḥ ye plutim na viduḥ kāmam teṣu tu viproṣya strīṣu iva ayam aham vadet . \\
\noindent \textbf{(Pas\_4.5)} abhivāde strīvat mā bhūma iti adhyeyam vyākaraṇam . \\
\noindent \textbf{(Pas\_4.5)} avidvāṃsaḥ \\
\noindent \textbf{(Pas\_4.6)} vibhaktim kurvanti . \\
\noindent \textbf{(Pas\_4.6)} yājñikāḥ paṭhanti : prayājāḥ savibhaktikāḥ kāryāḥ iti . \\
\noindent \textbf{(Pas\_4.6)} na ca antareṇa vyākaraṇam prayājāḥ savibhaktikāḥ śakyāḥ kartum. vibhaktim kurvanti \\
\noindent \textbf{(Pas\_4.7)} yaḥ vai imām . \\
\noindent \textbf{(Pas\_4.7)} yaḥ vai imām padaśaḥ svaraśaḥ akṣaraśaḥ vācam vidadhāti saḥ ārtvijīnaḥ . \\
\noindent \textbf{(Pas\_4.7)} ārtvijīnāḥ syāma iti adhyeyam vyākaraṇam . \\
\noindent \textbf{(Pas\_4.7)} yaḥ vai imām . \\
\noindent \textbf{(Pas\_4.8)} catvāri . \\
\noindent \textbf{(Pas\_4.8)} catvari śṛṅgā trayaḥ asya padā dve śīrṣe sapta hastāsaḥ asya tridhā baddhaḥ vṛṣabhaḥ roravīti mahaḥ devaḥ martyān a viveśa . \\
\noindent \textbf{(Pas\_4.8)} catvāri śṛṅgāni catvāri padajātāni nāmākhyātopasarganipātāḥ ca . \\
\noindent \textbf{(Pas\_4.8)} trayaḥ asya pādāḥ trayaḥ kālāḥ bhūtabhaviṣyadvartamānāḥ . \\
\noindent \textbf{(Pas\_4.8)} dve śīrṣe dvau śabdātmānau nityaḥ kāryaḥ ca . \\
\noindent \textbf{(Pas\_4.8)} sapta hastāsaḥ asya sapta vibhaktayaḥ . \\
\noindent \textbf{(Pas\_4.8)} tridhā baddhaḥ triṣu sthāneṣu baddhaḥ urasi kaṇṭhe śirasi iti . \\
\noindent \textbf{(Pas\_4.8)} vṛṣabhaḥ varṣaṇāt . \\
\noindent \textbf{(Pas\_4.8)} roravīti śabdam karoti . \\
\noindent \textbf{(Pas\_4.8)} kutaḥ etat . \\
\noindent \textbf{(Pas\_4.8)} rautiḥ śabdakarmā . \\
\noindent \textbf{(Pas\_4.8)} mahaḥ devaḥ martyān āviveśa iti . \\
\noindent \textbf{(Pas\_4.8)} mahān devaḥ śabdaḥ . \\
\noindent \textbf{(Pas\_4.8)} martyāḥ maraṇadharmāṇaḥ manuṣyāḥ . \\
\noindent \textbf{(Pas\_4.8)} tān āviveśa . \\
\noindent \textbf{(Pas\_4.8)} mahatā devena naḥ sāmyam yathā syāt iti adhyeyam vyākaraṇam . \\
\noindent \textbf{(Pas\_4.8)} aparaḥ āha : catvari vak parimitā padani tani viduḥ brāhmaṇa ye manīṣiṇaḥ guhā trīṇi nihitā na iṅgayanti turīyam vācaḥ manuṣyāḥ vadanti . \\
\noindent \textbf{(Pas\_4.8)} catvāri vāk parimitā padāni . \\
\noindent \textbf{(Pas\_4.8)} catvāri padajātāni nāmākhyātopasarganipātāḥ ca . \\
\noindent \textbf{(Pas\_4.8)} tāni viduḥ brāhmaṇāḥ ye manīṣiṇaḥ . \\
\noindent \textbf{(Pas\_4.8)} manasaḥ īṣiṇaḥ manīṣiṇaḥ . \\
\noindent \textbf{(Pas\_4.8)} guhā trīṇi nihitā na iṅgayanti . \\
\noindent \textbf{(Pas\_4.8)} guhāyām trīṇi nihitāni na iṅgayanti . \\
\noindent \textbf{(Pas\_4.8)} na ceṣṭante . \\
\noindent \textbf{(Pas\_4.8)} na nimiṣanti iti arthaḥ . \\
\noindent \textbf{(Pas\_4.8)} turīyam vācaḥ manuṣyāḥ vadanti . \\
\noindent \textbf{(Pas\_4.8)} turīyam ha vai etat vācaḥ yat manuṣyeṣu vartate . \\
\noindent \textbf{(Pas\_4.8)} caturtham iti arthaḥ . \\
\noindent \textbf{(Pas\_4.8)} catvāri . \\
\noindent \textbf{(Pas\_4.9)} uta tvaḥ . \\
\noindent \textbf{(Pas\_4.9)} uta tvaḥ paśyan na dadarśa vacam uta tvaḥ śrṇvan na śṛṇoti enām uto tvasmai tanvam visasre jāya iva patye uśatī suvasāḥ . \\
\noindent \textbf{(Pas\_4.9)} api khalu ekaḥ paśyan api na paśyati vācam . \\
\noindent \textbf{(Pas\_4.9)} api khalu ekaḥ śrṇvan api na śrṇoti enām . \\
\noindent \textbf{(Pas\_4.9)} avidvāṃsam āha ardham . \\
\noindent \textbf{(Pas\_4.9)} uto tvasmai tanvam visasre . \\
\noindent \textbf{(Pas\_4.9)} tanum vivṛṇute . \\
\noindent \textbf{(Pas\_4.9)} jāyā iva patye uśatī suvāsāḥ . \\
\noindent \textbf{(Pas\_4.9)} tad yathā jāyā patye kāmayamānā suvāsāḥ svam ātmānam vivṛṇute evam vāk vāgvide svātmānam vivṛṇute . \\
\noindent \textbf{(Pas\_4.9)} vāk naḥ vivṛṇuyāt ātmānam iti adhyeyam vyākaraṇam . \\
\noindent \textbf{(Pas\_4.9)} uta tvaḥ . \\
\noindent \textbf{(Pas\_4.10)} saktum iva . \\
\noindent \textbf{(Pas\_4.10)} saktum iva titaunā punantaḥ yatra dhīrāḥ manasā vacam akrata atrā sakhāyaḥ sakhyani jānate bhadra eṣām lakṣmīḥ nihitā adhi vāci . \\
\noindent \textbf{(Pas\_4.10)} saktuḥ sacateḥ durdhāvaḥ bhavati . \\
\noindent \textbf{(Pas\_4.10)} kasateḥ vā viparītāt vikasito bhavati . \\
\noindent \textbf{(Pas\_4.10)} titau paripavanam bhavati tatavat vā tunnavat vā . \\
\noindent \textbf{(Pas\_4.10)} dhīrāḥ dhyānavantaḥ manasā prajñānena vācam akrata vācam akṛṣata . \\
\noindent \textbf{(Pas\_4.10)} atrā sakhāyaḥ sakhyāni jānate . \\
\noindent \textbf{(Pas\_4.10)} sāyujyāni jānate . \\
\noindent \textbf{(Pas\_4.10)} kva . \\
\noindent \textbf{(Pas\_4.10)} yaḥ eṣaḥ durghaḥ mārgaḥ ekagamyaḥ vāgviṣayaḥ . \\
\noindent \textbf{(Pas\_4.10)} ke punaḥ te . \\
\noindent \textbf{(Pas\_4.10)} vaiyākaraṇāḥ . \\
\noindent \textbf{(Pas\_4.10)} kutaḥ etat . \\
\noindent \textbf{(Pas\_4.10)} bhadrā eṣām lakṣmīḥ nihitā adhi vāci . \\
\noindent \textbf{(Pas\_4.10)} eṣām vāci bhadrā lakṣmīḥ nihitā bhavati . \\
\noindent \textbf{(Pas\_4.10)} lakṣmīḥ lakṣaṇāt bhāsanāt parivṛḍhā bhavati . \\
\noindent \textbf{(Pas\_4.10)} saktum iva . \\
\noindent \textbf{(Pas\_4.11)} sārasvatīm. yājñikāḥ paṭhanti : āhitāgniḥ apaśabdam prayujya prāyaścittīyām sārasvatīm iṣṭim nirvapet iti . \\
\noindent \textbf{(Pas\_4.11)} prāyaścittīyāḥ mā bhūma iti adhyeyam vyākaraṇam . \\
\noindent \textbf{(Pas\_4.11)} sārasvatīm . \\
\noindent \textbf{(Pas\_4.12)} daśamyām putrasya . \\
\noindent \textbf{(Pas\_4.12)} yājñikāḥ paṭhanti : daśamyuttarakālam putrasya jātasya nāma vidadhyāt ghoṣavadādi antarantaḥstham avṛddham tripuruṣānūkam anaripratiṣṭhitam . \\
\noindent \textbf{(Pas\_4.12)} tat hi pratiṣṭhitatamam bhavati . \\
\noindent \textbf{(Pas\_4.12)} dvyakṣaram caturakṣaram vā nāma kṛtam kuryāt na taddhitam iti . \\
\noindent \textbf{(Pas\_4.12)} na ca antareṇa vyākaraṇam kṛtaḥ taddhitāḥ vā śakyāḥ vijñātum . \\
\noindent \textbf{(Pas\_4.12)} daśamyām putrasya . \\
\noindent \textbf{(Pas\_4.13)} sudevaḥ asi . \\
\noindent \textbf{(Pas\_4.13)} sudevaḥ asi varuṇa yasya te sapta sindhavaḥ anukṣaranti kākudam sūrmyam suṣiram iva . \\
\noindent \textbf{(Pas\_4.13)} sudevaḥ asi varuṇa satyadevaḥ asi yasya te sapta sindhavaḥ sapta vibhaktayaḥ . \\
\noindent \textbf{(Pas\_4.13)} anukṣaranti kākudam . \\
\noindent \textbf{(Pas\_4.13)} kākudam tālu . \\
\noindent \textbf{(Pas\_4.13)} kākuḥ jihvā sā asmin udyate iti kākudam . \\
\noindent \textbf{(Pas\_4.13)} sūrmyam suṣirām iva . \\
\noindent \textbf{(Pas\_4.13)} tad yathā śobhanām ūrmīm suṣirām agniḥ antaḥ praviśya dahati evam tava sapta sindhavaḥ sapta vibhaktayaḥ tālu anukṣaranti . \\
\noindent \textbf{(Pas\_4.13)} tena asi satyadevaḥ . \\
\noindent \textbf{(Pas\_4.13)} satyadevāḥ syāma iti adhyeyam vyākaraṇam . \\
\noindent \textbf{(Pas\_4.13)} sudevaḥ asi . \\
\noindent \textbf{(Pas\_5)} kim punaḥ idam vyākaranam eva adhijigāṃsamānebhyaḥ prayojanam anvākhyāyate na punaḥ anyat api kim cit . \\
\noindent \textbf{(Pas\_5)} om iti uktvā vṛttāntaśaḥ śam iti evamādīn śabdān paṭhanti . \\
\noindent \textbf{(Pas\_5)} purākalpe etat āsīt : saṃskārottarakālam brāhmaṇāḥ vyākaraṇam sma adhīyate . \\
\noindent \textbf{(Pas\_5)} tebhyaḥ tatra sthānakaraṇānupradānajñebhyaḥ vaidikāḥ śabdāḥ upadiśyante . \\
\noindent \textbf{(Pas\_5)} tat adyatve na tathā . \\
\noindent \textbf{(Pas\_5)} vedam adhītya tvaritāḥ vaktāraḥ bhavanti : vedāt naḥ vaidikāḥ śabdāḥ siddhāḥ lokāt ca laukikāḥ . \\
\noindent \textbf{(Pas\_5)} anarthakam vyākaraṇam iti . \\
\noindent \textbf{(Pas\_5)} tebhyaḥ vipratipannabuddhibhyaḥ adhyetṛbhyaḥ ācāryaḥ idam śāstram anvācaṣṭe : imāni prayojanāni adhyeyam vyākaraṇam iti . \\
\noindent \textbf{(Pas\_6)} uktaḥ śabdaḥ . \\
\noindent \textbf{(Pas\_6)} svarūpam api uktam . \\
\noindent \textbf{(Pas\_6)} prayojanāni api uktāni . \\
\noindent \textbf{(Pas\_6)} śabdānuśāsanam idānīm kartavyam . \\
\noindent \textbf{(Pas\_6)} tat katham kartavyam . \\
\noindent \textbf{(Pas\_6)} kim śabdopadeśaḥ kartavyaḥ āhosvit apaśabdopadeśaḥ āhosvit ubhayopadeśaḥ iti . \\
\noindent \textbf{(Pas\_6)} anyataropadeśena kṛtam syāt . \\
\noindent \textbf{(Pas\_6)} tat yathā bhakṣyaniyamena abhakṣyapratiṣedho gamyate . \\
\noindent \textbf{(Pas\_6)} pañca pañcanakhāḥ bhakṣyāḥ iti ukte gamyate etat : ataḥ anye abhakṣyāḥ iti . \\
\noindent \textbf{(Pas\_6)} abhakṣyapratiṣedhena vā bhakṣyaniyamaḥ . \\
\noindent \textbf{(Pas\_6)} tat yathā abhakṣyaḥ grāmyakukkuṭaḥ abhakṣyaḥ grāmyaśūkaraḥ iti ukte gamyate etat : āraṇyaḥ bhakṣyaḥ iti . \\
\noindent \textbf{(Pas\_6)} evam iha api : yadi tāvat śabdopadeśaḥ kriyate gauḥ iti etasmin upadiṣṭe gamyate etat : gāvyādayaḥ apaśabdāḥ iti . \\
\noindent \textbf{(Pas\_6)} atha apaśabdopadeśaḥ kriyate gāvyādiṣu upadiṣṭeṣu gamyate etat : gauḥ iti eṣaḥ śabdaḥ iti . \\
\noindent \textbf{(Pas\_6)} kim punaḥ atra jyāyaḥ . \\
\noindent \textbf{(Pas\_6)} laghutvāt śabdopadeśaḥ . \\
\noindent \textbf{(Pas\_6)} laghīyān śabdopadeśaḥ garīyān apaśabdopadeśaḥ . \\
\noindent \textbf{(Pas\_6)} ekaikasya śabdasya bahavaḥ apabhraṃśāḥ . \\
\noindent \textbf{(Pas\_6)} tat yathā . \\
\noindent \textbf{(Pas\_6)} gauḥ iti asya śabdasya gāvīgoṇīgotāgopotalikādayaḥ apabhraṃśāḥ . \\
\noindent \textbf{(Pas\_6)} iṣṭānvākhyānam khalu api bhavati . \\
\noindent \textbf{(Pas\_7)} atha etasmin śabdopadeśe sati kim śabdānām pratipattau pratipadapāṭhaḥ kartavyaḥ : gauḥ aśvaḥ puruṣaḥ hastī śakuniḥ mṛgaḥ brāhmaṇaḥ iti evamādayaḥ śabdāḥ paṭhitavyāḥ . \\
\noindent \textbf{(Pas\_7)} na iti āha . \\
\noindent \textbf{(Pas\_7)} anabhyupāyaḥ eṣaḥ śabdānām pratipattau pratipadapāṭhaḥ . \\
\noindent \textbf{(Pas\_7)} evam hi śrūyate : bṛhaspatiḥ indrāya divyam varṣasahasram pratipadoktānām śabdānām śabdapārāyaṇam provāca na antam jagāma . \\
\noindent \textbf{(Pas\_7)} bṛhaspatiḥ ca pravaktā indraḥ ca adhyetā divyam varṣasahasram adhyayanakālaḥ na ca antam jagāma . \\
\noindent \textbf{(Pas\_7)} kim punaḥ adyatve . \\
\noindent \textbf{(Pas\_7)} yaḥ sarvathā ciram jīvati saḥ varṣaśatam jīvati . \\
\noindent \textbf{(Pas\_7)} caturbhiḥ ca prakāraiḥ vidyā upayuktā bhavati āgamakālena svādhyāyakālena pravacanakālena vyavahārakālena iti . \\
\noindent \textbf{(Pas\_7)} tatra ca āgamakālena eva āyuḥ paryupayuktam syāt . \\
\noindent \textbf{(Pas\_7)} tasmāt anabhyupāyaḥ śabdānām pratipattau pratipadapāṭhaḥ . \\
\noindent \textbf{(Pas\_7)} katham tarhi ime śabdāḥ pratipattavyāḥ . \\
\noindent \textbf{(Pas\_7)} kim cit sāmanyaviśeṣavat lakṣaṇam pravartyam yena alpena yatnena mahataḥ mahataḥ śabdaughān pratipadyeran . \\
\noindent \textbf{(Pas\_7)} kim punaḥ tat . \\
\noindent \textbf{(Pas\_7)} utsargāpavādau . \\
\noindent \textbf{(Pas\_7)} kaḥ cit utsargaḥ kartavyaḥ kaḥ cit apavādaḥ . \\
\noindent \textbf{(Pas\_7)} kathañjātīyakaḥ punaḥ utsargaḥ kartavyaḥ kathañjātīyakaḥ apavādaḥ . \\
\noindent \textbf{(Pas\_7)} sāmanyena utsargaḥ kartavyaḥ . \\
\noindent \textbf{(Pas\_7)} tat yathā karmaṇi aṇ . \\
\noindent \textbf{(Pas\_7)} tasya viśeṣeṇa apavādaḥ . \\
\noindent \textbf{(Pas\_7)} tat yathā . \\
\noindent \textbf{(Pas\_7)} ātaḥ anupasarge kaḥ . \\
\noindent \textbf{(Pas\_8)} kim punaḥ ākṛtiḥ padārthaḥ āhosvit dravyam . \\
\noindent \textbf{(Pas\_8)} ubhayam iti āha . \\
\noindent \textbf{(Pas\_8)} katham jñāyate . \\
\noindent \textbf{(Pas\_8)} ubhayathā hi ācāryeṇa sūtrāṇi paṭhitāni . \\
\noindent \textbf{(Pas\_8)} ākṛtim padārtham matvā jātyākhyāyām ekasmin bahuvacanam anyatarasyām iti ucyate . \\
\noindent \textbf{(Pas\_8)} dravyam padārtham matvā sarūpāṇām ekaśeṣaḥ ekavibhaktau iti ekaśeṣaḥ ārabhyate . \\
\noindent \textbf{(Pas\_9)} kim punaḥ nityaḥ śabdaḥ āhosvit kāryaḥ . \\
\noindent \textbf{(Pas\_9)} saṅgrahe etat prādhānyena parīkṣitam nityaḥ vā syāt kāryaḥ vā iti . \\
\noindent \textbf{(Pas\_9)} tatra uktāḥ doṣāḥ prayojanāni api uktāni . \\
\noindent \textbf{(Pas\_9)} tatra tu eṣaḥ nirṇayaḥ yadi eva nityaḥ atha api kāryaḥ ubhayathā api lakṣaṇam pravartyam iti . \\
\noindent \textbf{(Pas\_10.1)} katham punaḥ idam bhagavataḥ pāṇineḥ ācāryasya lakṣaṇam pravṛttam . \\
\noindent \textbf{(Pas\_10.1)} siddhe śabdārthasambandhe . \\
\noindent \textbf{(Pas\_10.1)} siddhe śabde arthe sambandhe ca iti . \\
\noindent \textbf{(Pas\_10.1)} atha siddhaśabdasya kaḥ padārthaḥ . \\
\noindent \textbf{(Pas\_10.1)} nityaparyāyavācī siddhaśabdaḥ . \\
\noindent \textbf{(Pas\_10.1)} katham jñāyate . \\
\noindent \textbf{(Pas\_10.1)} yat kūṭastheṣu avicāliṣu bhāveṣu vartate . \\
\noindent \textbf{(Pas\_10.1)} tat yathā siddhā dyauḥ , siddhā pṛthivī siddham ākāśam iti . \\
\noindent \textbf{(Pas\_10.1)} nanu ca bhoḥ kāryeṣu api vartate . \\
\noindent \textbf{(Pas\_10.1)} tat yathā siddhaḥ odanaḥ , siddhaḥ sūpaḥ siddhā yavāgūḥ iti . \\
\noindent \textbf{(Pas\_10.1)} yāvatā kāryeṣu api vartate tatra kutaḥ etat nityaparyāyavācinaḥ grahaṇam na punaḥ kārye yaḥ siddhaśabdaḥ iti . \\
\noindent \textbf{(Pas\_10.1)} saṅgrahe tāvat kāryapratidvandvibhāvāt manyāmahe nityaparyāyavācinaḥ grahaṇam iti . \\
\noindent \textbf{(Pas\_10.1)} iha api tat eva . \\
\noindent \textbf{(Pas\_10.1)} atha vā santi ekapadāni api avadhāraṇāni . \\
\noindent \textbf{(Pas\_10.1)} tat yathā : abbhakṣaḥ vāyubhakṣaḥ iti . \\
\noindent \textbf{(Pas\_10.1)} apaḥ eva bhakṣayati vāyum eva bhakṣayati iti gamyate . \\
\noindent \textbf{(Pas\_10.1)} evam iha api siddhaḥ eva na sādhyaḥ iti . \\
\noindent \textbf{(Pas\_10.1)} atha vā pūrvapadalopaḥ atra draṣṭavyaḥ : atyantasiddhaḥ siddhaḥ iti . \\
\noindent \textbf{(Pas\_10.1)} tat yathā devadattaḥ dattaḥ , satyabhāmā bhāmā iti . \\
\noindent \textbf{(Pas\_10.1)} atha vā vyākhyānataḥ viśeṣapratipattiḥ na hi sandehāt alakṣaṇam iti nityaparyāyavācinaḥ grahaṇam iti vyākhyāsyāmaḥ . \\
\noindent \textbf{(Pas\_10.1)} kim punaḥ anena varṇyena . \\
\noindent \textbf{(Pas\_10.1)} kim na mahatā kaṇṭhena nityaśabdaḥ eva upāttaḥ yasmin upādīyamāne asandehaḥ syāt . \\
\noindent \textbf{(Pas\_10.1)} maṅgalārtham . \\
\noindent \textbf{(Pas\_10.1)} māṅgalikaḥ ācāryaḥ mahataḥ śāstraughasya maṅgalārtham siddhaśabdam āditaḥ prayuṅkte . \\
\noindent \textbf{(Pas\_10.1)} maṅgalādīni hi śāstrāṇi prathante vīrapuruṣakāṇi ca bhavanti āyuṣmatpuruṣakāṇi ca . \\
\noindent \textbf{(Pas\_10.1)} adhyetāraḥ ca siddhārthāḥ yathā syuḥ iti . \\
\noindent \textbf{(Pas\_10.1)} ayam khalu api nityaśabdaḥ na avaśyam kūṭastheṣu avicāliṣu bhāveṣu vartate . \\
\noindent \textbf{(Pas\_10.1)} kim tarhi . \\
\noindent \textbf{(Pas\_10.1)} ābhīkṣṇye api vartate . \\
\noindent \textbf{(Pas\_10.1)} tat yathā nityaprahasitaḥ nityaprajalpitaḥ iti . \\
\noindent \textbf{(Pas\_10.1)} yāvatā ābhīkṣṇye api vartate tatra api anyena eva arthaḥ syāt vyākhyānataḥ viśeṣapratipattiḥ na hi sandehāt alakṣaṇam iti . \\
\noindent \textbf{(Pas\_10.1)} paśyati tu ācāryaḥ maṅgalārthaḥ ca eva siddhaśabdaḥ āditaḥ prayuktaḥ bhaviṣyati śakṣyāmi ca enam nityaparyāyavācinam varṇayitum iti . \\
\noindent \textbf{(Pas\_10.1)} ataḥ siddhaśabdaḥ eva upāttaḥ na nityaśabdaḥ \\
\noindent \textbf{(Pas\_10.2)} atha kam punaḥ padārtham matvā eṣaḥ vigrahaḥ kriyate siddhe śabde arthe sambandhe ca iti . \\
\noindent \textbf{(Pas\_10.2)} ākṛtim iti āha . \\
\noindent \textbf{(Pas\_10.2)} kutaḥ etat . \\
\noindent \textbf{(Pas\_10.2)} ākṛtiḥ hi nityā . \\
\noindent \textbf{(Pas\_10.2)} dravyam anityam . \\
\noindent \textbf{(Pas\_10.2)} atha dravye padārthe katham vigrahaḥ kartavyaḥ . \\
\noindent \textbf{(Pas\_10.2)} siddhe śabde arthasambandhe ca iti . \\
\noindent \textbf{(Pas\_10.2)} nityaḥ hi arthavatām arthaiḥ abhisambandhaḥ . \\
\noindent \textbf{(Pas\_10.2)} atha vā dravye eva padārthe eṣaḥ vigrahaḥ nyāyyaḥ siddhe śabde arthe sambandhe ca iti. dravyam hi nityam ākṛtiḥ anityā . \\
\noindent \textbf{(Pas\_10.2)} katham jñāyate . \\
\noindent \textbf{(Pas\_10.2)} evam hi dṛśyate loke . \\
\noindent \textbf{(Pas\_10.2)} mṛt kayā cit ākṛtyā yuktā piṇḍaḥ bhavati . \\
\noindent \textbf{(Pas\_10.2)} piṇḍākṛtim upamṛdya ghaṭikāḥ kiryante . \\
\noindent \textbf{(Pas\_10.2)} ghaṭikākṛtim upamṛdya kuṇḍikāḥ kriyante . \\
\noindent \textbf{(Pas\_10.2)} tathā suvarṇam kayā cit ākṛtyā yuktam piṇḍaḥ bhavati . \\
\noindent \textbf{(Pas\_10.2)} piṇḍākṛtim upamṛdya rucakāḥ kriyante . \\
\noindent \textbf{(Pas\_10.2)} rucakākṛtim upamṛdya kaṭakāḥ kriyante . \\
\noindent \textbf{(Pas\_10.2)} kaṭakākṛtim upmṛdya svastikāḥ kriyante . \\
\noindent \textbf{(Pas\_10.2)} punaḥ āvṛttaḥ suvarṇapiṇḍaḥ punaḥ aparayā ākṛtyā yuktaḥ khadirāgārasavarṇe kuṇḍale bhavataḥ . \\
\noindent \textbf{(Pas\_10.2)} ākṛtiḥ anyā ca anyā ca bhavati dravyam punaḥ tad eva . \\
\noindent \textbf{(Pas\_10.2)} ākṛtyupamardena dravyam eva avaśiṣyate . \\
\noindent \textbf{(Pas\_10.2)} ākṛtau api padārthe eṣaḥ vigrahaḥ nyāyyaḥ siddhe śabde arthe sambandhe ca iti . \\
\noindent \textbf{(Pas\_10.2)} nanu ca uktam ākṛtiḥ anityā iti . \\
\noindent \textbf{(Pas\_10.2)} na etat asti . \\
\noindent \textbf{(Pas\_10.2)} nityā ākṛtiḥ . \\
\noindent \textbf{(Pas\_10.2)} katham . \\
\noindent \textbf{(Pas\_10.2)} na kva cit uparatā iti kṛtvā sarvatra uparatā bhavati . \\
\noindent \textbf{(Pas\_10.2)} dravyāntarasthā tu upalabhyate . \\
\noindent \textbf{(Pas\_10.2)} atha vā na idam eva nityalakṣaṇam dhruvam kūṭastham avicāli anapāyopajanavikāri anutpatti avṛddhi avyayayogi iti tan nityam iti . \\
\noindent \textbf{(Pas\_10.2)} tat api nityam yasmin tattvam na vihanyate . \\
\noindent \textbf{(Pas\_10.2)} kim punaḥ tattvam . \\
\noindent \textbf{(Pas\_10.2)} tadbhāvaḥ tattvam . \\
\noindent \textbf{(Pas\_10.2)} ākṛtau api tattvam na vihanyate . \\
\noindent \textbf{(Pas\_10.2)} atha vā kim naḥ etena idam nityam idam anityam iti . \\
\noindent \textbf{(Pas\_10.2)} yat nityam tam padārtham matvā eṣaḥ vigrahaḥ kriyate siddhe śabde arthe sambandhe ca iti . \\
\noindent \textbf{(Pas\_10.2)} katham punaḥ jñāyate siddhaḥ śabdaḥ arthaḥ sambandhaḥ ca iti . \\
\noindent \textbf{(Pas\_10.2)} lokataḥ . \\
\noindent \textbf{(Pas\_10.2)} yat loke artham upādāya śabdān prayuñjate . \\
\noindent \textbf{(Pas\_10.2)} na eṣām nirvṛttau yatnam kurvanti . \\
\noindent \textbf{(Pas\_10.2)} ye punaḥ kāryāḥ bhāvāḥ nirvṛttau tāvat teṣām yatnaḥ kriyate . \\
\noindent \textbf{(Pas\_10.2)} tat yathā . \\
\noindent \textbf{(Pas\_10.2)} ghaṭena kāryam kariṣyan kumbhakārakulam gatvā āha kuru ghaṭam . \\
\noindent \textbf{(Pas\_10.2)} kāryam anena kariṣyāmi iti . \\
\noindent \textbf{(Pas\_10.2)} na tadvat śabdān prayokṣyamāṇaḥ vaiyākaraṇakulam gatvā āha . \\
\noindent \textbf{(Pas\_10.2)} kuru śabdān . \\
\noindent \textbf{(Pas\_10.2)} prayokṣye iti . \\
\noindent \textbf{(Pas\_10.2)} tāvati eva artham upādāya śabdān prayuñjate . \\
\noindent \textbf{(Pas\_11)} yadi tarhi lokaḥ eṣu pramāṇam kim śāstreṇa kriyate . \\
\noindent \textbf{(Pas\_11)} lokataḥ arthaprayukte śabdaprayoge śāstreṇa dharmaniyamaḥ . \\
\noindent \textbf{(Pas\_11)} lokataḥ arthaprayukte śabdaprayoge śāstreṇa dharmaniyamaḥ kriyate . \\
\noindent \textbf{(Pas\_11)} kim idam dharmaniyamaḥ iti . \\
\noindent \textbf{(Pas\_11)} dharmāya niyamaḥ dharmaniyamaḥ dharmārthaḥ vā niyamaḥ dharmaniyamaḥ dharmaprayojanaḥ vā niyamaḥ dharmaniyamaḥ . \\
\noindent \textbf{(Pas\_11)} yathā laukikavaidikeṣu . \\
\noindent \textbf{(Pas\_11)} priyataddhitāḥ dākṣiṇātyāḥ . \\
\noindent \textbf{(Pas\_11)} yathā loke vede ca iti prayoktavye yathā laukikavaidikeṣu iti prayuñjate . \\
\noindent \textbf{(Pas\_11)} atha vā yuktaḥ eva taddhitārthaḥ . \\
\noindent \textbf{(Pas\_11)} yathā laukikeṣu vaidikeṣu ca kṛtānteṣu . \\
\noindent \textbf{(Pas\_11)} loke tāvat abhakṣyaḥ grāmyakukkuṭaḥ abhakṣyaḥ grāmyaśūkaraḥ iti ucyate . \\
\noindent \textbf{(Pas\_11)} bhakṣyam ca nāma kṣutpratīghātārtham upādīyate . \\
\noindent \textbf{(Pas\_11)} śakyam ca anena śvamāṃsādibhiḥ api kṣut pratihantum . \\
\noindent \textbf{(Pas\_11)} tatra niyamaḥ kriyate . \\
\noindent \textbf{(Pas\_11)} idam bhakṣyam . \\
\noindent \textbf{(Pas\_11)} idam abhakṣyam iti . \\
\noindent \textbf{(Pas\_11)} tathā khedāt strīṣu pravṛttiḥ bhavati . \\
\noindent \textbf{(Pas\_11)} samānaḥ ca khedavigamaḥ gamyāyām ca agamyāyām ca . \\
\noindent \textbf{(Pas\_11)} tatra niyamaḥ kriyate : iyam gamyā iyam agamyā iti . \\
\noindent \textbf{(Pas\_11)} vede khalu api payovrataḥ brāhmaṇaḥ yavāgūvrataḥ rājanyaḥ āmikṣāvrataḥ vaiśyaḥ iti ucyate . \\
\noindent \textbf{(Pas\_11)} vratam ca nāma abhyavahārārtham upādīyate . \\
\noindent \textbf{(Pas\_11)} śakyam ca anena śālimāṃsādīni api vratayitum . \\
\noindent \textbf{(Pas\_11)} tatra niyamaḥ kriyate . \\
\noindent \textbf{(Pas\_11)} tathā bailvaḥ khādiraḥ vā yūpaḥ syāt iti ucyate . \\
\noindent \textbf{(Pas\_11)} yūpaḥ ca nāma paśvanubandhārtham upādīyate . \\
\noindent \textbf{(Pas\_11)} śakyam ca anena kim cit eva kāṣṭham ucchritya anucchritya vā paśuḥ anubanddhum . \\
\noindent \textbf{(Pas\_11)} tatra niyamaḥ kriyate . \\
\noindent \textbf{(Pas\_11)} tathā agnau kapālāni adhiśritya abhimantrayate . \\
\noindent \textbf{(Pas\_11)} bhṛgūṇām aṅgirasām gharmasya tapasā tapyadhvam iti . \\
\noindent \textbf{(Pas\_11)} antareṇa api mantram agniḥ dahanakarmā kapālāni santāpayati . \\
\noindent \textbf{(Pas\_11)} tatra niyamaḥ kriyate . \\
\noindent \textbf{(Pas\_11)} evam kriyamāṇam abhyudayakāri bhavati iti . \\
\noindent \textbf{(Pas\_11)} evam iha api samānāyām arthagatau śabdena ca apaśabdena ca dharmaniyamaḥ kriyate . \\
\noindent \textbf{(Pas\_11)} śabdena eva arthaḥ abhidheyaḥ na apaśabdena iti . \\
\noindent \textbf{(Pas\_11)} evam kriyamāṇam abhyudayakāri bhavati iti . \\
\noindent \textbf{(Pas\_12)} asti aprayuktaḥ . \\
\noindent \textbf{(Pas\_12)} santi vai śabdāḥ aprayuktāḥ . \\
\noindent \textbf{(Pas\_12)} tat yathā ūṣa tera cakra peca iti . \\
\noindent \textbf{(Pas\_12)} kim ataḥ yat santi aprayuktāḥ . \\
\noindent \textbf{(Pas\_12)} prayogāt hi bhavān śabdānām sādhutvam adhyavasyati . \\
\noindent \textbf{(Pas\_12)} ye idānīm aprayuktāḥ na amī sādhavaḥ syuḥ . \\
\noindent \textbf{(Pas\_12)} idam vipratiṣiddham yat ucyate santi vai śabdāḥ aprayuktāḥ iti . \\
\noindent \textbf{(Pas\_12)} yadi santi na aprayuktāḥ . \\
\noindent \textbf{(Pas\_12)} atha aprayuktāḥ na santi . \\
\noindent \textbf{(Pas\_12)} santi ca aprayuktāḥ ca iti vipratiṣiddham . \\
\noindent \textbf{(Pas\_12)} prayuñjānaḥ eva khalu bhavān āha santi śabdāḥ aprayuktāḥ iti . \\
\noindent \textbf{(Pas\_12)} kaḥ ca idānīm anyaḥ bhavajjātīyakaḥ puruṣaḥ śabdānām prayoge sādhuḥ syāt . \\
\noindent \textbf{(Pas\_12)} na etat vipratiṣiddham . \\
\noindent \textbf{(Pas\_12)} santi iti tāvat brūmaḥ yat etān śāstravidaḥ śāstreṇa anuvidadhate . \\
\noindent \textbf{(Pas\_12)} aprayuktāḥ iti brūmaḥ yat loke aprayuktāḥ iti . \\
\noindent \textbf{(Pas\_12)} yat api ucyate kaḥ ca idānīm anyaḥ bhavajjātīyakaḥ puruṣaḥ śabdānām prayoge sādhuḥ syāt iti . \\
\noindent \textbf{(Pas\_12)} na brūmaḥ asmābhiḥ aprayuktāḥ iti . \\
\noindent \textbf{(Pas\_12)} kim tarhi . \\
\noindent \textbf{(Pas\_12)} loke aprayuktāḥ iti . \\
\noindent \textbf{(Pas\_12)} nanu ca bhavān api abhyantaraḥ loke . \\
\noindent \textbf{(Pas\_12)} abhyantaraḥ aham loke na tu aham lokaḥ . \\
\noindent \textbf{(Pas\_12)} asti aprayuktaḥ iti cet na arthe śabdaprayogāt . \\
\noindent \textbf{(Pas\_12)} asti aprayuktaḥ iti cet tat na . \\
\noindent \textbf{(Pas\_12)} kim kāraṇam . \\
\noindent \textbf{(Pas\_12)} arthe śabdaprayogāt . \\
\noindent \textbf{(Pas\_12)} arthe śabdāḥ prayujyante . \\
\noindent \textbf{(Pas\_12)} santi ca eṣām śabdānām arthāḥ yeṣu artheṣu prayujyante . \\
\noindent \textbf{(Pas\_12)} aprayogaḥ prayogānyatvāt . \\
\noindent \textbf{(Pas\_12)} aprayogaḥ khalu eṣāṃ śabdānām nyāyyaḥ . \\
\noindent \textbf{(Pas\_12)} kutaḥ . \\
\noindent \textbf{(Pas\_12)} prayogānyatvāt . \\
\noindent \textbf{(Pas\_12)} yat eteṣām śabdānām arthe anyān śabdān prayuñjate . \\
\noindent \textbf{(Pas\_12)} tat yathā . \\
\noindent \textbf{(Pas\_12)} ūṣa iti etasya śabdasya arthe kva yūyam uṣitāḥ . \\
\noindent \textbf{(Pas\_12)} tera iti asya arthe kim yūyam tīrṇāḥ . \\
\noindent \textbf{(Pas\_12)} cakra iti asya arthe kim yūyam kṛtavantaḥ . \\
\noindent \textbf{(Pas\_12)} peca iti asya arthe kim yūyam pakvavantaḥ iti . \\
\noindent \textbf{(Pas\_12)} aprayukte dīrghasattravat . \\
\noindent \textbf{(Pas\_12)} yadi api aprayuktāḥ avaśyam dīrghasattravat lakṣaṇena anuvidheyāḥ . \\
\noindent \textbf{(Pas\_12)} tat yathā . \\
\noindent \textbf{(Pas\_12)} dīrghasattrāṇi vārṣaśatikāni vārṣasahasrikāṇi ca . \\
\noindent \textbf{(Pas\_12)} na ca adyatve kaḥ cit api vyavaharati . \\
\noindent \textbf{(Pas\_12)} kevalam ṛṣisampradāyaḥ dharmaḥ iti kṛtvā yājñikāḥ śāstreṇa anuvidadhate . \\
\noindent \textbf{(Pas\_12)} sarve deśāntare . \\
\noindent \textbf{(Pas\_12)} sarve khalu api ete śabdāḥ deśāntare prayujyante . \\
\noindent \textbf{(Pas\_12)} na ca ete upalabhyante . \\
\noindent \textbf{(Pas\_12)} upalabdhau yatnaḥ kriyatām . \\
\noindent \textbf{(Pas\_12)} mahān hi śabdasya prayogaviṣayaḥ . \\
\noindent \textbf{(Pas\_12)} saptadvīpā vasumatī trayaḥ lokāḥ catvāraḥ vedāḥ sāṅgāḥ sarahasyāḥ bahudhā vibhinnāḥ ekaśatam adhvaryuśākhāḥ sahasravartmā sāmavedaḥ ekaviṃsatidhā bāhvṛcyam navadhā ātharvaṇaḥ vedaḥ vākovākyam itihāsaḥ purāṇam vaidyakam iti etāvān śabdasya prayogaviṣayaḥ . \\
\noindent \textbf{(Pas\_12)} etāvantam śabdasya prayogaviṣayam ananuniśamya santi aprayuktāḥ iti vacanam kevalam sāhasamātram . \\
\noindent \textbf{(Pas\_12)} etasmin atimahati śabdasya prayogaviṣaye te te śabdāḥ tatra tatra niyataviṣayāḥ dṛśyante . \\
\noindent \textbf{(Pas\_12)} tat yathā . \\
\noindent \textbf{(Pas\_12)} śavatiḥ gatikarmā kambojeṣu eva bhāṣitaḥ bhavati . \\
\noindent \textbf{(Pas\_12)} vikāre enam āryāḥ bhāṣante śavaḥ iti . \\
\noindent \textbf{(Pas\_12)} hammatiḥ surāṣṭreṣu raṃhatiḥ prācyamadhyeṣu gamim eva tu āryāḥ prayuñjate . \\
\noindent \textbf{(Pas\_12)} dātiḥ lavanārthe prācyeṣu dātram udīcyeṣu . \\
\noindent \textbf{(Pas\_12)} ye ca api ete bhavataḥ aprayuktāḥ abhimatāḥ śabdāḥ eteṣām api prayogaḥ dṛśyate . \\
\noindent \textbf{(Pas\_12)} kva . \\
\noindent \textbf{(Pas\_12)} vede . \\
\noindent \textbf{(Pas\_12)} yat vaḥ revatīḥ revatyam tat ūṣa . \\
\noindent \textbf{(Pas\_12)} yat me naraḥ śrutyam brahma cakra . \\
\noindent \textbf{(Pas\_12)} yatra naḥ cakra jarasam tanunām iti . \\
\noindent \textbf{(Pas\_13)} kim punaḥ śabdasya jñāne dharmaḥ āhosvit prayoge . \\
\noindent \textbf{(Pas\_13)} kaḥ ca atra viśeṣaḥ . \\
\noindent \textbf{(Pas\_13)} jñāne dharmaḥ iti cet tathā adharmaḥ . \\
\noindent \textbf{(Pas\_13)} jñāne dharmaḥ iti cet tathā adharmaḥ prāpnoti . \\
\noindent \textbf{(Pas\_13)} yaḥ hi śabdān jānāti apaśabdān api asau jānāti . \\
\noindent \textbf{(Pas\_13)} yathā eva śabdajñāne dharmaḥ evam apaśabdajñāne api adharmaḥ . \\
\noindent \textbf{(Pas\_13)} atha vā bhūyān adharmaḥ prāpnoti .bhūyāṃsaḥ apaśabdāḥ alpīyāṃsaḥ śabdāḥ . \\
\noindent \textbf{(Pas\_13)} ekaikasya śabdasya bahavaḥ apabhraṃśāḥ . \\
\noindent \textbf{(Pas\_13)} tat yathā . \\
\noindent \textbf{(Pas\_13)} gauḥ iti asya gāvī goṇī gotā gopotalikā iti evamādayaḥ apabhraṃśāḥ . \\
\noindent \textbf{(Pas\_13)} ācāre niyamaḥ . \\
\noindent \textbf{(Pas\_13)} ācāre punaḥ ṛṣiḥ niyamam vedayate . \\
\noindent \textbf{(Pas\_13)} te asurāḥ helayaḥ helayaḥ iti kurvantaḥ parābabhūvuḥ iti . \\
\noindent \textbf{(Pas\_13)} astu tarhi prayoge . \\
\noindent \textbf{(Pas\_13)} prayoge sarvalokasya . \\
\noindent \textbf{(Pas\_13)} yadi prayoge dharmaḥ sarvaḥ lokaḥ abhyudayena yujyeta . \\
\noindent \textbf{(Pas\_13)} kaḥ ca idānīm bhavataḥ matsaraḥ yadi sarvaḥ lokaḥ abhyudayena yujyeta . \\
\noindent \textbf{(Pas\_13)} na khalu kaḥ cit matsaraḥ . \\
\noindent \textbf{(Pas\_13)} prayatnānarthakyam tu bhavati . \\
\noindent \textbf{(Pas\_13)} phalavatā ca nāma prayatnena bhavitavyam na ca prayatnaḥ phalāt vyatirecyaḥ . \\
\noindent \textbf{(Pas\_13)} nanu ca ye kṛtaprayatnāḥ te sādhīyaḥ śabdān prayokṣyante . \\
\noindent \textbf{(Pas\_13)} te eva sādhīyaḥ abhyudayena yokṣyante . \\
\noindent \textbf{(Pas\_13)} vyatirekaḥ api vai lakṣyate . \\
\noindent \textbf{(Pas\_13)} dṛśyante hi kṛtaprayatnāḥ ca apravīṇāḥ akṛtaprayatnāḥ ca pravīṇāḥ . \\
\noindent \textbf{(Pas\_13)} tatra phalavyatirekaḥ api syāt . \\
\noindent \textbf{(Pas\_13)} evam tarhi na api jñāne eva dharmaḥ na api prayoge eva . \\
\noindent \textbf{(Pas\_13)} kim tarhi śāstrapūrvake prayoge abhyudayaḥ tat tulyam vedaśabdena . \\
\noindent \textbf{(Pas\_13)} śāstrapūrvakam yaḥ śabdān prayuṅkte saḥ abhyudayena yujyate . \\
\noindent \textbf{(Pas\_13)} tat tulyam vedaśabdena . \\
\noindent \textbf{(Pas\_13)} vedaśabdāḥ api evam abhivadanti . \\
\noindent \textbf{(Pas\_13)} yaḥ agniṣṭomena yajate yaḥ u ca enam evam veda . \\
\noindent \textbf{(Pas\_13)} yaḥ agnim nāciketam cinute yaḥ u ca enam evam veda . \\
\noindent \textbf{(Pas\_13)} aparaḥ āha : tat tulyam vedaśabdena iti . \\
\noindent \textbf{(Pas\_13)} yathā vedaśabdāḥ niyamapūrvam adhītāḥ phalavantaḥ bhavanti evam yaḥ śāstrapūrvakam śabdān prayuṅkte saḥ abhyudayena yujyate iti . \\
\noindent \textbf{(Pas\_13)} atha vā punaḥ astu jñāne eva dharmaḥ iti . \\
\noindent \textbf{(Pas\_13)} nanu ca uktam jñāne dharmaḥ iti cet tathā adharmaḥ iti . \\
\noindent \textbf{(Pas\_13)} na eṣaḥ doṣaḥ . \\
\noindent \textbf{(Pas\_13)} śabdapramāṇakāḥ vayam . \\
\noindent \textbf{(Pas\_13)} yat śabdaḥ āha tat asmākam pramāṇam . \\
\noindent \textbf{(Pas\_13)} śabdaḥ ca śabdajñāne dharmam āha na apaśabdajñāne adharmam . \\
\noindent \textbf{(Pas\_13)} yat ca punaḥ aśiṣṭāpratiṣiddham na eva tat doṣāya bhavati na abhyudayāya . \\
\noindent \textbf{(Pas\_13)} tat yathā . \\
\noindent \textbf{(Pas\_13)} hikkitahasitakaṇḍūyitāni na eva doṣāya bhavanti na api abhyudayāya . \\
\noindent \textbf{(Pas\_13)} atha vā abhyupāyaḥ eva apaśabdajñānam śabdajñāne . \\
\noindent \textbf{(Pas\_13)} yaḥ apaśabdān jānāti śabdān api asau jānāti . \\
\noindent \textbf{(Pas\_13)} tat evam jñāne dharmaḥ iti bruvataḥ arthāt āpannam bhavati apaśabdajñānapūrvake śabdajñāne dharmaḥ iti . \\
\noindent \textbf{(Pas\_13)} atha vā kūpakhānakavat etat bhavati . \\
\noindent \textbf{(Pas\_13)} tat yathā kūpakhānakaḥ khanan yadi api mṛdā pāṃsubhiḥ ca avakīrṇaḥ bhavati saḥ apsu sañjātāsu tataḥ eva tam guṇam āsādayati yena saḥ ca doṣaḥ nirhaṇyate bhūyasā ca abhyudayena yogaḥ bhavati evam iha api yadi api apaśabdajñāne adharmaḥ tathā api yaḥ tu asau śabdajñāne dharmaḥ tena saḥ ca doṣaḥ nirghāniṣyate bhūyasā ca abhyudayena yogaḥ bhaviṣyati . \\
\noindent \textbf{(Pas\_13)} yat api ucyate ācāre niyamaḥ iti yājñe karmaṇi saḥ niyamaḥ . \\
\noindent \textbf{(Pas\_13)} evam hi śrūyate . \\
\noindent \textbf{(Pas\_13)} yarvāṇaḥ tarvāṇaḥ nāma ṛṣayaḥ babhūvuḥ pratyakṣadharmāṇaḥ parāparajñāḥ viditaveditavyāḥ adhigatayāthātathyāḥ . \\
\noindent \textbf{(Pas\_13)} te tatrabhavantaḥ yat vā naḥ tat vā naḥ iti proyoktavye yar vā ṇaḥ tar vā ṇaḥ iti prayuñjate yājñe punaḥ karmaṇi na apabhāṣante . \\
\noindent \textbf{(Pas\_13)} taiḥ punaḥ asuraiḥ yājñe karmaṇi apabhāṣitam . \\
\noindent \textbf{(Pas\_13)} tataḥ te parābabhūtāḥ . \\
\noindent \textbf{(Pas\_14)} atha vyākaraṇam iti asya śabdasya kaḥ padārthaḥ . \\
\noindent \textbf{(Pas\_14)} sūtram . \\
\noindent \textbf{(Pas\_14)} sūtre vyākaraṇe ṣaṣṭhyarthaḥ anupapannaḥ . \\
\noindent \textbf{(Pas\_14)} sūtre vyākaraṇe ṣaṣṭhyarthaḥ na upapadyate vyākaraṇasya sūtram iti . \\
\noindent \textbf{(Pas\_14)} kim hi tat anyat sūtrāt vyākaraṇam yasya adaḥ sūtram syāt . \\
\noindent \textbf{(Pas\_14)} śabdāpratipattiḥ . \\
\noindent \textbf{(Pas\_14)} śabdānām ca apratipattiḥ prāpnoti vyākaraṇāt śabdān pratipadyāmahe iti . \\
\noindent \textbf{(Pas\_14)} na hi sūtrataḥ eva śabdān pratipadyante . \\
\noindent \textbf{(Pas\_14)} kim tarhi . \\
\noindent \textbf{(Pas\_14)} vyākhyānataḥ ca . \\
\noindent \textbf{(Pas\_14)} nanu ca tat eva sūtram vigṛhītam vyākhyānam bhavati . \\
\noindent \textbf{(Pas\_14)} na kevalāni carcāpadāni vyākhyanam vṛddhiḥ āt aic iti . \\
\noindent \textbf{(Pas\_14)} kim tarhi . \\
\noindent \textbf{(Pas\_14)} udāharaṇam pratyudāharaṇam vākyādhyāhāraḥ iti etat samuditam vyākhyānam bhavati . \\
\noindent \textbf{(Pas\_14)} evam tarhi śabdaḥ . \\
\noindent \textbf{(Pas\_14)} śabde lyuḍarthaḥ . \\
\noindent \textbf{(Pas\_14)} yadi śabdaḥ vyākaraṇam lyuḍarthaḥ na upapadyate vyākriyate anena iti vyākaraṇam . \\
\noindent \textbf{(Pas\_14)} na hi śabdena kim cit vyākriyate . \\
\noindent \textbf{(Pas\_14)} kena tarhi . \\
\noindent \textbf{(Pas\_14)} sūtreṇa . \\
\noindent \textbf{(Pas\_14)} bhave . \\
\noindent \textbf{(Pas\_14)} bhave ca taddhitaḥ na upapadyate vyākaraṇe bhavaḥ yogaḥ vaiyākaraṇaḥ iti . \\
\noindent \textbf{(Pas\_14)} na hi śabde bhavaḥ yogaḥ . \\
\noindent \textbf{(Pas\_14)} kva tarhi . \\
\noindent \textbf{(Pas\_14)} sūtre . \\
\noindent \textbf{(Pas\_14)} proktādayaḥ ca taddhitāḥ . \\
\noindent \textbf{(Pas\_14)} proktādayaḥ ca taddhitāḥ na upapadyante pāṇininā proktam pāṇinīyam , āpiśalam , kāśakṛtsnam iti . \\
\noindent \textbf{(Pas\_14)} na hi pāṇininā śabdāḥ proktāḥ . \\
\noindent \textbf{(Pas\_14)} kim tarhi . \\
\noindent \textbf{(Pas\_14)} sūtram . \\
\noindent \textbf{(Pas\_14)} kimartham idam ubhayam ucyate bhave proktādayaḥ ca taddhitāḥ iti na proktādayaḥ ca taddhitāḥ iti eva bhave api taddhitaḥ coditaḥ syāt . \\
\noindent \textbf{(Pas\_14)} purastāt idam ācāryeṇa dṛṣṭam bhave taddhitaḥ iti tat paṭhitam . \\
\noindent \textbf{(Pas\_14)} tataḥ uttarakālam idam dṛṣṭam proktādayaḥ ca taddhitāḥ iti tat api paṭhitam . \\
\noindent \textbf{(Pas\_14)} na ca idānīm ācāryāḥ sūtrāṇi kṛtvā nivartayanti . \\
\noindent \textbf{(Pas\_14)} ayam tāvat adoṣaḥ yat ucyate śabde lyuḍarthaḥ iti . \\
\noindent \textbf{(Pas\_14)} na avaśyam karaṇādhikaraṇayoḥ eva lyuṭ vidhīyate kim tarhi anyeṣu api kārakeṣu kṛtyalyuṭaḥ bahulam iti . \\
\noindent \textbf{(Pas\_14)} tat yathā praskandanam prapatanam iti . \\
\noindent \textbf{(Pas\_14)} atha vā śabdaiḥ api śabdāḥ vyākriyante . \\
\noindent \textbf{(Pas\_14)} tat yathā gauḥ iti ukte sarve sandehāḥ nivartante na aśvaḥ na gardabhaḥ iti . \\
\noindent \textbf{(Pas\_14)} ayam tarhi doṣaḥ bhave proktādayaḥ ca taddhitāḥ iti . \\
\noindent \textbf{(Pas\_14)} evam tarhi lakṣyalakṣaṇe vyākaraṇam . \\
\noindent \textbf{(Pas\_14)} lakṣyam ca lakṣaṇam ca etat samuditam vyākaraṇam bhavati . \\
\noindent \textbf{(Pas\_14)} kim punaḥ lakṣyam lakṣaṇam ca . \\
\noindent \textbf{(Pas\_14)} śabdaḥ lakṣyam sūtram lakṣaṇam . \\
\noindent \textbf{(Pas\_14)} evam api ayam doṣaḥ samudāye vyākaraṇaśabdaḥ pravṛttaḥ avayave na upapadyate . \\
\noindent \textbf{(Pas\_14)} sūtrāṇi ca adhīyānaḥ iṣyate vaiyākaraṇaḥ iti . \\
\noindent \textbf{(Pas\_14)} na eṣaḥ doṣaḥ . \\
\noindent \textbf{(Pas\_14)} samudāyeṣu hi śabdāḥ pravṛttāḥ avayaveṣu api vartante . \\
\noindent \textbf{(Pas\_14)} tat yathā pūrve pañcālāḥ , uttare pañcālāḥ , tailam bhuktam , ghṛtam bhuktam , śuklaḥ , nīlaḥ , kṛṣṇaḥ iti . \\
\noindent \textbf{(Pas\_14)} evam ayam samudāye vyākaraṇaśabdaḥ pravṛttaḥ avayave api vartate . \\
\noindent \textbf{(Pas\_14)} atha vā punaḥ astu sūtram . \\
\noindent \textbf{(Pas\_14)} nanu ca uktam sūtre vyākaraṇe ṣaṣṭhyarthaḥ anupapannaḥ iti . \\
\noindent \textbf{(Pas\_14)} na eṣa doṣaḥ . \\
\noindent \textbf{(Pas\_14)} vyapadeśivadbhāvena bhaviṣyati . \\
\noindent \textbf{(Pas\_14)} yat api ucyate śabdāpratipattiḥ iti na hi sūtrataḥ eva śabdān pratipadyante kim tarhi vyākhyānataḥ ca iti parihṛtam etat tat eva sūtram vigṛhītam vyākhyānam bhavati iti . \\
\noindent \textbf{(Pas\_14)} nanu ca uktam na kevalāni carcāpadāni vyākhyānam vṛddhiḥ āt aic iti kim tarhi udāharaṇam pratyudāharaṇam vākyādhyāhāraḥ iti etat samuditam vyākhyānam bhavati iti . \\
\noindent \textbf{(Pas\_14)} avijānataḥ etat evam bhavati . \\
\noindent \textbf{(Pas\_14)} sūtrataḥ eva hi śabdān pratipadyante . \\
\noindent \textbf{(Pas\_14)} ātaḥ ca sūtrataḥ eva yaḥ hi utsūtram kathayet na adaḥ gṛhyeta . \\
\noindent \textbf{(Pas\_15)} atha kimarthaḥ varṇānām upadeśaḥ . \\
\noindent \textbf{(Pas\_15)} vṛttisamavāyārthaḥ upadeśaḥ . vṛttisamavāyārthaḥ varṇānām upadeśaḥ kartavyaḥ . \\
\noindent \textbf{(Pas\_15)} kim idam vṛttisamavayārthaḥ iti . \\
\noindent \textbf{(Pas\_15)} vṛttaye samavāyaḥ vṛttisamavāyaḥ , vṛttyarthaḥ vā samavāyaḥ vṛttisamavāyaḥ , vṛttiprayojanaḥ vā vṛttisamavayaḥ . \\
\noindent \textbf{(Pas\_15)} kā punaḥ vṛttiḥ . \\
\noindent \textbf{(Pas\_15)} śāstrapravṛttiḥ . \\
\noindent \textbf{(Pas\_15)} atha kaḥ samavayaḥ . \\
\noindent \textbf{(Pas\_15)} varṇānām ānupūrvyeṇa sanniveśaḥ . \\
\noindent \textbf{(Pas\_15)} atha kaḥ upadeśaḥ . \\
\noindent \textbf{(Pas\_15)} uccāraṇam . \\
\noindent \textbf{(Pas\_15)} kutaḥ etat . \\
\noindent \textbf{(Pas\_15)} diśiḥ uccāraṇakriyaḥ . \\
\noindent \textbf{(Pas\_15)} uccārya hi varṇān āha : upadiṣṭāḥ ime varṇāḥ iti . \\
\noindent \textbf{(Pas\_15)} anubandhakaraṇārthaḥ ca . anubandhakaraṇārthaḥ ca varṇānām upadeśaḥ kartavyaḥ . \\
\noindent \textbf{(Pas\_15)} anubandhān āsaṅkṣyāmi iti . \\
\noindent \textbf{(Pas\_15)} na hi anupadiśya varṇān anubandhāḥ śakyāḥ āsaṅktum . \\
\noindent \textbf{(Pas\_15)} saḥ eṣaḥ varṇānām upadeśaḥ vṛttisamavāyārthaḥ ca anubandhakaraṇārthaḥ ca . \\
\noindent \textbf{(Pas\_15)} vṛttisamavāyaḥ ca anubandhakaraṇārthaḥ ca pratyāhārārtham . \\
\noindent \textbf{(Pas\_15)} pratyāhāraḥ vṛttyarthaḥ . \\
\noindent \textbf{(Pas\_15)} iṣṭabuddhyarthaḥ ca . \\
\noindent \textbf{(Pas\_15)} iṣṭabuddhyarthaḥ ca varṇānām upadeśaḥ . \\
\noindent \textbf{(Pas\_15)} iṣṭān varṇān bhotsye iti . \\
\noindent \textbf{(Pas\_15)} iṣṭabuddhyarthaḥ ca iti cet udāttānudāttasvaritānunāsikdīrghaplutānām api upadeśaḥ . iṣṭabuddhyarthaḥ ca iti cet udāttānudāttasvaritānunāsikdīrghaplutānām api upadeśaḥ kartavyaḥ . \\
\noindent \textbf{(Pas\_15)} evaṅguṇāḥ api hi varṇāḥ iṣyante . \\
\noindent \textbf{(Pas\_15)} ākṛtyupadeśāt siddham . \\
\noindent \textbf{(Pas\_15)} ākṛtyupadeśāt siddham etat . \\
\noindent \textbf{(Pas\_15)} avarṇākṛtiḥ upadiṣṭā sarvam avarṇakulam grahīṣyati . \\
\noindent \textbf{(Pas\_15)} tathā ivarṇakulākṛtiḥ . \\
\noindent \textbf{(Pas\_15)} tathā uvarṇakulākṛtiḥ . \\
\noindent \textbf{(Pas\_15)} ākṛtyupadeśāt siddham iti cet saṃvṛtādīnām pratiṣedhaḥ . \\
\noindent \textbf{(Pas\_15)} ākṛtyupadeśāt siddham iti cet saṃvṛtādīnām pratiṣedhaḥ vaktavyaḥ . \\
\noindent \textbf{(Pas\_15)} ke punaḥ saṃvṛtādayaḥ . \\
\noindent \textbf{(Pas\_15)} saṃvṛtaḥ kalaḥ dhmātaḥ eṇīkṛtaḥ ambūkṛtaḥ ardhakaḥ grastaḥ nirastaḥ pragītaḥ upagītaḥ kṣviṇṇaḥ romaśaḥ iti . \\
\noindent \textbf{(Pas\_15)} aparaḥ āha : grastam nirastam avilambitam nirhatam ambūkṛtam dhmātam atho vikampitam sandaṣṭam eṇīkṛtam ardhakam drutam vikīrṇam etāḥ svaradoṣabhāvanāḥ iti . \\
\noindent \textbf{(Pas\_15)} ataḥ anye vyañjanadoṣāḥ . \\
\noindent \textbf{(Pas\_15)} na eṣaḥ doṣaḥ . \\
\noindent \textbf{(Pas\_15)} gargādibidādipāṭhāt saṃvṛtādīnām nivṛttiḥ bhaviṣyati . \\
\noindent \textbf{(Pas\_15)} asti anyat gargādibidādipāṭhe prayojanam . \\
\noindent \textbf{(Pas\_15)} kim . \\
\noindent \textbf{(Pas\_15)} samudāyānām sādhutvam yathā syāt iti . \\
\noindent \textbf{(Pas\_15)} evam tarhi aṣṭādaśadhā bhinnām nivṛttakalādikām avarṇasya pratyāpattim vakṣyāmi . \\
\noindent \textbf{(Pas\_15)} sā tarhi vaktavyā . \\
\noindent \textbf{(Pas\_15)} liṅgārthā tu pratyāpattiḥ . \\
\noindent \textbf{(Pas\_15)} liṅgārtha sā tarhi bhaviṣyati . \\
\noindent \textbf{(Pas\_15)} tat tarhi vaktavyam . \\
\noindent \textbf{(Pas\_15)} yadi api etat ucyate atha vā etarhi anubandhaśatam na uccāryam itsañjñā ca na vaktavyā lopaḥ ca na vaktavyaḥ . \\
\noindent \textbf{(Pas\_15)} yat anubandhaiḥ kriyate tat kalādibhiḥ kariṣyati . \\
\noindent \textbf{(Pas\_15)} sidhyati evam apāṇinīyam tu bhavati . \\
\noindent \textbf{(Pas\_15)} yathānyāsam eva astu . \\
\noindent \textbf{(Pas\_15)} nanu ca uktam ākṛtyupadeśāt siddham iti cet saṃvṛtādīnām pratiṣedhaḥ iti .parihṛtam etat gargādibidādipāṭhāt saṃvṛtādīnām nivṛttiḥ bhaviṣyati . \\
\noindent \textbf{(Pas\_15)} nanu ca anyat gargādibidādipāṭhe prayojanam uktam . \\
\noindent \textbf{(Pas\_15)} kim . \\
\noindent \textbf{(Pas\_15)} samudāyānām sādhutvam yathā syāt iti . \\
\noindent \textbf{(Pas\_15)} evam tarhi ubhayam anena kriyate . \\
\noindent \textbf{(Pas\_15)} pāṭhaḥ ca eva viśeṣyate kalādayaḥ ca nivartyante . \\
\noindent \textbf{(Pas\_15)} katham punaḥ ekena yatnena ubhayam labhyam . \\
\noindent \textbf{(Pas\_15)} labhyam iti āha . \\
\noindent \textbf{(Pas\_15)} katham . \\
\noindent \textbf{(Pas\_15)} dvigatāḥ api hetavaḥ bhavanti . \\
\noindent \textbf{(Pas\_15)} tat yathā : āmrāḥ ca siktāḥ pitaraḥ ca prīṇitāḥ iti . \\
\noindent \textbf{(Pas\_15)} tathā vākyāni api ḍviṣṭhāni bhavanti . \\
\noindent \textbf{(Pas\_15)} śvetaḥ dhāvati , alambusānām yātā iti . \\
\noindent \textbf{(Pas\_15)} atha vā idam tāvat ayam praṣṭavyaḥ . \\
\noindent \textbf{(Pas\_15)} kve ime saṃvṛtādayaḥ śrūyeran iti . \\
\noindent \textbf{(Pas\_15)} āgameṣu . \\
\noindent \textbf{(Pas\_15)} āgamāḥ śuddhāḥ paṭhyante . \\
\noindent \textbf{(Pas\_15)} vikāreṣu tarhi . \\
\noindent \textbf{(Pas\_15)} vikārāḥ śuddhāḥ paṭhyante . \\
\noindent \textbf{(Pas\_15)} pratyayeṣu tarhi . \\
\noindent \textbf{(Pas\_15)} pratyayāḥ śuddhāḥ paṭhyante . \\
\noindent \textbf{(Pas\_15)} dhātuṣu tarhi . \\
\noindent \textbf{(Pas\_15)} dhātavaḥ api śuddhāḥ paṭhyante . \\
\noindent \textbf{(Pas\_15)} prātipadikeṣu tarhi . \\
\noindent \textbf{(Pas\_15)} prātipadikāni api śuddhāni paṭhyante . \\
\noindent \textbf{(Pas\_15)} yāni tarhi agrahaṇāni prātipadikāni . \\
\noindent \textbf{(Pas\_15)} eteṣām api svaravarṇānupūrvījñānārthaḥ upadeśaḥ kartavyaḥ . \\
\noindent \textbf{(Pas\_15)} śaśaḥ ṣaṣaḥ iti mā bhūt . \\
\noindent \textbf{(Pas\_15)} palāśaḥ palāṣaḥ iti mā bhūt . \\
\noindent \textbf{(Pas\_15)} mañcakaḥ mañjakaḥ iti mā bhūt . \\
\noindent \textbf{(Pas\_15)} āgamāḥ ca vikārāḥ ca pratyayāḥ saha dhātubhiḥ uccāryante tataḥ teṣu na ime prāptāḥ kalādayaḥ . \\
\noindent \textbf{(Śs\_1.1)} akārasya vivṛtopadeśaḥ ākāragrahaṇārthaḥ . \\
\noindent \textbf{(Śs\_1.1)} akārasya vivṛtopadeśaḥ kartavyaḥ . \\
\noindent \textbf{(Śs\_1.1)} kim prayojanam . \\
\noindent \textbf{(Śs\_1.1)} ākāragrahaṇārthaḥ . \\
\noindent \textbf{(Śs\_1.1)} akāraḥ savarṇagragaṇena ākāram api yathā gṛhṇīyāt . \\
\noindent \textbf{(Śs\_1.1)} kim ca kāraṇam na gṛhṇīyāt . \\
\noindent \textbf{(Śs\_1.1)} vivārabhedāt . \\
\noindent \textbf{(Śs\_1.1)} kim ucyate vivārabhedāt iti na punaḥ kālabhedād api . \\
\noindent \textbf{(Śs\_1.1)} yathā eva hi vivārabhinnaḥ evam kālabhinnaḥ api . \\
\noindent \textbf{(Śs\_1.1)} satyam etat . \\
\noindent \textbf{(Śs\_1.1)} vakṣyati tulyāsyaprayatnam savarṇam iti atra āsyagrahaṇasya prayojanam āsye yeṣām tulyaḥ deśaḥ prayatnaḥ ca te savarṇasañjñakāḥ bhavanti iti . \\
\noindent \textbf{(Śs\_1.1)} bāhyaḥ ca punaḥ āsyāt kālaḥ . \\
\noindent \textbf{(Śs\_1.1)} tena syāt eva kālabhinnasya grahaṇam na punaḥ vivārabhinnasya . \\
\noindent \textbf{(Śs\_1.1)} kim punaḥ idam vivṛtasya upadiśyamānasya prayojanam anvākhyāyate āhosvit saṃvṛtasya upadiśyamānasya vivṛtopadeśaḥ codyate . \\
\noindent \textbf{(Śs\_1.1)} vivṛtasya upadiśyamānasya prayojanam anvākhyāyate . \\
\noindent \textbf{(Śs\_1.1)} katham jñāyate . \\
\noindent \textbf{(Śs\_1.1)} yat ayam a* a iti akārasya vivṛtasya saṃvṛtatāpratyāpattim śāsti . \\
\noindent \textbf{(Śs\_1.1)} na etat asti jñāpakam . \\
\noindent \textbf{(Śs\_1.1)} asti hi anyat etasya vacane prayojanam . \\
\noindent \textbf{(Śs\_1.1)} kim . \\
\noindent \textbf{(Śs\_1.1)} atikhaṭvaḥ , atimālaḥ iti atra āntaryataḥ vivṛtasya vivṛtaḥ prāpnoti . \\
\noindent \textbf{(Śs\_1.1)} saṃvṛtaḥ syāt iti evamarthā pratyāpattiḥ . \\
\noindent \textbf{(Śs\_1.1)} na etat asti . \\
\noindent \textbf{(Śs\_1.1)} na eva loke na ca vede akāro vivṛtaḥ asti . \\
\noindent \textbf{(Śs\_1.1)} kaḥ tarhi . \\
\noindent \textbf{(Śs\_1.1)} saṃvṛtaḥ . \\
\noindent \textbf{(Śs\_1.1)} yaḥ asti saḥ bhaviṣyati . \\
\noindent \textbf{(Śs\_1.1)} tat etat pratyāpattivacanam jñāpakam eva bhaviṣyati vivṛtasya upadiśyamānasya prayojanam anvākhyāyate iti . \\
\noindent \textbf{(Śs\_1.1)} kaḥ punaḥ atra viśeṣaḥ vivṛtasya upadiśyamānasya prayojanam anvākhyāyeta saṃvṛtasya upadiśyamānasya vā vivṛtopadeśaḥ codyeta iti . \\
\noindent \textbf{(Śs\_1.1)} na khalu kaḥ cid viśeṣaḥ . \\
\noindent \textbf{(Śs\_1.1)} āhopuruṣikāmātram tu bhavān āha saṃvṛtasya upadiśyamānasya vivṛtopadeśaḥ codyate iti . \\
\noindent \textbf{(Śs\_1.1)} vayam tu brūmaḥ vivṛtasya upadiśyamānasya prayojanam anvākhyāyate iti. tasya vivṛtopadeśāt anyatra api vivṛtopadeśaḥ savarṇagrahaṇārthaḥ . tasya etasya ākṣarasamāmnāyikasya vivṛtopadeśāt anyatra api vivṛtopadeśaḥ kartavyaḥ . \\
\noindent \textbf{(Śs\_1.1)} kva anyatra . \\
\noindent \textbf{(Śs\_1.1)} dhātuprātipadikapratyayanipātasthasya . \\
\noindent \textbf{(Śs\_1.1)} kim prayojanam . \\
\noindent \textbf{(Śs\_1.1)} savarṇagrahaṇārthaḥ . \\
\noindent \textbf{(Śs\_1.1)} ākṣarasamāmnāyikena asya grahaṇam yathā syāt . \\
\noindent \textbf{(Śs\_1.1)} kim ca kāraṇam na syāt . \\
\noindent \textbf{(Śs\_1.1)} vivārabhedāt eva . \\
\noindent \textbf{(Śs\_1.1)} ācāryapravṛttiḥ jñāpayati bhavati ākṣarasamāmnāyikena dhātvādisthasya grahaṇam iti yat ayam akaḥ savarṇe dīrghaḥ iti pratyāhāre akaḥ grahaṇam karoti . \\
\noindent \textbf{(Śs\_1.1)} katham kṛtvā jñāpakam . \\
\noindent \textbf{(Śs\_1.1)} na hi dvayoḥ ākṣarasamāmnāyikayoḥ yugapat samavasthānam asti . \\
\noindent \textbf{(Śs\_1.1)} na etat asti jñāpakam . \\
\noindent \textbf{(Śs\_1.1)} asti hi anyat etasya vacane prayojanam . \\
\noindent \textbf{(Śs\_1.1)} kim. yasya ākṣarasamāmnāyikena grahaṇam asti tadartham etat syāt . \\
\noindent \textbf{(Śs\_1.1)} khaṭvāḍhakam mālāḍhakam iti . \\
\noindent \textbf{(Śs\_1.1)} sati prayojane na jñāpakam bhavati . \\
\noindent \textbf{(Śs\_1.1)} tasmāt vivṛtopadeśaḥ kartavyaḥ . \\
\noindent \textbf{(Śs\_1.1)} kaḥ eṣaḥ yatnaḥ codyate vivṛtopadeśaḥ nāma . \\
\noindent \textbf{(Śs\_1.1)} vivṛtaḥ vā upadiśyeta saṃvṛtaḥ vā kaḥ nu atra viśeṣaḥ . \\
\noindent \textbf{(Śs\_1.1)} saḥ eṣaḥ sarvaḥ evamarthaḥ yatnaḥ yāni etāni prātipadikāni agrahaṇāni teṣām etena abhyupāyena upadeśaḥ codyate . \\
\noindent \textbf{(Śs\_1.1)} tat guru bhavati . \\
\noindent \textbf{(Śs\_1.1)} tasmāt vaktavyam dhātvādisthaḥ ca vivṛtaḥ iti . \\
\noindent \textbf{(Śs\_1.1)} dīrghaplutavacane ca saṃvṛtanivṛttyarthaḥ . dīrghaplutavacane ca saṃvṛtanivṛttyarthaḥ vivṛtopadeśaḥ kartavyaḥ . \\
\noindent \textbf{(Śs\_1.1)} dīrghaplutau saṃvṛtau mā bhūtām iti . \\
\noindent \textbf{(Śs\_1.1)} vṛkṣābhyām devadattā iti . \\
\noindent \textbf{(Śs\_1.1)} na eva loke na ca vede dīrghaplutau saṃvṛtau staḥ . \\
\noindent \textbf{(Śs\_1.1)} kau tarhi . \\
\noindent \textbf{(Śs\_1.1)} vivṛtau . \\
\noindent \textbf{(Śs\_1.1)} yau staḥ tau bhaviṣyataḥ . \\
\noindent \textbf{(Śs\_1.1)} sthānī prakalpayet etau anusvāraḥ yathā yaṇam . \\
\noindent \textbf{(Śs\_1.1)} saṃvṛtaḥ sthānī saṃvṛtau dīrghaplutau prakalpayet anusvāraḥ yathā yaṇam . \\
\noindent \textbf{(Śs\_1.1)} tat yathā sayŚyantā savŚvatsaraḥ yalŚ lokam talŚ lokam iti . \\
\noindent \textbf{(Śs\_1.1)} ansvāraḥ sthānī yaṇam anunāsikam prakalpayati . \\
\noindent \textbf{(Śs\_1.1)} viṣamaḥ upanyāsaḥ . \\
\noindent \textbf{(Śs\_1.1)} yuktam yat sataḥ tatra prakḷptiḥ bhavati . \\
\noindent \textbf{(Śs\_1.1)} santi hi yaṇaḥ sānunāsikāḥ niranunāsikāḥ ca . \\
\noindent \textbf{(Śs\_1.1)} dīrghaplutau punaḥ na eva loke na ca veda saṃvṛtau staḥ . \\
\noindent \textbf{(Śs\_1.1)} kau tarhi . \\
\noindent \textbf{(Śs\_1.1)} vivṛtau . \\
\noindent \textbf{(Śs\_1.1)} yau staḥ tau bhaviṣyataḥ . \\
\noindent \textbf{(Śs\_1.1)} evam api kutaḥ etat tulyasthānau prayatnabhinnau bhaviṣyataḥ na punaḥ tulyaprayatnau sthānabhinnau syātām īkāraḥ ūkāraḥ vā iti . \\
\noindent \textbf{(Śs\_1.1)} vakṣyati sthāne antaratamaḥ iti atra sthāne iti vartamāne punaḥ sthānegrahaṇasya prayojanam . \\
\noindent \textbf{(Śs\_1.1)} yatra anekavidham āntaryam tatra sthānataḥ eva āntaryam balīyaḥ yathā syāt . \\
\noindent \textbf{(Śs\_1.2)} tatra anuvṛttinirdeśe savarṇāgrahaṇam anaṇtvāt . \\
\noindent \textbf{(Śs\_1.2)} tatra anuvṛttinirdeśe savarṇānām grahaṇam na prāpnoti . \\
\noindent \textbf{(Śs\_1.2)} asya cvau yasya īti ca . \\
\noindent \textbf{(Śs\_1.2)} kim kāraṇam . \\
\noindent \textbf{(Śs\_1.2)} anaṇtvāt . \\
\noindent \textbf{(Śs\_1.2)} na hi ete aṇaḥ ye anuvṛttau . \\
\noindent \textbf{(Śs\_1.2)} ke tarhi . \\
\noindent \textbf{(Śs\_1.2)} ye akṣarasamāmnaye upadiśyante . \\
\noindent \textbf{(Śs\_1.2)} ekatvāt akārasya siddham . \\
\noindent \textbf{(Śs\_1.2)} ekaḥ ayam akāraḥ yaḥ ca akṣarasamāmnaye yaḥ ca anuvṛttau yaḥ ca dhātvādisthaḥ . \\
\noindent \textbf{(Śs\_1.2)} anubandhasaṅkaraḥ tu . \\
\noindent \textbf{(Śs\_1.2)} anubandhasaṅkaraḥ tu prāpnoti . \\
\noindent \textbf{(Śs\_1.2)} karmaṇi aṇ , ātaḥ anupasarge kaḥ iti ke api ṇitkṛtam prāpnoti . \\
\noindent \textbf{(Śs\_1.2)} ekājanekājgrahaṇeṣu ca anupapattiḥ . \\
\noindent \textbf{(Śs\_1.2)} ekājanekājgrahaṇeṣu ca anupapattiḥ bhaviṣyati . \\
\noindent \textbf{(Śs\_1.2)} tatra kaḥ doṣaḥ . \\
\noindent \textbf{(Śs\_1.2)} kiriṇā giriṇā iti atra ekājlakṣaṇam antodāttatvam prāpnoti . \\
\noindent \textbf{(Śs\_1.2)} iha ca ghaṭena tarati ghaṭika iti dvyajlakṣaṇaḥ ṭhan na prāpnoti . \\
\noindent \textbf{(Śs\_1.2)} dravyavat ca upacārāḥ . \\
\noindent \textbf{(Śs\_1.2)} dravyavat ca upacārāḥ prāpnuvanti . \\
\noindent \textbf{(Śs\_1.2)} tat yathā . \\
\noindent \textbf{(Śs\_1.2)} dravyeṣu na ekena ghaṭena anekaḥ yugapat kāryam karoti . \\
\noindent \textbf{(Śs\_1.2)} evam imam akāram na anekaḥ yugapat uccārayet . \\
\noindent \textbf{(Śs\_1.2)} viṣayeṇa tu nānāliṅgakāraṇāt siddham . \\
\noindent \textbf{(Śs\_1.2)} yat ayam viṣaye viṣaye nānāliṅgam akāram karoti karmaṇi aṇ ātaḥ anupasarge kaḥ iti tena jñāyate nānubandhasaṅkaraḥ asti iti . \\
\noindent \textbf{(Śs\_1.2)} yadi hi syāt nānāliṅgakaraṇam anarthakam syāt . \\
\noindent \textbf{(Śs\_1.2)} ekam eva ayam sarvaguṇam uccārayet . \\
\noindent \textbf{(Śs\_1.2)} na etat asti jñapakam . \\
\noindent \textbf{(Śs\_1.2)} itsañjñāprakḷptyartham etat syāt . \\
\noindent \textbf{(Śs\_1.2)} na hi ayam anubandhaiḥ śalyakavat śakyaḥ upacetum . \\
\noindent \textbf{(Śs\_1.2)} itsañjñayām hi doṣaḥ syāt . \\
\noindent \textbf{(Śs\_1.2)} āyamya hi dvayoḥ itsañjñā syāt . \\
\noindent \textbf{(Śs\_1.2)} kayoḥ . \\
\noindent \textbf{(Śs\_1.2)} ādyantayoḥ . \\
\noindent \textbf{(Śs\_1.2)} evam tarhi viṣayeṇa tu punaḥ liṅgakaraṇāt siddham . \\
\noindent \textbf{(Śs\_1.2)} yat ayam viṣaye viṣaye punaḥ liṅgam akāram karoti prāk dīvyataḥ aṇ , śivādibhyaḥ aṇ iti tena jñāyate na anubandhasaṅkaraḥ asti iti . \\
\noindent \textbf{(Śs\_1.2)} yadi hi syāt punaḥ liṅgakaraṇam anarthakam syāt . \\
\noindent \textbf{(Śs\_1.2)} atha vā punaḥ astu viṣayeṇa tu nānāliṅgakāraṇāt siddham iti eva . \\
\noindent \textbf{(Śs\_1.2)} nanu ca uktam itsañjñāprakḷptyartham etat syāt iti . \\
\noindent \textbf{(Śs\_1.2)} na eṣa doṣaḥ . \\
\noindent \textbf{(Śs\_1.2)} lokataḥ etat siddham . \\
\noindent \textbf{(Śs\_1.2)} tat yathā : loke kaḥ cit devadattam āha : iha muṇḍo bhava , iha jaṭī bhava , iha śikhī bhava iti . \\
\noindent \textbf{(Śs\_1.2)} yalliṅgaḥ yatra ucyate talliṅgaḥ tatra upatiṣṭhate . \\
\noindent \textbf{(Śs\_1.2)} evam akāraḥ yalliṅgaḥ yatra ucyate talliṅgaḥ tatra upasthāsyate . \\
\noindent \textbf{(Śs\_1.2)} yat api ucyate ekājanekājgrahaṇeṣu ca anupapattiḥ iti . \\
\noindent \textbf{(Śs\_1.2)} ekājanekājgrahaṇeṣu ca āvṛttisaṅkhyānāt . \\
\noindent \textbf{(Śs\_1.2)} ekājanekājgrahaṇeṣu ca āvṛtteḥ saṅkhyānāt anekāctvam bhaviṣyati . \\
\noindent \textbf{(Śs\_1.2)} tat yathā saptadaśa sāmidhenyaḥ bhavanti iti triḥ prathamam anvāha triḥ uttamam iti āvṛttitaḥ saptadaśatvam bhavati . \\
\noindent \textbf{(Śs\_1.2)} evam iha api āvṛttitaḥ anekāctvam bhaviṣyati . \\
\noindent \textbf{(Śs\_1.2)} bhaved āvṛttitaḥ kāryam parihṛtam . \\
\noindent \textbf{(Śs\_1.2)} iha tu khalu kiriṇā giriṇā iti ekājlakṣaṇam antodāttatvam prāpnoti eva . \\
\noindent \textbf{(Śs\_1.2)} etat api siddham . \\
\noindent \textbf{(Śs\_1.2)} katham . \\
\noindent \textbf{(Śs\_1.2)} lokataḥ . \\
\noindent \textbf{(Śs\_1.2)} tat yathā loke ṛṣisahasram ekām kapilām ekaikaśaḥ sahasrakṛtvaḥ dattvā tayā sarve te sahasradakṣiṇaḥ saṃpannāḥ evam iha api anekāctvam bhaviṣyati . \\
\noindent \textbf{(Śs\_1.2)} yat api ucyate dravyavat ca upacārāḥ prāpnuvanti iti . \\
\noindent \textbf{(Śs\_1.2)} bhavet yat asambhavi kāryam tat na anekaḥ yugapat kuryāt . \\
\noindent \textbf{(Śs\_1.2)} yat tu khalu sambhavi kāryam anekaḥ api tat yugapat karoti . \\
\noindent \textbf{(Śs\_1.2)} tat yathā ghaṭasya darśanam sparśanam vā . \\
\noindent \textbf{(Śs\_1.2)} sambhavi ca idam kāryam akārasya uccāraṇam nāma anekaḥ api tat yugapat kariṣyati . \\
\noindent \textbf{(Śs\_1.2)} ānyabhāvyam tu kālaśabdavyavāyāt . \\
\noindent \textbf{(Śs\_1.2)} ānyabhāvyam tu akārasya . \\
\noindent \textbf{(Śs\_1.2)} kutaḥ . \\
\noindent \textbf{(Śs\_1.2)} kālaśabdavyavāyāt . \\
\noindent \textbf{(Śs\_1.2)} kālavyavāyāt śabdavyavāyāt ca . \\
\noindent \textbf{(Śs\_1.2)} kālavyāvāyāt : daṇḍa , agram . \\
\noindent \textbf{(Śs\_1.2)} śabdavyavāyāt : daṇḍaḥ . \\
\noindent \textbf{(Śs\_1.2)} na ca ekasya ātmanaḥ vyavāyena bhavitavyam . \\
\noindent \textbf{(Śs\_1.2)} bhavati cet bhavati ānyabhāvyam akārasya . \\
\noindent \textbf{(Śs\_1.2)} yugapat ca deśapṛthaktvadarśanāt . \\
\noindent \textbf{(Śs\_1.2)} yugapat ca deśapṛthaktvadarśanāt manyāmahe ānyabhāvyam akārasya iti . \\
\noindent \textbf{(Śs\_1.2)} yat ayam yugapat deśapṛthaktveṣu upalabhyate . \\
\noindent \textbf{(Śs\_1.2)} aśvaḥ , arkaḥ , arthaḥ iti . \\
\noindent \textbf{(Śs\_1.2)} na hi ekaḥ devadattaḥ yugapat srughne ca bhavati mathurāyām ca . \\
\noindent \textbf{(Śs\_1.2)} yadi punaḥ ime varṇāḥ śakunivat syuḥ . \\
\noindent \textbf{(Śs\_1.2)} tat yathā śakunayaḥ āśugamitvāt purastāt utpatitāḥ paścāt dṛśyante evam ayam akāraḥ da iti atra dṛṣṭaḥ ṇḍa iti atra dṛśyate . \\
\noindent \textbf{(Śs\_1.2)} na evam śakyam . \\
\noindent \textbf{(Śs\_1.2)} anityatvam evam syāt . \\
\noindent \textbf{(Śs\_1.2)} nityāḥ ca śabdāḥ . \\
\noindent \textbf{(Śs\_1.2)} nityeṣu ca śabdeṣu kūṭasthaiḥ avicālibhiḥ varṇaiḥ bhavitavyam anapāyopajanavikāribhiḥ . \\
\noindent \textbf{(Śs\_1.2)} yadi ca ayam da iti atra dṛṣṭaḥ ṇḍa iti atra dṛśyeta na ayam kūṭasthaḥ syāt . \\
\noindent \textbf{(Śs\_1.2)} yadi punaḥ ime varṇāḥ ādityavat syuḥ . \\
\noindent \textbf{(Śs\_1.2)} tat yathā ekaḥ ādityaḥ anekādhikaraṇsthaḥ yugapat deśapṛthaktveṣu upalabhyate . \\
\noindent \textbf{(Śs\_1.2)} viṣamaḥ upanyāsaḥ . \\
\noindent \textbf{(Śs\_1.2)} na ekaḥ draṣṭā ādityam anekādhikaraṇastham yugapat deśapṛthaktveṣu upalabhate . \\
\noindent \textbf{(Śs\_1.2)} akāram punaḥ upalabhate . \\
\noindent \textbf{(Śs\_1.2)} akāram api na upalabhate . \\
\noindent \textbf{(Śs\_1.2)} kim kāraṇam . \\
\noindent \textbf{(Śs\_1.2)} śrotropalabdhiḥ buddhinirgrāhyaḥ prayogeṇa abhijvalitaḥ ākāśadeśaḥ śabdaḥ ekam ca ākāśam . \\
\noindent \textbf{(Śs\_1.2)} ākāśadeśāḥ api bahavaḥ . \\
\noindent \textbf{(Śs\_1.2)} yāvatā bahavaḥ tasmāt ānyabhāvyam akārasya . \\
\noindent \textbf{(Śs\_1.2)} ākṛtigrahaṇāt siddham . \\
\noindent \textbf{(Śs\_1.2)} avarṇākṛtiḥ upadiṣṭā sarvam avarṇakulam grahīṣyati . \\
\noindent \textbf{(Śs\_1.2)} tathā ivarṇākṛtiḥ . \\
\noindent \textbf{(Śs\_1.2)} tathā uvarṇākṛtiḥ . \\
\noindent \textbf{(Śs\_1.2)} tadvat ca taparakaraṇam . \\
\noindent \textbf{(Śs\_1.2)} evam ca kṛtvā taparāḥ kriyante . \\
\noindent \textbf{(Śs\_1.2)} ākṛtigrahaṇena atiprasaktam iti . \\
\noindent \textbf{(Śs\_1.2)} nanu ca savarṇagrahaṇena atiprasaktam iti kṛtvā taparāḥ kriyeran . \\
\noindent \textbf{(Śs\_1.2)} pratyākhyāyate tat : savarṇe aṇgrahaṇam aparibhāṣyam ākṛtigrahaṇāt ananyatvāt ca iti . \\
\noindent \textbf{(Śs\_1.2)} halgrahaṇeṣu ca . \\
\noindent \textbf{(Śs\_1.2)} kim . \\
\noindent \textbf{(Śs\_1.2)} ākṛtigrahaṇāt siddham iti eva . \\
\noindent \textbf{(Śs\_1.2)} jhalo jhali : avāttām avāttam avātta yatra etat na asti : aṇ savarṇān gṛhṇāti iti . \\
\noindent \textbf{(Śs\_1.2)} rūpasāmanyāt vā . \\
\noindent \textbf{(Śs\_1.2)} rūpasāmānyāt vā siddham . \\
\noindent \textbf{(Śs\_1.2)} tat yathā : tān eva śāṭakān ācchādayāmaḥ ye mathurāyām , tān eva śālīn bhuñjmahe ye magadheṣu , tat eva idam bhavataḥ kārṣāpaṇam yat mathurāyām gṛhītam . \\
\noindent \textbf{(Śs\_1.2)} anyasmin ca anyasmin ca rūpasāmānyāt tat eva idam iti bhavati . \\
\noindent \textbf{(Śs\_1.2)} evam iha api rūpasāmānyāt siddham \\
\noindent \textbf{(Śs\_2)} ḷkāropadeśaḥ kimarthaḥ . \\
\noindent \textbf{(Śs\_2)} kim viśeṣeṇa ḷkāropadeśaḥ codyate na punaḥ anyeṣām api varṇānām upadeśaḥ codyate . \\
\noindent \textbf{(Śs\_2)} yadi kim cit anyeṣām api varṇānām upadeśe prayojanam asti ḷkāropadeśasya api tat bhavitum arhati . \\
\noindent \textbf{(Śs\_2)} kaḥ vā viśeṣaḥ . \\
\noindent \textbf{(Śs\_2)} ayam asti viśeṣaḥ . \\
\noindent \textbf{(Śs\_2)} asya hi ḷkārasya alpīyān ca eva prayogaviṣayaḥ yaḥ ca api prayogaviṣayaḥ saḥ api kḷpisthasya . \\
\noindent \textbf{(Śs\_2)} kḷpeḥ ca latvam asiddham . \\
\noindent \textbf{(Śs\_2)} tasya asiddhatvāt ṛkārasya eva ackāryāṇi bhaviṣyanti . \\
\noindent \textbf{(Śs\_2)} na arthaḥ ḷkāropadeśena . \\
\noindent \textbf{(Śs\_2)} ataḥ uttaram paṭhati : ḷkāropadeśaḥ yadṛcchāśaktijānukaraṇaplutyādyarthaḥ . ḷkāropadeśaḥ kriyate yadṛcchāśabdārthaḥ aśaktijānukaraṇārthaḥ plutyādyarthaḥ ca . \\
\noindent \textbf{(Śs\_2)} yadṛcchāśabdārthaḥ tāvat . \\
\noindent \textbf{(Śs\_2)} yadṛcchayā kaḥ cit ḷtakaḥ nāma tasmin ackāryāṇi yathā syuḥ . \\
\noindent \textbf{(Śs\_2)} dadhi ḷtaka dehi . \\
\noindent \textbf{(Śs\_2)} madhu ḷtaka dehi . \\
\noindent \textbf{(Śs\_2)} udaṅ ḷtakaḥ agamat . \\
\noindent \textbf{(Śs\_2)} pratyaṅ ḷtakaḥ agamat . \\
\noindent \textbf{(Śs\_2)} catuṣṭayī śabdānām pravṛttiḥ : jātiśabdāḥ guṇaśabdāḥ kriyāśabdāḥ yadṛcchāśabdāḥ caturthāḥ . \\
\noindent \textbf{(Śs\_2)} aśaktijānukaraṇārthaḥ . \\
\noindent \textbf{(Śs\_2)} aśaktyā kayā cit brāhmaṇyā ṛtakaḥ iti prayoktavye ḷtakaḥ iti prayuktam . \\
\noindent \textbf{(Śs\_2)} tasya anukaraṇam : brāhmaṇī ḷtakaḥ iti āha . \\
\noindent \textbf{(Śs\_2)} kumārī ḷtakaḥ iti āha iti . \\
\noindent \textbf{(Śs\_2)} plutādyarthaḥ ca ḷkāropadeśaḥ kartavyaḥ . \\
\noindent \textbf{(Śs\_2)} ke punaḥ plutādayaḥ . \\
\noindent \textbf{(Śs\_2)} plutidvirvacanasvaritāḥ : kḷptaśikha kḷpptaḥ , prakḷptaḥ . \\
\noindent \textbf{(Śs\_2)} plutyādiṣu kāryeṣu kḷpeḥ latvam siddham . \\
\noindent \textbf{(Śs\_2)} tasya siddhatvāt ackāryāṇi na sidhyanti . \\
\noindent \textbf{(Śs\_2)} tasmāt ḷkāropadeśaḥ kartavyaḥ . \\
\noindent \textbf{(Śs\_2)} na etāni santi prayojanāni . \\
\noindent \textbf{(Śs\_2)} nyāyyabhāvāt kalpanam sañjñādiṣu . nyāyyasya ṛtakaśabdasya bhāvāt kalpanam sañjñādiṣu sādhu manyante . \\
\noindent \textbf{(Śs\_2)} ṛtakaḥ eva asau na ḷtakaḥ iti . \\
\noindent \textbf{(Śs\_2)} aparaḥ āha : nyāyyaḥ ṛtakaśabdaḥ śāstrānvitaḥ asti . \\
\noindent \textbf{(Śs\_2)} saḥ kalpayitavyaḥ sādhuḥ sañjñādiṣu . \\
\noindent \textbf{(Śs\_2)} ṛtakaḥ eva asau na ḷtakaḥ . \\
\noindent \textbf{(Śs\_2)} ayam tarhi yadṛcchāśabdaḥ aparihāryaḥ . \\
\noindent \textbf{(Śs\_2)} ḷphiḍaḥ ḷphiḍḍaḥ . \\
\noindent \textbf{(Śs\_2)} eṣaḥ api ṛphiḍaḥ ṛphiḍḍaḥ ca . \\
\noindent \textbf{(Śs\_2)} katham . \\
\noindent \textbf{(Śs\_2)} artipravṛttiḥ ca eva loke lakṣyate phiḍiphiḍḍau auṇādikau pratyayau . \\
\noindent \textbf{(Śs\_2)} trayī ca śabdānām pravṛttiḥ . \\
\noindent \textbf{(Śs\_2)} jātiśabdāḥ guṇaśabdāḥ kriyāśabdāḥ iti . \\
\noindent \textbf{(Śs\_2)} na santi yadṛcchāśabdāḥ . \\
\noindent \textbf{(Śs\_2)} anyathā kṛtvā prayojanam uktam anyathā kṛtvā parihāraḥ . \\
\noindent \textbf{(Śs\_2)} santi yadṛcchāśabdāḥ iti kṛtvā prayojanam uktam . \\
\noindent \textbf{(Śs\_2)} na santi iti parihāraḥ . \\
\noindent \textbf{(Śs\_2)} samāne ca arthe śāstrānvitaḥ aśāstrānvitasya nivartakaḥ bhavati . \\
\noindent \textbf{(Śs\_2)} tat yathā . \\
\noindent \textbf{(Śs\_2)} devadattaśabdaḥ devadiṇṇaśabdaṃ nivartayati na gāvyādīn . \\
\noindent \textbf{(Śs\_2)} na eṣa doṣaḥ . \\
\noindent \textbf{(Śs\_2)} pakṣāntaraiḥ api parihārāḥ bhavanti . \\
\noindent \textbf{(Śs\_2)} anukaraṇam śiṣṭāśiṣṭāpratiṣiddheṣu yathā laukikavaidikeṣu . anukaraṇam hi śiṣṭasya sādhu bhavati . \\
\noindent \textbf{(Śs\_2)} aśiṣṭāpratiṣiddhasya vā na eva tat doṣāya bhavati na abhyudayāya . \\
\noindent \textbf{(Śs\_2)} yathā laukikavaidikeṣu . \\
\noindent \textbf{(Śs\_2)} yathā laukikeṣu vaidikeṣu ca kṛtānteṣu . \\
\noindent \textbf{(Śs\_2)} loke tāvat : yaḥ evam asau dadāti yaḥ evam asau yajate yaḥ evam asau adhīte iti tasya anukurvan dadyāt ca yajeta ca adhīyīta ca saḥ abhyudayena yujyate . \\
\noindent \textbf{(Śs\_2)} vede api : ye evam viśvasṛjaḥ sattrāṇi adhyāsate iti teṣām anukurvan tadvat sattrāṇi adhyāsīta saḥ api abhyudayena yujyate . \\
\noindent \textbf{(Śs\_2)} aśiṣṭāpratiṣiddham . \\
\noindent \textbf{(Śs\_2)} yaḥ evam asau hikkati yaḥ evam asau hasati yaḥ evam asau kaṇḍūyati iti tasya anukurvan hikket ca haset ca kaṇḍūyet ca na eva tat doṣāya syāt na abhyudayāya . \\
\noindent \textbf{(Śs\_2)} yaḥ tu khalu evam asau brāhmaṇam hanti evam asau surām pibati iti tasya anukurvan brāhmaṇam hanyāt surām vā pibet saḥ api manye patitaḥ syāt . \\
\noindent \textbf{(Śs\_2)} viṣamaḥ upanyāsaḥ . \\
\noindent \textbf{(Śs\_2)} yaḥ ca evam hanti yaḥ ca anuhanti ubhau tau hataḥ . \\
\noindent \textbf{(Śs\_2)} yaḥ ca pibati yaḥ ca anupibati ubhau tau pibataḥ . \\
\noindent \textbf{(Śs\_2)} yaḥ tu khalu evam asau brāhmaṇam hanti evam asau surām pibati iti tasya anukurvan snātānuliptaḥ mālyaguṇakaṇṭhaḥ kadalīstambham chindyāt payaḥ vā pibet na sa manye patitaḥ . \\
\noindent \textbf{(Śs\_2)} evam iha api yaḥ evam asau apaśabdam prayuṅkte iti tasya anukurvan apaśabdam prayuñjīta saḥ api apaśabdabhāk syāt . \\
\noindent \textbf{(Śs\_2)} ayam tu anyaḥ apaśabdapadārthakaḥ śabdaḥ yadarthaḥ upadeśaḥ kartavyaḥ . \\
\noindent \textbf{(Śs\_2)} na ca apaśabdapadārthakaḥ śabdaḥ apaśabdaḥ bhavati . \\
\noindent \textbf{(Śs\_2)} avaśyam ca etat evam vijñeyam . \\
\noindent \textbf{(Śs\_2)} yaḥ hi manyeta apaśabdapadārthakaḥ śabdaḥ apaśabdaḥ bhavati iti apaśabdaḥ iti eva tasya apaśabdaḥ syāt . \\
\noindent \textbf{(Śs\_2)} na ca eṣaḥ apaśabdaḥ . \\
\noindent \textbf{(Śs\_2)} ayam khalu api bhūyaḥ anukaraṇaśabdaḥ aparihāryaḥ yadarthaḥ upadeśaḥ kartavyaḥ . \\
\noindent \textbf{(Śs\_2)} sādhu ḷkāram adhīte . \\
\noindent \textbf{(Śs\_2)} madhu ḷkāram adhīte iti . \\
\noindent \textbf{(Śs\_2)} kvasthasya punaḥ etat anukaraṇam . \\
\noindent \textbf{(Śs\_2)} kḷpisthasya . \\
\noindent \textbf{(Śs\_2)} yadi kḷpisthasya kḷpeḥ ca latvam asiddham tasya asiddhatvāt ṛkāre eva ackāryāṇi bhaviṣyanti . \\
\noindent \textbf{(Śs\_2)} bhavet tadarthena na arthaḥ syāt . \\
\noindent \textbf{(Śs\_2)} ayam tu anyaḥ kḷpisthapadārthakaḥ śabdaḥ yadarthaḥ upadeśaḥ kartavyaḥ . \\
\noindent \textbf{(Śs\_2)} na kartavyaḥ . \\
\noindent \textbf{(Śs\_2)} idam avaśyam vaktavyam . \\
\noindent \textbf{(Śs\_2)} prakṛtivat anukaraṇam bhavati iti . \\
\noindent \textbf{(Śs\_2)} kim prayojanam . \\
\noindent \textbf{(Śs\_2)} dviḥ pacantu iti āha . \\
\noindent \textbf{(Śs\_2)} tiṅ atiṅaḥ iti nighātaḥ yathā syāt . \\
\noindent \textbf{(Śs\_2)} agnī iti āha . \\
\noindent \textbf{(Śs\_2)} īdūdet dvivacanam pragṛhyam iti pragṛhyasañjñā yathā syāt . \\
\noindent \textbf{(Śs\_2)} yadi prakṛtivat anukaraṇam bhavati iti ucyate apaśabdaḥ eva asau bhavati kumārīḷtakaḥ iti āha . \\
\noindent \textbf{(Śs\_2)} brahmaṇī ḷtakaḥ iti āha . \\
\noindent \textbf{(Śs\_2)} apaśabdaḥ hi asya prakṛtiḥ . \\
\noindent \textbf{(Śs\_2)} na cāpaśabdaḥ prakṛtiḥ . \\
\noindent \textbf{(Śs\_2)} na hi apaśabdāḥ upadiśyante . \\
\noindent \textbf{(Śs\_2)} na ca anupadiṣṭā prakṛtiḥ asti. ekadeśavikṛtasya ananyatvāt plutyādayaḥ . \\
\noindent \textbf{(Śs\_2)} ekadeśavikṛtam ananyavat bhavati iti plutyādayaḥ bhaviṣyanti . \\
\noindent \textbf{(Śs\_2)} yadi ekadeśavikṛtam ananyavat bhavati iti ucyate rājñaḥ ka ca rājakīyam allopaḥ anaḥ iti lopaḥ prāpnoti . \\
\noindent \textbf{(Śs\_2)} ekadeśavikṛtam ananyavat ṣaṣṭhīnirdiṣṭasya iti vakṣyāmi . \\
\noindent \textbf{(Śs\_2)} yadi ṣaṣṭhīnirdiṣṭasya iti ucyate kḷptaśikha iti plutaḥ na prāpnoti . \\
\noindent \textbf{(Śs\_2)} na hi atra ṛkāraḥ ṣaṣṭhīnirdiṣṭaḥ . \\
\noindent \textbf{(Śs\_2)} kaḥ tarhi . \\
\noindent \textbf{(Śs\_2)} rephaḥ . \\
\noindent \textbf{(Śs\_2)} ṛkāraḥ api atra ṣaṣṭhīnirdiṣṭaḥ . \\
\noindent \textbf{(Śs\_2)} katham . \\
\noindent \textbf{(Śs\_2)} avibhaktikaḥ nirdeśaḥ . \\
\noindent \textbf{(Śs\_2)} kṛpa uḥ raḥ laḥ kṛpo ro laḥ iti . \\
\noindent \textbf{(Śs\_2)} atha vā punaḥ astu aviśeṣeṇa . \\
\noindent \textbf{(Śs\_2)} nanu ca uktam . \\
\noindent \textbf{(Śs\_2)} rājñaḥ ka ca rājakīyam allopaḥ anaḥ iti lopaḥ prāpnoti iti . \\
\noindent \textbf{(Śs\_2)} vakṣyati etat . \\
\noindent \textbf{(Śs\_2)} śvādīnām prasāraṇe nakārāntagrahaṇam anakārāntapratiṣedhārtham iti . \\
\noindent \textbf{(Śs\_2)} tat prakṛtam uttaratra anuvartiṣyate . \\
\noindent \textbf{(Śs\_2)} allopaḥ anaḥ nakārāntasya iti . \\
\noindent \textbf{(Śs\_2)} iha tarhi kḷptaśikha iti anṛtaḥ iti pratiṣedhaḥ prāpnoti . \\
\noindent \textbf{(Śs\_2)} ravatpratiṣedhāt ca . ravatpratiṣedhāt ca etat sidhyati . \\
\noindent \textbf{(Śs\_2)} guroḥ aravataḥ iti vakṣyāmi . \\
\noindent \textbf{(Śs\_2)} yadi aravataḥ iti ucyate hotṛ-ṛkāra , hotṝṛkāra , atra na prāpnoti . \\
\noindent \textbf{(Śs\_2)} guroḥ aravataḥ hrasvasya iti vakṣyāmi . \\
\noindent \textbf{(Śs\_2)} saḥ eṣaḥ sūtrabhedena ḷkāraḥ plutyādyarthaḥ san pratyākhyāyate . \\
\noindent \textbf{(Śs\_2)} sā eṣā mahataḥ vaṃśastambāt laṭvā anukṛṣyate . \\
\noindent \textbf{(Śs\_3-4.1)} KA\_I,22.2-23.23 Ro\_I,79-84 {1/80}     idam vicāryate . \\
\noindent \textbf{(Śs\_3-4.1)} KA\_I,22.2-23.23 Ro\_I,79-84 {2/80}     imāni sandhyakṣarāṇi taparāṇi vā upadiśyeran . \\
\noindent \textbf{(Śs\_3-4.1)} KA\_I,22.2-23.23 Ro\_I,79-84 {3/80}     et , ot , ṅ . \\
\noindent \textbf{(Śs\_3-4.1)} KA\_I,22.2-23.23 Ro\_I,79-84 {4/80}     ait , aut , c iti . \\
\noindent \textbf{(Śs\_3-4.1)} KA\_I,22.2-23.23 Ro\_I,79-84 {5/80}     ataparāṇi vā yathānyāsam iti . \\
\noindent \textbf{(Śs\_3-4.1)} KA\_I,22.2-23.23 Ro\_I,79-84 {6/80}     kaḥ ca atra viśeṣaḥ . \\
\noindent \textbf{(Śs\_3-4.1)} KA\_I,22.2-23.23 Ro\_I,79-84 {7/80}     sandhyakṣareṣu taparopadeśaḥ cet taparoccāraṇam . \\
\noindent \textbf{(Śs\_3-4.1)} KA\_I,22.2-23.23 Ro\_I,79-84 {8/80}     sandhyakṣareṣu taparopadeśaḥ cet taparoccāraṇam kartavyam . \\
\noindent \textbf{(Śs\_3-4.1)} KA\_I,22.2-23.23 Ro\_I,79-84 {9/80}     plutyādiṣu ajvidhiḥ . \\
\noindent \textbf{(Śs\_3-4.1)} KA\_I,22.2-23.23 Ro\_I,79-84 {10/80}     plutyādiṣu ajāśrayaḥ vidhiḥ na sidhyati . \\
\noindent \textbf{(Śs\_3-4.1)} KA\_I,22.2-23.23 Ro\_I,79-84 {11/80}     gotrāta nautrāta iti atra anaci ca iti acaḥ uttarasya yaraḥ dve bhavataḥ iti dvirvacanam na prāpnoti . \\
\noindent \textbf{(Śs\_3-4.1)} KA\_I,22.2-23.23 Ro\_I,79-84 {12/80}     iha ca pratyaṅ aitikayana udaṅ aupagava iti aci iti ṅamuṭ na prāpnoti . \\
\noindent \textbf{(Śs\_3-4.1)} KA\_I,22.2-23.23 Ro\_I,79-84 {13/80}     plutasañjñā ca . \\
\noindent \textbf{(Śs\_3-4.1)} KA\_I,22.2-23.23 Ro\_I,79-84 {14/80}     plutasañjñā ca na sidhyati . \\
\noindent \textbf{(Śs\_3-4.1)} KA\_I,22.2-23.23 Ro\_I,79-84 {15/80}     aitikayana aupagava . \\
\noindent \textbf{(Śs\_3-4.1)} KA\_I,22.2-23.23 Ro\_I,79-84 {16/80}     ūkalaḥ ac hrasvadīrghaplutaḥ iti plutasañjñā na prāpnoti . \\
\noindent \textbf{(Śs\_3-4.1)} KA\_I,22.2-23.23 Ro\_I,79-84 {17/80}     santu tarhi ataparāṇi . \\
\noindent \textbf{(Śs\_3-4.1)} KA\_I,22.2-23.23 Ro\_I,79-84 {18/80}     atapare ecaḥ ik hrasvādeśe . \\
\noindent \textbf{(Śs\_3-4.1)} KA\_I,22.2-23.23 Ro\_I,79-84 {19/80}     yadi ataparāṇi ecaḥ ik hrasvādeśe iti vaktavyam . \\
\noindent \textbf{(Śs\_3-4.1)} KA\_I,22.2-23.23 Ro\_I,79-84 {20/80}     kim prayojanam . \\
\noindent \textbf{(Śs\_3-4.1)} KA\_I,22.2-23.23 Ro\_I,79-84 {21/80}     ecaḥ hrasvādeśaśāsaneṣu ardhaḥ ekāraḥ ardhaḥ okāraḥ vā mā bhūt iti . \\
\noindent \textbf{(Śs\_3-4.1)} KA\_I,22.2-23.23 Ro\_I,79-84 {22/80}     nanu ca yasya api taparāṇi tena api etat vaktavyam . \\
\noindent \textbf{(Śs\_3-4.1)} KA\_I,22.2-23.23 Ro\_I,79-84 {23/80}     imau aicau samāhāravarṇau mātrā avarṇasya mātrā ivarṇovarṇayoḥ . \\
\noindent \textbf{(Śs\_3-4.1)} KA\_I,22.2-23.23 Ro\_I,79-84 {24/80}     tayoḥ hrasvādeśaśāsaneṣu kadā cit avarṇaḥ syāt kadā cit ivarṇovarṇau . \\
\noindent \textbf{(Śs\_3-4.1)} KA\_I,22.2-23.23 Ro\_I,79-84 {25/80}     mā kadā cit avarṇam bhūt iti . \\
\noindent \textbf{(Śs\_3-4.1)} KA\_I,22.2-23.23 Ro\_I,79-84 {26/80}     pratyākhyāyate etat : aicoḥ ca uttarabhūyastvāt iti . \\
\noindent \textbf{(Śs\_3-4.1)} KA\_I,22.2-23.23 Ro\_I,79-84 {27/80}     yadi pratyākhyānapakṣaḥ idam api pratyākhyāyate : siddham eṅaḥ sasthānatvāt iti . \\
\noindent \textbf{(Śs\_3-4.1)} KA\_I,22.2-23.23 Ro\_I,79-84 {28/80}     nanu ca eṅaḥ sasthānatarau ardhaḥ ekāraḥ ardhaḥ okāraḥ ca . \\
\noindent \textbf{(Śs\_3-4.1)} KA\_I,22.2-23.23 Ro\_I,79-84 {29/80}     na tau staḥ . \\
\noindent \textbf{(Śs\_3-4.1)} KA\_I,22.2-23.23 Ro\_I,79-84 {30/80}     yadi hi tau syātām tau eva ayam upadiśet . \\
\noindent \textbf{(Śs\_3-4.1)} KA\_I,22.2-23.23 Ro\_I,79-84 {31/80}     nanu ca bhoḥ chandogānām sātyamugrirāṇāyanīyāḥ ardham ekāram ardham okāram ca adhīyate : sujāte eśvasūnṛte , adhvaryo odribhiḥ sutam , śukram te enyat yajatam te enyat iti . \\
\noindent \textbf{(Śs\_3-4.1)} KA\_I,22.2-23.23 Ro\_I,79-84 {32/80}     pārṣadakṛtiḥ eṣā tatrabhavatām . \\
\noindent \textbf{(Śs\_3-4.1)} KA\_I,22.2-23.23 Ro\_I,79-84 {33/80}     na eva hi loke na anyasmin vede ardhaḥ ekāraḥ ardhaḥ okāraḥ vā asti . \\
\noindent \textbf{(Śs\_3-4.1)} KA\_I,22.2-23.23 Ro\_I,79-84 {34/80}     ekādeśe dīrghagrahaṇam . \\
\noindent \textbf{(Śs\_3-4.1)} KA\_I,22.2-23.23 Ro\_I,79-84 {35/80}     ekādeśe dīrghagrahaṇam kartavyam . \\
\noindent \textbf{(Śs\_3-4.1)} KA\_I,22.2-23.23 Ro\_I,79-84 {36/80}     āt guṇaḥ dīrghaḥ . \\
\noindent \textbf{(Śs\_3-4.1)} KA\_I,22.2-23.23 Ro\_I,79-84 {37/80}     vṛddhiḥ eci dīrghaḥ iti . \\
\noindent \textbf{(Śs\_3-4.1)} KA\_I,22.2-23.23 Ro\_I,79-84 {38/80}     kim prayojanam . \\
\noindent \textbf{(Śs\_3-4.1)} KA\_I,22.2-23.23 Ro\_I,79-84 {39/80}     āntaryataḥ trimātracaturnātrāṇām sthāninām trimātracaturmātrāḥ ādeśāḥ mā bhūvan iti. khaṭvā , indraḥ khaṭvendraḥ , khaṭvā , udakam khaṭvodakam , khaṭvā , īṣā khaṭveṣā , khaṭvā , ūḍhā khaṭvoḍhā , khaṭvā , elakā khaṭvailakā , khaṭvā , odanaḥ khaṭvaudanaḥ , khaṭvā , aitikāyanaḥ khaṭvaitikāyanaḥ , khaṭvā , aupagavaḥ khaṭvaupagavaḥ . \\
\noindent \textbf{(Śs\_3-4.1)} KA\_I,22.2-23.23 Ro\_I,79-84 {40/80}     tat tarhi dīrghagrahaṇam kartavyam . \\
\noindent \textbf{(Śs\_3-4.1)} KA\_I,22.2-23.23 Ro\_I,79-84 {41/80}     na kartavyam . \\
\noindent \textbf{(Śs\_3-4.1)} KA\_I,22.2-23.23 Ro\_I,79-84 {42/80}     upariṣṭāt yogavibhāgaḥ kariṣyate . \\
\noindent \textbf{(Śs\_3-4.1)} KA\_I,22.2-23.23 Ro\_I,79-84 {43/80}     akaḥ savarṇe ekaḥ bhavati . \\
\noindent \textbf{(Śs\_3-4.1)} KA\_I,22.2-23.23 Ro\_I,79-84 {44/80}     tataḥ dīrghaḥ . \\
\noindent \textbf{(Śs\_3-4.1)} KA\_I,22.2-23.23 Ro\_I,79-84 {45/80}     dīrghaḥ ca sa bhavati yaḥ saḥ ekaḥ pūrvaparayoḥ iti evam nirdiṣṭaḥ iti . \\
\noindent \textbf{(Śs\_3-4.1)} KA\_I,22.2-23.23 Ro\_I,79-84 {46/80}     iha api tarhi prāpnoti . \\
\noindent \textbf{(Śs\_3-4.1)} KA\_I,22.2-23.23 Ro\_I,79-84 {47/80}     paśum , viddham , pacanti iti . \\
\noindent \textbf{(Śs\_3-4.1)} KA\_I,22.2-23.23 Ro\_I,79-84 {48/80}     na eṣaḥ doṣaḥ . \\
\noindent \textbf{(Śs\_3-4.1)} KA\_I,22.2-23.23 Ro\_I,79-84 {49/80}     iha tāvat paśum iti ami ekaḥ iti iyatā siddham . \\
\noindent \textbf{(Śs\_3-4.1)} KA\_I,22.2-23.23 Ro\_I,79-84 {50/80}     saḥ ayam evam siddhe sati yat pūrvagrahaṇam karoti tasya etat prayojanam yathājātīyakaḥ pūrvaḥ tathājātīyakaḥ ubhayoḥ yathā syāt iti . \\
\noindent \textbf{(Śs\_3-4.1)} KA\_I,22.2-23.23 Ro\_I,79-84 {51/80}     viddham iti pūrvaḥ iti eva anuvartate . \\
\noindent \textbf{(Śs\_3-4.1)} KA\_I,22.2-23.23 Ro\_I,79-84 {52/80}     atha vā ācāryapravṛttiḥ jñāpayati na anena samprasāraṇasya dīrghaḥ bhavati iti yat ayam halaḥ uttarasya samprasāraṇasya dīrghatvam śāsti . \\
\noindent \textbf{(Śs\_3-4.1)} KA\_I,22.2-23.23 Ro\_I,79-84 {53/80}     pacanti iti ataḥ guṇe paraḥ iti iyatā siddham . \\
\noindent \textbf{(Śs\_3-4.1)} KA\_I,22.2-23.23 Ro\_I,79-84 {54/80}     saḥ ayam evam siddhe sati yat rūpagrahaṇam karoti tasya etat prayojanam yathājātīyakam parasya rūpam tathājātīyakam ubhayoḥ yathā syāt iti . \\
\noindent \textbf{(Śs\_3-4.1)} KA\_I,22.2-23.23 Ro\_I,79-84 {55/80}     iha tarhi khaṭvarśyaḥ mālarśyaḥ iti dīrghavacanāt akāraḥ na ānantaryāt ekākaukārau na . \\
\noindent \textbf{(Śs\_3-4.1)} KA\_I,22.2-23.23 Ro\_I,79-84 {56/80}     tatra kaḥ doṣaḥ . \\
\noindent \textbf{(Śs\_3-4.1)} KA\_I,22.2-23.23 Ro\_I,79-84 {57/80}     vigṛhītasya śravaṇam prasajyeta . \\
\noindent \textbf{(Śs\_3-4.1)} KA\_I,22.2-23.23 Ro\_I,79-84 {58/80}     na brūmaḥ yatra kriyamāṇe doṣaḥ tatra kartavyam . \\
\noindent \textbf{(Śs\_3-4.1)} KA\_I,22.2-23.23 Ro\_I,79-84 {59/80}     kim tarhi . \\
\noindent \textbf{(Śs\_3-4.1)} KA\_I,22.2-23.23 Ro\_I,79-84 {60/80}     yatra kriyamāṇe na doṣaḥ tatra kartavyam iti . \\
\noindent \textbf{(Śs\_3-4.1)} KA\_I,22.2-23.23 Ro\_I,79-84 {61/80}     kva ca kriyamāṇe na doṣaḥ . \\
\noindent \textbf{(Śs\_3-4.1)} KA\_I,22.2-23.23 Ro\_I,79-84 {62/80}     sañjñāvidhau . \\
\noindent \textbf{(Śs\_3-4.1)} KA\_I,22.2-23.23 Ro\_I,79-84 {63/80}     vṛddhīḥ āt aic dīrghaḥ . \\
\noindent \textbf{(Śs\_3-4.1)} KA\_I,22.2-23.23 Ro\_I,79-84 {64/80}     at eṅ guṇaḥ dīrghaḥ iti . \\
\noindent \textbf{(Śs\_3-4.1)} KA\_I,22.2-23.23 Ro\_I,79-84 {65/80}     tat tarhi dīrghagrahaṇam kartavyam . \\
\noindent \textbf{(Śs\_3-4.1)} KA\_I,22.2-23.23 Ro\_I,79-84 {66/80}     na kartavyam . \\
\noindent \textbf{(Śs\_3-4.1)} KA\_I,22.2-23.23 Ro\_I,79-84 {67/80}     kasmāt eva āntaryataḥ trimātracaturnātrāṇām sthāninam trimātracaturmātrāḥ ādeśāḥ na bhavanti . \\
\noindent \textbf{(Śs\_3-4.1)} KA\_I,22.2-23.23 Ro\_I,79-84 {68/80}     tapare guṇavṛddhī . \\
\noindent \textbf{(Śs\_3-4.1)} KA\_I,22.2-23.23 Ro\_I,79-84 {69/80}     nanu ca bhoḥ taḥ paraḥ yasmāt saḥ ayam taparaḥ . \\
\noindent \textbf{(Śs\_3-4.1)} KA\_I,22.2-23.23 Ro\_I,79-84 {70/80}     na iti āha . \\
\noindent \textbf{(Śs\_3-4.1)} KA\_I,22.2-23.23 Ro\_I,79-84 {71/80}     tāt api paraḥ taparaḥ iti . \\
\noindent \textbf{(Śs\_3-4.1)} KA\_I,22.2-23.23 Ro\_I,79-84 {72/80}     yadi tāt api paraḥ taparaḥ ṝdoḥ ap iti iha eva syāt : yavaḥ stavaḥ . \\
\noindent \textbf{(Śs\_3-4.1)} KA\_I,22.2-23.23 Ro\_I,79-84 {73/80}     lavaḥ pavaḥ iti atra na syāt . \\
\noindent \textbf{(Śs\_3-4.1)} KA\_I,22.2-23.23 Ro\_I,79-84 {74/80}     na eṣaḥ takāraḥ . \\
\noindent \textbf{(Śs\_3-4.1)} KA\_I,22.2-23.23 Ro\_I,79-84 {75/80}     kaḥ tarhi . \\
\noindent \textbf{(Śs\_3-4.1)} KA\_I,22.2-23.23 Ro\_I,79-84 {76/80}     dakāraḥ . \\
\noindent \textbf{(Śs\_3-4.1)} KA\_I,22.2-23.23 Ro\_I,79-84 {77/80}     kim dakāre prayojanam . \\
\noindent \textbf{(Śs\_3-4.1)} KA\_I,22.2-23.23 Ro\_I,79-84 {78/80}     atha kim takāre prayojanam . \\
\noindent \textbf{(Śs\_3-4.1)} KA\_I,22.2-23.23 Ro\_I,79-84 {79/80}     yadi asandehārthaḥ takāraḥ dakāraḥ api . \\
\noindent \textbf{(Śs\_3-4.1)} KA\_I,22.2-23.23 Ro\_I,79-84 {80/80}     atha mukhasukhārthaḥ takāraḥ dakāraḥ api . \\
\noindent \textbf{(Śs\_3-4.2)} KA\_I,23.24-26.27 Ro\_I,84-93 {1/138}     idam vicāryate . \\
\noindent \textbf{(Śs\_3-4.2)} KA\_I,23.24-26.27 Ro\_I,84-93 {2/138}     ye ete varṇeṣu varṇaikadeśāḥ varṇāntarasamānākṛtayaḥ eteṣām avayavagrahaṇena grahaṇam syāt vā na vā iti . \\
\noindent \textbf{(Śs\_3-4.2)} KA\_I,23.24-26.27 Ro\_I,84-93 {3/138}     kutaḥ punaḥ iyam vicāraṇā . \\
\noindent \textbf{(Śs\_3-4.2)} KA\_I,23.24-26.27 Ro\_I,84-93 {4/138}     iha samudāyāḥ api upadiśyante avayavāḥ api . \\
\noindent \textbf{(Śs\_3-4.2)} KA\_I,23.24-26.27 Ro\_I,84-93 {5/138}     abhyantaraḥ ca samudāye avayavaḥ . \\
\noindent \textbf{(Śs\_3-4.2)} KA\_I,23.24-26.27 Ro\_I,84-93 {6/138}     tat yathā : vṛkṣaḥ pracalan saha avayavaiḥ pracalati . \\
\noindent \textbf{(Śs\_3-4.2)} KA\_I,23.24-26.27 Ro\_I,84-93 {7/138}     tatra samudāyasthasya avayavasya avayavagrahaṇena grahaṇam syāt vā na vā iti jāyate vicāraṇā . \\
\noindent \textbf{(Śs\_3-4.2)} KA\_I,23.24-26.27 Ro\_I,84-93 {8/138}     kaḥ ca atra viśeṣaḥ . \\
\noindent \textbf{(Śs\_3-4.2)} KA\_I,23.24-26.27 Ro\_I,84-93 {9/138}     varṇaikadeśāḥ varṇagrahaṇena cet sandhyakṣre samānākṣaravidhipratiṣedhaḥ . \\
\noindent \textbf{(Śs\_3-4.2)} KA\_I,23.24-26.27 Ro\_I,84-93 {10/138}     varṇaikadeśāḥ varṇagrahaṇena iti cet sandhyakṣare samānākṣarāśrayaḥ vidhiḥ prāpnoti . \\
\noindent \textbf{(Śs\_3-4.2)} KA\_I,23.24-26.27 Ro\_I,84-93 {11/138}     sa pratiṣedhyaḥ : agne , indram . \\
\noindent \textbf{(Śs\_3-4.2)} KA\_I,23.24-26.27 Ro\_I,84-93 {12/138}     vāyo , udakam . \\
\noindent \textbf{(Śs\_3-4.2)} KA\_I,23.24-26.27 Ro\_I,84-93 {13/138}     akaḥ savarṇe dīrghaḥ iti dīrghatvam prāpnoti . \\
\noindent \textbf{(Śs\_3-4.2)} KA\_I,23.24-26.27 Ro\_I,84-93 {14/138}     dīrghe hrasvavidhipratiṣedhaḥ . \\
\noindent \textbf{(Śs\_3-4.2)} KA\_I,23.24-26.27 Ro\_I,84-93 {15/138}     dīrghe hrasvākṣarāśrayaḥ vidhiḥ prāpnoti . \\
\noindent \textbf{(Śs\_3-4.2)} KA\_I,23.24-26.27 Ro\_I,84-93 {16/138}     sa pratiṣedhyaḥ . \\
\noindent \textbf{(Śs\_3-4.2)} KA\_I,23.24-26.27 Ro\_I,84-93 {17/138}     grāmaṇīḥ , ālūya , pralūya . \\
\noindent \textbf{(Śs\_3-4.2)} KA\_I,23.24-26.27 Ro\_I,84-93 {18/138}     hrasvasya piti kṛti tuk bhavati iti tuk prāpnoti . \\
\noindent \textbf{(Śs\_3-4.2)} KA\_I,23.24-26.27 Ro\_I,84-93 {19/138}     na eṣaḥ doṣaḥ . \\
\noindent \textbf{(Śs\_3-4.2)} KA\_I,23.24-26.27 Ro\_I,84-93 {20/138}     ācāryapravṛttiḥ jñāpayati na dīrghe hrasvāśrayaḥ vidhiḥ bhavati iti yat ayam dīrghāt che tukam śāsti . \\
\noindent \textbf{(Śs\_3-4.2)} KA\_I,23.24-26.27 Ro\_I,84-93 {21/138}     na etat asti jñāpakam . \\
\noindent \textbf{(Śs\_3-4.2)} KA\_I,23.24-26.27 Ro\_I,84-93 {22/138}     asti hi anyat etasya vacane prayojanam . \\
\noindent \textbf{(Śs\_3-4.2)} KA\_I,23.24-26.27 Ro\_I,84-93 {23/138}     kim . \\
\noindent \textbf{(Śs\_3-4.2)} KA\_I,23.24-26.27 Ro\_I,84-93 {24/138}     padāntāt vā iti vibhāṣām vakṣyāmi iti . \\
\noindent \textbf{(Śs\_3-4.2)} KA\_I,23.24-26.27 Ro\_I,84-93 {25/138}     yat tarhi yogavibhāgam karoti . \\
\noindent \textbf{(Śs\_3-4.2)} KA\_I,23.24-26.27 Ro\_I,84-93 {26/138}     itarathā hi dīrghāt padāntāt vā iti eva brūyāt . \\
\noindent \textbf{(Śs\_3-4.2)} KA\_I,23.24-26.27 Ro\_I,84-93 {27/138}     iha tarhi khaṭvābhiḥ , mālābhiḥ , ataḥ bihsaḥ ais , iti aisbhāvaḥ prāpnoti . \\
\noindent \textbf{(Śs\_3-4.2)} KA\_I,23.24-26.27 Ro\_I,84-93 {28/138}     taparakaraṇasāmarthyāt na bhaviṣyati . \\
\noindent \textbf{(Śs\_3-4.2)} KA\_I,23.24-26.27 Ro\_I,84-93 {29/138}     iha tarhi yātā vātā , ataḥ lopaḥ ārdhadhātuke iti akāralopaḥ prāpnoti . \\
\noindent \textbf{(Śs\_3-4.2)} KA\_I,23.24-26.27 Ro\_I,84-93 {30/138}     nanu ca atra api taparakaraṇasāmarthyāt eva na bhaviṣyati . \\
\noindent \textbf{(Śs\_3-4.2)} KA\_I,23.24-26.27 Ro\_I,84-93 {31/138}     asti hi anyat taparakaraṇe prayojanam . \\
\noindent \textbf{(Śs\_3-4.2)} KA\_I,23.24-26.27 Ro\_I,84-93 {32/138}     kim . \\
\noindent \textbf{(Śs\_3-4.2)} KA\_I,23.24-26.27 Ro\_I,84-93 {33/138}     sarvasya lopaḥ mā bhūt iti . \\
\noindent \textbf{(Śs\_3-4.2)} KA\_I,23.24-26.27 Ro\_I,84-93 {34/138}     atha kriyamāṇe api tapare parasya lope kṛte pūrvasya kasmāt na bhavati . \\
\noindent \textbf{(Śs\_3-4.2)} KA\_I,23.24-26.27 Ro\_I,84-93 {35/138}     paralopasya sthānivadbhāvāt asiddhatvāt ca . \\
\noindent \textbf{(Śs\_3-4.2)} KA\_I,23.24-26.27 Ro\_I,84-93 {36/138}     evam tarhi ācāryapravṛttiḥ jñāpayati nākārasthasya akārasya lopaḥ bhavati iti yat ayam ātaḥ anupasarge kaḥ iti kakāram anubandham karoti . \\
\noindent \textbf{(Śs\_3-4.2)} KA\_I,23.24-26.27 Ro\_I,84-93 {37/138}     katham kṛtvā jñāpakam . \\
\noindent \textbf{(Śs\_3-4.2)} KA\_I,23.24-26.27 Ro\_I,84-93 {38/138}     kitkaraṇe etat prayojanam kiti iti ākārlopaḥ yathā syāt iti . \\
\noindent \textbf{(Śs\_3-4.2)} KA\_I,23.24-26.27 Ro\_I,84-93 {39/138}     yadi ca ākārasthasya api akārlopaḥ syāt kitkaraṇam anarthakam syāt . \\
\noindent \textbf{(Śs\_3-4.2)} KA\_I,23.24-26.27 Ro\_I,84-93 {40/138}     parasya akārasya lope kṛte dvayoḥ akārayoḥ pararūpe hi siddham rūpam syāt: godaḥ , kambaladaḥ iti . \\
\noindent \textbf{(Śs\_3-4.2)} KA\_I,23.24-26.27 Ro\_I,84-93 {41/138}     paśyati tu ācāryaḥ nākārasthasya akārasya lopaḥ bhavati iti . \\
\noindent \textbf{(Śs\_3-4.2)} KA\_I,23.24-26.27 Ro\_I,84-93 {42/138}     ataḥ kakāram anubandham karoti . \\
\noindent \textbf{(Śs\_3-4.2)} KA\_I,23.24-26.27 Ro\_I,84-93 {43/138}     na etat asti jñāpakam . \\
\noindent \textbf{(Śs\_3-4.2)} KA\_I,23.24-26.27 Ro\_I,84-93 {44/138}     uttarārtham etat syāt . \\
\noindent \textbf{(Śs\_3-4.2)} KA\_I,23.24-26.27 Ro\_I,84-93 {45/138}     tundaśokayoḥ parimṛjāpanudoḥ iti . \\
\noindent \textbf{(Śs\_3-4.2)} KA\_I,23.24-26.27 Ro\_I,84-93 {46/138}     yat tarhi gāpoḥ ṭhak iti ananyārtham kakāram anubandham karoti . \\
\noindent \textbf{(Śs\_3-4.2)} KA\_I,23.24-26.27 Ro\_I,84-93 {47/138}     ekavarṇavat ca . \\
\noindent \textbf{(Śs\_3-4.2)} KA\_I,23.24-26.27 Ro\_I,84-93 {48/138}     ekavarṇavat ca dīrghaḥ bhavati iti vaktavyam . \\
\noindent \textbf{(Śs\_3-4.2)} KA\_I,23.24-26.27 Ro\_I,84-93 {49/138}     kim prayojanam . \\
\noindent \textbf{(Śs\_3-4.2)} KA\_I,23.24-26.27 Ro\_I,84-93 {50/138}     vācā tarati iti dvyajlakṣaṇaḥ ṭhan mā bhūt iti . \\
\noindent \textbf{(Śs\_3-4.2)} KA\_I,23.24-26.27 Ro\_I,84-93 {51/138}     iha ca vācaḥ nimittam , tasya nimittam saṃyogotpāttau iti dvyajlakṣaṇaḥ yat mā bhūt iti . \\
\noindent \textbf{(Śs\_3-4.2)} KA\_I,23.24-26.27 Ro\_I,84-93 {52/138}     atra api gonaugrahaṇam jñāpakam dīrghāt dvyajlakṣaṇaḥ vidhiḥ na bhavati iti . \\
\noindent \textbf{(Śs\_3-4.2)} KA\_I,23.24-26.27 Ro\_I,84-93 {53/138}     ayam tu sarveṣām parihāraḥ : na avyapavṛktasya avayave tadvidhiḥ yathā dravyeṣu . \\
\noindent \textbf{(Śs\_3-4.2)} KA\_I,23.24-26.27 Ro\_I,84-93 {54/138}     na avyapavṛktasya avayavsya avayavāśrayaḥ vidhiḥ bhavati yathā dravyeṣu . \\
\noindent \textbf{(Śs\_3-4.2)} KA\_I,23.24-26.27 Ro\_I,84-93 {55/138}     tat yathā dravyeṣu : saptadaśa sāmidhenyaḥ bhavanti iti na saptadaśāratnimātram kāṣṭham agnau abhyādhīyate . \\
\noindent \textbf{(Śs\_3-4.2)} KA\_I,23.24-26.27 Ro\_I,84-93 {56/138}     viṣamaḥ upanyāsaḥ . \\
\noindent \textbf{(Śs\_3-4.2)} KA\_I,23.24-26.27 Ro\_I,84-93 {57/138}     pratyṛcam ca eva hi tat karma codyate asambhavaḥ ca agnau vedyām ca . \\
\noindent \textbf{(Śs\_3-4.2)} KA\_I,23.24-26.27 Ro\_I,84-93 {58/138}     yathā tarhi saptadaśa prādeśamātrīḥ aśvatthīḥ samidhaḥ abhyādadhīta iti na saptadaśaprādeśamātram kāṣṭham abhyādhīyate . \\
\noindent \textbf{(Śs\_3-4.2)} KA\_I,23.24-26.27 Ro\_I,84-93 {59/138}     atra api pratipravaṇam ca etat karma codyate tulyaḥ ca asambhavaḥ agnau vedyām ca . \\
\noindent \textbf{(Śs\_3-4.2)} KA\_I,23.24-26.27 Ro\_I,84-93 {60/138}     yathā tarhi tailam na vikretavyam , māṃsam na vikretavyam iti . \\
\noindent \textbf{(Śs\_3-4.2)} KA\_I,23.24-26.27 Ro\_I,84-93 {61/138}     vyapavṛtkam ca na vikrīyate , avyapavṛktam ca gāvaḥ ca sarṣapāḥ ca vikriyante . \\
\noindent \textbf{(Śs\_3-4.2)} KA\_I,23.24-26.27 Ro\_I,84-93 {62/138}     tathā lomanakham spṛṣṭvā śaucam kartavyam iti , vyapavṛktam spṛṣṭvā niyogataḥ kartavyam avyapavṛkte kāmacāraḥ . \\
\noindent \textbf{(Śs\_3-4.2)} KA\_I,23.24-26.27 Ro\_I,84-93 {63/138}     yatra tarhi vyapavargaḥ asti . \\
\noindent \textbf{(Śs\_3-4.2)} KA\_I,23.24-26.27 Ro\_I,84-93 {64/138}     kva ca vyapavargaḥ asti . \\
\noindent \textbf{(Śs\_3-4.2)} KA\_I,23.24-26.27 Ro\_I,84-93 {65/138}     sandhyakṣareṣu . \\
\noindent \textbf{(Śs\_3-4.2)} KA\_I,23.24-26.27 Ro\_I,84-93 {66/138}     sandhyakṣareṣu vivṛtatvāt . \\
\noindent \textbf{(Śs\_3-4.2)} KA\_I,23.24-26.27 Ro\_I,84-93 {67/138}     yat atra avarṇam vivṛtataram tat anyasmāt avarṇāt . \\
\noindent \textbf{(Śs\_3-4.2)} KA\_I,23.24-26.27 Ro\_I,84-93 {68/138}     ye*api ivarṇovarṇe vivṛtatare te*anyābhyām ivarṇovarṇābhyām . \\
\noindent \textbf{(Śs\_3-4.2)} KA\_I,23.24-26.27 Ro\_I,84-93 {69/138}     athavā punaḥ na gṛhyante . \\
\noindent \textbf{(Śs\_3-4.2)} KA\_I,23.24-26.27 Ro\_I,84-93 {70/138}     agrahaṇam cet nuḍvidhilādeśavināmeṣu ṛkāragrahaṇam . \\
\noindent \textbf{(Śs\_3-4.2)} KA\_I,23.24-26.27 Ro\_I,84-93 {71/138}     agrahaṇam cet nuḍvidhilādeśavināmeṣu ṛkārasya grahaṇam kartavyam . \\
\noindent \textbf{(Śs\_3-4.2)} KA\_I,23.24-26.27 Ro\_I,84-93 {72/138}     tasmāt nuṭ dvihalaḥ , ṛkāre ca iti vaktavyam iha api yathā syāt : ānṛdhatuḥ , ānṛdhuḥ iti . \\
\noindent \textbf{(Śs\_3-4.2)} KA\_I,23.24-26.27 Ro\_I,84-93 {73/138}     yasya punaḥ gṛhyante dvihalaḥ iti eva tasya siddham . \\
\noindent \textbf{(Śs\_3-4.2)} KA\_I,23.24-26.27 Ro\_I,84-93 {74/138}     yasya api na gṛhyante tasya api eṣaḥ na doṣaḥ . \\
\noindent \textbf{(Śs\_3-4.2)} KA\_I,23.24-26.27 Ro\_I,84-93 {75/138}     dvihalgrahaṇam na kariṣyate . \\
\noindent \textbf{(Śs\_3-4.2)} KA\_I,23.24-26.27 Ro\_I,84-93 {76/138}     tasmāt nuṭ bhavati iti eva . \\
\noindent \textbf{(Śs\_3-4.2)} KA\_I,23.24-26.27 Ro\_I,84-93 {77/138}     yadi na kriyate āṭatuḥ , āṭuḥ iti atra api prāpnoti . \\
\noindent \textbf{(Śs\_3-4.2)} KA\_I,23.24-26.27 Ro\_I,84-93 {78/138}     aśnotigrahaṇam niyamārtham bhaviṣyati . \\
\noindent \textbf{(Śs\_3-4.2)} KA\_I,23.24-26.27 Ro\_I,84-93 {79/138}     aśnoteḥ eva avarṇopadhasya na anyasya avarṇopadhasya iti . \\
\noindent \textbf{(Śs\_3-4.2)} KA\_I,23.24-26.27 Ro\_I,84-93 {80/138}     lādeśe ca ṛkāragrahaṇam kartavyam . \\
\noindent \textbf{(Śs\_3-4.2)} KA\_I,23.24-26.27 Ro\_I,84-93 {81/138}     kṛpaḥ raḥ laḥ , ṛkārasya ca iti vaktavyam iha api yathā syāt : kḷptaḥ , kḷptavān iti . \\
\noindent \textbf{(Śs\_3-4.2)} KA\_I,23.24-26.27 Ro\_I,84-93 {82/138}     yasya punaḥ gṛhyante raḥ iti eva tasya siddham . \\
\noindent \textbf{(Śs\_3-4.2)} KA\_I,23.24-26.27 Ro\_I,84-93 {83/138}     yasya api na gṛhyante tasya api eṣaḥ na doṣaḥ . \\
\noindent \textbf{(Śs\_3-4.2)} KA\_I,23.24-26.27 Ro\_I,84-93 {84/138}     ṛkāraḥ api atra nirdiśyate . \\
\noindent \textbf{(Śs\_3-4.2)} KA\_I,23.24-26.27 Ro\_I,84-93 {85/138}     katham . \\
\noindent \textbf{(Śs\_3-4.2)} KA\_I,23.24-26.27 Ro\_I,84-93 {86/138}     avibhaktikaḥ nirdeśaḥ . \\
\noindent \textbf{(Śs\_3-4.2)} KA\_I,23.24-26.27 Ro\_I,84-93 {87/138}     kṛpa , uḥ , raḥ , laḥ kṛpo ro laḥ iti . \\
\noindent \textbf{(Śs\_3-4.2)} KA\_I,23.24-26.27 Ro\_I,84-93 {88/138}     atha vā ubhayataḥ sphoṭamātram nirdiśyate . \\
\noindent \textbf{(Śs\_3-4.2)} KA\_I,23.24-26.27 Ro\_I,84-93 {89/138}     raśruteḥ laśrutiḥ bhavati iti . \\
\noindent \textbf{(Śs\_3-4.2)} KA\_I,23.24-26.27 Ro\_I,84-93 {90/138}     vināme ṛkāragrahaṇam kartavyam . \\
\noindent \textbf{(Śs\_3-4.2)} KA\_I,23.24-26.27 Ro\_I,84-93 {91/138}     raṣābhyām naḥ ṇaḥ samānapade , ṛkārāt ca iti vaktayvam iha api yathā syāt : mātṝṇām , pitṝṇām iti . \\
\noindent \textbf{(Śs\_3-4.2)} KA\_I,23.24-26.27 Ro\_I,84-93 {92/138}     yasya punaḥ gṛhyante raṣābhyām iti eva tasya siddham . \\
\noindent \textbf{(Śs\_3-4.2)} KA\_I,23.24-26.27 Ro\_I,84-93 {93/138}     na sidhyati . \\
\noindent \textbf{(Śs\_3-4.2)} KA\_I,23.24-26.27 Ro\_I,84-93 {94/138}     yat tat rephāt param bhakteḥ tena vyavahitatvāt na prāpnoti . \\
\noindent \textbf{(Śs\_3-4.2)} KA\_I,23.24-26.27 Ro\_I,84-93 {95/138}     mā bhūt evam . \\
\noindent \textbf{(Śs\_3-4.2)} KA\_I,23.24-26.27 Ro\_I,84-93 {96/138}     aḍvyavāye iti eva siddham . \\
\noindent \textbf{(Śs\_3-4.2)} KA\_I,23.24-26.27 Ro\_I,84-93 {97/138}     na sidhyati . \\
\noindent \textbf{(Śs\_3-4.2)} KA\_I,23.24-26.27 Ro\_I,84-93 {98/138}     varṇaikadeśāḥ ke varṇagrahaṇena gṛhyante . \\
\noindent \textbf{(Śs\_3-4.2)} KA\_I,23.24-26.27 Ro\_I,84-93 {99/138}     ye vyapavṛktāḥ api varṇāḥ bhavanti . \\
\noindent \textbf{(Śs\_3-4.2)} KA\_I,23.24-26.27 Ro\_I,84-93 {100/138}     yat ca api rephāt param bhakteḥ na tat kva cit api vyapavṛktam dṛśyate . \\
\noindent \textbf{(Śs\_3-4.2)} KA\_I,23.24-26.27 Ro\_I,84-93 {101/138}     evam tarhi yogavibhāgaḥ kariṣyate . \\
\noindent \textbf{(Śs\_3-4.2)} KA\_I,23.24-26.27 Ro\_I,84-93 {102/138}     raṣābhyām naḥ ṇaḥ samānapade . \\
\noindent \textbf{(Śs\_3-4.2)} KA\_I,23.24-26.27 Ro\_I,84-93 {103/138}     tataḥ vyavāye . \\
\noindent \textbf{(Śs\_3-4.2)} KA\_I,23.24-26.27 Ro\_I,84-93 {104/138}     vyavāye ca raṣābhyām naḥ ṇaḥ bhavati iti . \\
\noindent \textbf{(Śs\_3-4.2)} KA\_I,23.24-26.27 Ro\_I,84-93 {105/138}     tataḥ aṭkupvāṅnumbhiḥ iti . \\
\noindent \textbf{(Śs\_3-4.2)} KA\_I,23.24-26.27 Ro\_I,84-93 {106/138}     idam idānīm kimartham . \\
\noindent \textbf{(Śs\_3-4.2)} KA\_I,23.24-26.27 Ro\_I,84-93 {107/138}     niyamārtham . \\
\noindent \textbf{(Śs\_3-4.2)} KA\_I,23.24-26.27 Ro\_I,84-93 {108/138}     etaiḥ eva ākṣarasamamnāyikaiḥ vyavāye na anyaiḥ iti . \\
\noindent \textbf{(Śs\_3-4.2)} KA\_I,23.24-26.27 Ro\_I,84-93 {109/138}     yasya api na gṛhyante tasya api eṣaḥ na doṣaḥ . \\
\noindent \textbf{(Śs\_3-4.2)} KA\_I,23.24-26.27 Ro\_I,84-93 {110/138}     ācāryapravṛttiḥ jñāpayati bhavati ṛkārāt naḥ ṇatvam iti yat ayam kṣubhādiṣu nṛnamanaśabdam paṭhati . \\
\noindent \textbf{(Śs\_3-4.2)} KA\_I,23.24-26.27 Ro\_I,84-93 {111/138}     na etat asti jñāpakam . \\
\noindent \textbf{(Śs\_3-4.2)} KA\_I,23.24-26.27 Ro\_I,84-93 {112/138}     vṛddhyartham etat syāt : nārnamaniḥ . \\
\noindent \textbf{(Śs\_3-4.2)} KA\_I,23.24-26.27 Ro\_I,84-93 {113/138}     yat tarhi tṛpnotiśabdam paṭhati . \\
\noindent \textbf{(Śs\_3-4.2)} KA\_I,23.24-26.27 Ro\_I,84-93 {114/138}     yat ca api nṛnamanaśabdam paṭhati . \\
\noindent \textbf{(Śs\_3-4.2)} KA\_I,23.24-26.27 Ro\_I,84-93 {115/138}     nanu ca uktam vṛddhyartham etat syāt iti . \\
\noindent \textbf{(Śs\_3-4.2)} KA\_I,23.24-26.27 Ro\_I,84-93 {116/138}     bahiraṅgā vṛddhiḥ . \\
\noindent \textbf{(Śs\_3-4.2)} KA\_I,23.24-26.27 Ro\_I,84-93 {117/138}     antaraṅgam ṇatvam . \\
\noindent \textbf{(Śs\_3-4.2)} KA\_I,23.24-26.27 Ro\_I,84-93 {118/138}     asiddham bahiraṅgam antaraṅge . \\
\noindent \textbf{(Śs\_3-4.2)} KA\_I,23.24-26.27 Ro\_I,84-93 {119/138}     atha vā upariṣṭāt yogavibhāgaḥ kariṣyate . \\
\noindent \textbf{(Śs\_3-4.2)} KA\_I,23.24-26.27 Ro\_I,84-93 {120/138}     ṛto naḥ ṇaḥ bhavati . \\
\noindent \textbf{(Śs\_3-4.2)} KA\_I,23.24-26.27 Ro\_I,84-93 {121/138}     tataḥ chandasi avagrahāt . \\
\noindent \textbf{(Śs\_3-4.2)} KA\_I,23.24-26.27 Ro\_I,84-93 {122/138}     ṛtaḥ iti eva . \\
\noindent \textbf{(Śs\_3-4.2)} KA\_I,23.24-26.27 Ro\_I,84-93 {123/138}     plutau aicaḥ idutau . \\
\noindent \textbf{(Śs\_3-4.2)} KA\_I,23.24-26.27 Ro\_I,84-93 {124/138}     etat ca vaktavyam . \\
\noindent \textbf{(Śs\_3-4.2)} KA\_I,23.24-26.27 Ro\_I,84-93 {125/138}     yasya punaḥ gṛhyante guroḥ ṭeḥ iti eva plutyā tasya siddham . \\
\noindent \textbf{(Śs\_3-4.2)} KA\_I,23.24-26.27 Ro\_I,84-93 {126/138}     yasya api na gṛhyante tasya api eṣaḥ na doṣaḥ . \\
\noindent \textbf{(Śs\_3-4.2)} KA\_I,23.24-26.27 Ro\_I,84-93 {127/138}     kriyate etat nyāse eva . \\
\noindent \textbf{(Śs\_3-4.2)} KA\_I,23.24-26.27 Ro\_I,84-93 {128/138}     tulyarūpe saṃyoge dvivyañjanavidhiḥ . \\
\noindent \textbf{(Śs\_3-4.2)} KA\_I,23.24-26.27 Ro\_I,84-93 {129/138}     tulyarūpe saṃyoge dvivyañjanāśrayaḥ vidhiḥ na sidhyati : kukkuṭaḥ , pippalaḥ , pittam iti . \\
\noindent \textbf{(Śs\_3-4.2)} KA\_I,23.24-26.27 Ro\_I,84-93 {130/138}     yasya punaḥ gṛhyante tasya dvau kakārau dvau pakārau dvau takārau . \\
\noindent \textbf{(Śs\_3-4.2)} KA\_I,23.24-26.27 Ro\_I,84-93 {131/138}     yasya api na gṛhyante tasya api dvau kakārau dvau pakārau dvau takārau . \\
\noindent \textbf{(Śs\_3-4.2)} KA\_I,23.24-26.27 Ro\_I,84-93 {132/138}     katham . \\
\noindent \textbf{(Śs\_3-4.2)} KA\_I,23.24-26.27 Ro\_I,84-93 {133/138}     mātrākālaḥ atra gamyate . \\
\noindent \textbf{(Śs\_3-4.2)} KA\_I,23.24-26.27 Ro\_I,84-93 {134/138}     na ca mātrikam vyañjanam asti . \\
\noindent \textbf{(Śs\_3-4.2)} KA\_I,23.24-26.27 Ro\_I,84-93 {135/138}     anupadiṣṭam sat katham śakyam vijñātum . \\
\noindent \textbf{(Śs\_3-4.2)} KA\_I,23.24-26.27 Ro\_I,84-93 {136/138}     yadi api tāvat atra etat śakyate vaktum yatra etat na asti aṇ savarṇān gṛhṇāti iti iha tu katham sayŚyantā savŚvatsaraḥ yalŚ lokam talŚ lokam iti yatra etat asti aṇ savarṇān gṛhṇāti iti . \\
\noindent \textbf{(Śs\_3-4.2)} KA\_I,23.24-26.27 Ro\_I,84-93 {137/138}     atra api mātrākalaḥ gṛhyate na ca mātrikam vyañjanam asti . \\
\noindent \textbf{(Śs\_3-4.2)} KA\_I,23.24-26.27 Ro\_I,84-93 {138/138}     anupadiṣṭam sat katham śakyam pratipattum . \\
\noindent \textbf{(Śs\_5.1)} sarve varṇāḥ sakṛt upadiṣṭāḥ . \\
\noindent \textbf{(Śs\_5.1)} ayam hakāraḥ dviḥ upadiśyate pūrvaḥ ca paraḥ ca . \\
\noindent \textbf{(Śs\_5.1)} yadi punaḥ pūrvaḥ eva upadiśyeta paraḥ eva vā . \\
\noindent \textbf{(Śs\_5.1)} kaḥ ca atra viśeṣaḥ . \\
\noindent \textbf{(Śs\_5.1)} hakārasya paropadeśe aḍgrahaṇeṣu hagrahaṇam . \\
\noindent \textbf{(Śs\_5.1)} hakārasya paropadeśe aḍgrahaṇeṣu hagrahaṇam kartavyam . \\
\noindent \textbf{(Śs\_5.1)} ātaḥ aṭi nityam , śaḥ chaḥ aṭi , dīrghāt aṭi samānapade . \\
\noindent \textbf{(Śs\_5.1)} hakāre ca iti vaktavyam iha api yathā syāt : mahāŚ hi saḥ . \\
\noindent \textbf{(Śs\_5.1)} uttve ca . uttve ca hakāragrahaṇam kartavyam . \\
\noindent \textbf{(Śs\_5.1)} ataḥ roḥ aplutāt aplute , haśi ca . \\
\noindent \textbf{(Śs\_5.1)} hakāre ca iti vaktavyam iha api yathā syāt : puruṣaḥ hasati , brāhmaṇaḥ hasati iti . \\
\noindent \textbf{(Śs\_5.1)} astu tarhi pūrvopadeśaḥ . \\
\noindent \textbf{(Śs\_5.1)} pūrvopadeśe kittvakseḍvidhayaḥ jhalgrahaṇāni ca . \\
\noindent \textbf{(Śs\_5.1)} yadi pūrvopadeśaḥ kittvam vidheyam . \\
\noindent \textbf{(Śs\_5.1)} snihitvā snehitvā sisnihiṣati sisnehiṣati . \\
\noindent \textbf{(Śs\_5.1)} ralaḥ vyupadhāt halādeḥ iti kittvam na prāpnoti . \\
\noindent \textbf{(Śs\_5.1)} ksavidhiḥ . \\
\noindent \textbf{(Śs\_5.1)} ksaḥ ca vidheyaḥ . \\
\noindent \textbf{(Śs\_5.1)} adhukṣat alikṣat . \\
\noindent \textbf{(Śs\_5.1)} śalaḥ igupadhāt aniṭaḥ ksaḥ iti ksaḥ na prāpnoti . \\
\noindent \textbf{(Śs\_5.1)} iḍvidhiḥ . \\
\noindent \textbf{(Śs\_5.1)} iṭ ca vidheyaḥ . \\
\noindent \textbf{(Śs\_5.1)} rudihi svapihi . \\
\noindent \textbf{(Śs\_5.1)} valādilakṣaṇaḥ iṭ na prāpnoti. jhalgrahaṇāni ca . \\
\noindent \textbf{(Śs\_5.1)} kim . \\
\noindent \textbf{(Śs\_5.1)} ahakārāṇi syuḥ . \\
\noindent \textbf{(Śs\_5.1)} tatra kaḥ doṣaḥ . \\
\noindent \textbf{(Śs\_5.1)} jhalaḥ jhali iti iha na syāt : adāgdhām adāgdham . \\
\noindent \textbf{(Śs\_5.1)} tasmāt pūrvaḥ ca upadeṣṭavyaḥ paraḥ ca . \\
\noindent \textbf{(Śs\_5.1)} yadi ca kim cit anyatra api upadeśe prayojanam asti tatra api upadeśaḥ kartavyaḥ . \\
\noindent \textbf{(Śs\_5.2)} idam vicāryate : ayam rephaḥ yakāravakārābhyām pūrvaḥ eva upadiśyeta ha ra ya vaṭ iti paraḥ eva vā yathānyāsam iti . \\
\noindent \textbf{(Śs\_5.2)} kaḥ ca atra viśeṣaḥ . \\
\noindent \textbf{(Śs\_5.2)} rephasya paropadeśe anunāsikadvirvacanaparasavarṇapratiṣedhaḥ . \\
\noindent \textbf{(Śs\_5.2)} rephasya paropadeśe anunāsikadvirvacanaparasavarṇānām pratiṣedhaḥ vaktavyaḥ . \\
\noindent \textbf{(Śs\_5.2)} anunāsikasya : svaḥ nayati , prātaḥ nayati iti yaraḥ anunāsike anunāsikaḥ vā iti anunāsikaḥ prāpnoti . \\
\noindent \textbf{(Śs\_5.2)} dvirvacanasya : bhadrahradaḥ , madrahradaḥ iti yaraḥ iti dvirvacanam prāpnoti . \\
\noindent \textbf{(Śs\_5.2)} parasavarṇasya : kuṇḍam rathena , vanam rathena . \\
\noindent \textbf{(Śs\_5.2)} anusvārasya yayi iti parasavarṇaḥ prāpnoti . \\
\noindent \textbf{(Śs\_5.2)} astu tarhi pūrvopadeśaḥ . \\
\noindent \textbf{(Śs\_5.2)} pūrvopadeśe kittvapratiṣedhaḥ vyalopavacanam ca . \\
\noindent \textbf{(Śs\_5.2)} yadi pūrvopadeśaḥ kittvam pratiṣedhyam . \\
\noindent \textbf{(Śs\_5.2)} devitvā dideviṣati . \\
\noindent \textbf{(Śs\_5.2)} ralaḥ vyupadhāt iti kittvam prāpnoti . \\
\noindent \textbf{(Śs\_5.2)} na eṣaḥ doṣaḥ . \\
\noindent \textbf{(Śs\_5.2)} na evam vijñāyate ralaḥ vyupadhāt iti . \\
\noindent \textbf{(Śs\_5.2)} kim tarhi . \\
\noindent \textbf{(Śs\_5.2)} ralaḥ avvyupadhāt iti . \\
\noindent \textbf{(Śs\_5.2)} kim idam avvyupadhāt iti . \\
\noindent \textbf{(Śs\_5.2)} avakārāntāt vyuvpadhāt avvyupadhāt iti . \\
\noindent \textbf{(Śs\_5.2)} vyalopavacanam ca . \\
\noindent \textbf{(Śs\_5.2)} vyoḥ ca lopaḥ vaktavyaḥ : gaudheraḥ , paceran yajeran . \\
\noindent \textbf{(Śs\_5.2)} jīveḥ radānuk : jīradānuḥ . \\
\noindent \textbf{(Śs\_5.2)} vali iti lopaḥ na prāpnoti . \\
\noindent \textbf{(Śs\_5.2)} na eṣaḥ doṣaḥ . \\
\noindent \textbf{(Śs\_5.2)} rephaḥ api atra nirdiśyate . \\
\noindent \textbf{(Śs\_5.2)} lopaḥ vyoḥ vali iti rephe ca vali ca iti . \\
\noindent \textbf{(Śs\_5.2)} atha vā punaḥ astu paropadeśaḥ . \\
\noindent \textbf{(Śs\_5.2)} nanu ca uktam rephasya paropadeśe anunāsikadvirvacanaparasavarṇapratiṣedhaḥ iti . \\
\noindent \textbf{(Śs\_5.2)} anunāsikaparasavarṇayoḥ tāvat pratiṣedhaḥ na vaktavyaḥ . \\
\noindent \textbf{(Śs\_5.2)} rephoṣmaṇām savarṇāḥ na santi . \\
\noindent \textbf{(Śs\_5.2)} dvirvacane api na imau rahau kāryiṇau dvirvacanasya . \\
\noindent \textbf{(Śs\_5.2)} kim tarhi . \\
\noindent \textbf{(Śs\_5.2)} nimittam imau rahau dvirvacanasya . \\
\noindent \textbf{(Śs\_5.2)} tat yathā . \\
\noindent \textbf{(Śs\_5.2)} brāhmaṇāḥ bhojyantām . \\
\noindent \textbf{(Śs\_5.2)} māṭharakauṇḍinyau pariveviṣṭām iti. na idānīm tau bhuñjāte . \\
\noindent \textbf{(Śs\_5.3)} idam vicāryate . \\
\noindent \textbf{(Śs\_5.3)} ime ayogavāhāḥ na kva cit upadiśyante śrūyante ca . \\
\noindent \textbf{(Śs\_5.3)} teṣām kāryārthaḥ upadeśaḥ kartavyaḥ . \\
\noindent \textbf{(Śs\_5.3)} ke punaḥ ayogavāhāḥ . \\
\noindent \textbf{(Śs\_5.3)} visarjanīyajihvāmūlīyopadhmānīyānusvārānunāsikyayamāḥ . \\
\noindent \textbf{(Śs\_5.3)} katham punaḥ ayogavāhāḥ . \\
\noindent \textbf{(Śs\_5.3)} yat ayuktāḥ vahanti anupadiṣṭāḥ ca śrūyante . \\
\noindent \textbf{(Śs\_5.3)} kva punaḥ eṣām upadeśaḥ kartavyaḥ . \\
\noindent \textbf{(Śs\_5.3)} ayogavāhānām aṭsu ṇatvam . \\
\noindent \textbf{(Śs\_5.3)} ayogavāhānām aṭsu upadeśaḥ kartavyaḥ . \\
\noindent \textbf{(Śs\_5.3)} kim prayojanam . \\
\noindent \textbf{(Śs\_5.3)} ṇatvam . \\
\noindent \textbf{(Śs\_5.3)} uraḥkeṇa , uraḥpeṇa : aḍvyavāye iti ṇatvam siddham bhavati . \\
\noindent \textbf{(Śs\_5.3)} śarṣu jaśbhāvaṣatve . \\
\noindent \textbf{(Śs\_5.3)} śarṣu upadeśaḥ kartavyaḥ . \\
\noindent \textbf{(Śs\_5.3)} kim prayojanam . \\
\noindent \textbf{(Śs\_5.3)} jaśbhāvaṣatve . \\
\noindent \textbf{(Śs\_5.3)} ayam ubjiḥ upadhmānīyopadhaḥ paṭhyate . \\
\noindent \textbf{(Śs\_5.3)} tasya jaśtve kṛte ubjitā ubjitum iti etat rūpam yathā syāt . \\
\noindent \textbf{(Śs\_5.3)} yadi ubjiḥ upadhmānīyopadhaḥ paṭhyate ubjijiṣati iti upadhmānīyādeḥ eva dvirvacanam prāpnoti. dakāropadhe punaḥ nandrāḥ saṃyogādayaḥ iti pratiṣedhasḥ siddhaḥ bhavati . \\
\noindent \textbf{(Śs\_5.3)} yadi dakāropadhaḥ paṭhyate kā rūpasiddhiḥ : ubjitā ubjitum iti . \\
\noindent \textbf{(Śs\_5.3)} asiddhe bhaḥ udjeḥ . \\
\noindent \textbf{(Śs\_5.3)} idam asti stoḥ ścunā ścuḥ iti . \\
\noindent \textbf{(Śs\_5.3)} tataḥ vakṣyāmi bhaḥ udjeḥ . \\
\noindent \textbf{(Śs\_5.3)} udjeḥ ścunā sannipāte bhaḥ bhavati iti . \\
\noindent \textbf{(Śs\_5.3)} tat tarhi vaktavyam . \\
\noindent \textbf{(Śs\_5.3)} na vaktavyam . \\
\noindent \textbf{(Śs\_5.3)} nipātanāt eva siddham . \\
\noindent \textbf{(Śs\_5.3)} kim nipātanam . \\
\noindent \textbf{(Śs\_5.3)} bhujanyubjau paṇyupatapayoḥ iti . \\
\noindent \textbf{(Śs\_5.3)} iha api tarhi prāpnoti abhyudgaḥ , samudgaḥ iti . \\
\noindent \textbf{(Śs\_5.3)} akutvaviṣaye tat nipātanam . \\
\noindent \textbf{(Śs\_5.3)} atha vā na etat ubjeḥ rūpam . \\
\noindent \textbf{(Śs\_5.3)} gameḥ dvyuparsargāt ḍaḥ vidhīyate . \\
\noindent \textbf{(Śs\_5.3)} abhyudgataḥ abhyudgaḥ . \\
\noindent \textbf{(Śs\_5.3)} samudgataḥ samudgaḥ iti . \\
\noindent \textbf{(Śs\_5.3)} ṣatvam ca prayojanam . \\
\noindent \textbf{(Śs\_5.3)} sarpiḥṣu dhanuḥṣu . \\
\noindent \textbf{(Śs\_5.3)} śarvyavāye iti ṣatvam siddham bhavati . \\
\noindent \textbf{(Śs\_5.3)} numvisarjanīyaśarvyavāye api iti visarjanīyagrahaṇam na kartavyam bhavati . \\
\noindent \textbf{(Śs\_5.3)} numaḥ ca api tarhi grahaṇam śakyam akartum . \\
\noindent \textbf{(Śs\_5.3)} katham sarpīṃṣi dhanūṃṣi . \\
\noindent \textbf{(Śs\_5.3)} anusvāre kṛte śarvyavāye iti eva siddham . \\
\noindent \textbf{(Śs\_5.3)} avaśyam numaḥ grahaṇam kartavyam anusvāraviśeṣaṇam numgrahaṇam numaḥ yaḥ anusvāraḥ tatra yathā syāt iha mā bhūt : puṃsu iti . \\
\noindent \textbf{(Śs\_5.3)} atha vā aviśeṣeṇa upadeśaḥ kartavyaḥ . \\
\noindent \textbf{(Śs\_5.3)} kim prayojanam . \\
\noindent \textbf{(Śs\_5.3)} aviśeṣeṇa saṃyogopadhāsañjñālontyadvirvacanasthānivadbhāvapratiṣedhāḥ . \\
\noindent \textbf{(Śs\_5.3)} aviśeṣeṇa saṃyogasañjñā prayojanam . \\
\noindent \textbf{(Śs\_5.3)} ūbjaka . \\
\noindent \textbf{(Śs\_5.3)} halaḥ anantarāḥ saṃyogaḥ iti saṃyogasañjñā saṃyoge guru iti gurusañjñā guroḥ iti plutaḥ bhavati . \\
\noindent \textbf{(Śs\_5.3)} upadhāsañjñā ca prayojanam . \\
\noindent \textbf{(Śs\_5.3)} duṣkṛtam , niṣkṛtam , niṣpītam , duṣpītam . \\
\noindent \textbf{(Śs\_5.3)} idudupadhasya ca apratyayasya iti ṣatvam siddham bhavati . \\
\noindent \textbf{(Śs\_5.3)} na etat asti prayojanam . \\
\noindent \textbf{(Śs\_5.3)} na idudupadhagrahaṇanena visarjanīyaḥ viśeṣyate . \\
\noindent \textbf{(Śs\_5.3)} kim tarhi . \\
\noindent \textbf{(Śs\_5.3)} sakāraḥ viśeṣyate . \\
\noindent \textbf{(Śs\_5.3)} idudupadhasya sakārasya yaḥ visarjanīyaḥ iti . \\
\noindent \textbf{(Śs\_5.3)} atha vā upadhāgrahaṇam na kariṣyate . \\
\noindent \textbf{(Śs\_5.3)} idudbhyām tu visarjanīyam viśeṣayiṣyāmaḥ . \\
\noindent \textbf{(Śs\_5.3)} idudbhyām uttarasya visarjanīyasya iti . \\
\noindent \textbf{(Śs\_5.3)} alaḥ antyavidhiḥ prayojanam . \\
\noindent \textbf{(Śs\_5.3)} vṛkṣaḥ tarati . \\
\noindent \textbf{(Śs\_5.3)} plakṣaḥ tarati . \\
\noindent \textbf{(Śs\_5.3)} alaḥ antyasya vidhayaḥ bhavanti iti alaḥ antyasya satvam siddham bhavati . \\
\noindent \textbf{(Śs\_5.3)} etat api na asti prayojanam . \\
\noindent \textbf{(Śs\_5.3)} nirdiśyamānasya ādeśāḥ bhavanti iti visarjanīyasya eva bhaviṣyati . \\
\noindent \textbf{(Śs\_5.3)} dvirvacanam prayojanam . \\
\noindent \textbf{(Śs\_5.3)} uraḥkaḥ , uraḥpaḥ . \\
\noindent \textbf{(Śs\_5.3)} anaci ca . \\
\noindent \textbf{(Śs\_5.3)} acaḥ uttarasya yaraḥ dve bhavataḥ iti dvirvacanam siddham bhavati . \\
\noindent \textbf{(Śs\_5.3)} sthānivadbhāvapratiṣedhaḥ ca prayojanam . \\
\noindent \textbf{(Śs\_5.3)} yathā iha bhavati uraḥkeṇa , uraḥpeṇa iti aḍvyavāye api iti ṇatvam evam iha api sthānivadbhāvāt prāpnoti vyūḍhoraskena mahoraskena iti . \\
\noindent \textbf{(Śs\_5.3)} tatra analvidhau iti pratiṣedhaḥ siddhaḥ bhavati . \\
\noindent \textbf{(Śs\_5.4)} kim punaḥ ime varṇāḥ arthavantaḥ āhosvit anarthakāḥ . \\
\noindent \textbf{(Śs\_5.4)} arthavantaḥ varṇāḥ dhātuprātipadikapratyayanipātānām ekavarṇānām arthadarśanāt . \\
\noindent \textbf{(Śs\_5.4)} dhātavaḥ ekavarṇāḥ arthavantaḥ dṛśyante : eti , adhyeti , adhīte iti . \\
\noindent \textbf{(Śs\_5.4)} prātipadikani ekavarṇāni arthavanti : ābhyām , ebhiḥ , eṣu . \\
\noindent \textbf{(Śs\_5.4)} pratyayāḥ ekavarṇāḥ arthavantaḥ : aupagavaḥ , kāpaṭavaḥ . \\
\noindent \textbf{(Śs\_5.4)} nipātāḥ ekavarṇāḥ arthavantaḥ : a*apehi . \\
\noindent \textbf{(Śs\_5.4)} i*indram paśya . \\
\noindent \textbf{(Śs\_5.4)} u*uttiṣṭha . \\
\noindent \textbf{(Śs\_5.4)} dhātuprātipadikapratyayanipātānām ekavarṇānām arthadarśanāt manyāmahe arthavantaḥ varṇāḥ iti . \\
\noindent \textbf{(Śs\_5.4)} varṇavyatyaye ca arthāntaragamanāt . \\
\noindent \textbf{(Śs\_5.4)} varṇavyatyaye ca arthāntaragamanāt manyāmahe arthavantaḥ varṇāḥ iti . \\
\noindent \textbf{(Śs\_5.4)} kūpaḥ , sūpaḥ , yūpaḥ iti . \\
\noindent \textbf{(Śs\_5.4)} kūpaḥ iti sakakāreṇa kaḥ cit arthaḥ gamyate . \\
\noindent \textbf{(Śs\_5.4)} sūpaḥ iti kakārāpāye sakāropajane ca arthāntaram gamyate . \\
\noindent \textbf{(Śs\_5.4)} yūpaḥ iti kakārasakārāpāye yakāropajane ca arthāntatam gamyate . \\
\noindent \textbf{(Śs\_5.4)} te manyāmahe : yaḥ kūpe kūpārthaḥ saḥ kakārasya . \\
\noindent \textbf{(Śs\_5.4)} yaḥ sūpe sūpārthaḥ saḥ sakārasya . \\
\noindent \textbf{(Śs\_5.4)} yaḥ yūpe yūpārthaḥ saḥ yakārasya iti . \\
\noindent \textbf{(Śs\_5.4)} varṇānupalabdhau ca anarthagateḥ . \\
\noindent \textbf{(Śs\_5.4)} varṇānupalabdhau ca anarthagateḥ manyāmahe arthavantaḥ varṇāḥ iti . \\
\noindent \textbf{(Śs\_5.4)} vṛkṣaḥ , ṛkṣaḥ , kāṇḍīraḥ , āṇḍīraḥ . \\
\noindent \textbf{(Śs\_5.4)} vṛkṣaḥ iti sakakāreṇa kaḥ cit arthaḥ gamyate . \\
\noindent \textbf{(Śs\_5.4)} ṛkṣaḥ iti vakārāpāye saḥ arthaḥ na gamyate . \\
\noindent \textbf{(Śs\_5.4)} kāṇḍīraḥ iti sakakāreṇa kaḥ cit arthaḥ gamyate . \\
\noindent \textbf{(Śs\_5.4)} āṇḍīraḥ iti kakārāpāye saḥ arthaḥ na gamyate . \\
\noindent \textbf{(Śs\_5.4)} kim tarhi ucyate anarthagateḥ iti . \\
\noindent \textbf{(Śs\_5.4)} na sādhīyaḥ hi atra arthasya gatiḥ bhavati . \\
\noindent \textbf{(Śs\_5.4)} evam tarhi idam paṭhitavyam syāt . \\
\noindent \textbf{(Śs\_5.4)} varṇānupalabdhau ca atadarthagateḥ . \\
\noindent \textbf{(Śs\_5.4)} kim idam atadarthagateḥ iti . \\
\noindent \textbf{(Śs\_5.4)} tasya arthaḥ tadarthaḥ . \\
\noindent \textbf{(Śs\_5.4)} tadarthasya gatiḥ tadarthagatiḥ . \\
\noindent \textbf{(Śs\_5.4)} na tadarthagatiḥ atadarthagatiḥ . \\
\noindent \textbf{(Śs\_5.4)} atadarthagateḥ iti . \\
\noindent \textbf{(Śs\_5.4)} atha vā saḥ arthaḥ tadarthaḥ . \\
\noindent \textbf{(Śs\_5.4)} tadarthasya gatiḥ tadarthagatiḥ . \\
\noindent \textbf{(Śs\_5.4)} na tadarthagatiḥ atadarthagatiḥ . \\
\noindent \textbf{(Śs\_5.4)} atadarthagateḥ iti . \\
\noindent \textbf{(Śs\_5.4)} saḥ tarhi tathā nirdeśaḥ kartavyaḥ . \\
\noindent \textbf{(Śs\_5.4)} na kartavyaḥ . \\
\noindent \textbf{(Śs\_5.4)} uttarapadalopaḥ atra draṣṭavyaḥ . \\
\noindent \textbf{(Śs\_5.4)} tat yathā: uṣṭramukham iva mukham asya uṣṭramukhaḥ , kharamukham iva mukham asya kharamukhaḥ . \\
\noindent \textbf{(Śs\_5.4)} evam atadarthagateḥ anarthagateḥ . \\
\noindent \textbf{(Śs\_5.4)} saṅghātārthavattvāt ca . \\
\noindent \textbf{(Śs\_5.4)} saṅghātārthavattvāt ca manyāmahe arthavantaḥ varṇāḥ iti . \\
\noindent \textbf{(Śs\_5.4)} yeṣām saṅghātāḥ arthavantaḥ avayavāḥ api teṣām arthavantaḥ . \\
\noindent \textbf{(Śs\_5.4)} yeṣām punaḥ avayavāḥ anarthakāḥ samudāyāḥ api teṣām anarthakāḥ . \\
\noindent \textbf{(Śs\_5.4)} tat yathā: ekaḥ cakṣuṣmān darśane samarthaḥ tatsamudāyaḥ ca śatam api samartham . \\
\noindent \textbf{(Śs\_5.4)} ekaḥ ca tilaḥ tailadāne samarthaḥ tatsamudāyaḥ ca khārī api samarthā . \\
\noindent \textbf{(Śs\_5.4)} yeṣām punaḥ avayavāḥ anarthakāḥ samudāyāḥ api teṣām anarthakāḥ . \\
\noindent \textbf{(Śs\_5.4)} tat yathā: ekaḥ andhaḥ darśane asamarthaḥ tatsamudāyaḥ ca śatam api asamartham . \\
\noindent \textbf{(Śs\_5.4)} ekā ca sikatā tailadāne asamarthā tatsamudāyaḥ ca khārīśatam api asamartham . \\
\noindent \textbf{(Śs\_5.4)} yadi tarhi ime varṇāḥ arthavantaḥ arthavatkṛtāni prāpnuvanti . \\
\noindent \textbf{(Śs\_5.4)} kāni . \\
\noindent \textbf{(Śs\_5.4)} arthavat prātipadikam iti prātipadikasañjñā prātipadikāt iti svādyutpattiḥ subantam padam iti padasañjñā . \\
\noindent \textbf{(Śs\_5.4)} tatra kaḥ doṣaḥ . \\
\noindent \textbf{(Śs\_5.4)} padasya iti nalopādīni prāpnuvanti . \\
\noindent \textbf{(Śs\_5.4)} dhanam , vanam iti . \\
\noindent \textbf{(Śs\_5.4)} saṅghātasya aikārthyāt subabhāvaḥ varṇāt . \\
\noindent \textbf{(Śs\_5.4)} saṅghātasya ekatvam arthaḥ . \\
\noindent \textbf{(Śs\_5.4)} tena varṇāt subutpattiḥ na bhaviṣyati . \\
\noindent \textbf{(Śs\_5.4)} anarthakāḥ tu prativarṇam arthānupalabdheḥ . \\
\noindent \textbf{(Śs\_5.4)} anarthakāḥ tu varṇāḥ . \\
\noindent \textbf{(Śs\_5.4)} kutaḥ . \\
\noindent \textbf{(Śs\_5.4)} prativarṇam arthānupalabdheḥ . \\
\noindent \textbf{(Śs\_5.4)} na hi prativarṇam arthāḥ upalabhyante . \\
\noindent \textbf{(Śs\_5.4)} kim idam prativarṇam iti . \\
\noindent \textbf{(Śs\_5.4)} varṇam varṇam prati prativarṇam . \\
\noindent \textbf{(Śs\_5.4)} varṇavyatyayāpāyopajanavikāreṣu arthadarśanāt . \\
\noindent \textbf{(Śs\_5.4)} varṇavyatyayāpāyopajanavikāreṣu arthadarśanāt manyāmahe anarthakāḥ varṇāḥ iti . \\
\noindent \textbf{(Śs\_5.4)} varṇavyatyaye: kṛteḥ tarkuḥ , kaseḥ sikatāḥ , hiṃseḥ siṃhaḥ . \\
\noindent \textbf{(Śs\_5.4)} varṇavyatyayaḥ na arthavyatyayaḥ . \\
\noindent \textbf{(Śs\_5.4)} apāyaḥ lopaḥ . \\
\noindent \textbf{(Śs\_5.4)} ghnanti , ghantu , aghnan . \\
\noindent \textbf{(Śs\_5.4)} varṇāpāyaḥ nārthāpāyaḥ . \\
\noindent \textbf{(Śs\_5.4)} upajanaḥ āgamaḥ . \\
\noindent \textbf{(Śs\_5.4)} lavitā , lavitum . \\
\noindent \textbf{(Śs\_5.4)} varṇopajanaḥ na arthopajanaḥ . \\
\noindent \textbf{(Śs\_5.4)} vikāraḥ ādeśaḥ . \\
\noindent \textbf{(Śs\_5.4)} ghātayati , ghātakaḥ . \\
\noindent \textbf{(Śs\_5.4)} varṇavikāraḥ na arthavikāraḥ . \\
\noindent \textbf{(Śs\_5.4)} yathā eva varṇavyatyayāpāyopajanavikārāḥ bhavanti tadvat arthavyatyayāpāyopajanavikāraiḥ bhavitavyam . \\
\noindent \textbf{(Śs\_5.4)} na ca iha tadvat . \\
\noindent \textbf{(Śs\_5.4)} ataḥ manyāmahe anarthakāḥ varṇāḥ iti . \\
\noindent \textbf{(Śs\_5.4)} ubhayam idam varṇeṣu uktam arthavantaḥ anarthakāḥ iti ca . \\
\noindent \textbf{(Śs\_5.4)} kim atra nyāyyam . \\
\noindent \textbf{(Śs\_5.4)} ubhayam iti āha . \\
\noindent \textbf{(Śs\_5.4)} kutaḥ . \\
\noindent \textbf{(Śs\_5.4)} svabhāvataḥ . \\
\noindent \textbf{(Śs\_5.4)} tat yathā: samānam īhamānānām adhīyānānām ca ke cit arthaiḥ yujyante apare na . \\
\noindent \textbf{(Śs\_5.4)} na ca idānīm kaḥ cit arthavān iti kṛtvā sarvaiḥ arthavadbhiḥ śakyam bhavitum , kaḥ cit anarthakaḥ iti kṛtvā sarvaiḥ anarthakaiḥ . \\
\noindent \textbf{(Śs\_5.4)} tatra kim asmābhiḥ śakyam kartum . \\
\noindent \textbf{(Śs\_5.4)} yat dhātupratyayaprātipadikanipātāḥ ekavarṇāḥ arthavantaḥ ataḥ anye anarthakāḥ iti svābhāvikam etat . \\
\noindent \textbf{(Śs\_5.4)} katham yaḥ eṣaḥ bhavatā varṇānām arthavattāyām hetuḥ upadiṣṭaḥ arthavantaḥ varṇāḥ dhātuprātipadikapratyayanipātānām ekavarṇānām arthadarśanāt varṇavyatyaye ca arthāntaragamanāt varṇānupalabdhau ca anarthagateḥ saṅghātārthavattvāt ca iti . \\
\noindent \textbf{(Śs\_5.4)} saṅghātāntarāṇi eva etāni evañjātīyakāni arthāntareṣu vartante: kūpaḥ , sūpaḥ , yūpaḥ iti . \\
\noindent \textbf{(Śs\_5.4)} yadi hi varṇavyatyayakṛtam arthāntaragamanam syāt bhūyiṣṭhaḥ kūpārthaḥ sūpe syāt sūpārthaḥ ca kūpe kūpārthaḥ ca yūpe yūpārthaḥ ca kūpe sūpārthaḥ ca yūpe yūpārthaḥ ca sūpe . \\
\noindent \textbf{(Śs\_5.4)} yataḥ tu khalu na kaḥ cit kūpasya vā sūpe sūpasya vā kūpe kūpasya vā yūpe yūpasya vā kūpe sūpasya vā yūpe yūpasya vā sūpe ataḥ manyāmahe saṅghātāntarāṇi eva etāni evañjātīyakāni arthāntareṣu vartante iti . \\
\noindent \textbf{(Śs\_5.4)} idam khalu api bhavatā varṇānām arthavattām bruvatā sādhīyaḥ anarthakatvam dyotitam . \\
\noindent \textbf{(Śs\_5.4)} yaḥ manyate yaḥ kūpe kūpārthaḥ saḥ kakārasya sūpe sūpārthaḥ saḥ sakārasya yūpe yūpārthaḥ saḥ yakārasya iti ūpaśabdaḥ tasya anarthakaḥ syāt . \\
\noindent \textbf{(Śs\_5.4)} tatra idam aparihṛtam saṅghātārthavattvāt ca . \\
\noindent \textbf{(Śs\_5.4)} etasya api prātipadikasñjñāyām vakṣyati . \\
\noindent \textbf{(Śs\_5.5)} a , i , uṇ ṛ , ḷk e , oṅ ai , auc . \\
\noindent \textbf{(Śs\_5.5)} pratyāhāre anubandhānām katham ajgrahaṇeṣu na . \\
\noindent \textbf{(Śs\_5.5)} ye ete akṣu pratyāhārārthāḥ anubandhāḥ kriyante eteṣām ajgrahaṇeṣu grahaṇam kasmāt na bhavati . \\
\noindent \textbf{(Śs\_5.5)} kim ca syāt . \\
\noindent \textbf{(Śs\_5.5)} dadhi ṇakārīyati madhu ṇakārīyati iti . \\
\noindent \textbf{(Śs\_5.5)} ikaḥ yaṇ aci iti yaṇādeśaḥ prasajyeta . \\
\noindent \textbf{(Śs\_5.5)} ācārāt . \\
\noindent \textbf{(Śs\_5.5)} kim idam ācārāt iti . \\
\noindent \textbf{(Śs\_5.5)} ācāryāṇām upacārāt . \\
\noindent \textbf{(Śs\_5.5)} na eteṣu ācāryāḥ ackāryāṇi kṛtavantaḥ . \\
\noindent \textbf{(Śs\_5.5)} apradhānatvāt . \\
\noindent \textbf{(Śs\_5.5)} apradhānatvāt ca . \\
\noindent \textbf{(Śs\_5.5)} na khalu api eteṣām akṣu prādhānyena upadeśaḥ kriyate . \\
\noindent \textbf{(Śs\_5.5)} kva tarhi . \\
\noindent \textbf{(Śs\_5.5)} halṣu . \\
\noindent \textbf{(Śs\_5.5)} kutaḥ etat . \\
\noindent \textbf{(Śs\_5.5)} eṣā hi ācāryasya śailī lakṣyate yat tulyajātīyān tulyajātīyeṣu upadiśati . \\
\noindent \textbf{(Śs\_5.5)} acaḥ akṣu halaḥ halṣu . \\
\noindent \textbf{(Śs\_5.5)} lopaḥ ca balavattaraḥ . \\
\noindent \textbf{(Śs\_5.5)} lopaḥ khalu api tāvat bhavati . \\
\noindent \textbf{(Śs\_5.5)} ūkālaḥ ac iti vā yogaḥ tatkālānām yathā bhavet acām grahaṇam ackāryam tena eteṣām na bhaviṣyati . \\
\noindent \textbf{(Śs\_5.5)} atha vā yogavibhāgaḥ kariṣyate . \\
\noindent \textbf{(Śs\_5.5)} ūkālaḥ ac . \\
\noindent \textbf{(Śs\_5.5)} u ū u3 iti evaṅkālaḥ ac bhavati . \\
\noindent \textbf{(Śs\_5.5)} tataḥ hrasvadīrghaplutaḥ . \\
\noindent \textbf{(Śs\_5.5)} hrasvadīrghaplutasañjñaḥ ca bhavati ūkālaḥ ac . \\
\noindent \textbf{(Śs\_5.5)} evam api kukkuṭaḥ iti atra api prāpnoti . \\
\noindent \textbf{(Śs\_5.5)} tasmāt pūrvoktaḥ eva parihāraḥ . \\
\noindent \textbf{(Śs\_5.5)} eṣaḥ eva arthaḥ. aparaḥ āha : hrasvādīnām vacanāt prāk yāvat tāvat eva yogaḥ astu ackāryāṇi yathā syuḥ tatkāleṣu akṣu kāryāṇi \\
\noindent \textbf{(Śs\_5.6)} atha kimartham antaḥsthānām aṇsu upadeśaḥ kriyate . \\
\noindent \textbf{(Śs\_5.6)} iha sayŚyŚyantā savŚvŚvatsaraḥ yalŚ lŚlokam talŚ lŚlokam iti parasavarṇasya asiddhatvāt anusvārasya eva dvirvacanam . \\
\noindent \textbf{(Śs\_5.6)} tatra parasya parasavarṇe kṛte tasya yaygrahaṇena grahaṇāt pūrvasya api parasavarṇaḥ yathā syāt . \\
\noindent \textbf{(Śs\_5.6)} na etat asti prayojanam . \\
\noindent \textbf{(Śs\_5.6)} vakṣyati etat . \\
\noindent \textbf{(Śs\_5.6)} dvirvacane parasavarṇatvam siddham vaktavyam iti . \\
\noindent \textbf{(Śs\_5.6)} yāvatā siddhatvam ucyate parasavarṇaḥ eva tāvad bhavati . \\
\noindent \textbf{(Śs\_5.6)} parasavarṇe tarhi kṛte tasya yargrahaṇeṇa grahaṇāt dvirvacanam yathā syāt . \\
\noindent \textbf{(Śs\_5.6)} mā bhūt dvirvacanam . \\
\noindent \textbf{(Śs\_5.6)} nanu ca bhedaḥ bhavati. sati dvirvacane triyakāram asati dvirvacane dviyakāram . \\
\noindent \textbf{(Śs\_5.6)} na asti bhedaḥ . \\
\noindent \textbf{(Śs\_5.6)} sati api dvirvacane dviyakāram eva . \\
\noindent \textbf{(Śs\_5.6)} katham . \\
\noindent \textbf{(Śs\_5.6)} halaḥ yamām yami lopaḥ iti evam ekasya lopena bhavitavyam . \\
\noindent \textbf{(Śs\_5.6)} evam api bhedaḥ . \\
\noindent \textbf{(Śs\_5.6)} sati dvirvacane kadā cit dviyakāram kadā cit triyakāram . \\
\noindent \textbf{(Śs\_5.6)} saḥ eṣaḥ katham bhedaḥ na syāt . \\
\noindent \textbf{(Śs\_5.6)} yadi nityaḥ lopaḥ syāt . \\
\noindent \textbf{(Śs\_5.6)} vibhāṣā ca saḥ lopaḥ . \\
\noindent \textbf{(Śs\_5.6)} yathā abhedaḥ tathā astu . \\
\noindent \textbf{(Śs\_5.6)} anuvartate vibhāṣā śaraḥ aci yat vārayati ayam dvitvam . \\
\noindent \textbf{(Śs\_5.6)} yat ayam śaraḥ aci iti dvirvacanapratiṣedham śāsti tat jñāpayati ācāryaḥ anuvartate vibhāṣā iti . \\
\noindent \textbf{(Śs\_5.6)} katham kṛtvā jñāpakam . \\
\noindent \textbf{(Śs\_5.6)} nitye hi tasya lope pratiṣedhārthaḥ na kaḥ cit syāt . \\
\noindent \textbf{(Śs\_5.6)} yadi nityaḥ lopaḥ syāt pratiṣedhavacanam anarthakam syāt . \\
\noindent \textbf{(Śs\_5.6)} astu atra dvirvacanam . \\
\noindent \textbf{(Śs\_5.6)} jharaḥ jhari savarṇe iti lopaḥ bhaviṣyati . \\
\noindent \textbf{(Śs\_5.6)} paśyati tu ācāryaḥ vibhāṣā saḥ lopaḥ iti . \\
\noindent \textbf{(Śs\_5.6)} tataḥ dvirvacanapratiṣedham śāsti . \\
\noindent \textbf{(Śs\_5.6)} na etat asti jñāpakam . \\
\noindent \textbf{(Śs\_5.6)} nitye api tasya lope saḥ pratiṣedhaḥ avaśyam vaktavyaḥ . \\
\noindent \textbf{(Śs\_5.6)} yat etat acaḥ rahābhyām iti dvirvacanam lopāpavādaḥ saḥ vijñāyate . \\
\noindent \textbf{(Śs\_5.6)} katham . \\
\noindent \textbf{(Śs\_5.6)} yaraḥ iti ucyate . \\
\noindent \textbf{(Śs\_5.6)} etāvantaḥ ca yaraḥ yat uta jharaḥ vā yamaḥ vā . \\
\noindent \textbf{(Śs\_5.6)} yadi ca atra nityaḥ lopaḥ syāt dvirvacanam anarthakam syāt . \\
\noindent \textbf{(Śs\_5.6)} kim tarhi tayoḥ yogayoḥ udāharaṇam . \\
\noindent \textbf{(Śs\_5.6)} yat akṛte dvirvacane trivyañjanaḥ saṃyogaḥ . \\
\noindent \textbf{(Śs\_5.6)} pratttam , avatttam , ādityyaḥ . \\
\noindent \textbf{(Śs\_5.6)} iha idānīm karttā , harttā iti dvirvacanasāmarthyāt lopaḥ na bhavati . \\
\noindent \textbf{(Śs\_5.6)} evam iha api lopaḥ na syāt:: karṣati varṣati iti . \\
\noindent \textbf{(Śs\_5.6)} tasmāt nitye api lope avaśyam saḥ pratiṣedhaḥ vaktavyaḥ . \\
\noindent \textbf{(Śs\_5.6)} tat etat atyantam sandigdham vartate ācāryāṇām vibhāṣā anuvartate na vā iti . \\
\noindent \textbf{(Śs\_6)} ayam ṇakāraḥ dviḥ anubadhyate pūrvaḥ ca paraḥ ca . \\
\noindent \textbf{(Śs\_6)} tatra aṇgrahaṇeṣu iṇgrahaṇeṣu ca sandehaḥ bhavati pūrveṇa vā syuḥ pareṇa vā iti . \\
\noindent \textbf{(Śs\_6)} katarasmin tāvat aṇgrahaṇe sandehaḥ . \\
\noindent \textbf{(Śs\_6)} ḍhralope pūrvasya dīrghaḥ aṇaḥ iti . \\
\noindent \textbf{(Śs\_6)} asandigdham pūrveṇa na pareṇa . \\
\noindent \textbf{(Śs\_6)} kutaḥ etat . \\
\noindent \textbf{(Śs\_6)} parābhāvāt . \\
\noindent \textbf{(Śs\_6)} na hi ḍhralope pare aṇaḥ santi . \\
\noindent \textbf{(Śs\_6)} nanu ca ayam asti : ātṛḍham āvṛḍham iti . \\
\noindent \textbf{(Śs\_6)} evam tarhi sāmarthyāt pūrveṇa na pareṇa . \\
\noindent \textbf{(Śs\_6)} yadi hi pareṇa syāt aṇgrahaṇam anarthakam syāt . \\
\noindent \textbf{(Śs\_6)} ḍhralope pūrvasya dīrghaḥ acaḥ iti eva brūyāt . \\
\noindent \textbf{(Śs\_6)} atha vā etat api na brūyāt . \\
\noindent \textbf{(Śs\_6)} acaḥ hi etat bhavati hrasvaḥ dīrghaḥ plutaḥ iti . \\
\noindent \textbf{(Śs\_6)} asmin tarhi aṇgrahaṇe sandehaḥ ke aṇaḥ iti . \\
\noindent \textbf{(Śs\_6)} asandigdham pūrveṇa na pareṇa . \\
\noindent \textbf{(Śs\_6)} kutaḥ etat . \\
\noindent \textbf{(Śs\_6)} parābhāvāt . \\
\noindent \textbf{(Śs\_6)} na hi ke pare aṇaḥ santi . \\
\noindent \textbf{(Śs\_6)} nanu ca ayam asti . \\
\noindent \textbf{(Śs\_6)} gokā naukā iti . \\
\noindent \textbf{(Śs\_6)} evam tarhi sāmarthyāt pūrveṇa na pareṇa . \\
\noindent \textbf{(Śs\_6)} yadi hi pareṇa syāt aṇgrahaṇam anarthakam syāt . \\
\noindent \textbf{(Śs\_6)} ke acaḥ iti eva brūyāt . \\
\noindent \textbf{(Śs\_6)} atha vā etat api na brūyāt . \\
\noindent \textbf{(Śs\_6)} acaḥ hi etat bhavati hrasvaḥ dīrghaḥ plutaḥ iti . \\
\noindent \textbf{(Śs\_6)} asmin tarhi aṇgrahaṇe sandehaḥ aṇaḥ apragṛhyasya anunāsikaḥ iti . \\
\noindent \textbf{(Śs\_6)} asandigdham pūrveṇa na pareṇa . \\
\noindent \textbf{(Śs\_6)} kutaḥ etat . \\
\noindent \textbf{(Śs\_6)} parābhāvāt . \\
\noindent \textbf{(Śs\_6)} na hi padāntāḥ pare aṇaḥ santi . \\
\noindent \textbf{(Śs\_6)} nanu ca ayam asti kartṛ hartṛ iti . \\
\noindent \textbf{(Śs\_6)} evam tarhi sāmarthyāt pūrveṇa na pareṇa . \\
\noindent \textbf{(Śs\_6)} yadi hi pareṇa syāt aṇgrahaṇam anarthakam syāt .acaḥ apragṛhyasya anunāsikaḥ iti eva brūyāt . \\
\noindent \textbf{(Śs\_6)} atha vā etat api na brūyāt . \\
\noindent \textbf{(Śs\_6)} acaḥ eva hi pragṛhyāḥ bhavanti . \\
\noindent \textbf{(Śs\_6)} asmin tarhi aṇgrahaṇe sandehaḥ uḥ aṇ raparaḥ iti . \\
\noindent \textbf{(Śs\_6)} asandigdham pūrveṇa na pareṇa . \\
\noindent \textbf{(Śs\_6)} kutaḥ etat . \\
\noindent \textbf{(Śs\_6)} parābhāvāt . \\
\noindent \textbf{(Śs\_6)} na hi uḥ sthāne pare aṇaḥ santi . \\
\noindent \textbf{(Śs\_6)} nanu ca ayam asti kartrartham hartrartham iti . \\
\noindent \textbf{(Śs\_6)} kim ca syāt . \\
\noindent \textbf{(Śs\_6)} yadi atra raparatvam syāt dvayoḥ rephayoḥ śravaṇam prasajyeta . \\
\noindent \textbf{(Śs\_6)} halaḥ yamām yami lopaḥ iti evam ekasya atra lopaḥ bhavati. vibhāṣā saḥ lopaḥ . \\
\noindent \textbf{(Śs\_6)} vibhāṣā śravaṇam prasajyeta . \\
\noindent \textbf{(Śs\_6)} ayam tarhi nityaḥ lopaḥ raḥ ri iti . \\
\noindent \textbf{(Śs\_6)} padāntasya iti evam saḥ . \\
\noindent \textbf{(Śs\_6)} na śakyaḥ padāntasya vijñātum . \\
\noindent \textbf{(Śs\_6)} iha hi lopaḥ na syāt jargṛdheḥ laṅ ajarghāḥ pāspardheḥ apāspāḥ iti . \\
\noindent \textbf{(Śs\_6)} iha tarhi mātṝṇām pitṝṇām iti raparatvam prasajyeta . \\
\noindent \textbf{(Śs\_6)} ācāryapravṛttiḥ jñāpayati na atra raparatvam bhavati iti yat ayam ṝtaḥ it dhātoḥ iti dhātugrahaṇam karoti . \\
\noindent \textbf{(Śs\_6)} katham kṛtvā jñāpakam . \\
\noindent \textbf{(Śs\_6)} dhātugrahaṇasya etat prayojanam . \\
\noindent \textbf{(Śs\_6)} iha mā bhūt : mātṝṇām , pitṝṇām iti . \\
\noindent \textbf{(Śs\_6)} yadi ca atra raparatvam syāt dhātugrahaṇam anarthakam syāt . \\
\noindent \textbf{(Śs\_6)} raparatve kṛte anantyatvāt ittvam na bhaviṣyati . \\
\noindent \textbf{(Śs\_6)} paśyati tu ācāryaḥ na atra raparatvam bhavati iti tataḥ dhātugrahaṇam karoti . \\
\noindent \textbf{(Śs\_6)} iha api tarhi ittvam na prāpnoti cikīrṣati jihīrṣati iti . \\
\noindent \textbf{(Śs\_6)} mā bhūt evam . \\
\noindent \textbf{(Śs\_6)} upadhāyāḥ ca iti evam bhaviṣyati . \\
\noindent \textbf{(Śs\_6)} iha api tarhi prāpnoti mātṝṇām , pitṝṇām iti . \\
\noindent \textbf{(Śs\_6)} tasmāt tatra dhātugrahaṇam kartavyam . \\
\noindent \textbf{(Śs\_6)} evam tarhi aṇgrahaṇasāmarthyāt pūrveṇa na pareṇa . \\
\noindent \textbf{(Śs\_6)} yadi pareṇa syāt aṇgrahaṇam anarthakam syāt . \\
\noindent \textbf{(Śs\_6)} uḥ ac raparaḥ iti eva brūyāt . \\
\noindent \textbf{(Śs\_6)} asmin tarhi aṇgrahaṇe sandehaḥ : aṇuditsavarṇasya ca apratyayaḥ iti . \\
\noindent \textbf{(Śs\_6)} asandigdham pareṇa na pūrveṇa iti. kutaḥ etat . \\
\noindent \textbf{(Śs\_6)} savarṇe aṇ taparam hi uḥ ṛt . \\
\noindent \textbf{(Śs\_6)} yat ayam uḥ ṛt iti ṛkāram taparam karoti tat jñāpayati ācāryaḥ pareṇa na pūrveṇa . \\
\noindent \textbf{(Śs\_6)} iṇgrahaṇeṣu tarhi sandehaḥ . \\
\noindent \textbf{(Śs\_6)} asandigdham pareṇa na pūrveṇa iti. kutaḥ etat . \\
\noindent \textbf{(Śs\_6)} yvoḥ anyatra pareṇa iṇ syāt . \\
\noindent \textbf{(Śs\_6)} yatra icchati pūrveṇa sammṛdya grahaṇam tatra karoti yvoḥ iti . \\
\noindent \textbf{(Śs\_6)} tat ca guru bhavati . \\
\noindent \textbf{(Śs\_6)} katham kṛtvā jñāpakam . \\
\noindent \textbf{(Śs\_6)} tatra vibhaktinirdeśe sammṛdya grahaṇe ardhacatasraḥ mātrāḥ . \\
\noindent \textbf{(Śs\_6)} pratyāhāragrahaṇe punaḥ tisraḥ mātrāḥ . \\
\noindent \textbf{(Śs\_6)} saḥ ayam evam laghīyasā nyāsena siddhe sati yat garīyāṃsam yatnam ārabhate tat jñāpayati ācāryaḥ pareṇa na pūrveṇa iti . \\
\noindent \textbf{(Śs\_6)} kim punaḥ varṇotsattau iva ṇakāraḥ dviḥ anubadhyate . \\
\noindent \textbf{(Śs\_6)} etat jñāpayati ācāryaḥ bhavati eṣā paribhāṣā vyākhyānataḥ viśeṣapratipattiḥ na hi sandehāt alakṣaṇam iti. aṇuditsavarṇam parihāya pūrveṇa aṇgrahaṇam pareṇa iṇgrahaṇam iti vyākhyāsyāmaḥ . \\
\noindent \textbf{(Śs\_7-8.1)} KA\_I,35.20-36.4 Ro\_I,115-116 {1/17}     kimartham imau mukhanāsikāvacanau varṇau ubhau api anubadhyete na ñakāra eva anubadhyeta . \\
\noindent \textbf{(Śs\_7-8.1)} KA\_I,35.20-36.4 Ro\_I,115-116 {2/17}     katham yāni makāreṇa grahaṇāni halaḥ yamām yami lopaḥ iti . \\
\noindent \textbf{(Śs\_7-8.1)} KA\_I,35.20-36.4 Ro\_I,115-116 {3/17}     santu ñakāreṇa halaḥ yañām yañi lopaḥ iti . \\
\noindent \textbf{(Śs\_7-8.1)} KA\_I,35.20-36.4 Ro\_I,115-116 {4/17}     na evam śakyam. jhakārabhakāraparayoḥ api hi jhakārabhakāryoḥ lopaḥ prasajyeta . \\
\noindent \textbf{(Śs\_7-8.1)} KA\_I,35.20-36.4 Ro\_I,115-116 {5/17}     na jhakārabhakārau jhakārabhakārayoḥ staḥ . \\
\noindent \textbf{(Śs\_7-8.1)} KA\_I,35.20-36.4 Ro\_I,115-116 {6/17}     katham pumaḥ khayyi ampare iti . \\
\noindent \textbf{(Śs\_7-8.1)} KA\_I,35.20-36.4 Ro\_I,115-116 {7/17}     etat api astu ñakāreṇa pumaḥ khayyi añpare iti . \\
\noindent \textbf{(Śs\_7-8.1)} KA\_I,35.20-36.4 Ro\_I,115-116 {8/17}     na evam śakyam . \\
\noindent \textbf{(Śs\_7-8.1)} KA\_I,35.20-36.4 Ro\_I,115-116 {9/17}     jhakārabhakārapare hi khayyi ruḥ prasajyeta . \\
\noindent \textbf{(Śs\_7-8.1)} KA\_I,35.20-36.4 Ro\_I,115-116 {10/17}     na jhakārabhakāraparaḥ khay asti .katham ṅamaḥ hrasvāt aci ṅamuṭ nityam iti . \\
\noindent \textbf{(Śs\_7-8.1)} KA\_I,35.20-36.4 Ro\_I,115-116 {11/17}     etat api astu ñakāreṇa ṅañaḥ hrasvāt aci ṅañuṭ nityam iti . \\
\noindent \textbf{(Śs\_7-8.1)} KA\_I,35.20-36.4 Ro\_I,115-116 {12/17}     na evam śakyam . \\
\noindent \textbf{(Śs\_7-8.1)} KA\_I,35.20-36.4 Ro\_I,115-116 {13/17}     jhakārabhakāraparayoḥ api hi padāntayoḥ jhakārabhakārau āgamau syātām . \\
\noindent \textbf{(Śs\_7-8.1)} KA\_I,35.20-36.4 Ro\_I,115-116 {14/17}     na jhakārabhakārau padāntau staḥ . \\
\noindent \textbf{(Śs\_7-8.1)} KA\_I,35.20-36.4 Ro\_I,115-116 {15/17}     evam api pañca āgamāḥ trayaḥ āgaminau vaiṣamyāt saṅkhyātānudeśaḥ na prāpnoti . \\
\noindent \textbf{(Śs\_7-8.1)} KA\_I,35.20-36.4 Ro\_I,115-116 {16/17}     santu tāvat yeṣām āgamānām āgaminaḥ santi . \\
\noindent \textbf{(Śs\_7-8.1)} KA\_I,35.20-36.4 Ro\_I,115-116 {17/17}     jhakārabhakārau padāntau na staḥ iti kṛtvā āgamau api na bhaviṣyataḥ . \\
\noindent \textbf{(Śs\_7-8.2)} KA\_I,36.5-11 Ro\_I,117 {1/6}     atha kim idam akṣaram iti . \\
\noindent \textbf{(Śs\_7-8.2)} KA\_I,36.5-11 Ro\_I,117 {2/6}     akṣaram na kṣaram vidyāt . na kṣīyate na kṣarati iti vā akṣaram . \\
\noindent \textbf{(Śs\_7-8.2)} KA\_I,36.5-11 Ro\_I,117 {3/6}     aśnoteḥ vā saraḥ akṣaram . \\
\noindent \textbf{(Śs\_7-8.2)} KA\_I,36.5-11 Ro\_I,117 {4/6}     aśnoteḥ vā punaḥ ayam auṇādikaḥ saranpratyayaḥ . \\
\noindent \textbf{(Śs\_7-8.2)} KA\_I,36.5-11 Ro\_I,117 {5/6}     varṇam vā āhuḥ pūrvasūtre . \\
\noindent \textbf{(Śs\_7-8.2)} KA\_I,36.5-11 Ro\_I,117 {6/6}     atha vā pūrvasūtre varṇasya akṣaram iti sañjñā kriyate . \\
\noindent \textbf{(Śs\_7-8.3)} KA\_I,36.12-18 Ro\_I,118-120 {1/7}     kimartham idam upadiśyate . \\
\noindent \textbf{(Śs\_7-8.3)} KA\_I,36.12-18 Ro\_I,118-120 {2/7}     atha kimartham idam upadiśyate . \\
\noindent \textbf{(Śs\_7-8.3)} KA\_I,36.12-18 Ro\_I,118-120 {3/7}     varṇajñānam vāgviṣayaḥ yatra brahma vartate . \\
\noindent \textbf{(Śs\_7-8.3)} KA\_I,36.12-18 Ro\_I,118-120 {4/7}     tadartham iṣṭadbuddhyartham laghvartham ca upadiśyate . \\
\noindent \textbf{(Śs\_7-8.3)} KA\_I,36.12-18 Ro\_I,118-120 {5/7}     saḥ ayam akṣarasamāmnāyaḥ vāksamāmnāyaḥ puṣpitaḥ phalitaḥ candratārakavat pratimaṇḍitaḥ veditavyaḥ brahmarāśiḥ . \\
\noindent \textbf{(Śs\_7-8.3)} KA\_I,36.12-18 Ro\_I,118-120 {6/7}     sarvavedapuṇyaphalāvāptiḥ ca asya jñāne bhavati . \\
\noindent \textbf{(Śs\_7-8.3)} KA\_I,36.12-18 Ro\_I,118-120 {7/7}     mātāpitarau ca asya svarge loke mahīyete . \\
\noindent \textbf{(P\_1,1.1.1)} kutvam kasmāt na bhavati coḥ kuḥ padasya iti . \\
\noindent \textbf{(P\_1,1.1.1)} bhatvāt . \\
\noindent \textbf{(P\_1,1.1.1)} katham bhasañjñā . \\
\noindent \textbf{(P\_1,1.1.1)} ayasmayādīni chandasi iti . \\
\noindent \textbf{(P\_1,1.1.1)} chandasi iti ucyate . \\
\noindent \textbf{(P\_1,1.1.1)} na ca idam chandaḥ . \\
\noindent \textbf{(P\_1,1.1.1)} chandovat sūtrāṇi bhavanti . \\
\noindent \textbf{(P\_1,1.1.1)} yadi bhasañjñā vṛddhiḥ ād aic at eṅ guṇaḥ iti jaśtvam api na prāpnoti . \\
\noindent \textbf{(P\_1,1.1.1)} ubhayasañjñāni api chandāṃsi dṛśyante . \\
\noindent \textbf{(P\_1,1.1.1)} tat yathā . \\
\noindent \textbf{(P\_1,1.1.1)} saḥ suṣṭubhā saḥ ṛkvatā gaṇena . \\
\noindent \textbf{(P\_1,1.1.1)} padatvāt kutvam bhatvāt jaśvtam na bhavati . \\
\noindent \textbf{(P\_1,1.1.1)} evam iha api bhatvāt kutvam na bhaviṣyati \\
\noindent \textbf{(P\_1,1.1.2)} kim punaḥ idam tadbhāvitagrahaṇam : vṛddhiḥ iti evam ye ākāraikāraukārāḥ bhāvyante teṣām grahaṇam āhosvit ādaijmātrasya . \\
\noindent \textbf{(P\_1,1.1.2)} kim ca ataḥ . \\
\noindent \textbf{(P\_1,1.1.2)} yadi tadbhāvitagrahaṇam śālīyaḥ mālīyaḥ iti vṛddhalakṣaṇaḥ chaḥ na prāpnoti . \\
\noindent \textbf{(P\_1,1.1.2)} āmramayam śālamayam vṛddhalakṣaṇaḥ mayaṭ na prāpnoti . \\
\noindent \textbf{(P\_1,1.1.2)} āmraguptāyaniḥ śālaguptayaniḥ vṛddhalakṣaṇaḥ phiñ na prāpnoti . \\
\noindent \textbf{(P\_1,1.1.2)} atha aijmātrasya grahaṇam sarvaḥ bhāsaḥ sarvabhāsaḥ iti uttarapadapadavṛddhau sarvam ca iti eṣaḥ vidhiḥ prāpnoti . \\
\noindent \textbf{(P\_1,1.1.2)} iha ca tāvatī bhāryā asya tāvadbhāryaḥ yāvadbhāryaḥ vṛddhinimittasya iti puṃvadbhāvapratiṣedhaḥ prāpnoti . \\
\noindent \textbf{(P\_1,1.1.2)} astu tarhi aijmātrasya . \\
\noindent \textbf{(P\_1,1.1.2)} nanu ca uktam sarvaḥ bhāsaḥ sarvabhāsaḥ iti uttarapadapadavṛddhau sarvam ca iti eṣaḥ vidhiḥ prāpnoti . \\
\noindent \textbf{(P\_1,1.1.2)} na eṣaḥ doṣaḥ . \\
\noindent \textbf{(P\_1,1.1.2)} na evam vijñāyate uttarapadasya vṛddhiḥ uttarapadavṛddhiḥ uttarapadavṛddhau iti . \\
\noindent \textbf{(P\_1,1.1.2)} katham tarhi . \\
\noindent \textbf{(P\_1,1.1.2)} uttarapadasya iti evam prakṛtya yā vṛddhiḥ tadvati uttarapade iti evam etat vijñāyate . \\
\noindent \textbf{(P\_1,1.1.2)} avaśyam ca etat evam vijñeyam . \\
\noindent \textbf{(P\_1,1.1.2)} tadbhāvitagrahaṇe sati api iha prasajyeta : sarvaḥ kārakaḥ sarvakārakaḥ iti . \\
\noindent \textbf{(P\_1,1.1.2)} yad api ucyate iha tāvatī bhāryā asya tāvadbhāryaḥ yāvadbhāryaḥ vṛddhinimittasya iti puṃvadbhāvapratiṣedhaḥ prāpnoti iti na eṣaḥ doṣaḥ . \\
\noindent \textbf{(P\_1,1.1.2)} na evam vijñāyate vṛddheḥ nimittam vṛddhinimittam vṛddhinimittasya iti . \\
\noindent \textbf{(P\_1,1.1.2)} katham tarhi . \\
\noindent \textbf{(P\_1,1.1.2)} vṛddheḥ nimittam yasmin saḥ ayam vṛddhinimittaḥ vṛddhinimittasya iti . \\
\noindent \textbf{(P\_1,1.1.2)} kim ca vṛddheḥ nimittam. yaḥ asau kakāraḥ ṇakāraḥ ñakāraḥ vā . \\
\noindent \textbf{(P\_1,1.1.2)} atha vā yaḥ kṛtsnāyāḥ vṛddheḥ nimittam . \\
\noindent \textbf{(P\_1,1.1.2)} kaḥ ca kṛtsnāyāḥ vṛddheḥ nimittam . \\
\noindent \textbf{(P\_1,1.1.2)} yaḥ trayāṇām ākāraikāraukārāṇām . \\
\noindent \textbf{(P\_1,1.1.3)} sañjñādhikāraḥ sañjñāsampratyayārthaḥ . \\
\noindent \textbf{(P\_1,1.1.3)} atha sañjñā iti prakṛtya vṛddhyādayaḥ śabdāḥ paṭhitavyāḥ . \\
\noindent \textbf{(P\_1,1.1.3)} kim prayojanam . \\
\noindent \textbf{(P\_1,1.1.3)} sañjñāsampratyayārthaḥ . \\
\noindent \textbf{(P\_1,1.1.3)} vṛddhyādīnām śabdānām sañjñā iti eṣaḥ sampratyayaḥ yathā syāt . \\
\noindent \textbf{(P\_1,1.1.3)} itarathā hi asampratyayaḥ yathā loke . \\
\noindent \textbf{(P\_1,1.1.3)} akriyamāṇe hi sañjñādhikāre vṛddhyādīnām sañjñā iti eṣaḥ sampratyayaḥ na syāt . \\
\noindent \textbf{(P\_1,1.1.3)} idam idānīm bahusūtram anarthakam syāt . \\
\noindent \textbf{(P\_1,1.1.3)} anarthakam iti āha . \\
\noindent \textbf{(P\_1,1.1.3)} katham . \\
\noindent \textbf{(P\_1,1.1.3)} yathā loke . \\
\noindent \textbf{(P\_1,1.1.3)} loke hi arthavanti ca anarthakāni ca vākyāni dṛśyante . \\
\noindent \textbf{(P\_1,1.1.3)} arthavanti tāvat : devadatta gām abhyāja śuklām daṇḍena . \\
\noindent \textbf{(P\_1,1.1.3)} devadatta gām abhyāja kṛṣṇām iti . \\
\noindent \textbf{(P\_1,1.1.3)} anarthakāni ca . \\
\noindent \textbf{(P\_1,1.1.3)} daśa dāḍimāni ṣaṭ apūpāḥ kuṇḍam ajājinam palalapiṇḍaḥ adhorukam etat kumāryāḥ sphaiyakṛtasya pitā pratiśīnaḥ iti . \\
\noindent \textbf{(P\_1,1.1.3)} sañjñāsañjñyasandehaḥ ca . \\
\noindent \textbf{(P\_1,1.1.3)} kriyamāṇe api sañjñādhikāre sañjñāsañjñinoḥ asandehaḥ vaktavyaḥ . \\
\noindent \textbf{(P\_1,1.1.3)} kutaḥ hi etat vṛddhiśabdaḥ sañjñā ādaicaḥ sañjñinaḥ iti . \\
\noindent \textbf{(P\_1,1.1.3)} na punaḥ ādaicaḥ sañjñā vṛddhiśabdaḥ sañjñī iti . \\
\noindent \textbf{(P\_1,1.1.3)} yat tāvat ucyate sañjñādhikāraḥ kartavyaḥ sañjñāsampratyayārthaḥ iti na kartavyaḥ . \\
\noindent \textbf{(P\_1,1.1.3)} ācāryācārāt sañjñāsiddhiḥ . \\
\noindent \textbf{(P\_1,1.1.3)} ācāryācārāt sañjñāsiddhiḥ bhaviṣyati . \\
\noindent \textbf{(P\_1,1.1.3)} kim idam ācāryācārāt . \\
\noindent \textbf{(P\_1,1.1.3)} ācāryāṇām upacārāt . \\
\noindent \textbf{(P\_1,1.1.3)} yathā laukikavaidikeṣu . \\
\noindent \textbf{(P\_1,1.1.3)} tat yathā laukikeṣu vaidikeṣu ca kṛtānteṣu . \\
\noindent \textbf{(P\_1,1.1.3)} loke tāvat : mātāpitarau putrasya jātasya saṃvṛte avakāśe nāma kurvāte devadattaḥ yajñadattaḥ iti . \\
\noindent \textbf{(P\_1,1.1.3)} tayoḥ upacārāt anye api jānanti iyam asya sañjñā iti . \\
\noindent \textbf{(P\_1,1.1.3)} vede : yājñikāḥ sañjñām kurvanti sphyaḥ yūpaḥ caṣālaḥ iti . \\
\noindent \textbf{(P\_1,1.1.3)} tatrabhavatām upacārāt anye api jānanti iyam asya sañjñā iti . \\
\noindent \textbf{(P\_1,1.1.3)} apare punaḥ sici vṛddhiḥ iti uktvā ākāraikāraukārān udāharanti . \\
\noindent \textbf{(P\_1,1.1.3)} te manyāmahe : yayā pratyāyyante sā sañjñā ye pratīyante te sañjñinaḥ iti . \\
\noindent \textbf{(P\_1,1.1.3)} yat api ucyate kriyamāṇe api sañjñādhikāre sañjñāsañjñinoḥ asandehaḥ vaktavyaḥ iti . \\
\noindent \textbf{(P\_1,1.1.3)} sañjñāsañjñyasandehaḥ . \\
\noindent \textbf{(P\_1,1.1.3)} sañjñāsañjñinoḥ asandehaḥ siddhaḥ . \\
\noindent \textbf{(P\_1,1.1.3)} kutaḥ . \\
\noindent \textbf{(P\_1,1.1.3)} ācāryācārāt eva . \\
\noindent \textbf{(P\_1,1.1.3)} uktaḥ ācāryācāraḥ . \\
\noindent \textbf{(P\_1,1.1.3)} anākṛtiḥ . \\
\noindent \textbf{(P\_1,1.1.3)} atha vā anākṛtiḥ sañjñā . \\
\noindent \textbf{(P\_1,1.1.3)} ākṛtimantaḥ sañjñinaḥ . \\
\noindent \textbf{(P\_1,1.1.3)} loke api hi ākrtimataḥ māṃsapiṇḍasya devadattaḥ iti sañjñā kriyate . \\
\noindent \textbf{(P\_1,1.1.3)} liṅgena vā . \\
\noindent \textbf{(P\_1,1.1.3)} atha vā kim cit liṅgam āsajya vakṣyāmi itthaṃliṅgā sañjñā iti . \\
\noindent \textbf{(P\_1,1.1.3)} vṛddhiśabde ca tat liṅgam kariṣyate na ādaicchabde . \\
\noindent \textbf{(P\_1,1.1.3)} idam tāvat ayuktam yat ucyate ācāryācārāt iti . \\
\noindent \textbf{(P\_1,1.1.3)} kim atra ayuktam . \\
\noindent \textbf{(P\_1,1.1.3)} tam eva upālabhya agamakam te sūtram iti tasya eva punaḥ pramāṇīkaraṇam iti etat ayuktam . \\
\noindent \textbf{(P\_1,1.1.3)} aparituṣyan khalu api bhavān anena parihāreṇa ākṛtiḥ liṅgena vā iti āha . \\
\noindent \textbf{(P\_1,1.1.3)} tat ca api vaktavyam . \\
\noindent \textbf{(P\_1,1.1.3)} yadi api etat ucyate atha vā etarhi itsañjñā na vaktavyā lopaḥ ca na vaktavyaḥ . \\
\noindent \textbf{(P\_1,1.1.3)} sañjñāliṅgam anubandheṣu kariṣyate . \\
\noindent \textbf{(P\_1,1.1.3)} na ca sañjñāyāḥ nivṛttiḥ ucyate . \\
\noindent \textbf{(P\_1,1.1.3)} svabhāvataḥ sañjñāḥ sañjñinaḥ pratyāyya nivartante . \\
\noindent \textbf{(P\_1,1.1.3)} tena anubandhānām api nivṛttiḥ bhaviṣyati . \\
\noindent \textbf{(P\_1,1.1.3)} sidhyati evam . \\
\noindent \textbf{(P\_1,1.1.3)} apāṇinīyam tu bhavati . \\
\noindent \textbf{(P\_1,1.1.3)} yathānyāsam eva astu . \\
\noindent \textbf{(P\_1,1.1.3)} nanu ca uktam sañjñādhikāraḥ sañjñāsampratyayārthaḥ itarathā hi asampratyayaḥ yathā loke iti . \\
\noindent \textbf{(P\_1,1.1.3)} na yathā loke tathā vyākaraṇe . \\
\noindent \textbf{(P\_1,1.1.3)} pramāṇabhūtaḥ ācāryaḥ darbhapavitrapāṇiḥ śucau avakāśe prāṅmukhaḥ upaviśya mahatā yatnena sūtram praṇayati sma . \\
\noindent \textbf{(P\_1,1.1.3)} tatra aśakyam varṇena api anarthakena bhavitum kim punaḥ iyatā sūtreṇa . \\
\noindent \textbf{(P\_1,1.1.3)} kim ataḥ yat aśakyam . \\
\noindent \textbf{(P\_1,1.1.3)} ataḥ sañjñāsañjñinau eva . \\
\noindent \textbf{(P\_1,1.1.3)} kutaḥ nu khalu etat sañjñāsañjñinau eva iti na punaḥ sādhvanuśāsane asmin śāstre sādhutvam anena kiryate . \\
\noindent \textbf{(P\_1,1.1.3)} kṛtam anayoḥ sādhutvam . \\
\noindent \textbf{(P\_1,1.1.3)} katham . \\
\noindent \textbf{(P\_1,1.1.3)} vṛdhiḥ asmai aviśeṣeṇa upadiṣṭaḥ prakṛtipāṭhe . \\
\noindent \textbf{(P\_1,1.1.3)} tasmāt ktinpratyayaḥ . \\
\noindent \textbf{(P\_1,1.1.3)} ādaicaḥ api akṣarasamāmnāye upadiṣṭāḥ . \\
\noindent \textbf{(P\_1,1.1.3)} prayoganiyamārtham tarhi idam syāt . \\
\noindent \textbf{(P\_1,1.1.3)} vṛddhiśabdāt paraḥ ādaicaḥ prayoktavyāḥ iti . \\
\noindent \textbf{(P\_1,1.1.3)} na iha prayoganiyamaḥ ārabhyate . \\
\noindent \textbf{(P\_1,1.1.3)} kim tarhi . \\
\noindent \textbf{(P\_1,1.1.3)} saṃskṛtya saṃskṛtya padāni utsṛjyante . \\
\noindent \textbf{(P\_1,1.1.3)} teṣām yatheṣṭham abhisambandhaḥ bhavati . \\
\noindent \textbf{(P\_1,1.1.3)} tat yathā : āhara pātram , pātram āhara iti . \\
\noindent \textbf{(P\_1,1.1.3)} ādeśāḥ tarhi ime syuḥ . \\
\noindent \textbf{(P\_1,1.1.3)} vṛddhiśabdasya ādaicaḥ . \\
\noindent \textbf{(P\_1,1.1.3)} ṣaṣṭhīnirdiṣṭasya ādeśāḥ ucyante na ca atra ṣaṣṭhīm paśyāmaḥ . \\
\noindent \textbf{(P\_1,1.1.3)} āgamāḥ tarhi ime syuḥ . \\
\noindent \textbf{(P\_1,1.1.3)} vṛddhiśabdasya ādaicaḥ āgamāḥ . \\
\noindent \textbf{(P\_1,1.1.3)} āgamāḥ api ṣaṣṭhīnirdiṣṭasya eva ucyante liṅgena ca . \\
\noindent \textbf{(P\_1,1.1.3)} na ca atra ṣaṣṭhīm na khalu api āgamaliṅgam paśyāmaḥ . \\
\noindent \textbf{(P\_1,1.1.3)} idam khalu api bhūyaḥ sāmanādhikaraṇyam ekavibhaktikatvam ca . \\
\noindent \textbf{(P\_1,1.1.3)} dvayoḥ ca etat bhavati . \\
\noindent \textbf{(P\_1,1.1.3)} kayoḥ . \\
\noindent \textbf{(P\_1,1.1.3)} viśeṣaṇaviśeṣyayoḥ vā sañjñāsañjñinoḥ vā . \\
\noindent \textbf{(P\_1,1.1.3)} tatra etat syāt viśeṣaṇaviśeṣye iti . \\
\noindent \textbf{(P\_1,1.1.3)} tat ca na . \\
\noindent \textbf{(P\_1,1.1.3)} dvayoḥ hi pratītpadārthakayoḥ loke viśeṣaṇaviśeṣyabhāvaḥ bhavati . \\
\noindent \textbf{(P\_1,1.1.3)} na ca ādaicchabdaḥ pratītapadārthakaḥ . \\
\noindent \textbf{(P\_1,1.1.3)} tasmāt sañjñāsañjñinau eva . \\
\noindent \textbf{(P\_1,1.1.3)} tatra tu etāvān sandedhaḥ kaḥ sañjñī kā sañjñā iti . \\
\noindent \textbf{(P\_1,1.1.3)} saḥ ca api kva sandehaḥ . \\
\noindent \textbf{(P\_1,1.1.3)} yatra ubhe samānākṣare . \\
\noindent \textbf{(P\_1,1.1.3)} yatra tu anyatarat laghu yat laghu sā sañjñā . \\
\noindent \textbf{(P\_1,1.1.3)} kutaḥ etat . \\
\noindent \textbf{(P\_1,1.1.3)} laghvartham hi sañjñākaraṇam . \\
\noindent \textbf{(P\_1,1.1.3)} tatra api ayam na avaśyam gurulaghutām eva upalakṣayitum arhati . \\
\noindent \textbf{(P\_1,1.1.3)} kim tarhi . \\
\noindent \textbf{(P\_1,1.1.3)} anākṛtitām api . \\
\noindent \textbf{(P\_1,1.1.3)} anākṛtiḥ sañjñā . \\
\noindent \textbf{(P\_1,1.1.3)} ākṛtimantaḥ sañjñinaḥ . \\
\noindent \textbf{(P\_1,1.1.3)} loke hi ākṛtimataḥ māṃsapiṇḍasya devadattaḥ iti sañjñā kriyate . \\
\noindent \textbf{(P\_1,1.1.3)} atha vā āvartinyaḥ sañjñāḥ bhavanti . \\
\noindent \textbf{(P\_1,1.1.3)} vṛddhiśabdaḥ ca āvartate na ādaicchabdaḥ . \\
\noindent \textbf{(P\_1,1.1.3)} tat yathā . \\
\noindent \textbf{(P\_1,1.1.3)} itaratra api devadattaśabdaḥ āvartate na māṃsapiṇḍaḥ . \\
\noindent \textbf{(P\_1,1.1.3)} atha vā pūrvoccāritaḥ sañjñī paroccāritā sañjñā . \\
\noindent \textbf{(P\_1,1.1.3)} kutaḥ etat . \\
\noindent \textbf{(P\_1,1.1.3)} sataḥ hi kāryiṇaḥ kāryeṇa bhavitavyam . \\
\noindent \textbf{(P\_1,1.1.3)} tat yathā . \\
\noindent \textbf{(P\_1,1.1.3)} itaratra api sataḥ māṃsapiṇḍasya devadattaḥ iti sañjñā kriyate . \\
\noindent \textbf{(P\_1,1.1.3)} katham vṛddhiḥ āt aic iti . \\
\noindent \textbf{(P\_1,1.1.3)} etat ekam ācāryasya maṅgalārtham mṛṣyatām . \\
\noindent \textbf{(P\_1,1.1.3)} māṅgalikaḥ ācāryaḥ mahataḥ śāstraughasya maṅgalārtham vṛddhiśabam āditaḥ prayuṅkte . \\
\noindent \textbf{(P\_1,1.1.3)} maṅgalādīni hi śāstrāṇi prathante vīrapuruṣakāṇi ca bhavanti āyuṣmatpuruṣakāṇi ca . \\
\noindent \textbf{(P\_1,1.1.3)} adhyetāraḥ ca siddhārthāḥ yathā syuḥ iti . \\
\noindent \textbf{(P\_1,1.1.3)} sarvatra eva hi vyākaraṇe pūrvoccāritaḥ sañjñī paroccāritā sañjñā . \\
\noindent \textbf{(P\_1,1.1.3)} at eṅ guṇaḥ iti yathā . \\
\noindent \textbf{(P\_1,1.1.3)} doṣavān khalu api sañjñādhikāraḥ . \\
\noindent \textbf{(P\_1,1.1.3)} aṣṭame api hi sañjñā kriyate tasya param āmreḍitam iti . \\
\noindent \textbf{(P\_1,1.1.3)} tatra api idam anuvartyam syāt . \\
\noindent \textbf{(P\_1,1.1.3)} atha vā asthāne ayam yatnaḥ kriyate . \\
\noindent \textbf{(P\_1,1.1.3)} na hi idam lokāt bhidyate . \\
\noindent \textbf{(P\_1,1.1.3)} yadi idam lokāt bhidyeta tataḥ yatnārham syāt . \\
\noindent \textbf{(P\_1,1.1.3)} tat yathā agojñāya kaḥ cit gām sakthani karṇe vā gṛhītvā upadiśati ayam gauḥ iti . \\
\noindent \textbf{(P\_1,1.1.3)} na ca asmai ācaṣṭe iyam asya sañjñā iti . \\
\noindent \textbf{(P\_1,1.1.3)} bhavati ca asya sampratyayaḥ . \\
\noindent \textbf{(P\_1,1.1.3)} tatra etat syāt kṛtaḥ pūrvaiḥ abhisambandhaḥ iti . \\
\noindent \textbf{(P\_1,1.1.3)} iha api kṛtaḥ pūrvaiḥ abhisambandhaḥ . \\
\noindent \textbf{(P\_1,1.1.3)} kaiḥ . \\
\noindent \textbf{(P\_1,1.1.3)} ācāryaiḥ . \\
\noindent \textbf{(P\_1,1.1.3)} tatra etat syāt . \\
\noindent \textbf{(P\_1,1.1.3)} yasmai samprati upadiśati tasya akṛtaḥ iti . \\
\noindent \textbf{(P\_1,1.1.3)} loke api yasmai samprati upadiśati tasya akṛtaḥ . \\
\noindent \textbf{(P\_1,1.1.3)} atha tatra kṛtaḥ iha api kṛtaḥ draṣṭavyaḥ . \\
\noindent \textbf{(P\_1,1.1.4)} sataḥ vṛddhyādiṣu sañjñābhāvāt tadāśraye itaretarāśrayatvāt asiddhiḥ . \\
\noindent \textbf{(P\_1,1.1.4)} sataḥ sañjñinaḥ sañjñābhāvāt sañjñāśraye sañjñini vṛddhyādiṣu itaretarāśrayatvāt aprasiddhiḥ . \\
\noindent \textbf{(P\_1,1.1.4)} kā iteretarāśrayatā . \\
\noindent \textbf{(P\_1,1.1.4)} satām ādaicām sañjñayā bhavitavyam sañjñayā ca ādaicaḥ bhāvyante . \\
\noindent \textbf{(P\_1,1.1.4)} tat itaretarāśrayam bhavati , itaretarāśrayāṇi ca kāryāṇi na prakalpante . \\
\noindent \textbf{(P\_1,1.1.4)} tat yathā . \\
\noindent \textbf{(P\_1,1.1.4)} nauḥ nāvi baddhā na itaratrāṇāya bhavati. nanu ca bhoḥ itaretarāśrayāṇi api kāryāṇi dṛśyante . \\
\noindent \textbf{(P\_1,1.1.4)} tat yathā . \\
\noindent \textbf{(P\_1,1.1.4)} nauḥ śakaṭam vahati śakaṭam ca nāvam vahati . \\
\noindent \textbf{(P\_1,1.1.4)} anyat api tatra kim cit bhavati jalam sthalam vā . \\
\noindent \textbf{(P\_1,1.1.4)} sthale śakaṭam nāvam vahati . \\
\noindent \textbf{(P\_1,1.1.4)} jale nauḥ śakaṭam vahati . \\
\noindent \textbf{(P\_1,1.1.4)} yathā tari triviṣṭabdhakam . \\
\noindent \textbf{(P\_1,1.1.4)} tatra api antataḥ sūtrakam bhavati . \\
\noindent \textbf{(P\_1,1.1.4)} idam punaḥ itaretarāśrayam eva . \\
\noindent \textbf{(P\_1,1.1.4)} siddham tu nityaśabdatvāt . \\
\noindent \textbf{(P\_1,1.1.4)} siddham etat . \\
\noindent \textbf{(P\_1,1.1.4)} katham . \\
\noindent \textbf{(P\_1,1.1.4)} nityaśabdatvāt . \\
\noindent \textbf{(P\_1,1.1.4)} nityāḥ śabdāḥ . \\
\noindent \textbf{(P\_1,1.1.4)} nityeṣu śabdeṣu satām ādaicām sañjñā kriyate . \\
\noindent \textbf{(P\_1,1.1.4)} na sañjñayā ādaicaḥ bhāvyante . \\
\noindent \textbf{(P\_1,1.1.4)} yadi tarhi nityāḥ śabdāḥ kimartham śāstram . \\
\noindent \textbf{(P\_1,1.1.4)} kimartham śāstram iti cet nivartakatvāt siddham . \\
\noindent \textbf{(P\_1,1.1.4)} nivartakam śāstram . \\
\noindent \textbf{(P\_1,1.1.4)} katham . \\
\noindent \textbf{(P\_1,1.1.4)} mṛjiḥ asmai aviśeṣeṇa upadiṣṭaḥ . \\
\noindent \textbf{(P\_1,1.1.4)} tasya sarvatra mṛjibuddhiḥ prasaktā . \\
\noindent \textbf{(P\_1,1.1.4)} tatra anena nivṛttiḥ kriyate . \\
\noindent \textbf{(P\_1,1.1.4)} mṛjeḥ akṅitsu pratyayeṣu mṛjiprasaṅge mārjiḥ sādhuḥ bhavati iti . \\
\noindent \textbf{(P\_1,1.1.5)} pratyekam vṛddhiguṇasañjñe bhavataḥ iti vaktavyam . \\
\noindent \textbf{(P\_1,1.1.5)} kim prayojanam . \\
\noindent \textbf{(P\_1,1.1.5)} samudāye mā bhūtām iti . \\
\noindent \textbf{(P\_1,1.1.5)} anyatra sahavacanāt samudāye sañjñāprasaṅgaḥ . \\
\noindent \textbf{(P\_1,1.1.5)} anyatra sahavacanāt samudāye vṛddhiguṇasañjñayoḥ aprasaṅgaḥ . \\
\noindent \textbf{(P\_1,1.1.5)} yatra icchati sahabhūtānām kāryam karoti tatra sahagrahaṇam . \\
\noindent \textbf{(P\_1,1.1.5)} tat yathā . \\
\noindent \textbf{(P\_1,1.1.5)} saha supā . \\
\noindent \textbf{(P\_1,1.1.5)} ubhe abhyastam saha iti . \\
\noindent \textbf{(P\_1,1.1.5)} pratyavayavam ca vākyaparisamāpteḥ . \\
\noindent \textbf{(P\_1,1.1.5)} pratyavayavam ca vākyaparisamāptiḥ dṛśyate . \\
\noindent \textbf{(P\_1,1.1.5)} tat yathā . \\
\noindent \textbf{(P\_1,1.1.5)} devadattayajñadattaviṣṇumitrāḥ bhojyantām iti . \\
\noindent \textbf{(P\_1,1.1.5)} na ca ucyate pratyekam iti . \\
\noindent \textbf{(P\_1,1.1.5)} pratyekam ca bhujiḥ parisamāpyate . \\
\noindent \textbf{(P\_1,1.1.5)} nanu ca ayam api asti dṛṣṭāntaḥ samudāye vākyaparisamāptiḥ iti . \\
\noindent \textbf{(P\_1,1.1.5)} tat yathā . \\
\noindent \textbf{(P\_1,1.1.5)} gargāḥ śatam daṇḍyantām iti . \\
\noindent \textbf{(P\_1,1.1.5)} arthinaḥ ca rājānaḥ hiraṇyena bhavanti na ca pratyekam daṇḍayanti . \\
\noindent \textbf{(P\_1,1.1.5)} sati etasmin dṛṣṭānte yadi tatra sahagrahaṇam kriyate iha api pratyekam iti vaktavyam . \\
\noindent \textbf{(P\_1,1.1.5)} atha tatra antareṇa sahagrahaṇam sahabhūtānām kāryam bhavati iha api na arthaḥ pratyekam iti vacanena . \\
\noindent \textbf{(P\_1,1.1.6)} atha kimartham ākāraḥ taparaḥ kriyate . \\
\noindent \textbf{(P\_1,1.1.6)} ākārasya taparakaraṇam savarṇārtham . \\
\noindent \textbf{(P\_1,1.1.6)} ākārasya taparakaraṇam kriyate . \\
\noindent \textbf{(P\_1,1.1.6)} kim prayojanam . \\
\noindent \textbf{(P\_1,1.1.6)} savarṇārtham . \\
\noindent \textbf{(P\_1,1.1.6)} taparaḥ tatkālasya iti tatkālānām grahaṇam yathā syāt . \\
\noindent \textbf{(P\_1,1.1.6)} keṣām . \\
\noindent \textbf{(P\_1,1.1.6)} udāttānundāttasvaritānām . \\
\noindent \textbf{(P\_1,1.1.6)} kim ca kāraṇam na syāt . \\
\noindent \textbf{(P\_1,1.1.6)} bhedakatvāt svarasya . \\
\noindent \textbf{(P\_1,1.1.6)} bhedakāḥ udāttādayaḥ . \\
\noindent \textbf{(P\_1,1.1.6)} katham punaḥ jñāyate bhedakāḥ udāttādayaḥ iti . \\
\noindent \textbf{(P\_1,1.1.6)} evam hi dṛśyate loke . \\
\noindent \textbf{(P\_1,1.1.6)} yaḥ udātte kartavye anudāttam karoti khaṇḍikopādhyāyaḥ tasmai capeṭām dadāti anyat tvam karoṣi iti . \\
\noindent \textbf{(P\_1,1.1.6)} asti prayojanam etat . \\
\noindent \textbf{(P\_1,1.1.6)} kim tarhi iti . \\
\noindent \textbf{(P\_1,1.1.6)} bhedakatvāt guṇasya iti vaktavyam . \\
\noindent \textbf{(P\_1,1.1.6)} kim prayojanam . \\
\noindent \textbf{(P\_1,1.1.6)} ānunāsikyam nāma guṇaḥ . \\
\noindent \textbf{(P\_1,1.1.6)} tadbhinnasya api yathā syāt . \\
\noindent \textbf{(P\_1,1.1.6)} kim ca kāraṇam na syāt . \\
\noindent \textbf{(P\_1,1.1.6)} bhedakatvāt guṇasya . \\
\noindent \textbf{(P\_1,1.1.6)} bhedakāḥ guṇāḥ . \\
\noindent \textbf{(P\_1,1.1.6)} katham punaḥ jñāyate bhedakāḥ guṇāḥ iti . \\
\noindent \textbf{(P\_1,1.1.6)} evam hi dṛśyate loke . \\
\noindent \textbf{(P\_1,1.1.6)} ekaḥ ayam ātmā udakam nāma . \\
\noindent \textbf{(P\_1,1.1.6)} tasya guṇabhedāt anyatvam bhavati : anyat idam śītam anyat idam ūṣṇam iti . \\
\noindent \textbf{(P\_1,1.1.6)} nanu ca bhoḥ abhedakāḥ api guṇāḥ dṛśyante . \\
\noindent \textbf{(P\_1,1.1.6)} tat yathā . \\
\noindent \textbf{(P\_1,1.1.6)} devadattaḥ muṇḍī api jaṭī api śikhī api svām ākhyām na jahāti . \\
\noindent \textbf{(P\_1,1.1.6)} tathā bālaḥ yuvā vṛddhaḥ vatsaḥ damyaḥ balīvardaḥ iti . \\
\noindent \textbf{(P\_1,1.1.6)} ubhayam idam guṇeṣu uktam bhedakāḥ abhedakāḥ iti . \\
\noindent \textbf{(P\_1,1.1.6)} kim punaḥ atra nyāyyam . \\
\noindent \textbf{(P\_1,1.1.6)} abhedakāḥ guṇāḥ iti eva nyāyyam . \\
\noindent \textbf{(P\_1,1.1.6)} kutaḥ etat . \\
\noindent \textbf{(P\_1,1.1.6)} yat ayam asthidadhisakthyakṣṇām anaṅ udāttaḥ iti udāttagrahaṇam karoti . \\
\noindent \textbf{(P\_1,1.1.6)} yadi bhedakāḥ guṇāḥ syuḥ udāttam eva uccārayet . \\
\noindent \textbf{(P\_1,1.1.6)} yadi tarhi abhedakāḥ guṇāḥ anudāttādeḥ antodāttāt ca yat ucyate tat svaritādeḥ svaritāntāt ca prāpnoti . \\
\noindent \textbf{(P\_1,1.1.6)} na eṣaḥ doṣaḥ . \\
\noindent \textbf{(P\_1,1.1.6)} āśrīyamāṇaḥ guṇaḥ bhedakaḥ bhavati . \\
\noindent \textbf{(P\_1,1.1.6)} tat yathā . \\
\noindent \textbf{(P\_1,1.1.6)} śuklam ālabheta . \\
\noindent \textbf{(P\_1,1.1.6)} kṛṣṇam ālabheta . \\
\noindent \textbf{(P\_1,1.1.6)} tatra yaḥ śukle ālabdhavye kṛṣṇam ālabheta na hi tena yathoktam kṛtam bhavati . \\
\noindent \textbf{(P\_1,1.1.6)} asandehārthaḥ tarhi takāraḥ . \\
\noindent \textbf{(P\_1,1.1.6)} aic iti ucyamane sandehaḥ syāt . \\
\noindent \textbf{(P\_1,1.1.6)} kim imau aicau eva āhosvit ākāraḥ api atra nirdiśyate iti . \\
\noindent \textbf{(P\_1,1.1.6)} sandehamātram etat bhavati . \\
\noindent \textbf{(P\_1,1.1.6)} sarvasandeheṣu ca idam upatiṣṭhate vyākhyānataḥ viśeṣapratipattiḥ na hi sandehāt alakṣaṇam iti . \\
\noindent \textbf{(P\_1,1.1.6)} trayāṇām grahaṇam iti vyākhyāsyāmaḥ . \\
\noindent \textbf{(P\_1,1.1.6)} anyatra api hi ayam evañjātīyakeṣu sandeheṣu na kam cid yatnam karoti . \\
\noindent \textbf{(P\_1,1.1.6)} tat yathā . \\
\noindent \textbf{(P\_1,1.1.6)} autaḥ amśasoḥ iti .idam tarhi prayojanam : āntaryataḥ trimātracaturnātrāṇām sthāninām trimātracaturmātrāḥ ādeśāḥ mā bhūvan iti : khaṭvā* indraḥ khaṭvendraḥ , khaṭvā* udakam khaṭvodakam , khaṭvā* īṣā khaṭveṣā khaṭvā* ūḍhā khaṭvoḍhā khaṭvā* elakā khaṭvailakā khaṭvā* odanaḥ , khaṭvaudanaḥ , khaṭvā* aitikāyanaḥ , khaṭvaitikāyanaḥ , khaṭvā* aupagavaḥ , khaṭvaupagavaḥ iti . \\
\noindent \textbf{(P\_1,1.1.6)} atha kriyamāṇe api takāre kasmāt eva trimātracaturnātrāṇām sthāninam trimātracaturmātrāḥ ādeśāḥ na bhavanti . \\
\noindent \textbf{(P\_1,1.1.6)} taparaḥ tatkālasya iti niyamāt . \\
\noindent \textbf{(P\_1,1.1.6)} nanu taḥ paraḥ yasmāt saḥ ayam taparaḥ . \\
\noindent \textbf{(P\_1,1.1.6)} na iti āha . \\
\noindent \textbf{(P\_1,1.1.6)} tāt api paraḥ taparaḥ iti . \\
\noindent \textbf{(P\_1,1.1.6)} yadi tāt api paraḥ taparaḥ ṝdoḥ ap iti iha eva syāt . \\
\noindent \textbf{(P\_1,1.1.6)} yavaḥ stavaḥ . \\
\noindent \textbf{(P\_1,1.1.6)} lavaḥ pavaḥ iti atra na syāt . \\
\noindent \textbf{(P\_1,1.1.6)} na eṣaḥ takāraḥ . \\
\noindent \textbf{(P\_1,1.1.6)} kaḥ tarhi . \\
\noindent \textbf{(P\_1,1.1.6)} dakāraḥ . \\
\noindent \textbf{(P\_1,1.1.6)} kim dakāre prayojanam . \\
\noindent \textbf{(P\_1,1.1.6)} atha kim takāre prayojanam . \\
\noindent \textbf{(P\_1,1.1.6)} yadi asandehārthaḥ takāraḥ dakāraḥ api . \\
\noindent \textbf{(P\_1,1.1.6)} atha mukhasukhārthaḥ takāraḥ dakāraḥ api . \\
\noindent \textbf{(P\_1,1.3.1)} iggrahaṇam kimartham . \\
\noindent \textbf{(P\_1,1.3.1)} iggrahaṇam ātsandhyakṣaravyañjananivṛttyartham . iggrahaṇam kriyate ākāranivṛttyartham sandhyakṣaranivṛttyartham vyañjananivṛttyartham ca. ākāranivṛttyartham tāvat . \\
\noindent \textbf{(P\_1,1.3.1)} yātā vātā . \\
\noindent \textbf{(P\_1,1.3.1)} ākārasya guṇaḥ prāpnoti . \\
\noindent \textbf{(P\_1,1.3.1)} iggrahaṇāt na bhavati . \\
\noindent \textbf{(P\_1,1.3.1)} sandhyakṣaranivṛttyartham . \\
\noindent \textbf{(P\_1,1.3.1)} glāyati mlāyati . \\
\noindent \textbf{(P\_1,1.3.1)} sandhyakṣarasya guṇaḥ prāpnoti . \\
\noindent \textbf{(P\_1,1.3.1)} iggrahaṇāt na bhavati . \\
\noindent \textbf{(P\_1,1.3.1)} vyañjananivṛttyartham . \\
\noindent \textbf{(P\_1,1.3.1)} umbhitā , umbhitum umbhitavyam . \\
\noindent \textbf{(P\_1,1.3.1)} vyañjanasya guṇaḥ prāpnoti . \\
\noindent \textbf{(P\_1,1.3.1)} iggrahaṇāt na bhavati . \\
\noindent \textbf{(P\_1,1.3.1)} ākāranivṛttyarthena tāvat nārthaḥ . \\
\noindent \textbf{(P\_1,1.3.1)} ācāryapravṛttiḥ jñāpayati na ākārasya guṇaḥ bhavati iti yat ayam ātaḥ anupasarge kaḥ iti kakāram anubandham karoti . \\
\noindent \textbf{(P\_1,1.3.1)} katham kṛtvā jñāpakam . \\
\noindent \textbf{(P\_1,1.3.1)} kitkaraṇe etat prayojanam kiti iti ākāralopaḥ yathā syāt . \\
\noindent \textbf{(P\_1,1.3.1)} yadi ca ākārasya guṇaḥ syāt kitkaraṇam anarthakam syāt . \\
\noindent \textbf{(P\_1,1.3.1)} guṇe kṛte dvayoḥ akārayoḥ pararūpeṇa siddham rūpam syāt godaḥ , kambaladaḥ iti . \\
\noindent \textbf{(P\_1,1.3.1)} paśyati tu ācāryaḥ na ākārasya guṇaḥ bhavati iti . \\
\noindent \textbf{(P\_1,1.3.1)} tataḥ kakāram anubandham karoti . \\
\noindent \textbf{(P\_1,1.3.1)} sandhyakṣarārthena api na arthaḥ . \\
\noindent \textbf{(P\_1,1.3.1)} upadeśasāmarthyāt sandhyakṣarasya guṇaḥ na bhaviṣyati . \\
\noindent \textbf{(P\_1,1.3.1)} vyañjananivṛttyarthena api na arthaḥ . \\
\noindent \textbf{(P\_1,1.3.1)} ācāryapravṛttiḥ jñāpayati na vyañjanasya guṇaḥ bhavati iti yat ayam janeḥ ḍam śāsti . \\
\noindent \textbf{(P\_1,1.3.1)} katham kṛtvā jñāpakam . \\
\noindent \textbf{(P\_1,1.3.1)} ḍitkaraṇe etat prayojanam ḍiti iti ṭilopaḥ yathā syāt . \\
\noindent \textbf{(P\_1,1.3.1)} yadi vyañjansya guṇaḥ syāt ḍitkaraṇam anarthakam syāt . \\
\noindent \textbf{(P\_1,1.3.1)} guṇe kṛte trayāṇām akārāṇām pararūpeṇa siddham rūpam syāt : upasarajaḥ , mandurajaḥ iti . \\
\noindent \textbf{(P\_1,1.3.1)} paśyati tu ācāryaḥ na vyañjanasya guṇaḥ bhavati iti . \\
\noindent \textbf{(P\_1,1.3.1)} tataḥ janeḥ ḍam śāsti. na etāni santi prayojanāni . \\
\noindent \textbf{(P\_1,1.3.1)} yat tāvat ucyate kitkaraṇam jñāpakam ākārasya guṇaḥ na bhavati iti . \\
\noindent \textbf{(P\_1,1.3.1)} uttarārtham etat syāt . \\
\noindent \textbf{(P\_1,1.3.1)} tundaśokayoḥ parimṛjāpanudoḥ iti . \\
\noindent \textbf{(P\_1,1.3.1)} yat tarhi gāpoḥ ṭhak iti ananyārtham kakāram anubandham karoti . \\
\noindent \textbf{(P\_1,1.3.1)} yat api ucyate upadeśasāmarthyāt sandhyakṣarasya guṇaḥ na bhavati iti . \\
\noindent \textbf{(P\_1,1.3.1)} yadi yat yat sandhyakṣarasya prāpnoti tat tat upadeśasāmarthyāt bādhyate āyādayaḥ api tarhi na prāpnuvanti . \\
\noindent \textbf{(P\_1,1.3.1)} na eṣaḥ doṣaḥ . \\
\noindent \textbf{(P\_1,1.3.1)} yam vidhim prati upadeśaḥ anarthakaḥ sa vidhiḥ bādhyate . \\
\noindent \textbf{(P\_1,1.3.1)} yasya tu vidheḥ nimittam eva na asau bādhyate . \\
\noindent \textbf{(P\_1,1.3.1)} guṇam ca prati upadeśaḥ anarthakaḥ āyādīnām punaḥ nimittam eva . \\
\noindent \textbf{(P\_1,1.3.1)} yat api ucyate janeḥ ḍavacanam jñāpakam na vyañjanasya guṇaḥ bhavati iti . \\
\noindent \textbf{(P\_1,1.3.1)} siddhe vidhiḥ ārabhyamāṇaḥ jñāpakārthaḥ bhavati . \\
\noindent \textbf{(P\_1,1.3.1)} na ca janeḥ guṇena sidhyati . \\
\noindent \textbf{(P\_1,1.3.1)} kutaḥ hi etat janeḥ guṇaḥ ucyamānaḥ akāraḥ bhavati na punaḥ ekāraḥ vā syāt okāraḥ vā iti . \\
\noindent \textbf{(P\_1,1.3.1)} āntaryataḥ ardhamātrikasya vyañjanasya mātrikaḥ akāraḥ bhaviṣyati . \\
\noindent \textbf{(P\_1,1.3.1)} evam api anunāsikaḥ prāpnoti . \\
\noindent \textbf{(P\_1,1.3.1)} pararūpeṇa śuddhaḥ bhaviṣyati . \\
\noindent \textbf{(P\_1,1.3.1)} evam tarhi gameḥ api ayam ḍaḥ vaktavyaḥ . \\
\noindent \textbf{(P\_1,1.3.1)} gameḥ ca guṇaḥ ucyamānaḥ āntaryataḥ okāraḥ prāpnoti . \\
\noindent \textbf{(P\_1,1.3.1)} tasmāt iggrahaṇam kartavyam . \\
\noindent \textbf{(P\_1,1.3.1)} yadi iggrahaṇam kriyate dyauḥ , panthāḥ , saḥ , imam ite ete api ikaḥ prāpnuvanti . \\
\noindent \textbf{(P\_1,1.3.1)} sañjñayā vidhāne niyamaḥ . \\
\noindent \textbf{(P\_1,1.3.1)} sañjñayā ye vidhīyante teṣu niyamaḥ . \\
\noindent \textbf{(P\_1,1.3.1)} kim vaktavyam etat . \\
\noindent \textbf{(P\_1,1.3.1)} na hi . \\
\noindent \textbf{(P\_1,1.3.1)} katham anucyamānam gaṃsyate . \\
\noindent \textbf{(P\_1,1.3.1)} guṇavṛddhigrahaṇasāmarthyāt . \\
\noindent \textbf{(P\_1,1.3.1)} katham punaḥ antareṇa guṇavṛddhigrahaṇam ikaḥ guṇavṛddhī syātām . \\
\noindent \textbf{(P\_1,1.3.1)} prakṛtam guṇavṛddhigrahaṇam anuvartate . \\
\noindent \textbf{(P\_1,1.3.1)} kva prakṛtam . \\
\noindent \textbf{(P\_1,1.3.1)} vṛddhiḥ āt aic at eṅ guṇaḥ iti . \\
\noindent \textbf{(P\_1,1.3.1)} yadi tat anuvartate at eṅ guṇaḥ vṛddhiḥ ca iti adeṅām api vṛddhisañjñā prāpnoti . \\
\noindent \textbf{(P\_1,1.3.1)} sambandham anuvartiṣyate . \\
\noindent \textbf{(P\_1,1.3.1)} vṛddhiḥ āt aic . \\
\noindent \textbf{(P\_1,1.3.1)} at eṅ guṇaḥ vṛddhiḥ āt aic . \\
\noindent \textbf{(P\_1,1.3.1)} tataḥ ikaḥ guṇavṛddhī iti . \\
\noindent \textbf{(P\_1,1.3.1)} guṇavṛddhigrahaṇam anuvartate ādaijgrahaṇam nivṛttam . \\
\noindent \textbf{(P\_1,1.3.1)} atha vā maṇḍūkagatayaḥ adhikārāḥ . \\
\noindent \textbf{(P\_1,1.3.1)} yathā maṇḍūkāḥ utplutya utplutya gacchanti tadvat adhikārāḥ . \\
\noindent \textbf{(P\_1,1.3.1)} atha vā ekayogaḥ kariṣyate vṛddhiḥ āt aic at eṅ guṇaḥ iti . \\
\noindent \textbf{(P\_1,1.3.1)} tataḥ iko guṇavṛddhī iti . \\
\noindent \textbf{(P\_1,1.3.1)} na ca ekayoge anuvṛttiḥ bhavati . \\
\noindent \textbf{(P\_1,1.3.1)} atha vā anyavacanāt cakārākaraṇāt prakṛtāpavādaḥ vijñāyate yathā utsargeṇa prasaktasya apavādaḥ bādhakaḥ bhavati . \\
\noindent \textbf{(P\_1,1.3.1)} anyasyāḥ sañjñāyāḥ vacanāt cakārasya anukarṣaṇārthasya akaraṇāt prakṛtāyāḥ vṛddhisañjñāyāḥ guṇasañjña bādhikā bhaviṣyati yathā utsargeṇa prasaktasya apavādaḥ bādhakaḥ bhavati . \\
\noindent \textbf{(P\_1,1.3.1)} atha vā vakṣyati etat . \\
\noindent \textbf{(P\_1,1.3.1)} anuvartante ca nāma vidhayaḥ na ca anuvartanāt eva bhavanti . \\
\noindent \textbf{(P\_1,1.3.1)} kim tarhi yatnāt bhavanti . \\
\noindent \textbf{(P\_1,1.3.1)} atha vā ubhayam nivṛttam . \\
\noindent \textbf{(P\_1,1.3.1)} tat apekṣiṣyāmahe . \\
\noindent \textbf{(P\_1,1.3.2)} kim punaḥ ayam alontyaśeṣaḥ āhosvit alontyāpavādaḥ . \\
\noindent \textbf{(P\_1,1.3.2)} katham ca ayam taccheṣaḥ syāt katham vā tadapavādaḥ . \\
\noindent \textbf{(P\_1,1.3.2)} yadi ekam vākyam tat ca idam ca alaḥ : antyasya vidhayaḥ bhavanti ikaḥ guṇavṛddhī* alaḥ antyasya iti tataḥ ayam taccheṣaḥ . \\
\noindent \textbf{(P\_1,1.3.2)} atha nānā vākyam : alaḥ antyasya vidhayaḥ bhavanti , ikaḥ guṇavṛddhī* antyasya ca anantyasya ca iti tataḥ ayam tadapavādaḥ . \\
\noindent \textbf{(P\_1,1.3.2)} kaḥ ca atra viśeṣaḥ . \\
\noindent \textbf{(P\_1,1.3.2)} vṛddhiguṇau alaḥ antyasya iti cet midipugantalaghūpadharcchidṛśikṣiprakṣudreṣu iggrahaṇam . \\
\noindent \textbf{(P\_1,1.3.2)} vṛddhiguṇau alaḥ antyasya iti cet midipugantalaghūpadharcchidṛśikṣiprakṣudreṣu iggrahaṇam kartavyam . \\
\noindent \textbf{(P\_1,1.3.2)} mideḥ guṇaḥ ikaḥ iti vaktavyam . \\
\noindent \textbf{(P\_1,1.3.2)} anantyatvāt hi na prāpnoti . \\
\noindent \textbf{(P\_1,1.3.2)} pugantalaghūpadhasya guṇaḥ ikaḥ iti vaktavyam . \\
\noindent \textbf{(P\_1,1.3.2)} anantyatvāt hi na prāpnoti . \\
\noindent \textbf{(P\_1,1.3.2)} ṛccheḥ liṭi guṇaḥ ikaḥ iti vaktavyam . \\
\noindent \textbf{(P\_1,1.3.2)} anantyatvāt hi na prāpnoti . \\
\noindent \textbf{(P\_1,1.3.2)} ṛdṛsaḥ aṅi guṇaḥ ikaḥ iti vaktavyam . \\
\noindent \textbf{(P\_1,1.3.2)} anantyatvāt hi na prāpnoti . \\
\noindent \textbf{(P\_1,1.3.2)} kṣiprakṣudrayoḥ guṇaḥ ikaḥ iti vaktavyam . \\
\noindent \textbf{(P\_1,1.3.2)} anantyatvāt hi na prāpnoti . \\
\noindent \textbf{(P\_1,1.3.2)} sarvādeśaprasaṅgaḥ ca anigantasya . \\
\noindent \textbf{(P\_1,1.3.2)} sarvādeśaḥ ca guṇaḥ ca anigantasya prāpnoti : yātā vātā . \\
\noindent \textbf{(P\_1,1.3.2)} kim kāraṇam . \\
\noindent \textbf{(P\_1,1.3.2)} alaḥ antyasya iti ṣaṣṭhī ca eva hi antyam ikam upasaṅkrāntā , aṅgasya iti ca sthānaṣaṣṭhī . \\
\noindent \textbf{(P\_1,1.3.2)} tat yat idānīm anigantam aṅgam tasya guṇaḥ sarvādeśaḥ prāpnoti . \\
\noindent \textbf{(P\_1,1.3.2)} na eṣaḥ doṣaḥ . \\
\noindent \textbf{(P\_1,1.3.2)} yathā eva hi alaḥ antyasya iti ṣaṣṭhī antyam ikam upasaṅkrāntā evam aṅgasya iti api sthānaṣaṣṭhī . \\
\noindent \textbf{(P\_1,1.3.2)} tat yad idānīm anigantam aṅgam , tatra ṣaṣṭhī eva na asti kutaḥ guṇaḥ kutaḥ sarvādeśaḥ . \\
\noindent \textbf{(P\_1,1.3.2)} evam tarhi na ayam doṣasamuccayaḥ . \\
\noindent \textbf{(P\_1,1.3.2)} kim tarhi . \\
\noindent \textbf{(P\_1,1.3.2)} pūrvāpekṣaḥ ayam doṣaḥ , hyarthe ca ayam caḥ paṭhitaḥ . \\
\noindent \textbf{(P\_1,1.3.2)} midipugantalaghūpadharcchidṛśikṣiprakṣudreṣu iggrahaṇam sarvādeśaprasaṅgaḥ hi anigantasya iti . \\
\noindent \textbf{(P\_1,1.3.2)} mideḥ guṇaḥ ikaḥ iti vacanāt antyasya na . \\
\noindent \textbf{(P\_1,1.3.2)} antyasya iti vacanāt ikaḥ na . \\
\noindent \textbf{(P\_1,1.3.2)} ucyate tu guṇaḥ . \\
\noindent \textbf{(P\_1,1.3.2)} saḥ sarvādeśaḥ prāpnoti . \\
\noindent \textbf{(P\_1,1.3.2)} evam sarvatra . \\
\noindent \textbf{(P\_1,1.3.2)} astu tarhi tadapavādaḥ . \\
\noindent \textbf{(P\_1,1.3.2)} igmātrasya iti cet jusisārvadhātukārdhadhātukahrasvādyoḥ guṇeṣu anantyapratiṣedhaḥ . \\
\noindent \textbf{(P\_1,1.3.2)} igmātrasya iti cet jusisārvadhātukārdhadhātukahrasvādyoḥ guṇeṣu anantyapratiṣedhaḥ vaktavyaḥ . \\
\noindent \textbf{(P\_1,1.3.2)} jusi guṇaḥ . \\
\noindent \textbf{(P\_1,1.3.2)} saḥ yathā iha bhavati : ajuhavuḥ , abibhayuḥ , evam anenijuḥ , paryaviviṣuḥ , atra api prāpnoti . \\
\noindent \textbf{(P\_1,1.3.2)} sārvadhātukārdhadhātukayoḥ guṇaḥ . \\
\noindent \textbf{(P\_1,1.3.2)} saḥ yathā iha bhavati : kartā hartā nayati tarati bhavati , evam īhitā , īhitum iti atra api prāpnoti . \\
\noindent \textbf{(P\_1,1.3.2)} hrasvasya guṇaḥ . \\
\noindent \textbf{(P\_1,1.3.2)} saḥ yathā iha bhavati : he agne he vāyo , evam he agnicit , he somasut iti atra api prāpnoti . \\
\noindent \textbf{(P\_1,1.3.2)} jasi guṇaḥ . \\
\noindent \textbf{(P\_1,1.3.2)} saḥ yathā iha bhavati agnayaḥ , vāyavaḥ iti evam agnicitaḥ , somasutaḥ iti atra api prāpnoti . \\
\noindent \textbf{(P\_1,1.3.2)} ṛto ṅisarvanāmasthānayoḥ guṇaḥ . \\
\noindent \textbf{(P\_1,1.3.2)} saḥ yathā iha bhavati kartari kartārau kartāraḥ iti evam sukṛti sukṛtau sukṛtaḥ iti atra api prāpnoti .gheḥ ṅiti guṇaḥ . \\
\noindent \textbf{(P\_1,1.3.2)} saḥ yathā iha bhavati agnaye vāyave evam agnicite somasute iti atra api prāpnoti . \\
\noindent \textbf{(P\_1,1.3.2)} oḥ guṇaḥ . \\
\noindent \textbf{(P\_1,1.3.2)} saḥ yathā iha bhavati bābhravyaḥ , māṇḍavyaḥ iti evam suśrut , sauśrutaḥ iti atra api prāpnoti . \\
\noindent \textbf{(P\_1,1.3.2)} na eṣaḥ doṣaḥ . \\
\noindent \textbf{(P\_1,1.3.2)} pugantalaghūpadhagrahaṇam anantyaniyamārtham . \\
\noindent \textbf{(P\_1,1.3.2)} pugantalaghūpadhagrahaṇam anantyaniyamārtham bhaviṣyati . \\
\noindent \textbf{(P\_1,1.3.2)} pugantalaghūpadhasya eva anantyasya na anyasya anantyasya iti . \\
\noindent \textbf{(P\_1,1.3.2)} prakṛtasya eṣaḥ niyamaḥ syāt . \\
\noindent \textbf{(P\_1,1.3.2)} kim ca prakṛtam . \\
\noindent \textbf{(P\_1,1.3.2)} sārvadhātukārdhadhātukayoḥ iti . \\
\noindent \textbf{(P\_1,1.3.2)} tena bhavet iha niyamāt na syāt īhitā , īhitum , īhitavyam iti . \\
\noindent \textbf{(P\_1,1.3.2)} hrasvādyoḥ guṇaḥ tu aniyataḥ . \\
\noindent \textbf{(P\_1,1.3.2)} saḥ anantyasya api prāpnoti . \\
\noindent \textbf{(P\_1,1.3.2)} atha api evam niyamaḥ syāt . \\
\noindent \textbf{(P\_1,1.3.2)} pugantalaghūpadhasya sārvadhātukārdhadhātukayoḥ eva iti . \\
\noindent \textbf{(P\_1,1.3.2)} evam api sārvadhātukārdhadhātukayoḥ guṇaḥ aniyataḥ . \\
\noindent \textbf{(P\_1,1.3.2)} saḥ anantyasya api prāpnoti : īhitā , īhitum īhitavyam iti . \\
\noindent \textbf{(P\_1,1.3.2)} atha api ubhayataḥ niyamaḥ syāt : pugantalaghūpadhasya eva sārvadhātukārdhadhātukayoḥ , sārvadhātukārdhadhātukayoḥ eva pugantalaghūpadhasya iti , evam api ayam jusi guṇaḥ aniyataḥ . \\
\noindent \textbf{(P\_1,1.3.2)} saḥ anantyasya api prāpnoti : anenijuḥ , paryaviviṣuḥ iti . \\
\noindent \textbf{(P\_1,1.3.2)} evam tarhi na ayam taccheṣaḥ na api tadapavādaḥ . \\
\noindent \textbf{(P\_1,1.3.2)} anyat eva idam paribhāṣāntaram asambaddham anayā paribhāṣayā . \\
\noindent \textbf{(P\_1,1.3.2)} paribhāṣāntaram iti ca matvā kroṣṭrīyāḥ paṭhanti : niyamāt ikaḥ guṇavṛddhī bhavataḥ vipratiṣedhena iti . \\
\noindent \textbf{(P\_1,1.3.2)} yadi ca ayam taccheṣaḥ syāt tena eva tasya ayuktaḥ vipratiṣedhaḥ . \\
\noindent \textbf{(P\_1,1.3.2)} atha api tadapavādaḥ utsargāpavādayoḥ api ayuktaḥ vipratiṣedhaḥ . \\
\noindent \textbf{(P\_1,1.3.2)} tatra niyamasya avakāśaḥ : rājñaḥ ka ca , rājakīyam . \\
\noindent \textbf{(P\_1,1.3.2)} ikaḥ guṇavṛddhī* iti asya avakāśaḥ : cayanam , cāyakaḥ , lavanam , lāvakaḥ iti . \\
\noindent \textbf{(P\_1,1.3.2)} iha ubhayam prāpnoti : medyati mārṣṭi iti . \\
\noindent \textbf{(P\_1,1.3.2)} ikaḥ guṇavṛddhī* iti etat bhavati vipratiṣedhena . \\
\noindent \textbf{(P\_1,1.3.2)} na eṣaḥ yuktaḥ vipratiṣedhaḥ . \\
\noindent \textbf{(P\_1,1.3.2)} vipratiṣedhe hi param iti ucyate , pūrvaḥ ca ayam yogaḥ paraḥ niyamaḥ . \\
\noindent \textbf{(P\_1,1.3.2)} iṣṭavācī paraśabdaḥ . \\
\noindent \textbf{(P\_1,1.3.2)} vipratiṣedhe param yat iṣṭam tat bhavati . \\
\noindent \textbf{(P\_1,1.3.2)} evam api ayuktaḥ vipratiṣedhaḥ . \\
\noindent \textbf{(P\_1,1.3.2)} dvikāryayogaḥ hi vipratiṣedhaḥ . \\
\noindent \textbf{(P\_1,1.3.2)} na ca atra ekaḥ dvikāryayuktaḥ . \\
\noindent \textbf{(P\_1,1.3.2)} na avaśyam dvikāryayogaḥ eva vipratiṣedhaḥ . \\
\noindent \textbf{(P\_1,1.3.2)} kim tarhi. asambhavaḥ api . \\
\noindent \textbf{(P\_1,1.3.2)} saḥ ca asti atra asambhavaḥ . \\
\noindent \textbf{(P\_1,1.3.2)} kaḥ asau asmbhavaḥ . \\
\noindent \textbf{(P\_1,1.3.2)} iha tāvat : vṛkṣebhyaḥ , plakṣebhyaḥ iti ekaḥ sthānī dvau ādeśau . \\
\noindent \textbf{(P\_1,1.3.2)} na ca asti sambhavaḥ yat ekasya sthāninaḥ dvau ādeśau syātām . \\
\noindent \textbf{(P\_1,1.3.2)} iha idānīm medyati medyataḥ medyanti iti dvau sthāninau ekaḥ ādeśaḥ . \\
\noindent \textbf{(P\_1,1.3.2)} na ca asti sambhavaḥ yat dvayoḥ sthāninoḥ ekaḥ ādeśaḥ syāt iti eṣaḥ asambhavaḥ . \\
\noindent \textbf{(P\_1,1.3.2)} satyam etasmin asambhave yuktaḥ vipratiṣedhaḥ . \\
\noindent \textbf{(P\_1,1.3.2)} evam api ayuktaḥ vipratiṣedhaḥ . \\
\noindent \textbf{(P\_1,1.3.2)} dvayoḥ hi sāvakāśayoḥ samavasthitayoḥ vipratiṣedhaḥ bhavati . \\
\noindent \textbf{(P\_1,1.3.2)} anavakāśaḥ ca ayam yogaḥ . \\
\noindent \textbf{(P\_1,1.3.2)} nanu ca idānīm eva asya avakāśaḥ prakḷptaḥ : cayanam , cāyakaḥ , lavanam , lāvakaḥ iti . \\
\noindent \textbf{(P\_1,1.3.2)} atra api niyamaḥ prāpnoti . \\
\noindent \textbf{(P\_1,1.3.2)} yāvatā na aprāpte niyame ayam yogaḥ ārabhyate ataḥ tadapavādaḥ ayam yogaḥ bhavati . \\
\noindent \textbf{(P\_1,1.3.2)} utsargāpavādayoḥ ca ayuktaḥ vipratiṣedhaḥ . \\
\noindent \textbf{(P\_1,1.3.2)} atha api katham cit ikaḥ guṇavṛddhī* iti asya avakāśaḥ syāt , evam api yathā iha vipratiṣedhāt ikaḥ guṇaḥ bhavati : medyati medyataḥ medyanti , evam iha api syāt : anenijuḥ , paryaveviṣuḥ iti . \\
\noindent \textbf{(P\_1,1.3.2)} evam tarhi vṛddhiḥ bhavati guṇaḥ bhavati iti yatra brūyāt ikaḥ iti etat tatra upasthitam draṣṭavyam . \\
\noindent \textbf{(P\_1,1.3.2)} kim kṛtam bhavati . \\
\noindent \textbf{(P\_1,1.3.2)} dvitīyā ṣaṣṭhī prāduḥ bhāvyate . \\
\noindent \textbf{(P\_1,1.3.2)} tatra kāmacāraḥ : gṛhyamāṇena vā ikam viśeṣayitum ikā vā gṛhyamāṇam . \\
\noindent \textbf{(P\_1,1.3.2)} yāvatā kāmacāraḥ iha tāvat : midipugantalaghūpadharcchidṛśikṣiprakṣudreṣu gṛhyamāṇena ikam viśeṣayiṣyāmaḥ : eteṣām yaḥ ik iti . \\
\noindent \textbf{(P\_1,1.3.2)} iha idānīm : jusisārvadhātukārdhadhātukahrasvādyoḥ guṇeṣu ikā gṛhyamāṇam viśeṣayiṣyāmaḥ : eteṣām guṇaḥ bhavati ikaḥ . \\
\noindent \textbf{(P\_1,1.3.2)} igantānām iti . \\
\noindent \textbf{(P\_1,1.3.2)} atha vā sarvatra eva sthānī nirdiśyate . \\
\noindent \textbf{(P\_1,1.3.2)} iha tāvat : mideḥ iti , avibhaktikaḥ nirdeśaḥ : mida , eḥ , mideḥ , mideḥ iti . \\
\noindent \textbf{(P\_1,1.3.2)} atha vā ṣaṣṭhīsamāsaḥ bhaviṣyati midaḥ iḥ , midiḥ , mideḥ iti . \\
\noindent \textbf{(P\_1,1.3.2)} pugantalaghūpadhasya iti na evam vijñāyate : pugantasya aṅgasya laghūpadhasya ca iti . \\
\noindent \textbf{(P\_1,1.3.2)} katham tarhi . \\
\noindent \textbf{(P\_1,1.3.2)} puki antaḥ pugantaḥ , laghvī upadhā laghūpadhā , pugantaḥ ca laghūpadhā ca pugantalaghūpadham pugantalaghūpadhasya iti . \\
\noindent \textbf{(P\_1,1.3.2)} avaśyam ca etat evam vijñeyam . \\
\noindent \textbf{(P\_1,1.3.2)} aṅgaviśeṣaṇe hi sati iha prasajyeta : bhinatti chinatti iti . \\
\noindent \textbf{(P\_1,1.3.2)} ṛccheḥ api praśliṣṭanirdeśaḥ ayam . \\
\noindent \textbf{(P\_1,1.3.2)} ṛcchati , ṛ , ṛ , ṝtām ṛcchatyṝtām iti . \\
\noindent \textbf{(P\_1,1.3.2)} dṛśeḥ api yogavibhāgaḥ kariṣyate . \\
\noindent \textbf{(P\_1,1.3.2)} uḥ aṅi guṇaḥ . \\
\noindent \textbf{(P\_1,1.3.2)} uḥ aṅi guṇaḥ bhavati . \\
\noindent \textbf{(P\_1,1.3.2)} tataḥ dṛśeḥ . \\
\noindent \textbf{(P\_1,1.3.2)} dṛśeḥ ca aṅi guṇaḥ bhavati . \\
\noindent \textbf{(P\_1,1.3.2)} uḥ iti eva. kṣiprakṣudrayoḥ api yaṇādiparam guṇa iti iyatā siddham . \\
\noindent \textbf{(P\_1,1.3.2)} saḥ ayam evam siddhe sati yat pūrvagrahaṇam karoti tasya etat prayojanam ikaḥ yathā syāt anikaḥ mā bhūt iti . \\
\noindent \textbf{(P\_1,1.3.3)} atha vṛddhigrahaṇam kimartham . \\
\noindent \textbf{(P\_1,1.3.3)} kim viśeṣeṇa vṛddhigrahaṇam codyate na punaḥ guṇagrahaṇam api . \\
\noindent \textbf{(P\_1,1.3.3)} yadi kim cit guṇagrahaṇasya prayojanam asti vṛddhigrahaṇasya api tat bhavitum arhati . \\
\noindent \textbf{(P\_1,1.3.3)} kaḥ vā viśeṣaḥ . \\
\noindent \textbf{(P\_1,1.3.3)} ayam asti viśeṣaḥ . \\
\noindent \textbf{(P\_1,1.3.3)} guṇavidhau na kva cit sthānī nirdiśyate . \\
\noindent \textbf{(P\_1,1.3.3)} tatra avaśyam sthāninirdeśārtham guṇagrahaṇam kartavyam . \\
\noindent \textbf{(P\_1,1.3.3)} vṛddhividhau punaḥ sarvatra eva sthānī nirdiśyate . \\
\noindent \textbf{(P\_1,1.3.3)} acaḥ ñṇiti ataḥ upadhāyāḥ taddhiteṣu acām ādeḥ iti . \\
\noindent \textbf{(P\_1,1.3.3)} ataḥ uttaram paṭhati . \\
\noindent \textbf{(P\_1,1.3.3)} vṛddhigrahaṇam uttarārtham . \\
\noindent \textbf{(P\_1,1.3.3)} vṛddhigrahaṇam kriyate uttarārtham . \\
\noindent \textbf{(P\_1,1.3.3)} kṅiti iti pratiṣedham vakṣyati . \\
\noindent \textbf{(P\_1,1.3.3)} saḥ vṛddheḥ api yathā syāt . \\
\noindent \textbf{(P\_1,1.3.3)} kaḥ ca idānīm kṅitpratyayeṣu vṛddheḥ prasaṅgaḥ yāvatā ñṇiti iti ucyate . \\
\noindent \textbf{(P\_1,1.3.3)} tat ca mṛjyartham . \\
\noindent \textbf{(P\_1,1.3.3)} mṛjeḥ vṛddhiḥ aviśeṣeṇa ucyate . \\
\noindent \textbf{(P\_1,1.3.3)} sā kṅiti mā bhūt : mṛṣṭaḥ , mṛṣṭavān iti . \\
\noindent \textbf{(P\_1,1.3.3)} ihārtham ca api mṛjyartham vṛddhigrahaṇam kartavyam . \\
\noindent \textbf{(P\_1,1.3.3)} mṛjeḥ vṛddhiḥ aviśeṣeṇa ucyate . \\
\noindent \textbf{(P\_1,1.3.3)} sā ikaḥ yathā syāt . \\
\noindent \textbf{(P\_1,1.3.3)} anikāḥ mā bhūt iti . \\
\noindent \textbf{(P\_1,1.3.3)} mṛjyartham iti cet yogavibhāgāt siddham . \\
\noindent \textbf{(P\_1,1.3.3)} mṛjyartham iti cet yogavibhāgaḥ kariṣyate . \\
\noindent \textbf{(P\_1,1.3.3)} mṛjeḥ vṛddhiḥ . \\
\noindent \textbf{(P\_1,1.3.3)} tataḥ ñṇiti . \\
\noindent \textbf{(P\_1,1.3.3)} ñiti ṇiti ca vṛddhiḥ bhavati . \\
\noindent \textbf{(P\_1,1.3.3)} acaḥ iti eva . \\
\noindent \textbf{(P\_1,1.3.3)} yadi acaḥ vṛddhiḥ ucyate nyamārṭ : aṭaḥ api vṛddhiḥ prāpnoti . \\
\noindent \textbf{(P\_1,1.3.3)} aṭi ca uktam . \\
\noindent \textbf{(P\_1,1.3.3)} kim uktam. anantyavikāre antyasadeśasya kāryam bhavati iti . \\
\noindent \textbf{(P\_1,1.3.3)} vṛddhipratiṣedhānupapattiḥ tu ikprakaraṇāt . \\
\noindent \textbf{(P\_1,1.3.3)} vṛddheḥ tu pratiṣedhaḥ na upapadyate . \\
\noindent \textbf{(P\_1,1.3.3)} kim kāraṇam . \\
\noindent \textbf{(P\_1,1.3.3)} ikprakaraṇāt . \\
\noindent \textbf{(P\_1,1.3.3)} iglakṣaṇayoḥ guṇavṛddhyoḥ pratiṣedhaḥ na ca evam sati mṛjeḥ iglakṣaṇā vṛddhiḥ bhavati . \\
\noindent \textbf{(P\_1,1.3.3)} tasmāt mṛjeḥ iglakṣaṇā vṛddhiḥ eṣitavyā . \\
\noindent \textbf{(P\_1,1.3.3)} evam tarhi iha anye vaiyākaraṇāḥ mṛjeḥ ajādau saṅkrame vibhāṣā vṛddhim ārabhante : parimṛjanti parimārjanti parimṛjantu parimārjantu parimamṛjatuḥ , parimamārjatuḥ ityādyartham . \\
\noindent \textbf{(P\_1,1.3.3)} tat iha api sādhyam . \\
\noindent \textbf{(P\_1,1.3.3)} tasmin sādhye yogavibhāgaḥ kariṣyate . \\
\noindent \textbf{(P\_1,1.3.3)} mṛjeḥ vṛddhiḥ acaḥ bhavati . \\
\noindent \textbf{(P\_1,1.3.3)} tataḥ aci kṅiti . \\
\noindent \textbf{(P\_1,1.3.3)} ajādau ca kṅiti mṛjeḥ vṛddhiḥ bhavati : parimārjanti parimārjantu . \\
\noindent \textbf{(P\_1,1.3.3)} kimartham idam . \\
\noindent \textbf{(P\_1,1.3.3)} niyamārtham : ajādau eva kṅiti na anyatra . \\
\noindent \textbf{(P\_1,1.3.3)} kva anyatra mā bhūt . \\
\noindent \textbf{(P\_1,1.3.3)} mṛṣṭaḥ , mṛṣṭavān iti . \\
\noindent \textbf{(P\_1,1.3.3)} tataḥ vā . \\
\noindent \textbf{(P\_1,1.3.3)} vā aci kṅiti mṛjeḥ vṛddhiḥ bhavati . \\
\noindent \textbf{(P\_1,1.3.3)} parimṛjanti , parimārjanti , parimamṛjatuḥ , parimamārjatuḥ iti . \\
\noindent \textbf{(P\_1,1.3.3)} ihārtham eva sijartham vṛddhigrahaṇam kartavyam . \\
\noindent \textbf{(P\_1,1.3.3)} sici vṛddhiḥ aviśeṣeṇa ucyate . \\
\noindent \textbf{(P\_1,1.3.3)} sā ikaḥ yathā syāt , anikaḥ mā bhūt iti . \\
\noindent \textbf{(P\_1,1.3.3)} kasya punaḥ anikaḥ prāpnoti . \\
\noindent \textbf{(P\_1,1.3.3)} akārasya : acikīrṣīt , ajihīrṣīt . \\
\noindent \textbf{(P\_1,1.3.3)} na etat asti . \\
\noindent \textbf{(P\_1,1.3.3)} lopaḥ atra bādhakaḥ bhaviṣyati . \\
\noindent \textbf{(P\_1,1.3.3)} ākārasya tarhi prāpnoti : ayāsīt , avāsīt . \\
\noindent \textbf{(P\_1,1.3.3)} na asti atra viśeṣaḥ satyām vṛddhau asatyām vā . \\
\noindent \textbf{(P\_1,1.3.3)} sandhyakṣarasya tarhi prāpnoti . \\
\noindent \textbf{(P\_1,1.3.3)} na eva sandhyakṣaram antyam asti . \\
\noindent \textbf{(P\_1,1.3.3)} nanu ca idam asti ḍhalope kṛte udavoḍhām udavoḍham udavoḍha iti . \\
\noindent \textbf{(P\_1,1.3.3)} na etat asti . \\
\noindent \textbf{(P\_1,1.3.3)} asiddhaḥ ḍhalopaḥ . \\
\noindent \textbf{(P\_1,1.3.3)} tasya asiddhatvāt na etat antyam bhavati . \\
\noindent \textbf{(P\_1,1.3.3)} vyañjanasya tarhi prāpnoti : abhaitsīt , acchaitsīt . \\
\noindent \textbf{(P\_1,1.3.3)} halantalakṣaṇā vṛddhiḥ bādhikā bhaviṣyati . \\
\noindent \textbf{(P\_1,1.3.3)} yatra tarhi sā pratiṣidhyate : akoṣīt , amoṣīt . \\
\noindent \textbf{(P\_1,1.3.3)} sici vṛddheḥ api eṣaḥ pratiṣedhaḥ . \\
\noindent \textbf{(P\_1,1.3.3)} katham . \\
\noindent \textbf{(P\_1,1.3.3)} lakṣaṇam hi nāma dhvanati bhramati muhūrtam api na avatiṣṭhate . \\
\noindent \textbf{(P\_1,1.3.3)} atha vā sici vṛddhiḥ parasmaipadeṣu iti sici vṛddhiḥ prāpnoti . \\
\noindent \textbf{(P\_1,1.3.3)} tasyāḥ halantalakṣaṇā vṛddhiḥ bādhikā . \\
\noindent \textbf{(P\_1,1.3.3)} tasyāḥ api na iṭi iti pratiṣedhaḥ . \\
\noindent \textbf{(P\_1,1.3.3)} asti punaḥ kva cid anyatra api apavāde pratiṣiddhe utsargaḥ api na bhavati . \\
\noindent \textbf{(P\_1,1.3.3)} asti iti āha . \\
\noindent \textbf{(P\_1,1.3.3)} sujāte* aśvasūnṛte , adhvaryo* adribhiḥ sutam , śukram te* anyat iti . \\
\noindent \textbf{(P\_1,1.3.3)} pūrvarūpatve pratiṣiddhe ayādayaḥ api na bhavanti . \\
\noindent \textbf{(P\_1,1.3.3)} uttarārtham eva tarhi sijartham vṛddhigrahaṇam kartavyam . \\
\noindent \textbf{(P\_1,1.3.3)} sici vṛddhiḥ aviśeṣeṇa ucyate . \\
\noindent \textbf{(P\_1,1.3.3)} sā kṅiti mā bhūt nyanuvīt , nyadhuvīt . \\
\noindent \textbf{(P\_1,1.3.3)} na etat asti prayojanam . \\
\noindent \textbf{(P\_1,1.3.3)} antaraṅgatvāt atra uvaṅādeśe kṛte anantyatvāt vṛddhiḥ na bhaviṣyati . \\
\noindent \textbf{(P\_1,1.3.3)} yadi tarhi sici antaraṅgam bhavati , akārṣīt , ahārṣīt : guṇe kṛte raparatve ca anantyatvāt vṛddhiḥ na prāpnoti . \\
\noindent \textbf{(P\_1,1.3.3)} mā bhūt evam . \\
\noindent \textbf{(P\_1,1.3.3)} halantasya iti evam bhaviṣyati . \\
\noindent \textbf{(P\_1,1.3.3)} iha tarhi : nyastārīt , vyadārīt : guṇe raparatve ca anantyatvāt vṛddhiḥ na prāpnoti . \\
\noindent \textbf{(P\_1,1.3.3)} halantalakṣaṇāyāḥ ca na iṭi iti pratiṣedhaḥ . \\
\noindent \textbf{(P\_1,1.3.3)} mā bhūt evam . \\
\noindent \textbf{(P\_1,1.3.3)} rlāntasya iti evam bhaviṣyati . \\
\noindent \textbf{(P\_1,1.3.3)} iha tarhi : alāvīt apāvīt : guṇe kṛte avādeśe ca anantyatvāt vṛddhiḥ na prāpnoti . \\
\noindent \textbf{(P\_1,1.3.3)} halantalakṣaṇāyāḥ ca na iṭi iti pratiṣedhaḥ . \\
\noindent \textbf{(P\_1,1.3.3)} mā bhūt evam . \\
\noindent \textbf{(P\_1,1.3.3)} rlāntasya iti evam bhaviṣyati . \\
\noindent \textbf{(P\_1,1.3.3)} rlāntasya iti ucyate na ca idam rlāntam . \\
\noindent \textbf{(P\_1,1.3.3)} rlāntasya iti atra vakāraḥ api nirdiśyate . \\
\noindent \textbf{(P\_1,1.3.3)} kim vakāraḥ na śrūyate . \\
\noindent \textbf{(P\_1,1.3.3)} luptanirdiṣtaḥ vakāraḥ . \\
\noindent \textbf{(P\_1,1.3.3)} yadi evam mā bhavān mavīt : atra api prāpnoti . \\
\noindent \textbf{(P\_1,1.3.3)} avimavyoḥ iti vakṣyāmi . \\
\noindent \textbf{(P\_1,1.3.3)} tat vaktavyam . \\
\noindent \textbf{(P\_1,1.3.3)} ṇiśvibhyām tau nimātavyau . \\
\noindent \textbf{(P\_1,1.3.3)} yadi api etat ucyate atha vā etarhi ṇiśvyoḥ pratiṣedhaḥ na vaktavyaḥ bhavati . \\
\noindent \textbf{(P\_1,1.3.3)} guṇe kṛte ayādeśe ca yāntānam na iti pratiṣedhaḥ bhaviṣyati . \\
\noindent \textbf{(P\_1,1.3.3)} evam tarhi ācāryapravṛttiḥ jñāpayati na sici antaraṅgam bhavati iti yat ayam ataḥ halādeḥ laghoḥ iti akāragrahaṇam karoti . \\
\noindent \textbf{(P\_1,1.3.3)} katham kṛtvā jñāpakam . \\
\noindent \textbf{(P\_1,1.3.3)} akāragrahaṇasya etat prayojanam iha mā bhūt : akoṣīt , amoṣīt . \\
\noindent \textbf{(P\_1,1.3.3)} yadi sici antaraṅgam syāt akāragrahaṇam anarthakam syāt . \\
\noindent \textbf{(P\_1,1.3.3)} guṇe kṛte alaghutvāt vṛddhiḥ na bhaviṣyati . \\
\noindent \textbf{(P\_1,1.3.3)} paśyati tu ācāryaḥ na sici antaraṅgam bhavati iti . \\
\noindent \textbf{(P\_1,1.3.3)} tataḥ akāragrahaṇam karoti . \\
\noindent \textbf{(P\_1,1.3.3)} na etat asti jñāpakam . \\
\noindent \textbf{(P\_1,1.3.3)} asti anyat etasya vacane prayojanam . \\
\noindent \textbf{(P\_1,1.3.3)} kim . \\
\noindent \textbf{(P\_1,1.3.3)} yatra guṇaḥ pratiṣidhyate tadartham etat syāt : nyakuṭīt , nyapuṭīt iti . \\
\noindent \textbf{(P\_1,1.3.3)} yat tarhi ṇiśvyoḥ pratiṣedham śāsti tena na iha antaraṅgam asti iti darśayati . \\
\noindent \textbf{(P\_1,1.3.3)} yat ca karoti akāragrahaṇam laghoḥ iti kṛte api . \\
\noindent \textbf{(P\_1,1.3.3)} tasmāt iglakṣaṇā vṛddhiḥ . tasmāt iglakṣaṇā vṛddhiḥ āstheyā . \\
\noindent \textbf{(P\_1,1.3.4)} ṣaṣṭhyāḥ sthāneyogatvāt ignivṛttiḥ . ṣaṣṭhyāḥ sthāneyogatvāt sarveṣām ikām nivṛttiḥ prāpnoti . \\
\noindent \textbf{(P\_1,1.3.4)} asya api prāpnoti : dadhi madhu . \\
\noindent \textbf{(P\_1,1.3.4)} punarvacanam idānīm kimartham syāt . \\
\noindent \textbf{(P\_1,1.3.4)} anyatarārtham punarvacanam . \\
\noindent \textbf{(P\_1,1.3.4)} anyatarārtham etat syāt . \\
\noindent \textbf{(P\_1,1.3.4)} sārvadhādukārdhadhātukayoḥ guṇaḥ eva iti . \\
\noindent \textbf{(P\_1,1.3.4)} prasāraṇe ca . \\
\noindent \textbf{(P\_1,1.3.4)} prasāraṇe ca sarveṣam yaṇām nivṛttiḥ prāpnoti . \\
\noindent \textbf{(P\_1,1.3.4)} asya api prāpnoti : yātā vātā . \\
\noindent \textbf{(P\_1,1.3.4)} punarvacanam idānīm kimartham syāt . \\
\noindent \textbf{(P\_1,1.3.4)} viṣayārtham punarvacanam . viṣayārtham etat syāt . \\
\noindent \textbf{(P\_1,1.3.4)} vacisvapiyajādīnām kiti eva iti . \\
\noindent \textbf{(P\_1,1.3.4)} uḥ aṇ rapare ca . \\
\noindent \textbf{(P\_1,1.3.4)} uḥ aṇ rapare ca sarvarkārāṇām nivṛttiḥ prāpnoti . \\
\noindent \textbf{(P\_1,1.3.4)} asya api prāpnoti kartṛ hartṛ . \\
\noindent \textbf{(P\_1,1.3.4)} siddham tu ṣaṣṭhyadhikāre vacanāt . \\
\noindent \textbf{(P\_1,1.3.4)} siddham etat . \\
\noindent \textbf{(P\_1,1.3.4)} katham . \\
\noindent \textbf{(P\_1,1.3.4)} ṣaṣṭhyadhikāre ime yogāḥ kartavyāḥ . \\
\noindent \textbf{(P\_1,1.3.4)} ekaḥ tāvat kriyate tatra eva . \\
\noindent \textbf{(P\_1,1.3.4)} imau api yogau ṣaṣṭhadhikāram anuvartiṣyete . \\
\noindent \textbf{(P\_1,1.3.4)} atha vā ṣaṣṭhadhikāre imau yogau apekṣiṣyāmahe . \\
\noindent \textbf{(P\_1,1.3.4)} atha vā idam tāvad ayam praṣṭavyaḥ . \\
\noindent \textbf{(P\_1,1.3.4)} sārvadhātukārdhadhātukayoḥ guṇaḥ bhavati iti iha kasmāt na bhavati : yātā vātā . \\
\noindent \textbf{(P\_1,1.3.4)} idam tatra apekṣiṣyate ikaḥ guṇavṛddhī* iti . \\
\noindent \textbf{(P\_1,1.3.4)} yathā eva tarhi idam tatra apekṣiṣyate evam iha api tad apekṣiṣyāmahe sārvadhātukārdhadhātukayoḥ ikaḥ guṇavṛddhī* iti. \\
\noindent \textbf{(P\_1,1.4.1)} dhātugrahaṇam kimartham . \\
\noindent \textbf{(P\_1,1.4.1)} iha mā bhūt: lūñ lavitā lavitum pūñ pavitā pavitum . \\
\noindent \textbf{(P\_1,1.4.1)} ārdhadhātuke iti kimartham . \\
\noindent \textbf{(P\_1,1.4.1)} tridhā baddhaḥ vṛṣabhaḥ roravīti . \\
\noindent \textbf{(P\_1,1.4.1)} kim punaḥ idam ārdhadhātukagrahaṇam lopaviśeṣaṇam : ārdhadhātukanimitte lope sati ye guṇavṛddhī prāpnutaḥ te na bhavataḥ iti , āhosvit guṇavṛddhiviśeṣaṇam ārdhadhātukagrahaṇam : dhātulope sati ārdhadhātukanimitte ye guṇavṛddhī prāpnutaḥ te na bhavataḥ iti . \\
\noindent \textbf{(P\_1,1.4.1)} kim ca ataḥ . \\
\noindent \textbf{(P\_1,1.4.1)} yadi lopaviśeṣaṇam upeddhaḥ preddhaḥ atra api prāpnoti . \\
\noindent \textbf{(P\_1,1.4.1)} atha guṇavṛddhiviśeṣaṇam knopayati iti atra api prāpnoti . \\
\noindent \textbf{(P\_1,1.4.1)} yathā icchasi tathā astu . \\
\noindent \textbf{(P\_1,1.4.1)} astu lopaviśeṣaṇam . \\
\noindent \textbf{(P\_1,1.4.1)} katham upeddhaḥ preddhaḥ iti . \\
\noindent \textbf{(P\_1,1.4.1)} bahiraṅgaḥ guṇaḥ antaraṅgaḥ pratiṣedhaḥ . \\
\noindent \textbf{(P\_1,1.4.1)} asiddham bahiraṅgam antaraṅge . \\
\noindent \textbf{(P\_1,1.4.1)} yadi evam na arthaḥ dhātugrahaṇena . \\
\noindent \textbf{(P\_1,1.4.1)} iha kasmāt na bhavati: lūñ lavitā lavitum . \\
\noindent \textbf{(P\_1,1.4.1)} ārdhadhātukanimitte lope pratiṣedhaḥ na ca eṣaḥ ārdhadhātukanimittaḥ lopaḥ . \\
\noindent \textbf{(P\_1,1.4.1)} atha vā punaḥ astu guṇavṛddhiviśeṣaṇam . \\
\noindent \textbf{(P\_1,1.4.1)} nanu ca uktam knopayati iti atra api prāpnoti iti . \\
\noindent \textbf{(P\_1,1.4.1)} na eṣaḥ doṣaḥ . \\
\noindent \textbf{(P\_1,1.4.1)} nipātanāt siddham . \\
\noindent \textbf{(P\_1,1.4.1)} kim nipātanam . \\
\noindent \textbf{(P\_1,1.4.1)} cele knopeḥ iti \\
\noindent \textbf{(P\_1,1.4.2)} parigaṇanam kartavyam . \\
\noindent \textbf{(P\_1,1.4.2)} yaṅyakkyavalope pratiṣedhaḥ . \\
\noindent \textbf{(P\_1,1.4.2)} yaṅyakkyavalope pratiṣedhaḥ vaktavyaḥ . \\
\noindent \textbf{(P\_1,1.4.2)} yaṅ: bebhiditā marīmṛjaḥ . \\
\noindent \textbf{(P\_1,1.4.2)} yak: kuṣubhitā magadhakaḥ . \\
\noindent \textbf{(P\_1,1.4.2)} kya: samidhitā dṛṣadakaḥ . \\
\noindent \textbf{(P\_1,1.4.2)} valope : jīradānuḥ . \\
\noindent \textbf{(P\_1,1.4.2)} kim prayojanam . \\
\noindent \textbf{(P\_1,1.4.2)} numlopasrivyanubandhalope apratiṣedhārtham . \\
\noindent \textbf{(P\_1,1.4.2)} numlope srivyanubandhalope ca pratiṣedhaḥ mā bhūt iti . \\
\noindent \textbf{(P\_1,1.4.2)} numlope: abhāji rāgaḥ upabarhaṇam . \\
\noindent \textbf{(P\_1,1.4.2)} sriveḥ : āsremāṇam . \\
\noindent \textbf{(P\_1,1.4.2)} anubandhalope : lūñ lavitā lavitum . \\
\noindent \textbf{(P\_1,1.4.2)} yadi parigaṇanam kriyate syadaḥ, praśrathaḥ, himaśrathaḥ iti atra api prāpnoti . \\
\noindent \textbf{(P\_1,1.4.2)} vakṣyati etat nipātanāt syadādiṣu iti . \\
\noindent \textbf{(P\_1,1.4.2)} tat tarhi parigaṇanam kartavyam . \\
\noindent \textbf{(P\_1,1.4.2)} na kartavyam . \\
\noindent \textbf{(P\_1,1.4.2)} numlope kasmāt na bhavati . \\
\noindent \textbf{(P\_1,1.4.2)} ikprakaraṇāt numlope vṛddhiḥ . \\
\noindent \textbf{(P\_1,1.4.2)} iglakṣaṇayoḥ guṇavṛddhyoḥ pratiṣedhaḥ na ca eṣā iglakṣaṇā vṛddhiḥ . \\
\noindent \textbf{(P\_1,1.4.2)} yadi iglakṣaṇayoḥ guṇavṛddhyoḥ pratiṣedhaḥ syadaḥ, praśrathaḥ, himaśrathaḥ iti atra na prāpnoti . \\
\noindent \textbf{(P\_1,1.4.2)} iha ca prāpnoti: avodaḥ, edhaḥ, odmaḥ iti . \\
\noindent \textbf{(P\_1,1.4.2)} nipātanāt syadādiṣu . nipātanāt syadādiṣu pratiṣedhaḥ bhaviṣyati na ca bhaviṣyati . \\
\noindent \textbf{(P\_1,1.4.2)} yadi iglakṣaṇayoḥ guṇavṛddhyoḥ pratiṣedhaḥ srivyanubandhalope katham sriveḥ āsremāṇam lūñ lavitā . \\
\noindent \textbf{(P\_1,1.4.2)} pratyayāśrayatvāt anyatra siddham . \\
\noindent \textbf{(P\_1,1.4.2)} ārdhadhātukanimitte lope pratiṣedhaḥ na ca eṣaḥ ārdhadhātukanimittaḥ lopaḥ . \\
\noindent \textbf{(P\_1,1.4.2)} yadi ārdhadhātukanimitte lope pratiṣedhaḥ jīradānuḥ atra na prāpnoti . \\
\noindent \textbf{(P\_1,1.4.2)} raki jyaḥ prasāraṇam . na etat jīveḥ rūpam . \\
\noindent \textbf{(P\_1,1.4.2)} raki etat jyaḥ prasāraṇam . \\
\noindent \textbf{(P\_1,1.4.2)} yāvatā ca idānīm raki jīveḥ api siddham bhavati . \\
\noindent \textbf{(P\_1,1.4.2)} katham upabarhaṇam . \\
\noindent \textbf{(P\_1,1.4.2)} bṛhiḥ prakṛtyantaram . \\
\noindent \textbf{(P\_1,1.4.2)} katham jñāyate bṛhiḥ prakṛtyantaram iti . \\
\noindent \textbf{(P\_1,1.4.2)} aci iti hi lopaḥ ucyate anajādau api dṛśyate: nibṛhyate . \\
\noindent \textbf{(P\_1,1.4.2)} aniṭi iti ca ucyate . \\
\noindent \textbf{(P\_1,1.4.2)} iḍādau api dṛśyate: nibarhitā nibarhitum iti . \\
\noindent \textbf{(P\_1,1.4.2)} ajādau api na dṛśyate: bṛṃhayati bṛṃhakaḥ . \\
\noindent \textbf{(P\_1,1.4.2)} tasmāt na arthaḥ parigaṇanena . \\
\noindent \textbf{(P\_1,1.4.2)} yadi parigaṇanam na kriyate bhedyate chedyate atra api prāpnoti . \\
\noindent \textbf{(P\_1,1.4.2)} na eṣaḥ doṣaḥ . \\
\noindent \textbf{(P\_1,1.4.2)} dhātulope iti na evam vijñāyate: dhātoḥ lopaḥ dhātulopaḥ, dhātulope iti . \\
\noindent \textbf{(P\_1,1.4.2)} katham tarhi . \\
\noindent \textbf{(P\_1,1.4.2)} dhātoḥ lopaḥ asmin tat idam dhātulopam, dhātulope iti . \\
\noindent \textbf{(P\_1,1.4.2)} tasmāt iglakṣaṇayoḥ guṇavṛddhyoḥ pratiṣedhaḥ . \\
\noindent \textbf{(P\_1,1.4.2)} yadi tarhi iglakṣaṇayoḥ guṇavṛddhyoḥ pratiṣedhaḥ pāpacakaḥ, pāpaṭhakaḥ, magadhakaḥ, dṛṣadakaḥ atra na prāpnoti . \\
\noindent \textbf{(P\_1,1.4.2)} allopasya sthānivatvāt . \\
\noindent \textbf{(P\_1,1.4.2)} akāralope kṛte tasya sthānivatvāt guṇavṛddhī na bhaviṣyataḥ . \\
\noindent \textbf{(P\_1,1.4.3)} anārambhaḥ vā . \\
\noindent \textbf{(P\_1,1.4.3)} anārambhaḥ vā punaḥ asya yogasya nyāyyaḥ . \\
\noindent \textbf{(P\_1,1.4.3)} katham bebhiditā, marīmṛjakaḥ, kuṣubhitā samidhitā iti . \\
\noindent \textbf{(P\_1,1.4.3)} atra api akāralope kṛte tasya sthānivatvāt guṇavṛddhī na bhaviṣyataḥ . \\
\noindent \textbf{(P\_1,1.4.3)} yatra tarhi sthānivadbhāvaḥ na asti tadartham ayam yogaḥ vaktavyaḥ . \\
\noindent \textbf{(P\_1,1.4.3)} kva ca sthānivadbhāvaḥ na asti . \\
\noindent \textbf{(P\_1,1.4.3)} yatra halacoḥ ādeśaḥ: loluvaḥ popuvaḥ marīmṛjaḥ sarīsṛpaḥ iti . \\
\noindent \textbf{(P\_1,1.4.3)} atra api akāralope kṛte tasya sthānivatvāt guṇavṛddhī na bhaviṣyataḥ . \\
\noindent \textbf{(P\_1,1.4.3)} luki kṛte na prāpnoti . \\
\noindent \textbf{(P\_1,1.4.3)} idam iha sampradhāryam: luk kriyatām allopaḥ iti kim atra kartavyam . \\
\noindent \textbf{(P\_1,1.4.3)} paratvāt allopaḥ . \\
\noindent \textbf{(P\_1,1.4.3)} nityaḥ luk . \\
\noindent \textbf{(P\_1,1.4.3)} kṛte api allope prāpnoti akṛte api prāpnoti . \\
\noindent \textbf{(P\_1,1.4.3)} luk api anityaḥ . \\
\noindent \textbf{(P\_1,1.4.3)} katham . \\
\noindent \textbf{(P\_1,1.4.3)} anyasya kṛte allope prāpnoti anyasya akṛte . \\
\noindent \textbf{(P\_1,1.4.3)} śabdāntarasya ca prāpnuvan vidhiḥ anityaḥ bhavati . \\
\noindent \textbf{(P\_1,1.4.3)} anavakāśaḥ tarhi luk. sāvakāśaḥ luk . \\
\noindent \textbf{(P\_1,1.4.3)} kaḥ avakāśaḥ . \\
\noindent \textbf{(P\_1,1.4.3)} avaśiṣṭaḥ . \\
\noindent \textbf{(P\_1,1.4.3)} atham katham cit anavakāśaḥ luk syāt evam api na doṣaḥ . \\
\noindent \textbf{(P\_1,1.4.3)} allope yogavibhāgaḥ kariṣyate : ataḥ lopaḥ . \\
\noindent \textbf{(P\_1,1.4.3)} tataḥ yasya : yasya ca lopaḥ bhavati . \\
\noindent \textbf{(P\_1,1.4.3)} ataḥ iti eva . \\
\noindent \textbf{(P\_1,1.4.3)} kimartham idam . \\
\noindent \textbf{(P\_1,1.4.3)} lukam vakṣyati tadbādhanārtham . \\
\noindent \textbf{(P\_1,1.4.3)} tato halaḥ . \\
\noindent \textbf{(P\_1,1.4.3)} halaḥ uttarasya ca yasya lopaḥ bhavati iti . \\
\noindent \textbf{(P\_1,1.4.3)} iha api paratvāt yogavibhāgāt va lopaḥ lukam bādheta: kṛṣṇaḥ nonāva vṛṣabhaḥ yadi idam . \\
\noindent \textbf{(P\_1,1.4.3)} nonūyateḥ nonāva . \\
\noindent \textbf{(P\_1,1.4.3)} samānāśrayaḥ luk lopena bādhyate . \\
\noindent \textbf{(P\_1,1.4.3)} kaḥ ca samānāśrayaḥ . \\
\noindent \textbf{(P\_1,1.4.3)} yaḥ pratyayāśrayaḥ . \\
\noindent \textbf{(P\_1,1.4.3)} atra ca prāk eva pratyayotpatteḥ luk bhavati . \\
\noindent \textbf{(P\_1,1.4.3)} katham syadaḥ, praśrathaḥ, himaśrathaḥ, jīradānuḥ, nikucitaḥ iti . \\
\noindent \textbf{(P\_1,1.4.3)} uktam śeṣe . \\
\noindent \textbf{(P\_1,1.4.3)} kim uktam . \\
\noindent \textbf{(P\_1,1.4.3)} nipātanāt syadādiṣu . \\
\noindent \textbf{(P\_1,1.4.3)} pratyayāśratvāt anyatra siddham . \\
\noindent \textbf{(P\_1,1.4.3)} raki jyaḥ prasāraṇam iti . \\
\noindent \textbf{(P\_1,1.4.3)} nikucite api uktam sannipātalakṣaṇaḥ vidhiḥ animittam tadvighātasya iti . \\
\noindent \textbf{(P\_1,1.5.1)} kṅiti pratiṣedhe tannimittagrahaṇam . \\
\noindent \textbf{(P\_1,1.5.1)} kṅiti pratiṣedhe tannimittagrahaṇam kartavyam . \\
\noindent \textbf{(P\_1,1.5.1)} kṅinnimitte ye guṇavṛddhī prāpnutaḥ te na bhavataḥ iti vaktavyam . \\
\noindent \textbf{(P\_1,1.5.1)} kim prayojanam . \\
\noindent \textbf{(P\_1,1.5.1)} upadhāroravītyartham . upadhārtham roravītyartham ca . \\
\noindent \textbf{(P\_1,1.5.1)} upadhārtham tāvat : bhinnaḥ , bhinnavān iti . \\
\noindent \textbf{(P\_1,1.5.1)} kim punaḥ kāraṇam na sidhyati . \\
\noindent \textbf{(P\_1,1.5.1)} kṅiti iti ucyate . \\
\noindent \textbf{(P\_1,1.5.1)} tena yatra kṅiti anantaraḥ guṇabhāvī asti tatra eva syāt: citam stutam . \\
\noindent \textbf{(P\_1,1.5.1)} iha tu na syāt: bhinnaḥ , bhinnavān iti . \\
\noindent \textbf{(P\_1,1.5.1)} nanu ca yasya guṇaḥ ucyate tat kṅitparatvena viśeṣayiṣyāmaḥ . \\
\noindent \textbf{(P\_1,1.5.1)} pugantalaghūpadhasya ca guṇaḥ ucyate tat ca atra kṅitparam . \\
\noindent \textbf{(P\_1,1.5.1)} pugantalaghūpadhasya iti na evam vijñāyate : pugantasya aṅgasya laghūpadhasya ca iti . \\
\noindent \textbf{(P\_1,1.5.1)} katham tarhi . \\
\noindent \textbf{(P\_1,1.5.1)} puki antaḥ pugantaḥ , laghvī upadhā laghūpadhā , pugantaḥ ca laghūpadhā ca pugantalaghūpadham , pugantalaghūpadhasya iti . \\
\noindent \textbf{(P\_1,1.5.1)} avaśyam ca etat evam vijñeyam . \\
\noindent \textbf{(P\_1,1.5.1)} aṅgaviśeṣaṇe hi sati iha prasajyeta : bhinatti chinatti iti . \\
\noindent \textbf{(P\_1,1.5.1)} roravītyartham ca . \\
\noindent \textbf{(P\_1,1.5.1)} tridhā baddhaḥ vṛṣabhaḥ roravīti iti . \\
\noindent \textbf{(P\_1,1.5.1)} yadi tannimittagrahaṇam kriyate śacaṅante doṣaḥ . \\
\noindent \textbf{(P\_1,1.5.1)} riyati piyati dhiyati . \\
\noindent \textbf{(P\_1,1.5.1)} prādudruvat prāsusruvat . \\
\noindent \textbf{(P\_1,1.5.1)} atra prāpnoti . \\
\noindent \textbf{(P\_1,1.5.1)} śacaṅantasya antaraṅgalakṣaṇatvāt . \\
\noindent \textbf{(P\_1,1.5.1)} antaraṅgalakṣaṇatvāt atra iyaṅuvaṅoḥ kṛtayoḥ anupadhātvāt guṇaḥ na bhaviṣyati . \\
\noindent \textbf{(P\_1,1.5.1)} evam kriyate ca idam tannimittagrahaṇam na ca kaḥ cit doṣaḥ bhavati . \\
\noindent \textbf{(P\_1,1.5.1)} imāni ca bhūyaḥ tannimittagrahaṇasya prayojanāni : hataḥ , hathaḥ , upoyate , auyata , lauyamāniḥ , pauyamāniḥ , nenikte iti . \\
\noindent \textbf{(P\_1,1.5.1)} na etāni santi prayojanāni . \\
\noindent \textbf{(P\_1,1.5.1)} iha tāvat hataḥ , hathaḥ iti . \\
\noindent \textbf{(P\_1,1.5.1)} prasaktasya anabhinirvṛttasya pratiṣedhena nivṛttiḥ śakyā kartum atra ca dhātūpadeśāvasthāyām eva akāraḥ . \\
\noindent \textbf{(P\_1,1.5.1)} iha ca upoyate , auyata , lauyamāniḥ , pauyamāniḥ iti . \\
\noindent \textbf{(P\_1,1.5.1)} bahiraṅge guṇavṛddhī antaraṅgaḥ pratiṣedhaḥ . \\
\noindent \textbf{(P\_1,1.5.1)} asiddham bahiraṅgam antaraṅge . \\
\noindent \textbf{(P\_1,1.5.1)} nenikte iti pareṇa rūpeṇa vyavahitatvāt na bhaviṣyati . \\
\noindent \textbf{(P\_1,1.5.2)} upadhārthena tāvat na arthaḥ . \\
\noindent \textbf{(P\_1,1.5.2)} dhātoḥ iti vartate . \\
\noindent \textbf{(P\_1,1.5.2)} dhātum kṅitparatvena viśeṣayiṣyāmaḥ . \\
\noindent \textbf{(P\_1,1.5.2)} yadi dhātuḥ viśeṣyate vikaraṇasya na prāpnoti : cinutaḥ , sunutaḥ , lunītaḥ , punītaḥ iti . \\
\noindent \textbf{(P\_1,1.5.2)} na eṣaḥ doṣaḥ . \\
\noindent \textbf{(P\_1,1.5.2)} vihitaviśeṣaṇam dhātugrahaṇam . \\
\noindent \textbf{(P\_1,1.5.2)} dhātoḥ yaḥ vihitaḥ iti . \\
\noindent \textbf{(P\_1,1.5.2)} dhātoḥ eva tarhi na prāpnoti . \\
\noindent \textbf{(P\_1,1.5.2)} na evam vijñāyate dhātoḥ vihitasya kṅiti iti . \\
\noindent \textbf{(P\_1,1.5.2)} katham tarhi . \\
\noindent \textbf{(P\_1,1.5.2)} dhātoḥ vihite kṅiti iti . \\
\noindent \textbf{(P\_1,1.5.2)} atha vā kāryakālam hi sañjñāparibhāṣam . \\
\noindent \textbf{(P\_1,1.5.2)} yatra kāryam tatra draṣṭavyam . \\
\noindent \textbf{(P\_1,1.5.2)} pugantalaghūpadhasya guṇaḥ bhavati iti upasthitam idam bhavati kṅiti na iti . \\
\noindent \textbf{(P\_1,1.5.2)} atha vā yat etasmin yoge kṅidgrahaṇam tad anavakāśam . \\
\noindent \textbf{(P\_1,1.5.2)} tasya anavakāśatvāt guṇavṛddhī na bhaviṣyataḥ . \\
\noindent \textbf{(P\_1,1.5.2)} atha vā ācāryapravṛttiḥ jñāpayati bhavati upadhālakṣaṇasya guṇasya pratiṣedhaḥ iti yat ayam trasigṛdhidhṛṣikṣipeḥ knuḥ ikaḥ jhal halantāt ca iti knusanau kitau karoti . \\
\noindent \textbf{(P\_1,1.5.2)} katham kṛtvā jñāpakam . \\
\noindent \textbf{(P\_1,1.5.2)} kitkaraṇe etat prayojanam guṇaḥ katham na syāt iti . \\
\noindent \textbf{(P\_1,1.5.2)} yadi ca atra guṇapratiṣedhaḥ na syāt kitkaraṇam anarthakam syāt . \\
\noindent \textbf{(P\_1,1.5.2)} paśyati tu ācāryaḥ bhavati upadhālakṣaṇasya guṇasya pratiṣedhaḥ iti . \\
\noindent \textbf{(P\_1,1.5.2)} tataḥ knusanau kitau karoti . \\
\noindent \textbf{(P\_1,1.5.2)} roravītyarthena api na arthaḥ . \\
\noindent \textbf{(P\_1,1.5.2)} kṅiti iti ucyate . \\
\noindent \textbf{(P\_1,1.5.2)} na ca atra kṅitam paśyāmaḥ . \\
\noindent \textbf{(P\_1,1.5.2)} pratyayalakṣaṇena prāpnoti . \\
\noindent \textbf{(P\_1,1.5.2)} na lumatā tasmin iti pratyayalakṣaṇapratiṣedhaḥ . \\
\noindent \textbf{(P\_1,1.5.2)} atha api na lumatā aṅgasya iti ucyate evam api na doṣaḥ . \\
\noindent \textbf{(P\_1,1.5.2)} katham . \\
\noindent \textbf{(P\_1,1.5.2)} na lumatā lupte aṅgādhikāraḥ pratinirdiśyate . \\
\noindent \textbf{(P\_1,1.5.2)} kim tarhi . \\
\noindent \textbf{(P\_1,1.5.2)} yaḥ asau lumatā lupyate tasmin yat aṅgam tasya yat kāryam tat na bhavati iti . \\
\noindent \textbf{(P\_1,1.5.2)} atha api āṅgādhikāraḥ pratinirdiśyate evam api na doṣaḥ . \\
\noindent \textbf{(P\_1,1.5.2)} katham . \\
\noindent \textbf{(P\_1,1.5.2)} kāryakālam hi sañjñāparibhāṣam yatra kāryam tatra draṣṭavyam . \\
\noindent \textbf{(P\_1,1.5.2)} sārvadhātukārdhadhātukayoḥ guṇaḥ bhavati iti upasthitam idam bhavati kṅiti na iti . \\
\noindent \textbf{(P\_1,1.5.2)} atha vā chāndasam etat . \\
\noindent \textbf{(P\_1,1.5.2)} dṛṣṭānuvidhiḥ ca chandasi bhavati . \\
\noindent \textbf{(P\_1,1.5.2)} atha vā bahiraṅgaḥ guṇaḥ antaraṅgaḥ pratiṣedhaḥ . \\
\noindent \textbf{(P\_1,1.5.2)} asiddham bahiraṅgam antaraṅge . \\
\noindent \textbf{(P\_1,1.5.2)} atha vā pūrvasmin yoge yad ārdhadhātukagrahaṇam tat anavakāśam . \\
\noindent \textbf{(P\_1,1.5.2)} tasya anavakāśatvāt guṇaḥ bhaviṣyati . \\
\noindent \textbf{(P\_1,1.5.3)} iha kasmāt na bhavati : laigavāyanaḥ , kāmayate . \\
\noindent \textbf{(P\_1,1.5.3)} taddhitakāmyoḥ ikprakaraṇāt . \\
\noindent \textbf{(P\_1,1.5.3)} iglakṣaṇayoḥ guṇavṛddhyoḥ pratiṣedhaḥ na ca ete iglakṣaṇe . \\
\noindent \textbf{(P\_1,1.5.3)} lakārasya ṅittvāt ādeśeṣu sthānivadbhāvaprasaṅgaḥ . \\
\noindent \textbf{(P\_1,1.5.3)} lakārasya ṅittvāt ādeśeṣu sthānivadbhāvaḥ prāpnoti : acinavam asunavam akaravam . \\
\noindent \textbf{(P\_1,1.5.3)} lakārasya ṅittvāt ādeśeṣu sthānivadbhāvaprasaṅgaḥ iti cet yāsuṭaḥ ṅidvacanāt siddham . yat ayam yāsuṭaḥ ṅidvacanam śāsti tat jñāpayati ācāryaḥ na ṅidādeśāḥ ṅitaḥ bhavanti iti . \\
\noindent \textbf{(P\_1,1.5.3)} yadi etat jñāpyate katham nityam ṅitaḥ itaḥ ca iti . \\
\noindent \textbf{(P\_1,1.5.3)} ṅitaḥ yat kāryam tat bhavati ṅiti yat kāryam tat na bhavati iti . \\
\noindent \textbf{(P\_1,1.5.3)} kim vaktavyam etat . \\
\noindent \textbf{(P\_1,1.5.3)} na hi . \\
\noindent \textbf{(P\_1,1.5.3)} katham anucyamānam gaṃsyate . \\
\noindent \textbf{(P\_1,1.5.3)} yāsuṭaḥ eva ṅidvacanāt . \\
\noindent \textbf{(P\_1,1.5.3)} aparyāptaḥ ca eva hi yāsuṭ samudāyasya ṅittve ṅitam ca enam karoti . \\
\noindent \textbf{(P\_1,1.5.3)} tasya etat prayojanam ṅitaḥ yat kāryam tat yathā syāt ṅiti yat kāryam tat mā bhūt iti . \\
\noindent \textbf{(P\_1,1.6)} kimartham idam udyate . \\
\noindent \textbf{(P\_1,1.6)} guṇavṛddhī mā bhūtām iti : ādīdhyanam ādīdhyakaḥ , āvevyanam āvevyakaḥ iti . \\
\noindent \textbf{(P\_1,1.6)} ayam yogaḥ śakyaḥ akartum . \\
\noindent \textbf{(P\_1,1.6)} katham . \\
\noindent \textbf{(P\_1,1.6)} dīdhīvevyoḥ chandoviṣayatvāt dṛṣṭānuvidhitvāt ca chandasaḥ adīdhet adīdhayuḥ iti ca guṇadarśanāt apratiṣedhaḥ . \\
\noindent \textbf{(P\_1,1.6)} dīdhīvevyoḥ chandoviṣayatvāt . \\
\noindent \textbf{(P\_1,1.6)} dīdhīvevyau chandoviṣayau . \\
\noindent \textbf{(P\_1,1.6)} dṛṣṭānuvidhitvāt ca chandasaḥ . \\
\noindent \textbf{(P\_1,1.6)} dṛṣṭānuvidhiḥ ca chandasi bhavati . \\
\noindent \textbf{(P\_1,1.6)} adīdhet , adīdhayuḥ iti ca guṇadarśanāt apratiṣedhaḥ . \\
\noindent \textbf{(P\_1,1.6)} anarthakaḥ pratiṣedhaḥ apratiṣedhaḥ . \\
\noindent \textbf{(P\_1,1.6)} prajapatiḥ vai yat kim cana manasā adīdhet . \\
\noindent \textbf{(P\_1,1.6)} hotraya vṛtaḥ kṛpayan adīdhet . \\
\noindent \textbf{(P\_1,1.6)} adīdhayuḥ dāśarājñe vṛtasaḥ . \\
\noindent \textbf{(P\_1,1.6)} bhavet idam yuktam udāharaṇam : adīdhet iti . \\
\noindent \textbf{(P\_1,1.6)} idam tu ayuktam : adīdhayuḥ iti . \\
\noindent \textbf{(P\_1,1.6)} ayam jusi guṇaḥ pratiṣedhaviṣaye [pratiṣedhaviṣayaḥ] ārabhyate . \\
\noindent \textbf{(P\_1,1.6)} saḥ yathā eva kṅiti na iti etam pratiṣedham bādhate evam imam api bādhate . \\
\noindent \textbf{(P\_1,1.6)} na eṣaḥ doṣaḥ . \\
\noindent \textbf{(P\_1,1.6)} jusi guṇaḥ pratiṣedhaviṣaye ārabhyamāṇaḥ tulyajātīyam pratiṣedham bādhate . \\
\noindent \textbf{(P\_1,1.6)} kaḥ ca tulyajātīyaḥ pratiṣedhaḥ . \\
\noindent \textbf{(P\_1,1.6)} yaḥ pratyayāśrayaḥ . \\
\noindent \textbf{(P\_1,1.6)} prakṛtyāśrayaḥ ca ayam . \\
\noindent \textbf{(P\_1,1.6)} atha vā yena na aprāpte tasya bādhanam bhavati . \\
\noindent \textbf{(P\_1,1.6)} na ca aprāpte kṅiti na iti etasmin pratiṣedhe jusi guṇaḥ ārabhyate . \\
\noindent \textbf{(P\_1,1.6)} asmin punaḥ prāpte ca aprāpte ca . \\
\noindent \textbf{(P\_1,1.6)} yadi tarhi ayam yogaḥ na ārabhyate katham dīdhyat iti . \\
\noindent \textbf{(P\_1,1.6)} dīdhyat iti śyanvyatyayena . \\
\noindent \textbf{(P\_1,1.6)} dīdhyat iti śyan eṣaḥ vyatyayena bhaviṣyati . \\
\noindent \textbf{(P\_1,1.6)} iṭaḥ ca api grahaṇam śakyam akartum . \\
\noindent \textbf{(P\_1,1.6)} katham akaṇiṣam araṇiṣam , kaṇitā śvaḥ , raṇitā śvaḥ iti . \\
\noindent \textbf{(P\_1,1.6)} ārdhadhātukasya iṭ valādeḥ iti atra iṭ iti vartamāne punaḥ iḍgrahaṇasya prayojanam iṭ eva yathā syāt yat anyat prāpnoti tat mā bhūt iti . \\
\noindent \textbf{(P\_1,1.6)} kim ca anyat prāpnoti . \\
\noindent \textbf{(P\_1,1.6)} guṇaḥ . \\
\noindent \textbf{(P\_1,1.6)} yadi niyamaḥ kriyate pipaṭhiṣateḥ apratyayaḥ pipaṭḥīḥ : dīrghatvam na prāpnoti . \\
\noindent \textbf{(P\_1,1.6)} na eṣaḥ doṣaḥ . \\
\noindent \textbf{(P\_1,1.6)} āṅgam yat kāryam tat niyamyate . \\
\noindent \textbf{(P\_1,1.6)} na ca etat āṅgam . \\
\noindent \textbf{(P\_1,1.6)} atha vā asiddham dīrghatvam . \\
\noindent \textbf{(P\_1,1.6)} tasya asiddhatvāt niyamaḥ na bhaviṣyati \\
\noindent \textbf{(P\_1,1.7.1)} anantarāḥ iti katham idam vijñāyate : avidyamānam antaram eṣām iti āhosvit avidyamānāḥ antarā eṣām iti . \\
\noindent \textbf{(P\_1,1.7.1)} kim ca ataḥ . \\
\noindent \textbf{(P\_1,1.7.1)} yadi vijñāyate avidyamānam antaram eṣām iti avagrahe saṃyogasañjñā na prāpnoti apsu iti ap-su iti . \\
\noindent \textbf{(P\_1,1.7.1)} vidyate hi atra antaram . \\
\noindent \textbf{(P\_1,1.7.1)} atha vijñāyate avidyamānāḥ antarā eṣām iti na doṣaḥ bhavati . \\
\noindent \textbf{(P\_1,1.7.1)} yathā na doṣaḥ tathā astu . \\
\noindent \textbf{(P\_1,1.7.1)} atha vā punaḥ astu avidyamānam antaram eṣām iti . \\
\noindent \textbf{(P\_1,1.7.1)} nanu ca uktam avagrahe saṃyogasañjñā na prāpnoti ap-su iti apsu iti . \\
\noindent \textbf{(P\_1,1.7.1)} vidyate hi atra antaram iti . \\
\noindent \textbf{(P\_1,1.7.1)} na eva doṣaḥ na prayojanam . \\
\noindent \textbf{(P\_1,1.7.2)} saṃyogasañjñāyām sahavacanam yathā anyatra . \\
\noindent \textbf{(P\_1,1.7.2)} saṃyogasañjñāyām sahavacanam kartavyam . \\
\noindent \textbf{(P\_1,1.7.2)} halaḥ anantarāḥ saṃyogaḥ saha iti vaktavyam . \\
\noindent \textbf{(P\_1,1.7.2)} kim prayojanam . \\
\noindent \textbf{(P\_1,1.7.2)} sahabhūtānām saṃyogasañjñā yathā syāt ekaikasya mā bhūt iti . \\
\noindent \textbf{(P\_1,1.7.2)} yathā anyatra . \\
\noindent \textbf{(P\_1,1.7.2)} tat yathā anyatra api yatra icchati sahabūtānām kāryam karoti tatra sahagrahaṇam . \\
\noindent \textbf{(P\_1,1.7.2)} tat yathā . \\
\noindent \textbf{(P\_1,1.7.2)} saha supā . \\
\noindent \textbf{(P\_1,1.7.2)} ubhe abhyastam saha iti . \\
\noindent \textbf{(P\_1,1.7.2)} kim ca syāt yadi ekaikasya halaḥ saṃyogasañjñā syāt . \\
\noindent \textbf{(P\_1,1.7.2)} iha niryāyāt , nirvāyāt , vā anyasya saṃyogādeḥ iti ettvam prasajyeta . \\
\noindent \textbf{(P\_1,1.7.2)} iha ca saṃhṛṣīṣṭa iti ṛtaḥ ca saṃyogādeḥ iti iṭ prasajyeta . \\
\noindent \textbf{(P\_1,1.7.2)} iha ca saṃhriyate iti guṇaḥ artisaṃyogādyoḥ iti guṇaḥ prasajyeta . \\
\noindent \textbf{(P\_1,1.7.2)} iha ca dṛṣat karoti samit karoti iti saṃyogāntasya lopaḥ prasajyeta . \\
\noindent \textbf{(P\_1,1.7.2)} iha ca śaktā vastā iti skoḥ saṃyogādyoḥ iti lopaḥ prasajyeta . \\
\noindent \textbf{(P\_1,1.7.2)} iha ca niryātaḥ , nirvātaḥ saṃyogādeḥ ātaḥ dhātoḥ yaṇvataḥ iti niṣṭhānatvam prasajyeta . \\
\noindent \textbf{(P\_1,1.7.2)} na eṣaḥ doṣaḥ . \\
\noindent \textbf{(P\_1,1.7.2)} yat tāvat ucyate iha tāvat niryāyāt , nirvāyāt vā anyasya saṃyogādeḥ iti ettvam prasajyeta iti . \\
\noindent \textbf{(P\_1,1.7.2)} na evam vijñāyate saṃyogaḥ ādiḥ yasya saḥ ayam saṃyogādiḥ , saṃyogādeḥ iti . \\
\noindent \textbf{(P\_1,1.7.2)} katham tarhi . \\
\noindent \textbf{(P\_1,1.7.2)} saṃyogau ādī yasya saḥ ayam saṃyogādiḥ , saṃyogādeḥ iti . \\
\noindent \textbf{(P\_1,1.7.2)} evam tāvat sarvam āṅgam parihṛtam . \\
\noindent \textbf{(P\_1,1.7.2)} yat api ucyate iha ca dṛṣat karoti samit karoti iti saṃyogāntasya lopaḥ prasajyeta iti . \\
\noindent \textbf{(P\_1,1.7.2)} na evam vijñāyate saṃyogaḥ antaḥ yasya tat idam saṃyogāntam , saṃyogāntasya iti . \\
\noindent \textbf{(P\_1,1.7.2)} katham tarhi . \\
\noindent \textbf{(P\_1,1.7.2)} saṃyogau antau asya tad idam saṃyogāntam , saṃyogāntasya iti . \\
\noindent \textbf{(P\_1,1.7.2)} yat api ucyate iha ca śaktā vastā iti skoḥ saṃyogādyoḥ iti lopaḥ prasajyeta iti . \\
\noindent \textbf{(P\_1,1.7.2)} na evam vijñāyate saṃyogau ādī saṃyodādī saṃyogādyoḥ iti . \\
\noindent \textbf{(P\_1,1.7.2)} katham tarhi . \\
\noindent \textbf{(P\_1,1.7.2)} saṃyogayoḥ ādī saṃyogādī saṃyogādyoḥ iti . \\
\noindent \textbf{(P\_1,1.7.2)} yat api ucyate iha ca niryātaḥ , nirvātaḥ saṃyogādeḥ ātaḥ dhātoḥ yaṇvataḥ iti niṣṭhānatvam prasajyeta iti . \\
\noindent \textbf{(P\_1,1.7.2)} na evam vijñāyate saṃyogaḥ ādiḥ yasya saḥ ayam saṃyogādiḥ , saṃyogādeḥ iti . \\
\noindent \textbf{(P\_1,1.7.2)} katham tarhi . \\
\noindent \textbf{(P\_1,1.7.2)} saṃyogau ādī yasya saḥ ayam saṃyogādiḥ , saṃyogādeḥ iti . \\
\noindent \textbf{(P\_1,1.7.2)} katham kṛtvā ekaikasya saṃyogasañjñā prāpnoti . \\
\noindent \textbf{(P\_1,1.7.2)} pratyekam vākyaparisamāptiḥ dṛṣṭā . \\
\noindent \textbf{(P\_1,1.7.2)} tat yathā . \\
\noindent \textbf{(P\_1,1.7.2)} vṛddhiguṇasañjñe pratyekam bhavataḥ . \\
\noindent \textbf{(P\_1,1.7.2)} nanu ca ayam api asti dṛṣṭāntaḥ : samudāye vākyaparisamāptiḥ iti . \\
\noindent \textbf{(P\_1,1.7.2)} tat yathā . \\
\noindent \textbf{(P\_1,1.7.2)} gargāḥ śatam daṇḍyantām iti . \\
\noindent \textbf{(P\_1,1.7.2)} arthinaḥ ca rājānaḥ hiraṇyena bhavanti na ca pratyekam daṇḍayanti . \\
\noindent \textbf{(P\_1,1.7.2)} sati etasmin dṛṣṭānte yadi tatra pratyekam iti ucyate iha api sahagrahaṇam kartavyam . \\
\noindent \textbf{(P\_1,1.7.2)} atha tatra antareṇa pratyekam iti vacanam pratyekam vṛddhiguṇasañjñe bhavataḥ iha api na arthaḥ sahagrahaṇena \\
\noindent \textbf{(P\_1,1.7.3)} atha yatra bahūnām ānantaryam kim tatra dvayoḥ dvayoḥ saṃyogasñjñā bhavati āhosvit aviśeṣeṇa . \\
\noindent \textbf{(P\_1,1.7.3)} kaḥ ca atra viśeṣaḥ . \\
\noindent \textbf{(P\_1,1.7.3)} samudāye saṃyogādilopaḥ masjeḥ . \\
\noindent \textbf{(P\_1,1.7.3)} samudāye saṃyogādilopaḥ masjeḥ na sidhyati . \\
\noindent \textbf{(P\_1,1.7.3)} maṅktā maṅktum . \\
\noindent \textbf{(P\_1,1.7.3)} iha ca nirgleyāt , nirglāyāt , nirmleyāt , nirmlāyāt : vā anyasya saṃyogādeḥ iti ettvam na prāpnoti . \\
\noindent \textbf{(P\_1,1.7.3)} iha ca saṃsvariṣīṣṭa iti ṛtaḥ ca saṃyogādeḥ iti iṭ na prāpnoti . \\
\noindent \textbf{(P\_1,1.7.3)} iha ca saṃsvaryate iti guṇaḥ artisaṃyogādyoḥ iti guṇaḥ na prāpnoti . \\
\noindent \textbf{(P\_1,1.7.3)} iha ca gomān karoti yavamān karoti iti saṃyogāntasya lopaḥ iti lopaḥ na prāpnoti . \\
\noindent \textbf{(P\_1,1.7.3)} iha ca nirglānaḥ , nirmlānaḥ iti saṃyogādeḥ ātaḥ dhātoḥ yaṇvataḥ iti niṣṭhānatvam na prāpnoti . \\
\noindent \textbf{(P\_1,1.7.3)} astu tarhi dvayoḥ dvayoḥ . \\
\noindent \textbf{(P\_1,1.7.3)} dvayoḥ haloḥ saṃyogaḥ iti cet dvirvacanam . dvayoḥ haloḥ saṃyogaḥ iti cet dvirvacanam na sidhyati . \\
\noindent \textbf{(P\_1,1.7.3)} indram icchati indrīyati . \\
\noindent \textbf{(P\_1,1.7.3)} indrīyateḥ san : indidrīyiṣati . \\
\noindent \textbf{(P\_1,1.7.3)} na ndrāḥ saṃyogādayaḥ iti dakārasya dvirvacanam na prāpnoti . \\
\noindent \textbf{(P\_1,1.7.3)} na vā ajvidheḥ . \\
\noindent \textbf{(P\_1,1.7.3)} na vā eṣaḥ doṣaḥ . \\
\noindent \textbf{(P\_1,1.7.3)} kim kāraṇam . \\
\noindent \textbf{(P\_1,1.7.3)} ajvidheḥ . \\
\noindent \textbf{(P\_1,1.7.3)} ndrāḥ saṃyogādayaḥ na dviḥ ucyante . \\
\noindent \textbf{(P\_1,1.7.3)} ajādeḥ iti vartate . \\
\noindent \textbf{(P\_1,1.7.3)} atha yadi eva bahūnām saṃyogasñjñā atha api dvayoḥ dvayoḥ kim gatam etat iyatā sūtreṇa āhosvit anyatarasmin pakṣe bhūyaḥ sūtram kartavyam . \\
\noindent \textbf{(P\_1,1.7.3)} gatam iti āha . \\
\noindent \textbf{(P\_1,1.7.3)} katham . \\
\noindent \textbf{(P\_1,1.7.3)} yada tāvat bahūnām saṃyogasñjñā tadā evam vigrahaḥ kariṣyate : avidyamānam antaram eṣām iti . \\
\noindent \textbf{(P\_1,1.7.3)} yadā dvayoḥ dvayoḥ tadā evam vigrahaḥ kariṣyate : avidyamānāḥ antarā eṣām iti . \\
\noindent \textbf{(P\_1,1.7.3)} dvayoḥ ca eva antarā kaḥ cit vidyate na vā . \\
\noindent \textbf{(P\_1,1.7.3)} evam api bahūnām eva prāpnoti . \\
\noindent \textbf{(P\_1,1.7.3)} yān hi bhavān ṣaṣṭhyā pratinirdiśati eteṣām anyena vyavāye na bhavitavyam . \\
\noindent \textbf{(P\_1,1.7.3)} astu tarhi samudāye sañjñā . \\
\noindent \textbf{(P\_1,1.7.3)} nanu ca uktam samudāye saṃyogādilopaḥ masjeḥ iti . \\
\noindent \textbf{(P\_1,1.7.3)} na eṣaḥ doṣaḥ . \\
\noindent \textbf{(P\_1,1.7.3)} vakṣyati etat . \\
\noindent \textbf{(P\_1,1.7.3)} antyāt pūrvaḥ masjeḥ mit anuṣaṅgasaṃyogādilopārtham iti . \\
\noindent \textbf{(P\_1,1.7.3)} atha vā aviśeṣeṇa saṃyogasañjñā vijñāsyate dvayoḥ api bahūnām api . \\
\noindent \textbf{(P\_1,1.7.3)} tatra dvayoḥ yā saṃyogasñjñā tadāśrayaḥ lopaḥ bhaviṣyati . \\
\noindent \textbf{(P\_1,1.7.3)} yat api ucyate iha ca nirgleyāt , nirglāyāt , nirmleyāt , nirmlāyāt : vā anyasya saṃyogādeḥ iti ettvam na prāpnoti iti aṅgena saṃyogādim viśeṣayiṣyāmaḥ . \\
\noindent \textbf{(P\_1,1.7.3)} aṅgasya saṃyogādeḥ iti . \\
\noindent \textbf{(P\_1,1.7.3)} evam tāvat sarvam āṅgam parihṛtam . \\
\noindent \textbf{(P\_1,1.7.3)} yat api ucyate iha ca gomān karoti yavamān karoti iti saṃyogāntasya lopaḥ iti lopaḥ na prāpnoti iti padena saṃyogāntam viśeṣayiṣyāmaḥ . \\
\noindent \textbf{(P\_1,1.7.3)} padasya saṃyogāntasya . \\
\noindent \textbf{(P\_1,1.7.3)} yat api ucyate iha ca nirglānaḥ , nirmlānaḥ iti saṃyogādeḥ ātaḥ dhātoḥ yaṇvataḥ iti niṣṭhānatvam na prāpnoti iti dhātuna saṃyogādim viśeṣayiṣyāmaḥ . \\
\noindent \textbf{(P\_1,1.7.3)} dhātoḥ saṃyogādeḥ iti \\
\noindent \textbf{(P\_1,1.7.4)} svarānantarhitavacanam . \\
\noindent \textbf{(P\_1,1.7.4)} svaraiḥ anantarhitāḥ halaḥ saṃyogasañjñāḥ bhavanti iti vaktavyam . \\
\noindent \textbf{(P\_1,1.7.4)} kim prayojanam . \\
\noindent \textbf{(P\_1,1.7.4)} vyavahitānām mā bhūt . \\
\noindent \textbf{(P\_1,1.7.4)} pacati panasam . \\
\noindent \textbf{(P\_1,1.7.4)} nanu ca anantarāḥ iti ucyate . \\
\noindent \textbf{(P\_1,1.7.4)} tena vyavahitānām na bhaviṣyati . \\
\noindent \textbf{(P\_1,1.7.4)} dṛṣṭam ānantaryam vyavahite . \\
\noindent \textbf{(P\_1,1.7.4)} vyavahite api anantaraśabdaḥ dṛśyate . \\
\noindent \textbf{(P\_1,1.7.4)} tat yathā : anantarau imau grāmau iti ucyate . \\
\noindent \textbf{(P\_1,1.7.4)} tayoḥ ca eva antarā nadyaḥ ca parvatāḥ ca bhavanti . \\
\noindent \textbf{(P\_1,1.7.4)} yadi tarhi vyavahite api anantaraśabdaḥ bhavati ānantaryavacanam idānīm kimartham syāt . \\
\noindent \textbf{(P\_1,1.7.4)} ānantaryavacanam kimartham iti cet ekapratiṣedhārtham . \\
\noindent \textbf{(P\_1,1.7.4)} ekasya halaḥ saṃyogasañjñā mā bhūt iti . \\
\noindent \textbf{(P\_1,1.7.4)} kim ca syāt yadi ekasya halaḥ saṃyogasañjñā syāt . \\
\noindent \textbf{(P\_1,1.7.4)} iyeṣa , uvoṣa . \\
\noindent \textbf{(P\_1,1.7.4)} ijādeḥ ca gurumataḥ anṛcchaḥ iti ām prasajyeta . \\
\noindent \textbf{(P\_1,1.7.4)} na vā atajjātīyavyavāyāt . \\
\noindent \textbf{(P\_1,1.7.4)} na vā eṣaḥ doṣaḥ . \\
\noindent \textbf{(P\_1,1.7.4)} kim kāraṇam . \\
\noindent \textbf{(P\_1,1.7.4)} atajjātīyasya vyavāyāt . \\
\noindent \textbf{(P\_1,1.7.4)} atajjātīyakam hi loke vyavadhāyakam bhavati . \\
\noindent \textbf{(P\_1,1.7.4)} katham punaḥ jñāyate atajjātīyakam loke vyavadhāyakam bhavati iti . \\
\noindent \textbf{(P\_1,1.7.4)} evam hi kam cit kaḥ cit pṛcchati . \\
\noindent \textbf{(P\_1,1.7.4)} anantare* ete brāhmaṇakule* iti . \\
\noindent \textbf{(P\_1,1.7.4)} saḥ āha . \\
\noindent \textbf{(P\_1,1.7.4)} na anantare . \\
\noindent \textbf{(P\_1,1.7.4)} vṛṣalakulam anayoḥ antarā iti . \\
\noindent \textbf{(P\_1,1.7.4)} kim punaḥ kāraṇam kva cit atajjātīyakam vyavadhāyakam bhavati kva cit na . \\
\noindent \textbf{(P\_1,1.7.4)} sarvatra eva hi atajjātīyakam vyavadhāyakam bhavati . \\
\noindent \textbf{(P\_1,1.7.4)} katham anantarau imau grāmau iti . \\
\noindent \textbf{(P\_1,1.7.4)} grāmaśabdaḥ ayam bahvarthaḥ . \\
\noindent \textbf{(P\_1,1.7.4)} asti eva śālāsamudāye vartate . \\
\noindent \textbf{(P\_1,1.7.4)} tat yathā grāmaḥ dagdhaḥ iti . \\
\noindent \textbf{(P\_1,1.7.4)} asti vāṭaparikṣepe vartate . \\
\noindent \textbf{(P\_1,1.7.4)} tat yathā grāmam praviṣṭaḥ . \\
\noindent \textbf{(P\_1,1.7.4)} asti manuṣyeṣu vartate . \\
\noindent \textbf{(P\_1,1.7.4)} tat yathā grāmaḥ gataḥ , grāmaḥ āgataḥ iti . \\
\noindent \textbf{(P\_1,1.7.4)} asti sāraṇyake sasīmake sasthaṇḍilake vartate . \\
\noindent \textbf{(P\_1,1.7.4)} tat yathā grāmaḥ labdhaḥ iti . \\
\noindent \textbf{(P\_1,1.7.4)} tat yaḥ sāraṇyake sasīmake sasthaṇḍilake vartate tam abhisamīkṣya etat prayujyate : anantarau imau grāmau iti . \\
\noindent \textbf{(P\_1,1.7.4)} sarvatra eva hi atajjātīyakam vyavadhāyakam bhavati . \\
\noindent \textbf{(P\_1,1.8.1)} kim idam mukhanāsikāvacanaḥ iti . \\
\noindent \textbf{(P\_1,1.8.1)} mukham ca nāsikā ca mukhanāsikam . \\
\noindent \textbf{(P\_1,1.8.1)} mukhanāsikam vacanam asya saḥ ayam mukhanāsikāvacanaḥ . \\
\noindent \textbf{(P\_1,1.8.1)} yadi evam mukhanāsikavacanaḥ iti prāpnoti . \\
\noindent \textbf{(P\_1,1.8.1)} nipātanāt dīrghatvam bhaviṣyati . \\
\noindent \textbf{(P\_1,1.8.1)} atha vā mukhanāsikam āvacanam asya saḥ ayam mukhanāsikāvacanaḥ . \\
\noindent \textbf{(P\_1,1.8.1)} kim idam āvacanam iti . \\
\noindent \textbf{(P\_1,1.8.1)} īṣadvacanam āvacanam . \\
\noindent \textbf{(P\_1,1.8.1)} kim cit mukhavacanam kim cit nāsikāvacanam . \\
\noindent \textbf{(P\_1,1.8.1)} mukhadvitīyā vā nāsikā vacanam asya saḥ ayam mukhanāsikāvacanaḥ . \\
\noindent \textbf{(P\_1,1.8.1)} mukhopasaṃhitā vā nāsikā vacanam asya saḥ ayam mukhanāsikāvacanaḥ . \\
\noindent \textbf{(P\_1,1.8.2)} atha mukhagrahaṇam kimartham . \\
\noindent \textbf{(P\_1,1.8.2)} nāsikāvacanaḥ anunāsikaḥ iti iyati ucyamāne yamānusvārāṇām eva prasajyeta . \\
\noindent \textbf{(P\_1,1.8.2)} mukhagrahaṇe punaḥ kriyamāṇe na doṣaḥ bhavati . \\
\noindent \textbf{(P\_1,1.8.2)} atha nāsikāgrahaṇam kimartham . \\
\noindent \textbf{(P\_1,1.8.2)} mukhavacanaḥ anunāsikaḥ iti iyati ucyamāne kacaṭatapānām eva prasajyeta . \\
\noindent \textbf{(P\_1,1.8.2)} nāsikāgrahaṇe punaḥ kriyamāṇe na doṣaḥ bhavati . \\
\noindent \textbf{(P\_1,1.8.2)} mukhagrahaṇam śakyam akartum . \\
\noindent \textbf{(P\_1,1.8.2)} kena idānīm ubhayavacanānām bhaviṣyati . \\
\noindent \textbf{(P\_1,1.8.2)} prāsādavāsinyāyena . \\
\noindent \textbf{(P\_1,1.8.2)} tat yathā . \\
\noindent \textbf{(P\_1,1.8.2)} ke cit prāsādavāsinaḥ ke cit bhūmivāsinaḥ ke cit ubhayavāsinaḥ . \\
\noindent \textbf{(P\_1,1.8.2)} ye prāsādavāsinaḥ gṛhyante te prāsādavāsigrahaṇena . \\
\noindent \textbf{(P\_1,1.8.2)} ye bhūmivāsinaḥ gṛhyante te bhūmivāsinyāyena . \\
\noindent \textbf{(P\_1,1.8.2)} ye ubhayavāsinaḥ gṛhyante te prāsādavāsigrahaṇena bhūmivāsinyāyena ca . \\
\noindent \textbf{(P\_1,1.8.2)} evam iha api ke cit mukhavacanāḥ ke cit nāsikāvacanāḥ ke cit ubhayavacanāḥ . \\
\noindent \textbf{(P\_1,1.8.2)} tatra ye mukhavacanāḥ gṛhyante te mukhagrahaṇena . \\
\noindent \textbf{(P\_1,1.8.2)} ye nāsikāvacanāḥ gṛhyante te nāsikāgrahaṇena . \\
\noindent \textbf{(P\_1,1.8.2)} ye ubhayavacanāḥ gṛhyante eva te mukhagrahaṇena nāsikāgrahaṇena ca. bhavet ubhayavacanānām siddham . \\
\noindent \textbf{(P\_1,1.8.2)} yamānusvārāṇām api prāpnoti . \\
\noindent \textbf{(P\_1,1.8.2)} na eva doṣaḥ na prayojanam . \\
\noindent \textbf{(P\_1,1.8.3)} itaretarāśrayam tu bhavati . \\
\noindent \textbf{(P\_1,1.8.3)} kā itaretarāśrayatā . \\
\noindent \textbf{(P\_1,1.8.3)} sataḥ anunāsikasya sañjñayā bhavitavyam sañjñayā ca nāma anunāsikaḥ bhāvyate . \\
\noindent \textbf{(P\_1,1.8.3)} tat itaretarāśrayam bhavati . \\
\noindent \textbf{(P\_1,1.8.3)} itaretarāśrayāṇi ca kāryāṇi na prakalpante . \\
\noindent \textbf{(P\_1,1.8.3)} anunāsikasañjñāyām itaretarāśraye uktam . \\
\noindent \textbf{(P\_1,1.8.3)} siddham tu nityaśabdatvāt iti . \\
\noindent \textbf{(P\_1,1.8.3)} nityāḥ śabdāḥ . \\
\noindent \textbf{(P\_1,1.8.3)} nityeṣu śabdeṣu sataḥ anunāsikasya sañjñā kriyate . \\
\noindent \textbf{(P\_1,1.8.3)} na sañjñayā anunāsikaḥ bhāvyate . \\
\noindent \textbf{(P\_1,1.8.3)} yadi tarhi nityāḥ śabdāḥ kimartham śāstram . \\
\noindent \textbf{(P\_1,1.8.3)} kimartham śāstram iti cet nivartakatvāt siddham . \\
\noindent \textbf{(P\_1,1.8.3)} nivartakam śāstram . \\
\noindent \textbf{(P\_1,1.8.3)} katham . \\
\noindent \textbf{(P\_1,1.8.3)} āṅ asmai aviśeṣeṇa upadiṣṭaḥ ananunāsikaḥ . \\
\noindent \textbf{(P\_1,1.8.3)} tasya sarvatra ananunāsikabuddhiḥ prasaktā . \\
\noindent \textbf{(P\_1,1.8.3)} tatra anena nivṛttiḥ kriyate . \\
\noindent \textbf{(P\_1,1.8.3)} chandasi aci parataḥ āṅaḥ ananunāsikasya prasaṅge anunāsikaḥ sādhuḥ bhavati iti . \\
\noindent \textbf{(P\_1,1.9.1)} tulayā sammitam tulyam . \\
\noindent \textbf{(P\_1,1.9.1)} āsyam ca prayatnaḥ ca āsyaprayatnam . \\
\noindent \textbf{(P\_1,1.9.1)} tulyāsyam tulyaprayatnam ca savarṇasañjñam bhavati . \\
\noindent \textbf{(P\_1,1.9.1)} kim punaḥ āsyam . \\
\noindent \textbf{(P\_1,1.9.1)} laukikam āsyam oṣṭhāt prabhṛti prāk kākalakāt . \\
\noindent \textbf{(P\_1,1.9.1)} katham punaḥ āsyam . \\
\noindent \textbf{(P\_1,1.9.1)} asyanti anena varṇān iti āsyam . \\
\noindent \textbf{(P\_1,1.9.1)} annam etat āsyandate iti vā āsyam . \\
\noindent \textbf{(P\_1,1.9.1)} atha kaḥ prayatnaḥ . \\
\noindent \textbf{(P\_1,1.9.1)} prayatanam prayatnaḥ . \\
\noindent \textbf{(P\_1,1.9.1)} prapūrvāt yatateḥ bhāvasādhanaḥ naṅpratyayaḥ . \\
\noindent \textbf{(P\_1,1.9.1)} yadi laukikam āsyam kim āsyopādāne prayojanam . \\
\noindent \textbf{(P\_1,1.9.1)} sarveṣām hi tat tulyam bhavati . \\
\noindent \textbf{(P\_1,1.9.1)} vakṣyati etat : prayatnaviśeṣaṇam āsyopādānam iti . \\
\noindent \textbf{(P\_1,1.9.2)} savarṇasañjñāyām bhinnadeśeṣu atiprasaṅgaḥ prayatnasāmānyāt . \\
\noindent \textbf{(P\_1,1.9.2)} savarṇasañjñāyām bhinnadeśeṣu atiprasaṅgaḥ bhavati jabagaḍadaśām . \\
\noindent \textbf{(P\_1,1.9.2)} kim kāraṇam . \\
\noindent \textbf{(P\_1,1.9.2)} prayatnasāmānyāt . \\
\noindent \textbf{(P\_1,1.9.2)} eteṣām hi samānaḥ prayatnaḥ . \\
\noindent \textbf{(P\_1,1.9.2)} siddham tu āsye tulyadeśaprayatnam savarṇam . \\
\noindent \textbf{(P\_1,1.9.2)} siddham etat. katham . \\
\noindent \textbf{(P\_1,1.9.2)} āsye yeṣām tulyaḥ deśaḥ yatnaḥ ca te savarṇasañjñāḥ bhavanti iti vaktavyam . \\
\noindent \textbf{(P\_1,1.9.2)} evam api kim āsyopādāne prayojanam . \\
\noindent \textbf{(P\_1,1.9.2)} sarveṣām hi tat tulyam . \\
\noindent \textbf{(P\_1,1.9.2)} prayatnaviśeṣaṇam āsyopādānam . \\
\noindent \textbf{(P\_1,1.9.2)} santi hi āsyāt bāhyāḥ prayatnāḥ . \\
\noindent \textbf{(P\_1,1.9.2)} te hāpitāḥ bhavanti . \\
\noindent \textbf{(P\_1,1.9.2)} teṣu satsu asatsu api savarṇasañjñā sidhyati . \\
\noindent \textbf{(P\_1,1.9.2)} ke punaḥ te . \\
\noindent \textbf{(P\_1,1.9.2)} vivārasaṃvārau śvāsanādau ghoṣavadaghoṣatā alpaprāṇatā mahāprāṇatā iti . \\
\noindent \textbf{(P\_1,1.9.2)} tatra vargāṇām prathamadvitīyāḥ vivṛtakaṇṭhāḥ śvāsānupradānāḥ aghoṣāḥ . \\
\noindent \textbf{(P\_1,1.9.2)} eke alpaprāṇāḥ apare mahāprāṇāḥ . \\
\noindent \textbf{(P\_1,1.9.2)} tṛtīyacaturthāḥ saṃvṛtakaṇṭhāḥ nādānupradānāḥ ghoṣavantaḥ . \\
\noindent \textbf{(P\_1,1.9.2)} eke alpaprāṇāḥ apare mahāprāṇāḥ . \\
\noindent \textbf{(P\_1,1.9.2)} yathā tṛtīyāḥ tathā pañcamāḥ ānunāsikyavarjam . \\
\noindent \textbf{(P\_1,1.9.2)} ānunāsikyam teṣām adhikaḥ guṇaḥ . \\
\noindent \textbf{(P\_1,1.9.2)} evam api avarṇasya savarṇasañjñā na prāpnoti . \\
\noindent \textbf{(P\_1,1.9.2)} kim kāraṇam . \\
\noindent \textbf{(P\_1,1.9.2)} bāhyam hi āsyāt sthānam avarṇasya . \\
\noindent \textbf{(P\_1,1.9.2)} sarvamukhasthānam avarṇam eke icchanti . \\
\noindent \textbf{(P\_1,1.9.2)} evam api vyapadeśaḥ na prakalpate : āsye yeṣām tulyaḥ deśaḥ iti . \\
\noindent \textbf{(P\_1,1.9.2)} vyapadeśivadbhāvena vyapadeśaḥ bhaviṣyati . \\
\noindent \textbf{(P\_1,1.9.2)} sidhyati . \\
\noindent \textbf{(P\_1,1.9.2)} sūtram tarhi bhidyate . \\
\noindent \textbf{(P\_1,1.9.2)} yathānyāsam eva astu . \\
\noindent \textbf{(P\_1,1.9.2)} nanu ca uktam savarṇasañjñayām bhinnadeśeṣu atiprasaṅgaḥ prayatnasāmānyāt iti . \\
\noindent \textbf{(P\_1,1.9.2)} na eṣaḥ doṣaḥ . \\
\noindent \textbf{(P\_1,1.9.2)} na hi laukikam āsyam . \\
\noindent \textbf{(P\_1,1.9.2)} kim tarhi . \\
\noindent \textbf{(P\_1,1.9.2)} taddhitāntam āsyam : āsye bhavam āsyam . \\
\noindent \textbf{(P\_1,1.9.2)} śarīrāvayavāt yat . \\
\noindent \textbf{(P\_1,1.9.2)} kim punaḥ āsye bhavam . \\
\noindent \textbf{(P\_1,1.9.2)} sthānam karaṇam ca . \\
\noindent \textbf{(P\_1,1.9.2)} evam api prayatnaḥ aviśeṣitaḥ bhavati . \\
\noindent \textbf{(P\_1,1.9.2)} prayatnaḥ ca viśeṣitaḥ . \\
\noindent \textbf{(P\_1,1.9.2)} katham . \\
\noindent \textbf{(P\_1,1.9.2)} na hi prayatanam prayatnaḥ . \\
\noindent \textbf{(P\_1,1.9.2)} kim tarhi . \\
\noindent \textbf{(P\_1,1.9.2)} prārambhaḥ yatnasya prayatnaḥ . \\
\noindent \textbf{(P\_1,1.9.2)} yadi prārambhaḥ yatnasya prayatnaḥ evam api avarṇasya eṅoḥ ca savarṇasañjñā prāpnoti . \\
\noindent \textbf{(P\_1,1.9.2)} praśliṣṭavarṇau etau . \\
\noindent \textbf{(P\_1,1.9.2)} avarṇasya tarhi aicoḥ ca savarṇasañjñā prāpnoti . \\
\noindent \textbf{(P\_1,1.9.2)} vivṛtatarāvarṇau etau . \\
\noindent \textbf{(P\_1,1.9.2)} etayoḥ eva tarhi mithaḥ savarṇasañjñā prāpnoti . \\
\noindent \textbf{(P\_1,1.9.2)} na etau tulyasthānau . \\
\noindent \textbf{(P\_1,1.9.2)} udāttādīnām tarhi savarṇasañjñā na prāpnoti . \\
\noindent \textbf{(P\_1,1.9.2)} abhedakāḥ udāttādayaḥ . \\
\noindent \textbf{(P\_1,1.9.2)} atha vā kim naḥ etena prārambhaḥ yatnasya prayatnaḥ iti .prayatanam eva prayatnaḥ . \\
\noindent \textbf{(P\_1,1.9.2)} tat eva ca taddhitāntam āsyam . \\
\noindent \textbf{(P\_1,1.9.2)} yat samānam tat āśrayiṣyāmaḥ . \\
\noindent \textbf{(P\_1,1.9.2)} kim sati bhede . \\
\noindent \textbf{(P\_1,1.9.2)} sati iti āha . \\
\noindent \textbf{(P\_1,1.9.2)} sati eva hi bhede savarṇasañjñayā bhavitavyam . \\
\noindent \textbf{(P\_1,1.9.2)} kutaḥ etat . \\
\noindent \textbf{(P\_1,1.9.2)} bhedādhiṣṭhānā hi savarṇasañjñā . \\
\noindent \textbf{(P\_1,1.9.2)} yadi hi yatra sarvam samānam tatra syāt savarṇasañjñāvacanam anarthakam syāt . \\
\noindent \textbf{(P\_1,1.9.2)} yadi tarhi sati bhede kim cit samānam iti kṛtva savarṇasañjñā bhaviṣyati śakārachakārayoḥ ṣakāraṭhakārahoḥ sakārathakārayoḥ savarṇasañjñā prāpnoti . \\
\noindent \textbf{(P\_1,1.9.2)} eteṣām hi sarvam anyat samānam karaṇavarjam . \\
\noindent \textbf{(P\_1,1.9.2)} evam tarhi prayatanam eva prayatnaḥ tat eva taddhitāntam āsyam na tu ayam dvandvaḥ : āsyam ca prayatnaḥ ca āsyaprayatnam iti . \\
\noindent \textbf{(P\_1,1.9.2)} kim tarhi . \\
\noindent \textbf{(P\_1,1.9.2)} tripadaḥ bahuvrīhiḥ : tulyaḥ āsye prayatnaḥ eṣām iti . \\
\noindent \textbf{(P\_1,1.9.2)} atha vā pūrvaḥ tatpuruṣaḥ tataḥ bahuvrīhiḥ : tulyaḥ āsye tulyāsyaḥ , tulyāsyaḥ prayatnaḥ eṣām iti . \\
\noindent \textbf{(P\_1,1.9.2)} atha vā paraḥ tatpuruṣaḥ tataḥ bahuvrīhiḥ : āsye yatnaḥ āsyayatnaḥ , tulyaḥ āsyayatnaḥ eṣām iti . \\
\noindent \textbf{(P\_1,1.9.3)} tasya . \\
\noindent \textbf{(P\_1,1.9.3)} tasya iti tu vaktavyam . \\
\noindent \textbf{(P\_1,1.9.3)} kim prayojanam . \\
\noindent \textbf{(P\_1,1.9.3)} yaḥ yasya tulyāsyaprayatnaḥ saḥ tasya savarṇasañjñaḥ yathā syāt . \\
\noindent \textbf{(P\_1,1.9.3)} anyasya tulyāsyaprayatnaḥ anyasya savarṇasañjñaḥ mā bhūt . \\
\noindent \textbf{(P\_1,1.9.3)} tasya avacanam vacanaprāmāṇyāt . \\
\noindent \textbf{(P\_1,1.9.3)} tasya iti na vaktavyam . \\
\noindent \textbf{(P\_1,1.9.3)} anyasya tulyāsyaprayatnaḥ anyasya savarṇasañjñaḥ kasmāt na bhavati . \\
\noindent \textbf{(P\_1,1.9.3)} vacanaprāmāṇyāt : savarṇasañjñāvacanasāmarthyāt . \\
\noindent \textbf{(P\_1,1.9.3)} yadi hi anyasya tulyāsyaprayatnaḥ saḥ anyasya savarṇasañjñaḥ syāt savarṇasañjñāvacanam anarthakam syāt . \\
\noindent \textbf{(P\_1,1.9.3)} sambandhiśabdaiḥ vā tulyam . \\
\noindent \textbf{(P\_1,1.9.3)} sambandhiśabdaiḥ vā punaḥ tulyam etat . \\
\noindent \textbf{(P\_1,1.9.3)} tat yathā sambandhiśabdāḥ : mātari vartitavyam , pitari śuśrūṣitavyam iti . \\
\noindent \textbf{(P\_1,1.9.3)} na ca ucyate svasyām mātari svasmin vā pitari iti sambandhāt ca etat gamyate yā yasya mātā yaḥ ca yasya pitā iti . \\
\noindent \textbf{(P\_1,1.9.3)} evam iha api tulyāsyaprayatnam savarṇam iti atra sambandiśabdau etau . \\
\noindent \textbf{(P\_1,1.9.3)} tatra sambandhāt etat gantavyam : yat prati yat tulyāsyaprayatnam tat prati tat savarṇasañjñam bhavati iti . \\
\noindent \textbf{(P\_1,1.9.4)} ṛkāraḷkārayoḥ savarṇavidhiḥ . \\
\noindent \textbf{(P\_1,1.9.4)} ṛkāraḷkārayoḥ savarṇasañjñā vidheyā . \\
\noindent \textbf{(P\_1,1.9.4)} hotṛ , ḷkāraḥ , hotṝḷkāraḥ . \\
\noindent \textbf{(P\_1,1.9.4)} kim prayojanam . \\
\noindent \textbf{(P\_1,1.9.4)} akaḥ savarṇe dīrghaḥ iti dīrghatvam yathā syāt . \\
\noindent \textbf{(P\_1,1.9.4)} na etat asti prayojanam . \\
\noindent \textbf{(P\_1,1.9.4)} vakṣyati etat . \\
\noindent \textbf{(P\_1,1.9.4)} savarṇadīrghatve ṛti , rṛvāvacanam ḷti , lḷvāvacanam iti . \\
\noindent \textbf{(P\_1,1.9.4)} tat savarṇe yathā syāt . \\
\noindent \textbf{(P\_1,1.9.4)} iha mā bhūt : dadhi , ḷkāraḥ , madhu , ḷkāraḥ iti . \\
\noindent \textbf{(P\_1,1.9.4)} yat etat savarṇadīrghatve ṛti iti etat ṛtaḥ iti vakṣyāmi . \\
\noindent \textbf{(P\_1,1.9.4)} tataḥ ḷti . \\
\noindent \textbf{(P\_1,1.9.4)} ḷti ca vā lḷ bhavati . \\
\noindent \textbf{(P\_1,1.9.4)} ṛtaḥ iti eva . \\
\noindent \textbf{(P\_1,1.9.4)} tat na vaktavyam bhavati . \\
\noindent \textbf{(P\_1,1.9.4)} avaśyam tat vaktavyam . \\
\noindent \textbf{(P\_1,1.9.4)} ūkālaḥ ac hrasvardīrghaplutasañjñaḥ bhavati iti ucyate . \\
\noindent \textbf{(P\_1,1.9.4)} na ca rṛkāraḥ lḷkāraḥ vā ac asti . \\
\noindent \textbf{(P\_1,1.9.4)} rṛkārasya , lḷkārasya ca actvam vakṣyāmi . \\
\noindent \textbf{(P\_1,1.9.4)} tat ca avaśyam vaktavyam plutaḥ yathā syāt : hotṛ , ṛkāraḥ hotṝkāraḥ , hotṛ3kāraḥ , hotṛ , ḷkāraḥ , hotḷkāraḥ , hotḷ3kāraḥ . \\
\noindent \textbf{(P\_1,1.9.4)} kim punaḥ atra jyāyaḥ . \\
\noindent \textbf{(P\_1,1.9.4)} savarṇasañjñāvacanam eva jyāyaḥ . \\
\noindent \textbf{(P\_1,1.9.4)} dīrghatvam ca eva hi siddham bhavati. api ca ṛkāragrahaṇe ḷkāragrahaṇam sannihitam bhavati . \\
\noindent \textbf{(P\_1,1.9.4)} yathā iha bhavati : ṛti akaḥ: khaṭva ṛśyaḥ , māla ṛśyaḥ idam api saṅgṛhītam bahavati : khaṭva , ḷkāraḥ, māla , ḷkāraḥ iti . \\
\noindent \textbf{(P\_1,1.9.4)} vā supi āpiśaleḥ : uparkārīyati , upārkārīyati , idam api siddham bhavati : upalkārīyati, upālkārīyati iti . \\
\noindent \textbf{(P\_1,1.9.4)} yadi tarhi ṛkāragrahaṇe ḷkāragrahaṇam sannihitam bhavati uḥ aṇ raparaḥ , ḷkārasya api raparatvam prāpnoti . \\
\noindent \textbf{(P\_1,1.9.4)} ḷkārasya laparatvam vakṣyāmi . \\
\noindent \textbf{(P\_1,1.9.4)} tat ca avaśyam vaktavyam asatyām savarṇasañjñāyām vidhyartham . \\
\noindent \textbf{(P\_1,1.9.4)} tat eva satyām rephabādhanārtham bhaviṣyati . \\
\noindent \textbf{(P\_1,1.9.4)} iha tarhi raṣābhyām naḥ ṇaḥ samānapade iti ṛkāragrahaṇam coditam mātṝṇām , pitṝṇām iti evamartham . \\
\noindent \textbf{(P\_1,1.9.4)} tat iha api prāpnoti : kḷpyamānam paśya iti . \\
\noindent \textbf{(P\_1,1.9.4)} atha asatyām api savarṇasañjñāyām iha kasmāt na bhavati : prakḷpyamānam paśya iti . \\
\noindent \textbf{(P\_1,1.9.4)} cuṭutulaśarvyavāye na iti vakṣyāmi . \\
\noindent \textbf{(P\_1,1.9.4)} aparaḥ āha : tribhiḥ ca madhyamaiḥ vargaiḥ laśasaiḥ ca vyavāye na iti vakṣyāmi iti . \\
\noindent \textbf{(P\_1,1.9.4)} varṇaikadeśāḥ ca varṇagrahaṇena gṛhyante iti yaḥ asau ḷkāre lakāraḥ tadāśrayaḥ pratiṣedhaḥ bhaviṣyati . \\
\noindent \textbf{(P\_1,1.9.4)} yadi evam na arthaḥ raṣābhyām ṇatve ṛkāragrahaṇena . \\
\noindent \textbf{(P\_1,1.9.4)} varṇaikadeśāḥ ca varṇagrahaṇena gṛhyante iti yaḥ asau ṛkāre rephaḥ tadāśrayam ṇatvam bhaviṣyati . \\
\noindent \textbf{(P\_1,1.10)} ajjhaloḥ pratiṣedhe śakārapratiṣedhaḥ ajjhaltvāt . \\
\noindent \textbf{(P\_1,1.10)} ajjhaloḥ pratiṣedhe śakārasya śakāreṇa savarṇasañjñāyāḥ pratiṣedhaḥ prāpnoti . \\
\noindent \textbf{(P\_1,1.10)} kim kāraṇam . \\
\noindent \textbf{(P\_1,1.10)} ajjhaltvāt . \\
\noindent \textbf{(P\_1,1.10)} ac ca eva hi śakāraḥ hal ca . \\
\noindent \textbf{(P\_1,1.10)} katham tāvat actvam . \\
\noindent \textbf{(P\_1,1.10)} ikāraḥ savarṇagrahaṇena śakāram api gṛhṇāti iti actvam . \\
\noindent \textbf{(P\_1,1.10)} halṣu upadeśāt haltvam . \\
\noindent \textbf{(P\_1,1.10)} tatra kaḥ doṣaḥ . \\
\noindent \textbf{(P\_1,1.10)} tatra savarṇalope doṣaḥ . \\
\noindent \textbf{(P\_1,1.10)} tatra savarṇalope doṣaḥ bhavati . \\
\noindent \textbf{(P\_1,1.10)} paraśśatāni kāryāṇi . \\
\noindent \textbf{(P\_1,1.10)} jharaḥ jhari savarṇe iti lopaḥ na prāpnoti . \\
\noindent \textbf{(P\_1,1.10)} siddham anactvāt . \\
\noindent \textbf{(P\_1,1.10)} siddham etat . \\
\noindent \textbf{(P\_1,1.10)} katham . \\
\noindent \textbf{(P\_1,1.10)} anactvāt . \\
\noindent \textbf{(P\_1,1.10)} katham anactvam . \\
\noindent \textbf{(P\_1,1.10)} spṛṣṭam sparśānām karaṇam . \\
\noindent \textbf{(P\_1,1.10)} īṣatspṛṣṭam antaḥsthānām . \\
\noindent \textbf{(P\_1,1.10)} vivṛtam ūṣmaṇām . \\
\noindent \textbf{(P\_1,1.10)} īṣat iti anuvartate . \\
\noindent \textbf{(P\_1,1.10)} svarāṇām vivṛtam . \\
\noindent \textbf{(P\_1,1.10)} īṣat iti nivṛttam . \\
\noindent \textbf{(P\_1,1.10)} vākyāparisamāpteḥ vā . \\
\noindent \textbf{(P\_1,1.10)} vākyāparisamāpteḥ vā siddham etat . \\
\noindent \textbf{(P\_1,1.10)} kim idam vākyāparisamāpteḥ iti . \\
\noindent \textbf{(P\_1,1.10)} varṇānām upadeśaḥ tāvat . \\
\noindent \textbf{(P\_1,1.10)} upadeśottarakālā itsañjñā . \\
\noindent \textbf{(P\_1,1.10)} itsañjñottarakālaḥ ādiḥ antyena saha itā iti pratyāhāraḥ . \\
\noindent \textbf{(P\_1,1.10)} pratyāhārottarakālā savarṇasañjñā . \\
\noindent \textbf{(P\_1,1.10)} savarṇasañjñottarakālam aṇ udit savarṇasya ca apratyayaḥ iti savarṇagrahaṇam . \\
\noindent \textbf{(P\_1,1.10)} etena sarveṇa samuditena vākyena anyatra savarṇānām grahaṇam bhavati . \\
\noindent \textbf{(P\_1,1.10)} ca ca atra ikāraḥ śakāram gṛhṇāti . \\
\noindent \textbf{(P\_1,1.10)} yathā eva tarhi ikāraḥ śakāram na gṛhṇāti evam īkāram api na gṛhṇīyāt . \\
\noindent \textbf{(P\_1,1.10)} tatra kaḥ doṣaḥ . \\
\noindent \textbf{(P\_1,1.10)} kumārī , īhate kumārīhate . \\
\noindent \textbf{(P\_1,1.10)} akaḥ savarṇadīrghatvam na prāpnoti . \\
\noindent \textbf{(P\_1,1.10)} na eṣaḥ doṣaḥ . \\
\noindent \textbf{(P\_1,1.10)} yat etat akaḥ savarṇe dīrghaḥ iti pratyāhāragrahaṇam tata ikāraḥ īkāram gṛhṇāti . \\
\noindent \textbf{(P\_1,1.10)} śakāram na gṛhṇāti . \\
\noindent \textbf{(P\_1,1.10)} aparaḥ āha : ajjhaloḥ pratiṣedhe śakārapratiṣedhaḥ ajhaltvāt . \\
\noindent \textbf{(P\_1,1.10)} ajjhaloḥ pratiṣedhe śakārasya śakāreṇa savarṇasañjñāyāḥ pratiṣedhaḥ prāpnoti . \\
\noindent \textbf{(P\_1,1.10)} kim kāraṇam . \\
\noindent \textbf{(P\_1,1.10)} ajjhaltvāt . \\
\noindent \textbf{(P\_1,1.10)} ac ca eva śakāraḥ hal ca . \\
\noindent \textbf{(P\_1,1.10)} katham tāvat actvam . \\
\noindent \textbf{(P\_1,1.10)} ikāraḥ savarṇagrahaṇena śakāram api gṛhṇāti iti actvam . \\
\noindent \textbf{(P\_1,1.10)} halṣu upadeśāt haltvam . \\
\noindent \textbf{(P\_1,1.10)} tatra kaḥ doṣaḥ . \\
\noindent \textbf{(P\_1,1.10)} tatra savarṇalope doṣaḥ . \\
\noindent \textbf{(P\_1,1.10)} tatra savarṇalope doṣaḥ bhavati . \\
\noindent \textbf{(P\_1,1.10)} paraśśatāni kāryāṇi . \\
\noindent \textbf{(P\_1,1.10)} jharaḥ jhari savarṇe iti lopaḥ na prāpnoti . \\
\noindent \textbf{(P\_1,1.10)} siddham anactvāt . \\
\noindent \textbf{(P\_1,1.10)} siddham etat . \\
\noindent \textbf{(P\_1,1.10)} katham . \\
\noindent \textbf{(P\_1,1.10)} anactvāt . \\
\noindent \textbf{(P\_1,1.10)} katham anactvam . \\
\noindent \textbf{(P\_1,1.10)} vākyāparisamāpteḥ vā . \\
\noindent \textbf{(P\_1,1.10)} uktā vākyāparisamāptiḥ . \\
\noindent \textbf{(P\_1,1.10)} asmin pakṣe vā iti etat asamarthitam bhavati . \\
\noindent \textbf{(P\_1,1.10)} etat ca samarthitam . \\
\noindent \textbf{(P\_1,1.10)} katham . \\
\noindent \textbf{(P\_1,1.10)} astu vā śakārasya śakāreṇa savarṇasañjñā mā vā bhūt . \\
\noindent \textbf{(P\_1,1.10)} nanu ca uktam : paraśśatāni kāryāṇi . \\
\noindent \textbf{(P\_1,1.10)} jharaḥ jhari savarṇe iti lopaḥ na prāpnoti iti . \\
\noindent \textbf{(P\_1,1.10)} mā bhūt lopaḥ . \\
\noindent \textbf{(P\_1,1.10)} nanu ca bhedaḥ bhavati . \\
\noindent \textbf{(P\_1,1.10)} sati lope dviśakāram asati lope triśakāram . \\
\noindent \textbf{(P\_1,1.10)} na asti bhedaḥ . \\
\noindent \textbf{(P\_1,1.10)} asati api lope dviśakāram eva . \\
\noindent \textbf{(P\_1,1.10)} katham . \\
\noindent \textbf{(P\_1,1.10)} vibhāṣā dvirvacanam . \\
\noindent \textbf{(P\_1,1.10)} evam api bhedaḥ . \\
\noindent \textbf{(P\_1,1.10)} asati lope kadā cit dviśakāram kadā cit triśakāram sati lope dviśakāram eva . \\
\noindent \textbf{(P\_1,1.10)} saḥ eṣaḥ katham bhedaḥ na syāt . \\
\noindent \textbf{(P\_1,1.10)} yadi nityaḥ lopaḥ syāt . \\
\noindent \textbf{(P\_1,1.10)} vibhāṣā tu saḥ lopaḥ . \\
\noindent \textbf{(P\_1,1.10)} yathā abhedaḥ tathā astu . \\
\noindent \textbf{(P\_1,1.11.1)} kimartham īdādīnām taparāṇām pragṛhyasañjñā ucyate . \\
\noindent \textbf{(P\_1,1.11.1)} taparaḥ tatkālasya iti tatkālānām savarṇānām grahaṇam yathā syāt . \\
\noindent \textbf{(P\_1,1.11.1)} keṣām . \\
\noindent \textbf{(P\_1,1.11.1)} udāttānudāttasvaritānām . \\
\noindent \textbf{(P\_1,1.11.1)} asti prayojanam etat . \\
\noindent \textbf{(P\_1,1.11.1)} kim tarhi iti . \\
\noindent \textbf{(P\_1,1.11.1)} plutānām tu pragṛhyasañjñā na prāpnoti . \\
\noindent \textbf{(P\_1,1.11.1)} kim kāraṇam . \\
\noindent \textbf{(P\_1,1.11.1)} atatkālatvāt . \\
\noindent \textbf{(P\_1,1.11.1)} na hi plutāḥ tatkālāḥ . \\
\noindent \textbf{(P\_1,1.11.1)} asiddhaḥ plutaḥ . \\
\noindent \textbf{(P\_1,1.11.1)} tasyāsiddhatvāt tatkālāḥ eva bhavanti . \\
\noindent \textbf{(P\_1,1.11.1)} siddhaḥ plutaḥ svarasandhiṣu . \\
\noindent \textbf{(P\_1,1.11.1)} katham jñāyate siddhaḥ plutaḥ svarasandhiṣu iti . \\
\noindent \textbf{(P\_1,1.11.1)} yat ayam plutapragṛhyāḥ aci iti plutasya prakṛtibhāvam śāsti . \\
\noindent \textbf{(P\_1,1.11.1)} katham kṛtvā jñāpakam . \\
\noindent \textbf{(P\_1,1.11.1)} sataḥ hi kāryiṇaḥ kāryeṇa bhavitavyam . \\
\noindent \textbf{(P\_1,1.11.1)} kim etasya jñāpane prayojanam . \\
\noindent \textbf{(P\_1,1.11.1)} aplutāt aplute iti etat na vaktavyam bhavati . \\
\noindent \textbf{(P\_1,1.11.1)} kim ataḥ yat siddhaḥ plutaḥ svarasandhiṣu . \\
\noindent \textbf{(P\_1,1.11.1)} sañjñāvidhau asiddhaḥ . \\
\noindent \textbf{(P\_1,1.11.1)} tasya asiddhatvāt tatkālāḥ eva bhavanti . \\
\noindent \textbf{(P\_1,1.11.1)} sañjñāvidhau ca siddhaḥ . \\
\noindent \textbf{(P\_1,1.11.1)} katham . \\
\noindent \textbf{(P\_1,1.11.1)} kāryakālam sañjñāparibhāṣam . \\
\noindent \textbf{(P\_1,1.11.1)} yatra kāryam tatra upasthitam draṣṭavyam . \\
\noindent \textbf{(P\_1,1.11.1)} pragṛhyaḥ prakṛtyā iti upasthitam idam bhavati īdūdet dvivacanam pragṛhyam iti . \\
\noindent \textbf{(P\_1,1.11.1)} kim punaḥ plutasya pragṛhyasañjñāvacane prayojanam . \\
\noindent \textbf{(P\_1,1.11.1)} pragṛhyāśrayaḥ prakṛtibhāvaḥ yathā syāt . \\
\noindent \textbf{(P\_1,1.11.1)} mā bhūt evam . \\
\noindent \textbf{(P\_1,1.11.1)} plutaḥ prakṛtyā iti evam bhaviṣyati . \\
\noindent \textbf{(P\_1,1.11.1)} na evam śakyam . \\
\noindent \textbf{(P\_1,1.11.1)} upasthite hi doṣaḥ syāt . \\
\noindent \textbf{(P\_1,1.11.1)} aplutavat upasthithe iti atra paṭhiṣyati hi ācāryaḥ : vadvacanam plutakāryapratiṣedhārtham , plutapratiṣedhe hi pragṛhyplutapratiṣedhaprasaṅgaḥ anyena vihitatvāt iti . \\
\noindent \textbf{(P\_1,1.11.1)} tasmāt plutasya pragṛhyasañjñā eṣitavyā pragṛhyāśrayaḥ prakṛtibhāvaḥ yathā syāt . \\
\noindent \textbf{(P\_1,1.11.1)} yadi punaḥ dīrghāṇām ataparāṇām pragṛhyasañjñā ucyeta . \\
\noindent \textbf{(P\_1,1.11.1)} evam api ekāraḥ eva ekaḥ savarṇān gṛhṇīyāt . \\
\noindent \textbf{(P\_1,1.11.1)} īkārokārau na gṛhṇīyātām . \\
\noindent \textbf{(P\_1,1.11.1)} kim kāraṇam . \\
\noindent \textbf{(P\_1,1.11.1)} anaṇtvāt . \\
\noindent \textbf{(P\_1,1.11.1)} yadi punaḥ hrasvānām ataparāṇām pragṛhyasañjñā ucyeta . \\
\noindent \textbf{(P\_1,1.11.1)} na evam śakyam . \\
\noindent \textbf{(P\_1,1.11.1)} iha api prasajyeta : akurvahi , atra akurvahi atra iti . \\
\noindent \textbf{(P\_1,1.11.1)} tasmāt dīrghāṇām eva taparāṇām pragṛhyasañjñā vaktavyā . \\
\noindent \textbf{(P\_1,1.11.1)} dīrghāṇām ca ucyamānā plutānām na prāpnoti . \\
\noindent \textbf{(P\_1,1.11.1)} evam tarhi kim naḥ etena yatnena yat siddhaḥ plutaḥ svarasandhiṣu iti . \\
\noindent \textbf{(P\_1,1.11.1)} asiddhaḥ plutaḥ . \\
\noindent \textbf{(P\_1,1.11.1)} tasya asiddhatvāt tatkālāḥ eva bhavanti iti . \\
\noindent \textbf{(P\_1,1.11.1)} katham yat tat jñāpakam uktam plutapragṛhyāḥ aci iti. plutabhāvī prakṛtyā iti evam etat vijñāyate . \\
\noindent \textbf{(P\_1,1.11.1)} katham yat tat prayojanam uktam . \\
\noindent \textbf{(P\_1,1.11.1)} kriyate tat nyāse eva aplutāt aplute iti . \\
\noindent \textbf{(P\_1,1.11.1)} evam api yat siddhe pragṛhyakāryam tat plutasya na prāpnoti . \\
\noindent \textbf{(P\_1,1.11.1)} aṇaḥ apragṛhyasya anunāsikaḥ iti . \\
\noindent \textbf{(P\_1,1.11.1)} evam tarhi kim naḥ etena kāryakālam sañjñāparibhāṣam iti . \\
\noindent \textbf{(P\_1,1.11.1)} yathoddeśam eva sañjñāparibhāṣam . \\
\noindent \textbf{(P\_1,1.11.1)} tatra ca asau asiddhaḥ . \\
\noindent \textbf{(P\_1,1.11.1)} tasyāsiddhatvāt tatkālāḥ eva bhavanti . \\
\noindent \textbf{(P\_1,1.11.2)} katham punaḥ idam vijñāyate : īdādayaḥ yat dvivacanam iti āhosvit īdādyantam yat dvivacanam iti . \\
\noindent \textbf{(P\_1,1.11.2)} kaḥ ca atra viśeṣaḥ . \\
\noindent \textbf{(P\_1,1.11.2)} īdādayaḥ dvivacanam pragṛhyāḥ iti cet antyasya vidhiḥ . \\
\noindent \textbf{(P\_1,1.11.2)} īdādayaḥ dvivacanam pragṛhyāḥ iti cet antyasya pragṛhyasañjñā vidheyā . \\
\noindent \textbf{(P\_1,1.11.2)} pacete* iti , pacethe* iti . \\
\noindent \textbf{(P\_1,1.11.2)} vacanāt bhaviṣyati . \\
\noindent \textbf{(P\_1,1.11.2)} asti vacane prayojanam . \\
\noindent \textbf{(P\_1,1.11.2)} kim . \\
\noindent \textbf{(P\_1,1.11.2)} khaṭve* iti , māle* iti . \\
\noindent \textbf{(P\_1,1.11.2)} astu tarhi īdādyantam yat dvivacanam iti . \\
\noindent \textbf{(P\_1,1.11.2)} īdādyantam iti cet ekasya vidhiḥ . \\
\noindent \textbf{(P\_1,1.11.2)} īdādyantam iti cet ekasya pragṛhyasañjñā vidheyā . \\
\noindent \textbf{(P\_1,1.11.2)} khaṭve* iti , māle* iti . \\
\noindent \textbf{(P\_1,1.11.2)} na vā ādyantatvāt . \\
\noindent \textbf{(P\_1,1.11.2)} na vā eṣaḥ doṣaḥ . \\
\noindent \textbf{(P\_1,1.11.2)} kim kāraṇam . \\
\noindent \textbf{(P\_1,1.11.2)} ādyantatvāt . \\
\noindent \textbf{(P\_1,1.11.2)} ādyantavat ekasmin iti ekasya api bhaviṣyati . \\
\noindent \textbf{(P\_1,1.11.2)} atha vā evam vakṣyāmi : īdādyantam yat dvivacanāntam iti . \\
\noindent \textbf{(P\_1,1.11.2)} īdādyantam dvivacanāntam iti cet luki pratiṣedhaḥ . \\
\noindent \textbf{(P\_1,1.11.2)} īdādyantam dvivacanāntam iti cet luki pratiṣedhaḥ vaktavyaḥ . \\
\noindent \textbf{(P\_1,1.11.2)} kumāryoḥ agāram , kumāryagāram vadhvoḥ agāram , vadhvagāram . \\
\noindent \textbf{(P\_1,1.11.2)} etat hi īdādyantam ca śrūyate dvivacanāntam ca bhavati pratyayalakṣaṇena . \\
\noindent \textbf{(P\_1,1.11.2)} saptamyām arthagrahaṇam jñāpakam pratyayalakṣaṇapratiṣedhasya . \\
\noindent \textbf{(P\_1,1.11.2)} yat ayam īdūtau ca saptamyarthe iti arthagrahaṇam karoti tat jñāpayati ācāryaḥ na pragṛhyasañjñāyām pratyayalakṣaṇam bhavati iti . \\
\noindent \textbf{(P\_1,1.11.2)} tat tarhi jñāpkārtham arthagrahaṇam kartavyam . \\
\noindent \textbf{(P\_1,1.11.2)} na kartavyam . \\
\noindent \textbf{(P\_1,1.11.2)} īdādibhiḥ dvivacanam viśeṣayiṣyāmaḥ īdādiviśiṣṭena ca dvivacanena tadantavidhiḥ bhaviṣyati . \\
\noindent \textbf{(P\_1,1.11.2)} īdādyantam yat dvivacanam tadantam īdādyantam iti . \\
\noindent \textbf{(P\_1,1.11.2)} evam api aśukle vastre śukle sampadyetām , śuklī āstām vastre* iti atra prāpnoti . \\
\noindent \textbf{(P\_1,1.11.2)} atra hi īdādi dvivacanam tadantam ca bhavati pratyayalakṣaṇena . \\
\noindent \textbf{(P\_1,1.11.2)} atra api akṛte śībhāve luk bhaviṣyati . \\
\noindent \textbf{(P\_1,1.11.2)} idam iha sampradhāryam . \\
\noindent \textbf{(P\_1,1.11.2)} luk kriyatām śībhāvaḥ iti kim atra kartavyam . \\
\noindent \textbf{(P\_1,1.11.2)} paratvāt śībhāvaḥ . \\
\noindent \textbf{(P\_1,1.11.2)} nityaḥ luk . \\
\noindent \textbf{(P\_1,1.11.2)} kṛte api śībhāve prāpnoti akṛte api prāpnoti . \\
\noindent \textbf{(P\_1,1.11.2)} anityaḥ luk . \\
\noindent \textbf{(P\_1,1.11.2)} anyasya kṛte śībhāve prāpnoti anyasya akṛte . \\
\noindent \textbf{(P\_1,1.11.2)} śabdāntarasya ca prāpnuvan vidhiḥ anityaḥ bhavati .śībhāvaḥ api anityaḥ . \\
\noindent \textbf{(P\_1,1.11.2)} na hi kṛte luki prāpnoti . \\
\noindent \textbf{(P\_1,1.11.2)} ubhayoḥ anityayoḥ paratvāt śībhāvaḥ śībhāve kṛte luk . \\
\noindent \textbf{(P\_1,1.11.2)} atha api katham cit nityaḥ luk syāt evam api doṣaḥ . \\
\noindent \textbf{(P\_1,1.11.2)} vakṣyati etat . \\
\noindent \textbf{(P\_1,1.11.2)} padasañjñāyām antavacanam anyatra sañjñāvidhau pratyayagrahaṇe tadantavidhipratiṣedhārtham iti. idam ca api pratyayagrahaṇam ayam ca api sañjñāvidhiḥ . \\
\noindent \textbf{(P\_1,1.11.2)} avaśyam khalu etasmin api pakṣe ādyantavadbhāvaḥ eṣitavyaḥ . \\
\noindent \textbf{(P\_1,1.11.2)} tasmāt astu saḥ eva madhyamaḥ pakṣaḥ . \\
\noindent \textbf{(P\_1,1.12)} māt pragṛhyasañjñāyām tasya asiddhatvāt ayāvekādeśapratiṣedhaḥ . \\
\noindent \textbf{(P\_1,1.12)} māt pragṛhyasañjñāyām tasya īttvasya ūttvasya ca asiddhatvāt ayāvekādeśāḥ prāpnuvanti . \\
\noindent \textbf{(P\_1,1.12)} teṣām pratiṣedhaḥ vaktavyaḥ . \\
\noindent \textbf{(P\_1,1.12)} amī* atra , amī* āsate , amū* atra , amū* āsāte . \\
\noindent \textbf{(P\_1,1.12)} nanu ca pragṛhyasañjñāvacanasāmarthyāt ayādayaḥ na bhaviṣyanti . \\
\noindent \textbf{(P\_1,1.12)} vacanārthaḥ hi siddhe . \\
\noindent \textbf{(P\_1,1.12)} na idam vacanāt labhyam . \\
\noindent \textbf{(P\_1,1.12)} asti hi anyat etasya vacane prayojanam . \\
\noindent \textbf{(P\_1,1.12)} kim . \\
\noindent \textbf{(P\_1,1.12)} yat siddhe pragṛhyasañjñākāryam tadartham etat syāt . \\
\noindent \textbf{(P\_1,1.12)} aṇaḥ apragṛhyasya anunāsikaḥ iti . \\
\noindent \textbf{(P\_1,1.12)} na ekam prayojanam yogārambham prayojayati . \\
\noindent \textbf{(P\_1,1.12)} yadi etāvat prayojanam syāt tatra eva ayam brūyāt aṇaḥ apragṛhyasya anunāsikaḥ adasaḥ na iti . \\
\noindent \textbf{(P\_1,1.12)} vipratiṣedhāt vā . \\
\noindent \textbf{(P\_1,1.12)} atha vā pragṛhyasañjñā kriyatām ayādayaḥ vā . \\
\noindent \textbf{(P\_1,1.12)} pragṛhyasañjñā bhaviṣyati vipratiṣedhena . \\
\noindent \textbf{(P\_1,1.12)} na eṣaḥ yuktaḥ vipratiṣedhaḥ . \\
\noindent \textbf{(P\_1,1.12)} vipratiṣedhe param iti ucyate . \\
\noindent \textbf{(P\_1,1.12)} pūrvā ca pragṛhyasañjñā pare ayādayaḥ . \\
\noindent \textbf{(P\_1,1.12)} parā pragṛhyasañjñā kariṣyate . \\
\noindent \textbf{(P\_1,1.12)} sūtraviparyāsaḥ kṛtaḥ bhavati . \\
\noindent \textbf{(P\_1,1.12)} evam tarhi parā eva pragṛhyasañjñā . \\
\noindent \textbf{(P\_1,1.12)} katham . \\
\noindent \textbf{(P\_1,1.12)} kāryakālam hi sañjñāparibhāṣam . \\
\noindent \textbf{(P\_1,1.12)} yatra kāryam tatra upasthitam draṣṭavyam . \\
\noindent \textbf{(P\_1,1.12)} pragṛhyaḥ prakṛtyā iti etat upasthitam bhavati adasaḥ māt iti . \\
\noindent \textbf{(P\_1,1.12)} evam api ayuktaḥ vipratiṣedhaḥ . \\
\noindent \textbf{(P\_1,1.12)} katham . \\
\noindent \textbf{(P\_1,1.12)} dvikāryayogaḥ hi vipratiṣedhaḥ . \\
\noindent \textbf{(P\_1,1.12)} na ca atra ekaḥ dvikāryayuktaḥ . \\
\noindent \textbf{(P\_1,1.12)} ecām ayādayaḥ . \\
\noindent \textbf{(P\_1,1.12)} īdūtoḥ pragṛhyasñjñā . \\
\noindent \textbf{(P\_1,1.12)} na avaśyam dvikāryayogaḥ eva vipratiṣedhaḥ . \\
\noindent \textbf{(P\_1,1.12)} kim tarhi . \\
\noindent \textbf{(P\_1,1.12)} asambhavaḥ api . \\
\noindent \textbf{(P\_1,1.12)} saḥ ca asti atra asambhavaḥ . \\
\noindent \textbf{(P\_1,1.12)} kaḥ asau asambhavaḥ . \\
\noindent \textbf{(P\_1,1.12)} pragṛhyasañjñā abhinirvartamānā ayādīn bādhate , ayādayaḥ abhinirvartamanāḥ pragṛhyasañjñānimittam vighnanti iti eṣaḥ asambhavaḥ . \\
\noindent \textbf{(P\_1,1.12)} sati asambhave yuktaḥ vipratiṣedhaḥ . \\
\noindent \textbf{(P\_1,1.12)} evam api ayuktaḥ vipratiṣedhaḥ . \\
\noindent \textbf{(P\_1,1.12)} satoḥ hi vipratiṣedhaḥ bhavati . \\
\noindent \textbf{(P\_1,1.12)} na ca atra īttvottve staḥ na api makāraḥ . \\
\noindent \textbf{(P\_1,1.12)} ubhayam asiddham . \\
\noindent \textbf{(P\_1,1.12)} āśrayāt siddhatvam ca yathā roḥ uttve . āśrayāt siddhatvam bhaviṣyati . \\
\noindent \textbf{(P\_1,1.12)} tat yathā ruḥ uttve āśrayāt siddhaḥ bhavati . \\
\noindent \textbf{(P\_1,1.12)} kim punaḥ kāraṇam ruḥ uttve āśrayāt siddhaḥ bhavati na punaḥ yatra eva ruḥ siddhaḥ tatra eva uttvam api ucyate . \\
\noindent \textbf{(P\_1,1.12)} na evam śakyam . \\
\noindent \textbf{(P\_1,1.12)} asiddhe hi uttve ādguṇāprasiddhiḥ . \\
\noindent \textbf{(P\_1,1.12)} asiddhe hi uttve ādguṇāprasiddhiḥ syāt . \\
\noindent \textbf{(P\_1,1.12)} vṛkṣaḥ atra , plakṣaḥ atra . \\
\noindent \textbf{(P\_1,1.12)} tasmāt tatra āśrayāt siddhatvam eṣitavyam . \\
\noindent \textbf{(P\_1,1.12)} tatra yathā āśrayāt siddham bhavati evam iha api bhaviṣyati . \\
\noindent \textbf{(P\_1,1.12)} atha vā pragṛhyasañjñāvacanasāmarthyāt ayādayaḥ ādeśāḥ na bhaviṣyanti . \\
\noindent \textbf{(P\_1,1.12)} atha vā yogavibhāgaḥ kariṣyate . \\
\noindent \textbf{(P\_1,1.12)} adasaḥ . \\
\noindent \textbf{(P\_1,1.12)} adasaḥ īdādayaḥ pragṛhyasañjñāḥ bhavanti . \\
\noindent \textbf{(P\_1,1.12)} tataḥ māt . \\
\noindent \textbf{(P\_1,1.12)} māt ca pare īdādayaḥ pragṛhyasañjñāḥ bhavanti . \\
\noindent \textbf{(P\_1,1.12)} adasaḥ iti eva . \\
\noindent \textbf{(P\_1,1.12)} kimarthaḥ yogavibhāgaḥ . \\
\noindent \textbf{(P\_1,1.12)} ekaḥ yat tat siddhe pragṛhyakāryam tadarthaḥ . \\
\noindent \textbf{(P\_1,1.12)} aparaḥ yat asiddhe . \\
\noindent \textbf{(P\_1,1.12)} iha api tarhi prāpnoti : amuyā , amuyoḥ iti . \\
\noindent \textbf{(P\_1,1.12)} kim ca syāt yadi pragṛhyasañjñā syāt . \\
\noindent \textbf{(P\_1,1.12)} pragṛhyāśrayaḥ prakṛtibhāvaḥ prasajyeta . \\
\noindent \textbf{(P\_1,1.12)} na eṣaḥ doṣaḥ . \\
\noindent \textbf{(P\_1,1.12)} padāntaprakaraṇe prakṛtibhāvaḥ . \\
\noindent \textbf{(P\_1,1.12)} na ca eṣaḥ padāntaḥ . \\
\noindent \textbf{(P\_1,1.12)} evam api amuke atra atra api prāpnoti . \\
\noindent \textbf{(P\_1,1.12)} dvivacanam iti vartate . \\
\noindent \textbf{(P\_1,1.12)} yadi dvivacanam iti vartate amī* atra iti na prāpnoti . \\
\noindent \textbf{(P\_1,1.12)} evam tarhi edantam iti nivṛttam . \\
\noindent \textbf{(P\_1,1.12)} atha vā āha ayam adasaḥ māt iti . \\
\noindent \textbf{(P\_1,1.12)} na ca īttvottve staḥ na api makāraḥ . \\
\noindent \textbf{(P\_1,1.12)} te evam vijñāsyāmaḥ mārthāt īdādyarthānām iti . \\
\noindent \textbf{(P\_1,1.12)} uktam vā . \\
\noindent \textbf{(P\_1,1.12)} kim uktam . \\
\noindent \textbf{(P\_1,1.12)} adasaḥ īttvottve svare bahiṣpadalakṣaṇe pragṛhyasañjāyām ca siddhe vaktavye iti . \\
\noindent \textbf{(P\_1,1.12)} tatra saki doṣaḥ . \\
\noindent \textbf{(P\_1,1.12)} tatra sakakāre doṣaḥ bhavati . \\
\noindent \textbf{(P\_1,1.12)} amuke atra . \\
\noindent \textbf{(P\_1,1.12)} na vā grahaṇaviśeṣaṇatvāt . \\
\noindent \textbf{(P\_1,1.12)} na vā eṣaḥ doṣaḥ . \\
\noindent \textbf{(P\_1,1.12)} kim kāraṇam . \\
\noindent \textbf{(P\_1,1.12)} grahaṇaviśeṣaṇatvāt . \\
\noindent \textbf{(P\_1,1.12)} na mādgrahaṇena īdādyantam viśeṣyate . \\
\noindent \textbf{(P\_1,1.12)} kim tarhi . \\
\noindent \textbf{(P\_1,1.12)} īdādayaḥ viśeṣyante . \\
\noindent \textbf{(P\_1,1.12)} māt pare ye īdādayaḥ iti . \\
\noindent \textbf{(P\_1,1.13)} iha kasmāt na bhavati : kāśe kuśe vaṃśe iti . \\
\noindent \textbf{(P\_1,1.13)} śe arthavadgrahaṇāt . \\
\noindent \textbf{(P\_1,1.13)} arthavataḥ śeśabdasya grahaṇam . \\
\noindent \textbf{(P\_1,1.13)} na ca ayam arthavān . \\
\noindent \textbf{(P\_1,1.13)} evam api hariśe babhruśe iti atra prāpnoti . \\
\noindent \textbf{(P\_1,1.13)} evam tarhi lakṣaṇapratipadoktayoḥ pratipadoktasya eva iti evam na bhaviṣyati . \\
\noindent \textbf{(P\_1,1.13)} atha vā punaḥ astu arthavadgrahaṇe na anarthakasya iti . \\
\noindent \textbf{(P\_1,1.13)} katham hariśe babhruśe iti . \\
\noindent \textbf{(P\_1,1.13)} ekaḥ atra vibhaktyarthena arthavān aparaḥ taddhitārthena . \\
\noindent \textbf{(P\_1,1.13)} samudāyaḥ anarthakaḥ . \\
\noindent \textbf{(P\_1,1.14)} nipātaḥ iti kimartham . \\
\noindent \textbf{(P\_1,1.14)} cakāra atra , jahāra atra . \\
\noindent \textbf{(P\_1,1.14)} ekāc iti kimartham . \\
\noindent \textbf{(P\_1,1.14)} pra idam brahma , pra idam kṣatram . \\
\noindent \textbf{(P\_1,1.14)} ekāc iti api ucyamāne atra api prāpnoti . \\
\noindent \textbf{(P\_1,1.14)} eṣaḥ api hi ekāc . \\
\noindent \textbf{(P\_1,1.14)} ekāc iti na ayam bahuvrīhiḥ : ekaḥ ac asmin saḥ ayam ekāc iti . \\
\noindent \textbf{(P\_1,1.14)} kim tarhi . \\
\noindent \textbf{(P\_1,1.14)} tatpuruṣaḥ ayam samānādhikaraṇaḥ : ekaḥ ac ekāc . \\
\noindent \textbf{(P\_1,1.14)} yadi tatpuruṣaḥ samānādhikaraṇaḥ na arthaḥ ekagrahaṇena . \\
\noindent \textbf{(P\_1,1.14)} iha kasmāt na bhavati : pra idam brahma , pra idam kṣatram . \\
\noindent \textbf{(P\_1,1.14)} ac eva yaḥ nipātaḥ iti evam vijñāsyate . \\
\noindent \textbf{(P\_1,1.14)} kim vaktavyam etat . \\
\noindent \textbf{(P\_1,1.14)} na hi . \\
\noindent \textbf{(P\_1,1.14)} katham anucyamānam gaṃsyate . \\
\noindent \textbf{(P\_1,1.14)} ajgrahaṇasāmarthyāt . \\
\noindent \textbf{(P\_1,1.14)} yadi hi yat ca ac ca anyat ca tatra syāt ajgrahaṇam anarthakam syāt . \\
\noindent \textbf{(P\_1,1.14)} asti anyat ajgrahaṇasya prayojanam . \\
\noindent \textbf{(P\_1,1.14)} kim. ajantasya yathā syāt . \\
\noindent \textbf{(P\_1,1.14)} halantasya mā bhūt . \\
\noindent \textbf{(P\_1,1.14)} na eva doṣaḥ na prayojanam . \\
\noindent \textbf{(P\_1,1.14)} evam api kutaḥ etat dvayoḥ paribhāṣayoḥ sāvakāśayoḥ samavasthitayoḥ ādyantavat ekasmin iti ca yena vidhiḥ tadantasya iti ca iyam iha paribhāṣā bhaviṣyati ādyantavat ekasmin iti iyam na bhaviṣyati yena vidhiḥ tadantasya iti . \\
\noindent \textbf{(P\_1,1.14)} ācāryapravṛttiḥ jñāpayati iyam iha paribhāṣā bhavati ādyantavat ekasmin iti iyam na bhavati yena vidhiḥ tadantasya iti yat ayam anāṅ iti pratiṣedham śāsti . \\
\noindent \textbf{(P\_1,1.14)} evam tarhi siddhe sati yat ajgrahaṇe kriyamāṇe ekagrahaṇam karoti tat jñāpayati ācāryaḥ anyatra varṇagrahaṇe jātigrahaṇam bhavati iti . \\
\noindent \textbf{(P\_1,1.14)} kim etasya jñāpane prayojanam . \\
\noindent \textbf{(P\_1,1.14)} dambheḥ halgrahaṇasya jātivācakatvāt siddham iti yat uktam tat upapannam bhavati . \\
\noindent \textbf{(P\_1,1.14)} anāṅ iti kimartham . \\
\noindent \textbf{(P\_1,1.14)} ā , udakāntāt odakāntāt . \\
\noindent \textbf{(P\_1,1.14)} iha kasmāt na bhavati: ā* evam nu manyase , ā* evam kila tat iti . \\
\noindent \textbf{(P\_1,1.14)} sānubandhakasya grahaṇam ananubandhakaḥ ca atra ākāraḥ . \\
\noindent \textbf{(P\_1,1.14)} kva punaḥ ayam sānubandhakaḥ kva niranubandhakaḥ . \\
\noindent \textbf{(P\_1,1.14)} īṣadarthe kriyāyoge maryādābhividhau ca yaḥ etam ātam ṅitam vidyāt vākyasmaraṇayoḥ aṅit \\
\noindent \textbf{(P\_1,1.15.1)} kim udāharaṇam . \\
\noindent \textbf{(P\_1,1.15.1)} āho* iti , utāho* iti . \\
\noindent \textbf{(P\_1,1.15.1)} na etat asti prayojanam . \\
\noindent \textbf{(P\_1,1.15.1)} nipātasamāhāraḥ ayam : āha , u : āho* iti, uta , āha , u : utāho* iti . \\
\noindent \textbf{(P\_1,1.15.1)} tatra nipātaḥ ekāc anāṅ iti eva siddham . \\
\noindent \textbf{(P\_1,1.15.1)} evam tarhi ekanipātāḥ ime . \\
\noindent \textbf{(P\_1,1.15.1)} atha vā pratiṣiddhārthaḥ ayam ārambhaḥ : o ṣu yātam marutaḥ , oṣu yātam bṛhatī śakvarī ca , o cit sakhāyam sakhya vavṛtyām . \\
\noindent \textbf{(P\_1,1.15.2)} otaḥ cvipratiṣedhaḥ . \\
\noindent \textbf{(P\_1,1.15.2)} odantaḥ nipātaḥ iti atra cvyantasya pratiṣedhaḥ vaktavyaḥ . \\
\noindent \textbf{(P\_1,1.15.2)} anadaḥ , adaḥ , abhavat : adobhavat , tirobhavat . \\
\noindent \textbf{(P\_1,1.15.2)} na vaktavyam . \\
\noindent \textbf{(P\_1,1.15.2)} lakṣaṇapratipadoktayoḥ pratipadoktasya eva iti evam na bhaviṣyati . \\
\noindent \textbf{(P\_1,1.15.2)} evam api agauḥ gauḥ sampadyate gobhavat : atra prāpnoti . \\
\noindent \textbf{(P\_1,1.15.2)} evam tarhi gauṇamukhyayoḥ mukhye kāryasamprayayaḥ iti . \\
\noindent \textbf{(P\_1,1.15.2)} tat yathā : gauḥ anubandhyaḥ ajaḥ agnīṣomīyaḥ iti na bāhīkaḥ anubadhyate . \\
\noindent \textbf{(P\_1,1.15.2)} katham tarhi bāhīke vṛddhyāttve bhavataḥ : gauḥ tiṣthati . \\
\noindent \textbf{(P\_1,1.15.2)} gām ānaya iti . \\
\noindent \textbf{(P\_1,1.15.2)} arthāśraye etat evam bhavati . \\
\noindent \textbf{(P\_1,1.15.2)} yat hi śabdāśrayam śabdamātre tat bhavati . \\
\noindent \textbf{(P\_1,1.15.2)} śabdāśraye ca vṛddhyāttve . \\
\noindent \textbf{(P\_1,1.17-18.1)} KA\_I,71.23-72.6 Ro\_I,233-234 {1/13}     iha kasmāt na bhavati : āho* iti , utāho* iti . \\
\noindent \textbf{(P\_1,1.17-18.1)} KA\_I,71.23-72.6 Ro\_I,233-234 {2/13}     uñaḥ iti ucyate . \\
\noindent \textbf{(P\_1,1.17-18.1)} KA\_I,71.23-72.6 Ro\_I,233-234 {3/13}     na ca atra uñam paśyāmaḥ . \\
\noindent \textbf{(P\_1,1.17-18.1)} KA\_I,71.23-72.6 Ro\_I,233-234 {4/13}     uñaḥ ayam anyena saha ekādeśaḥ uñgrahaṇena gṛhyate . \\
\noindent \textbf{(P\_1,1.17-18.1)} KA\_I,71.23-72.6 Ro\_I,233-234 {5/13}     ācāryapravṛttiḥ jñāpayati na uñekādeśaḥ uñgrahaṇena gṛhyate iti yat ayam ot iti odantasya nipātasya pragṛhyasañjñām śāsti . \\
\noindent \textbf{(P\_1,1.17-18.1)} KA\_I,71.23-72.6 Ro\_I,233-234 {6/13}     na etat asti jñāpakam . \\
\noindent \textbf{(P\_1,1.17-18.1)} KA\_I,71.23-72.6 Ro\_I,233-234 {7/13}     uktam etat pratiṣiddhārthaḥ ayam ārambhaḥ . \\
\noindent \textbf{(P\_1,1.17-18.1)} KA\_I,71.23-72.6 Ro\_I,233-234 {8/13}     doṣaḥ khalu api syāt yadi uñekādeśaḥ uñgrahaṇena na gṛhyeta : jānu , u . \\
\noindent \textbf{(P\_1,1.17-18.1)} KA\_I,71.23-72.6 Ro\_I,233-234 {9/13}     asya rujati jānū* asya rujati jānvasya rujati . \\
\noindent \textbf{(P\_1,1.17-18.1)} KA\_I,71.23-72.6 Ro\_I,233-234 {10/13}     mayaḥ uñaḥ vaḥ vā iti vatvam na syāt . \\
\noindent \textbf{(P\_1,1.17-18.1)} KA\_I,71.23-72.6 Ro\_I,233-234 {11/13}     evam tarhi ekanipātāḥ ime . \\
\noindent \textbf{(P\_1,1.17-18.1)} KA\_I,71.23-72.6 Ro\_I,233-234 {12/13}     atha vā dvau ukārau imau ekaḥ ananubandhakaḥ aparaḥ sānubandhakaḥ . \\
\noindent \textbf{(P\_1,1.17-18.1)} KA\_I,71.23-72.6 Ro\_I,233-234 {13/13}     tat yaḥ ananubandhakaḥ tasya eṣaḥ ekādeśaḥ . \\
\noindent \textbf{(P\_1,1.17-18.2)} KA\_I,72.7-13 Ro\_I,234-235 {1/10}     uñaḥ iti yogavibhāgaḥ . \\
\noindent \textbf{(P\_1,1.17-18.2)} KA\_I,72.7-13 Ro\_I,234-235 {2/10}     uñaḥ iti yogavibhāgaḥ kartavyaḥ . \\
\noindent \textbf{(P\_1,1.17-18.2)} KA\_I,72.7-13 Ro\_I,234-235 {3/10}     uñaḥ śākalyasya ācāryasya matena pragṛhyasañjñā bhavati . \\
\noindent \textbf{(P\_1,1.17-18.2)} KA\_I,72.7-13 Ro\_I,234-235 {4/10}     u* iti v iti . \\
\noindent \textbf{(P\_1,1.17-18.2)} KA\_I,72.7-13 Ro\_I,234-235 {5/10}     tataḥ uŚ . \\
\noindent \textbf{(P\_1,1.17-18.2)} KA\_I,72.7-13 Ro\_I,234-235 {6/10}     uñaḥ ūŚ iti ayam ādeśaḥ bhavati śākalyasya ācāryasya matena dīrghaḥ anunāsikaḥ pragṛhyasañjñakaḥ ca uŚ iti . \\
\noindent \textbf{(P\_1,1.17-18.2)} KA\_I,72.7-13 Ro\_I,234-235 {7/10}     kimarthaḥ yogavibhāgaḥ . \\
\noindent \textbf{(P\_1,1.17-18.2)} KA\_I,72.7-13 Ro\_I,234-235 {8/10}      uŚ vā śākalyasya . \\
\noindent \textbf{(P\_1,1.17-18.2)} KA\_I,72.7-13 Ro\_I,234-235 {9/10}     śākalyasya ācāryasya matena uŚ vibhāṣā yathā syāt : ūŚ iti , u* iti . \\
\noindent \textbf{(P\_1,1.17-18.2)} KA\_I,72.7-13 Ro\_I,234-235 {10/10}     anyeṣām ācāryāṇām matena v iti . \\
\noindent \textbf{(P\_1,1.19)} īdūtau saptamī iti eva . īdūtau saptamī iti eva siddham . \\
\noindent \textbf{(P\_1,1.19)} na arthaḥ arthagrahaṇena . \\
\noindent \textbf{(P\_1,1.19)} lupte arthagrahaṇāt bhavet . \\
\noindent \textbf{(P\_1,1.19)} luptāyam saptamyām pragṛhyasañjñā na prāpnoti . \\
\noindent \textbf{(P\_1,1.19)} kva . \\
\noindent \textbf{(P\_1,1.19)} somo gaurī adhi śritaḥ . \\
\noindent \textbf{(P\_1,1.19)} iṣyate ca atra api syāt iti . \\
\noindent \textbf{(P\_1,1.19)} tat ca antareṇa yatnam na sidhyati iti evamartham arthagrahaṇam . \\
\noindent \textbf{(P\_1,1.19)} na atra saptamī lupyate . \\
\noindent \textbf{(P\_1,1.19)} kim tarhi . \\
\noindent \textbf{(P\_1,1.19)} pūrvasavarṇaḥ atra bhavati . pūrvasya cet savarṇaḥ asau āḍāmbhāvaḥ prasajyate . \\
\noindent \textbf{(P\_1,1.19)} yadi pūrvasavarṇaḥ āṭ āmbhāvaḥ ca prāpnoti . \\
\noindent \textbf{(P\_1,1.19)} evam tarhi āha ayam īdūtau saptamī iti na sa asti saptamī īdūtau . \\
\noindent \textbf{(P\_1,1.19)} tatra vacanāt bhaviṣyati . \\
\noindent \textbf{(P\_1,1.19)} vacanāt yatra dīrghatvam . \\
\noindent \textbf{(P\_1,1.19)} na idam vacanāt labhyam. asti hi anyat etasya vacane prayojanam . \\
\noindent \textbf{(P\_1,1.19)} kim . \\
\noindent \textbf{(P\_1,1.19)} yatra saptamyāḥ dīrghatvam ucyate : dṛtim na śuṣkam sarasī śayānam iti . \\
\noindent \textbf{(P\_1,1.19)} sati prayojane iha na prāpnoti somo gaurī adhi śritaḥ iti . \\
\noindent \textbf{(P\_1,1.19)} tatra api sarasī yadi . \\
\noindent \textbf{(P\_1,1.19)} tatra api siddham . \\
\noindent \textbf{(P\_1,1.19)} katham . \\
\noindent \textbf{(P\_1,1.19)} yadi sarasīśabdasya pravṛttiḥ asti . \\
\noindent \textbf{(P\_1,1.19)} asti ca loke sarasīśabdasya pravṛttiḥ . \\
\noindent \textbf{(P\_1,1.19)} katham . \\
\noindent \textbf{(P\_1,1.19)} dakṣiṇāpathe hi mahānti sarāṃsi sarasyaḥ iti ucyante . \\
\noindent \textbf{(P\_1,1.19)} jñāpakam syāt tadantatve . \\
\noindent \textbf{(P\_1,1.19)} evam tarhi jñāpayati ācāryaḥ na pragṛhyasañjñāyām pratyayalakṣaṇam bhavati iti . \\
\noindent \textbf{(P\_1,1.19)} kim etasya jñāpane prayojanam . \\
\noindent \textbf{(P\_1,1.19)} kumāryoḥ agāram kumāryagāram , vadhvoḥ agāram vadhvagāram . \\
\noindent \textbf{(P\_1,1.19)} pratyayalakṣaṇena pragṛhyasañjñā na bhavati . \\
\noindent \textbf{(P\_1,1.19)} mā vā pūrvapadasya bhūt . atha vā pūrvapadasya mā bhūt iti evamartham arthagrahaṇam : vāpyām aśvaḥ vāpyaśvaḥ , nadyām ātiḥ nadyātiḥ . \\
\noindent \textbf{(P\_1,1.19)} atha kriyamāṇe api arthagrahaṇe kasmāt eva atra na bhavati . \\
\noindent \textbf{(P\_1,1.19)} jahatsvārthā vṛttiḥ iti . \\
\noindent \textbf{(P\_1,1.19)} atha ajahatsvārthāyām vṛttau doṣaḥ eva . \\
\noindent \textbf{(P\_1,1.19)} ajahatsvārthāyām ca na doṣaḥ . \\
\noindent \textbf{(P\_1,1.19)} samudāyārthaḥ abhidhīyate . \\
\noindent \textbf{(P\_1,1.19)} īdutau saptamī iti eva lupte arthagrahaṇāt bhavet pūrvasya cet savarṇaḥ asau āḍāmbhāvaḥ prasajyate vacanāt yatra dīrghatvam tatra api sarasī yadi jñāpakam syāt tadantatve mā vā pūrvapadasya bhūt . \\
\noindent \textbf{(P\_1,1.20.1)} ghusañjñāyām prakṛtigrahaṇam śidartham . ghusañjñāyām prakṛtigrahaṇam kartavyam . \\
\noindent \textbf{(P\_1,1.20.1)} dādhāprakṛtayaḥ ghusañjñā bhavanti iti vaktavyam . \\
\noindent \textbf{(P\_1,1.20.1)} kim prayojanam . \\
\noindent \textbf{(P\_1,1.20.1)} āttvabhūtānām iyam sañjñā kriyate . \\
\noindent \textbf{(P\_1,1.20.1)} sā āttvabhūtānām eva syāt anāttvabhūtānām na syāt . \\
\noindent \textbf{(P\_1,1.20.1)} nanu ca bhūyiṣthāni ghusañjñākāryāṇi ārdhadhātuke tatra ca ete āttvabhūtāḥ dṛśyante . \\
\noindent \textbf{(P\_1,1.20.1)} śidartham . \\
\noindent \textbf{(P\_1,1.20.1)} śidartham prakṛtigrahaṇam kartavyam . \\
\noindent \textbf{(P\_1,1.20.1)} śiti āttvam pratiṣidhyate tadartham : praṇidayate praṇidhayati iti . \\
\noindent \textbf{(P\_1,1.20.1)} bhāradvājīyāḥ paṭhanti ghusañjñāyām prakṛtigrahaṇam śidvikṛtārtham . \\
\noindent \textbf{(P\_1,1.20.1)} ghusañjñāyām prakṛtigrahaṇam kriyate . \\
\noindent \textbf{(P\_1,1.20.1)} kim prayojanam . \\
\noindent \textbf{(P\_1,1.20.1)} śidartham vikṛtārtham ca . \\
\noindent \textbf{(P\_1,1.20.1)} śiti udāhṛtam . \\
\noindent \textbf{(P\_1,1.20.1)} vikṛtārtham khalu api : praṇidātā praṇidhātā . \\
\noindent \textbf{(P\_1,1.20.1)} kim punaḥ kāraṇam na sidhyati . \\
\noindent \textbf{(P\_1,1.20.1)} lakṣaṇapratipadoktayoḥ pratipadoktasya eva iti pratipadam ye āttvabhūtāḥ teṣām eva syāt . \\
\noindent \textbf{(P\_1,1.20.1)} lakṣaṇena ye āttvabhūtāḥ teṣām na syāt . \\
\noindent \textbf{(P\_1,1.20.1)} atha kriyamāṇe api prakṛtigrahaṇe katham idam vijñāyate . \\
\noindent \textbf{(P\_1,1.20.1)} dādhāḥ prakṛtayaḥ āhosvit dādhām prakṛtayaḥ iti . \\
\noindent \textbf{(P\_1,1.20.1)} kim ca ataḥ . \\
\noindent \textbf{(P\_1,1.20.1)} yadi vijñāyate dādhāḥ prakṛtayaḥ iti saḥ eva doṣaḥ . \\
\noindent \textbf{(P\_1,1.20.1)} āttvabhūtānām eva syāt anāttvabhūtānām na syāt . \\
\noindent \textbf{(P\_1,1.20.1)} atha vijñāyate dādhām prakṛtayaḥ iti anāttvabhūtānām eva syāt āttvabhūtānām na syāt . \\
\noindent \textbf{(P\_1,1.20.1)} evam tarhi na evam vijñāyate dādhāḥ prakṛtayaḥ iti na api dādhām prakṛtayaḥ iti . \\
\noindent \textbf{(P\_1,1.20.1)} katham tarhi . \\
\noindent \textbf{(P\_1,1.20.1)} dādhāḥ ghusañjñāḥ bhavanti prakṛtayaḥ ca eṣām iti . \\
\noindent \textbf{(P\_1,1.20.1)} tat tarhi prakṛtigrahaṇam kartavyam . \\
\noindent \textbf{(P\_1,1.20.1)} na kartavyam . \\
\noindent \textbf{(P\_1,1.20.1)} idam prakṛtam arthagrahaṇam anuvartate . \\
\noindent \textbf{(P\_1,1.20.1)} kva prakṛtam . \\
\noindent \textbf{(P\_1,1.20.1)} īdūtau ca saptamyarthe iti . \\
\noindent \textbf{(P\_1,1.20.1)} tataḥ vakṣyāmi dādhāḥ ghu adāp . \\
\noindent \textbf{(P\_1,1.20.1)} arthe iti . \\
\noindent \textbf{(P\_1,1.20.1)} na evam śakyam . \\
\noindent \textbf{(P\_1,1.20.1)} dadātinā samānārthān rātirāsatidāśatimaṃhatiprīṇātiprabhṛtīn āhuḥ . \\
\noindent \textbf{(P\_1,1.20.1)} eteṣām api ghusañjñā prāpnoti . \\
\noindent \textbf{(P\_1,1.20.1)} tasmāt na evam śakyam . \\
\noindent \textbf{(P\_1,1.20.1)} na cet evam prakṛtigrahaṇam kartavyam . \\
\noindent \textbf{(P\_1,1.20.1)} śidarthena tāvat na arthaḥ prakṛtigrahaṇena . \\
\noindent \textbf{(P\_1,1.20.1)} avaśyam tatra mārtham prakṛtigrahaṇam kartavyam praṇimayate praṇyamayata iti evamartham . \\
\noindent \textbf{(P\_1,1.20.1)} tat purastāt apakrakṣyate : ghuprakṛtau māprakṛtau ca iti . \\
\noindent \textbf{(P\_1,1.20.1)} yadi prakṛtigrahaṇam kriyate praniminoti pranimīnāti atra api prāpnoti . \\
\noindent \textbf{(P\_1,1.20.1)} atha akriyamāṇe api prakṛtigrahaṇe iha kasmāt na bhavati : pranimātā pranimātum iti . \\
\noindent \textbf{(P\_1,1.20.1)} ākārāntasya ṅitaḥ grahaṇam vijñāsyate . \\
\noindent \textbf{(P\_1,1.20.1)} yathā eva tarhi akriyamāṇe prakṛtigrahaṇe ākārāntasya ṅitaḥ grahaṇam vijñāyate evam kriyamāṇe api prakṛtigrahaṇe ākārāntasya ṅitaḥ grahaṇam vijñāsyate . \\
\noindent \textbf{(P\_1,1.20.1)} vikṛtārthena ca api na arthaḥ . \\
\noindent \textbf{(P\_1,1.20.1)} doṣaḥ eva etasyāḥ paribhāṣāyāḥ lakṣaṇapratipadoktayoḥ pratipadoktasya eva iti gāmādāgrahaṇeṣu aviśeṣaḥ iti . \\
\noindent \textbf{(P\_1,1.20.2)} samānaśabdapratiṣedhaḥ . \\
\noindent \textbf{(P\_1,1.20.2)} samānaśabdānām pratiṣedhaḥ vaktavyaḥ : pranidārayati pranidhārayati . \\
\noindent \textbf{(P\_1,1.20.2)} dādhāḥ ghusañjñāḥ bhavanti iti ghusañjñā prāpnoti . \\
\noindent \textbf{(P\_1,1.20.2)} samānaśabdāpratiṣedhaḥ arthavadgrahaṇāt . \\
\noindent \textbf{(P\_1,1.20.2)} samānaśabdānām apratiṣedhaḥ . \\
\noindent \textbf{(P\_1,1.20.2)} anarthakaḥ pratiṣedhaḥ apratiṣedhaḥ . \\
\noindent \textbf{(P\_1,1.20.2)} ghusañjñā kasmāt na bhavati . \\
\noindent \textbf{(P\_1,1.20.2)} arthavadgrahaṇāt . \\
\noindent \textbf{(P\_1,1.20.2)} arthavatoḥ dādhoḥ grahaṇam . \\
\noindent \textbf{(P\_1,1.20.2)} na ca etau arthavantau . \\
\noindent \textbf{(P\_1,1.20.2)} anupasargāt vā . \\
\noindent \textbf{(P\_1,1.20.2)} atha vā yatkriyāyuktās prādayaḥ tam prati gatyupasargasañjñe bhavataḥ . \\
\noindent \textbf{(P\_1,1.20.2)} na ca etau dādhau prati kriyāyogaḥ . \\
\noindent \textbf{(P\_1,1.20.2)} yadi evam iha api tarhi na prāpnoti praṇidāpayati praṇidhāpayati . \\
\noindent \textbf{(P\_1,1.20.2)} atra api na etau dādhau arthavantau na api etau dādhau prati kriyāyogaḥ . \\
\noindent \textbf{(P\_1,1.20.2)} na vā arthavataḥ hi āgamaḥ tadguṇībhūtaḥ tadgrahaṇena gṛhyate yathā anyatra . \\
\noindent \textbf{(P\_1,1.20.2)} na vā eṣaḥ doṣaḥ . \\
\noindent \textbf{(P\_1,1.20.2)} kim kāraṇam . \\
\noindent \textbf{(P\_1,1.20.2)} arthavataḥ āgamaḥ tadguṇībhūtaḥ arthavadgrahaṇena gṛhyate yathā anyatra . \\
\noindent \textbf{(P\_1,1.20.2)} tat yathā . \\
\noindent \textbf{(P\_1,1.20.2)} anyatra api arthavataḥ āgamaḥ arthavadgrahaṇena gṛhyate . \\
\noindent \textbf{(P\_1,1.20.2)} kva anyatra . \\
\noindent \textbf{(P\_1,1.20.2)} lavitā cikīrṣitā iti . \\
\noindent \textbf{(P\_1,1.20.2)} yuktam punaḥ yat nityeṣu nāma śabdeṣu āgamaśāsanam syāt na nityeṣu śabdeṣu kūṭasthaiḥ avicālibhiḥ varṇaiḥ bhavitavyam anapāyopajanavikāribhiḥ . \\
\noindent \textbf{(P\_1,1.20.2)} āgamaḥ ca nāma apūrvaḥ śabdopajanaḥ . \\
\noindent \textbf{(P\_1,1.20.2)} atha yuktam yat nityeṣu śabdeṣu ādeśāḥ syuḥ . \\
\noindent \textbf{(P\_1,1.20.2)} bāḍham yuktam . \\
\noindent \textbf{(P\_1,1.20.2)} śabdāntaraiḥ iha bhavitavyam . \\
\noindent \textbf{(P\_1,1.20.2)} tatra śabdāntarāt śabdāntarasya pratipattiḥ yuktā . \\
\noindent \textbf{(P\_1,1.20.2)} ādeśāḥ tarhi ime bhaviṣyanti anāgamakānām sāgamakāḥ . \\
\noindent \textbf{(P\_1,1.20.2)} tat katham . \\
\noindent \textbf{(P\_1,1.20.2)} sarve sarvapadādeśāḥ dākṣīputrasya pāṇineḥ ekadeśavikāre hi nityatvam na upapadyate . \\
\noindent \textbf{(P\_1,1.20.3)} dīṅaḥ pratiṣedhaḥ sthāghvoḥ ittve . \\
\noindent \textbf{(P\_1,1.20.3)} dīṅaḥ pratiṣedhaḥ sthāghvoḥ ittve vaktavyaḥ . \\
\noindent \textbf{(P\_1,1.20.3)} upādāsta asya svaraḥ śikṣakasya iti . \\
\noindent \textbf{(P\_1,1.20.3)} mīnātiminoti iti āttve kṛte sthāghvoḥ it ca iti ittvam prāpnoti . \\
\noindent \textbf{(P\_1,1.20.3)} kutaḥ punaḥ ayam doṣaḥ jāyate . \\
\noindent \textbf{(P\_1,1.20.3)} kim prakṛtigrahaṇāt āhosvit rūpagrahaṇāt . \\
\noindent \textbf{(P\_1,1.20.3)} rūpagrahaṇāt iti āha . \\
\noindent \textbf{(P\_1,1.20.3)} iha khalu prakṛtigrahaṇāt doṣaḥ jāyate : upadidīṣate . \\
\noindent \textbf{(P\_1,1.20.3)} sani mīmāghurabhalabha iti . \\
\noindent \textbf{(P\_1,1.20.3)} na eṣaḥ doṣaḥ . \\
\noindent \textbf{(P\_1,1.20.3)} dāprakṛtiḥ iti ucyate . \\
\noindent \textbf{(P\_1,1.20.3)} na ca iyam dāprakṛtiḥ . \\
\noindent \textbf{(P\_1,1.20.3)} ākārāntānām ejantāḥ prakṛtayaḥ ejantānām api īkārāntāḥ . \\
\noindent \textbf{(P\_1,1.20.3)} na ca prakṛtiprakṛtiḥ prakṛtigrahaṇena gṛhyate . \\
\noindent \textbf{(P\_1,1.20.3)} saḥ tarhi pratiṣedhaḥ vaktavyaḥ . \\
\noindent \textbf{(P\_1,1.20.3)} na vaktavyaḥ . \\
\noindent \textbf{(P\_1,1.20.3)} ghusañjñā kasmāt na bhavati . \\
\noindent \textbf{(P\_1,1.20.3)} sannipātalakṣaṇaḥ vidhiḥ animittam tadvighātasya iti evam na bhaviṣyati . \\
\noindent \textbf{(P\_1,1.20.4)} dāppratiṣedhe na daipi anejantatvāt . \\
\noindent \textbf{(P\_1,1.20.4)} dāppratiṣedhe daipi pratiṣedhaḥ na prāpnoti : avadātam mukham . \\
\noindent \textbf{(P\_1,1.20.4)} nanu ca āttve kṛte bhaviṣyati . \\
\noindent \textbf{(P\_1,1.20.4)} tat hi āttvam na prāpnoti . \\
\noindent \textbf{(P\_1,1.20.4)} kim kāraṇam . \\
\noindent \textbf{(P\_1,1.20.4)} anejantatvāt . \\
\noindent \textbf{(P\_1,1.20.4)} siddham anubandhasya anekāntatvāt . siddham etat . \\
\noindent \textbf{(P\_1,1.20.4)} katham . \\
\noindent \textbf{(P\_1,1.20.4)} anubandhasya anekāntatvāt . \\
\noindent \textbf{(P\_1,1.20.4)} anekāntāḥ anubandhāḥ . \\
\noindent \textbf{(P\_1,1.20.4)} pitpratiṣedhāt vā . \\
\noindent \textbf{(P\_1,1.20.4)} atha vā dādhāḥ ghu apit iti vakṣyāmi . \\
\noindent \textbf{(P\_1,1.20.4)} tat ca avaśyam vaktavyam . \\
\noindent \textbf{(P\_1,1.20.4)} adāp iti hi ucyamāne iha api prasajyeta : praṇidāpayati iti . \\
\noindent \textbf{(P\_1,1.20.4)} śakyam tāvat anena adāp iti bruvatā bāntasya pratiṣedhaḥ vijñātum . \\
\noindent \textbf{(P\_1,1.20.4)} sūtram tarhi bhidyate . \\
\noindent \textbf{(P\_1,1.20.4)} yathānyāsam eva astu . \\
\noindent \textbf{(P\_1,1.20.4)} nanu ca uktam dappratiṣedhe na daipi iti . \\
\noindent \textbf{(P\_1,1.20.4)} parihṛtam etat siddham anubandhasya anekāntatvāt iti . \\
\noindent \textbf{(P\_1,1.20.4)} atha ekānteṣu doṣaḥ eva . \\
\noindent \textbf{(P\_1,1.20.4)} ekānteṣu ca na doṣaḥ . \\
\noindent \textbf{(P\_1,1.20.4)} āttve kṛte bhaviṣyati . \\
\noindent \textbf{(P\_1,1.20.4)} nanu ca uktam tat hi āttvam na prāpnoti . \\
\noindent \textbf{(P\_1,1.20.4)} kim kāraṇam . \\
\noindent \textbf{(P\_1,1.20.4)} anejantatvāt iti . \\
\noindent \textbf{(P\_1,1.20.4)} pakāralope kṛte bhaviṣyati . \\
\noindent \textbf{(P\_1,1.20.4)} na hi ayam tadā dāp bhavati . \\
\noindent \textbf{(P\_1,1.20.4)} bhūtapūrvagatyā bhaviṣyati . \\
\noindent \textbf{(P\_1,1.20.4)} etat ca atra yuktam yat sarveṣu eva sānubandhakagrahaṇeṣu bhūtapūrvagatiḥ vijñāyate . \\
\noindent \textbf{(P\_1,1.20.4)} anaimittikaḥ hi anubandhalopaḥ tāvati eva bhavati . \\
\noindent \textbf{(P\_1,1.20.4)} atha vā ācāryapravṛttiḥ jñāpayati na anubandhakṛtam anejantatvam iti yat ayam udīcām māṅaḥ vyatīhāre iti meṅaḥ sānubandhakasya āttvabhūtasya grahaṇam karoti . \\
\noindent \textbf{(P\_1,1.20.4)} atha vā dāp eva ayam na daip asti . \\
\noindent \textbf{(P\_1,1.20.4)} katham avadāyayati iti . \\
\noindent \textbf{(P\_1,1.20.4)} śyan vikaraṇaḥ bhaviṣyati . \\
\noindent \textbf{(P\_1,1.21.1)} kimartham idam ucyate . \\
\noindent \textbf{(P\_1,1.21.1)} sati anyasmin ādyantavadbhāvāt ekasmin ādyantavadvacanam . \\
\noindent \textbf{(P\_1,1.21.1)} sati anyasmin yasmāt pūrvam na asti param asti saḥ ādiḥ iti ucyate . \\
\noindent \textbf{(P\_1,1.21.1)} sati anyasmin yasmāt param na asti pūrvam asti saḥ antaḥ iti ucyate . \\
\noindent \textbf{(P\_1,1.21.1)} sati anyasmin ādyantavadbhāvāt etasmāt kāraṇāt ekasmin ādyantāpadiṣṭāni kāryāṇi na sidhyanti . \\
\noindent \textbf{(P\_1,1.21.1)} iṣyante ca syuḥ iti . \\
\noindent \textbf{(P\_1,1.21.1)} tāni antareṇa yatnam na sidhyanti iti ekasmin ādyantavadvacanam . \\
\noindent \textbf{(P\_1,1.21.1)} evamartham idam ucyate . \\
\noindent \textbf{(P\_1,1.21.1)} asti prayojanam etat . \\
\noindent \textbf{(P\_1,1.21.1)} kim tarhi iti . \\
\noindent \textbf{(P\_1,1.21.1)} tatra vyapadeśivadvacanam . \\
\noindent \textbf{(P\_1,1.21.1)} tatra vyapadeśivadbhāvaḥ vaktavyaḥ . \\
\noindent \textbf{(P\_1,1.21.1)} vyapadeśivat ekasmin kāryam bhavati iti vaktavyam . \\
\noindent \textbf{(P\_1,1.21.1)} kim prayojanam . \\
\noindent \textbf{(P\_1,1.21.1)} ekācaḥ dve prathamārtham . \\
\noindent \textbf{(P\_1,1.21.1)} vakṣyati ekācaḥ dve prathamasya iti bahuvrīhinirdeśaḥ iti . \\
\noindent \textbf{(P\_1,1.21.1)} tasmin kriyamāṇe iha eva : syāt papāca papāṭha . \\
\noindent \textbf{(P\_1,1.21.1)} iyāya , āra iti atra na syāt . \\
\noindent \textbf{(P\_1,1.21.1)} vyapdeśivat ekasmin kāryam bhavati iti atra api siddham bhavati . \\
\noindent \textbf{(P\_1,1.21.1)} ṣatve ca ādeśasampratyayārtham . \\
\noindent \textbf{(P\_1,1.21.1)} vakṣyati ādeśapratyayayoḥ iti avayavaṣaṣṭhī eva iti . \\
\noindent \textbf{(P\_1,1.21.1)} etasmin kriyamāṇe iha eva syāt : kariṣyati hariṣyati . \\
\noindent \textbf{(P\_1,1.21.1)} iha na syāt : indraḥ mā vakṣat , saḥ devan yakṣat . \\
\noindent \textbf{(P\_1,1.21.1)} vyapdeśivat ekasmin kāryam bhavati iti atra api siddham bhavati . \\
\noindent \textbf{(P\_1,1.21.1)} saḥ tarhi vyapadeśivadbhāvaḥ vaktavyaḥ . \\
\noindent \textbf{(P\_1,1.21.1)} na vaktavyaḥ . \\
\noindent \textbf{(P\_1,1.21.1)} avacanāt lokavijñānāt siddham . \\
\noindent \textbf{(P\_1,1.21.1)} antareṇa eva vacanam lokavijñānāt siddham etat . \\
\noindent \textbf{(P\_1,1.21.1)} tat yathā : loke śālāsamudāyaḥ grāmaḥ iti ucyate . \\
\noindent \textbf{(P\_1,1.21.1)} bhavati ca etat ekasmin api ekaśālaḥ grāmaḥ iti . \\
\noindent \textbf{(P\_1,1.21.1)} viṣamaḥ upanyāsaḥ . \\
\noindent \textbf{(P\_1,1.21.1)} grāmaśabdaḥ ayam bahvarthaḥ . \\
\noindent \textbf{(P\_1,1.21.1)} asti eva śālāsamudāye vartate . \\
\noindent \textbf{(P\_1,1.21.1)} tat yathā grāmaḥ dagdhaḥ iti . \\
\noindent \textbf{(P\_1,1.21.1)} asti vāṭaparikṣepe vartate . \\
\noindent \textbf{(P\_1,1.21.1)} tat yathā grāmam praviṣṭaḥ . \\
\noindent \textbf{(P\_1,1.21.1)} asti manuṣyeṣu vartate . \\
\noindent \textbf{(P\_1,1.21.1)} tat yathā . \\
\noindent \textbf{(P\_1,1.21.1)} grāmaḥ gataḥ , grāmaḥ āgataḥ iti . \\
\noindent \textbf{(P\_1,1.21.1)} asti sāraṇyake sasīmake sasthaṇḍilake vartate . \\
\noindent \textbf{(P\_1,1.21.1)} tat yathā grāmaḥ labdhaḥ iti . \\
\noindent \textbf{(P\_1,1.21.1)} tat yaḥ sāraṇyake sasīmake sasthaṇḍilake vartate tam abhisamīkṣya etat prayujyate : ekaśālaḥ grāmaḥ iti . \\
\noindent \textbf{(P\_1,1.21.1)} yathā tarhi varṇasamudāyaḥ padam padasamudāyaḥ ṛk ṛksamudāyaḥ sūktam iti ucyate . \\
\noindent \textbf{(P\_1,1.21.1)} bhavati ca etat ekasmin api ekavarṇam padam ekapadā ṛk ekarcam sūktam iti . \\
\noindent \textbf{(P\_1,1.21.1)} atra api arthena yuktaḥ vyapadeśaḥ . \\
\noindent \textbf{(P\_1,1.21.1)} padam nāma arthaḥ sūktam nama arthaḥ . \\
\noindent \textbf{(P\_1,1.21.1)} yathā tarhi bahuṣu putreṣu etat upapannam : bhavati ayam me jyeṣṭhaḥ ayam eva me madhyamaḥ ayam eva me kanīyān iti . \\
\noindent \textbf{(P\_1,1.21.1)} bhavati ca etat ekasmin api ayam eva me jyeṣṭhaḥ ayam me madhyamaḥ ayam me kanīyān iti . \\
\noindent \textbf{(P\_1,1.21.1)} tathā asūtāyām asoṣyamāṇāyām ca bhavati prathamagarbheṇa hatā iti . \\
\noindent \textbf{(P\_1,1.21.1)} tathā anetya anājigamiṣuḥ āha idam me prathamam āgamanam iti . \\
\noindent \textbf{(P\_1,1.21.1)} ādyantavadbhāvaḥ ca śakyaḥ avaktum . \\
\noindent \textbf{(P\_1,1.21.1)} katham . \\
\noindent \textbf{(P\_1,1.21.1)} apūrvānuttaralakṣaṇatvāt ādyantayoḥ siddham ekasmin . \\
\noindent \textbf{(P\_1,1.21.1)} apūrvalakṣaṇaḥ ādiḥ anuttaralakṣaṇaḥ antaḥ . \\
\noindent \textbf{(P\_1,1.21.1)} etat ca ekasmin api bhavati . \\
\noindent \textbf{(P\_1,1.21.1)} apūrvānuttaralakṣaṇatvāt etasmāt kāraṇāt ekasmin api ādyantāpadiṣṭani kāryāṇi bhaviṣyanti . \\
\noindent \textbf{(P\_1,1.21.1)} na arthaḥ ādyantavadbhāvena . \\
\noindent \textbf{(P\_1,1.21.1)} gonardīyaḥ tu āha satyam etat sati tu anyasmin iti . \\
\noindent \textbf{(P\_1,1.21.2)} kāni punaḥ asya yogasya prayojanāni . \\
\noindent \textbf{(P\_1,1.21.2)} ādivattve prayojanam pratyayañnidādyudāttatve . \\
\noindent \textbf{(P\_1,1.21.2)} pratyayasya ādiḥ udāttaḥ bhavati iti iha eva syāt : kartavyam , taittirīyaḥ . \\
\noindent \textbf{(P\_1,1.21.2)} aupagavaḥ , kāpaṭavaḥ iti atra na syāt . \\
\noindent \textbf{(P\_1,1.21.2)} ñniti ādiḥ nityam iti iha eva syāt : ahicumbakāyaniḥ , agniveśyaḥ . \\
\noindent \textbf{(P\_1,1.21.2)} gargyaḥ , kṛtiḥ iti atra na syāt . \\
\noindent \textbf{(P\_1,1.21.2)} valādeḥ ārdhadhātukasya iṭ . \\
\noindent \textbf{(P\_1,1.21.2)} valādeḥ ārdhadhātukasya iṭ prayojanam . \\
\noindent \textbf{(P\_1,1.21.2)} ārdhadhātukasya iṭ valādeḥ iha eva : syāt kariṣyati hariṣyati . \\
\noindent \textbf{(P\_1,1.21.2)} joṣiṣat , manidṣat iti atra na syāt . \\
\noindent \textbf{(P\_1,1.21.2)} yasmin vidhiḥ tadāditve . \\
\noindent \textbf{(P\_1,1.21.2)} yasmin vidhiḥ tadāditve prayojanam . \\
\noindent \textbf{(P\_1,1.21.2)} vakṣyati yasmin vidhiḥ tadādau algrahaṇe iti . \\
\noindent \textbf{(P\_1,1.21.2)} tasmin kriyamāṇe aci śnudhātubhruvām yvoḥ iyaṅuvaṅau iha eva syāt : śriyaḥ , bhruvaḥ . \\
\noindent \textbf{(P\_1,1.21.2)} śriyau bhruvau iti atra na syāt . \\
\noindent \textbf{(P\_1,1.21.2)} ajādyāṭtve . ajādyāṭtve prayojanam . \\
\noindent \textbf{(P\_1,1.21.2)} āṭ ajādīnām iha eva syāt : aihiṣṭa , aikṣiṣṭa . \\
\noindent \textbf{(P\_1,1.21.2)} ait , adhyaiṣṭa iti atra na syāt . \\
\noindent \textbf{(P\_1,1.21.2)} atha antavattve kāni prayojanāni . \\
\noindent \textbf{(P\_1,1.21.2)} antavat dvivacanāntapragṛhyatve . antavat dvivacanāntapragṛhyatve prayojanam . \\
\noindent \textbf{(P\_1,1.21.2)} īdūdet dvivacanam pragṛhyam iha eva syāt : pacete* iti pacethe* iti . \\
\noindent \textbf{(P\_1,1.21.2)} khaṭve* iti māle* iti iti atra na syāt . \\
\noindent \textbf{(P\_1,1.21.2)} mit acaḥ antyāt paraḥ . \\
\noindent \textbf{(P\_1,1.21.2)} mit acaḥ antyāt paraḥ prayojanam . \\
\noindent \textbf{(P\_1,1.21.2)} iha eva syāt : kuṇḍāni vanāni . \\
\noindent \textbf{(P\_1,1.21.2)} tāni yāni iti atra na syāt . \\
\noindent \textbf{(P\_1,1.21.2)} acaḥ antyādi ṭi . \\
\noindent \textbf{(P\_1,1.21.2)} acaḥ antyādi ṭi prayojanam . \\
\noindent \textbf{(P\_1,1.21.2)} ṭitaḥ ātmanepadānām ṭeḥ e iti iha eva syāt : kurvāte kurvāthe . \\
\noindent \textbf{(P\_1,1.21.2)} kurute kurve iti atra na syāt . \\
\noindent \textbf{(P\_1,1.21.2)} alaḥ antyasya . \\
\noindent \textbf{(P\_1,1.21.2)} alaḥ antyasya prayojanam . \\
\noindent \textbf{(P\_1,1.21.2)} ataḥ dīrghaḥ yañi supi ca iha eva syāt : ghaṭābhyam , paṭābhyām . \\
\noindent \textbf{(P\_1,1.21.2)} ābhyām iti atra na syāt . \\
\noindent \textbf{(P\_1,1.21.2)} yena vidiḥ tadantatve . \\
\noindent \textbf{(P\_1,1.21.2)} yena vidiḥ tadantatve prayojanam . \\
\noindent \textbf{(P\_1,1.21.2)} acaḥ yat iha eva syāt : ceyam , jeyam . \\
\noindent \textbf{(P\_1,1.21.2)} eyam adhyeyam iti atra na syāt . \\
\noindent \textbf{(P\_1,1.21.2)} ādyantavat ekasmin kāryam bhavati iti atra api siddham bhavati . \\
\noindent \textbf{(P\_1,1.22)} ghasañjñāyām nadītare pratiṣedhaḥ . \\
\noindent \textbf{(P\_1,1.22)} ghasañjñāyām nadītare pratiṣedhaḥ vaktavyaḥ . \\
\noindent \textbf{(P\_1,1.22)} nadyāḥ taraḥ nadītaraḥ iti . \\
\noindent \textbf{(P\_1,1.22)} ghasañjñāyām nadītare apratiṣedhaḥ . \\
\noindent \textbf{(P\_1,1.22)} anarthakaḥ pratiṣedhaḥ apratiṣedhaḥ . \\
\noindent \textbf{(P\_1,1.22)} ghasañjñā kasmāt na bhavati . \\
\noindent \textbf{(P\_1,1.22)} tarabgrahaṇam hi aupadeśikam . \\
\noindent \textbf{(P\_1,1.22)} aupadeśikasya tarapaḥ grahaṇam . \\
\noindent \textbf{(P\_1,1.22)} na ca eṣaḥ upadeśe tarapśabdaḥ . \\
\noindent \textbf{(P\_1,1.22)} kim vaktavyam etat . \\
\noindent \textbf{(P\_1,1.22)} na hi . \\
\noindent \textbf{(P\_1,1.22)} katham anucyamānam gaṃsyate . \\
\noindent \textbf{(P\_1,1.22)} iha hi vyākaraṇe sarveṣu eva sānubandhakeṣu grahaṇeṣu rūpam āśrīyate : yatra etat rūpam iti . \\
\noindent \textbf{(P\_1,1.22)} rūpanirgrahaḥ ca na antareṇa laukikam prayogam . \\
\noindent \textbf{(P\_1,1.22)} tasmin ca laukike prayoge sānubandhakānām prayogaḥ na asti iti kṛtvā dvitīyaḥ prayogaḥ upāsyate . \\
\noindent \textbf{(P\_1,1.22)} kaḥ asau . \\
\noindent \textbf{(P\_1,1.22)} upadeśaḥ nāma . \\
\noindent \textbf{(P\_1,1.22)} na ca eṣaḥ upadeśe tarapśabdaḥ . \\
\noindent \textbf{(P\_1,1.22)} atha vā astu asya ghasañjñā . \\
\noindent \textbf{(P\_1,1.22)} kaḥ doṣaḥ . \\
\noindent \textbf{(P\_1,1.22)} ghādiṣu nadyāḥ hrasvaḥ bhavati iti hrasvatvam prasajyeta . \\
\noindent \textbf{(P\_1,1.22)} samānādhikaraṇeṣu ghādiṣu iti evam tat . \\
\noindent \textbf{(P\_1,1.22)} yadā tarhi sā eva nadī saḥ eva taraḥ tadā prāpnoti . \\
\noindent \textbf{(P\_1,1.22)} strīliṅgeṣu eva ghādiṣu iti evam tat . \\
\noindent \textbf{(P\_1,1.22)} avaśyam ca etat evam vijñeyam . \\
\noindent \textbf{(P\_1,1.22)} samānādhikaraṇeṣu ghādiṣu iti ucyamāne iha prasajyeta mahiṣī rūpam iva brāhmaṇī rūpam iva . \\
\noindent \textbf{(P\_1,1.23.1)} saṅkhyāsañjñāyām saṅkhyāgrahaṇam . \\
\noindent \textbf{(P\_1,1.23.1)} saṅkhyāsañjñāyām saṅkhyāgrahaṇam kartavyam . \\
\noindent \textbf{(P\_1,1.23.1)} bahugaṇvatuḍatayaḥ saṅkhyāsañjñāḥ bhavanti . \\
\noindent \textbf{(P\_1,1.23.1)} saṅkhyā ca saṅkhyāsañjñā bhavati iti vaktavvyam . \\
\noindent \textbf{(P\_1,1.23.1)} kim prayojanam . \\
\noindent \textbf{(P\_1,1.23.1)} saṅkhyāsampratyayārtham . \\
\noindent \textbf{(P\_1,1.23.1)} ekādikāyāḥ saṅkhyāyāḥ saṅkhyāpradeśeṣu saṅkhyā iti eṣaḥ sampratyayaḥ yathā syāt . \\
\noindent \textbf{(P\_1,1.23.1)} nanu ca ekādikā saṅkhyā loke saṅkhyā iti pratītā . \\
\noindent \textbf{(P\_1,1.23.1)} tena asyāḥ saṅkhyāpradeśeṣu saṅkhyāsampratyayaḥ bhaviṣyati . \\
\noindent \textbf{(P\_1,1.23.1)} evam api kartavyam . \\
\noindent \textbf{(P\_1,1.23.1)} itarathā hi asampratyayaḥ akṛtrimatvāt yathā loke . \\
\noindent \textbf{(P\_1,1.23.1)} akriyamāṇe hi saṅkhyāgrahaṇe ekādikāyāḥ saṅkhyāyāḥ saṅkhyā iti sampratyayaḥ na syāt . \\
\noindent \textbf{(P\_1,1.23.1)} kim kāraṇam . \\
\noindent \textbf{(P\_1,1.23.1)} akṛtrimatvāt . \\
\noindent \textbf{(P\_1,1.23.1)} bahvādīnām kṛtrimā sañjñā . \\
\noindent \textbf{(P\_1,1.23.1)} kṛtrimākṛtrimayoḥ kṛtrime kāryasampratyayaḥ bhavati yathā loke . \\
\noindent \textbf{(P\_1,1.23.1)} tat yathā loke gopālakam ānaya kaṭajakam ānaya iti yasya eṣā sañjñā bhavati saḥ ānīyate na yaḥ gāḥ pālayati yaḥ vā kaṭe jātaḥ . \\
\noindent \textbf{(P\_1,1.23.1)} yadi tarhi kṛtrimākṛtrimayoḥ kṛtrime sampratyayaḥ bhavati nadīpaurṇamāsyāgrahāyaṇībhyaḥ iti atra api prasajyeta . \\
\noindent \textbf{(P\_1,1.23.1)} paurṇamāsyāgrahāyaṇīgrahaṇasāmarthyāt na bhaviṣyati . \\
\noindent \textbf{(P\_1,1.23.1)} tadviśeṣebhyaḥ tarhi prāpnoti : gaṅgā yamunā iti . \\
\noindent \textbf{(P\_1,1.23.1)} evam tarhi ācāryapravṛttiḥ jñāpayati na tadviśeṣebhyaḥ bhavati iti yat ayam vipāṭśabdam śaratprabhṛtiṣu paṭhati . \\
\noindent \textbf{(P\_1,1.23.1)} iha tarhi prāpnoti : nadībhiḥ ca iti . \\
\noindent \textbf{(P\_1,1.23.1)} bahuvacananirdeśāt na bhaviṣyati . \\
\noindent \textbf{(P\_1,1.23.1)} svarūpavidhiḥ tarhi prāpnoti . \\
\noindent \textbf{(P\_1,1.23.1)} bahuvacananirdeśāt eva na bhaviṣyati . \\
\noindent \textbf{(P\_1,1.23.1)} evam na ca idam akṛtam bhavati kṛtrimākṛtrimayoḥ kṛtrime sampratyayaḥ iti na ca kaḥ cit doṣaḥ . \\
\noindent \textbf{(P\_1,1.23.1)} uttarārtham ca . \\
\noindent \textbf{(P\_1,1.23.1)} uttarārtham ca saṅkhyāgrahaṇam kartavyam . \\
\noindent \textbf{(P\_1,1.23.1)} ṣṇāntā ṣaṭ . \\
\noindent \textbf{(P\_1,1.23.1)} ṣakāranakārāntāyāḥ saṅkhyāyāḥ ṣaṭsañjñā yathā syāt . \\
\noindent \textbf{(P\_1,1.23.1)} iha mā bhūt : pāmānaḥ , vipruṣaḥ iti . \\
\noindent \textbf{(P\_1,1.23.1)} ihārthena tāvat na arthaḥ saṅkhyāgrahaṇena . \\
\noindent \textbf{(P\_1,1.23.1)} nanu ca uktam itarathā hi asampratyayaḥ akṛtrimatvāt yathā loke iti . \\
\noindent \textbf{(P\_1,1.23.1)} na eṣaḥ doṣaḥ . \\
\noindent \textbf{(P\_1,1.23.1)} arthāt prakaraṇāt vā loke kṛtrimākṛtrimayoḥ kṛtrime sampratyayaḥ bhavati . \\
\noindent \textbf{(P\_1,1.23.1)} arthaḥ vā asya evaṃsaṅjñakena bhavati prakṛtam vā tatra bhavati idam evaṃsaṅjñakena kartavyam iti. ātaḥ ca arthāt prakaraṇāt vā . \\
\noindent \textbf{(P\_1,1.23.1)} aṅga hi bhavān grāmyam pāṃsurapādam aprakaraṇajñam āgatam bravītu gopālakam ānaya kaṭajakam ānaya iti . \\
\noindent \textbf{(P\_1,1.23.1)} ubhayagatiḥ tasya bhavati sādhīyaḥ vā yaṣṭihastam gamiṣyati . \\
\noindent \textbf{(P\_1,1.23.1)} yathā eva tarhi arthāt prakaraṇāt vā loke kṛtrimākṛtrimayoḥ kṛtrime sampratyayaḥ bhavati evam iha api prāpnoti . \\
\noindent \textbf{(P\_1,1.23.1)} jānāti hi asau bahvādīnām iyam sañjñā kṛtā iti . \\
\noindent \textbf{(P\_1,1.23.1)} na yathā loke tathā vyākaraṇe . \\
\noindent \textbf{(P\_1,1.23.1)} ubhayagatiḥ punaḥ iha bhavati . \\
\noindent \textbf{(P\_1,1.23.1)} anyatra api na avaśyam iha eva . \\
\noindent \textbf{(P\_1,1.23.1)} tat yathā : kartuḥ īpsitatamam karma iti kṛtrimā sañjñā . \\
\noindent \textbf{(P\_1,1.23.1)} karmapradeśeṣu ca ubhayagatiḥ bhavati . \\
\noindent \textbf{(P\_1,1.23.1)} karmaṇi dvitīyā iti kṛtrimasya grahaṇam kartari karmavyatihāre iti akṛtrimasya . \\
\noindent \textbf{(P\_1,1.23.1)} tathā sādhakatamam karaṇam iti kṛtrimā karaṇasañjñā . \\
\noindent \textbf{(P\_1,1.23.1)} karaṇapradeśeṣu ca ubhayagatiḥ bhavati . \\
\noindent \textbf{(P\_1,1.23.1)} kartṛkaraṇayoḥ tṛtīyā iti kṛtrimasya grahaṇam śabdavairakalahābhrakaṇvameghebhyaḥ karaṇe iti akṛtrimasya . \\
\noindent \textbf{(P\_1,1.23.1)} tathā ādhāraḥ adhikaraṇam iti kṛtrimā adhikaraṇasañjñā . \\
\noindent \textbf{(P\_1,1.23.1)} adhikaraṇepradeśeṣu ca ubhayagatiḥ bhavati . \\
\noindent \textbf{(P\_1,1.23.1)} saptamī adhikaraṇe ca iti kṛtrimasya grahaṇam vipratiṣiddham ca anadhikaraṇavāci iti akṛtrimasya . \\
\noindent \textbf{(P\_1,1.23.1)} atha vā na idam sañjñākaraṇam . \\
\noindent \textbf{(P\_1,1.23.1)} tadvadatideśaḥ ayam : bahugaṇavatuḍatayaḥ saṅkhyāvat bhavanti iti . \\
\noindent \textbf{(P\_1,1.23.1)} saḥ tarhi vatinirdeśaḥ kartavyaḥ . \\
\noindent \textbf{(P\_1,1.23.1)} na hi antareṇa vatim atideśaḥ gamyate . \\
\noindent \textbf{(P\_1,1.23.1)} antareṇa api vatim atideśaḥ gamyate . \\
\noindent \textbf{(P\_1,1.23.1)} tat yathā : eṣaḥ brahmadattaḥ . \\
\noindent \textbf{(P\_1,1.23.1)} abrahmadattam brahmadattaḥ iti āha . \\
\noindent \textbf{(P\_1,1.23.1)} te manyāmahe : brahmadattavat ayam bhavati iti . \\
\noindent \textbf{(P\_1,1.23.1)} evam iha api asaṅkhyām saṅkhyā iti āha . \\
\noindent \textbf{(P\_1,1.23.1)} saṅkhyāvat iti gamyate . \\
\noindent \textbf{(P\_1,1.23.1)} atha vā ācāryapravṛttiḥ jñāpayati bhavati ekādikāyāḥ saṅkhyāyāḥ saṅkhyāpradeśeṣu saṅkhyāsampratyayaḥ iti yat ayam saṅkhyāyāḥ atiśadantāyāḥ kan iti tiśadantāyāḥ pratiṣedham śāsti . \\
\noindent \textbf{(P\_1,1.23.1)} katham kṛtvā jñāpakam . \\
\noindent \textbf{(P\_1,1.23.1)} na hi kṛtrimā tyantā śadantā vā saṅkhyā asti . \\
\noindent \textbf{(P\_1,1.23.1)} nanu ca iyam asti ḍatiḥ . \\
\noindent \textbf{(P\_1,1.23.1)} yat tarhi śadantāyāḥ pratiṣedham śāsti . \\
\noindent \textbf{(P\_1,1.23.1)} yat ca api tyantāyāḥ pratiṣedham śāsti . \\
\noindent \textbf{(P\_1,1.23.1)} nanu ca uktam ḍatyartham etat syāt iti . \\
\noindent \textbf{(P\_1,1.23.1)} arthavadgrahaṇe na anarthakasya iti arthavataḥ tiśabdasya grahaṇam . \\
\noindent \textbf{(P\_1,1.23.1)} na ca ḍateḥ tiśabdaḥ arthavān. atha vā mahatī iyam sañjñā kriyate . \\
\noindent \textbf{(P\_1,1.23.1)} sañjñā ca nāma yataḥ na laghīyaḥ . \\
\noindent \textbf{(P\_1,1.23.1)} kutaḥ etat . \\
\noindent \textbf{(P\_1,1.23.1)} laghvartham hi sañjñākaraṇam . \\
\noindent \textbf{(P\_1,1.23.1)} tatra mahatyāḥ sañjñāyāḥ karaṇe etat prayojanam anvarthasañjñā yatha vijñāyeta . \\
\noindent \textbf{(P\_1,1.23.1)} saṅkhyāyate anayā saṅkhyā iti . \\
\noindent \textbf{(P\_1,1.23.1)} ekādikayā ca api saṅkhyāyate . \\
\noindent \textbf{(P\_1,1.23.1)} uttarārthena ca api na arthaḥ saṅkhyāgrahaṇena . \\
\noindent \textbf{(P\_1,1.23.1)} idam prakṛtam anuvartiṣyate . \\
\noindent \textbf{(P\_1,1.23.1)} idam vai sañjñārtham uttaratra ca sañjñiviśeṣaṇārthaḥ . \\
\noindent \textbf{(P\_1,1.23.1)} na ca anyārtham prakṛtam anyārtham bhavati . \\
\noindent \textbf{(P\_1,1.23.1)} na khalu api anyat prakṛtam anuvartanāt anyat bhavati . \\
\noindent \textbf{(P\_1,1.23.1)} na hi godhā sarpantī sarpaṇāt ahiḥ bhavati . \\
\noindent \textbf{(P\_1,1.23.1)} yat tāvat ucyate na ca anyārtham prakṛtam anyārtham bhavati iti anyārtham api prakṛtam anyārtham bhavati . \\
\noindent \textbf{(P\_1,1.23.1)} tat yathā : śālyartham kulyāḥ praṇīyante tābhyaḥ ca pāṇīyam pīyate upaśpṛśyate ca śālayaḥ ca bhāvyante . \\
\noindent \textbf{(P\_1,1.23.1)} yad api ucyate na khalu api anyat prakṛtam anuvartanāt anyat bhavati . \\
\noindent \textbf{(P\_1,1.23.1)} na hi godhā sarpantī sarpaṇāt ahiḥ bhavati iti . \\
\noindent \textbf{(P\_1,1.23.1)} bhavet dravyeṣu etat evam syāt . \\
\noindent \textbf{(P\_1,1.23.1)} śabdaḥ tu khalu yena yena viśeṣeṇa abhisambadhyate tasya tasya viśeṣakaḥ bhavati . \\
\noindent \textbf{(P\_1,1.23.1)} atha vā sāpekṣaḥ ayam nirdeśaḥ kriyate na ca anyat kim cit apekṣyam asti . \\
\noindent \textbf{(P\_1,1.23.1)} te saṅkhyām eva apekṣiṣyāmahe . \\
\noindent \textbf{(P\_1,1.23.2)} adhyardhagrahaṇam ca samāsakanvidhyartham . \\
\noindent \textbf{(P\_1,1.23.2)} adhyardhagrahaṇam ca kartavyam . \\
\noindent \textbf{(P\_1,1.23.2)} kim prayojanam . \\
\noindent \textbf{(P\_1,1.23.2)} samāsakanvidhyartham . \\
\noindent \textbf{(P\_1,1.23.2)} samāsavidhyartham kandvidhyartham ca . \\
\noindent \textbf{(P\_1,1.23.2)} samāsavidhyartham tāvat : adhyardhaśūrpam . \\
\noindent \textbf{(P\_1,1.23.2)} kanvidhyartham : adhyardhakam . \\
\noindent \textbf{(P\_1,1.23.2)} luki ca agrahaṇam . \\
\noindent \textbf{(P\_1,1.23.2)} luki ca adhyardhagrahaṇam na kartavyam bhavati : adhyardhapūrvadvigoḥ luk asañjñāyām iti . \\
\noindent \textbf{(P\_1,1.23.2)} dvigoḥ iti eva siddham . \\
\noindent \textbf{(P\_1,1.23.2)} ardhapūrvapadaḥ ca pūraṇapratyayāntaḥ . \\
\noindent \textbf{(P\_1,1.23.2)} ardhapūrvapadaḥ ca pūraṇapratyayāntaḥ saṅkhyāsañjñaḥ bhavati iti vaktavyam . \\
\noindent \textbf{(P\_1,1.23.2)} kim prayojanam . \\
\noindent \textbf{(P\_1,1.23.2)} samāsakanvidhyartham . \\
\noindent \textbf{(P\_1,1.23.2)} samāsavidhyartham kandvidhyartham ca . \\
\noindent \textbf{(P\_1,1.23.2)} samāsavidhyartham tāvat : ardhapañcamaśūrpam . \\
\noindent \textbf{(P\_1,1.23.2)} kanvidhyartham : ardhapañcamakam . \\
\noindent \textbf{(P\_1,1.23.2)} adhikagrahaṇam ca aluki samāsottarapadavṛddhyartham . \\
\noindent \textbf{(P\_1,1.23.2)} adhikagrahaṇam ca aluki kartavyam . \\
\noindent \textbf{(P\_1,1.23.2)} kim prayojanam . \\
\noindent \textbf{(P\_1,1.23.2)} samāsottarapadavṛddhyartham . \\
\noindent \textbf{(P\_1,1.23.2)} samāsavṛddhyartham uttaravṛddhyartam ca . \\
\noindent \textbf{(P\_1,1.23.2)} samāsavṛddhyartham tāvat : adhikaṣāṣṭhikaḥ , adhikasāptatikaḥ . \\
\noindent \textbf{(P\_1,1.23.2)} uttarapadavṛddhyartham adhikaṣāṣṭhikaḥ , adhikasāptatikaḥ . \\
\noindent \textbf{(P\_1,1.23.2)} aluki iti kim artham . \\
\noindent \textbf{(P\_1,1.23.2)} adhikaṣāṣṭhikaḥ , adhikasāptatikaḥ . \\
\noindent \textbf{(P\_1,1.23.2)} bahuvrīhau ca agrahaṇam . \\
\noindent \textbf{(P\_1,1.23.2)} bahuvrīhau ca adhikaśabdasya grahaṇam na kartavyam bhavati : saṅkhyayā avyayāsannādūrādhikasaṅkhyāḥ saṅkhyeye iti . \\
\noindent \textbf{(P\_1,1.23.2)} saṅkhyā iti eva siddham . \\
\noindent \textbf{(P\_1,1.23.2)} bahvādīnām agrahaṇam . \\
\noindent \textbf{(P\_1,1.23.2)} bahvādīnām grahaṇam śakyam akartum . \\
\noindent \textbf{(P\_1,1.23.2)} kena idānīm saṅkhyāpradeśeṣu saṅkhyasampratyayaḥ bhaviṣyati . \\
\noindent \textbf{(P\_1,1.23.2)} jñāpakāt siddham . \\
\noindent \textbf{(P\_1,1.23.2)} kim jñāpakam . \\
\noindent \textbf{(P\_1,1.23.2)} yat ayam vatoḥ iṭ vā iti saṅkhyāyāḥ vihitasya kanaḥ vatvantāt iṭam śāsti . \\
\noindent \textbf{(P\_1,1.23.2)} vatoḥ eva tat jñāpakam syāt . \\
\noindent \textbf{(P\_1,1.23.2)} na iti āha . \\
\noindent \textbf{(P\_1,1.23.2)} yogāpekṣam jñāpakam . \\
\noindent \textbf{(P\_1,1.24)} ṣaṭsañjñāyām upadeśavacanam . \\
\noindent \textbf{(P\_1,1.24)} ṣaṭsañjñāyām upadeśagrahaṇam kartavyam . \\
\noindent \textbf{(P\_1,1.24)} upadeśe ṣakāranakārāntā saṅkhyā ṣaṭsañjñā bhavati iti vaktavyam . \\
\noindent \textbf{(P\_1,1.24)} kim prayojanam . \\
\noindent \textbf{(P\_1,1.24)} śatādyaṣṭanoḥ numnuḍartham . \\
\noindent \textbf{(P\_1,1.24)} śatāni sahasrāṇi . \\
\noindent \textbf{(P\_1,1.24)} numi kṛte ṣṇāntā ṣaṭ iti ṣaṭsañjñā prāpnoti . \\
\noindent \textbf{(P\_1,1.24)} upadeśagrahaṇāt na bhavati . \\
\noindent \textbf{(P\_1,1.24)} aṣṭānām iti atra ātve kṛte ṣaṭsañjñā na prāpnoti . \\
\noindent \textbf{(P\_1,1.24)} upadeśagrahaṇāt bhavati . \\
\noindent \textbf{(P\_1,1.24)} uktam vā . \\
\noindent \textbf{(P\_1,1.24)} kim uktam . \\
\noindent \textbf{(P\_1,1.24)} iha tāvat śatāni sahasrāṇi iti . \\
\noindent \textbf{(P\_1,1.24)} sannipātalakṣaṇaḥ vidhiḥ animittam tadvighātasya iti . \\
\noindent \textbf{(P\_1,1.24)} aṣṭanaḥ api uktam . \\
\noindent \textbf{(P\_1,1.24)} kim uktam . \\
\noindent \textbf{(P\_1,1.24)} aṣṭanaḥ dīrghagrahaṇam ṣaṭsañjñājñāpakam ākārāntasya nuḍartham iti . \\
\noindent \textbf{(P\_1,1.24)} atha vā ākāraḥ api atra nirdiśyate . \\
\noindent \textbf{(P\_1,1.24)} ṣakārāntā nakārāntā ākārāntā ca saṅkhyā ṣaṭsañjñā bhavati iti . \\
\noindent \textbf{(P\_1,1.24)} iha api tarhi prāpnoti : sadhamadhaḥ dyumnaḥ ekāḥ taḥ ekāḥ iti . \\
\noindent \textbf{(P\_1,1.24)} na eṣaḥ doṣaḥ . \\
\noindent \textbf{(P\_1,1.24)} ekaśabdaḥ ayam bahvarthaḥ . \\
\noindent \textbf{(P\_1,1.24)} asti eva saṅkhyāpadam . \\
\noindent \textbf{(P\_1,1.24)} tat yathā : ekaḥ , dvau , bahavaḥ iti . \\
\noindent \textbf{(P\_1,1.24)} asti asahāyavācī . \\
\noindent \textbf{(P\_1,1.24)} tat yathā : ekāgnayaḥ , ekahalāni , ekākibhiḥ kṣudrakaiḥ jitam iti . \\
\noindent \textbf{(P\_1,1.24)} asahāyaiḥ iti arthaḥ . \\
\noindent \textbf{(P\_1,1.24)} asti anyārthe vartate . \\
\noindent \textbf{(P\_1,1.24)} tat yathā : prajam ekā rakṣati urjam ekā iti . \\
\noindent \textbf{(P\_1,1.24)} anyā iti arthaḥ . \\
\noindent \textbf{(P\_1,1.24)} sadhamadaḥ dyumnaḥ ekāḥ taḥ . \\
\noindent \textbf{(P\_1,1.24)} anyāḥ iti arthaḥ . \\
\noindent \textbf{(P\_1,1.24)} tat yaḥ anyārthe vartate tasya eṣaḥ prayogaḥ . \\
\noindent \textbf{(P\_1,1.24)} iha tarhi prāpnoti : dvabhyām iṣṭaye viṃśatya ca iti . \\
\noindent \textbf{(P\_1,1.24)} evam tarhi saptame yogavibhāgaḥ kariṣyate . \\
\noindent \textbf{(P\_1,1.24)} aṣṭābhyaḥ auś . \\
\noindent \textbf{(P\_1,1.24)} tataḥ ṣaḍbhyaḥ : ṣaḍbhyaḥ ca yat uktam aṣṭābhyaḥ api tat bhavati . \\
\noindent \textbf{(P\_1,1.24)} tataḥ luk : luk ca bhavati ṣaḍbhyaḥ iti . \\
\noindent \textbf{(P\_1,1.24)} atha vā upariṣṭāt yogavibhāgaḥ kariṣyate . \\
\noindent \textbf{(P\_1,1.24)} aṣṭanaḥ ā vibhaktau . \\
\noindent \textbf{(P\_1,1.24)} tataḥ rāyaḥ : rāyaḥ ca vibhaktau ākārādeśaḥ bhavati . \\
\noindent \textbf{(P\_1,1.24)} hali iti ubhayoḥ śeṣaḥ . \\
\noindent \textbf{(P\_1,1.24)} yadi evam priyāṣṭau priyāṣṭāḥ iti na sidhyati priyāṣṭānau priyāṣṭānaḥ iti ca prāpnoti . \\
\noindent \textbf{(P\_1,1.24)} yathālakṣaṇam aprayukte . \\
\noindent \textbf{(P\_1,1.25)} idam ḍatigrahaṇam dviḥ kriyate saṅkhyāsañjñāyām ṣaṭsañjñāyām ca . \\
\noindent \textbf{(P\_1,1.25)} ekam śakyam akartum . \\
\noindent \textbf{(P\_1,1.25)} katham . \\
\noindent \textbf{(P\_1,1.25)} yadi tāvat saṅkhyāsañjñāyām kriyate ṣaṭsañjñāyām na kariṣyate . \\
\noindent \textbf{(P\_1,1.25)} katham . \\
\noindent \textbf{(P\_1,1.25)} ṣṇāntā ṣaṭ iti atra ḍati iti anuvartiṣyate . \\
\noindent \textbf{(P\_1,1.25)} atha ṣaṭsañjñāyām kriyate saṅkhyāsañjñāyām na kariṣyate . \\
\noindent \textbf{(P\_1,1.25)} ḍati ca iti atra saṅkhyāsañjñā anuvartiṣyate . \\
\noindent \textbf{(P\_1,1.26)} niṣṭhāsañjñāyām samānaśabdapratiṣedhaḥ . \\
\noindent \textbf{(P\_1,1.26)} niṣṭhāsañjñāyām samānaśabdānām pratiṣedhaḥ kartavyaḥ . \\
\noindent \textbf{(P\_1,1.26)} lotaḥ gartaḥ iti . \\
\noindent \textbf{(P\_1,1.26)} niṣṭhāsañjñāyām samānaśabdāpratiṣedhaḥ . \\
\noindent \textbf{(P\_1,1.26)} niṣṭhāsañjñāyām samānaśabdapratiṣedhaḥ . \\
\noindent \textbf{(P\_1,1.26)} anarthakaḥ pratiṣedhaḥ apratiṣedhaḥ . \\
\noindent \textbf{(P\_1,1.26)} niṣṭhāsañjñā kasmāt na bhavati . \\
\noindent \textbf{(P\_1,1.26)} anubandhaḥ anyatvakaraḥ . \\
\noindent \textbf{(P\_1,1.26)} anubandhaḥ kriyate . \\
\noindent \textbf{(P\_1,1.26)} saḥ anyatvam kariṣyati . \\
\noindent \textbf{(P\_1,1.26)} anubandhaḥ anyatvakaraḥ iti cet na lopāt . \\
\noindent \textbf{(P\_1,1.26)} anubandhaḥ anyatvakaraḥ iti cet tat na . \\
\noindent \textbf{(P\_1,1.26)} kim kāraṇam . \\
\noindent \textbf{(P\_1,1.26)} lopāt . \\
\noindent \textbf{(P\_1,1.26)} lupyate atra anubandhaḥ . \\
\noindent \textbf{(P\_1,1.26)} lupte atra anubandhe na anyatvam bhaviṣyati . \\
\noindent \textbf{(P\_1,1.26)} tat yathā : katarat devadattasya gṛham . \\
\noindent \textbf{(P\_1,1.26)} adaḥ yatra asau kākaḥ iti . \\
\noindent \textbf{(P\_1,1.26)} utpatite kāke naṣṭam tat gṛham bhavati . \\
\noindent \textbf{(P\_1,1.26)} evam iha api lupte anubandhe naṣṭaḥ pratyayaḥ bhavati . \\
\noindent \textbf{(P\_1,1.26)} yadi api lupyate jānāti tu asau sānubandhakasya iyam sañjñā kṛtā iti . \\
\noindent \textbf{(P\_1,1.26)} tat yathā itaratra api : katarat devadattasya gṛham . \\
\noindent \textbf{(P\_1,1.26)} adaḥ yatra asau kākaḥ iti . \\
\noindent \textbf{(P\_1,1.26)} utpatite kāke yadi api naṣṭam tat gṛham bhavati antataḥ tam uddeśam jānāti . \\
\noindent \textbf{(P\_1,1.26)} siddhaviparyāsaḥ ca . \\
\noindent \textbf{(P\_1,1.26)} siddhaḥ ca viparyāsaḥ . \\
\noindent \textbf{(P\_1,1.26)} yadi api jānāti sandehaḥ tasya bhavati : ayam saḥ taśabdaḥ lotaḥ gartaḥ iti ayam saḥ taśabdaḥ lūnaḥ gīrṇaḥ iti . \\
\noindent \textbf{(P\_1,1.26)} tat yathā itaratra api : katarat devadattasya gṛham . \\
\noindent \textbf{(P\_1,1.26)} adaḥ yatra asau kākaḥ iti . \\
\noindent \textbf{(P\_1,1.26)} utpatite kāke yadi api naṣṭam tat gṛham bhavati antataḥ tam uddeśam jānāti . \\
\noindent \textbf{(P\_1,1.26)} sandehaḥ tu tasya bhavati : idam tat gṛham idam tat gṛham iti . \\
\noindent \textbf{(P\_1,1.26)} evam tarhi . \\
\noindent \textbf{(P\_1,1.26)} kārakakālaviśeṣāt siddham . \\
\noindent \textbf{(P\_1,1.26)} kārakakālaviśeṣau upādeyau . \\
\noindent \textbf{(P\_1,1.26)} bhūte yaḥ taśabdaḥ kartari karmaṇi bhāve ca iti . \\
\noindent \textbf{(P\_1,1.26)} tat yathā itaratra api . \\
\noindent \textbf{(P\_1,1.26)} yaḥ eṣaḥ manuṣyaḥ prekṣāpūrvakārī bhavati saḥ adhruveṇa nimittena dhruvam nimittam upādatte vedikām puṇḍarīkam vā . \\
\noindent \textbf{(P\_1,1.26)} evam api prākīrṣṭa iti atra prāpnoti . \\
\noindent \textbf{(P\_1,1.26)} luṅi sijādidarśanāt . \\
\noindent \textbf{(P\_1,1.26)} luṅi sijādidarśanāt na bhaviṣyati . \\
\noindent \textbf{(P\_1,1.26)} yatra tarhi sijādayaḥ na dṛśyante prābhitta iti . \\
\noindent \textbf{(P\_1,1.26)} dṛśyante atra api sijādayaḥ . \\
\noindent \textbf{(P\_1,1.26)} kim vaktavyam etat . \\
\noindent \textbf{(P\_1,1.26)} na hi . \\
\noindent \textbf{(P\_1,1.26)} katham anucyamānam gaṃsyate . \\
\noindent \textbf{(P\_1,1.26)} yathā eva ayam anupadiṣṭān kārakakālaviśeṣān avagacchati evam etat api avagantum arhati : yatra sijādayaḥ na iti . \\
\noindent \textbf{(P\_1,1.27.1)} sarvādīni iti kaḥ ayam samāsaḥ . \\
\noindent \textbf{(P\_1,1.27.1)} bahuvrīhiḥ iti āha . \\
\noindent \textbf{(P\_1,1.27.1)} kaḥ asya vigrahaḥ . \\
\noindent \textbf{(P\_1,1.27.1)} sarvaśabdaḥ ādiḥ yeṣām tāni imāni iti . \\
\noindent \textbf{(P\_1,1.27.1)} yadi evam sarvaśabdasya sarvanāmasañjñā na prāpnoti . \\
\noindent \textbf{(P\_1,1.27.1)} kim kāraṇam . \\
\noindent \textbf{(P\_1,1.27.1)} anyapadārthatvāt bahuvrīheḥ . \\
\noindent \textbf{(P\_1,1.27.1)} bahuvrīhiḥ ayam anyapadārthe vartate . \\
\noindent \textbf{(P\_1,1.27.1)} tena yat anyat sarvaśabdāt tasya sarvanāmasñjñā prāpnoti . \\
\noindent \textbf{(P\_1,1.27.1)} tat yathā citraguḥ ānīyatām iti ukte yasya tāḥ gāvaḥ bhavanti sa ānīyate na gāvaḥ . \\
\noindent \textbf{(P\_1,1.27.1)} na eṣaḥ doṣaḥ . \\
\noindent \textbf{(P\_1,1.27.1)} bhavati bahuvrīhau tadguṇasaṃvijñānam api . \\
\noindent \textbf{(P\_1,1.27.1)} tat yathā : citravāsam ānaya . \\
\noindent \textbf{(P\_1,1.27.1)} lohitoṣṇīṣāḥ ṛtvijaḥ pracaranti . \\
\noindent \textbf{(P\_1,1.27.1)} tadguṇaḥ ānīyate tadguṇāḥ ca pracaranti . \\
\noindent \textbf{(P\_1,1.27.2)} iha sarvanāmāni iti pūrvapadāt sañjñāyām agaḥ iti ṇatvam prāpnoti . \\
\noindent \textbf{(P\_1,1.27.2)} tasya pratiṣedhaḥ vaktavyaḥ . \\
\noindent \textbf{(P\_1,1.27.2)} sarvanāmasañjñāyām nipātanāt ṇatvābhāvaḥ . \\
\noindent \textbf{(P\_1,1.27.2)} sarvanāmasañjñāyām nipātanāt ṇatvam na bhaviṣyati . \\
\noindent \textbf{(P\_1,1.27.2)} kim etat nipātanam nāma . \\
\noindent \textbf{(P\_1,1.27.2)} atha kaḥ pratiṣedhaḥ nāma . \\
\noindent \textbf{(P\_1,1.27.2)} aviśeṣeṇa kim cit uktvā viśeṣeṇa na iti ucyate . \\
\noindent \textbf{(P\_1,1.27.2)} tatra vyaktam ācāryasya abhiprāyaḥ gamyate : idam na bhavati iti . \\
\noindent \textbf{(P\_1,1.27.2)} nipātanam api evañjātīyakam eva . \\
\noindent \textbf{(P\_1,1.27.2)} aviśeṣeṇa ṇatvam uktvā viśeṣeṇa nipātanam kriyate . \\
\noindent \textbf{(P\_1,1.27.2)} tatra vyaktam ācāryasya abhiprāyaḥ gamyate : idam na bhavati iti . \\
\noindent \textbf{(P\_1,1.27.2)} nanu ca nipātanāt ca aṇatvam syāt yathāprāptam ca ṇatvam . \\
\noindent \textbf{(P\_1,1.27.2)} kim anye api evam vidhayaḥ bhavanti . \\
\noindent \textbf{(P\_1,1.27.2)} iha ikaḥ yaṇ aci iti vacanāt ca yaṇ syāt yathāprāptaḥ ca ik śrūyeta . \\
\noindent \textbf{(P\_1,1.27.2)} na eṣaḥ doṣaḥ . \\
\noindent \textbf{(P\_1,1.27.2)} asti atra viśeṣaḥ . \\
\noindent \textbf{(P\_1,1.27.2)} ṣaṣṭhyā atra nirdeśaḥ kriyate . \\
\noindent \textbf{(P\_1,1.27.2)} ṣaṣṭhī ca punaḥ sthāninam nivartayati . \\
\noindent \textbf{(P\_1,1.27.2)} iha tarhi : kartari śap divādibhyaḥ śyan iti vacanāt ca śyan syāt yathāprāptaḥ ca śap śrūyeta . \\
\noindent \textbf{(P\_1,1.27.2)} na eṣaḥ doṣaḥ . \\
\noindent \textbf{(P\_1,1.27.2)} śabādeśāḥ śyanādayaḥ kariṣyante . \\
\noindent \textbf{(P\_1,1.27.2)} tat tarhi śapaḥ grahaṇam kartavyam . \\
\noindent \textbf{(P\_1,1.27.2)} na kartavyam . \\
\noindent \textbf{(P\_1,1.27.2)} prakṛtam anuvartate . \\
\noindent \textbf{(P\_1,1.27.2)} kva prakṛtam . \\
\noindent \textbf{(P\_1,1.27.2)} kartari śap iti . \\
\noindent \textbf{(P\_1,1.27.2)} tat vai prathamānirdiṣṭam ṣaṣṭhīnirdiṣṭena ca iha arthaḥ . \\
\noindent \textbf{(P\_1,1.27.2)} divādibhyaḥ iti eṣā pañcamī śap iti prathamāyāḥ ṣaṣṭhīm prakalpayiṣyati : tasmāt iti uttarasya . \\
\noindent \textbf{(P\_1,1.27.2)} pratyayavidhiḥ ayam . \\
\noindent \textbf{(P\_1,1.27.2)} na ca pratyayavidhau pañcamyaḥ prakalpikāḥ bhavanti . \\
\noindent \textbf{(P\_1,1.27.2)} na ayam pratyayavidhiḥ . \\
\noindent \textbf{(P\_1,1.27.2)} vihitaḥ pratyayaḥ prakṛtaḥ ca anuvartate . \\
\noindent \textbf{(P\_1,1.27.2)} iha tarhi : avyayasarvanāmnām akac prāk ṭeḥ iti . \\
\noindent \textbf{(P\_1,1.27.2)} vacanāt ca akac syāt yathāprāptaḥ ca kaḥ śrūyeta . \\
\noindent \textbf{(P\_1,1.27.2)} na eṣaḥ doṣaḥ . \\
\noindent \textbf{(P\_1,1.27.2)} na aprāpte hi ke akac ārabhyate . \\
\noindent \textbf{(P\_1,1.27.2)} saḥ bādhakaḥ bhaviṣyati . \\
\noindent \textbf{(P\_1,1.27.2)} nipātanam api evañjātīyakam eva . \\
\noindent \textbf{(P\_1,1.27.2)} na aprāpte ṇatve nipātanam ārabhyate . \\
\noindent \textbf{(P\_1,1.27.2)} tat bādhakam bhaviṣyati . \\
\noindent \textbf{(P\_1,1.27.2)} yadi tarhi nipātanāni api evañjātīyakāni bhavanti samaḥ tate doṣaḥ bhavati . \\
\noindent \textbf{(P\_1,1.27.2)} iha anye vaiyākaraṇāḥ samaḥ tate vibhāṣā lopam ārabhante : samaḥ hi tatayoḥ vā iti . \\
\noindent \textbf{(P\_1,1.27.2)} satatam , santatam , sahitam , saṃhitam iti . \\
\noindent \textbf{(P\_1,1.27.2)} iha punaḥ bhavān nipātanāt ca malopam icchati aparasparāḥ kriyāsātatye iti yathāprāptam ca alopam santatam iti . \\
\noindent \textbf{(P\_1,1.27.2)} etat na sidhyati . \\
\noindent \textbf{(P\_1,1.27.2)} kartavyaḥ atra yatnaḥ . \\
\noindent \textbf{(P\_1,1.27.2)} bādhakāni eva hi nipātanāni bhavanti . \\
\noindent \textbf{(P\_1,1.27.3)} sañjñopasarjanapratiṣedhaḥ . \\
\noindent \textbf{(P\_1,1.27.3)} sañjñopasarjanībhūtānām sarvādīnām pratiṣedhaḥ vaktavyaḥ . \\
\noindent \textbf{(P\_1,1.27.3)} sarvaḥ nāma kaḥ cit . \\
\noindent \textbf{(P\_1,1.27.3)} tasmai sarvāya dehi . \\
\noindent \textbf{(P\_1,1.27.3)} atisarvāya dehi . \\
\noindent \textbf{(P\_1,1.27.3)} saḥ katham kartavyaḥ . \\
\noindent \textbf{(P\_1,1.27.3)} pāṭhāt paryudāsaḥ paṭhitānām sañjñākaraṇam . \\
\noindent \textbf{(P\_1,1.27.3)} pāṭhāt eva paryudāsaḥ kartavyaḥ . \\
\noindent \textbf{(P\_1,1.27.3)} śuddhānām paṭhitānām sañjñā kartavyā . \\
\noindent \textbf{(P\_1,1.27.3)} sarvādīni sarvanāmasañjñāni bhavanti . \\
\noindent \textbf{(P\_1,1.27.3)} sañjñopasarjanībhūtāni na sarvādīni . \\
\noindent \textbf{(P\_1,1.27.3)} kim aviśeṣeṇa . \\
\noindent \textbf{(P\_1,1.27.3)} na iti āha . \\
\noindent \textbf{(P\_1,1.27.3)} viśeṣeṇa ca . \\
\noindent \textbf{(P\_1,1.27.3)} kim prayojanam . \\
\noindent \textbf{(P\_1,1.27.3)} sarvādyānantaryakāryārtham . \\
\noindent \textbf{(P\_1,1.27.3)} sarvādīnām ānantaryeṇa yat ucyate kāryam tat api sañjñopasarjanībhūtānām mā bhūt iti . \\
\noindent \textbf{(P\_1,1.27.3)} kim prayojanam . \\
\noindent \textbf{(P\_1,1.27.3)} prayojanam ḍatarādīnām adbhāve . \\
\noindent \textbf{(P\_1,1.27.3)} ḍatarādīnām adbhāve prayojanam . \\
\noindent \textbf{(P\_1,1.27.3)} atikrāntam idam bṛāhmaṇakulam katarat , atikataram brāhmaṇakulam iti . \\
\noindent \textbf{(P\_1,1.27.3)} tyadādividhau ca . \\
\noindent \textbf{(P\_1,1.27.3)} tyadādividhau ca prayojanam . \\
\noindent \textbf{(P\_1,1.27.3)} atikrāntaḥ ayam brāhmaṇaḥ tam atitat brāhmaṇaḥ iti . \\
\noindent \textbf{(P\_1,1.27.3)} sañjāpratiṣedhaḥ tāvat na vaktavyaḥ . \\
\noindent \textbf{(P\_1,1.27.3)} upariṣṭāt yogavibhāgaḥ kariṣyate . \\
\noindent \textbf{(P\_1,1.27.3)} pūrvaparāvaradakṣiṇottarāparādharāṇi vyavasthāyām . \\
\noindent \textbf{(P\_1,1.27.3)} tataḥ asañjñāyām iti . \\
\noindent \textbf{(P\_1,1.27.3)} sarvādīni iti evam yāni anukrāntāni asañjñāyām tāni draṣṭavyāni . \\
\noindent \textbf{(P\_1,1.27.3)} upasarjanapratiṣedhaḥ ca na kartavyaḥ . \\
\noindent \textbf{(P\_1,1.27.3)} anupasarjanāt iti eṣaḥ yogaḥ pratyākhyāyate . \\
\noindent \textbf{(P\_1,1.27.3)} tam evam abhisambhantsyāmaḥ : anupasarjana* a* at iti . \\
\noindent \textbf{(P\_1,1.27.3)} kim idam a* at iti . \\
\noindent \textbf{(P\_1,1.27.3)} akārātkārau śiṣyamāṇau anupasarjanasya draṣṭavyau . \\
\noindent \textbf{(P\_1,1.27.3)} yadi evam atiyuṣmat atyasmat iti na sidhyati . \\
\noindent \textbf{(P\_1,1.27.3)} praśliṣṭanirdeśaḥ ayam : anupasarjana* a* a* at iti . \\
\noindent \textbf{(P\_1,1.27.3)} akārāntāt akārātkārau śiṣyamāṇau anupasarjanasya draṣṭavyau . \\
\noindent \textbf{(P\_1,1.27.3)} atha vā aṅgādhikāre yat ucyate gṛhyamāṇavibhakteḥ tat bhavati . \\
\noindent \textbf{(P\_1,1.27.3)} yadi evam paramapañca paramasapta ṣaḍbhyaḥ luk iti luk na prāpnoti . \\
\noindent \textbf{(P\_1,1.27.3)} na eṣaḥ doṣaḥ . \\
\noindent \textbf{(P\_1,1.27.3)} ṣaṭpradhānaḥ eṣaḥ samāsaḥ . \\
\noindent \textbf{(P\_1,1.27.3)} iha tarhi priyasakthnā brāhmaṇena anaṅ na prāpnoti . \\
\noindent \textbf{(P\_1,1.27.3)} saptamīnirdiṣṭe yat ucyate prakṛtavibhaktau tat bhavati . \\
\noindent \textbf{(P\_1,1.27.3)} yadi evam atitat , atitadau , atitadaḥ iti atvam prāpnoti . \\
\noindent \textbf{(P\_1,1.27.3)} tat ca api vaktavyam . \\
\noindent \textbf{(P\_1,1.27.3)} na vaktavyam . \\
\noindent \textbf{(P\_1,1.27.3)} iha tāvat adḍ ḍatarādibhyaḥ pañcabhyaḥ iti pañcamī aṅgasya iti ṣaṣṭhī . \\
\noindent \textbf{(P\_1,1.27.3)} tatra aśakyam vivibhaktitvāt ḍatarādibhyaḥ iti pañcamyā aṅgam viśeṣayitum . \\
\noindent \textbf{(P\_1,1.27.3)} tatra kim anyat śakyam viśeṣayitum anyat ataḥ vihitāt pratyayāt . \\
\noindent \textbf{(P\_1,1.27.3)} ḍatarādibhyaḥ yaḥ vihitaḥ iti . \\
\noindent \textbf{(P\_1,1.27.3)} iha idānīm asthidadhisakhthyakṣṇām anaṅ udāttaḥ iti tyadādīnām aḥ bhavati iti asthyādīnām iti eṣā ṣaṣṭhī aṅgasya iti api tyadādīnām iti api ṣaṣṭhī aṅgasya iti api . \\
\noindent \textbf{(P\_1,1.27.3)} tatra kāmacāraḥ : gṛhyamāṇena vā vibhaktim viśeṣayitum aṅgena vā . \\
\noindent \textbf{(P\_1,1.27.3)} yāvatā kāmacāraḥ iha tāvat asthidadhisakhthyakṣṇām anaṅ udāttaḥ iti aṅgena vibhaktim viśeṣayiṣyāmaḥ asthyādibhiḥ anaṅam : aṅgasya vibhaktau anaṅ bhavati asthyādīnām iti . \\
\noindent \textbf{(P\_1,1.27.3)} iha idānīm tyadādīnām aḥ bhavati iti gṛhyamāṇena vibhaktim viśeṣayiṣyāmaḥ aṅgena akāram : tyadādīnām vibhaktau aḥ bhavati aṅgasya iti . \\
\noindent \textbf{(P\_1,1.27.3)} yadi evam atisaḥ : atvam na prāpnoti . \\
\noindent \textbf{(P\_1,1.27.3)} na eṣaḥ doṣaḥ . \\
\noindent \textbf{(P\_1,1.27.3)} tyadādipradhānaḥ eṣaḥ samāsaḥ . \\
\noindent \textbf{(P\_1,1.27.3)} atha vā na idam sañjñākaraṇam . \\
\noindent \textbf{(P\_1,1.27.3)} pāṭhaviśeṣaṇam idam : sarveṣām yāni nāmāni tāni sarvādīni . \\
\noindent \textbf{(P\_1,1.27.3)} sañjñopasarjane ca viśeṣe avatiṣṭhete . \\
\noindent \textbf{(P\_1,1.27.3)} yadi evam sañjñāśrayam yat kāryam tat na sidhyati : sarvanāmnaḥ smai , āmi sarvanāmnaḥ suṭ iti . \\
\noindent \textbf{(P\_1,1.27.3)} anvarthagrahaṇam tatra vijñāsyate : sarveṣām yat nāma tat sarvanāma . \\
\noindent \textbf{(P\_1,1.27.3)} sarvanāmnaḥ uttarasya ṅeḥ smai bhavati . \\
\noindent \textbf{(P\_1,1.27.3)} sarvanāmnaḥ uttarasya āmaḥ suṭ bhavati . \\
\noindent \textbf{(P\_1,1.27.3)} yadi evam sakalam , kṛtsnam , jagat iti atra api prāpnoti . \\
\noindent \textbf{(P\_1,1.27.3)} eteṣām ca api śabdānām ekaikasya saḥ saḥ viṣayaḥ . \\
\noindent \textbf{(P\_1,1.27.3)} tasmin tasmin viṣaye yaḥ yaḥ śabdaḥ vartate tasya tasya tasmin tasmin vartamānasya sarvanāmakāryam prāpnoti . \\
\noindent \textbf{(P\_1,1.27.3)} evam tarhi ubhayam anena kriyate . \\
\noindent \textbf{(P\_1,1.27.3)} pāṭhaḥ ca eva viśeṣyate sañjñā ca . \\
\noindent \textbf{(P\_1,1.27.3)} katham punaḥ ekena yatnena ubhayam labhyam . \\
\noindent \textbf{(P\_1,1.27.3)} labhyam iti āha . \\
\noindent \textbf{(P\_1,1.27.3)} katham. ekaśeṣanirdeśāt . \\
\noindent \textbf{(P\_1,1.27.3)} ekaśeṣanirdeśaḥ ayam : sarvādīni ca sarvādīni ca sarvādīni . \\
\noindent \textbf{(P\_1,1.27.3)} sarvanāmāni ca sarvanāmāni ca sarvanāmāni . \\
\noindent \textbf{(P\_1,1.27.3)} sarvādīni sarvanānasañjñāni bhavanti . \\
\noindent \textbf{(P\_1,1.27.3)} sarveṣām yāni ca nāmāni tāni sarvādīni . \\
\noindent \textbf{(P\_1,1.27.3)} sañjñopasarjane ca viśeṣe avatiṣṭhete . \\
\noindent \textbf{(P\_1,1.27.3)} atha vā mahatī iyam sañjñā kriyate . \\
\noindent \textbf{(P\_1,1.27.3)} sañjñā ca nāma yataḥ na laghīyaḥ . \\
\noindent \textbf{(P\_1,1.27.3)} kutaḥ etat . \\
\noindent \textbf{(P\_1,1.27.3)} laghvartham hi sañjñākaraṇam . \\
\noindent \textbf{(P\_1,1.27.3)} tatra mahatyāḥ sañjñāyāḥ karaṇe etat prayojanam anvarthasañjñā yathā vijñāyeta . \\
\noindent \textbf{(P\_1,1.27.3)} sarvādīni sarvanānasañjñāni bhavanti sarveṣām nāmāni iti ca ataḥ sarvanāmāni . \\
\noindent \textbf{(P\_1,1.27.3)} sañjñopasarjane ca viśeṣe avatiṣṭhete . \\
\noindent \textbf{(P\_1,1.27.4)} atha ubhasya sarvanāmatve kaḥ arthaḥ . \\
\noindent \textbf{(P\_1,1.27.4)} ubhasya sarvanāmatve akajarthaḥ . \\
\noindent \textbf{(P\_1,1.27.4)} ubhasya sarvanāmatve akajarthaḥ pāṭhaḥ kriyate : ubhakau . \\
\noindent \textbf{(P\_1,1.27.4)} kim ucyate akajarthaḥ iti na punaḥ anyāni api sarvanāmakāryāṇi . \\
\noindent \textbf{(P\_1,1.27.4)} anyābhāvaḥ dvivacanaṭābviṣayatvāt . \\
\noindent \textbf{(P\_1,1.27.4)} anyeṣām sarvanāmkāryāṇām abhāvaḥ . \\
\noindent \textbf{(P\_1,1.27.4)} kim kāraṇam . \\
\noindent \textbf{(P\_1,1.27.4)} dvivacanaṭābviṣayatvāt . \\
\noindent \textbf{(P\_1,1.27.4)} ubhaśabdaḥ ayam dvivacanaṭābviṣayaḥ . \\
\noindent \textbf{(P\_1,1.27.4)} anyāni ca sarvanāmakāryāṇi ekavacanabahuvacaneṣu ucyante . \\
\noindent \textbf{(P\_1,1.27.4)} yadā punaḥ ayam ubhaśabdaḥ dvivacanaṭābviṣayaḥ kaḥ idānīm asya anyatra bhavati . \\
\noindent \textbf{(P\_1,1.27.4)} ubhayaḥ anyatra . \\
\noindent \textbf{(P\_1,1.27.4)} ubhayaśabdaḥ asya anyatra bahvati . \\
\noindent \textbf{(P\_1,1.27.4)} ubhaye devamanuṣyāḥ , ubhayaḥ maṇiḥ iti . \\
\noindent \textbf{(P\_1,1.27.4)} kim ca syāt yadi atra akac na syāt . \\
\noindent \textbf{(P\_1,1.27.4)} kaḥ prasajyeta . \\
\noindent \textbf{(P\_1,1.27.4)} kaḥ ca idānīm kākacoḥ viśeṣaḥ . \\
\noindent \textbf{(P\_1,1.27.4)} ubhaśabdaḥ ayam dvivacanaṭābviṣayaḥ iti uktam . \\
\noindent \textbf{(P\_1,1.27.4)} tatra akaci sati akacaḥ tanmadhyapatitatvāt śakyate etat vaktum : dvivacanaparaḥ ayam iti . \\
\noindent \textbf{(P\_1,1.27.4)} ke punaḥ sati na ayam dvivacanaparaḥ syāt . \\
\noindent \textbf{(P\_1,1.27.4)} tatra dvivacanaparatā vaktavyā . \\
\noindent \textbf{(P\_1,1.27.4)} yathā eva tarhi ke sati na ayam dvivacanaparaḥ evam āpi api sati na ayam dvivacanaparaḥ syāt . \\
\noindent \textbf{(P\_1,1.27.4)} tatra api dvivacanaparatā vaktavyā . \\
\noindent \textbf{(P\_1,1.27.4)} avacanāt api tatparavijñānam . \\
\noindent \textbf{(P\_1,1.27.4)} antareṇa api vacanam āpi dvivacanaparaḥ ayam bhaviṣyati . \\
\noindent \textbf{(P\_1,1.27.4)} kim vaktavyam etat . \\
\noindent \textbf{(P\_1,1.27.4)} na hi . \\
\noindent \textbf{(P\_1,1.27.4)} katham anucyamānam gaṃsyate . \\
\noindent \textbf{(P\_1,1.27.4)} ekādeśe kṛte dvivacanaparaḥ ayam antādivadbhāvena . \\
\noindent \textbf{(P\_1,1.27.4)} avacanāt āpi tatparavijñānam iti cet ke api tulyam . \\
\noindent \textbf{(P\_1,1.27.4)} avacanāt āpi tatparavijñānam iti cet ke api antareṇa vacanam dvivacanaparaḥ bhaviṣyati . \\
\noindent \textbf{(P\_1,1.27.4)} katham . \\
\noindent \textbf{(P\_1,1.27.4)} svārthikāḥ pratyayāḥ prakṛtitaḥ aviśiṣṭāḥ bhavanti iti prakṛtigrahaṇena svārthikānām api grahaṇam bhavati . \\
\noindent \textbf{(P\_1,1.27.4)} atha bhavataḥ sarvanāmatve kāni projanāni . \\
\noindent \textbf{(P\_1,1.27.4)} bhavataḥ akaccheṣātvāni . \\
\noindent \textbf{(P\_1,1.27.4)} bhavataḥ akaccheṣātvāni prayojanāni . \\
\noindent \textbf{(P\_1,1.27.4)} akac : bhavakān . \\
\noindent \textbf{(P\_1,1.27.4)} śeṣaḥ : saḥ ca bhavān ca bhavantau . \\
\noindent \textbf{(P\_1,1.27.4)} ātvam : bhavādṛk iti . \\
\noindent \textbf{(P\_1,1.27.4)} kim punaḥ idam parigaṇanam āhosvit udāharaṇamātram . \\
\noindent \textbf{(P\_1,1.27.4)} udāharaṇamātram iti āha . \\
\noindent \textbf{(P\_1,1.27.4)} tṛtīyādayaḥ api hi iṣyante . \\
\noindent \textbf{(P\_1,1.27.4)} sarvanāmnaḥ tṛtīyā ca : bhavatā hetunā , bhavataḥ hetoḥ iti . \\
\noindent \textbf{(P\_1,1.28)} diggrahaṇam kimartham . \\
\noindent \textbf{(P\_1,1.28)} na bahuvrīhau iti pratiṣedham vakṣyati . \\
\noindent \textbf{(P\_1,1.28)} tatra na jñāyate kva vibhāṣā kva pratiṣedhaḥ iti . \\
\noindent \textbf{(P\_1,1.28)} diggrahaṇe punaḥ kriyamāṇe na doṣaḥ bhavati . \\
\noindent \textbf{(P\_1,1.28)} digupadiṣṭe vibhāṣā anyatra pratiṣedhaḥ . \\
\noindent \textbf{(P\_1,1.28)} atha samāsagrahaṇam kimartham. samāsaḥ eva yaḥ bahuvrīhiḥ tatra yathā syāt . \\
\noindent \textbf{(P\_1,1.28)} bahuvrīhivadbhāvena yaḥ bahuvrīhiḥ tatra mā bhūt iti : dakṣiṇadakṣiṇasyai dehi iti . \\
\noindent \textbf{(P\_1,1.28)} atha bahuvrīhigrahaṇam kimartham . \\
\noindent \textbf{(P\_1,1.28)} dvandve mā bhūt dakṣiṇottarapūrvāṇām iti . \\
\noindent \textbf{(P\_1,1.28)} na etat asti prayojanam . \\
\noindent \textbf{(P\_1,1.28)} dvandve ca iti pratiṣedhaḥ bhaviṣyati . \\
\noindent \textbf{(P\_1,1.28)} na aprāpte pratiṣedhe iyam paribhāṣā ārabhyate . \\
\noindent \textbf{(P\_1,1.28)} sā yathā eva bahuvrīhau iti etam pratiṣedham bādhate evam dvandve ca iti etam api bādheta . \\
\noindent \textbf{(P\_1,1.28)} na bādhate . \\
\noindent \textbf{(P\_1,1.28)} kim kāraṇam . \\
\noindent \textbf{(P\_1,1.28)} yena na aprāpte tasya bādhanam bhavati . \\
\noindent \textbf{(P\_1,1.28)} na ca aprāpte na bahuvrīhau iti etasmin pratiṣedhe iyam paribhāṣā ārabhyate . \\
\noindent \textbf{(P\_1,1.28)} dvandve ca iti etasmin punaḥ prāpte ca aprāpte ca . \\
\noindent \textbf{(P\_1,1.28)} atha vā purastāt apavādāḥ anantarān vidhīn bādhante iti evam iyam vibhāṣā na bahuvrīhau iti etam pratiṣedham bādhiṣyate dvandve ca iti etam pratiṣedham na bādhiṣyate . \\
\noindent \textbf{(P\_1,1.28)} atha vā idam tāvat ayam praṣṭavyaḥ . \\
\noindent \textbf{(P\_1,1.28)} iha kasmāt na bhavati : yā pūrvā sā uttarā asya unmugdhasya saḥ ayam pūrvottaraḥ unmugdhaḥ , tasmai pūrvottarāya dehi . \\
\noindent \textbf{(P\_1,1.28)} lakṣaṇapratipadoktayoḥ pratipadoktasya eva iti . \\
\noindent \textbf{(P\_1,1.28)} yadi evam na arthaḥ bahuvrīhigrahaṇena . \\
\noindent \textbf{(P\_1,1.28)} dvandve kasmāt na bhavati . \\
\noindent \textbf{(P\_1,1.28)} lakṣaṇapratipadoktayoḥ pratipadoktasya eva iti . \\
\noindent \textbf{(P\_1,1.28)} uttarārtham tarhi bahuvrīhigrahaṇam kartavyam . \\
\noindent \textbf{(P\_1,1.28)} na kartavyam . \\
\noindent \textbf{(P\_1,1.28)} kriyate tatra eva bahuvrīhau iti . \\
\noindent \textbf{(P\_1,1.28)} dvitīyam kartavyam . \\
\noindent \textbf{(P\_1,1.28)} bahuvrīhiḥ eva yaḥ bahuvrīhiḥ tatra yathā syāt . \\
\noindent \textbf{(P\_1,1.28)} bahuvrīhivadbhāvena yaḥ bahuvrīhiḥ tatra mā bhūt iti : ekaikasmai dehi . \\
\noindent \textbf{(P\_1,1.28)} etat api na asti prayojanam. samāse iti vartate . \\
\noindent \textbf{(P\_1,1.28)} tena bahuvrīhim viśeṣayiṣyāmaḥ : samāsaḥ yaḥ bahuvrīhiḥ iti . \\
\noindent \textbf{(P\_1,1.28)} idam tarhi prayojanam . \\
\noindent \textbf{(P\_1,1.28)} avayavabhūtasya api bahuvrīheḥ pratiṣedhaḥ yathā syāt . \\
\noindent \textbf{(P\_1,1.28)} iha mā bhūt vastram antaram eṣām te ime vastrāntarāḥ vasanam antaram eṣām te ime vasanāntarāḥ vastrāntarāḥ ca vasanāntarāḥ ca vastrāntaravasanāntarāḥ . \\
\noindent \textbf{(P\_1,1.29.1)} kim udāharaṇam . \\
\noindent \textbf{(P\_1,1.29.1)} priyaviśvāya . \\
\noindent \textbf{(P\_1,1.29.1)} na etat asti prayojanam . \\
\noindent \textbf{(P\_1,1.29.1)} sarvādyantasya bahuvrīheḥ pratiṣedhena bhavitavyam . \\
\noindent \textbf{(P\_1,1.29.1)} vakṣyati ca etat : bahuvrīhau sarvanāmasaṅkhyayoḥ upasaṅkhyānam iti . \\
\noindent \textbf{(P\_1,1.29.1)} tatra viśvapriyāya iti bhavitavyam . \\
\noindent \textbf{(P\_1,1.29.1)} idam tarhi : dvyanyāya tryanyāya . \\
\noindent \textbf{(P\_1,1.29.1)} nanu ca atra api sarvanāmnaḥ eva pūrvanipātena bhavitavyam . \\
\noindent \textbf{(P\_1,1.29.1)} na eṣaḥ doṣaḥ . \\
\noindent \textbf{(P\_1,1.29.1)} vakṣyati etat : saṅkhyāsarvanāmnoḥ yaḥ bahuvrīhiḥ paratvāt tatra saṅkhyāyāḥ pūrvanipātaḥ bhavati iti . \\
\noindent \textbf{(P\_1,1.29.1)} idam ca api udāharaṇam priyaviśvāya . \\
\noindent \textbf{(P\_1,1.29.1)} nanu ca uktam viśvapriyāya iti bhavitavyam iti . \\
\noindent \textbf{(P\_1,1.29.1)} vakṣyati etat : vā priyasya iti . \\
\noindent \textbf{(P\_1,1.29.1)} na khalu api avaśyam sarvādyantasya eva bahuvrīheḥ pratiṣedhena bhavitavyam . \\
\noindent \textbf{(P\_1,1.29.1)} kim tarhi . \\
\noindent \textbf{(P\_1,1.29.1)} asarvādyantasya api bhavitavyam . \\
\noindent \textbf{(P\_1,1.29.1)} kim prayojanam . \\
\noindent \textbf{(P\_1,1.29.1)} akac mā bhūt . \\
\noindent \textbf{(P\_1,1.29.1)} kim ca syāt yadi akac syāt . \\
\noindent \textbf{(P\_1,1.29.1)} kaḥ na syāt . \\
\noindent \textbf{(P\_1,1.29.1)} kaḥ ca idānīm kākacoḥ viśeṣaḥ . \\
\noindent \textbf{(P\_1,1.29.1)} vyañjanānteṣu viśeṣaḥ . \\
\noindent \textbf{(P\_1,1.29.1)} ahakam pitā asya makatpitṛkaḥ , tvakam pitā asya tvakatpitṛkaḥ iti prāpnoti , matkapitṛkaḥ tvatkapitṛkaḥ iti ca iṣyate . \\
\noindent \textbf{(P\_1,1.29.1)} katham punaḥ icchatā api bhavatā bahiraṅgena pratiṣedhena antaraṅgaḥ vidhiḥ śakhyaḥ bādhitum . \\
\noindent \textbf{(P\_1,1.29.1)} antaraṅgān api vidhīn bahiraṅgaḥ vidhiḥ bādhate gomatpriyaḥ iti yathā . \\
\noindent \textbf{(P\_1,1.29.1)} kriyate tatra yatnaḥ : pratyayottarapadayoḥ ca iti . \\
\noindent \textbf{(P\_1,1.29.1)} nanu ca iha api kriyate : na bahuvrīhau iti . \\
\noindent \textbf{(P\_1,1.29.1)} asti anyat etasya vacane prayojanam . \\
\noindent \textbf{(P\_1,1.29.1)} kim . \\
\noindent \textbf{(P\_1,1.29.1)} priyaviśvāya . \\
\noindent \textbf{(P\_1,1.29.1)} upasarjanapratiṣedhena api etat siddham . \\
\noindent \textbf{(P\_1,1.29.1)} ayam khalu api bahuvrīhiḥ asti eva prāthamakalpikaḥ yasmin aikapadyam aikasvaryam aikavibhaktikatvam ca . \\
\noindent \textbf{(P\_1,1.29.1)} asti tādarthyāt tācchabdyam : bahuvrīhyarthāni padāni bahuvrīhiḥ iti . \\
\noindent \textbf{(P\_1,1.29.1)} tat yat tādarthyāt tācchabdyam tasya idam grahaṇam . \\
\noindent \textbf{(P\_1,1.29.1)} gonardīyaḥ āha akacsvarau tu kartavyau pratyaṅgam muktasaṃśayau . tvakatpitṛkaḥ makatpitṛkaḥ iti eva bhavitavyam iti . \\
\noindent \textbf{(P\_1,1.29.2)} pratiṣedhe bhūtapūrvasya upasaṅkhyānam . \\
\noindent \textbf{(P\_1,1.29.2)} pratiṣedhe bhūtapūrvasya upasaṅkhyānam kartavyam . \\
\noindent \textbf{(P\_1,1.29.2)} āḍhyaḥ bhūtapūrvaḥ āḍhyapūrvaḥ , āḍhyapūrvāya dehi iti . \\
\noindent \textbf{(P\_1,1.29.2)} pratiṣedhe bhūtapūrvasya upasaṅkhyānānarthakyam pūrvādīnām vyavasthāyām iti vacanāt . \\
\noindent \textbf{(P\_1,1.29.2)} pratiṣedhe bhūtapūrvasya upasaṅkhyānam narthakam . \\
\noindent \textbf{(P\_1,1.29.2)} kim kāraṇam . \\
\noindent \textbf{(P\_1,1.29.2)} pūrvādīnām vyavasthāyām iti vacanāt . \\
\noindent \textbf{(P\_1,1.29.2)} pūrvādīnām vyavasthāyām sarvanāmasñjñā ucyate . \\
\noindent \textbf{(P\_1,1.29.2)} na ca atra vyavasthā gamyate . \\
\noindent \textbf{(P\_1,1.30)} samāse iti vartamāne punaḥ samāsagrahaṇam kimartham . \\
\noindent \textbf{(P\_1,1.30)} ayam tṛtīyāsamāsaḥ asti eva prāthamakalpikaḥ yasmin aikapadyam aikasvaryam aikavibhaktikatvam ca . \\
\noindent \textbf{(P\_1,1.30)} asti tādarthyāt tācchabdyam : tṛtīyāsamāsārthāni padāni tṛtīyāsamāsaḥ iti . \\
\noindent \textbf{(P\_1,1.30)} tat yat tādarthyāt tācchabdyam tasya idam grahaṇam . \\
\noindent \textbf{(P\_1,1.30)} atha vā samase iti vartamāne punaḥ samāsagrahaṇasya etat prayojanam : yogāṅgam yathā upajāyeta . \\
\noindent \textbf{(P\_1,1.30)} sati yogāṅge yogavibhāgaḥ kariṣyate . \\
\noindent \textbf{(P\_1,1.30)} tṛtīyā . \\
\noindent \textbf{(P\_1,1.30)} tṛtīyāsamāse sarvādīni sarvanāmasañjñāni na bhavanti . \\
\noindent \textbf{(P\_1,1.30)} māsapūrvāya dehi saṃvatsarapūrvāya dehi . \\
\noindent \textbf{(P\_1,1.30)} tataḥ asamāse . \\
\noindent \textbf{(P\_1,1.30)} asamāse ca tṛtīyāyāḥ sarvādīni sarvanāmasañjñāni na bhavanti . \\
\noindent \textbf{(P\_1,1.30)} māsena pūrvāya iti \\
\noindent \textbf{(P\_1,1.32)} jasaḥ kāryam prati vibhāṣā , akac hi na bhavati . \\
\noindent \textbf{(P\_1,1.34)} avarādīnām ca punaḥ sūtrapāṭhe grahaṇānarthakyam gaṇe paṭhitatvāt . \\
\noindent \textbf{(P\_1,1.34)} avarādīnām ca punaḥ sūtrapāṭhe grahaṇam anarthakam . \\
\noindent \textbf{(P\_1,1.34)} kim kāraṇam . \\
\noindent \textbf{(P\_1,1.34)} gaṇe paṭhitatvāt . \\
\noindent \textbf{(P\_1,1.34)} gaṇe hi etāni paṭhyante . \\
\noindent \textbf{(P\_1,1.34)} katham punaḥ jñāyate saḥ pūrvaḥ pāṭhaḥ ayam punaḥ pāṭhaḥ iti . \\
\noindent \textbf{(P\_1,1.34)} tāni hi pūrvādīni imāni avarādīni . \\
\noindent \textbf{(P\_1,1.34)} imāni api pūrvādīni . \\
\noindent \textbf{(P\_1,1.34)} evam tarhi ācāryapravṛttiḥ jñāpayati saḥ pūrvaḥ pāṭhaḥ ayam punaḥ pāṭhaḥ iti yat ayam pūrvādibhyaḥ navabhyaḥ vā iti navagrahaṇam karoti . \\
\noindent \textbf{(P\_1,1.34)} nava eva pūrvādīni . \\
\noindent \textbf{(P\_1,1.34)} idam tarhi prayojanam : vyavasthāyām asañjñāyām iti vakṣyāmi iti . \\
\noindent \textbf{(P\_1,1.34)} etat api na asti prayojanam . \\
\noindent \textbf{(P\_1,1.34)} evaṃviśiṣṭāni eva etāni gaṇe paṭhyante . \\
\noindent \textbf{(P\_1,1.34)} idam tarhi prayojanam dvyādiparyudāsena paryudāsaḥ mā bhūt iti . \\
\noindent \textbf{(P\_1,1.34)} etat api na asti prayojanam . \\
\noindent \textbf{(P\_1,1.34)} ācāryapravṛttiḥ jñāpayati na eṣām dvyādiparyudāsena paryudāsaḥ bhavati iti yat ayam pūrvatra asiddham iti nipātanam karoti . \\
\noindent \textbf{(P\_1,1.34)} vārttikakāraḥ ca paṭhati : jaśbhāvāt iti cet uttaratra abhāvāt apavādaprasaṅgaḥ iti . \\
\noindent \textbf{(P\_1,1.34)} idam tarhi prayojanam jasi vibhāṣām vakṣyāmi iti . \\
\noindent \textbf{(P\_1,1.35)} ākhyāgrahaṇam kimartham . \\
\noindent \textbf{(P\_1,1.35)} jñātidhanaparyāyavācī yaḥ svaśabdaḥ tasya yathā syāt . \\
\noindent \textbf{(P\_1,1.35)} iha mā bhūt : sve putrāḥ svāḥ putrāḥ sve gāvaḥ svāḥ gāvaḥ . \\
\noindent \textbf{(P\_1,1.36.1)} upasaṃvyānagrahaṇam anarthakam bahiryogeṇa kṛtatvāt . \\
\noindent \textbf{(P\_1,1.36.1)} upasaṃvyānagrahaṇam anarthakam . \\
\noindent \textbf{(P\_1,1.36.1)} kim kāraṇam . \\
\noindent \textbf{(P\_1,1.36.1)} bahiryogeṇa kṛtatvāt . \\
\noindent \textbf{(P\_1,1.36.1)} bahiryoge iti eva siddham . \\
\noindent \textbf{(P\_1,1.36.1)} na vā śāṭakayugādyartham . \\
\noindent \textbf{(P\_1,1.36.1)} na vā anarthakam . \\
\noindent \textbf{(P\_1,1.36.1)} kim kāraṇam . \\
\noindent \textbf{(P\_1,1.36.1)} śāṭakayugādyartham . \\
\noindent \textbf{(P\_1,1.36.1)} śāṭakayugādyartham tarhi idam vaktavyam yatra etat na jñāyate kim antarīyam kim uttarīyam iti. atra api yaḥ eṣaḥ manuṣyaḥ prekṣāpūrvakārī bhavati nirjñātam tasya bhavati idam antarīyam idam uttarīyam iti . \\
\noindent \textbf{(P\_1,1.36.2)} apuri iti vaktavyam . \\
\noindent \textbf{(P\_1,1.36.2)} iha mā bhūt : antarāyām puri vasati iti . \\
\noindent \textbf{(P\_1,1.36.2)} vāprakaraṇe tīyasya ṅitsu upasaṅkhyānam . vāprakaraṇe tīyasya ṅitsu upasaṅkhyānam kartavyam : dvitīyāyai dvitīyasyai tṛtīyāyai tṛtīyasyai . \\
\noindent \textbf{(P\_1,1.36.2)} vibhāṣā dvitīyātṛtīyābhyām iti etat na vaktavyam bhavati . \\
\noindent \textbf{(P\_1,1.36.2)} kim punaḥ atra jyāyaḥ . \\
\noindent \textbf{(P\_1,1.36.2)} upasaṅkhyānam eva atra jyāyaḥ . \\
\noindent \textbf{(P\_1,1.36.2)} idam api siddham bhavati : dvitīyāya dvitīyasmai tṛtīyāya tṛtīyasmai . \\
\noindent \textbf{(P\_1,1.37)} kimartham pṛthak grahaṇam svarādīnām kriyate na cādiṣu eva paṭhyeran . \\
\noindent \textbf{(P\_1,1.37)} cādīnām vai asattvavacanānām nipātasañjñā svarādīnām punaḥ sattvavacanānām asattvavacanānām ca . \\
\noindent \textbf{(P\_1,1.37)} atha kimartham ubhe sañjñe kriyete na nipātsañjñā eva syāt . \\
\noindent \textbf{(P\_1,1.37)} na evam śakyam . \\
\noindent \textbf{(P\_1,1.37)} nipātaḥ ekāc anāṅ iti pragṛhyasañjñā uktā . \\
\noindent \textbf{(P\_1,1.37)} sā svarādīnām api ekācām prasajyeta . \\
\noindent \textbf{(P\_1,1.37)} evam tarhi avyayasañjñā eva astu . \\
\noindent \textbf{(P\_1,1.37)} tat ca aśakyam . \\
\noindent \textbf{(P\_1,1.37)} vakṣyati etat : avyaye nañkunipātānām iti . \\
\noindent \textbf{(P\_1,1.37)} tat garīyasā nyāsena parigaṇanam kartavyam syāt . \\
\noindent \textbf{(P\_1,1.37)} tasmāt pṛthak grahaṇam kartavyam ubhe ca sañjñe kartavye . \\
\noindent \textbf{(P\_1,1.38.1)} asarvavibhaktau avibhaktinimittasya upasaṅkhyānam . \\
\noindent \textbf{(P\_1,1.38.1)} asarvavibhaktau avibhaktinimittasya upasaṅkhyānam kartavyam : nānā vinā . \\
\noindent \textbf{(P\_1,1.38.1)} kim punaḥ kāraṇam na sidhyati . \\
\noindent \textbf{(P\_1,1.38.1)} sarvavibhaktiḥ hi aviśeṣāt . sarvavibhaktiḥ hi eṣaḥ bhavati . \\
\noindent \textbf{(P\_1,1.38.1)} kim kāraṇam . \\
\noindent \textbf{(P\_1,1.38.1)} aviśeṣāt . \\
\noindent \textbf{(P\_1,1.38.1)} aviśeṣeṇa vihitatvāt . \\
\noindent \textbf{(P\_1,1.38.1)} tralādīnām ca upasaṅkhyānam . \\
\noindent \textbf{(P\_1,1.38.1)} tralādīnām ca upasaṅkhyānam kartavyam . \\
\noindent \textbf{(P\_1,1.38.1)} tatra yatra tataḥ yataḥ . \\
\noindent \textbf{(P\_1,1.38.1)} nanu ca viśeṣeṇa ete vidhīyante : pañcamyāḥ tasil saptamyāḥ tral iti . \\
\noindent \textbf{(P\_1,1.38.1)} vakṣyati etat : itarābhyaḥ api dṛśyante iti . \\
\noindent \textbf{(P\_1,1.38.2)} yadi punaḥ avibhaktiḥ śabdaḥ avyayasañjñaḥ bhavati iti ucyeta . \\
\noindent \textbf{(P\_1,1.38.2)} avibhaktau itaretarāśrayatvāt aprasiddhiḥ . \\
\noindent \textbf{(P\_1,1.38.2)} avibhaktau itaretarāśrayatvāt aprasiddhiḥ sañjñāyāḥ . \\
\noindent \textbf{(P\_1,1.38.2)} kā itaretarāśrayatā . \\
\noindent \textbf{(P\_1,1.38.2)} sati avibhaktitve sañjñayā bhavitavyam sañjñayā ca avibhaktitvam bhāvyate . \\
\noindent \textbf{(P\_1,1.38.2)} tat itaretarāśrayam bhavati , itaretarāśrayāṇi ca kāryāṇi na prakalpante . \\
\noindent \textbf{(P\_1,1.38.2)} aliṅgam asaṅkhyam iti vā . \\
\noindent \textbf{(P\_1,1.38.2)} atha vā aliṅgam asaṅkhyam avyayam iti vaktavyam . \\
\noindent \textbf{(P\_1,1.38.2)} evam api itaretarāśrayam eva bhavati . \\
\noindent \textbf{(P\_1,1.38.2)} kā itaretarāśrayatā . \\
\noindent \textbf{(P\_1,1.38.2)} sati aliṅgāsaṅkhyatve sañjñayā bhavitavyam sañjñayā ca aliṅgāsaṅkhyatvam bhāvyate . \\
\noindent \textbf{(P\_1,1.38.2)} tat itaretarāśrayam bhavati , itaretarāśrayāṇi ca kāryāṇi na prakalpante . \\
\noindent \textbf{(P\_1,1.38.2)} na idam vācanikam aliṅgatā asaṅkhyatā ca . \\
\noindent \textbf{(P\_1,1.38.2)} kim tarhi . \\
\noindent \textbf{(P\_1,1.38.2)} svābhāvikam etat . \\
\noindent \textbf{(P\_1,1.38.2)} tat yathā : samānam īhamānānām adhīyānānām ca ke cit arthaiḥ yujyante apare na . \\
\noindent \textbf{(P\_1,1.38.2)} tatra kim asmābhiḥ kartum śakyam . \\
\noindent \textbf{(P\_1,1.38.2)} svābhāvikam etat . \\
\noindent \textbf{(P\_1,1.38.2)} tat tarhi vaktavyam aliṅgam asaṅkhyam iti . \\
\noindent \textbf{(P\_1,1.38.2)} na vaktavyam . \\
\noindent \textbf{(P\_1,1.38.2)} siddham tu pāṭhāt . \\
\noindent \textbf{(P\_1,1.38.2)} pāṭhāt vā siddham etat. katham pāṭhaḥ kartavyaḥ . \\
\noindent \textbf{(P\_1,1.38.2)} tasilādayaḥ prāk pāsapaḥ , śasprabhṛtayaḥ prāk samāsāntebhyaḥ , māntaḥ , kṛtvorthaḥ , tasivatī , nānāñau iti . \\
\noindent \textbf{(P\_1,1.38.3)} atha vā punaḥ astu avibhaktiḥ śabdaḥ avyayasañjñaḥ bhavati iti eva . \\
\noindent \textbf{(P\_1,1.38.3)} nanu ca uktam avibhaktau itaretarāśrayatvāt aprasiddhiḥ iti . \\
\noindent \textbf{(P\_1,1.38.3)} na eṣaḥ doṣaḥ . \\
\noindent \textbf{(P\_1,1.38.3)} idam tāvat ayam praṣṭavyaḥ . \\
\noindent \textbf{(P\_1,1.38.3)} yadi api vaiyākaraṇāḥ vibhaktilopam ārabhamāṇāḥ avibhaktikān śabdāñ prayuñjate ye tu ete vaiyākaraṇebhyaḥ anye manuṣyāḥ katham te avibhaktikān śabdān prayuñjate iti . \\
\noindent \textbf{(P\_1,1.38.3)} abhijñāḥ ca punaḥ laukikāḥ ekatvādīnām arthānām . \\
\noindent \textbf{(P\_1,1.38.3)} ātaḥ ca abhijñāḥ : anyena hi vasnena ekam gām krīṇanti , anyena dvau , anyena trīn . \\
\noindent \textbf{(P\_1,1.38.3)} abhijñāḥ ca na ca prayuñjate . \\
\noindent \textbf{(P\_1,1.38.3)} tat etat evam sandṛśyatām : artharūpam etat evañjātīyakam yena atra vibhaktiḥ na bhavati iti . \\
\noindent \textbf{(P\_1,1.38.3)} tat ca api etat evam anugamyamānam dṛśyatām : kim cit avyayam vibhaktyarthapradhānam kim cit kriyāpradhānam . \\
\noindent \textbf{(P\_1,1.38.3)} uccaiḥ , nīcaiḥ iti vibhaktyarthapradhānam , hiruk pṛthak iti kriyāpradhānam . \\
\noindent \textbf{(P\_1,1.38.3)} taddhitaḥ ca api kaḥ cit vibhaktyarthapradhānaḥ kaḥ cit kriyāpradhānaḥ . \\
\noindent \textbf{(P\_1,1.38.3)} tatra yatra iti vibhaktyarthapradhānaḥ , nānā vina iti kriyāpradhānaḥ . \\
\noindent \textbf{(P\_1,1.38.3)} na ca etayoḥ arthayoḥ liṅgasaṅkhyābhyām yogaḥ asti . \\
\noindent \textbf{(P\_1,1.38.4)} atha api asarvavibhaktiḥ iti ucyate evam api na doṣaḥ . \\
\noindent \textbf{(P\_1,1.38.4)} katham . \\
\noindent \textbf{(P\_1,1.38.4)} idam ca api adyatve atibahu kriyate : ekasmin ekavacanam , dvayoḥ dvivacanam , bahuṣu bahuvacanam iti . \\
\noindent \textbf{(P\_1,1.38.4)} katham tarhi . \\
\noindent \textbf{(P\_1,1.38.4)} ekavacanam utsargaḥ kariṣyate . \\
\noindent \textbf{(P\_1,1.38.4)} tasya dvibahvoḥ arthayoḥ dvivacanabahuvacane bādhake bhaviṣyataḥ . \\
\noindent \textbf{(P\_1,1.38.4)} na ca api evam vigrahaḥ kariṣyate : na sarvāḥ asarvāḥ , asarvāḥ vibhaktayaḥ asmāt iti . \\
\noindent \textbf{(P\_1,1.38.4)} katham tarhi . \\
\noindent \textbf{(P\_1,1.38.4)} na sarvā asarvā , asarvā vibhaktiḥ asmāt iti . \\
\noindent \textbf{(P\_1,1.38.4)} trikam punaḥ vibhaktisañjñam . \\
\noindent \textbf{(P\_1,1.38.4)} evam gate kṛti api tulyam etat māntasya kāryam grahaṇam na tatra . \\
\noindent \textbf{(P\_1,1.38.4)} tataḥ pare ca abhimatāḥ kāryāḥ trayaḥ kṛdarthāḥ grahaṇena yogāḥ . \\
\noindent \textbf{(P\_1,1.38.4)} kṛttaddhitānām grahaṇam tu kāryam saṅkhyāviśeṣam hi abhiniśritāḥ ye . \\
\noindent \textbf{(P\_1,1.38.4)} teṣām pratiṣedhaḥ bhavati iti vaktavyam . \\
\noindent \textbf{(P\_1,1.38.4)} iha mā bhūt : ekaḥ , dvau , bahavaḥ iti . \\
\noindent \textbf{(P\_1,1.38.4)} tasmāt svarādigrahaṇam ca kāryam kṛttaddhitānām ca pāṭhe . \\
\noindent \textbf{(P\_1,1.38.5)} pāṭhena iyam avyayasañjñā kriyate . \\
\noindent \textbf{(P\_1,1.38.5)} sā iha na prāpnoti : paramoccaiḥ , paramanīcaiḥ iti . \\
\noindent \textbf{(P\_1,1.38.5)} tadantavidhinā bhaviṣyati . \\
\noindent \textbf{(P\_1,1.38.5)} iha api tarhi prāpnoti : atyuccaiḥ atyuccaisau atyuccaisaḥ iti . \\
\noindent \textbf{(P\_1,1.38.5)} upasarjanasya na iti pratiṣedhaḥ bhaviṣyati . \\
\noindent \textbf{(P\_1,1.38.5)} saḥ tarhi pratiṣedhaḥ vaktavyaḥ . \\
\noindent \textbf{(P\_1,1.38.5)} na vaktavyaḥ . \\
\noindent \textbf{(P\_1,1.38.5)} sarvanāmasañjñāyām prakṛtaḥ pratiṣedhaḥ iha anuvartiṣyate . \\
\noindent \textbf{(P\_1,1.38.5)} saḥ vai tatra pratyākhyāyate . \\
\noindent \textbf{(P\_1,1.38.5)} yathā saḥ tatra pratyākhyāyate iha api tathā śakyaḥ pratyākhyātum . \\
\noindent \textbf{(P\_1,1.38.5)} katham saḥ tatra pratyākhyāyate . \\
\noindent \textbf{(P\_1,1.38.5)} mahatī iyam sañjñā kriyate . \\
\noindent \textbf{(P\_1,1.38.5)} iyam api ca mahatī sañjñā kriyate . \\
\noindent \textbf{(P\_1,1.38.5)} sañjñā ca nāma yataḥ na laghīyaḥ . \\
\noindent \textbf{(P\_1,1.38.5)} kutaḥ etat . \\
\noindent \textbf{(P\_1,1.38.5)} laghvartham hi sañjñākaraṇam . \\
\noindent \textbf{(P\_1,1.38.5)} tatra mahatyāḥ sañjñāyāḥ karaṇe etat prayojanam anvarthasañjñā yathā vijñāyeta : na vyeti iti avyayam iti . \\
\noindent \textbf{(P\_1,1.38.5)} kva punaḥ na vyeti . \\
\noindent \textbf{(P\_1,1.38.5)} strīpuṃnapuṃsakāni sattvaguṇāḥ ekatvadvitvabahutvāni ca . \\
\noindent \textbf{(P\_1,1.38.5)} etān arthān ke cit viyanti ke cit na viyanti . \\
\noindent \textbf{(P\_1,1.38.5)} ye na viyanti tad avyayam . \\
\noindent \textbf{(P\_1,1.38.5)} sadṛśam triṣu liṅgeṣu sarvāsu ca vibhaktiṣu vacaneṣu ca sarveṣu yat na vyeti tat avyayam . \\
\noindent \textbf{(P\_1,1.39.1)} katham idam vijñāyate : kṛt yaḥ māntaḥ iti āhosvit kṛdantam yat māntam iti . \\
\noindent \textbf{(P\_1,1.39.1)} kim ca ataḥ . \\
\noindent \textbf{(P\_1,1.39.1)} yadi vijñāyate kṛt yaḥ māntaḥ iti kārayām cakāra hārayām cakāra iti atra na prāpnoti . \\
\noindent \textbf{(P\_1,1.39.1)} atha vijñāyate kṛdantam yat māntam iti pratāmau pratāmaḥ iti atra api prāpnoti . \\
\noindent \textbf{(P\_1,1.39.1)} yathā icchasi tathā astu . \\
\noindent \textbf{(P\_1,1.39.1)} astu tāvat kṛt yaḥ māntaḥ iti . \\
\noindent \textbf{(P\_1,1.39.1)} katham kārayām cakāra hārayām cakāra iti . \\
\noindent \textbf{(P\_1,1.39.1)} kim punaḥ atra avyayasañjñayā prārthyate . \\
\noindent \textbf{(P\_1,1.39.1)} avyayāt iti luk yathā syāt . \\
\noindent \textbf{(P\_1,1.39.1)} mā bhūt evam . \\
\noindent \textbf{(P\_1,1.39.1)} āmaḥ iti evam bhaviṣyati . \\
\noindent \textbf{(P\_1,1.39.1)} na sidhyati . \\
\noindent \textbf{(P\_1,1.39.1)} ligrahaṇam tatra anuvartate . \\
\noindent \textbf{(P\_1,1.39.1)} ligrahaṇam nivartiṣyate . \\
\noindent \textbf{(P\_1,1.39.1)} yadi nivartate pratyayamātrasya luk prāpnoti . \\
\noindent \textbf{(P\_1,1.39.1)} iṣyate ca pratyayamātrasya . \\
\noindent \textbf{(P\_1,1.39.1)} ātaḥ ca iṣyate . \\
\noindent \textbf{(P\_1,1.39.1)} evam hi āha : kṛñ ca anuprayujyate liṭi iti . \\
\noindent \textbf{(P\_1,1.39.1)} yadi ca pratyayamātrasya luk bhavati tataḥ etat upapannam bhavati . \\
\noindent \textbf{(P\_1,1.39.1)} atha vā punaḥ astu kṛdantam yat māntam iti . \\
\noindent \textbf{(P\_1,1.39.1)} katham pratāmau pratāmaḥ iti . \\
\noindent \textbf{(P\_1,1.39.1)} ācāryapravṛttiḥ jñāpayati na pratyayalakṣaṇena avyayasañjñā bhavati iti yat ayam praśānśabdam svarādiṣu paṭhati . \\
\noindent \textbf{(P\_1,1.39.2)} kṛt mejantaḥ ca anikārokāraprakṛtiḥ . \\
\noindent \textbf{(P\_1,1.39.2)} kṛt mejantaḥ ca anikārokāraprakṛtiḥ iti vaktavyam . \\
\noindent \textbf{(P\_1,1.39.2)} iha mā bhūt : ādhaye , ādheḥ , cikīrṣave , cikīrṣoḥ iti . \\
\noindent \textbf{(P\_1,1.39.2)} ananyaprakṛtiḥ iti vā . \\
\noindent \textbf{(P\_1,1.39.2)} atha vā ananyaprakṛtiḥ kṛt avyayasañjñaḥ bhavati iti vaktavyam . \\
\noindent \textbf{(P\_1,1.39.2)} kim punaḥ atra jyāyaḥ . \\
\noindent \textbf{(P\_1,1.39.2)} ananyaprakṛtivacanam eva jyāyaḥ . \\
\noindent \textbf{(P\_1,1.39.2)} idam api siddham bhavati : kumbhakārebhyaḥ , nagarakārebhyaḥ iti . \\
\noindent \textbf{(P\_1,1.39.2)} tat tarhi vaktavyam . \\
\noindent \textbf{(P\_1,1.39.2)} na vā sannipātalakṣaṇaḥ vidhiḥ animittam tadvighātasya . \\
\noindent \textbf{(P\_1,1.39.2)} na vā vaktavyam . \\
\noindent \textbf{(P\_1,1.39.2)} kim kāraṇam . \\
\noindent \textbf{(P\_1,1.39.2)} sannipātalakṣaṇaḥ vidhiḥ animittam tadvighātasya iti eṣā paribhāṣā kartavyā . \\
\noindent \textbf{(P\_1,1.39.2)} kaḥ punaḥ atra viśeṣaḥ eṣā vā paribhāṣā kriyeta ananyaprakṛtiḥ iti vā ucyeta . \\
\noindent \textbf{(P\_1,1.39.2)} avaśyam eṣā paribhāṣā kartavyā . \\
\noindent \textbf{(P\_1,1.39.2)} bahūni etasyāḥ paribhāṣāyāḥ prayojanāni . \\
\noindent \textbf{(P\_1,1.39.2)} kāni punaḥ tāni . \\
\noindent \textbf{(P\_1,1.39.2)} prayojanam hrasvatam tugvidheḥ grāmaṇikulam . \\
\noindent \textbf{(P\_1,1.39.2)} grāmaṇikulam , senānikulam iti atra hrasvatve kṛte hrasvasya piti kṛti tuk bhavati iti tuk prāpnoti . \\
\noindent \textbf{(P\_1,1.39.2)} sannipātalakṣaṇaḥ vidhiḥ animittam tadvighātasya iti na doṣaḥ bhavati . \\
\noindent \textbf{(P\_1,1.39.2)} na etat asti prayojanam . \\
\noindent \textbf{(P\_1,1.39.2)} bahiraṅgam hrasvatvam . \\
\noindent \textbf{(P\_1,1.39.2)} antaraṅgaḥ tuk . \\
\noindent \textbf{(P\_1,1.39.2)} asiddham bahiraṅgam antaraṅge . \\
\noindent \textbf{(P\_1,1.39.2)} nalopaḥ vṛtrahabhiḥ . \\
\noindent \textbf{(P\_1,1.39.2)} vṛtrhabhiḥ , bhrūṇhabhiḥ iti atra nalope kṛte hrasvasya piti kṛti tuk bhavati iti tuk prāpnoti . \\
\noindent \textbf{(P\_1,1.39.2)} sannipātalakṣaṇaḥ vidhiḥ animittam tadvighātasya iti na doṣaḥ bhavati . \\
\noindent \textbf{(P\_1,1.39.2)} etat api na asti prayojanam . \\
\noindent \textbf{(P\_1,1.39.2)} asiddhaḥ nalopaḥ . \\
\noindent \textbf{(P\_1,1.39.2)} tasya asiddhatvāt na bhaviṣyati . \\
\noindent \textbf{(P\_1,1.39.2)} udupadhatvam akittvasya nikucite . \\
\noindent \textbf{(P\_1,1.39.2)} udupadhatvam akittvasya animittam . \\
\noindent \textbf{(P\_1,1.39.2)} kva . \\
\noindent \textbf{(P\_1,1.39.2)} nikucite . \\
\noindent \textbf{(P\_1,1.39.2)} nikucitaḥ iti atra nalope kṛte udupadhāt bhāvādikarmaṇoḥ anyatarasyām iti akittvam prāpnoti . \\
\noindent \textbf{(P\_1,1.39.2)} sannipātalakṣaṇaḥ vidhiḥ animittam tadvighātasya iti na doṣaḥ bhavati . \\
\noindent \textbf{(P\_1,1.39.2)} etat api na asti prayojanam . \\
\noindent \textbf{(P\_1,1.39.2)} astu atra akittvam . \\
\noindent \textbf{(P\_1,1.39.2)} na dhātulope ārdhadhātuke iti pratiṣedhaḥ bhaviṣyati . \\
\noindent \textbf{(P\_1,1.39.2)} nābhāvaḥ yañi dīrghatvasya amunā . \\
\noindent \textbf{(P\_1,1.39.2)} nābhāvaḥ yañi dīrghatvasya asnimittam . \\
\noindent \textbf{(P\_1,1.39.2)} kva . \\
\noindent \textbf{(P\_1,1.39.2)} amunā . \\
\noindent \textbf{(P\_1,1.39.2)} nābhāve kṛte ataḥ dīrghaḥ yañi supi ca iti dīrghatvam prāpnoti . \\
\noindent \textbf{(P\_1,1.39.2)} sannipātalakṣaṇaḥ vidhiḥ animittam tadvighātasya iti na doṣaḥ bhavati . \\
\noindent \textbf{(P\_1,1.39.2)} etat api na asti prayojanam . \\
\noindent \textbf{(P\_1,1.39.2)} vakṣyati etat : na mu ṭādeśe iti . \\
\noindent \textbf{(P\_1,1.39.2)} āttvam kittvasya upādāsta . \\
\noindent \textbf{(P\_1,1.39.2)} āttvam kittvasya animittam . \\
\noindent \textbf{(P\_1,1.39.2)} kva . \\
\noindent \textbf{(P\_1,1.39.2)} upādāsta asya svaraḥ śikṣakasya iti . \\
\noindent \textbf{(P\_1,1.39.2)} āttve kṛte sthāghvoḥ it ca iti ittvam prāpnoti . \\
\noindent \textbf{(P\_1,1.39.2)} sannipātalakṣaṇaḥ vidhiḥ animittam tadvighātasya iti na doṣaḥ bhavati . \\
\noindent \textbf{(P\_1,1.39.2)} etat api na asti prayojanam . \\
\noindent \textbf{(P\_1,1.39.2)} uktam etat : dīṅaḥ pratiṣedhaḥ sthāghvoḥ ittve iti . \\
\noindent \textbf{(P\_1,1.39.2)} tisṛcatasṛtvam ṅībvidheḥ . \\
\noindent \textbf{(P\_1,1.39.2)} tisṛcatasṛtvam ṅībvidheḥ animittam . \\
\noindent \textbf{(P\_1,1.39.2)} tisraḥ tiṣṭhanti catasraḥ tiṣṭhanti . \\
\noindent \textbf{(P\_1,1.39.2)} tisṛcatasṛbhāve kṛte ṛnnebhyaḥ ṅīp iti ṅīp prāpnoti . \\
\noindent \textbf{(P\_1,1.39.2)} sannipātalakṣaṇaḥ vidhiḥ animittam tadvighātasya iti na doṣaḥ bhavati . \\
\noindent \textbf{(P\_1,1.39.2)} etat api na asti prayojanam . \\
\noindent \textbf{(P\_1,1.39.2)} ācāryapravṛttiḥ jñāpayati na tisṛcatasṛbhāve kṛte ṅīp bhavati iti yat ayam na tisṛcatasṛ iti nāmi dīrghatvapratiṣedham śāsti . \\
\noindent \textbf{(P\_1,1.39.2)} imāni tarhi prayojanāni : śatāni sahasrāṇi . \\
\noindent \textbf{(P\_1,1.39.2)} numi kṛte ṣṇāntā ṣaṭ it ṣaṭsañjñā prāpnoti . \\
\noindent \textbf{(P\_1,1.39.2)} sannipātalakṣaṇaḥ vidhiḥ animittam tadvighātasya iti na doṣaḥ bhavati . \\
\noindent \textbf{(P\_1,1.39.2)} śakaṭau paddhatau . \\
\noindent \textbf{(P\_1,1.39.2)} attve kṛte ataḥ iti ṭāp prāpnoti . \\
\noindent \textbf{(P\_1,1.39.2)} sannipātalakṣaṇaḥ vidhiḥ animittam tadvighātasya iti na doṣaḥ bhavati . \\
\noindent \textbf{(P\_1,1.39.2)} iyeṣa , uvoṣa . \\
\noindent \textbf{(P\_1,1.39.2)} guṇe kṛte ijādeḥ ca gurumataḥ anṛcchaḥ iti ām prāpnoti . \\
\noindent \textbf{(P\_1,1.39.2)} sannipātalakṣaṇaḥ vidhiḥ animittam tadvighātasya iti na doṣaḥ bhavati . \\
\noindent \textbf{(P\_1,1.39.2)} tasya doṣaḥ varṇāśrayaḥ pratyayaḥ varṇavicālasya . \\
\noindent \textbf{(P\_1,1.39.2)} tasya etasya lakṣaṇasya doṣaḥ varṇāśrayaḥ pratyayaḥ varṇavicālasya animittam syāt . \\
\noindent \textbf{(P\_1,1.39.2)} kva . \\
\noindent \textbf{(P\_1,1.39.2)} ata iñ : dākṣiḥ , plākṣiḥ . \\
\noindent \textbf{(P\_1,1.39.2)} na pratyayaḥ sannipātalakṣaṇaḥ . \\
\noindent \textbf{(P\_1,1.39.2)} aṅgasañjñā tarhi animittam syāt . \\
\noindent \textbf{(P\_1,1.39.2)} āttvam pugvidheḥ krāpayati . \\
\noindent \textbf{(P\_1,1.39.2)} āttvam pugvidheḥ animittam syāt . \\
\noindent \textbf{(P\_1,1.39.2)} kva . \\
\noindent \textbf{(P\_1,1.39.2)} krāpayati iti . \\
\noindent \textbf{(P\_1,1.39.2)} pug hrasvatvasya adīdapat . \\
\noindent \textbf{(P\_1,1.39.2)} puk hrasvatvasya animittam syāt . \\
\noindent \textbf{(P\_1,1.39.2)} kva . \\
\noindent \textbf{(P\_1,1.39.2)} adīdapat iti . \\
\noindent \textbf{(P\_1,1.39.2)} tyadādyakāraḥ ṭābvidheḥ . \\
\noindent \textbf{(P\_1,1.39.2)} tyadādyakāraḥ ṭābvidheḥ animittam syāt . \\
\noindent \textbf{(P\_1,1.39.2)} kva . \\
\noindent \textbf{(P\_1,1.39.2)} yā sā . \\
\noindent \textbf{(P\_1,1.39.2)} iḍvidhiḥ ākāralopasya papivān . \\
\noindent \textbf{(P\_1,1.39.2)} iḍvidhiḥ ākāralopasya animittam syāt . \\
\noindent \textbf{(P\_1,1.39.2)} kva . \\
\noindent \textbf{(P\_1,1.39.2)} papivān tasthivān iti . \\
\noindent \textbf{(P\_1,1.39.2)} matubvibhaktyudāttatvam pūrvanighātasya . \\
\noindent \textbf{(P\_1,1.39.2)} matubvibhaktyudāttatvam pūrvanighātasya animittam syāt . \\
\noindent \textbf{(P\_1,1.39.2)} kva . \\
\noindent \textbf{(P\_1,1.39.2)} agnimān vāyumān paramavācā paramavāce . \\
\noindent \textbf{(P\_1,1.39.2)} nadīhrasvratvam sambuddhilopasya . \\
\noindent \textbf{(P\_1,1.39.2)} nadīhrasvratvam sambuddhilopasya animittam syāt . \\
\noindent \textbf{(P\_1,1.39.2)} kva . \\
\noindent \textbf{(P\_1,1.39.2)} nadi kumāri kiśori brāhmaṇi brahmabandhu . \\
\noindent \textbf{(P\_1,1.39.2)} hrasvatve kṛte eṅhrasvāt sambuddheḥ iti lopaḥ na prāpnoti . \\
\noindent \textbf{(P\_1,1.39.2)} mā bhūt evam . \\
\noindent \textbf{(P\_1,1.39.2)} ṅyantāt iti evam bhaviṣyati . \\
\noindent \textbf{(P\_1,1.39.2)} na sidhyati . \\
\noindent \textbf{(P\_1,1.39.2)} dīrghāt iti ucyate . \\
\noindent \textbf{(P\_1,1.39.2)} hrasvāntāt ca na prāpnoti . \\
\noindent \textbf{(P\_1,1.39.2)} idam iha sampradhāryam : hrasvatvam kriyatām sambuddhilopaḥ iti kim atra kartavyam . \\
\noindent \textbf{(P\_1,1.39.2)} paratvāt hrasvatvam . \\
\noindent \textbf{(P\_1,1.39.2)} nityaḥ sambuddhilopaḥ . \\
\noindent \textbf{(P\_1,1.39.2)} kṛte api hrasvatve prāpnoti akṛte api . \\
\noindent \textbf{(P\_1,1.39.2)} anityaḥ sambuddhilopaḥ . \\
\noindent \textbf{(P\_1,1.39.2)} na hi kṛte hrasvatve prāpnoti . \\
\noindent \textbf{(P\_1,1.39.2)} kim kāraṇam . \\
\noindent \textbf{(P\_1,1.39.2)} sannipātalakṣaṇaḥ vidhiḥ animittam tadvighātasya iti . \\
\noindent \textbf{(P\_1,1.39.2)} ete doṣāḥ samāḥ bhūyāṃsaḥ vā . \\
\noindent \textbf{(P\_1,1.39.2)} tasmāt na arthaḥ anayā paribhāṣayā . \\
\noindent \textbf{(P\_1,1.39.2)} na hi doṣāḥ santi iti paribhāṣā na kartavyā lakṣaṇam vā na praṇeyam . \\
\noindent \textbf{(P\_1,1.39.2)} na hi bhikṣukāḥ santi iti sthālyaḥ na adhiśrīyante na ca mṛgāḥ santi iti yavāḥ na upyante . \\
\noindent \textbf{(P\_1,1.39.2)} doṣāḥ khalu api sākalyena parigaṇitāḥ prayojanānām udāharaṇamātram . \\
\noindent \textbf{(P\_1,1.39.2)} kutaḥ etat . \\
\noindent \textbf{(P\_1,1.39.2)} na hi doṣāṇām lakṣaṇam asti . \\
\noindent \textbf{(P\_1,1.39.2)} tasmāt yāni etasyāḥ paribhāṣāyāḥ prayojanāni tadartham eṣā paribhāṣā kartavyā pratividheyam ca doṣeṣu . \\
\noindent \textbf{(P\_1,1.41)} avyayībhāvasya avyayatve prayojanam lugmukhsvaropacārāḥ . \\
\noindent \textbf{(P\_1,1.41)} avyayībhāvasya avyayatve prayojanam kim . \\
\noindent \textbf{(P\_1,1.41)} lugmukhsvaropacārāḥ . \\
\noindent \textbf{(P\_1,1.41)} luk : upāgni pratyagni . \\
\noindent \textbf{(P\_1,1.41)} avyayāt iti luk siddhaḥ bhavati . \\
\noindent \textbf{(P\_1,1.41)} mukhasvaraḥ . \\
\noindent \textbf{(P\_1,1.41)} upāgnimukhaḥ , pratyagnimukhaḥ . \\
\noindent \textbf{(P\_1,1.41)} na avyayadikśabdgomahatsthūlapṛthuvatsebhyaḥ iti pratiṣedhaḥ siddhaḥ bhavati . \\
\noindent \textbf{(P\_1,1.41)} upacāraḥ : upapayaḥkāraḥ , upapayaḥkāmaḥ iti . \\
\noindent \textbf{(P\_1,1.41)} ataḥ kṛkamikaṃsakumbhapātrakuśākarṇīṣu anavyayasya iti pratiṣedhaḥ siddhaḥ bhavati . \\
\noindent \textbf{(P\_1,1.41)} kim punaḥ idam parigaṇanam āhosvit udāharaṇamātram . \\
\noindent \textbf{(P\_1,1.41)} parigaṇanam iti āha . \\
\noindent \textbf{(P\_1,1.41)} api khalu api āhuḥ . \\
\noindent \textbf{(P\_1,1.41)} yat anyat avyayībhāvasya avyayakṛtam prāpnoti tasya pratiṣedhaḥ vaktavyaḥ iti . \\
\noindent \textbf{(P\_1,1.41)} kim punaḥ tat . \\
\noindent \textbf{(P\_1,1.41)} parāṅgavadbhāvaḥ . \\
\noindent \textbf{(P\_1,1.41)} parāṅgavadbhāve avyayapratiṣedhaḥ coditaḥ uccaiḥ adhīyāna nīcaiḥ adhīyāna iti evamartham . \\
\noindent \textbf{(P\_1,1.41)} saḥ iha api prāpnoti : upāgni adhīyāna pratyagni adhīyāna . \\
\noindent \textbf{(P\_1,1.41)} akaci avyayagrahaṇam kriyate uccakaiḥ , nīcakaiḥ iti evamartham . \\
\noindent \textbf{(P\_1,1.41)} tat iha api prāpnoti : upāgnikam , pratyagnikam iti . \\
\noindent \textbf{(P\_1,1.41)} mumi avyayapratiṣedhaḥ ucyate doṣāmanyam ahaḥ , divāmanyā rātriḥ iti evamartham . \\
\noindent \textbf{(P\_1,1.41)} saḥ iha api prāpnoti : aupakumbhammanyaḥ , upamaṇikammanyaḥ . \\
\noindent \textbf{(P\_1,1.41)} asya cvau avyayapratiṣedhaḥ ucyate doṣābhūtam ahaḥ , divābhūtā rātriḥ iti evamartham . \\
\noindent \textbf{(P\_1,1.41)} saḥ iha api prāpnoti : upakumbhībhūtam upamaṇikībhūtam . \\
\noindent \textbf{(P\_1,1.41)} yadi parigaṇanam kriyate na arthaḥ avyayībhāvasya avyayasañjñayā . \\
\noindent \textbf{(P\_1,1.41)} katham yāni avyayībhāvasya avyayatve prayojanāni . \\
\noindent \textbf{(P\_1,1.41)} na etāni santi . \\
\noindent \textbf{(P\_1,1.41)} yat tāvat ucyate luk iti : ācāryapravṛttiḥ jñāpayati bhavati avyayībhāvāt luk iti yad ayam na avyayībhāvāt ataḥ iti pratiṣedham śāsti . \\
\noindent \textbf{(P\_1,1.41)} upacāraḥ : anuttarapadasthasya iti vartate . \\
\noindent \textbf{(P\_1,1.41)} tatra mukhasvaraḥ ekaḥ prayojayati . \\
\noindent \textbf{(P\_1,1.41)} na ca ekam prayojanam yogārambham prayojayati . \\
\noindent \textbf{(P\_1,1.41)} yadi etāvat prayojanam syāt tatra eva ayam brūyāt nāvyayāt avyayībhāvāt ca iti . \\
\noindent \textbf{(P\_1,1.42-43)} KA\_I,101.2-16 Ro\_I,320-322 {1/25}     śi sarvanāmasthānam suṭ anapuṃsakasya iti cet jasi śipratiṣedhaḥ . \\
\noindent \textbf{(P\_1,1.42-43)} KA\_I,101.2-16 Ro\_I,320-322 {2/25}     śi sarvanāmasthānam suṭ anapuṃsakasya iti cet jasi śeḥ pratiṣedhaḥ prāpnoti : kuṇḍāni tiṣṭhanti vanāni tiṣṭhanti . \\
\noindent \textbf{(P\_1,1.42-43)} KA\_I,101.2-16 Ro\_I,320-322 {3/25}     asamarthasamāsaḥ ca ayam draṣṭavyaḥ anapuṃsakasya iti . \\
\noindent \textbf{(P\_1,1.42-43)} KA\_I,101.2-16 Ro\_I,320-322 {4/25}     na hi nañaḥ napuṃsakena sāmarthyam. kena tarhi . \\
\noindent \textbf{(P\_1,1.42-43)} KA\_I,101.2-16 Ro\_I,320-322 {5/25}     bhavatinā : na bhavati napuṃsakasya iti . \\
\noindent \textbf{(P\_1,1.42-43)} KA\_I,101.2-16 Ro\_I,320-322 {6/25}     yat tāvat ucyate śi sarvanāmasthānam suṭ anapuṃsakasya iti cet jasi śipratiṣedhaḥ iti . \\
\noindent \textbf{(P\_1,1.42-43)} KA\_I,101.2-16 Ro\_I,320-322 {7/25}     na apratiṣedhāt . \\
\noindent \textbf{(P\_1,1.42-43)} KA\_I,101.2-16 Ro\_I,320-322 {8/25}     na ayam prasajyapratiṣedhaḥ : napuṃsakasya na iti . \\
\noindent \textbf{(P\_1,1.42-43)} KA\_I,101.2-16 Ro\_I,320-322 {9/25}     kim tarhi . \\
\noindent \textbf{(P\_1,1.42-43)} KA\_I,101.2-16 Ro\_I,320-322 {10/25}     paryudāsaḥ ayam : yat anyat napuṃsakāt iti . \\
\noindent \textbf{(P\_1,1.42-43)} KA\_I,101.2-16 Ro\_I,320-322 {11/25}     napuṃsake avyāpāraḥ . \\
\noindent \textbf{(P\_1,1.42-43)} KA\_I,101.2-16 Ro\_I,320-322 {12/25}     yadi kena cit prāpnoti tena bhaviṣyati . \\
\noindent \textbf{(P\_1,1.42-43)} KA\_I,101.2-16 Ro\_I,320-322 {13/25}     pūrveṇa ca prāpnoti . \\
\noindent \textbf{(P\_1,1.42-43)} KA\_I,101.2-16 Ro\_I,320-322 {14/25}     aprāpteḥ vā . \\
\noindent \textbf{(P\_1,1.42-43)} KA\_I,101.2-16 Ro\_I,320-322 {15/25}     atha vā anantarā yā prāptiḥ sā pratiṣidhyate . \\
\noindent \textbf{(P\_1,1.42-43)} KA\_I,101.2-16 Ro\_I,320-322 {16/25}     kutaḥ etat . \\
\noindent \textbf{(P\_1,1.42-43)} KA\_I,101.2-16 Ro\_I,320-322 {17/25}     anantarasya vidhiḥ vā bhavati pratiṣedhaḥ vā iti . \\
\noindent \textbf{(P\_1,1.42-43)} KA\_I,101.2-16 Ro\_I,320-322 {18/25}     pūrvā prāptiḥ apratiṣiddhā . \\
\noindent \textbf{(P\_1,1.42-43)} KA\_I,101.2-16 Ro\_I,320-322 {19/25}     tayā bhaviṣyati . \\
\noindent \textbf{(P\_1,1.42-43)} KA\_I,101.2-16 Ro\_I,320-322 {20/25}     nanu ca iyam prāptiḥ pūrvām prāptim bādhate . \\
\noindent \textbf{(P\_1,1.42-43)} KA\_I,101.2-16 Ro\_I,320-322 {21/25}     na utsahate pratiṣiddhā satī bādhitum . \\
\noindent \textbf{(P\_1,1.42-43)} KA\_I,101.2-16 Ro\_I,320-322 {22/25}     yat api ucyate . \\
\noindent \textbf{(P\_1,1.42-43)} KA\_I,101.2-16 Ro\_I,320-322 {23/25}     asamarthasamāsaḥ ca ayam draṣṭavyaḥ iti yadi api vaktavyaḥ atha vā etarhi bahūni prayojanāni . \\
\noindent \textbf{(P\_1,1.42-43)} KA\_I,101.2-16 Ro\_I,320-322 {24/25}     kāni . \\
\noindent \textbf{(P\_1,1.42-43)} KA\_I,101.2-16 Ro\_I,320-322 {25/25}     asūryampaśyāni mukhāni , apunargeyāḥ ślokāḥ , aśrāddhabhojī brāhmaṇaḥ iti . \\
\noindent \textbf{(P\_1,1.44.1)} na vā iti vibhāṣāyām arthasañjñākaraṇam . \\
\noindent \textbf{(P\_1,1.44.1)} na vā iti vibhāṣāyām arthasya sañjñā kartavyā . \\
\noindent \textbf{(P\_1,1.44.1)} navāśabdasya yaḥ arthaḥ tasya sañjñā bhavati iti vaktavyam . \\
\noindent \textbf{(P\_1,1.44.1)} kim prayojanam . \\
\noindent \textbf{(P\_1,1.44.1)} śabdasañjñāyām hi arthāsampratyayaḥ yathā anyatra . \\
\noindent \textbf{(P\_1,1.44.1)} śabdasañjñāyām hi satyām arthasya asampratyayaḥ syāt yathā anyatra . \\
\noindent \textbf{(P\_1,1.44.1)} anyatra api śabdasañjñāyām śabdasya sampratyayaḥ bhavati na arthasya . \\
\noindent \textbf{(P\_1,1.44.1)} kva anyatra . \\
\noindent \textbf{(P\_1,1.44.1)} dādhāḥ ghu adāp taraptamapau ghaḥ iti ghugrahaṇeṣu ghagrahaṇeṣu ca śabdasya sampratyayaḥ bhavati na arthasya . \\
\noindent \textbf{(P\_1,1.44.1)} tat tarhi vaktavyam . \\
\noindent \textbf{(P\_1,1.44.1)} na vaktavyam . \\
\noindent \textbf{(P\_1,1.44.1)} itikaraṇaḥ arthanirdeśārthaḥ . \\
\noindent \textbf{(P\_1,1.44.1)} itikaraṇaḥ kriyate . \\
\noindent \textbf{(P\_1,1.44.1)} saḥ arthanirdeśārthaḥ bhaviṣyati . \\
\noindent \textbf{(P\_1,1.44.1)} kim gatam etat itinā āhosvit śabādhikyāt arthādhikyam . \\
\noindent \textbf{(P\_1,1.44.1)} gatam iti āha . \\
\noindent \textbf{(P\_1,1.44.1)} kutaḥ . \\
\noindent \textbf{(P\_1,1.44.1)} lokataḥ . \\
\noindent \textbf{(P\_1,1.44.1)} tat yathā loke gauḥ ayam iti āha iti gośabdāt itikaraṇaḥ paraḥ prayujyamānaḥ gośabdam svasmāt padārthāt pracyāvayati . \\
\noindent \textbf{(P\_1,1.44.1)} saḥ asau svasmāt padārthāt pracyutaḥ yā asau arthapadārthakatā tasyāḥ śabdapadārthakaḥ sampadyate . \\
\noindent \textbf{(P\_1,1.44.1)} evam iha api navāśabdāt itikaraṇaḥ paraḥ prayujyamānaḥ navāśabdam svasmāt padārthāt pracyāvayati . \\
\noindent \textbf{(P\_1,1.44.1)} saḥ asau svasmāt padārthāt pracyutaḥ yā asau śabdapadārthakatā tasyāḥ laukikam artham sampratyāyayati . \\
\noindent \textbf{(P\_1,1.44.1)} na vā iti yat gamyate na vā iti yat pratīyate iti . \\
\noindent \textbf{(P\_1,1.44.2)} samānaśabdapratiṣedhaḥ . \\
\noindent \textbf{(P\_1,1.44.2)} samānaśabdānām pratiṣedhaḥ vaktavyaḥ : navā kuṇḍikā navā ghaṭikā iti . \\
\noindent \textbf{(P\_1,1.44.2)} kim ca syāt yadi eteṣām api vibhāṣāsañjñā syāt . \\
\noindent \textbf{(P\_1,1.44.2)} vibhāṣā diksamāse bahuvrīhau : dakṣiṇapūrvasyām śālāyām . \\
\noindent \textbf{(P\_1,1.44.2)} acirakṛtāyām sampratyayaḥ syāt . \\
\noindent \textbf{(P\_1,1.44.2)} na vā vidhipūrvakatvāt pratiṣedhasamapratyayaḥ yathā loke . \\
\noindent \textbf{(P\_1,1.44.2)} na vā eṣaḥ doṣaḥ . \\
\noindent \textbf{(P\_1,1.44.2)} kim kāraṇam . \\
\noindent \textbf{(P\_1,1.44.2)} vidhipūrvakatvāt . \\
\noindent \textbf{(P\_1,1.44.2)} vidhāya kim cit na vā iti ucyate . \\
\noindent \textbf{(P\_1,1.44.2)} tena pratiṣedhavācinaḥ samapratyayaḥ bhavati . \\
\noindent \textbf{(P\_1,1.44.2)} tat yathā loke : grāmaḥ bhavatā gantavyaḥ na vā . \\
\noindent \textbf{(P\_1,1.44.2)} na iti gamyate . \\
\noindent \textbf{(P\_1,1.44.2)} asti kārṇam yena loke pratiṣedhavācinaḥ samapratyayaḥ bhavati . \\
\noindent \textbf{(P\_1,1.44.2)} kim kāraṇam . \\
\noindent \textbf{(P\_1,1.44.2)} viliṅgam hi bhavān loke nirdeśam karoti . \\
\noindent \textbf{(P\_1,1.44.2)} āṅga hi samānaliṅgaḥ nirdeśaḥ kriyatām pratyagravācinaḥ sampratyayaḥ bhaviṣyati . \\
\noindent \textbf{(P\_1,1.44.2)} tat yathā : grāmaḥ bhavatā gantavyaḥ navaḥ . \\
\noindent \textbf{(P\_1,1.44.2)} pratyagraḥ iti gamyate . \\
\noindent \textbf{(P\_1,1.44.2)} etat ca eva na jānīmaḥ : kva cit vyākaraṇe samānaliṅgaḥ nirdeśaḥ kriyate iti . \\
\noindent \textbf{(P\_1,1.44.2)} api ca kāmacāraḥ prayoktuḥ śabdānām abhisambandhe . \\
\noindent \textbf{(P\_1,1.44.2)} tat yathā : yavāgūḥ bhavatā bhoktavyā navā . \\
\noindent \textbf{(P\_1,1.44.2)} yadā yavāgūśabdaḥ bhujinā abhisambadhyate bhujiḥ navāśabdena tadā pratiṣedhavācinaḥ sampratyayaḥ bhavati : yavāgūḥ bhavatā bhoktavyā navā . \\
\noindent \textbf{(P\_1,1.44.2)} na iti gamyate . \\
\noindent \textbf{(P\_1,1.44.2)} yadā yavāgūśabdaḥ navāśabdena abhisambadhyate na bhujinā tadā pratyagravācinaḥ sampratyayaḥ bhavati : yavāgūḥ navā bhavatā bhoktavyā . \\
\noindent \textbf{(P\_1,1.44.2)} pratyagrā iti gamyate . \\
\noindent \textbf{(P\_1,1.44.2)} na ca iha vayam vibhāṣāgrahaṇena sarvādīni abhisambadhnīmaḥ : diksamāse bahuvrīhau sarvādīni vibhāṣā bhavanti iti . \\
\noindent \textbf{(P\_1,1.44.2)} kim tarhi . \\
\noindent \textbf{(P\_1,1.44.2)} bhavatiḥ abhisambadhyate : diksamāse bahuvrīhau sarvādīni bhavanti vibhāṣā iti . \\
\noindent \textbf{(P\_1,1.44.3)} vidhyanityatvam anupapannam pratiṣedhasañjñākaraṇāt . \\
\noindent \textbf{(P\_1,1.44.3)} vidhyanityatvam na upapadyate : śuśāva , śuśuvatuḥ , śuśuvuḥ , śiśvāya , śiśviyatuḥ , śiśviyuḥ . \\
\noindent \textbf{(P\_1,1.44.3)} kim kāraṇam . \\
\noindent \textbf{(P\_1,1.44.3)} pratiṣedhasañjñākaraṇāt . \\
\noindent \textbf{(P\_1,1.44.3)} pratiṣedhasya iyam sañjñā kriyate . \\
\noindent \textbf{(P\_1,1.44.3)} tena vibhāṣāpradeśeṣu pratiṣedhasya eva sampratyayaḥ syāt . \\
\noindent \textbf{(P\_1,1.44.3)} siddham tu prasajyapratiṣedhāt . \\
\noindent \textbf{(P\_1,1.44.3)} siddham etat . \\
\noindent \textbf{(P\_1,1.44.3)} katham . \\
\noindent \textbf{(P\_1,1.44.3)} prasajyapratiṣedhāt . \\
\noindent \textbf{(P\_1,1.44.3)} prasajya kim cit na vā iti ucyate . \\
\noindent \textbf{(P\_1,1.44.3)} tena ubhayam bhaviṣyati . \\
\noindent \textbf{(P\_1,1.44.3)} vipratiṣiddham tu . \\
\noindent \textbf{(P\_1,1.44.3)} vipratiṣiddham tu bhavati . \\
\noindent \textbf{(P\_1,1.44.3)} atra na jñāyate : kena abhiprāyeṇa prasajati kena nivṛttim karoti iti . \\
\noindent \textbf{(P\_1,1.44.3)} na vā prasaṅgasāmarthyāt anyatra pratiṣedhaviṣayāt . \\
\noindent \textbf{(P\_1,1.44.3)} na vā eṣaḥ doṣaḥ . \\
\noindent \textbf{(P\_1,1.44.3)} kim kāraṇam . \\
\noindent \textbf{(P\_1,1.44.3)} prasaṅgasāmarthyāt . \\
\noindent \textbf{(P\_1,1.44.3)} prasaṅgasāmarthyāt ca vidhiḥ bhaviṣyati anyatra pratiṣedhaviṣayāt pratiṣedhasāmarthyāt ca pratiṣedhaḥ bhaviṣyati anyatra vidhiviṣayāt . \\
\noindent \textbf{(P\_1,1.44.3)} tat etat kva siddham bhavati . \\
\noindent \textbf{(P\_1,1.44.3)} yā aprāpte vibhāṣā . \\
\noindent \textbf{(P\_1,1.44.3)} yā hi prāpte kṛtasāmarthyaḥ tatra pūrveṇa vidhiḥ iti kṛtvā pratiṣedhasya eva sampratyayaḥ syāt . \\
\noindent \textbf{(P\_1,1.44.3)} etat api siddham . \\
\noindent \textbf{(P\_1,1.44.3)} katham . \\
\noindent \textbf{(P\_1,1.44.3)} vibhāṣā iti mahatīsañjñā kriyate . \\
\noindent \textbf{(P\_1,1.44.3)} sañjñā ca nāma yataḥ na laghīyaḥ . \\
\noindent \textbf{(P\_1,1.44.3)} kutaḥ etat . \\
\noindent \textbf{(P\_1,1.44.3)} laghvartham hi sañjñākaraṇam . \\
\noindent \textbf{(P\_1,1.44.3)} tatra mahatyāḥ sañjñāyāḥ karaṇe etat prayojanam ubhayoḥ sañjñā yathā vijñāyeta : na iti ca vā iti ca . \\
\noindent \textbf{(P\_1,1.44.3)} tatra yā tāvat aprāpte vibhāṣā tatra pratiṣedhyam na asti iti kṛtvā vā iti anena vikalpaḥ bhaviṣyati . \\
\noindent \textbf{(P\_1,1.44.3)} yā hi prāpte vibhāṣā tatra ubhayam upasthitam bhavati : na iti ca vā iti ca . \\
\noindent \textbf{(P\_1,1.44.3)} tatra na iti anena pratiṣiddhe vā iti anena vikalpaḥ bhaviṣyati . \\
\noindent \textbf{(P\_1,1.44.3)} evam api vipratiṣedhayoḥ yugapadvacanānupapattiḥ . \\
\noindent \textbf{(P\_1,1.44.3)} vipratiṣedhayoḥ yugapadvacanam na upapadyate : śuśāva śuśuvatuḥ śuśuvuḥ śiśvāya śiśviyatuḥ śiśviyuḥ . \\
\noindent \textbf{(P\_1,1.44.3)} kim kāraṇam . \\
\noindent \textbf{(P\_1,1.44.3)} bhavati iti cet na pratiṣedhaḥ . \\
\noindent \textbf{(P\_1,1.44.3)} bhavati iti cet pratiṣedhaḥ na prāpnoti . \\
\noindent \textbf{(P\_1,1.44.3)} na iti cet na vidhiḥ . \\
\noindent \textbf{(P\_1,1.44.3)} na it cet vidhiḥ na sidhyati . \\
\noindent \textbf{(P\_1,1.44.3)} siddham tu pūrvasya uttareṇa bādhitatvāt . \\
\noindent \textbf{(P\_1,1.44.3)} siddham etat. katham . \\
\noindent \textbf{(P\_1,1.44.3)} pūrvavidhim uttarvidhiḥ bādhate . \\
\noindent \textbf{(P\_1,1.44.3)} itikaraṇaḥ arthanirdeśāṛthaḥ iti uktam . \\
\noindent \textbf{(P\_1,1.44.4)} sādhvanuśāsane asmin yasya vibhāṣā tasya sādhutvam . \\
\noindent \textbf{(P\_1,1.44.4)} sādhvanuśāsane asmin śāstre yasya vibhāṣā kriyate saḥ vibhāṣā sādhuḥ syāt . \\
\noindent \textbf{(P\_1,1.44.4)} samāsaḥ ca eva hi vibhāṣā . \\
\noindent \textbf{(P\_1,1.44.4)} tena samāsasya eva vibhāṣā sādhutvam syāt . \\
\noindent \textbf{(P\_1,1.44.4)} astu . \\
\noindent \textbf{(P\_1,1.44.4)} yaḥ sādhuḥ saḥ prayokṣyate. asādhuḥ na prayokṣyate . \\
\noindent \textbf{(P\_1,1.44.4)} na ca eva hi kadā cit rājapuruṣaḥ iti asyām avasthāyām asādhutvam iṣyate . \\
\noindent \textbf{(P\_1,1.44.4)} api ca dvedhāpratipattiḥ . \\
\noindent \textbf{(P\_1,1.44.4)} dvaidham śabdānām apratipattiḥ . \\
\noindent \textbf{(P\_1,1.44.4)} icchāmaḥ ca punaḥ vibhāṣāpradeśeṣu dvaidham śabdānām pratipattiḥ syāt iti tat ca na sidhyati . \\
\noindent \textbf{(P\_1,1.44.4)} yasya punaḥ kāryāḥ śabdāḥ vibhāṣā asau samāsam nirvartayati . \\
\noindent \textbf{(P\_1,1.44.4)} yasya api nityāḥ śabdāḥ tasya api eṣaḥ na doṣaḥ . \\
\noindent \textbf{(P\_1,1.44.4)} katham . \\
\noindent \textbf{(P\_1,1.44.4)} na vibhāṣāgrahaṇena sādhutvam abhisambadhyate . \\
\noindent \textbf{(P\_1,1.44.4)} kim tarhi . \\
\noindent \textbf{(P\_1,1.44.4)} samāsasañjñā abhisambadhyate : samāsaḥ iti eṣā sañjñā vibhāṣā bhavati iti . \\
\noindent \textbf{(P\_1,1.44.4)} tat yathā : medhyaḥ paśuḥ vibhāṣitaḥ . \\
\noindent \textbf{(P\_1,1.44.4)} medhyaḥ anaḍvān vibhāṣitaḥ iti . \\
\noindent \textbf{(P\_1,1.44.4)} na etat vicāryate : anaḍvān na anaḍvān iti . \\
\noindent \textbf{(P\_1,1.44.4)} kim tarhi ālabdhavyaḥ na ālabdhavyaḥ iti . \\
\noindent \textbf{(P\_1,1.44.4)} kārye yugapadanvācayayaugapadyam . \\
\noindent \textbf{(P\_1,1.44.4)} kāryeṣu śabdeṣu yugapat anvācayena ca yat ucyate tasya yugapadvacanatā prāpnoti : tavyattavyānīyaraḥ , ḍhak ca maṇḍūkāt iti . \\
\noindent \textbf{(P\_1,1.44.4)} yasya punaḥ nityāḥ śabdāḥ prayuktānām asau sādhutvam anvācaṣṭe . \\
\noindent \textbf{(P\_1,1.44.4)} nanu ca yasya api kāryāḥ tasya api eṣaḥ na doṣaḥ . \\
\noindent \textbf{(P\_1,1.44.4)} katham . \\
\noindent \textbf{(P\_1,1.44.4)} pratyayaḥ paraḥ bhavati iti ucyate . \\
\noindent \textbf{(P\_1,1.44.4)} na ca ekasyāḥ prakṛteḥ anekasya pratyayasya yugapat paratvena sambhavaḥ asti . \\
\noindent \textbf{(P\_1,1.44.4)} na api brūmaḥ pratyayamālā prāpnoti . \\
\noindent \textbf{(P\_1,1.44.4)} kim tarhi . \\
\noindent \textbf{(P\_1,1.44.4)} kartavyam iti prayoktavye yugapat dvitīyasya tṛtīyasya ca prayogaḥ prāpnoti . \\
\noindent \textbf{(P\_1,1.44.4)} na eṣaḥ doṣaḥ . \\
\noindent \textbf{(P\_1,1.44.4)} arthagatyarthaḥ śabdaprayogaḥ . \\
\noindent \textbf{(P\_1,1.44.4)} artham sampratyāyayiṣyāmi iti śabdaḥ prayujyate . \\
\noindent \textbf{(P\_1,1.44.4)} tatra ekena uktatvāt tasya arthasya dvitīyasya prayogeṇa na bhavitavyam uktārthānām aprayogaḥ iti . \\
\noindent \textbf{(P\_1,1.44.4)} ācāryadeśaśīlane ca tadviṣayatā . \\
\noindent \textbf{(P\_1,1.44.4)} ācāryadeśaśīlanena yat ucyate tasya tadviṣayatā prāpnoti . \\
\noindent \textbf{(P\_1,1.44.4)} ikaḥ hrasvaḥ aṅyaḥ gālavasya prācām avṛddhāt phin bahulam iti gālavāḥ eva hrasvān prayuñjīran prākṣu ca eva hi phin syāt . \\
\noindent \textbf{(P\_1,1.44.4)} tat yathā : jamadagniḥ vai etat pañcamam avadānam avādyat tasmāt na ajāmadagnyaḥ pañcāvattam juhoti . \\
\noindent \textbf{(P\_1,1.44.4)} yasya punaḥ nityāḥ śabdāḥ gālavagrahaṇam tasya pūjārtham deśagrahaṇam ca kīrtyartham . \\
\noindent \textbf{(P\_1,1.44.4)} nanu ca yasya api kāryāḥ tasya api pūjārtham gālavagrahaṇam syāt deśagraham ca kīrtyartham . \\
\noindent \textbf{(P\_1,1.44.4)} tatkīrtane ca dvedhāpratipattiḥ . \\
\noindent \textbf{(P\_1,1.44.4)} tatkīrtane ca dvaidham śabdānām apratipattiḥ syāt . \\
\noindent \textbf{(P\_1,1.44.4)} icchāmaḥ ca punaḥ ācāryagrahaṇeṣu deśagrahaṇeṣu ca dvaidham śabdānām pratipattiḥ syāt iti tat ca na sidhyati . \\
\noindent \textbf{(P\_1,1.44.5)} aśiṣyaḥ vā viditatvāt . \\
\noindent \textbf{(P\_1,1.44.5)} aśiṣyaḥ vā punaḥ ayam yogaḥ . \\
\noindent \textbf{(P\_1,1.44.5)} kim kāraṇam . \\
\noindent \textbf{(P\_1,1.44.5)} viditatvāt . \\
\noindent \textbf{(P\_1,1.44.5)} yat anena yogena prārthyate tasya arthasya viditatvāt . \\
\noindent \textbf{(P\_1,1.44.5)} ye api hi etām sañjñām na ārabhante te api vibhāṣā iti ukte anityatvam avagacchanti . \\
\noindent \textbf{(P\_1,1.44.5)} yājñikāḥ khalu api sañjñām anārabhamāṇāḥ vibhāṣā iti ukte anityatvam avagacchanti . \\
\noindent \textbf{(P\_1,1.44.5)} tat yathā . \\
\noindent \textbf{(P\_1,1.44.5)} medhyaḥ paśuḥ vibhāṣitaḥ . \\
\noindent \textbf{(P\_1,1.44.5)} medhyaḥ anaḍvān vibhāṣitaḥ iti . \\
\noindent \textbf{(P\_1,1.44.5)} ālabdhavyaḥ na ālabdhavyaḥ iti gamyate . \\
\noindent \textbf{(P\_1,1.44.5)} ācāryaḥ khalu api sañjñām ārabhamāṇaḥ bhūyiṣṭham anyaiḥ api śabdaiḥ etam artham sampratyāyayati bahulam anyatarasyām ubhayathā vā ekeṣām iti . \\
\noindent \textbf{(P\_1,1.44.6)} aprāpte trisaṃśayāḥ . \\
\noindent \textbf{(P\_1,1.44.6)} itaḥ uttaram yāḥ vibhāṣāḥ anukramiṣyāmaḥ aprāpte tāḥ draṣṭavyāḥ . \\
\noindent \textbf{(P\_1,1.44.6)} trisaṃśayāḥ tu bhavanti : prāpte aprāpte ubhayatra vā iti . \\
\noindent \textbf{(P\_1,1.44.6)} dvandve ca vibhāṣā jasi : prāpte aprāpte ubhayatra vā iti sandehaḥ . \\
\noindent \textbf{(P\_1,1.44.6)} katham ca prāpte katham vā aprāpte katham vā ubhayatra . \\
\noindent \textbf{(P\_1,1.44.6)} ubhayaśabdaḥ sarvādiṣu paṭhyate tayapaḥ ca ayajādeśaḥ kriyate . \\
\noindent \textbf{(P\_1,1.44.6)} tena vā nitye prāpte anyatra vā aprāpte ubhayatra vā iti . \\
\noindent \textbf{(P\_1,1.44.6)} aprāpte . \\
\noindent \textbf{(P\_1,1.44.6)} ayac pratyayāntaram . \\
\noindent \textbf{(P\_1,1.44.6)} yadi pratyayāntaram ubhayī iti īkāraḥ na prāpnoti . \\
\noindent \textbf{(P\_1,1.44.6)} mā bhūt evam . \\
\noindent \textbf{(P\_1,1.44.6)} mātracaḥ iti evam bhaviṣyati . \\
\noindent \textbf{(P\_1,1.44.6)} katham . \\
\noindent \textbf{(P\_1,1.44.6)} mātrac iti na idam pratyayagrahaṇam . \\
\noindent \textbf{(P\_1,1.44.6)} kim tarhi . \\
\noindent \textbf{(P\_1,1.44.6)} pratyāhāragrahaṇam . \\
\noindent \textbf{(P\_1,1.44.6)} kva sanniviṣṭānām pratyāhāraḥ . \\
\noindent \textbf{(P\_1,1.44.6)} mātraśabdāt prabhṛti ā āyacaḥ cakārāt . \\
\noindent \textbf{(P\_1,1.44.6)} yadi pratyāhāragrahaṇam kati tiṣṭhanti : atra api prāpnoti . \\
\noindent \textbf{(P\_1,1.44.6)} ataḥ iti vartate . \\
\noindent \textbf{(P\_1,1.44.6)} evam api tailamātrā ghrtamātrā iti atra api prāpnoti . \\
\noindent \textbf{(P\_1,1.44.6)} sadṛśasya api asanniviṣṭasya na bhaviṣyati pratyāhāreṇa grahaṇam . \\
\noindent \textbf{(P\_1,1.44.6)} ūrṇoḥ vibhāṣā : prāpte aprāpte ubhayatra vā iti sandehaḥ . \\
\noindent \textbf{(P\_1,1.44.6)} katham ca prāpte katham vā aprāpte katham vā ubhayatra . \\
\noindent \textbf{(P\_1,1.44.6)} asaṃyogāt liṭ kit iti vā nitye prāpte anyatra vā aprāpte ubhayatra vā iti . \\
\noindent \textbf{(P\_1,1.44.6)} aprāpte . \\
\noindent \textbf{(P\_1,1.44.6)} anyat hi kittvam anyat ṅittvam . \\
\noindent \textbf{(P\_1,1.44.6)} ekam cet ṅitkitau . \\
\noindent \textbf{(P\_1,1.44.6)} yadi ekam ṅitkitau tataḥ asti sandehaḥ . \\
\noindent \textbf{(P\_1,1.44.6)} atha hi nānā na asti sandehaḥ . \\
\noindent \textbf{(P\_1,1.44.6)} yadi api nānā evam api sandehaḥ . \\
\noindent \textbf{(P\_1,1.44.6)} katham . \\
\noindent \textbf{(P\_1,1.44.6)} praurṇuvi iti . \\
\noindent \textbf{(P\_1,1.44.6)} sārvadhātukam apit iti vā nitye prāpte anyatra vā aprāpte ubhayatra vā iti . \\
\noindent \textbf{(P\_1,1.44.6)} aprāpte . \\
\noindent \textbf{(P\_1,1.44.6)} vibhāṣā upayamane : prāpte aprāpte ubhayatra vā iti sandehaḥ . \\
\noindent \textbf{(P\_1,1.44.6)} katham ca prāpte katham vā aprāpte katham vā ubhayatra . \\
\noindent \textbf{(P\_1,1.44.6)} gandhane iti vā nitye prāpte anyatra vā aprāpte ubhayatra vā iti . \\
\noindent \textbf{(P\_1,1.44.6)} aprāpte . \\
\noindent \textbf{(P\_1,1.44.6)} gandhane iti nivṛttam . \\
\noindent \textbf{(P\_1,1.44.6)} anupasargāt vā : prāpte aprāpte ubhayatra vā iti sandehaḥ . \\
\noindent \textbf{(P\_1,1.44.6)} katham ca prāpte katham vā aprāpte katham vā ubhayatra . \\
\noindent \textbf{(P\_1,1.44.6)} vṛttisargatāyaneṣu kramaḥ iti vā nitye prāpte anyatra vā aprāpte ubhayatra vā iti . \\
\noindent \textbf{(P\_1,1.44.6)} aprāpte . \\
\noindent \textbf{(P\_1,1.44.6)} vṛttyādiṣu iti nivṛttam . \\
\noindent \textbf{(P\_1,1.44.6)} vibhāṣā vṛkṣamṛgādīnām : prāpte aprāpte ubhayatra vā iti sandehaḥ . \\
\noindent \textbf{(P\_1,1.44.6)} katham ca prāpte katham vā aprāpte katham vā ubhayatra . \\
\noindent \textbf{(P\_1,1.44.6)} jātiḥ aprāṇinām iti vā nitye prāpte anyatra vā aprāpte ubhayatra vā iti . \\
\noindent \textbf{(P\_1,1.44.6)} aprāpte . \\
\noindent \textbf{(P\_1,1.44.6)} jātiḥ aprāṇinām iti nivṛttam . \\
\noindent \textbf{(P\_1,1.44.6)} uṣavidajāgṛbhyaḥ anyatarasyām : prāpte aprāpte ubhayatra vā iti sandehaḥ . \\
\noindent \textbf{(P\_1,1.44.6)} katham ca prāpte katham vā aprāpte katham vā ubhayatra . \\
\noindent \textbf{(P\_1,1.44.6)} pratyayāntāt iti vā nitye prāpte anyatra vā aprāpte ubhayatra vā iti . \\
\noindent \textbf{(P\_1,1.44.6)} aprāpte . \\
\noindent \textbf{(P\_1,1.44.6)} pratyayāntāḥ dhātvantarāṇi . \\
\noindent \textbf{(P\_1,1.44.6)} dīpādīnām vibhāṣā : prāpte aprāpte ubhayatra vā iti sandehaḥ . \\
\noindent \textbf{(P\_1,1.44.6)} katham ca prāpte katham vā aprāpte katham vā ubhayatra . \\
\noindent \textbf{(P\_1,1.44.6)} bhāvakarmaṇoḥ iti vā nitye prāpte anyatra vā aprāpte ubhayatra vā iti . \\
\noindent \textbf{(P\_1,1.44.6)} aprāpte . \\
\noindent \textbf{(P\_1,1.44.6)} kartari iti vartate . \\
\noindent \textbf{(P\_1,1.44.6)} evam api sandehaḥ : nyāyye vā kartari karmakartari vā iti . \\
\noindent \textbf{(P\_1,1.44.6)} na asti sandehaḥ . \\
\noindent \textbf{(P\_1,1.44.6)} sakarmakasya kartā karmavat bhavati akarmakāḥ ca dīpādayaḥ . \\
\noindent \textbf{(P\_1,1.44.6)} akarmakāḥ api vai sopasargāḥ sakarmakāḥ bhavanti . \\
\noindent \textbf{(P\_1,1.44.6)} karmāpadiṣṭāḥ vidhayaḥ karmasthabhāvakānām karmasthakriyāṇām ca bhavanti kartṛsthabhāvakāḥ ca dīpādayaḥ . \\
\noindent \textbf{(P\_1,1.44.6)} vibhāṣā agreprathamapūrveṣu : prāpte aprāpte ubhayatra vā iti sandehaḥ . \\
\noindent \textbf{(P\_1,1.44.6)} katham ca prāpte katham vā aprāpte katham vā ubhayatra . \\
\noindent \textbf{(P\_1,1.44.6)} ābhīkṣṇye iti vā nitye prāpte anyatra vā aprāpte ubhayatra vā iti . \\
\noindent \textbf{(P\_1,1.44.6)} aprāpte . \\
\noindent \textbf{(P\_1,1.44.6)} ābhīkṣṇye iti nivṛttam . \\
\noindent \textbf{(P\_1,1.44.6)} tṛnādīnām vibhāṣā : prāpte aprāpte ubhayatra vā iti sandehaḥ . \\
\noindent \textbf{(P\_1,1.44.6)} katham ca prāpte katham vā aprāpte katham vā ubhayatra . \\
\noindent \textbf{(P\_1,1.44.6)} ākrośe iti vā nitye prāpte anyatra vā aprāpte ubhayatra vā iti . \\
\noindent \textbf{(P\_1,1.44.6)} aprāpte . \\
\noindent \textbf{(P\_1,1.44.6)} ākrośe iti nivṛttam . \\
\noindent \textbf{(P\_1,1.44.6)} ekahalādau pūrayitavye anyatarasyām . \\
\noindent \textbf{(P\_1,1.44.6)} prāpte aprāpte ubhayatra vā iti sandehaḥ . \\
\noindent \textbf{(P\_1,1.44.6)} katham ca prāpte katham vā aprāpte katham vā ubhayatra . \\
\noindent \textbf{(P\_1,1.44.6)} udakasya udaḥ sañjñāyām iti vā nitye prāpte anyatra vā aprāpte ubhayatra vā iti . \\
\noindent \textbf{(P\_1,1.44.6)} aprāpte . \\
\noindent \textbf{(P\_1,1.44.6)} sañjñāyām iti nivṛttam . \\
\noindent \textbf{(P\_1,1.44.6)} śvādeḥ iñi padāntasya anyatarasyām . \\
\noindent \textbf{(P\_1,1.44.6)} prāpte aprāpte ubhayatra vā iti sandehaḥ . \\
\noindent \textbf{(P\_1,1.44.6)} katham ca prāpte katham vā aprāpte katham vā ubhayatra . \\
\noindent \textbf{(P\_1,1.44.6)} iñi iti vā nitye prāpte anyatra vā aprāpte ubhayatra vā iti . \\
\noindent \textbf{(P\_1,1.44.6)} aprāpte . \\
\noindent \textbf{(P\_1,1.44.6)} iñi iti nivṛttam . \\
\noindent \textbf{(P\_1,1.44.6)} sapūrvāyāḥ prathamāyāḥ vibhāṣā . \\
\noindent \textbf{(P\_1,1.44.6)} prāpte aprāpte ubhayatra vā iti sandehaḥ . \\
\noindent \textbf{(P\_1,1.44.6)} katham ca prāpte katham vā aprāpte katham vā ubhayatra . \\
\noindent \textbf{(P\_1,1.44.6)} cādibhiḥ yoge iti vā nitye prāpte anyatra vā aprāpte ubhayatra vā iti . \\
\noindent \textbf{(P\_1,1.44.6)} aprāpte . \\
\noindent \textbf{(P\_1,1.44.6)} cādibhiḥ yoge iti nivṛttam . \\
\noindent \textbf{(P\_1,1.44.6)} graḥ yaṅi aci vibhāṣā . \\
\noindent \textbf{(P\_1,1.44.6)} prāpte aprāpte ubhayatra vā iti sandehaḥ . \\
\noindent \textbf{(P\_1,1.44.6)} katham ca prāpte katham vā aprāpte katham vā ubhayatra . \\
\noindent \textbf{(P\_1,1.44.6)} yaṅi iti vā nitye prāpte anyatra vā aprāpte ubhayatra vā iti . \\
\noindent \textbf{(P\_1,1.44.6)} aprāpte . \\
\noindent \textbf{(P\_1,1.44.6)} aprāpte . \\
\noindent \textbf{(P\_1,1.44.6)} yaṅi iti nivṛttam \\
\noindent \textbf{(P\_1,1.44.7)} prāpte ca . \\
\noindent \textbf{(P\_1,1.44.7)} itaḥ uttaram yāḥ vibhāṣāḥ anukramiṣyāmaḥ prāpte tāḥ draṣṭavyāḥ . \\
\noindent \textbf{(P\_1,1.44.7)} trisaṃśayāḥ tu bhavanti : prāpte aprāpte ubhayatra vā iti . \\
\noindent \textbf{(P\_1,1.44.7)} vibhāṣā vipralāpe : prāpte aprāpte ubhayatra vā iti sandehaḥ . \\
\noindent \textbf{(P\_1,1.44.7)} katham ca prāpte katham vā aprāpte katham vā ubhayatra . \\
\noindent \textbf{(P\_1,1.44.7)} vyaktavācām iti vā nitye prāpte anyatra vā aprāpte ubhayatra vā iti . \\
\noindent \textbf{(P\_1,1.44.7)} prāpte . \\
\noindent \textbf{(P\_1,1.44.7)} vyaktavācām iti hi vartate . \\
\noindent \textbf{(P\_1,1.44.7)} vibhāṣā upapadena pratīyamāne : prāpte aprāpte ubhayatra vā iti sandehaḥ . \\
\noindent \textbf{(P\_1,1.44.7)} katham ca prāpte katham vā aprāpte katham vā ubhayatra . \\
\noindent \textbf{(P\_1,1.44.7)} svaritañitaḥ iti vā nitye prāpte anyatra vā aprāpte ubhayatra vā iti . \\
\noindent \textbf{(P\_1,1.44.7)} prāpte . \\
\noindent \textbf{(P\_1,1.44.7)} svaritañitaḥ iti hi vartate . \\
\noindent \textbf{(P\_1,1.44.7)} tiraḥ antardhau vibhāṣā kṛñi : prāpte aprāpte ubhayatra vā iti sandehaḥ . \\
\noindent \textbf{(P\_1,1.44.7)} katham ca prāpte katham vā aprāpte katham vā ubhayatra . \\
\noindent \textbf{(P\_1,1.44.7)} antardhau iti vā nitye prāpte anyatra vā aprāpte ubhayatra vā iti . \\
\noindent \textbf{(P\_1,1.44.7)} prāpte . \\
\noindent \textbf{(P\_1,1.44.7)} antardhau iti hi vartate . \\
\noindent \textbf{(P\_1,1.44.7)} adhiḥ īśvare vibhāṣā kṛñi : prāpte aprāpte ubhayatra vā iti sandehaḥ . \\
\noindent \textbf{(P\_1,1.44.7)} katham ca prāpte katham vā aprāpte katham vā ubhayatra . \\
\noindent \textbf{(P\_1,1.44.7)} īśvare iti vā nitye prāpte anyatra vā aprāpte ubhayatra vā iti . \\
\noindent \textbf{(P\_1,1.44.7)} prāpte . \\
\noindent \textbf{(P\_1,1.44.7)} īśvare iti hi vartate . \\
\noindent \textbf{(P\_1,1.44.7)} divaḥ tadarthasya vibhāṣā upasarge : prāpte aprāpte ubhayatra vā iti sandehaḥ . \\
\noindent \textbf{(P\_1,1.44.7)} katham ca prāpte katham vā aprāpte katham vā ubhayatra . \\
\noindent \textbf{(P\_1,1.44.7)} tadarthasya iti vā nitye prāpte anyatra vā aprāpte ubhayatra vā iti . \\
\noindent \textbf{(P\_1,1.44.7)} prāpte . \\
\noindent \textbf{(P\_1,1.44.7)} tadarthasya iti vartate \\
\noindent \textbf{(P\_1,1.44.8)} ubhayatra ca . \\
\noindent \textbf{(P\_1,1.44.8)} itaḥ uttaram yāḥ vibhāṣāḥ anukramiṣyāmaḥ ubhayatra tāḥ draṣṭavyāḥ . \\
\noindent \textbf{(P\_1,1.44.8)} trisaṃśayāḥ tu bhavanti : prāpte aprāpte ubhayatra vā iti . \\
\noindent \textbf{(P\_1,1.44.8)} hṛkroḥ anyatarasyām : prāpte aprāpte ubhayatra vā iti sandehaḥ . \\
\noindent \textbf{(P\_1,1.44.8)} katham ca prāpte katham vā aprāpte katham vā ubhayatra . \\
\noindent \textbf{(P\_1,1.44.8)} gatibuddhipratyavasānārthaśabdakarmākarmakāṇām iti vā nitye prāpte anyatra vā aprāpte ubhayatra vā iti . \\
\noindent \textbf{(P\_1,1.44.8)} ubhayatra . \\
\noindent \textbf{(P\_1,1.44.8)} prāpte tāvat : abhyavahārayati saindhavān , abhyavahārayati saindhavaiḥ , vikārayati saindhavān , vikārayati saindhavaiḥ . \\
\noindent \textbf{(P\_1,1.44.8)} aprāpte : harati bhāram devadattaḥ . \\
\noindent \textbf{(P\_1,1.44.8)} hārayati bhāram devadattam , hārayati bhāram devadattena . \\
\noindent \textbf{(P\_1,1.44.8)} karoti kaṭam devadattaḥ . \\
\noindent \textbf{(P\_1,1.44.8)} kārayati kaṭam devadattam , kārayati kaṭam devadattena . \\
\noindent \textbf{(P\_1,1.44.8)} na yadi vibhāṣā sākāṅkṣe : prāpte aprāpte ubhayatra vā iti sandehaḥ . \\
\noindent \textbf{(P\_1,1.44.8)} katham ca prāpte katham vā aprāpte katham vā ubhayatra .yadi iti vā nitye prāpte anyatra vā aprāpte ubhayatra vā iti . \\
\noindent \textbf{(P\_1,1.44.8)} ubhayatra . \\
\noindent \textbf{(P\_1,1.44.8)} prāpte tāvat : abhijānāsi devadatta yat kaśmīreṣu vatsyāmaḥ , yat kaśmīreṣu avasāma , yat tatra odanān bhokṣyāmahe , yat tatra odanān abhuñjmahi . \\
\noindent \textbf{(P\_1,1.44.8)} aprāpte : abhijānāsi devadatta kaśmīrān gamiṣyāmaḥ , kaśmīrān agacchāma , tatra odanām bhokṣyāmahe , tatra odanān abhuñjmahi . \\
\noindent \textbf{(P\_1,1.44.8)} vibhāṣā śveḥ : prāpte aprāpte ubhayatra vā iti sandehaḥ . \\
\noindent \textbf{(P\_1,1.44.8)} katham ca prāpte katham vā aprāpte katham vā ubhayatra .kiti iti vā nitye prāpte anyatra vā aprāpte ubhayatra vā iti . \\
\noindent \textbf{(P\_1,1.44.8)} ubhayatra . \\
\noindent \textbf{(P\_1,1.44.8)} prāpte tāvat : śuśuvatuḥ , śuśuvuḥ , śiśviyatuḥ , śiśviyuḥ . \\
\noindent \textbf{(P\_1,1.44.8)} aprāpte : śuśāva śuśavitha śiśvāya śiśvayitha . \\
\noindent \textbf{(P\_1,1.44.8)} vibhāṣā saṅghuṣāsvanām : sampūrvāt ghuṣeḥ prāpte aprāpte ubhayatra vā iti sandehaḥ . \\
\noindent \textbf{(P\_1,1.44.8)} katham ca prāpte katham vā aprāpte katham vā ubhayatra . \\
\noindent \textbf{(P\_1,1.44.8)} ghuṣiḥ aviśabdane iti vā nitye prāpte anyatra vā aprāpte ubhayatra vā iti . \\
\noindent \textbf{(P\_1,1.44.8)} ubhayatra . \\
\noindent \textbf{(P\_1,1.44.8)} prāpte tāvat : saṅghuṣṭā rajjuḥ , saṅghuṣitā rajjuḥ . \\
\noindent \textbf{(P\_1,1.44.8)} aprāpte : saṅghuṣṭam vākyam , saṅghuṣitam vākyam . \\
\noindent \textbf{(P\_1,1.44.8)} āṅpūrvāt svaneḥ prāpte aprāpte ubhayatra vā iti sandehaḥ . \\
\noindent \textbf{(P\_1,1.44.8)} katham ca prāpte katham vā aprāpte katham vā ubhayatra . \\
\noindent \textbf{(P\_1,1.44.8)} manasi iti vā nitye prāpte anyatra vā aprāpte ubhayatra vā iti . \\
\noindent \textbf{(P\_1,1.44.8)} ubhayatra . \\
\noindent \textbf{(P\_1,1.44.8)} prāpte tāvat : āsvāntam manaḥ , āsvanitam manaḥ . \\
\noindent \textbf{(P\_1,1.44.8)} aprāpte : āsvāntaḥ devadattaḥ , āsvanitaḥ devadattaḥ iti . \\
\noindent \textbf{(P\_1,1.45)} fkim iyam vākyasya samprasāraṇasañjñā kriyate : ik yaṇaḥ iti etat vākyam samprasāraṇasañjñam bhavati iti , āhosvit varṇasya : ik yaḥ yaṇaḥ sthāne saḥ samprasāraṇasañjñaḥ bhavati iti . \\
\noindent \textbf{(P\_1,1.45)} kaḥ ca atra viśeṣaḥ . \\
\noindent \textbf{(P\_1,1.45)} samprasāraṇasañjñāyām vākyasañjñā cet varṇavidhiḥ . \\
\noindent \textbf{(P\_1,1.45)} samprasāraṇasañjñāyām vākyasañjñā cet varṇavidhiḥ na sidhyati : samprasāraṇāt paraḥ pūrvaḥ bhavati , samprasāraṇasya dīrghaḥ bhavati iti . \\
\noindent \textbf{(P\_1,1.45)} na hi vākyasya samprasāraṇasañjñāyām satyām eṣaḥ nirdeśaḥ upapadyate na api etayoḥ kāryayoḥ sambhavaḥ asti . \\
\noindent \textbf{(P\_1,1.45)} astu tarhi varṇasya . \\
\noindent \textbf{(P\_1,1.45)} varṇasañjñā cet nirvṛttiḥ . varṇasañjñā cet nirvṛttiḥ na sidhyati : ṣyaṅaḥ samprasāraṇam iti . \\
\noindent \textbf{(P\_1,1.45)} saḥ eva hi tāvat ik durlabhaḥ yasya sañjñā kriyate . \\
\noindent \textbf{(P\_1,1.45)} atha api katham cit labhyeta kena asu yaṇaḥ sthāne syāt . \\
\noindent \textbf{(P\_1,1.45)} anena eva hi asau vyavasthāpyate . \\
\noindent \textbf{(P\_1,1.45)} tat etat itaretarāśrayam bhavati , itaretarāśrayāṇi ca kāryāṇi na prakalpante . \\
\noindent \textbf{(P\_1,1.45)} vibhaktiviśeṣanirdeśaḥ tu jñāpakaḥ ubhayasañjñātvasya . \\
\noindent \textbf{(P\_1,1.45)} yat ayam vibhaktiviśeṣaiḥ nirdeśam karoti samprasāraṇāt paraḥ pūrvaḥ bhavati samprasāraṇasya dīrghaḥ bhavati ṣyaṅaḥ samprasāraṇam iti tena jñāyate ubhayoḥ sañjñā bhavati iti . \\
\noindent \textbf{(P\_1,1.45)} yat tāvat āha samprasāraṇāt paraḥ pūrvaḥ bhavati samprasāraṇasya dīrghaḥ bhavati iti tena jñāyate varṇasya bhavati iti . \\
\noindent \textbf{(P\_1,1.45)} yat api āha ṣyaṅaḥ samprasāraṇam iti tena jñāyate vākyasya api sañjñā bhavati iti . \\
\noindent \textbf{(P\_1,1.45)} atha vā punaḥ astu vākyasya eva . \\
\noindent \textbf{(P\_1,1.45)} nanu ca uktam samprasāraṇasañjñāyām vākyasañjñā cet varṇavidhiḥ iti . \\
\noindent \textbf{(P\_1,1.45)} na eṣaḥ doṣaḥ . \\
\noindent \textbf{(P\_1,1.45)} yathā kākāt jātaḥ kākaḥ , śyenāt jātaḥ śyenaḥ evam samprasāraṇāt jātam samprasāraṇam . \\
\noindent \textbf{(P\_1,1.45)} yat tat samprasāraṇāt jātam samprasāraṇam tasmāt paraḥ pūrvaḥ bhavati tasya dīrghaḥ bhavati iti . \\
\noindent \textbf{(P\_1,1.45)} atha vā dṛśyante hi vākyeṣu vākyaikadeśān prayuñjānāḥ padeṣu ca padaikadeśān . \\
\noindent \textbf{(P\_1,1.45)} vākyeṣu tāvat vākyaikadeśān : praviśa piṇḍīm , praviśa tarparṇam . \\
\noindent \textbf{(P\_1,1.45)} padeṣu padaikadeśān : devadattaḥ dattaḥ , satyabhāmā bhāmā iti . \\
\noindent \textbf{(P\_1,1.45)} evam iha api samprasāraṇanirvṛttāt samprasāraṇanirvṛttasya iti etasya vākyasya arthe samprasāraṇāt samprasāraṇasya iti vākyaikadeśaḥ prayujyate . \\
\noindent \textbf{(P\_1,1.45)} tena nirvṛttasya vidhim vijñāsyāmaḥ . \\
\noindent \textbf{(P\_1,1.45)} samprasāraṇanirvṛttāt samprasāraṇanirvṛttasya iti . \\
\noindent \textbf{(P\_1,1.45)} atha vā āha ayam samprasāraṇāt paraḥ pūrvaḥ bhavati samprasāraṇasya dīrghaḥ bhavati iti . \\
\noindent \textbf{(P\_1,1.45)} na ca vākyasya samprasāraṇasañjñāyām satyām eṣaḥ nirdeśaḥ upapadyate na api etayoḥ kāryayoḥ sambhavaḥ asti . \\
\noindent \textbf{(P\_1,1.45)} tatra vacanāt bhaviṣyati . \\
\noindent \textbf{(P\_1,1.45)} atha vā punaḥ astu varṇasya . \\
\noindent \textbf{(P\_1,1.45)} nanu ca uktam varṇasañjñā cet nirvṛttiḥ iti . \\
\noindent \textbf{(P\_1,1.45)} na eṣaḥ doṣaḥ . \\
\noindent \textbf{(P\_1,1.45)} itaretarāśrayamātram etat coditam . \\
\noindent \textbf{(P\_1,1.45)} sarvāṇi ca itaretarāśrayāṇi ekatvena parihṛtāni siddham tu nityaśabdatvāt iti . \\
\noindent \textbf{(P\_1,1.45)} na idam tulyam anyaiḥ itaretarāśrayaiḥ . \\
\noindent \textbf{(P\_1,1.45)} na hi tatra kim cit ucyate asya sthāne ye ākāraikāraukārāḥ bhāvyante te vṛddhisañjñāḥ bhavanti iti . \\
\noindent \textbf{(P\_1,1.45)} iha punaḥ ucyate ik yaḥ yaṇaḥ sthāne saḥ samprasāraṇasañjñaḥ bhavati iti . \\
\noindent \textbf{(P\_1,1.45)} evam tarhi bhāvinī iyam sañjñā vijñāsyate . \\
\noindent \textbf{(P\_1,1.45)} tat yathā : kaḥ cit kam cit tantuvāyam āha : asya sūtrasya śāṭakam vaya iti . \\
\noindent \textbf{(P\_1,1.45)} saḥ paśyati : yadi śāṭakaḥ na vātavyaḥ atha vātavyaḥ na śāṭakaḥ . \\
\noindent \textbf{(P\_1,1.45)} śāṭakaḥ vātavyaḥ iti vipratiṣiddham . \\
\noindent \textbf{(P\_1,1.45)} bhāvinī khalu asya sañjñā abhipretā . \\
\noindent \textbf{(P\_1,1.45)} saḥ manye vātavyaḥ yasmin ute śāṭakaḥ iti etat bhavati iti . \\
\noindent \textbf{(P\_1,1.45)} evam iha api saḥ yaṇaḥ sthāne bhavati yasya abhinirvṛttasya samprasāraṇam iti eṣā sañjñā bhaviṣyati . \\
\noindent \textbf{(P\_1,1.45)} atha vā ijādiyajādipravṛttiḥ ca eva hi loke lakṣyate . \\
\noindent \textbf{(P\_1,1.45)} yajādyupadeśāt tu ijādinivṛttiḥ prasaktā . \\
\noindent \textbf{(P\_1,1.45)} prayuñjate ca punaḥ lokāḥ iṣṭam uptam iti . \\
\noindent \textbf{(P\_1,1.45)} te manyāmahe : asya yaṇaḥ sthāne imam ikam prayuñjate iti . \\
\noindent \textbf{(P\_1,1.45)} tatra tasya asādhvabhimatasya śāstreṇa sādhutvam avasthāpyate : kiti sādhuḥ bhavati ṅiti sādhuḥ bhavati iti . \\
\noindent \textbf{(P\_1,1.46.1)} samāsanirdeśaḥ ayam . \\
\noindent \textbf{(P\_1,1.46.1)} tatra na jñāyate kaḥ ādiḥ kaḥ antaḥ iti . \\
\noindent \textbf{(P\_1,1.46.1)} tat yathā : ajāvidhanau devadattayajñadattau iti ukte na jñāyate kasya ajāḥ dhanam kasya avayaḥ iti . \\
\noindent \textbf{(P\_1,1.46.1)} yadi api tāvat loke eṣaḥ dṛṣṭāntaḥ dṛṣṭāntasya api puruṣārambhaḥ nivartakaḥ bhavati . \\
\noindent \textbf{(P\_1,1.46.1)} asti ca iha kaḥ cit puruṣārambhaḥ . \\
\noindent \textbf{(P\_1,1.46.1)} asti iti āha . \\
\noindent \textbf{(P\_1,1.46.1)} kaḥ . \\
\noindent \textbf{(P\_1,1.46.1)} saṅkhyātanudeśaḥ nāma . \\
\noindent \textbf{(P\_1,1.46.2)} kau punaḥ ṭakitau ādyantau bhavataḥ . \\
\noindent \textbf{(P\_1,1.46.2)} āgamau iti āha . \\
\noindent \textbf{(P\_1,1.46.2)} yuktam punaḥ yat nityeṣu nāma śabdeṣu āgamaśāsanam syāt na nityeṣu śabdeṣu kūṭasthaiḥ avicālibhiḥ varṇaiḥ bhavitavyam anapāyopajanavikāribhiḥ . \\
\noindent \textbf{(P\_1,1.46.2)} āgamaḥ ca nāma apūrvaḥ śabdopajanaḥ . \\
\noindent \textbf{(P\_1,1.46.2)} atha yuktam yat nityeṣu śabdeṣu ādeśāḥ syuḥ . \\
\noindent \textbf{(P\_1,1.46.2)} bāḍham yuktam . \\
\noindent \textbf{(P\_1,1.46.2)} śabdāntaraiḥ iha bhavitavyam . \\
\noindent \textbf{(P\_1,1.46.2)} tatra śabdāntarāt śabdāntarasya pratipattiḥ yuktā . \\
\noindent \textbf{(P\_1,1.46.2)} ādeśāḥ tarhi ime bhaviṣyanti anāgamakānām sāgamakāḥ . \\
\noindent \textbf{(P\_1,1.46.2)} tat katham . \\
\noindent \textbf{(P\_1,1.46.2)} sañjñādhikāraḥ ayam . \\
\noindent \textbf{(P\_1,1.46.2)} ādyantau ca iha saṅkīrtyete . \\
\noindent \textbf{(P\_1,1.46.2)} ṭakārkakārau itau udāhriyete . \\
\noindent \textbf{(P\_1,1.46.2)} tatra ādyantayoḥ ṭakārakakārau itau sañjñe bhaviṣyataḥ . \\
\noindent \textbf{(P\_1,1.46.2)} tatra ārdhadhātukasya iṭ valādeḥ iti upasthitam idam bhavati : ādiḥ iti . \\
\noindent \textbf{(P\_1,1.46.2)} tena ikārādiḥ ādeśaḥ bhaviṣyati . \\
\noindent \textbf{(P\_1,1.46.2)} etāvat iha sūtram iṭ iti . \\
\noindent \textbf{(P\_1,1.46.2)} katham punaḥ iyatā sūtreṇa ikārādiḥ ādeśaḥ labhyaḥ . \\
\noindent \textbf{(P\_1,1.46.2)} labhyaḥ iti āha . \\
\noindent \textbf{(P\_1,1.46.2)} katham. bahuvrīhinirdeśāt . \\
\noindent \textbf{(P\_1,1.46.2)} bahuvrīhinirdeśaḥ ayam : ikāraḥ ādiḥ asya iti . \\
\noindent \textbf{(P\_1,1.46.2)} yadi api tāvat atra etat śakyate vakutm iha katham : luṅlaṅlṛṅkṣu aṭ udāttaḥ iti yatra aśakyam udāttagrahaṇena akāraḥ viśeṣayitum . \\
\noindent \textbf{(P\_1,1.46.2)} tatra kaḥ doṣaḥ . \\
\noindent \textbf{(P\_1,1.46.2)} aṅgasya udāttatvam prasajyeta . \\
\noindent \textbf{(P\_1,1.46.2)} na eṣaḥ doṣaḥ . \\
\noindent \textbf{(P\_1,1.46.2)} tripadaḥ ayam bahuvrīhiḥ . \\
\noindent \textbf{(P\_1,1.46.2)} tatra vākye eva udāttagrahaṇena akāraḥ viśeṣyate : akāraḥ udāttaḥ ādiḥ asya iti . \\
\noindent \textbf{(P\_1,1.46.2)} yatra tarhi anuvṛttyā etat bhavati : āṭ ajādīnām iti . \\
\noindent \textbf{(P\_1,1.46.2)} vakṣyati etat : ajādīnām aṭā siddham iti . \\
\noindent \textbf{(P\_1,1.46.2)} atha vā yat tāvat ayam sāmānyena śaknoti upadeṣṭum tat tāvat upadiśati prakṛtim tataḥ valādi ārdhadhātukam tataḥ paścāt ikāram . \\
\noindent \textbf{(P\_1,1.46.2)} tena ayam viśeṣeṇa śabdāntaram samudāyam pratipadyate . \\
\noindent \textbf{(P\_1,1.46.2)} tat yathā khadiraburburayoḥ : khadiraburburau gaurakāṇḍau sūkṣmaparṇau . \\
\noindent \textbf{(P\_1,1.46.2)} tataḥ paścāt āha kaṇṭakavān khadiraḥ iti . \\
\noindent \textbf{(P\_1,1.46.2)} tena asau viśeṣeṇa dravyāntaram samudāyam pratipadyate . \\
\noindent \textbf{(P\_1,1.46.2)} atha vā etayā ānupūrvyā ayam śabdāntaram upadiśati : prakṛtim tataḥ valādi ārdhadhātukam tataḥ paścāt ikāram yasmin tasya āgamabuddhiḥ bhavati . \\
\noindent \textbf{(P\_1,1.46.3)} ṭakitoḥ ādyantavidhāne pratyayapratiṣedhaḥ . \\
\noindent \textbf{(P\_1,1.46.3)} ṭakitoḥ ādyantavidhāne pratyayasya pratiṣedhaḥ vaktavyaḥ . \\
\noindent \textbf{(P\_1,1.46.3)} pratyayaḥ ādiḥ antaḥ vā mā bhūt : careḥ ṭaḥ ātaḥ anupasarge kaḥ iti . \\
\noindent \textbf{(P\_1,1.46.3)} paravacanāt siddham . \\
\noindent \textbf{(P\_1,1.46.3)} paravacanāt pratyayaḥ ādiḥ antaḥ vā na bhaviṣyati . \\
\noindent \textbf{(P\_1,1.46.3)} paravacanāt siddham iti cet na apavādatvāt . \\
\noindent \textbf{(P\_1,1.46.3)} paravacanāt siddham iti cet na . \\
\noindent \textbf{(P\_1,1.46.3)} kim kāraṇam . \\
\noindent \textbf{(P\_1,1.46.3)} apavādatvāt . \\
\noindent \textbf{(P\_1,1.46.3)} apavādaḥ ayam yogaḥ . \\
\noindent \textbf{(P\_1,1.46.3)} tat yathā mit acaḥ antyāt paraḥ iti eṣaḥ yogaḥ sthāneyogatvasya pratyayaparatvasya ca apavādaḥ . \\
\noindent \textbf{(P\_1,1.46.3)} viṣamaḥ upanyāsaḥ . \\
\noindent \textbf{(P\_1,1.46.3)} yuktam tatra yat anavakāśam mitkaraṇam sthāneyogatvam pratyayaparatvam ca bādhate . \\
\noindent \textbf{(P\_1,1.46.3)} iha punaḥ ubhayam sāvakāsam . \\
\noindent \textbf{(P\_1,1.46.3)} kaḥ avakāśaḥ . \\
\noindent \textbf{(P\_1,1.46.3)} ṭitkaraṇasya avakāśaḥ : ṭitaḥ iti īkāraḥ yathā syāt . \\
\noindent \textbf{(P\_1,1.46.3)} kitkaraṇasya avakāśaḥ : kiti iti ākāralopaḥ yathā syāt . \\
\noindent \textbf{(P\_1,1.46.3)} prayojanam nāma tat vaktavyam yat niyogataḥ syāt . \\
\noindent \textbf{(P\_1,1.46.3)} yadi ca ayam niyogataḥ paraḥ syāt tataḥ etat prayojanam syāt . \\
\noindent \textbf{(P\_1,1.46.3)} kutaḥ nu khalu etat ṭitkaraṇāt ayam paraḥ bhaviṣyati na punaḥ ādiḥ iti kitkaraṇāt ca paraḥ bhaviṣyati na punaḥ antaḥ iti . \\
\noindent \textbf{(P\_1,1.46.3)} ṭitaḥ khalu api eṣaḥ parihāraḥ yatra na asti sambhavaḥ yat paraḥ ca syāt ādiḥ ca . \\
\noindent \textbf{(P\_1,1.46.3)} kitaḥ tu aparihāraḥ . \\
\noindent \textbf{(P\_1,1.46.3)} asti hi sambhavaḥ yat paraḥ ca syāt antaḥ ca . \\
\noindent \textbf{(P\_1,1.46.3)} tatra kaḥ doṣaḥ . \\
\noindent \textbf{(P\_1,1.46.3)} upasarge ghoḥ kiḥ : ādhyoḥ , pradhyoḥ . \\
\noindent \textbf{(P\_1,1.46.3)} noṅdhātvoḥ iti pratiṣedhaḥ prasajyeta . \\
\noindent \textbf{(P\_1,1.46.3)} ṭitaḥ ca api aparihāraḥ . \\
\noindent \textbf{(P\_1,1.46.3)} syāt eva hi ayam ṭitkaraṇāt ādiḥ na punaḥ paraḥ . \\
\noindent \textbf{(P\_1,1.46.3)} kva tarhi idānīm idam syāt : ṭitaḥ īkāraḥ bhavati iti . \\
\noindent \textbf{(P\_1,1.46.3)} yaḥ ubhayavān : gāpoḥ ṭak iti . \\
\noindent \textbf{(P\_1,1.46.3)} siddham tu ṣaṣṭhyadhikāre vacanāt . \\
\noindent \textbf{(P\_1,1.46.3)} siddham etat . \\
\noindent \textbf{(P\_1,1.46.3)} katham . \\
\noindent \textbf{(P\_1,1.46.3)} ṣaṣṭhyadhikāre ayam yogaḥ karatvyaḥ : ādyantau ṭakitau ṣaṣṭhīnirdiṣṭasya iti . \\
\noindent \textbf{(P\_1,1.46.3)} ādyantayoḥ vā ṣaṣthyarthatvāt tadabhāve asampratyayaḥ . \\
\noindent \textbf{(P\_1,1.46.3)} ādyantayoḥ vā ṣaṣthyarthatvāt ṣaṣṭhyāḥ abhāve asampratyayaḥ . \\
\noindent \textbf{(P\_1,1.46.3)} ādiḥ antaḥ vā na bhaviṣyati . \\
\noindent \textbf{(P\_1,1.46.3)} yuktam punaḥ yat śabdanimittakaḥ nāma arthaḥ syāt na arthanimittakena śabdena bhavitavyam . \\
\noindent \textbf{(P\_1,1.46.3)} arthanimittakaḥ eva śabdaḥ . \\
\noindent \textbf{(P\_1,1.46.3)} tat katham . \\
\noindent \textbf{(P\_1,1.46.3)} ādyantau ṣaṣṭhyarthau . \\
\noindent \textbf{(P\_1,1.46.3)} na ca atra ṣaṣṭhīm paśyāmaḥ . \\
\noindent \textbf{(P\_1,1.46.3)} te manyāmahe : ādyantau eva atra na staḥ . \\
\noindent \textbf{(P\_1,1.46.3)} tayoḥ abhāve ṣaṣṭhī api na bhavati iti . \\
\noindent \textbf{(P\_1,1.47.1)} kimartham idam ucyate . \\
\noindent \textbf{(P\_1,1.47.1)} mit acaḥ antyāt paraḥ iti sthānaparapratyayāpavādaḥ . \\
\noindent \textbf{(P\_1,1.47.1)} mit acaḥ antyāt paraḥ iti ucyate sthāneyogatvasya pratyayaparatvasya ca apavādaḥ . \\
\noindent \textbf{(P\_1,1.47.1)} sthāneyogatvasya tāvat : kuṇḍāni vanāni payāṃsi yaśāṃsi . \\
\noindent \textbf{(P\_1,1.47.1)} pratyayaparatvasya : bhinatti chinatti . \\
\noindent \textbf{(P\_1,1.47.1)} bhavet idam yuktam udāharaṇam kuṇḍāni vanāni yatra na asti sambhavaḥ yat ayam acaḥ anytāt paraḥ ca syāt sthāne ca iti . \\
\noindent \textbf{(P\_1,1.47.1)} idam tu ayuktam payāṃsi yaśāṃsi . \\
\noindent \textbf{(P\_1,1.47.1)} asti hi sambhavaḥ yat acaḥ anytāt paraḥ ca syāt sthāne ca . \\
\noindent \textbf{(P\_1,1.47.1)} etat api yuktam . \\
\noindent \textbf{(P\_1,1.47.1)} katham . \\
\noindent \textbf{(P\_1,1.47.1)} na eva īśvaraḥ ājñāpayati na api dharmasūtrakārāḥ paṭhanti apavādaiḥ utsargāḥ bādhyantām iti . \\
\noindent \textbf{(P\_1,1.47.1)} kim tarhi . \\
\noindent \textbf{(P\_1,1.47.1)} laukikaḥ ayam dṛṣṭāntaḥ . \\
\noindent \textbf{(P\_1,1.47.1)} loke hi sati api sambhave bādhanam bhavati . \\
\noindent \textbf{(P\_1,1.47.1)} tat yathā : dadhi brāhmaṇebhyaḥ dīyatām takram kauṇḍinyāya iti sati api sambhave dadhidānasya takradānam nivartakam bhavati . \\
\noindent \textbf{(P\_1,1.47.1)} evam iha api sati api sambhave acām antyāt paratvam ṣaṣṭhīsthāneyogatvam bādhiṣyate . \\
\noindent \textbf{(P\_1,1.47.2)} antyāt pūrvaḥ masjeḥ anuṣaṅgasaṃyogādilopārtham . \\
\noindent \textbf{(P\_1,1.47.2)} antyāt pūrvaḥ masjeḥ mit vaktavyaḥ . \\
\noindent \textbf{(P\_1,1.47.2)} kim prayojanam . \\
\noindent \textbf{(P\_1,1.47.2)} anuṣaṅgasaṃyogādilopārtham . \\
\noindent \textbf{(P\_1,1.47.2)} anuṣaṅgalopārtham saṃyogādilopārtham ca . \\
\noindent \textbf{(P\_1,1.47.2)} anuṣaṅgalopārtham tāvat : magnaḥ , magnavān . \\
\noindent \textbf{(P\_1,1.47.2)} saṃyogādilopārtham maṅktā maṅktum , maṅktavyam . \\
\noindent \textbf{(P\_1,1.47.2)} bharjimarcyoḥ ca . \\
\noindent \textbf{(P\_1,1.47.2)} bharjimarcyoḥ ca antyāt pūrvaḥ mit vaktavyaḥ . \\
\noindent \textbf{(P\_1,1.47.2)} bharūjā marīcayaḥ iti . \\
\noindent \textbf{(P\_1,1.47.2)} saḥ tarhi vaktavyaḥ . \\
\noindent \textbf{(P\_1,1.47.2)} na vaktavyaḥ . \\
\noindent \textbf{(P\_1,1.47.2)} nipātanāt siddham . \\
\noindent \textbf{(P\_1,1.47.2)} kim nipātanam . \\
\noindent \textbf{(P\_1,1.47.2)} bharūjāśabdaḥ aṅgulyādiṣu paṭhyate marīciśabdaḥ bāhvādiṣu . \\
\noindent \textbf{(P\_1,1.47.3)} kim punaḥ ayam pūrvāntaḥ āhosvit parādiḥ āhosvit abhaktaḥ . \\
\noindent \textbf{(P\_1,1.47.3)} katham ca ayam pūrvāntaḥ syāt katham vā parādiḥ katham vā abhaktaḥ . \\
\noindent \textbf{(P\_1,1.47.3)} yadi antaḥ iti vartate tataḥ pūrvāntaḥ . \\
\noindent \textbf{(P\_1,1.47.3)} atha ādiḥ iti vartate tataḥ parādiḥ . \\
\noindent \textbf{(P\_1,1.47.3)} atha ubhayam nivṛttam tataḥ abhaktaḥ . \\
\noindent \textbf{(P\_1,1.47.3)} kaḥ ca atra viśeṣaḥ . \\
\noindent \textbf{(P\_1,1.47.3)} abhakte dīrghanalopasvaraṇatvānusvāraśībhāvāḥ . \\
\noindent \textbf{(P\_1,1.47.3)} yadi abhaktaḥ dīrghatvam na prāpnoti : kuṇḍāni vanāni . \\
\noindent \textbf{(P\_1,1.47.3)} nopadhāyāḥ sarvanāmasthāne ca asambuddhau iti dīrghatvam na prāpnoti . \\
\noindent \textbf{(P\_1,1.47.3)} dīrgha . \\
\noindent \textbf{(P\_1,1.47.3)} nalopa : nalopaḥ ca na sidhyati : agne trī te vajinā trī sadhasthā , ta tā piṇḍānām . \\
\noindent \textbf{(P\_1,1.47.3)} nalopaḥ prātipadikāntasya iti nalopaḥ na prāpnoti . \\
\noindent \textbf{(P\_1,1.47.3)} nalopa . \\
\noindent \textbf{(P\_1,1.47.3)} svara : svaraḥ ca na sidhyati : sarvāṇi jyotīṃṣi . \\
\noindent \textbf{(P\_1,1.47.3)} sarvasya supi iti ādyudāttatvam na prāpnoti . \\
\noindent \textbf{(P\_1,1.47.3)} svara . \\
\noindent \textbf{(P\_1,1.47.3)} ṇatva : ṇatvam ca na sidhyati : māṣavāpāṇi vrīhivāpāṇi . \\
\noindent \textbf{(P\_1,1.47.3)} pūrvānte prātipadikāntanakārasya iti siddham , parādau vibhaktinakārasya , abhakte numaḥ grahaṇam kartavyam . \\
\noindent \textbf{(P\_1,1.47.3)} na kartavyam . \\
\noindent \textbf{(P\_1,1.47.3)} kriyate nyāse eva : prātipadikāntanumvibhaktiṣu iti . \\
\noindent \textbf{(P\_1,1.47.3)} ṇatva . \\
\noindent \textbf{(P\_1,1.47.3)} anusvāra : anusvāraḥ ca na sidhyati : dviṣantapaḥ , parantapaḥ . \\
\noindent \textbf{(P\_1,1.47.3)} maḥ anusvāraḥ hali iti anusvāraḥ na prāpnoti . \\
\noindent \textbf{(P\_1,1.47.3)} mā bhūt evam . \\
\noindent \textbf{(P\_1,1.47.3)} naḥ ca apadāntasya jhali iti evam bhaviṣyati . \\
\noindent \textbf{(P\_1,1.47.3)} yaḥ tarhi na jhalparaḥ : vahaṃlihaḥ gauḥ , abhraṃlihaḥ vāyuḥ . \\
\noindent \textbf{(P\_1,1.47.3)} anusvāra . \\
\noindent \textbf{(P\_1,1.47.3)} śībhāva : śībhāvaḥ ca na sidhyati : trapuṇī jatunī tumburuṇī . \\
\noindent \textbf{(P\_1,1.47.3)} napuṃsakāt uttarasya auṅaḥ śībhāvaḥ bhavati iti śībhāvaḥ na prāpnoti . \\
\noindent \textbf{(P\_1,1.47.3)} śībhāva . \\
\noindent \textbf{(P\_1,1.47.3)} evam tarhi parādiḥ kariṣyate . \\
\noindent \textbf{(P\_1,1.47.3)} parādau guṇavṛddhyauttvadīrghanalopānusvāraśībhāvenakārapratiṣedhaḥ . \\
\noindent \textbf{(P\_1,1.47.3)} yadi parādiḥ guṇaḥ pratiṣedhyaḥ : trapuṇe jatune tumburuṇe . \\
\noindent \textbf{(P\_1,1.47.3)} gheḥ ṅiti iti guṇaḥ prāpnoti . \\
\noindent \textbf{(P\_1,1.47.3)} guṇa . \\
\noindent \textbf{(P\_1,1.47.3)} vṛddhi : vṛddhiḥ pratiṣedhyā : atisakhīni brāhmaṇakulāni . \\
\noindent \textbf{(P\_1,1.47.3)} sakhyuḥ asambuddhau iti ṇittve acaḥ ñṇiti iti vṛddhiḥ prāpnoti . \\
\noindent \textbf{(P\_1,1.47.3)} vṛddhi . \\
\noindent \textbf{(P\_1,1.47.3)} auttva : auttvam ca pratiṣedhyam : trapuṇi jatuni tumburuṇi . \\
\noindent \textbf{(P\_1,1.47.3)} idudbhyām aut at ca gheḥ iti auttvam prāpnoti . \\
\noindent \textbf{(P\_1,1.47.3)} auttva . \\
\noindent \textbf{(P\_1,1.47.3)} dīrgha : dīrghatvam ca na sidhyati : kuṇḍāni vanāni . \\
\noindent \textbf{(P\_1,1.47.3)} nopadhāyāḥ sarvanāmasthāne ca asambuddhau iti dīrghatvam na prāpnoti . \\
\noindent \textbf{(P\_1,1.47.3)} mā bhūt evam . \\
\noindent \textbf{(P\_1,1.47.3)} ataḥ dīrghaḥ yañi supi ca iti evam bhaviṣyati . \\
\noindent \textbf{(P\_1,1.47.3)} iha tarhi : asthīni dadhīni priyasakhīni brāhmaṇakulāni . \\
\noindent \textbf{(P\_1,1.47.3)} dīrgha . \\
\noindent \textbf{(P\_1,1.47.3)} nalopa : nalopaḥ ca na sidhyati : agne trī te vajinā trī sadhasthā , ta tā piṇḍānām . \\
\noindent \textbf{(P\_1,1.47.3)} nalopaḥ prātipadikāntasya iti nalopaḥ na prāpnoti . \\
\noindent \textbf{(P\_1,1.47.3)} nalopa . \\
\noindent \textbf{(P\_1,1.47.3)} anusvāra : anusvāraḥ ca na sidhyati : dviṣantapaḥ , parantapaḥ . \\
\noindent \textbf{(P\_1,1.47.3)} maḥ anusvāraḥ hali iti anusvāraḥ na prāpnoti . \\
\noindent \textbf{(P\_1,1.47.3)} mā bhūt evam . \\
\noindent \textbf{(P\_1,1.47.3)} naḥ ca apadāntasya jhali iti evam bhaviṣyati . \\
\noindent \textbf{(P\_1,1.47.3)} yaḥ tarhi na jhalparaḥ : vahaṃlihaḥ gauḥ , abhraṃlihaḥ vāyuḥ . \\
\noindent \textbf{(P\_1,1.47.3)} anusvāra . \\
\noindent \textbf{(P\_1,1.47.3)} śībhāvenakārapratiṣedhaḥ : śībhāve nakārasya pratiṣedhaḥ vaktavyaḥ : trapuṇī jatunī tumburuṇī . \\
\noindent \textbf{(P\_1,1.47.3)} sanumkasya śībhāvaḥ prāpnoti . \\
\noindent \textbf{(P\_1,1.47.3)} na eṣaḥ doṣaḥ . \\
\noindent \textbf{(P\_1,1.47.3)} nirdiśyamānasya ādeśāḥ bhavanti iti evam na bhaviṣyati . \\
\noindent \textbf{(P\_1,1.47.3)} yaḥ tarhi nirdiśyate tasya na prāpnoti . \\
\noindent \textbf{(P\_1,1.47.3)} kasmāt . \\
\noindent \textbf{(P\_1,1.47.3)} numā vyavahitatvāt . \\
\noindent \textbf{(P\_1,1.47.3)} evam tarhi pūrvāntaḥ kariṣyate . \\
\noindent \textbf{(P\_1,1.47.3)} pūrvānte napuṃsakopasarjanahrasvatvam dvigusvaraḥ ca . \\
\noindent \textbf{(P\_1,1.47.3)} yadi pūrvantaḥ kriyate napuṃsakopasarjanahrasvatvam dvigusvaraḥ ca na sidhyati . \\
\noindent \textbf{(P\_1,1.47.3)} napuṃsakopasarjanahrasvatvam : ārāśastriṇī dhānāśaṣkulinī niṣkauśāmbinī nirvārāṇasinī . \\
\noindent \textbf{(P\_1,1.47.3)} dvigusvara : pañcāratninī daśāratninī . \\
\noindent \textbf{(P\_1,1.47.3)} numi kṛte anantyatvāt ete vidhayaḥ na prāpnunvanti . \\
\noindent \textbf{(P\_1,1.47.3)} na vā bahiraṅgalakṣaṇatvāt . \\
\noindent \textbf{(P\_1,1.47.3)} na vā eṣaḥ doṣaḥ . \\
\noindent \textbf{(P\_1,1.47.3)} kim kāraṇam . \\
\noindent \textbf{(P\_1,1.47.3)} bahiraṅgalakṣaṇatvāt . \\
\noindent \textbf{(P\_1,1.47.3)} bahiraṅgaḥ num , antaraṅgāḥ ete vidhayaḥ . \\
\noindent \textbf{(P\_1,1.47.3)} asiddham bahiraṅgam antaraṅge . \\
\noindent \textbf{(P\_1,1.47.3)} dvigusvare bhūyān parihāraḥ : saṅghātabhaktaḥ asau na utsahate avayavasya igantatām vihantum iti kṛtvā dvigusvaraḥ bhaviṣyati . \\
\noindent \textbf{(P\_1,1.48)} kimartham idam ucyate . \\
\noindent \textbf{(P\_1,1.48)} ecaḥ ik savarṇākāranivṛttyartham . \\
\noindent \textbf{(P\_1,1.48)} ecaḥ ik bhavati iti ucyate savarṇanivṛttyartham akāranivṛttyartham ca . \\
\noindent \textbf{(P\_1,1.48)} savarṇnivṛttyartham tāvat : eṅaḥ hrasvaśāsaneṣu ardhaḥ ekāraḥ ardhaḥ okāraḥ vā mā bhūt iti . \\
\noindent \textbf{(P\_1,1.48)} akāranivṛttyartham ca . \\
\noindent \textbf{(P\_1,1.48)} imau aicau samāhāravarṇau . \\
\noindent \textbf{(P\_1,1.48)} mātrā avarṇasya mātrā ivarṇovarṇayoḥ . \\
\noindent \textbf{(P\_1,1.48)} tayoḥ hrasvaśāsaneṣu kadā cit avarṇaḥ syāt kadā cit ivarṇovarṇau . \\
\noindent \textbf{(P\_1,1.48)} mā kadā cit avarṇam bhūt iti evamartham idam ucyate . \\
\noindent \textbf{(P\_1,1.48)} asti prayojanam etat . \\
\noindent \textbf{(P\_1,1.48)} kim tarhi iti . \\
\noindent \textbf{(P\_1,1.48)} dīrghaprasaṅgaḥ . \\
\noindent \textbf{(P\_1,1.48)} dīrghāḥ tu ikaḥ prāpnuvanti . \\
\noindent \textbf{(P\_1,1.48)} kim kāraṇam . \\
\noindent \textbf{(P\_1,1.48)} sthāne antaratamaḥ bhavati iti . \\
\noindent \textbf{(P\_1,1.48)} nanu ca hrasvādeśe iti ucyate . \\
\noindent \textbf{(P\_1,1.48)} tena dīrghāḥ na bhaviṣyanti . \\
\noindent \textbf{(P\_1,1.48)} viṣyārtham etat syāt . \\
\noindent \textbf{(P\_1,1.48)} ecaḥ hrasvaprasaṅge ik bhavati iti . \\
\noindent \textbf{(P\_1,1.48)} dīrghāprasaṅgaḥ tu nivartakatvāt . \\
\noindent \textbf{(P\_1,1.48)} dīrghāṇām tu ikām aprasaṅgaḥ . \\
\noindent \textbf{(P\_1,1.48)} kim kāraṇam . \\
\noindent \textbf{(P\_1,1.48)} nivartakatvāt . \\
\noindent \textbf{(P\_1,1.48)} na anena ikaḥ nirvartyante . \\
\noindent \textbf{(P\_1,1.48)} kim tarhi . \\
\noindent \textbf{(P\_1,1.48)} anikaḥ nivartyante . \\
\noindent \textbf{(P\_1,1.48)} siddhāḥ eva hrasvāḥ ikaḥ ca anikaḥ ca . \\
\noindent \textbf{(P\_1,1.48)} tatra anena anikaḥ nivartyante . \\
\noindent \textbf{(P\_1,1.48)} savarṇanivṛttyarthena tāvat na arthaḥ . \\
\noindent \textbf{(P\_1,1.48)} siddham eṅaḥ sasthānatvāt . \\
\noindent \textbf{(P\_1,1.48)} siddham etat . \\
\noindent \textbf{(P\_1,1.48)} katham . \\
\noindent \textbf{(P\_1,1.48)} eṅaḥ sasthānatvāt ikārokārau bhaviṣyataḥ . \\
\noindent \textbf{(P\_1,1.48)} ardhaḥ ekāraḥ aradhaḥ okāraḥ vā na bhaviṣyati . \\
\noindent \textbf{(P\_1,1.48)} nanu ca eṅaḥ sasthānatarau ardhaḥ ekāraukārau . \\
\noindent \textbf{(P\_1,1.48)} na tau staḥ . \\
\noindent \textbf{(P\_1,1.48)} yadi hi tau syātām tau eva ayam upadiśet . \\
\noindent \textbf{(P\_1,1.48)} nanu ca bhoḥ chandogānām sātyamugrirāṇāyanīyāḥ ardham ekāram ardham okāram ca adhīyate : sujāte eśvasūnṛte , adhvaryo odribhiḥ sutam , śukram te enyat yajatam te enyat iti . \\
\noindent \textbf{(P\_1,1.48)} pārṣadakṛtiḥ eṣā tatrabhavatām . \\
\noindent \textbf{(P\_1,1.48)} na eva loke na anyasmin vede ardhaḥ ekāraḥ ardhaḥ okāraḥ vā asti . \\
\noindent \textbf{(P\_1,1.48)} akāranivṛttyarthena api na arthaḥ . \\
\noindent \textbf{(P\_1,1.48)} aicoḥ ca uttarabhūyastvāt . \\
\noindent \textbf{(P\_1,1.48)} aicoḥ ca uttarabhūyastvāt avarṇaḥ na bhaviṣyati . \\
\noindent \textbf{(P\_1,1.48)} bhūyasī mātrā ivarṇovarṇayoḥ alpīyasī avarṇasya . \\
\noindent \textbf{(P\_1,1.48)} bhūyasaḥ eva grahaṇāni bhaviṣyanti . \\
\noindent \textbf{(P\_1,1.48)} tat yathā brāhmaṇagrāmaḥ ānīyatām iti ucyate tatra ca avarataḥ pañcakārukī bhavati . \\
\noindent \textbf{(P\_1,1.49.1)} kim idam sthāneyogā iti . \\
\noindent \textbf{(P\_1,1.49.1)} sthāne yogaḥ asyāḥ sā iyam sthāneyogā . \\
\noindent \textbf{(P\_1,1.49.1)} saptamyalopaḥ nipātanāt . \\
\noindent \textbf{(P\_1,1.49.1)} tṛtiyāyā vā etvam : sthānena yogaḥ asyāḥ sā iyam sthāneyogā . \\
\noindent \textbf{(P\_1,1.49.2)} kimartham punaḥ idam ucyate . \\
\noindent \textbf{(P\_1,1.49.2)} ṣaṣṭhyāḥ sthāneyogavacanam niyamārtham . \\
\noindent \textbf{(P\_1,1.49.2)} niyamāṛthaḥ ayam ārambhaḥ . \\
\noindent \textbf{(P\_1,1.49.2)} ekaśatam ṣaṣṭhyarthāḥ yāvantaḥ vā te sarve ṣaṣṭhyām uccāritāyām prāpnuvanti . \\
\noindent \textbf{(P\_1,1.49.2)} iṣyate ca vyākaraṇe yā ṣaṣṭhī sā sthāneyogā eva syāt iti . \\
\noindent \textbf{(P\_1,1.49.2)} tat ca antareṇa yatnam na sidhyati iti ṣaṣṭhyāḥ sthāneyogavacanam niyamārtham . \\
\noindent \textbf{(P\_1,1.49.2)} evamartham idam ucyate . \\
\noindent \textbf{(P\_1,1.49.2)} asti prayojanam etat . \\
\noindent \textbf{(P\_1,1.49.2)} kim tarhi iti . \\
\noindent \textbf{(P\_1,1.49.2)} avayavaṣaṣṭhyādiṣu atiprasaṅgaḥ śāsaḥ gohaḥ iti . \\
\noindent \textbf{(P\_1,1.49.2)} avayavaṣaṣṭhyādayaḥ tu na sidhyanti . \\
\noindent \textbf{(P\_1,1.49.2)} tatra kaḥ doṣaḥ . \\
\noindent \textbf{(P\_1,1.49.2)} śāsaḥ it aṅhaloḥ iti śāseḥ ca antyasya syāt upadhāmātrasya ca . \\
\noindent \textbf{(P\_1,1.49.2)} ūt upadhāyāḥ gohaḥ iti gohaḥ ca antyasya syāt upadhāmātrasya ca . \\
\noindent \textbf{(P\_1,1.49.2)} avayavaṣaṣṭhyādīnām ca aprāptiḥ yogasya asandigdhatvāt . \\
\noindent \textbf{(P\_1,1.49.2)} avayavaṣaṣṭhyādīnām ca niyamasya aprāptiḥ . \\
\noindent \textbf{(P\_1,1.49.2)} kim kāraṇam . \\
\noindent \textbf{(P\_1,1.49.2)} yogasya asandigdhatvāt . \\
\noindent \textbf{(P\_1,1.49.2)} sandehe niyamaḥ na ca avayavaṣaṣṭhyādiṣu sandehaḥ . \\
\noindent \textbf{(P\_1,1.49.2)} kim vaktavyam etat . \\
\noindent \textbf{(P\_1,1.49.2)} na hi . \\
\noindent \textbf{(P\_1,1.49.2)} katham anucyamānam gaṃsyate. laukikaḥ ayam dṛṣṭāṇtaḥ . \\
\noindent \textbf{(P\_1,1.49.2)} tat yathā loke : kam cit kaḥ cit pṛcchati : grāmāntaram gamiṣyāmi panthānam me bhavān upadiśatu iti . \\
\noindent \textbf{(P\_1,1.49.2)} saḥ tasmai ācaṣṭe . \\
\noindent \textbf{(P\_1,1.49.2)} amuṣmin avakāśe hastadakṣiṇaḥ grahītavyaḥ amuṣmin avakāśe hastavāmaḥ iti . \\
\noindent \textbf{(P\_1,1.49.2)} yaḥ tu atra tiryakpathaḥ bhavati na tasmin sandehaḥ iti kṛtvā na asau upadiśyate . \\
\noindent \textbf{(P\_1,1.49.2)} evam iha api sandehe niyamaḥ na ca avayavaṣaṣṭhyādiṣu sandehaḥ . \\
\noindent \textbf{(P\_1,1.49.2)} atha vā sthāne ayogā sthāneyogā kim idam ayogā iti . \\
\noindent \textbf{(P\_1,1.49.2)} avyaktayogā ayogā . \\
\noindent \textbf{(P\_1,1.49.2)} atha vā yogavatī yogā . \\
\noindent \textbf{(P\_1,1.49.2)} kā punaḥ yogavatī . \\
\noindent \textbf{(P\_1,1.49.2)} yasyāḥ bahavaḥ yogāḥ . \\
\noindent \textbf{(P\_1,1.49.2)} kutaḥ etat . \\
\noindent \textbf{(P\_1,1.49.2)} bhūmni hi matup bhavati . \\
\noindent \textbf{(P\_1,1.49.2)} viśiṣṭā vā ṣaṣṭhī sthāneyogā . \\
\noindent \textbf{(P\_1,1.49.2)} atha vā kim cid liṅgam āsajya vakṣyāmi : itthaṃliṅgā ṣaṣṭhī sthāneyogā bhavati iti . \\
\noindent \textbf{(P\_1,1.49.2)} na tat liṅgam avayavaṣaṣṭhyādiṣu kariṣyate . \\
\noindent \textbf{(P\_1,1.49.2)} yadi evam śāsaḥ it aṅhaloḥ śā hau śāsigrahaṇam kartavyam sthāneyogārtham liṅgam āsaṅkṣyāmi iti . \\
\noindent \textbf{(P\_1,1.49.2)} na kartavyam . \\
\noindent \textbf{(P\_1,1.49.2)} yat eva adaḥ purastāt avayavaṣaṣṭhyartham prakṛtam etat uttaratra anuvṛttam sat sthāneyogārtham bhaviṣyati . \\
\noindent \textbf{(P\_1,1.49.2)} katham . \\
\noindent \textbf{(P\_1,1.49.2)} adhikāraḥ nāma triprakāraḥ . \\
\noindent \textbf{(P\_1,1.49.2)} kaḥ cit ekadeśasthaḥ sarvam śāstram abhijvalayati yathā pradīpaḥ supravijvalitaḥ sarvam veśma abhijvalayati . \\
\noindent \textbf{(P\_1,1.49.2)} aparaḥ adhikāraḥ yathā rajjvā ayasā vā baddham kāṣṭham anukṛṣyate tadvat anukṛṣyate cakāreṇa . \\
\noindent \textbf{(P\_1,1.49.2)} aparaḥ adhikāraḥ pratiyogam tasya anirdeśārthaḥ iti yoge yoge upatiṣṭhate . \\
\noindent \textbf{(P\_1,1.49.2)} tat yadā eṣaḥ pakṣaḥ adhikāraḥ pratiyogam tasya anirdeśārthaḥ iti tadā hi yat eva adaḥ purastāt avayavaṣaṣṭhyartham prakaṛtam etat uttaratra anuvṛttam sat sthāneyogārtham bhaviṣyati . \\
\noindent \textbf{(P\_1,1.49.2)} sampratyayamātram etat bhavati . \\
\noindent \textbf{(P\_1,1.49.2)} na hi anuccārya śabdam liṅgam śakyam āsaṅktum . \\
\noindent \textbf{(P\_1,1.49.2)} evam tarhi ādeśe tat liṅgam kariṣyate tat prakṛtim āskantsyati . \\
\noindent \textbf{(P\_1,1.49.2)} yadi niyamaḥ kriyate yatra ekā ṣaṣṭhī anekam ca viśeṣyam tatra na sidhyati : aṅgasya , halaḥ , aṇaḥ , samprasāraṇasya iti . \\
\noindent \textbf{(P\_1,1.49.2)} hal api viśeṣyaḥ aṇ api viśeṣyaḥ samprasāraṇam api viśeṣyam . \\
\noindent \textbf{(P\_1,1.49.2)} asati punaḥ niyame kāmacāraḥ ekayā ṣaṣthyā anekam viśeṣayitum . \\
\noindent \textbf{(P\_1,1.49.2)} tat yathā . \\
\noindent \textbf{(P\_1,1.49.2)} devadattasya putraḥ pāṇiḥ kambalaḥ iti . \\
\noindent \textbf{(P\_1,1.49.2)} tasmāt na arthaḥ niyamena . \\
\noindent \textbf{(P\_1,1.49.2)} nanu ca uktam ekaśatam ṣaṣṭhyarthāḥ yāvantaḥ vā te sarve ṣaṣṭhyām uccāritāyām prāpnuvanti iti . \\
\noindent \textbf{(P\_1,1.49.2)} na eṣaḥ doṣaḥ . \\
\noindent \textbf{(P\_1,1.49.2)} yadi api loke bahavaḥ abhisambandhāḥ ārthāḥ yaunāḥ maukhāḥ srauvāḥ ca śabdasya tu śabdena kaḥ anyaḥ abhisambandhaḥ bhavitum arhati anyat ataḥ sthānāt . \\
\noindent \textbf{(P\_1,1.49.2)} śabdasya api śabdena anantarādayaḥ abhisambandhāḥ . \\
\noindent \textbf{(P\_1,1.49.2)} asteḥ bhūḥ bhavati iti sandehaḥ : sthāne anantare samīpe iti . \\
\noindent \textbf{(P\_1,1.49.2)} sandehamātram etat bhavati . \\
\noindent \textbf{(P\_1,1.49.2)} sarvasandeheṣu ca idam upatiṣṭhate : vyākhyānataḥ viśeṣapratipattiḥ na hi sandehāt alakṣaṇam iti . \\
\noindent \textbf{(P\_1,1.49.2)} sthāne iti vyākhyāsyāmaḥ . \\
\noindent \textbf{(P\_1,1.49.2)} na tarhi idānīm ayam yogaḥ vaktavyaḥ . \\
\noindent \textbf{(P\_1,1.49.2)} vaktavyaḥ ca . \\
\noindent \textbf{(P\_1,1.49.2)} kim prayojanam . \\
\noindent \textbf{(P\_1,1.49.2)} ṣaṣṭhyantam sthānena yathā yujyeta yataḥ ṣaṣṭhī uccāritā . \\
\noindent \textbf{(P\_1,1.49.2)} kim kṛtam bhavati . \\
\noindent \textbf{(P\_1,1.49.2)} nirdiśyamānasya ādeśāḥ bhavanti iti eṣā paribhāṣā na kartavyā bhavati . \\
\noindent \textbf{(P\_1,1.50.1)} kim udāharaṇam . \\
\noindent \textbf{(P\_1,1.50.1)} ikaḥ yaṇ aci : dadhi atra madhu atra : tālusthānasya tālusthānaḥ oṣṭhasthānasya oṣṭhasthānaḥ yathā syāt . \\
\noindent \textbf{(P\_1,1.50.1)} na etat asti . \\
\noindent \textbf{(P\_1,1.50.1)} saṅkhyātānudeśena api etat siddham . \\
\noindent \textbf{(P\_1,1.50.1)} idam tarhi : tasthasthamipām tāmtamtāmaḥ iti ekārthasya ekārthaḥ dvyarthasya dvyarthaḥ bahvarthasya bahuvarthaḥ yathā syāt . \\
\noindent \textbf{(P\_1,1.50.1)} nanu ca etat api saṅkhyātānudeśena eva siddham . \\
\noindent \textbf{(P\_1,1.50.1)} idam tarhi : akaḥ savarṇe dīrghaḥ iti : daṇḍāgram , kṣupāgram , dadhi indraḥ , madhu uṣṭraḥ iti : kaṇṭhasthānayoḥ kaṇṭhasthānaḥ tālusthānayoḥ tālusthānaḥ oṣṭhasthānayoḥ oṣṭhasthānaḥ yathā syāt iti . \\
\noindent \textbf{(P\_1,1.50.1)} atha sthāne iti vartamāne punaḥ sthānagrahaṇam kimartham . \\
\noindent \textbf{(P\_1,1.50.1)} yatra anekavidham āntaryam tatra sthānataḥ eva āntaryam balīyaḥ yathā syāt . \\
\noindent \textbf{(P\_1,1.50.1)} kim punaḥ tat . \\
\noindent \textbf{(P\_1,1.50.1)} cetā stotā : pramāṇataḥ akāraḥ guṇaḥ prāpnoti sthānataḥ ekāraukārau . \\
\noindent \textbf{(P\_1,1.50.1)} punaḥ sthānagrahaṇāt ekāraukārau bhavataḥ . \\
\noindent \textbf{(P\_1,1.50.1)} atha tamabgrahaṇam kimartham . \\
\noindent \textbf{(P\_1,1.50.1)} jhayaḥ haḥ anyatarasyām iti atra soṣmaṇaḥ soṣmāṇaḥ iti dvitīyāḥ prasaktāḥ nādavataḥ nādavantaḥ iti tṛtīyāḥ . \\
\noindent \textbf{(P\_1,1.50.1)} tamapgrahaṇena soṣmāṇaḥ nādavantaḥ ca te bhavanti caturthāḥ : vāg ghasati triṣṭub bhasati iti . \\
\noindent \textbf{(P\_1,1.50.2)} kimartham punaḥ idam ucyate . \\
\noindent \textbf{(P\_1,1.50.2)} sthāninaḥ ekatvanirdeśāt anekādeśanirdeśāt ca sarvaprasaṅgaḥ tasmāt sthānentaratamavacanam . sthānī ekatvena nirdiśyate : akaḥ iti , anekaḥ ca punaḥ ādeśaḥ pratinirdiśyate : dīrghaḥ iti . \\
\noindent \textbf{(P\_1,1.50.2)} sthāninaḥ ekatvanirdeśāt anekādeśanirdeśāt ca sarvaprasaṅgaḥ . \\
\noindent \textbf{(P\_1,1.50.2)} sarve sarvatra prāpnuvanti . \\
\noindent \textbf{(P\_1,1.50.2)} iṣyate ca antaratamāḥ eva syuḥ iti . \\
\noindent \textbf{(P\_1,1.50.2)} tat ca antareṇa yatnam na sidhyati . \\
\noindent \textbf{(P\_1,1.50.2)} tasmāt sthāne antaratamaḥ iti vacanam niyamārtham . \\
\noindent \textbf{(P\_1,1.50.2)} evamartham idam ucyate . \\
\noindent \textbf{(P\_1,1.50.2)} asti prayojanam etat . \\
\noindent \textbf{(P\_1,1.50.2)} kim tarhi iti . \\
\noindent \textbf{(P\_1,1.50.2)} yathā punaḥ iyam antaratamanirvṛtttiḥ sā kim prakṛtitaḥ bhavati : sthānini antaratame ṣaṣṭhī , āhosvit ādeśataḥ : sthāne prāpyamāṇānām antaratamaḥ ādeśaḥ bhavati iti . \\
\noindent \textbf{(P\_1,1.50.2)} kutaḥ punaḥ iyam vicāraṇā . \\
\noindent \textbf{(P\_1,1.50.2)} ubhayathā api tulyā saṃhitā : sthānentaratama , uraṇ raparaḥ iti . \\
\noindent \textbf{(P\_1,1.50.2)} kim ca ataḥ . \\
\noindent \textbf{(P\_1,1.50.2)} yadi prakṛtitaḥ : ikaḥ yaṇ aci : yaṇām ye antaratamāḥ ikaḥ tatra ṣaṣṭhī , yatra ṣaṣṭhī tatra ādeśāḥ bhavanti iti iha eva syāt : dadhi atra madhu atra . \\
\noindent \textbf{(P\_1,1.50.2)} kumārī atra brahmabandhvartham iti atra na syāt . \\
\noindent \textbf{(P\_1,1.50.2)} ādeśataḥ punaḥ antaratamanirvṛttau satyām sarvatra ṣaṣṭhī , yatra ṣaṣṭhī tatra ādeśāḥ bhavanti iti sarvatra siddham bhavati . \\
\noindent \textbf{(P\_1,1.50.2)} tathā ikaḥ guṇavṛddhī : guṇavṛddhyoḥ ye antaratamāḥ ikaḥ tatra ṣaṣṭhī , yatra ṣaṣṭhī tatra ādeśāḥ bhavanti iti iha eva syāt : netā lavitā nāyakaḥ lāvakaḥ . \\
\noindent \textbf{(P\_1,1.50.2)} cetā stotā cāyakaḥ stāvakaḥ iti atra na syāt . \\
\noindent \textbf{(P\_1,1.50.2)} ādeśataḥ punaḥ antaratamanirvṛttau satyām sarvatra ṣaṣṭhī , yatra ṣaṣṭhī tatra ādeśāḥ bhavanti iti sarvatra siddham bhavati . \\
\noindent \textbf{(P\_1,1.50.2)} tathā ṛvarṇasya guṇavṛddhiprasaṅge guṇavṛddhyoḥ yat antaratamam ṛvarṇam tatra ṣaṣṭhī , yatra ṣaṣṭhī tatra ādeśāḥ bhavanti iti iha eva syāt : kartā hartā , āstārakaḥ , nipārakaḥ . \\
\noindent \textbf{(P\_1,1.50.2)} āstaritā niparitā kārakaḥ , hārakaḥ iti atra na syāt . \\
\noindent \textbf{(P\_1,1.50.2)} ādeśataḥ punaḥ antaratamanirvṛttau satyām sarvatra ṣaṣṭhī , yatra ṣaṣṭhī tatra ādeśāḥ bhavanti iti sarvatra siddham bhavati . \\
\noindent \textbf{(P\_1,1.50.2)} atha ādeśataḥ antaratamanirvṛttau satyām ayam doṣaḥ : vāntaḥ yi pratyaye : sthāninirdeśaḥ kartavyaḥ . \\
\noindent \textbf{(P\_1,1.50.2)} okāraukārayoḥ iti vaktavyam ekāraikārayoḥ mā bhūt iti . \\
\noindent \textbf{(P\_1,1.50.2)} prakṛtitaḥ punaḥ antaratamanirvṛttau satyām vāntādeśasya yā antaratamā prakṛtiḥ tatra ṣaṣṭhī , yatra ṣaṣṭhī tatra ādeśāḥ bhavanti iti antareṇa sthāninirdeśam siddham bhavati . \\
\noindent \textbf{(P\_1,1.50.2)} ādeśataḥ api antaratamanirvṛttau satyām na doṣaḥ . \\
\noindent \textbf{(P\_1,1.50.2)} katham . \\
\noindent \textbf{(P\_1,1.50.2)} vāntagrahaṇam na kariṣyate . \\
\noindent \textbf{(P\_1,1.50.2)} yi pratyaye ecaḥ ayādayaḥ bhavanti iti eva. yadi na kriyate ceyam , jeyam iti atra api prāpnoti . \\
\noindent \textbf{(P\_1,1.50.2)} kṣayyajayyau śakyārthe iti etat niyamārtham bhaviṣyati : kṣijyoḥ eva ecaḥ iti . \\
\noindent \textbf{(P\_1,1.50.2)} tayoḥ tarhi śakyārthāt anyatra api prāpnoti : kṣeyam pāpam jeyaḥ vṛṣalaḥ iti . \\
\noindent \textbf{(P\_1,1.50.2)} ubhayataḥ niyamaḥ vijñāsyate : kṣijyoḥ eva ecaḥ anayoḥ ca śakyāṛthe eva iti . \\
\noindent \textbf{(P\_1,1.50.2)} iha api tarhi niyamāt na prāpnoti : lavyam , pavyam avaśyalāvyam avaśyapāvyam iti . \\
\noindent \textbf{(P\_1,1.50.2)} tulyajātīyasya niyamaḥ . \\
\noindent \textbf{(P\_1,1.50.2)} kaḥ ca tulyajātīyaḥ . \\
\noindent \textbf{(P\_1,1.50.2)} yathājātīyakaḥ kṣijyoḥ ec . \\
\noindent \textbf{(P\_1,1.50.2)} kathañjātīyakaḥ kṣijyoḥ ec . \\
\noindent \textbf{(P\_1,1.50.2)} ekāraḥ . \\
\noindent \textbf{(P\_1,1.50.2)} evam api rāyam icchati raiyati atra api prāpnoti . \\
\noindent \textbf{(P\_1,1.50.2)} rāyiḥ chāndasaḥ . \\
\noindent \textbf{(P\_1,1.50.2)} dṛṣṭānuvidhiḥ chandasi bhavati . \\
\noindent \textbf{(P\_1,1.50.2)} ūdupadhayāḥ gohaḥ : ādeśataḥ antaratamanirvṛttau satyām upadhāgrahaṇam kartavyam . \\
\noindent \textbf{(P\_1,1.50.2)} prakṛtitaḥ punaḥ antaratamanirvṛttau satyām ūkārasya gohaḥ yā antaratamā prakṛtiḥ tatra ṣaṣṭhī , yatra ṣaṣṭhī tatra ādeśāḥ bhavanti iti antareṇa upadhāgrahaṇam siddham bhavati . \\
\noindent \textbf{(P\_1,1.50.2)} ādeśataḥ api antaratamanirvṛttau satyām na doṣaḥ . \\
\noindent \textbf{(P\_1,1.50.2)} kriyate etat nyāse eva . \\
\noindent \textbf{(P\_1,1.50.2)} radābhyām niṣṭhātaḥ naḥ pūrvasya ca daḥ : ādeśataḥ antaratamanirvṛttau satyām takāragrahaṇam kartavyam . \\
\noindent \textbf{(P\_1,1.50.2)} prakṛtitaḥ punaḥ antaratamanirvṛttau satyām nakārasya niṣṭhāyām yā antaratamā prakṛtiḥ tatra ṣaṣṭhī , yatra ṣaṣṭhī tatra ādeśāḥ bhavanti iti antareṇa takāragrahaṇam siddham bhavati . \\
\noindent \textbf{(P\_1,1.50.2)} ādeśataḥ api antaratamanirvṛttau satyām na doṣaḥ . \\
\noindent \textbf{(P\_1,1.50.2)} kriyate etat nyāse eva . \\
\noindent \textbf{(P\_1,1.50.3)} kim punaḥ idam nirvartakam : antaratamāḥ anena nirvartyante , āhosvit pratipādakam : anyena nirvṛttānām anena pratipattiḥ . \\
\noindent \textbf{(P\_1,1.50.3)} kaḥ ca atra viśeṣaḥ . \\
\noindent \textbf{(P\_1,1.50.3)} sthāne antaratamanirvatake sthāninivṛttiḥ . \\
\noindent \textbf{(P\_1,1.50.3)} sthāne antaratamanirvatake sarvasthāninām nivṛttiḥ prāpnoti . \\
\noindent \textbf{(P\_1,1.50.3)} asya api prāpnoti : dadhi madhu . \\
\noindent \textbf{(P\_1,1.50.3)} astu . \\
\noindent \textbf{(P\_1,1.50.3)} na kaḥ cit anyaḥ ādeśaḥ pratinirdiśyate . \\
\noindent \textbf{(P\_1,1.50.3)} tatra āntaryataḥ dadhiśabdasya dadhiśabdaḥ eva madhuśabdasya madhuśabdaḥ eva ādeśaḥ bhaviṣyati . \\
\noindent \textbf{(P\_1,1.50.3)} yadi ca evam kva cit vairūpyam tatra doṣaḥ syāt : bisam bisam , musalam musalam . \\
\noindent \textbf{(P\_1,1.50.3)} iṇkoḥ iti ṣatvam prāpnoti . \\
\noindent \textbf{(P\_1,1.50.3)} api ca iṣṭā vyavasthā na prakalpeta . \\
\noindent \textbf{(P\_1,1.50.3)} tat yathā tapte bhrāṣṭre tilāḥ kṣiptāḥ muhūrtam api na avatiṣṭhante evam ime varṇāḥ muhūrtam api na avatiṣṭheran . \\
\noindent \textbf{(P\_1,1.50.3)} astu tarhi pratipādakam : anyena nirvṛttānām anena pratipattiḥ . \\
\noindent \textbf{(P\_1,1.50.3)} nirvṛttapratipattau nirvṛttiḥ . \\
\noindent \textbf{(P\_1,1.50.3)} nirvṛttapratipattau nirvṛttiḥ na sidhyati . \\
\noindent \textbf{(P\_1,1.50.3)} sarve sarvatra prāpnuvanti . \\
\noindent \textbf{(P\_1,1.50.3)} kim tarhi ucyate nirvṛttiḥ na sidhyati iti . \\
\noindent \textbf{(P\_1,1.50.3)} na sādhīyaḥ nirvṛttiḥ siddhā bhavati . \\
\noindent \textbf{(P\_1,1.50.3)} na brūmaḥ nirvṛttiḥ na sidhyati iti . \\
\noindent \textbf{(P\_1,1.50.3)} kim tarhi . \\
\noindent \textbf{(P\_1,1.50.3)} iṣṭā vyavasthā na prakalpeta . \\
\noindent \textbf{(P\_1,1.50.3)} na sarve sarvatra iṣyante . \\
\noindent \textbf{(P\_1,1.50.3)} idam idānīm kimartham syāt . \\
\noindent \textbf{(P\_1,1.50.3)} anarthakam ca . \\
\noindent \textbf{(P\_1,1.50.3)} anarthakam etat syāt . \\
\noindent \textbf{(P\_1,1.50.3)} yaḥ hi bhuktavantam brūyāt mā bhukthāḥ iti kim tena kṛtam syāt . \\
\noindent \textbf{(P\_1,1.50.3)} uktam vā . \\
\noindent \textbf{(P\_1,1.50.3)} kim uktam . \\
\noindent \textbf{(P\_1,1.50.3)} siddham tu ṣaṣṭhyadhikāre vacanāt iti . \\
\noindent \textbf{(P\_1,1.50.3)} ṣaṣṭhyadhikāre ayam yogaḥ kartvyaḥ : sthāne antaratamaḥ ṣaṣṭhīnirdiṣṭasya . \\
\noindent \textbf{(P\_1,1.50.4)} pratyātmavacanam ca . \\
\noindent \textbf{(P\_1,1.50.4)} pratyātmam iti ca vaktavyam . \\
\noindent \textbf{(P\_1,1.50.4)} kim prayojanam . \\
\noindent \textbf{(P\_1,1.50.4)} yaḥ yasya antaratamaḥ sa tasya sthāne yathā syāt anyasya antaratamaḥ anyasya sthāne mā bhūt iti . \\
\noindent \textbf{(P\_1,1.50.4)} pratyātmavacanam aśiṣyam svabhāvasiddhatvāt . \\
\noindent \textbf{(P\_1,1.50.4)} pratyātmavacanam aśiṣyam . \\
\noindent \textbf{(P\_1,1.50.4)} kim kāraṇam . \\
\noindent \textbf{(P\_1,1.50.4)} svabhāvasiddhatvāt . \\
\noindent \textbf{(P\_1,1.50.4)} svabhāvataḥ etat siddham . \\
\noindent \textbf{(P\_1,1.50.4)} tat yathā : samājeṣu samāśeṣu samavāyeṣu ca āsyatām iti ucyate . \\
\noindent \textbf{(P\_1,1.50.4)} na ca ucyate pratyātmam iti pratyātmam ca āsate . \\
\noindent \textbf{(P\_1,1.50.4)} antaratamavacanam ca . \\
\noindent \textbf{(P\_1,1.50.4)} antaratamavacanam ca aśiṣyam . \\
\noindent \textbf{(P\_1,1.50.4)} yogaḥ ca api ayam aśiṣyaḥ . \\
\noindent \textbf{(P\_1,1.50.4)} kutaḥ . \\
\noindent \textbf{(P\_1,1.50.4)} svabhāvasiddhatvāt eva . \\
\noindent \textbf{(P\_1,1.50.4)} tat yathā : samājeṣu samāśeṣu samavāyeṣu ca āsyatām iti ukte na eva kṛśāḥ kṛśaiḥ saha āsata na pāṇḍavaḥ pāṇḍubhiḥ . \\
\noindent \textbf{(P\_1,1.50.4)} yeṣām eva kim cit arthakṛtam āntaryam taiḥ eva saḥ āsate . \\
\noindent \textbf{(P\_1,1.50.4)} tathā gāvaḥ divasam caritavatyaḥ yaḥ yasyāḥ prasavaḥ bhavati tena saha śerate . \\
\noindent \textbf{(P\_1,1.50.4)} tathā yāni etāni goyuktakāni saṅghuṣṭakāni bhavanti tāni anyonyam paśyanti sabdam kurvanti . \\
\noindent \textbf{(P\_1,1.50.4)} evam tāvat cetanāvatsu . \\
\noindent \textbf{(P\_1,1.50.4)} acetaneṣu api . \\
\noindent \textbf{(P\_1,1.50.4)} loṣṭaḥ kṣiptaḥ bāhuvegam gatvā na eva tiryak gacchati na ūrdhvam ārohati pṛthivīvikāraḥ pṛthivīm gacchati āntaryataḥ . \\
\noindent \textbf{(P\_1,1.50.4)} tathā yā etāḥ āntarikṣyaḥ sūkṣmāḥ āpaḥ tāsām vikāraḥ dhūmaḥ saḥ ākāśadeśe nivāte na eva tiryak gacchati na avāk avarohati abvikāraḥ apaḥ eva gacchati āntaryataḥ . \\
\noindent \textbf{(P\_1,1.50.4)} tathā jyotiṣaḥ vikāraḥ arciḥ ākāśadeśe nivāte suprajvalitaḥ na eva tiryak gacchati na avāk avarohati jyotiṣaḥ vikāraḥ jyotiḥ eva gacchati āntaryataḥ . \\
\noindent \textbf{(P\_1,1.50.5)} vyañjanasvaravyatikrame ca tatkālaprasaṅgaḥ . \\
\noindent \textbf{(P\_1,1.50.5)} vyañjanavyatikrame svaravyatikrame ca tatkālatā prāpnoti . \\
\noindent \textbf{(P\_1,1.50.5)} vyañjanavyatikrame : iṣṭam uptam . \\
\noindent \textbf{(P\_1,1.50.5)} āntaryataḥ ardhamātrikasya vyañjanasya ardhamātrikaḥ ik prāpnoti . \\
\noindent \textbf{(P\_1,1.50.5)} na eva loke na ca vede ardhamātrikaḥ ik asti . \\
\noindent \textbf{(P\_1,1.50.5)} kaḥ tarhi . \\
\noindent \textbf{(P\_1,1.50.5)} mātrikaḥ . \\
\noindent \textbf{(P\_1,1.50.5)} yaḥ asti saḥ bhaviṣyati . \\
\noindent \textbf{(P\_1,1.50.5)} svaravyatikrame : dadhi atra madhu atra kumārī atra brahmabandhvartham iti . \\
\noindent \textbf{(P\_1,1.50.5)} āntaryataḥ mātrikasya dvimātrikasya ikaḥ mātrikaḥ dvimātrikaḥ vā yaṇ prāpnoti . \\
\noindent \textbf{(P\_1,1.50.5)} na eva loke na ca vede mātrikaḥ dvimātrikaḥ vā yaṇ asti .kaḥ tarhi . \\
\noindent \textbf{(P\_1,1.50.5)} ardhamātrikaḥ . \\
\noindent \textbf{(P\_1,1.50.5)} yaḥ asti saḥ bhaviṣyati . \\
\noindent \textbf{(P\_1,1.50.5)} akṣu ca anekavarṇādeśeṣu . \\
\noindent \textbf{(P\_1,1.50.5)} akṣu ca anekavarṇādeśeṣu tatkālatā prāpnoti . \\
\noindent \textbf{(P\_1,1.50.5)} idamaḥ iś : āntaryataḥ ardhtṛtīyamātrasya idamaḥ sthāne ardhtṛtīyamātram ivarṇam prāpnoti . \\
\noindent \textbf{(P\_1,1.50.5)} na eṣaḥ doṣaḥ . \\
\noindent \textbf{(P\_1,1.50.5)} bhāvyamānena savarṇānām grahaṇam na iti evam na bhaviṣyati . \\
\noindent \textbf{(P\_1,1.50.5)} guṇavṛddhyejbhāveṣu ca . \\
\noindent \textbf{(P\_1,1.50.5)} guṇavṛddhyejbhāveṣu ca tatkālatā prāpnoti : khaṭvā indraḥ khaṭvendraḥ khaṭvā udakam khaṭvodakam khaṭvā īṣā khaṭveṣā khaṭvā ūḍhā khaṭvoḍhā khaṭvā elakā khaṭvailakā khaṭvā odanaḥ khaṭvaudanaḥ , khaṭvā aitikayanaḥ khaṭvaitikāyanaḥ , khaṭvā aupagavaḥ khaṭvaupagavaḥ iti . \\
\noindent \textbf{(P\_1,1.50.5)} āntaryataḥ trimātrcaturmātrāṇām sthāninām trimātracaturmātrāḥ ādeśāḥ prāpnuvanti . \\
\noindent \textbf{(P\_1,1.50.5)} na eṣaḥ doṣāḥ . \\
\noindent \textbf{(P\_1,1.50.5)} tapare guṇavṛddhī . \\
\noindent \textbf{(P\_1,1.50.5)} nanu ca taḥ paraḥ yasmāt saḥ ayam taparaḥ . \\
\noindent \textbf{(P\_1,1.50.5)} na iti āha . \\
\noindent \textbf{(P\_1,1.50.5)} tāt api paraḥ taparaḥ . \\
\noindent \textbf{(P\_1,1.50.5)} yadi tāt api paraḥ taparaḥ ṝdoḥ ap iti iha eva syāt : yavaḥ , stavaḥ . \\
\noindent \textbf{(P\_1,1.50.5)} lavaḥ , pavaḥ iti atra na syāt . \\
\noindent \textbf{(P\_1,1.50.5)} na eṣaḥ takāraḥ . \\
\noindent \textbf{(P\_1,1.50.5)} kaḥ tarhi . \\
\noindent \textbf{(P\_1,1.50.5)} dakāraḥ . \\
\noindent \textbf{(P\_1,1.50.5)} kim dakāre prayojanam . \\
\noindent \textbf{(P\_1,1.50.5)} atha kim takāre prayojanam . \\
\noindent \textbf{(P\_1,1.50.5)} yadi asandehārthaḥ takāraḥ dakāraḥ api . \\
\noindent \textbf{(P\_1,1.50.5)} atha mukhasukhārthaḥ takāraḥ dakāraḥ api . \\
\noindent \textbf{(P\_1,1.50.5)} ejbhāve : kurvāte kurvāthe . \\
\noindent \textbf{(P\_1,1.50.5)} āntaryataḥ ardhatṛtīyamātrasya ṭisañjñakasya ardhatṛtīyamātraḥ eḥ prāpnoti . \\
\noindent \textbf{(P\_1,1.50.5)} na eva loke na ca vede ardhatṛtīyamātraḥ eḥ asti . \\
\noindent \textbf{(P\_1,1.50.5)} ṛvarṇasya guṇavṛddhiprasaṅge sarvaprasaṅgaḥ aviśeṣāt . \\
\noindent \textbf{(P\_1,1.50.5)} ṛvarṇasya guṇavṛddhiprasaṅge sarvaprasaṅgaḥ . \\
\noindent \textbf{(P\_1,1.50.5)} sarve guṇavṛddhisañjñakāḥ ṛvarṇasya sthāne prāpnuvanti . \\
\noindent \textbf{(P\_1,1.50.5)} kim kāraṇam . \\
\noindent \textbf{(P\_1,1.50.5)} aviśeṣāt . \\
\noindent \textbf{(P\_1,1.50.5)} na hi kaḥ cit viśeṣaḥ upādīyate evañjātīyakaḥ guṇavṛddhisañjñakaḥ ṛvarṇasya sthāne bhavati iti . \\
\noindent \textbf{(P\_1,1.50.5)} anupādīyamāne viśeṣe sarvaprasaṅgaḥ . \\
\noindent \textbf{(P\_1,1.50.5)} na vā ṛvarṇasya sthāne raparaprasaṅgāt avarṇasya āntaryam . \\
\noindent \textbf{(P\_1,1.50.5)} na vā eṣaḥ doṣaḥ . \\
\noindent \textbf{(P\_1,1.50.5)} kim kāraṇam . \\
\noindent \textbf{(P\_1,1.50.5)} ṛvarṇasya sthāne raparaprasaṅgāt . \\
\noindent \textbf{(P\_1,1.50.5)} uḥ sthāne aṇ prasajyamānaḥ eva raparaḥ bhavati iti ucyate . \\
\noindent \textbf{(P\_1,1.50.5)} tatra ṛvarṇasya āntaryataḥ rephavataḥ rephavān akāraḥ eva antaratamaḥ bhavati . \\
\noindent \textbf{(P\_1,1.50.5)} sarvādeśaprasaṅgaḥ tu anekāltvāt . \\
\noindent \textbf{(P\_1,1.50.5)} sarvādeśaprasaṅgaḥ tu guṇavṛddhisañjñakaḥ ṛvarṇasya prāpnoti . \\
\noindent \textbf{(P\_1,1.50.5)} kim kāraṇam . \\
\noindent \textbf{(P\_1,1.50.5)} anekāltvāt . \\
\noindent \textbf{(P\_1,1.50.5)} anekāl śit sarvasya iti . \\
\noindent \textbf{(P\_1,1.50.5)} na vā anekāltvasya tadāśrayatvāt ṛvarṇādeśasya avighātaḥ . \\
\noindent \textbf{(P\_1,1.50.5)} na vā eṣaḥ doṣaḥ . \\
\noindent \textbf{(P\_1,1.50.5)} kim kāraṇam . \\
\noindent \textbf{(P\_1,1.50.5)} anekāltvasya tadāśrayatvāt . \\
\noindent \textbf{(P\_1,1.50.5)} yadā ayam uḥ sthāne tadā anekāl . \\
\noindent \textbf{(P\_1,1.50.5)} anekāltvasya tadāśrayatvāt ṛvarṇādeśasya vighātaḥ na bhaviṣyati . \\
\noindent \textbf{(P\_1,1.50.5)} athavā anāntaryam eva etayoḥ āntaryam . \\
\noindent \textbf{(P\_1,1.50.5)} ekasya api antaratamā prakṛtiḥ na asti aparasya api antaratamaḥ ādeśaḥ na asti . \\
\noindent \textbf{(P\_1,1.50.5)} etat eva etayoḥ āntaryam . \\
\noindent \textbf{(P\_1,1.50.5)} samprayogaḥ vā naṣṭāśvadagdharathavat . \\
\noindent \textbf{(P\_1,1.50.5)} atha vā naṣṭāśvadagdharathavat samprayogaḥ bhavati . \\
\noindent \textbf{(P\_1,1.50.5)} tat yathā : tava aśvaḥ naṣṭaḥ mama api rathaḥ dagdhaḥ . \\
\noindent \textbf{(P\_1,1.50.5)} ubhau samprayujyāvahai iti . \\
\noindent \textbf{(P\_1,1.50.5)} evam iha api : tava api antaratamā prakṛtiḥ na asti mama api antaratamaḥ ādeśaḥ na asti . \\
\noindent \textbf{(P\_1,1.50.5)} astu nau samprayogaḥ iti . \\
\noindent \textbf{(P\_1,1.50.5)} viṣamaḥ upanyāsaḥ . \\
\noindent \textbf{(P\_1,1.50.5)} cetanāvatsu arthāt prakaraṇāt vā loke samprayogaḥ bhavati . \\
\noindent \textbf{(P\_1,1.50.5)} varṇāḥ ca punaḥ acetanāḥ . \\
\noindent \textbf{(P\_1,1.50.5)} tatra kiṅkṛtaḥ samprayogaḥ . \\
\noindent \textbf{(P\_1,1.50.5)} yadi api varṇāḥ acetanāḥ yaḥ tu asau prayuṅkte saḥ cetanāvān . \\
\noindent \textbf{(P\_1,1.50.5)} ejavarṇayoḥ ādeśe avarṇam sthāninaḥ avarṇapradhānatvāt . \\
\noindent \textbf{(P\_1,1.50.5)} ejavarṇayoḥ ādeśe avarṇam prāpnoti :khaṭvā elakā , mālā aupagavaḥ . \\
\noindent \textbf{(P\_1,1.50.5)} kim kāraṇam . \\
\noindent \textbf{(P\_1,1.50.5)} sthāninaḥ avarṇapradhānatvāt . \\
\noindent \textbf{(P\_1,1.50.5)} sthānī hi atra avarṇapradhānaḥ . \\
\noindent \textbf{(P\_1,1.50.5)} siddham tu ubhayāntaryāt . \\
\noindent \textbf{(P\_1,1.50.5)} siddham etat . \\
\noindent \textbf{(P\_1,1.50.5)} katham . \\
\noindent \textbf{(P\_1,1.50.5)} ubhayoḥ yaḥ antaratamaḥ tena bhavitavyam . \\
\noindent \textbf{(P\_1,1.50.5)} na ca avarṇam ubhayoḥ antaratamam . \\
\noindent \textbf{(P\_1,1.51.1)} kim idam uraṇraparavacanam anyanivṛttyartham : uḥ sthāne aṇ eva bhavati raparaḥ ca iti , āhosvit raparatvam anena vidhīyate : uḥ sthāne aṇ ca anaṇ ca aṇ tu raparaḥ eva . \\
\noindent \textbf{(P\_1,1.51.1)} kaḥ ca atra viśeṣaḥ . \\
\noindent \textbf{(P\_1,1.51.1)} uraṇraparavacanam anyanivṛttyartham cet udāttādiṣu doṣaḥ . \\
\noindent \textbf{(P\_1,1.51.1)} uraṇraparavacanam anyanivṛttyartham cet udāttādiṣu doṣaḥ bhavati . \\
\noindent \textbf{(P\_1,1.51.1)} ke punaḥ udāttādayaḥ . \\
\noindent \textbf{(P\_1,1.51.1)} udāttānudāttasvaritānunāsikāḥ . \\
\noindent \textbf{(P\_1,1.51.1)} kṛtiḥ , hṛtiḥ , kṛtam , hṛtam , prakṛtam , prahṛtam nṝŚḥ pāhi . \\
\noindent \textbf{(P\_1,1.51.1)} astu tarhi uḥ sthāne aṇ ca anaṇ ca aṇ tu raparaḥ iti . \\
\noindent \textbf{(P\_1,1.51.1)} yaḥ uḥ sthāne saḥ raparaḥ iti cet guṇavṛddhyoḥ avarṇāpratipattiḥ . \\
\noindent \textbf{(P\_1,1.51.1)} yaḥ uḥ sthāne saḥ raparaḥ iti cet guṇavṛddhyoḥ avarṇāpratipattiḥ . \\
\noindent \textbf{(P\_1,1.51.1)} kartā hartā vārṣagaṇyaḥ . \\
\noindent \textbf{(P\_1,1.51.1)} kim hi sādhīyaḥ ṛvarṇasya asavarṇe yat avarṇam syāt na punaḥ eṅaicau . \\
\noindent \textbf{(P\_1,1.51.1)} pūrvasmin api pakṣe eṣaḥ doṣaḥ . \\
\noindent \textbf{(P\_1,1.51.1)} kim hi sādhīyaḥ tatra api ṛvarṇasya asavarṇe yat avarṇam syāt na punaḥ ivarṇovarṇau . \\
\noindent \textbf{(P\_1,1.51.1)} atha matam etat uḥ sthāne aṇaḥ ca anaṇaḥ ca prasaṅge aṇ eva bhavati raparaḥ ca iti siddhā pūrvasmin pakṣe avarṇasya pratipattiḥ . \\
\noindent \textbf{(P\_1,1.51.1)} yat tu tat uktam udāttādiṣu doṣaḥ bhavati iti iha saḥ doṣaḥ jāyate . \\
\noindent \textbf{(P\_1,1.51.1)} na jāyate . \\
\noindent \textbf{(P\_1,1.51.1)} jāyate saḥ doṣaḥ . \\
\noindent \textbf{(P\_1,1.51.1)} katham . \\
\noindent \textbf{(P\_1,1.51.1)} udāttaḥ iti anena aṇaḥ api pratinirdiśyante anaṇaḥ api . \\
\noindent \textbf{(P\_1,1.51.1)} yadi api pratinirdiśyante na tu prāpnuvanti . \\
\noindent \textbf{(P\_1,1.51.1)} kim kāraṇam . \\
\noindent \textbf{(P\_1,1.51.1)} sthāne antaratamaḥ bhavati . \\
\noindent \textbf{(P\_1,1.51.1)} kutaḥ nu khalu dvayoḥ paribhāṣayoḥ sāvakāśayoḥ samavasthitayoḥ sthāne antaratamaḥ uḥ aṇ raparaḥ iti ca sthāne antaratamaḥ iti anayā paribhāṣayā vyavasthā bhaviṣyati na punaḥ uḥ aṇ raparaḥ iti . \\
\noindent \textbf{(P\_1,1.51.1)} ataḥ kim . \\
\noindent \textbf{(P\_1,1.51.1)} ataḥ eṣaḥ doṣaḥ jāyate : udāttādiṣu doṣaḥ iti . \\
\noindent \textbf{(P\_1,1.51.1)} ye ca api ete ṛvarṇasya sthāne pratipadam ādeśāḥ ucyante teṣu raparatvam na prāpnoti : ṝtaḥ it dhātoḥ ut oṣṭhyapūrvasya iti . \\
\noindent \textbf{(P\_1,1.51.1)} siddham tu prasaṅge raparatvāt . \\
\noindent \textbf{(P\_1,1.51.1)} siddham etat . \\
\noindent \textbf{(P\_1,1.51.1)} katham . \\
\noindent \textbf{(P\_1,1.51.1)} prasaṅge raparatvāt . \\
\noindent \textbf{(P\_1,1.51.1)} uḥ sthāne aṇ prasajyamānaḥ eva raparaḥ bhavati iti . \\
\noindent \textbf{(P\_1,1.51.1)} kim vaktavyam etat . \\
\noindent \textbf{(P\_1,1.51.1)} na hi . \\
\noindent \textbf{(P\_1,1.51.1)} katham anucyamānam gaṃsyate . \\
\noindent \textbf{(P\_1,1.51.1)} sthāne iti vartate sthānaśabdaḥ ca prasaṅgavācī . \\
\noindent \textbf{(P\_1,1.51.1)} yadi evam ādeśaḥ aviśeṣitaḥ bhavati . \\
\noindent \textbf{(P\_1,1.51.1)} ādeśaḥ ca viśeṣitaḥ . \\
\noindent \textbf{(P\_1,1.51.1)} katham . \\
\noindent \textbf{(P\_1,1.51.1)} dvitīyam sthānagrahaṇam prakṛtam anuvartate. tatra evam abhisambandhaḥ kariṣyate : uḥ sthāne aṇ sthāne iti . \\
\noindent \textbf{(P\_1,1.51.1)} uḥ prasaṅge aṇ prasajyamānaḥ eva raparaḥ bhavati . \\
\noindent \textbf{(P\_1,1.51.2)} atha aṇgrahaṇam kimartham na uḥ raparaḥ iti eva ucyeta . \\
\noindent \textbf{(P\_1,1.51.2)} uḥ raparaḥ iti iyati ucyamāne kaḥ idānīm raparaḥ syāt . \\
\noindent \textbf{(P\_1,1.51.2)} yaḥ uḥ sthāne bhavati . \\
\noindent \textbf{(P\_1,1.51.2)} kaḥ ca uḥ sthāne bhavati . \\
\noindent \textbf{(P\_1,1.51.2)} ādeśaḥ . \\
\noindent \textbf{(P\_1,1.51.2)} ādeśaḥ raparaḥ iti cet rīrividhiṣu raparapratiṣedhaḥ . \\
\noindent \textbf{(P\_1,1.51.2)} ādeśaḥ raparaḥ iti cet rīrividhiṣu raparatvasya pratiṣedhaḥ vaktavyaḥ . \\
\noindent \textbf{(P\_1,1.51.2)} ke punaḥ rīrividhayaḥ . \\
\noindent \textbf{(P\_1,1.51.2)} akaṅlopānaṅanaṅrīṅriṅādeśāḥ . \\
\noindent \textbf{(P\_1,1.51.2)} akaṅ: saudhātakiḥ . \\
\noindent \textbf{(P\_1,1.51.2)} lopaḥ : paitṛṣvaseyaḥ . \\
\noindent \textbf{(P\_1,1.51.2)} ānaṅ : hotāpotārau . \\
\noindent \textbf{(P\_1,1.51.2)} anaṅ : kartā hartā . \\
\noindent \textbf{(P\_1,1.51.2)} rīṅ : mātrīyati pitrīyati . \\
\noindent \textbf{(P\_1,1.51.2)} riṅ : kriyate hriyate . \\
\noindent \textbf{(P\_1,1.51.2)} udāttādiṣu ca . \\
\noindent \textbf{(P\_1,1.51.2)} kim . \\
\noindent \textbf{(P\_1,1.51.2)} raparatvasya pratiṣedhaḥ vaktavyaḥ . \\
\noindent \textbf{(P\_1,1.51.2)} kṛtiḥ , hṛtiḥ , kṛtam , hṛtam , prakṛtam , prahṛtam nṝŚḥ pāhi . \\
\noindent \textbf{(P\_1,1.51.2)} tasmāt aṇgrahaṇam kartavyam . \\
\noindent \textbf{(P\_1,1.51.3)} ekādeśasya upasaṅkhyānam . \\
\noindent \textbf{(P\_1,1.51.3)} ekādeśasya upasaṅkhyānam kartavyam : khaṭvarśyaḥ , mālarśyaḥ . \\
\noindent \textbf{(P\_1,1.51.3)} kim punaḥ kāraṇam na sidhyati . \\
\noindent \textbf{(P\_1,1.51.3)} uḥ sthāne aṇ prasajyamānaḥ eva raparaḥ bhavati iti ucyate na ca ayam uḥ eva sthāne aṇ śiṣyate . \\
\noindent \textbf{(P\_1,1.51.3)} kim tarhi . \\
\noindent \textbf{(P\_1,1.51.3)} uḥ ca anyasya ca . \\
\noindent \textbf{(P\_1,1.51.3)} avayavagrahaṇāt siddham . \\
\noindent \textbf{(P\_1,1.51.3)} yat atra ṛvarṇam tadāśrayam raparatvam bhaviṣyati . \\
\noindent \textbf{(P\_1,1.51.3)} tat yathā māṣāḥ na bhoktavyāḥ iti miśrāḥ api na bhujyante . \\
\noindent \textbf{(P\_1,1.51.3)} avayavagrahaṇāt siddham iti cet ādeśe rāntapratiṣedhaḥ . \\
\noindent \textbf{(P\_1,1.51.3)} avyayavagrahaṇāt siddham iti cet ādeśe rāntasya pratiṣedhaḥ vaktavyaḥ : hotāpotārau . \\
\noindent \textbf{(P\_1,1.51.3)} yathā eva uḥ ca anyasya ca sthāne aṇ raparaḥ bhavati evam yaḥ uḥ sthāne aṇ ca anaṇ ca saḥ api raparaḥ syāt . \\
\noindent \textbf{(P\_1,1.51.3)} yadi punaḥ ṛvarṇāntasya sthāninaḥ raparatvam ucyeta : khaṭvarśyaḥ , mālarśyaḥ . \\
\noindent \textbf{(P\_1,1.51.3)} na evam śakyam . \\
\noindent \textbf{(P\_1,1.51.3)} iha hi doṣaḥ syāt : kartā hartā kirati girati . \\
\noindent \textbf{(P\_1,1.51.3)} ṛvarṇāntasya iti ucyate . \\
\noindent \textbf{(P\_1,1.51.3)} na ca etat ṛvarṇāntam . \\
\noindent \textbf{(P\_1,1.51.3)} nanu ca etat api vyapadeśivadbhāvena ṛvarṇāntam . \\
\noindent \textbf{(P\_1,1.51.3)} arthavatā vyapadeśivadbhāvaḥ na ca eṣaḥ arthavān . \\
\noindent \textbf{(P\_1,1.51.3)} tasmāt na evam śakyam . \\
\noindent \textbf{(P\_1,1.51.3)} na cet evam upasaṅkhyānam kartavyam . \\
\noindent \textbf{(P\_1,1.51.3)} iha ca raparatvapratiṣedhaḥ vaktavyaḥ : mātuḥ , pituḥ iti . \\
\noindent \textbf{(P\_1,1.51.3)} ubhayam na vaktavyam . \\
\noindent \textbf{(P\_1,1.51.3)} katham . \\
\noindent \textbf{(P\_1,1.51.3)} iha yaḥ dvayoḥ ṣaṣṭhīnirdiṣṭayoḥ prasaṅge bhavati labhate asau anyatarataḥ vyapadeśam . \\
\noindent \textbf{(P\_1,1.51.3)} tat yathā devadattasya putraḥ , devadattāyāḥ putraḥ iti . \\
\noindent \textbf{(P\_1,1.51.3)} katham mātuḥ pituḥ iti . \\
\noindent \textbf{(P\_1,1.51.3)} astu atra raparatvam . \\
\noindent \textbf{(P\_1,1.51.3)} kā rūpasiddiḥ . \\
\noindent \textbf{(P\_1,1.51.3)} rāt sasya iti sakārasya lopaḥ rephasya visarjanīyaḥ . \\
\noindent \textbf{(P\_1,1.51.3)} na evam śakyam . \\
\noindent \textbf{(P\_1,1.51.3)} iha hi mātuḥ karoti , pituḥ karoti iti apratyayavisarjanīyasya iti ṣatvam prasajyeta . \\
\noindent \textbf{(P\_1,1.51.3)} apratyayasvisarjanīyasya iti ucyate . \\
\noindent \textbf{(P\_1,1.51.3)} pratyayavisarjanīyaḥ ca ayam . \\
\noindent \textbf{(P\_1,1.51.3)} lupyate atra pratyayaḥ rāt sasya iti . \\
\noindent \textbf{(P\_1,1.51.3)} evam tarhi bhrātuṣputragrahaṇam jñāpakam ekādeśanimittāt ṣatvapratiṣedhasya . \\
\noindent \textbf{(P\_1,1.51.3)} yat ayam kaskādiṣu bhrātuṣputraśabdam paṭhati tat jñāpayati ācāryaḥ na ekādeśanimitttāt ṣatvam bhavati iti. \\
\noindent \textbf{(P\_1,1.51.4)} kim punaḥ ayam pūrvāntaḥ āhosvit parādiḥ āhosvit abhaktaḥ . \\
\noindent \textbf{(P\_1,1.51.4)} katham ca ayam pūrvāntaḥ syāt katham vā parādiḥ katham vā abhaktaḥ . \\
\noindent \textbf{(P\_1,1.51.4)} yadi antaḥ iti vartate tataḥ pūrvāntaḥ . \\
\noindent \textbf{(P\_1,1.51.4)} atha ādiḥ iti vartate tataḥ parādiḥ . \\
\noindent \textbf{(P\_1,1.51.4)} atha ubhayam nivṛttam tataḥ abhaktaḥ . \\
\noindent \textbf{(P\_1,1.51.4)} kaḥ ca atra viśeṣaḥ . \\
\noindent \textbf{(P\_1,1.51.4)} abhakte dīrghalatvayagabhyastasvarahalādiśeṣavisarjanīyapratiṣedhaḥ pratyayāvyavasthā ca . \\
\noindent \textbf{(P\_1,1.51.4)} yadi abhaktaḥ dīrghatvam na prāpnoti : gīḥ , pūḥ . \\
\noindent \textbf{(P\_1,1.51.4)} rephavakārāntasya dhātoḥ iti dīrghatvam na prāpnoti . \\
\noindent \textbf{(P\_1,1.51.4)} kim punaḥ kāraṇam rephavakārābhyām dhātuḥ viśeṣyate na punaḥ padam viśeṣyate rephavakārāntasya padasya iti . \\
\noindent \textbf{(P\_1,1.51.4)} na evam śakyam . \\
\noindent \textbf{(P\_1,1.51.4)} iha api prasajyeta : agniḥ , vāyuḥ iti . \\
\noindent \textbf{(P\_1,1.51.4)} evam tarhi rephavakārābhyām padam viśeṣayiṣyāmaḥ dhātunā ikam : rephavakārāntasya padasya ikaḥ dhātoḥ iti . \\
\noindent \textbf{(P\_1,1.51.4)} evam api priyam grāmaṇi kulam asya priyagrāmaṇiḥ , priyasenāniḥ atra api prāpnoti . \\
\noindent \textbf{(P\_1,1.51.4)} tasmāt dhātuḥ eva viśeṣyate . \\
\noindent \textbf{(P\_1,1.51.4)} dhātau ca viśeṣyamāṇe iha dīrghatvam na prāpnoti : gīḥ , pūḥ . \\
\noindent \textbf{(P\_1,1.51.4)} dīrgha . \\
\noindent \textbf{(P\_1,1.51.4)} latva : latvam ca na sidhyati : nijegilyate . \\
\noindent \textbf{(P\_1,1.51.4)} graḥ yaṅi iti latvam na prāpnoti . \\
\noindent \textbf{(P\_1,1.51.4)} na eṣaḥ doṣaḥ . \\
\noindent \textbf{(P\_1,1.51.4)} graḥ iti anantarayogā eṣā ṣaṣṭhī . \\
\noindent \textbf{(P\_1,1.51.4)} evam api svaḥ jegilyate iti atra api prāpnoti . \\
\noindent \textbf{(P\_1,1.51.4)} evam tarhi yaṅā ānantaryam viśeṣayiṣyāmaḥ . \\
\noindent \textbf{(P\_1,1.51.4)} atha vā graḥ iti pañcamī . \\
\noindent \textbf{(P\_1,1.51.4)} latva . \\
\noindent \textbf{(P\_1,1.51.4)} yaksvara : yaksvaraḥ ca na sidhyati . \\
\noindent \textbf{(P\_1,1.51.4)} gīryate svayam eva , puryate svayam eva . \\
\noindent \textbf{(P\_1,1.51.4)} acaḥ kartṛyaki iti eṣaḥ svaraḥ na prāpnoti repheṇa vyavahitatvāt . \\
\noindent \textbf{(P\_1,1.51.4)} na eṣaḥ doṣaḥ . \\
\noindent \textbf{(P\_1,1.51.4)} svaravidhau vyañjanam avidyamānavat iti na asti vyavadhānam . \\
\noindent \textbf{(P\_1,1.51.4)} yaksvara . \\
\noindent \textbf{(P\_1,1.51.4)} abhyastasvara : abhyastasvaraḥ ca na sidhyati : ma hi sma te piparuḥ , ma hi sma te bibharuḥ . \\
\noindent \textbf{(P\_1,1.51.4)} abhyastānām ādiḥ udāttaḥ bhavati ajādau lasārvadhātuke iti eṣaḥ svaraḥ na prāpnoti repheṇa vyavahitatvāt . \\
\noindent \textbf{(P\_1,1.51.4)} na eṣaḥ doṣaḥ . \\
\noindent \textbf{(P\_1,1.51.4)} svaravidhau vyañjamam avidyamānavat iti na asti vyavadhānam . \\
\noindent \textbf{(P\_1,1.51.4)} abhyastasvara . \\
\noindent \textbf{(P\_1,1.51.4)} halādiśeṣa : halādiśeṣaḥ ca na sidhyati : vavṛte vavṛdhe . \\
\noindent \textbf{(P\_1,1.51.4)} abhyāsasya iti halādiśeṣaḥ na prāpnoti . \\
\noindent \textbf{(P\_1,1.51.4)} halādiśeṣa . \\
\noindent \textbf{(P\_1,1.51.4)} visarjanīya : visarjanīyasya ca pratiṣedhaḥ vaktavyaḥ : nārkuṭaḥ , nārpatyaḥ . \\
\noindent \textbf{(P\_1,1.51.4)} kharavasānayoḥ visarjanīyaḥ iti visarjanīyaḥ prāpnoti . \\
\noindent \textbf{(P\_1,1.51.4)} visarjanīya . \\
\noindent \textbf{(P\_1,1.51.4)} pratyayāvyavasthā : pratyaye vyavasthā na prakalpate : kirataḥ , girataḥ . \\
\noindent \textbf{(P\_1,1.51.4)} rephaḥ api abhaktaḥ pratyayaḥ api . \\
\noindent \textbf{(P\_1,1.51.4)} tatra vyavasthā na prakalpate . \\
\noindent \textbf{(P\_1,1.51.4)} evam tarhi pūrvāntaḥ kariṣyate . \\
\noindent \textbf{(P\_1,1.51.4)} pūrvānte rvavadhāraṇam visarjanīyapratiṣedhaḥ yaksvaraḥ ca . \\
\noindent \textbf{(P\_1,1.51.4)} yadi pūrvāntaḥ roḥ avadhāraṇam kartavyam : roḥ supi . \\
\noindent \textbf{(P\_1,1.51.4)} roḥ eva supi na anyasya rephasya : sarpiṣṣu dhanuṣṣu . \\
\noindent \textbf{(P\_1,1.51.4)} iha mā bhūt : gīrṣu pūrṣu . \\
\noindent \textbf{(P\_1,1.51.4)} parādau api sati avadhāraṇam kartavyam caturṣu iti evam artham . \\
\noindent \textbf{(P\_1,1.51.4)} visarjanīyapratiṣedhaḥ : visarjanīyasya ca pratiṣedhaḥ vaktavyaḥ : nārkuṭaḥ , nārpatyaḥ . \\
\noindent \textbf{(P\_1,1.51.4)} kharavasānayoḥ visarjanīyaḥ iti visarjanīyaḥ prāpnoti . \\
\noindent \textbf{(P\_1,1.51.4)} parādau api visarjanīyasya pratiṣedhaḥ vaktavyaḥ nārkalpiḥ iti evamartham . \\
\noindent \textbf{(P\_1,1.51.4)} kalpipadasaṅghātabhaktaḥ asau na utsahate avayavasya padāntatām vihantum iti kṛtvā visarjanīyaḥ prāpnoti . \\
\noindent \textbf{(P\_1,1.51.4)} yaksvaraḥ : yaksvaraḥ ca na sidhyati : gīryate svayam eva , puryate svayam eva . \\
\noindent \textbf{(P\_1,1.51.4)} acaḥ kartṛyaki iti eṣaḥ svaraḥ na prāpnoti . \\
\noindent \textbf{(P\_1,1.51.4)} na eṣaḥ doṣaḥ . \\
\noindent \textbf{(P\_1,1.51.4)} upadeśe iti vartate . \\
\noindent \textbf{(P\_1,1.51.4)} atha vā punaḥ astu parādiḥ . \\
\noindent \textbf{(P\_1,1.51.4)} parādau akāralopautvapukpratiṣedhaḥ caṅi upadhāhrasvatvam iṭaḥ avyavasthā abhyāsalopaḥ abhyastatādisvaraḥ dīrghatvam ca . \\
\noindent \textbf{(P\_1,1.51.4)} yadi parādiḥ akāralopaḥ pratiṣedhyaḥ : kartā hartā : ataḥ lopaḥ ārdhadhātuke iti akāralopaḥ prāpnoti . \\
\noindent \textbf{(P\_1,1.51.4)} na eṣaḥ doṣaḥ . \\
\noindent \textbf{(P\_1,1.51.4)} upadeśe iti vartate . \\
\noindent \textbf{(P\_1,1.51.4)} yadi upadeśe iti vartate dhinutaḥ , kṛṇutaḥ atra lopaḥ na prāpnoti . \\
\noindent \textbf{(P\_1,1.51.4)} na upadeśagrahaṇena prakṛtiḥ abhisambadhyate . \\
\noindent \textbf{(P\_1,1.51.4)} kim tarhi . \\
\noindent \textbf{(P\_1,1.51.4)} ārdhadhātukam abhisambadhyate . \\
\noindent \textbf{(P\_1,1.51.4)} ārdhadhātukopadeśe yat akārāntam iti . \\
\noindent \textbf{(P\_1,1.51.4)} akāralopa . \\
\noindent \textbf{(P\_1,1.51.4)} autva : autvam ca pratiṣedhyam : cakāra jahāra . \\
\noindent \textbf{(P\_1,1.51.4)} ātaḥ au ṇalaḥ iti autvam prāpnoti . \\
\noindent \textbf{(P\_1,1.51.4)} na eṣaḥ doṣaḥ . \\
\noindent \textbf{(P\_1,1.51.4)} nirdiśyamānasya ādeśāḥ bhavanti iti evam na bhaviṣyati . \\
\noindent \textbf{(P\_1,1.51.4)} yaḥ tarhi nirdiśyate tasya kasmāt na bhavati . \\
\noindent \textbf{(P\_1,1.51.4)} rephena vyavahitatvāt . \\
\noindent \textbf{(P\_1,1.51.4)} autva . \\
\noindent \textbf{(P\_1,1.51.4)} pukpratiṣedhaḥ : puk ca pratiṣedhyaḥ : kārayati hārayati . \\
\noindent \textbf{(P\_1,1.51.4)} ātām puk iti puk prāpnoti . \\
\noindent \textbf{(P\_1,1.51.4)} pukpratiṣedhaḥ . \\
\noindent \textbf{(P\_1,1.51.4)} caṅi upadhāhrasvatvam ca na sidhyati : acīkarat ajīharat . \\
\noindent \textbf{(P\_1,1.51.4)} ṇau caṅi upadhāyāḥ hrasvaḥ iti hrasvatvam na prāpnoti . \\
\noindent \textbf{(P\_1,1.51.4)} caṅi upadhāhrasvatvam . \\
\noindent \textbf{(P\_1,1.51.4)} iṭaḥ avyavasthā : iṭaḥ ca vyavasthā na prakalpate : āstaritā niparitā . \\
\noindent \textbf{(P\_1,1.51.4)} iṭ api parādiḥ rephaḥ api . \\
\noindent \textbf{(P\_1,1.51.4)} tatra vyavasthā na prakalpate . \\
\noindent \textbf{(P\_1,1.51.4)} iṭaḥ avyavasthā . \\
\noindent \textbf{(P\_1,1.51.4)} abhyāsalopaḥ : abhyāsalopaḥ ca vaktavyaḥ : vavṛte vavṛdhe . \\
\noindent \textbf{(P\_1,1.51.4)} abhyāsasya iti halādiśeṣaḥ na prāpnoti . \\
\noindent \textbf{(P\_1,1.51.4)} abhyāsalopaḥ . \\
\noindent \textbf{(P\_1,1.51.4)} abhyastasvara : abhyastasvaraḥ ca na sidhyati : ma hi sma te piparuḥ , ma hi sma te bibharuḥ . \\
\noindent \textbf{(P\_1,1.51.4)} abhyastānām ādiḥ udāttaḥ bhavati ajādau lasārvadhātuke iti eṣaḥ svaraḥ na prāpnoti . \\
\noindent \textbf{(P\_1,1.51.4)} abhyastasvara . \\
\noindent \textbf{(P\_1,1.51.4)} tādisvaraḥ : tādisvaraḥ ca na sidhyati : prakartā prakartum , prahartā prahartum . \\
\noindent \textbf{(P\_1,1.51.4)} tādau ca niti kṛti atau iti eṣaḥ svaraḥ na prāpnoti . \\
\noindent \textbf{(P\_1,1.51.4)} na eṣaḥ doṣaḥ . \\
\noindent \textbf{(P\_1,1.51.4)} uktam etat : kṛdupadeśe vā tādyartham iḍartham iti . \\
\noindent \textbf{(P\_1,1.51.4)} tādisvaraḥ . \\
\noindent \textbf{(P\_1,1.51.4)} dīrghatvam : dīrghatvam ca na sidhyati : gīḥ , pūḥ . \\
\noindent \textbf{(P\_1,1.51.4)} rephavakārāntasya dhātoḥ iti dīrghatvam na prāpnoti . \\
\noindent \textbf{(P\_1,1.52.1)} kim idam algrahaṇam antyaviśeṣaṇam āhosvit ādeśaviśeṣaṇam . \\
\noindent \textbf{(P\_1,1.52.1)} kim ca ataḥ . \\
\noindent \textbf{(P\_1,1.52.1)} yadi antyaviśeṣaṇam ādeśaḥ aviśeṣitaḥ bhavati . \\
\noindent \textbf{(P\_1,1.52.1)} tatra kaḥ doṣaḥ . \\
\noindent \textbf{(P\_1,1.52.1)} anekāl api ādeśaḥ antyasya prasajyeta . \\
\noindent \textbf{(P\_1,1.52.1)} yadi punaḥ al antyasya iti ucyeta . \\
\noindent \textbf{(P\_1,1.52.1)} tatra ayam api arthaḥ anekāl śit sarvasya iti etat na vaktavyam bhavati . \\
\noindent \textbf{(P\_1,1.52.1)} idam niyamārtham bhaviṣyati : al eva antyasya bhavati na anyaḥ iti . \\
\noindent \textbf{(P\_1,1.52.1)} evam api antyaḥ aviśeṣitaḥ bhavati . \\
\noindent \textbf{(P\_1,1.52.1)} tatra kaḥ doṣaḥ . \\
\noindent \textbf{(P\_1,1.52.1)} vākyasya api padasya api antyasya prasajyeta . \\
\noindent \textbf{(P\_1,1.52.1)} yadi khalu api eṣaḥ abhiprāyaḥ tat na kriyeta iti antyaviśeṣaṇe api sati tat na kariṣyate . \\
\noindent \textbf{(P\_1,1.52.1)} katham . \\
\noindent \textbf{(P\_1,1.52.1)} ṅit ca alaḥ antyasya iti etat niyamārtham bhaviṣyati : ṅit eva anekāl antyasya bhavati na anyaḥ iti . \\
\noindent \textbf{(P\_1,1.52.2)} kimartham punaḥ idam ucyate . \\
\noindent \textbf{(P\_1,1.52.2)} alaḥ antyasya iti sthāne vijñātasya anusaṃhāraḥ . \\
\noindent \textbf{(P\_1,1.52.2)} alaḥ antyasya iti sthāne vijñātasya anusaṃhāraḥ kriyate sthāne prasaktasya . \\
\noindent \textbf{(P\_1,1.52.2)} itarathā hi aniṣṭaprasaṅgaḥ . \\
\noindent \textbf{(P\_1,1.52.2)} itarathā hi aniṣṭaprasaṅgaḥ prasajyeta . \\
\noindent \textbf{(P\_1,1.52.2)} ṭitkinmitaḥ api antyasya syuḥ . \\
\noindent \textbf{(P\_1,1.52.2)} yadi punaḥ ayam yogaśeṣaḥ vijñāyeta . \\
\noindent \textbf{(P\_1,1.52.2)} yogaśeṣe ca . \\
\noindent \textbf{(P\_1,1.52.2)} kim . \\
\noindent \textbf{(P\_1,1.52.2)} aniṣṭam prasajyete . \\
\noindent \textbf{(P\_1,1.52.2)} ṭitkinmitaḥ api antyasya syuḥ . \\
\noindent \textbf{(P\_1,1.52.2)} tasmāt suṣṭhu ucyate : alaḥ antyasya iti sthāne vijñātasya anusaṃhāraḥ itarathā hi aniṣṭaprasaṅgaḥ iti . \\
\noindent \textbf{(P\_1,1.53)} tātaṅ antyasya kasmāt na bhavati . \\
\noindent \textbf{(P\_1,1.53)} ṅit ca alaḥ antyasya iti prāpnoti . \\
\noindent \textbf{(P\_1,1.53)} tātaṅi ṅitkaraṇasya sāvakāśatvāt vipratiṣedhāt sarvādeśaḥ . \\
\noindent \textbf{(P\_1,1.53)} tātaṅi ṅitkaraṇasya sāvakāśam . \\
\noindent \textbf{(P\_1,1.53)} kaḥ avakāśaḥ . \\
\noindent \textbf{(P\_1,1.53)} guṇavṛddhipratiṣedhārthaḥ ṅakāraḥ . \\
\noindent \textbf{(P\_1,1.53)} tātaṅi ṅitkaraṇasya sāvakāśatvāt vipratiṣedhāt sarvādeśaḥ bhaviṣyati . \\
\noindent \textbf{(P\_1,1.53)} prayojanam nāma tat vaktavyam yat niyogataḥ syāt . \\
\noindent \textbf{(P\_1,1.53)} yadi ca ayam niyogataḥ sarvādeśaḥ syāt tataḥ etat prayojanam syāt . \\
\noindent \textbf{(P\_1,1.53)} kutaḥ nu khalu etat ṅitkaraṇāt ayam sarvādeśaḥ bhaviṣyati na punaḥ antyasya syāt iti . \\
\noindent \textbf{(P\_1,1.53)} evam tarhi etat eva jñāpayati na tātaṅ antyasya sthāne bhavati iti yat etam ṅitam karoti . \\
\noindent \textbf{(P\_1,1.53)} itarathā hi loṭaḥ eruprakaraṇe eva brūyāt tihyoḥ tāt āśiṣi anyatarasyām iti . \\
\noindent \textbf{(P\_1,1.54)} alaḥ antyasya adeḥ parasya anekāl śit sarvasya iti apavādavipratiṣedhāt sarvādeśaḥ . \\
\noindent \textbf{(P\_1,1.54)} alaḥ antyasya iti utsargaḥ . \\
\noindent \textbf{(P\_1,1.54)} tasya ādeḥ parasya anekālśit sarvasya iti apavādau . \\
\noindent \textbf{(P\_1,1.54)} apavādavipratiṣedhāt tu sarvādeśaḥ bhaviṣyati . \\
\noindent \textbf{(P\_1,1.54)} ādeḥ parasya iti asya avakāśaḥ dvyantarupasargebhyaḥ apaḥ īt : dvīpam anvīpam . \\
\noindent \textbf{(P\_1,1.54)} anekālśit sarvasya iti asya avakāśaḥ asteḥ bhūḥ : bhavitā bhavitum . \\
\noindent \textbf{(P\_1,1.54)} iha ubhayam prāpnoti : ataḥ bhisaḥ ais . \\
\noindent \textbf{(P\_1,1.54)} anekālśit sarvasya iti etat bhavati vipratiṣedhena . \\
\noindent \textbf{(P\_1,1.54)} śit sarvasya iti asya avakāśaḥ idamaḥ iś : itaḥ , iha . \\
\noindent \textbf{(P\_1,1.54)} ādeḥ parasya iti asya avakāśaḥ saḥ eva . \\
\noindent \textbf{(P\_1,1.54)} iha ubhayam prāpnoti : aṣṭābhyaḥ auś . \\
\noindent \textbf{(P\_1,1.54)} śit sarvasya iti etat bhavati vipratiṣedhena . \\
\noindent \textbf{(P\_1,1.55)} śit sarvasya iti kim udāharaṇam . \\
\noindent \textbf{(P\_1,1.55)} idamaḥ iś : itaḥ , iha . \\
\noindent \textbf{(P\_1,1.55)} na etat asti prayojanam . \\
\noindent \textbf{(P\_1,1.55)} śitkaraṇāt eva atra sarvādeśaḥ bhaviṣyati . \\
\noindent \textbf{(P\_1,1.55)} idam tarhi : aṣṭābhyaḥ auś . \\
\noindent \textbf{(P\_1,1.55)} nanu ca atra api śitkaraṇāt eva sarvādeśaḥ bhaviṣyati . \\
\noindent \textbf{(P\_1,1.55)} idam tarhi : jasaḥ śī jaśśasoḥ śiḥ . \\
\noindent \textbf{(P\_1,1.55)} nanu ca atra api śitkaraṇāt eva sarvādeśaḥ bhaviṣyati . \\
\noindent \textbf{(P\_1,1.55)} asti anyat śitkaraṇe prayojanam . \\
\noindent \textbf{(P\_1,1.55)} kim . \\
\noindent \textbf{(P\_1,1.55)} viśeṣaṇārthaḥ . \\
\noindent \textbf{(P\_1,1.55)} kva viśeṣaṇārthena arthaḥ . \\
\noindent \textbf{(P\_1,1.55)} śi sarvanāmasthānam vibhāṣā ṅiśyoḥ iti . \\
\noindent \textbf{(P\_1,1.55)} śit sarvasya iti śakyam akartum . \\
\noindent \textbf{(P\_1,1.55)} katham . \\
\noindent \textbf{(P\_1,1.55)} antyasya ayam sthāne bhavan na pratyayaḥ syāt . \\
\noindent \textbf{(P\_1,1.55)} asatyām pratyayasañjñāyām itsañjñā na syāt . \\
\noindent \textbf{(P\_1,1.55)} asatyām itsañjñāyām lopaḥ na syāt . \\
\noindent \textbf{(P\_1,1.55)} asati lope anekāl . \\
\noindent \textbf{(P\_1,1.55)} yadā anekāl tadā sarvādeśaḥ . \\
\noindent \textbf{(P\_1,1.55)} yadā sarvādeśaḥ tada pratyayaḥ . \\
\noindent \textbf{(P\_1,1.55)} yadā pratyayaḥ tadā itsañjñā . \\
\noindent \textbf{(P\_1,1.55)} yadā itsañjñā tadā lopaḥ . \\
\noindent \textbf{(P\_1,1.55)} evam tarhi siddhe sati yat śit sarvasya iti āha tat jñāpayati ācāryaḥ asti eṣā paribhāṣā : na anubandhakṛtam anekāltvam bhavati iti . \\
\noindent \textbf{(P\_1,1.55)} kim etasya jñāpane prayojanam . \\
\noindent \textbf{(P\_1,1.55)} tatra asarūpasarvādeśāppratiṣedheṣu pṛthaktvanirdeśaḥ anākārāntatvāt iti uktam . \\
\noindent \textbf{(P\_1,1.55)} tat na vaktavyam bhavati iti . \\
\noindent \textbf{(P\_1,1.56.1)} vatkaraṇam kimartham . \\
\noindent \textbf{(P\_1,1.56.1)} sthānī ādeśaḥ analvidhau iti iyati ucyamāne sañjñādhikaraḥ ayam tatra sthānī ādeśasya sañjñā syāt . \\
\noindent \textbf{(P\_1,1.56.1)} tatra kaḥ doṣaḥ . \\
\noindent \textbf{(P\_1,1.56.1)} āṅaḥ yamahanaḥ ātmanepadam bhavati iti vadheḥ eva syāt . \\
\noindent \textbf{(P\_1,1.56.1)} hanteḥ na syāt . \\
\noindent \textbf{(P\_1,1.56.1)} vatkaraṇe punaḥ kriyamāṇe na doṣaḥ bhavati . \\
\noindent \textbf{(P\_1,1.56.1)} sthānikāryam ādeśe atidiśyate guruvat guruputraḥ iti yathā . \\
\noindent \textbf{(P\_1,1.56.1)} atha ādeśagrahaṇam kimartham . \\
\noindent \textbf{(P\_1,1.56.1)} sthānivat analvidhau iti iyati ucyamāne kaḥ idānīm sthānivat syāt . \\
\noindent \textbf{(P\_1,1.56.1)} yaḥ sthāne bhavati . \\
\noindent \textbf{(P\_1,1.56.1)} kaḥ ca sthāne bhavati . \\
\noindent \textbf{(P\_1,1.56.1)} ādeśaḥ . \\
\noindent \textbf{(P\_1,1.56.1)} idam tarhi prayojanam ādeśamātram sthānivat yathā syāt . \\
\noindent \textbf{(P\_1,1.56.1)} ekadeśavikṛtasya upasaṅkhyānam codayiṣyati . \\
\noindent \textbf{(P\_1,1.56.1)} tat na vaktavyam bhavati . \\
\noindent \textbf{(P\_1,1.56.1)} atha vidhigrahaṇam kimartham. sarvavibhaktyantaḥ samāsaḥ yathā vijñāyeta : alaḥ parasya vidhiḥ alvidhiḥ , alaḥ vidhiḥ alvidhiḥ , ali vidhiḥ alvidhiḥ , alā vidhiḥ alvidhiḥ iti . \\
\noindent \textbf{(P\_1,1.56.1)} na etat asti prayojanam . \\
\noindent \textbf{(P\_1,1.56.1)} prātipadikarnirdeśaḥ ayam . \\
\noindent \textbf{(P\_1,1.56.1)} prātipadikarnirdeśāḥ ca arthatantrāḥ bhavanti . \\
\noindent \textbf{(P\_1,1.56.1)} na kāṃ cit prādhānyena vibhaktim āśrayanti . \\
\noindent \textbf{(P\_1,1.56.1)} tatra prātipadikārthe nirdiṣṭe yām yām vibhaktim āśrayitum buddhiḥ upajāyate sā sā āśrayitavyā . \\
\noindent \textbf{(P\_1,1.56.1)} idam tarhi prayojanam : uttarapadalopaḥ yathā vijñāyeta : alam āśrayate alāśṛayaḥ , alāśrayaḥ vidhiḥ alvidhiḥ iti . \\
\noindent \textbf{(P\_1,1.56.1)} yatra prādhānyena al āśrīyate tatra eva pratiṣedhaḥ syāt . \\
\noindent \textbf{(P\_1,1.56.1)} yatra viśeṣaṇatvena al āśrīyate tatra pratiṣedhaḥ na syāt . \\
\noindent \textbf{(P\_1,1.56.1)} kim prayojanam . \\
\noindent \textbf{(P\_1,1.56.1)} pradīvya prasīvya iti valādilakṣaṇaḥ iṭ mā bhūt iti . \\
\noindent \textbf{(P\_1,1.56.2)} kimartham punaḥ idam ucyate . \\
\noindent \textbf{(P\_1,1.56.2)} sthānyādeśapṛthaktvāt ādeśe sthānivadanudeśaḥ guruvat guruputre iti yathā . \\
\noindent \textbf{(P\_1,1.56.2)} anyaḥ sthānī anyaḥ ādeśaḥ . \\
\noindent \textbf{(P\_1,1.56.2)} sthānyādeśapṛthaktvāt etasmāt kāraṇāt sthānikāryam ādeśe na prāpnoti . \\
\noindent \textbf{(P\_1,1.56.2)} tatra kaḥ doṣaḥ . \\
\noindent \textbf{(P\_1,1.56.2)} āṅaḥ yamahanaḥ ātmanepadam bhavati iti hanteḥ eva syāt vadheḥ na syāt . \\
\noindent \textbf{(P\_1,1.56.2)} iṣyate ca vadheḥ api syāt iti . \\
\noindent \textbf{(P\_1,1.56.2)} tat ca antareṇa yatnam na sidhyati . \\
\noindent \textbf{(P\_1,1.56.2)} tasmāt . \\
\noindent \textbf{(P\_1,1.56.2)} sthānivadanudeśaḥ . \\
\noindent \textbf{(P\_1,1.56.2)} evamartham idam ucyate . \\
\noindent \textbf{(P\_1,1.56.2)} guruvat guruputraḥ iti yathā . \\
\noindent \textbf{(P\_1,1.56.2)} tat yathā guruvat asmin guruputre vartitavyam iti gurau yat kāryam tat guruputre atidiśyate , evam iha api sthānikāryam ādeśe atidiśyate . \\
\noindent \textbf{(P\_1,1.56.2)} na etat asti prayojanam . \\
\noindent \textbf{(P\_1,1.56.2)} lokataḥ etat siddham . \\
\noindent \textbf{(P\_1,1.56.2)} tat yathā loke yaḥ yasya prasaṅge bhavati labhate asau tatkāryāṇi . \\
\noindent \textbf{(P\_1,1.56.2)} tat yathā upādhyāyasya śiṣyaḥ yājyakulāni gatvā agrāsanādīni labhate . \\
\noindent \textbf{(P\_1,1.56.2)} yadi api tāvat loke eṣaḥ dṛṣṭāntaḥ dṛṣṭāntasya api tu puruṣārambhaḥ nivartakaḥ bhavati . \\
\noindent \textbf{(P\_1,1.56.2)} asti ca iha kaḥ cit puruṣārambhaḥ . \\
\noindent \textbf{(P\_1,1.56.2)} asti iti āha . \\
\noindent \textbf{(P\_1,1.56.2)} kaḥ . \\
\noindent \textbf{(P\_1,1.56.2)} svarūpavidhiḥ . \\
\noindent \textbf{(P\_1,1.56.2)} hanteḥ ātmanepadam ucyamānam hanteḥ eva syāt vadheḥ na syāt . \\
\noindent \textbf{(P\_1,1.56.2)} evam tarhi ācāryapravṛttiḥ jñāpayati sthānivat ādeśaḥ bhavati iti yat ayam yuṣmadasmadoḥ anādeśe iti ādeśapratiṣedham śāsti . \\
\noindent \textbf{(P\_1,1.56.2)} katham kṛtvā jñāpakam . \\
\noindent \textbf{(P\_1,1.56.2)} yuṣmadasmadoḥ vibhaktau kāryam ucyamānam kaḥ prasaṅgaḥ yat ādeśe syāt . \\
\noindent \textbf{(P\_1,1.56.2)} paśyati tu ācāryaḥ sthānivat ādeśaḥ bhavati iti . \\
\noindent \textbf{(P\_1,1.56.2)} ataḥ ādeśe pratiṣedham śāsti . \\
\noindent \textbf{(P\_1,1.56.2)} idam tarhi prayojanam : analvidhau iti pratiṣedham vakṣyāmi iti , iha mā bhūt : dyauḥ , panthāḥ , saḥ iti . \\
\noindent \textbf{(P\_1,1.56.2)} etat api na asti prayojanam . \\
\noindent \textbf{(P\_1,1.56.2)} ācāryapravṛttiḥ jñāpayati alvidhau sthānivadbhāvaḥ na bhavati iti yat ayam adaḥ jagdhiḥ lyap ti kiti iti ti kiti iti eva siddhe lyabgrahaṇam karoti . \\
\noindent \textbf{(P\_1,1.56.2)} tasmāt na arthaḥ anena yogena . \\
\noindent \textbf{(P\_1,1.56.3)} ārabhyamāṇe api etasmin yoge alvidhau pratiṣedhe aviśeṣaṇe aprāptiḥ tasya adarśanāt . \\
\noindent \textbf{(P\_1,1.56.3)} alvidhau pratiṣedhe asati api viśeṣaṇe samāśrīyamaṇe asati tasmin viśeṣaṇe aprāptiḥ vidheḥ : pradīvya prasīvya . \\
\noindent \textbf{(P\_1,1.56.3)} kim kāraṇam . \\
\noindent \textbf{(P\_1,1.56.3)} tasya adarśanāt . \\
\noindent \textbf{(P\_1,1.56.3)} valādeḥ iti ucyate na ca atra valādim paśyāmaḥ . \\
\noindent \textbf{(P\_1,1.56.3)} nanu ca evamarthaḥ eva ayam yatnaḥ kriyate : anyasya kāryam ucyamānam anyasya yathā syāt iti . \\
\noindent \textbf{(P\_1,1.56.3)} satyam evamarthaḥ na tu prāpnoti . \\
\noindent \textbf{(P\_1,1.56.3)} kim kāraṇam . \\
\noindent \textbf{(P\_1,1.56.3)} sāmānyātideśe viśeṣānatideśaḥ . sāmanye hi atidiśyamāne viśeṣaḥ na atidiṣṭaḥ bhavati . \\
\noindent \textbf{(P\_1,1.56.3)} tat yathā : brahmaṇavat asmin kṣatriye vartitavyam iti sāmānyam yat brāhmaṇakāryam tat kṣatriye atidiśyate . \\
\noindent \textbf{(P\_1,1.56.3)} yat viśiṣṭam māṭhare kauṇḍinye vā na tat atidiśyate . \\
\noindent \textbf{(P\_1,1.56.3)} evam iha api sāmānyam yat pratyayakāryam tat atidiśyate yat viśiṣṭam valādeḥ iti na tat atidiśyate . \\
\noindent \textbf{(P\_1,1.56.3)} yadi evam agrahīt iti iṭaḥ īṭi iti sicaḥ lopaḥ na prāpnoti . \\
\noindent \textbf{(P\_1,1.56.3)} analvidhau iti punaḥ ucyamāne iha api pratiṣedhaḥ bhaviṣyati : pradīvya prasīvya iti . \\
\noindent \textbf{(P\_1,1.56.3)} viśiṣṭam hi eṣaḥ alam āśrayate valam nāma . \\
\noindent \textbf{(P\_1,1.56.3)} iha ca pratiṣedhaḥ na bhaviṣyati : agrahīt iti . \\
\noindent \textbf{(P\_1,1.56.3)} viśiṣṭam hi eṣaḥ analam āśrayati iṭam nāma . \\
\noindent \textbf{(P\_1,1.56.3)} yadi tarhi sāmānyam api atidiśyate viśeṣaḥ ca sati āśraye vidhiḥ iṣṭaḥ . \\
\noindent \textbf{(P\_1,1.56.3)} sati ca valāditve iṭā bhavitavyam : aruditām aruditam arudita . \\
\noindent \textbf{(P\_1,1.56.3)} kim ataḥ yat sati bhavitavyam . \\
\noindent \textbf{(P\_1,1.56.3)} pratiṣedhaḥ tu prāpnoti alvidhitvāt . \\
\noindent \textbf{(P\_1,1.56.3)} pratiṣedhaḥ tu prāpnoti . \\
\noindent \textbf{(P\_1,1.56.3)} kim kāraṇam . \\
\noindent \textbf{(P\_1,1.56.3)} alvidhitvāt . \\
\noindent \textbf{(P\_1,1.56.3)} alvidhiḥ ayam bhavati . \\
\noindent \textbf{(P\_1,1.56.3)} tatra analvidhau iti pratiṣedhaḥ prāpnoti . \\
\noindent \textbf{(P\_1,1.56.3)} na vā ānudeśikasya pratiṣedhāt itareṇa bhāvaḥ . \\
\noindent \textbf{(P\_1,1.56.3)} na vā eṣaḥ doṣaḥ . \\
\noindent \textbf{(P\_1,1.56.3)} kim kāraṇam . \\
\noindent \textbf{(P\_1,1.56.3)} ānudeśikasya pratiṣedhāt . \\
\noindent \textbf{(P\_1,1.56.3)} astu atra ānudeśikasya valāditvasya pratiṣedhaḥ . \\
\noindent \textbf{(P\_1,1.56.3)} svāśrayam atra valāditvam bhaviṣyati . \\
\noindent \textbf{(P\_1,1.56.3)} na etat vivadāmahe valādiḥ na valādiḥ iti . \\
\noindent \textbf{(P\_1,1.56.3)} kim tarhi . \\
\noindent \textbf{(P\_1,1.56.3)} sthānivadbhāvāt sārvadhātukatvam eṣitavyam . \\
\noindent \textbf{(P\_1,1.56.3)} tatra analvidhau iti pratiṣedhaḥ prāpnoti . \\
\noindent \textbf{(P\_1,1.56.4)} kim punaḥ ādeśini ali āśrīyamāṇe pratiṣedhaḥ bhavati āhosvit aviśeṣeṇa ādeśe ādeśini ca . \\
\noindent \textbf{(P\_1,1.56.4)} kaḥ ca atra viśeṣaḥ . \\
\noindent \textbf{(P\_1,1.56.4)} ādeśyalvidhipratiṣedhe kuruvadhapibām guṇavṛddhipratiṣedhaḥ . \\
\noindent \textbf{(P\_1,1.56.4)} ādeśyalvidhipratiṣedhe kuruvadhapibām guṇavṛddhipratiṣedhaḥ vaktavyaḥ . \\
\noindent \textbf{(P\_1,1.56.4)} kuru iti atra sthānivadbhāvāt aṅgasañjñā śvāśrayam ca laghūpadhatvam . \\
\noindent \textbf{(P\_1,1.56.4)} tatra laghūpadhaguṇaḥ prāpnoti . \\
\noindent \textbf{(P\_1,1.56.4)} vadhakam iti atra sthānivadbhāvāt aṅgasañjñā śvāśrayam ca adupadhatvam . \\
\noindent \textbf{(P\_1,1.56.4)} tatra vṛddhiḥ prāpnoti . \\
\noindent \textbf{(P\_1,1.56.4)} piba iti atra sthānivadbhāvāt aṅgasañjñā śvāśrayam ca laghūpadhatvam . \\
\noindent \textbf{(P\_1,1.56.4)} tatra guṇaḥ prāpnoti . \\
\noindent \textbf{(P\_1,1.56.4)} astu tarhi aviśeṣeṇa ādeśe ādeśini ca . \\
\noindent \textbf{(P\_1,1.56.4)} ādeśyādeśe iti cet suptiṅkṛdatidiṣṭeṣu upasaṅkhyānam . \\
\noindent \textbf{(P\_1,1.56.4)} ādeśyādeśe iti cet suptiṅkṛdatidiṣṭeṣu upasaṅkhyānam kartavyam . \\
\noindent \textbf{(P\_1,1.56.4)} sup : vṛkṣāya plakṣāya . \\
\noindent \textbf{(P\_1,1.56.4)} sthānivadbhāvāt supsañjñā svāśrayam ca yañāditvam . \\
\noindent \textbf{(P\_1,1.56.4)} tatra pratiṣedhaḥ prāpnoti . \\
\noindent \textbf{(P\_1,1.56.4)} sup . \\
\noindent \textbf{(P\_1,1.56.4)} tiṅ : aruditām aruditam arudita . \\
\noindent \textbf{(P\_1,1.56.4)} sthānivadbhāvāt sārvadhātukasañjñā svāśrayam ca valāditvam . \\
\noindent \textbf{(P\_1,1.56.4)} tatra pratiṣedhaḥ prāpnoti . \\
\noindent \textbf{(P\_1,1.56.4)} tiṅ . \\
\noindent \textbf{(P\_1,1.56.4)} kṛdatidiṣṭam : bhuvanam , suvanam , dhuvanam . \\
\noindent \textbf{(P\_1,1.56.4)} sthānivadbhāvāt pratyayasañjñā svāśrayam ca ajāditvam . \\
\noindent \textbf{(P\_1,1.56.4)} tatra pratiṣedhaḥ prāpnoti . \\
\noindent \textbf{(P\_1,1.56.4)} kim punaḥ atra jyāyaḥ . \\
\noindent \textbf{(P\_1,1.56.4)} ādeśini ali āśrīyamāṇe pratiṣedhaḥ iti jyāyaḥ . \\
\noindent \textbf{(P\_1,1.56.4)} kutaḥ etat . \\
\noindent \textbf{(P\_1,1.56.4)} tathā hi ayam viśiṣṭam sthānikāryam ādeśe atidiśati guruvat guruputre iti yathā . \\
\noindent \textbf{(P\_1,1.56.4)} tat yathā : guruvat guruputre vartitavyam anyatra ucchiṣṭabhojanāt pādopasaṅgrahaṇāc ca iti . \\
\noindent \textbf{(P\_1,1.56.4)} yadi ca guruputraḥ api guruḥ bhavati tat api kartavyam . \\
\noindent \textbf{(P\_1,1.56.4)} astu tarhi ādeśini ali āśrīyamāṇe pratiṣedhaḥ . \\
\noindent \textbf{(P\_1,1.56.4)} nanu ca uktam ādeśyalvidhipratiṣedhe kuruvadhapibām guṇavṛddhipratiṣedhaḥ iti . \\
\noindent \textbf{(P\_1,1.56.4)} na eṣaḥ doṣaḥ . \\
\noindent \textbf{(P\_1,1.56.4)} karotau taparakaraṇanirdeśāt siddham . \\
\noindent \textbf{(P\_1,1.56.4)} pibatiḥ adantaḥ . \\
\noindent \textbf{(P\_1,1.56.4)} vadhakam iti na ayam ṇvul . \\
\noindent \textbf{(P\_1,1.56.4)} anyaḥ ayam akaśabdaḥ kit auṇādikaḥ rucakaḥ iti yathā . \\
\noindent \textbf{(P\_1,1.56.5)} ekadeśavikṛtasya upasaṅkhyānam . \\
\noindent \textbf{(P\_1,1.56.5)} ekadeśavikṛtasya upasaṅkhyānam kartavyam . \\
\noindent \textbf{(P\_1,1.56.5)} kim prayojanam . \\
\noindent \textbf{(P\_1,1.56.5)} pacatu pacantu . \\
\noindent \textbf{(P\_1,1.56.5)} tiṅgrahaṇena grahaṇam yathā syāt . \\
\noindent \textbf{(P\_1,1.56.5)} ekadeśavikṛtasya ananyatvāt siddham . \\
\noindent \textbf{(P\_1,1.56.5)} ekadeśavikṛtam ananyavat bhavati iti tiṅgrahaṇena grahaṇam bhaviṣyati . \\
\noindent \textbf{(P\_1,1.56.5)} tat yatha : śvā karṇe vā pucche vā chinne śvā eva bhavati na aśvaḥ na gardabhaḥ iti . \\
\noindent \textbf{(P\_1,1.56.5)} anityatvavijñānam tu tasmāt upasaṅkhyanam . \\
\noindent \textbf{(P\_1,1.56.5)} anityatvavijñānam tu bhavati . \\
\noindent \textbf{(P\_1,1.56.5)} nityāḥ śabdāḥ . \\
\noindent \textbf{(P\_1,1.56.5)} nityeṣu nāma śabdeṣu kūṭasthaiḥ avicālibhiḥ varṇaiḥ bhavitavyam anapāyopajanavikāribhiḥ . \\
\noindent \textbf{(P\_1,1.56.5)} tatra saḥ eva ayam vikṛtaḥ ca etat nityeṣu na upapadyate . \\
\noindent \textbf{(P\_1,1.56.5)} tasmāt upasaṅkhyanam kartavyam . \\
\noindent \textbf{(P\_1,1.56.5)} bhāradvājīyāḥ paṭhanti : ekadeśavikṛteṣu upasaṅkhyānam . \\
\noindent \textbf{(P\_1,1.56.5)} ekadeśavikṛteṣu upasaṅkhyānam kartavyam . \\
\noindent \textbf{(P\_1,1.56.5)} kim prayojanam . \\
\noindent \textbf{(P\_1,1.56.5)} pacatu pacantu : tiṅgrahaṇena grahaṇam yathā syāt . \\
\noindent \textbf{(P\_1,1.56.5)} kim ca kāraṇam na syāt . \\
\noindent \textbf{(P\_1,1.56.5)} anādeśatvāt . \\
\noindent \textbf{(P\_1,1.56.5)} ādeśaḥ sthānivat iti ucyate , na ca ime ādeśāḥ . \\
\noindent \textbf{(P\_1,1.56.5)} rūpānyatvāt ca . \\
\noindent \textbf{(P\_1,1.56.5)} anyat khalu api rūpam pacati iti anyat pacatu iti . \\
\noindent \textbf{(P\_1,1.56.5)} ime api ādeśāḥ . \\
\noindent \textbf{(P\_1,1.56.5)} katham . \\
\noindent \textbf{(P\_1,1.56.5)} ādiśyate yaḥ saḥ ādeśaḥ . \\
\noindent \textbf{(P\_1,1.56.5)} ime ca api ādiśyante . \\
\noindent \textbf{(P\_1,1.56.5)} ādeśaḥ sthānivat iti cet na anāśritatvāt . \\
\noindent \textbf{(P\_1,1.56.5)} ādeśaḥ sthānivat iti cet tat na . \\
\noindent \textbf{(P\_1,1.56.5)} kim kāraṇam . \\
\noindent \textbf{(P\_1,1.56.5)} anāśritatvāt . \\
\noindent \textbf{(P\_1,1.56.5)} yaḥ atra ādeśaḥ na asau āśrīyate yaḥ ca āśrīyate na asau ādeśaḥ . \\
\noindent \textbf{(P\_1,1.56.5)} na etat mantavyam : samudāye āśrīyamāṇe avayavaḥ na āśrīyate iti . \\
\noindent \textbf{(P\_1,1.56.5)} abhyantaraḥ hi samudāyasya avayavaḥ . \\
\noindent \textbf{(P\_1,1.56.5)} tat yathā : vṛkṣaḥ pracalan saha avayavaiḥ pracalati . \\
\noindent \textbf{(P\_1,1.56.5)} āśrayaḥ iti cet alvidhiprasaṅgaḥ . āśrayaḥ iti cet alvidhiḥ ayam bhavati . \\
\noindent \textbf{(P\_1,1.56.5)} tatra analvidhau iti pratiṣedhaḥ prāpnoti . \\
\noindent \textbf{(P\_1,1.56.5)} na eṣaḥ doṣaḥ . \\
\noindent \textbf{(P\_1,1.56.5)} na evam sati kaḥ cit api analvidhiḥ syāt . \\
\noindent \textbf{(P\_1,1.56.5)} ucyate ca idam analvidhau iti . \\
\noindent \textbf{(P\_1,1.56.5)} tatra prakarṣagatiḥ vijñāsyate : sādhīyaḥ yaḥ alvidhiḥ iti . \\
\noindent \textbf{(P\_1,1.56.5)} kaḥ ca sādhīyaḥ alvidhiḥ . \\
\noindent \textbf{(P\_1,1.56.5)} yatra prādhānyena al āśrīyate . \\
\noindent \textbf{(P\_1,1.56.5)} yatra nāntarīyakaḥ al āśrīyate na asau alvidhiḥ . \\
\noindent \textbf{(P\_1,1.56.5)} atha vā uktam ādeśagrahaṇasya prayojanam : ādeśamātram sthānivat yathā syāt iti . \\
\noindent \textbf{(P\_1,1.56.6)} anupapannam sthānyādeśatvam nityatvāt . \\
\noindent \textbf{(P\_1,1.56.6)} sthānī ādeśaḥ iti etat nityeṣu śabdeṣu na upapadyate . \\
\noindent \textbf{(P\_1,1.56.6)} kim kāraṇam . \\
\noindent \textbf{(P\_1,1.56.6)} nityatvāt . \\
\noindent \textbf{(P\_1,1.56.6)} sthānī hi nām yaḥ bhūtvā na bhavati . \\
\noindent \textbf{(P\_1,1.56.6)} ādeśaḥ hi nāma yaḥ abhūtvā bhavati . \\
\noindent \textbf{(P\_1,1.56.6)} etat ca nityeṣu śabdeṣu na upapadyate yat sataḥ nāma vināśaḥ syāt asataḥ vā prādurbhāvaḥ iti . \\
\noindent \textbf{(P\_1,1.56.6)} siddham tu yathā laukikavaidikeṣu abhūtapūrve api sthānaśabdaprayogāt . \\
\noindent \textbf{(P\_1,1.56.6)} siddham etat . \\
\noindent \textbf{(P\_1,1.56.6)} katham . \\
\noindent \textbf{(P\_1,1.56.6)} yathā laukikeṣu vaidikeṣu ca kṛtānteṣu abhūtapūrve api sthānaśabdaḥ vartate . \\
\noindent \textbf{(P\_1,1.56.6)} loke tāvat : upādhyāyasya sthāne śiṣyaḥ iti ucyate na ca tatra upādhyāyaḥ bhūtapūrvaḥ bhavati . \\
\noindent \textbf{(P\_1,1.56.6)} vede api : somasya sthāne pūtīkatṛṇāni abhiṣuṇuyāt iti ucyate na ca tatra somaḥ bhūtapūrvaḥ bhavati . \\
\noindent \textbf{(P\_1,1.56.6)} kāryavipariṇāmāt vā siddham . \\
\noindent \textbf{(P\_1,1.56.6)} atha vā kāryavipariṇāmāt siddham etat . \\
\noindent \textbf{(P\_1,1.56.6)} kim idam kāryavipariṇāmāt iti . \\
\noindent \textbf{(P\_1,1.56.6)} kāryā buddhiḥ . \\
\noindent \textbf{(P\_1,1.56.6)} sā vipariṇamyate . \\
\noindent \textbf{(P\_1,1.56.6)} nanu ca kāryāvipariṇāmāt iti bhavitavyam . \\
\noindent \textbf{(P\_1,1.56.6)} santi ca eva hi auttarpadikāni hrasvatvāni . \\
\noindent \textbf{(P\_1,1.56.6)} api ca buddhiḥ sampratyayaḥ iti anarthāntaram . \\
\noindent \textbf{(P\_1,1.56.6)} kāryā buddhiḥ kāryaḥ sampratyayaḥ kāryasya sampratyayasya vipariṇāmaḥ kāryavipariṇāmaḥ kāryavipariṇāmāt iti . \\
\noindent \textbf{(P\_1,1.56.6)} parihārantaram eva idam matvā paṭhitam . \\
\noindent \textbf{(P\_1,1.56.6)} katham ca idam parihārāntaram syāt . \\
\noindent \textbf{(P\_1,1.56.6)} yadi bhūtapūrve sthānaśabdaḥ vartate . \\
\noindent \textbf{(P\_1,1.56.6)} bhūtapūrve ca api sthānaśabdaḥ vartate . \\
\noindent \textbf{(P\_1,1.56.6)} katham . \\
\noindent \textbf{(P\_1,1.56.6)} buddhyā . \\
\noindent \textbf{(P\_1,1.56.6)} tat yathā kaḥ cit kasmai cit upadiśati prācīnam grāmāt āmrāḥ iti . \\
\noindent \textbf{(P\_1,1.56.6)} tasya sarvatra āmrabuddhiḥ prasaktā . \\
\noindent \textbf{(P\_1,1.56.6)} tataḥ paścāt aha ye kṣīriṇaḥ avarohavantaḥ pṛthuparṇāḥ te nyagrodhāḥ iti. saḥ tatra āmrabuddhyāḥ nyagrodhabuddhim pratipadyate . \\
\noindent \textbf{(P\_1,1.56.6)} saḥ tataḥ paśyati buddhyā āmrān ca apakṛṣyamāṇān nyagrodhān ca ādhīyamānān . \\
\noindent \textbf{(P\_1,1.56.6)} nityāḥ eva ca svasmin viṣaye āmrāḥ nityāḥ ca nyagrodhāḥ . \\
\noindent \textbf{(P\_1,1.56.6)} buddhiḥ tu asya vipariṇamyate . \\
\noindent \textbf{(P\_1,1.56.6)} evam iha api astiḥ asmai aviśeṣeṇa upadiṣṭaḥ . \\
\noindent \textbf{(P\_1,1.56.6)} tasya sarvatra astibuddhiḥ prasaktā . \\
\noindent \textbf{(P\_1,1.56.6)} saḥ asteḥ bhūḥ iti astibuddhyāḥ bhavatibuddhim pratipadyate . \\
\noindent \textbf{(P\_1,1.56.6)} saḥ tataḥ paśyati buddhyā astim ca apakṛṣyamāṇam bhavatim ca ādhīyamānam . \\
\noindent \textbf{(P\_1,1.56.6)} nityaḥ eva svasmin viṣaye astiḥ nityaḥ bhavatiḥ . \\
\noindent \textbf{(P\_1,1.56.6)} buddhiḥ tu asya vipariṇamyate . \\
\noindent \textbf{(P\_1,1.56.7)} apavādaprasaṅgaḥ tu sthānivattvāt . \\
\noindent \textbf{(P\_1,1.56.7)} apavāde utsargakṛtam ca prāpnoti . \\
\noindent \textbf{(P\_1,1.56.7)} karmaṇi aṇ ātaḥ anupasarge kaḥ iti ke api aṇi kṛtam prāpnoti . \\
\noindent \textbf{(P\_1,1.56.7)} kim kāraṇam . \\
\noindent \textbf{(P\_1,1.56.7)} sthānivattvāt . \\
\noindent \textbf{(P\_1,1.56.7)} uktam vā . \\
\noindent \textbf{(P\_1,1.56.7)} kim uktam . \\
\noindent \textbf{(P\_1,1.56.7)} viṣayeṇa tu nānāliṅgakaraṇāt siddham iti . \\
\noindent \textbf{(P\_1,1.56.7)} atha vā . \\
\noindent \textbf{(P\_1,1.56.7)} siddham tu ṣaṣṭhīnirdiṣṭasya sthānivadvacanāt . \\
\noindent \textbf{(P\_1,1.56.7)} siddham etat . \\
\noindent \textbf{(P\_1,1.56.7)} katham . \\
\noindent \textbf{(P\_1,1.56.7)} ṣaṣṭhīnirdiṣṭasya ādeśaḥ sthānivat iti vaktavyam . \\
\noindent \textbf{(P\_1,1.56.7)} tat tarhi ṣaṣṭhīnirdiṣṭagrahaṇam kartavyam . \\
\noindent \textbf{(P\_1,1.56.7)} na kartavyam . \\
\noindent \textbf{(P\_1,1.56.7)} prakṛtam anuvartate . \\
\noindent \textbf{(P\_1,1.56.7)} kva prakṛtam . \\
\noindent \textbf{(P\_1,1.56.7)} ṣaṣṭhī sthāneyogā iti . \\
\noindent \textbf{(P\_1,1.56.7)} atha vā ācāryapravṛttiḥ jñāpayati na apavāde utsargakṛtam bhavati iti yat ayam śyanādīnām kān cit śitaḥ karoti : śyan , śnam , śnā , śaḥ , śnuḥ iti . \\
\noindent \textbf{(P\_1,1.56.8)} tasya doṣaḥ tayādeśe ubhayapratiṣedhaḥ . tasya etasya lakṣaṇasya doṣaḥ : tayādeśe ubhayapratiṣedhaḥ vaktavyaḥ : ubhaye devamanuṣyāḥ . \\
\noindent \textbf{(P\_1,1.56.8)} tayapaḥ grahaṇena grahaṇāt jasi vibhāṣā prāpnoti . \\
\noindent \textbf{(P\_1,1.56.8)} na eṣaḥ doṣaḥ . \\
\noindent \textbf{(P\_1,1.56.8)} ayac pratyayāntaram . \\
\noindent \textbf{(P\_1,1.56.8)} yadi pratyayāntaram ubhayī iti īkāraḥ na prāpnoti . \\
\noindent \textbf{(P\_1,1.56.8)} mā bhūt evam . \\
\noindent \textbf{(P\_1,1.56.8)} mātracaḥ iti evam bhaviṣyati . \\
\noindent \textbf{(P\_1,1.56.8)} katham . \\
\noindent \textbf{(P\_1,1.56.8)} mātrac iti na idam pratyayagrahaṇam . \\
\noindent \textbf{(P\_1,1.56.8)} kim tarhi . \\
\noindent \textbf{(P\_1,1.56.8)} pratyāhāragrahaṇam . \\
\noindent \textbf{(P\_1,1.56.8)} kva sanniviṣṭānām pratyāhāraḥ . \\
\noindent \textbf{(P\_1,1.56.8)} mātraśabdāt prabhṛti ā āyacaḥ cakārāt . \\
\noindent \textbf{(P\_1,1.56.8)} yadi pratyāhāragrahaṇam kati tiṣṭhanti atra api prāpnoti . \\
\noindent \textbf{(P\_1,1.56.8)} ataḥ iti vartate . \\
\noindent \textbf{(P\_1,1.56.8)} evam api tailamātrā ghrtamātrā iti atra api prāpnoti . \\
\noindent \textbf{(P\_1,1.56.8)} sadṛśasya api asanniviṣṭasya na bhaviṣyati pratyāhāreṇa grahaṇam . \\
\noindent \textbf{(P\_1,1.56.8)} jātyakhyāyām vacanātideśe sthānivadbhāvapratiṣedhaḥ . \\
\noindent \textbf{(P\_1,1.56.8)} jātyakhyāyām vacanātideśe sthānivadbhāvasya pratiṣedhaḥ vaktavyaḥ . \\
\noindent \textbf{(P\_1,1.56.8)} vrīhibhyaḥ āgataḥ iti atra gheḥ ṅiti it guṇaḥ prāpnoti . \\
\noindent \textbf{(P\_1,1.56.8)} na eṣaḥ doṣaḥ . \\
\noindent \textbf{(P\_1,1.56.8)} uktam etat : arthātideśāt siddham iti . \\
\noindent \textbf{(P\_1,1.56.8)} ṅyābgrahaṇe adīrghaḥ . ṅyābgrahaṇe adīrghaḥ ādeśaḥ na sthānivat iti vaktavyam . \\
\noindent \textbf{(P\_1,1.56.8)} kim prayojanam . \\
\noindent \textbf{(P\_1,1.56.8)} niṣkauśāmbiḥ , atikhaṭvaḥ . \\
\noindent \textbf{(P\_1,1.56.8)} ṅyābgrahaṇena grahaṇāt sulopaḥ mā bhūt iti . \\
\noindent \textbf{(P\_1,1.56.8)} nanu ca dīrghāt iti ucyate . \\
\noindent \textbf{(P\_1,1.56.8)} tat na vaktavyam bhavati . \\
\noindent \textbf{(P\_1,1.56.8)} kim punaḥ atra jyāyaḥ . \\
\noindent \textbf{(P\_1,1.56.8)} sthānivatpratiṣedhaḥ eva jyāyān . \\
\noindent \textbf{(P\_1,1.56.8)} idam api siddham bhavati : atikhaṭvāya atimālāya . \\
\noindent \textbf{(P\_1,1.56.8)} yāṭ āpaḥ iti yāṭ na bhavati . \\
\noindent \textbf{(P\_1,1.56.8)} atha idānīm asati api sthānivadbhāve dīrghatve kṛte pit ca asau bhūtapūrvaḥ iti kṛtvā yāṭ āpaḥ iti yāṭ kasmāt na bhavati . \\
\noindent \textbf{(P\_1,1.56.8)} lakṣaṇapratipadoktayoḥ pratipadoktasya eva iti . \\
\noindent \textbf{(P\_1,1.56.8)} nanu ca idānīm sati api sthānivadbhāve etayā paribhāṣayā śakyam iha upasthātum . \\
\noindent \textbf{(P\_1,1.56.8)} na iti āha . \\
\noindent \textbf{(P\_1,1.56.8)} na hi idānīm kva cit api sthānivadbhāvaḥ syāt . \\
\noindent \textbf{(P\_1,1.56.8)} tat tarhi vaktavyam . \\
\noindent \textbf{(P\_1,1.56.8)} na vaktavyam. praśliṣṭanirdeśāt siddham . \\
\noindent \textbf{(P\_1,1.56.8)} praśliṣṭanirdeśaḥ ayam : ṅī* ī* īkārāntāt ā* āp ākārāntāt iti . \\
\noindent \textbf{(P\_1,1.56.8)} āhibhuvoḥ īṭpratiṣedhaḥ . \\
\noindent \textbf{(P\_1,1.56.8)} āhibhuvoḥ īṭpratiṣedhaḥ vaktavyaḥ : āttha abhūt . \\
\noindent \textbf{(P\_1,1.56.8)} astibrūgrahaṇena grahaṇāt īṭ prāpnoti . \\
\noindent \textbf{(P\_1,1.56.8)} āheḥ tāvat na vaktavyaḥ . \\
\noindent \textbf{(P\_1,1.56.8)} ācāryapravṛttiḥ jñāpayati na āheḥ īṭ bhavati iti yat ayam āhaḥ thaḥ iti jhalādiprakaraṇe thatvam śāsti . \\
\noindent \textbf{(P\_1,1.56.8)} na etat asti prayojanam. asti hi anyat etasya vacane prayojanam . \\
\noindent \textbf{(P\_1,1.56.8)} kim . \\
\noindent \textbf{(P\_1,1.56.8)} bhūtapūrvagatiḥ yathā vijñāyeta : jhalādiḥ yaḥ bhūtapūrvaḥ iti . \\
\noindent \textbf{(P\_1,1.56.8)} yadi evam thavacanam anarthakam syāt . \\
\noindent \textbf{(P\_1,1.56.8)} āthim eva ayam uccārayet : bruvaḥ pañcānām āditaḥ āthaḥ bruvaḥ iti . \\
\noindent \textbf{(P\_1,1.56.8)} bhavateḥ ca api na vaktavyaḥ . \\
\noindent \textbf{(P\_1,1.56.8)} astisicaḥ apṛkte iti dvisakārakaḥ nirdeśaḥ : asteḥ sakārāntāt iti . \\
\noindent \textbf{(P\_1,1.56.8)} vadhyādeśe vṛddhitatvapratiṣedhaḥ . \\
\noindent \textbf{(P\_1,1.56.8)} vadhyādeśe vṛddhitatvapratiṣedhaḥ vaktavyaḥ : vadhakam puṣkaram iti . \\
\noindent \textbf{(P\_1,1.56.8)} sthānivadbhāvāt vṛddhitatve prāpnutaḥ . \\
\noindent \textbf{(P\_1,1.56.8)} na eṣaḥ doṣaḥ . \\
\noindent \textbf{(P\_1,1.56.8)} uktam etat : na ayam ṇvul . \\
\noindent \textbf{(P\_1,1.56.8)} anyaḥ ayam akaśabdaḥ kit auṇādikaḥ rucakaḥ iti yathā . \\
\noindent \textbf{(P\_1,1.56.8)} iḍvidhiḥ ca . \\
\noindent \textbf{(P\_1,1.56.8)} iḍvidheyaḥ : āvadhiṣīṣṭa . \\
\noindent \textbf{(P\_1,1.56.8)} ekācaḥ upadeśe anudāttāt iti pratiṣedhaḥ prāpnoti . \\
\noindent \textbf{(P\_1,1.56.8)} na eṣaḥ doṣaḥ . \\
\noindent \textbf{(P\_1,1.56.8)} ādyudāttanipātanam kariṣyate . \\
\noindent \textbf{(P\_1,1.56.8)} sa nipātanasvaraḥ prakṛtisvarasya bādhakaḥ bhaviṣyati . \\
\noindent \textbf{(P\_1,1.56.8)} evam api upadeśivadbhāvaḥ vaktavyaḥ . \\
\noindent \textbf{(P\_1,1.56.8)} yathā eva hi nipātanasvaraḥ prakṛtisvaram bādhate evam pratyayasvaram api bādheta : āvadhiṣīṣṭa iti . \\
\noindent \textbf{(P\_1,1.56.8)} na eṣaḥ doṣaḥ . \\
\noindent \textbf{(P\_1,1.56.8)} ārdhadhātukīyāḥ sāmānyena bhavanti anavasthiteṣu pratyayeṣu . \\
\noindent \textbf{(P\_1,1.56.8)} tatra ārdhadhātukasāmānye vadhibhāve kṛte sati śiṣṭatvāt pratyayasvaraḥ bhaviṣyati . \\
\noindent \textbf{(P\_1,1.56.8)} ākārāntāt nukṣukpratiṣedhaḥ . \\
\noindent \textbf{(P\_1,1.56.8)} ākārāntāt nukṣukoḥ pratiṣedhaḥ vaktavyaḥ : vilāpayati bhāpayate . \\
\noindent \textbf{(P\_1,1.56.8)} lībhīgrahaṇena grahaṇāt nukṣukau prāpnutaḥ . \\
\noindent \textbf{(P\_1,1.56.8)} lībhiyoḥ praśliṣṭanirdeśāt siddham . \\
\noindent \textbf{(P\_1,1.56.8)} lībhiyoḥ praśliṣṭanirdeśaḥ ayam : lī* ī* īkārāntasya bhī* ī* īkārāntasya ca iti . \\
\noindent \textbf{(P\_1,1.56.8)} loḍādeśe śābhāvajabhāvadhitvahilopaittvapratiṣedhaḥ . \\
\noindent \textbf{(P\_1,1.56.8)} loḍādeśe eṣām pratiṣedhaḥ vaktavyaḥ : śiṣṭāt , hatāt , bhintāt , kurutāt , stāt . \\
\noindent \textbf{(P\_1,1.56.8)} loḍādeśe kṛte śābhāvaḥ jabhāvaḥ dhitvam hilopaḥ ettvam iti ete vidhayaḥ prāpnuvanti . \\
\noindent \textbf{(P\_1,1.56.8)} na eṣaḥ doṣaḥ . \\
\noindent \textbf{(P\_1,1.56.8)} idam iha sampradhāryam : loḍādeśaḥ kriyatām ete vidhayaḥ iti kim atra kartavyam . \\
\noindent \textbf{(P\_1,1.56.8)} paratvāt loḍādeśaḥ . \\
\noindent \textbf{(P\_1,1.56.8)} atha idānīm loḍādeśe kṛte punaḥprasaṅgavijñānāt kasmāt ete vidhayaḥ na bhavanti . \\
\noindent \textbf{(P\_1,1.56.8)} sakṛdgatau vipratiṣedhe yat bādhitam tat bādhitam eva iti kṛtvā . \\
\noindent \textbf{(P\_1,1.56.8)} trayādeśe srantapratiṣedhaḥ . \\
\noindent \textbf{(P\_1,1.56.8)} trayādeśe srantasya pratiṣedhaḥ vaktavyaḥ : tisṛṇām . \\
\noindent \textbf{(P\_1,1.56.8)} tisṛbhāve kṛte treḥ trayaḥ iti trayādeśaḥ prāpnoti . \\
\noindent \textbf{(P\_1,1.56.8)} na eṣaḥ doṣaḥ . \\
\noindent \textbf{(P\_1,1.56.8)} idam iha sampradhāryam : tisṛbhāvaḥ kriyatām trayādeśaḥ iti kim atra kartavyam . \\
\noindent \textbf{(P\_1,1.56.8)} paratvāt tisṛbhāvaḥ . \\
\noindent \textbf{(P\_1,1.56.8)} atha idānīm tisṛbhāve kṛte punaḥprasaṅgavijñānāt trayādeśaḥ kasmāt na bhavati . \\
\noindent \textbf{(P\_1,1.56.8)} sakṛdgatau vipratiṣedhe yat bādhitam tat bādhitam eva iti . \\
\noindent \textbf{(P\_1,1.56.8)} āmvidhau ca . \\
\noindent \textbf{(P\_1,1.56.8)} āmvidhau ca srantasya pratiṣedhaḥ vaktavyaḥ : catasraḥ tiṣṭhanti . \\
\noindent \textbf{(P\_1,1.56.8)} catasṛbhāve kṛte caturanaḍuhoḥ ām udāttaḥ iti ām prāpnoti . \\
\noindent \textbf{(P\_1,1.56.8)} na eṣaḥ doṣaḥ . \\
\noindent \textbf{(P\_1,1.56.8)} idam iha sampradhāryam : catasṛbhāvaḥ kriyatām caturanaḍuhoḥ ām udāttaḥ iti ām iti kim atra kartavyam . \\
\noindent \textbf{(P\_1,1.56.8)} paratvāt catasṛbhāvaḥ . \\
\noindent \textbf{(P\_1,1.56.8)} atha idānīm catasṛbhāve kṛte punaḥprasaṅgavijñānāt ām kasmāt na bhavati . \\
\noindent \textbf{(P\_1,1.56.8)} sakṛdgatau vipratiṣedhe yat bādhitam tat bādhitam eva iti . \\
\noindent \textbf{(P\_1,1.56.8)} svare vasvādeśe . \\
\noindent \textbf{(P\_1,1.56.8)} svare vasvādeśe pratiṣedhaḥ vaktavyaḥ : viduṣaḥ paśya . \\
\noindent \textbf{(P\_1,1.56.8)} śatuḥ anumaḥ nadyajādī antodāttāt iti eṣaḥ svaraḥ prāpnoti . \\
\noindent \textbf{(P\_1,1.56.8)} na eṣaḥ doṣaḥ . \\
\noindent \textbf{(P\_1,1.56.8)} anumaḥ iti pratiṣedhaḥ bhaviṣyati . \\
\noindent \textbf{(P\_1,1.56.8)} anumaḥ iti ucyate na ca atra numam paśyāmaḥ . \\
\noindent \textbf{(P\_1,1.56.8)} anumaḥ iti na idam āgamagrahaṇam . \\
\noindent \textbf{(P\_1,1.56.8)} kim tarhi . \\
\noindent \textbf{(P\_1,1.56.8)} pratyāhāragrahaṇam . \\
\noindent \textbf{(P\_1,1.56.8)} kva sanniviṣṭānām pratyāhāraḥ . \\
\noindent \textbf{(P\_1,1.56.8)} ukārāt prabhṛti ā numaḥ makārāt . \\
\noindent \textbf{(P\_1,1.56.8)} yadi pratyāhāragrahaṇam lunata punata atra api prāpnoti . \\
\noindent \textbf{(P\_1,1.56.8)} anumgrahaṇena na śatrantam viśeṣyate . \\
\noindent \textbf{(P\_1,1.56.8)} kim tarhi . \\
\noindent \textbf{(P\_1,1.56.8)} śatā eva viśeṣyate : śatā yaḥ anumkaḥ iti . \\
\noindent \textbf{(P\_1,1.56.8)} avaśyam ca etat evam vijñeyam . \\
\noindent \textbf{(P\_1,1.56.8)} āgamagrahaṇe hi sati iha prasajyeta : muñcatā muñcataḥ iti . \\
\noindent \textbf{(P\_1,1.56.8)} goḥ pūrvaṇittvātvasvareṣu . \\
\noindent \textbf{(P\_1,1.56.8)} goḥ pūrvaṇittvātvasvareṣu pratiṣedhaḥ vaktavyaḥ : citragvagram , śabalagvagram . \\
\noindent \textbf{(P\_1,1.56.8)} sarvatra vibhāṣā goḥ iti vibhāṣā pūrvatvam prāpnoti . \\
\noindent \textbf{(P\_1,1.56.8)} na eṣaḥ doṣaḥ . \\
\noindent \textbf{(P\_1,1.56.8)} eṅaḥ iti vartate . \\
\noindent \textbf{(P\_1,1.56.8)} tatra analvidhau iti pratiṣedhaḥ bhaviṣyati . \\
\noindent \textbf{(P\_1,1.56.8)} evam api he citrago agram atra prāpnoti . \\
\noindent \textbf{(P\_1,1.56.8)} ṇittvam : citraguḥ , citragū citragavaḥ . \\
\noindent \textbf{(P\_1,1.56.8)} goto ṇit iti ṇittvam prāpnoti . \\
\noindent \textbf{(P\_1,1.56.8)} ātvam : citragum paśya śabalagum paśya . \\
\noindent \textbf{(P\_1,1.56.8)} ā otaḥ iti ātvam prāpnoti . \\
\noindent \textbf{(P\_1,1.56.8)} na eṣaḥ doṣaḥ . \\
\noindent \textbf{(P\_1,1.56.8)} taparakaraṇāt siddham . \\
\noindent \textbf{(P\_1,1.56.8)} taparakaraṇasāmārthyāt ṇittvātve na bhaviṣyataḥ . \\
\noindent \textbf{(P\_1,1.56.8)} svara : bahugumān . \\
\noindent \textbf{(P\_1,1.56.8)} na gośvansāvavarṇa iti pratiṣedhaḥ prāpnoti . \\
\noindent \textbf{(P\_1,1.56.8)} karotipibyoḥ pratiṣedhaḥ . \\
\noindent \textbf{(P\_1,1.56.8)} karotipibyoḥ pratiṣedhaḥ vaktavyaḥ : kuru piba iti . \\
\noindent \textbf{(P\_1,1.56.8)} sthānivadbhāvāt laghūpadhaguṇaḥ prāpnoti . \\
\noindent \textbf{(P\_1,1.56.8)} uktam vā . \\
\noindent \textbf{(P\_1,1.56.8)} kim uktam . \\
\noindent \textbf{(P\_1,1.56.8)} karotau taparakaraṇanirdeśāt siddham , pibatiḥ adantaḥ iti . \\
\noindent \textbf{(P\_1,1.57.1)} acaḥ iti kimartham . \\
\noindent \textbf{(P\_1,1.57.1)} praśnaḥ , dyūtvā , ākrāṣṭām āgatya . \\
\noindent \textbf{(P\_1,1.57.1)} praśnaḥ , viśnaḥ iti atra chakārasya śakāraḥ paranimittakaḥ . \\
\noindent \textbf{(P\_1,1.57.1)} tasya sthānivadbhāvāt che ca iti tuk prāpnoti . \\
\noindent \textbf{(P\_1,1.57.1)} acaḥ iti vacanāt na bhavati . \\
\noindent \textbf{(P\_1,1.57.1)} na etat asti prayojanam . \\
\noindent \textbf{(P\_1,1.57.1)} kriyamāṇe api vai ajgrahaṇe avaśyam atra tugabhāve yatnaḥ kartavyaḥ . \\
\noindent \textbf{(P\_1,1.57.1)} antaraṅgatvāt hi tuk prāpnoti . \\
\noindent \textbf{(P\_1,1.57.1)} idam tarhi : dyūtvā . \\
\noindent \textbf{(P\_1,1.57.1)} vakārasya ūṭh paranimittakaḥ . \\
\noindent \textbf{(P\_1,1.57.1)} tasya sthānivadbhāvāt aci iti yaṇādeśaḥ na prāpnoti . \\
\noindent \textbf{(P\_1,1.57.1)} acaḥ iti vacanāt bhavati . \\
\noindent \textbf{(P\_1,1.57.1)} etat api na asti prayojanam . \\
\noindent \textbf{(P\_1,1.57.1)} svāśrayam atra actvam bhaviṣyati . \\
\noindent \textbf{(P\_1,1.57.1)} atha vā yaḥ atra ādeśaḥ na asau āśrīyate yaḥ ca āśrīyate na asau ādeśaḥ . \\
\noindent \textbf{(P\_1,1.57.1)} idam tarhi prayojanam : ākrāṣṭām . \\
\noindent \textbf{(P\_1,1.57.1)} sicaḥ lopaḥ paranimittakaḥ . \\
\noindent \textbf{(P\_1,1.57.1)} tasya sthānivadbhāvāt ṣaḍhoḥ kaḥ si iti katvam prāpnoti . \\
\noindent \textbf{(P\_1,1.57.1)} acaḥ iti vacanāt bhavati . \\
\noindent \textbf{(P\_1,1.57.1)} etat api na asti prayojanam .vakṣyati etat : pūrvatrāsiddhe na sthānivat iti . \\
\noindent \textbf{(P\_1,1.57.1)} idam tarhi prayojanam : āgatya , abhigatya . \\
\noindent \textbf{(P\_1,1.57.1)} anunāsikalopaḥ paranimittakaḥ . \\
\noindent \textbf{(P\_1,1.57.1)} tasya sthānivadbhāvāt hrasvasya iti tuk na prāpnoti . \\
\noindent \textbf{(P\_1,1.57.1)} acaḥ iti vacanāt bhavati . \\
\noindent \textbf{(P\_1,1.57.1)} atha parasmin iti kimartham . \\
\noindent \textbf{(P\_1,1.57.1)} yuvajāniḥ , dvipadikā , vaiyāghrapadyaḥ , ādīdhye . \\
\noindent \textbf{(P\_1,1.57.1)} yuvajāniḥ , vadhūjāniḥ iti : jāyāyāḥ niṅ na paranimittakaḥ . \\
\noindent \textbf{(P\_1,1.57.1)} tasya sthānivadbhāvāt vali iti yalopaḥ na prāpnoti . \\
\noindent \textbf{(P\_1,1.57.1)} parasmin iti vacanāt bhavati . \\
\noindent \textbf{(P\_1,1.57.1)} na etat asti prayojanam . \\
\noindent \textbf{(P\_1,1.57.1)} svāśrayam atra valtam bhaviṣyati . \\
\noindent \textbf{(P\_1,1.57.1)} atha vā yaḥ atra ādeśaḥ na asau āśrīyate yaḥ ca āśrīyate na asau ādeśaḥ . \\
\noindent \textbf{(P\_1,1.57.1)} idam tarhi prayojanam : dvipadikā tripadikā . \\
\noindent \textbf{(P\_1,1.57.1)} pādasya lopaḥ na paranimittakaḥ . \\
\noindent \textbf{(P\_1,1.57.1)} tasya sthānivadbhāvāt padbhāvaḥ na prāpnoti . \\
\noindent \textbf{(P\_1,1.57.1)} parasmin iti vacanāt bhavati . \\
\noindent \textbf{(P\_1,1.57.1)} etat api na asti prayojanam . \\
\noindent \textbf{(P\_1,1.57.1)} punarlopavacanasāmarthyāt sthānivadbhāvaḥ na bhaviṣyati . \\
\noindent \textbf{(P\_1,1.57.1)} idam tarhi prayojanam : vaiyāghrapadyaḥ . \\
\noindent \textbf{(P\_1,1.57.1)} nanu ca atra api punarvacanasāmarthyāt eva na bhaviṣyati . \\
\noindent \textbf{(P\_1,1.57.1)} asti hi anyat punarlopavacane prayojanam . \\
\noindent \textbf{(P\_1,1.57.1)} kim . \\
\noindent \textbf{(P\_1,1.57.1)} yatra bhasañjñā na : vyāghrapāt , śyenapāt iti . \\
\noindent \textbf{(P\_1,1.57.1)} idam ca api udāharaṇam : ādīdhye , āvevye . \\
\noindent \textbf{(P\_1,1.57.1)} ikārasya ekāraḥ na paranimittakaḥ . \\
\noindent \textbf{(P\_1,1.57.1)} tasya sthānivadbhāvāt yīvarṇayoḥ dīdhīvevyoḥ iti lopaḥ prāpnoti . \\
\noindent \textbf{(P\_1,1.57.1)} parasmin iti vacanāt bhavati . \\
\noindent \textbf{(P\_1,1.57.1)} atha pūrvavidhau iti kim artham . \\
\noindent \textbf{(P\_1,1.57.1)} he gauḥ , bābhravīyāḥ , naidheyaḥ . \\
\noindent \textbf{(P\_1,1.57.1)} he gauḥ iti aukāraḥ paranimittakaḥ . \\
\noindent \textbf{(P\_1,1.57.1)} tasya sthānivadbhāvāt eṅhrasvāt sambuddheḥ iti lopaḥ prāpnoti . \\
\noindent \textbf{(P\_1,1.57.1)} pūrvavidhau iti vacanāt na bhavati . \\
\noindent \textbf{(P\_1,1.57.1)} na etat asti prayojanam . \\
\noindent \textbf{(P\_1,1.57.1)} ācāryapravṛttiḥ jñāpayati na sambuddhilope sthānivadbhāvaḥ bhavati iti yat ayam eṅhrasvāt sambuddheḥ iti eṅgrahaṇam karoti . \\
\noindent \textbf{(P\_1,1.57.1)} na etat asti jñāpakam . \\
\noindent \textbf{(P\_1,1.57.1)} gortham etat syāt . \\
\noindent \textbf{(P\_1,1.57.1)} yat tarhi pratyāhāragrahaṇam karoti . \\
\noindent \textbf{(P\_1,1.57.1)} itarathā hi ohrasvāt iti eva brūyāt . \\
\noindent \textbf{(P\_1,1.57.1)} idam tarhi prayojanam : bābhravīyāḥ , mādhavīyāḥ . \\
\noindent \textbf{(P\_1,1.57.1)} vāntādeśaḥ paranimittakaḥ . \\
\noindent \textbf{(P\_1,1.57.1)} tasya sthānivadbhāvāt halaḥ taddhitasya iti yalopaḥ na prāpnoti . \\
\noindent \textbf{(P\_1,1.57.1)} pūrvavidhau iti vacanāt na bhavati . \\
\noindent \textbf{(P\_1,1.57.1)} etat api na asti prayojanam . \\
\noindent \textbf{(P\_1,1.57.1)} svāśrayam atra haltvam bhaviṣyati . \\
\noindent \textbf{(P\_1,1.57.1)} atha vā yaḥ atra ādeśaḥ na asau āśrīyate yaḥ ca āśrīyate na asau ādeśaḥ . \\
\noindent \textbf{(P\_1,1.57.1)} idam tarhi prayojanam : naidheyaḥ . \\
\noindent \textbf{(P\_1,1.57.1)} ākāralopaḥ paranimittakaḥ . \\
\noindent \textbf{(P\_1,1.57.1)} tasya sthānivadbhāvāt dvyajlakṣaṇaḥ ḍhak na prāpnoti . \\
\noindent \textbf{(P\_1,1.57.1)} pūrvavidhau iti vacanāt na bhavati . \\
\noindent \textbf{(P\_1,1.57.1)} atha vidhigrahaṇam kimartham . \\
\noindent \textbf{(P\_1,1.57.1)} sarvavibhaktyantaḥ samāsaḥ yathā vijñāyeta : pūrvasya vidhiḥ pūrvavidhiḥ , pūrvasmāt vidhiḥ pūrvavidhiḥ iti . \\
\noindent \textbf{(P\_1,1.57.1)} kāni punaḥ pūrvasmāt vidhau sthānivadbhāvasya prayojanāni . \\
\noindent \textbf{(P\_1,1.57.1)} bebhiditā , māthitikaḥ , apīpacan . \\
\noindent \textbf{(P\_1,1.57.1)} bebhiditā , cecchiditā iti akāralope kṛte ekājlakṣaṇaḥ iṭpratiṣedhaḥ prāpnoti . \\
\noindent \textbf{(P\_1,1.57.1)} sthānivadbhāvāt na bhavati . \\
\noindent \textbf{(P\_1,1.57.1)} māthitikaḥ iti akāralope kṛte tāntāt kaḥ iti kādeśaḥ prāpnoti . \\
\noindent \textbf{(P\_1,1.57.1)} sthānivadbhāvāt na bhavati . \\
\noindent \textbf{(P\_1,1.57.1)} apīpacan iti ekādeśe kṛte abhyastāt jheḥ jus bhavati iti jusbhāvaḥ prāpnoti . \\
\noindent \textbf{(P\_1,1.57.1)} sthānivadbhāvāt na bhavati . \\
\noindent \textbf{(P\_1,1.57.1)} na etāni santi prayojanāni . \\
\noindent \textbf{(P\_1,1.57.1)} kutaḥ . \\
\noindent \textbf{(P\_1,1.57.1)} prātipadikarnirdeśaḥ ayam . \\
\noindent \textbf{(P\_1,1.57.1)} prātipadikarnirdeśāḥ ca arthatantrāḥ bhavanti . \\
\noindent \textbf{(P\_1,1.57.1)} na kāṃ cit prādhānyena vibhaktim āśrayanti . \\
\noindent \textbf{(P\_1,1.57.1)} tatra prātipadikārthe nirdiṣṭe yām yām vibhaktim āśrayitum buddhiḥ upajāyate sā sā āśrayitavyā . \\
\noindent \textbf{(P\_1,1.57.1)} idam tarhi prayojanam : vidhimātre sthānivat yathā syāt anāśrīyamāṇāyām api prakṛtau : vāyvoḥ , adhvaryvoḥ . \\
\noindent \textbf{(P\_1,1.57.1)} lopaḥ vyoḥ vali iti yalapaḥ mā bhūt iti . \\
\noindent \textbf{(P\_1,1.57.1)} asti prayojanam etat . \\
\noindent \textbf{(P\_1,1.57.1)} kim tarhi iti . \\
\noindent \textbf{(P\_1,1.57.1)} aparavidhau iti tu vaktavyam . \\
\noindent \textbf{(P\_1,1.57.1)} kim prayojanam . \\
\noindent \textbf{(P\_1,1.57.1)} svavidhau api sthānivadbhāvaḥ yathā syāt . \\
\noindent \textbf{(P\_1,1.57.1)} kāni punaḥ svavidhau sthānivadbhāvasya prayojanāni . \\
\noindent \textbf{(P\_1,1.57.1)} āyan , āsan , dhinvanti kṛṇvanti dadhi atra , madhu atra cakratuḥ , cakruḥ . \\
\noindent \textbf{(P\_1,1.57.1)} iha tāvat : āyan , āsan iti iṇastyoḥ yaṇlopayoḥ kṛtayoḥ anajāditvāt āṭ ajādīnām iti āṭ na prāpnoti . \\
\noindent \textbf{(P\_1,1.57.1)} sthānivadbhāvāt bhavati . \\
\noindent \textbf{(P\_1,1.57.1)} dhinvanti kṛṇvanti iti yaṇādeśe kṛte valādilakṣaṇaḥ iṭ prāpnoti . \\
\noindent \textbf{(P\_1,1.57.1)} sthānivadbhāvāt na bhavati . \\
\noindent \textbf{(P\_1,1.57.1)} dadhi atra madhu atra iti yaṇādeśe kṛte saṃyogāntalopaḥ prāpnoti . \\
\noindent \textbf{(P\_1,1.57.1)} sthānivadbhāvāt na bhavati . \\
\noindent \textbf{(P\_1,1.57.1)} cakratuḥ , cakruḥ iti yaṇādeśe kṛte anactvāt dvirvacanam na prāpnoti . \\
\noindent \textbf{(P\_1,1.57.1)} sthānivadbhāvāt bhavati . \\
\noindent \textbf{(P\_1,1.57.1)} yadi tarhi svavidhau api sthānivadbhāvaḥ bhavati dvābhyām , deyam , lavanam atra api prāpnoti . \\
\noindent \textbf{(P\_1,1.57.1)} dvābhyām iti atra atvasya sthānivadbhāvāt dīrghatvam na prāpnoti . \\
\noindent \textbf{(P\_1,1.57.1)} deyam iti īttvasya sthānivadbhāvāt guṇaḥ na prāpnoti . \\
\noindent \textbf{(P\_1,1.57.1)} lavanam iti guṇasya sthānivadbhāvāt avādeśaḥ na prāpnoti . \\
\noindent \textbf{(P\_1,1.57.1)} na eṣaḥ doṣaḥ . \\
\noindent \textbf{(P\_1,1.57.1)} svāśrayāḥ atra ete vidhayaḥ bhaviṣyanti . \\
\noindent \textbf{(P\_1,1.57.1)} tat tarhi vaktavyam aparavidhau iti . \\
\noindent \textbf{(P\_1,1.57.1)} na vaktavyam . \\
\noindent \textbf{(P\_1,1.57.1)} pūrvavidhau iti eva siddham . \\
\noindent \textbf{(P\_1,1.57.1)} katham . \\
\noindent \textbf{(P\_1,1.57.1)} na pūrvgrahaṇena ādeśaḥ abhisambadhyate : ajādeśaḥ paranimittakaḥ pūrvasya vidhim prati sthānivat bhavati . \\
\noindent \textbf{(P\_1,1.57.1)} kutaḥ pūrvasya . \\
\noindent \textbf{(P\_1,1.57.1)} ādeśāt iti . \\
\noindent \textbf{(P\_1,1.57.1)} kim tarhi . \\
\noindent \textbf{(P\_1,1.57.1)} nimittam abhisambadhyate : ajādeśaḥ paranimittakaḥ pūrvasya vidhim prati sthānivat bhavati . \\
\noindent \textbf{(P\_1,1.57.1)} kutaḥ pūrvasya . \\
\noindent \textbf{(P\_1,1.57.1)} nimittāt iti . \\
\noindent \textbf{(P\_1,1.57.1)} atha nimitte abhisambadhyamāne yat tat asya yogasya mūrdhābhiṣiktam udāharaṇam tat api saṅgṛhītam bhavati . \\
\noindent \textbf{(P\_1,1.57.1)} kim punaḥ tat . \\
\noindent \textbf{(P\_1,1.57.1)} paṭvyā mṛdvyā iti . \\
\noindent \textbf{(P\_1,1.57.1)} bāḍham saṅgṛhītam . \\
\noindent \textbf{(P\_1,1.57.1)} nanu ca īkārayaṇā vyavahitatvāt na asau nimittāt pūrvaḥ bhavati . \\
\noindent \textbf{(P\_1,1.57.1)} vyavahite api pūrvaśabdaḥ vartate . \\
\noindent \textbf{(P\_1,1.57.1)} tat yathā : pūrvam mathurāyāḥ pāṭaliputram iti . \\
\noindent \textbf{(P\_1,1.57.1)} atha vā ādeśaḥ eva abhisambadhyate . \\
\noindent \textbf{(P\_1,1.57.1)} katham yāni svavidhau sthānivadbhāvasya prayojanāni . \\
\noindent \textbf{(P\_1,1.57.1)} na etāni santi . \\
\noindent \textbf{(P\_1,1.57.1)} iha tāvat āyan , āsan , dhinvanti kṛṇvanti iti . \\
\noindent \textbf{(P\_1,1.57.1)} ayam vidhiśabdaḥ asti eva karmasādhanaḥ : vidhīyate vidhiḥ . \\
\noindent \textbf{(P\_1,1.57.1)} asti bhāvasādhanaḥ : vidhānam vidhiḥ iti . \\
\noindent \textbf{(P\_1,1.57.1)} karmasādhanasya vidhiśabdasya upādāne na sarvam iṣṭam saṅgṛhītam iti kṛtvā bhāvasādhanasya vidhiśabdasya upādānam vijñāsyate : pūrvasya vidhānam prati pūrvasya bhāvam prati pūrvaḥ syāt iti sthānivat bhavati iti evam āṭ bhaviṣyati iṭ ca na bhaviṣyati . \\
\noindent \textbf{(P\_1,1.57.1)} dadhi atra madhu atra cakratuḥ cakruḥ iti parihāram vakṣyati \\
\noindent \textbf{(P\_1,1.57.2)} kāni punaḥ asya yogasya prayojanāni . \\
\noindent \textbf{(P\_1,1.57.2)} stoṣyāmi aham pādikam audavāhim tataḥ śvobhūte śātanīm pātanīm ca . \\
\noindent \textbf{(P\_1,1.57.2)} netārau āgacchatam dhāraṇim rāvaṇim ca tataḥ paścāt sraṃsyate dhvaṃsyate ca . \\
\noindent \textbf{(P\_1,1.57.2)} iha tāvat pādikam audavāhim śātanīm pātanīm dhāraṇim rāvaṇim iti akāralope kṛte padbhāvaḥ ūṭh allopaḥ ṭilopaḥ iti ete vidhayaḥ prāpnuvanti . \\
\noindent \textbf{(P\_1,1.57.2)} sthānivadbhāvāt na bhavanti . \\
\noindent \textbf{(P\_1,1.57.2)} sraṃsyate dhvaṃsyate : ṇilope kṛte aniditām halaḥ upadhāyāḥ kṅiti iti nalopaḥ prāpnoti . \\
\noindent \textbf{(P\_1,1.57.2)} sthānivadbhāvāt na bhavati . \\
\noindent \textbf{(P\_1,1.57.2)} na etāni santi prayojanāni . \\
\noindent \textbf{(P\_1,1.57.2)} asiddhavat atra ā bhāt iti anena api etāni siddhāni . \\
\noindent \textbf{(P\_1,1.57.2)} idam tarhi prayojanam : yājyate vāpyate . \\
\noindent \textbf{(P\_1,1.57.2)} ṇilope kṛte yajādīnām kiti iti samprasāraṇam prāpnoti . \\
\noindent \textbf{(P\_1,1.57.2)} sthānivadbhāvāt na bhavati . \\
\noindent \textbf{(P\_1,1.57.2)} etat api na asti prayojanam . \\
\noindent \textbf{(P\_1,1.57.2)} yajādibhiḥ atra kitam viśeṣayiṣyāmaḥ yajādīnām yaḥ kit iti . \\
\noindent \textbf{(P\_1,1.57.2)} kaḥ ca yajādīnām kit . \\
\noindent \textbf{(P\_1,1.57.2)} yajādibhyaḥ yaḥ vihitaḥ iti . \\
\noindent \textbf{(P\_1,1.57.2)} idam tarhi prayojanam : paṭvyā mṛdvyā iti . \\
\noindent \textbf{(P\_1,1.57.2)} parasya yaṇādeśe kṛte pūrvasya na prāpnoti īkārayaṇā vyavahitatvāt . \\
\noindent \textbf{(P\_1,1.57.2)} sthānivadbhāvāt bhavati . \\
\noindent \textbf{(P\_1,1.57.2)} kim punaḥ kāraṇam parasya tāvat bhavati na punaḥ pūrvasya . \\
\noindent \textbf{(P\_1,1.57.2)} nityatvāt . \\
\noindent \textbf{(P\_1,1.57.2)} nityaḥ parayaṇādeśaḥ . \\
\noindent \textbf{(P\_1,1.57.2)} kṛte api pūrvayaṇādeśe prāpnoti akṛte api prāpnoti . \\
\noindent \textbf{(P\_1,1.57.2)} nityatvāt parayaṇādeśe kṛte pūrvasya na prāpnoti . \\
\noindent \textbf{(P\_1,1.57.2)} sthānivadbhāvāt bhavati . \\
\noindent \textbf{(P\_1,1.57.2)} etat api na asti prayojanam . \\
\noindent \textbf{(P\_1,1.57.2)} asiddham bahiraṅgalakṣaṇam antaraṅgalakṣaṇe iti asiddhatvāt bahiraṅgalakṣaṇasya parayaṇādeśasya antaraṅgalakṣaṇaḥ pūrvayaṇādeśaḥ bhaviṣyati . \\
\noindent \textbf{(P\_1,1.57.2)} avaśyam ca eṣā paribhāṣā āśrayitavyā svarārtham kartrya hartrya iti udāttayaṇaḥ halpūrvāt iti eṣaḥ svaraḥ yathā syāt . \\
\noindent \textbf{(P\_1,1.57.2)} anena api siddhaḥ svaraḥ . \\
\noindent \textbf{(P\_1,1.57.2)} katham . \\
\noindent \textbf{(P\_1,1.57.2)} ārabhyamāṇe nityaḥ asau . \\
\noindent \textbf{(P\_1,1.57.2)} ārabhyamāṇe tu asmin yoge nityaḥ pūrvayaṇādeśaḥ . \\
\noindent \textbf{(P\_1,1.57.2)} kṛte api parayaṇādeśe prāpnoti akṛte api . \\
\noindent \textbf{(P\_1,1.57.2)} parayaṇādeśaḥ api nityaḥ . \\
\noindent \textbf{(P\_1,1.57.2)} kṛte api pūrvayaṇādeśe prāpnoti akṛte api . \\
\noindent \textbf{(P\_1,1.57.2)} paraḥ ca asau vyavasthā . \\
\noindent \textbf{(P\_1,1.57.2)} vyavasthayā ca asau paraḥ . \\
\noindent \textbf{(P\_1,1.57.2)} yugapatsambhavaḥ na asti . \\
\noindent \textbf{(P\_1,1.57.2)} na ca asti yaugapadyena sambhavaḥ . \\
\noindent \textbf{(P\_1,1.57.2)} katham ca sidhyati . \\
\noindent \textbf{(P\_1,1.57.2)} bahiraṅgeṇa sidhyati . \\
\noindent \textbf{(P\_1,1.57.2)} asiddham bahiraṅgalakṣaṇam antaraṅgalakṣaṇe iti anena sidhyati . \\
\noindent \textbf{(P\_1,1.57.2)} evam tarhi yaḥ atra udāttayaṇ tadāśrayaḥ svaraḥ bhaviṣyati . \\
\noindent \textbf{(P\_1,1.57.2)} īkārayaṇā vyavahitatvāt na prāpnoti . \\
\noindent \textbf{(P\_1,1.57.2)} svaravidhau vyañjanam avidyamānavat bhavati iti na asti vyavadhānam . \\
\noindent \textbf{(P\_1,1.57.2)} sā tarhi eṣā paribhāṣā kartavyā . \\
\noindent \textbf{(P\_1,1.57.2)} nanu ca iyam api kartavyā : asiddham bahiraṅgalakṣaṇam antaraṅgalakṣaṇe iti . \\
\noindent \textbf{(P\_1,1.57.2)} bahuprayojanā eṣā paribhāṣā . \\
\noindent \textbf{(P\_1,1.57.2)} avaśyam eṣā kartavyā . \\
\noindent \textbf{(P\_1,1.57.2)} sā ca api eṣā lokataḥ siddhā . \\
\noindent \textbf{(P\_1,1.57.2)} katham . \\
\noindent \textbf{(P\_1,1.57.2)} pratyaṅgavartī lokaḥ lakṣyate . \\
\noindent \textbf{(P\_1,1.57.2)} tat yathā : puruṣaḥ ayam prātaḥ utthāya yāni asya pratiśarīram kāryāṇi tāni tāvat karoti tataḥ suhṛdām tataḥ sambandhinām . \\
\noindent \textbf{(P\_1,1.57.2)} prātipadikam ca api upadiṣṭam sāmānyabhūte arthe vartate . \\
\noindent \textbf{(P\_1,1.57.2)} sāmanye vartamānasya vyaktiḥ upajāyate . \\
\noindent \textbf{(P\_1,1.57.2)} vyaktasya sataḥ liṅgasaṅkhyābhyām anvitasya bāhyena arthena yogaḥ bhavati . \\
\noindent \textbf{(P\_1,1.57.2)} yayā eva ānupūrvyā arthānām prādurbhāvaḥ tayā eva śabdānām api tadvat kāryaiḥ api bhavitavyam . \\
\noindent \textbf{(P\_1,1.57.2)} imāni tarhi prayojanāni : paṭayati , avadhīt , bahukhaṭvakaḥ . \\
\noindent \textbf{(P\_1,1.57.2)} paṭayati laghayati iti ṭilope kṛte ataḥ upadhāyāḥ iti vṛddhiḥ prāpnoti . \\
\noindent \textbf{(P\_1,1.57.2)} sthānivadbhāvāt na bhavati . \\
\noindent \textbf{(P\_1,1.57.2)} avadhīt iti akāralope kṛte ataḥ halādeḥ laghoḥ iti vibhāṣā vṛddhiḥ prāpnoti . \\
\noindent \textbf{(P\_1,1.57.2)} sthānivadbhāvāt na bhavati . \\
\noindent \textbf{(P\_1,1.57.2)} bahukhaṭvakaḥ it āpaḥ anyatarasyām hrasvatve kṛte hrasvānte antyāt pūrvam iti eṣaḥ svaraḥ prāpnoti . \\
\noindent \textbf{(P\_1,1.57.2)} sthānivadbhāvāt na bhavati . \\
\noindent \textbf{(P\_1,1.57.3)} iha vaiyākaraṇaḥ , sauvaśvaḥ iti yvoḥ sthānivadbhāvāt āyāvau prāpnutaḥ . \\
\noindent \textbf{(P\_1,1.57.3)} tayoḥ pratiṣedhaḥ vaktavyaḥ . \\
\noindent \textbf{(P\_1,1.57.3)} acaḥ pūrvavijñānāt aicoḥ siddham . \\
\noindent \textbf{(P\_1,1.57.3)} yaḥ anādiṣṭāt acaḥ pūrvaḥ tasya vidhim prati sthānivadbhāvaḥ . \\
\noindent \textbf{(P\_1,1.57.3)} ādiṣṭāt ca eṣaḥ acaḥ pūrvaḥ . \\
\noindent \textbf{(P\_1,1.57.3)} kim vaktavyam etat . \\
\noindent \textbf{(P\_1,1.57.3)} na hi . \\
\noindent \textbf{(P\_1,1.57.3)} katham anucyamānam gaṃsyate . \\
\noindent \textbf{(P\_1,1.57.3)} acaḥ iti pañcamī : acaḥ pūrvasya . \\
\noindent \textbf{(P\_1,1.57.3)} yadi evam ādeśaḥ aviśeṣitaḥ bhavati . \\
\noindent \textbf{(P\_1,1.57.3)} ādeśaḥ ca viśeṣitaḥ . \\
\noindent \textbf{(P\_1,1.57.3)} katham . \\
\noindent \textbf{(P\_1,1.57.3)} na brūmaḥ yat ṣaṣṭhīnirdiṣṭam ajgrahaṇam tat pañcamīnirdiṣṭam kartavyam . \\
\noindent \textbf{(P\_1,1.57.3)} kim tarhi anyat kartavyam . \\
\noindent \textbf{(P\_1,1.57.3)} anyat ca na kartavyam . \\
\noindent \textbf{(P\_1,1.57.3)} yat eva adaḥ ṣaṣṭhīnirdiṣṭam ajgrahaṇam tasya dikśabdaiḥ yoge pañcamī bhavati : ajādeśaḥ paranimittakaḥ pūrvasya vidhim prati sthānivat bhavati . \\
\noindent \textbf{(P\_1,1.57.3)} kutaḥ pūrvasya . \\
\noindent \textbf{(P\_1,1.57.3)} acaḥ iti . \\
\noindent \textbf{(P\_1,1.57.3)} tat yathā ādeśaḥ prathamānirdiṣṭaḥ . \\
\noindent \textbf{(P\_1,1.57.3)} tasya dikśabdaiḥ yoge pañcamī bhavati : ajādeśaḥ paranimittakaḥ pūrvasya vidhim prati sthānivat bhavati . \\
\noindent \textbf{(P\_1,1.57.3)} kutaḥ pūrvasya . \\
\noindent \textbf{(P\_1,1.57.3)} ādeśāt iti . \\
\noindent \textbf{(P\_1,1.57.4)} tatra ādeśalakṣaṇapratiṣedhaḥ . \\
\noindent \textbf{(P\_1,1.57.4)} tatra ādeśalakṣaṇam kāryam prāpnoti . \\
\noindent \textbf{(P\_1,1.57.4)} tasya pratiṣedhaḥ vaktavyaḥ : vāyvoḥ , adhvaryvoḥ . \\
\noindent \textbf{(P\_1,1.57.4)} lopaḥ vyoḥ vali iti yalopaḥ prāpnoti . \\
\noindent \textbf{(P\_1,1.57.4)} asiddhavacanāt siddham . \\
\noindent \textbf{(P\_1,1.57.4)} ajādeśaḥ paranimittakaḥ pūrvasya vidhim prati asiddhaḥ bhavati iti vaktavyam . \\
\noindent \textbf{(P\_1,1.57.4)} asiddhavacanāt siddham iti cet utsargalakṣaṇānām anudeśaḥ . \\
\noindent \textbf{(P\_1,1.57.4)} asiddhavacanāt siddham iti cet utsargalakṣaṇānām anudeśaḥ kartavyaḥ : paṭvyā mrdvyā iti . \\
\noindent \textbf{(P\_1,1.57.4)} nanu ca etat api asiddhavacanāt siddham . \\
\noindent \textbf{(P\_1,1.57.4)} asiddhavacanāt siddham iti cet na anyasya asiddhavacanāt anyasya bhāvaḥ . \\
\noindent \textbf{(P\_1,1.57.4)} asiddhavacanāt siddham iti cet tat na . \\
\noindent \textbf{(P\_1,1.57.4)} kim kāraṇam . \\
\noindent \textbf{(P\_1,1.57.4)} na anyasya asiddhavacanāt anyasya bhāvaḥ . \\
\noindent \textbf{(P\_1,1.57.4)} na hi anyasya asiddhavacanāt anyasya prādurbhāvaḥ bhavati . \\
\noindent \textbf{(P\_1,1.57.4)} na hi devadattasya hantari hate devadattasya prādurbhāvaḥ bhavati . \\
\noindent \textbf{(P\_1,1.57.4)} tasmāt sthānivadvacanam asiddhatvam ca . \\
\noindent \textbf{(P\_1,1.57.4)} tasmāt sthānivadbhāvaḥ vaktavyaḥ asiddhatvam ca . \\
\noindent \textbf{(P\_1,1.57.4)} paṭvyā mṛdvyā iti atra sthānivadbhāvaḥ . \\
\noindent \textbf{(P\_1,1.57.4)} vāyvoḥ , adhvaryvoḥ iti asiddhatvam . \\
\noindent \textbf{(P\_1,1.57.4)} uktam vā . \\
\noindent \textbf{(P\_1,1.57.4)} kim uktam . \\
\noindent \textbf{(P\_1,1.57.4)} sthānivadvacanānarthakyam śāstrāsiddhatvāt iti . \\
\noindent \textbf{(P\_1,1.57.4)} viṣamaḥ upanyāsaḥ . \\
\noindent \textbf{(P\_1,1.57.4)} yuktam tatra yat ekādeśaśāstram tukśāstre asiddham syāt : anyat anyasmin . \\
\noindent \textbf{(P\_1,1.57.4)} iha punaḥ na yuktam . \\
\noindent \textbf{(P\_1,1.57.4)} katham hi tad eva nāma tasmin asiddham syāt . \\
\noindent \textbf{(P\_1,1.57.4)} tad eva ca api tasmin asiddham bhavati . \\
\noindent \textbf{(P\_1,1.57.4)} vakṣyati hi ācāryaḥ : ciṇaḥ luki tagrahaṇānarthakyam saṅghātasya apratyayatvāt talopasya ca asiddhatvāt iti . \\
\noindent \textbf{(P\_1,1.57.4)} ciṇaḥ luk ciṇaḥ luki eva asiddhaḥ bhavati . \\
\noindent \textbf{(P\_1,1.57.4)} kāmam atidiśyatām vā sat ca asat ca api na iha bhāraḥ asti . \\
\noindent \textbf{(P\_1,1.57.4)} kalpyaḥ hi vākyaśeṣaḥ vākyam vaktari adhīnam hi . \\
\noindent \textbf{(P\_1,1.57.4)} atha vā vatinirdeśaḥ ayam . \\
\noindent \textbf{(P\_1,1.57.4)} kāmacāraḥ ca vatinirdeśe vākyaśeṣam samarthayitum . \\
\noindent \textbf{(P\_1,1.57.4)} tat yathā . \\
\noindent \textbf{(P\_1,1.57.4)} uśīnaravat madreṣu yavāḥ . \\
\noindent \textbf{(P\_1,1.57.4)} santi na santi iti . \\
\noindent \textbf{(P\_1,1.57.4)} mātṛvat asyāḥ kalāḥ . \\
\noindent \textbf{(P\_1,1.57.4)} santi na santi . \\
\noindent \textbf{(P\_1,1.57.4)} evam iha api sthānivat bhavati sthānivat na bhavati iti vākyaśeṣam samarthayiṣyāmahe . \\
\noindent \textbf{(P\_1,1.57.4)} iha tāvat paṭvyā mṛdvyā iti yathā sthānini yaṇādeśaḥ bhavati evam ādeśe api . \\
\noindent \textbf{(P\_1,1.57.4)} iha idānīm vāyvoḥ adhvaryvoḥ iti yathā sthānini yalopaḥ na bhavati evam ādeśe api na bhavati . \\
\noindent \textbf{(P\_1,1.57.5)} kim punaḥ anantarasya vidhim prati sthānivadbhāvaḥ āhosvit pūrvamātrasya . \\
\noindent \textbf{(P\_1,1.57.5)} kaḥ ca atra viśeṣaḥ . \\
\noindent \textbf{(P\_1,1.57.5)} anantarasya cet ekānanudāttadvigusvaragatinighāteṣu upasaṅkhyānam . \\
\noindent \textbf{(P\_1,1.57.5)} anantarasya cet ekānanudāttadvigusvaragatinighāteṣu upasaṅkhyānam kartavyam . \\
\noindent \textbf{(P\_1,1.57.5)} ekānanudātta : lunīhi atra punīhi atra . \\
\noindent \textbf{(P\_1,1.57.5)} anudāttam padam ekavarjam iti eṣaḥ svaraḥ na prāpnoti . \\
\noindent \textbf{(P\_1,1.57.5)} dvigusvara : pañcāratnyaḥ , daśāratnyaḥ . \\
\noindent \textbf{(P\_1,1.57.5)} igantakāla iti eṣaḥ svaraḥ na prāpnoti . \\
\noindent \textbf{(P\_1,1.57.5)} gatinighāta : yat pralunīhi atra , yat prapunīhi atra . \\
\noindent \textbf{(P\_1,1.57.5)} tiṅi codāttavati iti eṣaḥ svaraḥ na prāpnoti . \\
\noindent \textbf{(P\_1,1.57.5)} astu tarhi pūrvamātrasya . \\
\noindent \textbf{(P\_1,1.57.5)} pūrvamātrasya iti cet upadhāhrasvatvam . pūrvamātrasya iti cet upadhāhrasvatvam vaktavyam : vāditavantam prayojitavān : avīvadat vīṇām parivādakena . \\
\noindent \textbf{(P\_1,1.57.5)} kim punaḥ kāraṇam na sidhyati . \\
\noindent \textbf{(P\_1,1.57.5)} yaḥ asau ṇau ṇiḥ lupyate tasya sthānivadbhāvāt hrasvatvam na prāpnoti . \\
\noindent \textbf{(P\_1,1.57.5)} gurusañjñā ca . \\
\noindent \textbf{(P\_1,1.57.5)} gurusañjñā ca na sidhyati : śleṣmā3ghna pittā3ghna dā3dhyaśva mā3dhvaśva . \\
\noindent \textbf{(P\_1,1.57.5)} halaḥ anantarāḥ saṃyogaḥ iti saṃyogasañjñā . \\
\noindent \textbf{(P\_1,1.57.5)} saṃyoge guru iti gurusañjñā . \\
\noindent \textbf{(P\_1,1.57.5)} guroḥ iti plutaḥ na prāpnoti . \\
\noindent \textbf{(P\_1,1.57.5)} nanu ca yasya api anantarasya vidhim prati sthānivadbhāvaḥ tasya api anantaralakṣaṇaḥ vidhiḥ saṃyogasañjñā vidheyā . \\
\noindent \textbf{(P\_1,1.57.5)} na vā saṃyogasya apūrvavidhitvāt . \\
\noindent \textbf{(P\_1,1.57.5)} na vā eṣaḥ doṣaḥ . \\
\noindent \textbf{(P\_1,1.57.5)} kim kāraṇam . \\
\noindent \textbf{(P\_1,1.57.5)} saṃyogasya apūrvavidhitvāt . \\
\noindent \textbf{(P\_1,1.57.5)} na pūrvavidhiḥ saṃyogaḥ . \\
\noindent \textbf{(P\_1,1.57.5)} kim tarhi . \\
\noindent \textbf{(P\_1,1.57.5)} pūrvaparavidhiḥ saṃyogaḥ . \\
\noindent \textbf{(P\_1,1.57.5)} ekādeśasya upasaṅkhyānam . \\
\noindent \textbf{(P\_1,1.57.5)} ekādeśasya upasaṅkhyānam kartavyam : śrāyasau gaumatau cāturau , ānaḍuhau pāde , udavāhe . \\
\noindent \textbf{(P\_1,1.57.5)} ekādeśe kṛte numāmau padbhāvaḥ ūṭh iti ete vidhayaḥ prāpnuvanti . \\
\noindent \textbf{(P\_1,1.57.5)} kim punaḥ kāraṇam na sidhyati . \\
\noindent \textbf{(P\_1,1.57.5)} ubhayanimittatvāt . \\
\noindent \textbf{(P\_1,1.57.5)} ajādeśaḥ paranimittakaḥ iti ucyate ubhayanimittaḥ ca ayam . \\
\noindent \textbf{(P\_1,1.57.5)} ubhayādeśatvāt ca . \\
\noindent \textbf{(P\_1,1.57.5)} acaḥ ādeśaḥ ici ucyate acoḥ ca ayam ādeśaḥ . \\
\noindent \textbf{(P\_1,1.57.5)} na eṣaḥ doṣaḥ . \\
\noindent \textbf{(P\_1,1.57.5)} yat tāvat ucyate ubhayanimittatvāt iti : iha yasya grāme nagare vā anekam kāryam bhavati śaknoti asau tataḥ anyatarat vyapadeṣṭum : gurunimittam vasāmaḥ . \\
\noindent \textbf{(P\_1,1.57.5)} adhyayananimittam vasāmaḥ iti . \\
\noindent \textbf{(P\_1,1.57.5)} yat api ucyate ubhayādeśatvāt ca iti .:iha yaḥ dvayoḥ ṣaṣṭhīnirdiṣṭayoḥ prasaṅge bhavati labhate asau anyatarataḥ vyapadeśam . \\
\noindent \textbf{(P\_1,1.57.5)} tat yathā devadattasya putraḥ , devadattāyāḥ putraḥ iti. . \\
\noindent \textbf{(P\_1,1.57.6)} atha halacoḥ ādeśaḥ sthānivat bhavati utāho na . \\
\noindent \textbf{(P\_1,1.57.6)} kaḥ ca atra viśeṣaḥ . \\
\noindent \textbf{(P\_1,1.57.6)} halacoḥ ādeśaḥ sthānivat iti cet viṃśateḥ tilopaḥ ekādeśaḥ . \\
\noindent \textbf{(P\_1,1.57.6)} halacoḥ ādeśaḥ sthānivat iti cet viṃśateḥ tilope ekādeśaḥ vaktavyaḥ : viṃśakaḥ , viṃśam śatam , viṃśaḥ . \\
\noindent \textbf{(P\_1,1.57.6)} sthūlādīnām yaṇādilope avādeśaḥ . \\
\noindent \textbf{(P\_1,1.57.6)} sthūlādīnām yaṇādilope kṛte avādeśaḥ vaktavyaḥ : sthavīyān , davīyān . \\
\noindent \textbf{(P\_1,1.57.6)} kekayimitrayvoḥ iyādeśe etvam . \\
\noindent \textbf{(P\_1,1.57.6)} kekayimitrayvoḥ iyādeśe etvam na sidhyati : kaikeyaḥ , maitreyaḥ . \\
\noindent \textbf{(P\_1,1.57.6)} aci iti etvam na sidhyati . \\
\noindent \textbf{(P\_1,1.57.6)} uttarapadalope ca . \\
\noindent \textbf{(P\_1,1.57.6)} uttarapadalope ca doṣaḥ bhavati : dadhyupasiktāḥ saktavaḥ dadhisaktavaḥ . \\
\noindent \textbf{(P\_1,1.57.6)} aci iti yaṇādeśaḥ prāpnoti . \\
\noindent \textbf{(P\_1,1.57.6)} yaṅlope yaṇiyaṅuvaṅaḥ . \\
\noindent \textbf{(P\_1,1.57.6)} yaṅlope yaṇiyaṅuvaṅaḥ na sidhyanti : cecyaḥ , nenyaḥ , cekriyaḥ , loluvaḥ , popuvaḥ . \\
\noindent \textbf{(P\_1,1.57.6)} aci iti yaṇiyaṅuvaṅaḥ na sidhyanti . \\
\noindent \textbf{(P\_1,1.57.6)} astu tarhi na sthānivat . \\
\noindent \textbf{(P\_1,1.57.6)} asthānivattve yaṅlope guṇavṛddhipratiṣedhaḥ . asthānivattve yaṅlope guṇavṛddhipratiṣedhaḥ vaktavyaḥ : loluvaḥ , popuvaḥ , sarīsṛpaḥ , marīmṛjaḥ iti . \\
\noindent \textbf{(P\_1,1.57.6)} na eṣaḥ doṣaḥ . \\
\noindent \textbf{(P\_1,1.57.6)} na dhātulope ārdhadhātuke iti pratiṣedhaḥ bhaviṣyati . \\
\noindent \textbf{(P\_1,1.57.7)} kim punaḥ āśrīyamāṇāyām prakṛtau sthānivat bhavati āhosvit aviśeṣeṇa . \\
\noindent \textbf{(P\_1,1.57.7)} kaḥ ca atra viśeṣaḥ . \\
\noindent \textbf{(P\_1,1.57.7)} aviśeṣeṇa sthānivat iti cet lopayaṇādeśe guruvidhiḥ . \\
\noindent \textbf{(P\_1,1.57.7)} aviśeṣeṇa sthānivat iti cet lopayaṇādeśe guruvidhiḥ na sidhyati : śleṣmā3ghna pittā3ghna dā3dhyaśva mā3dhvaśva . \\
\noindent \textbf{(P\_1,1.57.7)} halaḥ anantarāḥ saṃyogaḥ iti saṃyogasañjñā saṃyoge guru iti gurusañjñā guroḥ iti plutaḥ na prāpnoti . \\
\noindent \textbf{(P\_1,1.57.7)} dvirvacanādayaḥ ca pratiṣedhe . \\
\noindent \textbf{(P\_1,1.57.7)} dvirvacanādayaḥ ca pratiṣedhe vaktavyāḥ : dvirvacanavareyalopa iti . \\
\noindent \textbf{(P\_1,1.57.7)} ksalope lugvacanam . \\
\noindent \textbf{(P\_1,1.57.7)} ksalope luk vaktavyaḥ : adugdha , adugdhāḥ : luk vā duhadihalihaguhām ātmanepade dantye iti . \\
\noindent \textbf{(P\_1,1.57.7)} hanteḥ ghatvam . \\
\noindent \textbf{(P\_1,1.57.7)} hanteḥ ca ghatvam vaktavyam : ghnanti ghnantu , aghnan . \\
\noindent \textbf{(P\_1,1.57.7)} astu tarhi āśrīyamāṇāyām prakṛtau iti . \\
\noindent \textbf{(P\_1,1.57.7)} grahaṇeṣu sthānivat iti cet jagdhyādiṣu ādeśapratiṣedhaḥ . \\
\noindent \textbf{(P\_1,1.57.7)} grahaṇeṣu sthānivat iti cet jagdhyādiṣu ādeśasya pratiṣedhaḥ vaktavyaḥ : nirādya samādya . \\
\noindent \textbf{(P\_1,1.57.7)} adaḥ jagdhiḥ lyap ti kiti iti jagdhibhāvaḥ prāpnoti . \\
\noindent \textbf{(P\_1,1.57.7)} yaṇādeśe yulopetvānunāsikāttvapratiṣedhaḥ . \\
\noindent \textbf{(P\_1,1.57.7)} yaṇādeśe yulopetvānunāsikāttvānām pratiṣedhaḥ vaktavyaḥ . \\
\noindent \textbf{(P\_1,1.57.7)} yalopa : vāyvoḥ , adhvaryvoḥ . \\
\noindent \textbf{(P\_1,1.57.7)} lopaḥ vyoḥ vali iti yalopaḥ prāpnoti . \\
\noindent \textbf{(P\_1,1.57.7)} ulopa : akurvi* āśām akurvy āśām . \\
\noindent \textbf{(P\_1,1.57.7)} nityam karoteḥ ye ca iti ukāralopaḥ prāpnoti . \\
\noindent \textbf{(P\_1,1.57.7)} ītva : aluni* āśām aluny āśām . \\
\noindent \textbf{(P\_1,1.57.7)} ī hali aghoḥ iti ītvam prāpnoti . \\
\noindent \textbf{(P\_1,1.57.7)} anunāsikāttva : ajajñi* āśām ajajñy āśām . \\
\noindent \textbf{(P\_1,1.57.7)} ye vibhāṣā iti anunāsikāttvam prāpnoti . \\
\noindent \textbf{(P\_1,1.57.7)} rāyātvapratiṣedhaḥ ca . \\
\noindent \textbf{(P\_1,1.57.7)} rāyaḥ ātvasya ca pratiṣedhaḥ vaktavyaḥ : rāyi* āśām rāyy āśām . \\
\noindent \textbf{(P\_1,1.57.7)} rāyaḥ hali iti ātvam prāpnoti . \\
\noindent \textbf{(P\_1,1.57.7)} dīrghe yalopapratiṣedhaḥ . \\
\noindent \textbf{(P\_1,1.57.7)} dīrghe yalopasya pratiṣedhaḥ vaktavyaḥ : saurye nāma himavataḥ śrṅge tadvān sauryī himavān iti sau ināśraye dīrghatve kṛte īti yalopaḥ prāpnoti . \\
\noindent \textbf{(P\_1,1.57.7)} ataḥ dīrghe yalopavacanam . \\
\noindent \textbf{(P\_1,1.57.7)} ataḥ dīrghe yalopaḥ vaktavyaḥ : gārgābhyām , vātsābhyām . \\
\noindent \textbf{(P\_1,1.57.7)} dīrghe kṛte āpatyasya ca taddhite anāti iti pratiṣedhaḥ prāpnoti . \\
\noindent \textbf{(P\_1,1.57.7)} na eṣaḥ doṣaḥ . \\
\noindent \textbf{(P\_1,1.57.7)} āśrīyate tatra prakṛtiḥ : taddhite iti . \\
\noindent \textbf{(P\_1,1.57.7)} sarveṣām eṣām parihāraḥ : uktam vidhigrahaṇasya prayojanam vidhimātre sthānivat yathā syāt anāśrīyamāṇāyām api prakṛtau iti . \\
\noindent \textbf{(P\_1,1.57.7)} atha vā punaḥ astu aviśeṣeṇa sthānivat iti . \\
\noindent \textbf{(P\_1,1.57.7)} nanu ca uktam aviśeṣeṇa sthānivat iti cet lopayaṇādeśe guruvidhiḥ dvirvacanādayaḥ ca pratiṣedhe , ksalope lugvacanam , hanteḥ ghatvam iti . \\
\noindent \textbf{(P\_1,1.57.7)} na eṣaḥ doṣaḥ . \\
\noindent \textbf{(P\_1,1.57.7)} yat tāvat ucyate aviśeṣeṇa sthānivat iti cet lopayaṇādeśe guruvidhiḥ iti : uktam etat : na vā saṃyogasya apūrvavidhitvāt iti . \\
\noindent \textbf{(P\_1,1.57.7)} yat api ucyate dvirvacanādayaḥ ca pratiṣedhe vaktavyāḥ iti : ucyante nyāse eva . \\
\noindent \textbf{(P\_1,1.57.7)} ksalope lugvacanam iti : kriyate nyāse eva . \\
\noindent \textbf{(P\_1,1.57.7)} hanteḥ ghatvam iti . \\
\noindent \textbf{(P\_1,1.57.7)} saptame parihāram vakṣyati . \\
\noindent \textbf{(P\_1,1.58.1)} padāntavidhim prati na sthānivat iti ucyate . \\
\noindent \textbf{(P\_1,1.58.1)} tatra vetasvān iti ruḥ prāpnoti . \\
\noindent \textbf{(P\_1,1.58.1)} na eṣaḥ doṣaḥ . \\
\noindent \textbf{(P\_1,1.58.1)} bhasañjñā atra bādhikā bhaviṣyati : tasau matvarthe iti . \\
\noindent \textbf{(P\_1,1.58.1)} akārāntam etat bhasañjñām prati . \\
\noindent \textbf{(P\_1,1.58.1)} padasañjñām prati sakārāntam . \\
\noindent \textbf{(P\_1,1.58.1)} nanu ca evam vijñāsyate : yaḥ samprati padāntaḥ iti . \\
\noindent \textbf{(P\_1,1.58.1)} karmasādhanasya vidhiśabdasya upādāne etat evam syāt . \\
\noindent \textbf{(P\_1,1.58.1)} ayam ca vidhiśabdaḥ asti eva karmasādhanaḥ : vidhīyate vidhiḥ . \\
\noindent \textbf{(P\_1,1.58.1)} asti bhāvasādhanaḥ : vidhānam vidhiḥ iti . \\
\noindent \textbf{(P\_1,1.58.1)} tatra bhāvasādhanasya upādāne eṣaḥ doṣaḥ bhavati . \\
\noindent \textbf{(P\_1,1.58.1)} iha ca : brahmabandhvā brahmabandhvai : dhakārasya jaśtvam prāpnoti . \\
\noindent \textbf{(P\_1,1.58.1)} asti punaḥ kim cit bhāvasādhanasya vidhiśabdasya upādāne sati iṣṭam saṅgṛhītam āhosvit doṣāntam eva . \\
\noindent \textbf{(P\_1,1.58.1)} asti iti āha . \\
\noindent \textbf{(P\_1,1.58.1)} iha kāni santi yāni santi kau staḥ , yau staḥ iti yaḥ asau padāntaḥ yakāraḥ vakāraḥ vā śrūyeta saḥ na śrūyate . \\
\noindent \textbf{(P\_1,1.58.1)} ṣaḍikaḥ ca api siddhaḥ bhavati . \\
\noindent \textbf{(P\_1,1.58.1)} vācikaḥ tu na sidhyati . \\
\noindent \textbf{(P\_1,1.58.1)} astu tarhi karmasādhanaḥ . \\
\noindent \textbf{(P\_1,1.58.1)} yadi karmasādhanaḥ ṣaḍikaḥ na sidhyati . \\
\noindent \textbf{(P\_1,1.58.1)} astu tarhi bhāvasādhanaḥ . \\
\noindent \textbf{(P\_1,1.58.1)} vācikaḥ na sidhyati . \\
\noindent \textbf{(P\_1,1.58.1)} vācikaṣaḍikau na saṃvadete . \\
\noindent \textbf{(P\_1,1.58.1)} kartavyaḥ atra yatnaḥ . \\
\noindent \textbf{(P\_1,1.58.1)} katham brahmabandhvā brahmabandhvai. ubhayataḥ āśraye na antādivat iti . \\
\noindent \textbf{(P\_1,1.58.1)} katham vetasvān . \\
\noindent \textbf{(P\_1,1.58.1)} na evam vijñāyate : padasya antaḥ padāntaḥ padantavidhim prati iti . \\
\noindent \textbf{(P\_1,1.58.1)} katham tarhi . \\
\noindent \textbf{(P\_1,1.58.1)} pade antaḥ padāntaḥ padāntavidhim prati iti . \\
\noindent \textbf{(P\_1,1.58.1)} atha vā yathā eva anyāni api padakāryāṇi upaplavante rutvam jaśtvam ca evam idam api padakāryam upaploṣyate . \\
\noindent \textbf{(P\_1,1.58.1)} kim. bhasañjñā nāma . \\
\noindent \textbf{(P\_1,1.58.1)} vare yalopavidhim prati na sthānivat bhavati iti ucyate . \\
\noindent \textbf{(P\_1,1.58.1)} tatra te apsu yāyāvaraḥ pravapeta piṇḍān avarṇalopavidhim prati sthānivat syāt . \\
\noindent \textbf{(P\_1,1.58.1)} na eṣaḥ doṣaḥ . \\
\noindent \textbf{(P\_1,1.58.1)} na evam vijñāyate : vare yalopavidhim prati na sthānivat bhavati iti . \\
\noindent \textbf{(P\_1,1.58.1)} katham tarhi . \\
\noindent \textbf{(P\_1,1.58.1)} vare ayalopavidhim prati iti . \\
\noindent \textbf{(P\_1,1.58.1)} kim idam ayalopavidhim prati iti . \\
\noindent \textbf{(P\_1,1.58.1)} avarṇalopavidhim prati yalopavidhim ca prati iti . \\
\noindent \textbf{(P\_1,1.58.1)} atha vā yogavibhāgaḥ kariṣyate : vare luptam na sthānivat . \\
\noindent \textbf{(P\_1,1.58.1)} tataḥ yalopavidhim ca prati na sthānivat iti . \\
\noindent \textbf{(P\_1,1.58.1)} yalope kim udāharaṇam . \\
\noindent \textbf{(P\_1,1.58.1)} kaṇḍūyateḥ apratyayaḥ kaṇḍūḥ iti . \\
\noindent \textbf{(P\_1,1.58.1)} na etat asti . \\
\noindent \textbf{(P\_1,1.58.1)} kvau luptam na sthānivat . \\
\noindent \textbf{(P\_1,1.58.1)} idam tarhi : saurī balākā . \\
\noindent \textbf{(P\_1,1.58.1)} na etat asti . \\
\noindent \textbf{(P\_1,1.58.1)} upadhātvavidhim prati na sthānivat . \\
\noindent \textbf{(P\_1,1.58.1)} idam tarhi prayojanam : ādityaḥ . \\
\noindent \textbf{(P\_1,1.58.1)} na etat asti . \\
\noindent \textbf{(P\_1,1.58.1)} pūrvatrāsiddhe na sthānivat . \\
\noindent \textbf{(P\_1,1.58.1)} idam tarhi : kaṇḍūtiḥ , valgūtiḥ . \\
\noindent \textbf{(P\_1,1.58.1)} na etat asti prayojanam . \\
\noindent \textbf{(P\_1,1.58.1)} kaṇḍūyā valgūyā iti bhavitavyam . \\
\noindent \textbf{(P\_1,1.58.1)} idam tarhi : kaṇḍūyateḥ ktic : brāhmaṇakaṇḍūtiḥ , kṣatriyakaṇḍūtiḥ . \\
\noindent \textbf{(P\_1,1.58.2)} pratiṣedhe svaradīrghayalopeṣu lopājādeśaḥ na sthānivat . \\
\noindent \textbf{(P\_1,1.58.2)} pratiṣedhe svaradīrghayalopeṣu lopājādeśaḥ na sthānivat iti vaktavyam . \\
\noindent \textbf{(P\_1,1.58.2)} svara : ākarṣikaḥ , cikīrṣakaḥ , jihīrṣakaḥ . \\
\noindent \textbf{(P\_1,1.58.2)} yaḥ hi anyaḥ ādeśaḥ sthānivat eva asau bhavati : pañcāratnyaḥ , daśāratnyaḥ . \\
\noindent \textbf{(P\_1,1.58.2)} svara . \\
\noindent \textbf{(P\_1,1.58.2)} dīrgha : pratidīvnā pratidīvne . \\
\noindent \textbf{(P\_1,1.58.2)} yaḥ hi anyaḥ ādeśaḥ sthānivat eva asau bhavati : kiryoḥ , giryoḥ . \\
\noindent \textbf{(P\_1,1.58.2)} dīrgha . \\
\noindent \textbf{(P\_1,1.58.2)} yalopa : brāhmaṇakaṇḍūtiḥ , kṣatriyakaṇḍūtiḥ . \\
\noindent \textbf{(P\_1,1.58.2)} yaḥ hi anyaḥ ādeśaḥ sthānivat eva asau bhavati : vāyvoḥ , adhvaryvoḥ iti . \\
\noindent \textbf{(P\_1,1.58.2)} tat tarhi vaktavyam . \\
\noindent \textbf{(P\_1,1.58.2)} na vaktavyam . \\
\noindent \textbf{(P\_1,1.58.2)} iha hi lopaḥ api prakṛtaḥ ādeśaḥ api . \\
\noindent \textbf{(P\_1,1.58.2)} vidhigrahaṇam api prakṛtam anuvartate . \\
\noindent \textbf{(P\_1,1.58.2)} dīrghādayaḥ api nirdiśyante . \\
\noindent \textbf{(P\_1,1.58.2)} kevalam atra abhisambandhamātram kartavyam : svaradīrghayalopavidhiṣu lopājādeśaḥ na sthānivat iti . \\
\noindent \textbf{(P\_1,1.58.2)} ānupūrvyeṇa sanniviṣṭānām yatheṣṭam abhisambandhaḥ śakyate kartum . \\
\noindent \textbf{(P\_1,1.58.2)} na ca etani ānupūrvyeṇa sanniviṣṭāni . \\
\noindent \textbf{(P\_1,1.58.2)} anānupūrvyeṇa api sanniviṣṭānām yatheṣtam abhisambandhaḥ bhavati . \\
\noindent \textbf{(P\_1,1.58.2)} tat yathā : anaḍvāham udahāri yā tvam harasi śirasā kumbham bhagini sācīnam abhidhāvantam adrākṣīḥ iti . \\
\noindent \textbf{(P\_1,1.58.2)} tasya yatheṣtam abhisambandhaḥ bhavati : udahāri bhagini yā tvam kumbham harasi śirasā anaḍvāham sācīnam abhidhāvantam adrākṣīḥ iti . \\
\noindent \textbf{(P\_1,1.58.3)} kvilugupadhātvacaṅparanirhrāsakutveṣu upasaṅkhyānam . \\
\noindent \textbf{(P\_1,1.58.3)} kvilugupadhātvacaṅparanirhrāsakutveṣu upasaṅkhyānam kartavyam . \\
\noindent \textbf{(P\_1,1.58.3)} kvau kim udāharaṇam . \\
\noindent \textbf{(P\_1,1.58.3)} kaṇḍūyateḥ apratyayaḥ kaṇḍūḥ iti . \\
\noindent \textbf{(P\_1,1.58.3)} na etat asti . \\
\noindent \textbf{(P\_1,1.58.3)} yalopavidhim prati na sthānivat . \\
\noindent \textbf{(P\_1,1.58.3)} idam tarhi : pipaṭhiṣateḥ apratyayaḥ pipaṭhīḥ . \\
\noindent \textbf{(P\_1,1.58.3)} na etat asti . \\
\noindent \textbf{(P\_1,1.58.3)} dīrghatvam prati na sthānivat . \\
\noindent \textbf{(P\_1,1.58.3)} idam tarhi : lāvayateḥ lauḥ , pāvayateḥ pauḥ . \\
\noindent \textbf{(P\_1,1.58.3)} na etat asti . \\
\noindent \textbf{(P\_1,1.58.3)} akṛtvā vṛddhyāvādeśau ṇilopaḥ . \\
\noindent \textbf{(P\_1,1.58.3)} pratyayalakṣaṇena vṛddhiḥ bhaviṣyati . \\
\noindent \textbf{(P\_1,1.58.3)} idam tarhi : lavam ācaṣṭe lavayati . \\
\noindent \textbf{(P\_1,1.58.3)} lavayateḥ apratyayaḥ lauḥ , pauḥ . \\
\noindent \textbf{(P\_1,1.58.3)} sthānivadbhāvāt ṇeḥ ūṭh na prāpnoti . \\
\noindent \textbf{(P\_1,1.58.3)} kvau luptam na sthānivat iti bhavati . \\
\noindent \textbf{(P\_1,1.58.3)} evam api na sidhyati . \\
\noindent \textbf{(P\_1,1.58.3)} katham . \\
\noindent \textbf{(P\_1,1.58.3)} kvau ṇilopaḥ ṇau akāralopaḥ . \\
\noindent \textbf{(P\_1,1.58.3)} tasya sthānivadbhāvāt ūṭh na prāpnoti . \\
\noindent \textbf{(P\_1,1.58.3)} na eṣaḥ doṣaḥ . \\
\noindent \textbf{(P\_1,1.58.3)} na evam vijñāyate : kvau luptam na sthānivat iti . \\
\noindent \textbf{(P\_1,1.58.3)} katham tarhi . \\
\noindent \textbf{(P\_1,1.58.3)} kvau vidhim prati na sthānivat . \\
\noindent \textbf{(P\_1,1.58.3)} luki kim udāharaṇam . \\
\noindent \textbf{(P\_1,1.58.3)} bimbam , badaram . \\
\noindent \textbf{(P\_1,1.58.3)} na etat asti . \\
\noindent \textbf{(P\_1,1.58.3)} puṃvadbhāvena api etat siddham . \\
\noindent \textbf{(P\_1,1.58.3)} idam tarhi : āmalakam . \\
\noindent \textbf{(P\_1,1.58.3)} etat api na asti . \\
\noindent \textbf{(P\_1,1.58.3)} vakṣyati etat : phale lugvacanānarthakyam prakṛtyantaratvāt iti . \\
\noindent \textbf{(P\_1,1.58.3)} idam tarhi : pañcabhiḥ paṭvībhiḥ krītaḥ pañcapaṭuḥ , daśapaṭuḥ iti . \\
\noindent \textbf{(P\_1,1.58.3)} nanu ca etat api puṃvadbhāvena eva siddham . \\
\noindent \textbf{(P\_1,1.58.3)} katham puṃvadbhāvaḥ . \\
\noindent \textbf{(P\_1,1.58.3)} bhasya aḍhe taddhite puṃvat bhavati iti . \\
\noindent \textbf{(P\_1,1.58.3)} bhasya iti ucyate . \\
\noindent \textbf{(P\_1,1.58.3)} yajādau ca bham bhavati na ca atra yajādim paśyāmaḥ . \\
\noindent \textbf{(P\_1,1.58.3)} pratyayalakṣaṇena yajādiḥ . \\
\noindent \textbf{(P\_1,1.58.3)} varṇāśraye na asti pratyayalakṣaṇam . \\
\noindent \textbf{(P\_1,1.58.3)} evam tarhi ṭhakchasoḥ ca iti evam bhaviṣyati . \\
\noindent \textbf{(P\_1,1.58.3)} ṭakchasoḥ ca iti ucyate . \\
\noindent \textbf{(P\_1,1.58.3)} na ca atra ṭakchasau paśyāmaḥ . \\
\noindent \textbf{(P\_1,1.58.3)} pratyayalakṣaṇena . \\
\noindent \textbf{(P\_1,1.58.3)} na lumatā tasmin iti pratyayalakṣaṇasya pratiṣedhaḥ . \\
\noindent \textbf{(P\_1,1.58.3)} na khalu api ṭhak eva krītapratyayaḥ krītādyarthāḥ eva vā taddhitāḥ . \\
\noindent \textbf{(P\_1,1.58.3)} kim tarhi . \\
\noindent \textbf{(P\_1,1.58.3)} anye api taddhitāḥ ye lukam prayojayanti : pañcendrāṇyaḥ devatāḥ asya iti pañcendraḥ , daśendraḥ , pañcāgniḥ , daśāgniḥ . \\
\noindent \textbf{(P\_1,1.58.3)} upadhātve kim udāharaṇam . \\
\noindent \textbf{(P\_1,1.58.3)} pipaṭhiṣateḥ apratyayaḥ pipaṭhīḥ iti . \\
\noindent \textbf{(P\_1,1.58.3)} na etat asti . \\
\noindent \textbf{(P\_1,1.58.3)} dīrghavidhim prati na sthānivat . \\
\noindent \textbf{(P\_1,1.58.3)} idam tarhi: saurī balākā . \\
\noindent \textbf{(P\_1,1.58.3)} na etat asti . \\
\noindent \textbf{(P\_1,1.58.3)} yalopavidhim prati na sthānivat . \\
\noindent \textbf{(P\_1,1.58.3)} idam tarhi : pārikhīyaḥ . \\
\noindent \textbf{(P\_1,1.58.3)} caṅparanirhrāse ca upasaṅkhyanam kartavyam . \\
\noindent \textbf{(P\_1,1.58.3)} vāditavantam prayojitavān : avīvadat vīṇām parivādakena . \\
\noindent \textbf{(P\_1,1.58.3)} kim punaḥ kāraṇam na sidhyati . \\
\noindent \textbf{(P\_1,1.58.3)} yaḥ asau ṇau ṇiḥ lupyate tasya sthānivadbhāvāt hrasvatvam na prāpnoti . \\
\noindent \textbf{(P\_1,1.58.3)} nanu ca etat api upadhātvavidhim prati na sthānivat iti eva siddham . \\
\noindent \textbf{(P\_1,1.58.3)} viśeṣe etat vaktavyam . \\
\noindent \textbf{(P\_1,1.58.3)} kva . \\
\noindent \textbf{(P\_1,1.58.3)} pratyayavidhau iti . \\
\noindent \textbf{(P\_1,1.58.3)} iha mā bhūt : paṭayati laghayati iti . \\
\noindent \textbf{(P\_1,1.58.3)} kutve ca upasaṅkhyanam kartavyam . \\
\noindent \textbf{(P\_1,1.58.3)} arcayateḥ arkaḥ , marcayateḥ markaḥ . \\
\noindent \textbf{(P\_1,1.58.3)} na etat ghañantam . \\
\noindent \textbf{(P\_1,1.58.3)} auṇādikaḥ eṣaḥ kaśabdaḥ . \\
\noindent \textbf{(P\_1,1.58.3)} tasmin āṣṭamikam kutvam . \\
\noindent \textbf{(P\_1,1.58.3)} etat api ṇicā vyavahitatvāt na prāpnoti . \\
\noindent \textbf{(P\_1,1.58.4)} pūrvatrāsiddhe ca . \\
\noindent \textbf{(P\_1,1.58.4)} pūrvatrāsiddhe ca na sthānivat iti vaktavyam . \\
\noindent \textbf{(P\_1,1.58.4)} kim prayojanam . \\
\noindent \textbf{(P\_1,1.58.4)} prayojanam ksalopaḥ salope ksalopaḥ salope prayojanam : adugdha , adugdhāḥ . \\
\noindent \textbf{(P\_1,1.58.4)} luk vā duhadihalihaguhām ātmanepade dantye iti luggrahaṇam na kartavyam . \\
\noindent \textbf{(P\_1,1.58.4)} dadhaḥ ākāralope ādicaturthatve . \\
\noindent \textbf{(P\_1,1.58.4)} dadhaḥ ākāralope ādicaturthatve prayojanam : dhatse dhaddhve dhaddhvam iti . \\
\noindent \textbf{(P\_1,1.58.4)} dadhaḥ tathoḥ ca iti cakāraḥ na kartavyaḥ bhavati . \\
\noindent \textbf{(P\_1,1.58.4)} halaḥ yamām yami lope . \\
\noindent \textbf{(P\_1,1.58.4)} halaḥ yamām yami lope prayojanam : ādityaḥ . \\
\noindent \textbf{(P\_1,1.58.4)} halaḥ yamām yami lopaḥ siddhaḥ bhavati . \\
\noindent \textbf{(P\_1,1.58.4)} allopaṇilopau saṃyogāntalopaprabhṛtiṣu . \\
\noindent \textbf{(P\_1,1.58.4)} allopaṇilopau saṃyogāntalopaprabhṛtiṣu prayojanam : pāpacyateḥ pāpaktiḥ , yāyajyateḥ yāyaṣṭiḥ , pācayateḥ pāktiḥ , yājayateḥ yāṣṭiḥ . \\
\noindent \textbf{(P\_1,1.58.4)} dvirvacanādīni ca . \\
\noindent \textbf{(P\_1,1.58.4)} dvirvacanādīni ca na paṭhitavyāni bhavanti . \\
\noindent \textbf{(P\_1,1.58.4)} pūrvatrāsiddhena eva siddhāni bhavanti . \\
\noindent \textbf{(P\_1,1.58.4)} kim aviśeṣeṇa . \\
\noindent \textbf{(P\_1,1.58.4)} na iti āha . \\
\noindent \textbf{(P\_1,1.58.4)} vareyalopasvaravarjam . \\
\noindent \textbf{(P\_1,1.58.4)} vareyalopam svaram ca varjayitvā . \\
\noindent \textbf{(P\_1,1.58.4)} tasya doṣaḥ saṃyogādilopalatvaṇatveṣu . \\
\noindent \textbf{(P\_1,1.58.4)} tasya etasya lakṣaṇasya doṣaḥ saṃyogādilopalatvaṇatveṣu . \\
\noindent \textbf{(P\_1,1.58.4)} saṃyogādilopa : kākyartham , vāsyartham . \\
\noindent \textbf{(P\_1,1.58.4)} skoḥ saṃyogādyoḥ ante ca iti lopaḥ prāpnoti . \\
\noindent \textbf{(P\_1,1.58.4)} latvam : nigāryate nigālyate . \\
\noindent \textbf{(P\_1,1.58.4)} aci vibhāṣā iti latvam na prāpnoti . \\
\noindent \textbf{(P\_1,1.58.4)} ṇatvam : māṣavapanī vrīhivāpanī . \\
\noindent \textbf{(P\_1,1.58.4)} prātipadikāntasya iti ṇatvam prāpnoti . \\
\noindent \textbf{(P\_1,1.59.1)} ādeśe sthānivadanudeśāt tadvataḥ dvirvacanam . \\
\noindent \textbf{(P\_1,1.59.1)} ādeśe sthānivadanudeśāt tadvataḥ . \\
\noindent \textbf{(P\_1,1.59.1)} kiṃvataḥ . \\
\noindent \textbf{(P\_1,1.59.1)} ādeśavataḥ dvirvacanam prāpnoti . \\
\noindent \textbf{(P\_1,1.59.1)} tata kaḥ doṣaḥ . \\
\noindent \textbf{(P\_1,1.59.1)} tatra abhyāsarūpam . \\
\noindent \textbf{(P\_1,1.59.1)} tatra abhyāsarūpam na sidhyati : cakratuḥ , cakruḥ iti . \\
\noindent \textbf{(P\_1,1.59.1)} ajgrahaṇam tu jñāpakam rūpasthānivadbhāvasya . \\
\noindent \textbf{(P\_1,1.59.1)} yat ayam ajgrahaṇam karoti tat jñāpayati ācāryaḥ rūpam sthānivat bhavati iti . \\
\noindent \textbf{(P\_1,1.59.1)} katham kṛtvā jñāpakam . \\
\noindent \textbf{(P\_1,1.59.1)} ajgrahaṇasya etat prayojanam : iha mā bhūt : jeghrīyate , dedhmīyate iti . \\
\noindent \textbf{(P\_1,1.59.1)} yadi rūpam sthānivat bhavati tataḥ ajgrahaṇam arthavat bhavati . \\
\noindent \textbf{(P\_1,1.59.1)} atha hi kāryam na arthaḥ ajgrahaṇena . \\
\noindent \textbf{(P\_1,1.59.1)} bhavati eva atra dvirvacanam . \\
\noindent \textbf{(P\_1,1.59.2)} tatra gāṅpratiṣedhaḥ . tatra gāṅaḥ pratiṣedhaḥ vaktavyaḥ : adhijage . \\
\noindent \textbf{(P\_1,1.59.2)} ivarṇābhyāsatā prāpnoti . \\
\noindent \textbf{(P\_1,1.59.2)} na vaktavyaḥ . \\
\noindent \textbf{(P\_1,1.59.2)} gāṅ liṭi iti dvilakārakaḥ nirdeśaḥ : liṭi lakārādau iti . \\
\noindent \textbf{(P\_1,1.59.2)} kṝtyejantadivādināmadhātuṣu abhyāsarūpam . \\
\noindent \textbf{(P\_1,1.59.2)} kṝtyejantadivādināmadhātuṣu abhyāsarūpam na sidhyati . \\
\noindent \textbf{(P\_1,1.59.2)} kṝti : acikīrtat . \\
\noindent \textbf{(P\_1,1.59.2)} kṝti . \\
\noindent \textbf{(P\_1,1.59.2)} ejanta : jagle mamle . \\
\noindent \textbf{(P\_1,1.59.2)} ejanta . \\
\noindent \textbf{(P\_1,1.59.2)} divādi : dudyūṣati susyūṣati . \\
\noindent \textbf{(P\_1,1.59.2)} divādi . \\
\noindent \textbf{(P\_1,1.59.2)} nāmadhātu : bhavanam icchati bhavanīyati bhavanīyateḥ san : bibhavanīyiṣati . \\
\noindent \textbf{(P\_1,1.59.2)} evam tarhi pratyaye iti vakṣyāmi . \\
\noindent \textbf{(P\_1,1.59.2)} pratyaye iti cet kṝtyejantanamadhātuṣu abhyāsarūpam . \\
\noindent \textbf{(P\_1,1.59.2)} pratyaye iti cet kṝtyejantanamadhātuṣu abhyāsarūpam na sidhyati . \\
\noindent \textbf{(P\_1,1.59.2)} divādayaḥ eke parihṛtāḥ . \\
\noindent \textbf{(P\_1,1.59.2)} evam tarhi dvirvacananimitte aci ajādeśaḥ sthānivat iti vakṣyāmi . \\
\noindent \textbf{(P\_1,1.59.2)} saḥ tarhi nimittaśabdaḥ upādeyaḥ . \\
\noindent \textbf{(P\_1,1.59.2)} na hi antareṇa nimittaśabdam nimittārthaḥ gamyate . \\
\noindent \textbf{(P\_1,1.59.2)} antareṇa api nimittaśabdam nimittārthaḥ gamyate . \\
\noindent \textbf{(P\_1,1.59.2)} tat yathā : dadhitrapusam pratyakṣaḥ jvaraḥ . \\
\noindent \textbf{(P\_1,1.59.2)} jvaranimittam iti gamyate . \\
\noindent \textbf{(P\_1,1.59.2)} naḍvalodakam pādarogaḥ . \\
\noindent \textbf{(P\_1,1.59.2)} pādaroganimittam iti gamyate . \\
\noindent \textbf{(P\_1,1.59.2)} ayuḥ ghṛtam . \\
\noindent \textbf{(P\_1,1.59.2)} āyuṣaḥ nimittam iti gamyate . \\
\noindent \textbf{(P\_1,1.59.2)} atha vā akāraḥ matvarthīyaḥ : dvirvacanam asmin asti saḥ ayam dvirvacanaḥ , dvirvacane iti . \\
\noindent \textbf{(P\_1,1.59.2)} evam api na jñāyate kiyantam asau kālam sthānivat bhavati iti . \\
\noindent \textbf{(P\_1,1.59.2)} yaḥ punaḥ āha dvirvacane kartavye iti kṛte tasya dvirvacane sthānivat na bhaviṣyati . \\
\noindent \textbf{(P\_1,1.59.2)} evam tarhi pratiṣedhaḥ prakṛtaḥ . \\
\noindent \textbf{(P\_1,1.59.2)} saḥ anuvartiṣyate . \\
\noindent \textbf{(P\_1,1.59.2)} kva prakṛtaḥ . \\
\noindent \textbf{(P\_1,1.59.2)} na padāntadvirvacana iti . \\
\noindent \textbf{(P\_1,1.59.2)} dvirvacananimitte aci ajādeśaḥ na bhavati iti . \\
\noindent \textbf{(P\_1,1.59.2)} evam api na jñāyate kiyantam asau kālam na bhavati iti . \\
\noindent \textbf{(P\_1,1.59.2)} yaḥ punaḥ āha dvirvacane kartavye iti kṛte tasya dvirvacane ajādeśaḥ bhaviṣyati . \\
\noindent \textbf{(P\_1,1.59.2)} evam tarhi ubhayam anena kriyate : pratyayaḥ ca viśeṣyate dvirvacanam ca . \\
\noindent \textbf{(P\_1,1.59.2)} katham punaḥ ekena yatnena ubhayam labhyam . \\
\noindent \textbf{(P\_1,1.59.2)} labhyam iti āha . \\
\noindent \textbf{(P\_1,1.59.2)} katham . \\
\noindent \textbf{(P\_1,1.59.2)} ekaśeṣanirdeśāt . \\
\noindent \textbf{(P\_1,1.59.2)} ekaśeṣanirdeśaḥ ayam : dvirvacanam ca dvirvacanam ca dvirvacanam . \\
\noindent \textbf{(P\_1,1.59.2)} dvirvacane ca kartavye dvirvacane aci pratyaye iti dvirvacananimitte aci sthānivat bhavati . \\
\noindent \textbf{(P\_1,1.59.2)} dvirvacananimitte aci sthānivat iti cet ṇau sthānivadvacanam . dvirvacananimitte aci sthānivat iti cet ṇau sthānivadbhāvaḥ vaktavyaḥ : avanunāvayiṣati , avacukṣāvayiṣati . \\
\noindent \textbf{(P\_1,1.59.2)} na vaktavyaḥ . \\
\noindent \textbf{(P\_1,1.59.2)} oḥ puyaṇjiṣu vacanam jñāpakam ṇau sthānivadbhāvasya . \\
\noindent \textbf{(P\_1,1.59.2)} yat ayam puyaṇji apare iti āha tat jñāpayati ācāryaḥ bhavati ṇau sthānivat iti . \\
\noindent \textbf{(P\_1,1.59.2)} yadi etat jñāpyate acīkīrtat atra api prāpnoti . \\
\noindent \textbf{(P\_1,1.59.2)} tulyajātīyasya jñāpakam . \\
\noindent \textbf{(P\_1,1.59.2)} kaḥ ca tulyajātīyaḥ . \\
\noindent \textbf{(P\_1,1.59.2)} yathājātīyakāḥ puyaṇjayaḥ . \\
\noindent \textbf{(P\_1,1.59.2)} kathañjātīyakāḥ ca ete . \\
\noindent \textbf{(P\_1,1.59.2)} avarṇaparāḥ . \\
\noindent \textbf{(P\_1,1.59.2)} katham jagle mamle . \\
\noindent \textbf{(P\_1,1.59.2)} anaimittikam āttvam śiti tu pratiṣedhaḥ . \\
\noindent \textbf{(P\_1,1.59.3)} kāni punaḥ asya yogasya prayojanāni . \\
\noindent \textbf{(P\_1,1.59.3)} papatuḥ , papuḥ , tasthatuḥ , tasthuḥ , jagmatuḥ , jagmuḥ , āṭitat , āśiśat , cakratuḥ , cakruḥ iti . \\
\noindent \textbf{(P\_1,1.59.3)} āllopopadhālopaṇilopayaṇādeśeṣu kṛteṣu anackatvāt dvirvacanam na prāpnoti . \\
\noindent \textbf{(P\_1,1.59.3)} sthānivadbhāvāt bhavati . \\
\noindent \textbf{(P\_1,1.59.3)} na etāni santi prayojanāni . \\
\noindent \textbf{(P\_1,1.59.3)} pūrvavipratiṣedhena api etāni siddhāni . \\
\noindent \textbf{(P\_1,1.59.3)} katham . \\
\noindent \textbf{(P\_1,1.59.3)} vakṣyati hi ācāryaḥ : dvirvacanam yaṇayavāyāvādeśāllopopadhālopakikinoruttvebhyaḥ iti . \\
\noindent \textbf{(P\_1,1.59.3)} saḥ pūrvavipratiṣedhaḥ na paṭhitavyaḥ bhavati . \\
\noindent \textbf{(P\_1,1.59.3)} kim punaḥ atra jyāyaḥ . \\
\noindent \textbf{(P\_1,1.59.3)} sthānivadbhāvaḥ eva jyāyān . \\
\noindent \textbf{(P\_1,1.59.3)} pūrvavipratiṣedhe hi sati idam vaktavyam syāt : odaudādeśasya ut bhavati cuṭutuśarādeḥ abhyāsasya iti . \\
\noindent \textbf{(P\_1,1.59.3)} nanu ca tvayā api ittvam vaktavyam . \\
\noindent \textbf{(P\_1,1.59.3)} parārtham mama bhaviṣyati : sani ataḥ it bhavati iti . \\
\noindent \textbf{(P\_1,1.59.3)} mama api tarhi uttvam parārtham bhaviṣyati : utparasya ataḥ ti ca iti . \\
\noindent \textbf{(P\_1,1.59.3)} ittvam api tvayā vaktavyam yat samānāśrayam tadartham : utpipaviṣate saṃyiyaviṣati iti evamartham . \\
\noindent \textbf{(P\_1,1.59.3)} tasmāt sthānivat iti eṣaḥ eva pakṣaḥ jyāyān . \\
\noindent \textbf{(P\_1,1.60)} arthasya sañjñā kartavyā śabdasya mā bhūt iti . \\
\noindent \textbf{(P\_1,1.60)} itaretarāśrayam ca bhavati . \\
\noindent \textbf{(P\_1,1.60)} kā itaretarāśrayatā . \\
\noindent \textbf{(P\_1,1.60)} sataḥ adarśanasya sañjñayā bhavitavyam sañjñaya ca adarśanam bhāvyate . \\
\noindent \textbf{(P\_1,1.60)} tat etat itaretarāśrayam bhavati . \\
\noindent \textbf{(P\_1,1.60)} itaretarāśrayāṇi ca kāryāṇi na prakalpante . \\
\noindent \textbf{(P\_1,1.60)} lopasañjñāyām arthasatoḥ uktam . \\
\noindent \textbf{(P\_1,1.60)} kim uktam . \\
\noindent \textbf{(P\_1,1.60)} arthasya tāvat uktam : itikaraṇaḥ arthanirdeśārthaḥ iti . \\
\noindent \textbf{(P\_1,1.60)} sataḥ api uktam : siddham tu nityaśabdatvāt iti . \\
\noindent \textbf{(P\_1,1.60)} nityāḥ śabdāḥ . \\
\noindent \textbf{(P\_1,1.60)} nityeṣu ca śabdeṣu sataḥ adarśanasya sañjñā kriyate . \\
\noindent \textbf{(P\_1,1.60)} na sañjñayā adarśanam bhāvyate . \\
\noindent \textbf{(P\_1,1.60)} sarvaprasaṅgaḥ tu sarvasya anyatra adṛṣṭatvāt . \\
\noindent \textbf{(P\_1,1.60)} sarvaprasaṅgaḥ tu bhavati . \\
\noindent \textbf{(P\_1,1.60)} sarvasya adarśanasya lopasañjñā prāpnoti . \\
\noindent \textbf{(P\_1,1.60)} kim kāraṇam . \\
\noindent \textbf{(P\_1,1.60)} sarvasya anyatra adṛṣṭatvāt . \\
\noindent \textbf{(P\_1,1.60)} sarvaḥ hi śabdaḥ yaḥ yasya prayogaviṣayaḥ saḥ tataḥ anyatra na dṛśyate . \\
\noindent \textbf{(P\_1,1.60)} trapu jatu iti atra aṇaḥ adarśanam . \\
\noindent \textbf{(P\_1,1.60)} tatra adarśanam lopaḥ iti lopasañjñā prāpnoti . \\
\noindent \textbf{(P\_1,1.60)} tatra kaḥ doṣaḥ . \\
\noindent \textbf{(P\_1,1.60)} tatra pratyayalakṣaṇapratiṣedhaḥ . \\
\noindent \textbf{(P\_1,1.60)} tatra pratyayalakṣaṇam kāryam prāpnoti . \\
\noindent \textbf{(P\_1,1.60)} tasya pratiṣedhaḥ vaktavyaḥ . \\
\noindent \textbf{(P\_1,1.60)} acaḥ ñṇiti iti vṛddhiḥ prāpnoti . \\
\noindent \textbf{(P\_1,1.60)} na eṣaḥ doṣaḥ . \\
\noindent \textbf{(P\_1,1.60)} ñṇiti aṅgasya acaḥ vṛddhiḥ ucyate . \\
\noindent \textbf{(P\_1,1.60)} yasmāt pratyayavidhiḥ tadādi pratyaye aṅgam bhavati . \\
\noindent \textbf{(P\_1,1.60)} yasmāt ca atra pratyayavidhiḥ na tat pratyaye parataḥ yat ca pratyaye parataḥ na tasmāt pratyayavidhiḥ . \\
\noindent \textbf{(P\_1,1.60)} kvipaḥ tarhi adarśanam . \\
\noindent \textbf{(P\_1,1.60)} tatra adarśanam lopaḥ iti lopasañjñā prāpnoti . \\
\noindent \textbf{(P\_1,1.60)} tatra kaḥ doṣaḥ . \\
\noindent \textbf{(P\_1,1.60)} tatra pratyayalakṣaṇapratiṣedhaḥ . \\
\noindent \textbf{(P\_1,1.60)} tatra pratyayalakṣaṇam kāryam prāpnoti . \\
\noindent \textbf{(P\_1,1.60)} tasya pratiṣedhaḥ vaktavyaḥ . \\
\noindent \textbf{(P\_1,1.60)} hrasvasya piti kṛti tuk bhavati iti tuk prāpnoti . \\
\noindent \textbf{(P\_1,1.60)} siddham tu prasaktādarśanasya lopasañjñitvāt . \\
\noindent \textbf{(P\_1,1.60)} siddham etat . \\
\noindent \textbf{(P\_1,1.60)} katham . \\
\noindent \textbf{(P\_1,1.60)} prasaktādarśanam lopasañjñam bhavati iti vaktavyam . \\
\noindent \textbf{(P\_1,1.60)} yadi prasaktādarśanam lopasañjñam bhavati iti ucyate grāmaṇīḥ , senānīḥ : atra vṛddhiḥ prāpnoti . \\
\noindent \textbf{(P\_1,1.60)} prasaktādarśanam lopasañjñam bhavati ṣaṣṭhīnirdiṣṭasya . \\
\noindent \textbf{(P\_1,1.60)} yadi ṣaṣṭhīnirdiṣṭasya iti ucyate cāhalope eva iti avadhāraṇe cādilope vibhāṣā iti atra lopasañjñā na prāpnoti . \\
\noindent \textbf{(P\_1,1.60)} atha prasaktādarśanam lopasañjñam bhavati iti ucyamāne katham iva etat sidhyati . \\
\noindent \textbf{(P\_1,1.60)} kaḥ śabdasya prasaṅgaḥ . \\
\noindent \textbf{(P\_1,1.60)} yatra gamyate ca arthaḥ na ca prayujyate . \\
\noindent \textbf{(P\_1,1.60)} astu tarhi prasaktādarśanam lopasañjñam bhavati iti eva . \\
\noindent \textbf{(P\_1,1.60)} katham grāmaṇīḥ , senānīḥ . \\
\noindent \textbf{(P\_1,1.60)} yaḥ atra aṇaḥ prasaṅgaḥ kvipā asau bādhyate . \\
\noindent \textbf{(P\_1,1.61)} pratyayagrahaṇam kimartham . \\
\noindent \textbf{(P\_1,1.61)} lumati pratyayagrahaṇam apratyayasañjñāpratiṣedhārtham . lumati pratyayagrahaṇam kriyate apratyayasya etāḥ sañjñāḥ mā bhūvan iti . \\
\noindent \textbf{(P\_1,1.61)} kim prayojanam . \\
\noindent \textbf{(P\_1,1.61)} prayojanam taddhitaluki kaṃsīyaparaśavyayoḥ luki ca goprakṛtinivṛttyartham . taddhitaluki gonivṛttyartham kaṃsīyaparaśavyayoḥ ca luki prakṛtinivṛttyartham . \\
\noindent \textbf{(P\_1,1.61)} luk taddhitaluki iti goḥ api luk prāpnoti . \\
\noindent \textbf{(P\_1,1.61)} pratyayagrahaṇāt na bhavati . \\
\noindent \textbf{(P\_1,1.61)} kaṃsīyaparaśavyayoḥ yañañau luk ca iti prakṛteḥ api luk prāpnoti . \\
\noindent \textbf{(P\_1,1.61)} pratyayagrahaṇāt na bhavati . \\
\noindent \textbf{(P\_1,1.61)} gonivṛttyarthena tāvat na arthaḥ . \\
\noindent \textbf{(P\_1,1.61)} yogavibhāgāt siddham . \\
\noindent \textbf{(P\_1,1.61)} yogavibhāgaḥ kariṣyate : goḥ upasarjanasya . \\
\noindent \textbf{(P\_1,1.61)} gontasya prātipadikasya upasarjanasya hrasvaḥ bhavati . \\
\noindent \textbf{(P\_1,1.61)} tataḥ striyāḥ . \\
\noindent \textbf{(P\_1,1.61)} strīpratyayāntasya prātipadikasya upasarjanasya hrasvaḥ bhavati . \\
\noindent \textbf{(P\_1,1.61)} tataḥ luk taddhitaluki iti . \\
\noindent \textbf{(P\_1,1.61)} striyāḥ iti vartate . \\
\noindent \textbf{(P\_1,1.61)} goḥ iti nivṛttam . \\
\noindent \textbf{(P\_1,1.61)} kaṃsīyaparaśavyayoḥ viśiṣṭanirdeśāt siddham . \\
\noindent \textbf{(P\_1,1.61)} kaṃsīyaparaśavyayoḥ api viśiṣṭanirdeśḥ kartavyaḥ : kaṃsīyaparaśavyayoḥ yañañau bhavataḥ chayatoḥ ca luk bhavati iti . \\
\noindent \textbf{(P\_1,1.61)} saḥ ca avaśyam viśiṣṭanirdeśaḥ kartavyaḥ kriyamāṇe api vai pratyayagrahaṇe ukārasaśabdayoḥ mā bhūt iti : kameḥ saḥ kaṃsaḥ . \\
\noindent \textbf{(P\_1,1.61)} parān śṛṇāti iti paraśuḥ iti . \\
\noindent \textbf{(P\_1,1.61)} na eṣaḥ doṣaḥ . \\
\noindent \textbf{(P\_1,1.61)} uṇādayaḥ avyutpannāni prātipadikāni . \\
\noindent \textbf{(P\_1,1.61)} saḥ eṣaḥ ananyārthaḥ viśiṣṭanirdeśaḥ kartavyaḥ pratyayagrahaṇam vā kartavyam . \\
\noindent \textbf{(P\_1,1.61)} uktam vā . \\
\noindent \textbf{(P\_1,1.61)} kim uktam . \\
\noindent \textbf{(P\_1,1.61)} ṅyāpprātipadikagrahaṇam āṅgabhapadasañjñārtham yacchayoḥ ca lugartham iti . \\
\noindent \textbf{(P\_1,1.61)} ṣaṣṭhīnirdeśārtham tu . \\
\noindent \textbf{(P\_1,1.61)} ṣaṣṭhīnirdeśārtham tarhi pratyayagrahaṇam kartavyam . \\
\noindent \textbf{(P\_1,1.61)} anirdeśe hi ṣaṣṭhyarthāprasiddhiḥ . \\
\noindent \textbf{(P\_1,1.61)} akriyamāṇe hi pratyayagrahaṇe ṣaṣṭhyarthasya aprasiddhiḥ syāt . \\
\noindent \textbf{(P\_1,1.61)} kasya . \\
\noindent \textbf{(P\_1,1.61)} sthāneyogatvasya . \\
\noindent \textbf{(P\_1,1.61)} kva punaḥ iha ṣaṣṭhīnirdeśārthena arthaḥ pratyayagrahaṇena yāvatā sarvatra eva ṣaṣṭhī uccāryate : aṇiñoḥ tadrājasya yañañoḥ śapaḥ iti . \\
\noindent \textbf{(P\_1,1.61)} iha na kā cit ṣaṣṭhī : janapade lup iti . \\
\noindent \textbf{(P\_1,1.61)} atra api prakṛtam pratyayagrahaṇam anuvartate . \\
\noindent \textbf{(P\_1,1.61)} kva prakṛtam . \\
\noindent \textbf{(P\_1,1.61)} pratyayaḥ paraḥ ca iti . \\
\noindent \textbf{(P\_1,1.61)} tat vai prathamānirdiṣṭam ṣaṣṭhīnirdiṣṭena ca iha arthaḥ . \\
\noindent \textbf{(P\_1,1.61)} ṅyāpprātipadikāt iti eṣā pañcamī pratyayaḥ iti prathamāyāḥ ṣaṣṭhīm prakalpayiṣyati tasmāt iti uttarasya . \\
\noindent \textbf{(P\_1,1.61)} pratyayavidhiḥ ayam . \\
\noindent \textbf{(P\_1,1.61)} na ca pratyayavidhau pañcamyaḥ prakalpikāḥ bhavanti . \\
\noindent \textbf{(P\_1,1.61)} na ayam pratyayavidhiḥ . \\
\noindent \textbf{(P\_1,1.61)} vihitaḥ pratyayaḥ prakṛtaḥ ca anuvartate . \\
\noindent \textbf{(P\_1,1.61)} sarvādeśārtham vā vacanaprāmāṇyāt . \\
\noindent \textbf{(P\_1,1.61)} sarvādeśārtham tarhi pratyayagrahaṇam kartavyam . \\
\noindent \textbf{(P\_1,1.61)} lukślulupaḥ sarvādeśāḥ yathā syuḥ . \\
\noindent \textbf{(P\_1,1.61)} atha kriyamāṇe api pratyayagrahaṇe katham iva lukślulupaḥ sarvādeśāḥ labhyāḥ . \\
\noindent \textbf{(P\_1,1.61)} vacanaprāmāṇyāt : pratyayagrahaṇasāmāṛthyāt . \\
\noindent \textbf{(P\_1,1.61)} etat api na asti prayojanam . \\
\noindent \textbf{(P\_1,1.61)} ācāryapravṛttiḥ jñāpayati lukślulupaḥ sarvādeśāḥ bhavanti iti yat ayam luk vā duhadihalihaguhām ātmanepade dantye iti lope kṛte lukam śāsti . \\
\noindent \textbf{(P\_1,1.61)} uttarārtham tu . \\
\noindent \textbf{(P\_1,1.61)} uttarārtham tarhi pratyayagrahaṇam kartavyam . \\
\noindent \textbf{(P\_1,1.61)} na kartavyam. kriyate tatra eva : pratyayalope pratyayalakṣaṇam iti . \\
\noindent \textbf{(P\_1,1.61)} dvitīyam kartavyam kṛtsnapratyayalope pratyayalakṣaṇam yathā syāt . \\
\noindent \textbf{(P\_1,1.61)} ekadeśalope mā bhūt iti : āghnīta sam rāyaspoṣeṇa gmīya iti . \\
\noindent \textbf{(P\_1,1.62.1)} pratyayagrahaṇam kimartham. lope pratyayalakṣaṇam iti iyati ucyamāne saurathī vahatī iti gurūpottamalakṣaṇaḥ ṣyaṅ prasajyeta . \\
\noindent \textbf{(P\_1,1.62.1)} na eṣaḥ doṣaḥ . \\
\noindent \textbf{(P\_1,1.62.1)} na evam vijñāyate : lope pratyayalakṣaṇam pratyayasya prādurbhāvaḥ iti . \\
\noindent \textbf{(P\_1,1.62.1)} katham tarhi . \\
\noindent \textbf{(P\_1,1.62.1)} pratyayaḥ lakṣaṇam yasya kāryasya tat lupte api bhavati iti . \\
\noindent \textbf{(P\_1,1.62.1)} idam tarhi prayojanam : sati pratyaye yat prāpnoti tat pratyayalakṣanena yathā syāt . \\
\noindent \textbf{(P\_1,1.62.1)} lopottarakālam yat prāpnoti tat pratyayalakṣaṇena mā bhūt iti . \\
\noindent \textbf{(P\_1,1.62.1)} kim prayojanam . \\
\noindent \textbf{(P\_1,1.62.1)} grāmaṇikulam , senānikulam : auttarapadike hrasvatve kṛte hrasvasya piti kṛti tuk bhavati iti tuk prāpnoti . \\
\noindent \textbf{(P\_1,1.62.1)} saḥ mā bhūt iti . \\
\noindent \textbf{(P\_1,1.62.1)} yadi tarhi yat sati pratyaye prāpnoti tat pratyayalakṣanena bhavati . \\
\noindent \textbf{(P\_1,1.62.1)} lopottarakālam yat prāpnoti tat na bhavati jagat , janagat iti atra tuk na prāpnoti . \\
\noindent \textbf{(P\_1,1.62.1)} lopottarakalaḥ hi atra tuk āgamaḥ . \\
\noindent \textbf{(P\_1,1.62.1)} tasmāt na arthaḥ evamarthena pratyayagrahaṇena . \\
\noindent \textbf{(P\_1,1.62.1)} kasmāt na bhavati grāmaṇikulam , senānikulam . \\
\noindent \textbf{(P\_1,1.62.1)} bahiraṅgam hrasvatvam . \\
\noindent \textbf{(P\_1,1.62.1)} antaraṅgaḥ tuk . \\
\noindent \textbf{(P\_1,1.62.1)} asiddham bahiraṅgam antaraṅge . \\
\noindent \textbf{(P\_1,1.62.1)} idam tarhi prayojanam : kṛtsnapratyayalope pratyayalakṣaṇam yathā syāt . \\
\noindent \textbf{(P\_1,1.62.1)} ekadeśalope mā bhūt iti : āghnīta sam rāyaspoṣeṇa gmīya iti . \\
\noindent \textbf{(P\_1,1.62.1)} pūrvasmin api yoge pratyayagrahaṇasya etat prayojanam uktam. anyatarat śakyam akartum . \\
\noindent \textbf{(P\_1,1.62.1)} atha dvitīyam pratyayagrahaṇam kimartham . \\
\noindent \textbf{(P\_1,1.62.1)} pratyayalakṣaṇam yathā syāt varṇalakṣaṇam mā bhūt iti : gave hitam gohitam , rāyaḥ kulam raikulam iti . \\
\noindent \textbf{(P\_1,1.62.2)} kimartham punaḥ idam ucyate . \\
\noindent \textbf{(P\_1,1.62.2)} pratyayalope pratyayalakṣaṇavacanam sadanvākhyānāt śāstrasya . \\
\noindent \textbf{(P\_1,1.62.2)} pratyayalope pratyayalakṣaṇam iti ucyate sadanvākhyānāt śāstrasya . \\
\noindent \textbf{(P\_1,1.62.2)} sat śāstreṇa anvākhyāyate sataḥ vā śāstram anvyākhāyakam bhavati . \\
\noindent \textbf{(P\_1,1.62.2)} sadanvākhyānāt śāstrasya ugidacām sarvanāmasthāne adhātoḥ iti iha : eva syāt gomantau yavamantau . \\
\noindent \textbf{(P\_1,1.62.2)} gomān yavamān iti atra na syāt . \\
\noindent \textbf{(P\_1,1.62.2)} iṣyate ca syāt iti . \\
\noindent \textbf{(P\_1,1.62.2)} tat ca antareṇa yatnam na sidhyati . \\
\noindent \textbf{(P\_1,1.62.2)} ataḥ pratyayalope pratyayalakṣaṇavacanam . \\
\noindent \textbf{(P\_1,1.62.2)} evamartham idam ucyate . \\
\noindent \textbf{(P\_1,1.62.2)} asti prayojanam etat . \\
\noindent \textbf{(P\_1,1.62.2)} kim tarhi iti. luki upasaṅkhyānam . \\
\noindent \textbf{(P\_1,1.62.2)} luki upasaṅkhyānam kartavyam : pañca sapta . \\
\noindent \textbf{(P\_1,1.62.2)} kim punaḥ kāraṇam na sidhyati . \\
\noindent \textbf{(P\_1,1.62.2)} lope hi vidhānam . \\
\noindent \textbf{(P\_1,1.62.2)} lope hi pratyayalakṣaṇam vidhīyate . \\
\noindent \textbf{(P\_1,1.62.2)} tena luki na prāpnoti . \\
\noindent \textbf{(P\_1,1.62.2)} na vā adarśanasya lopasañjñitvāt . \\
\noindent \textbf{(P\_1,1.62.2)} na vā kartavyam . \\
\noindent \textbf{(P\_1,1.62.2)} kim kāraṇam . \\
\noindent \textbf{(P\_1,1.62.2)} adarśanasya lopasañjñitvāt . \\
\noindent \textbf{(P\_1,1.62.2)} adarśanam lopasañjñam iti ucyate . \\
\noindent \textbf{(P\_1,1.62.2)} lumatsañjñāḥ ca adarśanasya kriyante . \\
\noindent \textbf{(P\_1,1.62.2)} tena luki api bhaviṣyati . \\
\noindent \textbf{(P\_1,1.62.2)} yadi evam . \\
\noindent \textbf{(P\_1,1.62.2)} pratyayādarśanam tu lumatsañjñam . \\
\noindent \textbf{(P\_1,1.62.2)} pratyayādarśanam tu lumatsañjñam api prāpnoti . \\
\noindent \textbf{(P\_1,1.62.2)} tatra kaḥ doṣaḥ . \\
\noindent \textbf{(P\_1,1.62.2)} tatra luki śluvidhipratiṣedhaḥ . \\
\noindent \textbf{(P\_1,1.62.2)} tatra luki śluvidhiḥ api prāpnoti . \\
\noindent \textbf{(P\_1,1.62.2)} saḥ pratiṣedhyaḥ : atti hanti . \\
\noindent \textbf{(P\_1,1.62.2)} ślau iti dvirvacanam prāpnoti . \\
\noindent \textbf{(P\_1,1.62.2)} na vā pṛthaksañjñākaraṇāt . \\
\noindent \textbf{(P\_1,1.62.2)} na vā eṣaḥ doṣaḥ . \\
\noindent \textbf{(P\_1,1.62.2)} kim kāraṇam . \\
\noindent \textbf{(P\_1,1.62.2)} pṛthaksañjñākaraṇāt . \\
\noindent \textbf{(P\_1,1.62.2)} pṛthaksañjñākaraṇasāmarthyāt luki śluvidhiḥ na bhaviṣyati . \\
\noindent \textbf{(P\_1,1.62.2)} tasmāt adarśanasāmānyāt lopasañjñā lumatsañjñāḥ avagāhate . \\
\noindent \textbf{(P\_1,1.62.2)} yathā eva tarhi adarśanasāmānyāt lopasañjñā lumatsañjñāḥ avagāhate evam lumatsañjñāḥ api lopasañjñām avagāheran . \\
\noindent \textbf{(P\_1,1.62.2)} tatra kaḥ doṣaḥ . \\
\noindent \textbf{(P\_1,1.62.2)} agomatī gomatī sampannā gomatībhūtā : luk taddhitaluki iti ṅīpaḥ luk prasajyeta . \\
\noindent \textbf{(P\_1,1.62.2)} nanu ca atra api pṛthaksañjñākaraṇāt iti eva siddham . \\
\noindent \textbf{(P\_1,1.62.2)} yathā eva tarhi pṛthaksañjñākaraṇasāmarthyāt lumatsañjñāḥ lopasañjñām na avagāhante evam lopasañjñā api lumatsañjñāḥ na avagāheta . \\
\noindent \textbf{(P\_1,1.62.2)} tatra saḥ eva doṣaḥ : luki upasaṅkhyānam iti . \\
\noindent \textbf{(P\_1,1.62.2)} asti anyat lopasañjñāyāḥ pṛthaksañjñākaraṇe prayojanam . \\
\noindent \textbf{(P\_1,1.62.2)} kim. lumatsañjñāsu yat ucyate tat lopamātre mā bhūt iti . \\
\noindent \textbf{(P\_1,1.62.2)} lumati pratiṣedhāt vā . \\
\noindent \textbf{(P\_1,1.62.2)} atha vā yat ayam na lumatā aṅgasya iti pratṣedham śāsti tat jñāpayati ācāryaḥ bhavati luki pratyayalakṣaṇam iti . \\
\noindent \textbf{(P\_1,1.62.3)} sataḥ nimittābhāvāt padasañjñābhāvaḥ . san pratyayaḥ yeṣām kāryāṇām animittam : rājñaḥ puruṣaḥ iti saḥ luptaḥ api animittam syāt: rājapuruṣaḥ iti . \\
\noindent \textbf{(P\_1,1.62.3)} astu tasyāḥ animittam yā svādau padam iti padasañjñā yā tu subantam padam iti padasañjñā sā bhaviṣyati . \\
\noindent \textbf{(P\_1,1.62.3)} sati etatpratyaye āsīt : anayā bhaviṣyati anayā na bhaviṣyati iti . \\
\noindent \textbf{(P\_1,1.62.3)} lupte idānīm pratyaye yāvataḥ eva avadheḥ svādau padam iti padasañjñā tāvataḥ eva avadheḥ subantam padam iti . \\
\noindent \textbf{(P\_1,1.62.3)} asti ca pratyayalakṣaṇena yajādiparatā iti kṛtvā bhasañjñā prāpnoti . \\
\noindent \textbf{(P\_1,1.62.3)} tugdīrghatvayoḥ ca vipratiṣedhānupapattiḥ ekayogalakṣaṇatvāt parivīḥ iti . \\
\noindent \textbf{(P\_1,1.62.3)} tugdīrghatvayoḥ ca vipratiṣedhaḥ na upapadyate . \\
\noindent \textbf{(P\_1,1.62.3)} kva . \\
\noindent \textbf{(P\_1,1.62.3)} parivīḥ iti . \\
\noindent \textbf{(P\_1,1.62.3)} kim kāraṇam . \\
\noindent \textbf{(P\_1,1.62.3)} ekayogalakṣaṇatvāt . \\
\noindent \textbf{(P\_1,1.62.3)} ekayogalakṣaṇe tugdīrghatve . \\
\noindent \textbf{(P\_1,1.62.3)} iha lupte pratyaye sarvāṇi pratyayāśrayāṇi kāryāṇi paryavapannāni bhavanti . \\
\noindent \textbf{(P\_1,1.62.3)} tāni etāni pratyutthāpyante . \\
\noindent \textbf{(P\_1,1.62.3)} anena eva tuk anena eva ca dīrghatvam iti . \\
\noindent \textbf{(P\_1,1.62.3)} tat etat ekayogalakṣaṇam bhavati . \\
\noindent \textbf{(P\_1,1.62.3)} ekayogalakṣaṇāni ca na prakalpante . \\
\noindent \textbf{(P\_1,1.62.3)} siddham tu sthānisañjñānudeśāt ānyabhāvyasya . \\
\noindent \textbf{(P\_1,1.62.3)} siddham etat . \\
\noindent \textbf{(P\_1,1.62.3)} katham . \\
\noindent \textbf{(P\_1,1.62.3)} sthānisañjñā anyabhūtasya bhavati iti vaktavyam . \\
\noindent \textbf{(P\_1,1.62.3)} kim kṛtam bhavati . \\
\noindent \textbf{(P\_1,1.62.3)} sattāmātram anena kriyate . \\
\noindent \textbf{(P\_1,1.62.3)} yathāprāpte tugdīrghatve bhaviṣyataḥ . \\
\noindent \textbf{(P\_1,1.62.3)} tat vaktavyam bhavati . \\
\noindent \textbf{(P\_1,1.62.3)} yadi api etat ucyate atha vā etarhi sthānivadbhāvaḥ na ārabhyate . \\
\noindent \textbf{(P\_1,1.62.3)} sthānisañjñā anyabhūtasya analvidhau iti vakṣyāmi . \\
\noindent \textbf{(P\_1,1.62.3)} yadi evam āṅaḥ yamahanaḥ ātmanepadam bhavati iti hanteḥ eva syāt vadheḥ na syāt . \\
\noindent \textbf{(P\_1,1.62.3)} na hi kā cit hanteḥ sañjñā asti yā vadheḥ atidiśyeta . \\
\noindent \textbf{(P\_1,1.62.3)} hanteḥ api sañjñā asti . \\
\noindent \textbf{(P\_1,1.62.3)} kā . \\
\noindent \textbf{(P\_1,1.62.3)} hantiḥ eva . \\
\noindent \textbf{(P\_1,1.62.3)} katham . \\
\noindent \textbf{(P\_1,1.62.3)} svam rūpam śabdasya aśabdasañjñā iti vacanāt svam rūpam śabdasya sañjñā bhavati iti hanteḥ api hantiḥ sañjñā bhaviṣyati . \\
\noindent \textbf{(P\_1,1.62.3)} bhasañjñāṅīpṣphagorātveṣu ca siddham . \\
\noindent \textbf{(P\_1,1.62.3)} bhasañjñāṅīpṣphagorātveṣu ca siddham bhavati . \\
\noindent \textbf{(P\_1,1.62.3)} bhasañjñā : rājñaḥ puruṣaḥ rājapuruṣaḥ . \\
\noindent \textbf{(P\_1,1.62.3)} pratyayalakṣaṇena yaci bham iti bhāsañjñā prāpnoti . \\
\noindent \textbf{(P\_1,1.62.3)} sthānisañjñā anyabhūtasya analvidhau iti vacanāt na bhavati . \\
\noindent \textbf{(P\_1,1.62.3)} ṅīp : citrāyām jātā citrā . \\
\noindent \textbf{(P\_1,1.62.3)} pratyayalakṣaṇena aṇantāt īkāraḥ prāpnoti . \\
\noindent \textbf{(P\_1,1.62.3)} sthānisañjñā anyabhūtasya analvidhau iti vacanāt na bhavati . \\
\noindent \textbf{(P\_1,1.62.3)} ṣpha : vataṅḍī . \\
\noindent \textbf{(P\_1,1.62.3)} pratyayalakṣaṇena yañantāt iti ṣphaḥ prāpnoti . \\
\noindent \textbf{(P\_1,1.62.3)} sthānisañjñā anyabhūtasya analvidhau iti vacanāt na bhavati . \\
\noindent \textbf{(P\_1,1.62.3)} goḥ ātvam . \\
\noindent \textbf{(P\_1,1.62.3)} gām icchati gavyati . \\
\noindent \textbf{(P\_1,1.62.3)} pratyayalakṣaṇena ami ā otaḥ amśasoḥ iti ātvam prāpnoti . \\
\noindent \textbf{(P\_1,1.62.3)} sthānisañjñā anyabhūtasya analvidhau iti vacanāt na bhavati . \\
\noindent \textbf{(P\_1,1.62.3)} tasya doṣaḥ ṅaunakāralopettvemvidhayaḥ . tasya etasya lakṣaṇasya doṣaḥ ṅaunakāralopaḥ . \\
\noindent \textbf{(P\_1,1.62.3)} ārdre carman lohite carman . \\
\noindent \textbf{(P\_1,1.62.3)} pratyayalakṣaṇena yaci bham iti bhasañjñā siddhā bhavati . \\
\noindent \textbf{(P\_1,1.62.3)} sthānisañjñā anyabhūtasya analvidhau iti vacanāt na prāpnoti . \\
\noindent \textbf{(P\_1,1.62.3)} ittvam : āśīḥ . \\
\noindent \textbf{(P\_1,1.62.3)} pratyayalakṣaṇena hali iti itvam siddham bhavati . \\
\noindent \textbf{(P\_1,1.62.3)} sthānisañjñā anyabhūtasya analvidhau iti vacanāt na prāpnoti . \\
\noindent \textbf{(P\_1,1.62.3)} im : atṛṇet . \\
\noindent \textbf{(P\_1,1.62.3)} pratyayalakṣaṇena hali iti ittvam siddham bhavati . \\
\noindent \textbf{(P\_1,1.62.3)} sthānisañjñā anyabhūtasya analvidhau iti vacanāt na prāpnoti . \\
\noindent \textbf{(P\_1,1.62.3)} sūtram ca bhidyate . \\
\noindent \textbf{(P\_1,1.62.3)} yathānyāsam eva astu . \\
\noindent \textbf{(P\_1,1.62.3)} nanu ca uktam sataḥ nimittābhāvāt padasañjñābhāvaḥ tugdīrghatvayoḥ ca vipratiṣedhānupapattiḥ ekayogalakṣaṇatvāt parivīḥ iti . \\
\noindent \textbf{(P\_1,1.62.3)} na eṣaḥ doṣaḥ . \\
\noindent \textbf{(P\_1,1.62.3)} vakṣyati atra parihāram . \\
\noindent \textbf{(P\_1,1.62.3)} iha api parivīḥ iti śāstraparavipratiṣedhena paratvāt dīrghatvam bhaviṣyati . \\
\noindent \textbf{(P\_1,1.62.4)} kāni punaḥ asya yogasya prayojanāni . \\
\noindent \textbf{(P\_1,1.62.4)} prayojanam apṛktaśilope num amāmau guṇavṛddhidīrghatvemaḍāṭśnamvidhayaḥ . \\
\noindent \textbf{(P\_1,1.62.4)} apṛktalope śilope ca kṛte num amāmau guṇavṛddhī dīrghatvam imaḍāṭau śnamvidhiḥ iti prayojanāni . \\
\noindent \textbf{(P\_1,1.62.4)} num : agne trī te vajinā trī sadhasthā , ta tā piṇḍānām . \\
\noindent \textbf{(P\_1,1.62.4)} num . \\
\noindent \textbf{(P\_1,1.62.4)} amāmau : he anaḍvan , anaḍvān . \\
\noindent \textbf{(P\_1,1.62.4)} guṇaḥ : adhok , aleṭ . \\
\noindent \textbf{(P\_1,1.62.4)} vṛddhiḥ : ni amārṭ . \\
\noindent \textbf{(P\_1,1.62.4)} dīrghatvam : agne trī te vajinā trī sadhasthā , ta tā piṇḍānām . \\
\noindent \textbf{(P\_1,1.62.4)} im : atṛṇeṭ . \\
\noindent \textbf{(P\_1,1.62.4)} aḍāṭau : adhok , aleṭ , aiyaḥ , aunaḥ . \\
\noindent \textbf{(P\_1,1.62.4)} śnamvidhiḥ : abhinaḥ atra , acchinaḥ atra . \\
\noindent \textbf{(P\_1,1.62.4)} apṛktaśilopayoḥ kṛtayoḥ ete vidhayaḥ na prāpnuvanti . \\
\noindent \textbf{(P\_1,1.62.4)} pratyayalakṣaṇena bhavanti . \\
\noindent \textbf{(P\_1,1.62.4)} na etāni santi prayojanāni . \\
\noindent \textbf{(P\_1,1.62.4)} sthānivadbhāvena api etāni siddhāni . \\
\noindent \textbf{(P\_1,1.62.4)} na sidhyanti . \\
\noindent \textbf{(P\_1,1.62.4)} ādeśaḥ sthānivat iti ucyate . \\
\noindent \textbf{(P\_1,1.62.4)} na ca lopaḥ ādeśaḥ . \\
\noindent \textbf{(P\_1,1.62.4)} lopaḥ api ādeśaḥ . \\
\noindent \textbf{(P\_1,1.62.4)} katham . \\
\noindent \textbf{(P\_1,1.62.4)} ādiśyate yaḥ saḥ ādeśaḥ . \\
\noindent \textbf{(P\_1,1.62.4)} lopaḥ api ādiśyate . \\
\noindent \textbf{(P\_1,1.62.4)} doṣaḥ khalu api syāt yadi lopaḥ na ādeśaḥ syāt . \\
\noindent \textbf{(P\_1,1.62.4)} iha acaḥ parasmin pūrvavidhau iti etasya bhūyiṣṭhāni lope udāharaṇāni tāni na syuḥ . \\
\noindent \textbf{(P\_1,1.62.4)} yatra tarhi sthānivadbhāvaḥ na asti tadartham ayam yogaḥ vaktavyaḥ . \\
\noindent \textbf{(P\_1,1.62.4)} kva ca sthānivadbhāvaḥ na asti . \\
\noindent \textbf{(P\_1,1.62.4)} yaḥ alvidhiḥ . \\
\noindent \textbf{(P\_1,1.62.4)} kim prayojanam . \\
\noindent \textbf{(P\_1,1.62.4)} prayojanam ṅaunakārlopettvemvidhayaḥ . \\
\noindent \textbf{(P\_1,1.62.4)} bhasañjñāṅīpṣphgorātveṣu doṣaḥ . \\
\noindent \textbf{(P\_1,1.62.4)} bhasañjñāṅīpṣphgorātveṣu doṣaḥ bhavati. bhasañjñāyām tāvat na doṣaḥ . \\
\noindent \textbf{(P\_1,1.62.4)} ācāryapravṛttiḥ jñāpayati na pratyayalakṣaṇena bhasañjñā bhavati iti yat ayam na ṅisambuddhyoḥ iti ṅau pratiṣedham śāsti . \\
\noindent \textbf{(P\_1,1.62.4)} ṅīpi api : na evam vijñāyate : aṇantāt akārāntāt . \\
\noindent \textbf{(P\_1,1.62.4)} katham tarhi . \\
\noindent \textbf{(P\_1,1.62.4)} aṇ yaḥ akāraḥ iti . \\
\noindent \textbf{(P\_1,1.62.4)} ṣphe api : na evam vijñāyate : yañantāt akārantāt iti . \\
\noindent \textbf{(P\_1,1.62.4)} katham tarhi . \\
\noindent \textbf{(P\_1,1.62.4)} yañ yaḥ akāraḥ iti . \\
\noindent \textbf{(P\_1,1.62.4)} goḥ ātve api : na evam vijñāyate : ami aci iti . \\
\noindent \textbf{(P\_1,1.62.4)} katham tarhi . \\
\noindent \textbf{(P\_1,1.62.4)} aci ami iti . \\
\noindent \textbf{(P\_1,1.62.4)} prayojanāni api tarhi tāni na santi . \\
\noindent \textbf{(P\_1,1.62.4)} yat tāvat ucyate ṅaunakārlopaḥ iti kriyate etat nyāse eva : na ṅisambuddhyoḥ iti . \\
\noindent \textbf{(P\_1,1.62.4)} ittvam api . \\
\noindent \textbf{(P\_1,1.62.4)} vakṣyati etat : śāsaḥ ittve āśāsaḥ kvau iti . \\
\noindent \textbf{(P\_1,1.62.4)} imvidhiḥ api : hali iti nivṛttam . \\
\noindent \textbf{(P\_1,1.62.4)} yadi hali iti nivṛttam tṛṇahāni atra api prāpnoti . \\
\noindent \textbf{(P\_1,1.62.4)} evam tarhi aci na iti api anuvartiṣyate . \\
\noindent \textbf{(P\_1,1.62.4)} na tarhi idānīm ayam yogaḥ vaktavyaḥ . \\
\noindent \textbf{(P\_1,1.62.4)} vaktavyaḥ ca . \\
\noindent \textbf{(P\_1,1.62.4)} kim prayojanam . \\
\noindent \textbf{(P\_1,1.62.4)} pratyayam gṛhītvā yat ucyate tat pratyayalakṣaṇena yathā syāt śabdam gṛhītvā yat ucyate tat pratyayalakṣaṇena mā bhūt iti . \\
\noindent \textbf{(P\_1,1.62.4)} kim prayojanam . \\
\noindent \textbf{(P\_1,1.62.4)} śobhanāḥ dṛṣadaḥ asya sudṛṣat brāhmaṇaḥ . \\
\noindent \textbf{(P\_1,1.62.4)} soḥ manasī* alomoṣasī* iti eṣaḥ svaraḥ mā bhūt iti . \\
\noindent \textbf{(P\_1,1.63.1)} lumati pratiṣedhe ekapadasvarasya upasaṅkhyānam . \\
\noindent \textbf{(P\_1,1.63.1)} lumati pratiṣedhe ekapadasvarasya upasaṅkhyānam kartavyam. ekapadasvare ca lumatā lupte pratyayalakṣaṇam na bhavati iti vaktavyam . \\
\noindent \textbf{(P\_1,1.63.1)} kim aviśeṣeṇa . \\
\noindent \textbf{(P\_1,1.63.1)} na iti āha . \\
\noindent \textbf{(P\_1,1.63.1)} sarvāmantritasijluksvaravarjam . \\
\noindent \textbf{(P\_1,1.63.1)} sarvasvaram āmantritasvaravam sijluksvaram ca varjayitvā . \\
\noindent \textbf{(P\_1,1.63.1)} sarvasvara : sarvastomaḥ , sarvapṛṣṭhaḥ : sarvasya supi iti ādyudāttatvam yathā syāt . \\
\noindent \textbf{(P\_1,1.63.1)} āmantritasvara : sarpiḥ āgaccha , sapta āgacchata : āmantritasya ca iti ādyudāttatvam yathā syāt . \\
\noindent \textbf{(P\_1,1.63.1)} sijluksvara : ma hi datām , ma hi dhatām : ādiḥ sicaḥ anyatarasyām iti eṣaḥ svaraḥ yathā syāt . \\
\noindent \textbf{(P\_1,1.63.1)} kim prayojanam . \\
\noindent \textbf{(P\_1,1.63.1)} prayojanam ñinikilluki svarāḥ . \\
\noindent \textbf{(P\_1,1.63.1)} ñinikitsvarāḥ luki prayojayanti . \\
\noindent \textbf{(P\_1,1.63.1)} gargaḥ , vatsaḥ , bidaḥ , urvaḥ , uṣṭragrīvaḥ , vāmarajjuḥ : ñniti iti ādyudāttatvam mā bhūt iti . \\
\noindent \textbf{(P\_1,1.63.1)} iha ca : atrayaḥ : kitaḥ iti antodāttatvam mā bhūt iti . \\
\noindent \textbf{(P\_1,1.63.1)} pathimathoḥ sarvanāmasthāne . \\
\noindent \textbf{(P\_1,1.63.1)} pathimathoḥ sarvanāmasthāne luki prayojanam . \\
\noindent \textbf{(P\_1,1.63.1)} pathipriyaḥ , mathipriyaḥ : pathimathoḥ sarvanāmasthāne iti eṣaḥ svaraḥ mā bhūt iti . \\
\noindent \textbf{(P\_1,1.63.1)} ahnaḥ ravidhau . \\
\noindent \textbf{(P\_1,1.63.1)} ahnaḥ ravidhau lumatā lupte pratyayalakṣaṇam na bhavati iti vaktavyam . \\
\noindent \textbf{(P\_1,1.63.1)} ahaḥ dadati , ahaḥ bhuṅkte : raḥ asupi iti pratyayalakṣaṇena pratiṣedhaḥ mā bhūt iti . \\
\noindent \textbf{(P\_1,1.63.2)} uttarapadatve ca apadādividhau . uttarapadatve ca apadādividhau lumatā lupte pratyayalakṣaṇam na bhavati iti vaktavyam . \\
\noindent \textbf{(P\_1,1.63.2)} paramavācā paramavāce paramagoduhā paramagoduhe paramaśvalihā paramaśvalihe : padasya iti pratyayalakṣaṇena kutvādīni mā bhūvan iti . \\
\noindent \textbf{(P\_1,1.63.2)} apadādividhau iti kimartham . \\
\noindent \textbf{(P\_1,1.63.2)} dadhisecau dadhisecaḥ : sātpadādyoḥ iti pratiṣedhaḥ yathā syāt . \\
\noindent \textbf{(P\_1,1.63.2)} yadi apadādividhau iti ucyate uttarapadādhikāraḥ na prakalpeta . \\
\noindent \textbf{(P\_1,1.63.2)} tatra kaḥ doṣaḥ . \\
\noindent \textbf{(P\_1,1.63.2)} karṇaḥ varṇalakṣaṇāt iti evamādiḥ vidhiḥ na sidhyati . \\
\noindent \textbf{(P\_1,1.63.2)} yadi punaḥ nalopādividhau plutyante lumatā lupte pratyayalakṣaṇam na bhavati iti ucyeta . \\
\noindent \textbf{(P\_1,1.63.2)} na evam śakyam . \\
\noindent \textbf{(P\_1,1.63.2)} iha hi : rājakumāryau rājakumāryaḥ iti śākalam prasajyeta . \\
\noindent \textbf{(P\_1,1.63.2)} na eṣaḥ doṣaḥ . \\
\noindent \textbf{(P\_1,1.63.2)} yat etat siti śākalam na iti etat pratyaye śākalam na iti vakṣyāmi . \\
\noindent \textbf{(P\_1,1.63.2)} yadi pratyaye śākalam na iti ucyate dadhi adhunā madhu adhunā : atra api na prasajyeta . \\
\noindent \textbf{(P\_1,1.63.2)} pratyaye śākalam na bhavati . \\
\noindent \textbf{(P\_1,1.63.2)} kasmin . \\
\noindent \textbf{(P\_1,1.63.2)} yasmāt yaḥ pratyayaḥ vihitaḥ iti . \\
\noindent \textbf{(P\_1,1.63.2)} iha tarhi paramadivā paramadive : diva ut iti uttvam prāpnoti iti. astu tarhi aviśeṣeṇa . \\
\noindent \textbf{(P\_1,1.63.2)} nanu ca uktam uttarapadādhikāraḥ na prakalpeta iti . \\
\noindent \textbf{(P\_1,1.63.2)} vacanāt uttarapadādhikāraḥ bhaviṣyati . \\
\noindent \textbf{(P\_1,1.63.2)} tat tarhi vaktavyam . \\
\noindent \textbf{(P\_1,1.63.2)} na vaktavyam . \\
\noindent \textbf{(P\_1,1.63.2)} anuvṛttiḥ kariṣyate . \\
\noindent \textbf{(P\_1,1.63.2)} idam asti : yasmāt pratayayavidhiḥ tadādi pratyaye aṅgam , suptiṅantam padam . \\
\noindent \textbf{(P\_1,1.63.2)} yasmāt suptiṅvidhiḥ tadādi subantam ca . \\
\noindent \textbf{(P\_1,1.63.2)} naḥ kye . \\
\noindent \textbf{(P\_1,1.63.2)} nāntam kye padasañjñam bhavati yasmāt kyavidhiḥ subantam ca . \\
\noindent \textbf{(P\_1,1.63.2)} siti ca . \\
\noindent \textbf{(P\_1,1.63.2)} siti ca pūrvam padasañjñam bhavati yasmāt sidvidhiḥ tadādi subantam ca . \\
\noindent \textbf{(P\_1,1.63.2)} svādiṣu asarvanāmasthāne . \\
\noindent \textbf{(P\_1,1.63.2)} svādiṣu asarvanāmasthāne pūrvam padasañjñam bhavati yasmāt svādividhiḥ tadādi subantam ca . \\
\noindent \textbf{(P\_1,1.63.2)} yaci bham . \\
\noindent \textbf{(P\_1,1.63.2)} yajādipratyaye pūrvam padasañjñam bhavati yasmāt yajādividhiḥ tadādi subantam ca . \\
\noindent \textbf{(P\_1,1.63.2)} iha tarhi : paramavāk : asarvanāmasthāne iti pratiṣedhaḥ prāpnoti . \\
\noindent \textbf{(P\_1,1.63.2)} astu tasyāḥ pratiṣedhaḥ yā svādau padam iti padasañjñā yā tu subantam padam iti padasañjñā sā bhaviṣyati . \\
\noindent \textbf{(P\_1,1.63.2)} sati etatpratyaye āsīt anayā bhaviṣyati anayā na bhaviṣyati iti . \\
\noindent \textbf{(P\_1,1.63.2)} lupte idānīm pratyaye yāvataḥ eva avadheḥ svādau padam iti padasañjñā tāvataḥ eva avadheḥ subantam padam iti . \\
\noindent \textbf{(P\_1,1.63.2)} asti ca pratyayalakṣaṇena sarvanāmasthānaparatā iti kṛtvā pratiṣedhāḥ ca balīyāṃsaḥ bhavanti iti pratiṣedhaḥ prāpnoti . \\
\noindent \textbf{(P\_1,1.63.2)} na apratiṣedhāt . \\
\noindent \textbf{(P\_1,1.63.2)} na ayam prasajyapratiṣedhaḥ : sarvanāmasthāne na iti . \\
\noindent \textbf{(P\_1,1.63.2)} kim tarhi . \\
\noindent \textbf{(P\_1,1.63.2)} paryudāsaḥ ayam : yat anyat sarvanāmasthānāt iti . \\
\noindent \textbf{(P\_1,1.63.2)} sarvanāmasthāne avyāpāraḥ . \\
\noindent \textbf{(P\_1,1.63.2)} yadi kena cit prāpnoti tena bhaviṣyati . \\
\noindent \textbf{(P\_1,1.63.2)} pūrveṇa ca prāpnoti . \\
\noindent \textbf{(P\_1,1.63.2)} aprāpteḥ vā . \\
\noindent \textbf{(P\_1,1.63.2)} atha vā anantarā ya prāptiḥ sā pratiṣidhyate . \\
\noindent \textbf{(P\_1,1.63.2)} kutaḥ etat . \\
\noindent \textbf{(P\_1,1.63.2)} anantarasya vidhiḥ vā bhavati pratiṣedhaḥ vā iti . \\
\noindent \textbf{(P\_1,1.63.2)} pūrvā prāptiḥ apratiṣiddhā tayā bhaviṣyati . \\
\noindent \textbf{(P\_1,1.63.2)} nanu ca iyam prāptiḥ pūrvām prāptim bādhate . \\
\noindent \textbf{(P\_1,1.63.2)} na utsahate pratiṣiddhā satī bādhitum . \\
\noindent \textbf{(P\_1,1.63.2)} yadi evam paramavācau paramavācaḥ iti suptiṅantam padam iti padasañjñā prāpnoti . \\
\noindent \textbf{(P\_1,1.63.2)} evam tarhi yogavibhāgaḥ kariṣyate . \\
\noindent \textbf{(P\_1,1.63.2)} svādiṣu pūrvam padasañjñam bhavati . \\
\noindent \textbf{(P\_1,1.63.2)} tataḥ sarvanāmasthāne ayaci pūrvam padasañjñam bhavati . \\
\noindent \textbf{(P\_1,1.63.2)} tataḥ bham . \\
\noindent \textbf{(P\_1,1.63.2)} bhasañjñam bhavati yajādau asarvanāmasthane iti . \\
\noindent \textbf{(P\_1,1.63.2)} yadi tarhi sau api padam bhavati , ecaḥ plutādhikāre padāntagrahaṇam codayiṣyati iha mā bhūt : bhadram karoṣi gauḥ iti , tasmin kriyamāṇe api bhaviṣyati . \\
\noindent \textbf{(P\_1,1.63.2)} vākyapadayoḥ antyasya iti evam tat . \\
\noindent \textbf{(P\_1,1.63.2)} iha tarhi : dadhisecau dadhisecaḥ : sātpadādyoḥ iti padādilakṣaṇaḥ pratiṣedhaḥ na prāpnoti . \\
\noindent \textbf{(P\_1,1.63.2)} mā bhūt evam : padasya ādiḥ padādiḥ , padādeḥ na iti . \\
\noindent \textbf{(P\_1,1.63.2)} katham tarhi . \\
\noindent \textbf{(P\_1,1.63.2)} padāt ādiḥ padādiḥ , padādeḥ na iti evam bhaviṣyati . \\
\noindent \textbf{(P\_1,1.63.2)} na evam śakyam . \\
\noindent \textbf{(P\_1,1.63.2)} iha api prasajyeta : ṛkṣu vākṣu tvakṣu kumārīṣu kiśorīṣu iti . \\
\noindent \textbf{(P\_1,1.63.2)} sātpratiṣedhaḥ jñāpakaḥ svādiṣu padatvena yeṣām padasañjñā na tebhyaḥ pratiṣedhaḥ bhavati iti . \\
\noindent \textbf{(P\_1,1.63.2)} iha tarhi : bahusecau , bahusecaḥ : bahuc ayam pratyayaḥ . \\
\noindent \textbf{(P\_1,1.63.2)} atra padāt ādiḥ padādiḥ , padādeḥ na iti ucyamāne api na sidhyati . \\
\noindent \textbf{(P\_1,1.63.2)} evam tarhi uttarapadatve ca padādividhau lumatā lupte pratyayalakṣaṇam bhavati iti vakṣyāmi . \\
\noindent \textbf{(P\_1,1.63.2)} tat niyamārtham bhaviṣyati : padādividhau eva na padāntavidhau iti . \\
\noindent \textbf{(P\_1,1.63.2)} katham bahusecau bahusecaḥ . \\
\noindent \textbf{(P\_1,1.63.2)} bahucpūrvasya ca padādividhau na padāntavidhau iti . \\
\noindent \textbf{(P\_1,1.63.2)} dvandve antyasya . \\
\noindent \textbf{(P\_1,1.63.2)} dvandve antyasyalumatā lupte pratyayalakṣaṇam na bhavati iti vaktavyam . \\
\noindent \textbf{(P\_1,1.63.2)} vāksraktvacam . \\
\noindent \textbf{(P\_1,1.63.3)} iha abhūvan iti pratyayalakṣaṇena jusbhāvaḥ prāpnoti . \\
\noindent \textbf{(P\_1,1.63.3)} sicaḥ usaḥ aprasaṅgaḥ ākāraprakaraṇāt . \\
\noindent \textbf{(P\_1,1.63.3)} sicaḥ usaḥ aprasaṅgaḥ . \\
\noindent \textbf{(P\_1,1.63.3)} kim kāraṇam . \\
\noindent \textbf{(P\_1,1.63.3)} ākāraprakaraṇāt . \\
\noindent \textbf{(P\_1,1.63.3)} ātaḥ iti etat niyamārtham bhaviṣyati : ātaḥ eva sijlugantāt na anyasmāt sijlugantāt iti . \\
\noindent \textbf{(P\_1,1.63.3)} iha : iti yuṣmatputraḥ dadāti , iti asmatputraḥ dadāti iti atra yuṣmadasmadoḥ ṣaṣṭhīcaturthīdvitīyāsthayoḥ vāmnāvau iti vāmnāvādayaḥ prāpnuvanti . \\
\noindent \textbf{(P\_1,1.63.3)} yuṣmadasmadoḥ sthagrahaṇāt . \\
\noindent \textbf{(P\_1,1.63.3)} sthagrahaṇam tatra kriyate . \\
\noindent \textbf{(P\_1,1.63.3)} tat śrūyamāṇavibhaktiviśeṣaṇam vijñāsyate . \\
\noindent \textbf{(P\_1,1.63.3)} asti anyat sthagrahaṇasya prayojanam . \\
\noindent \textbf{(P\_1,1.63.3)} kim . \\
\noindent \textbf{(P\_1,1.63.3)} savibhaktikasya vāmnāvādayaḥ yathā syuḥ iti . \\
\noindent \textbf{(P\_1,1.63.3)} na etat asti prayojanam . \\
\noindent \textbf{(P\_1,1.63.3)} padasya iti vartate vibhaktyantam ca padam . \\
\noindent \textbf{(P\_1,1.63.3)} tatra antareṇa api sthagrahaṇam savibhaktikasya eva grahaṇam bhaviṣyati . \\
\noindent \textbf{(P\_1,1.63.3)} bhavet siddham yatra vibhaktyantam padam . \\
\noindent \textbf{(P\_1,1.63.3)} yatra tu khalu vibhaktau padam tatra na sidhyati : grāmaḥ vām dīyate , grāmaḥ nau dīyate janapadaḥ vām dīyate , janapadaḥ nau dīyate . \\
\noindent \textbf{(P\_1,1.63.3)} sarvagrahaṇam api prakṛtam anuvartate . \\
\noindent \textbf{(P\_1,1.63.3)} tena savibhaktikasya eva bhaviṣyati . \\
\noindent \textbf{(P\_1,1.63.3)} iha : cakṣuṣkāmam yājayām cakāra iti tiṅ atiṅaḥ iti . \\
\noindent \textbf{(P\_1,1.63.3)} tasya ca nighātaḥ tasmāt ca anighātaḥ prāpnoti . \\
\noindent \textbf{(P\_1,1.63.3)} āmi lilopāt tasya ca anighātaḥ tasmāt ca nighātaḥ . \\
\noindent \textbf{(P\_1,1.63.3)} āmi lilopāt tasya ca anighātaḥ tasmāt ca nighātaḥ siddhaḥ bhaviṣyati . \\
\noindent \textbf{(P\_1,1.63.3)} aṅgādhikāre iṭaḥ vidhipratiṣedhau . \\
\noindent \textbf{(P\_1,1.63.3)} aṅgādhikāre iṭaḥ vidhipratiṣedhau na sidhyataḥ : jigamiṣa saṃvivṛtsa . \\
\noindent \textbf{(P\_1,1.63.3)} aṅgasya iti iṭaḥ vidhipratiṣedhau na prāpnutaḥ . \\
\noindent \textbf{(P\_1,1.63.3)} krameḥ dīrghatvam ca . \\
\noindent \textbf{(P\_1,1.63.3)} kim ca . \\
\noindent \textbf{(P\_1,1.63.3)} iṭaḥ ca vidhipratiṣedhau . \\
\noindent \textbf{(P\_1,1.63.3)} na iti āha . \\
\noindent \textbf{(P\_1,1.63.3)} adeśe ayam caḥ paṭhitaḥ . \\
\noindent \textbf{(P\_1,1.63.3)} krameḥ ca dīrghatvam : utkrāma saṅkrāma iti . \\
\noindent \textbf{(P\_1,1.63.4)} iha kim cit aṅgādhikāre lumatā lupte pratyayalakṣaṇena bhavati kim cit ca anyatra na bhavati . \\
\noindent \textbf{(P\_1,1.63.4)} yadi punaḥ na lumatā tasmin iti ucyeta . \\
\noindent \textbf{(P\_1,1.63.4)} atha na lumatā tasmin iti ucyamāne kim siddham etat bhavati iṭaḥ vidhipratiṣedhau krameḥ dīrghatvam ca . \\
\noindent \textbf{(P\_1,1.63.4)} bāḍham siddham . \\
\noindent \textbf{(P\_1,1.63.4)} na iṭaḥ ividhipratiṣedhau parasmaipadeṣu iti ucyate . \\
\noindent \textbf{(P\_1,1.63.4)} katham tarhi . \\
\noindent \textbf{(P\_1,1.63.4)} sakārādau iti . \\
\noindent \textbf{(P\_1,1.63.4)} tadviśeṣaṇam parasmaipadagrahaṇam . \\
\noindent \textbf{(P\_1,1.63.4)} na khalu api krameḥ dīrghatvam parasmaipadeṣu iti ucyate . \\
\noindent \textbf{(P\_1,1.63.4)} katham tarhi . \\
\noindent \textbf{(P\_1,1.63.4)} śiti iti . \\
\noindent \textbf{(P\_1,1.63.4)} tadviśeṣaṇam parasmaipadagrahaṇam . \\
\noindent \textbf{(P\_1,1.63.4)} na lumatā tasmin iti cet haniṇiṅādeśāḥ talope . \\
\noindent \textbf{(P\_1,1.63.4)} na lumatā tasmin iti cet haniṇiṅādeśāḥ talope na sidhyanti : avadhi bhavatā dasyuḥ , agāyi bhavatā grāmaḥ , adhyagāyi bhavatā anuvākaḥ . \\
\noindent \textbf{(P\_1,1.63.4)} talope kṛte luṅi iti haniṇiṅādeśāḥ na prāpnuvanti . \\
\noindent \textbf{(P\_1,1.63.4)} na eṣaḥ doṣaḥ . \\
\noindent \textbf{(P\_1,1.63.4)} na luṅi iti haniṇiṅādeśāḥ ucyante . \\
\noindent \textbf{(P\_1,1.63.4)} kim tarhi . \\
\noindent \textbf{(P\_1,1.63.4)} ārdhadhātuke iti . \\
\noindent \textbf{(P\_1,1.63.4)} tadviśeṣaṇam luṅgrahaṇam . \\
\noindent \textbf{(P\_1,1.63.4)} iha ca : sarvastomaḥ , sarvapṛṣṭhaḥ sarvasya supi iti ādyudāttatvam na prāpnoti . \\
\noindent \textbf{(P\_1,1.63.4)} tat ca api vaktavyam . \\
\noindent \textbf{(P\_1,1.63.4)} na vaktavyam . \\
\noindent \textbf{(P\_1,1.63.4)} na lumatā aṅgasya iti eva siddham . \\
\noindent \textbf{(P\_1,1.63.4)} katham . \\
\noindent \textbf{(P\_1,1.63.4)} na lumatā lupte aṅgādhikāraḥ pratinirdiśyate . \\
\noindent \textbf{(P\_1,1.63.4)} kim tarhi . \\
\noindent \textbf{(P\_1,1.63.4)} yaḥ asau lumatā lupyate tasmin yat aṅgam tasya yat kāryam tat na bhavati . \\
\noindent \textbf{(P\_1,1.63.4)} evam api sarvasvaraḥ na sidhyati . \\
\noindent \textbf{(P\_1,1.63.4)} kartavyaḥ atra yatnaḥ . \\
\noindent \textbf{(P\_1,1.65.1)} kim idam algrahaṇam antyaviśeṣaṇam . \\
\noindent \textbf{(P\_1,1.65.1)} evam bhavitum arhati . \\
\noindent \textbf{(P\_1,1.65.1)} upadhāsañjñāyām algrahaṇam antyanirdeśaḥ cet saṅghātapratiṣedhaḥ . \\
\noindent \textbf{(P\_1,1.65.1)} upadhāsañjñāyām algrahaṇam antyanirdeśaḥ cet saṅghātasya pratiṣedhaḥ vaktavyaḥ . \\
\noindent \textbf{(P\_1,1.65.1)} saṅghātasya upadhāsañjñā prāpnoti . \\
\noindent \textbf{(P\_1,1.65.1)} tatra kaḥ doṣaḥ . \\
\noindent \textbf{(P\_1,1.65.1)} śāsaḥ it aṅhaloḥ : śiṣṭvā śiṣṭaḥ : saṅghātasya ittvam prāpnoti . \\
\noindent \textbf{(P\_1,1.65.1)} yadi punaḥ al antyāt iti ucyeta . \\
\noindent \textbf{(P\_1,1.65.1)} evam api antyaḥ aviśeṣitaḥ bhavati . \\
\noindent \textbf{(P\_1,1.65.1)} tatra kaḥ doṣaḥ . \\
\noindent \textbf{(P\_1,1.65.1)} saṅghātāt api pūrvasya upadhāsañjñā prasajyeta . \\
\noindent \textbf{(P\_1,1.65.1)} tatra kaḥ doṣaḥ . \\
\noindent \textbf{(P\_1,1.65.1)} śāsaḥ it aṅhaloḥ : śiṣṭaḥ , śiṣṭavān : śakārasya ittvam prasajyeta . \\
\noindent \textbf{(P\_1,1.65.1)} sūtram ca bhidyate . \\
\noindent \textbf{(P\_1,1.65.1)} yathānyāsam eva astu . \\
\noindent \textbf{(P\_1,1.65.1)} nanu ca uktam upadhāsañjñāyām algrahaṇam antyanirdeśaḥ cet saṅghātapratiṣedhaḥ iti . \\
\noindent \textbf{(P\_1,1.65.1)} na eṣaḥ doṣaḥ . \\
\noindent \textbf{(P\_1,1.65.1)} antyavijñānāt siddham . \\
\noindent \textbf{(P\_1,1.65.1)} siddham etat . \\
\noindent \textbf{(P\_1,1.65.1)} katham . \\
\noindent \textbf{(P\_1,1.65.1)} alaḥ antyasya vidhayaḥ bhavanti iti antyasya bhaviṣyati . \\
\noindent \textbf{(P\_1,1.65.2)} antyavijñānāt siddham iti cet na anarthake alontyavidhiḥ anabhyāsavikāre . \\
\noindent \textbf{(P\_1,1.65.2)} antyavijñānāt siddham iti cet tat na . \\
\noindent \textbf{(P\_1,1.65.2)} kim kāraṇam . \\
\noindent \textbf{(P\_1,1.65.2)} na anarthake alontyavidhiḥ anabhyāsavikāre . \\
\noindent \textbf{(P\_1,1.65.2)} anarthake alontyavidhiḥ na iti eṣā paribhāṣā kartavyā . \\
\noindent \textbf{(P\_1,1.65.2)} kim aviśeṣeṇa . \\
\noindent \textbf{(P\_1,1.65.2)} na iti āha . \\
\noindent \textbf{(P\_1,1.65.2)} anabhyāsavikāre . \\
\noindent \textbf{(P\_1,1.65.2)} abhyāsavikārān varjayitvā . \\
\noindent \textbf{(P\_1,1.65.2)} bhṛñām it , artipipartyoḥ ca iti . \\
\noindent \textbf{(P\_1,1.65.2)} kāni etasyāḥ paribhāṣāyāḥ prayojanāni . \\
\noindent \textbf{(P\_1,1.65.2)} prayojanam avyaktānukaraṇasya ataḥ itau . \\
\noindent \textbf{(P\_1,1.65.2)} antyasya prāpnoti . \\
\noindent \textbf{(P\_1,1.65.2)} anarthake alontyavidhiḥ na bhavati iti na doṣaḥ bhavati . \\
\noindent \textbf{(P\_1,1.65.2)} na etat asti prayojanam . \\
\noindent \textbf{(P\_1,1.65.2)} ācāryapravṛttiḥ jñāpayati na antyasya pararūpam bhavati iti yat ayam na āmreḍitasya antyasya tu vā iti āha . \\
\noindent \textbf{(P\_1,1.65.2)} ghvasoḥ et hau abhyāsalopaḥ ca . \\
\noindent \textbf{(P\_1,1.65.2)} ghvasoḥ et hau abhyāsalopaḥ ca iti antyasya prāpnoti . \\
\noindent \textbf{(P\_1,1.65.2)} anarthake alontyavidhiḥ na bhavati iti na doṣaḥ bhavati . \\
\noindent \textbf{(P\_1,1.65.2)} etat api na asti prayojanam . \\
\noindent \textbf{(P\_1,1.65.2)} punarlopavacanasāmarthyāt sarvasya bhaviṣyati . \\
\noindent \textbf{(P\_1,1.65.2)} atha vā śit lopaḥ kariṣyate . \\
\noindent \textbf{(P\_1,1.65.2)} saḥ śit sarvasya iti sarvādeśaḥ bhaviṣyati . \\
\noindent \textbf{(P\_1,1.65.2)} saḥ tarhi śakāraḥ kartavyaḥ . \\
\noindent \textbf{(P\_1,1.65.2)} na kartavyaḥ . \\
\noindent \textbf{(P\_1,1.65.2)} kriyate nyāse eva . \\
\noindent \textbf{(P\_1,1.65.2)} dviśakārakaḥ nirdeśaḥ : ghvasoḥ et hau abhyāsalopaśśca iti . \\
\noindent \textbf{(P\_1,1.65.2)} āpi lopaḥ akaḥ anaci . \\
\noindent \textbf{(P\_1,1.65.2)} tiṣṭhati sūtram . \\
\noindent \textbf{(P\_1,1.65.2)} anyathā vyākhyāyate : āpi hali lopaḥ iti antyasya prāpnoti . \\
\noindent \textbf{(P\_1,1.65.2)} anarthake alontyavidhiḥ na bhavati iti na doṣaḥ bhavati . \\
\noindent \textbf{(P\_1,1.65.2)} etat api na asti prayojanam . \\
\noindent \textbf{(P\_1,1.65.2)} anaḥ eva lopam vakṣyāmi . \\
\noindent \textbf{(P\_1,1.65.2)} tat anaḥ grahaṇam kartavyam . \\
\noindent \textbf{(P\_1,1.65.2)} na kartavyam . \\
\noindent \textbf{(P\_1,1.65.2)} prakṛtam anuvartate . \\
\noindent \textbf{(P\_1,1.65.2)} kva prakṛtam . \\
\noindent \textbf{(P\_1,1.65.2)} an āpi akaḥ iti . \\
\noindent \textbf{(P\_1,1.65.2)} tat vai prathamānirdiṣṭam . \\
\noindent \textbf{(P\_1,1.65.2)} ṣaṣṭhīnirdiṣtena ca iha arthaḥ . \\
\noindent \textbf{(P\_1,1.65.2)} hali iti eṣā saptamī an iti prathamāyāḥ ṣaṣṭhīm prakalpayiṣyati : tasmin iti nirdiṣṭe pūrvasya iti . \\
\noindent \textbf{(P\_1,1.65.2)} atra lopaḥ abhyāsasya . \\
\noindent \textbf{(P\_1,1.65.2)} atra lopaḥ abhyāsasya iti antyasya prāpnoti . \\
\noindent \textbf{(P\_1,1.65.2)} anarthake alontyavidhiḥ na bhavati iti na doṣaḥ bhavati . \\
\noindent \textbf{(P\_1,1.65.2)} etat api na asti prayojanam . \\
\noindent \textbf{(P\_1,1.65.2)} atragrahaṇasāmarthyāt sarvasya bhaviṣyati . \\
\noindent \textbf{(P\_1,1.65.2)} asti anyat atragrahaṇasya prayojanam . \\
\noindent \textbf{(P\_1,1.65.2)} kim . \\
\noindent \textbf{(P\_1,1.65.2)} sanadhikāraḥ apekṣyate , iha mā bhūt : dadhau dadau . \\
\noindent \textbf{(P\_1,1.65.2)} antareṇa api atragrahaṇam sanadhikāram apekṣiṣyāmahe . \\
\noindent \textbf{(P\_1,1.65.2)} san tarhi sakārādiḥ apekṣyate sani sakārādau iti , iha mā bhūt : jijñāpayiṣati . \\
\noindent \textbf{(P\_1,1.65.2)} antareṇa api atragrahaṇam sanam sakārādim apekṣiṣyāmahe . \\
\noindent \textbf{(P\_1,1.65.2)} prakṛtayaḥ tarhi apekṣyante . \\
\noindent \textbf{(P\_1,1.65.2)} etāsām prakṛtīnām lopaḥ yathā syāt , iha mā bhūt : pipakṣati yiyakṣati . \\
\noindent \textbf{(P\_1,1.65.2)} antareṇa api atragrahaṇam etāḥ prakṛtīḥ apekṣiṣyāmahe . \\
\noindent \textbf{(P\_1,1.65.2)} viṣayaḥ tarhi apekṣyate . \\
\noindent \textbf{(P\_1,1.65.2)} mucaḥ akarmakasya guṇaḥ vā iti iha mā bhūt : mumukṣati gām iti . \\
\noindent \textbf{(P\_1,1.65.2)} antareṇa api atragrahaṇam viṣayam apekṣiṣyāmahe . \\
\noindent \textbf{(P\_1,1.65.2)} katham . \\
\noindent \textbf{(P\_1,1.65.2)} akarmakasya iti ucyate . \\
\noindent \textbf{(P\_1,1.65.2)} tena yatra eva ayam muciḥ akarmakaḥ tatra eva bhaviṣyati . \\
\noindent \textbf{(P\_1,1.65.2)} tasmāt na arthaḥ anayā paribhāṣayā . \\
\noindent \textbf{(P\_1,1.65.3)} alaḥ antyāt pūrvaḥ al upadhā iti vā . \\
\noindent \textbf{(P\_1,1.65.3)} atha vā vyaktam eva pathitavyam alaḥ antyāt pūrvaḥ al upadhāsañjñaḥ bhavati iti . \\
\noindent \textbf{(P\_1,1.65.3)} tat tarhi vaktavyam . \\
\noindent \textbf{(P\_1,1.65.3)} na vaktavyam . \\
\noindent \textbf{(P\_1,1.65.3)} avacanāt lokavijñānāt siddham . \\
\noindent \textbf{(P\_1,1.65.3)} antareṇa api vacanam lokavijñānāt siddham etat . \\
\noindent \textbf{(P\_1,1.65.3)} katham . \\
\noindent \textbf{(P\_1,1.65.3)} loke amīṣām brāhmaṇānām antyāt pūrvaḥ ānīyatām iti ukte yathājātīyakaḥ antyaḥ tathājātīyakaḥ antyāt pūrvaḥ ānīyate . \\
\noindent \textbf{(P\_1,1.66-67.1)} KA\_I,171.18-172.17 Ro\_I,507-511 {1/42}     kim udāharaṇam . \\
\noindent \textbf{(P\_1,1.66-67.1)} KA\_I,171.18-172.17 Ro\_I,507-511 {2/42}     iha tāvat : tasmin iti nirdiṣṭe pūrvasya iti : ikaḥ yaṇ aci : dadhi atra madhu atra . \\
\noindent \textbf{(P\_1,1.66-67.1)} KA\_I,171.18-172.17 Ro\_I,507-511 {3/42}     iha : tasmāt iti uttarasya iti : dvayantarupasargebhyaḥ apaḥ īt : dvīpam antarīpam samīpam . \\
\noindent \textbf{(P\_1,1.66-67.1)} KA\_I,171.18-172.17 Ro\_I,507-511 {4/42}     anyathājātīyakena śabdena nirdeśaḥ kriyate anyathājātīyakaḥ udāhriyate . \\
\noindent \textbf{(P\_1,1.66-67.1)} KA\_I,171.18-172.17 Ro\_I,507-511 {5/42}     kim punaḥ udāharaṇam . \\
\noindent \textbf{(P\_1,1.66-67.1)} KA\_I,171.18-172.17 Ro\_I,507-511 {6/42}     iha tāvat : tasmin iti nirdiṣṭe pūrvasya iti : tasmin aṇi ca yuṣmākāsmākau iti . \\
\noindent \textbf{(P\_1,1.66-67.1)} KA\_I,171.18-172.17 Ro\_I,507-511 {7/42}     tasmāt iti uttarasya iti : tasmāt śasaḥ naḥ puṃsi iti . \\
\noindent \textbf{(P\_1,1.66-67.1)} KA\_I,171.18-172.17 Ro\_I,507-511 {8/42}     idam ca api udāharaṇam : ikaḥ yaṇ aci dvyantarupasargebhyaḥ apaḥ īt iti . \\
\noindent \textbf{(P\_1,1.66-67.1)} KA\_I,171.18-172.17 Ro\_I,507-511 {9/42}     katham . \\
\noindent \textbf{(P\_1,1.66-67.1)} KA\_I,171.18-172.17 Ro\_I,507-511 {10/42}     sarvanāmnā ayam nirdeśaḥ kriyate sarvanāma ca sāmānyavāci . \\
\noindent \textbf{(P\_1,1.66-67.1)} KA\_I,171.18-172.17 Ro\_I,507-511 {11/42}     tatra sāmānye nirdiṣṭe viśeṣāḥ api udāharaṇāni bhavanti . \\
\noindent \textbf{(P\_1,1.66-67.1)} KA\_I,171.18-172.17 Ro\_I,507-511 {12/42}     kim punaḥ sāmānyam kaḥ vā viśeṣaḥ . \\
\noindent \textbf{(P\_1,1.66-67.1)} KA\_I,171.18-172.17 Ro\_I,507-511 {13/42}     gauḥ sāmānyam kṛṣṇaḥ viśeṣaḥ . \\
\noindent \textbf{(P\_1,1.66-67.1)} KA\_I,171.18-172.17 Ro\_I,507-511 {14/42}     na tarhi idānīm kṛṣṇaḥ sāmānyam bhavati gauḥ viśeṣaḥ bhavati . \\
\noindent \textbf{(P\_1,1.66-67.1)} KA\_I,171.18-172.17 Ro\_I,507-511 {15/42}     bhavati ca . \\
\noindent \textbf{(P\_1,1.66-67.1)} KA\_I,171.18-172.17 Ro\_I,507-511 {16/42}     yadi sāmānyam api viśeṣaḥ viśeṣaḥ api sāmānyam sāmānyaviśeṣau na prakalpete . \\
\noindent \textbf{(P\_1,1.66-67.1)} KA\_I,171.18-172.17 Ro\_I,507-511 {17/42}     prakalpete ca . \\
\noindent \textbf{(P\_1,1.66-67.1)} KA\_I,171.18-172.17 Ro\_I,507-511 {18/42}     katham . \\
\noindent \textbf{(P\_1,1.66-67.1)} KA\_I,171.18-172.17 Ro\_I,507-511 {19/42}     vivakṣātaḥ . \\
\noindent \textbf{(P\_1,1.66-67.1)} KA\_I,171.18-172.17 Ro\_I,507-511 {20/42}     yadā asya gauḥ sāmānyena vivakṣitaḥ bhavati kṛṣṇaḥ viśeṣatvena tadā gauḥ sāmānyam kṛṣṇaḥ viśeṣaḥ . \\
\noindent \textbf{(P\_1,1.66-67.1)} KA\_I,171.18-172.17 Ro\_I,507-511 {21/42}     yadā kṛṣṇaḥ sāmānyena vivakṣitaḥ bhavati gauḥ viśeṣatvena tadā kṛṣṇaḥ sāmānyam kṛṣṇaḥ viśeṣaḥ . \\
\noindent \textbf{(P\_1,1.66-67.1)} KA\_I,171.18-172.17 Ro\_I,507-511 {22/42}     aparaḥ āha : prakalpete ca . \\
\noindent \textbf{(P\_1,1.66-67.1)} KA\_I,171.18-172.17 Ro\_I,507-511 {23/42}     katham . \\
\noindent \textbf{(P\_1,1.66-67.1)} KA\_I,171.18-172.17 Ro\_I,507-511 {24/42}     pitāputravat . \\
\noindent \textbf{(P\_1,1.66-67.1)} KA\_I,171.18-172.17 Ro\_I,507-511 {25/42}     tat yathā saḥ eva kam cit prati pitā bhavati kam cit prati putraḥ bhavati evam iha api saḥ eva kam cit prati sāmānyam kam cit prati viśeṣaḥ . \\
\noindent \textbf{(P\_1,1.66-67.1)} KA\_I,171.18-172.17 Ro\_I,507-511 {26/42}     ete khalu api nairdeśikānām vārttatarakāḥ bhavanti ye sarvanāmnā nirdeśāḥ kriyante . \\
\noindent \textbf{(P\_1,1.66-67.1)} KA\_I,171.18-172.17 Ro\_I,507-511 {27/42}     etaiḥ hi bahutarakam vyāpyate . \\
\noindent \textbf{(P\_1,1.66-67.1)} KA\_I,171.18-172.17 Ro\_I,507-511 {28/42}     atha kimartham upasargeṇa nirdeśaḥ kriyate . \\
\noindent \textbf{(P\_1,1.66-67.1)} KA\_I,171.18-172.17 Ro\_I,507-511 {29/42}     śabde saptamyā nirdiṣṭe pūrvasya kāryam yathā syāt arthe mā bhūt : janapade atiśāyane iti . \\
\noindent \textbf{(P\_1,1.66-67.1)} KA\_I,171.18-172.17 Ro\_I,507-511 {30/42}     kim gatam etat upasargeṇa āhosvit śabdādhikyāt arthādhikyam . \\
\noindent \textbf{(P\_1,1.66-67.1)} KA\_I,171.18-172.17 Ro\_I,507-511 {31/42}     gatam iti āha . \\
\noindent \textbf{(P\_1,1.66-67.1)} KA\_I,171.18-172.17 Ro\_I,507-511 {32/42}     katham . \\
\noindent \textbf{(P\_1,1.66-67.1)} KA\_I,171.18-172.17 Ro\_I,507-511 {33/42}     niḥ ayam bahirbhāve vartate . \\
\noindent \textbf{(P\_1,1.66-67.1)} KA\_I,171.18-172.17 Ro\_I,507-511 {34/42}     tat yathā : niṣkrāntaḥ deśāt nirdeśaḥ . \\
\noindent \textbf{(P\_1,1.66-67.1)} KA\_I,171.18-172.17 Ro\_I,507-511 {35/42}     bahirdeśaḥ iti gamyate . \\
\noindent \textbf{(P\_1,1.66-67.1)} KA\_I,171.18-172.17 Ro\_I,507-511 {36/42}     śabdaḥ ca śabdāt bahirbhūtaḥ arthaḥ abahirbhūtaḥ . \\
\noindent \textbf{(P\_1,1.66-67.1)} KA\_I,171.18-172.17 Ro\_I,507-511 {37/42}     atha nirdiṣṭagrahaṇam kimartham . \\
\noindent \textbf{(P\_1,1.66-67.1)} KA\_I,171.18-172.17 Ro\_I,507-511 {38/42}     nirdiṣṭagrahaṇam ānantaryārtham . \\
\noindent \textbf{(P\_1,1.66-67.1)} KA\_I,171.18-172.17 Ro\_I,507-511 {39/42}     nirdiṣṭagrahaṇam kriyate ānantaryārtham . \\
\noindent \textbf{(P\_1,1.66-67.1)} KA\_I,171.18-172.17 Ro\_I,507-511 {40/42}     ānantaryamātre kāryam yathā syāt . \\
\noindent \textbf{(P\_1,1.66-67.1)} KA\_I,171.18-172.17 Ro\_I,507-511 {41/42}     ikaḥ yaṇ aci : dadhi atra madhu atra . \\
\noindent \textbf{(P\_1,1.66-67.1)} KA\_I,171.18-172.17 Ro\_I,507-511 {42/42}     iha mā bhūt :samidhau samidhaḥ , dṛṣadau dṛṣadaḥ . \\
\noindent \textbf{(P\_1,1.66-67.2)} KA\_I,172.19-174.5 Ro\_I,511-515 {1/50}     kimartham punaḥ idam ucyate . \\
\noindent \textbf{(P\_1,1.66-67.2)} KA\_I,172.19-174.5 Ro\_I,511-515 {2/50}     tasmin tasmāt iti pūrvottarayoḥ yogayoḥ aviśeṣāt niyamārtham vacanam dadhi udakam pacati odanam . \\
\noindent \textbf{(P\_1,1.66-67.2)} KA\_I,172.19-174.5 Ro\_I,511-515 {3/50}     tasmin tasmāt iti pūrvottarayoḥ yogayoḥ aviśeṣāt niyamārthaḥ ayam ārambhaḥ . \\
\noindent \textbf{(P\_1,1.66-67.2)} KA\_I,172.19-174.5 Ro\_I,511-515 {4/50}     grāme devadattaḥ . \\
\noindent \textbf{(P\_1,1.66-67.2)} KA\_I,172.19-174.5 Ro\_I,511-515 {5/50}     pūrvaḥ paraḥ iti sandehaḥ . \\
\noindent \textbf{(P\_1,1.66-67.2)} KA\_I,172.19-174.5 Ro\_I,511-515 {6/50}     grāmāt devadattaḥ . \\
\noindent \textbf{(P\_1,1.66-67.2)} KA\_I,172.19-174.5 Ro\_I,511-515 {7/50}     pūrvaḥ paraḥ iti sandehaḥ . \\
\noindent \textbf{(P\_1,1.66-67.2)} KA\_I,172.19-174.5 Ro\_I,511-515 {8/50}     evam iha api : ikaḥ yaṇ aci . \\
\noindent \textbf{(P\_1,1.66-67.2)} KA\_I,172.19-174.5 Ro\_I,511-515 {9/50}     dadhi udakam , pacati odanam . \\
\noindent \textbf{(P\_1,1.66-67.2)} KA\_I,172.19-174.5 Ro\_I,511-515 {10/50}     ubhau ikau ubhau acau . \\
\noindent \textbf{(P\_1,1.66-67.2)} KA\_I,172.19-174.5 Ro\_I,511-515 {11/50}     aci pūrvasya aci parasya iti sandehaḥ . \\
\noindent \textbf{(P\_1,1.66-67.2)} KA\_I,172.19-174.5 Ro\_I,511-515 {12/50}     tiṅ atiṅaḥ iti atiṅaḥ pūrvasya atiṅaḥ parasya iti sandehaḥ . \\
\noindent \textbf{(P\_1,1.66-67.2)} KA\_I,172.19-174.5 Ro\_I,511-515 {13/50}     iṣyate ca atra aci pūrvasya syāt , atiṅaḥ parasya iti . \\
\noindent \textbf{(P\_1,1.66-67.2)} KA\_I,172.19-174.5 Ro\_I,511-515 {14/50}     tat ca antareṇa yatnam na sidhyati iti niyamārtham vacanam . \\
\noindent \textbf{(P\_1,1.66-67.2)} KA\_I,172.19-174.5 Ro\_I,511-515 {15/50}     asti prayojanam etat . \\
\noindent \textbf{(P\_1,1.66-67.2)} KA\_I,172.19-174.5 Ro\_I,511-515 {16/50}     kim tarhi iti . \\
\noindent \textbf{(P\_1,1.66-67.2)} KA\_I,172.19-174.5 Ro\_I,511-515 {17/50}     atha yatra ubhayam nirdiśyate kim tatra pūrvasya kāryam bhavati āhosvit parasya iti . \\
\noindent \textbf{(P\_1,1.66-67.2)} KA\_I,172.19-174.5 Ro\_I,511-515 {18/50}     ubhayanirdeśe vipratiṣedhāt pañcamīnirdeśaḥ . \\
\noindent \textbf{(P\_1,1.66-67.2)} KA\_I,172.19-174.5 Ro\_I,511-515 {19/50}     ubhayanirdeśe vipratiṣedhāt pañcamīnirdeśaḥ bhaviṣyati . \\
\noindent \textbf{(P\_1,1.66-67.2)} KA\_I,172.19-174.5 Ro\_I,511-515 {20/50}     kim prayojanam . \\
\noindent \textbf{(P\_1,1.66-67.2)} KA\_I,172.19-174.5 Ro\_I,511-515 {21/50}     prayojanam ataḥ lasārvadhātukanudāttatve . \\
\noindent \textbf{(P\_1,1.66-67.2)} KA\_I,172.19-174.5 Ro\_I,511-515 {22/50}     vakṣyati tāsyādibhyaḥ anudāttatve saptamīnirdeśaḥ abhyastasijarthaḥ iti . \\
\noindent \textbf{(P\_1,1.66-67.2)} KA\_I,172.19-174.5 Ro\_I,511-515 {23/50}     tasmin kriyamāṇe tāsyādibhyaḥ parasya lasārvadhātukasya lasārvadhātuke parataḥ tāsyādīnām iti sandehaḥ . \\
\noindent \textbf{(P\_1,1.66-67.2)} KA\_I,172.19-174.5 Ro\_I,511-515 {24/50}     tāsyādibhyaḥ parasya lasārvadhātukasya . \\
\noindent \textbf{(P\_1,1.66-67.2)} KA\_I,172.19-174.5 Ro\_I,511-515 {25/50}     bahoḥ iṣṭhādīnām ādilopaḥ . \\
\noindent \textbf{(P\_1,1.66-67.2)} KA\_I,172.19-174.5 Ro\_I,511-515 {26/50}     bahoḥ uttareṣām iṣṭhemeyasām iṣṭhemayaḥsu parataḥ bahoḥ iti sandehaḥ . \\
\noindent \textbf{(P\_1,1.66-67.2)} KA\_I,172.19-174.5 Ro\_I,511-515 {27/50}     bahoḥ uttareṣām iṣṭhemeyasām . \\
\noindent \textbf{(P\_1,1.66-67.2)} KA\_I,172.19-174.5 Ro\_I,511-515 {28/50}     gotaḥ ṇit . \\
\noindent \textbf{(P\_1,1.66-67.2)} KA\_I,172.19-174.5 Ro\_I,511-515 {29/50}     gotaḥ parasya sarvanāmasthānasya sarvanāmasthāne parataḥ gotaḥ iti sandehaḥ . \\
\noindent \textbf{(P\_1,1.66-67.2)} KA\_I,172.19-174.5 Ro\_I,511-515 {30/50}     gotaḥ parasya sarvanāmasthānasya . \\
\noindent \textbf{(P\_1,1.66-67.2)} KA\_I,172.19-174.5 Ro\_I,511-515 {31/50}     rudādibhyaḥ sārvadhātuke . \\
\noindent \textbf{(P\_1,1.66-67.2)} KA\_I,172.19-174.5 Ro\_I,511-515 {32/50}     rudādibhyaḥ parasya sārvadhātukasya sārvadhātuke parataḥ rudādīnām iti sandehaḥ . \\
\noindent \textbf{(P\_1,1.66-67.2)} KA\_I,172.19-174.5 Ro\_I,511-515 {33/50}     rudādibhyaḥ parasya sārvadhātukasya . \\
\noindent \textbf{(P\_1,1.66-67.2)} KA\_I,172.19-174.5 Ro\_I,511-515 {34/50}     āne muk īt āsaḥ . \\
\noindent \textbf{(P\_1,1.66-67.2)} KA\_I,172.19-174.5 Ro\_I,511-515 {35/50}     āsaḥ uttarasya ānasya , āne parataḥ āsaḥ iti sandehaḥ . \\
\noindent \textbf{(P\_1,1.66-67.2)} KA\_I,172.19-174.5 Ro\_I,511-515 {36/50}     āsaḥ uttarasya ānasya . \\
\noindent \textbf{(P\_1,1.66-67.2)} KA\_I,172.19-174.5 Ro\_I,511-515 {37/50}     āmi sarvanāmnaḥ suṭ . \\
\noindent \textbf{(P\_1,1.66-67.2)} KA\_I,172.19-174.5 Ro\_I,511-515 {38/50}     sarvanāmnaḥ uttarasya āmaḥ āmi parataḥ sarvanāmnaḥ iti sandehaḥ . \\
\noindent \textbf{(P\_1,1.66-67.2)} KA\_I,172.19-174.5 Ro\_I,511-515 {39/50}     sarvanāmnaḥ uttarasya . \\
\noindent \textbf{(P\_1,1.66-67.2)} KA\_I,172.19-174.5 Ro\_I,511-515 {40/50}     gheḥ ṅiti āṭ nadyāḥ . \\
\noindent \textbf{(P\_1,1.66-67.2)} KA\_I,172.19-174.5 Ro\_I,511-515 {41/50}     nadyāḥ uttareṣām ṅitām ṅitsu parataḥ nadyāḥ iti sandehaḥ . \\
\noindent \textbf{(P\_1,1.66-67.2)} KA\_I,172.19-174.5 Ro\_I,511-515 {42/50}     nadyāḥ uttareṣām ṅitām . \\
\noindent \textbf{(P\_1,1.66-67.2)} KA\_I,172.19-174.5 Ro\_I,511-515 {43/50}     yāṭ āpaḥ . \\
\noindent \textbf{(P\_1,1.66-67.2)} KA\_I,172.19-174.5 Ro\_I,511-515 {44/50}      āpaḥ uttarasya ṅitaḥ ṅiti parataḥ āpaḥ iti sandehaḥ . \\
\noindent \textbf{(P\_1,1.66-67.2)} KA\_I,172.19-174.5 Ro\_I,511-515 {45/50}     āpaḥ uttarasya ṅitaḥ . ṅamaḥ hrasvāt aci ṅamuṭ nityam . \\
\noindent \textbf{(P\_1,1.66-67.2)} KA\_I,172.19-174.5 Ro\_I,511-515 {46/50}     ṅamaḥ uttarasya acaḥ aci parataḥ ṅamaḥ iti sandehaḥ . \\
\noindent \textbf{(P\_1,1.66-67.2)} KA\_I,172.19-174.5 Ro\_I,511-515 {47/50}     ṅamaḥ uttarasya acaḥ . \\
\noindent \textbf{(P\_1,1.66-67.2)} KA\_I,172.19-174.5 Ro\_I,511-515 {48/50}     vibhaktiviśeṣanirdeśānavakāśatvāt avipratiṣedhaḥ . \\
\noindent \textbf{(P\_1,1.66-67.2)} KA\_I,172.19-174.5 Ro\_I,511-515 {49/50}     vibhaktiviśeṣanirdeśasya anavakāśatvāt ayuktaḥ ayam vipratiṣedhaḥ . \\
\noindent \textbf{(P\_1,1.66-67.2)} KA\_I,172.19-174.5 Ro\_I,511-515 {50/50}     sarvatra eva atra kṛtasāmarthyā saptamī akṛtasāmārthyā pañcamī iti kṛtvā pañcamīnirdeśaḥ bhaviṣyati . \\
\noindent \textbf{(P\_1,1.66-67.3)} KA\_I,174.6-175.18 Ro\_I,515-518 {1/62}     yathārtham vā ṣaṣṭhīnirdeśaḥ . \\
\noindent \textbf{(P\_1,1.66-67.3)} KA\_I,174.6-175.18 Ro\_I,515-518 {2/62}     yathārtham vā ṣaṣṭhīnirdeśaḥ kartavyaḥ . \\
\noindent \textbf{(P\_1,1.66-67.3)} KA\_I,174.6-175.18 Ro\_I,515-518 {3/62}     yatra pūrvasya kāryam iṣyate tatra pūrvasya ṣaṣṭhī kartavyā . \\
\noindent \textbf{(P\_1,1.66-67.3)} KA\_I,174.6-175.18 Ro\_I,515-518 {4/62}     yatra parasya kāryam iṣyate tatra parasya ṣaṣṭhī kartavyā . \\
\noindent \textbf{(P\_1,1.66-67.3)} KA\_I,174.6-175.18 Ro\_I,515-518 {5/62}     saḥ tarhi tathā nirdeśaḥ kartavyaḥ . \\
\noindent \textbf{(P\_1,1.66-67.3)} KA\_I,174.6-175.18 Ro\_I,515-518 {6/62}     na kartavyaḥ . \\
\noindent \textbf{(P\_1,1.66-67.3)} KA\_I,174.6-175.18 Ro\_I,515-518 {7/62}     anena eva prakḷptiḥ bhaviṣyati : tasmin iti nirdiṣṭe pūrvasya ṣaṣṭhī . \\
\noindent \textbf{(P\_1,1.66-67.3)} KA\_I,174.6-175.18 Ro\_I,515-518 {8/62}     tasmāt iti nirdiṣṭe parasya ṣaṣṭhī . \\
\noindent \textbf{(P\_1,1.66-67.3)} KA\_I,174.6-175.18 Ro\_I,515-518 {9/62}     tat tarhi ṣaṣṭhīgrahaṇam kartavyam . \\
\noindent \textbf{(P\_1,1.66-67.3)} KA\_I,174.6-175.18 Ro\_I,515-518 {10/62}     na kartavyam . \\
\noindent \textbf{(P\_1,1.66-67.3)} KA\_I,174.6-175.18 Ro\_I,515-518 {11/62}     prakṛtam anuvartate . \\
\noindent \textbf{(P\_1,1.66-67.3)} KA\_I,174.6-175.18 Ro\_I,515-518 {12/62}     kva prakṛtam . \\
\noindent \textbf{(P\_1,1.66-67.3)} KA\_I,174.6-175.18 Ro\_I,515-518 {13/62}     ṣaṣṭhī sthāneyogā iti . \\
\noindent \textbf{(P\_1,1.66-67.3)} KA\_I,174.6-175.18 Ro\_I,515-518 {14/62}     prakalpakam iti cet niyamābhāvaḥ . \\
\noindent \textbf{(P\_1,1.66-67.3)} KA\_I,174.6-175.18 Ro\_I,515-518 {15/62}     prakalpakam iti cet niyamasya abhāvaḥ . \\
\noindent \textbf{(P\_1,1.66-67.3)} KA\_I,174.6-175.18 Ro\_I,515-518 {16/62}     uktam ca etat : niyamārthaḥ ayam ārambhaḥ iti . \\
\noindent \textbf{(P\_1,1.66-67.3)} KA\_I,174.6-175.18 Ro\_I,515-518 {17/62}     pratyayavidhau khalu api pañcamyāḥ prakalpikāḥ syuḥ . \\
\noindent \textbf{(P\_1,1.66-67.3)} KA\_I,174.6-175.18 Ro\_I,515-518 {18/62}     tatra kaḥ doṣaḥ . \\
\noindent \textbf{(P\_1,1.66-67.3)} KA\_I,174.6-175.18 Ro\_I,515-518 {19/62}     guptijkibhyaḥ san iti eṣā pañcamī san iti prathamāyāḥ ṣaṣṭhīm prakalpayet tasmāt iti uttarasya iti . \\
\noindent \textbf{(P\_1,1.66-67.3)} KA\_I,174.6-175.18 Ro\_I,515-518 {20/62}     astu . \\
\noindent \textbf{(P\_1,1.66-67.3)} KA\_I,174.6-175.18 Ro\_I,515-518 {21/62}     na kaḥ cit ādeśaḥ pratinirdiśyate . \\
\noindent \textbf{(P\_1,1.66-67.3)} KA\_I,174.6-175.18 Ro\_I,515-518 {22/62}     tatra āntaryataḥ sanaḥ san eva bhaviṣyati . \\
\noindent \textbf{(P\_1,1.66-67.3)} KA\_I,174.6-175.18 Ro\_I,515-518 {23/62}     na evam śakyam . \\
\noindent \textbf{(P\_1,1.66-67.3)} KA\_I,174.6-175.18 Ro\_I,515-518 {24/62}     itsañjñā na prakalpeta . \\
\noindent \textbf{(P\_1,1.66-67.3)} KA\_I,174.6-175.18 Ro\_I,515-518 {25/62}     upadeśe iti itsañjñā ucyate . \\
\noindent \textbf{(P\_1,1.66-67.3)} KA\_I,174.6-175.18 Ro\_I,515-518 {26/62}     prakṛtivikārāvyavasthā ca . \\
\noindent \textbf{(P\_1,1.66-67.3)} KA\_I,174.6-175.18 Ro\_I,515-518 {27/62}     prakṛtivikārayoḥ ca vyavasthā na prakalpeta . \\
\noindent \textbf{(P\_1,1.66-67.3)} KA\_I,174.6-175.18 Ro\_I,515-518 {28/62}     ikaḥ yaṇ aci : aci iti eṣā saptamī yaṇ iti prathamāyāḥ ṣaṣṭhīm prakalpayet tasmin iti nirdiṣṭe pūrvasya iti . \\
\noindent \textbf{(P\_1,1.66-67.3)} KA\_I,174.6-175.18 Ro\_I,515-518 {29/62}     saptamīpañcamyoḥ ca bhāvāt ubhayatra ṣaṣṭhīprakḷptiḥ tatra ubhayakāryaprasaṅgaḥ . \\
\noindent \textbf{(P\_1,1.66-67.3)} KA\_I,174.6-175.18 Ro\_I,515-518 {30/62}     saptamīpañcamyoḥ ca bhāvāt ubhayatra eva ṣaṣṭhī prāpnoti . \\
\noindent \textbf{(P\_1,1.66-67.3)} KA\_I,174.6-175.18 Ro\_I,515-518 {31/62}     tāsyādibhyaḥ iti eṣā pañcamī lasārvadhātuke iti asyāḥ saptamyāḥ ṣaṣṭhīm prakalpayet tasmāt iti uttarasya iti . \\
\noindent \textbf{(P\_1,1.66-67.3)} KA\_I,174.6-175.18 Ro\_I,515-518 {32/62}     tathā lasārvadhātuke iti eṣā saptamī tāsyādibhyaḥ iti pañcamyāḥ ṣaṣṭhīm prakalpayet tasmin iti nirdiṣṭe pūrvasya iti . \\
\noindent \textbf{(P\_1,1.66-67.3)} KA\_I,174.6-175.18 Ro\_I,515-518 {33/62}     tatra kaḥ doṣaḥ . \\
\noindent \textbf{(P\_1,1.66-67.3)} KA\_I,174.6-175.18 Ro\_I,515-518 {34/62}     ubhayoḥ kāryam tatra prāpnoti . \\
\noindent \textbf{(P\_1,1.66-67.3)} KA\_I,174.6-175.18 Ro\_I,515-518 {35/62}     na eṣaḥ doṣaḥ . \\
\noindent \textbf{(P\_1,1.66-67.3)} KA\_I,174.6-175.18 Ro\_I,515-518 {36/62}     yat tāvat ucyate : prakalpakam iti cet niyamābhāvaḥ iti . \\
\noindent \textbf{(P\_1,1.66-67.3)} KA\_I,174.6-175.18 Ro\_I,515-518 {37/62}     mā bhūt niyamaḥ . \\
\noindent \textbf{(P\_1,1.66-67.3)} KA\_I,174.6-175.18 Ro\_I,515-518 {38/62}     saptamīnirdiṣṭe pūrvasya ṣaṣṭhī prakalpyate pañcamīnirdiṣṭe parasya . \\
\noindent \textbf{(P\_1,1.66-67.3)} KA\_I,174.6-175.18 Ro\_I,515-518 {39/62}     yāvatā saptamīnirdiṣṭe pūrvasya ṣaṣṭhī prakalpyate evam pañcamīnirdiṣṭe parasya . \\
\noindent \textbf{(P\_1,1.66-67.3)} KA\_I,174.6-175.18 Ro\_I,515-518 {40/62}     na utsahate saptamīnirdiṣṭe parasya kāryam bhavitum na api pañcamīnirdiṣṭe pūrvasya . \\
\noindent \textbf{(P\_1,1.66-67.3)} KA\_I,174.6-175.18 Ro\_I,515-518 {41/62}     yat api ucyate : pratyayavidhau khalu api pañcamyāḥ prakalpikāḥ syuḥ iti . \\
\noindent \textbf{(P\_1,1.66-67.3)} KA\_I,174.6-175.18 Ro\_I,515-518 {42/62}     santu prakalpikāḥ . \\
\noindent \textbf{(P\_1,1.66-67.3)} KA\_I,174.6-175.18 Ro\_I,515-518 {43/62}     nanu ca uktam guptijkibhyaḥ san iti eṣā pañcamī san iti prathamāyāḥ ṣaṣṭhīm prakalpayet tasmāt iti uttarasya iti . \\
\noindent \textbf{(P\_1,1.66-67.3)} KA\_I,174.6-175.18 Ro\_I,515-518 {44/62}     parihṛtam etat : na kaḥ cit ādeśaḥ pratinirdiśyate . \\
\noindent \textbf{(P\_1,1.66-67.3)} KA\_I,174.6-175.18 Ro\_I,515-518 {45/62}     tatra āntaryataḥ sanaḥ san eva bhaviṣyati iti . \\
\noindent \textbf{(P\_1,1.66-67.3)} KA\_I,174.6-175.18 Ro\_I,515-518 {46/62}     nanu ca uktam : na evam śakyam . \\
\noindent \textbf{(P\_1,1.66-67.3)} KA\_I,174.6-175.18 Ro\_I,515-518 {47/62}     itsañjñā na prakalpeta . \\
\noindent \textbf{(P\_1,1.66-67.3)} KA\_I,174.6-175.18 Ro\_I,515-518 {48/62}     upadeśe iti itsañjñā ucyate iti . \\
\noindent \textbf{(P\_1,1.66-67.3)} KA\_I,174.6-175.18 Ro\_I,515-518 {49/62}     syāt eṣaḥ doṣaḥ yadi itsañjñā ādeśam pratīkṣeta . \\
\noindent \textbf{(P\_1,1.66-67.3)} KA\_I,174.6-175.18 Ro\_I,515-518 {50/62}     tatra khalu kṛtāyām itsañjñāyām lope ca kṛte ādeśaḥ bhaviṣyati . \\
\noindent \textbf{(P\_1,1.66-67.3)} KA\_I,174.6-175.18 Ro\_I,515-518 {51/62}     upadeśe iti hi itsañjñā ucyate . \\
\noindent \textbf{(P\_1,1.66-67.3)} KA\_I,174.6-175.18 Ro\_I,515-518 {52/62}     atha vā na anutpanne sani prakḷptyā bhavitavyam . \\
\noindent \textbf{(P\_1,1.66-67.3)} KA\_I,174.6-175.18 Ro\_I,515-518 {53/62}     yadā ca utpannaḥ san tadā kṛtasāmarthyā pañcamī iti kṛtvā prakḷptiḥ na bhaviṣyati . \\
\noindent \textbf{(P\_1,1.66-67.3)} KA\_I,174.6-175.18 Ro\_I,515-518 {54/62}     yat api ucyate : prakṛtivikārāvyavasthā ca iti . \\
\noindent \textbf{(P\_1,1.66-67.3)} KA\_I,174.6-175.18 Ro\_I,515-518 {55/62}     tatra api kṛtā prakṛtau ṣaṣṭhī ikaḥ iti vikṛtau prathamā yaṇ iti . \\
\noindent \textbf{(P\_1,1.66-67.3)} KA\_I,174.6-175.18 Ro\_I,515-518 {56/62}     yatra ca nāma sautrī ṣaṣṭhī na asti tatra prakḷptyā bhavitavyam . \\
\noindent \textbf{(P\_1,1.66-67.3)} KA\_I,174.6-175.18 Ro\_I,515-518 {57/62}     atha vā astu tāvat ikaḥ yaṇ aci iti yatra nāma sautrī ṣaṣṭhī . \\
\noindent \textbf{(P\_1,1.66-67.3)} KA\_I,174.6-175.18 Ro\_I,515-518 {58/62}     yadi ca idānīm aci iti eṣā saptamī yaṇ iti prathamāyāḥ ṣaṣṭhīm prakalpayet tasmin iti nirdiṣṭe pūrvasya iti astu . \\
\noindent \textbf{(P\_1,1.66-67.3)} KA\_I,174.6-175.18 Ro\_I,515-518 {59/62}     na kaḥ cit anyaḥ ādeśaḥ pratinirdiśyate . \\
\noindent \textbf{(P\_1,1.66-67.3)} KA\_I,174.6-175.18 Ro\_I,515-518 {60/62}     tatra āntaryataḥ yaṇaḥ yaṇ eva bhaviṣyati . \\
\noindent \textbf{(P\_1,1.66-67.3)} KA\_I,174.6-175.18 Ro\_I,515-518 {61/62}     yat api ucyate : saptamīpañcamyoḥ ca bhāvāt ubhayatra ṣaṣṭhīprakḷptiḥ tatra ubhayakāryaprasaṅgaḥ iti . \\
\noindent \textbf{(P\_1,1.66-67.3)} KA\_I,174.6-175.18 Ro\_I,515-518 {62/62}     ācāryapravṛttiḥ jñāpayati na ubhe yugapat prakalpike bhavataḥ iti yat ayam ekaḥ pūrvaparayoḥ iti pūrvagrahaṇam karoti . \\
\noindent \textbf{(P\_1,1.68.1)} rūpagrahaṇam kim artham na svam śabdasya aśabdasañjñā bhavati iti eva rūpam śabasya sañjñā bhaviṣyati . \\
\noindent \textbf{(P\_1,1.68.1)} na hi anyat svam śabdasya asti anyat ataḥ rūpāt . \\
\noindent \textbf{(P\_1,1.68.1)} evam tarhi siddhe sati yat rūpagrahaṇam karoti tat jñāpayati ācāryaḥ asti anyat rūpāt svam śabdasya iti . \\
\noindent \textbf{(P\_1,1.68.1)} kim punaḥ tat . \\
\noindent \textbf{(P\_1,1.68.1)} arthaḥ . \\
\noindent \textbf{(P\_1,1.68.1)} kim etasya jñāpane prayojanam . \\
\noindent \textbf{(P\_1,1.68.1)} arthavadgrahaṇe na anarthakasya iti eṣā paribhāṣā na kartavyā bhavati . \\
\noindent \textbf{(P\_1,1.68.2)} kimartham punaḥ idam ucyate . \\
\noindent \textbf{(P\_1,1.68.2)} śabdena arthagateḥ arthasya asambhavāt tadvācinaḥ sañjñāpratiṣedhārtham svaṃrūpavacanam . śabdena uccāritena arthaḥ gamyate . \\
\noindent \textbf{(P\_1,1.68.2)} gām ānaya dadhi aśāna iti arthaḥ ānīyate arthaḥ ca bhujyate . \\
\noindent \textbf{(P\_1,1.68.2)} arthasya asambhavāt . \\
\noindent \textbf{(P\_1,1.68.2)} iha vyākaraṇe arthe kāryasya asambhavaḥ . \\
\noindent \textbf{(P\_1,1.68.2)} agneḥ ḍak iti : na śakyate aṅgārebhyaḥ paraḥ ḍhak kartum . \\
\noindent \textbf{(P\_1,1.68.2)} śabdena arthagateḥ arthasya asambhavāt yāvantaḥ tadvācinaḥ śabdāḥ tāvadbhyaḥ sarvebhyaḥ utpattiḥ prāpnoti . \\
\noindent \textbf{(P\_1,1.68.2)} iṣyate ca tasmāt eva syāt iti . \\
\noindent \textbf{(P\_1,1.68.2)} tat ca antareṇa yatnam na sidhyati iti tadvācinaḥ sañjñāpratiṣedhārtham svaṃrūpavacanam . \\
\noindent \textbf{(P\_1,1.68.2)} evamartham idam ucyate . \\
\noindent \textbf{(P\_1,1.68.2)} na vā śabdapūrvakaḥ hi arthe sampratyayaḥ tasmāt arthanivṛttiḥ . \\
\noindent \textbf{(P\_1,1.68.2)} na vā etat prayojanam asti . \\
\noindent \textbf{(P\_1,1.68.2)} kim kāraṇam . \\
\noindent \textbf{(P\_1,1.68.2)} śabdapūrvakaḥ hi arthe sampratyayaḥ . \\
\noindent \textbf{(P\_1,1.68.2)} śabdapūrvakaḥ hi arthasya sampratyayaḥ . \\
\noindent \textbf{(P\_1,1.68.2)} ātaḥ ca śabdapūrvakaḥ : yaḥ api hi asau āhūyate nāmnā nāma yadā anena na upalabdham bhavati tada pṛcchati kim bhavān āha iti . \\
\noindent \textbf{(P\_1,1.68.2)} śabdapūrvakaḥ ca arthasya sampratyayaḥ iha ca vyākaraṇe śabde kāryasya sambhavaḥ arthe asambhavaḥ . \\
\noindent \textbf{(P\_1,1.68.2)} tasmāt arthanivṛttiḥ bhaviṣyati . \\
\noindent \textbf{(P\_1,1.68.2)} idam tarhi prayojanam aśabdasañjñā iti vakṣyāmi iti . \\
\noindent \textbf{(P\_1,1.68.2)} iha mā bhūt : dādhāḥ ghu adāp taraptamapau ghaḥ iti . \\
\noindent \textbf{(P\_1,1.68.2)} sañjñāpratiṣedhānarthakyam vacanaprāmāṇyāt . \\
\noindent \textbf{(P\_1,1.68.2)} sañjñāpratiṣedhaḥ ca anarthakaḥ . \\
\noindent \textbf{(P\_1,1.68.2)} śabdasañjñāyām svarūpavidhiḥ kasmāt na bhavati . \\
\noindent \textbf{(P\_1,1.68.2)} vacanaprāmāṇyāt . \\
\noindent \textbf{(P\_1,1.68.2)} śabdasañjñāvacanasāmarthyāt . \\
\noindent \textbf{(P\_1,1.68.2)} nanu ca vacanaprāmāṇyāt sañjñinām sampratyayaḥ syāt svarūpagrahaṇāt ca sañjñāyāḥ . \\
\noindent \textbf{(P\_1,1.68.2)} etat api na asti prayojanam . \\
\noindent \textbf{(P\_1,1.68.2)} ācāryapravṛttiḥ jñāpayati śabdasañjñāyām na svarūpavidhiḥ bhavati iti yat ayam ṣṇāntā ṣaṭ iti ṣakārāntāyāḥ saṅkhyāyāḥ ṣaṭsañjñām śāsti . \\
\noindent \textbf{(P\_1,1.68.2)} itarathā hi vacanaprāmāṇyāt nakārāntāyāḥ saṅkhyāyāḥ sampratyayaḥ syāt svarūpagrahaṇāt ca ṣakārāntāyāḥ . \\
\noindent \textbf{(P\_1,1.68.2)} na etat asti prayojanam . \\
\noindent \textbf{(P\_1,1.68.2)} na hi ṣakārāntā sañjñā . \\
\noindent \textbf{(P\_1,1.68.2)} kā tarhi . \\
\noindent \textbf{(P\_1,1.68.2)} ḍakārāntā . \\
\noindent \textbf{(P\_1,1.68.2)} asiddham jaśtvam . \\
\noindent \textbf{(P\_1,1.68.2)} tasya asiddhatvāt ṣakārāntā . \\
\noindent \textbf{(P\_1,1.68.2)} mantrādyartham tarhi idam vaktavyam . \\
\noindent \textbf{(P\_1,1.68.2)} mantre , ṛci yajuṣi iti yat ucyate tat mantraśabde ṛkśabde ca yajuḥśabde ca mā bhūt . \\
\noindent \textbf{(P\_1,1.68.2)} mantrādyartham iti cet śāstrasāmarthyāt arthagateḥ siddham . \\
\noindent \textbf{(P\_1,1.68.2)} mantrādyartham iti cet na . \\
\noindent \textbf{(P\_1,1.68.2)} kim kāraṇam . \\
\noindent \textbf{(P\_1,1.68.2)} śāstrasāmarthyāt arthasya gatiḥ bhaviṣyati . \\
\noindent \textbf{(P\_1,1.68.2)} mantre , ṛci yajuṣi iti yat ucyate tat mantraśabde ṛkśabde ca yajuḥśabde ca tasya kāryasya sambhavaḥ na asti iti kṛtvā mantrādisahacaritaḥ yaḥ arthaḥ tasya gatiḥ bhaviṣyati sāhacaryāt . \\
\noindent \textbf{(P\_1,1.68.3)} sit tadviśeṣāṇām vṛkṣādyartham . \\
\noindent \textbf{(P\_1,1.68.3)} sinnirdeśaḥ kartavyaḥ . \\
\noindent \textbf{(P\_1,1.68.3)} tataḥ vaktavyam : tadviśeṣāṇām grahaṇam bhavati iti . \\
\noindent \textbf{(P\_1,1.68.3)} kim prayojanam . \\
\noindent \textbf{(P\_1,1.68.3)} vṛkṣādyartham . \\
\noindent \textbf{(P\_1,1.68.3)} vibhāṣā vṛkṣamṛga iti : plakṣanyagrodham , plakṣanyagrodhāḥ . \\
\noindent \textbf{(P\_1,1.68.3)} pit paryāyavacanasya ca svādyartham . \\
\noindent \textbf{(P\_1,1.68.3)} pinnirdeśaḥ kartavyaḥ . \\
\noindent \textbf{(P\_1,1.68.3)} tataḥ vaktavyam : paryāyavacanasya tadviśeṣāṇām ca grahaṇam bhavati svasya ca rūpasya iti . \\
\noindent \textbf{(P\_1,1.68.3)} kim prayojanam . \\
\noindent \textbf{(P\_1,1.68.3)} svādyartham . \\
\noindent \textbf{(P\_1,1.68.3)} sve puṣaḥ : svapoṣam puṣyati raipoṣam , vidyāpoṣam , gopoṣam aśvapoṣam . \\
\noindent \textbf{(P\_1,1.68.3)} jit paryāyavacanasya eva rājādyartham . \\
\noindent \textbf{(P\_1,1.68.3)} jinnirdeśaḥ kartavyaḥ . \\
\noindent \textbf{(P\_1,1.68.3)} tataḥ vaktavyam paryāyavacanasya eva grahaṇam bhavati . \\
\noindent \textbf{(P\_1,1.68.3)} kim prayojanam . \\
\noindent \textbf{(P\_1,1.68.3)} rājādyartham . \\
\noindent \textbf{(P\_1,1.68.3)} sabhā rājāmanuṣyapūrvā : inasabham īśvarasabham . \\
\noindent \textbf{(P\_1,1.68.3)} tasya eva na bhavati : rājasabhā . \\
\noindent \textbf{(P\_1,1.68.3)} tadviśeṣāṇām ca na bhavati : puṣyamitrasabhā candraguptasabhā . \\
\noindent \textbf{(P\_1,1.68.3)} jhit tasya ca tadviśeṣāṇām ca matsyādyartham . \\
\noindent \textbf{(P\_1,1.68.3)} jhinnirdeśaḥ kartavyaḥ . \\
\noindent \textbf{(P\_1,1.68.3)} tataḥ vaktavyam : tasya ca grahaṇam bhavati tadviśeṣāṇām ca iti . \\
\noindent \textbf{(P\_1,1.68.3)} kim prayojanam . \\
\noindent \textbf{(P\_1,1.68.3)} matsyādyartham . \\
\noindent \textbf{(P\_1,1.68.3)} pakṣimatsyamṛgān hanti : mātsyikaḥ . \\
\noindent \textbf{(P\_1,1.68.3)} tadviśeṣāṇām : śāpharikaḥ , śākulikaḥ . \\
\noindent \textbf{(P\_1,1.68.3)} paryāyavacanānām na bhavati : ajihmān hanti iti . \\
\noindent \textbf{(P\_1,1.68.3)} asya ekasya paryāyavacanasya iṣyate : mīnān hanti mainikaḥ . \\
\noindent \textbf{(P\_1,1.69.1)} apratyayaḥ iti kimartham . \\
\noindent \textbf{(P\_1,1.69.1)} sanāśaṃsabhikṣaḥ uḥ , a sāmpratike . \\
\noindent \textbf{(P\_1,1.69.1)} atyalpam idam ucyate : apratyayaḥ iti . \\
\noindent \textbf{(P\_1,1.69.1)} apratyayādeśaṭitkinmitaḥ iti vaktavyam . \\
\noindent \textbf{(P\_1,1.69.1)} pratyaye udāhṛtam . \\
\noindent \textbf{(P\_1,1.69.1)} ādeśe : idamaḥ iś : iha , itaḥ . \\
\noindent \textbf{(P\_1,1.69.1)} ṭiti . \\
\noindent \textbf{(P\_1,1.69.1)} lavitā lavitum . \\
\noindent \textbf{(P\_1,1.69.1)} kiti . \\
\noindent \textbf{(P\_1,1.69.1)} babhūva . \\
\noindent \textbf{(P\_1,1.69.1)} miti . \\
\noindent \textbf{(P\_1,1.69.1)} he anaḍvan . \\
\noindent \textbf{(P\_1,1.69.1)} ṭitaḥ parihāraḥ . \\
\noindent \textbf{(P\_1,1.69.1)} ācāryapravṛttiḥ jñāpayati na ṭitaḥ savarṇānām grahaṇam bhavati iti yat ayam grahaḥ aliṭi dīrghatvam śāsti . \\
\noindent \textbf{(P\_1,1.69.1)} na etat asti jñāpakam . \\
\noindent \textbf{(P\_1,1.69.1)} niyamārtham etat syāt : grahaḥ aliṭi dīrghaḥ eva iti . \\
\noindent \textbf{(P\_1,1.69.1)} yat tarhi vṝtaḥ vā iti vibhāṣām śāsti . \\
\noindent \textbf{(P\_1,1.69.1)} sarveṣām eva parihāraḥ : bhāvyamānena savarṇānām grahaṇam na iti evam bhaviṣyati . \\
\noindent \textbf{(P\_1,1.69.1)} pratyaye bhūyān parihāraḥ : anabhidhānāt pratyayaḥ savarṇān na grahīṣyati . \\
\noindent \textbf{(P\_1,1.69.1)} yān hi pratyayaḥ savarṇagrahaṇena gṛhṇīyāt na taiḥ arthasya abhidhānam syāt . \\
\noindent \textbf{(P\_1,1.69.1)} anabhidhānāt na bhaviṣyati . \\
\noindent \textbf{(P\_1,1.69.1)} idam tarhi prayojanam : iha ke cit pratīyante ke cit pratyāyyante . \\
\noindent \textbf{(P\_1,1.69.1)} hrasvāḥ pratīyante ḍīrghāḥ pratyāyyante . \\
\noindent \textbf{(P\_1,1.69.1)} yāvat brūyāt pratyāyyamānena savarṇānam grahaṇam na iti tāvat apratyayaḥ iti . \\
\noindent \textbf{(P\_1,1.69.1)} kam punaḥ dīrghaḥ savarṇagrahaṇena gṛhṇīyāt . \\
\noindent \textbf{(P\_1,1.69.1)} hrasvam . \\
\noindent \textbf{(P\_1,1.69.1)} yatnādhikyāt na grahīṣyati . \\
\noindent \textbf{(P\_1,1.69.1)} plutam tarhi gṛhṇīyāt . \\
\noindent \textbf{(P\_1,1.69.1)} anaṇtvāt na grahīṣyati . \\
\noindent \textbf{(P\_1,1.69.1)} evam tarhi siddhe sati yat apratyayaḥ iti pratiṣedham śāsti tat jñāpayati ācāryaḥ bhavati eṣā paribhāṣā : bhāvyamānena savarṇānām grahaṇam na iti . \\
\noindent \textbf{(P\_1,1.69.2)} kimartham punaḥ idam ucyate . \\
\noindent \textbf{(P\_1,1.69.2)} aṇ savarṇasya iti svarānunāsikyakālabhedāt . \\
\noindent \textbf{(P\_1,1.69.2)} aṇ savarṇasya iti ucyate . \\
\noindent \textbf{(P\_1,1.69.2)} svarabhedāt ānunāsikyabhedāt kālabhedāt ca aṇ savarṇān na gṛhṇīyāt . \\
\noindent \textbf{(P\_1,1.69.2)} iṣyate ca savarṇagrahaṇam syāt iti . \\
\noindent \textbf{(P\_1,1.69.2)} tat ca antareṇa yatnam na sidhyati iti evamartham idam ucyate . \\
\noindent \textbf{(P\_1,1.69.2)} asti prayojanam etat . \\
\noindent \textbf{(P\_1,1.69.2)} kim tarhi iti . \\
\noindent \textbf{(P\_1,1.69.2)} tatra pratyāhāragrahaṇe savarṇāgrahaṇam anupadeśāt . \\
\noindent \textbf{(P\_1,1.69.2)} tatra pratyāhāragrahaṇe savarṇānām grahaṇam na prāpnoti : akaḥ savarṇe dīrghaḥ iti . \\
\noindent \textbf{(P\_1,1.69.2)} kim kāraṇam . \\
\noindent \textbf{(P\_1,1.69.2)} anupadeśāt . \\
\noindent \textbf{(P\_1,1.69.2)} yathājātīyakānām sañjñā kṛtā tathājātīyakānām sampratyāyikā syāt . \\
\noindent \textbf{(P\_1,1.69.2)} hrasvānām ca kriyate . \\
\noindent \textbf{(P\_1,1.69.2)} hrasvānām eva sampratyāyikā syāt dīrghānām na syāt . \\
\noindent \textbf{(P\_1,1.69.2)} nanu ca hrasvāḥ pratīyamānāḥ dīrghān sampratyāyayiṣyanti . \\
\noindent \textbf{(P\_1,1.69.2)} hrasvasampratyayāt iti cet uccāryamāṇasampratyāyakatvāt śabdasya avacanam . \\
\noindent \textbf{(P\_1,1.69.2)} hrasvasampratyayāt iti cet uccāryamāṇaḥ śabdaḥ sampratyāyakaḥ bhavati na sampratīyamanaḥ . \\
\noindent \textbf{(P\_1,1.69.2)} tat yathā ṛk iti ukte sampāṭhamātram gamyate na asyāḥ arthaḥ gamyate . \\
\noindent \textbf{(P\_1,1.69.2)} evam tarhi varṇapāṭhe eva upadeśaḥ kariṣyate . \\
\noindent \textbf{(P\_1,1.69.2)} varṇapāṭhe upadeśaḥ iti cet avarkālatvāt paribhāṣāyāḥ anupadeśaḥ . varṇapāṭhe upadeśaḥ iti cet avarkālatvāt paribhāṣāyāḥ anupadeśaḥ . \\
\noindent \textbf{(P\_1,1.69.2)} kim parā sūtrāt kriyate iti ataḥ avarakālā . \\
\noindent \textbf{(P\_1,1.69.2)} na iti āha . \\
\noindent \textbf{(P\_1,1.69.2)} sarvathā avarakālā eva . \\
\noindent \textbf{(P\_1,1.69.2)} varṇānām upadeśaḥ tāvat . \\
\noindent \textbf{(P\_1,1.69.2)} upadeśottarakālaḥ ādiḥ antyena saha itā iti pratyāhāraḥ . \\
\noindent \textbf{(P\_1,1.69.2)} pratyāhārottarakālā savarṇasañjñā . \\
\noindent \textbf{(P\_1,1.69.2)} savarṇasañjñottarakālam aṇudit savarṇasya ca apratyayaḥ iti . \\
\noindent \textbf{(P\_1,1.69.2)} sā eṣā upadeśottarakālā avarakālā satī varṇānām utpattau nimittatvāya kalpayiṣyate iti tat na . \\
\noindent \textbf{(P\_1,1.69.2)} tasmāt upadeśaḥ . \\
\noindent \textbf{(P\_1,1.69.2)} tasmāt upadeśaḥ kartavyaḥ . \\
\noindent \textbf{(P\_1,1.69.2)} tatra anuvṛttinirdeśe savarṇāgrahaṇam anaṇtvāt . \\
\noindent \textbf{(P\_1,1.69.2)} tatra anuvṛttinirdeśe savarṇānām grahaṇam na prāpnoti : asya cvau yasya īti ca . \\
\noindent \textbf{(P\_1,1.69.2)} kim kāraṇam . \\
\noindent \textbf{(P\_1,1.69.2)} anaṇtvāt . \\
\noindent \textbf{(P\_1,1.69.2)} na hi ete aṇaḥ ye anuvṛttinirdeśe . \\
\noindent \textbf{(P\_1,1.69.2)} ke tarhi . \\
\noindent \textbf{(P\_1,1.69.2)} ye akṣarasamāmnāye upadiśyante . \\
\noindent \textbf{(P\_1,1.69.2)} evam tarhi anaṇtvāt anuvṛttau na anupadeśāt ca pratyāhāre na . \\
\noindent \textbf{(P\_1,1.69.2)} ucyate ca idam aṇ savarṇān gṛhṇāti iti . \\
\noindent \textbf{(P\_1,1.69.2)} tatra vacanāt bhaviṣyati . \\
\noindent \textbf{(P\_1,1.69.2)} vacanāt yatra tat na asti . \\
\noindent \textbf{(P\_1,1.69.2)} na idam vacanāt labhyam . \\
\noindent \textbf{(P\_1,1.69.2)} asti hi anyat etasya vacane prayojanam . \\
\noindent \textbf{(P\_1,1.69.2)} kim . \\
\noindent \textbf{(P\_1,1.69.2)} ye ete pratyāhārāṇām āditaḥ varṇāḥ taiḥ savarṇānām grahaṇam yathā syāt . \\
\noindent \textbf{(P\_1,1.69.3)} evam tarhi savarṇe aṇgrahaṇam aparibhāṣyam ākṛtigrahaṇāt . \\
\noindent \textbf{(P\_1,1.69.3)} savarṇe aṇgrahaṇam aparibhāṣyam . \\
\noindent \textbf{(P\_1,1.69.3)} kutaḥ . \\
\noindent \textbf{(P\_1,1.69.3)} ākṛtigrahaṇāt . \\
\noindent \textbf{(P\_1,1.69.3)} avarṇākṛtiḥ upadiṣṭā sā sarvam avarṇakulam grahīṣyati . \\
\noindent \textbf{(P\_1,1.69.3)} tathā ivarṇakulākṛtiḥ . \\
\noindent \textbf{(P\_1,1.69.3)} tathā uvarṇakulākṛtiḥ . \\
\noindent \textbf{(P\_1,1.69.3)} nanu ca anyā ākṛtiḥ akārasya ākārasya ca . \\
\noindent \textbf{(P\_1,1.69.3)} ananyatvāt ca . \\
\noindent \textbf{(P\_1,1.69.3)} ananyākṛtiḥ akārasya ākārasya ca . \\
\noindent \textbf{(P\_1,1.69.3)} anekāntaḥ hi ananyatvakaraḥ . \\
\noindent \textbf{(P\_1,1.69.3)} yaḥ hi anekāntena bhedaḥ na asau anyatvam karoti . \\
\noindent \textbf{(P\_1,1.69.3)} tat yathā : na yaḥ goḥ ca goḥ ca bhedaḥ saḥ anyatvam karoti . \\
\noindent \textbf{(P\_1,1.69.3)} yaḥ tu khalu goḥ ca aśvasya ca bhedaḥ saḥ anyatvam karoti . \\
\noindent \textbf{(P\_1,1.69.3)} aparaḥ āha : savarṇe aṇgrahaṇam aparibhāṣyam . \\
\noindent \textbf{(P\_1,1.69.3)} ākṛtigrahaṇāt ananyatvam . \\
\noindent \textbf{(P\_1,1.69.3)} savarṇe aṇgrahaṇam aparibhāṣyam . \\
\noindent \textbf{(P\_1,1.69.3)} ākṛtigrahaṇāt ananyatvam bhaviṣyati . \\
\noindent \textbf{(P\_1,1.69.3)} ananyākṛtiḥ akārasya ākārasya ca . \\
\noindent \textbf{(P\_1,1.69.3)} anekāntaḥ hi ananyatvakaraḥ . \\
\noindent \textbf{(P\_1,1.69.3)} yaḥ hi anekāntena bhedaḥ na asau anyatvam karoti . \\
\noindent \textbf{(P\_1,1.69.3)} tat yathā : na yaḥ goḥ ca goḥ ca bhedaḥ saḥ anyatvam karoti . \\
\noindent \textbf{(P\_1,1.69.3)} yaḥ tu khalu goḥ ca aśvasya ca bhedaḥ saḥ anyatvam karoti . \\
\noindent \textbf{(P\_1,1.69.3)} tadvat ca halgrahaṇeṣu . \\
\noindent \textbf{(P\_1,1.69.3)} evam ca kṛtvā ca halgrahaṇeṣu siddham bhavati . \\
\noindent \textbf{(P\_1,1.69.3)} jhalaḥ jhali : avāttām avāttam avātta yatra etat na asti aṇ savarṇān gṛhṇāti iti . \\
\noindent \textbf{(P\_1,1.69.3)} anekāntaḥ hi ananyatvakaraḥ iti uktārtham . \\
\noindent \textbf{(P\_1,1.69.3)} drutavilambitayoḥ ca anupadeśāt . \\
\noindent \textbf{(P\_1,1.69.3)} drutavilambitayoḥ ca anupadeśāt manyāmahe ākṛtigrahaṇāt siddham iti . \\
\noindent \textbf{(P\_1,1.69.3)} yat ayam kasyām cit vṛttau varṇān upadiśya sarvatra kṛtī bhavati . \\
\noindent \textbf{(P\_1,1.69.3)} asti prayojanam etat . \\
\noindent \textbf{(P\_1,1.69.3)} kim tarhi iti . \\
\noindent \textbf{(P\_1,1.69.3)} vṛttipṛthaktvam tu na upapadyate . \\
\noindent \textbf{(P\_1,1.69.3)} vṛtteḥ tu pṛthaktvam na upapadyate . \\
\noindent \textbf{(P\_1,1.69.3)} tasmāt tatra taparanirdeśāt siddham . \\
\noindent \textbf{(P\_1,1.69.3)} tasmāt tatra taparanirdeśaḥ kartavyaḥ . \\
\noindent \textbf{(P\_1,1.69.3)} na kartavyaḥ . \\
\noindent \textbf{(P\_1,1.69.3)} kriyate etat nyāse eva : ataḥ bhisaḥ ais iti . \\
\noindent \textbf{(P\_1,1.70.1)} ayuktaḥ ayam nirdeśaḥ . \\
\noindent \textbf{(P\_1,1.70.1)} tat iti anena kālaḥ pratinirdiśyate tat iti ayam ca varṇaḥ . \\
\noindent \textbf{(P\_1,1.70.1)} tatra ayuktam varṇasya kālena saha sāmanādhikaraṇyam . \\
\noindent \textbf{(P\_1,1.70.1)} katham tarhi nirdeśaḥ kartavyaḥ . \\
\noindent \textbf{(P\_1,1.70.1)} tatkālakālasya iti . \\
\noindent \textbf{(P\_1,1.70.1)} kim idam tatkālakālasya iti . \\
\noindent \textbf{(P\_1,1.70.1)} tasya kālaḥ tatkālaḥ , tatkālaḥ kālaḥ yasya saḥ ayam tatkālakālaḥ , tatkālakālasya iti . \\
\noindent \textbf{(P\_1,1.70.1)} saḥ tarhi tathā nirdeśaḥ kartavyaḥ . \\
\noindent \textbf{(P\_1,1.70.1)} na kartavyaḥ . \\
\noindent \textbf{(P\_1,1.70.1)} uttarapadalopaḥ atra draṣṭavyaḥ . \\
\noindent \textbf{(P\_1,1.70.1)} tat yathā uṣṭramukham iva mukham asya uṣṭramukhaḥ , kharamukhaḥ . \\
\noindent \textbf{(P\_1,1.70.1)} evam tatkālakālaḥ tatkālaḥ , tatkālasya iti . \\
\noindent \textbf{(P\_1,1.70.1)} atha vā sāhacaryāt tācchabdyam bhaviṣyati . \\
\noindent \textbf{(P\_1,1.70.1)} kālasahacaritaḥ varṇaḥ api kālaḥ eva. \\
\noindent \textbf{(P\_1,1.70.2)} kim punaḥ idam niyamārtham āhosvit prāpakam . \\
\noindent \textbf{(P\_1,1.70.2)} katham ca niyamārtham syāt katham vā prāpakam . \\
\noindent \textbf{(P\_1,1.70.2)} yadi atra aṇgrahaṇam anuvartate tataḥ niyamārtham . \\
\noindent \textbf{(P\_1,1.70.2)} atha nivṛttam tataḥ prāpakam . \\
\noindent \textbf{(P\_1,1.70.2)} kaḥ ca atra viśeṣaḥ . \\
\noindent \textbf{(P\_1,1.70.2)} taparaḥ tatkālasya iti niyamāṛtham iti cet dīrghagrahaṇe svarabhinnāgrahaṇam . \\
\noindent \textbf{(P\_1,1.70.2)} taparaḥ tatkālasya iti niyamāṛtham iti cet dīrghagrahaṇe svarabhinnānām grahaṇam na prāpnoti . \\
\noindent \textbf{(P\_1,1.70.2)} keṣām . \\
\noindent \textbf{(P\_1,1.70.2)} udāttānudāttasvaritānām . \\
\noindent \textbf{(P\_1,1.70.2)} astu tarhi prāpakam . \\
\noindent \textbf{(P\_1,1.70.2)} prāpakam iti cet hrasvagrahaṇe dīrghaplutapratiṣedhaḥ . \\
\noindent \textbf{(P\_1,1.70.2)} prāpakam iti cet hrasvagrahaṇe dīrghaplutayoḥ tu pratiṣedhaḥ vaktavyaḥ . \\
\noindent \textbf{(P\_1,1.70.2)} vipratiṣedhāt siddham . \\
\noindent \textbf{(P\_1,1.70.2)} aṇ savarṇān gṛhṇāti iti etat astu taparaḥ tatkālasya iti vā . \\
\noindent \textbf{(P\_1,1.70.2)} taparaḥ tatkālasya iti etat bhavati vipratiṣedhena . \\
\noindent \textbf{(P\_1,1.70.2)} aṇ savarṇān gṛhṇāti iti asya avakāśaḥ hrasvāḥ ataparāḥ aṇaḥ . \\
\noindent \textbf{(P\_1,1.70.2)} taparaḥ tatkālasya iti asya avakāśaḥ dīrghāḥ taparāḥ . \\
\noindent \textbf{(P\_1,1.70.2)} hrasveṣu tapareṣu ubhayam prāpnoti . \\
\noindent \textbf{(P\_1,1.70.2)} taparaḥ tatkālasya iti etat bhavati vipratiṣedhena . \\
\noindent \textbf{(P\_1,1.70.2)} yadi evam drutāyām taparakaraṇe madhyamavilambitayoḥ upasaṅkhyānam kālabhedāt . drutāyām taparakaraṇe madhyamavilambitayoḥ upasaṅkhyānam kartavyam tathā madhyamāyām drutavilambitayoḥ tathā vilambitāyām drutamadhyamayoḥ . \\
\noindent \textbf{(P\_1,1.70.2)} kim punaḥ kāraṇam na sidhyati . \\
\noindent \textbf{(P\_1,1.70.2)} kālabhedāt . \\
\noindent \textbf{(P\_1,1.70.2)} ye hi drutāyām vṛttau varṇāḥ tribhāgādhikāḥ te madhyamāyām . \\
\noindent \textbf{(P\_1,1.70.2)} ye madhyamāyām varṇāḥ tribhāgādhikāḥ te vilambitāyām . \\
\noindent \textbf{(P\_1,1.70.2)} siddham tu avasthitāḥ varṇāḥ vaktuḥ cirāciravacanāt vṛttayaḥ viśiṣyante . siddham etat . \\
\noindent \textbf{(P\_1,1.70.2)} katham . \\
\noindent \textbf{(P\_1,1.70.2)} avasthitāḥ varṇāḥ drutamadhyamavilambitāsu . \\
\noindent \textbf{(P\_1,1.70.2)} kiṅkṛtaḥ tu vṛttiviśeṣaḥ . \\
\noindent \textbf{(P\_1,1.70.2)} vaktuḥ cirāciravacanāt vṛttayaḥ viśiṣyante . \\
\noindent \textbf{(P\_1,1.70.2)} vaktā kaḥ cit āśvabhidhāyī bhavati , āśu varṇān abhidhatte . \\
\noindent \textbf{(P\_1,1.70.2)} kaḥ cit cireṇa kaḥ cit ciratareṇa . \\
\noindent \textbf{(P\_1,1.70.2)} tat yathā : tam eva adhvānam kaḥ cit āśu gacchati kaḥ cit cireṇa gacchati kaḥ cit ciratareṇa gacchati . \\
\noindent \textbf{(P\_1,1.70.2)} rathikaḥ āśu gacchati āśvikaḥ cireṇa padātiḥ ciratareṇa . \\
\noindent \textbf{(P\_1,1.70.2)} viṣamaḥ upanyāsaḥ . \\
\noindent \textbf{(P\_1,1.70.2)} adhikaraṇam atra adhvā vrajikriyāyāḥ . \\
\noindent \textbf{(P\_1,1.70.2)} tatra ayuktam yat adhikaraṇasya vṛddhihrāsau syātām . \\
\noindent \textbf{(P\_1,1.70.2)} evam tarhi sphoṭaḥ śabdaḥ dhvaniḥ śabdaguṇaḥ . \\
\noindent \textbf{(P\_1,1.70.2)} katham . \\
\noindent \textbf{(P\_1,1.70.2)} bheryāghātavat . \\
\noindent \textbf{(P\_1,1.70.2)} tat yathā bheryāghātaḥ bherīm āhatya kaḥ cit viṃśati padāni gacchati kaḥ cit triṃśat kaḥ cit catvāriṃśat . \\
\noindent \textbf{(P\_1,1.70.2)} sphoṭaḥ ca tāvān eva bhavati . \\
\noindent \textbf{(P\_1,1.70.2)} dhvanikṛtā vṛddhiḥ . \\
\noindent \textbf{(P\_1,1.70.2)} dhvaniḥ sphoṭaḥ ca śabdānām dhvaniḥ tu khalu lakṣyate | alpaḥ mahān ca keṣām cit ubhayam tat svabhāvataḥ . \\
\noindent \textbf{(P\_1,1.71)} ādiḥ antyena saha iti asampratyayaḥ sañjñinaḥ anirdeśāt . \\
\noindent \textbf{(P\_1,1.71)} ādiḥ antyena saha iti asampratyayaḥ . \\
\noindent \textbf{(P\_1,1.71)} kim kāraṇam . \\
\noindent \textbf{(P\_1,1.71)} sañjñinaḥ anirdeśāt . \\
\noindent \textbf{(P\_1,1.71)} na hi sañjñinaḥ nirdiśyante . \\
\noindent \textbf{(P\_1,1.71)} siddham tu ādiḥ itā saha tanmadhyasya iti vacanāt . \\
\noindent \textbf{(P\_1,1.71)} siddham etat . \\
\noindent \textbf{(P\_1,1.71)} katham . \\
\noindent \textbf{(P\_1,1.71)} ādiḥ antyena saha itā gṛhyamāṇaḥ svasya ca rūpasya grāhakaḥ tanmadhyānām ca iti vaktavyam . \\
\noindent \textbf{(P\_1,1.71)} sambandhiśabdaiḥ vā tulyam . \\
\noindent \textbf{(P\_1,1.71)} sambandhiśabdaiḥ vā tulyam etat . \\
\noindent \textbf{(P\_1,1.71)} tat yathā sambandhiśabdāḥ : mātari vartitavyam , pitari śuśrūṣitavyam iti . \\
\noindent \textbf{(P\_1,1.71)} na ca ucyate svasyām mātari svasmin pitari iti sambandhāt ca gamyate yā yasya mātā yaḥ ca yasya pitā iti . \\
\noindent \textbf{(P\_1,1.71)} evam iha api ādiḥ antyaḥ iti sambandhiśabdau etau . \\
\noindent \textbf{(P\_1,1.71)} tatra sambandhāt etat gantavyam : yam prati ādiḥ antyaḥ iti ca bhavati tasya grahaṇam bhavati svasya ca rūpasya iti . \\
\noindent \textbf{(P\_1,1.72.1)} iha kasmāt na bhavati : ikaḥ yaṇ aci : dadhi atra madhu atra . \\
\noindent \textbf{(P\_1,1.72.1)} astu . \\
\noindent \textbf{(P\_1,1.72.1)} alaḥ antyasya vidhayaḥ bhavanti iti antyasya bhaviṣyati . \\
\noindent \textbf{(P\_1,1.72.1)} na evam śakyam . \\
\noindent \textbf{(P\_1,1.72.1)} ye anekālaḥ ādeśāḥ teṣu doṣaḥ syāt : ecaḥ ayavāyāvaḥ iti . \\
\noindent \textbf{(P\_1,1.72.1)} na eṣaḥ doṣaḥ . \\
\noindent \textbf{(P\_1,1.72.1)} yathā eva prakṛtitaḥ tadantavidhiḥ bhavati evam ādeśataḥ api bhaviṣyati . \\
\noindent \textbf{(P\_1,1.72.1)} tatra ejantasya ayādyantā ādeśāḥ bhaviṣyanti . \\
\noindent \textbf{(P\_1,1.72.1)} yadi ca evam kva cit vairūpyam tatra doṣaḥ syāt . \\
\noindent \textbf{(P\_1,1.72.1)} api ca antaraṅgabahiraṅge na prakalpyeyātām . \\
\noindent \textbf{(P\_1,1.72.1)} tatra kaḥ doṣaḥ . \\
\noindent \textbf{(P\_1,1.72.1)} syonaḥ , syonā : antaraṅgalakṣaṇasya yaṇādeśasya bahiraṅgalakṣaṇaḥ guṇaḥ bādhakaḥ prasajyeta . \\
\noindent \textbf{(P\_1,1.72.1)} ūnaśabdam hi āśritya yaṇādeśaḥ naśabdam āśritya guṇaḥ . \\
\noindent \textbf{(P\_1,1.72.1)} alvidhiḥ ca na prakalpeta : dyauḥ , panthāḥ , saḥ iti . \\
\noindent \textbf{(P\_1,1.72.1)} tasmāt prakṛte tadantavidhiḥ iti vaktavyam . \\
\noindent \textbf{(P\_1,1.72.1)} na vaktavyam . \\
\noindent \textbf{(P\_1,1.72.1)} yena iti karaṇe eṣā tṛtīyā anyena ca anyasya vidhiḥ bhavati . \\
\noindent \textbf{(P\_1,1.72.1)} tat yathā : devadattasya samāśam śarāvaiḥ odanena ca yajñadattaḥ pratividhatte , tathā saṅgrāmam hastyaśvarathapadātibhiḥ . \\
\noindent \textbf{(P\_1,1.72.1)} evam iha api acā dhātoḥ yatam vidhatte . \\
\noindent \textbf{(P\_1,1.72.1)} akāreṇa prātipadikasya iñam vidhatte . \\
\noindent \textbf{(P\_1,1.72.2)} yena vidhiḥ tadantasya iti cet grahaṇopādhīnām tadantopādhiprasaṅgaḥ . \\
\noindent \textbf{(P\_1,1.72.2)} yena vidhiḥ tadantasya iti cet grahaṇopādhīnām tadantopādhitāprasaṅgaḥ . \\
\noindent \textbf{(P\_1,1.72.2)} ye grahaṇopādhayaḥ te api tadantopādhayaḥ syuḥ . \\
\noindent \textbf{(P\_1,1.72.2)} tatra kaḥ doṣaḥ . \\
\noindent \textbf{(P\_1,1.72.2)} utaḥ ca pratyayāt asaṃyogapūrvāt iti asaṃyogapūrvagrahaṇam ukārāntviśeṣaṇam syāt . \\
\noindent \textbf{(P\_1,1.72.2)} tatra kaḥ doṣaḥ . \\
\noindent \textbf{(P\_1,1.72.2)} asaṃyogapūrvagrahaṇena iha eva paryudāsaḥ syāt:: akṣṇuhi takṣṇuhi iti . \\
\noindent \textbf{(P\_1,1.72.2)} iha na syāt : āpnuhi śaknuhi iti . \\
\noindent \textbf{(P\_1,1.72.2)} tathā ut oṣthyapūrvasya iti oṣṭhyapūrvagrahaṇam ṝkārāntaviśeṣaṇam syāt . \\
\noindent \textbf{(P\_1,1.72.2)} tatra kaḥ doṣaḥ . \\
\noindent \textbf{(P\_1,1.72.2)} oṣṭhyapūrvagrahaṇena iha ca prasajyeta : saṅkīrṇam iti . \\
\noindent \textbf{(P\_1,1.72.2)} iha ca na syāt : nipūrtāḥ piṇḍāḥ iti . \\
\noindent \textbf{(P\_1,1.72.2)} siddham tu viśeṣaṇaviśeṣyayoḥ yatheṣṭatvāt . \\
\noindent \textbf{(P\_1,1.72.2)} siddham etat . \\
\noindent \textbf{(P\_1,1.72.2)} katham . \\
\noindent \textbf{(P\_1,1.72.2)} yatheṣṭam viśeṣaṇaviśeṣyayoḥ yogaḥ bhavati . \\
\noindent \textbf{(P\_1,1.72.2)} yāvatā yatheṣṭam iha tāvat : utaḥ ca pratyayāt asaṃyogapūrvāt iti na asaṃyogapūrvagrahaṇena ukārāntam viśeṣyate . \\
\noindent \textbf{(P\_1,1.72.2)} kim tarhi . \\
\noindent \textbf{(P\_1,1.72.2)} ukāraḥ eva viśeṣyate : ukāraḥ yaḥ asaṃyogapūrvaḥ tadantāt pratyayāt iti . \\
\noindent \textbf{(P\_1,1.72.2)} tathā ut oṣthyapūrvasya iti na oṣṭhapūrvagrahaṇena ṝkārāntam viśeṣyate . \\
\noindent \textbf{(P\_1,1.72.2)} kim tarhi . \\
\noindent \textbf{(P\_1,1.72.2)} ṝkāraḥ eva viśeṣyate : ṝkāraḥ yaḥ oṣṭhyapūrvaḥ tadantasya dhātoḥ iti . \\
\noindent \textbf{(P\_1,1.72.3)} samāsapratyayavidhau pratiṣedhaḥ . \\
\noindent \textbf{(P\_1,1.72.3)} samāsavidhau pratyayavidhau ca pratiṣedhaḥ vaktavyaḥ . \\
\noindent \textbf{(P\_1,1.72.3)} samāsavidhau tāvat : dvitīyā śritādibhiḥ samasyate : kaṣṭaśritaḥ , narakaśritaḥ . \\
\noindent \textbf{(P\_1,1.72.3)} kaṣṭam paramaśrita iti atra mā bhūt . \\
\noindent \textbf{(P\_1,1.72.3)} pratyayavidhau : naḍasya apatyam nāḍāyanaḥ . \\
\noindent \textbf{(P\_1,1.72.3)} iha na bhavati : sūtranaḍasya apatyam sautranāḍiḥ . \\
\noindent \textbf{(P\_1,1.72.3)} kim aviśeṣeṇa . \\
\noindent \textbf{(P\_1,1.72.3)} na iti āha . \\
\noindent \textbf{(P\_1,1.72.3)} ugidvarṇagrahaṇavarjam . \\
\noindent \textbf{(P\_1,1.72.3)} ugidgrahaṇam varṇagrahaṇam ca varjayitvā . \\
\noindent \textbf{(P\_1,1.72.3)} ugidgrahaṇam : bhavatī , atibhavatī mahatī , atimahatī . \\
\noindent \textbf{(P\_1,1.72.3)} varṇagrahaṇam : ataḥ iñ : dākṣiḥ , plākṣiḥ . \\
\noindent \textbf{(P\_1,1.72.3)} asti ca idānīm kaḥ cit kevalaḥ akāraḥ prātipadikam yadarthaḥ vidhiḥ syāt . \\
\noindent \textbf{(P\_1,1.72.3)} asti iti āha . \\
\noindent \textbf{(P\_1,1.72.3)} atateḥ ḍaḥ : aḥ , tasya apatyam : ataḥ iñ iḥ . \\
\noindent \textbf{(P\_1,1.72.3)} akacśnamvataḥ sarvanāmāvyayadhātuvidhau upasaṅkhyānam . akacvataḥ sarvanāmāvyayavidhau śnamvataḥ dhātuvidhau upasaṅkhyānam kartavyam . \\
\noindent \textbf{(P\_1,1.72.3)} akacvataḥ : sarvake viśvake . \\
\noindent \textbf{(P\_1,1.72.3)} avyayavidhau : uccakaiḥ nīcakaiḥ . \\
\noindent \textbf{(P\_1,1.72.3)} śnamvataḥ : bhinatti chinatti . \\
\noindent \textbf{(P\_1,1.72.3)} kim punaḥ kāraṇam na sidhyati . \\
\noindent \textbf{(P\_1,1.72.3)} iha tasya vā grahaṇam bhavati tadantasya vā . \\
\noindent \textbf{(P\_1,1.72.3)} na ca idam tat na api tadantam . \\
\noindent \textbf{(P\_1,1.72.3)} siddham tu tadantāntavacanāt . \\
\noindent \textbf{(P\_1,1.72.3)} siddham etat . \\
\noindent \textbf{(P\_1,1.72.3)} katham . \\
\noindent \textbf{(P\_1,1.72.3)} tadantāntavacanāt . \\
\noindent \textbf{(P\_1,1.72.3)} tadantāntasya iti vaktavyam . \\
\noindent \textbf{(P\_1,1.72.3)} kim idam tadantāntasya iti . \\
\noindent \textbf{(P\_1,1.72.3)} tasya antaḥ tadantaḥ , tadantaḥ antaḥ yasya tat idam tadantāntam , tadantāntasya iti . \\
\noindent \textbf{(P\_1,1.72.3)} saḥ tarhi tathā nirdeśaḥ kartavyaḥ . \\
\noindent \textbf{(P\_1,1.72.3)} na kartavyaḥ . \\
\noindent \textbf{(P\_1,1.72.3)} uttarapadalopaḥ atra draṣṭavyaḥ . \\
\noindent \textbf{(P\_1,1.72.3)} tat yathā : uṣṭramukham iva mukham asya : uṣṭramukhaḥ , kharamukhaḥ . \\
\noindent \textbf{(P\_1,1.72.3)} evam iha api tadantaḥ antaḥ yasya tadantasya iti . \\
\noindent \textbf{(P\_1,1.72.3)} tadekadeśavijñānāt vā siddham . tadekadeśavijñānāt vā punaḥ siddham etat . \\
\noindent \textbf{(P\_1,1.72.3)} tadekadeśabhūtaḥ tadgrahaṇena gṛhyate . \\
\noindent \textbf{(P\_1,1.72.3)} tat yathā gaṅgā yamunā devadattā iti . \\
\noindent \textbf{(P\_1,1.72.3)} anekā nadī gaṅgām yamunām ca praviṣṭā gaṅgāyamunāgrahaṇena gṛhyate . \\
\noindent \textbf{(P\_1,1.72.3)} tathā devadattāsthaḥ garbhaḥ devadattāgrahaṇena gṛhyate . \\
\noindent \textbf{(P\_1,1.72.3)} viṣamaḥ upanyāsaḥ . \\
\noindent \textbf{(P\_1,1.72.3)} iha ke cit śabdāḥ aktaparimāṇānām arthānām vācakāḥ bhavanti ye ete saṅkhyāśabdāḥ parimāṇaśabdāḥ ca . \\
\noindent \textbf{(P\_1,1.72.3)} pañca sapta iti : ekena api apāye na bhavanti . \\
\noindent \textbf{(P\_1,1.72.3)} droṇaḥ khārī āḍhakam iti : na eva adhike bhavanti na nyūne . \\
\noindent \textbf{(P\_1,1.72.3)} ke cit yāvat eva tat bhavati tāvat eva āhuḥ ye ete jātiśabdāḥ guṇaśabdāḥ ca . \\
\noindent \textbf{(P\_1,1.72.3)} tailam ghṛtam iti : khāryām api bhavanti droṇe api . \\
\noindent \textbf{(P\_1,1.72.3)} śuklaḥ nīlaḥ kṛṣṇaḥ iti : himavati api bhavati vaṭakaṇikāmātre api dravye . \\
\noindent \textbf{(P\_1,1.72.3)} imāḥ ca api sañjñāḥ aktaparimāṇānām arthānām kriyante . \\
\noindent \textbf{(P\_1,1.72.3)} tāḥ kena adhikasya syuḥ . \\
\noindent \textbf{(P\_1,1.72.3)} evam tarhi ācāryapravṛttiḥ jñāpayati tadekadeśabhūtam tadgrahaṇena gṛhyate iti yat ayam na idamadasoḥ akoḥ iti sakakārayoḥ idamadasoḥ pratiṣedham śāsti . \\
\noindent \textbf{(P\_1,1.72.3)} katham kṛtvā jñāpakam . \\
\noindent \textbf{(P\_1,1.72.3)} idamadasoḥ kāryam ucyamānam kaḥ prasaṅgaḥ yat sakakārayoḥ syāt . \\
\noindent \textbf{(P\_1,1.72.3)} paśyati tu ācāryaḥ tadekadeśabhūtam tadgrahaṇena gṛhyate iti . \\
\noindent \textbf{(P\_1,1.72.3)} tataḥ sakakārayoḥ pratiṣedham śāsti . \\
\noindent \textbf{(P\_1,1.72.4)} kāni punaḥ asya yogasya prayojanāni . \\
\noindent \textbf{(P\_1,1.72.4)} prayojanam sarvanāmāvyayasañjñāyām . \\
\noindent \textbf{(P\_1,1.72.4)} sarvanāmāvyayasañjñāyām prayojanam : sarve paramasarve viśve paramaviśve , uccaiḥ , paramoccaiḥ , nīcaiḥ , paramanīcaiḥ iti . \\
\noindent \textbf{(P\_1,1.72.4)} upapadavidhau bhayāḍhyādigrahaṇam . \\
\noindent \textbf{(P\_1,1.72.4)} upapadavidhau bhayāḍhyādigrahaṇam prayojanam : bhayaṅkaraḥ , abhayaṅkaraḥ , āḍhyaṅkaraṇam , khāḍyaṅkaraṇam . \\
\noindent \textbf{(P\_1,1.72.4)} ṅībvidhau ugidgrahaṇam . \\
\noindent \textbf{(P\_1,1.72.4)} ṅībvidhau ugidgrahaṇam prayojanam : bhavatī , atibhavatī mahatī , atimahatī . \\
\noindent \textbf{(P\_1,1.72.4)} pratiṣedhe svasrādigrahaṇam . \\
\noindent \textbf{(P\_1,1.72.4)} pratiṣedhe svasrādigrahaṇam prayojanam : svasā paramasvasā duhitā paramaduhitā . \\
\noindent \textbf{(P\_1,1.72.4)} aparimāṇabistādigrahaṇam ca pratiṣedhe . \\
\noindent \textbf{(P\_1,1.72.4)} aparimāṇabistādigrahaṇam ca pratiṣedhe prayojanam . \\
\noindent \textbf{(P\_1,1.72.4)} aparimāṇabistācitakambalebhyaḥ na taddhitaluki : dvibistā dviparamabistā tribistā triparamabistā dvyācitā dviparamācitā . \\
\noindent \textbf{(P\_1,1.72.4)} diti . \\
\noindent \textbf{(P\_1,1.72.4)} ditigrahaṇam ca prayojanam . \\
\noindent \textbf{(P\_1,1.72.4)} diteḥ apatyam daityaḥ , aditeḥ apatyam ādityaḥ . \\
\noindent \textbf{(P\_1,1.72.4)} dityadityāditya iti aditigrahaṇam na kartavyam bhavati . \\
\noindent \textbf{(P\_1,1.72.4)} roṇyāḥ aṇ . \\
\noindent \textbf{(P\_1,1.72.4)} roṇyāḥ aṇgrahaṇam ca prayojanam : ājakaroṇaḥ , saiṃhakaroṇaḥ . \\
\noindent \textbf{(P\_1,1.72.4)} tasya ca . \\
\noindent \textbf{(P\_1,1.72.4)} tasya ca iti vaktavyam : rauṇaḥ . \\
\noindent \textbf{(P\_1,1.72.4)} kim punaḥ kāraṇam na sidhyati . \\
\noindent \textbf{(P\_1,1.72.4)} tadantāt ca tadantavidhinā siddham kevalāt ca vyapdeśivadbhāvena . \\
\noindent \textbf{(P\_1,1.72.4)} vyapdeśivadbhāvaḥ aprātipadikena . \\
\noindent \textbf{(P\_1,1.72.4)} kim punaḥ kāraṇam vyapdeśivadbhāvaḥ aprātipadikena . \\
\noindent \textbf{(P\_1,1.72.4)} iha : sūtrāntāt ṭhak bhavati daśāntāt ḍaḥ bhavati iti : kevalāt utpattiḥ mā bhūt iti . \\
\noindent \textbf{(P\_1,1.72.4)} na etat asti prayojanam . \\
\noindent \textbf{(P\_1,1.72.4)} siddham atra tadantāt ca tadantavidhinā kevalāt ca vyapdeśivadbhāvena . \\
\noindent \textbf{(P\_1,1.72.4)} saḥ ayam evam siddhe sati yat antagrahaṇam karoti tat jñāpayati ācāryaḥ sūtrāntāt eva daśāntāt eva iti . \\
\noindent \textbf{(P\_1,1.72.4)} na atra tadantāt utpattiḥ prāpnoti . \\
\noindent \textbf{(P\_1,1.72.4)} idānīm eva hi uktam : samāsapratyayavidhau pratiṣedhaḥ iti . \\
\noindent \textbf{(P\_1,1.72.4)} sā tarhi eṣā paribhāṣā kartavyā . \\
\noindent \textbf{(P\_1,1.72.4)} na kartavyā . \\
\noindent \textbf{(P\_1,1.72.4)} ācāryapravṛttiḥ jñāpayati vyapdeśivadbhāvaḥ aprātipadikena iti yat ayam pūrvāt iniḥ sapūrvāt ca iti āha . \\
\noindent \textbf{(P\_1,1.72.4)} na etat asti prayojanam . \\
\noindent \textbf{(P\_1,1.72.4)} asti hi anyat etasya vacane prayojanam . \\
\noindent \textbf{(P\_1,1.72.4)} kim . \\
\noindent \textbf{(P\_1,1.72.4)} sapūrvāt pūrvāt inim vakṣyāmi iti . \\
\noindent \textbf{(P\_1,1.72.4)} yat tarhi yogavibhāgam karoti . \\
\noindent \textbf{(P\_1,1.72.4)} itarathā hi pūrvāt sapūrvāt iti eva brūyāt . \\
\noindent \textbf{(P\_1,1.72.4)} kim punaḥ ayam asya eva śeṣaḥ : tasya ca iti . \\
\noindent \textbf{(P\_1,1.72.4)} na iti āha . \\
\noindent \textbf{(P\_1,1.72.4)} yat ca anukrāntam yat ca anukraṃsayte sarvasya eva śeṣaḥ tasya ca iti . \\
\noindent \textbf{(P\_1,1.72.4)} rathasītāhalebhyaḥ yadvidhau . \\
\noindent \textbf{(P\_1,1.72.4)} rathasītāhalebhyaḥ yadvidhau prayojanam : rathyaḥ , paramarathyaḥ , sītyam , paramasītyam , halyā paramahalyā . \\
\noindent \textbf{(P\_1,1.72.4)} susarvārdhadikśabdebhyaḥ janapadasya . \\
\noindent \textbf{(P\_1,1.72.4)} susarvārdhadikśabdebhyaḥ janapadasya prayojanam : supāñcālakaḥ , sumāgadhakaḥ . \\
\noindent \textbf{(P\_1,1.72.4)} su . \\
\noindent \textbf{(P\_1,1.72.4)} sarva : sarvapāñcālakaḥ , sarvamāgadhakaḥ . \\
\noindent \textbf{(P\_1,1.72.4)} sarva . \\
\noindent \textbf{(P\_1,1.72.4)} ardha : ardhapāñcālakaḥ , ardhamāgadhakaḥ . \\
\noindent \textbf{(P\_1,1.72.4)} ardha . \\
\noindent \textbf{(P\_1,1.72.4)} dikśabda : pūrvapāñcālakaḥ , pūrvamāgadhakaḥ . \\
\noindent \textbf{(P\_1,1.72.4)} ṛtoḥ vṛddhimadvidhau avayavānām . \\
\noindent \textbf{(P\_1,1.72.4)} ṛtoḥ vṛddhimadvidhau avayavānām prayojanam : pūrvaśāradam , aparaśāradam , pūrvanaidāgham , aparanaidāgham . \\
\noindent \textbf{(P\_1,1.72.4)} ṭhañvidhau saṅkhyāyāḥ . \\
\noindent \textbf{(P\_1,1.72.4)} ṭhañvidhau saṅkhyāyāḥ prayojanam : dviṣāṣṭikam , pañcaṣāṣṭikam . \\
\noindent \textbf{(P\_1,1.72.4)} dharmāt nañaḥ . \\
\noindent \textbf{(P\_1,1.72.4)} dharmāt nañaḥ prayojanam : dharmam carati dhārmikaḥ , adharmam carati ādharmikaḥ . \\
\noindent \textbf{(P\_1,1.72.4)} adharmāt ca iti na vaktavyam bhavati \\
\noindent \textbf{(P\_1,1.72.5)} padāṅgādhikāre tasya ca taduttarapadasya ca . \\
\noindent \textbf{(P\_1,1.72.5)} padāṅgādhikāre tasya ca taduttarapadasya ca iti vaktavyam . \\
\noindent \textbf{(P\_1,1.72.5)} padādhikāre kim prayojanam . \\
\noindent \textbf{(P\_1,1.72.5)} prayojanam iṣṭikeṣīkāmālānām citatūlabhāriṣu : iṣṭakacitam cinvīta , pakveṣṭikcitam cinvīta , iṣīkatūlena muñjeṣīkatūlena mālabhāriṇī kanyā , utpalamālabhāriṇī kanyā . \\
\noindent \textbf{(P\_1,1.72.5)} aṅgādhikāre kim prayojanam . \\
\noindent \textbf{(P\_1,1.72.5)} mahadapsvasṛtṝṇām dīrghavidhau . \\
\noindent \textbf{(P\_1,1.72.5)} mahadapsvasṛtṝṇām dīrghavidhau prayojanam : mahān , paramamahān . \\
\noindent \textbf{(P\_1,1.72.5)} mahat . \\
\noindent \textbf{(P\_1,1.72.5)} ap : āpaḥ tiṣṭhanti , svāpaḥ tiṣṭhanti . \\
\noindent \textbf{(P\_1,1.72.5)} ap . \\
\noindent \textbf{(P\_1,1.72.5)} svasṛ : svasā svasārau svasāraḥ , paramasvasā paramasvasārau paramasvasāraḥ . \\
\noindent \textbf{(P\_1,1.72.5)} svasṛ . \\
\noindent \textbf{(P\_1,1.72.5)} naptṛ : naptā naptārau naptāraḥ . \\
\noindent \textbf{(P\_1,1.72.5)} evam paramanaptā paramanaptārau paramanaptāraḥ . \\
\noindent \textbf{(P\_1,1.72.5)} padyuṣmadasmadasthyādyanḍuhaḥ num . padbhāvaḥ prayojanam : divpadaḥ paśya . \\
\noindent \textbf{(P\_1,1.72.5)} asti ca idānīm kaḥ cit kevalaḥ pācchabhdaḥ yadarthaḥ vidhiḥ syāt . \\
\noindent \textbf{(P\_1,1.72.5)} na asti iti āha . \\
\noindent \textbf{(P\_1,1.72.5)} evam tarhi aṅgādhikāre prayojanam na asti iti kṛtvā padādhikārasya idam prayojanam uktam : himakāṣihatiṣu ca : yathā patkāṣiṇau patkāṣiṇaḥ evam paramapatkāṣiṇau paramapatkāṣiṇaḥ . \\
\noindent \textbf{(P\_1,1.72.5)} yadi tarhi padādhikāre pādasya tadantavidhiḥ bhavati pādasya pada ājyatigopahateṣu : yathā iha bhavati : pādena upahatam padopahatam atra api syāt : digdhapādena upahatam digdhapādopahatam . \\
\noindent \textbf{(P\_1,1.72.5)} evam tarhi aṅgādhikāre eva prayojanam . \\
\noindent \textbf{(P\_1,1.72.5)} nanu ca uktam na asti kevalaḥ pācchabdaḥ iti . \\
\noindent \textbf{(P\_1,1.72.5)} ayam asti pādayateḥ apratyayaḥ pāt : padā pade . \\
\noindent \textbf{(P\_1,1.72.5)} pad . \\
\noindent \textbf{(P\_1,1.72.5)} yuṣmat asmat : yūyam , vayam atiyūyam ativayam . \\
\noindent \textbf{(P\_1,1.72.5)} asthyādi : asthnā dadhnā sakthnā parmāsthnā paramadadhnā paramasakthnā . \\
\noindent \textbf{(P\_1,1.72.5)} anaḍuhaḥ num : anaḍvān , paramānaḍvān . \\
\noindent \textbf{(P\_1,1.72.5)} dyupathimathipuṅgosakhicaturanaḍuttrigrahaṇam . \\
\noindent \textbf{(P\_1,1.72.5)} dyupathimathipuṅgosakhicaturanaḍuttrigrahaṇam prayojanam : dyauḥ , sudyauḥ , panthāḥ , supanthāḥ , manthāḥ , sumanthāḥ , paramamanthāḥ , pumān paramapumān , gauḥ , sugauḥ , sakhā sakhāyau sakhāyaḥ , susakhā susakhāyau susakhāyaḥ , paramasakhā paramasakhāyau paramasakhāyaḥ , catvāraḥ paramacatvāraḥ , anaḍvāhaḥ , parmānaḍvāhaḥ , trayāṇām , paramatrayāṇām . \\
\noindent \textbf{(P\_1,1.72.5)} tyadādividhibhastrādistrīgrahaṇam ca . \\
\noindent \textbf{(P\_1,1.72.5)} tyadādividhibhastrādistrīgrahaṇam ca prayojanam : saḥ , atisaḥ , bhastrakā bhastrikā nirbhastrakā nirbhastrikā bahubhastrakā bahubhastrikā . \\
\noindent \textbf{(P\_1,1.72.5)} strīgrahaṇam ca prayojanam . \\
\noindent \textbf{(P\_1,1.72.5)} striyau striyaḥ rājastriyau rājastriyaḥ . \\
\noindent \textbf{(P\_1,1.72.5)} varṇagrahaṇam ca sarvatra . \\
\noindent \textbf{(P\_1,1.72.5)} varṇagrahaṇam ca sarvatra prayojanam . \\
\noindent \textbf{(P\_1,1.72.5)} kva sarvatra . \\
\noindent \textbf{(P\_1,1.72.5)} aṅgādhikāre ca anyatra ca . \\
\noindent \textbf{(P\_1,1.72.5)} anyatra udāhṛtam . \\
\noindent \textbf{(P\_1,1.72.5)} āngādhikāre : ataḥ dīrghaḥ yañi supi ca : iha eva syāt : ābhyām . \\
\noindent \textbf{(P\_1,1.72.5)} ghaṭābhyām iti atra na syāt . \\
\noindent \textbf{(P\_1,1.72.5)} pratyayagrahaṇam ca apañcamyāḥ . pratyayagrahaṇam ca apañcamyāḥ prayojanam : yañiñoḥ phak bhavati . \\
\noindent \textbf{(P\_1,1.72.5)} gārgyāyaṇaḥ vātsyāyanaḥ paramagārgyāyaṇaḥ paramavātsyāyanaḥ . \\
\noindent \textbf{(P\_1,1.72.5)} apañcamyāḥ iti kimartham . \\
\noindent \textbf{(P\_1,1.72.5)} dṛṣattīrṇā pariṣattīrṇā . \\
\noindent \textbf{(P\_1,1.72.5)} alā eva anarthakena na anyena anarthakena iti vaktavyam . \\
\noindent \textbf{(P\_1,1.72.5)} kim prayojanam . \\
\noindent \textbf{(P\_1,1.72.5)} hangrahaṇe plīhangrahaṇam mā bhūt . \\
\noindent \textbf{(P\_1,1.72.5)} udgrahaṇe garmudgrahaṇam . \\
\noindent \textbf{(P\_1,1.72.5)} strīgrahaṇe śastrīgrahaṇam . \\
\noindent \textbf{(P\_1,1.72.5)} saṅgrahaṇe pāyasam karoti iti mā bhūt . \\
\noindent \textbf{(P\_1,1.72.5)} kimartham idam ucyate na padāṅgādhikāre tasya ca taduttarapadasya ca iti eva siddham . \\
\noindent \textbf{(P\_1,1.72.5)} na ca idam tat na api taduttarapadam . \\
\noindent \textbf{(P\_1,1.72.5)} tat na vaktavyam bhavati . \\
\noindent \textbf{(P\_1,1.72.5)} kim punaḥ atra jyāyaḥ . \\
\noindent \textbf{(P\_1,1.72.5)} tadantavidhiḥ eva jyāyān . \\
\noindent \textbf{(P\_1,1.72.5)} idam api siddham bhavati : paramātimahān . \\
\noindent \textbf{(P\_1,1.72.5)} etat hi na eva tat na api taduttarapadam . \\
\noindent \textbf{(P\_1,1.72.5)} aninasmangrahaṇāni ca arthavatā ca anarthakena ca tadantavidhim prayojayanti . \\
\noindent \textbf{(P\_1,1.72.5)} an : rājñā iti arthavatā sāmnā iti anarthakena . \\
\noindent \textbf{(P\_1,1.72.5)} an . \\
\noindent \textbf{(P\_1,1.72.5)} in : daṇḍī* iti arthavatā vāgmī* iti anarthakena . \\
\noindent \textbf{(P\_1,1.72.5)} in . \\
\noindent \textbf{(P\_1,1.72.5)} as : supayāḥ iti arthavatā susrotāḥ iti anarthakena . \\
\noindent \textbf{(P\_1,1.72.5)} as . \\
\noindent \textbf{(P\_1,1.72.5)} man : suśarmā iti arthavatā suprathimā iti anarthakena . \\
\noindent \textbf{(P\_1,1.72.6)} yasmin vidhiḥ tadādau algrahaṇe . \\
\noindent \textbf{(P\_1,1.72.6)} algrahaṇeṣu yasmin vidhiḥ tadādau iti vaktavyam . \\
\noindent \textbf{(P\_1,1.72.6)} kim prayojanam . \\
\noindent \textbf{(P\_1,1.72.6)} aci śnudhātubhruvām yvoḥ iyaṅuvaṅau iti iha eva syāt : śriyau bhruvau . \\
\noindent \textbf{(P\_1,1.72.6)} śriyaḥ , bhruvaḥ iti atra na syāt . \\
\noindent \textbf{(P\_1,1.73.1)} vṛddhigrahaṇam kimartham. yasya acām ādiḥ tat vṛddham iti iyati ucyamāne dāttāḥ , rākṣitāḥ atra api prasajyeta . \\
\noindent \textbf{(P\_1,1.73.1)} vṛddhigrahaṇe punaḥ kriyamāṇe na doṣaḥ bhavati . \\
\noindent \textbf{(P\_1,1.73.1)} atha yasyagrahaṇam kimartham . \\
\noindent \textbf{(P\_1,1.73.1)} yasya iti vyapadeśāya . \\
\noindent \textbf{(P\_1,1.73.1)} atha ajgrahaṇam kimartham . \\
\noindent \textbf{(P\_1,1.73.1)} vṛddhiḥ yasya ādiḥ tat vṛddham iti iyati ucyamāne iha eva syāt : aitikāyanīyāḥ , aupagavīyāḥ . \\
\noindent \textbf{(P\_1,1.73.1)} iha na syāt : gārgīyāḥ , vātsīyāḥ iti . \\
\noindent \textbf{(P\_1,1.73.1)} ajgrahaṇe punaḥ kriyamāṇe na doṣaḥ bhavati . \\
\noindent \textbf{(P\_1,1.73.1)} atha ādigrahaṇam kimartham . \\
\noindent \textbf{(P\_1,1.73.1)} vṛddhiḥ yasya acām tat vṛddham iti iyati ucyamāne sabhāsannayane bhavaḥ sābhasannayanaḥ iti atra prasajyeta . \\
\noindent \textbf{(P\_1,1.73.1)} ādigrahaṇe punaḥ kriyamāṇe na doṣaḥ bhavati . \\
\noindent \textbf{(P\_1,1.73.1)} vṛddhasañjñāyām ajasanniveśāt anāditvam . \\
\noindent \textbf{(P\_1,1.73.1)} vṛddhasañjñāyām ajasanniveśāt ādiḥ iti etat na upapadyate . \\
\noindent \textbf{(P\_1,1.73.1)} na hi acām sanniveśaḥ asti . \\
\noindent \textbf{(P\_1,1.73.1)} nanu ca evam vijñāyate : ac eva ādiḥ ajādiḥ . \\
\noindent \textbf{(P\_1,1.73.1)} na evam śakyam . \\
\noindent \textbf{(P\_1,1.73.1)} iha eva prasajyeta : aupagavīyāḥ . \\
\noindent \textbf{(P\_1,1.73.1)} iha na syāt : gārgīyāḥ iti . \\
\noindent \textbf{(P\_1,1.73.1)} ekāntāditvam tarhi vijñāyate . \\
\noindent \textbf{(P\_1,1.73.1)} ekāntāditve ca sarvaprasaṅgaḥ . \\
\noindent \textbf{(P\_1,1.73.1)} iha api prasajyeta : sabhāsannayane bhavaḥ sābhasannayanaḥ iti . \\
\noindent \textbf{(P\_1,1.73.1)} siddham ajākṛtinirdeśāt . \\
\noindent \textbf{(P\_1,1.73.1)} siddham etat . \\
\noindent \textbf{(P\_1,1.73.1)} katham . \\
\noindent \textbf{(P\_1,1.73.1)} ajākṛtiḥ nirdiśyate . \\
\noindent \textbf{(P\_1,1.73.1)} evam api vyañjanaiḥ vyavahitatvāt na prāpnoti . \\
\noindent \textbf{(P\_1,1.73.1)} vyañjanasya avidyamānatvam yathā anyatra . \\
\noindent \textbf{(P\_1,1.73.1)} vyañjanasya avidyamānavadbhāvaḥ vaktavyaḥ yathā anyatra api bhavati vyañjanasya avidyamānavadbhāvaḥ . \\
\noindent \textbf{(P\_1,1.73.1)} kva anyatra . \\
\noindent \textbf{(P\_1,1.73.1)} svare . \\
\noindent \textbf{(P\_1,1.73.2)} vā nāmadheyasya . \\
\noindent \textbf{(P\_1,1.73.2)} vā nāmadheyasya vṛddhasañjñā vaktavyā : devadattīyāḥ , daivadattāḥ , yajñadattīyāḥ , yājñadattāḥ . \\
\noindent \textbf{(P\_1,1.73.2)} gotrottarapadasya ca . \\
\noindent \textbf{(P\_1,1.73.2)} gotrottarapadasya ca vṛddhasañjñā vaktavyā : kambalacārāyaṇīyāḥ , odanapāṇinīyāḥ , ghṛtarauḍhīyāḥ . \\
\noindent \textbf{(P\_1,1.73.2)} gotrāntāt vā asamastavat . \\
\noindent \textbf{(P\_1,1.73.2)} gotrāntāt vā asamastavat pratyayaḥ bhavati iti vaktavyam : etāni eva udāharaṇāni . \\
\noindent \textbf{(P\_1,1.73.2)} kim aviśeṣeṇa . \\
\noindent \textbf{(P\_1,1.73.2)} na iti āha . \\
\noindent \textbf{(P\_1,1.73.2)} jihvākātyaharitakātyavarjam . \\
\noindent \textbf{(P\_1,1.73.2)} jihvākātyam haritakātyam ca varjayitvā : jaihavākātāḥ , hāritakātāḥ . \\
\noindent \textbf{(P\_1,1.73.2)} kim punaḥ atra jyāyaḥ . \\
\noindent \textbf{(P\_1,1.73.2)} gotrāntāt vā asamastavat iti eva jyāyaḥ . \\
\noindent \textbf{(P\_1,1.73.2)} idam api siddham bhavati : piṅgalakāṇvasya chāttrāḥ paiṅgalakāṇvāḥ . \\
\noindent \textbf{(P\_1,1.74)} yasyācāmādigrahaṇam anuvartate utāho na . \\
\noindent \textbf{(P\_1,1.74)} kim ca ataḥ . \\
\noindent \textbf{(P\_1,1.74)} yadi anuvartate iha ca prasajyeta : tvatputrasya chāttrāḥ tvātputrāḥ , mātputrāḥ iha ca na syāt : tvadīyaḥ , madīyaḥ iti . \\
\noindent \textbf{(P\_1,1.74)} atha nivṛttam eṅ prācām deśe yasyācāmādigrahaṇam kartavyam . \\
\noindent \textbf{(P\_1,1.74)} evam tarhi anuvartate . \\
\noindent \textbf{(P\_1,1.74)} katham tvāputrāḥ , mātputrāḥ iti . \\
\noindent \textbf{(P\_1,1.74)} sambandham anuvartiṣyate . \\
\noindent \textbf{(P\_1,1.74)} vṛddhiḥ yasya acām ādiḥ tat vṛddham . \\
\noindent \textbf{(P\_1,1.74)} tyadādīni ca vṛddhasañjñāni bhavanti . \\
\noindent \textbf{(P\_1,1.74)} vṛddhiḥ yasya acām ādiḥ tat vṛddham . \\
\noindent \textbf{(P\_1,1.74)} eṅ prācām deśe . \\
\noindent \textbf{(P\_1,1.74)} yasyācāmādigrahaṇam anuvartate . \\
\noindent \textbf{(P\_1,1.74)} vṛddhigrahaṇam nivṛttam . \\
\noindent \textbf{(P\_1,1.74)} tat yathā kaḥ cit kāntāre samupasthite sārtham upādatte . \\
\noindent \textbf{(P\_1,1.74)} saḥ yadā niṣkāntārībhūtaḥ bhavati tadā sārtham jahāti \\
\noindent \textbf{(P\_1,1.75)} eṅ prācām deśe śaiṣikeṣu iti vaktavyam : saipurikī saipurikā skaunagarikī skaunagarikā iti . \\
\noindent \textbf{(P\_1,2.1.1)} ṅitkidvacane tayoḥ abhāvāt aprasiddhiḥ . \\
\noindent \textbf{(P\_1,2.1.1)} ṅitkidvacane tayoḥ abhāvāt , ṅakārakakārayoḥ abhāvāt , ṅittvakittvayoḥ aprasiddhiḥ . \\
\noindent \textbf{(P\_1,2.1.1)} satā hi abhisambandhaḥ śakyate kartum na ca atra ṅakārakakārau itau paśyāmaḥ . \\
\noindent \textbf{(P\_1,2.1.1)} tat yathā citraguḥ devadattaḥ iti : yasya tāḥ gāvaḥ santi saḥ eva tābhyām śabdābhyām śakyate abhisambandhum . \\
\noindent \textbf{(P\_1,2.1.1)} bhāvyete tarhi anena . \\
\noindent \textbf{(P\_1,2.1.1)} gāṅkuṭādibhyaḥ añṇit ṅit bhavati iti . \\
\noindent \textbf{(P\_1,2.1.1)} asaṃyogāt liṭ kit bhavati iti . \\
\noindent \textbf{(P\_1,2.1.1)} bhavati iti cet ādeśapratiṣedhaḥ . \\
\noindent \textbf{(P\_1,2.1.1)} bhavati iti cet ādeśasya pratiṣedhaḥ vaktavyaḥ . \\
\noindent \textbf{(P\_1,2.1.1)} ṅakārakakārau itau ādeśau prāpnutaḥ . \\
\noindent \textbf{(P\_1,2.1.1)} katham punaḥ itsañjñaḥ nāma ādeśaḥ syāt . \\
\noindent \textbf{(P\_1,2.1.1)} kim hi vacanāt na bhavati . \\
\noindent \textbf{(P\_1,2.1.1)} evam tarhi ṣaṣṭhīnirdiṣṭasya ādeśāḥ ucyante na ca atra ṣaṣṭhīm paśyāmaḥ . \\
\noindent \textbf{(P\_1,2.1.1)} gāṅkuṭādibhyaḥ iti eṣā pañcamī añṇit iti prathamāyāḥ ṣaṣṭhīm prakalpayiṣyati tasmāt iti uttarasya iti . \\
\noindent \textbf{(P\_1,2.1.1)} sañjñākaraṇam tarhi idam . \\
\noindent \textbf{(P\_1,2.1.1)} gāṅkuṭādibhyaḥ añṇit ṅitsañjñaḥ bhavati iti . \\
\noindent \textbf{(P\_1,2.1.1)} asaṃyogāt liṭ kitsañjñaḥ bhavati iti . \\
\noindent \textbf{(P\_1,2.1.1)} sañjñākaraṇe kṅidgrahaṇe asampratyayaḥ śabdabhedāt . \\
\noindent \textbf{(P\_1,2.1.1)} sañjñākaraṇe kṅidgrahaṇe asampratyayaḥ syāt . \\
\noindent \textbf{(P\_1,2.1.1)} kim kāraṇam . \\
\noindent \textbf{(P\_1,2.1.1)} śabdabhedāt . \\
\noindent \textbf{(P\_1,2.1.1)} anyaḥ hi śabdaḥ kṅiti iti anyaḥ kiti iti ṅiti iti ca . \\
\noindent \textbf{(P\_1,2.1.1)} tathā kidgrahaṇeṣu ṅidgrahaṇeṣu ca anayoḥ eva sampratyayaḥ syāt . \\
\noindent \textbf{(P\_1,2.1.1)} tadvadatideśaḥ tarhi ayam : gāṅkuṭādibhyaḥ añṇit ṅidvat bhavati iti . \\
\noindent \textbf{(P\_1,2.1.1)} asaṃyogāt liṭ kidvat bhavati iti . \\
\noindent \textbf{(P\_1,2.1.1)} saḥ tarhi vatinirdeśaḥ kartavyaḥ . \\
\noindent \textbf{(P\_1,2.1.1)} na hi antareṇa vatim atideśaḥ gamyate . \\
\noindent \textbf{(P\_1,2.1.1)} antareṇa api vatim atideśaḥ gamyate . \\
\noindent \textbf{(P\_1,2.1.1)} tat yathā : eṣaḥ brahmadattaḥ . \\
\noindent \textbf{(P\_1,2.1.1)} abrahmadattam brahmadattaḥ iti āha . \\
\noindent \textbf{(P\_1,2.1.1)} te manyāmahe brahmadattavat ayam bhavati iti . \\
\noindent \textbf{(P\_1,2.1.1)} evam iha api aṅitam ṅit iti āha . \\
\noindent \textbf{(P\_1,2.1.1)} ṅidvat iti gamyate . \\
\noindent \textbf{(P\_1,2.1.1)} akitam kit iti āha . \\
\noindent \textbf{(P\_1,2.1.1)} kidvat iti gamyate . \\
\noindent \textbf{(P\_1,2.1.1)} tadvadatideśe akidvidhiprasaṅgaḥ . tadvadatideśe akidvidhiḥ api prāpnoti . \\
\noindent \textbf{(P\_1,2.1.1)} sṛjidṛśoḥ jhali am akiti : sisṛkṣati didṛkṣate : akillakṣaṇaḥ amāgamaḥ prāpnoti . \\
\noindent \textbf{(P\_1,2.1.1)} siddham tu prasajyapratiṣedhāt . \\
\noindent \textbf{(P\_1,2.1.1)} siddham etat . \\
\noindent \textbf{(P\_1,2.1.1)} katham . \\
\noindent \textbf{(P\_1,2.1.1)} prasajya ayam pratiṣedhaḥ kriyate : kiti na iti . \\
\noindent \textbf{(P\_1,2.1.1)} sarvatra sanantāt ātmanepadapratiṣedhaḥ . \\
\noindent \textbf{(P\_1,2.1.1)} sarveṣu pakṣeṣu sanantāt ātmanepadam prāpnoti . \\
\noindent \textbf{(P\_1,2.1.1)} uccukuṭiṣati nicukuṭiṣati : ṅiti iti ātmanepadam prāpnoti . \\
\noindent \textbf{(P\_1,2.1.1)} tasya pratiṣedhaḥ vaktavyaḥ . \\
\noindent \textbf{(P\_1,2.1.1)} siddham tu pūrvasya kāryātideśāt . \\
\noindent \textbf{(P\_1,2.1.1)} siddham etat . \\
\noindent \textbf{(P\_1,2.1.1)} katham . \\
\noindent \textbf{(P\_1,2.1.1)} pūrvasya yat kāryam tat atidiśyate . \\
\noindent \textbf{(P\_1,2.1.1)} kim vaktavyam etat . \\
\noindent \textbf{(P\_1,2.1.1)} na hi . \\
\noindent \textbf{(P\_1,2.1.1)} katham anucyamānam gaṃsyate . \\
\noindent \textbf{(P\_1,2.1.1)} saptamyarthe api vatiḥ bhavati . \\
\noindent \textbf{(P\_1,2.1.1)} tat yathā : mathurāyām iva mathurāvat pāṭaliputre iva pāṭaliputravat evam ṅiti iva ṅidvat . \\
\noindent \textbf{(P\_1,2.1.2)} atha kimartham pṛthak ṅitkitau kriyete na sarvam kit eva vā syāt ṅit eva vā . \\
\noindent \textbf{(P\_1,2.1.2)} pṛthaganubandhatve prayojanam vacisvapiyajādīnām asamprasāraṇam sārvadhātukacaṅādiṣu . \\
\noindent \textbf{(P\_1,2.1.2)} pṛthaganubandhatve prayojanam vacisvapiyajādīnām asamprasāraṇam sārvadhātuke caṅādiṣu ca. sārvadhātuke prayojanam : yathā iha bhavati suptaḥ , suptavān iti evam svapitaḥ , svapithaḥ : atra api prāpnoti . \\
\noindent \textbf{(P\_1,2.1.2)} caṅādiṣu prayojanam . \\
\noindent \textbf{(P\_1,2.1.2)} ke punaḥ caṅādayaḥ . \\
\noindent \textbf{(P\_1,2.1.2)} caṅaṅnajiṅṅvanibathaṅnaṅaḥ . \\
\noindent \textbf{(P\_1,2.1.2)} caṅ : yathā iha bhavati śūnaḥ , śūnavān iti evam aśiśviyat : atra api prāpnoti . \\
\noindent \textbf{(P\_1,2.1.2)} aṅ : yathā iha bhavati śūnaḥ , uktaḥ iti evam aśvat , avocat : atra api prāpnoti . \\
\noindent \textbf{(P\_1,2.1.2)} najiṅ : yathā iha bhavati suptaḥ iti evam svapnak : atra api prāpnoti . \\
\noindent \textbf{(P\_1,2.1.2)} ṅvanip : yathā iha bhavati iṣṭaḥ iti evam yajvā : atra api prāpnoti . \\
\noindent \textbf{(P\_1,2.1.2)} athaṅ : yathā iha bhavati uṣitaḥ iti evam āvasathaḥ : atra api prāpnoti . \\
\noindent \textbf{(P\_1,2.1.2)} naṅ : yathā iha bhavati iṣṭam evam yajñaḥ : atra api prāpnoti . \\
\noindent \textbf{(P\_1,2.1.2)} jāgraḥ aguṇavidhiḥ . \\
\noindent \textbf{(P\_1,2.1.2)} jāgarteḥ aguṇavidhiḥ prayojanam . \\
\noindent \textbf{(P\_1,2.1.2)} yathā iha bhavati jāgṛtaḥ , jāgṛthaḥ iti aṅiti iti paryudāsaḥ evam jāgaritaḥ, jāgaritavān iti atra api prāpnoti . \\
\noindent \textbf{(P\_1,2.1.2)} aparaḥ āha : jāgraḥ guṇavidhiḥ . \\
\noindent \textbf{(P\_1,2.1.2)} jāgarteḥ guṇavidhiḥ prayojanam . \\
\noindent \textbf{(P\_1,2.1.2)} yathā iha bhavati jāgaritaḥ , jāgaritavān evam jāgṛtaḥ jāgṛthaḥ iti atra api prāpnoti . \\
\noindent \textbf{(P\_1,2.1.2)} kuṭādīnām iṭpratiṣedhaḥ . \\
\noindent \textbf{(P\_1,2.1.2)} kuṭādīnām iṭpratiṣedhaḥ prayojanam . \\
\noindent \textbf{(P\_1,2.1.2)} yathā iha bhavati lūtvā pūtvā śryukaḥ kiti iti iṭpratiṣedhaḥ evam nuvitā dhuvitā : atra api prāpnoti . \\
\noindent \textbf{(P\_1,2.1.2)} ktvāyām kitpratiṣedhaḥ ca . \\
\noindent \textbf{(P\_1,2.1.2)} ktvāyām kitpratiṣedhaḥ ca prayojanam . \\
\noindent \textbf{(P\_1,2.1.2)} kim ca . \\
\noindent \textbf{(P\_1,2.1.2)} iṭpratiṣedhaḥ ca . \\
\noindent \textbf{(P\_1,2.1.2)} na iti āha . \\
\noindent \textbf{(P\_1,2.1.2)} adeśe ayam caḥ paṭhitaḥ . \\
\noindent \textbf{(P\_1,2.1.2)} ktvāyām ca kitpratiṣedhaḥ iti . \\
\noindent \textbf{(P\_1,2.1.2)} yathā iha bhavati devitvā sevitvā na ktvā seṭ iti pratiṣedhaḥ evam kuṭitvā puṭitvā : atra api prāpnoti . \\
\noindent \textbf{(P\_1,2.1.2)} atha vā deśe eva ayam caḥ paṭhitaḥ . \\
\noindent \textbf{(P\_1,2.1.2)} ktvāyām kitpratiṣedhaḥ ca iṭpratiṣedhaḥ ca . \\
\noindent \textbf{(P\_1,2.1.2)} iṭpratiṣedhaḥ : yathā iha bhavati lūtvā pūtvā śryukaḥ kiti iti iṭpratiṣedhaḥ evam nuvitvā dhuvitvā : atra api prāpnoti . \\
\noindent \textbf{(P\_1,2.1.2)} syāt etat prayojanam yadi atra niyogataḥ ātideśikena ṅittvena aupadeśikam kittvam bādhyeta . \\
\noindent \textbf{(P\_1,2.1.2)} sati api tu ṅittve kit eva eṣaḥ . \\
\noindent \textbf{(P\_1,2.1.2)} tasmāt nūtvā dhūtvā iti eva bhavitavyam . \\
\noindent \textbf{(P\_1,2.4.1)} sārvadhātukagrahaṇam kimartham . \\
\noindent \textbf{(P\_1,2.4.1)} apit iti iyati ucyamāne ārdhadhātukasya api apitaḥ ṅittvam prasajyeta : kartā hartā . \\
\noindent \textbf{(P\_1,2.4.1)} na eṣaḥ doṣaḥ . \\
\noindent \textbf{(P\_1,2.4.1)} ācāryapravṛttiḥ jñāpayati na anena ārdhadhātukasya ṅittvam bhavati iti yat ayam ārdhadhātukīyān kān cit ṅitaḥ karoti : caṅaṅnajiṅṅvanibathaṅnaṅaḥ . \\
\noindent \textbf{(P\_1,2.4.1)} sārvadhātuke api etat jñāpakam syāt . \\
\noindent \textbf{(P\_1,2.4.1)} na iti āha . \\
\noindent \textbf{(P\_1,2.4.1)} tulyajātīyasya jñāpakam . \\
\noindent \textbf{(P\_1,2.4.1)} kaḥ ca tulyajātīyaḥ . \\
\noindent \textbf{(P\_1,2.4.1)} yathājātīyakāḥ caṅaṅnajiṅṅvanibathaṅnaṅaḥ . \\
\noindent \textbf{(P\_1,2.4.1)} kathañjātīyakāḥ ca ete . \\
\noindent \textbf{(P\_1,2.4.1)} ārdhadhātukāḥ . \\
\noindent \textbf{(P\_1,2.4.1)} yadi etat asti tulyajātīyasya jñāpakam iti caṅaṅau luṅvikaraṇānam jñāpakau syātām najiṅ vartamānakālānām ṅvanip bhūtakālānām athaṅśabdaḥ auṇādikānām naṅśabdaḥ ghañarthānām . \\
\noindent \textbf{(P\_1,2.4.1)} tasmāt sārvadhātukagrahaṇam kartavyam . \\
\noindent \textbf{(P\_1,2.4.2)} kim punaḥ ayam paryudāsaḥ : yat anyat pitaḥ , āhosvit prasajya ayam pratiṣedhaḥ : pit na iti . \\
\noindent \textbf{(P\_1,2.4.2)} kaḥ ca atra viśeṣaḥ . \\
\noindent \textbf{(P\_1,2.4.2)} apit ṅit iti cet śabdekādeśapratiṣedhaḥ ādivattvāt . \\
\noindent \textbf{(P\_1,2.4.2)} apit ṅit iti cet śabdekādeśe pratiṣedhaḥ vaktavyaḥ : cyavante plavante . \\
\noindent \textbf{(P\_1,2.4.2)} kim kāraṇam . \\
\noindent \textbf{(P\_1,2.4.2)} ādivattvāt . \\
\noindent \textbf{(P\_1,2.4.2)} pidapitoḥ ekādeśaḥ apitaḥ ādivat syāt . \\
\noindent \textbf{(P\_1,2.4.2)} asti anyat pitaḥ iti kṛtvā ṅittvam prāpnoti . \\
\noindent \textbf{(P\_1,2.4.2)} astu tarhi prasajyapratiṣedhaḥ : pit na iti . \\
\noindent \textbf{(P\_1,2.4.2)} na pit ṅit iti cet uttamaikādeśapratiṣedhaḥ . \\
\noindent \textbf{(P\_1,2.4.2)} pit na iti cet uttamaikādeśe pratiṣedhaḥ prāpnoti : tudāni likhāni . \\
\noindent \textbf{(P\_1,2.4.2)} kim kāraṇam . \\
\noindent \textbf{(P\_1,2.4.2)} ādivattvāt eva . \\
\noindent \textbf{(P\_1,2.4.2)} pidapitoḥ ekādeśaḥ pitaḥ ādivat syāt . \\
\noindent \textbf{(P\_1,2.4.2)} tatra pit na iti pratiṣedhaḥ prāpnoti . \\
\noindent \textbf{(P\_1,2.4.2)} yathā icchasi tathā astu . \\
\noindent \textbf{(P\_1,2.4.2)} nanu ca uktam ubhayathā api doṣaḥ iti . \\
\noindent \textbf{(P\_1,2.4.2)} ubhayathā api na doṣaḥ . \\
\noindent \textbf{(P\_1,2.4.2)} ekādeśaḥ pūrvavidhau sthānivat iti sthānivadbhāvāt vyavadhānam . \\
\noindent \textbf{(P\_1,2.5)} ṛdupadhebhyaḥ liṭaḥ kittvam guṇāt vipratiṣedhena . \\
\noindent \textbf{(P\_1,2.5)} ṛdupadhebhyaḥ liṭaḥ kittvam guṇāt bhavati vipratiṣedhena : vavṛte vavṛdhe . \\
\noindent \textbf{(P\_1,2.5)} uktam vā . \\
\noindent \textbf{(P\_1,2.5)} kim uktam . \\
\noindent \textbf{(P\_1,2.5)} na vā ksasya anavakāśatvāt apavādaḥ guṇasya iti . \\
\noindent \textbf{(P\_1,2.5)} viṣamaḥ upanyāsaḥ . \\
\noindent \textbf{(P\_1,2.5)} yuktam tatra yat anavakāśam kitkaraṇam guṇam bādhate . \\
\noindent \textbf{(P\_1,2.5)} iha punaḥ ubhayam sāvakāśam . \\
\noindent \textbf{(P\_1,2.5)} kitkaraṇasya avakāśaḥ: ījatuḥ ījuḥ . \\
\noindent \textbf{(P\_1,2.5)} guṇasya avakāśaḥ : vartitvā vardhitvā . \\
\noindent \textbf{(P\_1,2.5)} iha ubhayam prāpnoti : vavṛte vavṛdhe . \\
\noindent \textbf{(P\_1,2.5)} paratvāt guṇaḥ prāpnoti . \\
\noindent \textbf{(P\_1,2.5)} idam tarhi uktam : iṣṭavācī paraśabdaḥ . \\
\noindent \textbf{(P\_1,2.5)} vipratiṣedhe param yat iṣṭam tat bhavati itI . \\
\noindent \textbf{(P\_1,2.6)} kimartham idam ucyate . \\
\noindent \textbf{(P\_1,2.6)} indheḥ samyogārtham vacanam bhavateḥ pidartham . \\
\noindent \textbf{(P\_1,2.6)} ayam yogaḥ śakyaḥ avaktum . \\
\noindent \textbf{(P\_1,2.6)} katham . \\
\noindent \textbf{(P\_1,2.6)} indheḥ chandoviṣayatvāt bhuvaḥ vukaḥ nityatvāt tābhyām kidvacanānarthakyam . \\
\noindent \textbf{(P\_1,2.6)} indheḥ chandoviṣayaḥ liṭ . \\
\noindent \textbf{(P\_1,2.6)} na hi antareṇa chandaḥ indheḥ anantaraḥ liṭ labhyaḥ . \\
\noindent \textbf{(P\_1,2.6)} āmā bhāṣāyām bhavitavyam . \\
\noindent \textbf{(P\_1,2.6)} bhuvaḥ vukaḥ nityatvāt . \\
\noindent \textbf{(P\_1,2.6)} bhavateḥ api nityaḥ vuk . \\
\noindent \textbf{(P\_1,2.6)} kṛte api prāpnoti akṛte api . \\
\noindent \textbf{(P\_1,2.6)} tābhyām kidvacanānarthakyam . \\
\noindent \textbf{(P\_1,2.6)} tābhyām indhibhavitibhyām kidvacanam anarthakam . \\
\noindent \textbf{(P\_1,2.7)} kimartham mṛḍādibhyaḥ parasya ktvaḥ kittvam ucyate . \\
\noindent \textbf{(P\_1,2.7)} kit eva hi ktvā . \\
\noindent \textbf{(P\_1,2.7)} na ktvā seṭ iti pratiṣedhaḥ prāpnoti tadbādhanārtham . \\
\noindent \textbf{(P\_1,2.7)} yadi tarhi mṛḍādibhyaḥ parasya ktvaḥ kittvam ucyate na arthaḥ na ktvā seṭ iti anena kittvapratiṣedhena . \\
\noindent \textbf{(P\_1,2.7)} idam niyamārtham bhaviṣyati : mṛḍādibhyaḥ eva parasya ktvaḥ kittvam bhavati na anyebhyaḥ iti . \\
\noindent \textbf{(P\_1,2.7)} yadi niyamaḥ kriyate iha api tarhi niyamān na prāpnoti : lūtvā pūtvā . \\
\noindent \textbf{(P\_1,2.7)} atra api akittvam prāpnoti . \\
\noindent \textbf{(P\_1,2.7)} tulyajātīyasya niyamaḥ . \\
\noindent \textbf{(P\_1,2.7)} kaḥ ca tulyajātīyaḥ . \\
\noindent \textbf{(P\_1,2.7)} yathājātīyakaḥ mṛḍādibhyaḥ paraḥ ktvā . \\
\noindent \textbf{(P\_1,2.7)} kathañjātīyakaḥ mṛḍādibhyaḥ paraḥ ktvā . \\
\noindent \textbf{(P\_1,2.7)} seṭ . \\
\noindent \textbf{(P\_1,2.7)} evam api asti atra kaḥ cit vibhāṣiteṭ . \\
\noindent \textbf{(P\_1,2.7)} saḥ aniṭām niyāmakaḥ syāt . \\
\noindent \textbf{(P\_1,2.7)} astu tāvat ye seṭaḥ teṣām grahaṇam niyamāṛtham . \\
\noindent \textbf{(P\_1,2.7)} yaḥ idānīm vibhāṣiteṭ tasya grahaṇam vidhyartham bhaviṣyati . \\
\noindent \textbf{(P\_1,2.8)} svapipracchyoḥ sanartham grahaṇam kit eva hi ktvā . \\
\noindent \textbf{(P\_1,2.9)} kimartham ikaḥ parasya sanaḥ kittvam ucyate . \\
\noindent \textbf{(P\_1,2.9)} ikaḥ kittvam guṇaḥ mā bhūt . \\
\noindent \textbf{(P\_1,2.9)} ikaḥ kittvam ucyate guṇaḥ mā bhūt iti : cicīṣati tuṣṭūṣati . \\
\noindent \textbf{(P\_1,2.9)} na etat asti prayojanam . \\
\noindent \textbf{(P\_1,2.9)} dīrghārambhāt . \\
\noindent \textbf{(P\_1,2.9)} dīrghatvam atra bādhakam bhaviṣyati . \\
\noindent \textbf{(P\_1,2.9)} kṛte bhavet . \\
\noindent \textbf{(P\_1,2.9)} kṛte khalu dīrghatve guṇaḥ prāpnoti . \\
\noindent \textbf{(P\_1,2.9)} anarthakam tu . \\
\noindent \textbf{(P\_1,2.9)} anarthakam evam sati dīrghatvam syāt . \\
\noindent \textbf{(P\_1,2.9)} na anarthakam . \\
\noindent \textbf{(P\_1,2.9)} hrasvārtham . \\
\noindent \textbf{(P\_1,2.9)} hrasvānām dīrghavacanasāmarthyāt guṇaḥ na bhaviṣyati . \\
\noindent \textbf{(P\_1,2.9)} bhavet hrasvānām dīrghavacanasāmarthyāt guṇaḥ na syāt . \\
\noindent \textbf{(P\_1,2.9)} dīrghāṇām tu prasajyate . \\
\noindent \textbf{(P\_1,2.9)} dīrghāṇām tu khalu guṇaḥ prāpnoti . \\
\noindent \textbf{(P\_1,2.9)} nanu ca dīrghāṇām api dīrghavacanasāmarthyāt guṇaḥ na bhaviṣyati . \\
\noindent \textbf{(P\_1,2.9)} na dīrghāṇām dīrghāḥ prāpnuvanti . \\
\noindent \textbf{(P\_1,2.9)} kim kāraṇam . \\
\noindent \textbf{(P\_1,2.9)} na hi bhuktavān punaḥ bhuṅkte na ca kṛtaśmaśruḥ punaḥ śmaśrūni kārayati . \\
\noindent \textbf{(P\_1,2.9)} nanu ca punaḥ pravṛttiḥ api dṛṣṭā . \\
\noindent \textbf{(P\_1,2.9)} bhuktavān ca punaḥ bhuṅkte kṛtaśmaśruḥ ca punaḥ śmaśrūni kārayati . \\
\noindent \textbf{(P\_1,2.9)} sāmarthyāt hi punaḥ bhāvyam . \\
\noindent \textbf{(P\_1,2.9)} sāmarthyāt tatra punaḥ pravṛttiḥ bhavati bhojanaviśeṣāt śilpiviśeṣāt vā . \\
\noindent \textbf{(P\_1,2.9)} dīrghāṇām punaḥ dīrghatvavacane na kim cit prayojanam asti . \\
\noindent \textbf{(P\_1,2.9)} akṛtakāri khalu api śāstram agnivat . \\
\noindent \textbf{(P\_1,2.9)} tat yathā agniḥ yad adagdham tat dahati . \\
\noindent \textbf{(P\_1,2.9)} dīrghāṇām api dīrghavacane etat prayojanam guṇaḥ mā bhūt iti . \\
\noindent \textbf{(P\_1,2.9)} kṛtakāri khalu api śāstram parjanyavat . \\
\noindent \textbf{(P\_1,2.9)} tat yathā parjanyaḥ yāvat ūnam pūrṇam ca sarvam abhivarṣati . \\
\noindent \textbf{(P\_1,2.9)} yathā eva tarhi dīrghavacanasāmarthyāt guṇaḥ na bhavati evam ṝdittvam api na prāpnoti : cikīrṣati jihīrṣati iti. ṝdittvam dīrghasaṃśrayam . \\
\noindent \textbf{(P\_1,2.9)} na akṛte dīrghe ṝdittvam prāpnoti . \\
\noindent \textbf{(P\_1,2.9)} kim kāraṇam . \\
\noindent \textbf{(P\_1,2.9)} ṝtaḥ iti ucyate . \\
\noindent \textbf{(P\_1,2.9)} bhavet hrasvānām na akṛte dīrghe ṝdittvam syāt dīrghāṇām tu khalu akṛte api dīrghatve ṝdittvam prāpnoti . \\
\noindent \textbf{(P\_1,2.9)} dīrghāṇām na akṛte dīrghe . \\
\noindent \textbf{(P\_1,2.9)} dīrghāṇām api na akṛte dīrghe ṝdittvam prāpnoti . \\
\noindent \textbf{(P\_1,2.9)} yadā dīrghatvena guṇaḥ bādhitaḥ tataḥ uttarakālam ṝdittvam bhavati . \\
\noindent \textbf{(P\_1,2.9)} ṇilopaḥ tu prayojanam . \\
\noindent \textbf{(P\_1,2.9)} idam tarhi prayojanam : ṇilopaḥ yathā syāt iti : jñīpsati . \\
\noindent \textbf{(P\_1,2.9)} kva astāḥ kva nipatitāḥ kva kittvam kva ṇilopaḥ . \\
\noindent \textbf{(P\_1,2.9)} kaḥ vā abhisambandhaḥ yat sati kittve ṇilopaḥ syāt asati na syāt . \\
\noindent \textbf{(P\_1,2.9)} eṣaḥ abhisambandhaḥ yat sati kittve sāvakāśam dīrghatvam paratvāt ṇilopaḥ bādhate . \\
\noindent \textbf{(P\_1,2.9)} asati punaḥ kittve anavakāśam dīrghatvam yathā eva guṇam bādhate evam ṇilopam api bādheta . \\
\noindent \textbf{(P\_1,2.9)} tatra ṇilopasya avakāśaḥ : kāraṇā hāraṇā . \\
\noindent \textbf{(P\_1,2.9)} dīrghatvasya avakāśaḥ : cicīṣati tuṣṭūṣati . \\
\noindent \textbf{(P\_1,2.9)} iha ubhayam prāpnoti : jñīpsati . \\
\noindent \textbf{(P\_1,2.9)} paratvāt ṇilopaḥ . \\
\noindent \textbf{(P\_1,2.9)} asati api kittve sāvakāśam dīrghatvam . \\
\noindent \textbf{(P\_1,2.9)} kaḥ avakāśaḥ . \\
\noindent \textbf{(P\_1,2.9)} isbhāvaḥ . \\
\noindent \textbf{(P\_1,2.9)} nimitsati pramitsati . \\
\noindent \textbf{(P\_1,2.9)} mīnātiminotyoḥ dīrghatve kṛte mīgrahaṇena grahaṇam yathā syāt . \\
\noindent \textbf{(P\_1,2.9)} yathā eva tarhi asati kittve sāvakāśam dīrghatvam paratvāt ṇilopaḥ bādhate evam guṇaḥ api bādheta . \\
\noindent \textbf{(P\_1,2.9)} tasmāt kittvam vaktavyam . \\
\noindent \textbf{(P\_1,2.9)} ikaḥ kittvam guṇaḥ mā bhūt dīrghārambhāt kṛte bhavet | anarthakam tu hrasvārtham dīrghāṇām tu prasajyate || sāmarthyāt hi punaḥ bhāvyam ṝdittvam dīrghasaṃśrayam dīrghāṇām na akṛte dīrghe ṇilopaḥ tu prayojanam . \\
\noindent \textbf{(P\_1,2.10)} ayuktaḥ ayam nirdeśaḥ . \\
\noindent \textbf{(P\_1,2.10)} katham hi ikaḥ nāma hal antaḥ syāt anyasya anyaḥ . \\
\noindent \textbf{(P\_1,2.10)} katham tarhi nirdeśaḥ kartavyaḥ . \\
\noindent \textbf{(P\_1,2.10)} igvataḥ halaḥ iti . \\
\noindent \textbf{(P\_1,2.10)} yadi evam yiyakṣati atra api prāpnoti . \\
\noindent \textbf{(P\_1,2.10)} evam tarhi igupadhāt halantāt iti vakṣyāmi . \\
\noindent \textbf{(P\_1,2.10)} evam api dambheḥ na prāpnoti . \\
\noindent \textbf{(P\_1,2.10)} sūtram ca bhidyate . \\
\noindent \textbf{(P\_1,2.10)} yathānyāsam eva astu . \\
\noindent \textbf{(P\_1,2.10)} nanu ca uktam ayuktaḥ ayam nirdeśaḥ iti . \\
\noindent \textbf{(P\_1,2.10)} na ayuktaḥ . \\
\noindent \textbf{(P\_1,2.10)} antaśabdaḥ ayam asti eva avayavavācī . \\
\noindent \textbf{(P\_1,2.10)} tat yathā vastrāntaḥ , vasanāntaḥ : vastrāvayavaḥ , vasanāvayavaḥ iti gamyate . \\
\noindent \textbf{(P\_1,2.10)} asti sāmīpye vartate . \\
\noindent \textbf{(P\_1,2.10)} tat yathā udakāntam gataḥ iti . \\
\noindent \textbf{(P\_1,2.10)} udakasamīpam gataḥ iti gamyate . \\
\noindent \textbf{(P\_1,2.10)} tat yaḥ sāmīpye vartate tasya idam grahaṇam . \\
\noindent \textbf{(P\_1,2.10)} evam api dambheḥ na sidhyati . \\
\noindent \textbf{(P\_1,2.10)} yaḥ atra iksamīpe hal na tasmāt uttaraḥ san . \\
\noindent \textbf{(P\_1,2.10)} yasmāt uttaraḥ san na asau iksamīpe hal . \\
\noindent \textbf{(P\_1,2.10)} evam tarhi dambheḥ halgrahaṇasya jātivācakatvāt siddham . \\
\noindent \textbf{(P\_1,2.10)} haljātiḥ nirdiśyate : ikaḥ uttarā yā haljātiḥ iti . \\
\noindent \textbf{(P\_1,2.11)} katham idam vijñāyate : ātmanepadam yau liṅsicau iti āhosvit ātmanepadeṣu parataḥ yau liṅsicau iti . \\
\noindent \textbf{(P\_1,2.11)} kim ca ataḥ. yadi vijñāyate ātmanepadam yau liṅsicau iti liṅ viśeṣitaḥ sic aviśeṣitaḥ . \\
\noindent \textbf{(P\_1,2.11)} atha vijñāyate ātmanepadeṣu parataḥ yau liṅsicau iti sic viśeṣitaḥ liṅ aviśeṣitaḥ . \\
\noindent \textbf{(P\_1,2.11)} yathā icchasi tathā astu . \\
\noindent \textbf{(P\_1,2.11)} astu tāvat ātmanepadam yau liṅsicau iti . \\
\noindent \textbf{(P\_1,2.11)} nanu ca uktam . \\
\noindent \textbf{(P\_1,2.11)} liṅ viśeṣitaḥ sic aviśeṣitaḥ iti . \\
\noindent \textbf{(P\_1,2.11)} sic ca viśeṣitaḥ . \\
\noindent \textbf{(P\_1,2.11)} katham . \\
\noindent \textbf{(P\_1,2.11)} ātmanepadam sic na asti iti kṛtvā ātmanepadapare sici kāryam vijñāsyate . \\
\noindent \textbf{(P\_1,2.11)} atha vā punaḥ astu ātmanepadeṣu parataḥ yau liṅsicau iti . \\
\noindent \textbf{(P\_1,2.11)} nanu ca uktam sic viśeṣitaḥ liṅ aviśeṣitaḥ iti . \\
\noindent \textbf{(P\_1,2.11)} liṅ ca viśeṣitaḥ . \\
\noindent \textbf{(P\_1,2.11)} katham . \\
\noindent \textbf{(P\_1,2.11)} ātmanepadeṣu parataḥ liṅ na asti iti kṛtvā ātmanepade liṅi kāryam vijñāsyate . \\
\noindent \textbf{(P\_1,2.11)} na eva vā punaḥ arthaḥ liṅviśeṣaṇena ātmanepadagrahaṇena . \\
\noindent \textbf{(P\_1,2.11)} kim kāraṇam . \\
\noindent \textbf{(P\_1,2.11)} jhal iti vartate . \\
\noindent \textbf{(P\_1,2.11)} ātmanepadeṣu ca eva liṅ jhalādiḥ na parasmaipadeṣu . \\
\noindent \textbf{(P\_1,2.11)} tat etat sijviśeṣaṇam ātmanepadagrahaṇam . \\
\noindent \textbf{(P\_1,2.11)} atha sijviśeṣaṇe ātmanepadagrahaṇe sati kim prayojanam . \\
\noindent \textbf{(P\_1,2.11)} iha mā bhūt : ayākṣīt , avātsīt . \\
\noindent \textbf{(P\_1,2.11)} na etat asti . \\
\noindent \textbf{(P\_1,2.11)} ikaḥ iti vartate . \\
\noindent \textbf{(P\_1,2.11)} evam api anaiṣīt , acaiṣīt : atra api prāpnoti . \\
\noindent \textbf{(P\_1,2.11)} etat api na asti prayojanam . \\
\noindent \textbf{(P\_1,2.11)} halantāt iti vartate . \\
\noindent \textbf{(P\_1,2.11)} evam api akoṣīt , amoṣīt : atra api prāpnoti . \\
\noindent \textbf{(P\_1,2.11)} na etat asti prayojanam . \\
\noindent \textbf{(P\_1,2.11)} jhal iti vartate . \\
\noindent \textbf{(P\_1,2.11)} evam api abhaitsīt , acchaitsīt : atra api prāpnoti . \\
\noindent \textbf{(P\_1,2.11)} na etat asti . \\
\noindent \textbf{(P\_1,2.11)} iglakṣaṇayoḥ guṇavṛddhyoḥ pratiṣedhaḥ na ca eṣā iglakṣaṇā vṛddhiḥ . \\
\noindent \textbf{(P\_1,2.11)} idam tarhi prayojanam : iha mā bhūt : adrākṣīt , asrākṣīt . \\
\noindent \textbf{(P\_1,2.11)} kim ca syāt . \\
\noindent \textbf{(P\_1,2.11)} akillakṣaṇaḥ amāgamaḥ na syāt . \\
\noindent \textbf{(P\_1,2.17)} it ca kasya takārettvam . \\
\noindent \textbf{(P\_1,2.17)} kasya hetoḥ ikāraḥ taparaḥ kriyate . \\
\noindent \textbf{(P\_1,2.17)} dīrghaḥ mā bhūt . \\
\noindent \textbf{(P\_1,2.17)} dīrghaḥ mā bhūt iti . \\
\noindent \textbf{(P\_1,2.17)} ṛte api saḥ . \\
\noindent \textbf{(P\_1,2.17)} antareṇa api ārambham siddhaḥ atra dīrghaḥ : ghumāsthāgāpājahāti iti . \\
\noindent \textbf{(P\_1,2.17)} anantare plutaḥ mā bhūt . \\
\noindent \textbf{(P\_1,2.17)} idam tarhi prayojanam : anantare plutaḥ mā bhūt iti . \\
\noindent \textbf{(P\_1,2.17)} kutaḥ nu khalu etat anantarārthe ārambhe hrasvaḥ bhaviṣyati na punaḥ plutaḥ iti . \\
\noindent \textbf{(P\_1,2.17)} plutaḥ ca viṣaye smṛtaḥ . \\
\noindent \textbf{(P\_1,2.17)} viṣaye khalu plutaḥ ucyate . \\
\noindent \textbf{(P\_1,2.17)} yadā ca saḥ viṣayaḥ bhavitavyam eva tadā plutena . \\
\noindent \textbf{(P\_1,2.17)} it ca kasya takārettvam dīrghaḥ mā bhūt ṛte api saḥ | anantare plutaḥ mā bhūt plutaḥ ca viṣaye smṛtaḥ \\
\noindent \textbf{(P\_1,2.18)} na seṭ iti kṛte akittve . \\
\noindent \textbf{(P\_1,2.18)} na seṭ iti eva siddham . \\
\noindent \textbf{(P\_1,2.18)} na arthaḥ ktvāgrahaṇena . \\
\noindent \textbf{(P\_1,2.18)} niṣṭhāyam api tarhi prāpnoti : gudhitaḥ gudhitavān iti . \\
\noindent \textbf{(P\_1,2.18)} niṣṭhāyām avadhāraṇāt . \\
\noindent \textbf{(P\_1,2.18)} niṣṭhāyām avadhāraṇāt na bhaviṣyati . \\
\noindent \textbf{(P\_1,2.18)} kim avadhāraṇam . \\
\noindent \textbf{(P\_1,2.18)} niṣṭhā śīṅsvidimidikṣvididhṛṣaḥ iti . \\
\noindent \textbf{(P\_1,2.18)} parokṣāyām tarhi prāpnoti . \\
\noindent \textbf{(P\_1,2.18)} kim ca syāt . \\
\noindent \textbf{(P\_1,2.18)} papiva papima : kṅiti iti ākāralopaḥ na syāt . \\
\noindent \textbf{(P\_1,2.18)} mā bhūt evam . \\
\noindent \textbf{(P\_1,2.18)} iṭi iti evam bhaviṣyati . \\
\noindent \textbf{(P\_1,2.18)} idam tarhi : jagmiva jaghniva . \\
\noindent \textbf{(P\_1,2.18)} kṅiti iti upadhālopaḥ na syāt . \\
\noindent \textbf{(P\_1,2.18)} jñāpakāt na parokṣāyām . \\
\noindent \textbf{(P\_1,2.18)} jñāpakāt parokṣāyām na bhaviṣyati . \\
\noindent \textbf{(P\_1,2.18)} kim jñāpakam . \\
\noindent \textbf{(P\_1,2.18)} sani jhalgrahaṇam viduḥ . \\
\noindent \textbf{(P\_1,2.18)} yat ayam ikaḥ jhal iti jhalgrahaṇam karoti tat jñāpayati ācāryaḥ aupadeśikasya kittvasya pratiṣedhaḥ na ātideśikasya iti . \\
\noindent \textbf{(P\_1,2.18)} katham kṛtvā jñāpakam . \\
\noindent \textbf{(P\_1,2.18)} jhalgrahaṇasya etat prayojanam : iha mā bhūt : śiśayiṣate iti . \\
\noindent \textbf{(P\_1,2.18)} yadi ca atra ātideśikasya kittvasya pratiṣedhaḥ syāt jhalgrahaṇam anarthakam syāt . \\
\noindent \textbf{(P\_1,2.18)} astu atra kittvam . \\
\noindent \textbf{(P\_1,2.18)} na seṭ iti pratiṣedhaḥ bhaviṣyati . \\
\noindent \textbf{(P\_1,2.18)} paśyati tu ācāryaḥ aupadeśikasya kittvasya pratiṣedhaḥ na ātideśikasya iti . \\
\noindent \textbf{(P\_1,2.18)} tataḥ jhalgrahaṇam karoti . \\
\noindent \textbf{(P\_1,2.18)} na etat asti prayojanam . \\
\noindent \textbf{(P\_1,2.18)} uttarārtham etat syāt . \\
\noindent \textbf{(P\_1,2.18)} sthāghvoḥ it ca jhalādau yathā syāt . \\
\noindent \textbf{(P\_1,2.18)} iha mā bhūt : upāsthāyiṣātām upāsthāyiṣata . \\
\noindent \textbf{(P\_1,2.18)} ittvam kitsanniyogena . \\
\noindent \textbf{(P\_1,2.18)} kitsanniyogena ittvam ucyate . \\
\noindent \textbf{(P\_1,2.18)} tena asati kittve ittvam na bhaviṣyati . \\
\noindent \textbf{(P\_1,2.18)} reṇa tulyam sudhīvani . \\
\noindent \textbf{(P\_1,2.18)} tat yathā sudhīvā supīvā : ṅīpsanniyogena raḥ ucyamānaḥ asati ṅīpi na bhavati . \\
\noindent \textbf{(P\_1,2.18)} atha vā astu atra ittvam . \\
\noindent \textbf{(P\_1,2.18)} kā rūpasiddhiḥ . \\
\noindent \textbf{(P\_1,2.18)} vṛddhau kṛtāyām āyādeśaḥ bhaviṣyati . \\
\noindent \textbf{(P\_1,2.18)} vasvartham . \\
\noindent \textbf{(P\_1,2.18)} vasvartham tarhi ktvāgrahaṇam kartavyam . \\
\noindent \textbf{(P\_1,2.18)} vasau hi aupadeśikam kittvam . \\
\noindent \textbf{(P\_1,2.18)} kim ca syāt . \\
\noindent \textbf{(P\_1,2.18)} papivān papimān : kṅiti iti ākāralopaḥ na syāt . \\
\noindent \textbf{(P\_1,2.18)} mā bhūt evam . \\
\noindent \textbf{(P\_1,2.18)} iṭi iti evam bhaviṣyati . \\
\noindent \textbf{(P\_1,2.18)} idam tarhi : jagmivān , jaghnivān : kṅiti iti upadhālopaḥ na syāt . \\
\noindent \textbf{(P\_1,2.18)} kidatīdeśāt . \\
\noindent \textbf{(P\_1,2.18)} astu atra aupadeśikasya kittvasya pratiṣedhaḥ . \\
\noindent \textbf{(P\_1,2.18)} ātideśikasya kittvam bhaviṣyati . \\
\noindent \textbf{(P\_1,2.18)} yatra tarhi tat pratiṣidhyate : añjeḥ ājivān iti . \\
\noindent \textbf{(P\_1,2.18)} evam tarhi chāndasḥ kvasuḥ liṭ ca chandasi sārvadhātukam api bhavati . \\
\noindent \textbf{(P\_1,2.18)} tatra sārvadhātukam apit ṅit bhavati iti ṅiti upadhālopaḥ bhaviṣyati . \\
\noindent \textbf{(P\_1,2.18)} nigṛhītiḥ . \\
\noindent \textbf{(P\_1,2.18)} nigṛhītiḥ prayojanam . \\
\noindent \textbf{(P\_1,2.18)} idam tarhi prayojanam : iha mā bhūt : nigṛhītiḥ , upasannihitiḥ , nikucitiḥ . \\
\noindent \textbf{(P\_1,2.18)} tat tarhi ktvāgrahaṇam kartavyam . \\
\noindent \textbf{(P\_1,2.18)} na kartavyam . \\
\noindent \textbf{(P\_1,2.18)} ktvā ca vigrahāt . \\
\noindent \textbf{(P\_1,2.18)} upariṣṭāt yogavibhāgaḥ kariṣyate : na seṭ niṣṭhā śīṅsvidimidikṣvididhṛṣaḥ . \\
\noindent \textbf{(P\_1,2.18)} mṛṣaḥ titikṣāyām . \\
\noindent \textbf{(P\_1,2.18)} udupadhāt bhāvādikarmaṇoḥ anyatarasyām . \\
\noindent \textbf{(P\_1,2.18)} tataḥ pūṅaḥ . \\
\noindent \textbf{(P\_1,2.18)} pūṅaḥ niṣṭhā seṭ na kit bhavati . \\
\noindent \textbf{(P\_1,2.18)} tataḥ ktvā . \\
\noindent \textbf{(P\_1,2.18)} ktvā ca seṭ na kit bhavati . \\
\noindent \textbf{(P\_1,2.18)} pūṅ iti nivṛttam . \\
\noindent \textbf{(P\_1,2.18)} na seṭ iti kṛte akittve niṣṭhāyām avadhāraṇāt | jñāpakāt na parokṣāyām sani jhalgrahaṇam viduḥ || ittvam kitsanniyogena reṇa tulyam sudhīvani | vasvartham kidatīdeśāt nigṛhītiḥ ktvā ca vigrahāt || \\
\noindent \textbf{(P\_1,2.21)} iha kasmāt na bhavati . \\
\noindent \textbf{(P\_1,2.21)} gudhitaḥ gudhitavān iti . \\
\noindent \textbf{(P\_1,2.21)} udupadhāt śapaḥ . \\
\noindent \textbf{(P\_1,2.21)} śabvikaraṇebhyaḥ iṣyate . \\
\noindent \textbf{(P\_1,2.22)} pūṅaḥ ktvāniṣṭhayoḥ iṭi vāprasaṅgaḥ seṭprakaraṇāt . \\
\noindent \textbf{(P\_1,2.22)} pūṅaḥ ktvāniṣṭhayoḥ iṭi vibhāṣā prāpnoti . \\
\noindent \textbf{(P\_1,2.22)} kim kāraṇam . \\
\noindent \textbf{(P\_1,2.22)} seṭprakaraṇāt . \\
\noindent \textbf{(P\_1,2.22)} seṭ iti vartate . \\
\noindent \textbf{(P\_1,2.22)} na vā seṭtvasya akidāśrayatvāt aniṭi vā kittvam . \\
\noindent \textbf{(P\_1,2.22)} na vā eṣaḥ doṣaḥ . \\
\noindent \textbf{(P\_1,2.22)} kim kāraṇam . \\
\noindent \textbf{(P\_1,2.22)} seṭtvasya akidāśrayatvāt . \\
\noindent \textbf{(P\_1,2.22)} akidāśrayam seṭtvam . \\
\noindent \textbf{(P\_1,2.22)} yadā akittvam tadā iṭā bhavitavyam . \\
\noindent \textbf{(P\_1,2.22)} seṭtvasya akidāśrayatvāt aniṭi eva vibhāṣā kittvam bhaviṣyati . \\
\noindent \textbf{(P\_1,2.22)} iḍvidhau pūṅaḥ grahaṇam kriyate . \\
\noindent \textbf{(P\_1,2.22)} tena vacanāt iṭ seṭprakaraṇāt ca iṭi eva vibhāṣā kittvam prāpnoti . \\
\noindent \textbf{(P\_1,2.22)} iḍvidhau hi agrahaṇam . \\
\noindent \textbf{(P\_1,2.22)} iḍvidhau hi pūṅaḥ grahaṇam na kartavyam bhavati . \\
\noindent \textbf{(P\_1,2.22)} bhāradvājīyāḥ paṭhanti : nityam akittvam iḍādyoḥ . ktvāgrahaṇam uttarārtham . nityam akittvam iḍādyoḥ siddham . \\
\noindent \textbf{(P\_1,2.22)} katham. vibhāṣāmadhye ayam yogaḥ kriyate . \\
\noindent \textbf{(P\_1,2.22)} vibhāṣāmadhye ca ye vidhayaḥ te nityāḥ bhavanti . \\
\noindent \textbf{(P\_1,2.22)} kimartham tarhi ktvāgrahaṇam . \\
\noindent \textbf{(P\_1,2.22)} ktvāgrahaṇam uttarārtham . uttarārtham ktvāgrahaṇam kriyate : nopadhāt thapāntāt vā vañciluñcyṛtaḥ ca iti . \\
\noindent \textbf{(P\_1,2.25)} kāśyapagrahaṇam kimartham. kāśyapagrahaṇam pūjārtham . \\
\noindent \textbf{(P\_1,2.25)} vā iti eva hi vartate . \\
\noindent \textbf{(P\_1,2.26)} kim idam ralaḥ ktvāsanoḥ kittvam vidhīyate āhosvit pratiṣidhyate . \\
\noindent \textbf{(P\_1,2.26)} kim ca ataḥ . \\
\noindent \textbf{(P\_1,2.26)} yadi vidhīyate ktvāgrahaṇam anarthakam . \\
\noindent \textbf{(P\_1,2.26)} kit eva hi ktvā . \\
\noindent \textbf{(P\_1,2.26)} atha pratiṣidhyate sangrahaṇam anarthakam . \\
\noindent \textbf{(P\_1,2.26)} akit eva hi san . \\
\noindent \textbf{(P\_1,2.26)} ataḥ uttaram paṭhati : ralaḥ ktvāsanoḥ kittvam . \\
\noindent \textbf{(P\_1,2.26)} ralaḥ ktvāsanoḥ kittvam vidhīyate . \\
\noindent \textbf{(P\_1,2.26)} nanu ca uktam : ktvāgrahaṇam anarthakam . \\
\noindent \textbf{(P\_1,2.26)} kit eva hi ktvā iti . \\
\noindent \textbf{(P\_1,2.26)} na anarthakam . \\
\noindent \textbf{(P\_1,2.26)} na ktvā seṭ iti pratiṣedhaḥ prāpnoti . \\
\noindent \textbf{(P\_1,2.26)} tadbādhanārtham . \\
\noindent \textbf{(P\_1,2.27.1)} ayuktaḥ ayam nirdeśaḥ . \\
\noindent \textbf{(P\_1,2.27.1)} ū iti anena kālaḥ pratinirdiśyate ū iti ayam ca varṇaḥ . \\
\noindent \textbf{(P\_1,2.27.1)} tatra ayuktam varṇasya kālena saha sāmanādhikaraṇyam . \\
\noindent \textbf{(P\_1,2.27.1)} katham tarhi nirdeśaḥ kartavyaḥ . \\
\noindent \textbf{(P\_1,2.27.1)} ūkālakālasya iti . \\
\noindent \textbf{(P\_1,2.27.1)} kim idam ūkālakālasya iti . \\
\noindent \textbf{(P\_1,2.27.1)} ū iti etasya kālaḥ ūkālaḥ . \\
\noindent \textbf{(P\_1,2.27.1)} ūkālaḥ kālaḥ asya ūkālakālaḥ iti . \\
\noindent \textbf{(P\_1,2.27.1)} saḥ tarhi tathā nirdeśaḥ kartavyaḥ . \\
\noindent \textbf{(P\_1,2.27.1)} na kartavyaḥ . \\
\noindent \textbf{(P\_1,2.27.1)} uttarapadalopaḥ atra draṣṭavyaḥ . \\
\noindent \textbf{(P\_1,2.27.1)} tat yathā uṣṭramukham mukham asya uṣṭramukhaḥ , kharamukhaḥ evam ūkālakālaḥ ūkālaḥ iti . \\
\noindent \textbf{(P\_1,2.27.1)} atha vā sāhacaryāt tācchabdyam bhaviṣyati . \\
\noindent \textbf{(P\_1,2.27.1)} kālasahacaritaḥ varṇaḥ api kālaḥ eva . \\
\noindent \textbf{(P\_1,2.27.2)} hrasvādiṣu samasaṅkhyāprasiddhiḥ nirdeśavaiṣamyāt . \\
\noindent \textbf{(P\_1,2.27.2)} hrasvādiṣu samasaṅkhyatvasya aprasiddhiḥ . \\
\noindent \textbf{(P\_1,2.27.2)} kim kāraṇam . \\
\noindent \textbf{(P\_1,2.27.2)} nirdeśavaiṣamyāt . \\
\noindent \textbf{(P\_1,2.27.2)} tisraḥ sañjñāḥ ekā sañjñī . \\
\noindent \textbf{(P\_1,2.27.2)} vaiṣamyāt saṅkhyātānudeśaḥ na prāpnoti . \\
\noindent \textbf{(P\_1,2.27.2)} siddham tu samasṅkhyatvāt . \\
\noindent \textbf{(P\_1,2.27.2)} siddham etat . \\
\noindent \textbf{(P\_1,2.27.2)} katham . \\
\noindent \textbf{(P\_1,2.27.2)} samasaṅkhyatvāt . \\
\noindent \textbf{(P\_1,2.27.2)} katham samasaṅkhyatvam . \\
\noindent \textbf{(P\_1,2.27.2)} trayāṇām hi vikāranirdeśaḥ . \\
\noindent \textbf{(P\_1,2.27.2)} trayāṇām ayam praśliṣṭanirdeśaḥ . \\
\noindent \textbf{(P\_1,2.27.2)} katham punaḥ jñāyate trayāṇām ayam praśliṣṭanirdeśaḥ iti . \\
\noindent \textbf{(P\_1,2.27.2)} tisṛṇām sañjñānām karaṇasāmarthyāt . \\
\noindent \textbf{(P\_1,2.27.2)} yadi api tāvat tisṛṇām sañjñānām karaṇasāmarthyāt jñāyate trayāṇām ayam praśliṣṭanirdeśaḥ iti kutaḥ tu etat : etena ānupūrvyeṇa sanniviṣṭānām sañjñāḥ bhaviṣyanti iti . \\
\noindent \textbf{(P\_1,2.27.2)} ādau mātrikaḥ tataḥ dvimātraḥ tataḥ trimātraḥ iti . \\
\noindent \textbf{(P\_1,2.27.2)} na punaḥ mātrikaḥ madhye vā ante vā syāt tathā dvimātraḥ ādau vā ante vā syāt tathā trimātraḥ ādau vā madhye vā syāt . \\
\noindent \textbf{(P\_1,2.27.2)} ayam tāvat trimātraḥ aśakyaḥ ādau vā madhye vā kartum . \\
\noindent \textbf{(P\_1,2.27.2)} kutaḥ . \\
\noindent \textbf{(P\_1,2.27.2)} plutāśrayaḥ hi prakṛtibhāvaḥ prasajyeta . \\
\noindent \textbf{(P\_1,2.27.2)} mātrikadvimātrikayoḥ api ghyantam pūrvam nipatati iti mātrikasya pūrvanipātaḥ bhaviṣyati . \\
\noindent \textbf{(P\_1,2.27.2)} yat tāvat ucyate ayam tāvat trimātraḥ aśakyaḥ ādau vā madhye vā kartum . \\
\noindent \textbf{(P\_1,2.27.2)} plutāśrayaḥ hi prakṛtibhāvaḥ prasajyeta iti . \\
\noindent \textbf{(P\_1,2.27.2)} plutāśrayaḥ prakṛtibhāvaḥ plutasañjñā ca anena eva . \\
\noindent \textbf{(P\_1,2.27.2)} yadi ca trimātraḥ ādau vā madhye vā syāt plutasañjñā eva asya na syāt kutaḥ pratkṛtibhāvaḥ . \\
\noindent \textbf{(P\_1,2.27.2)} yat api ucyate mātrikadvimātrikayoḥ api ghyantam pūrvam nipatati iti mātrikasya pūrvanipātaḥ bhaviṣyati iti . \\
\noindent \textbf{(P\_1,2.27.2)} hrasvāśrayā hi ghisañjñā hrasvasañjñā ca anena eva. yadi ca mātrikaḥ madhye vā ante vā syāt hrasvasañjñā eva asya na syāt kutaḥ ghisañjñā kutaḥ pūrvanipātaḥ . \\
\noindent \textbf{(P\_1,2.27.2)} evam eṣā vyavasthā na prakalpate . \\
\noindent \textbf{(P\_1,2.27.2)} evam tarhi ācāryapravṛttiḥ jñāpayati na mātrikaḥ ante bhavati iti yat ayam vibhāṣā pṛṣṭaprativacane heḥ iti mātrikasya plutam śāsti . \\
\noindent \textbf{(P\_1,2.27.2)} katham kṛtvā jñāpakam . \\
\noindent \textbf{(P\_1,2.27.2)} yaḥ ante saḥ plutasañjñakaḥ . \\
\noindent \textbf{(P\_1,2.27.2)} yadi ca mātrikaḥ ante syāt plutasañjñā asya syāt. tatra mātrākālasya mātrākālavacanam anarthakam syāt . \\
\noindent \textbf{(P\_1,2.27.2)} madhye tarhi syāt iti . \\
\noindent \textbf{(P\_1,2.27.2)} atra api ācāryapravṛttiḥ jñāpayati na mātrikaḥ madhye bhavati iti yat ayam ataḥ dīrghaḥ yañi supi ca it dīrghatvam śāsti . \\
\noindent \textbf{(P\_1,2.27.2)} katham kṛtvā jñāpakam . \\
\noindent \textbf{(P\_1,2.27.2)} yaḥ madhye saḥ dīrghasañjñakaḥ . \\
\noindent \textbf{(P\_1,2.27.2)} yadi ca mātrikaḥ madhye syāt dīrghasañjñā asya syāt . \\
\noindent \textbf{(P\_1,2.27.2)} tatra mātrākālasya mātrākālavacanam anarthakam syāt . \\
\noindent \textbf{(P\_1,2.27.2)} dvimātraḥ tarhi ante syāt . \\
\noindent \textbf{(P\_1,2.27.2)} atra api ācāryapravṛttiḥ jñāpayati na dvimātraḥ ante bhavati iti yat ayam om abhyādāne iti dvimātrikasya plutam śāsti . \\
\noindent \textbf{(P\_1,2.27.2)} katham kṛtvā jñāpakam . \\
\noindent \textbf{(P\_1,2.27.2)} yaḥ ante saḥ plutasañjñakaḥ . \\
\noindent \textbf{(P\_1,2.27.2)} yadi ca dvimātrikaḥ ante syāt plutasañjñā asya syāt . \\
\noindent \textbf{(P\_1,2.27.2)} tatra dvimātrakālasya dvimātrakālavacanam anarthakam syāt . \\
\noindent \textbf{(P\_1,2.27.2)} mātrikeṇa ca asya pūrvanipātaḥ bādhitaḥ iti kṛtvā kva anyatra utsahate bhavitum anyat ataḥ madhyāt . \\
\noindent \textbf{(P\_1,2.27.2)} evam eṣā vyavasthā prakḷptā . \\
\noindent \textbf{(P\_1,2.27.2)} bhavet vyavasthā prakḷptā . \\
\noindent \textbf{(P\_1,2.27.2)} dīrghaplutayoḥ tu pūrvasañjñāprasaṅgaḥ . dīrghaplutayoḥ api pūrvasañjñā prāpnoti . \\
\noindent \textbf{(P\_1,2.27.2)} kā . \\
\noindent \textbf{(P\_1,2.27.2)} hrasvasañjñā . \\
\noindent \textbf{(P\_1,2.27.2)} kim kāraṇam . \\
\noindent \textbf{(P\_1,2.27.2)} aṇ savarṇān gṛhṇāti iti . \\
\noindent \textbf{(P\_1,2.27.2)} siddham tu taparanirdeśāt . \\
\noindent \textbf{(P\_1,2.27.2)} siddham etat . \\
\noindent \textbf{(P\_1,2.27.2)} katham . \\
\noindent \textbf{(P\_1,2.27.2)} taparanirdeśaḥ kartavyaḥ : udūkālaḥ iti . \\
\noindent \textbf{(P\_1,2.27.2)} yadi evam drutāyām taparakaraṇe madhyamavilambitayoḥ upasaṅkhyānam kālabhedāt . \\
\noindent \textbf{(P\_1,2.27.2)} drutādiṣu ca uktam . \\
\noindent \textbf{(P\_1,2.27.2)} kim utktam . \\
\noindent \textbf{(P\_1,2.27.2)} siddham tu . \\
\noindent \textbf{(P\_1,2.27.2)} avasthitāḥ varṇāḥ vaktuḥ cirāciravacanāt vṛttayaḥ viśiṣyante iti . \\
\noindent \textbf{(P\_1,2.27.2)} saḥ tarhi taparanirdeśaḥ kartavyaḥ . \\
\noindent \textbf{(P\_1,2.27.2)} na kartavyaḥ . \\
\noindent \textbf{(P\_1,2.27.2)} iha kālagrahaṇam kriyate . \\
\noindent \textbf{(P\_1,2.27.2)} yāvat ca taparakaraṇam tāvat kālagrahaṇam . \\
\noindent \textbf{(P\_1,2.27.2)} pratyekam ca kālaśabdaḥ parisamāpyate : ukālaḥ ūkālaḥ ū3kālaḥ iti . \\
\noindent \textbf{(P\_1,2.27.2)} atha vā ekasañjñādhikāre ayam yogaḥ kartavyaḥ . \\
\noindent \textbf{(P\_1,2.27.2)} tatra ekā sañjñā bhavati yā parā anavakāśā ca iti evam hi dīrghaplutayoḥ pūrvasañjñā na bhaviṣyati . \\
\noindent \textbf{(P\_1,2.27.2)} atha vā svam rūpam śabdasya aśabdasañjñā iti ayam yogaḥ pratyākhyāyate . \\
\noindent \textbf{(P\_1,2.27.2)} tatra yat etat aśabdasañjñā iti etat yayā vibhaktyā nirdiśyamānam arthavat bhavati tayā nirdiṣṭam anuvartiṣyate : aṇudit savarṇasya ca apratyayaḥ aśabdasañjñāyām iti . \\
\noindent \textbf{(P\_1,2.27.2)} atha vā hrasvasañjñāvacanasāmarthyāt dīrghaplutayoḥ pūrvasañjñā na bhaviṣyati . \\
\noindent \textbf{(P\_1,2.27.2)} nanu ca idam prayojanam syāt : sañjñayā vidhāne niyamam vakṣyāmi iti . \\
\noindent \textbf{(P\_1,2.27.2)} hrasvasañjñayā yat ucyate tat acaḥ sthāne yathā syāt iti . \\
\noindent \textbf{(P\_1,2.27.2)} syāt etat prayojanam yadi kiñcitkarāṇi hrasvaśāsanāni syuḥ . \\
\noindent \textbf{(P\_1,2.27.2)} yataḥ tu khalu yāvat ajgrahaṇam tāvat hrasvagrahaṇam ataḥ akiñcitkarāṇi hrasvaśāsanāni . \\
\noindent \textbf{(P\_1,2.27.2)} idam tarhi prayojanam : ecaḥ ik hrasvādeśe iti vakṣyāmi iti . \\
\noindent \textbf{(P\_1,2.27.2)} anucyamāne hi etasmin hrasvapradeśeṣu ecaḥ ik bhavati iti vaktavyam syāt . \\
\noindent \textbf{(P\_1,2.27.2)} hrasvaḥ napuṃsake prātipadikasya ecaḥ ik bhavati iti . \\
\noindent \textbf{(P\_1,2.27.2)} ṇau caṅi upadhāyāḥ hrasvaḥ ecaḥ ik bhavati iti . \\
\noindent \textbf{(P\_1,2.27.2)} hrasvaḥ halādiḥ śeṣaḥ ecaḥ ik bhavati iti . \\
\noindent \textbf{(P\_1,2.27.2)} sañjñā ca nāma yataḥ na laghīyaḥ . \\
\noindent \textbf{(P\_1,2.27.2)} kutaḥ etat . \\
\noindent \textbf{(P\_1,2.27.2)} laghvartham hi sañjñākaraṇam . \\
\noindent \textbf{(P\_1,2.27.2)} laghīyaḥ ca triḥ hrasvapradeśeṣu ecaḥ ik bhavati iti na punaḥ sañjñākaraṇam . \\
\noindent \textbf{(P\_1,2.27.2)} triḥ hrasvapradeśeṣu ecaḥ ik bhavati iti ṣaṭ grahaṇāni . \\
\noindent \textbf{(P\_1,2.27.2)} sañjñākaraṇe punaḥ aṣṭau . \\
\noindent \textbf{(P\_1,2.27.2)} hrasvasañjñā vaktavyā . \\
\noindent \textbf{(P\_1,2.27.2)} triḥ hrasvapradeśeṣu hrasvagrahaṇam kartavyam hrasvaḥ hrasvaḥ hrasvaḥ iti . \\
\noindent \textbf{(P\_1,2.27.2)} ecaḥ ik hrasvādeśe iti . \\
\noindent \textbf{(P\_1,2.27.2)} saḥ ayam laghīyasā nyāsena siddhe sati yat garīyāṃsam yatnam ārabhate tasya etat prayojanam dīrghaplutayoḥ tu pūrvasañjñā mā bhūt iti . \\
\noindent \textbf{(P\_1,2.28.1)} kim ayam alontyaśeṣaḥ āhosvit alontyāpavādaḥ . \\
\noindent \textbf{(P\_1,2.28.1)} katham ca ayam taccheṣaḥ syāt katham vā tadapavādaḥ . \\
\noindent \textbf{(P\_1,2.28.1)} yadi ekam vākyam tat ca idam ca : alaḥ antyasya vidhayaḥ bhavanti acaḥ hrasvadīrghaplutāḥ antyasya iti tataḥ ayam taccheśaḥ . \\
\noindent \textbf{(P\_1,2.28.1)} atha nānā vākyam : alaḥ antyasya vidhayaḥ bhavanti , acaḥ hrasvadīrghaplutāḥ antyasya anantyasya ca iti tataḥ ayam tadapavādaḥ . \\
\noindent \textbf{(P\_1,2.28.1)} kaḥ ca atra viśeṣaḥ . \\
\noindent \textbf{(P\_1,2.28.1)} hrasvādividhiḥ alaḥ antyasya iti cet vacipracchiśamādiprabhṛtihanigamidīrgheṣu ajgrahaṇam . \\
\noindent \textbf{(P\_1,2.28.1)} hrasvādividhiḥ alaḥ antyasya iti cet vacipracchiśamādiprabhṛtihanigamidīrgheṣu ajgrahaṇam kartavyam . \\
\noindent \textbf{(P\_1,2.28.1)} vacipracchyoḥ dīrghaḥ acaḥ iti vaktavyam . \\
\noindent \textbf{(P\_1,2.28.1)} anantyatvāt na prāpnoti . \\
\noindent \textbf{(P\_1,2.28.1)} śamādīnām dīrghaḥ acaḥ iti vaktavyam . \\
\noindent \textbf{(P\_1,2.28.1)} anantyatvāt na prāpnoti . \\
\noindent \textbf{(P\_1,2.28.1)} hanigamyoḥ dīrghaḥ acaḥ iti vaktavyam . \\
\noindent \textbf{(P\_1,2.28.1)} anantyatvāt na prāpnoti . \\
\noindent \textbf{(P\_1,2.28.1)} astu tarhi tadapavādaḥ . \\
\noindent \textbf{(P\_1,2.28.1)} acaḥ cet napuṃsakahrasvākṛtsārvadhātukanāmidīrgheṣu anantyapratiṣedhaḥ . \\
\noindent \textbf{(P\_1,2.28.1)} acaḥ cet napuṃsakahrasvākṛtsārvadhātukanāmidīrgheṣu anantyapratiṣedhaḥ vaktavyaḥ . \\
\noindent \textbf{(P\_1,2.28.1)} hrasvaḥ napuṃsake prātipadikasya : yathā iha bhavati : rai : atiri nau : atinau evam suvāk brāhmaṇakulam iti atra api prāpnoti . \\
\noindent \textbf{(P\_1,2.28.1)} akṛtsārvadhātukayoḥ dīrghaḥ : yathā iha bhavati : cīyate stūyate evam bhidyate atra api prāpnoti . \\
\noindent \textbf{(P\_1,2.28.1)} nāmi dīrghaḥ : yathā iha bhavati : agnīnām , vāyūnām evam atra api prāpnoti : ṣaṇṇām . \\
\noindent \textbf{(P\_1,2.28.1)} na eṣaḥ doṣaḥ . \\
\noindent \textbf{(P\_1,2.28.1)} nopadhyāyāḥ iti etat niyamārtham bhaviṣyati . \\
\noindent \textbf{(P\_1,2.28.1)} prakṛtasya eṣaḥ niyamaḥ syāt. kim ca prakṛtam . \\
\noindent \textbf{(P\_1,2.28.1)} nāmi iti . \\
\noindent \textbf{(P\_1,2.28.1)} tena bhavet iha niyamāt na syāt : ṣaṇṇām . \\
\noindent \textbf{(P\_1,2.28.1)} anyate tanyate atra api prāpnoti . \\
\noindent \textbf{(P\_1,2.28.1)} atha api evam niyamaḥ syāt : nopadhāyāḥ nāmi eva iti evam api bhavet iha niyamāt na syāt : anyate tanyate . \\
\noindent \textbf{(P\_1,2.28.1)} ṣaṇṇām iti atra prāpnoti . \\
\noindent \textbf{(P\_1,2.28.1)} atha api ubhayataḥ niyamaḥ syāt : nopadhāyāḥ eva nāmi nāmi eva nopadadhyāyāḥ iti evam api bhidyate suvāk brāhmaṇakulam ita atra api prāpnoti . \\
\noindent \textbf{(P\_1,2.28.1)} evam tarhi hrasvaḥ dīrghaḥ plutaḥ iti yatra brūyāt acaḥ iti etat tatra upasthitam draṣṭavyam . \\
\noindent \textbf{(P\_1,2.28.1)} kim kṛtam bhavati . \\
\noindent \textbf{(P\_1,2.28.1)} dvitīyā ṣaṣṭhī prāduḥ bhāvyate . \\
\noindent \textbf{(P\_1,2.28.1)} tatra kāmacāraḥ : gṛhyamāṇena vā acam viśeṣayitum acā vā gṛhyamāṇaḥ . \\
\noindent \textbf{(P\_1,2.28.1)} yāvatā kāmacāraḥ iha tāvat vacipracchiśamādiprabhṛtihanigamidīrgheṣu gṛhyamāṇena acam viśeṣayiṣyāmaḥ : eṣām acaḥ dīrghaḥ bhavati iti . \\
\noindent \textbf{(P\_1,2.28.1)} iha idānīm napuṃsakahrasvākṛtsārvadhātukanāmidīrgheṣu acā gṛhyamāṇam viśeṣayiṣyāmaḥ : napuṃsakasya hrasvaḥ bhavati acaḥ . \\
\noindent \textbf{(P\_1,2.28.1)} ajantasya iti . \\
\noindent \textbf{(P\_1,2.28.1)} akṛtsārvadhātukayoḥ dīrghaḥ acaḥ . \\
\noindent \textbf{(P\_1,2.28.1)} ajantasya iti . \\
\noindent \textbf{(P\_1,2.28.1)} nāmi dīrghaḥ bhavati acaḥ . \\
\noindent \textbf{(P\_1,2.28.1)} ajantasya iti . \\
\noindent \textbf{(P\_1,2.28.2)} iha kasmāt na bhavati : dyauḥ , panthāḥ , saḥ iti . \\
\noindent \textbf{(P\_1,2.28.2)} sañjñayā vidhāne niyamaḥ . \\
\noindent \textbf{(P\_1,2.28.2)} sañjñayā ye vidhīyante teṣu niyamaḥ . \\
\noindent \textbf{(P\_1,2.28.2)} kim vaktavyam etat . \\
\noindent \textbf{(P\_1,2.28.2)} na hi . \\
\noindent \textbf{(P\_1,2.28.2)} katham anucyamānam gaṃsyate . \\
\noindent \textbf{(P\_1,2.28.2)} ac iti vartate . \\
\noindent \textbf{(P\_1,2.28.2)} tatra evam abhisambandhaḥ kariṣyate : acaḥ ac bhavati hrasvaḥ dīrghaḥ plutaḥ iti evam bhāvyamanaḥ iti . \\
\noindent \textbf{(P\_1,2.28.2)} atha pūrvasmin yoge ajgrahaṇe sati kim prayojanam . \\
\noindent \textbf{(P\_1,2.28.2)} ajgrahaṇam saṃyogācsamudāyanivṛttyartham . \\
\noindent \textbf{(P\_1,2.28.2)} ajgrahaṇam kriyate saṃyognivṛttyartham acsamudāyanivṛttyartham ca . \\
\noindent \textbf{(P\_1,2.28.2)} saṃyognivṛttyartham tāvat : pratakṣya prarakṣya . \\
\noindent \textbf{(P\_1,2.28.2)} hrasvasya piti kṛti tuk iti tuk mā bhūt iti . \\
\noindent \textbf{(P\_1,2.28.2)} acsamudāyanivṛttyartham : titaucchatram , titaucchāyā . \\
\noindent \textbf{(P\_1,2.28.2)} dīrghāt padāntāt vā iti vibhāṣā mā bhūt . \\
\noindent \textbf{(P\_1,2.29-30.1)} KA\_I,206.14-25 Ro\_II,43-45 {1/23}     kim ṣaṣṭhīnirdiṣṭam ajgrahaṇam anuvartate utāho na . \\
\noindent \textbf{(P\_1,2.29-30.1)} KA\_I,206.14-25 Ro\_II,43-45 {2/23}     kim ca ataḥ . \\
\noindent \textbf{(P\_1,2.29-30.1)} KA\_I,206.14-25 Ro\_II,43-45 {3/23}     yadi anuvartate halsvaraprāptau vyañjanam avidyamānavat iti paribhāṣā na prakalpate . \\
\noindent \textbf{(P\_1,2.29-30.1)} KA\_I,206.14-25 Ro\_II,43-45 {4/23}     katham halaḥ nāma svaraprāptiḥ syāt . \\
\noindent \textbf{(P\_1,2.29-30.1)} KA\_I,206.14-25 Ro\_II,43-45 {5/23}     evam tarhi nivṛttam . \\
\noindent \textbf{(P\_1,2.29-30.1)} KA\_I,206.14-25 Ro\_II,43-45 {6/23}     bahūni etasyāḥ paribhāṣāyāḥ prayojanāni . \\
\noindent \textbf{(P\_1,2.29-30.1)} KA\_I,206.14-25 Ro\_II,43-45 {7/23}     atha prathamānirdiṣṭam ajgrahaṇam anuvartate utāho na . \\
\noindent \textbf{(P\_1,2.29-30.1)} KA\_I,206.14-25 Ro\_II,43-45 {8/23}     kim ca arthaḥ anuvṛttyā . \\
\noindent \textbf{(P\_1,2.29-30.1)} KA\_I,206.14-25 Ro\_II,43-45 {9/23}     bāḍham arthaḥ yadi ete vyañjanasya api guṇāḥ lakṣyante . \\
\noindent \textbf{(P\_1,2.29-30.1)} KA\_I,206.14-25 Ro\_II,43-45 {10/23}     nanu ca pratyakṣam upalabhyante : iṣe tvā ūrje tvā . \\
\noindent \textbf{(P\_1,2.29-30.1)} KA\_I,206.14-25 Ro\_II,43-45 {11/23}     na ete vyañjanasya guṇāḥ . \\
\noindent \textbf{(P\_1,2.29-30.1)} KA\_I,206.14-25 Ro\_II,43-45 {12/23}     acaḥ ete guṇāḥ . \\
\noindent \textbf{(P\_1,2.29-30.1)} KA\_I,206.14-25 Ro\_II,43-45 {13/23}     tatsāmīpyāt tu vyañjanam api tadguṇam upalabhyate . \\
\noindent \textbf{(P\_1,2.29-30.1)} KA\_I,206.14-25 Ro\_II,43-45 {14/23}     tat yathā . \\
\noindent \textbf{(P\_1,2.29-30.1)} KA\_I,206.14-25 Ro\_II,43-45 {15/23}     dvayoḥ raktayoḥ vastrayoḥ madhye śuklam vastram tadguṇam upalabhyate badarapiṭake riktakaḥ lohakaṃsaḥ tadguṇaḥ upalabhyate . \\
\noindent \textbf{(P\_1,2.29-30.1)} KA\_I,206.14-25 Ro\_II,43-45 {16/23}     kutaḥ nu khalu etat acaḥ ete guṇāḥ . \\
\noindent \textbf{(P\_1,2.29-30.1)} KA\_I,206.14-25 Ro\_II,43-45 {17/23}     tatsāmīpyāt tu vyañjanam api tadguṇam upalabhyate iti . \\
\noindent \textbf{(P\_1,2.29-30.1)} KA\_I,206.14-25 Ro\_II,43-45 {18/23}     na punaḥ vyañjanasya ete guṇāḥ syuḥ . \\
\noindent \textbf{(P\_1,2.29-30.1)} KA\_I,206.14-25 Ro\_II,43-45 {19/23}     tatsāmīpyāt tu ac api tadguṇaḥ upalabhyate iti . \\
\noindent \textbf{(P\_1,2.29-30.1)} KA\_I,206.14-25 Ro\_II,43-45 {20/23}     antareṇa api vyañjanam acaḥ eva ete guṇāḥ lakṣyante . \\
\noindent \textbf{(P\_1,2.29-30.1)} KA\_I,206.14-25 Ro\_II,43-45 {21/23}     na punaḥ antareṇa acam vyañjanasya uccāraṇam api bhavati . \\
\noindent \textbf{(P\_1,2.29-30.1)} KA\_I,206.14-25 Ro\_II,43-45 {22/23}     anvartham khalu api nirvacanam : svayam rājante svarāḥ . \\
\noindent \textbf{(P\_1,2.29-30.1)} KA\_I,206.14-25 Ro\_II,43-45 {23/23}     anvak bhavati vyañjanam iti . \\
\noindent \textbf{(P\_1,2.29-30.2)} KA\_I,207.1-17 Ro\_II,45-46 {1/28}     uccanīcasya anavasthitatvāt sañjñāprasiddhiḥ . \\
\noindent \textbf{(P\_1,2.29-30.2)} KA\_I,207.1-17 Ro\_II,45-46 {2/28}     idam uccanīcam anavasthitapadarthakam . \\
\noindent \textbf{(P\_1,2.29-30.2)} KA\_I,207.1-17 Ro\_II,45-46 {3/28}     tat eva hi kam cit prati uccaiḥ bhavati kam cit prati nīcaiḥ . \\
\noindent \textbf{(P\_1,2.29-30.2)} KA\_I,207.1-17 Ro\_II,45-46 {4/28}     evam kam cit kaḥ cit adhīyānam āha : kim uccaiḥ rorūyase . \\
\noindent \textbf{(P\_1,2.29-30.2)} KA\_I,207.1-17 Ro\_II,45-46 {5/28}     atha nīcaiḥ vartatām iti . \\
\noindent \textbf{(P\_1,2.29-30.2)} KA\_I,207.1-17 Ro\_II,45-46 {6/28}     tam eva tathā adhīyānam aparaḥ āha : kim antardantakena adhīṣe . \\
\noindent \textbf{(P\_1,2.29-30.2)} KA\_I,207.1-17 Ro\_II,45-46 {7/28}     uccaiḥ vartatām iti . \\
\noindent \textbf{(P\_1,2.29-30.2)} KA\_I,207.1-17 Ro\_II,45-46 {8/28}     evam uccanīcam anavasthitapadarthakam . \\
\noindent \textbf{(P\_1,2.29-30.2)} KA\_I,207.1-17 Ro\_II,45-46 {9/28}     tasya anavasthānāt sañjñāyāḥ aprasiddhiḥ . \\
\noindent \textbf{(P\_1,2.29-30.2)} KA\_I,207.1-17 Ro\_II,45-46 {10/28}     evam tarhi lakṣaṇam kariṣyate : āyāmaḥ dāruṇyam aṇutā khasya iti uccaiḥkarāṇi śabdasya . \\
\noindent \textbf{(P\_1,2.29-30.2)} KA\_I,207.1-17 Ro\_II,45-46 {11/28}     āyāmaḥ gātrāṇām nigrahaḥ . \\
\noindent \textbf{(P\_1,2.29-30.2)} KA\_I,207.1-17 Ro\_II,45-46 {12/28}     dāruṇyam svarasya dāruṇatā rūkṣatā . \\
\noindent \textbf{(P\_1,2.29-30.2)} KA\_I,207.1-17 Ro\_II,45-46 {13/28}     aṇutā khasya kaṇṭhasya saṃvṛtatā . \\
\noindent \textbf{(P\_1,2.29-30.2)} KA\_I,207.1-17 Ro\_II,45-46 {14/28}     uccaiḥkarāṇi śabdasya . \\
\noindent \textbf{(P\_1,2.29-30.2)} KA\_I,207.1-17 Ro\_II,45-46 {15/28}     atha nīcaiḥkarāṇi śabdasya . \\
\noindent \textbf{(P\_1,2.29-30.2)} KA\_I,207.1-17 Ro\_II,45-46 {16/28}     anvavasargaḥ mārdavam urutā khasya iti nīcaiḥkarāṇi śabdasya . \\
\noindent \textbf{(P\_1,2.29-30.2)} KA\_I,207.1-17 Ro\_II,45-46 {17/28}     anvavasargaḥ gātrāṇām śithilatā . \\
\noindent \textbf{(P\_1,2.29-30.2)} KA\_I,207.1-17 Ro\_II,45-46 {18/28}     mārdavam svarasya mṛdutā snigdhatā . \\
\noindent \textbf{(P\_1,2.29-30.2)} KA\_I,207.1-17 Ro\_II,45-46 {19/28}     urutā khasya mahattā kaṇṭhasya . \\
\noindent \textbf{(P\_1,2.29-30.2)} KA\_I,207.1-17 Ro\_II,45-46 {20/28}     iti nīcaiḥkarāṇi śabdasya . \\
\noindent \textbf{(P\_1,2.29-30.2)} KA\_I,207.1-17 Ro\_II,45-46 {21/28}     etat api anaikāntikam . \\
\noindent \textbf{(P\_1,2.29-30.2)} KA\_I,207.1-17 Ro\_II,45-46 {22/28}     yat alpaprāṇasya sarvoccaiḥ tat mahāprāṇasya sarvanīcaiḥ . \\
\noindent \textbf{(P\_1,2.29-30.2)} KA\_I,207.1-17 Ro\_II,45-46 {23/28}     siddham tu samānaprakramavacanāt . \\
\noindent \textbf{(P\_1,2.29-30.2)} KA\_I,207.1-17 Ro\_II,45-46 {24/28}     siddham etat . \\
\noindent \textbf{(P\_1,2.29-30.2)} KA\_I,207.1-17 Ro\_II,45-46 {25/28}     katham . \\
\noindent \textbf{(P\_1,2.29-30.2)} KA\_I,207.1-17 Ro\_II,45-46 {26/28}      samāne prakrame iti vaktavyam . \\
\noindent \textbf{(P\_1,2.29-30.2)} KA\_I,207.1-17 Ro\_II,45-46 {27/28}     kaḥ punaḥ prakramaḥ . \\
\noindent \textbf{(P\_1,2.29-30.2)} KA\_I,207.1-17 Ro\_II,45-46 {28/28}     uraḥ kaṇṭhaḥ śiraḥ iti . \\
\noindent \textbf{(P\_1,2.31)} samāhāraḥ svaritaḥ iti ucyate . \\
\noindent \textbf{(P\_1,2.31)} kasya samāhāraḥ svaritasañjñaḥ bhavati . \\
\noindent \textbf{(P\_1,2.31)} acoḥ iti āha . \\
\noindent \textbf{(P\_1,2.31)} samāhāraḥ acoḥ cet na abhāvāt . \\
\noindent \textbf{(P\_1,2.31)} samāhāraḥ acoḥ cet tat na . \\
\noindent \textbf{(P\_1,2.31)} kim kāraṇam . \\
\noindent \textbf{(P\_1,2.31)} abhāvāt . \\
\noindent \textbf{(P\_1,2.31)} na hi acoḥ samāhāraḥ asti . \\
\noindent \textbf{(P\_1,2.31)} nanu ayam asti gāṅgenūpe iti . \\
\noindent \textbf{(P\_1,2.31)} na eṣaḥ acoḥ samāhāraḥ . \\
\noindent \textbf{(P\_1,2.31)} anyaḥ ayam udāttānudāttayoḥ sthāne ekaḥ ādiśyate . \\
\noindent \textbf{(P\_1,2.31)} evam tarhi guṇayoḥ . \\
\noindent \textbf{(P\_1,2.31)} guṇayoḥ cet na acprakaraṇāt . \\
\noindent \textbf{(P\_1,2.31)} guṇayoḥ samahāraḥ iti cet tat na . \\
\noindent \textbf{(P\_1,2.31)} kim kāraṇam . \\
\noindent \textbf{(P\_1,2.31)} acprakaraṇāt . \\
\noindent \textbf{(P\_1,2.31)} ac iti vartate . \\
\noindent \textbf{(P\_1,2.31)} siddham tu acsamudāyasya abhāvāt tadguṇe sampratyayaḥ . \\
\noindent \textbf{(P\_1,2.31)} siddham etat . \\
\noindent \textbf{(P\_1,2.31)} katham . \\
\noindent \textbf{(P\_1,2.31)} acsamudāyaḥ na asti iti kṛtvā tadguṇasya acaḥ samāhāraguṇasya sampratyayaḥ bhaviṣyati . \\
\noindent \textbf{(P\_1,2.31)} katham punaḥ samāhāraḥ iti anena ac śakyaḥ pratinirdeṣṭum . \\
\noindent \textbf{(P\_1,2.31)} matublopaḥ atra draṣṭavyaḥ . \\
\noindent \textbf{(P\_1,2.31)} tat yathā puṣpakāḥ eṣām te puṣpakāḥ , kālakāḥ eṣām te kālakāḥ iti evam samāhāravān samāhāraḥ . \\
\noindent \textbf{(P\_1,2.31)} atha vā akāraḥ matvarthīyaḥ . \\
\noindent \textbf{(P\_1,2.31)} tat yathā tundaḥ , ghāṭaḥ iti . \\
\noindent \textbf{(P\_1,2.31)} yadi evam traisvaryam na prakalpate . \\
\noindent \textbf{(P\_1,2.31)} tatra kaḥ doṣaḥ . \\
\noindent \textbf{(P\_1,2.31)} traisvaryam adhīmahe iti etat na upapadyate . \\
\noindent \textbf{(P\_1,2.31)} na etat guṇāpekṣam . \\
\noindent \textbf{(P\_1,2.31)} kim tarhi . \\
\noindent \textbf{(P\_1,2.31)} ajapekṣam . \\
\noindent \textbf{(P\_1,2.31)} traisvaryam adhīmahe : triprakāraiḥ ajbhiḥ adhīmahe kaiḥ cit udāttaguṇaiḥ kaiḥ cit anudāttaguṇaiḥ kaiḥ cit ubhayaguṇaiḥ . \\
\noindent \textbf{(P\_1,2.31)} tat yathā : śuklaguṇaḥ śuklaḥ kṛṣṇaguṇaḥ kṛṣṇaḥ . \\
\noindent \textbf{(P\_1,2.31)} yaḥ idānīm ubhayaguṇaḥ saḥ tṛtīyām ākhyām labhate kalmāṣaḥ iti vā sāraṅgaḥ iti vā . \\
\noindent \textbf{(P\_1,2.31)} evam iha api udāttaguṇaḥ udāttaḥ anudāttaguṇaḥ anudāttaḥ . \\
\noindent \textbf{(P\_1,2.31)} yaḥ idānīm ubhayavān sa tṛtīyām ākhyām labhate svaritaḥ iti . \\
\noindent \textbf{(P\_1,2.32.1)} ardhahrasvam iti ucyate . \\
\noindent \textbf{(P\_1,2.32.1)} tatra dīrghaplutayoḥ na prāpnoti . \\
\noindent \textbf{(P\_1,2.32.1)} kanyā śaktike3 śaktike . \\
\noindent \textbf{(P\_1,2.32.1)} na eṣaḥ doṣaḥ . \\
\noindent \textbf{(P\_1,2.32.1)} mātracaḥ atra lopaḥ draṣṭavyaḥ . \\
\noindent \textbf{(P\_1,2.32.1)} ardhahrasvamātram ardhahrasvam iti . \\
\noindent \textbf{(P\_1,2.32.1)} kimartham idam ucyate . \\
\noindent \textbf{(P\_1,2.32.1)} āmiśrībhūtam iva idam bhavati . \\
\noindent \textbf{(P\_1,2.32.1)} tat yathā : kṣīrodake sampṛkte* āmiśrībhūtatvāt na jñāyate : kiyat kṣīram kiyat udakam kasmin avakāśe kṣīram kasmin avakāśe udakam iti . \\
\noindent \textbf{(P\_1,2.32.1)} evam iha api āmiśrībhūtatvāt na jñāyate : kiyat udāttam kiyat anudāttam kasmin avakāśe udāttam kasmin avakāśe anudāttam iti . \\
\noindent \textbf{(P\_1,2.32.1)} tat ācāryaḥ suhṛt bhūtvā anvācaṣṭe : iyat udāttam iyat anudāttam asmin avakāśe udāttam asmin avakāśe anudāttam iti . \\
\noindent \textbf{(P\_1,2.32.1)} yadi ayam evam suhṛt kim anyāni api evañjātīyakāni na upadiśati . \\
\noindent \textbf{(P\_1,2.32.1)} kāni punaḥ tāni . \\
\noindent \textbf{(P\_1,2.32.1)} sthānakaraṇānupradānāni . \\
\noindent \textbf{(P\_1,2.32.1)} vyākaraṇam nāma iyam uttarā vidyā . \\
\noindent \textbf{(P\_1,2.32.1)} saḥ asau chandaḥśāstreṣu abhivinītaḥ upalabdhyā avagantum utsahate . \\
\noindent \textbf{(P\_1,2.32.1)} yadi evam na arthaḥ anena . \\
\noindent \textbf{(P\_1,2.32.1)} idam api upalabdhyā gamiṣyati . \\
\noindent \textbf{(P\_1,2.32.1)} sañjñākaraṇam tarhi idam : tasya svaritasya āditaḥ ardhahrasvam udāttasañjñam iti . \\
\noindent \textbf{(P\_1,2.32.1)} kim kṛtam bhavati . \\
\noindent \textbf{(P\_1,2.32.1)} triḥ udāttapradeśeṣu svaritagrahaṇam na kartavyam bhavati : udāttasvaritaparasya sannataraḥ , udāttasvaritayoḥ yaṇaḥ svaritaḥ anudāttasya na udāttasvaritodayam iti . \\
\noindent \textbf{(P\_1,2.32.1)} sañjñākaraṇam hi nāma yataḥ na laghīyaḥ . \\
\noindent \textbf{(P\_1,2.32.1)} kutaḥ etat . \\
\noindent \textbf{(P\_1,2.32.1)} laghvartham hi sañjñākaraṇam . \\
\noindent \textbf{(P\_1,2.32.1)} laghīyaḥ ca triḥ udāttapradeśeṣu svaritagrahaṇam na punaḥ sañjñākaraṇam . \\
\noindent \textbf{(P\_1,2.32.1)} triḥ udāttapradeśeṣu svaritagrahaṇe nava akṣarāṇi sañjñākaraṇe punaḥ ekādaśa . \\
\noindent \textbf{(P\_1,2.32.1)} evam tarhi ubhayam anena kriyate anvākhyānam ca sañjñā ca . \\
\noindent \textbf{(P\_1,2.32.1)} katham punaḥ ekena yatnena ubhayam labhyam . \\
\noindent \textbf{(P\_1,2.32.1)} labhyam iti āha . \\
\noindent \textbf{(P\_1,2.32.1)} katham . \\
\noindent \textbf{(P\_1,2.32.1)} anvarthagrahaṇam vijñāsyate : tasya svaritasya āditaḥ ardhahrasvram udāttasañjñam bhavati iti . \\
\noindent \textbf{(P\_1,2.32.1)} ūrdhvam āttam iti ca ataḥ udāttam . \\
\noindent \textbf{(P\_1,2.32.1)} yadi tarhi sañjñākaraṇam udāttādeḥ yat ucyate tat svaritādeḥ api prāpnoti . \\
\noindent \textbf{(P\_1,2.32.1)} anvākhyānam eva tarhi idam mandabuddheḥ . \\
\noindent \textbf{(P\_1,2.32.2)} svaritasyārdhahrasvodāttāt ā udāttasvaritaparasyasannatarāt ūrdhvam udāttādanudāttasyasvaritāt kāryam svaritāt iti siddhyartham . \\
\noindent \textbf{(P\_1,2.32.2)} svaritasyārdhahrasvodāttāt ā udāttasvaritaparasyasannataraḥ iti etasmāt sūtrāt idam sūtrakāṇḍam ūrdhvam udāttāt anudāttasya svaritaḥ iti ataḥ kāryam . \\
\noindent \textbf{(P\_1,2.32.2)} kim prayojanam . \\
\noindent \textbf{(P\_1,2.32.2)} svaritāt iti siddhyartham . \\
\noindent \textbf{(P\_1,2.32.2)} svaritāt iti siddhiḥ yathā syāt . \\
\noindent \textbf{(P\_1,2.32.2)} svaritāt saṃhitāyām anudāttānām iti : imam me gaṅge yamune sarasvati śutudri . \\
\noindent \textbf{(P\_1,2.32.2)} kva tarhi syāt . \\
\noindent \textbf{(P\_1,2.32.2)} yaḥ siddhaḥ svaritaḥ : kāryam devadattayajñadattau . \\
\noindent \textbf{(P\_1,2.32.2)} svaritodāttārtham ca . \\
\noindent \textbf{(P\_1,2.32.2)} svaritodāttārtham ca tatra eva kartavyam . \\
\noindent \textbf{(P\_1,2.32.2)} na subrahmaṇyāyām svaritasya tu udāttaḥ : indra agaccha . \\
\noindent \textbf{(P\_1,2.32.2)} kva tarhi syāt . \\
\noindent \textbf{(P\_1,2.32.2)} yaḥ siddhaḥ svaritaḥ : subrahmaṇyom indra agaccha . \\
\noindent \textbf{(P\_1,2.32.2)} svaritodāttāt ca asvaritārtham . \\
\noindent \textbf{(P\_1,2.32.2)} svaritodāttāt ca asvaritārtham tatra eva kartavyam . \\
\noindent \textbf{(P\_1,2.32.2)} indra agaccha harivaḥ agaccha . \\
\noindent \textbf{(P\_1,2.32.2)} svaritaparasannatarārtham ca . \\
\noindent \textbf{(P\_1,2.32.2)} svaritaparasannatarārtham ca tatra eva kartavyam . \\
\noindent \textbf{(P\_1,2.32.2)} udāttasvaritaparasya sannataraḥ : maṇavaka jaṭilakādhyāpaka nyaṅ . \\
\noindent \textbf{(P\_1,2.32.2)} kva tarhi syāt . \\
\noindent \textbf{(P\_1,2.32.2)} yaḥ siddhaḥ svaritaḥ : maṇavaka jaṭilakābhirūpaka kva . \\
\noindent \textbf{(P\_1,2.32.2)} tat tarhi vaktavyam . \\
\noindent \textbf{(P\_1,2.32.2)} na vaktavyam . \\
\noindent \textbf{(P\_1,2.32.2)} devabrahmaṇoḥ anudāttavacanam jñāpakam svaritāt iti siddhatvasya . \\
\noindent \textbf{(P\_1,2.32.2)} devabrahmaṇoḥ anudāttavacanam jñāpakam siddhaḥ iha svaritaḥ iti . \\
\noindent \textbf{(P\_1,2.32.2)} yadi etat jñāpyate svaritodāttaparasya anudāttasya svaritatvam prāpnoti . \\
\noindent \textbf{(P\_1,2.32.2)} na brūmaḥ devabrahmaṇoḥ anudāttavacanam jñāpakam siddhaḥ iha svaritaḥ iti . \\
\noindent \textbf{(P\_1,2.32.2)} kim tarhi . \\
\noindent \textbf{(P\_1,2.32.2)} param etat kāṇḍam iti . \\
\noindent \textbf{(P\_1,2.33.1)} kim idam pāribhāṣikyāḥ sambuddheḥ grahaṇam : ekavacanam sambuddhiḥ āhosvit anvarthagrahaṇam : sambodhanam sambuddhiḥ iti . \\
\noindent \textbf{(P\_1,2.33.1)} kim ca ataḥ . \\
\noindent \textbf{(P\_1,2.33.1)} yadi pāribhāṣikyāḥ devāḥ brahmāṇaḥ iti atra na prāpnoti . \\
\noindent \textbf{(P\_1,2.33.1)} atha anvarthagrahaṇam na doṣaḥ . \\
\noindent \textbf{(P\_1,2.33.1)} yathā na doṣaḥ tathā astu . \\
\noindent \textbf{(P\_1,2.33.2)} kim punaḥ iyam ekaśrutiḥ udāttā āhosvit anudāttā . \\
\noindent \textbf{(P\_1,2.33.2)} na udāttā . \\
\noindent \textbf{(P\_1,2.33.2)} katham jñāyate . \\
\noindent \textbf{(P\_1,2.33.2)} yat ayam uccaistarām vā vaṣaṭkāraḥ iti āha . \\
\noindent \textbf{(P\_1,2.33.2)} katham kṛtvā jñāpakam . \\
\noindent \textbf{(P\_1,2.33.2)} atantram taranirdeśaḥ . \\
\noindent \textbf{(P\_1,2.33.2)} yāvat uccaiḥ tāvat uccaistarām iti . \\
\noindent \textbf{(P\_1,2.33.2)} yadi tarhi na udāttā anudāttā . \\
\noindent \textbf{(P\_1,2.33.2)} anudāttā ca na . \\
\noindent \textbf{(P\_1,2.33.2)} katham jñāyate . \\
\noindent \textbf{(P\_1,2.33.2)} yat ayam udāttasvaritaparasya sannataraḥ iti āha . \\
\noindent \textbf{(P\_1,2.33.2)} katham kṛtvā jñāpakam . \\
\noindent \textbf{(P\_1,2.33.2)} atantram taranirdeśaḥ . \\
\noindent \textbf{(P\_1,2.33.2)} yāvat sannaḥ tāvat sannataraḥ iti . \\
\noindent \textbf{(P\_1,2.33.2)} sā eṣā jñāpakābhyām udāttānudāttayoḥ madhyam ekaśrutiḥ antarālam hriyate . \\
\noindent \textbf{(P\_1,2.33.2)} aparaḥ āha : kim punaḥ iyam ekaśrutiḥ udāttā uta anudāttā . \\
\noindent \textbf{(P\_1,2.33.2)} udāttā . \\
\noindent \textbf{(P\_1,2.33.2)} katham jñāyate . \\
\noindent \textbf{(P\_1,2.33.2)} yat ayam uccaistarām vā vaṣaṭkāraḥ iti āha . \\
\noindent \textbf{(P\_1,2.33.2)} katham kṛtvā jñāpakam . \\
\noindent \textbf{(P\_1,2.33.2)} tantram taranirdeśaḥ . \\
\noindent \textbf{(P\_1,2.33.2)} uccaiḥ dṛṣṭvā uccaistarām iti etat bhavati . \\
\noindent \textbf{(P\_1,2.33.2)} yadi tarhi udāttā na anudāttā . \\
\noindent \textbf{(P\_1,2.33.2)} anudāttā ca . \\
\noindent \textbf{(P\_1,2.33.2)} katham jñāyate . \\
\noindent \textbf{(P\_1,2.33.2)} yat ayam udāttasvaritaparasya sannataraḥ iti āha . \\
\noindent \textbf{(P\_1,2.33.2)} katham kṛtvā jñāpakam . \\
\noindent \textbf{(P\_1,2.33.2)} tantram taranirdeśaḥ . \\
\noindent \textbf{(P\_1,2.33.2)} sannam dṛṣṭvā sannataraḥ iti etat bhavati . \\
\noindent \textbf{(P\_1,2.33.2)} te ete tantre taranirdeśe sapta svarāḥ bhavanti : udāttaḥ , udāttataraḥ , anudāttaḥ , anudāttaraḥ , svaritaḥ , svarite yaḥ udāttaḥ saḥ anyena viśiṣtaḥ , ekaśrutiḥ saptamaḥ . \\
\noindent \textbf{(P\_1,2.37)} subrahmaṇyāyam okāraḥ udāttaḥ . \\
\noindent \textbf{(P\_1,2.37)} subrahmaṇyāyam okāraḥ udāttaḥ bhavati : subrahmaṇyom . \\
\noindent \textbf{(P\_1,2.37)} ākāraḥ ākhyāte parādiḥ ca . \\
\noindent \textbf{(P\_1,2.37)} ākāraḥ ākhyāte parādiḥ ca udāttaḥ bhavati : indra agaccha . \\
\noindent \textbf{(P\_1,2.37)} harivaḥ agaccha . \\
\noindent \textbf{(P\_1,2.37)} vākyādau ca dve dve . \\
\noindent \textbf{(P\_1,2.37)} vākyādau ca dve dve udātte bhavataḥ : indra agaccha . \\
\noindent \textbf{(P\_1,2.37)} harivaḥ agaccha . \\
\noindent \textbf{(P\_1,2.37)} maghavanvarjam . agaccha maghavan . \\
\noindent \textbf{(P\_1,2.37)} sutyāparāṇām antaḥ . \\
\noindent \textbf{(P\_1,2.37)} sutyāparāṇām antaḥ udāttaḥ bhavati : dvyahe sutyam . \\
\noindent \textbf{(P\_1,2.37)} tryahe sutyam . \\
\noindent \textbf{(P\_1,2.37)} asau iti antaḥ . \\
\noindent \textbf{(P\_1,2.37)} asau iti antaḥ udāttaḥ bhavati : gārgyaḥ yajate . \\
\noindent \textbf{(P\_1,2.37)} vātsyaḥ yajate . \\
\noindent \textbf{(P\_1,2.37)} amuṣya iti antaḥ . amuṣya iti antaḥ : dākṣeḥ pita yajate . \\
\noindent \textbf{(P\_1,2.37)} syāntasya upottamam ca . \\
\noindent \textbf{(P\_1,2.37)} syāntasya upottamam udāttam bhavati antaḥ ca . \\
\noindent \textbf{(P\_1,2.37)} gārgyasya pita yajate . \\
\noindent \textbf{(P\_1,2.37)} vātsyasya pita yajate . \\
\noindent \textbf{(P\_1,2.37)} vā nāmadheyasya . \\
\noindent \textbf{(P\_1,2.37)} vā nāmadheyasya syāntasya upottamam udāttam bhavati : devadattasya pita yajate . \\
\noindent \textbf{(P\_1,2.37)} devadattasya pita yajate . \\
\noindent \textbf{(P\_1,2.38)} devabrahmaṇoḥ anudāttatvam eke . \\
\noindent \textbf{(P\_1,2.38)} devabrahmaṇoḥ anudāttatvam eke icchanti : devāḥ brahmāṇaḥ . \\
\noindent \textbf{(P\_1,2.38)} devāḥ brahmāṇaḥ . \\
\noindent \textbf{(P\_1,2.39)} svaritāt saṃhitāyām anudāttānām iti cet dvyekayoḥ aikaśrutyavacanam . \\
\noindent \textbf{(P\_1,2.39)} svaritāt saṃhitāyām anudāttānām iti cet dvyekayoḥ aikaśrutyam vaktavyam : agniveśyaḥ pacati . \\
\noindent \textbf{(P\_1,2.39)} kim punaḥ kāraṇam na sidhyati . \\
\noindent \textbf{(P\_1,2.39)} bahuvacanena nirdeśaḥ kriyate . \\
\noindent \textbf{(P\_1,2.39)} tena bahūnām aikaśrutyam syāt dvyekayoḥ na syāt . \\
\noindent \textbf{(P\_1,2.39)} na eṣaḥ doṣaḥ . \\
\noindent \textbf{(P\_1,2.39)} na atra nirdeśaḥ tantram . \\
\noindent \textbf{(P\_1,2.39)} katham punaḥ tena eva nirdeśaḥ kriyate tat ca atantram syāt . \\
\noindent \textbf{(P\_1,2.39)} tatkārī ca bhavān taddveṣī ca . \\
\noindent \textbf{(P\_1,2.39)} nāntarīyakatvāt atra bahuvacanena nirdeśaḥ kriyate : avaśyam kayā cit vibhaktyā kena cit vacanena nirdeśaḥ kartavyaḥ iti . \\
\noindent \textbf{(P\_1,2.39)} tat yathā : kaḥ cit annārthī śālikalāpam sapalālam satuṣam āharati nāntayīyakatvāt . \\
\noindent \textbf{(P\_1,2.39)} saḥ yāvat ādeyam tāvat ādāya tuṣapalālāni utsṛjati . \\
\noindent \textbf{(P\_1,2.39)} tathā kaḥ cit māṃsārthī matsyān sakaṇṭakān saśakalān āharati nāntayīyakatvāt . \\
\noindent \textbf{(P\_1,2.39)} saḥ yāvat ādeyam tāvat ādāya śakalakaṇṭakān utsṛjati . \\
\noindent \textbf{(P\_1,2.39)} evam iha api nāntarīyakatvāt bahuvacanena nirdeśaḥ kriyate . \\
\noindent \textbf{(P\_1,2.39)} aviśeṣeṇa aikaśrutyam . \\
\noindent \textbf{(P\_1,2.39)} aviśeṣeṇa aikaśrutyam iti cet vyavahitānām aprasiddhiḥ . \\
\noindent \textbf{(P\_1,2.39)} aviśeṣeṇa aikaśrutyam iti cet vyavahitānām aikaśrutyam na prāpnoti : imam me gaṅge yamune sarasvati śutudri . \\
\noindent \textbf{(P\_1,2.39)} anekam api iti tu vacanāt siddham . \\
\noindent \textbf{(P\_1,2.39)} anekam api ekam api svaritāt param saṃhitāyām ekaśruti bhavati iti vaktavyam . \\
\noindent \textbf{(P\_1,2.39)} sidhyati . \\
\noindent \textbf{(P\_1,2.39)} sūtram tarhi bhidyate . \\
\noindent \textbf{(P\_1,2.39)} yathānyāsam eva astu . \\
\noindent \textbf{(P\_1,2.39)} nanu ca uktam : svaritāt saṃhitāyām anudāttānām iti cet dvyekayoḥ aikaśrutyavacanam . \\
\noindent \textbf{(P\_1,2.39)} aviśeṣeṇa aikaśrutyam vyavahitānām aprasiddhiḥ iti . \\
\noindent \textbf{(P\_1,2.39)} na eṣaḥ doṣaḥ . \\
\noindent \textbf{(P\_1,2.39)} katham . \\
\noindent \textbf{(P\_1,2.39)} ekaśeṣanirdeśaḥ ayam anudāttasya ca : anudāttayoḥ ca anudāttānām ca anudāttānām iti . \\
\noindent \textbf{(P\_1,2.39)} evam api ṣaṭprabhṛtīnām eva prāpnoti . \\
\noindent \textbf{(P\_1,2.39)} ṣaṭprabhṛtiṣu ekaśeṣaḥ parisamāpyate . \\
\noindent \textbf{(P\_1,2.39)} pratyekam vākyaparisamāptiḥ dṛṣṭā iti dvyekayoḥ api bhaviṣyati . \\
\noindent \textbf{(P\_1,2.41)} apṛktasañjñāyām halgrahaṇam svādilope halaḥ agrahaṇārtham . \\
\noindent \textbf{(P\_1,2.41)} apṛktasañjñāyām halgrahaṇam kartavyam . \\
\noindent \textbf{(P\_1,2.41)} ekahal pratyayaḥ apṛktasañjñaḥ bhavati iti vaktavyam . \\
\noindent \textbf{(P\_1,2.41)} kim prayojanam . \\
\noindent \textbf{(P\_1,2.41)} svādilope halaḥ agrahaṇārtham . \\
\noindent \textbf{(P\_1,2.41)} svādilope halaḥ grahaṇam na kartavyam bhavati : halṅyābbhyaḥ dīrghāt sutisi apṛktam hal iti apṛktasya iti eva siddham . \\
\noindent \textbf{(P\_1,2.41)} aṇiñoḥ lugartham algrahaṇam . \\
\noindent \textbf{(P\_1,2.41)} aṇiñoḥ lugartham algrahaṇam kartavyam . \\
\noindent \textbf{(P\_1,2.41)} kim prayojanam . \\
\noindent \textbf{(P\_1,2.41)} aṇiñoḥ luki grahaṇam na kartavyam bhavati : ṇyakṣatriyārṣañitaḥ yūni luk aṇiñoḥ iti apṛktasya iti eva siddham . \\
\noindent \textbf{(P\_1,2.41)} aṇiñoḥ lugartham iti cet ṇe atiprasaṅgaḥ . \\
\noindent \textbf{(P\_1,2.41)} aṇiñoḥ lugartham iti cet ṇe atiprasaṅgaḥ bhavati . \\
\noindent \textbf{(P\_1,2.41)} iha api prāpnoti : phāṇṭāhṛteḥ apatyam māṇavakaḥ pḥāṇṭāhṛtaḥ iti . \\
\noindent \textbf{(P\_1,2.41)} ṇavacanasāmārthyāt na bhaviṣyati . \\
\noindent \textbf{(P\_1,2.41)} vacanaprāmāṇyāt iti cet phagnivṛttyartham vacanam . \\
\noindent \textbf{(P\_1,2.41)} vacanaprāmāṇyāt iti cet phagnivṛttyartham etat syāt : phak ataḥ mā bhūt iti . \\
\noindent \textbf{(P\_1,2.41)} pailādiṣu vacanāt siddham . \\
\noindent \textbf{(P\_1,2.41)} yadi etāvat prayojanam syāt pailādiṣu eva pāṭham kurvīta . \\
\noindent \textbf{(P\_1,2.41)} tatra pāṭhāt anyeṣām api phakaḥ nivṛttiḥ bhavati . \\
\noindent \textbf{(P\_1,2.41)} evam siddhe sati yat ayam ṇam śāsti tat jñāpayati ācāryaḥ na asya luk bhavati iti . \\
\noindent \textbf{(P\_1,2.41)} tāni etāni trīṇi grahaṇāni bhavanti . \\
\noindent \textbf{(P\_1,2.41)} apṛktasañjñāyām halgrahaṇam kartavyam . \\
\noindent \textbf{(P\_1,2.41)} svādilope halaḥ grahaṇam na kartavyam . \\
\noindent \textbf{(P\_1,2.41)} aṇiñoḥ luki grahaṇam kartavyam . \\
\noindent \textbf{(P\_1,2.41)} algrahaṇe api vai kriyamāṇe tāni eva trīṇi grahaṇāni bhavanti . \\
\noindent \textbf{(P\_1,2.41)} apṛktasañjñāyām algrahaṇam kartavyam . \\
\noindent \textbf{(P\_1,2.41)} svādilope halaḥ grahaṇam kartavyam . \\
\noindent \textbf{(P\_1,2.41)} aṇiñoḥ luki grahaṇam na kartavyam bhavati . \\
\noindent \textbf{(P\_1,2.41)} apṛktagrahaṇam kartavyam . \\
\noindent \textbf{(P\_1,2.41)} tatra na asti lāghavakṛtaḥ viśeṣaḥ . \\
\noindent \textbf{(P\_1,2.41)} ayam asti viśeṣaḥ : algrahaṇe kriyamāṇe ekagrahaṇam na kariṣyate . \\
\noindent \textbf{(P\_1,2.41)} kasmāt na bhavati darviḥ , jāgṛviḥ . \\
\noindent \textbf{(P\_1,2.41)} al eva yaḥ pratyayaḥ . \\
\noindent \textbf{(P\_1,2.41)} kim vaktavyam etat . \\
\noindent \textbf{(P\_1,2.41)} na hi . \\
\noindent \textbf{(P\_1,2.41)} katham anucyamānam gaṃsyate . \\
\noindent \textbf{(P\_1,2.41)} algrahaṇasāmarthyāt . \\
\noindent \textbf{(P\_1,2.41)} yadi yaḥ al ca anyaḥ ca tatra syāt algrahaṇam anarthakam syāt . \\
\noindent \textbf{(P\_1,2.41)} halgrahaṇe api kriyamāṇe ekagrahaṇam na kariṣyate . \\
\noindent \textbf{(P\_1,2.41)} kasmāt na bhavati darviḥ jāgṛviḥ . \\
\noindent \textbf{(P\_1,2.41)} hal eva yaḥ pratyayaḥ . \\
\noindent \textbf{(P\_1,2.41)} kim vaktavyam etat . \\
\noindent \textbf{(P\_1,2.41)} na hi . \\
\noindent \textbf{(P\_1,2.41)} katham anucyamānam gaṃsyate . \\
\noindent \textbf{(P\_1,2.41)} halgrahaṇasāmarthyāt . \\
\noindent \textbf{(P\_1,2.41)} yadi yaḥ hal ca anyaḥ ca tatra syāt halgrahaṇam anarthakam syāt . \\
\noindent \textbf{(P\_1,2.41)} asti anyat halgrahaṇasya prayojanam . \\
\noindent \textbf{(P\_1,2.41)} kim . \\
\noindent \textbf{(P\_1,2.41)} halantasya yathā syāt alantasya mā bhūt iti . \\
\noindent \textbf{(P\_1,2.41)} evam tarhi siddhe sati yat algrahaṇe kriyamāṇe ekagrahaṇam karoti tat jñāpayati ācāryaḥ anyatra varṇagrahaṇe jātigrahaṇam bhavati iti . \\
\noindent \textbf{(P\_1,2.41)} kim etasya jñāpane prayojanam . \\
\noindent \textbf{(P\_1,2.41)} dambheḥ halgrahaṇasya jātivācakatvāt siddham iti uktam . \\
\noindent \textbf{(P\_1,2.41)} tat upapannam bhavati . \\
\noindent \textbf{(P\_1,2.42)} tatpuruṣaḥ samānādhikaraṇaḥ karmadhārayaḥ iti cet samāsaikārthatvāt aprasiddhiḥ . \\
\noindent \textbf{(P\_1,2.42)} tatpuruṣaḥ samānādhikaraṇaḥ karmadhārayaḥ iti cet samāsasya ekārthatvāt sañjñāyāḥ aprasiddhiḥ . \\
\noindent \textbf{(P\_1,2.42)} ekaḥ ayam arthaḥ tatpuruṣaḥ nāma anekārthāśrayam ca sāmānādhikaraṇyam . \\
\noindent \textbf{(P\_1,2.42)} siddham tu padasāmānādhikaraṇyāt . \\
\noindent \textbf{(P\_1,2.42)} siddham etat . \\
\noindent \textbf{(P\_1,2.42)} katham . \\
\noindent \textbf{(P\_1,2.42)} tatpuruṣaḥ samānādhikaraṇapadaḥ karmadhārayasañjñaḥ bhavati iti vaktavyam . \\
\noindent \textbf{(P\_1,2.42)} sidhyati . \\
\noindent \textbf{(P\_1,2.42)} sūtram tarhi bhidyate . \\
\noindent \textbf{(P\_1,2.42)} yathānyāsam eva astu . \\
\noindent \textbf{(P\_1,2.42)} nanu ca uktam tatpuruṣaḥ samānādhikaraṇaḥ karmadhārayaḥ iti cet samāsaikārthatvāt aprasiddhiḥ iti . \\
\noindent \textbf{(P\_1,2.42)} na eṣaḥ doṣaḥ . \\
\noindent \textbf{(P\_1,2.42)} ayam tatpuruṣaḥ asti prāthamakalpikaḥ yasmin aikapadyam aikasvaryam ekavibhaktikatvam ca . \\
\noindent \textbf{(P\_1,2.42)} asti tādarthyāt tācchabdyam ; tatpuruṣārthāni padāni tatpuruṣaḥ iti . \\
\noindent \textbf{(P\_1,2.42)} tat yaḥ tādarthyāt tācchabdyam tasya iha grahaṇam . \\
\noindent \textbf{(P\_1,2.43.1)} prathamānirdiṣṭam samāse upasajanam iti cet anirdeśāt prathamāyāḥ samāse sañjñāprasiddhiḥ .prathamānirdiṣṭam samāse upasajanam iti cet anirdeśāt prathamāyāḥ samāse sañjñāyāḥ aprasiddhiḥ . \\
\noindent \textbf{(P\_1,2.43.1)} na hi kaṣṭādīnām samāse prathamām paśyāmaḥ . \\
\noindent \textbf{(P\_1,2.43.1)} siddham tu samāsavidhāne vacanāt . \\
\noindent \textbf{(P\_1,2.43.1)} siddham etat . \\
\noindent \textbf{(P\_1,2.43.1)} katham . \\
\noindent \textbf{(P\_1,2.43.1)} samāsavidhāne prathamānirdiṣṭam upasarjanasañjñam bhavati iti vaktavyam . \\
\noindent \textbf{(P\_1,2.43.1)} tat tarhi vaktavyam . \\
\noindent \textbf{(P\_1,2.43.1)} na vā tādarthyāt tācchabdyam . \\
\noindent \textbf{(P\_1,2.43.1)} na vā vaktavyam . \\
\noindent \textbf{(P\_1,2.43.1)} kim kāraṇam . \\
\noindent \textbf{(P\_1,2.43.1)} tādarthyāt tācchabdyam bhavati. samāsārtham śāstram samāsaḥ iti . \\
\noindent \textbf{(P\_1,2.43.2)} yasya vidhau prathamānirdeśaḥ tataḥ anyatra api upasarjanasañjñāprasaṅgaḥ . \\
\noindent \textbf{(P\_1,2.43.2)} yasya vidhau prathamānirdeśaḥ kriyate tataḥ anyatra api tasya upasarjanasañjñā prāpnoti : rajñaḥ kumārīm rājakumārīm śritaḥ . \\
\noindent \textbf{(P\_1,2.43.2)} śritādisamāse dvitīyāntam prathamānirdiṣṭam . \\
\noindent \textbf{(P\_1,2.43.2)} tasya ṣaṣṭhīsamāse api upasarjanasañjñā prāpnoti . \\
\noindent \textbf{(P\_1,2.43.2)} siddham tu yasya vidhau tam prati iti vacanāt . \\
\noindent \textbf{(P\_1,2.43.2)} siddham etat . \\
\noindent \textbf{(P\_1,2.43.2)} katham . \\
\noindent \textbf{(P\_1,2.43.2)} yasya vidhau yat prathamānirdiṣṭam tam prati tat upasarjanasañjñam bhavati iti vaktavyam . \\
\noindent \textbf{(P\_1,2.43.2)} tat tarhi vaktavyam . \\
\noindent \textbf{(P\_1,2.43.2)} na vaktavyam . \\
\noindent \textbf{(P\_1,2.43.2)} upasarjanam iti mahatī sañjñā kriyate . \\
\noindent \textbf{(P\_1,2.43.2)} sañjñā ca nāma yataḥ na laghīyaḥ . \\
\noindent \textbf{(P\_1,2.43.2)} kutaḥ etat . \\
\noindent \textbf{(P\_1,2.43.2)} laghvartham hi sañjñākaraṇam . \\
\noindent \textbf{(P\_1,2.43.2)} tatra mahatyāḥ sañjñāyāḥ karaṇe etat prayojanam anvarthasañjñā yathā vijñāyeta . \\
\noindent \textbf{(P\_1,2.43.2)} apradhānam upasarjanam iti . \\
\noindent \textbf{(P\_1,2.43.2)} pradhānam upasarjanam iti ca sambandhiśabdau etau . \\
\noindent \textbf{(P\_1,2.43.2)} tatra sambandhāt etat gantavyam : yam prati yat apradhānam tam prati tat upasarjanasñjñam bhavati iti . \\
\noindent \textbf{(P\_1,2.43.2)} atha yatra dve ṣaṣṭhyante kasmāt tatra pradhānasya upasarjanasañjñā na bhavati : rājñaḥ puruṣasya rājapuruṣasya iti . \\
\noindent \textbf{(P\_1,2.43.2)} ṣaṣṭhyantayoḥ upasarjanatve uktam . \\
\noindent \textbf{(P\_1,2.43.2)} kim uktam . \\
\noindent \textbf{(P\_1,2.43.2)} ṣaṣṭhyantayoḥ samāse arthābhedāt pradhānasya apūrvanipātaḥ iti . \\
\noindent \textbf{(P\_1,2.43.2)} evam na ca idam akṛtam bhavet upsarjanam pūrvam iti arthaḥ ca abhinnaḥ iti kṛtvā pradhānasya pūrvanipātaḥ na bhaviṣyati . \\
\noindent \textbf{(P\_1,2.43.2)} yadi api tāvat etat upasarjanakāryam parihṛtam idam aparam prāpnoti . \\
\noindent \textbf{(P\_1,2.43.2)} rājñaḥ kumāryāḥ rājakumāryāḥ . \\
\noindent \textbf{(P\_1,2.43.2)} gostriyoḥ upasarjanasya iti hrasvatvam prāpnoti . \\
\noindent \textbf{(P\_1,2.43.2)} uktam vā . \\
\noindent \textbf{(P\_1,2.43.2)} kim uktam . \\
\noindent \textbf{(P\_1,2.43.2)} paravat liṅgam iti śabdaśabdārthau iti . \\
\noindent \textbf{(P\_1,2.43.2)} tatra aupadeśikasya hrasvatvam . \\
\noindent \textbf{(P\_1,2.43.2)} ātideśikasya śravaṇam bhaviṣyati . \\
\noindent \textbf{(P\_1,2.44.1)} dvitīyādīnām api anena upasarjanasañjñā prāpnoti . \\
\noindent \textbf{(P\_1,2.44.1)} tatra kaḥ doṣaḥ . \\
\noindent \textbf{(P\_1,2.44.1)} tatra apūrvanipāte iti pratiṣedhaḥ prasajyeta . \\
\noindent \textbf{(P\_1,2.44.1)} na apratiṣedhāt . \\
\noindent \textbf{(P\_1,2.44.1)} na ayam prasajyapratiṣedhaḥ : pūrvanipāte na iti . \\
\noindent \textbf{(P\_1,2.44.1)} kim tarhi . \\
\noindent \textbf{(P\_1,2.44.1)} paryudāsaḥ ayam : yat anyat pūrvanipātāt iti . \\
\noindent \textbf{(P\_1,2.44.1)} pūrvanipāte avyāpāraḥ . \\
\noindent \textbf{(P\_1,2.44.1)} yadi kena cit prāpnoti tena bhaviṣyati . \\
\noindent \textbf{(P\_1,2.44.1)} pūrveṇa ca prāpnoti . \\
\noindent \textbf{(P\_1,2.44.1)} tena bhaviṣyati . \\
\noindent \textbf{(P\_1,2.44.1)} aprāpteḥ vā . \\
\noindent \textbf{(P\_1,2.44.1)} atha vā anantarā yā prāptiḥ sā pratiṣidhyate . \\
\noindent \textbf{(P\_1,2.44.1)} kutaḥ etat . \\
\noindent \textbf{(P\_1,2.44.1)} anantarasya vidhiḥ vā bhavati pratiṣedhaḥ vā iti \\
\noindent \textbf{(P\_1,2.44.2)} ekavibhaktau aṣaṣṭhyantavacanam . \\
\noindent \textbf{(P\_1,2.44.2)} ekavibhaktau aṣaṣṭhyantānām iti vaktavyam . \\
\noindent \textbf{(P\_1,2.44.2)} iha mā bhūt : ardham pippalyāḥ ardhapippalī iti . \\
\noindent \textbf{(P\_1,2.44.2)} uktam vā . \\
\noindent \textbf{(P\_1,2.44.2)} kim uktam . \\
\noindent \textbf{(P\_1,2.44.2)} paravat liṅgam iti śabdaśardārthau iti . \\
\noindent \textbf{(P\_1,2.44.2)} tatra aupadeśikasya hrasvatvam . \\
\noindent \textbf{(P\_1,2.44.2)} ātideśikasya śravaṇam bhaviṣyati . \\
\noindent \textbf{(P\_1,2.44.3)} kāni punaḥ asya yogasya prayojanāni . \\
\noindent \textbf{(P\_1,2.44.3)} prayojanam dviguprāptāpannālampūrvopasargāḥ ktārthe . \\
\noindent \textbf{(P\_1,2.44.3)} dviguḥ : pañcabhiḥ gobhiḥ krītaḥ pañcaguḥ . \\
\noindent \textbf{(P\_1,2.44.3)} prāptāpanna : prāptaḥ jivikām prāptajīvikaḥ . \\
\noindent \textbf{(P\_1,2.44.3)} āpannaḥ jīvikām āpannajīvikaḥ . \\
\noindent \textbf{(P\_1,2.44.3)} alampūrva : alam kumāryai alaṅkumāriḥ . \\
\noindent \textbf{(P\_1,2.44.3)} upasargāḥ ktārthe : niṣkauśāmbiḥ nirvārāṇasiḥ . \\
\noindent \textbf{(P\_1,2.45.1)} arthavat iti vyapadeśāya : varṇānām ca mā bhūt iti . \\
\noindent \textbf{(P\_1,2.45.1)} kim ca syāt . \\
\noindent \textbf{(P\_1,2.45.1)} vanam , dhanam iti nalopaḥ prātipadikāntasya iti nalopaḥ prasajyeta . \\
\noindent \textbf{(P\_1,2.45.1)} adhātuḥ iti kimartham . \\
\noindent \textbf{(P\_1,2.45.1)} ahan vṛtram iti . \\
\noindent \textbf{(P\_1,2.45.1)} adhātuḥ iti śakyam akartum . \\
\noindent \textbf{(P\_1,2.45.1)} kasmāt na bhavati ahan vṛtram iti . \\
\noindent \textbf{(P\_1,2.45.1)} ācāryapravṛttiḥ jñāpayati na dhātoḥ prātipadikasañjñā bhavati iti yat ayam supaḥ dhātuprātipadikayoḥ iti dhātugrahaṇam karoti . \\
\noindent \textbf{(P\_1,2.45.1)} na etat asti jñāpakam . \\
\noindent \textbf{(P\_1,2.45.1)} pratiṣiddhārtham etat syāt : api kākaḥ śyenāyate iti . \\
\noindent \textbf{(P\_1,2.45.1)} apratyayaḥ iti kimartham . \\
\noindent \textbf{(P\_1,2.45.1)} kāṇḍe kuḍye . \\
\noindent \textbf{(P\_1,2.45.1)} apratyayaḥ iti śakyam akartum . \\
\noindent \textbf{(P\_1,2.45.1)} kasmāt na bhavati kāṇḍe kuḍye* iti . \\
\noindent \textbf{(P\_1,2.45.1)} kṛttaddhitagrahaṇam niyamārtham bhaviṣyati : kṛttaddhitāntasya eva pratyayāntasya prātipadikasañjñā bhavati na anyasya iti . \\
\noindent \textbf{(P\_1,2.45.2)} arthavati anekapadaprasaṅgaḥ . \\
\noindent \textbf{(P\_1,2.45.2)} arthavati prātipadikasañjñāyām anekasya api padasya prātipadikasañjñā prāpnoti : daśa dāḍimāni ṣaṭ apūpāḥ kuṇḍam ajājinam palalapiṇḍaḥ adhorukam etat kumāryāḥ sphaiyakṛtasya pitā pratiśīnaḥ iti . \\
\noindent \textbf{(P\_1,2.45.2)} samudāyaḥ anarthakaḥ . \\
\noindent \textbf{(P\_1,2.45.2)} samudāyaḥ anarthakaḥ iti cet avayavārthavattvāt samudāyārthavattvam yathā loke . \\
\noindent \textbf{(P\_1,2.45.2)} samudāyaḥ anarthakaḥ iti cet avayavaiḥ arthavadbhiḥ samudāyāḥ api arthavantaḥ bhavanti yathā loke . \\
\noindent \textbf{(P\_1,2.45.2)} tat yathā loke āḍhyam idam nagaram , gomat idam nagaram iti ucyate na ca tatra sarve āḍhyāḥ bhavanti sarve vā gomantaḥ . \\
\noindent \textbf{(P\_1,2.45.2)} yathā loke iti ucyate loke ca avayavāḥ eva arthavantaḥ na samudāyaḥ . \\
\noindent \textbf{(P\_1,2.45.2)} ātaḥ ca avayavāḥ eva arthavantaḥ na samudāyaḥ . \\
\noindent \textbf{(P\_1,2.45.2)} yasya hi tat dravyam bhavati saḥ tena kāryam karoti yasya ca gāvaḥ santi saḥ tāsām kṣīram ghṛtam ca upabhuṅkte . \\
\noindent \textbf{(P\_1,2.45.2)} anyaiḥ etat draṣṭum api aśakyam . \\
\noindent \textbf{(P\_1,2.45.2)} kā tarhi iyam vācoyuktiḥ : āḍhyam idam nagaram , gomat idam iti . \\
\noindent \textbf{(P\_1,2.45.2)} eṣā eṣā vācoyuktiḥ : iha tāvat āḍhyam idam nagaram iti akāraḥ matvarthīyaḥ : āḍhyāḥ asmin santi iti tat idam āḍhyam iti . \\
\noindent \textbf{(P\_1,2.45.2)} gomat idam iti matvantāt matvarthīyaḥ lupyate . \\
\noindent \textbf{(P\_1,2.45.2)} evam api vākyapratiṣedhaḥ arthavattvāt . \\
\noindent \textbf{(P\_1,2.45.2)} vākyasya prātipadikasñjñāyāḥ pratiṣedhaḥ vaktavyaḥ : devadatta gām abhyāja śuklām . \\
\noindent \textbf{(P\_1,2.45.2)} devadatta gām abhyāja kṛṣṇām iti . \\
\noindent \textbf{(P\_1,2.45.2)} kim kāraṇam . \\
\noindent \textbf{(P\_1,2.45.2)} arthavattvāt . \\
\noindent \textbf{(P\_1,2.45.2)} arthavat hi etat vākyam bhavati . \\
\noindent \textbf{(P\_1,2.45.2)} na vai padārthāt anyasya arthasya upalabdhiḥ bhavati vākye . \\
\noindent \textbf{(P\_1,2.45.2)} padārthāt anyasya anupalabdhiḥ iti cet padārthābhisambandhasya upalabdhiḥ . \\
\noindent \textbf{(P\_1,2.45.2)} padārthāt anyasya anupalabdhiḥ iti cet evam ucyate : padārthābhisambandhasya upalabdhiḥ bhavati vākye . \\
\noindent \textbf{(P\_1,2.45.2)} iha devadatta iti ukte kartā nirdiṣṭaḥ karma kriyāguṇau ca anirdiṣṭau . \\
\noindent \textbf{(P\_1,2.45.2)} gām iti ukte karma nirdiṣṭam kartā kriyāguṇau ca anirdiṣṭau . \\
\noindent \textbf{(P\_1,2.45.2)} abhyāja iti ukte kriyā nirdiṣṭā kartṛkarmaṇī guṇaḥ ca anirdiṣṭaḥ . \\
\noindent \textbf{(P\_1,2.45.2)} śuklām iti ukte guṇaḥ nirdiṣṭaḥ kartṛkarmaṇī kriyā ca anirdiṣṭā . \\
\noindent \textbf{(P\_1,2.45.2)} iha idānīm devadatta gām abhyāja śuklām iti ukte sarvam nirdiṣṭam bhavati : devadattaḥ eva kartā na anyaḥ . \\
\noindent \textbf{(P\_1,2.45.2)} gauḥ eva karma na anyat . \\
\noindent \textbf{(P\_1,2.45.2)} abhyājiḥ eva kriyā na anyā . \\
\noindent \textbf{(P\_1,2.45.2)} śuklām eva na kṛṣṇām iti . \\
\noindent \textbf{(P\_1,2.45.2)} eteṣām padānām sāmānye vartamānānām yadviśeṣe avasthānam saḥ vākyārthaḥ . \\
\noindent \textbf{(P\_1,2.45.2)} tasmāt pratiṣedhaḥ . \\
\noindent \textbf{(P\_1,2.45.2)} tasmāt pratiṣedhaḥ vaktavyaḥ . \\
\noindent \textbf{(P\_1,2.45.2)} na vaktavyaḥ . \\
\noindent \textbf{(P\_1,2.45.2)} arthavatsamudāyānām samāsagrahaṇam niyamārtham . \\
\noindent \textbf{(P\_1,2.45.2)} arthavatsamudāyānām samāsagrahaṇam niyamārtham bhaviṣyati : samāsaḥ eva arthavatām samudāyāyānām prātipadikasañjñaḥ bhavati na anyaḥ iti . \\
\noindent \textbf{(P\_1,2.45.2)} yadi niyamaḥ kriyate prakṛtipratyayasamudāyasya prātipadikasañjñā na prāpnoti : bahupaṭavaḥ , uccakaiḥ iti . \\
\noindent \textbf{(P\_1,2.45.2)} kim punaḥ atra prātipadikasañjñayā prārthyate . \\
\noindent \textbf{(P\_1,2.45.2)} prātipadikāt iti svādyutpattiḥ yathā syāt . \\
\noindent \textbf{(P\_1,2.45.2)} na eṣaḥ doṣaḥ . \\
\noindent \textbf{(P\_1,2.45.2)} yathā eva atra aprātipadikatvāt svādyutpattiḥ na bhavati evam luk api na bhaviṣyati . \\
\noindent \textbf{(P\_1,2.45.2)} tatra yā eva antarvartinī vibhaktiḥ tasyāḥ eva śravaṇam bhaviṣyati . \\
\noindent \textbf{(P\_1,2.45.2)} na evam śakyam . \\
\noindent \textbf{(P\_1,2.45.2)} svare hi doṣaḥ syāt . \\
\noindent \textbf{(P\_1,2.45.2)} bahupaṭavaḥ iti evam svaraḥ prasajyeta bahupaṭavaḥ iti ca iṣyate . \\
\noindent \textbf{(P\_1,2.45.2)} paṭhiṣyati hi ācāryaḥ : citaḥ saprakṛteḥ bahvakajartham iti . \\
\noindent \textbf{(P\_1,2.45.2)} tasyām punaḥ luptāyam yā anyā vibhaktiḥ utpadyate tasyāḥ prakṛtyanekadeśatvāt antodāttatvam na bhaviṣyati . \\
\noindent \textbf{(P\_1,2.45.2)} evam tarhi ācāryapravṛttiḥ jñāpayati bhavati prakṛtipratyayasamudāyasya prātipadikasañjñā iti yat ayam aprayayaḥ iti pratiṣedham śāsti . \\
\noindent \textbf{(P\_1,2.45.2)} saḥ ca tadantapratiṣedhaḥ . \\
\noindent \textbf{(P\_1,2.45.2)} saḥ tarhi jñāpakārthaḥ pratyayapratiṣedhaḥ vaktavyaḥ . \\
\noindent \textbf{(P\_1,2.45.2)} nanu ca ayam prāptyarthaḥ api vaktavyaḥ . \\
\noindent \textbf{(P\_1,2.45.2)} na arthaḥ prāptyarthena . \\
\noindent \textbf{(P\_1,2.45.2)} kṛttaddhitagrahaṇam niyamārtham bhaviṣyati : kṛttaddhitāntasya eva pratyayāntasya prātipadikasañjñā bhaviṣyati na anyasya pratyayāntasya iti . \\
\noindent \textbf{(P\_1,2.45.2)} saḥ eṣaḥ ananyārthaḥ pratyayapratiṣedhaḥ vaktavyaḥ prakṛtipratyayasamudāyasya vā prātipadikasañjñā vaktavyā . \\
\noindent \textbf{(P\_1,2.45.2)} ubhayam na vaktavyam . \\
\noindent \textbf{(P\_1,2.45.2)} tulyajātīyasya niyamaḥ . \\
\noindent \textbf{(P\_1,2.45.2)} kaḥ ca tulyajātīyaḥ . \\
\noindent \textbf{(P\_1,2.45.2)} yathājātīyakānām samāsaḥ . \\
\noindent \textbf{(P\_1,2.45.2)} kathañjātīyakānām samāsaḥ . \\
\noindent \textbf{(P\_1,2.45.2)} subantānām . \\
\noindent \textbf{(P\_1,2.45.2)} suptiṅsamudāyasya tarhi prātipadikasañjñā prāpnoti . \\
\noindent \textbf{(P\_1,2.45.2)} suptiṅsamudāyasya prātipadikasañjñā ārabhyate : jahi karmaṇā bahulam ābhīkṣṇye kartāram ca abhidadhāti iti . \\
\noindent \textbf{(P\_1,2.45.2)} tat niyamārtham bhaviṣyati : etasya eva suptiṅsamudāyasya prātipadikasañjñā bhavati na anyasya iti . \\
\noindent \textbf{(P\_1,2.45.2)} tiṅsamudāyasya tarhi prātipadikasañjñā prāpnoti . \\
\noindent \textbf{(P\_1,2.45.2)} tiṅsamudāyasya api prātipadikasañjñā ārabhyate : ākhyātam ākhyātena kriyāsātatye iti . \\
\noindent \textbf{(P\_1,2.45.2)} tat niyamārtham bhaviṣyati : etasya eva tiṅsamudāyasya prātipadikasañjñā bhavati na anyasya iti . \\
\noindent \textbf{(P\_1,2.45.3)} arthavattā na upapadyate kevalena avacanāt . \\
\noindent \textbf{(P\_1,2.45.3)} arthavattā na upapadyate vṛkṣaśabdasya . \\
\noindent \textbf{(P\_1,2.45.3)} kim kāraṇam . \\
\noindent \textbf{(P\_1,2.45.3)} kevalena avacanāt . \\
\noindent \textbf{(P\_1,2.45.3)} na kevalena vṛkṣaśabdena arthaḥ gamyate . \\
\noindent \textbf{(P\_1,2.45.3)} kena tarhi . \\
\noindent \textbf{(P\_1,2.45.3)} sapratyayakena . \\
\noindent \textbf{(P\_1,2.45.3)} na vā pratyayena nityasambandhāt kevalasya aprayogaḥ . \\
\noindent \textbf{(P\_1,2.45.3)} na vā eṣaḥ doṣaḥ . \\
\noindent \textbf{(P\_1,2.45.3)} kim kāraṇam . \\
\noindent \textbf{(P\_1,2.45.3)} pratyayena nityasambandhāt . \\
\noindent \textbf{(P\_1,2.45.3)} nityasambandhau etau arthau prakṛtiḥ pratyayaḥ iti . \\
\noindent \textbf{(P\_1,2.45.3)} pratyayena nityasambandhāt kevalasya aprayogaḥ na bhaviṣyati . \\
\noindent \textbf{(P\_1,2.45.3)} anyat bhavān pṛṣṭaḥ anyat ācaṣṭe . \\
\noindent \textbf{(P\_1,2.45.3)} āmrān pṛṣṭaḥ kovidārān ācaṣṭe . \\
\noindent \textbf{(P\_1,2.45.3)} arthavattā na upapadyate kevalena avacanāt iti bhavān asmābhiḥ coditaḥ kevalasya aprayoge hetum āha . \\
\noindent \textbf{(P\_1,2.45.3)} evam ca kila nāma kṛtvā codyate : samudāyasya arthe prayogāt avayavānām aprasiddhiḥ iti . \\
\noindent \textbf{(P\_1,2.45.3)} siddham tu anvayavyatirekābhyām . \\
\noindent \textbf{(P\_1,2.45.3)} siddham etat . \\
\noindent \textbf{(P\_1,2.45.3)} katham . \\
\noindent \textbf{(P\_1,2.45.3)} anvayāt vyatirekāt ca . \\
\noindent \textbf{(P\_1,2.45.3)} kaḥ asau anvayaḥ vyatirekaḥ vā . \\
\noindent \textbf{(P\_1,2.45.3)} iha vṛkṣaḥ iti ukte kaḥ cit śabdaḥ śrūyate : vṛkṣaśabdaḥ akārāntaḥ sakārāntaḥ ca pratyayaḥ . \\
\noindent \textbf{(P\_1,2.45.3)} arthaḥ api kaḥ cit gamyate : mūlaskandhaphalapalāśavān ekatvam ca . \\
\noindent \textbf{(P\_1,2.45.3)} vṛkṣau iti ukte kaḥ cit śabdaḥ hīyate kaḥ cit upajāyate kaḥ cit anvayī : sakāraḥ hīyate , aukāraḥ upajāyate vṛkṣaśabdaḥ akārāntaḥ anvayī . \\
\noindent \textbf{(P\_1,2.45.3)} arthaḥ api kaḥ cit hīyate kaḥ cit upajāyate kaḥ cit anvayī : ekatvam hīyate dvitvam upajāyate mūlaskandhaphalapalāśavān anvayī . \\
\noindent \textbf{(P\_1,2.45.3)} te manyāmahe : yaḥ śabdaḥ hīyate tasya asau arthaḥ yaḥ arthaḥ hīyate . \\
\noindent \textbf{(P\_1,2.45.3)} yaḥ śabdaḥ upajāyate tasya asau arthaḥ yaḥ arthaḥ upajāyate . \\
\noindent \textbf{(P\_1,2.45.3)} yaḥ śabdaḥ anvayī tasya asau arthaḥ yaḥ arthaḥ anvayī. viṣamaḥ upanyāsaḥ . \\
\noindent \textbf{(P\_1,2.45.3)} bahavaḥ hi śabdāḥ ekārthāḥ bhavanti . \\
\noindent \textbf{(P\_1,2.45.3)} tat yathā : indraḥ śakraḥ puruhūtaḥ purandaraḥ , kanduḥ koṣṭhaḥ kuśūlaḥ iti . \\
\noindent \textbf{(P\_1,2.45.3)} ekaḥ ca śabdaḥ bahvarthaḥ . \\
\noindent \textbf{(P\_1,2.45.3)} tat yathā : akṣāḥ pādāḥ māṣāḥ iti . \\
\noindent \textbf{(P\_1,2.45.3)} ataḥ kim na sādhīyaḥ arthavattā siddhā bhavati . \\
\noindent \textbf{(P\_1,2.45.3)} na brūmaḥ arthavattā na sidhyati iti .varṇitā arthavattā anvayavyatirekābhyām eva . \\
\noindent \textbf{(P\_1,2.45.3)} tatra kutaḥ etat : ayam prakṛtyarthaḥ ayam pratyayārthaḥ iti na punaḥ prakṛtiḥ eva ubhau arthau brūyāt pratyayaḥ eva vā . \\
\noindent \textbf{(P\_1,2.45.3)} sāmānyaśabdāḥ ete evam syuḥ . \\
\noindent \textbf{(P\_1,2.45.3)} sāmānyaśabdāḥ ca na antareṇa viśeṣam prakaraṇam vā viśeṣeṣu avatiṣṭhante . \\
\noindent \textbf{(P\_1,2.45.3)} yataḥ tu niyogataḥ vṛkṣaḥ iti ukte svabhāvataḥ kasmin cid arthe pratītiḥ upajāyate ataḥ manyāmahe na ime sāmānyaśabdāḥ iti . \\
\noindent \textbf{(P\_1,2.45.3)} na cet sāmānyaśabdāḥ prakṛtiḥ prakṛtyarthe vartate pratyayaḥ pratyayārthe . \\
\noindent \textbf{(P\_1,2.45.4)} kim punaḥ ime varṇāḥ arthavantaḥ āhosvit anarthakāḥ . \\
\noindent \textbf{(P\_1,2.45.4)} varṇasya arthavadanarthakatve uktam . \\
\noindent \textbf{(P\_1,2.45.4)} kim uktam . \\
\noindent \textbf{(P\_1,2.45.4)} arthavantaḥ varṇāḥ dhātuprātipadikapratyayanipātānām ekavarṇānām arthadarśanāt . \\
\noindent \textbf{(P\_1,2.45.4)} varṇavyatyaye ca arthāntaragamanāt . \\
\noindent \textbf{(P\_1,2.45.4)} varṇānupalabdhau ca anarthagateḥ . \\
\noindent \textbf{(P\_1,2.45.4)} saṅghātārthavattvāt ca . \\
\noindent \textbf{(P\_1,2.45.4)} saṅghātasya aikārthyāt subabhāvaḥ varṇāt . \\
\noindent \textbf{(P\_1,2.45.4)} anarthakāḥ tu prativarṇam arthānupalabdheḥ . \\
\noindent \textbf{(P\_1,2.45.4)} varṇavyatyayāpāyopajanavikāreṣu arthadarśanāt iti . \\
\noindent \textbf{(P\_1,2.45.4)} tatra idam aparihṛtam : saṅghātārthavattvāt ca iti . \\
\noindent \textbf{(P\_1,2.45.4)} tasya parihāraḥ . \\
\noindent \textbf{(P\_1,2.45.4)} saṅghātārthavattvāt ca iti cet dṛṣṭaḥ hiatadarthena guṇena guṇinaḥ arthabhāvaḥ . \\
\noindent \textbf{(P\_1,2.45.4)} saṅghātārthavattvāt ca iti cet dṛśyate hi punaḥ atadarthena guṇena guṇinaḥ arthabhāvaḥ . \\
\noindent \textbf{(P\_1,2.45.4)} tat yathā . \\
\noindent \textbf{(P\_1,2.45.4)} ekaḥ tantuḥ tvaktrāṇe asamarthaḥ tatsamudāyaḥ ca kambalaḥ samarthaḥ . \\
\noindent \textbf{(P\_1,2.45.4)} ekaḥ taṇḍulaḥ kṣutpratighāte asamarthaḥ tatsamudāyaḥ ca vardhatikam samarthaḥ . \\
\noindent \textbf{(P\_1,2.45.4)} ekaḥ ca balvajaḥ bandhane asamarthaḥ tatsamudāyaḥ ca rajjuḥ samarthā bhavati . \\
\noindent \textbf{(P\_1,2.45.4)} viṣamaḥ upanyāsaḥ . \\
\noindent \textbf{(P\_1,2.45.4)} bhavati hi tatra yā ca yāvatī ca arthamātrā . \\
\noindent \textbf{(P\_1,2.45.4)} bhavati hi kim cit prati ekaḥ tantuḥ tvaktrāṇe samarthaḥ ekaḥ ca taṇḍulaḥ kṣutpratighāte samarthaḥ ekaḥ ca balvajaḥ bandhane samarthaḥ . \\
\noindent \textbf{(P\_1,2.45.4)} ime punaḥ varṇāḥ atyantāya eva anarthakāḥ . \\
\noindent \textbf{(P\_1,2.45.4)} yathā tarhi rathāṅgāni vihṛtāni pratyekam vrajikriyām prati asamarthāni bhavanti tatsamudāyaḥ ca rathaḥ samarthaḥ evam eṣām varṇānām samudāyāḥ arthavantaḥ avayavāḥ anarthakāḥ iti . \\
\noindent \textbf{(P\_1,2.45.5)} nipātasya anarthakasya prātipadikatvam . \\
\noindent \textbf{(P\_1,2.45.5)} nipātasya anarthakasya prātipadikasañjñā vaktavyā . \\
\noindent \textbf{(P\_1,2.45.5)} khañjati nikhañjati lambate pralambate . \\
\noindent \textbf{(P\_1,2.45.5)} kim punaḥ atra prātipadikasañjñayā prārthyate . \\
\noindent \textbf{(P\_1,2.45.5)} prātipadikāt iti svādyutpattiḥ , subantam padam iti padasañjñā , padasya padāt iti nighātaḥ yathā syāt . \\
\noindent \textbf{(P\_1,2.45.5)} na etat asti prayojanam . \\
\noindent \textbf{(P\_1,2.45.5)} satyām api prātipadikasañjñāyām svādyutpattiḥ na prāpnoti . \\
\noindent \textbf{(P\_1,2.45.5)} kim kāraṇam . \\
\noindent \textbf{(P\_1,2.45.5)} na hi prātipadikasañjñāyām eva svādyutpattiḥ pratibaddhā . \\
\noindent \textbf{(P\_1,2.45.5)} kim tarhi . \\
\noindent \textbf{(P\_1,2.45.5)} ekatvādiṣu artheṣu svādayaḥ vidhīyante na ca eṣām ekatvādayaḥ santi . \\
\noindent \textbf{(P\_1,2.45.5)} na eṣaḥ doṣaḥ . \\
\noindent \textbf{(P\_1,2.45.5)} aviśeṣeṇa utpadyante . \\
\noindent \textbf{(P\_1,2.45.5)} utpannānām niyamaḥ kriyate . \\
\noindent \textbf{(P\_1,2.45.5)} atha vā prakṛtārthān apekṣya niyamaḥ . \\
\noindent \textbf{(P\_1,2.45.5)} ke ca prakṛtāḥ . \\
\noindent \textbf{(P\_1,2.45.5)} ekatvādayaḥ . \\
\noindent \textbf{(P\_1,2.45.5)} ekasmin eva arthe ekavacanam na dvayoḥ na bahuṣu . \\
\noindent \textbf{(P\_1,2.45.5)} dvayoḥ eva dvivacanam na ekasmin na bahuṣu . \\
\noindent \textbf{(P\_1,2.45.5)} bahuṣu eva artheṣu bahuvacanam na ekasmin na dvayoḥ iti . \\
\noindent \textbf{(P\_1,2.45.5)} atha vā ācāryapravṛttiḥ jñāpayati anarthakānām api eteṣām bhavati arthavatkṛtam iti yat ayam adhiparī* anarthakau iti anarthakayoḥ gatyupasargasañjābādhikām karmapravacanīyasañjñām śāsti \\
\noindent \textbf{(P\_1,2.45.6)} kim punaḥ ayam paryudāsaḥ : yat anyat pratyayāt āhosvit prasajya ayam pratiṣedhaḥ : pratyayaḥ na iti . \\
\noindent \textbf{(P\_1,2.45.6)} kaḥ ca atra viśeṣaḥ . \\
\noindent \textbf{(P\_1,2.45.6)} apratyayaḥ iti cet tibekādeśe pratiṣedhaḥ antavattvāt . \\
\noindent \textbf{(P\_1,2.45.6)} apratyayaḥ iti cet tibekādeśe pratiṣedhaḥ vaktavyaḥ : kāṇḍe kuḍye . \\
\noindent \textbf{(P\_1,2.45.6)} kim kāraṇam . \\
\noindent \textbf{(P\_1,2.45.6)} antavattvāt . \\
\noindent \textbf{(P\_1,2.45.6)} tibatipoḥ ekādeśaḥ atipaḥ antavat syāt . \\
\noindent \textbf{(P\_1,2.45.6)} asti anyat tipaḥ iti kṛtvā prātipadikasañjñā prāpnoti . \\
\noindent \textbf{(P\_1,2.45.6)} astu tarhi prasajyapratiṣedhaḥ : pratyayaḥ na iti . \\
\noindent \textbf{(P\_1,2.45.6)} na pratyayaḥ iti cet ūṅekādeśe pratiṣedhaḥ ādivattvāt . \\
\noindent \textbf{(P\_1,2.45.6)} na pratyayaḥ iti cet ūṅekādeśe pratiṣedhaḥ prāpnoti : brahmabandhūḥ . \\
\noindent \textbf{(P\_1,2.45.6)} kim kāraṇam . \\
\noindent \textbf{(P\_1,2.45.6)} ādivattvāt . \\
\noindent \textbf{(P\_1,2.45.6)} pratyayāpratyayayoḥ pratyayasya ādivat syāt . \\
\noindent \textbf{(P\_1,2.45.6)} tatra pratyayaḥ na iti pratiṣedhaḥ prāpnoti . \\
\noindent \textbf{(P\_1,2.45.6)} na eṣaḥ doṣaḥ . \\
\noindent \textbf{(P\_1,2.45.6)} ācāryapravṛttiḥ jñāpayati utpadyante ūṅantāt svādayaḥ iti yat ayam na ūṅdhātvoḥ iti vibhaktisvarasya pratiṣedham śāśti . \\
\noindent \textbf{(P\_1,2.45.6)} atha vā dve hi atra prātipadikasañjñe : avayavasya api samudāyasya api . \\
\noindent \textbf{(P\_1,2.45.6)} tatra avayavasya yā prātipadikasañjñā tayā antavadbhāvāt svādyutpattiḥ bhaviṣyati . \\
\noindent \textbf{(P\_1,2.45.6)} sublope ca pratyayalakṣaṇatvāt . \\
\noindent \textbf{(P\_1,2.45.6)} sublope ca pratyayalakṣaṇena pratiṣedhaḥ prāpnoti : rājā takṣā . \\
\noindent \textbf{(P\_1,2.45.6)} pratyayalakṣaṇena pratyayaḥ na iti pratiṣedhaḥ prāpnoti . \\
\noindent \textbf{(P\_1,2.45.6)} na eṣaḥ doṣaḥ . \\
\noindent \textbf{(P\_1,2.45.6)} ācāryapravṛttiḥ jñāpayati na pratyayalakṣaṇena pratiṣedhaḥ bhavati iti yat ayam na ṅisambuddhyoḥ iti pratiṣedham śāsti . \\
\noindent \textbf{(P\_1,2.45.6)} atha vā punaḥ astu paryudāsaḥ . \\
\noindent \textbf{(P\_1,2.45.6)} nanu ca uktam : apratyayaḥ iti cet tibekādeśe pratiṣedhaḥ antavattvāt iti . \\
\noindent \textbf{(P\_1,2.45.6)} prasajyapratiṣedhe api eṣaḥ doṣaḥ . \\
\noindent \textbf{(P\_1,2.45.6)} dve hi atra prātipadikasañjñe : avayavasya api samudāyasya api . \\
\noindent \textbf{(P\_1,2.45.6)} gṛhyate ca prātipadikāprātipadikayoḥ ekādeśaḥ prātipadikagrahaṇena . \\
\noindent \textbf{(P\_1,2.45.6)} tasmāt ubhābhyām api vaktavyam syāt : hrasvaḥ napuṃsake yat tasya iti . \\
\noindent \textbf{(P\_1,2.45.6)} kim ca napuṃsake . \\
\noindent \textbf{(P\_1,2.45.6)} napuṃsakam yasya guṇaḥ . \\
\noindent \textbf{(P\_1,2.45.6)} kasya ca napuṃsakam guṇaḥ . \\
\noindent \textbf{(P\_1,2.45.6)} prātipadikasya . \\
\noindent \textbf{(P\_1,2.46)} samāsagrahaṇam kimartham . \\
\noindent \textbf{(P\_1,2.46)} samāsagrahaṇe uktam . \\
\noindent \textbf{(P\_1,2.46)} kim uktam . \\
\noindent \textbf{(P\_1,2.46)} arthavatsamudāyānām samāsagrahaṇam niyamārtham iti . \\
\noindent \textbf{(P\_1,2.47.1)} prātipadikagrahaṇam kimartham . \\
\noindent \textbf{(P\_1,2.47.1)} napuṃsakahrasvatve prātipadikagrahaṇam tibnivṛttyartham . napuṃsakahrasvatve prātipadikagrahaṇam kriyate tibnivṛttyartham . \\
\noindent \textbf{(P\_1,2.47.1)} tibantasya hrasvatvam mā bhūt : kāṇḍe kuḍye . \\
\noindent \textbf{(P\_1,2.47.1)} ramate brāhmaṇakulam . \\
\noindent \textbf{(P\_1,2.47.1)} avyayapratiṣedhaḥ . \\
\noindent \textbf{(P\_1,2.47.1)} avyayānām pratiṣedhaḥ vaktavyaḥ : doṣā brāhmaṇakulam divā brāhmaṇakulam iti . \\
\noindent \textbf{(P\_1,2.47.1)} saḥ tarhi vaktavyaḥ . \\
\noindent \textbf{(P\_1,2.47.1)} na vaktavyaḥ . \\
\noindent \textbf{(P\_1,2.47.1)} na atra avyayam napuṃsake vartate . \\
\noindent \textbf{(P\_1,2.47.1)} kim tarhi . \\
\noindent \textbf{(P\_1,2.47.1)} adhikaraṇam atra avyayam napuṃsakasya . \\
\noindent \textbf{(P\_1,2.47.1)} iha tarhi prāpnoti : kāṇḍībhūtam vṛṣalakulam , kuḍyībhūtam vṛṣalakulam iti . \\
\noindent \textbf{(P\_1,2.47.1)} na vā liṅgābhāvāt . \\
\noindent \textbf{(P\_1,2.47.1)} na vā vaktavyam . \\
\noindent \textbf{(P\_1,2.47.1)} kim kāraṇam . \\
\noindent \textbf{(P\_1,2.47.1)} liṅgābhāvāt . \\
\noindent \textbf{(P\_1,2.47.1)} aliṅgam avyayam . \\
\noindent \textbf{(P\_1,2.47.1)} kim punaḥ ayam avyayasya eva parihāraḥ āhosvit tibantasya api . \\
\noindent \textbf{(P\_1,2.47.1)} tibantasya api iti āha . \\
\noindent \textbf{(P\_1,2.47.1)} katham . \\
\noindent \textbf{(P\_1,2.47.1)} avyayam hi kim cit vibhaktyarthapradhānam kim cit kriyāpradhānam . \\
\noindent \textbf{(P\_1,2.47.1)} uccaiḥ, nīcaiḥ iti vibhaktyarthapradhānam , hiruk pṛthak iti kriyāpradhānam . \\
\noindent \textbf{(P\_1,2.47.1)} tibantam ca api kim cit vibhaktyarthapradhānam kim cit kriyāpradhānam . \\
\noindent \textbf{(P\_1,2.47.1)} kāṇḍe kuḍye* iti vibhaktarthyapradhānam , ramate brāhmaṇakulam iti kriyāpradhānam . \\
\noindent \textbf{(P\_1,2.47.1)} na ca etayoḥ arthayoḥ liṅgasaṅkhyābhyām yogaḥ asti . \\
\noindent \textbf{(P\_1,2.47.1)} avaśyam ca etat evam vijñeyam . \\
\noindent \textbf{(P\_1,2.47.1)} kriyamāṇe api hi prātipadikagrahaṇe iha prasajyeta : kāṇḍe kuḍye . \\
\noindent \textbf{(P\_1,2.47.1)} dve hi atra prātipadikasañjñe avayavasya api samudāyasya api . \\
\noindent \textbf{(P\_1,2.47.1)} gṛhyate ca prātipadikāprātipadikayoḥ ekādeśaḥ prātipadikagrahaṇena . \\
\noindent \textbf{(P\_1,2.47.1)} tasmāt ubhābhyām api vaktavyam syāt : hrasvaḥ napuṃsake yat tasya iti . \\
\noindent \textbf{(P\_1,2.47.1)} kim ca napuṃsake . \\
\noindent \textbf{(P\_1,2.47.1)} napuṃsakam yasya guṇaḥ . \\
\noindent \textbf{(P\_1,2.47.1)} kasya ca napuṃsakam guṇaḥ . \\
\noindent \textbf{(P\_1,2.47.1)} prātipadikasya . \\
\noindent \textbf{(P\_1,2.47.2)} yañekādeśadīrghaittveṣu pratiṣedhaḥ . \\
\noindent \textbf{(P\_1,2.47.2)} yañekādeśadīrghaittveṣu pratiṣedhaḥ vaktavyaḥ : yugavaratrāya yugavaratrārtham , yugavaratrebhyaḥ . \\
\noindent \textbf{(P\_1,2.47.2)} yañekādeśadīrghaittveṣu bahiraṅgalakṣaṇatvāt siddham . \\
\noindent \textbf{(P\_1,2.47.2)} bahiraṅgāḥ ete vidhayaḥ . \\
\noindent \textbf{(P\_1,2.47.2)} antaraṅgam hrasvatvam . \\
\noindent \textbf{(P\_1,2.47.2)} asiddham bahiraṅgam antaraṅge \\
\noindent \textbf{(P\_1,2.48.1)} upasarjanahrasvatve ca . \\
\noindent \textbf{(P\_1,2.48.1)} upasarjanahrasvatve ca . \\
\noindent \textbf{(P\_1,2.48.1)} kim . \\
\noindent \textbf{(P\_1,2.48.1)} yañekādeśadīrghaittveṣu pratiṣedhaḥ vaktavyaḥ : atikhaṭvāya atikhaṭvārtham atikhaṭvebhyaḥ . \\
\noindent \textbf{(P\_1,2.48.1)} upasarjanahrasvatve ca . \\
\noindent \textbf{(P\_1,2.48.1)} kim . \\
\noindent \textbf{(P\_1,2.48.1)} bahiraṅgalakṣaṇatvāt siddham iti eva . \\
\noindent \textbf{(P\_1,2.48.1)} bahiraṅgāḥ ete vidhayaḥ . \\
\noindent \textbf{(P\_1,2.48.1)} antaraṅgam hrasvatvam . \\
\noindent \textbf{(P\_1,2.48.1)} asiddham bahiraṅgam antaraṅge \\
\noindent \textbf{(P\_1,2.48.2)} goṭāṅgrahaṇam kṛnnivṛttyartham . \\
\noindent \textbf{(P\_1,2.48.2)} goṭāṅgrahaṇam kartavyam . \\
\noindent \textbf{(P\_1,2.48.2)} kim idam ṭāṅ iti . \\
\noindent \textbf{(P\_1,2.48.2)} pratyāhāragrahaṇam . \\
\noindent \textbf{(P\_1,2.48.2)} kva sanniviṣṭānām pratyāhāraḥ . \\
\noindent \textbf{(P\_1,2.48.2)} ṭāpaḥ prabhṛti ā ṣyaṅaḥ ṅakārāt . \\
\noindent \textbf{(P\_1,2.48.2)} kim prayojanam . \\
\noindent \textbf{(P\_1,2.48.2)} kṛnnivṛttyartham . \\
\noindent \textbf{(P\_1,2.48.2)} kṛtstriyāḥ dhātustriyāḥ ca hrasvatvam mā bhūt iti : atitantrīḥ , atiśrīḥ , atilakṣmīḥ iti . \\
\noindent \textbf{(P\_1,2.48.2)} tat tarhi vaktavyam . \\
\noindent \textbf{(P\_1,2.48.2)} na vaktavyam . \\
\noindent \textbf{(P\_1,2.48.2)} strīgrahaṇam svaryate . \\
\noindent \textbf{(P\_1,2.48.2)} tatra svaritena adhikāragatiḥ bhavati . \\
\noindent \textbf{(P\_1,2.48.2)} striyām iti evam prakṛtya ye vihitāḥ teṣām grahaṇam vijñāsyate . \\
\noindent \textbf{(P\_1,2.48.2)} svaritena adhikāragatiḥ bhavati iti na doṣaḥ bhavati . \\
\noindent \textbf{(P\_1,2.48.2)} yadi evam pratyayagrahaṇam idam bhavati . \\
\noindent \textbf{(P\_1,2.48.2)} tatra pratyayagrahaṇe yasmāt saḥ tadādeḥ grahaṇam bhavati iti iha na prāpnoti : atirājakumāriḥ , atisenānīkumāriḥ iti . \\
\noindent \textbf{(P\_1,2.48.2)} astrīpratyayena iti evam tat . \\
\noindent \textbf{(P\_1,2.48.2)} īyasaḥ bahuvrīhau puṃvadvacanam . \\
\noindent \textbf{(P\_1,2.48.2)} īyasaḥ bahuvrīhau puṃvadbhāvaḥ vaktavyaḥ . \\
\noindent \textbf{(P\_1,2.48.2)} bahvyaḥ śreyasyaḥ asya bahuśreyasī vidyamānaśreyasī . \\
\noindent \textbf{(P\_1,2.48.2)} pūrvapadasya ca pratiṣedhaḥ gosamāsanivṛttyartham . \\
\noindent \textbf{(P\_1,2.48.2)} pūrvapadasya ca pratiṣedhaḥ vaktavyaḥ . \\
\noindent \textbf{(P\_1,2.48.2)} kim prayojanam . \\
\noindent \textbf{(P\_1,2.48.2)} gosamāsanivṛttyartham . \\
\noindent \textbf{(P\_1,2.48.2)} gonivṛttyartham samāsanivṛttyartham ca . \\
\noindent \textbf{(P\_1,2.48.2)} gonivṛttyartham tāvat : gokulam , gokṣīram , gopālakaḥ iti . \\
\noindent \textbf{(P\_1,2.48.2)} samāsanivṛttyartham : rājakumārīputraḥ , senānīkumārīputraḥ iti . \\
\noindent \textbf{(P\_1,2.48.2)} kim ucyate samāsanivṛttyartham iti na punaḥ asamāsaḥ api kim cit pūrvapadam yadarthaḥ pratiṣedhaḥ syāt . \\
\noindent \textbf{(P\_1,2.48.2)} stryantasya prātipadikasya upasarjanasya hrasvaḥ bhavati iti ucyate na ca antareṇa samāsam stryantam prātipadikam upasarjanam asti . \\
\noindent \textbf{(P\_1,2.48.2)} nanu ca idam asti : khaṭvāpādaḥ , mālāpādaḥ iti . \\
\noindent \textbf{(P\_1,2.48.2)} ekādeśe kṛte antādivadbhāvāt prāpnoti . \\
\noindent \textbf{(P\_1,2.48.2)} ubhayataḥ āśraye na antādivat . \\
\noindent \textbf{(P\_1,2.48.2)} gonivṛttyarthena tāvat na arthaḥ . \\
\noindent \textbf{(P\_1,2.48.2)} gontasya prātipadikasya upasarjanasya hrasvaḥ bhavati iti ucyate na ca etat gontam . \\
\noindent \textbf{(P\_1,2.48.2)} nanu ca etat api vyapadeśivadbhāvena gontam . \\
\noindent \textbf{(P\_1,2.48.2)} vyapadeśivadbhāvaḥ aprātipadikena . \\
\noindent \textbf{(P\_1,2.48.2)} samāsanivṛttyarthena ca api na arthaḥ . \\
\noindent \textbf{(P\_1,2.48.2)} stryantasya prātipadikasya upasarjanasya hrasvaḥ bhavati iti ucyate . \\
\noindent \textbf{(P\_1,2.48.2)} pradhānam upasarjanam iti ca sambandhiśabdau etau . \\
\noindent \textbf{(P\_1,2.48.2)} tatra sambandhāt etat gantavyam : yam prati yat apradhānam tasya cet saḥ antaḥ bhavati iti . \\
\noindent \textbf{(P\_1,2.48.2)} avaśyam ca etat evam vijñeyam . \\
\noindent \textbf{(P\_1,2.48.2)} ucyamāne api hi pratiṣedhe iha prasajyeta : pañca kumāryaḥ priyāḥ asya pañcakumārīpriyaḥ , daśakumārīpriyaḥ iti . \\
\noindent \textbf{(P\_1,2.48.2)} kapi ca . \\
\noindent \textbf{(P\_1,2.48.2)} kapi ca pratiṣedhaḥ vaktavyaḥ : bahukumārīkaḥ , bahuvṛṣalīkaḥ . \\
\noindent \textbf{(P\_1,2.48.2)} dvandve ca . \\
\noindent \textbf{(P\_1,2.48.2)} dvandve ca pratiṣedhaḥ vaktavyaḥ : kukkuṭamayūryau . \\
\noindent \textbf{(P\_1,2.48.2)} uktam vā . \\
\noindent \textbf{(P\_1,2.48.2)} kim uktam . \\
\noindent \textbf{(P\_1,2.48.2)} kapi tāvat uktam : na kapi iti pratiṣedhaḥ iti . \\
\noindent \textbf{(P\_1,2.48.2)} na etat asti uktam . \\
\noindent \textbf{(P\_1,2.48.2)} ke aṇaḥ iti yā hrasvaprāptiḥ tasyāḥ pratiṣedhaḥ . \\
\noindent \textbf{(P\_1,2.48.2)} kutaḥ etat . \\
\noindent \textbf{(P\_1,2.48.2)} anantarasya vidhiḥ vā bhavati pratiṣedhaḥ vā iti . \\
\noindent \textbf{(P\_1,2.48.2)} avaśyam ca etat evam vijñeyam . \\
\noindent \textbf{(P\_1,2.48.2)} yaḥ hi manyate yā ca yāvatīca hrasvaprāptiḥ tasyāḥ sarvasyāḥ pratiṣedhaḥ iti iha api tasya pratiṣedhaḥ prasajyeta : priyam grāmaṇi brāhmaṇakulam asya priyagrāmaṇikaḥ . \\
\noindent \textbf{(P\_1,2.48.2)} idam tarhi uktam : kapi kṛte anantyatvāt hrasvatvam na bhaviṣyati . \\
\noindent \textbf{(P\_1,2.48.2)} idam iha sampradhāryam : kap kriyatām hrasvatvam iti . \\
\noindent \textbf{(P\_1,2.48.2)} kim atra kartavyam . \\
\noindent \textbf{(P\_1,2.48.2)} paratvāt kap . \\
\noindent \textbf{(P\_1,2.48.2)} antaraṅgam hrasvatvam . \\
\noindent \textbf{(P\_1,2.48.2)} antaraṅgataraḥ kap . \\
\noindent \textbf{(P\_1,2.48.2)} nanu ca ayam kap samāsāntaḥ ici ucyate . \\
\noindent \textbf{(P\_1,2.48.2)} tādarthyāt tācchabdyam bhaviṣyati . \\
\noindent \textbf{(P\_1,2.48.2)} yeṣām padānām samāsaḥ na tāvat teṣām anyat bhavati . \\
\noindent \textbf{(P\_1,2.48.2)} kapam tāvat pratīkṣate . \\
\noindent \textbf{(P\_1,2.48.2)} dvandve api uktam . \\
\noindent \textbf{(P\_1,2.48.2)} kim uktam . \\
\noindent \textbf{(P\_1,2.48.2)} paravat liṅgam iti śabdaśabdārthau iti . \\
\noindent \textbf{(P\_1,2.48.2)} tatra aupadeśikasya hrasvatvam ātideśikasya śravaṇam bhaviṣyati . \\
\noindent \textbf{(P\_1,2.49)} taddhitaluki avantyādīnām pratiṣedhaḥ . \\
\noindent \textbf{(P\_1,2.49)} taddhitaluki avantyādīnām pratiṣedhaḥ vaktavyaḥ : avantī kuntī kurūḥ . \\
\noindent \textbf{(P\_1,2.49)} taddhitaluki avantyādīnām apratiṣedhaḥ alukparatvāt . \\
\noindent \textbf{(P\_1,2.49)} taddhitaluki avantyādīnām apratiṣedhaḥ . \\
\noindent \textbf{(P\_1,2.49)} anarthakaḥ pratiṣedhaḥ apratiṣedhaḥ . \\
\noindent \textbf{(P\_1,2.49)} luk kasmāt na bhavati . \\
\noindent \textbf{(P\_1,2.49)} alukparatvāt . \\
\noindent \textbf{(P\_1,2.49)} luki iti ucyate . \\
\noindent \textbf{(P\_1,2.49)} na ca atra lukam paśyāmaḥ . \\
\noindent \textbf{(P\_1,2.49)} luki iti na eṣā parasaptamī śakyā vijñātum . \\
\noindent \textbf{(P\_1,2.49)} na hi lukā paurvāparyam asti . \\
\noindent \textbf{(P\_1,2.49)} kā tarhi . \\
\noindent \textbf{(P\_1,2.49)} satsaptamī : luki sati iti . \\
\noindent \textbf{(P\_1,2.49)} satsaptamī cet prāpnoti . \\
\noindent \textbf{(P\_1,2.49)} evam tarhi idam iha vyapadeśyam sat ācāryaḥ na vyapadiśati . \\
\noindent \textbf{(P\_1,2.49)} kim . \\
\noindent \textbf{(P\_1,2.49)} upasarjanasya iti vartate . \\
\noindent \textbf{(P\_1,2.49)} na ca jātiḥ upasarjanam . \\
\noindent \textbf{(P\_1,2.50)} it goṇyāḥ na iti vaktavyam . goṇyāḥ na iti eva vaktavyam . \\
\noindent \textbf{(P\_1,2.50)} na arthaḥ ittvena . \\
\noindent \textbf{(P\_1,2.50)} kā rūpasiddhiḥ : pañcagoṇiḥ , daśagoṇiḥ . \\
\noindent \textbf{(P\_1,2.50)} hrasvatā hi vidhīyate . \\
\noindent \textbf{(P\_1,2.50)} hrasvatvam atra vidhīyate : gostriyoḥ upasarjanasya iti . \\
\noindent \textbf{(P\_1,2.50)} iti vā vacane tāvat . \\
\noindent \textbf{(P\_1,2.50)} it iti vā ucyeta na iti vā kaḥ nu atra viśeṣaḥ . \\
\noindent \textbf{(P\_1,2.50)} mātrārtham vā kṛtam bhavet . atha vā mātrārtham idam vaktavyam : goṇīmātram idam goṇiḥ . \\
\noindent \textbf{(P\_1,2.50)} aparaḥ āha : goṇyāḥ ittvam prakaraṇāt . \\
\noindent \textbf{(P\_1,2.50)} aśiṣyam goṇyāḥ ittvam . \\
\noindent \textbf{(P\_1,2.50)} kim kāraṇam . \\
\noindent \textbf{(P\_1,2.50)} prakaraṇāt . \\
\noindent \textbf{(P\_1,2.50)} prakṛtam hrasvatvam . \\
\noindent \textbf{(P\_1,2.50)} hrasvaḥ iti vartate . \\
\noindent \textbf{(P\_1,2.50)} nanu sūcyāḥ . \\
\noindent \textbf{(P\_1,2.50)} sūcyādyartham atha api vā . \\
\noindent \textbf{(P\_1,2.50)} sūcyādyartham idam draṣṭavyam : pañcasūciḥ , daśasūcīḥ . \\
\noindent \textbf{(P\_1,2.50)} it goṇyāḥ na iti vaktavyam hrasvatā hi vidhīyate | iti vā vacane tāvat . \\
\noindent \textbf{(P\_1,2.50)} mātrārtham vā kṛtam bhavet || goṇyāḥ ittvam prakaraṇāt . \\
\noindent \textbf{(P\_1,2.50)} sūcyādyartham atha api vā . \\
\noindent \textbf{(P\_1,2.51.1)} vyaktivacane iti kimartham . \\
\noindent \textbf{(P\_1,2.51.1)} śirīṣāṇām adūrabhavaḥ grāmaḥ śirīṣāḥ . \\
\noindent \textbf{(P\_1,2.51.1)} tasya grāmasya vanam śirīṣavanam . \\
\noindent \textbf{(P\_1,2.51.1)} kim ca syāt . \\
\noindent \textbf{(P\_1,2.51.1)} vibhāṣā oṣadhivanaspatibhyaḥ iti ṇatvam prasajyeta . \\
\noindent \textbf{(P\_1,2.51.1)} aparaḥ āha : kaṭubadaryāḥ adūrabhavaḥ grāmaḥ kaṭubadarī . \\
\noindent \textbf{(P\_1,2.51.1)} ṣaṣṭhī yuktavadbhāvena mā bhūt iti . \\
\noindent \textbf{(P\_1,2.51.1)} atha vyaktivacane iti api ucyamāne kasmāt eva atra na bhavati . \\
\noindent \textbf{(P\_1,2.51.1)} ṣaṣṭhī api hi vacanam . \\
\noindent \textbf{(P\_1,2.51.1)} na idam pāribhāṣikasya vacanasya grahaṇam . \\
\noindent \textbf{(P\_1,2.51.1)} kim tarhi . \\
\noindent \textbf{(P\_1,2.51.1)} anvarthagrahaṇam : ucyate vacanam iti . \\
\noindent \textbf{(P\_1,2.51.1)} evam api ṣaṣṭhī prāpnoti . \\
\noindent \textbf{(P\_1,2.51.1)} ṣaṣṭhī api hi ucyate . \\
\noindent \textbf{(P\_1,2.51.1)} lupā uktatvāt tasya arthasya dvitīyasya prayogeṇa na bhavitavyam uktārthānām aprayogaḥ iti . \\
\noindent \textbf{(P\_1,2.51.1)} ātideśikī tarhi prāpnoti . \\
\noindent \textbf{(P\_1,2.51.1)} evam tarhi prāk api vr̥tteḥ yuktam vṛttam ca api . iha yāvatā yuktam vaktuḥ ca kāmacāraḥ prāk vṛtteḥ liṅgasaṅkhye ye . \\
\noindent \textbf{(P\_1,2.51.1)} prāk api vṛtteḥ yuktam vanaspatibhiḥ nagaram vṛttam ca api yuktam vanaspatibhiḥ nagaram . \\
\noindent \textbf{(P\_1,2.51.1)} vṛtte ca yuktavadbhāvaḥ vidhīyate . \\
\noindent \textbf{(P\_1,2.51.1)} kāmacāraḥ ca prayoktuḥ prāk vṛtteḥ ye liṅgasaṅkhye te* atideṣṭum vṛttasya vā ye liṅgasaṅkhye . \\
\noindent \textbf{(P\_1,2.51.1)} yāvatā kārmacāraḥ vṛttasya ye liṅgasaṅkhye te* atidiśyete na prāk vṛtteḥ ye . \\
\noindent \textbf{(P\_1,2.51.2)} kimartham punaḥ idam ucyate . \\
\noindent \textbf{(P\_1,2.51.2)} anyatra abhidheyavyaktivacanabhāvāt lupi yuktavadanudeśaḥ . \\
\noindent \textbf{(P\_1,2.51.2)} anyatra abhidheyavat liṅgavacanāni bhavanti . \\
\noindent \textbf{(P\_1,2.51.2)} kva anyatra . \\
\noindent \textbf{(P\_1,2.51.2)} luki : lavaṇaḥ supaḥ , lavaṇā yavāguḥ , lavaṇam śākam iti . \\
\noindent \textbf{(P\_1,2.51.2)} anyatra abhidheyavat liṅgavacanāni bhavanti luki. iha api anyatra abhidheyavat liṅgavacanāni prāpnuvanti . \\
\noindent \textbf{(P\_1,2.51.2)} iṣyante ca abhidhānavat syuḥ iti . \\
\noindent \textbf{(P\_1,2.51.2)} tat ca antareṇa yatnam na sidhyati iti lupi yuktavadanudeśaḥ . \\
\noindent \textbf{(P\_1,2.51.2)} evamartham idam ucyate . \\
\noindent \textbf{(P\_1,2.51.2)} asti prayojanam etat . \\
\noindent \textbf{(P\_1,2.51.2)} kim tarhi iti . \\
\noindent \textbf{(P\_1,2.51.2)} lupaḥ adarśanasañjñitvāt arthagatiḥ na upapadyate . \\
\noindent \textbf{(P\_1,2.51.2)} lup nāma iyam adarśanasya sañjñā kriyate . \\
\noindent \textbf{(P\_1,2.51.2)} na ca adarśanasya liṅgasaṅkhye śakyete* atideṣṭum . \\
\noindent \textbf{(P\_1,2.51.2)} lupaḥ adarśanasañjñitvāt arthagatiḥ na upapadyate . \\
\noindent \textbf{(P\_1,2.51.2)} na vā adarśanasya aśakyatvāt arthagatiḥ sāhacaryāt . \\
\noindent \textbf{(P\_1,2.51.2)} na vā eṣaḥ doṣaḥ . \\
\noindent \textbf{(P\_1,2.51.2)} kim kāraṇam . \\
\noindent \textbf{(P\_1,2.51.2)} adarśanasya aśakyatvāt . \\
\noindent \textbf{(P\_1,2.51.2)} adarśanasya liṅgasaṅkhye* aśakye* atideṣṭum iti kṛtvā adarśanasahacaritaḥ yaḥ arthaḥ tasya gatiḥ bhaviṣyati sāhacaryāt . \\
\noindent \textbf{(P\_1,2.51.2)} yogābhāvāt ca anyasya . \\
\noindent \textbf{(P\_1,2.51.2)} adarśanena ca yogaḥ na asti iti kṛtvā adarśanasahacaritaḥ yaḥ arthaḥ tasya gatiḥ bhaviṣyati sāhacaryāt \\
\noindent \textbf{(P\_1,2.51.3)} samāse uttarapadasya bahuvacanasya lupaḥ . \\
\noindent \textbf{(P\_1,2.51.3)} samāse uttarapadasya bahuvacanasya lupaḥ yuktavadbhāvaḥ vaktavyaḥ : madhurāpañcālāḥ . \\
\noindent \textbf{(P\_1,2.51.3)} kim prayojanam . \\
\noindent \textbf{(P\_1,2.51.3)} niyamārtham . \\
\noindent \textbf{(P\_1,2.51.3)} samāse uttarapadasya eva . \\
\noindent \textbf{(P\_1,2.51.3)} kva mā bhūt . \\
\noindent \textbf{(P\_1,2.51.3)} pañcālamadhure* iti . \\
\noindent \textbf{(P\_1,2.52.1)} katham idam vijñāyate : jātiḥ yat viśeṣaṇam iti āhosvit jāteḥ yāni viśeṣaṇāni iti . \\
\noindent \textbf{(P\_1,2.52.1)} kim ca ataḥ . \\
\noindent \textbf{(P\_1,2.52.1)} yadi vijñāyate jātiḥ yat viśeṣaṇam iti siddham pañcālāḥ janapadaḥ iti . \\
\noindent \textbf{(P\_1,2.52.1)} subhikṣaḥ sampannapānīyaḥ bahumālyaphalaḥ iti na sidhyati . \\
\noindent \textbf{(P\_1,2.52.1)} atha vijñāyate jāteḥ yāni viśeṣaṇāni iti siddham subhikṣaḥ sampannapānīyaḥ bahumālyaphalaḥ iti . \\
\noindent \textbf{(P\_1,2.52.1)} pañcālāḥ janapadaḥ iti na sidhyati . \\
\noindent \textbf{(P\_1,2.52.1)} evam tarhi na evam vijñāyate jātiḥ yat viśeṣaṇam iti na api jāteḥ yāni viśeṣaṇāni iti . \\
\noindent \textbf{(P\_1,2.52.1)} katham tarhi . \\
\noindent \textbf{(P\_1,2.52.1)} viśeṣaṇānām yuktavadbhāvaḥ bhavati ā jātiprayogāt . \\
\noindent \textbf{(P\_1,2.52.2)} kimartham punaḥ idam ucyate . \\
\noindent \textbf{(P\_1,2.52.2)} viśeṣaṇānām vacanam jātinivṛttyartham . \\
\noindent \textbf{(P\_1,2.52.2)} jātinivṛttyarthaḥ ayam ārambhaḥ . \\
\noindent \textbf{(P\_1,2.52.2)} kim ucyate jātinivṛttyarthaḥ iti na punaḥ viśeṣaṇānām api yuktavadbhāvaḥ yathā syāt iti . \\
\noindent \textbf{(P\_1,2.52.2)} samānādhikaraṇatvāt siddham . \\
\noindent \textbf{(P\_1,2.52.2)} samānādhikaraṇatvāt viśeṣaṇānām yuktavadbhāvaḥ bhaviṣyati . \\
\noindent \textbf{(P\_1,2.52.2)} yadi evam na arthaḥ anena . \\
\noindent \textbf{(P\_1,2.52.2)} lupaḥ anyatra api jāteḥ yuktavadbhāvaḥ na bhavati . \\
\noindent \textbf{(P\_1,2.52.2)} kva anyatra . \\
\noindent \textbf{(P\_1,2.52.2)} badarī sūkṣmakaṇṭakā madhurā vṛkṣaḥ iti . \\
\noindent \textbf{(P\_1,2.52.2)} kim punaḥ kāraṇam anyatra api jāteḥ yuktavadbhāvaḥ na bhavati . \\
\noindent \textbf{(P\_1,2.52.2)} āviṣṭaliṅgā jātiḥ yat liṅgam upādāya pravartate utpattiprabhṛti ā vināśāt na tat liṅgam jahāti . \\
\noindent \textbf{(P\_1,2.52.2)} na tarhi idānīm ayam yogaḥ vaktavyaḥ . \\
\noindent \textbf{(P\_1,2.52.2)} vaktavyaḥ ca . \\
\noindent \textbf{(P\_1,2.52.2)} kim prayojanam . \\
\noindent \textbf{(P\_1,2.52.2)} idam tatra tatra ucyate guṇavacanānām śabdānām āśrayataḥ liṅgavacanāni bhavanti iti . \\
\noindent \textbf{(P\_1,2.52.2)} tat anena kriyate \\
\noindent \textbf{(P\_1,2.52.3)} harītakyādiṣu vyaktiḥ . harītakyādiṣu vyaktiḥ bhavati yuktavadbhāvena : harītakyāḥ phalāni harītakyaḥ phalāni . \\
\noindent \textbf{(P\_1,2.52.3)} khalatikādiṣu vacanam . \\
\noindent \textbf{(P\_1,2.52.3)} khalatikādiṣu vacanam bhavati yuktavadbhāvena : khalatikasya parvatasya adūrabhavāni vanāni khalatikam vanāni . \\
\noindent \textbf{(P\_1,2.52.3)} manuṣyalupi pratiṣedhaḥ . \\
\noindent \textbf{(P\_1,2.52.3)} manuṣyalupi pratiṣedhaḥ vaktavyaḥ : cañcā abhirūpaḥ , vadhrikā darśanīyaḥ . \\
\noindent \textbf{(P\_1,2.53)} kim yāḥ etāḥ kṛtrimāḥ ṭighubhādisañjñāḥ tatprāmāṇyāt aśiṣyam . \\
\noindent \textbf{(P\_1,2.53)} na iti āha . \\
\noindent \textbf{(P\_1,2.53)} sañjñānam sañjñā . \\
\noindent \textbf{(P\_1,2.58)} idam ayuktam vartate . \\
\noindent \textbf{(P\_1,2.58)} kim atra ayuktam . \\
\noindent \textbf{(P\_1,2.58)} bahavaḥ te arthāḥ . \\
\noindent \textbf{(P\_1,2.58)} tatra yuktam bahuvacanam . \\
\noindent \textbf{(P\_1,2.58)} tat yat ekavacane śāsitavye bahuvacanam śiṣyate etat ayuktam . \\
\noindent \textbf{(P\_1,2.58)} bahuṣu ekavacanam iti nāma vaktavyam . \\
\noindent \textbf{(P\_1,2.58)} ataḥ uttaram paṭhati : jātyākhyāyām sāmānyābhidhānāt aikārthyam . jātyākhyāyām sāmānyābhidhānāt aikārthyam bhaviṣyati . \\
\noindent \textbf{(P\_1,2.58)} yat tat vrīhau vrīhitvam yave yavatvam gārgye gārgyatvam tat ekam tac ca vivakṣitam . \\
\noindent \textbf{(P\_1,2.58)} tasya ekatvāt ekavacanam eva prāpnoti . \\
\noindent \textbf{(P\_1,2.58)} iṣyate ca bahuvacanam syāt iti . \\
\noindent \textbf{(P\_1,2.58)} tat ca antareṇa yatnam na sidhyati iti jātyākhyāyam ekasmin bahuvacanam . \\
\noindent \textbf{(P\_1,2.58)} evamartham idam ucyate . \\
\noindent \textbf{(P\_1,2.58)} asti prayojanam etat . \\
\noindent \textbf{(P\_1,2.58)} kim tarhi iti . \\
\noindent \textbf{(P\_1,2.58)} tatra ekavacanādeśe uktam . \\
\noindent \textbf{(P\_1,2.58)} kim uktam . \\
\noindent \textbf{(P\_1,2.58)} vrīhibhyaḥ āgataḥ iti atra gheḥ ṅiti iti guṇaḥ prāpnoti iti . \\
\noindent \textbf{(P\_1,2.58)} na eṣaḥ doṣaḥ . \\
\noindent \textbf{(P\_1,2.58)} arthātideśāt siddham . arthātideśaḥ ayam . \\
\noindent \textbf{(P\_1,2.58)} na idam pāribhāṣikasya vacanasya grahaṇam . \\
\noindent \textbf{(P\_1,2.58)} kim tarhi . \\
\noindent \textbf{(P\_1,2.58)} anvarthagrahaṇam : ucyate vacanam . \\
\noindent \textbf{(P\_1,2.58)} bahūnām arthānām vacanam bahuvacanam iti . \\
\noindent \textbf{(P\_1,2.58)} yāvat brūyāt ekaḥ arthaḥ bahuvat bhavati iti tāvat ekasmin bahuvacanam iti . \\
\noindent \textbf{(P\_1,2.58)} saṅkhyāprayoge pratiṣedhaḥ . saṅkhyāprayoge pratiṣedhaḥ vaktavyaḥ . \\
\noindent \textbf{(P\_1,2.58)} ekaḥ vrīhiḥ sampannaḥ subhikṣam karoti . \\
\noindent \textbf{(P\_1,2.58)} asmadaḥ nāmayuvapratyayayoḥ ca . asmadaḥ nāmaprayoge yuvapratyayaprayoge ca pratiṣedhaḥ vaktavyaḥ . \\
\noindent \textbf{(P\_1,2.58)} nāmaprayoge : aham devadattaḥ bravīmi . \\
\noindent \textbf{(P\_1,2.58)} aham yajñadattaḥ bravīmi . \\
\noindent \textbf{(P\_1,2.58)} yuvapratyayaprayoge : ahaṃ gārgyāyaṇaḥ bravīmi . \\
\noindent \textbf{(P\_1,2.58)} aham vātsyāyanaḥ bravīmi . \\
\noindent \textbf{(P\_1,2.58)} yuvagrahaṇena nārthaḥ . \\
\noindent \textbf{(P\_1,2.58)} asmadaḥ nāmapratyayaprayoge na iti eva . \\
\noindent \textbf{(P\_1,2.58)} idam api siddham bhavati : aham gārgyaḥ bravīmi . \\
\noindent \textbf{(P\_1,2.58)} aham vātsyaḥ bravīmi . \\
\noindent \textbf{(P\_1,2.58)} aparaḥ āha : asmadaḥ saviśeṣaṇasya prayoge na iti eva . \\
\noindent \textbf{(P\_1,2.58)} idam api siddham bhavati : aham paṭuḥ bravīmi . \\
\noindent \textbf{(P\_1,2.58)} aham paṇḍitaḥ bravīmi . \\
\noindent \textbf{(P\_1,2.58)} aśiṣyam vā bahuvat pṛthakttvābhidhānāt . \\
\noindent \textbf{(P\_1,2.58)} aśiṣyaḥ vā bahuvadbhāvaḥ . \\
\noindent \textbf{(P\_1,2.58)} kim kāraṇam . \\
\noindent \textbf{(P\_1,2.58)} pṛthaktvābhidhānāt . \\
\noindent \textbf{(P\_1,2.58)} pṛthaktvena hi dravyāṇi abhidhīyante . \\
\noindent \textbf{(P\_1,2.58)} bahavaḥ te arthāḥ . \\
\noindent \textbf{(P\_1,2.58)} tatra yuktam bahuvacanam . \\
\noindent \textbf{(P\_1,2.58)} kim ucyate pṛthaktvābhidhānāt iti yāvatā idānīm eva uktam : jātyākhyāyām sāmānyābhidhānāt aikārthyam iti . \\
\noindent \textbf{(P\_1,2.58)} jātiśabdena hi dravyābhidhānam . jātiśabdena hi dravyam api abhidhīyate jātiḥ api . \\
\noindent \textbf{(P\_1,2.58)} katham punaḥ jñāyate jātiśabdena dravyam api abhidhīyate iti . \\
\noindent \textbf{(P\_1,2.58)} evam hi kaḥ cit mahati gomaṇḍale gopālakam āsīnam pṛcchati : asti atra kām cid gām paśyasi iti . \\
\noindent \textbf{(P\_1,2.58)} saḥ paśyati : paśyati ca ayam gāḥ pṛcchati ca kām cid atra gām paśyasi iti . \\
\noindent \textbf{(P\_1,2.58)} nūnam asya dravyam vivakṣitam iti . \\
\noindent \textbf{(P\_1,2.58)} tat yadā dravyābhidhānam tadā bahuvacanam bhaviṣyati . \\
\noindent \textbf{(P\_1,2.58)} yadā sāmānyābhidhānam tadā ekavacanam bhaviṣyati . \\
\noindent \textbf{(P\_1,2.59)} ayam api yogaḥ śakyaḥ avaktum . \\
\noindent \textbf{(P\_1,2.59)} katham aham bravīmi , āvām brūvaḥ , vayam brūmaḥ . \\
\noindent \textbf{(P\_1,2.59)} imāni indriyāṇi kadā cit svātantryeṇa vivakṣitāni bhavanti . \\
\noindent \textbf{(P\_1,2.59)} tat yathā : idam me akṣi suṣṭhu paśyati . \\
\noindent \textbf{(P\_1,2.59)} ayam me karṇaḥ suṣṭhu śṛṇoti iti . \\
\noindent \textbf{(P\_1,2.59)} kadā cit pāratantryeṇa : anena akṣṇā suṣṭhu paśyāmi . \\
\noindent \textbf{(P\_1,2.59)} anena karṇena suṣṭhu śṛṇomi iti . \\
\noindent \textbf{(P\_1,2.59)} tat yadā svātantryeṇa vivakṣā tadā bahuvacanam bhaviṣyati . \\
\noindent \textbf{(P\_1,2.59)} yadā pāratantryeṇa tadā ekavacanadvivacane bhaviṣyataḥ . \\
\noindent \textbf{(P\_1,2.60)} ayam api yogaḥ śakyaḥ avaktum . \\
\noindent \textbf{(P\_1,2.60)} katham udite pūrve phalgunyau , uditāḥ pūrvāḥ phalgunyaḥ , udite pūrve proṣṭhapade , uditāḥ pūrvāḥ proṣṭhapadāḥ . \\
\noindent \textbf{(P\_1,2.60)} phalgunīsampīpagate candramasi phalgunīśabdaḥ vartate . \\
\noindent \textbf{(P\_1,2.60)} bahavaḥ te arthāḥ . \\
\noindent \textbf{(P\_1,2.60)} tatra yuktam bahuvacanam . \\
\noindent \textbf{(P\_1,2.60)} yadā tayoḥ eva abhidhānam tadā dvivacanam bhaviṣyati . \\
\noindent \textbf{(P\_1,2.61-62)} KA\_I,231.10-12 Ro\_II,110 {1/4}     imau api yogau śakyau avaktum . \\
\noindent \textbf{(P\_1,2.61-62)} KA\_I,231.10-12 Ro\_II,110 {2/4}     katham . \\
\noindent \textbf{(P\_1,2.61-62)} KA\_I,231.10-12 Ro\_II,110 {3/4}     punarvasuviśākhayoḥ supām sulukpūrvasavarṇa iti siddham . \\
\noindent \textbf{(P\_1,2.61-62)} KA\_I,231.10-12 Ro\_II,110 {4/4}     punarvasuviśākhayoḥ supām sulukpūrvasavarṇa iti eva siddham \\
\noindent \textbf{(P\_1,2.63)} tiṣyapunarvasvoḥ iti kimartham . \\
\noindent \textbf{(P\_1,2.63)} kṛttikārohiṇyaḥ . \\
\noindent \textbf{(P\_1,2.63)} nakṣatra iti kimartham . \\
\noindent \textbf{(P\_1,2.63)} tiṣyaḥ ca māṇavakaḥ punarvasū maṇavakau tiṣyapunarvasavaḥ . \\
\noindent \textbf{(P\_1,2.63)} atha nakṣatre iti vartamāne punaḥ nakṣatragrahaṇam kimartham . \\
\noindent \textbf{(P\_1,2.63)} ayam tiṣyapunarvasuśabdaḥ asti eva jyotiṣi vartate . \\
\noindent \textbf{(P\_1,2.63)} asti ca kālavācī . \\
\noindent \textbf{(P\_1,2.63)} tat yathā : bahavaḥ tiṣyapunarvasavaḥ atikrāntāḥ . \\
\noindent \textbf{(P\_1,2.63)} katareṇa tiṣyeṇa gataḥ iti . \\
\noindent \textbf{(P\_1,2.63)} tat yaḥ jyotiṣi vartate tasya idam grahaṇam . \\
\noindent \textbf{(P\_1,2.63)} atha vā nakṣatre iti vartamāne punaḥ nakṣatragrahaṇasya etat prayojanam : videśastham api tiṣyapunarvasvoḥ kāryam tat api nakṣatrasya eva yathā syāt : tiṣyapuṣyayoḥ nakṣatrāṇi yalopaḥ vaktavyaḥ iti nakṣatragrahaṇam na kartavyam bhavati . \\
\noindent \textbf{(P\_1,2.63)} atha vā atha vā nakṣatre iti vartamāne punaḥ nakṣatra grahaṇasya etat prayojanam : tiṣyapunarvasuparyāyavācinām api yathā syāt : puṣyapunarvasū sidhyapunarvasū . \\
\noindent \textbf{(P\_1,2.63)} atha dvandve iti kimartham . \\
\noindent \textbf{(P\_1,2.63)} yaḥ tiṣyaḥ tau punarvasū yeṣām te ime tiṣyapunarvasavaḥ unmugdhāḥ . \\
\noindent \textbf{(P\_1,2.63)} bahuvacanasya iti kimartham . \\
\noindent \textbf{(P\_1,2.63)} uditam tiṣyapunarvasū . \\
\noindent \textbf{(P\_1,2.63)} katham ca atra ekavacanam . \\
\noindent \textbf{(P\_1,2.63)} jātidvandvaḥ ekavat bhavati iti . \\
\noindent \textbf{(P\_1,2.63)} aprāṇinām iti pratiṣedhaḥ prāpnoti . \\
\noindent \textbf{(P\_1,2.63)} evam tarhi siddhe sati yat bahuvacanagrahaṇam karoti tat jñāpayati ācāryaḥ : sarvaḥ dvandvaḥ vibhāṣā ekavat bhavati iti . \\
\noindent \textbf{(P\_1,2.63)} kim etasya jñāpane prayojanam . \\
\noindent \textbf{(P\_1,2.63)} bābhravaśālaṅkāyanam bābhravaśālaṅkāyanāḥ iti etat siddham bhavati . \\
\noindent \textbf{(P\_1,2.63)} atha vā na atra bhavantaḥ prāṇināḥ . \\
\noindent \textbf{(P\_1,2.63)} prāṇāḥ eva atra bhavantaḥ . \\
\noindent \textbf{(P\_1,2.64.1)} rūpagrahaṇam kimartham . \\
\noindent \textbf{(P\_1,2.64.1)} samānānām ekaśeṣa ekavibhaktau iti iyati ucyamāne yatra eva sarvam samānam śabdaḥ arthaḥ ca tatra eva syāt : vṛkṣāḥ , plakṣāḥ iti . \\
\noindent \textbf{(P\_1,2.64.1)} iha na syāt : akṣāḥ ,. pādāḥ , māṣāḥ iti . \\
\noindent \textbf{(P\_1,2.64.1)} rūpagrahaṇe punaḥ kriyamāṇe na doṣaḥ bhavati . \\
\noindent \textbf{(P\_1,2.64.1)} rūpam nimittatvena āśrīyate śrutau ca rūpagrahaṇam . \\
\noindent \textbf{(P\_1,2.64.1)} atha ekagrahaṇam kimartham . \\
\noindent \textbf{(P\_1,2.64.1)} sarūpāṇām śeṣaḥ ekavibhaktau iti iyati ucyamāne dvibahvoḥ api śeṣaḥ prasajyeta . \\
\noindent \textbf{(P\_1,2.64.1)} ekagrahaṇe punaḥ kriyamāṇe na doṣaḥ bhavati . \\
\noindent \textbf{(P\_1,2.64.1)} atha śeṣagrahaṇam kimartham . \\
\noindent \textbf{(P\_1,2.64.1)} sarūpāṇām ekaḥ ekavibhaktau iti iyati ucyamāne ādeśaḥ ayam vijñāyeta . \\
\noindent \textbf{(P\_1,2.64.1)} tatra kaḥ doṣaḥ . \\
\noindent \textbf{(P\_1,2.64.1)} aśvaḥ ca asvaḥ ca aśvau : āntaryataḥ dvyudāttavataḥ sthāninaḥ dvyudāttavān ādeśaḥ prasajyeta . \\
\noindent \textbf{(P\_1,2.64.1)} lopyalopitā ca na prakalpeta . \\
\noindent \textbf{(P\_1,2.64.1)} tatra kaḥ doṣaḥ . \\
\noindent \textbf{(P\_1,2.64.1)} gargāḥ , vatsāḥ , bidāḥ , urvāḥ . \\
\noindent \textbf{(P\_1,2.64.1)} añ yaḥ bahuṣu yañ yaḥ bahuṣu iti ucyamānaḥ luk na prāpnoti . \\
\noindent \textbf{(P\_1,2.64.1)} mā bhūt evam . \\
\noindent \textbf{(P\_1,2.64.1)} añantam yat bahuṣu yañantam yat bahuṣu iti evam bhaviṣyati . \\
\noindent \textbf{(P\_1,2.64.1)} na evam śakyam . \\
\noindent \textbf{(P\_1,2.64.1)} iha hi doṣaḥ syāt. : kāśyapapratikṛtayaḥ kāśyapāḥ iti . \\
\noindent \textbf{(P\_1,2.64.1)} ekavibhaktau iti kimartham . \\
\noindent \textbf{(P\_1,2.64.1)} payaḥ payaḥ jarayati . \\
\noindent \textbf{(P\_1,2.64.1)} vāsaḥ vāsaḥ chādayati . \\
\noindent \textbf{(P\_1,2.64.1)} brāhmaṇābhyām ca kṛtam brāhmaṇābhyām ca dehi iti . \\
\noindent \textbf{(P\_1,2.64.2)} kimartham punaḥ idam ucyate . \\
\noindent \textbf{(P\_1,2.64.2)} pratyartham śabdaniveśāt na ekena anekasya abhidhānam . pratyartham śabdāḥ abhiniviśante . \\
\noindent \textbf{(P\_1,2.64.2)} kim idam pratyartham iti . \\
\noindent \textbf{(P\_1,2.64.2)} artham artham prati pratyartham . \\
\noindent \textbf{(P\_1,2.64.2)} pratyartham śābdaniveśāt etasmāt kāraṇāt na ekena śabdena anekasya arthasya abhidhānam prāpnoti . \\
\noindent \textbf{(P\_1,2.64.2)} tatra kaḥ doṣaḥ . \\
\noindent \textbf{(P\_1,2.64.2)} tatra anekārthābhidhāne anekaśabdatvam . \\
\noindent \textbf{(P\_1,2.64.2)} tatra anekārthābhidhāne anekaśabdatvam prāpnoti . \\
\noindent \textbf{(P\_1,2.64.2)} iṣyate ca ekena api anekasya abhidhānam syāt iti . \\
\noindent \textbf{(P\_1,2.64.2)} tat ca antareṇa yatnam na sidhyati . \\
\noindent \textbf{(P\_1,2.64.2)} tasmāt ekaśeṣaḥ . \\
\noindent \textbf{(P\_1,2.64.2)} evamartham idam ucyate . \\
\noindent \textbf{(P\_1,2.64.2)} asti prayojanam etat . \\
\noindent \textbf{(P\_1,2.64.2)} kim tarhi iti . \\
\noindent \textbf{(P\_1,2.64.2)} kim idam pratyartham śabdāḥ abhiniveśante iti etam dṛṣṭāntam āsthāya sarūpāṇām ekaśeṣaḥ ārabhyate na punaḥ apratyartham śabdāḥ abhiniviśante iti etam dṛṣṭāntam āsthāya virūpāṇām anekaśeṣaḥ ārabhyate . \\
\noindent \textbf{(P\_1,2.64.2)} tatra etat syāt : laghīyasī sarūpanivṛttirḥ garīyasī virūpapratipattiḥ iti . \\
\noindent \textbf{(P\_1,2.64.2)} tat ca na . \\
\noindent \textbf{(P\_1,2.64.2)} laghīyasī virūpapratipattiḥ . \\
\noindent \textbf{(P\_1,2.64.2)} kim kāraṇam . \\
\noindent \textbf{(P\_1,2.64.2)} yatra hi bahūnām sarūpāṇām ekaḥ śiṣyate tatra avarataḥ dvayoḥ sarūpayoḥ nivṛttiḥ vaktavyā syāt . \\
\noindent \textbf{(P\_1,2.64.2)} evam api etasmin sati kim cit ācāryaḥ sukaratarakam manyate . \\
\noindent \textbf{(P\_1,2.64.2)} sukaratarakam ca ekaśeṣārambham manyate . \\
\noindent \textbf{(P\_1,2.64.3)} kim punaḥ ayam ekavibhaktau ekaśeṣaḥ bhavati . \\
\noindent \textbf{(P\_1,2.64.3)} evam bhavitum arhati . \\
\noindent \textbf{(P\_1,2.64.3)} ekavibhaktau iti cet na abhāvād vibhakteḥ . ekavibhaktau iti cet tat na . \\
\noindent \textbf{(P\_1,2.64.3)} kim kāraṇam . \\
\noindent \textbf{(P\_1,2.64.3)} abhāvāt vibhakteḥ . \\
\noindent \textbf{(P\_1,2.64.3)} na hi samudāyāt parā vibhaktiḥ asti . \\
\noindent \textbf{(P\_1,2.64.3)} kim kāraṇam . \\
\noindent \textbf{(P\_1,2.64.3)} aprātipadikatvāt . \\
\noindent \textbf{(P\_1,2.64.3)} nanu ca arthavat prātipadikam iti prātipadikasañjñā bhaviṣyati . \\
\noindent \textbf{(P\_1,2.64.3)} niyamāt na prāpnoti . \\
\noindent \textbf{(P\_1,2.64.3)} arthavatsamudayānām samāsagrahaṇam niyamārtham iti . \\
\noindent \textbf{(P\_1,2.64.3)} yadi punaḥ pṛthak sarveṣām vibhaktiparāṇām ekaśeṣaḥ ucyeta . \\
\noindent \textbf{(P\_1,2.64.3)} pṛthak sarveṣām iti cet ekaśeṣe pṛthak vibhaktyupalabdhiḥ tadāśrayatvāt . \\
\noindent \textbf{(P\_1,2.64.3)} pṛthak sarveṣām iti cet ekaśeṣe pṛthak vibhaktyupalabdhiḥ prāpnoti . \\
\noindent \textbf{(P\_1,2.64.3)} kim ucyate ekaśeṣe pṛthak vibhaktyupalabdhiḥ iti yāvatā samayaḥ kṛtaḥ : na kevalā prakṛtiḥ prayoktavyā na kevalaḥ pratyayaḥ iti . \\
\noindent \textbf{(P\_1,2.64.3)} tadāśrayatvāt prāpnoti . \\
\noindent \textbf{(P\_1,2.64.3)} yatra hi prakṛtinimittā pratyayanivṛttiḥ tatra apratyayikāyāḥ prakṛteḥ prayogaḥ bhavati agnicit somasut iti yathā . \\
\noindent \textbf{(P\_1,2.64.3)} yatra ca pratyayanimittā prakṛtinivṛttiḥ tatra aprakṛtikasya pratyayasya prayogaḥ bhavati adhunā , iyān iti yathā . \\
\noindent \textbf{(P\_1,2.64.3)} astu saṃyogāntalopena siddham . \\
\noindent \textbf{(P\_1,2.64.3)} kutaḥ nu khalu etat parayoḥ vṛkṣaśabdayoḥ nivṛttiḥ bhaviṣyati na punaḥ pūrvayoḥ iti . \\
\noindent \textbf{(P\_1,2.64.3)} tatra etat syāt : pūrvanivṛttav api satyām saṃyogādilopena siddham iti . \\
\noindent \textbf{(P\_1,2.64.3)} na sidhyati . \\
\noindent \textbf{(P\_1,2.64.3)} tatra avarataḥ dvayoḥ sakārayoḥ śravaṇam prasajyeta. yatra ca saṃyogāntalopaḥ na asti tatra ca na sidhyati . \\
\noindent \textbf{(P\_1,2.64.3)} kva ca saṃyogāntalopaḥ na asti . \\
\noindent \textbf{(P\_1,2.64.3)} dvivacanabahuvacanayoḥ . \\
\noindent \textbf{(P\_1,2.64.3)} yadi punaḥ samāse ekaśeṣaḥ ucyeta . \\
\noindent \textbf{(P\_1,2.64.3)} kim kṛtam bhavati . \\
\noindent \textbf{(P\_1,2.64.3)} kaḥ cit vacanalopaḥ parihṛtaḥ bhavati . \\
\noindent \textbf{(P\_1,2.64.3)} tat tarhi samāsagrahaṇam kartavyam . \\
\noindent \textbf{(P\_1,2.64.3)} na kartavyam . \\
\noindent \textbf{(P\_1,2.64.3)} prakṛtam anuvartate . \\
\noindent \textbf{(P\_1,2.64.3)} kva prakṛtam . \\
\noindent \textbf{(P\_1,2.64.3)} tiṣyapunarvasvoḥ nakṣatradvandve bahuvacanasya dvivacanam nityam iti . \\
\noindent \textbf{(P\_1,2.64.3)} samāse iti cet svarasamāsānteṣu doṣaḥ . \\
\noindent \textbf{(P\_1,2.64.3)} samāse iti cet svarasamāsānteṣu doṣaḥ bhavati . \\
\noindent \textbf{(P\_1,2.64.3)} svara : aśvaḥ ca aśvaḥ ca aśvau. samāsāntodāttatve kṛte ekaśeṣaḥ prāpnoti . \\
\noindent \textbf{(P\_1,2.64.3)} idam iha sampradhāryam : samāsāntodāttattvam kriyatām ekaśeṣaḥ iti . \\
\noindent \textbf{(P\_1,2.64.3)} kim atra kartavyam . \\
\noindent \textbf{(P\_1,2.64.3)} paratvāt samāsāntodāttatvam . \\
\noindent \textbf{(P\_1,2.64.3)} samāsāntodāttatve ca doṣaḥ bhavati . \\
\noindent \textbf{(P\_1,2.64.3)} svara . \\
\noindent \textbf{(P\_1,2.64.3)} samāsānta : ṛk ca ṛk ca ṛcau . \\
\noindent \textbf{(P\_1,2.64.3)} samāsānte kṛte asārūpyāt ekaśeṣaḥ na prāpnoti . \\
\noindent \textbf{(P\_1,2.64.3)} idam iha sampradhāryam : samāsāntaḥ kriyatām ekaśeṣaḥ iti . \\
\noindent \textbf{(P\_1,2.64.3)} kim atra kartavyam . \\
\noindent \textbf{(P\_1,2.64.3)} paratvāt samāsāntaḥ . \\
\noindent \textbf{(P\_1,2.64.3)} samāsānte ca doṣaḥ bhavati . \\
\noindent \textbf{(P\_1,2.64.3)} aṅgāśraye ca ekaśeṣavacanam . aṅgāśraye ca kārye ekaśeṣaḥ vaktavyaḥ . \\
\noindent \textbf{(P\_1,2.64.3)} svasā ca svasārau ca svasāraḥ . \\
\noindent \textbf{(P\_1,2.64.3)} aṅgāśraye kṛte asārūpyāt ekaśeṣaḥ na prāpnoti . \\
\noindent \textbf{(P\_1,2.64.3)} idam iha sampradhāryam : aṅgāśrayam kriyatām ekaśeṣaḥ iti . \\
\noindent \textbf{(P\_1,2.64.3)} kim atra kartavyam . \\
\noindent \textbf{(P\_1,2.64.3)} paratvāt aṅgāśrayam . \\
\noindent \textbf{(P\_1,2.64.3)} tiṅsamāse tiṅsamāsavacanam . \\
\noindent \textbf{(P\_1,2.64.3)} tiṅsamāse tiṅsamāsaḥ vaktavyaḥ . \\
\noindent \textbf{(P\_1,2.64.3)} ekam tiṅgrahaṇam anarthakam . \\
\noindent \textbf{(P\_1,2.64.3)} samāse tiṅsamāsaḥ iti eva siddham . \\
\noindent \textbf{(P\_1,2.64.3)} na anarthakam . \\
\noindent \textbf{(P\_1,2.64.3)} tiṅsamāse prakṛte tiṅsamāsaḥ vaktavyaḥ . \\
\noindent \textbf{(P\_1,2.64.3)} tiṅvidhipratiṣedhaḥ ca . \\
\noindent \textbf{(P\_1,2.64.3)} tiṅ ca kaḥ cit vidheyaḥ kaḥ cit pratiṣedhyaḥ . \\
\noindent \textbf{(P\_1,2.64.3)} pacati ca pacati ca pacataḥ : taḥśabdaḥ vidheyaḥ tiśabdaḥ pratiṣedhyaḥ . \\
\noindent \textbf{(P\_1,2.64.3)} yadi punaḥ asamāse ekaśeṣaḥ ucyeta . \\
\noindent \textbf{(P\_1,2.64.3)} asamāse vacanalopaḥ . \\
\noindent \textbf{(P\_1,2.64.3)} yadi asamāse vacanalopaḥ vaktavyaḥ . \\
\noindent \textbf{(P\_1,2.64.3)} nanu ca utpatatā eva vacanalopam coditāḥ smaḥ . \\
\noindent \textbf{(P\_1,2.64.3)} dvivacanabahuvacanavidhim dvandvapratiṣedham ca vakṣyati tadartham punaḥ codyate . \\
\noindent \textbf{(P\_1,2.64.3)} dvivacanabahuvacanavidhiḥ . \\
\noindent \textbf{(P\_1,2.64.3)} dvivacanabahuvacanāni vidheyāni : vṛkṣaḥ ca vṛkṣaḥ ca vṛkṣau , vṛkṣaḥ ca vṛkṣaḥ ca vṛkṣaḥ ca vṛkṣāḥ iti . \\
\noindent \textbf{(P\_1,2.64.3)} dvandvapratiṣedhaḥ ca . \\
\noindent \textbf{(P\_1,2.64.3)} dvandvasya ca pratiṣedhaḥ vaktavyaḥ : vṛkṣaḥ ca vṛkṣaḥ ca vṛkṣau , vṛkṣaḥ ca vṛkṣaḥ ca vṛkṣaḥ ca vṛkṣāḥ iti . \\
\noindent \textbf{(P\_1,2.64.3)} cārthe dvandvaḥ iti dvandvaḥ prāpnoti . \\
\noindent \textbf{(P\_1,2.64.3)} na eṣaḥ doṣaḥ . \\
\noindent \textbf{(P\_1,2.64.3)} anavakāśaḥ ekaśeṣaḥ dvandvam bādhiṣyate . \\
\noindent \textbf{(P\_1,2.64.3)} sāvakāśaḥ ekaśeṣaḥ . \\
\noindent \textbf{(P\_1,2.64.3)} kaḥ avakāśaḥ . \\
\noindent \textbf{(P\_1,2.64.3)} tiṅantāni avakāśaḥ . \\
\noindent \textbf{(P\_1,2.64.3)} yadi punaḥ pṛthak sarveṣām vibhaktyantānām ekaśeṣaḥ ucyeta . \\
\noindent \textbf{(P\_1,2.64.3)} kim kṛtam bhavati . \\
\noindent \textbf{(P\_1,2.64.3)} kaḥ cit vacanalopaḥ parihṛtaḥ bhavati . \\
\noindent \textbf{(P\_1,2.64.3)} vibhaktyantānām ekaśeṣe vibhaktyantānām ekaśeṣe vibhaktyantānām eva tu nivṛttiḥ bhavati . \\
\noindent \textbf{(P\_1,2.64.3)} ekavibhaktyantānām iti tu pṛthagvibhaktipratiṣedhārtham . \\
\noindent \textbf{(P\_1,2.64.3)} ekavibhaktyantānām iti tu vaktavyam . \\
\noindent \textbf{(P\_1,2.64.3)} kim prayojanam . \\
\noindent \textbf{(P\_1,2.64.3)} pṛthagvibhaktipratiṣedhārtham . \\
\noindent \textbf{(P\_1,2.64.3)} pṛthagvibhaktyantānām mā bhūt : brāhmaṇābhyām ca kṛtam brāhmaṇābhyām ca dehi . \\
\noindent \textbf{(P\_1,2.64.3)} na vā arthavipratiṣedhāt yugapadvacanābhāvaḥ . \\
\noindent \textbf{(P\_1,2.64.3)} na vā eṣaḥ doṣaḥ . \\
\noindent \textbf{(P\_1,2.64.3)} kim kāraṇam . \\
\noindent \textbf{(P\_1,2.64.3)} arthavipratiṣedhāt . \\
\noindent \textbf{(P\_1,2.64.3)} vipratiṣiddhau etau arthau kartā saṃpradānam iti aśakyau yugapat nirdeṣṭum . \\
\noindent \textbf{(P\_1,2.64.3)} tayoḥ vipratiṣiddhatvāt yugapadvacanam na bhaviṣyati . \\
\noindent \textbf{(P\_1,2.64.3)} anekārthāśrayaḥ ca punaḥ ekaśeṣaḥ . \\
\noindent \textbf{(P\_1,2.64.3)} anekam artham sampratyāyayiṣyāmi iti ekaśeṣaḥ ārabhyate . \\
\noindent \textbf{(P\_1,2.64.3)} tasmāt na ekaśabdatvam . \\
\noindent \textbf{(P\_1,2.64.3)} tasmāt ekaśabdatvam na bhaviṣyati . \\
\noindent \textbf{(P\_1,2.64.3)} ayam tarhi doṣaḥ : kaḥ cit vacanalopaḥ dvivacanabahuvacanavidhiḥ dvandvapratiṣedhaḥ ca iti . \\
\noindent \textbf{(P\_1,2.64.3)} yadi punaḥ prātipadikānām ekaśeṣaḥ ucyeta . \\
\noindent \textbf{(P\_1,2.64.3)} kim kṛtam bhavati . \\
\noindent \textbf{(P\_1,2.64.3)} vacanalopaḥ parihṛtaḥ bhavati . \\
\noindent \textbf{(P\_1,2.64.3)} prātipadikānām ekaśeṣe mātṛmātroḥ pratiṣedhaḥ sarūpatvāt . prātipadikānām ekaśeṣe mātṛmātroḥ pratiṣedhaḥ vaktavyaḥ : mātā ca janayitrī mātārau ca dhānyasya mātṛmātāraḥ . \\
\noindent \textbf{(P\_1,2.64.3)} kim kāraṇam . \\
\noindent \textbf{(P\_1,2.64.3)} sarūpatvāt . \\
\noindent \textbf{(P\_1,2.64.3)} sarūpāṇi hi etāni prātipadikāni . \\
\noindent \textbf{(P\_1,2.64.3)} kim ucyate prātipadikānām ekaśeṣe mātṛmātroḥ pratiṣedhaḥ vaktavyaḥ iti na punaḥ yasya api vibhaktyantānām ekaśeṣaḥ tena api mātṛmātroḥ pratiṣedhaḥ vaktavyaḥ syāt . \\
\noindent \textbf{(P\_1,2.64.3)} tasya api hi etāni kva cit vibhaktyantāni sarūpāṇi : mātṛbhyām ca mātṛbhyāṃ ca iti . \\
\noindent \textbf{(P\_1,2.64.3)} atha matam etat vibhaktyantānām sārūpye bhavitavyam eva ekaśeṣeṇa iti prātipadikānām eva ekaśeṣe doṣaḥ bhavati . \\
\noindent \textbf{(P\_1,2.64.3)} evam ca kṛtvā codyate . \\
\noindent \textbf{(P\_1,2.64.3)} haritahariṇaśyetaśyenarohitarohiṇānām striyām upasaṅkhyānam . haritahariṇaśyetaśyenarohitarohiṇānām striyām upasaṅkhyānam kartavyam . \\
\noindent \textbf{(P\_1,2.64.3)} haritasya strī hariṇī hariṇasya api hariṇī , hariṇī ca hariṇī ca hariṇyau . \\
\noindent \textbf{(P\_1,2.64.3)} śyetasya strī śyenī śyenasya api śyenī , śyenī ca śyenī ca śyenyau . \\
\noindent \textbf{(P\_1,2.64.3)} rohitasya strī rohiṇī rohiṇasya api rohiṇī , rohiṇī ca rohiṇī ca rohiṇyau . \\
\noindent \textbf{(P\_1,2.64.3)} na vā padasya arthe prayogāt . \\
\noindent \textbf{(P\_1,2.64.3)} na vā eṣaḥ doṣaḥ . \\
\noindent \textbf{(P\_1,2.64.3)} kim kāraṇam . \\
\noindent \textbf{(P\_1,2.64.3)} padasya arthe prayogāt . \\
\noindent \textbf{(P\_1,2.64.3)} padam arthe prayujyate vibhaktyantam ca padam . \\
\noindent \textbf{(P\_1,2.64.3)} rūpam ca iha āśrīyate . \\
\noindent \textbf{(P\_1,2.64.3)} rūpanirgrahaḥ ca śabdasya na antareṇa laukikam prayogam . \\
\noindent \textbf{(P\_1,2.64.3)} tasmin ca laukike prayoge sarūpāṇi etāni . \\
\noindent \textbf{(P\_1,2.64.3)} aparaḥ āha : na vā padasya arthe prayogāt . \\
\noindent \textbf{(P\_1,2.64.3)} na vā eṣaḥ pakṣaḥ eva asti prātipadikānām ekaśeṣaḥ iti . \\
\noindent \textbf{(P\_1,2.64.3)} kim kārāṇam . \\
\noindent \textbf{(P\_1,2.64.3)} padasya arthe prayogāt . \\
\noindent \textbf{(P\_1,2.64.3)} padam arthe prayujyate vibhaktyantam ca padam . \\
\noindent \textbf{(P\_1,2.64.3)} rūpam ca iha āśrīyate rūpanirgrahaḥ ca śabdasya na antareṇa laukikam prayogam . \\
\noindent \textbf{(P\_1,2.64.3)} tasmin ca laukike prayoge prātipadikānām prayogaḥ na asti . \\
\noindent \textbf{(P\_1,2.64.3)} atha anena pakṣeṇa arthaḥ syāt : prātipadikānām ekaśeṣaḥ iti . \\
\noindent \textbf{(P\_1,2.64.3)} bāḍham arthaḥ . \\
\noindent \textbf{(P\_1,2.64.3)} kim vaktavyam etat . \\
\noindent \textbf{(P\_1,2.64.3)} na hi . \\
\noindent \textbf{(P\_1,2.64.3)} katham anucyamānam gaṃsyate . \\
\noindent \textbf{(P\_1,2.64.3)} etena eva abhihitam sūtreṇa sarūpāṇām ekaśeṣaḥ ekavibhaktau iti . \\
\noindent \textbf{(P\_1,2.64.3)} katham . \\
\noindent \textbf{(P\_1,2.64.3)} vibhaktiḥ sārūpyeṇa āśrīyate . \\
\noindent \textbf{(P\_1,2.64.3)} anaimittikaḥ ekaśeṣaḥ . \\
\noindent \textbf{(P\_1,2.64.3)} ekavibhaktau yāni sarūpāṇi teṣām ekaśeṣaḥ bhavati . \\
\noindent \textbf{(P\_1,2.64.3)} kva . \\
\noindent \textbf{(P\_1,2.64.3)} yatra vā tatra vā iti . \\
\noindent \textbf{(P\_1,2.64.3)} atha anena pakṣeṇa arthaḥ syāt : vibhaktyantānām ekaśeṣaḥ iti . \\
\noindent \textbf{(P\_1,2.64.3)} bāḍham arthaḥ . \\
\noindent \textbf{(P\_1,2.64.3)} kim vaktavyam etat . \\
\noindent \textbf{(P\_1,2.64.3)} na hi . \\
\noindent \textbf{(P\_1,2.64.3)} katham anucyamānam gaṃsyate . \\
\noindent \textbf{(P\_1,2.64.3)} etat api etena eva abhihitam sūtreṇa sarūpāṇām ekaśeṣaḥ ekavibhaktau iti . \\
\noindent \textbf{(P\_1,2.64.3)} katham . \\
\noindent \textbf{(P\_1,2.64.3)} na idam pāribhāṣikyāḥ vibhakteḥ grahaṇam . \\
\noindent \textbf{(P\_1,2.64.3)} kim tarhi . \\
\noindent \textbf{(P\_1,2.64.3)} anvarthagrahaṇam : vibhāgaḥ vibhaktiḥ iti . \\
\noindent \textbf{(P\_1,2.64.3)} ekavibhāge yāni sarūpāṇi teṣām ekaśeṣaḥ bhavati iti . \\
\noindent \textbf{(P\_1,2.64.3)} nanu ca uktaṃ : kaḥ cit vacanalopaḥ dvivacanabahuvacanavidhiḥ dvandvapratiṣedhaḥ ca iti . \\
\noindent \textbf{(P\_1,2.64.3)} na eṣaḥ doṣaḥ . \\
\noindent \textbf{(P\_1,2.64.3)} yat tāvat ucyate kaḥ cit vacanalopaḥ dvivacanabahuvacanavidhiḥiti . \\
\noindent \textbf{(P\_1,2.64.3)} sahavivakṣāyām ekaśeṣaḥ . \\
\noindent \textbf{(P\_1,2.64.3)} yugapadvivakṣāyām ekaśeṣeṇa bhavitavyam . \\
\noindent \textbf{(P\_1,2.64.3)} na tarhi idānīm idam bhavati : vṛkṣaḥ ca vṛkṣaḥ ca vṛkṣau , vṛkṣaḥ ca vṛkṣaḥ ca vṛkṣaḥ ca vṛkṣāḥ iti . \\
\noindent \textbf{(P\_1,2.64.3)} na etat sahavivakṣāyām bhavati . \\
\noindent \textbf{(P\_1,2.64.3)} atha api nidarśayitum buddhiḥ evam nidarśayitavyam : vṛkṣau ca vṛkṣau ca vṛkṣau , vṛkṣāḥ ca vṛkṣāḥ ca vṛkṣāḥ ca vṛkṣāḥ iti . \\
\noindent \textbf{(P\_1,2.64.3)} yat api ucyate dvandvapratiṣedhaḥ ca vaktavyaḥ iti . \\
\noindent \textbf{(P\_1,2.64.3)} na eṣaḥ doṣaḥ . \\
\noindent \textbf{(P\_1,2.64.3)} anavakāśaḥ ekaśeṣḥ dvandvam bādhiṣyate . \\
\noindent \textbf{(P\_1,2.64.3)} nanu ca uktam sāvakāśaḥ ekaśeṣaḥ . \\
\noindent \textbf{(P\_1,2.64.3)} kaḥ avakāśaḥ . \\
\noindent \textbf{(P\_1,2.64.3)} tiṅantāni avakāśaḥ iti . \\
\noindent \textbf{(P\_1,2.64.3)} na tiṅantāni ekaśeṣārambham prayojayanti . \\
\noindent \textbf{(P\_1,2.64.3)} kim kāṛaṇam . \\
\noindent \textbf{(P\_1,2.64.3)} yathājātīyakānām dvitīyasya padasya prayoge sāmarthyam asti tathājātīyakānām ekaśeṣaḥ . \\
\noindent \textbf{(P\_1,2.64.3)} na ca tiṅantānām dvitīyasya padasya prayoge sāmarthyam asti . \\
\noindent \textbf{(P\_1,2.64.3)} kim kāraṇam . \\
\noindent \textbf{(P\_1,2.64.3)} ekā hi kriyā . \\
\noindent \textbf{(P\_1,2.64.3)} ekena uktatvāt tasya arthasya dvitīyasya prayogeṇa na bhavitavyam uktārthānām aprayogaḥ iti . \\
\noindent \textbf{(P\_1,2.64.3)} yadi tarhi ekā kriyā dvivacanabahuvacanāni na sidhyanti : pacataḥ pacanti . \\
\noindent \textbf{(P\_1,2.64.3)} na etāni kriyāpekṣāṇi . \\
\noindent \textbf{(P\_1,2.64.3)} kim tarhi . \\
\noindent \textbf{(P\_1,2.64.3)} sādhanāpekṣāṇi . \\
\noindent \textbf{(P\_1,2.64.3)} atha vā punaḥ astu ekavibhaktau iti . \\
\noindent \textbf{(P\_1,2.64.3)} nanu ca uktam ekavibhaktau iti cet na abhāvāt vibhakteḥ iti . \\
\noindent \textbf{(P\_1,2.64.3)} na eṣaḥ doṣaḥ . \\
\noindent \textbf{(P\_1,2.64.3)} parihṛtam etat : arthavat prātipadikam iti prātipadikasañjñā bhaviṣyati iti . \\
\noindent \textbf{(P\_1,2.64.3)} nanu ca uktam niyamāt na prāpnoti arthavatsamudāyānām samāsagrahaṇam niyamārtham iti . \\
\noindent \textbf{(P\_1,2.64.3)} na eṣaḥ doṣaḥ . \\
\noindent \textbf{(P\_1,2.64.3)} tulyajātīyasya niyamaḥ . \\
\noindent \textbf{(P\_1,2.64.3)} kaḥ ca tulyajātīyaḥ . \\
\noindent \textbf{(P\_1,2.64.3)} yathājātīyakānām samāsaḥ . \\
\noindent \textbf{(P\_1,2.64.3)} kathañjātīyakānām samāsaḥ . \\
\noindent \textbf{(P\_1,2.64.3)} subantānām \\
\noindent \textbf{(P\_1,2.64.4)} sarvatra apatyādiṣu upasaṅkhyānam . \\
\noindent \textbf{(P\_1,2.64.4)} sarveṣu pakṣeṣu apatyādiṣu upasaṅkhyānam kartavyam : bhikṣāṇām samūhaḥ bhaikṣam iti . \\
\noindent \textbf{(P\_1,2.64.4)} sarvatra iti ucyate prātipadikāṇām ca ekaśeṣe siddham . \\
\noindent \textbf{(P\_1,2.64.4)} apatyādiṣu iti ucyate bahavaḥ ca apatyādayaḥ : gargasya apatyam bahavaḥ gargāḥ . \\
\noindent \textbf{(P\_1,2.64.4)} ekā prakṛtiḥ bahavaḥ ca yañaḥ . \\
\noindent \textbf{(P\_1,2.64.4)} asārūpyāt ekaśeṣaḥ na prāpnoti . \\
\noindent \textbf{(P\_1,2.64.4)} nanu ca yathā eva bahavaḥ yañaḥ evam prakṛtayaḥ api bahvyaḥ syuḥ . \\
\noindent \textbf{(P\_1,2.64.4)} na evam śakyam . \\
\noindent \textbf{(P\_1,2.64.4)} iha hi doṣaḥ syāt : gargāḥ , vatsāḥ , bidāḥ , urvāḥ iti . \\
\noindent \textbf{(P\_1,2.64.4)} añ yaḥ bahuṣu yañ yaḥ bahuṣu iti ucyamānaḥ luk na prāpnoti . \\
\noindent \textbf{(P\_1,2.64.4)} mā bhūt evam . \\
\noindent \textbf{(P\_1,2.64.4)} añantam yat bahuṣu yañantam yat bahuṣu iti evam bhaviṣyati . \\
\noindent \textbf{(P\_1,2.64.4)} nanu ca uktam : na evam śakyam . \\
\noindent \textbf{(P\_1,2.64.4)} iha hi doṣaḥ syāt : kāśyapapratikṛtayaḥ kāśyapāḥ iti . \\
\noindent \textbf{(P\_1,2.64.4)} na eṣaḥ doṣaḥ . \\
\noindent \textbf{(P\_1,2.64.4)} laukikasya tatra gotrasya grahaṇam na ca etat laukikam gotram . \\
\noindent \textbf{(P\_1,2.64.4)} atha vā punaḥ astu ekā prakṛtiḥ bahavaḥ ca yañaḥ . \\
\noindent \textbf{(P\_1,2.64.4)} nanu ca uktam : asārūpyāt ekaśeṣaḥ na prāpnoti iti . \\
\noindent \textbf{(P\_1,2.64.4)} siddham tu samānārthānām ekaśeṣavacanāt . \\
\noindent \textbf{(P\_1,2.64.4)} siddham etat . \\
\noindent \textbf{(P\_1,2.64.4)} katham . \\
\noindent \textbf{(P\_1,2.64.4)} samānārthānām ekaśeṣaḥ bhavati iti vaktavyam . \\
\noindent \textbf{(P\_1,2.64.4)} yadi samānārthānām ekaśeṣaḥ ucyate katham akṣāḥ , pādāḥ , māṣāḥ iti . \\
\noindent \textbf{(P\_1,2.64.4)} nānārthānām api sarūpāṇām . \\
\noindent \textbf{(P\_1,2.64.4)} nānārthānām api sarūpāṇām ekaśeṣḥ vaktavyaḥ . \\
\noindent \textbf{(P\_1,2.64.4)} ekārthānām api virūpāṇām . \\
\noindent \textbf{(P\_1,2.64.4)} ekārthānām api virūpāṇām ekaśeṣaḥ vaktavyaḥ : vakradaṇḍaḥ ca kuṭiladaṇḍaḥ ca vakradaṇḍau kuṭiladaṇḍāu iti vā . \\
\noindent \textbf{(P\_1,2.64.4)} svarabhinnānām yasya uttarasvaravidhiḥ . \\
\noindent \textbf{(P\_1,2.64.4)} svarabhinnānām yasya uttarasvaravidhiḥ tasya ekaśeṣaḥ vaktavyaḥ . \\
\noindent \textbf{(P\_1,2.64.4)} akṣaḥ ca akṣaḥ ca akṣau , mīmaṃsakaḥ ca mīmāṃsakaḥ ca mīmaṃsakau . \\
\noindent \textbf{(P\_1,2.64.5)} iha kasmāt na bhavati : ekaḥ ca ekaḥ ca , dvau ca dvau ca iti . \\
\noindent \textbf{(P\_1,2.64.5)} saṅkhyāyāḥ arthāsampratyayāt anyapadārthatvāt ca anekaśeṣaḥ . \\
\noindent \textbf{(P\_1,2.64.5)} saṅkhyāyāḥ arthāsampratyayāt ekaśeṣaḥ na bhaviṣyati . \\
\noindent \textbf{(P\_1,2.64.5)} na hi ekau iti anena arthaḥ gamyate . \\
\noindent \textbf{(P\_1,2.64.5)} anyapadārthatvāt ca saṅkhyāyāḥ ekaśeṣaḥ na bhaviṣyati . \\
\noindent \textbf{(P\_1,2.64.5)} ekaḥ ca ekaḥ ca iti asya dvau iti arthaḥ . \\
\noindent \textbf{(P\_1,2.64.5)} dvau ca dvau ca iti asya catvāraḥ iti arthaḥ . \\
\noindent \textbf{(P\_1,2.64.5)} na etau staḥ parihārau . \\
\noindent \textbf{(P\_1,2.64.5)} yat tāvat ucyate saṅkhyāyāḥ arthāsampratyayāt iti . \\
\noindent \textbf{(P\_1,2.64.5)} arthāsampratyaye api ekaśeṣaḥ bhavati . \\
\noindent \textbf{(P\_1,2.64.5)} tat yathā . \\
\noindent \textbf{(P\_1,2.64.5)} gārgyaḥ ca gārgyāyaṇaḥ ca gārgyau . \\
\noindent \textbf{(P\_1,2.64.5)} na ca ucyate vṛddhayuvānau iti bhavati ca ekaśeṣaḥ . \\
\noindent \textbf{(P\_1,2.64.5)} yat api ucyate : anyapadārthatvāt ca iti . \\
\noindent \textbf{(P\_1,2.64.5)} anyapadārthe api ekaśeṣaḥ bhavati . \\
\noindent \textbf{(P\_1,2.64.5)} tat yathā : viṃśatiḥ ca viṃśatiḥ ca viṃśatī iti . \\
\noindent \textbf{(P\_1,2.64.5)} tayoḥ catvāriṃśat iti arthaḥ . \\
\noindent \textbf{(P\_1,2.64.5)} evam tarhi na imau pṛthak parihārau . \\
\noindent \textbf{(P\_1,2.64.5)} ekaparihāraḥ ayam : saṅkhyāyāḥ arthāsampratyayāt anyapadārthatvāt ca iti . \\
\noindent \textbf{(P\_1,2.64.5)} yatra hi arthāsampratyayaḥ eva vā anyapadārthatā eva vā bhavati tatra ekaśeṣaḥ gārgyau viṃśatī iti yathā . \\
\noindent \textbf{(P\_1,2.64.5)} atha vā na ime ekaśeṣaśabdāḥ . \\
\noindent \textbf{(P\_1,2.64.5)} yadi tarhi na ime ekaśeṣaśabdāḥ samudāyaśabdāḥ tarhi bhavanti . \\
\noindent \textbf{(P\_1,2.64.5)} tatra kaḥ doṣaḥ . \\
\noindent \textbf{(P\_1,2.64.5)} ekavacanam prāpnoti . \\
\noindent \textbf{(P\_1,2.64.5)} ekārthāḥ hi samudāyāḥ bhavanti . \\
\noindent \textbf{(P\_1,2.64.5)} tat yathā yūtham , śatam , vanam iti . \\
\noindent \textbf{(P\_1,2.64.5)} santu tarhi ekaśeṣaśabdāḥ . \\
\noindent \textbf{(P\_1,2.64.5)} kiṅkṛtam sārūpyam . \\
\noindent \textbf{(P\_1,2.64.5)} anyonyakṛtam sārūpyam . \\
\noindent \textbf{(P\_1,2.64.5)} santi punaḥ ke cit anye api śabdāḥ yeṣām anyonyakṛtaḥ bhāvaḥ . \\
\noindent \textbf{(P\_1,2.64.5)} santi iti āha . \\
\noindent \textbf{(P\_1,2.64.5)} tad yathā mātā pitā bhrātā iti . \\
\noindent \textbf{(P\_1,2.64.5)} viṣamaḥ upanyāsaḥ . \\
\noindent \textbf{(P\_1,2.64.5)} sakṛt ete śabdāḥ pravṛttāḥ apāyeṣu api vartante . \\
\noindent \textbf{(P\_1,2.64.5)} iha punaḥ ekena api apāye na bhavati catvāraḥ iti . \\
\noindent \textbf{(P\_1,2.64.5)} anyat idānīm etat ucyate sakṛt ete śabdāḥ pravṛttāḥ apāyeṣu api vartante iti . \\
\noindent \textbf{(P\_1,2.64.5)} yat tu bhavān asmān codayati santi punaḥ ke cit anye api śabdāḥ yeṣām anyonyakṛtḥ bhāvaḥ iti tatra ete asmābhiḥ upanyastāḥ . \\
\noindent \textbf{(P\_1,2.64.5)} tatra etat bhavān āha sakṛt ete śabdāḥ pravṛttāḥ apāyeṣu api vartante iti . \\
\noindent \textbf{(P\_1,2.64.5)} etat ca vārttam . \\
\noindent \textbf{(P\_1,2.64.5)} ekaikaḥ na udyantum bhāram śaknoti yat katham tatra | ekaikaḥ kartā syāt sarve vā syuḥ katham yuktam || kāraṇam udyamanam cet na udyacchati ca antareṇa tat tulyam | tasmāt pṛthak pṛthak te kartāraḥ savyapekṣāḥ tu || \\
\noindent \textbf{(P\_1,2.64.6)} prathamamadhyamottamānām ekaśeṣaḥ sarūpatvāt . \\
\noindent \textbf{(P\_1,2.64.6)} prathamamadhyamottamānām ekaśeṣaḥ vaktavyaḥ : pacati ca pacasi ca pacathaḥ , pacasi ca pacāmi ca pacāvaḥ , pacati ca pacasi ca pacāmi ca pacāmaḥ . \\
\noindent \textbf{(P\_1,2.64.6)} kim punaḥ kāraṇam na sidhyati . \\
\noindent \textbf{(P\_1,2.64.6)} asarūpatvāt . \\
\noindent \textbf{(P\_1,2.64.7)} dvivacanabahuvacanāprasiddhiḥ ca ekārthatvāt . dvivacanabahuvacanayoḥ ca aprasiddhiḥ . \\
\noindent \textbf{(P\_1,2.64.7)} kim kāraṇam . \\
\noindent \textbf{(P\_1,2.64.7)} ekārthatvāt . \\
\noindent \textbf{(P\_1,2.64.7)} ekaḥ ayam avaśiṣyate . \\
\noindent \textbf{(P\_1,2.64.7)} tena anena tadarthena bhavitavyam . \\
\noindent \textbf{(P\_1,2.64.7)} kimarthena . \\
\noindent \textbf{(P\_1,2.64.7)} yadarthaḥ ekaḥ . \\
\noindent \textbf{(P\_1,2.64.7)} kimarthaḥ ca ekaḥ . \\
\noindent \textbf{(P\_1,2.64.7)} ekaḥ ekārthaḥ . \\
\noindent \textbf{(P\_1,2.64.7)} na aikārthyam . \\
\noindent \textbf{(P\_1,2.64.7)} na ayam ekārthaḥ . \\
\noindent \textbf{(P\_1,2.64.7)} kim tarhi . \\
\noindent \textbf{(P\_1,2.64.7)} dvyarthaḥ bahvarthaḥ ca . \\
\noindent \textbf{(P\_1,2.64.7)} na aikārthyam iti cet ārambhānarthakyam . \\
\noindent \textbf{(P\_1,2.64.7)} na aikārthyam iti cet ekaśeṣārambhaḥ anarthakaḥ syāt . \\
\noindent \textbf{(P\_1,2.64.7)} iha hi śabdasya svābhāvikī vā anekārthatā syāt vācanikī vā . \\
\noindent \textbf{(P\_1,2.64.7)} tat yadi tāvat svābhāvikī aśiṣyaḥ ekaśeṣaḥ ekena uktatvāt . aśiṣyaḥ ekaśeṣaḥ . \\
\noindent \textbf{(P\_1,2.64.7)} kim kāraṇam . \\
\noindent \textbf{(P\_1,2.64.7)} ekena uktatvāt tasya arthasya dvitīyasya prayogeṇa na bhavitavyam uktārthānām aprayogaḥ iti . \\
\noindent \textbf{(P\_1,2.64.7)} atha vācanikī tat vaktavyam : ekaḥ ayam aviśiṣyate saḥ ca dvyarthaḥ bhavati bahvarthaḥ ca iti . \\
\noindent \textbf{(P\_1,2.64.7)} na vaktavyam . \\
\noindent \textbf{(P\_1,2.64.7)} siddham ekaśeṣaḥ iti eva . \\
\noindent \textbf{(P\_1,2.64.7)} katham punaḥ ekaḥ ayam aviśiṣyate iti anena dvyarthatā bahvarthatā vā śakyā labdhum . \\
\noindent \textbf{(P\_1,2.64.7)} tat ca ekaśeṣakṛtam . \\
\noindent \textbf{(P\_1,2.64.7)} na hi antareṇa tadvācinaḥ śabdasya prayogam tasya arthasya gatiḥ bhavati . \\
\noindent \textbf{(P\_1,2.64.7)} paśyāmaḥ ca punaḥ antareṇa api tadvācinaḥ śabdasya prayogam tasya arthasya gatiḥ bhavati iti agnicit somasut iti yathā . \\
\noindent \textbf{(P\_1,2.64.7)} te manyāmahe : lopakṛtam etat yena atra antareṇa api tadvācinaḥ śabdasya prayogam tasya arthasya gatiḥ bhavatiti . \\
\noindent \textbf{(P\_1,2.64.7)} evam iha api ekaśeṣakṛtam etat yena atra ekaḥ ayam avaśiṣyate iti anena dvyarthatā bahvarthatā vā bhavati . \\
\noindent \textbf{(P\_1,2.64.7)} ucyeta tarhi na tu gamyeta . \\
\noindent \textbf{(P\_1,2.64.7)} yaḥ hi gām aśvaḥ iti brūyāt aśvam vā gauḥ iti na jātu cit sampratyayaḥ syāt . \\
\noindent \textbf{(P\_1,2.64.7)} tena anekārthābhidhāne yatnam kurvatā avaśyam lokaḥ pṛṣṭhataḥ anugantavyaḥ : keṣu artheṣu laukikāḥ kān śabdān prayuñjate iti . \\
\noindent \textbf{(P\_1,2.64.7)} loke ca ekasmin vṛkṣaḥ iti prayuñjate dvayoḥ vṛkṣau iti bahuṣu vṛkṣāḥ iti . \\
\noindent \textbf{(P\_1,2.64.7)} yadi tarhi lokaḥ avaśyam śabdeṣu pramāṇam kimartham ekaśeṣaḥ ārabhyate . \\
\noindent \textbf{(P\_1,2.64.7)} atha kimartham lopaḥ ārabhyate . \\
\noindent \textbf{(P\_1,2.64.7)} pratyayalakṣaṇam ācāryaḥ prārthayamānaḥ lopam ārabhate . \\
\noindent \textbf{(P\_1,2.64.7)} ekaśeṣārambhe punaḥ asya na kim cit prayojanam asti . \\
\noindent \textbf{(P\_1,2.64.7)} nanu ca uktam : pratyartham śabdaniveśāt na ekena anekasya abhidhānam iti . \\
\noindent \textbf{(P\_1,2.64.7)} yadi ca ekena śabdena anekasya arthasya abhidhānam syāt na pratyartham śabdaniveśaḥ kṛtaḥ syāt . \\
\noindent \textbf{(P\_1,2.64.7)} pratyartham śabdaniveśāt ekena anekasya abhidhānāt apratyartham iti cet tat api pratyartham eva . \\
\noindent \textbf{(P\_1,2.64.7)} pratyartham śabdaniveśāt ekena anekasyābhidhānāt apratyartham iti cet evam ucyate : yat api ekena anekasya abhidhānam bhavati tat api pratyartham eva . \\
\noindent \textbf{(P\_1,2.64.7)} yat api hi arthau arthau prati tat api pratyartham eva . \\
\noindent \textbf{(P\_1,2.64.7)} yat api hi arthān arthān prati tat api pratyartham eva . \\
\noindent \textbf{(P\_1,2.64.7)} yāvatām abhidhānam tāvatām prayogaḥ nyāyyaḥ . \\
\noindent \textbf{(P\_1,2.64.7)} yāvatām arthānām abhidhānam bhavati tāvatām śabdānām prayogaḥ iti eṣaḥ pakṣaḥ nyāyyaḥ . \\
\noindent \textbf{(P\_1,2.64.7)} yāvatām abhidhānam tāvatām prayogaḥ nyāyyaḥ iti cet ekena api anekasya abhidhānam . \\
\noindent \textbf{(P\_1,2.64.7)} yāvatām abhidhānam tāvatām prayogaḥ nyāyyaḥ iti cet evam ucyate : eṣaḥ api nyāyyaḥ eva yat api ekena api anekasya abhidhānam bhavati . \\
\noindent \textbf{(P\_1,2.64.7)} yadi tarhi ekena anekasya abhidhānam bhavati plakṣanyagrodhau : ekena uktatvāt aparasya prayogaḥ anupapannaḥ . \\
\noindent \textbf{(P\_1,2.64.7)} ekena uktatvāt tasya arthasya aparasya prayogeṇa na bhavitavyam . \\
\noindent \textbf{(P\_1,2.64.7)} kim kāraṇam . \\
\noindent \textbf{(P\_1,2.64.7)} uktārthānām aprayogaḥ iti . \\
\noindent \textbf{(P\_1,2.64.7)} ekena uktatvāt aparasya prayogaḥ anupapannaḥ iti cet anuktatvāt plakṣeṇa nyagrodhasya nyagrodhaprayogaḥ . ekena uktatvāt aparasya prayogaḥ anupapannaḥ iti cet anuktaḥ plakṣeṇa nyagrodhārthaḥ iti kṛtvā nyagrodhaśabdaḥ prayujyate . \\
\noindent \textbf{(P\_1,2.64.7)} katham anuktaḥ yāvatā idānīm eva uktam ekena api anekasya abhidhānam bhavati iti . \\
\noindent \textbf{(P\_1,2.64.7)} sarūpāṇām ekena api anekasya abhidhānam bhavati na virūpāṇām . \\
\noindent \textbf{(P\_1,2.64.7)} kim punaḥ kāraṇam sarūpāṇām ekena api anekasya abhidhānam bhavati na punaḥ virūpāṇām . \\
\noindent \textbf{(P\_1,2.64.7)} abhidhānam punaḥ svābhāvikam . \\
\noindent \textbf{(P\_1,2.64.7)} svābhāvikam abhidhānam . \\
\noindent \textbf{(P\_1,2.64.7)} ubhayadarśanāt ca . \\
\noindent \textbf{(P\_1,2.64.7)} ubhayam khalu api dṛśyate : virūpāṇām api ekena anekasya abhidhānam bhavati . \\
\noindent \textbf{(P\_1,2.64.7)} tat yathā : dyavā ha kṣamā . \\
\noindent \textbf{(P\_1,2.64.7)} dyavā cit asmai pṛthivī namete iti . \\
\noindent \textbf{(P\_1,2.64.7)} virūpāṇām kila nāma ekena anekasya abhidhānam syāt kim punaḥ sarūpāṇām . \\
\noindent \textbf{(P\_1,2.64.8)} ākṛtyabhidhānāt vā ekam vibhaktau vājapyāyanaḥ . \\
\noindent \textbf{(P\_1,2.64.8)} ākṛtyabhidhānāt vā ekam śabdam vibhaktau vājapyāyanaḥ ācāryaḥ nyāyyam manyate : ekā ākṛtiḥ sā ca abhidhīyate iti . \\
\noindent \textbf{(P\_1,2.64.8)} katham punaḥ jñāyate ekā ākṛtiḥ sā ca abhidhīyate iti . \\
\noindent \textbf{(P\_1,2.64.8)} prakhyāviśeṣāt . na hi gauḥ iti ukte viśeṣaḥ prakhyāyate śuklā nīlā kapilā kapotikā iti . \\
\noindent \textbf{(P\_1,2.64.8)} yadi api tāvat prakhyāviśeṣāt jñāyate ekā ākṛtiḥ iti kutaḥ tu etat sā abhidhīyate iti . \\
\noindent \textbf{(P\_1,2.64.8)} avyapavargagateḥ ca . \\
\noindent \textbf{(P\_1,2.64.8)} avyapavargagateḥ ca manyāmahe ākṛtiḥ abhidhīyate iti . \\
\noindent \textbf{(P\_1,2.64.8)} na hi gauḥ iti ukte vyapavargaḥ gamyate śuklā nīlā kapilā kapotikā iti . \\
\noindent \textbf{(P\_1,2.64.8)} jñāyate ca ekopadiṣṭam . jñāyate khalu api ekopadiṣṭam . \\
\noindent \textbf{(P\_1,2.64.8)} gauḥ asya kadā cit upadiṣṭaḥ bhavati . \\
\noindent \textbf{(P\_1,2.64.8)} saḥ tam anyasmin deśe anyasmin kāle anyasyām ca vayovasthāyām dṛṣṭvā jānāti ayam gauḥ iti . \\
\noindent \textbf{(P\_1,2.64.8)} kaḥ punaḥ asya viśeṣaḥ prakhyāviśeṣāt iti ataḥ . \\
\noindent \textbf{(P\_1,2.64.8)} tasya eva upodbalakam etat : prakhyāviśeṣāt jñāyate ca ekopadiṣṭam iti . \\
\noindent \textbf{(P\_1,2.64.8)} dharmaśāstram ca tathā . evam ca kṛtvā dharmaśāstram pravṛttam : brāhmaṇaḥ na hantavyaḥ . \\
\noindent \textbf{(P\_1,2.64.8)} surā na peyā iti . \\
\noindent \textbf{(P\_1,2.64.8)} brāhmaṇamātram na hanyate surāmātram ca na pīyate . \\
\noindent \textbf{(P\_1,2.64.8)} yadi dravyam padārthaḥ syāt ekam brāhmaṇam ahatvā ekām ca surām apītvā anyatra kāmacāraḥ syāt . \\
\noindent \textbf{(P\_1,2.64.8)} kaḥ punaḥ asya viśeṣaḥ avyapavargagateḥ ca iti ataḥ . \\
\noindent \textbf{(P\_1,2.64.8)} tasya eva upodbalakam etat : avyapavargagateḥ ca dharmaśāstram ca tathā iti . \\
\noindent \textbf{(P\_1,2.64.8)} asti ca ekam anekādhikaraṇastham yugapat . \\
\noindent \textbf{(P\_1,2.64.8)} asti khalu api ekam anekādhikaraṇastham yugapat upalabhyate . \\
\noindent \textbf{(P\_1,2.64.8)} kim . \\
\noindent \textbf{(P\_1,2.64.8)} ādityaḥ . \\
\noindent \textbf{(P\_1,2.64.8)} tad yathā ekaḥ ādityaḥ anekādhikaraṇasthaḥ yugapat upalabhyate . \\
\noindent \textbf{(P\_1,2.64.8)} viṣamaḥ upanyāsaḥ . \\
\noindent \textbf{(P\_1,2.64.8)} na ekaḥ draṣṭā ādityam anekādhikaraṇastham yugapat upalabhate . \\
\noindent \textbf{(P\_1,2.64.8)} evam tarhi itīndravat viṣayaḥ . tat yathā ekaḥ indraḥ anekasmin kratuśate āhūtaḥ yugapat sarvatra bhavati evam ākṛtiḥ api yugapat sarvatra bhaviṣyati . \\
\noindent \textbf{(P\_1,2.64.8)} avaśyam ca etat evam vijñeyam ekam anekādhikaraṇastham yugapat upalabhyate iti. na ekam anekādhikaraṇastham yugapat iti cet tathā ekaśeṣe . \\
\noindent \textbf{(P\_1,2.64.8)} yaḥ hi manyate na ekam anekādhikaraṇastham yugapad upalabhyate iti ekaśeṣe tasya doṣaḥ syāt . \\
\noindent \textbf{(P\_1,2.64.8)} ekaśeṣe api na ekaḥ vṛkṣaśabdaḥ anekam artham yugapat abhidadhīta . \\
\noindent \textbf{(P\_1,2.64.8)} avaśyam ca etat evam vijñeyam ākṛtiḥ abhidhīyate iti . \\
\noindent \textbf{(P\_1,2.64.8)} dravyābhidhāne hi ākṛtyasampratyayaḥ . \\
\noindent \textbf{(P\_1,2.64.8)} dravyābhidhāne sati ākṛteḥ asampratyayaḥ syāt . \\
\noindent \textbf{(P\_1,2.64.8)} tatra kaḥ doṣaḥ . \\
\noindent \textbf{(P\_1,2.64.8)} tatra asarvadravyagatiḥ . \\
\noindent \textbf{(P\_1,2.64.8)} tatra asarvadravyagatiḥ prāpnoti . \\
\noindent \textbf{(P\_1,2.64.8)} asarvadravyagatau kaḥ doṣaḥ . \\
\noindent \textbf{(P\_1,2.64.8)} gauḥ anubandhyaḥ ajaḥ agnīṣomīyaḥ iti : ekaḥ śāstroktam kurvīta aparaḥ aśāstroktam . \\
\noindent \textbf{(P\_1,2.64.8)} aśāstrokte ca kriyamāṇe viguṇam karma bhavati . \\
\noindent \textbf{(P\_1,2.64.8)} viguṇe ca karmaṇi phalānavāptiḥ . \\
\noindent \textbf{(P\_1,2.64.8)} nanu ca yasya api ākṛtiḥ padārthaḥ tasya api yadi anavayavena codyate na ca anubadhyate viguṇam karma bhavati . \\
\noindent \textbf{(P\_1,2.64.8)} viguṇe ca karmaṇi phalānavāptiḥ . \\
\noindent \textbf{(P\_1,2.64.8)} ekā ākṛtiḥ iti ca pratijñā hīyeta . \\
\noindent \textbf{(P\_1,2.64.8)} yat ca asya pakṣasya upādāne prayojanam ekaśeṣaḥ na vaktavyaḥ iti saḥ ca idānīm vaktavyaḥ bhavati . \\
\noindent \textbf{(P\_1,2.64.8)} evam tarhi anavayavena codyate pratyekam ca parisamāpyate yathā ādityaḥ . \\
\noindent \textbf{(P\_1,2.64.8)} nanu ca yasya api dravyam padārthaḥ tasya api anavayavena codyate pratyekam ca parisamāpyate . \\
\noindent \textbf{(P\_1,2.64.8)} ekaśeṣaḥ tvayā vaktavyaḥ . \\
\noindent \textbf{(P\_1,2.64.8)} tvayā api tarhi dvivacanabahuvacanāni sādhyāni . \\
\noindent \textbf{(P\_1,2.64.8)} codanāyām ca ekasya upādhivṛtteḥ . \\
\noindent \textbf{(P\_1,2.64.8)} codanāyām ca ekasya upādhivṛtteḥ manyāmahe ākṛtiḥ abhidhīyate iti . \\
\noindent \textbf{(P\_1,2.64.8)} āgneyam aṣṭākapālam nirvapet : ekam nirupya dvitīyas tṛtīyaḥ ca nirupyate . \\
\noindent \textbf{(P\_1,2.64.8)} yadi ca dravyam padārthaḥ syāt ekam nirupya dvitīyasya tṛtīyasya ca nirvapaṇam na prakalpeta . \\
\noindent \textbf{(P\_1,2.64.8)} kaḥ punaḥ etayoḥ jāticodanayoḥ viśeṣaḥ . \\
\noindent \textbf{(P\_1,2.64.8)} ekā nirvṛttena aparā nirvartyena . \\
\noindent \textbf{(P\_1,2.64.9)} dravyābhidhānam vyāḍiḥ . dravyābhidhānam vyāḍiḥ ācāryaḥ nyāyyam manyate : dravyam abhidhīyate iti . \\
\noindent \textbf{(P\_1,2.64.9)} tathā ca liṅgavacanasiddhiḥ . \\
\noindent \textbf{(P\_1,2.64.9)} evam ca kṛtvā liṅgavacanāni siddhāni bhavanti : brāhmaṇī brāhmaṇaḥ , brāhmaṇau brāhmaṇāḥ iti . \\
\noindent \textbf{(P\_1,2.64.9)} codanāsu ca tasya ārambhāt . \\
\noindent \textbf{(P\_1,2.64.9)} codanāsu ca tasya ārambhāt manyāmahe dravyam abhidhīyate iti . \\
\noindent \textbf{(P\_1,2.64.9)} gauḥ anubandhyaḥ ajaḥ agnīṣomīyaḥ iti : ākṛtau coditāyām dravye ārambhaṇālambhanaprokṣaṇaviśasanādīni kriyante . \\
\noindent \textbf{(P\_1,2.64.9)} na ca ekam anekādhikaraṇastham yugapat . \\
\noindent \textbf{(P\_1,2.64.9)} na khalu api ekam anekādhikaraṇastham yugapat upalabhyate . \\
\noindent \textbf{(P\_1,2.64.9)} na hi ekaḥ devadattaḥ yugapat srughne bhavati mathurāyām ca . \\
\noindent \textbf{(P\_1,2.64.9)} vināśe prādurbhāve ca sarvam tathā syāt . \\
\noindent \textbf{(P\_1,2.64.9)} kim . \\
\noindent \textbf{(P\_1,2.64.9)} vinaśyet ca prāduḥ ṣyāt ca . \\
\noindent \textbf{(P\_1,2.64.9)} śvā mṛtaḥ iti śvā nāma loke na pracaret . \\
\noindent \textbf{(P\_1,2.64.9)} gauḥ jātaḥ iti sarvam gobhūtam anavakāśam syāt . \\
\noindent \textbf{(P\_1,2.64.9)} asti ca vairūpyam . \\
\noindent \textbf{(P\_1,2.64.9)} asti khalu api vairūpyam : gauḥ ca gauḥ ca khaṇḍaḥ muṇḍaḥ iti . \\
\noindent \textbf{(P\_1,2.64.9)} tathā ca vigrahaḥ . \\
\noindent \textbf{(P\_1,2.64.9)} evam ca kṛtvā vigrahaḥ upapannaḥ bhavati : gauḥ ca gauḥ ca iti . \\
\noindent \textbf{(P\_1,2.64.9)} vyartheṣu ca muktasaṃśayam . \\
\noindent \textbf{(P\_1,2.64.9)} vyartheṣu ca muktasaṃśayam bhavati . \\
\noindent \textbf{(P\_1,2.64.9)} ākṛtau api padārthe ekaśeṣaḥ vaktavyaḥ : akṣāḥ , pādāḥ , māṣāḥ iti . \\
\noindent \textbf{(P\_1,2.64.10)} liṅgavacanasiddhiḥ guṇasya anityatvāt . \\
\noindent \textbf{(P\_1,2.64.10)} liṅgavacanāni siddhāni bhavanti . \\
\noindent \textbf{(P\_1,2.64.10)} kutaḥ . \\
\noindent \textbf{(P\_1,2.64.10)} guṇasya anityatvāt . \\
\noindent \textbf{(P\_1,2.64.10)} anityāḥ guṇāḥ apāyinaḥ upāyinaḥ ca . \\
\noindent \textbf{(P\_1,2.64.10)} kim ye ete śuklādayaḥ . \\
\noindent \textbf{(P\_1,2.64.10)} na iti āha . \\
\noindent \textbf{(P\_1,2.64.10)} strīpuṃnapuṃsakāni sattvaguṇāḥ ekatvadvitvabahutvāni ca . \\
\noindent \textbf{(P\_1,2.64.10)} kadā cit ākṛtiḥ ekatvena yujyate kadā cit dvitvena kadā cit bahutvena kadā cit strītvena kadā cit puṃstvena kadācit napuṃsakatvena . \\
\noindent \textbf{(P\_1,2.64.10)} bhavet liṅgaparihāraḥ upapannaḥ vacanaparihāraḥ tu na upapadyate . \\
\noindent \textbf{(P\_1,2.64.10)} yadi hi kadā cit ākṛtiḥ ekatvena yujyate kadā cit dvitvena kadā cit bahutvena ekā ākṛtiḥ iti pratijñā hīyeta . \\
\noindent \textbf{(P\_1,2.64.10)} yat ca asya pakṣasya upādāne prayojanam uktam ekaśeṣaḥ na vaktavyaḥ iti saḥ ca idānīm vaktavyaḥ bhavati . \\
\noindent \textbf{(P\_1,2.64.10)} evam tarhi liṅgavacanasiddhiḥ guṇavivakṣānityatvāt . \\
\noindent \textbf{(P\_1,2.64.10)} liṅgavacanāni siddhāni bhavanti . \\
\noindent \textbf{(P\_1,2.64.10)} kutaḥ . \\
\noindent \textbf{(P\_1,2.64.10)} guṇavivakṣāyāḥ anityatvāt . \\
\noindent \textbf{(P\_1,2.64.10)} anityā guṇavivakṣā . \\
\noindent \textbf{(P\_1,2.64.10)} kadā cit ākṛtiḥ ekatvena vivakṣitā bhavati kadā cit dvitvena kadā cit bahutvena kadā cit strītvena kadā cit puṃstvena kadā cit napuṃsakatvena . \\
\noindent \textbf{(P\_1,2.64.10)} bhavet liṅgaparihāraḥ upapannaḥ vacanaparihāraḥ tu na upapadyate . \\
\noindent \textbf{(P\_1,2.64.10)} yadi kadā cit ākṛtiḥ ekatvena vivakṣitā bhavati kadā cit dvitvena kadā cit bahutvena ekā ākṛtiḥ iti pratijñā hīyeta . \\
\noindent \textbf{(P\_1,2.64.10)} yat ca asya pakṣasya upādāne prayojanam uktam ekaśeṣaḥ na vaktavyaḥ iti saḥ ca idānīm vaktavyaḥ bhavati . \\
\noindent \textbf{(P\_1,2.64.10)} liṅgaparihāraḥ ca api na upapadyate . \\
\noindent \textbf{(P\_1,2.64.10)} kim kāraṇam . \\
\noindent \textbf{(P\_1,2.64.10)} āviṣṭaliṅgā jātiḥ yat liṅgam upādāya pravartate utpattiprabhṛti ā vināśāt tat liṅgam na jahāti . \\
\noindent \textbf{(P\_1,2.64.10)} tasmāt na vaiyākaraṇaiḥ śakyam laukikam liṅgam āsthātum . \\
\noindent \textbf{(P\_1,2.64.10)} avaśyam kaḥ cit svakṛtāntaḥ āstheyaḥ . \\
\noindent \textbf{(P\_1,2.64.10)} kaḥ asau svakṛtāntaḥ . \\
\noindent \textbf{(P\_1,2.64.10)} saṃstyānaprasavau liṅgam . \\
\noindent \textbf{(P\_1,2.64.10)} saṃstyānaprasavau liṅgam āstheyau . \\
\noindent \textbf{(P\_1,2.64.10)} kim idam saṃstyānaprasavau iti . \\
\noindent \textbf{(P\_1,2.64.10)} saṃstyāne styāyateḥ ḍraṭ : strī . sūteḥ sap prasave pumān . \\
\noindent \textbf{(P\_1,2.64.10)} nanu ca loke api styāyateḥ eva strī sūteḥ ca pumān . \\
\noindent \textbf{(P\_1,2.64.10)} adhikaraṇasādhanā loke strī : styāyati asyām garbhaḥ iti . \\
\noindent \textbf{(P\_1,2.64.10)} kartṛsādhanaḥ ca pumān : sūte pumān iti . \\
\noindent \textbf{(P\_1,2.64.10)} iha punaḥ ubhayam bhāvasādhanam : styānam pravṛttiḥ ca . \\
\noindent \textbf{(P\_1,2.64.10)} kasya punaḥ styānam strī pravṛttiḥ vā pumān . \\
\noindent \textbf{(P\_1,2.64.10)} guṇānām . \\
\noindent \textbf{(P\_1,2.64.10)} keṣām . \\
\noindent \textbf{(P\_1,2.64.10)} śabdasparśarūparasagandhānām . \\
\noindent \textbf{(P\_1,2.64.10)} sarvāḥ ca punaḥ mūrtayaḥ evamātmikāḥ saṃstyānaprasavaguṇāḥ śabdasparśarūparasagandhavatyaḥ . \\
\noindent \textbf{(P\_1,2.64.10)} yatra alpīyāṃsaḥ guṇāḥ tatra avarataḥ trayaḥ : śabdaḥ sparśaḥ rūpam iti . \\
\noindent \textbf{(P\_1,2.64.10)} rasagandhau na sarvatra . \\
\noindent \textbf{(P\_1,2.64.10)} pravṛttiḥ khalu api nityā . \\
\noindent \textbf{(P\_1,2.64.10)} na hi iha kaḥ cit api svasmin ātmani muhūrtam api avatiṣṭhate . \\
\noindent \textbf{(P\_1,2.64.10)} vardhate yāvat anena vardhitavyam apacayena vā yujyate . \\
\noindent \textbf{(P\_1,2.64.10)} tat ca ubhayam sarvatra . \\
\noindent \textbf{(P\_1,2.64.10)} yadi ubhayam sarvatra kutaḥ vyavasthā . \\
\noindent \textbf{(P\_1,2.64.10)} vivakṣātaḥ. saṃstyānavivakṣāyām strī prasavavivakṣāyām pumān ubhayoḥ api avivakṣāyām napuṃsakam . \\
\noindent \textbf{(P\_1,2.64.10)} tatra liṅgavacanasiddhiḥ guṇavivakṣānityatvāt iti liṅgaparihāraḥ upapannaḥ . \\
\noindent \textbf{(P\_1,2.64.10)} vacanaparihāraḥ tu na upapadyate . \\
\noindent \textbf{(P\_1,2.64.10)} vacanaparihāraḥ ca api upapannaḥ . \\
\noindent \textbf{(P\_1,2.64.10)} idam tāvat ayam praṣṭavyaḥ : atha yasya dravyam padārthaḥ katham tasya ekavacanadvivacanabahuvacanāni bhavanti iti . \\
\noindent \textbf{(P\_1,2.64.10)} evam saḥ vakṣyati : ekasmin ekavacanam dvayoḥ dvivacanam bahuṣu bahuvacanam iti . \\
\noindent \textbf{(P\_1,2.64.10)} yadi tasya api vācanikāni na svābhāvikāni aham api evam vakṣyāmi : ekasmin ekavacanam dvayoḥ dvivacanam bahuṣu bahuvacanam iti . \\
\noindent \textbf{(P\_1,2.64.10)} na hi ākṛtipadārthikasya dravyam na padārthaḥ dvavyapadārthikasya vā ākṛtiḥ na padārthaḥ . \\
\noindent \textbf{(P\_1,2.64.10)} ubhayoḥ ubhayam padārthaḥ . \\
\noindent \textbf{(P\_1,2.64.10)} kasya cit tu kim cit pradhānabhūtam kim cit guṇabhūtam . \\
\noindent \textbf{(P\_1,2.64.10)} ākṛtipadārthikasya ākṛtiḥ pradhānabhūtā dravyam guṇabhūtam . \\
\noindent \textbf{(P\_1,2.64.10)} dravyapadārthikasya dravyam pradhānabhūtam ākṛtiḥ guṇabhūtā . \\
\noindent \textbf{(P\_1,2.64.10)} guṇavacanavat vā . \\
\noindent \textbf{(P\_1,2.64.10)} guṇavacanavat vā liṅgavacanāni bhaviṣyanti . \\
\noindent \textbf{(P\_1,2.64.10)} tat yathā guṇavacanānām śabdānām āśrayataḥ liṅgavacanāni bhavanti : śuklam vastram , śuklā śāṭī śuklaḥ kambalaḥ , śuklau kambalau śuklāḥ kambalāḥ iti . \\
\noindent \textbf{(P\_1,2.64.10)} yat asau dravyam śritaḥ bhavati guṇaḥ tasya yat liṅgam vacanam ca tat guṇasya api bhavati . \\
\noindent \textbf{(P\_1,2.64.10)} evam iha api yat asau dravyam śritā ākṛtiḥ tasya yat liṅgam vacanam ca tat ākṛteḥ api bhaviṣyati . \\
\noindent \textbf{(P\_1,2.64.10)} adhikaraṇagatiḥ sāhacaryāt . \\
\noindent \textbf{(P\_1,2.64.10)} ākṛtau ārambhaṇādīnām sambhavaḥ na asti iti kṛtvā ākṛtisahacarite dravye ārambhaṇādīni bhaviṣyanti . \\
\noindent \textbf{(P\_1,2.64.10)} na ca ekam anekādhikaraṇastham yugapat iti ādityavat viṣayaḥ . \\
\noindent \textbf{(P\_1,2.64.10)} na khalu api ekam anekādhikaraṇastham yugapat upalabhyate iti ādityavat viṣayaḥ bhaviṣyati . \\
\noindent \textbf{(P\_1,2.64.10)} tat yathā ekaḥ ādityaḥ anekādhikaraṇasthaḥ yugapat upalabhyate . \\
\noindent \textbf{(P\_1,2.64.10)} viṣamaḥ upanyāsaḥ . \\
\noindent \textbf{(P\_1,2.64.10)} na ekaḥ draṣṭā anekādhikaraṇastham ādityam yugapat upalabhate . \\
\noindent \textbf{(P\_1,2.64.10)} evam tarhi itīndravat viṣayaḥ . \\
\noindent \textbf{(P\_1,2.64.10)} tad yathā ekaḥ indraḥ anekasmin kratuśate āhūtaḥ yugapat sarvatra bhavati evam ākṛtiḥ yugapat sarvatra bhaviṣyati . \\
\noindent \textbf{(P\_1,2.64.10)} avināśaḥ anāśritatvāt . \\
\noindent \textbf{(P\_1,2.64.10)} dravyavināśe ākṛteḥ avināśaḥ . \\
\noindent \textbf{(P\_1,2.64.10)} kutaḥ . \\
\noindent \textbf{(P\_1,2.64.10)} anāśritatvāt . \\
\noindent \textbf{(P\_1,2.64.10)} anāśritā ākṛtiḥ dravyam . \\
\noindent \textbf{(P\_1,2.64.10)} kim ucyate anāśritatvāt iti yat idānīm eva uktam adhikaraṇagatiḥ sāhacaryāt iti . \\
\noindent \textbf{(P\_1,2.64.10)} evam tarhi avināśaḥ anaikātmyāt . \\
\noindent \textbf{(P\_1,2.64.10)} dravyavināśe ākṛteḥ avināśaḥ . \\
\noindent \textbf{(P\_1,2.64.10)} kutaḥ . \\
\noindent \textbf{(P\_1,2.64.10)} anaikātmyāt . \\
\noindent \textbf{(P\_1,2.64.10)} anekaḥ ātmā ākṛteḥ dravyasya ca . \\
\noindent \textbf{(P\_1,2.64.10)} tat yathā vṛkṣasthaḥ avatānaḥ vṛkṣe chinne api na vinaśyati . \\
\noindent \textbf{(P\_1,2.64.10)} vairūpyavigrahau dravyabhedāt . vairūpyavigrahau api dravyabhedāt bhaviṣyataḥ . \\
\noindent \textbf{(P\_1,2.64.10)} vyartheṣu ca sāmānyāt siddham . \\
\noindent \textbf{(P\_1,2.64.10)} vibhinnārtheṣu ca sāmānyāt siddham sarvam . \\
\noindent \textbf{(P\_1,2.64.10)} aśnoteḥ akṣaḥ . \\
\noindent \textbf{(P\_1,2.64.10)} padyateḥ pādaḥ . \\
\noindent \textbf{(P\_1,2.64.10)} mimīteḥ māṣaḥ . \\
\noindent \textbf{(P\_1,2.64.10)} tatra kriyāsāmānyāt siddham . \\
\noindent \textbf{(P\_1,2.64.10)} aparaḥ tu āha . \\
\noindent \textbf{(P\_1,2.64.10)} purākalpe etat āsīt ṣoḍaśa māṣāḥ kārṣāpaṇam ṣoḍaśaphalāḥca māṣaśambaṭyaḥ . \\
\noindent \textbf{(P\_1,2.64.10)} tatra saṃkhyāsāmānyāt siddham . \\
\noindent \textbf{(P\_1,2.65)} iha kasmāt na bhavati : ajaḥ ca barkaraḥ ca , aśvaḥ ca kiśoraḥ ca , uṣṭraḥ ca karabhaḥ ca iti . \\
\noindent \textbf{(P\_1,2.65)} tallakṣaṇaḥ cet eva viśeṣaḥ iti ucyate na ca atra tallakṣaṇaḥ eva viśeṣaḥ . \\
\noindent \textbf{(P\_1,2.65)} tallakṣaṇaḥ eva viśeṣaḥ yat samānāyām ākṛtau śabdabhedaḥ . \\
\noindent \textbf{(P\_1,2.66.1)} idam sarveṣu strīgrahaṇeṣu vicāryate : strīgrahaṇe strīpratyayagrahaṇam vā syāt stryarthagrahaṇam vā strīśabdagrahaṇam vā iti . \\
\noindent \textbf{(P\_1,2.66.1)} kim ca ataḥ . \\
\noindent \textbf{(P\_1,2.66.1)} yadi pratyayagrahaṇam vā śabdagrahaṇam vā gārgī ca gārgyāyaṇau ca gargāḥ : kena yaśabdaḥ na śrūyeta . \\
\noindent \textbf{(P\_1,2.66.1)} astriyām iti hi luk ucyate . \\
\noindent \textbf{(P\_1,2.66.1)} iha ca gārgī ca gārgyāyaṇau ca gargān paśya : tasmāt śasaḥ naḥ puṃsi iti natvam na prāpnoti . \\
\noindent \textbf{(P\_1,2.66.1)} atha arthagrahaṇam na doṣaḥ bhavati . \\
\noindent \textbf{(P\_1,2.66.1)} yathā na doṣaḥ tathā astu \\
\noindent \textbf{(P\_1,2.66.2)} iha kasmāt na bhavati : ajā ca barkaraḥ ca , vaḍavā ca kiśoraḥ ca , uṣṭrī ca karabhaḥ ca iti . \\
\noindent \textbf{(P\_1,2.66.2)} tallakṣaṇaḥ cet eva viśeṣaḥ iti ucyate . \\
\noindent \textbf{(P\_1,2.66.2)} na ca atra tallakṣaṇaḥ eva viśeṣaḥ . \\
\noindent \textbf{(P\_1,2.66.2)} tallakṣaṇaḥ eva viśeṣaḥ yat samānāyām ākṛtau śabdabhedaḥ . \\
\noindent \textbf{(P\_1,2.67)} iha kasmāt na bhavati : haṃsaḥ ca varaṭā ca kacchapaḥ ca ḍulī ca , rśyaḥ ca rohit ca iti . \\
\noindent \textbf{(P\_1,2.67)} tallakṣaṇaḥ cet eva viśeṣaḥ iti ucyate . \\
\noindent \textbf{(P\_1,2.67)} na ca atra tallakṣaṇaḥ eva viśeṣaḥ . \\
\noindent \textbf{(P\_1,2.67)} tallakṣaṇaḥ eva viśeṣaḥ yat samānāyām ākṛtau śabdabhedaḥ . \\
\noindent \textbf{(P\_1,2.68.1)} kimartham idam ucyate na pumān striyā iti eva siddham . \\
\noindent \textbf{(P\_1,2.68.1)} na sidhyati . \\
\noindent \textbf{(P\_1,2.68.1)} tallakṣaṇaḥ cet eva viśeṣaḥ iti ucyate . \\
\noindent \textbf{(P\_1,2.68.1)} na ca atra tallakṣaṇaḥ eva viśeṣaḥ . \\
\noindent \textbf{(P\_1,2.68.1)} tallakṣaṇaḥ eva viśeṣaḥ yat samānāyām ākṛtau śabdabhedaḥ . \\
\noindent \textbf{(P\_1,2.68.1)} evam tarhi siddhe sati yat imam yogam śāsti tat jñāpayati ācāryaḥ : yatra ūrdhvam prakṛteḥ tallakṣaṇaḥ eva viśeṣaḥ tatra ekaśeṣaḥ bhavati iti . \\
\noindent \textbf{(P\_1,2.68.1)} kim etasya jñāpane prayojanam . \\
\noindent \textbf{(P\_1,2.68.1)} haṃsaḥ ca varaṭā ca , kacchapaḥ ca ḍulī ca , rśyaḥ ca rohit ca iti atra ekaśeṣaḥ na bhavati . \\
\noindent \textbf{(P\_1,2.68.1)} pūrvayoḥ yogayoḥ bhūyān parihāraḥ . \\
\noindent \textbf{(P\_1,2.68.1)} yāvat brūyāt gotram yūnā iti tāvat vṛddhaḥ yūnā iti . \\
\noindent \textbf{(P\_1,2.68.1)} pūrvasūtre gotrasya vṛddham iti sañjñā kriyate . \\
\noindent \textbf{(P\_1,2.68.2)} asarūpāṇām yuvasthavirastrīpuṃsānām viśeṣasya avivakṣitatvāt sāmānyasya ca vivakṣitatvāt siddham . asarūpāṇām yuvasthavirastrīpuṃsānām viśeṣaḥ ca avivakṣitaḥ sāmānyam ca vivakṣitam . \\
\noindent \textbf{(P\_1,2.68.2)} viśeṣasya avivakṣitatvāt sāmānyasya ca vivakṣitatvāt sarūpāṇām ekaśeṣaḥ ekavibhaktau iti eva siddham . \\
\noindent \textbf{(P\_1,2.68.2)} pumān striyā iha kasmāt na bhavati : brāhmaṇavatsā ca brāhmaṇīvatsaḥ ca iti . \\
\noindent \textbf{(P\_1,2.68.2)} brāhmaṇavatsābrāhmaṇīvatsayoḥ vibhaktiparasya viśeṣavācakatvāt anekaśeṣaḥ . \\
\noindent \textbf{(P\_1,2.68.2)} brāhmaṇavatsābrāhmaṇīvatsayoḥ liṅgasya vibhaktiparasya viśeṣavācakatvāt ekaśeṣaḥ na bhaviṣyati . \\
\noindent \textbf{(P\_1,2.68.2)} yatra liṅgam vibhaktiparam eva viśeṣavācakam tatra ekaśeṣaḥ bhavati . \\
\noindent \textbf{(P\_1,2.68.2)} na atra liṅgam vibhaktiparam eva viśeṣavācakam . \\
\noindent \textbf{(P\_1,2.68.2)} yadi tarhi yatra liṅgam vibhaktiparam eva viśeṣavācakam tatra ekaśeṣaḥ bhavati iha na prāpnoti : kārakaḥ ca kārikā ca kārakau . \\
\noindent \textbf{(P\_1,2.68.2)} na hi atra liṅgam vibhaktiparam eva viśeṣavācakam . \\
\noindent \textbf{(P\_1,2.68.2)} katham punaḥ idam vijñāyate : śabdaḥ yā strī tallakṣaṇaḥ cet eva viśeṣaḥ iti āhosvit arthaḥ yā strī tallakṣaṇaḥ cet eva viśeṣaḥ iti . \\
\noindent \textbf{(P\_1,2.68.2)} kim ca ataḥ . \\
\noindent \textbf{(P\_1,2.68.2)} yadi vijñāyate śabdaḥ yā strī tallakṣaṇaḥ cet eva viśeṣaḥ iti siddham kārakaḥ ca kārikā ca kārakau . \\
\noindent \textbf{(P\_1,2.68.2)} idam tu na sidhyati : gomān ca gomatī ca gomantau . \\
\noindent \textbf{(P\_1,2.68.2)} atha vijñāyate arthaḥ yā strī tallakṣaṇaḥ cet eva viśeṣaḥ iti siddham gomān ca gomatī ca gomantau . \\
\noindent \textbf{(P\_1,2.68.2)} idam tu na sidhyati : kārakaḥ ca kārikā ca kārakau . \\
\noindent \textbf{(P\_1,2.68.2)} ubhayathā api paṭuḥ ca paṭvī ca paṭū* iti etat na sidhyati . \\
\noindent \textbf{(P\_1,2.68.2)} evam tarhi na evam vijñāyate śabdaḥ yā strī tallakṣaṇaḥ cet eva viśeṣaḥ iti na api arthaḥ yā strī tallakṣaṇaḥ cet eva viśeṣaḥ iti . \\
\noindent \textbf{(P\_1,2.68.2)} katham tarhi . \\
\noindent \textbf{(P\_1,2.68.2)} śabdārthau yā strī tatsadbhāvena ca tallakṣaṇaḥ viśeṣaḥ āśrīyate . \\
\noindent \textbf{(P\_1,2.68.2)} evam ca kṛtvā iha api prāptiḥ : brāhmaṇavatsā ca brāhmaṇīvatsaḥ ca iti . \\
\noindent \textbf{(P\_1,2.68.2)} evam tarhi idam iha vyapadeśyam sat ācāryaḥ na vyapadiśati . \\
\noindent \textbf{(P\_1,2.68.2)} kim . \\
\noindent \textbf{(P\_1,2.68.2)} tat iti anuvartate . \\
\noindent \textbf{(P\_1,2.68.2)} tat iti anena prakṛtau strīpuṃsau pratinirdiśyete . \\
\noindent \textbf{(P\_1,2.68.2)} kau ca prakṛtau . \\
\noindent \textbf{(P\_1,2.68.2)} pradhāne . \\
\noindent \textbf{(P\_1,2.68.2)} pradhānam yā śabdastrī pradhānam yā arthastrī iti . \\
\noindent \textbf{(P\_1,2.69)} ayam yogaḥ śakyaḥ avaktum . \\
\noindent \textbf{(P\_1,2.69)} katham śuklaḥ ca kambalaḥ śuklam ca vastram tat idam śuklam , te* ime śukle , śuklaḥ ca kambalaḥ śuklā ca bṛhatikā śuklam ca vastram tat idam śuklam , tāni imani śuklāni . \\
\noindent \textbf{(P\_1,2.69)} pradhāne kāryasampratyayāt śeṣaḥ . \\
\noindent \textbf{(P\_1,2.69)} pradhāne kāryasampratyayāt śeṣaḥ bhaviṣyati . \\
\noindent \textbf{(P\_1,2.69)} kim ca pradhānam . \\
\noindent \textbf{(P\_1,2.69)} napuṃsakam . \\
\noindent \textbf{(P\_1,2.69)} katham punaḥ jñāyate napuṃsakam pradhānam iti . \\
\noindent \textbf{(P\_1,2.69)} evam hi dṛśyate loke : anirjñāte arthe guṇasandehe ca napuṃsakaliṅgam prayujyate . \\
\noindent \textbf{(P\_1,2.69)} kim jātam iti ucyate . \\
\noindent \textbf{(P\_1,2.69)} dvayam ca eva hi jāyate strī vā pumān vā . \\
\noindent \textbf{(P\_1,2.69)} tathā vidūre avyaktam ārūpam dṛṣṭvā vaktāraḥ bhavanti mahiṣīrūpam iva brāhmaṇīrūpam iva . \\
\noindent \textbf{(P\_1,2.69)} pradhāne kāryasampratyayāt napuṃsakasya śeṣaḥ bhaviṣyati . \\
\noindent \textbf{(P\_1,2.69)} idam tarhi prayojanam : ekavat ca asya anyatarasyām iti vakṣyāmi iti . \\
\noindent \textbf{(P\_1,2.69)} etat api na asti prayojanam . \\
\noindent \textbf{(P\_1,2.69)} ākṛtivācitvāt ekavacanam . \\
\noindent \textbf{(P\_1,2.69)} ākṛtivācitvāt ekavacanam bhaviṣyati . \\
\noindent \textbf{(P\_1,2.69)} yadā dravyābhidhānam tadā dvivacanabahuvacane bhaviṣyataḥ . \\
\noindent \textbf{(P\_1,2.68,} 70-71) KA\_I,250.13-251.7 Ro\_II,168-169 {1/27}     kimartham idam ucyate na pumān striyā iti eva siddham . \\
\noindent \textbf{(P\_1,2.68,} 70-71) KA\_I,250.13-251.7 Ro\_II,168-169 {2/27}     bhrātṛputrapitṛśvaśurāṇām kāraṇāt dravye śabdaniveśaḥ . \\
\noindent \textbf{(P\_1,2.68,} 70-71) KA\_I,250.13-251.7 Ro\_II,168-169 {3/27}     bhrātṛputrapitṛśvaśurāṇām kāraṇāt dravye śabdaniveśaḥ bhavati . \\
\noindent \textbf{(P\_1,2.68,} 70-71) KA\_I,250.13-251.7 Ro\_II,168-169 {4/27}     bhrātṛputrapitṛśvaśurāṇām kāraṇād dravye śabdaniveśaḥ iti cet tulyakāraṇatvāt siddham . \\
\noindent \textbf{(P\_1,2.68,} 70-71) KA\_I,250.13-251.7 Ro\_II,168-169 {5/27}     yadi tāvat bibharti iti bhrātā svasari api etat bhavati . \\
\noindent \textbf{(P\_1,2.68,} 70-71) KA\_I,250.13-251.7 Ro\_II,168-169 {6/27}     tathā yadi punāti prīṇāti iti vā putraḥ duhitari api etat bhavati . \\
\noindent \textbf{(P\_1,2.68,} 70-71) KA\_I,250.13-251.7 Ro\_II,168-169 {7/27}     tathā yadi pāti pālayati iti vā pitā mātari api etat bhavati . \\
\noindent \textbf{(P\_1,2.68,} 70-71) KA\_I,250.13-251.7 Ro\_II,168-169 {8/27}     tathā yadi āśu āptavyaḥ śvaśuraḥ śvaśrvām api etat bhavati . \\
\noindent \textbf{(P\_1,2.68,} 70-71) KA\_I,250.13-251.7 Ro\_II,168-169 {9/27}     darśanam vai hetuḥ . \\
\noindent \textbf{(P\_1,2.68,} 70-71) KA\_I,250.13-251.7 Ro\_II,168-169 {10/27}     na hi svasari bhrātṛśabdaḥ dṛśyate . \\
\noindent \textbf{(P\_1,2.68,} 70-71) KA\_I,250.13-251.7 Ro\_II,168-169 {11/27}     darśanam hetuḥ iti cet tulyam . \\
\noindent \textbf{(P\_1,2.68,} 70-71) KA\_I,250.13-251.7 Ro\_II,168-169 {12/27}     darśanam hetuḥ iti cet tulyam etat bhavati . \\
\noindent \textbf{(P\_1,2.68,} 70-71) KA\_I,250.13-251.7 Ro\_II,168-169 {13/27}     svasari api bhrātṛśabdaḥ dṛśyatām . \\
\noindent \textbf{(P\_1,2.68,} 70-71) KA\_I,250.13-251.7 Ro\_II,168-169 {14/27}     tulyam hi kāraṇam . \\
\noindent \textbf{(P\_1,2.68,} 70-71) KA\_I,250.13-251.7 Ro\_II,168-169 {15/27}     na vai eṣaḥ loke sampratyayaḥ . \\
\noindent \textbf{(P\_1,2.68,} 70-71) KA\_I,250.13-251.7 Ro\_II,168-169 {16/27}     na hi loke bhrātā ānīyatām iti ukte svasā ānīyate . \\
\noindent \textbf{(P\_1,2.68,} 70-71) KA\_I,250.13-251.7 Ro\_II,168-169 {17/27}     tadviṣayam ca . \\
\noindent \textbf{(P\_1,2.68,} 70-71) KA\_I,250.13-251.7 Ro\_II,168-169 {18/27}     tadviṣayam ca etat draṣṭavyam bhavati : svasari bhrātṛtvam . \\
\noindent \textbf{(P\_1,2.68,} 70-71) KA\_I,250.13-251.7 Ro\_II,168-169 {19/27}     kiṃviṣayam . \\
\noindent \textbf{(P\_1,2.68,} 70-71) KA\_I,250.13-251.7 Ro\_II,168-169 {20/27}     ekaśeṣaviṣayam . \\
\noindent \textbf{(P\_1,2.68,} 70-71) KA\_I,250.13-251.7 Ro\_II,168-169 {21/27}     yuktam punaḥ yat niyataviṣayāḥ śabdāḥ syuḥ . \\
\noindent \textbf{(P\_1,2.68,} 70-71) KA\_I,250.13-251.7 Ro\_II,168-169 {22/27}     bāḍham yuktam . \\
\noindent \textbf{(P\_1,2.68,} 70-71) KA\_I,250.13-251.7 Ro\_II,168-169 {23/27}     anyatra api tadviṣayadarśanāt . \\
\noindent \textbf{(P\_1,2.68,} 70-71) KA\_I,250.13-251.7 Ro\_II,168-169 {24/27}     anyatra api tadviṣayāḥ śabdāḥ dṛśyante . \\
\noindent \textbf{(P\_1,2.68,} 70-71) KA\_I,250.13-251.7 Ro\_II,168-169 {25/27}     tat yathā : samāne rakte varṇe gauḥ lohitaḥ iti bhavati aśvaḥ śoṇaḥ iti . \\
\noindent \textbf{(P\_1,2.68,} 70-71) KA\_I,250.13-251.7 Ro\_II,168-169 {26/27}     samāne ca kāle varṇe gauḥ kṛṣṇaḥ iti bhavati aśvaḥ hemaḥ iti . \\
\noindent \textbf{(P\_1,2.68,} 70-71) KA\_I,250.13-251.7 Ro\_II,168-169 {27/27}     samāne ca śukle varṇe gauḥ śvetaḥ iti bhavati aśvaḥ karkaḥ iti . \\
\noindent \textbf{(P\_1,2.72.1)} tyadāditaḥ śeṣe punnapuṃsakataḥ liṅgavacanāni . \\
\noindent \textbf{(P\_1,2.72.1)} tyadāditaḥ śeṣe punnapuṃsakataḥ liṅgavacanāni bhavanti . \\
\noindent \textbf{(P\_1,2.72.1)} sā ca devadattaḥ ca tau sā ca kuṇḍe ca tāni . \\
\noindent \textbf{(P\_1,2.72.1)} advandvatatpuruṣaviśeṣaṇānām . \\
\noindent \textbf{(P\_1,2.72.1)} advandvatatpuruṣaviśeṣaṇānām iti vaktavyam . \\
\noindent \textbf{(P\_1,2.72.1)} iha mā bhūt . \\
\noindent \textbf{(P\_1,2.72.1)} saḥ ca kukkuṭaḥ sā ca mayūrī kukkuṭamayūryau te . \\
\noindent \textbf{(P\_1,2.72.1)} ardham pippalyāḥ tat ardhapippalī ca sā ardhapippalyau te . \\
\noindent \textbf{(P\_1,2.72.2)} ayam api yogaḥ śakyaḥ avaktum . \\
\noindent \textbf{(P\_1,2.72.2)} katham . \\
\noindent \textbf{(P\_1,2.72.2)} tyadādīnām sāmānyārthatvāt . \\
\noindent \textbf{(P\_1,2.72.2)} tyadādīnām sāmānyam arthaḥ . \\
\noindent \textbf{(P\_1,2.72.2)} ātaḥ ca sāmānyam . \\
\noindent \textbf{(P\_1,2.72.2)} devadatte api hi saḥ iti etat bhavati yajñadatte api . \\
\noindent \textbf{(P\_1,2.72.2)} tyadādīnām sāmānyārthatvāt śeṣaḥ bhaviṣyati . \\
\noindent \textbf{(P\_1,2.72.2)} idam tarhi prayojanam : parasya śeṣam vakṣyāmi iti . \\
\noindent \textbf{(P\_1,2.72.2)} parasya ca ubhayavācitvāt . \\
\noindent \textbf{(P\_1,2.72.2)} ubhayavāci param . \\
\noindent \textbf{(P\_1,2.72.2)} pūrvaśeṣadarśanāt ca . \\
\noindent \textbf{(P\_1,2.72.2)} pūrvasya khalu api śeṣaḥ dṛśyate : saḥ ca yaḥ ca tau ānaya , yau ānaya iti . \\
\noindent \textbf{(P\_1,2.72.2)} idam tarhi prayojanam :dvandvaḥ mā bhūt iti . \\
\noindent \textbf{(P\_1,2.72.2)} etat api na asti prayojanam . \\
\noindent \textbf{(P\_1,2.72.2)} sāmānyaviśeṣavācinoḥ ca dvandvābhāvāt siddham . sāmānyaviśeṣavācinoḥ ca dvandvaḥ na bhavati iti vaktavyam . \\
\noindent \textbf{(P\_1,2.72.2)} yadi sāmānyaviśeṣavācinoḥ dvandvaḥ na bhavati iti ucyate śūdrābhīram , gobalīvardam , tṛṇolapam iti na sidhyati . \\
\noindent \textbf{(P\_1,2.72.2)} na eṣaḥ doṣaḥ . \\
\noindent \textbf{(P\_1,2.72.2)} iha tāvat śūdrābhīram iti : ābhīrāḥ jātyantarāṇi . \\
\noindent \textbf{(P\_1,2.72.2)} gobalīvardam iti : gāvaḥ utkālitapuṃskāḥ vāhāya ca vikrayāya ca . \\
\noindent \textbf{(P\_1,2.72.2)} striyaḥ eva avaśiṣyante . \\
\noindent \textbf{(P\_1,2.72.2)} tṛṇolapam iti : apām ulapam iti nāmadheyam . \\
\noindent \textbf{(P\_1,2.72.2)} tat tarhi vaktavyam . \\
\noindent \textbf{(P\_1,2.72.2)} na vaktavyam . \\
\noindent \textbf{(P\_1,2.72.2)} sāmānyena uktatvāt viśeṣasya prayogaḥ na bhaviṣyati . \\
\noindent \textbf{(P\_1,2.72.2)} sāmānyena uktatvāt tasya arthasya viśeṣasya prayogeṇa na bhavitavyam . \\
\noindent \textbf{(P\_1,2.72.2)} kim kāraṇam . \\
\noindent \textbf{(P\_1,2.72.2)} uktārthānām aprayogaḥ iti . \\
\noindent \textbf{(P\_1,2.72.2)} na tarhi idānīm idam bhavati : tam brāhmaṇam ānaya gārgyam iti . \\
\noindent \textbf{(P\_1,2.72.2)} bhavati yadā niyogataḥ tasya eva ānayanam bhavati . \\
\noindent \textbf{(P\_1,2.72.2)} evam tarhi yena eva khalu api hetunā etat vākyam bhavati tam brāhmaṇam ānaya gārgyam iti tena eva hetunā vṛttiḥ api prāpnoti . \\
\noindent \textbf{(P\_1,2.72.2)} tasmāt sāmānyaviśeṣavācinoḥ dvandvaḥ na bhavati iti vaktavyam . \\
\noindent \textbf{(P\_1,2.73)} ayam api yogaḥ śakyaḥ avaktum . \\
\noindent \textbf{(P\_1,2.73)} katham gāvaḥ imāḥ caranti , ajāḥ imāḥ caranti . \\
\noindent \textbf{(P\_1,2.73)} gāvaḥ utkālitapuṃskāḥ vāhāya ca vikrayāya ca . \\
\noindent \textbf{(P\_1,2.73)} striyaḥ eva avaśiṣyante . \\
\noindent \textbf{(P\_1,2.73)} idam tarhi prayojanam : grāmyeṣu iti vakṣyāmi iti . \\
\noindent \textbf{(P\_1,2.73)} iha mā bhūt : nyaṅkavaḥ ime , śūkarāḥ ime iti . \\
\noindent \textbf{(P\_1,2.73)} kaḥ punaḥ arhati agrāmyāṇām puṃsaḥ utkālayitum ye grahītum aśakyāḥ . \\
\noindent \textbf{(P\_1,2.73)} kutaḥ eva vāhāya ca vikrayāya ca . \\
\noindent \textbf{(P\_1,2.73)} idam tarhi prayojanam : paśuṣu iti vakṣyāmi iti . \\
\noindent \textbf{(P\_1,2.73)} iha mā bhūt : brāhmaṇāḥ ime , vṛṣalāḥ ime . \\
\noindent \textbf{(P\_1,2.73)} kaḥ punaḥ arhati apaśūnām puṃsaḥ utkālayitum ye aśakyāḥ vāhāya ca vikrayāya ca . \\
\noindent \textbf{(P\_1,2.73)} idam tarhi prayojanam : saṅgheṣu iti vakṣyāmi iti . \\
\noindent \textbf{(P\_1,2.73)} iha mā bhūt : etau gāvaḥ carataḥ . \\
\noindent \textbf{(P\_1,2.73)} kaḥ punaḥ arhati nirjñāte arthe anyathā prayoktum . \\
\noindent \textbf{(P\_1,2.73)} idam tarhi prayojanam : ataruṇeṣu iti vakṣyāmi iti . \\
\noindent \textbf{(P\_1,2.73)} iha mā bhūt : uraṇakāḥ ime , barkarāḥ ime iti . \\
\noindent \textbf{(P\_1,2.73)} kaḥ punaḥ arhati taruṇānām puṃsaḥ utkālayitum ye aśakyāḥ vāhāya ca vikrayāya ca . \\
\noindent \textbf{(P\_1,2.73)} anekaśapheṣu iti vaktavyam iha mā bhūt : aśvāḥ caranti . \\
\noindent \textbf{(P\_1,2.73)} gardabhāḥ caranti iti . \\
\noindent \textbf{(P\_1,3.1.1)} kutaḥ ayam vakāraḥ . \\
\noindent \textbf{(P\_1,3.1.1)} yadi tāvat saṃhitayā nirdeśaḥ kriyate bhvādayaḥ iti bhavitavyam . \\
\noindent \textbf{(P\_1,3.1.1)} atha asaṃhitayā bhū-ādayaḥ iti bhavitavyam . \\
\noindent \textbf{(P\_1,3.1.1)} ataḥ uttaram paṭhati : bhūvādīnām vakāraḥ ayam maṅgalārthaḥ prayujyate . \\
\noindent \textbf{(P\_1,3.1.1)} māṅgalikaḥ ācāryaḥ mahataḥ śāstraughasya maṅgalārtham vakāram āgamam prayuṅkte . \\
\noindent \textbf{(P\_1,3.1.1)} maṅgalādīni maṅgalamadhyāni maṅgalāntāni hi śāstrāṇi prathante vīrapuruṣāṇi ca bhavanti āyuṣmatpuruṣāṇi ca . \\
\noindent \textbf{(P\_1,3.1.1)} adhyetāraḥ ca siddhārthāḥ yathā syuḥ iti . \\
\noindent \textbf{(P\_1,3.1.1)} atha ādigrahaṇam kimartham . \\
\noindent \textbf{(P\_1,3.1.1)} yadi tāvat paṭhyante na arthaḥ ādigrahaṇena . \\
\noindent \textbf{(P\_1,3.1.1)} anyatra api hi ayam paṭhan ādigrahaṇam na karoti . \\
\noindent \textbf{(P\_1,3.1.1)} kva anyatra . \\
\noindent \textbf{(P\_1,3.1.1)} mṛḍamṛdagudhakuṣakliśavadavasaḥ ktvā iti . \\
\noindent \textbf{(P\_1,3.1.1)} atha na paṭhyante natarām arthaḥ ādigrahaṇena . \\
\noindent \textbf{(P\_1,3.1.1)} na hi apaṭhitāḥ śakyāḥ ādigrahaṇena viśeṣayitum . \\
\noindent \textbf{(P\_1,3.1.1)} evam tarhi siddhe sati yat ādigrahaṇam karoti tat jñāpayati ācāryaḥ asti ca pāṭhaḥ bāhyaḥ ca sūtrāt iti . \\
\noindent \textbf{(P\_1,3.1.1)} kim etasya jñāpane prayojanam . \\
\noindent \textbf{(P\_1,3.1.1)} pāṭhena dhātusañjñā iti etat upapannam bhavati . \\
\noindent \textbf{(P\_1,3.1.1)} pāṭhena dhātusañjñāyām samānaśabdapratiṣedhaḥ . \\
\noindent \textbf{(P\_1,3.1.1)} pāṭhena dhātusañjñāyām samānaśabdānām pratiṣedhaḥ vaktavyaḥ . \\
\noindent \textbf{(P\_1,3.1.1)} yā iti dhātuḥ yā iti ābantaḥ . \\
\noindent \textbf{(P\_1,3.1.1)} vā iti dhātuḥ vā iti nipātaḥ . \\
\noindent \textbf{(P\_1,3.1.1)} nu iti dhātuḥ nu iti pratyayaḥ ca nipātaḥ ca . \\
\noindent \textbf{(P\_1,3.1.1)} div iti dhātuḥ div iti prātipadikam . \\
\noindent \textbf{(P\_1,3.1.1)} kim ca syāt yadi eteṣām api dhātusañjñā syāt . \\
\noindent \textbf{(P\_1,3.1.1)} dhātoḥ iti tavyādīnām utpattiḥ prasajyeta . \\
\noindent \textbf{(P\_1,3.1.1)} na eṣaḥ doṣaḥ . \\
\noindent \textbf{(P\_1,3.1.1)} sādhane tavyādayaḥ vidhīyante sādhanam ca kriyāyāḥ . \\
\noindent \textbf{(P\_1,3.1.1)} kriyābhāvāt sādhanābhāvaḥ . \\
\noindent \textbf{(P\_1,3.1.1)} sādhanābhāvāt satyām api dhātusañjñāyām tavyādayaḥ na bhaviṣyanti . \\
\noindent \textbf{(P\_1,3.1.1)} iha tarhi : yāḥ paśya : ātaḥ dhātoḥ iti lopaḥ prasajyeta . \\
\noindent \textbf{(P\_1,3.1.1)} na eṣaḥ doṣaḥ . \\
\noindent \textbf{(P\_1,3.1.1)} anāpaḥ iti evam saḥ . \\
\noindent \textbf{(P\_1,3.1.1)} asya tarhi vāśabdasya nipātasya adhātuḥ iti prātipadikasañjñāyāḥ pratiṣedhaḥ prasajyeta . \\
\noindent \textbf{(P\_1,3.1.1)} aprātipadikatvāt svādyutpattiḥ na syāt . \\
\noindent \textbf{(P\_1,3.1.1)} na eṣaḥ doṣaḥ . \\
\noindent \textbf{(P\_1,3.1.1)} nipātasya anarthakasya prātipadikatvam coditam . \\
\noindent \textbf{(P\_1,3.1.1)} tatra anarthakagrahaṇam na kariṣyate : nipātaḥ prātipadikam iti eva . \\
\noindent \textbf{(P\_1,3.1.1)} iha tarhi : trasnū iti : aci śnudhātubhruvām yvoḥ iyaṅuvaṅau iti uvaṅādeśaḥ prasajyeta . \\
\noindent \textbf{(P\_1,3.1.1)} na eṣaḥ doṣaḥ . \\
\noindent \textbf{(P\_1,3.1.1)} ācāryapravṛttiḥ jñāpayati na pratyayasya uvaṅādeśaḥ bhavati iti yat ayam tatra śnugrahaṇam karoti . \\
\noindent \textbf{(P\_1,3.1.1)} asya tarhi divśabdasya adhātuḥ iti prātipadikasañjñāyāḥ pratiṣedhaḥ prasajyeta . \\
\noindent \textbf{(P\_1,3.1.1)} aprātipadikatvāt svādyutpattiḥ na syāt . \\
\noindent \textbf{(P\_1,3.1.1)} na eṣaḥ doṣaḥ . \\
\noindent \textbf{(P\_1,3.1.1)} ācāryapravṛttiḥ jñāpayati utpadyante divśabdāt svādayaḥ iti yat ayam divaḥ sau auttvam śāsti . \\
\noindent \textbf{(P\_1,3.1.1)} na etat asti jñāpakam . \\
\noindent \textbf{(P\_1,3.1.1)} asti hi anyat etasya vacane prayojanam . \\
\noindent \textbf{(P\_1,3.1.1)} kim . \\
\noindent \textbf{(P\_1,3.1.1)} divśabdaḥ yat prātipadikam tadartham etat syāt : akṣadyūḥ iti . \\
\noindent \textbf{(P\_1,3.1.1)} na vai atra iṣyate . \\
\noindent \textbf{(P\_1,3.1.1)} aniṣṭam ca prāpnoti iṣṭam ca na sidhyati . \\
\noindent \textbf{(P\_1,3.1.1)} evam tarhi ananubandhakagrahaṇe na sānubandhakasya iti evam etasya na bhaviṣyati . \\
\noindent \textbf{(P\_1,3.1.1)} evam api ananubandhakaḥ divśabdaḥ na asti iti kṛtvā sānubandhakasya grahaṇam vijñāsyate . \\
\noindent \textbf{(P\_1,3.1.1)} parimāṇagrahaṇam ca . \\
\noindent \textbf{(P\_1,3.1.1)} parimāṇagrahaṇam ca kartavyam . \\
\noindent \textbf{(P\_1,3.1.1)} iyān avadhiḥ dhātusañjñaḥ bhavati iti vaktavyam kutaḥ hi etat bhūśabdaḥ dhātusañjñaḥ bhaviṣyati na punaḥ bhvedhśabdaḥ iti \\
\noindent \textbf{(P\_1,3.1.2)} yadi punaḥ kriyāvacanaḥ dhātuḥ iti etat lakṣaṇam kriyeta . \\
\noindent \textbf{(P\_1,3.1.2)} kā punaḥ kriyā . \\
\noindent \textbf{(P\_1,3.1.2)} īhā . \\
\noindent \textbf{(P\_1,3.1.2)} kā punaḥ īhā . \\
\noindent \textbf{(P\_1,3.1.2)} ceṣṭā . \\
\noindent \textbf{(P\_1,3.1.2)} kā punaḥ ceṣṭā . \\
\noindent \textbf{(P\_1,3.1.2)} vyāpāraḥ . \\
\noindent \textbf{(P\_1,3.1.2)} sarvathā bhavān śabdena eva śabdān ācaṣṭe . \\
\noindent \textbf{(P\_1,3.1.2)} na kim cid arthajātam nidarśayati : evañjātīyikā kriyā iti . \\
\noindent \textbf{(P\_1,3.1.2)} kriyā nāma iyam atyantāparidṛṣṭā . \\
\noindent \textbf{(P\_1,3.1.2)} aśakyā kriyā piṇḍībhūtā nidarśayitum yathā garbhaḥ nirluṭhitaḥ . \\
\noindent \textbf{(P\_1,3.1.2)} sā asau anumānagamyā . \\
\noindent \textbf{(P\_1,3.1.2)} kaḥ asau anumānaḥ . \\
\noindent \textbf{(P\_1,3.1.2)} iha sarveṣu sādhaneṣu sannihiteṣu kadā cit pacati iti etat bhavati kadācit na bhavati . \\
\noindent \textbf{(P\_1,3.1.2)} yasmin sādhane sannihite pacati iti etat bhavati sā nūnam kriyā . \\
\noindent \textbf{(P\_1,3.1.2)} atha vā yayā devadattaḥ iha bhūtvā pāṭaliputre bhavati sā nūnam kriyā . \\
\noindent \textbf{(P\_1,3.1.2)} katham punaḥ jñāyate kriyāvacanāḥ pacādayaḥ iti . \\
\noindent \textbf{(P\_1,3.1.2)} yat eṣām karotinā sāmānādhikaraṇyam : kim karoti . \\
\noindent \textbf{(P\_1,3.1.2)} pacati . \\
\noindent \textbf{(P\_1,3.1.2)} kim kariṣyati . \\
\noindent \textbf{(P\_1,3.1.2)} pakṣyati . \\
\noindent \textbf{(P\_1,3.1.2)} kim akārṣīt . \\
\noindent \textbf{(P\_1,3.1.2)} apākṣīt iti . \\
\noindent \textbf{(P\_1,3.1.2)} tatra kriyāvacane upasargapratyayapratiṣedhaḥ . \\
\noindent \textbf{(P\_1,3.1.2)} kriyāvacane dhātau upasargapratyayayoḥ pratiṣedhaḥ vaktavyaḥ . \\
\noindent \textbf{(P\_1,3.1.2)} pacati prapacati . \\
\noindent \textbf{(P\_1,3.1.2)} kim punaḥ kāraṇam prāpnoti . \\
\noindent \textbf{(P\_1,3.1.2)} saṅghātena arthagateḥ . \\
\noindent \textbf{(P\_1,3.1.2)} saṅghātena hi arthaḥ gamyate saprakṛtikena sapratyayakena sopasargeṇa ca . \\
\noindent \textbf{(P\_1,3.1.2)} astibhavatividyatīnām dhātutvam . \\
\noindent \textbf{(P\_1,3.1.2)} astibhavatividyatīnām dhātusañjñā vaktavyā . \\
\noindent \textbf{(P\_1,3.1.2)} yathā hi bhavatā karotinā pacādīnām sāmānādhikaraṇyam nidarśitam na tathā astyādīnām nidarśyate . \\
\noindent \textbf{(P\_1,3.1.2)} na hi bhavati kim karoti asti iti . \\
\noindent \textbf{(P\_1,3.1.2)} pratyayārthasya avyatirekāt prakṛtyantareṣu . \\
\noindent \textbf{(P\_1,3.1.2)} pratyayārthasya avyatirekāt prakṛtyantareṣu manyāmahe dhātuḥ eva kriyām āha iti . \\
\noindent \textbf{(P\_1,3.1.2)} pacati paṭhati : prakṛtyarthaḥ anyaḥ ca anyaḥ ca . \\
\noindent \textbf{(P\_1,3.1.2)} pratyayārthaḥ saḥ eva . \\
\noindent \textbf{(P\_1,3.1.2)} dhātoḥ ca arthābhedāt pratyayāntareṣu . \\
\noindent \textbf{(P\_1,3.1.2)} dhātoḥ ca arthābhedāt pratyayāntareṣu manyāmahe dhātuḥ eva kriyām āha iti . \\
\noindent \textbf{(P\_1,3.1.2)} paktā pacanam pākaḥ iti : pratyayārthaḥ anyaḥ ca anyaḥ ca bhavati . \\
\noindent \textbf{(P\_1,3.1.2)} prakṛtyarthaḥ saḥ eva. katham punaḥ jñāyate ayam prakṛtyarthaḥ ayam pratyayārthaḥ iti . \\
\noindent \textbf{(P\_1,3.1.2)} siddham tu anvayavyatirekābhyām . \\
\noindent \textbf{(P\_1,3.1.2)} anvayāt vyatirekāt ca . \\
\noindent \textbf{(P\_1,3.1.2)} kaḥ asau anvayaḥ vyatirekaḥ vā . \\
\noindent \textbf{(P\_1,3.1.2)} iha pacati iti ukte kaḥ cit śabdaḥ śrūyate : pacśabdaḥ cakārāntaḥ atiśabdaḥ ca pratyayaḥ . \\
\noindent \textbf{(P\_1,3.1.2)} arthaḥ api kaḥ cit gamyate : viklittiḥ kartṛtvam ekatvam ca . \\
\noindent \textbf{(P\_1,3.1.2)} paṭhati iti ukte kaḥ cit śabdaḥ hīyate kaḥ cit upajāyate kaḥ cit anvayī : pacśabdaḥ hīyate paṭhśabdaḥ upajāyate atiśabdaḥ anvayī . \\
\noindent \textbf{(P\_1,3.1.2)} arthaḥ api kaḥ cit hīyate kaḥ cit upajāyate kaḥ cit anvayī : viklittiḥ hīyate paṭhikriyā upajāyate kartṛtvam ca ekatvam ca anvayī . \\
\noindent \textbf{(P\_1,3.1.2)} te manyāmahe : yaḥ śabdaḥ hīyate tasya asau arthaḥ yaḥ arthaḥ hīyate . \\
\noindent \textbf{(P\_1,3.1.2)} yaḥ śabdaḥ upajāyate tasya asau arthaḥ yaḥ arthaḥ upajāyate . \\
\noindent \textbf{(P\_1,3.1.2)} yaḥ śabdaḥ anvayī tasya asau arthaḥ yaḥ arthaḥ anvayī . \\
\noindent \textbf{(P\_1,3.1.2)} viṣamaḥ upanyāsaḥ . \\
\noindent \textbf{(P\_1,3.1.2)} bahavaḥ hi śabdāḥ ekārthāḥ bhavanti . \\
\noindent \textbf{(P\_1,3.1.2)} tat yathā : indraḥ śakraḥ puruhūtaḥ purandaraḥ , kanduḥ koṣṭhaḥ kuśūlaḥ iti . \\
\noindent \textbf{(P\_1,3.1.2)} ekaḥ ca śabdaḥ bahvarthaḥ . \\
\noindent \textbf{(P\_1,3.1.2)} tat yathā : akṣāḥ pādāḥ māṣāḥ iti . \\
\noindent \textbf{(P\_1,3.1.2)} ataḥ kim na sādhīyaḥ arthavattā siddhā bhavati . \\
\noindent \textbf{(P\_1,3.1.2)} na api brūmaḥ arthavattā na sidhyati iti .varṇitā arthavattā anvayavyatirekābhyām eva . \\
\noindent \textbf{(P\_1,3.1.2)} tatra kutaḥ etat : ayam prakṛtyarthaḥ ayam pratyayārthaḥ iti na punaḥ prakṛtiḥ eva ubhau arthau brūyāt pratyayaḥ eva vā . \\
\noindent \textbf{(P\_1,3.1.2)} sāmānyaśabdāḥ ete evam syuḥ . \\
\noindent \textbf{(P\_1,3.1.2)} sāmānyaśabdāḥ ca na antareṇa prakaraṇam viśeṣam vā viśeṣeṣu avatiṣṭhante . \\
\noindent \textbf{(P\_1,3.1.2)} yataḥ tu khalu niyogataḥ pacati iti ukte svabhāvataḥ kasmin cit viśeṣe pacatiśabdaḥ vartate ataḥ manyāmahe na ime sāmānyaśabdāḥ iti . \\
\noindent \textbf{(P\_1,3.1.2)} na cet sāmānyaśabdāḥ prakṛtiḥ prakṛtyarthe vartate pratyayaḥ pratyayārthe . \\
\noindent \textbf{(P\_1,3.1.2)} kriyāviśeṣakaḥ upasargaḥ . \\
\noindent \textbf{(P\_1,3.1.2)} pacati iti kriyā gamyate . \\
\noindent \textbf{(P\_1,3.1.2)} tām praḥ viśinaṣṭi . \\
\noindent \textbf{(P\_1,3.1.2)} yadi api tāvat atra etat śakyate vaktum yatra dhātuḥ upasargam vyabhicarati yatra na khalu tam vyabhicarati tatra katham : adhyeti , adhīte iti . \\
\noindent \textbf{(P\_1,3.1.2)} yadi api atra dhātuḥ upasargam na vyabhicarati upasargaḥ tu dhātum vyabhicarati . \\
\noindent \textbf{(P\_1,3.1.2)} te manyāmahe : yaḥ eva asya adheḥ anyatra arthaḥ sa iha api iti . \\
\noindent \textbf{(P\_1,3.1.2)} kaḥ punaḥ anyatra adheḥ arthaḥ . \\
\noindent \textbf{(P\_1,3.1.2)} adhiḥ uparibhāve vartate . \\
\noindent \textbf{(P\_1,3.1.2)} iha tarhi vyaktam arthāntaram gamyate : tiṣṭhati pratiṣṭhate iti . \\
\noindent \textbf{(P\_1,3.1.2)} tiṣṭhati iti vrajikriyāyāḥ nivṛttiḥ pratiṣṭhate iti vrajikriyā gamyate . \\
\noindent \textbf{(P\_1,3.1.2)} te manyāmahe upasargakṛtam etat yena atra vrajikriyā gamyate . \\
\noindent \textbf{(P\_1,3.1.2)} praḥ ayam dṛṣṭāpacāraḥ ādikarmaṇi vartate . \\
\noindent \textbf{(P\_1,3.1.2)} na ca idam na asti bahvarthāḥ api dhātavaḥ bhavanti iti . \\
\noindent \textbf{(P\_1,3.1.2)} tat yathā : vapiḥ prakiraṇe ḍṛṣṭaḥ chedane api vartate : keśaśmaśru vapati iti . \\
\noindent \textbf{(P\_1,3.1.2)} īḍiḥ stuticodanāyācñāsu dṛṣṭaḥ preraṇe api vartate : agniḥ vai itaḥ vṛṣṭim īṭṭe marutaḥ amutaḥ cyāvayanti iti . \\
\noindent \textbf{(P\_1,3.1.2)} karotiḥ abhūtaprādurbhāve dṛṣṭaḥ nirmalīkaraṇe api vartate : pṛṣṭham kuru pācau , kuru . \\
\noindent \textbf{(P\_1,3.1.2)} unmṛdāna iti gamyate . \\
\noindent \textbf{(P\_1,3.1.2)} nikṣepaṇe ca api vartate : kaṭe kuru , ghaṭe kuru , aśmānam itaḥ kuru . \\
\noindent \textbf{(P\_1,3.1.2)} sthāpaya iti gamyate . \\
\noindent \textbf{(P\_1,3.1.2)} evam iha api tiṣṭhatiḥ eva vrajikriyām āha tiṣṭhatiḥ eva vrajikriyāyāḥ nivṛttim . \\
\noindent \textbf{(P\_1,3.1.2)} ayam tarhi doṣaḥ : astibhavatividyatīnām dhātutvam iti . \\
\noindent \textbf{(P\_1,3.1.3)} yadi punaḥ bhāvavacanaḥ dhātuḥ iti evam lakṣaṇam kriyeta . \\
\noindent \textbf{(P\_1,3.1.3)} katham punaḥ jñāyate bhāvavacanāḥ pacādayaḥ iti . \\
\noindent \textbf{(P\_1,3.1.3)} yat eṣām bhavatinā sāmānādhikaraṇyam : bhavati pacati , bhavati pakṣyati , bhavati apākṣīt iti . \\
\noindent \textbf{(P\_1,3.1.3)} kaḥ punaḥ bhāvaḥ . \\
\noindent \textbf{(P\_1,3.1.3)} bhavateḥ svapadārthaḥ bhavanam bhāvaḥ iti . \\
\noindent \textbf{(P\_1,3.1.3)} yadi bhavateḥ svapadārthaḥ bhavanam bhāvaḥ vipratiṣiddhānām dhātusañjñā na prāpnoti : bhedaḥ , chedaḥ . \\
\noindent \textbf{(P\_1,3.1.3)} anyaḥ hi bhāvaḥ anyaḥ hi abhāvaḥ . \\
\noindent \textbf{(P\_1,3.1.3)} ātaḥ ca anyaḥ bhāvaḥ anyaḥ abhāvaḥ iti . \\
\noindent \textbf{(P\_1,3.1.3)} yaḥ hi yasya bhāvam icchati saḥ na tasya abhāvam yasya ca abhāvam na tasya bhāvam . \\
\noindent \textbf{(P\_1,3.1.3)} pacādīnām ca dhātusañjñā na prāpnoti . \\
\noindent \textbf{(P\_1,3.1.3)} yathā hi bhavatā kriyāvacane dhātau karotinā pacādīnām sāmānādhikaraṇyam nidarśitam na tathā bhāvavacane dhātau nidarśyate . \\
\noindent \textbf{(P\_1,3.1.3)} karotiḥ pacācīnām sarvān kālān sarvān puruṣān sarvāṇi ca vacanāni anuvartate . \\
\noindent \textbf{(P\_1,3.1.3)} bhavatiḥ punaḥ vartamānakālam ca eva ekatvam ca . \\
\noindent \textbf{(P\_1,3.1.3)} kā tarhi iyam vācoyuktiḥ : bhavati pacati , bhavati pakṣyati , bhavati apākṣīt iti . \\
\noindent \textbf{(P\_1,3.1.3)} eṣā eṣā vācoyuktiḥ : pacādayaḥ kriyāḥ bhavatikriyāyāḥ kartryaḥ bhavanti iti . \\
\noindent \textbf{(P\_1,3.1.3)} yadi api tāvat atra etat śakyate vaktum yatra anyā ca anyā ca kriyā yatra khalu sā eva kriyā tatra katham : bhavet api bhavet , syāt api syāt iti . \\
\noindent \textbf{(P\_1,3.1.3)} atra api anyatvam asti . \\
\noindent \textbf{(P\_1,3.1.3)} kutaḥ . \\
\noindent \textbf{(P\_1,3.1.3)} kālabhedāt sādhanabhedāt ca . \\
\noindent \textbf{(P\_1,3.1.3)} ekasya atra bhavateḥ bhavatiḥ sādhanam sarvakālaḥ ca pratyayaḥ . \\
\noindent \textbf{(P\_1,3.1.3)} aparasya bāhyam sādhanam vartamānakālaḥ ca pratyayaḥ . \\
\noindent \textbf{(P\_1,3.1.3)} yāvatā atra api anyatvam asti pacādayaḥ ca kriyāḥ bhavatikriyāyāḥ kartryaḥ bhavanti iti astu ayam kartṛsādhanaḥ : bhavati iti bhāvaḥ iti . \\
\noindent \textbf{(P\_1,3.1.3)} kim kṛtam bhavati . \\
\noindent \textbf{(P\_1,3.1.3)} vipratiṣiddhānām dhātusañjñā siddhā bhavati . \\
\noindent \textbf{(P\_1,3.1.3)} bhavet vipratiṣiddhānām dhātusañjñā siddhā syāt prātipadikānām api prāpnoti : vṛkṣaḥ , plakṣaḥ iti . \\
\noindent \textbf{(P\_1,3.1.3)} kim kāraṇam . \\
\noindent \textbf{(P\_1,3.1.3)} etāni api hi bhavanti . \\
\noindent \textbf{(P\_1,3.1.3)} evam tarhi karmasādhanaḥ bhaviṣyati : bhāvyate yaḥ saḥ bhāvaḥ iti . \\
\noindent \textbf{(P\_1,3.1.3)} kriyā ca eva hi bhāvyate svabhāvasiddham tu dravyam . \\
\noindent \textbf{(P\_1,3.1.3)} evam api bhavet keṣām cit na syāt yāni na bhāvyante . \\
\noindent \textbf{(P\_1,3.1.3)} ye tu ete sambandhiśabdāḥ teṣām prāpnoti : mātā pitā bhrātā iti . \\
\noindent \textbf{(P\_1,3.1.3)} sarvathā vayam prātipadikaparyudāsāt na mucyāmahe . \\
\noindent \textbf{(P\_1,3.1.3)} paṭhiṣyati hi ācāryaḥ : bhūvādipāṭhaḥ prātipadikāṇapayatyādinivṛttyarthaḥ iti . \\
\noindent \textbf{(P\_1,3.1.3)} yāvatā paṭhiṣyati pacādayaḥ ca kriyāḥ bhavatikriyāyāḥ kartryaḥ bhavanti iti astu ayam kartṛsādhanaḥ : bhavati iti bhāvaḥ . \\
\noindent \textbf{(P\_1,3.1.3)} kim vaktavyam etat . \\
\noindent \textbf{(P\_1,3.1.3)} na hi . \\
\noindent \textbf{(P\_1,3.1.3)} katham anucyamānam gaṃsyate . \\
\noindent \textbf{(P\_1,3.1.3)} etena eva abhihitam sūtreṇa bhūvādayaḥ dhātavaḥ iti . \\
\noindent \textbf{(P\_1,3.1.3)} katham . \\
\noindent \textbf{(P\_1,3.1.3)} na idam ādigrahaṇam . \\
\noindent \textbf{(P\_1,3.1.3)} vadeḥ ayam auṇādikaḥ iñ kartṛsādhanaḥ : bhuvam vadanti iti bhūvādayaḥ iti . \\
\noindent \textbf{(P\_1,3.1.3)} bhāvavacane tadarthapratyayapratiṣedhaḥ . \\
\noindent \textbf{(P\_1,3.1.3)} bhāvavacane dhātau tadarthasya pratyayasya pratiṣedhaḥ vaktavyaḥ : śiśye iti . \\
\noindent \textbf{(P\_1,3.1.3)} kim ca syāt . \\
\noindent \textbf{(P\_1,3.1.3)} aśiti iti āttvam prasajyeta . \\
\noindent \textbf{(P\_1,3.1.3)} tat hi dhātoḥ vihitam . \\
\noindent \textbf{(P\_1,3.1.3)} itaretarāśrayam ca pratyaye bhāvavacanatvam tasmāt ca pratyayaḥ . \\
\noindent \textbf{(P\_1,3.1.3)} itaretarāśrayam ca bhavati . \\
\noindent \textbf{(P\_1,3.1.3)} kā itaretarāśrayatā . \\
\noindent \textbf{(P\_1,3.1.3)} pratyaye bhāvavacanatvam tasmāt ca pratyayaḥ . \\
\noindent \textbf{(P\_1,3.1.3)} utpanne hi pratyaye bhāvavacanatvam gamyate saḥ ca tāvat bhāvavacanāt utpannaḥ . \\
\noindent \textbf{(P\_1,3.1.3)} tat etat itaretarāśrayam bhavati . \\
\noindent \textbf{(P\_1,3.1.3)} itaretarāśrayāṇi ca na prakalpante . \\
\noindent \textbf{(P\_1,3.1.3)} siddham tu nityaśabdatvāt anāśritya bhāvavacanatvam pratyayaḥ . \\
\noindent \textbf{(P\_1,3.1.3)} siddham etat . \\
\noindent \textbf{(P\_1,3.1.3)} katham . \\
\noindent \textbf{(P\_1,3.1.3)} nityāḥ śabdāḥ . \\
\noindent \textbf{(P\_1,3.1.3)} nityeṣu ca śabdeṣu anāśritya bhāvavacanatvam pratyayaḥ utpadyate . \\
\noindent \textbf{(P\_1,3.1.3)} prathamabhāvagrahaṇam ca . \\
\noindent \textbf{(P\_1,3.1.3)} prathamabhāvagrahaṇam ca kartavyam . \\
\noindent \textbf{(P\_1,3.1.3)} prathamam yaḥ bhāvam āha iti . \\
\noindent \textbf{(P\_1,3.1.3)} kutaḥ punaḥ prāthamyam . \\
\noindent \textbf{(P\_1,3.1.3)} kim śabdataḥ āhosvit arthataḥ . \\
\noindent \textbf{(P\_1,3.1.3)} kim ca ataḥ . \\
\noindent \textbf{(P\_1,3.1.3)} yadi śabdataḥ sanādīnām dhātusañjñā na prāpnoti : putrīyati vastrīyati iti . \\
\noindent \textbf{(P\_1,3.1.3)} atha arthataḥ siddhā sanādīnām dhātusañjñā saḥ eva tu doṣaḥ bhavati : bhāvavacane tadarthapratyayapratiṣedhaḥ iti . \\
\noindent \textbf{(P\_1,3.1.3)} evam tarhi na eva arthataḥ na eva śabdataḥ . \\
\noindent \textbf{(P\_1,3.1.3)} kim tarhi . \\
\noindent \textbf{(P\_1,3.1.3)} abhidhānataḥ . \\
\noindent \textbf{(P\_1,3.1.3)} sumadhyame abhidhāne yaḥ prathamam bhāvam āha. \\
\noindent \textbf{(P\_1,3.1.4)} iha ye eva bhāvavacane dhātau doṣāḥ te eva kriyāvacane api . \\
\noindent \textbf{(P\_1,3.1.4)} tatra te eva parihārāḥ . \\
\noindent \textbf{(P\_1,3.1.4)} tatra idam aparihṛtam : astibhavatividyatīnām dhātutvam iti . \\
\noindent \textbf{(P\_1,3.1.4)} tasya parihāraḥ . \\
\noindent \textbf{(P\_1,3.1.4)} kām punaḥ kriyām bhavān matvā āha astibhavatividyatīnām dhātusañjñā na prāpnoti iti . \\
\noindent \textbf{(P\_1,3.1.4)} kim yat tat devadattaḥ kaṃsapātryām pāṇinā odanam bhuṅkte iti . \\
\noindent \textbf{(P\_1,3.1.4)} na brūmaḥ kārakāṇi kriyā iti . \\
\noindent \textbf{(P\_1,3.1.4)} kim tarhi . \\
\noindent \textbf{(P\_1,3.1.4)} kārakāṇām pravṛttiviśeṣaḥ kriyā . \\
\noindent \textbf{(P\_1,3.1.4)} anyathā ca kārakāṇi śuṣkaudane pravartante anyathā ca māṃsaudane . \\
\noindent \textbf{(P\_1,3.1.4)} yadi evam siddhā astibhavatividyatīnām dhātusañjñā . \\
\noindent \textbf{(P\_1,3.1.4)} anyathā hi kārakāṇi astau pravartante anyathā hi mriyatau . \\
\noindent \textbf{(P\_1,3.1.4)} ṣaṭ bhāvavikārāḥ iti ha sma āha bhagavān vārṣyāyaṇiḥ : jāyate asti vipariṇamate vardhate apakṣīyate vinaśyati iti . \\
\noindent \textbf{(P\_1,3.1.4)} sarvathā sthitaḥ iti atra dhātusañjñā na prāpnoti . \\
\noindent \textbf{(P\_1,3.1.4)} bāhyaḥ hi ebhyaḥ tiṣṭhatiḥ . \\
\noindent \textbf{(P\_1,3.1.4)} evam tarhi kriyāyāḥ kriyā nivartikā bhavati dravyam dravyasya nivartakam . \\
\noindent \textbf{(P\_1,3.1.4)} evam hi kaḥ cit kam cit pṛcchati . \\
\noindent \textbf{(P\_1,3.1.4)} kimavasthaḥ devadattasya vyādhiḥ iti . \\
\noindent \textbf{(P\_1,3.1.4)} saḥ āha : vardhate iti . \\
\noindent \textbf{(P\_1,3.1.4)} aparaḥ āha : apakṣīyate iti . \\
\noindent \textbf{(P\_1,3.1.4)} aparaḥ āha : sthitaḥ iti . \\
\noindent \textbf{(P\_1,3.1.4)} sthitaḥ iti ukte vardhateḥ ca apakṣīyateḥ ca nivṛttiḥ bhavati . \\
\noindent \textbf{(P\_1,3.1.4)} atha vā na antareṇa kriyām bhūtabhaviṣyadvartamānāḥ kālāḥ vyajyante . \\
\noindent \textbf{(P\_1,3.1.4)} astyādibhiḥ ca api bhūtabhaviṣyadvartamānāḥ kālāḥ vyajyante . \\
\noindent \textbf{(P\_1,3.1.4)} atha vā na anyat prṣṭena anyat ākhyeyam . \\
\noindent \textbf{(P\_1,3.1.4)} tena na bhaviṣyati kim karoti asti iti . \\
\noindent \textbf{(P\_1,3.1.5)} atha yadi eva kriyāvacanaḥ dhātuḥ iti eṣaḥ pakṣaḥ atha api bhāvavacanaḥ dhātuḥ kim gatam etat iyatā sūtreṇa āhosvit anyatarasmin pakṣe bhūyaḥ sūtram kartavyam . \\
\noindent \textbf{(P\_1,3.1.5)} gatam iti āha . \\
\noindent \textbf{(P\_1,3.1.5)} katham . \\
\noindent \textbf{(P\_1,3.1.5)} ayam ādiśabdaḥ asti eva vyavasthāyām vartate . \\
\noindent \textbf{(P\_1,3.1.5)} tat yathā : devadattādīn samupaviṣṭān āha : devadattādayaḥ ānīyantām iti . \\
\noindent \textbf{(P\_1,3.1.5)} te utthāpya ānīyante . \\
\noindent \textbf{(P\_1,3.1.5)} asti prakāre vartate . \\
\noindent \textbf{(P\_1,3.1.5)} tat yathā : devadattādayaḥ āḍhyāḥ abhirūpāḥ darśanīyāḥ pakṣavantaḥ . \\
\noindent \textbf{(P\_1,3.1.5)} devadattaprakārāḥ iti gamyate . \\
\noindent \textbf{(P\_1,3.1.5)} pratyekam ca ādiśabdaḥ parisamāpyate . \\
\noindent \textbf{(P\_1,3.1.5)} bhvādayaḥ iti ca vādayaḥ iti ca . \\
\noindent \textbf{(P\_1,3.1.5)} tat yadā tāvat kriyāvacanaḥ dhātuḥ iti eṣaḥ pakṣaḥ tadā bhū iti atra yaḥ ādiśabdaḥ saḥ vyavasthāyām vartate vā iti atra yaḥ ādiśabdaḥ saḥ prakāre . \\
\noindent \textbf{(P\_1,3.1.5)} bhū iti evamādayaḥ vā iti evamprakārāḥ iti . \\
\noindent \textbf{(P\_1,3.1.5)} yadā tu bhāvavacanaḥ dhātuḥ iti eṣaḥ pakṣaḥ tadā vā iti atra yaḥ ādiśabdaḥ saḥ vyavasthāyām bhū iti atra yaḥ ādiśabdaḥ saḥ prakāre . \\
\noindent \textbf{(P\_1,3.1.5)} vā iti evamādayaḥ bhū iti evamprakārāḥ iti . \\
\noindent \textbf{(P\_1,3.1.5)} yadi tarhi lakṣaṇam kriyate na idānīm pāṭhaḥ kartavyaḥ . \\
\noindent \textbf{(P\_1,3.1.5)} kartavyaḥ ca . \\
\noindent \textbf{(P\_1,3.1.5)} kim prayojanam . \\
\noindent \textbf{(P\_1,3.1.5)} bhūvādipāṭhaḥ prātipadikāṇapayatyādinivṛttyarthaḥ . \\
\noindent \textbf{(P\_1,3.1.5)} bhūvādipāṭhaḥ kartavyaḥ . \\
\noindent \textbf{(P\_1,3.1.5)} kim prayojanam . \\
\noindent \textbf{(P\_1,3.1.5)} prātipadikāṇapayatyādinivṛttyarthaḥ . \\
\noindent \textbf{(P\_1,3.1.5)} prātipadikanivṛttyarthaḥ āṇapayatyādinivṛttyarthaḥ ca . \\
\noindent \textbf{(P\_1,3.1.5)} ke punaḥ āṇapayatyādayaḥ . \\
\noindent \textbf{(P\_1,3.1.5)} āṇapayati vaṭṭati vaḍḍhati iti . \\
\noindent \textbf{(P\_1,3.1.5)} svarānubandhajñāpanāya ca . \\
\noindent \textbf{(P\_1,3.1.5)} svarānubandhajñāpanāya ca pāṭhaḥ kartavyaḥ : svarān anubandhān ca jñāsyāmi iti . \\
\noindent \textbf{(P\_1,3.1.5)} na hi antareṇa pātham svarāḥ anubandhāḥ vā śakyāḥ vijñātum . \\
\noindent \textbf{(P\_1,3.1.5)} ye tu ete nyāyyavikaraṇāḥ udāttāḥ ananubandhakāḥ paṭhyante eteṣām pāṭhaḥ śakyaḥ akartum . \\
\noindent \textbf{(P\_1,3.1.5)} eteṣām api avaśyam āṇapayatyādinivṛttyarthaḥ pāṭhaḥ kartavyaḥ . \\
\noindent \textbf{(P\_1,3.1.5)} na kartavyaḥ . \\
\noindent \textbf{(P\_1,3.1.5)} śiṣṭaprayogāt āṇapayatyādīnām nivṛttiḥ bhaviṣyati . \\
\noindent \textbf{(P\_1,3.1.5)} saḥ ca avaśyam śiṣṭaprayogaḥ upāsyaḥ ye api paṭhyante teṣām api viparyāsanivṛttyarthaḥ . \\
\noindent \textbf{(P\_1,3.1.5)} loke hi kṛṣyarthe kasim prayuñjate dṛśyarthe ca diśim . \\
\noindent \textbf{(P\_1,3.2.1)} upadeśe iti kimartham . \\
\noindent \textbf{(P\_1,3.2.1)} abhre āŚ apaḥ : uddeśe yaḥ anunāsikaḥ tasya mā bhūt iti . \\
\noindent \textbf{(P\_1,3.2.1)} kaḥ punaḥ uddeśopadeśayoḥ viśeṣaḥ . \\
\noindent \textbf{(P\_1,3.2.1)} pratyakṣam ākhyānam upadeśaḥ , guṇaiḥ prāpaṇam uddeśaḥ . \\
\noindent \textbf{(P\_1,3.2.1)} pratyakṣam tāvat ākhyānam upadeśaḥ . \\
\noindent \textbf{(P\_1,3.2.1)} tat yathā : agojñāya kaḥ cit gām sakhthani karṇe vā gṛhītvā upadiśati : ayam gauḥ iti . \\
\noindent \textbf{(P\_1,3.2.1)} saḥ pratyakṣam ākhyātam āha : upadiṣṭaḥ me gauḥ iti . \\
\noindent \textbf{(P\_1,3.2.1)} guṇaiḥ prāpaṇam uddeśaḥ . \\
\noindent \textbf{(P\_1,3.2.1)} tat yathā : kaḥ cit kam cit āha : devadattam me bhavān uddiśatu iti . \\
\noindent \textbf{(P\_1,3.2.1)} saḥ ihasthaḥ pāṭaliputrastham devadattam uddiśati : aṅgadī kuṇḍalī kirīṭī vyūḍhoraskaḥ vṛttabāhuḥ lohitākṣaḥ tuṅganāsaḥ citrābharaṇaḥ īdṛśaḥ devadattaḥ iti . \\
\noindent \textbf{(P\_1,3.2.1)} saḥ guṇaiḥ prāpyamāṇam āha : uddiṣṭaḥ me devadattaḥ iti . \\
\noindent \textbf{(P\_1,3.2.2)} itsañjñāyām sarvaprasaṅgaḥ aviśeṣāt . \\
\noindent \textbf{(P\_1,3.2.2)} itsañjñāyām sarvaprasaṅgaḥ . \\
\noindent \textbf{(P\_1,3.2.2)} sarvasya anunāsikasya itsañjñā prāpnoti . \\
\noindent \textbf{(P\_1,3.2.2)} asya api prāpnoti : abhre āŚ apaḥ . \\
\noindent \textbf{(P\_1,3.2.2)} kim kāraṇam . \\
\noindent \textbf{(P\_1,3.2.2)} aviśeṣāt . \\
\noindent \textbf{(P\_1,3.2.2)} na hi kaḥ cit viśeṣaḥ : upādīyate evañjātīyakasya anunāsikasya itsañjñā bhavati iti . \\
\noindent \textbf{(P\_1,3.2.2)} anupādīyamāne viśeṣe sarvaprasaṅgaḥ . \\
\noindent \textbf{(P\_1,3.2.2)} kim ucyate anupādīyamāne viśeṣe iti . \\
\noindent \textbf{(P\_1,3.2.2)} katham na nāma upādīyate yadā upadeśe iti ucyate . \\
\noindent \textbf{(P\_1,3.2.2)} lakṣaṇena hi upadeśaḥ . \\
\noindent \textbf{(P\_1,3.2.2)} saṅkīrṇau uddeśopadeśau . \\
\noindent \textbf{(P\_1,3.2.2)} pratyakṣam ākhyānam uddeśaḥ guṇaiḥ ca prāpaṇam upadeśaḥ . \\
\noindent \textbf{(P\_1,3.2.2)} pratyakṣam tāvat ākhyānam uddeśaḥ . \\
\noindent \textbf{(P\_1,3.2.2)} tat yathā : kaḥ cit kam cit āha : anuvākam me bhavān uddiśatu iti . \\
\noindent \textbf{(P\_1,3.2.2)} saḥ tasmai ācaṣṭe : iṣetvakam adhīṣva . \\
\noindent \textbf{(P\_1,3.2.2)} śannodevīyam adhīṣva iti . \\
\noindent \textbf{(P\_1,3.2.2)} saḥ pratyakṣam ākhyātam āha : uddiṣṭaḥ me anuvākaḥ . \\
\noindent \textbf{(P\_1,3.2.2)} tam adhyeṣye iti . \\
\noindent \textbf{(P\_1,3.2.2)} guṇaiḥ ca prāpaṇam upadeśaḥ . \\
\noindent \textbf{(P\_1,3.2.2)} tat yathā : kaḥ cit kam cit āha : grāmantaram gamiṣyāmi . \\
\noindent \textbf{(P\_1,3.2.2)} panthānam me bhavān upadiśatu iti . \\
\noindent \textbf{(P\_1,3.2.2)} saḥ tasmai ācaṣṭe : amuṣmin avakāśe hastadakṣiṇaḥ grahītavyaḥ , amuṣmin hastavāmaḥ iti . \\
\noindent \textbf{(P\_1,3.2.2)} saḥ guṇaiḥ prāpyamāṇam āha : upadiṣṭaḥ me panthāḥ iti . \\
\noindent \textbf{(P\_1,3.2.2)} evam etau saṅkīrṇau uddeśopadeśau . \\
\noindent \textbf{(P\_1,3.2.2)} evam tarhi itkāryābhāvāt itsañjñā na bhaviṣyati . \\
\noindent \textbf{(P\_1,3.2.2)} nanu ca lopaḥ eva itkāryam syāt . \\
\noindent \textbf{(P\_1,3.2.2)} akāryam lopaḥ . \\
\noindent \textbf{(P\_1,3.2.2)} iha hi śabdasya dvyarthaḥ upadeśaḥ . \\
\noindent \textbf{(P\_1,3.2.2)} kāryārthaḥ vā bhavati upadeśaḥ śravaṇārthaḥ vā . \\
\noindent \textbf{(P\_1,3.2.2)} kāryam ca iha na asti . \\
\noindent \textbf{(P\_1,3.2.2)} kārye ca asati yadi śravaṇam api na syāt upadeśaḥ anarthakaḥ syāt . \\
\noindent \textbf{(P\_1,3.2.2)} idam asti itkāryam : abhre āŚ aṭitaḥ : anantaralakṣaṇāyām itsañjñāyām satyām āditaḥ ca iti iṭpratiṣedhaḥ prasajyeta . \\
\noindent \textbf{(P\_1,3.2.2)} siddham tu upadeśane anunāsikavacanāt . \\
\noindent \textbf{(P\_1,3.2.2)} siddham etat katham . \\
\noindent \textbf{(P\_1,3.2.2)} upadeśane yaḥ anunāsikaḥ saḥ itsañjñaḥ bhavati iti vaktavyam . \\
\noindent \textbf{(P\_1,3.2.2)} kim punaḥ upadeśanam . \\
\noindent \textbf{(P\_1,3.2.2)} śāstram . \\
\noindent \textbf{(P\_1,3.2.2)} sidhyati . \\
\noindent \textbf{(P\_1,3.2.2)} sūtram tarhi bhidyate . \\
\noindent \textbf{(P\_1,3.2.2)} yathānyāsam eva astu . \\
\noindent \textbf{(P\_1,3.2.2)} nanu ca uktam itsañjñāyām sarvaprasaṅgaḥ aviśeṣāt iti . \\
\noindent \textbf{(P\_1,3.2.2)} na eṣaḥ doṣaḥ . \\
\noindent \textbf{(P\_1,3.2.2)} upadeśaḥ iti ghañ ayam karaṇasādhanaḥ . \\
\noindent \textbf{(P\_1,3.2.2)} na sidhyati . \\
\noindent \textbf{(P\_1,3.2.2)} paratvāt lyuṭ prāpnoti . \\
\noindent \textbf{(P\_1,3.2.2)} na brūmaḥ akartari ca kārake sañjñāyām iti . \\
\noindent \textbf{(P\_1,3.2.2)} kim tarhi . \\
\noindent \textbf{(P\_1,3.2.2)} halaḥ ca iti . \\
\noindent \textbf{(P\_1,3.2.2)} tatra api sañjñāyām iti vartate . \\
\noindent \textbf{(P\_1,3.2.2)} na ca eṣā sañjñā . \\
\noindent \textbf{(P\_1,3.2.2)} prāyavacanāt asañjñāyām api bhaviṣyati . \\
\noindent \textbf{(P\_1,3.2.2)} prāyavacanāt sañjñāyām eva syāt vā na vā . \\
\noindent \textbf{(P\_1,3.2.2)} na hi upādheḥ upādhiḥ bhavati viśeṣaṇasya vā viśeṣaṇam . \\
\noindent \textbf{(P\_1,3.2.2)} yadi na upādheḥ upādhiḥ bhavati viśeṣaṇasya vā viśeṣaṇam kalyāṇyādīnām inaṅ kulaṭāyāḥ vā inaṅ vibhāṣā na prāpnoti . \\
\noindent \textbf{(P\_1,3.2.2)} inaṅ eva atra pradhānam . \\
\noindent \textbf{(P\_1,3.2.2)} vihitaḥ pratyayaḥ prakṛtaḥ ca anuvartate . \\
\noindent \textbf{(P\_1,3.2.2)} iha tarhi : vākinādīnām kuk ca putrāt anyatarasyām iti kuk vibhāṣā na prāpnoti . \\
\noindent \textbf{(P\_1,3.2.2)} atra api kuk eva pradhānam . \\
\noindent \textbf{(P\_1,3.2.2)} vihitaḥ pratyayaḥ prakṛtaḥ ca anuvartate . \\
\noindent \textbf{(P\_1,3.2.2)} evam na ca idam akṛtam bhavati na upādheḥ upādhiḥ bhavati viśeṣaṇasya vā viśeṣaṇam iti na ca kaḥ cit doṣaḥ bhavati . \\
\noindent \textbf{(P\_1,3.2.2)} evam ca kṛtvā ghañ na prāpnoti . \\
\noindent \textbf{(P\_1,3.2.2)} evam tarhi kṛtyalyuṭaḥ bahulam iti evam atra ghañ bhaviṣyati . \\
\noindent \textbf{(P\_1,3.3.1)} halantye sarvprasaṅgaḥ sarvāntyatvāt . \\
\noindent \textbf{(P\_1,3.3.1)} halantye sarvprasaṅgaḥ . \\
\noindent \textbf{(P\_1,3.3.1)} sarvasya halaḥ itsañjñā prāpnoti . \\
\noindent \textbf{(P\_1,3.3.1)} kim kāraṇam . \\
\noindent \textbf{(P\_1,3.3.1)} sarvāntyatvāt . \\
\noindent \textbf{(P\_1,3.3.1)} sarvaḥ hi hal tam tam avadhim prati antyaḥ bhavati . \\
\noindent \textbf{(P\_1,3.3.1)} siddham tu vyavasitāntyatvāt . \\
\noindent \textbf{(P\_1,3.3.1)} siddham etat . \\
\noindent \textbf{(P\_1,3.3.1)} katham . \\
\noindent \textbf{(P\_1,3.3.1)} vyavasitāntyatvāt . \\
\noindent \textbf{(P\_1,3.3.1)} vyavasitāntyaḥ hal itsañjñaḥ bhavati iti vaktavyam . \\
\noindent \textbf{(P\_1,3.3.1)} ke punaḥ vyavasitāḥ . \\
\noindent \textbf{(P\_1,3.3.1)} dhātuprātipadikapratyayanipātāgamādeśāḥ . \\
\noindent \textbf{(P\_1,3.3.1)} sidhyati . \\
\noindent \textbf{(P\_1,3.3.1)} sūtram tarhi bhidyate . \\
\noindent \textbf{(P\_1,3.3.1)} yathānyāsam eva astu . \\
\noindent \textbf{(P\_1,3.3.1)} nanu ca uktam halantye sarvprasaṅgaḥ sarvāntyatvāt iti . \\
\noindent \textbf{(P\_1,3.3.1)} na eṣaḥ doṣaḥ . \\
\noindent \textbf{(P\_1,3.3.1)} āha ayam hal antyam itsañjñam bhavati iti . \\
\noindent \textbf{(P\_1,3.3.1)} sarvaḥ ca hal tam tam avadhim prati antyaḥ bhavati . \\
\noindent \textbf{(P\_1,3.3.1)} tatra prakarṣagatiḥ vijñāsyate : sādhīyaḥ yaḥ antyaḥ iti . \\
\noindent \textbf{(P\_1,3.3.1)} kaḥ ca sādhīyaḥ . \\
\noindent \textbf{(P\_1,3.3.1)} yaḥ vyavasitāntyaḥ . \\
\noindent \textbf{(P\_1,3.3.1)} atha vā sāpekṣaḥ ayam nirdeśaḥ kriyate . \\
\noindent \textbf{(P\_1,3.3.1)} na ca anyat kim cit apekṣyam asti . \\
\noindent \textbf{(P\_1,3.3.1)} te vyavasitam eva apekṣiṣyāmahe . \\
\noindent \textbf{(P\_1,3.3.2)} lakārasya anubandhājñāpitatvāt halgrahaṇāprasiddhiḥ . \\
\noindent \textbf{(P\_1,3.3.2)} lakārasya anubandhatvena ajñāpitatvāt halgrahaṇāprasiddhiḥ . \\
\noindent \textbf{(P\_1,3.3.2)} hal antyam itsañjñam bhavati iti ucyate . \\
\noindent \textbf{(P\_1,3.3.2)} lakārasya eva tāvat itsañjñā na prāpnoti . \\
\noindent \textbf{(P\_1,3.3.2)} siddham tu lakāranirdeśāt . \\
\noindent \textbf{(P\_1,3.3.2)} siddham etat . \\
\noindent \textbf{(P\_1,3.3.2)} katham . \\
\noindent \textbf{(P\_1,3.3.2)} lakāranirdeśaḥ kartavyaḥ . \\
\noindent \textbf{(P\_1,3.3.2)} hal antyam itsañjñam bhavati lakāraḥ ca iti vaktavyam . \\
\noindent \textbf{(P\_1,3.3.2)} ekaśeṣanirdeśāt vā . \\
\noindent \textbf{(P\_1,3.3.2)} atha vā ekaśeṣanirdeśaḥ ayam . \\
\noindent \textbf{(P\_1,3.3.2)} hal ca hal ca hal . \\
\noindent \textbf{(P\_1,3.3.2)} hal antyam itsañjñam bhavati iti . \\
\noindent \textbf{(P\_1,3.3.2)} atha vā ḷkārasya eva idam guṇabhūtasya grahaṇam . \\
\noindent \textbf{(P\_1,3.3.2)} tatra upadeśe ac anunāsika it iti itsañjñā bhaviṣyati . \\
\noindent \textbf{(P\_1,3.3.2)} atha vā ācāryapravṛttiḥ jñāpayati bhavati lakārasya itsañjñā iti yat ayam ṇalam litam karoti . \\
\noindent \textbf{(P\_1,3.3.3)} prātipadikapratiṣedhaḥ akṛttaddhite . \\
\noindent \textbf{(P\_1,3.3.3)} akṛttaddhitāntasya prātipadikasya pratiṣedhaḥ vaktavyaḥ . \\
\noindent \textbf{(P\_1,3.3.3)} udaśvit śakrṭ iti . \\
\noindent \textbf{(P\_1,3.3.3)} akṛttaddhitāntasya iti kimartham . \\
\noindent \textbf{(P\_1,3.3.3)} kumbhakāraḥ nagarakāraḥ aupagavaḥ kāpaṭavaḥ . \\
\noindent \textbf{(P\_1,3.3.3)} idarthābhāvāt siddham . \\
\noindent \textbf{(P\_1,3.3.3)} itkāryābhāvāt atra itsañjña na bhaviṣyati . \\
\noindent \textbf{(P\_1,3.3.3)} idam asti itkāryam titsvaritam iti svaritatvam yathā syāt . \\
\noindent \textbf{(P\_1,3.3.3)} na etat asti . \\
\noindent \textbf{(P\_1,3.3.3)} pratyayagrahaṇam tatra codayiṣyati . \\
\noindent \textbf{(P\_1,3.3.3)} idam tarhi : rājā takṣā . \\
\noindent \textbf{(P\_1,3.3.3)} ñniti ādyudāttatvam yathā syāt . \\
\noindent \textbf{(P\_1,3.3.3)} ñniti iti ucyate . \\
\noindent \textbf{(P\_1,3.3.3)} tatra vyapavargābhāvāt na bhaviṣyati . \\
\noindent \textbf{(P\_1,3.3.3)} idam tarhi svaḥ . \\
\noindent \textbf{(P\_1,3.3.3)} upottamam riti eṣaḥ svaraḥ yathā syāt . \\
\noindent \textbf{(P\_1,3.3.3)} svaritakaraṇasāmarthyāt na bhaviṣyati . \\
\noindent \textbf{(P\_1,3.3.3)} nyaṅsvarau svritau iti . \\
\noindent \textbf{(P\_1,3.3.3)} iha tarhi antaḥ . \\
\noindent \textbf{(P\_1,3.3.3)} uttamaśabdaḥ triprabhṛtiṣu vartate . \\
\noindent \textbf{(P\_1,3.3.3)} na ca atra triprabhṛtayaḥ santi . \\
\noindent \textbf{(P\_1,3.3.3)} iha tarhi sanutaḥ . \\
\noindent \textbf{(P\_1,3.3.3)} upottamam riti iti eṣaḥ svaraḥ yathā syāt . \\
\noindent \textbf{(P\_1,3.3.3)} antodāttanipātanam kariṣyate . \\
\noindent \textbf{(P\_1,3.3.3)} saḥ ca nipātasvaraḥ ritsvarasya bādhakaḥ bhaviṣyati . \\
\noindent \textbf{(P\_1,3.3.3)} etat ca atra yuktam yat itkāryābhāvāt itsañjñā na syāt . \\
\noindent \textbf{(P\_1,3.3.3)} yatra itkāryam bhavati bhavati tatra itsañjñā . \\
\noindent \textbf{(P\_1,3.3.3)} tat yathā āgastyakauṇḍinyayoḥ agastikuṇḍinac . \\
\noindent \textbf{(P\_1,3.4)} vibhaktau tavargapratiṣedhaḥ ataddhite . \\
\noindent \textbf{(P\_1,3.4)} vibhaktau tavargapratiṣedhaḥ ataddhite iti vaktavyam . \\
\noindent \textbf{(P\_1,3.4)} iha mā bhūt . \\
\noindent \textbf{(P\_1,3.4)} kimaḥ at kve prepsan dīpyase kva ardhamāsāḥ iti . \\
\noindent \textbf{(P\_1,3.4)} saḥ tarhi pratiṣedhaḥ vaktavyaḥ . \\
\noindent \textbf{(P\_1,3.4)} na vaktavyaḥ . \\
\noindent \textbf{(P\_1,3.4)} ācāryapravṛttiḥ jñāpāyati na vibhaktau taddhite pratiṣedhaḥ bhavati iti yat ayam idamaḥ thamuḥ iti makārasye itsañjñāparitrāṇārtham ukāram anubandham karoti . \\
\noindent \textbf{(P\_1,3.4)} yadi etat jñāpyate idānīm iti atra api prāpnoti . \\
\noindent \textbf{(P\_1,3.4)} itkāryābhāvāt atra itsañjñā na bhaviṣyati . \\
\noindent \textbf{(P\_1,3.4)} idam asti itkāryam mit acaḥ antyāt paraḥ iti acām antyāt paraḥ yathā syāt . \\
\noindent \textbf{(P\_1,3.4)} iśbhāve kṛte na asti viśeṣaḥ mit acaḥ antyāt paraḥ iti vā paratve pratyayaḥ paraḥ iti vā . \\
\noindent \textbf{(P\_1,3.4)} saḥ eva tāvat iśbhāvaḥ na prāpnoti . \\
\noindent \textbf{(P\_1,3.4)} kim kāraṇam . \\
\noindent \textbf{(P\_1,3.4)} prāk diśaḥ pratyayeṣu iti ucyate . \\
\noindent \textbf{(P\_1,3.4)} kaḥ punaḥ arhati iśbhāvam prāg diśaḥ pratyayeṣu vaktum . \\
\noindent \textbf{(P\_1,3.4)} kim tarhi . \\
\noindent \textbf{(P\_1,3.4)} prāk diśaḥ artheṣu iśbhāvaḥ kiṃsarvanāmabahubhyaḥ advyādibhyaḥ pratyayotpattiḥ . \\
\noindent \textbf{(P\_1,3.4)} evam tarhi tadaḥ api ayam vaktavyaḥ . \\
\noindent \textbf{(P\_1,3.4)} tadaḥ ca mit acaḥ antyāt paratvena na sidhyati . \\
\noindent \textbf{(P\_1,3.4)} nanu ca atra api atve kṛte na asti viśeṣaḥ mit acaḥ antyāt paraḥ iti vā paratve pratyayaḥ paraḥ iti vā . \\
\noindent \textbf{(P\_1,3.4)} tat hi attvam na prāpnoti . \\
\noindent \textbf{(P\_1,3.4)} kim kāraṇam . \\
\noindent \textbf{(P\_1,3.4)} vibhaktau iti ucyate . \\
\noindent \textbf{(P\_1,3.4)} evam tarhi yakārāntaḥ dānīm kariṣyate . \\
\noindent \textbf{(P\_1,3.4)} kim yakāraḥ na śrūyate . \\
\noindent \textbf{(P\_1,3.4)} luptanirdiṣṭaḥ yakāraḥ . \\
\noindent \textbf{(P\_1,3.7.1)} cuñcupcaṇapoḥ cakārapratiṣedhaḥ . \\
\noindent \textbf{(P\_1,3.7.1)} cuñcupcaṇapoḥ cakārasya pratiṣedhaḥ vaktavyaḥ . \\
\noindent \textbf{(P\_1,3.7.1)} keśacuñcuḥ keśacaṇaḥ . \\
\noindent \textbf{(P\_1,3.7.1)} itkāryābhāvāt siddham . \\
\noindent \textbf{(P\_1,3.7.1)} itkāryābhāvāt atra itsañjñā na bhaviṣyati . \\
\noindent \textbf{(P\_1,3.7.1)} idam asti itkāryam citaḥ antaḥ udāttaḥ bhavati iti antodāttatvam yathā syāt . \\
\noindent \textbf{(P\_1,3.7.1)} pitkaraṇam idānīm kimartham syāt . \\
\noindent \textbf{(P\_1,3.7.1)} pitkaraṇam kimartham iti cet paryāyārtham . \\
\noindent \textbf{(P\_1,3.7.1)} pitkaraṇam kimartham iti cet paryāyārtham etat syāt . \\
\noindent \textbf{(P\_1,3.7.1)} evam tarhi yakārādī cuñcupcaṇapau . \\
\noindent \textbf{(P\_1,3.7.1)} kim yakāraḥ na śrūyate . \\
\noindent \textbf{(P\_1,3.7.1)} luptanirdiṣṭaḥ yakāraḥ . \\
\noindent \textbf{(P\_1,3.7.2)} iraḥ upasaṅkhyānam . \\
\noindent \textbf{(P\_1,3.7.2)} iraḥ upasaṅkhyānam kartavyam : rudhir : arudhat , arautsīt . \\
\noindent \textbf{(P\_1,3.7.2)} avayavagrahaṇāt siddham . \\
\noindent \textbf{(P\_1,3.7.2)} rephasya atra halantyam iti itsañjñā bhaviṣyati ikārasya upadeśe ac anunāsikaḥ iti . \\
\noindent \textbf{(P\_1,3.7.2)} avayavagrahaṇāt iti cet ididvidhiprasaṅgaḥ . \\
\noindent \textbf{(P\_1,3.7.2)} avayavagrahaṇāt iti cet ididvidhiprasaṅgaḥ prāpnoti . \\
\noindent \textbf{(P\_1,3.7.2)} bhettā chettā . \\
\noindent \textbf{(P\_1,3.7.2)} iditaḥ num dhātoḥ iti num prāpnoti . \\
\noindent \textbf{(P\_1,3.7.2)} yadi punaḥ ayam ididvidhiḥ kumbhīdhānyanyāyena vijñāyeta . \\
\noindent \textbf{(P\_1,3.7.2)} tat yathā . \\
\noindent \textbf{(P\_1,3.7.2)} kumbhīdhānyaḥ śrotriyaḥ iti ucyate . \\
\noindent \textbf{(P\_1,3.7.2)} yasya kumbhyām eva dhānyam saḥ kumbhīdhānyaḥ . \\
\noindent \textbf{(P\_1,3.7.2)} yasya punaḥ kumbhyām ca anyatra ca na asau kumbhīdhānyaḥ . \\
\noindent \textbf{(P\_1,3.7.2)} na ayam ididvidhiḥ kumbhīdhānyanyāyena śakyaḥ vijñātum . \\
\noindent \textbf{(P\_1,3.7.2)} iha hi doṣaḥ syāt . \\
\noindent \textbf{(P\_1,3.7.2)} ṭunadi nandathuḥ iti . \\
\noindent \textbf{(P\_1,3.7.2)} evam tarhi na evam vijñāyate ikāraḥ it yasya saḥ ayam idit tasya iditaḥ iti . \\
\noindent \textbf{(P\_1,3.7.2)} katham tarhi . \\
\noindent \textbf{(P\_1,3.7.2)} ikāraḥ eva it idit ididantasya iti . \\
\noindent \textbf{(P\_1,3.7.2)} atha vā ṝkārasya eva idam irtvabhūtasya grahaṇam . \\
\noindent \textbf{(P\_1,3.7.2)} tatra upadeśe ac anunāsikaḥ it iti itsañjñā bhaviṣyati . \\
\noindent \textbf{(P\_1,3.7.2)} atha vā ācāryapravṛttiḥ jñāpayati na evañjātīyakānām ididvidhiḥ bhavati iti yat ayam iritaḥ kān cit numanuṣaktān paṭhati . \\
\noindent \textbf{(P\_1,3.7.2)} ubundir niśāmane . \\
\noindent \textbf{(P\_1,3.7.2)} skandir gatiśoṣaṇayoḥ . \\
\noindent \textbf{(P\_1,3.7.2)} atha vā ācāryapravṛttiḥ jñāpayati irśabdasya itsañjñā bhavati iti yat ayam iritaḥ vā iti āha . \\
\noindent \textbf{(P\_1,3.7.2)} atha vā ante iti vartate . \\
\noindent \textbf{(P\_1,3.9.1)} tasyagrahaṇam kimartham . \\
\noindent \textbf{(P\_1,3.9.1)} itsañjñakaḥ pratinirdiśyate . \\
\noindent \textbf{(P\_1,3.9.1)} na etat asti prayojanam . \\
\noindent \textbf{(P\_1,3.9.1)} prakṛtam it iti vartate . \\
\noindent \textbf{(P\_1,3.9.1)} kva prakṛtam . \\
\noindent \textbf{(P\_1,3.9.1)} upadeśe ac anunāsikaḥ it iti . \\
\noindent \textbf{(P\_1,3.9.1)} tat vai prathamānirdiṣṭam ṣaṣṭhīnirdiṣṭena ca iha arthaḥ . \\
\noindent \textbf{(P\_1,3.9.1)} arthāt vibhaktivipariṇāmaḥ bhaviṣyati . \\
\noindent \textbf{(P\_1,3.9.1)} tat yathā . \\
\noindent \textbf{(P\_1,3.9.1)} uccāni devadattasya gṛhāṇi . \\
\noindent \textbf{(P\_1,3.9.1)} āmantrayasva enam . \\
\noindent \textbf{(P\_1,3.9.1)} devadattam iti gamyate . \\
\noindent \textbf{(P\_1,3.9.1)} devadattasya gāvaḥ aśvāḥ hiraṇyam iti . \\
\noindent \textbf{(P\_1,3.9.1)} āḍhyaḥ vaidhaveyaḥ . \\
\noindent \textbf{(P\_1,3.9.1)} devadattaḥ iti gamyate . \\
\noindent \textbf{(P\_1,3.9.1)} purastāt ṣaṣṭhīnirdiṣṭam sat arthāt dvitīyānirdiṣṭam prathamānirdiṣṭam ca bhavati . \\
\noindent \textbf{(P\_1,3.9.1)} evam iha api purastāt prathamānirdiṣṭam sat arthāt ṣaṣṭhīnirdiṣṭam bhaviṣyati . \\
\noindent \textbf{(P\_1,3.9.1)} idam tarhi prayojanam . \\
\noindent \textbf{(P\_1,3.9.1)} ye anekālaḥ itsañjñāḥ teṣām lopaḥ sarvādeśaḥ yathā syāt . \\
\noindent \textbf{(P\_1,3.9.1)} atha kriyamāṇe api ca tasyagrahaṇe katham iva lopaḥ sarvādeśaḥ labhyaḥ . \\
\noindent \textbf{(P\_1,3.9.1)} labhyaḥ iti āha . \\
\noindent \textbf{(P\_1,3.9.1)} kutaḥ . \\
\noindent \textbf{(P\_1,3.9.1)} vacanaprāmāṇyāt . \\
\noindent \textbf{(P\_1,3.9.1)} tasyagrahaṇasāmarthyāt . \\
\noindent \textbf{(P\_1,3.9.2)} itaḥ lope ṇalktvāniṣṭhāsu upasaṅkhyānam itpratiṣedhāt . \\
\noindent \textbf{(P\_1,3.9.2)} itaḥ lope ṇalktvāniṣṭhāsu upasaṅkhyānam kartavyam . \\
\noindent \textbf{(P\_1,3.9.2)} ṇal . \\
\noindent \textbf{(P\_1,3.9.2)} aham papaca . \\
\noindent \textbf{(P\_1,3.9.2)} ktvā . \\
\noindent \textbf{(P\_1,3.9.2)} devitvā sevitvā . \\
\noindent \textbf{(P\_1,3.9.2)} niṣṭhā . \\
\noindent \textbf{(P\_1,3.9.2)} śayitaḥ śayitavān . \\
\noindent \textbf{(P\_1,3.9.2)} kim punaḥ kāraṇam na sidhyati . \\
\noindent \textbf{(P\_1,3.9.2)} itpratiṣedhāt . \\
\noindent \textbf{(P\_1,3.9.2)} pratiṣidhyate atra itsañjñā . \\
\noindent \textbf{(P\_1,3.9.2)} ṇal uttamaḥ ṇit vā bhavati . \\
\noindent \textbf{(P\_1,3.9.2)} ktvā seṭ na kit bhavati . \\
\noindent \textbf{(P\_1,3.9.2)} niṣṭhā seṭ na kit bhavati iti . \\
\noindent \textbf{(P\_1,3.9.2)} siddham tu ṇalādīnām grahaṇapratiṣedhāt . siddham etat . \\
\noindent \textbf{(P\_1,3.9.2)} katham . \\
\noindent \textbf{(P\_1,3.9.2)} ṇalādīnām grahaṇāni pratiṣidhante . \\
\noindent \textbf{(P\_1,3.9.2)} ṇal uttamaḥ vā ṇidgrahaṇena gṛhyate . \\
\noindent \textbf{(P\_1,3.9.2)} ktvā seṭ na kidgrahaṇena gṛhyate . \\
\noindent \textbf{(P\_1,3.9.2)} niṣṭhā seṭ na kidgrahaṇena gṛhyate iti . \\
\noindent \textbf{(P\_1,3.9.2)} nirdiṣṭalopāt vā . \\
\noindent \textbf{(P\_1,3.9.2)} nirdiṣṭalopāt vā siddham eva . \\
\noindent \textbf{(P\_1,3.9.2)} atha vā nirdiṣṭasya ayam lopaḥ kriyate . \\
\noindent \textbf{(P\_1,3.9.2)} tasmāt siddham etat . \\
\noindent \textbf{(P\_1,3.9.2)} tatra tusmānām pratiṣedhaḥ . \\
\noindent \textbf{(P\_1,3.9.2)} tatra tusmānām pratiṣedhaḥ vaktavyaḥ . \\
\noindent \textbf{(P\_1,3.9.2)} tasmāt tasmin yasmāt yasmin vṛkṣāḥ plakṣāḥ acinavam asunavam akaravam . \\
\noindent \textbf{(P\_1,3.9.2)} na vā uccāraṇasāmarthyāt . \\
\noindent \textbf{(P\_1,3.9.2)} na vā vaktavyaḥ . \\
\noindent \textbf{(P\_1,3.9.2)} kim kāraṇam . \\
\noindent \textbf{(P\_1,3.9.2)} uccāraṇasāmarthyāt atra lopaḥ na bhaviṣyati . \\
\noindent \textbf{(P\_1,3.9.2)} anubandhalope bhāvābhāvayoḥ vipratiṣedhāt aprasiddhiḥ . anubandhalope bhāvābhāvayoḥ virodhāt aprasiddhiḥ . \\
\noindent \textbf{(P\_1,3.9.2)} na jñāyate kena abhiprāyeṇa prasajati kena nivṛttim karoti iti . \\
\noindent \textbf{(P\_1,3.9.2)} siddham tu apavādanyāyena . \\
\noindent \textbf{(P\_1,3.9.2)} siddham etat . \\
\noindent \textbf{(P\_1,3.9.2)} katham . \\
\noindent \textbf{(P\_1,3.9.2)} apavādanyāyena . \\
\noindent \textbf{(P\_1,3.9.2)} kim punaḥ iha tathā yathā utsargāpavādau . \\
\noindent \textbf{(P\_1,3.9.2)} bhāvaḥ hi kāryāṛthaḥ nanyārthaḥ lopaḥ . \\
\noindent \textbf{(P\_1,3.9.2)} kāryam kariṣyāmi iti anubandhaḥ āsajyate kāryād anyan mā bhūt iti lopaḥ . \\
\noindent \textbf{(P\_1,3.9.3)} atha yasya anubandhaḥ āsajyate kim saḥ tasya ekāntaḥ bhavati āhosvit anekāntaḥ . \\
\noindent \textbf{(P\_1,3.9.3)} ekāntaḥ tatra upalabdheḥ . \\
\noindent \textbf{(P\_1,3.9.3)} ekāntaḥ iti āha . \\
\noindent \textbf{(P\_1,3.9.3)} kutaḥ . \\
\noindent \textbf{(P\_1,3.9.3)} tatra upalabdheḥ . \\
\noindent \textbf{(P\_1,3.9.3)} tatrasthaḥ hi asau upalabhyate . \\
\noindent \textbf{(P\_1,3.9.3)} tat yathā vṛkṣasthā śākhā vṛkṣaikāntā upalabhyate . \\
\noindent \textbf{(P\_1,3.9.3)} tatra asarūpasarvādeśadāppratiṣedhe pṛthaktvanirdeśaḥ anākārāntatvāt . \\
\noindent \textbf{(P\_1,3.9.3)} tatra asarūpavidhau doṣaḥ bhavati . \\
\noindent \textbf{(P\_1,3.9.3)} karmaṇi aṇ ātaḥ anupasarge kaḥ iti . \\
\noindent \textbf{(P\_1,3.9.3)} kaṇviṣaye aṇ api prāpnoti . \\
\noindent \textbf{(P\_1,3.9.3)} sarvādeśe ca doṣaḥ bhavati . \\
\noindent \textbf{(P\_1,3.9.3)} divaḥ aut sarvādeśaḥ prāpnoti . \\
\noindent \textbf{(P\_1,3.9.3)} dāppratiṣedhe pṛthaktvanirdeśaḥ kartavyaḥ . \\
\noindent \textbf{(P\_1,3.9.3)} adābdaipau iti vaktavyam . \\
\noindent \textbf{(P\_1,3.9.3)} kim punaḥ kāraṇam na sidhyati . \\
\noindent \textbf{(P\_1,3.9.3)} anākārāntatvāt . \\
\noindent \textbf{(P\_1,3.9.3)} nanu ca āttve kṛte bhaviṣyati . \\
\noindent \textbf{(P\_1,3.9.3)} tat hi āttvam na prāpnoti . \\
\noindent \textbf{(P\_1,3.9.3)} kim kāraṇam . \\
\noindent \textbf{(P\_1,3.9.3)} anejantatvāt . \\
\noindent \textbf{(P\_1,3.9.3)} astu tarhi anekāntaḥ . \\
\noindent \textbf{(P\_1,3.9.3)} anekānte vṛttiviśeṣaḥ . \\
\noindent \textbf{(P\_1,3.9.3)} yadi anekāntaḥ vṛttiviśeṣaḥ na sidhyati . \\
\noindent \textbf{(P\_1,3.9.3)} kiti ṇiti iti kāryāṇi na sidhyanti . \\
\noindent \textbf{(P\_1,3.9.3)} kim hi saḥ tasya it bhavati yena itkṛtam syāt . \\
\noindent \textbf{(P\_1,3.9.3)} evam tarhi anantaraḥ . \\
\noindent \textbf{(P\_1,3.9.3)} anantaraḥ iti cet pūrvaparayoḥ itkṛtaprasaṅgaḥ . \\
\noindent \textbf{(P\_1,3.9.3)} anantaraḥ iti cet pūrvaparayoḥ itkṛtam prāpnoti . \\
\noindent \textbf{(P\_1,3.9.3)} vuñchaṇ . \\
\noindent \textbf{(P\_1,3.9.3)} siddham tu vyavasitapāṭhāt . \\
\noindent \textbf{(P\_1,3.9.3)} siddham etat . \\
\noindent \textbf{(P\_1,3.9.3)} katham . \\
\noindent \textbf{(P\_1,3.9.3)} vyavasitapāṭhaḥ kartavyaḥ . \\
\noindent \textbf{(P\_1,3.9.3)} vuñ chaṇ . \\
\noindent \textbf{(P\_1,3.9.3)} saḥ ca avaśyam vyavasitapāṭhaḥ kartavyaḥ . \\
\noindent \textbf{(P\_1,3.9.3)} itarathā hi ekānte api sandehaḥ . \\
\noindent \textbf{(P\_1,3.9.3)} akriyamāṇe vyavasitapāṭhe ekānte api sandehaḥ syāt . \\
\noindent \textbf{(P\_1,3.9.3)} tatra na jñāyate kim ayam pūrvasya bhavati āhosvit parasya iti . \\
\noindent \textbf{(P\_1,3.9.3)} sandehamātram etat bhavati . \\
\noindent \textbf{(P\_1,3.9.3)} sarvasandeheṣu ca idam upatiṣṭhate vyākhyānataḥ viśeṣapratipattiḥ na hi sandehāt alakṣaṇam iti . \\
\noindent \textbf{(P\_1,3.9.3)} pūrvasya iti vyākhyāsyamaḥ . \\
\noindent \textbf{(P\_1,3.9.3)} vṛttāt vā . \\
\noindent \textbf{(P\_1,3.9.3)} vṛttāt vā punaḥ siddham etat . \\
\noindent \textbf{(P\_1,3.9.3)} vṛddhimantam ādyudāttam dṛṣṭvā ñit iti vyavaseyam . \\
\noindent \textbf{(P\_1,3.9.3)} antodāttam dṛṣṭvā kit iti . \\
\noindent \textbf{(P\_1,3.9.3)} yuktam punaḥ yat vṛttinimittakaḥ anubandhaḥ syāt na anubandhanimittakena nāma vṛttena bhavitavyam . \\
\noindent \textbf{(P\_1,3.9.3)} vṛttinimittakaḥ eva anubandhaḥ . \\
\noindent \textbf{(P\_1,3.9.3)} vṛttijñaḥ hi ācāryaḥ anubandhān āsajati . \\
\noindent \textbf{(P\_1,3.9.3)} ubhayam idam anubandheṣu uktam ekāntāḥ anekāntāḥ iti . \\
\noindent \textbf{(P\_1,3.9.3)} kim atra nyāyyam . \\
\noindent \textbf{(P\_1,3.9.3)} ekāntāḥ iti nyāyyam . \\
\noindent \textbf{(P\_1,3.9.3)} kutaḥ etat . \\
\noindent \textbf{(P\_1,3.9.3)} atra hi hetuḥ vyapadiṣṭaḥ . \\
\noindent \textbf{(P\_1,3.9.3)} yat ca nāma sahetukam tat nyāyyam . \\
\noindent \textbf{(P\_1,3.9.3)} nanu ca uktam tatra asarūpasarvādeśadāppratiṣedhe pṛthaktvanirdeśaḥ anākārāntatvāt iti . \\
\noindent \textbf{(P\_1,3.9.3)} asarūpavidhau tāvat na doṣaḥ . \\
\noindent \textbf{(P\_1,3.9.3)} ācāryapravṛttiḥ jñāpayati na anubandhakṛtam asārūpyam bhavati iti yat ayam dadātidadhātyoḥ vibhāṣā iti vibhāṣā śam śāsti . \\
\noindent \textbf{(P\_1,3.9.3)} yat api uktam sarvādeśe iti . \\
\noindent \textbf{(P\_1,3.9.3)} atra api ācāryapravṛttiḥ jñāpayati na anubandhakṛtam anekāltvam bhavati iti yat ayam śit sarvasya iti āha . \\
\noindent \textbf{(P\_1,3.9.3)} yat api uktam dāppratiṣedhe pṛthaktvanirdeśaḥ kartavyaḥ iti . \\
\noindent \textbf{(P\_1,3.9.3)} na kartavyaḥ . \\
\noindent \textbf{(P\_1,3.9.3)} ācāryapravṛttiḥ jñāpayati na anubandhakṛtam anejantatvam bhavati iti yat ayam udīcām māṅaḥ vyatīhāre iti meṅaḥ sānubandhakasya āttvabhūtasya grahaṇam karoti . \\
\noindent \textbf{(P\_1,3.10.1)} kim iha udāharaṇam . \\
\noindent \textbf{(P\_1,3.10.1)} ikaḥ yaṇ aci . \\
\noindent \textbf{(P\_1,3.10.1)} dadhi atra madhu atra . \\
\noindent \textbf{(P\_1,3.10.1)} na etat asti . \\
\noindent \textbf{(P\_1,3.10.1)} sthāne antaratamena api etat siddham . \\
\noindent \textbf{(P\_1,3.10.1)} kutaḥ āntaryam . \\
\noindent \textbf{(P\_1,3.10.1)} tālusthānasya tālusthānaḥ oṣṭhasthānasya oṣṭhasthānaḥ bhaviṣyati iti . \\
\noindent \textbf{(P\_1,3.10.1)} idam tarhi . \\
\noindent \textbf{(P\_1,3.10.1)} tasthasthamipām tāmtamtāmaḥ iti . \\
\noindent \textbf{(P\_1,3.10.1)} nanu ca etat api sthāne antaratamena eva siddham . \\
\noindent \textbf{(P\_1,3.10.1)} kutaḥ āntaryam . \\
\noindent \textbf{(P\_1,3.10.1)} ekārthasya ekārthaḥ dvyarthasya dvyarthaḥ bahvarthasya bahvarthaḥ bhaviṣyati iti . \\
\noindent \textbf{(P\_1,3.10.1)} idam tarhi tūdīśalātruavarmatīkūcavārāt ḍhakchaṇḍhañyakaḥ iti . \\
\noindent \textbf{(P\_1,3.10.2)} kimartham punaḥ idam ucyate . \\
\noindent \textbf{(P\_1,3.10.2)} sañjñāsamāsanirdeśāt sarvaprasaṅgaḥ anudeśasya yathāsaṅkhyavacanam niyamārtham . sañjñayā samāsaiḥ ca nirdeśāḥ kriyante . \\
\noindent \textbf{(P\_1,3.10.2)} sañjñayā tāvat . \\
\noindent \textbf{(P\_1,3.10.2)} parasmaipadānām ṇalatususthalathusaṇalvamāḥ iti . \\
\noindent \textbf{(P\_1,3.10.2)} samāsaiḥ . \\
\noindent \textbf{(P\_1,3.10.2)} tūdīśalāturavarmatīkūcavārāt ḍhakchaṇḍhañyakaḥ iti . \\
\noindent \textbf{(P\_1,3.10.2)} sañjñāsamāsanirdeśāt sarvaprasaṅgaḥ anudeśasya yathāsaṅkhyavacanam niyamārtham . \\
\noindent \textbf{(P\_1,3.10.2)} sarvasya uddeśasya sarvaḥ anudeśaḥ prāpnoti . \\
\noindent \textbf{(P\_1,3.10.2)} iṣyate ca samasaṅkhyam yathā syāt iti . \\
\noindent \textbf{(P\_1,3.10.2)} tat ca antareṇa yatnam na sidhyati iti tatra yathāsaṅkhyavacanam niyamārtham . \\
\noindent \textbf{(P\_1,3.10.2)} evamartham idam ucyate . \\
\noindent \textbf{(P\_1,3.10.2)} kim punaḥ kāraṇam sañjñayā samāsaiḥ ca nirdeśāḥ kriyante . \\
\noindent \textbf{(P\_1,3.10.2)} sañjñāsamāsanirdeśaḥ lpṛthak vibhaktisañjñyanuccāraṇārthaḥ . sañjñayā samāsaiḥ ca nirdeśāḥ kriyantepṛthak vibhaktīḥ sañjñinaḥ ca mā uccicīram iti . \\
\noindent \textbf{(P\_1,3.10.2)} prakaraṇe ca sarvasampratyayārthaḥ . \\
\noindent \textbf{(P\_1,3.10.2)} prakaraṇe ca sarveṣām sampratyayaḥ yathā syāt . \\
\noindent \textbf{(P\_1,3.10.2)} vidaḥ laṭaḥ vā iti . \\
\noindent \textbf{(P\_1,3.10.3)} kim punaḥ śabdataḥ sāmye saṅkhyātānudeśaḥ bhavati āhosvit arthataḥ . \\
\noindent \textbf{(P\_1,3.10.3)} kaḥ ca atra viśeṣaḥ . \\
\noindent \textbf{(P\_1,3.10.3)} saṅkhyāsāmyam śabdataḥ cet ṇalādayaḥ parasmaipadānām ḍāraurasaḥ prathamasya ayavāyāvaḥ ecaḥ iti anirdeśaḥ . \\
\noindent \textbf{(P\_1,3.10.3)} agamakaḥ nirdeśaḥ anirdeśaḥ . \\
\noindent \textbf{(P\_1,3.10.3)} parasmaipadānām ṇalatususthalathusaṇalvamāḥ iti ṇalādayaḥ bahavaḥ parasmaipadānām iti ekaḥ śabdaḥ . \\
\noindent \textbf{(P\_1,3.10.3)} vaiṣamyāt saṅkhyātānudeśaḥ na prāpnoti . \\
\noindent \textbf{(P\_1,3.10.3)} ḍāraurasaḥ prathamasya . \\
\noindent \textbf{(P\_1,3.10.3)} ḍāraurasaḥ bahavaḥ prathamasya iti ekaḥ śabdaḥ . \\
\noindent \textbf{(P\_1,3.10.3)} vaiṣamyāt saṅkhyātānudeśaḥ na prāpnoti . \\
\noindent \textbf{(P\_1,3.10.3)} ecaḥ ayavāyāvaḥ . \\
\noindent \textbf{(P\_1,3.10.3)} ayavāyāvaḥ bahavaḥ ecaḥ iti ekaḥ śabdaḥ . \\
\noindent \textbf{(P\_1,3.10.3)} vaiṣamyāt saṅkhyātānudeśaḥ na prāpnoti . \\
\noindent \textbf{(P\_1,3.10.3)} astu tarhi arthataḥ . \\
\noindent \textbf{(P\_1,3.10.3)} arthataḥ cet lṛluṭornandyarīhaṇasindhutakṣaśilādiṣu doṣaḥ . \\
\noindent \textbf{(P\_1,3.10.3)} lṛluṭornandyarīhaṇasindhutakṣaśilādiṣu doṣaḥ bhavati . \\
\noindent \textbf{(P\_1,3.10.3)} syatāsīlṛluṭoḥ . \\
\noindent \textbf{(P\_1,3.10.3)} syatāsī dvau lṛluṭoḥ iti asya trayaḥ arthāḥ . \\
\noindent \textbf{(P\_1,3.10.3)} vaiṣamyāt saṅkhyātānudeśaḥ na prāpnoti . \\
\noindent \textbf{(P\_1,3.10.3)} nandigrahipacādibhyaḥ lyuṇinyacaḥ . \\
\noindent \textbf{(P\_1,3.10.3)} nandyādayaḥ bahavaḥ lyuṇinyacaḥ trayaḥ . \\
\noindent \textbf{(P\_1,3.10.3)} vaiṣamyāt saṅkhyātānudeśaḥ na prāpnoti . \\
\noindent \textbf{(P\_1,3.10.3)} arīhaṇādayaḥ bahavaḥ vuñādayaḥ saptadaśa . \\
\noindent \textbf{(P\_1,3.10.3)} vaiṣamyāt saṅkhyātānudeśaḥ na prāpnoti . \\
\noindent \textbf{(P\_1,3.10.3)} sindhutakṣaśilādibhyaḥ aṇañau . \\
\noindent \textbf{(P\_1,3.10.3)} sindhutakṣaśilādayaḥ bahavaḥ aṇañau dvau . \\
\noindent \textbf{(P\_1,3.10.3)} vaiṣamyāt saṅkhyātānudeśaḥ na prāpnoti . \\
\noindent \textbf{(P\_1,3.10.3)} ātmanepadavidhiniṣṭhāsārvadhātukadvigrahaṇeṣu . \\
\noindent \textbf{(P\_1,3.10.3)} ātmanepadavidhiniṣṭhāsārvadhātukadvigrahaṇeṣu ca doṣaḥ bhavati . \\
\noindent \textbf{(P\_1,3.10.3)} ātmanepadavidhiḥ ca na sidhyati . \\
\noindent \textbf{(P\_1,3.10.3)} anudāttaṅitaḥ ātmanepadam . \\
\noindent \textbf{(P\_1,3.10.3)} anudāttaṅitau dvau ātmanepadam iti asya dvau arthau . \\
\noindent \textbf{(P\_1,3.10.3)} tatra saṅkhyātānudeśaḥ prāpnoti . \\
\noindent \textbf{(P\_1,3.10.3)} niṣṭhā . \\
\noindent \textbf{(P\_1,3.10.3)} radābhyām niṣṭhātaḥ naḥ pūrvasya ca daḥ iti . \\
\noindent \textbf{(P\_1,3.10.3)} rephadakārau dvau niṣṭhā iti asya dvau arthau . \\
\noindent \textbf{(P\_1,3.10.3)} tatra saṅkhyātānudeśaḥ prāpnoti . \\
\noindent \textbf{(P\_1,3.10.3)} sārvadhātukadvigrahaṇeṣu ca doṣaḥ bhavati . \\
\noindent \textbf{(P\_1,3.10.3)} śnasoḥ allopaḥ śnamastī dvau sārvadhātukam iti asya dvau arthau . \\
\noindent \textbf{(P\_1,3.10.3)} tatra saṅkhyātānudeśaḥ prāpnoti . \\
\noindent \textbf{(P\_1,3.10.3)} eṅaḥ pūrvatve pratiṣedhaḥ . \\
\noindent \textbf{(P\_1,3.10.3)} eṅaḥ pūrvatve pratiṣedhaḥ vaktavyaḥ . \\
\noindent \textbf{(P\_1,3.10.3)} eṅaḥ padāntāt ati ṅasiṅasoḥ ca . \\
\noindent \textbf{(P\_1,3.10.3)} ṅasiṅasau dvau eṅ iti asya dvau arthau . \\
\noindent \textbf{(P\_1,3.10.3)} tatra saṅkhyātānudeśaḥ prāpnoti . \\
\noindent \textbf{(P\_1,3.10.3)} astu tarhi śabdataḥ . \\
\noindent \textbf{(P\_1,3.10.3)} nanu ca uktam saṅkhyāsāmyam śabdataḥ cet ṇalādayaḥ parasmaipadānām ḍāraurasaḥ prathamasya ayavāyāvaḥ ecaḥ iti anirdeśaḥ iti . \\
\noindent \textbf{(P\_1,3.10.3)} na eṣaḥ doṣāḥ . \\
\noindent \textbf{(P\_1,3.10.3)} sthāne antaratamaḥ iti anena vyavasthā bhaviṣyati . \\
\noindent \textbf{(P\_1,3.10.3)} kutaḥ āntaryam . \\
\noindent \textbf{(P\_1,3.10.3)} ekārthasya ekārthaḥ dvyarthasya dvyarthaḥ bahvarthasya bahvarthaḥ . \\
\noindent \textbf{(P\_1,3.10.3)} saṃvṛtāvarṇasya saṃvṛtāvarṇaḥ vivṛtāvarṇasya vivṛtāvarṇaḥ . \\
\noindent \textbf{(P\_1,3.10.3)} atiprasaṅgaḥ guṇavṛddhipratiṣedhe kṅiti . \\
\noindent \textbf{(P\_1,3.10.3)} atiprasaṅgaḥ bhavati guṇavṛddhipratiṣedhe kṅiti . \\
\noindent \textbf{(P\_1,3.10.3)} guṇavṛddhī dve kṅitau dvau . \\
\noindent \textbf{(P\_1,3.10.3)} tatra saṅkhyātānudeśaḥ prāpnoti . \\
\noindent \textbf{(P\_1,3.10.3)} na eṣaḥ doṣaḥ . \\
\noindent \textbf{(P\_1,3.10.3)} gakāraḥ api atra nirdiśyate . \\
\noindent \textbf{(P\_1,3.10.3)} tat gakāragrahaṇam api kartavyam . \\
\noindent \textbf{(P\_1,3.10.3)} na kartavyam . \\
\noindent \textbf{(P\_1,3.10.3)} kriyate nyāse eva . \\
\noindent \textbf{(P\_1,3.10.3)} kakāre gakāraḥ cartvabhūtaḥ nirdiśyate . \\
\noindent \textbf{(P\_1,3.10.3)} giti kiti ṅiti iti . \\
\noindent \textbf{(P\_1,3.10.3)} udi kūle rujivahoḥ . \\
\noindent \textbf{(P\_1,3.10.3)} udikūle dve rujivahau dvau . \\
\noindent \textbf{(P\_1,3.10.3)} tatra saṅkhyātānudeśaḥ prāpnoti . \\
\noindent \textbf{(P\_1,3.10.3)} na eṣaḥ doṣaḥ . \\
\noindent \textbf{(P\_1,3.10.3)} na udiḥ upapadam . \\
\noindent \textbf{(P\_1,3.10.3)} kim tarhi . \\
\noindent \textbf{(P\_1,3.10.3)} viśeṣaṇam rujivahoḥ . \\
\noindent \textbf{(P\_1,3.10.3)} utpūrvābhyām rujivahibhyām kūle upapade iti . \\
\noindent \textbf{(P\_1,3.10.3)} tacchīlādiṣu dhātutrigrahaṇeṣu . \\
\noindent \textbf{(P\_1,3.10.3)} tacchīlādiṣu dhātutrigrahaṇeṣu doṣaḥ bhavati . \\
\noindent \textbf{(P\_1,3.10.3)} vidibhidicchideḥ kurac . \\
\noindent \textbf{(P\_1,3.10.3)} vidibhidicchidayaḥ trayaḥ tacchīlādayaḥ trayaḥ . \\
\noindent \textbf{(P\_1,3.10.3)} tatra saṅkhyātānudeśaḥ prāpnoti . \\
\noindent \textbf{(P\_1,3.10.3)} ghañādiṣu dvigrahaṇeṣu . \\
\noindent \textbf{(P\_1,3.10.3)} ghañādiṣu dvigrahaṇeṣu doṣaḥ bhavati . \\
\noindent \textbf{(P\_1,3.10.3)} nirabhyoḥ pūlvoḥ . \\
\noindent \textbf{(P\_1,3.10.3)} nirabhī dvau pūlvau dvau . \\
\noindent \textbf{(P\_1,3.10.3)} tatra saṅkhyātānudeśaḥ prāpnoti . \\
\noindent \textbf{(P\_1,3.10.3)} na eṣaḥ doṣaḥ . \\
\noindent \textbf{(P\_1,3.10.3)} iṣyate ca atra saṅkhyātānudeśaḥ : niṣpāvaḥ , abhilāvaḥ iti . \\
\noindent \textbf{(P\_1,3.10.3)} evam tarhi akartari ca kārake bhāve ca iti dvau pūlvau ca dvau . \\
\noindent \textbf{(P\_1,3.10.3)} tatra saṅkhyātānudeśaḥ prāpnoti . \\
\noindent \textbf{(P\_1,3.10.3)} ave tṝstroḥ karaṇādhikaraṇayoḥ . \\
\noindent \textbf{(P\_1,3.10.3)} tṝstrau dvau karaṇādhikaraṇe dve . \\
\noindent \textbf{(P\_1,3.10.3)} tatra saṅkhyātānudeśaḥ prāpnoti . \\
\noindent \textbf{(P\_1,3.10.3)} kartṛkarmaṇoḥ ca bhūkṛñoḥ . \\
\noindent \textbf{(P\_1,3.10.3)} kartṛkarmaṇī dve bhūkṛñau dvau . \\
\noindent \textbf{(P\_1,3.10.3)} tatra saṅkhyātānudeśaḥ prāpnoti . \\
\noindent \textbf{(P\_1,3.10.3)} anavakḷptyamarṣayoḥ akiṃvṛtte api . \\
\noindent \textbf{(P\_1,3.10.3)} anavakḷptyamarṣau dvau kiṃvṛttākiṃvṛtte dve . \\
\noindent \textbf{(P\_1,3.10.3)} tatra saṅkhyātānudeśaḥ prāpnoti . \\
\noindent \textbf{(P\_1,3.10.3)} kṛbhvoḥ ktvāṇamulau . \\
\noindent \textbf{(P\_1,3.10.3)} kṛbhvau dvau ktvāṇamulau dvau . \\
\noindent \textbf{(P\_1,3.10.3)} tatra saṅkhyātānudeśaḥ prāpnoti . \\
\noindent \textbf{(P\_1,3.10.3)} adhīyānaviduṣoḥ chandobrāhmaṇāni . \\
\noindent \textbf{(P\_1,3.10.3)} chandobrāhmaṇāni iti dve adhīte veda iti ca dvau . \\
\noindent \textbf{(P\_1,3.10.3)} tatra saṅkhyātānudeśaḥ prāpnoti . \\
\noindent \textbf{(P\_1,3.10.3)} ropadhetoḥ pathidūtayoḥ . \\
\noindent \textbf{(P\_1,3.10.3)} ropadhetoḥ prācām tat gacchati pathidūtayoḥ . \\
\noindent \textbf{(P\_1,3.10.3)} ropadhetau dvau pathidūtau dvau . \\
\noindent \textbf{(P\_1,3.10.3)} tatra saṅkhyātānudeśaḥ prāpnoti . \\
\noindent \textbf{(P\_1,3.10.3)} tatra bhavataḥ tasya vyākhyānaḥ kratuyajñebhyaḥ . \\
\noindent \textbf{(P\_1,3.10.3)} tatra bhavatastasyavyākhyānau dvau kratuyajñau dvau . \\
\noindent \textbf{(P\_1,3.10.3)} tatra saṅkhyātānudeśaḥ prāpnoti . \\
\noindent \textbf{(P\_1,3.10.3)} saṅghādiṣu añprabhṛtayaḥ ṣaṅghādiṣu añprabhṛtayaḥ saṅkhyātānudeśena na sidhyanti . \\
\noindent \textbf{(P\_1,3.10.3)} na eṣaḥ doṣaḥ . \\
\noindent \textbf{(P\_1,3.10.3)} ghoṣagrahaṇam atra kartavyam . \\
\noindent \textbf{(P\_1,3.10.3)} veśoyaśāadeḥ bhagāt yalkhau . \\
\noindent \textbf{(P\_1,3.10.3)} veśoyaśāadī dvau yalkhau dvau . \\
\noindent \textbf{(P\_1,3.10.3)} tatra saṅkhyātānudeśaḥ prāpnoti . \\
\noindent \textbf{(P\_1,3.10.3)} ṅasiṅasoḥ khyatyāt parasya . \\
\noindent \textbf{(P\_1,3.10.3)} ṅasiṅasau dvau khyatyau dvau . \\
\noindent \textbf{(P\_1,3.10.3)} tatra saṅkhyātānudeśaḥ prāpnoti . \\
\noindent \textbf{(P\_1,3.10.3)} na vā samānayogavacanāt . \\
\noindent \textbf{(P\_1,3.10.3)} na vā eṣaḥ doṣaḥ . \\
\noindent \textbf{(P\_1,3.10.3)} kim kāraṇam . \\
\noindent \textbf{(P\_1,3.10.3)} samānayogavacanāt . \\
\noindent \textbf{(P\_1,3.10.3)} samānayoge saṅkhyātānudeśam vakṣyāmi . \\
\noindent \textbf{(P\_1,3.10.3)} tasya doṣaḥ vidaḥ laṭaḥ vā . \\
\noindent \textbf{(P\_1,3.10.3)} tasya etasya lakṣaṇasya doṣaḥ vidaḥ laṭaḥ vā iti saṅkhyātānudeśaḥ na prāpnoti . \\
\noindent \textbf{(P\_1,3.10.3)} dhmādhetṭoḥ nāḍīmuṣṭyoḥ ca . \\
\noindent \textbf{(P\_1,3.10.3)} dhmādhetṭoḥ nāḍīmuṣṭyoḥ ca saṅkhyātānudeśaḥ na prāpnoti . \\
\noindent \textbf{(P\_1,3.10.3)} khalagorathāt initrakaṭyacaḥ ca . \\
\noindent \textbf{(P\_1,3.10.3)} saṅkhyātānudeśaḥ na prāpnoti . \\
\noindent \textbf{(P\_1,3.10.3)} sindhvapakarābhyām kan aṇañau ca . \\
\noindent \textbf{(P\_1,3.10.3)} saṅkhyātānudeśaḥ na prāpnoti . \\
\noindent \textbf{(P\_1,3.10.3)} yuṣmadasmadoḥ ca ādeśāḥ . \\
\noindent \textbf{(P\_1,3.10.3)} yuṣmadasmadoḥ ca ādeśāḥ saṅkhyātānudeśena na sidhyanti . \\
\noindent \textbf{(P\_1,3.10.3)} tasmāt yasmin pakṣe alpīyāṃsaḥ doṣāḥ tām āsthāya pratividheyam doṣeṣu . \\
\noindent \textbf{(P\_1,3.10.3)} atha vā evam vakṣyāmi . \\
\noindent \textbf{(P\_1,3.10.3)} yathāsaṅkhyam anudeśaḥ samānām svaritena . \\
\noindent \textbf{(P\_1,3.10.3)} tataḥ adhikāraḥ . \\
\noindent \textbf{(P\_1,3.10.3)} adhikāraḥ ca bhavati svaritena iti . \\
\noindent \textbf{(P\_1,3.10.3)} evam api svaritam dṛṣṭvā sandehaḥ syāt . \\
\noindent \textbf{(P\_1,3.10.3)} na jñāyate kim ayam samasaṅkhyārthaḥ āhosvit adhikārārthaḥ iti . \\
\noindent \textbf{(P\_1,3.10.3)} sandehamātram etat bhavati . \\
\noindent \textbf{(P\_1,3.10.3)} sarvasandeheṣu ca idam upatiṣṭhate vyākhyānataḥ viśeṣapratipattiḥ na hi sandehāt alakṣaṇam iti samasaṅkhyārthaḥ iti vyākhyāsyāmaḥ . \\
\noindent \textbf{(P\_1,3.11.1)} kimartham idam ucyate . \\
\noindent \textbf{(P\_1,3.11.1)} adhikāraḥ pratiyogam tasya anirdeśārthaḥ . \\
\noindent \textbf{(P\_1,3.11.1)} adhikāraḥ kriyate pratiyogam tasya anirdeśārthaḥ iti . \\
\noindent \textbf{(P\_1,3.11.1)} kim idam pratiyogam iti . \\
\noindent \textbf{(P\_1,3.11.1)} yogam yogam prati pratiyogam . \\
\noindent \textbf{(P\_1,3.11.1)} yoge yoge tasya grahaṇam mā kārṣam iti . \\
\noindent \textbf{(P\_1,3.11.1)} kim gatam etat iyatā sūtreṇa . \\
\noindent \textbf{(P\_1,3.11.1)} gatam iti āha . \\
\noindent \textbf{(P\_1,3.11.1)} kutaḥ . \\
\noindent \textbf{(P\_1,3.11.1)} lokataḥ . \\
\noindent \textbf{(P\_1,3.11.1)} tat yathā loke adhikṛtaḥ asau grāme adhikṛtaḥ asau nagare iti ucyate yaḥ yatra vyāpāram gacchati . \\
\noindent \textbf{(P\_1,3.11.1)} śabdena ca api adhikṛtena kaḥ anyaḥ vyāpāraḥ śakyaḥ avagantum anyat ataḥ yoge yoge upasthānāt . \\
\noindent \textbf{(P\_1,3.11.1)} na vā nirdiśyamānādhikṛtatvāt yathā loke . \\
\noindent \textbf{(P\_1,3.11.1)} na vā etat prayojanam . \\
\noindent \textbf{(P\_1,3.11.1)} kim kāraṇam . \\
\noindent \textbf{(P\_1,3.11.1)} nirdiśyamānādhikṛtatvāt yathā loke . \\
\noindent \textbf{(P\_1,3.11.1)} nirdiśyamānam adhikṛtam gamyate . \\
\noindent \textbf{(P\_1,3.11.1)} tat yathā . \\
\noindent \textbf{(P\_1,3.11.1)} devadattāya gauḥ dīyatām yajñadattāya viṣṇumitrāya iti . \\
\noindent \textbf{(P\_1,3.11.1)} gauḥ iti gamyate . \\
\noindent \textbf{(P\_1,3.11.1)} evam iha api padarujaviśaspṛśaḥ ghañ sṛ sthire bhāve . \\
\noindent \textbf{(P\_1,3.11.1)} ghañ iti gamyate . \\
\noindent \textbf{(P\_1,3.11.1)} anyanirdeśaḥ tu nivartakaḥ tasmāt paribhāṣā . \\
\noindent \textbf{(P\_1,3.11.1)} anyanirdeśaḥ tu loke nivartakaḥ bhavati . \\
\noindent \textbf{(P\_1,3.11.1)} tat yathā . \\
\noindent \textbf{(P\_1,3.11.1)} devadattāya gauḥ dīyatām yajñadattāya kambalaḥ viṣṇumitrāya ca iti . \\
\noindent \textbf{(P\_1,3.11.1)} kambalaḥ gonivartakaḥ bhavati . \\
\noindent \textbf{(P\_1,3.11.1)} evam iha api abhividhau bhāve inuṇ ghañaḥ nivartakaḥ syāt . \\
\noindent \textbf{(P\_1,3.11.1)} tasmāt paribhāṣā kartavyā . \\
\noindent \textbf{(P\_1,3.11.2)} adhikāraparimāṇājñānam tu . \\
\noindent \textbf{(P\_1,3.11.2)} adhikāraparimāṇājñānam tu bhavati . \\
\noindent \textbf{(P\_1,3.11.2)} na jñāyate kiyantam avadhim adhikāraḥ anuvartate iti . \\
\noindent \textbf{(P\_1,3.11.2)} adhikāraparimāṇajñānārtham tu . \\
\noindent \textbf{(P\_1,3.11.2)} adhikāraparimāṇajñānārtham eva tarhi ayam yogaḥ vaktavyaḥ . \\
\noindent \textbf{(P\_1,3.11.2)} adhikāraparimāṇam jñāsyāmi iti . \\
\noindent \textbf{(P\_1,3.11.2)} katham punaḥ svaritena adhikāraḥ iti anena adhikāraparimāṇam śakyam vijñātum . \\
\noindent \textbf{(P\_1,3.11.2)} evam vakṣyāmi svarite na adhikāraḥ iti . \\
\noindent \textbf{(P\_1,3.11.2)} svaritam dṛṣṭvā adhikāraḥ na bhavati iti . \\
\noindent \textbf{(P\_1,3.11.2)} kena idānīm adhikāraḥ bhaviṣyati . \\
\noindent \textbf{(P\_1,3.11.2)} laukikaḥ adhikāraḥ . \\
\noindent \textbf{(P\_1,3.11.2)} na adhikāraḥ iti cet uktam . \\
\noindent \textbf{(P\_1,3.11.2)} kim uktam . \\
\noindent \textbf{(P\_1,3.11.2)} anyanirdeśaḥ tu nivartakaḥ tasmāt paribhāṣā iti . \\
\noindent \textbf{(P\_1,3.11.2)} adhikārārtham eva tarhi ayam yogaḥ vaktavyaḥ . \\
\noindent \textbf{(P\_1,3.11.2)} nanu ca uktam adhikāraparimāṇājñānam tu iti . \\
\noindent \textbf{(P\_1,3.11.2)} yāvatithaḥ al anubandhaḥ tāvataḥ yogān iti vacanāt siddham . \\
\noindent \textbf{(P\_1,3.11.2)} yāvatithaḥ al anubadhyate tāvataḥ yogān adhikāraḥ anuvartate iti vaktavyam . \\
\noindent \textbf{(P\_1,3.11.2)} atha idānīm yatra alpīyāṃsaḥ alaḥ bhūyasaḥ ca yogān adhikāraḥ anuvartate katham tatra kartavyam . \\
\noindent \textbf{(P\_1,3.11.2)} bhūyasi prāgvacanam . \\
\noindent \textbf{(P\_1,3.11.2)} bhūyasi prāgvacanam kartavyam . \\
\noindent \textbf{(P\_1,3.11.2)} prāk amutaḥ iti vaktavyam . \\
\noindent \textbf{(P\_1,3.11.2)} tat tarhi vaktavyam . \\
\noindent \textbf{(P\_1,3.11.2)} na vaktavyam . \\
\noindent \textbf{(P\_1,3.11.2)} sandehamātram etat bhavati . \\
\noindent \textbf{(P\_1,3.11.2)} sarvasandeheṣu ca idam upatiṣṭhate vyākhyānataḥ viśeṣapratipattiḥ na hi sandehāt alakṣaṇam iti . \\
\noindent \textbf{(P\_1,3.11.2)} prāk amutaḥ iti vyākhyāsyāmaḥ . \\
\noindent \textbf{(P\_1,3.11.2)} yadi evam na arthaḥ anena . \\
\noindent \textbf{(P\_1,3.11.2)} kena idānīm adhikāraḥ bhaviṣyati . \\
\noindent \textbf{(P\_1,3.11.2)} laukikaḥ adhikāraḥ . \\
\noindent \textbf{(P\_1,3.11.2)} nanu ca uktam na adhikāraḥ iti cet uktam . \\
\noindent \textbf{(P\_1,3.11.2)} kim uktam . \\
\noindent \textbf{(P\_1,3.11.2)} anyanirdeśaḥ tu nivartakaḥ tasmāt paribhāṣā . \\
\noindent \textbf{(P\_1,3.11.2)} sandehamātram etat bhavati . \\
\noindent \textbf{(P\_1,3.11.2)} sarvasandeheṣu ca idam upatiṣṭhate vyākhyānataḥ viśeṣapratipattiḥ na hi sandehāt alakṣaṇam iti . \\
\noindent \textbf{(P\_1,3.11.2)} inuṇ ghañ iti sandehe ghañ iti vyākhyāsyāmaḥ . \\
\noindent \textbf{(P\_1,3.11.3)} na tarhi idānīm ayam yogaḥ vaktavyaḥ . \\
\noindent \textbf{(P\_1,3.11.3)} vaktavyaḥ ca . \\
\noindent \textbf{(P\_1,3.11.3)} kim prayojanam . \\
\noindent \textbf{(P\_1,3.11.3)} svaritena adhikāragatiḥ yathā vijñāyeta , adhikam kāryam , adhikaḥ kāraḥ . \\
\noindent \textbf{(P\_1,3.11.3)} adhikāragatiḥ : gostriyoḥ upasarjanam iti atra goṭāṅgrahaṇam coditam . \\
\noindent \textbf{(P\_1,3.11.3)} tat na kartavyam bhavati . \\
\noindent \textbf{(P\_1,3.11.3)} strīgrahaṇam svarayiṣyate . \\
\noindent \textbf{(P\_1,3.11.3)} svaritena adhikāragatiḥ bhavati iti striyām iti evam prakṛtya ye pratyayāḥ vihitāḥ teṣām grahaṇam vijñāsyate . \\
\noindent \textbf{(P\_1,3.11.3)} tatra svaritena adhikāragatiḥ bhavati iti na doṣaḥ bhavati . \\
\noindent \textbf{(P\_1,3.11.3)} adhikam kāryam : apādānam ācāryaḥ kim nyāyyam manyate . \\
\noindent \textbf{(P\_1,3.11.3)} yatra prāpya nivṛttiḥ . \\
\noindent \textbf{(P\_1,3.11.3)} tena iha eva syāt : grāmāt āgacchati . \\
\noindent \textbf{(P\_1,3.11.3)} nagarāt āgacchati . \\
\noindent \textbf{(P\_1,3.11.3)} sāṅkāśyakebhyaḥ pāṭaliputrakāḥ abhirūpatarāḥ iti atra na syāt . \\
\noindent \textbf{(P\_1,3.11.3)} svaritena adhikarm kāryam bhavati iti atra api siddham bhavati . \\
\noindent \textbf{(P\_1,3.11.3)} tathā adhikaraṇam ācāryaḥ kim nyāyyam manyate . \\
\noindent \textbf{(P\_1,3.11.3)} yatra kṛtsnaḥ ādhārātmā vyāptaḥ bhavati . \\
\noindent \textbf{(P\_1,3.11.3)} tena iha eva syāt : tileṣu tailam . \\
\noindent \textbf{(P\_1,3.11.3)} dadhni sarpiḥ iti . \\
\noindent \textbf{(P\_1,3.11.3)} gaṅgāyām gāvaḥ . \\
\noindent \textbf{(P\_1,3.11.3)} kūpe gargakulam iti atra na syāt . \\
\noindent \textbf{(P\_1,3.11.3)} svaritena adhikarm kāryam bhavati iti atra api siddham bhavati . \\
\noindent \textbf{(P\_1,3.11.3)} adhikam kāryam . \\
\noindent \textbf{(P\_1,3.11.3)} adhikaḥ kāraḥ : pūrvavipratiṣedhāḥ na paṭhitavyāḥ bhavanti . \\
\noindent \textbf{(P\_1,3.11.3)} guṇavṛddhyauttvatṛjvadbhāvebhyaḥ num pūrvavipratiṣiddham numaciratṛjvadbhāvebhyaḥ nuṭ iti . \\
\noindent \textbf{(P\_1,3.11.3)} numnuṭau svarayiṣyete . \\
\noindent \textbf{(P\_1,3.11.3)} tatra svaritena adhikaḥ kāraḥ bhavati iti numnuṭau bhaviṣyataḥ . \\
\noindent \textbf{(P\_1,3.11.3)} katham punaḥ adhikaḥ kāraḥ iti anena pūrvavipratiṣedhāḥ śakya na paṭhitum . \\
\noindent \textbf{(P\_1,3.11.3)} lokataḥ . \\
\noindent \textbf{(P\_1,3.11.3)} tat yathā loke adhikam ayam kāram karoti iti ucyate yaḥ ayam durbalaḥ san balavadbhiḥ saha bhāram vahati . \\
\noindent \textbf{(P\_1,3.11.3)} evam iha api adhikam ayam kāram karoti iti ucyate yaḥ ayam pūrvaḥ san param bādhate . \\
\noindent \textbf{(P\_1,3.11.3)} adhikāragatiḥ stryarthā viśeṣāya adhikam kāryam . \\
\noindent \textbf{(P\_1,3.11.3)} atha yaḥ anyaḥ adhikaḥ kāraḥ pūrvavipratiṣedhārthaḥ saḥ . \\
\noindent \textbf{(P\_1,3.12.1)} vikaraṇebhyaḥ pratiṣedhaḥ vaktavyaḥ . \\
\noindent \textbf{(P\_1,3.12.1)} cinutaḥ sunutaḥ lunītaḥ punītaḥ . \\
\noindent \textbf{(P\_1,3.12.1)} ṅitaḥ iti ātmanepadam prāpnoti . \\
\noindent \textbf{(P\_1,3.12.1)} na eṣaḥ doṣaḥ . \\
\noindent \textbf{(P\_1,3.12.1)} na evam vijñāyate ṅakāraḥ it asya saḥ ayam ṅit ṅitaḥ iti . \\
\noindent \textbf{(P\_1,3.12.1)} katham tarhi . \\
\noindent \textbf{(P\_1,3.12.1)} ṅakāraḥ eva it ṅit ṅitaḥ . \\
\noindent \textbf{(P\_1,3.12.1)} atha vā upadeśe iti vartate . \\
\noindent \textbf{(P\_1,3.12.1)} atha vā uktam etat siddham tu pūrvasya kāryātideśāt iti . \\
\noindent \textbf{(P\_1,3.12.1)} sarvathā caṅaṅbhyām prāpnoti . \\
\noindent \textbf{(P\_1,3.12.1)} evam tarhi dhātoḥ iti vartate . \\
\noindent \textbf{(P\_1,3.12.1)} kva prakṛtam . \\
\noindent \textbf{(P\_1,3.12.1)} bhūvādayaḥ dhātavaḥ iti . \\
\noindent \textbf{(P\_1,3.12.1)} tat vai prathamānirdiṣṭam pañcamīnirdiṣṭena ca iha arthaḥ . \\
\noindent \textbf{(P\_1,3.12.1)} arthāt vibhaktivipariṇāmaḥ bhaviṣyati . \\
\noindent \textbf{(P\_1,3.12.1)} tat yathā . \\
\noindent \textbf{(P\_1,3.12.1)} uccāni devadattasya gṛhāṇi . \\
\noindent \textbf{(P\_1,3.12.1)} āmantrayasva enam . \\
\noindent \textbf{(P\_1,3.12.1)} devadattam iti gamyate . \\
\noindent \textbf{(P\_1,3.12.1)} devadattasya gāvaḥ aśvāḥ hiraṇyam iti . \\
\noindent \textbf{(P\_1,3.12.1)} āḍhyaḥ vaidhaveyaḥ . \\
\noindent \textbf{(P\_1,3.12.1)} devadattaḥ iti gamyate . \\
\noindent \textbf{(P\_1,3.12.1)} purastāt ṣaṣṭhīnirdiṣṭam sat arthāt dvitīyānirdiṣṭam prathamānirdiṣṭam ca bhavati . \\
\noindent \textbf{(P\_1,3.12.1)} evam iha api purastāt prathamānirdiṣṭam sat arthāt pañcamīnirdiṣṭam bhaviṣyati . \\
\noindent \textbf{(P\_1,3.12.2)} kimartham punaḥ idam ucyate . \\
\noindent \textbf{(P\_1,3.12.2)} ātmanepadavacanam niyamārtham . \\
\noindent \textbf{(P\_1,3.12.2)} niyamārthaḥ ayam ārambhaḥ . \\
\noindent \textbf{(P\_1,3.12.2)} kim ucyate niyamārthaḥ ayam iti na punaḥ vidhyarthaḥ api syāt . \\
\noindent \textbf{(P\_1,3.12.2)} lavidhānāt vihitam . \\
\noindent \textbf{(P\_1,3.12.2)} lavidhānāt hi ātmanepadam parasmaipadam ca vihitam . \\
\noindent \textbf{(P\_1,3.12.2)} asti prayojanam etat . \\
\noindent \textbf{(P\_1,3.12.2)} kim tarhi iti . \\
\noindent \textbf{(P\_1,3.12.2)} vikaraṇaiḥ tu vyavahitatvāt niyamaḥ na prāpnoti . \\
\noindent \textbf{(P\_1,3.12.2)} idam iha sampradhāryam . \\
\noindent \textbf{(P\_1,3.12.2)} vikaraṇāḥ kriyantām niyamaḥ iti . \\
\noindent \textbf{(P\_1,3.12.2)} kim atra kartavyam . \\
\noindent \textbf{(P\_1,3.12.2)} paratvāt vikaraṇāḥ . \\
\noindent \textbf{(P\_1,3.12.2)} nityāḥ khalu api vikaraṇāḥ . \\
\noindent \textbf{(P\_1,3.12.2)} kṛte api niyame prāpnuvanti akṛte api prāpnuvanti . \\
\noindent \textbf{(P\_1,3.12.2)} nityatvāt paratvāt ca vikaraṇeṣu kṛteṣu vikaraṇaiḥ vyavahitatvāt niyamaḥ na prāpnoti . \\
\noindent \textbf{(P\_1,3.12.2)} na eṣaḥ doṣaḥ . \\
\noindent \textbf{(P\_1,3.12.2)} anavakāśaḥ niyamaḥ . \\
\noindent \textbf{(P\_1,3.12.2)} sāvakāśaḥ . \\
\noindent \textbf{(P\_1,3.12.2)} kaḥ avakāśaḥ . \\
\noindent \textbf{(P\_1,3.12.2)} ye ete lugvikaraṇāḥ śluvikaraṇāḥ liṅliṭau ca . \\
\noindent \textbf{(P\_1,3.12.2)} yadi punaḥ iyam paribhāṣā vijñāyeta . \\
\noindent \textbf{(P\_1,3.12.2)} kim kṛtam bhavati . \\
\noindent \textbf{(P\_1,3.12.2)} kāryakālam sañjñāparibhāṣam yatra kāryam tatra draṣṭavyam . \\
\noindent \textbf{(P\_1,3.12.2)} lasya tibādayaḥ bhavanti iti upasthitam idam bhavati anudāttaṅitaḥ ātmanepadam śeṣāt kartari parasmaipadam iti . \\
\noindent \textbf{(P\_1,3.12.2)} evam api itaretarāśrayam bhavati . \\
\noindent \textbf{(P\_1,3.12.2)} kā itaretarāśrayatā . \\
\noindent \textbf{(P\_1,3.12.2)} abhinirvṛttānām lasya sthāne tibādīnām ātmanepadaparasmaipadasañjñayā bhavitavyam sañjñayā ca tibādayaḥ bhāvyante . \\
\noindent \textbf{(P\_1,3.12.2)} tat itaretarāśrayam bhavati . \\
\noindent \textbf{(P\_1,3.12.2)} itaretarāśrayāṇi kāryāṇi ca na prakalpante . \\
\noindent \textbf{(P\_1,3.12.2)} parasmaipadeṣu tāvat na itaretarāśrayam bhavati . \\
\noindent \textbf{(P\_1,3.12.2)} parasmaipadānukramaṇam na kariṣyate . \\
\noindent \textbf{(P\_1,3.12.2)} avaśyam kartavyam anuparābhyām kṛñaḥ iti evamartham . \\
\noindent \textbf{(P\_1,3.12.2)} nanu ca etat api ātmanepadānukramaṇe eva kariṣye . \\
\noindent \textbf{(P\_1,3.12.2)} svaritañitaḥ kartrabhipraye kriyāphale ātmanepadam bhavati kartari . \\
\noindent \textbf{(P\_1,3.12.2)} anuparābhyām kṛñaḥ na iti . \\
\noindent \textbf{(P\_1,3.12.2)} ātmanepadeṣu ca api na itaretarāśrayam bhavati . \\
\noindent \textbf{(P\_1,3.12.2)} katham . \\
\noindent \textbf{(P\_1,3.12.2)} bhāvinī sañjñā vijñāsyate sūtraśāṭakavat . \\
\noindent \textbf{(P\_1,3.12.2)} tat yathā : kaḥ cit kam cit tantuvāyam āha : asya sūtrasya śāṭakam vaya iti . \\
\noindent \textbf{(P\_1,3.12.2)} saḥ paśyati . \\
\noindent \textbf{(P\_1,3.12.2)} yadi śāṭakaḥ na vātavyaḥ atha vātavyaḥ na śāṭakaḥ . \\
\noindent \textbf{(P\_1,3.12.2)} śāṭakaḥ vātavyaḥ iti vipratiṣiddham. bhāvinī khalu asya sañjñā abhipretā . \\
\noindent \textbf{(P\_1,3.12.2)} saḥ manye vātavyaḥ yasmin ute śāṭakaḥ iti etat bhavati iti . \\
\noindent \textbf{(P\_1,3.12.2)} evam iha api saḥ lasya sthāne kartavyaḥ lyasya abhinirvṛttasya ātmanepadam iti eṣā sañjñā bhaviṣyati . \\
\noindent \textbf{(P\_1,3.12.2)} atha vā punaḥ astu niyamaḥ . \\
\noindent \textbf{(P\_1,3.12.2)} nanu ca utkam vikaraṇaiḥ tu vyavahitatvāt niyamaḥ na prāpnoti . \\
\noindent \textbf{(P\_1,3.12.2)} na eṣaḥ doṣaḥ . \\
\noindent \textbf{(P\_1,3.12.2)} ācāryapravṛttiḥ jñāpayati vikaraṇebhyaḥ niyamaḥ balīyān iti yat ayam vikaraṇavidhau ātmanepadaparasmaipadāni āśrayati . \\
\noindent \textbf{(P\_1,3.12.2)} puṣādidyutādḷrditaḥ parasmaipadeṣu ātmanepadeṣu anyatarasyām iti . \\
\noindent \textbf{(P\_1,3.12.2)} na etat asti jñāpakam . \\
\noindent \textbf{(P\_1,3.12.2)} abhinirvṛttāni hi lasya sthāne ātmanepadāni parasmaipadāni ca . \\
\noindent \textbf{(P\_1,3.12.2)} yat tarhi anupasargāt vā iti vibhāṣām śāsti . \\
\noindent \textbf{(P\_1,3.12.3)} kim punaḥ ayam pratyayaniyamaḥ : anudāttaṅitaḥ eva ātmanepadam bhavati , bhāvakarmaṇoḥ eva ātmanepadam bhavati iti . \\
\noindent \textbf{(P\_1,3.12.3)} āhosvit prakṛtyarthaniyamaḥ : anudāttaṅitaḥ ātmanepadam eva , bhāvakarmaṇoḥ ātmanepadam eva . \\
\noindent \textbf{(P\_1,3.12.3)} kaḥ ca atra viśeṣaḥ . \\
\noindent \textbf{(P\_1,3.12.3)} tatra pratyayaniyame śeṣavacanam parasmaipadasya anivṛttatvāt . \\
\noindent \textbf{(P\_1,3.12.3)} tatra pratyayaniyame śeṣagrahaṇam kartavyam parasmaipadaniyamārtham . \\
\noindent \textbf{(P\_1,3.12.3)} śeṣāt kartari parasmaipadam iti . \\
\noindent \textbf{(P\_1,3.12.3)} kim kāraṇam . \\
\noindent \textbf{(P\_1,3.12.3)} parasmaipadasya anivṛttatvāt . \\
\noindent \textbf{(P\_1,3.12.3)} pratyayāḥ niyatāḥ prakṛtyarthau aniyatau . \\
\noindent \textbf{(P\_1,3.12.3)} tatra parasmaipadam prāpnoti . \\
\noindent \textbf{(P\_1,3.12.3)} tatra śeṣagrahaṇam kartavyam parasmaipadaniyamārtham . \\
\noindent \textbf{(P\_1,3.12.3)} śeṣāt eva parasmaipadam bhavati na anyataḥ iti . \\
\noindent \textbf{(P\_1,3.12.3)} kyaṣaḥ ātmanepadavacanam tasya anyatra niyamāt . \\
\noindent \textbf{(P\_1,3.12.3)} kyaṣaḥ ātmanepadam vaktavyam . \\
\noindent \textbf{(P\_1,3.12.3)} lohitāyati lohitāyate . \\
\noindent \textbf{(P\_1,3.12.3)} kim punaḥ kāraṇam na sidhyati . \\
\noindent \textbf{(P\_1,3.12.3)} tasya anyatra niyamāt . \\
\noindent \textbf{(P\_1,3.12.3)} tat hi anyatra niyamyate . \\
\noindent \textbf{(P\_1,3.12.3)} ucyate ca na ca prāpnoti . \\
\noindent \textbf{(P\_1,3.12.3)} tat vacanāt bhaviṣyati . \\
\noindent \textbf{(P\_1,3.12.3)} astu tarhi prakṛtyarthaniyamaḥ . \\
\noindent \textbf{(P\_1,3.12.3)} prakṛtyarthaniyame anyābhāvaḥ . \\
\noindent \textbf{(P\_1,3.12.3)} prakṛtyarthaniyame anyeṣām pratyayānām abhāvaḥ . \\
\noindent \textbf{(P\_1,3.12.3)} anudāttaṅitaḥ tṛjādayaḥ na prāpnuvanti . \\
\noindent \textbf{(P\_1,3.12.3)} na eṣaḥ doṣaḥ . \\
\noindent \textbf{(P\_1,3.12.3)} anavakāśāḥ tṛjādayaḥ ucyante ca . \\
\noindent \textbf{(P\_1,3.12.3)} te vacanāt bhaviṣyanti . \\
\noindent \textbf{(P\_1,3.12.3)} sāvakāśāḥ tṛjādayaḥ . \\
\noindent \textbf{(P\_1,3.12.3)} kaḥ avakāśaḥ . \\
\noindent \textbf{(P\_1,3.12.3)} parasmaipadinaḥ avakāśaḥ . \\
\noindent \textbf{(P\_1,3.12.3)} tatra api niyamāt na prāpnuvanti . \\
\noindent \textbf{(P\_1,3.12.3)} tavyādayaḥ tarhi bhāvakarmaṇoḥ niyamāt na prāpnuvanti . \\
\noindent \textbf{(P\_1,3.12.3)} tavyādayaḥ api anavakāśāḥ . \\
\noindent \textbf{(P\_1,3.12.3)} te vacanāt bhaviṣyanti . \\
\noindent \textbf{(P\_1,3.12.3)} ciṇ tarhi bhāvakarmaṇoḥ niyamāt na prāpnoti . \\
\noindent \textbf{(P\_1,3.12.3)} ciṇ api vacanāt bhaviṣyati . \\
\noindent \textbf{(P\_1,3.12.3)} ghañ tarhi bhāvakarmaṇoḥ niyamāt na prāpnoti . \\
\noindent \textbf{(P\_1,3.12.3)} tatra api prakṛtam karmagrahaṇam . \\
\noindent \textbf{(P\_1,3.12.3)} kva prakṛtam . \\
\noindent \textbf{(P\_1,3.12.3)} aṇ karmaṇi ca iti . \\
\noindent \textbf{(P\_1,3.12.3)} tat vai tatra upapadaviśeṣaṇam abhidheyaviśeṣaṇena ca iha arthaḥ . \\
\noindent \textbf{(P\_1,3.12.3)} na ca anyārtham prakṛtam anyārtham bhavati . \\
\noindent \textbf{(P\_1,3.12.3)} na khalu api anyat prakṛtam anuvartanāt anyat bhavati . \\
\noindent \textbf{(P\_1,3.12.3)} na hi godhā sarpantī sarpaṇāt ahiḥ bhavati . \\
\noindent \textbf{(P\_1,3.12.3)} yat tāvat ucyate na ca anyārtham prakṛtam anyārtham bhavati iti anyārtham api prakṛtam anyārtham bhavati ṭat yathā . \\
\noindent \textbf{(P\_1,3.12.3)} śālyartham kulyāḥ praṇīyante tābhyaḥ ca pāṇīyam pīyate upaśpṛśyate ca śālayaḥ ca bhāvyante . \\
\noindent \textbf{(P\_1,3.12.3)} yad api ucyate na khalu api anyat prakṛtam anuvartanāt anyat bhavati . \\
\noindent \textbf{(P\_1,3.12.3)} na hi godhā sarpantī sarpaṇāt ahiḥ bhavati iti . \\
\noindent \textbf{(P\_1,3.12.3)} bhavet dravyeṣu etat evam syāt . \\
\noindent \textbf{(P\_1,3.12.3)} śabdaḥ tu khalu yena yena viśeṣeṇa abhisambadhyate tasya tasya viśeṣakaḥ bhavati . \\
\noindent \textbf{(P\_1,3.12.3)} śeṣavacanam ca . \\
\noindent \textbf{(P\_1,3.12.3)} śeṣagrahaṇam ca kartavyam . \\
\noindent \textbf{(P\_1,3.12.3)} śeṣāt kartari parasmaipadam iti . \\
\noindent \textbf{(P\_1,3.12.3)} kim prayojanam . \\
\noindent \textbf{(P\_1,3.12.3)} śeṣaniyamārtham . \\
\noindent \textbf{(P\_1,3.12.3)} prakṛtarthau niyatau . \\
\noindent \textbf{(P\_1,3.12.3)} pratyayāḥ aniyatāḥ . \\
\noindent \textbf{(P\_1,3.12.3)} te śeṣe api prāpnuvanti . \\
\noindent \textbf{(P\_1,3.12.3)} tatra śeṣagrahaṇam kartavyam . \\
\noindent \textbf{(P\_1,3.12.3)} śeṣāt kartari parasmaipadam eva na anyat iti . \\
\noindent \textbf{(P\_1,3.12.3)} kartari ca ātmanepadaviṣaye parasmaipadapratiṣedhārtham . \\
\noindent \textbf{(P\_1,3.12.3)} kartari ca ātmanepadaviṣaye parasmaipadapratiṣedhārtham dvitīyam śeṣagrahaṇam kartavyam . \\
\noindent \textbf{(P\_1,3.12.3)} śeṣāt śeṣe iti vaktavyam . \\
\noindent \textbf{(P\_1,3.12.3)} iha mā bhūt . \\
\noindent \textbf{(P\_1,3.12.3)} bhidyate kuśūlaḥ svayam eva iti . \\
\noindent \textbf{(P\_1,3.12.3)} katarasmin pakṣe ayam doṣaḥ . \\
\noindent \textbf{(P\_1,3.12.3)} prakṛtyarthaniyame . \\
\noindent \textbf{(P\_1,3.12.3)} prakṛtyarthaniyame tāvat na doṣaḥ . \\
\noindent \textbf{(P\_1,3.12.3)} prakṛtyarthau niyatau . \\
\noindent \textbf{(P\_1,3.12.3)} pratyayāḥ aniyatāḥ . \\
\noindent \textbf{(P\_1,3.12.3)} tatra na arthaḥ kartṛgrahaṇena kartṛgrahaṇāt ca doṣaḥ . \\
\noindent \textbf{(P\_1,3.12.3)} pratyayaniyame tarhi ayam doṣaḥ . \\
\noindent \textbf{(P\_1,3.12.3)} pratyayāḥ niyatāḥ . \\
\noindent \textbf{(P\_1,3.12.3)} prakṛtyarthau aniyatau . \\
\noindent \textbf{(P\_1,3.12.3)} tatra kartṛgrahaṇam kartavyam bhāvakarmaṇoḥ nivṛttyartham . \\
\noindent \textbf{(P\_1,3.12.3)} kartṛgrahaṇāt ca eṣaḥ doṣaḥ . \\
\noindent \textbf{(P\_1,3.12.3)} prakṛtyarthaniyame śeṣagrahaṇam śakyam akartum . \\
\noindent \textbf{(P\_1,3.12.3)} katham . \\
\noindent \textbf{(P\_1,3.12.3)} prakṛtayarthau niyatau . \\
\noindent \textbf{(P\_1,3.12.3)} pratyayāḥ aniyatāḥ . \\
\noindent \textbf{(P\_1,3.12.3)} tataḥ vakṣyāmi parasmaipadam bhavati iti . \\
\noindent \textbf{(P\_1,3.12.3)} tat niyamārtham bhaviṣyati . \\
\noindent \textbf{(P\_1,3.12.3)} yatra parasmaipadam ca anyat ca prāpnoti tatra parasmaipadam eva bhavati iti . \\
\noindent \textbf{(P\_1,3.12.3)} tat tarhi pratyayaniyame dvitīyam śeṣagrahaṇam kartavyam . \\
\noindent \textbf{(P\_1,3.12.3)} na kartavyam . \\
\noindent \textbf{(P\_1,3.12.3)} yogavibhāgaḥ kariṣyate . \\
\noindent \textbf{(P\_1,3.12.3)} anudāttaṅitaḥ ātmanepadam . \\
\noindent \textbf{(P\_1,3.12.3)} tataḥ bhāvakarmaṇoḥ . \\
\noindent \textbf{(P\_1,3.12.3)} tataḥ kartari . \\
\noindent \textbf{(P\_1,3.12.3)} kartari ca ātmanepadam bhavati bhāvakarmaṇoḥ . \\
\noindent \textbf{(P\_1,3.12.3)} tataḥ karkmavyatihāre . \\
\noindent \textbf{(P\_1,3.12.3)} kartari iti eva . \\
\noindent \textbf{(P\_1,3.12.3)} bhāvakarmaṇoḥ iti nivṛttam . \\
\noindent \textbf{(P\_1,3.12.3)} yathā eva tarhi karmaṇi kartari bhavati evam bhāve api kartari prāpnoti . \\
\noindent \textbf{(P\_1,3.12.3)} eti jīvantam ānandaḥ . \\
\noindent \textbf{(P\_1,3.12.3)} na asya kim cit rujati iti . \\
\noindent \textbf{(P\_1,3.12.3)} dvitīyaḥ yogavibhāgaḥ kariṣyate . \\
\noindent \textbf{(P\_1,3.12.3)} anudāttaṅitaḥ ātmanepadam . \\
\noindent \textbf{(P\_1,3.12.3)} tataḥ bhāve . \\
\noindent \textbf{(P\_1,3.12.3)} tataḥ karmaṇi . \\
\noindent \textbf{(P\_1,3.12.3)} karmaṇi ca ātmanepadam bhavati . \\
\noindent \textbf{(P\_1,3.12.3)} tataḥ kartari . \\
\noindent \textbf{(P\_1,3.12.3)} kartari ca ātmanepadam bhavati . \\
\noindent \textbf{(P\_1,3.12.3)} karmaṇi iti anuvartate . \\
\noindent \textbf{(P\_1,3.12.3)} bhāve iti nivṛttam . \\
\noindent \textbf{(P\_1,3.12.3)} tataḥ karmavyatihāre . \\
\noindent \textbf{(P\_1,3.12.3)} kartari iti eva . \\
\noindent \textbf{(P\_1,3.12.3)} karmaṇi iti nivṛttam . \\
\noindent \textbf{(P\_1,3.12.3)} evam api śeṣagrahaṇam kartavyam anuparābhyām kṛñaḥ iti evamartham . \\
\noindent \textbf{(P\_1,3.12.3)} iha mā bhūt anukriyate svayam eva . \\
\noindent \textbf{(P\_1,3.12.3)} parākriyate svayam eva . \\
\noindent \textbf{(P\_1,3.12.3)} nanu ca etat api yogavibhāgāt eva siddham . \\
\noindent \textbf{(P\_1,3.12.3)} na sidhyati . \\
\noindent \textbf{(P\_1,3.12.3)} anantarā yā prāptiḥ sā yogavibhāgena śakyā bādhitum . \\
\noindent \textbf{(P\_1,3.12.3)} kutaḥ etat . \\
\noindent \textbf{(P\_1,3.12.3)} anantarasya vidhiḥ vā bhavati pratiṣedhaḥ vā iti . \\
\noindent \textbf{(P\_1,3.12.3)} parā prāptiḥ apratiṣiddhā . \\
\noindent \textbf{(P\_1,3.12.3)} tayā bhaviṣyati . \\
\noindent \textbf{(P\_1,3.12.3)} nanu ca iyam prāptiḥ pūrvām prāptim bādhate . \\
\noindent \textbf{(P\_1,3.12.3)} na utsahate pratiṣiddhā satī bādhitum . \\
\noindent \textbf{(P\_1,3.12.3)} evam tarhi kartari karmavyatihāre iti atra kartṛgrahaṇam pratyākhyāyate . \\
\noindent \textbf{(P\_1,3.12.3)} tat prakṛtam uttaratra anuvartiṣyate . \\
\noindent \textbf{(P\_1,3.12.3)} śeṣāt kartari kartari iti . \\
\noindent \textbf{(P\_1,3.12.3)} kimartham idam kartari kartari iti . \\
\noindent \textbf{(P\_1,3.12.3)} kartā eva yaḥ kartā tatra yathā syāt . \\
\noindent \textbf{(P\_1,3.12.3)} kartā ca anyaḥ ca yaḥ kartā tatra mā bhūt iti . \\
\noindent \textbf{(P\_1,3.12.3)} tataḥ anuparābhyām kṛñaḥ . \\
\noindent \textbf{(P\_1,3.12.3)} kartari kartari iti eva . \\
\noindent \textbf{(P\_1,3.14.1)} kriyāvyatirhāre iti vaktavyam . \\
\noindent \textbf{(P\_1,3.14.1)} karmavyatirhāre iti ucyamāne iha prasajyeta devadattasya dhānyam vyatilunanti iti iha ca na syāt vyatilunate vyatipunate iti . \\
\noindent \textbf{(P\_1,3.14.1)} tat tarhi vaktavyam . \\
\noindent \textbf{(P\_1,3.14.1)} na vaktavyam . \\
\noindent \textbf{(P\_1,3.14.1)} kriyām hi loke karma iti upacaranti . \\
\noindent \textbf{(P\_1,3.14.1)} kām kriyām kariṣyasi . \\
\noindent \textbf{(P\_1,3.14.1)} kim karma kariṣyasi iti . \\
\noindent \textbf{(P\_1,3.14.1)} evam api kartavyam . \\
\noindent \textbf{(P\_1,3.14.1)} kṛtrimākṛtrimayoḥ kṛtrime sampratyayaḥ bhavati . \\
\noindent \textbf{(P\_1,3.14.1)} kriyā api kṛtrimam karma . \\
\noindent \textbf{(P\_1,3.14.1)} na sidhyati . \\
\noindent \textbf{(P\_1,3.14.1)} kartuḥ īpsitatamam karma iti ucyate . \\
\noindent \textbf{(P\_1,3.14.1)} katham ca kriyā nāma kriyepsitatamā syāt . \\
\noindent \textbf{(P\_1,3.14.1)} kriyā api kriyepsitatamā bhavati . \\
\noindent \textbf{(P\_1,3.14.1)} kayā kriyayā . \\
\noindent \textbf{(P\_1,3.14.1)} sampaśyatikriyayā prārthayatikriyayā adhyavasyatikriyayā vā . \\
\noindent \textbf{(P\_1,3.14.1)} iha yaḥ eṣaḥ manuṣyaḥ prekṣāpūrvakārī bhavati salḥ buddhyā tāvat kam cit artham sampaśyati . \\
\noindent \textbf{(P\_1,3.14.1)} sandṛṣṭe prārthanā prārthite adhyavasāyaḥ adhyavasāye ārambhaḥ ārambhe nirvṛttiḥ nirvṛttau phalāvaptiḥ . \\
\noindent \textbf{(P\_1,3.14.1)} evam kriyā api kṛtrimam karma . \\
\noindent \textbf{(P\_1,3.14.1)} evam api ubhayoḥ kṛtrimākṛtrimayoḥ ubhayagatiḥ prasajyeta . \\
\noindent \textbf{(P\_1,3.14.1)} tasmāt kriyāvyatihāre iti vaktavyam . \\
\noindent \textbf{(P\_1,3.14.1)} na vaktavyam . \\
\noindent \textbf{(P\_1,3.14.1)} iha kartari vyatihāre iti iyatā siddham . \\
\noindent \textbf{(P\_1,3.14.1)} saḥ ayam evam siddhe sati yat karmagrahaṇam karoti tasya etat prayojanam kriyāvyatihāre yathā syāt karmavyatihāre mā bhūt iti . \\
\noindent \textbf{(P\_1,3.14.2)} atha kartṛgrahaṇam kimartham . \\
\noindent \textbf{(P\_1,3.14.2)} karmavyatihārādiṣu kartṛgrahaṇam bhāvakarmanivṛttyartham . \\
\noindent \textbf{(P\_1,3.14.2)} karmavyatihārādiṣu kartṛgrahaṇam kriyate bhāvakarmaṇoḥ anena ātmanepadam mā bhūt iti . \\
\noindent \textbf{(P\_1,3.14.2)} itarathā hi tatra pratiṣedhe bhāvakarmaṇoḥ pratiṣedhaḥ . \\
\noindent \textbf{(P\_1,3.14.2)} akriyamāṇe kartṛgrahaṇe bhāvakarmaṇoḥ api ātmanepadam prasajyeta . \\
\noindent \textbf{(P\_1,3.14.2)} tatra kaḥ doṣaḥ . \\
\noindent \textbf{(P\_1,3.14.2)} tatra pratiṣedhe bhāvakarmaṇoḥ pratiṣedhaḥ . \\
\noindent \textbf{(P\_1,3.14.2)} tatra kaḥ doṣaḥ . \\
\noindent \textbf{(P\_1,3.14.2)} tatra pratiṣedhe bhāvakarmaṇoḥ api anena ātmanepadasya pratiṣedhaḥ prasajyeta . \\
\noindent \textbf{(P\_1,3.14.2)} vyatigamyante grāmāḥ vyatihanyante dasyavaḥ iti . \\
\noindent \textbf{(P\_1,3.14.2)} na vā anantarasya pratiṣedhāt . \\
\noindent \textbf{(P\_1,3.14.2)} na vā eṣaḥ doṣaḥ . \\
\noindent \textbf{(P\_1,3.14.2)} kim kāraṇam . \\
\noindent \textbf{(P\_1,3.14.2)} anantarasya pratiṣedhāt . \\
\noindent \textbf{(P\_1,3.14.2)} anantaram yat ātmanepadavidhānam tasya pratiṣedhāt . \\
\noindent \textbf{(P\_1,3.14.2)} kutaḥ etat . \\
\noindent \textbf{(P\_1,3.14.2)} anantarasya vidhiḥ vā bhavati pratiṣedhaḥ vā iti . \\
\noindent \textbf{(P\_1,3.14.2)} pūrvā prāptiḥ apratiṣiddhā tayā bhaviṣyati . \\
\noindent \textbf{(P\_1,3.14.2)} nanu ca iyam prāptiḥ pūrvām prāptim bādhate . \\
\noindent \textbf{(P\_1,3.14.2)} na utsahate pratiṣiddhā satī bādhitum . \\
\noindent \textbf{(P\_1,3.14.2)} uttarārtham tarhi kartṛgrahaṇam kartavyam . \\
\noindent \textbf{(P\_1,3.14.2)} na kartavyam . \\
\noindent \textbf{(P\_1,3.14.2)} kriyate tatra eva śeṣāt kartari parasmaipadam iti . \\
\noindent \textbf{(P\_1,3.14.2)} dvitīyam kartṛgrahaṇam kartavyam . \\
\noindent \textbf{(P\_1,3.14.2)} kim prayojanam . \\
\noindent \textbf{(P\_1,3.14.2)} kartā eva yaḥ kartā tatra yathā syāt . \\
\noindent \textbf{(P\_1,3.14.2)} karta ca anyaḥ ca yaḥ kartā tatra mā bhūt iti . \\
\noindent \textbf{(P\_1,3.15)} pratiṣedhe hasādīnām upasaṅkhyānam . \\
\noindent \textbf{(P\_1,3.15)} pratiṣedhe hasādīnām upasaṅkhyānam kartavyam . \\
\noindent \textbf{(P\_1,3.15)} vyatihasanti vayatijalpanti vyatipaṭhanti . \\
\noindent \textbf{(P\_1,3.15)} harivahyoḥ apratiṣedhaḥ . \\
\noindent \textbf{(P\_1,3.15)} harivahyoḥ apratiṣedhaḥ bhavati iti vaktavyam . \\
\noindent \textbf{(P\_1,3.15)} sampraharante rājānaḥ . \\
\noindent \textbf{(P\_1,3.15)} saṃvivahante gargaiḥ iti . \\
\noindent \textbf{(P\_1,3.15)} na vahiḥ gatyarthaḥ . \\
\noindent \textbf{(P\_1,3.15)} deśāntaraprāpaṇakriyaḥ vahiḥ . \\
\noindent \textbf{(P\_1,3.16)} prarasparopapadāt ca . \\
\noindent \textbf{(P\_1,3.16)} prarasparopapadāt ca iti vaktavyam . \\
\noindent \textbf{(P\_1,3.16)} parasparasya vyatilunanti . \\
\noindent \textbf{(P\_1,3.16)} parasparasya vyatipunanti . \\
\noindent \textbf{(P\_1,3.19)} upasargagrahaṇam kartavyam . \\
\noindent \textbf{(P\_1,3.19)} parā jayati senā iti . \\
\noindent \textbf{(P\_1,3.19)} tat tarhi vaktavyam . \\
\noindent \textbf{(P\_1,3.19)} na vaktavyam . \\
\noindent \textbf{(P\_1,3.19)} yadi api tāvat ayam parāśabdaḥ dṛṣṭāpacāraḥ upasargaḥ ca anupasargaḥ ca ayam tu khalu viśabdaḥ adṛṣṭāpacāraḥ upasargaḥ eva . \\
\noindent \textbf{(P\_1,3.19)} tasya asya kaḥ dvitīyaḥ sahāyaḥ bhavitum arhati anyat ataḥ upasargāt . \\
\noindent \textbf{(P\_1,3.19)} tat yathā asya goḥ dvitīyena arthaḥ iti gauḥ eva upādīyate na aśvaḥ na gardabhaḥ . \\
\noindent \textbf{(P\_1,3.20)} āṅaḥ daḥ avyasanakriyasya . \\
\noindent \textbf{(P\_1,3.20)} āṅaḥ daḥ avyasanakriyasya iti vaktavyam . \\
\noindent \textbf{(P\_1,3.20)} iha api yathā syāt . \\
\noindent \textbf{(P\_1,3.20)} vipādikām vyādadāti . \\
\noindent \textbf{(P\_1,3.20)} kūlam vyādadāti iti . \\
\noindent \textbf{(P\_1,3.20)} tat tarhi vaktavyam . \\
\noindent \textbf{(P\_1,3.20)} na vaktavyam . \\
\noindent \textbf{(P\_1,3.20)} iha āṅaḥ daḥ anāsye iti iyatā siddham . \\
\noindent \textbf{(P\_1,3.20)} saḥ ayam evam siddhe sati yat viharaṇagrahaṇam karoti tasya etat prayojanam āsyaviharaṇasamānakriyāt api yathā syāt . \\
\noindent \textbf{(P\_1,3.20)} yathājātīyakā ca āsyaviharaṇakriyā tathājātīyakā atra api . \\
\noindent \textbf{(P\_1,3.20)} svāṅgakarmāt ca . \\
\noindent \textbf{(P\_1,3.20)} svāṅgakarmāt ca iti vaktavyam . \\
\noindent \textbf{(P\_1,3.20)} iha mā bhūt . \\
\noindent \textbf{(P\_1,3.20)} vyādadate pipīlikāḥ pataṅgamukham iti . \\
\noindent \textbf{(P\_1,3.21)} upasargagrahaṇam kartavyam . \\
\noindent \textbf{(P\_1,3.21)} iha mā bhūt . \\
\noindent \textbf{(P\_1,3.21)} anu krīḍati māṇavakam . \\
\noindent \textbf{(P\_1,3.21)} samaḥ akūjane . \\
\noindent \textbf{(P\_1,3.21)} samaḥ akūjane iti vaktavyam . \\
\noindent \textbf{(P\_1,3.21)} iha mā bhūt . \\
\noindent \textbf{(P\_1,3.21)} saṅkrīḍanti śakaṭāni . \\
\noindent \textbf{(P\_1,3.21)} āgameḥ kṣamāyām . \\
\noindent \textbf{(P\_1,3.21)} āgameḥ kṣamāyām upasaṅkhyanam kartavyam . \\
\noindent \textbf{(P\_1,3.21)} āgamayasva tāvat māṇavka . \\
\noindent \textbf{(P\_1,3.21)} śikṣeḥ jijñāsāyām . \\
\noindent \textbf{(P\_1,3.21)} śikṣeḥ jijñāsāyām upasaṅkhyanam kartavyam . \\
\noindent \textbf{(P\_1,3.21)} vidyāsu śikṣate . \\
\noindent \textbf{(P\_1,3.21)} dhanuṣi śikṣate . \\
\noindent \textbf{(P\_1,3.21)} kirateḥ harṣajīvikākulāyakaraṇeṣu . \\
\noindent \textbf{(P\_1,3.21)} kirateḥ harṣajīvikākulāyakaraṇeṣu upasaṅkhyanam kartavyam . \\
\noindent \textbf{(P\_1,3.21)} apaskirate vṛṣabhaḥ hṛṣṭaḥ . \\
\noindent \textbf{(P\_1,3.21)} apaskirate kukkuṭaḥ bhakṣārthī . \\
\noindent \textbf{(P\_1,3.21)} apaskirate śvā āśrayārthī . \\
\noindent \textbf{(P\_1,3.21)} harateḥ gatatācchīlye . \\
\noindent \textbf{(P\_1,3.21)} harateḥ gatatācchīlye upasaṅkhyanam kartavyam . \\
\noindent \textbf{(P\_1,3.21)} paitṛkam aśvāḥ anuharante . \\
\noindent \textbf{(P\_1,3.21)} mātṛkam gāvaḥ anuharante . \\
\noindent \textbf{(P\_1,3.21)} āṅi nupracchyoḥ . \\
\noindent \textbf{(P\_1,3.21)} āṅi nupracchyoḥ upasaṅkhyanam kartavyam . \\
\noindent \textbf{(P\_1,3.21)} ānute śṛgālaḥ . \\
\noindent \textbf{(P\_1,3.21)} āpṛcchate gurum . \\
\noindent \textbf{(P\_1,3.21)} āśiṣi nāthaḥ . \\
\noindent \textbf{(P\_1,3.21)} āśiṣi nāthaḥ upasaṅkhyanam kartavyam . \\
\noindent \textbf{(P\_1,3.21)} sarpiṣaḥ nāthate . \\
\noindent \textbf{(P\_1,3.21)} madhunaḥ nāthate . \\
\noindent \textbf{(P\_1,3.21)} śapaḥ upalambhane . \\
\noindent \textbf{(P\_1,3.21)} śapaḥ upalambhane upasaṅkhyanam kartavyam . \\
\noindent \textbf{(P\_1,3.21)} devadattāya śapate . \\
\noindent \textbf{(P\_1,3.21)} yajñadattāya śapate . \\
\noindent \textbf{(P\_1,3.22)} āṅaḥ sthaḥ pratijñāne . \\
\noindent \textbf{(P\_1,3.22)} āṅaḥ sthaḥ pratijñāne iti vaktavyam . \\
\noindent \textbf{(P\_1,3.22)} astim sakāram ātiṣṭhate . \\
\noindent \textbf{(P\_1,3.22)} āgamau guṇavṛddhī ātiṣṭhate . \\
\noindent \textbf{(P\_1,3.22)} vikārau guṇavṛddhī ātiṣṭhate . \\
\noindent \textbf{(P\_1,3.24)} udaḥ īhāyām . \\
\noindent \textbf{(P\_1,3.24)} udaḥ īhāyām iti vaktavyam . \\
\noindent \textbf{(P\_1,3.24)} iha mā bhūt . \\
\noindent \textbf{(P\_1,3.24)} uttiṣṭhati senā iti . \\
\noindent \textbf{(P\_1,3.25)} upāt pūjāsaṅgatakaraṇayoḥ . \\
\noindent \textbf{(P\_1,3.25)} upāt pūjāsaṅgatakaraṇayoḥ iti vaktavyam . \\
\noindent \textbf{(P\_1,3.25)} ādityam upatiṣṭhate . \\
\noindent \textbf{(P\_1,3.25)} candramasam upatiṣṭhate . \\
\noindent \textbf{(P\_1,3.25)} saṅgatakaraṇe . \\
\noindent \textbf{(P\_1,3.25)} rathikān upatiṣṭhate . \\
\noindent \textbf{(P\_1,3.25)} aśvārohān upatiṣṭhate . \\
\noindent \textbf{(P\_1,3.25)} bahūnām api acittānām ekaḥ bhavati cittavān . \\
\noindent \textbf{(P\_1,3.25)} paśya vānarasainye asmin yat arkam upatiṣṭhate . \\
\noindent \textbf{(P\_1,3.25)} mā evam maṃsthāḥ sacittaḥ ayam eṣaḥ api yathā vayam . \\
\noindent \textbf{(P\_1,3.25)} etat api asya kāpeyam yat arkam upatiṣṭhati . \\
\noindent \textbf{(P\_1,3.25)} aparaḥ āha : upāt devapūjāsaṅgatakaraṇamitrakaraṇapathiṣu iti vaktavyam . \\
\noindent \textbf{(P\_1,3.25)} saṅgatakaraṇe udāhṛtam . \\
\noindent \textbf{(P\_1,3.25)} mitrakaraṇe . \\
\noindent \textbf{(P\_1,3.25)} rathikān upatiṣṭhate . \\
\noindent \textbf{(P\_1,3.25)} aśvārohān upatiṣṭhate . \\
\noindent \textbf{(P\_1,3.25)} pathi. ayam panthāḥ srughnam upatiṣṭhate . \\
\noindent \textbf{(P\_1,3.25)} ayam panthāḥ sāketam upatiṣthate . \\
\noindent \textbf{(P\_1,3.25)} vā lipsāyām . \\
\noindent \textbf{(P\_1,3.25)} vā lipsāyām iti vaktavyam . \\
\noindent \textbf{(P\_1,3.25)} bhikṣukaḥ brāhmaṇakulam upatiṣṭhate . \\
\noindent \textbf{(P\_1,3.25)} bhikṣukaḥ brāhmaṇakulam upatiṣṭhati vā \\
\noindent \textbf{(P\_1,3.27)} akarmakāt iti eva . \\
\noindent \textbf{(P\_1,3.27)} ut tapati suvarṇam suvarṇakāraḥ . \\
\noindent \textbf{(P\_1,3.27)} svāṅgakarmakāt ca . \\
\noindent \textbf{(P\_1,3.27)} svāṅgakarmakāt ca iti vaktavyam . \\
\noindent \textbf{(P\_1,3.27)} uttapate pāṇī . \\
\noindent \textbf{(P\_1,3.27)} vitapate pāṇī . \\
\noindent \textbf{(P\_1,3.27)} uttapate pṛṣṭham . \\
\noindent \textbf{(P\_1,3.27)} vitapate pṛṣṭham . \\
\noindent \textbf{(P\_1,3.27)} atha udbhibhyām iti atra kim pratyudāhriyate . \\
\noindent \textbf{(P\_1,3.27)} niṣṭapyate iti . \\
\noindent \textbf{(P\_1,3.27)} kim punaḥ kāraṇam ātmanepadam eva udāhriyate na parasmaipadam pratyudāhāryam syāt . \\
\noindent \textbf{(P\_1,3.27)} tapiḥ ayam akarmakaḥ . \\
\noindent \textbf{(P\_1,3.27)} akarmakāḥ ca api sopasargāḥ sakarmakāḥ bhavanti . \\
\noindent \textbf{(P\_1,3.27)} na ca antareṇa karmakartāram sakarmakāḥ akarmakāḥ bhavanti . \\
\noindent \textbf{(P\_1,3.27)} yat ucyate na ca antareṇa karmakartāram sakarmakāḥ akarmakāḥ bhavanti iti antareṇa api karmakartāram sakarmakāḥ akarmakāḥ bhavanti . \\
\noindent \textbf{(P\_1,3.27)} tat yathā . \\
\noindent \textbf{(P\_1,3.27)} nadīvahati iti akarmakaḥ . \\
\noindent \textbf{(P\_1,3.27)} bhāram vahati iti sakarmakaḥ . \\
\noindent \textbf{(P\_1,3.27)} tasmāt niṣṭapati iti pratyudāhartavyam . \\
\noindent \textbf{(P\_1,3.28)} akarmakāt iti eva . \\
\noindent \textbf{(P\_1,3.28)} āyacchati rajjum kūpāt . \\
\noindent \textbf{(P\_1,3.28)} āhanti vṛṣalam pādena . \\
\noindent \textbf{(P\_1,3.28)} svāṅgakarmakāt ca . \\
\noindent \textbf{(P\_1,3.28)} svāṅgakarmakāt ca iti vaktavyam . \\
\noindent \textbf{(P\_1,3.28)} āyacchate pāṇī . \\
\noindent \textbf{(P\_1,3.28)} āhate udaram iti . \\
\noindent \textbf{(P\_1,3.29)} samaḥ gamyādiṣu vidipracchisvaratīnām upasaṅkhyānam . \\
\noindent \textbf{(P\_1,3.29)} samaḥ gamyādiṣu vidipracchisvaratīnām upasaṅkhyānam kartavyam . \\
\noindent \textbf{(P\_1,3.29)} saṃvitte sampṛcchate saṃsvarate . \\
\noindent \textbf{(P\_1,3.29)} artiśrudṛśibhyaḥ ca . \\
\noindent \textbf{(P\_1,3.29)} artiśrudṛśibhyaḥ ca iti vaktavyam . \\
\noindent \textbf{(P\_1,3.29)} ma samṛta mā samṛṣātām mā samṛṣata . \\
\noindent \textbf{(P\_1,3.29)} arti . \\
\noindent \textbf{(P\_1,3.29)} śru . \\
\noindent \textbf{(P\_1,3.29)} saṃśrṇute . \\
\noindent \textbf{(P\_1,3.29)} dṛśi . \\
\noindent \textbf{(P\_1,3.29)} sampaśyate . \\
\noindent \textbf{(P\_1,3.29)} upasargāt asyatyūhyoḥ vāvacanam . \\
\noindent \textbf{(P\_1,3.29)} upasargāt astyūhyoḥ vā iti vaktavyam . \\
\noindent \textbf{(P\_1,3.29)} nirasyati nirasyate . \\
\noindent \textbf{(P\_1,3.29)} samūhati samūhate . \\
\noindent \textbf{(P\_1,3.40)} jyotiṣām udgamane. jyotiṣām udgamane iti vaktavyam . \\
\noindent \textbf{(P\_1,3.40)} iha mā bhūt . \\
\noindent \textbf{(P\_1,3.40)} ākrāmati dhūmaḥ harmyatalam iti . \\
\noindent \textbf{(P\_1,3.48)} vyaktavācām iti kimartham . \\
\noindent \textbf{(P\_1,3.48)} varatanu sampravadanti kukkuṭāḥ . \\
\noindent \textbf{(P\_1,3.48)} vyaktavācām iti ucyamāne api atra prāpnoti . \\
\noindent \textbf{(P\_1,3.48)} ete api hi vyaktavācaḥ . \\
\noindent \textbf{(P\_1,3.48)} ātaḥ ca vyaktavācaḥ kukkuṭena udite ucyate kukkuṭaḥ vadati iti . \\
\noindent \textbf{(P\_1,3.48)} evam tarhi vyaktavācām iti ucyate . \\
\noindent \textbf{(P\_1,3.48)} sarve eva hi vyaktavācaḥ . \\
\noindent \textbf{(P\_1,3.48)} tatra prakarṣagatiḥ vijñāsyate . \\
\noindent \textbf{(P\_1,3.48)} sādhīyaḥ ye vyaktavācaḥ iti . \\
\noindent \textbf{(P\_1,3.48)} ke ca sādhīyaḥ . \\
\noindent \textbf{(P\_1,3.48)} yeṣām vāci akārādayaḥ varṇāḥ vyajyante . \\
\noindent \textbf{(P\_1,3.48)} na ca eteṣām vāci akārādayaḥ varṇāḥ vyajyante . \\
\noindent \textbf{(P\_1,3.48)} eteṣām api vāci akārādayaḥ varṇāḥ vyajyante . \\
\noindent \textbf{(P\_1,3.48)} ātaḥ ca vyajyante evam hi āhuḥ kukkuṭāḥ kukkuṭ iti . \\
\noindent \textbf{(P\_1,3.48)} na evam te āhuḥ . \\
\noindent \textbf{(P\_1,3.48)} anukaraṇam etat teṣām . \\
\noindent \textbf{(P\_1,3.48)} atha vā na evam vijñāyate vyaktā vāk yeṣām te ime vyaktavācaḥ iti . \\
\noindent \textbf{(P\_1,3.48)} katham tarhi . \\
\noindent \textbf{(P\_1,3.48)} vyaktā vāci varṇāḥ yeṣām te ime vyaktavācaḥ iti . \\
\noindent \textbf{(P\_1,3.51)} avād graḥ girateḥ . \\
\noindent \textbf{(P\_1,3.51)} avād graḥ iti atra girateḥ iti vaktavyam . \\
\noindent \textbf{(P\_1,3.51)} gṛṇāteḥ mā bhūt . \\
\noindent \textbf{(P\_1,3.51)} tat tarhi vaktavyam . \\
\noindent \textbf{(P\_1,3.51)} na vaktavyam prayogābhāvāt . \\
\noindent \textbf{(P\_1,3.51)} avāt graḥ iti ucyate na ca āpūrvasya gṛṇāteḥ prayogaḥ asti . \\
\noindent \textbf{(P\_1,3.54)} tṛtīyāyuktāt iti kimartham . \\
\noindent \textbf{(P\_1,3.54)} ubhau lokau cañcarasi imam ca amum ca devala . \\
\noindent \textbf{(P\_1,3.54)} tṛtīyāyuktāt iti ucyamāne api atra prāpnoti . \\
\noindent \textbf{(P\_1,3.54)} atra api hi tṛtīyayā yogaḥ . \\
\noindent \textbf{(P\_1,3.54)} evam tarhi tṛtīyāyuktāt iti ucyate sarvatra ca tṛtīyayā yogaḥ . \\
\noindent \textbf{(P\_1,3.54)} tatra prakarṣagatiḥ vijñāsyate : sādhīyaḥ yatra tṛtīyayā yogaḥ iti . \\
\noindent \textbf{(P\_1,3.54)} kva ca sādhīyaḥ . \\
\noindent \textbf{(P\_1,3.54)} yatra tṛtīyayā yogaḥ śrūyate . \\
\noindent \textbf{(P\_1,3.55)} sā cet tṛtīyā caturthyarthe iti ucyate . \\
\noindent \textbf{(P\_1,3.55)} katham nāma tṛtīyā caturthyarthe syāt . \\
\noindent \textbf{(P\_1,3.55)} evam tarhi aśiṣṭavyavahāre anena tṛtīyā ca vidhīyate ātmanepadam ca . \\
\noindent \textbf{(P\_1,3.55)} dāsyā samprayacchate . \\
\noindent \textbf{(P\_1,3.55)} vṛṣalyā sampracchate . \\
\noindent \textbf{(P\_1,3.55)} yaḥ hi śiṣṭavyavahāraḥ brāhmaṇībhyaḥ samprayacchati iti eva tatra bhavitavyam . \\
\noindent \textbf{(P\_1,3.55)} yadi evam na arthaḥ anena yogena . \\
\noindent \textbf{(P\_1,3.55)} kena idānīm tṛtīyā bhaviṣyati ātmanepadam ca . \\
\noindent \textbf{(P\_1,3.55)} sahayukte tṛtīyā syāt vyatihāre taṅaḥ vidhiḥ . \\
\noindent \textbf{(P\_1,3.55)} sahayukte apradhāne iti eva tṛtīyā bhaviṣyati . \\
\noindent \textbf{(P\_1,3.55)} kartari karmavyatihāre iti ātmanepadam . \\
\noindent \textbf{(P\_1,3.56)} iha kasmāt na bhavati . \\
\noindent \textbf{(P\_1,3.56)} svam śāṭakāntam upayacchati iti . \\
\noindent \textbf{(P\_1,3.56)} asvam yadā svam karoti tadā bhavitavyam . \\
\noindent \textbf{(P\_1,3.56)} yadi evam svīkaraṇe iti prāpnoti . \\
\noindent \textbf{(P\_1,3.56)} vicitrāḥ taddhitavṛttayaḥ . \\
\noindent \textbf{(P\_1,3.56)} na ataḥ taddhitaḥ utpadyate . \\
\noindent \textbf{(P\_1,3.58)} anoḥ jñaḥ pratiṣedhe sakarmakavacanam . \\
\noindent \textbf{(P\_1,3.58)} anoḥ jñaḥ pratiṣedhe sakarmakagrahaṇam kartavyam . \\
\noindent \textbf{(P\_1,3.58)} iha ma bhūt . \\
\noindent \textbf{(P\_1,3.58)} auṣadhasya anujijñāsate iti . \\
\noindent \textbf{(P\_1,3.58)} na vā akarmakasya uttareṇa vidhānāt . \\
\noindent \textbf{(P\_1,3.58)} na vā kartavyam . \\
\noindent \textbf{(P\_1,3.58)} kim kāraṇam . \\
\noindent \textbf{(P\_1,3.58)} akarmakasya uttareṇa vidhānāt . \\
\noindent \textbf{(P\_1,3.58)} akarmakāt janāteḥ uttareṇa yogena ātmanepadam vidhīyate pūrvavat sanaḥ iti . \\
\noindent \textbf{(P\_1,3.58)} pratiṣedhaḥ pūrvasya ca . \\
\noindent \textbf{(P\_1,3.58)} pūrvasya ca ayam pratiṣedhaḥ . \\
\noindent \textbf{(P\_1,3.58)} saḥ ca sakarmakārthaḥ ārambhaḥ . \\
\noindent \textbf{(P\_1,3.58)} katham punaḥ jñāyate pūrvaysa ayam pratiṣedhaḥ iti . \\
\noindent \textbf{(P\_1,3.58)} anantarasya vidhiḥ vā bhavati pratiṣedhaḥ vā iti . \\
\noindent \textbf{(P\_1,3.58)} katham punaḥ jñāyate sakarmakārthaḥ ārambhaḥ iti . \\
\noindent \textbf{(P\_1,3.58)} akarmakāt jānāteḥ sanaḥ ātmanepadavacane prayojanam na asti iti kṛtvā sakarmakārthaḥ vijñāyate . \\
\noindent \textbf{(P\_1,3.60.1)} śadeḥ śitaḥ parasmaipadāśrayatvāt ātmanepadābhāvaḥ . \\
\noindent \textbf{(P\_1,3.60.1)} śadeḥ śitaḥ parasmaipadāśrayatvāt ātmanepadasya abhāvaḥ . \\
\noindent \textbf{(P\_1,3.60.1)} śīyate śīyete śīyante . \\
\noindent \textbf{(P\_1,3.60.1)} kim ca bhoḥ śadeḥ śit parasmaipadeṣu iti ucyate . \\
\noindent \textbf{(P\_1,3.60.1)} na khalu parasmaipadeṣu iti ucyate parasmaipadeṣu tu vijñāyate . \\
\noindent \textbf{(P\_1,3.60.1)} katham anudāttaṅitaḥ ātmanepadam bhāvakarmaṇoḥ ātmanepadam iti etau dvau yogau uktvā śeṣāt kartari parasmaipadam ucyate . \\
\noindent \textbf{(P\_1,3.60.1)} evam na ca parasmaipadeṣu ucyate parasmaipadeṣu ca vijñāyate . \\
\noindent \textbf{(P\_1,3.60.1)} kaḥ punaḥ arhati etau dvau yogau uktvā śeṣāt kartari parasmaipadam vaktum . \\
\noindent \textbf{(P\_1,3.60.1)} kim tarhi . \\
\noindent \textbf{(P\_1,3.60.1)} aviśeṣeṇa sarvam ātmanepadaprakaraṇam anukramya śeṣāt kartari parasmaipadam iti ucyate . \\
\noindent \textbf{(P\_1,3.60.1)} evam api parasmaipadāśrayaḥ bhavati . \\
\noindent \textbf{(P\_1,3.60.1)} katham . \\
\noindent \textbf{(P\_1,3.60.1)} idam tāvad ayam praṣṭavyaḥ . \\
\noindent \textbf{(P\_1,3.60.1)} yadi idam na ucyeta kim iha syāt iti . \\
\noindent \textbf{(P\_1,3.60.1)} parasmaipadam iti āha . \\
\noindent \textbf{(P\_1,3.60.1)} parasmaipadam iti cet parasmaipadāśrayaḥ bhavati . \\
\noindent \textbf{(P\_1,3.60.1)} siddham tu laḍādīnām ātmanepadavacanam . \\
\noindent \textbf{(P\_1,3.60.1)} siddham etat . \\
\noindent \textbf{(P\_1,3.60.1)} katham . \\
\noindent \textbf{(P\_1,3.60.1)} śadeḥ laḍādīnām ātmanepadam bhavati iti vaktavyam . \\
\noindent \textbf{(P\_1,3.60.1)} sidhyati . \\
\noindent \textbf{(P\_1,3.60.1)} sūtram tarhi bhidyate . \\
\noindent \textbf{(P\_1,3.60.1)} yathānyāsam eva astu . \\
\noindent \textbf{(P\_1,3.60.1)} nanu ca uktam śadeḥ śitaḥ parasmaipadāśrayatvāt ātmanepadābhāvaḥ iti . \\
\noindent \textbf{(P\_1,3.60.1)} na eṣaḥ doṣaḥ . \\
\noindent \textbf{(P\_1,3.60.1)} śitaḥ iti na eṣā pañcamī . \\
\noindent \textbf{(P\_1,3.60.1)} kā tarhi . \\
\noindent \textbf{(P\_1,3.60.1)} sambandhaṣaṣṭhī . \\
\noindent \textbf{(P\_1,3.60.1)} śitaḥ yaḥ śadiḥ . \\
\noindent \textbf{(P\_1,3.60.1)} kaḥ ca śitaḥ śadiḥ . \\
\noindent \textbf{(P\_1,3.60.1)} prakṛtiḥ . \\
\noindent \textbf{(P\_1,3.60.1)} śadeḥ śitprakṛteḥ iti . \\
\noindent \textbf{(P\_1,3.60.1)} atha vā āha ayam śadeḥ śitaḥ iti na ca śadiḥ śit asti . \\
\noindent \textbf{(P\_1,3.60.1)} te evam vijñāsyāmaḥ śadeḥ śidviṣayāt iti . \\
\noindent \textbf{(P\_1,3.60.1)} atha vā yadi api tāvat etat anyatra bhavati vikaraṇebhyaḥ niyamaḥ balīyān iti iha etat na asti . \\
\noindent \textbf{(P\_1,3.60.1)} vikaraṇaḥ hi iha āśrīyate śitaḥ iti . \\
\noindent \textbf{(P\_1,3.60.2)} upasargapūrvaniyame aḍvyavāye upasaṅkhyānam . \\
\noindent \textbf{(P\_1,3.60.2)} upasargapūrvasya niyame aḍvyavāye upasaṅkhyānam kartavyam . \\
\noindent \textbf{(P\_1,3.60.2)} nyaviśata vyakrīṇīta . \\
\noindent \textbf{(P\_1,3.60.2)} kim punaḥ kāraṇāt na sidhyati . \\
\noindent \textbf{(P\_1,3.60.2)} aṭā vyavahitatvāt . \\
\noindent \textbf{(P\_1,3.60.2)} nanu ca ayam aṭ dhātubhaktaḥ dhātugrahaṇena grahīṣyate . \\
\noindent \textbf{(P\_1,3.60.2)} na sidhyati . \\
\noindent \textbf{(P\_1,3.60.2)} aṅgasya hi aṭ ucyate vikaraṇāntam ca aṅgam . \\
\noindent \textbf{(P\_1,3.60.2)} saḥ asau saṅghātabhaktaḥ na śakyaḥ dhātugrahaṇena grahītum . \\
\noindent \textbf{(P\_1,3.60.2)} evam tarhi idam iha sampradhāryam : aṭ kriyatām vikaraṇaḥ iti . \\
\noindent \textbf{(P\_1,3.60.2)} kim atra kartavyam . \\
\noindent \textbf{(P\_1,3.60.2)} paratvāt aṭ āgamaḥ . \\
\noindent \textbf{(P\_1,3.60.2)} nityāḥ vikaraṇāḥ . \\
\noindent \textbf{(P\_1,3.60.2)} kṛte api aṭi prāpnuvanti akṛte api prāpnuvanti . \\
\noindent \textbf{(P\_1,3.60.2)} aṭ api nityaḥ . \\
\noindent \textbf{(P\_1,3.60.2)} kṛteṣu api vikaraṇeṣu prāpnoti akṛteṣu api prāpnoti . \\
\noindent \textbf{(P\_1,3.60.2)} anityaḥ aṭ . \\
\noindent \textbf{(P\_1,3.60.2)} anyasya kṛteṣu vikaraṇeṣu prāpnoti anyasya akṛteṣu . \\
\noindent \textbf{(P\_1,3.60.2)} śabdāntarasya ca prāpnuvan vidhiḥ anityaḥ bhavati . \\
\noindent \textbf{(P\_1,3.60.2)} evam tarhi idam iha sampradhāryam . \\
\noindent \textbf{(P\_1,3.60.2)} aṭ kriyatām lādeśaḥ iti . \\
\noindent \textbf{(P\_1,3.60.2)} kim atra kartavyam . \\
\noindent \textbf{(P\_1,3.60.2)} paratvāt aṭ āgamaḥ . \\
\noindent \textbf{(P\_1,3.60.2)} nityaḥ lādeśaḥ . \\
\noindent \textbf{(P\_1,3.60.2)} kṛte api aṭi prāpnoti akṛte api prāpnoti . \\
\noindent \textbf{(P\_1,3.60.2)} nityatvāt lādeśasya ātmanepade eva aḍāgamaḥ bhaviṣyati . \\
\noindent \textbf{(P\_1,3.60.2)} nityatvāt lādeśasya ātmanepade aṭ āgamaḥ iti cet aṭaḥ nityanimittatvāt ātmanepadābhāvaḥ . \\
\noindent \textbf{(P\_1,3.60.2)} nityatvāt lādeśasya ātmanepade eva aḍāgamaḥ iti cet evam ucyate . \\
\noindent \textbf{(P\_1,3.60.2)} aṭ api nityanimittaḥ . \\
\noindent \textbf{(P\_1,3.60.2)} kṛte api ladeśe prāpnoti akṛte api prāpnoti . \\
\noindent \textbf{(P\_1,3.60.2)} aṭaḥ nityanimittatvāt ātmanepadābhāvaḥ . \\
\noindent \textbf{(P\_1,3.60.2)} tasmāt upasaṅkhyānam . \\
\noindent \textbf{(P\_1,3.60.2)} tasmāt upasaṅkhyānam kartavyam . \\
\noindent \textbf{(P\_1,3.60.2)} na kartavyam . \\
\noindent \textbf{(P\_1,3.60.2)} antaraṅgaḥ tarhi lādeśaḥ . \\
\noindent \textbf{(P\_1,3.60.2)} na etat vivadāmahe antaraṅgaḥ na antaraṅgaḥ iti . \\
\noindent \textbf{(P\_1,3.60.2)} astu ayam nityaḥ antaraṅgaḥ ca . \\
\noindent \textbf{(P\_1,3.60.2)} atra khalu lādeśe kṛte trīṇi kāryāṇi yugapat prāpnuvanti : vikaraṇāḥ aṭ āgamaḥ niyamaḥ iti . \\
\noindent \textbf{(P\_1,3.60.2)} tat yadi sarvataḥ niyamaḥ labhyeta kṛtam syāt . \\
\noindent \textbf{(P\_1,3.60.2)} tat tu na labhyam . \\
\noindent \textbf{(P\_1,3.60.2)} atha api vikaraṇāt aṭ iti aṭ labhyeta evam api kṛtam syāt . \\
\noindent \textbf{(P\_1,3.60.2)} tat tu na labhyam . \\
\noindent \textbf{(P\_1,3.60.2)} kim kāraṇam . \\
\noindent \textbf{(P\_1,3.60.2)} āṅgāt pūrvam vikaraṇāḥ eṣitavyāḥ tarataḥ , taranti iti evamartham . \\
\noindent \textbf{(P\_1,3.60.2)} aḍāḍbhyām api anyat āṅgam pūrvam eṣitavyam upārcchati iti evamartham . \\
\noindent \textbf{(P\_1,3.60.2)} tatra hi āṭi kṛte sāṭkasya ṛcchibhāvaḥ prāpnoti . \\
\noindent \textbf{(P\_1,3.60.2)} nanu ca ṛcchibhāve kṛte śabdāntarasya akṛtaḥ āṭ iti kṛtvā punaḥ āṭ bhaviṣyati . \\
\noindent \textbf{(P\_1,3.60.2)} punaḥ ṛcchibhāvaḥ punaḥ āṭ iti cakrakam avyavasthā prāpnoti . \\
\noindent \textbf{(P\_1,3.60.2)} na eṣaḥ doṣaḥ . \\
\noindent \textbf{(P\_1,3.60.2)} yat tāvat ucyate āṅgāt pūrvam vikaraṇāḥ eṣitavyāḥ tarataḥ taranti iti evamartham iti . \\
\noindent \textbf{(P\_1,3.60.2)} bhavet siddham yatra vikaraṇāḥ nityāḥ āṅgam anityam tatra āṅgāt pūrvam vikaraṇāḥ syuḥ . \\
\noindent \textbf{(P\_1,3.60.2)} yatra tu khalu ubhayam nityam paratvāt tatra āṅgam tāvat bhavati . \\
\noindent \textbf{(P\_1,3.60.2)} yat api ucyate aḍāḍbhyām api anyat āṅgam pūrvam eṣitavyam upārcchati iti evamartham iti . \\
\noindent \textbf{(P\_1,3.60.2)} astu atra āṭ . \\
\noindent \textbf{(P\_1,3.60.2)} āṭi kṛte sāṭkasya ṛcchibhāve kṛte śabdāntarasya akṛtaḥ āṭ iti kṛtvā punaḥ āṭ bhaviṣyati . \\
\noindent \textbf{(P\_1,3.60.2)} nanu ca uktam punaḥ ṛcchibhāvaḥ punaḥ āṭ iti cakrakam avyavasthā prāpnoti . \\
\noindent \textbf{(P\_1,3.60.2)} na eṣaḥ doṣaḥ . \\
\noindent \textbf{(P\_1,3.60.2)} cakrakeṣu iṣṭataḥ vyavasthā . \\
\noindent \textbf{(P\_1,3.60.2)} atha vā neḥ iti na eṣā pañcamī . \\
\noindent \textbf{(P\_1,3.60.2)} kā tarhi . \\
\noindent \textbf{(P\_1,3.60.2)} viśeṣaṇaṣaṣṭhī . \\
\noindent \textbf{(P\_1,3.60.2)} neḥ yaḥ viśiḥ . \\
\noindent \textbf{(P\_1,3.60.2)} kaḥ ca neḥ viśiḥ . \\
\noindent \textbf{(P\_1,3.60.2)} viśeṣyaḥ . \\
\noindent \textbf{(P\_1,3.60.2)} vyavahitaḥ ca api śakyate viśeṣayitum . \\
\noindent \textbf{(P\_1,3.60.2)} atha vā niḥ api padam viśiḥ api padam .padavidhiḥ ca samarthāmām . \\
\noindent \textbf{(P\_1,3.60.2)} vyavahite api sāmarthyam bhavati . \\
\noindent \textbf{(P\_1,3.62.1)} kim idam pūrvagrahaṇam sanapekṣam . \\
\noindent \textbf{(P\_1,3.62.1)} prāk sanaḥ yebhyaḥ ātmanepadm uktam tebhyaḥ sanantebhyaḥ api bhavati iti . \\
\noindent \textbf{(P\_1,3.62.1)} āhosvit yogāpekṣam . \\
\noindent \textbf{(P\_1,3.62.1)} prāk etasmāt yogāt yebhyaḥ ātmanepadam uktam tebhyaḥ sanantebhyaḥ ātmanepadam bhavati iti . \\
\noindent \textbf{(P\_1,3.62.1)} kim ca ataḥ . \\
\noindent \textbf{(P\_1,3.62.1)} yadi sanapekṣam nimittam aviśeṣitam bhavati . \\
\noindent \textbf{(P\_1,3.62.1)} pūrvavat sanaḥ na jñāyate kimantāt bhavitavyam . \\
\noindent \textbf{(P\_1,3.62.1)} atha yogāpekṣam uttaratra vidhiḥ na prakalpate . \\
\noindent \textbf{(P\_1,3.62.1)} bubhukṣate upayuyukṣate iti . \\
\noindent \textbf{(P\_1,3.62.1)} yathā icchasi tathā astu . \\
\noindent \textbf{(P\_1,3.62.1)} astu tāvat sanapekṣam . \\
\noindent \textbf{(P\_1,3.62.1)} nanu ca uktam nimittam aviśeṣitam bhavati . \\
\noindent \textbf{(P\_1,3.62.1)} nimittam ca viśeṣitam . \\
\noindent \textbf{(P\_1,3.62.1)} katham . \\
\noindent \textbf{(P\_1,3.62.1)} sanam eva atra nimittatvena apekṣiṣyāmahe . \\
\noindent \textbf{(P\_1,3.62.1)} pūrvavat sanaḥ ātmanepadam bhavati . \\
\noindent \textbf{(P\_1,3.62.1)} kutaḥ sanaḥ iti . \\
\noindent \textbf{(P\_1,3.62.1)} atha vā punaḥ astu yogāpekṣam . \\
\noindent \textbf{(P\_1,3.62.1)} nanu ca uktam uttaratra vidhiḥ na prakalpate . \\
\noindent \textbf{(P\_1,3.62.1)} vidhiḥ ca prakḷptaḥ . \\
\noindent \textbf{(P\_1,3.62.1)} katham . \\
\noindent \textbf{(P\_1,3.62.1)} uttaratra api pūrvavat sanaḥ iti eva anuvartiṣyate . \\
\noindent \textbf{(P\_1,3.62.2)} kimartham punaḥ idam ucyate . \\
\noindent \textbf{(P\_1,3.62.2)} pūrvavat sanaḥ iti śadimriyatyartham . \\
\noindent \textbf{(P\_1,3.62.2)} śadimriyatyarthaḥ ayam ārambhaḥ . \\
\noindent \textbf{(P\_1,3.62.2)} śadimriyatibhyām sanantābhyām ātmanepadam mā bhūt iti . \\
\noindent \textbf{(P\_1,3.62.2)} itarathā hi tābhyām sanantābhyām ātmanepadapratiṣedhaḥ . \\
\noindent \textbf{(P\_1,3.62.2)} itarathā hi anucyamāne asmin tābhyām sanantābhyām ātmanepadasya pratiṣedhaḥ vaktavyaḥ syāt . \\
\noindent \textbf{(P\_1,3.62.2)} śiśitsati mumūrṣati . \\
\noindent \textbf{(P\_1,3.62.2)} katham punaḥ pūrvavat sanaḥ iti anena śadimriyatibhyām sanantābhyām ātmanepadasya pratiṣedhaḥ śakyaḥ vijñātum . \\
\noindent \textbf{(P\_1,3.62.2)} vatinirdeśaḥ ayam kāmacāraḥ ca vatinirdeśe vākyaśeṣam samarthayitum . \\
\noindent \textbf{(P\_1,3.62.2)} tat yathā : uśīnaravat madreṣu yavāḥ . \\
\noindent \textbf{(P\_1,3.62.2)} santi na santi iti . \\
\noindent \textbf{(P\_1,3.62.2)} mātṛvat asyāḥ kalāḥ . \\
\noindent \textbf{(P\_1,3.62.2)} santi na santi . \\
\noindent \textbf{(P\_1,3.62.2)} evam iha api pūrvavat bhavati na bhavati iti . \\
\noindent \textbf{(P\_1,3.62.2)} na bhavati iti vākyaśeṣam samarthayiṣyāmahe . \\
\noindent \textbf{(P\_1,3.62.2)} yathā pūrvayogayoḥ sanantābhyām ātmanepadam na bhavati evam iha api śadimriyatibhyām sanantābhyām ātmanepadam na bhavati iti . \\
\noindent \textbf{(P\_1,3.62.2)} yadi tarhi śadimriyatyarthaḥ ayam ārambhaḥ vidhiḥ na prakalpate . \\
\noindent \textbf{(P\_1,3.62.2)} āsisiṣate śiśayiṣate . \\
\noindent \textbf{(P\_1,3.62.2)} atha vidhyarthaḥ śadimriyatibhyām sanantābhyām ātmanepadam prāpnoti . \\
\noindent \textbf{(P\_1,3.62.2)} yathā icchasi tathā astu . \\
\noindent \textbf{(P\_1,3.62.2)} astu tāvat pratiṣedhārthaḥ . \\
\noindent \textbf{(P\_1,3.62.2)} nanu ca uktam vidhiḥ na prakalpate iti . \\
\noindent \textbf{(P\_1,3.62.2)} vidhiḥ ca prakḷptaḥ . \\
\noindent \textbf{(P\_1,3.62.2)} katham . \\
\noindent \textbf{(P\_1,3.62.2)} etat eva jñāpayati sanantāt ātmanepadam bhavati iti yat ayam śadimriyatibhyām sanantābhyām ātmanepadasya pratiṣedham śāsti . \\
\noindent \textbf{(P\_1,3.62.2)} atha vā punaḥ astu vidhyarthaḥ . \\
\noindent \textbf{(P\_1,3.62.2)} nanu ca uktam śadimriyatibhyām sanantābhyām ātmanepadam prāpnoti iti . \\
\noindent \textbf{(P\_1,3.62.2)} na eṣaḥ doṣaḥ . \\
\noindent \textbf{(P\_1,3.62.2)} prakṛtam sanaḥ na iti anuvartiṣyate . \\
\noindent \textbf{(P\_1,3.62.2)} kva prakṛtam . \\
\noindent \textbf{(P\_1,3.62.2)} jñāśrusmṛdṛśām sanaḥ na anoḥ jñaḥ . \\
\noindent \textbf{(P\_1,3.62.2)} sakarmakāt sanaḥ na . \\
\noindent \textbf{(P\_1,3.62.2)} pratyāṅbhyām śruvaḥ sanaḥ na . \\
\noindent \textbf{(P\_1,3.62.2)} śadeḥ śitaḥ sanaḥ na . \\
\noindent \textbf{(P\_1,3.62.2)} mriyateḥ luṅliṅoḥ ca sanaḥ na iti . \\
\noindent \textbf{(P\_1,3.62.2)} iha idānīm pūrvavat sanaḥ iti sanaḥ iti vartate na iti nivṛttam . \\
\noindent \textbf{(P\_1,3.62.2)} evam ca kṛtvā saḥ api adoṣaḥ bhavati yat uktam nimittam aviśeṣitam bhavati iti . \\
\noindent \textbf{(P\_1,3.62.2)} na eva vā punaḥ atra śadimriyatibhyām sanantābhyām ātmanepadam prāpnoti . \\
\noindent \textbf{(P\_1,3.62.2)} kim kāraṇam . \\
\noindent \textbf{(P\_1,3.62.2)} śadeḥ śitaḥ iti ucyate na ca śadiḥ eva ātmanepadasya nimittam . \\
\noindent \textbf{(P\_1,3.62.2)} kim tarhi . \\
\noindent \textbf{(P\_1,3.62.2)} śit api nimittam . \\
\noindent \textbf{(P\_1,3.62.2)} atha api śadiḥ eva śitparaḥ tu nimittam . \\
\noindent \textbf{(P\_1,3.62.2)} na ca ayam sanparaḥ śitparaḥ bhavati . \\
\noindent \textbf{(P\_1,3.62.2)} yatra tarhi śit na āśrīyate mriyateḥ luṅliṅoḥ ca iti . \\
\noindent \textbf{(P\_1,3.62.2)} atra api na mriyatiḥ eva ātmanepadasya nimittam . \\
\noindent \textbf{(P\_1,3.62.2)} kim tarhi . \\
\noindent \textbf{(P\_1,3.62.2)} luṅliṅau api nimittam . \\
\noindent \textbf{(P\_1,3.62.2)} atha api mriyatiḥ eva luṅliṅparaḥ tu nimittam . \\
\noindent \textbf{(P\_1,3.62.2)} na ca ayam sanparaḥ luṅliṅparaḥ bhavati . \\
\noindent \textbf{(P\_1,3.62.3)} kim punaḥ pūrvasya yat ātmanepadadarśanam tat sanantasya api atidiśyate . \\
\noindent \textbf{(P\_1,3.62.3)} evam bhavitum arhati . \\
\noindent \textbf{(P\_1,3.62.3)} pūrvasya ātmanepadadarśanāt sanantāt ātmanepadabhāvaḥ iti cet gupādiṣu aprasiddhiḥ . \\
\noindent \textbf{(P\_1,3.62.3)} pūrvasya ātmanepadadarśanāt sanantāt ātmanepadabhāvaḥ iti cet gupādiṣu aprasiddhiḥ . \\
\noindent \textbf{(P\_1,3.62.3)} gupādīnām na prāpnoti . \\
\noindent \textbf{(P\_1,3.62.3)} jugupsate mīmāṃsate iti . \\
\noindent \textbf{(P\_1,3.62.3)} na hi etebhyaḥ prāk sanaḥ ātmanepadam na api parasmaipadam paśyāmaḥ . \\
\noindent \textbf{(P\_1,3.62.3)} siddham tu pūrvasya liṅgātideśāt . \\
\noindent \textbf{(P\_1,3.62.3)} siddham etat . \\
\noindent \textbf{(P\_1,3.62.3)} katham . \\
\noindent \textbf{(P\_1,3.62.3)} pūrvasya yat ātmanepadaliṅgam tat sanantasya api atidiśyate . \\
\noindent \textbf{(P\_1,3.62.3)} kṛñādiṣu tu liṅgapratiṣedhaḥ . \\
\noindent \textbf{(P\_1,3.62.3)} kṛñādiṣu tu liṅgapratiṣedhaḥ vaktavyaḥ . \\
\noindent \textbf{(P\_1,3.62.3)} anucikīrṣati parācikīrṣati iti . \\
\noindent \textbf{(P\_1,3.62.3)} astu tarhi prāk sanaḥ yebhyaḥ ātmanepadam dṛṣṭam tebhyaḥ sanantebhyaḥ api bahvati iti . \\
\noindent \textbf{(P\_1,3.62.3)} nanu ca uktam pūrvasya ātmanepadadarśanāt sanantāt ātmanepadabhāvaḥ iti cet gupādiṣu aprasiddhiḥ iti . \\
\noindent \textbf{(P\_1,3.62.3)} na eṣaḥ doṣaḥ . \\
\noindent \textbf{(P\_1,3.62.3)} anubandhakaraṇasāmarthyāt bhaviṣyati . \\
\noindent \textbf{(P\_1,3.62.3)} atha vā avayave kṛtam liṅgam samudāyasya viśeṣakam bhavati . \\
\noindent \textbf{(P\_1,3.62.3)} tat yathā goḥ sakthani karṇe vā kṛtam liṅgam samudāyasya viśeṣakam bhavati . \\
\noindent \textbf{(P\_1,3.62.3)} yadi avayave kṛtam liṅgam samudāyasya viśeṣakam bhavati jugupsayati mīmāṃsāyati iti atra api prāpnoti . \\
\noindent \textbf{(P\_1,3.62.3)} na eṣaḥ doṣaḥ . \\
\noindent \textbf{(P\_1,3.62.3)} avayave kṛtam liṅgam kasya samudāyasya viśeṣakam bhavati . \\
\noindent \textbf{(P\_1,3.62.3)} yam samudāyam yaḥ avayavaḥ na vyabhicarati . \\
\noindent \textbf{(P\_1,3.62.3)} sanam ca na vyabhicarati . \\
\noindent \textbf{(P\_1,3.62.3)} ṇicam punaḥ vyabhicarati . \\
\noindent \textbf{(P\_1,3.62.3)} tat yathā goḥ sakthani karṇe vā kṛtam liṅgam goḥ eva viśeṣakam bhavati na gomaṇḍalasya . \\
\noindent \textbf{(P\_1,3.62.4)} pratyayagrahaṇam ṇijartham . \\
\noindent \textbf{(P\_1,3.62.4)} pratyayasya grahaṇam kartavyam . \\
\noindent \textbf{(P\_1,3.62.4)} pūrvavat pratyayāt iti vaktavyam . \\
\noindent \textbf{(P\_1,3.62.4)} kim prayojanam . \\
\noindent \textbf{(P\_1,3.62.4)} ṇijartham . \\
\noindent \textbf{(P\_1,3.62.4)} ṇijantāt api yathā syāt iti . \\
\noindent \textbf{(P\_1,3.62.4)} ākusmayate vikusmayate hṛṇīyate mahīyate iti. tatra kaḥ doṣaḥ . \\
\noindent \textbf{(P\_1,3.62.4)} tatra hetumaṇṇicaḥ pratiṣedhaḥ . \\
\noindent \textbf{(P\_1,3.62.4)} tatra hetumaṇṇicaḥ pratiṣedhaḥ vaktavyaḥ . \\
\noindent \textbf{(P\_1,3.62.4)} āsayati śāyayati . \\
\noindent \textbf{(P\_1,3.62.4)} sūtram ca bhidyate . \\
\noindent \textbf{(P\_1,3.62.4)} yathānyāsam eva astu . \\
\noindent \textbf{(P\_1,3.62.4)} katham ākusmayate vikusmayate hṛṇīyate mahīyate iti. anubandhakaraṇasāmarthyāt bhaviṣyati . \\
\noindent \textbf{(P\_1,3.62.4)} atha vā avayave kṛtam liṅgam samudāyasya viśeṣakam bhavati . \\
\noindent \textbf{(P\_1,3.62.4)} tat yathā goḥ sakthani karṇe vā kṛtam liṅgam samudāyasya viśeṣakam bhavati . \\
\noindent \textbf{(P\_1,3.62.4)} yadi avayave kṛtam liṅgam samudāyasya viśeṣakam bhavati hṛṇīyayati mahīyayati atra api prāpnoti . \\
\noindent \textbf{(P\_1,3.62.4)} avayave kṛtam liṅgam kasya samudāyasya viśeṣakam bhavati . \\
\noindent \textbf{(P\_1,3.62.4)} yam samudāyam yaḥ avayavaḥ na vyabhicarati . \\
\noindent \textbf{(P\_1,3.62.4)} yakam ca na vyabhicarati . \\
\noindent \textbf{(P\_1,3.62.4)} ṇicam tu vyabhicarati . \\
\noindent \textbf{(P\_1,3.62.4)} tat yathā goḥ sakthani karṇe vā kṛtam liṅgam goḥ eva viśeṣakam bhavati na gomaṇḍalasya . \\
\noindent \textbf{(P\_1,3.63)} kṛñgrahaṇam kimartham . \\
\noindent \textbf{(P\_1,3.63)} iha mā bhūt . \\
\noindent \textbf{(P\_1,3.63)} īhāmāsa īhāmāsatuḥ īhāmāsuḥ . \\
\noindent \textbf{(P\_1,3.63)} katham ca atra asteḥ anuprayogaḥ bhavati . \\
\noindent \textbf{(P\_1,3.63)} pratyāhāragrahaṇam tatra vijñāyate . \\
\noindent \textbf{(P\_1,3.63)} katham punaḥ jñāyate tatra pratyāhāragrahaṇam iti . \\
\noindent \textbf{(P\_1,3.63)} iha kṛñgrahaṇāt . \\
\noindent \textbf{(P\_1,3.63)} iha kasmāt pratyāhāragrahaṇam na bhavati . \\
\noindent \textbf{(P\_1,3.63)} iha eva kṛñgrahaṇāt . \\
\noindent \textbf{(P\_1,3.63)} atha iha kasmāt na bhavati . \\
\noindent \textbf{(P\_1,3.63)} udumbhām cakāra udubjām cakāra . \\
\noindent \textbf{(P\_1,3.63)} nanu ca āmpratyayavat iti ucyate na ca atra āmpratyayāt ātmanepadam paśyāmaḥ . \\
\noindent \textbf{(P\_1,3.63)} na brūmaḥ anena iti . \\
\noindent \textbf{(P\_1,3.63)} kim tarhi . \\
\noindent \textbf{(P\_1,3.63)} svaritañitaḥ kartrabhiprāye kriyāphale ātmanepadam bhavati iti . \\
\noindent \textbf{(P\_1,3.63)} na eṣaḥ doṣaḥ . \\
\noindent \textbf{(P\_1,3.63)} iha niyamārtham bhaviṣyati . \\
\noindent \textbf{(P\_1,3.63)} āmpratyayavat eva iti . \\
\noindent \textbf{(P\_1,3.63)} yadi niyamārtham vidhiḥ na prakalpate . \\
\noindent \textbf{(P\_1,3.63)} īhām cakre ūhām cakre iti . \\
\noindent \textbf{(P\_1,3.63)} vidhiḥ ca prakḷptaḥ . \\
\noindent \textbf{(P\_1,3.63)} katham pūrvavat iti vartate . \\
\noindent \textbf{(P\_1,3.63)} āmpratyayavat pūrvavat ca iti . \\
\noindent \textbf{(P\_1,3.64)} svarādyupasṛṣṭāt iti vaktavyam . \\
\noindent \textbf{(P\_1,3.64)} udyuṅkte anuyuṅkte . \\
\noindent \textbf{(P\_1,3.64)} aparaḥ āha : svarādyantopasṛṣṭāt iti vaktayam . \\
\noindent \textbf{(P\_1,3.64)} prayuṅkte niyuṅkte niniyuṅkte . \\
\noindent \textbf{(P\_1,3.65)} kimartham videśasthasya grahaṇam kriyate na samaḥ gamādiṣu eva ucyeta . \\
\noindent \textbf{(P\_1,3.65)} samaḥ kṣṇuvaḥ sakarmakārtham . \\
\noindent \textbf{(P\_1,3.65)} sakarmakārthaḥ ayam ārambhaḥ . \\
\noindent \textbf{(P\_1,3.65)} akarmakāt iti hi tatra anuvartate \\
\noindent \textbf{(P\_1,3.66)} anavanakauṭilyayoḥ iti vaktavyam . \\
\noindent \textbf{(P\_1,3.66)} iha api yathā syāt . \\
\noindent \textbf{(P\_1,3.66)} prabhujati vāsasī nibhujati jānuśirasī iti . \\
\noindent \textbf{(P\_1,3.66)} tat tarhi vaktavyam . \\
\noindent \textbf{(P\_1,3.66)} na vaktavyam . \\
\noindent \textbf{(P\_1,3.66)} yasya bhujeḥ avanam anavanam ca arthaḥ tasya grahaṇam . \\
\noindent \textbf{(P\_1,3.66)} na ca asya bhujeḥ avanam anavanam ca arthaḥ \\
\noindent \textbf{(P\_1,3.67.1)} ṇeḥ ātmanepadavidhāne aṇyantasya karmaṇaḥ tatra upalabdhiḥ ṇeḥ ātmanepadavidhāne aṇyantasya yat karma yadā ṇyante tat eva karma bhavati tadā ātmanepadam bhavati iti vaktavyam . \\
\noindent \textbf{(P\_1,3.67.1)} itarathā hi sarvaprasaṅgaḥ . \\
\noindent \textbf{(P\_1,3.67.1)} itarathā hi sarvatra prasaṅgaḥ syāt . \\
\noindent \textbf{(P\_1,3.67.1)} iha api prasajyeta : ārohanti hastinam hastipakāḥ . \\
\noindent \textbf{(P\_1,3.67.1)} ārohamāṇaḥ hastīsthalam ārohayati manuṣyān . \\
\noindent \textbf{(P\_1,3.67.1)} tat tarhi vaktavyam . \\
\noindent \textbf{(P\_1,3.67.1)} na vaktavyam . \\
\noindent \textbf{(P\_1,3.67.1)} kasmāt na bhavati : ārohanti hastinam hastipakāḥ . \\
\noindent \textbf{(P\_1,3.67.1)} ārohamāṇaḥ hastīsthalam ārohayati manuṣyān iti . \\
\noindent \textbf{(P\_1,3.67.1)} evam vakṣyāmi . \\
\noindent \textbf{(P\_1,3.67.1)} ṇeḥ ātmanepadam bhavati . \\
\noindent \textbf{(P\_1,3.67.1)} tataḥ aṇau yat karma ṇau cet . \\
\noindent \textbf{(P\_1,3.67.1)} aṇyante yat karma ṇau yadi tat eva karma bhavati . \\
\noindent \textbf{(P\_1,3.67.1)} tataḥ saḥ kartā . \\
\noindent \textbf{(P\_1,3.67.1)} kartā cet saḥ bhavati ṇau iti . \\
\noindent \textbf{(P\_1,3.67.1)} yadi evam karmakāryam bhavati . \\
\noindent \textbf{(P\_1,3.67.1)} tatra karmakartṛtvāt siddham . \\
\noindent \textbf{(P\_1,3.67.1)} karmakartṛtvāt siddham iti cet yakciṇoḥ nivṛttyartham vacanam . \\
\noindent \textbf{(P\_1,3.67.1)} karmakartṛtvāt siddham iti cet yakciṇoḥ nivṛttyartham idam vaktavyam . \\
\noindent \textbf{(P\_1,3.67.1)} karmāpadiṣṭau yakciṇau mā bhūtām iti . \\
\noindent \textbf{(P\_1,3.67.1)} na vā yakciṇoḥ pratiṣedhāt . \\
\noindent \textbf{(P\_1,3.67.1)} na vā eṣaḥ doṣaḥ . \\
\noindent \textbf{(P\_1,3.67.1)} kim kāraṇam . \\
\noindent \textbf{(P\_1,3.67.1)} yakciṇoḥ pratiṣedhāt . \\
\noindent \textbf{(P\_1,3.67.1)} pratiṣidhyete atra yakciṇau . \\
\noindent \textbf{(P\_1,3.67.1)} yakciṇoḥ pratiṣedhe hetumaṇṇiśribrūñām upasaṅkhyānam iti . \\
\noindent \textbf{(P\_1,3.67.1)} yaḥ tarhi na hetumaṇṇic tadartham idam vaktavyam . \\
\noindent \textbf{(P\_1,3.67.1)} tasya karmāpadiṣṭau yakciṇau mā bhūtām iti : utpucchayate puccham svayam eva . \\
\noindent \textbf{(P\_1,3.67.1)} udapuppucchata puccham svayam eva . \\
\noindent \textbf{(P\_1,3.67.1)} atra api yathā bhāradvājīyāḥ paṭhanti tathā bhavitavyam pratiṣedhena : yakciṇoḥ pratiṣedhe ṇiśrigranthibrūñātmanepadākarmakāṇām upsaṅkhyānam iti . \\
\noindent \textbf{(P\_1,3.67.1)} saḥ ca avaśyam pratiṣedhaḥ āśrayitavyaḥ . \\
\noindent \textbf{(P\_1,3.67.1)} itarathā hi yatra niyamaḥ tataḥ anyatra pratiṣedhaḥ . \\
\noindent \textbf{(P\_1,3.67.1)} anucyamāne hi etasmin yatra niyamaḥ tataḥ anyatra tena yakciṇoḥ pratiṣedhaḥ vaktavyaḥ syāt . \\
\noindent \textbf{(P\_1,3.67.1)} gaṇayati gaṇam gopālakaḥ . \\
\noindent \textbf{(P\_1,3.67.1)} gaṇayati gaṇaḥ svayam eva . \\
\noindent \textbf{(P\_1,3.67.1)} ātmanepadasya ca . \\
\noindent \textbf{(P\_1,3.67.1)} ātmanepadasya ca pratiṣedhaḥ vaktavyaḥ : gaṇayati gaṇaḥ svayam eva . \\
\noindent \textbf{(P\_1,3.67.1)} ātmanepadapratiṣedhārtham tu . \\
\noindent \textbf{(P\_1,3.67.1)} ātmanepadapratiṣedhārtham idam vaktavyam . \\
\noindent \textbf{(P\_1,3.67.1)} gaṇayati gaṇaḥ svayam eva . \\
\noindent \textbf{(P\_1,3.67.1)} iṣyate eva atra ātmanepadam . \\
\noindent \textbf{(P\_1,3.67.1)} kim iṣyate eva āhosvit prāpnoti api . \\
\noindent \textbf{(P\_1,3.67.1)} iṣyate ca prāpnoti ca . \\
\noindent \textbf{(P\_1,3.67.1)} katham . \\
\noindent \textbf{(P\_1,3.67.1)} aṇau it kasya idam ṇeḥ grahaṇam . \\
\noindent \textbf{(P\_1,3.67.1)} yamāt ṇeḥ prāk karma kartā vā vidyate . \\
\noindent \textbf{(P\_1,3.67.1)} na ca etasmāt ṇeḥ prāk karma kartā vā vidyate . \\
\noindent \textbf{(P\_1,3.67.1)} idam tarhi prayojanam : anādhyāne iti vakṣyāmi iti . \\
\noindent \textbf{(P\_1,3.67.1)} iha mā bhūt . \\
\noindent \textbf{(P\_1,3.67.1)} smarati vanagulmasya kokilaḥ . \\
\noindent \textbf{(P\_1,3.67.1)} smarayati enam vanagulmaḥ svayam eva . \\
\noindent \textbf{(P\_1,3.67.1)} etat api na asti prayojanam . \\
\noindent \textbf{(P\_1,3.67.1)} karmāpadiṣṭāḥ vidhayaḥ karmasthabhāvakānām karmasthakriyāṇām bhavanti kartṛsthabhāvakaḥ ca ayam . \\
\noindent \textbf{(P\_1,3.67.1)} evam tarhi siddhe sati yat anādhyāne iti pratiṣedham śāsti tat jñāpayati ācāryaḥ bhavati evañjātīyakānām ātmanepadam iti . \\
\noindent \textbf{(P\_1,3.67.1)} kim etasya jñāpane prayojanam . \\
\noindent \textbf{(P\_1,3.67.1)} paśyanti bhṛtyāḥ rājānam . \\
\noindent \textbf{(P\_1,3.67.1)} darśayate bhṛtyān rājā . \\
\noindent \textbf{(P\_1,3.67.1)} darśayate bhṛtyaiḥ rājā . \\
\noindent \textbf{(P\_1,3.67.1)} atra ātmanepadam siddham bhavati . \\
\noindent \textbf{(P\_1,3.67.2)} ātmanaḥ karmatve pratiṣedhaḥ . \\
\noindent \textbf{(P\_1,3.67.2)} ātmanaḥ karmatve pratiṣedhaḥ vaktavyaḥ . \\
\noindent \textbf{(P\_1,3.67.2)} hanti ātmānam . \\
\noindent \textbf{(P\_1,3.67.2)} ghātayati ātmā iti . \\
\noindent \textbf{(P\_1,3.67.2)} saḥ tarhi vaktavyaḥ . \\
\noindent \textbf{(P\_1,3.67.2)} na vā ṇyante anyasya kartṛtvāt . \\
\noindent \textbf{(P\_1,3.67.2)} na vā vaktavyaḥ . \\
\noindent \textbf{(P\_1,3.67.2)} kim kāraṇam . \\
\noindent \textbf{(P\_1,3.67.2)} ṇyante anyasya kartṛtvāt . \\
\noindent \textbf{(P\_1,3.67.2)} anyat atra aṇyante karma anyaḥ ṇyantasya kartā . \\
\noindent \textbf{(P\_1,3.67.2)} katham . \\
\noindent \textbf{(P\_1,3.67.2)} dvau ātmanau antarātmā śarīrātmā ca . \\
\noindent \textbf{(P\_1,3.67.2)} antarātmā tat karma karoti yena śarīrātmā sukhaduḥkhe anubhavati . \\
\noindent \textbf{(P\_1,3.67.2)} śarīrātmā tat karma karoti yena antarātmā sukhaduḥkhe anubhavati iti . \\
\noindent \textbf{(P\_1,3.72)} svaritañitaḥ iti kimartham . \\
\noindent \textbf{(P\_1,3.72)} yāti vāti drāti psāti . \\
\noindent \textbf{(P\_1,3.72)} svaritañitaḥ iti śakyam akartum . \\
\noindent \textbf{(P\_1,3.72)} kasmāt na bhavati yāti vāti drāti psāti iti . \\
\noindent \textbf{(P\_1,3.72)} kartrabhiprāye kriyāphale iti ucyate sarveṣām ca kartrabhiprāyam kriyāphalam asti . \\
\noindent \textbf{(P\_1,3.72)} te evam vijñāsyāmaḥ . \\
\noindent \textbf{(P\_1,3.72)} yeṣām kartrabhiprāyam akartrabhiprayam ca kriyāphalam asti tebhyaḥ ātmanepadam bhavati iti . \\
\noindent \textbf{(P\_1,3.72)} na ca eteṣām kartrabhiprāyam akartrabhiprayam ca kriyāphalam asti . \\
\noindent \textbf{(P\_1,3.72)} tathājātīyakāḥ khalu ācāryeṇa svaritañitaḥ paṭhitāḥ yeṣām kartrabhiprāyam akartrabhiprayam ca kriyāphalam asti . \\
\noindent \textbf{(P\_1,3.72)} atha abhiprāyagrahaṇam kimartham . \\
\noindent \textbf{(P\_1,3.72)} svaritañitaḥ kartrāye kriyāphale iti iyati ucyamāne yam eva samprati eti kriyāphalam tatra eva syāt . \\
\noindent \textbf{(P\_1,3.72)} lūñ lunīte pūñ punīte . \\
\noindent \textbf{(P\_1,3.72)} iha na syāt . \\
\noindent \textbf{(P\_1,3.72)} yaj yajate vap vapate . \\
\noindent \textbf{(P\_1,3.72)} abhiprayagrahaṇe punaḥ kriyamāṇe na doṣaḥ bhavati . \\
\noindent \textbf{(P\_1,3.72)} abhiḥ ābhimukhye vartate pra ādikarmaṇi . \\
\noindent \textbf{(P\_1,3.72)} tena yam ca abhipraiti yam ca abhipraiṣyati yam ca abhiprāgāt tatra sarvatra ābhimukhyamātre siddham bhavati . \\
\noindent \textbf{(P\_1,3.72)} kartrabhiprāye kriyāphale iti kimartham . \\
\noindent \textbf{(P\_1,3.72)} pacanti bhaktakārāḥ . \\
\noindent \textbf{(P\_1,3.72)} kurvanti karmakārāḥ . \\
\noindent \textbf{(P\_1,3.72)} yajanti yājakāḥ . \\
\noindent \textbf{(P\_1,3.72)} kartrabhiprāye kriyāphale iti ucyamāne api atra prāpnoti . \\
\noindent \textbf{(P\_1,3.72)} atra api hi kriyāphalam kartāram abhipraiti . \\
\noindent \textbf{(P\_1,3.72)} yājakāḥ yajanti gāḥ lapsyāmahe iti . \\
\noindent \textbf{(P\_1,3.72)} karmakarāḥ kurvanti pādikam ahaḥ lapsyāmahe iti . \\
\noindent \textbf{(P\_1,3.72)} evam tarhi kartrabhipraye kriyāphale iti ucyate sarvatra ca kartāram kriyāphalam abhipraiti . \\
\noindent \textbf{(P\_1,3.72)} tatra prakarṣagatiḥ vijñāsyate . \\
\noindent \textbf{(P\_1,3.72)} sādhīyaḥ yatra kartāram kriyāphalam abhipraiti iti . \\
\noindent \textbf{(P\_1,3.72)} na ca antareṇa yajim yajiphalam vapim va vapiphalam labhante . \\
\noindent \textbf{(P\_1,3.72)} yājakāḥ punaḥ antareṇa api yajim gāḥ labhante bhṛtakāḥ ca pādikam iti . \\
\noindent \textbf{(P\_1,3.78)} śeṣavacanam pañcamyā cet arthe pratiṣedhaḥ . \\
\noindent \textbf{(P\_1,3.78)} śeṣavacanam pañcamyā cet arthe pratiṣedhaḥ vaktavyaḥ . \\
\noindent \textbf{(P\_1,3.78)} bhidyate kuśūlaḥ svayam eva . \\
\noindent \textbf{(P\_1,3.78)} chidyate rajjuḥ svayam eva . \\
\noindent \textbf{(P\_1,3.78)} evam tarhi śeṣe iti vakṣyāmi . \\
\noindent \textbf{(P\_1,3.78)} saptamyā cet prakṛteḥ . saptamyā cet prakṛteḥ pratiṣedhaḥ vaktavyaḥ . \\
\noindent \textbf{(P\_1,3.78)} āste śete cyavante plavante . \\
\noindent \textbf{(P\_1,3.78)} siddham tu ubhayanirdeśāt . \\
\noindent \textbf{(P\_1,3.78)} siddham etat . \\
\noindent \textbf{(P\_1,3.78)} katham . \\
\noindent \textbf{(P\_1,3.78)} ubhayanirdeśaḥ kartavyaḥ . \\
\noindent \textbf{(P\_1,3.78)} śeṣāt śeṣe iti vaktavyam . \\
\noindent \textbf{(P\_1,3.78)} kartṛgrahaṇam idānīm kimartham syāt . \\
\noindent \textbf{(P\_1,3.78)} kartṛgrahaṇam anuparādyartham . \\
\noindent \textbf{(P\_1,3.78)} anuparādyartham etat syāt . \\
\noindent \textbf{(P\_1,3.78)} iha mā bhūt . \\
\noindent \textbf{(P\_1,3.78)} anukriyate svayam eva . \\
\noindent \textbf{(P\_1,3.78)} parākriyate svayam eva iti . \\
\noindent \textbf{(P\_1,3.78)} sidhyati . \\
\noindent \textbf{(P\_1,3.78)} sutram tarhi bhidyate . \\
\noindent \textbf{(P\_1,3.78)} yathānyāsam eva astu . \\
\noindent \textbf{(P\_1,3.78)} nanu ca uktam śeṣavacanam pañcamyā cet arthe pratiṣedhaḥ iti . \\
\noindent \textbf{(P\_1,3.78)} na eṣaḥ doṣaḥ . \\
\noindent \textbf{(P\_1,3.78)} kartari karmavyatihāre iti atra kartṛgrahaṇam pratyākhyāyate . \\
\noindent \textbf{(P\_1,3.78)} tat prakṛtam iha anuvartiṣyate . \\
\noindent \textbf{(P\_1,3.78)} śeṣāt kartari kartari iti . \\
\noindent \textbf{(P\_1,3.78)} kim idam kartari kartari iti . \\
\noindent \textbf{(P\_1,3.78)} kartā eva yaḥ kartā tatra yathā syāt . \\
\noindent \textbf{(P\_1,3.78)} kartā ca anyaḥ ca yaḥ kartā tatra mā bhūt iti . \\
\noindent \textbf{(P\_1,3.79)} kimartham idam ucyate . \\
\noindent \textbf{(P\_1,3.79)} parasmaipadapratiṣedhāt kṛñādiṣu vidhānam . \\
\noindent \textbf{(P\_1,3.79)} parasmaipadapratiṣedhāt kṛñādiṣu parasmaipadam vidhīyate . \\
\noindent \textbf{(P\_1,3.79)} pratiṣidhyate tatra parasmaipadam svaritañitaḥ kartrabhiprāye kriyāphale ātmanepadam bhavati iti . \\
\noindent \textbf{(P\_1,3.79)} asti prayojanam etat . \\
\noindent \textbf{(P\_1,3.79)} kim tarhi iti . \\
\noindent \textbf{(P\_1,3.79)} tatra ātmanepadapratiṣedhaḥ apratiṣiddhatvāt . \\
\noindent \textbf{(P\_1,3.79)} tatra ātmanepadasya pratiṣedhaḥ vaktavyaḥ . \\
\noindent \textbf{(P\_1,3.79)} kim kāraṇam . \\
\noindent \textbf{(P\_1,3.79)} apratiṣiddhatvāt . \\
\noindent \textbf{(P\_1,3.79)} na hi ātmanepadam pratiṣidhyate . \\
\noindent \textbf{(P\_1,3.79)} kim tarhi . \\
\noindent \textbf{(P\_1,3.79)} parasmaipadam anena vidhīyate . \\
\noindent \textbf{(P\_1,3.79)} na vā dyutādibhyaḥ vāvacanāt . \\
\noindent \textbf{(P\_1,3.79)} na vā eṣaḥ doṣaḥ . \\
\noindent \textbf{(P\_1,3.79)} kim kāraṇam . \\
\noindent \textbf{(P\_1,3.79)} dyutādibhyaḥ vāvacanāt . \\
\noindent \textbf{(P\_1,3.79)} yat ayam dyutādibhyaḥ vāvacanam karoti tat jñāpayati ācāryaḥ na parasmaipadaviṣaye ātmanepadam bhavati iti . \\
\noindent \textbf{(P\_1,3.79)} ātmanepadaniyame vā pratiṣedhaḥ . \\
\noindent \textbf{(P\_1,3.79)} ātmanepadaniyame vā pratiṣedhaḥ vaktavyaḥ . \\
\noindent \textbf{(P\_1,3.79)} svaritañitaḥ kartrabhiprāye kriyāphale ātmanepadam bhavati kartari . \\
\noindent \textbf{(P\_1,3.79)} anuparābhyām kṛñaḥ na iti . \\
\noindent \textbf{(P\_1,3.79)} sidhyati . \\
\noindent \textbf{(P\_1,3.79)} sūtram tarhi bhidyate . \\
\noindent \textbf{(P\_1,3.79)} yathānyāsam eva astu . \\
\noindent \textbf{(P\_1,3.79)} nanu ca uktam tatra ātmanepadapratiṣedhaḥ apratiṣiddhatvāt iti . \\
\noindent \textbf{(P\_1,3.79)} parihṛtam etat na vā dyutādibhyaḥ vāvacanāt . \\
\noindent \textbf{(P\_1,3.79)} atha vā idam tāvat ayam praṣṭavyaḥ . \\
\noindent \textbf{(P\_1,3.79)} svaritañitaḥ kartrabhiprāye kriyāphale ātmanepadam bhavati iti parasmaipadam kasmāt na bhavati . \\
\noindent \textbf{(P\_1,3.79)} ātmanepadena bādhyate . \\
\noindent \textbf{(P\_1,3.79)} yathā eva tarhi ātmanepadena parasmaipadam bādhyate evam parasmaipadena ātmanepadam bādhiṣyate . \\
\noindent \textbf{(P\_1,3.86)} budhādiṣu ye akarmakāḥ teṣām grahaṇam kimartham . \\
\noindent \textbf{(P\_1,3.86)} sakarmakārtham acittavatkartṛkārtham vā . \\
\noindent \textbf{(P\_1,3.88)} aṇau akarmakāt iti curādiṇicaḥ ṇyantāt parasmaipadavacanam . \\
\noindent \textbf{(P\_1,3.88)} aṇau akarmakāt iti curādiṇicaḥ ṇyantāt parasmaipadam vaktavyam . \\
\noindent \textbf{(P\_1,3.88)} iha api yathā syāt : cetayamāṇam prayojayati cetayati iti . \\
\noindent \textbf{(P\_1,3.88)} yadi tarhi atra api iṣyate aṇigrahaṇam idānīm kimartham syāt . \\
\noindent \textbf{(P\_1,3.88)} akarmakagrahaṇam aṇyantaviśeṣaṇam yathā vijñāyeta . \\
\noindent \textbf{(P\_1,3.88)} atha akriyamāṇe aṇigrahaṇam kasya akarmakgrahaṇam viśeṣaṇam syāt . \\
\noindent \textbf{(P\_1,3.88)} ṇeḥ iti vartate . \\
\noindent \textbf{(P\_1,3.88)} ṇyantaviśeṣaṇam . \\
\noindent \textbf{(P\_1,3.88)} tatra kaḥ doṣaḥ . \\
\noindent \textbf{(P\_1,3.88)} iha eva syāt : cetayamānam prayojayati cetayati iti . \\
\noindent \textbf{(P\_1,3.88)} iha na syāt : āsayati śāyayati iti . \\
\noindent \textbf{(P\_1,3.88)} siddham tu atasmin ṇau iti vacanāt . \\
\noindent \textbf{(P\_1,3.88)} siddham etat . \\
\noindent \textbf{(P\_1,3.88)} katham . \\
\noindent \textbf{(P\_1,3.88)} atasmin ṇau yaḥ akarmakaḥ tatra iti vaktavyam . \\
\noindent \textbf{(P\_1,3.88)} sidhyati . \\
\noindent \textbf{(P\_1,3.88)} sūtram tarhi bhidyate . \\
\noindent \textbf{(P\_1,3.88)} yathānyāsam eva astu . \\
\noindent \textbf{(P\_1,3.88)} nanu ca uktam aṇav akarmakāt iti curādiṇicaḥ ṇyantāt parasmaipadavacanam iti . \\
\noindent \textbf{(P\_1,3.88)} na eṣaḥ doṣaḥ . \\
\noindent \textbf{(P\_1,3.88)} aṇau iti kasya idam ṇeḥ grahaṇam . \\
\noindent \textbf{(P\_1,3.88)} yasmāṇ ṇeḥ prāk karma kartā vā vidyate . \\
\noindent \textbf{(P\_1,3.88)} na ca etasmāṇ ṇeaḥ prāk karma kartā vā vidyate . \\
\noindent \textbf{(P\_1,3.89)} pādiṣu dheṭaḥ upasaṅkhyānam . \\
\noindent \textbf{(P\_1,3.89)} pādiṣu dheṭaḥ upasaṅkhyānam kartavyam . \\
\noindent \textbf{(P\_1,3.89)} dhāpayate śiśumeka samīcī . \\
\noindent \textbf{(P\_1,3.93)} kimarthaḥ cakāraḥ . \\
\noindent \textbf{(P\_1,3.93)} syasanoḥ iti etat anukṛṣyate . \\
\noindent \textbf{(P\_1,3.93)} yadi tarhi na antareṇa cakāram anuṅrttiḥ bhavati dyudbhyaḥ luṅi iti atra api cakāraḥ kartavyaḥ vibhāṣā iti anukarṣaṇārthaḥ . \\
\noindent \textbf{(P\_1,3.93)} atha idānīm antareṇa api cakāram anuvṛttiḥ bhavati iha api na arthaḥ cakāreṇa . \\
\noindent \textbf{(P\_1,3.93)} evam sarve cakārāḥ pratyākhyāyante. \\
\noindent \textbf{(P\_1,4.1.1)} kimartham idam ucyate . \\
\noindent \textbf{(P\_1,4.1.1)} anyatra sañjñāsamāveśānniyamārtham vacanam . \\
\noindent \textbf{(P\_1,4.1.1)} anyatra sañjñāsamāveśaḥ bhavati . \\
\noindent \textbf{(P\_1,4.1.1)} kvānyatra . \\
\noindent \textbf{(P\_1,4.1.1)} loke vyākaraṇe ca . \\
\noindent \textbf{(P\_1,4.1.1)} loke tāvat . \\
\noindent \textbf{(P\_1,4.1.1)} indraḥ śakraḥ puruhūtaḥ purandaraḥ . \\
\noindent \textbf{(P\_1,4.1.1)} kanduḥ koṣṭhaḥ kuśūlaḥ iti . \\
\noindent \textbf{(P\_1,4.1.1)} ekasya dravyasya bahvyaḥ sañjñāḥ bhavanti . \\
\noindent \textbf{(P\_1,4.1.1)} vyākaraṇe api kartavyam hartavyam iti atra pratyayakṛtkṛtyasañjñānām samāveśaḥ bhavati . \\
\noindent \textbf{(P\_1,4.1.1)} pāñcālaḥ vaidehaḥ vaidarbhaḥ iti atra pratyayataddhitatadrājasañjñānāṃ samāveśaḥ bhavati . \\
\noindent \textbf{(P\_1,4.1.1)} anyatra sañjñāsamāveśāt etasmāt kāraṇāt ā kaḍārāt api sañjñānām samāveśaḥ prāpnoti . \\
\noindent \textbf{(P\_1,4.1.1)} iṣyate ca ekā eva sañjñā syāt iti . \\
\noindent \textbf{(P\_1,4.1.1)} tat ca antareṇa yatnam na sidhyati iti niyamārtham vacanam . \\
\noindent \textbf{(P\_1,4.1.1)} evamartham idam ucyate . \\
\noindent \textbf{(P\_1,4.1.1)} asti prayojanam etat . \\
\noindent \textbf{(P\_1,4.1.1)} kim tarhi iti . \\
\noindent \textbf{(P\_1,4.1.2)} katham tvetatsūtram paṭhitavyam . \\
\noindent \textbf{(P\_1,4.1.2)} kim ā kaḍārāt ekā sañjñā iti āhosvit prāk kaḍārāt param kāryam iti . \\
\noindent \textbf{(P\_1,4.1.2)} kutaḥ punaḥ ayam sandehaḥ . \\
\noindent \textbf{(P\_1,4.1.2)} ubhayathā hi ācāryeṇa śiṣyāḥ sūtram pratipāditāḥ : kecit ā kaḍārāt ekā sañjñā iti , kecit prāk kaḍārāt param kāryam iti . \\
\noindent \textbf{(P\_1,4.1.2)} kaḥ ca atra viśeṣaḥ . \\
\noindent \textbf{(P\_1,4.1.2)} tatra ekasañjñādhikāre tadvacanam . tatra ekasañjñādhikāre tat vaktavyam . \\
\noindent \textbf{(P\_1,4.1.2)} kim . \\
\noindent \textbf{(P\_1,4.1.2)} ekā sañjñā bhavati iti . \\
\noindent \textbf{(P\_1,4.1.2)} nanu ca yasya api paraṅkāryatvam tena api paragrahaṇam kartavyam . \\
\noindent \textbf{(P\_1,4.1.2)} parārtham mama bhaviṣyati . \\
\noindent \textbf{(P\_1,4.1.2)} vipratiṣedhe ca iti . \\
\noindent \textbf{(P\_1,4.1.2)} mama api tarhi ekagrahaṇam parārtham bhaviṣyati . \\
\noindent \textbf{(P\_1,4.1.2)} sarūpāṇām ekaśeṣaḥ ekavibhaktau iti . \\
\noindent \textbf{(P\_1,4.1.2)} sañjñādhikāraḥ ca ayam . \\
\noindent \textbf{(P\_1,4.1.2)} tatra kim anyat śakyam vijñātum anyat ataḥ sañjñāyāḥ . \\
\noindent \textbf{(P\_1,4.1.2)} tatra etāvat vācyam . \\
\noindent \textbf{(P\_1,4.1.2)} ā kaḍārāt ekā . \\
\noindent \textbf{(P\_1,4.1.2)} kim . \\
\noindent \textbf{(P\_1,4.1.2)} ekā sañjñā bhavati iti . \\
\noindent \textbf{(P\_1,4.1.2)} aṅgasañjñayā bhapadasañjñayoḥ asamāveśaḥ . \\
\noindent \textbf{(P\_1,4.1.2)} aṇgasañjñayā bhapadasañjñayoḥ samāveśaḥ na prāpnoti . \\
\noindent \textbf{(P\_1,4.1.2)} sārpiṣkaḥ bārhiṣkaḥ yājuṣkaḥ dhānuṣkaḥ . \\
\noindent \textbf{(P\_1,4.1.2)} bābhravyaḥ māṇḍavya iti . \\
\noindent \textbf{(P\_1,4.1.2)} anavakāśe bhapadasañjñe aṅgasañjñāṃ bādheyātām . \\
\noindent \textbf{(P\_1,4.1.2)} paravacane hi niyamānupapatteḥ ubhayasañjñābhāvaḥ . \\
\noindent \textbf{(P\_1,4.1.2)} yasya punaḥ paraṅkāryatvam niyamānupapatteḥ tasya ubhayoḥ sañjñayoḥ bhāvaḥ siddhaḥ . \\
\noindent \textbf{(P\_1,4.1.2)} katham . \\
\noindent \textbf{(P\_1,4.1.2)} pūrve tasya bhapadasañjñe parā aṅgasañjñā . \\
\noindent \textbf{(P\_1,4.1.2)} katham . \\
\noindent \textbf{(P\_1,4.1.2)} evam sa vakṣyati . \\
\noindent \textbf{(P\_1,4.1.2)} yasmātpratyayavidhiḥ tadādi suptiṅantaṃ padam naḥ kye siti ca . \\
\noindent \textbf{(P\_1,4.1.2)} svādiṣu asarvanāmasthāne yaci bham . \\
\noindent \textbf{(P\_1,4.1.2)} tasya ante pratyaye aṅgamiti . \\
\noindent \textbf{(P\_1,4.1.2)} tatra ārambhasāmarthyāc ca bhapadasañjñe paraṅkāryatvāt ca aṅgasañjñā bhaviṣyati . \\
\noindent \textbf{(P\_1,4.1.2)} nanu ca yasya api ekasañjñādhikāraḥ tasya api aṅgasañjñāpūrvike bhapadasañjñe . \\
\noindent \textbf{(P\_1,4.1.2)} katham .anuvṛttiḥ kriyate . \\
\noindent \textbf{(P\_1,4.1.2)} paryāyaḥ prasajyeta. ekā sañjñā iti vacanāt na asti yaugapadyena saṃbhavaḥ . \\
\noindent \textbf{(P\_1,4.1.2)} karmadhārayatve tatpuruṣagrahaṇam . karmadhārayatve tatpuruṣagrahaṇam kartavyam . \\
\noindent \textbf{(P\_1,4.1.2)} tatpuruṣaḥ samānādhikaraṇaḥ karmadhārayaḥ iti . \\
\noindent \textbf{(P\_1,4.1.2)} ekasañjñādhikāraḥ iti coditam . \\
\noindent \textbf{(P\_1,4.1.2)} akriyamāṇe hi anavakāśā karmadhārayasañjñā tatpuruṣasañjñām bādheta . \\
\noindent \textbf{(P\_1,4.1.2)} paravacane hi niyamānupapatteḥ ubhayasañjñābhāvaḥ . \\
\noindent \textbf{(P\_1,4.1.2)} yasya punaḥ paraṅkāryatvam niyamānupapatteḥ tasya ubhayoḥ sañjñayoḥ bhāvaḥ siddhaḥ . \\
\noindent \textbf{(P\_1,4.1.2)} katham . \\
\noindent \textbf{(P\_1,4.1.2)} pūrvā tasya karmadhārayasañjñā parā tatpuruṣasañjñā . \\
\noindent \textbf{(P\_1,4.1.2)} katham . \\
\noindent \textbf{(P\_1,4.1.2)} evam sa vakṣyati . \\
\noindent \textbf{(P\_1,4.1.2)} pūrvakālaikasarvajaratpurāṇanavakevalāḥ samānādhikaraṇena karmadhārayaḥ iti . \\
\noindent \textbf{(P\_1,4.1.2)} evam sarvam karmadhārayaprakaraṇam anukramya tasya ante śritādiḥ tatpuruṣaḥ iti . \\
\noindent \textbf{(P\_1,4.1.2)} tatra ārambhasāmarthyāt ca karmadhārayasañjñā paraṅkāryatvāt ca tatpuruṣasañjñā bhaviṣyati . \\
\noindent \textbf{(P\_1,4.1.2)} nanu ca yasya api ekasañjñādhikāraḥ tasya api tatpuruṣasañjñāpūrvikā karmadhārayasañjñā . \\
\noindent \textbf{(P\_1,4.1.2)} katham . \\
\noindent \textbf{(P\_1,4.1.2)} anuvṛttiḥ kriyate . \\
\noindent \textbf{(P\_1,4.1.2)} paryāyaḥ prasajyeta . \\
\noindent \textbf{(P\_1,4.1.2)} ekā sañjñā iti vacanāt na asti yaugapadyena sambhavaḥ . \\
\noindent \textbf{(P\_1,4.1.2)} tatpuruṣatve dvigucagrahaṇam . \\
\noindent \textbf{(P\_1,4.1.2)} tatpuruṣatve dvigucagrahaṇam kartavyam . \\
\noindent \textbf{(P\_1,4.1.2)} tatpuruṣaḥ dviguḥ ca iti cakāraḥ kartavyaḥ . \\
\noindent \textbf{(P\_1,4.1.2)} akriyamāṇe hi cakāre anavakāśā dvigusañjñā tatpuruṣasañjñām bādheta . \\
\noindent \textbf{(P\_1,4.1.2)} paravacane hi niyamānupapatteḥ ubhayasañjñābhāvaḥ . \\
\noindent \textbf{(P\_1,4.1.2)} yasya punaḥ paraṅkāryatvam niyamānupappateḥ tasya ubhayoḥ sañjñayoḥ bhāvaḥ siddhaḥ . \\
\noindent \textbf{(P\_1,4.1.2)} katham . \\
\noindent \textbf{(P\_1,4.1.2)} pūrvā tasya dvigusañjñā parā tatpuruṣasañjñā. katham . \\
\noindent \textbf{(P\_1,4.1.2)} evaṃ sa vakṣyati . \\
\noindent \textbf{(P\_1,4.1.2)} taddhitārthottarapadasamāhāre ca saṅkhyāpūrvaḥ dviguḥ iti . \\
\noindent \textbf{(P\_1,4.1.2)} evam sarvam dviguprakaraṇam anukramya tasya ante śritādiḥ tatpuruṣaḥ iti . \\
\noindent \textbf{(P\_1,4.1.2)} tatra ārambhasāmarthyāt ca dvigusañjñā paraṅkāryatvāt ca tatpuruṣasañjñā bhaviṣyati . \\
\noindent \textbf{(P\_1,4.1.2)} nanu ca yasya api ekasañjñādhikāraḥ tasya api tatpuruṣasañjñāpūrvikā dvigusañjñā . \\
\noindent \textbf{(P\_1,4.1.2)} katham . \\
\noindent \textbf{(P\_1,4.1.2)} anuvṛttiḥ kriyate . \\
\noindent \textbf{(P\_1,4.1.2)} paryāyaḥ prasajyeta . \\
\noindent \textbf{(P\_1,4.1.2)} ekā sañjñā iti vacanāt na asti yaugapadyena saṃbhavaḥ . \\
\noindent \textbf{(P\_1,4.1.2)} gatidivaḥkarmahetumatsu cagrahaṇam . \\
\noindent \textbf{(P\_1,4.1.2)} gatidivaḥkarmahetumatsu cagrahaṇam kartavyam . \\
\noindent \textbf{(P\_1,4.1.2)} upasargāḥ kriyāyoge gatiśca iti cakāraḥ kartavyaḥ . \\
\noindent \textbf{(P\_1,4.1.2)} akriyamāṇe hi cakāre anavakāśaḥ pasargasañjñā gatisañjñām bādheta . \\
\noindent \textbf{(P\_1,4.1.2)} paravacane hi niyamānupapatteḥ ubhayasañjñābhāvaḥ . \\
\noindent \textbf{(P\_1,4.1.2)} yasya punaḥ paraṅkāryatvam niyamānupapatteḥ tasya ubhayoḥ sañjñayoḥ bhāvaḥ siddhaḥ . \\
\noindent \textbf{(P\_1,4.1.2)} katham . \\
\noindent \textbf{(P\_1,4.1.2)} pūrvā tasya upasargasañjñā parā gatisañjñā . \\
\noindent \textbf{(P\_1,4.1.2)} atra ārambhasāmarthyāt ca upasargasañjñā paraṅkāryatvāt ca gatisañjñā bhaviṣyati . \\
\noindent \textbf{(P\_1,4.1.2)} nanu ca yasya api ekasañjñādhikāraḥ tasya upasargasañjñāpūrvikā gatisañjñā . \\
\noindent \textbf{(P\_1,4.1.2)} katham . \\
\noindent \textbf{(P\_1,4.1.2)} anuvṛttiḥ kriyate . \\
\noindent \textbf{(P\_1,4.1.2)} paryāyaḥ prasajyeta . \\
\noindent \textbf{(P\_1,4.1.2)} ekā sañjñā iti vacanāt na asti yaugapadyena saṃbhavaḥ . \\
\noindent \textbf{(P\_1,4.1.2)} gatisañjñā api anavakāśā sā vacanāt bhaviṣyati . \\
\noindent \textbf{(P\_1,4.1.2)} sāvakāśā gatisañjñā . \\
\noindent \textbf{(P\_1,4.1.2)} kaḥ avakāśaḥ . \\
\noindent \textbf{(P\_1,4.1.2)} ūryādīni avakāśaḥ . \\
\noindent \textbf{(P\_1,4.1.2)} prādīnāṃ yā gatisañjñā sā anavakāśā . \\
\noindent \textbf{(P\_1,4.1.2)} gati . \\
\noindent \textbf{(P\_1,4.1.2)} divaḥ karma . \\
\noindent \textbf{(P\_1,4.1.2)} sādhakatamam karaṇam divaḥ karma ca iti cakāraḥ kartavyaḥ . \\
\noindent \textbf{(P\_1,4.1.2)} akriyamāṇe hi cakāre anavakāśā karmasañjñā karaṇasañjñām bādheta . \\
\noindent \textbf{(P\_1,4.1.2)} paravacane hi niyamānupapatteḥ ubhayasañjñābhāvaḥ . \\
\noindent \textbf{(P\_1,4.1.2)} yasya punaḥ paraṅkāryatvam niyamānupapatteḥ tasya ubhayoḥ sañjñayoḥ bhāvaḥ siddhaḥ . \\
\noindent \textbf{(P\_1,4.1.2)} katham . \\
\noindent \textbf{(P\_1,4.1.2)} pūrvā tasya karmasañjñā parā karaṇasañjñā . \\
\noindent \textbf{(P\_1,4.1.2)} katham . \\
\noindent \textbf{(P\_1,4.1.2)} evam sa vakṣyati . \\
\noindent \textbf{(P\_1,4.1.2)} divaḥ sādhakatamam karma . \\
\noindent \textbf{(P\_1,4.1.2)} tataḥ karaṇam . \\
\noindent \textbf{(P\_1,4.1.2)} karaṇasañjñām ca bhavati sādhakatamam . \\
\noindent \textbf{(P\_1,4.1.2)} diva iti nivṛttam . \\
\noindent \textbf{(P\_1,4.1.2)} tatra ārambhasāmarthyāt ca karmasañjñā paraṅkāryatvāt ca karaṇasañjñā bhaviṣyati . \\
\noindent \textbf{(P\_1,4.1.2)} nanu ca yasya api ekasañjñādhikāraḥtasya api karaṇasañjñāpūrvikā karmasañjñā . \\
\noindent \textbf{(P\_1,4.1.2)} katham . \\
\noindent \textbf{(P\_1,4.1.2)} anuvṛttiḥ kriyate . \\
\noindent \textbf{(P\_1,4.1.2)} paryāyaḥ prasajyeta . \\
\noindent \textbf{(P\_1,4.1.2)} ekā sañjñā iti vacanāt na asti yaugapadyena saṃbhavaḥ . \\
\noindent \textbf{(P\_1,4.1.2)} divaḥ karma . \\
\noindent \textbf{(P\_1,4.1.2)} hetumat . \\
\noindent \textbf{(P\_1,4.1.2)} svatantraḥ kartā tatprayojako hetuḥ ca iti cakāraḥ kartavyaḥ . \\
\noindent \textbf{(P\_1,4.1.2)} akriyamāṇe hi cakāre anavakāśā hetusañjñā kartṛsañjñām bādheta . \\
\noindent \textbf{(P\_1,4.1.2)} paravacane hi niyamānupapatteḥ ubhayasañjñābhāvaḥ . \\
\noindent \textbf{(P\_1,4.1.2)} yasya punaḥ paraṅkāryatvam niyamānupapatteḥ tasya ubhayoḥ sañjñayorbhāvaḥ siddhaḥ . \\
\noindent \textbf{(P\_1,4.1.2)} katham . \\
\noindent \textbf{(P\_1,4.1.2)} pūrvā tasya hetusañjñā parā kartṛsañjñā . \\
\noindent \textbf{(P\_1,4.1.2)} katham . \\
\noindent \textbf{(P\_1,4.1.2)} evam sa vakṣyati . \\
\noindent \textbf{(P\_1,4.1.2)} svatantraḥ prayojakaḥ hetuḥ iti . \\
\noindent \textbf{(P\_1,4.1.2)} tataḥ kartā . \\
\noindent \textbf{(P\_1,4.1.2)} kartṛsañjñaḥ ca bhavati svatantraḥ . \\
\noindent \textbf{(P\_1,4.1.2)} prayojakaḥ iti nivṛttam . \\
\noindent \textbf{(P\_1,4.1.2)} tatra ārambhasāmarthyāt ca hetusañjñā paraṅkāryatvāt ca kartṛsañjñā bhaviṣyati . \\
\noindent \textbf{(P\_1,4.1.2)} nanu ca yasya api ekasañjñādhikāraḥ tasya api kartṛsañjñāpūrvikā hetusañjñā . \\
\noindent \textbf{(P\_1,4.1.2)} katham . \\
\noindent \textbf{(P\_1,4.1.2)} anuvṛttiḥ kriyate . \\
\noindent \textbf{(P\_1,4.1.2)} paryāyaḥ prasajyeta . \\
\noindent \textbf{(P\_1,4.1.2)} ekā sañjñā iti vacanāt na asti yaugapadyena saṃbhavaḥ . \\
\noindent \textbf{(P\_1,4.1.2)} gurulaghusañjñe nadīghisañjñe . \\
\noindent \textbf{(P\_1,4.1.2)} gurulaghusañjñe nadīghisañjñe bādheyātām . \\
\noindent \textbf{(P\_1,4.1.2)} gārgibandhuḥ vātsībandhuḥ vaitram viviniyya . \\
\noindent \textbf{(P\_1,4.1.2)} paravacane hi niyamānupatteḥ ubhayasañjñābhāvaḥ . \\
\noindent \textbf{(P\_1,4.1.2)} yasya punaḥ paraṅkāryatvam niyamānupatteḥ tasya ubhayoḥ sañjñayoḥ bhāvaḥ siddhaḥ . \\
\noindent \textbf{(P\_1,4.1.2)} katham . \\
\noindent \textbf{(P\_1,4.1.2)} pūrve tasya nadīghisañjñe pare gurulaghusañjñe . \\
\noindent \textbf{(P\_1,4.1.2)} tatra ārambhasāmarthyāt ca nadīghisañjñe paraṅkāryatvāt ca gurulaghusañjñe bhaviṣyataḥ . \\
\noindent \textbf{(P\_1,4.1.2)} nanu ca yasya api ekasañjñādhikāraḥ tasya api nadīghisaṃghisañjñāpūrvike gurulaghusañjñe . \\
\noindent \textbf{(P\_1,4.1.2)} katham . \\
\noindent \textbf{(P\_1,4.1.2)} anuvṛttiḥ kriyate . \\
\noindent \textbf{(P\_1,4.1.2)} paryāyaḥ prasajyeta . \\
\noindent \textbf{(P\_1,4.1.2)} ekā sañjñā iti vacanāt na asti yaugapadyena sambhavaḥ . \\
\noindent \textbf{(P\_1,4.1.2)} parasmaipadasañjñām puruṣasañjñā . parasmaipadasañjñām puruṣasañjñā bādheta . \\
\noindent \textbf{(P\_1,4.1.2)} paravacane hi niyamānupapatteḥ ubhayasañjñābhāvaḥ . \\
\noindent \textbf{(P\_1,4.1.2)} yasya punaḥ paraṅkāryatvam niyamānupapatteḥ tasyobhayoḥ sañjñayoḥ bhāvaḥ siddhaḥ . \\
\noindent \textbf{(P\_1,4.1.2)} katham . \\
\noindent \textbf{(P\_1,4.1.2)} pūrvā tasya puruṣasañjñā parā parasmaipadasañjñā . \\
\noindent \textbf{(P\_1,4.1.2)} katham . \\
\noindent \textbf{(P\_1,4.1.2)} evaṃ sa vakṣyati . \\
\noindent \textbf{(P\_1,4.1.2)} tiṅaḥ trīṇi trīṇi prathamamadhyottamāḥ iti . \\
\noindent \textbf{(P\_1,4.1.2)} evam sarvam puruṣaniyamam anukramya tasya ante laḥ parasmaipadam iti . \\
\noindent \textbf{(P\_1,4.1.2)} tatra ārambhasāmārthyāt ca puruṣasañjñā paraṅkāryatvāt ca parasmaipadasañjñā bhaviṣyati . \\
\noindent \textbf{(P\_1,4.1.2)} nanu ca yasya api ekasañjñādhikaraḥ tasya api parasmaipadasañjñāpūrvikā puruṣasañjñā . \\
\noindent \textbf{(P\_1,4.1.2)} katham . \\
\noindent \textbf{(P\_1,4.1.2)} anuvṛttiḥ kriyate . \\
\noindent \textbf{(P\_1,4.1.2)} paryāyaḥ prasajyeta. ekā sañjñā iti vacanāt na asti yaugapadyena sambhavaḥ . \\
\noindent \textbf{(P\_1,4.1.2)} parasmaipadasañjñā api anavakāśā sā vacanāt bhaviṣyati . \\
\noindent \textbf{(P\_1,4.1.2)} sāvakāśā parasmaipadasañjñā . \\
\noindent \textbf{(P\_1,4.1.2)} kaḥ avakāśaḥ . \\
\noindent \textbf{(P\_1,4.1.2)} śatṛkvasū avakāśaḥ . \\
\noindent \textbf{(P\_1,4.1.3)} paravacane siti padaṃ bham . \\
\noindent \textbf{(P\_1,4.1.3)} paravacane siti padam bhasañjñamapi prāpnoti . \\
\noindent \textbf{(P\_1,4.1.3)} ayam te yoniḥ ṛtviyaḥ . \\
\noindent \textbf{(P\_1,4.1.3)} prajam vindāma ṛtvijayām . \\
\noindent \textbf{(P\_1,4.1.3)} ārambhasāmarthyāt ca padasañjñā paraṅkāryatvāt ca bhasañjñā prāpnoti . \\
\noindent \textbf{(P\_1,4.1.3)} gatibuddhyādīnām ṇyantānām karma kartṛsañjñam . \\
\noindent \textbf{(P\_1,4.1.3)} gatibuddhyādīāṃ ṇyantānām karma kartṛsañjñam api prāpnoti . \\
\noindent \textbf{(P\_1,4.1.3)} ārambhasāmarthyāt ca karmasañjñā padasañjñā paraṅkāryatvāt ca kartṛsañjñā prāpnoti . \\
\noindent \textbf{(P\_1,4.1.3)} na eṣaḥ doṣaḥ . \\
\noindent \textbf{(P\_1,4.1.3)} ācāryapravṛttiḥ jñāpayati na karmasañjñāyām kartṛsañjñā bhavati iti yat ayam hṛkroḥ anyatarasyām iti antarasyāṅgrahaṇam karoti . \\
\noindent \textbf{(P\_1,4.1.3)} śeṣavacanam ca ghisañjñānivṛttyartham . \\
\noindent \textbf{(P\_1,4.1.3)} śeṣagrahaṇam ca kartavyam . \\
\noindent \textbf{(P\_1,4.1.3)} śeṣaḥ ghi asakhi iti . \\
\noindent \textbf{(P\_1,4.1.3)} kim prayojanam . \\
\noindent \textbf{(P\_1,4.1.3)} ghisāñjñānivṛttyartham . \\
\noindent \textbf{(P\_1,4.1.3)} nadīsañjñāyām ghisañjñā mā bhūt iti . \\
\noindent \textbf{(P\_1,4.1.3)} śakaṭyai paddhatyai buddhayai dhenvai . \\
\noindent \textbf{(P\_1,4.1.3)} itarathā hi paraṅkāryatvāt ca ghisañjñā ārambhasāmarthyāt ca ṅiti hrasvaḥ ca iti nadīsañjñā . \\
\noindent \textbf{(P\_1,4.1.3)} na vā asambhavāt . \\
\noindent \textbf{(P\_1,4.1.3)} na vā asambhavāt . \\
\noindent \textbf{(P\_1,4.1.3)} na vā kartavyam . \\
\noindent \textbf{(P\_1,4.1.3)} nadīsañjñāyām ghisañjñā kasmāt na bhavati . \\
\noindent \textbf{(P\_1,4.1.3)} asambhavāt . \\
\noindent \textbf{(P\_1,4.1.3)} kaḥ asau asambhavaḥ . \\
\noindent \textbf{(P\_1,4.1.3)} hrasvalakṣaṇā hi nadīsañjñā ghisañjñāyām ca guṇaḥ . \\
\noindent \textbf{(P\_1,4.1.3)} hrasvalakṣaṇā hi nadīsañjñā ghisañjñāyām ca guṇena bhavitavyam . \\
\noindent \textbf{(P\_1,4.1.3)} tatra vacanaprāmāṇyāt nadīsañjñāyām ghisañjñābhāvaḥ . \\
\noindent \textbf{(P\_1,4.1.3)} tatra vacanaprāmāṇyāt nadīsañjñāyām ghisañjñā na bhaviṣyati . \\
\noindent \textbf{(P\_1,4.1.3)} kim kāraṇam . \\
\noindent \textbf{(P\_1,4.1.3)} āśrayābhāvāt . \\
\noindent \textbf{(P\_1,4.1.3)} āśrayābhāvāt nadīsañjñāyām ghisañjñānivṛttiḥ iti cet yaṇādeśābhāvaḥ . \\
\noindent \textbf{(P\_1,4.1.3)} āśrayābhāvāt nadīsañjñāyām ghisañjñānivṛttiḥ iti cet evam ucyate yaṇādeśaḥ api na prāpnoti . \\
\noindent \textbf{(P\_1,4.1.3)} na eṣaḥ doṣaḥ . \\
\noindent \textbf{(P\_1,4.1.3)} nadyāśrayatvāt yaṇādeśasya hrasvasya nadīsañjñābhāvaḥ . \\
\noindent \textbf{(P\_1,4.1.3)} nadyāśrayaḥ yaṇādeśaḥ . \\
\noindent \textbf{(P\_1,4.1.3)} yadā nadīsañjñayā ghisañjñā bādhitā tata uttarakālam yaṇādeśena bhavitavyam . \\
\noindent \textbf{(P\_1,4.1.3)} nadyāśrayatvāt yaṇādeśasya hrasvasya nadīsañjñā bhaviṣyati . \\
\noindent \textbf{(P\_1,4.1.3)} bahuvrīhyartham tu . bahvrīhipratiṣedhārtham tu śeṣagrahaṇam kartavyam . \\
\noindent \textbf{(P\_1,4.1.3)} śeṣaḥ bavuvrīhiḥ iti . \\
\noindent \textbf{(P\_1,4.1.3)} kiṃ prayojanam . \\
\noindent \textbf{(P\_1,4.1.3)} prayojanamavyayībhāvopamānadvigukṛllopeṣu . avyayībhāve . \\
\noindent \textbf{(P\_1,4.1.3)} unmattagaṅgam lohitagaṅgam . \\
\noindent \textbf{(P\_1,4.1.3)} upamāne . \\
\noindent \textbf{(P\_1,4.1.3)} śastrīśyāmā kumudaśyenī . \\
\noindent \textbf{(P\_1,4.1.3)} dvigu . \\
\noindent \textbf{(P\_1,4.1.3)} pañcagavam daśagavam . \\
\noindent \textbf{(P\_1,4.1.3)} kṛllope . \\
\noindent \textbf{(P\_1,4.1.3)} niṣkauśāmbiḥ nirvārāṇasiḥ . \\
\noindent \textbf{(P\_1,4.1.3)} tatra śeṣavacanāt doṣaḥ saṅkhyāsamānādhikaraṇañsamāseṣu bahuvrīhipratiṣedhaḥ . \\
\noindent \textbf{(P\_1,4.1.3)} tatra śeṣavacanāt doṣaḥ bhavati . \\
\noindent \textbf{(P\_1,4.1.3)} saṅkhyāsamānādhikaraṇañsamāseṣu bahuvrīheḥ pratiṣedhaḥ prāpnoti . \\
\noindent \textbf{(P\_1,4.1.3)} saṅkhyā . \\
\noindent \textbf{(P\_1,4.1.3)} dvīrāvatīkaḥ deśaḥ trīrāvatīkaḥ deśaḥ . \\
\noindent \textbf{(P\_1,4.1.3)} samānādhikaraṇa . \\
\noindent \textbf{(P\_1,4.1.3)} vīrpuruṣakaḥ grāmaḥ . \\
\noindent \textbf{(P\_1,4.1.3)} nañsamāse . \\
\noindent \textbf{(P\_1,4.1.3)} abrāhmaṇakaḥ deśaḥ avṛṣalakaḥ deśaḥ . \\
\noindent \textbf{(P\_1,4.1.3)} kṛllope ca śeṣavacanāt prādibhiḥ na bahuvrīhiḥ . \\
\noindent \textbf{(P\_1,4.1.3)} kṛllope ca śeṣavacanāt prādibhiḥ na prāpnoti . \\
\noindent \textbf{(P\_1,4.1.3)} prapatitaparṇaḥ praparṇakaḥ prapatitapalāśaḥ prapalāśakaḥ iti . \\
\noindent \textbf{(P\_1,4.1.3)} atha ekasañjñādhikāre katham sidhyati . \\
\noindent \textbf{(P\_1,4.1.3)} ekasañjñādhikāre vipratiṣedhād bahuvrīhiḥ . \\
\noindent \textbf{(P\_1,4.1.3)} ekasañjñādhikāre vipratiṣedhāt bahuvrīhiḥ bhaviṣyati . \\
\noindent \textbf{(P\_1,4.1.3)} ekasañjñādhikāre vipratiṣedhāt bahuvrīhiḥ iti cet ktārthe pratiṣedhaḥ . \\
\noindent \textbf{(P\_1,4.1.3)} ekasañjñādhikāre vipratiṣedhāt bahuvrīhiḥ iti cet ktārthe pratiṣedhaḥ vaktavyaḥ . \\
\noindent \textbf{(P\_1,4.1.3)} niṣkauśāmbiḥ nirvārāṇasiḥ . \\
\noindent \textbf{(P\_1,4.1.3)} tatpuruṣaḥ atra bādhakaḥ bhaviṣyati . \\
\noindent \textbf{(P\_1,4.1.3)} tatpuruṣaḥ iti cet anyatra ktārthāt pratiṣedhaḥ . tatpuruṣaḥ iti cet anyatra ktārthāt pratiṣedhaḥ vaktavyaḥ . \\
\noindent \textbf{(P\_1,4.1.3)} prapatitaparṇaḥ praparṇakaḥ pratitatpalāśaḥ prapalāśakaḥ iti . \\
\noindent \textbf{(P\_1,4.1.3)} siddham tu prādīnām ktārthe tatpuruṣavacanāt . \\
\noindent \textbf{(P\_1,4.1.3)} siddhametat . \\
\noindent \textbf{(P\_1,4.1.3)} katham . \\
\noindent \textbf{(P\_1,4.1.3)} prādīnām ktārthe tatpuruṣaḥ bhavati iti vaktavyam . \\
\noindent \textbf{(P\_1,4.1.4)} kāni punaḥ asya yogasya prayojanāni . \\
\noindent \textbf{(P\_1,4.1.4)} prayojanam hrasvasañjñām dīrghaplutau . \\
\noindent \textbf{(P\_1,4.1.4)} hrasvasañjñām dīrghaplutasañjñe bādhete . \\
\noindent \textbf{(P\_1,4.1.4)} tiṅsārvadhātukam liṅliṭoḥ ārdhadhātukam . \\
\noindent \textbf{(P\_1,4.1.4)} tiṅaḥ sārvadhātukasañjñām liṅliṭoḥārdhadhātusañjñā bādhate . \\
\noindent \textbf{(P\_1,4.1.4)} apatyam vṛddham yuvā . \\
\noindent \textbf{(P\_1,4.1.4)} apatyam vṛddham yavusañjñā bādhate . \\
\noindent \textbf{(P\_1,4.1.4)} ghim nadī . \\
\noindent \textbf{(P\_1,4.1.4)} ghisañjñām nadīsañjñā bādhate . \\
\noindent \textbf{(P\_1,4.1.4)} laghu guru . \\
\noindent \textbf{(P\_1,4.1.4)} laghusañjñām gurusañjñā bādhate . \\
\noindent \textbf{(P\_1,4.1.4)} padam bham . \\
\noindent \textbf{(P\_1,4.1.4)} padasañjñām bhasañjñā bādhate . \\
\noindent \textbf{(P\_1,4.1.4)} apādādnam uttarāṇi dhanuṣā vidhyati kaṃsapātryām bhuṅkte gām dogdhi dhanuḥ vidhyati iti . \\
\noindent \textbf{(P\_1,4.1.4)} apādānasañjñām uttarāṇi kārakāṇi bādhante . \\
\noindent \textbf{(P\_1,4.1.4)} kva . \\
\noindent \textbf{(P\_1,4.1.4)} dhanuṣā vidhyati . \\
\noindent \textbf{(P\_1,4.1.4)} kaṃsapātryām bhuṅkte . \\
\noindent \textbf{(P\_1,4.1.4)} gām dogdhi . \\
\noindent \textbf{(P\_1,4.1.4)} dhanuḥ vidhyati . \\
\noindent \textbf{(P\_1,4.1.4)} dhanuṣā vidhyati iti apāyayuktatvāt ca dhruvamapāye apādānam iti apādānasañjñā prāpnoti sādhakatamam karaṇam iti ca karaṇasañjñā . \\
\noindent \textbf{(P\_1,4.1.4)} karaṇasañjñā parā sā bhavati . \\
\noindent \textbf{(P\_1,4.1.4)} kaṃsapātryām bhuṅkte iti atra apāyuktatvāt ca dhruvamapāye apādānamiti apādānasañjñā prāpnoti ādhāraḥ adhikaraṇam iti ca adhikaraṇasañjñā . \\
\noindent \textbf{(P\_1,4.1.4)} adhikaraṇasañjñā parā sā bhavati . \\
\noindent \textbf{(P\_1,4.1.4)} gām dogdhi iti atra apāyuktatvāt ca apādānasañjñā prāpnoti karturīpsitatamam karma . \\
\noindent \textbf{(P\_1,4.1.4)} iti ca karmasañjñā . \\
\noindent \textbf{(P\_1,4.1.4)} karmasañjñā parā sā bhavati . \\
\noindent \textbf{(P\_1,4.1.4)} dhanuḥ vidhyati iti atra apāyayuktatvāt ca apādānasañjñā prāpnoti svatantraḥ kartā iti ca . \\
\noindent \textbf{(P\_1,4.1.4)} kartṛsañjñā parā sā bhavati . \\
\noindent \textbf{(P\_1,4.1.4)} krudhadruhoḥ upsṛṣṭayoḥ karma saṃpradānam . krudhadruhoḥ upsṛṣṭayoḥ karmasañjñā saṃpradānasañjñām bādhate . \\
\noindent \textbf{(P\_1,4.1.4)} karaṇam parāṇi sādhu asiḥ chinatti . \\
\noindent \textbf{(P\_1,4.1.4)} karaṇasañjñām parāṇi āṇi bādhante . \\
\noindent \textbf{(P\_1,4.1.4)} kva . \\
\noindent \textbf{(P\_1,4.1.4)} dhanuḥ vidhyati . \\
\noindent \textbf{(P\_1,4.1.4)} asiḥ chinatti iti . \\
\noindent \textbf{(P\_1,4.1.4)} adhikaraṇam karma geham praviśati . adhikaraṇasañjñām karmasañjñā bādhate . \\
\noindent \textbf{(P\_1,4.1.4)} kva . \\
\noindent \textbf{(P\_1,4.1.4)} geham praviśati iti . \\
\noindent \textbf{(P\_1,4.1.4)} adhikaraṇam kartā sthālī pacati . \\
\noindent \textbf{(P\_1,4.1.4)} adhikaraṇasañjñām karmasañjñā bādhate . \\
\noindent \textbf{(P\_1,4.1.4)} kva . \\
\noindent \textbf{(P\_1,4.1.4)} sthālī pacati iti . \\
\noindent \textbf{(P\_1,4.1.4)} adhyupasṛṣṭam karma . adhyupasṛṣṭam karma adhikaraṇasañjñām bādhate . \\
\noindent \textbf{(P\_1,4.1.4)} gatyupasargasañjñe karmapravacanīyasañjñā . \\
\noindent \textbf{(P\_1,4.1.4)} gatyupasargasañjñe karmapravacanīyasañjñā bādhate . \\
\noindent \textbf{(P\_1,4.1.4)} parasmaipadam ātmanepadam . parasmaipadasañjñām ātmanepadasañjñā bādhate . \\
\noindent \textbf{(P\_1,4.1.4)} samāsasañjñāḥ ca . \\
\noindent \textbf{(P\_1,4.1.4)} samāsasañjñāḥ ca yāḥ yāḥ parāḥ anavakāśāḥ ca tāḥ tāḥ pūrvāḥ sāvakāśāḥ ca bādhante . \\
\noindent \textbf{(P\_1,4.1.5)} arthavat prātipadikam . arthavat prātipadikasañjñam bhavati . \\
\noindent \textbf{(P\_1,4.1.5)} guṇavacanam ca . \\
\noindent \textbf{(P\_1,4.1.5)} guṇavacanasañjñam ca bhavati arthavat . \\
\noindent \textbf{(P\_1,4.1.5)} samāsakṛttaddhitāvyayasarvanāma asarvaliṅgā jātiḥ . \\
\noindent \textbf{(P\_1,4.1.5)} 41||samāsa | samāsasañjñā ca vaktavyā . \\
\noindent \textbf{(P\_1,4.1.5)} kṛt . \\
\noindent \textbf{(P\_1,4.1.5)} kṛtsañjñā ca vaktavyā . \\
\noindent \textbf{(P\_1,4.1.5)} taddhita . \\
\noindent \textbf{(P\_1,4.1.5)} taddhitasañjñā ca vaktavyā . \\
\noindent \textbf{(P\_1,4.1.5)} avyaya . \\
\noindent \textbf{(P\_1,4.1.5)} avyayasañjñā ca vaktavyā . \\
\noindent \textbf{(P\_1,4.1.5)} sarvanāma . \\
\noindent \textbf{(P\_1,4.1.5)} sarvanāmasañjñā ca vaktavyā . \\
\noindent \textbf{(P\_1,4.1.5)} asarvaliṅgā jātiḥ iti etat ca vaktavyam . \\
\noindent \textbf{(P\_1,4.1.5)} saṅkhyā . \\
\noindent \textbf{(P\_1,4.1.5)} saṅkhyāsañjñā ca vaktavyā . \\
\noindent \textbf{(P\_1,4.1.5)} ḍu ca . \\
\noindent \textbf{(P\_1,4.1.5)} ḍusañjñā ca vaktavyā . \\
\noindent \textbf{(P\_1,4.1.5)} kā punaḥ ḍusañjñā . \\
\noindent \textbf{(P\_1,4.1.5)} ṣaṭsañjñā . \\
\noindent \textbf{(P\_1,4.1.5)} ekadravyopaniveśinī sañjñā . \\
\noindent \textbf{(P\_1,4.1.5)} ekadravyopaniveśinī sañjñā iti etat ca vaktavyam . \\
\noindent \textbf{(P\_1,4.1.5)} kimartham idamucyate . \\
\noindent \textbf{(P\_1,4.1.5)} yathānyāse eva bhūyiṣṭhāḥ sañjñāḥ kriyante . \\
\noindent \textbf{(P\_1,4.1.5)} santi ca eva atra kāḥ cit apūrvāḥ sañjñāḥ . \\
\noindent \textbf{(P\_1,4.1.5)} api ca etena ānupūrvyeṇa sanniviṣṭānām bādhanam yathā syāt . \\
\noindent \textbf{(P\_1,4.1.5)} guṇavacanasañjñāyāḥ ca etābhiḥ bādhanam yathā syāt iti . \\
\noindent \textbf{(P\_1,4.2.1)} vipratiṣedhaḥ iti kaḥ ayam śabdaḥ . \\
\noindent \textbf{(P\_1,4.2.1)} vipratipūrvāt siddheḥ karmavyatihāre karmavyatihāre ghañ . \\
\noindent \textbf{(P\_1,4.2.1)} itaretarapratiṣedhaḥ vipratiṣedhaḥ . \\
\noindent \textbf{(P\_1,4.2.1)} anyonyapratiṣedhaḥ vipratiṣedhaḥ . \\
\noindent \textbf{(P\_1,4.2.1)} kaḥ punaḥ vipratiṣedhaḥ . \\
\noindent \textbf{(P\_1,4.2.1)} dvau prasaṅgau anyārthau ekasmin saḥ vipratiṣedhaḥ . \\
\noindent \textbf{(P\_1,4.2.1)} dvau prasaṇgau yadā anyārthau bhavataḥ ekasmin ca yugapat prāptnutaḥ saḥ vipratiṣedhaḥ . \\
\noindent \textbf{(P\_1,4.2.1)} kva punaḥ anyārthau kva ca ekasmin yugapat prāpnutaḥ . \\
\noindent \textbf{(P\_1,4.2.1)} vṛkṣābhyām , vṛkṣeṣu iti anyārthau vṛkṣebhyaḥ iti atra yugapat prāpnutaḥ . \\
\noindent \textbf{(P\_1,4.2.1)} kim ca syāt . \\
\noindent \textbf{(P\_1,4.2.1)} ekasmin yugapadasambhavāt pūrvaparaprāpteḥ ubhayaprasaṅgaḥ . \\
\noindent \textbf{(P\_1,4.2.1)} ekasmin yugapadasambhavātpūrvasyāḥ ca parasyāḥ ca prāpteḥ ubhayaprasaṅgaḥ . \\
\noindent \textbf{(P\_1,4.2.1)} idam vipratiṣiddham yat ucyate ekasmin yugapadasambhavātpūrvaparaprāpteḥ ubhayaprasaṅgaḥ iti . \\
\noindent \textbf{(P\_1,4.2.1)} katham hi ekasmin ca nāma yugapadasambhavaḥ syāt pūrvasyāḥ ca parasyāḥca prāpteḥ ubhayaprasaṅgaḥ ca syāt . \\
\noindent \textbf{(P\_1,4.2.1)} na etat vipratiṣiddham . \\
\noindent \textbf{(P\_1,4.2.1)} yat ucyate ekasmin yugapadasambhavāt iti kāryayoḥ yugapadasambhavaḥ śāstrayoḥ ubhayaprasaṅgaḥ . \\
\noindent \textbf{(P\_1,4.2.1)} tṛjādibhiḥ tulyam . tṛjādibhiḥ tulyam paryāyaḥ prāpnoti . \\
\noindent \textbf{(P\_1,4.2.1)} tat yathā tṛjādayaḥ paryāyeṇa bhavanti . \\
\noindent \textbf{(P\_1,4.2.1)} kim punaḥ kāraṇam tṛjādayaḥ paryāyeṇa bhavanti . \\
\noindent \textbf{(P\_1,4.2.1)} anavayaprasaṅgāt pratipadam vidheḥ ca . \\
\noindent \textbf{(P\_1,4.2.1)} anavayavena prasajyante pratipadam ca vidhīyante . \\
\noindent \textbf{(P\_1,4.2.1)} apratipattiḥ vā ubhayoḥ tulyabalatvāt . \\
\noindent \textbf{(P\_1,4.2.1)} apratipattiḥ vā punaḥ ubhayoḥ śāstrayoḥ syāt . \\
\noindent \textbf{(P\_1,4.2.1)} kim kāraṇam . \\
\noindent \textbf{(P\_1,4.2.1)} tulyabalatvāt . \\
\noindent \textbf{(P\_1,4.2.1)} tulyabale hi ubhe śāstre . \\
\noindent \textbf{(P\_1,4.2.1)} tat yathā . \\
\noindent \textbf{(P\_1,4.2.1)} dvayoḥ tulyabalayoḥ ekaḥ preṣyaḥ bhavati . \\
\noindent \textbf{(P\_1,4.2.1)} saḥ tayoḥ paryāyeṇa kāryam karoti . \\
\noindent \textbf{(P\_1,4.2.1)} yadā tam ubhau yugapat preṣayataḥ nānādikṣu ca kārye bhavataḥ tadā yadi asau avirodhārthī bhavati tataḥ ubhayoḥ na karoti . \\
\noindent \textbf{(P\_1,4.2.1)} kim punaḥ kāraṇam ubhayoḥ na karoti . \\
\noindent \textbf{(P\_1,4.2.1)} yaugapadyāsambhavāt . \\
\noindent \textbf{(P\_1,4.2.1)} na asti yaugapadyena sambhavaḥ . \\
\noindent \textbf{(P\_1,4.2.1)} tatra pratipattyartham vacanam . \\
\noindent \textbf{(P\_1,4.2.1)} tatra pratipattyartham idam vaktavyam . \\
\noindent \textbf{(P\_1,4.2.1)} tavyadādīnām tu aprasiddhiḥ . \\
\noindent \textbf{(P\_1,4.2.1)} tavyadādīnām tu kāryasya aprasiddhiḥ . \\
\noindent \textbf{(P\_1,4.2.1)} na hi kim cit tavyadādiṣu niyamakāri śāstram ārabhyate yena tavyadādayaḥ syuḥ . \\
\noindent \textbf{(P\_1,4.2.1)} yaḥ ca bhavatā hetuḥ vyapadiṣṭaḥ apratipattiḥ vā ubhayoḥ tulyabalatvāt iti tulyaḥ sa tavyadādiṣu . \\
\noindent \textbf{(P\_1,4.2.1)} na eṣaḥ doṣaḥ . \\
\noindent \textbf{(P\_1,4.2.1)} anavakāśāḥ tavyadādayaḥ ucyante ca . \\
\noindent \textbf{(P\_1,4.2.1)} te vacanāt bhaviṣyanti . \\
\noindent \textbf{(P\_1,4.2.1)} yaḥ ca bhavatā hetuḥ vyapadiṣṭaḥ tṛjādibhiḥ tulyam paryāyaḥ prāpnoti iti tulyaḥ sa tavyadādiṣu . \\
\noindent \textbf{(P\_1,4.2.1)} etāvat iha sūtram vipratiṣedhe param iti . \\
\noindent \textbf{(P\_1,4.2.1)} paṭhiṣyati hi ācāryaḥ sakṛdgatau virpratiṣedhe yat bādhitam tat bādhitam eva iti . \\
\noindent \textbf{(P\_1,4.2.1)} punaḥ ca paṭhiṣyati punaḥ prasaṅgavijñānāt siddham iti . \\
\noindent \textbf{(P\_1,4.2.1)} kim punaḥ iyatā sūtreṇa ubhayam labhyam . \\
\noindent \textbf{(P\_1,4.2.1)} labhyam iti āha . \\
\noindent \textbf{(P\_1,4.2.1)} katham . \\
\noindent \textbf{(P\_1,4.2.1)} iha bhavatā dvau hetū vyapadiṣṭau . \\
\noindent \textbf{(P\_1,4.2.1)} tṛjādibhiḥ tulyam paryāyaḥ prāpnoti iti ca apratipattiḥ vā ubhayoḥ tulyabalatvāt iti ca . \\
\noindent \textbf{(P\_1,4.2.1)} tat yadā tāvat eṣaḥ hetuḥ tṛjādibhiḥ tulyam paryāyaḥ prāpnoti iti tadā vipratiṣeḍhe param iti anena kim kriyate . \\
\noindent \textbf{(P\_1,4.2.1)} niyamaḥ . \\
\noindent \textbf{(P\_1,4.2.1)} vipratiṣedhe param eva bhavati iti . \\
\noindent \textbf{(P\_1,4.2.1)} tadā etat upapannam bhavati sakṛdgatau vipratiṣedhe yat bādhitam tat bādhitam eva iti . \\
\noindent \textbf{(P\_1,4.2.1)} yadā tu eṣaḥ hetuḥ apratipattiḥ vā ubhayoḥ tulyabalatvāt iti tadā vipratiṣedhe param iti anena kim kriyate . \\
\noindent \textbf{(P\_1,4.2.1)} dvāram . \\
\noindent \textbf{(P\_1,4.2.1)} vipratiṣedhe param tāvat bhavati tasmin krṭe yadi pūrvam api prāpnoti tat api bhavati . \\
\noindent \textbf{(P\_1,4.2.1)} tadā etat upapannam bhavati punaḥprasaṅgavijñānāt siddham iti . \\
\noindent \textbf{(P\_1,4.2.1)} vipratiṣedhe param iti uktvā aṅgādhikāre pūrvam iti vaktavyam . \\
\noindent \textbf{(P\_1,4.2.1)} kim kṛtam bhavati . \\
\noindent \textbf{(P\_1,4.2.1)} pūrvavipratiṣedhāḥ na paṭhitavyāḥ bhavanti . \\
\noindent \textbf{(P\_1,4.2.1)} guṇavṛddhyautvatṛjvadbhāvebhyaḥ num pūrvavipratiṣiddham . \\
\noindent \textbf{(P\_1,4.2.1)} numaciratṛjvadbhāvebhyaḥ nuṭ iti . \\
\noindent \textbf{(P\_1,4.2.1)} katham ye paravipratiṣedhāḥ . \\
\noindent \textbf{(P\_1,4.2.1)} ittvottvābhyām guṇavṛddhī bhavataḥ vipratiṣedhena iti . \\
\noindent \textbf{(P\_1,4.2.1)} sūtram ca bhidyate . \\
\noindent \textbf{(P\_1,4.2.1)} yathānyāsam eva astu . \\
\noindent \textbf{(P\_1,4.2.1)} katham ye pūrvavipratiṣedhāḥ . \\
\noindent \textbf{(P\_1,4.2.1)} vipratiṣedhe param iti eva siddham . \\
\noindent \textbf{(P\_1,4.2.1)} katham . \\
\noindent \textbf{(P\_1,4.2.1)} paraśabdaḥ ayam bahvarthaḥ . \\
\noindent \textbf{(P\_1,4.2.1)} asti eva vyavasthāyām vartate . \\
\noindent \textbf{(P\_1,4.2.1)} tat yathā pūrvaḥ paraḥ iti . \\
\noindent \textbf{(P\_1,4.2.1)} asti anyārthe vartate . \\
\noindent \textbf{(P\_1,4.2.1)} paraputraḥ parabhāryā . \\
\noindent \textbf{(P\_1,4.2.1)} anyaputraḥ anyabhāryā iti gamyate . \\
\noindent \textbf{(P\_1,4.2.1)} asti prādhānye vartate . \\
\noindent \textbf{(P\_1,4.2.1)} tat yathā param iyaṃ brāhmāṇī asmin kuṭumbe . \\
\noindent \textbf{(P\_1,4.2.1)} pradhānam iti gamyate . \\
\noindent \textbf{(P\_1,4.2.1)} asti iṣṭavācī paraśabdaḥ . \\
\noindent \textbf{(P\_1,4.2.1)} tat yathā . \\
\noindent \textbf{(P\_1,4.2.1)} param dhāma gataḥ iti . \\
\noindent \textbf{(P\_1,4.2.1)} iṣṭam dhāma iti gamyate . \\
\noindent \textbf{(P\_1,4.2.1)} tat yaḥ iṣṭavācī paraśabdaḥ tasya idam grahaṇam . \\
\noindent \textbf{(P\_1,4.2.1)} vipratiṣedhe param yat iṣṭam tat bhavati . \\
\noindent \textbf{(P\_1,4.2.2)} antaraṅgam ca . \\
\noindent \textbf{(P\_1,4.2.2)} antaraṅgam ca balīyaḥ bhavati iti vaktavyam . \\
\noindent \textbf{(P\_1,4.2.2)} kim prayojanam . \\
\noindent \textbf{(P\_1,4.2.2)} prayojanam yaṇekādeśettvottvāni guṇavṛddhidvirvacanāllopasvarebhyaḥ . \\
\noindent \textbf{(P\_1,4.2.2)} guṇāt yaṇādeśaḥ : syonaḥ , syonā . \\
\noindent \textbf{(P\_1,4.2.2)} guṇaḥ ca prāpnoti yaṇādeḥ ca . \\
\noindent \textbf{(P\_1,4.2.2)} paratvāt guṇaḥ syāt . \\
\noindent \textbf{(P\_1,4.2.2)} yaṇādeśaḥ bhavatyantaraṅgataḥ . \\
\noindent \textbf{(P\_1,4.2.2)} vṛddheḥ yaṇādeśaḥ . \\
\noindent \textbf{(P\_1,4.2.2)} dyaukāmiḥ syaukāmiḥ . \\
\noindent \textbf{(P\_1,4.2.2)} vṛddhiḥ ca prāpnoti yaṇādeḥ ca . \\
\noindent \textbf{(P\_1,4.2.2)} paratvāt vṛddhiḥ syāt . \\
\noindent \textbf{(P\_1,4.2.2)} yaṇādeśaḥ bhavati antaraṅgataḥ . \\
\noindent \textbf{(P\_1,4.2.2)} dvirvacanāt yaṇādeśaḥ . \\
\noindent \textbf{(P\_1,4.2.2)} dudyūṣati susyūṣati . \\
\noindent \textbf{(P\_1,4.2.2)} dvirvacanam ca prāpnoti yaṇādeḥ ca . \\
\noindent \textbf{(P\_1,4.2.2)} nityatvāt dvirvacanam syāt . \\
\noindent \textbf{(P\_1,4.2.2)} yaṇādeśaḥ bhavati antaraṅgataḥ . \\
\noindent \textbf{(P\_1,4.2.2)} allopasya ca yaṇādeśasya ca na asti sampradhāraṇā . \\
\noindent \textbf{(P\_1,4.2.2)} svarāt yaṇādeśaḥ . \\
\noindent \textbf{(P\_1,4.2.2)} dyaukāmiḥ syaukāmiḥ . \\
\noindent \textbf{(P\_1,4.2.2)} svaraḥ ca prāpnoti yaṇādeḥ ca . \\
\noindent \textbf{(P\_1,4.2.2)} paratvāt svaraḥ syāt . \\
\noindent \textbf{(P\_1,4.2.2)} yaṇādeśaḥ bhavati antaraṅgataḥ . \\
\noindent \textbf{(P\_1,4.2.2)} guṇāt ekādeśaḥ . \\
\noindent \textbf{(P\_1,4.2.2)} kādraveyaḥ mantram apaśyat . \\
\noindent \textbf{(P\_1,4.2.2)} guṇaḥ ca prāpnoti ekādeśaḥ ca . \\
\noindent \textbf{(P\_1,4.2.2)} paratvāt guṇaḥ syāt . \\
\noindent \textbf{(P\_1,4.2.2)} ekādeśaḥ bhavati antaraṅgataḥ . \\
\noindent \textbf{(P\_1,4.2.2)} vṛddheḥ ekādeśaḥ . \\
\noindent \textbf{(P\_1,4.2.2)} vaikṣmāṇiḥ sausthitiḥ . \\
\noindent \textbf{(P\_1,4.2.2)} vṛddhiḥ ca prāpnoti ekādeśaḥ ca . \\
\noindent \textbf{(P\_1,4.2.2)} paratvāt vṛddhiḥ syāt . \\
\noindent \textbf{(P\_1,4.2.2)} ekādeśaḥ bhavati antaraṅgataḥ . \\
\noindent \textbf{(P\_1,4.2.2)} dvirvacanāt ekādeśaḥ . \\
\noindent \textbf{(P\_1,4.2.2)} jñāyā odanaḥ jñaudanaḥ . \\
\noindent \textbf{(P\_1,4.2.2)} jñaudanam icchati jñaudanīyati . \\
\noindent \textbf{(P\_1,4.2.2)} jñaudanīyateḥ san jujñaudanīyiṣati . \\
\noindent \textbf{(P\_1,4.2.2)} dvirvacanam ca prāpnoti ekādeśaḥ ca . \\
\noindent \textbf{(P\_1,4.2.2)} nityatvāt dvirvacanam syāt . \\
\noindent \textbf{(P\_1,4.2.2)} ekādeśaḥ bhavati antaraṅgataḥ . \\
\noindent \textbf{(P\_1,4.2.2)} allopāt ekādeśaḥ . \\
\noindent \textbf{(P\_1,4.2.2)} śunā śune . \\
\noindent \textbf{(P\_1,4.2.2)} allopaḥ ca prāpnoti ekādeśaḥ ca . \\
\noindent \textbf{(P\_1,4.2.2)} paratvāt allopaḥ syāt . \\
\noindent \textbf{(P\_1,4.2.2)} ekādeśaḥ bhavati antaraṅgataḥ . \\
\noindent \textbf{(P\_1,4.2.2)} na etat asti prayojanam . \\
\noindent \textbf{(P\_1,4.2.2)} na asti atra viśeṣaḥ allopena vā nivṛttau satyām pūrvatvena vā . \\
\noindent \textbf{(P\_1,4.2.2)} ayam asti viśeṣaḥ . \\
\noindent \textbf{(P\_1,4.2.2)} allopena nivṛttau satyām udāttanivṛttisvaraḥ prasajyeta . \\
\noindent \textbf{(P\_1,4.2.2)} na atra udāttanivṛttisvaraḥ prāpnoti . \\
\noindent \textbf{(P\_1,4.2.2)} kim kāraṇam . \\
\noindent \textbf{(P\_1,4.2.2)} na gośvansāvavarṇa iti pratiṣedhāt . \\
\noindent \textbf{(P\_1,4.2.2)} na eṣaḥ udāttanivṛttisvarasya pratiṣedhaḥ . \\
\noindent \textbf{(P\_1,4.2.2)} kasya tarhi . \\
\noindent \textbf{(P\_1,4.2.2)} tṛtīyādisvarasya . \\
\noindent \textbf{(P\_1,4.2.2)} yatra tarhi tṛtīyādisvaraḥ na asti . \\
\noindent \textbf{(P\_1,4.2.2)} śunaḥ paśya iti . \\
\noindent \textbf{(P\_1,4.2.2)} evam tarhi na lākṣaṇikasya pratiṣedham śiṣmaḥ . \\
\noindent \textbf{(P\_1,4.2.2)} kim tarhi . \\
\noindent \textbf{(P\_1,4.2.2)} yena kena cit lakṣaṇena prāptasya vibhaktisvarasya pratiṣedhaḥ . \\
\noindent \textbf{(P\_1,4.2.2)} yatra tarhi vibhaktiḥ na asti . \\
\noindent \textbf{(P\_1,4.2.2)} bahhuśunī iti . \\
\noindent \textbf{(P\_1,4.2.2)} yadi punaḥ ayam udāttanivṛttisvarasya api pratiṣedhaḥ vijñāyeta . \\
\noindent \textbf{(P\_1,4.2.2)} na evam śakyam . \\
\noindent \textbf{(P\_1,4.2.2)} iha api prasajyeta : kumārī iti . \\
\noindent \textbf{(P\_1,4.2.2)} evam tarhi ācāryapravṛttiḥ jñāpayati nodāttanivṛttisvaraḥ śuni avatarati iti yat ayam śvanśabdam gaurādiṣu paṭhati . \\
\noindent \textbf{(P\_1,4.2.2)} antodāttārtham yatnam karoti . \\
\noindent \textbf{(P\_1,4.2.2)} siddham hi syān ṅīpā eva . \\
\noindent \textbf{(P\_1,4.2.2)} svarāt ekādekādeśaḥ . \\
\noindent \textbf{(P\_1,4.2.2)} sautthitiḥ vaikṣmāṇiḥ . \\
\noindent \textbf{(P\_1,4.2.2)} svaraḥ ca prāpnoti ekādeśaḥ ca . \\
\noindent \textbf{(P\_1,4.2.2)} paratvātsvaraḥ syāt . \\
\noindent \textbf{(P\_1,4.2.2)} ekādeśaḥ bhavati antaraṅgataḥ . \\
\noindent \textbf{(P\_1,4.2.2)} guṇasya ca ittvottvayoḥ ca na asti sampradhāraṇā . \\
\noindent \textbf{(P\_1,4.2.2)} vṛddheḥ ittvottve . \\
\noindent \textbf{(P\_1,4.2.2)} staurṇiḥ paurtiḥ . \\
\noindent \textbf{(P\_1,4.2.2)} vṛddhiḥ ca prāpnoti ittvottve ca . \\
\noindent \textbf{(P\_1,4.2.2)} paratvāt vṛddhiḥ syāt . \\
\noindent \textbf{(P\_1,4.2.2)} ittvottve bhavataḥ antaraṅgataḥ . \\
\noindent \textbf{(P\_1,4.2.2)} dvirvacanāt ittvottve . \\
\noindent \textbf{(P\_1,4.2.2)} ātestīryate āpopūryate . \\
\noindent \textbf{(P\_1,4.2.2)} dvirvacanam ca prāpnoti ittvottve ca . \\
\noindent \textbf{(P\_1,4.2.2)} nityatvāt dvirvacanam syāt . \\
\noindent \textbf{(P\_1,4.2.2)} ittvottve bhavataḥ antaraṅgataḥ . \\
\noindent \textbf{(P\_1,4.2.2)} allopasya ca ittvottvayoḥ ca na asti sampradhāraṇā . \\
\noindent \textbf{(P\_1,4.2.2)} svare nāsti viśeṣaḥ . \\
\noindent \textbf{(P\_1,4.2.2)} iṇṅiśīnām āt guṇaḥ savarṇadīrghatvāt . \\
\noindent \textbf{(P\_1,4.2.2)} iṇṅiśīnām āt guṇaḥ savarṇadīrghatvāt prayojanam . \\
\noindent \textbf{(P\_1,4.2.2)} ayaje indram avape indram . \\
\noindent \textbf{(P\_1,4.2.2)} vṛkṣe indram plakṣe indram . \\
\noindent \textbf{(P\_1,4.2.2)} ye indram te indram . \\
\noindent \textbf{(P\_1,4.2.2)} āt guṇaḥ ca prāpnoti savarṇadīrghatvam ca . \\
\noindent \textbf{(P\_1,4.2.2)} paratvāt savaraṇadīrghatvam syāt . \\
\noindent \textbf{(P\_1,4.2.2)} āt guṇaḥ bhavati antaraṅgataḥ . \\
\noindent \textbf{(P\_1,4.2.2)} na vā savarṇadīrghatvasya anavakāśatvāt . \\
\noindent \textbf{(P\_1,4.2.2)} na vā etat antaraṅgeṇa api sidhyati . \\
\noindent \textbf{(P\_1,4.2.2)} kim kāraṇam . \\
\noindent \textbf{(P\_1,4.2.2)} savarṇadīrghatvasya anavakāśatvāt . \\
\noindent \textbf{(P\_1,4.2.2)} anavakāśam savarṇadīrghatvam āt guṇam bādheta . \\
\noindent \textbf{(P\_1,4.2.2)} na etat antaraṅge asti anavakāśam param iti . \\
\noindent \textbf{(P\_1,4.2.2)} iha api syonaḥ , syonā iti śakyam vaktum na vā paratvāt guṇasya iti . \\
\noindent \textbf{(P\_1,4.2.2)} ūṅāpoḥ ekādeśaḥ ītvalopābhyām . \\
\noindent \textbf{(P\_1,4.2.2)} ūṅāpoḥ ekādeśaḥ ītvalopābhyām bhavati antaraṅgataḥ prayojanam . \\
\noindent \textbf{(P\_1,4.2.2)} ītvāt ekādeśaḥ . \\
\noindent \textbf{(P\_1,4.2.2)} khaṭvīyati mālīyati . \\
\noindent \textbf{(P\_1,4.2.2)} ītvam ca prāpnoti ekādeśaḥ ca . \\
\noindent \textbf{(P\_1,4.2.2)} paratvāt ītvam syāt . \\
\noindent \textbf{(P\_1,4.2.2)} ekādeśaḥ bhavati antaraṅgataḥ . \\
\noindent \textbf{(P\_1,4.2.2)} lopāt ekādeśaḥ . \\
\noindent \textbf{(P\_1,4.2.2)} kāmaṇḍaleyaḥ bhādrabāheyaḥ . \\
\noindent \textbf{(P\_1,4.2.2)} lopaḥ ca prāpnoti ekādeśaḥ ca . \\
\noindent \textbf{(P\_1,4.2.2)} paratvāt lopaḥ syāt . \\
\noindent \textbf{(P\_1,4.2.2)} ekādeśaḥ bhavati antaraṅgataḥ . \\
\noindent \textbf{(P\_1,4.2.2)} atha kimartham ītvalopābhyām iti ucyate na lopetvābhyāmiti eva ucyeta . \\
\noindent \textbf{(P\_1,4.2.2)} saṅkhyātānudeśaḥ mā bhūt iti . \\
\noindent \textbf{(P\_1,4.2.2)} āpaḥ api ekādeśaḥ lope prayojayati . \\
\noindent \textbf{(P\_1,4.2.2)} cauḍiḥ bālākiḥ . \\
\noindent \textbf{(P\_1,4.2.2)} āttvanapuṃsakopasarjanahrasvatvāni ayavāyāvekādeśatugvidhibhyaḥ . āttvanapuṃsakopasarjanahrasvatvāni ayavāyāvekādeśatugvidhibhyaḥ bhavanti antaraṅgataḥ . \\
\noindent \textbf{(P\_1,4.2.2)} veñ vānīyam śo śānīyam glai glānīyam mlai mlānīyam glācchattram clācchatram . \\
\noindent \textbf{(P\_1,4.2.2)} āttvam ca prāpnoti ete ca vidhayaḥ . \\
\noindent \textbf{(P\_1,4.2.2)} paratvāt ete vidhayaḥ syuḥ . \\
\noindent \textbf{(P\_1,4.2.2)} āttvam bhavati antaraṅgataḥ . \\
\noindent \textbf{(P\_1,4.2.2)} napuṃsakopasarjanahrasvatvam ca prayojanam . \\
\noindent \textbf{(P\_1,4.2.2)} atiri atra atinu atra atiricchattram atinucchatram ārāśastri idam dhānāśaṣkuli idam niṣkauśāmbi idam nirvārāṇasi idam niṣkauśāmbicchatram nirvārāṇāsicchatram . \\
\noindent \textbf{(P\_1,4.2.2)} napuṃsakopasarjanahrasvatvam ca prāpnoti ete ca vidhayaḥ . \\
\noindent \textbf{(P\_1,4.2.2)} paratvāt ete vidhayaḥ syuḥ . \\
\noindent \textbf{(P\_1,4.2.2)} napuṃsakopasarjanahrasvatvam bhavati antaraṅgataḥ . \\
\noindent \textbf{(P\_1,4.2.2)} tuk yaṇekādeśaguṇavṛddhyauttvadīrghatvetvamumetttvarīvidhibhyaḥ . yaṇekādeśaguṇavṛddhyauttvadīrghatvetvamumetttvarīvidhibhyaḥ tuk bhavati antaraṅgataḥ . \\
\noindent \textbf{(P\_1,4.2.2)} yaṇādeśāt . \\
\noindent \textbf{(P\_1,4.2.2)} agnicit atra somasut atra . \\
\noindent \textbf{(P\_1,4.2.2)} ekādeśāt . \\
\noindent \textbf{(P\_1,4.2.2)} agnicit idam somasut udakam . \\
\noindent \textbf{(P\_1,4.2.2)} guṇāt . \\
\noindent \textbf{(P\_1,4.2.2)} agnicite somasute . \\
\noindent \textbf{(P\_1,4.2.2)} vṛddheḥ . \\
\noindent \textbf{(P\_1,4.2.2)} praṛcchakaḥ prārcchakaḥ . \\
\noindent \textbf{(P\_1,4.2.2)} auttvāt . \\
\noindent \textbf{(P\_1,4.2.2)} agniciti somasuti . \\
\noindent \textbf{(P\_1,4.2.2)} dīrghatvāt . \\
\noindent \textbf{(P\_1,4.2.2)} jagadbhyām janagadbhyām . \\
\noindent \textbf{(P\_1,4.2.2)} ītvāt . \\
\noindent \textbf{(P\_1,4.2.2)} jagatyati janagatyati . \\
\noindent \textbf{(P\_1,4.2.2)} mumaḥ . \\
\noindent \textbf{(P\_1,4.2.2)} agnicinmanyaḥ somasunmanyaḥ . \\
\noindent \textbf{(P\_1,4.2.2)} etvāt . \\
\noindent \textbf{(P\_1,4.2.2)} jagadbhyaḥ janagadbhyaḥ . \\
\noindent \textbf{(P\_1,4.2.2)} rīvidheḥ . \\
\noindent \textbf{(P\_1,4.2.2)} sukṛtyati pāpakrṭyati . \\
\noindent \textbf{(P\_1,4.2.2)} anaṅānaṅbhyām ca iti vaktavyam . \\
\noindent \textbf{(P\_1,4.2.2)} sukṛt sukṛtdduṣkrṭau . \\
\noindent \textbf{(P\_1,4.2.2)} tuk ca prāpnoti ete ca vidhayaḥ . \\
\noindent \textbf{(P\_1,4.2.2)} paratvāt ete vidhayaḥ syuḥ . \\
\noindent \textbf{(P\_1,4.2.2)} tuk bhavati antaraṅgataḥ . \\
\noindent \textbf{(P\_1,4.2.2)} iyaṅādeśaḥ guṇāt . \\
\noindent \textbf{(P\_1,4.2.2)} iyaṅādeśaḥ guṇāt bhavati antaraṅgataḥ prayojanam . \\
\noindent \textbf{(P\_1,4.2.2)} dhiyati riyati . \\
\noindent \textbf{(P\_1,4.2.2)} iyaṅādeśaḥ ca prāpnoti guṇaḥ ca . \\
\noindent \textbf{(P\_1,4.2.2)} paratvāt guṇaḥ syāt . \\
\noindent \textbf{(P\_1,4.2.2)} iyaṅādeśaḥ bhavati antaraṇgataḥ . \\
\noindent \textbf{(P\_1,4.2.2)} uvaṅādeśaḥ ca iti vaktavyam . \\
\noindent \textbf{(P\_1,4.2.2)} prādudruvat prāsusruvat . \\
\noindent \textbf{(P\_1,4.2.2)} śveḥ samprasāraṇapūrvatvam yaṇādeśāt . \\
\noindent \textbf{(P\_1,4.2.2)} śveḥ samprasāraṇapūrvatvam yaṇādeśāt bhavati antaraṇgataḥ prayojanam . \\
\noindent \textbf{(P\_1,4.2.2)} śuśuvatuḥ śuśuvuḥ . \\
\noindent \textbf{(P\_1,4.2.2)} pūrvatvam ca prāpnoti yaṇādeśaḥ ca . \\
\noindent \textbf{(P\_1,4.2.2)} paratvāt yaṇādeśaḥ syāt . \\
\noindent \textbf{(P\_1,4.2.2)} pūrvatvam bhavati antaraṇgataḥ . \\
\noindent \textbf{(P\_1,4.2.2)} hvaḥ ākāralopāt . \\
\noindent \textbf{(P\_1,4.2.2)} hvaḥ ākāralopāt pūrvatvam bhavati antaraṅgataḥ prayojanam . \\
\noindent \textbf{(P\_1,4.2.2)} juhuvatuḥ juhuvuḥ . \\
\noindent \textbf{(P\_1,4.2.2)} pūrvatvam ca prāpnoti ākāralopaḥ ca paratvāt ākāralopaḥ syāt . \\
\noindent \textbf{(P\_1,4.2.2)} pūrvatvam bhavati antaraṅgataḥ . \\
\noindent \textbf{(P\_1,4.2.2)} svaraḥ lopāt . \\
\noindent \textbf{(P\_1,4.2.2)} svaraḥ lopāt bhavati antaraṅgataḥ prayojanam . \\
\noindent \textbf{(P\_1,4.2.2)} aupagavī saudāmanī . \\
\noindent \textbf{(P\_1,4.2.2)} svaraḥ ca prāpnoti lopaḥ ca . \\
\noindent \textbf{(P\_1,4.2.2)} paratvāt lopaḥ syāt . \\
\noindent \textbf{(P\_1,4.2.2)} svaraḥ bhavati antaraṅgataḥ . \\
\noindent \textbf{(P\_1,4.2.2)} pratyayavidhiḥ ekādeśāt . \\
\noindent \textbf{(P\_1,4.2.2)} pratyayavidhiḥ ekādeśāt bhavati antaraṅgataḥ prayojanam . \\
\noindent \textbf{(P\_1,4.2.2)} agniḥ indraḥ vāyuḥ udakam . \\
\noindent \textbf{(P\_1,4.2.2)} pratyayavidhiḥ ca prāpnoti ekādeśaḥ ca . \\
\noindent \textbf{(P\_1,4.2.2)} paratvāt ekādeśaḥ syāt . \\
\noindent \textbf{(P\_1,4.2.2)} | pratyayavidhiḥ bhavati antaraṅgataḥ . \\
\noindent \textbf{(P\_1,4.2.2)} yaṇādeśāt ca iti vaktavyam . \\
\noindent \textbf{(P\_1,4.2.2)} agniḥ atra vāyuḥ atra . \\
\noindent \textbf{(P\_1,4.2.2)} lādeśaḥ varṇavidheḥ . lādeśaḥ varṇavidheḥ bhavati antaraṅgataḥ prayojanam . \\
\noindent \textbf{(P\_1,4.2.2)} pacatu atra paṭhtu atra . \\
\noindent \textbf{(P\_1,4.2.2)} lādeśaḥ ca prāpnoti yaṇādeśaḥ ca .paratvāt yaṇādeśaḥ syāt . \\
\noindent \textbf{(P\_1,4.2.2)} lādeśaḥ bhavati antaraṅgataḥ . \\
\noindent \textbf{(P\_1,4.2.2)} tatpuruṣāntodāttatvam pūrvapadaprakṛtisvarāt . \\
\noindent \textbf{(P\_1,4.2.2)} tatpuruṣāntodāttatvam pūrvapadaprakṛtisvarāt bhavati antaraṅgataḥ prayojanam . \\
\noindent \textbf{(P\_1,4.2.2)} pūrvaśālāpriyaḥ aparaśālāpriyaḥ ṭatpuruṣāntodāttatvam ca prāpnoti pūrvapadaprakṛtisvaratvam ca . \\
\noindent \textbf{(P\_1,4.2.2)} paratvāt pūrvapadaprakṛtisvaratvam syāt . \\
\noindent \textbf{(P\_1,4.2.2)} tatpuruṣāntodāttatvam bhavati antaraṅgataḥ . \\
\noindent \textbf{(P\_1,4.2.2)} etāni asyāḥ paribhāṣāyāḥ prayojanāni yadartham eṣā paribhāṣā kartavyā . \\
\noindent \textbf{(P\_1,4.2.3)} yadi santi prayojanāni iti eṣā paribhāṣā kriyate nanu ca iyam api kartavyā asiddham bahiraṅgalakṣaṇam antaraṅgalakṣaṇe iti . \\
\noindent \textbf{(P\_1,4.2.3)} kim prayojanam . \\
\noindent \textbf{(P\_1,4.2.3)} pacāvedam pacāmedam . \\
\noindent \textbf{(P\_1,4.2.3)} asiddhatvāt bahiraṅgalakṣaṇasya guṇasya antaraṇgalakṣaṇam aitvam mā bhūt iti . \\
\noindent \textbf{(P\_1,4.2.3)} ubhe tarhi kartavye . \\
\noindent \textbf{(P\_1,4.2.3)} na iti āha . \\
\noindent \textbf{(P\_1,4.2.3)} anayā eva siddham . \\
\noindent \textbf{(P\_1,4.2.3)} iha api syonaḥ syonā iti asiddhatvāt bahiraṅgalakṣaṇasya guṇasya antaraṅgalakṣaṇaḥ yaṇādeśo bhaviṣyati . \\
\noindent \textbf{(P\_1,4.2.3)} yadi asiddham bahiraṅgalakṣaṇam antaraṅgalakṣaṇe iti ucyate akṣadyūḥ hiraṇyadyūḥ asiddhatvāt asiddhatvāt bahiraṅgalakṣaṇasya ūṭhaḥ antaraṅgalakṣaṇaḥ yaṇādeśaḥ na prāpnoti . \\
\noindent \textbf{(P\_1,4.2.3)} na eṣaḥ doṣaḥ . \\
\noindent \textbf{(P\_1,4.2.3)} asiddham bahiraṅgalakṣaṇam antaraṅgalakṣaṇe iti uktvā tataḥ vakṣyāmi na ajānantarye bahiṣṭvaprakḷptiḥ iti . \\
\noindent \textbf{(P\_1,4.2.3)} sā tarhi eṣā paribhāṣā kartavyā . \\
\noindent \textbf{(P\_1,4.2.3)} na kartavyā . \\
\noindent \textbf{(P\_1,4.2.3)} ācāryapravṛttiḥ jñāpayati bhavati eṣā paribhāṣā iti yat ayam ṣatvatukoḥ asiddhaḥ iti āha . \\
\noindent \textbf{(P\_1,4.2.3)} iyam tarhi paribhāṣā kartavyā asiddham bahiraṅgalakṣaṇamantaraṅgalakṣaṇe iti . \\
\noindent \textbf{(P\_1,4.2.3)} eṣā ca na kartavyā . \\
\noindent \textbf{(P\_1,4.2.3)} ācāryapravṛttiḥ jñāpayati bhavati eṣā paribhāṣā iti yat ayam vāhaḥ ūṭh iti ūṭham śāsti . \\
\noindent \textbf{(P\_1,4.2.3)} tasya doṣaḥ pūrvapadottarapadayoḥ vṛddhisvarau ekādeśāt . \\
\noindent \textbf{(P\_1,4.2.3)} tasya etasya lakṣaṇasya doṣaḥ pūrvapadottarapadayoḥ vṛddhisvarau ekādeśāt antaraṅgataḥ abhinirvṛttāt na prāpnutaḥ . \\
\noindent \textbf{(P\_1,4.2.3)} pūrvaiṣukāmaśamaḥ aparaiṣukāmaśamaḥ guḍodakam tilodakam . \\
\noindent \textbf{(P\_1,4.2.3)} udake akevale iti pūrvottarapadayoḥ vyapavargābhāvāt na syāt . \\
\noindent \textbf{(P\_1,4.2.3)} na eṣaḥ doṣaḥ . \\
\noindent \textbf{(P\_1,4.2.3)} ācāryapravṛttiḥ jñāpayati pūrvottarapadayoḥ tāvat kāryam bhavati na ekādeśaḥ iti yat ayam na indrasya parasya iti pratiṣedham śāsti . \\
\noindent \textbf{(P\_1,4.2.3)} katham kṛtvā jñāpakam . \\
\noindent \textbf{(P\_1,4.2.3)} indre dvau acau . \\
\noindent \textbf{(P\_1,4.2.3)} tatra ekaḥ yasya īti ca iti lopena hriyate aparaḥ ekādeśena . \\
\noindent \textbf{(P\_1,4.2.3)} tataḥ anackaḥ indraḥ sampannaḥ . \\
\noindent \textbf{(P\_1,4.2.3)} tatra kaḥ prasaṅgaḥ vṛddheḥ . \\
\noindent \textbf{(P\_1,4.2.3)} paśyati tu ācāryaḥ pūrvapadottarapadyoḥ tāvatkāryam bhavati na ekādeśaḥ iti tataḥ na indrasya parasya iti pratiṣedham śāsti . \\
\noindent \textbf{(P\_1,4.2.3)} yaṇādeśāt iyuvau . yaṇādeśāt iyuvau antaraṅgataḥ abhinirvṛttāt na prāpnutaḥ . \\
\noindent \textbf{(P\_1,4.2.3)} vaiyākaraṇaḥ sauvaśvaḥ iti . \\
\noindent \textbf{(P\_1,4.2.3)} lakṣaṇam hi bhavati yvoḥ vṛddhiprasaṅge iyuvau bhavataḥ iti . \\
\noindent \textbf{(P\_1,4.2.3)} na eṣaḥ doṣaḥ . \\
\noindent \textbf{(P\_1,4.2.3)} anavakāśau iyuvau . \\
\noindent \textbf{(P\_1,4.2.3)} aci iti ucyate . \\
\noindent \textbf{(P\_1,4.2.3)} kim punaḥ kāraṇam aci ti ucyate . \\
\noindent \textbf{(P\_1,4.2.3)} iha mā bhūtām . \\
\noindent \textbf{(P\_1,4.2.3)} aitikāyanaḥ aupagavaḥ iti . \\
\noindent \textbf{(P\_1,4.2.3)} stām atra iyuvau lopaḥ vyoḥ vali iti lopaḥ bhaviṣyati . \\
\noindent \textbf{(P\_1,4.2.3)} yatra tarhi lopaḥ na asti . \\
\noindent \textbf{(P\_1,4.2.3)} praiyamedhaḥ praiyamgavaḥ iti . \\
\noindent \textbf{(P\_1,4.2.3)} usi pararūpāt ca . \\
\noindent \textbf{(P\_1,4.2.3)} usi pararūpāt ca antaraṅgataḥ abhinirvṛttāt iyādeśaḥ na prāpnoti . \\
\noindent \textbf{(P\_1,4.2.3)} paceyuḥ yajeyuḥ . \\
\noindent \textbf{(P\_1,4.2.3)} na eṣaḥ doṣaḥ . \\
\noindent \textbf{(P\_1,4.2.3)} na evam vijñāyate yā iti etasya iy bhavati iti . \\
\noindent \textbf{(P\_1,4.2.3)} katham tarhi . \\
\noindent \textbf{(P\_1,4.2.3)} yās iti etasya iy bhavati iti . \\
\noindent \textbf{(P\_1,4.2.3)} luk lopayaṇayavāyāvekādeśebhyaḥ . \\
\noindent \textbf{(P\_1,4.2.3)} lopayaṇayavāyāvekādeśebhyaḥ luk balīyān iti vaktavyam . \\
\noindent \textbf{(P\_1,4.2.3)} lopāt . \\
\noindent \textbf{(P\_1,4.2.3)} gomān priyaḥ asya gomatpriyaḥ yavamatpriyaḥ . \\
\noindent \textbf{(P\_1,4.2.3)} gomān iva ācarati gomatyate yavamatyate . \\
\noindent \textbf{(P\_1,4.2.3)} yaṇādeśāt . \\
\noindent \textbf{(P\_1,4.2.3)} grāmaṇyaḥ kulam grāmaṇikulam senānyaḥ kulam senānikulam . \\
\noindent \textbf{(P\_1,4.2.3)} ayavāyāvekādeśebhyaḥ . \\
\noindent \textbf{(P\_1,4.2.3)} gave hitam gohitam rāyaḥ kulam raikulam nāvaḥ kulam naukulam vṛkādbhayam vṛkabhayam . \\
\noindent \textbf{(P\_1,4.2.3)} luk ca prāpnoti ete ca vidhayaḥ . \\
\noindent \textbf{(P\_1,4.2.3)} paratvāt ete vidhayaḥ syuḥ . \\
\noindent \textbf{(P\_1,4.2.3)} luk balīyān iti vaktavyam luk yathā syāt . \\
\noindent \textbf{(P\_1,4.3.1)} yū iti kimartham . \\
\noindent \textbf{(P\_1,4.3.1)} khaṭvā mālā . \\
\noindent \textbf{(P\_1,4.3.1)} kim ca syāt . \\
\noindent \textbf{(P\_1,4.3.1)} khaṭvābandhuḥ mālābandhuḥ . \\
\noindent \textbf{(P\_1,4.3.1)} nadī bandhuni iti eṣaḥ svaraḥ prasajyeta . \\
\noindent \textbf{(P\_1,4.3.1)} iha ca bahukhaṭvakaḥ iti nadyṛtaḥ ca iti nityaḥ kap prasajyeta . \\
\noindent \textbf{(P\_1,4.3.1)} na eṣaḥ doṣaḥ . \\
\noindent \textbf{(P\_1,4.3.1)} ācāryapravṛttiḥ jñāpayati na āpaḥ nadīsañjñā bhavati iti yat ayam ṅeḥ ām nadyāmnībhyaḥ iti pṛthak ābgrahaṇam karoti . \\
\noindent \textbf{(P\_1,4.3.1)} iha tarhi mātre mātuḥ iti āṭ nadyāḥ iti āṭ prasajyeta . \\
\noindent \textbf{(P\_1,4.3.1)} kim punaḥ idam dīrghayoḥ grahaṇam āhosvid hrasvayoḥ . \\
\noindent \textbf{(P\_1,4.3.1)} kim cātaḥ . \\
\noindent \textbf{(P\_1,4.3.1)} yadi dīrghayoḥ grahaṇam yū iti nirdeśaḥ na upapadyate . \\
\noindent \textbf{(P\_1,4.3.1)} dīrghāt hi pūrvasavarṇaḥ pratiṣidhyate . \\
\noindent \textbf{(P\_1,4.3.1)} uttaratra ca viśeṣaṇam na prakalpeta yū hrasvau iti . \\
\noindent \textbf{(P\_1,4.3.1)} yadi yū na hrasvau . \\
\noindent \textbf{(P\_1,4.3.1)} atha hrasvau na yū . \\
\noindent \textbf{(P\_1,4.3.1)} yū hrasvau ceti vipratiṣiddham . \\
\noindent \textbf{(P\_1,4.3.1)} atha hrasvayoḥ he śakaṭe atra api prasajyeta . \\
\noindent \textbf{(P\_1,4.3.1)} na eṣaḥ doṣaḥ . \\
\noindent \textbf{(P\_1,4.3.1)} avaśyam atra vibhāṣā nadīsañjñā eṣitavyā . \\
\noindent \textbf{(P\_1,4.3.1)} ubhayam hi iṣyate : he śakaṭi he śakaṭe iti . \\
\noindent \textbf{(P\_1,4.3.1)} iha tarhi śakaṭibandhuḥ iti nadī bandhuni iti eṣaḥ svaraḥ prasajyeta . \\
\noindent \textbf{(P\_1,4.3.1)} iha ca bahuśakaṭiḥ iti nadyṛtaḥ ceti nityaḥ kap prasajyeta . \\
\noindent \textbf{(P\_1,4.3.1)} na eṣaḥ doṣaḥ . \\
\noindent \textbf{(P\_1,4.3.1)} ṅiti hrasvaḥ ca iti ayam niyamārthaḥ bhaviṣyati . \\
\noindent \textbf{(P\_1,4.3.1)} ṅiti eva yū hrasvau nadīsañjñau bhavataḥ na nyatra iti . \\
\noindent \textbf{(P\_1,4.3.1)} kaimarthakyāt niyamaḥ bhavati . \\
\noindent \textbf{(P\_1,4.3.1)} vidheyam na asti iti kṛtvā . \\
\noindent \textbf{(P\_1,4.3.1)} iha ca asti vidheyam . \\
\noindent \textbf{(P\_1,4.3.1)} kim . \\
\noindent \textbf{(P\_1,4.3.1)} nityā nadīsañjñā prāptā sā vibhāṣā vidheyā . \\
\noindent \textbf{(P\_1,4.3.1)} tatra apūrvaḥ vidhiḥ astu niyamaḥ astu iti apūrvaḥ eva vidhiḥ bhaviṣyati na niyamaḥ . \\
\noindent \textbf{(P\_1,4.3.1)} atha ayam nityaḥ yogaḥ syāt prakalpeta niyamaḥ . \\
\noindent \textbf{(P\_1,4.3.1)} bāḍham prakalpeta . \\
\noindent \textbf{(P\_1,4.3.1)} nityaḥ tarthi bhaviṣyati | tat katham . \\
\noindent \textbf{(P\_1,4.3.1)} yogavibhāgaḥ kariṣyate . \\
\noindent \textbf{(P\_1,4.3.1)} idam asti . \\
\noindent \textbf{(P\_1,4.3.1)} yū stryākhyau nadī na iyaṅuvaṅsthānau astrī vāmi . \\
\noindent \textbf{(P\_1,4.3.1)} tataḥ ṅiti . \\
\noindent \textbf{(P\_1,4.3.1)} ṅiti ca iyaṅuvaṅsthānau yū vā astrī nadīsañjñau na bhavataḥ . \\
\noindent \textbf{(P\_1,4.3.1)} tataḥ hrasvau . \\
\noindent \textbf{(P\_1,4.3.1)} hrasvau ca yū stryākhyau ṅiti nadīsañjñau bhavataḥ . \\
\noindent \textbf{(P\_1,4.3.1)} iyaṅuvaṅsthānau vā na iti ca nivṛttam . \\
\noindent \textbf{(P\_1,4.3.1)} yadi evam śakaṭaye atra guṇaḥ na prāpnoti . \\
\noindent \textbf{(P\_1,4.3.1)} dvitīyaḥ yogavibhāgaḥ kariṣyate . \\
\noindent \textbf{(P\_1,4.3.1)} śeṣagrahaṇam na kariṣyate . \\
\noindent \textbf{(P\_1,4.3.1)} katham . \\
\noindent \textbf{(P\_1,4.3.1)} idam asti . \\
\noindent \textbf{(P\_1,4.3.1)} yū stryākhyau nadī na iyaṅuvaṅsthānau astrī vāmi . \\
\noindent \textbf{(P\_1,4.3.1)} tataḥ ṅiti . \\
\noindent \textbf{(P\_1,4.3.1)} ṅiti ca iyaṅuvaṅsthānau yū vā astrī nadīsañjñau na bhavataḥ . \\
\noindent \textbf{(P\_1,4.3.1)} tataḥ hrasvau . \\
\noindent \textbf{(P\_1,4.3.1)} hrasvau ca yū stryākhyau ṅiti nadīsañjñau bhavataḥ . \\
\noindent \textbf{(P\_1,4.3.1)} tataḥ hrasvau . \\
\noindent \textbf{(P\_1,4.3.1)} hrasvau ca yū stryākhyau ṅiti nadīsañjñau bhavataḥ . \\
\noindent \textbf{(P\_1,4.3.1)} iyaṅuvaṅsthānau vā na iti ca nivṛttam . \\
\noindent \textbf{(P\_1,4.3.1)} tataḥ ghi . \\
\noindent \textbf{(P\_1,4.3.1)} ghisañjñau ca bhavataḥ stryākhyau yū hrasvau ṅiti . \\
\noindent \textbf{(P\_1,4.3.1)} tataḥ asakhi . \\
\noindent \textbf{(P\_1,4.3.1)} sakhivarjitau ca yū hrasvau ghisañjñau bhavataḥ . \\
\noindent \textbf{(P\_1,4.3.1)} stryākhyau ṅiti iti ca nivṛttam . \\
\noindent \textbf{(P\_1,4.3.1)} yadi tarhi śeṣagrahaṇam na kriyate na arthaḥ ekena api yogavibhāgena . \\
\noindent \textbf{(P\_1,4.3.1)} aviśeṣeṇa nadīsañjñā utsargaḥ . \\
\noindent \textbf{(P\_1,4.3.1)} tasyāḥ hrasvayoḥ ghisañjñā bādhikā . \\
\noindent \textbf{(P\_1,4.3.1)} tasyām nityāyām prāptāyām ṅiti vibhāṣā ārabhyate . \\
\noindent \textbf{(P\_1,4.3.1)} atha vā punaḥ astu dīrghayoḥ . \\
\noindent \textbf{(P\_1,4.3.1)} nanu ca uktam nirdeśaḥ na upapadyate . \\
\noindent \textbf{(P\_1,4.3.1)} dīrghāt hi pūrvasavarṇaḥ pratiṣidhyate . \\
\noindent \textbf{(P\_1,4.3.1)} vā chandasi iti evam bhaviṣyati . \\
\noindent \textbf{(P\_1,4.3.1)} chandasi iti ucyate na ca idam chandaḥ . \\
\noindent \textbf{(P\_1,4.3.1)} chandovat sūtrāṇi bhavanti iti . \\
\noindent \textbf{(P\_1,4.3.1)} yat api ucyate uttaratra viśeṣeṇam na prakalpeta yū hrasvau iti . \\
\noindent \textbf{(P\_1,4.3.1)} yadi yū na hrasvau atha hrasvau na yū . \\
\noindent \textbf{(P\_1,4.3.1)} yū hrasvau iti vipratiṣiddham iti . \\
\noindent \textbf{(P\_1,4.3.1)} na etat vipratiṣiddham . \\
\noindent \textbf{(P\_1,4.3.1)} āha ayam yū hrasvau iti . \\
\noindent \textbf{(P\_1,4.3.1)} yadi yū na hrasvau . \\
\noindent \textbf{(P\_1,4.3.1)} atha hrasvau na yū . \\
\noindent \textbf{(P\_1,4.3.1)} te evam vijñāsyāmaḥ yvoḥ yau hrasvau iti . \\
\noindent \textbf{(P\_1,4.3.1)} kau ca yvoḥ hrasvau . \\
\noindent \textbf{(P\_1,4.3.1)} savarṇau . \\
\noindent \textbf{(P\_1,4.3.1)} atha stryākhyau iti kaḥ ayam śabdaḥ . \\
\noindent \textbf{(P\_1,4.3.1)} striyam ācakṣāte stryākhyau . \\
\noindent \textbf{(P\_1,4.3.1)} yadi evam stryākhyāyau iti prāpnoti . \\
\noindent \textbf{(P\_1,4.3.1)} anupasarge hi kaḥ vidhīyate . \\
\noindent \textbf{(P\_1,4.3.1)} na tarhi idānīm idam bhavati yasmin daśa sahasrāṇi putre jāte gavām dadau . \\
\noindent \textbf{(P\_1,4.3.1)} brāhmaṇebhyaḥ priyākhyebhyaḥ saḥ ayam uñchena jīvati . \\
\noindent \textbf{(P\_1,4.3.1)} chandovat kavayaḥ kurvanti . \\
\noindent \textbf{(P\_1,4.3.1)} na hi eṣā iṣṭiḥ . \\
\noindent \textbf{(P\_1,4.3.1)} evam tarhi karmasādhanaḥ bhaviṣyati : striyām ākhyāyete stryākhyau . \\
\noindent \textbf{(P\_1,4.3.1)} yadi karmasādhanaḥ kṛtstriyāḥ dhātustriyāḥ ca na sidhyati . \\
\noindent \textbf{(P\_1,4.3.1)} tantryai lakṣmyai śriyai bhruvai . \\
\noindent \textbf{(P\_1,4.3.1)} evam tarhi bahuvrīhiḥ bhaviṣyati . \\
\noindent \textbf{(P\_1,4.3.1)} striyām ākhyā anayoḥ stryākhyau . \\
\noindent \textbf{(P\_1,4.3.1)} evam api kṛtstriyāḥ dhātustriyāḥ ca na sidhyati . \\
\noindent \textbf{(P\_1,4.3.1)} tantryai lakṣmyai śriyai bhruvai . \\
\noindent \textbf{(P\_1,4.3.1)} evam tarhi vic bhaviṣyati . \\
\noindent \textbf{(P\_1,4.3.1)} atha vā punaḥ astu kaḥ eva . \\
\noindent \textbf{(P\_1,4.3.1)} striyam ācakṣāte stryākhyau iti . \\
\noindent \textbf{(P\_1,4.3.1)} nanu ca uktam stryākhyāyau iti prāpnoti . \\
\noindent \textbf{(P\_1,4.3.1)} anupasarge hi kaḥ vidhīyate . \\
\noindent \textbf{(P\_1,4.3.1)} mūlavibhujādipāṭhāt kaḥ bhaviṣyati . \\
\noindent \textbf{(P\_1,4.3.1)} evam ca kṛtvā saḥ api adoṣaḥ bhavati yat uktam yasmin daśa sahasrāṇi putre jāte gavām dadau . \\
\noindent \textbf{(P\_1,4.3.1)} brāhmaṇebhyaḥ priyākhyebhyaḥ saḥ ayam uñchena jīvati . \\
\noindent \textbf{(P\_1,4.3.2)} atha ākhyāgrahaṇam kimartham . \\
\noindent \textbf{(P\_1,4.3.2)} nadīsañjñāyām ākhyāgrahaṇam strīviṣayārtham . nadīsañjñāyām ākhyāgrahaṇam strīviṣayārtham . \\
\noindent \textbf{(P\_1,4.3.2)} strīviṣayau eva yau nityam tayoḥ eva nadīsañjñā yathā syāt . \\
\noindent \textbf{(P\_1,4.3.2)} iha mā bhūt grāmaṇye senānye striyai iti . \\
\noindent \textbf{(P\_1,4.3.2)} prathamaliṅgagrahaṇam ca . \\
\noindent \textbf{(P\_1,4.3.2)} prathamaliṅgagrahaṇam ca kartavyam . \\
\noindent \textbf{(P\_1,4.3.2)} prathamaliṅge yau stryākhau iti vaktavyam . \\
\noindent \textbf{(P\_1,4.3.2)} kim prayojanam . \\
\noindent \textbf{(P\_1,4.3.2)} prayojanam kviblupsamāsāḥ . kumāryai brāhmaṇāya . \\
\noindent \textbf{(P\_1,4.3.2)} lup . \\
\noindent \textbf{(P\_1,4.3.2)} kharakuṭyai brāhmaṇāya . \\
\noindent \textbf{(P\_1,4.3.2)} atitantryai brāhmaṇāya atilakṣmyai brāhmaṇāya . \\
\noindent \textbf{(P\_1,4.3.2)} tat tarhi vaktavyam . \\
\noindent \textbf{(P\_1,4.3.2)} na vaktavyam . \\
\noindent \textbf{(P\_1,4.3.2)} avayavastrīviṣayatvāt siddham . \\
\noindent \textbf{(P\_1,4.3.2)} avayavaḥ atra strīviṣayaḥ tadāśrayā nadīsañjñā bhaviṣyati . \\
\noindent \textbf{(P\_1,4.3.2)} avayavastrīviṣayatvāt siddham iti cet iyaṅuvaṅsthānapratiṣedhe yaṇsthānapratiṣedhaprasaṅgaḥ avayavasya iyaṅuvaṅsthānatvāt . \\
\noindent \textbf{(P\_1,4.3.2)} avayavastrīviṣayatvāt siddham iti cet iyaṅuvaṅsthānapratiṣedhe yaṇsthānayoḥ api yvoḥ pratiṣedhaḥ prasajyeta . \\
\noindent \textbf{(P\_1,4.3.2)} ādhyai pradhyai brāhmaṇyai . \\
\noindent \textbf{(P\_1,4.3.2)} kim kāraṇam . \\
\noindent \textbf{(P\_1,4.3.2)} avayavasya iyaṅuvaṅsthānatvāt . \\
\noindent \textbf{(P\_1,4.3.2)} avayavaḥ atra iyaṅuvaṅsthānaḥ . \\
\noindent \textbf{(P\_1,4.3.2)} siddham tvaṅgarūpagrahaṇāt yasya aṅgasya iyuvau tatpratiṣedhāt . \\
\noindent \textbf{(P\_1,4.3.2)} siddham etat . \\
\noindent \textbf{(P\_1,4.3.2)} katham . \\
\noindent \textbf{(P\_1,4.3.2)} aṅgarūpam gṛhyate . \\
\noindent \textbf{(P\_1,4.3.2)} yasya aṅgasya iyuvau bhavataḥ tasya idam grahaṇam . \\
\noindent \textbf{(P\_1,4.3.2)} na ca etasya aṅgasya iyuvau bhavataḥ . \\
\noindent \textbf{(P\_1,4.3.2)} hrasveyuvsthānapravṛttau ca strīvacane . \\
\noindent \textbf{(P\_1,4.3.2)} hrasvau ca iyuvsthānau ca pravṛttau ca prāk ca pravṛtteḥ strīvacanau eva nadīsañjñau bhavataḥ iti vaktavyam . \\
\noindent \textbf{(P\_1,4.3.2)} śakaṭyai atiśakaṭyai brāmaṇāṇyai . \\
\noindent \textbf{(P\_1,4.3.2)} kva mā bhūt . \\
\noindent \textbf{(P\_1,4.3.2)} śakaṭaye atiśakaṭaye brāhmaṇāya . \\
\noindent \textbf{(P\_1,4.3.2)} dhenvai atidhenvai brāhmaṇyai . \\
\noindent \textbf{(P\_1,4.3.2)} kva mā bhūt . \\
\noindent \textbf{(P\_1,4.3.2)} dhenave atidhenave brāhmaṇāya . \\
\noindent \textbf{(P\_1,4.3.2)} śriyai atiśriyai brāhmaṇyai . \\
\noindent \textbf{(P\_1,4.3.2)} kva mā bhūt . \\
\noindent \textbf{(P\_1,4.3.2)} śriye atiśriye brāhmaṇāya . \\
\noindent \textbf{(P\_1,4.3.2)} bhruvai atibhruvai brāhmaṇyai . \\
\noindent \textbf{(P\_1,4.3.2)} kva mā bhūt . \\
\noindent \textbf{(P\_1,4.3.2)} bhruve atibhruve brāhmaṇāya . \\
\noindent \textbf{(P\_1,4.3.2)} aparaḥ āha : hrasvau ca iyuvsthānau ca pravṛttau api strīvacanau eva nadīsañjñau bhavataḥ iti vaktavyam : śakaṭyai , atiśakaṭyai brāmaṇāṇyai . \\
\noindent \textbf{(P\_1,4.3.2)} kva mā bhūt . \\
\noindent \textbf{(P\_1,4.3.2)} śakaṭaye atiśakaṭaye brāhmaṇāya . \\
\noindent \textbf{(P\_1,4.3.2)} dhenvai atidhenvai brāhmaṇyai . \\
\noindent \textbf{(P\_1,4.3.2)} kva mā bhūt . \\
\noindent \textbf{(P\_1,4.3.2)} dhenave atidhenave brāhmaṇāya . \\
\noindent \textbf{(P\_1,4.3.2)} śriyai atiśriyai brāhmaṇyai . \\
\noindent \textbf{(P\_1,4.3.2)} kva mā bhūt . \\
\noindent \textbf{(P\_1,4.3.2)} śriye atiśriye brāhmaṇāya . \\
\noindent \textbf{(P\_1,4.3.2)} bhruvai atibhruvai brāhmaṇyai . \\
\noindent \textbf{(P\_1,4.3.2)} kva mā bhūt . \\
\noindent \textbf{(P\_1,4.3.2)} bhruve atibhruve brāhmaṇāya . \\
\noindent \textbf{(P\_1,4.3.2)} kimartham punaḥ idam ucyate . \\
\noindent \textbf{(P\_1,4.3.2)} prathamaliṅgagrahaṇam coditam . \\
\noindent \textbf{(P\_1,4.3.2)} tat dveṣyam vijānīyāt : sarvam etat vikalpate iti . \\
\noindent \textbf{(P\_1,4.3.2)} tat ācāryaḥ suhṛt bhūtvā anvācaṣṭe hrasvau ca iyuvsthānau ca pravṛttau ca prāk ca pravṛtteḥ strīvacanau eva iti . \\
\noindent \textbf{(P\_1,4.9)} yogavibhāgaḥ kartavyaḥ . \\
\noindent \textbf{(P\_1,4.9)} ṣaṣthīyuktaḥ chandasi . \\
\noindent \textbf{(P\_1,4.9)} ṣaṣthīyuktaḥ patiśabdaḥ chandasi ghisañjñaḥ bhavati . \\
\noindent \textbf{(P\_1,4.9)} tataḥ vā . \\
\noindent \textbf{(P\_1,4.9)} vā chandasi sarve vidhayo bhavati . \\
\noindent \textbf{(P\_1,4.9)} supām vyatyayaḥ . \\
\noindent \textbf{(P\_1,4.9)} tiṅām vyatyayaḥ . \\
\noindent \textbf{(P\_1,4.9)} varṇavyatyayaḥ . \\
\noindent \textbf{(P\_1,4.9)} liṅgavyatyayaḥ . \\
\noindent \textbf{(P\_1,4.9)} kālavyatyayaḥ . \\
\noindent \textbf{(P\_1,4.9)} puruṣavyatyayaḥ . \\
\noindent \textbf{(P\_1,4.9)} ātmanepadavyatyayaḥ . \\
\noindent \textbf{(P\_1,4.9)} parasmaipadavyatyayaḥ . \\
\noindent \textbf{(P\_1,4.9)} supām vyatyayaḥ . \\
\noindent \textbf{(P\_1,4.9)} yukta māta asīt dhuri dakṣiṇāyāḥ . \\
\noindent \textbf{(P\_1,4.9)} dakṣiṇāyām iti prāpte . \\
\noindent \textbf{(P\_1,4.9)} tiṅām vyatyayaḥ . \\
\noindent \textbf{(P\_1,4.9)} caṣalam ye aśvayūpaya takṣati . \\
\noindent \textbf{(P\_1,4.9)} takṣanti iti prāpte . \\
\noindent \textbf{(P\_1,4.9)} varṇavyatyayaḥ . \\
\noindent \textbf{(P\_1,4.9)} triṣṭubhaujaḥ śubhitam ugravīram . \\
\noindent \textbf{(P\_1,4.9)} suhitamiti prāpte . \\
\noindent \textbf{(P\_1,4.9)} liṅgavyatyayaḥ . \\
\noindent \textbf{(P\_1,4.9)} madhoḥ gṛhṇāti madhoḥ tṛptāḥ iva āsate . \\
\noindent \textbf{(P\_1,4.9)} madhunaḥ iti prāpte . \\
\noindent \textbf{(P\_1,4.9)} kālavyatyayaḥ . \\
\noindent \textbf{(P\_1,4.9)} śvaḥ agnīn ādhāsyamānena śvaḥ somena yakṣamāṇena . \\
\noindent \textbf{(P\_1,4.9)} śvaḥ ādhātā śvaḥ yaṣṭā iti prāpte . \\
\noindent \textbf{(P\_1,4.9)} puruṣavyatyayaḥ . \\
\noindent \textbf{(P\_1,4.9)} adhā saḥ vīraiḥ daśabhiḥ viyūyāḥ . \\
\noindent \textbf{(P\_1,4.9)} viyūyāt iti prāpte . \\
\noindent \textbf{(P\_1,4.9)} ātmanepadavyatyayaḥ . \\
\noindent \textbf{(P\_1,4.9)} brahmacāriṇam icchate . \\
\noindent \textbf{(P\_1,4.9)} icchati iti prāpte . \\
\noindent \textbf{(P\_1,4.9)} parasmaipadavyatyayaḥ . \\
\noindent \textbf{(P\_1,4.9)} pratīpam anyaḥ ūrmiḥ yudhyati | anvīpam anyaḥ ūrmiḥ yudhyati . \\
\noindent \textbf{(P\_1,4.9)} yudhyate iti prāpte . \\
\noindent \textbf{(P\_1,4.13.1)} yasmāt iti vyapadeśāya . \\
\noindent \textbf{(P\_1,4.13.1)} atha pratyayagrahaṇam kimartham . \\
\noindent \textbf{(P\_1,4.13.1)} yasmāt vidhiḥ tadādi pratyaye aṅgam iti iyati ucyamāne strī iyatī strīyati iti atra api prasajyeta . \\
\noindent \textbf{(P\_1,4.13.1)} pratyayagrahaṇe punaḥ kriyamāṇe na doṣaḥ bhavati . \\
\noindent \textbf{(P\_1,4.13.1)} atha vidhigrahaṇam kimartham . \\
\noindent \textbf{(P\_1,4.13.1)} yasmāt pratyayaḥ tadādi pratyaye aṅgam iti iyati ucyamāne dadhi adhunā madhu adhunā atrāpi prasajyeta . \\
\noindent \textbf{(P\_1,4.13.1)} vidhigrahaṇeṇa punaḥ kriyamāṇe na doṣaḥ bhavati . \\
\noindent \textbf{(P\_1,4.13.1)} tat etat pratyayagrahaṇena vidhigrahaṇena ca samuditena kriyate sanniyogaḥ . \\
\noindent \textbf{(P\_1,4.13.1)} yasmāt yaḥ pratyayaḥ vidhīyate tadādi tasmin aṅgasañjñam bhavati iti . \\
\noindent \textbf{(P\_1,4.13.1)} atha tadādigrahaṇam kimartham . \\
\noindent \textbf{(P\_1,4.13.1)} aṅgasañjñāyām tadādivacanam syādinumartham . aṅgasañjñāyām tadādigrahaṇam kriyate syādyartham numartham ca . \\
\noindent \textbf{(P\_1,4.13.1)} syādyartham tāvat . \\
\noindent \textbf{(P\_1,4.13.1)} kariṣyāvaḥ kariṣyāmaḥ . \\
\noindent \textbf{(P\_1,4.13.1)} numartham . \\
\noindent \textbf{(P\_1,4.13.1)} kuṇḍāni vanāni . \\
\noindent \textbf{(P\_1,4.13.1)} mitsuṭoḥ upasamkhyānam . \\
\noindent \textbf{(P\_1,4.13.1)} mitvataḥ suḍvataḥ ca pasamkhyānam kartavyam . \\
\noindent \textbf{(P\_1,4.13.1)} mitvataḥ . \\
\noindent \textbf{(P\_1,4.13.1)} bhinatti chinatti abhinat acchinat . \\
\noindent \textbf{(P\_1,4.13.1)} suḍvataḥ . \\
\noindent \textbf{(P\_1,4.13.1)} sañcarakastu sañcaskaruḥ . \\
\noindent \textbf{(P\_1,4.13.1)} kim punaḥ kāraṇam na sidhyati . \\
\noindent \textbf{(P\_1,4.13.1)} suṭaḥ bahiraṅgatvāt . \\
\noindent \textbf{(P\_1,4.13.1)} bahiraṅgaḥ suṭ . \\
\noindent \textbf{(P\_1,4.13.1)} antaraṅgaḥ guṇaḥ . \\
\noindent \textbf{(P\_1,4.13.1)} asiddham bahiraṅgam antaraṅge . \\
\noindent \textbf{(P\_1,4.13.1)} vakṣyati etat saṃyogādeḥ guṇavidhāne saṃyogopadhagrahaṇam kṛñartham . \\
\noindent \textbf{(P\_1,4.13.1)} yadi saṃyogopadhagrahaṇam kriyate na arthaḥ saṃyogādigrahaṇena . \\
\noindent \textbf{(P\_1,4.13.1)} iha api sasvaratuḥ sasvaruḥ iti saṃyogopadhasya iti eva siddham . \\
\noindent \textbf{(P\_1,4.13.1)} bhavet evamarthena na arthaḥ . \\
\noindent \textbf{(P\_1,4.13.1)} idam tu na sidhyati sañcakaratuḥ sañcaskaruḥ . \\
\noindent \textbf{(P\_1,4.13.1)} kim punaḥ kāraṇam na sidhyati . \\
\noindent \textbf{(P\_1,4.13.1)} iha tasya vā grahaṇam bhavati tadādeḥ vā na cedam tat na api tadādi . \\
\noindent \textbf{(P\_1,4.13.1)} siddham tu tadādyādivacanāt . \\
\noindent \textbf{(P\_1,4.13.1)} siddham etat . \\
\noindent \textbf{(P\_1,4.13.1)} katham . \\
\noindent \textbf{(P\_1,4.13.1)} tadādyādi aṅgasañjñam bhavati iti vaktavyam . \\
\noindent \textbf{(P\_1,4.13.1)} kim idam tadādyādi iti . \\
\noindent \textbf{(P\_1,4.13.1)} tasya ādiḥ tadādiḥ , tadādiḥ ādiḥ yasya tadidam tadādyādi iti . \\
\noindent \textbf{(P\_1,4.13.1)} saḥ tarhi tathā nirdeśaḥ kartavyaḥ . \\
\noindent \textbf{(P\_1,4.13.1)} na kartavyaḥ . \\
\noindent \textbf{(P\_1,4.13.1)} uttarapadalopaḥ atra draṣṭavyaḥ . \\
\noindent \textbf{(P\_1,4.13.1)} tat yathā : uṣṭramukham iva mukham asya uṣṭramukhaḥ , kharamukhaḥ , evam tadādyādi tadādi iti . \\
\noindent \textbf{(P\_1,4.13.1)} tadekadeśavijñānāt vā siddham . tadekadeśavijñānāt vā siddham etat . \\
\noindent \textbf{(P\_1,4.13.1)} tadekadeśabhūtam tadgrahaṇena gṛhyate. tad yathā . \\
\noindent \textbf{(P\_1,4.13.1)} gaṅgā yamunā devadattā iti . \\
\noindent \textbf{(P\_1,4.13.1)} anekā nadī gaṅgām yamunām ca praviṣṭā gaṅgāyamunāgrahaṇena gṛhyate . \\
\noindent \textbf{(P\_1,4.13.1)} tathā devadattāsthaḥ garbhaḥ devadattāgrahaṇena gṛhyate . \\
\noindent \textbf{(P\_1,4.13.1)} viṣamaḥ upanyāsaḥ . \\
\noindent \textbf{(P\_1,4.13.1)} iha ke cit śabdāḥ aktaparimāṇānām arthānām vācakāḥ bhavanti ye ete saṅkhyāśabdāḥ parimāṇaśabdāḥ ca . \\
\noindent \textbf{(P\_1,4.13.1)} pañca sapta iti : ekena api apāye na bhavanti . \\
\noindent \textbf{(P\_1,4.13.1)} droṇaḥ khārī āḍhakam iti : naivā adhike bhavanti na ca nyūne . \\
\noindent \textbf{(P\_1,4.13.1)} ke cit yāvat eva tat bhavati tāvat eva āhuḥ ye ete jātiśabdāḥ guṇaśabdāḥ ca . \\
\noindent \textbf{(P\_1,4.13.1)} tailam ghṛtam iti : khāryām api bhavanti droṇe api . \\
\noindent \textbf{(P\_1,4.13.1)} śuklaḥ nīlaḥ kṛṣṇaḥ iti : himavati api bhavati vaṭakaṇikāmātre api dravye . \\
\noindent \textbf{(P\_1,4.13.1)} aṅgasañjñā ca api aktaparimāṇānām kriyate . \\
\noindent \textbf{(P\_1,4.13.1)} sā kena adhikasya syāt . \\
\noindent \textbf{(P\_1,4.13.1)} evam tarhi ācāryapravṛttiḥ jñāpayati tadekadeśabhūtam tadgrahaṇena gṛhyate iti yat ayam na idamadasoḥ akoḥ iti sakakārayoḥ pratiṣedham śāsti . \\
\noindent \textbf{(P\_1,4.13.1)} katham kṛtvā jñāpakam . \\
\noindent \textbf{(P\_1,4.13.1)} idamadasoḥ kāryam ucyamānam kaḥ prasaṅgo yat sakakārayoḥ syāt . \\
\noindent \textbf{(P\_1,4.13.1)} paśyati tu ācāryaḥ tadekadeśabhūtam tadgrahaṇena gṛhyate iti tataḥ sakakārayoḥ pratiṣedham śāsti . \\
\noindent \textbf{(P\_1,4.13.2)} atha dvitīyam pratyayagrahaṇam kimartham . \\
\noindent \textbf{(P\_1,4.13.2)} pratyayagrahaṇam padādau aprasaṅgārtham . \\
\noindent \textbf{(P\_1,4.13.2)} pratyayagrahaṇam kriyate padādau aṅgasañjñā mā bhūt iti . \\
\noindent \textbf{(P\_1,4.13.2)} kim ca syāt . \\
\noindent \textbf{(P\_1,4.13.2)} stryartham , śryartham , bhvartham : aṅgasya iti iyaṅuvaṅau syātām . \\
\noindent \textbf{(P\_1,4.13.2)} parimāṇārtham ca . \\
\noindent \textbf{(P\_1,4.13.2)} parimāṇārtham ca dvitīyam pratyayagrahaṇam kriyate . \\
\noindent \textbf{(P\_1,4.13.2)} yasmāt pratyayavidhiḥ tadādi aṅgam iti iyati ucyamāne dāśatayasya api aṅgasañjñā prasajyeta . \\
\noindent \textbf{(P\_1,4.13.2)} tat tarhi kartavyam . \\
\noindent \textbf{(P\_1,4.13.2)} na kartavyam . \\
\noindent \textbf{(P\_1,4.13.2)} kena idānīm aṅgakāryam bhaviṣyati . \\
\noindent \textbf{(P\_1,4.13.2)} pratyaye iti prakṛtya aṅgakāryam adhyeṣye . \\
\noindent \textbf{(P\_1,4.13.2)} pratyaye iti prakṛtya aṅgakāryam adhīṣe prākarot upaihiṣṭa upasargāt pūrvau aḍāṭau prāpnutaḥ . \\
\noindent \textbf{(P\_1,4.13.2)} siddham tu pratyayagrahaṇe yasmāt saḥ tadāditadantavijñānāt . \\
\noindent \textbf{(P\_1,4.13.2)} siddham etat . \\
\noindent \textbf{(P\_1,4.13.2)} katham . \\
\noindent \textbf{(P\_1,4.13.2)} pratyayagrahaṇe yasmāt saḥ pratyayaḥ vihitaḥ tadādeḥ tadantasya ca grahaṇam bhavati iti eṣā paribhāṣā kartavyā . \\
\noindent \textbf{(P\_1,4.13.2)} kaḥ punaḥ atra viśeṣaḥ eṣā paribhāṣā kriyeta pratyayagrahaṇam vā . \\
\noindent \textbf{(P\_1,4.13.2)} avaśyam eṣā paribhāṣā kartavyā . \\
\noindent \textbf{(P\_1,4.13.2)} bahūni etasyāḥ paribhāṣāyāḥ prayojanāni . \\
\noindent \textbf{(P\_1,4.13.2)} prayojanam dhātuprātipadikapratyayasamāsataddhitavidhisvarāḥ . \\
\noindent \textbf{(P\_1,4.13.2)} dhātu . \\
\noindent \textbf{(P\_1,4.13.2)} devadattaḥ cikīrṣati . \\
\noindent \textbf{(P\_1,4.13.2)} saṅghātasya dhātusañjñā prāpnoti . \\
\noindent \textbf{(P\_1,4.13.2)} prātipadika . \\
\noindent \textbf{(P\_1,4.13.2)} devadattḥ gārgyaḥ . \\
\noindent \textbf{(P\_1,4.13.2)} saṅghātasya prātipadikasañjñā prāpnoti . \\
\noindent \textbf{(P\_1,4.13.2)} pratyaya . \\
\noindent \textbf{(P\_1,4.13.2)} mahāntam putram icchati . \\
\noindent \textbf{(P\_1,4.13.2)} samghātāt pratyayotpattiḥ prāpnoti . \\
\noindent \textbf{(P\_1,4.13.2)} samāsa . \\
\noindent \textbf{(P\_1,4.13.2)} ṛddhasya rājñaḥ puruṣaḥ . \\
\noindent \textbf{(P\_1,4.13.2)} samghātasya samāsasañjñā prāpnoti . \\
\noindent \textbf{(P\_1,4.13.2)} taddhitavidhi . \\
\noindent \textbf{(P\_1,4.13.2)} devadattaḥ gārgyāyaṇaḥ . \\
\noindent \textbf{(P\_1,4.13.2)} samghātāt taddhitotpattiḥ prāpnoti . \\
\noindent \textbf{(P\_1,4.13.2)} svara . \\
\noindent \textbf{(P\_1,4.13.2)} devadattaḥ gārgyaḥ . \\
\noindent \textbf{(P\_1,4.13.2)} samghātsya ñnityādiḥ nityam iti ādyudāttatvam prāpnoti . \\
\noindent \textbf{(P\_1,4.13.2)} pratyayagrahaṇe yasmāt sa tadādeḥ tadantasya grahaṇam bhavati iti na doṣaḥ bhavati . \\
\noindent \textbf{(P\_1,4.13.2)} sā tarhi eṣā paribhāṣā kartavyā . \\
\noindent \textbf{(P\_1,4.13.2)} na kartavyā . \\
\noindent \textbf{(P\_1,4.13.2)} evam vakṣyāmi : yasmāt pratyayavidhiḥ tadādi pratyaye gṛhyamāṇe gṛhyate . \\
\noindent \textbf{(P\_1,4.13.2)} tataḥ aṅgam . \\
\noindent \textbf{(P\_1,4.13.2)} aṅgasañjñam ca bhavati yasmā tpratyayavidhiḥ tadādi pratyaye . \\
\noindent \textbf{(P\_1,4.13.3)} yadi pratyayagrahaṇe yasmāt saḥ tadādeḥ grahaṇam bhavati iti ucyate avatapenakulasthitam te etat udakeviśīrṇam te etat sagatikena sanakulena ca samāsaḥ na prāpnoti . \\
\noindent \textbf{(P\_1,4.13.3)} evam tarhi pratyayagrahaṇe yasmāt saḥ tadādeḥ grahaṇam bhavati iti uktvā tataḥ vakṣyāmi : kṛdgrahaṇe gatikārakapūrvasya api . \\
\noindent \textbf{(P\_1,4.13.3)} kṛdgrahaṇe gatikārakapūrvasya api grahaṇam bhavati ti eṣā paribhāṣā kartavyā . \\
\noindent \textbf{(P\_1,4.13.3)} kāni etasyāḥ paribhāṣāyāḥ prayojanāni . \\
\noindent \textbf{(P\_1,4.13.3)} prayojanam samāsataddhitavidhisvarāḥ . \\
\noindent \textbf{(P\_1,4.13.3)} samāsa . \\
\noindent \textbf{(P\_1,4.13.3)} avatapenakulasthitam te etat udakeviśīrṇam te etat sagatikena sanakulena ca samāsaḥ siddhaḥ bhavati . \\
\noindent \textbf{(P\_1,4.13.3)} samāsa . \\
\noindent \textbf{(P\_1,4.13.3)} taddhitavidhi . \\
\noindent \textbf{(P\_1,4.13.3)} sāṅkūṭinam vyāvakrośī . \\
\noindent \textbf{(P\_1,4.13.3)} samghātāt taddhitopattiḥ siddhā bhavati . \\
\noindent \textbf{(P\_1,4.13.3)} taddhitavidhi . \\
\noindent \textbf{(P\_1,4.13.3)} svara . \\
\noindent \textbf{(P\_1,4.13.3)} dūrāt āgataḥ dūrādāgataḥ iti . \\
\noindent \textbf{(P\_1,4.13.3)} antaḥ thāthaghañktājabitrakāṇām iti eṣaḥ svaraḥ siddhaḥ bhavati . \\
\noindent \textbf{(P\_1,4.13.3)} kṛdgrahaṇe gatikārakapūrvasya api grahaṇam bhavati iti na doṣaḥ bhavati . \\
\noindent \textbf{(P\_1,4.13.3)} sā tarhi eṣā paribhāṣā kartavyā . \\
\noindent \textbf{(P\_1,4.13.3)} na kartavyā . \\
\noindent \textbf{(P\_1,4.13.3)} ācāryapravṛttiḥ jñāpayati bhavati eṣā paribhāṣa iti yat ayam gatiḥ anantaraḥ iti anantaragrahaṇam karoti . \\
\noindent \textbf{(P\_1,4.14)} antagrahaṇam kimartham na suptiṅ padam iti eva ucyate . \\
\noindent \textbf{(P\_1,4.14)} kena idānīm tadantānām bhaviṣyati . \\
\noindent \textbf{(P\_1,4.14)} tadantavidhinā . \\
\noindent \textbf{(P\_1,4.14)} ataḥ uttaram paṭhati . \\
\noindent \textbf{(P\_1,4.14)} padasañjñāyām antavacanam anyatra sañjñāvidhau pratyayagrahaṇe tadantavidhipratiṣedhārtham . \\
\noindent \textbf{(P\_1,4.14)} padasañjñāyām antagrahaṇam kriyate jñāpakārtham . \\
\noindent \textbf{(P\_1,4.14)} kim jñāpyam . \\
\noindent \textbf{(P\_1,4.14)} etat jñāpayati ācāryaḥ anyatra sañjñāvidhau pratyayagrahaṇe tadantavidhiḥ na bhavati iti . \\
\noindent \textbf{(P\_1,4.14)} kim etasya jñāpane prayojanam . \\
\noindent \textbf{(P\_1,4.14)} taraptamau ghaḥ . \\
\noindent \textbf{(P\_1,4.14)} taraptamabantasya ghasañjñā na bhavati . \\
\noindent \textbf{(P\_1,4.14)} kim ca syāt . \\
\noindent \textbf{(P\_1,4.14)} kumārī gauritarā . \\
\noindent \textbf{(P\_1,4.14)} ghādiṣu nadyāḥ hrasvaḥ bhavati iti hrasvatvam prasajyeta . \\
\noindent \textbf{(P\_1,4.14)} yadi etat jñāpyate sanādyantāḥ dhātavaḥ iti antagrahaṇam kartavyam . \\
\noindent \textbf{(P\_1,4.14)} kṛttaddhitasamāsāḥ ca iti antagrahaṇam kartavyam . \\
\noindent \textbf{(P\_1,4.14)} idam tṛtīyam jñāpakārtham . \\
\noindent \textbf{(P\_1,4.14)} dve tāvat kriyete nyāse eva . \\
\noindent \textbf{(P\_1,4.14)} yat api ucyate kṛttaddhitasamāsāḥ ca iti antagrahaṇam kartavyam iti . \\
\noindent \textbf{(P\_1,4.14)} na kartavyam . \\
\noindent \textbf{(P\_1,4.14)} arthavat iti vartate kṛttaddhāntam ca eva arthavat na kevalāḥ kṛtaḥ taddhitāḥ vā . \\
\noindent \textbf{(P\_1,4.15)} kimartham idam ucyate na subantam padam iti eva siddham . \\
\noindent \textbf{(P\_1,4.15)} niyamārthaḥ ayamārambhaḥ . \\
\noindent \textbf{(P\_1,4.15)} nāntameva kye padasañjñam bhavati na anyat . \\
\noindent \textbf{(P\_1,4.15)} kva mā bhūt . \\
\noindent \textbf{(P\_1,4.15)} vācyati srucyati \\
\noindent \textbf{(P\_1,4.17)} asarvanāmasthāne iti ucyate . \\
\noindent \textbf{(P\_1,4.17)} tatra te rājā takṣā asarvanāmasthāne iti padasañjñāyāḥ pratiṣedhaḥ prasajyeta . \\
\noindent \textbf{(P\_1,4.17)} na apratiṣedhāt . \\
\noindent \textbf{(P\_1,4.17)} na ayam prasajyapratiṣedhaḥ sarvanāmasthāne na iti . \\
\noindent \textbf{(P\_1,4.17)} kim tarhi . \\
\noindent \textbf{(P\_1,4.17)} paryudāsaḥ ayam yat anyat sarvanāmasthānāt iti . \\
\noindent \textbf{(P\_1,4.17)} sarvanāmasthāne avyāpāraḥ . \\
\noindent \textbf{(P\_1,4.17)} yadi kena citprāpnoti tena bhaviṣyati . \\
\noindent \textbf{(P\_1,4.17)} pūrveṇa ca prāpnoti . \\
\noindent \textbf{(P\_1,4.17)} aprāpteḥ vā . \\
\noindent \textbf{(P\_1,4.17)} atha vā anantarā yā prāptiḥ sā pratiṣidhyate . \\
\noindent \textbf{(P\_1,4.17)} kutaḥ etat . \\
\noindent \textbf{(P\_1,4.17)} antarasya vidhiḥ vā bhavati pratiṣedhaḥ vā iti . \\
\noindent \textbf{(P\_1,4.17)} pūrvā prāptiḥ apratiṣiddhā tayā bhaviṣyati . \\
\noindent \textbf{(P\_1,4.17)} nanu ca iyam prāptiḥ pūrvām prāptim bādhate . \\
\noindent \textbf{(P\_1,4.17)} na utsahate pratiṣiddhā satī bādhitum . \\
\noindent \textbf{(P\_1,4.17)} atha vā yogavibhāgaḥ kariṣyate . \\
\noindent \textbf{(P\_1,4.17)} svādiṣu pūrvam padasañjñam bhavati . \\
\noindent \textbf{(P\_1,4.17)} tataḥ sarvanāmasthāne ayacipūrvam padasañjñam bhavati . \\
\noindent \textbf{(P\_1,4.17)} tato bham . \\
\noindent \textbf{(P\_1,4.17)} bhasañjñam ca bhavati yajādau asarvanāmasthāne iti . \\
\noindent \textbf{(P\_1,4.17)} yadi tarhi sau api padam bhavati ecaḥ plutavikāre padāntagrahaṇam coditam iha mā bhūt bhadram karoṣi gauḥ iti tasmin kriyamāṇe api prāpnoti . \\
\noindent \textbf{(P\_1,4.17)} vākyapadayoḥ antyasya iti evam tat . \\
\noindent \textbf{(P\_1,4.17)} bhuvadvadbhyaḥ dhārayadbhyaḥ etayoḥ padasañjñā vaktavyā . \\
\noindent \textbf{(P\_1,4.17)} bhuvadvadbhyaḥ dhārayadvadbhyaḥ . \\
\noindent \textbf{(P\_1,4.18)} bhasañjñāyām uttarapadalope ṣaṣaḥ pratiṣedhaḥ . \\
\noindent \textbf{(P\_1,4.18)} bhasañjñāyām uttarapadalope ṣaṣaḥ pratiṣedhaḥ vaktavyaḥ . \\
\noindent \textbf{(P\_1,4.18)} anukampitaḥ ṣaḍaḍguliḥ ṣaḍikaḥ . \\
\noindent \textbf{(P\_1,4.18)} siddham acaḥ sthānivattvāt . \\
\noindent \textbf{(P\_1,4.18)} siddham etat . \\
\noindent \textbf{(P\_1,4.18)} katham . \\
\noindent \textbf{(P\_1,4.18)} acaḥ sthānivadbhāvāt bhasañjñā na bhaviṣyati . \\
\noindent \textbf{(P\_1,4.18)} iha api tarhi prāpnoti . \\
\noindent \textbf{(P\_1,4.18)} vāgāśīrdattaḥ vācikaḥ iti . \\
\noindent \textbf{(P\_1,4.18)} vakṣyati etat : siddham ekākṣarapūrvapadānām uttarapadalopavacanāt iti . \\
\noindent \textbf{(P\_1,4.18)} iha api tarhi prāpnoti ṣaḍaḍguliḥ ṣaḍikaḥ iti . \\
\noindent \textbf{(P\_1,4.18)} vakṣyati etat : ṣaṣaḥ ṭhājādivacanāt siddham iti . \\
\noindent \textbf{(P\_1,4.18)} nabhoṅgiromanuṣām vati upasamkhyānam . \\
\noindent \textbf{(P\_1,4.18)} nabhoṅgiromanuṣām vati upasamkhyānam kartavyam . \\
\noindent \textbf{(P\_1,4.18)} nabhasvat aṅgirasvat manuṣvat . \\
\noindent \textbf{(P\_1,4.18)} vṛṣaṇ vasvaśvayoḥ . \\
\noindent \textbf{(P\_1,4.18)} vṛṣaṇ iti etasya vasvaśayoḥ bhasañjñā vaktavyā . \\
\noindent \textbf{(P\_1,4.18)} vṛṣaṇvasuḥ vṛṣaṇaśvasya yat śiraḥ vṛṣaṇaśvasya mene . \\
\noindent \textbf{(P\_1,4.19)} arthagrahaṇam kimartham na tasau matau iti eva ucyeta . \\
\noindent \textbf{(P\_1,4.19)} tasau matau iti iyati ucyamāne ihaiva syāt payasvān yaśasvān . \\
\noindent \textbf{(P\_1,4.19)} iha na syāt payasvī yaśasvī . \\
\noindent \textbf{(P\_1,4.19)} arthagrahaṇe punaḥ kriyamāṇe matupi ca siddham bhavati yaḥ ca nyaḥ tena samānārthaḥ tasmin ca . \\
\noindent \textbf{(P\_1,4.19)} yadi arthagrahaṇam kriyate payasvān yaśasvān atra na prāpnoti . \\
\noindent \textbf{(P\_1,4.19)} kim kāraṇam . \\
\noindent \textbf{(P\_1,4.19)} na hi matup matvarthe vartate . \\
\noindent \textbf{(P\_1,4.19)} matup api matvarthe vartate . \\
\noindent \textbf{(P\_1,4.19)} tat yathā devadattaśālāyām brāhmaṇā ānīyantām iti ukte yadi devadattaḥ pi brāhmaṇaḥ bhavati saḥ api ānīyate . \\
\noindent \textbf{(P\_1,4.20)} ubhayasañjñānyapi iti vaktavyam . \\
\noindent \textbf{(P\_1,4.20)} saḥ suṣṭhubhā saḥ ṛkvatā gaṇena . \\
\noindent \textbf{(P\_1,4.21.1)} bahuṣu bahuvacanam iti ucyate . \\
\noindent \textbf{(P\_1,4.21.1)} keṣu bahuṣu . \\
\noindent \textbf{(P\_1,4.21.1)} artheṣu . \\
\noindent \textbf{(P\_1,4.21.1)} yadi evam vṛkṣaḥ plakṣaḥ atra api prāpnoti . \\
\noindent \textbf{(P\_1,4.21.1)} bahavaḥ te arthāḥ mūlam skandhaḥ phalam palāśam iti . \\
\noindent \textbf{(P\_1,4.21.1)} evam tarhi ekavacanam dvivacanam bahuvacanamiti śabdasañjñāḥ etāḥ . \\
\noindent \textbf{(P\_1,4.21.1)} yeṣu artheṣu svādayaḥ vidhīyante teṣu bahuṣu . \\
\noindent \textbf{(P\_1,4.21.1)} keṣu ca artheṣu svādayaḥ vidhīyante . \\
\noindent \textbf{(P\_1,4.21.1)} karmādiṣu . \\
\noindent \textbf{(P\_1,4.21.1)} na vai karmādayaḥ vibhaktyarthāḥ . \\
\noindent \textbf{(P\_1,4.21.1)} ke tarhi . \\
\noindent \textbf{(P\_1,4.21.1)} ekatvādayaḥ . \\
\noindent \textbf{(P\_1,4.21.1)} ekatvādiṣu api vai vibhakyartheṣu avaśyam karmādayaḥ nimittatvena upādeyāḥ . \\
\noindent \textbf{(P\_1,4.21.1)} karmaṇaḥ ekatve karmaṇaḥ dvitve karmaṇaḥ bahutve iti . \\
\noindent \textbf{(P\_1,4.21.1)} saḥ tarhi tathā nirdeśaḥ kartavyaḥ . \\
\noindent \textbf{(P\_1,4.21.1)} na hi antareṇa bhāvapratyayam guṇapradhānaḥ bhavati nirdeśaḥ . \\
\noindent \textbf{(P\_1,4.21.1)} iha ca : iti eke manyante , tat eke manyante iti paratvāt ekavacanam prāpnoti . \\
\noindent \textbf{(P\_1,4.21.1)} bahuṣu bahuvacanam iti eṣaḥ yogaḥ paraḥ kariṣyate . \\
\noindent \textbf{(P\_1,4.21.1)} sūtraviparyāsaḥ kṛtaḥ bhavati . \\
\noindent \textbf{(P\_1,4.21.1)} iha ca : bahuḥ odanaḥ , bahuḥ sūpaḥ iti paratvāt bahuvacanam prāpnoti . \\
\noindent \textbf{(P\_1,4.21.1)} na eṣaḥ doṣaḥ . \\
\noindent \textbf{(P\_1,4.21.1)} yat tāvat ucyate na hi antareṇa bhāvaprayayam guṇapradhānaḥ bhavati nirdeśaḥ iti tan na . \\
\noindent \textbf{(P\_1,4.21.1)} antareṇa api bhāvapratyayam guṇapradhānaḥ bhavati nirdeśaḥ . \\
\noindent \textbf{(P\_1,4.21.1)} katham . \\
\noindent \textbf{(P\_1,4.21.1)} iha kadā cit guṇaḥ guṇiviśeṣakaḥ bhavati . \\
\noindent \textbf{(P\_1,4.21.1)} tat yathā paṭaḥ śuklaḥ iti . \\
\noindent \textbf{(P\_1,4.21.1)} kadā cit ca guṇinā guṇaḥ vypadiśyate . \\
\noindent \textbf{(P\_1,4.21.1)} paṭhasya śuklaḥ iti . \\
\noindent \textbf{(P\_1,4.21.1)} tat yadā tāvat guṇaḥ guṇiviśeṣakaḥ bhavati paṭaḥ śuklaḥ iti tadā sāmānādhikaraṇyam guṇaguṇinoḥ . \\
\noindent \textbf{(P\_1,4.21.1)} tadā na antareṇa bhāvapratyayam guṇapradhānaḥ bhavati nirdeśaḥ . \\
\noindent \textbf{(P\_1,4.21.1)} yadā tu guṇinā guṇaḥ vyapadiśyate paṭasya śuklaḥ iti svapradhānaḥ tadā guṇaḥ bhavati . \\
\noindent \textbf{(P\_1,4.21.1)} tadā dravye ṣaṣṭhī . \\
\noindent \textbf{(P\_1,4.21.1)} tadā antareṇa bhāvapratyayam guaṇapradhānaḥ bhavati nirdeśaḥ . \\
\noindent \textbf{(P\_1,4.21.1)} na ca iha vayam ekatvādibhiḥ karmādīn viśeṣayiṣyāmaḥ . \\
\noindent \textbf{(P\_1,4.21.1)} kim tarhi . \\
\noindent \textbf{(P\_1,4.21.1)} karmādibhiḥ ekatvādīn viśeṣayiṣyāmaḥ . \\
\noindent \textbf{(P\_1,4.21.1)} katham . \\
\noindent \textbf{(P\_1,4.21.1)} ekasmin ekavacanam . \\
\noindent \textbf{(P\_1,4.21.1)} kasyaikasmin . \\
\noindent \textbf{(P\_1,4.21.1)} karmaṇaḥ . \\
\noindent \textbf{(P\_1,4.21.1)} dvayoḥ dvivacanam . \\
\noindent \textbf{(P\_1,4.21.1)} kayoḥ dvayoḥ . \\
\noindent \textbf{(P\_1,4.21.1)} karmaṇoḥ . \\
\noindent \textbf{(P\_1,4.21.1)} bahuṣu bahuvacanam . \\
\noindent \textbf{(P\_1,4.21.1)} keṣām bahuṣu . \\
\noindent \textbf{(P\_1,4.21.1)} karmaṇām iti . \\
\noindent \textbf{(P\_1,4.21.1)} katham bahuṣu bahuvacanamiti . \\
\noindent \textbf{(P\_1,4.21.1)} etat eva jñāpayati ācāryaḥ nānādhikaraṇavācī yaḥ bahuśabdaḥ tasya idam grahaṇam na vaipulyavācinaḥ iti . \\
\noindent \textbf{(P\_1,4.21.1)} kim etasya jñāpane prayojanam . \\
\noindent \textbf{(P\_1,4.21.1)} yat uktam bahuḥ odanaḥ bahuḥ sūpaḥ iti paratvāt bahuvacanam prāpnoti iti sa doṣaḥ na bhavati . \\
\noindent \textbf{(P\_1,4.21.1)} yat apyucyate iti eke manyante tat eke manyanta iti paratvāt ekavacanam prāpnoti iti na eṣaḥ doṣaḥ . \\
\noindent \textbf{(P\_1,4.21.1)} ekaśabdaḥ ayam bahvarthaḥ . \\
\noindent \textbf{(P\_1,4.21.1)} asti eva samkhyāvācī . \\
\noindent \textbf{(P\_1,4.21.1)} tat yathā ekaḥ dvau bahavaḥ iti . \\
\noindent \textbf{(P\_1,4.21.1)} asti asahāyavācī . \\
\noindent \textbf{(P\_1,4.21.1)} tat yathā ekāgnayaḥ ekahalāni ekākibhiḥ kṣudrakaiḥ jitam iti . \\
\noindent \textbf{(P\_1,4.21.1)} asti anyārthe vartate . \\
\noindent \textbf{(P\_1,4.21.1)} tat yathā sadhamādaḥ dyumnaḥ ekāḥ tāḥ . \\
\noindent \textbf{(P\_1,4.21.1)} anyāḥ iti arthaḥ . \\
\noindent \textbf{(P\_1,4.21.1)} tat yaḥ anyārthe vartate tasya eṣaḥ prayogaḥ . \\
\noindent \textbf{(P\_1,4.21.2)} kimartham punaḥ idam ucyate . \\
\noindent \textbf{(P\_1,4.21.2)} suptiṅām aviśeṣavidhānāt dṛṣtaviprayogatvāt ca niyamārtham vacanam . supaḥ aviśeṣeṇa prātipadikamātrāt vidhīyante . \\
\noindent \textbf{(P\_1,4.21.2)} tiṅaḥ aviśeṣeṇa dhātumātrāt vidhīyante . \\
\noindent \textbf{(P\_1,4.21.2)} tatra etat syāt yadi apyaviśeṣeṇa vidhīyante na eva viprayogaḥ lakṣyate iti . \\
\noindent \textbf{(P\_1,4.21.2)} dṛṣṭaviprayogatvāt ca . \\
\noindent \textbf{(P\_1,4.21.2)} dṛśyate khalu api viprayogaḥ . \\
\noindent \textbf{(P\_1,4.21.2)} tadyathā : akṣīṇi me darśanīyāni , pādāḥ me sukumārāḥ iti . \\
\noindent \textbf{(P\_1,4.21.2)} suptiṅoḥ aviśeṣavidhānāt dṛṣṭaviprayogatvāt ca vyatikaraḥ prāpnoti . \\
\noindent \textbf{(P\_1,4.21.2)} iṣyate ca avyatikaraḥ syāt iti . \\
\noindent \textbf{(P\_1,4.21.2)} tat ca antareṇa yatnam na sidhyati iti niyamārtham vacanam . \\
\noindent \textbf{(P\_1,4.21.2)} evamartham idam ucyate . \\
\noindent \textbf{(P\_1,4.21.2)} atha etasmin niyamārthe sati kim punaḥ ayam pratyayaniyamaḥ : ekasmin eva ekavacanam , dvayoḥ eva dvivacanam , bahuṣu eva bahuvacanam iti . \\
\noindent \textbf{(P\_1,4.21.2)} āhosvit arthaniyamaḥ : ekasmin ekavacanam eva , dvayoḥ dvivacanam eva , bahuṣu bahuvacanam eva iti . \\
\noindent \textbf{(P\_1,4.21.2)} kaḥ ca atra viśeṣaḥ . \\
\noindent \textbf{(P\_1,4.21.2)} tatra pratyayaniyame avyayānām padasañjñābhāvaḥ asubantatvāt . \\
\noindent \textbf{(P\_1,4.21.2)} tatra pratyayaniyame avyayānām padasañjñā na prāpnoti . \\
\noindent \textbf{(P\_1,4.21.2)} uccaiḥ nīcaiḥ iti . \\
\noindent \textbf{(P\_1,4.21.2)} kim kāraṇam . \\
\noindent \textbf{(P\_1,4.21.2)} asubantatvāt . \\
\noindent \textbf{(P\_1,4.21.2)} arthaniyame siddham . \\
\noindent \textbf{(P\_1,4.21.2)} arthaniyame siddham bhavati . \\
\noindent \textbf{(P\_1,4.21.2)} astu arthaniyamaḥ . \\
\noindent \textbf{(P\_1,4.21.2)} atha vā punaḥ astu pratyayaniyamaḥ . \\
\noindent \textbf{(P\_1,4.21.2)} nanu ca uktam : tatra pratyayaniyame avyayānām padasañjñābhāvaḥ asubantatvāt iti . \\
\noindent \textbf{(P\_1,4.21.2)} na eṣaḥ doṣaḥ . \\
\noindent \textbf{(P\_1,4.21.2)} supām karmādayaḥ api arthāḥ saṅkhyā ca eva tathā tiṅām . supām saṅkhyā ca eva arthaḥ karmādayaḥ ca . \\
\noindent \textbf{(P\_1,4.21.2)} tathā tiṅām . \\
\noindent \textbf{(P\_1,4.21.2)} prasiddhaḥ niyamaḥ tatra . \\
\noindent \textbf{(P\_1,4.21.2)} prasiddhaḥ tatra niyamaḥ . \\
\noindent \textbf{(P\_1,4.21.2)} niyamaḥ prakṛteṣu vā . \\
\noindent \textbf{(P\_1,4.21.2)} atha vā prakṛtān arthān apekṣya niyamaḥ . \\
\noindent \textbf{(P\_1,4.21.2)} ke ca prakṛtāḥ . \\
\noindent \textbf{(P\_1,4.21.2)} ekatvādayaḥ . \\
\noindent \textbf{(P\_1,4.21.2)} ekasmin eva ekavacanam na dvayoḥ na bahuṣu . \\
\noindent \textbf{(P\_1,4.21.2)} dvayoḥ eva dvivacanam naikasmin na bahuṣu . \\
\noindent \textbf{(P\_1,4.21.2)} bahuṣu eva bahuvacanam na dvayoḥ na ekasmin iti . \\
\noindent \textbf{(P\_1,4.21.2)} atha vā ācāryapravṛttiḥ jñāpayati utpadyante avyayebhyaḥ svādayaḥ iti yat ayam avyayāt āpsupaḥ iti avyayāt lukam śāsti . \\
\noindent \textbf{(P\_1,4.23.1)} kim idam kārake iti . \\
\noindent \textbf{(P\_1,4.23.1)} sañjñānirdeśaḥ . \\
\noindent \textbf{(P\_1,4.23.1)} kim vaktavyam etat . \\
\noindent \textbf{(P\_1,4.23.1)} na hi . \\
\noindent \textbf{(P\_1,4.23.1)} katham anucyamānam gaṃsyate . \\
\noindent \textbf{(P\_1,4.23.1)} iha hi vyākaraṇe ye vā ete loke pratītapadārthakāḥ śabdāḥ taiḥ nirdeśāḥ kriyante paśuḥ apatyam devatā iti yāḥ vā etāḥ kṛtrimāḥ ṭighughabhasañjñāḥ tābhiḥ . \\
\noindent \textbf{(P\_1,4.23.1)} na ca ayam loke dhruvādīnām pratītapadārthakaḥ śabdaḥ na khalu api kṛtrimā sañjñā anyatra avidhānāt . \\
\noindent \textbf{(P\_1,4.23.1)} sañjñādhikāraḥ ca ayam . \\
\noindent \textbf{(P\_1,4.23.1)} tatra kim anyat śakyam vijñātum anyat ataḥ sañjñāyāḥ . \\
\noindent \textbf{(P\_1,4.23.1)} kārake iti sañjñānirdeśaḥ cetsañjñinaḥ api nirdeśaḥ . \\
\noindent \textbf{(P\_1,4.23.1)} kārake iti sañjñānirdeśaḥ cetsañjñinaḥ api nirdeśaḥ kartavyaḥ . \\
\noindent \textbf{(P\_1,4.23.1)} sādhakam nirvartakam kārakasañjñam bhavati iti vaktavyam . \\
\noindent \textbf{(P\_1,4.23.1)} itarathā hi aniṣṭaprasaṅgaḥ grāmasya samīpāt āgacchati iti akārakasya . itarathā hi aniṣṭam prasajyeta . \\
\noindent \textbf{(P\_1,4.23.1)} akārakasya api apādādanasañjñā prasajyeta . \\
\noindent \textbf{(P\_1,4.23.1)} kva . \\
\noindent \textbf{(P\_1,4.23.1)} grāmasya samīpāt āgacchati iti . \\
\noindent \textbf{(P\_1,4.23.1)} na eṣaḥ doṣaḥ . \\
\noindent \textbf{(P\_1,4.23.1)} na atra grāmaḥ apāyayuktaḥ . \\
\noindent \textbf{(P\_1,4.23.1)} kim tarhi . \\
\noindent \textbf{(P\_1,4.23.1)} samīpam . \\
\noindent \textbf{(P\_1,4.23.1)} yadā ca grāmaḥ apāyayuktaḥ bhavati bhavati tadā apādānasañjñā . \\
\noindent \textbf{(P\_1,4.23.1)} tat yathā grāmāt āgacchati iti . \\
\noindent \textbf{(P\_1,4.23.1)} karmasañjñāprasaṅgaḥ akathitasya brāhmaṇasya putram panthānam pṛcchati iti . \\
\noindent \textbf{(P\_1,4.23.1)} karmasañjñā ca prāpnoti akathitasya. kva . \\
\noindent \textbf{(P\_1,4.23.1)} brāhmaṇasya putram panthānam pṛcchati iti . \\
\noindent \textbf{(P\_1,4.23.1)} na eṣaḥ doṣaḥ . \\
\noindent \textbf{(P\_1,4.23.1)} ayam akathitaśabdaḥ asti eva saṅkīrtite vartate . \\
\noindent \textbf{(P\_1,4.23.1)} tat yathā kaḥ cit kam cit sañcakṣya āha asau atra akathitaḥ . \\
\noindent \textbf{(P\_1,4.23.1)} asaṃkīrtitaḥ iti gamyate . \\
\noindent \textbf{(P\_1,4.23.1)} asti aprādhānye vartate . \\
\noindent \textbf{(P\_1,4.23.1)} tat yathā akathitaḥ asau grāme akathitaḥ asau nagare iti ucyate yaḥ yatra apradhānaḥ bhavati . \\
\noindent \textbf{(P\_1,4.23.1)} tat yadā aprādhānye akathitśabdaḥ vartate tadā eṣaḥ doṣaḥ karmasañjñāprasaṅgaḥ akathitasya brāhmaṇasya putram panthānam pṛcchati iti . \\
\noindent \textbf{(P\_1,4.23.1)} apādānam ca vṛkṣasya parṇam patati iti . \\
\noindent \textbf{(P\_1,4.23.1)} apādānasañjñā ca prāpnoti . \\
\noindent \textbf{(P\_1,4.23.1)} kva . \\
\noindent \textbf{(P\_1,4.23.1)} vṛkṣasya parṇam patati . \\
\noindent \textbf{(P\_1,4.23.1)} kuḍyasya piṇḍaḥ patati iti . \\
\noindent \textbf{(P\_1,4.23.1)} nā vā apāyasyāvivakṣitatvāt . \\
\noindent \textbf{(P\_1,4.23.1)} na vā eṣaḥ doṣaḥ . \\
\noindent \textbf{(P\_1,4.23.1)} kim kāraṇam . \\
\noindent \textbf{(P\_1,4.23.1)} apāyasya avivakṣitatvāt . \\
\noindent \textbf{(P\_1,4.23.1)} na atra apāyaḥ vivakṣitaḥ . \\
\noindent \textbf{(P\_1,4.23.1)} kim tarhi . \\
\noindent \textbf{(P\_1,4.23.1)} sambandhaḥ . \\
\noindent \textbf{(P\_1,4.23.1)} yadā ca apāyaḥ vivakṣitaḥ bhavati bhavati tadā apādānasañjñā . \\
\noindent \textbf{(P\_1,4.23.1)} tat yathā . \\
\noindent \textbf{(P\_1,4.23.1)} vṛkṣāt parṇam patati iti . \\
\noindent \textbf{(P\_1,4.23.1)} sambandhastu tadā na vivakṣitaḥ bhavati . \\
\noindent \textbf{(P\_1,4.23.1)} na jñāyate kaṅkasya vā kurarasya vā iti . \\
\noindent \textbf{(P\_1,4.23.2)} ayam tarhi doṣaḥ karmasañjñāprasaṅgaḥ ca akathitasya brāhmaṇasya putram panthānam pṛcchati iti . \\
\noindent \textbf{(P\_1,4.23.2)} na eṣaḥ doṣaḥ . \\
\noindent \textbf{(P\_1,4.23.2)} kārakaḥ iti mahatī sañjñā kriyate . \\
\noindent \textbf{(P\_1,4.23.2)} sañjñā ca nāma yataḥ na laghīyaḥ . \\
\noindent \textbf{(P\_1,4.23.2)} kutaḥ etat . \\
\noindent \textbf{(P\_1,4.23.2)} laghvartham hi sañjñākāraṇam . \\
\noindent \textbf{(P\_1,4.23.2)} tatra mahatyāḥ sañjñāyāḥ karaṇe etatprayojanam anvarthasañjñā yathā vijñāyeta . \\
\noindent \textbf{(P\_1,4.23.2)} karoti iti kārakam iti . \\
\noindent \textbf{(P\_1,4.23.2)} anvartham iti cet akartari kartṛśabdānupapattiḥ . anvartham iti cet akartari kartṛśabdaḥ na upapadyate . \\
\noindent \textbf{(P\_1,4.23.2)} karaṇam kārakam adhikaraṇam kārakam iti . \\
\noindent \textbf{(P\_1,4.23.2)} siddham tu pratikārakam kriyābhedāt pacādīnām karaṇādhikaraṇayoḥ kartṛbhāvaḥ . \\
\noindent \textbf{(P\_1,4.23.2)} siddhaḥ karaṇādhikaraṇayoḥ kartṛbhāvaḥ . \\
\noindent \textbf{(P\_1,4.23.2)} kutaḥ . \\
\noindent \textbf{(P\_1,4.23.2)} pratikārakam kriyābhedāt pacādīnām . \\
\noindent \textbf{(P\_1,4.23.2)} pacādīnām hi pratikārakam kriyā bhidyate . \\
\noindent \textbf{(P\_1,4.23.2)} kim idam pratikārakam iti . \\
\noindent \textbf{(P\_1,4.23.2)} kārakam kārakam prati pratikārakam . \\
\noindent \textbf{(P\_1,4.23.2)} kaḥ asau pratikārakam kriyābhedaḥ pacādīnām . \\
\noindent \textbf{(P\_1,4.23.2)} adhiśrayaṇodakāsecanataṇḍulāvapanaidhopadakarṣaṇikriyāḥ pradhānasya kartuḥ pākaḥ . \\
\noindent \textbf{(P\_1,4.23.2)} adhiśrayaṇodakāsecanataṇḍulāvapanaidhopakarṣaṇādikriyāḥ kurvan eva devadattaḥ pacati iti ucyate . \\
\noindent \textbf{(P\_1,4.23.2)} tatra tadā paciḥ vartate . \\
\noindent \textbf{(P\_1,4.23.2)} eṣaḥ pradhānakartuḥ pākaḥ . \\
\noindent \textbf{(P\_1,4.23.2)} etat pradhānakartuḥ kartṛtvam . \\
\noindent \textbf{(P\_1,4.23.2)} droṇam pacati āḍhakam pacati ti sambhavanakriyā dhāraṇakriyā ca adhikaraṇasya pākaḥ . \\
\noindent \textbf{(P\_1,4.23.2)} droṇam pacati āḍhakam pacati iti sambhavanakriyām dhāraṇakriyām ca kurvatī sthālī pacati iti ucatye . \\
\noindent \textbf{(P\_1,4.23.2)} tatra tadā paciḥ vartate . \\
\noindent \textbf{(P\_1,4.23.2)} eṣaḥ dhikaraṇasya pākaḥ . \\
\noindent \textbf{(P\_1,4.23.2)} etat adhikaraṇasya kartṛtvam . \\
\noindent \textbf{(P\_1,4.23.2)} edhāḥ pakṣyanti ā viklitteḥ jvaliṣyanti iti jvalvanakriyā karaṇasya pākaḥ . \\
\noindent \textbf{(P\_1,4.23.2)} edhāḥ pakṣyanti ā viklitteḥ jvaliṣyanti iti kurvanti kāṣṭhāni pacanti iti ucyante . \\
\noindent \textbf{(P\_1,4.23.2)} tatra tadā paciḥ vartate . \\
\noindent \textbf{(P\_1,4.23.2)} eṣaḥ karaṇasya pākaḥ .etat karaṇasya kartṛtvam . \\
\noindent \textbf{(P\_1,4.23.2)} udyamananipātanāni kartuḥ chidikriyā . \\
\noindent \textbf{(P\_1,4.23.2)} udyamananipātanāni kurvan devadattaḥ chinatti iti ucyate . \\
\noindent \textbf{(P\_1,4.23.2)} tatra tadā chidiḥ vartate . \\
\noindent \textbf{(P\_1,4.23.2)} eṣaḥ pradhānakartuḥ chedaḥ . \\
\noindent \textbf{(P\_1,4.23.2)} etat pradhānakartuḥ kartṛtvam . \\
\noindent \textbf{(P\_1,4.23.2)} yat tat na tṛṇena tatparaśoḥ chedanam . yat tat samāne udyamane nipātane ca paraśunā chidyate na tṛṇena tat paraśoḥ chedanam . \\
\noindent \textbf{(P\_1,4.23.2)} avaśyam ca etat evam vijñeyam . \\
\noindent \textbf{(P\_1,4.23.2)} itarathā hi asitṛṇayoḥ chedane aviśeṣaḥ syāt . \\
\noindent \textbf{(P\_1,4.23.2)} yaḥ hi manyata udyamananipātanāt eva etat bhavati chinatti iti asitṛṇayoḥ chedane na tasya viśeṣaḥ syāt . \\
\noindent \textbf{(P\_1,4.23.2)} yat asinā chidyate tṛṇena api tat chidyeta . \\
\noindent \textbf{(P\_1,4.23.2)} apādānādīnām tu aprasiddhiḥ . \\
\noindent \textbf{(P\_1,4.23.2)} apādānādīnām kartṛtvasya aprasiddhiḥ . \\
\noindent \textbf{(P\_1,4.23.2)} yathā hi bhavatā karaṇādhikaraṇayoḥ kartṛtvam nirdarśitam na tathā apādānādīnām kartṛtvam nidarśyate . \\
\noindent \textbf{(P\_1,4.23.2)} na vā svatantraparatantratvāt tayoḥ paryāyeṇa vacanam vacanāśrayā ca sañjñā . \\
\noindent \textbf{(P\_1,4.23.2)} na vā eṣaḥ doṣaḥ . \\
\noindent \textbf{(P\_1,4.23.2)} kim kāraṇam . \\
\noindent \textbf{(P\_1,4.23.2)} svatantraparatantratvāt . \\
\noindent \textbf{(P\_1,4.23.2)} sarvatra eva atra svātantryam pāratantryam ca vivakṣitam . \\
\noindent \textbf{(P\_1,4.23.2)} tayoḥ paryāyeṇa vacanam . \\
\noindent \textbf{(P\_1,4.23.2)} tayoḥ svātantryapāratantryayoḥ paryāyeṇa vacanam bhaviṣyati . \\
\noindent \textbf{(P\_1,4.23.2)} vacanāśrayā ca sañjñā bhaviṣyati . \\
\noindent \textbf{(P\_1,4.23.2)} tat yathā . \\
\noindent \textbf{(P\_1,4.23.2)} balāhakāt vidyotate . \\
\noindent \textbf{(P\_1,4.23.2)} balāhake vidyotate . \\
\noindent \textbf{(P\_1,4.23.2)} balāhakaḥ vidyotate iti . \\
\noindent \textbf{(P\_1,4.23.2)} kim tarhi ucyate apādānādīnām tu aprasiddhiḥ iti . \\
\noindent \textbf{(P\_1,4.23.2)} evam tarhi na brūmaḥ apādānādīnām kartṛtvasya aprasiddhiḥ iti . \\
\noindent \textbf{(P\_1,4.23.2)} paryāptam karaṇādhikaraṇayoḥ kartṛtvam nidarśitam apādānīnām kartṛtvanirdarśanāya . \\
\noindent \textbf{(P\_1,4.23.2)} paryāptaḥ hyekaḥ pulākaḥ sthālyāḥ nirdarśanāya . \\
\noindent \textbf{(P\_1,4.23.2)} kim tarhi . \\
\noindent \textbf{(P\_1,4.23.2)} sañjñāyāḥ aprasiddhiḥ . \\
\noindent \textbf{(P\_1,4.23.2)} yāvatā sarvatra eva atra svātantryam vidyate pāratantryam ca tatra paratvāt kartṛsañjñā eva prāpnoti . \\
\noindent \textbf{(P\_1,4.23.2)} atra api na vā svatantryaparatantratvāt tayoḥ paryāyeṇa vacanam vacanāśrayā ca sañjñā iti eva . \\
\noindent \textbf{(P\_1,4.23.2)} yathā punaḥ idam sthālyāḥ svātantryam nidarśitam sambhavanakriyām dhāraṇakriyām ca kurvatī sthālī svatantrā iti kva idānīm paratantrā syāt . \\
\noindent \textbf{(P\_1,4.23.2)} yat tat prakṣālanam parivartanam vā . \\
\noindent \textbf{(P\_1,4.23.2)} na vai evamartham sthālī upādiyate prakṣālanam parivartanam ca kariṣyāmi iti . \\
\noindent \textbf{(P\_1,4.23.2)} katham tarhi . \\
\noindent \textbf{(P\_1,4.23.2)} sambhavanakriyām dhāraṇakriyām ca kariṣyati iti . \\
\noindent \textbf{(P\_1,4.23.2)} tatra ca sau svatantrā . \\
\noindent \textbf{(P\_1,4.23.2)} kva idānīm paratantrā . \\
\noindent \textbf{(P\_1,4.23.2)} evam tarhi sthālīsthe yatne kathyamāne sthālī svatantrā kartṛsthe yatne kathyamāne paratantrā . \\
\noindent \textbf{(P\_1,4.23.2)} nanu ca bhoḥ kartṛsthe api vai yatne kathyamāne sthālī sambhavanakriyām dhāraṇakriyām ca karoti . \\
\noindent \textbf{(P\_1,4.23.2)} tatra asau svatantrā . \\
\noindent \textbf{(P\_1,4.23.2)} kva idānīm paratantrā . \\
\noindent \textbf{(P\_1,4.23.2)} evam tarhi pradhānena samavāye sthālī paratantrā vyavāye svatantrā . \\
\noindent \textbf{(P\_1,4.23.2)} tat yathā amātyādīnām rājñā saha samavāye pāratantryam vyavāye svātantryam . \\
\noindent \textbf{(P\_1,4.23.2)} kim punaḥ pradhānam . \\
\noindent \textbf{(P\_1,4.23.2)} kartā .katham punaḥ jñāyate kartā pradhānam iti . \\
\noindent \textbf{(P\_1,4.23.2)} yat sarveṣu sādhaneṣu samnihiteṣu kartā pravartayitā bhavati . \\
\noindent \textbf{(P\_1,4.23.2)} nanu ca bhoḥ pradhānena api vai samavāye sthālyāḥ anenārthaḥ adhikaraṇam kārakam iti . \\
\noindent \textbf{(P\_1,4.23.2)} na hi kārakam iti anena adhikaraṇatvam uktam adhikaraṇamiti vā kārakatvam . \\
\noindent \textbf{(P\_1,4.23.2)} ubhau ca anyonyaviśeṣakau bhavataḥ . \\
\noindent \textbf{(P\_1,4.23.2)} katham . \\
\noindent \textbf{(P\_1,4.23.2)} ekadravyasamavāyitvāt . \\
\noindent \textbf{(P\_1,4.23.2)} tat yathā gārgyaḥ devadattaḥ iti . \\
\noindent \textbf{(P\_1,4.23.2)} na hi gārgyaḥ iti anena devadattatvam uktam devadattaḥ iti anena vā gārgyatvam . \\
\noindent \textbf{(P\_1,4.23.2)} ubhau ca anyonyaviśeṣakau bhavataḥ ekadravyasamavāyitvāt . \\
\noindent \textbf{(P\_1,4.23.2)} evam tarhi sāmānyabhūtā kriyā vartate tasyāḥ nirvartakam kārakam . \\
\noindent \textbf{(P\_1,4.23.2)} atha vā yāvat brūyāt kriyāyāmiti tāvat kārake iti . \\
\noindent \textbf{(P\_1,4.23.2)} evam ca kṛtvā nirdeśaḥ upapannaḥ bhavati kārake iti . \\
\noindent \textbf{(P\_1,4.23.2)} itarathā hi kārakeṣu iti brūyāt . \\
\noindent \textbf{(P\_1,4.24.1)} dhruvam iti kimartham . \\
\noindent \textbf{(P\_1,4.24.1)} grāmāt āgacchati śakaṭena . \\
\noindent \textbf{(P\_1,4.24.1)} na etat asti . \\
\noindent \textbf{(P\_1,4.24.1)} karaṇasañjñā atra bādhikā bhaviṣyati . \\
\noindent \textbf{(P\_1,4.24.1)} idam tarhi grāmāt āgacchan kaṃsapātryām pāṇinā odanam bhuṅkte iti . \\
\noindent \textbf{(P\_1,4.24.1)} atra api adhikaraṇasañjñā bādhikā bhaviṣyati . \\
\noindent \textbf{(P\_1,4.24.1)} idam tarhi vṛkṣasya parṇam patati . \\
\noindent \textbf{(P\_1,4.24.1)} kuḍyasya piṇḍaḥ patiti iti . \\
\noindent \textbf{(P\_1,4.24.1)} jugupsāvirāmapramādārthānām upasamkhyānam . jugupsāvirāmapramādārthānām upasamkhyānam kartavyam . \\
\noindent \textbf{(P\_1,4.24.1)} jugupsā . \\
\noindent \textbf{(P\_1,4.24.1)} adharmāt jugupsate. adharmāt bībhatsate . \\
\noindent \textbf{(P\_1,4.24.1)} virāma | dharmāt viramati. dharmāt nivartate . \\
\noindent \textbf{(P\_1,4.24.1)} pramāda . \\
\noindent \textbf{(P\_1,4.24.1)} dharmāt pramādyati . \\
\noindent \textbf{(P\_1,4.24.1)} dharmāt muhyati . \\
\noindent \textbf{(P\_1,4.24.1)} iha ca upasamkhyānam kartavyam . \\
\noindent \textbf{(P\_1,4.24.1)} sāmkāśyakebhyaḥ pāṭaliputrakāḥ abhirūpatarāḥ iti . \\
\noindent \textbf{(P\_1,4.24.1)} tat tarhi idam vaktavyam . \\
\noindent \textbf{(P\_1,4.24.1)} na vaktavyam . \\
\noindent \textbf{(P\_1,4.24.1)} iha tāvat adharmāt jugupsate adharmāt bībhatsate iti . \\
\noindent \textbf{(P\_1,4.24.1)} yaḥ eṣaḥ manuṣyaḥ prekṣāpūrvakārī bhavati saḥ paśyati duḥkhaḥ adharmaḥ na anena kṛtyam asti iti . \\
\noindent \textbf{(P\_1,4.24.1)} saḥ buddhyā samprāpya nivartate . \\
\noindent \textbf{(P\_1,4.24.1)} tatra dhruvamapāye apādānam iti eva siddham . \\
\noindent \textbf{(P\_1,4.24.1)} iha ca dharmāt viramati dharmāt nivartate iti dharmāt pramādyati dharmāt muhyati iti . \\
\noindent \textbf{(P\_1,4.24.1)} yaḥ eṣaḥ manuṣyaḥ sambhinnabuddhiḥ bhavati saḥ paśyati na idam kim cit dharmaḥ nāma na enam kariṣyāmi iti . \\
\noindent \textbf{(P\_1,4.24.1)} saḥ buddhyā samprāpya nivartate . \\
\noindent \textbf{(P\_1,4.24.1)} tatra dhruvamapāye apādānam iti eva siddham . \\
\noindent \textbf{(P\_1,4.24.1)} iha ca sāmkāśyakebhyaḥ pāṭaliputrakāḥ abhirūpatarāḥ iti . \\
\noindent \textbf{(P\_1,4.24.1)} yaḥ taiḥ sāmyam gatavān bhavati saḥ etatprayuṅkte . \\
\noindent \textbf{(P\_1,4.24.1)} gatiyukteṣu apādānasañjñā na papadyate adhruvatvāt . gatiyukteṣu apādānasañjñā na upapadyate . \\
\noindent \textbf{(P\_1,4.24.1)} aśvāt trastāt patitaḥ . \\
\noindent \textbf{(P\_1,4.24.1)} rathāt pravītāt patitaḥ . \\
\noindent \textbf{(P\_1,4.24.1)} sārthāt gacchataḥ hīnaḥ iti . \\
\noindent \textbf{(P\_1,4.24.1)} kim kāraṇam . \\
\noindent \textbf{(P\_1,4.24.1)} adhruvatvāt . \\
\noindent \textbf{(P\_1,4.24.1)} na vā adhrauvyasya avivikṣitatvāt . \\
\noindent \textbf{(P\_1,4.24.1)} na vā eṣaḥ doṣaḥ . \\
\noindent \textbf{(P\_1,4.24.1)} kim kāraṇam . \\
\noindent \textbf{(P\_1,4.24.1)} adhrauvyasya avivakṣitatvāt . \\
\noindent \textbf{(P\_1,4.24.1)} na atra adhrauvyam vivakṣitam . \\
\noindent \textbf{(P\_1,4.24.1)} kim tarhi . \\
\noindent \textbf{(P\_1,4.24.1)} dhrauvyam . \\
\noindent \textbf{(P\_1,4.24.1)} iha tāvat aśvāt trastāt patitaḥ iti . \\
\noindent \textbf{(P\_1,4.24.1)} yat tadaśve aśvatvam āśugāmitvam tat dhruvam tat ca vivakṣitam . \\
\noindent \textbf{(P\_1,4.24.1)} rathāt pravītāt patitaḥ iti yat tat rathe rathatvam ramante asmin rathaḥ iti tat dhruvam tat ca vivakṣitam . \\
\noindent \textbf{(P\_1,4.24.1)} sārthāt gacchataḥ hīnaḥ iti yat tatsārthe sārthatvam sahārthībhāvaḥ tat dhruvam tat ca vivakṣiktam . \\
\noindent \textbf{(P\_1,4.24.1)} yadi api tāvat atra etat śakyate vaktum ye tu ete atyantagatiyuktāḥ tatra katham . \\
\noindent \textbf{(P\_1,4.24.1)} dhāvataḥ patitaḥ . \\
\noindent \textbf{(P\_1,4.24.1)} tvaramāṇāt patitaḥ iti . \\
\noindent \textbf{(P\_1,4.24.1)} atra api na vā adhrauvyasya avivakṣitatvāt iti eva siddham . \\
\noindent \textbf{(P\_1,4.24.1)} katham punaḥ sataḥ nāma avivakṣā syāt . \\
\noindent \textbf{(P\_1,4.24.1)} sataḥ api avivakṣā bhavati . \\
\noindent \textbf{(P\_1,4.24.1)} tat yathā alomikā eḍakā . \\
\noindent \textbf{(P\_1,4.24.1)} anudarā kanya iti . \\
\noindent \textbf{(P\_1,4.24.1)} asataḥ ca vivakṣā bhavati . \\
\noindent \textbf{(P\_1,4.24.1)} samudraḥ kuṇḍikā . \\
\noindent \textbf{(P\_1,4.24.1)} vindhyaḥ vardhitakam iti \\
\noindent \textbf{(P\_1,4.24.2)} ayam yogaḥ śakyaḥ avaktum . \\
\noindent \textbf{(P\_1,4.24.2)} katham vṛkebhyaḥ bibheti dasyubhyaḥ bibheti caurebhyaḥ trāyate dasyubhyaḥ trāyate iti . \\
\noindent \textbf{(P\_1,4.24.2)} iha tāvat vṛkebhyaḥ bibheti dasyubhyaḥ bibheti iti . \\
\noindent \textbf{(P\_1,4.24.2)} yaḥ eṣaḥ manuṣyaḥ prekṣāpūrvakārī bhavati saḥ paśyati yadi mām vṛkāḥ paśyanti dhruvaḥ me mṛtyuḥ iti . \\
\noindent \textbf{(P\_1,4.24.2)} saḥ buddhyā samprāpya nivartate . \\
\noindent \textbf{(P\_1,4.24.2)} tatra dhruvam apāye apādānam iti eva siddham . \\
\noindent \textbf{(P\_1,4.24.2)} iha caurebhyaḥ trāyate dasyubhyaḥ trāyate iti . \\
\noindent \textbf{(P\_1,4.24.2)} yaḥ eṣaḥ manuṣyaḥ prekṣāpūrvakārī bhavati saḥ paśyati yadi imam paśyanti dhruvam asya vadhabandhaparikleśāḥ iti . \\
\noindent \textbf{(P\_1,4.24.2)} saḥ buddhyā samprāpya nivartate . \\
\noindent \textbf{(P\_1,4.24.2)} tatra dhruvam apāye apādānam iti eva siddham . \\
\noindent \textbf{(P\_1,4.26)} ayam api yogaḥ śakyaḥ avaktum . \\
\noindent \textbf{(P\_1,4.26)} katham adhyayanāt parājayate iti . \\
\noindent \textbf{(P\_1,4.26)} yaḥ eṣaḥ manuṣyaḥ prekṣāpūrvakārī bhavati saḥ paśyati duḥkham adhyayanam durdharam ca guravaḥ ca durupacārāḥ iti . \\
\noindent \textbf{(P\_1,4.26)} saḥ buddhyā samprāpya nivartate . \\
\noindent \textbf{(P\_1,4.26)} tatra dhruvam apāye apādānam iti eva siddham . \\
\noindent \textbf{(P\_1,4.27)} kim udāharaṇam . \\
\noindent \textbf{(P\_1,4.27)} māṣebhyaḥ gāḥ vārayati . \\
\noindent \textbf{(P\_1,4.27)} bhaved yasya māṣāḥ na gāvaḥ tasya māṣāḥ īpsitāḥ syuḥ . \\
\noindent \textbf{(P\_1,4.27)} yasya tu khalu gāvaḥ na māṣāḥ katham tasya māṣāḥ īpsitāḥ syuḥ . \\
\noindent \textbf{(P\_1,4.27)} tasya api māṣāḥ eva īpsitāḥ . \\
\noindent \textbf{(P\_1,4.27)} ātaḥ ca īpsitāḥ yavebhyḥ gāḥ vārayati . \\
\noindent \textbf{(P\_1,4.27)} iha kūpāt andham vārayati iti kūpe apādānasañjñā na prāpnoti . \\
\noindent \textbf{(P\_1,4.27)} na hi tasya kūpaḥ īpsitaḥ . \\
\noindent \textbf{(P\_1,4.27)} kaḥ tarhi . \\
\noindent \textbf{(P\_1,4.27)} andhaḥ . \\
\noindent \textbf{(P\_1,4.27)} tasya api kūpaḥ eva īpsitaḥ . \\
\noindent \textbf{(P\_1,4.27)} paśyati ayam andhaḥ kūpam mā prāpat iti . \\
\noindent \textbf{(P\_1,4.27)} atha vā yathā eva asya anyatra apaśyataḥ īpsā evam kūpe api . \\
\noindent \textbf{(P\_1,4.27)} iha agneḥ māṇavakam vārayati iti māṇavake apādānasañjñā prāpnoti . \\
\noindent \textbf{(P\_1,4.27)} karmasañjñātra bādhikā bhaviṣyati . \\
\noindent \textbf{(P\_1,4.27)} agnau api tarhi bādhikā syāt . \\
\noindent \textbf{(P\_1,4.27)} tasmāt vaktavyam karmaṇaḥ yat īpsitam iti īpsitepsitam iti vā . \\
\noindent \textbf{(P\_1,4.27)} vāraṇartheṣu karmagrahaṇānarthakyam kartuḥ īpsitatamam karma iti vacanāt . \\
\noindent \textbf{(P\_1,4.27)} vāraṇārtheṣu karmagrahaṇam anarthakam . \\
\noindent \textbf{(P\_1,4.27)} kim kāraṇam . \\
\noindent \textbf{(P\_1,4.27)} kartuḥ īpsitatamam karma iti vacanāt . \\
\noindent \textbf{(P\_1,4.27)} kartuḥ īpsitatamam karma iti eva siddham . \\
\noindent \textbf{(P\_1,4.27)} ayam api yogaḥ śakyaḥ avaktum . \\
\noindent \textbf{(P\_1,4.27)} katham māṣebhyaḥ gāḥ vārayati iti . \\
\noindent \textbf{(P\_1,4.27)} paśyati ayam yadi imāḥ gāvaḥ tatra gacchanti dhruvam sasyavināśaḥ sasyavināśe adharmaḥ ca eva rājabhayam ca . \\
\noindent \textbf{(P\_1,4.27)} saḥ buddhyā samprāpya nivartate . \\
\noindent \textbf{(P\_1,4.27)} tatra dhruvam apāye apādānam iti eva siddham . \\
\noindent \textbf{(P\_1,4.28)} ayam api yogaḥ śakyaḥ avaktum . \\
\noindent \textbf{(P\_1,4.28)} katham upādhyāyāt antardhatte iti . \\
\noindent \textbf{(P\_1,4.28)} paśyati ayam yadi mām upādhyāyaḥ paśyati dhruvam preṣaṇam upālambhaḥ vā iti . \\
\noindent \textbf{(P\_1,4.28)} saḥ buddhyā samprāpya nivartate . \\
\noindent \textbf{(P\_1,4.28)} tatra dhruvam apāye apādānam iti eva siddham . \\
\noindent \textbf{(P\_1,4.29)} upayoge iti kimartham . \\
\noindent \textbf{(P\_1,4.29)} naṭasya śṛṇoti . \\
\noindent \textbf{(P\_1,4.29)} granthikasya śṛṇoti . \\
\noindent \textbf{(P\_1,4.29)} upayoge iti ucyamāne api atra prāpnoti . \\
\noindent \textbf{(P\_1,4.29)} eṣaḥ api hi upayogaḥ . \\
\noindent \textbf{(P\_1,4.29)} ātaḥ ca upayogaḥ yat ārambhakāḥ raṅgam gacchanti naṭasya śroṣyāmaḥ , granthikasya śroṣyāmaḥ iti . \\
\noindent \textbf{(P\_1,4.29)} evam tarhi upayoge iti ucyate sarvaḥ ca upayogaḥ . \\
\noindent \textbf{(P\_1,4.29)} tatra prakarṣagatiḥ vijñāsyate : sādhīyaḥ yaḥ upayogaḥ iti . \\
\noindent \textbf{(P\_1,4.29)} kaḥ ca sādhīyaḥ . \\
\noindent \textbf{(P\_1,4.29)} yaḥ granthārthayoḥ . \\
\noindent \textbf{(P\_1,4.29)} atha vā upayogaḥ kaḥ bhavitum arhati . \\
\noindent \textbf{(P\_1,4.29)} yaḥ niyamapūrvakaḥ . \\
\noindent \textbf{(P\_1,4.29)} tat yathā upayuktāḥ māṇavakāḥ iti ucyante ye ete niyamapūrvakam adhītavantaḥ bhavanti . \\
\noindent \textbf{(P\_1,4.29)} kim punaḥ ākhyātā anupayoge kārakam āhosvit akārakam . \\
\noindent \textbf{(P\_1,4.29)} kaḥ ca atra viśeṣaḥ . \\
\noindent \textbf{(P\_1,4.29)} ākhyātā anupayoge kārakam iti cet akathitvāt karmasañjñāprasaṅgaḥ . ākhyātā anupayoge kārakam iti cet akathitvāt karmasañjñā prāpnoti . \\
\noindent \textbf{(P\_1,4.29)} astu tarhi akārakam . \\
\noindent \textbf{(P\_1,4.29)} akārakam iti cet upayogavacanānarthakyam . \\
\noindent \textbf{(P\_1,4.29)} yadi akārakam upayogavacanam anarthakam . \\
\noindent \textbf{(P\_1,4.29)} astu tarhi kārakam . \\
\noindent \textbf{(P\_1,4.29)} nanu ca uktam ākhyātā anupayoge kārakam iti cet akathitatvāt karmasañjñāprasaṅgaḥ iti . \\
\noindent \textbf{(P\_1,4.29)} na eṣaḥ doṣaḥ . \\
\noindent \textbf{(P\_1,4.29)} parigaṇanam tatra kriyate . \\
\noindent \textbf{(P\_1,4.29)} duhiyācirudhipracchibhikṣiciñām iti . \\
\noindent \textbf{(P\_1,4.29)} ayam api yogaḥ śakyaḥ avaktum . \\
\noindent \textbf{(P\_1,4.29)} katham upādhyāyāt adhīte iti . \\
\noindent \textbf{(P\_1,4.29)} apakrāmati tasmāt tadadhyayanam . \\
\noindent \textbf{(P\_1,4.29)} yadi apakrāmati kim na atyantāya apakrāmati . \\
\noindent \textbf{(P\_1,4.29)} sattatatvāt . \\
\noindent \textbf{(P\_1,4.29)} atha vā jyotirvat jñānāni bhavanti . \\
\noindent \textbf{(P\_1,4.30)} ayam api yogaḥ śakyaḥ avaktum . \\
\noindent \textbf{(P\_1,4.30)} katham gomayāt vṛścikaḥ jāyate . \\
\noindent \textbf{(P\_1,4.30)} golomāvilomabhyaḥ durvāḥ jāyante iti . \\
\noindent \textbf{(P\_1,4.30)} apakrāmanti tāḥ tebhyaḥ . \\
\noindent \textbf{(P\_1,4.30)} yadi apakrāmati kim na atyantāya apakrāmati . \\
\noindent \textbf{(P\_1,4.30)} santatatvāt . \\
\noindent \textbf{(P\_1,4.30)} atha vā anyāḥ canyāḥ ca prādurbhavanti . \\
\noindent \textbf{(P\_1,4.31)} ayam api yogaḥ śakyaḥ avaktum . \\
\noindent \textbf{(P\_1,4.31)} katham himavataḥ gaṅgā prabhavati iti . \\
\noindent \textbf{(P\_1,4.31)} apakrāmanti tāḥ tasmāt āpaḥ . \\
\noindent \textbf{(P\_1,4.31)} yadi apakrāmati kim na atyantāya apakrāmati . \\
\noindent \textbf{(P\_1,4.31)} santatatvāt . \\
\noindent \textbf{(P\_1,4.31)} atha vā anyāḥ canyāḥ ca prādurbhavanti . \\
\noindent \textbf{(P\_1,4.32.1)} karmagrahaṇam kimartham . \\
\noindent \textbf{(P\_1,4.32.1)} yam abhipraiti saḥ sampradānam iti iyati ucyamāne karmaṇaḥ eva sampradānasañjñā prasajyeta . \\
\noindent \textbf{(P\_1,4.32.1)} karmagrahaṇe punaḥ kriyamāṇe na doṣaḥ bhavati . \\
\noindent \textbf{(P\_1,4.32.1)} karma nimittatvena āśrīyate . \\
\noindent \textbf{(P\_1,4.32.1)} atha yamsagrahaṇam kimartham . \\
\noindent \textbf{(P\_1,4.32.1)} karmaṇā abhipraiti sampradānam iti iyati ucyamāne abhiprayataḥ eva sampradānasañjñā prasajyeta . \\
\noindent \textbf{(P\_1,4.32.1)} yamsagrahaṇe punaḥ kriyamāṇe na doṣaḥ bhavati . \\
\noindent \textbf{(P\_1,4.32.1)} yamsagrahaṇāt abhiprayataḥ sampradānasañjñā nirbhajyate . \\
\noindent \textbf{(P\_1,4.32.1)} atha abhipragrahaṇam kimartham . \\
\noindent \textbf{(P\_1,4.32.1)} karmaṇā yam eti sa sampradānam iti iyati ucyamāne yam eva sampratyeti tatra eva syāt . \\
\noindent \textbf{(P\_1,4.32.1)} upādhyāyāya gām dadāti iti . \\
\noindent \textbf{(P\_1,4.32.1)} iha na syāt . \\
\noindent \textbf{(P\_1,4.32.1)} upādhyāyāya gām adāt . \\
\noindent \textbf{(P\_1,4.32.1)} upādhyāyāya gām dāsyati iti . \\
\noindent \textbf{(P\_1,4.32.1)} abhipragrahaṇe punaḥ kriyamāṇe na doṣaḥ bhavati . \\
\noindent \textbf{(P\_1,4.32.1)} abhiḥ ābhimukhye vartate praśabdaḥ ādikarmaṇi . \\
\noindent \textbf{(P\_1,4.32.1)} tena yam ca abhipraiti yam ca abhipraiṣyati yam ca abhiprāgād ābhimukhyamātre sarvatra siddham bhavati . \\
\noindent \textbf{(P\_1,4.32.2)} kriyāgrahaṇam api kartavyam . \\
\noindent \textbf{(P\_1,4.32.2)} iha api yathā syāt . \\
\noindent \textbf{(P\_1,4.32.2)} śrāddhāya nigarhate . \\
\noindent \textbf{(P\_1,4.32.2)} yuddhāya sannahyate . \\
\noindent \textbf{(P\_1,4.32.2)} patye śete iti . \\
\noindent \textbf{(P\_1,4.32.2)} tat tarhi vaktavyam . \\
\noindent \textbf{(P\_1,4.32.2)} na vaktavyam . \\
\noindent \textbf{(P\_1,4.32.2)} katham . \\
\noindent \textbf{(P\_1,4.32.2)} kriyām hi loke karma iti upacaranti . \\
\noindent \textbf{(P\_1,4.32.2)} kām kriyām kariṣyasi . \\
\noindent \textbf{(P\_1,4.32.2)} kim karma kariṣyasi iti . \\
\noindent \textbf{(P\_1,4.32.2)} evam api kartavyam . \\
\noindent \textbf{(P\_1,4.32.2)} kṛtrimākṛtrimayoḥ kṛtrime sampratyayaḥ bhavati . \\
\noindent \textbf{(P\_1,4.32.2)} kriyā api kṛtrimam karma . \\
\noindent \textbf{(P\_1,4.32.2)} na sidhyati . \\
\noindent \textbf{(P\_1,4.32.2)} kartuḥ īpsitatamam karma iti ucyate katham ca nāma kriyayā kriyā īpsitatamā syāt . \\
\noindent \textbf{(P\_1,4.32.2)} kriyā api kriyayā īpsitatamā bhavati . \\
\noindent \textbf{(P\_1,4.32.2)} kayā kriyayā . \\
\noindent \textbf{(P\_1,4.32.2)} sandarśanakriyayā vā prārthayatikriyayā vā adhyavasyatikriyayā vā . \\
\noindent \textbf{(P\_1,4.32.2)} iha yaḥ eṣaḥ manuṣyaḥ prekṣāpūrvakārī bhavati saḥ buddhyā tāvat kamcidartham sampaśyati sandṛṣṭe prārthanā prārthanāyām adhavasāyaḥ adhyavasāye ārambhaḥ ārambhe nirvṛttiḥ nirvṛttau phalāvāptiḥ . \\
\noindent \textbf{(P\_1,4.32.2)} evam kriyā api kṛtrimam karma . \\
\noindent \textbf{(P\_1,4.32.2)} evam api karmaṇaḥ karaṇasañjñā vaktavyā sampradānasya ca karmasañjñā . \\
\noindent \textbf{(P\_1,4.32.2)} paśunā rudram yajate . \\
\noindent \textbf{(P\_1,4.32.2)} paśum rudrāya dadāti iti arthaḥ . \\
\noindent \textbf{(P\_1,4.32.2)} agnau kila paśuḥ prakṣipyate tat rudrāya pahriyate iti . \\
\noindent \textbf{(P\_1,4.37)} kimete ekārthāḥ āhosvit nānārthāḥ . \\
\noindent \textbf{(P\_1,4.37)} kim ca ataḥ . \\
\noindent \textbf{(P\_1,4.37)} yadi ekārthāḥ kimartham pṛthak nirdiśyante . \\
\noindent \textbf{(P\_1,4.37)} atha nānārthāḥ katham kupinā śakyante viśeṣayitum . \\
\noindent \textbf{(P\_1,4.37)} evam tarhi nānārthāḥ kupau tu eṣām sāmānyam asti . \\
\noindent \textbf{(P\_1,4.37)} na hi akupitaḥ krudhyate na vā akupitaḥ druhyati na vā akupitaḥ īrṣyati na vā akupitaḥ asūyati . \\
\noindent \textbf{(P\_1,4.42)} tamagrahaṇam kimartham na sādhakam karaṇam iti eva ucyeta ṣādhakam karaṇam iti iyati ucyamāne sarveṣām kārakāṇām karaṇasañjñā prasajyeta . \\
\noindent \textbf{(P\_1,4.42)} sarvāṇi hi kārakāṇi sādhakāni . \\
\noindent \textbf{(P\_1,4.42)} tamagrahaṇe punaḥ kriyamāṇe na doṣaḥ bhavati . \\
\noindent \textbf{(P\_1,4.42)} na etat asti prayojanam . \\
\noindent \textbf{(P\_1,4.42)} pūrvāḥ tāvat sañjñāḥ apavādatvāt bādhikāḥ bhaviṣyanti parāḥ paratvāt ca anavakāśatvāt ca . \\
\noindent \textbf{(P\_1,4.42)} iha tarhi dhanuṣā vidhyati apāyayuktatvāt ca apādānasañjñā sādhakatvāt ca karaṇasañjñā prāpnoti . \\
\noindent \textbf{(P\_1,4.42)} tamagrahaṇe punaḥ kriyamāṇe na doṣaḥ bhavati . \\
\noindent \textbf{(P\_1,4.42)} evam tarhi lokataḥ etat siddham . \\
\noindent \textbf{(P\_1,4.42)} tat yathā . \\
\noindent \textbf{(P\_1,4.42)} loke abhirūpāya udakamāneyam abhirūpāya kanyā deyā iti na ca anabhirūpe pravṛttiḥ asti . \\
\noindent \textbf{(P\_1,4.42)} tatra abhirūpatamāya iti gamyate . \\
\noindent \textbf{(P\_1,4.42)} evam iha api sādhakam karaṇam iti ucyate sarvāṇi ca kārakāṇi sādhakāni na ca asādhake pravṛttiḥ asti . \\
\noindent \textbf{(P\_1,4.42)} tatra sādhakatamam iti vijñāsyate . \\
\noindent \textbf{(P\_1,4.42)} evam tarhi siddhe sati yat tamagrahaṇam karoti tat jñāpayati ācāryaḥ kārakasañjñāyām taratamayogaḥ na bhavati iti . \\
\noindent \textbf{(P\_1,4.42)} kim etasya jñāpane prayojanam . \\
\noindent \textbf{(P\_1,4.42)} apādānam ācāryaḥ kim nyāyyam manyate . \\
\noindent \textbf{(P\_1,4.42)} yatra samprāpya nivṛttiḥ . \\
\noindent \textbf{(P\_1,4.42)} tena iha eva syāt grāmāt āgacchati nagarāt āgacchati iti . \\
\noindent \textbf{(P\_1,4.42)} sāṅkāśyakebhyaḥ pāṭaliputrakāḥ abhirūpatarāḥ iti atra na syāt . \\
\noindent \textbf{(P\_1,4.42)} kārakasañjñāyām taratamayogaḥ na bhavati iti atra api siddham bhavati . \\
\noindent \textbf{(P\_1,4.42)} tathā ādhāram ācāryaḥ kim nyāyyam manyate . \\
\noindent \textbf{(P\_1,4.42)} yatra kṛtsnaḥ ādhārātmā vyāptaḥ bhavati . \\
\noindent \textbf{(P\_1,4.42)} tena iha eva syāt tileṣu tailam dadhni sarpiḥ iti . \\
\noindent \textbf{(P\_1,4.42)} gaṅgāyām gāvaḥ kūpe gargarkulam iti atra na syāt . \\
\noindent \textbf{(P\_1,4.42)} kārakasañjñāyām taratamayogaḥ na bhavati iti atra api siddham bhavati . \\
\noindent \textbf{(P\_1,4.48)} vaseḥ aśyarthasya pratiṣedhaḥ . vaseḥ aśyarthasya pratiṣedhaḥ vaktavyaḥ . \\
\noindent \textbf{(P\_1,4.48)} grāma upavasati iti . \\
\noindent \textbf{(P\_1,4.48)} saḥ tarhi vaktavyaḥ . \\
\noindent \textbf{(P\_1,4.48)} na vaktavyaḥ . \\
\noindent \textbf{(P\_1,4.48)} na atra upapūrvasya vaseḥ grāmaḥ adhikaraṇam . \\
\noindent \textbf{(P\_1,4.48)} kasya tarhi . \\
\noindent \textbf{(P\_1,4.48)} anupasargasya . \\
\noindent \textbf{(P\_1,4.48)} grāme asau vasan trirātram upavasati iti . \\
\noindent \textbf{(P\_1,4.49.1)} tamagrahaṇam kimartham . \\
\noindent \textbf{(P\_1,4.49.1)} kartuḥ īpsitam karma iti iyati ucyamāne iha: agneḥ māṇavakam vārayati iti māṇavake apādānasañjñā prasajyeta . \\
\noindent \textbf{(P\_1,4.49.1)} na eṣaḥ doṣaḥ . \\
\noindent \textbf{(P\_1,4.49.1)} karmasañjñā tatra bādhikā bhaviṣyati . \\
\noindent \textbf{(P\_1,4.49.1)} agnau api tarhi bādhikā syāt . \\
\noindent \textbf{(P\_1,4.49.1)} iha punaḥ tamagrahaṇe kriyamāṇe tat upapannam bhavati yat uktam vāraṇārtheṣu karmagrahaṇānarthakyam kartuḥ īpsitatamam karma iti vacanāt iti . \\
\noindent \textbf{(P\_1,4.49.2)} iha ucyate odanam pacati iti . \\
\noindent \textbf{(P\_1,4.49.2)} yadi odanaḥ pacyeta dravyāntarama bhinirvarteta . \\
\noindent \textbf{(P\_1,4.49.2)} na eṣaḥ doṣaḥ . \\
\noindent \textbf{(P\_1,4.49.2)} tādarthyāt tācchabdyam bhaviṣyati . \\
\noindent \textbf{(P\_1,4.49.2)} odanārthāḥ taṇḍulāḥ odanaḥ iti . \\
\noindent \textbf{(P\_1,4.49.2)} atha iha katham bhavitavyam taṇḍulān odanam pacati iti āhosvit taṇḍulānām odanam pacati iti . \\
\noindent \textbf{(P\_1,4.49.2)} ubhyathā api bhavitavyam . \\
\noindent \textbf{(P\_1,4.49.2)} katham . \\
\noindent \textbf{(P\_1,4.49.2)} iha hi taṇḍulān odanam pacati iti dvyarthaḥ paciḥ . \\
\noindent \textbf{(P\_1,4.49.2)} taṇḍulān pacan odanam nirvartayati iti . \\
\noindent \textbf{(P\_1,4.49.2)} iha idanīm taṇḍulānām odanam pacati iti dvyarthaḥ ca eva paciḥ vikārayoge ca ṣaṣṭhī . \\
\noindent \textbf{(P\_1,4.49.2)} taṇḍulavikāram odanam nirvartayati iti . \\
\noindent \textbf{(P\_1,4.49.2)} iha kaḥ cit kam cidāmantrayate siddham bhujyatām iti . \\
\noindent \textbf{(P\_1,4.49.2)} saḥ āmantrayamāṇaḥ āha prabhūtam bhuktam asmābhiḥ iti . \\
\noindent \textbf{(P\_1,4.49.2)} āmantrayamāṇaḥ āha dadhi khalu bhaviṣyati payaḥ khalu bhaviṣyati . \\
\noindent \textbf{(P\_1,4.49.2)} āmantryamāṇaḥ āha dadhnā khalu bhuñjīya payasā khalu bhuñjīya iti . \\
\noindent \textbf{(P\_1,4.49.2)} atra karmasañjñā prāpnoti . \\
\noindent \textbf{(P\_1,4.49.2)} tat hi tasya īpsitatamam bhavati . \\
\noindent \textbf{(P\_1,4.49.2)} tasya api odanaḥ eva eva īpsitatamaḥ na tu guṇeṣu asya anurodhaḥ . \\
\noindent \textbf{(P\_1,4.49.2)} tat yathā bhuñjīya aham odanam yadi mṛduviśadaḥ syāt iti evam iha api dadhiguṇam odanam bhuñjīya payoguṇamodanam bhuñjīya iti . \\
\noindent \textbf{(P\_1,4.49.3)} īpsitasya karmasañjñāyām nirvṛttasya kārakatve karmasañjñāprasaṅgaḥ kriyepsitatvāt . \\
\noindent \textbf{(P\_1,4.49.3)} īpsitasya karmasañjñāyām nirvṛttasya kārakatve karmasañjñā na prāpnoti . \\
\noindent \textbf{(P\_1,4.49.3)} guḍam bhakṣayati iti . \\
\noindent \textbf{(P\_1,4.49.3)} kim kāraṇam . \\
\noindent \textbf{(P\_1,4.49.3)} kriyepsitatvāt . \\
\noindent \textbf{(P\_1,4.49.3)} kriyā tasya īpsitā . \\
\noindent \textbf{(P\_1,4.49.3)} na vā ubhayepsitatvāt . \\
\noindent \textbf{(P\_1,4.49.3)} na vā eṣaḥ doṣaḥ . \\
\noindent \textbf{(P\_1,4.49.3)} kim kāraṇam . \\
\noindent \textbf{(P\_1,4.49.3)} ubhayepsitatvāt . \\
\noindent \textbf{(P\_1,4.49.3)} ubhayam tasya īpsitam . \\
\noindent \textbf{(P\_1,4.49.3)} ātaḥ ca ubhayam yasya hi guḍabhakṣaṇe buddhiḥ prasaktā bhavati na asau loṣṭam bhakṣayitvā kṛtī bhavati . \\
\noindent \textbf{(P\_1,4.49.3)} yadi api tāvat atra etat śakyate vaktum ye tu ete rājakarmiṇaḥ manuṣyāḥ teṣām kaḥ cit kam cit āha kaṭam kuru iti . \\
\noindent \textbf{(P\_1,4.49.3)} sa āha na aham kaṭam kariṣyāmi ghaṭaḥ mayā āhṛtaḥ iti . \\
\noindent \textbf{(P\_1,4.49.3)} tasya kriyāmātram īpsitam . \\
\noindent \textbf{(P\_1,4.49.3)} yadi api tasya kriyāmātram īpsitam yaḥ tu asau preṣayati tasya ubhayam īpsitam iti . \\
\noindent \textbf{(P\_1,4.50)} kim udāharaṇam . \\
\noindent \textbf{(P\_1,4.50)} viṣam bhakṣayati iti . \\
\noindent \textbf{(P\_1,4.50)} na etat asti . \\
\noindent \textbf{(P\_1,4.50)} pūrveṇa api etat sidhyati . \\
\noindent \textbf{(P\_1,4.50)} na sidhyati . \\
\noindent \textbf{(P\_1,4.50)} kartuḥ īpsitatamam karma iti ucyate kasya ca nāma viṣabhakṣaṇam īpsitam syāt . \\
\noindent \textbf{(P\_1,4.50)} viṣabhakṣaṇam api kasya cit īpsitam bhavati . \\
\noindent \textbf{(P\_1,4.50)} katham . \\
\noindent \textbf{(P\_1,4.50)} iha yaḥ eṣaḥ manuṣyaḥ duḥkhārtaḥ bhavati saḥ anyāni duḥkhāni anuniśamya viṣabhakṣaṇam eva jyāyaḥ manyate . \\
\noindent \textbf{(P\_1,4.50)} ātaḥ ca īpsitam yat tat bhakṣayati . \\
\noindent \textbf{(P\_1,4.50)} yat tarthi anyat kariṣyāmi iti anyat karoti tat udāharaṇam . \\
\noindent \textbf{(P\_1,4.50)} kim punaḥ tat . \\
\noindent \textbf{(P\_1,4.50)} grāmāntaram ayam gacchan caurān paśyati ahim laṅghayati kaṇṭakān mṛdnāti . \\
\noindent \textbf{(P\_1,4.50)} iha īpsitasya api karmasañjñā ārabhyate anīpsitasya api . \\
\noindent \textbf{(P\_1,4.50)} yat idānīm na eva īpsitamam na api anīpsitam tatra katham bhavitavyam . \\
\noindent \textbf{(P\_1,4.50)} grāmāntaram ayam gacchan vṛkṣamūlāni upasarpati kuḍyamūlāni upasarpati iti . \\
\noindent \textbf{(P\_1,4.50)} atra api siddham . \\
\noindent \textbf{(P\_1,4.50)} katham . \\
\noindent \textbf{(P\_1,4.50)} anīpsitam iti na ayam prasajyapratiṣedhaḥ īpsitam na iti . \\
\noindent \textbf{(P\_1,4.50)} kim tarhi . \\
\noindent \textbf{(P\_1,4.50)} paryudāsaḥ ayam yat anyat īpsitāt tat anīpsitam iti . \\
\noindent \textbf{(P\_1,4.50)} anyat ca etat īpsitāt yat na eva īpsitam na api anīpsitam iti . \\
\noindent \textbf{(P\_1,4.51.1)} kena akathitam . \\
\noindent \textbf{(P\_1,4.51.1)} apādānādibhiḥ viśeṣakathābhiḥ . \\
\noindent \textbf{(P\_1,4.51.1)} kim udāharaṇam . \\
\noindent \textbf{(P\_1,4.51.1)} duhiyācirudhiprachibhikṣiciñām upayoganimittam apūrvavidhau bruviśāsiguṇena ca yatsacate tat akīrtitam ācaritam kavinā . \\
\noindent \textbf{(P\_1,4.51.1)} duhi : gām dogdhi payaḥ . \\
\noindent \textbf{(P\_1,4.51.1)} na etat asti . \\
\noindent \textbf{(P\_1,4.51.1)} kathitā atra pūrvā apādansañjñā . \\
\noindent \textbf{(P\_1,4.51.1)} duhi . \\
\noindent \textbf{(P\_1,4.51.1)} yāci : idam tarhi pauravam gām yācate iti . \\
\noindent \textbf{(P\_1,4.51.1)} na etat asti . \\
\noindent \textbf{(P\_1,4.51.1)} kathitā atra pūrvā apādansañjñā . \\
\noindent \textbf{(P\_1,4.51.1)} na yācanāt eva apāyaḥ bhavati . \\
\noindent \textbf{(P\_1,4.51.1)} yācitaḥ asau yadi dadāti tataḥ apāyena yujyate . \\
\noindent \textbf{(P\_1,4.51.1)} yāci . \\
\noindent \textbf{(P\_1,4.51.1)} rudhi : anvavaruṇaddhi gām vrajam . \\
\noindent \textbf{(P\_1,4.51.1)} na etat asti . \\
\noindent \textbf{(P\_1,4.51.1)} kathitā atra pūrvā adhikaraṇasañjñā . \\
\noindent \textbf{(P\_1,4.51.1)} rudhi . \\
\noindent \textbf{(P\_1,4.51.1)} pracchi : māṇavakam panthānam pṛcchati . \\
\noindent \textbf{(P\_1,4.51.1)} na etat asti . \\
\noindent \textbf{(P\_1,4.51.1)} kathitā atra pūrvā apādānasañjñā . \\
\noindent \textbf{(P\_1,4.51.1)} na praśnāt eva apāyaḥ bhavati . \\
\noindent \textbf{(P\_1,4.51.1)} pṛṣtaḥ asau yadi ācāṣṭe tataḥ apāyena yujyate . \\
\noindent \textbf{(P\_1,4.51.1)} pracchi. bhikṣi : pauravam gām bhikṣate . \\
\noindent \textbf{(P\_1,4.51.1)} na etat asti . \\
\noindent \textbf{(P\_1,4.51.1)} kathitā atra pūrvā apādānasañjñā . \\
\noindent \textbf{(P\_1,4.51.1)} na bhikṣaṇāt eva apāyaḥ bhavati . \\
\noindent \textbf{(P\_1,4.51.1)} bhikṣitaḥ asau yadi dadāti tataḥ apāyena yujyate . \\
\noindent \textbf{(P\_1,4.51.1)} bhikṣi . \\
\noindent \textbf{(P\_1,4.51.1)} ciñ : vṛkṣam avacinoti phalāni . \\
\noindent \textbf{(P\_1,4.51.1)} na etat asti . \\
\noindent \textbf{(P\_1,4.51.1)} kathitā atra pūrvā apādānasañjñā . \\
\noindent \textbf{(P\_1,4.51.1)} bruviśāsiguṇena ca yat sacate tat akīrtitam ācaritam kavinā . \\
\noindent \textbf{(P\_1,4.51.1)} bruviśāsiguṇena ca yat sacate sambadhyate tat ca dāharaṇam . \\
\noindent \textbf{(P\_1,4.51.1)} kim punaḥ tat . \\
\noindent \textbf{(P\_1,4.51.1)} putram brūte dharmam . \\
\noindent \textbf{(P\_1,4.51.1)} putram anuśāsti dharmam iti . \\
\noindent \textbf{(P\_1,4.51.1)} na etat asti . \\
\noindent \textbf{(P\_1,4.51.1)} kathitā atra pūrvā sampradānasañjñā . \\
\noindent \textbf{(P\_1,4.51.1)} tasmāt trīṇi eva udāharaṇāni . \\
\noindent \textbf{(P\_1,4.51.1)} pauravam gām yācate . \\
\noindent \textbf{(P\_1,4.51.1)} māṇavakam panthānam pṛcchati . \\
\noindent \textbf{(P\_1,4.51.1)} pauravam gām bhikṣate iti . \\
\noindent \textbf{(P\_1,4.51.2)} atha ye dhātūnām dvikarmakāḥ teṣām kim kathite lādayaḥ bhavanti āhosvit akathite . \\
\noindent \textbf{(P\_1,4.51.2)} kathite lādayaḥ . \\
\noindent \textbf{(P\_1,4.51.2)} kathite lādibhiḥ abhihite guṇamkarmaṇi kā kartavyā . kathite lādayaḥ cet syuḥ ṣaṣṭḥīm kuryāt tadā guṇe . \\
\noindent \textbf{(P\_1,4.51.2)} kathite lādayaḥ cet syuḥ ṣaṣṭhī guṇakarmaṇi tadā kartavyā . \\
\noindent \textbf{(P\_1,4.51.2)} duhyate goḥ payaḥ . \\
\noindent \textbf{(P\_1,4.51.2)} yācyate pauravasya kambalaḥ iti . \\
\noindent \textbf{(P\_1,4.51.2)} katham . \\
\noindent \textbf{(P\_1,4.51.2)} akārakam hyakathitatvāt . \\
\noindent \textbf{(P\_1,4.51.2)} akārakam hi etat bhavati . \\
\noindent \textbf{(P\_1,4.51.2)} kim kāraṇam . \\
\noindent \textbf{(P\_1,4.51.2)} akathitatvāt . \\
\noindent \textbf{(P\_1,4.51.2)} kārakam cet tu na akathā . \\
\noindent \textbf{(P\_1,4.51.2)} atha kārakam na akathitam . \\
\noindent \textbf{(P\_1,4.51.2)} atha kārake sati kā kartavyā . \\
\noindent \textbf{(P\_1,4.51.2)} kārakam cet vijānīyāt yām yām manyeta sā bhavet . \\
\noindent \textbf{(P\_1,4.51.2)} kārakam cet vijānātīyāt yā yā prāpnoti sā kartavyā . \\
\noindent \textbf{(P\_1,4.51.2)} duhyate goḥ payaḥ . \\
\noindent \textbf{(P\_1,4.51.2)} yācyate pauravāt kambalaḥ iti . \\
\noindent \textbf{(P\_1,4.51.2)} kathite abhihite tvavidhiḥ tvamatiḥ guṇakarmaṇi lādividhiḥ sapare . \\
\noindent \textbf{(P\_1,4.51.2)} kathite lādibhiḥ abhihite tvavidhiḥ eṣaḥ bhavati . \\
\noindent \textbf{(P\_1,4.51.2)} kim idam tvavidhiḥ iti . \\
\noindent \textbf{(P\_1,4.51.2)} tava vidhiḥ tvavidhiḥ . \\
\noindent \textbf{(P\_1,4.51.2)} tvamatiḥ . \\
\noindent \textbf{(P\_1,4.51.2)} kimidam tvamatiḥ iti . \\
\noindent \textbf{(P\_1,4.51.2)} tava matiḥ tvamatiḥ . \\
\noindent \textbf{(P\_1,4.51.2)} na evam anye manyante . \\
\noindent \textbf{(P\_1,4.51.2)} katham tarhi anye manyante . \\
\noindent \textbf{(P\_1,4.51.2)} guṇakarmaṇi lādividhiḥ sapare . \\
\noindent \textbf{(P\_1,4.51.2)} guṇakarmaṇi lādividhiyaḥ bhavanti saha pareṇa yogena . \\
\noindent \textbf{(P\_1,4.51.2)} gatibuddhipratyavasānārthaśabdakarmākarmakāṇām aṇikartā saḥ ṇau iti . \\
\noindent \textbf{(P\_1,4.51.2)} dhruvaceṣṭitayuktiṣu ca api aguṇe tat analpamateḥ vacanam smarata . dhruvayuktiṣu ceṣṭitayuktiṣu ca api aguṇe karmaṇi lādayaḥ bhavanti . \\
\noindent \textbf{(P\_1,4.51.2)} tat analpmateḥ ācāryasya vacanam smaryatām . \\
\noindent \textbf{(P\_1,4.51.2)} aparaḥ āha : pradhānakarmaṇi ākhyeye lādīn āhuḥ dvikarmaṇām . pradhānakarmaṇi abhidheye dvikarmaṇām dhātūnām karmaṇi lādayaḥ bhavanti iti vaktavyam . \\
\noindent \textbf{(P\_1,4.51.2)} ajām nayati grāmam . \\
\noindent \textbf{(P\_1,4.51.2)} ajā nīyate grāmam . \\
\noindent \textbf{(P\_1,4.51.2)} ajā nītā grāmam iti . \\
\noindent \textbf{(P\_1,4.51.2)} apradhāne duhādīnām . apradhāne duhādīnām karmaṇi lādayaḥ bhavanti iti vaktavyam . \\
\noindent \textbf{(P\_1,4.51.2)} duhyate gauḥ payaḥ . \\
\noindent \textbf{(P\_1,4.51.2)} ṇyante kartuḥ ca karmaṇaḥ . lādayaḥ bhavanti iti vaktavyam . \\
\noindent \textbf{(P\_1,4.51.2)} gamyate devadattaḥ grāmam yajñadattena . \\
\noindent \textbf{(P\_1,4.51.2)} ke punaḥ dhātūnām dvikarmakāḥ . \\
\noindent \textbf{(P\_1,4.51.2)} nīvahyoḥ harateḥ ca api gatyarthānām tathā eva ca dvikarmakeṣu grahaṇam draṣṭavyam iti niścayaḥ . \\
\noindent \textbf{(P\_1,4.51.2)} ajām nayati grāmam . \\
\noindent \textbf{(P\_1,4.51.2)} bhāram vahati grāmam . \\
\noindent \textbf{(P\_1,4.51.2)} bhāram harati grāmam . \\
\noindent \textbf{(P\_1,4.51.2)} gatyarthānām . \\
\noindent \textbf{(P\_1,4.51.2)} gamayati devadattam grāmam . \\
\noindent \textbf{(P\_1,4.51.2)} yāpayati devadattam grāmam . \\
\noindent \textbf{(P\_1,4.51.2)} siddham vā apo anyakarmaṇaḥ . \\
\noindent \textbf{(P\_1,4.51.2)} siddham vā punaḥ etat bhavati . \\
\noindent \textbf{(P\_1,4.51.2)} kutaḥ . \\
\noindent \textbf{(P\_1,4.51.2)} anyakarmaṇaḥ . \\
\noindent \textbf{(P\_1,4.51.2)} anyasya atra ajā karma anyasya grāmaḥ . \\
\noindent \textbf{(P\_1,4.51.2)} ajām asau gṛhītvā grāmam nayati . \\
\noindent \textbf{(P\_1,4.51.2)} anyakarma iti cet brūyāt lādīnām avidhiḥ bhavet . anyakarma iti cet brūyāt lādīnām avidhiḥ ayam bhavet . \\
\noindent \textbf{(P\_1,4.51.2)} ajā nīyate grāmam iti . \\
\noindent \textbf{(P\_1,4.51.2)} parasādhane utpadyamānena lena ajāyāḥ abhidhānam na prāpnoti . \\
\noindent \textbf{(P\_1,4.51.3)} kālabhāvādhvagantavyāḥ karmasañjñā hi arkarmaṇām . \\
\noindent \textbf{(P\_1,4.51.3)} kālabhāvādhvagantavyāḥ akarmakāṇām dhātūnām karmasañjñāḥ bhavanti iti vaktavyam . \\
\noindent \textbf{(P\_1,4.51.3)} kāla . \\
\noindent \textbf{(P\_1,4.51.3)} māsam āste . \\
\noindent \textbf{(P\_1,4.51.3)} māsam svapiti . \\
\noindent \textbf{(P\_1,4.51.3)} bhāva . \\
\noindent \textbf{(P\_1,4.51.3)} godoham āste . \\
\noindent \textbf{(P\_1,4.51.3)} godoham svapiti . \\
\noindent \textbf{(P\_1,4.51.3)} adhvagantavya . \\
\noindent \textbf{(P\_1,4.51.3)} krośam āste . \\
\noindent \textbf{(P\_1,4.51.3)} krośam svapiti . \\
\noindent \textbf{(P\_1,4.51.3)} deśaḥ ca akarmaṇām karmasañjñaḥ bhavati iti vaktavyam . \\
\noindent \textbf{(P\_1,4.51.3)} kurūn svapiti . \\
\noindent \textbf{(P\_1,4.51.3)} pañcālān svapiti . \\
\noindent \textbf{(P\_1,4.51.3)} viparītam tu yat karma tat kalma kavayaḥ viduḥ . kimidam kalma iti . \\
\noindent \textbf{(P\_1,4.51.3)} aparisamāptam karma kalma . \\
\noindent \textbf{(P\_1,4.51.3)} na vā asmin sarvāṇi karmakāryāṇi kriyante . \\
\noindent \textbf{(P\_1,4.51.3)} kim tarhi . \\
\noindent \textbf{(P\_1,4.51.3)} dvitīyā eva . \\
\noindent \textbf{(P\_1,4.51.3)} yasmin tu karmaṇi upajāyate anyat dhātvarthayogā api ca yatra ṣaṣṭhī tat karma kalma iti ca kalma na uktam dhātoḥ hi vṛttiḥ na ralatvataḥ asti . \\
\noindent \textbf{(P\_1,4.51.3)} etena karmasañjñā sarvā siddhā bhavati kathitena . \\
\noindent \textbf{(P\_1,4.51.3)} tatra īpsitasya kim syāt prayojanam karmasañjñāyāḥ . \\
\noindent \textbf{(P\_1,4.51.3)} yat tu kathitam purastāt īpsitatayuktam ca tasya siddhyartham . \\
\noindent \textbf{(P\_1,4.51.3)} īpsitam eva tu yat syāt tasya bhaviṣyati kathitena . \\
\noindent \textbf{(P\_1,4.51.3)} atha iha katham bhavitavyam . \\
\noindent \textbf{(P\_1,4.51.3)} netā aśvasya srughnam iti āhosvit netā aśvasya srughnasya iti . \\
\noindent \textbf{(P\_1,4.51.3)} ubhayathā goṇikāputraḥ . \\
\noindent \textbf{(P\_1,4.52.1)} śabdakarma iti katham idam vijñāyate . \\
\noindent \textbf{(P\_1,4.52.1)} śabdaḥ yeṣām kriyā iti āhosvit śabdaḥ yeṣām karma iti . \\
\noindent \textbf{(P\_1,4.52.1)} kaḥ ca atra viśeṣaḥ . \\
\noindent \textbf{(P\_1,4.52.1)} śabdakarmanirdeśe śabdakriyāṇām iti cet hvayatyādīnām pratiṣedhaḥ . śabdakarmanirdeśe śabdakriyāṇāmiti ced hvayadādīnām pratiṣedhaḥ vaktavyaḥ . \\
\noindent \textbf{(P\_1,4.52.1)} ke punaḥ hvayatādayaḥ . \\
\noindent \textbf{(P\_1,4.52.1)} hvayati krandati śabdāyate . \\
\noindent \textbf{(P\_1,4.52.1)} hvayati devadattaḥ . \\
\noindent \textbf{(P\_1,4.52.1)} hvāyayati devadattena . \\
\noindent \textbf{(P\_1,4.52.1)} krandati devadattaḥ . \\
\noindent \textbf{(P\_1,4.52.1)} krandayati devadattena . \\
\noindent \textbf{(P\_1,4.52.1)} śabdāyate devadattaḥ . \\
\noindent \textbf{(P\_1,4.52.1)} śabdāyayati devadattena iti . \\
\noindent \textbf{(P\_1,4.52.1)} śṛṇotyādīn ām ca upasamkhyānam aśabdakriyatvāt . \\
\noindent \textbf{(P\_1,4.52.1)} śṛṇotyādīnām ca upasamkhyānam kartavyam . \\
\noindent \textbf{(P\_1,4.52.1)} ke punaḥ śṛṇotyādayaḥ . \\
\noindent \textbf{(P\_1,4.52.1)} śṛṇoti vijānāti upalabhate . \\
\noindent \textbf{(P\_1,4.52.1)} śṛṇoti devadattaḥ . \\
\noindent \textbf{(P\_1,4.52.1)} śrāvayati devadattam . \\
\noindent \textbf{(P\_1,4.52.1)} vijānāti devadattaḥ . \\
\noindent \textbf{(P\_1,4.52.1)} vijñāpayati devadattam . \\
\noindent \textbf{(P\_1,4.52.1)} upalabhate devadattaḥ . \\
\noindent \textbf{(P\_1,4.52.1)} upalambhayati devadattam . \\
\noindent \textbf{(P\_1,4.52.1)} kim punaḥ kāraṇam na sidhyati . \\
\noindent \textbf{(P\_1,4.52.1)} aśabdakriyatvād . \\
\noindent \textbf{(P\_1,4.52.1)} astu tarhi śabdaḥ yeṣām karma iti . \\
\noindent \textbf{(P\_1,4.52.1)} śabdakarmaṇaḥ iti cet jalpatiprabhṛtīnām upasamkhyānam . śabdakarmaṇa iti cet jalpatiprabhṛtīnāmupasaṅkhyānam kartavyam . \\
\noindent \textbf{(P\_1,4.52.1)} ke punaḥ jalpatiprabhṛatayaḥ . \\
\noindent \textbf{(P\_1,4.52.1)} jalpati vilapati ābhāṣate . \\
\noindent \textbf{(P\_1,4.52.1)} jalpati devadattaḥ . \\
\noindent \textbf{(P\_1,4.52.1)} jalpayati devadattam . \\
\noindent \textbf{(P\_1,4.52.1)} vilapati devadattaḥ . \\
\noindent \textbf{(P\_1,4.52.1)} vilāpayati devadattam . \\
\noindent \textbf{(P\_1,4.52.1)} ābhāṣate devadattaḥ .ābhāṣayati devadattam . \\
\noindent \textbf{(P\_1,4.52.1)} dṛśeḥ sarvatra . \\
\noindent \textbf{(P\_1,4.52.1)} dṛśeḥ sarvatra upasaṅkhyānam kartavyam . \\
\noindent \textbf{(P\_1,4.52.1)} paśyati rūpatarkaḥ kārṣāpaṇam . \\
\noindent \textbf{(P\_1,4.52.1)} darśayati rūpatarkam kārṣāpaṇam . \\
\noindent \textbf{(P\_1,4.52.2)} adikhādinīvahīnām pratiṣedhaḥ . \\
\noindent \textbf{(P\_1,4.52.2)} adikhādinīvahīnām pratiṣedhaḥ vaktavyaḥ . \\
\noindent \textbf{(P\_1,4.52.2)} atti devadattaḥ . \\
\noindent \textbf{(P\_1,4.52.2)} ādayate devadattena . \\
\noindent \textbf{(P\_1,4.52.2)} aparaḥ āha : sarvam eva pratyavasānakāryam adeḥ na bhavati iti vaktavyam , parasmaipadam api . \\
\noindent \textbf{(P\_1,4.52.2)} idam ekam iṣyate : ktaḥ adhikaraṇe ca drauvyagatipratyavasānārthebhyaḥ : idam eṣām jagdham . \\
\noindent \textbf{(P\_1,4.52.2)} khādi . \\
\noindent \textbf{(P\_1,4.52.2)} khādati devadattaḥ . \\
\noindent \textbf{(P\_1,4.52.2)} khādayati devadattena . \\
\noindent \textbf{(P\_1,4.52.2)} nī . \\
\noindent \textbf{(P\_1,4.52.2)} nayati devadattaḥ . \\
\noindent \textbf{(P\_1,4.52.2)} nāyayati devadattena . \\
\noindent \textbf{(P\_1,4.52.2)} vaheraniyantṛkartṛkasya . \\
\noindent \textbf{(P\_1,4.52.2)} vaheḥ aniyantṛkartṛkasya iti vaktavyam . \\
\noindent \textbf{(P\_1,4.52.2)} vahati bhāram devadattaḥ . \\
\noindent \textbf{(P\_1,4.52.2)} vāhayati bhāram devadattena . \\
\noindent \textbf{(P\_1,4.52.2)} aniyantṛkartṛkasya iti kimartham . \\
\noindent \textbf{(P\_1,4.52.2)} vahanti yavān balīvardāḥ . \\
\noindent \textbf{(P\_1,4.52.2)} vāhayanti balīvardān yavān . \\
\noindent \textbf{(P\_1,4.52.2)} bhakṣeḥ ahiṃsārthasya . \\
\noindent \textbf{(P\_1,4.52.2)} bhakṣeḥ ahimsārthasya iti vaktavyam . \\
\noindent \textbf{(P\_1,4.52.2)} bhakṣayati piṇḍīm devadattaḥ . \\
\noindent \textbf{(P\_1,4.52.2)} bhakṣayati piṇḍīm devadattena . \\
\noindent \textbf{(P\_1,4.52.2)} ahimsārthasya iti kimartham . \\
\noindent \textbf{(P\_1,4.52.2)} bhakṣayanti yavān balīvardāḥ . \\
\noindent \textbf{(P\_1,4.52.2)} bhakṣayanti balīvardān yavān . \\
\noindent \textbf{(P\_1,4.52.3)} akarmakagrahaṇe kālakarmakāṇām upasaṅkhyānam . \\
\noindent \textbf{(P\_1,4.52.3)} akarmakagrahaṇe kālakarmakāṇām upasaṅkhyānam kartavyam . \\
\noindent \textbf{(P\_1,4.52.3)} māsam āste devadattaḥ . \\
\noindent \textbf{(P\_1,4.52.3)} māsam āsayati devadattam . \\
\noindent \textbf{(P\_1,4.52.3)} māsam śete devadattaḥ . \\
\noindent \textbf{(P\_1,4.52.3)} māsam śāyayati devadattam . \\
\noindent \textbf{(P\_1,4.52.3)} siddham tu kālakarmakāṇām akarmakavadvacanāt . \\
\noindent \textbf{(P\_1,4.52.3)} siddham etat . \\
\noindent \textbf{(P\_1,4.52.3)} katham . \\
\noindent \textbf{(P\_1,4.52.3)} kālakarmakāḥ akarmakavat bhavanti iti vaktavyam . \\
\noindent \textbf{(P\_1,4.52.3)} tat tarhi vaktavyam . \\
\noindent \textbf{(P\_1,4.52.3)} na vaktavyam . \\
\noindent \textbf{(P\_1,4.52.3)} akarmakāṇām iti ucyate na ca ke cit kadā cit kālabhāvādhvabhiḥ akarmakāḥ . \\
\noindent \textbf{(P\_1,4.52.3)} te evam vijñāsyāmaḥ . \\
\noindent \textbf{(P\_1,4.52.3)} kva cit ye akarmakāḥ iti . \\
\noindent \textbf{(P\_1,4.52.3)} atha vā yena karmaṇā sakarmkāḥ ca akarmakāḥ ca bhavanti tena akarmakāṇām . \\
\noindent \textbf{(P\_1,4.52.3)} na ca etena karmaṇā kaḥ cit api akarmakaḥ . \\
\noindent \textbf{(P\_1,4.52.3)} atha vā yat karma bhavati na ca bhavati tena karmakāṇām . \\
\noindent \textbf{(P\_1,4.52.3)} na ca etat karma kva cit api na bhavati . \\
\noindent \textbf{(P\_1,4.53)} hṛkroḥ vāvacane abhivādidṛśyoḥ ātmanepade upasaṅkhyānam . \\
\noindent \textbf{(P\_1,4.53)} hṛkrorvāvacane abhivādidṛśoḥ ātmanepade upasaṅkhyānam kartavyam . \\
\noindent \textbf{(P\_1,4.53)} abhivadati gurum devadattaḥ . \\
\noindent \textbf{(P\_1,4.53)} abhivādayate gurum devadattam . \\
\noindent \textbf{(P\_1,4.53)} abhivādayate gurum devadattena . \\
\noindent \textbf{(P\_1,4.53)} paśyanti bhṛtyāḥ rājānam . \\
\noindent \textbf{(P\_1,4.53)} darśayate bhṛtyān rājā . \\
\noindent \textbf{(P\_1,4.53)} darśayate bhṛtyaiḥ rājā . \\
\noindent \textbf{(P\_1,4.53)} katham ca atra ātmanepadam . \\
\noindent \textbf{(P\_1,4.53)} ekasya ṇeḥ aṇau i ti aparasya ṇicaḥ ca iti . \\
\noindent \textbf{(P\_1,4.54.1)} kim yasya svam tantram saḥ svatantraḥ . \\
\noindent \textbf{(P\_1,4.54.1)} kim ca ataḥ . \\
\noindent \textbf{(P\_1,4.54.1)} tantuvāye prāpnoti . \\
\noindent \textbf{(P\_1,4.54.1)} na eṣaḥ doṣaḥ . \\
\noindent \textbf{(P\_1,4.54.1)} ayam tantraśabdaḥ asti eva vitāne vartate . \\
\noindent \textbf{(P\_1,4.54.1)} tat yathā : āstīrṇam tantram . \\
\noindent \textbf{(P\_1,4.54.1)} pretam tantram . \\
\noindent \textbf{(P\_1,4.54.1)} vitānaḥ iti gamyate . \\
\noindent \textbf{(P\_1,4.54.1)} asti prādhānye vartate . \\
\noindent \textbf{(P\_1,4.54.1)} tat yathā svatantraḥ asau brāhmaṇaḥ iti ucyate . \\
\noindent \textbf{(P\_1,4.54.1)} svapradhānaḥ iti gamyate . \\
\noindent \textbf{(P\_1,4.54.1)} tat yaḥ prādhānye vartate tantraśabdaḥ tasya idam grahaṇam . \\
\noindent \textbf{(P\_1,4.54.2)} svatantrasya kartṛsañjñāyām hetumati upasaṅkhyānam asvatrantvāt . \\
\noindent \textbf{(P\_1,4.54.2)} svatantrasya kartṛsañjñāyām hetumati upasaṅkhyānam kartavyam . \\
\noindent \textbf{(P\_1,4.54.2)} pācayati odanam devadattaḥ yajñadattena iti . \\
\noindent \textbf{(P\_1,4.54.2)} kim punaḥ kāraṇam na sidhyati . \\
\noindent \textbf{(P\_1,4.54.2)} asvatantratvāt . \\
\noindent \textbf{(P\_1,4.54.2)} na vā svātrantryāt itarathā hi akurvati api kārayati iti syāt . \\
\noindent \textbf{(P\_1,4.54.2)} na vā kartavyam . \\
\noindent \textbf{(P\_1,4.54.2)} kim kāraṇam . \\
\noindent \textbf{(P\_1,4.54.2)} svātantryāt . \\
\noindent \textbf{(P\_1,4.54.2)} svatantraḥ asau bhavati . \\
\noindent \textbf{(P\_1,4.54.2)} itarathā hi akurvati api kārayati iti syāt . \\
\noindent \textbf{(P\_1,4.54.2)} yaḥ hi manyate na asau svatantraḥ akurvati api tasya kārayati iti etat syāt . \\
\noindent \textbf{(P\_1,4.54.2)} na akurvat i iti cet svatantraḥ . \\
\noindent \textbf{(P\_1,4.54.2)} na cet akurvati tasmin kārayati iti etat bhavati svatantraḥ asau bhavati . \\
\noindent \textbf{(P\_1,4.54.2)} śakyam tāvat anena upasamkhyānam kurvatā vaktum kurvan svatantraḥ akurvan na iti . \\
\noindent \textbf{(P\_1,4.54.2)} sādhīyaḥ jñāpakam bhavati . \\
\noindent \textbf{(P\_1,4.54.2)} preṣite ca kila ayam kriyām ca akriyām ca dṛṣṭvā adhyavasyati kurvan svatantraḥ akurvan na iti . \\
\noindent \textbf{(P\_1,4.54.2)} yadi ca preṣitaḥ asau na karoti svatantraḥ asau bhavati iti . \\
\noindent \textbf{(P\_1,4.55)} praiṣe asvatantraprayojakatvāt hetusañjñāprasiddhiḥ . \\
\noindent \textbf{(P\_1,4.55)} praiṣe asvatantraprayojakatvāt hetusañjñāyāḥ aprasiddhiḥ . \\
\noindent \textbf{(P\_1,4.55)} svatantraprayojakaḥ hetusañjñaḥ bhavati iti ucyate . \\
\noindent \textbf{(P\_1,4.55)} na ca asau svatantram prayojayati . \\
\noindent \textbf{(P\_1,4.55)} svatantratvāt siddham . \\
\noindent \textbf{(P\_1,4.55)} siddham etat . \\
\noindent \textbf{(P\_1,4.55)} katham . \\
\noindent \textbf{(P\_1,4.55)} svatantratvāt . \\
\noindent \textbf{(P\_1,4.55)} svatantram asau prayojayati . \\
\noindent \textbf{(P\_1,4.55)} svatantratvāt siddham iti cet svatantraparatantratvam vipratiṣiddam . \\
\noindent \textbf{(P\_1,4.55)} yadi svatantraḥ na prayojyaḥ atha prayojyaḥ na svatantraḥ prayojyaḥ svatantraḥ ca iti vipratiṣiddham . \\
\noindent \textbf{(P\_1,4.55)} uktam vā . \\
\noindent \textbf{(P\_1,4.55)} kim uktam . \\
\noindent \textbf{(P\_1,4.55)} ekam tāvat uktam na vā svātrantryāt itarathā hi akurvati api kārayati iti syāt iti . \\
\noindent \textbf{(P\_1,4.55)} aparam uktam . \\
\noindent \textbf{(P\_1,4.55)} na vā sāmānyakṛtatvāt hetutaḥ hi aviśiṣṭam svatantraprayojakatvāt aprayojakaḥ iti cet muktamasamśayena tulyam iti . \\
\noindent \textbf{(P\_1,4.56)} kimartham rephādhikaḥ īśvaraśabdaḥ gṛhyate . \\
\noindent \textbf{(P\_1,4.56)} rīśvarāt vīśvarāt mā bhūt . \\
\noindent \textbf{(P\_1,4.56)} rīśvarāt iti ucyate vīśvarāt mā bhūt . \\
\noindent \textbf{(P\_1,4.56)} śaki ṇamulkamulau īśvare tosunkasunau iti . \\
\noindent \textbf{(P\_1,4.56)} na etat asti prayojanam . \\
\noindent \textbf{(P\_1,4.56)} ācāryapravṛttiḥ jñāpayati anantaraḥ yaḥ īśvaraśabdaḥ tasya grahaṇam iti yat ayam kṛt mejantaḥ iti kṛtaḥ māntasya ejantasya ca avyayasañjñām śāsti . \\
\noindent \textbf{(P\_1,4.56)} kṛt mejantaḥ paraḥ api saḥ . \\
\noindent \textbf{(P\_1,4.56)} paraḥ api etasmāt kṛt māntaḥ ejantaḥ ca asti . \\
\noindent \textbf{(P\_1,4.56)} tadartham etat syāt . \\
\noindent \textbf{(P\_1,4.56)} yat tarhi avyayībhāvasya avyayasañjñām śāsti tat jñāpayati ācāryaḥ nantaraḥ yaḥ īśvaraśabdaḥ tasya grahaṇam iti . \\
\noindent \textbf{(P\_1,4.56)} samāseṣu avyayībhāvaḥ . \\
\noindent \textbf{(P\_1,4.56)} samāsasya etat jñāpakam syāt . \\
\noindent \textbf{(P\_1,4.56)} avyayībhāvaḥ eva samāsaḥ avyayasañjñaḥ bhavati na anyaḥ iti . \\
\noindent \textbf{(P\_1,4.56)} evam tarhi lokataḥ etat siddham . \\
\noindent \textbf{(P\_1,4.56)} tat yathā loke ā vanāntāt ā udakāntāt priyam pānthaman uvrajet iti yaḥ eva prathamaḥ vanāntaḥ udakāntaḥ ca tataḥ nuvrajati . \\
\noindent \textbf{(P\_1,4.56)} laukikam ca ativartate . \\
\noindent \textbf{(P\_1,4.56)} dvitīyam ca tṛtīyam ca vanāntam udakāntam vā anuvrajati . \\
\noindent \textbf{(P\_1,4.56)} tasmāt rephādikaḥ īśvaraśabdaḥ grahītavyaḥ . \\
\noindent \textbf{(P\_1,4.56)} atha prāgvacanam kimartham . \\
\noindent \textbf{(P\_1,4.56)} prāgvacanam sañjñānivṛttyartham . \\
\noindent \textbf{(P\_1,4.56)} prāgvacanam kriyate nipātasañjñāyāḥ anivṛttiḥ yathā syāt . \\
\noindent \textbf{(P\_1,4.56)} akriyamāṇe hi prāgvacane anavakāśāḥ gatyupasargakarmapravacanīyasañjñāḥ nipātasañjñām bādheran . \\
\noindent \textbf{(P\_1,4.56)} tāḥ mā bādhiṣata iti prāgvacanam kriyate . \\
\noindent \textbf{(P\_1,4.56)} atha kriyamāṇe api prāgvacane yāvatā anavakāśāḥ etāḥ sañjñāḥ kasmāt eva na bādhante . \\
\noindent \textbf{(P\_1,4.56)} kriyamāṇe hi prāgvacane satyām nipātasañjñāyām etāḥ avayavasañjñāḥ ārabhyante . \\
\noindent \textbf{(P\_1,4.56)} tatra vacanāt samāveśaḥ bhavati . \\
\noindent \textbf{(P\_1,4.57)} ayam sattvaśabdaḥ asti eva dravyapadārthakaḥ . \\
\noindent \textbf{(P\_1,4.57)} tat yathā sattvam ayam brāhmaṇaḥ sattvamiyam brāhmaṇī iti . \\
\noindent \textbf{(P\_1,4.57)} asti kriyāpadārthakaḥ . \\
\noindent \textbf{(P\_1,4.57)} sadbhāvaḥ sattvam iti . \\
\noindent \textbf{(P\_1,4.57)} kasya idam grahaṇam . \\
\noindent \textbf{(P\_1,4.57)} dravyapadārthakasya . \\
\noindent \textbf{(P\_1,4.57)} kutaḥ etat . \\
\noindent \textbf{(P\_1,4.57)} evam hi kṛtvā vidhiḥ ca siddhaḥ bhavati pratiṣedhaḥ ca . \\
\noindent \textbf{(P\_1,4.57)} kim punaḥ ayam paryudāsaḥ . \\
\noindent \textbf{(P\_1,4.57)} yat anyat sattvavacanāt iti . \\
\noindent \textbf{(P\_1,4.57)} āhosvit prasajya ayam pratiṣedhaḥ . \\
\noindent \textbf{(P\_1,4.57)} sattvavacane na iti . \\
\noindent \textbf{(P\_1,4.57)} kim ca ataḥ . \\
\noindent \textbf{(P\_1,4.57)} yadi paryudāsaḥ vipraḥ iti atra api prāpnoti . \\
\noindent \textbf{(P\_1,4.57)} kriyādravyavacanaḥ ayam samghāto dravyāt anyaśca vidhinā āśrīyate . \\
\noindent \textbf{(P\_1,4.57)} asti ca prādibhiḥ sāmānyam iti kṛtvā tadantavidhinā nipātasañjñā prāpnoti . \\
\noindent \textbf{(P\_1,4.57)} atha prasajyapratiṣedhaḥ na doṣaḥ bhavati . \\
\noindent \textbf{(P\_1,4.57)} yathā na doṣaḥ tathā astu . \\
\noindent \textbf{(P\_1,4.58-59.1)} KA\_I,341.11-18 Ro\_II,444 {1/9}     prādayaḥ iti yogavibhāgaḥ . \\
\noindent \textbf{(P\_1,4.58-59.1)} KA\_I,341.11-18 Ro\_II,444 {2/9}     prādayaḥ iti yogavibhāgaḥ kartavyaḥ . \\
\noindent \textbf{(P\_1,4.58-59.1)} KA\_I,341.11-18 Ro\_II,444 {3/9}     prādayaḥ sattvavacanāḥ nipātasañjñāḥ bhavanti . \\
\noindent \textbf{(P\_1,4.58-59.1)} KA\_I,341.11-18 Ro\_II,444 {4/9}     tataḥ upasargāḥ kriyāyoge iti . \\
\noindent \textbf{(P\_1,4.58-59.1)} KA\_I,341.11-18 Ro\_II,444 {5/9}     kimarthaḥ yogavibhāgaḥ . \\
\noindent \textbf{(P\_1,4.58-59.1)} KA\_I,341.11-18 Ro\_II,444 {6/9}     nipātasañjñārthaḥ . \\
\noindent \textbf{(P\_1,4.58-59.1)} KA\_I,341.11-18 Ro\_II,444 {7/9}     nipātasañjñā yathā syāt . \\
\noindent \textbf{(P\_1,4.58-59.1)} KA\_I,341.11-18 Ro\_II,444 {8/9}     ekayoge hi nipātasañjñābhāvaḥ . ekayoge hi sati nipātasañjñāyā abhāvaḥ syāt . \\
\noindent \textbf{(P\_1,4.58-59.1)} KA\_I,341.11-18 Ro\_II,444 {9/9}     yasmin eva viśeṣe gatyupasargakarmapravacanīyasañjñāḥ tasmin eva viśeṣe nipātasañjñā syāt . \\
\noindent \textbf{(P\_1,4.58-59.2)} KA\_I,341.19-23 Ro\_II,445 {1/7}     marucchabdasya upsaṅkhyānam . \\
\noindent \textbf{(P\_1,4.58-59.2)} KA\_I,341.19-23 Ro\_II,445 {2/7}     marucchabdasya upsaṅkhyānam kartavyam . \\
\noindent \textbf{(P\_1,4.58-59.2)} KA\_I,341.19-23 Ro\_II,445 {3/7}     maruddatto marutyaḥ . \\
\noindent \textbf{(P\_1,4.58-59.2)} KA\_I,341.19-23 Ro\_II,445 {4/7}     aca upasargāt iti tattvam yathā syāt . \\
\noindent \textbf{(P\_1,4.58-59.2)} KA\_I,341.19-23 Ro\_II,445 {5/7}     śracchabdasya upasamkhyānam . \\
\noindent \textbf{(P\_1,4.58-59.2)} KA\_I,341.19-23 Ro\_II,445 {6/7}     śracchabdasya upasamkhyānam kartavyam . \\
\noindent \textbf{(P\_1,4.58-59.2)} KA\_I,341.19-23 Ro\_II,445 {7/7}     śraddhā . \\
\noindent \textbf{(P\_1,4.60.1)} kārikāśabdasya . \\
\noindent \textbf{(P\_1,4.60.1)} kārikāśabdasya upasaṅkhyānam kartavyam . \\
\noindent \textbf{(P\_1,4.60.1)} kārikākṛtya . \\
\noindent \textbf{(P\_1,4.60.1)} punaścanasau chandasi . \\
\noindent \textbf{(P\_1,4.60.1)} punaścanasau chandasi gatisañjñau bhavataḥ iti vaktavyam . \\
\noindent \textbf{(P\_1,4.60.1)} punarutsyūtam vāsaḥ deyam . \\
\noindent \textbf{(P\_1,4.60.1)} punarniṣkṛtaḥ rathaḥ . \\
\noindent \textbf{(P\_1,4.60.1)} uśik dūtaḥ canohitaḥ . \\
\noindent \textbf{(P\_1,4.60.2)} gatyupasargasañjñāḥ kriyāyoge yatkriyāyuktāḥ tam prati iti vacanam . \\
\noindent \textbf{(P\_1,4.60.2)} gatyupasargasañjñāḥ kriyāyoge yatkriyāyuktāḥ tam prati gatyupasargasañjñāḥ bhavanti iti vaktavyam . \\
\noindent \textbf{(P\_1,4.60.2)} kim prayojanam . \\
\noindent \textbf{(P\_1,4.60.2)} prayojanam ghañ ṣaṭvaṇatve . \\
\noindent \textbf{(P\_1,4.60.2)} ghañ . \\
\noindent \textbf{(P\_1,4.60.2)} pravṛddhaḥ bhāvaḥ prabhāvaḥ . \\
\noindent \textbf{(P\_1,4.60.2)} anupasarge iti pratiṣedhaḥ mā bhūt . \\
\noindent \textbf{(P\_1,4.60.2)} ṣatvam . \\
\noindent \textbf{(P\_1,4.60.2)} vigatāḥ secakāḥ asmāt grāmāt visecakaḥ grāmaḥ . \\
\noindent \textbf{(P\_1,4.60.2)} upasargāt iti ṣatvam mā bhūt . \\
\noindent \textbf{(P\_1,4.60.2)} ṇatvam . \\
\noindent \textbf{(P\_1,4.60.2)} pragatāḥ nāyakāḥ asmāt grāmāt pranāyakaḥ grāmaḥ . \\
\noindent \textbf{(P\_1,4.60.2)} upasargāditi ṇatvam mā bhūt . \\
\noindent \textbf{(P\_1,4.60.2)} vṛddhividhau ca dhātugrahaṇānarthakyam . vṛddhividhau ca dhātugrahaṇam anarthakam . \\
\noindent \textbf{(P\_1,4.60.2)} upasargāt ṛti dhātau iti . \\
\noindent \textbf{(P\_1,4.60.2)} tatra dhātugrahaṇasya etat prayojanam iha mā bhūt prarṣabham vanam iti . \\
\noindent \textbf{(P\_1,4.60.2)} kriyamāṇe ca api dhātugrahaṇe prarcchaka iti atra prāpnoti . \\
\noindent \textbf{(P\_1,4.60.2)} yatkriyāyuktāḥ tam prati iti vacanāt na bhavati . \\
\noindent \textbf{(P\_1,4.60.2)} vadvidhnabhāvābīttvasvāṅgādisvaraṇatveṣu doṣaḥ bhavati . \\
\noindent \textbf{(P\_1,4.60.2)} vadvidhi. yat udvataḥ nivataḥ yāsi bapsat . \\
\noindent \textbf{(P\_1,4.60.2)} vadvidhi . \\
\noindent \textbf{(P\_1,4.60.2)} nasbhāva . \\
\noindent \textbf{(P\_1,4.60.2)} praṇasam mukham unnasam mukham . \\
\noindent \textbf{(P\_1,4.60.2)} nasbhāva . \\
\noindent \textbf{(P\_1,4.60.2)} abīttva .prepam parepam . \\
\noindent \textbf{(P\_1,4.60.2)} abittva . \\
\noindent \textbf{(P\_1,4.60.2)} svāṅgādisvara . \\
\noindent \textbf{(P\_1,4.60.2)} prasphik prodaraḥ . \\
\noindent \textbf{(P\_1,4.60.2)} svāṅgādisvara . \\
\noindent \textbf{(P\_1,4.60.2)} ṇatva . \\
\noindent \textbf{(P\_1,4.60.2)} pra ṇaḥ śūdraḥ pra ṇaḥ ācāryaḥ pra ṇaḥ rājā pra ṇaḥ vṛtrahā . \\
\noindent \textbf{(P\_1,4.60.2)} upasargāt iti ete vidhayaḥ na prāpnuvanti . \\
\noindent \textbf{(P\_1,4.60.2)} vadvidhinasbhāvabīttvasvāṅgadisvaraṇatveṣu vacanaprāmāṇyāt siddham . \\
\noindent \textbf{(P\_1,4.60.2)} anavakāśāḥ ete vidhayaḥ . \\
\noindent \textbf{(P\_1,4.60.2)} te vacanaprāmāṇyāt bhaviṣyanti . \\
\noindent \textbf{(P\_1,4.60.2)} suduroḥ pratiṣedhaḥ numvidhitatvaṣatvaṇatveṣu . \\
\noindent \textbf{(P\_1,4.60.2)} suduroḥ pratiṣedhaḥ numvidhitatvaṣatvaṇatveṣu vaktavyaḥ . \\
\noindent \textbf{(P\_1,4.60.2)} numvidhi : sulabham durlabham . \\
\noindent \textbf{(P\_1,4.60.2)} upasargāt iti num mā bhūt iti . \\
\noindent \textbf{(P\_1,4.60.2)} na sudurbhyām kevalābhyām . \\
\noindent \textbf{(P\_1,4.60.2)} iti etat na vaktavyam bhavati . \\
\noindent \textbf{(P\_1,4.60.2)} na etat asti prayojanam . \\
\noindent \textbf{(P\_1,4.60.2)} kriyate etat nyāse eva . \\
\noindent \textbf{(P\_1,4.60.2)} tatvam . \\
\noindent \textbf{(P\_1,4.60.2)} sudattam . \\
\noindent \textbf{(P\_1,4.60.2)} acaḥ upasargāt taḥ iti tatvam mā bhūt iti . \\
\noindent \textbf{(P\_1,4.60.2)} ṣatvam . \\
\noindent \textbf{(P\_1,4.60.2)} susiktam ghaṭaśatena sustutam ślokaśatena . \\
\noindent \textbf{(P\_1,4.60.2)} upasargāt iti ṣatvam mā bhūtiti . \\
\noindent \textbf{(P\_1,4.60.2)} suḥ pūjāyām iti etat na vaktavyam bhavati . \\
\noindent \textbf{(P\_1,4.60.2)} na etat asti prayojanam . \\
\noindent \textbf{(P\_1,4.60.2)} kriyata etat nyāse eva . \\
\noindent \textbf{(P\_1,4.60.2)} ṇatvam . \\
\noindent \textbf{(P\_1,4.60.2)} durnayam durnītamiti . \\
\noindent \textbf{(P\_1,4.60.2)} upasargāt iti ṇatvam mā bhūt iti . \\
\noindent \textbf{(P\_1,4.61)} kṛbhvastiyoge iti vaktavyam . \\
\noindent \textbf{(P\_1,4.61)} iha eva yathā syāt . \\
\noindent \textbf{(P\_1,4.61)} ūrīkṛtya ūrībhūya . \\
\noindent \textbf{(P\_1,4.61)} iha mā bhūt . \\
\noindent \textbf{(P\_1,4.61)} ūrī paktvā . \\
\noindent \textbf{(P\_1,4.61)} tat tarhi vaktavyam . \\
\noindent \textbf{(P\_1,4.61)} na vaktavyam . \\
\noindent \textbf{(P\_1,4.61)} kriyāyoge iti anuvartate na ca anyayā kriyayā ūryādicviḍācām yogaḥ asti . \\
\noindent \textbf{(P\_1,4.62.1)} katham idam vijñāyate . \\
\noindent \textbf{(P\_1,4.62.1)} iteḥ param itiparam na itiparam anitiparam iti āhosvit itiḥ paro yasmāt tat idam itiparam na itiparam anitiparamiti . \\
\noindent \textbf{(P\_1,4.62.1)} kim ca ataḥ . \\
\noindent \textbf{(P\_1,4.62.1)} yadi vijñāyata iteḥ param itiparam na itiparam anitiparam iti khāṭ iti kṛtvā niraṣṭhīvat iti atra prāpnoti . \\
\noindent \textbf{(P\_1,4.62.1)} atha itiḥ paro yasmāt tat idam itiparam na itiparam anitiparamiti śrauṣaṭ vauṣaṭ iti kṛtvā niraṣṭhīvat iti atra prāpnoti . \\
\noindent \textbf{(P\_1,4.62.1)} astu tāvat itiḥ paro yasmāt tat idam itiparam na itiparam anitiparamiti . \\
\noindent \textbf{(P\_1,4.62.1)} nanu ca uktam śrauṣaṭ vauṣaṭ iti kṛtvā niraṣṭhīvat iti atra prāpnoti iti . \\
\noindent \textbf{(P\_1,4.62.1)} na eṣaḥ doṣaḥ . \\
\noindent \textbf{(P\_1,4.62.1)} idam tāvat ayam praṣṭavyaḥ . \\
\noindent \textbf{(P\_1,4.62.1)} atha iha te prāk dhātoḥ iti katham gatimātrasya pūrvaprayogḥ bhavati . \\
\noindent \textbf{(P\_1,4.62.1)} upoddharati iti . \\
\noindent \textbf{(P\_1,4.62.1)} gatyākṛtiḥ pratinirdiśyate . \\
\noindent \textbf{(P\_1,4.62.1)} iha api tarhi anukaraṇākṛtiḥ nirdiśyate . \\
\noindent \textbf{(P\_1,4.62.2)} kimartham idam ucyate . \\
\noindent \textbf{(P\_1,4.62.2)} anukaraṇasya itikaraṇaparatvapratiṣedhaḥ aniṣṭaśabdanivṛttyarthaḥ . \\
\noindent \textbf{(P\_1,4.62.2)} anukaraṇasya itikaraṇaparatvapratiṣedhaḥ ucyate . \\
\noindent \textbf{(P\_1,4.62.2)} kim prayojanam . \\
\noindent \textbf{(P\_1,4.62.2)} aniṣṭaśabdanivṛttyarthaḥ . \\
\noindent \textbf{(P\_1,4.62.2)} aniṣṭhaśabdatā mā bhūt iti . \\
\noindent \textbf{(P\_1,4.62.2)} idam vicārayiṣyati teprāgdhātuvacanam prayoganiyamārtham vā syāt sañjñāniyamārtham vā iti . \\
\noindent \textbf{(P\_1,4.62.2)} tat yadā prayoganiyamārtham tadā aniṣṭhaśabdanivṛttyartham idam vaktavyam . \\
\noindent \textbf{(P\_1,4.62.2)} yadā hi sañjñāniyamārtham tadā na doṣaḥ bhavati . \\
\noindent \textbf{(P\_1,4.63)} idam atibahu kriyate ādare anādare sat asat iti . \\
\noindent \textbf{(P\_1,4.63)} ādāre sat iti eva siddham . \\
\noindent \textbf{(P\_1,4.63)} katham asatkṛtya iti . \\
\noindent \textbf{(P\_1,4.63)} tadantividhinā bhaviṣyati . \\
\noindent \textbf{(P\_1,4.63)} kena idānīm anādare bhaviṣyati . \\
\noindent \textbf{(P\_1,4.63)} nañā ādarapratiṣedham vijñāsyāmaḥ . \\
\noindent \textbf{(P\_1,4.63)} nādare anādare iti . \\
\noindent \textbf{(P\_1,4.63)} na evam śakyam . \\
\noindent \textbf{(P\_1,4.63)} ādaraprasaṅge eva hi syāt . \\
\noindent \textbf{(P\_1,4.63)} anādaraprasaṅge na syāt . \\
\noindent \textbf{(P\_1,4.63)} anādaragrahaṇe punaḥ kriyamāṇe bahuvrīhiḥ ayam vijñāyate . \\
\noindent \textbf{(P\_1,4.63)} avidyamānādare anādare iti . \\
\noindent \textbf{(P\_1,4.63)} tasmāt anādaragrahaṇam kartavyam . \\
\noindent \textbf{(P\_1,4.63)} asataḥ tu tadantavidhinā siddham . \\
\noindent \textbf{(P\_1,4.65)} antaḥśabdasya āṅkividhisamāsaṇatveṣu pasaṅkhyānam . \\
\noindent \textbf{(P\_1,4.65)} antaḥśabdasya āṅkividhisamāsaṇatveṣu pasaṅkhyānam kartavyam . \\
\noindent \textbf{(P\_1,4.65)} aṅ . \\
\noindent \textbf{(P\_1,4.65)} antardhā . \\
\noindent \textbf{(P\_1,4.65)} kividhiḥ . \\
\noindent \textbf{(P\_1,4.65)} antardhiḥ . \\
\noindent \textbf{(P\_1,4.65)} samāsaḥ . \\
\noindent \textbf{(P\_1,4.65)} antarhatya . \\
\noindent \textbf{(P\_1,4.65)} ṇatvam . \\
\noindent \textbf{(P\_1,4.65)} antarhaṇyāt gobhyo gāḥ . \\
\noindent \textbf{(P\_1,4.74)} sākṣātprabhṛtiṣu cvyarthagrahaṇam . \\
\noindent \textbf{(P\_1,4.74)} sākṣātprabhṛtiṣu cvyarthagrahaṇam kartavyam . \\
\noindent \textbf{(P\_1,4.74)} asākṣātsākṣātkṛtvā sākṣātkṛtya . \\
\noindent \textbf{(P\_1,4.74)} yadā hi sākṣāt eva kim cit kriyate tadā mā bhūt iti . \\
\noindent \textbf{(P\_1,4.74)} makārāntatvam ca gatisañjñāsanniyuktam . \\
\noindent \textbf{(P\_1,4.74)} makārāntatvam ca gatisañjñāsanniyogena vaktavyam . \\
\noindent \textbf{(P\_1,4.74)} lavaṇaṅkṛtya . \\
\noindent \textbf{(P\_1,4.74)} tatra cvipratiṣedhaḥ . \\
\noindent \textbf{(P\_1,4.74)} tatra cvyantasya pratiṣedhaḥ vaktavyaḥ . \\
\noindent \textbf{(P\_1,4.74)} lavaṇīkṛtya . \\
\noindent \textbf{(P\_1,4.74)} na vā pūrveṇa kṛtatvāt ṇa vā vaktavyam . \\
\noindent \textbf{(P\_1,4.74)} kim kāraṇam . \\
\noindent \textbf{(P\_1,4.74)} pūrveṇa kṛtatvāt . \\
\noindent \textbf{(P\_1,4.74)} astu anena vibhāṣā . \\
\noindent \textbf{(P\_1,4.74)} pūrveṇa nityaḥ bhaviṣyati . \\
\noindent \textbf{(P\_1,4.74)} idam tarhi prayojanam . \\
\noindent \textbf{(P\_1,4.74)} makārāntatvam ca gatisañjñāsanniyuktam iti uktam . \\
\noindent \textbf{(P\_1,4.74)} tat cvyantasya mā bhūt iti . \\
\noindent \textbf{(P\_1,4.74)} etat api na asti prayojanam . \\
\noindent \textbf{(P\_1,4.74)} lavaṇaśabdasya ayam vibhāṣā lavaṇamśabda ādeśaḥ kriyate . \\
\noindent \textbf{(P\_1,4.74)} yadi ca lavaṇī śabdasya api vibhāṣā lavaṇamśabdaḥ ādeśaḥ bhavati na kim cid duṣyati . \\
\noindent \textbf{(P\_1,4.74)} traiśabdyam ca ha sādhyam . \\
\noindent \textbf{(P\_1,4.74)} | tacca evam sati siddham bhavati iti . \\
\noindent \textbf{(P\_1,4.80)} kimidam prāgdhātuvacanam prayoganiyamārtham : ete prāk eva dhātoḥ prayoktavyāḥ . \\
\noindent \textbf{(P\_1,4.80)} āhosvit sañjñāniyamārtham : ete prāk ca akprāk ca prayoktavyāḥ , prāk prayujyamānānām gatisañjñā bhavati iti . \\
\noindent \textbf{(P\_1,4.80)} kaḥ ca atra viśeṣaḥ . \\
\noindent \textbf{(P\_1,4.80)} prāgdhātuvacanam prayoganiyamārtham iti cet anukaraṇasya itikaraṇaparapratiṣedhaḥ aniṣṭaśabdanivṛttyarthaḥ . \\
\noindent \textbf{(P\_1,4.80)} prāgdhātuvacanam prayoganiyamārtham iti cet anukaraṇasya itikaraṇaparapratiṣedhaḥ vaktavyaḥ . \\
\noindent \textbf{(P\_1,4.80)} kim prayojanam . \\
\noindent \textbf{(P\_1,4.80)} aniṣṭaśabdanivṛttyarthaḥ . \\
\noindent \textbf{(P\_1,4.80)} aniṣṭaśabatā mā bhūt iti . \\
\noindent \textbf{(P\_1,4.80)} chandasi paravyavahitavacanam ca . \\
\noindent \textbf{(P\_1,4.80)} chandasi pare api vyavahitāḥ ca iti vaktavyam . \\
\noindent \textbf{(P\_1,4.80)} sañjñāniyame siddham . \\
\noindent \textbf{(P\_1,4.80)} sañjñāniyame siddham etat bhavati . \\
\noindent \textbf{(P\_1,4.80)} astu tarhi sañjñāniyamaḥ . \\
\noindent \textbf{(P\_1,4.80)} ubhayoḥ anarthakam vacanam aniṣṭādarśanāt . \\
\noindent \textbf{(P\_1,4.80)} ubhayoḥ api pakṣayoḥ vacanamanarthakam . \\
\noindent \textbf{(P\_1,4.80)} kim kāraṇam . \\
\noindent \textbf{(P\_1,4.80)} aniṣṭādarśanāt . \\
\noindent \textbf{(P\_1,4.80)} na hi kaḥ citprapacati iti prayoktavye pacatipra iti prayuṅkte . \\
\noindent \textbf{(P\_1,4.80)} yadi ca aniṣṭam dṛśyeta tataḥ yatnārham syāt . \\
\noindent \textbf{(P\_1,4.80)} upasarjanasannipāte tu pūrvaparavyavasthārtham . \\
\noindent \textbf{(P\_1,4.80)} upasarjanasannipāte tu pūrvaparavyavasthārtham etat vaktavyam . \\
\noindent \textbf{(P\_1,4.80)} ṛṣabham kūlamudrujam ṛṣabham kūlamudvaham . \\
\noindent \textbf{(P\_1,4.80)} atra gateḥ prāk dhātoḥ prayogaḥ yathā syāt . \\
\noindent \textbf{(P\_1,4.80)} yadi upasarjanasannipāte pūrvaparavyavasthārtham idam ucyate sukaṭamkarāṇi vīraṇān i iti atra gateḥ prāk dhātoḥ prayogaḥ prāpnoti . \\
\noindent \textbf{(P\_1,4.80)} ācāryapravṛttiḥ jñāpayati na atra gateḥ prākprayogaḥ bhavati iti yat ayam īṣadduḥsuṣu kṛcchrākṛcchārtheṣu khal iti khakāram anubandham karoti . \\
\noindent \textbf{(P\_1,4.80)} katham kṛtvā jñāpakam . \\
\noindent \textbf{(P\_1,4.80)} khitkaraṇe etat prayojanam khiti iti mum yathā syāt iti . \\
\noindent \textbf{(P\_1,4.80)} yadi ca atra gateḥ prākprayogaḥ syāt khitkaraṇam anarthakam syāt . \\
\noindent \textbf{(P\_1,4.80)} astu atra mum . \\
\noindent \textbf{(P\_1,4.80)} anavyayasya iti pratiṣedhaḥ bhaviṣyati . \\
\noindent \textbf{(P\_1,4.80)} paśyati tu ācāryaḥ na atra gateḥ prāk dhatoḥ prayogaḥ bhavati iti tataḥ khakāram anubandham karoti . \\
\noindent \textbf{(P\_1,4.80)} na etat asti jñāpakam . \\
\noindent \textbf{(P\_1,4.80)} yadi api atra gateḥ prākprayogaḥ syāt syātevātra mumāgamaḥ . \\
\noindent \textbf{(P\_1,4.80)} katham . \\
\noindent \textbf{(P\_1,4.80)} kṛdgrahaṇe gatikārakapūrvasya api grahaṇam bhavati iti . \\
\noindent \textbf{(P\_1,4.80)} tasmāt na arthaḥ evamarthena prāgdhātuvacanena . \\
\noindent \textbf{(P\_1,4.80)} katham ṛṣabham kūlamudrujam ṛṣabham kūlamudvaham . \\
\noindent \textbf{(P\_1,4.80)} na eṣaḥ doṣaḥ . \\
\noindent \textbf{(P\_1,4.80)} na eṣaḥ udiḥ upapadam . \\
\noindent \textbf{(P\_1,4.80)} kim tarhi . \\
\noindent \textbf{(P\_1,4.80)} viśeṣaṇam . \\
\noindent \textbf{(P\_1,4.80)} udi kūle rujivahoḥ . \\
\noindent \textbf{(P\_1,4.80)} utpūrvābhyām rujivahibhyām kūle upapade iti . \\
\noindent \textbf{(P\_1,4.83)} kimartham mahatī sañjñā kriyate . \\
\noindent \textbf{(P\_1,4.83)} anvarthasañjñā yathā vijñāyeta . \\
\noindent \textbf{(P\_1,4.83)} karma proktavantaḥ karmapravacanīyāḥ iti . \\
\noindent \textbf{(P\_1,4.83)} ke punaḥ karma proktavantaḥ . \\
\noindent \textbf{(P\_1,4.83)} ye samprati kriyām na āhuḥ . \\
\noindent \textbf{(P\_1,4.83)} ke ca samprati kriyām na āhuḥ . \\
\noindent \textbf{(P\_1,4.83)} ye aprayujyamānasya kriyām āhuḥ te karmapravacanīyāḥ . \\
\noindent \textbf{(P\_1,4.84)} kimartham idam ucyate . \\
\noindent \textbf{(P\_1,4.84)} karmapravacanīyasañjñā yathā syāt . \\
\noindent \textbf{(P\_1,4.84)} gatyupasargasañjñe mā bhūtām iti . \\
\noindent \textbf{(P\_1,4.84)} kim ca syāt . \\
\noindent \textbf{(P\_1,4.84)} śākalyasya samhitām anu prāvarṣat : gatiḥ gatau iti nighātaḥ prasajyeta . \\
\noindent \textbf{(P\_1,4.84)} yadi evam veḥ api karmapravacanīyasañjñā vaktavyā . \\
\noindent \textbf{(P\_1,4.84)} veḥ api nighātaḥ na iṣyate : prādeśam prādeśam viparilikhati . \\
\noindent \textbf{(P\_1,4.84)} asti atra viśeṣaḥ . \\
\noindent \textbf{(P\_1,4.84)} na atra veḥ likhim prati kriyāyogaḥ . \\
\noindent \textbf{(P\_1,4.84)} kim tarhi . \\
\noindent \textbf{(P\_1,4.84)} aprayujyamānam . \\
\noindent \textbf{(P\_1,4.84)} prādeśam prādeśam vimāya parilikhati iti . \\
\noindent \textbf{(P\_1,4.84)} yadi evam anoḥ api karmapravacanīyasañjñayā na arthaḥ . \\
\noindent \textbf{(P\_1,4.84)} anoḥ api hi na vṛṣim prati kriyāyogaḥ . \\
\noindent \textbf{(P\_1,4.84)} kim tarhi aprayujyamānam . \\
\noindent \textbf{(P\_1,4.84)} śākalyena sukṛtām samhitām anuviśamya devaḥ prāvarṣat . \\
\noindent \textbf{(P\_1,4.84)} idam tarhi prayojanam dvitīyā yathā syāt karmapravacanīyayukte dvitīyā iti . \\
\noindent \textbf{(P\_1,4.84)} ataḥ uttaram paṭhati . \\
\noindent \textbf{(P\_1,4.84)} anurlakṣaṇevacanānarthakyam sāmānyakṛtatvāt . \\
\noindent \textbf{(P\_1,4.84)} anurlakṣaṇevacanārthakyam . \\
\noindent \textbf{(P\_1,4.84)} kim kāraṇam . \\
\noindent \textbf{(P\_1,4.84)} sāmānyakṛtatvāt . \\
\noindent \textbf{(P\_1,4.84)} sāmānyena eva atra karmapravacanīyasañjñā bhaviṣyati lakṣaṇetthambhūtākhyānabhāgavīpsāsu pratiparyanavaḥ iti . \\
\noindent \textbf{(P\_1,4.84)} hetvartham tu vacanam . \\
\noindent \textbf{(P\_1,4.84)} hetvartham idam vaktavyam . \\
\noindent \textbf{(P\_1,4.84)} hetuḥ śākalyasya samhitā varṣasya na lakṣaṇam . \\
\noindent \textbf{(P\_1,4.84)} kim vaktavyam etat . \\
\noindent \textbf{(P\_1,4.84)} na hi . \\
\noindent \textbf{(P\_1,4.84)} katham anucyamānam gaṃsyate . \\
\noindent \textbf{(P\_1,4.84)} lakṣaṇam hi nāma saḥ bhavati yena punaḥ punaḥ lakṣyate na yaḥ sakṛdapi nimittatvāya kalpate . \\
\noindent \textbf{(P\_1,4.84)} sakṛt ca asau śākalyena sukṛtām samhitām anuśimya devaḥ prāvarṣat . \\
\noindent \textbf{(P\_1,4.84)} saḥ tarhi tathā . \\
\noindent \textbf{(P\_1,4.84)} nirdeśaḥ kartavyaḥ anuḥ hetau iti . \\
\noindent \textbf{(P\_1,4.84)} atha idānīm lakṣaṇena hetuḥ api vyāptaḥ na arthaḥ anena . \\
\noindent \textbf{(P\_1,4.84)} lakṣaṇena hetuḥ api vyāptaḥ . \\
\noindent \textbf{(P\_1,4.84)} na hi avaśyam tat eva lakṣaṇam bhavati yena punaḥ punaḥ lakṣyate . \\
\noindent \textbf{(P\_1,4.84)} kim tarhi . \\
\noindent \textbf{(P\_1,4.84)} yat sakṛt api nimittvāya kalpate tat api lakṣaṇam bhavati . \\
\noindent \textbf{(P\_1,4.84)} tat yathā api bhavān kamaṇḍulapāṇim chātram adrakṣīt iti . \\
\noindent \textbf{(P\_1,4.84)} sakṛt āsau kamaṇḍalupāṇiḥ chātraḥ dṛṣṭaḥ tasya tat eva lakṣaṇam bhavati . \\
\noindent \textbf{(P\_1,4.84)} tat eva tarhi prayojanam dvitīyā yathā syāt karmapravacanīyayukte dvitīyā iti . \\
\noindent \textbf{(P\_1,4.84)} etat api na asti prayojanam . \\
\noindent \textbf{(P\_1,4.84)} siddhā atra dvitīyā karmapravacanīyayukte iti eva . \\
\noindent \textbf{(P\_1,4.84)} na sidhyati . \\
\noindent \textbf{(P\_1,4.84)} paratvāt hetutvāśrayā tṛtīyā prāpnoti . \\
\noindent \textbf{(P\_1,4.89)} āṅ maryādābhividhyoḥ iti vaktavyam . \\
\noindent \textbf{(P\_1,4.89)} iha api yathā syāt ākumāram yaśaḥ pāṇineḥ iti . \\
\noindent \textbf{(P\_1,4.89)} tat tarhi vaktavyam . \\
\noindent \textbf{(P\_1,4.89)} na vaktavyam . \\
\noindent \textbf{(P\_1,4.89)} maryādāvacane iti eva siddham . \\
\noindent \textbf{(P\_1,4.89)} eṣā asya yaśasaḥ maryādā . \\
\noindent \textbf{(P\_1,4.90)} kasya lakṣaṇadayaḥ arthāḥ nirdiśyante . \\
\noindent \textbf{(P\_1,4.90)} vṛkṣādīnām . \\
\noindent \textbf{(P\_1,4.90)} kimartham punaḥ idam ucyate . \\
\noindent \textbf{(P\_1,4.90)} karmapravacanīyasañjñā yathā syāt . \\
\noindent \textbf{(P\_1,4.90)} gatyupasargasañjñe mā bhūtām iti . \\
\noindent \textbf{(P\_1,4.90)} na etat asti prayojanam . \\
\noindent \textbf{(P\_1,4.90)} yatkriyāyuktāḥ tam prati gatyupasargasañjñe bhavataḥ na ca vṛkṣādīn prati kriyāyogaḥ . \\
\noindent \textbf{(P\_1,4.90)} idam tarhi prayojanam dvitīyā yathā syāt karmapravacanīyayukte dvitīyā iti . \\
\noindent \textbf{(P\_1,4.90)} vṛkṣam prati vidyotate . \\
\noindent \textbf{(P\_1,4.90)} vṛkṣamanu vidyotate iti . \\
\noindent \textbf{(P\_1,4.93)} kimartham adhiparyoḥ anarthakayoḥ karmapravacanīyasañjñā ucyate . \\
\noindent \textbf{(P\_1,4.93)} karmapravacanīyasañjñā yathā syāt . \\
\noindent \textbf{(P\_1,4.93)} gatyupasargasañjñe mā bhūtām iti . \\
\noindent \textbf{(P\_1,4.93)} na etat asti prayojanam . \\
\noindent \textbf{(P\_1,4.93)} yatkriyāyuktāḥ tam prati gatyupasargasañjñe bhavataḥ anarthakau ca imau . \\
\noindent \textbf{(P\_1,4.93)} idam tarhi prayojanam pañcamī yathā syāt pañcamī apāṅparibhiḥ iti . \\
\noindent \textbf{(P\_1,4.93)} kutaḥ paryāgamyata iti . \\
\noindent \textbf{(P\_1,4.93)} siddhā atra pañcamī apādāne iti eva . \\
\noindent \textbf{(P\_1,4.93)} ātaḥ ca apādānapañcamī eṣā . \\
\noindent \textbf{(P\_1,4.93)} yatra api adhiśabdena yoge pañcamī na vidhīyate tatra api śrūyate . \\
\noindent \textbf{(P\_1,4.93)} kutaḥ adhyāgamyata iti . \\
\noindent \textbf{(P\_1,4.93)} evam tarhi siddhe sati yat anarthakayoḥ gatyupasargasañjñābādhikām karmapravacanīyasañjñām śāsti tat jñāpayati ācāryaḥ anarthakānām api eṣām bhavati arthavatkṛtam iti . \\
\noindent \textbf{(P\_1,4.93)} kim etasya jñāpane prayojanam . \\
\noindent \textbf{(P\_1,4.93)} nipātasya anarthakasya prātipadikatvam coditam . \\
\noindent \textbf{(P\_1,4.93)} tat na vaktavyam bhavati . \\
\noindent \textbf{(P\_1,4.93)} atha vā na eva imau anarthakau . \\
\noindent \textbf{(P\_1,4.93)} kim tarhi anarthakau iti ucyate . \\
\noindent \textbf{(P\_1,4.93)} anarthāntarvācinau anarthakau . \\
\noindent \textbf{(P\_1,4.93)} dhātunā uktām kriyām āhatuḥ . \\
\noindent \textbf{(P\_1,4.93)} tad aviśiṣṭam bhavati yathā śaṅkhe payaḥ . \\
\noindent \textbf{(P\_1,4.93)} yadi evam dhātunā uktatvāt tasyārthasya upasargaprayogo na prāpnoti uktārthānām aprayogaḥ iti . \\
\noindent \textbf{(P\_1,4.93)} uktārthānāmapi prayogaḥ dṛśyate . \\
\noindent \textbf{(P\_1,4.93)} tat yathā apūpau dvau ānaya . \\
\noindent \textbf{(P\_1,4.93)} brāhmaṇau dvau anaya iti . \\
\noindent \textbf{(P\_1,4.96)} iha kasmāt na bhavati . \\
\noindent \textbf{(P\_1,4.96)} sarpiṣaḥ api syāt . \\
\noindent \textbf{(P\_1,4.96)} gomūtrasya api syāt . \\
\noindent \textbf{(P\_1,4.96)} kim ca syāt . \\
\noindent \textbf{(P\_1,4.96)} dvitīyā api prasajyeta karmapravacanīyayukte dvitīyā iti . \\
\noindent \textbf{(P\_1,4.96)} na eṣaḥ doṣaḥ . \\
\noindent \textbf{(P\_1,4.96)} na ime apyarthāḥ nirdiśyante . \\
\noindent \textbf{(P\_1,4.96)} kim tarhi . \\
\noindent \textbf{(P\_1,4.96)} parapadārthāḥ ime nirdiśyante . \\
\noindent \textbf{(P\_1,4.96)} eteṣu artheṣu yat padam vartate tat prati apiḥ karmapravacanīyasañjñaḥ bhavati iti . \\
\noindent \textbf{(P\_1,4.96)} atha vā yat atra karmapravacanīyayuktam na adaḥ prayujyate . \\
\noindent \textbf{(P\_1,4.96)} kim punaḥ tat . \\
\noindent \textbf{(P\_1,4.96)} binduḥ . \\
\noindent \textbf{(P\_1,4.96)} bindoḥ tarhi kasmāt na bhavati . \\
\noindent \textbf{(P\_1,4.96)} upapadavibhakteḥ kārakavibhaktiḥ balīyasī iti prathamā bhaviṣyati iti . \\
\noindent \textbf{(P\_1,4.97)} adhirīśvaravacane uktam . \\
\noindent \textbf{(P\_1,4.97)} kim uktam . \\
\noindent \textbf{(P\_1,4.97)} yasya ca īśvaravacanam iti kartṛnirdeśaḥ cet avacanāt siddham . \\
\noindent \textbf{(P\_1,4.97)} prathamānupapattiḥ tu . \\
\noindent \textbf{(P\_1,4.97)} svavacanāt siddham iti . \\
\noindent \textbf{(P\_1,4.97)} adhiḥ svam prati karmapravacanīyasañjñaḥ bhavati iti vaktavyam . \\
\noindent \textbf{(P\_1,4.99)} lādeśe parasmaipadagrahaṇam puruṣabādhitatvāt . \\
\noindent \textbf{(P\_1,4.99)} lādeśe parasmaipadagrahaṇam kartavyam . \\
\noindent \textbf{(P\_1,4.99)} kim kāraṇam . \\
\noindent \textbf{(P\_1,4.99)} puruṣabādhitatvāt . \\
\noindent \textbf{(P\_1,4.99)} iha vacane hi sañjñābādhanam . \\
\noindent \textbf{(P\_1,4.99)} iha hi kriyamāṇe anavakāśā puruṣasañjñā parasmaipadasañjñām bādheta . \\
\noindent \textbf{(P\_1,4.99)} parasmaipadasañjñā api anavakāśā . \\
\noindent \textbf{(P\_1,4.99)} sā vacanāt bhaviṣyati . \\
\noindent \textbf{(P\_1,4.99)} sāvakāśā parasamaipadasañjñā . \\
\noindent \textbf{(P\_1,4.99)} kaḥ vakāśaḥ . \\
\noindent \textbf{(P\_1,4.99)} śatṛkkvasū avakāśaḥ . \\
\noindent \textbf{(P\_1,4.99)} sici vṛddhau tu parasmaipadagrahaṇam jñāpakam puruṣābādhakatvasya . \\
\noindent \textbf{(P\_1,4.99)} yat ayam sici vṛddhiḥ parasmaipadeṣu iti parasmaipadagrahaṇam karoti tat jñāpayati ācāryaḥ na puruṣasañjñā parasmaipadasañjñām bādhate iti . \\
\noindent \textbf{(P\_1,4.101)} prathamamadhyamottamasañjñāyām ātmanepadagrahaṇam samasaṅkhyārtham . prathamamadhyamottamasañjñāyām ātmanepadagrahaṇam kartavyam . \\
\noindent \textbf{(P\_1,4.101)} ātmanepadānām ca prathamamadhyamottamasañjñāḥ bhavanti iti vaktavyam . \\
\noindent \textbf{(P\_1,4.101)} kim prayojanam . \\
\noindent \textbf{(P\_1,4.101)} samasaṅkhyārtham . \\
\noindent \textbf{(P\_1,4.101)} saṅkhyātānudeśaḥ yathā syāt . \\
\noindent \textbf{(P\_1,4.101)} akriyamāṇe hi ātmanepadagrahaṇe tisraḥ sañjñāḥ ṣaṭ sañjñinaḥ . \\
\noindent \textbf{(P\_1,4.101)} vaiṣamyāt saṅkhyātānudeśaḥ na prāpnoti . \\
\noindent \textbf{(P\_1,4.101)} kriyamāṇe api ca ātmanepadagrahaṇe ānupūrvyavacanam ca . \\
\noindent \textbf{(P\_1,4.101)} ānupūrvyavacanam ca kartavyam . \\
\noindent \textbf{(P\_1,4.101)} akriyamaṇe hi kasya cit eva trikasya prathamasañjñā syāt kasya cit eva madhyamasañjñā kasya cit eva uttamasañjñā . \\
\noindent \textbf{(P\_1,4.101)} na vaikaśeṣanirdeśāt . \\
\noindent \textbf{(P\_1,4.101)} yat tāvat ucyate ātmanepadagrahaṇam kartavyam samasaṅkhyārtham iti . \\
\noindent \textbf{(P\_1,4.101)} tat na kartavyam . \\
\noindent \textbf{(P\_1,4.101)} sañjñāḥ api ṣat eva nirdiśyante . \\
\noindent \textbf{(P\_1,4.101)} katham . \\
\noindent \textbf{(P\_1,4.101)} ekaśeṣanirdesāt . \\
\noindent \textbf{(P\_1,4.101)} ekaśeṣanirdeśaḥ ayam . \\
\noindent \textbf{(P\_1,4.101)} atha etasmin ekaśeṣanirdeśe sati kim ayam kṛtaikaśeṣāṇām dvandvaḥ . \\
\noindent \textbf{(P\_1,4.101)} prathamaḥ ca prathamaḥ ca prathamau madhyamaḥ ca madhyamaḥ ca madhyamau uttamaḥ ca uttamaḥ ca uttamau prathamau ca madhyamau ca uttamau ca prathamamadhyamottamāḥ iti . \\
\noindent \textbf{(P\_1,4.101)} āhosvit kṛtadvandvānām ekaśeṣaḥ . \\
\noindent \textbf{(P\_1,4.101)} prathamau ca madhyamaḥ ca uttamaḥ ca prathamamadhyamottamāḥ . \\
\noindent \textbf{(P\_1,4.101)} prathamamadhyamotttamāḥ ca prathamamadhyamottamāḥ ca prathamamadhyamottamāḥ iti . \\
\noindent \textbf{(P\_1,4.101)} kim ca ataḥ . \\
\noindent \textbf{(P\_1,4.101)} yadi kṛtaikeśeṣāṇām dvandvaḥ prathamamadhyamayoḥ prathamasañjñā prāpnoti . \\
\noindent \textbf{(P\_1,4.101)} uttamaprathamayoḥ madhyamasañjñā prāpnoti . \\
\noindent \textbf{(P\_1,4.101)} madhyamottamayoḥ uttamasañjñā prāpnoti . \\
\noindent \textbf{(P\_1,4.101)} atha kṛtadvandvānāmekaśeṣo na doṣaḥ bhavati . \\
\noindent \textbf{(P\_1,4.101)} yathā na doṣaḥ tathā astu . \\
\noindent \textbf{(P\_1,4.101)} kim punaḥ atra nyāyyam . \\
\noindent \textbf{(P\_1,4.101)} ubhayam iti āha . \\
\noindent \textbf{(P\_1,4.101)} ubhayam hi dṛśyate . \\
\noindent \textbf{(P\_1,4.101)} tat yathā . \\
\noindent \textbf{(P\_1,4.101)} bahu śaktikiṭakam bahūni śaktikiṭakāni bahu sthālīpiṭharam bahūni sthālīpiṭharāṇi . \\
\noindent \textbf{(P\_1,4.101)} yat api ucyate kriyamāṇe api ātmanepadagrahaṇe ānurpūrvyavacanam kartavyam iti . \\
\noindent \textbf{(P\_1,4.101)} na kartavyam . \\
\noindent \textbf{(P\_1,4.101)} lokataḥ etat siddham . \\
\noindent \textbf{(P\_1,4.101)} tat yathā loke vihavyasya dvābhyām dvābhyām agniḥ upstheyaḥ iti . \\
\noindent \textbf{(P\_1,4.101)} na ca ucyate ānupūrvyeṇa iti ānupūrvyeṇa ca upasthīyata iti . \\
\noindent \textbf{(P\_1,4.104)} trīṇi trīṇi iti anuvartate utāho na . \\
\noindent \textbf{(P\_1,4.104)} kim ca ataḥ . \\
\noindent \textbf{(P\_1,4.104)} yadi anuvartate aṣṭhanaḥ ā vibhaktau iti ātvam na prāpnoti . \\
\noindent \textbf{(P\_1,4.104)} atha nivṛttam prathamayoḥ pūrvasavarṇaḥ iti atra pratyayayoḥ eva grahaṇam prāpnoti . \\
\noindent \textbf{(P\_1,4.104)} yathā icchasi tathā astu . \\
\noindent \textbf{(P\_1,4.104)} astu tāvat anuvartate iti . \\
\noindent \textbf{(P\_1,4.104)} nanu ca uktam aṣṭhana ā vibhaktau iti ātvam na prāpnoti iti . \\
\noindent \textbf{(P\_1,4.104)} vacanāt bhaviṣyati . \\
\noindent \textbf{(P\_1,4.104)} atha vā punaḥ astu nivṛttam . \\
\noindent \textbf{(P\_1,4.104)} nanu ca uktam prathamayoḥ pūrvasavarṇaḥ iti atra pratyayayoḥ eva grahaṇam prāpnoti iti . \\
\noindent \textbf{(P\_1,4.104)} na eṣaḥ doṣaḥ . \\
\noindent \textbf{(P\_1,4.104)} aci iti anuvartate . \\
\noindent \textbf{(P\_1,4.104)} na cājādī prathamau pratyayau staḥ . \\
\noindent \textbf{(P\_1,4.104)} nanu ca evam vijñāyate ajādī yau prathamau ajādīnām vā yau prathamau iti . \\
\noindent \textbf{(P\_1,4.104)} yat tarhi tasmāt śasaḥ naḥ pumsi iti anukrāntam pūrvasavarṇadīrgham pratinirdiśati tat jñāpayati ācāryaḥ vibhaktyoḥ grahaṇam iti . \\
\noindent \textbf{(P\_1,4.104)} atha vā vacanagrahaṇam eva kuryāt . \\
\noindent \textbf{(P\_1,4.104)} aujasoḥ pūrvasavarṇaḥ iti . \\
\noindent \textbf{(P\_1,4.105,} 107-108.1) KA\_I,351.13-353.27 Ro\_II,471-476 {1/100}     kimartham idam ucyate . \\
\noindent \textbf{(P\_1,4.105,} 107-108.1) KA\_I,351.13-353.27 Ro\_II,471-476 {2/100}     yuṣmadasmaccheṣavacanam niyamārtham . \\
\noindent \textbf{(P\_1,4.105,} 107-108.1) KA\_I,351.13-353.27 Ro\_II,471-476 {3/100}     niyamārthaḥ ayam ārambhaḥ . \\
\noindent \textbf{(P\_1,4.105,} 107-108.1) KA\_I,351.13-353.27 Ro\_II,471-476 {4/100}     atha etasmin niyamārthe vijñāyamāne kim ayam upapadaniyamaḥ yuṣmadi madhyamaḥ eva asmadi uttamaḥ eva āhosvit puruṣaniyamaḥ yuṣmadi eva madhyamaḥ asmadi eva uttamaḥ iti . \\
\noindent \textbf{(P\_1,4.105,} 107-108.1) KA\_I,351.13-353.27 Ro\_II,471-476 {5/100}     kim ca ataḥ . \\
\noindent \textbf{(P\_1,4.105,} 107-108.1) KA\_I,351.13-353.27 Ro\_II,471-476 {6/100}     yadi puruṣaniyamaḥ śeṣagrahaṇam kartavyam śeṣe prathamaḥ iti . \\
\noindent \textbf{(P\_1,4.105,} 107-108.1) KA\_I,351.13-353.27 Ro\_II,471-476 {7/100}     kim kāraṇam . \\
\noindent \textbf{(P\_1,4.105,} 107-108.1) KA\_I,351.13-353.27 Ro\_II,471-476 {8/100}     madhyamottamau niyatau yuṣmadasmadī aniyate . \\
\noindent \textbf{(P\_1,4.105,} 107-108.1) KA\_I,351.13-353.27 Ro\_II,471-476 {9/100}     tatra prathamaḥ api prāpnoti . \\
\noindent \textbf{(P\_1,4.105,} 107-108.1) KA\_I,351.13-353.27 Ro\_II,471-476 {10/100}     tatra śeṣagrahaṇam kartavyam prathamaniyamārtham . \\
\noindent \textbf{(P\_1,4.105,} 107-108.1) KA\_I,351.13-353.27 Ro\_II,471-476 {11/100}     śeṣe eva prathamaḥ bhavati na anyatra iti . \\
\noindent \textbf{(P\_1,4.105,} 107-108.1) KA\_I,351.13-353.27 Ro\_II,471-476 {12/100}     atha api upapadaniyamaḥ evam api śeṣagrahaṇam kartavyam śeṣe prathamaḥ iti . \\
\noindent \textbf{(P\_1,4.105,} 107-108.1) KA\_I,351.13-353.27 Ro\_II,471-476 {13/100}     yuṣmadasmadī niyate madhyamottamau aniyatau tau śeṣe api prāpnutaḥ . \\
\noindent \textbf{(P\_1,4.105,} 107-108.1) KA\_I,351.13-353.27 Ro\_II,471-476 {14/100}     tatra śeṣagrahaṇam kartavyam śeṣaniyamārtham . \\
\noindent \textbf{(P\_1,4.105,} 107-108.1) KA\_I,351.13-353.27 Ro\_II,471-476 {15/100}     śeṣe prathamaḥ eva bhavati na anyaḥ iti . \\
\noindent \textbf{(P\_1,4.105,} 107-108.1) KA\_I,351.13-353.27 Ro\_II,471-476 {16/100}     upapadaniyame śeṣagrahaṇam śakyam akartum . \\
\noindent \textbf{(P\_1,4.105,} 107-108.1) KA\_I,351.13-353.27 Ro\_II,471-476 {17/100}     katham . \\
\noindent \textbf{(P\_1,4.105,} 107-108.1) KA\_I,351.13-353.27 Ro\_II,471-476 {18/100}     yuṣmadasmadī niyate madhyamottamau aniyatau tau śeṣe api prāpnutaḥ . \\
\noindent \textbf{(P\_1,4.105,} 107-108.1) KA\_I,351.13-353.27 Ro\_II,471-476 {19/100}     tataḥ vakṣyāmi prathamaḥ bhavati iti . \\
\noindent \textbf{(P\_1,4.105,} 107-108.1) KA\_I,351.13-353.27 Ro\_II,471-476 {20/100}     tat niyamārtham bhaviṣyati . \\
\noindent \textbf{(P\_1,4.105,} 107-108.1) KA\_I,351.13-353.27 Ro\_II,471-476 {21/100}     yatra prathamaḥ ca anyaḥ ca prāpnoti tatra prathamaḥ bhavati iti . \\
\noindent \textbf{(P\_1,4.105,} 107-108.1) KA\_I,351.13-353.27 Ro\_II,471-476 {22/100}     tatra yuṣmadasmadanyeṣu prathamapratiṣedhaḥ śeṣatvāt . \\
\noindent \textbf{(P\_1,4.105,} 107-108.1) KA\_I,351.13-353.27 Ro\_II,471-476 {23/100}     tatra yuṣmadasmadanyeṣu prathamasya pratiṣedhaḥ vaktavyaḥ . \\
\noindent \textbf{(P\_1,4.105,} 107-108.1) KA\_I,351.13-353.27 Ro\_II,471-476 {24/100}     tvam ca devadattaḥ ca pacathaḥ . \\
\noindent \textbf{(P\_1,4.105,} 107-108.1) KA\_I,351.13-353.27 Ro\_II,471-476 {25/100}     aham ca devadattaḥ ca pacāvaḥ . \\
\noindent \textbf{(P\_1,4.105,} 107-108.1) KA\_I,351.13-353.27 Ro\_II,471-476 {26/100}     kim kāraṇam . \\
\noindent \textbf{(P\_1,4.105,} 107-108.1) KA\_I,351.13-353.27 Ro\_II,471-476 {27/100}     śeṣatvāt . \\
\noindent \textbf{(P\_1,4.105,} 107-108.1) KA\_I,351.13-353.27 Ro\_II,471-476 {28/100}     śeṣe prathamaḥ iti prathamaḥ prāpnoti . \\
\noindent \textbf{(P\_1,4.105,} 107-108.1) KA\_I,351.13-353.27 Ro\_II,471-476 {29/100}     siddham tu yuṣmadasmadoḥ pratiṣedhāt . \\
\noindent \textbf{(P\_1,4.105,} 107-108.1) KA\_I,351.13-353.27 Ro\_II,471-476 {30/100}     siddham etat . \\
\noindent \textbf{(P\_1,4.105,} 107-108.1) KA\_I,351.13-353.27 Ro\_II,471-476 {31/100}     katham . \\
\noindent \textbf{(P\_1,4.105,} 107-108.1) KA\_I,351.13-353.27 Ro\_II,471-476 {32/100}     yuṣmadasmadoḥ pratiṣedhāt . \\
\noindent \textbf{(P\_1,4.105,} 107-108.1) KA\_I,351.13-353.27 Ro\_II,471-476 {33/100}     śeṣe prathamḥ yuṣmadasmadoḥ na iti vaktavyam . \\
\noindent \textbf{(P\_1,4.105,} 107-108.1) KA\_I,351.13-353.27 Ro\_II,471-476 {34/100}     yuṣmadi madhyamāt asmadi uttamaḥ pratiṣedhena . \\
\noindent \textbf{(P\_1,4.105,} 107-108.1) KA\_I,351.13-353.27 Ro\_II,471-476 {35/100}     yuṣmadi madhyamāt asmadi uttamaḥ iti etat bhavati vipratiṣedhena . \\
\noindent \textbf{(P\_1,4.105,} 107-108.1) KA\_I,351.13-353.27 Ro\_II,471-476 {36/100}     yuṣmadi madhyamaḥ iti asya avakāśaḥ tvam pacasi . \\
\noindent \textbf{(P\_1,4.105,} 107-108.1) KA\_I,351.13-353.27 Ro\_II,471-476 {37/100}     asmadi uttamaḥ iti asya avakāśaḥ aham pacāmi . \\
\noindent \textbf{(P\_1,4.105,} 107-108.1) KA\_I,351.13-353.27 Ro\_II,471-476 {38/100}     iha ubhayam prāpnoti tvam ca aham ca pacāvaḥ . \\
\noindent \textbf{(P\_1,4.105,} 107-108.1) KA\_I,351.13-353.27 Ro\_II,471-476 {39/100}     asmadi uttamaḥ iti etat bhavati virpratiṣedhena . \\
\noindent \textbf{(P\_1,4.105,} 107-108.1) KA\_I,351.13-353.27 Ro\_II,471-476 {40/100}     saḥ tarhi vipratiṣedhaḥ vaktvayaḥ . \\
\noindent \textbf{(P\_1,4.105,} 107-108.1) KA\_I,351.13-353.27 Ro\_II,471-476 {41/100}     na vaktavyaḥ . \\
\noindent \textbf{(P\_1,4.105,} 107-108.1) KA\_I,351.13-353.27 Ro\_II,471-476 {42/100}     tyadādīnām yat yat param tat tat śiṣyate iti evam asmadaḥ śeṣaḥ bhaviṣyati . \\
\noindent \textbf{(P\_1,4.105,} 107-108.1) KA\_I,351.13-353.27 Ro\_II,471-476 {43/100}     tatra asmadi uttamaḥ iti eva siddham . \\
\noindent \textbf{(P\_1,4.105,} 107-108.1) KA\_I,351.13-353.27 Ro\_II,471-476 {44/100}     anekaśeṣabhāvārtham tu . anekaśeṣabhāvārtham tu saḥ vipratiṣedhaḥ vaktavyaḥ . \\
\noindent \textbf{(P\_1,4.105,} 107-108.1) KA\_I,351.13-353.27 Ro\_II,471-476 {45/100}     yadā ca ekaśeṣaḥ na . \\
\noindent \textbf{(P\_1,4.105,} 107-108.1) KA\_I,351.13-353.27 Ro\_II,471-476 {46/100}     kadā ca ekaśeṣaḥ na . \\
\noindent \textbf{(P\_1,4.105,} 107-108.1) KA\_I,351.13-353.27 Ro\_II,471-476 {47/100}     sahavivakṣāyām ekaśeṣaḥ . \\
\noindent \textbf{(P\_1,4.105,} 107-108.1) KA\_I,351.13-353.27 Ro\_II,471-476 {48/100}     yadā na sahavivakṣā tada ekaśeṣaḥ na asti . \\
\noindent \textbf{(P\_1,4.105,} 107-108.1) KA\_I,351.13-353.27 Ro\_II,471-476 {49/100}     na vā yuṣmadasmadoḥ anekaśeṣabhāvāt tadadhikaraṇānām api anekaśeṣabhāvāt avipratiṣedhaḥ . \\
\noindent \textbf{(P\_1,4.105,} 107-108.1) KA\_I,351.13-353.27 Ro\_II,471-476 {50/100}     na vā arthaḥ vipratiṣedhena . \\
\noindent \textbf{(P\_1,4.105,} 107-108.1) KA\_I,351.13-353.27 Ro\_II,471-476 {51/100}     kim kāraṇam . \\
\noindent \textbf{(P\_1,4.105,} 107-108.1) KA\_I,351.13-353.27 Ro\_II,471-476 {52/100}     yuṣmadasmadoḥ anekaśeṣabhāvāt tadadhikaraṇānām api yuṣmadasmadadhikaraṇānām api ekaśeṣena na bhavitavyam . \\
\noindent \textbf{(P\_1,4.105,} 107-108.1) KA\_I,351.13-353.27 Ro\_II,471-476 {53/100}     tvam ca aham ca pacasi pacāmi ca iti . \\
\noindent \textbf{(P\_1,4.105,} 107-108.1) KA\_I,351.13-353.27 Ro\_II,471-476 {54/100}     kriyāpṛthaktve ca dravyapṛthaktvadarśanam anumānam uttaratra anekaśeṣabhāvasya . \\
\noindent \textbf{(P\_1,4.105,} 107-108.1) KA\_I,351.13-353.27 Ro\_II,471-476 {55/100}     kriyāpṛthaktve ca dravyapṛthaktvam dṛśyate . \\
\noindent \textbf{(P\_1,4.105,} 107-108.1) KA\_I,351.13-353.27 Ro\_II,471-476 {56/100}     tat yathā pacasi pacāmi ca tvam ca aham ca iti . \\
\noindent \textbf{(P\_1,4.105,} 107-108.1) KA\_I,351.13-353.27 Ro\_II,471-476 {57/100}     tat anumānam uttarayoḥ api kriyayoḥ ekaśeṣaḥ na bhavati iti . \\
\noindent \textbf{(P\_1,4.105,} 107-108.1) KA\_I,351.13-353.27 Ro\_II,471-476 {58/100}     evam ca kṛtvā saḥ pi adoṣaḥ bhavati yat uktam tatra yuṣmadasmadanyeṣu prathamapratiṣedhaḥ śeṣatvāt iti . \\
\noindent \textbf{(P\_1,4.105,} 107-108.1) KA\_I,351.13-353.27 Ro\_II,471-476 {59/100}     tatra api hi evam bhavitavyam tvam ca devadattaḥ ca pacasi pacati ca . \\
\noindent \textbf{(P\_1,4.105,} 107-108.1) KA\_I,351.13-353.27 Ro\_II,471-476 {60/100}     aham ca devadattaḥ ca pacāmi pacati ca iti . \\
\noindent \textbf{(P\_1,4.105,} 107-108.1) KA\_I,351.13-353.27 Ro\_II,471-476 {61/100}     yat tāvat ucyate na vā yuṣmadasmadoḥ anekaśeṣabhāvāt tadadhikaraṇānām api anekaśeṣabhāvāt avipratiṣedhaḥ iti . \\
\noindent \textbf{(P\_1,4.105,} 107-108.1) KA\_I,351.13-353.27 Ro\_II,471-476 {62/100}     dṛśyate hi yuṣmadasmadoḥ cānekaśeṣaḥ tadadhikaraṇānām ca ekaśeṣaḥ . \\
\noindent \textbf{(P\_1,4.105,} 107-108.1) KA\_I,351.13-353.27 Ro\_II,471-476 {63/100}     tat yathā tvam ca aham ca vṛttrahan ubhau samprayujyāvahai iti . \\
\noindent \textbf{(P\_1,4.105,} 107-108.1) KA\_I,351.13-353.27 Ro\_II,471-476 {64/100}     yat api ucyate kriyāpṛthaktve ca dravyapṛthaktvadarśanam anumānam uttaratra anekaśeṣabhāvasya iti . \\
\noindent \textbf{(P\_1,4.105,} 107-108.1) KA\_I,351.13-353.27 Ro\_II,471-476 {65/100}      kriyāpṛthaktve khalu api dravyaikaśeṣaḥ bhavati iti dṛśyate . \\
\noindent \textbf{(P\_1,4.105,} 107-108.1) KA\_I,351.13-353.27 Ro\_II,471-476 {66/100}     tat yathā akṣāḥ bhajyantām bhakṣyantām dīvyantām iti . \\
\noindent \textbf{(P\_1,4.105,} 107-108.1) KA\_I,351.13-353.27 Ro\_II,471-476 {67/100}     evam ca kṛtvā saḥ api doṣo bhavati yat uktam tatra yuṣmadasmadanyeṣu prathamapratiṣedhaḥ śeṣatvāt iti . \\
\noindent \textbf{(P\_1,4.105,} 107-108.1) KA\_I,351.13-353.27 Ro\_II,471-476 {68/100}     na eṣaḥ doṣaḥ . \\
\noindent \textbf{(P\_1,4.105,} 107-108.1) KA\_I,351.13-353.27 Ro\_II,471-476 {69/100}     parihṛtam etat siddham tu yuṣmadasmadoḥ pratiṣedhāt iti . \\
\noindent \textbf{(P\_1,4.105,} 107-108.1) KA\_I,351.13-353.27 Ro\_II,471-476 {70/100}     saḥ tarhi pratiṣedhaḥ vaktavyaḥ . \\
\noindent \textbf{(P\_1,4.105,} 107-108.1) KA\_I,351.13-353.27 Ro\_II,471-476 {71/100}     na vaktavyaḥ . \\
\noindent \textbf{(P\_1,4.105,} 107-108.1) KA\_I,351.13-353.27 Ro\_II,471-476 {72/100}     śeṣe prathamaḥ vidhīyate . \\
\noindent \textbf{(P\_1,4.105,} 107-108.1) KA\_I,351.13-353.27 Ro\_II,471-476 {73/100}      na hi śeṣaḥ ca anyaḥ ca śeṣagrahaṇena gṛhyate . \\
\noindent \textbf{(P\_1,4.105,} 107-108.1) KA\_I,351.13-353.27 Ro\_II,471-476 {74/100}     bhavet prathamaḥ na syān . \\
\noindent \textbf{(P\_1,4.105,} 107-108.1) KA\_I,351.13-353.27 Ro\_II,471-476 {75/100}     madhyamottamau api na prāpnutaḥ . \\
\noindent \textbf{(P\_1,4.105,} 107-108.1) KA\_I,351.13-353.27 Ro\_II,471-476 {76/100}     kim kāraṇam . \\
\noindent \textbf{(P\_1,4.105,} 107-108.1) KA\_I,351.13-353.27 Ro\_II,471-476 {77/100}      yuṣmadasmadoḥ upapadayoḥ madhyamottamau ucyete . \\
\noindent \textbf{(P\_1,4.105,} 107-108.1) KA\_I,351.13-353.27 Ro\_II,471-476 {78/100}      na ca yuṣmadasmadī anyaḥ ca yuṣmadasmadgrahaṇena gṛhyate . \\
\noindent \textbf{(P\_1,4.105,} 107-108.1) KA\_I,351.13-353.27 Ro\_II,471-476 {79/100}     yat atra yuṣmat yat ca asmat tattadāśrayau madhyamottamau bhaviṣyataḥ . \\
\noindent \textbf{(P\_1,4.105,} 107-108.1) KA\_I,351.13-353.27 Ro\_II,471-476 {80/100}     yathā eva tarhi yat atra yuṣmat yat ca smat tadāśrayau madhyamottamau bhavataḥ evam yaḥ atra śeṣaḥ tadāśrayaḥ prathamaḥ prāpnoti . \\
\noindent \textbf{(P\_1,4.105,} 107-108.1) KA\_I,351.13-353.27 Ro\_II,471-476 {81/100}     evam tarhi śeṣe upapade prathamaḥ vidhīyate . \\
\noindent \textbf{(P\_1,4.105,} 107-108.1) KA\_I,351.13-353.27 Ro\_II,471-476 {82/100}     upoccāri padam upapadam . \\
\noindent \textbf{(P\_1,4.105,} 107-108.1) KA\_I,351.13-353.27 Ro\_II,471-476 {83/100}     yat ca atra upoccāri na saḥ śeṣaḥ yaḥ ca śeṣaḥ na tat upoccāri . \\
\noindent \textbf{(P\_1,4.105,} 107-108.1) KA\_I,351.13-353.27 Ro\_II,471-476 {84/100}     bhavet prathamaḥ na syāt . \\
\noindent \textbf{(P\_1,4.105,} 107-108.1) KA\_I,351.13-353.27 Ro\_II,471-476 {85/100}     madhyamottamau api na prāpnutaḥ . \\
\noindent \textbf{(P\_1,4.105,} 107-108.1) KA\_I,351.13-353.27 Ro\_II,471-476 {86/100}     kim kāraṇam . \\
\noindent \textbf{(P\_1,4.105,} 107-108.1) KA\_I,351.13-353.27 Ro\_II,471-476 {87/100}     yuṣmadasmadoḥ upapadayoḥ madhyamottamau ucyete . \\
\noindent \textbf{(P\_1,4.105,} 107-108.1) KA\_I,351.13-353.27 Ro\_II,471-476 {88/100}     upoccāri padam upapadam . \\
\noindent \textbf{(P\_1,4.105,} 107-108.1) KA\_I,351.13-353.27 Ro\_II,471-476 {89/100}     yat ca atra upoccāri na te yuṣmadasmadī ye ca yuṣmadasmadī na tat upoccāri . \\
\noindent \textbf{(P\_1,4.105,} 107-108.1) KA\_I,351.13-353.27 Ro\_II,471-476 {90/100}     evam tarhi śeṣeṇa sāmānādhikaraṇye prathamaḥ vidhīyate . \\
\noindent \textbf{(P\_1,4.105,} 107-108.1) KA\_I,351.13-353.27 Ro\_II,471-476 {91/100}     na ca atra śeṣeṇa eva sāmānādhikaraṇyam . \\
\noindent \textbf{(P\_1,4.105,} 107-108.1) KA\_I,351.13-353.27 Ro\_II,471-476 {92/100}     bhavet prathamaḥ na syāt . \\
\noindent \textbf{(P\_1,4.105,} 107-108.1) KA\_I,351.13-353.27 Ro\_II,471-476 {93/100}     madhyamottamau api na prāpnutaḥ . \\
\noindent \textbf{(P\_1,4.105,} 107-108.1) KA\_I,351.13-353.27 Ro\_II,471-476 {94/100}     kim kāraṇam . \\
\noindent \textbf{(P\_1,4.105,} 107-108.1) KA\_I,351.13-353.27 Ro\_II,471-476 {95/100}     yuṣmadasmadbhyām sāmānādhikaraṇye madhyamottamau ucyete na ca atra yuṣmadasmadbhyām eva sāmānādhikaraṇyam . \\
\noindent \textbf{(P\_1,4.105,} 107-108.1) KA\_I,351.13-353.27 Ro\_II,471-476 {96/100}     evam tarhi tyadādīni sarvaiḥ nityam iti evam atra yuṣmadasmadoḥ śeṣaḥ bhaviṣyati . \\
\noindent \textbf{(P\_1,4.105,} 107-108.1) KA\_I,351.13-353.27 Ro\_II,471-476 {97/100}     tatra yuṣmadi madhyamaḥ asmadi uttamaḥ iti eva siddham . \\
\noindent \textbf{(P\_1,4.105,} 107-108.1) KA\_I,351.13-353.27 Ro\_II,471-476 {98/100}     na sidhyati . \\
\noindent \textbf{(P\_1,4.105,} 107-108.1) KA\_I,351.13-353.27 Ro\_II,471-476 {99/100}     sthānini api iti prathamaḥ prāpnoti . \\
\noindent \textbf{(P\_1,4.105,} 107-108.1) KA\_I,351.13-353.27 Ro\_II,471-476 {100/100}     tyadādīnām khalu api yat yat param tat tat śiṣyate iti yadā bhavataḥ śeṣaḥ tadā prathamaḥ prāpnoti \\
\noindent \textbf{(P\_1,4.105,} 107-108.2) KA\_I,353.27-354. 15 Ro\_II,477-478 {1/30}     yuṣmadi madhyamaḥ asmadi uttamaḥ iti eva ucyate. tau iha na prāpnutaḥ : paramatvam pacasi . \\
\noindent \textbf{(P\_1,4.105,} 107-108.2) KA\_I,353.27-354. 15 Ro\_II,477-478 {2/30}      paramāham pacāmi iti . \\
\noindent \textbf{(P\_1,4.105,} 107-108.2) KA\_I,353.27-354. 15 Ro\_II,477-478 {3/30}     tadantavidhinā bhaviṣyati . \\
\noindent \textbf{(P\_1,4.105,} 107-108.2) KA\_I,353.27-354. 15 Ro\_II,477-478 {4/30}     iha api tarhi tadantavidhinā prāpnutaḥ : atitvam pacati . \\
\noindent \textbf{(P\_1,4.105,} 107-108.2) KA\_I,353.27-354. 15 Ro\_II,477-478 {5/30}     atyaham pacati iti . \\
\noindent \textbf{(P\_1,4.105,} 107-108.2) KA\_I,353.27-354. 15 Ro\_II,477-478 {6/30}     ye ca api ete samānādhikaraṇavṛttayaḥ taddhitāḥ tatra ca madhyamottamau na prāpnutaḥ : tvattaraḥ pacasi mattaraḥ pacāmi iti . \\
\noindent \textbf{(P\_1,4.105,} 107-108.2) KA\_I,353.27-354. 15 Ro\_II,477-478 {7/30}     tvadrūpaḥ pacasi madrūpaḥ pacāmi iti . \\
\noindent \textbf{(P\_1,4.105,} 107-108.2) KA\_I,353.27-354. 15 Ro\_II,477-478 {8/30}     tvatkalpaḥ pacasi . \\
\noindent \textbf{(P\_1,4.105,} 107-108.2) KA\_I,353.27-354. 15 Ro\_II,477-478 {9/30}     matkalpaḥ pacāmi iti. evam tarhi yuṣmadvati asmadvati iti evam bhaviṣyati . \\
\noindent \textbf{(P\_1,4.105,} 107-108.2) KA\_I,353.27-354. 15 Ro\_II,477-478 {10/30}     iha api tarhi prāpnutaḥ : atitvam pacati . \\
\noindent \textbf{(P\_1,4.105,} 107-108.2) KA\_I,353.27-354. 15 Ro\_II,477-478 {11/30}     atyaham pacati iti . \\
\noindent \textbf{(P\_1,4.105,} 107-108.2) KA\_I,353.27-354. 15 Ro\_II,477-478 {12/30}     evam tarhi yuṣmadi sādhane asmadi sādhane iti evam bhaviṣyati . \\
\noindent \textbf{(P\_1,4.105,} 107-108.2) KA\_I,353.27-354. 15 Ro\_II,477-478 {13/30}     evam ca kṛtvā saḥ api adoṣaḥ bhavati yat uktam tatra yuṣmadasmadanyeṣu prathamapratiṣedhaḥ śeṣatvāt iti . \\
\noindent \textbf{(P\_1,4.105,} 107-108.2) KA\_I,353.27-354. 15 Ro\_II,477-478 {14/30}     atha vā prathamaḥ utsargaḥ kariṣyate . \\
\noindent \textbf{(P\_1,4.105,} 107-108.2) KA\_I,353.27-354. 15 Ro\_II,477-478 {15/30}     tasya yuṣmadasmadoḥ upapadayoḥ madhyamottamau apavādau bhaviṣyataḥ . \\
\noindent \textbf{(P\_1,4.105,} 107-108.2) KA\_I,353.27-354. 15 Ro\_II,477-478 {16/30}     tatra yuṣmadgandhaḥ ca asmadgandhaḥ ca asti iti kṛtvā madhyamottamau bhaviṣyataḥ . \\
\noindent \textbf{(P\_1,4.105,} 107-108.2) KA\_I,353.27-354. 15 Ro\_II,477-478 {17/30}     atha iha katham bhavitavyam . \\
\noindent \textbf{(P\_1,4.105,} 107-108.2) KA\_I,353.27-354. 15 Ro\_II,477-478 {18/30}     atvam tvam sampadyate tvadbhavati madbhavati iti . \\
\noindent \textbf{(P\_1,4.105,} 107-108.2) KA\_I,353.27-354. 15 Ro\_II,477-478 {19/30}     āhosvit tvadbhavasi madbhavāmi iti . \\
\noindent \textbf{(P\_1,4.105,} 107-108.2) KA\_I,353.27-354. 15 Ro\_II,477-478 {20/30}     tvadbhavati madbhavati iti evam bhavitavyam . \\
\noindent \textbf{(P\_1,4.105,} 107-108.2) KA\_I,353.27-354. 15 Ro\_II,477-478 {21/30}     madhyamottamau kasmāt na bhavataḥ . \\
\noindent \textbf{(P\_1,4.105,} 107-108.2) KA\_I,353.27-354. 15 Ro\_II,477-478 {22/30}     gauṇamukhyayoḥ mukhye sampratyayaḥ bhavati . \\
\noindent \textbf{(P\_1,4.105,} 107-108.2) KA\_I,353.27-354. 15 Ro\_II,477-478 {23/30}     tat yathā . \\
\noindent \textbf{(P\_1,4.105,} 107-108.2) KA\_I,353.27-354. 15 Ro\_II,477-478 {24/30}     gauḥ anubandhyaḥ ajaḥ agnīṣhomīyaḥ iti na bāhīkaḥ anubadhyate . \\
\noindent \textbf{(P\_1,4.105,} 107-108.2) KA\_I,353.27-354. 15 Ro\_II,477-478 {25/30}     katham tarhi bāhīke vṛddhyāttve bhavataḥ . \\
\noindent \textbf{(P\_1,4.105,} 107-108.2) KA\_I,353.27-354. 15 Ro\_II,477-478 {26/30}     gauḥ tiṣṭhati . \\
\noindent \textbf{(P\_1,4.105,} 107-108.2) KA\_I,353.27-354. 15 Ro\_II,477-478 {27/30}     gām ānaya iti . \\
\noindent \textbf{(P\_1,4.105,} 107-108.2) KA\_I,353.27-354. 15 Ro\_II,477-478 {28/30}     arthāśraye etat evam bhavati . \\
\noindent \textbf{(P\_1,4.105,} 107-108.2) KA\_I,353.27-354. 15 Ro\_II,477-478 {29/30}     yat hi śabdāśrayam śabdamātre tat bhavati . \\
\noindent \textbf{(P\_1,4.105,} 107-108.2) KA\_I,353.27-354. 15 Ro\_II,477-478 {30/30}     śabdāśraye ca vṛddhyāttve . \\
\noindent \textbf{(P\_1,4.109)} paraḥ sannikarṣaḥ samhitā cet adrutāyām asaṃhitam . paraḥ sannikarṣaḥ saṃhitā cet adrutāyām vṛttau saṃhitāsañjñā na prāpnoti . \\
\noindent \textbf{(P\_1,4.109)} drutāyām eva hi paraḥ sannikarṣo varṇānām na adrutāyām . \\
\noindent \textbf{(P\_1,4.109)} tulyaḥ samnikarṣaḥ . \\
\noindent \textbf{(P\_1,4.109)} tulyaḥ samnikarṣaḥ varṇānām drutamadhyamavilimbitāsu vṛttiṣu . \\
\noindent \textbf{(P\_1,4.109)} kiṅkṛtaḥ tarthi viśeṣaḥ . \\
\noindent \textbf{(P\_1,4.109)} varṇakālabhūyastvam tu . varṇānām tu kālabhūyastvam . \\
\noindent \textbf{(P\_1,4.109)} tat yathā . \\
\noindent \textbf{(P\_1,4.109)} hastimaśakayoḥ tulyaḥ sannikarṣaḥ prāṇibhūyastvam tu . \\
\noindent \textbf{(P\_1,4.109)} yadi evam drutāyām taparakaraṇe madhyamavilimbitayoḥ upasaṅkhyānam kālabhedāt . \\
\noindent \textbf{(P\_1,4.109)} drutāyām taparakaraṇe madhyamavilimbitayoḥ upasaṅkhyānam kartavyam . \\
\noindent \textbf{(P\_1,4.109)} kim kāraṇam . \\
\noindent \textbf{(P\_1,4.109)} kālabhedāt . \\
\noindent \textbf{(P\_1,4.109)} ye drutāyām vṛttau varṇāḥ tribhāgādhikāḥ te madhyamāyām ye madhyamāyām vṛttau varṇāḥ tribhāgādhikāste vilimbitāyām . \\
\noindent \textbf{(P\_1,4.109)} uktam vā . \\
\noindent \textbf{(P\_1,4.109)} kim uktam . \\
\noindent \textbf{(P\_1,4.109)} siddham tu avasthitāḥ varṇāḥ vaktuḥ cirāciravacanāt vṛttayaḥ viśiṣyante iti . \\
\noindent \textbf{(P\_1,4.109)} atha vā śabdāvirāmaḥ saṃhitā iti etat lakṣaṇam kariṣyate . \\
\noindent \textbf{(P\_1,4.109)} śabdāvirāme prativarṇam avasānam . \\
\noindent \textbf{(P\_1,4.109)} śabdāvirāme prativarṇam avasānasañjñā prāpnoti . \\
\noindent \textbf{(P\_1,4.109)} kim idam prativarṇam iti . \\
\noindent \textbf{(P\_1,4.109)} varṇam varṇam prati prativarṇam . \\
\noindent \textbf{(P\_1,4.109)} yena eva yatnena ekaḥ varṇaḥ uccyāryate vicchinne varṇe upasaṃhṛtya tam anyam upādāya dvitīyaḥ prayujyate tathā tṛtīyaḥ tathā caturthaḥ . \\
\noindent \textbf{(P\_1,4.109)} evam tarhi anavakāśā saṃhitāsañjñā avasānasañjñām bādhiṣyate . \\
\noindent \textbf{(P\_1,4.109)} atha vā avasānasañjñāyām prakarṣagatiḥ vijñāsyate : sādhīyaḥ yaḥ virāmaḥ iti . \\
\noindent \textbf{(P\_1,4.109)} kaḥ ca sādhīyaḥ . \\
\noindent \textbf{(P\_1,4.109)} yaḥ śabdārthayoḥ virāmaḥ . \\
\noindent \textbf{(P\_1,4.109)} atha vā hrādāvirāmaḥ saṃhā iti etat lakṣaṇam kariṣyate . \\
\noindent \textbf{(P\_1,4.109)} hrādāvirāme sparśāghoṣasaṃyoge asannidhānāt asaṃhitam . \\
\noindent \textbf{(P\_1,4.109)} hrādāvirāme sparśānām aghoṣāṇām saṃyoge asamnidhānāt saṃhitāsañjñā na prāpnoti . \\
\noindent \textbf{(P\_1,4.109)} kukkuṭaḥ pippakā pittam iti . \\
\noindent \textbf{(P\_1,4.109)} kim ucyate saṃyoge iti . \\
\noindent \textbf{(P\_1,4.109)} atha yatra ekaḥ pacati iti ekaḥ pūrvaparayoḥ hrādena pracchādyate . \\
\noindent \textbf{(P\_1,4.109)} tad yathā . \\
\noindent \textbf{(P\_1,4.109)} dvayoḥ raktayoḥ vastrayoḥ madhye śuklam vastram tadguṇam upalabhyate . \\
\noindent \textbf{(P\_1,4.109)} badarapiṭake riktakaḥ lohakaṃsaḥ tadguṇaḥ upalabhyate . \\
\noindent \textbf{(P\_1,4.109)} ekena tulyaḥ sannidhiḥ . \\
\noindent \textbf{(P\_1,4.109)} yathā ekaḥ varṇaḥ hrādena pracchādyate evam anekaḥ api . \\
\noindent \textbf{(P\_1,4.109)} atha vā paurvāparyam akālavyapetam saṃhitā iti etat lakṣaṇam kariṣyate. paurvāparyam akālavyapetam saṃhitā cet pūrvāparābhāvāt asaṃhitam . paurvāparyam akālavyapetam saṃhitā cet pūrvāparābhāvāt saṃhitāsañjñā na prāpnoti . \\
\noindent \textbf{(P\_1,4.109)} na hi varṇānām paurvāparyam asti . \\
\noindent \textbf{(P\_1,4.109)} kim kāraṇam . \\
\noindent \textbf{(P\_1,4.109)} ekaikavarṇavartitvāt vācaḥ uccaritapradhvaṃsitvāt ca varṇānām . \\
\noindent \textbf{(P\_1,4.109)} ekaikavarṇavartinī vāk . \\
\noindent \textbf{(P\_1,4.109)} na dvau yugapat uccārayati . \\
\noindent \textbf{(P\_1,4.109)} gauḥ iti yāvat gakāre vāk vartate na aukāre na visarjanīye . \\
\noindent \textbf{(P\_1,4.109)} yāvat aukāre na gakāre na visarjanīye . \\
\noindent \textbf{(P\_1,4.109)} yāvat visarjanīye na gakāre na aukāre . \\
\noindent \textbf{(P\_1,4.109)} uccaritapradhvaṃsitvāt . \\
\noindent \textbf{(P\_1,4.109)} uccaritapradhvaṃsinaḥ khalu api varṇāḥ . \\
\noindent \textbf{(P\_1,4.109)} uccaritaḥ pradhvastaḥ . \\
\noindent \textbf{(P\_1,4.109)} atha aparaḥ prayujyate . \\
\noindent \textbf{(P\_1,4.109)} na varṇaḥ varṇasya sahāyaḥ . \\
\noindent \textbf{(P\_1,4.109)} evam tarhi buddhau kṛtvā sarvāḥ ceṣṭhāḥ kartā dhīraḥ tatvannītiḥ śabdena arthān vācyān dṛṣṭvā buddhau kuryāt paurvāparyam . buddhiviṣayam eva śabdānām paurvāparyam . \\
\noindent \textbf{(P\_1,4.109)} iha yaḥ eṣaḥ manuṣyaḥ prekṣāpūrvakārī bhavati saḥ paśyati asamin arthe ayam śabdaḥ prayoktvayaḥ smin tāvat śabde ayam tāvat varṇaḥ tataḥ ayam tataḥ ayam iti . \\
\noindent \textbf{(P\_1,4.110)} idam vicāryate abhāvaḥ avasānalakṣaṇam syād virāmaḥ vā iti . \\
\noindent \textbf{(P\_1,4.110)} kaḥ ca atra viśeṣaḥ . \\
\noindent \textbf{(P\_1,4.110)} abhāve avasānalakṣaṇe uparyabhāvavacanam . abhāveāvasānalakṣaṇe uparyabhāvagrahaṇam kartavyam . \\
\noindent \textbf{(P\_1,4.110)} upari yaḥ abhāvaḥ iti vaktavyam . \\
\noindent \textbf{(P\_1,4.110)} purastāt api hi śabdasya abhāvaḥ tatra mā bhūt iti . \\
\noindent \textbf{(P\_1,4.110)} kim ca syāt . \\
\noindent \textbf{(P\_1,4.110)} rasaḥ rathaḥ . \\
\noindent \textbf{(P\_1,4.110)} kharavasānayorvisarjanīyaḥ iti visarjanīyaḥ prasajyeta . \\
\noindent \textbf{(P\_1,4.110)} astu tarhi virāmaḥ . \\
\noindent \textbf{(P\_1,4.110)} virāme virāmavacanam . \\
\noindent \textbf{(P\_1,4.110)} yasya virāmaḥ virāmagrahaṇam tena kartavyam . \\
\noindent \textbf{(P\_1,4.110)} nanu ca yasya api abhāvaḥ tena api abhāvagrahaṇam kartavyam . \\
\noindent \textbf{(P\_1,4.110)} parārtham mama bhaviṣyati . \\
\noindent \textbf{(P\_1,4.110)} abhāvaḥ lopaḥ . \\
\noindent \textbf{(P\_1,4.110)} tataḥ avasānam ca iti . \\
\noindent \textbf{(P\_1,4.110)} mama api tarhi virāmagrahaṇam parārtham bhaviṣyati . \\
\noindent \textbf{(P\_1,4.110)} virāmaḥ lopaḥ avasānam ca iti . \\
\noindent \textbf{(P\_1,4.110)} upari yaḥ virāmaḥ iti vaktavyam . \\
\noindent \textbf{(P\_1,4.110)} purastāt api śabdasya virāmaḥ tatra mā bhūt iti . \\
\noindent \textbf{(P\_1,4.110)} kim ca syāt . \\
\noindent \textbf{(P\_1,4.110)} rasaḥ rathaḥ . \\
\noindent \textbf{(P\_1,4.110)} kharavasānayorvisarjanīyaḥ iti visarjanīyaḥ prasajyeta . \\
\noindent \textbf{(P\_1,4.110)} ārambhapūrvakaḥ mama virāmaḥ . \\
\noindent \textbf{(P\_1,4.110)} atha vā na idam avasānalakṣaṇam vicāryate . \\
\noindent \textbf{(P\_1,4.110)} kim tarhi . \\
\noindent \textbf{(P\_1,4.110)} sañjñī . \\
\noindent \textbf{(P\_1,4.110)} abhāvaḥ vasānasañjñī syāt virāmaḥ vā iti . \\
\noindent \textbf{(P\_1,4.110)} kaḥ ca atra viśeṣaḥ . \\
\noindent \textbf{(P\_1,4.110)} abhāve avasānasañjñini uparyabhāvavacanam . \\
\noindent \textbf{(P\_1,4.110)} abhāve avasānasañjñin yuparyabhāvagrahaṇam kartavyam . \\
\noindent \textbf{(P\_1,4.110)} upari yaḥ bhāvaḥ iti vaktavyam . \\
\noindent \textbf{(P\_1,4.110)} purastāt api hi śabdasya abhāvaḥ tatra mā bhūt iti . \\
\noindent \textbf{(P\_1,4.110)} kim ca syāt . \\
\noindent \textbf{(P\_1,4.110)} rasaḥ rathaḥ . \\
\noindent \textbf{(P\_1,4.110)} kharavasānayorvisarjanīyaḥ iti visarjanīyaḥ prasajyeta . \\
\noindent \textbf{(P\_1,4.110)} astu tarhi virāmaḥ avāsanam . \\
\noindent \textbf{(P\_1,4.110)} virāme virāmavacanam . \\
\noindent \textbf{(P\_1,4.110)} yasya virāmaḥ tena virāmagrahaṇam kartavyam . \\
\noindent \textbf{(P\_1,4.110)} nanu ca yasya api abhāvaḥ tena api abhāvagrahaṇam kartavyam . \\
\noindent \textbf{(P\_1,4.110)} parārtham mama bhaviṣyati . \\
\noindent \textbf{(P\_1,4.110)} abhāvaḥ lopaḥ . \\
\noindent \textbf{(P\_1,4.110)} tataḥ avasānam ca iti . \\
\noindent \textbf{(P\_1,4.110)} mama api tarhi virāmagrahaṇam parārtham bhaviṣyati . \\
\noindent \textbf{(P\_1,4.110)} virāmaḥ lopaḥ avasānam ca iti . \\
\noindent \textbf{(P\_1,4.110)} upari yaḥ virāmaḥ iti vaktavyam . \\
\noindent \textbf{(P\_1,4.110)} nanu ca yasya api abhāvaḥ tena api abhāvagrahaṇam kartavyam . \\
\noindent \textbf{(P\_1,4.110)} parārtham mama bhaviṣyati . \\
\noindent \textbf{(P\_1,4.110)} abhāvaḥ lopaḥ . \\
\noindent \textbf{(P\_1,4.110)} tataḥ avasānam ca iti . \\
\noindent \textbf{(P\_1,4.110)} mama api tarhi virāmagrahaṇam parārtham bhaviṣyati . \\
\noindent \textbf{(P\_1,4.110)} virāmaḥ lopaḥ avasānam ca iti . \\
\noindent \textbf{(P\_1,4.110)} upari yaḥ virāmaḥ iti vaktavyam . \\
\noindent \textbf{(P\_1,4.110)} nanu ca uktam ārambhapūrvakaḥ iti . \\
\noindent \textbf{(P\_1,4.110)} na avaśyam ayam ramiḥ pravṛttau eva vartate . \\
\noindent \textbf{(P\_1,4.110)} kim tarhi . \\
\noindent \textbf{(P\_1,4.110)} apravṛttau api . \\
\noindent \textbf{(P\_1,4.110)} tat yathā . \\
\noindent \textbf{(P\_1,4.110)} uparatāni asmin kule vratāni . \\
\noindent \textbf{(P\_1,4.110)} uparataḥ svādhyāyaḥ iti . \\
\noindent \textbf{(P\_1,4.110)} na ca tatra svādhyāyaḥ bhūtapūrvaḥ bhavati na api vratāni . \\
\noindent \textbf{(P\_1,4.110)} bhāvāvirāmabhāvitvāt śabdasya vasānalakṣaṇam na . \\
\noindent \textbf{(P\_1,4.110)} bhāvāvirāmabhāvitvāt śabdasya avasānalakṣaṇam na upapadyate . \\
\noindent \textbf{(P\_1,4.110)} kim idam bhāvāvirāmabhāvitvāt iti . \\
\noindent \textbf{(P\_1,4.110)} bhāvasya avirāmaḥ bhāvāvirāmaḥ bhāvāvirāmeṇa bhavati iti bhāvāvirāmabhāvī bhāvāvirāmabhāvinaḥ bhāvo bhāvāvirāmabhāvitvam . \\
\noindent \textbf{(P\_1,4.110)} aparaḥ āha . \\
\noindent \textbf{(P\_1,4.110)} bhāvabhāvitvādavirāmabhāvitvāt ca śabdasya avasānalakṣaṇam na upapadyate iti . \\
\noindent \textbf{(P\_1,4.110)} tatparaḥ iti vā varṇasya avasānam . virāmaparaḥ varṇaḥ vasānasañjñaḥ bhavati iti vaktavyam . \\
\noindent \textbf{(P\_1,4.110)} varṇaḥ antyaḥ vā avasānam . \\
\noindent \textbf{(P\_1,4.110)} atha vā vyaktam eva paṭhitavyam antyaḥ varṇaḥ vasānasañjñaḥ bhavati iti . \\
\noindent \textbf{(P\_1,4.110)} tat tarhi vaktavyam . \\
\noindent \textbf{(P\_1,4.110)} na vaktavyam . \\
\noindent \textbf{(P\_1,4.110)} saṃhitāvasānayoḥ lokaviditatvāt siddham . \\
\noindent \textbf{(P\_1,4.110)} samhitā avasānam iti lokaviditau etau arthau . \\
\noindent \textbf{(P\_1,4.110)} evam hi kaḥ cit kam cid adhīyānam āha : śannodevīyam samhitayā adhīṣva iti . \\
\noindent \textbf{(P\_1,4.110)} saḥ tatra paramasannikarṣam adhīte . \\
\noindent \textbf{(P\_1,4.110)} aparaḥ āha : kena vasyasi iti . \\
\noindent \textbf{(P\_1,4.110)} saḥ āha : akāreṇa ikāreṇa ukāreṇa iti . \\
\noindent \textbf{(P\_1,4.110)} evam etau lokaviditatau arthau . \\
\noindent \textbf{(P\_1,4.110)} tayoḥ lokaviditatvāt siddham iti . \\
\noindent \textbf{(P\_2,1.1.1)} vidhiḥ iti kaḥ ayam śabdaḥ . \\
\noindent \textbf{(P\_2,1.1.1)} vipūrvāt dhāñaḥ karmasādhanaḥ ikāraḥ . \\
\noindent \textbf{(P\_2,1.1.1)} vidhīyate vidhiḥ iti . \\
\noindent \textbf{(P\_2,1.1.1)} kim punaḥ vidhīyate . \\
\noindent \textbf{(P\_2,1.1.1)} samāsaḥ vibhaktividhānam parāṅgavadbhāvaḥ ca . \\
\noindent \textbf{(P\_2,1.1.1)} kim punaḥ ayam adhikāraḥ āhosvit paribhāṣā . \\
\noindent \textbf{(P\_2,1.1.1)} kaḥ punaḥ adhikāraparibhāṣayoḥ viśeṣaḥ . \\
\noindent \textbf{(P\_2,1.1.1)} adhikāraḥ pratiyogam tasya anirdeśārthaḥ iti yoge yoge upatiṣṭhate . \\
\noindent \textbf{(P\_2,1.1.1)} paribhāṣā punaḥ ekadeśasthā satī sarvam śāstram abhijvalayati pradīpavat . \\
\noindent \textbf{(P\_2,1.1.1)} tat yathā pradīpaḥ suprajvalitaḥ ekadeśasthaḥ sarvam veśma abhijvalayati . \\
\noindent \textbf{(P\_2,1.1.1)} kaḥ punaḥ atra prayatnaviśeṣaḥ . \\
\noindent \textbf{(P\_2,1.1.1)} adhikāre sati svarayitavyam paribhāṣāyām punaḥ satyām sarvam apekṣyam . \\
\noindent \textbf{(P\_2,1.1.1)} tathā idam aparam dvaitam bhavati . \\
\noindent \textbf{(P\_2,1.1.1)} ekārthībhāvaḥ vā sāmarthyam syāt vyapekṣā vā iti . \\
\noindent \textbf{(P\_2,1.1.1)} tatra ekārthībhāve sāmarthye adhikāre ca sati samāsaḥ ekaḥ saṅgṛhītaḥ bhavati bibhaktividhānam parāṅgavadbhāvaḥ ca asaṅgṛhītaḥ . \\
\noindent \textbf{(P\_2,1.1.1)} vyapekṣāyām punaḥ sāmarthye adhikāre ca sati bibhaktividhānam parāṅgavadbhāvaḥ ca saṅgṛhītaḥ samāsaḥ tu ekaḥ asaṅgṛhītaḥ . \\
\noindent \textbf{(P\_2,1.1.1)} anyatra khalu api samarthagrahaṇāni yuktagrahaṇāni ca kartavyāni bhavanti . \\
\noindent \textbf{(P\_2,1.1.1)} kva anyatra . \\
\noindent \textbf{(P\_2,1.1.1)} isusoḥ sāmarthye na cavāhāhaivayukte iti . \\
\noindent \textbf{(P\_2,1.1.1)} vyapekṣāyām punaḥ sāmarthye paribhāṣāyām ca satyām yāvān vyākaraṇe padagandhaḥ asti saḥ sarvaḥ saṅgṛhītaḥ bhavati samāsaḥ tu ekaḥ asaṅgṛhītaḥ . \\
\noindent \textbf{(P\_2,1.1.1)} tatra ekārthībhāvaḥ sāmarthyam paribhāṣā ca iti evam sūtram abhinnatarakam bhavati . \\
\noindent \textbf{(P\_2,1.1.1)} evam api kva cit akartavyam samarthagrahaṇam kriyate kva cit ca kartavyam na kriyate . \\
\noindent \textbf{(P\_2,1.1.1)} akartavyam tāvat kriyate samarthānām prathamāt vā iti . \\
\noindent \textbf{(P\_2,1.1.1)} kartavyam ca na kriyate karmaṇi aṇ samarthāt iti . \\
\noindent \textbf{(P\_2,1.1.1)} nanu ca gamyate tatra sāmarthyam . \\
\noindent \textbf{(P\_2,1.1.1)} kumbhakāraḥ nagarakāraḥ iti . \\
\noindent \textbf{(P\_2,1.1.1)} satyam gamyate utpanne tu pratyaye . \\
\noindent \textbf{(P\_2,1.1.1)} saḥ eva tāvat samarthāt utpādyaḥ . \\
\noindent \textbf{(P\_2,1.1.2)} atha samarthagrahaṇam kimartham . \\
\noindent \textbf{(P\_2,1.1.2)} vakṣyati dvitīyā śritādibhiḥ samasyate . \\
\noindent \textbf{(P\_2,1.1.2)} kaṣṭaśritaḥ narakaśritaḥ iti . \\
\noindent \textbf{(P\_2,1.1.2)} samarthagrahaṇam kimartham . \\
\noindent \textbf{(P\_2,1.1.2)} paśya devadatta kaṣṭam . \\
\noindent \textbf{(P\_2,1.1.2)} śritaḥ viṣṇumitraḥ gurukulam . \\
\noindent \textbf{(P\_2,1.1.2)} tṛtīyā tatkṛtārthena guṇavacanena . \\
\noindent \textbf{(P\_2,1.1.2)} śaṅkulākhaṇḍaḥ kirikāṇaḥ . \\
\noindent \textbf{(P\_2,1.1.2)} samarthagrahaṇam kimartham . \\
\noindent \textbf{(P\_2,1.1.2)} tiṣṭha tvam śaṅkulayā . \\
\noindent \textbf{(P\_2,1.1.2)} khaṇḍaḥ dhāvati musalena . \\
\noindent \textbf{(P\_2,1.1.2)} caturthī tadarthārthabalihitasukharakṣitaiḥ . \\
\noindent \textbf{(P\_2,1.1.2)} gohitam aśrahitam . \\
\noindent \textbf{(P\_2,1.1.2)} samarthagrahaṇam kimartham . \\
\noindent \textbf{(P\_2,1.1.2)} sukham gobhyaḥ . \\
\noindent \textbf{(P\_2,1.1.2)} hitam devadattāya . \\
\noindent \textbf{(P\_2,1.1.2)} pañcamī bhayena . \\
\noindent \textbf{(P\_2,1.1.2)} vṛkabhayam dasyubhayam corabhayam . \\
\noindent \textbf{(P\_2,1.1.2)} samarthagrahaṇam kimartham . \\
\noindent \textbf{(P\_2,1.1.2)} gaccha tvam mā vṛkebhyaḥ . \\
\noindent \textbf{(P\_2,1.1.2)} bhayam devadattasya yajñadattāt . \\
\noindent \textbf{(P\_2,1.1.2)} ṣaṣṭhī subantena samasyate : rājapuruṣaḥ , brāhmaṇakambalaḥ . \\
\noindent \textbf{(P\_2,1.1.2)} samarthagrahaṇam kimartham . \\
\noindent \textbf{(P\_2,1.1.2)} bhāryā rājñaḥ . \\
\noindent \textbf{(P\_2,1.1.2)} puruṣaḥ devadattasya . \\
\noindent \textbf{(P\_2,1.1.2)} saptamī śauṇḍaiḥ : akṣaśauṇḍaḥ , strīśauṇḍaḥ . \\
\noindent \textbf{(P\_2,1.1.2)} samarthagrahaṇam kimartham . \\
\noindent \textbf{(P\_2,1.1.2)} kuśalaḥ devadattaḥ akṣeṣu . \\
\noindent \textbf{(P\_2,1.1.2)} śauṇḍaḥ pibati pānāgāre . \\
\noindent \textbf{(P\_2,1.1.2)} atha kriyamāṇe api samarthagrahaṇe iha kasmāt na bhavati mahat kaṣṭam śritaḥ iti . \\
\noindent \textbf{(P\_2,1.1.2)} na vā bhavati mahākaṣṭaśritaḥ iti . \\
\noindent \textbf{(P\_2,1.1.2)} bhavati yadā etat vākyam bhavati : mahat kaṣṭam mahākaṣṭam , mahākaṣṭam śritaḥ mahākaṣṭaśritaḥ iti . \\
\noindent \textbf{(P\_2,1.1.2)} yadā tu etat vākyam bhavati : mahat kaṣṭam śritaḥ iti tadā na bhavitavyam tadā ca prapnoti . \\
\noindent \textbf{(P\_2,1.1.2)} tadā kasmāt na bhavati . \\
\noindent \textbf{(P\_2,1.1.2)} kasya kasmāt na bhavati . \\
\noindent \textbf{(P\_2,1.1.2)} kim dvayoḥ āhosvit bahūnām . \\
\noindent \textbf{(P\_2,1.1.2)} bahūnām kasmāt na bhavati . \\
\noindent \textbf{(P\_2,1.1.2)} sup supā iti vartate . \\
\noindent \textbf{(P\_2,1.1.2)} nanu ca bhoḥ ākṛtau śāstrāṇi pravartante . \\
\noindent \textbf{(P\_2,1.1.2)} tat yathā prātipadikāt iti vartamāne anyasmāt ca anyasmāt ca prātipadikāt utpattiḥ bhavati . \\
\noindent \textbf{(P\_2,1.1.2)} satyam etat . \\
\noindent \textbf{(P\_2,1.1.2)} ākṛtiḥ tu pratyekam parisamāpyate . \\
\noindent \textbf{(P\_2,1.1.2)} yāvati etat parisamāpyate prātipadikāt iti tāvataḥ utpattyā bhavitavyam . \\
\noindent \textbf{(P\_2,1.1.2)} pratyekam ca etat parisamāpyate na samudāye . \\
\noindent \textbf{(P\_2,1.1.2)} evam iha api yāvati etat parisamāpyate sup supā iti tāvataḥ samāsena bhavitavyam . \\
\noindent \textbf{(P\_2,1.1.2)} dvayoḥ dvayoḥ ca etat parisamāpyate na bahuṣu . \\
\noindent \textbf{(P\_2,1.1.2)} dvayoḥ tarhi kasmāt na bhavati . \\
\noindent \textbf{(P\_2,1.1.2)} asāmarthyāt . \\
\noindent \textbf{(P\_2,1.1.2)} katham asāmarthyam . \\
\noindent \textbf{(P\_2,1.1.2)} sāpekṣam asamartham bhavati iti . \\
\noindent \textbf{(P\_2,1.1.2)} yadi sāpekṣam asamartham bhavati iti ucyate rājapuruṣaḥ abhirūpaḥ rājapuruṣaḥ darśanīyaḥ atra vṛttiḥ na prāpnoti . \\
\noindent \textbf{(P\_2,1.1.2)} na eṣaḥ doṣaḥ . \\
\noindent \textbf{(P\_2,1.1.2)} pradhānam atra sāpekṣam . \\
\noindent \textbf{(P\_2,1.1.2)} bhavati ca pradhānasya sāpekṣasya api samāsaḥ . \\
\noindent \textbf{(P\_2,1.1.2)} yatra tarhi apradhānam sāpekṣam bhavati tatra te vṛttiḥ na prāpnoti : devadattasya gurukulam , devadattasya guruputraḥ , devadattasya dāsabhāryā iti . \\
\noindent \textbf{(P\_2,1.1.2)} na eṣaḥ doṣaḥ . \\
\noindent \textbf{(P\_2,1.1.2)} samudāyapekṣā atra ṣaṣṭḥī sarvam gurukulam apekṣate . \\
\noindent \textbf{(P\_2,1.1.2)} yatra tarhi na samudāyapekṣā ṣaṣṭḥī tatra vṛttiḥ na prāpnoti : kim odanaḥ śālīnām . \\
\noindent \textbf{(P\_2,1.1.2)} saktvāḍhakam āpaṇīyānām . \\
\noindent \textbf{(P\_2,1.1.2)} kutaḥ bhavān pāṭaliputrakaḥ . \\
\noindent \textbf{(P\_2,1.1.2)} iha ca api : devadattasya gurukulam , devadattasya guruputraḥ , devadattasya dāsabhāryā iti : yadi eṣā samudāyapekṣā ṣaṣṭḥī syāt na etat niyogataḥ gamyeta devadattasya yaḥ guruḥ tasya yaḥ putraḥ iti . \\
\noindent \textbf{(P\_2,1.1.2)} kim tarhi . \\
\noindent \textbf{(P\_2,1.1.2)} anyasya api guruputraḥ devadattasya kim cit iti eṣaḥ arthaḥ gamyeta . \\
\noindent \textbf{(P\_2,1.1.2)} yataḥ tu niyogataḥ devadattasya yaḥ guruḥ tasya yaḥ putraḥ iti eṣaḥ arthaḥ gamyate ataḥ manyāmahe na samudāyapekṣā ṣaṣṭḥī iti . \\
\noindent \textbf{(P\_2,1.1.2)} anyatra khalu api samarthagrahaṇe sāpekṣasya api kāryam bhavati . \\
\noindent \textbf{(P\_2,1.1.2)} kva anyatra . \\
\noindent \textbf{(P\_2,1.1.2)} isusoḥ sāmarthye . \\
\noindent \textbf{(P\_2,1.1.2)} brāhamaṇasya sarpiḥ karoti iti . \\
\noindent \textbf{(P\_2,1.1.2)} tasmāt na aeta śakyak vaktum sāpekṣam asamartham bhavati iti . \\
\noindent \textbf{(P\_2,1.1.2)} vṛttiḥ tarhi kasmāt na bhavati mahat kaṣṭam śritaḥ iti . \\
\noindent \textbf{(P\_2,1.1.2)} saviśeṣaṇānām vṛttiḥ na vṛttasya vā viśeṣaṇam na prayujyate iti vaktavyam . \\
\noindent \textbf{(P\_2,1.1.2)} yadi saviśeṣaṇānām vṛttiḥ na vṛttasya vā viśeṣaṇam na prayujyate iti ucyate devadattasya gurukulam devadattasya guruputraḥ devadattasya dāsabhāryā iti atra vṛttiḥ na prāpnoti . \\
\noindent \textbf{(P\_2,1.1.2)} agurukulaputrādīnām iti vaktavyam . \\
\noindent \textbf{(P\_2,1.1.2)} tat tarhi vaktavyam saviśeṣaṇānām vṛttiḥ na vṛttasya vā viśeṣaṇam na prayujyate agurukulaputrādīnām iti . \\
\noindent \textbf{(P\_2,1.1.2)} na vaktavyam . \\
\noindent \textbf{(P\_2,1.1.2)} vṛttiḥ tarhi kasmāt na bhavati . \\
\noindent \textbf{(P\_2,1.1.2)} agamakatvāt . \\
\noindent \textbf{(P\_2,1.1.2)} iha samānārthena vākyena bhavitavyam samāsena ca . \\
\noindent \textbf{(P\_2,1.1.2)} yaḥ ca iha arthaḥ vākyena gamyate mahat kaṣṭam śritaḥ iti na jātu cit samāsena asau gamyate mahat kaṣṭaśritaḥ iti . \\
\noindent \textbf{(P\_2,1.1.2)} etasmāt hetoḥ brūmaḥ agamakatvāt iti . \\
\noindent \textbf{(P\_2,1.1.2)} na brūmaḥ apaśabdaḥ syāt iti . \\
\noindent \textbf{(P\_2,1.1.2)} yatra gamakaḥ bhavati bhavati tatra vṛttiḥ . \\
\noindent \textbf{(P\_2,1.1.2)} tat yathā : devadattasya gurukulam , devadattasya guruputraḥ , devadattasya dāsabhāryā iti . \\
\noindent \textbf{(P\_2,1.1.2)} yadi agamakatvam hetuḥ na arthaḥ samarthagrahaṇena . \\
\noindent \textbf{(P\_2,1.1.2)} iha api bhāryā rājñaḥ puruṣaḥ devadattasya iti yaḥ arthaḥ vākyena gamyate na asau jātu cit samāsena asau gamyate bhārya rājapuruṣaḥ devadattasya iti . \\
\noindent \textbf{(P\_2,1.1.2)} tasmāt na arthaḥ samarthagrahaṇena . \\
\noindent \textbf{(P\_2,1.1.2)} idam tarhi prayojanam . \\
\noindent \textbf{(P\_2,1.1.2)} asti asamarthasamāsaḥ nañsamāsaḥ gamakaḥ . \\
\noindent \textbf{(P\_2,1.1.2)} tasya sādhutvam mā bhūt . \\
\noindent \textbf{(P\_2,1.1.2)} akiñcit kurvāṇam amāṣam haramāṇam agādhāt utsṛṣṭam iti . \\
\noindent \textbf{(P\_2,1.1.2)} etat api na asti prayojanam . \\
\noindent \textbf{(P\_2,1.1.2)} avaśyam kasya cit nañsamāsasya gamakasya sādhutvam vaktavyam . \\
\noindent \textbf{(P\_2,1.1.2)} asūryampaśyāni mukhāni apunargeyāḥ ślokāḥ aśrāddhabhojī alavaṇabhojī brāhmaṇaḥ . \\
\noindent \textbf{(P\_2,1.1.2)} suṭ anapuṃsakasya etat niyamārtham bhaviṣyati . \\
\noindent \textbf{(P\_2,1.1.2)} etasya eva asamarthasamāsasya nañsamāsasya gamakasya sādhutvam bhavati na anyasya iti . \\
\noindent \textbf{(P\_2,1.1.2)} tasmān na arthaḥ samarthagrahaṇena . \\
\noindent \textbf{(P\_2,1.1.3)} atha kriyamāṇe api samarthagrahaṇe samartham iti ucyate kim samartham nāma . \\
\noindent \textbf{(P\_2,1.1.3)} pṛthagarthānām ekārthībhāvaḥ samarthavacanam . \\
\noindent \textbf{(P\_2,1.1.3)} pṛthagarthānām padānam ekārthībhāvaḥ samartham iti ucyate . \\
\noindent \textbf{(P\_2,1.1.3)} vākye pṛthagarthāni rājñaḥ puruṣaḥ iti . \\
\noindent \textbf{(P\_2,1.1.3)} samāse punaḥ ekārthāni rājapuruṣaḥ iti . \\
\noindent \textbf{(P\_2,1.1.3)} kim ucyate pṛthagarthāni iti yāvatā rājñaḥ puruṣaḥ ānīyatām iti ukte rājapuruṣaḥ iti ca saḥ eva . \\
\noindent \textbf{(P\_2,1.1.3)} na api brūmaḥ anyasya ānayanam bhavati iti . \\
\noindent \textbf{(P\_2,1.1.3)} kaḥ tarhi ekārthībhāvakṛtaḥ viśeṣaḥ . \\
\noindent \textbf{(P\_2,1.1.3)} subalopaḥ vyavadhānam yatheṣṭam anyatareṇa abhisambandhaḥ svaraḥ . \\
\noindent \textbf{(P\_2,1.1.3)} supaḥ alopaḥ bhavati vākye . \\
\noindent \textbf{(P\_2,1.1.3)} rājñaḥ puruṣaḥ iti . \\
\noindent \textbf{(P\_2,1.1.3)} samāse punaḥ na bhavati . \\
\noindent \textbf{(P\_2,1.1.3)} rājapuruṣaḥ iti . \\
\noindent \textbf{(P\_2,1.1.3)} vyavadhānam ca bhavati vākye . \\
\noindent \textbf{(P\_2,1.1.3)} rājñaḥ ṛddhasya puruṣaḥ iti . \\
\noindent \textbf{(P\_2,1.1.3)} samāse na bhavati . \\
\noindent \textbf{(P\_2,1.1.3)} rājapuruṣaḥ iti . \\
\noindent \textbf{(P\_2,1.1.3)} yatheṣṭam anyatareṇa abhisambandhaḥ bhavati vākye . \\
\noindent \textbf{(P\_2,1.1.3)} rajñaḥ puruṣaḥ puruṣaḥ rājñaḥ iti . \\
\noindent \textbf{(P\_2,1.1.3)} samāse na bhavati . \\
\noindent \textbf{(P\_2,1.1.3)} rājapuruṣaḥ iti . \\
\noindent \textbf{(P\_2,1.1.3)} dvau svarau bhavataḥ vākye . \\
\noindent \textbf{(P\_2,1.1.3)} rajñaḥ puruṣaḥ . \\
\noindent \textbf{(P\_2,1.1.3)} samāse punaḥ ekaḥ eva . \\
\noindent \textbf{(P\_2,1.1.3)} rājapuruṣaḥ iti . \\
\noindent \textbf{(P\_2,1.1.3)} na ete ekārthībhāvakṛtāḥ viśeṣāḥ . \\
\noindent \textbf{(P\_2,1.1.3)} kim tarhi . \\
\noindent \textbf{(P\_2,1.1.3)} vācanikāni etāni . \\
\noindent \textbf{(P\_2,1.1.3)} āha hi bhagavān supaḥ dhātuprātipadikayoḥ upasarjanam pūrvam samāsasya antaḥ udāttaḥ bhavati iti . \\
\noindent \textbf{(P\_2,1.1.3)} ime tarhi ekārthībhāvkṛtāḥ viśeṣāḥ . \\
\noindent \textbf{(P\_2,1.1.3)} saṅkhyāviśeṣaḥ vyaktābhidānam lkupasarjanaviśeṣaṇam cayogaḥ iti . \\
\noindent \textbf{(P\_2,1.1.3)} saṅkhyāviśeṣaḥ bhavati vākye . \\
\noindent \textbf{(P\_2,1.1.3)} rājñaḥ puruṣaḥ rājñoḥ puruṣaḥ rājñām puruṣaḥ iti . \\
\noindent \textbf{(P\_2,1.1.3)} samāse na bhavati . \\
\noindent \textbf{(P\_2,1.1.3)} rājapuruṣaḥ iti . \\
\noindent \textbf{(P\_2,1.1.3)} asti kāraṇam yena etat evam bhavati . \\
\noindent \textbf{(P\_2,1.1.3)} kim kāraṇam . \\
\noindent \textbf{(P\_2,1.1.3)} yaḥ asau viśeṣavācī śabdaḥ tadasānnidhyāt . \\
\noindent \textbf{(P\_2,1.1.3)} aṅga hi bhavān tam uccārayatu gaṃsyate saḥ viśeṣaḥ . \\
\noindent \textbf{(P\_2,1.1.3)} nanu ca na etena evam bhavitavyam . \\
\noindent \textbf{(P\_2,1.1.3)} na hi śabdakṛtena nāma arthena bhavitavyam . \\
\noindent \textbf{(P\_2,1.1.3)} arthakṛtena nāma śabdena bhavitavyam . \\
\noindent \textbf{(P\_2,1.1.3)} tat etat evam dṛśyatām . \\
\noindent \textbf{(P\_2,1.1.3)} artharūpam eva etat evañjātīyakam yena atra viśeṣaḥ na gamyate iti . \\
\noindent \textbf{(P\_2,1.1.3)} avaśyam ca etat evam vijñeyam . \\
\noindent \textbf{(P\_2,1.1.3)} yaḥ hi manyate yaḥ asau viśeṣavācī śabdaḥ tadasānnidhyāt atra viśeṣaḥ na gamyate iti iha tasya viśeṣaḥ gamyeta : apsucaraḥ goṣucaraḥ varṣāsujaḥ iti . \\
\noindent \textbf{(P\_2,1.1.3)} vyaktābhidhānam bhavati vākye . \\
\noindent \textbf{(P\_2,1.1.3)} brāhmaṇasya kambalaḥ tiṣṭhati iti . \\
\noindent \textbf{(P\_2,1.1.3)} samāse punaḥ avyaktam . \\
\noindent \textbf{(P\_2,1.1.3)} brāhmaṇakambalaḥ tiṣṭhati iti . \\
\noindent \textbf{(P\_2,1.1.3)} sandehaḥ bhavati sambuddhiḥ syāt ṣaṣṭhīsamāsaḥ vā iti . \\
\noindent \textbf{(P\_2,1.1.3)} eṣaḥ api aviśeṣaḥ . \\
\noindent \textbf{(P\_2,1.1.3)} bhavati hi kim cit vākye avyaktam tat ca samāse vyaktam . \\
\noindent \textbf{(P\_2,1.1.3)} vākye tāvat avyaktam . \\
\noindent \textbf{(P\_2,1.1.3)} ardham paśoḥ devadattasya iti . \\
\noindent \textbf{(P\_2,1.1.3)} sandehaḥ bhavati paśuguṇasya vā devadattasya yat ardham artha vā yaḥ asau sañjñībhūtaḥ paśuḥ nāma tasya yat ardham iti . \\
\noindent \textbf{(P\_2,1.1.3)} tat ca samāse vyaktam bhavati . \\
\noindent \textbf{(P\_2,1.1.3)} ardhhapaśuḥ devadattasya iti . \\
\noindent \textbf{(P\_2,1.1.3)} upasarjanaviśeṣaṇam bhavati vākye . \\
\noindent \textbf{(P\_2,1.1.3)} ṛddhasya rājñaḥ puruṣaḥ iti . \\
\noindent \textbf{(P\_2,1.1.3)} samāse na bhavati . \\
\noindent \textbf{(P\_2,1.1.3)} rājapuruṣaḥ iti . \\
\noindent \textbf{(P\_2,1.1.3)} eṣaḥ api adoṣaḥ . \\
\noindent \textbf{(P\_2,1.1.3)} samāse api upasarjanaviśeṣaṇam bhavati . \\
\noindent \textbf{(P\_2,1.1.3)} tat yathā devadattasya gurukulam devadattasya guruputraḥ devadattasya dāsabhāryā iti . \\
\noindent \textbf{(P\_2,1.1.3)} cayogaḥ bhavati vākye . \\
\noindent \textbf{(P\_2,1.1.3)} svacayogaḥ svāmicayogaḥ ca . \\
\noindent \textbf{(P\_2,1.1.3)} svacayogaḥ rājñaḥ gauḥ ca aśvaḥ ca puruṣaḥ ca iti . \\
\noindent \textbf{(P\_2,1.1.3)} samāse na bhavati . \\
\noindent \textbf{(P\_2,1.1.3)} rājñaḥ gavāśvapuruṣāḥ iti . \\
\noindent \textbf{(P\_2,1.1.3)} svāmicayogaḥ devadattasya ca yajñadattasya ca viṣṇumitrasya ca gauḥ iti . \\
\noindent \textbf{(P\_2,1.1.3)} samāse na bhavati . \\
\noindent \textbf{(P\_2,1.1.3)} devadattayajñadattaviṣṇumitrāṇām gauḥ iti . \\
\noindent \textbf{(P\_2,1.1.3)} atha etasmin ekārthībhāvkṛte viśeṣe kim svābhāvikam śabdaiḥ arthābhidhānam āhosvit vācanikam . \\
\noindent \textbf{(P\_2,1.1.3)} svābhāvikam iti āha . \\
\noindent \textbf{(P\_2,1.1.3)} kutaḥ etat . \\
\noindent \textbf{(P\_2,1.1.3)} arthānādeśāt . \\
\noindent \textbf{(P\_2,1.1.3)} na hi arthāḥ ādiśyante . \\
\noindent \textbf{(P\_2,1.1.3)} katham punaḥ arthān ādiśan evam brūyāt na arthāḥ ādiśyante it . \\
\noindent \textbf{(P\_2,1.1.3)} yat āha bhavān anekam anyapadārthe cārthe dvandvaḥ apatye rakte nirvṛtte iti . \\
\noindent \textbf{(P\_2,1.1.3)} na etāni arthādeśanāni . \\
\noindent \textbf{(P\_2,1.1.3)} svabhāvataḥ eteṣām śabdānām eteṣu artheṣu abhiniviṣṭānām nimittatvena anvākhyānam kriyate . \\
\noindent \textbf{(P\_2,1.1.3)} tat yathā . \\
\noindent \textbf{(P\_2,1.1.3)} kūpe hastadakṣiṇaḥ panthāḥ . \\
\noindent \textbf{(P\_2,1.1.3)} abhre candramasam paśya iti . \\
\noindent \textbf{(P\_2,1.1.3)} svabhāvataḥ tatrasthasya pathaḥ candramasaḥ ca nimittatvena anvākhyānam kriyate . \\
\noindent \textbf{(P\_2,1.1.3)} evam iha api cārthe yaḥ saḥ dvandvasamāsaḥ anyapadārthaḥ yaḥ saḥ bahuvrīhiḥ iti . \\
\noindent \textbf{(P\_2,1.1.3)} kim puna kāraṇam na ādiśyante . \\
\noindent \textbf{(P\_2,1.1.3)} tat ca laghvartham . \\
\noindent \textbf{(P\_2,1.1.3)} laghvartham hi arthāḥ na ādiśyante . \\
\noindent \textbf{(P\_2,1.1.3)} avaśyam hi anena arthān ādiśatā kena cit śabdena nirdeśaḥ kartavyaḥ syāt . \\
\noindent \textbf{(P\_2,1.1.3)} tasya ca tāvat kena kṛtaḥ yena asau kriyate . \\
\noindent \textbf{(P\_2,1.1.3)} atha tasya kena cit kṛtaḥ tasya kena kṛtaḥ iti anavasthā . \\
\noindent \textbf{(P\_2,1.1.3)} asambhavaḥ khalu api ādeśaḥ tasya . \\
\noindent \textbf{(P\_2,1.1.3)} kaḥ hi nāma samarthaḥ dhātuprātipadikapratyayanipātānām arthān ādeṣṭum . \\
\noindent \textbf{(P\_2,1.1.3)} na ca etat mantavyam pratyayārthe nirdiṣṭe prakṛtyarthaḥ anirdiṣṭaḥ iti . \\
\noindent \textbf{(P\_2,1.1.3)} bhavati hi guṇābhidhāne guṇinaḥ sampratyayaḥ . \\
\noindent \textbf{(P\_2,1.1.3)} tat yathā śuklaḥ kṛṣṇaḥ iti . \\
\noindent \textbf{(P\_2,1.1.3)} viṣamaḥ upanyāsaḥ . \\
\noindent \textbf{(P\_2,1.1.3)} sāmānyaśabdāḥ ete evam syuḥ . \\
\noindent \textbf{(P\_2,1.1.3)} sāmanyaśabdāḥ ca na antareṇa viśeṣam prakaraṇam vā viśeṣeṣu avatiṣṭhante . \\
\noindent \textbf{(P\_2,1.1.3)} yataḥ tu khalu niyogataḥ vṛkṣaḥ iti ukte svabhāvataḥ kasmin cit eva viśeṣe vṛkṣaśabdaḥ vartate ataḥ manyāmahe na ime sāmānyaśabdāḥ iti . \\
\noindent \textbf{(P\_2,1.1.3)} na cet sāmānyaśabdāḥ prakṛtiḥ prakṛtyarthe vartate pratyayaḥ pratyayārthe vartate . \\
\noindent \textbf{(P\_2,1.1.3)} apravṛttiḥ khalu api arthādeśanasya . \\
\noindent \textbf{(P\_2,1.1.3)} bahavaḥ hi śabdāḥ yeṣām arthāḥ na vijñāyante . \\
\noindent \textbf{(P\_2,1.1.3)} jarbharī turpharītū . \\
\noindent \textbf{(P\_2,1.1.3)} antareṇa khalu api śabdaprayogam bahavaḥ arthāḥ gamyante akṣinikocaiḥ pāṇivihāraiḥ ca . \\
\noindent \textbf{(P\_2,1.1.3)} na khalu api nirjñātasya arthasya anvyākhyane kim cit prayojanam asti . \\
\noindent \textbf{(P\_2,1.1.3)} yaḥ hi brūyāt purastāt ādityaḥ udeti paścāt astam eti madhuraḥ guḍaḥ kaṭukam śṛṅgaveram iti kim tena kṛtam syāt . \\
\noindent \textbf{(P\_2,1.1.4)} vāvacanānarthakyam ca svabhāvasiddhatvāt . \\
\noindent \textbf{(P\_2,1.1.4)} vāvacanānarthakyam . \\
\noindent \textbf{(P\_2,1.1.4)} kim kāraṇam . \\
\noindent \textbf{(P\_2,1.1.4)} svabhāvasiddhatvāt . \\
\noindent \textbf{(P\_2,1.1.4)} iha dvau pakṣau vṛttipakṣaḥ avṛttipakṣaḥ ca . \\
\noindent \textbf{(P\_2,1.1.4)} svabhāvataḥ ca etat bhavati vākyam ca samāsaḥ ca . \\
\noindent \textbf{(P\_2,1.1.4)} tatra svābhāvike vṛttiviṣaye nitye samāse prapte vāvacanena kim anyat śakyam abhisambandhum anyat ataḥ sañjñāyāḥ . \\
\noindent \textbf{(P\_2,1.1.4)} na ca sañjñāyāḥ bhāvābhāvau iṣyete . \\
\noindent \textbf{(P\_2,1.1.4)} tasmāt na arthaḥ vā vacanena \\
\noindent \textbf{(P\_2,1.1.5)} atha ye vṛttim vartayanti kim te āhuḥ . \\
\noindent \textbf{(P\_2,1.1.5)} parārthābhidhānam vṛttiḥ iti āhuḥ . \\
\noindent \textbf{(P\_2,1.1.5)} atha teṣām evam bruvatām kim jahatsvārthā vṛttiḥ āhosvit ajahatsvārthā . \\
\noindent \textbf{(P\_2,1.1.5)} kim ca ataḥ . \\
\noindent \textbf{(P\_2,1.1.5)} yadi jahatsvārthā vṛttiḥ rājapuruṣam ānaya iti ukte puruṣamātrasya ānayanam prāpnoti aupagavam ānaya iti ukte apatyamātrasya . \\
\noindent \textbf{(P\_2,1.1.5)} atha ajahatsvārthā vṛttiḥ ubhayoḥ vidyamānasvārthayoḥ dvayoḥ dvivacanam iti dvivacanam prāpnoti . \\
\noindent \textbf{(P\_2,1.1.5)} kā punaḥ vṛttiḥ nyāyyā . \\
\noindent \textbf{(P\_2,1.1.5)} jahatsvārthā . \\
\noindent \textbf{(P\_2,1.1.5)} yuktam punaḥ yat jahatsvārthā nāma vṛttiḥ syāt . \\
\noindent \textbf{(P\_2,1.1.5)} bāḍham yuktam . \\
\noindent \textbf{(P\_2,1.1.5)} evam hi dṛśyate loke . \\
\noindent \textbf{(P\_2,1.1.5)} puruṣaḥ ayam parakarmaṇi pravartamānaḥ svam karma jahāti . \\
\noindent \textbf{(P\_2,1.1.5)} tat yathā takṣā rājakarmaṇi vartamānaḥ svam karma jahāti . \\
\noindent \textbf{(P\_2,1.1.5)} evam yuktam yat rājā puruṣārthe vartamānaḥ svam artham jahyāt upaguḥ ca apartyārthe vartamānaḥ svam artham jahyāt . \\
\noindent \textbf{(P\_2,1.1.5)} nanu ca uktam rājapuruṣam ānaya iti ukte puruṣamātrasya ānayanam prāpnoti aupagavam ānaya iti ukte apatyamātrasya iti . \\
\noindent \textbf{(P\_2,1.1.5)} na eṣaḥ doṣaḥ . \\
\noindent \textbf{(P\_2,1.1.5)} jahat api asau svārtham na atyantāya jahāti . \\
\noindent \textbf{(P\_2,1.1.5)} yaḥ parārthavirodhī svārthaḥ tam jahāti . \\
\noindent \textbf{(P\_2,1.1.5)} tat yathā takṣā rājakarmaṇi vartamānaḥ svam takṣakarma jahāti na hikkitahasitakaṇḍūyitāni . \\
\noindent \textbf{(P\_2,1.1.5)} na ca ayam arthaḥ parārthavirodhī viśeṣaṇam nāma . \\
\noindent \textbf{(P\_2,1.1.5)} tasmāt na hāsyati . \\
\noindent \textbf{(P\_2,1.1.5)} atha vā anvayāt viśeṣaṇam bhaviṣyati . \\
\noindent \textbf{(P\_2,1.1.5)} tat yathā . \\
\noindent \textbf{(P\_2,1.1.5)} ghṛtaghaṭaḥ tailaghaṭaḥ iti niṣikte ghṛte taile vā anvyayāt viśeṣaṇam bhavati ayam ghṛtaghaṭaḥ ayam tailaghaṭaḥ iti . \\
\noindent \textbf{(P\_2,1.1.5)} viṣamaḥ upanyāsaḥ . \\
\noindent \textbf{(P\_2,1.1.5)} bhavati hi tatra yā ca yāvatī ca arthamātrā . \\
\noindent \textbf{(P\_2,1.1.5)} aṅga hi bhavān agnau niṣṭapya ghṛtaghaṭam tṛṇakūrcena prakṣālayatu . \\
\noindent \textbf{(P\_2,1.1.5)} na gaṃsyate saḥ viśeṣaḥ . \\
\noindent \textbf{(P\_2,1.1.5)} yathā tarhi mallikāpuṭaḥ campakaputaḥ iti niṣkīrṇāsu api sumanaḥsu anvayāt viśeṣaṇam bhavati ayam mallikapuṭaḥ ayam campakapuṭaḥ iti . \\
\noindent \textbf{(P\_2,1.1.5)} atha vā samarthādhikāraḥ ayam vṛttau kriyate . \\
\noindent \textbf{(P\_2,1.1.5)} sāmartham nama bhedaḥ saṃsargaḥ vā . \\
\noindent \textbf{(P\_2,1.1.5)} aparaḥ āha : bhedasaṃsargau vā sāmarthyam iti . \\
\noindent \textbf{(P\_2,1.1.5)} kaḥ punaḥ bhedaḥ saṃsargaḥ vā . \\
\noindent \textbf{(P\_2,1.1.5)} iha rājñaḥ iti ukte sarvam svam prasaktam puruṣaḥ iti ukte sarvaḥ svāmi prasaktaḥ . \\
\noindent \textbf{(P\_2,1.1.5)} iha idānīm rājapuruṣaḥ iti ukte rājā puruṣam nivartayati anyebhyaḥ svāmibhyaḥ puruṣaḥ api rājānam anyebhyaḥ svebhyaḥ . \\
\noindent \textbf{(P\_2,1.1.5)} evam etasmin ubhayataḥ vyavacchinne yadi jahāti kāmam jahātu . \\
\noindent \textbf{(P\_2,1.1.5)} na jātu cit puruṣamātrasya ānayanam bhaviṣyati . \\
\noindent \textbf{(P\_2,1.1.5)} atha vā punaḥ astu ajahatsvārthā vṛttiḥ . \\
\noindent \textbf{(P\_2,1.1.5)} yuktam punaḥ yat ajahatsvārthā nāma vṛttiḥ syāt . \\
\noindent \textbf{(P\_2,1.1.5)} bāḍham yuktam . \\
\noindent \textbf{(P\_2,1.1.5)} evam hi dṛśyate loke . \\
\noindent \textbf{(P\_2,1.1.5)} bhikṣukaḥ ayam dvitīyām bhikṣām āsādya pūrvām na jahāti sañcayāya pravartate . \\
\noindent \textbf{(P\_2,1.1.5)} nanu ca uktam ubhayoḥ vidyamānasvārthayoḥ dvayoḥ dvivacanam iti dvivacanam prāpnoti iti . \\
\noindent \textbf{(P\_2,1.1.5)} kasyāḥ punaḥ dvivacanam prāpnoti . \\
\noindent \textbf{(P\_2,1.1.5)} prathāmāyāḥ . \\
\noindent \textbf{(P\_2,1.1.5)} na prathamāsamarthaḥ rājā . \\
\noindent \textbf{(P\_2,1.1.5)} ṣaṣṭhyāḥ tarhi . \\
\noindent \textbf{(P\_2,1.1.5)} na ṣaṣṭhīsamarthaḥ puruṣaḥ . \\
\noindent \textbf{(P\_2,1.1.5)} prathamāyāḥ eva tarhi prāpnoti . \\
\noindent \textbf{(P\_2,1.1.5)} nanu ca uktam na prathamāsamarthaḥ rājā iti . \\
\noindent \textbf{(P\_2,1.1.5)} abhinihitaḥ saḥ saḥ arthaḥ antarbhūtaḥ prātipadikārthaḥ sampannaḥ . \\
\noindent \textbf{(P\_2,1.1.5)} tatra prātipadikārthe prathamā iti prathamāyāḥ eva dvivacanam prāpnoti . \\
\noindent \textbf{(P\_2,1.1.5)} saṅghātasya aikārthyāt na avayavasaṅkhyātaḥ subutpattiḥ . \\
\noindent \textbf{(P\_2,1.1.5)} saṅghātasya ekatvam arthaḥ . \\
\noindent \textbf{(P\_2,1.1.5)} tena avayavasaṅkhyātaḥ subutpattiḥ na bhaviṣyati . \\
\noindent \textbf{(P\_2,1.1.5)} parasparavyapekṣā sāmarthyem eke . \\
\noindent \textbf{(P\_2,1.1.5)} parasparavyapekṣā sāmarthyem eke icchanti . \\
\noindent \textbf{(P\_2,1.1.5)} kā punaḥ śabdayoḥ vyapekṣā . \\
\noindent \textbf{(P\_2,1.1.5)} na brūmaḥ śabdayoḥ iti . \\
\noindent \textbf{(P\_2,1.1.5)} kim tarhi . \\
\noindent \textbf{(P\_2,1.1.5)} arthayoḥ . \\
\noindent \textbf{(P\_2,1.1.5)} iha rājñaḥ puruṣaḥ iti ukte rājā puruṣam apekṣate mama ayam iti . \\
\noindent \textbf{(P\_2,1.1.5)} puruṣaḥ api rājānam apekṣate aham asya iti . \\
\noindent \textbf{(P\_2,1.1.5)} tayoḥ abhisambandhasya ṣaṣṭhī vācikā bhavati . \\
\noindent \textbf{(P\_2,1.1.5)} tathā kaṣṭam śritaḥ iti kriyākārakayoḥ abhisambandhasya dvitīyā vācikā bhavati . \\
\noindent \textbf{(P\_2,1.1.6)} atha yadi eva ekārthībhāvaḥ sāmarthyam atha api vyapekṣā sāmarthyam kim gatam etat iyatā sūtreṇa āhosvit anyatarasmin pakṣe bhūyaḥ sūtram kartavyam . \\
\noindent \textbf{(P\_2,1.1.6)} gatam iti āha . \\
\noindent \textbf{(P\_2,1.1.6)} katham . \\
\noindent \textbf{(P\_2,1.1.6)} samaḥ ayam arthaśabdena saha samāsaḥ . \\
\noindent \textbf{(P\_2,1.1.6)} sam ca upasargaḥ . \\
\noindent \textbf{(P\_2,1.1.6)} upasargāḥ ca punaḥ evamātmakāḥ yatra kaḥ cit kriyāvācī śabdaḥ prayujyate tatra kriyāviśeṣam āhuḥ . \\
\noindent \textbf{(P\_2,1.1.6)} na ca iha kaḥ cit kriyāvācī śabdaḥ prayujyate yena samaḥ sāmarthyam syāt . \\
\noindent \textbf{(P\_2,1.1.6)} tatra prayogāt etat gantavyam nūnam atra kaḥ cit prayogārhaḥ śabdaḥ na prayujyate yena samaḥ sāmarthyam iti . \\
\noindent \textbf{(P\_2,1.1.6)} tat yathā . \\
\noindent \textbf{(P\_2,1.1.6)} dhūmam dṛṣṭvā agniḥ atra iti gamyate triviṣṭabdhakam ca dṛṣṭvā parivrājakaḥ iti . \\
\noindent \textbf{(P\_2,1.1.6)} kaḥ punaḥ asau prayogārhaḥ śabdaḥ . \\
\noindent \textbf{(P\_2,1.1.6)} ucyate . \\
\noindent \textbf{(P\_2,1.1.6)} saṅgatārtham samartham saṃsṛṣṭārtham samartham samprekṣitam artham samartham sambaddhārtham samartham iti . \\
\noindent \textbf{(P\_2,1.1.6)} tat yadā tāvat ekārthībhāvaḥ sāmarthyam tadā evam vigrahaḥ kariṣyate saṅgatārthaḥ saṃsṛṣṭārthaḥ samarthaḥ iti . \\
\noindent \textbf{(P\_2,1.1.6)} tat yathā saṅgatam ghṛtam saṅgatam tailam iti ucyate . \\
\noindent \textbf{(P\_2,1.1.6)} ekībhūtam iti gamyate . \\
\noindent \textbf{(P\_2,1.1.6)} yadā vyapekṣā sāmarthyam tadā evam vigrahaḥ kariṣyate samprekṣitārthaḥ samarthaḥ sambaddhārthaḥ samarthaḥ iti . \\
\noindent \textbf{(P\_2,1.1.6)} kaḥ punaḥ iha badhnātyarthaḥ . \\
\noindent \textbf{(P\_2,1.1.6)} sambaddhaḥ iti ucyate yaḥ rajjvā ayasā vā kīle vyatiṣiktaḥ bhavati . \\
\noindent \textbf{(P\_2,1.1.6)} na avaśyam badhnātiḥ vyatiṣaṅge eva vartate . \\
\noindent \textbf{(P\_2,1.1.6)} kim tarhi . \\
\noindent \textbf{(P\_2,1.1.6)} ahānau api vartate . \\
\noindent \textbf{(P\_2,1.1.6)} tat yathā sambaddhau imau damyau iti ucyete yau anyonyam na jahītaḥ . \\
\noindent \textbf{(P\_2,1.1.6)} atha vā bhavati ca evañjātīyakeṣu badhnātiḥ vartate . \\
\noindent \textbf{(P\_2,1.1.6)} tat yathā . \\
\noindent \textbf{(P\_2,1.1.6)} asti naḥ gargaiḥ sambandhaḥ . \\
\noindent \textbf{(P\_2,1.1.6)} asti naḥ vatsaiḥ sambandhaḥ . \\
\noindent \textbf{(P\_2,1.1.6)} saṃyogaḥ iti arthaḥ . \\
\noindent \textbf{(P\_2,1.1.6)} atha etasmin vyapekṣāyām sāmarthye yaḥ asau ekārthībhāvkṛtaḥ viśeṣaḥ sa vaktavyaḥ . \\
\noindent \textbf{(P\_2,1.1.6)} tatra nānākārakāt nighātayuṣmadasmadādeśepratiṣedhaḥ . \\
\noindent \textbf{(P\_2,1.1.6)} tatra etasmin vyapekṣāyām sāmarthye nānākārakāt nighātayuṣmadasmadādeśāḥ prāpnuvanti . \\
\noindent \textbf{(P\_2,1.1.6)} teṣām pratiṣedhaḥ vaktavyaḥ . \\
\noindent \textbf{(P\_2,1.1.6)} nighātaḥ . \\
\noindent \textbf{(P\_2,1.1.6)} ayam daṇḍaḥ . \\
\noindent \textbf{(P\_2,1.1.6)} hara anena . \\
\noindent \textbf{(P\_2,1.1.6)} asti daṇḍasya harateḥ ca vyapekṣā iti kṛtvā nighātaḥ prāpnoti . \\
\noindent \textbf{(P\_2,1.1.6)} yuṣmadasmadādeśāḥ . \\
\noindent \textbf{(P\_2,1.1.6)} odanam paca . \\
\noindent \textbf{(P\_2,1.1.6)} tava bhaviṣyati . \\
\noindent \textbf{(P\_2,1.1.6)} odanam paca mama bhaviṣyati . \\
\noindent \textbf{(P\_2,1.1.6)} asti odanasya yuṣmadasmadoḥ ca vyapekṣā iti kṛtvā vāmnāvādayaḥ prāpnuvanti . \\
\noindent \textbf{(P\_2,1.1.6)} teṣām pratiṣedhaḥ vaktavyaḥ . \\
\noindent \textbf{(P\_2,1.1.6)} kim ucyate nānākārakāt iti yadā tena eva āsajya hriyate . \\
\noindent \textbf{(P\_2,1.1.6)} na api brūmaḥ anyena āsajya hriyate iti . \\
\noindent \textbf{(P\_2,1.1.6)} kim tarhi . \\
\noindent \textbf{(P\_2,1.1.6)} śabdapramāṇakāḥ vayam . \\
\noindent \textbf{(P\_2,1.1.6)} yat śabdaḥ āha tat asmākam pramāṇam . \\
\noindent \textbf{(P\_2,1.1.6)} śabdaḥ ca iha sattām āha . \\
\noindent \textbf{(P\_2,1.1.6)} ayam daṇḍaḥ . \\
\noindent \textbf{(P\_2,1.1.6)} asti iti gamyate . \\
\noindent \textbf{(P\_2,1.1.6)} saḥ daṇḍaḥ kartā bhūtvā anyena śabdena abhisambadhyamānaḥ karaṇam sampadyate . \\
\noindent \textbf{(P\_2,1.1.6)} tat yathā . \\
\noindent \textbf{(P\_2,1.1.6)} kaḥ cit kam cit pṛcchati . \\
\noindent \textbf{(P\_2,1.1.6)} kva devadattaḥ iti . \\
\noindent \textbf{(P\_2,1.1.6)} saḥ tasmai ācaṣṭe . \\
\noindent \textbf{(P\_2,1.1.6)} asau vṛkṣe iti . \\
\noindent \textbf{(P\_2,1.1.6)} katarasmin . \\
\noindent \textbf{(P\_2,1.1.6)} yaḥ tiṣṭhati iti . \\
\noindent \textbf{(P\_2,1.1.6)} saḥ vṛkṣaḥ adhikaraṇam bhūtvā anyena śabdena abhisambadhyamānaḥ kartā sampadyate . \\
\noindent \textbf{(P\_2,1.1.6)} pracaye samāsapratiṣedhaḥ . \\
\noindent \textbf{(P\_2,1.1.6)} pracaye samāsapratiṣedhaḥ vaktavyaḥ . \\
\noindent \textbf{(P\_2,1.1.6)} rājñaḥ gauḥ ca aśvaḥ ca puruṣaḥ ca rājagavāśvapuruṣāḥ iti . \\
\noindent \textbf{(P\_2,1.1.6)} samarthatarāṇām vā . \\
\noindent \textbf{(P\_2,1.1.6)} samarthatarāṇām vā padānām samāsaḥ bhaviṣyati . \\
\noindent \textbf{(P\_2,1.1.6)} kāni punaḥ samarthatarāṇi . \\
\noindent \textbf{(P\_2,1.1.6)} yāni dvandvabhāvīni . \\
\noindent \textbf{(P\_2,1.1.6)} kutaḥ etat . \\
\noindent \textbf{(P\_2,1.1.6)} eṣām hi āśutarā vṛttiḥ prāpnoti . \\
\noindent \textbf{(P\_2,1.1.6)} tat yathā samarthataraḥ ayam māṇavakaḥ adhyayanāya iti ucyate . \\
\noindent \textbf{(P\_2,1.1.6)} āśrutaragranthaḥ iti gamyate . \\
\noindent \textbf{(P\_2,1.1.6)} aparaḥ āha : samarthatarāṇām vā padānām samāsaḥ bhaviṣyati . \\
\noindent \textbf{(P\_2,1.1.6)} kāni punaḥ samarthatarāṇi . \\
\noindent \textbf{(P\_2,1.1.6)} yāni dvandvabhāvīni . \\
\noindent \textbf{(P\_2,1.1.6)} kutaḥ etat . \\
\noindent \textbf{(P\_2,1.1.6)} etāni samānavibhaktīni anyavibhaktiḥ rājā . \\
\noindent \textbf{(P\_2,1.1.6)} bhavati viśeṣaḥ svasmin bhrātari pitṛvyaputre ca . \\
\noindent \textbf{(P\_2,1.1.6)} samudāyasāmarthyāt vā siddham ṣamudāyasāmarthyāt vā punaḥ siddham etat . \\
\noindent \textbf{(P\_2,1.1.6)} samudāyena rājñaḥ sāmarthyam bhavati na avayavena . \\
\noindent \textbf{(P\_2,1.1.6)} aparaḥ āha . \\
\noindent \textbf{(P\_2,1.1.6)} samarthatarāṇām vā samudāyasāmarthyāt . \\
\noindent \textbf{(P\_2,1.1.6)} samarthatarāṇām vā padānām samāsaḥ bhavati . \\
\noindent \textbf{(P\_2,1.1.6)} kutaḥ etat . \\
\noindent \textbf{(P\_2,1.1.6)} samudāyasāmarthyāt eva . \\
\noindent \textbf{(P\_2,1.1.6)} asmin pakṣe vā iti etat asamarthitam bhavati . \\
\noindent \textbf{(P\_2,1.1.6)} etat ca samarthitam . \\
\noindent \textbf{(P\_2,1.1.6)} katham . \\
\noindent \textbf{(P\_2,1.1.6)} na eva vā punaḥ atra rājñaḥ aśvapuruṣau apekṣamāṇasya gavā saha samāsaḥ bhavati . \\
\noindent \textbf{(P\_2,1.1.6)} kim tarhi . \\
\noindent \textbf{(P\_2,1.1.6)} goḥ rājānam apekṣamāṇasya āsvapuruṣābhyām samāsaḥ samāsaḥ bhavati . \\
\noindent \textbf{(P\_2,1.1.6)} pradhānam atra tada gauḥ bhavati . \\
\noindent \textbf{(P\_2,1.1.6)} bhavati ca pradhānasya sāpekṣasya api samāsaḥ . \\
\noindent \textbf{(P\_2,1.1.7)} ākhyātam sāvyayakārakaviśeṣaṇam vākyam . \\
\noindent \textbf{(P\_2,1.1.7)} ākhyātam sāvyayam sakārakam sakārakaviśeṣaṇam vākyasañjñam bhavati vaktavyam . \\
\noindent \textbf{(P\_2,1.1.7)} sāvyayam . \\
\noindent \textbf{(P\_2,1.1.7)} uccaiḥ paṭhati . \\
\noindent \textbf{(P\_2,1.1.7)} nīcaiḥ paṭhati . \\
\noindent \textbf{(P\_2,1.1.7)} sakārakam . \\
\noindent \textbf{(P\_2,1.1.7)} odanam pacati . \\
\noindent \textbf{(P\_2,1.1.7)} sakārakaviśeṣaṇam . \\
\noindent \textbf{(P\_2,1.1.7)} odanam mṛdu viśadam pacati . \\
\noindent \textbf{(P\_2,1.1.7)} sakriyāviśeṣaṇam ca iti vaktavyam . \\
\noindent \textbf{(P\_2,1.1.7)} suṣṭhu pacati . \\
\noindent \textbf{(P\_2,1.1.7)} duṣṭhu pacati . \\
\noindent \textbf{(P\_2,1.1.7)} aparaḥ āha : ākhyātam saviśeṣaṇam iti eva . \\
\noindent \textbf{(P\_2,1.1.7)} sarvāṇi hi etāni kriyāviśeṣaṇāni . \\
\noindent \textbf{(P\_2,1.1.7)} ekatiṅ . \\
\noindent \textbf{(P\_2,1.1.7)} ekatiṅ vākyasañjñam bhavati vaktavyam . \\
\noindent \textbf{(P\_2,1.1.7)} brūhi brūhi . \\
\noindent \textbf{(P\_2,1.1.7)} samānavākye nighātayuṣmadasmadādeśāḥ . \\
\noindent \textbf{(P\_2,1.1.7)} samānavākye iti prakṛtya nighātayuṣmadasmadādeśāḥ vaktavyāḥ . \\
\noindent \textbf{(P\_2,1.1.7)} kim prayojanam . \\
\noindent \textbf{(P\_2,1.1.7)} nānāvākye mā bhūvan nighātādayaḥ iti . \\
\noindent \textbf{(P\_2,1.1.7)} ayam daṇḍaḥ . \\
\noindent \textbf{(P\_2,1.1.7)} hara anena . \\
\noindent \textbf{(P\_2,1.1.7)} odanam paca . \\
\noindent \textbf{(P\_2,1.1.7)} tava bhaviṣyati . \\
\noindent \textbf{(P\_2,1.1.7)} odanam paca . \\
\noindent \textbf{(P\_2,1.1.7)} mama bhaviṣyati . \\
\noindent \textbf{(P\_2,1.1.7)} yoge pratiṣedhaḥ cādibhiḥ . \\
\noindent \textbf{(P\_2,1.1.7)} cādibhiḥ yoge pratiṣedhaḥ vaktavyaḥ . \\
\noindent \textbf{(P\_2,1.1.7)} grāmaḥ tava ca svam mama ca svam . \\
\noindent \textbf{(P\_2,1.1.7)} kimartham icam ucyate . \\
\noindent \textbf{(P\_2,1.1.7)} yathānyāsam eva cādibhiḥ yoge pratiṣedhaḥ ucyate . \\
\noindent \textbf{(P\_2,1.1.7)} idam adya apūrvam kriyate vākyasañjñā samānavākyādhikāraḥ ca . \\
\noindent \textbf{(P\_2,1.1.7)} tat dveṣyam vijānīyāt : sarvam etat vikalpate iti . \\
\noindent \textbf{(P\_2,1.1.7)} tat ācāryaḥ suhṛt bhūtvā anvācaṣṭe cādibhiḥ yoge yathānyāsam eva bhavati iti . \\
\noindent \textbf{(P\_2,1.1.7)} samarthanighāte hi samānādhikaraṇayuktayukteṣu upasaṅkhyānam asamarthatvāt . \\
\noindent \textbf{(P\_2,1.1.7)} samarthanighāte hi samānādhikaraṇayuktayukteṣu upasaṅkhyānam kartavyam syāt . \\
\noindent \textbf{(P\_2,1.1.7)} samānādhikaraṇe . \\
\noindent \textbf{(P\_2,1.1.7)} paṭave te dāsyāmi . \\
\noindent \textbf{(P\_2,1.1.7)} mrdave te dāsyāmi . \\
\noindent \textbf{(P\_2,1.1.7)} samānādhikaraṇe . \\
\noindent \textbf{(P\_2,1.1.7)} yuktayukte . \\
\noindent \textbf{(P\_2,1.1.7)} nadyāḥ tiṣṭhati kūle . \\
\noindent \textbf{(P\_2,1.1.7)} vṛkṣasya lambate śākhā . \\
\noindent \textbf{(P\_2,1.1.7)} śālīnām te odanam dadāmi . \\
\noindent \textbf{(P\_2,1.1.7)} śālīnām me odanam dadāti . \\
\noindent \textbf{(P\_2,1.1.7)} kim punaḥ kāraṇam na sidhyati . \\
\noindent \textbf{(P\_2,1.1.7)} asamarthatvāt . \\
\noindent \textbf{(P\_2,1.1.7)} rājagavīkṣīre dvisamāsaprasaṅgaḥ dviṣaṣṭhībhāvāt . \\
\noindent \textbf{(P\_2,1.1.7)} rājagavīkṣīre dvisamāsaprasaṅgaḥ . \\
\noindent \textbf{(P\_2,1.1.7)} kim kāraṇam . \\
\noindent \textbf{(P\_2,1.1.7)} dviṣaṣṭhībhāvāt . \\
\noindent \textbf{(P\_2,1.1.7)} dve hi atra ṣaṣṭhyau . \\
\noindent \textbf{(P\_2,1.1.7)} rājñaḥ goḥ kṣīram iti . \\
\noindent \textbf{(P\_2,1.1.7)} kim ucyate dvisamāsaprasaṅgaḥ iti yāvatā sup saha supā iti vartate . \\
\noindent \textbf{(P\_2,1.1.7)} dvisamāsaprasaṅgaḥ iti na evam vijñāyate dvayoḥ subantayoḥ samāsaprasaṅgaḥ dvisamāsaprasaṅgaḥ iti . \\
\noindent \textbf{(P\_2,1.1.7)} katham tarhi . \\
\noindent \textbf{(P\_2,1.1.7)} dviprakārasya samāsasya prasaṅgaḥ dvisamāsaprasaṅgaḥ iti . \\
\noindent \textbf{(P\_2,1.1.7)} rājagokṣīram iti api prāpnoti na ca evam bhavitavyam . \\
\noindent \textbf{(P\_2,1.1.7)} bhavitavyam ca yadā etat vākyam bhavati goḥ kṣīram gokṣīram rājñaḥ gokṣīram rājagokṣīram iti . \\
\noindent \textbf{(P\_2,1.1.7)} yadā tu etat vākyam bhavati rājñaḥ goḥ kṣīram iti tadā na bhavitavyam tadā ca prāpnoti . \\
\noindent \textbf{(P\_2,1.1.7)} tadā kasmāt na bhavati . \\
\noindent \textbf{(P\_2,1.1.7)} siddham tu rājaviśiṣṭāyāḥ goḥ kṣīreṇa sāmarthyāt . \\
\noindent \textbf{(P\_2,1.1.7)} siddham etat . \\
\noindent \textbf{(P\_2,1.1.7)} katham . \\
\noindent \textbf{(P\_2,1.1.7)} rājaviśiṣṭāyāḥ goḥ kṣīreṇa saha samāsaḥ bhavati na kevalāyāḥ . \\
\noindent \textbf{(P\_2,1.1.7)} kim vaktavyam etat . \\
\noindent \textbf{(P\_2,1.1.7)} na hi . \\
\noindent \textbf{(P\_2,1.1.7)} katham anucyamānam gaṃsyate . \\
\noindent \textbf{(P\_2,1.1.7)} yathā eva ayam gavi yatate na kṣīramātreṇa santoṣam karoti evam rājani api yatate . \\
\noindent \textbf{(P\_2,1.1.7)} rājñaḥ yā gauḥ tasyāḥ yat kṣīram iti . \\
\noindent \textbf{(P\_2,1.1.7)} na eva vā punaḥ atra goḥ rājānam apekṣamāṇāyāḥ kṣīreṇa saha samāsaḥ prāpnoti . \\
\noindent \textbf{(P\_2,1.1.7)} kim kāraṇam . \\
\noindent \textbf{(P\_2,1.1.7)} asāmarthyāt . \\
\noindent \textbf{(P\_2,1.1.7)} katham asāmarthyam . \\
\noindent \textbf{(P\_2,1.1.7)} sāpekṣam asamartham bhavati iti . \\
\noindent \textbf{(P\_2,1.1.7)} katham tarhi goḥ kṣīram apekṣamāṇāyāḥ rājñā saha samāsaḥ bhavati . \\
\noindent \textbf{(P\_2,1.1.7)} pradhānam atra tada gauḥ bhavati . \\
\noindent \textbf{(P\_2,1.1.7)} bhavati ca pradhānasya sāpekṣasya api samāsaḥ \\
\noindent \textbf{(P\_2,1.1.8)} atha kimartham padavidhau samarthādhikāraḥ kriyate . \\
\noindent \textbf{(P\_2,1.1.8)} padavidhau samarthavacanam varṇāśraye śāstre ānantaryavijñānāt. padavidhau samarthādhikāraḥ kriyate varṇāśraye śāstre ānantaryamātre kāryam yathā vijñāyeta . \\
\noindent \textbf{(P\_2,1.1.8)} tiṣṭhatu dadhi aśāna tvam śākena . \\
\noindent \textbf{(P\_2,1.1.8)} tiṣṭhatu kumārī chatram hara devadatta iti . \\
\noindent \textbf{(P\_2,1.1.8)} samarthādhikārasya vidheyasāmānādhikaraṇyāt nirdeśānarthakyam . \\
\noindent \textbf{(P\_2,1.1.8)} samarthādhikāraḥ ayam vidheyena samānādhikaraṇaḥ . \\
\noindent \textbf{(P\_2,1.1.8)} kim ca vidheyam . \\
\noindent \textbf{(P\_2,1.1.8)} samāsaḥ . \\
\noindent \textbf{(P\_2,1.1.8)} yāvat brūyāt samarthaḥ samāsaḥ iti tāvat samarthaḥ padavidhiḥ . \\
\noindent \textbf{(P\_2,1.1.8)} na ca rājapuruṣaḥ iti etasyām avasthāyām samarthādhikāreṇa kim cit api śakyam pravartayitum nivartayitum vā . \\
\noindent \textbf{(P\_2,1.1.8)} samarthādhikārasya vidheyasāmānādhikaraṇyāt nirdeśaḥ anarthakaḥ . \\
\noindent \textbf{(P\_2,1.1.8)} siddham tu samarthānām iti vacanāt . \\
\noindent \textbf{(P\_2,1.1.8)} siddham etat . \\
\noindent \textbf{(P\_2,1.1.8)} katham . \\
\noindent \textbf{(P\_2,1.1.8)} samarthānām padānām vidhiḥ iti vaktavyam . \\
\noindent \textbf{(P\_2,1.1.8)} evam api dvyekayoḥ na prāpnoti . \\
\noindent \textbf{(P\_2,1.1.8)} ekaśeṣanirdeśāt vā . \\
\noindent \textbf{(P\_2,1.1.8)} atha vā ekaśeṣanirdeśaḥ ayam. samarthasya ca samarthayoḥ ca samarthānām ca samarthānām iti . \\
\noindent \textbf{(P\_2,1.1.8)} evam api ṣaṭprabhṛtīnām eva prāpnoti . \\
\noindent \textbf{(P\_2,1.1.8)} ṣaṭprabhṛtiṣu hi ekaśeṣaḥ parisamāpyate . \\
\noindent \textbf{(P\_2,1.1.8)} na eṣaḥ doṣaḥ . \\
\noindent \textbf{(P\_2,1.1.8)} pratyekam vākyaparisamāptiḥ dṛṣṭā iti dvyekayoḥ api bhaviṣyati . \\
\noindent \textbf{(P\_2,1.1.8)} evam api vivibhaktīnām na prāpnoti . \\
\noindent \textbf{(P\_2,1.1.8)} samarthāt samarthe padāt pade iti . \\
\noindent \textbf{(P\_2,1.1.8)} evam tarhi samarthapadayoḥ vidhiśabdena sarvavibhaktyantaḥ samāsaḥ : samarthasya vidhiḥ samarthavidhiḥ , samarthayoḥ vidhiḥ samarthavidhiḥ , samarthāt vidhiḥ samarthavidhiḥ , samarthe vidhiḥ samarthvidhiḥ . \\
\noindent \textbf{(P\_2,1.1.8)} padasya vidhiḥ padavidhiḥ , padayoḥ vidhiḥ padavidhiḥ , padānām vidhiḥ padavidhiḥ , padāt vidhiḥ padavidhiḥ , pade vidhiḥ padavidhiḥ . \\
\noindent \textbf{(P\_2,1.1.8)} samarthavidhiḥ ca samarthavidhiḥ ca samarthavidhiḥ ca samarthavidhiḥ ca samarthavidhiḥ ca samarthavidhiḥ . \\
\noindent \textbf{(P\_2,1.1.8)} padavidhiḥ ca padavidhiḥ ca padavidhiḥ ca padavidhiḥ ca padavidhiḥ ca padavidhiḥ . \\
\noindent \textbf{(P\_2,1.1.8)} samarthavidhiḥ ca padavidhiḥ ca samarthaḥ padavidhiḥ . \\
\noindent \textbf{(P\_2,1.1.8)} pūrvaḥ samāsaḥ uttarapadalopī yādṛcchikīvibhaktiḥ . \\
\noindent \textbf{(P\_2,1.1.9)} samānādhikaraṇeṣu upasaṅkhyānam asamarthatvāt . \\
\noindent \textbf{(P\_2,1.1.9)} samānādhikaraṇeṣu upasaṅkhyānam kartavyam . \\
\noindent \textbf{(P\_2,1.1.9)} vīraḥ pūruṣaḥ vīrapuruṣaḥ . \\
\noindent \textbf{(P\_2,1.1.9)} kim punaḥ kāraṇam na sidhyati . \\
\noindent \textbf{(P\_2,1.1.9)} asamarthatvāt . \\
\noindent \textbf{(P\_2,1.1.9)} katham asamarthatvam . \\
\noindent \textbf{(P\_2,1.1.9)} dravyam padārthaḥ iti cet. yadi dravyam padārthaḥ na bhavati tadā sāmarthyam . \\
\noindent \textbf{(P\_2,1.1.9)} atha hi guṇaḥ padārthaḥ bhavati tadā sāmarthyam . \\
\noindent \textbf{(P\_2,1.1.9)} anyaḥ hi vīratvam guṇaḥ anyaḥ hi puruṣatvam . \\
\noindent \textbf{(P\_2,1.1.9)} na anyatvam asti iti iyatā sāmarthyam bhavati . \\
\noindent \textbf{(P\_2,1.1.9)} anyaḥ hi devadattaḥ gobhyaḥ aśvebhyaḥ ca na ca tasya etāvatā sāmartham bhavati . \\
\noindent \textbf{(P\_2,1.1.9)} kaḥ vā viśeṣaḥ yat guṇe padārthe sāmarthyam syāt dravye ca na syāt . \\
\noindent \textbf{(P\_2,1.1.9)} eṣaḥ viśeṣaḥ . \\
\noindent \textbf{(P\_2,1.1.9)} ekam tayoḥ adhikaraṇam anyaḥ ca vīratvam guṇaḥ anyaḥ puruṣatvam . \\
\noindent \textbf{(P\_2,1.1.9)} dravyapadārthikasya api tarhi guṇabhedāt sāmarthyam bhaviṣyati . \\
\noindent \textbf{(P\_2,1.1.9)} aśakyaḥ dravyapadārthikena dravyasya guṇakṛtaḥ upakāraḥ pratijñātum . \\
\noindent \textbf{(P\_2,1.1.9)} nanu ca abhyantaraḥ asau bhavati . \\
\noindent \textbf{(P\_2,1.1.9)} yadi api abhyantaraḥ na tu gamyate . \\
\noindent \textbf{(P\_2,1.1.9)} na hi guḍaḥ iti ukte madhuratvam gamyate śṛṅgaveram iti vā kaṭukatvam . \\
\noindent \textbf{(P\_2,1.1.9)} guṇapadārthikena api tarhi aśakyaḥ guṇasya dravyakṛtaḥ upakāraḥ pratijñātum . \\
\noindent \textbf{(P\_2,1.1.9)} atha guṇapadārthikaḥ pratijānīte dravyapadārthikaḥ api kasmāt na pratijānīte . \\
\noindent \textbf{(P\_2,1.1.9)} evam anayoḥ sāmarthyam syāt vā na vā . \\
\noindent \textbf{(P\_2,1.1.9)} kva ca tāvat idam syāt samānādhikaraṇena iti . \\
\noindent \textbf{(P\_2,1.1.9)} yatra sarvam samāman : indraḥ śakraḥ puruhūtaḥ purandaraḥ . \\
\noindent \textbf{(P\_2,1.1.9)} kanduḥ koṣṭhaḥ kuśūlaḥ iti . \\
\noindent \textbf{(P\_2,1.1.9)} na evañjātīyakānām samāsena bhavitavyam pratyayena vā utpattavyam . \\
\noindent \textbf{(P\_2,1.1.9)} kim kāraṇam . \\
\noindent \textbf{(P\_2,1.1.9)} arthagatyarthaḥ śabdaprayogaḥ . \\
\noindent \textbf{(P\_2,1.1.9)} artham pratyāyayiṣyāmi iti śabdaḥ prayujyate . \\
\noindent \textbf{(P\_2,1.1.9)} tatra ekena uktatvāt tasya arthasya dvitīyasya prayogeṇa na bhavitavyam . \\
\noindent \textbf{(P\_2,1.1.9)} kim kāraṇam . \\
\noindent \textbf{(P\_2,1.1.9)} uktārthānām aprayogaḥ iti . \\
\noindent \textbf{(P\_2,1.1.9)} na tarhi idānīm idam bhavati bhṛtyabharaṇīyaḥ iti . \\
\noindent \textbf{(P\_2,1.1.9)} na etau samānārthau . \\
\noindent \textbf{(P\_2,1.1.9)} ekaḥ atra śakyārthe kṛtyaḥ bhavati aparaḥ arhatyarthe : śakyaḥ bhartum bhṛtyaḥ . \\
\noindent \textbf{(P\_2,1.1.9)} arhati bhṛtim bharaṇīyaḥ . \\
\noindent \textbf{(P\_2,1.1.9)} bhṛtyaḥ bharaṇīyaḥ bhṛtyabharaṇīyaḥ iti . \\
\noindent \textbf{(P\_2,1.1.9)} yadi tarhi yatra kim cit samānam kaḥ cit viśeṣaḥ tatra bhavitavyam iha api tarhi prāpnoti . \\
\noindent \textbf{(P\_2,1.1.9)} darśanīyāyāḥ mātā darśanīyamātā iti . \\
\noindent \textbf{(P\_2,1.1.9)} atra api kim cit samānam kaḥ cit viśeṣaḥ . \\
\noindent \textbf{(P\_2,1.1.9)} kim punaḥ tat . \\
\noindent \textbf{(P\_2,1.1.9)} sadbhāvānyabhāvau . \\
\noindent \textbf{(P\_2,1.1.9)} na kva cit sadbhāvānyabhāvau na staḥ ucyate ca samānādhikaraṇena iti . \\
\noindent \textbf{(P\_2,1.1.9)} tatra prakarṣagatiḥ vijñāsyate : yatra sādhīyaḥ sāmānādhikaraṇyam . \\
\noindent \textbf{(P\_2,1.1.9)} kva ca sādhīyaḥ sāmānādhikaraṇyam . \\
\noindent \textbf{(P\_2,1.1.9)} yatra sarvam samānam sadbhāvānyabhāvau dravyam ca. atha vā samānādhikaraṇena iti tat samānam āśrīyate yat samānam bhavati na ca bhavati . \\
\noindent \textbf{(P\_2,1.1.9)} na ca etat samānam kva cit api na bhavati . \\
\noindent \textbf{(P\_2,1.1.9)} atha vā yāvat brūyāt samānadravyeṇa iti tāvat samānādhikaraṇena iti . \\
\noindent \textbf{(P\_2,1.1.9)} dravyam hi loke adhikaraṇam iti ucyate . \\
\noindent \textbf{(P\_2,1.1.9)} tat yathā . \\
\noindent \textbf{(P\_2,1.1.9)} ekasmin dravye vyuditam . \\
\noindent \textbf{(P\_2,1.1.9)} ekasmin adhikaraṇe vyuditam iti . \\
\noindent \textbf{(P\_2,1.1.9)} tathā vyākaraṇe vipratiṣiddham ca anadhikaraṇavāci iti adravyavāci iti gamyate . \\
\noindent \textbf{(P\_2,1.1.9)} evam api idam avaśyam kartavyam samānādhikaraṇam asamarthavat bhavati iti . \\
\noindent \textbf{(P\_2,1.1.9)} kim prayojanam . \\
\noindent \textbf{(P\_2,1.1.9)} sarpiḥ kālakam yajuḥ pītakam iti evamartham . \\
\noindent \textbf{(P\_2,1.1.9)} yadi samānādhikaraṇam asamarthavat bhavati iti ucyate sarpiḥ pīyate yajuḥ kriyate iti atra ṣatvam na prāpnoti . \\
\noindent \textbf{(P\_2,1.1.9)} adhātvabhihitam iti evam tat . \\
\noindent \textbf{(P\_2,1.1.9)} evam ca kṛtvā samānādhikaraṇeṣu upasaṅkhyānam kartavyam . \\
\noindent \textbf{(P\_2,1.1.9)} vīraḥ pūruṣaḥ vīrapuruṣaḥ . \\
\noindent \textbf{(P\_2,1.1.9)} kim kāraṇam asamarthatvāt . \\
\noindent \textbf{(P\_2,1.1.9)} na vā vacanaprāmāṇyāt . \\
\noindent \textbf{(P\_2,1.1.9)} na vā kartavyam . \\
\noindent \textbf{(P\_2,1.1.9)} kim kāraṇam . \\
\noindent \textbf{(P\_2,1.1.9)} vacanaprāmāṇyāt . \\
\noindent \textbf{(P\_2,1.1.9)} vacanaprāmāṇyāt atra samāsaḥ bhaviṣyati . \\
\noindent \textbf{(P\_2,1.1.9)} kim vacanaprāmāṇyam . \\
\noindent \textbf{(P\_2,1.1.9)} samānamadhyamadhyamavīrāḥ ca iti . \\
\noindent \textbf{(P\_2,1.1.9)} luptākhyāteṣu ca . \\
\noindent \textbf{(P\_2,1.1.9)} luptākhyāteṣu ca upasaṅkhyānam kartavyam . \\
\noindent \textbf{(P\_2,1.1.9)} niṣkauśāmbiḥ nirvārāṇasiḥ . \\
\noindent \textbf{(P\_2,1.1.9)} luptākhyāteṣu ca . \\
\noindent \textbf{(P\_2,1.1.9)} kim . \\
\noindent \textbf{(P\_2,1.1.9)} vacanaprāmāṇyāt . \\
\noindent \textbf{(P\_2,1.1.9)} kim vacanaprāmāṇyam . \\
\noindent \textbf{(P\_2,1.1.9)} kugatiprādayaḥ ca iti . \\
\noindent \textbf{(P\_2,1.1.9)} asti anyat etasya vacane prayojanam . \\
\noindent \textbf{(P\_2,1.1.9)} kim . \\
\noindent \textbf{(P\_2,1.1.9)} surājā atirājā iti . \\
\noindent \textbf{(P\_2,1.1.9)} na brūmaḥ vṛttisūtravacanaprāmāṇyāt iti . \\
\noindent \textbf{(P\_2,1.1.9)} kim tarhi . \\
\noindent \textbf{(P\_2,1.1.9)} vārttikavacanaprāmāṇyāt iti . \\
\noindent \textbf{(P\_2,1.1.9)} siddham tu kvāṅksvatidurgativacanāt prādayaḥ ktārthe iti . \\
\noindent \textbf{(P\_2,1.1.9)} tadarthagateḥ vā . \\
\noindent \textbf{(P\_2,1.1.9)} tadarthagateḥ vā punaḥ siddham etat . \\
\noindent \textbf{(P\_2,1.1.9)} kim idam tadarthagateḥ iti . \\
\noindent \textbf{(P\_2,1.1.9)} tasya arthaḥ tadarthaḥ tadarthasya gatiḥ tadarthagatiḥ tadarthagateḥ iti . \\
\noindent \textbf{(P\_2,1.1.9)} yasya arthasya kauśāmbyā sāmarthyam saḥ nisā ucyate . \\
\noindent \textbf{(P\_2,1.1.9)} atha vā saḥ arthaḥ tadarthaḥ tadarthasya gatiḥ tadarthagatiḥ tadarthagateḥ iti . \\
\noindent \textbf{(P\_2,1.1.9)} yaḥ arthaḥ kauśāmbyā samarthaḥ saḥ nisā ucyate . \\
\noindent \textbf{(P\_2,1.1.10)} atha yatra bahūnām samāsaprasaṅgaḥ kim tatra dvayoḥ dvayoḥ samāsaḥ bhavati āhosvit aviśeṣeṇa . \\
\noindent \textbf{(P\_2,1.1.10)} kaḥ ca atra viśeṣaḥ . \\
\noindent \textbf{(P\_2,1.1.10)} samāsaḥ dvayoḥ dvayoḥ cet dvandve anekagrahaṇam . \\
\noindent \textbf{(P\_2,1.1.10)} samāsaḥ dvayoḥ dvayoḥ cet dvandve anekagrahaṇam kartavyam . \\
\noindent \textbf{(P\_2,1.1.10)} carthe dvandvaḥ . \\
\noindent \textbf{(P\_2,1.1.10)} anekam iti vaktavyam iha api yathā syāt . \\
\noindent \textbf{(P\_2,1.1.10)} plakṣanyagrodhakhadirapalāśāḥ iti . \\
\noindent \textbf{(P\_2,1.1.10)} na eṣaḥ doṣaḥ . \\
\noindent \textbf{(P\_2,1.1.10)} atra api dvayoḥ dvayoḥ samāsaḥ bhaviṣyati . \\
\noindent \textbf{(P\_2,1.1.10)} dvayoḥ dvayoḥ samāsaḥ iti cet na bahuṣu dvitvābhāvāt . \\
\noindent \textbf{(P\_2,1.1.10)} dvayoḥ dvayoḥ samāsaḥ iti cet tat na . \\
\noindent \textbf{(P\_2,1.1.10)} kim kāraṇam . \\
\noindent \textbf{(P\_2,1.1.10)} bahuṣu dvitvābhāvāt . \\
\noindent \textbf{(P\_2,1.1.10)} na bahuṣu dvitvam asti . \\
\noindent \textbf{(P\_2,1.1.10)} na avaśyam evam vigrahaḥ kartavyaḥ : plakṣaḥ ca nyagrodhaḥ ca khadiraḥ ca palāśaḥ ca iti . \\
\noindent \textbf{(P\_2,1.1.10)} kim tarhi evam vighrahaḥ kariṣyate : plakṣaḥ ca nyagrodhaḥ ca plakṣanyagrodhau khadiraḥ ca palāśaḥ ca khadirapalāśau plakṣanyagrodhau ca khadirapalāśau plakṣanyagrodhakhadirapalāśāḥ iti . \\
\noindent \textbf{(P\_2,1.1.10)} hotṛpotṛneṣṭodgātāraḥ tarhi na sidhyanti . \\
\noindent \textbf{(P\_2,1.1.10)} hotāpotāneṣṭodgātāraḥ iti prāpnoti . \\
\noindent \textbf{(P\_2,1.1.10)} na ca evam bhavitavyam . \\
\noindent \textbf{(P\_2,1.1.10)} bhavitavyam ca yadā evam vigrahaḥ kriyate hotā ca potā ca hotāpotārau neṣṭā ca udgātā ca neṣṭodgātārau hotāpotārau ca neṣṭodgātārau ca hotāpotāneṣṭodgātāraḥ iti . \\
\noindent \textbf{(P\_2,1.1.10)} hotṛpotṛneṣṭodgātāraḥ tu na sidhyanti . \\
\noindent \textbf{(P\_2,1.1.10)} samāsāntapratiṣedhaḥ ca . \\
\noindent \textbf{(P\_2,1.1.10)} samāsāntasya ca pratiṣedhaḥ vaktavyaḥ . \\
\noindent \textbf{(P\_2,1.1.10)} vāktvaksrugdṛṣadam iti . \\
\noindent \textbf{(P\_2,1.1.10)} vāktvacasrugdṛṣadam iti prāpnoti . \\
\noindent \textbf{(P\_2,1.1.10)} na eṣaḥ doṣaḥ . \\
\noindent \textbf{(P\_2,1.1.10)} atra api pareṇa pareṇa saha samāsaḥ bhaviṣyati . \\
\noindent \textbf{(P\_2,1.1.10)} sruk ca dṛṣadam ca srugdṛṣadam tvak ca srugdṛṣadam ca tvaksrugdṛṣadam vāk ca tvaksrugdṛṣadam ca vāktvaksrugdṛṣadam iti . \\
\noindent \textbf{(P\_2,1.1.10)} hotṛpotṛneṣṭodgātāraḥ evam tarhi sidhyanti . \\
\noindent \textbf{(P\_2,1.1.10)} iha ca susukṣmajaṭakeśena sunatājivāsanā samantaśitirandhreṇa dvayoḥ vṛttau na sidhyati . \\
\noindent \textbf{(P\_2,1.1.10)} astu tarhi aviśeṣeṇa . \\
\noindent \textbf{(P\_2,1.1.10)} aviśeṣeṇa bahuvrīhau anekapadaprasaṅgaḥ . \\
\noindent \textbf{(P\_2,1.1.10)} yadi aviśeṣeṇa bahuvrīhau anekapadaprasaṅgaḥ . \\
\noindent \textbf{(P\_2,1.1.10)} tatra kaḥ doṣaḥ . \\
\noindent \textbf{(P\_2,1.1.10)} tatra svarasamāsāntapuṃvadbhāveṣu doṣaḥ . \\
\noindent \textbf{(P\_2,1.1.10)} tatra svarasamāsāntapuṃvadbhāveṣu doṣaḥ bhavati . \\
\noindent \textbf{(P\_2,1.1.10)} svara . \\
\noindent \textbf{(P\_2,1.1.10)} pūrvaśālāpriyaḥ aparaśālāpriyaḥ . \\
\noindent \textbf{(P\_2,1.1.10)} svara . \\
\noindent \textbf{(P\_2,1.1.10)} samāsānta . \\
\noindent \textbf{(P\_2,1.1.10)} pañcagavapriyaḥ . \\
\noindent \textbf{(P\_2,1.1.10)} samāsānta . \\
\noindent \textbf{(P\_2,1.1.10)} puṃvadbhāva . \\
\noindent \textbf{(P\_2,1.1.10)} khādiretaraśamyam rauravetarśamyam . \\
\noindent \textbf{(P\_2,1.1.10)} na vā avayavatatpuruṣatvāt . na vā eṣaḥ doṣaḥ . \\
\noindent \textbf{(P\_2,1.1.10)} kim kāraṇam . \\
\noindent \textbf{(P\_2,1.1.10)} avayavatatpuruṣatvāt . \\
\noindent \textbf{(P\_2,1.1.10)} avayavaḥ atra tatpuruṣasañjñaḥ tadāśrayau samāsāntapuṃvadbhāvau bhaviṣyataḥ . \\
\noindent \textbf{(P\_2,1.1.10)} svaraḥ katham . \\
\noindent \textbf{(P\_2,1.1.10)} tasya antodāttatvam vipratiṣedhāt . \\
\noindent \textbf{(P\_2,1.1.10)} antodāttatvam kriyatām pūrvapadaprakṛtisvaraḥ iti antodāttatvam bhavati vipratiṣedhena . \\
\noindent \textbf{(P\_2,1.1.10)} na eṣaḥ yuktaḥ vipratiṣedhaḥ . \\
\noindent \textbf{(P\_2,1.1.10)} vipratiṣedhe param iti ucyate . \\
\noindent \textbf{(P\_2,1.1.10)} pūrvam ca antodāttatvam param pūrvapadaprakṛtisvaratvam . \\
\noindent \textbf{(P\_2,1.1.10)} na paravipratiṣedham brūmaḥ . \\
\noindent \textbf{(P\_2,1.1.10)} kim tarhi . \\
\noindent \textbf{(P\_2,1.1.10)} antaraṅgavipratiṣedham . \\
\noindent \textbf{(P\_2,1.1.10)} nimittisvarabalīyastvāt vā . \\
\noindent \textbf{(P\_2,1.1.10)} atha vā nimittasvarāt nimittisvaraḥ balīyān iti vaktavyam . \\
\noindent \textbf{(P\_2,1.1.10)} kim punaḥ nimittam kaḥ vā nimittī . \\
\noindent \textbf{(P\_2,1.1.10)} bahuvrīhiḥ nimittam tatpuruṣaḥ nimittī . \\
\noindent \textbf{(P\_2,1.1.10)} tat tarhi vaktavyam nimittasvarāt nimittisvaraḥ balīyān iti . \\
\noindent \textbf{(P\_2,1.1.10)} na vaktavyam . \\
\noindent \textbf{(P\_2,1.1.10)} ekaśitipātsvaravacanam tu jñāpakam nimittisvarabalīyastvasya . \\
\noindent \textbf{(P\_2,1.1.10)} yat ayam yuktārohyādiṣu ekaśitipacchabdam paṭhati tat jñāpayati ācāryaḥ nimittasvarāt nimittisvaraḥ balīyān iti . \\
\noindent \textbf{(P\_2,1.1.10)} kaḥ punaḥ arhati yuktārohyādiṣu ekaśitipacchabdam paṭhitum . \\
\noindent \textbf{(P\_2,1.1.10)} evam kila nāma paṭhyate ekaḥ śitiḥ ekaśitiḥ ekaḥ śitiḥ pādaḥ yasya iti . \\
\noindent \textbf{(P\_2,1.1.10)} tat ca na . \\
\noindent \textbf{(P\_2,1.1.10)} evam vigrahaḥ kariṣyate . \\
\noindent \textbf{(P\_2,1.1.10)} ekaḥ śitiḥ eṣu te ime ekaśitayaḥ ekaśitayaḥ pādāḥ yasya iti . \\
\noindent \textbf{(P\_2,1.1.10)} atha api evam vigrahaḥ kriyate ekaḥ śitiḥ ekaśitiḥ ekaḥ śitiḥ pādaḥ yasya iti evam api na arthaḥ pāṭhena . \\
\noindent \textbf{(P\_2,1.1.10)} igante dvigau iti eṣaḥ svaraḥ atra bādhakaḥ bhaviṣyati . \\
\noindent \textbf{(P\_2,1.1.10)} asya tarhi bahuvrīhyavayavasya tatpuruṣañjñā prāpnoti susukṣmajaṭakeśena sunatājivāsanā samantaśitirandhreṇa iti . \\
\noindent \textbf{(P\_2,1.1.10)} tatra kaḥ doṣaḥ . \\
\noindent \textbf{(P\_2,1.1.10)} tasya antodāttatvam vipratiṣedhāt iti antodāttatvam syāt vipratiṣedhena . \\
\noindent \textbf{(P\_2,1.1.10)} na eṣaḥ doṣaḥ . \\
\noindent \textbf{(P\_2,1.1.10)} na idam bahuvrīhyavayavasya tatpuruṣasya lakṣaṇam ārabhyate . \\
\noindent \textbf{(P\_2,1.1.10)} kim tarhi . \\
\noindent \textbf{(P\_2,1.1.10)} yasya bahuvrīhyavayavasya tatpuruṣasya tat lakṣaṇam asti tasya antodāttatvam bhaviṣyati vipratiṣedhena . \\
\noindent \textbf{(P\_2,1.1.10)} nanu ca asya api asti kim viśeṣaṇam viśeṣyeṇa bahulam iti. bahulavacanāt na bhaviṣyati . \\
\noindent \textbf{(P\_2,1.1.10)} asya tarhi bahuvrīhyavayavasya tatpuruṣañjñā prāpnoti . \\
\noindent \textbf{(P\_2,1.1.10)} adhikaṣaṣṭivarṣaḥ iti . \\
\noindent \textbf{(P\_2,1.1.10)} tatra kaḥ doṣaḥ . \\
\noindent \textbf{(P\_2,1.1.10)} tasya antodāttatvam vipratiṣedhāt iti antodāttatvam syāt vipratiṣedhena . \\
\noindent \textbf{(P\_2,1.1.10)} na eṣaḥ doṣaḥ . \\
\noindent \textbf{(P\_2,1.1.10)} igante dvigau iti eṣaḥ svaraḥ bhaviṣyati . \\
\noindent \textbf{(P\_2,1.1.10)} yaḥ tarhi na igantaḥ adhikaśatavarṣaḥ iti . \\
\noindent \textbf{(P\_2,1.1.10)} iha ca api adhikaṣaṣṭivarṣaḥ iti samāsantaḥ prāpnoti . \\
\noindent \textbf{(P\_2,1.1.10)} ḍacprakaraṇe saṅkhyāyāḥ tatpuruṣasya upasaṅkhyānam nistriṃśādyartham iti . \\
\noindent \textbf{(P\_2,1.1.10)} na eṣaḥ doṣaḥ . \\
\noindent \textbf{(P\_2,1.1.10)} avyayādeḥ iti evam tat . \\
\noindent \textbf{(P\_2,1.1.10)} kim punaḥ kāraṇam avyayādeḥ iti evam tat . \\
\noindent \textbf{(P\_2,1.1.10)} iha mā bhūt gotriṃśat gocatvāriṃśat iti . \\
\noindent \textbf{(P\_2,1.1.10)} bahuvrīhisañjñā tarhi prāpnoti. saṅkhyayā avyayāsannādūrādhikasaṅhyāḥ saṅkhyeye iti . \\
\noindent \textbf{(P\_2,1.1.10)} na saṅkhyām saṅkhyeye vartayiṣyāmaḥ . \\
\noindent \textbf{(P\_2,1.1.10)} katham . \\
\noindent \textbf{(P\_2,1.1.10)} evam vigrahaḥ kariṣyate adhikā ṣaṣṭiḥ varṣāṇām asya iti . \\
\noindent \textbf{(P\_2,1.1.10)} yathā tarhi saḥ yogaḥ pratyākhyāyate tathā pūrveṇa prāpnoti . \\
\noindent \textbf{(P\_2,1.1.10)} katham ca saḥ yogaḥ pratyākhyāyate . \\
\noindent \textbf{(P\_2,1.1.10)} aśiṣyaḥ saṅkhyottarapadaḥ saṅkhyā iva abhidhāyitvāt iti . \\
\noindent \textbf{(P\_2,1.1.10)} pratyākhyāte tasmin yoge saṅkhyām saṅkhyeye vartayiṣyāmaḥ . \\
\noindent \textbf{(P\_2,1.1.10)} tatra evam vigrahaḥ kariṣyate adhikā ṣaṣṭiḥ varṣāṇi asya iti . \\
\noindent \textbf{(P\_2,1.1.10)} sarvatha vayam adhikaṣaṣṭivarṣāt na mucyāmahe . \\
\noindent \textbf{(P\_2,1.1.10)} katham . \\
\noindent \textbf{(P\_2,1.1.10)} yāvatā saḥ ca yogaḥ pratyākhyāyate ayam ca vigrahaḥ asti adhikā ṣaṣṭiḥ varṣāṇām asya iti . \\
\noindent \textbf{(P\_2,1.1.10)} yat tu tat uktam adhikaṣaṣṭivarṣaḥ na sidhyati iti saḥ siddhaḥ bhavati . \\
\noindent \textbf{(P\_2,1.1.10)} katham . \\
\noindent \textbf{(P\_2,1.1.10)} yāvatā saḥ ca yogaḥ pratyākhyāyate ayam ca vigrahaḥ asti adhikā ṣaṣṭiḥ varṣāṇi asya iti . \\
\noindent \textbf{(P\_2,1.1.10)} adhikaśatavarṣaḥ tu na sidhyati . \\
\noindent \textbf{(P\_2,1.1.10)} kartavyaḥ atra yatnaḥ \\
\noindent \textbf{(P\_2,1.2)} sup iti kimartham . \\
\noindent \textbf{(P\_2,1.2)} karoṣi aṭan . \\
\noindent \textbf{(P\_2,1.2)} na etat asti . \\
\noindent \textbf{(P\_2,1.2)} asāmarthyāt atra na bhaviṣyati . \\
\noindent \textbf{(P\_2,1.2)} katham asāmarthyam . \\
\noindent \textbf{(P\_2,1.2)} samānādhikaraṇam asamarthavat bhavati iti . \\
\noindent \textbf{(P\_2,1.2)} idam tarhi . \\
\noindent \textbf{(P\_2,1.2)} pīḍye pīḍyamāna . \\
\noindent \textbf{(P\_2,1.2)} idam ca api udāharaṇam karoṣi aṭan . \\
\noindent \textbf{(P\_2,1.2)} nanu ca uktam asāmarthyāt atra na bhaviṣyati . \\
\noindent \textbf{(P\_2,1.2)} katham asāmarthyam . \\
\noindent \textbf{(P\_2,1.2)} samānādhikaraṇam asamarthavat bhavati iti . \\
\noindent \textbf{(P\_2,1.2)} na eṣaḥ doṣaḥ . \\
\noindent \textbf{(P\_2,1.2)} adhātvabhihitam iti evam tat . \\
\noindent \textbf{(P\_2,1.2)} āmantritasya parāṅgvadbhāve ṣaṣṭhyāmantritakārakavacanam . \\
\noindent \textbf{(P\_2,1.2)} āmantritasya parāṅgvadbhāve ṣaṣṭhyantam āmantritakārakam ca parasya aṅgavat bhavati iti vaktavyam . \\
\noindent \textbf{(P\_2,1.2)} ṣaṣthyantam tāvat . \\
\noindent \textbf{(P\_2,1.2)} madrāṇām rājan magadhānām rājan . \\
\noindent \textbf{(P\_2,1.2)} āmantritakārakam . \\
\noindent \textbf{(P\_2,1.2)} kuṇḍena aṭan . \\
\noindent \textbf{(P\_2,1.2)} na asti atra viśeṣaḥ parāṅgavadbhāve sati asati vā . \\
\noindent \textbf{(P\_2,1.2)} idam tarhi . \\
\noindent \textbf{(P\_2,1.2)} paraśunā vṛścan . \\
\noindent \textbf{(P\_2,1.2)} tannimittagrahaṇam vā . \\
\noindent \textbf{(P\_2,1.2)} tannimittagrahaṇam vā kartavyam. āmantritanimittam parasya aṅgavat bhavati iti vaktavyam . \\
\noindent \textbf{(P\_2,1.2)} tat ca avaśyam anyatarat vaktavyam . \\
\noindent \textbf{(P\_2,1.2)} avacane hi subantamātraprasaṅgaḥ . \\
\noindent \textbf{(P\_2,1.2)} anucyamāne hi etasmin subantramātrasya parāṅgavadbhāvaḥ prāpnoti . \\
\noindent \textbf{(P\_2,1.2)} asya api prasajyeta . \\
\noindent \textbf{(P\_2,1.2)} kṣtreṇa agne svāyuḥ saṃrabhasya mitreṇa agne mitradheye yatasva . \\
\noindent \textbf{(P\_2,1.2)} kim punaḥ atra jyāyaḥ . \\
\noindent \textbf{(P\_2,1.2)} tannimittagrahaṇam eva jyāyaḥ . \\
\noindent \textbf{(P\_2,1.2)} idam api siddham bhavati . \\
\noindent \textbf{(P\_2,1.2)} goṣu svāmin aśveṣu svāmin . \\
\noindent \textbf{(P\_2,1.2)} etat hi na eva ṣaṣthyantam na api āmantritakārakam . \\
\noindent \textbf{(P\_2,1.2)} subantasya parāṅgavadbhāve samānādhikaraṇasya upasaṅkhyanam ananantaratvāt . \\
\noindent \textbf{(P\_2,1.2)} subantasya parāṅgavadbhāve samānādhikaraṇasya upasaṅkhyanam kartavyam . \\
\noindent \textbf{(P\_2,1.2)} tīkṣṇayā sūcyā sīvyan tīkṣṇena paraśunā vṛścan . \\
\noindent \textbf{(P\_2,1.2)} kim punaḥ kāraṇam na sidhyati . \\
\noindent \textbf{(P\_2,1.2)} ananantaratvāt . \\
\noindent \textbf{(P\_2,1.2)} nanu ca parasya parāṅgavadbhāve kṛte pūrvasya api bhaviṣyati . \\
\noindent \textbf{(P\_2,1.2)} svare avadhāraṇāt ca . \\
\noindent \textbf{(P\_2,1.2)} svare avadhāraṇāt ca na sidhyati . \\
\noindent \textbf{(P\_2,1.2)} svare avadhāraṇam kriyate na ānantarye . \\
\noindent \textbf{(P\_2,1.2)} param api chandasi . \\
\noindent \textbf{(P\_2,1.2)} param api chandasi pūrvasya aṅgavat bhavati iti vaktavyam . \\
\noindent \textbf{(P\_2,1.2)} ā te pitaḥ marutām sumnam etu . \\
\noindent \textbf{(P\_2,1.2)} prati tvā duhitaḥ divaḥ . \\
\noindent \textbf{(P\_2,1.2)} vṛṇīṣva duhitaḥ divaḥ . \\
\noindent \textbf{(P\_2,1.2)} avyayapratiṣedhaḥ ca . \\
\noindent \textbf{(P\_2,1.2)} avyayānām ca pratiṣedhaḥ vaktavyaḥ . \\
\noindent \textbf{(P\_2,1.2)} uccaiḥ adhīyāna nīcaiḥ adhīyāna . \\
\noindent \textbf{(P\_2,1.2)} anavyayībhāvasya . \\
\noindent \textbf{(P\_2,1.2)} anavyayībhāvasya iti vaktavyam . \\
\noindent \textbf{(P\_2,1.2)} iha mā bhūt . \\
\noindent \textbf{(P\_2,1.2)} upāgni adhīyāna pratyagni adhīyāna . \\
\noindent \textbf{(P\_2,1.2)} atha kimartham svare avadhāraṇam kriyate . \\
\noindent \textbf{(P\_2,1.2)} svare avadhāraṇam subalopārtham . \\
\noindent \textbf{(P\_2,1.2)} svare avadhāraṇam kriyate sublaḥ mā bhūt iti . \\
\noindent \textbf{(P\_2,1.2)} paraśunā vṛścan . \\
\noindent \textbf{(P\_2,1.2)} na vā subantaikāntatvāt . \\
\noindent \textbf{(P\_2,1.2)} na vā kartavyam . \\
\noindent \textbf{(P\_2,1.2)} kim kāraṇam . \\
\noindent \textbf{(P\_2,1.2)} subantaikāntatvāt . \\
\noindent \textbf{(P\_2,1.2)} subantaikāntaḥ parāṅgavadbhāvaḥ bhavati . \\
\noindent \textbf{(P\_2,1.2)} prātipadikaikāntaḥ tu sublope . \\
\noindent \textbf{(P\_2,1.2)} prātipadikaikāntaḥ tu bhavati sublope kṛte . \\
\noindent \textbf{(P\_2,1.2)} pratyayalakṣaṇena subantaikāntatā syāt . \\
\noindent \textbf{(P\_2,1.2)} tasmāt svare avadhāraṇam na kartavyam subalopārtham . \\
\noindent \textbf{(P\_2,1.2)} prātipadikasthāyāḥ supaḥ luk ucyate . \\
\noindent \textbf{(P\_2,1.2)} tasmāt svaragrahaṇena na arthaḥ . \\
\noindent \textbf{(P\_2,1.2)} idam tarhi prayojanam ṣatvaṇatve mā bhūtām iti . \\
\noindent \textbf{(P\_2,1.2)} kūpe siñcan carma naman iti . \\
\noindent \textbf{(P\_2,1.2)} etat api na asti prayojanam . \\
\noindent \textbf{(P\_2,1.2)} iha tāvat kūpe siñcan iti svāśrayam padāditvam bhaviṣyati . \\
\noindent \textbf{(P\_2,1.2)} carma naman iti pūrvapadāt sañjñāyām agaḥ iti etasmāt niyamāt na bhaviṣyati . \\
\noindent \textbf{(P\_2,1.2)} nanu ca samāse etat bhavati pūrvapadam uttarapadam iti . \\
\noindent \textbf{(P\_2,1.2)} na iti āha . \\
\noindent \textbf{(P\_2,1.2)} aviśeṣeṇa etat bhavati . \\
\noindent \textbf{(P\_2,1.2)} pūrvam padam pūrvapadam uttaram padam uttarapadam iti . \\
\noindent \textbf{(P\_2,1.3)} prāgvacanam kimartham . \\
\noindent \textbf{(P\_2,1.3)} prāgvacanam sañjñānivṛttyartham . \\
\noindent \textbf{(P\_2,1.3)} prāgvacanam kriyate samāsasañjñāyāḥ anivṛttiḥ yathā syāt . \\
\noindent \textbf{(P\_2,1.3)} akriyamāṇe hi prāgvacane anavakāśāḥ avyayībhāvādayaḥ sañjñāḥ samāsasañjñām bādheran . \\
\noindent \textbf{(P\_2,1.3)} tāḥ mā bādhiṣata iti prāgvacanam kriyate . \\
\noindent \textbf{(P\_2,1.3)} atha kriyamāṇe api prāgvacane yāvatā anavakāśāḥ avyayībhāvādayaḥ sañjñāḥ kasmāt eva na bādhante . \\
\noindent \textbf{(P\_2,1.3)} kriyamāṇe hi prāgvacane satyām samāsasañjñāyām etāḥ avayavasañjñāḥ ārabhyante . \\
\noindent \textbf{(P\_2,1.3)} tatra vacanasamāveśaḥ bhaviṣyati . \\
\noindent \textbf{(P\_2,1.3)} samāsasañjñā api anavakāśā . \\
\noindent \textbf{(P\_2,1.3)} sā vacanāt bhaviṣyati . \\
\noindent \textbf{(P\_2,1.3)} sāvakāśā samāsasañjñā . \\
\noindent \textbf{(P\_2,1.3)} kaḥ avakāśaḥ . \\
\noindent \textbf{(P\_2,1.3)} vispaṣṭādīni avakāśaḥ . \\
\noindent \textbf{(P\_2,1.3)} vispaṣṭam paṭuḥ vispaṣṭapaṭuḥ iti . \\
\noindent \textbf{(P\_2,1.3)} na eṣaḥ asti avakāśaḥ . \\
\noindent \textbf{(P\_2,1.3)} eṣā hi ācāryasya śailī lakṣyate yena eva avayavakāryam bhavati tena eva samudāyakāryam api bhavati . \\
\noindent \textbf{(P\_2,1.3)} yena eva avayavakāryam svaraḥ tena eva samudākāryam api samāsaḥ bhaviṣyati . \\
\noindent \textbf{(P\_2,1.3)} vispaṣṭādīni guṇavacaneṣu iti . \\
\noindent \textbf{(P\_2,1.3)} idam tarhi kākatālīyam ajākṛpāṇīyam . \\
\noindent \textbf{(P\_2,1.3)} atra api yena eva avayavakāryam pratyayotpattiḥ kriyate tena eva samudākāryam samāsasañjñā bhaviṣyati . \\
\noindent \textbf{(P\_2,1.3)} samāsāt ca tadviṣayāt . \\
\noindent \textbf{(P\_2,1.3)} idam tarhi punārājaḥ punargavaḥ . \\
\noindent \textbf{(P\_2,1.3)} atra api avaśyam tatpuruṣasañjñā vaktavya tatpuruṣāśrayaḥ samāsāntaḥ yathā syāt . \\
\noindent \textbf{(P\_2,1.3)} idam tarhi . \\
\noindent \textbf{(P\_2,1.3)} punarādheyam . \\
\noindent \textbf{(P\_2,1.3)} atra api avaśyam gatisañjñā vaktavyā gatikārakopapadāt kṛt iti eṣaḥ svaraḥ yathā syāt . \\
\noindent \textbf{(P\_2,1.3)} idam tarhi punarutsyūtam vāsaḥ deyam . \\
\noindent \textbf{(P\_2,1.3)} atra api avaśyam gatisañjñā vaktavyā gatiḥ gatau iti nighātaḥ yathā syāt . \\
\noindent \textbf{(P\_2,1.3)} yadi tat na asti punaścanasau chandasi iti. sati tasmin tena eva siddham . \\
\noindent \textbf{(P\_2,1.3)} evam api ekā sañjñā iti vacanāt na asti yaugapadyena sambhavaḥ . \\
\noindent \textbf{(P\_2,1.3)} paryāyaḥ prasajyeta . \\
\noindent \textbf{(P\_2,1.3)} tasmāt prāgvacanam kartavyam . \\
\noindent \textbf{(P\_2,1.4)} sahavacanam kimartham . \\
\noindent \textbf{(P\_2,1.4)} sahavacanam pṛthak asamāsārtham . \\
\noindent \textbf{(P\_2,1.4)} sahagrahaṇam kriyate sahabhūtayoḥ samāsañjñā yathā syāt ekaikasya mā bhūt iti . \\
\noindent \textbf{(P\_2,1.4)} kim ca syāt . \\
\noindent \textbf{(P\_2,1.4)} yadi ekaikasya samāsañjñā syāt iha ṛkpādaḥ iti samāsāntaḥ prasajyeta . \\
\noindent \textbf{(P\_2,1.4)} iha ca rājāśvaḥ iti dvau svarau syātām . \\
\noindent \textbf{(P\_2,1.4)} katham ca kṛtvā ekaikasya sañjñā prāpnoti . \\
\noindent \textbf{(P\_2,1.4)} pratyekam vākyaparisamāptiḥ dṛṣṭā . \\
\noindent \textbf{(P\_2,1.4)} tat yathā vṛddhiguṇasañjñe pratyekam bhavataḥ . \\
\noindent \textbf{(P\_2,1.4)} nanu ca ayam api asti dṛṣṭāntaḥ samudāye vākyaparisamāptiḥ iti . \\
\noindent \textbf{(P\_2,1.4)} tat yathā gargāḥ śatam daṇḍyantām iti . \\
\noindent \textbf{(P\_2,1.4)} arthinaḥ ca rājānaḥ hiraṇyena bhavanti na ca pratyekam daṇḍayanti . \\
\noindent \textbf{(P\_2,1.4)} sati etasmin dṛṣṭānte yadi tatra pratyekam iti ucyate iha api sahagrahaṇam kartavyam . \\
\noindent \textbf{(P\_2,1.4)} atha tatra antareṇa pratyekam iti vacanam pratyekam guṇavṛddhisañjñe bhavataḥ iha api na arthaḥ sahagrahaṇena . \\
\noindent \textbf{(P\_2,1.4)} evam tarhi siddhe sati yat sahagrahaṇam karoti tasya etat prayojanam yogāṅgam yathā vijñāyeta . \\
\noindent \textbf{(P\_2,1.4)} sati ca yogāṅge yogavibhāgaḥ kariṣyate . \\
\noindent \textbf{(P\_2,1.4)} saha sup samasyate . \\
\noindent \textbf{(P\_2,1.4)} kena saha . \\
\noindent \textbf{(P\_2,1.4)} samarthena. anuvyācalat anuprāviśat . \\
\noindent \textbf{(P\_2,1.4)} tataḥ supā . \\
\noindent \textbf{(P\_2,1.4)} supā ca saha sup samasyate . \\
\noindent \textbf{(P\_2,1.4)} adhikāraḥ ca lakṣaṇam ca . \\
\noindent \textbf{(P\_2,1.4)} yasya samāsasya anyat lakṣaṇam na asti idam tasya lakṣaṇam bhaviṣyati . \\
\noindent \textbf{(P\_2,1.4)} punarutsyūtam vāsaḥ deyam punarniṣkṛtaḥ rathaḥ iti . \\
\noindent \textbf{(P\_2,1.4)} ivena vibhaktyalopaḥ pūrvapadaprakṛtisvaratvam ca . \\
\noindent \textbf{(P\_2,1.4)} ivena saha samāsaḥ vibhaktyalopaḥ pūrvapadaprakṛtisvaratvam ca vaktavyam . \\
\noindent \textbf{(P\_2,1.4)} vāsasīiva kanye iva . \\
\noindent \textbf{(P\_2,1.5)} kimartham mahatī sañjñā kriyate . \\
\noindent \textbf{(P\_2,1.5)} anvarthasañjñā yathā vijñayeta . \\
\noindent \textbf{(P\_2,1.5)} anavyayam avyayam bhavati iti avyayībhāvaḥ . \\
\noindent \textbf{(P\_2,1.5)} avyayībhāvaḥ ca samāsaḥ avyayasañjñaḥ bhavati iti etat na vaktavyam bhavati . \\
\noindent \textbf{(P\_2,1.6)} iha kasmāt na bhavati . \\
\noindent \textbf{(P\_2,1.6)} sumadrāḥ sumagadhāḥ saputraḥ sacchātraḥ iti . \\
\noindent \textbf{(P\_2,1.6)} samṛddhau sākalye iti ca prāpnoti . \\
\noindent \textbf{(P\_2,1.6)} na eṣaḥ doṣaḥ . \\
\noindent \textbf{(P\_2,1.6)} iha kaḥ cit samāsaḥ pūrvapadārthapradhānaḥ , kaḥ cit uttarapadārthapradhānaḥ , kaḥ cit anyapadārthapradhānaḥ , kaḥ cit ubhayapadārthapradhānaḥ . \\
\noindent \textbf{(P\_2,1.6)} pūrvapadārthapradhānaḥ avyayībhāvaḥ , uttarapadārthapradhānaḥ tatpuruṣaḥ , anyapadārthapradhānaḥ bahuvrīhiḥ ubhayapadārthapradhānaḥ dvandvaḥ . \\
\noindent \textbf{(P\_2,1.6)} na ca atra pūrvapadārthaprādhānyam gamyate . \\
\noindent \textbf{(P\_2,1.6)} atha vā na ime samāsārthāḥ nirdiśyante . \\
\noindent \textbf{(P\_2,1.6)} kim tarhi . \\
\noindent \textbf{(P\_2,1.6)} avyayārthāḥ nirdiśyante ime . \\
\noindent \textbf{(P\_2,1.6)} eteṣu artheṣu yat avyayam vartate tat subantena samasyate iti . \\
\noindent \textbf{(P\_2,1.7)} asādṛśye iti kimartham . \\
\noindent \textbf{(P\_2,1.7)} yathā devadattaḥ tathā yajñadattaḥ iti . \\
\noindent \textbf{(P\_2,1.7)} asādṛśye iti ucyate . \\
\noindent \textbf{(P\_2,1.7)} tatra idam na sidhyati : yathāśakti yathābalam iti . \\
\noindent \textbf{(P\_2,1.7)} kim kāraṇam . \\
\noindent \textbf{(P\_2,1.7)} yathā iti ayam prakāravacane thāl saḥ ca sādṛśye vartate . \\
\noindent \textbf{(P\_2,1.7)} na eṣaḥ doṣaḥ . \\
\noindent \textbf{(P\_2,1.7)} ayam yathāśabdaḥ asti eva avyutpannam prātipadikam vīpsāvācī . \\
\noindent \textbf{(P\_2,1.7)} asti prakāravacane thāl . \\
\noindent \textbf{(P\_2,1.7)} tat yat avyutpannam prātipadikam vīpsāvāci tasya idam grahaṇam . \\
\noindent \textbf{(P\_2,1.7)} atha yaḥ prakāravacane thāl tasya grahaṇam kasmāt na bhavati . \\
\noindent \textbf{(P\_2,1.7)} pūrveṇa prāpnoti sādṛśyasampatti iti . \\
\noindent \textbf{(P\_2,1.7)} pratiṣedhavacanasāmarthyāt na bhaviṣyati . \\
\noindent \textbf{(P\_2,1.9)} sup iti vartamāne punaḥ subgrahaṇam kimartham . \\
\noindent \textbf{(P\_2,1.9)} avyayam iti evam tat abhūt submātre yathā syāt . \\
\noindent \textbf{(P\_2,1.9)} māṣaprati sūpaprati odanaprati . \\
\noindent \textbf{(P\_2,1.10)} akṣādayaḥ tṛtīyāntāḥ pūrvoktasya yathā na tat . akṣādayaḥ tṛtīyāntāḥ pariṇā saha samasyante iti vaktavyam . \\
\noindent \textbf{(P\_2,1.10)} pūrvoktasya yathā na tat . \\
\noindent \textbf{(P\_2,1.10)} ayathājātīyake dyotye . \\
\noindent \textbf{(P\_2,1.10)} akṣeṇa na tathā vṛttam yathā pūrvam iti akṣapari śalākāpari . \\
\noindent \textbf{(P\_2,1.10)} ekatve akṣaśalākayoḥ . \\
\noindent \textbf{(P\_2,1.10)} akṣaśalākayoḥ ca ekavacanāntayoḥ iti vaktavyam . \\
\noindent \textbf{(P\_2,1.10)} iha mā bhūt . \\
\noindent \textbf{(P\_2,1.10)} akṣābhyām vṛttam akṣaiḥ vṛttam . \\
\noindent \textbf{(P\_2,1.10)} kitavavyavahāre ca . \\
\noindent \textbf{(P\_2,1.10)} kitavavyavahāre iti vaktavyam . \\
\noindent \textbf{(P\_2,1.10)} iha mā bhūt . \\
\noindent \textbf{(P\_2,1.10)} akṣeṇa idam na vṛttam śakaṭena yathā pūrvam . \\
\noindent \textbf{(P\_2,1.11-12)} KA\_I,380.7-12 Ro\_II,574-575 {1/12}     yogavibhāgaḥ kartavyaḥ . \\
\noindent \textbf{(P\_2,1.11-12)} KA\_I,380.7-12 Ro\_II,574-575 {2/12}     vibhāṣā iti ayam adhikāraḥ. tataḥ apaparibahirañcavaḥ pañcamyā iti . \\
\noindent \textbf{(P\_2,1.11-12)} KA\_I,380.7-12 Ro\_II,574-575 {3/12}     pañcamīgrahaṇam śakyam akartum . \\
\noindent \textbf{(P\_2,1.11-12)} KA\_I,380.7-12 Ro\_II,574-575 {4/12}     katham . \\
\noindent \textbf{(P\_2,1.11-12)} KA\_I,380.7-12 Ro\_II,574-575 {5/12}     subantena iti vartate etaiḥ ca karmapravacanīyaiḥ yoge pañcamī vidhīyate . \\
\noindent \textbf{(P\_2,1.11-12)} KA\_I,380.7-12 Ro\_II,574-575 {6/12}     tatra antareṇa api pañcamīgrahaṇam pañcamyantena eva samāsaḥ bhaviṣyati . \\
\noindent \textbf{(P\_2,1.11-12)} KA\_I,380.7-12 Ro\_II,574-575 {7/12}     idam tarhi prayojanam . \\
\noindent \textbf{(P\_2,1.11-12)} KA\_I,380.7-12 Ro\_II,574-575 {8/12}     bahiḥśabdena yoge pañcamī na vidhīyate . \\
\noindent \textbf{(P\_2,1.11-12)} KA\_I,380.7-12 Ro\_II,574-575 {9/12}     tatra api yathā syāt iti . \\
\noindent \textbf{(P\_2,1.11-12)} KA\_I,380.7-12 Ro\_II,574-575 {10/12}     bahirgrāmāt . \\
\noindent \textbf{(P\_2,1.11-12)} KA\_I,380.7-12 Ro\_II,574-575 {11/12}     atha kriyamāṇe api pañcamīgrahaṇe yāvatā bahiḥśabdena yoge pañcamī na vidhīyate katham eva etat sidhyati . \\
\noindent \textbf{(P\_2,1.11-12)} KA\_I,380.7-12 Ro\_II,574-575 {12/12}     pañcamīgrahaṇasāmarthyāt . \\
\noindent \textbf{(P\_2,1.13)} maryādābhividhigrahaṇam śakyam akartum . \\
\noindent \textbf{(P\_2,1.13)} katham . \\
\noindent \textbf{(P\_2,1.13)} pañcamyantena iti vartate āṅā ca karmapravacanīyayukte pañcamī vidhīyate . \\
\noindent \textbf{(P\_2,1.13)} etayoḥ ca eva arthayoḥ āṅ karmapravacanīyasañjñaḥ bhavati na anyatra . \\
\noindent \textbf{(P\_2,1.16)} kim udāharaṇam . \\
\noindent \textbf{(P\_2,1.16)} anugaṅgam hāstinapuram anugaṅgam vārāṇasī anuśoṇam pāṭaliputram . \\
\noindent \textbf{(P\_2,1.16)} yasya ca āyāmaḥ iti ucyate gaṅgā ca api āyatā vārāṇasī api āyatā . \\
\noindent \textbf{(P\_2,1.16)} tatra kutaḥ etat gaṅgayā saha samāsaḥ bhaviṣyati na punaḥ vārāṇasyā iti . \\
\noindent \textbf{(P\_2,1.16)} evarm tarhi lakṣaṇena iti vartate gaṅgā ca eva hi lakṣaṇam na vārāṇasī . \\
\noindent \textbf{(P\_2,1.16)} atha vā yasya ca āyāmaḥ iti ucyate gaṅgā ca api āyatā vārāṇasī api āyatā . \\
\noindent \textbf{(P\_2,1.16)} tatra prakarṣagatiḥ vijñāsyate : sādhīyaḥ yasya āyāmaḥ iti . \\
\noindent \textbf{(P\_2,1.16)} sādhīyaḥ ca gaṅgāyāḥ na vārāṇasyāḥ . \\
\noindent \textbf{(P\_2,1.17)} kimarthaḥ cakāraḥ . \\
\noindent \textbf{(P\_2,1.17)} evakārārthaḥ . \\
\noindent \textbf{(P\_2,1.17)} tiṣṭhadguprabhṛtīni eva . \\
\noindent \textbf{(P\_2,1.17)} kva mā bhūt . \\
\noindent \textbf{(P\_2,1.17)} paramam tiṣṭhadgu . \\
\noindent \textbf{(P\_2,1.17)} tiṣṭhadgu kālaviśeṣe . \\
\noindent \textbf{(P\_2,1.17)} tiṣṭhadgu kālaviśeṣe iti vaktavyam . \\
\noindent \textbf{(P\_2,1.17)} tiṣṭhanti gāvaḥ asmin kāle tiṣṭhadgu . \\
\noindent \textbf{(P\_2,1.17)} vahadgu . \\
\noindent \textbf{(P\_2,1.17)} khaleyavādīni prathamāntāni anyapadārthe . \\
\noindent \textbf{(P\_2,1.17)} khaleyavādīni prathamāntāni anyapadārthe samasyante . \\
\noindent \textbf{(P\_2,1.17)} khaleyavam khalebusam lūnayavam lūyamānayavam pūtayavam pūyamānayavam . \\
\noindent \textbf{(P\_2,1.18)} vāvacanam kimartham . \\
\noindent \textbf{(P\_2,1.18)} vibhāṣā samāsaḥ yatha syāt . \\
\noindent \textbf{(P\_2,1.18)} samāsena mukte vākyam api yathā syāt . \\
\noindent \textbf{(P\_2,1.18)} pāram gaṅgāyāḥ iti . \\
\noindent \textbf{(P\_2,1.18)} na etat asti prayojanam . \\
\noindent \textbf{(P\_2,1.18)} prakṛtā mahāvibhāṣā . \\
\noindent \textbf{(P\_2,1.18)} tayā vākyam api bhaviṣyati . \\
\noindent \textbf{(P\_2,1.18)} idam tarhi prayojanam avyayībhāvena mukte ṣaṣṭhīsamāsaḥ yathā syāt . \\
\noindent \textbf{(P\_2,1.18)} gaṅgāpāram iti . \\
\noindent \textbf{(P\_2,1.18)} etat api na asti prayojanam . \\
\noindent \textbf{(P\_2,1.18)} ayam api vibhāṣā ṣaṣṭhīsamāsaḥ api . \\
\noindent \textbf{(P\_2,1.18)} tau ubhau vacanāt bhaviṣyataḥ . \\
\noindent \textbf{(P\_2,1.18)} ataḥ uttaram paṭhati . \\
\noindent \textbf{(P\_2,1.18)} pāre madhye ṣaṣṭhyā vāvacanam . \\
\noindent \textbf{(P\_2,1.18)} pāre madhye ṣaṣṭhyā vā iti vaktavyam . \\
\noindent \textbf{(P\_2,1.18)} avacane hi ṣaṣṭhīsamāsābhāvaḥ yathā ekadeśipradhāne . \\
\noindent \textbf{(P\_2,1.18)} akriyamāṇe hi vāvacane ṣaṣṭhīsamāsasya abhāvaḥ syāt yathā ekadeśipradhāne . \\
\noindent \textbf{(P\_2,1.18)} tat yatha ekadeśisamāsena mukte ṣaṣṭhīsamāsaḥ na bhavati . \\
\noindent \textbf{(P\_2,1.18)} kim punaḥ kāraṇam ekadeśisamāsena mukte ṣaṣṭhīsamāsaḥ na bhavati . \\
\noindent \textbf{(P\_2,1.18)} samāsataddhitānām vṛttiḥ vibhāṣā . \\
\noindent \textbf{(P\_2,1.18)} vṛttiviṣaye nityaḥ apavādaḥ . \\
\noindent \textbf{(P\_2,1.18)} iha punaḥ vāvacane kriyamāṇe ekayā vṛttiḥ vibhāṣā aparayā vṛttiviṣaye vibhāṣāpavādaḥ . \\
\noindent \textbf{(P\_2,1.18)} ekārāntanipātanam ca . \\
\noindent \textbf{(P\_2,1.18)} ekārāntanipātanam ca kartavyam . \\
\noindent \textbf{(P\_2,1.18)} pāregaṅgam iti . \\
\noindent \textbf{(P\_2,1.18)} na kartavyam . \\
\noindent \textbf{(P\_2,1.18)} saptamyāḥ alukā siddham . \\
\noindent \textbf{(P\_2,1.18)} bhavet siddham yadā saptamī yadā tu anyāḥ vibhaktayaḥ tadā na sidhyati . \\
\noindent \textbf{(P\_2,1.20)} nadībhiḥ saṅkhyāsamāse anyapadārthe pratiṣedhaḥ . \\
\noindent \textbf{(P\_2,1.20)} nadībhiḥ saṅkhyāsamāse anyapadārthe pratiṣedhaḥ vaktavyaḥ . \\
\noindent \textbf{(P\_2,1.20)} dvīrāvatīkaḥ deśaḥ trīrāvatīkaḥ deśaḥ . \\
\noindent \textbf{(P\_2,1.20)} nadībhiḥ saṅkhyā iti prāpnoti . \\
\noindent \textbf{(P\_2,1.20)} na vaktavyaḥ . \\
\noindent \textbf{(P\_2,1.20)} iha kaḥ cit samāsaḥ pūrvapadārthapradhānaḥ , kaḥ cit uttarapadārthapradhānaḥ , kaḥ cit anyapadārthapradhānaḥ , kaḥ cit ubhayapadārthapradhānaḥ . \\
\noindent \textbf{(P\_2,1.20)} pūrvapadārthapradhānaḥ avyayībhāvaḥ , uttarapadārthapradhānaḥ tatpuruṣaḥ , anyapadārthapradhānaḥ bahuvrīhiḥ , ubhayapadārthapradhānaḥ dvandvaḥ. na ca atra pūrvapadārthaprādhānyam gamyate . \\
\noindent \textbf{(P\_2,1.20)} nanu ca yat yena ucyate saḥ tasya arthaḥ bhavati . \\
\noindent \textbf{(P\_2,1.20)} atra ca vayam etābhyām padābhyām etam artham ucyamānam paśyāmaḥ . \\
\noindent \textbf{(P\_2,1.20)} etat eva ca jānīmaḥ yat yena ucyate saḥ tasya arthaḥ iti . \\
\noindent \textbf{(P\_2,1.20)} api ca anyapadārthatā na prakalpeta . \\
\noindent \textbf{(P\_2,1.20)} citraguḥ śabalaguḥ iti . \\
\noindent \textbf{(P\_2,1.20)} kim kāraṇam . \\
\noindent \textbf{(P\_2,1.20)} atra api hi vayam etābhyām śabdābhyām etam artham ucyamānam paśyāmaḥ . \\
\noindent \textbf{(P\_2,1.20)} yadi api atra etābhyām śabdābhyām eṣaḥ arthaḥ ucyate anyapadārthaḥ api tu gamyate . \\
\noindent \textbf{(P\_2,1.20)} tatra anyapadārthāśrayaḥ bahuvrīhiḥ bhaviṣyati . \\
\noindent \textbf{(P\_2,1.20)} iha api tarhi anyapadārthaḥ gamyate svapadārthaḥ api tu gamyate . \\
\noindent \textbf{(P\_2,1.20)} tatra svapadārthāśrayaḥ avyayībhāvaḥ prāpnoti . \\
\noindent \textbf{(P\_2,1.20)} evam tarhi idam iha sampradhāryam . \\
\noindent \textbf{(P\_2,1.20)} avyayībhāvaḥ kriyatām bahuvrīhiḥ iti . \\
\noindent \textbf{(P\_2,1.20)} bahuvrīhiḥ bhaviṣyati vipratiṣedhena . \\
\noindent \textbf{(P\_2,1.20)} bhavet ekasañjñādhikāre siddham paraṅkāryatve tu na sidhyati . \\
\noindent \textbf{(P\_2,1.20)} ārambhasāmarthyāt avyayībhāvaḥ prāpnoti paraṅkāryatvāt ca bahuvrīhiḥ . \\
\noindent \textbf{(P\_2,1.20)} paraṅkāryatve ca na doṣaḥ . \\
\noindent \textbf{(P\_2,1.20)} nadībhiḥ saṅkhyāyāḥ samāhāre avyayībhāvaḥ vaktavyaḥ . \\
\noindent \textbf{(P\_2,1.20)} saḥ ca avaśyam vaktavyaḥ . \\
\noindent \textbf{(P\_2,1.20)} sarvam ekanadītare . \\
\noindent \textbf{(P\_2,1.23)} dvigoḥ tatpuruṣatve kāni prayojanāni . \\
\noindent \textbf{(P\_2,1.23)} dvigoḥ tatpuruṣatve samāsāntāḥ prayojanam . \\
\noindent \textbf{(P\_2,1.23)} pañcagavam daśagavam pañcarājam daśarājam . \\
\noindent \textbf{(P\_2,1.24)} śritādiṣu gamigāmyādīnām upasaṅkhyānam . \\
\noindent \textbf{(P\_2,1.24)} śritādiṣu gamigāmyādīnām upasaṅkhyānam kartavyam . \\
\noindent \textbf{(P\_2,1.24)} grāmam gamī gramagamī gramam gāmī grāmagāmī . \\
\noindent \textbf{(P\_2,1.24)} śritādibhiḥ ahīne dvitīyāsamāsavacanānarthakyam bahuvrīhikṛtatvāt . \\
\noindent \textbf{(P\_2,1.24)} śritādibhiḥ ahīnavācinyāḥ dvitīyāyāḥ samāsavacanam anarthakam . \\
\noindent \textbf{(P\_2,1.24)} kim kāraṇam . \\
\noindent \textbf{(P\_2,1.24)} bahuvrīhikṛtatvāt . \\
\noindent \textbf{(P\_2,1.24)} iha hi yaḥ kaṣṭam śritaḥ kaṣṭam anena śritam bhavati . \\
\noindent \textbf{(P\_2,1.24)} tatra bahuvrīhiṇā siddham . \\
\noindent \textbf{(P\_2,1.24)} ahīne dvītīyāsvaravacanānarthakyam ca . \\
\noindent \textbf{(P\_2,1.24)} ahīne dvitīyā pūrvapadam prakṛtisvaram bhavati iti etat svaravacanam anarthakam . \\
\noindent \textbf{(P\_2,1.24)} kim kāraṇam . \\
\noindent \textbf{(P\_2,1.24)} bahuvrīhikṛtatvāt eva . \\
\noindent \textbf{(P\_2,1.24)} jātisvaraprasaṅgaḥ tu . \\
\noindent \textbf{(P\_2,1.24)} jātisvaraḥ tu prāpnoti . \\
\noindent \textbf{(P\_2,1.24)} grāmatataḥ araṇyagataḥ . \\
\noindent \textbf{(P\_2,1.24)} jātikālasukhādibhyaḥ anācchādanāt ktaḥ akṛtamitapratipannāḥ iti . \\
\noindent \textbf{(P\_2,1.24)} tatra jātādiṣu vāvacanāt siddham . \\
\noindent \textbf{(P\_2,1.24)} yat etat vā jāte iti etat vā jātādiṣu iti vakṣyāmi . \\
\noindent \textbf{(P\_2,1.24)} ime jātādayaḥ bhaviṣyanti . \\
\noindent \textbf{(P\_2,1.24)} nanu ca bhedaḥ bhavati . \\
\noindent \textbf{(P\_2,1.24)} bahuvrīhau sati samāsāntodāttatvena api bhavitavyam pūrvapadaprakṛtisvaratvena api tatpuruṣatve sati pūrvapadaprakṛtisvaratvena eva . \\
\noindent \textbf{(P\_2,1.24)} na asti bhedaḥ . \\
\noindent \textbf{(P\_2,1.24)} yaḥ api tatpuruṣam ārabhate na tasya daṇḍavāritaḥ bahuvrīhiḥ . \\
\noindent \textbf{(P\_2,1.24)} tatra tatpuruṣe sati dvau samāsau dvau svarau . \\
\noindent \textbf{(P\_2,1.24)} bahuvrīhau sati ekaḥ samāsaḥ dvisvaratvam . \\
\noindent \textbf{(P\_2,1.24)} evam tarhi siddhe sati yat tatpuruṣam śāsti tat jñāpayati ācāryaḥ samāne arthe kevalam vigrahabhedāt yatra tatpuruṣaḥ prāpnoti bahuvrīhiḥ ca tatra tatpuruṣaḥ bhavati iti . \\
\noindent \textbf{(P\_2,1.24)} kim etasya jñāpane prayojanam . \\
\noindent \textbf{(P\_2,1.24)} rājñaḥ sakhā rājasakhaḥ . \\
\noindent \textbf{(P\_2,1.24)} rājā sakhā asya iti bahuvrīhiḥ na bhavati . \\
\noindent \textbf{(P\_2,1.24)} na etat jñāpakasādhyam apavādaiḥ utsargāḥ bādhyante iti . \\
\noindent \textbf{(P\_2,1.24)} bādhakena anena bhavitavyam sāmānyavihitasya viśeṣavihitena . \\
\noindent \textbf{(P\_2,1.24)} atha na sāmānyavihitaḥ . \\
\noindent \textbf{(P\_2,1.24)} yat ucyate bahuvrīhikṛtatvāt iti etat ayuktam . \\
\noindent \textbf{(P\_2,1.24)} asti khalu api viśeṣaḥ bahuvrīheḥ tatpuruṣasya ca . \\
\noindent \textbf{(P\_2,1.24)} kim śabdakṛtaḥ atha arthakṛtaḥ . \\
\noindent \textbf{(P\_2,1.24)} śabdakṛtaḥ va arthakṛtaḥ ca . \\
\noindent \textbf{(P\_2,1.24)} śabdakṛtaḥ tāvat . \\
\noindent \textbf{(P\_2,1.24)} bahuvrīhau sati kapā bhavitavyam . \\
\noindent \textbf{(P\_2,1.24)} tatpuruṣe sati na bhavitavyam . \\
\noindent \textbf{(P\_2,1.24)} arthakṛtaḥ . \\
\noindent \textbf{(P\_2,1.24)} tatpuruṣe sati ruhādīnām ktaḥ kartari bhavati dhātvarthasya anapavarge . \\
\noindent \textbf{(P\_2,1.24)} ārūḍhaḥ vṛkṣam devadattaḥ iti . \\
\noindent \textbf{(P\_2,1.24)} bahuvrīhau vyapavṛkte karmaṇi bhavati . \\
\noindent \textbf{(P\_2,1.24)} ārūḍhaḥ vṛkṣaḥ devadattena iti . \\
\noindent \textbf{(P\_2,1.24)} anyathājātīyakaḥ khalu api pratyakṣeṇa arthasampratyayaḥ anyathājātīyakaḥ sambandhāt . \\
\noindent \textbf{(P\_2,1.24)} rājñaḥ sakhā rājasakhā . \\
\noindent \textbf{(P\_2,1.24)} sambandhāt etat gantavyam nūnam rāja api asya sakhā iti . \\
\noindent \textbf{(P\_2,1.24)} ubhayam khalu api iṣyate : svasti somasakhā punaḥ ehi . \\
\noindent \textbf{(P\_2,1.24)} gavāṅsakhaḥ iti . \\
\noindent \textbf{(P\_2,1.26)} kim udāharaṇam . \\
\noindent \textbf{(P\_2,1.26)} khaṭvārūḍhaḥ jālmaḥ . \\
\noindent \textbf{(P\_2,1.26)} kṣepe iti ucyate . \\
\noindent \textbf{(P\_2,1.26)} kaḥ kṣepaḥ nāma . \\
\noindent \textbf{(P\_2,1.26)} adhītya snātvā gurubhiḥ anujñātena khaṭvā āroḍhavyā . \\
\noindent \textbf{(P\_2,1.26)} yaḥ idānīm ataḥ anyatha karoti saḥ khaṭvārūḍhaḥ ayam jālmaḥ . \\
\noindent \textbf{(P\_2,1.26)} na ativratavān iti . \\
\noindent \textbf{(P\_2,1.29)} atyantasaṃyoge samāsasya aviśeṣavacanāt ktena samāsavacanānarthakyam . \\
\noindent \textbf{(P\_2,1.29)} atyantasaṃyoge samāsasya aviśeṣavacanāt ktāntena ca aktāntena ca kālāḥ ktāntena iti samāsavacanam anarthakam . \\
\noindent \textbf{(P\_2,1.29)} atyantasaṃyoge iti eva siddham . \\
\noindent \textbf{(P\_2,1.29)} anatyantasaṃyogārtham tu . \\
\noindent \textbf{(P\_2,1.29)} anatyantasaṃyogārtham tarhi idam vaktavyam . \\
\noindent \textbf{(P\_2,1.29)} ṣaṭ muhūrtāḥ carācarāḥ . \\
\noindent \textbf{(P\_2,1.29)} te kadā cit ahaḥ gacchanti kadā cit rātrim . \\
\noindent \textbf{(P\_2,1.29)} tat ucyate ahargatāḥ rātrigatāḥ iti . \\
\noindent \textbf{(P\_2,1.29)} na etat asti . \\
\noindent \textbf{(P\_2,1.29)} gatagrahaṇāt api etat siddham . \\
\noindent \textbf{(P\_2,1.29)} idam tarhi . \\
\noindent \textbf{(P\_2,1.29)} aharatisṛtāḥ rātryatisṛtāḥ māsapramitaḥ candramāḥ . \\
\noindent \textbf{(P\_2,1.30)} tatkṛtārthena iti kimartham . \\
\noindent \textbf{(P\_2,1.30)} dadhnā paṭuḥ ghṛtena paṭuḥ . \\
\noindent \textbf{(P\_2,1.30)} na etat asti . \\
\noindent \textbf{(P\_2,1.30)} asāmarthyāt atra na bhaviṣyati . \\
\noindent \textbf{(P\_2,1.30)} katham asāmarthyam . \\
\noindent \textbf{(P\_2,1.30)} sāpekṣam asamartham bhavati iti . \\
\noindent \textbf{(P\_2,1.30)} na hi dadhnaḥ paṭunā sāmarthyam . \\
\noindent \textbf{(P\_2,1.30)} kena tarhi . \\
\noindent \textbf{(P\_2,1.30)} bhujinā . \\
\noindent \textbf{(P\_2,1.30)} dadhnā bhuṅkte paṭuḥ iti . \\
\noindent \textbf{(P\_2,1.30)} iha api tarhi na prāpnoti . \\
\noindent \textbf{(P\_2,1.30)} śaṅkulākhaṇḍaḥ kirikāṇaḥ iti . \\
\noindent \textbf{(P\_2,1.30)} atra api na śaṅkulāyāḥ khaṇḍena sāmarthyam . \\
\noindent \textbf{(P\_2,1.30)} kena tarhi . \\
\noindent \textbf{(P\_2,1.30)} karotinā . \\
\noindent \textbf{(P\_2,1.30)} śaṅkulayā kṛtaḥ khaṇḍaḥ iti . \\
\noindent \textbf{(P\_2,1.30)} vacanāt bhaviṣyati . \\
\noindent \textbf{(P\_2,1.30)} iha api vacanāt bhaviṣyati dadhnā paṭuḥ ghṛtena paṭuḥ iti . \\
\noindent \textbf{(P\_2,1.30)} tasmāt tatkṛtārthagrahaṇam kartavyam . \\
\noindent \textbf{(P\_2,1.30)} guṇavacanena iti kimartham . \\
\noindent \textbf{(P\_2,1.30)} gobhiḥ vapāvān dhānyena dhanavān . \\
\noindent \textbf{(P\_2,1.30)} kim punaḥ iha udāharaṇam . \\
\noindent \textbf{(P\_2,1.30)} śaṅkulākhaṇḍaḥ devadattaḥ iti . \\
\noindent \textbf{(P\_2,1.30)} katham punaḥ guṇavacanena samāsaḥ ucyamānaḥ dravyavacanena syāt . \\
\noindent \textbf{(P\_2,1.30)} iha tṛtīyā tatkṛtārthena guṇena iti iyatā siddham . \\
\noindent \textbf{(P\_2,1.30)} saḥ ayam evam siddhe sati yat vacanagrahaṇam karoti tasya etat prayojanam evam yathā vijñāyeta guṇam uktavatā guṇavacanena iti . \\
\noindent \textbf{(P\_2,1.30)} katham punaḥ ayam guṇavacanaḥ san dravyavacanaḥ sampadyate . \\
\noindent \textbf{(P\_2,1.30)} ārabhyate tatra matublopaḥ guṇavacanebhyaḥ matupaḥ luk iti . \\
\noindent \textbf{(P\_2,1.30)} tat yathā śuklaguṇaḥ śuklaḥ kṛṣṇaguṇaḥ kṛṣṇaḥ evam khaṇḍaguṇaḥ khaṇḍaḥ . \\
\noindent \textbf{(P\_2,1.30)} yadi evam na arthaḥ kṛtārthagrahaṇena . \\
\noindent \textbf{(P\_2,1.30)} bhavati hi śaṅkulāyāḥ khaṇḍena sāmarthyam . \\
\noindent \textbf{(P\_2,1.30)} asāmarthyāt ca atra na bhaviṣyati dadhnā paṭuḥ ghṛtena paṭuḥ iti . \\
\noindent \textbf{(P\_2,1.30)} tasmat na arthaḥ tatkṛtārthagrahaṇena . \\
\noindent \textbf{(P\_2,1.30)} tṛtīyāsamāse arthagrahaṇam anarthakam arthagatiḥ hi avacanāt . \\
\noindent \textbf{(P\_2,1.30)} tṛtīyāsamāse arthagrahaṇam anarthakam . \\
\noindent \textbf{(P\_2,1.30)} kim kāraṇam . \\
\noindent \textbf{(P\_2,1.30)} arthagatiḥ hi avacanāt . \\
\noindent \textbf{(P\_2,1.30)} antareṇa api vacanam arthagatiḥ bhaviṣyati . \\
\noindent \textbf{(P\_2,1.30)} nirdeśyam iti cet tṛtīyārthanirdeśaḥ api . \\
\noindent \textbf{(P\_2,1.30)} atha evam api nirdeśaḥ kartavyaḥ iti cet tṛtīyārthanirdeśaḥ api kartavyaḥ syāt . \\
\noindent \textbf{(P\_2,1.30)} tṛtīyā tadarthakṛtārthena iti vaktavyam . \\
\noindent \textbf{(P\_2,1.30)} tat tarhi vaktavyam . \\
\noindent \textbf{(P\_2,1.30)} na vaktavyam . \\
\noindent \textbf{(P\_2,1.30)} na ayam arthanirdeśaḥ . \\
\noindent \textbf{(P\_2,1.30)} kim tarhi . \\
\noindent \textbf{(P\_2,1.30)} yogāṅgam idam nirdiśyate . \\
\noindent \textbf{(P\_2,1.30)} sati ca yogāṅge yogavibhāgaḥ kariṣyate . \\
\noindent \textbf{(P\_2,1.30)} tṛtīyā tatkṛtena guṇavacanena samasyate . \\
\noindent \textbf{(P\_2,1.30)} tataḥ arthena . \\
\noindent \textbf{(P\_2,1.30)} arthaśabdena ca tṛtīyā samasyate . \\
\noindent \textbf{(P\_2,1.30)} dhānyāṛthaḥ vasanārthaḥ . \\
\noindent \textbf{(P\_2,1.30)} pūrvasadṛśasamonārtha iti arthagrahaṇam na kartavyam bhavati . \\
\noindent \textbf{(P\_2,1.31)} pūrvādiṣu avarasya upasaṅkhyānam . \\
\noindent \textbf{(P\_2,1.31)} pūrvādiṣu avarasya upasaṅkhyānam . \\
\noindent \textbf{(P\_2,1.31)} māsāvaraḥ ayam saṃvatsarāvaraḥ ayam . \\
\noindent \textbf{(P\_2,1.31)} sadṛśagrahaṇe uktam . \\
\noindent \textbf{(P\_2,1.31)} kim uktam . \\
\noindent \textbf{(P\_2,1.31)} sadrśagrahaṇam anarthakam tṛtīyāsamāsavacanāt . \\
\noindent \textbf{(P\_2,1.31)} ṣaṣṭhyartham iti cet tṛtīyāsamāsavacanānarthakyam iti . \\
\noindent \textbf{(P\_2,1.32)} kartṛkaraṇe kṛtā ktena . \\
\noindent \textbf{(P\_2,1.32)} kartṛkaraṇe kṛtā ktena iti vaktavyam. ahihataḥ nakhanirbhinnaḥ dātralūnaḥ paraśucchinnaḥ . \\
\noindent \textbf{(P\_2,1.32)} kṛtā ktena iti kimartham . \\
\noindent \textbf{(P\_2,1.32)} iha mā bhūt . \\
\noindent \textbf{(P\_2,1.32)} dātreṇa lūnavān paraśunā chinnavān . \\
\noindent \textbf{(P\_2,1.32)} tat tarhi vaktavyam . \\
\noindent \textbf{(P\_2,1.32)} na vaktavyam . \\
\noindent \textbf{(P\_2,1.32)} bahulavacanāt siddham . \\
\noindent \textbf{(P\_2,1.33)} kṛtryaiḥ adhikārthavacane anyatra api dṛśyate .kṛtryaiḥ adhikārthavacane anyatra api dṛśyate iti vaktavyam . \\
\noindent \textbf{(P\_2,1.33)} busopendhyam tṛṇopendhyam ghanaghātyam . \\
\noindent \textbf{(P\_2,1.33)} sādhanam kṛtā iti vā pādahārakādyartham . \\
\noindent \textbf{(P\_2,1.33)} atha vā sādhanam kṛtā saha samasyate iti vaktavyam . \\
\noindent \textbf{(P\_2,1.33)} kim prayojanam . \\
\noindent \textbf{(P\_2,1.33)} pādahārakādyartham . \\
\noindent \textbf{(P\_2,1.33)} pādābhyām hriyate pādahārakaḥ gale copyate galecopakaḥ . \\
\noindent \textbf{(P\_2,1.34-35)} KA\_I,386.18-388.4 Ro\_II,595-597 {1/37}     annena vyañjanam bhakṣyeṇa miśrīkaraṇam iti asamarthasamāsaḥ . \\
\noindent \textbf{(P\_2,1.34-35)} KA\_I,386.18-388.4 Ro\_II,595-597 {2/37}     annena vyañjanam bhakṣyeṇa miśrīkaraṇam iti asamarthasamāsaḥ ayam draṣṭavyaḥ . \\
\noindent \textbf{(P\_2,1.34-35)} KA\_I,386.18-388.4 Ro\_II,595-597 {3/37}     kim kāraṇam . \\
\noindent \textbf{(P\_2,1.34-35)} KA\_I,386.18-388.4 Ro\_II,595-597 {4/37}     kārakāṇām kriyayā sāmarthyāt . \\
\noindent \textbf{(P\_2,1.34-35)} KA\_I,386.18-388.4 Ro\_II,595-597 {5/37}     kārakāṇām kriyayā sāmarthyam bhavati na teṣām anyonyena . \\
\noindent \textbf{(P\_2,1.34-35)} KA\_I,386.18-388.4 Ro\_II,595-597 {6/37}     tat yathā niśrayaṇyā dvābhyām kāṣṭhābhyām sāmarthyam na teṣām anyonyena . \\
\noindent \textbf{(P\_2,1.34-35)} KA\_I,386.18-388.4 Ro\_II,595-597 {7/37}     evam tarhi āha ayam annena vyañjanam bhakṣyeṇa miśrīkaraṇam iti na ca asti sāmarthyam . \\
\noindent \textbf{(P\_2,1.34-35)} KA\_I,386.18-388.4 Ro\_II,595-597 {8/37}     tatra vacanāt samāsaḥ bhaviṣyati . \\
\noindent \textbf{(P\_2,1.34-35)} KA\_I,386.18-388.4 Ro\_II,595-597 {9/37}     vacanaprāmāṇyāt iti cet nānākārakāṇām pratiṣedhaḥ . vacanaprāmāṇyāt iti cet nānākārakāṇām pratiṣedhaḥ vaktavyaḥ . \\
\noindent \textbf{(P\_2,1.34-35)} KA\_I,386.18-388.4 Ro\_II,595-597 {10/37}     tiṣṭhatu dadhnā odanaḥ bhujyate devadattena . \\
\noindent \textbf{(P\_2,1.34-35)} KA\_I,386.18-388.4 Ro\_II,595-597 {11/37}     siddham tu samānādhikaraṇādhikāre ktaḥ tṛtīyāpūrvapadaḥ uttarapadalopaḥ ca ṣiddham etat . \\
\noindent \textbf{(P\_2,1.34-35)} KA\_I,386.18-388.4 Ro\_II,595-597 {12/37}     katham . \\
\noindent \textbf{(P\_2,1.34-35)} KA\_I,386.18-388.4 Ro\_II,595-597 {13/37}     samānādhikaraṇādhikāre vaktavyam ktaḥ tṛtīyāpūrvapadaḥ samasyatesupā uttarapadasya ca lopaḥ bhavati iti . \\
\noindent \textbf{(P\_2,1.34-35)} KA\_I,386.18-388.4 Ro\_II,595-597 {14/37}     dadhnā upasiktaḥ dadhyupasiktaḥ dadhyupasiktaḥ odanaḥ dadhyodanaḥ guḍena saṃsṛṣṭāḥ guḍasaṃsṛṣṭāḥ , guḍasaṃsṛṣṭāḥ dhānāḥ guḍadhānāḥ . \\
\noindent \textbf{(P\_2,1.34-35)} KA\_I,386.18-388.4 Ro\_II,595-597 {15/37}     ṣaṣṭhīsamāsaḥ ca yuktapūrṇāntaḥ . \\
\noindent \textbf{(P\_2,1.34-35)} KA\_I,386.18-388.4 Ro\_II,595-597 {16/37}     ṣaṣṭhīsamāsaḥ ca yuktapūrṇāntaḥ samasyate uttarapadasya ca lopaḥ vaktavyaḥ . \\
\noindent \textbf{(P\_2,1.34-35)} KA\_I,386.18-388.4 Ro\_II,595-597 {17/37}     aśvānām yuktaḥ aśvayuktaḥ aśvayuktaḥ rathaḥ aśvarathaḥ . \\
\noindent \textbf{(P\_2,1.34-35)} KA\_I,386.18-388.4 Ro\_II,595-597 {18/37}     dadhnaḥ pūrṇaḥ dadipūrṇaḥ dadhipūrṇaḥ ghaṭaḥ dadhighaṭaḥ . \\
\noindent \textbf{(P\_2,1.34-35)} KA\_I,386.18-388.4 Ro\_II,595-597 {19/37}     tat tarhi bahu vaktavyam . \\
\noindent \textbf{(P\_2,1.34-35)} KA\_I,386.18-388.4 Ro\_II,595-597 {20/37}     na vā asamāse adarśanāt . \\
\noindent \textbf{(P\_2,1.34-35)} KA\_I,386.18-388.4 Ro\_II,595-597 {21/37}     na vā vaktavyam . \\
\noindent \textbf{(P\_2,1.34-35)} KA\_I,386.18-388.4 Ro\_II,595-597 {22/37}     kim kāraṇam . \\
\noindent \textbf{(P\_2,1.34-35)} KA\_I,386.18-388.4 Ro\_II,595-597 {23/37}     asamāse adarśanāt . \\
\noindent \textbf{(P\_2,1.34-35)} KA\_I,386.18-388.4 Ro\_II,595-597 {24/37}     yat hi asamāse dṛśyate samāse ca na dṛśyate tat lopārambham prayojayati . \\
\noindent \textbf{(P\_2,1.34-35)} KA\_I,386.18-388.4 Ro\_II,595-597 {25/37}     na ca asamāse upasiktaśabdaḥ saṃsṛṣṭaśabdaḥ pūrṇaśabdaḥ vā dṛśyate . \\
\noindent \textbf{(P\_2,1.34-35)} KA\_I,386.18-388.4 Ro\_II,595-597 {26/37}     katham tarhi sāmarthyam gamyate . \\
\noindent \textbf{(P\_2,1.34-35)} KA\_I,386.18-388.4 Ro\_II,595-597 {27/37}     yuktārthasampratyayāt ca sāmarthyam . \\
\noindent \textbf{(P\_2,1.34-35)} KA\_I,386.18-388.4 Ro\_II,595-597 {28/37}     dadhnā yuktārthatā sampratīyate . \\
\noindent \textbf{(P\_2,1.34-35)} KA\_I,386.18-388.4 Ro\_II,595-597 {29/37}     katham punaḥ jñāyate dadhnā yuktārthatā sampratīyate iti . \\
\noindent \textbf{(P\_2,1.34-35)} KA\_I,386.18-388.4 Ro\_II,595-597 {30/37}     sampratyayāt ca tadarthādhyavasānam . \\
\noindent \textbf{(P\_2,1.34-35)} KA\_I,386.18-388.4 Ro\_II,595-597 {31/37}     sampratyayāt ca tadarthaḥ adhyavasīyate . \\
\noindent \textbf{(P\_2,1.34-35)} KA\_I,386.18-388.4 Ro\_II,595-597 {32/37}     avaśyam ca etat evam vijñeyam . \\
\noindent \textbf{(P\_2,1.34-35)} KA\_I,386.18-388.4 Ro\_II,595-597 {33/37}     sampratīyamānārthalope hi anavasthā .yaḥ hi manyate sampratīyamānārthānām śabdānām lopaḥ bhavati iti anavasthā tasya lopasya syāt . \\
\noindent \textbf{(P\_2,1.34-35)} KA\_I,386.18-388.4 Ro\_II,595-597 {34/37}     dadhi iti ukte bahavaḥ arthāḥ gamyante mandakam uttarakam nilīnakam iti tadvācinām śabdānām lopaḥ vaktavyaḥ syāt . \\
\noindent \textbf{(P\_2,1.34-35)} KA\_I,386.18-388.4 Ro\_II,595-597 {35/37}     tathā guḍaḥ iti ukte madhuraśabdasya śṛṅgaveram iti ukte ca kaṭuśabdasya . \\
\noindent \textbf{(P\_2,1.34-35)} KA\_I,386.18-388.4 Ro\_II,595-597 {36/37}     antareṇa khalu api śabdaprayogam bahavaḥ arthāḥ gamyante akṣinikocaiḥ pāṇivihāraiḥ ca . \\
\noindent \textbf{(P\_2,1.34-35)} KA\_I,386.18-388.4 Ro\_II,595-597 {37/37}     tadvācinām śabdānām lopaḥ vaktavyaḥ syāt . \\
\noindent \textbf{(P\_2,1.36)} kim caturthyantasya tadarthamātreṇa samāsaḥ bhavati . \\
\noindent \textbf{(P\_2,1.36)} evam bhavitum arhati . \\
\noindent \textbf{(P\_2,1.36)} caturthī tadarthamātreṇa cet sarvaprasaṅgaḥ aviśeṣāt . \\
\noindent \textbf{(P\_2,1.36)} caturthī tadarthamātreṇa cet sarvaprasaṅgaḥ sarvasya caturthyantasya tadarthamātreṇa saha samāsaḥ prāpnoti . \\
\noindent \textbf{(P\_2,1.36)} anena api prāpnoti . \\
\noindent \textbf{(P\_2,1.36)} randhanāya sthālī avahananāya ulūkhalam iti . \\
\noindent \textbf{(P\_2,1.36)} kim kāraṇam . \\
\noindent \textbf{(P\_2,1.36)} aviśeṣāt . \\
\noindent \textbf{(P\_2,1.36)} na hi kaḥ cit viśeṣaḥ upādīyate evañjātīyakasya caturthyantasya tadarthena saha samāsaḥ bhavati iti . \\
\noindent \textbf{(P\_2,1.36)} anupādīyamane viśeṣe sarvaprasaṅgaḥ . \\
\noindent \textbf{(P\_2,1.36)} balirakṣitābhyām ca anarthakam vacanam . \\
\noindent \textbf{(P\_2,1.36)} balirakṣitābhyām ca samāsavacanam anarthakam . \\
\noindent \textbf{(P\_2,1.36)} yaḥ hi mahārājāya baliḥ mahārājārthaḥ saḥ bhavati . \\
\noindent \textbf{(P\_2,1.36)} tatra tadarthaḥ iti eva siddham . \\
\noindent \textbf{(P\_2,1.36)} yadi punaḥ vikṛtiḥ caturthyantā prakṛtyā saha samasyate iti etat lakṣaṇam kriyeta . \\
\noindent \textbf{(P\_2,1.36)} vikṛtiḥ prakṛtyā iti cet aśvaghāsādīnām upasaṅkhyānam . \\
\noindent \textbf{(P\_2,1.36)} vikṛtiḥ prakṛtyā iti cet aśvaghāsādīnām upasaṅkhyānam kartavyam . \\
\noindent \textbf{(P\_2,1.36)} aśvaghāsaḥ śvaśrūsuram hastividhā iti . \\
\noindent \textbf{(P\_2,1.36)} arthena nityasamāsavacanam . \\
\noindent \textbf{(P\_2,1.36)} arthśabdena nityasamāsaḥ vaktavyaḥ . \\
\noindent \textbf{(P\_2,1.36)} brāhmaṇārtham kṣatriyārtham . \\
\noindent \textbf{(P\_2,1.36)} kim vikṛtiḥ caturthyantā prakṛtyā saha samasyate iti ataḥ arthena nityasamāsaḥ vaktavyaḥ . \\
\noindent \textbf{(P\_2,1.36)} na iti āha sarvathā arthena nityasamāsaḥ vaktavyaḥ vigrahaḥ mā bhūt iti . \\
\noindent \textbf{(P\_2,1.36)} sarvaliṅgatā ca . \\
\noindent \textbf{(P\_2,1.36)} sarvaliṅgatā ca vaktavyā . \\
\noindent \textbf{(P\_2,1.36)} brāhmaṇārtham payaḥ brāhmaṇārthaḥ sūpaḥ brāhmaṇārthā yavāgūḥ iti . \\
\noindent \textbf{(P\_2,1.36)} kim arthena nityasamāsaḥ ucyate iti ataḥ sarvaliṅgatā vaktavyā . \\
\noindent \textbf{(P\_2,1.36)} na iti āha . \\
\noindent \textbf{(P\_2,1.36)} sarvathā sarvaliṅgatā vaktavyā . \\
\noindent \textbf{(P\_2,1.36)} kim kāraṇam . \\
\noindent \textbf{(P\_2,1.36)} arthaśabdaḥ ayam puṃliṅgaḥ uttarapadārthapradhānaḥ ca tatpuruṣaḥ . \\
\noindent \textbf{(P\_2,1.36)} tena puṃliṅgasya eva samāsasya abhidhānam syāt strīnapuṃsakaliṅgasya na syāt . \\
\noindent \textbf{(P\_2,1.36)} tat tarhi bahu vaktavyam . \\
\noindent \textbf{(P\_2,1.36)} vikṛtiḥ prakṛtyā iti vaktavyam . \\
\noindent \textbf{(P\_2,1.36)} aśvaghāsādīnām upasaṅkhyānam kartavyam . \\
\noindent \textbf{(P\_2,1.36)} arthena nityasamāsaḥ vaktavyaḥ . \\
\noindent \textbf{(P\_2,1.36)} sarvaliṅgatā ca vaktavyā . \\
\noindent \textbf{(P\_2,1.36)} na vaktavyam . \\
\noindent \textbf{(P\_2,1.36)} yat tāvat ucyate vikṛtiḥ prakṛtyā iti vaktavyam . \\
\noindent \textbf{(P\_2,1.36)} na vaktavyam . \\
\noindent \textbf{(P\_2,1.36)} ācāryapravṛttiḥ jñāyapati vikṛtiḥ caturthyantā prakṛtyā saha samasyate iti yat ayam balirakiṣitagrahaṇam karoti . \\
\noindent \textbf{(P\_2,1.36)} katham kṛtvā jñāpakam . \\
\noindent \textbf{(P\_2,1.36)} yathājātīyakānām samāse balirakṣitagrahaṇena arthaḥ tathājātīyakānām samāsaḥ . \\
\noindent \textbf{(P\_2,1.36)} yadi ca vikṛtiḥ caturthyantā prakṛtyā saha samasyate na tadarthamātreṇa tataḥ balirakiṣitagrahaṇam arthavat bhavati . \\
\noindent \textbf{(P\_2,1.36)} yat api ucyate aśvaghāsādīnām upasaṅkhyānam kartavyam iti . \\
\noindent \textbf{(P\_2,1.36)} na kartavyam . \\
\noindent \textbf{(P\_2,1.36)} aśvaghāsādayaḥ ṣaṣṭhīsamāsāḥ bhaviṣyanti . \\
\noindent \textbf{(P\_2,1.36)} yat hi yadartham bhavati ayam api tatra abhisambandhaḥ bhavati asya idam iti . \\
\noindent \textbf{(P\_2,1.36)} tat yathā guroḥ idam gurvartham iti . \\
\noindent \textbf{(P\_2,1.36)} nanu ca svarabhedaḥ bhavati . \\
\noindent \textbf{(P\_2,1.36)} caturthīsamāse sati pūrvapadaprakṛtisvaratvena bhavitavyam ṣaṣṭhīsamāse punaḥ antodāttatvena ṇa asti bhedaḥ . \\
\noindent \textbf{(P\_2,1.36)} caturthīsamāse api sati antodāttatvena eva bhavitavyam . \\
\noindent \textbf{(P\_2,1.36)} katham . \\
\noindent \textbf{(P\_2,1.36)} ācāryapravṛttiḥ jñāpayati vikṛtiḥ caturthyantā prakṛtisvarā bhavati na caturthīmātram iti yat ayam caturthī tadarthe arthe kte ca iti arthagrahaṇam ktagrahaṇam ca karoti . \\
\noindent \textbf{(P\_2,1.36)} katham kṛtvā jñāpakam . \\
\noindent \textbf{(P\_2,1.36)} yathājātīyakānām arthagrahaṇena ktagrahaṇena ca arthaḥ tathājātīyakānām prakṛtisvaratvam . \\
\noindent \textbf{(P\_2,1.36)} yadi ca vikṛtiḥ caturthyantā prakṛtyā bhavati na caturthīmātram tataḥ arthagrahaṇam ktagrahaṇam ca arthavat bhavati .yat api ucyate arthena nityasamāsaḥ vaktavyaḥ iti . \\
\noindent \textbf{(P\_2,1.36)} na vaktavyaḥ . \\
\noindent \textbf{(P\_2,1.36)} sarthappratyayaḥ kariṣyate . \\
\noindent \textbf{(P\_2,1.36)} kim kṛtam bhavati . \\
\noindent \textbf{(P\_2,1.36)} na ca eva hi kadā cit vigrahaḥ bhavati . \\
\noindent \textbf{(P\_2,1.36)} api ca sarvaliṅgatā siddhā bhavati . \\
\noindent \textbf{(P\_2,1.36)} yadi sarthappratayaḥ kriyate itsañjñā na prāpnoti . \\
\noindent \textbf{(P\_2,1.36)} atha api katham cit itsañjñā syāt evam api śryartham bhvartham iti aṅgasya iti iyaṅuvaṅau syātām . \\
\noindent \textbf{(P\_2,1.36)} evam tarhi bahuvrīhiḥ bhavaiṣyati . \\
\noindent \textbf{(P\_2,1.36)} kim kṛtam bhavati . \\
\noindent \textbf{(P\_2,1.36)} bhavati vai kaḥ cit asvapadavigrahaḥ bahuvrīhiḥ . \\
\noindent \textbf{(P\_2,1.36)} tat yathā śobhanam mukham asyāḥ sumukhī iti . \\
\noindent \textbf{(P\_2,1.36)} na evam śakyam . \\
\noindent \textbf{(P\_2,1.36)} iha hi mahadartham iti āttvakapau prasajyetām . \\
\noindent \textbf{(P\_2,1.36)} evam tarhi tadarthasya uttarapadasya arthaśabdaḥ ādeśaḥ kariṣyate . \\
\noindent \textbf{(P\_2,1.36)} kim kṛtam bhavati . \\
\noindent \textbf{(P\_2,1.36)} na ca eva kadā cit ādeśena vigrahaḥ bhavati . \\
\noindent \textbf{(P\_2,1.36)} api ca sarvaliṅgatā siddhā bhavati . \\
\noindent \textbf{(P\_2,1.36)} tat tarhi vaktavyam . \\
\noindent \textbf{(P\_2,1.36)} na vaktavyam . \\
\noindent \textbf{(P\_2,1.36)} yogavibhāgaḥ kariṣyate . \\
\noindent \textbf{(P\_2,1.36)} caturthī subantena saha samasyate . \\
\noindent \textbf{(P\_2,1.36)} tataḥ tadarthārtha . \\
\noindent \textbf{(P\_2,1.36)} tadarthasya uttarapadasya arthaśabdaḥ ādeśaḥ bhavati . \\
\noindent \textbf{(P\_2,1.36)} iha api tarhi samāsaḥ prāpnoti chātrāya rucitam chātrāya svaditam iti . \\
\noindent \textbf{(P\_2,1.36)} ācāryapravṛttiḥ jñāpayati tādarthye ya caturthī sā samasyate na caturthīmātram iti yat ayam hitasukhagrahaṇam karoti . \\
\noindent \textbf{(P\_2,1.36)} katham kṛtvā jñāpakam . \\
\noindent \textbf{(P\_2,1.36)} yathājātīyakānām samāse hitasukhagrahaṇena arthaḥ tathājātīyakānam samāsaḥ . \\
\noindent \textbf{(P\_2,1.36)} yadi ca tādarthye yā caturthī sā samasyate na caturthīmātram tataḥ hitasukhagrahaṇam arthavat bhavati . \\
\noindent \textbf{(P\_2,1.36)} iha api tarhi tadarthasya uttarapadasya arthaśabdaḥ ādeśaḥ prāpnoti . \\
\noindent \textbf{(P\_2,1.36)} yūpāya dāru yūpadāru rathadāru . \\
\noindent \textbf{(P\_2,1.36)} vāvacanam vidhāsyate . \\
\noindent \textbf{(P\_2,1.36)} iha api tarhi vibhāṣā prāprnoti . \\
\noindent \textbf{(P\_2,1.36)} brāhmaṇāṛtham kṣatriyārtham iti . \\
\noindent \textbf{(P\_2,1.36)} evam tarhi ācāryapravṛttiḥ jñāpayati prakṛtivikṛtyoḥ yaḥ samāsaḥ tatra tadarthasya uttarapadasya arthaśabdaḥ ādeśaḥ bhavati anyatra nityaḥ iti yat ayam balihitagrahaṇam karoti . \\
\noindent \textbf{(P\_2,1.36)} evam tarhi udakārthaḥ vīvadhaḥ . \\
\noindent \textbf{(P\_2,1.36)} sthānivadbhāvāt udabhāvaḥ prāpnoti . \\
\noindent \textbf{(P\_2,1.36)} tasmāt na evam śakyam . \\
\noindent \textbf{(P\_2,1.36)} na cet evam arthena nityasamāsaḥ vaktavyaḥ sarvaliṅgatā ca . \\
\noindent \textbf{(P\_2,1.36)} na eṣaḥ doṣaḥ . \\
\noindent \textbf{(P\_2,1.36)} idam tāvat ayam praṣṭavyaḥ . \\
\noindent \textbf{(P\_2,1.36)} atha iha brāhmaṇebhyaḥ iti kā eṣā caturthī . \\
\noindent \textbf{(P\_2,1.36)} tādarthye iti āha . \\
\noindent \textbf{(P\_2,1.36)} yadi tādarthye caturthī arthaśabdasya prayogeṇa na bhavitavyam uktārthānām aprayogaḥ iti . \\
\noindent \textbf{(P\_2,1.36)} samāsaḥ api tarhi na prāpnoti . \\
\noindent \textbf{(P\_2,1.36)} vacanāt samāsaḥ bhaviṣyati . \\
\noindent \textbf{(P\_2,1.36)} yat api ucyate sarvaliṅgatā ca vaktavyā iti . \\
\noindent \textbf{(P\_2,1.36)} na vaktavyā . \\
\noindent \textbf{(P\_2,1.36)} liṅgam aśiṣyam lokāśrayatvāt liṅgasya . \\
\noindent \textbf{(P\_2,1.37)} atyalpam idam ucyate bhayena iti . \\
\noindent \textbf{(P\_2,1.37)} bhayabhītabhītibhībhiḥ iti vaktavyam . \\
\noindent \textbf{(P\_2,1.37)} vṛkāt bhayam vṛkabhayam vṛkāt bhītaḥ vṛkabhītaḥ vṛkāt bhītiḥ vṛkabhītiḥ vṛkāt bhīḥ vṛkabhīḥ iti . \\
\noindent \textbf{(P\_2,1.37)} aparaḥ āha : bhayanirgatajugupsubhiḥ iti vaktavyam : vṛkabhayam grāmanirgataḥ adharmajugupsuḥ iti . \\
\noindent \textbf{(P\_2,1.40)} śauṇḍādibhiḥ iti vaktavyam . \\
\noindent \textbf{(P\_2,1.40)} iha api yathā syāt . \\
\noindent \textbf{(P\_2,1.40)} akṣadhūrtaḥ strīdhūrtaḥ akṣakitavaḥ strīkitavaḥ iti . \\
\noindent \textbf{(P\_2,1.40)} tat tarhi vaktavyam . \\
\noindent \textbf{(P\_2,1.40)} na vaktavyam . \\
\noindent \textbf{(P\_2,1.40)} bahuvacananirdeśāt śauṇḍādibhiḥ iti vijñāsyate . \\
\noindent \textbf{(P\_2,1.42)} dhvāṅkṣeṇa iti arthagrahaṇam . \\
\noindent \textbf{(P\_2,1.42)} dhvāṅkṣeṇa iti arthagrahaṇam kartavyam . \\
\noindent \textbf{(P\_2,1.42)} iha api yathā syāt . \\
\noindent \textbf{(P\_2,1.42)} tīrthakākaḥ iti . \\
\noindent \textbf{(P\_2,1.42)} kṣepe iti ucyate . \\
\noindent \textbf{(P\_2,1.42)} kaḥ iha kṣepaḥ nāma . \\
\noindent \textbf{(P\_2,1.42)} yathā tīrthe kākāḥ na ciram sthātāraḥ bhavanti evam yaḥ gurukulāni gatvā na ciram tiṣṭhati sa ucyate tīrthakākaḥ iti . \\
\noindent \textbf{(P\_2,1.43)} kṛtyaiḥ niyoge yadgrahaṇam . \\
\noindent \textbf{(P\_2,1.43)} kṛtyaiḥ niyoge iti vaktavyam . \\
\noindent \textbf{(P\_2,1.43)} iha api yathā syāt . \\
\noindent \textbf{(P\_2,1.43)} pūrvāḥṇegeyam sāma prātaḥ adhyeyaḥ anuvākaḥ iti . \\
\noindent \textbf{(P\_2,1.43)} tat tarhi vaktavyam . \\
\noindent \textbf{(P\_2,1.43)} na vaktavyam . \\
\noindent \textbf{(P\_2,1.43)} ṛṇe iti eva siddham . \\
\noindent \textbf{(P\_2,1.43)} iha yat yasya niyogataḥ kāryam ṛṇam tasya tat bhavati . \\
\noindent \textbf{(P\_2,1.43)} tataḥ ṛṇe iti eva siddham . \\
\noindent \textbf{(P\_2,1.43)} yagrahaṇam ca kartavyam . \\
\noindent \textbf{(P\_2,1.43)} iha mā bhūt . \\
\noindent \textbf{(P\_2,1.43)} pūrvāhṇe dātavyā bhikṣā iti . \\
\noindent \textbf{(P\_2,1.47)} kim udāharaṇam . \\
\noindent \textbf{(P\_2,1.47)} avataptenakulasthitam te etat . \\
\noindent \textbf{(P\_2,1.47)} kṣepe iti ucyate . \\
\noindent \textbf{(P\_2,1.47)} kaḥ iha kṣepaḥ nāma . \\
\noindent \textbf{(P\_2,1.47)} yathā avatapte nakulāḥ na ciram sthātāraḥ bhavanti evam kāryāṇi ārabhya yaḥ na ciram tiṣṭhati sa ucyate avataptenakulasthitam te etat iti . \\
\noindent \textbf{(P\_2,1.47)} kṣepe saptamyantam ktāntena saha samasyate iti ucyate . \\
\noindent \textbf{(P\_2,1.47)} tatra sagatikena sanakulena ca samāsaḥ na prāpnoti . \\
\noindent \textbf{(P\_2,1.47)} kṣepe gatikārakapūrve uktam . kim uktam . \\
\noindent \textbf{(P\_2,1.47)} kṛdgrahaṇe gatikārakapūrvasya api iti . \\
\noindent \textbf{(P\_2,1.48)} kimarthaḥ cakāraḥ . \\
\noindent \textbf{(P\_2,1.48)} evakārārthaḥ . \\
\noindent \textbf{(P\_2,1.48)} pātresamitādayaḥ eva . \\
\noindent \textbf{(P\_2,1.48)} kva mā bhūt . \\
\noindent \textbf{(P\_2,1.48)} paramam pātresamitāḥ iti . \\
\noindent \textbf{(P\_2,1.49)} iha kasmāt avyayībhāvaḥ na bhavati . \\
\noindent \textbf{(P\_2,1.49)} ekā nadī ekanadī . \\
\noindent \textbf{(P\_2,1.49)} nadībhiḥ saṅkhyā iti prāpnoti . \\
\noindent \textbf{(P\_2,1.49)} iha kaḥ cit samāsaḥ pūrvapadārthapradhānaḥ , kaḥ cit uttarapadārthapradhānaḥ , kaḥ cit anyapadārthapradhānaḥ , kaḥ cit ubhayapadārthapradhānaḥ . \\
\noindent \textbf{(P\_2,1.49)} pūrvapadārthapradhānaḥ avyayībhāvaḥ , uttarapadārthapradhānaḥ tatpuruṣaḥ , anyapadārthapradhānaḥ bahuvrīhiḥ , ubhayapadārthapradhānaḥ dvandvaḥ. na ca atra pūrvapadārthaprādhānyam gamyate . \\
\noindent \textbf{(P\_2,1.49)} athavā avyayībhāvaḥ kriyatām bahuvrīhiḥ iti . \\
\noindent \textbf{(P\_2,1.49)} bahuvrīhiḥ bhaviṣyati vipratiṣedhena . \\
\noindent \textbf{(P\_2,1.49)} bhavet ekasañjñādhikāre siddham paraṅkāryatve tu na sidhyati . \\
\noindent \textbf{(P\_2,1.49)} ārambhasāmarthyāt ca avyayībhāvaḥ prāpnoti paraṅkāryatvāt ca bahuvrīhiḥ . \\
\noindent \textbf{(P\_2,1.49)} paraṅkāryatve ca na doṣaḥ . \\
\noindent \textbf{(P\_2,1.49)} nadībhiḥ saṅkhyāyāḥ samāhāre avyayībhāvaḥ vaktavyaḥ . \\
\noindent \textbf{(P\_2,1.49)} saḥ ca avaśyam vaktavyaḥ . \\
\noindent \textbf{(P\_2,1.49)} sarvam ekanadītare . \\
\noindent \textbf{(P\_2,1.51.1)} samāhāraḥ iti kaḥ ayam śabdaḥ . \\
\noindent \textbf{(P\_2,1.51.1)} samāṅpūrvāt harateḥ sarmasādhane ghañ . \\
\noindent \textbf{(P\_2,1.51.1)} samāhriyate samāhāraḥ iti . \\
\noindent \textbf{(P\_2,1.51.1)} yadi karmasādhanaḥ pañca kumāryaḥ samahṛtāḥ pañcakumāri daśakumāri gostriyoḥ upasarjanasya iti hrasvatvam na prāpnoti dviguḥ ekavacanam iti etat ca vaktavyam . \\
\noindent \textbf{(P\_2,1.51.1)} evam tarhi bhāvasādhanaḥ bhaviṣyati . \\
\noindent \textbf{(P\_2,1.51.1)} samāharaṇam samāhāraḥ . \\
\noindent \textbf{(P\_2,1.51.1)} atha bhāvasādhane sati kim abhidhīyate . \\
\noindent \textbf{(P\_2,1.51.1)} yat tat auttarādharyam . \\
\noindent \textbf{(P\_2,1.51.1)} kaḥ punaḥ gavām samāhāraḥ . \\
\noindent \textbf{(P\_2,1.51.1)} yat tat arjanam krayaṇam bhiṣaṇam aparaharaṇam vā . \\
\noindent \textbf{(P\_2,1.51.1)} yadi evam vikṣipteṣu pūleṣu goṣu carantīṣu na sidhyati . \\
\noindent \textbf{(P\_2,1.51.1)} evam tarhi samabhyāśīkaraṇam samāhāraḥ . \\
\noindent \textbf{(P\_2,1.51.1)} evam api pañcagrāmī ṣaṇṇagarī tripurī iti na sidhyati . \\
\noindent \textbf{(P\_2,1.51.1)} kim kāraṇam . \\
\noindent \textbf{(P\_2,1.51.1)} sam ekatvavācī āṅ ābhimukhye vartate haratiḥ deśāntaraprāpaṇe . \\
\noindent \textbf{(P\_2,1.51.1)} na avaśyam haratiḥ deśāntaraprāpaṇe eva vartate . \\
\noindent \textbf{(P\_2,1.51.1)} kim tarhi . \\
\noindent \textbf{(P\_2,1.51.1)} sādṛśye api vartate . \\
\noindent \textbf{(P\_2,1.51.1)} tat yathā mātuḥ anuharati pituḥ anuharati . \\
\noindent \textbf{(P\_2,1.51.1)} atha vā pañcagrāmī ṣaṇṇagarī tripurī iti na eva idam iyati eva avatiṣṭhate . \\
\noindent \textbf{(P\_2,1.51.1)} avaśyam asau tataḥ kim cit ākāṅkṣati kriyām vā guṇam vā . \\
\noindent \textbf{(P\_2,1.51.1)} yat ākāṅkṣata tat ekam sa ca samāhāraḥ . \\
\noindent \textbf{(P\_2,1.51.1)} ayam tarhi bhāvasādhane sati doṣaḥ . \\
\noindent \textbf{(P\_2,1.51.1)} pañcapūlī ānīyatām iti bhāvānayane codite dravyānayanam na prāpoti . \\
\noindent \textbf{(P\_2,1.51.1)} na eṣaḥ doṣaḥ . \\
\noindent \textbf{(P\_2,1.51.1)} iha tāvat ayam praṣṭavyaḥ . \\
\noindent \textbf{(P\_2,1.51.1)} atha iha gauḥ anubandhyaḥ ajaḥ agnīṣomīyaḥ iti katham ākṛtau coditāyām dravye ārambhaṇālambhanaprokṣaṇaviśasanāni kriyante . \\
\noindent \textbf{(P\_2,1.51.1)} asambhavāt . \\
\noindent \textbf{(P\_2,1.51.1)} ākṛtau ārambhaṇādīnām sambhavaḥ na asti iti kṛtvā ākṛtisahacarite dravye ārambhaṇādīni kriyante . \\
\noindent \textbf{(P\_2,1.51.1)} idam api evañjātīyakam eva . \\
\noindent \textbf{(P\_2,1.51.1)} asambhavāt bhāvānayanasya dravyānayanam bhaviṣyati . \\
\noindent \textbf{(P\_2,1.51.1)} atha vā avyatirekāt dravyākṛtyoḥ . \\
\noindent \textbf{(P\_2,1.51.2)} kim punaḥ dvigusañjñā pratyayottarapadayoḥ bhavati . \\
\noindent \textbf{(P\_2,1.51.2)} evam bhavitum arhati . \\
\noindent \textbf{(P\_2,1.51.2)} dvigusañjñā pratyayottarapadayoḥ cet itaretarāśrayatvāt aprasiddhiḥ . \\
\noindent \textbf{(P\_2,1.51.2)} dvigusañjñā pratyayottarapadayoḥ cet itaretarāśrayatvāt aprasiddhiḥ . \\
\noindent \textbf{(P\_2,1.51.2)} kā etaretarāśrayatā . \\
\noindent \textbf{(P\_2,1.51.2)} dvigunimitte pratyayottarapade pratyayottarapadanimittā ca dvigusañjñā . \\
\noindent \textbf{(P\_2,1.51.2)} tat etat itaretarāśrayam . \\
\noindent \textbf{(P\_2,1.51.2)} itaretarāśrayāṇi ca na prakalpante . \\
\noindent \textbf{(P\_2,1.51.2)} evam tarhi arthe it vakṣyāmi . \\
\noindent \textbf{(P\_2,1.51.2)} arthe cet taddhitānutpattiḥ bahuvrīhivat . \\
\noindent \textbf{(P\_2,1.51.2)} arthe cet taddhitotpattiḥ na prāpnoti . \\
\noindent \textbf{(P\_2,1.51.2)} pāñcanāpitiḥ , dvimāturaḥ , traimāturaḥ . \\
\noindent \textbf{(P\_2,1.51.2)} kim kāraṇam . \\
\noindent \textbf{(P\_2,1.51.2)} dvigunā uktatvāt bahuvrīhivat . \\
\noindent \textbf{(P\_2,1.51.2)} tat yathā citraguḥ iti bahuvrīhiṇoktatvāt matvarthasya matvarthīyaḥ na bhavati . \\
\noindent \textbf{(P\_2,1.51.2)} evam tarhi samāsataddhitavidhau iti vakṣyāmi . \\
\noindent \textbf{(P\_2,1.51.2)} samāsataddhitavidhau iti cet anyatra samāsasañjñābhāvaḥ . \\
\noindent \textbf{(P\_2,1.51.2)} samāsataddhitavidhau iti cet anyatra samāsasañjñā na prāpnoti . \\
\noindent \textbf{(P\_2,1.51.2)} kva anyatra . \\
\noindent \textbf{(P\_2,1.51.2)} svare . \\
\noindent \textbf{(P\_2,1.51.2)} pañcāratniḥ , daśāratniḥ . \\
\noindent \textbf{(P\_2,1.51.2)} igante dvigau iti eṣaḥ svaraḥ na prāpnoti . \\
\noindent \textbf{(P\_2,1.51.2)} siddham tu pratyayottarapadayoḥ ca iti vacanāt . \\
\noindent \textbf{(P\_2,1.51.2)} siddham etat . \\
\noindent \textbf{(P\_2,1.51.2)} katham . \\
\noindent \textbf{(P\_2,1.51.2)} pratyayottarapadayoḥ ca iti vacanāt . \\
\noindent \textbf{(P\_2,1.51.2)} pratyayottarapadayoḥ dvigusañjñā bhavati iti vaktavyam . \\
\noindent \textbf{(P\_2,1.51.2)} nanu ca uktam dvigusañjñā pratyayottarapadayoḥ cet itaretarāśrayatvāt aprasiddhiḥ iti . \\
\noindent \textbf{(P\_2,1.51.2)} na eṣaḥ doṣaḥ . \\
\noindent \textbf{(P\_2,1.51.2)} itaretarāśrayamātram etat coditam sarvāṇi ca itaretarāśrayāṇi ekatvena parihṛtāni siddham tu nityaśabdatvāt iti . \\
\noindent \textbf{(P\_2,1.51.2)} na idam tulyam anyaiḥ itaretarāśrayaiḥ . \\
\noindent \textbf{(P\_2,1.51.2)} na hi sañjñā nityā . \\
\noindent \textbf{(P\_2,1.51.2)} evam tarhi bhāvinī sañjñā vijñāsyate . \\
\noindent \textbf{(P\_2,1.51.2)} tat yathā : kaḥ cit kam cit tantuvāyam āha : asya sūtrasya śāṭakam vaya iti . \\
\noindent \textbf{(P\_2,1.51.2)} saḥ paśyati . \\
\noindent \textbf{(P\_2,1.51.2)} yadi śāṭakaḥ na vātavyaḥ atha vātavyaḥ na śāṭakaḥ . \\
\noindent \textbf{(P\_2,1.51.2)} śāṭakaḥ vātavyaḥ ca iti vipratiṣiddham . \\
\noindent \textbf{(P\_2,1.51.2)} bhāvinī khalu asya sañjñā abhipretā . \\
\noindent \textbf{(P\_2,1.51.2)} saḥ manye vātavyaḥ yasmin ute śāṭakaḥ iti etat bhavati iti . \\
\noindent \textbf{(P\_2,1.51.2)} evam iha api tasmin dviguḥ bhavati yasya abhinirvṛttasya pratyaya uttarapadam iti ca ete sañjñe bhaviṣyataḥ . \\
\noindent \textbf{(P\_2,1.51.2)} atha vā punaḥ astu arthe iti . \\
\noindent \textbf{(P\_2,1.51.2)} nanu ca uktam arthe cet taddhitānutpattiḥ bahuvrīhivat iti . \\
\noindent \textbf{(P\_2,1.51.2)} na eṣaḥ doṣaḥ . \\
\noindent \textbf{(P\_2,1.51.2)} na avaśyam arthaśabdaḥ abhidheye eva vartate . \\
\noindent \textbf{(P\_2,1.51.2)} kim tarhi . \\
\noindent \textbf{(P\_2,1.51.2)} syādarthe api vartate . \\
\noindent \textbf{(P\_2,1.51.2)} tat yathā . \\
\noindent \textbf{(P\_2,1.51.2)} dārārtham ghaṭāmahe . \\
\noindent \textbf{(P\_2,1.51.2)} dhanārtham bhikṣāmahe . \\
\noindent \textbf{(P\_2,1.51.2)} dārāḥ naḥ syuḥ . \\
\noindent \textbf{(P\_2,1.51.2)} dhanāni naḥ syuḥ iti . \\
\noindent \textbf{(P\_2,1.51.2)} evam iha api taddhitārthe dviguḥ bhavati taddhitaḥ syāt iti . \\
\noindent \textbf{(P\_2,1.51.2)} dvigoḥ vā lugvacanam jñāpakam taddhitotpatteḥ . \\
\noindent \textbf{(P\_2,1.51.2)} atha vā yat ayam dvigoḥ luk anapatye iti dvigoḥ uttarasya taddhitasya lukam śāsti tat jñāpayati ācāryaḥ utpadyate dvigoḥ taddhitaḥ iti. \\
\noindent \textbf{(P\_2,1.51.3)} samāhārasamūhayoḥ aviśeṣāt samāhāragrahaṇānarthakyam taddhitārthena kṛtatvāt . \\
\noindent \textbf{(P\_2,1.51.3)} samāhāraḥ samūhaḥ iti aviśiṣtau etau arthau . \\
\noindent \textbf{(P\_2,1.51.3)} samāhārasamūhayoḥ aviśeṣāt samāhāragrahaṇam anarthakam . \\
\noindent \textbf{(P\_2,1.51.3)} kim kāraṇam . \\
\noindent \textbf{(P\_2,1.51.3)} taddhitārthe kṛtatvāt . \\
\noindent \textbf{(P\_2,1.51.3)} taddhitārthe dviguḥ iti evam atra dviguḥ bhaviṣyati . \\
\noindent \textbf{(P\_2,1.51.3)} yadi taddhitārthe dviguḥ iti evam atra dviguḥ bhavati taddhitotpattiḥ prāpnoti . \\
\noindent \textbf{(P\_2,1.51.3)} utpadyatām . \\
\noindent \textbf{(P\_2,1.51.3)} luk bhaviṣyati . \\
\noindent \textbf{(P\_2,1.51.3)} lukkṛtāni prāpnuvanti . \\
\noindent \textbf{(P\_2,1.51.3)} kāni . \\
\noindent \textbf{(P\_2,1.51.3)} pañcapūlī daśapūlī . \\
\noindent \textbf{(P\_2,1.51.3)} aparimāṇabistācitakambalebhyaḥ na taddhitaluki iti pratiṣedhaḥ prāpnoti . \\
\noindent \textbf{(P\_2,1.51.3)} pañcagavam daśagavam . \\
\noindent \textbf{(P\_2,1.51.3)} goḥ ataddhitaluki it ṭac na prapnoti . \\
\noindent \textbf{(P\_2,1.51.3)} na eṣaḥ doṣaḥ . \\
\noindent \textbf{(P\_2,1.51.3)} aviśeṣeṇa dvigoḥ ṅīp bhavati iti uktvā samāhāre iti vakṣyāmi . \\
\noindent \textbf{(P\_2,1.51.3)} tat niyamārtham bhaviṣyati . \\
\noindent \textbf{(P\_2,1.51.3)} samāhāre eva na anyatra iti . \\
\noindent \textbf{(P\_2,1.51.3)} goḥ akāraḥ dvigoḥ samāhāre . \\
\noindent \textbf{(P\_2,1.51.3)} aviśeṣeṇa goḥ ṭac bhavati iti uktvā dvigoḥ samāhāre iti vakṣyāmi . \\
\noindent \textbf{(P\_2,1.51.3)} tat niyamārtham bhaviṣyati . \\
\noindent \textbf{(P\_2,1.51.3)} samāhāre eva na anyatra iti . \\
\noindent \textbf{(P\_2,1.51.3)} abhidhānārtham tu . \\
\noindent \textbf{(P\_2,1.51.3)} abhidhānārtham tu samāhāragrahaṇam kartavyam . \\
\noindent \textbf{(P\_2,1.51.3)} samāhāreṇa abhidhānam yathā syāt taddhitārthena mā bhūt iti . \\
\noindent \textbf{(P\_2,1.51.3)} kim ca syāt . \\
\noindent \textbf{(P\_2,1.51.3)} taddhitotpattiḥ prasajyeta . \\
\noindent \textbf{(P\_2,1.51.3)} utpadyatām . \\
\noindent \textbf{(P\_2,1.51.3)} luk bhaviṣyati . \\
\noindent \textbf{(P\_2,1.51.3)} lukkṛtāni prāpnuvanti . \\
\noindent \textbf{(P\_2,1.51.3)} sarvāṇi parihṛtāni . \\
\noindent \textbf{(P\_2,1.51.3)} na sarvāṇi parihṛtāni . \\
\noindent \textbf{(P\_2,1.51.3)} pañcakumāri daśakumāri . \\
\noindent \textbf{(P\_2,1.51.3)} lik taddhitaluki iti ṅīpaḥ luk prasajyeta . \\
\noindent \textbf{(P\_2,1.51.3)} dvandvatatpuruṣayoḥ uttarapade nityasamāsavacanam . \\
\noindent \textbf{(P\_2,1.51.3)} dvandvatatpuruṣayoḥ uttarapade nityasamāsaḥ vaktavyaḥ . \\
\noindent \textbf{(P\_2,1.51.3)} vāgdṛṣadapriyaḥ chatropānahapriyaḥ pañcagavapriyaḥ daśagavapriyaḥ . \\
\noindent \textbf{(P\_2,1.51.3)} kim prayojanam . \\
\noindent \textbf{(P\_2,1.51.3)} samudāyavṛttau avayavānām mā kadā cit avṛttiḥ bhūt iti . \\
\noindent \textbf{(P\_2,1.51.3)} tat tarhi vaktavyam . \\
\noindent \textbf{(P\_2,1.51.3)} na vaktavyam . \\
\noindent \textbf{(P\_2,1.51.3)} iha dvau pakṣau vṛttipakṣaḥ avṛttipakṣaḥ ca . \\
\noindent \textbf{(P\_2,1.51.3)} yadā vṛttipakṣaḥ tadā sarveṣām eva vṛttiḥ . \\
\noindent \textbf{(P\_2,1.51.3)} yadā tu avṛttiḥ tadā sarveṣām avṛttiḥ . \\
\noindent \textbf{(P\_2,1.51.3)} uttarapadena parimāṇina dvigoḥ samāsavacanam . \\
\noindent \textbf{(P\_2,1.51.3)} uttarapadena parimāṇina dvigoḥ samāsaḥ vaktavyaḥ . \\
\noindent \textbf{(P\_2,1.51.3)} dvimāsajātaḥ trimāsajātaḥ . \\
\noindent \textbf{(P\_2,1.51.3)} kim punaḥ kāraṇam na sidhyati . \\
\noindent \textbf{(P\_2,1.51.3)} sup supā iti vartate . \\
\noindent \textbf{(P\_2,1.51.3)} evam tarhi idam syāt : dvau māsau dvimāsam , dvimāsam jātasya iti . \\
\noindent \textbf{(P\_2,1.51.3)} na evam śakyam . \\
\noindent \textbf{(P\_2,1.51.3)} svare hi doṣaḥ syāt . \\
\noindent \textbf{(P\_2,1.51.3)} dvimāsajātaḥ iti prāpnoti dvimāsajātaḥ iti ca iṣyate . \\
\noindent \textbf{(P\_2,1.51.3)} dvyāhnajātaḥ ca na sidhyati . \\
\noindent \textbf{(P\_2,1.51.3)} dvyahajāta iti prāpnoti na ca evam bhavitavyam . \\
\noindent \textbf{(P\_2,1.51.3)} bhavitavyam ca yadā samāhāre dviguḥ . \\
\noindent \textbf{(P\_2,1.51.3)} dvyahnajātaḥ tu na sidhyati . \\
\noindent \textbf{(P\_2,1.51.3)} kim ucyate parimāṇinā iti na punaḥ anyatra api . \\
\noindent \textbf{(P\_2,1.51.3)} pañcagavapriyaḥ daśagavapriyaḥ . \\
\noindent \textbf{(P\_2,1.51.3)} anyatra samudāyabahuvrīhitvāt uttarapadaprasiddhiḥ . \\
\noindent \textbf{(P\_2,1.51.3)} anyatra samudāyabaḥ huvrīhisañjñḥ . \\
\noindent \textbf{(P\_2,1.51.3)} anyatra samudāyabahuvrīhitvāt uttarapadam prasiddham . \\
\noindent \textbf{(P\_2,1.51.3)} uttarapade prasiddhe uttarapade iti dviguḥ bhaviṣyati . \\
\noindent \textbf{(P\_2,1.51.3)} sarvatra matvarthe pratiṣedhaḥ . \\
\noindent \textbf{(P\_2,1.51.3)} sarveṣu pakṣeṣu dvigusañjñāyāḥ matvarthe pratiṣedhaḥ vaktavyaḥ . \\
\noindent \textbf{(P\_2,1.51.3)} kim prayojanam . \\
\noindent \textbf{(P\_2,1.51.3)} pañcakhaṭvā daśakhaṭvā . \\
\noindent \textbf{(P\_2,1.51.3)} dvigoḥ iti īkāraḥ mā bhūt . \\
\noindent \textbf{(P\_2,1.51.3)} pañcaguḥ daśaguḥ . \\
\noindent \textbf{(P\_2,1.51.3)} goḥ ataddhitaluki iti ṭac mā bhūt iti . \\
\noindent \textbf{(P\_2,1.52)} kim anantare yoge saṅkhyāpūrvaḥ saḥ dvigusañjñaḥ āhosvit pūrvamātre . \\
\noindent \textbf{(P\_2,1.52)} kim ca ataḥ . \\
\noindent \textbf{(P\_2,1.52)} yadi anantare yoge ekaśāṭī dvigoḥ iti īkāraḥ na prāpnoti . \\
\noindent \textbf{(P\_2,1.52)} atha pūrvamātre akabhikṣā atra api prāpnoti . \\
\noindent \textbf{(P\_2,1.52)} astu anantare . \\
\noindent \textbf{(P\_2,1.52)} kamam ekaśāṭī . \\
\noindent \textbf{(P\_2,1.52)} īkārāntena samāsaḥ bhaviṣyati . \\
\noindent \textbf{(P\_2,1.52)} ekā śāṭī ekaśāṭī . \\
\noindent \textbf{(P\_2,1.52)} iha tarhi ekāpūpī dvigoḥ iti īkāraḥ na prāpnoti . \\
\noindent \textbf{(P\_2,1.52)} astu tarhi pūrvamātre. katham ekabhikṣā . \\
\noindent \textbf{(P\_2,1.52)} ṭābantena samāsaḥ bhaviṣyati . \\
\noindent \textbf{(P\_2,1.52)} ekā bhikṣā ekabhikṣā . \\
\noindent \textbf{(P\_2,1.52)} iha tarhi saptarṣayaḥ igante dvigau iti eṣaḥ svaraḥ prāpnoti . \\
\noindent \textbf{(P\_2,1.52)} astu tarhi anantare . \\
\noindent \textbf{(P\_2,1.52)} katham ekāpūpī . \\
\noindent \textbf{(P\_2,1.52)} samāhāre iti eva siddham . \\
\noindent \textbf{(P\_2,1.52)} kaḥ punaḥ atra samāhāraḥ . \\
\noindent \textbf{(P\_2,1.52)} yat taddānam sambhramaḥ vā . \\
\noindent \textbf{(P\_2,1.52)} iha tarhi pañcahotāraḥ daśahotāraḥ igante dvigau iti eṣaḥ svaraḥ na prapnoti . \\
\noindent \textbf{(P\_2,1.52)} astu tarhi pūrvamātre . \\
\noindent \textbf{(P\_2,1.52)} katham saptarṣayaḥ . \\
\noindent \textbf{(P\_2,1.52)} antodāttaprakaraṇe tricakrādīnām chandasi iti evam etat siddham . \\
\noindent \textbf{(P\_2,1.52)} atha vā punaḥ astu anantare . \\
\noindent \textbf{(P\_2,1.52)} katham pañcahotāraḥ daśahotāraḥ . \\
\noindent \textbf{(P\_2,1.52)} ādyudāttaprakaraṇe divodāsādīnām chandasi iti eva siddham . \\
\noindent \textbf{(P\_2,1.53)} kim udāharaṇm . \\
\noindent \textbf{(P\_2,1.53)} vaiyākaraṇakhasūciḥ . \\
\noindent \textbf{(P\_2,1.53)} kim vyākaraṇam kutsitam āhosvit vaiyākaraṇaḥ . \\
\noindent \textbf{(P\_2,1.53)} vaiyākaraṇaḥ kutsitaḥ . \\
\noindent \textbf{(P\_2,1.53)} tasmin kutsite tatstham api kutsitam bhavati . \\
\noindent \textbf{(P\_2,1.55)} upamānāni iti ucyate . \\
\noindent \textbf{(P\_2,1.55)} kāni punaḥ upamānāni . \\
\noindent \textbf{(P\_2,1.55)} kim yat eva upamānam tat eva upameyam āhosvit anyat upamānam anyat upameyam . \\
\noindent \textbf{(P\_2,1.55)} kim ca ataḥ . \\
\noindent \textbf{(P\_2,1.55)} yadi yat eva upamānam tat eva upameyam kaḥ iha upamārthaḥ gauḥ iva gauḥ iti . \\
\noindent \textbf{(P\_2,1.55)} atha anyat eva upamānam anyat upameyam kaḥ iha upamārthaḥ gauḥ iva aśvaḥ iti . \\
\noindent \textbf{(P\_2,1.55)} evam tarhi yatra kim cit sāmānyam kaḥ cit viśeṣaḥ tatra upamānopameye bhavataḥ . \\
\noindent \textbf{(P\_2,1.55)} kim vaktavyam etat . \\
\noindent \textbf{(P\_2,1.55)} na hi . \\
\noindent \textbf{(P\_2,1.55)} katham anucyamānam gaṃsyate . \\
\noindent \textbf{(P\_2,1.55)} mānam hi nāma anirjñātajñānārtham upādīyate anirjñātam artham jñāsyāmi iti . \\
\noindent \textbf{(P\_2,1.55)} tat samīpe yat na atyantāya mimīte tat upamānam . \\
\noindent \textbf{(P\_2,1.55)} gauḥ iva gavayaḥ iti . \\
\noindent \textbf{(P\_2,1.55)} gauḥ nirjñātaḥ gavayaḥ anirjñātaḥ . \\
\noindent \textbf{(P\_2,1.55)} kāmam tarhi anena eva hetunā yasya gavayaḥ nirjñātaḥ syāt gauḥ anirjñātaḥ tena kartavyam syāt gavayaḥ iva gauḥ iti. bāḍham kartavyam . \\
\noindent \textbf{(P\_2,1.55)} kim punaḥ iha udāharaṇam . \\
\noindent \textbf{(P\_2,1.55)} śastrīśyāmā . \\
\noindent \textbf{(P\_2,1.55)} kva punaḥ ayam śyāmāśabdaḥ vartate . \\
\noindent \textbf{(P\_2,1.55)} śatryām iti āha . \\
\noindent \textbf{(P\_2,1.55)} kena idānīm devadattā abhidhīyate . \\
\noindent \textbf{(P\_2,1.55)} samāsena. yadi evam śastrīśyāmo devadattaḥ iti na sidhyati . \\
\noindent \textbf{(P\_2,1.55)} upasarjanasya iti hrasvatvam bhaviṣyati . \\
\noindent \textbf{(P\_2,1.55)} yadi tarhi upasarjanāni api evañjātīyakāni bhavanti tittirikalmāṣī kumbhakapālalohinī anupasarjanalakṣaṇaḥ īkāraḥ na prāpnoti . \\
\noindent \textbf{(P\_2,1.55)} evam tarhi śastryām eva śastrīśabdaḥ vartate devadattāyām śyāmāśabdaḥ . \\
\noindent \textbf{(P\_2,1.55)} evam api guṇaḥ anirdiṣṭaḥ bhavati . \\
\noindent \textbf{(P\_2,1.55)} bahavaḥ śastryām guṇāḥ tīkṣṇā sūkṣmā pṛthuḥ iti . \\
\noindent \textbf{(P\_2,1.55)} anirdiśyamānasya api guṇasya bhavati loke sampratyayaḥ . \\
\noindent \textbf{(P\_2,1.55)} tat yathā candramukhī devadattā iti . \\
\noindent \textbf{(P\_2,1.55)} bahavaḥ candre guṇāḥ yā ca asau priyadarśanatā sā gamyate . \\
\noindent \textbf{(P\_2,1.55)} evam api samānādhikaraṇena iti vartate . \\
\noindent \textbf{(P\_2,1.55)} vyadhikaraṇatvāt samāsaḥ na prāpnoti . \\
\noindent \textbf{(P\_2,1.55)} kim hi vacanāt na bhavati . \\
\noindent \textbf{(P\_2,1.55)} yadi api tāvat vacanāt samāsaḥ syāt iha tu khalu mṛgī iva capalā mṛgacapalā samānādhikaraṇalakṣaṇaḥ puṃvadbhāvaḥ na prāpnoti . \\
\noindent \textbf{(P\_2,1.55)} evam tarhi tasyām eva ubhayam vartate . \\
\noindent \textbf{(P\_2,1.55)} etat ca atra yuktam yat tasyām eva ubhayam vartate iti . \\
\noindent \textbf{(P\_2,1.55)} itarathā hi bahu apekṣyam syāt . \\
\noindent \textbf{(P\_2,1.55)} yadi tāvat evam vigrahaḥ kriyate śastrī iva śyāmā devadattā iti śastryām śyāmā iti etat apekṣyam . \\
\noindent \textbf{(P\_2,1.55)} atha api evam vigrahaḥ kriyate yathā sāstrīśyāmā tadvat iyam devadattā iti evam api devadattāyām śyāmā iti apekṣyam syāt . \\
\noindent \textbf{(P\_2,1.55)} evam api guṇaḥ anirdiṣṭaḥ bhavati . \\
\noindent \textbf{(P\_2,1.55)} bahavaḥ śastryām guṇāḥ tīkṣṇā sūkṣmā pṛthuḥ iti . \\
\noindent \textbf{(P\_2,1.55)} anirdiśyamānasya api guṇasya bhavati loke sampratyayaḥ . \\
\noindent \textbf{(P\_2,1.55)} tat yathā candramukhī devadattā iti . \\
\noindent \textbf{(P\_2,1.55)} bahavaḥ candre guṇāḥ yā ca asau priyadarśanatā sā gamyate . \\
\noindent \textbf{(P\_2,1.55)} upamānasamāse guṇavacanasya viśeṣabhāktvāt sāmanyavacanāprasiddhiḥ . \\
\noindent \textbf{(P\_2,1.55)} upamānasamāse guṇavacanasya viśeṣabhāktvāt sāmanyavacanasya aprasiddhiḥ syāt . \\
\noindent \textbf{(P\_2,1.55)} śastrīśyāmā . \\
\noindent \textbf{(P\_2,1.55)} śyāmāśabdaḥ ayam śastrīśabdena abhisambadhyamānaḥ viśeṣavacanaḥ sampadyate . \\
\noindent \textbf{(P\_2,1.55)} tatra sāmānyavacanaiḥ iti samāsaḥ na prāpnoti . \\
\noindent \textbf{(P\_2,1.55)} na vā śyāmatvasyo uhhayatra bhāvāt tadvācaktvāt ca śabdasya sāmānyavacanaprasiddhiḥ . \\
\noindent \textbf{(P\_2,1.55)} na vā eṣaḥ doṣaḥ . \\
\noindent \textbf{(P\_2,1.55)} kim kāraṇam . \\
\noindent \textbf{(P\_2,1.55)} śyāmatvasyo uhhayatra bhāvāt . \\
\noindent \textbf{(P\_2,1.55)} ubhayatra eva śyāmatvam asti śastryām devadattāyām ca . \\
\noindent \textbf{(P\_2,1.55)} tadvācaktvāt ca śabdasya . \\
\noindent \textbf{(P\_2,1.55)} sāmānyavacanaprasiddhiḥ tadvācakaḥ ca atra śyāmāśabdaḥ prayujyate . \\
\noindent \textbf{(P\_2,1.55)} kimvācakaḥ . \\
\noindent \textbf{(P\_2,1.55)} ubhayavācakaḥ . \\
\noindent \textbf{(P\_2,1.55)} śyāmatvasya ubhayatra bhāvāt tadvācakatvāt ca śabdasya sāmānyavacanam prasiddham . \\
\noindent \textbf{(P\_2,1.55)} sāmānyavacane prasiddhe sāmānyavacanaiḥ iti samāsaḥ bhaviṣyati . \\
\noindent \textbf{(P\_2,1.55)} na ca avaśyam saḥ eva sāmānyavacanaḥ yaḥ bahūnām sāmānyam āha. dvayoḥ api sāmānyam āha saḥ api sāmānyavacanaḥ eva . \\
\noindent \textbf{(P\_2,1.55)} atha vā sāmānyavacanaiḥ iti ucyate . \\
\noindent \textbf{(P\_2,1.55)} sarvaḥ ca śabdaḥ anyena śabdena abhisambadhyamānaḥ viśeṣavacanaḥ sampadyate . \\
\noindent \textbf{(P\_2,1.55)} te evam vijñāsyāmaḥ prāk abhisambandhāt sāmānyavacanaḥ iti . \\
\noindent \textbf{(P\_2,1.56)} sāmānyāprayoge iti kimartham . \\
\noindent \textbf{(P\_2,1.56)} iha mā bhūt . \\
\noindent \textbf{(P\_2,1.56)} puruṣaḥ ayam vyāghraḥ iva śūraḥ . \\
\noindent \textbf{(P\_2,1.56)} puruṣaḥ ayam vyāghraḥ iva balavān . \\
\noindent \textbf{(P\_2,1.56)} sāmānyāprayoge iti śakyam akartum . \\
\noindent \textbf{(P\_2,1.56)} kasmāt na bhavati puruṣaḥ ayam vyāghraḥ iva śūraḥ . \\
\noindent \textbf{(P\_2,1.56)} puruṣaḥ ayam vyāghraḥ iva balavān iti . \\
\noindent \textbf{(P\_2,1.56)} asāmarthyāt . \\
\noindent \textbf{(P\_2,1.56)} katham asāmarthyam . \\
\noindent \textbf{(P\_2,1.56)} sāpekṣam asamartham bhavati iti . \\
\noindent \textbf{(P\_2,1.56)} evam tarhi siddhe sati yat sāmānyāprayoge iti pratiṣedham śāsti tat jñāpayati ācāryaḥ bhavati vai pradhānasya sāpekṣasya api samāsaḥ iti . \\
\noindent \textbf{(P\_2,1.56)} kim etasya jñāpane prayojanam . \\
\noindent \textbf{(P\_2,1.56)} rājapuruṣaḥ abhirūpaḥ rājapuruṣaḥ darśanīyaḥ atra vṛttiḥ siddhā bhavati . \\
\noindent \textbf{(P\_2,1.57)} viśeṣaṇaviśeṣyayoḥ ubhayaviśeṣaṇatvāt ubhayoḥ ca viśeṣyatvāt upasarjanāprasiddhiḥ . \\
\noindent \textbf{(P\_2,1.57)} viśeṣaṇaviśeṣyayoḥ ubhayaviśeṣaṇatvāt ubhayoḥ ca viśeṣyatvāt upasarjansya aprasiddhiḥ . \\
\noindent \textbf{(P\_2,1.57)} kṛṣṇatilāḥ iti kṛṣṇaśabdaḥ ayam tilaśabdena abhisambadhyamānaḥ viśeṣaṇavacanaḥ sampadyate . \\
\noindent \textbf{(P\_2,1.57)} tathā tilaśabdaḥ kṛṣṇaśabdena abhisambadhyamānaḥ viśeṣaṇavacanaḥ sampadyate . \\
\noindent \textbf{(P\_2,1.57)} tat ubhayam viśeṣaṇam bhavati ubhayam ca viśeṣyam . \\
\noindent \textbf{(P\_2,1.57)} viśeṣaṇaviśeṣyayoḥ ubhayaviśeṣaṇatvāt ubhayoḥ ca viśeṣyatvāt upasarjansya aprasiddhiḥ . \\
\noindent \textbf{(P\_2,1.57)} na vā anyatarasya pradhānabhāvāt tadviśeṣakatvāt ca aparasya upasarjanaprasiddhiḥ . \\
\noindent \textbf{(P\_2,1.57)} na vā eṣaḥ doṣaḥ . \\
\noindent \textbf{(P\_2,1.57)} kim kāraṇam . \\
\noindent \textbf{(P\_2,1.57)} anyatarasya pradhānabhāvāt . \\
\noindent \textbf{(P\_2,1.57)} anyatarat atra pradhānam . \\
\noindent \textbf{(P\_2,1.57)} tadviśeṣakatvāt ca aparasya . \\
\noindent \textbf{(P\_2,1.57)} tadviśeṣakam ca aparam . \\
\noindent \textbf{(P\_2,1.57)} anyatarasya pradhānabhāvāt tadviśeṣakatvāt ca aparasya upasarjanasañjñā bhaviṣyati . \\
\noindent \textbf{(P\_2,1.57)} yadā asya tilāḥ prādhānyena vivakṣitāḥ bhavanti kṛṣṇaḥ viśeṣaṇatvena tadā tilāḥ pradhānam kṛṣṇaḥ viśeṣaṇam . \\
\noindent \textbf{(P\_2,1.57)} kāmam tarhi anena eva hetunā yasya kṛṣṇāḥ prādhānyena vivakṣitāḥ bhavanti tilāḥ viśeṣaṇatvena tena kartavyam tilakṛṣṇāḥ iti . \\
\noindent \textbf{(P\_2,1.57)} na kartavyam . \\
\noindent \textbf{(P\_2,1.57)} na hi ayam dvandvaḥ tilāḥ ca kṛṣṇāḥ ca iti . \\
\noindent \textbf{(P\_2,1.57)} na khalu api ṣaṣṭhīsamāsaḥ tilānām kṛṣṇāḥ iti . \\
\noindent \textbf{(P\_2,1.57)} kim tarhi . \\
\noindent \textbf{(P\_2,1.57)} dvau imau pradhānaśabdau ekasmin arthe yugapat avarudhyete . \\
\noindent \textbf{(P\_2,1.57)} na ca dvayoḥ pradhānaśabdayoḥ ekasmin arthe yugapat avarudhyamānayoḥ kim cit api prayojanam asti . \\
\noindent \textbf{(P\_2,1.57)} tatra prayogāt etat gantavyam . \\
\noindent \textbf{(P\_2,1.57)} nūnam atra anyatarat pradhānam tadviśeṣakam ca aparam iti . \\
\noindent \textbf{(P\_2,1.57)} tatra tu etāvān sandehaḥ kim pradhānam kim viśeṣaṇam iti . \\
\noindent \textbf{(P\_2,1.57)} saḥ ca api kva sandehaḥ . \\
\noindent \textbf{(P\_2,1.57)} yatra ubhau guṇaśabdau . \\
\noindent \textbf{(P\_2,1.57)} tat yathā kuñjakhañjakaḥ khañjakubjakaḥ iti . \\
\noindent \textbf{(P\_2,1.57)} yatra hi anyatarat dravyam anyataraḥ guṇaḥ tatra yat dravyam tat pradhānam . \\
\noindent \textbf{(P\_2,1.57)} tat yathā śuklam ālabheta kṛṣṇam ālabheta iti na piṣṭapiṇḍīm ālabhya kṛtī bhavati . \\
\noindent \textbf{(P\_2,1.57)} avaśyam tadguṇam dravyam ākāṅkṣati . \\
\noindent \textbf{(P\_2,1.57)} katham tarhi imau dvau pradhānaśabdau ekasmin arthe yugapat avarudhyete vṛkṣaḥ śiṃśipā iti . \\
\noindent \textbf{(P\_2,1.57)} na etayoḥ āvaśyakaḥ samāveśaḥ . \\
\noindent \textbf{(P\_2,1.57)} na hi avṛkṣaḥ śiṃśipā asti . \\
\noindent \textbf{(P\_2,1.58)} atha kimartham uttaratra evamādi anukramaṇam kriyate na viśeṣaṇam viśeṣyeṇa bahulam iti eva siddham . \\
\noindent \textbf{(P\_2,1.58)} bahulavacanasya akṛtsnatvāt uttaratrānukramaṇasāmarthyam . \\
\noindent \textbf{(P\_2,1.58)} akṛtsnam bahulavacanam iti uttaratra anukramaṇam kriyate . \\
\noindent \textbf{(P\_2,1.58)} yadi akṛtsnam yat anena kṛtam akṛtam tat . \\
\noindent \textbf{(P\_2,1.58)} evam tarhi na brūmaḥ akṛtsnam iti . \\
\noindent \textbf{(P\_2,1.58)} kṛtsnam ca kārakam ca sādhakam ca nirvartakam ca . \\
\noindent \textbf{(P\_2,1.58)} yat ca anena kṛtam suktṛtam tat . \\
\noindent \textbf{(P\_2,1.58)} kimartham tarhi evamādi anukramaṇam kriyate . \\
\noindent \textbf{(P\_2,1.58)} udāharaṇabhūyastvāt . \\
\noindent \textbf{(P\_2,1.58)} te khalu api vidhayaḥ suparigṛhītāḥ bhavanti yeṣu lakṣaṇam prapañcaḥ ca . \\
\noindent \textbf{(P\_2,1.58)} kevalam lakṣaṇam kevalaḥ prapañcaḥ vā na tathā kārakam bhavati . \\
\noindent \textbf{(P\_2,1.58)} avaśyam khalu asmābhiḥ idam vaktavyam bahulam anyatarasyām ubhayathā vā ekeṣām iti . \\
\noindent \textbf{(P\_2,1.58)} sarvavedapāṛiṣadam hi idam śāstram . \\
\noindent \textbf{(P\_2,1.58)} tatra na ekaḥ panthāḥ śakyaḥ āsthātum \\
\noindent \textbf{(P\_2,1.59)} śreṇyādayaḥ paṭhyante . \\
\noindent \textbf{(P\_2,1.59)} kṛtādiḥ ākṛtigaṇaḥ . \\
\noindent \textbf{(P\_2,1.59)} śreṇyādiṣu cvyarthavacanam . \\
\noindent \textbf{(P\_2,1.59)} śreṇyādiṣu cvyarthagrahaṇam kartavyam . \\
\noindent \textbf{(P\_2,1.59)} aśreṇayaḥ śreṇayaḥ kṛtāḥ śreṇikṛtāḥ . \\
\noindent \textbf{(P\_2,1.59)} yadā hi śreṇayaḥ eva kim cit kriyante tadā mā bhūt . \\
\noindent \textbf{(P\_2,1.59)} anyatra ayam cvyarthagrahaṇeṣu cvyantasya pratiṣedham śāsti . \\
\noindent \textbf{(P\_2,1.59)} tat iha na tathā . \\
\noindent \textbf{(P\_2,1.59)} kim kāraṇam . \\
\noindent \textbf{(P\_2,1.59)} anyatra pūrvam cvyantakāryam param cvyarthakāryam . \\
\noindent \textbf{(P\_2,1.59)} iha punaḥ pūrvam cyvarthakāryam param cvyantakāryam iti . \\
\noindent \textbf{(P\_2,1.60)} nañviśiṣṭe samānaprakṛtivacanam . \\
\noindent \textbf{(P\_2,1.60)} nañviśiṣṭe samānaprakṛtigrahaṇam kartavyam . \\
\noindent \textbf{(P\_2,1.60)} iha mā būt . \\
\noindent \textbf{(P\_2,1.60)} siddham ca abhuktam ca iti . \\
\noindent \textbf{(P\_2,1.60)} anañ iti ca pratiṣedhaḥ kartavyaḥ . \\
\noindent \textbf{(P\_2,1.60)} iha mā bhūt . \\
\noindent \textbf{(P\_2,1.60)} kartavyam akṛtam iti . \\
\noindent \textbf{(P\_2,1.60)} nuḍiḍadhikena ca . \\
\noindent \textbf{(P\_2,1.60)} nuḍiḍadhikena ca samāsaḥ vaktavyaḥ . \\
\noindent \textbf{(P\_2,1.60)} iha api yathā syāt . \\
\noindent \textbf{(P\_2,1.60)} aśitānaśitena jīvati . \\
\noindent \textbf{(P\_2,1.60)} kliṣṭākliśitena jīvati . \\
\noindent \textbf{(P\_2,1.60)} kim ucyate samānaprakṛtigrahaṇam kartavyam iti yadā nañviśiṣṭena iti ucyate . \\
\noindent \textbf{(P\_2,1.60)} na ca atra nañkṛtaḥ eva viśeṣaḥ . \\
\noindent \textbf{(P\_2,1.60)} kim tarhi . \\
\noindent \textbf{(P\_2,1.60)} prakṛtikṛtaḥ api . \\
\noindent \textbf{(P\_2,1.60)} ayam viśiṣṭaśabdaḥ asti eva avadhāraṇe vartate . \\
\noindent \textbf{(P\_2,1.60)} tat yathā . \\
\noindent \textbf{(P\_2,1.60)} devadattayajñadattau āḍhyau abhirūpau darśanīyau pakṣavantau devadattaḥ tu yajñadattāt svādhyāyena viśiṣṭaḥ . \\
\noindent \textbf{(P\_2,1.60)} svādhyāyena eva iti gamyate . \\
\noindent \textbf{(P\_2,1.60)} anye guṇāḥ samāḥ bhavanti . \\
\noindent \textbf{(P\_2,1.60)} asti ādhikye vartate . \\
\noindent \textbf{(P\_2,1.60)} tat yathā . \\
\noindent \textbf{(P\_2,1.60)} devadattayajñadattau āḍhyau abhirūpau darśanīyau pakṣavantau devadattaḥ tu yajñadattāt svādhyāyena viśiṣṭaḥ . \\
\noindent \textbf{(P\_2,1.60)} svādhyāyena adhikaḥ anye guṇāḥ avivakṣitāḥ bhavanti . \\
\noindent \textbf{(P\_2,1.60)} tat yadā tāvat avadhāraṇe viśiṣṭaśabdaḥ tadā na eva arthaḥ samānaprakṛtigrahaṇena . \\
\noindent \textbf{(P\_2,1.60)} na iha bhaviṣyati . \\
\noindent \textbf{(P\_2,1.60)} siddham ca abhuktam ca iti . \\
\noindent \textbf{(P\_2,1.60)} na api anañ iti pratiṣedhena . \\
\noindent \textbf{(P\_2,1.60)} na iha bhaviṣyati kartavyam akṛtam iti . \\
\noindent \textbf{(P\_2,1.60)} nuḍiḍadhikena api tu tadā samāsaḥ na prāpnoti . \\
\noindent \textbf{(P\_2,1.60)} yadā ādhikye viśiṣṭaśabdaḥ tadā samānaprakṛtigrahaṇam kartavyam . \\
\noindent \textbf{(P\_2,1.60)} iha mā bhūt ṣiddham ca abhuktam ca iti . \\
\noindent \textbf{(P\_2,1.60)} anañ iti ca pratiṣedhaḥ kartavyaḥ . \\
\noindent \textbf{(P\_2,1.60)} iha mā bhūt . \\
\noindent \textbf{(P\_2,1.60)} kartavyam akṛtam iti . \\
\noindent \textbf{(P\_2,1.60)} nuḍiḍadhikena api tu samāsadḥ siddhaḥ bhavati . \\
\noindent \textbf{(P\_2,1.60)} tatra ādhikye viśiṣṭagrahaṇam matvā samānaprakṛtigrahaṇam codyate . \\
\noindent \textbf{(P\_2,1.60)} avadhāraṇam nañā cet nuḍiḍviśiṣṭena na prakalpeta . \\
\noindent \textbf{(P\_2,1.60)} atha cet adhikavivakṣā kāryam tulyaprakṛtikena iti . \\
\noindent \textbf{(P\_2,1.60)} kṛtāpakṛtādīnām ca upasaṅkhyānam . \\
\noindent \textbf{(P\_2,1.60)} kṛtāpakṛtādīnām ca upasaṅkhyānam . \\
\noindent \textbf{(P\_2,1.60)} kṛtāpakṛtam bhuktavibhuktam pītavipītam . \\
\noindent \textbf{(P\_2,1.60)} siddham tu ktena visamāptau anañ . \\
\noindent \textbf{(P\_2,1.60)} siddham etat . \\
\noindent \textbf{(P\_2,1.60)} katham. ktāntena kriyāvisamāptau anañ ktāntam samasyate iti vaktavyam . \\
\noindent \textbf{(P\_2,1.60)} gatapratyāgatādīnām ca upasaṅkhyānam . \\
\noindent \textbf{(P\_2,1.60)} gatapratyāgatādīnām ca upasaṅkhyānam kartavyam . \\
\noindent \textbf{(P\_2,1.60)} gatapratyāgatam yātānuyātam puṭāpuṭikā krayākrayikā phalāphalikā mānonmānikā . \\
\noindent \textbf{(P\_2,1.67)} ayuktaḥ ayam nirdeśaḥ . \\
\noindent \textbf{(P\_2,1.67)} samānādhikaraṇena iti vartate . \\
\noindent \textbf{(P\_2,1.67)} kaḥ prasaṅgaḥ yad vyadhikaraṇānām samāsaḥ syāt . \\
\noindent \textbf{(P\_2,1.67)} evam tarhi jñāpayati ācāryaḥ yathājātīyakam uktam uttarapadam tathājātīyakena pūrvapadena samasyate iti . \\
\noindent \textbf{(P\_2,1.67)} kim etasya jñāpane prayojanam . \\
\noindent \textbf{(P\_2,1.67)} prātipadikagrahaṇe liṅgaviśiṣṭasya api grahaṇam bhavati iti eṣā paribhāṣā na kartavyā bhavati . \\
\noindent \textbf{(P\_2,1.69.1)} idam vicāryate : varṇena tṛtīyāsamāsaḥ vā syāt : kṛṣṇena sāraṅgaḥ kṛṣṇasāraṅgaḥ samānādhikaraṇena vā : kṛṣṇaḥ sāraṅgaḥ kṛṣṇasāraṅgaḥ iti . \\
\noindent \textbf{(P\_2,1.69.1)} kaḥ ca atra viśeṣaḥ . \\
\noindent \textbf{(P\_2,1.69.1)} varṇena tṛtīyāsamāsaḥ etapratiṣedhe varṇagrahaṇam . \\
\noindent \textbf{(P\_2,1.69.1)} varṇena tṛtīyāsamāsaḥ etapratiṣedhe varṇagrahaṇam kartavyam . \\
\noindent \textbf{(P\_2,1.69.1)} tṛtīyā pūrvapadam prakṛtisvaram bhavati . \\
\noindent \textbf{(P\_2,1.69.1)} anete varṇaḥ iti vaktavyam . \\
\noindent \textbf{(P\_2,1.69.1)} atha dvitīyena varṇagrahaṇena etaviśeṣaṇena arthaḥ . \\
\noindent \textbf{(P\_2,1.69.1)} bāḍham arthaḥ yadi avarṇa etaśabdaḥ asti . \\
\noindent \textbf{(P\_2,1.69.1)} nanu ca ayam asti : ā* itaḥ etaḥ , kṛṣṇetaḥ , lohitetaḥ iti . \\
\noindent \textbf{(P\_2,1.69.1)} na arthaḥ evamarthena varṇagrahaṇena . \\
\noindent \textbf{(P\_2,1.69.1)} yadi tāvat ayam kartari ktaḥ tṛtīyā karmaṇi iti anena svareṇa bhavitavyam . \\
\noindent \textbf{(P\_2,1.69.1)} atha api kartari paratvāt kṛtsvareṇa bhavitavyam . \\
\noindent \textbf{(P\_2,1.69.1)} atha samānādhikaraṇaḥ . \\
\noindent \textbf{(P\_2,1.69.1)} samānādhikaraṇe dviḥ varṇagrahaṇam . \\
\noindent \textbf{(P\_2,1.69.1)} samānādhikaraṇe dviḥ varṇagrahaṇam kartavyam . \\
\noindent \textbf{(P\_2,1.69.1)} varṇaḥ varṇeṣu anete iti vaktavyam . \\
\noindent \textbf{(P\_2,1.69.1)} ekam varṇagrahaṇam kartavyam iha mā bhūt . \\
\noindent \textbf{(P\_2,1.69.1)} paramaśuklaḥ paramakṛṣṇaḥ iti . \\
\noindent \textbf{(P\_2,1.69.1)} dvitīyam varṇagrahaṇam kartavyam iha mā bhūt . \\
\noindent \textbf{(P\_2,1.69.1)} kṛṣṇatilāḥ iti . \\
\noindent \textbf{(P\_2,1.69.1)} ekam varṇagrahaṇam anakrthakam . \\
\noindent \textbf{(P\_2,1.69.1)} anyataratra kasmāt na bhavati . \\
\noindent \textbf{(P\_2,1.69.1)} lakṣaṇapratipadikoktayoḥ pratipadoktasya eva iti . \\
\noindent \textbf{(P\_2,1.69.1)} evam sati . \\
\noindent \textbf{(P\_2,1.69.1)} tāni etāni trīṇi varṇagrahaṇāni bhavanti samāsavidhau dve svaravidhau ca ekam . \\
\noindent \textbf{(P\_2,1.69.1)} yasya api tṛtīyāsamāsaḥ tasya api tāni eva trīṇi varṇagrahaṇāni bhavanti samāsavidhau dve svaravidhau ca ekam . \\
\noindent \textbf{(P\_2,1.69.1)} sāmānyena mama tṛtīyāsamāsaḥ bhaviṣyati tṛtīyā tatkṛtārthena guṇavacanena iti . \\
\noindent \textbf{(P\_2,1.69.1)} avaśyam varṇena pratipadam samāsaḥ vaktavyaḥ yatra tena na sidhyati tadartham . \\
\noindent \textbf{(P\_2,1.69.1)} kva ca tena na sidhyati . \\
\noindent \textbf{(P\_2,1.69.1)} śukababhruḥ haritababhruḥ iti . \\
\noindent \textbf{(P\_2,1.69.1)} tathā ca sati tāni eva trīṇi varṇagrahaṇāni bhavanti samāsavidhau dve svaravidhau ca ekam . \\
\noindent \textbf{(P\_2,1.69.1)} atha samānādhikaraṇaḥ sāmānyena siddhaḥ syāt . \\
\noindent \textbf{(P\_2,1.69.1)} bāḍham siddhaḥ . \\
\noindent \textbf{(P\_2,1.69.1)} katham . \\
\noindent \textbf{(P\_2,1.69.1)} viśeṣaṇam viśeṣyeṇa bahulam iti . \\
\noindent \textbf{(P\_2,1.69.1)} evam api dve varṇagrahaṇe kartavye svaravidhau eva pratipadoktasya abhāvāt . \\
\noindent \textbf{(P\_2,1.69.1)} tasmāt samānādhikaraṇaḥ iti eṣaḥ pakṣaḥ jyāyān . \\
\noindent \textbf{(P\_2,1.69.2)} samānādhikaraṇādhikāre pradhānopasarjanānām param param vipratiṣedhena . \\
\noindent \textbf{(P\_2,1.69.2)} samānādhikaraṇādhikāre pradhānopasarjanānām param param bhavati vipratiṣedhena . \\
\noindent \textbf{(P\_2,1.69.2)} pradhānānām pradhānam upasarjanānām upasarjanam . \\
\noindent \textbf{(P\_2,1.69.2)} pradhānānām tāvat pradhānam . \\
\noindent \textbf{(P\_2,1.69.2)} vṛdārakanāgakuñjaraiḥ pūjyamānam iti asya avakāśaḥ govṛndārakaḥ aśvavṛndārakaḥ . \\
\noindent \textbf{(P\_2,1.69.2)} poṭāyuvatīnām avakāśaḥ ibhyayuvatiḥ āḍhyayuvatiḥ . \\
\noindent \textbf{(P\_2,1.69.2)} iha ubhayam prāpnoti . \\
\noindent \textbf{(P\_2,1.69.2)} nāgayuvatiḥ vṛndārakayuvatiḥ . \\
\noindent \textbf{(P\_2,1.69.2)} pradhānānām param bhavati vipratiṣedhena . \\
\noindent \textbf{(P\_2,1.69.2)} upasarjanānām param upasarjanam . \\
\noindent \textbf{(P\_2,1.69.2)} sanmahatparamotkṛṣṭāḥ iti asya avakāśaḥ sadgavaḥ sadaśvaḥ . \\
\noindent \textbf{(P\_2,1.69.2)} kṛtyatulyākhyā ajātyā iti asya avakāśaḥ tulyaśvetaḥ tulyakṛṣṇaḥ . \\
\noindent \textbf{(P\_2,1.69.2)} iha ubhayam prāpnoti : tulyasat tulyamahān . \\
\noindent \textbf{(P\_2,1.69.2)} upasarjanānām param upasarjanam bhavati vipratiṣedhena . \\
\noindent \textbf{(P\_2,1.69.2)} samānādhikaraṇasamāsāt bahuvrīhiḥ . samānādhikaraṇasamāsāt bahuvrīhiḥ bhavati vipratiṣedhena . \\
\noindent \textbf{(P\_2,1.69.2)} samānādhikaraṇasamāsasya avakāśaḥ vīraḥ puruṣaḥ vīrapuruṣaḥ . \\
\noindent \textbf{(P\_2,1.69.2)} bahuvrīheḥ avakāśaḥ kaṇṭhekālaḥ . \\
\noindent \textbf{(P\_2,1.69.2)} iha ubhayam prāpnoti : vīrapuruṣakaḥ grāmaḥ . \\
\noindent \textbf{(P\_2,1.69.2)} bahuvrīhiḥ bhavati vipratiṣedhena . \\
\noindent \textbf{(P\_2,1.69.2)} kadā cit karmadhārayaḥ sarvadhanādyarthaḥ . \\
\noindent \textbf{(P\_2,1.69.2)} kadā cit karmadhārayaḥ bhavati bahuvrīheḥ . \\
\noindent \textbf{(P\_2,1.69.2)} kim prayojanam . \\
\noindent \textbf{(P\_2,1.69.2)} sarvadhanādyarthaḥ . \\
\noindent \textbf{(P\_2,1.69.2)} sarvadhanī sarvabījī sarvakeśī naṭaḥ gaurakharavat vanam gauramṛgavat vanam kṛṣṇasarpavān valmīkaḥ lohitaśālimān grāmaḥ . \\
\noindent \textbf{(P\_2,1.69.2)} kim prayojanam . \\
\noindent \textbf{(P\_2,1.69.2)} karmadhārayaprakṛtibhiḥ matvarthīyaiḥ abhidhānam yathā syāt . \\
\noindent \textbf{(P\_2,1.69.2)} kim ca kāraṇam na syāt . \\
\noindent \textbf{(P\_2,1.69.2)} bahuvrīhiṇā uktatvāt matvarthasya . \\
\noindent \textbf{(P\_2,1.69.2)} yadi uktatvam hetuḥ karmadhārayeṇa api uktatvāt na prāpnoti . \\
\noindent \textbf{(P\_2,1.69.2)} na khalu api sañjñāśrayaḥ matvarthīyaḥ . \\
\noindent \textbf{(P\_2,1.69.2)} kim tarhi . \\
\noindent \textbf{(P\_2,1.69.2)} arthāśrayaḥ . \\
\noindent \textbf{(P\_2,1.69.2)} saḥ yathā eva bahuvrīhiṇā uktatvāt na bhavati evam karmadhārayeṇa api uktatvāt na bhaviṣyati . \\
\noindent \textbf{(P\_2,1.69.2)} evam tarhi idam syāt : sarvāṇi dhanāni sarvadhanāni sarvadhanāni asya saniti sarvadhanī . \\
\noindent \textbf{(P\_2,1.69.2)} na evam śakyam . \\
\noindent \textbf{(P\_2,1.69.2)} nityam evam sati karmadhārayaḥ syāt . \\
\noindent \textbf{(P\_2,1.69.2)} tatra yat uktam kadā cit karmadhārayaḥ iti etat ayuktam . \\
\noindent \textbf{(P\_2,1.69.2)} evam tarhi bhavati vai kim cit ācāryāḥ kāryavat buddhim kṛtvā paṭhanti kāryāḥ śabdāḥ iti . \\
\noindent \textbf{(P\_2,1.69.2)} tadvat idam paṭhitam samānādhikaraṇasamādāt bahuvrīhiḥ kartavyaḥ kadā cit karmadhārayaḥ sarvadhanādyarthaḥ iti . \\
\noindent \textbf{(P\_2,1.69.2)} yad ucyate samānādhikaraṇasamāsāt bahuvrīhiḥ bhavati vipratiṣedhena iti na eṣaḥ yuktaḥ vipratiṣedhaḥ . \\
\noindent \textbf{(P\_2,1.69.2)} antaraṅgaḥ karmadhārayaḥ . \\
\noindent \textbf{(P\_2,1.69.2)} kā antaraṅgatā . \\
\noindent \textbf{(P\_2,1.69.2)} svapadārthe karmadhārayaḥ anyapadārthe bahuvrīhiḥ . \\
\noindent \textbf{(P\_2,1.69.2)} astu . \\
\noindent \textbf{(P\_2,1.69.2)} vibhāṣā karmadhārayaḥ . \\
\noindent \textbf{(P\_2,1.69.2)} yadā na karmadhārayaḥ tadā bahuvrīhiḥ bhaviṣyati . \\
\noindent \textbf{(P\_2,1.69.2)} evam api yadi atra kadā cit karmadhārayaḥ bhavati karmadhārayaprakṛtibhiḥ matvarthīyaiḥ abhidhānam prāpnoti . \\
\noindent \textbf{(P\_2,1.69.2)} sarvaḥ ca ayam evamarthaḥ yatnaḥ karmadhārayaprakṛtibhiḥ matvarthīyaiḥ abhidhānam mā bhūt iti . \\
\noindent \textbf{(P\_2,1.69.2)} evam tarhi na idam tasya yogasya udāharaṇam vipratiṣedhe param iti . \\
\noindent \textbf{(P\_2,1.69.2)} kim tarhi . \\
\noindent \textbf{(P\_2,1.69.2)} iṣṭiḥ iyam paṭhitā . \\
\noindent \textbf{(P\_2,1.69.2)} samānādhikaraṇasamāsāt bahuvrīhiḥ iṣṭaḥ kadā cit karmadhārayaḥ sarvadhanādyarthaḥ iti . \\
\noindent \textbf{(P\_2,1.69.2)} yadi iṣṭiḥ paṭhitā na arthaḥ anena . \\
\noindent \textbf{(P\_2,1.69.2)} iha hi sarve manuṣyāḥ alpena yatnena mahataḥ arthān ākāṅkṣanti . \\
\noindent \textbf{(P\_2,1.69.2)} ekena māṣeṇa śatasahasram . \\
\noindent \textbf{(P\_2,1.69.2)} ekena kuddālakena khārīsahasram . \\
\noindent \textbf{(P\_2,1.69.2)} tatra karmadhārayaprakṛtibhiḥ matvarthīyaiḥ abhidhānam astu bahuvrīhiṇā iti bahuvrīhiṇā bhaviṣyati laghutvāt . \\
\noindent \textbf{(P\_2,1.69.2)} katham sarvadhanī sarvabījī sarvakeśī naṭaḥ iti . \\
\noindent \textbf{(P\_2,1.69.2)} iniprakaraṇe sarvādeḥ inim vakṣyāmi . \\
\noindent \textbf{(P\_2,1.69.2)} tat ca avaśyam vaktavyam ṭhanaḥ bādhanārtham . \\
\noindent \textbf{(P\_2,1.69.2)} katham gaurakharavat vanam gauramṛgavat vanam kṛṣṇasarpavān valmīkaḥ lohitaśālimān grāmaḥ . \\
\noindent \textbf{(P\_2,1.69.2)} asti atra viśeṣaḥ . \\
\noindent \textbf{(P\_2,1.69.2)} jātyā atra abhisambandhaḥ kriyate . \\
\noindent \textbf{(P\_2,1.69.2)} kṛṣṇasarpaḥ nāma sarpajātiḥ sā asmin valmīke asti . \\
\noindent \textbf{(P\_2,1.69.2)} yadā hi antareṇa jātim tadvatām abhisambandhaḥ kriyate kṛṣṇasarpaḥ valmīkaḥ iti evam tadā bhaviṣyati . \\
\noindent \textbf{(P\_2,1.69.2)} pūrvapadātiśaye ātiśāyikāt bahuvrīhiḥ sūkṣmavastratarādyarthaḥ . \\
\noindent \textbf{(P\_2,1.69.2)} pūrvapadātiśaye ātiśāyikāt bahuvrīhiḥ bhavati vipratiṣedhena . \\
\noindent \textbf{(P\_2,1.69.2)} kim prayojanam . \\
\noindent \textbf{(P\_2,1.69.2)} sūkṣmavastratarādyarthaḥ . \\
\noindent \textbf{(P\_2,1.69.2)} ātiśāyikasya avakāśaḥ paṭutaraḥ paṭutamaḥ . \\
\noindent \textbf{(P\_2,1.69.2)} bahuvrīheḥ avakāśaḥ citraguḥ śabalaguḥ . \\
\noindent \textbf{(P\_2,1.69.2)} iha ubhayam prāpnoti sūkṣmavastrataraḥ tīkṣṇśṛṅgataraḥ . \\
\noindent \textbf{(P\_2,1.69.2)} bahuvrīhiḥ bhavati vipratiṣedhena . \\
\noindent \textbf{(P\_2,1.69.2)} na eṣaḥ yuktaḥ vipratiṣedhaḥ . \\
\noindent \textbf{(P\_2,1.69.2)} virpatiṣedhe param iti ucyate . \\
\noindent \textbf{(P\_2,1.69.2)} pūrvaḥ ca bahuvrīhiḥ paraḥ ātiśāyikaḥ . \\
\noindent \textbf{(P\_2,1.69.2)} iṣṭavācī paraśabdaḥ . \\
\noindent \textbf{(P\_2,1.69.2)} vipratiṣedhe param yat iṣṭam tat bhavati . \\
\noindent \textbf{(P\_2,1.69.2)} evam api ayuktaḥ . \\
\noindent \textbf{(P\_2,1.69.2)} antaraṅgaḥ ātiśāiyikaḥ . \\
\noindent \textbf{(P\_2,1.69.2)} kā antaraṅgatā . \\
\noindent \textbf{(P\_2,1.69.2)} ṅyāpprātipadikāt ātiśāyikaḥ subantānām bahuvrīhiḥ . \\
\noindent \textbf{(P\_2,1.69.2)} ātiśāyikaḥ api na antaraṅgaḥ . \\
\noindent \textbf{(P\_2,1.69.2)} katham . \\
\noindent \textbf{(P\_2,1.69.2)} samarthāt taddhitaḥ utpadyate sāmarthyam ca subantenta . \\
\noindent \textbf{(P\_2,1.69.2)} evam api antaraṅgaḥ . \\
\noindent \textbf{(P\_2,1.69.2)} katham . \\
\noindent \textbf{(P\_2,1.69.2)} svapadārthe ātiśāyikaḥ anyapadārthe bahuvrīhiḥ . \\
\noindent \textbf{(P\_2,1.69.2)} evam api na antaraṅgaḥ . \\
\noindent \textbf{(P\_2,1.69.2)} katham . \\
\noindent \textbf{(P\_2,1.69.2)} spardhāyām ātiśāyikaḥ bhavati . \\
\noindent \textbf{(P\_2,1.69.2)} na ca antareṇa pratiyoginam spardhā bhavati . \\
\noindent \textbf{(P\_2,1.69.2)} na eva vā atra ātiśāyikaḥ prāpnoti . \\
\noindent \textbf{(P\_2,1.69.2)} kim kāraṇam . \\
\noindent \textbf{(P\_2,1.69.2)} asāmarthyāt . \\
\noindent \textbf{(P\_2,1.69.2)} katham asāmarthyam. sāpekṣam asamartham bhavati iti . \\
\noindent \textbf{(P\_2,1.69.2)} yāvatā vastrāṇi tadvantam apekṣante tadvantam ca apekṣya vastrāṇām vastraiḥ yugapat spardhā bhavati . \\
\noindent \textbf{(P\_2,1.69.2)} nanu ca ayam ātiśāyikaḥ evamātmakaḥ satyām vyapekṣāyām vidhīyate . \\
\noindent \textbf{(P\_2,1.69.2)} satyam evamātmakaḥ yām ca na anatareṇa vyapekṣām pravṛttiḥ tasyam satyām bhavitavyam . \\
\noindent \textbf{(P\_2,1.69.2)} kām ca na antareṇa vyapekṣām ātiśāyikasya pravṛttiḥ . \\
\noindent \textbf{(P\_2,1.69.2)} yā hi pratiyoginam prati vyapekṣā . \\
\noindent \textbf{(P\_2,1.69.2)} yā hi tadvantam prati na tasyām bhavitavyam . \\
\noindent \textbf{(P\_2,1.69.2)} bahuvrīhiḥ api tarhi na prāpnoti . \\
\noindent \textbf{(P\_2,1.69.2)} kim kāraṇam . \\
\noindent \textbf{(P\_2,1.69.2)} asāmarthyāt eva . \\
\noindent \textbf{(P\_2,1.69.2)} katham asāmarthyam . \\
\noindent \textbf{(P\_2,1.69.2)} sāpekṣam asamartham bhavati iti . \\
\noindent \textbf{(P\_2,1.69.2)} yāvatā vastrāṇi vastrāntarāṇi apekṣante tadvatā ca abhisambandhaḥ . \\
\noindent \textbf{(P\_2,1.69.2)} evam tarhi na idam tasya yogasya udāharaṇam vipratiṣedhe param iti . \\
\noindent \textbf{(P\_2,1.69.2)} kim tarhi . \\
\noindent \textbf{(P\_2,1.69.2)} iṣṭiḥ iyam paṭhitā . \\
\noindent \textbf{(P\_2,1.69.2)} pūrvapadātiśaye ātiśāyikāt bahuvrīhiḥ iṣṭaḥ : sūkṣmavastratarādyarthaḥ iti . \\
\noindent \textbf{(P\_2,1.69.2)} yadi iṣṭiḥ iyam paṭhitā na arthaḥ anena . \\
\noindent \textbf{(P\_2,1.69.2)} katham yā eṣā yuktiḥ uktā : yāvatā vastrāṇi vastrāntarāṇi apekṣante tadvatā ca abhisambandhaḥ iti . \\
\noindent \textbf{(P\_2,1.69.2)} yadā hi antareṇa vastrāṇām vastraiḥ yugapat spardhām tadvatā ca abhisambandhaḥ kriyate niṣpratidvandvaḥ tadā bahuvrīhiḥ . \\
\noindent \textbf{(P\_2,1.69.2)} bahuvrīheḥ ātiśāyikaḥ . \\
\noindent \textbf{(P\_2,1.69.2)} na tarhi idānīm idam bhavati : sūkṣmataravastraḥ iti . \\
\noindent \textbf{(P\_2,1.69.2)} bhavati . \\
\noindent \textbf{(P\_2,1.69.2)} yadā antareṇa tadvantam vastrāṇām vastraiḥ yugapat spardhā niṣpratidvandvaḥ tadā ātiśāyikaḥ . \\
\noindent \textbf{(P\_2,1.69.2)} katham punaḥ anyasya prakarṣeṇa anyasya prakarṣaḥ syāt . \\
\noindent \textbf{(P\_2,1.69.2)} na eva anyasya prakarṣeṇa anyasya prakarṣeṇa bhavitavyam . \\
\noindent \textbf{(P\_2,1.69.2)} yathā eva ayam dravyeṣu yatate vastrāṇi me syuḥ iti evam guṇeṣu api yatate sūkṣmatarāṇi me syuḥ iti . \\
\noindent \textbf{(P\_2,1.69.2)} na atra ātiśāyikaḥ prāpnoti . \\
\noindent \textbf{(P\_2,1.69.2)} kim kāraṇam . \\
\noindent \textbf{(P\_2,1.69.2)} guṇavacanāt iti ucyate . \\
\noindent \textbf{(P\_2,1.69.2)} na ca samāsaḥ guṇavacanaḥ . \\
\noindent \textbf{(P\_2,1.69.2)} samāsaḥ api guṇavacanaḥ . \\
\noindent \textbf{(P\_2,1.69.2)} katham . \\
\noindent \textbf{(P\_2,1.69.2)} ajahatsvārthā vṛttiḥ . \\
\noindent \textbf{(P\_2,1.69.2)} atha jahatsvārthāyām tu doṣaḥ eva . \\
\noindent \textbf{(P\_2,1.69.2)} jahatsvāṛthāyām api na doṣaḥ . \\
\noindent \textbf{(P\_2,1.69.2)} bhavati bahuvrīhau tadguṇasaṃvijñānam api . \\
\noindent \textbf{(P\_2,1.69.2)} tat yathā . \\
\noindent \textbf{(P\_2,1.69.2)} śuklavāsasam ānaya . \\
\noindent \textbf{(P\_2,1.69.2)} lohitoṣṇīṣāḥ ṛtvijaḥ pracaranti iti . \\
\noindent \textbf{(P\_2,1.69.2)} tatguṇaḥ ānīyate tadguṇāḥ ca pracaranti . \\
\noindent \textbf{(P\_2,1.69.2)} uttarapadārthātiśaye ātiśāyikaḥ bahuvrīheḥ bahvāḍhyatarādyarthaḥ . \\
\noindent \textbf{(P\_2,1.69.2)} uttarapadārthātiśaye ātiśāyikaḥ bahuvrīheḥ bhavati vipratiṣedhena . \\
\noindent \textbf{(P\_2,1.69.2)} kim prayojanam . \\
\noindent \textbf{(P\_2,1.69.2)} bahvāḍhyatarādyarthaḥ . \\
\noindent \textbf{(P\_2,1.69.2)} bahvāḍhyataraḥ bahusukumārataraḥ . \\
\noindent \textbf{(P\_2,1.69.2)} kaḥ punaḥ atra viśeṣaḥ bahuvrīheḥ vā ātiśāyikaḥ syāt ātiśāyikāntena vā bahuvrīhiḥ . \\
\noindent \textbf{(P\_2,1.69.2)} svarakapoḥ viśeṣaḥ . \\
\noindent \textbf{(P\_2,1.69.2)} yadi atra ātiśāyikāt bahuvrīhiḥ syāt bahvāḍyataraḥ evam svaraḥ prasajyeta bahvāḍhyataraḥ iti ca iṣyate . \\
\noindent \textbf{(P\_2,1.69.2)} bahvāḍhyakataraḥ iti ca prāpnoti bahvāḍhyatarakaḥ iti ca iṣyate . \\
\noindent \textbf{(P\_2,1.69.2)} samānādhikaraṇādhikāre śākapārthivādīnām upasaṅkhyānam uttarapadalopaḥ ca . \\
\noindent \textbf{(P\_2,1.69.2)} samānādhikaraṇādhikāre śākapārthivādīnām upasaṅkhyānam kartavyam uttarapadalopaḥ ca vaktavyaḥ . \\
\noindent \textbf{(P\_2,1.69.2)} śākabhojī pārthivaḥ śākapārthivaḥ . \\
\noindent \textbf{(P\_2,1.69.2)} kutapavāsaḥ sauśrutaḥ kutapasauśrutraḥ . \\
\noindent \textbf{(P\_2,1.69.2)} ajāpaṇyaḥ taulvaliḥ ajātaulvaliḥ . \\
\noindent \textbf{(P\_2,1.69.2)} yaṣṭipradhānaḥ maudgalyaḥ yaṣṭimaudgalyaḥ . \\
\noindent \textbf{(P\_2,1.71)} catuṣpāt jātiḥ iti vaktavyam . \\
\noindent \textbf{(P\_2,1.71)} iha mā bhūt . \\
\noindent \textbf{(P\_2,1.71)} kālākṣīgarbhiṇī svastimatī garbhiṇī . \\
\noindent \textbf{(P\_2,1.72)} kimarthaḥ cakāraḥ . \\
\noindent \textbf{(P\_2,1.72)} evakārārthaḥ . \\
\noindent \textbf{(P\_2,1.72)} mayūravyaṃsakādayaḥ eva . \\
\noindent \textbf{(P\_2,1.72)} kva mā bhūt . \\
\noindent \textbf{(P\_2,1.72)} paramaḥ mayūravyaṃsakaḥ iti . \\
\noindent \textbf{(P\_2,2.2)} iha kasmāt na bhavati : grāmārdhaḥ , nagarārdhaḥ iti . \\
\noindent \textbf{(P\_2,2.2)} ardhaśabdasya napuṃsakaliṅgasya idam grahaṇam puṃliṅgaḥ ca ayam ardhaśabdaḥ . \\
\noindent \textbf{(P\_2,2.2)} kva punaḥ ayam napuṃsakaliṅgaḥ kva puṃliṅgaḥ . \\
\noindent \textbf{(P\_2,2.2)} samapravibhāge napuṃsakaliṅgaḥ , avayavavācī puṃliṅgaḥ . \\
\noindent \textbf{(P\_2,2.2)} iha kasmāt na bhavati : ardham pippalīnām iti . \\
\noindent \textbf{(P\_2,2.2)} na vā bhavati ardhapippalyaḥ iti . \\
\noindent \textbf{(P\_2,2.2)} bhavati yadā khaṇḍasamuccayaḥ : ardhapippalī ca ardhapippalī ca ardhapippalī ca ardhapippalyaḥ iti . \\
\noindent \textbf{(P\_2,2.2)} yada tu etat vākyam bhavati ardham pippalīnām iti tadā na bhavitavyam . \\
\noindent \textbf{(P\_2,2.2)} tadā kasmāt na bhavati . \\
\noindent \textbf{(P\_2,2.2)} ekādkhikaraṇe iti vartate . \\
\noindent \textbf{(P\_2,2.2)} na tarhi idānīm idam bhavati : ardharāśiḥ iti . \\
\noindent \textbf{(P\_2,2.2)} bhavati . \\
\noindent \textbf{(P\_2,2.2)} ekam etat adhikaraṇam yaḥ asau rāśiḥ nāma . \\
\noindent \textbf{(P\_2,2.3)} anyatarasyāṅgrahaṇam kimartham . \\
\noindent \textbf{(P\_2,2.3)} anyatarasyām samāsaḥ yathā syāt . \\
\noindent \textbf{(P\_2,2.3)} samāsena mukte vākyam api yathā syāt . \\
\noindent \textbf{(P\_2,2.3)} dvitīyam bhikṣāyāḥ iti . \\
\noindent \textbf{(P\_2,2.3)} na etat asti prayojanam . \\
\noindent \textbf{(P\_2,2.3)} prakṛtā mahāvibhāṣā . \\
\noindent \textbf{(P\_2,2.3)} tayā vākyam api bhaviṣyati . \\
\noindent \textbf{(P\_2,2.3)} idam tarhi prayojanam . \\
\noindent \textbf{(P\_2,2.3)} ekadeśisamāsena mukte ṣaṣṭhīsamāsaḥ api yathā syāt . \\
\noindent \textbf{(P\_2,2.3)} bhikṣādvitīyam iti . \\
\noindent \textbf{(P\_2,2.3)} etat api na asti prayojanam . \\
\noindent \textbf{(P\_2,2.3)} ayam api vibhāṣā ṣaṣṭhīsamāsaḥ api . \\
\noindent \textbf{(P\_2,2.3)} tau ubhau vacanāt bhaviṣyataḥ . \\
\noindent \textbf{(P\_2,2.3)} ataḥ uttaram paṭhati . \\
\noindent \textbf{(P\_2,2.3)} dvitīyādīnām vibhāṣāprakaraṇe vibhāṣāvacanam jñāpakam avayavavidhāne sāmānyavidhānābhāvasya . \\
\noindent \textbf{(P\_2,2.3)} dvitīyādīnām vibhāṣāprakaraṇe vibhāṣāvacanam kriyate jñāpārtham . \\
\noindent \textbf{(P\_2,2.3)} kim jñāpyate . \\
\noindent \textbf{(P\_2,2.3)} etat jñāpayati ācāryaḥ . \\
\noindent \textbf{(P\_2,2.3)} avayavavidhau sāmānyavidhiḥ na bhavati iti . \\
\noindent \textbf{(P\_2,2.3)} kim etasya jñāpane prayojanam . \\
\noindent \textbf{(P\_2,2.3)} bhinatti chinatti . \\
\noindent \textbf{(P\_2,2.3)} śnami kṛte śap na bhavati . \\
\noindent \textbf{(P\_2,2.3)} na etat asti prayojanam . \\
\noindent \textbf{(P\_2,2.3)} śabādeśāḥ śyanādayaḥ kariṣyante . \\
\noindent \textbf{(P\_2,2.3)} tat tarhi śapaḥ grahaṇam kartavyam . \\
\noindent \textbf{(P\_2,2.3)} na kartavyam .prakṛtam anuvartate . \\
\noindent \textbf{(P\_2,2.3)} kva prakṛtam . \\
\noindent \textbf{(P\_2,2.3)} kartari śap iti . \\
\noindent \textbf{(P\_2,2.3)} tat vai prathamānirdiṣṭam ṣaṣṭhīnirdiṣṭena ca iha arthaḥ . \\
\noindent \textbf{(P\_2,2.3)} rudhādibhyaḥ iti eṣā pañcamī śap iti prathamāyāḥ ṣaṣthīm prakalpayiṣyati tasmāt iti uttarasya iti . \\
\noindent \textbf{(P\_2,2.3)} pratyayavidhiḥ ayam na ca pratyayavidhau pañcamyaḥ prakalpikāḥ bhavanti . \\
\noindent \textbf{(P\_2,2.3)} na ayam pratyayavidhiḥ . \\
\noindent \textbf{(P\_2,2.3)} vihitaḥ pratyayaḥ prakṛtaḥ ca anuvartate . \\
\noindent \textbf{(P\_2,2.3)} evam tarhi jñāpayati ācāryaḥ yatra utsargāpavādam vibhāṣā tatra apavādena mukte utsargaḥ na bhavati iti . \\
\noindent \textbf{(P\_2,2.3)} kim etasya jñāpane prayojanam . \\
\noindent \textbf{(P\_2,2.3)} dikpūrvapadāt ṅīp . \\
\noindent \textbf{(P\_2,2.3)} prāṅmukhī prāṅmukhā pratyaṅmukhī pratyaṅmukhā . \\
\noindent \textbf{(P\_2,2.3)} ṅīpa mukte ṅīṣ na bhavati . \\
\noindent \textbf{(P\_2,2.3)} na etat asti prayojanam . \\
\noindent \textbf{(P\_2,2.3)} vakṣyati etat . \\
\noindent \textbf{(P\_2,2.3)} dikpūrvapadāt ṅīṣaḥ anudāttatvam ṅībvidhāne hi anyatra api ṅīṣviṣayāt ṅīpprasaṅgaḥ iti . \\
\noindent \textbf{(P\_2,2.3)} idam tarhi prayojanam . \\
\noindent \textbf{(P\_2,2.3)} ardhapippalī ardhakośātakī . \\
\noindent \textbf{(P\_2,2.3)} ekadeśisamāsena mukte ṣaṣṭhīsamāsaḥ na bhavati . \\
\noindent \textbf{(P\_2,2.3)} unmattagaṅgam lohitagaṅgam . \\
\noindent \textbf{(P\_2,2.3)} avyayībhāvena mukte bahuvrīhiḥ na bhavati . \\
\noindent \textbf{(P\_2,2.3)} dākṣiḥ plākṣiḥ . \\
\noindent \textbf{(P\_2,2.3)} iñā mukte aṇ na bhavati . \\
\noindent \textbf{(P\_2,2.3)} yadi etat jñāpyate upagoḥ apatyam aupagavaḥ . \\
\noindent \textbf{(P\_2,2.3)} taddhitena mukte upagvapatyam iti na sidhyati . \\
\noindent \textbf{(P\_2,2.3)} asti atra viśeṣaḥ . \\
\noindent \textbf{(P\_2,2.3)} dve hi atra vibhāṣā . \\
\noindent \textbf{(P\_2,2.3)} daivayajñiśaucivṛkṣisātyamugrikāṇṭheviddhibhyaḥ anyatarasyām iti samarthānām prathamāt vā iti ca . \\
\noindent \textbf{(P\_2,2.3)} tatra ekaya vṛttiḥ bhaviṣyati aparaya vṛttiviṣaye vibhāṣapavādaḥ . \\
\noindent \textbf{(P\_2,2.3)} kriyamāṇe api vai anyatarasyāṅgrahaṇe ṣaṣṭhīsamāsaḥ na prāpnoti . \\
\noindent \textbf{(P\_2,2.3)} kim kāraṇam . \\
\noindent \textbf{(P\_2,2.3)} pūraṇena iti pratiṣedhāt . \\
\noindent \textbf{(P\_2,2.3)} na etat pūraṇāntam . \\
\noindent \textbf{(P\_2,2.3)} anā etat paryavapannam . \\
\noindent \textbf{(P\_2,2.3)} etat api pūraṇāntam eva . \\
\noindent \textbf{(P\_2,2.3)} katha. pūraṇam nāma arthaḥ tam āha tīyaśabdaḥ . \\
\noindent \textbf{(P\_2,2.3)} ataḥ pūraṇam . \\
\noindent \textbf{(P\_2,2.3)} yaḥ asau pūraṇāntāt svārthe bhāge an saḥ api pūraṇam eva . \\
\noindent \textbf{(P\_2,2.3)} evam tarhi anyatarasyāṅgrahaṇasāmarthyāt ṣaṣṭhīsamāsaḥ api bhaviṣyati . \\
\noindent \textbf{(P\_2,2.4)} kimarthaḥ cakāraḥ . \\
\noindent \textbf{(P\_2,2.4)} anukaraṣaṇārthaḥ . \\
\noindent \textbf{(P\_2,2.4)} anyatarasyām iti etat anukṛṣyate . \\
\noindent \textbf{(P\_2,2.4)} kim prayojanam . \\
\noindent \textbf{(P\_2,2.4)} anyatarasyām samāsaḥ yathā syāt . \\
\noindent \textbf{(P\_2,2.4)} samāsena mukte vākyam api yatha syāt . \\
\noindent \textbf{(P\_2,2.4)} jīvikām prāptaḥ iti . \\
\noindent \textbf{(P\_2,2.4)} na etat asti prayojanam . \\
\noindent \textbf{(P\_2,2.4)} prakṛtā mahāvibhāṣā . \\
\noindent \textbf{(P\_2,2.4)} tayā vākyam bhaviṣyati . \\
\noindent \textbf{(P\_2,2.4)} idam tarhi prayojanam . \\
\noindent \textbf{(P\_2,2.4)} dvitīyāsamāsaḥ api yatha syāt . \\
\noindent \textbf{(P\_2,2.4)} jīvikāprāptaḥ iti . \\
\noindent \textbf{(P\_2,2.4)} etat api na asti prayojanam. ayam api ucyate dvitīyāsamāsaḥ api . \\
\noindent \textbf{(P\_2,2.4)} tat ubhayam vacanāt bhaviṣyati . \\
\noindent \textbf{(P\_2,2.4)} evam tarhi na ayam anukarṣaṇārthaḥ cakāraḥ . \\
\noindent \textbf{(P\_2,2.4)} kim tarhi . \\
\noindent \textbf{(P\_2,2.4)} atvam anena vidhīyate . \\
\noindent \textbf{(P\_2,2.4)} prāptāpanne dvitīyāntena saha samasyete atvam ca bhavati prāptapannayoḥ iti . \\
\noindent \textbf{(P\_2,2.4)} prāptā jīvikām prāptajīvikā āpannā jīvikām āpannajivikā . \\
\noindent \textbf{(P\_2,2.5.1)} kimpradhānaḥ ayam samāsaḥ . \\
\noindent \textbf{(P\_2,2.5.1)} uttarapadārthapradhānaḥ . \\
\noindent \textbf{(P\_2,2.5.1)} yadi uttarapadārthapradhānaḥ sadharmaṇā anena anyaiḥ uttarapadārthapradhānaiḥ bhavitavyam . \\
\noindent \textbf{(P\_2,2.5.1)} anyeṣu ca uttarapadārthapradhāneṣu yā eva asau antarvartinī vibhaktiḥ tasyāḥ samāse api śravaṇam bhavati : rājñaḥ puruṣaḥ rājapuruṣaḥ iti . \\
\noindent \textbf{(P\_2,2.5.1)} iha punaḥ vākye ṣaṣṭhī samāse prathamā . \\
\noindent \textbf{(P\_2,2.5.1)} kena etat evam bhavati . \\
\noindent \textbf{(P\_2,2.5.1)} yaḥ asau māsajātayoḥ abhisambandhaḥ saḥ samāse nivartate . \\
\noindent \textbf{(P\_2,2.5.1)} abhihitaḥ saḥ arthaḥ antarbhūtaḥ prātipadikārthaḥ sampannaḥ . \\
\noindent \textbf{(P\_2,2.5.1)} tatra prātipadikārthe prathamā iti prathamā bhavati . \\
\noindent \textbf{(P\_2,2.5.1)} na tarhi idānīm idam bhavati : māsajātasya iti . \\
\noindent \textbf{(P\_2,2.5.1)} bhavati . \\
\noindent \textbf{(P\_2,2.5.1)} bāhyam artham apekṣya ṣaṣṭhī . \\
\noindent \textbf{(P\_2,2.5.2)} kālasya yena samāsaḥ tasya aparimāṇitvāt anirdeśaḥ . \\
\noindent \textbf{(P\_2,2.5.2)} kālasya yena samāsaḥ saḥ aparimāṇī . \\
\noindent \textbf{(P\_2,2.5.2)} tasya aparimāṇitvāt anirdeśaḥ . \\
\noindent \textbf{(P\_2,2.5.2)} agamakaḥ nirdeśaḥ anirdeśaḥ . \\
\noindent \textbf{(P\_2,2.5.2)} na hi jātasya māsaḥ parimāṇam . \\
\noindent \textbf{(P\_2,2.5.2)} kasya tarhi . \\
\noindent \textbf{(P\_2,2.5.2)} triṃśadrātrasya . \\
\noindent \textbf{(P\_2,2.5.2)} tat yathā . \\
\noindent \textbf{(P\_2,2.5.2)} droṇaḥ badarāṇām devadattasya iti . \\
\noindent \textbf{(P\_2,2.5.2)} na devadattasya droṇaḥ parimāṇam . \\
\noindent \textbf{(P\_2,2.5.2)} kasya tarhi . \\
\noindent \textbf{(P\_2,2.5.2)} badarāṇām . \\
\noindent \textbf{(P\_2,2.5.2)} siddham tu kālaparimāṇam yasya sa kālaḥ tena . \\
\noindent \textbf{(P\_2,2.5.2)} siddham etat . \\
\noindent \textbf{(P\_2,2.5.2)} katham . \\
\noindent \textbf{(P\_2,2.5.2)} kālaparimāṇam yasya sa kālaḥ tena samasyate iti vaktavyam . \\
\noindent \textbf{(P\_2,2.5.2)} sidhyati . \\
\noindent \textbf{(P\_2,2.5.2)} sūtram tarhi bhidyate . \\
\noindent \textbf{(P\_2,2.5.2)} yathānyāsam eva astu . \\
\noindent \textbf{(P\_2,2.5.2)} nanu ca uktam kālasya yena samāsaḥ tasya aparimāṇitvāt anirdeśaḥ iti . \\
\noindent \textbf{(P\_2,2.5.2)} kam punaḥ kālam matvā bhavān āha kālasya yena samāsaḥ tasya aparimāṇitvāt anirdeśaḥ iti . \\
\noindent \textbf{(P\_2,2.5.2)} yena mūrtīnām upacayāḥ ca apacayāḥ ca lakṣyante tam kālam āhuḥ . \\
\noindent \textbf{(P\_2,2.5.2)} tasya eva kayā cit kriyayā yuktasya ahaḥ iti ca bhavati rātriḥ iti ca . \\
\noindent \textbf{(P\_2,2.5.2)} kayā kriyayā . \\
\noindent \textbf{(P\_2,2.5.2)} ādityagatyā . \\
\noindent \textbf{(P\_2,2.5.2)} tayā eva asakṛt āvṛttayā māsaḥ iti bhavati saṃvatsaraḥ iti ca . \\
\noindent \textbf{(P\_2,2.5.2)} yadi evam jātasya māsaḥ parimāṇam . \\
\noindent \textbf{(P\_2,2.5.2)} ekavacanadvigoḥ ca upasaṅkhyānam . ekavacanāntānām iti vaktavyam . \\
\noindent \textbf{(P\_2,2.5.2)} iha mā bhūt . \\
\noindent \textbf{(P\_2,2.5.2)} māsau jātasya māsāḥ jātasya iti . \\
\noindent \textbf{(P\_2,2.5.2)} dvigoḥ ca iti vaktavyam iha api yathā syāt : dvimāsajātaḥ , trimāsajātaḥ . \\
\noindent \textbf{(P\_2,2.5.2)} uktam vā . \\
\noindent \textbf{(P\_2,2.5.2)} kim uktam . \\
\noindent \textbf{(P\_2,2.5.2)} ekavacane tāvat uktam anabhidhānāt iti . \\
\noindent \textbf{(P\_2,2.5.2)} dvigoḥ kim uktam . \\
\noindent \textbf{(P\_2,2.5.2)} uttarapadena parimāṇina dvigoḥ samāsavacanam iti . \\
\noindent \textbf{(P\_2,2.6)} kimpradhānaḥ ayam samāsaḥ . \\
\noindent \textbf{(P\_2,2.6)} uttarapadārthapradhānaḥ . \\
\noindent \textbf{(P\_2,2.6)} yadi uttarapadārthapradhānaḥ abrāhmaṇam ānaya iti ukte brāhmaṇamātrasya ānayanam prāpnoti . \\
\noindent \textbf{(P\_2,2.6)} anyapadārthapradhānaḥ tarhi bhaviṣyati . \\
\noindent \textbf{(P\_2,2.6)} yadi anyapadārthapradhānaḥ , avarṣā hemantaḥ iti hemantasya yat liṅgam vacanam ca tat samāsasya api prāpnoti . \\
\noindent \textbf{(P\_2,2.6)} pūrvapadārthapradhānaḥ tarhi bhaviṣyati . \\
\noindent \textbf{(P\_2,2.6)} yadi pūrvapadārthapradhānaḥ avyayasañjñā prāpnoti : avyayam hi asya pūrvapadam iti . \\
\noindent \textbf{(P\_2,2.6)} na eṣaḥ doṣaḥ . \\
\noindent \textbf{(P\_2,2.6)} pāṭhena avyayasañjñā kriyate . \\
\noindent \textbf{(P\_2,2.6)} na ca nañsamāsaḥ tatra paṭhyate . \\
\noindent \textbf{(P\_2,2.6)} yadi api nañsamāsaḥ na paṭhyate nañ tu paṭhyate . \\
\noindent \textbf{(P\_2,2.6)} pāṭhena api avyayasañjñāyām satyām abhideheyavat liṅgavacanāni bhavanti . \\
\noindent \textbf{(P\_2,2.6)} yaḥ ca iha arthaḥ abhidhīyate na tasya liṅgasaṅkhyābhyām yogaḥ asti . \\
\noindent \textbf{(P\_2,2.6)} na idam vācanikam aliṅgatā asaṅkhyatā va . \\
\noindent \textbf{(P\_2,2.6)} kim tarhi . \\
\noindent \textbf{(P\_2,2.6)} svābhāvikam etat . \\
\noindent \textbf{(P\_2,2.6)} tat yathā: samānam īhamānānām adhīyānānām ca ke cit arthaiḥ yujyante apare na . \\
\noindent \textbf{(P\_2,2.6)} na ca idānīm kaḥ cit arthavān iti kṛtvā sarvaiḥ arthavadbhiḥ śakyam bhavitum kaḥ cit anarthakaḥ iti kṛtvā sarvaiḥ anarthakaiḥ . \\
\noindent \textbf{(P\_2,2.6)} tatra kim asmābhiḥ śakyam kartum . \\
\noindent \textbf{(P\_2,2.6)} yat nañaḥ prāk samāsāt liṅgasaṅkhyābhyām yogaḥ na asti samāse ca bhavati svābhāvikam etat . \\
\noindent \textbf{(P\_2,2.6)} atha vā āśrayataḥ liṅgavacanāni bhaviṣyanti . \\
\noindent \textbf{(P\_2,2.6)} guṅavacanānām hi śabdānām āśrayataḥ liṅgavacanāni bhavanti . \\
\noindent \textbf{(P\_2,2.6)} tat yathā : śuklam vastram , śuklā śāṭī śuklaḥ kambalaḥ , śuklau kambalau śuklāḥ kambalāḥ iti . \\
\noindent \textbf{(P\_2,2.6)} yat asau dravyam śritaḥ bhavati guṇaḥ tasya yat liṅgam vacanam ca tat guṇasya api bhavati . \\
\noindent \textbf{(P\_2,2.6)} evam iha api yat asau dravyam śritaḥ bhavati samāsaḥ tasya yat liṅgam vacanam ca tat samāsasya api bhaviṣyati . \\
\noindent \textbf{(P\_2,2.6)} atha vā punaḥ astu uttarapadārthapradhānaḥ . \\
\noindent \textbf{(P\_2,2.6)} nanu ca uktam abrāhmaṇam ānaya iti ukte brāhmaṇamātrasya ānayanam prāpnoti iti . \\
\noindent \textbf{(P\_2,2.6)} na eṣaḥ doṣaḥ . \\
\noindent \textbf{(P\_2,2.6)} idam tāvat ayam praṣṭavyaḥ : atha iha rājapuruṣam ānaya iti ukte puruṣamātrasya ānayanam kasmāt na bhavati . \\
\noindent \textbf{(P\_2,2.6)} asti atra viśeṣaḥ . \\
\noindent \textbf{(P\_2,2.6)} rājā viśeṣakaḥ prayujyate . \\
\noindent \textbf{(P\_2,2.6)} tena viśiṣṭasya ānayanam bhavati . \\
\noindent \textbf{(P\_2,2.6)} iha api tarhi nañ viśeṣakaḥ prayujyate . \\
\noindent \textbf{(P\_2,2.6)} tena nañviśiṣṭasya ānayanam bhaviṣyati . \\
\noindent \textbf{(P\_2,2.6)} kaḥ punaḥ asau . \\
\noindent \textbf{(P\_2,2.6)} nivṛttapadārthakaḥ . \\
\noindent \textbf{(P\_2,2.6)} yadā punaḥ asya padārthaḥ nivartate kim svābhāvikī nivṛttiḥ āhosvit vācanikī . \\
\noindent \textbf{(P\_2,2.6)} kim ca ataḥ . \\
\noindent \textbf{(P\_2,2.6)} yadi svābhāvikī kim nañ prayujyamānaḥ karoti . \\
\noindent \textbf{(P\_2,2.6)} atha vācanikī tat vaktavyam : nañ prayujyamānaḥ padārtham nivartayati iti . \\
\noindent \textbf{(P\_2,2.6)} evam tarhi svābhāvikī nivṛttiḥ . \\
\noindent \textbf{(P\_2,2.6)} nanu ca uktam kim nañ prayujyamānaḥ karoti iti . \\
\noindent \textbf{(P\_2,2.6)} nañ prayujyamānaḥ padārtham nivartayati katham . \\
\noindent \textbf{(P\_2,2.6)} kīlapratikīlavat . \\
\noindent \textbf{(P\_2,2.6)} tat yathā kīlaḥ āhanyamānaḥ pratikīlam nirhanti . \\
\noindent \textbf{(P\_2,2.6)} yadi etat nañaḥ māhātmyam syāt na jātu cit rājānaḥ hastyaśvam bibhṛyuḥ . \\
\noindent \textbf{(P\_2,2.6)} na iti eva rājānaḥ brūyuḥ . \\
\noindent \textbf{(P\_2,2.6)} evam tarhi svābhāvikī nivṛttiḥ . \\
\noindent \textbf{(P\_2,2.6)} nanu ca uktam kim nañ prayujyamānaḥ karoti iti . \\
\noindent \textbf{(P\_2,2.6)} nañnimittā tu upalabdhiḥ . \\
\noindent \textbf{(P\_2,2.6)} tat yathā samandhakāre dravyāṇām samavasthitānām pradīpnimittam darśanam na ca teṣām pradīpaḥ nirvartakaḥ bhavati . \\
\noindent \textbf{(P\_2,2.6)} yadi punaḥ ayam nivṛttapadārthakaḥ kimartham brāhmaṇaśabdaḥ prayujyate . \\
\noindent \textbf{(P\_2,2.6)} evam yathā vijñayeta asya padārthaḥ nivartate iti . \\
\noindent \textbf{(P\_2,2.6)} na iti hi ukte sandehaḥ syāt kasya padārthaḥ nivartate iti . \\
\noindent \textbf{(P\_2,2.6)} tatra asandehārtham brāhmaṇaśabdaḥ prayujyate . \\
\noindent \textbf{(P\_2,2.6)} evam vā etat . \\
\noindent \textbf{(P\_2,2.6)} atha vā sarve ete śabdāḥ guṇasamudāyeṣu vartante brāhmaṇaḥ kṣatriyaḥ vaiśyaḥ śūdraḥ iti . tapaḥ śrutam ca yoniḥ ca iti etad brāhmaṇakārakam . \\
\noindent \textbf{(P\_2,2.6)} tapaḥśtrutābhyām yaḥ hīnaḥ jātibrāhmaṇaḥ eva saḥ . tathā gauraḥ śucyācāraḥ piṅgalaḥ kapilakeśaḥ iti etān api abhyantarān brāhmaṇye guṇān kurvanti . \\
\noindent \textbf{(P\_2,2.6)} samudāyeṣu ca vṛttāḥ śabdāḥ avayaveṣu api vartante. tad yathā . \\
\noindent \textbf{(P\_2,2.6)} pūrve pañcālāḥ uttare pañcālāḥ tailam bhuktam ghṛtaṃ bhuktam śuklaḥ nīlaḥ kapilaḥ kṛṣṇaḥ iti . \\
\noindent \textbf{(P\_2,2.6)} evam ayam samudāye brāhmaṇaśabdaḥ pravṛttaḥ avayaveṣu api vartate jātihīne guṇahīne ca . \\
\noindent \textbf{(P\_2,2.6)} guṇahīne tāvat. abrāhmaṇaḥ ayam yaḥ tiṣṭhan mūtrayati . \\
\noindent \textbf{(P\_2,2.6)} abrāhmaṇaḥ ayaṃ yaḥ gacchan bhakṣayati . \\
\noindent \textbf{(P\_2,2.6)} jātihīne sandehāt durupadeśāt ca brāhmaṇaśabdaḥ vartate . \\
\noindent \textbf{(P\_2,2.6)} sandehāt tāvat : gauram śucyācāraṃ piṅgalam kapilakeśam dṛṣṭvā adhyavasyati brāhmaṇaḥ ayam iti . \\
\noindent \textbf{(P\_2,2.6)} tataḥ paścāt upalabhate na ayaṃ brāhmaṇaḥ abrāhmaṇaḥ ayam iti . \\
\noindent \textbf{(P\_2,2.6)} tatra sandehāt ca brāhmaṇaśabdaḥ vartate jātikṛtā ca arthasya nivṛttiḥ . \\
\noindent \textbf{(P\_2,2.6)} durupadeśāt : durupadiṣṭam asya bhavati amuṣmin avakāśe brāhmaṇaḥ tam ānaya iti . \\
\noindent \textbf{(P\_2,2.6)} sa tatra gatvā yam paśyati tam adhyavasyati brāhmaṇaḥ ayam iti . \\
\noindent \textbf{(P\_2,2.6)} tataḥ paścāt upalabhate na ayaṃ brāhmaṇaḥ abrāhmaṇaḥ ayam iti . \\
\noindent \textbf{(P\_2,2.6)} tatra durupadeśāt ca brāhmaṇaśabdaḥ vartate jātikṛtā ca arthasya nivṛttiḥ . \\
\noindent \textbf{(P\_2,2.6)} ātaḥ ca sandehāt durupadeśāt vā . \\
\noindent \textbf{(P\_2,2.6)} na hi ayam kālam māṣarāśivarṇam āpaṇe āsīnam dṛṣṭvā adhyavasyati brāhmaṇaḥ ayam iti . \\
\noindent \textbf{(P\_2,2.6)} nirjñātam tasya bhavati . \\
\noindent \textbf{(P\_2,2.6)} idam khalu api bhūyaḥ uttarapadārthaprādhānye sati saṅgṛhītam bhavati . \\
\noindent \textbf{(P\_2,2.6)} kim . \\
\noindent \textbf{(P\_2,2.6)} anekam iti . \\
\noindent \textbf{(P\_2,2.6)} kim atra saṅgṛhītam . \\
\noindent \textbf{(P\_2,2.6)} ekavacanam . \\
\noindent \textbf{(P\_2,2.6)} katham punaḥ ekasya pratiṣedhena anekasya sampratyayaḥ syāt . \\
\noindent \textbf{(P\_2,2.6)} prasajya ayam kriyāguṇau tataḥ paścāt nivṛttim karoti . \\
\noindent \textbf{(P\_2,2.6)} tat yathā : āsaya śāyaya bhojaya anekam iti . \\
\noindent \textbf{(P\_2,2.6)} yadi api tāvat atra etat śakyate vaktum yatra kriyāguṇau prasajyete yatra khalu na prasajyete tatra katham : anekaḥ tiṣṭhati iti . \\
\noindent \textbf{(P\_2,2.6)} bhavati ca evañjātīyakānām api ekasya pratiṣedhena bahūnām sampratyayaḥ . \\
\noindent \textbf{(P\_2,2.6)} tat yathā na naḥ ekam priyam na naḥ ekam sukham iti . \\
\noindent \textbf{(P\_2,2.6)} iha abrāhmaṇatvam abrāhmaṇatā paratvāt tvatalau prāpnutaḥ . \\
\noindent \textbf{(P\_2,2.6)} tatra kaḥ doṣaḥ . \\
\noindent \textbf{(P\_2,2.6)} svare hi doṣaḥ syāt . \\
\noindent \textbf{(P\_2,2.6)} abrāhmaṇatvam iti evam svaraḥ prasajyeta . \\
\noindent \textbf{(P\_2,2.6)} abrāhmaṇatvam iti ca iṣyate . \\
\noindent \textbf{(P\_2,2.6)} nañsamāse bhāvavacane uktam . \\
\noindent \textbf{(P\_2,2.6)} kim uktam . \\
\noindent \textbf{(P\_2,2.6)} tvatalbhyām nañsamāsaḥ pūrvapratiṣiddham svarsiddhyartham iti . \\
\noindent \textbf{(P\_2,2.7)} īṣat guṇavacanena .īṣat guṇavacanena iti vaktavyam . \\
\noindent \textbf{(P\_2,2.7)} akṛtā iti ucyamāne iha ca prasajyeta . \\
\noindent \textbf{(P\_2,2.7)} īṣat gārgyaḥ iti . \\
\noindent \textbf{(P\_2,2.7)} iha ca na syāt . \\
\noindent \textbf{(P\_2,2.7)} īṣatkaḍāraḥ . \\
\noindent \textbf{(P\_2,2.8)} kṛdyogā ca . \\
\noindent \textbf{(P\_2,2.8)} kṛdyogā ca ṣaṣṭhī samasyate iti vaktavyam . \\
\noindent \textbf{(P\_2,2.8)} idhmapravraścanaḥ palāśaśātanaḥ . \\
\noindent \textbf{(P\_2,2.8)} kimartham idam ucyate . \\
\noindent \textbf{(P\_2,2.8)} pratipadavidhānā ca ṣaṣṭhī na samasyate iti vakṣyati . \\
\noindent \textbf{(P\_2,2.8)} tasya ayam purastāt apakarṣaḥ . \\
\noindent \textbf{(P\_2,2.8)} kā punaḥ ṣaṣṭhīpratipadavidhānā kā kṛdyogā . \\
\noindent \textbf{(P\_2,2.8)} sarvā ṣaṣṭhī pratipadavidhānā śeṣalakṣaṇām varjayitvā . \\
\noindent \textbf{(P\_2,2.8)} kartṛkarmaṇoḥ kṛti iti yā ṣaṣṭhī sā kṛdyogā . \\
\noindent \textbf{(P\_2,2.8)} tatsthaiḥ ca guṇaiḥ . \\
\noindent \textbf{(P\_2,2.8)} tatsthaiḥ ca guṇaiḥ ṣaṣthīguṇaiḥ ṣaṣṭhī samasyate iti vaktavyam . \\
\noindent \textbf{(P\_2,2.8)} brāhmaṇavarṇaḥ candanagandhaḥ paṭahaśabdaḥ nadhīghoṣaḥ . \\
\noindent \textbf{(P\_2,2.8)} na tu tadviśeṣaṇaiḥ . \\
\noindent \textbf{(P\_2,2.8)} na tu tadviśeṣaṇaiḥ iti vaktavyam . \\
\noindent \textbf{(P\_2,2.8)} iha mā bhūt . \\
\noindent \textbf{(P\_2,2.8)} ghṛtasya tīvraḥ candanasya mṛduḥ iti . \\
\noindent \textbf{(P\_2,2.8)} kimartham idam ucyate . \\
\noindent \textbf{(P\_2,2.8)} guṇena iti pratiṣedham vakṣyati . \\
\noindent \textbf{(P\_2,2.8)} tasya ayam purastāt apakarṣaḥ . \\
\noindent \textbf{(P\_2,2.8)} kim kāraṇam guṇena na iti ucyate na punaḥ guṇavacanena iti ucyate . \\
\noindent \textbf{(P\_2,2.8)} na evam śakyam . \\
\noindent \textbf{(P\_2,2.8)} iha hi na syāt . \\
\noindent \textbf{(P\_2,2.8)} kākasya kārṣṇyam kaṇṭakasya taikṣṇyam balākāyāḥ śauklyam iti . \\
\noindent \textbf{(P\_2,2.8)} etat eva tasmin yoge udāharaṇam . \\
\noindent \textbf{(P\_2,2.8)} yat vai brāhmaṇasya śuklāḥ vṛṣalasya kṛṣṇāḥ iti asāmarthyāt atra na bhaviṣyati . \\
\noindent \textbf{(P\_2,2.8)} katham asāmarthyam . \\
\noindent \textbf{(P\_2,2.8)} sāpekṣam asamartham bhavati iti . \\
\noindent \textbf{(P\_2,2.8)} dravyam atra apekṣyate dantāḥ . \\
\noindent \textbf{(P\_2,2.8)} tasmāt guṇena na iti vaktavyam . \\
\noindent \textbf{(P\_2,2.8)} guṇena na iti ucyamāne tatsthaiḥ ca guṇaiḥ iti vaktavyam . \\
\noindent \textbf{(P\_2,2.8)} tatsthaiḥ ca guṇaiḥ iti ucyamāne na tu tadviśeṣaṇaiḥ iti vaktavyam . \\
\noindent \textbf{(P\_2,2.10)} pratipadavidhānā ca . \\
\noindent \textbf{(P\_2,2.10)} pratipadavidhānā ca ṣaṣṭhī na samasyate iti vaktavyam . \\
\noindent \textbf{(P\_2,2.10)} iha mā bhūt . \\
\noindent \textbf{(P\_2,2.10)} sarpiṣaḥ jñānam madhunaḥ jñānam iti . \\
\noindent \textbf{(P\_2,2.11)} guṇe kim udāharaṇam . \\
\noindent \textbf{(P\_2,2.11)} brāhmaṇasya śuklāḥ vṛṣalasya kṛṣṇāḥ iti . \\
\noindent \textbf{(P\_2,2.11)} na etat asti prayojanam . \\
\noindent \textbf{(P\_2,2.11)} asāmarthyāt atra na bhaviṣyati . \\
\noindent \textbf{(P\_2,2.11)} katham asāmarthyam . \\
\noindent \textbf{(P\_2,2.11)} sāpekṣam asamartham bhavati iti . \\
\noindent \textbf{(P\_2,2.11)} dravyam atra apekṣyate dantāḥ . \\
\noindent \textbf{(P\_2,2.11)} idam tarhi kākasya kārṣṇyam kaṇṭakasya taikṣṇyam balākāyāḥ śauklyam iti . \\
\noindent \textbf{(P\_2,2.11)} idam api udāharaṇam brāhmaṇasya śuklāḥ vṛṣalasya kṛṣṇāḥ iti . \\
\noindent \textbf{(P\_2,2.11)} nanu ca uktam . \\
\noindent \textbf{(P\_2,2.11)} asāmarthyāt atra na bhaviṣyati . \\
\noindent \textbf{(P\_2,2.11)} katham asāmarthyam . \\
\noindent \textbf{(P\_2,2.11)} sāpekṣam asamartham bhavati iti . \\
\noindent \textbf{(P\_2,2.11)} dravyam atra apekṣyate dantāḥ iti . \\
\noindent \textbf{(P\_2,2.11)} na eṣaḥ doṣaḥ . \\
\noindent \textbf{(P\_2,2.11)} bhavati vai kasya cit arthāt prakaraṇāt vā apekṣyam nirjñātam tadā vṛttiḥ prāpnoti . \\
\noindent \textbf{(P\_2,2.11)} sati kim udāharaṇam . \\
\noindent \textbf{(P\_2,2.11)} brahmaṇasya pakṣyan brāhmaṇasya pakṣyamāṇaḥ . \\
\noindent \textbf{(P\_2,2.11)} na etat asti . \\
\noindent \textbf{(P\_2,2.11)} pratiṣidhyate atra ṣaṣṭhī laprayoge na iti . \\
\noindent \textbf{(P\_2,2.11)} yā ca śrūyate eṣā bāhyam artham apekṣya bhavati . \\
\noindent \textbf{(P\_2,2.11)} tatra asmārthyāt na bhaviṣyati . \\
\noindent \textbf{(P\_2,2.11)} katham asāmarthyam . \\
\noindent \textbf{(P\_2,2.11)} sāpekṣam asamartham bhavati iti . \\
\noindent \textbf{(P\_2,2.11)} dravyam atra apekṣyate odanaḥ . \\
\noindent \textbf{(P\_2,2.11)} idam tarhi caurasya dviṣan vṛṣalasya dviṣan . \\
\noindent \textbf{(P\_2,2.11)} nanu ca atra api pratiṣidhyate . \\
\noindent \textbf{(P\_2,2.11)} vakṣyati etat dviṣaḥ śatuḥ vāvacanam iti . \\
\noindent \textbf{(P\_2,2.11)} avyaye kim udāharaṇam . \\
\noindent \textbf{(P\_2,2.11)} brāhmaṇasya uccaiḥ vṛṣalasya nīcaiḥ iti . \\
\noindent \textbf{(P\_2,2.11)} na etat asti . \\
\noindent \textbf{(P\_2,2.11)} asāmarthyāt atra na bhaviṣyati . \\
\noindent \textbf{(P\_2,2.11)} katham asāmarthyam . \\
\noindent \textbf{(P\_2,2.11)} sāpekṣam asamartham bhavati iti . \\
\noindent \textbf{(P\_2,2.11)} dravyam atra apekṣyate āsanam . \\
\noindent \textbf{(P\_2,2.11)} idan tarhi brāhmaṇasya kṛṭvā vṛṣalasya kṛtvā iti . \\
\noindent \textbf{(P\_2,2.11)} etat api na asti . \\
\noindent \textbf{(P\_2,2.11)} pratiṣidhyate tatra ṣaṣṭhī avyayaprayoge na iti . \\
\noindent \textbf{(P\_2,2.11)} yā ca śrūyate eṣā bāhyam artham apekṣya bhavati . \\
\noindent \textbf{(P\_2,2.11)} tatra asmārthyāt na bhaviṣyati . \\
\noindent \textbf{(P\_2,2.11)} katham asāmarthyam . \\
\noindent \textbf{(P\_2,2.11)} sāpekṣam asamartham bhavati iti . \\
\noindent \textbf{(P\_2,2.11)} dravyam atra apekṣyate kaṭaḥ . \\
\noindent \textbf{(P\_2,2.11)} idam tarhi . \\
\noindent \textbf{(P\_2,2.11)} purā sūryasya udetoḥ ādheyaḥ . \\
\noindent \textbf{(P\_2,2.11)} purā vatsānām apākartoḥ . \\
\noindent \textbf{(P\_2,2.11)} nanu ca atra api pratiṣidhyate avyayam iti kṛtvā . \\
\noindent \textbf{(P\_2,2.11)} vakṣyati etat . \\
\noindent \textbf{(P\_2,2.11)} avyayapratiṣedhe tosunkasunoḥ apratiṣedhaḥ iti . \\
\noindent \textbf{(P\_2,2.11)} samānādhikaraṇe kim udāharaṇam . \\
\noindent \textbf{(P\_2,2.11)} rājñaḥ pāṭaliputrakasya śukasya mārāvidasya pāṇineḥ sūtrakārasya . \\
\noindent \textbf{(P\_2,2.11)} na etat asti . \\
\noindent \textbf{(P\_2,2.11)} asāmarthyāt atra na bhaviṣyati . \\
\noindent \textbf{(P\_2,2.11)} katham asāmarthyam . \\
\noindent \textbf{(P\_2,2.11)} samānādhikaraṇam asamarthavat bhavati iti . \\
\noindent \textbf{(P\_2,2.11)} idam tarhi . \\
\noindent \textbf{(P\_2,2.11)} sarpiṣaḥ pīyamānaḥ yajuṣaḥ kriyamāṇasya iti . \\
\noindent \textbf{(P\_2,2.11)} nanu ca atra api asāmarthyāt eva na bhaviṣyati . \\
\noindent \textbf{(P\_2,2.11)} katham asāmarthyam . \\
\noindent \textbf{(P\_2,2.11)} samānādhikaraṇam asamarthavat bhavati iti . \\
\noindent \textbf{(P\_2,2.11)} adhātvabhihitam iti evam tat . \\
\noindent \textbf{(P\_2,2.14)} katham idam vijñāyate karmaṇi yā ṣaṣṭhī sā na samasyate iti āhosvit karmaṇi yaḥ ktaḥ iti . \\
\noindent \textbf{(P\_2,2.14)} kutaḥ sandehaḥ . \\
\noindent \textbf{(P\_2,2.14)} ubhayam prakṛtam . \\
\noindent \textbf{(P\_2,2.14)} tatra anyatarat śakyam viśeṣayitum . \\
\noindent \textbf{(P\_2,2.14)} kaḥ ca atra viśeṣaḥ . \\
\noindent \textbf{(P\_2,2.14)} karmaṇi iti ṣaṣṭhīnirdeśaḥ cet akartari kṛtā samāsavacanam . \\
\noindent \textbf{(P\_2,2.14)} karmaṇi iti ṣaṣṭhīnirdeśaḥ cet akartari kṛtā samāsaḥ vaktavyaḥ . \\
\noindent \textbf{(P\_2,2.14)} idhmapravraścanaḥ palāśaśātanaḥ . \\
\noindent \textbf{(P\_2,2.14)} tṛkakābhyam ca anarthakaḥ pratiṣedhaḥ . \\
\noindent \textbf{(P\_2,2.14)} tṛkakābhyam ca anarthakaḥ pratiṣedhaḥ . \\
\noindent \textbf{(P\_2,2.14)} apām sraṣṭā . \\
\noindent \textbf{(P\_2,2.14)} karmaṇi iti eva siddham . \\
\noindent \textbf{(P\_2,2.14)} astu tarhi karmaṇi yaḥ ktaḥ iti . \\
\noindent \textbf{(P\_2,2.14)} kim udāharaṇam . \\
\noindent \textbf{(P\_2,2.14)} brāhmaṇasya bhuktam vṛṣalasya pītam iti . \\
\noindent \textbf{(P\_2,2.14)} ktanirdeśe asamarthatvāt apratiṣedhaḥ . \\
\noindent \textbf{(P\_2,2.14)} ktanirdeśe asamarthatvāt apratiṣedhaḥ . \\
\noindent \textbf{(P\_2,2.14)} anarthakaḥ pratiṣedhaḥ apratiṣedhaḥ . \\
\noindent \textbf{(P\_2,2.14)} samāsaḥ kasmāt na bhavati . \\
\noindent \textbf{(P\_2,2.14)} asāmarthyāt . \\
\noindent \textbf{(P\_2,2.14)} katham asāmarthyam . \\
\noindent \textbf{(P\_2,2.14)} sāpekṣam asamartham bhavati iti . \\
\noindent \textbf{(P\_2,2.14)} dravyam atra apekṣyate odanaḥ . \\
\noindent \textbf{(P\_2,2.14)} pratiṣedhyam iti cet kartari api pratiṣedhaḥ .atha evam sati pratiṣedhaḥ kartavyaḥ iti dṛśyate kartari api pratiṣedhaḥ vaktavyaḥ syāt . \\
\noindent \textbf{(P\_2,2.14)} brāhmaṇasya gataḥ brāhmaṇasya yātaḥ iti . \\
\noindent \textbf{(P\_2,2.14)} pūjāyām ca pratiṣedhānarthakyam . \\
\noindent \textbf{(P\_2,2.14)} pūjāyām ca pratiṣedhaḥ anarthaḥ . \\
\noindent \textbf{(P\_2,2.14)} rājñām pūjitaḥ . \\
\noindent \textbf{(P\_2,2.14)} karmaṇi iti eva siddham . \\
\noindent \textbf{(P\_2,2.14)} tasmāt ubhayaprāptau karmaṇi ṣaṣṭhyāḥ pratiṣedhaḥ . \\
\noindent \textbf{(P\_2,2.14)} tasmāt ubhayaprāptau karmaṇi iti evam yā ṣaṣṭhī tasyāḥ pratiṣedhaḥ vaktavyaḥ . \\
\noindent \textbf{(P\_2,2.14)} saḥ tarhi vaktavyaḥ . \\
\noindent \textbf{(P\_2,2.14)} na vaktavyaḥ ityarthe ayam caḥ paṭhitaḥ . \\
\noindent \textbf{(P\_2,2.14)} kamaṇi ca . \\
\noindent \textbf{(P\_2,2.14)} karmaṇi iti evam yā ṣaṣṭhī iti . \\
\noindent \textbf{(P\_2,2.17)} kim iha nityagrahaṇena abhisambadhyate vidhiḥ āhosvit pratiṣedhaḥ . \\
\noindent \textbf{(P\_2,2.17)} vidhiḥ iti āha . \\
\noindent \textbf{(P\_2,2.17)} kutaḥ etat . \\
\noindent \textbf{(P\_2,2.17)} vidhiḥ hi vibhāṣā nityaḥ pratiṣedhaḥ . \\
\noindent \textbf{(P\_2,2.18)} prādiprasaṅge karmapravacanīyapratiṣedhaḥ . prādiprasaṅge karmapravacanīyānām pratiṣedhaḥ vaktavyaḥ . \\
\noindent \textbf{(P\_2,2.18)} vṛkṣam prati vidyotate vidyut . \\
\noindent \textbf{(P\_2,2.18)} sādhuḥ devadattaḥ mātaram prati . \\
\noindent \textbf{(P\_2,2.18)} vyavetapratiṣedhaḥ ca . \\
\noindent \textbf{(P\_2,2.18)} vyavetānām ca pratiṣedhaḥ vaktavyaḥ . \\
\noindent \textbf{(P\_2,2.18)} a mandraiḥ indra haribhiḥ yahi mayuraromabhiḥ . \\
\noindent \textbf{(P\_2,2.18)} siddham tu kvāṅsvatidurgativacanāt . \\
\noindent \textbf{(P\_2,2.18)} siddham etat . \\
\noindent \textbf{(P\_2,2.18)} katham . \\
\noindent \textbf{(P\_2,2.18)} kvāṅsvatidurgatayaḥ samasyante iti vaktavyam . \\
\noindent \textbf{(P\_2,2.18)} ku . \\
\noindent \textbf{(P\_2,2.18)} kubrāhmaṇaḥ kuvṛṣalaḥ . \\
\noindent \textbf{(P\_2,2.18)} āṅ . \\
\noindent \textbf{(P\_2,2.18)} ākaḍāraḥ āpiṅgalaḥ . \\
\noindent \textbf{(P\_2,2.18)} su . \\
\noindent \textbf{(P\_2,2.18)} subrāhmaṇaḥ suvṛṣalaḥ . \\
\noindent \textbf{(P\_2,2.18)} at . \\
\noindent \textbf{(P\_2,2.18)} atibrāhmaṇaḥ ativṛṣalaḥ . \\
\noindent \textbf{(P\_2,2.18)} dur . \\
\noindent \textbf{(P\_2,2.18)} durbrāhmaṇaḥ . \\
\noindent \textbf{(P\_2,2.18)} gati . \\
\noindent \textbf{(P\_2,2.18)} prakārakaḥ praṇāyakaḥ prasecakaḥ ūrīkṛtya ūrīkṛtam . \\
\noindent \textbf{(P\_2,2.18)} prādayaḥ ktārthe . \\
\noindent \textbf{(P\_2,2.18)} prādayaḥ ktārthe samasyante iti vaktavyam . \\
\noindent \textbf{(P\_2,2.18)} pragataḥ ācāryaḥ prācāryaḥ prāntevāsī prapitāmahaḥ . \\
\noindent \textbf{(P\_2,2.18)} etat eva ca saunāgaiḥ vistaratarakeṇa paṭhitam . \\
\noindent \textbf{(P\_2,2.18)} svatī pūjāyām . \\
\noindent \textbf{(P\_2,2.18)} svatīpūjāyām iti vaktavyam . \\
\noindent \textbf{(P\_2,2.18)} surājā atirājā . \\
\noindent \textbf{(P\_2,2.18)} duḥ nindāyām . \\
\noindent \textbf{(P\_2,2.18)} duḥ nindāyām iti vaktavyam . \\
\noindent \textbf{(P\_2,2.18)} duṣkulam durgavaḥ . \\
\noindent \textbf{(P\_2,2.18)} āṅ īṣadarthe . \\
\noindent \textbf{(P\_2,2.18)} āṅ īṣadarthe iti vaktavyam . \\
\noindent \textbf{(P\_2,2.18)} ākaḍāraḥ āpiṅgalaḥ . \\
\noindent \textbf{(P\_2,2.18)} kuḥ pāpārthe . \\
\noindent \textbf{(P\_2,2.18)} kuḥ pāpārtheiti vaktavyam . \\
\noindent \textbf{(P\_2,2.18)} kubrāhmaṇaḥ kuvṛṣalaḥ . \\
\noindent \textbf{(P\_2,2.18)} prādayaḥ gatādyarthe prathamayā . \\
\noindent \textbf{(P\_2,2.18)} prādayaḥ gatādyarthe prathamayā samasyante iti vaktavyam . \\
\noindent \textbf{(P\_2,2.18)} pragataḥ ācāryaḥ prācāryaḥ prāntevāsī prapitāmahaḥ . \\
\noindent \textbf{(P\_2,2.18)} atyādayaḥ krāntādyarthe dvitīyayā . \\
\noindent \textbf{(P\_2,2.18)} atyādayaḥ krāntādyarthe dvitīyayā samasyante iti vaktavyam . \\
\noindent \textbf{(P\_2,2.18)} atikrāntaḥ khaṭvām atikhaṭvaḥ atimālaḥ . \\
\noindent \textbf{(P\_2,2.18)} avādayaḥ kruṣṭādyarthe tṛtīyayā . \\
\noindent \textbf{(P\_2,2.18)} avādayaḥ kruṣṭādyarthe tṛtīyayā samasyante iti vaktavyam . \\
\noindent \textbf{(P\_2,2.18)} avakruṣṭaḥ kokilayā avakokilaḥ vasantaḥ . \\
\noindent \textbf{(P\_2,2.18)} paryādayaḥ glānādyarthe caturthyā . \\
\noindent \textbf{(P\_2,2.18)} paryādayaḥ glānādyarthe caturthyā samasyante iti vaktavyam . \\
\noindent \textbf{(P\_2,2.18)} pariglānaḥ adhyayanāya paryadhayanaḥ . \\
\noindent \textbf{(P\_2,2.18)} nirādayaḥ krāntādyarthe pañcamyā . \\
\noindent \textbf{(P\_2,2.18)} nirādayaḥ krāntādyarthe pañcamyā samasyante iti vaktavyam . \\
\noindent \textbf{(P\_2,2.18)} niṣkrāntaḥ kauśāmbyāḥ niṣkauśāmbiḥ nirvārāṇasiḥ . \\
\noindent \textbf{(P\_2,2.18)} avyayam pravṛddhādibhiḥ . \\
\noindent \textbf{(P\_2,2.18)} avyayam pravṛddhādibhiḥ samasyate iti vaktavyam . \\
\noindent \textbf{(P\_2,2.18)} punaḥprvavṛddham barhiḥ bhavati punarṇavam punaḥsukham . \\
\noindent \textbf{(P\_2,2.18)} ivena vibhaktyantalopaḥ pūrvpadaprakṛtisvaratvam ca . \\
\noindent \textbf{(P\_2,2.18)} vāsasīiva kanyeiva . \\
\noindent \textbf{(P\_2,2.18)} udāttavatā tiṅā gatimatā ca avyayam samasyate iti vaktavyam . \\
\noindent \textbf{(P\_2,2.18)} anuvyacalat anuprāviśat yat pariyanti . \\
\noindent \textbf{(P\_2,2.19)} atiṅ iti kimartham . \\
\noindent \textbf{(P\_2,2.19)} kārakaḥ vrajati . \\
\noindent \textbf{(P\_2,2.19)} hārakaḥ vrajati . \\
\noindent \textbf{(P\_2,2.19)} atiṅ iti śakyam akartum . \\
\noindent \textbf{(P\_2,2.19)} kasmāt na bhavati . \\
\noindent \textbf{(P\_2,2.19)} kārakaḥ vrajati . \\
\noindent \textbf{(P\_2,2.19)} hārakaḥ vrajati iti . \\
\noindent \textbf{(P\_2,2.19)} sup supā iti vartate . \\
\noindent \textbf{(P\_2,2.19)} ataḥ uttaram paṭhati . \\
\noindent \textbf{(P\_2,2.19)} upapadam atiṅ iti tadarthapratiṣedhaḥ . \\
\noindent \textbf{(P\_2,2.19)} upapadam atiṅ iti tadarthasya ayam pratiṣedhaḥ vaktavyaḥ . \\
\noindent \textbf{(P\_2,2.19)} kasya . \\
\noindent \textbf{(P\_2,2.19)} tiṅarthasya . \\
\noindent \textbf{(P\_2,2.19)} kaḥ punaḥ tiṅarthaḥ . \\
\noindent \textbf{(P\_2,2.19)} kriyā . \\
\noindent \textbf{(P\_2,2.19)} kriyāpratiṣedhaḥ vā . \\
\noindent \textbf{(P\_2,2.19)} atha vā vyaktam eva idam paṭhitavyam upapadam akriyayā iti . \\
\noindent \textbf{(P\_2,2.19)} atha akriyayā iti kim pratyudāhriyate . \\
\noindent \textbf{(P\_2,2.19)} kārakaḥ gataḥ hārakaḥ gataḥ . \\
\noindent \textbf{(P\_2,2.19)} na etat kriyāvāci . \\
\noindent \textbf{(P\_2,2.19)} kim tarhi. dravyavāci . \\
\noindent \textbf{(P\_2,2.19)} idam tarhi kārakasya gatiḥ kārakasya vrajyā . \\
\noindent \textbf{(P\_2,2.19)} etat api dravyvāci . \\
\noindent \textbf{(P\_2,2.19)} katham . \\
\noindent \textbf{(P\_2,2.19)} kṛdabhihitaḥ bhāvaḥ dravyavat bhavati iti . \\
\noindent \textbf{(P\_2,2.19)} evam tarhi siddhe sati yat atiṅ iti pratiṣedham śāsti tat jñāpayati ācāryaḥ anayoḥ yogayoḥ nivṛttam sup supā iti . \\
\noindent \textbf{(P\_2,2.19)} kim etasya jñāpane prayojanam . \\
\noindent \textbf{(P\_2,2.19)} gatikārakopapadānām kṛdbhiḥ samāsaḥ bhavati iti eṣā paribhāṣā na kartavyā bhavati . \\
\noindent \textbf{(P\_2,2.19)} yadi etat jñāpyate kena idānīm samāsaḥ bhaviṣyati . \\
\noindent \textbf{(P\_2,2.19)} samarthena . \\
\noindent \textbf{(P\_2,2.19)} yadi evam dhātūpasargayoḥ api samāsaḥ prāpnoti . \\
\noindent \textbf{(P\_2,2.19)} pūrvam dhātuḥ upasargeṇa yujyate paścāt sādhanena iti . \\
\noindent \textbf{(P\_2,2.19)} na etat asti . \\
\noindent \textbf{(P\_2,2.19)} pūrvam dhātuḥ sādhanena yujyate paścāt upasargeṇa . \\
\noindent \textbf{(P\_2,2.19)} sādhanam hi kriyām nirvartayati . \\
\noindent \textbf{(P\_2,2.19)} tām upasargaḥ viśinaṣṭi . \\
\noindent \textbf{(P\_2,2.19)} abhinirvṛttasya ca arthasya upasargeṇa viśeṣaḥ śakyaḥ vaktum . \\
\noindent \textbf{(P\_2,2.19)} ṣaṣṭhīsamāsāt upasargasamāsaḥ vipratiṣedhena . \\
\noindent \textbf{(P\_2,2.19)} ṣaṣṭhīsamāsāt upasargasamāsaḥ vipratiṣedhena . \\
\noindent \textbf{(P\_2,2.19)} ṣaṣṭhīsamāsasya avakāśaḥ rājñaḥ puruṣaḥ rājapuruṣaḥ . \\
\noindent \textbf{(P\_2,2.19)} upapadasamāsasya avakāśaḥ stamberamaḥ karṇejapaḥ . \\
\noindent \textbf{(P\_2,2.19)} iha ubhayam prāpnoti . \\
\noindent \textbf{(P\_2,2.19)} kumbhakāraḥ nagarakāraḥ . \\
\noindent \textbf{(P\_2,2.19)} upapadasamāsaḥ bhavati vipratiṣedhena . \\
\noindent \textbf{(P\_2,2.19)} na vā ṣaṣṭhīsamāsābhāvād upapadasamāsaḥ . na vā arthaḥ vipratiṣedhena . \\
\noindent \textbf{(P\_2,2.19)} kim kāraṇam . \\
\noindent \textbf{(P\_2,2.19)} na vā ṣaṣṭhīsamāsābhāvād upapadasamāsaḥ bhaviṣyati . \\
\noindent \textbf{(P\_2,2.19)} katham . \\
\noindent \textbf{(P\_2,2.19)} gatikārakopadānām kṛdbhiḥ saha samāsavanacam prāk subutpatteḥ iti vacanāt . \\
\noindent \textbf{(P\_2,2.19)} atha vā vibhāṣā ṣaṣṭhīsamāsaḥ . \\
\noindent \textbf{(P\_2,2.19)} yadā na ṣaṣṭhīsamāsaḥ tadā upapadasamāsaḥ bhaviṣyati . \\
\noindent \textbf{(P\_2,2.19)} anena eva yathā syāt tena mā bhūt iti . \\
\noindent \textbf{(P\_2,2.19)} kaḥ ca atra viśeṣaḥ tena vā syāt anena vā . \\
\noindent \textbf{(P\_2,2.19)} upapadasamāsaḥ nityasamāsaḥ ṣaṣṭhīsamāsaḥ punaḥ vibhāṣā . \\
\noindent \textbf{(P\_2,2.19)} nanu ca nityam yaḥ samāsaḥ saḥ nityasamāsaḥ . \\
\noindent \textbf{(P\_2,2.19)} yasya vigrahaḥ na asti . \\
\noindent \textbf{(P\_2,2.19)} na iti āha . \\
\noindent \textbf{(P\_2,2.19)} nityādhikāre yaḥ samāsaḥ saḥ nityasamāsaḥ . \\
\noindent \textbf{(P\_2,2.19)} na evam śakyam . \\
\noindent \textbf{(P\_2,2.19)} avyayībhāvasya hi anityasamāsatā prasajyeta . \\
\noindent \textbf{(P\_2,2.19)} tasmāt nityaḥ samāsaḥ nityasamāsaḥ . \\
\noindent \textbf{(P\_2,2.19)} yasya vigrahaḥ na asti . \\
\noindent \textbf{(P\_2,2.20)} evakāraḥ kimarthaḥ . \\
\noindent \textbf{(P\_2,2.20)} niyamārthaḥ . \\
\noindent \textbf{(P\_2,2.20)} na etat asti prayojanam . \\
\noindent \textbf{(P\_2,2.20)} siddhe vidhiḥ ārabhyamāṇaḥ antareṇa api evakāram niyamārthaḥ bhaviṣyati . \\
\noindent \textbf{(P\_2,2.20)} iṣṭataḥ avadhāraṇārthaḥ tarhi bhaviṣyati . \\
\noindent \textbf{(P\_2,2.20)} yathā evam vijñāyeta : amā eva avyayena iti . \\
\noindent \textbf{(P\_2,2.20)} mā evam vijñāyi : amā avyayena eva iti . \\
\noindent \textbf{(P\_2,2.20)} asti ca idānīm anavyayam amśabdaḥ yadarthaḥ vidhiḥ syāt . \\
\noindent \textbf{(P\_2,2.20)} asti iti āha . \\
\noindent \textbf{(P\_2,2.20)} khaśayam brāhmaṇakulam iti . \\
\noindent \textbf{(P\_2,2.20)} na etat asti prayojanam . \\
\noindent \textbf{(P\_2,2.20)} antaraṅgatvāt atra samāsaḥ bhaviṣyati . \\
\noindent \textbf{(P\_2,2.20)} idam tarhi prayojanam . \\
\noindent \textbf{(P\_2,2.20)} amā eva yat tulyavidhānam upapadam tatra eva yathā syāt . \\
\noindent \textbf{(P\_2,2.20)} amā ca anyena ca yat tulyavidhānam upapadam tatra mā bhūt iti . \\
\noindent \textbf{(P\_2,2.20)} agre bhojam agre bhuktvā . \\
\noindent \textbf{(P\_2,2.20)} agrādiṣu aprāptavidheḥ samāsapratiṣedham codayiṣyati . \\
\noindent \textbf{(P\_2,2.20)} saḥ na vaktavyaḥ bhavati . \\
\noindent \textbf{(P\_2,2.23)} śeṣaḥ iti ucyate . \\
\noindent \textbf{(P\_2,2.23)} kaḥ śeṣaḥ nāma . \\
\noindent \textbf{(P\_2,2.23)} yeṣām padānām anuktaḥ samāsaḥ saḥ śeṣaḥ . \\
\noindent \textbf{(P\_2,2.23)} śeṣavacanam padataḥ cet na abhāvāt . \\
\noindent \textbf{(P\_2,2.23)} śeṣavacanam padataḥ cet tat na . \\
\noindent \textbf{(P\_2,2.23)} kim kāraṇam . \\
\noindent \textbf{(P\_2,2.23)} abhāvāt . \\
\noindent \textbf{(P\_2,2.23)} na hi santi tāni padāni yeṣām padānām anuktaḥ samāsaḥ . \\
\noindent \textbf{(P\_2,2.23)} arthataḥ tarhi śeṣagrahaṇam . \\
\noindent \textbf{(P\_2,2.23)} yeṣu artheṣu anuktaḥ samāsaḥ saḥ śeṣaḥ . \\
\noindent \textbf{(P\_2,2.23)} arthataḥ cet aviśiṣṭam . arthataḥ cet aviśiṣṭam etat bhavati . \\
\noindent \textbf{(P\_2,2.23)} kutaḥ . \\
\noindent \textbf{(P\_2,2.23)} padataḥ . \\
\noindent \textbf{(P\_2,2.23)} na hi santi te arthāḥ yeṣu anuktaḥ samāsaḥ . \\
\noindent \textbf{(P\_2,2.23)} trikataḥ tarhi śeṣagrahaṇam . \\
\noindent \textbf{(P\_2,2.23)} yasya trikasya anuktaḥ samāsaḥ saḥ śeṣaḥ . \\
\noindent \textbf{(P\_2,2.23)} kasya ca anuktaḥ . \\
\noindent \textbf{(P\_2,2.23)} prathamāyāḥ . \\
\noindent \textbf{(P\_2,2.24.1)} padagrahaṇam kimartham. anekam anyārthe iti iyati ucyamāne vakyārthe api bahuvrīhiḥ syāt . \\
\noindent \textbf{(P\_2,2.24.1)} yathā me mātā tathā me pitā susnātam bhoḥ iti . \\
\noindent \textbf{(P\_2,2.24.1)} padagrahaṇe punaḥ kriyamāṇe na doṣaḥ bhavati . \\
\noindent \textbf{(P\_2,2.24.1)} atha anyagrahaṇam kimartham . \\
\noindent \textbf{(P\_2,2.24.1)} anekam padārthe iti iyati ucyamāne svapadārthe api bahurvīhiḥ syāt . \\
\noindent \textbf{(P\_2,2.24.1)} rājapuruṣaḥ takṣapuruṣaḥ iti . \\
\noindent \textbf{(P\_2,2.24.1)} na etat asti prayojanam . \\
\noindent \textbf{(P\_2,2.24.1)} tatpuruṣaḥ svapadārthe bādhakaḥ bhaviṣyati . \\
\noindent \textbf{(P\_2,2.24.1)} bhavet ekasañjñādhikāre siddham . \\
\noindent \textbf{(P\_2,2.24.1)} paraṅkāryatve tu na sidhyati . \\
\noindent \textbf{(P\_2,2.24.1)} ārambhasāmarthyāt ca tatpuruṣaḥ paraṅkāryatvāt ca bahuvrīhiḥ prāpnoti . \\
\noindent \textbf{(P\_2,2.24.1)} paraṅkāryatve ca na doṣaḥ . \\
\noindent \textbf{(P\_2,2.24.1)} śeṣaḥ iti vartate . \\
\noindent \textbf{(P\_2,2.24.1)} śeṣatvāt na bhaviṣyati . \\
\noindent \textbf{(P\_2,2.24.1)} śeṣavacane uktam . \\
\noindent \textbf{(P\_2,2.24.1)} kim uktam . \\
\noindent \textbf{(P\_2,2.24.1)} tatra śeṣavacanāt doṣaḥ saṅkhyāsamānādhikaraṇanañsamāseṣu bahuvrīhipratiṣedhaḥ iti . \\
\noindent \textbf{(P\_2,2.24.1)} atha ekasaṅjñādhikāre na arthaḥ anyagrahaṇena . \\
\noindent \textbf{(P\_2,2.24.1)} ekasaṅjñādhikāre ca kartavyam . \\
\noindent \textbf{(P\_2,2.24.1)} akriyamāṇe hi anyagrahaṇe yathā eva tatpuruṣaḥ svapadārthe bahuvrīhim bādhate evam anyapadārthe api bādheta . \\
\noindent \textbf{(P\_2,2.24.1)} atha anekagrahaṇam kimartham. anyapadārthe iti iyati ucyamāne ekasya api padasya bahuvrīhiḥ syāt . \\
\noindent \textbf{(P\_2,2.24.1)} sarpiṣaḥ api syāt . \\
\noindent \textbf{(P\_2,2.24.1)} madhunaḥ api syāt . \\
\noindent \textbf{(P\_2,2.24.1)} gomūtrasya api syāt . \\
\noindent \textbf{(P\_2,2.24.1)} na etat asti prayojanam . \\
\noindent \textbf{(P\_2,2.24.1)} sup supā iti vartate . \\
\noindent \textbf{(P\_2,2.24.1)} idam tarhi prayojanam . \\
\noindent \textbf{(P\_2,2.24.1)} bahūnām api samāsaḥ yathā syāt . \\
\noindent \textbf{(P\_2,2.24.1)} susūkṣmajaṭakeśena sunatājinavāsasā . \\
\noindent \textbf{(P\_2,2.24.1)} uttarārtham ca anekagrahaṇam kartavyam cārthe dvandvaḥ anekam iti . \\
\noindent \textbf{(P\_2,2.24.1)} iha api yathā syāt . \\
\noindent \textbf{(P\_2,2.24.1)} plakṣanyagrodhakhadirapalāśāḥ iti . \\
\noindent \textbf{(P\_2,2.24.1)} etat api na asti prayojanam . \\
\noindent \textbf{(P\_2,2.24.1)} ācāryapravṛttiḥ jñāpayati bahūnām api samāsaḥ bhavati iti yat ayam uttarapade dvigum śāsti . \\
\noindent \textbf{(P\_2,2.24.1)} tatpuruṣaḥ api tarhi bahūnām prāpnoti . \\
\noindent \textbf{(P\_2,2.24.1)} grahaṇena tatpuruṣaḥ ucyate . \\
\noindent \textbf{(P\_2,2.24.1)} tena bahūnām na bhaviṣyati . \\
\noindent \textbf{(P\_2,2.24.1)} ataḥ uttaram paṭhati . \\
\noindent \textbf{(P\_2,2.24.1)} anekavacanam upasarjanārtham . \\
\noindent \textbf{(P\_2,2.24.1)} anekagrahaṇam kriyate upasarjanārtham . \\
\noindent \textbf{(P\_2,2.24.1)} prathamānirdiṣṭam samāse upasarjanam iti anekasya supaḥ upasarjanasañjñā yathā syāt . \\
\noindent \textbf{(P\_2,2.24.1)} citraguḥ śabalaguḥ iti . \\
\noindent \textbf{(P\_2,2.24.1)} na vā ekavibhaktitvāt . \\
\noindent \textbf{(P\_2,2.24.1)} na vā etat api prayojanam asti . \\
\noindent \textbf{(P\_2,2.24.1)} kim kāraṇam . \\
\noindent \textbf{(P\_2,2.24.1)} ekavibhaktitvāt . \\
\noindent \textbf{(P\_2,2.24.1)} ekavibhakti ca apūrvnipāte iti upasarjanasañjñā bhaviṣyati . \\
\noindent \textbf{(P\_2,2.24.1)} citraguḥ śabalaguḥ iti . \\
\noindent \textbf{(P\_2,2.24.1)} citrāḥ yasya gāvaḥ citraguḥ tiṣṭhati . \\
\noindent \textbf{(P\_2,2.24.1)} citrāḥ yasya gāvaḥ citragum paśya . \\
\noindent \textbf{(P\_2,2.24.1)} citrāḥ yasya gāvaḥ citraguṇā kṛtam . \\
\noindent \textbf{(P\_2,2.24.1)} citrāḥ yasya gāvaḥ citragave dehi . \\
\noindent \textbf{(P\_2,2.24.1)} citrāḥ yasya gāvaḥ citragoḥ ānaya . \\
\noindent \textbf{(P\_2,2.24.1)} citrāḥ yasya gāvaḥ citragoḥ svam . \\
\noindent \textbf{(P\_2,2.24.1)} citrāḥ yasya gāvaḥ citragau nidhehi . \\
\noindent \textbf{(P\_2,2.24.1)} citrāḥ yasya gāvaḥ he citrago iti . \\
\noindent \textbf{(P\_2,2.24.1)} yadi tarhi yataḥ kutaḥ cit eva kim cit padam adhyāhṛtya ekavibhaktyā yogaḥ kriyate etat api ekavibhaktiyuktam bhavati iha api prāpnoti . \\
\noindent \textbf{(P\_2,2.24.1)} rājakumārī takṣakumārī . \\
\noindent \textbf{(P\_2,2.24.1)} rājñaḥ yā kumārī rājakumārī tiṣṭhati . \\
\noindent \textbf{(P\_2,2.24.1)} rājñaḥ yā kumārī rājakumārīm paśya . \\
\noindent \textbf{(P\_2,2.24.1)} rājñaḥ yā kumārī rājakumāryā kṛtam . \\
\noindent \textbf{(P\_2,2.24.1)} rājñaḥ yā kumārī rājakumāryai dehi . \\
\noindent \textbf{(P\_2,2.24.1)} rājñaḥ yā kumārī rājakumāryāḥ ānaya . \\
\noindent \textbf{(P\_2,2.24.1)} rājñaḥ yā kumārī rājakumāryāḥ svam . \\
\noindent \textbf{(P\_2,2.24.1)} rājñaḥ yā kumārī rājakumāryām nidhehi . \\
\noindent \textbf{(P\_2,2.24.1)} rājñaḥ yā kumārī he rājakumāri iti . \\
\noindent \textbf{(P\_2,2.24.1)} kim vaktavyam etat . \\
\noindent \textbf{(P\_2,2.24.1)} na hi . \\
\noindent \textbf{(P\_2,2.24.1)} katham anucyamānam gaṃsyate . \\
\noindent \textbf{(P\_2,2.24.1)} ekagrahaṇasāmarthyāt . \\
\noindent \textbf{(P\_2,2.24.1)} yadi hi yat ekavibhaktiyuktam ca anekavibhaktiyuktam ca tatra syāt ekagrahaṇam anarthakam syāt . \\
\noindent \textbf{(P\_2,2.24.1)} vibhaktiyuktam ca apūrvanipāte iti eva brūyāt . \\
\noindent \textbf{(P\_2,2.24.2)} padārthābhidhāne anuprayogānupapattiḥ abhihitatvāt . \\
\noindent \textbf{(P\_2,2.24.2)} padārthasya abhidhāne anuprayogasya anupapattiḥ . \\
\noindent \textbf{(P\_2,2.24.2)} citraguḥ devadattaḥ iti . \\
\noindent \textbf{(P\_2,2.24.2)} kim kāraṇam . \\
\noindent \textbf{(P\_2,2.24.2)} abhihitatvāt . \\
\noindent \textbf{(P\_2,2.24.2)} citraguśabdena abhihitaḥ saḥ arthaḥ iti kṛtvā anuprayogaḥ na prāpnoti . \\
\noindent \textbf{(P\_2,2.24.2)} na vā anabhihitatvāt . \\
\noindent \textbf{(P\_2,2.24.2)} na vā eṣaḥ doṣaḥ . \\
\noindent \textbf{(P\_2,2.24.2)} kim kāraṇam . \\
\noindent \textbf{(P\_2,2.24.2)} anabhihitatvāt . \\
\noindent \textbf{(P\_2,2.24.2)} citraguśabdena anabhihitaḥ saḥ arthaḥ iti kṛtvā anuprayogaḥ bhaviṣyati . \\
\noindent \textbf{(P\_2,2.24.2)} katham anabhihitaḥ yāvatā idānīm eva uktam padārthābhidhāne anuprayogānupapattiḥ abhihitatvāt iti . \\
\noindent \textbf{(P\_2,2.24.2)} sāmānyābhidhāne hi viśeṣānabhidhānam . \\
\noindent \textbf{(P\_2,2.24.2)} sāmānye hi abhidhīyamāne viśeṣaḥ anabhhitaḥ bhavati . \\
\noindent \textbf{(P\_2,2.24.2)} tatra avaśyam viśeṣārthinā viśeṣaḥ anuprayoktavyaḥ . \\
\noindent \textbf{(P\_2,2.24.2)} citraguḥ . \\
\noindent \textbf{(P\_2,2.24.2)} kaḥ . \\
\noindent \textbf{(P\_2,2.24.2)} devadattaḥ iti. bhavet siddham yadā sāmānye vṛttiḥ . \\
\noindent \textbf{(P\_2,2.24.2)} yadā tu khalu viśeṣe vṛttiḥ tadā na sidhyati . \\
\noindent \textbf{(P\_2,2.24.2)} citrā gāvaḥ devadattasya citraguḥ devadattaḥ iti . \\
\noindent \textbf{(P\_2,2.24.2)} tat api siddham . \\
\noindent \textbf{(P\_2,2.24.2)} katham . \\
\noindent \textbf{(P\_2,2.24.2)} na idam ubhayam yugapat bhavati vākyam ca samāsaḥ ca . \\
\noindent \textbf{(P\_2,2.24.2)} yadā vākyam tadā na samāsaḥ . \\
\noindent \textbf{(P\_2,2.24.2)} yadā samāsaḥ tadā na vakyam . \\
\noindent \textbf{(P\_2,2.24.2)} yadā samāsaḥ tadā sāmānye vṛttiḥ . \\
\noindent \textbf{(P\_2,2.24.2)} tatra avaśyam viśeṣārthinā viśeṣaḥ anuprayoktavyaḥ . \\
\noindent \textbf{(P\_2,2.24.2)} citraguḥ . \\
\noindent \textbf{(P\_2,2.24.2)} kaḥ . \\
\noindent \textbf{(P\_2,2.24.2)} devadattaḥ iti. sāmānyasya eva tarhi anuprayogaḥ na prāpnoti . \\
\noindent \textbf{(P\_2,2.24.2)} citragu tat . \\
\noindent \textbf{(P\_2,2.24.2)} citragu kim cit . \\
\noindent \textbf{(P\_2,2.24.2)} citragu sarvam iti . \\
\noindent \textbf{(P\_2,2.24.2)} sāmānyam api yathā viśeṣaḥ tadvat . \\
\noindent \textbf{(P\_2,2.24.2)} citragu iti ukte sandehaḥ syāt . \\
\noindent \textbf{(P\_2,2.24.2)} sarvam vā viśvam vā iti . \\
\noindent \textbf{(P\_2,2.24.2)} tatra avaśyam sandehanivṛttyartham viśeṣārthinā viśeṣaḥ anuprayoktavyaḥ . \\
\noindent \textbf{(P\_2,2.24.2)} atha vā vibhaktyarthaḥ abhidīyate . \\
\noindent \textbf{(P\_2,2.24.2)} etat ca atra yuktam yat vibhaktyarthaḥ abhidhīyate . \\
\noindent \textbf{(P\_2,2.24.2)} tatra hi sarvapaścāt padam vartate asya iti . \\
\noindent \textbf{(P\_2,2.24.2)} vibhaktyarthābhidhāne adravyasya liṅgasaṅkhyopacārānupapattiḥ . \\
\noindent \textbf{(P\_2,2.24.2)} vibhaktyarthābhidhāne adravyasya liṅgasaṅkhyābhyām upacāraḥ anupapannaḥ . \\
\noindent \textbf{(P\_2,2.24.2)} bahuyavam bahuyavā bahuyavaḥ bahuyavau bahuhavāḥ iti . \\
\noindent \textbf{(P\_2,2.24.2)} aparaḥ āha : vibhaktyarthābhidhāne adravyasya liṅgasaṅkhyopacārānupapattiḥ vibhaktyarthābhidhāne dravyasya ye liṅgasaṅkhye tābhyām vibhaktyarthasya upacāraḥ anupapannaḥ . \\
\noindent \textbf{(P\_2,2.24.2)} bahuyavam bahuyavāḥ bahuyavaḥ bahuyavau bahuhavāḥ iti . \\
\noindent \textbf{(P\_2,2.24.2)} katham hi anyasya liṅgasaṅkhyābhyām anyasya upacāraḥ syāt . \\
\noindent \textbf{(P\_2,2.24.2)} siddham tu yathā guṇavacaneṣu . \\
\noindent \textbf{(P\_2,2.24.2)} siddham etat. katham . \\
\noindent \textbf{(P\_2,2.24.2)} yathā guṇavacaneṣu . \\
\noindent \textbf{(P\_2,2.24.2)} guṇavacaneṣu uktam : guṇavacanānām śabdānām āśrayataḥ liṅgavacanāni bhavanti iti . \\
\noindent \textbf{(P\_2,2.24.2)} tat yathā śuklam vastram śuklā śāṭī śuklaḥ kambalaḥ śuklau kambalau śuklāḥ kambalāḥ iti . \\
\noindent \textbf{(P\_2,2.24.2)} yat asau dravyam śritaḥ bhavati guṇaḥ tasya yat liṅgam vacanam ca tat guṇasya api bhavati . \\
\noindent \textbf{(P\_2,2.24.2)} evam iha api yat asau dravyam śritaḥ vibhaktyarthaḥ tasya yat liṅgam vacanam ca tat samāsasya api bhaviṣyati . \\
\noindent \textbf{(P\_2,2.24.2)} yadi tarhi vibhaktyarthaḥ abhidhīyate kṛtsnaḥ padārthaḥ katham abhihitaḥ bhavati sadravyaḥ saliṅgaḥ sasaṅkhyaḥ ca . \\
\noindent \textbf{(P\_2,2.24.2)} arthagrahaṇasāmarthyāt . \\
\noindent \textbf{(P\_2,2.24.2)} iha anekam anyapade iti iyatā siddham . \\
\noindent \textbf{(P\_2,2.24.2)} katham punaḥ pade nāma vṛttiḥ syāt . \\
\noindent \textbf{(P\_2,2.24.2)} śabdaḥ hi eṣaḥ . \\
\noindent \textbf{(P\_2,2.24.2)} śabde asambhavāt arthe kāryam vijñāsyate . \\
\noindent \textbf{(P\_2,2.24.2)} saḥ ayam evam siddhe sati yat arthagrahaṇam karoti tasya etat prayojanam kṛtsnaḥ padārthaḥ yathā abhidhīyeta sadravyaḥ saliṅgaḥ sasaṅkhyaḥ ca iti . \\
\noindent \textbf{(P\_2,2.24.2)} yadi tarhi kṛtsnaḥ padārthaḥ abhidhīyate laiṅgāḥ sāṅkhyāḥ ca vidhayaḥ na sidhyanti . \\
\noindent \textbf{(P\_2,2.24.2)} uktam vā . \\
\noindent \textbf{(P\_2,2.24.2)} kim uktam . \\
\noindent \textbf{(P\_2,2.24.2)} liṅgeṣu tāvat . \\
\noindent \textbf{(P\_2,2.24.2)} siddham tu striyāḥ prātipadikaviśeṣaṇatvāt svārthe ṭābādayaḥ iti . \\
\noindent \textbf{(P\_2,2.24.2)} sāṅkhyeṣu api uktam karmādīnām anuktāḥ ekatvādayaḥ iti kṛtvā sāṅkhyāḥ bhaviṣyanti . \\
\noindent \textbf{(P\_2,2.24.2)} prathamā tarhi na prāpnoti . \\
\noindent \textbf{(P\_2,2.24.2)} samayāt bhaviṣyati . \\
\noindent \textbf{(P\_2,2.24.2)} yadi sāmayikī na niyogataḥ anyāḥ kasmāt na bhavanti . \\
\noindent \textbf{(P\_2,2.24.2)} karmādīnām abhāvāt . \\
\noindent \textbf{(P\_2,2.24.2)} ṣaṣṭhī tarhi prāpnoti . \\
\noindent \textbf{(P\_2,2.24.2)} śeṣalakṣaṇā ṣaṣṭhī . \\
\noindent \textbf{(P\_2,2.24.2)} aśeṣatvāt na bhaviṣyati . \\
\noindent \textbf{(P\_2,2.24.2)} evam api vyatikaraḥ . \\
\noindent \textbf{(P\_2,2.24.2)} ekasmin api dvivacanabahuvacane prāpnutaḥ dvayoḥ api ekavacanabahuvacane bahuṣu api ekavacanadvivacane . \\
\noindent \textbf{(P\_2,2.24.2)} arthataḥ vyavasthā bhaviṣyati . \\
\noindent \textbf{(P\_2,2.24.2)} atha vā saṅkhyā nāma iyam parapradhānā . \\
\noindent \textbf{(P\_2,2.24.2)} saṅkhyeam anyā viśeṣyam . \\
\noindent \textbf{(P\_2,2.24.2)} yadi ca atra prathamā na syāt saṅkhyeyam aviśeṣitam syāt . \\
\noindent \textbf{(P\_2,2.24.2)} atha vā vakṣyati etat . \\
\noindent \textbf{(P\_2,2.24.2)} tatra vacanagrahaṇasya prayojanam ukteṣu api ekatvādiṣu prathamā yathā syāt iti . \\
\noindent \textbf{(P\_2,2.24.2)} evam api ṣaṣṭhī prāpnoti . \\
\noindent \textbf{(P\_2,2.24.2)} kim kāraṇam . \\
\noindent \textbf{(P\_2,2.24.2)} vyabhicarati eva hi ayam samāsaḥ liṅgasaṅkhye . \\
\noindent \textbf{(P\_2,2.24.2)} ṣaṣthyartham punaḥ na vyabhicarati . \\
\noindent \textbf{(P\_2,2.24.2)} abhihitaḥ saḥ arthaḥ antarbhūtaḥ prātipadikārthaḥ sampannaḥ . \\
\noindent \textbf{(P\_2,2.24.2)} tatra prātipadikārthe prathamā iti prathamā bhaviṣyati . \\
\noindent \textbf{(P\_2,2.24.2)} na tarhi idānīm idam bhavati : citragoḥ devadattasya . \\
\noindent \textbf{(P\_2,2.24.2)} bhavati . \\
\noindent \textbf{(P\_2,2.24.2)} bāhyam artham apekṣya ṣaṣṭhī . \\
\noindent \textbf{(P\_2,2.24.3)} parigaṇanam kartavyam . \\
\noindent \textbf{(P\_2,2.24.3)} bahuvrīhiḥ samānādhikaraṇānām . \\
\noindent \textbf{(P\_2,2.24.3)} samānādhikaraṇānām bahuvrīhiḥ vaktavyaḥ . \\
\noindent \textbf{(P\_2,2.24.3)} kim prayojanam . \\
\noindent \textbf{(P\_2,2.24.3)} vyadhikaraṇānām mā bhūt iti . \\
\noindent \textbf{(P\_2,2.24.3)} pañcabhiḥ bhuktam asya iti . \\
\noindent \textbf{(P\_2,2.24.3)} avyayānām ca . \\
\noindent \textbf{(P\_2,2.24.3)} avyayānām bahuvrīhiḥ vaktavyaḥ . \\
\noindent \textbf{(P\_2,2.24.3)} uccairmukhaḥ nīcairmukhaḥ . \\
\noindent \textbf{(P\_2,2.24.3)} saptamyupamānapūrvapadasya uttarapadalopaḥ ca . \\
\noindent \textbf{(P\_2,2.24.3)} saptamīpūrvasya upamānapūrvasya ca bahuvrīhiḥ vaktavyaḥ uttarapadasya ca lopaḥ vaktavyaḥ . \\
\noindent \textbf{(P\_2,2.24.3)} kaṇṭhesthaḥ kālaḥ asya kaṇṭhekālaḥ uṣṭramukham iva mukham asya uṣṭramukhaḥ kharamukhaḥ . \\
\noindent \textbf{(P\_2,2.24.3)} samudāyavikāraṣaṣṭhyāḥ ca . \\
\noindent \textbf{(P\_2,2.24.3)} samudāyaṣaṣṭhyāḥ vikāraṣaṣṭhyāḥ ca bahuvrīhiḥ vaktavyaḥ uttarapadasya ca lopaḥ vaktavyaḥ . \\
\noindent \textbf{(P\_2,2.24.3)} keśānām samāhāraḥ cūḍā asya keśacūḍaḥ suvarṇasya vikāraḥ alaṅkāraḥ asya suvarṇālaṅkāraḥ . \\
\noindent \textbf{(P\_2,2.24.3)} prādibhyaḥ dhātujasya vā . \\
\noindent \textbf{(P\_2,2.24.3)} prādibhyaḥ dhātujasya bahuvrīhiḥ vaktavyaḥ uttarapadasya ca vā lopaḥ vaktavyaḥ . \\
\noindent \textbf{(P\_2,2.24.3)} prapatitaparṇaḥ praparṇaḥ prapatitapalāśaḥ prapalāśaḥ . \\
\noindent \textbf{(P\_2,2.24.3)} nañaḥ astyarthānām . \\
\noindent \textbf{(P\_2,2.24.3)} nañaḥ astyarthānām bahuvrīhiḥ vaktavyaḥ uttarapadasya ca vā lopaḥ vaktavyaḥ . \\
\noindent \textbf{(P\_2,2.24.3)} avidyamānaputraḥ aputraḥ avidyamānabhāryaḥ abhāryaḥ . \\
\noindent \textbf{(P\_2,2.24.3)} tat tarhi bahu vaktavyam . \\
\noindent \textbf{(P\_2,2.24.3)} na vā anabhidhānāt asamānādhikaraṇe sañjñābhāvaḥ . \\
\noindent \textbf{(P\_2,2.24.3)} na vā vaktavyam . \\
\noindent \textbf{(P\_2,2.24.3)} asamānādhikaraṇānām bahuvrīhiḥ kasmāt na bhavati : pañcabhiḥ bhuktam asya iti . \\
\noindent \textbf{(P\_2,2.24.3)} anabhidhānāt . \\
\noindent \textbf{(P\_2,2.24.3)} tat ca avaśyam anabhidhānam āśrayitavyam . \\
\noindent \textbf{(P\_2,2.24.3)} kriyamāṇe api vai parigaṇane yatra abhidhānam na asti na bhavati tatra bahuvrīhiḥ . \\
\noindent \textbf{(P\_2,2.24.3)} tat yathā pañca bhuktavantaḥ asya iti . \\
\noindent \textbf{(P\_2,2.24.3)} atha etasmin sati anabhidhāne yadi vṛttiparigaṇanam kriyate vartiparigaṇanam api kartavyam . \\
\noindent \textbf{(P\_2,2.24.3)} tat katham kartavyam . \\
\noindent \textbf{(P\_2,2.24.3)} arthaniyame matvarthagrahaṇam . \\
\noindent \textbf{(P\_2,2.24.3)} arthaniyame matvarthagrahaṇam kartavyam . \\
\noindent \textbf{(P\_2,2.24.3)} matvarthe yaḥ saḥ bahuvrīhiḥ iti vaktavyam . \\
\noindent \textbf{(P\_2,2.24.3)} iha mā bhūt : kaṣṭam śritam anena iti . \\
\noindent \textbf{(P\_2,2.24.3)} tathā ca uttarasya vacanārthaḥ . \\
\noindent \textbf{(P\_2,2.24.3)} evam ca kṛtvā uttarasya yogasya vacanārthaḥ upapannaḥ bhavati . \\
\noindent \textbf{(P\_2,2.24.3)} ke cit tāvat āhuḥ : yat vṛttisūtre iti . \\
\noindent \textbf{(P\_2,2.24.3)} saṅkhyāvyayāsannādūrādhikasaṅkhyāḥ saṅkhyeye iti . \\
\noindent \textbf{(P\_2,2.24.3)} aparaḥ āha : yat vārttike iti . \\
\noindent \textbf{(P\_2,2.24.3)} karmavacanena aprathmāyāḥ . \\
\noindent \textbf{(P\_2,2.24.3)} karmavacanena aprathmāyāḥ bahuvrīhiḥ vaktavyaḥ . \\
\noindent \textbf{(P\_2,2.24.3)} ūḍhaḥ rathaḥ anena ūḍharathaḥ anaḍvān upahṛtaḥ paśuḥ rudrāya upahṛtapaśuḥ rudraḥ uddhṛtaḥ odanaḥ sthālyāḥ uddhṛtaudanā sthālī . \\
\noindent \textbf{(P\_2,2.24.3)} yadi karmavacanena iti ucyate kartṛvacanena katham . \\
\noindent \textbf{(P\_2,2.24.3)} prāptam udakam grāmam prāptodakaḥ grāmaḥ āgatāḥ atithayaḥ grāmam āgatātithiḥ grāmaḥ . \\
\noindent \textbf{(P\_2,2.24.3)} kartṛvacanena api . \\
\noindent \textbf{(P\_2,2.24.3)} kartṛvacanena api iti vaktavyam . \\
\noindent \textbf{(P\_2,2.24.3)} aprathamāyāḥ iti kimartham . \\
\noindent \textbf{(P\_2,2.24.3)} vṛṣṭe deve gataḥ . \\
\noindent \textbf{(P\_2,2.24.3)} aprathamāyāḥ iti ucyamāne iha kasmāt na bhavati . \\
\noindent \textbf{(P\_2,2.24.3)} vṛṣṭe deve gatam paśya iti . \\
\noindent \textbf{(P\_2,2.24.3)} bahiraṅgā atra aprathamā . \\
\noindent \textbf{(P\_2,2.24.3)} subadhikāre astikṣīrādivacanam . \\
\noindent \textbf{(P\_2,2.24.3)} subadhikāre astikṣīrādīnām upasaṅkhyānam kartavyam . \\
\noindent \textbf{(P\_2,2.24.3)} astikṣīrā brāhmaṇī . \\
\noindent \textbf{(P\_2,2.24.3)} tat tarhi vaktavyam . \\
\noindent \textbf{(P\_2,2.24.3)} na vā avyayatvāt . \\
\noindent \textbf{(P\_2,2.24.3)} na vā vaktavyam . \\
\noindent \textbf{(P\_2,2.24.3)} kim kāraṇam . \\
\noindent \textbf{(P\_2,2.24.3)} avyayatvāt . \\
\noindent \textbf{(P\_2,2.24.3)} avyayaḥ ayam astiśabdaḥ . \\
\noindent \textbf{(P\_2,2.24.3)} na eṣaḥ asteḥ laṭ . \\
\noindent \textbf{(P\_2,2.24.3)} katham avyayatvam . \\
\noindent \textbf{(P\_2,2.24.3)} upasargavibhaktisvarapratirūpakāḥ ca nipātasañjñāḥ bhavanti iti nipātasñjñā . \\
\noindent \textbf{(P\_2,2.24.3)} nipātaḥ avyayam iti avyayasañjñā . \\
\noindent \textbf{(P\_2,2.24.4)} atha kiṃsabrahmacārī iti kaḥ ayam samāsaḥ . \\
\noindent \textbf{(P\_2,2.24.4)} bahuvrīhiḥ iti aha . \\
\noindent \textbf{(P\_2,2.24.4)} kaḥ asya vigrahaḥ . \\
\noindent \textbf{(P\_2,2.24.4)} ke sabrahmacāriṇaḥ asya iti . \\
\noindent \textbf{(P\_2,2.24.4)} yadi evam kaṭhaḥ iti prativacanam na upapadyate . \\
\noindent \textbf{(P\_2,2.24.4)} na hi anyat pṛṣṭena anyat ākhyāyate . \\
\noindent \textbf{(P\_2,2.24.4)} evam tarhi evam vigrahaḥ kariṣyate : keṣām sabrahmacārī kiṃsbrahmacārī iti . \\
\noindent \textbf{(P\_2,2.24.4)} prativacanam ca eva na upapadyate svare ca doṣaḥ bhavati . \\
\noindent \textbf{(P\_2,2.24.4)} kiṃsabrahmacārī iti evam svaraḥ prasajyeta . \\
\noindent \textbf{(P\_2,2.24.4)} kiṃsabrahmacārī iti ca iṣyate . \\
\noindent \textbf{(P\_2,2.24.4)} evam tarhi evam vigrahaḥ kariṣyate . \\
\noindent \textbf{(P\_2,2.24.4)} kaḥ sabrahmacārī kiṃsabrahmacārī iti . \\
\noindent \textbf{(P\_2,2.24.4)} bhavet prativacanam upapannam svare tu doṣaḥ bhavati . \\
\noindent \textbf{(P\_2,2.24.4)} evam tarhi evam vigrahaḥ kariṣyate . \\
\noindent \textbf{(P\_2,2.24.4)} kaḥ sabrahmacārī tava kiṃsabrahmacārī tvam iti . \\
\noindent \textbf{(P\_2,2.24.4)} atha vā punaḥ astu evam vigrahaḥ : ke sabrahmacāriṇaḥ asya iti . \\
\noindent \textbf{(P\_2,2.24.4)} nanu ca uktam kaṭhaḥ iti prativacanam na upapadyate . \\
\noindent \textbf{(P\_2,2.24.4)} na eṣaḥ doṣaḥ . \\
\noindent \textbf{(P\_2,2.24.4)} agnaukaravāṇinyāyena bhaviṣyati . \\
\noindent \textbf{(P\_2,2.24.4)} tat yathā . \\
\noindent \textbf{(P\_2,2.24.4)} kaḥ cit kam cit āha . \\
\noindent \textbf{(P\_2,2.24.4)} agnau karavāṇi iti . \\
\noindent \textbf{(P\_2,2.24.4)} kuru iti kartari anujñāte karma api anujñātam bhavati . \\
\noindent \textbf{(P\_2,2.24.4)} aparaḥ āha : agnau kariṣyate iti . \\
\noindent \textbf{(P\_2,2.24.4)} kriyatām iti karmaṇi anujñāte kartā api anujñātaḥ bhavati . \\
\noindent \textbf{(P\_2,2.24.4)} yathā eva khalu api ke sabrahmacāriṇaḥ asya iti kaṭhāḥ iti ukte sambandhāt etat gamyate . \\
\noindent \textbf{(P\_2,2.24.4)} nūnam saḥ api kaṭha iti . \\
\noindent \textbf{(P\_2,2.24.4)} evam kaṭhaḥ iti ukte sambandhāt etat gantavyam syāt . \\
\noindent \textbf{(P\_2,2.24.4)} nūnam te api kaṭhāḥ iti . \\
\noindent \textbf{(P\_2,2.24.4)} na khalu api te śakyāḥ samāsena pratinirdeṣṭum . \\
\noindent \textbf{(P\_2,2.24.4)} upasarjanam he te bhavanti . \\
\noindent \textbf{(P\_2,2.24.4)} atha arthatṛtīyāḥ iti kaḥ ayam samāsaḥ . \\
\noindent \textbf{(P\_2,2.24.4)} bahuvrīhiḥ iti āha . \\
\noindent \textbf{(P\_2,2.24.4)} kaḥ asya vigrahaḥ . \\
\noindent \textbf{(P\_2,2.24.4)} ardham tṛtīyam eṣām iti . \\
\noindent \textbf{(P\_2,2.24.4)} kaḥ samāsārthaḥ . \\
\noindent \textbf{(P\_2,2.24.4)} samāsārthaḥ na upapadyate . \\
\noindent \textbf{(P\_2,2.24.4)} anyapadārthaḥ hi nāma saḥ bhavati . \\
\noindent \textbf{(P\_2,2.24.4)} yeṣām padānām samāsaḥ tataḥ anyasya padasya arthaḥ anyapadārthaḥ . \\
\noindent \textbf{(P\_2,2.24.4)} evam tarhi evam vigrahaḥ kariṣyate . \\
\noindent \textbf{(P\_2,2.24.4)} ardham tṛtīyam anayoḥ iti . \\
\noindent \textbf{(P\_2,2.24.4)} evam api kaḥ ṣaṣṭhyarthaḥ . \\
\noindent \textbf{(P\_2,2.24.4)} ṣaṣṭhyarthaḥ na upapadyate . \\
\noindent \textbf{(P\_2,2.24.4)} kim hi tayoḥ ardham bhavati . \\
\noindent \textbf{(P\_2,2.24.4)} astu tari evam vigrahaḥ ardham tṛtīyam eṣām iti . \\
\noindent \textbf{(P\_2,2.24.4)} nanu ca uktam samāsārthaḥ na upapadyate iti . \\
\noindent \textbf{(P\_2,2.24.4)} na eṣaḥ doṣaḥ . \\
\noindent \textbf{(P\_2,2.24.4)} avayavena vigrahaḥ samudāyaḥ samāsārthaḥ . \\
\noindent \textbf{(P\_2,2.24.4)} yadi avayavena vigrahaḥ samudāyaḥ samāsārthaḥ asidvitīyaḥ anusasāra pāṇḍavam . \\
\noindent \textbf{(P\_2,2.24.4)} saṅkarṣaṇadvitīyasya balam kṛṣṇasya vardhatām iti. dvayoḥ dvivacanam prāpnoti . \\
\noindent \textbf{(P\_2,2.24.4)} astu tarhi ayam eva vigrahaḥ ardham tṛtīyam anayoḥ . \\
\noindent \textbf{(P\_2,2.24.4)} nanu ca uktam . \\
\noindent \textbf{(P\_2,2.24.4)} ṣaṣṭhyarthaḥ na upapadyate iti . \\
\noindent \textbf{(P\_2,2.24.4)} na eṣaḥ doṣaḥ . \\
\noindent \textbf{(P\_2,2.24.4)} idam tāvat ayam praṣṭavyaḥ . \\
\noindent \textbf{(P\_2,2.24.4)} atha iha devadattasya bhrātā iti kaḥ ṣaṣṭhyarthaḥ . \\
\noindent \textbf{(P\_2,2.24.4)} tatra etat syāt . \\
\noindent \textbf{(P\_2,2.24.4)} ekasmāt prādurbhāvaḥ iti . \\
\noindent \textbf{(P\_2,2.24.4)} etat ca vārtam . \\
\noindent \textbf{(P\_2,2.24.4)} tat yathā . \\
\noindent \textbf{(P\_2,2.24.4)} sārthikanam ekapratiśraye uṣitānām prātaḥ utthāya pratiṣṭhamānānām na kaḥ cit parasparam sambandhaḥ bhavati . \\
\noindent \textbf{(P\_2,2.24.4)} evañjātīyakam bhrātṛtvam nāma . \\
\noindent \textbf{(P\_2,2.24.4)} atra cet yuktaḥ ṣaṣṭhyarthaḥ dṛśyate iha api yuktaḥ dṛśyatām . \\
\noindent \textbf{(P\_2,2.24.4)} iha tarhi ardhatṛtīyāḥ ānīyantām iti ukte ardhasya ānayanam na prāpnoti . \\
\noindent \textbf{(P\_2,2.24.4)} astu tarhi ayam eva vigrahaḥ ardham tṛtīyam eṣām iti . \\
\noindent \textbf{(P\_2,2.24.4)} nanu ca uktam anusasāra pāṇḍavam . \\
\noindent \textbf{(P\_2,2.24.4)} saṅkarṣaṇadvitīyasya balam kṛṣṇasya vardhatām iti. dvayoḥ dvivacanam prāpnoti iti . \\
\noindent \textbf{(P\_2,2.24.4)} na eṣaḥ doṣaḥ . \\
\noindent \textbf{(P\_2,2.24.4)} ayam tīyantaḥ śabdaḥ asti eva pūraṇam . \\
\noindent \textbf{(P\_2,2.24.4)} asti sahāyavācī . \\
\noindent \textbf{(P\_2,2.24.4)} tat yaḥ sahāyavācī tasya idam grahaṇam . \\
\noindent \textbf{(P\_2,2.24.4)} asidvitīyaḥ asisahāyaḥ iti gamyate . \\
\noindent \textbf{(P\_2,2.24.4)} evam api ardhatṛtīyāḥ iti ekasmin ekavacanam iti ekavacanam prāpnoti . \\
\noindent \textbf{(P\_2,2.24.4)} ekārthāḥ hi samudāyāḥ bhavanti . \\
\noindent \textbf{(P\_2,2.24.4)} tat yathā śatam yūtham vanam iti . \\
\noindent \textbf{(P\_2,2.24.4)} astu tarhi ayam eva vigrahaḥ ardham tṛtīyam anayoḥ iti . \\
\noindent \textbf{(P\_2,2.24.4)} nanu ca uktam ardhatṛtīyāḥ ānīyantām iti ukte ardhasya ānayanam na prāpnoti iti . \\
\noindent \textbf{(P\_2,2.24.4)} na eṣaḥ doṣaḥ . \\
\noindent \textbf{(P\_2,2.24.4)} bhavati bahuvrīhau tadguṇasaṃvijñānam api . \\
\noindent \textbf{(P\_2,2.24.4)} tat yathā . \\
\noindent \textbf{(P\_2,2.24.4)} śuklavāsasam ānaya . \\
\noindent \textbf{(P\_2,2.24.4)} lohitoṣṇīṣāḥ ṛtvijaḥ pracaranti iti . \\
\noindent \textbf{(P\_2,2.24.4)} tadguṇaḥ ānīyate tadguṇāḥ ca pracaranti . \\
\noindent \textbf{(P\_2,2.24.4)} atha vā punaḥ astu ayam eva vigrahaḥ ardham tṛtīyam eṣām iti . \\
\noindent \textbf{(P\_2,2.24.4)} nanu ca uktam ekavacanam prāpnoti iti . \\
\noindent \textbf{(P\_2,2.24.4)} na eṣaḥ doṣaḥ . \\
\noindent \textbf{(P\_2,2.24.4)} saṅkhyā nāma iyam parapradhānā . \\
\noindent \textbf{(P\_2,2.24.4)} saṅkhyeyam anayā viśeṣyam . \\
\noindent \textbf{(P\_2,2.24.4)} yadi ca atra ekavacanam syāt saṅkhyeyam aviśeṣitam syāt . \\
\noindent \textbf{(P\_2,2.24.4)} iha tarhi ardhatṛtīyāḥ droṇāḥ iti ayam droṇaśabdaḥ samudāye pravṛttaḥ avayave na upapadyate . \\
\noindent \textbf{(P\_2,2.24.4)} na eṣaḥ doṣaḥ . \\
\noindent \textbf{(P\_2,2.24.4)} samudāyeṣu api śabdāḥ pravṛttāḥ avayaveṣu api vartante. tad yathā . \\
\noindent \textbf{(P\_2,2.24.4)} pūrve pañcālāḥ uttare pañcālāḥ tailam bhuktam ghṛtaṃ bhuktam śuklaḥ nīlaḥ kapilaḥ kṛṣṇaḥ iti . \\
\noindent \textbf{(P\_2,2.24.4)} evam ayam samudāye droṇaśabdaḥ pravṛttaḥ avayaveṣu api vartati . \\
\noindent \textbf{(P\_2,2.24.4)} kāmam tarhi anena eva hetunā yadā dvau droṇau ardhārḍhakam ca kartavyam ardhatṛtīyāḥ droṇāḥ iti . \\
\noindent \textbf{(P\_2,2.24.4)} na kartavyam . \\
\noindent \textbf{(P\_2,2.24.4)} samudāyeṣu api hi śabdāḥ pravṛttāḥ avayaveṣu api vartante. keṣu avayaveṣu . \\
\noindent \textbf{(P\_2,2.24.4)} yaḥ avayavaḥ tam samudāyam na vyabhicarati . \\
\noindent \textbf{(P\_2,2.24.4)} kam ca samudāyam na vyabhicarati . \\
\noindent \textbf{(P\_2,2.24.4)} ardhdroṇaḥ droṇam . \\
\noindent \textbf{(P\_2,2.24.4)} ardhāḍhakam punaḥ vyabhicarati . \\
\noindent \textbf{(P\_2,2.25)} dvitrāḥ tricaturāḥ iti kaḥ ayam samāsaḥ . \\
\noindent \textbf{(P\_2,2.25)} bahuvrīhiḥ iti āha . \\
\noindent \textbf{(P\_2,2.25)} kaḥ asya vigrahaḥ . \\
\noindent \textbf{(P\_2,2.25)} dvau vā trayaḥ vā iti . \\
\noindent \textbf{(P\_2,2.25)} bhavet yadā bahūnām ānayanam tadā bahuvacanam upapannam yadā tu khalu dvau ānīyete tadā na sidhyati . \\
\noindent \textbf{(P\_2,2.25)} tadā api sidhyati . \\
\noindent \textbf{(P\_2,2.25)} katham . \\
\noindent \textbf{(P\_2,2.25)} ke cit tāvat āhuḥ : anirjñāte arthe bahuvacanam prayoktavyam iti . \\
\noindent \textbf{(P\_2,2.25)} tat yathā : kati bhavataḥ putrāḥ . \\
\noindent \textbf{(P\_2,2.25)} kati bhavataḥ bhāryāḥ iti . \\
\noindent \textbf{(P\_2,2.25)} aparaḥ āha : dvau vā iti ukte trayaḥ vā iti gamyate . \\
\noindent \textbf{(P\_2,2.25)} trayaḥ vā iti ukte dvau vā iti gamyate . \\
\noindent \textbf{(P\_2,2.25)} sā eṣā pañcādhiṣṭhānā vāk . \\
\noindent \textbf{(P\_2,2.25)} atra yuktam bahuvacanam . \\
\noindent \textbf{(P\_2,2.25)} atha dvidaśāḥ tridaśāḥ iti kaḥ ayam samāsaḥ . \\
\noindent \textbf{(P\_2,2.25)} bahuvrīhiḥ iti āha . \\
\noindent \textbf{(P\_2,2.25)} kaḥ asya vigrahaḥ . \\
\noindent \textbf{(P\_2,2.25)} dviḥ daśa dviśaśāḥ iti . \\
\noindent \textbf{(P\_2,2.25)} saṅkhyāsamāse sujantatvāt saṅkhyāprasiddhiḥ . \\
\noindent \textbf{(P\_2,2.25)} saṅkhyāsamāse sujantatvāt saṅkhyā iti aprasiddhiḥ . \\
\noindent \textbf{(P\_2,2.25)} na hi sujantā saṅkhyā asti . \\
\noindent \textbf{(P\_2,2.25)} evam tarhi evam vigrahaḥ kariṣyate . \\
\noindent \textbf{(P\_2,2.25)} dvau daśatau dvidaśāḥ iti . \\
\noindent \textbf{(P\_2,2.25)} evam api atkārāntatvāt saṅkhyā iti aprasiddhiḥ . \\
\noindent \textbf{(P\_2,2.25)} na hi atkārāntā saṅkhyā asti . \\
\noindent \textbf{(P\_2,2.25)} astu tarhi ayam eva vigrahaḥ dviḥ daśa dviśaśāḥ iti . \\
\noindent \textbf{(P\_2,2.25)} nanu ca uktam saṅkhyāsamāse sujantatvāt saṅkhyā iti aprasiddhiḥ iti . \\
\noindent \textbf{(P\_2,2.25)} na vā asujantatvāt . \\
\noindent \textbf{(P\_2,2.25)} na vā eṣaḥ doṣaḥ . \\
\noindent \textbf{(P\_2,2.25)} kim kāraṇam . \\
\noindent \textbf{(P\_2,2.25)} asujantatvāt . \\
\noindent \textbf{(P\_2,2.25)} sujantā iti ucyate . \\
\noindent \textbf{(P\_2,2.25)} na ca atra sujantam paśyāmaḥ . \\
\noindent \textbf{(P\_2,2.25)} kim punaḥ kāraṇam vākye suc dṛśyate samāse tu na dṛśyate . \\
\noindent \textbf{(P\_2,2.25)} sujabhāvaḥ ahihitārthatvāt samāse . \\
\noindent \textbf{(P\_2,2.25)} samāse sucaḥ abhāvaḥ . \\
\noindent \textbf{(P\_2,2.25)} kim kāraṇam . \\
\noindent \textbf{(P\_2,2.25)} ahihitārthatvāt . \\
\noindent \textbf{(P\_2,2.25)} abhihitaḥ sujarthaḥ samāsena iti kṛtvā samāse suc na bhaviṣyati .kim ca bhoḥ sujarthe iti samāsaḥ ucyate . \\
\noindent \textbf{(P\_2,2.25)} na khalu sujarthe iti ucyate gamyate tu sujarthaḥ . \\
\noindent \textbf{(P\_2,2.25)} katham . \\
\noindent \textbf{(P\_2,2.25)} yāvatā saṅkhyeyaḥ yaḥ saṅkhyayā saṅkhyāyate saḥ ca kriyābhyāvṛttyarthaḥ . \\
\noindent \textbf{(P\_2,2.25)} saḥ ca uktaḥ samāsena iti kṛtvā samāse suc na bhaviṣyati . \\
\noindent \textbf{(P\_2,2.25)} aśiṣyaḥ saṅkhyottarapadaḥ saṅkhyeyavābhidhyāyitvāt . \\
\noindent \textbf{(P\_2,2.25)} aśiṣyaḥ saṅkhyottarapadaḥ bahuvrīhiḥ . \\
\noindent \textbf{(P\_2,2.25)} kim kāraṇam . \\
\noindent \textbf{(P\_2,2.25)} saṅkhyeyavābhidhyāyitvāt . \\
\noindent \textbf{(P\_2,2.25)} saṅkhyeyam vārthaḥ ca abhidīyate . \\
\noindent \textbf{(P\_2,2.25)} tatra anyapadārthe iti eva siddham . \\
\noindent \textbf{(P\_2,2.25)} bhavet siddham adhikaviṃśāḥ adhikatriṃśāḥ iti yatra etat vicāryate . \\
\noindent \textbf{(P\_2,2.25)} viśatyādayaḥ daśadarthe vā syuḥ parimāṇini vā iti . \\
\noindent \textbf{(P\_2,2.25)} idam tu na sidhyati adhikadaśāḥ iti yatra niyogataḥ saṅkhyeye eva vartate . \\
\noindent \textbf{(P\_2,2.25)} atha upadaśāḥ iti kaḥ ayam samāsaḥ . \\
\noindent \textbf{(P\_2,2.25)} bahuvrīhiḥ iti āha . \\
\noindent \textbf{(P\_2,2.25)} kaḥ asya vigrahaḥ . \\
\noindent \textbf{(P\_2,2.25)} daśānām samīpe upadaśāḥ iti . \\
\noindent \textbf{(P\_2,2.25)} kasya punaḥ sāmīpyam arthaḥ . \\
\noindent \textbf{(P\_2,2.25)} upasya . \\
\noindent \textbf{(P\_2,2.25)} yadi evam na anyapadārthaḥ bhavati . \\
\noindent \textbf{(P\_2,2.25)} tatra prathānirdiṣṭam saṅkhyāgrahaṇam śakyam akartum . \\
\noindent \textbf{(P\_2,2.25)} matvarthe vā pūrvasya vidhānāt . \\
\noindent \textbf{(P\_2,2.25)} atha vā matvarthe pūrvaḥ yogaḥ . \\
\noindent \textbf{(P\_2,2.25)} amatvarthaḥ ayam ārambhaḥ . \\
\noindent \textbf{(P\_2,2.25)} kababhāvārtham vā . \\
\noindent \textbf{(P\_2,2.25)} atha va kap mā bhūt iti . \\
\noindent \textbf{(P\_2,2.26,} 28) KA\_I,428.19-429.16 Ro\_II,725-727 {1/33}     diksamāsasahayogayoḥ ca antarālapradhānābhidhānāt . \\
\noindent \textbf{(P\_2,2.26,} 28) KA\_I,428.19-429.16 Ro\_II,725-727 {2/33}     diksamāsasahayogayoḥ ca aśiṣyaḥ bahuvrīhiḥ . \\
\noindent \textbf{(P\_2,2.26,} 28) KA\_I,428.19-429.16 Ro\_II,725-727 {3/33}     kim kāraṇam . \\
\noindent \textbf{(P\_2,2.26,} 28) KA\_I,428.19-429.16 Ro\_II,725-727 {4/33}     antarālapradhānābhidhānāt . \\
\noindent \textbf{(P\_2,2.26,} 28) KA\_I,428.19-429.16 Ro\_II,725-727 {5/33}     diksamāse sahayoge ca antarālam pradhānam ca abhidhīyate . \\
\noindent \textbf{(P\_2,2.26,} 28) KA\_I,428.19-429.16 Ro\_II,725-727 {6/33}     tatra anyapadārthe iti eva siddham . \\
\noindent \textbf{(P\_2,2.26,} 28) KA\_I,428.19-429.16 Ro\_II,725-727 {7/33}     yadi evam dakṣiṇapūrvā dik samānādhikaraṇalakṣaṇaḥ puṃvadbhāvaḥ na prāpnoti . \\
\noindent \textbf{(P\_2,2.26,} 28) KA\_I,428.19-429.16 Ro\_II,725-727 {8/33}     adya punaḥ iyam sā eva dakṣiṇā sā eva pūrvā iti kṛtvā samānādhikaraṇalakṣaṇaḥ puṃvadbhāvaḥ siddhaḥ bhavati . \\
\noindent \textbf{(P\_2,2.26,} 28) KA\_I,428.19-429.16 Ro\_II,725-727 {9/33}     na sidhyati . \\
\noindent \textbf{(P\_2,2.26,} 28) KA\_I,428.19-429.16 Ro\_II,725-727 {10/33}     bhāṣitapuṃskasya puṃvadbhāvaḥ . \\
\noindent \textbf{(P\_2,2.26,} 28) KA\_I,428.19-429.16 Ro\_II,725-727 {11/33}     na ca etau bhāṣitapuṃskau . \\
\noindent \textbf{(P\_2,2.26,} 28) KA\_I,428.19-429.16 Ro\_II,725-727 {12/33}     nanu ca bhoḥ dakṣiṇaśabdaḥ pūrvaśabdaḥ ca puṃsi bhāṣyete . \\
\noindent \textbf{(P\_2,2.26,} 28) KA\_I,428.19-429.16 Ro\_II,725-727 {13/33}     samānāyām ākṛtau yat bhāṣitapuṃskam . \\
\noindent \textbf{(P\_2,2.26,} 28) KA\_I,428.19-429.16 Ro\_II,725-727 {14/33}     ākṛtyantare ca etau bhāṣitapuṃskau . \\
\noindent \textbf{(P\_2,2.26,} 28) KA\_I,428.19-429.16 Ro\_II,725-727 {15/33}     dakṣiṇā pūrvā iti dikśabdau . \\
\noindent \textbf{(P\_2,2.26,} 28) KA\_I,428.19-429.16 Ro\_II,725-727 {16/33}     dakṣiṇaḥ pūrvaḥ iti vyavasthāśabdau . \\
\noindent \textbf{(P\_2,2.26,} 28) KA\_I,428.19-429.16 Ro\_II,725-727 {17/33}     yadi punaḥ dikśabdāḥ api vyavasthāśabdāḥ syuḥ . \\
\noindent \textbf{(P\_2,2.26,} 28) KA\_I,428.19-429.16 Ro\_II,725-727 {18/33}     katham yāni digapadiṣṭāni kāryāṇi . \\
\noindent \textbf{(P\_2,2.26,} 28) KA\_I,428.19-429.16 Ro\_II,725-727 {19/33}     yadā diśaḥ vyavasthām vakṣyanti . \\
\noindent \textbf{(P\_2,2.26,} 28) KA\_I,428.19-429.16 Ro\_II,725-727 {20/33}     yadi tari yaḥ yaḥ diśi vartate saḥ saḥ dikśabdaḥ ramaṇīyādiṣu atiprasaṅgaḥ bhavati . \\
\noindent \textbf{(P\_2,2.26,} 28) KA\_I,428.19-429.16 Ro\_II,725-727 {21/33}     ramaṇīyā dik śobhanā dik iti . \\
\noindent \textbf{(P\_2,2.26,} 28) KA\_I,428.19-429.16 Ro\_II,725-727 {22/33}     atha matam etat diśi dṛṣṭaḥ digdṛṣṭaḥ digdiṣṭaḥ śabdaḥ dikśabdaḥ diśam yaḥ na vyabhicarati iti ramaṇīyādiṣu atiprasaṅgaḥ na bhavati . \\
\noindent \textbf{(P\_2,2.26,} 28) KA\_I,428.19-429.16 Ro\_II,725-727 {23/33}     puṃvadbhāvaḥ tu prāpnoti . \\
\noindent \textbf{(P\_2,2.26,} 28) KA\_I,428.19-429.16 Ro\_II,725-727 {24/33}     evam tarhi sarvanāmnaḥ vṛttimātre puṃvadbhāvaḥ vaktavyaḥ dakṣiṇottarapūrvāṇām iti evamartham . \\
\noindent \textbf{(P\_2,2.26,} 28) KA\_I,428.19-429.16 Ro\_II,725-727 {25/33}     evam ca kṛtvā dik diksamāsasahayogayoḥ ca antarālapradhānābhidhānāt iti eva . \\
\noindent \textbf{(P\_2,2.26,} 28) KA\_I,428.19-429.16 Ro\_II,725-727 {26/33}     nanu ca uktam dakṣiṇapūrvā dik samānādhikaraṇalakṣaṇaḥ puṃvadbhāvaḥ na prāpnoti iti . \\
\noindent \textbf{(P\_2,2.26,} 28) KA\_I,428.19-429.16 Ro\_II,725-727 {27/33}     na eṣaḥ doṣaḥ . \\
\noindent \textbf{(P\_2,2.26,} 28) KA\_I,428.19-429.16 Ro\_II,725-727 {28/33}     sarvanāmnaḥ vṛttimātre puṃvadbhāvena parihṛtam . \\
\noindent \textbf{(P\_2,2.26,} 28) KA\_I,428.19-429.16 Ro\_II,725-727 {29/33}     matvarthe vā pūrvasya vidhānāt . \\
\noindent \textbf{(P\_2,2.26,} 28) KA\_I,428.19-429.16 Ro\_II,725-727 {30/33}     atha vā matvarthe pūrvaḥ yogaḥ . \\
\noindent \textbf{(P\_2,2.26,} 28) KA\_I,428.19-429.16 Ro\_II,725-727 {31/33}     amatvarthaḥ ayam ārambhaḥ . \\
\noindent \textbf{(P\_2,2.26,} 28) KA\_I,428.19-429.16 Ro\_II,725-727 {32/33}     kababhāvārtham vā . \\
\noindent \textbf{(P\_2,2.26,} 28) KA\_I,428.19-429.16 Ro\_II,725-727 {33/33}     atha va kap mā bhūt iti . \\
\noindent \textbf{(P\_2,2.27)} tṛtīyāsaptamyanteṣu ca kriyābhidhānāt . \\
\noindent \textbf{(P\_2,2.27)} tṛtīyāsaptamyanteṣu ca kriyābhidhānāt aśiṣyaḥ bahuvrīhiḥ . \\
\noindent \textbf{(P\_2,2.27)} kim kāraṇam . \\
\noindent \textbf{(P\_2,2.27)} kriyābhidhānāt . \\
\noindent \textbf{(P\_2,2.27)} kriyā abhidhīyate . \\
\noindent \textbf{(P\_2,2.27)} tatra anyapadārthe iti eva siddham . \\
\noindent \textbf{(P\_2,2.27)} na vā ekaśeṣapratiṣedhārtham . \\
\noindent \textbf{(P\_2,2.27)} na vā aśiṣyaḥ . \\
\noindent \textbf{(P\_2,2.27)} kim kāraṇam . \\
\noindent \textbf{(P\_2,2.27)} ekaśeṣapratiṣedhārtham idam vaktavyam . \\
\noindent \textbf{(P\_2,2.27)} pūrvadīrghārtham ca . \\
\noindent \textbf{(P\_2,2.27)} pūrvadīrghārtham ca idam vaktavyam . \\
\noindent \textbf{(P\_2,2.27)} keśākeśi . \\
\noindent \textbf{(P\_2,2.27)} syāt etat prayojanam yadi niyogataḥ asya anena eva dīrghatvam syāt . \\
\noindent \textbf{(P\_2,2.27)} atha idānīm anyeṣām api dṛśyate iti dīrghatvam na prayojanam bhavati . \\
\noindent \textbf{(P\_2,2.27)} matvarthe vā pūrvasya vidhānāt . \\
\noindent \textbf{(P\_2,2.27)} atha va matvarthe pūrvaḥ yogaḥ . \\
\noindent \textbf{(P\_2,2.27)} amatvarthaḥ ayam ārambhaḥ . \\
\noindent \textbf{(P\_2,2.27)} kababhāvārtham vā . \\
\noindent \textbf{(P\_2,2.27)} atha va kap mā bhūt iti . \\
\noindent \textbf{(P\_2,2.29.1)} cārthe iti ucyate caḥ ca avyayam . \\
\noindent \textbf{(P\_2,2.29.1)} tena samāsasya avyayasañjñā prāpnoti . \\
\noindent \textbf{(P\_2,2.29.1)} na eṣaḥ doṣaḥ . \\
\noindent \textbf{(P\_2,2.29.1)} pāṭhena avyayasañjñā kriyate . \\
\noindent \textbf{(P\_2,2.29.1)} na ca samāsaḥ tatra paṭhyate . \\
\noindent \textbf{(P\_2,2.29.1)} pāṭhena api avyayasañjñāyām satyām abhideheyavat liṅgavacanāni bhavanti . \\
\noindent \textbf{(P\_2,2.29.1)} yaḥ ca iha arthaḥ abhidhīyate na tasya liṅgasaṅkhyābhyām yogaḥ asti . \\
\noindent \textbf{(P\_2,2.29.1)} na idam vācanikam aliṅgatā asaṅkhyatā va . \\
\noindent \textbf{(P\_2,2.29.1)} kim tarhi . \\
\noindent \textbf{(P\_2,2.29.1)} svābhāvikam etat . \\
\noindent \textbf{(P\_2,2.29.1)} tat yathā : samānam īhamānānām adhīyānānām ca ke cit arthaiḥ yujyante apare na . \\
\noindent \textbf{(P\_2,2.29.1)} na ca idānīm kaḥ cit arthavān iti kṛtvā sarvaiḥ arthavadbhiḥ śakyam bhavitum kaḥ cit anarthakaḥ iti kṛtvā sarvaiḥ anarthakaiḥ . \\
\noindent \textbf{(P\_2,2.29.1)} tatra kim asmābhiḥ śakyam kartum . \\
\noindent \textbf{(P\_2,2.29.1)} yat prāk samāsāt cārthasya liṅgasaṅkhyābhyām yogaḥ na asti samāse ca bhavati svābhāvikam etat . \\
\noindent \textbf{(P\_2,2.29.1)} atha vā āśrayataḥ liṅgavacanāni bhaviṣyanti . \\
\noindent \textbf{(P\_2,2.29.1)} guṇavacanānām hi śabdānām āśrayataḥ liṅgavacanāni bhavanti . \\
\noindent \textbf{(P\_2,2.29.1)} tat yathā śuklam vastram , śuklā śāṭī śuklaḥ kambalaḥ , śuklau kambalau śuklāḥ kambalāḥ iti . \\
\noindent \textbf{(P\_2,2.29.1)} yat asau dravyam śritaḥ bhavati guṇaḥ tasya yat liṅgam vacanam ca tat guṇasya api bhavati . \\
\noindent \textbf{(P\_2,2.29.1)} evam iha api yat asau dravyam śritaḥ bhavati samāsaḥ tasya yat liṅgam vacanam ca tat samāsasya api bhaviṣyati . \\
\noindent \textbf{(P\_2,2.29.1)} atha iha kasmāt na bhavati . \\
\noindent \textbf{(P\_2,2.29.1)} yājñikaḥ ca ayam vaiyākaraṇaḥ ca . \\
\noindent \textbf{(P\_2,2.29.1)} kaṭhaḥ ca ayam bahvṛcaḥ ca . \\
\noindent \textbf{(P\_2,2.29.1)} aukthikaḥ ca ayam mīmāṃsakaḥ ca iti . \\
\noindent \textbf{(P\_2,2.29.1)} śeṣaḥ iti vartate . \\
\noindent \textbf{(P\_2,2.29.1)} aśeṣatvāt na bhaviṣyati . \\
\noindent \textbf{(P\_2,2.29.1)} yadi śeṣaḥ iti vartate upāsnātam sthūlasiktam tūṣṇīṅgaṅgam mahāhradam droṇam cet aśakaḥ gantum mā tvā tāptām kṛtākṛte iti etat na sidhyati . \\
\noindent \textbf{(P\_2,2.29.1)} na eṣaḥ doṣaḥ . \\
\noindent \textbf{(P\_2,2.29.1)} anyat hi kṛtam anyat akṛtam . \\
\noindent \textbf{(P\_2,2.29.2)} cārthe dvandvavacane asamāse api cārthasampratyayāt aniṣṭaprasaṅgaḥ . \\
\noindent \textbf{(P\_2,2.29.2)} cārthe dvandvavacane asamāse api cārthasampratyayāt aniṣṭam prāpnoti . \\
\noindent \textbf{(P\_2,2.29.2)} ahaḥ ahaḥ nayamānaḥ gām aśvam puruṣam paśum vaivasvataḥ na tṛpyati surāyāḥ iva durmadī indraḥ tvaṣṭā varuṇaḥ vāyuḥ ādityaḥ iti . \\
\noindent \textbf{(P\_2,2.29.2)} siddham tu yugapadadhikaraṇavacane dvandvavacanāt . \\
\noindent \textbf{(P\_2,2.29.2)} siddham etat . \\
\noindent \textbf{(P\_2,2.29.2)} katham . \\
\noindent \textbf{(P\_2,2.29.2)} yugapadadhikaraṇavacane dvandvaḥ bhavati iti vaktavyam . \\
\noindent \textbf{(P\_2,2.29.2)} tatra puṃvadbhāvapratiṣedhaḥ . \\
\noindent \textbf{(P\_2,2.29.2)} tatra etasmin lakṣaṇe puṃvadbhāvasya pratiṣedhaḥ vaktavyaḥ . \\
\noindent \textbf{(P\_2,2.29.2)} paṭvīmṛdvyau . \\
\noindent \textbf{(P\_2,2.29.2)} samānādhikaraṇalakṣaṇaḥ puṃvadbhāvaḥ prāpnoti . \\
\noindent \textbf{(P\_2,2.29.2)} vipratiṣiddheṣu ca anupapattiḥ . \\
\noindent \textbf{(P\_2,2.29.2)} vipratiṣiddheṣu yugapadadhikaraṇavacatāyāḥ anupapattiḥ . \\
\noindent \textbf{(P\_2,2.29.2)} śītoṣṇe sukhaduḥkhe jananamaraṇe . \\
\noindent \textbf{(P\_2,2.29.2)} kim kāraṇam . \\
\noindent \textbf{(P\_2,2.29.2)} sukhapratighātena hi duḥkham duḥkapratighātena ca sukham . \\
\noindent \textbf{(P\_2,2.29.2)} yat tāvat ucyate tatra puṃvadbhāvapratiṣedhaḥ iti . \\
\noindent \textbf{(P\_2,2.29.2)} idam tāvat ayam praṣṭavyaḥ : atha iha kasmāt na bhavati . \\
\noindent \textbf{(P\_2,2.29.2)} darśanīyāyāḥ mātā darśanīyāmātā iti . \\
\noindent \textbf{(P\_2,2.29.2)} atha matam etat prāk samāsāt yatra sāmānādhikaraṇyam tatra puṃvadbhāvaḥ bhavati iti iha api na doṣaḥ bhavati . \\
\noindent \textbf{(P\_2,2.29.2)} yad api ucyate vipratiṣiddheṣu ca anupapattiḥ iti . \\
\noindent \textbf{(P\_2,2.29.2)} sarve eva hi śabdāḥ vipratiṣiddhāḥ . \\
\noindent \textbf{(P\_2,2.29.2)} iha api plakṣanyagrodhau iti plakṣaśabdaḥ prayujyamānaḥ plakṣārtham sampratyāyayati nyagrodhārtham nivartayati . \\
\noindent \textbf{(P\_2,2.29.2)} nyagrodhaśabdaḥ prayujyamānaḥ nyagrodhārtham sampratyāyayati plakṣārtham nivartayati . \\
\noindent \textbf{(P\_2,2.29.2)} atra cet yuktā yugapat adhikaraṇvacanatā dṛśyate iha api yuktā dṛśyatām . \\
\noindent \textbf{(P\_2,2.29.2)} evam api śabdapaurvāparyaprayogāt arthapaurvāparyābhidhānam . \\
\noindent \textbf{(P\_2,2.29.2)} śabdapaurvāparyaprayogāt arthapaurvāparyābhidhānam prāpnoti . \\
\noindent \textbf{(P\_2,2.29.2)} ataḥ kim . \\
\noindent \textbf{(P\_2,2.29.2)} yugapatadhikaraṇavacanatāyāḥ anupapattiḥ . \\
\noindent \textbf{(P\_2,2.29.2)} plakṣanyagrodhau plakṣanyagrodhāḥ iti . \\
\noindent \textbf{(P\_2,2.29.2)} yatha eva hi śabdānām paurvāparyam tadvat arthānām api bhavitavyam . \\
\noindent \textbf{(P\_2,2.29.2)} śabdapaurvāparyaprayogāt arthapaurvāparyābhidhānam iti cet dvivacanabahuvacanānupapattiḥ . \\
\noindent \textbf{(P\_2,2.29.2)} śabdapaurvāparyaprayogāt arthapaurvāparyābhidhānam iti cet dvivacanabahuvacanānupapattiḥ : plakṣanyagrodhau plakṣanyagrodhāḥ iti . \\
\noindent \textbf{(P\_2,2.29.2)} plakṣaśabdaḥ sārthakaḥ nivṛttaḥ nyagrodhaśabdaḥ upasthitaḥ ekārthaḥ tasya ekārthatvāt ekavacanam eva prāpnoti . \\
\noindent \textbf{(P\_2,2.29.2)} vigrahe tu yugapadvacanam jñāpakam yugapadvacanasya . \\
\noindent \textbf{(P\_2,2.29.2)} vigrahe khalu api yugapadvacanatā dṛśyate : dyavā ha kṣamā . \\
\noindent \textbf{(P\_2,2.29.2)} dyavā cit asmai pṛthivī namete iti . \\
\noindent \textbf{(P\_2,2.29.2)} kim etat . \\
\noindent \textbf{(P\_2,2.29.2)} yugapadadhikaraṇavacanatāyāḥ upodbalakam . \\
\noindent \textbf{(P\_2,2.29.2)} vigrahe kila nāma yugapadadhikaraṇavacanatā syāt kim punaḥ samāse . \\
\noindent \textbf{(P\_2,2.29.2)} samudāyāt siddham . \\
\noindent \textbf{(P\_2,2.29.2)} samudāyāt siddham etat . \\
\noindent \textbf{(P\_2,2.29.2)} kim etat samudāyāt siddham iti . \\
\noindent \textbf{(P\_2,2.29.2)} dvivacanabahuvacanaprasiddhiḥ iti coditam . \\
\noindent \textbf{(P\_2,2.29.2)} tasya ayam parihāraḥ . \\
\noindent \textbf{(P\_2,2.29.2)} samudāyāt siddham iti cet na ekārthatvāt samudāyasya . \\
\noindent \textbf{(P\_2,2.29.2)} samudāyāt siddham iti cet tat na . \\
\noindent \textbf{(P\_2,2.29.2)} kim kāraṇam . \\
\noindent \textbf{(P\_2,2.29.2)} ekārthatvāt samudāyasya . \\
\noindent \textbf{(P\_2,2.29.2)} ekārthāḥ hi samudāyāḥ bhavanti . \\
\noindent \textbf{(P\_2,2.29.2)} tat yathā śatam yūtham vanam iti . \\
\noindent \textbf{(P\_2,2.29.2)} na aikārthyam . \\
\noindent \textbf{(P\_2,2.29.2)} na ayam ekārthaḥ . \\
\noindent \textbf{(P\_2,2.29.2)} kim tarhi . \\
\noindent \textbf{(P\_2,2.29.2)} dvyarthaḥ bahvarthaḥ ca . \\
\noindent \textbf{(P\_2,2.29.2)} plakṣaḥ api dvyarthaḥ nyagrodhaḥ api dvyarthaḥ . \\
\noindent \textbf{(P\_2,2.29.2)} yadi tarhi plakṣaḥ api dvyarthaḥ nyagrodhaḥ api dvyarthaḥ tayoḥ anekārthatvāt bahuvacanaprasaṅgaḥ . \\
\noindent \textbf{(P\_2,2.29.2)} tayoḥ anekārthatvāt bahuṣu bahuvacanam iti bahuvacanam prāpnoti . \\
\noindent \textbf{(P\_2,2.29.2)} tayoḥ anekārthatvāt bahuvacanaprasaṅgaḥ iti cet na bahutvābhāvāt . \\
\noindent \textbf{(P\_2,2.29.2)} tayoḥ anekārthatvāt bahuvacanaprasaṅgaḥ iti cet tat na . \\
\noindent \textbf{(P\_2,2.29.2)} kim kāraṇam . \\
\noindent \textbf{(P\_2,2.29.2)} bahutvābhāvāt . \\
\noindent \textbf{(P\_2,2.29.2)} na atra bahutvam asti . \\
\noindent \textbf{(P\_2,2.29.2)} kim ucyate bahutvābhāvāt iti yāvatā idānīm eva uktam plakṣaḥ api dvyarthaḥ nyagrodhaḥ api dvyarthaḥ iti . \\
\noindent \textbf{(P\_2,2.29.2)} yābhyām eva atra ekaḥ dvyarthaḥ tābhyām eva aparaḥ api . \\
\noindent \textbf{(P\_2,2.29.2)} yadi evam anyavācakena anyasya vacanānupapattiḥ . \\
\noindent \textbf{(P\_2,2.29.2)} anyavācakena śabdena anyasya vacanam na upapadyate . \\
\noindent \textbf{(P\_2,2.29.2)} anyavācakena anyasya vacanānupapattiḥ iti cet plakṣasya nyagrodhatvāt nyagrodhasya plakṣatvāt svaśabdena abhidhānam . \\
\noindent \textbf{(P\_2,2.29.2)} anyavācakena anyasya vacanānupapattiḥ iti cet ucyate tat na . \\
\noindent \textbf{(P\_2,2.29.2)} kim kāraṇam . \\
\noindent \textbf{(P\_2,2.29.2)} plakṣasya nyagrodhatvāt nyagrodhasya plakṣatvāt svaśabdena abhidhānam . \\
\noindent \textbf{(P\_2,2.29.2)} plakṣaḥ api nyagrodhaḥ nyagrodhaḥ api plakṣaḥ . \\
\noindent \textbf{(P\_2,2.29.2)} katham punaḥ plakṣaḥ api nyagrodhaḥ nyagrodhaḥ api plakṣaḥ syāt yāvatā kāraṇāt dravye śabdaniveśaḥ . \\
\noindent \textbf{(P\_2,2.29.2)} kāraṇāt dravye śabdaniveśaḥ iti cet tulyakāraṇatvāt siddham . \\
\noindent \textbf{(P\_2,2.29.2)} kāraṇāt dravye śabdaniveśaḥ iti cet evam ucyate : tat na tulyakāraṇatvāt siddham . \\
\noindent \textbf{(P\_2,2.29.2)} tulyam hi kāraṇam . \\
\noindent \textbf{(P\_2,2.29.2)} yadi tāvat prakṣarati iti plakṣaḥ syān nyagrodhe api etat bhavati . \\
\noindent \textbf{(P\_2,2.29.2)} tathā yadi nyak rohati iti nyagrodhaḥ plakṣe api etat bhavati . \\
\noindent \textbf{(P\_2,2.29.2)} darśanam vai hetuḥ na ca nyagrodhe plakṣaśabdaḥ dṛśyate . \\
\noindent \textbf{(P\_2,2.29.2)} darśanam hetuḥ iti cet tulyam . \\
\noindent \textbf{(P\_2,2.29.2)} darśanam hetuḥ iti cet tulyam etat bhavati . \\
\noindent \textbf{(P\_2,2.29.2)} plakṣe api nyagrodhaśabdaḥ dṛśyatām . \\
\noindent \textbf{(P\_2,2.29.2)} tulyam hi kāraṇam . \\
\noindent \textbf{(P\_2,2.29.2)} na vai loke eṣaḥ sampratyayaḥ bhavati . \\
\noindent \textbf{(P\_2,2.29.2)} na hi plakṣaḥ ānīyatām iti ukte nyragrodhaḥ ānīyate . \\
\noindent \textbf{(P\_2,2.29.2)} tadviṣayam ca . \\
\noindent \textbf{(P\_2,2.29.2)} tadviṣayam ca etat draṣṭavyam plakṣasya nyagrodhatvam . \\
\noindent \textbf{(P\_2,2.29.2)} kiṃviṣayam . \\
\noindent \textbf{(P\_2,2.29.2)} dvandvaviṣayam . \\
\noindent \textbf{(P\_2,2.29.2)} yuktam punaḥ yat niyataviṣayāḥ nāma śabdāḥ syuḥ . \\
\noindent \textbf{(P\_2,2.29.2)} bāḍham yuktam . \\
\noindent \textbf{(P\_2,2.29.2)} anyatra api tadviṣayadarśanāt . \\
\noindent \textbf{(P\_2,2.29.2)} anyatra api hi niyataviṣayāḥ śabdāḥ dṛśyante . \\
\noindent \textbf{(P\_2,2.29.2)} tat yathā : samāne rakte varṇe gauḥ lohitaḥ iti bhavati āsvaḥ śoṇaḥ iti . \\
\noindent \textbf{(P\_2,2.29.2)} samāne ca kāle varṇe gauḥ kṛṣṇaḥ iti bhavati aśvaḥ hemaḥ iti . \\
\noindent \textbf{(P\_2,2.29.2)} samāne ca śukle varṇe gauḥ śvetaḥ iti bhavati aśvaḥ karkaḥ iti . \\
\noindent \textbf{(P\_2,2.29.2)} yadi tarhi plakṣaḥ api nyagraodhaḥ nyagrodhaḥ api plakṣaḥ ekena uktatvāt aparasya prayogaḥ anupapannaḥ . \\
\noindent \textbf{(P\_2,2.29.2)} ekena uktatvāt tasya arthasya aparasya prayogaḥ na upapadyate . \\
\noindent \textbf{(P\_2,2.29.2)} plakṣeṇa nyagrodhasya nyagrodhaprayogaḥ . \\
\noindent \textbf{(P\_2,2.29.2)} ekena uktatvāt aparasya prayogaḥ anupapannaḥ iti cet anuktatvāt plakṣeṇa nyagrodhasya nyagrodhaprayogaḥ . ekena uktatvāt aparasya prayogaḥ anupapannaḥ iti cet tat na . \\
\noindent \textbf{(P\_2,2.29.2)} kim kāraṇam . \\
\noindent \textbf{(P\_2,2.29.2)} anuktatvāt plakṣeṇa nyagrodhasya nyagrodhaprayogaḥ . \\
\noindent \textbf{(P\_2,2.29.2)} anuktaḥ plakṣeṇa nyagrodhārthaḥ iti kṛtvā nyagrodhaśabdaḥ prayujyate . \\
\noindent \textbf{(P\_2,2.29.2)} katham anuktaḥ yāvatā idānīm eva uktam plakṣaḥ api nyagrodhaḥ nyagrodhaḥ api plakṣaḥ iti . \\
\noindent \textbf{(P\_2,2.29.2)} sahabhūtau etau anyonyasya artham āhatuḥ na pṛthagbhūtau . \\
\noindent \textbf{(P\_2,2.29.2)} kim punaḥ kāraṇam sahabhūtau etau anyonyasya artham āhatuḥ na pṛthagbhūtau . \\
\noindent \textbf{(P\_2,2.29.2)} abhidhānam punaḥ svābhāvikam . \\
\noindent \textbf{(P\_2,2.29.2)} svābhāvikam abhidhānam . \\
\noindent \textbf{(P\_2,2.29.2)} atha vā iha kau cit prāthamakalpikau plakṣanyagrodhau kau cit kriyayā vā guṇena va plakṣaḥ iva ayam plakṣaḥ , nyagrodhaḥ iva ayam nyagrodhaḥ iti . \\
\noindent \textbf{(P\_2,2.29.2)} tatra plakṣau iti ukte sandehaḥ syāt : kim imau plakṣau āhosvit plakṣanyagrodhau iti . \\
\noindent \textbf{(P\_2,2.29.2)} tatra asandehārtham nyagrodhaśabdaḥ prayujyate . \\
\noindent \textbf{(P\_2,2.29.2)} iyam yugapadadhikaraṇavacanata nāma duḥkhā ca durupapādā ca . \\
\noindent \textbf{(P\_2,2.29.2)} yat ca api asyā nibandhanam uktam dyāvā ha kṣāmā iti tat api chāndasam . \\
\noindent \textbf{(P\_2,2.29.2)} tatra supām supaḥ bhavanti iti eva siddham . \\
\noindent \textbf{(P\_2,2.29.2)} sūtram ca bhidyate . \\
\noindent \textbf{(P\_2,2.29.2)} yathānyāsam eva astu . \\
\noindent \textbf{(P\_2,2.29.2)} nanu ca uktam cārthe dvandvavacane asamāse api cārthasampratyayāt aniṣṭaprasaṅgaḥ iti . \\
\noindent \textbf{(P\_2,2.29.2)} na eṣaḥ doṣaḥ . \\
\noindent \textbf{(P\_2,2.29.2)} iha ce dvandve iti iyatā siddham . \\
\noindent \textbf{(P\_2,2.29.2)} katham punaḥ ce nāma vṛttiḥ syāt . \\
\noindent \textbf{(P\_2,2.29.2)} śabdaḥ hi eṣaḥ . \\
\noindent \textbf{(P\_2,2.29.2)} śabde asambhavāt arthe kāryam vijñāsyate . \\
\noindent \textbf{(P\_2,2.29.2)} saḥ ayam evam siddhe sati yat arthagrahaṇam karoti tasya etat prayojanam evam yathā vijñāyeta cena kṛtaḥ artaḥ cārthaḥ iti . \\
\noindent \textbf{(P\_2,2.29.2)} kaḥ punaḥ cena kṛtaḥ arthaḥ . \\
\noindent \textbf{(P\_2,2.29.2)} samuccayaḥ anvācayaḥ itaretarayogaḥ samāhāraḥ iti . \\
\noindent \textbf{(P\_2,2.29.2)} samuccayaḥ . \\
\noindent \textbf{(P\_2,2.29.2)} plakṣaḥ ca iti ukte gamyate etat nyagrodhaḥ ca iti . \\
\noindent \textbf{(P\_2,2.29.2)} anvācayaḥ . \\
\noindent \textbf{(P\_2,2.29.2)} plakṣaḥ ca iti ukte gamyate etat sāpaekṣaḥ ayam prayujyate iti . \\
\noindent \textbf{(P\_2,2.29.2)} itaretarayogaḥ . \\
\noindent \textbf{(P\_2,2.29.2)} plakṣaḥ ca nyagrodhaḥ ca iti ukte gamyate etat plakṣaḥ api nyagrodhasahāyaḥ nyagrodhaḥ api plakṣasahāyaḥ iti . \\
\noindent \textbf{(P\_2,2.29.2)} samāhāre api kriyate plakṣanyagrodham iti . \\
\noindent \textbf{(P\_2,2.29.2)} tatra ayam api arthaḥ dvandvaikavadbhāvaḥ na paṭhitavyaḥ bhavati . \\
\noindent \textbf{(P\_2,2.29.2)} samāhārasya ekatvāt eva siddham . \\
\noindent \textbf{(P\_2,2.29.3)} ekādaśa dvādaśa iti kaḥ ayam samāsaḥ . \\
\noindent \textbf{(P\_2,2.29.3)} ekādīnām daśādibhiḥ dvandvaḥ . \\
\noindent \textbf{(P\_2,2.29.3)} ekādīnām daśādibhiḥ dvandvaḥ samāsaḥ . \\
\noindent \textbf{(P\_2,2.29.3)} ekādīnām daśādibhiḥ dvandvaḥ iti cet viṃśatyādiṣu vacanaprasaṅgaḥ . \\
\noindent \textbf{(P\_2,2.29.3)} ekādīnām daśādibhiḥ dvandvaḥ iti cet viṃśatyādiṣu vacanam prāpnoti . \\
\noindent \textbf{(P\_2,2.29.3)} ekaviṃśatiḥ dvāviṃśatiḥ . \\
\noindent \textbf{(P\_2,2.29.3)} siddham tu adhikāntā saṅkhya saṅkhyayā samānādhikaraṇādhikāre adhikalopaḥ ca . siddham etat . \\
\noindent \textbf{(P\_2,2.29.3)} katham . \\
\noindent \textbf{(P\_2,2.29.3)} samānādhikaraṇādhikāre vaktavyam adhikāntā saṅkhya saṅkhyayā saha samasyate adhikaśabdasya ca lopaḥ bhavati iti . \\
\noindent \textbf{(P\_2,2.29.3)} ekādhikā viṃśatiḥ ekaviṃśatiḥ dvyadhikā viṃśatiḥ dvāviṃśatiḥ . \\
\noindent \textbf{(P\_2,2.29.3)} yadi samānādhikaraṇaḥ svaraḥ na sidhyati . \\
\noindent \textbf{(P\_2,2.29.3)} yat hi tat saṅkhyā pūrvapadam prakṛtisvaram bhavati iti dvandve iti tat . \\
\noindent \textbf{(P\_2,2.29.3)} kim punaḥ kāraṇam dvandve iti evam tat . \\
\noindent \textbf{(P\_2,2.29.3)} iha mā bhūt śatasahasram iti . \\
\noindent \textbf{(P\_2,2.29.3)} astu tarhi dvandvaḥ . \\
\noindent \textbf{(P\_2,2.29.3)} nanu ca uktam ekādīnām daśādibhiḥ dvandvaḥ iti cet viṃśatyādiṣu vacanaprasaṅgaḥ iti. na eṣaḥ doṣaḥ . \\
\noindent \textbf{(P\_2,2.29.3)} sarvaḥ dvandvaḥ vibhāṣā ekavat bhavati . \\
\noindent \textbf{(P\_2,2.29.3)} yadā tarhi ekavacanam tadā napuṃsakaliṅgam prāpnoti . \\
\noindent \textbf{(P\_2,2.29.3)} liṅgam aśiṣyam lokāśrayatvāt liṅgasya . \\
\noindent \textbf{(P\_2,2.30)} kimartham idam ucyate . \\
\noindent \textbf{(P\_2,2.30)} upasarjanasya pūrvavacanam paraprayoganivṛttyartham . \\
\noindent \textbf{(P\_2,2.30)} upasarjanasya pūrvavacanam kriyate paraprayogaḥ mā bhūt iti . \\
\noindent \textbf{(P\_2,2.30)} na vā aniṣṭadarśanāt . \\
\noindent \textbf{(P\_2,2.30)} na vā etat prayojanam asti . \\
\noindent \textbf{(P\_2,2.30)} kim kāraṇam . \\
\noindent \textbf{(P\_2,2.30)} aniṣṭadarśanāt . \\
\noindent \textbf{(P\_2,2.30)} na hi kim cit aniṣṭam dṛśyate . \\
\noindent \textbf{(P\_2,2.30)} na hi kaḥ cit rājapuruṣaḥ iti prayoktavye puruṣarājaḥ iti prayuṅkte . \\
\noindent \textbf{(P\_2,2.30)} yadi ca aniṣṭam drśyete tataḥ yatnārtham syāt . \\
\noindent \textbf{(P\_2,2.30)} atha yatra dve ṣaṣṭhyante bhavataḥ kasmāt tatra pradhānasya pūrvanipātaḥ na bhavati . \\
\noindent \textbf{(P\_2,2.30)} rājñaḥ puruṣasya rājapuruṣasya iti . \\
\noindent \textbf{(P\_2,2.30)} ṣaṣṭhyantayoḥ samāse arthābhedāt pradhānasya apūrvanipātaḥ . \\
\noindent \textbf{(P\_2,2.30)} ṣaṣṭhyantayoḥ samāse arthābhedāt pradhānasya pūrvanipātaḥ na bhaviṣyati . \\
\noindent \textbf{(P\_2,2.30)} evam na ca idam akṛtam bhavati upasarjanam pūrvam iti arthaḥ ca abhinnaḥ iti kṛtvā pradhānasya pūrvanipātaḥ na bhaviṣyati . \\
\noindent \textbf{(P\_2,2.34.1)} kim ayam tantram taranirdeśaḥ āhosvit atantram . \\
\noindent \textbf{(P\_2,2.34.1)} kim ca ataḥ . \\
\noindent \textbf{(P\_2,2.34.1)} yadi tantram dvayoḥ niyamaḥ bahuṣu aniyamaḥ . \\
\noindent \textbf{(P\_2,2.34.1)} tatra kaḥ doṣaḥ . \\
\noindent \textbf{(P\_2,2.34.1)} śaṅkhadundubhivīṅānām iti na sidhyati . \\
\noindent \textbf{(P\_2,2.34.1)} dundubhiśabdasya api pūrvnipātaḥ prāpnoti . \\
\noindent \textbf{(P\_2,2.34.1)} atha atantram mṛdaṅgaśaṅkhatūṇavāḥ pṛthak nadanti saṃsadi. prāsāde dhanapatirāmakeśavānām iti etat na sidhyati . \\
\noindent \textbf{(P\_2,2.34.1)} yathā icchasi tathā astu . \\
\noindent \textbf{(P\_2,2.34.1)} astu tāvat tantram . \\
\noindent \textbf{(P\_2,2.34.1)} nanu ca uktam dvayoḥ niyamaḥ bahuṣu aniyamaḥ iti . \\
\noindent \textbf{(P\_2,2.34.1)} tatra śaṅkhadundubhivīṅānām iti na sidhyati . \\
\noindent \textbf{(P\_2,2.34.1)} dundubhiśabdasya api pūrvnipātaḥ prāpnoti iti . \\
\noindent \textbf{(P\_2,2.34.1)} na eṣaḥ doṣaḥ . \\
\noindent \textbf{(P\_2,2.34.1)} yat etat alpāctaram iti tat alpāc iti vakṣyāmi . \\
\noindent \textbf{(P\_2,2.34.1)} atha vā punaḥ astu atantram . \\
\noindent \textbf{(P\_2,2.34.1)} nanu ca uktam mṛdaṅgaśaṅkhatūṇavāḥ pṛthak nadanti saṃsadi. prāsāde dhanapatirāmakeśavānām iti etat na sidhyati iti . \\
\noindent \textbf{(P\_2,2.34.1)} atantre taranirdeśe śaṅkhatūṇavayoḥ mṛdaṅgena samāsaḥ . \\
\noindent \textbf{(P\_2,2.34.1)} atantre taranirdeśe śaṅkhatūṇavayoḥ mṛdaṅgena samāsaḥ kariṣyate . \\
\noindent \textbf{(P\_2,2.34.1)} śaṅkhaḥ ca tūṇavaḥ ca śaṇkhatūṇavau . \\
\noindent \textbf{(P\_2,2.34.1)} mṛdaṅgaḥ ca śaṇkhatūṇavau ca mṛdaṅgaśaṅkhatūṇavāḥ . \\
\noindent \textbf{(P\_2,2.34.1)} rāmaḥ ca keśavaḥ ca rāmakeśavau dhanapatiḥ ca rāmakeśavau ca dhanapatirāmakeśavāḥ teṣām dhanapatirāmakeśavānām iti . \\
\noindent \textbf{(P\_2,2.34.1)} atha yatra bahūnām pūrvanipātaprasaṅgaḥ kim tatra ekasya niyamaḥ bhavati ahosvit aviśeṣeṇa . \\
\noindent \textbf{(P\_2,2.34.1)} anekaprāptau ekasya niyamaḥ aniyamaḥ śeṣeṣu . \\
\noindent \textbf{(P\_2,2.34.1)} anekaprāptau ekasya niyamaḥ aniyamaḥ śeṣeṣu . \\
\noindent \textbf{(P\_2,2.34.1)} paṭumṛduśuklāḥ paṭuśuklamṛdavaḥ iti . \\
\noindent \textbf{(P\_2,2.34.2)} ṛtunakṣatrāṇām ānupūrvyeṇa samānākṣarāṇām . \\
\noindent \textbf{(P\_2,2.34.2)} ṛtunakṣatrāṇām ānupūrvyeṇa samānākṣarāṇām pūrvanipātaḥ vaktavyaḥ . \\
\noindent \textbf{(P\_2,2.34.2)} śiśiravasantau udagayanasthau kṛttikārohiṇyaḥ . \\
\noindent \textbf{(P\_2,2.34.2)} abhyarhitam . \\
\noindent \textbf{(P\_2,2.34.2)} abhyarhitam pūrvam nipatati iti vaktavyam . \\
\noindent \textbf{(P\_2,2.34.2)} mātāpitarau śraddhāmedhe . \\
\noindent \textbf{(P\_2,2.34.2)} laghvakṣaram . \\
\noindent \textbf{(P\_2,2.34.2)} laghvakṣaram pūrvam nipatati iti vaktavyam . \\
\noindent \textbf{(P\_2,2.34.2)} kuśakāśam śaraśīryam . \\
\noindent \textbf{(P\_2,2.34.2)} aparaḥ āha : sarvatra eva abhyarhitam pūrvam nipatati iti vaktavyam . \\
\noindent \textbf{(P\_2,2.34.2)} laghvakṣarāt api iti . \\
\noindent \textbf{(P\_2,2.34.2)} śraddhātapasī dīkṣātapasī . \\
\noindent \textbf{(P\_2,2.34.2)} varṇānām ānupūrvyeṇa . \\
\noindent \textbf{(P\_2,2.34.2)} varṇānām ānupūrvyeṇa pūrvanipātaḥ bhavati iti vaktavyam . \\
\noindent \textbf{(P\_2,2.34.2)} brāhmaṇakṣatriyaviṭśūdrāḥ . \\
\noindent \textbf{(P\_2,2.34.2)} bhrātuḥ ca jyāyasaḥ . \\
\noindent \textbf{(P\_2,2.34.2)} bhrātuḥ ca jyāyasaḥ pūrvanipātaḥ bhavati iti vaktavyam . \\
\noindent \textbf{(P\_2,2.34.2)} yudhiṣṭhirārjunau . \\
\noindent \textbf{(P\_2,2.34.2)} saṅkhyāyāḥ alpīyasaḥ . \\
\noindent \textbf{(P\_2,2.34.2)} saṅkhyāyāḥ alpīyasaḥ pūrvanipātaḥ vaktavyaḥ . \\
\noindent \textbf{(P\_2,2.34.2)} ekādaśa dvādaśa . \\
\noindent \textbf{(P\_2,2.34.2)} dharmādiṣu ubhayam . \\
\noindent \textbf{(P\_2,2.34.2)} dharmādiṣu ubhayam pūrvam nipatati iti vaktavyam . \\
\noindent \textbf{(P\_2,2.34.2)} dharmārthau arthadharmau kāmārthau arthakāmau guṇavṛddhī vṛddhiguṇau ādyantau antādī \\
\noindent \textbf{(P\_2,2.35)} bahuvrīhau sarvanāmasaṅkhyayoḥ upasaṅkhyānam . \\
\noindent \textbf{(P\_2,2.35)} bahuvrīhau sarvanāmasaṅkhyayoḥ upasaṅkhyānam kartavyam . \\
\noindent \textbf{(P\_2,2.35)} viśvadevaḥ viśvayasāḥ dviputraḥ dvibhāryaḥ . \\
\noindent \textbf{(P\_2,2.35)} atha yatra saṅkhyāsarvanāmnoḥ eva bahurvīhiḥ kasya tatra pūrvanipātena bhavitavyam . \\
\noindent \textbf{(P\_2,2.35)} paratvāt saṅkhyāyāḥ : dvyanyāya tryanyāya . \\
\noindent \textbf{(P\_2,2.35)} vā priyasya . \\
\noindent \textbf{(P\_2,2.35)} vā priyasya pūrvanipātaḥ vaktavyaḥ . \\
\noindent \textbf{(P\_2,2.35)} priyaguḍaḥ guḍapriyaḥ . \\
\noindent \textbf{(P\_2,2.35)} saptamyāḥ pūrvanipāte gaḍvādibhyaḥ paravacanam . saptamyāḥ pūrvanipāte gaḍvādibhyaḥ parā saptamībhavati iti vaktavyam . \\
\noindent \textbf{(P\_2,2.35)} gaḍukaṇṭhaḥ gaḍuśirāḥ . \\
\noindent \textbf{(P\_2,2.36)} niṣṭhāyāḥ pūrvanipāte jātikālasukhādibhyaḥ paravacanam . \\
\noindent \textbf{(P\_2,2.36)} niṣṭhāyāḥ pūrvanipāte jātikālasukhādibhyaḥ parā niṣṭhā bhavati iti vaktavyam . \\
\noindent \textbf{(P\_2,2.36)} śārṅgajagdhī palāṇḍubhakṣitī māsajātā saṃvatsarajātā sukhajātā duḥkhajātā . \\
\noindent \textbf{(P\_2,2.36)} na vā uttarapadasya antodāttavacanam jñāpakam parabhāvasya . \\
\noindent \textbf{(P\_2,2.36)} na vā vaktavyam . \\
\noindent \textbf{(P\_2,2.36)} kim kāraṇam . \\
\noindent \textbf{(P\_2,2.36)} uttarapadasya antodāttavacanam jñāpakam parabhāvasya . \\
\noindent \textbf{(P\_2,2.36)} yat ayam jātikālasukhādibhyaḥ parasyāḥ niṣṭhāyāḥ uttarapadasya antodāttatvam śāsti tat jñāpayati ācāryaḥ parā atra niṣṭhā bhavati iti . \\
\noindent \textbf{(P\_2,2.36)} pratiṣedhe tu pūrvanipātaprasaṅgaḥ tasmāt rājadantādiṣu pāṭhaḥ . \\
\noindent \textbf{(P\_2,2.36)} pratiṣedhe tu pūrvanipātaḥ prāpnoti . \\
\noindent \textbf{(P\_2,2.36)} akṛtamitapratipannāḥ iti . \\
\noindent \textbf{(P\_2,2.36)} tasmāt rājadantādiṣu pāṭhaḥ kartavyaḥ . \\
\noindent \textbf{(P\_2,2.36)} na kartavyaḥ . \\
\noindent \textbf{(P\_2,2.36)} atra api pratiṣedhavacanam jñāpakam parā niṣṭhā bhavati iti . \\
\noindent \textbf{(P\_2,2.36)} praharaṇārthebhyaḥ ca . \\
\noindent \textbf{(P\_2,2.36)} praharaṇārthebhyaḥ ca pare niṣṭhāsaptamyau bhavataḥ iti vaktavyam . \\
\noindent \textbf{(P\_2,2.36)} asyudyataḥ musalodyataḥ asipāṇiḥ daṇḍapāṇiḥ . \\
\noindent \textbf{(P\_2,2.36)} dvandve ghi ajādyantam vipratiṣedhena . \\
\noindent \textbf{(P\_2,2.36)} dvandve ghi iti asmāt ajādyantam iti etat bhavati vipratiṣedhena . \\
\noindent \textbf{(P\_2,2.36)} dvandve ghi iti asya avakāśaḥ paṭuguptau . \\
\noindent \textbf{(P\_2,2.36)} ajādyadantam iti asya avakāśaḥ uṣṭrakharau . \\
\noindent \textbf{(P\_2,2.36)} iha ubhayam prāpnoti indrāgnī . \\
\noindent \textbf{(P\_2,2.36)} ajādyadantam iti etat bhavati vipratiṣedhena . \\
\noindent \textbf{(P\_2,2.36)} ubhābhyām alpāctaram . \\
\noindent \textbf{(P\_2,2.36)} ubhābhyām alpāctaram iti etat bhavati . \\
\noindent \textbf{(P\_2,2.36)} dvandve ghi iti asya avakāśaḥ paṭuguptau . \\
\noindent \textbf{(P\_2,2.36)} alpāctaram iti asya avakāśaḥ vāgdṛṣadau . \\
\noindent \textbf{(P\_2,2.36)} iha ubhayam prāpnoti vāgagnī . \\
\noindent \textbf{(P\_2,2.36)} alpāctaram iti etat bhavati vipratiṣedhena . \\
\noindent \textbf{(P\_2,2.36)} ajādyadantam iti asya avakāśaḥ uṣṭrakharau . \\
\noindent \textbf{(P\_2,2.36)} alpāctaram iti asya avakāśaḥ saḥ eva . \\
\noindent \textbf{(P\_2,2.36)} iha ubhayam prāpnoti vāgindrau . \\
\noindent \textbf{(P\_2,2.36)} alpāctaram iti etat bhavati vipratiṣedhena . \\
\noindent \textbf{(P\_2,2.38)} kaḍārādayaḥ iti vaktavyam iha api yathā syāt . \\
\noindent \textbf{(P\_2,2.38)} gaḍulaśāṇḍilyaḥ śāṇḍilyagaḍulaḥ khaṇḍavātsyaḥ vatsyakaṇḍaḥ . \\
\noindent \textbf{(P\_2,2.38)} tat tarhi vaktavyam . \\
\noindent \textbf{(P\_2,2.38)} na vaktavyam . \\
\noindent \textbf{(P\_2,2.38)} bahuvacananirdeśāt kaḍārādayaḥ iti vijñāsyate . \\
\noindent \textbf{(P\_2,3.1.1)} anabhihite iti ucyate . \\
\noindent \textbf{(P\_2,3.1.1)} kim idam anabhihitam nāma . \\
\noindent \textbf{(P\_2,3.1.1)} uktam nirdiṣṭam abhihitam iti anarthāntaram . \\
\noindent \textbf{(P\_2,3.1.1)} yāvat brūyāt anukte anirdiṣṭe iti tāvat anabhihite iti . \\
\noindent \textbf{(P\_2,3.1.1)} anabhihitavacanam anarthakam anyatra api vihitasya abhāvāt abhihite . \\
\noindent \textbf{(P\_2,3.1.1)} anabhihitavacanam anarthakam . \\
\noindent \textbf{(P\_2,3.1.1)} kim kāraṇam . \\
\noindent \textbf{(P\_2,3.1.1)} anyatra api vihitasya abhāvāt abhihite . \\
\noindent \textbf{(P\_2,3.1.1)} anyatra api abhihite vihitam na bhavati . \\
\noindent \textbf{(P\_2,3.1.1)} kva anyatra . \\
\noindent \textbf{(P\_2,3.1.1)} citraguḥ śabalaguḥ . \\
\noindent \textbf{(P\_2,3.1.1)} bahuvrīhiṇā uktatvāt matvarthasya matvarthīyaḥ na bhavati . \\
\noindent \textbf{(P\_2,3.1.1)} gargāḥ vatsāḥ vidāḥ urvāḥ . \\
\noindent \textbf{(P\_2,3.1.1)} yañañbhyām uktatvāt apatyārthasya nyāyyotpattiḥ na bhavati . \\
\noindent \textbf{(P\_2,3.1.1)} saptaparṇaḥ aṣṭāpadamiti . \\
\noindent \textbf{(P\_2,3.1.1)} samāsena uktatvāt vīpsāyāḥ dvirvacanam na bhavati . \\
\noindent \textbf{(P\_2,3.1.1)} yat tāvat ucyate citraguḥ śabalguḥ bahuvrīhiṇā uktatvāt matvarthasya matvarthīyaḥ na bhavati iti . \\
\noindent \textbf{(P\_2,3.1.1)} astinā sāmānādhikaraṇye matup vidhīyate . \\
\noindent \textbf{(P\_2,3.1.1)} na ca atra astinā sāmānādhikaraṇyam . \\
\noindent \textbf{(P\_2,3.1.1)} yat api ucyate gargāḥ vatsāḥ vidāḥ urvāḥ yañañbhyām uktatvāt apatyārthasya nyāyyotpattiḥ na bhavati iti . \\
\noindent \textbf{(P\_2,3.1.1)} samarthānām prathamāt vā iti vartate . \\
\noindent \textbf{(P\_2,3.1.1)} na ca etat samarthānām prathamam . \\
\noindent \textbf{(P\_2,3.1.1)} kiṃ tarhi . \\
\noindent \textbf{(P\_2,3.1.1)} dvitīyam arthamupasaṃkrāntam . \\
\noindent \textbf{(P\_2,3.1.1)} yat api ucyate saptaparṇaḥ aṣṭāpadam iti samāsena uktatvāt vīpsāyāḥ dvirvacanam na bhavati iti . \\
\noindent \textbf{(P\_2,3.1.1)} yat atra vīpsāyuktam na adaḥ prayujyate . \\
\noindent \textbf{(P\_2,3.1.1)} kim punaḥ tat . \\
\noindent \textbf{(P\_2,3.1.1)} parvaṇi parvaṇi sapta parṇāni asya . \\
\noindent \textbf{(P\_2,3.1.1)} paṅktau paṅktau aṣṭau padāni iti . \\
\noindent \textbf{(P\_2,3.1.1)} śnambahujakakṣu tarhi . \\
\noindent \textbf{(P\_2,3.1.1)} śnam : bhinatti chinatti . \\
\noindent \textbf{(P\_2,3.1.1)} śnamā uktatvāt kartṛtvasya kartari śap na bhavati . \\
\noindent \textbf{(P\_2,3.1.1)} bahuc : bahukṛtam , bahubhinnam iti . \\
\noindent \textbf{(P\_2,3.1.1)} bahucā uktatvāt īṣadasmāpteḥ kalpabādayaḥ na bhavanti iti . \\
\noindent \textbf{(P\_2,3.1.1)} akac : uccakaiḥ , nīcakaiḥ iti . \\
\noindent \textbf{(P\_2,3.1.1)} akacā uktatvāt kutsādīnām kādayaḥ na bhavanti . \\
\noindent \textbf{(P\_2,3.1.1)} nanu ca śnambahujakacaḥ apavādāḥ te apavādatvāt bādhakāḥ bhaviṣyanti . \\
\noindent \textbf{(P\_2,3.1.1)} śnambahujakakṣu nānādeśatvāt utsargāpratiṣedhaḥ . \\
\noindent \textbf{(P\_2,3.1.1)} samānadeśaiḥ apavādaiḥ utsargāṇām bādhanam bhavati . \\
\noindent \textbf{(P\_2,3.1.1)} nānādeśatvāt na prāpnoti . \\
\noindent \textbf{(P\_2,3.1.1)} kim punaḥ iha akartavyaḥ anabhihitādhikāraḥ kriyate āhosvit anyatra kartavyaḥ na kriyate . \\
\noindent \textbf{(P\_2,3.1.1)} iha akartavyaḥ kriyate . \\
\noindent \textbf{(P\_2,3.1.1)} eṣaḥ eva hi nyāyyaḥ pakṣaḥ yat abhihite vihitam na syāt . \\
\noindent \textbf{(P\_2,3.1.1)} anabhitaḥ tu vibhaktyarthaḥ tasmāt anabhihitavacanam . anabhihitaḥ tu vibhaktyarthaḥ . \\
\noindent \textbf{(P\_2,3.1.1)} kaḥ punaḥ vibhaktyarthaḥ . \\
\noindent \textbf{(P\_2,3.1.1)} ekatvādayaḥ vibhaktyarthāḥ teṣu anabhihiteṣu karmādayaḥ bhihitāḥ vibhaktīnām utpattau nimittatvāya mā bhūvan iti . \\
\noindent \textbf{(P\_2,3.1.1)} tasmāt anabhihitavacanam . \\
\noindent \textbf{(P\_2,3.1.1)} tasmāt anabhihitādhikāraḥ kriyate . \\
\noindent \textbf{(P\_2,3.1.1)} avaśyam ca etat evam vijñeyam ekatvādayaḥ vibhaktyarthāḥ iti . \\
\noindent \textbf{(P\_2,3.1.1)} abhihite prathamābhāvaḥ . yaḥ hi manyate karmādayaḥ vibhaktyarthāḥ teṣu abhihiteṣu sāmarthyāt me vibhaktīnām utpattiḥ na bhaviṣyati iti prathamā tasya na prāpnoti . \\
\noindent \textbf{(P\_2,3.1.1)} kva . \\
\noindent \textbf{(P\_2,3.1.1)} vṛkṣaḥ plakṣaḥ .kiṃ kāraṇam . \\
\noindent \textbf{(P\_2,3.1.1)} prātipadikena uktaḥ prātipadikārthaḥ iti . \\
\noindent \textbf{(P\_2,3.1.1)} na kva cit prātipadikena anuktaḥ prātipadikārthaḥ ucyate ca prathamā . \\
\noindent \textbf{(P\_2,3.1.1)} sā vacanāt bhaviṣyati . \\
\noindent \textbf{(P\_2,3.1.1)} tava eva tu khalu eṣaḥ doṣaḥ yasya te ekatvādayaḥ vibhaktyarthāḥ abhihite prathamābhāvaḥ iti . \\
\noindent \textbf{(P\_2,3.1.1)} prathamā te na prāpnoti . \\
\noindent \textbf{(P\_2,3.1.1)} kva . \\
\noindent \textbf{(P\_2,3.1.1)} pacati odanam devadattaḥ iti . \\
\noindent \textbf{(P\_2,3.1.1)} kim kāraṇam . \\
\noindent \textbf{(P\_2,3.1.1)} tiṅā uktāḥ ekatvādayaḥ iti . \\
\noindent \textbf{(P\_2,3.1.1)} anabhihitādhikāram ca tvam karoṣi parigaṇanam ca . \\
\noindent \textbf{(P\_2,3.1.1)} na kva cit tiṅā ekatvādīnām anabhidhānam ucyate ca prathamā . \\
\noindent \textbf{(P\_2,3.1.1)} sā vacanāt bhaviṣyati . \\
\noindent \textbf{(P\_2,3.1.1)} nanu ca iha anabhidhānam vṛkṣaḥ plakṣaḥ iti . \\
\noindent \textbf{(P\_2,3.1.1)} atra api abhidhānam asti . \\
\noindent \textbf{(P\_2,3.1.1)} katham . \\
\noindent \textbf{(P\_2,3.1.1)} vakṣyati etat : astiḥ bhavantīparaḥ prathamapuruṣaḥ aprayujyamānaḥ api asti iti . \\
\noindent \textbf{(P\_2,3.1.1)} vṛkṣaḥ plakṣaḥ . \\
\noindent \textbf{(P\_2,3.1.1)} asti iti gamyate . \\
\noindent \textbf{(P\_2,3.1.1)} tava eva tu khalu eṣaḥ doṣaḥ yasya te karmādayaḥ vibhaktyārthāḥ abhihite prathamābhāvaḥ iti . \\
\noindent \textbf{(P\_2,3.1.1)} prathamā te prāpnoti . \\
\noindent \textbf{(P\_2,3.1.1)} kva . \\
\noindent \textbf{(P\_2,3.1.1)} kaṭam karoti bhīṣmam udāram śobhanam darśanīyam iti . \\
\noindent \textbf{(P\_2,3.1.1)} kaṭaśabdāt utpadyamānayā dvitīyayā abhihitam karma iti kṛtvā bhīṣādibhyaḥ dvitīyā na prāpnoti . \\
\noindent \textbf{(P\_2,3.1.1)} kā tarhi prāpnoti . \\
\noindent \textbf{(P\_2,3.1.1)} prathamā . \\
\noindent \textbf{(P\_2,3.1.1)} tat yathā . \\
\noindent \textbf{(P\_2,3.1.1)} kṛtaḥ kaṭaḥ bhīṣmaḥ udāraḥ śobhanaḥ darśanīyaḥ iti . \\
\noindent \textbf{(P\_2,3.1.1)} karoteḥ utpadyamānena ktena abhihitam karma iti kṛtvā bhīṣmādibhyaḥ dvitīyā na bhavati . \\
\noindent \textbf{(P\_2,3.1.1)} kā tarhi . \\
\noindent \textbf{(P\_2,3.1.1)} prathamā bhavati . \\
\noindent \textbf{(P\_2,3.1.1)} na eṣaḥ doṣaḥ . \\
\noindent \textbf{(P\_2,3.1.1)} na hi mama anabhihitādhikāraḥ asti na api parigaṇanam . \\
\noindent \textbf{(P\_2,3.1.1)} sāmarthyāt me vibhaktīnām utpattiḥ bhaviṣyati . \\
\noindent \textbf{(P\_2,3.1.1)} asti ca sāmarthyam . \\
\noindent \textbf{(P\_2,3.1.1)} kim . \\
\noindent \textbf{(P\_2,3.1.1)} karmaviśeṣaḥ vaktavyaḥ . \\
\noindent \textbf{(P\_2,3.1.1)} atha vā kaṭaḥ api karma bhīṣmādayaḥ api . \\
\noindent \textbf{(P\_2,3.1.1)} tatra karmaṇi iti eva siddham . \\
\noindent \textbf{(P\_2,3.1.1)} atha vā kaṭaḥ eva karma tat sāmānādhikaraṇyāt bhīṣmādibhyaḥ dvitīyā bhaviṣyati . \\
\noindent \textbf{(P\_2,3.1.1)} asti khalvapi viśeṣaḥ kaṭaṃ karoti bhīṣmamudāram śobhanam darśanīyam iti ca kṛtaḥ kaṭo bhīṣmaḥ udāraḥ śobhanaḥ darśanīyaḥ iti ca . \\
\noindent \textbf{(P\_2,3.1.1)} karoteḥ utpadyamānaḥ ktaḥ anavayavena sarvam karma abhidhatte . \\
\noindent \textbf{(P\_2,3.1.1)} kaṭaśabdāt punaḥ utpadyamānayā dvitīyayā yat kaṭastham karma tat śakyamabhidhātum na hi karmaviśeṣaḥ . \\
\noindent \textbf{(P\_2,3.1.1)} tava eva tu khalu eṣaḥ doṣaḥ yasya te ekatvādayaḥ vibhaktyarthāḥ abhihite prathamābhāvaḥ iti . \\
\noindent \textbf{(P\_2,3.1.1)} prathamā te na prāpnoti . \\
\noindent \textbf{(P\_2,3.1.1)} kva . \\
\noindent \textbf{(P\_2,3.1.1)} ekaḥ dvau bahavaḥ iti . \\
\noindent \textbf{(P\_2,3.1.1)} kim kāraṇam . \\
\noindent \textbf{(P\_2,3.1.1)} prātipadikena uktāḥ ekatvādayaḥ iti . \\
\noindent \textbf{(P\_2,3.1.1)} karmādiṣu api vai vibhaktyartheṣu avaśyam ekatvādayaḥ nimittatvena upādeyāḥ . \\
\noindent \textbf{(P\_2,3.1.1)} karmaṇaḥ evatve karmaṇaḥ dvitve karmaṇaḥ bahutve iti . \\
\noindent \textbf{(P\_2,3.1.1)} na ca ekatvādīnām ekatvādayaḥ santi . \\
\noindent \textbf{(P\_2,3.1.1)} atha santi mama api santi . \\
\noindent \textbf{(P\_2,3.1.1)} teṣu anabhihiteṣu prathamā bhaviṣyati . \\
\noindent \textbf{(P\_2,3.1.1)} atha vā ubhayavacanāḥ hyete . \\
\noindent \textbf{(P\_2,3.1.1)} dravyam ca āhuḥ guṇam ca . \\
\noindent \textbf{(P\_2,3.1.1)} yatsthaḥ asau guṇaḥ tasya anuktāḥ ekatvādayaḥ iti kṛtvā prathamā bhaviṣyati . \\
\noindent \textbf{(P\_2,3.1.1)} atha vā saṅkhyā nāma iyam parapradhānā . \\
\noindent \textbf{(P\_2,3.1.1)} saṃkhyeyam anayā viśeṣyam . \\
\noindent \textbf{(P\_2,3.1.1)} yadi ca atra prathamā na syāt saṅkhyeyam aviśeṣitam syāt . \\
\noindent \textbf{(P\_2,3.1.1)} atha vā vakṣyati tatra vacanagrahaṇasya prayojanam ukteṣu api ekatvādiṣu prathamā yathā syāt iti . \\
\noindent \textbf{(P\_2,3.1.1)} atha vā samayāt bhaviṣyati . \\
\noindent \textbf{(P\_2,3.1.1)} yadi sāmayakī na niyogataḥ anyāḥ kasmāt na bhavanti . \\
\noindent \textbf{(P\_2,3.1.1)} karmādīnām abhāvāt . \\
\noindent \textbf{(P\_2,3.1.1)} ṣaṣṭhī tarhi prāpnoti . \\
\noindent \textbf{(P\_2,3.1.1)} śeṣalakṣaṇā ṣaṣṭhī aśeṣatvāt na bhaviṣyati . \\
\noindent \textbf{(P\_2,3.1.1)} evam api vyatikaraḥ prāpnoti . \\
\noindent \textbf{(P\_2,3.1.1)} ekasmin api dvivacanabahuvacane prāpnutaḥ . \\
\noindent \textbf{(P\_2,3.1.1)} dvayoḥ api ekavavacanabahuvacane prāpnutaḥ . \\
\noindent \textbf{(P\_2,3.1.1)} bahuṣu api ekavacanadvivacane prāpnutaḥ . \\
\noindent \textbf{(P\_2,3.1.1)} arthataḥ vyavasthā bhaviṣyati . \\
\noindent \textbf{(P\_2,3.1.2)} parigaṇanaṃ kartavyam . \\
\noindent \textbf{(P\_2,3.1.2)} tiṅkṛttaddhitasamāsaiḥ parisaṅkhyānam . \\
\noindent \textbf{(P\_2,3.1.2)} tiṅkṛttaddhitasamāsaiḥ parisaṅkhyānam kartavyam . \\
\noindent \textbf{(P\_2,3.1.2)} tiṅ . \\
\noindent \textbf{(P\_2,3.1.2)} kriyate kaṭaḥ . \\
\noindent \textbf{(P\_2,3.1.2)} kṛt . \\
\noindent \textbf{(P\_2,3.1.2)} kṛtaḥ kaṭaḥ . \\
\noindent \textbf{(P\_2,3.1.2)} taddhita . \\
\noindent \textbf{(P\_2,3.1.2)} aupagavaḥ kāpaṭavaḥ . \\
\noindent \textbf{(P\_2,3.1.2)} samāsa . \\
\noindent \textbf{(P\_2,3.1.2)} citraguḥ śabalaguḥ . \\
\noindent \textbf{(P\_2,3.1.2)} utsarge hi prātipadikasāmānādhikaraṇye vibhaktivacanam . utsarge hi prātipadikasāmānādhikaraṇye vibhaktiḥ vaktavyā . \\
\noindent \textbf{(P\_2,3.1.2)} kva . \\
\noindent \textbf{(P\_2,3.1.2)} kaṭam karoti bhīṣmam udāram śobhanam darśanīyam iti . \\
\noindent \textbf{(P\_2,3.1.2)} kaṭaśabdāt utpadyamānayā dvitīyayā abhihitam karma iti kṛtvā bhīṣmādibhyaḥ dvitīyā na prāpnoti . \\
\noindent \textbf{(P\_2,3.1.2)} kā tarhi syāt . \\
\noindent \textbf{(P\_2,3.1.2)} ṣaṣṭhī . \\
\noindent \textbf{(P\_2,3.1.2)} śeṣalakṣaṇā ṣaṣṭhī . \\
\noindent \textbf{(P\_2,3.1.2)} aśeṣatvāt na bhaviṣyati . \\
\noindent \textbf{(P\_2,3.1.2)} anyāḥ api na prāpnuvanti . \\
\noindent \textbf{(P\_2,3.1.2)} kim kāraṇam . \\
\noindent \textbf{(P\_2,3.1.2)} karmādīn āmabhāvāt . \\
\noindent \textbf{(P\_2,3.1.2)} samayaśca kṛtaḥ ne kevalā prakṛtiḥ proktavyā na kevalaḥ pratyayaḥ iti . \\
\noindent \textbf{(P\_2,3.1.2)} na cānyā utpadyamānā etam abhisambandham utsahante vaktum iti kṛtvā dvitīyā bhaviṣyati . \\
\noindent \textbf{(P\_2,3.1.2)} atha vā kaṭaḥ api karma bhīṣmādayaḥ api . \\
\noindent \textbf{(P\_2,3.1.2)} tatra karmaṇi iti eva siddham . \\
\noindent \textbf{(P\_2,3.1.2)} atha vā kaṭaḥ eva karma . \\
\noindent \textbf{(P\_2,3.1.2)} tatsāmānādhikaraṇyāt bhīṣmādibhyaḥ dvitīyā bhaviṣyati . \\
\noindent \textbf{(P\_2,3.1.2)} tasmāt na arthaḥ parigaṇanena . \\
\noindent \textbf{(P\_2,3.1.3)} dvayoḥ kriyayoḥ kārake anyatareṇa abhihite vibhaktyabhāvaprasaṅgaḥ . dvayoḥ kriyayoḥ kārake anyatareṇa abhihite vibhaktiḥ na prāpnoti . \\
\noindent \textbf{(P\_2,3.1.3)} kva . \\
\noindent \textbf{(P\_2,3.1.3)} prāsāde āste , śayane āste iti . \\
\noindent \textbf{(P\_2,3.1.3)} kim kāraṇam . \\
\noindent \textbf{(P\_2,3.1.3)} sadipratyayena abhihitam adhikaraṇam iti kṛtvā saptamī na prāpnoti . \\
\noindent \textbf{(P\_2,3.1.3)} na vā anyatareṇa anabhidhānāt . \\
\noindent \textbf{(P\_2,3.1.3)} na vā eṣaḥ doṣaḥ . \\
\noindent \textbf{(P\_2,3.1.3)} kim kāraṇam . \\
\noindent \textbf{(P\_2,3.1.3)} anyatareṇa anabhidhānāt . \\
\noindent \textbf{(P\_2,3.1.3)} anyatareṇa atra anabhidhānam . \\
\noindent \textbf{(P\_2,3.1.3)} sadipratyayena bhidhānam āsipratyayena anabhidhānam . \\
\noindent \textbf{(P\_2,3.1.3)} yataḥ anabhidhānam tadāśrayā saptamī bhaviṣyati . \\
\noindent \textbf{(P\_2,3.1.3)} kutaḥ na khalu etat sati abhidhāne ca anabhidhāne ca anabhihitāśrayā saptamī bhaviṣyati na punaḥ abhihitāśrayaḥ pratiṣedhaḥ iti . \\
\noindent \textbf{(P\_2,3.1.3)} anabhihite hi vidhānam . \\
\noindent \textbf{(P\_2,3.1.3)} anabhihite hi saptamī vidhīyate na abhihite pratiṣedhaḥ . \\
\noindent \textbf{(P\_2,3.1.3)} yadi api tāvat atra etat śakyate vaktum yatra anyā ca anyā ca kriyā yatra tu khalu sā eva kriyā tatra katham . \\
\noindent \textbf{(P\_2,3.1.3)} āsane āste . \\
\noindent \textbf{(P\_2,3.1.3)} śayane śete iti . \\
\noindent \textbf{(P\_2,3.1.3)} atra api anyatvam asti . \\
\noindent \textbf{(P\_2,3.1.3)} kutaḥ . \\
\noindent \textbf{(P\_2,3.1.3)} kālabhedāt sādhanabhedāt ca . \\
\noindent \textbf{(P\_2,3.1.3)} ekasya atra āseḥ āsiḥ sādhanam sarvakālaḥ ca pratyayaḥ . \\
\noindent \textbf{(P\_2,3.1.3)} aparasya bāhyam sādhanam vartamānakālaḥ ca pratyayaḥ . \\
\noindent \textbf{(P\_2,3.1.3)} kim punaḥ dravyam sādhanam āhosvit guṇaḥ . \\
\noindent \textbf{(P\_2,3.1.3)} kim ca ataḥ . \\
\noindent \textbf{(P\_2,3.1.3)} yadi dravyam sādhanam na etat anyat bhavati abhihitāt . \\
\noindent \textbf{(P\_2,3.1.3)} atha hi guṇaḥ sādhanam bhavati etat anyat abhihitāt . \\
\noindent \textbf{(P\_2,3.1.3)} anyaḥ hi sadiguṇaḥ anyaḥ ca āsiguṇaḥ . \\
\noindent \textbf{(P\_2,3.1.3)} kiṃ punaḥ sādhanam nyāyyam . \\
\noindent \textbf{(P\_2,3.1.3)} guṇaḥ iti āha . \\
\noindent \textbf{(P\_2,3.1.3)} katham jñāyate . \\
\noindent \textbf{(P\_2,3.1.3)} evam hi kaḥ cit kam cit pṛcchati . \\
\noindent \textbf{(P\_2,3.1.3)} kva devadattaḥ iti . \\
\noindent \textbf{(P\_2,3.1.3)} saḥ tasmai ācaṣṭe . \\
\noindent \textbf{(P\_2,3.1.3)} asau vṛkṣe iti . \\
\noindent \textbf{(P\_2,3.1.3)} katarasmin . \\
\noindent \textbf{(P\_2,3.1.3)} yaḥ tiṣṭhati iti . \\
\noindent \textbf{(P\_2,3.1.3)} saḥ vṛkṣaḥ adhikaraṇam bhūtvā anyena śabdena abhisambadhyamānaḥ kartā sampadyate . \\
\noindent \textbf{(P\_2,3.1.3)} dravye punaḥ sādhane sati yat karma karma eva syāt yat karaṇam karaṇam eva yat adhikaraṇam adhikaraṇam eva . \\
\noindent \textbf{(P\_2,3.1.4)} anabhihitavacanam anarthakam prathamāvidhānasya anavakāśatvāt . \\
\noindent \textbf{(P\_2,3.1.4)} anabhihitavacanam anarthakam . \\
\noindent \textbf{(P\_2,3.1.4)} kim kāraṇam . \\
\noindent \textbf{(P\_2,3.1.4)} prathamāvidhānasya anavakāśatvāt . \\
\noindent \textbf{(P\_2,3.1.4)} anavakāśā prathamā . \\
\noindent \textbf{(P\_2,3.1.4)} sā vacanāt bhaviṣyati. sāvakāśā prathamā . \\
\noindent \textbf{(P\_2,3.1.4)} kaḥ avakāśaḥ . \\
\noindent \textbf{(P\_2,3.1.4)} akārakam . \\
\noindent \textbf{(P\_2,3.1.4)} vṛkṣaḥ plakṣaḥ iti . \\
\noindent \textbf{(P\_2,3.1.4)} avakāśaḥ akārakam iti cet na astiḥ bhavantīparaḥ prathamapuruṣaḥ aprayujyamānaḥ api asti . \\
\noindent \textbf{(P\_2,3.1.4)} avakāśaḥ akārakam iti cet tat na . \\
\noindent \textbf{(P\_2,3.1.4)} kim kāraṇam . \\
\noindent \textbf{(P\_2,3.1.4)} astiḥ bhavantīparaḥ prathamapuruṣaḥ aprayujyamānaḥ api asti iti gamyate . \\
\noindent \textbf{(P\_2,3.1.4)} vṛkṣaḥ plakṣaḥ . \\
\noindent \textbf{(P\_2,3.1.4)} asti iti gamyate. vipratiṣedhāt vā prathamābhāvaḥ . atha vā dvitīyādayaḥ kriyantām prathamā vā iti prathamā bhaviṣyati vipratiṣedhena . \\
\noindent \textbf{(P\_2,3.1.4)} dvitīyādīnām avakāśaḥ kaṭam karoti bhīṣmam udāram śobhanamdarśanīyam iti . \\
\noindent \textbf{(P\_2,3.1.4)} prathamāyāḥ avakāśaḥ akārakam vṛkṣaḥ plakṣaḥ iti . \\
\noindent \textbf{(P\_2,3.1.4)} iha ubhayam prāpnoti . \\
\noindent \textbf{(P\_2,3.1.4)} kṛtaḥ kaṭaḥ bhīṣmaḥ udāraḥ śobhanaḥ darśanīyaḥ iti . \\
\noindent \textbf{(P\_2,3.1.4)} prathamā bhaviṣyati vipratiṣedhena . \\
\noindent \textbf{(P\_2,3.1.4)} na sidhyati . \\
\noindent \textbf{(P\_2,3.1.4)} paratvāt ṣaṣṭhī prāpnoti . \\
\noindent \textbf{(P\_2,3.1.4)} śeṣalakṣaṇā ṣaṣṭhī aśeṣatvāt na bhaviṣyati . \\
\noindent \textbf{(P\_2,3.1.4)} kṛtprayoge tu param vidhānam ṣaṣṭhyāḥ tatpratiṣedhārtham . \\
\noindent \textbf{(P\_2,3.1.4)} kṛtprayoge tu paratvā tṣaṣṭhī prāpnoti . \\
\noindent \textbf{(P\_2,3.1.4)} tatpratiṣedhārtham anabhihitādhikāraḥ kartavyaḥ . \\
\noindent \textbf{(P\_2,3.1.4)} kartavyaḥ kaṭaḥ iti . \\
\noindent \textbf{(P\_2,3.1.4)} saḥ katham kartavyaḥ . \\
\noindent \textbf{(P\_2,3.1.4)} yadi ekatvādayaḥ vibhaktyarthāḥ . \\
\noindent \textbf{(P\_2,3.1.4)} atha hi karmādayaḥ vibhaktyarthāḥ na arthaḥ anabhihitādhikāreṇa . \\
\noindent \textbf{(P\_2,3.2)} samayānikaṣāhāyogeṣu upasaṅkhyānam . samayānikaṣāhāyogeṣu upasaṅkhyānam kartavyam . \\
\noindent \textbf{(P\_2,3.2)} samayā grāmam nikaṣā grāmam . \\
\noindent \textbf{(P\_2,3.2)} hāyoge . \\
\noindent \textbf{(P\_2,3.2)} hā devadattam . \\
\noindent \textbf{(P\_2,3.2)} hā yajñadattam . \\
\noindent \textbf{(P\_2,3.2)} aparaḥ āha : dvitīyāvidhāne abhitaḥparitaḥsamayānikaṣādhyadhidhigyogeṣu upasaṅkhyānam . \\
\noindent \textbf{(P\_2,3.2)} dvitīyāvidhāne abhitaḥparitaḥsamayānikaṣādhyadhidhigyogeṣu upasaṅkhyānam kartavyam . \\
\noindent \textbf{(P\_2,3.2)} abhitaḥ grāmam paritaḥ grāmam . \\
\noindent \textbf{(P\_2,3.2)} samayā grāmam . \\
\noindent \textbf{(P\_2,3.2)} nikaṣā grāmam . \\
\noindent \textbf{(P\_2,3.2)} adhi adhi grāmam . \\
\noindent \textbf{(P\_2,3.2)} dhik jālmam dhik vṛṣalam . \\
\noindent \textbf{(P\_2,3.2)} aparaḥ āha . \\
\noindent \textbf{(P\_2,3.2)} ubhasarvatasoḥ kāryā dhiguparyādiṣu triṣu dvitīyā āmreḍitānteṣu tataḥ anyatra api dṛśyate . \\
\noindent \textbf{(P\_2,3.2)} ubhaya sarva iti etābhyām tasantābhyām dvitīyā vaktavyā . \\
\noindent \textbf{(P\_2,3.2)} ubhayataḥ grāmam sarvataḥ grāmam . \\
\noindent \textbf{(P\_2,3.2)} dhigyoge . \\
\noindent \textbf{(P\_2,3.2)} dhik jālmam dhik vṛṣalam . \\
\noindent \textbf{(P\_2,3.2)} uparyādiṣu triṣu āmreḍitānteṣu dvitīyā vaktavyā . \\
\noindent \textbf{(P\_2,3.2)} upari upari grāmam . \\
\noindent \textbf{(P\_2,3.2)} adhi adhi grāmam . \\
\noindent \textbf{(P\_2,3.2)} adhaḥ adhaḥ grāmam . \\
\noindent \textbf{(P\_2,3.2)} tataḥ anyatra api dṛśyate . \\
\noindent \textbf{(P\_2,3.2)} na devadattam pratibhāti kim cit . \\
\noindent \textbf{(P\_2,3.2)} bubhukṣitam na pratibhāti kim cit . \\
\noindent \textbf{(P\_2,3.3)} kimartham idam ucyate . \\
\noindent \textbf{(P\_2,3.3)} tṛtīyā yathā syāt . \\
\noindent \textbf{(P\_2,3.3)} atha dvitīyā siddhā . \\
\noindent \textbf{(P\_2,3.3)} siddhā karmaṇi iti eva . \\
\noindent \textbf{(P\_2,3.3)} tṛtīyā api siddhā . \\
\noindent \textbf{(P\_2,3.3)} katham . \\
\noindent \textbf{(P\_2,3.3)} supām supaḥ bhavanti iti eva . \\
\noindent \textbf{(P\_2,3.3)} asati etasmin supām supaḥ bhavanti iti tṛtīyārthaḥ ayam ārambhaḥ . \\
\noindent \textbf{(P\_2,3.3)} yavāgvā agnihotraṃ juhoti . \\
\noindent \textbf{(P\_2,3.3)} evam tarhi tṛtīyā api siddhā . \\
\noindent \textbf{(P\_2,3.3)} katham . \\
\noindent \textbf{(P\_2,3.3)} kartṛkaraṇayoḥ iti eva . \\
\noindent \textbf{(P\_2,3.3)} ayam agnihotraśabdaḥ asti eva jyotiṣi vartate . \\
\noindent \textbf{(P\_2,3.3)} tadyathā . \\
\noindent \textbf{(P\_2,3.3)} agnihotram prajvalayati iti . \\
\noindent \textbf{(P\_2,3.3)} asti haviṣi vartate . \\
\noindent \textbf{(P\_2,3.3)} tat yathā . \\
\noindent \textbf{(P\_2,3.3)} agnihotram juhoti iti . \\
\noindent \textbf{(P\_2,3.3)} juhotiḥ ca asti eva prakṣepaṇe vartate asti prīṇātyarthe vartate . \\
\noindent \textbf{(P\_2,3.3)} tat yathā tāvat yavāgūśabdāt tṛtīyā tadā agnihotraśabdaḥ jyotiṣi vartate juhotiḥ ca prīṇātyarthe . \\
\noindent \textbf{(P\_2,3.3)} tat yathā . \\
\noindent \textbf{(P\_2,3.3)} yavāgvā agnihotram juhoti . \\
\noindent \textbf{(P\_2,3.3)} agniṃ prīṇāti . \\
\noindent \textbf{(P\_2,3.3)} yadā yavāgūśabdāt dvitīyā tadā agnihotraśabdaḥ haviṣi vartate juhotiḥ ca prakṣepaṇe . \\
\noindent \textbf{(P\_2,3.3)} tat yathā . \\
\noindent \textbf{(P\_2,3.3)} yavāgūm agnihotram juhoti . \\
\noindent \textbf{(P\_2,3.3)} yavāgūm haviḥ agnau prakṣipati . \\
\noindent \textbf{(P\_2,3.4)} iha kasmāt na bhavati . \\
\noindent \textbf{(P\_2,3.4)} kim te bābhravaśālaṅkāyanānām antareṇa gatena iti . \\
\noindent \textbf{(P\_2,3.4)} lakṣaṇapratipadoktayoḥ pratipadoktasya eva iti . \\
\noindent \textbf{(P\_2,3.4)} atha vā yadi api tāvat ayam antareṇaśabdaḥ dṛṣṭāpacāraḥ nipātaḥ ca anipātaḥ ca ayaṃ tu khalu antarāśabdaḥ adṛṣṭāpacāraḥ nipātaḥ eva . \\
\noindent \textbf{(P\_2,3.4)} tasya asya kaḥ anyaḥ dvitīyaḥ sahāyaḥ bhavitum arhati anyat ataḥ nipātāt . \\
\noindent \textbf{(P\_2,3.4)} tat yathā . \\
\noindent \textbf{(P\_2,3.4)} asya goḥ dvitīyena arthaḥ iti gauḥ eva ānīyate na aśvaḥ na gadarbhaḥ . \\
\noindent \textbf{(P\_2,3.4)} antarāntareṇayuktānām apradhānavacanam . \\
\noindent \textbf{(P\_2,3.4)} antarāntareṇayuktānāmapradhānagrahaṇam vaktavyam . \\
\noindent \textbf{(P\_2,3.4)} apradhāne dvīiyā bhavati iti vaktavyam . \\
\noindent \textbf{(P\_2,3.4)} antarā tvām ca mām ca kamaṇḍaluḥ iti . \\
\noindent \textbf{(P\_2,3.4)} kamaṇḍaloḥ dvitīyā mā bhūt iti . \\
\noindent \textbf{(P\_2,3.4)} kaḥ punaḥ etābhyām kamaṇḍaloḥ yogaḥ . \\
\noindent \textbf{(P\_2,3.4)} yat tat tvām ca mām ca antarā tat kamaṇḍaloḥ sthānam . \\
\noindent \textbf{(P\_2,3.4)} tatt arhi vaktavyam . \\
\noindent \textbf{(P\_2,3.4)} na vaktavyam . \\
\noindent \textbf{(P\_2,3.4)} kamaṇḍaloḥ dvitīyā kasmāt na bhavati . \\
\noindent \textbf{(P\_2,3.4)} upapadavibhakteḥ kārakavibhaktiḥ balīyasī iti prathamā bhaviṣyati . \\
\noindent \textbf{(P\_2,3.5)} atyantasaṃyoge karmavat lādyartham . \\
\noindent \textbf{(P\_2,3.5)} atyantasaṃyoge kālādhvānau karmavat bhavataḥ iti vaktavyam . \\
\noindent \textbf{(P\_2,3.5)} kim prayojanam . \\
\noindent \textbf{(P\_2,3.5)} lādyartham . \\
\noindent \textbf{(P\_2,3.5)} lādibhiḥ abhidhānam yathā syāt . \\
\noindent \textbf{(P\_2,3.5)} āsyate māsaḥ . \\
\noindent \textbf{(P\_2,3.5)} śayyate krośaḥ . \\
\noindent \textbf{(P\_2,3.5)} atha vatkaraṇam kimartham . \\
\noindent \textbf{(P\_2,3.5)} svāśrayam api yathā syāt . \\
\noindent \textbf{(P\_2,3.5)} āsyate māsam . \\
\noindent \textbf{(P\_2,3.5)} śayyate krośam . \\
\noindent \textbf{(P\_2,3.5)} akarmakāṇām bhāve laḥ bhavati iti bhāve laḥ yathā syāt . \\
\noindent \textbf{(P\_2,3.5)} tat tarhi vaktavyam . \\
\noindent \textbf{(P\_2,3.5)} na vaktavyam . \\
\noindent \textbf{(P\_2,3.5)} prākṛtameva etat karma yathā kaṭam karoti śakaṭam karoti iti . \\
\noindent \textbf{(P\_2,3.5)} evam manyate . \\
\noindent \textbf{(P\_2,3.5)} yatra kaḥ citkriyākṛtaḥ viśeṣaḥ upajāyate tat nyāyyam karma iti . \\
\noindent \textbf{(P\_2,3.5)} na ca iha kaḥ cit kriyākṛtaḥ viśeṣaḥ upajāyate . \\
\noindent \textbf{(P\_2,3.5)} na evam śakyam . \\
\noindent \textbf{(P\_2,3.5)} iha api na syāt . \\
\noindent \textbf{(P\_2,3.5)} ādityam paśyati . \\
\noindent \textbf{(P\_2,3.5)} himavantam śṛṇoti . \\
\noindent \textbf{(P\_2,3.5)} grāmam gacchati . \\
\noindent \textbf{(P\_2,3.5)} tasmāt prākṛtameva etat karma yathā kaṭam karoti śakaṭam karoti iti . \\
\noindent \textbf{(P\_2,3.5)} yadi tarhi prākṛtam eva etat karma akarmakāṇām bhāve laḥ bhavati iti bhāve laḥ na prāpnoti. āsyate māsam devadattena iti . \\
\noindent \textbf{(P\_2,3.5)} tat tarhi vaktavyam . \\
\noindent \textbf{(P\_2,3.5)} na vaktavyam . \\
\noindent \textbf{(P\_2,3.5)} akarmakāṇām iti ucyate na ca ke cit kālabhāvādhvabhiḥ akarmakāḥ . \\
\noindent \textbf{(P\_2,3.5)} te evam vijñāsyāmaḥ . \\
\noindent \textbf{(P\_2,3.5)} kva cit ye akarmakāḥ iti . \\
\noindent \textbf{(P\_2,3.5)} atha vā yena karmaṇā sakarmakāḥ ca akarmakāḥ ca bhavanti tena akarmakāṇām . \\
\noindent \textbf{(P\_2,3.5)} na ca etena karmaṇā kaḥ cid api akarmakaḥ . \\
\noindent \textbf{(P\_2,3.5)} atha vā yat karma bhavati na ca bhavati tena karmakāṇām . \\
\noindent \textbf{(P\_2,3.5)} na ca etat karma kva cit api na bhavati . \\
\noindent \textbf{(P\_2,3.5)} na tarhi idānīm idam sūtram vaktavyam . \\
\noindent \textbf{(P\_2,3.5)} vaktavyam ca . \\
\noindent \textbf{(P\_2,3.5)} kim prayojanam . \\
\noindent \textbf{(P\_2,3.5)} yatra akriyayā atyantasaṃyogaḥ tadartham . \\
\noindent \textbf{(P\_2,3.5)} krośam kuṭilā nadī . \\
\noindent \textbf{(P\_2,3.5)} krośam ramaṇīyā vanarājiḥ . \\
\noindent \textbf{(P\_2,3.6)} kriyāparvarge iti vaktavyam . \\
\noindent \textbf{(P\_2,3.6)} sādhanāpavarge mā bhūt . \\
\noindent \textbf{(P\_2,3.6)} māsam adhītaḥ anuvāko na ca anena gṛhītaḥ iti . \\
\noindent \textbf{(P\_2,3.7)} kriyāmadhye iti vaktavyam . \\
\noindent \textbf{(P\_2,3.7)} iha api yathā syāt . \\
\noindent \textbf{(P\_2,3.7)} adya devadattaḥ bhuktvā dvyahāt bhoktā dvyahe bhoktā . \\
\noindent \textbf{(P\_2,3.7)} kārakamadhye iti iyati ucyamāne iha eva syāt : ihasthaḥ ayam iṣvāsaḥ krośāt lakṣyam vidhyati krośe lakṣyam vidhyati . \\
\noindent \textbf{(P\_2,3.7)} yam ca vidhyati yataḥ ca vidhyati ubhayoḥ tanmadhyam bhavati . \\
\noindent \textbf{(P\_2,3.7)} tat tarhi vaktavyam . \\
\noindent \textbf{(P\_2,3.7)} na vaktavyam . \\
\noindent \textbf{(P\_2,3.7)} na antareṇa sādhanam kriyāyāḥ pravṛttiḥ bhavati . \\
\noindent \textbf{(P\_2,3.7)} kriyāmadhyam cet kārakamadhyam api bhavati tatra kārakamadhye iti eva siddham . \\
\noindent \textbf{(P\_2,3.8)} karmapravacanīyayukte pratyādibhiḥ ca lakṣaṇādiṣu upasaṅkhyānam saptamīpañcamyoḥ pratiṣedhārtham . \\
\noindent \textbf{(P\_2,3.8)} karmapravacanīyayukte pratyādibhiḥ ca lakṣaṇādiṣu upasaṅkhyānam kartavyam . \\
\noindent \textbf{(P\_2,3.8)} vṛkṣam prati vidyotate vidyut . \\
\noindent \textbf{(P\_2,3.8)} vṛkṣam pari . \\
\noindent \textbf{(P\_2,3.8)} vṛkṣamanu . \\
\noindent \textbf{(P\_2,3.8)} sādhuḥ devadattaḥ mātaram prati . \\
\noindent \textbf{(P\_2,3.8)} mātaram pari . \\
\noindent \textbf{(P\_2,3.8)} mātaram anu . \\
\noindent \textbf{(P\_2,3.8)} kim prayojanam . \\
\noindent \textbf{(P\_2,3.8)} saptamīpañcamyoḥ pratiṣedhārtham . \\
\noindent \textbf{(P\_2,3.8)} saptamīpañcamyau mā bhūtām iti . \\
\noindent \textbf{(P\_2,3.8)} sādhunipuṇābhyām arcāyām saptamī iti saptamī . \\
\noindent \textbf{(P\_2,3.8)} pañcamī apāṅparibhiḥ iti pañcamī . \\
\noindent \textbf{(P\_2,3.8)} tatra ayam api arthaḥ aprateḥ iti na vaktavyam bhavati . \\
\noindent \textbf{(P\_2,3.8)} tat tarhi vaktavyam . \\
\noindent \textbf{(P\_2,3.8)} na vaktavyam . \\
\noindent \textbf{(P\_2,3.8)} uktaṃ vā . \\
\noindent \textbf{(P\_2,3.8)} kim uktam . \\
\noindent \textbf{(P\_2,3.8)} ekatra tāvat uktam aprateḥ iti . \\
\noindent \textbf{(P\_2,3.8)} itaratra api yadi api tāvat ayam pariḥ dṛṣṭāpacāraḥ varjane ca avarjane ca ayaṃ khalu apaśabdaḥ adṛṣṭāpacāraḥ varjanārthaḥ eva . \\
\noindent \textbf{(P\_2,3.8)} tasya kaḥ anyaḥ dvitīyaḥ sahāyaḥ bhavitum arhati anyat ataḥ varjanārthāt . \\
\noindent \textbf{(P\_2,3.8)} tat yathā . \\
\noindent \textbf{(P\_2,3.8)} asya goḥ dvitīyena arthaḥ iti gauḥ eva ānīyate na aśvaḥ na gadarbhaḥ . \\
\noindent \textbf{(P\_2,3.9)} katham idam vijñāyate . \\
\noindent \textbf{(P\_2,3.9)} yasya ca aiśvaryam īśvaratā īśvarabhāvaḥ tasmāt karmapravacanīyayuktāt iti . \\
\noindent \textbf{(P\_2,3.9)} āhosvit yasya svasya īśvaraḥ tasmāt karmapravacanīyayuktāt iti . \\
\noindent \textbf{(P\_2,3.9)} kaḥ ca atra viśeṣaḥ . \\
\noindent \textbf{(P\_2,3.9)} yasya ca īśvaravacanam iti kartṛnirdeśaḥ cet avacanāt siddham . \\
\noindent \textbf{(P\_2,3.9)} yasya ca īśvaravacanam iti kartṛnirdeśaḥ cet antareṇa vacanam siddham . \\
\noindent \textbf{(P\_2,3.9)} adhi brahmadatte pañcālāḥ . \\
\noindent \textbf{(P\_2,3.9)} ādhṛtāḥ te tasmin bhavanti . \\
\noindent \textbf{(P\_2,3.9)} satyam evam etat . \\
\noindent \textbf{(P\_2,3.9)} nityam parigrahītavyam parigrahītradhīnam bhavati . \\
\noindent \textbf{(P\_2,3.9)} prathamānupapattiḥ tu . prathamā na upapadyate . \\
\noindent \textbf{(P\_2,3.9)} kutaḥ . \\
\noindent \textbf{(P\_2,3.9)} pañcālebhyaḥ . \\
\noindent \textbf{(P\_2,3.9)} kā tarhi syāt . \\
\noindent \textbf{(P\_2,3.9)} ṣaṣṭhīsaptamyau . \\
\noindent \textbf{(P\_2,3.9)} svāmīśvarādhipati iti . \\
\noindent \textbf{(P\_2,3.9)} na tatra adhiśabdaḥ paṭhyate . \\
\noindent \textbf{(P\_2,3.9)} yadi api na paṭhyate adhiḥ īśvaravācī . \\
\noindent \textbf{(P\_2,3.9)} na tatra paryāyavacanānām grahaṇam . \\
\noindent \textbf{(P\_2,3.9)} katham jñāyate . \\
\noindent \textbf{(P\_2,3.9)} yat ayam kasya cit paryāyavacanasya grahaṇam karoti : adhipatidāyāda iti . \\
\noindent \textbf{(P\_2,3.9)} ṣaṣṭhī tarhi prāpnoti . \\
\noindent \textbf{(P\_2,3.9)} śeṣalakṣaṇā ṣaṣṭhī aśeṣatvāt na bhaviṣyati . \\
\noindent \textbf{(P\_2,3.9)} dvitīyā tarhi prāpnoti karmapravacanīyayukte dvitīyā iti . \\
\noindent \textbf{(P\_2,3.9)} saptamyā uktatvāt tasya abhisambandhasya dvitīyā na bhaviṣyati . \\
\noindent \textbf{(P\_2,3.9)} bhavet yaḥ adheḥ brahamadattasya ca abhisambandhaḥ saḥ saptamyā uktaḥ syāt . \\
\noindent \textbf{(P\_2,3.9)} yaḥ tu khalu adheḥ pañcālānām ca abhisambandhaḥ tatra dvitīyā prāpnoti . \\
\noindent \textbf{(P\_2,3.9)} svavacanāt siddham . astu yasya svasya īśvaraḥ tasmāt karmapravacanīyayuktāt iti . \\
\noindent \textbf{(P\_2,3.9)} evam api antareṇa vacanam siddham . \\
\noindent \textbf{(P\_2,3.9)} adhi brahmadattaḥ pañcāleṣu . \\
\noindent \textbf{(P\_2,3.9)} ādhṛtaḥ sa teṣu bhavati . \\
\noindent \textbf{(P\_2,3.9)} satyam evam etat . \\
\noindent \textbf{(P\_2,3.9)} nityam parigrahītā parigrahītavyādhīnaḥ bhavati . \\
\noindent \textbf{(P\_2,3.9)} prathamānupapattiḥ tu . \\
\noindent \textbf{(P\_2,3.9)} prathamā na upapadyate . \\
\noindent \textbf{(P\_2,3.9)} kutaḥ . \\
\noindent \textbf{(P\_2,3.9)} brahmadattāt . \\
\noindent \textbf{(P\_2,3.9)} kā tarhi syāt . \\
\noindent \textbf{(P\_2,3.9)} ṣaṣṭhīsaptamyau . \\
\noindent \textbf{(P\_2,3.9)} svāmīśvarādhipati iti . \\
\noindent \textbf{(P\_2,3.9)} na tatra adhiśabdaḥ paṭhyate . \\
\noindent \textbf{(P\_2,3.9)} yadi api na paṭhyate adhiḥ īśvaravācī . \\
\noindent \textbf{(P\_2,3.9)} na tatra paryāyavacanānām grahaṇam . \\
\noindent \textbf{(P\_2,3.9)} katham jñāyate . \\
\noindent \textbf{(P\_2,3.9)} yat ayam kasya cit paryāyavacanasya grahaṇam karoti . \\
\noindent \textbf{(P\_2,3.9)} adhipatidāyāda iti . \\
\noindent \textbf{(P\_2,3.9)} ṣaṣṭhī tarhi prāpnoti . \\
\noindent \textbf{(P\_2,3.9)} śeṣalakṣaṇā ṣaṣṭhī aśeṣatvāt na bhaviṣyati . \\
\noindent \textbf{(P\_2,3.9)} dvitīyā tarhi prāpnoti karmapravacanīyayukte dvitīyā iti . \\
\noindent \textbf{(P\_2,3.9)} saptamyā uktatvāt tasya abhisambandhasya dvitīyā na bhaviṣyati . \\
\noindent \textbf{(P\_2,3.9)} bhavet yaḥ dheḥ pañcālānām ca abhisambandhaḥ saḥ saptamyā uktaḥ syāt yaḥ tu khalu adheḥ brahmadattasya ca abhisaṃbandhaḥ tatra dvitīyā prāpnoti . \\
\noindent \textbf{(P\_2,3.9)} evaṃ tarhi svavacanāt siddham . \\
\noindent \textbf{(P\_2,3.9)} adhiḥ svam prati karmapravacanīyasaṃjñaḥ bhavati iti vaktavyam . \\
\noindent \textbf{(P\_2,3.9)} evam api yadā brahmadatte adhikaraṇe saptamī tadā pañcālebhyaḥ dvitīyā prāpnoti karmapravacanīyayukte dvitīya iti . \\
\noindent \textbf{(P\_2,3.9)} upapadavibhakteḥ kārakavibhaktiḥ balīyasī iti prathamā bhaviṣyati \\
\noindent \textbf{(P\_2,3.12)} adhvani arthagrahaṇam . adhvani arthagrahaṇam kartavyam . \\
\noindent \textbf{(P\_2,3.12)} iha api mā bhūt . \\
\noindent \textbf{(P\_2,3.12)} panthānam gacchati . \\
\noindent \textbf{(P\_2,3.12)} vīvadham gacchat i iti . \\
\noindent \textbf{(P\_2,3.12)} āsthitapratiṣedhaḥ ca . \\
\noindent \textbf{(P\_2,3.12)} āsthitapratiṣedhaḥ ca ayam vaktavyaḥ . \\
\noindent \textbf{(P\_2,3.12)} yaḥ hi utpathena panthānam gacchati pathe gacchati iti eva tatra bhavitavyam . \\
\noindent \textbf{(P\_2,3.12)} kim artham punaḥ idam ucyate . \\
\noindent \textbf{(P\_2,3.12)} caturthī yathā syāt . \\
\noindent \textbf{(P\_2,3.12)} atha dvitīyā siddhā . \\
\noindent \textbf{(P\_2,3.12)} siddhā karmaṇi iti eva . \\
\noindent \textbf{(P\_2,3.12)} caturthī api siddhā . \\
\noindent \textbf{(P\_2,3.12)} | katham . \\
\noindent \textbf{(P\_2,3.12)} sampradāne iti eva . \\
\noindent \textbf{(P\_2,3.12)} na sidhyati . \\
\noindent \textbf{(P\_2,3.12)} karmaṇā yam abhipraiti saḥ saṃpradānam iti ucyate . \\
\noindent \textbf{(P\_2,3.12)} kriyayā ca asau grāmam abhipraiti . \\
\noindent \textbf{(P\_2,3.12)} kayā kriyayā . \\
\noindent \textbf{(P\_2,3.12)} gamikriyayā . \\
\noindent \textbf{(P\_2,3.12)} kriyāgrahaṇam api tatra codyate . \\
\noindent \textbf{(P\_2,3.12)} ceṣṭāyām anadhvani striyam gacchati ajām nayati iti atiprasaṅgaḥ . \\
\noindent \textbf{(P\_2,3.12)} ceṣṭāyām anadhvani striyam gacchati ajām nayati iti atiprasaṅgaḥ bhavati . \\
\noindent \textbf{(P\_2,3.12)} siddham tu asamprāptavacanāt . \\
\noindent \textbf{(P\_2,3.12)} siddham etat . \\
\noindent \textbf{(P\_2,3.12)} katham . \\
\noindent \textbf{(P\_2,3.12)} asamprāpte karmaṇi dvitīyācaturthyau bhavataḥ iti vaktavyam . \\
\noindent \textbf{(P\_2,3.12)} adhvanaḥ ca anapavādaḥ . \\
\noindent \textbf{(P\_2,3.12)} evam ca kṛtvā anadhvani iti etat api na vaktavyam bhavati . \\
\noindent \textbf{(P\_2,3.12)} samprāptam hi etat karma adhvānam gacchati iti . \\
\noindent \textbf{(P\_2,3.13)} caturthīvidhāne tādarthye upasaṅkhyānam . caturthīvidhāne tādarthye upasaṅkhyānam kartavyam . \\
\noindent \textbf{(P\_2,3.13)} yūpāya dāru kuṇḍalāya hiraṇyam . \\
\noindent \textbf{(P\_2,3.13)} kim idam tādarthyam iti . \\
\noindent \textbf{(P\_2,3.13)} tadarthasya bhāvaḥ tādarthyam . \\
\noindent \textbf{(P\_2,3.13)} tadartham punaḥ kim . \\
\noindent \textbf{(P\_2,3.13)} sarvanāmnaḥ ayam caturthyantasya arthaśabdena saha samāsaḥ . \\
\noindent \textbf{(P\_2,3.13)} katham ca atra caturthī . \\
\noindent \textbf{(P\_2,3.13)} anena eva . \\
\noindent \textbf{(P\_2,3.13)} yadi evam itaretarāśrayam bhavati . \\
\noindent \textbf{(P\_2,3.13)} kā itaretarāśrayatā . \\
\noindent \textbf{(P\_2,3.13)} nirdeśottarakālam caturthyā bhavitavyam caturthyā ca nirdeśaḥ tat itaretarāśrayam bhavati . \\
\noindent \textbf{(P\_2,3.13)} itaretarāśrayāṇi ca na prakalpante . \\
\noindent \textbf{(P\_2,3.13)} tat tarhi vaktavyam . \\
\noindent \textbf{(P\_2,3.13)} na vaktavyam . \\
\noindent \textbf{(P\_2,3.13)} ācāryapravṛttiḥ jñāpayati bhavati arthaśabdena yoge caturthī iti yat ayam carturthī tadarthārtha iti caturthyantasya arthaśabdena saha samāsam śāsti . \\
\noindent \textbf{(P\_2,3.13)} na khalu api avaśyaṃ caturthyantasya eva arthaśabdena saha samāsaḥ bhavati . \\
\noindent \textbf{(P\_2,3.13)} kim tarhi . \\
\noindent \textbf{(P\_2,3.13)} ṣaṣṭhyantasya api bhavati . \\
\noindent \textbf{(P\_2,3.13)} tat yathā . \\
\noindent \textbf{(P\_2,3.13)} guroḥ idam gurvartham iti . \\
\noindent \textbf{(P\_2,3.13)} yadi tādarthye upasaṅkhyānam kriyate na arthaḥ sampradānagrahaṇena . \\
\noindent \textbf{(P\_2,3.13)} yaḥ api hi upādhyāyāya gauḥ dīyate upādhyāyārthaḥ saḥ bhavati . \\
\noindent \textbf{(P\_2,3.13)} tatra tādarthye iti eva siddham . \\
\noindent \textbf{(P\_2,3.13)} avaśyam saṃpradānagrahaṇam kartavyam yā anyena lakṣaṇena sampradānasañjñā tadartham . \\
\noindent \textbf{(P\_2,3.13)} chātrāya rucitam . \\
\noindent \textbf{(P\_2,3.13)} chātrāya svaditam iti . \\
\noindent \textbf{(P\_2,3.13)} tat tarhi upasaṅkhyānam kartavyam . \\
\noindent \textbf{(P\_2,3.13)} na kartavyam . \\
\noindent \textbf{(P\_2,3.13)} ācāryapravṛttiḥ jñāpayati bhavati tādarthye caturthī iti yat ayam caturthī tadarthārtha iti caturthyantasya tadarthena saha samāsam śāsti . \\
\noindent \textbf{(P\_2,3.13)} kḷpi sampadyamāne . \\
\noindent \textbf{(P\_2,3.13)} kḷpi sampadyamāne caturthī vaktavyā . \\
\noindent \textbf{(P\_2,3.13)} mūtrāya kalpate yavāgūḥ . \\
\noindent \textbf{(P\_2,3.13)} uccārāya kalpate yavānnam iti . \\
\noindent \textbf{(P\_2,3.13)} utpātena jñāpyamāne . \\
\noindent \textbf{(P\_2,3.13)} utpātena jñāpyamāne caturthī vaktavyā . \\
\noindent \textbf{(P\_2,3.13)} vātāya kapilā vidyut ātapāya atilohinī pītā bhavati sasyāya durbhikṣāya sitā bhavet . \\
\noindent \textbf{(P\_2,3.13)} māṃsaudanāya vyāharati mṛgaḥ . \\
\noindent \textbf{(P\_2,3.13)} hitayoge ca . \\
\noindent \textbf{(P\_2,3.13)} hitayoge caturthī vaktavyā . \\
\noindent \textbf{(P\_2,3.13)} | hitam arocakine hitam āmayāvine . \\
\noindent \textbf{(P\_2,3.16)} svastiyoge caturthī kuśalārthaiḥ āśiṣi vāvidhānāt . \\
\noindent \textbf{(P\_2,3.16)} svastiyoge caturthī kuśalārthaiḥ āśiṣi vāvidhānāt bhavati vipratiṣedhena . \\
\noindent \textbf{(P\_2,3.16)} svastiyoge caturthyāḥ avakāśaḥ svasti jālmāya svasti vṛṣalāya . \\
\noindent \textbf{(P\_2,3.16)} kuśalārthaiḥ āśiṣi vāvidhānasya avakāśaḥ anye kuśalārthāḥ . \\
\noindent \textbf{(P\_2,3.16)} kuśalam devattāya kuśalam devadattasya . \\
\noindent \textbf{(P\_2,3.16)} iha ubhayam prāpnoti . \\
\noindent \textbf{(P\_2,3.16)} svasti gobhyaḥ svasti brāhmaṇebhyaḥ iti . \\
\noindent \textbf{(P\_2,3.16)} caturthī bhavati vipratiṣedhena . \\
\noindent \textbf{(P\_2,3.16)} alamiti paryāptyarthagrahaṇam . \\
\noindent \textbf{(P\_2,3.16)} alamiti paryāptyarthagrahaṇam kartavyam . \\
\noindent \textbf{(P\_2,3.16)} iha mā bhūt . \\
\noindent \textbf{(P\_2,3.16)} alaṅkurute kanyām iti . \\
\noindent \textbf{(P\_2,3.16)} aparaḥ āha : alam iti paryāptyarthagrahaṇam kartavyam . \\
\noindent \textbf{(P\_2,3.16)} iha api yathā syāt . \\
\noindent \textbf{(P\_2,3.16)} alam mallaḥ mallāya . \\
\noindent \textbf{(P\_2,3.16)} prabhuḥmallaḥ mallāya . \\
\noindent \textbf{(P\_2,3.16)} prabhavati mallaḥ mallāya iti . \\
\noindent \textbf{(P\_2,3.17)} aprāṇiṣu iti ucyate . \\
\noindent \textbf{(P\_2,3.17)} tatra idam na sidhyati : na tvā śvānam manye , na tvā śune manye iti . \\
\noindent \textbf{(P\_2,3.17)} evam tarhi yogavibhāgaḥ kariṣyate . \\
\noindent \textbf{(P\_2,3.17)} manyakarmaṇi anādare vibhāṣā . \\
\noindent \textbf{(P\_2,3.17)} tataḥ aprāṇiṣu . \\
\noindent \textbf{(P\_2,3.17)} aprāṇiṣu ca vibhāṣā iti . \\
\noindent \textbf{(P\_2,3.17)} iha api tarhi prāpnoti : na tvā kākam manye , na tvā śukam manye iti . \\
\noindent \textbf{(P\_2,3.17)} yat etat aprāṇiṣu iti etat anāvādiṣu iti vakṣyāmi . \\
\noindent \textbf{(P\_2,3.17)} ime ca nāvādayaḥ bhaviṣyanti . \\
\noindent \textbf{(P\_2,3.17)} na tvā nāvam manye yāvat tīrṇam na nāvyam . \\
\noindent \textbf{(P\_2,3.17)} na tvā annam manye yāvat bhuktam na śrāddham . \\
\noindent \textbf{(P\_2,3.17)} atra yeṣu prāṇiṣu na iṣyate te nāvādayaḥ bhaviṣyanti . \\
\noindent \textbf{(P\_2,3.17)} manyakarmaṇi prakṛṣyakutsitagrahaṇam . manyakarmaṇi prakṛṣyakutsitagrahaṇam kartavyam iha mā bhūt : tvām tṛṇam manye iti . \\
\noindent \textbf{(P\_2,3.18)} tṛtīyāvidhāne prakṛtyādibhyaḥ upasaṅkhyānam . \\
\noindent \textbf{(P\_2,3.18)} tṛtīyāvidhāne prakṛtyādibhyaḥ upasaṅkhyānam kartavyam . \\
\noindent \textbf{(P\_2,3.18)} tṛtīyāvidhāne prakṛtyādibhyaḥ upasaṅkhyānam . \\
\noindent \textbf{(P\_2,3.18)} prakṛtyā abhirūpaḥ prakṛtyā darśanīyaḥ . \\
\noindent \textbf{(P\_2,3.18)} prāyeṇa yājñikāḥ prāyeṇa vaiyākaraṇāḥ . \\
\noindent \textbf{(P\_2,3.18)} māṭharaḥ asmi gotreṇa. gārgyaḥ asmi gotreṇa . \\
\noindent \textbf{(P\_2,3.18)} samena dhāvati . \\
\noindent \textbf{(P\_2,3.18)} viṣameṇa dhāvati . \\
\noindent \textbf{(P\_2,3.18)} dvidroṇena dhānyam krīṇāti . \\
\noindent \textbf{(P\_2,3.18)} tridroṇena dhānyam krīṇāti . \\
\noindent \textbf{(P\_2,3.18)} pañcakena paśūn krīṇāti . \\
\noindent \textbf{(P\_2,3.18)} sāhasreṇa aśvān krīṇāti. tat tarhi vaktavyam . \\
\noindent \textbf{(P\_2,3.18)} na vaktavyam . \\
\noindent \textbf{(P\_2,3.18)} kartṛkaraṇayoḥ tṛtīyā iti eva siddham . \\
\noindent \textbf{(P\_2,3.18)} iha tāvat prakṛtyā abhirūpaḥ prakṛtyā darśanīyaḥ iti prakṛtikṛtam tasya ābhirūpyam . \\
\noindent \textbf{(P\_2,3.18)} prāyeṇa yājñikāḥ prāyeṇa vaiyākaraṇāḥ iti . \\
\noindent \textbf{(P\_2,3.18)} eṣaḥ tatra prāyaḥ yena te adhīyate . \\
\noindent \textbf{(P\_2,3.18)} māṭharaḥ asmi gotreṇa. gārgyaḥ asmi gotreṇa iti . \\
\noindent \textbf{(P\_2,3.18)} etena aham sañjñāye . \\
\noindent \textbf{(P\_2,3.18)} samena dhāvati . \\
\noindent \textbf{(P\_2,3.18)} viṣameṇa dhāvati . \\
\noindent \textbf{(P\_2,3.18)} idam atra prayoktavyam sat na prayujyate samena pathā dhāvati viṣameṇa pathā dhāvatīti . \\
\noindent \textbf{(P\_2,3.18)} dvidroṇena dhānyam krīṇāti . \\
\noindent \textbf{(P\_2,3.18)} tridroṇena dhānyam krīṇāti . \\
\noindent \textbf{(P\_2,3.18)} tādarthyāt tācchabdyam . \\
\noindent \textbf{(P\_2,3.18)} dvidroṇārtham dvidroṇam . \\
\noindent \textbf{(P\_2,3.18)} dvidroṇena hiraṇyena dhānyam krīṇāti iti . \\
\noindent \textbf{(P\_2,3.18)} pañcakena paśūn krīṇāti iti . \\
\noindent \textbf{(P\_2,3.18)} atra api tādarthyāt tācchabdyam . \\
\noindent \textbf{(P\_2,3.18)} pañcapaśvarthaḥ pañcakaḥ . \\
\noindent \textbf{(P\_2,3.18)} pañcakena paśūn krīṇāti iti . \\
\noindent \textbf{(P\_2,3.18)} sāhasreṇa aśvān krīṇāti iti. sahasraparimāṇam sāhasram . \\
\noindent \textbf{(P\_2,3.18)} sāhasreṇa hiraṇyena aśvān krīṇāti iti. \\
\noindent \textbf{(P\_2,3.19)} kim udāharaṇam . \\
\noindent \textbf{(P\_2,3.19)} tilaiḥ saha māṣān vapati iti . \\
\noindent \textbf{(P\_2,3.19)} na etat asti . \\
\noindent \textbf{(P\_2,3.19)} tilaiḥ miśrīkṛtya māṣāḥ upyante . \\
\noindent \textbf{(P\_2,3.19)} tatra karaṇe iti eva siddham . \\
\noindent \textbf{(P\_2,3.19)} idam tarhi . \\
\noindent \textbf{(P\_2,3.19)} putreṇa saha āgataḥ devadattaḥ iti . \\
\noindent \textbf{(P\_2,3.19)} apradhāne kartari tṛtīyā yathā syāt . \\
\noindent \textbf{(P\_2,3.19)} etat api na asti prayojanam . \\
\noindent \textbf{(P\_2,3.19)} pradhāne kartari lādayaḥ bhavanti iti pradhānakartā ktena abhidhīyate yaḥ ca apradhānam siddhā tatra kartari iti eva tṛtīyā . \\
\noindent \textbf{(P\_2,3.19)} idam tarhi . \\
\noindent \textbf{(P\_2,3.19)} putreṇa saha āgamanam devadattasya iti . \\
\noindent \textbf{(P\_2,3.19)} ṣaṣṭhī atra bādhikā bhaviṣyati . \\
\noindent \textbf{(P\_2,3.19)} idam tarhi . \\
\noindent \textbf{(P\_2,3.19)} putreṇa saha sthūlaḥ . \\
\noindent \textbf{(P\_2,3.19)} putreṇa saha piṅgalaḥ iti . \\
\noindent \textbf{(P\_2,3.19)} idam ca api udāharaṇam tilaiḥ saha māṣān vapati iti . \\
\noindent \textbf{(P\_2,3.19)} nanu ca uktam tilaiḥ miśrīkṛtya māṣāḥ upyante . \\
\noindent \textbf{(P\_2,3.19)} tatra karaṇe iti eva siddham iti . \\
\noindent \textbf{(P\_2,3.19)} bhavet siddham yadā tilaiḥ miśrīkṛtya upyeran . \\
\noindent \textbf{(P\_2,3.19)} yadā tu khalu kasya cin māṣabījāvāpaḥ upasthitaḥ tadartham ca kṣetram upārjitam tatra anyat api kiṃ cid upyate yadi bhaviṣyati bhaviṣyati iti tadā na sidhyati . \\
\noindent \textbf{(P\_2,3.19)} sahayukte apradhānavacanam anarthakam upapadavibhakteḥ kārakavibhaktibalīyastvāt anyatra api . \\
\noindent \textbf{(P\_2,3.19)} sahayukte apradhānavacanam anarthakam . \\
\noindent \textbf{(P\_2,3.19)} kiṃ kāraṇam . \\
\noindent \textbf{(P\_2,3.19)} upapadavibhakteḥ kārakavibhaktibalīyastvāt . \\
\noindent \textbf{(P\_2,3.19)} anyatra api kārakavibhaktirbalīyasī iti prathamā bhaviṣyati . \\
\noindent \textbf{(P\_2,3.19)} kva anyatra . \\
\noindent \textbf{(P\_2,3.19)} gāḥ svāmī vrajati iti . \\
\noindent \textbf{(P\_2,3.20)} iha kasmāt na bhavati . \\
\noindent \textbf{(P\_2,3.20)} akṣi kāṇam asya iti . \\
\noindent \textbf{(P\_2,3.20)} aṅgāt vikṛtāt tadvikārataḥ cet aṅginaḥ vacanam . \\
\noindent \textbf{(P\_2,3.20)} aṅgāt vikṛtāt tṛtīyā vaktavyā tena eva cet vikāreṇa aṅgī dyotyate iti vaktavyam . \\
\noindent \textbf{(P\_2,3.20)} tat tarhi vaktavyam . \\
\noindent \textbf{(P\_2,3.20)} na vaktavyam . \\
\noindent \textbf{(P\_2,3.20)} aṅgaśabdaḥ ayam samudāyaśabdaḥ yena iti ca karaṇe eṣā tṛtīyā . \\
\noindent \textbf{(P\_2,3.20)} yena avayavena samudāyaḥ aṅgī dyotyate tasmin bhavitavyam na ca etena avayavena samudāyaḥ dyotyate . \\
\noindent \textbf{(P\_2,3.21)} itthambhūtalakṣaṇe tatsthe pratiṣedhaḥ . \\
\noindent \textbf{(P\_2,3.21)} itthambhūtalakṣaṇe tatsthe pratiṣedhaḥ vaktavyaḥ . \\
\noindent \textbf{(P\_2,3.21)} api bhavānkamaṇḍalupāṇim chātrama drākṣīt iti . \\
\noindent \textbf{(P\_2,3.21)} na vā itthambhūtasya lakṣaṇena apṛthagbhāvāt . \\
\noindent \textbf{(P\_2,3.21)} na vā vaktavyam . \\
\noindent \textbf{(P\_2,3.21)} kim kāraṇam . \\
\noindent \textbf{(P\_2,3.21)} itthambhūtasya lakṣaṇena apṛthagbhāvāt . \\
\noindent \textbf{(P\_2,3.21)} yatra itthambhūtasya pṛthagbhūtam lakṣaṇam tatra bhavitavyam . \\
\noindent \textbf{(P\_2,3.21)} na ca atra itthambhūtasya pṛthagbhūtam lakṣaṇam . \\
\noindent \textbf{(P\_2,3.21)} kim vaktavyam etat . \\
\noindent \textbf{(P\_2,3.21)} na hi . \\
\noindent \textbf{(P\_2,3.21)} katham anucyamānam gaṃsyate . \\
\noindent \textbf{(P\_2,3.21)} tathā hi ayam prādhānyena lakṣaṇam pratinirdiśati . \\
\noindent \textbf{(P\_2,3.21)} itthambhūtasya lakṣaṇam itthaṃbhūtalakṣaṇam tasmin nitthaṃbhūtalakṣaṇe iti . \\
\noindent \textbf{(P\_2,3.22)} sañjñaḥ kṛtprayoge ṣaṣṭhī vipratiṣedhena . sañjñaḥ anyatarasyām karmaṇi iti etasmāt kṛprayoge ṣaṣṭhī bhavati vipratiṣedhena . \\
\noindent \textbf{(P\_2,3.22)} sañjñaḥ anyatarasyām iti asya avakāśaḥ . \\
\noindent \textbf{(P\_2,3.22)} mātaram sañjānīte . \\
\noindent \textbf{(P\_2,3.22)} mātrā sañjānīte . \\
\noindent \textbf{(P\_2,3.22)} kṛtprayoge ṣaṣṭhyāḥ avakāśaḥ . \\
\noindent \textbf{(P\_2,3.22)} idhmapravraścanaḥ palāśaśātanaḥ . \\
\noindent \textbf{(P\_2,3.22)} iha ubhayam prāpnoti . \\
\noindent \textbf{(P\_2,3.22)} mātuḥ sañjñātā . \\
\noindent \textbf{(P\_2,3.22)} pituḥ sañjñātā iti . \\
\noindent \textbf{(P\_2,3.22)} ṣaṣṭhī bhavati vipratiṣedhena . \\
\noindent \textbf{(P\_2,3.22)} upapadavibhakteḥ ca upapadavibhaktiḥ . upapadavibhakteḥ ca upapadavibhaktiḥ bhavati vipratiṣedhena | anyārāditarertadikśabdāñcūttarapadājāhiyukte iti asya avakāśaḥ . \\
\noindent \textbf{(P\_2,3.22)} anyaḥ devadattāt . \\
\noindent \textbf{(P\_2,3.22)} svāmīśvarādhipatidāyādasākṣipratibhūprasūtaiḥ ca iti asya avakāśaḥ . \\
\noindent \textbf{(P\_2,3.22)} goṣu svāmī . \\
\noindent \textbf{(P\_2,3.22)} gavāṃ svāmī . \\
\noindent \textbf{(P\_2,3.22)} iha ubhayam prāpnoti . \\
\noindent \textbf{(P\_2,3.22)} anyaḥ goṣu svāmī . \\
\noindent \textbf{(P\_2,3.22)} anyaḥ gavāṃ svāmī iti . \\
\noindent \textbf{(P\_2,3.22)} svāmīśvarādhipati iti etat bhavati vipratiṣedhena . \\
\noindent \textbf{(P\_2,3.22)} na eṣaḥ yuktḥ vipratiṣedhaḥ . \\
\noindent \textbf{(P\_2,3.22)} na hi atra gāvaḥ anyayuktāḥ . \\
\noindent \textbf{(P\_2,3.22)} kaḥ tarhi . \\
\noindent \textbf{(P\_2,3.22)} svāmī . \\
\noindent \textbf{(P\_2,3.22)} evam tarhi tulyārthaḥ atulopamābhyām tṛtīyā anyatarasyām iti asya avakāśaḥ . \\
\noindent \textbf{(P\_2,3.22)} tulyaḥ devadattasya . \\
\noindent \textbf{(P\_2,3.22)} tulyaḥ devadattena iti . \\
\noindent \textbf{(P\_2,3.22)} svāmīśvarādhipati iti asya vakāśaḥ saḥ eva . \\
\noindent \textbf{(P\_2,3.22)} iha ubhayam prāpnoti . \\
\noindent \textbf{(P\_2,3.22)} tulyaḥ gobhiḥ svāmī . \\
\noindent \textbf{(P\_2,3.22)} tulyaḥ gavāṃ svāmī iti . \\
\noindent \textbf{(P\_2,3.22)} tulyārthaḥ ratulopamābhyām iti etat bhavati vipratiṣedhena . \\
\noindent \textbf{(P\_2,3.23)} nimittakāraṇahetuṣu sarvāsām prāyadarśanam . nimittakāraṇahetuṣu sarvā vibhaktayaḥ prāyeṇa dṛśyante iti vaktavyam . \\
\noindent \textbf{(P\_2,3.23)} kim nimittam vasati . \\
\noindent \textbf{(P\_2,3.23)} kena nimittena vasati . \\
\noindent \textbf{(P\_2,3.23)} kasmai nimittāya vasati . \\
\noindent \textbf{(P\_2,3.23)} kasmāt nimittāt vasati . \\
\noindent \textbf{(P\_2,3.23)} kasya nimittasya vasati . \\
\noindent \textbf{(P\_2,3.23)} kasmin nimitte vasati . \\
\noindent \textbf{(P\_2,3.23)} kim kāraṇam vasati . \\
\noindent \textbf{(P\_2,3.23)} kena kāraṇena vasati . \\
\noindent \textbf{(P\_2,3.23)} kasmai kāraṇāya vasati . \\
\noindent \textbf{(P\_2,3.23)} kasmāt kāraṇāt vasati . \\
\noindent \textbf{(P\_2,3.23)} kasya kāraṇasya vasati . \\
\noindent \textbf{(P\_2,3.23)} kasmin kāraṇe vasati . \\
\noindent \textbf{(P\_2,3.23)} kaḥ hetuḥ vasati . \\
\noindent \textbf{(P\_2,3.23)} kam hetum vasati . \\
\noindent \textbf{(P\_2,3.23)} kena hetunā vasati . \\
\noindent \textbf{(P\_2,3.23)} kasmai hetave vasati . \\
\noindent \textbf{(P\_2,3.23)} kasmāt hetoḥ vasati . \\
\noindent \textbf{(P\_2,3.23)} kasya hetoḥ vasati . \\
\noindent \textbf{(P\_2,3.23)} kasmin hetau vasati . \\
\noindent \textbf{(P\_2,3.28)} pañcamīvidhāne lyablope karmaṇi upasaṅkhyānam . \\
\noindent \textbf{(P\_2,3.28)} pañcamīvidhāne lyablope karmaṇi upasaṅkhyānam kartavyam . \\
\noindent \textbf{(P\_2,3.28)} prāsādam āruhya prekṣate . \\
\noindent \textbf{(P\_2,3.28)} prāsādātprekṣate . \\
\noindent \textbf{(P\_2,3.28)} adhikaraṇe ca . \\
\noindent \textbf{(P\_2,3.28)} adhikaraṇe ca upasaṅkhyānam kartavyam . \\
\noindent \textbf{(P\_2,3.28)} āsanāt prekṣate . \\
\noindent \textbf{(P\_2,3.28)} śayanāt prekṣate . \\
\noindent \textbf{(P\_2,3.28)} praśnākhyānayoḥ ca . \\
\noindent \textbf{(P\_2,3.28)} praśnākhyānayoḥ ca pañcamī vaktavyā . \\
\noindent \textbf{(P\_2,3.28)} kutaḥ bhavān . \\
\noindent \textbf{(P\_2,3.28)} pāṭaliputrāt . \\
\noindent \textbf{(P\_2,3.28)} yataḥ ca adhvakālanirmāṇam . yataḥ ca adhvakālanirmāṇam tatra pañcamī vaktavyā . \\
\noindent \textbf{(P\_2,3.28)} gavīdhumataḥ sāvakāśyaīm catvāri yojanāni . \\
\noindent \textbf{(P\_2,3.28)} kārtikyāḥ āgrahāyaṇī māse . \\
\noindent \textbf{(P\_2,3.28)} tadyuktāt kāle saptamī . tadyuktāt kāle saptamī vaktavyā . \\
\noindent \textbf{(P\_2,3.28)} kārtikyāḥ āgrahāyaṇī māse . \\
\noindent \textbf{(P\_2,3.28)} adhvanaḥ prathamā ca . adhvanaḥ prathamā ca saptamī ca vaktayvā . \\
\noindent \textbf{(P\_2,3.28)} gavīdhumataḥ sāvakāśyaīm catvāri yojanāni caturṣu yojaneṣu . \\
\noindent \textbf{(P\_2,3.28)} tat tarhi idaṃ bahu vaktavyam . \\
\noindent \textbf{(P\_2,3.28)} na vaktavyam . \\
\noindent \textbf{(P\_2,3.28)} apādāne iti eva siddham . \\
\noindent \textbf{(P\_2,3.28)} iha tāvat prāsādāt prekṣate . \\
\noindent \textbf{(P\_2,3.28)} śayanāt prekṣate iti . \\
\noindent \textbf{(P\_2,3.28)} apakrāmati tat tasmāt darśanam . \\
\noindent \textbf{(P\_2,3.28)} yadi apakrāmati kim na atyantāya apakrāmati . \\
\noindent \textbf{(P\_2,3.28)} santatatvāt . \\
\noindent \textbf{(P\_2,3.28)} atha vā anyānyaprādurbhāvā . \\
\noindent \textbf{(P\_2,3.28)} praśnākhyānayoḥ ca pañcamī vaktavyā iti . \\
\noindent \textbf{(P\_2,3.28)} idam atra prayoktavyam sat na prayujyate . \\
\noindent \textbf{(P\_2,3.28)} kutaḥ bhavān āgacchati . \\
\noindent \textbf{(P\_2,3.28)} pāṭaliputrāt āgacchami iti . \\
\noindent \textbf{(P\_2,3.28)} yataḥ ca adhvakālanirmāṇam tatra pañcamī vaktavyā iti . \\
\noindent \textbf{(P\_2,3.28)} idam atra prayoktavyam sat na prayujyate gavīdhumataḥ niḥsṛtya sāṅkāśyam catvāri yojanāni . \\
\noindent \textbf{(P\_2,3.28)} kārtikyāḥ āgrahāyaṇī māse iti . \\
\noindent \textbf{(P\_2,3.28)} idam atra prayoktavyam sat na prayujyate . \\
\noindent \textbf{(P\_2,3.28)} kārttikyāḥ prabhṛti āgrahāyaṇī māsa iti . \\
\noindent \textbf{(P\_2,3.28)} tadyuktāt kāle saptamī vaktavyā iti . \\
\noindent \textbf{(P\_2,3.28)} idam atra prayoktavyam sat na prayujyate . \\
\noindent \textbf{(P\_2,3.28)} kārttikyāḥ āgrahāyaṇī gate māse iti . \\
\noindent \textbf{(P\_2,3.28)} adhvanaḥ prathamā ca saptamī ca iti . \\
\noindent \textbf{(P\_2,3.28)} idam atra prayoktavyam sat na prayujyate . \\
\noindent \textbf{(P\_2,3.28)} gavīdhumato niḥsṛtya yadā catvāri yojanāni gatāni bhavanti tataḥ sāṅkāśyam . \\
\noindent \textbf{(P\_2,3.28)} caturṣu yojaneṣu gateṣu sāṅkāśyam iti . \\
\noindent \textbf{(P\_2,3.29)} añcūttarapadagrahaṇam kimartham na dikśabdaiḥ yoge iti eva siddham . \\
\noindent \textbf{(P\_2,3.29)} ṣaṣṭhī atasarthapratyayena iti vakṣyati . \\
\noindent \textbf{(P\_2,3.29)} tasya ayam purastāt apakarṣaḥ . \\
\noindent \textbf{(P\_2,3.30)} arthagrahaṇam kimartham . \\
\noindent \textbf{(P\_2,3.30)} ṣaṣṭhī ataspratyayena iti ucyamāne iha eva syāt . \\
\noindent \textbf{(P\_2,3.30)} dakṣiṇato grāmasya uttarato grāmasya iti . \\
\noindent \textbf{(P\_2,3.30)} iha na syāt . \\
\noindent \textbf{(P\_2,3.30)} upari grāmasya upariṣṭāt grāmasya iti . \\
\noindent \textbf{(P\_2,3.30)} arthagrahaṇe punaḥ kriyamāṇe ataspratyayena ca siddham bhavati yaḥ ca anyaḥ tena samānārthaḥ . \\
\noindent \textbf{(P\_2,3.30)} atha pratyayagrahaṇam kimartham . \\
\noindent \textbf{(P\_2,3.30)} iha mā bhūt . \\
\noindent \textbf{(P\_2,3.30)} prāk grāmāt pratyak grāmāt . \\
\noindent \textbf{(P\_2,3.30)} añcūttarapadasya api etat prayojanam uktam . \\
\noindent \textbf{(P\_2,3.30)} tatra anyatarat śakyam akartum . \\
\noindent \textbf{(P\_2,3.32)} pṛthagādiṣu pañcamīvidhānam . \\
\noindent \textbf{(P\_2,3.32)} pṛthagādiṣu pañcamīvidheyā . \\
\noindent \textbf{(P\_2,3.32)} pṛthak devadattāt . \\
\noindent \textbf{(P\_2,3.32)} kimartham na prakṛtam pañcamīgrahaṇam anuvartate . \\
\noindent \textbf{(P\_2,3.32)} kva prakṛtam . \\
\noindent \textbf{(P\_2,3.32)} apādāne pañcamī iti. anadhikārāt . \\
\noindent \textbf{(P\_2,3.32)} anadhikāraḥ saḥ . \\
\noindent \textbf{(P\_2,3.32)} adhikāre hi dvitīyāṣaṣṭhīviṣaye pratiṣedhaḥ . \\
\noindent \textbf{(P\_2,3.32)} adhikāre hi dvitīyāṣaṣṭhīviṣaye pratiṣedhaḥ vaktavyaḥ syāt . \\
\noindent \textbf{(P\_2,3.32)} dakṣiṇena grāmam , dakṣiṇataḥ grāmasya . \\
\noindent \textbf{(P\_2,3.32)} evam tarhi anyatarasyāṅgrahaṇasāmarthyāt pañcamī bhaviṣyati . \\
\noindent \textbf{(P\_2,3.32)} asti anyat anyatarasyāṅgrahaṇasya prayojanam . \\
\noindent \textbf{(P\_2,3.32)} kim . \\
\noindent \textbf{(P\_2,3.32)} yasyām na aprāptāyām tṛtīyā ārabhyate sā yathā syāt . \\
\noindent \textbf{(P\_2,3.32)} kasyām ca na aprāptāyām . \\
\noindent \textbf{(P\_2,3.32)} antataḥ ṣaṣṭhyām . \\
\noindent \textbf{(P\_2,3.32)} tat tarhi vaktavyam . \\
\noindent \textbf{(P\_2,3.32)} na vaktavyam . \\
\noindent \textbf{(P\_2,3.32)} prakṛtam anuvartate . \\
\noindent \textbf{(P\_2,3.32)} kva prakṛtam . \\
\noindent \textbf{(P\_2,3.32)} apādāne pañcamī iti . \\
\noindent \textbf{(P\_2,3.32)} nanu ca uktam anadhikāraḥ saḥ adhikāre hi dvitīyāṣaṣṭhīviṣaye pratiṣedhaḥ iti . \\
\noindent \textbf{(P\_2,3.32)} evam tarhi sambandham anuvartiṣyate . \\
\noindent \textbf{(P\_2,3.32)} apādāne pañcamī . \\
\noindent \textbf{(P\_2,3.32)} anyārāditarartedikśabdāñcūttarapadājāhiyukte pañcamī . \\
\noindent \textbf{(P\_2,3.32)} ṣaṣṭhī atasarthapratyayena anyārādibhiḥ yoge pañcamī . \\
\noindent \textbf{(P\_2,3.32)} enapā dvitīyā anyārādibhiryoge pañcamī . \\
\noindent \textbf{(P\_2,3.32)} pṛthagvinānānābhiḥ tṛtīyā anyatarasyām . \\
\noindent \textbf{(P\_2,3.32)} pañcamīgrahaṇam anuvartate anyārādibhiḥ yoge iti nivṛttam . \\
\noindent \textbf{(P\_2,3.32)} atha vā maṇḍūkagatayaḥ adhikārāḥ . \\
\noindent \textbf{(P\_2,3.32)} tat. yathā maṇḍūkāḥ utplutya utplutya gacchanti tadvat adhikārāḥ . \\
\noindent \textbf{(P\_2,3.32)} atha vā anyavacanāt cakārākaraṇāt prakṛtasya apavādaḥ vijñāyate yathā utsargeṇa prasaktasya . \\
\noindent \textbf{(P\_2,3.32)} anyasyā vibhakteḥ vacanāt cakārasya anukarṣaṇārthasya akaraṇāt prakṛtāyaḥ pañcamyāḥ dvitīyāṣaṣṭhyau bādhike bhaviṣyataḥ yathā utsargeṇa prasaktasya apavādaḥ bādhakaḥ bhavati . \\
\noindent \textbf{(P\_2,3.32)} atha vā vakṣyati etat . \\
\noindent \textbf{(P\_2,3.32)} anuvartante ca nāma vidhayaḥ . \\
\noindent \textbf{(P\_2,3.32)} na ca anuvartanāt eva bhavanti . \\
\noindent \textbf{(P\_2,3.32)} kim tarhi . \\
\noindent \textbf{(P\_2,3.32)} yatnāt bhavanti iti . \\
\noindent \textbf{(P\_2,3.35)} dūrāntikārthebhyaḥ pañcamīvidhāne tadyuktāt pañcamīpratiṣedhaḥ . \\
\noindent \textbf{(P\_2,3.35)} dūrāntikārthebhyaḥ pañcamīvidhāne tadyuktātpañcamyāḥ pratiṣedhaḥ vaktavyaḥ . \\
\noindent \textbf{(P\_2,3.35)} dūrād grāmasya . \\
\noindent \textbf{(P\_2,3.35)} na vā tatra api darśanāt apratiṣedhaḥ . \\
\noindent \textbf{(P\_2,3.35)} na vā tatra api darśanāt pañcamyāḥ pratiṣedhaḥ anarthakaḥ . \\
\noindent \textbf{(P\_2,3.35)} tatra api pañcamī dṛśyate . \\
\noindent \textbf{(P\_2,3.35)} dūrāt āvasathāt mūtram dūrāt pādāvasecanam dūrāt ca bhāvyam dasyubhyaḥ dūrāt ca kupitāt guroḥ . \\
\noindent \textbf{(P\_2,3.36)} saptamīvidhāne ktasya inviṣayasya karmaṇi upsaṅkhyānam . \\
\noindent \textbf{(P\_2,3.36)} saptamīvidhāne ktasya inviṣayasya karmaṇi upsaṅkhyānam vaktavyam . \\
\noindent \textbf{(P\_2,3.36)} adhītī vyākaraṇe . \\
\noindent \textbf{(P\_2,3.36)} parigaṇitī yājñikye . \\
\noindent \textbf{(P\_2,3.36)} āmnātī cchandasi . \\
\noindent \textbf{(P\_2,3.36)} sādhvasādhuprayoge ca . \\
\noindent \textbf{(P\_2,3.36)} sādhvasādhuprayoge ca saptamī vaktavyā . \\
\noindent \textbf{(P\_2,3.36)} sādhuḥ devadattaḥ mātari . \\
\noindent \textbf{(P\_2,3.36)} asādhuḥ pitari . \\
\noindent \textbf{(P\_2,3.36)} kārakārhāṇām ca kārakatve . \\
\noindent \textbf{(P\_2,3.36)} kārakārhāṇām ca kārakatve saptamī vaktavyā . \\
\noindent \textbf{(P\_2,3.36)} ṛddheṣu bhuñjāneṣu daridrāḥ āsate . \\
\noindent \textbf{(P\_2,3.36)} brāhmaṇeṣu taratsu vṛṣalāḥ āsate . \\
\noindent \textbf{(P\_2,3.36)} akārakārhāṇām cākārakatve . \\
\noindent \textbf{(P\_2,3.36)} akārakārhāṇām cākārakatve saptamī vaktavyā . \\
\noindent \textbf{(P\_2,3.36)} mūrkheṣu āsīneṣu vṛddhāḥ bhuñjate . \\
\noindent \textbf{(P\_2,3.36)} vṛṣaleṣu āsīneṣu brāhmaṇāḥ taranti . \\
\noindent \textbf{(P\_2,3.36)} tadviparyāse ca . \\
\noindent \textbf{(P\_2,3.36)} tadviparyāse ca saptamī vaktavyā . \\
\noindent \textbf{(P\_2,3.36)} ṛddheṣu āsīneṣu mūrkhāḥ bhuñjate . \\
\noindent \textbf{(P\_2,3.36)} brāhmaṇeṣu āsīneṣu vṛṣalāḥ taranti . \\
\noindent \textbf{(P\_2,3.36)} nimittāt karmasaṃyoge . \\
\noindent \textbf{(P\_2,3.36)} nimittātkarmasaṃyoge saptamī vaktavyā . \\
\noindent \textbf{(P\_2,3.36)} carmaṇi dvīpinam hanti . \\
\noindent \textbf{(P\_2,3.36)} dantayoḥ hanti kuñjaram . \\
\noindent \textbf{(P\_2,3.36)} keṣeṣu camarīm hanti . \\
\noindent \textbf{(P\_2,3.36)} sīmni puṣkalakaḥ hataḥ . \\
\noindent \textbf{(P\_2,3.37)} bhāvalakṣaṇe saptamīvidhāne abhāvalakṣaṇe upasaṅkhyānam . \\
\noindent \textbf{(P\_2,3.37)} bhāvalakṣaṇe saptamīvidhāne abhāvalakṣaṇe upasaṅkhyānam kartavyam . \\
\noindent \textbf{(P\_2,3.37)} agniṣu hūyamāneṣu prasthitaḥ huteṣu āgataḥ . \\
\noindent \textbf{(P\_2,3.37)} goṣu duhyamānāsu prasthitaḥ dugdhāsu āgataḥ . \\
\noindent \textbf{(P\_2,3.37)} kim punaḥ kāraṇam na sidhyati . \\
\noindent \textbf{(P\_2,3.37)} lakṣaṇam hi nāma tat bhavati yena punaḥ punaḥ lakṣyate . \\
\noindent \textbf{(P\_2,3.37)} sakṛt ca asau katham cit agniṣu hūyamāneṣu prasthitaḥ huteṣu āgataḥ goṣu duhyamānāsu prasthitaḥ dugdhāsu āgataḥ . \\
\noindent \textbf{(P\_2,3.37)} siddham tu bhāvapravṛttau yasya bhāvārambhavacanāt . \\
\noindent \textbf{(P\_2,3.37)} siddhametat . \\
\noindent \textbf{(P\_2,3.37)} katham . \\
\noindent \textbf{(P\_2,3.37)} yasya bhāvapravṛttau dvitīyaḥ bhāvaḥ ārabhyate tatra saptamī vaktavyā . \\
\noindent \textbf{(P\_2,3.37)} sidhyati . \\
\noindent \textbf{(P\_2,3.37)} sūtram tarhi bhidyate . \\
\noindent \textbf{(P\_2,3.37)} yathānyāsam eva astu . \\
\noindent \textbf{(P\_2,3.37)} nanu ca uktam bhāvalakṣaṇe saptamīvidhāne abhāvalakṣaṇe upasaṅkhyānam iti . \\
\noindent \textbf{(P\_2,3.37)} na eṣaḥ doṣaḥ . \\
\noindent \textbf{(P\_2,3.37)} na khalu avaśyam tat eva lakṣaṇam bhavati yena punaḥ punaḥ lakṣyate . \\
\noindent \textbf{(P\_2,3.37)} sakṛt api yat nimittatvāya kalpate tat api lakṣaṇam bhavati . \\
\noindent \textbf{(P\_2,3.37)} tat yathā . \\
\noindent \textbf{(P\_2,3.37)} api bhavān kamaṇḍalupāṇim chātram adrākṣīt iti . \\
\noindent \textbf{(P\_2,3.37)} sakṛt asau kamaṇḍalupāṇiḥ dṛṣṭaḥ chātraḥ . \\
\noindent \textbf{(P\_2,3.37)} tasya tat eva lakṣaṇam bhavati . \\
\noindent \textbf{(P\_2,3.42)} iha kasmāt na bhavati . \\
\noindent \textbf{(P\_2,3.42)} kṛṣṇā gavām sampannakṣīratamā iti . \\
\noindent \textbf{(P\_2,3.42)} vibhakte iti ucyate . \\
\noindent \textbf{(P\_2,3.42)} na ca etat vibhaktam . \\
\noindent \textbf{(P\_2,3.42)} vibhaktametat . \\
\noindent \textbf{(P\_2,3.42)} gobhyaḥ kṛṣṇā vibhajyate . \\
\noindent \textbf{(P\_2,3.42)} vibhaktam eva yat nityam tatra bhavitavyam . \\
\noindent \textbf{(P\_2,3.42)} na ca etat nityam vibhaktam . \\
\noindent \textbf{(P\_2,3.42)} kim vaktavyam etat . \\
\noindent \textbf{(P\_2,3.42)} na hi . \\
\noindent \textbf{(P\_2,3.42)} katham anucyamānam gaṃsyate . \\
\noindent \textbf{(P\_2,3.42)} vibhaktagrahaṇasāmarthyāt . \\
\noindent \textbf{(P\_2,3.42)} yadi hi yat vibhaktam ca avibhaktam ca tatra syāt vibhāktagrahaṇam anarthakam syāt . \\
\noindent \textbf{(P\_2,3.43)} apratyādibhiḥ iti vaktavyam . \\
\noindent \textbf{(P\_2,3.43)} iha api yathā syāt . \\
\noindent \textbf{(P\_2,3.43)} sādhuḥ devadattaḥ mātaram pari . \\
\noindent \textbf{(P\_2,3.43)} mātaram anu . \\
\noindent \textbf{(P\_2,3.44)} prasitaḥ iti ucyate kaḥ prasitaḥ nāma . \\
\noindent \textbf{(P\_2,3.44)} yaḥ tatra nityam pratibaddhaḥ . \\
\noindent \textbf{(P\_2,3.44)} kutaḥ etat . \\
\noindent \textbf{(P\_2,3.44)} sinotiḥ ayam badhnātyarthe vartate . \\
\noindent \textbf{(P\_2,3.44)} baddhaḥ iva asau tatra bhavati . \\
\noindent \textbf{(P\_2,3.45)} iha kasmāt na bhavati . \\
\noindent \textbf{(P\_2,3.45)} adya puṣyaḥ . \\
\noindent \textbf{(P\_2,3.45)} adya maghā iti . \\
\noindent \textbf{(P\_2,3.45)} adhikaraṇe iti vartate . \\
\noindent \textbf{(P\_2,3.46.1)} prātipadikagrahaṇam kimartham . \\
\noindent \textbf{(P\_2,3.46.1)} uccaiḥ nīcaiḥ iti āpi yathā syāt . \\
\noindent \textbf{(P\_2,3.46.1)} kim punaḥ atra prathamayā prārthyate . \\
\noindent \textbf{(P\_2,3.46.1)} padatvam . \\
\noindent \textbf{(P\_2,3.46.1)} na etat asti .ṣaṣṭhyā atra padatvam bhaviṣyati . \\
\noindent \textbf{(P\_2,3.46.1)} idam tarhi prayojanam . \\
\noindent \textbf{(P\_2,3.46.1)} grāmaḥ ucaiḥ te svam . \\
\noindent \textbf{(P\_2,3.46.1)} grāmaḥ uccaiḥ tava svam . \\
\noindent \textbf{(P\_2,3.46.1)} sapūrvāyāḥ prathamāyāḥ vibhāṣā iti eṣaḥ vidhiḥ yathā syāt . \\
\noindent \textbf{(P\_2,3.46.1)} atha liṅgagrahaṇam kimartham . \\
\noindent \textbf{(P\_2,3.46.1)} strī pumān napuṃsakam iti ata api yathā syāt . \\
\noindent \textbf{(P\_2,3.46.1)} na etat asti prayojanam . \\
\noindent \textbf{(P\_2,3.46.1)} eṣaḥ eva atra prātipadikārthaḥ . \\
\noindent \textbf{(P\_2,3.46.1)} idam tarhi . \\
\noindent \textbf{(P\_2,3.46.1)} kumārī vṛkṣaḥ kuṇḍam iti . \\
\noindent \textbf{(P\_2,3.46.1)} atha parimāṇagrahaṇam kimartham . \\
\noindent \textbf{(P\_2,3.46.1)} droṇaḥ khārī āḍhakam iti atra api yathā syāt . \\
\noindent \textbf{(P\_2,3.46.1)} atha vacanagrahaṇam kimartham . \\
\noindent \textbf{(P\_2,3.46.1)} iha samudāye vākyaparisamāptiḥ dṛśyate . \\
\noindent \textbf{(P\_2,3.46.1)} tat yathā . \\
\noindent \textbf{(P\_2,3.46.1)} gargāḥ śatam daṇḍyantām iti . \\
\noindent \textbf{(P\_2,3.46.1)} arthinaḥ ca rājānaḥ hiraṇyena bhavanti na ca pratyekam daṇḍayanti . \\
\noindent \textbf{(P\_2,3.46.1)} sati etasmin dṛṣṭānte yatra etāni samuditāni bhavanti tatra eva syāt . \\
\noindent \textbf{(P\_2,3.46.1)} droṇaḥ khārī āḍhakam iti . \\
\noindent \textbf{(P\_2,3.46.1)} iha na syāt . \\
\noindent \textbf{(P\_2,3.46.1)} kumārī vṛkṣaḥ kuṇḍam iti . \\
\noindent \textbf{(P\_2,3.46.1)} na etat asti prayojanam . \\
\noindent \textbf{(P\_2,3.46.1)} pratyekam api vākyaparisamāptiḥ dṛśyate . \\
\noindent \textbf{(P\_2,3.46.1)} tat yathā vṛddhiguṇasañjñe pratyekam bhavataḥ . \\
\noindent \textbf{(P\_2,3.46.1)} idam tarhi prayojanam ukteṣu api ekatvādiṣu prathamā yathā syāt . \\
\noindent \textbf{(P\_2,3.46.1)} ekaḥ dvau bahavaḥ iti . \\
\noindent \textbf{(P\_2,3.46.1)} atha mātragrahaṇam kimartham . \\
\noindent \textbf{(P\_2,3.46.1)} etanmātre eva prathamā yathā syāt karmādiviśiṣṭe mā bhūt iti . \\
\noindent \textbf{(P\_2,3.46.1)} kaṭam karoti . \\
\noindent \textbf{(P\_2,3.46.1)} na etat asti prayojanam . \\
\noindent \textbf{(P\_2,3.46.1)} karmādiṣu dvitīyādyāḥ vibhaktayaḥ tāḥ karmādiviśiṣṭe bādhikāḥ bhaviṣyanti . \\
\noindent \textbf{(P\_2,3.46.1)} atha vā ācāryapravṛttiḥ jñāpayati na karmādiviśiṣṭe prathamā bhavati iti yat ayam sambodhane prathamām śāsti . \\
\noindent \textbf{(P\_2,3.46.1)} na etat asti jñāpakam . \\
\noindent \textbf{(P\_2,3.46.1)} asti hi anyat etasya vacane prayojanam . \\
\noindent \textbf{(P\_2,3.46.1)} kim . \\
\noindent \textbf{(P\_2,3.46.1)} sā āmantritam iti vakṣyāmi iti . \\
\noindent \textbf{(P\_2,3.46.1)} yat tarhi yogavibhāgam karoti . \\
\noindent \textbf{(P\_2,3.46.1)} itarathā hi sambodhane āmantritam iti eva brūyāt . \\
\noindent \textbf{(P\_2,3.46.1)} idam tarhi ukteṣu api ekatvādiṣu prathamā yathā syāt . \\
\noindent \textbf{(P\_2,3.46.1)} ekaḥ dvau bahavaḥ iti . \\
\noindent \textbf{(P\_2,3.46.1)} vacanagrahaṇasya api etat prayojanam uktam . \\
\noindent \textbf{(P\_2,3.46.1)} anyatarat śakyam akartum . \\
\noindent \textbf{(P\_2,3.46.2)} prātipadikārthaliṅgaparimāṇavacanamātre prathamālakṣaṇe padasāmānādhikaraṇye upasaṅkhyānam adhikatvāt . \\
\noindent \textbf{(P\_2,3.46.2)} prātipadikārthaliṅgaparimāṇavacanamātre prathamālakṣaṇe padasāmānādhikaraṇye upasaṅkhyānam kartavyam . \\
\noindent \textbf{(P\_2,3.46.2)} vīraḥ puruṣaḥ . \\
\noindent \textbf{(P\_2,3.46.2)} kim punaḥ kāraṇam na sidhyati . \\
\noindent \textbf{(P\_2,3.46.2)} adhikatvāt . \\
\noindent \textbf{(P\_2,3.46.2)} vyatiriktaḥ prātipadikārthaḥ iti kṛtvā prathamā na prāpnoti . \\
\noindent \textbf{(P\_2,3.46.2)} katham vyatiriktiḥ . \\
\noindent \textbf{(P\_2,3.46.2)} puruṣe vīratvam . \\
\noindent \textbf{(P\_2,3.46.2)} na vā vākyārthatvāt . \\
\noindent \textbf{(P\_2,3.46.2)} na vā vaktavyam . \\
\noindent \textbf{(P\_2,3.46.2)} kim kāraṇam . \\
\noindent \textbf{(P\_2,3.46.2)} vākyārthatvāt . \\
\noindent \textbf{(P\_2,3.46.2)} yat atra ādhikyam vākyārthaḥ saḥ . \\
\noindent \textbf{(P\_2,3.46.2)} atha vā abhihite prathamā iti etat lakṣaṇam kariyṣyate . \\
\noindent \textbf{(P\_2,3.46.2)} abhihitalakṣaṇāyām anabhihite prathamāvidhiḥ . abhihitalakṣaṇāyāmanabhihite prathamā vidheyā . \\
\noindent \textbf{(P\_2,3.46.2)} vṛkṣaḥ plakṣaḥ iti . \\
\noindent \textbf{(P\_2,3.46.2)} uktam vā . kim uktam . \\
\noindent \textbf{(P\_2,3.46.2)} astiḥ bhavantīparaḥ prathamapuruṣaḥ aprayujyamānaḥ api asti iti . \\
\noindent \textbf{(P\_2,3.46.2)} vṛkṣaḥ plakṣaḥ . \\
\noindent \textbf{(P\_2,3.46.2)} asti iti gamyate . \\
\noindent \textbf{(P\_2,3.46.2)} abhihitānabhihite prathamābhāvaḥ . abhihitānabhhite prathamā prāpnoti . \\
\noindent \textbf{(P\_2,3.46.2)} kva . \\
\noindent \textbf{(P\_2,3.46.2)} prāsāde āste . \\
\noindent \textbf{(P\_2,3.46.2)} śayane āste . \\
\noindent \textbf{(P\_2,3.46.2)} sadipratyayena abhihitam adhikaraṇam iti kṛtvā prathamā prāpnoti . \\
\noindent \textbf{(P\_2,3.46.2)} evam tarhi tiṅsamānādhikaraṇe prathamā iti etat lakṣaṇaṃ kariṣyate . \\
\noindent \textbf{(P\_2,3.46.2)} tiṅsamānādhikaraṇe iti cet tiṅaḥ aprayoge prathamāvidhiḥ . tiṅsamānādhikaraṇe iti cet tiṅaḥ aprayoge prathamā vidheyā . \\
\noindent \textbf{(P\_2,3.46.2)} vṛkṣaḥ plakṣa iti . \\
\noindent \textbf{(P\_2,3.46.2)} uktam pūrveṇa . \\
\noindent \textbf{(P\_2,3.46.2)} kim uktam . \\
\noindent \textbf{(P\_2,3.46.2)} astiḥ bhavantīparaḥ prathamapuruṣaḥ aprayujyamānaḥ api asti iti . \\
\noindent \textbf{(P\_2,3.46.2)} vṛkṣaḥ plakṣaḥ . \\
\noindent \textbf{(P\_2,3.46.2)} asti iti gamyate . \\
\noindent \textbf{(P\_2,3.46.2)} śatṛśānacoḥ ca nimittabhāvāt tiṅaḥ abhāvaḥ tayoḥ apavādatvāt . \\
\noindent \textbf{(P\_2,3.46.2)} śatṛśānacoḥ ca nimittabhāvāt tiṅaḥ abhāvaḥ . \\
\noindent \textbf{(P\_2,3.46.2)} kva . \\
\noindent \textbf{(P\_2,3.46.2)} pacati odanam devadattaḥ iti . \\
\noindent \textbf{(P\_2,3.46.2)} kim kāraṇam . \\
\noindent \textbf{(P\_2,3.46.2)} tayoḥ apavādatvāt . \\
\noindent \textbf{(P\_2,3.46.2)} śatṛśānacau tiṅapavādau . \\
\noindent \textbf{(P\_2,3.46.2)} tau ca atra bādhakau . \\
\noindent \textbf{(P\_2,3.46.2)} na ca apavādaviṣaye utsargaḥ abhiniviśate . \\
\noindent \textbf{(P\_2,3.46.2)} pūrvam hi apavādāḥ abhiniviśante paścāt utsargaḥ . \\
\noindent \textbf{(P\_2,3.46.2)} prakalpya vā apavādaviṣayam tataḥ utsargaḥ abhiniviśate . \\
\noindent \textbf{(P\_2,3.46.2)} tat na tāvat atra kadā cit tiṅādeśo bhavati . \\
\noindent \textbf{(P\_2,3.46.2)} apavādau tāvat śatṛśānacau pratīkṣate . \\
\noindent \textbf{(P\_2,3.46.2)} pākṣikaḥ eṣaḥ doṣaḥ . \\
\noindent \textbf{(P\_2,3.46.2)} katarasmin pakṣe . \\
\noindent \textbf{(P\_2,3.46.2)} śatṛśānacoḥ dvaitam bhavati . \\
\noindent \textbf{(P\_2,3.46.2)} aprathamā vā vidhinā āśrīyate prathamā vā pratiṣedhena iti . \\
\noindent \textbf{(P\_2,3.46.2)} vibhaktiniyame ca api dvaitam bhavati . \\
\noindent \textbf{(P\_2,3.46.2)} vibhaktiniyamaḥ vā syāt arthaniyamaḥ vā iti . \\
\noindent \textbf{(P\_2,3.46.2)} tat yadā tāvat arthaniyamaḥ aprathamā ca vidhinā āśrīyate tadā eṣa doṣaḥ bhavati . \\
\noindent \textbf{(P\_2,3.46.2)} yadā hi vibhaktiniyamaḥ yadi eva aprathamā vidhinā āśrīyate atha api prathamā pratiṣedhena na tadā doṣaḥ bhavati \\
\noindent \textbf{(P\_2,3.50)} śeṣe iti ucyate. kaḥ śeṣaḥ nāma . \\
\noindent \textbf{(P\_2,3.50)} karmādibhyaḥ ye anye arthāḥ saḥ śeṣaḥ . \\
\noindent \textbf{(P\_2,3.50)} yadi evam śeṣaḥ na prakalpate . \\
\noindent \textbf{(P\_2,3.50)} na hi karmādibhyaḥ anye arthāḥ santi . \\
\noindent \textbf{(P\_2,3.50)} iha tāvat rājñaḥ puruṣaḥ iti rājā kartā puruṣaḥ sampradānam . \\
\noindent \textbf{(P\_2,3.50)} vṛkṣasya śākhā iti vṛkṣaḥ śākhyāyāḥ adhikaraṇam . \\
\noindent \textbf{(P\_2,3.50)} tathā yat etat svam nāma caturbhiḥ etat prakāraiḥ bhavati krayaṇāt apaharaṇāt yāñcāyāḥ vinimayāt iti . \\
\noindent \textbf{(P\_2,3.50)} atra ca sarvatra karmādayaḥ santi. evam tarhi karmādīnām avivakṣā śeṣaḥ . \\
\noindent \textbf{(P\_2,3.50)} katham punaḥ sataḥ nāma avāvivakṣā syāt . \\
\noindent \textbf{(P\_2,3.50)} sataḥ api avivakṣā bhavati . \\
\noindent \textbf{(P\_2,3.50)} tat yathā . \\
\noindent \textbf{(P\_2,3.50)} alomikā eḍakā . \\
\noindent \textbf{(P\_2,3.50)} anudarā kanyā iti . \\
\noindent \textbf{(P\_2,3.50)} asataḥ ca vivakṣā bhavati . \\
\noindent \textbf{(P\_2,3.50)} samudraḥ kuṇḍikā . \\
\noindent \textbf{(P\_2,3.50)} vindhyaḥ vardhitakam iti . \\
\noindent \textbf{(P\_2,3.50)} kimartham punaḥ śeṣagrahaṇam . \\
\noindent \textbf{(P\_2,3.50)} pratyayāvadhāraṇāt śeṣavacanam . pratyayāvadhāraṇāt śeṣavacanam kartavyam . \\
\noindent \textbf{(P\_2,3.50)} pratyayāḥ niyatāḥ arthāḥ aniyatāḥ tatra ṣaṣṭhī prāpnoti . \\
\noindent \textbf{(P\_2,3.50)} tatra śeṣagrahaṇam kartavyam ṣaṣṭhīniyamārtham . \\
\noindent \textbf{(P\_2,3.50)} śeṣe eva ṣaṣṭhī bhavati na anyatra iti . \\
\noindent \textbf{(P\_2,3.50)} arthāvadhāraṇāt vā . \\
\noindent \textbf{(P\_2,3.50)} atha vā arthāḥ niyatāḥ pratyayāḥ aniyatāḥ te śeṣe api prāpnuvanti . \\
\noindent \textbf{(P\_2,3.50)} tatra śeṣagrahaṇam kartavyam śeṣaniyamārtham . \\
\noindent \textbf{(P\_2,3.50)} śeṣe ṣaṣṭhī eva bhavati na anyā iti . \\
\noindent \textbf{(P\_2,3.50)} arthaniyame śeṣagrahaṇam śakyam akartum . \\
\noindent \textbf{(P\_2,3.50)} katham . \\
\noindent \textbf{(P\_2,3.50)} arthāḥ niyatāḥ pratyayāḥ aniyatāḥ . \\
\noindent \textbf{(P\_2,3.50)} tataḥ vakṣyāmi ṣaṣṭhī bhavati iti . \\
\noindent \textbf{(P\_2,3.50)} tat niyamārtham bhaviṣyati . \\
\noindent \textbf{(P\_2,3.50)} yatra ṣaṣṭhī ca anyā ca prāpnoti ṣaṣṭhī eva tatra bhavati iti . \\
\noindent \textbf{(P\_2,3.50)} ṣaṣṭhī śeṣe iti cet viśeṣyasya pratiṣedhaḥ . \\
\noindent \textbf{(P\_2,3.50)} ṣaṣṭhī śeṣe iti cet viśeṣyasya pratiṣedhaḥ vaktavyaḥ . \\
\noindent \textbf{(P\_2,3.50)} rājñaḥ puruṣaḥ iti atra rājā viśeṣaṇam puruṣaḥ viśeṣyaḥ . \\
\noindent \textbf{(P\_2,3.50)} tatra prātipadikārthaḥ vyatiriktaḥ iti kṛtvā prathamā na prāpnoti . \\
\noindent \textbf{(P\_2,3.50)} tatra ṣaṣṭhī syāt . \\
\noindent \textbf{(P\_2,3.50)} tasyāḥ pratiṣedhaḥ vaktavyaḥ . \\
\noindent \textbf{(P\_2,3.50)} tatra prathamāvidhiḥ . \\
\noindent \textbf{(P\_2,3.50)} tatra ṣaṣṭhīm pratiṣidhya prathamā vidheyā . \\
\noindent \textbf{(P\_2,3.50)} rājñaḥ puruṣaḥ iti . \\
\noindent \textbf{(P\_2,3.50)} uktam pūrveṇa . \\
\noindent \textbf{(P\_2,3.50)} kimuktam . \\
\noindent \textbf{(P\_2,3.50)} na vā vākyārthatvāt iti . \\
\noindent \textbf{(P\_2,3.50)} yadatrādikhyam vākyārthaḥ saḥ . \\
\noindent \textbf{(P\_2,3.50)} kutaḥ nu khalu etat puruṣe yat ādikhyam saḥ vākyārthaḥ iti na punaḥ rājani yat ādhikyam saḥ vākyārthaḥ syāt . \\
\noindent \textbf{(P\_2,3.50)} antareṇa api puruṣaśabdaprayogam rājani saḥ arthaḥ gamyate . \\
\noindent \textbf{(P\_2,3.50)} na punaḥ antareṇa rājaśabdaprayogam puruṣe saḥ arthaḥ gamyate . \\
\noindent \textbf{(P\_2,3.50)} asti kāraṇam yena etat evam bhavati . \\
\noindent \textbf{(P\_2,3.50)} kim kāraṇam . \\
\noindent \textbf{(P\_2,3.50)} rājaśabdāt hi bhavān ṣaṣṭhīm uccārayati. aṅga hi bhavān puruṣaśabdāt api uccārayatu gaṃsyate saḥ arthaḥ . \\
\noindent \textbf{(P\_2,3.50)} nanu ca na etena evam bhavitavyam . \\
\noindent \textbf{(P\_2,3.50)} na hi śabdakṛtena nāma arthena bhavitavyam . \\
\noindent \textbf{(P\_2,3.50)} arthakṛtena nāma śabdena bhavitavyam . \\
\noindent \textbf{(P\_2,3.50)} tat etat evam dṛśyatām : artharūpam eva etat evañjātīyakam yena atra antareṇa api puruṣaśabdaprayogam rājani saḥ arthaḥ gamyate . \\
\noindent \textbf{(P\_2,3.50)} kim punaḥ tat . \\
\noindent \textbf{(P\_2,3.50)} svāmitvam . \\
\noindent \textbf{(P\_2,3.50)} kiṅkṛtam punaḥ tat . \\
\noindent \textbf{(P\_2,3.50)} svakṛtam . \\
\noindent \textbf{(P\_2,3.50)} tat yathā : prātipadikārthānām kriyākṛtāḥ viśeṣāḥ upajāyante tatkṛtāḥ ca ākhyāḥ prādurbhavanti karma karaṇam apādānaṃ sampradānam adhikaraṇam iti . \\
\noindent \textbf{(P\_2,3.50)} tāḥ ca punaḥ vibhaktīnām utpattau kadā cit nimittatvena upādīyante kadā cit na . \\
\noindent \textbf{(P\_2,3.50)} kadā ca vibhaktīnām utpattau nimittatvena upādīyante . \\
\noindent \textbf{(P\_2,3.50)} yadā vyabhicaranti prātipadikārtham . \\
\noindent \textbf{(P\_2,3.50)} yadā hi na vyabhicaranti ākhyābhūtāḥ eva tadā bhavanti karma karaṇam apādānam sampradānam adhikaraṇam iti . \\
\noindent \textbf{(P\_2,3.50)} yathā eva tarhi rājani svakṛtam svāmitvam tatra ṣaṣṭhī evam puruṣe api svāmikṛtam svatvam . \\
\noindent \textbf{(P\_2,3.50)} tatra ṣaṣṭhī prāpnoti . \\
\noindent \textbf{(P\_2,3.50)} rājaśabdāt utpadyamānayā ṣaṣṭhyā abhihitaḥ saḥ arthaḥ iti kṛtvā puruṣaśabdāt ṣaṣṭhī na bhaviṣyati . \\
\noindent \textbf{(P\_2,3.50)} na tarhi idānīm idam bhavati puruṣasya rājā iti . \\
\noindent \textbf{(P\_2,3.50)} bhavati . \\
\noindent \textbf{(P\_2,3.50)} rājaśabdāt tu tadā prathamā . \\
\noindent \textbf{(P\_2,3.50)} na tarhi idānīm idam bhavati : rājñaḥ puruṣasya iti . \\
\noindent \textbf{(P\_2,3.50)} bhavati . \\
\noindent \textbf{(P\_2,3.50)} bāhyam artham abhisamīkṣya . \\
\noindent \textbf{(P\_2,3.52)} karmādiṣu akarmakavadvacanam . karmādiṣu akarmakavadbhāvaḥ vaktavyaḥ . \\
\noindent \textbf{(P\_2,3.52)} kim prayojanam . \\
\noindent \textbf{(P\_2,3.52)} akarmakāṇām bhāve laḥ bhavati . \\
\noindent \textbf{(P\_2,3.52)} bhāve laḥ yathā syāt . \\
\noindent \textbf{(P\_2,3.52)} mātuḥ smaryate . \\
\noindent \textbf{(P\_2,3.52)} pituḥ smaryate . \\
\noindent \textbf{(P\_2,3.52)} atha vatkaraṇam kimartham . \\
\noindent \textbf{(P\_2,3.52)} svāśrayam api yathā syāt . \\
\noindent \textbf{(P\_2,3.52)} mātā smaryate . \\
\noindent \textbf{(P\_2,3.52)} pitā smaryate iti . \\
\noindent \textbf{(P\_2,3.52)} karmābhidhāne hi liṅgavacanānupapattiḥ . \\
\noindent \textbf{(P\_2,3.52)} karmābhidhāne hi sati liṅgavacanayoḥ anupapattiḥ syāt . \\
\noindent \textbf{(P\_2,3.52)} mātuḥ smṛtam . \\
\noindent \textbf{(P\_2,3.52)} mātroḥ smṛtam . \\
\noindent \textbf{(P\_2,3.52)} māṭṝṇām smṛtam iti . \\
\noindent \textbf{(P\_2,3.52)} mātuḥ yat liṅgam vacanam ca tat smṛtaśabdasya api prāpnoti . \\
\noindent \textbf{(P\_2,3.52)} ṣaṣṭhīprasaṇgaḥ ca . \\
\noindent \textbf{(P\_2,3.52)} ṣaṣṭhī ca prāpnoti . \\
\noindent \textbf{(P\_2,3.52)} kutaḥ . \\
\noindent \textbf{(P\_2,3.52)} smṛtaśabdāt . \\
\noindent \textbf{(P\_2,3.52)} mātuḥ sāmānādhikaraṇyāt ṣaṣṭhī prāpnoti . \\
\noindent \textbf{(P\_2,3.52)} aparaḥ āha : ṣaṣṭhīprasaṅgaḥ ca . \\
\noindent \textbf{(P\_2,3.52)} ṣaṣṭhī ca prasaṅktavyā . \\
\noindent \textbf{(P\_2,3.52)} kutaḥ . \\
\noindent \textbf{(P\_2,3.52)} mātṛśabdāt . \\
\noindent \textbf{(P\_2,3.52)} smṛtaśabdena bhihitam karma iti kṛtvā ṣaṣṭhī na prāpnoti . \\
\noindent \textbf{(P\_2,3.52)} tat tarhi vaktavyam . \\
\noindent \textbf{(P\_2,3.52)} na vaktavyam . \\
\noindent \textbf{(P\_2,3.52)} avivakṣite karmaṇi ṣaṣṭhī bhavati . \\
\noindent \textbf{(P\_2,3.52)} kim vaktavyam etat . \\
\noindent \textbf{(P\_2,3.52)} na hi . \\
\noindent \textbf{(P\_2,3.52)} katham anucyamānam gaṃsyate . \\
\noindent \textbf{(P\_2,3.52)} śeṣe iti vartate . \\
\noindent \textbf{(P\_2,3.52)} śeṣaḥ ca kaḥ . \\
\noindent \textbf{(P\_2,3.52)} karmādīnām avivakṣā śeṣaḥ . \\
\noindent \textbf{(P\_2,3.52)} yadā karma vivakṣitam bhavati tadā ṣaṣṭhī na bhavati . \\
\noindent \textbf{(P\_2,3.52)} tat yathā . \\
\noindent \textbf{(P\_2,3.52)} smarāmi aham mātaram . \\
\noindent \textbf{(P\_2,3.52)} smarāmi aham pitaram iti . \\
\noindent \textbf{(P\_2,3.54)} ajvarisantāpyoḥ iti vaktavyam . \\
\noindent \textbf{(P\_2,3.54)} iha api yathā syāt . \\
\noindent \textbf{(P\_2,3.54)} cauram santāpayati . \\
\noindent \textbf{(P\_2,3.54)} vṛṣalam santāpayati . \\
\noindent \textbf{(P\_2,3.54)} atha kimartham bhāvavacanānām iti ucyate yāvatā rujārthāḥ bhāvavacanāḥ eva bhavanti . \\
\noindent \textbf{(P\_2,3.54)} bhāvakartṛkāt yathā syāt . \\
\noindent \textbf{(P\_2,3.54)} iha mā bhūt . \\
\noindent \textbf{(P\_2,3.54)} nadī kūlāni rujati iti . \\
\noindent \textbf{(P\_2,3.60)} kim udāharaṇam . \\
\noindent \textbf{(P\_2,3.60)} gām ghnanti . \\
\noindent \textbf{(P\_2,3.60)} gām pradīvyanti . \\
\noindent \textbf{(P\_2,3.60)} gām sabhāsadbhyaḥ upaharanti . \\
\noindent \textbf{(P\_2,3.60)} na etat asti . \\
\noindent \textbf{(P\_2,3.60)} pūrveṇa api etat siddham . \\
\noindent \textbf{(P\_2,3.60)} idam tarhi . \\
\noindent \textbf{(P\_2,3.60)} gāmasya tadahaḥ sabhāyām dīvyeyuḥ . \\
\noindent \textbf{(P\_2,3.61)} haviṣaḥ aprasthitasya . \\
\noindent \textbf{(P\_2,3.61)} haviṣaḥ aprasthitasya iti vaktavyam . \\
\noindent \textbf{(P\_2,3.61)} indrāgnibhyām chāgam haviḥ vapām medaḥ prasthitam preṣya . \\
\noindent \textbf{(P\_2,3.62)} ṣaṣṭhyarthe caturthīvacanam . \\
\noindent \textbf{(P\_2,3.62)} ṣaṣṭhyarthe caturthī vaktavyā . \\
\noindent \textbf{(P\_2,3.62)} yā kharveṇa pibati tasyai kharvaḥ tisraḥ rātrīḥ . \\
\noindent \textbf{(P\_2,3.62)} tasyāḥ iti prāpte . \\
\noindent \textbf{(P\_2,3.62)} yaḥ tataḥ jāyate saḥ bhiśastaḥ yām araṇye tasyai stenaḥ yām parācīm tasyai hrītamukhī apagagalbhaḥ yā snāti tasyai apsu mārukaḥ yā abhyaṅkte tasyai duścarmā yā pralikhate tasyai khalatiḥ apamārī yā āṅkte tasyai kāṇaḥ yā dataḥ dhāvate tasyai śyāvadan yā nakhani nikṛntate tasyai kunakhṝ yā kṛṇatti tasyai klībaḥ yā rajjum sṛjati tasyai udbandhukaḥ yā parṇena pibati tasyai unmādukaḥ jāyate . \\
\noindent \textbf{(P\_2,3.62)} ahalyāyai jāraḥ . \\
\noindent \textbf{(P\_2,3.62)} manāyyai tantuḥ . \\
\noindent \textbf{(P\_2,3.62)} tat tarhi vaktavyam . \\
\noindent \textbf{(P\_2,3.62)} na vaktavyam . \\
\noindent \textbf{(P\_2,3.62)} yogavibhāgāt siddham . \\
\noindent \textbf{(P\_2,3.62)} caturthī . \\
\noindent \textbf{(P\_2,3.62)} tataḥ arthe bahulam chandasi iti . \\
\noindent \textbf{(P\_2,3.65)} kṛdgrahaṇam kimartham . \\
\noindent \textbf{(P\_2,3.65)} iha mā bhūt . \\
\noindent \textbf{(P\_2,3.65)} pacati odanam devadattaḥ iti . \\
\noindent \textbf{(P\_2,3.65)} kartṛkarmaṇoḥ ṣaṣṭhīvidhāne kṛdgrahaṇānarthakyam lapratiṣedhāt . \\
\noindent \textbf{(P\_2,3.65)} kartṛkarmaṇoḥ ṣaṣṭhīvidhāne kṛdgrahaṇam anarthakam . \\
\noindent \textbf{(P\_2,3.65)} kim kāraṇam . \\
\noindent \textbf{(P\_2,3.65)} lapratiṣedhāt . \\
\noindent \textbf{(P\_2,3.65)} pratiṣidhyate tatra ṣaṣthī laprayoge na iti . \\
\noindent \textbf{(P\_2,3.65)} tasya karmakartrartham tarhi kṛdgrahaṇam kartavyam . \\
\noindent \textbf{(P\_2,3.65)} kṛtaḥ ye kartṛkarmaṇī tatra yathā syāt . \\
\noindent \textbf{(P\_2,3.65)} anyasya ye kartṛkarmaṇī tatra mā bhūt iti . \\
\noindent \textbf{(P\_2,3.65)} na etat asti prayojanam . \\
\noindent \textbf{(P\_2,3.65)} dhātoḥ hi dvaye pratyayāḥ vidhīyante tiṅaḥ ca kṛtaḥ ca . \\
\noindent \textbf{(P\_2,3.65)} tatra kṛtprayoge iṣyate tiṅprayoge pratiṣidhyate . \\
\noindent \textbf{(P\_2,3.65)} na brūmaḥ ihārtham tasya karmakartrartham kṛdgrahaṇam kartavyam iti . \\
\noindent \textbf{(P\_2,3.65)} kim tarhi . \\
\noindent \textbf{(P\_2,3.65)} uttarārtham . \\
\noindent \textbf{(P\_2,3.65)} avyayaprayoge na iti ṣaṣṭhyāḥ pratiṣedham vakṣyati . \\
\noindent \textbf{(P\_2,3.65)} saḥ kṛtaḥ avyayasya ye kartṛkarmaṇī tatra yathā syāt . \\
\noindent \textbf{(P\_2,3.65)} akṛtaḥ avyayasya ye kartṛkarmaṇī tatra mā bhūt iti . \\
\noindent \textbf{(P\_2,3.65)} uccaiḥ kaṭānām sraṣṭā iti . \\
\noindent \textbf{(P\_2,3.65)} tasya karmakartrartham iti cet pratiṣedhe api tadantakarmakartṛtvāt siddham . \\
\noindent \textbf{(P\_2,3.65)} kṛtaḥ ete kartṛkarmaṇī na avyayasya . \\
\noindent \textbf{(P\_2,3.65)} adhikaraṇam atra avyayam . \\
\noindent \textbf{(P\_2,3.65)} idam tarhi prayojanam . \\
\noindent \textbf{(P\_2,3.65)} ubhayaprāptau karmaṇi ṣaṣṭhyāḥ pratiṣedham vakṣyati . \\
\noindent \textbf{(P\_2,3.65)} saḥ kṛtaḥ ye kartṛkarmaṇī tatra yathā syāt . \\
\noindent \textbf{(P\_2,3.65)} kṛtoḥ ye kartṛkarmaṇī tatra mā bhūt iti . \\
\noindent \textbf{(P\_2,3.65)} āścaryam idam vṛttam odanasya ca nāma pākaḥ brāhmaṇānām ca prādurbhāvaḥ iti . \\
\noindent \textbf{(P\_2,3.65)} atha kriyamāṇe api kṛdgrahaṇe kasmāt eva atra na bhavati . \\
\noindent \textbf{(P\_2,3.65)} ubhayaprāptau iti na evam vijñāyate ubhayoḥ prāptiḥ ubhayaprāptiḥ ubhayaprāptau iti . \\
\noindent \textbf{(P\_2,3.65)} katham tarhi . \\
\noindent \textbf{(P\_2,3.65)} ubhayoḥ prāptiḥ yasmin kṛti saḥ ayam ubhayaprāptiḥ kṛt ubhayaprāptau iti . \\
\noindent \textbf{(P\_2,3.65)} atha vā kṛtaḥ ye kartṛkarmaṇī tatra yathā syāt . \\
\noindent \textbf{(P\_2,3.65)} taddhitasya ye kartṛkarmaṇī tatra mā bhūt iti . \\
\noindent \textbf{(P\_2,3.65)} kṛtapūrvī kaṭam . \\
\noindent \textbf{(P\_2,3.65)} bhuktapūrvī odanam iti . \\
\noindent \textbf{(P\_2,3.65)} nanu ca vākyena eva anena na bhavitavyam . \\
\noindent \textbf{(P\_2,3.65)} dvitīyayā tāvat na bhavitavyam . \\
\noindent \textbf{(P\_2,3.65)} kim kāraṇam . \\
\noindent \textbf{(P\_2,3.65)} ktena abhihitam karma iti kṛtvā . \\
\noindent \textbf{(P\_2,3.65)} inipratyayena ca api na utpattavyam . \\
\noindent \textbf{(P\_2,3.65)} kim kāraṇam . \\
\noindent \textbf{(P\_2,3.65)} asāmarthyāt . \\
\noindent \textbf{(P\_2,3.65)} katham asamārthyam . \\
\noindent \textbf{(P\_2,3.65)} sāpekṣam asamartham bhavati iti . \\
\noindent \textbf{(P\_2,3.65)} yat tāvat ucyate dvitīyayā tāvat na bhavitavyam . \\
\noindent \textbf{(P\_2,3.65)} kim kāraṇam . \\
\noindent \textbf{(P\_2,3.65)} ktena abhihitam karma iti kṛtvā iti . \\
\noindent \textbf{(P\_2,3.65)} yaḥ asau kṛtakaṭayoḥ abhisaṃbandhaḥ saḥ utpanne pratyaye nivartate . \\
\noindent \textbf{(P\_2,3.65)} asti ca karoteḥ kaṭena sāmarthyam iti kṛtvā dvitīyā bhaviṣyati . \\
\noindent \textbf{(P\_2,3.65)} yat api ucyate inipratyayena ca api na utpattavyam . \\
\noindent \textbf{(P\_2,3.65)} kim kāraṇam . \\
\noindent \textbf{(P\_2,3.65)} asāmarthyāt . \\
\noindent \textbf{(P\_2,3.65)} katham asamārthyam . \\
\noindent \textbf{(P\_2,3.65)} sāpekṣam asamartham bhavati iti . \\
\noindent \textbf{(P\_2,3.65)} na idam ubhayam yugapat bhavati vākyam ca pratyayaḥ ca . \\
\noindent \textbf{(P\_2,3.65)} yadā vākyam na tadā pratyayaḥ . \\
\noindent \textbf{(P\_2,3.65)} yadā pratyayaḥ sāmānyena tadā vṛttiḥ . \\
\noindent \textbf{(P\_2,3.65)} tatra avaśyaṃ viśeṣārthinā viśeṣaḥ anuprayoktavyaḥ . \\
\noindent \textbf{(P\_2,3.65)} kṛtapūrvī . \\
\noindent \textbf{(P\_2,3.65)} kim . \\
\noindent \textbf{(P\_2,3.65)} kaṭam . \\
\noindent \textbf{(P\_2,3.65)} bhuktapūrvī . \\
\noindent \textbf{(P\_2,3.65)} kim . \\
\noindent \textbf{(P\_2,3.65)} odanam iti . \\
\noindent \textbf{(P\_2,3.65)} atha vā idam prayojanam kartṛbhūtapūrvamātrāt api ṣaṣṭhīyathā syāt . \\
\noindent \textbf{(P\_2,3.65)} bhedikā devadattasya yajñadattasya kāṣṭhānām iti . \\
\noindent \textbf{(P\_2,3.66)} ubhayaprāptau karmaṇi ṣaṣṭhyāḥ pratiṣedhe akādiprayoge apratiṣedhaḥ . \\
\noindent \textbf{(P\_2,3.66)} ubhayaprāptau karmaṇi ṣaṣṭhyāḥ pratiṣedhe akādiprayoge pratiṣedhaḥ na bhavati iti vaktavyam . \\
\noindent \textbf{(P\_2,3.66)} bhedikā devadattasya kāṣṭhānām . \\
\noindent \textbf{(P\_2,3.66)} cikīrṣā viṣṇumitrasya kaṭasya . \\
\noindent \textbf{(P\_2,3.66)} aparaḥ āha : akākārayoḥ prayoge pratiṣedhaḥ na iti vaktavyam . \\
\noindent \textbf{(P\_2,3.66)} śeṣe vibhāṣā . \\
\noindent \textbf{(P\_2,3.66)} śobhanā khalu pāṇineḥ sūtrasya kṛtiḥ . \\
\noindent \textbf{(P\_2,3.66)} śobhanā khalu pāṇininā sūtrasya kṛtiḥ . \\
\noindent \textbf{(P\_2,3.66)} śobhanā khalu dākṣāyaṇasya saṅgrahasya kṛtiḥ . \\
\noindent \textbf{(P\_2,3.66)} śobhanā khalu dākṣāyeṇa saṅgrahasya kṛtiḥ iti . \\
\noindent \textbf{(P\_2,3.67)} ktasya ca vartamāne nāpuṃsake bhāve upasaṅkhyānam . \\
\noindent \textbf{(P\_2,3.67)} ktasya ca vartamāne nāpuṃsake bhāve upasaṅkhyānam kartavyam : chāttrasya hasitam , naṭasya bhuktam , mayūrasya nṛttam , kokilasya vyāhṛtam iti . \\
\noindent \textbf{(P\_2,3.67)} śeṣavijñānāt siddham . śeṣalakṣaṇā atra ṣaṣṭhī bhaviṣyati . \\
\noindent \textbf{(P\_2,3.67)} śeṣaḥ iti ucyate . \\
\noindent \textbf{(P\_2,3.67)} kaḥ ca śeṣaḥ . \\
\noindent \textbf{(P\_2,3.67)} karmādīnām avivakṣā śeṣaḥ . \\
\noindent \textbf{(P\_2,3.67)} katham punaḥ sataḥ nāma avivakṣā syāt yadā chātraḥ hasati , naṭaḥ bhuṅkte , mayūraḥ nṛtyati , kokilaḥ vyāharati . \\
\noindent \textbf{(P\_2,3.67)} sataḥ api avivakṣā bhavati . \\
\noindent \textbf{(P\_2,3.67)} tat yathā : alomikā eḍakā , anudarā kanyā iti . \\
\noindent \textbf{(P\_2,3.67)} asataḥ ca vivakṣā bhavati . \\
\noindent \textbf{(P\_2,3.67)} samudraḥ kuṇḍikā . \\
\noindent \textbf{(P\_2,3.67)} vindhyaḥ vardhitakam iti . \\
\noindent \textbf{(P\_2,3.67)} yadi evam uttaratra cātuḥśabdyam prāpnoti . \\
\noindent \textbf{(P\_2,3.67)} idam aheḥ sṛptam , iha ahinā sṛptam , iha ahiḥ sṛptaḥ , iha aheḥ sṛptam , grāmasya pārśve grāmasya madhye iti . \\
\noindent \textbf{(P\_2,3.67)} iṣyate eva cātuḥśabdyam . \\
\noindent \textbf{(P\_2,3.69)} lādeśe salliḍgrahaṇam kikinoḥ pratiṣedhārtham . \\
\noindent \textbf{(P\_2,3.69)} lādeśe salliḍgagrahaṇam kartavyam . \\
\noindent \textbf{(P\_2,3.69)} salliṭoḥ prayoge na iti vaktavyam . \\
\noindent \textbf{(P\_2,3.69)} kim prayojanam . \\
\noindent \textbf{(P\_2,3.69)} kikinoḥ pratiṣedhārtham . \\
\noindent \textbf{(P\_2,3.69)} kikinoḥ api prayoge pratiṣedhaḥ yathā syāt . \\
\noindent \textbf{(P\_2,3.69)} | papiḥ somaṃ dadiḥ gāḥ . \\
\noindent \textbf{(P\_2,3.69)} kim punaḥ kāraṇam na sidhyati . \\
\noindent \textbf{(P\_2,3.69)} tayoḥ alādeśatvāt . \\
\noindent \textbf{(P\_2,3.69)} na hi tau lādeśau . \\
\noindent \textbf{(P\_2,3.69)} atha tau lādeśau syātām syāt pratiṣedhaḥ . \\
\noindent \textbf{(P\_2,3.69)} bāḍham syāt . \\
\noindent \textbf{(P\_2,3.69)} lādeśau tarhi bhaviṣyataḥ . \\
\noindent \textbf{(P\_2,3.69)} tat katham . \\
\noindent \textbf{(P\_2,3.69)} ādṛgamahanajanaḥ kikinau liṭ ca iti liḍvat iti vakṣyāmi . \\
\noindent \textbf{(P\_2,3.69)} saḥ tarhi vatinirdeśaḥ kartavyaḥ . \\
\noindent \textbf{(P\_2,3.69)} na hi antareṇa vatim atideśaḥ gamyate . \\
\noindent \textbf{(P\_2,3.69)} antareṇa api vatim atideśaḥ gamyate . \\
\noindent \textbf{(P\_2,3.69)} tat yathā . \\
\noindent \textbf{(P\_2,3.69)} eṣaḥ brahmadattaḥ . \\
\noindent \textbf{(P\_2,3.69)} abrahmadattam brahmadattaḥ iti āha . \\
\noindent \textbf{(P\_2,3.69)} te manyāmahe : brahmadattavat ayam bhavati iti . \\
\noindent \textbf{(P\_2,3.69)} evam iha api aliṭam liṭ iti āha . \\
\noindent \textbf{(P\_2,3.69)} liḍvat iti vijñāsyate . \\
\noindent \textbf{(P\_2,3.69)} ukāraprayoge na iti vaktavyam . \\
\noindent \textbf{(P\_2,3.69)} kaṭam cikīrṣuḥ . \\
\noindent \textbf{(P\_2,3.69)} odanam bubhukṣuḥ . \\
\noindent \textbf{(P\_2,3.69)} tat tarhi vaktavyam . \\
\noindent \textbf{(P\_2,3.69)} na vaktavyam . \\
\noindent \textbf{(P\_2,3.69)} ukāraḥ api atra nirdiśyate . \\
\noindent \textbf{(P\_2,3.69)} katham . \\
\noindent \textbf{(P\_2,3.69)} praśliṣṭanirdeśaḥ ayam . \\
\noindent \textbf{(P\_2,3.69)} u uka ūka la ūka loka iti . \\
\noindent \textbf{(P\_2,3.69)} ukapratiṣedhe kameḥ bhāṣāyām apratiṣedhaḥ . \\
\noindent \textbf{(P\_2,3.69)} ukapratiṣedhe kameḥ bhāṣāyām pratiṣedhaḥ na bhavati iti vaktavyam . \\
\noindent \textbf{(P\_2,3.69)} dasyāḥ kāmukaḥ . \\
\noindent \textbf{(P\_2,3.69)} vṛṣalyāḥ kāmukaḥ . \\
\noindent \textbf{(P\_2,3.69)} avyayapratiṣedhe tosunkasunoḥ apratiṣedhaḥ . avyayapratiṣedhe tosunkasunoḥ pratiṣedhaḥ na bhavati iti vaktavyam . \\
\noindent \textbf{(P\_2,3.69)} purā sūryasya udetoḥ ādheyaḥ . \\
\noindent \textbf{(P\_2,3.69)} purā vatsānām apākartoḥ . \\
\noindent \textbf{(P\_2,3.69)} purā krūrasya visṛpaḥ virapśin . \\
\noindent \textbf{(P\_2,3.69)} śānaṃścānaśśatṝṛṇām upasaṅkhyānam . \\
\noindent \textbf{(P\_2,3.69)} śānaṃścānaśśatṝṛṇām upasaṅkhyānam kartavyam . \\
\noindent \textbf{(P\_2,3.69)} somam pavamānaḥ . \\
\noindent \textbf{(P\_2,3.69)} naḍam āghnānaḥ . \\
\noindent \textbf{(P\_2,3.69)} adhīyan pārāyaṇam . \\
\noindent \textbf{(P\_2,3.69)} laprayoge na iti pratiṣedhaḥ na prāpnoti . \\
\noindent \textbf{(P\_2,3.69)} mā bhūt evam . \\
\noindent \textbf{(P\_2,3.69)} tṛn iti evam bhaviṣyati . \\
\noindent \textbf{(P\_2,3.69)} katham . \\
\noindent \textbf{(P\_2,3.69)} tṛn iti na idam pratyayagrahaṇam . \\
\noindent \textbf{(P\_2,3.69)} kim tarhi . \\
\noindent \textbf{(P\_2,3.69)} pratyāhāragrahaṇam . \\
\noindent \textbf{(P\_2,3.69)} kva saṃniviṣṭānām pratyāhāraḥ . \\
\noindent \textbf{(P\_2,3.69)} laṭaḥ śatṛ iti ataḥ prabhṛti ā tṛnaḥ nakārāt . \\
\noindent \textbf{(P\_2,3.69)} yadi pratyāhāragrahaṇam caurasya dviṣan vṛṣalasya dviṣan atra api prāpnoti . \\
\noindent \textbf{(P\_2,3.69)} dviṣaḥ śatuḥ vāvacanam . dviṣaḥ śatuḥ vā iti vaktavyam . \\
\noindent \textbf{(P\_2,3.69)} tat ca avaśyaṃ vaktavyam pratyayagrhaṇe sati pratiṣedhārtham . \\
\noindent \textbf{(P\_2,3.69)} tat eva pratyāhāragrahaṇe sati vidhyartham bhaviṣyati . \\
\noindent \textbf{(P\_2,3.70)} akasya bhaviṣyati . \\
\noindent \textbf{(P\_2,3.70)} akasya bhaviṣyati iti vaktavyam . \\
\noindent \textbf{(P\_2,3.70)} yavān lāvakaḥ vrajati . \\
\noindent \textbf{(P\_2,3.70)} odanam bhojakaḥ vrajati . \\
\noindent \textbf{(P\_2,3.70)} saktūnpāyakaḥ vrajati . \\
\noindent \textbf{(P\_2,3.70)} inaḥ ādhamarṇye ca . tataḥ inaḥ ādhamarṇye ca bhaviṣyati ca iti vaktavyam . \\
\noindent \textbf{(P\_2,3.70)} śatam dāyī . \\
\noindent \textbf{(P\_2,3.70)} sahasram dāyī . \\
\noindent \textbf{(P\_2,3.70)} grāmam gāmī . \\
\noindent \textbf{(P\_2,3.71)} kartṛgrahaṇam kimartham . \\
\noindent \textbf{(P\_2,3.71)} karmaṇi mā bhūt iti . \\
\noindent \textbf{(P\_2,3.71)} na etat asti prayojanam . \\
\noindent \textbf{(P\_2,3.71)} bhāvakarmaṇoḥ kṛtyāḥ vidhīyante ṭatra kṛtyaiḥ abhihitatvāt karmaṇi ṣaṣṭhī na bhaviṣyati . \\
\noindent \textbf{(P\_2,3.71)} ataḥ uttaram paṭhati . \\
\noindent \textbf{(P\_2,3.71)} bhavyādīnām karmaṇaḥ anabhidhānāt kṛtyānām kartṛgrahaṇam . \\
\noindent \textbf{(P\_2,3.71)} bhavyādīnām karma kṛtyaiḥ anabhitam . \\
\noindent \textbf{(P\_2,3.71)} geyaḥ māṇavakaḥ sāmnām . \\
\noindent \textbf{(P\_2,3.71)} bhavyādīnam karmaṇaḥ anabhidhānāt kṛtyānām kartṛgrahaṇam kriyate . \\
\noindent \textbf{(P\_2,3.71)} kim ucyate bhavyādīnām karma kṛtyaiḥ anabhitam iti . \\
\noindent \textbf{(P\_2,3.71)} na iha api anabhihitaṃ bhavati . \\
\noindent \textbf{(P\_2,3.71)} ākraṣṭavyā grāṃam śākhā iti . \\
\noindent \textbf{(P\_2,3.71)} evam tarhi yogavibhāgaḥ kariṣyate . \\
\noindent \textbf{(P\_2,3.71)} kṛtyānām . \\
\noindent \textbf{(P\_2,3.71)} kṛtyānām prayoge ṣaṣṭhī na bhavati iti . \\
\noindent \textbf{(P\_2,3.71)} kim udāharaṇam . \\
\noindent \textbf{(P\_2,3.71)} grāmam ākraṣṭavyā śākhā . \\
\noindent \textbf{(P\_2,3.71)} tataḥ kartari vā iti . \\
\noindent \textbf{(P\_2,3.71)} iha api tarhi prāpnoti . \\
\noindent \textbf{(P\_2,3.71)} geyaḥ māṇavakaḥ sāmnām iti . \\
\noindent \textbf{(P\_2,3.71)} ubhayaprāptau iti vartate . \\
\noindent \textbf{(P\_2,3.71)} nanu ca ubhayaprāptiḥ eṣā . \\
\noindent \textbf{(P\_2,3.71)} geyaḥ māṇavakaḥ sāmnām iti ca geyāni māṇavakena sāmāni iti ca bhavati . \\
\noindent \textbf{(P\_2,3.71)} ubhayaprāptiḥ nāma sā bhavati yatra ubhayasya yugapatprasaṅgaḥ atra ca yadā karmaṇi na tadā kartari yadā kartari na tadā karmaṇi iti. \\
\noindent \textbf{(P\_2,4.1)} kimartham idam ucyate . \\
\noindent \textbf{(P\_2,4.1)} pratyadhikaraṇam vacanotpatteḥ saṅkhyāsāmānādhikaraṇyāt ca dvigoḥ ekavacanavidhānam . \\
\noindent \textbf{(P\_2,4.1)} pratyadhikaraṇam vacanotpattiḥ bhavati . \\
\noindent \textbf{(P\_2,4.1)} kim idam pratyadhikaraṇam iti . \\
\noindent \textbf{(P\_2,4.1)} adhikaraṇam adhikaraṇam prati pratyadhikaraṇam . \\
\noindent \textbf{(P\_2,4.1)} saṅkhyāsāmānādhikaraṇyāt ca . \\
\noindent \textbf{(P\_2,4.1)} saṅkhyayā bahvarthayā ca asya sāmānādhikaraṇyam . \\
\noindent \textbf{(P\_2,4.1)} pratyadhikaraṇam vacanotpatteḥ saṅkhyāsāmānādhikaraṇyāt ca bahuṣu bahuvacanam iti bahuvacanam prāpnoti . \\
\noindent \textbf{(P\_2,4.1)} iṣyate ca ekavacanam syāt iti . \\
\noindent \textbf{(P\_2,4.1)} tat ca antareṇa yatnam na sidhyati iti dvigoḥ ekavacanavidhānam . \\
\noindent \textbf{(P\_2,4.1)} evamartham idam ucyate . \\
\noindent \textbf{(P\_2,4.1)} asti prayojanam etat . \\
\noindent \textbf{(P\_2,4.1)} kim tarhi iti . \\
\noindent \textbf{(P\_2,4.1)} tatra anuprayogasya ekavacanābhāvaḥ advigutvāt . \\
\noindent \textbf{(P\_2,4.1)} tatra anuprayogasya ekavacanam na prāpnoti : pañcapūlī iyam iti . \\
\noindent \textbf{(P\_2,4.1)} kim kāraṇam . \\
\noindent \textbf{(P\_2,4.1)} advigutvāt . \\
\noindent \textbf{(P\_2,4.1)} dvigoḥ ekavacanm iti ucyate . \\
\noindent \textbf{(P\_2,4.1)} na ca atra anuprayogaḥ dvigusañjñaḥ . \\
\noindent \textbf{(P\_2,4.1)} siddham tu dvigvarthasya ekavadvacanāt . \\
\noindent \textbf{(P\_2,4.1)} siddham etat . \\
\noindent \textbf{(P\_2,4.1)} katham . \\
\noindent \textbf{(P\_2,4.1)} dvigvarthaḥ ekavat bhavati iti vaktavyam . \\
\noindent \textbf{(P\_2,4.1)} tat tarhi vaktavyam . \\
\noindent \textbf{(P\_2,4.1)} na vaktavyam . \\
\noindent \textbf{(P\_2,4.1)} na idam pāṛibhāṣikasya vacanasya grahaṇam . \\
\noindent \textbf{(P\_2,4.1)} kim tarhi . \\
\noindent \textbf{(P\_2,4.1)} anvarthagrahaṇam . \\
\noindent \textbf{(P\_2,4.1)} ucyate vacanam . \\
\noindent \textbf{(P\_2,4.1)} ekasya arthasya vacanam ekavacanam . \\
\noindent \textbf{(P\_2,4.1)} ekaśeṣapratiṣedhaḥ ca . \\
\noindent \textbf{(P\_2,4.1)} ekaśeṣasya ca pratiṣedhaḥ vaktavyaḥ . \\
\noindent \textbf{(P\_2,4.1)} pañcapūlī ca pañcapūlī ca pañcapūlī ca pañcapūlyaḥ . \\
\noindent \textbf{(P\_2,4.1)} na vā anyasya anekatvāt . \\
\noindent \textbf{(P\_2,4.1)} na vā vaktavyaḥ . \\
\noindent \textbf{(P\_2,4.1)} kim kāraṇam . \\
\noindent \textbf{(P\_2,4.1)} anyasya anekatvāt . \\
\noindent \textbf{(P\_2,4.1)} na etat dvigoḥ anekatvam . \\
\noindent \textbf{(P\_2,4.1)} kasya tarhi . \\
\noindent \textbf{(P\_2,4.1)} dvigvarthasamudāyasya . \\
\noindent \textbf{(P\_2,4.1)} samāhāragrahaṇam ca taddhitārthapratiṣedhārtham . \\
\noindent \textbf{(P\_2,4.1)} samāhāragrahaṇam ca kartavyam . \\
\noindent \textbf{(P\_2,4.1)} kim prayojanam . \\
\noindent \textbf{(P\_2,4.1)} taddhitārthapratiṣedhārtham . \\
\noindent \textbf{(P\_2,4.1)} taddhitārthe yaḥ dviguḥ tasya mā bhūt iti . \\
\noindent \textbf{(P\_2,4.1)} pañcakapālau pañcakapālāḥ iti . \\
\noindent \textbf{(P\_2,4.1)} kim punaḥ ayam pañcakapālaśabdaḥ pratyekam parisamāpyate āhosvit samudāye vartate . \\
\noindent \textbf{(P\_2,4.1)} kim ca ataḥ . \\
\noindent \textbf{(P\_2,4.1)} yadi pratyekam parisamāpyate purastāt eva coditam parihṛtam ca. atha samudāye vartate . \\
\noindent \textbf{(P\_2,4.1)} na vā samāhāraikatvāt . \\
\noindent \textbf{(P\_2,4.1)} na vā etat samāhāraikatvāt api sidhyati . \\
\noindent \textbf{(P\_2,4.1)} evam tarhi pratyekam parisamāpyate . \\
\noindent \textbf{(P\_2,4.1)} purastāt eva coditam parihṛtam ca. aparaḥ āha : na vā samāhāraikatvāt . \\
\noindent \textbf{(P\_2,4.1)} na vā yogārambheṇa eva arthaḥ . \\
\noindent \textbf{(P\_2,4.1)} kim kāraṇam . \\
\noindent \textbf{(P\_2,4.1)} samāhāraikatvāt . \\
\noindent \textbf{(P\_2,4.1)} ekaḥ ayam samāhāraḥ nāma . \\
\noindent \textbf{(P\_2,4.1)} tasya ekatvāt ekavacanam bhaviṣyati . \\
\noindent \textbf{(P\_2,4.2)} prāṇitūryasenāṅgānām tatpūrvapadottarapadagrahaṇam . \\
\noindent \textbf{(P\_2,4.2)} prāṇitūryasenāṅgānām tatpūrvapadottarapadagrahaṇam kartavyam . \\
\noindent \textbf{(P\_2,4.2)} prāṇyaṅgānām prāṇyaṅaiḥ iti vaktavyam . \\
\noindent \textbf{(P\_2,4.2)} tūryāṅgāṇāṃ tūryāṅgaiḥ . \\
\noindent \textbf{(P\_2,4.2)} senāṅgānām senāṅgaiḥ iti . \\
\noindent \textbf{(P\_2,4.2)} kim prayojanam . \\
\noindent \textbf{(P\_2,4.2)} vyatikaraḥ mā bhūt iti . \\
\noindent \textbf{(P\_2,4.2)} tat tarhi vaktavyam . \\
\noindent \textbf{(P\_2,4.2)} yogavibhāgāt siddham . \\
\noindent \textbf{(P\_2,4.2)} yogavibhāgaḥ kariṣyate . \\
\noindent \textbf{(P\_2,4.2)} dvandvaḥ ca prāṇyaṅgānām . \\
\noindent \textbf{(P\_2,4.2)} tataḥ tūryāṅgāṇām . \\
\noindent \textbf{(P\_2,4.2)} tataḥ senāṅgānām iti . \\
\noindent \textbf{(P\_2,4.2)} saḥ tarhi yogavibhāgaḥ kartavyaḥ . \\
\noindent \textbf{(P\_2,4.2)} na kartavyaḥ . \\
\noindent \textbf{(P\_2,4.2)} pratyekam aṅgaśabdaḥ parisamāpyate . \\
\noindent \textbf{(P\_2,4.3)} iha kasmāt na bhavati . \\
\noindent \textbf{(P\_2,4.3)} nandantu kaṭhakālāpāḥ . \\
\noindent \textbf{(P\_2,4.3)} vardhantām kaṭhakauthumāḥ . \\
\noindent \textbf{(P\_2,4.3)} stheṇoḥ . \\
\noindent \textbf{(P\_2,4.3)} stheṇoḥ iti vaktavyam . \\
\noindent \textbf{(P\_2,4.3)} evam api tiṣṭhantu kaṭhakālāpāḥ iti atra api prāpnoti . \\
\noindent \textbf{(P\_2,4.3)} adyatanyām ca . \\
\noindent \textbf{(P\_2,4.3)} adyatanyām ca iti vaktavyam . \\
\noindent \textbf{(P\_2,4.3)} udagāt kaṭhakāpālam . \\
\noindent \textbf{(P\_2,4.3)} pratyaṣṭhāt kaṭhakauthumam . \\
\noindent \textbf{(P\_2,4.3)} udagāt maudapaippalādam . \\
\noindent \textbf{(P\_2,4.7)} grāmapratiṣedhe nagarapratiṣedhaḥ . \\
\noindent \textbf{(P\_2,4.7)} agrāmāḥ iti atra anagarāṇām iti vaktavyam . \\
\noindent \textbf{(P\_2,4.7)} iha mā bhūt . \\
\noindent \textbf{(P\_2,4.7)} mathurāpāṭaliputram iti . \\
\noindent \textbf{(P\_2,4.7)} ubhayataḥ ca grāmāṇām . \\
\noindent \textbf{(P\_2,4.7)} ubhayataḥ ca grāmāṇām pratiṣedhaḥ vaktavyaḥ . \\
\noindent \textbf{(P\_2,4.7)} iha mā bhūt . \\
\noindent \textbf{(P\_2,4.7)} śauryam ca ketavatā ca śauryaketavate . \\
\noindent \textbf{(P\_2,4.7)} jāmbavam ca śālukinī ca jāmbavaśālukinyau \\
\noindent \textbf{(P\_2,4.8)} kṣudrjantavaḥ iti ucyate . \\
\noindent \textbf{(P\_2,4.8)} ke punaḥ kṣudrajantavaḥ . \\
\noindent \textbf{(P\_2,4.8)} kṣottavyāḥ jantavaḥ . \\
\noindent \textbf{(P\_2,4.8)} yadi evam yūkālikṣam kīṭapipīlikam iti na sidhyati . \\
\noindent \textbf{(P\_2,4.8)} evam tarhi anathikāḥ kṣudrajantavaḥ . \\
\noindent \textbf{(P\_2,4.8)} atha vā yeṣām na śoṇitam te kṣudrajantavaḥ . \\
\noindent \textbf{(P\_2,4.8)} atha vā yeṣām ā sahasrāt añjaliḥ na pūryate te kṣudrajantavaḥ . \\
\noindent \textbf{(P\_2,4.8)} atha vā yeṣām gocarmamātram na patitaḥ bhavati te kṣudrajantavaḥ . \\
\noindent \textbf{(P\_2,4.8)} atha va nakulaparyantāḥ kṣudrajantavaḥ . \\
\noindent \textbf{(P\_2,4.9)} kimarthaḥ cakāraḥ . \\
\noindent \textbf{(P\_2,4.9)} evakārārthaḥ cakāraḥ . \\
\noindent \textbf{(P\_2,4.9)} yeṣām virodhaḥ śāśvatikaḥ teṣām dvandve ekavacanam yathā syāt . \\
\noindent \textbf{(P\_2,4.9)} anyat yat prāpnoti tat mā bhūt iti . \\
\noindent \textbf{(P\_2,4.9)} kim ca anyat prāpnoti . \\
\noindent \textbf{(P\_2,4.9)} paśuśakunidvandve virodhinām pūrvavipratiṣiddham iti uktam saḥ pūrvavipratiṣedhaḥ na vaktavyaḥ bhavati . \\
\noindent \textbf{(P\_2,4.10)} aniravasitānām iti ucyate . \\
\noindent \textbf{(P\_2,4.10)} kutaḥ aniravasitānām . \\
\noindent \textbf{(P\_2,4.10)} āryāvartāt aniravasitānām . \\
\noindent \textbf{(P\_2,4.10)} kaḥ punaḥ āryāvartaḥ . \\
\noindent \textbf{(P\_2,4.10)} prāg ādarśāt pratyak kālakavanāt dakṣiṇena himavantam uttareṇa pāriyātram . \\
\noindent \textbf{(P\_2,4.10)} yadi evam kiṣkindhagandikam śakayavanam śauryakrauñcam iti na sidhyati . \\
\noindent \textbf{(P\_2,4.10)} evam tarhi āryanivāsāt aniravasitānām . \\
\noindent \textbf{(P\_2,4.10)} kaḥ punaḥ āryanivāsaḥ . \\
\noindent \textbf{(P\_2,4.10)} grāmaḥ ghoṣaḥ nagaram saṃvāhaḥ iti . \\
\noindent \textbf{(P\_2,4.10)} evam api ye ete mahāntaḥ saṃstyāyāḥ teṣu abhyantarāḥ caṇḍālāḥ mṛtapāḥ ca vasanti tatra caṇḍālamṛtapāḥ iti na sidhyati . \\
\noindent \textbf{(P\_2,4.10)} evam tarhi yājñāt karmaṇaḥ aniravasitānām . \\
\noindent \textbf{(P\_2,4.10)} evam api takṣāyaskāram rajakatantuvāyam iti na sidhyati . \\
\noindent \textbf{(P\_2,4.10)} evam tarhi pātrāt aniravasitānām . \\
\noindent \textbf{(P\_2,4.10)} yaiḥ bhukte pātram saṃskāreṇa śudhyati te aniravasitāḥ . \\
\noindent \textbf{(P\_2,4.10)} yaiḥ bhukte saṃskāreṇa api na śudhyati te niravasitāḥ . \\
\noindent \textbf{(P\_2,4.11)} gavāśvaprabhṛtiṣu yathoccāritam dvandvavṛttam . \\
\noindent \textbf{(P\_2,4.11)} avāśvaprabhṛtiṣu yathoccāritam dvandvavṛttam draṣṭavyam . \\
\noindent \textbf{(P\_2,4.11)} gavāśvam gavāvikam gavaiḍakam . \\
\noindent \textbf{(P\_2,4.12)} bahuprakṛtiḥ phalasenāvanaspatimṛgaśakuntkṣudrajantudhānyatṛṇānām . \\
\noindent \textbf{(P\_2,4.12)} phalasenāvanaspatimṛgaśakuntkṣudrajantudhānyatṛṇānām dvandvaḥ vibhāṣā ekavat bhavati bahuprakṛtiḥ iti vaktavyam . \\
\noindent \textbf{(P\_2,4.12)} phala. badarāmalkam badarāmalakani . \\
\noindent \textbf{(P\_2,4.12)} senā . \\
\noindent \textbf{(P\_2,4.12)} hastyaśvam hastyaśvāḥ . \\
\noindent \textbf{(P\_2,4.12)} vanaspati . \\
\noindent \textbf{(P\_2,4.12)} plakṣanyagrodham plakṣanyarodhāḥ . \\
\noindent \textbf{(P\_2,4.12)} mṛga . \\
\noindent \textbf{(P\_2,4.12)} rurupṛṣatam rurupṛṣatāḥ . \\
\noindent \textbf{(P\_2,4.12)} śakunta . \\
\noindent \textbf{(P\_2,4.12)} haṃsacakravākam haṃsacakravākāḥ . \\
\noindent \textbf{(P\_2,4.12)} kṣudrajantu . \\
\noindent \textbf{(P\_2,4.12)} yūkālikṣam yūkālikṣāḥ . \\
\noindent \textbf{(P\_2,4.12)} dhānya . \\
\noindent \textbf{(P\_2,4.12)} vrīhiyavam vrīhiyavāḥ māṣatilam māṣatilāḥ . \\
\noindent \textbf{(P\_2,4.12)} tṛṇa . \\
\noindent \textbf{(P\_2,4.12)} kuśakāsam kuśakāśāḥ śaraśīryam śaraśīryāḥ . \\
\noindent \textbf{(P\_2,4.12)} kim prayojanam . \\
\noindent \textbf{(P\_2,4.12)} bahuprakṛtiḥ eva yathā syāt . \\
\noindent \textbf{(P\_2,4.12)} kva mā bhūt . \\
\noindent \textbf{(P\_2,4.12)} badarāmalake tiṣṭhataḥ . \\
\noindent \textbf{(P\_2,4.12)} kim punaḥ anena yā prāptiḥ sā niyamyate āhosvit aviśeṣeṇa . \\
\noindent \textbf{(P\_2,4.12)} kim ca ataḥ . \\
\noindent \textbf{(P\_2,4.12)} yadi anena yā prāptiḥ sā niyamyate plakṣanyagrodhau jātiḥ aprāṇinām iti nityaḥ dvandvaikavadbhāvaḥ prāpnoti . \\
\noindent \textbf{(P\_2,4.12)} atha aviśeṣeṇa na doṣaḥ bhavati . \\
\noindent \textbf{(P\_2,4.12)} yathā na doṣaḥ tathāu astu . \\
\noindent \textbf{(P\_2,4.12)} paśuśakunidvandve virodhinām pūrvavipratiṣiddham . paśuśakunidvandve yeṣām ca virodhaḥ śāśvatikaḥ iti etat bhavati pūrvavipratiṣedhena . \\
\noindent \textbf{(P\_2,4.12)} paśuśakunidvandvasya avakāśaḥ mahājorabhram mahājorabhrāḥ haṃsacakravākam haṃsacakravākāḥ . \\
\noindent \textbf{(P\_2,4.12)} yeṣām ca virodhaḥ iti asya avakāśaḥ śramaṇabrāhmaṇam . \\
\noindent \textbf{(P\_2,4.12)} iha ubhayam prāpnoti . \\
\noindent \textbf{(P\_2,4.12)} kākolūkam śvaśṛgālam iti . \\
\noindent \textbf{(P\_2,4.12)} yeṣām ca virodhaḥ iti etat bhavati pūrvavipratiṣedhena . \\
\noindent \textbf{(P\_2,4.12)} saḥ tarhi pūrvavipratiṣedhaḥ vaktavyaḥ . \\
\noindent \textbf{(P\_2,4.12)} na vaktavyaḥ . \\
\noindent \textbf{(P\_2,4.12)} uktam tatra cakārakaraṇasya prayojanam yeṣām ca virodhaḥ śāśvatikaḥ teṣām dvandve ekavacanam yathā syāt . \\
\noindent \textbf{(P\_2,4.12)} anyat yat prāpnoti tat mā bhūt iti . \\
\noindent \textbf{(P\_2,4.12)} aśvavaḍavayoḥ pūrvaliṅgatvāt paśudvandvanapuṃsakam . \\
\noindent \textbf{(P\_2,4.12)} aśvavaḍavayoḥ pūrvaliṅgatvāt paśudvandvanapuṃsakam bhavati pūrvavipratiṣedhena . \\
\noindent \textbf{(P\_2,4.12)} aśvavaḍavayoḥ pūrvaliṅgatvasya avakāśaḥ . \\
\noindent \textbf{(P\_2,4.12)} vibhāṣā paśudvandvanapuṃsakam . \\
\noindent \textbf{(P\_2,4.12)} yadā na paśudvandvanapuṃsakam saḥ avakāśaḥ . \\
\noindent \textbf{(P\_2,4.12)} aśvavaḍavau . \\
\noindent \textbf{(P\_2,4.12)} paśudvandvanapuṃsakasya avakāśaḥ anye paśudvandvāḥ . \\
\noindent \textbf{(P\_2,4.12)} mahājorabhram mahājorabhrāḥ . \\
\noindent \textbf{(P\_2,4.12)} paśudvandvanapuṃsakaprasaṅge ubhayam prāpnoti . \\
\noindent \textbf{(P\_2,4.12)} aśvavaḍavam . \\
\noindent \textbf{(P\_2,4.12)} paśudvandvanapuṃsakam bhavati pūrvavipratiṣedhena . \\
\noindent \textbf{(P\_2,4.12)} saḥ tarhi pūrvavipratiṣedhaḥ vaktavyaḥ . \\
\noindent \textbf{(P\_2,4.12)} na vaktavyaḥ . \\
\noindent \textbf{(P\_2,4.12)} pratipadavidhānāt siddham . pratipadam atra napuṃsakam vidhīyate . \\
\noindent \textbf{(P\_2,4.12)} aśvavaḍavapūrvāpara iti . \\
\noindent \textbf{(P\_2,4.12)} ekavacanam anarthakam samāhāraikatvāt . \\
\noindent \textbf{(P\_2,4.12)} ekavadbhāvaḥ anarthakaḥ . \\
\noindent \textbf{(P\_2,4.12)} kim kāraṇam . \\
\noindent \textbf{(P\_2,4.12)} samāhāraikatvāt . \\
\noindent \textbf{(P\_2,4.12)} ekaḥ ayam arthaḥ samāhāraḥ nāma . \\
\noindent \textbf{(P\_2,4.12)} tasya ekatvāt ekavacanam bhaviṣyati . \\
\noindent \textbf{(P\_2,4.12)} idam tarhi prayojanam . \\
\noindent \textbf{(P\_2,4.12)} etat jñāsyāmi . \\
\noindent \textbf{(P\_2,4.12)} iha nityaḥ vidhiḥ iha vibhāṣā iti . \\
\noindent \textbf{(P\_2,4.12)} na etat asti prayojanam . \\
\noindent \textbf{(P\_2,4.12)} ācāryapravṛttiḥ jñāpayati sarvaḥ dvandvaḥ vibhāṣā ekavat bhavati iti yat ayam tiṣyapunarvasvoḥ nakṣatradvandve bahuvacanasya dvivacanam nityam iti āha . \\
\noindent \textbf{(P\_2,4.12)} idam tarhi prayojanam . \\
\noindent \textbf{(P\_2,4.12)} saḥ napuṃsakam iti vakṣyāmi iti . \\
\noindent \textbf{(P\_2,4.12)} etat api na asti prayojanam . \\
\noindent \textbf{(P\_2,4.12)} liṅgam aśiṣyam lokāśrayatvāt liṅgasya . \\
\noindent \textbf{(P\_2,4.12)} na tarhi idānīm idam vaktavyam . \\
\noindent \textbf{(P\_2,4.12)} vaktavyam ca . \\
\noindent \textbf{(P\_2,4.12)} kim prayojanam . \\
\noindent \textbf{(P\_2,4.12)} pūrvatra nityārtham uttaratra vyabhicārārtham vibhāṣā vṛkṣamṛga iti . \\
\noindent \textbf{(P\_2,4.16)} kim udāharaṇam . \\
\noindent \textbf{(P\_2,4.16)} upadaśam pāṇipādam upadaśāḥ pāṇipādāḥ . \\
\noindent \textbf{(P\_2,4.16)} na etat asti prayojanam . \\
\noindent \textbf{(P\_2,4.16)} ayam dvandvaikavadbhāvaḥ ārabhyate . \\
\noindent \textbf{(P\_2,4.16)} tatra kaḥ prasaṅgaḥ yat anuprayogasya syāt . \\
\noindent \textbf{(P\_2,4.16)} evam tarhi avyayasya saṅkhayā avyayībhāvaḥ api ārabhyate bahuvrīhiḥ api . \\
\noindent \textbf{(P\_2,4.16)} tat yadā tāvat ekavacanam tadā avyayībhāvaḥ anuprayujyate ekārthasya ekārthaḥ iti . \\
\noindent \textbf{(P\_2,4.16)} yadā bahuvacanam tadā bahuvrīhiḥ anuprayujyate bahvarthasya bahvarthaḥ iti . \\
\noindent \textbf{(P\_2,4.19)} kimartham idam ucyate . \\
\noindent \textbf{(P\_2,4.19)} sañjñāyām kanthośīnareṣu iti vakṣyati . \\
\noindent \textbf{(P\_2,4.19)} tat atatpuruṣasya nañsamāsasya karmadhārayasya vā mā bhūt iti . \\
\noindent \textbf{(P\_2,4.19)} na etat asti prayojanam . \\
\noindent \textbf{(P\_2,4.19)} na hi sañjñāyām kanthāntaḥ uśīnareṣu atatpuruṣaḥ nañsamāsaḥ karmadhārayaḥ vā asti . \\
\noindent \textbf{(P\_2,4.19)} uttarārtham tarhi . \\
\noindent \textbf{(P\_2,4.19)} upajñopakramam tadādyācikhyāsāyām iti vakṣyati . \\
\noindent \textbf{(P\_2,4.19)} tat atatpuruṣasya nañsamāsasya karmadhārayasya vā mā bhūt iti . \\
\noindent \textbf{(P\_2,4.19)} etat api na asti prayojanam . \\
\noindent \textbf{(P\_2,4.19)} na hi tadādyācikhyāsāyām upajñopakramāntaḥ atatpuruṣaḥ nañsamāsaḥ karmadhārayaḥ vā asti . \\
\noindent \textbf{(P\_2,4.19)} uttarārtham eva tarhi . \\
\noindent \textbf{(P\_2,4.19)} chāyā bāhulye iti vakṣyati . \\
\noindent \textbf{(P\_2,4.19)} tat atatpuruṣasya nañsamāsasya karmadhārayasya vā mā bhūt iti . \\
\noindent \textbf{(P\_2,4.19)} etat api na asti prayojanam . \\
\noindent \textbf{(P\_2,4.19)} na hi chāyāntaḥ bāhulye atatpuruṣaḥ nañsamāsaḥ karmadhārayaḥ vā asti . \\
\noindent \textbf{(P\_2,4.19)} uttarārtham eva tarhi . \\
\noindent \textbf{(P\_2,4.19)} sabhā rājāmanuṣyapūrvā aśālā ca iti vakṣyati . \\
\noindent \textbf{(P\_2,4.19)} tat atatpuruṣasya nañsamāsasya karmadhārayasya vā mā bhūt iti . \\
\noindent \textbf{(P\_2,4.19)} etat api na asti prayojanam . \\
\noindent \textbf{(P\_2,4.19)} na hi sabhāntaḥ aśālāyām atatpuruṣaḥ nañsamāsaḥ karmadhārayaḥ vā asti . \\
\noindent \textbf{(P\_2,4.19)} idam tarhi . \\
\noindent \textbf{(P\_2,4.19)} vibhāṣā senāsurā iti vakṣyati . \\
\noindent \textbf{(P\_2,4.19)} tat atatpuruṣasya nañsamāsasya karmadhārayasya vā mā bhūt iti . \\
\noindent \textbf{(P\_2,4.19)} dṛḍhasenaḥ rājā . \\
\noindent \textbf{(P\_2,4.19)} anañ iti kimartham . \\
\noindent \textbf{(P\_2,4.19)} asenā . \\
\noindent \textbf{(P\_2,4.19)} akarmadhārayaḥ iti kimartham . \\
\noindent \textbf{(P\_2,4.19)} paramsenā uttamasenā . \\
\noindent \textbf{(P\_2,4.26)} kimartham idam ucyate . \\
\noindent \textbf{(P\_2,4.26)} dvandvaḥ ayam ubhayapadārthapradhānaḥ . \\
\noindent \textbf{(P\_2,4.26)} tatra kadā cit pūrvapadasya yat liṅgam tat samāsasya api syāt kadā cit uttarapadasya . \\
\noindent \textbf{(P\_2,4.26)} iṣyate ca parasya yat liṅgam tat samāsasya syāt iti . \\
\noindent \textbf{(P\_2,4.26)} tat ca antareṇa yatnam na sidhyati iti paravat liṅgam dvandvatatpuruṣayoḥ iti . \\
\noindent \textbf{(P\_2,4.26)} evamartham idam ucyate . \\
\noindent \textbf{(P\_2,4.26)} tatpuruṣaḥ ca kaḥ prayojayati . \\
\noindent \textbf{(P\_2,4.26)} yaḥ pūrvapadārthapradhānaḥ ekadeśisamāsaḥ ardhapippalī iti . \\
\noindent \textbf{(P\_2,4.26)} yaḥ hi uttarapadārthapradhānaḥ daivakṛtam tasya paravat liṅgam . \\
\noindent \textbf{(P\_2,4.26)} paravat liṅgam dvandvatatpuruṣayoḥ iti cet prāptāpannālampūrvagatisamāseṣu pratiṣedhaḥ . \\
\noindent \textbf{(P\_2,4.26)} paravat liṅgam dvandvatatpuruṣayoḥ iti cet prāptāpannālampūrvagatisamāseṣu pratiṣedhaḥ vaktavyaḥ . \\
\noindent \textbf{(P\_2,4.26)} prāptaḥ jīvikām praptajīvikaḥ āpannaḥ jīvikām apannajīvikaḥ . \\
\noindent \textbf{(P\_2,4.26)} alampūrvaḥ . \\
\noindent \textbf{(P\_2,4.26)} alam jīvikāyāḥ alamjīvikaḥ . \\
\noindent \textbf{(P\_2,4.26)} gatisamāsa . \\
\noindent \textbf{(P\_2,4.26)} niṣkauśāmbiḥ nirvārāṇasiḥ . \\
\noindent \textbf{(P\_2,4.26)} pūrvapadasya ca . \\
\noindent \textbf{(P\_2,4.26)} pūrvapadasya ca pratiṣedhaḥ vaktavyaḥ . \\
\noindent \textbf{(P\_2,4.26)} mayūrīkukkuṭau . \\
\noindent \textbf{(P\_2,4.26)} yadi punaḥ yathājātīyakam parasya liṅgam tathājātīyakam samāsāt anyat atidiśyeta . \\
\noindent \textbf{(P\_2,4.26)} samāsāt anyat liṅgam iti cet aśvavaḍavayoḥ ṭāblugvacanam . \\
\noindent \textbf{(P\_2,4.26)} samāsāt anyat liṅgam iti cet aśvavaḍavayoḥ ṭāpaḥ luk vaktavyaḥ . \\
\noindent \textbf{(P\_2,4.26)} aśvavaḍavau . \\
\noindent \textbf{(P\_2,4.26)} nipātanāt siddham . \\
\noindent \textbf{(P\_2,4.26)} nipātanāt siddham etat . \\
\noindent \textbf{(P\_2,4.26)} kim nipātanam . \\
\noindent \textbf{(P\_2,4.26)} āśvavaḍavapūrvāpara iti . \\
\noindent \textbf{(P\_2,4.26)} upasarjanahrasvatvam vā . \\
\noindent \textbf{(P\_2,4.26)} atha vā upasarjanasya iti hrasvatvam bhaviṣyati . \\
\noindent \textbf{(P\_2,4.26)} iha api tarhi prāpnoti . \\
\noindent \textbf{(P\_2,4.26)} kukkuṭamayūryau . \\
\noindent \textbf{(P\_2,4.26)} astu . \\
\noindent \textbf{(P\_2,4.26)} paravat liṅgam iti śabdaśabdārthau . \\
\noindent \textbf{(P\_2,4.26)} paravat liṅgam iti śabdaśabdārthau atidiśyete . \\
\noindent \textbf{(P\_2,4.26)} tatra aupadeśikasya hrasvatvam ātideśikasya śravaṇam bhaviṣyati . \\
\noindent \textbf{(P\_2,4.26)} idam tarhi . \\
\noindent \textbf{(P\_2,4.26)} dattā ca kārīṣagandhyā ca dattākārīṣagandhye dattā ca gārgyāyaṇī dattāgārgyāyaṇyau . \\
\noindent \textbf{(P\_2,4.26)} dvau ṣyaṅau dvau ṣphau ca prāpnutaḥ . \\
\noindent \textbf{(P\_2,4.26)} stām .pūmvadbhāvena ekasya nivṛttiḥ bhaviṣyati . \\
\noindent \textbf{(P\_2,4.26)} idam tarhi . \\
\noindent \textbf{(P\_2,4.26)} dattā ca yuvatiḥ ca dattāyuvatī . \\
\noindent \textbf{(P\_2,4.26)} dvau tiśabdau prāpnutaḥ . \\
\noindent \textbf{(P\_2,4.26)} tasmāt na etat śakyam vaktum śabdaśabdārthau atidiśyete iti . \\
\noindent \textbf{(P\_2,4.26)} nanu ca uktam samāsāt anyat liṅgam iti cet aśvavaḍavayoḥ ṭāblugvacanam iti . \\
\noindent \textbf{(P\_2,4.26)} parihṛtam etat : nipātanāt siddham iti . \\
\noindent \textbf{(P\_2,4.26)} atha vā na evam vijñāyate parasya eva paravat iti . \\
\noindent \textbf{(P\_2,4.26)} katham tarhi . \\
\noindent \textbf{(P\_2,4.26)} parasya iva paravat iti . \\
\noindent \textbf{(P\_2,4.26)} yathājātīyakam parasya liṅgam tathājātīyakam samāsasya atidiśyate . \\
\noindent \textbf{(P\_2,4.26)} atha pūrvapadasya na pratiṣidhyate prāptādiṣu katham . \\
\noindent \textbf{(P\_2,4.26)} prāptādiṣu ca ekadeśigrahaṇāt siddham . \\
\noindent \textbf{(P\_2,4.26)} dvandvaikadeśinoḥ iti vakṣyāmi . \\
\noindent \textbf{(P\_2,4.26)} tat ekadeśigrahaṇam kartavyam . \\
\noindent \textbf{(P\_2,4.26)} na kartavyam . \\
\noindent \textbf{(P\_2,4.26)} ekadeśisamāsaḥ na ārapsyate . \\
\noindent \textbf{(P\_2,4.26)} katham ardhapippalī iti . \\
\noindent \textbf{(P\_2,4.26)} samānādhikaraṇasamāsaḥ bhaviṣyati . \\
\noindent \textbf{(P\_2,4.26)} ardham ca sā pippalī ca ardhapippalī iti . \\
\noindent \textbf{(P\_2,4.26)} na sidhyati . \\
\noindent \textbf{(P\_2,4.26)} paratvāt ṣaṣṭhīsamāsaḥ prāpnoti . \\
\noindent \textbf{(P\_2,4.26)} adya punaḥ ayam ekadeśisamāsaḥ ārabhyamāṇaḥ ṣaṣṭhīsamāsam bādhate . \\
\noindent \textbf{(P\_2,4.26)} iṣyate ca ṣaṣṭhīsamāsaḥ api . \\
\noindent \textbf{(P\_2,4.26)} tat yathā . \\
\noindent \textbf{(P\_2,4.26)} apūpārdham mayā bhakṣitam . \\
\noindent \textbf{(P\_2,4.26)} grāmārdham mayā labdham iti . \\
\noindent \textbf{(P\_2,4.26)} evam pippalyardham iti bhavitavyam . \\
\noindent \textbf{(P\_2,4.26)} katham ardhapippalī iti . \\
\noindent \textbf{(P\_2,4.26)} samānādhikaraṇaḥ bhaviṣyati . \\
\noindent \textbf{(P\_2,4.29)} anuvākādayaḥ puṃsi . \\
\noindent \textbf{(P\_2,4.29)} anuvākādayaḥ puṃsi bhāṣyante iti vaktavyam . \\
\noindent \textbf{(P\_2,4.29)} anuvākaḥ śamyuvākaḥ sūktavākaḥ . \\
\noindent \textbf{(P\_2,4.30)} puṇyasudinābhyām ahnaḥ napuṃsakatvam vaktavyam . \\
\noindent \textbf{(P\_2,4.30)} puṇyāham sudināham . \\
\noindent \textbf{(P\_2,4.30)} pathaḥ saṅkhyāvyayādeḥ . \\
\noindent \textbf{(P\_2,4.30)} pathaḥ saṅkhyāvyayādeḥ iti vaktavyam . \\
\noindent \textbf{(P\_2,4.30)} dvipatham tripatham catuṣpatham utpatham vipatham . \\
\noindent \textbf{(P\_2,4.30)} dviguḥ ca . \\
\noindent \textbf{(P\_2,4.30)} dviguḥ ca samāsaḥ napuṃsakaliṅgaḥ bhavati iti vaktavyam . \\
\noindent \textbf{(P\_2,4.30)} pañcagavam daśagavam . \\
\noindent \textbf{(P\_2,4.30)} akārāntottarapadaḥ dviguḥ striyām bhāṣyate iti vaktavyam . \\
\noindent \textbf{(P\_2,4.30)} pañcapūlī daśapūlī . \\
\noindent \textbf{(P\_2,4.30)} vā ābantaḥ . \\
\noindent \textbf{(P\_2,4.30)} vā ābantaḥ striyām bhāṣyate iti vaktavyam . \\
\noindent \textbf{(P\_2,4.30)} pañcakhaṭvam pañcakhaṭvī daśakhaṭvam daśakhaṭvī . \\
\noindent \textbf{(P\_2,4.30)} anaḥ nalopaḥ ca vā ca striyām bhāṣyate iti vaktavyam . \\
\noindent \textbf{(P\_2,4.30)} pañcatakṣam pañcatakṣī daśatakṣam daśatakṣī . \\
\noindent \textbf{(P\_2,4.30)} pātrādibhyaḥ pratiṣedhaḥ vaktavyaḥ . \\
\noindent \textbf{(P\_2,4.30)} dvipātram pañcapātram . \\
\noindent \textbf{(P\_2,4.31)} ardharcādayaḥ iti vaktavyam . \\
\noindent \textbf{(P\_2,4.31)} ardharcam ardharcaḥ kārṣāpaṇam kārṣāpaṇaḥ gomayam gomayaḥ sāram sāraḥ . \\
\noindent \textbf{(P\_2,4.31)} tat tarhi vaktavyam . \\
\noindent \textbf{(P\_2,4.31)} na vaktavyam . \\
\noindent \textbf{(P\_2,4.31)} bahuvacananirdeśāt ādyarthaḥ gamyate . \\
\noindent \textbf{(P\_2,4.32.1)} anvādeśe samānādhikaraṇagrahaṇam . \\
\noindent \textbf{(P\_2,4.32.1)} anvādeśe samānādhikaraṇagrahaṇam kartavyam . \\
\noindent \textbf{(P\_2,4.32.1)} kim prayojanam . \\
\noindent \textbf{(P\_2,4.32.1)} devadattam bhojaya imam ca iti aprasaṅgārtham . \\
\noindent \textbf{(P\_2,4.32.1)} iha mā bhūt . \\
\noindent \textbf{(P\_2,4.32.1)} devadattam bhojaya imam ca yajña dattam bhojaya iti . \\
\noindent \textbf{(P\_2,4.32.1)} anvādeśaḥ ca kathitānukathanamātram . \\
\noindent \textbf{(P\_2,4.32.1)} anvādeśaḥ ca kathitānukathanamātram draṣṭavyam . \\
\noindent \textbf{(P\_2,4.32.1)} tat dveṣyam vijānīyāt : idamā kathitam idamā yadā anukathyate iti . \\
\noindent \textbf{(P\_2,4.32.1)} tat ācāryaḥ suhṛt bhūtvā ācaṣṭe : anvādeśaḥ ca kathitānukathanamātram draṣṭavyam iti . \\
\noindent \textbf{(P\_2,4.32.2)} atha kimartham aśādeśaḥ kriyate na tṛtīyādiṣu iti eva ucyeta . \\
\noindent \textbf{(P\_2,4.32.2)} tatra ṭāyām osi ca enena bhavitavyam . \\
\noindent \textbf{(P\_2,4.32.2)} anyāḥ sarvāḥ halādayaḥ vibhaktayaḥ . \\
\noindent \textbf{(P\_2,4.32.2)} tatra idrūpalope kṛte kevalam idamaḥ anudāttatvam vaktavyam . \\
\noindent \textbf{(P\_2,4.32.2)} ataḥ uttaram paṭhati . \\
\noindent \textbf{(P\_2,4.32.2)} ādeśavacanam sākackārtham . \\
\noindent \textbf{(P\_2,4.32.2)} ādeśavacanam sākackārtham kriyate . \\
\noindent \textbf{(P\_2,4.32.2)} sākackasya api ādeśaḥ yathā syāt . \\
\noindent \textbf{(P\_2,4.32.2)} imakābhyām chātrābhyām rātriḥ adhītā atho ābhyām api adhītam . \\
\noindent \textbf{(P\_2,4.32.2)} atha kimartham śitkaraṇam . \\
\noindent \textbf{(P\_2,4.32.2)} śitkaraṇam sarvādeśārtham . śitkaraṇam kriyate sarvādeśārtham . \\
\noindent \textbf{(P\_2,4.32.2)} śit sarvasya iti sarvādeśaḥ yathā syāt : imakābhyām chātrābhyām rātriḥ adhītā atho* ābhyām api adhītam iti . \\
\noindent \textbf{(P\_2,4.32.2)} akriyamāṇe hi śitkaraṇe alaḥ antyasya vidhayaḥ bhavanti iti antyasya prasajyeta . \\
\noindent \textbf{(P\_2,4.32.2)} na vā antyasya vikāravacanānarthakyāt . \\
\noindent \textbf{(P\_2,4.32.2)} na vā kartavyam . \\
\noindent \textbf{(P\_2,4.32.2)} kim kāraṇam . \\
\noindent \textbf{(P\_2,4.32.2)} antyasya vikāravacanānarthakyāt . \\
\noindent \textbf{(P\_2,4.32.2)} akārasya akāravacane prayojanam na asti iti kṛtvā antareṇa api śakāram sarvādeśaḥ bhaviṣyati . \\
\noindent \textbf{(P\_2,4.32.2)} arthavat tu ādeśapratiṣedhāṛtham . \\
\noindent \textbf{(P\_2,4.32.2)} arthavat tu asya akāravacanam . \\
\noindent \textbf{(P\_2,4.32.2)} kaḥ arthaḥ . \\
\noindent \textbf{(P\_2,4.32.2)} ādeśapratiṣedhāṛtham . \\
\noindent \textbf{(P\_2,4.32.2)} ye anye akārasya ādeśāḥ prāpnuvanti tadbādhanārtham . \\
\noindent \textbf{(P\_2,4.32.2)} tat yathā . \\
\noindent \textbf{(P\_2,4.32.2)} maḥ rāji samaḥ kvau iti : makārasya makāravacane prayojanam na asti iti kṛtvā anusvārādayaḥ bādhyante . \\
\noindent \textbf{(P\_2,4.32.2)} tasmāt śitkaraṇam . tasmāt śakāraḥ kartavyaḥ . \\
\noindent \textbf{(P\_2,4.32.2)} na kartavyaḥ . \\
\noindent \textbf{(P\_2,4.32.2)} praśliṣṭanirdeśaḥ ayam . \\
\noindent \textbf{(P\_2,4.32.2)} a* a iti . \\
\noindent \textbf{(P\_2,4.32.2)} anekālśit sarvasya iti sarvādeśaḥ bhaviṣyati . \\
\noindent \textbf{(P\_2,4.32.2)} atha vā vicitrāḥ taddhitavṛttayaḥ . \\
\noindent \textbf{(P\_2,4.32.2)} na anvādeśe akac utpatsyate . \\
\noindent \textbf{(P\_2,4.33)} kimartham tratasoḥ anudāttatvam ucyate . \\
\noindent \textbf{(P\_2,4.33)} udāttau mā bhūtām iti . \\
\noindent \textbf{(P\_2,4.33)} na etat asti prayojanam . \\
\noindent \textbf{(P\_2,4.33)} litsvare kṛte nighāte etadaḥ anudāttatvena siddham . \\
\noindent \textbf{(P\_2,4.33)} idam iha sampradhāryam . \\
\noindent \textbf{(P\_2,4.33)} anudāttatvam kriyatām litsvaraḥ iti . \\
\noindent \textbf{(P\_2,4.33)} kim atra kartavyam . \\
\noindent \textbf{(P\_2,4.33)} paratvāt litsvaraḥ . \\
\noindent \textbf{(P\_2,4.33)} nityatvāt anudāttatvam . \\
\noindent \textbf{(P\_2,4.33)} kṛte api litsvare prāptnoti akṛte api . \\
\noindent \textbf{(P\_2,4.33)} tatra nityatvāt anudāttatve kṛte liti pūrvaḥ udāttabhāvī na asti iti kṛtvā yathāprāptaḥ pratyayasvaraḥ prasajyeta . \\
\noindent \textbf{(P\_2,4.33)} tat yathā goṣpadapram vṛṣṭaḥ devaḥ iti ūlope kṛte pūrvaḥ udāttabhāvī na asti iti kṛtvā yathāprāptaḥ pratyayasvaraḥ bhavati . \\
\noindent \textbf{(P\_2,4.33)} tasmāt tratasoḥ anudāttatvam vaktavyam . \\
\noindent \textbf{(P\_2,4.34)} kasya ayam enaḥ vidhīyate . \\
\noindent \textbf{(P\_2,4.34)} etadaḥ prāpnoti . \\
\noindent \textbf{(P\_2,4.34)} idamaḥ api tu iṣyate . \\
\noindent \textbf{(P\_2,4.34)} tat idamaḥ grahaṇam kartavyam . \\
\noindent \textbf{(P\_2,4.34)} na kartavyam . \\
\noindent \textbf{(P\_2,4.34)} prakṛtam anuvartate . \\
\noindent \textbf{(P\_2,4.34)} kva prakṛtam . \\
\noindent \textbf{(P\_2,4.34)} idamaḥ anvādeśe aś anudāttaḥ tṛtīyādau iti . \\
\noindent \textbf{(P\_2,4.34)} yadi tat anuvartate etadaḥ tratasoḥ tratasau ca anudāttau iti idamaḥ ca iti idamaḥ api prāpnoti . \\
\noindent \textbf{(P\_2,4.34)} na eṣaḥ doṣaḥ . \\
\noindent \textbf{(P\_2,4.34)} sambandham anuvartiṣyate . \\
\noindent \textbf{(P\_2,4.34)} idamaḥ anvādeśe aś anudāttaḥ tṛtīyādau . \\
\noindent \textbf{(P\_2,4.34)} etadaḥ tratasoḥ tratasau ca anudāttau idamaḥ anvādeśe aś anudāttaḥ tṛtīyādau aś bhavati . \\
\noindent \textbf{(P\_2,4.34)} tataḥ dvitīyāṭaussu enaḥ idamaḥ etadaḥ ca . \\
\noindent \textbf{(P\_2,4.34)} tṛtīyādau iti nivṛttam . \\
\noindent \textbf{(P\_2,4.34)} atha vā maṇḍūkagatayaḥ adhikārāḥ . \\
\noindent \textbf{(P\_2,4.34)} tat yathā maṇḍūkāḥ utplutya utplutya gacchanti tadvat adhikārāḥ . \\
\noindent \textbf{(P\_2,4.34)} atha vā ekayogaḥ kariṣyate . \\
\noindent \textbf{(P\_2,4.34)} idamaḥ anvādeśe aś anudāttaḥ tṛtīyādau iti etadaḥ tratasoḥ tratasau ca anudāttau . \\
\noindent \textbf{(P\_2,4.34)} tataḥ dvitīyāṭaussu enaḥ idamaḥ etadaḥ ca . \\
\noindent \textbf{(P\_2,4.34)} atha vā ubhayam nivṛttam . \\
\noindent \textbf{(P\_2,4.34)} tat apekṣiṣyāmahe . \\
\noindent \textbf{(P\_2,4.34)} enat iti napuṃsakaikavacane . \\
\noindent \textbf{(P\_2,4.34)} enat iti napuṃsakaikavacanekartavyam . \\
\noindent \textbf{(P\_2,4.34)} kuṇḍam ānaya prakṣālaya enat parivartaya enat . \\
\noindent \textbf{(P\_2,4.34)} yadi enat kriyate enaḥ na kartavyaḥ . \\
\noindent \textbf{(P\_2,4.34)} kā rūpasiddhiḥ : atho enam atho ene atho enān iti ṭyadādyatvena siddham . \\
\noindent \textbf{(P\_2,4.34)} yadi evam enaśritakaḥ na sidhyati . \\
\noindent \textbf{(P\_2,4.34)} enacchritakaḥ iti pāpnoti . \\
\noindent \textbf{(P\_2,4.34)} yathālakṣaṇam aprayukte . \\
\noindent \textbf{(P\_2,4.35)} jagdhyādiṣu ārdhadhātukāśrayatvāt sati tasmin vidhānam . \\
\noindent \textbf{(P\_2,4.35)} jagdhyādiṣu ārdhadhātukāśrayatvāt sati tasmin ārdhadhātuke jagdhyādibhiḥ bhavitavyam . \\
\noindent \textbf{(P\_2,4.35)} kim ataḥ yat sati bhavitavyam . \\
\noindent \textbf{(P\_2,4.35)} tatra utsargalakṣaṇapratiṣedhaḥ . \\
\noindent \textbf{(P\_2,4.35)} tatra utsargalakṣaṇam kāryam prāpnoti . \\
\noindent \textbf{(P\_2,4.35)} tasya pratiṣedhaḥ vaktavyaḥ . \\
\noindent \textbf{(P\_2,4.35)} bhavyam praveyam ākhyeyam . \\
\noindent \textbf{(P\_2,4.35)} ṇyati avasthite aniṣṭe pratyaye ādeśaḥ syāt . \\
\noindent \textbf{(P\_2,4.35)} ṇyataḥ śravaṇam prasajyeta . \\
\noindent \textbf{(P\_2,4.35)} na eṣaḥ doṣaḥ . \\
\noindent \textbf{(P\_2,4.35)} sāmānyāśrayatvāt viśeṣasya anāśrayaḥ . \\
\noindent \textbf{(P\_2,4.35)} sāmānyena hi āśrīyamāṇe viśeṣaḥ na āśritaḥ bhavati . \\
\noindent \textbf{(P\_2,4.35)} tatra ārdhadhātukasāmānye jagdhyādiṣu kṛteṣu yaḥ yataḥ pratyayaḥ prāpnoti saḥ tataḥ bhaviṣyati . \\
\noindent \textbf{(P\_2,4.35)} sāmānyāśrayatvāt viśeṣasya anāśrayaḥ iti cet uvarṇākārāntebhyaḥ ṇyadvidhiprasaṅgaḥ . \\
\noindent \textbf{(P\_2,4.35)} sāmānyāśrayatvāt viśeṣasya anāśrayaḥ iti cet uvarṇākārāntebhyaḥ ṇyat prāpnoti . \\
\noindent \textbf{(P\_2,4.35)} lavyam pavyam iti . \\
\noindent \textbf{(P\_2,4.35)} ārdhadhātukasāmānye guṇe kṛti yi pratyayasāmānya ca vāntādeśe halantāt iti ṇyat prāpnoti . \\
\noindent \textbf{(P\_2,4.35)} iha ca ditsyam dhitsyam ārdhadhātukasāmānye akāralope kṛte halantāt iti ṇyat prāpnoti . \\
\noindent \textbf{(P\_2,4.35)} paurvāparyābhāt ca sāmānyena anupapattiḥ . \\
\noindent \textbf{(P\_2,4.35)} paurvāparyābhāt ca sāmānyena jagdhyādīnām anupapattiḥ . \\
\noindent \textbf{(P\_2,4.35)} na hi sāmānyena paurvāparyam asti . \\
\noindent \textbf{(P\_2,4.35)} siddham tu sārvadhātuke pratiṣedhāt . \\
\noindent \textbf{(P\_2,4.35)} siddham etat . \\
\noindent \textbf{(P\_2,4.35)} katham . \\
\noindent \textbf{(P\_2,4.35)} aviśeṣeṇa jagdhyādīn uktvā sārvadhātuke na iti pratiṣedham vakṣyāmi . \\
\noindent \textbf{(P\_2,4.35)} sidhyati . \\
\noindent \textbf{(P\_2,4.35)} sūtram tarhi bhidyate . \\
\noindent \textbf{(P\_2,4.35)} yathānyāsam eva astu . \\
\noindent \textbf{(P\_2,4.35)} nanu ca uktam jagdhyādiṣu ārdhadhātukāśrayatvāt sati tasmin vidhānam iti . \\
\noindent \textbf{(P\_2,4.35)} parihṛtam etat sāmānyāśrayatvāt viśeṣasya anāśrayaḥ iti . \\
\noindent \textbf{(P\_2,4.35)} nanu ca uktam sāmānyāśrayatvāt viśeṣasya anāśrayaḥ iti cet uvarṇākārāntebhyaḥ ṇyadvidhiprasaṅgaḥ iti . \\
\noindent \textbf{(P\_2,4.35)} na eṣaḥ doṣaḥ . \\
\noindent \textbf{(P\_2,4.35)} vakṣyati tatra ajgrahaṇasya prayojanam ajantabhūtapūrvamātrāt api yathā syāt iti . \\
\noindent \textbf{(P\_2,4.35)} yat api ucyate paurvāparyābhāt ca sāmānyena anupapattiḥ iti . \\
\noindent \textbf{(P\_2,4.35)} arthasiddhiḥ eva eṣā yat sāmānyena paurvāparyam na asti . \\
\noindent \textbf{(P\_2,4.35)} asati paurvāparye viṣayasaptamī vijñāsyate . \\
\noindent \textbf{(P\_2,4.35)} ārdhadhātukaviṣaye iti . \\
\noindent \textbf{(P\_2,4.35)} atha vā ārdhadhātukāsu iti vakṣyāmi . \\
\noindent \textbf{(P\_2,4.35)} kāsu ārdhadhātukāsu . \\
\noindent \textbf{(P\_2,4.35)} uktiṣu yuktiṣu rūḍhiṣu pratītiṣu śrutiṣu sañjñāsu \\
\noindent \textbf{(P\_2,4.36)} lyabgrahaṇam kimartham na ti kiti iti eva siddham . \\
\noindent \textbf{(P\_2,4.36)} lyapi kṛte na prāpnoti . \\
\noindent \textbf{(P\_2,4.36)} idam iha sampradhāryam . \\
\noindent \textbf{(P\_2,4.36)} lyap kriyatām ādeśaḥ iti . \\
\noindent \textbf{(P\_2,4.36)} kim atra kartavyam . \\
\noindent \textbf{(P\_2,4.36)} paratvāt lyap . \\
\noindent \textbf{(P\_2,4.36)} antaraṅgaḥ ādeśaḥ . \\
\noindent \textbf{(P\_2,4.36)} evam tarhi siddhe sati yat lyabgrahaṇam karoti tat jñāpayati ācāryaḥ antaraṅgān api vidhīn bahiraṅgaḥ lyap bādhate iti . \\
\noindent \textbf{(P\_2,4.36)} kim etasya jñāpane prayojanam . \\
\noindent \textbf{(P\_2,4.36)} lyabadeśe upadeśivadvacanam anādiṣṭārtham bahiraṅgalakṣaṇatvāt iti vakṣyati . \\
\noindent \textbf{(P\_2,4.36)} tat na vaktavyam bhavati . \\
\noindent \textbf{(P\_2,4.36)} jagdhiḥ vidhiḥ lyapi yat tat akasmāt siddham asti kiti iti vidhānāt . \\
\noindent \textbf{(P\_2,4.36)} hiprabhṛtīn tu sadā bahiraṅgaḥ lyaP bharati iti kṛtam tat u viddhi . \\
\noindent \textbf{(P\_2,4.36)} eṣaḥ eva arthaḥ jagdhau siddhe antaraṅgatvāt ti kiti iti lyap ucyate . \\
\noindent \textbf{(P\_2,4.36)} jñāpayati antaraṅgāṇām lyapā bhavati bādhanam . \\
\noindent \textbf{(P\_2,4.37)} ghasḷbhāve aci upasaṅkhyānam . ghasḷbhāve aci upasaṅkhyānam kartavyam . \\
\noindent \textbf{(P\_2,4.37)} prātti iti praghasaḥ . \\
\noindent \textbf{(P\_2,4.42-43)} KA\_I,485. 2-5 Ro\_II,874 {1/6}     kimayam vadhiḥ vyañjantaḥ āhosvit adantaḥ . \\
\noindent \textbf{(P\_2,4.42-43)} KA\_I,485. 2-5 Ro\_II,874 {2/6}     kim ca ataḥ . \\
\noindent \textbf{(P\_2,4.42-43)} KA\_I,485. 2-5 Ro\_II,874 {3/6}     yadi vyañjanāntaḥ vadhau vyañjanānte uktam . kim uktam . \\
\noindent \textbf{(P\_2,4.42-43)} KA\_I,485. 2-5 Ro\_II,874 {4/6}     vadhyādeśe vṛddhitatvapratiṣedhaḥ iḍvidhiḥ ca iti . \\
\noindent \textbf{(P\_2,4.42-43)} KA\_I,485. 2-5 Ro\_II,874 {5/6}     atha adantaḥ na doṣaḥ bhavati . \\
\noindent \textbf{(P\_2,4.42-43)} KA\_I,485. 2-5 Ro\_II,874 {6/6}     yathā na doṣaḥ tathā astu . \\
\noindent \textbf{(P\_2,4.45)} iṇvat ikaḥ . \\
\noindent \textbf{(P\_2,4.45)} iṇvat ikaḥ iti vaktavyam . \\
\noindent \textbf{(P\_2,4.45)} iha api yathā syāt . \\
\noindent \textbf{(P\_2,4.45)} adhyagāt adhyagātām . \\
\noindent \textbf{(P\_2,4.46)} iṇvat ikaḥ iti eva . \\
\noindent \textbf{(P\_2,4.46)} adhigamayati adhigamayataḥ aghigamayanti . \\
\noindent \textbf{(P\_2,4.47)} iṇvat ikaḥ iti eva . \\
\noindent \textbf{(P\_2,4.47)} aghijigamiṣati adhijigamiśataḥ adhijigamiṣanti . \\
\noindent \textbf{(P\_2,4.49)} ṅitkaraṇam kimartham . \\
\noindent \textbf{(P\_2,4.49)} gāṅi anubandhakaraṇam viśeṣaṇāṛtham . \\
\noindent \textbf{(P\_2,4.49)} gāṅi anubandhakaraṇam kriyate viśeṣaṇāṛtham . \\
\noindent \textbf{(P\_2,4.49)} kva viśeṣaṇārthena arthaḥ . \\
\noindent \textbf{(P\_2,4.49)} gāṅkuṭādibhyaḥ añṇit ṅit iti . \\
\noindent \textbf{(P\_2,4.49)} gākuṭādibhyaḥ añṇit ṅit iti iyati ucyamāne iṇādeśasya api prasajyeta . \\
\noindent \textbf{(P\_2,4.49)} jñāpakam vā sānubandhakasya ādeśavacane itkāryābhāvasya . \\
\noindent \textbf{(P\_2,4.49)} atha vā etat jñāpayati ācāryaḥ sānubandhakasya ādeśe itkāryam na bhavati iti . \\
\noindent \textbf{(P\_2,4.49)} kim etasya jñāpane prayojanam . \\
\noindent \textbf{(P\_2,4.49)} prayojanam cakṣiṅaḥ khyāñ . ṅitaḥ iti ātmanepadam na bhavati . \\
\noindent \textbf{(P\_2,4.49)} laṭaḥ śatṛśānacau . laṭaḥ śatṛśānacau prayojanam . \\
\noindent \textbf{(P\_2,4.49)} pacamānaḥ yajamānaḥ iti . \\
\noindent \textbf{(P\_2,4.49)} ṭitaḥ iti etvam na bhavati . \\
\noindent \textbf{(P\_2,4.49)} yuvoḥ anākau . \\
\noindent \textbf{(P\_2,4.49)} yuvoḥ anākau ca prayojanam . \\
\noindent \textbf{(P\_2,4.49)} nandanaḥ kārakaḥ nandanā kārikā iti . \\
\noindent \textbf{(P\_2,4.49)} ugillakṣaṇau ṅībnumau na bhavataḥ . \\
\noindent \textbf{(P\_2,4.49)} meḥ ca ananubandhakasya amvacanam . \\
\noindent \textbf{(P\_2,4.49)} meḥ ca ananubandhakasya amvaktavyaḥ . \\
\noindent \textbf{(P\_2,4.49)} acinavam akaravam asunavam . \\
\noindent \textbf{(P\_2,4.49)} atyalpam idam ucyate . \\
\noindent \textbf{(P\_2,4.49)} tiptibmipām iti vaktavyam . \\
\noindent \textbf{(P\_2,4.49)} iha api yathā syāt : veda vettha . \\
\noindent \textbf{(P\_2,4.49)} asya jñāpakasya santi doṣāḥ santi prayojanāni . \\
\noindent \textbf{(P\_2,4.49)} doṣāḥ samāḥ bhūyāṃsaḥ vā . \\
\noindent \textbf{(P\_2,4.49)} tasmāt na arthaḥ anena jñāpakena . \\
\noindent \textbf{(P\_2,4.49)} katham yāni prayojanāni . \\
\noindent \textbf{(P\_2,4.49)} na etāni santi . \\
\noindent \textbf{(P\_2,4.49)} iha tāvat . \\
\noindent \textbf{(P\_2,4.49)} cakṣiṅaḥ khyāñ iti . \\
\noindent \textbf{(P\_2,4.49)} ñitkaraṇasāmarthyāt vibhāṣā ātmanepadam bhaviṣyati . \\
\noindent \textbf{(P\_2,4.49)} laṭaḥ śatṛśānacau iti . \\
\noindent \textbf{(P\_2,4.49)} vakṣyati etat . \\
\noindent \textbf{(P\_2,4.49)} prakṛtānām ātmanepadānām etvam bhavati iti . \\
\noindent \textbf{(P\_2,4.49)} yuvoḥ anākau iti . \\
\noindent \textbf{(P\_2,4.49)} vakṣyati etat . \\
\noindent \textbf{(P\_2,4.49)} siddham tu yuvoḥ ananunāsikatvāt iti . \\
\noindent \textbf{(P\_2,4.54.1)} kim ayam kaśādiḥ āhośvit khayādiḥ . \\
\noindent \textbf{(P\_2,4.54.1)} cakṣiṅaḥ kśāñkhyāñau . \\
\noindent \textbf{(P\_2,4.54.1)} cakiṅaḥ khyāñ kaśādiḥ khayādiḥ ca . \\
\noindent \textbf{(P\_2,4.54.1)} khaśādiḥ vā . \\
\noindent \textbf{(P\_2,4.54.1)} atha vā khaśādiḥ bhaviṣyati . \\
\noindent \textbf{(P\_2,4.54.1)} kena idānīm kaśādiḥ bhaviṣyati . \\
\noindent \textbf{(P\_2,4.54.1)} cartvena . \\
\noindent \textbf{(P\_2,4.54.1)} atha khayādiḥ katham . \\
\noindent \textbf{(P\_2,4.54.1)} asiddhe śasya yavacanam vibhāṣā . \\
\noindent \textbf{(P\_2,4.54.1)} asiddhe śasya vibhāṣā yatvam vaktavyam . \\
\noindent \textbf{(P\_2,4.54.1)} kim prayojanam . \\
\noindent \textbf{(P\_2,4.54.1)} prayojanam sauprakhye vuñvidhiḥ . \\
\noindent \textbf{(P\_2,4.54.1)} sauprakhyaḥ iti yopadhalakṣaṇaḥ vuñvidhiḥ na bhavati . \\
\noindent \textbf{(P\_2,4.54.1)} sauprakhyīyaḥ . \\
\noindent \textbf{(P\_2,4.54.1)} vṛddhāt chaḥ bhavati . \\
\noindent \textbf{(P\_2,4.54.1)} niṣṭhānatvam ākhyāte . \\
\noindent \textbf{(P\_2,4.54.1)} ākhyātaḥ iti niṣṭhānatvam na bhavati . \\
\noindent \textbf{(P\_2,4.54.1)} ruvidhiḥ puṅkhyāne . \\
\noindent \textbf{(P\_2,4.54.1)} puṅkhyānam iti ruvidhiḥ na bhavati . \\
\noindent \textbf{(P\_2,4.54.1)} ṇatvam paryākhyāte . \\
\noindent \textbf{(P\_2,4.54.1)} paryākhyānam iti ṇatvam na bhavati . \\
\noindent \textbf{(P\_2,4.54.1)} sasthānatvam namaḥ khyātre . \\
\noindent \textbf{(P\_2,4.54.1)} namaḥ khyātre iti sasthānatvam na bhavati. \\
\noindent \textbf{(P\_2,4.54.2)} varjane pratiṣedhaḥ . varjane pratiṣedhaḥ vaktavyaḥ . \\
\noindent \textbf{(P\_2,4.54.2)} avasañcakṣyāḥ parisañcakṣyāḥ . \\
\noindent \textbf{(P\_2,4.54.2)} asanayoḥ ca . \\
\noindent \textbf{(P\_2,4.54.2)} asanayoḥ ca pratiṣedhaḥ vaktavyaḥ . \\
\noindent \textbf{(P\_2,4.54.2)} nṛcakṣāḥ rakṣaḥ . \\
\noindent \textbf{(P\_2,4.54.2)} vicakṣaṇaḥ iti . \\
\noindent \textbf{(P\_2,4.54.2)} bahulam taṇi . \\
\noindent \textbf{(P\_2,4.54.2)} bahulam taṇi iti vaktavyam . \\
\noindent \textbf{(P\_2,4.54.2)} kim idam taṇi iti . \\
\noindent \textbf{(P\_2,4.54.2)} sañjñāchandasoḥ grahaṇam . \\
\noindent \textbf{(P\_2,4.54.2)} kim prayojanam . \\
\noindent \textbf{(P\_2,4.54.2)} annavadhakagātravicakṣaṇājirādyartham . \\
\noindent \textbf{(P\_2,4.54.2)} anna . \\
\noindent \textbf{(P\_2,4.54.2)} annam . \\
\noindent \textbf{(P\_2,4.54.2)} vadhaka . \\
\noindent \textbf{(P\_2,4.54.2)} vadhakam . \\
\noindent \textbf{(P\_2,4.54.2)} gātra . \\
\noindent \textbf{(P\_2,4.54.2)} gātram paśya . \\
\noindent \textbf{(P\_2,4.54.2)} vicakṣaṇa . \\
\noindent \textbf{(P\_2,4.54.2)} vicakṣaṇaḥ . \\
\noindent \textbf{(P\_2,4.54.2)} ajira . \\
\noindent \textbf{(P\_2,4.54.2)} ajire tiṣṭhati . \\
\noindent \textbf{(P\_2,4.56)} ghañapoḥ pratiṣedhe kyapaḥ upasaṅkhyānam . \\
\noindent \textbf{(P\_2,4.56)} ghañapoḥ pratiṣedhe kyapaḥ upasaṅkhyānam kartavyam . \\
\noindent \textbf{(P\_2,4.56)} iha api yathā syāt . \\
\noindent \textbf{(P\_2,4.56)} samajanam samajyā iti . \\
\noindent \textbf{(P\_2,4.56)} tat tarhi vaktavyam . \\
\noindent \textbf{(P\_2,4.56)} na vaktavyam . \\
\noindent \textbf{(P\_2,4.56)} api iti eva bhaviṣyati . \\
\noindent \textbf{(P\_2,4.56)} katham . \\
\noindent \textbf{(P\_2,4.56)} api iti na idam pratyayagrahaṇam . \\
\noindent \textbf{(P\_2,4.56)} kim tarhi . \\
\noindent \textbf{(P\_2,4.56)} pratyāhāragrahaṇam . \\
\noindent \textbf{(P\_2,4.56)} kva sanniviṣṭānām pratyāhāraḥ . \\
\noindent \textbf{(P\_2,4.56)} apaḥ akātāt prabhṛti ā kyapaḥ pakārāt . \\
\noindent \textbf{(P\_2,4.56)} yadi pratyāhāragrahaṇam saṃvītiḥ na sidhyati . \\
\noindent \textbf{(P\_2,4.56)} evam tarhi na arthaḥ uapsaṅkhyānena na api ghañnapoḥ pratiṣedhena . \\
\noindent \textbf{(P\_2,4.56)} idam asti . \\
\noindent \textbf{(P\_2,4.56)} cakṣiṅaḥ khyāñ vā liṭi iti . \\
\noindent \textbf{(P\_2,4.56)} tataḥ vakṣyāmi . \\
\noindent \textbf{(P\_2,4.56)} ajeḥ vī bhavati vā vyavasthitavibhāṣā ca iti . \\
\noindent \textbf{(P\_2,4.56)} tena iha ca bhaviṣyati : pravetā pravetum pravītaḥ rathaḥ , saṃvītiḥ iti . \\
\noindent \textbf{(P\_2,4.56)} iha ca na bhaviṣyati : samājaḥ , udājaḥ , samajaḥ , udajaḥ , samajanam udajanam , samajyā iti . \\
\noindent \textbf{(P\_2,4.56)} tatra ayam api arthaḥ . \\
\noindent \textbf{(P\_2,4.56)} idam api siddham bhavati : prājitā iti . \\
\noindent \textbf{(P\_2,4.56)} kim ca bhoḥ iṣyate etat rūpam . \\
\noindent \textbf{(P\_2,4.56)} bāḍham iṣyate . \\
\noindent \textbf{(P\_2,4.56)} evam hi kaḥ cit vaiyākaraṇaḥ āha . \\
\noindent \textbf{(P\_2,4.56)} kaḥ asya rathasya pravetā iti . \\
\noindent \textbf{(P\_2,4.56)} sūtaḥ āha . \\
\noindent \textbf{(P\_2,4.56)} āyuṣman aham prājitā iti . \\
\noindent \textbf{(P\_2,4.56)} vaiyākaraṇaḥ āha . \\
\noindent \textbf{(P\_2,4.56)} apaśabdaḥ iti . \\
\noindent \textbf{(P\_2,4.56)} sūtaḥ āha . \\
\noindent \textbf{(P\_2,4.56)} prāpitjñaḥ devānām priyaḥ na tu iṣṭajñaḥ . \\
\noindent \textbf{(P\_2,4.56)} iṣyate etat rūpam iti . \\
\noindent \textbf{(P\_2,4.56)} vaiyākaraṇaḥ āha . \\
\noindent \textbf{(P\_2,4.56)} āho khalu anena durutena bādhyāmahe iti . \\
\noindent \textbf{(P\_2,4.56)} sūtaḥ āha . \\
\noindent \textbf{(P\_2,4.56)} na khalu veñaḥ sūtaḥ . \\
\noindent \textbf{(P\_2,4.56)} suvateḥ eva sūtaḥ . \\
\noindent \textbf{(P\_2,4.56)} yadi suvateḥ kutsā prayoktavyā . \\
\noindent \textbf{(P\_2,4.56)} duḥsūtena iti vaktavyam . \\
\noindent \textbf{(P\_2,4.56)} na tarhi idānīm idam vā yau iti vaktavyam . \\
\noindent \textbf{(P\_2,4.56)} vaktavyam ca . \\
\noindent \textbf{(P\_2,4.56)} kim prayojanam . \\
\noindent \textbf{(P\_2,4.56)} na iyam vibhāṣā . \\
\noindent \textbf{(P\_2,4.56)} kim tarhi . \\
\noindent \textbf{(P\_2,4.56)} ādeśaḥ ayam vidhīyate . \\
\noindent \textbf{(P\_2,4.56)} vā iti ayam ādeśaḥ bhavati ajeḥ yau parataḥ . \\
\noindent \textbf{(P\_2,4.56)} vāyuḥ iti . \\
\noindent \textbf{(P\_2,4.58)} aṇiñoḥ luki tadrājāt yuvapratyayasya upasaṅkhyānam . aṇiñoḥ luki tadrājāt yuvapratyayasya upasaṅkhyānam kartavyam . \\
\noindent \textbf{(P\_2,4.58)} baudhiḥ pitā baudhīḥ putraḥ audumbariḥ pitā audumbariḥ putraḥ . \\
\noindent \textbf{(P\_2,4.58)} aparaḥ āha : aṇiñoḥ luki kṣatriyagotramātrāt yuvapratyayasya upasaṅkhyānam kartavyam iti . \\
\noindent \textbf{(P\_2,4.58)} jābāliḥ pitā jābāliḥ putraḥ . \\
\noindent \textbf{(P\_2,4.58)} aparaḥ āha . \\
\noindent \textbf{(P\_2,4.58)} abrāhmaṇagotramātrāt yuvapratyayasya upasaṅkhyānam kartavyam iti . \\
\noindent \textbf{(P\_2,4.58)} kim prayojanam . \\
\noindent \textbf{(P\_2,4.58)} idam api siddham bhavati . \\
\noindent \textbf{(P\_2,4.58)} bhāṇḍijaṅghiḥ pitā bhāṇḍijaṅghiḥ putraḥ kārṇakharakiḥ pitā kārṇakharakiḥ putraḥ . \\
\noindent \textbf{(P\_2,4.62)} tadrājādīnām luki samāsabahutve pratiṣedhaḥ . \\
\noindent \textbf{(P\_2,4.62)} tadrājādīnām luki samāsabahutve pratiṣedhaḥ . \\
\noindent \textbf{(P\_2,4.62)} priyaḥ āṅgaḥ eṣām te ime priyāṅgāḥ . \\
\noindent \textbf{(P\_2,4.62)} priyaḥ vāṅgaḥ eṣām te ime priyavāṅgāḥ iti . \\
\noindent \textbf{(P\_2,4.62)} kim ucyate samāsabahutve pratiṣedhaḥ iti yāvatā tena eva cet kṛtam bahutvam iti ucyate na ca atra tena eva kṛtam bahutvam . \\
\noindent \textbf{(P\_2,4.62)} bhavati vai kim cit ācāryāḥ kriyamāṇam api codayanti . \\
\noindent \textbf{(P\_2,4.62)} tat vā kartavyam tena eva cet bahutvam iti samāsabahutve vā pratiṣedhaḥ vaktavyaḥ iti . \\
\noindent \textbf{(P\_2,4.62)} abahutve ca lugvacanam . \\
\noindent \textbf{(P\_2,4.62)} abahutve ca luk vaktavyaḥ . \\
\noindent \textbf{(P\_2,4.62)} atikrāntaḥ aṅgān atyaṅgaḥ . \\
\noindent \textbf{(P\_2,4.62)} bahuvacane parataḥ yaḥ tadrājaḥ iti evam kṛtvā codyate . \\
\noindent \textbf{(P\_2,4.62)} atha kimartham punaḥ idam na bahuvacane iti eva siddham . \\
\noindent \textbf{(P\_2,4.62)} na sidhyati . \\
\noindent \textbf{(P\_2,4.62)} bahuvacane iti ucyate na ca atra bahuvacanam paśyāmaḥ . \\
\noindent \textbf{(P\_2,4.62)} pratyayalakṣaṇena bhaviṣyati . \\
\noindent \textbf{(P\_2,4.62)} na lumatā tasmin iti pratyayalakṣaṇasya pratiṣedhaḥ . \\
\noindent \textbf{(P\_2,4.62)} na lumatā āṅgasya iti vakṣyāmi . \\
\noindent \textbf{(P\_2,4.62)} na evam śakyam . \\
\noindent \textbf{(P\_2,4.62)} iha hi doṣaḥ syāt . \\
\noindent \textbf{(P\_2,4.62)} pañcabhiḥ gārgībhiḥ krītaḥ paṭaḥ pañcagārgyaḥ daśagārgyaḥ iti . \\
\noindent \textbf{(P\_2,4.62)} dvandve abahuṣu lugvacanam . \\
\noindent \textbf{(P\_2,4.62)} dvandve abahuṣu luk vaktavyaḥ . \\
\noindent \textbf{(P\_2,4.62)} gargavatsavājāḥ iti . \\
\noindent \textbf{(P\_2,4.62)} iha ca luk vaktavyaḥ . \\
\noindent \textbf{(P\_2,4.62)} gargebhyaḥ āgatam gargarūpyam gargamayam iti . \\
\noindent \textbf{(P\_2,4.62)} iha ca atrayaḥ iti udāttanivṛttisvaraḥ prāpnoti . \\
\noindent \textbf{(P\_2,4.62)} siddham tu pratyayārthabahutve lugvacanāt . \\
\noindent \textbf{(P\_2,4.62)} siddham etat . \\
\noindent \textbf{(P\_2,4.62)} katham . \\
\noindent \textbf{(P\_2,4.62)} pratyayārthabahutve luk vaktavyaḥ . \\
\noindent \textbf{(P\_2,4.62)} yadi pratyayārthabahutve luk ucyate astriyām iti vaktavyam . \\
\noindent \textbf{(P\_2,4.62)} iha mā bhūt : āṅgyaḥ striyaḥ , vāṅgyaḥ striyaḥ iti . \\
\noindent \textbf{(P\_2,4.62)} yasya punaḥ bahuvacane parataḥ luk ucyate tasya īkāreṇa vyavahitatvāt na bhaviṣyati . \\
\noindent \textbf{(P\_2,4.62)} yasya api bahuvacane parataḥ luk ucyate tena api astriyām iti vaktavyam āmbaṣṭhyāḥ striyaḥ sauvīryāḥ striyaḥ iti evamartham . \\
\noindent \textbf{(P\_2,4.62)} atra api cāpā vyavadhānam . \\
\noindent \textbf{(P\_2,4.62)} ekādeśe kṛte na asti vyavhadānam . \\
\noindent \textbf{(P\_2,4.62)} ekādeśaḥ pūvavidhau sthānivat bhavati iti sthānivadbhāvāt vyavadhānam eva . \\
\noindent \textbf{(P\_2,4.62)} dvandve abahuṣu lugvacanam . \\
\noindent \textbf{(P\_2,4.62)} dvandve abahuṣu luk vaktavyaḥ . \\
\noindent \textbf{(P\_2,4.62)} gargavatsavājāḥ iti . \\
\noindent \textbf{(P\_2,4.62)} gotrasya bahuṣu lopinaḥ bahuvacanāntasya pravṛttau dvyekayoḥ aluk . \\
\noindent \textbf{(P\_2,4.62)} gotrasya bahuṣu lopinaḥ bahuvacanāntasya pravṛttau dvyekayoḥ aluk vaktavyaḥ . \\
\noindent \textbf{(P\_2,4.62)} bidānām apatyam māṇavakaḥ baidaḥ baidau . \\
\noindent \textbf{(P\_2,4.62)} kimartham idam na aci iti eva aluk siddhaḥ . \\
\noindent \textbf{(P\_2,4.62)} aci iti ucyate . \\
\noindent \textbf{(P\_2,4.62)} na ca atra ajādim paśyāmaḥ . \\
\noindent \textbf{(P\_2,4.62)} pratyayalakṣaṇena . \\
\noindent \textbf{(P\_2,4.62)} varṇāśraye na asti pratyayalakṣaṇam . \\
\noindent \textbf{(P\_2,4.62)} ekavacanadvivacanāntasya pravṛttau bahuṣu lopaḥ yūni . \\
\noindent \textbf{(P\_2,4.62)} ekavacanadvivacanāntasya pravṛttau bahuṣu lopaḥ yūni vaktavyaḥ . \\
\noindent \textbf{(P\_2,4.62)} baidasya apatyam bahavaḥ māṇavakāḥ bidāḥ baidayoḥ vā bidāḥ . \\
\noindent \textbf{(P\_2,4.62)} añ yaḥ bahuṣu yañ yaḥ bahuṣu iti ucyamānaḥ luk na prāpnoti . \\
\noindent \textbf{(P\_2,4.62)} mā bhūt evam . \\
\noindent \textbf{(P\_2,4.62)} añantam yat bahuṣu yañantam yat bahuṣu iti evam bhaviṣyati . \\
\noindent \textbf{(P\_2,4.62)} na evam śakyam . \\
\noindent \textbf{(P\_2,4.62)} iha hi doṣaḥ syāt . \\
\noindent \textbf{(P\_2,4.62)} kāśyapapratikṛtayaḥ kāśyapāḥ iti . \\
\noindent \textbf{(P\_2,4.62)} na vā sarveṣām dvandve bahvarthatvāt . \\
\noindent \textbf{(P\_2,4.62)} na vā eṣaḥ doṣaḥ . \\
\noindent \textbf{(P\_2,4.62)} kim kāraṇam . \\
\noindent \textbf{(P\_2,4.62)} sarveṣām dvandve bahvarthatvāt . \\
\noindent \textbf{(P\_2,4.62)} sarvāṇi dvandve bahvarthāni . \\
\noindent \textbf{(P\_2,4.62)} katham . \\
\noindent \textbf{(P\_2,4.62)} yugapat adhikaraṇavivakṣāyām dvandvaḥ bhavati . \\
\noindent \textbf{(P\_2,4.62)} tataḥ ayam āha yasya bahuvacane parataḥ luk . \\
\noindent \textbf{(P\_2,4.62)} yadi sarvāṇi dvandve bahvarthāni aham api idam acodyam codye . \\
\noindent \textbf{(P\_2,4.62)} dvandve abahuṣu lugvacanam iti . \\
\noindent \textbf{(P\_2,4.62)} mama api sarvatra bahuvacanam param bhavati . \\
\noindent \textbf{(P\_2,4.62)} luke kṛte na prāpnoti . \\
\noindent \textbf{(P\_2,4.62)} pratyayalakṣaṇena . \\
\noindent \textbf{(P\_2,4.62)} na lumatā tasmin iti pratyayalakṣaṇasya pratiṣedhaḥ . \\
\noindent \textbf{(P\_2,4.62)} na lumatā āṅgasya iti vakṣyāmi . \\
\noindent \textbf{(P\_2,4.62)} nanu ca uktam na evam śakyam . \\
\noindent \textbf{(P\_2,4.62)} iha hi doṣaḥ syāt . \\
\noindent \textbf{(P\_2,4.62)} pañcabhiḥ gārgībhiḥ krītaḥ paṭaḥ pañcagārgyaḥ daśagārgyaḥ iti . \\
\noindent \textbf{(P\_2,4.62)} iṣtam eva etat saṅgṛhītam . \\
\noindent \textbf{(P\_2,4.62)} pañcagargaḥ daśagargaḥ iti eva bhavitavyam . \\
\noindent \textbf{(P\_2,4.62)} tathā idam aparam acodyam codye . \\
\noindent \textbf{(P\_2,4.62)} gargarūpyam gargamayam . \\
\noindent \textbf{(P\_2,4.62)} atra api bahuvacane iti eva siddham katham . \\
\noindent \textbf{(P\_2,4.62)} samarthāt taddhitaḥ utpadyate . \\
\noindent \textbf{(P\_2,4.62)} sāmarthyam ca subantena . \\
\noindent \textbf{(P\_2,4.62)} tataḥ ayam āha yasya prtayayārthabahutve luk . \\
\noindent \textbf{(P\_2,4.62)} yadi samarthāt taddhitaḥ utpadyate aham api idam acodyam codye . \\
\noindent \textbf{(P\_2,4.62)} gotrasya bahuṣu lopinaḥ bahuvacanāntasya pravṛttau dvyekayoḥ aluk iti . \\
\noindent \textbf{(P\_2,4.62)} katham . \\
\noindent \textbf{(P\_2,4.62)} yasya api bahuvacane parataḥ luk tena api atra aluk vaktavyaḥ . \\
\noindent \textbf{(P\_2,4.62)} tasya api hi atra bahuvacanam param bhavati . \\
\noindent \textbf{(P\_2,4.62)} na vaktavyaḥ . \\
\noindent \textbf{(P\_2,4.62)} aci iti evam aluk siddhaḥ . \\
\noindent \textbf{(P\_2,4.62)} aci iti ucyate . \\
\noindent \textbf{(P\_2,4.62)} na ca atra ajādim paśyāmaḥ . \\
\noindent \textbf{(P\_2,4.62)} nanu ca uktam pratyayalakṣaṇena . \\
\noindent \textbf{(P\_2,4.62)} varṇāśraye na asti pratyayalakṣaṇam iti . \\
\noindent \textbf{(P\_2,4.62)} yadi vā kāni cit varṇāśrayāṇi api pratyayalakṣaṇena bhavanti tathā idam api bhaviṣyati . \\
\noindent \textbf{(P\_2,4.62)} atha vā aviśeṣeṇa alukam uktvā hali na iti vakṣyāmi . \\
\noindent \textbf{(P\_2,4.62)} yadi aviśeṣeṇa alukam uktvā hali na iti ucyate bidānām apatyam bahavaḥ māṇavakāḥ bidāḥ atra api aluk prāpnoti . \\
\noindent \textbf{(P\_2,4.62)} astu . \\
\noindent \textbf{(P\_2,4.62)} punaḥ asya yuvabahutve vartamānasya luk bhaviṣyati . \\
\noindent \textbf{(P\_2,4.62)} punaḥ aluk kasmāt na bhavati . \\
\noindent \textbf{(P\_2,4.62)} samarthānām prathamasya gotrapratyayāntasya aluk ucyate na ca tat samarthānām prathamam gotrapratyayāntam . \\
\noindent \textbf{(P\_2,4.62)} kim tarhi . \\
\noindent \textbf{(P\_2,4.62)} dvitīyam artham upasaṅkrāntam . \\
\noindent \textbf{(P\_2,4.62)} avaśyam ca etat evam vijñeyam atribharadvājikā vasiṣṭhakaśyapikā bhṛgvaṅgirasikā kutsakuśikā iti evamartham . \\
\noindent \textbf{(P\_2,4.62)} gargabhārgavikāgrahaṇam vā kriyate . \\
\noindent \textbf{(P\_2,4.62)} tat niyamārtham bhaviṣyati . \\
\noindent \textbf{(P\_2,4.62)} etasya eva dvitīyam artham upasaṅkrāntasya aluk bhavati na anyasya iti . \\
\noindent \textbf{(P\_2,4.62)} yat api ucyate ekavacanadvivacanāntasya pravṛttau bahuṣu lopaḥ yūni vaktavyaḥ iti . \\
\noindent \textbf{(P\_2,4.62)} mā bhūt evam añ yaḥ bahuṣu yañ yaḥ bahuṣu iti . \\
\noindent \textbf{(P\_2,4.62)} añantam yatbahuṣu yañantam yat bahuṣu iti evam bhaviṣyati . \\
\noindent \textbf{(P\_2,4.62)} nanu ca uktam na evam śakyam . \\
\noindent \textbf{(P\_2,4.62)} iha hi doṣaḥ syāt . \\
\noindent \textbf{(P\_2,4.62)} kāśyapapratikṛtayaḥ kāśyapāḥ iti . \\
\noindent \textbf{(P\_2,4.62)} na eṣaḥ doṣaḥ . \\
\noindent \textbf{(P\_2,4.62)} laukikasya tatra gotrasya grahaṇam . \\
\noindent \textbf{(P\_2,4.62)} na ca etat laukikam gotram . \\
\noindent \textbf{(P\_2,4.62)} yasya bahuvacane parataḥ luk samāsabahutve tena pratiṣedhaḥ vaktavyaḥ tena eva cet kṛtam bahutvam iti vā vaktavyam . \\
\noindent \textbf{(P\_2,4.62)} yasya pratyayārthabahutve luk tena astriyām iti vaktavyam . \\
\noindent \textbf{(P\_2,4.62)} yasya bahuvacane parataḥ luk tasya ayam adhikaḥ doṣaḥ atraḥ iti udāttanivṛttisvaraḥ prāpnoti . \\
\noindent \textbf{(P\_2,4.62)} tasmāt pratyayārthabahutve luk iti eṣaḥ pakṣaḥ jyāyān . \\
\noindent \textbf{(P\_2,4.62)} atha iha katham bhavitavyam . \\
\noindent \textbf{(P\_2,4.62)} gārgī ca bātsyaḥ ca iti . \\
\noindent \textbf{(P\_2,4.62)} yadi tāvat astri vidhinā āśrīyate asti atra astrī iti kṛtvā bhavitavyam lukā . \\
\noindent \textbf{(P\_2,4.62)} atha strī pratiṣedhena āśrīyate asti atra strī iti kṛtvā bhavitavyam pratiṣedhena . \\
\noindent \textbf{(P\_2,4.62)} kim punaḥ atra arthasatattvam . \\
\noindent \textbf{(P\_2,4.62)} devāḥ etat jñātum arhanti . \\
\noindent \textbf{(P\_2,4.62)} atha yaḥ lopyalopinām samāsaḥ tatra katham bhavitavyam . \\
\noindent \textbf{(P\_2,4.62)} ubhayam hi dṛśyate . \\
\noindent \textbf{(P\_2,4.62)} śaradvat śunakadarbhāt bhruguvat sāgrāyaṇeṣu na udāttasvaritodayam agārgyakāśyapagālavānām iti . \\
\noindent \textbf{(P\_2,4.64)} yañādīnām ekadvayoḥ vā tatpuruṣe ṣaṣṭhyāḥ upasaṅkhyānam . \\
\noindent \textbf{(P\_2,4.64)} yañādīnām ekadvayoḥ vā tatpuruṣe ṣaṣṭhyāḥ upasaṅkhyānam kartavyam . \\
\noindent \textbf{(P\_2,4.64)} gārgyasya kulam gārgyakulam gargakulam vā . \\
\noindent \textbf{(P\_2,4.64)} gārgyayoḥ kulam gārgyakulam gargakulam vā . \\
\noindent \textbf{(P\_2,4.64)} baidasya kulam baidakulam bidakulam vā . \\
\noindent \textbf{(P\_2,4.64)} baidayoḥ kulam baidakulam bidakulam vā . \\
\noindent \textbf{(P\_2,4.64)} yañādīnām iti kimartham . \\
\noindent \textbf{(P\_2,4.64)} āṅgasya kulam āṅgakulam. āṅgayoḥ kulam āṅgakulam . \\
\noindent \textbf{(P\_2,4.64)} ekadvayoḥ iti kimartham . \\
\noindent \textbf{(P\_2,4.64)} gargāṇām kulam gargakulam . \\
\noindent \textbf{(P\_2,4.64)} tatpuruṣe iti kimartham. gārgyasya samīpam upagārgyam . \\
\noindent \textbf{(P\_2,4.64)} ṣaṣṭhyāḥ iti kim . \\
\noindent \textbf{(P\_2,4.64)} śobhanagārgyaḥ paramagārgyaḥ . \\
\noindent \textbf{(P\_2,4.66)} kim ayam samuccayaḥ . \\
\noindent \textbf{(P\_2,4.66)} prākṣu bharateṣu ca iti . \\
\noindent \textbf{(P\_2,4.66)} āhosvit bharataviśeṣaṇam prāggrahaṇam . \\
\noindent \textbf{(P\_2,4.66)} prāñcaḥ ye bharatāḥ iti . \\
\noindent \textbf{(P\_2,4.66)} kim ca ataḥ . \\
\noindent \textbf{(P\_2,4.66)} yadi samuccayaḥ bharatagrahaṇam anarthakam . \\
\noindent \textbf{(P\_2,4.66)} na hi anyatra bharatā santi . \\
\noindent \textbf{(P\_2,4.66)} atha prāggrahaṇam bharataviśeṣaṇam prāggrahaṇam anarthakam . \\
\noindent \textbf{(P\_2,4.66)} na hi aprāñcaḥ bharatāḥ santi . \\
\noindent \textbf{(P\_2,4.66)} evam tarhi samuccayaḥ . \\
\noindent \textbf{(P\_2,4.66)} nanu ca uktam bharatagrahaṇam anarthakam . \\
\noindent \textbf{(P\_2,4.66)} na hi anyatra bharatā santi iti . \\
\noindent \textbf{(P\_2,4.66)} na anarthakam. jñāpakārtham . \\
\noindent \textbf{(P\_2,4.66)} kim jñāpyate . \\
\noindent \textbf{(P\_2,4.66)} etat jñāpayati ācāryaḥ anyatra prāggrahaṇe bharatagrahaṇam na bhavati iti . \\
\noindent \textbf{(P\_2,4.66)} kim etasya jñāpane prayojanam . \\
\noindent \textbf{(P\_2,4.66)} iñaḥ prācām bharatgrahaṇam na bhavati . \\
\noindent \textbf{(P\_2,4.66)} auddālakiḥ pitā auddālakāyanaḥ putraḥ iti . \\
\noindent \textbf{(P\_2,4.67)} gopavanādipratiṣedhaḥ prāk haritādibhyaḥ . \\
\noindent \textbf{(P\_2,4.67)} gopavanādipratiṣedhaḥ prāk haritādibhyaḥ draṣṭavyaḥ . \\
\noindent \textbf{(P\_2,4.67)} hāritaḥ hāritau bahuṣu haritāḥ . \\
\noindent \textbf{(P\_2,4.69)} kimartham advandve iti ucyate . \\
\noindent \textbf{(P\_2,4.69)} dvandve mā bhūt iti . \\
\noindent \textbf{(P\_2,4.69)} na etat asti prayojanam . \\
\noindent \textbf{(P\_2,4.69)} iṣyate eva dvandve : bhraṣṭakakapiṣṭhalāḥ bhrāṣṭakikāpiṣthalayaḥ iti . \\
\noindent \textbf{(P\_2,4.69)} ataḥ uttaram paṭhati advandve iti dvandvādhikāranivṛttyartham . \\
\noindent \textbf{(P\_2,4.69)} advandve iti ucyate dvandvādhikāranivṛttyartham . \\
\noindent \textbf{(P\_2,4.69)} dvandvādhikāraḥ nivartate . \\
\noindent \textbf{(P\_2,4.69)} tasmin nivṛtte aviśeṣeṇa dvandve ca advandve ca bhaviṣyati . \\
\noindent \textbf{(P\_2,4.70)} āgastyakauṇḍinyayoḥ prakṛtinipātanam . āgastyakauṇḍinyayoḥ prakṛtinipātanam kartavyam . \\
\noindent \textbf{(P\_2,4.70)} agastikauṇḍinac iti etau prakṛtyādeśau bhavataḥ iti vaktavyam . \\
\noindent \textbf{(P\_2,4.70)} kim prayojanam . \\
\noindent \textbf{(P\_2,4.70)} lukpratiṣedhe vṛddhyartham . lukpratiṣedhe vṛddhiḥ yathā syāt . \\
\noindent \textbf{(P\_2,4.70)} pratyayāntanipātane hi vṛddhyabhāvaḥ . \\
\noindent \textbf{(P\_2,4.70)} pratyayāntanipātane hi sati vṛddhyabhāvaḥ syāt . \\
\noindent \textbf{(P\_2,4.70)} āgastīyāḥ kauṇḍinyāḥ iti . \\
\noindent \textbf{(P\_2,4.70)} yadi prakrṭinipātanam kriyate kena idānīm pratyayasya lopaḥ bhaviṣyati . \\
\noindent \textbf{(P\_2,4.70)} adhikārāt pratyayalopaḥ . \\
\noindent \textbf{(P\_2,4.70)} adhikārāt pratyayalopaḥ bhaviṣyati . \\
\noindent \textbf{(P\_2,4.70)} tat tarhi prakṛtinipātanam kartavyam . \\
\noindent \textbf{(P\_2,4.70)} na kartavyam . \\
\noindent \textbf{(P\_2,4.70)} yogavibhāgaḥ kariṣyate . \\
\noindent \textbf{(P\_2,4.70)} āgastyakauṇḍinyayoḥ bahuṣu luk bhavati . \\
\noindent \textbf{(P\_2,4.70)} tataḥ agastikuṇḍinac iti etau prakṛtyādeśau bhavataḥ āgastyakauṇḍinyayoḥ iti . \\
\noindent \textbf{(P\_2,4.70)} evam api pratyayāntayoḥ eva prāpnoti . \\
\noindent \textbf{(P\_2,4.70)} pratyayāntāt hi bhavān ṣaṣṭhīm uccārayati . \\
\noindent \textbf{(P\_2,4.70)} āgastyakauṇḍinyayoḥ iti . \\
\noindent \textbf{(P\_2,4.70)} na eṣaḥ doṣaḥ . \\
\noindent \textbf{(P\_2,4.70)} yathā hi paribhāṣitam pratyayasya lukślulupaḥ bhavanti iti pratyayasya eva bhaviṣyati . \\
\noindent \textbf{(P\_2,4.70)} avaśiṣṭasya ādeśau bhaviṣyataḥ . \\
\noindent \textbf{(P\_2,4.74)} ūtaḥ aci .ūtaḥ aci iti vaktavyam . \\
\noindent \textbf{(P\_2,4.74)} iha mā bhūt . \\
\noindent \textbf{(P\_2,4.74)} sanīsrasaḥ danīdhvasaḥ iti . \\
\noindent \textbf{(P\_2,4.74)} atha ūtaḥ iti ucyamāne iha kasmāt na bhavati . \\
\noindent \textbf{(P\_2,4.74)} yoyūyaḥ rorūvaḥ . \\
\noindent \textbf{(P\_2,4.74)} vihitaviśeṣaṇam ūkārāntagrahaṇan . \\
\noindent \textbf{(P\_2,4.74)} ūkārāntāt yaḥ vihitaḥ iti . \\
\noindent \textbf{(P\_2,4.74)} tat tarhi vaktavyam . \\
\noindent \textbf{(P\_2,4.74)} na vaktavyam . \\
\noindent \textbf{(P\_2,4.74)} iṣṭam eva etat saṅgṛhītam . \\
\noindent \textbf{(P\_2,4.74)} sanīsraṃsaḥ danīdhvaṃsaḥ iti eva bhavitavyam . \\
\noindent \textbf{(P\_2,4.77)} gāpoḥ grahaṇe iṇpibatyoḥ grahaṇam . gāpoḥ grahaṇe iṇpibatyoḥ grahaṇam kartavyam . \\
\noindent \textbf{(P\_2,4.77)} iṇaḥ yaḥ gāśabdaḥ pibateḥ yaḥ pāśabdaḥ iti vaktavyam . \\
\noindent \textbf{(P\_2,4.77)} iha mā bhūt . \\
\noindent \textbf{(P\_2,4.77)} agāsīt naṭaḥ . \\
\noindent \textbf{(P\_2,4.77)} apāsīt dhanam iti . \\
\noindent \textbf{(P\_2,4.77)} tat tarhi vaktavyam . \\
\noindent \textbf{(P\_2,4.77)} na vaktavyam . \\
\noindent \textbf{(P\_2,4.77)} iṇaḥ grahaṇe tāvat vārtam . \\
\noindent \textbf{(P\_2,4.77)} nirdeśāt eva idam vyaktam lugvikaraṇasya grahaṇam iti . \\
\noindent \textbf{(P\_2,4.77)} pāgrahaṇe ca api vārtam . \\
\noindent \textbf{(P\_2,4.77)} vaktavyam eva etat sarvatra eva pāgrahaṇe alugvikaraṇasya grahaṇam iti . \\
\noindent \textbf{(P\_2,4.79)} tathāsoḥ ātmanepadavacanam . \\
\noindent \textbf{(P\_2,4.79)} tathāsoḥ ātmanepadasya grahaṇam kartavyam . \\
\noindent \textbf{(P\_2,4.79)} ātmanepadam yau tathāsau iti vaktavyam . \\
\noindent \textbf{(P\_2,4.79)} ekavacanagrahaṇam vā . \\
\noindent \textbf{(P\_2,4.79)} atha vā ekavacane ye tathāsī iti vaktavyam . \\
\noindent \textbf{(P\_2,4.79)} tat ca avaśyam anyatarat kartavyam . \\
\noindent \textbf{(P\_2,4.79)} avacane hi aniṣṭaprasaṅgaḥ . \\
\noindent \textbf{(P\_2,4.79)} anucyamāne hi etasmin aniṣṭam prasajyeta . \\
\noindent \textbf{(P\_2,4.79)} ataniṣṭa yūpam . \\
\noindent \textbf{(P\_2,4.79)} asaniṣṭa yūpam iti . \\
\noindent \textbf{(P\_2,4.79)} tat tarhi vaktavyam . \\
\noindent \textbf{(P\_2,4.79)} na vaktavyam . \\
\noindent \textbf{(P\_2,4.79)} yadi api tāvat ayam taśabdaḥ dṛṣṭāpacāraḥ asti ātmanepadam asti eva parasmaipadam asti ekavacanam asti bahuvacanam ayam khalu thāsśabdaḥ adṛṣṭāpacāraḥ ātmanepadam ekavacanam eva . \\
\noindent \textbf{(P\_2,4.79)} tasya asya kaḥ anyaḥ sahāyaḥ bhavitum arhati anyat ataḥ ātmanepadāt ekavacanāt ca . \\
\noindent \textbf{(P\_2,4.79)} tat yathā asyas goḥ dvitīyena arthaḥ iti . \\
\noindent \textbf{(P\_2,4.79)} gauḥ eva ānīyate na aśvaḥ na gardabhaḥ . \\
\noindent \textbf{(P\_2,4.81.1)} āmaḥ leḥ lope luṅloṭoḥ upasaṅkhyānam . \\
\noindent \textbf{(P\_2,4.81.1)} āmaḥ leḥ lope luṅloṭoḥ upasaṅkhyānam kartavyam : tām baijavāpayaḥ vidām akran . \\
\noindent \textbf{(P\_2,4.81.1)} atra bhavantaḥ vidām kurvantu. tat tarhi vaktavyam . \\
\noindent \textbf{(P\_2,4.81.1)} na vaktavyam . \\
\noindent \textbf{(P\_2,4.81.1)} ligrahaṇam nivartiṣyate . \\
\noindent \textbf{(P\_2,4.81.1)} yadi nivartate pratyayamātrasya prāpnoti . \\
\noindent \textbf{(P\_2,4.81.1)} iṣyate ca pratyayamātrasya . \\
\noindent \textbf{(P\_2,4.81.1)} ātaḥ ca iṣyate . \\
\noindent \textbf{(P\_2,4.81.1)} evam hi āha : kṛñ ca anuprayujyate liṭi iti . \\
\noindent \textbf{(P\_2,4.81.1)} yadi ca pratyayamātrasya luk bhavati tataḥ etat upapannam bhavati . \\
\noindent \textbf{(P\_2,4.81.1)} āmantebhyaḥ ṇalaḥ pratiṣedhaḥ . \\
\noindent \textbf{(P\_2,4.81.1)} āmantebhyaḥ ṇalaḥ pratiṣedhaḥ vaktavyaḥ . \\
\noindent \textbf{(P\_2,4.81.1)} śaśāma tatāma . \\
\noindent \textbf{(P\_2,4.81.1)} vṛddhau kṛtyāyām āmaḥ iti luk prāpnoti . \\
\noindent \textbf{(P\_2,4.81.1)} āmantebhyaḥ arthavadgrahaṇāt ṇalaḥ apratiṣedhaḥ . \\
\noindent \textbf{(P\_2,4.81.1)} āmantebhyaḥ ṇalaḥ apratiṣedhaḥ . \\
\noindent \textbf{(P\_2,4.81.1)} anarthakaḥ pratiṣedhaḥ apratiṣedhaḥ . \\
\noindent \textbf{(P\_2,4.81.1)} luk kasmāt na bhavati : śaśāma tatāma iti . \\
\noindent \textbf{(P\_2,4.81.1)} arthavataḥ āmśabdasya grahaṇam . \\
\noindent \textbf{(P\_2,4.81.1)} na ca eṣaḥ arthavān . \\
\noindent \textbf{(P\_2,4.81.1)} āmantebhyaḥ arthavadgrahaṇāt ṇalaḥ apratiṣedhaḥ iti cet amaḥ pratiṣedhaḥ . āmantebhyaḥ arthavadgrahaṇāt ṇalaḥ apratiṣedhaḥ iti cet amaḥ pratiṣedhaḥ vaktavyaḥ . \\
\noindent \textbf{(P\_2,4.81.1)} āma . \\
\noindent \textbf{(P\_2,4.81.1)} uktam vā . \\
\noindent \textbf{(P\_2,4.81.1)} kim uktam . \\
\noindent \textbf{(P\_2,4.81.1)} sannipātalakṣaṇaḥ vidhiḥ animittam tadvighātasya iti . \\
\noindent \textbf{(P\_2,4.81.2)} kim punaḥ luk ādeśānām apavādaḥ āhosvit kṛteṣu ādeśeṣu bhavati . \\
\noindent \textbf{(P\_2,4.81.2)} luk ādeśāpavādaḥ . luk ādeśānām apavādaḥ . \\
\noindent \textbf{(P\_2,4.81.2)} tiṅkṛtābhāvaḥ tu . \\
\noindent \textbf{(P\_2,4.81.2)} tiṅkṛtasya tu abhāvaḥ . \\
\noindent \textbf{(P\_2,4.81.2)} kasya . \\
\noindent \textbf{(P\_2,4.81.2)} padatvasya . \\
\noindent \textbf{(P\_2,4.81.2)} subantapadatvāt siddham . \\
\noindent \textbf{(P\_2,4.81.2)} subantam padam iti padasañjñā bhaviṣyati . \\
\noindent \textbf{(P\_2,4.81.2)} katham svādyutpattiḥ . \\
\noindent \textbf{(P\_2,4.81.2)} lakārasya kṛttvāt prātipadikatvam tadāśrayam pratyayavidhānam . lakāraḥ kṛt . \\
\noindent \textbf{(P\_2,4.81.2)} kṛt prātipadikam iti prātipadikasañjñā . \\
\noindent \textbf{(P\_2,4.81.2)} tadāśrayam pratyayavidhānam . \\
\noindent \textbf{(P\_2,4.81.2)} prātipadikāśrayatvāt svādyutpattiḥ bhaviṣyati . \\
\noindent \textbf{(P\_2,4.81.2)} supaḥ śravaṇam prāpnoti . \\
\noindent \textbf{(P\_2,4.81.2)} avyayāt iti luk bhaviṣyati . \\
\noindent \textbf{(P\_2,4.81.2)} katham avyayatvam . \\
\noindent \textbf{(P\_2,4.81.2)} avyayatvam makārāntatvāt . \\
\noindent \textbf{(P\_2,4.81.2)} kṛdantam māntam avyayasañjñam bhavati iti avyayasañjñā bhaviṣyati . \\
\noindent \textbf{(P\_2,4.81.2)} svaraḥ katham . \\
\noindent \textbf{(P\_2,4.81.2)} yat prakārayām cakāra . \\
\noindent \textbf{(P\_2,4.81.2)} svaraḥ kṛdantaprakṛtisvaratvāt . \\
\noindent \textbf{(P\_2,4.81.2)} kṛdantam uttarapadam prakṛtisvaram bhavati iti eṣaḥ svaraḥ bhaviṣyati . \\
\noindent \textbf{(P\_2,4.81.2)} tathā ca nighātānighātasiddhiḥ . \\
\noindent \textbf{(P\_2,4.81.2)} tathā ca nighātānighātasiddhiḥ bhavati . \\
\noindent \textbf{(P\_2,4.81.2)} cakṣuṣkāmam yājayām cakāra . \\
\noindent \textbf{(P\_2,4.81.2)} tiṅ atiṅaḥ iti tasya ca anighātaḥ . \\
\noindent \textbf{(P\_2,4.81.2)} tasmāt ca nighātaḥ siddhaḥ bhavati . \\
\noindent \textbf{(P\_2,4.81.2)} nañā tu samāsaprasaṅgaḥ . \\
\noindent \textbf{(P\_2,4.81.2)} nañā tu samāsaḥ prāpnoti . \\
\noindent \textbf{(P\_2,4.81.2)} na kārayam na hārayām . \\
\noindent \textbf{(P\_2,4.81.2)} nañ subantena saha samasyate iti samāsaḥ prāpnoti . \\
\noindent \textbf{(P\_2,4.81.2)} uktam vā . \\
\noindent \textbf{(P\_2,4.81.2)} kim uktam . \\
\noindent \textbf{(P\_2,4.81.2)} asāmarthyāt iti . \\
\noindent \textbf{(P\_2,4.81.2)} na atra nañaḥ āmantena sāmarthyam . \\
\noindent \textbf{(P\_2,4.81.2)} kena tarhi . \\
\noindent \textbf{(P\_2,4.81.2)} tiṅantena . \\
\noindent \textbf{(P\_2,4.81.2)} na cakāra kārayām . \\
\noindent \textbf{(P\_2,4.81.2)} na cakāra hārayām iti . \\
\noindent \textbf{(P\_2,4.82)} avyayāt āpaḥ lugvacanānarthakyam liṅgābhāvāt .avyayāt āpaḥ lugvacanam anarthakam . \\
\noindent \textbf{(P\_2,4.82)} kim kāraṇam . \\
\noindent \textbf{(P\_2,4.82)} liṅgābhāvāt . \\
\noindent \textbf{(P\_2,4.82)} aliṅgam avyayam . \\
\noindent \textbf{(P\_2,4.82)} kim idam bhavān supaḥ lukam mṛṣyati āpaḥ lukam na mṛṣyati . \\
\noindent \textbf{(P\_2,4.82)} yathā eva hi aliṅgam avyayam evam asaṅkhyam api . \\
\noindent \textbf{(P\_2,4.82)} satyam etat . \\
\noindent \textbf{(P\_2,4.82)} pratyayalakṣaṇam ācāryaḥ prārthayamānaḥ supaḥ lukam mṛṣyati . \\
\noindent \textbf{(P\_2,4.82)} āpaḥ punaḥ asya luki sati na kim cit api prayojanam asti . \\
\noindent \textbf{(P\_2,4.82)} ucyamāne api etasmin svādyutpattiḥ na prāpnoti . \\
\noindent \textbf{(P\_2,4.82)} kim kāraṇam . \\
\noindent \textbf{(P\_2,4.82)} ekatvādīnām abhāvāt . \\
\noindent \textbf{(P\_2,4.82)} ekatvādiṣu artheṣu svādayaḥ vidhīyante . \\
\noindent \textbf{(P\_2,4.82)} na ca eṣām ekatvādayaḥ santi . \\
\noindent \textbf{(P\_2,4.82)} aviśeṣeṇa utpadyante . \\
\noindent \textbf{(P\_2,4.82)} utpannānām niyamaḥ kriyate . \\
\noindent \textbf{(P\_2,4.82)} atha vā prakṛtān arthān apekṣya niyamaḥ . \\
\noindent \textbf{(P\_2,4.82)} ke ca prakṛtāḥ . \\
\noindent \textbf{(P\_2,4.82)} ekatvādayaḥ . \\
\noindent \textbf{(P\_2,4.82)} ekasmin ekavacanam na dvayoḥ na bahuṣu . \\
\noindent \textbf{(P\_2,4.82)} dvayoḥ eva dvivacanam na ekasmin na bahuṣu . \\
\noindent \textbf{(P\_2,4.82)} bahuṣu eva bahuvacanam na ekasmin na dvayoḥ iti . \\
\noindent \textbf{(P\_2,4.82)} atha vā ācāryapravṛttiḥ jñāpayati utpadyante avyayebhyaḥ svādayaḥ iti yat ayam avyayāt āpasupaḥ iti sublukam śāsti . \\
\noindent \textbf{(P\_2,4.83.1)} na avyayībhāvāt ataḥ iti yogavyavasānam . \\
\noindent \textbf{(P\_2,4.83.1)} na avyayībhāvāt ataḥ iti yogaḥ vyavaseyaḥ . \\
\noindent \textbf{(P\_2,4.83.1)} na avyayībhāvāt akārāntāt supaḥ luk bhavati . \\
\noindent \textbf{(P\_2,4.83.1)} tataḥ am tu apañcamyāḥ iti . \\
\noindent \textbf{(P\_2,4.83.1)} kimarthaḥ yogavibhāgaḥ . \\
\noindent \textbf{(P\_2,4.83.1)} pañcamyāḥ ampratiṣedhārtham . pañcamyāḥ amaḥ pratiṣedhaḥ yathā syāt . \\
\noindent \textbf{(P\_2,4.83.1)} ekayoge hi ubhayoḥ pratiṣedhaḥ . \\
\noindent \textbf{(P\_2,4.83.1)} ekayoge hi sati ubhayoḥ pratiṣedhaḥ syāt amaḥ alukaḥ ca . \\
\noindent \textbf{(P\_2,4.83.1)} saḥ tarhi yogavibhāgaḥ kartavyaḥ . \\
\noindent \textbf{(P\_2,4.83.1)} na kartavyaḥ . \\
\noindent \textbf{(P\_2,4.83.1)} tuḥ niyāmakaḥ . \\
\noindent \textbf{(P\_2,4.83.1)} tuḥ kriyate . \\
\noindent \textbf{(P\_2,4.83.1)} sa niyāmakaḥ bhaviṣyati . \\
\noindent \textbf{(P\_2,4.83.1)} am eva apañcamyāḥ iti . \\
\noindent \textbf{(P\_2,4.83.2)} ami pañcamīpratiṣedhe apādānagrahaṇam . \\
\noindent \textbf{(P\_2,4.83.2)} ami pañcamīpratiṣedhe apādānagrahaṇam kartavyam . \\
\noindent \textbf{(P\_2,4.83.2)} apādānapañcamyāḥ iti vaktavyam . \\
\noindent \textbf{(P\_2,4.83.2)} kim prayojanam . \\
\noindent \textbf{(P\_2,4.83.2)} karmapravacanīyayukte apratiṣedhārtham . karmapravacanīyayukte mā bhūt . \\
\noindent \textbf{(P\_2,4.83.2)} āpāṭaliputram vṛṣṭaḥ devaḥ . \\
\noindent \textbf{(P\_2,4.83.2)} na vā uttarapadasya karmapravacanīyayogāt samāsāt pañcamyabhāvaḥ . \\
\noindent \textbf{(P\_2,4.83.2)} na vā vaktavyam . \\
\noindent \textbf{(P\_2,4.83.2)} kim kāraṇam . \\
\noindent \textbf{(P\_2,4.83.2)} uttarapadam atra karmapravacanīyayuktam . \\
\noindent \textbf{(P\_2,4.83.2)} uttarapadasya karmapravacanīyayogāt samāsāt pañcamī na bhaviṣyati . \\
\noindent \textbf{(P\_2,4.83.2)} yadā ca samāsaḥ karmapravacanīyayuktaḥ bhavati tadā pratiṣedhaḥ . \\
\noindent \textbf{(P\_2,4.83.2)} tat yathā . \\
\noindent \textbf{(P\_2,4.83.2)} ā upakumbhāt ā upamaṇikāt iti . \\
\noindent \textbf{(P\_2,4.84)} saptamyāḥ ṛddhinadīsamāsasaṅkhyāvayavebhyaḥ nityam . saptamyāḥ ṛddhinadīsamāsasaṅkhyāvayavebhyaḥ nityam iti vaktavyam . \\
\noindent \textbf{(P\_2,4.84)} ṛddhi . \\
\noindent \textbf{(P\_2,4.84)} sumaram sumagadham . \\
\noindent \textbf{(P\_2,4.84)} nadīsamāsaḥ . \\
\noindent \textbf{(P\_2,4.84)} unmattagaṅgam lohitagaṅgam . \\
\noindent \textbf{(P\_2,4.84)} saṅkhyāvayava . \\
\noindent \textbf{(P\_2,4.84)} ekaviṃsatibhāradvājam tripañcāśatgautamam . \\
\noindent \textbf{(P\_2,4.85.1)} ṭitām ṭeḥ evidheḥ luṭaḥ ḍāraurasaḥ pūrvavipratiṣiddham . \\
\noindent \textbf{(P\_2,4.85.1)} ṭitām ṭeḥ evidheḥ luṭaḥ ḍāraurasaḥ bhavanti pūrvavipratiṣedhena . \\
\noindent \textbf{(P\_2,4.85.1)} ṭeḥ etvasya avakāśaḥ pacate pacete pacante . \\
\noindent \textbf{(P\_2,4.85.1)} ḍāraurasām avakāśaḥ śvaḥ kartā śvaḥ kartārau śvaḥ kartāraḥ . \\
\noindent \textbf{(P\_2,4.85.1)} iha ubhayam prāpnoti . \\
\noindent \textbf{(P\_2,4.85.1)} śvaḥ adhetā śvaḥ adhyetārau śvaḥ adhyetāraḥ . \\
\noindent \textbf{(P\_2,4.85.1)} ḍāraurasaḥ bhavanti pūrvavipratiṣedhena . \\
\noindent \textbf{(P\_2,4.85.1)} saḥ tarhi pūrvavipratiṣedhaḥ vaktavyaḥ . \\
\noindent \textbf{(P\_2,4.85.1)} na vaktavyaḥ . \\
\noindent \textbf{(P\_2,4.85.1)} ātmanepadānām ca iti vacanāt siddham . \\
\noindent \textbf{(P\_2,4.85.1)} ātmanepadānām ca ḍāraurasaḥ bhavanti iti vaktavyam . \\
\noindent \textbf{(P\_2,4.85.1)} tat ca samasaṅkhyārtham . tat ca avaśyam ātmanepadagrahaṇam kartavyam samasaṅkhyārtham . \\
\noindent \textbf{(P\_2,4.85.1)} saṅkhyātanudeśaḥ yathā syāt . \\
\noindent \textbf{(P\_2,4.85.1)} akriyamāṇe hi ātmanepadagrahaṇe śaṭ sthāninaḥ trayaḥ ādeśāḥ . \\
\noindent \textbf{(P\_2,4.85.1)} vaiṣamyāt saṅkhyātānudeśaḥ na prāpnoti . \\
\noindent \textbf{(P\_2,4.85.1)} pūrvavipratiṣedhārthena tāvat na arthaḥ ātmanepadagrahaṇena . \\
\noindent \textbf{(P\_2,4.85.1)} idam iha sampradhāryam . \\
\noindent \textbf{(P\_2,4.85.1)} ḍāraurasaḥ kriyantām ṭeḥ etvam iti . \\
\noindent \textbf{(P\_2,4.85.1)} kim atra kartavyam . \\
\noindent \textbf{(P\_2,4.85.1)} paratvāt etvam . \\
\noindent \textbf{(P\_2,4.85.1)} nityāḥ ḍāraurasaḥ . \\
\noindent \textbf{(P\_2,4.85.1)} kṛte api etve prapnuvanti akṛte api prāpnuvanti . \\
\noindent \textbf{(P\_2,4.85.1)} ṭeḥ etvam api nityam . \\
\noindent \textbf{(P\_2,4.85.1)} kṛteṣu api ḍāraurassu prāpnoti akṛteṣu api prāpnoti . \\
\noindent \textbf{(P\_2,4.85.1)} anityam etvam . \\
\noindent \textbf{(P\_2,4.85.1)} anyasya kṛteṣu ḍāraurassu prāpnoti anyasya akṛteṣu . \\
\noindent \textbf{(P\_2,4.85.1)} śabdāntarasya ca prāpnuvan vidhiḥ anityaḥ bhavati . \\
\noindent \textbf{(P\_2,4.85.1)} ḍāraurasaḥ api anityāḥ . \\
\noindent \textbf{(P\_2,4.85.1)} anyasya kṛte etve prāpnuvanti anyasya akṛte . \\
\noindent \textbf{(P\_2,4.85.1)} śabdāntarasya prāpnuvantaḥ anityā bhavanti . \\
\noindent \textbf{(P\_2,4.85.1)} ubhayoḥ anityayoḥ paratvāt etvam . \\
\noindent \textbf{(P\_2,4.85.1)} etve kṛte punaḥprasaṅgavijñānāt ḍāraurasaḥ bhaviṣyanti . \\
\noindent \textbf{(P\_2,4.85.1)} samasaṅkhyārthena ca api na arthaḥ ātmanepadagrahaṇena . \\
\noindent \textbf{(P\_2,4.85.1)} sthāne antaramena vyavasthā bhaviṣyati . \\
\noindent \textbf{(P\_2,4.85.1)} kutaḥ āntaryam . \\
\noindent \textbf{(P\_2,4.85.1)} arthataḥ . \\
\noindent \textbf{(P\_2,4.85.1)} ekārthasya ekārthaḥ dvyarthasya dvyarthaḥ bahvarthasya bahvarthaḥ . \\
\noindent \textbf{(P\_2,4.85.1)} atha vā ādeśāḥ api ṣaṭ eva nirdiśyante . \\
\noindent \textbf{(P\_2,4.85.1)} katham . \\
\noindent \textbf{(P\_2,4.85.1)} ekaśeṣanirdeśāt . \\
\noindent \textbf{(P\_2,4.85.1)} ekaśeṣanirdeśaḥ ayam . \\
\noindent \textbf{(P\_2,4.85.1)} atha etasmin ekaśeṣanirdeśe sati kim ayam kṛtaikaśeṣāṇām dvandvaḥ . \\
\noindent \textbf{(P\_2,4.85.1)} ḍā ca ḍā ca ḍā rau ca rau ca rau raḥ ca raḥ ca raḥ. ḍā ca rau ca raḥ ca ḍāraurasaḥ iti . \\
\noindent \textbf{(P\_2,4.85.1)} āhosvit kṛtadvandvānām ekaśeṣaḥ . \\
\noindent \textbf{(P\_2,4.85.1)} ḍā ca rau ca raḥ ca ḍāraurasaḥ . \\
\noindent \textbf{(P\_2,4.85.1)} ḍāraurasaḥ ca ḍāraurasaḥ ca ḍāraurasaḥ iti . \\
\noindent \textbf{(P\_2,4.85.1)} kim ca ataḥ . \\
\noindent \textbf{(P\_2,4.85.1)} yadi kṛtaikaśeṣāṇām dvandvaḥ aniṣṭaḥ samasaṅkhyaḥ prāpnoti ekavacanadvivacanayoḥ ḍā prāpnoti . \\
\noindent \textbf{(P\_2,4.85.1)} bahuvacanaikavacanayoḥ rau prāpnoti dvivacanabahuvacanayoḥ ca raḥ prāpnoti . \\
\noindent \textbf{(P\_2,4.85.1)} atha kṛtadvandvānām ekaśeṣaḥ na doṣaḥ bhavati . \\
\noindent \textbf{(P\_2,4.85.1)} yathā na doṣaḥ tathā astu . \\
\noindent \textbf{(P\_2,4.85.1)} kim punaḥ atra jyāyaḥ . \\
\noindent \textbf{(P\_2,4.85.1)} ubhayam iti āha . \\
\noindent \textbf{(P\_2,4.85.1)} ubhayam hi dṛśyate . \\
\noindent \textbf{(P\_2,4.85.1)} bahu śaktikiṭakam bahūni śaktikiṭakāni bahu sthālīpiṭharam bahūni sthālīpiṭharāni . \\
\noindent \textbf{(P\_2,4.85.1)} ḍāraurasaḥ kṛte ṭeḥ e yathāt dvitvam prasāraṇe samasṅkhyena na artha asti . \\
\noindent \textbf{(P\_2,4.85.1)} siddham sthāne arthataḥ antarāḥ . \\
\noindent \textbf{(P\_2,4.85.1)} āntaryataḥ vyavasthā . \\
\noindent \textbf{(P\_2,4.85.1)} trayaḥ eva ime bhavantu sarveṣām . \\
\noindent \textbf{(P\_2,4.85.1)} ṭeḥ etvam ca paratvāt kṛte api tasmin ime santu . \\
\noindent \textbf{(P\_2,4.85.2)} ḍāvikārasya śitkaraṇam sarvādeśārtham . \\
\noindent \textbf{(P\_2,4.85.2)} ḍāvikāraḥ śit kartavyaḥ . \\
\noindent \textbf{(P\_2,4.85.2)} kim prayojanam . \\
\noindent \textbf{(P\_2,4.85.2)} sarvādeśārtham . \\
\noindent \textbf{(P\_2,4.85.2)} śit sarvasya iti sarvādeśaḥ yathā syāt . \\
\noindent \textbf{(P\_2,4.85.2)} akriyamāṇe hi śakāre alaḥ antyasya vidhayaḥ bhavanti iti antasya prasajyeta . \\
\noindent \textbf{(P\_2,4.85.2)} nighātaprasaṅgaḥ tu . \\
\noindent \textbf{(P\_2,4.85.2)} nighātaḥ tu prapnoti . \\
\noindent \textbf{(P\_2,4.85.2)} śvaḥ kartā . \\
\noindent \textbf{(P\_2,4.85.2)} tāseḥ param lasārvadhātukam anudāttam bhavati iti eṣaḥ svaraḥ prāpnoti . \\
\noindent \textbf{(P\_2,4.85.2)} yat tāvat ucyate ḍāvikārasya śitkaraṇam sarvādeśārtham iti . \\
\noindent \textbf{(P\_2,4.85.2)} siddham alaḥ antyavikārāt . \\
\noindent \textbf{(P\_2,4.85.2)} siddham etat . \\
\noindent \textbf{(P\_2,4.85.2)} katham .alaḥ antyavikārāt . \\
\noindent \textbf{(P\_2,4.85.2)} astu ayam alaḥ antyasya . \\
\noindent \textbf{(P\_2,4.85.2)} kā rūpasiddhiḥ : kartā . \\
\noindent \textbf{(P\_2,4.85.2)} ḍiti ṭeḥ lopāt lopaḥ . \\
\noindent \textbf{(P\_2,4.85.2)} ḍiti ṭeḥ lopena lopaḥ bhaviṣyati . \\
\noindent \textbf{(P\_2,4.85.2)} abhatvāt na prāpnoti . \\
\noindent \textbf{(P\_2,4.85.2)} ḍitkaraṇasāmarthyāt bhaviṣyati . \\
\noindent \textbf{(P\_2,4.85.2)} anittvāt vā . \\
\noindent \textbf{(P\_2,4.85.2)} atha vā anittvāt etat siddham . \\
\noindent \textbf{(P\_2,4.85.2)} kim idam anittvāt iti . \\
\noindent \textbf{(P\_2,4.85.2)} antyasya ayam sthāne bhavan na pratyayaḥ syāt . \\
\noindent \textbf{(P\_2,4.85.2)} asatyām pratyayasañjñāyām itsañjñā na . \\
\noindent \textbf{(P\_2,4.85.2)} asatyām itsañjñāyām lopaḥ na . \\
\noindent \textbf{(P\_2,4.85.2)} asati lope anekāl . \\
\noindent \textbf{(P\_2,4.85.2)} yada anekāl tadā sarvādeśaḥ . \\
\noindent \textbf{(P\_2,4.85.2)} yadā sarvādeśaḥ tadā pratyayaḥ . \\
\noindent \textbf{(P\_2,4.85.2)} yadā pratyayaḥ tadā itsañjñā . \\
\noindent \textbf{(P\_2,4.85.2)} yadā itsañjñā tadā lopaḥ . \\
\noindent \textbf{(P\_2,4.85.2)} praśliṣṭanirdeśāt vā . \\
\noindent \textbf{(P\_2,4.85.2)} atha vā praśliṣṭanirdeśaḥ ayam . \\
\noindent \textbf{(P\_2,4.85.2)} ḍā ā ḍā . \\
\noindent \textbf{(P\_2,4.85.2)} saḥ anekālśit sarvasya iti sarvādeśaḥ bhaviṣyati . \\
\noindent \textbf{(P\_2,4.85.2)} yadā tarhi ayam antyasya sthāne bhavati tadā tiṅgrahaṇena grahaṇam na prāpnoti . \\
\noindent \textbf{(P\_2,4.85.2)} tiṅgrahaṇam ekadeśavikṛtasya ananyatvāt . \\
\noindent \textbf{(P\_2,4.85.2)} ekadeśavikṛtam ananyavat bhavati iti tiṅgrahaṇena grahaṇam bhaviṣyati svaraḥ katham . \\
\noindent \textbf{(P\_2,4.85.2)} svare vipratiṣedhāt siddham . \\
\noindent \textbf{(P\_2,4.85.2)} ḍāraurasaḥ kriyantām anudāttatvam iti kim atra kartavyam . \\
\noindent \textbf{(P\_2,4.85.2)} paratvāt anudāttatvam . \\
\noindent \textbf{(P\_2,4.85.2)} nityāḥ ḍāraurasaḥ . \\
\noindent \textbf{(P\_2,4.85.2)} kṛte api anudāttatve prapnuvanti akṛte api prapnuvanti . \\
\noindent \textbf{(P\_2,4.85.2)} anudāttatvam api nityam . \\
\noindent \textbf{(P\_2,4.85.2)} kṛteṣu kṛteṣu api ḍāraurassu prāpnoti akṛteṣu api prāpnoti . \\
\noindent \textbf{(P\_2,4.85.2)} anityam anudāttatvam . \\
\noindent \textbf{(P\_2,4.85.2)} anyasya kṛteṣu ḍāraurassu prāpnoti anyasya akṛteṣu . \\
\noindent \textbf{(P\_2,4.85.2)} śabdāntarasya ca prāpnuvan vidhiḥ anityaḥ bhavati . \\
\noindent \textbf{(P\_2,4.85.2)} ḍāraurasaḥ api anityāḥ . \\
\noindent \textbf{(P\_2,4.85.2)} anyathāsvarasya kṛte anudāttatve prāpnuvanti anyathāsvarasya akṛte . \\
\noindent \textbf{(P\_2,4.85.2)} svarabhinnasya ca prāpnuvantaḥ anityāḥ bhavanti . \\
\noindent \textbf{(P\_2,4.85.2)} ubhayoḥ anityayoḥ paratvāt anudāttatvam . \\
\noindent \textbf{(P\_2,4.85.2)} anudāttatve kṛte punaḥprasaṅgavijñātāt ḍāraurasaḥ . \\
\noindent \textbf{(P\_2,4.85.2)} ṭilope udāttanivṛttisvareṇa siddham . \\
\noindent \textbf{(P\_2,4.85.2)} na sidhyati . \\
\noindent \textbf{(P\_2,4.85.2)} kim kāraṇam . \\
\noindent \textbf{(P\_2,4.85.2)} antaraṅgatvāt ḍāraurasaḥ . \\
\noindent \textbf{(P\_2,4.85.2)} tatra antaraṅgatvāt ḍāraurassu kṛteṣu anudāttatvam kriyatām ṭilopaḥ iti kim atra kartavyam . \\
\noindent \textbf{(P\_2,4.85.2)} paratvāt ṭilopena bhavitavyam . \\
\noindent \textbf{(P\_2,4.85.2)} evam tarhi svare vipratiṣedhāt siddham . \\
\noindent \textbf{(P\_2,4.85.2)} nyāyyaḥ eva ayam svare vipratiṣedhaḥ . \\
\noindent \textbf{(P\_2,4.85.2)} idam iha sampradhāryam . \\
\noindent \textbf{(P\_2,4.85.2)} anudāttatvam kriyatām udāttanivṛttisvaraḥ iti . \\
\noindent \textbf{(P\_2,4.85.2)} kim atra kartavyam . \\
\noindent \textbf{(P\_2,4.85.2)} paratvāt anudāttatvam . \\
\noindent \textbf{(P\_2,4.85.2)} anudāttatve kṛte punaḥprasaṅgavijñānāt udāttanivṛttisvaraḥ bhaviṣyati . \\
\noindent \textbf{(P\_2,4.85.2)} tat etat kva siddham bhavati . \\
\noindent \textbf{(P\_2,4.85.2)} yat pit vacanam . \\
\noindent \textbf{(P\_2,4.85.2)} yat apit vacanam tatra na sidhyati . \\
\noindent \textbf{(P\_2,4.85.2)} tatra api siddham . \\
\noindent \textbf{(P\_2,4.85.2)} katham idam adya lasārvadhātukānudāttatvam pratyayasvarasya apavādaḥ . \\
\noindent \textbf{(P\_2,4.85.2)} na ca apavādaviṣaye utsargaḥ abhiniviśate . \\
\noindent \textbf{(P\_2,4.85.2)} pūrvam hi apavādāḥ abhiniviśante paścāt utsargāḥ . \\
\noindent \textbf{(P\_2,4.85.2)} prakalpya vā apavādaviṣayam tataḥ utsargaḥ abhiniviśate . \\
\noindent \textbf{(P\_2,4.85.2)} tat na tāvat kadā cit pratyayasvaraḥ bhavati . \\
\noindent \textbf{(P\_2,4.85.2)} apavādam lasārvadhātukanudāttatvam pratīkṣate . \\
\noindent \textbf{(P\_2,4.85.2)} tatra anudāttatvam kriyatām lopaḥ iti . \\
\noindent \textbf{(P\_2,4.85.2)} yadi api paratvāt lopaḥ saḥ asau avidyamānodāttatve anudātte udāttaḥ lupyate . \\
\noindent \textbf{(P\_2,4.85.2)} pratyayasvarāpavādaḥ lasārvadhātukānudāttam . \\
\noindent \textbf{(P\_2,4.85.2)} tena tatra na prasaktaḥ pratyayasvaraḥ kadā cit . \\
\noindent \textbf{(P\_2,4.85.2)} pratyayasvaraḥ tu tāseḥ vṛttisanniyogaśiṣṭaḥ . \\
\noindent \textbf{(P\_2,4.85.2)} tena ca api asau udāttaḥ lopsyate . \\
\noindent \textbf{(P\_2,4.85.2)} tathā na doṣaḥ . \\
\noindent \textbf{(P\_3,1.1)} adhikāreṇa iyam pratyayasañjñā kriyate . \\
\noindent \textbf{(P\_3,1.1)} sā prakṛtyupapadopādhīnām api prāpnoti . \\
\noindent \textbf{(P\_3,1.1)} tasyāḥ pratiṣedhaḥ vaktavyaḥ . \\
\noindent \textbf{(P\_3,1.1)} prakṛti . \\
\noindent \textbf{(P\_3,1.1)} guptijkibhyaḥ san . \\
\noindent \textbf{(P\_3,1.1)} upapada . \\
\noindent \textbf{(P\_3,1.1)} stambakarṇayoḥ ramajapoḥ . \\
\noindent \textbf{(P\_3,1.1)} upādhi . \\
\noindent \textbf{(P\_3,1.1)} harateḥ dṛtināthayoḥ paśau . \\
\noindent \textbf{(P\_3,1.1)} eteṣām pratiṣedhaḥ vaktavyaḥ . \\
\noindent \textbf{(P\_3,1.1)} kim ca syāt yadi eteṣām api pratyayasañjñā syāt . \\
\noindent \textbf{(P\_3,1.1)} paratvam ādyudāttatvam aṅgasañjñā iti ete vidhayaḥ prasajyeran . \\
\noindent \textbf{(P\_3,1.1)} ataḥ uttaram paṭhati . \\
\noindent \textbf{(P\_3,1.1)} pratyayādhikāre prakṛtyupapadopādhīnām apratiṣedhaḥ . \\
\noindent \textbf{(P\_3,1.1)} adhikāreṇa api pratyayasañjñāyām satyām prakṛtyupapadopādhīnām apratiṣedhaḥ . \\
\noindent \textbf{(P\_3,1.1)} anarthakaḥ pratiṣedhaḥ apratiṣedhaḥ . \\
\noindent \textbf{(P\_3,1.1)} pratyayasañjñā kasmāt na bhavati . \\
\noindent \textbf{(P\_3,1.1)} nimittasya nimittikāryārthatvāt anyatra api . \\
\noindent \textbf{(P\_3,1.1)} nimittāni hi nimittikāryārthāni bhavanti . \\
\noindent \textbf{(P\_3,1.1)} kim punaḥ nimittam kaḥ vā nimittī . \\
\noindent \textbf{(P\_3,1.1)} prakṛtyupapaopādhayaḥ nimittam pratyayaḥ nimittī . \\
\noindent \textbf{(P\_3,1.1)} anyatra api ca eṣaḥ nyāyaḥ dṛṣṭaḥ . \\
\noindent \textbf{(P\_3,1.1)} kva anyatra . \\
\noindent \textbf{(P\_3,1.1)} loke . \\
\noindent \textbf{(P\_3,1.1)} tat yathā . \\
\noindent \textbf{(P\_3,1.1)} bahuṣu āsīneṣu kaḥ cit kam cit pṛcchati . \\
\noindent \textbf{(P\_3,1.1)} katamaḥ devadattaḥ . \\
\noindent \textbf{(P\_3,1.1)} kataraḥ yajñadattaḥ iti . \\
\noindent \textbf{(P\_3,1.1)} saḥ tasmai ācaṣṭe . \\
\noindent \textbf{(P\_3,1.1)} yaḥ aśve yaḥ pīṭhe iti ukte nimittasya nimittikāryārthatvāt adhyavasyati ayam devadattaḥ ayam yajñdatta iti . \\
\noindent \textbf{(P\_3,1.1)} na idānīm aśvasya pīṭhasya vā devadattaḥ iti sañjñā bhavati . \\
\noindent \textbf{(P\_3,1.1)} kim punaḥ nimittam kaḥ vā nimittī . \\
\noindent \textbf{(P\_3,1.1)} nirjñātaḥ arthaḥ nimittam anirjñātārthaḥ nimittī . \\
\noindent \textbf{(P\_3,1.1)} iha ca pratyayaḥ anirjñātaḥ prakṛtyupapadopādhayaḥ nirjñātāḥ . \\
\noindent \textbf{(P\_3,1.1)} kva . \\
\noindent \textbf{(P\_3,1.1)} dhātūpadeśe prātipadikopadeśe ca . \\
\noindent \textbf{(P\_3,1.1)} te nirjñātāḥ nimittatvena upādīyante . \\
\noindent \textbf{(P\_3,1.1)} pradhāne kāryasampratyayāt vā siddham . atha vā pradhāne kāryasampratyayaḥ bhaviṣyati . \\
\noindent \textbf{(P\_3,1.1)} kim ca pradhānam . \\
\noindent \textbf{(P\_3,1.1)} pratyayaḥ . \\
\noindent \textbf{(P\_3,1.1)} tat yathā . \\
\noindent \textbf{(P\_3,1.1)} bahuṣu yātsu kaḥ cit kam cit pṛcchati . \\
\noindent \textbf{(P\_3,1.1)} kaḥ yāti iti . \\
\noindent \textbf{(P\_3,1.1)} saḥ āha rājā iti . \\
\noindent \textbf{(P\_3,1.1)} rājā iti ukte pradhāne kāryasampratyayāt yaḥ pṛcchati yaḥ ca ācaṣṭe ubhayoḥ sampratyayaḥ bhavati . \\
\noindent \textbf{(P\_3,1.1)} kiṅkṛtam punaḥ prādhānyam . \\
\noindent \textbf{(P\_3,1.1)} arthakṛtam . \\
\noindent \textbf{(P\_3,1.1)} yathā punaḥ loke arthakṛtam prādhānyam śabdasya idānīm kiṅkṛtam prādhānyam . \\
\noindent \textbf{(P\_3,1.1)} śabdasya apūrvopadeśaḥ prādhānyam . \\
\noindent \textbf{(P\_3,1.1)} yasya apūrvopadeśaḥ saḥ pradhānam . \\
\noindent \textbf{(P\_3,1.1)} prakṛtyupapadopādhayaḥ ca upadiṣṭāḥ . \\
\noindent \textbf{(P\_3,1.1)} kva . \\
\noindent \textbf{(P\_3,1.1)} dhātūpadeśe prātipadikopadeśe ca . \\
\noindent \textbf{(P\_3,1.1)} yadi eva nimittasya nimittikāryārthatvāt atha api pradhāne kāryasampratyayāt prakṛtyupapadopādhīnām na bhavati vikārāgamānām tu prāpnoti . \\
\noindent \textbf{(P\_3,1.1)} hanaḥ ta ca . \\
\noindent \textbf{(P\_3,1.1)} trapujatunoḥ ṣuk iti . \\
\noindent \textbf{(P\_3,1.1)} eteṣām hi apūrvopadeśāt prādhānyam . \\
\noindent \textbf{(P\_3,1.1)} nimittinaḥ ca ete . \\
\noindent \textbf{(P\_3,1.1)} vikārāgameṣu ca paravijñānāt . \\
\noindent \textbf{(P\_3,1.1)} vikārāgameṣu ca paravijñānāt pratyayasañjñā na bhaviṣyati . \\
\noindent \textbf{(P\_3,1.1)} pratyayaḥ paraḥ bahvati iti ucyate . \\
\noindent \textbf{(P\_3,1.1)} na ca vikārāgamāḥ pare sambhavanti . \\
\noindent \textbf{(P\_3,1.1)} kim punaḥ kāraṇam samāne apūrvopadeśe pratyayaḥ paraḥ vikārāgamāḥ na pare . \\
\noindent \textbf{(P\_3,1.1)} ṣaṣṭhīnirdiṣṭasya ca tadyuktatvāt . \\
\noindent \textbf{(P\_3,1.1)} ṣaṣṭhīnirdiṣtam vikārāgamayuktam pañcamīnirdṣṭāt ca pratyayaḥ vidhīyate . \\
\noindent \textbf{(P\_3,1.1)} pratyayavidhānānupapattiḥ tu . \\
\noindent \textbf{(P\_3,1.1)} pratyayavidhiḥ tu na upapapdyate . \\
\noindent \textbf{(P\_3,1.1)} kva . \\
\noindent \textbf{(P\_3,1.1)} yatra vikārāgamāḥ vidhīyante . \\
\noindent \textbf{(P\_3,1.1)} hanaḥ ta ca . \\
\noindent \textbf{(P\_3,1.1)} tarpujatunoḥ ṣuk . \\
\noindent \textbf{(P\_3,1.1)} kim punaḥ kāraṇam na sidhyati . \\
\noindent \textbf{(P\_3,1.1)} vikārāgamayuktatvāt apañcamīnirdiṣṭatvāt ca . \\
\noindent \textbf{(P\_3,1.1)} tasmāt tatra pañcamīnirdeśāt siddham . \\
\noindent \textbf{(P\_3,1.1)} tasmāt tatra pañcamīnirdeśaḥ kartavyaḥ . \\
\noindent \textbf{(P\_3,1.1)} na kartavyaḥ . \\
\noindent \textbf{(P\_3,1.1)} iha tāvat hanaḥ te iti . \\
\noindent \textbf{(P\_3,1.1)} dhātoḥ iti vartate. iha trapujatunoḥ ṣuk iti . \\
\noindent \textbf{(P\_3,1.1)} prātipadikāt iti vartate . \\
\noindent \textbf{(P\_3,1.1)} yadi evam hanaḥ ta ca dhātoḥ kyap bhavati iti dhātumātrāt kyap prāpnoti . \\
\noindent \textbf{(P\_3,1.1)} na eṣaḥ doṣaḥ . \\
\noindent \textbf{(P\_3,1.1)} ācāryapravṛttiḥ jñāpayati na dhātumātrāt kyap bhavati iti yat ayam etistuśasvṛdṛjuṣaḥ kyap iti parigaṇanam karoti . \\
\noindent \textbf{(P\_3,1.1)} atha vā hantim eva atra dhātugrahaṇena abhisambhantsyāmaḥ . \\
\noindent \textbf{(P\_3,1.1)} hanaḥ taḥ bhavati . \\
\noindent \textbf{(P\_3,1.1)} dhātoḥ kyap bhavati . \\
\noindent \textbf{(P\_3,1.1)} kasmāt . \\
\noindent \textbf{(P\_3,1.1)} hanteḥ iti . \\
\noindent \textbf{(P\_3,1.1)} arthāśrayatvāt vā . \\
\noindent \textbf{(P\_3,1.1)} atha vā arthāśrayaḥ pratyayavidhiḥ . \\
\noindent \textbf{(P\_3,1.1)} yaḥ tam artham sampratyāyayati saḥ pratyayaḥ . \\
\noindent \textbf{(P\_3,1.1)} kim vaktavyam etat . \\
\noindent \textbf{(P\_3,1.1)} na hi . \\
\noindent \textbf{(P\_3,1.1)} katham anucyamānam gaṃsyate . \\
\noindent \textbf{(P\_3,1.1)} pratyayaḥ iti mahatī sañjñā kriyate . \\
\noindent \textbf{(P\_3,1.1)} sañjñā ca nāma yataḥ na laghīyaḥ . \\
\noindent \textbf{(P\_3,1.1)} kutaḥ etat . \\
\noindent \textbf{(P\_3,1.1)} laghvartham hi sañjñākaraṇam . \\
\noindent \textbf{(P\_3,1.1)} tatra mahatyāḥ sañjñāyāḥ karaṇe etat prayojanam anvarthasañjñā yathā vijñāyeta . \\
\noindent \textbf{(P\_3,1.1)} pratyāyayiti iti pratyayaḥ . \\
\noindent \textbf{(P\_3,1.1)} yadi pratyāyayiti iti pratyayaḥ avikādīnām pratyayasañjñā na prāpnoti . \\
\noindent \textbf{(P\_3,1.1)} na hi te kim cit pratyāyayanti . \\
\noindent \textbf{(P\_3,1.1)} evam tarhi pratyāyyate pratyayaḥ iti . \\
\noindent \textbf{(P\_3,1.1)} evam api sanādīnām na prāpnoti . \\
\noindent \textbf{(P\_3,1.1)} evam tari ubhayasādhanaḥ ayam kartṛsādhanaḥ karmasādhanaḥ ca . \\
\noindent \textbf{(P\_3,1.1)} evam api kutaḥ etat samāne apūrvopadeśe trāpuṣam jātuṣam iti atra akāraḥ tam artham sampratyāyayati na punaḥ ṣakāraḥ iti . \\
\noindent \textbf{(P\_3,1.1)} anyatra api akāreṇa tasya arthasya vacanāt manyāmahe akāraḥ tam artham sampratyāyayatina ṣakāraḥ iti . \\
\noindent \textbf{(P\_3,1.1)} kva anyatra . \\
\noindent \textbf{(P\_3,1.1)} bilvādibhyaḥ aṇ . \\
\noindent \textbf{(P\_3,1.1)} bailvaḥ . \\
\noindent \textbf{(P\_3,1.2)} kimartham idam ucyate . \\
\noindent \textbf{(P\_3,1.2)} paraḥ yathā syāt . \\
\noindent \textbf{(P\_3,1.2)} pūrvaḥ mā bhūt iti . \\
\noindent \textbf{(P\_3,1.2)} na etat asti prayojanam . \\
\noindent \textbf{(P\_3,1.2)} yam icchati pūrvam āha tam : vibhāṣā supaḥ bahuc purastāt tu iti . \\
\noindent \textbf{(P\_3,1.2)} madhye tarhi mā bhūt iti . \\
\noindent \textbf{(P\_3,1.2)} madhye api yam icchati āha tam : avyayasarvanāmnām akac prāk ṭeḥ iti . \\
\noindent \textbf{(P\_3,1.2)} yaḥ idānīm anyaḥ pratyayaḥ śeṣaḥ saḥ antareṇa vacanam paraḥ eva bhaviṣyati iti nā arthaḥ paravacanena . \\
\noindent \textbf{(P\_3,1.2)} evam api yeṣām eva pratyayānām deśaḥ niyamyate te eva niyatadeśāḥ syuḥ . \\
\noindent \textbf{(P\_3,1.2)} yaḥ idānīm aniyatadeśaḥ saḥ kadā cit pūrvaḥ kadā cit paraḥ kadā cit madhye syāt . \\
\noindent \textbf{(P\_3,1.2)} tat yathā mātuḥ vatsaḥ kadā cit agrataḥ kadā cit pṛṣṭhataḥ kadā cit pārśvataḥ bhavati . \\
\noindent \textbf{(P\_3,1.2)} paraḥ eva yathā syāt iti evamartham paravacanam . \\
\noindent \textbf{(P\_3,1.2)} paravacanam anarthakam pañcamīnirdiṣṭatvāt parasya . \\
\noindent \textbf{(P\_3,1.2)} paragrahaṇam anarthakam . \\
\noindent \textbf{(P\_3,1.2)} kim kāraṇam . \\
\noindent \textbf{(P\_3,1.2)} pañcamīnirdiṣṭatvāt parasya kāryam ucyate . \\
\noindent \textbf{(P\_3,1.2)} tat yathā dvyantarupasargebhyaḥ apaḥ īt . \\
\noindent \textbf{(P\_3,1.2)} viṣamaḥ upanyāsaḥ . \\
\noindent \textbf{(P\_3,1.2)} sataḥ tatra parasya kāryam ucyate . \\
\noindent \textbf{(P\_3,1.2)} iha idānīm kasya sataḥ parasya kāryam bhavitum arhati . \\
\noindent \textbf{(P\_3,1.2)} iha api sataḥ eva. katham . \\
\noindent \textbf{(P\_3,1.2)} paratvam svābhāvikam . \\
\noindent \textbf{(P\_3,1.2)} atha vācanike paratve sati arthaḥ syāt paragrahaṇena . \\
\noindent \textbf{(P\_3,1.2)} vācanike ca na arthaḥ . \\
\noindent \textbf{(P\_3,1.2)} etat hi tasya parasya kāryam yat asau paraḥ syāt . \\
\noindent \textbf{(P\_3,1.2)} atha vā yat asya parasya sataḥ sañjñā syāt . \\
\noindent \textbf{(P\_3,1.2)} yatra tarhi pañcamī na asti tadartham ayam yogaḥ vaktavyaḥ . \\
\noindent \textbf{(P\_3,1.2)} kva ca pañcamī na asti . \\
\noindent \textbf{(P\_3,1.2)} yatra vikārāgamāḥ śiṣyante . \\
\noindent \textbf{(P\_3,1.2)} kva ca vikārāgamāḥ śiṣyante . \\
\noindent \textbf{(P\_3,1.2)} hanaḥ ta ca . \\
\noindent \textbf{(P\_3,1.2)} trapujatunoḥ ṣuk iti . \\
\noindent \textbf{(P\_3,1.2)} vikārāgameṣu ca uktam . \\
\noindent \textbf{(P\_3,1.2)} kim uktam . \\
\noindent \textbf{(P\_3,1.2)} pratyayavidhānānupapattiḥ tu . \\
\noindent \textbf{(P\_3,1.2)} tasmāt tatra pañcamīnirdeśāt siddham iti . \\
\noindent \textbf{(P\_3,1.2)} atyantāparadṛṣṭānām vā parabhūtalopārtham . \\
\noindent \textbf{(P\_3,1.2)} atyantāparadṛṣṭānām tarhi parabhūtalopārtham paragrahaṇam kartavyam . \\
\noindent \textbf{(P\_3,1.2)} ye ete atyantāparadṛṣṭāḥ kvibādayaḥ lupyante teṣām parabhūtānām lopaḥ yathā syāt . \\
\noindent \textbf{(P\_3,1.2)} aparabhūtānām mā bhūt . \\
\noindent \textbf{(P\_3,1.2)} kim punaḥ atyantāparadṛṣṭānām parabhūtalopavacane prayojanam . \\
\noindent \textbf{(P\_3,1.2)} kiti ṇiti iti kāryāṇi yathā syuḥ iti . \\
\noindent \textbf{(P\_3,1.2)} etat api na asti prayojanam . \\
\noindent \textbf{(P\_3,1.2)} ācāryapravṛttiḥ jñāpayati atyantāparadṛṣṭāḥ parabhūtāḥ lupyante iti yat ayam teṣu kādīn anubandhān āsajati . \\
\noindent \textbf{(P\_3,1.2)} katham kṛtvā jñāpakam . \\
\noindent \textbf{(P\_3,1.2)} anubandhāsañjane etat prayojanam kiti ṇiti iti kāryāṇi yathā syuḥ iti . \\
\noindent \textbf{(P\_3,1.2)} yadi ca atra atyantāparadṛṣṭāḥ parabhūtāḥ lupyantetataḥ anubandhāsañjanam arthavat bhavati . \\
\noindent \textbf{(P\_3,1.2)} prayoganiyamārtham vā . \\
\noindent \textbf{(P\_3,1.2)} prayoganiyamārtham tarhi paragrahaṇam kartavyam . \\
\noindent \textbf{(P\_3,1.2)} parabhūtānām prayogaḥ yathā syāt . \\
\noindent \textbf{(P\_3,1.2)} aparabhūtānām mā bhūt iti . \\
\noindent \textbf{(P\_3,1.2)} asti punaḥ kim cit aniṣṭam yadarthaḥ niyamaḥ syāt . \\
\noindent \textbf{(P\_3,1.2)} asti iti āha . \\
\noindent \textbf{(P\_3,1.2)} prakṛteḥ arthābhidhāne pratyayādarśanāt . prakṛteḥ arthābhidhāne apratyayikāḥ dṛśyante . \\
\noindent \textbf{(P\_3,1.2)} kva saḥ devadattaḥ kva saḥ yajñdattaḥ babhruḥ maṇḍuḥ lamakaḥ iti . \\
\noindent \textbf{(P\_3,1.2)} bābhravyaḥ māṇḍavyaḥ lāmakāyanaḥ iti prayoktavye babhruḥ maṇḍuḥ lamakaḥ iti prayujyate . \\
\noindent \textbf{(P\_3,1.2)} dvayasajādīnām ca kevaladṛṣṭatvāt . \\
\noindent \textbf{(P\_3,1.2)} dvayasajādīnām ca kevalānām prayogaḥ dṛśyate . \\
\noindent \textbf{(P\_3,1.2)} kim asya dvayasam . \\
\noindent \textbf{(P\_3,1.2)} kim asya mātram . \\
\noindent \textbf{(P\_3,1.2)} kā adya tithī iti . \\
\noindent \textbf{(P\_3,1.2)} dvayasajādayaḥ vai vṛttijasadṛśāḥ avṛttijāḥ yathā bahuḥ tathā . \\
\noindent \textbf{(P\_3,1.2)} vāvacane ca anutpattyartham . vāvacane ca anutpattyartham paragrahaṇam kartavyam . \\
\noindent \textbf{(P\_3,1.2)} vā vacanena anutpattiḥ yathā syāt . \\
\noindent \textbf{(P\_3,1.2)} atha kriyamāṇe api vai paragrahaṇe katham iva vāvacanena anutpattiḥ labhyā . \\
\noindent \textbf{(P\_3,1.2)} kriyamāṇe paragrahaṇe vāvacanena vā paraḥ iti etat abhisambadhyate . \\
\noindent \textbf{(P\_3,1.2)} akriyamāṇe punaḥ paragrahaṇe vāvacanena kim anyat śakyam abhisambandhum anyat ataḥ sañjñāyāḥ . \\
\noindent \textbf{(P\_3,1.2)} na ca sañjñāyāḥ bhāvābhāvau iṣyete . \\
\noindent \textbf{(P\_3,1.2)} vāvacane ca uktam . \\
\noindent \textbf{(P\_3,1.2)} kim uktam . \\
\noindent \textbf{(P\_3,1.2)} vāvacanānarthakyam ca tatra nityatvāt sanaḥ iti . \\
\noindent \textbf{(P\_3,1.2)} prayoganiyamārtham eva tarhi paragrahaṇam kartavyam . \\
\noindent \textbf{(P\_3,1.2)} atha etasmin prayoganiyame sati kim ayam pratyayahiyamaḥ . \\
\noindent \textbf{(P\_3,1.2)} prakṛtiparaḥ eva pratyayaḥ prayoktavyaḥ aprakṛtiparaḥ na iti . \\
\noindent \textbf{(P\_3,1.2)} āhosvit prakṛtiniyamaḥ . \\
\noindent \textbf{(P\_3,1.2)} pratyayaparā eva prakṛtiḥ prayoktavyā apratyayā na iti . \\
\noindent \textbf{(P\_3,1.2)} kaḥ ca atra viśeṣaḥ . \\
\noindent \textbf{(P\_3,1.2)} tatra pratyayaniyame prkṛtiniyamābhāvaḥ . \\
\noindent \textbf{(P\_3,1.2)} tatra pratyayaniyame sati prkṛtiniyamaḥ na prāpnoti . \\
\noindent \textbf{(P\_3,1.2)} apratyayikāyāḥ prakṛteḥ prayogaḥ prāpnoti . \\
\noindent \textbf{(P\_3,1.2)} kva saḥ devadattaḥ kva saḥ yajñdattaḥ babhruḥ maṇḍuḥ lamakaḥ iti . \\
\noindent \textbf{(P\_3,1.2)} astu tarhi prakṛtiniyamaḥ . \\
\noindent \textbf{(P\_3,1.2)} prakṛtiniyame pratyayāniyamaḥ . \\
\noindent \textbf{(P\_3,1.2)} prakṛtiniyame sati pratyayasya niyamaḥ na prāpnoti . \\
\noindent \textbf{(P\_3,1.2)} kim asya dvayasam . \\
\noindent \textbf{(P\_3,1.2)} kim asya mātram . \\
\noindent \textbf{(P\_3,1.2)} kā adya tithī iti . \\
\noindent \textbf{(P\_3,1.2)} aprakṛtikasya pratyayasya prayogaḥ prāpnoti . \\
\noindent \textbf{(P\_3,1.2)} siddham tu ubhayaniyamāt . \\
\noindent \textbf{(P\_3,1.2)} siddham etat . \\
\noindent \textbf{(P\_3,1.2)} katham . \\
\noindent \textbf{(P\_3,1.2)} ubhayaniyamāt . \\
\noindent \textbf{(P\_3,1.2)} ubhayaniyamaḥ ayam . \\
\noindent \textbf{(P\_3,1.2)} prakṛtiparaḥ eva pratyayaḥ prayoktavyaḥ pratyayaparā eva ca prakṛtiḥ iti . \\
\noindent \textbf{(P\_3,1.2)} kim vaktavyam etat . \\
\noindent \textbf{(P\_3,1.2)} na hi . \\
\noindent \textbf{(P\_3,1.2)} katham anucyamānam gaṃsyate . \\
\noindent \textbf{(P\_3,1.2)} paragrahaṇasāmarthyāt . \\
\noindent \textbf{(P\_3,1.2)} antareṇa api paragrahaṇam syāt ayam paraḥ . \\
\noindent \textbf{(P\_3,1.2)} paraḥ eva yathā syāt iti evamartham paragrahaṇam . \\
\noindent \textbf{(P\_3,1.3.1)} kimartham idam ucyate . \\
\noindent \textbf{(P\_3,1.3.1)} ādyudāttaḥ yathā syāt . \\
\noindent \textbf{(P\_3,1.3.1)} antodāttaḥ mā bhūt iti . \\
\noindent \textbf{(P\_3,1.3.1)} na etat asti prayojanam . \\
\noindent \textbf{(P\_3,1.3.1)} yam icchati antodāttam karoti tatra cakāram anubandham āha ca citaḥ antaḥ udāttaḥ iti . \\
\noindent \textbf{(P\_3,1.3.1)} madhyodāttaḥ tarhi mā bhūt iti . \\
\noindent \textbf{(P\_3,1.3.1)} madyodāttam yam icchati tatra repham anubandham karoti āha ca upottamam riti iti . \\
\noindent \textbf{(P\_3,1.3.1)} anudāttaḥ tarhi mā bhūt iti . \\
\noindent \textbf{(P\_3,1.3.1)} anudāttam api yam icchati tatra pakāram anubandham karoti āha ca anudāttau suppitau iti . \\
\noindent \textbf{(P\_3,1.3.1)} svaritaḥ tarhi mā bhūt iti . \\
\noindent \textbf{(P\_3,1.3.1)} svaritam api yam icchati karoti tatra takāram anubandham āha ca tit svaritam iti . \\
\noindent \textbf{(P\_3,1.3.1)} yaḥ idānīm ataḥ anyaḥ pratyayaḥ śeṣaḥ saḥ antareṇa api vacanam ādyudāttaḥ eva bhaviṣyati iti na arthaḥ ādyudāttavacanena . \\
\noindent \textbf{(P\_3,1.3.1)} evam api yeṣām eva pratyayānām svaraḥ niyamyate te eva niyatasvarāḥ syuḥ . \\
\noindent \textbf{(P\_3,1.3.1)} yaḥ idānīm aniyatasvaraḥ saḥ kadā cit ādyudāttaḥ kadā cit antodāttaḥ kadā cit madhyodāttaḥ kadā cit anudāttaḥ kacā cit svaritaḥ syāt . \\
\noindent \textbf{(P\_3,1.3.1)} ādyudāttaḥ eva yathā syāt iti evam artham idam ucyate . \\
\noindent \textbf{(P\_3,1.3.2)} atha kimartham pratyayasañjñāsanniyogena ādyudāttatvam ucyate anudāttatvam ca na yatra eva anyaḥ svaraḥ tatra eva ayam ucyeta . \\
\noindent \textbf{(P\_3,1.3.2)} ñniti ādiḥ nityam pratyayasya ca . \\
\noindent \textbf{(P\_3,1.3.2)} adupadeśāt lasārvadhātukam anudāttam suppitau ca iti . \\
\noindent \textbf{(P\_3,1.3.2)} tatra ayam api arthaḥ dviḥ ādyudāttagrahaṇam dviḥ ca anudāttagrahaṇam na kartavyam bhavati . \\
\noindent \textbf{(P\_3,1.3.2)} prakṛtam anuvartate . \\
\noindent \textbf{(P\_3,1.3.2)} ataḥ uttaram paṭhati : ādyudāttatvasya pratyayasañjñāsanniyoge prayojanam yasya sañjñākaraṇam tasya ādyudāttārtham . \\
\noindent \textbf{(P\_3,1.3.2)} ādyudāttatvasya pratyayasañjñāsanniyogakaraṇe etat prayojanam yasya sañjñākaraṇam tasya ādyudāttatvam yathā syāt . \\
\noindent \textbf{(P\_3,1.3.2)} asanniyoge hi yasmāt saḥ tadādeḥ ādyudāttatvam tadantasya ca anudāttatvam . \\
\noindent \textbf{(P\_3,1.3.2)} akriyamāṇe hi pratyayasañjñāsanniyogena ādyudāttatve pratyayagrahaṇe yasmāt saḥ tadādeḥ grahaṇam bhavati iti tadādeḥ ādyudāttatvam prasajyeta tadantasya ca anudāttatvam . \\
\noindent \textbf{(P\_3,1.3.2)} atha kriyamāṇe api pratyayasañjñāsanniyogena ādyudāttatve anudāttatve ca kasmāt eva tadādeḥ ādyudāttatvam na bhavati tadantasya ca anudāttatvam . \\
\noindent \textbf{(P\_3,1.3.2)} utpannaḥ pratyayaḥ pratyayāśrayāṇām kāryāṇām nimittam bhavati na utpadyamānaḥ . \\
\noindent \textbf{(P\_3,1.3.2)} tat yathā ghaṭaḥ kṛtaḥ ghaṭāśrayāṇām kāryāṇām nimittam bhavati na kriyamāṇaḥ . \\
\noindent \textbf{(P\_3,1.3.2)} na vā prakṛteḥ ādyudāttavacanam jñāpakam tadādeḥ agrahaṇasya . \\
\noindent \textbf{(P\_3,1.3.2)} na vā eṣaḥ doṣaḥ . \\
\noindent \textbf{(P\_3,1.3.2)} kim kāraṇam . \\
\noindent \textbf{(P\_3,1.3.2)} yat ayam ñniti ādiḥ nityam iti prakṛteḥ ādyudāttatvam śāsti tat jñāpayati ācāryaḥ na tadādeḥ ādyudāttatvam bhavati iti . \\
\noindent \textbf{(P\_3,1.3.2)} tadantasya tarhi anudāttatvam prāpnoti . \\
\noindent \textbf{(P\_3,1.3.2)} prakṛtisvarasya ca vidhānasāmarthyāt pratyayasvarābhāvaḥ . yat ayam dhātoḥ antaḥ prātipadikasya antaḥ iti prakṛteḥ antodāttatvam śāsti tat jñāpayati ācāryaḥ na tadantasya anudāttatvam bhavati iti . \\
\noindent \textbf{(P\_3,1.3.2)} katham kṛtvā jñāpakam . \\
\noindent \textbf{(P\_3,1.3.2)} yatra hi anudāttaḥpratyayaḥ prakṛtisvaraḥ tat prayojayati . \\
\noindent \textbf{(P\_3,1.3.2)} āgamānudāttārtham vā . \\
\noindent \textbf{(P\_3,1.3.2)} āgamānudāttārtham tarhi pratyayasañjñāsanniyogena ādyudāttatvam ucyate . \\
\noindent \textbf{(P\_3,1.3.2)} pratyayasañjñāsanniyogena ādyudāttatve kṛte āgamāḥ anudāttāḥ yathā syuḥ iti . \\
\noindent \textbf{(P\_3,1.3.2)} na vā āgamasya anudāttavacanāt . \\
\noindent \textbf{(P\_3,1.3.2)} na vā etat api prayojanam asti . \\
\noindent \textbf{(P\_3,1.3.2)} kim kāraṇam . \\
\noindent \textbf{(P\_3,1.3.2)} āgamasya anudāttavacanāt . \\
\noindent \textbf{(P\_3,1.3.2)} āgamāḥ anudāttāḥ bhavanti iti vakṣyāmi . \\
\noindent \textbf{(P\_3,1.3.2)} ke punaḥ āgamāḥ anudāttatvam prayojayanti . \\
\noindent \textbf{(P\_3,1.3.2)} iṭ . \\
\noindent \textbf{(P\_3,1.3.2)} lavitā . \\
\noindent \textbf{(P\_3,1.3.2)} iṭ tāvat na prayojayati . \\
\noindent \textbf{(P\_3,1.3.2)} idam iha sampradhāryam . \\
\noindent \textbf{(P\_3,1.3.2)} iṭ kriyatām ādyudāttatvam iti . \\
\noindent \textbf{(P\_3,1.3.2)} kim atra kartavyam . \\
\noindent \textbf{(P\_3,1.3.2)} paratvāt iḍāgamaḥ . \\
\noindent \textbf{(P\_3,1.3.2)} nityam ādyudāttatvam . \\
\noindent \textbf{(P\_3,1.3.2)} kṛte api iṭi prāpnoti akṛte api prāpnoti . \\
\noindent \textbf{(P\_3,1.3.2)} iṭ api nityaḥ . \\
\noindent \textbf{(P\_3,1.3.2)} kṛte api ādyudāttatve prāpnoti akṛte api prāpnoti . \\
\noindent \textbf{(P\_3,1.3.2)} anityaḥ iṭ . \\
\noindent \textbf{(P\_3,1.3.2)} anyathāsvarasya kṛte ādyudāttatve prapnoti anyathāsvarasya akṛte . \\
\noindent \textbf{(P\_3,1.3.2)} svarabhinnasya ca prāpnuvan vidhiḥ anityaḥ bhavati . \\
\noindent \textbf{(P\_3,1.3.2)} ādyudāttatvam api anityam . \\
\noindent \textbf{(P\_3,1.3.2)} anyasya kṛte iṭi prāpnoti anyasya akṛte . \\
\noindent \textbf{(P\_3,1.3.2)} śabdāntarasya ca prāpnuvan vidhiḥ anityaḥ bhavati . \\
\noindent \textbf{(P\_3,1.3.2)} ubhayoḥ anityayoḥ paratvāt iḍāgamaḥ . \\
\noindent \textbf{(P\_3,1.3.2)} antaraṅgam tarhi ādyudāttatvam . \\
\noindent \textbf{(P\_3,1.3.2)} kā antaraṅgatā . \\
\noindent \textbf{(P\_3,1.3.2)} utpattisanniyogena ādyudāttatvam ucyate . \\
\noindent \textbf{(P\_3,1.3.2)} utpanne pratyaye prakṛtipratyayau āśritya aṅgasya iḍāgamaḥ . \\
\noindent \textbf{(P\_3,1.3.2)} ādyudāttatvam api na antaraṅgam yāvatā pratyaye āśrīyamāṇe prakṛtiḥ api āśritā bhavati . \\
\noindent \textbf{(P\_3,1.3.2)} antaraṅgam eva ādyudāttatvam . \\
\noindent \textbf{(P\_3,1.3.2)} katham . \\
\noindent \textbf{(P\_3,1.3.2)} idānīm eva hi uktam na pratyayasvaravidhau tadādividhiḥ bhavati iti . \\
\noindent \textbf{(P\_3,1.3.2)} sīyuṭ tarhi prayojayati . \\
\noindent \textbf{(P\_3,1.3.2)} avacane hi sīyuḍādeḥ ādyudāttatvam . akriyamāṇe hi āgamānudāttatve kriyamāṇe api pratyayasañjñāsanniyogena ādyudāttatve sīyuḍādeḥ liṅaḥ ādyudāttatvam prasajyeta . \\
\noindent \textbf{(P\_3,1.3.2)} laviṣīya paviṣīya . \\
\noindent \textbf{(P\_3,1.3.2)} tat tarhi vaktavyam āgamāḥ anudāttāḥ bhavanti iti . \\
\noindent \textbf{(P\_3,1.3.2)} na vaktavyam . \\
\noindent \textbf{(P\_3,1.3.2)} ācāryapravṛttiḥ jñāpayati āgamāḥ anudāttāḥ bhavanti iti yat ayam yāsuṭ parasamaipadeṣu udāttaḥ ṅit ca iti āha . \\
\noindent \textbf{(P\_3,1.3.2)} na etat asti jñāpakam vakṣyati etat . \\
\noindent \textbf{(P\_3,1.3.2)} yāsuṭaḥ ṅidvacanam pidartham udāttavacanam ca iti . \\
\noindent \textbf{(P\_3,1.3.2)} śakyam anena vaktum : yāsuṭ parasmaipadeṣu bhavati apit ca liṅ bhavati iti . \\
\noindent \textbf{(P\_3,1.3.2)} saḥ ayam evam laghīyasā nyāsena siddhe sati yat garīyāṃsam yatnam ārabhate tat jñapayati ācāryaḥ āgamāḥ anudāttāḥ bhavanti iti . \\
\noindent \textbf{(P\_3,1.3.2)} śakyam idam labdhum . \\
\noindent \textbf{(P\_3,1.3.2)} yadi eva vacanāt atha api jñāpakāt āgamāḥ anudāttāḥ bhavanti . \\
\noindent \textbf{(P\_3,1.3.2)} āgamaiḥ tu vyavahitatvāt ādyudāttatvam na prāpnoti . \\
\noindent \textbf{(P\_3,1.3.2)} āgamāḥ avidyamānavat bhavanti iti vakṣyāmi . \\
\noindent \textbf{(P\_3,1.3.2)} yadi āgamāḥ avidyamānavat bhavanti iti ucyate lavitā avādeśaḥ na prāpnoti . \\
\noindent \textbf{(P\_3,1.3.2)} svaravidhau iti vakṣyāmi . \\
\noindent \textbf{(P\_3,1.3.2)} evam api lavitā udāttāt anudāttasya svaritaḥ iti svaritaḥ na prāpnoti . \\
\noindent \textbf{(P\_3,1.3.2)} ṣāṣṭhike svare iti vakṣyāmi . \\
\noindent \textbf{(P\_3,1.3.2)} evam api śikṣitaḥ niṣṭhā ca dvyac anāt it eṣaḥ svaraḥ prāpnoti . \\
\noindent \textbf{(P\_3,1.3.2)} pratyayasvaravidhau iti vakṣyāmi . \\
\noindent \textbf{(P\_3,1.3.2)} tat tarhi vaktavyam avidyamānavat bhavanti iti . \\
\noindent \textbf{(P\_3,1.3.2)} na vaktavyam . \\
\noindent \textbf{(P\_3,1.3.2)} ācāryapravṛttiḥ jñāpayati āgamāḥ avidyamānavat bhavanti iti yat ayam yāsuṭ parasamaipadeṣu udāttaḥ ṅit ca iti āha . \\
\noindent \textbf{(P\_3,1.3.2)} na etat asti jñāpakam . \\
\noindent \textbf{(P\_3,1.3.2)} vakṣyati etat . \\
\noindent \textbf{(P\_3,1.3.2)} yāsuṭaḥ ṅidvacanam pidartham udāttavacanam ca iti . \\
\noindent \textbf{(P\_3,1.3.2)} śakyam anena vaktum . \\
\noindent \textbf{(P\_3,1.3.2)} yāsuṭ parasmaipadeṣu bhavati apit ca liṅ bhavati iti . \\
\noindent \textbf{(P\_3,1.3.2)} saḥ ayam evam laghīyasā nyāsena siddhe sati yat garīyāṃsam yatnam ārabhate tat jñapayati ācāryaḥ āgamāḥ avidyamānavat bhavanti iti . \\
\noindent \textbf{(P\_3,1.3.2)} ādyudāttasya vā lopārtham . \\
\noindent \textbf{(P\_3,1.3.2)} ādyudāttasya tarhi lopārtham pratyayasañjñāsanniyogena ādyudāttatvam ucyate . \\
\noindent \textbf{(P\_3,1.3.2)} pratyayasañjñāsanniyogena ādyudāttatve kṛte udāttanivṛttisvaraḥ siddhaḥ bhavati : sraughnī māthurī . \\
\noindent \textbf{(P\_3,1.3.2)} atra hi paratvāt lopaḥ pratyaysvaram bādheta . \\
\noindent \textbf{(P\_3,1.3.2)} na vā bahiraṅgalakṣaṇatvāt . \\
\noindent \textbf{(P\_3,1.3.2)} na vā etat prayojayati . \\
\noindent \textbf{(P\_3,1.3.2)} kim kāraṇam . \\
\noindent \textbf{(P\_3,1.3.2)} bahiraṅgalakṣaṇatvāt . \\
\noindent \textbf{(P\_3,1.3.2)} bahiraṅgalakṣaṇaḥ lopaḥ antaraṅgalakṣaṇaḥ svaraḥ . \\
\noindent \textbf{(P\_3,1.3.2)} asiddham bahiraṅgam antaraṅge . \\
\noindent \textbf{(P\_3,1.3.2)} avaśyam ca eṣā paribhāṣā āśrayitavyā . \\
\noindent \textbf{(P\_3,1.3.2)} avacane hi ñinnitkitsu atiprasaṅgaḥ . \\
\noindent \textbf{(P\_3,1.3.2)} anāśrīyamāṇāyām asyām paribhāṣāyām kriyamāṇe api pratyayasañjñāsanniyogena ādyudāttatveñinnitkitsu atiprasaṅgaḥ syāt . \\
\noindent \textbf{(P\_3,1.3.2)} autsī kaṃsikī ātreyī iti . \\
\noindent \textbf{(P\_3,1.3.2)} atra hi paratvāt lopaḥ ñinnitkitsvarān bādheta . \\
\noindent \textbf{(P\_3,1.3.2)} na eṣaḥ doṣaḥ . \\
\noindent \textbf{(P\_3,1.3.2)} ñinnitkitsvarāḥ pratyaysvarāpavādāḥ . \\
\noindent \textbf{(P\_3,1.3.2)} na ca apavādaviṣaye utsargaḥ bhiniviśate . \\
\noindent \textbf{(P\_3,1.3.2)} pūrvam hi apavādāḥ abhiniviśante paścāt utsargaḥ . \\
\noindent \textbf{(P\_3,1.3.2)} prakalpya vā apavādaviṣayam tataḥ utsargaḥ abhiniviśate . \\
\noindent \textbf{(P\_3,1.3.2)} na tāvat atra kadā cit pratyayādyudāttatvam bhavati . \\
\noindent \textbf{(P\_3,1.3.2)} apavādān ñinnitkitsvarān pratīkṣate . \\
\noindent \textbf{(P\_3,1.3.2)} kaṃsikyām bhūyān apahāraḥ . \\
\noindent \textbf{(P\_3,1.3.2)} anyasya atra udāttatvam anyasya lopaḥ . \\
\noindent \textbf{(P\_3,1.3.2)} ādeḥ udāttatvam antyasya lopaḥ . \\
\noindent \textbf{(P\_3,1.3.2)} idam tarhi ātreyī iti . \\
\noindent \textbf{(P\_3,1.3.2)} atra hi paratvāt lopaḥ kitsvaram bādheta . \\
\noindent \textbf{(P\_3,1.3.2)} tasmāt eṣā paribhāṣā āśrayitavyā . \\
\noindent \textbf{(P\_3,1.3.2)} etasyām ca satyām śakyam pratyayasanniyogena ādyudāttatvam avaktum . \\
\noindent \textbf{(P\_3,1.3.3)} pratyayādyudāttatvāt dhātoḥ antaḥ . \\
\noindent \textbf{(P\_3,1.3.3)} pratyayādyudāttatvāt dhātoḥ antaḥ iti etat bhavati vipratiṣedhena . \\
\noindent \textbf{(P\_3,1.3.3)} pratyayādyudāttatvasya avakāśaḥ yatra anudāttā prakṛtiḥ . \\
\noindent \textbf{(P\_3,1.3.3)} samatvam simatvam . \\
\noindent \textbf{(P\_3,1.3.3)} dhātoḥ antaḥ iti asya avakāśaḥ yatra anudāttaḥ pratyayaḥ . \\
\noindent \textbf{(P\_3,1.3.3)} pacati paṭhati . \\
\noindent \textbf{(P\_3,1.3.3)} iha ubhayam prāpnoti . \\
\noindent \textbf{(P\_3,1.3.3)} gopāyati dhapāyati . \\
\noindent \textbf{(P\_3,1.3.3)} dhātoḥ antaḥ iti etat bhavati vipratiṣedhena. pitsvarāt titsvaraḥ ṭāpi . pitsvarāt titsvaraḥ ṭāpi bhavati vipratiṣedhena . \\
\noindent \textbf{(P\_3,1.3.3)} pitsvarasya avakāśaḥ . \\
\noindent \textbf{(P\_3,1.3.3)} pacati paṭhati . \\
\noindent \textbf{(P\_3,1.3.3)} titsvarasya avakāśaḥ . \\
\noindent \textbf{(P\_3,1.3.3)} kāryam hāryam . \\
\noindent \textbf{(P\_3,1.3.3)} iha ubhayam prāpnoti . \\
\noindent \textbf{(P\_3,1.3.3)} kāryā hāryā . \\
\noindent \textbf{(P\_3,1.3.3)} titsvaraḥ bhavati vipratiṣedhena . \\
\noindent \textbf{(P\_3,1.3.3)} citsvaraḥ cāpi pitsvarāt . citsvaraḥ cāpi pitsvarāt bhavati vipratiṣedhena . \\
\noindent \textbf{(P\_3,1.3.3)} citsvarasya avakāśaḥ . \\
\noindent \textbf{(P\_3,1.3.3)} calanaḥ copanaḥ . \\
\noindent \textbf{(P\_3,1.3.3)} pitsvarasya saḥ eva . \\
\noindent \textbf{(P\_3,1.3.3)} iha ubhayam prāpnoti . \\
\noindent \textbf{(P\_3,1.3.3)} āmbaṣṭhyā sauvīryā . \\
\noindent \textbf{(P\_3,1.3.3)} citsvaraḥ bhavati vipratiṣedhena . \\
\noindent \textbf{(P\_3,1.3.3)} na vā ādyutāttasya pratyayasañjñāsanniyogāt . na vā arthaḥ vipratiṣedhena . \\
\noindent \textbf{(P\_3,1.3.3)} kim kāraṇam . \\
\noindent \textbf{(P\_3,1.3.3)} ādyutāttasya pratyayasañjñāsanniyogāt . \\
\noindent \textbf{(P\_3,1.3.3)} pratyayasañjñāsanniyogena ādyudāttatve kṛte satiśiṣṭatvāt dhātusvaraḥ bhaviṣyati . \\
\noindent \textbf{(P\_3,1.3.3)} ayam ca api ayuktaḥ vipratiṣedhaḥ pitsvarasya titsvarasya ca . \\
\noindent \textbf{(P\_3,1.3.3)} kim kāraṇam . \\
\noindent \textbf{(P\_3,1.3.3)} ṭāpi svaritenaikādeśaḥ . \\
\noindent \textbf{(P\_3,1.3.3)} ṭāpi svaritena ekādeśaḥ bhavati . \\
\noindent \textbf{(P\_3,1.3.3)} idam iha sampradhāryam . \\
\noindent \textbf{(P\_3,1.3.3)} svaritatvam kriyatām ekādeśaḥ iti . \\
\noindent \textbf{(P\_3,1.3.3)} kim atra kartavyam . \\
\noindent \textbf{(P\_3,1.3.3)} paratvāt svaritatvam . \\
\noindent \textbf{(P\_3,1.3.3)} nityaḥ ekādeśaḥ . \\
\noindent \textbf{(P\_3,1.3.3)} kṛte api svaritatve prāpnoti akṛte api prāpnoti . \\
\noindent \textbf{(P\_3,1.3.3)} . svaritatvam api nityam . \\
\noindent \textbf{(P\_3,1.3.3)} kṛte api ekādeśe prāpnoti akṛte api . \\
\noindent \textbf{(P\_3,1.3.3)} anityam svaritatvam . \\
\noindent \textbf{(P\_3,1.3.3)} anyasya kṛte ekādeśe prāpnoti anyasya akṛte . \\
\noindent \textbf{(P\_3,1.3.3)} śabdāntarasya ca prāpnuvan vidhiḥ anityaḥ bhavati . \\
\noindent \textbf{(P\_3,1.3.3)} ekādeśaḥ api anityaḥ . \\
\noindent \textbf{(P\_3,1.3.3)} anyathāsvarasya kṛte svaritatve prāpnoti anyathāsvarasya akṛte . \\
\noindent \textbf{(P\_3,1.3.3)} svarabhinnasya ca prāpnuvan vidhiḥ anityaḥ bhavati . \\
\noindent \textbf{(P\_3,1.3.3)} antaraṅgaḥ tarhi ekādeśaḥ . \\
\noindent \textbf{(P\_3,1.3.3)} kā antaraṅgatā . \\
\noindent \textbf{(P\_3,1.3.3)} varṇau āśritya ekādeśaḥ padasya svaritatvam . \\
\noindent \textbf{(P\_3,1.3.3)} svaritatvam api antaraṅgam . \\
\noindent \textbf{(P\_3,1.3.3)} katham . \\
\noindent \textbf{(P\_3,1.3.3)} vakṣyati etat . \\
\noindent \textbf{(P\_3,1.3.3)} padagrahaṇam parimāṇārtham iti . \\
\noindent \textbf{(P\_3,1.3.3)} ubhayoḥ antaraṅgayoḥ paratvāt svaritatvam . \\
\noindent \textbf{(P\_3,1.3.3)} svaritatve kṛte āntaryataḥ svaridānudāttayoḥ svaritaḥ bhaviṣyati . \\
\noindent \textbf{(P\_3,1.3.3)} ayam ca api ayuktaḥ vipratiṣedhaḥ pitsvarasya citsvarasya ca . \\
\noindent \textbf{(P\_3,1.3.3)} kim kāraṇam . \\
\noindent \textbf{(P\_3,1.3.3)} cāpi citkaraṇāt . \\
\noindent \textbf{(P\_3,1.3.3)} cāpi citkaraṇasāmarthyāt antodāttatvam bhaviṣyati . \\
\noindent \textbf{(P\_3,1.5)} gupādiṣu anubandhakaraṇam kimartham . \\
\noindent \textbf{(P\_3,1.5)} gupādiṣu anubandhakaraṇam ātmanepadārtham . gupādiṣu anubandhāḥ kriyante ātmanepadam yathā syāt . \\
\noindent \textbf{(P\_3,1.5)} kriyamāṇeṣu api anubandheṣu ātmanepadam na eva prāpnoti . \\
\noindent \textbf{(P\_3,1.5)} kim kāraṇam . \\
\noindent \textbf{(P\_3,1.5)} sanā vyavahitatvāt . \\
\noindent \textbf{(P\_3,1.5)} pūrvavat sanaḥ iti evam bhaviṣyati . \\
\noindent \textbf{(P\_3,1.5)} pūrvavat sanaḥ iti ucyate . \\
\noindent \textbf{(P\_3,1.5)} na ce etebhyaḥ prāk sanaḥ ātmanepadam na api parasmaipadam paśyāmaḥ . \\
\noindent \textbf{(P\_3,1.5)} evam tarhi anubandhakaraṇasāmarthyāt bhaviṣyati . \\
\noindent \textbf{(P\_3,1.5)} atha vā avayave kṛtam liṅgam samudāyasya viśeṣakam bhaviṣyati . \\
\noindent \textbf{(P\_3,1.5)} tat yathā goḥ sakthani karṇe vā kṛtam liṅgam goḥ viśeṣakam bhavati . \\
\noindent \textbf{(P\_3,1.5)} yadi avayave kṛtam liṅgam samudāyasya viśeṣakam bhavati jugupsayati mīmāṃsayati iti atra api prāpnoti . \\
\noindent \textbf{(P\_3,1.5)} avayave kṛtam liṅgam kasya samudāyasya viśeṣakam bhavati . \\
\noindent \textbf{(P\_3,1.5)} yam samudāyam yaḥ avayavaḥ na vyabhicarati . \\
\noindent \textbf{(P\_3,1.5)} sanam ca na vhabhicarati . \\
\noindent \textbf{(P\_3,1.5)} ṇicam punaḥ vyabhicarati . \\
\noindent \textbf{(P\_3,1.5)} tat yathā tat yathā goḥ sakthani karṇe vā kṛtam liṅgam goḥ viśeṣakam bhavati na gomaṇḍalasya . \\
\noindent \textbf{(P\_3,1.6)} abhyāsadīrghatve avarṇasya dīrghaprasaṅgaḥ . \\
\noindent \textbf{(P\_3,1.6)} abhyāsadīrghatve avarṇasya dīrghatvam prāpnoti . \\
\noindent \textbf{(P\_3,1.6)} mīmāṃsate . \\
\noindent \textbf{(P\_3,1.6)} nanu ce ittve kṛte dīrghatvam bhaviṣyati . \\
\noindent \textbf{(P\_3,1.6)} katham punaḥ utpattisanniyogena dīrghatvam ucyamānam ittvam pratīkṣate . \\
\noindent \textbf{(P\_3,1.6)} atha katham abhyāsam pratīkṣate . \\
\noindent \textbf{(P\_3,1.6)} vacanāt abhyāsam pratīkṣate . \\
\noindent \textbf{(P\_3,1.6)} ittvam punaḥ na pratīkṣate . \\
\noindent \textbf{(P\_3,1.6)} na vā abhyāsavikāreṣu apavādasya utsargābādhakatvāt . \\
\noindent \textbf{(P\_3,1.6)} na vā eṣaḥ doṣaḥ . \\
\noindent \textbf{(P\_3,1.6)} kim kāraṇam . \\
\noindent \textbf{(P\_3,1.6)} abhyāsavikāreṣu apavādasya utsargābādhakatvāt . \\
\noindent \textbf{(P\_3,1.6)} abhyāsavikāreṣu apavādāḥ utsargān na bādhante iti evam dīrghatvam ucyamānam ittvam na bādhiṣyate . \\
\noindent \textbf{(P\_3,1.6)} atha vā mānbadhadānśanbhyaḥ ī ca abhyāsasya iti vakṣyāmi . \\
\noindent \textbf{(P\_3,1.6)} evam api halādiśeṣāpavādaḥ īkāraḥ prāpnoti . \\
\noindent \textbf{(P\_3,1.6)} ī ca acaḥ iti vakṣyāmi . \\
\noindent \textbf{(P\_3,1.6)} atha vā mānbadhadānśanbhyaḥ dīrghaḥ ca itaḥ abhyāsasya iti vakṣyāmi . \\
\noindent \textbf{(P\_3,1.6)} sidhyati . \\
\noindent \textbf{(P\_3,1.6)} sūtram tarhi bhidyate . \\
\noindent \textbf{(P\_3,1.6)} yathānyāsam eva astu . \\
\noindent \textbf{(P\_3,1.6)} nanu ca uktam abhyāsadīrghatve avarṇasya dīrghaprasaṅgaḥ iti . \\
\noindent \textbf{(P\_3,1.6)} parihṛtam etat na vā abhyāsavikāreṣu apavādasya utsargābādhakatvāt iti . \\
\noindent \textbf{(P\_3,1.6)} atha vā na evam vijñāyate dīrghaḥ ca abhyāsasya iti . \\
\noindent \textbf{(P\_3,1.6)} katham tarhi . \\
\noindent \textbf{(P\_3,1.6)} dīrghaḥ ca ābhyāsasya iti . \\
\noindent \textbf{(P\_3,1.6)} kim idam ābhyāsasya iti . \\
\noindent \textbf{(P\_3,1.6)} abhyāsavikāraḥ ābhyāsaḥ tasya iti . \\
\noindent \textbf{(P\_3,1.7.1)} dhātoḥ iti kimartham . \\
\noindent \textbf{(P\_3,1.7.1)} prakartum aicchat prācikīrṣat . \\
\noindent \textbf{(P\_3,1.7.1)} sopasargāt mā bhūt . \\
\noindent \textbf{(P\_3,1.7.1)} karmagrahaṇāt sanvidhau dhātugrahaṇānarthakyam . \\
\noindent \textbf{(P\_3,1.7.1)} karmagrahaṇāt sanvidhau dhātugrahaṇam anarthakam . \\
\noindent \textbf{(P\_3,1.7.1)} karmaṇaḥ samānakartṛkāt icchāyām vā sambhavati iti eva dhātoḥ utpattiḥ bhaviṣyati . \\
\noindent \textbf{(P\_3,1.7.1)} soparsargam vai karma . \\
\noindent \textbf{(P\_3,1.7.1)} tataḥ utpattiḥ prāpnoti . \\
\noindent \textbf{(P\_3,1.7.1)} sopasargam karma iti cet karmaviśeṣakatvāt upasargasya anupasargam karma . \\
\noindent \textbf{(P\_3,1.7.1)} sopasargam karma iti cet karmaviśeṣakaḥ upasargaḥ . \\
\noindent \textbf{(P\_3,1.7.1)} anupasargam hi karma . \\
\noindent \textbf{(P\_3,1.7.1)} avaśyam ca etat evam vijñeyam anupasargam karma iti . \\
\noindent \textbf{(P\_3,1.7.1)} sopasargasya hi karmatve dhātvadhikāre api sanaḥ avidhānam akarmatvāt . \\
\noindent \textbf{(P\_3,1.7.1)} yaḥ hi manyate sopasargam karma iti kriyamāṇe api tasya dhātugrahaṇe sanaḥ avidhiḥ syāt . \\
\noindent \textbf{(P\_3,1.7.1)} kim kāraṇam . \\
\noindent \textbf{(P\_3,1.7.1)} akarmatvāt . \\
\noindent \textbf{(P\_3,1.7.1)} idam tarhi prayojanam . \\
\noindent \textbf{(P\_3,1.7.1)} subantāt utpattiḥ mā bhūt . \\
\noindent \textbf{(P\_3,1.7.1)} subantāt ca aprasaṅgaḥ kyajādīnām apavādatvāt . \\
\noindent \textbf{(P\_3,1.7.1)} subantāt ca sanaḥ aprasaṅgaḥ . \\
\noindent \textbf{(P\_3,1.7.1)} kim kāraṇam . \\
\noindent \textbf{(P\_3,1.7.1)} kyajādīnām apavādatvāt . \\
\noindent \textbf{(P\_3,1.7.1)} subantāt kyajādayaḥ vidhīyante . \\
\noindent \textbf{(P\_3,1.7.1)} te apavādatvāt bādhakāḥ bhaviṣyanti . \\
\noindent \textbf{(P\_3,1.7.1)} anabhidhānāt vā . \\
\noindent \textbf{(P\_3,1.7.1)} atha vā anabhidhānāt subantāt utpattiḥ na bhaviṣyati . \\
\noindent \textbf{(P\_3,1.7.1)} na hi subantāt utpadyamānena sanā icchāyā abhidhānam syāt . \\
\noindent \textbf{(P\_3,1.7.1)} anabhidhānāt tataḥ utpattiḥ na bhaviṣyati . \\
\noindent \textbf{(P\_3,1.7.1)} iyam tāvat agatikā gatiḥ yat ucyate anabhidhānāt iti . \\
\noindent \textbf{(P\_3,1.7.1)} yat api ucyate subantāt ca aprasaṅgaḥ kyajādīnām apavādatvāt iti . \\
\noindent \textbf{(P\_3,1.7.1)} bhavet kasmāt cit aprasaṅgaḥ syāt ātmecchāyām . \\
\noindent \textbf{(P\_3,1.7.1)} parecchāyām tu prāpnoti : rājñaḥ putram icchati iti . \\
\noindent \textbf{(P\_3,1.7.1)} evam tarhi idam iha vyapadeśyam sat ācāryaḥ na vyapadiśati . \\
\noindent \textbf{(P\_3,1.7.1)} kim . \\
\noindent \textbf{(P\_3,1.7.1)} samānakartṛkāt iti ucyate . \\
\noindent \textbf{(P\_3,1.7.1)} na ca subantasya samānaḥ kartā asti . \\
\noindent \textbf{(P\_3,1.7.1)} evam api bhavet kasmāt cit aprasaṅgaḥ yasya kartā na asti . \\
\noindent \textbf{(P\_3,1.7.1)} iha tu prāpnoti : āsitum icchati śayitum icchati . \\
\noindent \textbf{(P\_3,1.7.1)} icchāyām arthe san vidhīyate icchārtheṣu ca tumun . \\
\noindent \textbf{(P\_3,1.7.1)} tatra tumunā uktatatvāt tasya arthasya san na bhaviṣyati . \\
\noindent \textbf{(P\_3,1.7.1)} evam api iha prāpnoti : āsanam icchati śayanam icchati iti . \\
\noindent \textbf{(P\_3,1.7.1)} iha yaḥ viśeṣaḥ upādhiḥ vā upādīyate dyotye tasmin tena bhavitavyam . \\
\noindent \textbf{(P\_3,1.7.1)} yaḥ ca iha arthaḥ gamyate āsitum icchati śayitum icchati svayam tām kriyām kartum icchati iti na asau iha gamyate āsanam icchati śayanam icchati iti . \\
\noindent \textbf{(P\_3,1.7.1)} anyasya api āsanam icchati iti eṣaḥ api arthaḥ gamyate . \\
\noindent \textbf{(P\_3,1.7.1)} avaśyam ca etat evam vijñeyam . \\
\noindent \textbf{(P\_3,1.7.1)} yaḥ hi manyate adyotye tasmin tena bhavitavyam iti kriyamāṇe api tasya dhātugrahaṇe iha prasajyeta : saṅgatam icchati devadattaḥ yajñadattena iti . \\
\noindent \textbf{(P\_3,1.7.1)} karmasamānakartṛkagrahaṇānarthakyam ca icchābhidhāne pratyayavidhānāt . \\
\noindent \textbf{(P\_3,1.7.1)} karmasamānakartṛkagrahaṇam ca anarthakam . \\
\noindent \textbf{(P\_3,1.7.1)} kim kāraṇam . \\
\noindent \textbf{(P\_3,1.7.1)} icchābhidhāne pratyayavidhānāt . \\
\noindent \textbf{(P\_3,1.7.1)} icchāyām abhidheyāyām san vidhīyate . \\
\noindent \textbf{(P\_3,1.7.1)} akarmaṇaḥ hi asamānakartṛkāt vā anabhidhānam . icchāyām abhidheyāyām san vidhīyate . \\
\noindent \textbf{(P\_3,1.7.1)} na ca akarmaṇaḥ asamānakartṛkāt vā utpadyamānena sanā icchāyā abhidhānam syāt . \\
\noindent \textbf{(P\_3,1.7.1)} anabhidhānāt tataḥ utpattiḥ na bhaviṣyati . \\
\noindent \textbf{(P\_3,1.7.1)} aṅgaparimāṇārtham tu . \\
\noindent \textbf{(P\_3,1.7.1)} aṅgaparimāṇārtham tarhi anyatarat kartavyam karmagrahaṇam dhātugrahaṇam vā . \\
\noindent \textbf{(P\_3,1.7.1)} aṅgaparimāṇam jñāsyāmi iti . \\
\noindent \textbf{(P\_3,1.7.1)} kim punaḥ atra jyāyaḥ . \\
\noindent \textbf{(P\_3,1.7.1)} dhātugrahaṇam eva jyāyaḥ . \\
\noindent \textbf{(P\_3,1.7.1)} aṅgaparimāṇam ca eva vijñātam bhavati . \\
\noindent \textbf{(P\_3,1.7.1)} api ca dhātoḥ vihitaḥ pratyayaḥ śeṣaḥ ārdhadhātukasañjñaḥ bhavati iti sanaḥ ārdhadhātukasañjñā siddhā bhavati . \\
\noindent \textbf{(P\_3,1.7.1)} yat ca api etat uktam karmagrahaṇāt sanvidhau dhātugrahaṇānarthakyam sopasargam karma iti cet karmaviśeṣakatvāt upasargasya anupasargam karma sopasargasya hi karmatve dhātvadhikāre api sanaḥ avidhānam akarmatvāt iti svapakṣaḥ anena varṇitaḥ . \\
\noindent \textbf{(P\_3,1.7.1)} yuktam iha draṣṭavyam kim nyāyyam karma iti . \\
\noindent \textbf{(P\_3,1.7.1)} etat ca atra yuktam yat sopasargam karma syāt . \\
\noindent \textbf{(P\_3,1.7.1)} nanu ca uktam sopasargasya hi karmatve dhātvadhikāre api sanaḥ avidhānam akarmatvāt iti . \\
\noindent \textbf{(P\_3,1.7.1)} na eṣaḥ doṣaḥ . \\
\noindent \textbf{(P\_3,1.7.1)} karmaṇaḥ iti na eṣā dhātusamānādhikaraṇā pañcamī . \\
\noindent \textbf{(P\_3,1.7.1)} karmaṇaḥ dhātoḥ iti . \\
\noindent \textbf{(P\_3,1.7.1)} kim tarhi . \\
\noindent \textbf{(P\_3,1.7.1)} avayavayogā eṣā ṣaṣṭhī . \\
\noindent \textbf{(P\_3,1.7.1)} karmaṇaḥ yaḥ dhātuḥ avayavaḥ . \\
\noindent \textbf{(P\_3,1.7.1)} yadi avayavayogā eṣā ṣaṣṭhīkevalāt utpattiḥ na prāpnoti . \\
\noindent \textbf{(P\_3,1.7.1)} cikīrṣati jihīrṣati iti . \\
\noindent \textbf{(P\_3,1.7.1)} eṣaḥ api vyapadeśivadbhāvena karmaṇaḥ dhātuḥ avayaḥ bhavati . \\
\noindent \textbf{(P\_3,1.7.1)} kāmam tarhi anena eva hetunā kyac api kartavyaḥ . \\
\noindent \textbf{(P\_3,1.7.1)} mahāntam putram icchati . \\
\noindent \textbf{(P\_3,1.7.1)} karmaṇaḥ yat subantam avayayaḥ iti . \\
\noindent \textbf{(P\_3,1.7.1)} na kartavyaḥ . \\
\noindent \textbf{(P\_3,1.7.1)} asāmarthyāt na bhaviṣyati . \\
\noindent \textbf{(P\_3,1.7.1)} katham asāmarthyam . \\
\noindent \textbf{(P\_3,1.7.1)} sāpekṣam asamartham bhavati iti . \\
\noindent \textbf{(P\_3,1.7.1)} vāvacanānarthakyam ca tatra nityatvāt sanaḥ . \\
\noindent \textbf{(P\_3,1.7.1)} vāvacanam ca anarthakam . \\
\noindent \textbf{(P\_3,1.7.1)} kim kāraṇam . \\
\noindent \textbf{(P\_3,1.7.1)} tatra nityatvāt sanaḥ . \\
\noindent \textbf{(P\_3,1.7.1)} iha hi dvau pakṣau vṛttipakṣaḥ avṛttipakṣaḥ ca . \\
\noindent \textbf{(P\_3,1.7.1)} svabhāvataḥ ca etat bhavati vākyam ca pratyayaḥ ca . \\
\noindent \textbf{(P\_3,1.7.1)} tatra svābhāvike vṛttiviṣaye nitye pratyaye prāpte vāvacanena kim anyat śakyam abhisambandhum anyat ataḥ sañjñāyāḥ . \\
\noindent \textbf{(P\_3,1.7.1)} na ca sañjñāyāḥ bhāvābhāvau iṣyete . \\
\noindent \textbf{(P\_3,1.7.1)} tasmāt na arthaḥ vāvacanena . \\
\noindent \textbf{(P\_3,1.7.2)} tumunantāt vā tasya ca lugvacanam . \\
\noindent \textbf{(P\_3,1.7.2)} tumunantāt vā san vaktavyaḥ tasya ca tumunaḥ luk vaktavyaḥ . \\
\noindent \textbf{(P\_3,1.7.2)} kartum icchati cikīrṣati . \\
\noindent \textbf{(P\_3,1.7.2)} liṅuttamāt vā . \\
\noindent \textbf{(P\_3,1.7.2)} liṅuttamāt vā san vaktavyaḥ tasya ca liṅaḥ luk vaktavyaḥ . \\
\noindent \textbf{(P\_3,1.7.2)} kuryām iti icchati cikīrṣati . \\
\noindent \textbf{(P\_3,1.7.2)} āśaṅkāyām acetaneṣu upasaṅkhyānam . \\
\noindent \textbf{(P\_3,1.7.2)} āśaṅkāyām acetaneṣu upasaṅkhyānam kartavyam . \\
\noindent \textbf{(P\_3,1.7.2)} aśmā luluṭhiṣate . \\
\noindent \textbf{(P\_3,1.7.2)} kūlam pipatiṣati iti . \\
\noindent \textbf{(P\_3,1.7.2)} kim punaḥ kāraṇam na sidhyati . \\
\noindent \textbf{(P\_3,1.7.2)} evam manyate . \\
\noindent \textbf{(P\_3,1.7.2)} cetanāvataḥ etat bhavati icchā iti . \\
\noindent \textbf{(P\_3,1.7.2)} kūlam ca acetanam . \\
\noindent \textbf{(P\_3,1.7.2)} acetanagrahaṇena na arthaḥ . \\
\noindent \textbf{(P\_3,1.7.2)} āśaṅkāyām iti eva . \\
\noindent \textbf{(P\_3,1.7.2)} idam api siddham bhavati . \\
\noindent \textbf{(P\_3,1.7.2)} śvā mumūrṣati . \\
\noindent \textbf{(P\_3,1.7.2)} na vā tulyakāraṇatvāt icchayāḥ hi pravṛttitaḥ upalabdhiḥ . \\
\noindent \textbf{(P\_3,1.7.2)} na vā kartavyam . \\
\noindent \textbf{(P\_3,1.7.2)} kim kāraṇam . \\
\noindent \textbf{(P\_3,1.7.2)} tulyakāraṇatvāt . \\
\noindent \textbf{(P\_3,1.7.2)} tulyam hi kāraṇam cetanāvati devadatte kūle ca acetane . \\
\noindent \textbf{(P\_3,1.7.2)} kim kāraṇam . \\
\noindent \textbf{(P\_3,1.7.2)} icchayāḥ hi pravṛttitaḥ upalabdhiḥ . \\
\noindent \textbf{(P\_3,1.7.2)} icchayāḥ hi pravṛttitaḥ upalabdhiḥ bhavati . \\
\noindent \textbf{(P\_3,1.7.2)} yaḥ api asu kaṭam cikīrṣuḥ bhavati na asau āghoṣayati . \\
\noindent \textbf{(P\_3,1.7.2)} kaṭam kariṣyāmi iti . \\
\noindent \textbf{(P\_3,1.7.2)} kim tarhi . \\
\noindent \textbf{(P\_3,1.7.2)} sannaddham rajjukīlakpūlapāṇim dṛṣṭvā tataḥ icchā gamyate . \\
\noindent \textbf{(P\_3,1.7.2)} kūlasya api pipatiṣataḥ loṣṭāḥ śīryante bhidā jāyante deśāt deśāntaram upasaṅkrāmati . \\
\noindent \textbf{(P\_3,1.7.2)} śvānaḥ khalu api mumūrṣavaḥ ekāntaśīlāḥ śūnākṣāḥ ca bhavanti . \\
\noindent \textbf{(P\_3,1.7.2)} upamānāt vā siddham . \\
\noindent \textbf{(P\_3,1.7.2)} upamānāt vā siddham etat . \\
\noindent \textbf{(P\_3,1.7.2)} katham . \\
\noindent \textbf{(P\_3,1.7.2)} luluṭhiṣate iva luluṭhiṣate . \\
\noindent \textbf{(P\_3,1.7.2)} pipatiṣati iva pipatiṣati . \\
\noindent \textbf{(P\_3,1.7.2)} na tiṅantena upamānam asti . \\
\noindent \textbf{(P\_3,1.7.2)} evam tarhi icchā iva icchā . \\
\noindent \textbf{(P\_3,1.7.2)} sarvasya vā cetanāvattvāt . \\
\noindent \textbf{(P\_3,1.7.2)} atha vā sarvam cetanāvat . \\
\noindent \textbf{(P\_3,1.7.2)} evam hi āha . \\
\noindent \textbf{(P\_3,1.7.2)} kaṃsakāḥ sarpanti . \\
\noindent \textbf{(P\_3,1.7.2)} śirīṣaḥ adhaḥ svapiti . \\
\noindent \textbf{(P\_3,1.7.2)} suvarcalā ādityam anu paryeti . \\
\noindent \textbf{(P\_3,1.7.2)} āskanda kapilaka iti ukte tṛṇam āskandati . \\
\noindent \textbf{(P\_3,1.7.2)} ayaskāntam ayaḥ saṅkrāmati . \\
\noindent \textbf{(P\_3,1.7.2)} ṛṣiḥ paṭhati śrṇota grāvāṇaḥ \\
\noindent \textbf{(P\_3,1.7.3)} ime iṣavaḥ bahavaḥ paṭhyante . \\
\noindent \textbf{(P\_3,1.7.3)} tatra na jñāyate kasya ayam arthe san vidhīyate iti . \\
\noindent \textbf{(P\_3,1.7.3)} iṣeḥ chatvabhāvinaḥ . \\
\noindent \textbf{(P\_3,1.7.3)} yadi evam kartum anvicchati kartum anveṣaṇā atra api prāpnoti . \\
\noindent \textbf{(P\_3,1.7.3)} evam tarhi yasya striyām icchā iti etat rūpam nipātyate . \\
\noindent \textbf{(P\_3,1.7.3)} kasya ca etat nipātyate . \\
\noindent \textbf{(P\_3,1.7.3)} kāntikarmaṇaḥ . \\
\noindent \textbf{(P\_3,1.7.3)} atha iha grāmam gantum icchati iti kasya kim karma . \\
\noindent \textbf{(P\_3,1.7.3)} iṣeḥ ubhe karmaṇī . \\
\noindent \textbf{(P\_3,1.7.3)} yadi evam grāmam gantum icchati grāmāya gantum icchati iti gatyarthakarmaṇi dvitīyācaturthyau na prāpnutaḥ . \\
\noindent \textbf{(P\_3,1.7.3)} evam tarhi gameḥ grāmaḥ karma iṣeḥ gamiḥ karma . \\
\noindent \textbf{(P\_3,1.7.3)} evam api iṣyate grāmaḥ gantum iti parasādhane utpadyamānena lena grāmasya abhidhānam na prāpnoti . \\
\noindent \textbf{(P\_3,1.7.3)} evam tarhi gameḥ grāmaḥ karma iṣeḥ ubhe karmaṇī . \\
\noindent \textbf{(P\_3,1.7.3)} atha sanantāt sanā bhavitavyam : cikīrṣitum icchati jihīrṣitum icchati iti . \\
\noindent \textbf{(P\_3,1.7.3)} na bhavitavyam . \\
\noindent \textbf{(P\_3,1.7.3)} kim kāraṇam . \\
\noindent \textbf{(P\_3,1.7.3)} arthagatyarthaḥ śabdaprayogaḥ . \\
\noindent \textbf{(P\_3,1.7.3)} artham sampratyāyayiṣyāmi iti śabdaḥ prayujyate . \\
\noindent \textbf{(P\_3,1.7.3)} tatra ekena uktatvāt tasya arthasya aparasya prayogeṇa na bhavitavyam . \\
\noindent \textbf{(P\_3,1.7.3)} kim kāraṇam . \\
\noindent \textbf{(P\_3,1.7.3)} uktārthānām aprayogaḥ . \\
\noindent \textbf{(P\_3,1.7.3)} na tarhi idānīm idam bhavati : eṣitum icchati eṣiṣiṣati iti . \\
\noindent \textbf{(P\_3,1.7.3)} asti atra viśeṣaḥ . \\
\noindent \textbf{(P\_3,1.7.3)} ekasya atra iṣeḥ iṣiḥ sādhanam vartamānakālaḥ ca pratyayaḥ . \\
\noindent \textbf{(P\_3,1.7.3)} aparasya bāhyam sādhanam sarvakālaḥ ca pratyayaḥ . \\
\noindent \textbf{(P\_3,1.7.3)} iha api tarhi ekasya iṣeḥ karotiviṣiṣṭaḥ iṣiḥ sādhanam vartamānakālaḥ ca pratyayaḥ . \\
\noindent \textbf{(P\_3,1.7.3)} aparasya bāhyam sādhanam sarvakālaḥ ca pratyayaḥ . \\
\noindent \textbf{(P\_3,1.7.3)} yena eva khalu api hetunā etat vākyam bhavati cikīrṣitum icchati jihīrṣitum icchati iti tena eva hetunā vṛttiḥ api prāpnoti . \\
\noindent \textbf{(P\_3,1.7.3)} tasmāt sanantāt sanaḥ pratiṣedhaḥ vaktavyaḥ . \\
\noindent \textbf{(P\_3,1.7.3)} tam ca api bruvatā iṣisanaḥ iti vaktavyam . \\
\noindent \textbf{(P\_3,1.7.3)} bhavati hi jugupsiṣate mīmāṃsiṣate iti . \\
\noindent \textbf{(P\_3,1.7.3)} śaiṣikāt matubarthīyāt śaiṣikaḥ matubarthikaḥ sarūpaḥ pratyayaḥ na iṣṭaḥ . \\
\noindent \textbf{(P\_3,1.7.3)} sanantāt na san iṣyate . \\
\noindent \textbf{(P\_3,1.8.1)} kimarthaḥ cakāraḥ . \\
\noindent \textbf{(P\_3,1.8.1)} svarārthaḥ . \\
\noindent \textbf{(P\_3,1.8.1)} citaḥ antaḥ udāttaḥ bhavati iti antodāttatvam yathā syāt . \\
\noindent \textbf{(P\_3,1.8.1)} na etat asti prayojanam . \\
\noindent \textbf{(P\_3,1.8.1)} ekāc ayam . \\
\noindent \textbf{(P\_3,1.8.1)} tatra na arthaḥ svarārthena cakāreṇa anubandhena . \\
\noindent \textbf{(P\_3,1.8.1)} pratyayasvareṇa eva siddham . \\
\noindent \textbf{(P\_3,1.8.1)} viśeṣaṇārthaḥ tarhi . \\
\noindent \textbf{(P\_3,1.8.1)} kva viśeṣaṇārthena arthaḥ . \\
\noindent \textbf{(P\_3,1.8.1)} asya cvau kyaci ca iti . \\
\noindent \textbf{(P\_3,1.8.1)} kye ca iti ucyamāne api kākaḥ śyenāyate atra api prasayjeta . \\
\noindent \textbf{(P\_3,1.8.1)} na etat asti . \\
\noindent \textbf{(P\_3,1.8.1)} tadanubandhakagrahaṇe atadanubandhakasya grahaṇam na iti evam etasya na bhaviṣyati . \\
\noindent \textbf{(P\_3,1.8.1)} sāmānyagrahaṇāvighātārthaḥ tarhi . \\
\noindent \textbf{(P\_3,1.8.1)} kva ca sāmānyagrahaṇāvighātārthena arthaḥ . \\
\noindent \textbf{(P\_3,1.8.1)} naḥ kye iti . \\
\noindent \textbf{(P\_3,1.8.1)} atha ātmangrahaṇam kimartham . \\
\noindent \textbf{(P\_3,1.8.1)} ātmecchāyām yathā syāt . \\
\noindent \textbf{(P\_3,1.8.1)} parecchāyām mā bhūt iti . \\
\noindent \textbf{(P\_3,1.8.1)} rājñaḥ putram icchati iti . \\
\noindent \textbf{(P\_3,1.8.1)} kriyamāṇe api ātmagrahaṇe parecchāyām prāpnoti . \\
\noindent \textbf{(P\_3,1.8.1)} kim kāraṇam . \\
\noindent \textbf{(P\_3,1.8.1)} ātmanaḥ iti iyam kartari ṣaṣṭhī . \\
\noindent \textbf{(P\_3,1.8.1)} icchā iti akāraḥ bhāve . \\
\noindent \textbf{(P\_3,1.8.1)} saḥ yadi eva ātmanaḥ icchā atha api parasya ātmecchā eva asau bhavati . \\
\noindent \textbf{(P\_3,1.8.1)} na ātmagrahaṇena icchā abhisambadhyate . \\
\noindent \textbf{(P\_3,1.8.1)} kim tarhi . \\
\noindent \textbf{(P\_3,1.8.1)} subantam abhisambadhyate . \\
\noindent \textbf{(P\_3,1.8.1)} ātmanaḥ yat subantam iti . \\
\noindent \textbf{(P\_3,1.8.1)} yadi ātmagrahaṇam kriyate chandasi parecchāyām na prāpnoti . \\
\noindent \textbf{(P\_3,1.8.1)} ma tvā vṛkāḥ aghāyavaḥ vidan . \\
\noindent \textbf{(P\_3,1.8.1)} tasmāt na arthaḥ ātmagrahaṇena . \\
\noindent \textbf{(P\_3,1.8.1)} iha kasmāt na bhavati : rājñaḥ putram icchati iti . \\
\noindent \textbf{(P\_3,1.8.1)} asāmarthyāt . \\
\noindent \textbf{(P\_3,1.8.1)} katham asāmarthyam . \\
\noindent \textbf{(P\_3,1.8.1)} sāpekṣam asamartham bhavati iti . \\
\noindent \textbf{(P\_3,1.8.1)} chandasi api tarhi na prāpnoti . \\
\noindent \textbf{(P\_3,1.8.1)} ma tvā vṛkāḥ aghāyavaḥ vidan . \\
\noindent \textbf{(P\_3,1.8.1)} asti atra viśeṣaḥ . \\
\noindent \textbf{(P\_3,1.8.1)} antareṇa api atra tṛtīyasya padasya prayogam parecchā gamyate . \\
\noindent \textbf{(P\_3,1.8.1)} katham punaḥ antareṇa api atra tṛtīyasya padasya prayogam parecchā gamyate . \\
\noindent \textbf{(P\_3,1.8.1)} te ca eva vṛkāḥ evamātmakaḥ hiṃsrāḥ . \\
\noindent \textbf{(P\_3,1.8.1)} kaḥ ca ātmanaḥ agham eṣitum arhati . \\
\noindent \textbf{(P\_3,1.8.1)} ataḥ antareṇa api atra tṛtīyasya padasya prayogam parecchā gamyate . \\
\noindent \textbf{(P\_3,1.8.1)} yathā eva tarhi chandasi aghaśabdāt parecchāyām khyac bhavati evam bhāṣāyām api prāpnoti . \\
\noindent \textbf{(P\_3,1.8.1)} agham icchati iti . \\
\noindent \textbf{(P\_3,1.8.1)} tasmāt ātmagrahaṇam kartavyam . \\
\noindent \textbf{(P\_3,1.8.1)} chandasi katham . \\
\noindent \textbf{(P\_3,1.8.1)} ācāryapravṛttiḥ jñāpayati bhavati chandasi aghaśabdāt parecchāyām kyac iti yat ayam aśvāghasyāt iti kyaci pratkṛte ītvabādhanārtham ākāram śāsti . \\
\noindent \textbf{(P\_3,1.8.1)} atha subgrahaṇam kimartham . \\
\noindent \textbf{(P\_3,1.8.1)} subantāt utpattiḥ yatha syāt . \\
\noindent \textbf{(P\_3,1.8.1)} prātipadikāt mā bhūt iti . \\
\noindent \textbf{(P\_3,1.8.1)} na etat asti prayojanam . \\
\noindent \textbf{(P\_3,1.8.1)} na asti atra viśeṣaḥ subantāt utpattau satyām prātipadikāt vā ayam asti viśeṣaḥ . \\
\noindent \textbf{(P\_3,1.8.1)} subantāt utpattau satyām padasañjñā siddhā bhavati . \\
\noindent \textbf{(P\_3,1.8.1)} prātipadikāt utpattau satyām padasañjñā na prāpnoti . \\
\noindent \textbf{(P\_3,1.8.1)} nanu ca prātipadikāt utpattau satyām padasañjñā siddhā . \\
\noindent \textbf{(P\_3,1.8.1)} katham . \\
\noindent \textbf{(P\_3,1.8.1)} ārabhyate naḥ kye iti . \\
\noindent \textbf{(P\_3,1.8.1)} tat ca avaśyam kartavyam subantāt utpattau satyām niyamārtham . \\
\noindent \textbf{(P\_3,1.8.1)} tat eva prātipadikāt utpattau satyām vidhyartham bhaviṣyati . \\
\noindent \textbf{(P\_3,1.8.1)} idam tarhi prayojanam . \\
\noindent \textbf{(P\_3,1.8.1)} subantāt utpattiḥ yatha syāt . \\
\noindent \textbf{(P\_3,1.8.1)} dhātoḥ mā bhūt iti . \\
\noindent \textbf{(P\_3,1.8.1)} etat api na asti prayojanam . \\
\noindent \textbf{(P\_3,1.8.1)} dhātoḥ san vidhīyate . \\
\noindent \textbf{(P\_3,1.8.1)} saḥ bādhakaḥ bhaviṣyati . \\
\noindent \textbf{(P\_3,1.8.1)} anavakāśāḥ hi vidhayaḥ bādhakāḥ bhavanti sāvakāśaḥ ca san . \\
\noindent \textbf{(P\_3,1.8.1)} kaḥ avakāśaḥ . \\
\noindent \textbf{(P\_3,1.8.1)} parecchā . \\
\noindent \textbf{(P\_3,1.8.1)} na parecchāyām sanā bhavitavyam . \\
\noindent \textbf{(P\_3,1.8.1)} kim kāraṇam . \\
\noindent \textbf{(P\_3,1.8.1)} samānakartṛkāt iti ucyate . \\
\noindent \textbf{(P\_3,1.8.1)} yāvat ca iha ātmagrahaṇam tāvat tatra samānakartṛkagrahaṇam . \\
\noindent \textbf{(P\_3,1.8.1)} idam tarhi prayojanam . \\
\noindent \textbf{(P\_3,1.8.1)} subantāt utpattiḥ yatha syāt . \\
\noindent \textbf{(P\_3,1.8.1)} vākyāt māt bhūt iti . \\
\noindent \textbf{(P\_3,1.8.1)} mahāntam putram icchati iti . \\
\noindent \textbf{(P\_3,1.8.1)} na vā bhavati mahāputrīyati iti . \\
\noindent \textbf{(P\_3,1.8.1)} bhavati yadā etat vākyam bhavati . \\
\noindent \textbf{(P\_3,1.8.1)} mahān putraḥ mahāputraḥ . \\
\noindent \textbf{(P\_3,1.8.1)} mahāputram icchati mahāputrīyati iti . \\
\noindent \textbf{(P\_3,1.8.1)} yadā tu etat vākyam bhavati mahāntam putram icchati iti tadā na bhavitavyam tadā ca prāpnoti . \\
\noindent \textbf{(P\_3,1.8.1)} tadā mā bhūt iti . \\
\noindent \textbf{(P\_3,1.8.2)} atha kriyamāṇe api subgrahaṇe kasmāt eva atra na bhavati . \\
\noindent \textbf{(P\_3,1.8.2)} subantam hi etat vākyam . \\
\noindent \textbf{(P\_3,1.8.2)} na etat subantam . \\
\noindent \textbf{(P\_3,1.8.2)} katham . \\
\noindent \textbf{(P\_3,1.8.2)} pratyayagrahaṇe yasmāt tadādeḥ grahaṇam bhavati iti . \\
\noindent \textbf{(P\_3,1.8.2)} atha yat atra subantam tasmāt utpattiḥ kasmāt na bhavati . \\
\noindent \textbf{(P\_3,1.8.2)} samānādhikaraṇānām sarvatra avṛttiḥ ayogāt ekena . \\
\noindent \textbf{(P\_3,1.8.2)} samānādhikaraṇānām sarvatra eva vṛttiḥ na bhavati . \\
\noindent \textbf{(P\_3,1.8.2)} kva sarvatra . \\
\noindent \textbf{(P\_3,1.8.2)} samāsavidhau pratyayavidhau . \\
\noindent \textbf{(P\_3,1.8.2)} samāsavidhau tāvat . \\
\noindent \textbf{(P\_3,1.8.2)} ṛddhasya rājñaḥ puruṣaḥ . \\
\noindent \textbf{(P\_3,1.8.2)} mahat kaṣṭam śritaḥ iti . \\
\noindent \textbf{(P\_3,1.8.2)} pratyayavidhau . \\
\noindent \textbf{(P\_3,1.8.2)} ṛddhasya upagoḥ apatyam . \\
\noindent \textbf{(P\_3,1.8.2)} mahāntam putram icchati . \\
\noindent \textbf{(P\_3,1.8.2)} iti . \\
\noindent \textbf{(P\_3,1.8.2)} kim punaḥ kāraṇam samānādhikaraṇānām sarvatra vṛttiḥ na bhavati . \\
\noindent \textbf{(P\_3,1.8.2)} ayogāt ekena . \\
\noindent \textbf{(P\_3,1.8.2)} na hi ekena padena yogaḥ bhavati . \\
\noindent \textbf{(P\_3,1.8.2)} iha tāvat ṛddhasya rājñaḥ puruṣaḥ iti ṣaṣṭhyantena subantena sāmarthye sati samāsaḥ vidhīyate . \\
\noindent \textbf{(P\_3,1.8.2)} yat ca atra ṣaṣthyantam na tasya subantena sāmarthyam . \\
\noindent \textbf{(P\_3,1.8.2)} yasya ca sāmarthyam na tat ṣaṣṭhyantam . \\
\noindent \textbf{(P\_3,1.8.2)} vākyam tat . \\
\noindent \textbf{(P\_3,1.8.2)} ṛddhasya upagoḥ apatyam iti ca . \\
\noindent \textbf{(P\_3,1.8.2)} ṣaṣṭhīsamarthāt apatyena yoge pratyayaḥ vidhīyate . \\
\noindent \textbf{(P\_3,1.8.2)} yat ca atra ṣaṣṭḥīsamartham na tasya apatatyena yogaḥ yasya ca aptatyena yogaḥ na tat ṣaṣṭhyantam . \\
\noindent \textbf{(P\_3,1.8.2)} vākyam tat . \\
\noindent \textbf{(P\_3,1.8.2)} samānādhikaraṇānām iti ucyate . \\
\noindent \textbf{(P\_3,1.8.2)} atha vyadhikaraṇānām katham . \\
\noindent \textbf{(P\_3,1.8.2)} rājñaḥ putram icchati iti . \\
\noindent \textbf{(P\_3,1.8.2)} evam tarhi idam paṭhitavyam . \\
\noindent \textbf{(P\_3,1.8.2)} saviśeṣaṇānām sarvatra avṛttiḥ ayogāt ekena . \\
\noindent \textbf{(P\_3,1.8.2)} dvitīyānupapattiḥ tu . \\
\noindent \textbf{(P\_3,1.8.2)} dvitīyā tu na upapadyate . \\
\noindent \textbf{(P\_3,1.8.2)} mahāntam putram icchati iti . \\
\noindent \textbf{(P\_3,1.8.2)} kim kāraṇam . \\
\noindent \textbf{(P\_3,1.8.2)} na putraḥ iṣikarma . \\
\noindent \textbf{(P\_3,1.8.2)} yadi putraḥ na iṣikarma na ca avaśyam dvitīyā eva . \\
\noindent \textbf{(P\_3,1.8.2)} kim tarhi . \\
\noindent \textbf{(P\_3,1.8.2)} sarvāḥ dvitīyādayaḥ vibhaktayaḥ . \\
\noindent \textbf{(P\_3,1.8.2)} mahatā putreṇa kṛtam . \\
\noindent \textbf{(P\_3,1.8.2)} mahate putrāya dehi . \\
\noindent \textbf{(P\_3,1.8.2)} mahaḥ putrāt ānaya . \\
\noindent \textbf{(P\_3,1.8.2)} mahataḥ putrasya svam . \\
\noindent \textbf{(P\_3,1.8.2)} mahati putre nidhehi . \\
\noindent \textbf{(P\_3,1.8.2)} tasmāt na evam śakyam vaktum na putraḥ iṣikarma iti . \\
\noindent \textbf{(P\_3,1.8.2)} putra eva iṣikarma . \\
\noindent \textbf{(P\_3,1.8.2)} tatsāmānādhikaraṇyāt dvitīyādayaḥ bhaviṣyanti . \\
\noindent \textbf{(P\_3,1.8.2)} vṛttiḥ tarhi kasmāt na bhavati . \\
\noindent \textbf{(P\_3,1.8.2)} saviśeṣaṇānām vṛttiḥ na vṛttasya vā viśeṣaṇam na prayujyate iti vaktavyam . \\
\noindent \textbf{(P\_3,1.8.2)} yadi saviśeṣaṇānām vṛttiḥ na vṛttasya vā viśeṣaṇam na prayujyate iti ucyate muṇḍayati māṇavakam iti atra vṛttiḥ na prāpnoti . \\
\noindent \textbf{(P\_3,1.8.2)} amuṇḍādīnām iti vaktavyam . \\
\noindent \textbf{(P\_3,1.8.2)} tat tarhi vaktavyam saviśeṣaṇānām vṛttiḥ na vṛttasya vā viśeṣaṇam na prayujyate amuṇḍādīnām iti . \\
\noindent \textbf{(P\_3,1.8.2)} na vaktavyam . \\
\noindent \textbf{(P\_3,1.8.2)} vṛttiḥ kasmāt na bhavati mahāntam putram icchati iti . \\
\noindent \textbf{(P\_3,1.8.2)} agamakatvāt . \\
\noindent \textbf{(P\_3,1.8.2)} iha samānārthena vākyena bhavitavyam pratyayāntena ca . \\
\noindent \textbf{(P\_3,1.8.2)} yaḥ ca iha arthaḥ vākyena gamyate mahāntam putram icchati iti na asau jātu cit pratyayāntena gamyate mahāntam putrīyati iti . \\
\noindent \textbf{(P\_3,1.8.2)} etasmāt hetoḥ brūmaḥ agamakatvāt iti . \\
\noindent \textbf{(P\_3,1.8.2)} na brūmaḥ apaśabdaḥ syāt iti . \\
\noindent \textbf{(P\_3,1.8.2)} yatra ca gamakatvam bhavati tatra vṛttiḥ . \\
\noindent \textbf{(P\_3,1.8.2)} tat yathā muṇḍayati māṇavakam iti . \\
\noindent \textbf{(P\_3,1.8.3)} atha asya kyajantasya kāni sādhanāni bhavanti . \\
\noindent \textbf{(P\_3,1.8.3)} bhāvaḥ kartā ca . \\
\noindent \textbf{(P\_3,1.8.3)} atha karma . \\
\noindent \textbf{(P\_3,1.8.3)} na asti karma . \\
\noindent \textbf{(P\_3,1.8.3)} nanu ca ayam iṣiḥ sakarmakaḥ yasya ayam arthe kyac vidhīyate . \\
\noindent \textbf{(P\_3,1.8.3)} abhihitam tat karma antarbhūtam dhātvarthaḥ sampannaḥ . \\
\noindent \textbf{(P\_3,1.8.3)} na ca idānīm anyat karma asti yena sakarmakaḥ syāt . \\
\noindent \textbf{(P\_3,1.8.3)} katham tarhi ayam sakarmakaḥ bhavati aputram putram iva ācarati putrīyati māṇavakam iti . \\
\noindent \textbf{(P\_3,1.8.3)} asti atra viśeṣaḥ . \\
\noindent \textbf{(P\_3,1.8.3)} dve hi atra karmaṇī upamānakarma upameyakarma ca . \\
\noindent \textbf{(P\_3,1.8.3)} upamānakarma antarbhūtam . \\
\noindent \textbf{(P\_3,1.8.3)} upameyena karmaṇā sakarmakaḥ bhavati . \\
\noindent \textbf{(P\_3,1.8.3)} tat yathā . \\
\noindent \textbf{(P\_3,1.8.3)} api kākaḥ śyenāyate iti atra dvau kartārau upamānakartā ca upameyakartā ca. upamānakartā antarbhūtaḥ . \\
\noindent \textbf{(P\_3,1.8.3)} upemeyakartrā sakrtṛkaḥ bhavati . \\
\noindent \textbf{(P\_3,1.8.3)} ayam tarhi katham sakarmakaḥ bhavati . \\
\noindent \textbf{(P\_3,1.8.3)} muṇḍayati māṇavakam iti . \\
\noindent \textbf{(P\_3,1.8.3)} atra api dve karmaṇī sāmānyakarma viśeṣakarma ca . \\
\noindent \textbf{(P\_3,1.8.3)} sāmānyakarma antarbhūtam . \\
\noindent \textbf{(P\_3,1.8.3)} viśeṣakarmaṇā sakarmakaḥ bhavati . \\
\noindent \textbf{(P\_3,1.8.3)} nanu ca vṛttyā eva atra na bhavitavyam . \\
\noindent \textbf{(P\_3,1.8.3)} kim kāraṇam . \\
\noindent \textbf{(P\_3,1.8.3)} asāmarthyāt . \\
\noindent \textbf{(P\_3,1.8.3)} katham asāmarthyam . \\
\noindent \textbf{(P\_3,1.8.3)} sāpekṣam asamartham bhavati iti . \\
\noindent \textbf{(P\_3,1.8.3)} na eṣaḥ doṣaḥ . \\
\noindent \textbf{(P\_3,1.8.3)} na atra ubhau karotiyuktau muṇḍaḥ māṇavakaḥ ca . \\
\noindent \textbf{(P\_3,1.8.3)} na hi māṇavakaḥ kriyate . \\
\noindent \textbf{(P\_3,1.8.3)} yadā ca ubhau karotiyuktau bhavataḥ na bhavati tadā vṛttiḥ . \\
\noindent \textbf{(P\_3,1.8.3)} tat yathā balīvardam karoti muṇḍam ca enam karoti iti . \\
\noindent \textbf{(P\_3,1.8.3)} kāmam tarhi anena eva hetunā kyac api kartavyaḥ māṇavakam muṇḍam icchati iti . \\
\noindent \textbf{(P\_3,1.8.3)} na ubhau iṣiyuktau iti . \\
\noindent \textbf{(P\_3,1.8.3)} na kartavyaḥ . \\
\noindent \textbf{(P\_3,1.8.3)} ubhau atra iṣiyuktau muṇḍaḥ māṇavakaḥ ca . \\
\noindent \textbf{(P\_3,1.8.3)} katham . \\
\noindent \textbf{(P\_3,1.8.3)} na hi asau mauṇḍyamātreṇa santoṣam karoti . \\
\noindent \textbf{(P\_3,1.8.3)} māṇavakastham asau mauṇḍyam icchati . \\
\noindent \textbf{(P\_3,1.8.3)} iha api tarhi na prāpnoti muṇḍayati māṇavakam iti . \\
\noindent \textbf{(P\_3,1.8.3)} atra api hi ubhau karotiyukta muṇḍaḥ māṇavakaḥ ca . \\
\noindent \textbf{(P\_3,1.8.3)} na hi asau mauṇḍyamātreṇa santoṣam karoti . \\
\noindent \textbf{(P\_3,1.8.3)} māṇavakastham asu mauṇḍyam nirvartayati . \\
\noindent \textbf{(P\_3,1.8.3)} evam tarhi muṇḍādayaḥ guaṇavacanāḥ . \\
\noindent \textbf{(P\_3,1.8.3)} guṇavacanāḥ ca sāpekṣāḥ . \\
\noindent \textbf{(P\_3,1.8.3)} vacanāt sāpekṣāṇām api vṛttiḥ bhaviṣyati . \\
\noindent \textbf{(P\_3,1.8.3)} atha vā dhātavaḥ eva muṇḍādayaḥ . \\
\noindent \textbf{(P\_3,1.8.3)} na na eva hi arthāḥ ādiśyante kriyāvacanatā ca gamyate . \\
\noindent \textbf{(P\_3,1.8.3)} atha vā na idam ubhayam yugapat bhavati vākyam ca pratyayaḥ ca . \\
\noindent \textbf{(P\_3,1.8.3)} yadā vākyam na tadā pratyayaḥ . \\
\noindent \textbf{(P\_3,1.8.3)} yadā pratyayaḥ sāmānyena tadā vṛttiḥ . \\
\noindent \textbf{(P\_3,1.8.3)} tatra avśyam viśeṣārthinā viśeṣaḥ anuprayoktavyaḥ . \\
\noindent \textbf{(P\_3,1.8.3)} muṇḍayati . \\
\noindent \textbf{(P\_3,1.8.3)} kam . \\
\noindent \textbf{(P\_3,1.8.3)} māṇavakam iti . \\
\noindent \textbf{(P\_3,1.8.3)} muṇḍaviśiṣṭena vā karotina tam āptum icchati . \\
\noindent \textbf{(P\_3,1.8.3)} atha vā uktam etat . \\
\noindent \textbf{(P\_3,1.8.3)} na atra vyāpāraḥ anugantavyaḥ iti . \\
\noindent \textbf{(P\_3,1.8.3)} gamaktvāt iha vṛttiḥ bhaviṣyati . \\
\noindent \textbf{(P\_3,1.8.3)} muṇḍayati māṇavakam iti . \\
\noindent \textbf{(P\_3,1.8.3)} atha iha kyacā bhavitavyam . \\
\noindent \textbf{(P\_3,1.8.3)} iṣṭaḥ putraḥ . \\
\noindent \textbf{(P\_3,1.8.3)} iṣyate putraḥ iti . \\
\noindent \textbf{(P\_3,1.8.3)} ke cit tāvat āhuḥ na bhavitavyam it . \\
\noindent \textbf{(P\_3,1.8.3)} kim kāraṇam . \\
\noindent \textbf{(P\_3,1.8.3)} svaśabdena uktatvāt iti . \\
\noindent \textbf{(P\_3,1.8.3)} apare āhuḥ : bhavitavyam iti . \\
\noindent \textbf{(P\_3,1.8.3)} kim kāraṇam . \\
\noindent \textbf{(P\_3,1.8.3)} dhātvarthe ayam kyac vidhīyate . \\
\noindent \textbf{(P\_3,1.8.3)} saḥ ca dhātvarthaḥ kena cit eva śabdena nirdeṣṭavyaḥ iti . \\
\noindent \textbf{(P\_3,1.8.3)} ihabhavantaḥ tu āhuḥ na bhavitavyam iti . \\
\noindent \textbf{(P\_3,1.8.3)} kim kāraṇam . \\
\noindent \textbf{(P\_3,1.8.3)} iha samānārthena vākyena bhavitavyam pratyayāntena ca . \\
\noindent \textbf{(P\_3,1.8.3)} yaḥ ca iha arthaḥ vākyena gamyate iṣṭaḥ putraḥ iṣyate putraḥ iti na asau jātu cit pratyayāntena gamyate . \\
\noindent \textbf{(P\_3,1.8.4)} kyaci māntāvyayapratiṣedhaḥ . \\
\noindent \textbf{(P\_3,1.8.4)} kyaci māntāvyayānām pratiṣedhaḥ vaktavyaḥ . \\
\noindent \textbf{(P\_3,1.8.4)} iha māt bhūt . \\
\noindent \textbf{(P\_3,1.8.4)} idam icchati . \\
\noindent \textbf{(P\_3,1.8.4)} kim icchati . \\
\noindent \textbf{(P\_3,1.8.4)} uccaiḥ icchati . \\
\noindent \textbf{(P\_3,1.8.4)} nīcaiḥ icchati . \\
\noindent \textbf{(P\_3,1.8.4)} gosamānākṣaranāntāt iti eke . \\
\noindent \textbf{(P\_3,1.8.4)} gām icchati gavyati . \\
\noindent \textbf{(P\_3,1.8.4)} samānākṣarāt . \\
\noindent \textbf{(P\_3,1.8.4)} dadhīyati madhati kartrīyati hartrīyati . \\
\noindent \textbf{(P\_3,1.8.4)} nāntāt . \\
\noindent \textbf{(P\_3,1.8.4)} rājīyati takṣīyati . \\
\noindent \textbf{(P\_3,1.9)} kimarthaḥ cakāraḥ . \\
\noindent \textbf{(P\_3,1.9)} svarārthaḥ . \\
\noindent \textbf{(P\_3,1.9)} citaḥ antaḥ udāttaḥ bhavati iti antodāttatvam yathā syāt . \\
\noindent \textbf{(P\_3,1.9)} na etat asti prayojanam . \\
\noindent \textbf{(P\_3,1.9)} dhātusvareṇa api etat siddham . \\
\noindent \textbf{(P\_3,1.9)} kakārasya tarhi itsañjñāparitrāṇārthaḥ āditaḥ cakāraḥ kartavyaḥ . \\
\noindent \textbf{(P\_3,1.9)} ataḥ uttaram paṭhati . \\
\noindent \textbf{(P\_3,1.9)} kāmyacaḥ citkaraṇānarthakyam kasya idarthābhāvāt . \\
\noindent \textbf{(P\_3,1.9)} kāmyacaḥ citkaraṇam anarthakam . \\
\noindent \textbf{(P\_3,1.9)} kakārasya tarhi itsañjñā kasmāt na bhavati . \\
\noindent \textbf{(P\_3,1.9)} idarthābhāvāt . \\
\noindent \textbf{(P\_3,1.9)} itkāryābhāvāt atra itsañjñā na bhaviṣyati . \\
\noindent \textbf{(P\_3,1.9)} nanu ca lopaḥ eva itkāryam . \\
\noindent \textbf{(P\_3,1.9)} akāryam lopaḥ . \\
\noindent \textbf{(P\_3,1.9)} iha hi śabdasya kāryārthaḥ vā bhavati upadeśaḥ śravaṇārthaḥ vā . \\
\noindent \textbf{(P\_3,1.9)} karyam ca iha na asti . \\
\noindent \textbf{(P\_3,1.9)} kārye asati yadi śravaṇam api na syāt upadeśaḥ anarthakaḥ syāt . \\
\noindent \textbf{(P\_3,1.9)} idam tarhi itkāryam . \\
\noindent \textbf{(P\_3,1.9)} agnicitkamyati . \\
\noindent \textbf{(P\_3,1.9)} kiti iti guṇapratiṣedhaḥ yathā syāt . \\
\noindent \textbf{(P\_3,1.9)} na etat asti prayojanam . \\
\noindent \textbf{(P\_3,1.9)} sārvadhātukārdhadhātukayoḥ aṅgasya guṇaḥ ucyate . \\
\noindent \textbf{(P\_3,1.9)} dhātoḥ ca vihitaḥ pratyayaḥ śeṣaḥ ārdhadhātukasañjñām labhate . \\
\noindent \textbf{(P\_3,1.9)} na ca ayam dhātoḥ vidhīyate . \\
\noindent \textbf{(P\_3,1.9)} idam tarhi . \\
\noindent \textbf{(P\_3,1.9)} upayaṭkāmyati . \\
\noindent \textbf{(P\_3,1.9)} kiti iti samprasāraṇam yathā syāt .. etat api na asti prayojanam . \\
\noindent \textbf{(P\_3,1.9)} yajādibhiḥ atra kitam viśeṣayiṣyāmaḥ . \\
\noindent \textbf{(P\_3,1.9)} yajādīnām yaḥ kit iti . \\
\noindent \textbf{(P\_3,1.9)} kaḥ ca yajādīnām kit . \\
\noindent \textbf{(P\_3,1.9)} yajādibhyaḥ yaḥ vihitaḥ iti . \\
\noindent \textbf{(P\_3,1.9)} atha api katham cit itkāryam syāt . \\
\noindent \textbf{(P\_3,1.9)} evam api na doṣaḥ . \\
\noindent \textbf{(P\_3,1.9)} kriyate nyāse eva dvicakārakaḥ nirdeśaḥ . \\
\noindent \textbf{(P\_3,1.9)} supaḥ ātmanaḥ kyac ckāmyat ca iti . \\
\noindent \textbf{(P\_3,1.9)} atha vā chāndasam etat . \\
\noindent \textbf{(P\_3,1.9)} dṛṣṭānuvidhiḥ chandasi bhavati . \\
\noindent \textbf{(P\_3,1.9)} na ca atra samprasāraṇam dṛśyate . \\
\noindent \textbf{(P\_3,1.10)} adhikaraṇāt ca . \\
\noindent \textbf{(P\_3,1.10)} adhikaraṇāt ca iti vaktavyam . \\
\noindent \textbf{(P\_3,1.10)} prāsādayati kuṭyām kuṭīyati prāsāde iti atra api yathā syāt . \\
\noindent \textbf{(P\_3,1.11.1)} salopasanniyogena ayam kyaṅ vidhīyate . \\
\noindent \textbf{(P\_3,1.11.1)} tena yatra eva salopaḥ tatra eva syāt . \\
\noindent \textbf{(P\_3,1.11.1)} payāyate . \\
\noindent \textbf{(P\_3,1.11.1)} iha na syāt . \\
\noindent \textbf{(P\_3,1.11.1)} api kākaḥ śyenāyate . \\
\noindent \textbf{(P\_3,1.11.1)} na eṣaḥ doṣaḥ . \\
\noindent \textbf{(P\_3,1.11.1)} pradhānśiṣṭaḥ kyaṅ . \\
\noindent \textbf{(P\_3,1.11.1)} anvācayaśiṣṭaḥ salopaḥ . \\
\noindent \textbf{(P\_3,1.11.1)} yatra ca sakāram paśyasi iti . \\
\noindent \textbf{(P\_3,1.11.1)} tat yatha . \\
\noindent \textbf{(P\_3,1.11.1)} kaḥ cit uktaḥ grāme bhikṣām cara devadattam ca ānaya iti . \\
\noindent \textbf{(P\_3,1.11.1)} saḥ grāme bhikṣām carati . \\
\noindent \textbf{(P\_3,1.11.1)} yadi devadattam paśyati tam api ānayati . \\
\noindent \textbf{(P\_3,1.11.1)} salopaḥ vā . \\
\noindent \textbf{(P\_3,1.11.1)} salopaḥ vā iti vaktavyam . \\
\noindent \textbf{(P\_3,1.11.1)} payāyate payasyate . \\
\noindent \textbf{(P\_3,1.11.1)} ojopsarasoḥ nityam . \\
\noindent \textbf{(P\_3,1.11.1)} ojopsarasoḥ nityam salopaḥ vaktavyaḥ . \\
\noindent \textbf{(P\_3,1.11.1)} ojāyamānam yaḥ ahim jaghāna . \\
\noindent \textbf{(P\_3,1.11.1)} apsarāyate . \\
\noindent \textbf{(P\_3,1.11.1)} apraḥ āha salopaḥ apsarasaḥ eva . \\
\noindent \textbf{(P\_3,1.11.1)} payasyate iti eva bhavitavyam iti . \\
\noindent \textbf{(P\_3,1.11.1)} katham ojāyamānam yaḥ ahim jaghana iti . \\
\noindent \textbf{(P\_3,1.11.1)} chāndasaḥ prayogaḥ . \\
\noindent \textbf{(P\_3,1.11.1)} chandasi ca dṛṣṭānuvidhiḥ vidhīyate . \\
\noindent \textbf{(P\_3,1.11.2)} ācāre galbhaklībahoḍebhyaḥ kvip vā . \\
\noindent \textbf{(P\_3,1.11.2)} ācāre galbhaklībahoḍebhyaḥ kvip vā vaktavyaḥ . \\
\noindent \textbf{(P\_3,1.11.2)} avagalbhate avagalbhāyate . \\
\noindent \textbf{(P\_3,1.11.2)} klība . \\
\noindent \textbf{(P\_3,1.11.2)} viklībate viklībāyate . \\
\noindent \textbf{(P\_3,1.11.2)} klība . \\
\noindent \textbf{(P\_3,1.11.2)} hoḍa . \\
\noindent \textbf{(P\_3,1.11.2)} vihoḍate vihoḍāyate . \\
\noindent \textbf{(P\_3,1.11.2)} kim prayojanam . \\
\noindent \textbf{(P\_3,1.11.2)} kriyāvacanatā yathā syāt . \\
\noindent \textbf{(P\_3,1.11.2)} na etat asti prayojanam . \\
\noindent \textbf{(P\_3,1.11.2)} dhātavaḥ eva galbhādayaḥ . \\
\noindent \textbf{(P\_3,1.11.2)} na ca eva hi arthāḥ ādiśyante kriyāvacanatā ca gamyate . \\
\noindent \textbf{(P\_3,1.11.2)} idam tarhi prayojanam . \\
\noindent \textbf{(P\_3,1.11.2)} avagalbhā viklībā vihoḍā . \\
\noindent \textbf{(P\_3,1.11.2)} a pratyayāt iti akāraḥ yathā syāt . \\
\noindent \textbf{(P\_3,1.11.2)} mā bhūt evam . \\
\noindent \textbf{(P\_3,1.11.2)} guroḥ ca halaḥ iti evam bhaviṣyati . \\
\noindent \textbf{(P\_3,1.11.2)} idam tarhi . \\
\noindent \textbf{(P\_3,1.11.2)} avagalbhām cakre . \\
\noindent \textbf{(P\_3,1.11.2)} viklībām cakre . \\
\noindent \textbf{(P\_3,1.11.2)} vihoḍām cakre . \\
\noindent \textbf{(P\_3,1.11.2)} kāspratyayāt ām amantre iti ām yathā syāt . \\
\noindent \textbf{(P\_3,1.11.2)} aparaḥ āha : sarvaprātipadikebhyaḥ ācāre kvip vaktavyaḥ aśvati gardabhati iti evamartham . \\
\noindent \textbf{(P\_3,1.11.2)} na tarhi idānīm galbhādyanukramaṇam kartavyam . \\
\noindent \textbf{(P\_3,1.11.2)} kartavyam ca . \\
\noindent \textbf{(P\_3,1.11.2)} kim prayojanam . \\
\noindent \textbf{(P\_3,1.11.2)} ātmanepadārtham anubandhān āsaṅkṣyāmi iti . \\
\noindent \textbf{(P\_3,1.11.2)} galbha klība hoḍa . \\
\noindent \textbf{(P\_3,1.12.1)} halaḥ lopasanniyogena ayam kyaṅ vidhīyate . \\
\noindent \textbf{(P\_3,1.12.1)} tena yatra eva halaḥ lopaḥ tatra eva prasajyeta . \\
\noindent \textbf{(P\_3,1.12.1)} na eṣaḥ doṣaḥ . \\
\noindent \textbf{(P\_3,1.12.1)} pradhānaśiṣṭaḥ kyaṅ . \\
\noindent \textbf{(P\_3,1.12.1)} anvācayaśiṣṭaḥ halaḥ lopaḥ . \\
\noindent \textbf{(P\_3,1.12.1)} yatra ca halam paśyasi iti . \\
\noindent \textbf{(P\_3,1.12.2)} bhṛśādiṣu abhūtatadbhāvagrahaṇam . \\
\noindent \textbf{(P\_3,1.12.2)} bhṛśādiṣu abhūtatadbhāvagrahaṇam kartavyam . \\
\noindent \textbf{(P\_3,1.12.2)} iha mā bhūt . \\
\noindent \textbf{(P\_3,1.12.2)} kva divā bhṛśāḥ bhavanti iti . \\
\noindent \textbf{(P\_3,1.12.2)} cvipratiṣedhānarthakyam ca bhavatyarthe kyaṅvacanāt . cvipratiṣedhaḥ ca anarthakaḥ . \\
\noindent \textbf{(P\_3,1.12.2)} kim kāraṇam . \\
\noindent \textbf{(P\_3,1.12.2)} bhavatyarthe kyaṅvacanāt . \\
\noindent \textbf{(P\_3,1.12.2)} bhavatyarthe hi kyaṅ vidhīyate . \\
\noindent \textbf{(P\_3,1.12.2)} bhavatiyoge cvividhānam . bhavatinā yoge cviḥ vidhīyate . \\
\noindent \textbf{(P\_3,1.12.2)} tatra cvinā uktatvāt tasya arthasya kyaṅ na bhaviṣyati . \\
\noindent \textbf{(P\_3,1.12.2)} ḍājantāt api tarhi na prāpnoti . \\
\noindent \textbf{(P\_3,1.12.2)} paṭapaṭāyate . \\
\noindent \textbf{(P\_3,1.12.2)} ḍāc api hi bhavatinā yoge vidhīyate . \\
\noindent \textbf{(P\_3,1.12.2)} bhavatyarthe kyaṣ . \\
\noindent \textbf{(P\_3,1.12.2)} ḍāci vacanaprāmāṇyāt . \\
\noindent \textbf{(P\_3,1.12.2)} ḍāci vacanaprāmāṇyāt bhaviṣyati . \\
\noindent \textbf{(P\_3,1.12.2)} kim vacanaprāmāṇyam . \\
\noindent \textbf{(P\_3,1.12.2)} lohitādiḍājbhyaḥ kyaṣ iti . \\
\noindent \textbf{(P\_3,1.12.2)} iha kim cit akriyamāṇam codyate kim cit kriyamāṇam pratyākhyāyate . \\
\noindent \textbf{(P\_3,1.12.2)} saḥ sūtrabhedaḥ kṛtaḥ bhavati . \\
\noindent \textbf{(P\_3,1.12.2)} yathānyāsam eva astu . \\
\noindent \textbf{(P\_3,1.12.2)} nanu ca uktam iha kasmāt na bhavati kva divā bhṛśāḥ bhavanti iti . \\
\noindent \textbf{(P\_3,1.12.2)} nañivayuktam anyasadṛśādhikaraṇe tathā hi arthagatiḥ . \\
\noindent \textbf{(P\_3,1.12.2)} nañyuktam ivayuktam va yat kim cit iha dṛśyate tatra anyasmin tatsadṛśe kāryam vijñāyate . \\
\noindent \textbf{(P\_3,1.12.2)} tathā hi arthaḥ gamyate . \\
\noindent \textbf{(P\_3,1.12.2)} abrāhmaṇam ānaya iti ukte brāhmaṇasadṛśaḥ ānīyate . \\
\noindent \textbf{(P\_3,1.12.2)} na asau loṣṭam ānīya kṛtī bhavati . \\
\noindent \textbf{(P\_3,1.12.2)} evam iha api acveḥ iti cvipratiṣedhāt anyasmin acvyante cvisadṛśe kāryam vijñāsyate . \\
\noindent \textbf{(P\_3,1.12.2)} kim ca ataḥ anyat advyantam cvisadṛsam . \\
\noindent \textbf{(P\_3,1.12.2)} abhūtatadbhāvaḥ . \\
\noindent \textbf{(P\_3,1.12.3)} iha kāḥ cit prakṛtayaḥ sopasargāḥ paṭhyante : abhimanas , sumanas , unmanas , durmanas . \\
\noindent \textbf{(P\_3,1.12.3)} tatra vicāryate : bhṛśādiṣu upasargaḥ pratyayārthaviśeṣaṇam vā syāt : abhibhavatau subhavatau udbhavatau durbhavatau iti . \\
\noindent \textbf{(P\_3,1.12.3)} prakṛtyarthaviśeṣaṇam vā . \\
\noindent \textbf{(P\_3,1.12.3)} abhimanasśabdāt sumanasśabdāt unmanasśabdāt durmanasśabdāt iti . \\
\noindent \textbf{(P\_3,1.12.3)} yuktam punaḥ idam vicārayatum . \\
\noindent \textbf{(P\_3,1.12.3)} nanu tena asandigdhena prakṛtyarthaviśeṣaṇam bhavitavyam yāvatā prāk prakṛteḥ paṭhyante . \\
\noindent \textbf{(P\_3,1.12.3)} yadi hi pratyayārthaviśeṣaṇam syāt prāk bhavateḥ paṭhyeran . \\
\noindent \textbf{(P\_3,1.12.3)} na ime śakyāḥ prāk bhavateḥ paṭhitum . \\
\noindent \textbf{(P\_3,1.12.3)} evam viśiṣṭe hi pratyayāṛthe bhṛśādimātrāt utpattiḥ prasajyeta . \\
\noindent \textbf{(P\_3,1.12.3)} tasmāt na evam śakyam kartum . \\
\noindent \textbf{(P\_3,1.12.3)} na cet evam jāyate vicāraṇā . \\
\noindent \textbf{(P\_3,1.12.3)} kaḥ ca atra viśeṣaḥ . \\
\noindent \textbf{(P\_3,1.12.3)} bhṛśādiṣu upasargaḥ pratyayārthaviśeṣaṇam iti cet svare doṣaḥ . \\
\noindent \textbf{(P\_3,1.12.3)} bhṛśādiṣu upasargaḥ pratyayārthaviśeṣaṇam iti cet svare doṣaḥ bhavati . \\
\noindent \textbf{(P\_3,1.12.3)} abhimanāyate . \\
\noindent \textbf{(P\_3,1.12.3)} tiṅ atiṅaḥ iti nighātaḥ prasajyate . \\
\noindent \textbf{(P\_3,1.12.3)} astu tarhi prakṛtyarthaviśeṣaṇam . \\
\noindent \textbf{(P\_3,1.12.3)} sopasargāt iti cet aṭi doṣaḥ . sopasargāt iti cet aṭi doṣaḥ bhavati . \\
\noindent \textbf{(P\_3,1.12.3)} svamanayata iti . \\
\noindent \textbf{(P\_3,1.12.3)} atyalpam idam ucyate : aṭi doṣaḥ bhavati iti . \\
\noindent \textbf{(P\_3,1.12.3)} aḍlyavdvirvacaneṣu iti vaktavyam . \\
\noindent \textbf{(P\_3,1.12.3)} aṭi : udāhṛtam . \\
\noindent \textbf{(P\_3,1.12.3)} lyapi : sumanāyya . \\
\noindent \textbf{(P\_3,1.12.3)} dvirvacane : abhimimanāyiṣate . \\
\noindent \textbf{(P\_3,1.12.3)} na eṣaḥ doṣaḥ . \\
\noindent \textbf{(P\_3,1.12.3)} avaśyam saṅgrāmayateḥ sopasargāt utpattiḥ vaktavyā asaṅgrāmayata śūraḥ iti evamartham . \\
\noindent \textbf{(P\_3,1.12.3)} tat niyamārtham bhaviṣyati . \\
\noindent \textbf{(P\_3,1.12.3)} saṅgrāmayateḥ eva sopasargāt na anyasmāt sopasargāt iti . \\
\noindent \textbf{(P\_3,1.12.3)} yadi niyamaḥ kriyate svaraḥ na sidhyati . \\
\noindent \textbf{(P\_3,1.12.3)} evam tarhi bhṛśādiṣu upasargasya parāṅgavadbhāvam vakṣyāmi . \\
\noindent \textbf{(P\_3,1.12.3)} yadi parāṅgavadbhāvaḥ ucyate aḍlyavdvirvacanāni na sidhyanti . \\
\noindent \textbf{(P\_3,1.12.3)} svaravidhau iti vakṣyāmi . \\
\noindent \textbf{(P\_3,1.12.3)} evam ca kṛtvā astu pratyayārthaviśeṣaṇam . \\
\noindent \textbf{(P\_3,1.12.3)} nanu ca uktam bhṛśādiṣu upasargaḥ pratyayārthaviśeṣaṇam iti cet svare doṣaḥ iti . \\
\noindent \textbf{(P\_3,1.12.3)} svare parāṅgavadbhāvena parihṛtam . \\
\noindent \textbf{(P\_3,1.12.3)} ayam tarhi pratyayārthaviśeṣaṇe sati doṣaḥ . \\
\noindent \textbf{(P\_3,1.12.3)} kyaṅā uktatvāt tasya arthasya upasargasya prayogaḥ na prāpnoti . \\
\noindent \textbf{(P\_3,1.12.3)} kim kāraṇam . \\
\noindent \textbf{(P\_3,1.12.3)} uktārthānām aprayogaḥ iti . \\
\noindent \textbf{(P\_3,1.12.3)} tat yathā . \\
\noindent \textbf{(P\_3,1.12.3)} api kākaḥ śyenāyate iti kyaṅā uktatvāt ācārārthasya āṅaḥ prayogaḥ na bhavati . \\
\noindent \textbf{(P\_3,1.12.3)} asti atra viśeṣaḥ . \\
\noindent \textbf{(P\_3,1.12.3)} ekena atra viśiṣṭe pratyayārthe pratyayaḥ utpadyate iha punaḥ anekena . \\
\noindent \textbf{(P\_3,1.12.3)} tatra manāyate iti ukte sandehaḥ syāt abhibhavatau subhavatau durbhavatau iti . \\
\noindent \textbf{(P\_3,1.12.3)} tatra asandehārtham upasargaḥ prayujyate . \\
\noindent \textbf{(P\_3,1.12.3)} yatra tarhi ekena . \\
\noindent \textbf{(P\_3,1.12.3)} utpucchayate . \\
\noindent \textbf{(P\_3,1.12.3)} atra api anekena . \\
\noindent \textbf{(P\_3,1.12.3)} pucchāt udasane pucchāt vyasane pucchāt paryasane iti . \\
\noindent \textbf{(P\_3,1.13.1)} kimarthaḥ kakāraḥ . \\
\noindent \textbf{(P\_3,1.13.1)} kṅiti iti guṇapratiṣedhaḥ yathā syāt . \\
\noindent \textbf{(P\_3,1.13.1)} na etat asti prayojanam . \\
\noindent \textbf{(P\_3,1.13.1)} sārvadhātukārdhadhātukayoḥ aṅgasya guṇaḥ ucyate . \\
\noindent \textbf{(P\_3,1.13.1)} dhātoḥ ca vihitaḥ pratyayaḥ śeṣaḥ ārdhadhātukasañjñām labhate . \\
\noindent \textbf{(P\_3,1.13.1)} na ca ayam dhātoḥ vidhīyate . \\
\noindent \textbf{(P\_3,1.13.1)} lohitādīni prātipadikāni . \\
\noindent \textbf{(P\_3,1.13.1)} sāmānyagrahaṇārthaḥ tarhi . \\
\noindent \textbf{(P\_3,1.13.1)} kva sāmānyagrahaṇena arthaḥ . \\
\noindent \textbf{(P\_3,1.13.1)} naḥ kye iti . \\
\noindent \textbf{(P\_3,1.13.1)} na ayam nāntāt vidhīyate . \\
\noindent \textbf{(P\_3,1.13.1)} iha tarhi . \\
\noindent \textbf{(P\_3,1.13.1)} yasya halaḥ kyasya vibhāṣā iti . \\
\noindent \textbf{(P\_3,1.13.1)} na ayam halantāt vidhīyate . \\
\noindent \textbf{(P\_3,1.13.1)} iha tarhi . \\
\noindent \textbf{(P\_3,1.13.1)} āpatyayasya ca taddhite anāti kyacvyoḥ ca iti . \\
\noindent \textbf{(P\_3,1.13.1)} na ayam āpatyāt vidhīyate . \\
\noindent \textbf{(P\_3,1.13.1)} iha tarhi . \\
\noindent \textbf{(P\_3,1.13.1)} kyāt chandasi iti . \\
\noindent \textbf{(P\_3,1.13.1)} yāt chandasi iti etāvat vaktavyam caraṇyūḥ turaṇyuḥ bhuraṇyuḥ iti evamartham . \\
\noindent \textbf{(P\_3,1.13.1)} idam tarhi prayojanam . \\
\noindent \textbf{(P\_3,1.13.1)} yat tat akṛtyakāre iti dīrghatvam tatra kṅidgrahaṇam anuvartate . \\
\noindent \textbf{(P\_3,1.13.1)} tat iha api yathā syāt . \\
\noindent \textbf{(P\_3,1.13.1)} lohitāyate . \\
\noindent \textbf{(P\_3,1.13.1)} kim punaḥ kāraṇam tatra kṅidgrahaṇam anuvartate . \\
\noindent \textbf{(P\_3,1.13.1)} iha mā bhūt . \\
\noindent \textbf{(P\_3,1.13.1)} uruyā dhṛṣṇuyā iti . \\
\noindent \textbf{(P\_3,1.13.1)} yadi kṅidgrahaṇam anuvartate itryam iti pituḥ rīṅbhāvaḥ na prāpnoti . \\
\noindent \textbf{(P\_3,1.13.1)} rīṅbhāve kṅidgrahaṇam nivartiṣyate . \\
\noindent \textbf{(P\_3,1.13.1)} yadi nivartate katham asūyā vasūyā ca yamāmahe . \\
\noindent \textbf{(P\_3,1.13.1)} asūyateḥ asūyā vasūyateḥ vasūyā . \\
\noindent \textbf{(P\_3,1.13.1)} atha vā chāndasam etat . \\
\noindent \textbf{(P\_3,1.13.1)} dṛṣṭānuvidhiḥ ca chandasi bhavati iti . \\
\noindent \textbf{(P\_3,1.13.1)} yadi chāndasatvam hetuḥ na arthaḥ kṅidgrahaṇena anuvartamānena . \\
\noindent \textbf{(P\_3,1.13.1)} kasmāt na bhavati uruyā dhṛṣṇuyā iti . \\
\noindent \textbf{(P\_3,1.13.1)} chāndasatvāt . \\
\noindent \textbf{(P\_3,1.13.1)} atha vā astu atra dīrghatvam . \\
\noindent \textbf{(P\_3,1.13.1)} chāndasam hrasvatvam bhaviṣyati . \\
\noindent \textbf{(P\_3,1.13.1)} tat yathā upagāyantu mām patnayaḥ garbhiṇayaḥ yuvatayaḥ iti . \\
\noindent \textbf{(P\_3,1.13.1)} atha kimarthaḥ ṣakāraḥ . \\
\noindent \textbf{(P\_3,1.13.1)} viśeṣaṇārthaḥ . \\
\noindent \textbf{(P\_3,1.13.1)} kva viśeṣaṇārthena arthaḥ . \\
\noindent \textbf{(P\_3,1.13.1)} vā kyaṣaḥ iti . \\
\noindent \textbf{(P\_3,1.13.1)} vā yāt iti hi ucyamāne ataḥ api prasajyeta . \\
\noindent \textbf{(P\_3,1.13.1)} na etat asi prayojanam . \\
\noindent \textbf{(P\_3,1.13.1)} parasmaipadam iti ucyate . \\
\noindent \textbf{(P\_3,1.13.1)} na ca ataḥ parasmaipadam na api ātmanepadam paśyāmaḥ . \\
\noindent \textbf{(P\_3,1.13.1)} sāmānyagrahaṇāvighātārthaḥ tarhi bhaviṣyati . \\
\noindent \textbf{(P\_3,1.13.1)} kva sāmānyagrahaṇāvighātārthena arthaḥ . \\
\noindent \textbf{(P\_3,1.13.1)} kyāt chandasi iti . \\
\noindent \textbf{(P\_3,1.13.1)} yāt chandasi iti evam vaktavyam caraṇyūḥ turaṇyuḥ bhuraṇyuḥ iti evamartham . \\
\noindent \textbf{(P\_3,1.13.2)} lohitaḍājbhyaḥ kyaṣvacanam . \\
\noindent \textbf{(P\_3,1.13.2)} lohitaḍājbhyaḥ kyaṣ vaktavyaḥ . \\
\noindent \textbf{(P\_3,1.13.2)} lohitāyati lohitāyate paṭapaṭāyati paṭapaṭāyate . \\
\noindent \textbf{(P\_3,1.13.2)} atha anyāni lohitādīni . \\
\noindent \textbf{(P\_3,1.13.2)} bhṛśādiṣu itarāṇi . \\
\noindent \textbf{(P\_3,1.13.2)} bhṛśādiṣu itarāṇi paṭhitavyāni . \\
\noindent \textbf{(P\_3,1.13.2)} kim prayojanam . \\
\noindent \textbf{(P\_3,1.13.2)} ṅitaḥ iti ātmanepadam yathā syāt iti . \\
\noindent \textbf{(P\_3,1.14)} kaṣṭāya iti kim nipātyate . \\
\noindent \textbf{(P\_3,1.14)} kaṣṭaśabdāt caturthīsamarthāt kramaṇe anārjave kyṅ nipātyate . \\
\noindent \textbf{(P\_3,1.14)} kaṣṭāya karmaṇe krāmati kaṣṭāyate . \\
\noindent \textbf{(P\_3,1.14)} atyalpam idam ucyate : kaṣṭāya iti . \\
\noindent \textbf{(P\_3,1.14)} sattrakakṣakaṣṭagahanebhyaḥ kaṇvacikīrṣāyām . \\
\noindent \textbf{(P\_3,1.14)} sattrakakṣakaṣṭagahanebhyaḥ kaṇvacikīrṣāyām iti vaktavyam . \\
\noindent \textbf{(P\_3,1.14)} sattrāyate . \\
\noindent \textbf{(P\_3,1.14)} sattra. kakaṣa. kakṣāyate . \\
\noindent \textbf{(P\_3,1.14)} kaṣṭa . \\
\noindent \textbf{(P\_3,1.14)} kaṣṭāyate . \\
\noindent \textbf{(P\_3,1.14)} kaṣṭa . \\
\noindent \textbf{(P\_3,1.14)} gahana . \\
\noindent \textbf{(P\_3,1.14)} gahanāyate . \\
\noindent \textbf{(P\_3,1.14)} aparaḥ āha:. sattrādibhyaḥ caturthyantebhyaḥ kramaṇe anārjave kyaṅ vaktavyaḥ . \\
\noindent \textbf{(P\_3,1.14)} etāni eva udāharaṇāni . \\
\noindent \textbf{(P\_3,1.14)} sattrādibhyaḥ iti kimartham . \\
\noindent \textbf{(P\_3,1.14)} kuṭilāya krāmati anuvākāya . \\
\noindent \textbf{(P\_3,1.14)} caturthyantebhyaḥ iti kimartham . \\
\noindent \textbf{(P\_3,1.14)} ajaḥ kaṣṭam krāmati . \\
\noindent \textbf{(P\_3,1.14)} tat tarhi vaktavyam . \\
\noindent \textbf{(P\_3,1.14)} na vaktavyam . \\
\noindent \textbf{(P\_3,1.14)} na etat pratyayāntanipātanam . \\
\noindent \textbf{(P\_3,1.14)} kim tarhi . \\
\noindent \textbf{(P\_3,1.14)} tādarthye eṣā caturthī . \\
\noindent \textbf{(P\_3,1.14)} kaṣṭāya yat prātipadikam . \\
\noindent \textbf{(P\_3,1.14)} kaṣṭārthe yat prātipadikam iti . \\
\noindent \textbf{(P\_3,1.15)} romanthe iti ucyate . \\
\noindent \textbf{(P\_3,1.15)} kaḥ romanthaḥ nāma . \\
\noindent \textbf{(P\_3,1.15)} udgīrṇasya vā avagīrṇasya vā manthaḥ romanthaḥ iti . \\
\noindent \textbf{(P\_3,1.15)} yadi evam hanucalane iti vaktavyam . \\
\noindent \textbf{(P\_3,1.15)} iha mā bhūt . \\
\noindent \textbf{(P\_3,1.15)} kīṭaḥ romatham vartayati . \\
\noindent \textbf{(P\_3,1.15)} tat tarhi vaktavyam . \\
\noindent \textbf{(P\_3,1.15)} na vaktavyam . \\
\noindent \textbf{(P\_3,1.15)} kasmāt na bhavati . \\
\noindent \textbf{(P\_3,1.15)} kīṭaḥ romatham vartayati iti . \\
\noindent \textbf{(P\_3,1.15)} anabhidhānāt . \\
\noindent \textbf{(P\_3,1.15)} tapasaḥ parasmaipadam ca . \\
\noindent \textbf{(P\_3,1.15)} tapasaḥ parasmaipadam ca iti vaktavyam . \\
\noindent \textbf{(P\_3,1.15)} tapaḥ carati tapasyati . \\
\noindent \textbf{(P\_3,1.15)} katham tapasyate lokajigīṣuḥ agneḥ . \\
\noindent \textbf{(P\_3,1.15)} chāndasatvāt bhaviṣyati . \\
\noindent \textbf{(P\_3,1.16)} phenāt ca iti vaktavyam . \\
\noindent \textbf{(P\_3,1.16)} phenāyate . \\
\noindent \textbf{(P\_3,1.17)} aṭāṭtāśīkākoṭāpoṭāsoṭāpruṣṭāpluṣṭāgrahaṇam kartavyam . \\
\noindent \textbf{(P\_3,1.17)} aṭā . \\
\noindent \textbf{(P\_3,1.17)} aṭāyate . \\
\noindent \textbf{(P\_3,1.17)} aṭṭā . \\
\noindent \textbf{(P\_3,1.17)} aṭṭāyate . \\
\noindent \textbf{(P\_3,1.17)} śīkā . \\
\noindent \textbf{(P\_3,1.17)} śīkāyate . \\
\noindent \textbf{(P\_3,1.17)} koṭā . \\
\noindent \textbf{(P\_3,1.17)} koṭāyate . \\
\noindent \textbf{(P\_3,1.17)} poṭā . \\
\noindent \textbf{(P\_3,1.17)} poṭāyate . \\
\noindent \textbf{(P\_3,1.17)} soṭā . \\
\noindent \textbf{(P\_3,1.17)} soṭāyate . \\
\noindent \textbf{(P\_3,1.17)} pruṣṭā . \\
\noindent \textbf{(P\_3,1.17)} pruṣṭayate . \\
\noindent \textbf{(P\_3,1.17)} pluṣṭā . \\
\noindent \textbf{(P\_3,1.17)} pluṣṭāyate . \\
\noindent \textbf{(P\_3,1.17)} sudinadurdinābhyām ca . \\
\noindent \textbf{(P\_3,1.17)} sudinadurdinābhyām ca iti vaktavyam . \\
\noindent \textbf{(P\_3,1.17)} sudināyate . \\
\noindent \textbf{(P\_3,1.17)} durdināyate . \\
\noindent \textbf{(P\_3,1.17)} nīhārāt ca . \\
\noindent \textbf{(P\_3,1.17)} nīhārāt ca iti vaktavyam . \\
\noindent \textbf{(P\_3,1.17)} nīhārāyate . \\
\noindent \textbf{(P\_3,1.18)} kartṛvedanāyām iti kimartham . \\
\noindent \textbf{(P\_3,1.18)} iha mā bhūt . \\
\noindent \textbf{(P\_3,1.18)} sukham vedayate prasādhakaḥ devadattasya . \\
\noindent \textbf{(P\_3,1.18)} kartṛvedanāyām iti ucyamāne api atra prāpnoti . \\
\noindent \textbf{(P\_3,1.18)} kim kāraṇam . \\
\noindent \textbf{(P\_3,1.18)} kartuḥ iti iyam kartari ṣaṣṭhī . \\
\noindent \textbf{(P\_3,1.18)} vedanāyām iti ca anaḥ bhāve . \\
\noindent \textbf{(P\_3,1.18)} saḥ yadi eva ātmanaḥ vedayate atha api parasya kartṛvedanā eva asau bhavati . \\
\noindent \textbf{(P\_3,1.18)} na kartṛgrahaṇena vedanā abhisambadhyate . \\
\noindent \textbf{(P\_3,1.18)} kim tarhi . \\
\noindent \textbf{(P\_3,1.18)} sukhādīni abhisambadhyante . \\
\noindent \textbf{(P\_3,1.18)} kartuḥ yāni sukhādīni . \\
\noindent \textbf{(P\_3,1.19.1)} namasaḥ kyaci dvitīyānupapattiḥ . \\
\noindent \textbf{(P\_3,1.19.1)} namasaḥ kyaci dvitīyā na upapadyate . \\
\noindent \textbf{(P\_3,1.19.1)} namasyati devān . \\
\noindent \textbf{(P\_3,1.19.1)} kim kāraṇam . \\
\noindent \textbf{(P\_3,1.19.1)} namaḥśabdena yoge caturthī vidhīyate . \\
\noindent \textbf{(P\_3,1.19.1)} sā prāpnoti . \\
\noindent \textbf{(P\_3,1.19.1)} prakṛtyantaratvāt siddham . \\
\noindent \textbf{(P\_3,1.19.1)} namaḥśabdena yoge caturthī vidhīyate namasyatiśabdaḥ ca ayam . \\
\noindent \textbf{(P\_3,1.19.1)} nanu ca namasyatiśabde namaḥśabdaḥ asti . \\
\noindent \textbf{(P\_3,1.19.1)} tena yoge prāpnoti . \\
\noindent \textbf{(P\_3,1.19.1)} na eṣaḥ doṣaḥ . \\
\noindent \textbf{(P\_3,1.19.1)} arthavataḥ namaḥsābdasya grahaṇam . \\
\noindent \textbf{(P\_3,1.19.1)} na ca namasyatiśabde namaḥśabdaḥ arthavān . \\
\noindent \textbf{(P\_3,1.19.1)} atha vā upapadavibhakteḥ kārakavibhaktiḥ balīyasī iti dvitīyā bhaviṣyati . \\
\noindent \textbf{(P\_3,1.19.2)} kyajādiṣu pratyayārthanirdeśaḥ . \\
\noindent \textbf{(P\_3,1.19.2)} kyajādiṣu pratyayārthanirdeśaḥ kartavyaḥ . \\
\noindent \textbf{(P\_3,1.19.2)} namasaḥ pūjāyām . \\
\noindent \textbf{(P\_3,1.19.2)} varivasaḥ paricaryāyām . \\
\noindent \textbf{(P\_3,1.19.2)} citraṅaḥ āścarye . \\
\noindent \textbf{(P\_3,1.19.2)} bhāṇḍāt samācayane . \\
\noindent \textbf{(P\_3,1.19.2)} cīvarāt arjane paridhāne vā . \\
\noindent \textbf{(P\_3,1.19.2)} pucchāt udasane vysasane ca iti . \\
\noindent \textbf{(P\_3,1.19.2)} kim prayojanam . \\
\noindent \textbf{(P\_3,1.19.2)} kriyāvacanatā yathā syāt . \\
\noindent \textbf{(P\_3,1.19.2)} na etat asti prayojanam . \\
\noindent \textbf{(P\_3,1.19.2)} ācāryapravṛttiḥ jñāpayati kriyāvacanāḥ kyajādayaḥ iti yat ayam sanādyantāḥ dhātavaḥ iti dhātusañjñām śāsti . \\
\noindent \textbf{(P\_3,1.19.2)} dhātusañjñāvacane etat prayojanam : dhātoḥ iti tavyadādīnām utpattiḥ yathā syāt . \\
\noindent \textbf{(P\_3,1.19.2)} yadi ca atra kriyāvacanatā na syāt dhātusañjñāvacanam anarthakam syāt . \\
\noindent \textbf{(P\_3,1.19.2)} satyām api dhātusañjñāyām tavyadādayaḥ na syuḥ . \\
\noindent \textbf{(P\_3,1.19.2)} kim kāraṇam . \\
\noindent \textbf{(P\_3,1.19.2)} sādhane tāvyādayaḥ vidhīyante sādhanam ca kriyāyāḥ . \\
\noindent \textbf{(P\_3,1.19.2)} kriyābhāvāt sādhanābhāvaḥ . \\
\noindent \textbf{(P\_3,1.19.2)} sādhanābhāvāt satyām api dhātusañjñāyām tavyadādayaḥ na syuḥ . \\
\noindent \textbf{(P\_3,1.19.2)} paśyati tu ācāryaḥ kriyāvacanāḥ kyajādayaḥ iti tataḥ sanādyantāḥ dhātavaḥ iti dhātusañjñām śāsti . \\
\noindent \textbf{(P\_3,1.19.2)} nanu ca idam prayojanam syāt . \\
\noindent \textbf{(P\_3,1.19.2)} parasādhane utpattim vakṣyāmi iti . \\
\noindent \textbf{(P\_3,1.19.2)} na parasādhane utpattyā bhavitavyam . \\
\noindent \textbf{(P\_3,1.19.2)} kim kāraṇam . \\
\noindent \textbf{(P\_3,1.19.2)} sādhanam iti sambandhiśabdaḥ ayam . \\
\noindent \textbf{(P\_3,1.19.2)} sambandhiśabdāḥ ca punaḥ evamātmakāḥ yat uta sambandhinam ākṣipanti . \\
\noindent \textbf{(P\_3,1.19.2)} tat yathā . \\
\noindent \textbf{(P\_3,1.19.2)} mātari vartitatvyam , pitari śuśrūṣitavyam iti . \\
\noindent \textbf{(P\_3,1.19.2)} na ca ucyate svasyām mātari svasmin vā pitari iti , sambandhāt ca etat gamyate yā yasya mātā yaḥ ca yasya pitā iti . \\
\noindent \textbf{(P\_3,1.19.2)} evam iha api sambandhāt etat gantavyam yasya dhātoḥ yat sādhanam iti . \\
\noindent \textbf{(P\_3,1.19.2)} atha vā dhātavaḥ eva kyajādayaḥ . \\
\noindent \textbf{(P\_3,1.19.2)} na ca eva hi arthāḥ ādiśyante . \\
\noindent \textbf{(P\_3,1.19.2)} kriyāvacanatā ca gamyate . \\
\noindent \textbf{(P\_3,1.19.2)} kaḥ khalu api pacādīnām kriyāvacanatve yatnam karoti . \\
\noindent \textbf{(P\_3,1.19.2)} yena eva khalu api hetunā pacādayaḥ kriyāvacanāḥ tena eva kyajādayaḥ api . \\
\noindent \textbf{(P\_3,1.19.2)} evamartham ācāryaḥ citrayati . \\
\noindent \textbf{(P\_3,1.19.2)} kva cit arthān ādiśati kva cit na . \\
\noindent \textbf{(P\_3,1.19.2)} evam api arthādeśanam kartavyam . \\
\noindent \textbf{(P\_3,1.19.2)} katham ime abudhāḥ budhyeran iti . \\
\noindent \textbf{(P\_3,1.19.2)} atha vā śakyam ādeśanam akartum . \\
\noindent \textbf{(P\_3,1.19.2)} katham . \\
\noindent \textbf{(P\_3,1.19.2)} karaṇe iti vartate . \\
\noindent \textbf{(P\_3,1.19.2)} karaṇam ca karoteḥ karotiḥ ca kriyāsāmānye vartate . \\
\noindent \textbf{(P\_3,1.21)} imau halikalī staḥ ikārāntau . \\
\noindent \textbf{(P\_3,1.21)} asti halaśabdaḥ kalaśabdaḥ ca akārāntaḥ . \\
\noindent \textbf{(P\_3,1.21)} kayoḥ idam grahaṇam . \\
\noindent \textbf{(P\_3,1.21)} yau ikārāntau tayoḥ atvam nipātyate . \\
\noindent \textbf{(P\_3,1.21)} kim prayojanam . \\
\noindent \textbf{(P\_3,1.21)} halikalyoḥ atvanipātanam sanvadbhāvapratiṣedhārtham . \\
\noindent \textbf{(P\_3,1.21)} halikalyoḥ atvanipātanam kriyate sanvadbhāvaḥ mā bhūt iti . \\
\noindent \textbf{(P\_3,1.21)} ajahalat acakalat . \\
\noindent \textbf{(P\_3,1.21)} na etat asti prayojanam . \\
\noindent \textbf{(P\_3,1.21)} ikāralope kṛte aglopinām na iti pratiṣedhaḥ bhaviṣyati . \\
\noindent \textbf{(P\_3,1.21)} vṛddhau kṛtāyām lopaḥ . \\
\noindent \textbf{(P\_3,1.21)} tat na aglopi aṅgam bhavati . \\
\noindent \textbf{(P\_3,1.21)} idam iha sampradhāryam . \\
\noindent \textbf{(P\_3,1.21)} vṛddhiḥ kriyatām aglopaḥ iti . \\
\noindent \textbf{(P\_3,1.21)} kim atra kartavyam . \\
\noindent \textbf{(P\_3,1.21)} paratvāt vṛddhiḥ . \\
\noindent \textbf{(P\_3,1.21)} nityaḥ lopaḥ . \\
\noindent \textbf{(P\_3,1.21)} kṛtāyām api vṛddhau prāpnoti akṛtāyām api prāpnoti . \\
\noindent \textbf{(P\_3,1.21)} anityaḥ lopaḥ . \\
\noindent \textbf{(P\_3,1.21)} anyasya kṛtāyām vṛddhau prāpnoti anyasya akṛtāyām . \\
\noindent \textbf{(P\_3,1.21)} śabdāntarasya ca prāpnuvan vidhiḥ anityaḥ bhavati . \\
\noindent \textbf{(P\_3,1.21)} vṛddhiḥ api anityā . \\
\noindent \textbf{(P\_3,1.21)} anyasya kṛte lope prāpnoti anyasya akṛte . \\
\noindent \textbf{(P\_3,1.21)} śabdāntarasya ca prāpnuvan vidhiḥ anityaḥ bhavati . \\
\noindent \textbf{(P\_3,1.21)} ubhayoḥ anityayoḥ paratvāt vṛddhiḥ . \\
\noindent \textbf{(P\_3,1.21)} vṛddhau kṛtāyām lopaḥ . \\
\noindent \textbf{(P\_3,1.21)} tat na aglopi aṅgam bhavati . \\
\noindent \textbf{(P\_3,1.21)} atve punaḥ sati vṛddhiḥ kriyatām lopaḥ iti yadi api paratvāt vṛddhiḥ vṛddhau kṛtāyām api ak eva lupyate . \\
\noindent \textbf{(P\_3,1.21)} tasmāt suṣṭhu ucyate halikalyoḥ atvanipātanam sanvadbhāvapratiṣedhārtham . \\
\noindent \textbf{(P\_3,1.22.1)} samabhihāraḥ iti kaḥ ayam śabdaḥ . \\
\noindent \textbf{(P\_3,1.22.1)} samabhipūrvāt harateḥ bhāvasādhanaḥ ghañ . \\
\noindent \textbf{(P\_3,1.22.1)} samabhiharaṇam samabhihāraḥ . \\
\noindent \textbf{(P\_3,1.22.1)} tat yatha puṣpābhihāraḥ phalābhihāraḥ iti . \\
\noindent \textbf{(P\_3,1.22.1)} viṣamaḥ upanyāsaḥ . \\
\noindent \textbf{(P\_3,1.22.1)} bahvyaḥ hi tāḥ sumanasaḥ . \\
\noindent \textbf{(P\_3,1.22.1)} tatra yuktaḥ samabhihāraḥ . \\
\noindent \textbf{(P\_3,1.22.1)} iha punaḥ ekā kriyā . \\
\noindent \textbf{(P\_3,1.22.1)} yadi api ekā sāmānyakriyā avayavakriyāḥ tu bahvyaḥ adhiśrayaṇodakāsecanataṇḍulāvapanaidhkopakarṣaṇkriyāḥ . \\
\noindent \textbf{(P\_3,1.22.1)} tāḥ kaḥ cit kārtsnyena karoti kaḥ cit akārtsnyena . \\
\noindent \textbf{(P\_3,1.22.1)} yaḥ kārtsnyena karoti saḥ ucyate pāpacyate iti . \\
\noindent \textbf{(P\_3,1.22.1)} punaḥ punaḥ vā pacati pāpacyate iti . \\
\noindent \textbf{(P\_3,1.22.2)} atha dhātugrahaṇam kimartham . \\
\noindent \textbf{(P\_3,1.22.2)} iha mā bhūt prāṭati bhṛśam iti . \\
\noindent \textbf{(P\_3,1.22.2)} ataḥ uttaram paṭhati . \\
\noindent \textbf{(P\_3,1.22.2)} yaṅvidhau dhātugrahaṇe uktam . \\
\noindent \textbf{(P\_3,1.22.2)} kim uktam . \\
\noindent \textbf{(P\_3,1.22.2)} tatra tāvat uktam karmagrahaṇāt sanvidhau dhātugrahaṇānarthakyam . \\
\noindent \textbf{(P\_3,1.22.2)} sopasargam karma iti cet karmaviśeṣakatvāt upasargasya anupasargam karma . \\
\noindent \textbf{(P\_3,1.22.2)} sopasargasya hi karmatve dhātvadhikāre api sanaḥ avidhānam akarmatvāt iti . \\
\noindent \textbf{(P\_3,1.22.2)} evam iha api kriyāsamabhihāragrahaṇāt yaṅvidhau dhātugrahaṇānarthakyam . \\
\noindent \textbf{(P\_3,1.22.2)} sopasargaḥ kriyāsamabhihāraḥ iti cet kriyāsamabhihāraviśeṣakatvāt upasargasya anupasargaḥ kriyāsamabhihāraḥ . \\
\noindent \textbf{(P\_3,1.22.2)} sopasargasya hi kriyāsamahibhāratve dhātvadhikāre api yaṅaḥ avidhānam akriyāsamabhihāratvāt iti . \\
\noindent \textbf{(P\_3,1.22.2)} atha ekājjhalādigrahaṇam kimartham . \\
\noindent \textbf{(P\_3,1.22.2)} iha mā bhūt : jāgarti bhṛśam . \\
\noindent \textbf{(P\_3,1.22.2)} īkṣate bhṛśam . \\
\noindent \textbf{(P\_3,1.22.2)} ekājjhalādigrahaṇe ca . \\
\noindent \textbf{(P\_3,1.22.2)} ekājjhalādigrahaṇe ca uktam . \\
\noindent \textbf{(P\_3,1.22.2)} kim uktam . \\
\noindent \textbf{(P\_3,1.22.2)} tatra tāvat uktam karmasamānakartṛkagrahaṇānarthakyam ca icchābhidhāne pratyayavidhānāt . \\
\noindent \textbf{(P\_3,1.22.2)} akarmaṇaḥ hi asamānakartṛkāt vā anabhidhānam iti . \\
\noindent \textbf{(P\_3,1.22.2)} iha api ekājjhalādigrahaṇānarthakyam kriyāsamabhihāre yaṅvacanāt anekācaḥ ahalādeḥ hi anabhidhānam iti . \\
\noindent \textbf{(P\_3,1.22.2)} tat ca avaśyam anabhidānam āśrayitavyam . \\
\noindent \textbf{(P\_3,1.22.2)} kriyamāṇe api hi ekājjhalādigrahaṇe yatra ekācaḥ halādeḥ ca utpadyamānena yaṅā arthasya abhidhānam na bhavati na bhavati tatra utpattiḥ . \\
\noindent \textbf{(P\_3,1.22.2)} tat yathā . \\
\noindent \textbf{(P\_3,1.22.2)} bhṛśam śobhate . \\
\noindent \textbf{(P\_3,1.22.2)} bhṛśam rocate . \\
\noindent \textbf{(P\_3,1.22.2)} yatra ca anekācaḥ ahalādeḥ vo utpadyamānena yaṅā arthasya abhidhānam bhavati bhavati tatra utpattiḥ . \\
\noindent \textbf{(P\_3,1.22.2)} tat yathā . \\
\noindent \textbf{(P\_3,1.22.2)} aṭāṭyate arāryate aśāśyate sosūcyate sosūtryate momūtryate . \\
\noindent \textbf{(P\_3,1.22.3)} ūrṇoteḥ ca upasaṅkhyānam . \\
\noindent \textbf{(P\_3,1.22.3)} ūrṇoteḥ ca upasaṅkhyānam kartavyam . \\
\noindent \textbf{(P\_3,1.22.3)} prorṇonūyate . \\
\noindent \textbf{(P\_3,1.22.3)} atyalpam idam ucyate : ūrṇoteḥ iti . \\
\noindent \textbf{(P\_3,1.22.3)} sūcisūtrimūtryaṭyartyaśyūrṇugrahaṇam yaṅvidhau anekājahalādyartham . \\
\noindent \textbf{(P\_3,1.22.3)} sūcisūtrimūtryaṭyartyaśyūrṇotīnām grahaṇam kartavyam . \\
\noindent \textbf{(P\_3,1.22.3)} kim prayojanam . \\
\noindent \textbf{(P\_3,1.22.3)} yaṅvidhau anekājahalādyartham . \\
\noindent \textbf{(P\_3,1.22.3)} sosūcyate sosūtryate momūtryate aṭāṭyate arāryate aśāśyate prorṇonūyate . \\
\noindent \textbf{(P\_3,1.22.3)} vācyaḥ ūrṇorṇuvadbhāvaḥ yaṅprasiddhiḥ prayojanam . \\
\noindent \textbf{(P\_3,1.22.3)} āmaḥ ca pratiṣedhārtham ekācaḥ ca iḍupagrahāt . \\
\noindent \textbf{(P\_3,1.22.4)} kriyāsamabhihāre yaṅaḥ vipratiṣedhena loḍvidhānam . \\
\noindent \textbf{(P\_3,1.22.4)} kriyāsamabhihāre loṭ bhavati yaṅaḥ vipratiṣedhena . \\
\noindent \textbf{(P\_3,1.22.4)} kriyāsamabhihāre yaṅ bhavati iti asya avakāśaḥ dhātuḥ yaḥ ekāc halādiḥ kriyāsamabhihāre vartate adhātusambandhaḥ : lolūyate . \\
\noindent \textbf{(P\_3,1.22.4)} loṭaḥ avakāśaḥ dhātuḥ yaḥ anekāc ahalādiḥ kriyāsamabhihāre vartate dhātusambandhaḥ : saḥ bhavān jāgṛhi jāgṛhi iti eva ayam jāgarti . \\
\noindent \textbf{(P\_3,1.22.4)} saḥ bhavān īhasva īhasva iti eva ayam īhate . \\
\noindent \textbf{(P\_3,1.22.4)} dhātuḥ yaḥ ekāc halādiḥ kriyāsamabhihāre vartate dhātusambandhaḥ ca tasmāt ubhayam prāpnoti : saḥ bhavān lunīhi lunīhi iti eva ayam lunāti . \\
\noindent \textbf{(P\_3,1.22.4)} loṭ bhavati vipratiṣedhena . \\
\noindent \textbf{(P\_3,1.22.4)} na tarhi idānīm idam bhavati : saḥ bhavān lolūyasva lolūyasva iti eva ayam lolūyate . \\
\noindent \textbf{(P\_3,1.22.4)} bhavati ca . \\
\noindent \textbf{(P\_3,1.22.4)} na vā nānārthatvāt . \\
\noindent \textbf{(P\_3,1.22.4)} kartṛkarmaṇoḥ hi lavidhānam kriyāviśeṣe svārthe yaṅ . \\
\noindent \textbf{(P\_3,1.22.4)} na vā arthaḥ vipratiṣedhena . \\
\noindent \textbf{(P\_3,1.22.4)} kim kāraṇam . \\
\noindent \textbf{(P\_3,1.22.4)} nānārthatvāt . \\
\noindent \textbf{(P\_3,1.22.4)} kā nānārthatā . \\
\noindent \textbf{(P\_3,1.22.4)} kartṛkarmaṇoḥ hi lavidhānam . \\
\noindent \textbf{(P\_3,1.22.4)} kartṛkarmaṇoḥ hi loṭ vidhīyate . \\
\noindent \textbf{(P\_3,1.22.4)} kriyāviśeṣe svārthe yaṅ . \\
\noindent \textbf{(P\_3,1.22.4)} tatra antaraṅgatvāt yaṅā bhavitavyam . \\
\noindent \textbf{(P\_3,1.22.4)} na tarhi idam idānīm bhavati . \\
\noindent \textbf{(P\_3,1.22.4)} saḥ bhavān lunīhi lunīhi iti eva ayam lunāti . \\
\noindent \textbf{(P\_3,1.22.4)} bhavati ca . \\
\noindent \textbf{(P\_3,1.22.4)} vibhāṣā yaṅ . \\
\noindent \textbf{(P\_3,1.22.4)} yadā na yaṅ tadā loṭ . \\
\noindent \textbf{(P\_3,1.24)} uttarayoḥ vigraheṇa viśeṣāsampratyayāt nityagrahaṇānarthakyam . \\
\noindent \textbf{(P\_3,1.24)} uttarayoḥ yogayoḥ vigraheṇa viśeṣāsampratyayāt nityagrahaṇānarthakyam . \\
\noindent \textbf{(P\_3,1.24)} na hi kuṭilam krāmati iti caṅkramyate iti gamyate . \\
\noindent \textbf{(P\_3,1.24)} athe etebhyaḥ kriyāsamabhihāre yaṅā bhavitavyam . \\
\noindent \textbf{(P\_3,1.24)} kriyāsamabhihāre ca na etebhyaḥ . \\
\noindent \textbf{(P\_3,1.24)} kriyāsamabhihāre ca na etebhyaḥ yaṅā bhavitavyam . \\
\noindent \textbf{(P\_3,1.24)} bhṛśam japati brāhmaṇaḥ . \\
\noindent \textbf{(P\_3,1.24)} bhṛśam samidaḥ dahati iti eva. \\
\noindent \textbf{(P\_3,1.25)} satyāpa iti kim nipātyate . \\
\noindent \textbf{(P\_3,1.25)} satyasya kṛñi āpuk ca . satyasya kṛñi āpuk ca nipātyate ṇic ca . \\
\noindent \textbf{(P\_3,1.25)} satyam karoti satyāpayati . \\
\noindent \textbf{(P\_3,1.25)} atyalpam idam ucyate . \\
\noindent \textbf{(P\_3,1.25)} ṇividau arthavedasatyānām apuk ca . \\
\noindent \textbf{(P\_3,1.25)} ṇividau arthavedasatyānām apuk ca iti vaktavyam . \\
\noindent \textbf{(P\_3,1.25)} arthāpayati vedāpayati satyāpayati . \\
\noindent \textbf{(P\_3,1.25)} yadi āpuk kriyate ṭilopaḥ prāpnoti . \\
\noindent \textbf{(P\_3,1.25)} evam tarhi puk kariṣyate . \\
\noindent \textbf{(P\_3,1.25)} evam api ṭilopaḥ prāpnoti . \\
\noindent \textbf{(P\_3,1.25)} evam tarhi āk kariṣyate . \\
\noindent \textbf{(P\_3,1.25)} evam api ṭilopaḥ prāpnoti . \\
\noindent \textbf{(P\_3,1.25)} evam tarhi ak kariṣyate . \\
\noindent \textbf{(P\_3,1.25)} evam api anākārāntatvāt puk na prāpnoti . \\
\noindent \textbf{(P\_3,1.25)} evam tarhi apuṭ kariṣyate . \\
\noindent \textbf{(P\_3,1.25)} atha vā punaḥ astu āpuk eva . \\
\noindent \textbf{(P\_3,1.25)} nanu ca uktam . \\
\noindent \textbf{(P\_3,1.25)} ṭilopaḥ prāpnoti iti . \\
\noindent \textbf{(P\_3,1.25)} āpugvacanasāmarthyāt na bhaviṣyati . \\
\noindent \textbf{(P\_3,1.25)} atha vā punaḥ astu puk eva . \\
\noindent \textbf{(P\_3,1.25)} nanu ca uktam evam api ṭilopaḥ prāpnoti iti . \\
\noindent \textbf{(P\_3,1.25)} pugvacanasāmarthyāt na bhaviṣyati . \\
\noindent \textbf{(P\_3,1.25)} atha vā punaḥ astu āk eva . \\
\noindent \textbf{(P\_3,1.25)} nanu ca uktam evam api ṭilopaḥ prāpnoti iti . \\
\noindent \textbf{(P\_3,1.25)} āgvacanasāmarthyāt na bhaviṣyati . \\
\noindent \textbf{(P\_3,1.26.1)} katham idam vijñāyate . \\
\noindent \textbf{(P\_3,1.26.1)} hetumati abhidheye ṇic bhavati iti . \\
\noindent \textbf{(P\_3,1.26.1)} āhosvit hetumati yaḥ dhātuḥ vartate iti . \\
\noindent \textbf{(P\_3,1.26.1)} yuktam punaḥ idam vicārayitum . \\
\noindent \textbf{(P\_3,1.26.1)} nanu anena asandigdhena pratyayārthaviśeṣaṇena bhavitavyam yāvatā hetumati iti ucyate . \\
\noindent \textbf{(P\_3,1.26.1)} yadi hi prakṛtarthaviśeṣaṇam syāt hetumataḥ iti evam brūyāt . \\
\noindent \textbf{(P\_3,1.26.1)} na etat asti . \\
\noindent \textbf{(P\_3,1.26.1)} bhavanti iha hi viṣayasaptamyaḥ api . \\
\noindent \textbf{(P\_3,1.26.1)} tat yathā . \\
\noindent \textbf{(P\_3,1.26.1)} pramāṇe yat prātipadikam vartate striyām yat prātipadikam vartate iti . \\
\noindent \textbf{(P\_3,1.26.1)} evam iha api hetumati abhidheye ṇic bhavati hetumati yaḥ dhātuḥ vartate iti jāyate vicāraṇā . \\
\noindent \textbf{(P\_3,1.26.1)} ata uttaram paṭhati . \\
\noindent \textbf{(P\_3,1.26.1)} hetumati iti kārakopādānam pratyayārthaparigrahārtham yathā tanūkaraṇe takṣaḥ . \\
\noindent \textbf{(P\_3,1.26.1)} hetumati iti kārakam upādīyate . \\
\noindent \textbf{(P\_3,1.26.1)} kim prayojanam . \\
\noindent \textbf{(P\_3,1.26.1)} pratyayārthaparigrahārtham . \\
\noindent \textbf{(P\_3,1.26.1)} evam sati pratyayārthaḥ suparigṛhītaḥ bhavati . \\
\noindent \textbf{(P\_3,1.26.1)} yathā tanūkaraṇe takṣaḥ iti tanūkaraṇam upādīyate . \\
\noindent \textbf{(P\_3,1.26.1)} yadi tarhi tadvat prakṛtyarthaviśeṣaṇam bhavati . \\
\noindent \textbf{(P\_3,1.26.1)} prakṛtyarthaviśeṣaṇam hi tat tatra vijñāyate . \\
\noindent \textbf{(P\_3,1.26.1)} tanūkaraṇakriyāyām takṣaḥ iti . \\
\noindent \textbf{(P\_3,1.26.1)} astu prakṛtyarthaviśeṣaṇam . \\
\noindent \textbf{(P\_3,1.26.1)} kaḥ doṣaḥ . \\
\noindent \textbf{(P\_3,1.26.1)} iha hi uktaḥ karoti preṣitaḥ karoti iti ṇic prāpnoti . \\
\noindent \textbf{(P\_3,1.26.1)} pratyayārthaviśeṣaṇe punaḥ sati na eṣaḥ doṣaḥ . \\
\noindent \textbf{(P\_3,1.26.1)} svaśabdena uktatvāt na bhaviṣyati . \\
\noindent \textbf{(P\_3,1.26.1)} prakṛtyarthaviśeṣaṇe api sati na eṣaḥ doṣaḥ . \\
\noindent \textbf{(P\_3,1.26.1)} yatra na antareṇa śabdam arthasya gatiḥ bhavati tatra śabdaḥ prayujyate . \\
\noindent \textbf{(P\_3,1.26.1)} yatra hi antareṇa api śabdam arthasya gatiḥ bhavati na tatra śabdaḥ prayujyate . \\
\noindent \textbf{(P\_3,1.26.1)} iha tarhi pācayati odanam devadattaḥ yajñadattena iti ubhayoḥ kartroḥ lena abhidhānam prāpnoti . \\
\noindent \textbf{(P\_3,1.26.1)} pratyayārthaviśeṣaṇe punaḥ sati na doṣaḥ . \\
\noindent \textbf{(P\_3,1.26.1)} pradhānakartari lādayaḥ bhavanti iti pradhānakartā lena abhidhīyate . \\
\noindent \textbf{(P\_3,1.26.1)} yaḥ ca apradhānam siddha tatra kartari iti eva tṛtīyā . \\
\noindent \textbf{(P\_3,1.26.1)} iha ca gamitaḥ grāmam devadattaḥ yajñadattanea iti avyatiriktaḥ gatyarthaḥ iti kṛtvā gatyarthānām kartari iti kartari ktaḥ prāpnoti . \\
\noindent \textbf{(P\_3,1.26.1)} iha ca vyatibhedayante vyaticchedayante iti avyatiriktaḥ hiṃsārthaḥ iti kṛtvā na gatihiṃsārthebhyaḥ iti pratiṣedhaḥ prāpnoti . \\
\noindent \textbf{(P\_3,1.26.1)} astu tarhi pratyayārthaviśeṣaṇam . \\
\noindent \textbf{(P\_3,1.26.1)} yadi pratyayārthaviśeṣaṇam pācayati odanam devadattaḥ yajñadattena iti prayojye kartari karmasañjñā prāpnoti . \\
\noindent \textbf{(P\_3,1.26.1)} bhavati hi tasya tasmin īpsā . \\
\noindent \textbf{(P\_3,1.26.1)} iha ca grāmam gamayati grāmāya gamayati iti vyatiriktaḥ gatyarthaḥ iti kṛtvā gatyarthakarmaṇi dvitīyācaturthyau na prāpnutaḥ . \\
\noindent \textbf{(P\_3,1.26.1)} iha ca edhodakasya upaskārayati iti vyatiriktaḥ karotyarthaḥ iti kṛtvā kṛñaḥ pratiyatne iti ṣaṣṭhī na prāpnoti . \\
\noindent \textbf{(P\_3,1.26.1)} iha ca bhedikā devadattasya yajñadattasya kāṣṭhānām iti prayojye kartari ṣaṣṭhī na prāpnoti . \\
\noindent \textbf{(P\_3,1.26.1)} iha ca abhiṣāvayati pariṣāvayati iti vyatiriktaḥ sunotyarthaḥ iti kṛtvā upasargāt sunotyādīnām iti ṣatvam na prāpnoti . \\
\noindent \textbf{(P\_3,1.26.1)} na eṣaḥ doṣaḥ . \\
\noindent \textbf{(P\_3,1.26.1)} yat tāvat ucyate pācayati odanam devadattaḥ yajñadattena iti prayojye kartari karmasañjñā prāpnoti iti . \\
\noindent \textbf{(P\_3,1.26.1)} gatibuddhipratyavasānārthaśabdakarmākarmakāṇām aṇi iti etat niyamārtham bhaviṣyati . \\
\noindent \textbf{(P\_3,1.26.1)} eteṣām eva aṇyantānām yaḥ kartā saḥ ṇau karmasañjñaḥ bhavati na anyeṣām iti . \\
\noindent \textbf{(P\_3,1.26.1)} yat api ucyate iha ca grāmam gamayati grāmāya gamayati iti vyatiriktaḥ gatyarthaḥ iti kṛtvā gatyarthakarmaṇi dvitīyācaturthyau na prāpnutaḥ iti . \\
\noindent \textbf{(P\_3,1.26.1)} na asau evam preṣyate gaccha grāmam iti . \\
\noindent \textbf{(P\_3,1.26.1)} katham tarhi . \\
\noindent \textbf{(P\_3,1.26.1)} sādhanaviśiṣṭām asau kriyām preṣyate . \\
\noindent \textbf{(P\_3,1.26.1)} grāmam gaccha . \\
\noindent \textbf{(P\_3,1.26.1)} grāmāya gaccha iti . \\
\noindent \textbf{(P\_3,1.26.1)} yat api ucyate iha ca edhodakasya upaskārayati iti vyatiriktaḥ karotyarthaḥ iti kṛtvā kṛñaḥ pratiyatne iti ṣaṣṭhī na prāpnoti iti . \\
\noindent \textbf{(P\_3,1.26.1)} na asau evam preṣyate upaskuruṣva edhodakasya iti . \\
\noindent \textbf{(P\_3,1.26.1)} katham tarhi . \\
\noindent \textbf{(P\_3,1.26.1)} sādhanaviśiṣṭām asau kriyām preṣyate . \\
\noindent \textbf{(P\_3,1.26.1)} edhodakasya upaskuruṣva iti . \\
\noindent \textbf{(P\_3,1.26.1)} yat api ucyate iha ca bhedikā devadattasya yajñadattasya kāṣṭhānām iti prayojye kartari ṣaṣṭhī na prāpnoti iti . \\
\noindent \textbf{(P\_3,1.26.1)} uktam tatra kṛdgrahaṇasya prayojanam kartṛbhūtapūrvamātre api ṣaṣṭhī yathā syāt iti . \\
\noindent \textbf{(P\_3,1.26.1)} yat api ucyate iha ca abhiṣāvayati pariṣāvayati iti vyatiriktaḥ sunotyarthaḥ iti kṛtvā upasargāt sunotyādīnām iti ṣatvam na prāpnoti iti . \\
\noindent \textbf{(P\_3,1.26.1)} na asau evam preṣyate sunu abhi iti . \\
\noindent \textbf{(P\_3,1.26.1)} katham tarhi upasargaviśiṣṭām asau kriyām preṣyate . \\
\noindent \textbf{(P\_3,1.26.1)} abhiṣunu iti . \\
\noindent \textbf{(P\_3,1.26.1)} yuktam punaḥ idam vicārayitum . \\
\noindent \textbf{(P\_3,1.26.1)} nanu anena asandigdhena pratyayārthaviśeṣaṇena bhavitavyam yāvatā vyaktam arthāntaram gamyate pacati pācayati iti ca . \\
\noindent \textbf{(P\_3,1.26.1)} bāḍham yuktam . \\
\noindent \textbf{(P\_3,1.26.1)} iha paceḥ kaḥ pradhānārthaḥ . \\
\noindent \textbf{(P\_3,1.26.1)} yā asau taṇḍulānām viklittiḥ . \\
\noindent \textbf{(P\_3,1.26.1)} atha idānīm tadabhisandhipūrvakam preṣaṇam adhyeṣaṇam vā . \\
\noindent \textbf{(P\_3,1.26.1)} yuktam yat sarvam pacyarthaḥ syāt . \\
\noindent \textbf{(P\_3,1.26.2)} hetunirdeśaḥ ca nimittamātram bikṣādiṣu darśanāt . \\
\noindent \textbf{(P\_3,1.26.2)} hetunirdeśaḥ ca nimittamātram draṣṭavyam . \\
\noindent \textbf{(P\_3,1.26.2)} yāvat brūyāt nimittam kāraṇam iti tāvt hetuḥ iti . \\
\noindent \textbf{(P\_3,1.26.2)} kim prayojanam . \\
\noindent \textbf{(P\_3,1.26.2)} bikṣādiṣu darśanāt . \\
\noindent \textbf{(P\_3,1.26.2)} bhikṣādiṣu hi ṇic dṛśyate . \\
\noindent \textbf{(P\_3,1.26.2)} bhikṣāḥ vāsayanti .kāriṣaḥ agniḥ adhyāpayati iti . \\
\noindent \textbf{(P\_3,1.26.2)} kim punaḥ kāraṇam pāribhāṣike hetau na sidhyati . \\
\noindent \textbf{(P\_3,1.26.2)} evam manyate . \\
\noindent \textbf{(P\_3,1.26.2)} cetanāvataḥ etat bhavati preṣaṇam adhyeṣaṇam ca iti . \\
\noindent \textbf{(P\_3,1.26.2)} bhikṣāḥ ca acetanāḥ . \\
\noindent \textbf{(P\_3,1.26.2)} na eṣaḥ doṣaḥ . \\
\noindent \textbf{(P\_3,1.26.2)} na avaśyam saḥ eva vāsam prayojayati yaḥ āha uṣyatām iti . \\
\noindent \textbf{(P\_3,1.26.2)} tūṣṇīm āsīnaḥ yaḥ tatsamarthāni ācarati saḥ api vāsam prayojayati . \\
\noindent \textbf{(P\_3,1.26.2)} bhikṣāḥ ca api pracurāḥ vyañjanavatyaḥ labhyamānāḥ vāsam prayojayanti . \\
\noindent \textbf{(P\_3,1.26.2)} tathā kārīṣaḥ agniḥ nirvāte ekānte suprajvalitaḥ adhyayanam prayojayati . \\
\noindent \textbf{(P\_3,1.26.3)} iha kaḥ cit kam cit āha . \\
\noindent \textbf{(P\_3,1.26.3)} pṛcchatu mā bhavān . \\
\noindent \textbf{(P\_3,1.26.3)} anuyuṅktām mā bhavān iti . \\
\noindent \textbf{(P\_3,1.26.3)} atra ṇic kasmāt na bhavati . \\
\noindent \textbf{(P\_3,1.26.3)} akartṛtvāt . \\
\noindent \textbf{(P\_3,1.26.3)} na hi asau samprati pṛcchati . \\
\noindent \textbf{(P\_3,1.26.3)} tūṣṇīm āste . \\
\noindent \textbf{(P\_3,1.26.3)} kim ca bhoḥ vartamānakālāyāḥ eva kriyāyāḥ kartrā bhavitavyam na bhūtabhaviṣyatkālāyāḥ . \\
\noindent \textbf{(P\_3,1.26.3)} bhūtabhaviṣyatkālāyāḥ api bhavitavyam . \\
\noindent \textbf{(P\_3,1.26.3)} abhisambandhaḥ tatra kriyate . \\
\noindent \textbf{(P\_3,1.26.3)} imām kriyām akārṣīt . \\
\noindent \textbf{(P\_3,1.26.3)} imām kriyām kariṣyati iti . \\
\noindent \textbf{(P\_3,1.26.3)} iha punaḥ na kaḥ cit abhisambandhaḥ kriyate na ca asau samprati pṛcchati . \\
\noindent \textbf{(P\_3,1.26.3)} tūṣṇīm āste . \\
\noindent \textbf{(P\_3,1.26.3)} yadi tarhi kartā na asti katham tarhi kartṛpratyayena loṭā abhidhīyate . \\
\noindent \textbf{(P\_3,1.26.3)} atham katham asmin apṛcchati ayam pracchiḥ vartate . \\
\noindent \textbf{(P\_3,1.26.3)} abhisambandhaḥ tatra kriyate . \\
\noindent \textbf{(P\_3,1.26.3)} imām kriyām kuru iti . \\
\noindent \textbf{(P\_3,1.26.3)} kartrā api tarhi abhisambandhaḥ kriyate . \\
\noindent \textbf{(P\_3,1.26.3)} katham . \\
\noindent \textbf{(P\_3,1.26.3)} kartā ca asyāḥ kriyāyāḥ bhava iti . \\
\noindent \textbf{(P\_3,1.26.3)} evam na ca kartā kartṛpratyayena ca loṭā abhidhīyate . \\
\noindent \textbf{(P\_3,1.26.3)} atha api katham cit kartā syāt . \\
\noindent \textbf{(P\_3,1.26.3)} evam api na doṣaḥ . \\
\noindent \textbf{(P\_3,1.26.3)} loṭā uktatvāt preṣaṇasya ṇic na bhaviṣyati . \\
\noindent \textbf{(P\_3,1.26.3)} vidhīyante hi eteṣu artheṣu praiṣādiṣu loḍādayaḥ . \\
\noindent \textbf{(P\_3,1.26.3)} yatra ca dvitīyaḥ prayojyaḥ arthaḥ bhavati bhavati tatra ṇic . \\
\noindent \textbf{(P\_3,1.26.3)} tat yathā āsaya śāyaya iti . \\
\noindent \textbf{(P\_3,1.26.4)} kṛṣyādiṣu ca anutpattiḥ . \\
\noindent \textbf{(P\_3,1.26.4)} kṛṣyādiṣu ca anutpattiḥ vaktavyā . \\
\noindent \textbf{(P\_3,1.26.4)} ekānte tūṣṇīm āsīnaḥ ucyate pañcabhiḥ halaiḥ kṛṣati iti . \\
\noindent \textbf{(P\_3,1.26.4)} tatra bhavitavyam . \\
\noindent \textbf{(P\_3,1.26.4)} pañcabhiḥ halaiḥ karṣayati iti . \\
\noindent \textbf{(P\_3,1.26.4)} kṛṣyādiṣu ca anutpattiḥ nānākriyāṇām kṛṣyarthatvāt . \\
\noindent \textbf{(P\_3,1.26.4)} kṛṣyādiṣu ca anutpattiḥ siddhā . \\
\noindent \textbf{(P\_3,1.26.4)} kutaḥ . \\
\noindent \textbf{(P\_3,1.26.4)} nānākriyāṇām kṛṣyarthatvāt . \\
\noindent \textbf{(P\_3,1.26.4)} nānākriyāḥ kṛṣeḥ arthāḥ . \\
\noindent \textbf{(P\_3,1.26.4)} na avaśyam kṛṣiḥ vilekhane eva vartate . \\
\noindent \textbf{(P\_3,1.26.4)} kim tarhi. pratividhāne api vartate . \\
\noindent \textbf{(P\_3,1.26.4)} yat asau bhaktabījabalīvardaiḥ pratividhānam karoti saḥ kṛṣyarthaḥ . \\
\noindent \textbf{(P\_3,1.26.4)} ātaḥ ca pratividhāne vartate . \\
\noindent \textbf{(P\_3,1.26.4)} yadahaḥ eva asau na pratividhatte tadahaḥ tat karma na pravartate . \\
\noindent \textbf{(P\_3,1.26.4)} yajyādiṣu ca aviparyāsaḥ . \\
\noindent \textbf{(P\_3,1.26.4)} yajyādiṣu ca aviparyāsaḥ vaktavyaḥ . \\
\noindent \textbf{(P\_3,1.26.4)} puṣyamitraḥ yajate . \\
\noindent \textbf{(P\_3,1.26.4)} yājakāḥ yājayanti iti . \\
\noindent \textbf{(P\_3,1.26.4)} tatra bhavitavyam . \\
\noindent \textbf{(P\_3,1.26.4)} puṣyamitraḥ yājayate . \\
\noindent \textbf{(P\_3,1.26.4)} yājakāḥ yajanti iti . \\
\noindent \textbf{(P\_3,1.26.4)} yajyādiṣu ca aviparyāsaḥ nānākriyāṇām yajyarthatvāt . \\
\noindent \textbf{(P\_3,1.26.4)} yajyādiṣu ca aviparyāsaḥ siddhaḥ . \\
\noindent \textbf{(P\_3,1.26.4)} kutaḥ . \\
\noindent \textbf{(P\_3,1.26.4)} nānākriyāṇām yajyarthatvāt . \\
\noindent \textbf{(P\_3,1.26.4)} nānākriyāḥ yajeḥ arthāḥ . \\
\noindent \textbf{(P\_3,1.26.4)} na avaśyam yajiḥ haviṣprakṣepaṇe eva vartate . \\
\noindent \textbf{(P\_3,1.26.4)} kim tarhi . \\
\noindent \textbf{(P\_3,1.26.4)} tyāge api vartate . \\
\noindent \textbf{(P\_3,1.26.4)} aho yajate iti ucyate yaḥ suṣṭhu tyāgam karoti . \\
\noindent \textbf{(P\_3,1.26.4)} tam ca puṣyamitraḥ karoti . \\
\noindent \textbf{(P\_3,1.26.4)} yājakāḥ prayojayanti . \\
\noindent \textbf{(P\_3,1.26.5)} tat karoti iti upasaṅkhyānam sūtrayatyādyartham . \\
\noindent \textbf{(P\_3,1.26.5)} tat karoti iti upasaṅkhyānam kartavyam . \\
\noindent \textbf{(P\_3,1.26.5)} kim prayojanam . \\
\noindent \textbf{(P\_3,1.26.5)} sūtrayatyādyartham . \\
\noindent \textbf{(P\_3,1.26.5)} sūtram karoti . \\
\noindent \textbf{(P\_3,1.26.5)} sūtrayati . \\
\noindent \textbf{(P\_3,1.26.5)} iha vyākaraṇasya sūtram karoti . \\
\noindent \textbf{(P\_3,1.26.5)} vyākaraṇam sūtrayati iti . \\
\noindent \textbf{(P\_3,1.26.5)} vākye ṣaṣṭhī utpanne ca pratyaye dvitīyā . \\
\noindent \textbf{(P\_3,1.26.5)} kena etat evam bhavati . \\
\noindent \textbf{(P\_3,1.26.5)} yaḥ asau sūtravyākaraṇayoḥ abhisambandhaḥ saḥ utpanne pratyaye nivartate . \\
\noindent \textbf{(P\_3,1.26.5)} asti ca karoteḥ vyākaraṇena sāmarthyam iti kṛtvā dvitīyā bhaviṣyati . \\
\noindent \textbf{(P\_3,1.26.6)} ākhyānāt kṛtaḥ tat ācaṣṭe iti kṛlluk prakṛtipratyāpattiḥ prakṛtivat ca kārakam . ākhyānāt kṛdantāt tat ācaṣṭe iti etasmin arthe kṛlluk prakṛtipratyāpattiḥ prakṛtivat ca kārakam bhavati iti vaktavyam . \\
\noindent \textbf{(P\_3,1.26.6)} kaṃsavadham ācaṣṭe kaṃsam ghātayati . \\
\noindent \textbf{(P\_3,1.26.6)} balibandham ācaṣṭe balim bandhayati . \\
\noindent \textbf{(P\_3,1.26.6)} ākhyānāt ca pratiṣedhaḥ . \\
\noindent \textbf{(P\_3,1.26.6)} ākhyānaśabdāt ca pratiṣedhaḥ vaktavyaḥ . \\
\noindent \textbf{(P\_3,1.26.6)} ākhyānam ācaṣṭe . \\
\noindent \textbf{(P\_3,1.26.6)} kim punaḥ yāni etāni sañjñābhūtāni ākhyānāni tataḥ utpattyā bhavitavyam āhosvit kriyānvākhyānamātrāt . \\
\noindent \textbf{(P\_3,1.26.6)} kim ca ataḥ . \\
\noindent \textbf{(P\_3,1.26.6)} yadi sañjñābhūtebhyaḥ iha na prāpnoti . \\
\noindent \textbf{(P\_3,1.26.6)} rājāgamanam ācaṣṭe ṛajānam āgamayati . \\
\noindent \textbf{(P\_3,1.26.6)} atha kriyānvākhyānamātrāt na doṣaḥ bhavati . \\
\noindent \textbf{(P\_3,1.26.6)} yathā na doṣaḥ tathā astu . \\
\noindent \textbf{(P\_3,1.26.6)} dṛśyarthānām ca pravṛttau . \\
\noindent \textbf{(P\_3,1.26.6)} dṛśyarthānām ca pravṛttau kṛdantāt ṇic vaktayaḥ tat ācaṣṭe iti etasmin arthe kṛlluk prakṛtipratyāpattiḥ prakṛtivat ca kārakam bhavati iti . \\
\noindent \textbf{(P\_3,1.26.6)} mṛgaramaṇam ācaṣṭe mṛgān ramayati iti . \\
\noindent \textbf{(P\_3,1.26.6)} dṛśyarthānām iti kimartham . \\
\noindent \textbf{(P\_3,1.26.6)} yadā hi grāme mṛgaramaṇam ācaṣṭe mṛgaramaṇam ācaṣṭe iti eva tadā bhavati iti . \\
\noindent \textbf{(P\_3,1.26.6)} āṅlopaḥ ca kālātyantasaṃyoge maryādāyām . \\
\noindent \textbf{(P\_3,1.26.6)} kālātyantasaṃyoge maryādayām kṛdantāt ṇic vaktayaḥ tat ācaṣṭe iti etasmin arthe āṅlopaḥ ca kṛlluk prakṛtipratyāpattiḥ prakṛtivat ca kārakam bhavati iti . \\
\noindent \textbf{(P\_3,1.26.6)} ārātrimvivāsam ācaṣṭe rātrim vivāsayati iti . \\
\noindent \textbf{(P\_3,1.26.6)} citrīkaraṇe prāpi . \\
\noindent \textbf{(P\_3,1.26.6)} citrīkaraṇe prāpyarthe kṛdantāt ṇic vaktayaḥ kṛlluk prakṛtipratyāpattiḥ prakṛtivat ca kārakam bhavati iti . \\
\noindent \textbf{(P\_3,1.26.6)} ujjayinyāḥ prasthitaḥ māhiṣmatyām suryodgamanam sambhāvayate sūryam udgamayati . \\
\noindent \textbf{(P\_3,1.26.6)} nakṣatrayoge jñi . \\
\noindent \textbf{(P\_3,1.26.6)} nakṣatrayoge jānātyarthe kṛdantāt ṇic vaktayaḥ kṛlluk prakṛtipratyāpattiḥ prakṛtivat ca kārakam bhavati iti . \\
\noindent \textbf{(P\_3,1.26.6)} puṣyayogam jānāti puṣyeṇa yojayati . \\
\noindent \textbf{(P\_3,1.26.6)} maghābhiḥ yojayati . \\
\noindent \textbf{(P\_3,1.26.6)} tat tarhi bahu vaktavyam . \\
\noindent \textbf{(P\_3,1.26.6)} na vā sāmānyakṛtatvāt hetutaḥ hi aviśiṣṭam . \\
\noindent \textbf{(P\_3,1.26.6)} na vā vaktavyam . \\
\noindent \textbf{(P\_3,1.26.6)} kim kāraṇam . \\
\noindent \textbf{(P\_3,1.26.6)} sāmānyakṛtatvāt . \\
\noindent \textbf{(P\_3,1.26.6)} sāmānyena eva atra ṇic bhaviṣyati . \\
\noindent \textbf{(P\_3,1.26.6)} hetumati iti . \\
\noindent \textbf{(P\_3,1.26.6)} kim kāraṇam . \\
\noindent \textbf{(P\_3,1.26.6)} hetutaḥ hi aviśiṣṭam . \\
\noindent \textbf{(P\_3,1.26.6)} hetutaḥ hi aviśiṣṭam bhavati . \\
\noindent \textbf{(P\_3,1.26.6)} tulyā hi hetutā devadatte ca āditye ca . \\
\noindent \textbf{(P\_3,1.26.6)} na sidhyati . \\
\noindent \textbf{(P\_3,1.26.6)} svatantraprayojakaḥ hetusañjñaḥ bhavati iti ucyate . \\
\noindent \textbf{(P\_3,1.26.6)} na ca asau ādityam prayojayati . \\
\noindent \textbf{(P\_3,1.26.6)} svatantraprayojakatvāt aprayojakaḥ iti cet muktasaṃśayena tulyam . \\
\noindent \textbf{(P\_3,1.26.6)} yam bhavān svatrantraprayojakam muktasaṃśayam nyāyyam manyate pācayati odanam devadattaḥ yajñadattena iti tena etat tulyam . \\
\noindent \textbf{(P\_3,1.26.6)} katham . \\
\noindent \textbf{(P\_3,1.26.6)} pravṛttiḥ hi ubhayatra anapekṣya . \\
\noindent \textbf{(P\_3,1.26.6)} pravṛttiḥ hi ubhayatra anapekṣya eva kim cit bhavati devadatte ca āditye ca . \\
\noindent \textbf{(P\_3,1.26.6)} na iha kaḥ cit paraḥ anugrahītavyaḥ iti pravartate . \\
\noindent \textbf{(P\_3,1.26.6)} sarve ime svabhūtyartham pravartante . \\
\noindent \textbf{(P\_3,1.26.6)} ye tāvat ete guruśuśrūṣavaḥ te api svabhūtyartham eva pravartante pāralaukikam ca naḥ bhaviṣyati iha ca naḥ prītaḥ guruḥ adhyāpayiṣyati iti . \\
\noindent \textbf{(P\_3,1.26.6)} tathā yat etat dāsakarmakaram nāma ete api svabhūtyartham eva pravartantebhaktam celam ca lapsyāmahe paribhāṣāḥ ca na naḥ bhaviṣyanti iti . \\
\noindent \textbf{(P\_3,1.26.6)} tathā ye ete śilpinaḥ nāme te api svabhūtyartham eva pravartante vetanam ca lapsyāmahe mitrāṇi ca naḥ bhaviṣyanti iti . \\
\noindent \textbf{(P\_3,1.26.6)} evam eteṣu sarveṣu svabhūtyartham pravartamāneṣu kurvataḥ prayojakaḥ iti cet tulyam . \\
\noindent \textbf{(P\_3,1.26.6)} yadi kaḥ cit kurvataḥ prayojakaḥ nāma bhavati tena etat tulyam . \\
\noindent \textbf{(P\_3,1.26.6)} yadi tarhi sarve ime svabhūtyartham pravartantekaḥ prayojyārthaḥ . \\
\noindent \textbf{(P\_3,1.26.6)} yat abhiprāyeṣu sajjante . \\
\noindent \textbf{(P\_3,1.26.6)} īdṛśau vadhrau kuru . \\
\noindent \textbf{(P\_3,1.26.6)} īdṛśau paṭukau kuru . \\
\noindent \textbf{(P\_3,1.26.6)} ādityaḥ ca asya abhiprāye sajjate . \\
\noindent \textbf{(P\_3,1.26.6)} eṣaḥ tasya abhiprāyaḥ . \\
\noindent \textbf{(P\_3,1.26.6)} ujjayinyāḥ prasthitaḥ māhiṣmatyām suryodgamanam sambhāvayeya iti . \\
\noindent \textbf{(P\_3,1.26.6)} tam ca asya abhiprāyam ādityaḥ nirvartayati . \\
\noindent \textbf{(P\_3,1.26.6)} bhavet iha vartamānakālatā yuktā . \\
\noindent \textbf{(P\_3,1.26.6)} ujjayinyāḥ prasthitaḥ māhiṣmatyām suryodgamanam sambhāvayate sūryam udgamayati iti . \\
\noindent \textbf{(P\_3,1.26.6)} tatrasthasya hi tasya ādityaḥ udeti . \\
\noindent \textbf{(P\_3,1.26.6)} iha tu katham vartamānakālatām kaṃsam ghātayati balim bandhayati iti cirahate kaṃse cirabaddhe ca balau . \\
\noindent \textbf{(P\_3,1.26.6)} atra api yuktā . \\
\noindent \textbf{(P\_3,1.26.6)} katham . \\
\noindent \textbf{(P\_3,1.26.6)} ye tāvat ete śobhikāḥ nāma ete pratyakṣam kaṃsam ghātayanti pratyakṣam ca balim bandhayanti iti . \\
\noindent \textbf{(P\_3,1.26.6)} citreṣu katham . \\
\noindent \textbf{(P\_3,1.26.6)} citreṣu api udgūrṇāḥ nipatitāḥ ca prahārāḥ dṛśyante kaṃsakarṣaṇyaḥ ca . \\
\noindent \textbf{(P\_3,1.26.6)} granthikeṣu katham yatra śabdagaḍumātram lakṣyate . \\
\noindent \textbf{(P\_3,1.26.6)} te api hi teṣām utpattiprabhṛti ā vināśāt ṛddhīḥ vyācakṣāṇāḥ sataḥ buddhiviṣayān prakāśayanti . \\
\noindent \textbf{(P\_3,1.26.6)} ātaḥ ca sataḥ vyāmiśrāḥ hi dṛśyante . \\
\noindent \textbf{(P\_3,1.26.6)} ke cit kaṃsabhaktāḥ bhavanti ke cit vāsudevabhaktāḥ . \\
\noindent \textbf{(P\_3,1.26.6)} varṇānyatvam khalu api puṣyanti . \\
\noindent \textbf{(P\_3,1.26.6)} ke cit raktamukhāḥ bhavanti ke cit kālamukhāḥ . \\
\noindent \textbf{(P\_3,1.26.6)} traikālyam khalu api loke lakṣyate . \\
\noindent \textbf{(P\_3,1.26.6)} gaccha hanyate kaṃsaḥ . \\
\noindent \textbf{(P\_3,1.26.6)} gaccha ghāniṣyate kaṃsaḥ . \\
\noindent \textbf{(P\_3,1.26.6)} kim gatena hataḥ kaṃsaḥ iti . \\
\noindent \textbf{(P\_3,1.27)} kimarthaḥ kakāraḥ . \\
\noindent \textbf{(P\_3,1.27)} kṅiti iti guṇapratiṣedhaḥ yathā syāt . \\
\noindent \textbf{(P\_3,1.27)} na etat asti prayojanam . \\
\noindent \textbf{(P\_3,1.27)} sārvadhātukārdhadhātukayoḥ aṅgasya guṇaḥ ucyate . \\
\noindent \textbf{(P\_3,1.27)} dhātoḥ ca vihitaḥ pratyayaḥ śeṣaḥ ārdhadhātukasañjñām labhate . \\
\noindent \textbf{(P\_3,1.27)} na ca ayam dhātoḥ vidhīyate . \\
\noindent \textbf{(P\_3,1.27)} kaṇḍvādīni hi prātipadikāni . \\
\noindent \textbf{(P\_3,1.27)} kaṇḍvādibhyaḥ vāvacanam . \\
\noindent \textbf{(P\_3,1.27)} kaṇḍvādibhyaḥ vā iti vaktavyam . \\
\noindent \textbf{(P\_3,1.27)} avacane hi nityapratyayatvam . \\
\noindent \textbf{(P\_3,1.27)} akriyamāṇe hi vāvacane nityaḥ pratyayavidhiḥ prasajyeta . \\
\noindent \textbf{(P\_3,1.27)} tatra kaḥ doṣaḥ . \\
\noindent \textbf{(P\_3,1.27)} tatra dhātuvidhitukpratiṣedhaḥ . \\
\noindent \textbf{(P\_3,1.27)} tatra dhātuvidheḥ tukaḥ ca pratiṣedhaḥ vaktavyaḥ syāt . \\
\noindent \textbf{(P\_3,1.27)} kaṇḍvau kaṇḍvaḥ . \\
\noindent \textbf{(P\_3,1.27)} aci śnudhātubhruvām yvoḥ iyaṅuvaṅau iti uvaṅadeśaḥ prasajyeta . \\
\noindent \textbf{(P\_3,1.27)} iha ca kaṇḍvā kaṇḍve na ūṅdhātvoḥ iti pratiṣedhaḥ prasajyeta . \\
\noindent \textbf{(P\_3,1.27)} tuk ca pratiṣedhyaḥ . \\
\noindent \textbf{(P\_3,1.27)} valguḥ mantuḥ iti . \\
\noindent \textbf{(P\_3,1.27)} hrasvasya piti kṛti tuk prāpnoti . \\
\noindent \textbf{(P\_3,1.27)} hrasvayalopau ca vaktavyau . \\
\noindent \textbf{(P\_3,1.27)} hrasvayalopau ca vaktavyau syātām . \\
\noindent \textbf{(P\_3,1.27)} valguḥ mantuḥ iti . \\
\noindent \textbf{(P\_3,1.27)} kimartham idam na hrasvaḥ eva ayam . \\
\noindent \textbf{(P\_3,1.27)} antaraṅgatvāt akṛdyakāre iti dīrghatvam prāpnoti . \\
\noindent \textbf{(P\_3,1.27)} yalopaḥ . \\
\noindent \textbf{(P\_3,1.27)} yalopaḥ ca vaktavyaḥ . \\
\noindent \textbf{(P\_3,1.27)} kaṇḍūḥ valguḥ mantuḥ iti . \\
\noindent \textbf{(P\_3,1.27)} kimartham idam na vali iti eva siddham . \\
\noindent \textbf{(P\_3,1.27)} vali iti ucyate . \\
\noindent \textbf{(P\_3,1.27)} na ca atra valim paśyāmaḥ . \\
\noindent \textbf{(P\_3,1.27)} nanu cal kvip valādiḥ . \\
\noindent \textbf{(P\_3,1.27)} kviblope kṛte valādyabhāvāt na prāpnoti . \\
\noindent \textbf{(P\_3,1.27)} idam iha sampradhāryam . \\
\noindent \textbf{(P\_3,1.27)} kviblopaḥ kriyatām vali lopaḥ iti . \\
\noindent \textbf{(P\_3,1.27)} kim atra kartavyam . \\
\noindent \textbf{(P\_3,1.27)} paratvāt kviblopaḥ . \\
\noindent \textbf{(P\_3,1.27)} nityaḥ khalu api kviblopaḥ . \\
\noindent \textbf{(P\_3,1.27)} kṛte api yalope prāpnoti akṛte api prāpnoti . \\
\noindent \textbf{(P\_3,1.27)} nityatvāt paratvāt ca kvilope kṛte valādyabhāvāt na prāpnoti . \\
\noindent \textbf{(P\_3,1.27)} evam tarhi pratyayalakṣaṇena bhaviṣyati . \\
\noindent \textbf{(P\_3,1.27)} varṇāśraye na asti pratyayalakṣaṇam . \\
\noindent \textbf{(P\_3,1.27)} atha kriyamāṇe api vāvacane yadā yagantāt kvip tadā ete doṣāḥ kasmāt na bhavanti . \\
\noindent \textbf{(P\_3,1.27)} na etebhyaḥ tadā kvip drakṣyate . \\
\noindent \textbf{(P\_3,1.27)} kim kāraṇam . \\
\noindent \textbf{(P\_3,1.27)} anyebhyaḥ api dṛśyate iti ucyate . \\
\noindent \textbf{(P\_3,1.27)} na ca etebhyaḥ tadā kvip dṛśyate . \\
\noindent \textbf{(P\_3,1.27)} yathā eva tarhi kriyamāṇe vāvacane anyebhyaḥ api dṛśyate iti evam atra kvip na bhavati evam akriyamāṇe api na bhaviṣyati . \\
\noindent \textbf{(P\_3,1.27)} avaśyam etebhyaḥ tadā kvip eṣitavyaḥ . \\
\noindent \textbf{(P\_3,1.27)} kim prayojanam . \\
\noindent \textbf{(P\_3,1.27)} etāni rūpāṇi yathā syuḥ iti . \\
\noindent \textbf{(P\_3,1.27)} tat tarhi vāvacanam kartavyam . \\
\noindent \textbf{(P\_3,1.27)} na kartavyam . \\
\noindent \textbf{(P\_3,1.27)} ubhayam kaṇḍvādīni dhātavaḥ ca prātipadikāni ca . \\
\noindent \textbf{(P\_3,1.27)} ātaḥ ca ubhayam . \\
\noindent \textbf{(P\_3,1.27)} kaṇḍūyati iti kriyām kurvāṇe prayujyate asti me kaṇḍūḥ iti vedanāmātrasya sānnidhye . \\
\noindent \textbf{(P\_3,1.27)} aparaḥ āha : dhātuprakaraṇāt dhātuḥ kasya āsañjanāt api . \\
\noindent \textbf{(P\_3,1.27)} āha ca ayam imam dīrgham . \\
\noindent \textbf{(P\_3,1.27)} manye dhātuḥ vibhāṣitaḥ . \\
\noindent \textbf{(P\_3,1.30)} kimarthaḥ ayam ṇakāraḥ . \\
\noindent \textbf{(P\_3,1.30)} vṛddhyarthaḥ . \\
\noindent \textbf{(P\_3,1.30)} ñṇiti iti vṛddhiḥ yathā syāt . \\
\noindent \textbf{(P\_3,1.30)} kriyamāṇe api vai ṇakāre vṛddhiḥ na prāpnoti . \\
\noindent \textbf{(P\_3,1.30)} kim kāraṇam . \\
\noindent \textbf{(P\_3,1.30)} kṅiti ca iti pratiṣedhāt . \\
\noindent \textbf{(P\_3,1.30)} ṇitkaraṇasāmarthyāt bhaviṣyati . \\
\noindent \textbf{(P\_3,1.30)} ataḥ uttaram paṭhati . \\
\noindent \textbf{(P\_3,1.30)} ṇiṅi ṇitkaraṇasya sāvakāśatvāt vṛddhipratiṣedhaprasaṅgaḥ . \\
\noindent \textbf{(P\_3,1.30)} ṇiṅi ṇitkaraṇam sāvakāśam . \\
\noindent \textbf{(P\_3,1.30)} kaḥ avakāśaḥ . \\
\noindent \textbf{(P\_3,1.30)} sāmānyagrahaṇārthaḥ ṇakāraḥ . \\
\noindent \textbf{(P\_3,1.30)} kva sāmānyagrahaṇārthena arthaḥ . \\
\noindent \textbf{(P\_3,1.30)} ṇeḥ aniṭi iti . \\
\noindent \textbf{(P\_3,1.30)} ṇiṅi ṇitkaraṇasya sāvakāśatvāt vṛddhipratiṣedhaḥ prāpnoti . \\
\noindent \textbf{(P\_3,1.30)} ṅitkaraṇam api tarhi sāvakāśam . \\
\noindent \textbf{(P\_3,1.30)} kaḥ avakāśaḥ . \\
\noindent \textbf{(P\_3,1.30)} sāmānyagrahaṇāvighātārthaḥ ṅakāraḥ . \\
\noindent \textbf{(P\_3,1.30)} kva sāmānyagrahaṇāvighātārthena arthaḥ . \\
\noindent \textbf{(P\_3,1.30)} atra eva . \\
\noindent \textbf{(P\_3,1.30)} śakyaḥ atra sāmānyagrahaṇāvighātārthaḥ anyaḥ anubandhaḥ āsaṅktum . \\
\noindent \textbf{(P\_3,1.30)} tatra ṅakārānurodhāt vṛddhipratiṣedhaḥ prāpnoti . \\
\noindent \textbf{(P\_3,1.30)} avaśayam atra ātmanepadārthaḥ ṅakāraḥ anubandhaḥ āsaṅktavyaḥ ṅitaḥ iti ātmanepadam yathā syāt . \\
\noindent \textbf{(P\_3,1.30)} evam ubhayoḥ sāvakaśayoḥ pratiṣedhabalīyastvāt pratiṣedhaḥ prāpnoti . \\
\noindent \textbf{(P\_3,1.30)} evam tarhi ācāryapravṛttiḥ jñāpayati na kameḥ vṛddhipratiṣedhaḥ bhavati iti yat ayam na kamyamicamām iti mitsañjñāyā pratiṣedham śāsti . \\
\noindent \textbf{(P\_3,1.30)} mitpratiṣedhasya ca arthavattvāt . \\
\noindent \textbf{(P\_3,1.30)} mitpratiṣedhasya ca arthavattvāt pratiṣedhaḥ prāpnoti . \\
\noindent \textbf{(P\_3,1.30)} arthavān mitpratiṣedhaḥ . \\
\noindent \textbf{(P\_3,1.30)} kaḥ arthaḥ . \\
\noindent \textbf{(P\_3,1.30)} ṇiṅantasya ṇici yā vṛddhiḥ tasyāḥ hrasvatvam mā bhūt iti . \\
\noindent \textbf{(P\_3,1.30)} nanu etasyāḥ api kṅiti ca iti pratiṣedhena bhavitavyam . \\
\noindent \textbf{(P\_3,1.30)} na bhavitavyam . \\
\noindent \textbf{(P\_3,1.30)} uktam etat kṅiti pratiṣedhe tannimittagrahaṇam iti . \\
\noindent \textbf{(P\_3,1.30)} evam tarhi na ṇiṅantasya ṇici yā vṛddhiḥ tasyāḥ hrasvatvam prāpnoti . \\
\noindent \textbf{(P\_3,1.30)} kim kāraṇam . \\
\noindent \textbf{(P\_3,1.30)} ṇiṅā vyavahitatvāt . \\
\noindent \textbf{(P\_3,1.30)} lope kṛte na asti vyavadhānam . \\
\noindent \textbf{(P\_3,1.30)} sthānivadbhāvāt vyavadhānam eva . \\
\noindent \textbf{(P\_3,1.30)} ṇiṅi eva tarhi mā bhūt iti . \\
\noindent \textbf{(P\_3,1.30)} ṇiṅi ca na prāpnoti . \\
\noindent \textbf{(P\_3,1.30)} kim kāraṇam . \\
\noindent \textbf{(P\_3,1.30)} asiddham bahiraṅgalakṣaṇam antaraṅgalakṣaṇe iti . \\
\noindent \textbf{(P\_3,1.30)} na eva vā punaḥ ṇiṅantasya ṇici vṛddhiḥ prāpnoti . \\
\noindent \textbf{(P\_3,1.30)} kim kāraṇam . \\
\noindent \textbf{(P\_3,1.30)} ṇiṅā vyavahitatvāt . \\
\noindent \textbf{(P\_3,1.30)} lope kṛte na asti vyavadhānam . \\
\noindent \textbf{(P\_3,1.30)} sthānivadbhāvāt vyavadhānam eva . \\
\noindent \textbf{(P\_3,1.30)} idam tarhi prayojanam . \\
\noindent \textbf{(P\_3,1.30)} yat tat ciṇṇamuloḥ dīrghaḥ anyatarasyām iti dīrghatvam tat kameḥ ṇiṅi mā bhūt iti . \\
\noindent \textbf{(P\_3,1.30)} kim punaḥ kāraṇam tatra dīrghaḥ anyatarasyām iti ucyate . \\
\noindent \textbf{(P\_3,1.30)} na hrasvaḥ anyatarasyām iti eva ucyeta . \\
\noindent \textbf{(P\_3,1.30)} yathāprāptam ca api kameḥ hrasvatvam eva . \\
\noindent \textbf{(P\_3,1.30)} tatra ayam api arthaḥ . \\
\noindent \textbf{(P\_3,1.30)} hrasvagrahaṇam na kartavyam bhavati . \\
\noindent \textbf{(P\_3,1.30)} prakṛtam anuvartate . \\
\noindent \textbf{(P\_3,1.30)} kva prakṛtam . \\
\noindent \textbf{(P\_3,1.30)} mitām hrasvaḥ iti . \\
\noindent \textbf{(P\_3,1.30)} kā rūpasiddhiḥ : aśami aśāmi śamam śamam śāmam śāmam . \\
\noindent \textbf{(P\_3,1.30)} vṛddhyā siddham . \\
\noindent \textbf{(P\_3,1.30)} na sidhyati . \\
\noindent \textbf{(P\_3,1.30)} na sidhyati . \\
\noindent \textbf{(P\_3,1.30)} na udāttopadeśasya māntasya anācameḥ iti vṛddhipratiṣedhaḥ prāpnoti . \\
\noindent \textbf{(P\_3,1.30)} ciṇkṛtoḥ saḥ pratiṣedhaḥ na ṇici . \\
\noindent \textbf{(P\_3,1.30)} idam tarhi . \\
\noindent \textbf{(P\_3,1.30)} ajani ajāni janam janam jānam jānam . \\
\noindent \textbf{(P\_3,1.30)} janivadhyoḥ ca iti vṛddhipratiṣedhaḥ prāpnoti . \\
\noindent \textbf{(P\_3,1.30)} saḥ api ciṇkṛtoḥ eva . \\
\noindent \textbf{(P\_3,1.30)} ṇijvyavahiteṣu tarhi yaṅlope ca upasaṅkhyānam kartavyam syāt . \\
\noindent \textbf{(P\_3,1.30)} śamayantam prayojitavān aśami aśāmi śamam śamam śāmam śāmam . \\
\noindent \textbf{(P\_3,1.30)} śaṃśamayateḥ aśaṃsami aśaṃśāmi śaṃśamam śaṃśamam śaṃśāmam śaṃśāmam . \\
\noindent \textbf{(P\_3,1.30)} kim punaḥ kāraṇam na sidhyati . \\
\noindent \textbf{(P\_3,1.30)} ciṇṇamulpare ṇau mitām aṅgānām hrasvaḥ bhavati iti ucyate . \\
\noindent \textbf{(P\_3,1.30)} yaḥ ca atra ṇiḥ ciṇṇamulparaḥ na tasmin mit aṅgam yasmin ca mit aṅgam na asau ṇiḥ ṇamulparaḥ . \\
\noindent \textbf{(P\_3,1.30)} ṇilope kṛte ciṇṇamulparaḥ . \\
\noindent \textbf{(P\_3,1.30)} sthānivadbhāvāt na ciṇṇamulparaḥ . \\
\noindent \textbf{(P\_3,1.30)} atha dīrghaḥ anyatarasyām iti ucyamāne yāvatā sthānivadbhāvaḥ katham eva etat sidhyati . \\
\noindent \textbf{(P\_3,1.30)} etat idānīm dīrghagrahaṇasya prayojanam . \\
\noindent \textbf{(P\_3,1.30)} dīrghavidhim prati ajādeśaḥ na sthānivat iti sthānivadbhāvapratiṣedhaḥ siddhaḥ bhavati . \\
\noindent \textbf{(P\_3,1.30)} yadā khalu api āyādayaḥ ārdhadhātuke vā bhavanti tadā ṇici ṇiṅ na bhavati . \\
\noindent \textbf{(P\_3,1.30)} tadartham ca mitpratiṣedhaḥ syāt . \\
\noindent \textbf{(P\_3,1.30)} tasmāt pratiṣedhaḥ prāpnoti . \\
\noindent \textbf{(P\_3,1.30)} uktam vā . \\
\noindent \textbf{(P\_3,1.30)} kim uktam . \\
\noindent \textbf{(P\_3,1.30)} taddhitakāmyoḥ ikprakaraṇāt iti . \\
\noindent \textbf{(P\_3,1.31)} katham idam vijñāyate . \\
\noindent \textbf{(P\_3,1.31)} āyādibhyaḥ yat ārdhadhātukam tasmin avasthite vā āyādīnām nivṛttiḥ bhavati . \\
\noindent \textbf{(P\_3,1.31)} āhosvit āyādiprakṛteḥ yat ārdhadhātukam tasmin avasthite vā āyādīnām utpattiḥ bhavati iti . \\
\noindent \textbf{(P\_3,1.31)} kim gatam etat iyatā sūtreṇā āhosvit anyatarasmin pakṣe bhūyaḥ sūtram kartavyam . \\
\noindent \textbf{(P\_3,1.31)} gatam iti āha . \\
\noindent \textbf{(P\_3,1.31)} katham . \\
\noindent \textbf{(P\_3,1.31)} yadā tāvat āyādibhyaḥ yat ārdhadhātukam tasmin avasthite vā āyādīnām nivṛttiḥ bhavati iti tadā aviśeṣeṇa sarvam āyādiprakaraṇam anukramya āyādayaḥ ārdhadhātuke vā iti ucyate . \\
\noindent \textbf{(P\_3,1.31)} yadā api āyādiprakṛteḥ yat ārdhadhātukam tasmin avasthite vā āyādīnām utpattiḥ bhavati iti tadā ekam vākyam tat ca idam ca . \\
\noindent \textbf{(P\_3,1.31)} gupūdhūpavicchipaṇipanibhyaḥ āyaḥ ārdhadhātuke vā . \\
\noindent \textbf{(P\_3,1.31)} ṛteḥ īyaṅ ārdhadhātuke vā . \\
\noindent \textbf{(P\_3,1.31)} kameḥ ṇiṅ ārdhadhātuke vā iti . \\
\noindent \textbf{(P\_3,1.31)} kaḥ ca atra viśeṣaḥ . \\
\noindent \textbf{(P\_3,1.31)} āyādibhyaḥ yat ārdhadhātukam āyādiprakṛteḥ yat ārdhadhātukam iti ca ubhayathā aniṣṭaprasaṅgaḥ . āyādibhyaḥ yat ārdhadhātukam āyādiprakṛteḥ yat ārdhadhātukam iti ca ubhayathā aniṣṭam prāpnoti . \\
\noindent \textbf{(P\_3,1.31)} yadi vijñāyate āyādibhyaḥ yat ārdhadhātukam tasmin avasthite vā āyādīnām nivṛttiḥ bhavati iti guptiḥ jugopa iti ca iṣṭam na sidhyati idam ca aniṣṭam prāpnoti . \\
\noindent \textbf{(P\_3,1.31)} gopām cakāra gopā iti ca . \\
\noindent \textbf{(P\_3,1.31)} idam tāvat iṣṭam siddham bhavati . \\
\noindent \textbf{(P\_3,1.31)} gopāyām cakāra gopāya iti . \\
\noindent \textbf{(P\_3,1.31)} atha vijñāyate āyādiprakṛteḥ yat ārdhadhātukam tasmin avasthite vā āyādīnām utpattiḥ bhavati iti guptiḥ jugopa iti ca iṣṭam siddham bhavati . \\
\noindent \textbf{(P\_3,1.31)} idam ca aniṣṭam na prāpnoti . \\
\noindent \textbf{(P\_3,1.31)} gopāyām cakāra gopāya iti . \\
\noindent \textbf{(P\_3,1.31)} idam tu iṣṭam na sidhyati . \\
\noindent \textbf{(P\_3,1.31)} gopayām cakāra gopāya iti . \\
\noindent \textbf{(P\_3,1.31)} idam tāvat iṣṭam sidhyati . \\
\noindent \textbf{(P\_3,1.31)} gopayām cakāra iti . \\
\noindent \textbf{(P\_3,1.31)} katham . \\
\noindent \textbf{(P\_3,1.31)} astu atra āyādiprakṛteḥ yat ārdhadhātukam liṭ . \\
\noindent \textbf{(P\_3,1.31)} tasmin avasthite vā āyādayaḥ . \\
\noindent \textbf{(P\_3,1.31)} ām madhye patiṣyati yathā vikaraṇāḥ tadvat . \\
\noindent \textbf{(P\_3,1.31)} idam tarhi iṣṭam na sidhyati gopāyā iti . \\
\noindent \textbf{(P\_3,1.31)} siddham tu sārvadhātuke nityavacanāt anāśritya vāvidhānam . \\
\noindent \textbf{(P\_3,1.31)} siddham etat . \\
\noindent \textbf{(P\_3,1.31)} katham. aviśeṣeṇa āyādīnām vāvidhānam uktvā sārvadhātuke nityam iti vakṣyāmi . \\
\noindent \textbf{(P\_3,1.31)} syādibalīyastvam tu vipratiṣedhena tulyanimittatvāt . \\
\noindent \textbf{(P\_3,1.31)} syādibhiḥ tu āyādīnām bādhanam prāpnoti vipratiṣedhena . \\
\noindent \textbf{(P\_3,1.31)} kim kāraṇam . \\
\noindent \textbf{(P\_3,1.31)} tulyanimittatvāt . \\
\noindent \textbf{(P\_3,1.31)} tulyam nimittam syādīnām āyādīnām ca . \\
\noindent \textbf{(P\_3,1.31)} syādīnām avakāśaḥ kariṣyati hariṣyati . \\
\noindent \textbf{(P\_3,1.31)} āyādīnām avakāśaḥ gopāyati dhūpāyati . \\
\noindent \textbf{(P\_3,1.31)} iha ubhayam prāpnoti . \\
\noindent \textbf{(P\_3,1.31)} gopāyiṣyati dhūpāyiṣyati iti . \\
\noindent \textbf{(P\_3,1.31)} paratvāt syādayaḥ prāpnuvanti . \\
\noindent \textbf{(P\_3,1.31)} na vā āyādividhānasya anavakāśatvāt . \\
\noindent \textbf{(P\_3,1.31)} na vā eṣaḥ doṣaḥ . \\
\noindent \textbf{(P\_3,1.31)} kim kāraṇam . \\
\noindent \textbf{(P\_3,1.31)} āyādividhānasya anavakāśatvāt . \\
\noindent \textbf{(P\_3,1.31)} anavakāśāḥ āyādayaḥ ucyante ca . \\
\noindent \textbf{(P\_3,1.31)} te vacanāt bhaviṣyanti . \\
\noindent \textbf{(P\_3,1.31)} nanu ca idānīm eva avakāśaḥ prakḷptaḥ gopāyati dhūpāyati iti . \\
\noindent \textbf{(P\_3,1.31)} atra api śap syādiḥ bhavati . \\
\noindent \textbf{(P\_3,1.31)} yadi api atra api bhavati na tu atra asti viśeṣaḥ sati vā śapi asati vā . \\
\noindent \textbf{(P\_3,1.31)} anyat idānīm etat ucyate na asti viśeṣaḥ iti . \\
\noindent \textbf{(P\_3,1.31)} yat tu tat uktam āyādīnām syādibhiḥ avyāptaḥ avakāśaḥ it sa na asti avakāśaḥ . \\
\noindent \textbf{(P\_3,1.31)} avaśyam khalu api atra śap syādiḥ eṣitavyaḥ . \\
\noindent \textbf{(P\_3,1.31)} kim kāraṇam . \\
\noindent \textbf{(P\_3,1.31)} gopāyantī dhūpāyantī iti : śapśyanoḥ nityam iti num yathā syāt iti . \\
\noindent \textbf{(P\_3,1.31)} yadi tarhi anavakāśāḥ āyādayaḥ āyādibhiḥ syādīnām bādhanam prāpnoti . \\
\noindent \textbf{(P\_3,1.31)} yathā punaḥ ayam sūtrebhedena parihāraḥ yadi punaḥ śapi nityam iti ucyeta . \\
\noindent \textbf{(P\_3,1.31)} sidhyati . \\
\noindent \textbf{(P\_3,1.31)} sūtram tarhi bhidyate . \\
\noindent \textbf{(P\_3,1.31)} yathānyāsam eva astu . \\
\noindent \textbf{(P\_3,1.31)} nanu ca uktam āyādibhyaḥ yat ārdhadhātukam āyādiprakṛteḥ yat ārdhadhātukam iti ca ubhayathā aniṣṭaprasaṅgaḥ iti . \\
\noindent \textbf{(P\_3,1.31)} na eṣaḥ doṣaḥ . \\
\noindent \textbf{(P\_3,1.31)} ārdhadhātuke iti na eṣā parasaptamī . \\
\noindent \textbf{(P\_3,1.31)} kā tarhi . \\
\noindent \textbf{(P\_3,1.31)} viṣayasaptamī . \\
\noindent \textbf{(P\_3,1.31)} ārdhadhātukaviṣaye iti . \\
\noindent \textbf{(P\_3,1.31)} tatra ārdhadhātukaviṣaye āyādiprakṛteḥ āyādiṣu kṛteṣu yaḥ yataḥ pratyayaḥ prāpnoti saḥ tataḥ bhaviṣyati . \\
\noindent \textbf{(P\_3,1.32)} antagrahaṇam kimartham na sanādayaḥ dhātavaḥ iti eva ucyeta . \\
\noindent \textbf{(P\_3,1.32)} kena idānīm tadantānām bhaviṣyati . \\
\noindent \textbf{(P\_3,1.32)} tadantavidhinā . \\
\noindent \textbf{(P\_3,1.32)} ataḥ uttaram paṭhati . \\
\noindent \textbf{(P\_3,1.32)} sanādiṣu antagrahaṇe uktam . \\
\noindent \textbf{(P\_3,1.32)} kim uktam . \\
\noindent \textbf{(P\_3,1.32)} padasañjñāyām antagrahaṇam anyatra sañjñāvidhau pratyayagrahaṇe tadantavidhipratiṣedhārtham iti . \\
\noindent \textbf{(P\_3,1.32)} idam ca api pratyayagrahaṇam . \\
\noindent \textbf{(P\_3,1.32)} ayam ca api sañjñāvidhiḥ . \\
\noindent \textbf{(P\_3,1.32)} kimartham punaḥ idam ucyate na bhūvādayaḥ dhātavaḥ iti eva siddham . \\
\noindent \textbf{(P\_3,1.32)} na sidhyati . \\
\noindent \textbf{(P\_3,1.32)} pāṭhena dhātusañjñā kriyate na ca ime tatra paṭhyante . \\
\noindent \textbf{(P\_3,1.32)} katham tarhi anyeṣām apaṭhyamānānām dhātusañjñā bhavati : asteḥ bhūḥ . \\
\noindent \textbf{(P\_3,1.32)} bruvaḥ vaciḥ . \\
\noindent \textbf{(P\_3,1.32)} cakṣiṅaḥ khyāñ iti . \\
\noindent \textbf{(P\_3,1.32)} yadi api ete tatra na paṭhyante prakṛtayaḥ tu eṣām tatra paṭhyante . \\
\noindent \textbf{(P\_3,1.32)} tatra sthānivadbhāvāt siddham . \\
\noindent \textbf{(P\_3,1.32)} ime api tarhi yadi api tatra na paṭhyante yeṣām tu arthāḥ ādiśyante te tatra paṭhyante . \\
\noindent \textbf{(P\_3,1.32)} tatra sthānivadbhāvāt siddham . \\
\noindent \textbf{(P\_3,1.32)} na sidhyati . \\
\noindent \textbf{(P\_3,1.32)} ādeśaḥ sthānivat bhavati iti ucyate . \\
\noindent \textbf{(P\_3,1.32)} na ca ime ādeśāḥ . \\
\noindent \textbf{(P\_3,1.32)} ime api ādeśāḥ . \\
\noindent \textbf{(P\_3,1.32)} katham . \\
\noindent \textbf{(P\_3,1.32)} ādiśyate yaḥ saḥ ādeśaḥ . \\
\noindent \textbf{(P\_3,1.32)} ime ca api ādiśyante . \\
\noindent \textbf{(P\_3,1.32)} evam api ṣaṣṭhīnirdiṣṭasya ādeśāḥ sthānivat bhavanti iti ucyate . \\
\noindent \textbf{(P\_3,1.32)} na ce ime ṣaṣṭhīnirdiṣṭasya ādeśāḥ . \\
\noindent \textbf{(P\_3,1.32)} ṣaṣṭhīgrahaṇam nivartiṣyate . \\
\noindent \textbf{(P\_3,1.32)} yadi nivartate apavāde utsargakṛtam prāpnoti . \\
\noindent \textbf{(P\_3,1.32)} karmaṇi aṇ ātaḥ anupasarge kaḥ iti ke api aṇkṛtam prāpnoti . \\
\noindent \textbf{(P\_3,1.32)} na eṣaḥ doṣaḥ . \\
\noindent \textbf{(P\_3,1.32)} ācāryapravṛttiḥ jñāpayati na apavāde utsargakṛtam bhavati iti yat ayam śyanādīn kān cit śitaḥ karoti . \\
\noindent \textbf{(P\_3,1.32)} śnam śnā śnuḥ iti. \\
\noindent \textbf{(P\_3,1.33)} ime vikaraṇāḥ paṭhyante . \\
\noindent \textbf{(P\_3,1.33)} tatra na jñāyate kaḥ utsargaḥ kaḥ apavādaḥ iti . \\
\noindent \textbf{(P\_3,1.33)} tatra vaktyam : ayam utsargaḥ ayam apavādaḥ iti . \\
\noindent \textbf{(P\_3,1.33)} ime brūmaḥ . \\
\noindent \textbf{(P\_3,1.33)} yak utsargaḥ . \\
\noindent \textbf{(P\_3,1.33)} apavādaḥ śabdādiḥ syādayaḥ ca . \\
\noindent \textbf{(P\_3,1.33)} yadi evam apavādavipratiṣedhāt śabādibādhanam . \\
\noindent \textbf{(P\_3,1.33)} apavādvipratiṣedhāt śabādibhiḥ syādīnām bādhanam prāpnoti . \\
\noindent \textbf{(P\_3,1.33)} śabādīnām avakāśaḥ pacati yajati . \\
\noindent \textbf{(P\_3,1.33)} syādīnām avakāśaḥ pakṣyate yakṣyate . \\
\noindent \textbf{(P\_3,1.33)} iha ubhayam prāpnoti . \\
\noindent \textbf{(P\_3,1.33)} pakṣyati yakṣyati . \\
\noindent \textbf{(P\_3,1.33)} paratvāt śabādayaḥ prāpnuvanti . \\
\noindent \textbf{(P\_3,1.33)} apavādaḥ nāma anekalakṣaṇaprasaṅgaḥ . \\
\noindent \textbf{(P\_3,1.33)} apavādaḥ nāma bhavati yatra anekalakṣaṇaprasaṅgaḥ . \\
\noindent \textbf{(P\_3,1.33)} tatra bhāvakarmaṇoḥ yak vidhīyate kartari śap . \\
\noindent \textbf{(P\_3,1.33)} kaḥ prasaṅgaḥ yat bhāvakarmaṇoḥ yakam kartari śabādayaḥ bādheran . \\
\noindent \textbf{(P\_3,1.33)} evam tarhi yakśapau utsargau . \\
\noindent \textbf{(P\_3,1.33)} apavādāḥ śyanādaya syādayaḥ ca . \\
\noindent \textbf{(P\_3,1.33)} apavādavipratiṣedhāt śyanādibādhanam . \\
\noindent \textbf{(P\_3,1.33)} apavādvipratiṣedhāt śyanādibhiḥ syādīnām bādhanam prāpnoti . \\
\noindent \textbf{(P\_3,1.33)} śyanādīnām avakāśaḥ dīvyati sīvyati . \\
\noindent \textbf{(P\_3,1.33)} syādīnām avakāśaḥ pakṣyati yakṣyati . \\
\noindent \textbf{(P\_3,1.33)} iha ubhayam prāpnoti . \\
\noindent \textbf{(P\_3,1.33)} deviṣyati seviṣyati . \\
\noindent \textbf{(P\_3,1.33)} paratvāt śyanādayaḥ prāpnuvanti . \\
\noindent \textbf{(P\_3,1.33)} na eṣaḥ doṣaḥ . \\
\noindent \textbf{(P\_3,1.33)} śabādeśāḥ śyanādayaḥ kariṣyante . \\
\noindent \textbf{(P\_3,1.33)} śap ca syādibhiḥ bādhyate . \\
\noindent \textbf{(P\_3,1.33)} tatra divādibhyaḥ syādiviṣaye śap eva na asti kutaḥ śyanādayaḥ . \\
\noindent \textbf{(P\_3,1.33)} tat tarhi śapaḥ grahaṇam kartavyam . \\
\noindent \textbf{(P\_3,1.33)} na kartavyam . \\
\noindent \textbf{(P\_3,1.33)} prakṛtam anuvartate . \\
\noindent \textbf{(P\_3,1.33)} kva prakṛtam . \\
\noindent \textbf{(P\_3,1.33)} kartari śap iti . \\
\noindent \textbf{(P\_3,1.33)} tat vai prathamānirdiṣṭam ṣaṣṭhīnirdiṣṭena ca iha arthaḥ . \\
\noindent \textbf{(P\_3,1.33)} divādibhyaḥ iti eṣā pañcamī śap iti prathamāyāḥ ṣaṣṭhīm prakalpayiṣyati tasmāt iti uttarasya iti . \\
\noindent \textbf{(P\_3,1.33)} pratyayavidhiḥ ayam . \\
\noindent \textbf{(P\_3,1.33)} na ca pratyayavidhau pañcamyaḥ prakalpikāḥ bhavanti . \\
\noindent \textbf{(P\_3,1.33)} na ayam pratyayavidhiḥ . \\
\noindent \textbf{(P\_3,1.33)} vihitaḥ pratyayaḥ . \\
\noindent \textbf{(P\_3,1.33)} prakṛtaḥ ca anuvartate . \\
\noindent \textbf{(P\_3,1.33)} atha vā anuvṛttiḥ kariṣyate . \\
\noindent \textbf{(P\_3,1.33)} sārvadhātuke yak syatāsī lṛluṭoḥ cli luṅi cleḥ sic bhavati . \\
\noindent \textbf{(P\_3,1.33)} kartari śap syatāsī lṛluṭoḥ cli luṅi cleḥ sic bhavati . \\
\noindent \textbf{(P\_3,1.33)} divādibhyaḥ śyan syatāsī lṛluṭoḥ cli luṅi cleḥ sic bhavati . \\
\noindent \textbf{(P\_3,1.33)} atha vā antaraṅgāḥ syādayaḥ . \\
\noindent \textbf{(P\_3,1.33)} kā antaraṅgatā . \\
\noindent \textbf{(P\_3,1.33)} lāvasthāyām eva syādayaḥ . \\
\noindent \textbf{(P\_3,1.33)} sārvadhātuke śyanādayaḥ . \\
\noindent \textbf{(P\_3,1.34.1)} siP utsargaḥ chandasi . \\
\noindent \textbf{(P\_3,1.34.1)} sip utsargaḥ chandasi kartavyaḥ . \\
\noindent \textbf{(P\_3,1.34.1)} sanādyante neṣatvādyarthaḥ . \\
\noindent \textbf{(P\_3,1.34.1)} sanādyante ca kartavyaḥ . \\
\noindent \textbf{(P\_3,1.34.1)} kim prajojanam . \\
\noindent \textbf{(P\_3,1.34.1)} neṣatvādyarthaḥ . \\
\noindent \textbf{(P\_3,1.34.1)} indraḥ naḥ tena neṣatu . \\
\noindent \textbf{(P\_3,1.34.1)} gā vaḥ neṣṭāt . \\
\noindent \textbf{(P\_3,1.34.1)} prakṛtyantaratvāt siddham . \\
\noindent \textbf{(P\_3,1.34.1)} prakṛtyantaratvāt siddham etat . \\
\noindent \textbf{(P\_3,1.34.1)} prkṛtyantaram neṣatiḥ . \\
\noindent \textbf{(P\_3,1.34.1)} neṣatu neṣṭāt iti darśanāt . \\
\noindent \textbf{(P\_3,1.34.1)} neṣatu neṣṭāt iti dṛśyate . \\
\noindent \textbf{(P\_3,1.34.2)} atha kimarthaḥ pakāraḥ . \\
\noindent \textbf{(P\_3,1.34.2)} svarārthaḥ . \\
\noindent \textbf{(P\_3,1.34.2)} anudāttau suppitau iti eṣaḥ svaraḥ yathā syāt . \\
\noindent \textbf{(P\_3,1.34.2)} pitkaraṇānarthakyam ca anackatvāt . \\
\noindent \textbf{(P\_3,1.34.2)} pitkaraṇam ca anarthakam . \\
\noindent \textbf{(P\_3,1.34.2)} kim kāraṇam . \\
\noindent \textbf{(P\_3,1.34.2)} anackatvāt . \\
\noindent \textbf{(P\_3,1.34.2)} anackaḥ ayam . \\
\noindent \textbf{(P\_3,1.34.2)} tatra na arthaḥ svarārthena pakāreṇa anubandhena . \\
\noindent \textbf{(P\_3,1.34.2)} iṭi kṛte sāckaḥ bhaviṣyati . \\
\noindent \textbf{(P\_3,1.34.2)} iṭaḥ anudāttārtham iti cet āgamānudāttatvāt siddham . āgamānudāttatvena iṭaḥ anudāttatvam bhaviṣyati . \\
\noindent \textbf{(P\_3,1.34.2)} evam tarhi sap ayam kartavyaḥ . \\
\noindent \textbf{(P\_3,1.34.2)} kim prayojanam . \\
\noindent \textbf{(P\_3,1.34.2)} yat eva yāsiṣīṣṭhāḥ . \\
\noindent \textbf{(P\_3,1.34.2)} ekājlakṣaṇaḥ iṭpratiṣedhaḥ mā bhūt iti . \\
\noindent \textbf{(P\_3,1.34.2)} kva ayam akāraḥ śrūyate . \\
\noindent \textbf{(P\_3,1.34.2)} na kva cit śrūyate . \\
\noindent \textbf{(P\_3,1.34.2)} lopaḥ asya bhaviṣyati ataḥ lopaḥ ārdhadhātuke iti . \\
\noindent \textbf{(P\_3,1.34.2)} yadi na kva cit śrūyate na arthaḥ svarārthena pakāreṇa anubandhena . \\
\noindent \textbf{(P\_3,1.34.2)} evam api kartavyaḥ eva . \\
\noindent \textbf{(P\_3,1.34.2)} kim prayojanam . \\
\noindent \textbf{(P\_3,1.34.2)} anudāttasya lopaḥ yathā syāt . \\
\noindent \textbf{(P\_3,1.34.2)} udāttasya mā bhūt iti . \\
\noindent \textbf{(P\_3,1.34.2)} kim ca syāt . \\
\noindent \textbf{(P\_3,1.34.2)} udāttanivṛttisvaraḥ prasajyeta . \\
\noindent \textbf{(P\_3,1.34.2)} siP bahulam chandasi ṇit . \\
\noindent \textbf{(P\_3,1.34.2)} sip bahulam chandasi ṇit vaktavyaḥ . \\
\noindent \textbf{(P\_3,1.34.2)} savitā dharmam dāviṣat . \\
\noindent \textbf{(P\_3,1.34.2)} pra ṇaḥ āyūṃṣi tāriṣat . \\
\noindent \textbf{(P\_3,1.35)} kāsgrahaṇe cakāsaḥ upasaṅkhyānam . \\
\noindent \textbf{(P\_3,1.35)} kāsgrahaṇe cakāsaḥ upasaṅkhyānam kartavyam . \\
\noindent \textbf{(P\_3,1.35)} cakāsām cakāra . \\
\noindent \textbf{(P\_3,1.35)} na kartavyam . \\
\noindent \textbf{(P\_3,1.35)} cakāspratayayāt iti vakṣyāmi . \\
\noindent \textbf{(P\_3,1.35)} cakāsgrahaṇe kāsaḥ upasaṅkhyānam kartavyam . \\
\noindent \textbf{(P\_3,1.35)} kāsām cakre . \\
\noindent \textbf{(P\_3,1.35)} sūtram ca bhidyate . \\
\noindent \textbf{(P\_3,1.35)} yathānyāsam eva astu . \\
\noindent \textbf{(P\_3,1.35)} nanu ca uktam kāsgrahaṇe cakāsaḥ upasaṅkhyānam iti . \\
\noindent \textbf{(P\_3,1.35)} na eṣaḥ doṣaḥ . \\
\noindent \textbf{(P\_3,1.35)} cakāsśabde kāsśabdaḥ asti . \\
\noindent \textbf{(P\_3,1.35)} tatra kāspratyayāt iti eva siddham . \\
\noindent \textbf{(P\_3,1.35)} na sidhyati . \\
\noindent \textbf{(P\_3,1.35)} kim kāraṇam . \\
\noindent \textbf{(P\_3,1.35)} arthavataḥ kāsśabdasya grahaṇam . \\
\noindent \textbf{(P\_3,1.35)} na ca cakāsśabde kāsśabdaḥ arthavān . \\
\noindent \textbf{(P\_3,1.35)} evam tarhi kāsi anekācaḥ iti vaktavyam . \\
\noindent \textbf{(P\_3,1.35)} kim prayojanam . \\
\noindent \textbf{(P\_3,1.35)} culumpādyartham . \\
\noindent \textbf{(P\_3,1.35)} culumpām cakāra daridrām cakāra . \\
\noindent \textbf{(P\_3,1.36.1)} gurumataḥ āmvidhāne liṇnimittāt pratiṣedhaḥ . \\
\noindent \textbf{(P\_3,1.36.1)} gurumataḥ āmvidhāne liṇnimittāt pratiṣedhaḥ vaktavyaḥ . \\
\noindent \textbf{(P\_3,1.36.1)} iyeṣa uvoṣa . \\
\noindent \textbf{(P\_3,1.36.1)} guṇe kṛte ijādeḥ ca gurumataḥ anṛcchaḥ iti ām prāpnoti . \\
\noindent \textbf{(P\_3,1.36.1)} gurumadvacanam idānīm kimartham syāt . \\
\noindent \textbf{(P\_3,1.36.1)} gurumadvacanam kimartham iti cet ṇali uttame yajādipratiṣedhāṛtham . \\
\noindent \textbf{(P\_3,1.36.1)} gurumadvacanam kimartham iti cet ṇali uttame yajādīnām mā bhūt iti . \\
\noindent \textbf{(P\_3,1.36.1)} iyaja aham uvapa aham . \\
\noindent \textbf{(P\_3,1.36.1)} upadeśavacanāt siddham . \\
\noindent \textbf{(P\_3,1.36.1)} upadeśe gurumataḥ iti vaktavyam . \\
\noindent \textbf{(P\_3,1.36.1)} yadi upadeśagrahaṇam kriyate uccheḥ ām vaktavyaḥ . \\
\noindent \textbf{(P\_3,1.36.1)} vyucchām cakāra iti . \\
\noindent \textbf{(P\_3,1.36.1)} ṛcchipratiṣedhaḥ jñāpakaḥ uccheḥ āmbhāvasya . \\
\noindent \textbf{(P\_3,1.36.1)} yat ayam anṛcchaḥ iti pratiṣedham śāsti tat jñāpayati ācāryaḥ tugnimittā yasya gurumattā bhavati tasmāt ām iti . \\
\noindent \textbf{(P\_3,1.36.1)} sa tarhi jñāpakārthaḥ ṛcchipratiṣedhaḥ vaktavyaḥ . \\
\noindent \textbf{(P\_3,1.36.1)} nanu ca avaśyam prāptyarthaḥ api vaktavyaḥ . \\
\noindent \textbf{(P\_3,1.36.1)} na arthaḥ prāptyarthena . \\
\noindent \textbf{(P\_3,1.36.1)} ṛcchatyṝtām iti ṛccheḥ liṭi guṇavacanam jñāpakam na ṛccheḥ liṭi ām bhavati iti . \\
\noindent \textbf{(P\_3,1.36.1)} na etat asti jñāpakam . \\
\noindent \textbf{(P\_3,1.36.1)} artyartham etat syāt . \\
\noindent \textbf{(P\_3,1.36.1)} katham punaḥ ṛccheḥ liṭi guṇaḥ ucyamānaḥ artyarthaḥ śakyaḥ vijñātum . \\
\noindent \textbf{(P\_3,1.36.1)} sāmarthyāt . \\
\noindent \textbf{(P\_3,1.36.1)} ṛcchiḥ liṭi na asti iti kṛtvā prakṛtyartham vijñāyate . \\
\noindent \textbf{(P\_3,1.36.1)} tat yathā . \\
\noindent \textbf{(P\_3,1.36.1)} tiṣṭhateḥ it jighrateḥ vā iti caṅi tiṣṭhatijighratī na staḥ iti kṛtvā prakṛtyartham vijñāyate . \\
\noindent \textbf{(P\_3,1.36.1)} kim punaḥ arteḥ guṇavacane prayojanam . \\
\noindent \textbf{(P\_3,1.36.1)} āratuḥ āruḥ etat rūpam yathā syāt . \\
\noindent \textbf{(P\_3,1.36.1)} kim punaḥ kāraṇam na sidhyati . \\
\noindent \textbf{(P\_3,1.36.1)} dvirvacane kṛte savarṇadīrghatve ca yadi tāvat dhātugrahaṇena grahaṇam ṝkārāntānām liṭi guṇaḥ bhavati iti guṇe kṛte raparate aratuḥ aruḥ iti etat rūpam prasajyeta . \\
\noindent \textbf{(P\_3,1.36.1)} atha abhyāsagrahaṇena grahaṇam uḥ attvam raparatvam halādiśeṣaḥ ataḥ ādeḥ iti dīrghatvam ātaḥ lopaḥ iṭi ca iti ākāralopaḥ atuḥ uḥ iti vacanam eva śrūyeta . \\
\noindent \textbf{(P\_3,1.36.1)} guṇa punaḥ sati guṇe kṛte raparatve ca dvirvacanam ataḥ ādeḥ iti dīrghatvam . \\
\noindent \textbf{(P\_3,1.36.1)} tataḥ siddham bhavati yathā āṭatuḥ āṭuḥ iti . \\
\noindent \textbf{(P\_3,1.36.1)} kim punaḥ savarṇadīrghatvam tāvat bhavati na punaḥ uḥ attvam . \\
\noindent \textbf{(P\_3,1.36.1)} paratvāt uḥ attvena bhavitavyam . \\
\noindent \textbf{(P\_3,1.36.1)} antaraṅgatvāt . \\
\noindent \textbf{(P\_3,1.36.1)} antaraṅgam savarṇadīrghatvam . \\
\noindent \textbf{(P\_3,1.36.1)} bahiraṅgam uḥ attvam . \\
\noindent \textbf{(P\_3,1.36.1)} kā antaraṅgatā . \\
\noindent \textbf{(P\_3,1.36.1)} varṇau āśritya savarṇadīrghatvam . \\
\noindent \textbf{(P\_3,1.36.1)} aṅgasya uḥ attvam . \\
\noindent \textbf{(P\_3,1.36.1)} uḥ attvam api antaraṅgam . \\
\noindent \textbf{(P\_3,1.36.1)} katham . \\
\noindent \textbf{(P\_3,1.36.1)} vakṣyati etat . \\
\noindent \textbf{(P\_3,1.36.1)} prāk abhyāsavikārebhyaḥ aṅgādhikāraḥ iti . \\
\noindent \textbf{(P\_3,1.36.1)} ubhayoḥ antaraṅgayoḥ paratvāt uḥ attvam . \\
\noindent \textbf{(P\_3,1.36.1)} uḥ attve kṛte raparatvam halādiśeṣaḥ ataḥ ādeḥ iti dīrghatvam parasya rūpasya yaṇādeśaḥ . \\
\noindent \textbf{(P\_3,1.36.1)} siddham bhavati āratuḥ āruḥ iti . \\
\noindent \textbf{(P\_3,1.36.1)} atha api katham cit arteḥ liṭi guṇena arthaḥ syāt . \\
\noindent \textbf{(P\_3,1.36.1)} evam api na doṣaḥ . \\
\noindent \textbf{(P\_3,1.36.1)} ṛcchatyṝtām iti ṛkāraḥ api nirdiśyate . \\
\noindent \textbf{(P\_3,1.36.1)} katham . \\
\noindent \textbf{(P\_3,1.36.1)} ayam . \\
\noindent \textbf{(P\_3,1.36.1)} ṛcchati ṛ ṛtām ṛcchatyṝtām iti . \\
\noindent \textbf{(P\_3,1.36.1)} iha api tarhi prāpnoti . \\
\noindent \textbf{(P\_3,1.36.1)} cakratuḥ cakruḥ iti . \\
\noindent \textbf{(P\_3,1.36.1)} saṃyogādigrahaṇam niyamārtham bhaviṣyati . \\
\noindent \textbf{(P\_3,1.36.1)} saṃyogādeḥ eva akevalasya na anyasya akevalasya iti . \\
\noindent \textbf{(P\_3,1.36.1)} tat etat antareṇa arteḥ liṭi guṇavacanam rūpam siddham antareṇa ca ṛcchigrahaṇam arteḥ liṭi guṇaḥ siddhaḥ . \\
\noindent \textbf{(P\_3,1.36.1)} saḥ eṣaḥ ananyārthaḥ ṛcchipratiṣedhaḥ vaktavyaḥ uccheḥ vā ām vaktavyaḥ . \\
\noindent \textbf{(P\_3,1.36.1)} ubhayam na vaktavyam . \\
\noindent \textbf{(P\_3,1.36.1)} upadeśagrahaṇam na kariṣyate . \\
\noindent \textbf{(P\_3,1.36.1)} kasmāt na bhavati iyeṣa uvoṣa . \\
\noindent \textbf{(P\_3,1.36.1)} uktam vā .kim uktam . \\
\noindent \textbf{(P\_3,1.36.1)} sannipātalakṣaṇaḥ vidhiḥ animittam tadvighātasya iti . \\
\noindent \textbf{(P\_3,1.36.2)} ūrṇoteḥ ca upasaṅkhyānam . \\
\noindent \textbf{(P\_3,1.36.2)} ūrṇoteḥ ca upasaṅkhyānam kartavyam . \\
\noindent \textbf{(P\_3,1.36.2)} prorṇunāva . \\
\noindent \textbf{(P\_3,1.36.2)} na vaktavyam . \\
\noindent \textbf{(P\_3,1.36.2)} vācyaḥ ūrṇoḥ ṇuvadbhāvaḥ . \\
\noindent \textbf{(P\_3,1.36.2)} yaṅprasiddhiḥ prayojanam . \\
\noindent \textbf{(P\_3,1.36.2)} āmaḥ ca pratiṣedhārtham . \\
\noindent \textbf{(P\_3,1.36.2)} ekācaḥ ca iḍupagrahāt . atha vā ukāraḥ api atra nirdiśyate . \\
\noindent \textbf{(P\_3,1.36.2)} katham . \\
\noindent \textbf{(P\_3,1.36.2)} avibhaktikaḥ nirdeśaḥ . \\
\noindent \textbf{(P\_3,1.36.2)} anṛccha u anṛccho dayāyāsaḥ ca iti . \\
\noindent \textbf{(P\_3,1.38)} videḥ ām kit . \\
\noindent \textbf{(P\_3,1.38)} videḥ ām kit vaktavyaḥ . \\
\noindent \textbf{(P\_3,1.38)} vidām cakāra . \\
\noindent \textbf{(P\_3,1.38)} na vaktavyaḥ . \\
\noindent \textbf{(P\_3,1.38)} vidiḥ akārāntaḥ . \\
\noindent \textbf{(P\_3,1.38)} yadi akārāntaḥ vetti iti guṇaḥ na sidhyati . \\
\noindent \textbf{(P\_3,1.38)} liṭsanniyogena . \\
\noindent \textbf{(P\_3,1.38)} evam api viveda iti na sidhyati . \\
\noindent \textbf{(P\_3,1.38)} evam tarhi āmsanniyogena . \\
\noindent \textbf{(P\_3,1.38)} bhāradvājīyāḥ paṭhanti . \\
\noindent \textbf{(P\_3,1.38)} videḥ ām kit nipātanāt vā aguṇatvam iti . \\
\noindent \textbf{(P\_3,1.39)} śluvadatideśe kim prayojanam . \\
\noindent \textbf{(P\_3,1.39)} śluvadatideśe prayojanam dvitvettve . \\
\noindent \textbf{(P\_3,1.39)} bibharām cakāra . \\
\noindent \textbf{(P\_3,1.40)} kimartham idam ucyate . \\
\noindent \textbf{(P\_3,1.40)} anuprayogaḥ yathā syāt . \\
\noindent \textbf{(P\_3,1.40)} na etat asti prayojanam . \\
\noindent \textbf{(P\_3,1.40)} āmantam avyaktapadārthakam . \\
\noindent \textbf{(P\_3,1.40)} tena aparisamāptaḥ arthaḥ iti kṛtvā anuprayogaḥ bhaviṣyati . \\
\noindent \textbf{(P\_3,1.40)} ataḥ uttaram paṭhati . \\
\noindent \textbf{(P\_3,1.40)} kṛñaḥ anuprayogavacanam astibhūpratiṣedhārtham . \\
\noindent \textbf{(P\_3,1.40)} kṛñaḥ anuprayogavacanam kriyate astibhūpratiṣedhārtham . \\
\noindent \textbf{(P\_3,1.40)} astibhuvoḥ anuprayogaḥ mā bhūt iti . \\
\noindent \textbf{(P\_3,1.40)} ātmanepadavidhyartham ca . \\
\noindent \textbf{(P\_3,1.40)} ātmanepadavidhyartham ca kṛñaḥ anuprayogavacanam kriyate . \\
\noindent \textbf{(P\_3,1.40)} ātmanepadam yathā syāt . \\
\noindent \textbf{(P\_3,1.40)} ucyamāne api etasmin avaśyam ātmanepadārthaḥ yatnaḥ kartavyaḥ . \\
\noindent \textbf{(P\_3,1.40)} astibhūpratiṣedhārthena ca api na arthaḥ . \\
\noindent \textbf{(P\_3,1.40)} iṣṭaḥ sarvānuprayogaḥ . \\
\noindent \textbf{(P\_3,1.40)} sarveṣām eva kṛbhvastīnām anuprayogaḥ iṣyate . \\
\noindent \textbf{(P\_3,1.40)} kim iṣyate eva āhosvit prāpnoti api . \\
\noindent \textbf{(P\_3,1.40)} iṣyate ca prāpnoti ca . \\
\noindent \textbf{(P\_3,1.40)} katham . \\
\noindent \textbf{(P\_3,1.40)} kṛñ iti na etat dhātugrahaṇam . \\
\noindent \textbf{(P\_3,1.40)} kim tarhi . \\
\noindent \textbf{(P\_3,1.40)} pratyāhāragrahaṇam . \\
\noindent \textbf{(P\_3,1.40)} kva sanniviṣṭānām pratyāhāraḥ . \\
\noindent \textbf{(P\_3,1.40)} kṛbhvastiyoge iti ataḥ prabhṛti ā kṛñaḥ ñakārāt . \\
\noindent \textbf{(P\_3,1.40)} sarvānuprayogaḥ iti cet aśiṣyam arthābhāvāt . \\
\noindent \textbf{(P\_3,1.40)} sarvānuprayogaḥ iti cet aśiṣyam kṛñaḥ anuprayogavacanam . \\
\noindent \textbf{(P\_3,1.40)} kim kāraṇam . \\
\noindent \textbf{(P\_3,1.40)} arthābhāvāt . \\
\noindent \textbf{(P\_3,1.40)} āmantam avyaktapadārthakam . \\
\noindent \textbf{(P\_3,1.40)} tena aparisamāptaḥ arthaḥ iti kṛtvā anuprayogaḥ bhaviṣyati . \\
\noindent \textbf{(P\_3,1.40)} idam tarhi prayojanam . \\
\noindent \textbf{(P\_3,1.40)} kṛbhvastīnām eva anuprayogaḥ yathā syāt pacādīnām mā bhūt iti . \\
\noindent \textbf{(P\_3,1.40)} etat api na asti prayojanam . \\
\noindent \textbf{(P\_3,1.40)} arthābhāvāt ca anyasya . \\
\noindent \textbf{(P\_3,1.40)} arthābhāvāt ca anyasya siddham . \\
\noindent \textbf{(P\_3,1.40)} kṛbhvastayaḥ kriyāsāmānyavācinaḥ . \\
\noindent \textbf{(P\_3,1.40)} kriyāviśeṣavācinaḥ pacādayaḥ . \\
\noindent \textbf{(P\_3,1.40)} na ca sāmānyavācinoḥ eva viśeṣavācinoḥ eva va prayogaḥ bhavati . \\
\noindent \textbf{(P\_3,1.40)} tatra viśeṣavācinaḥ utpattiḥ . \\
\noindent \textbf{(P\_3,1.40)} sāmānyavācinaḥ anuprayokṣyante . \\
\noindent \textbf{(P\_3,1.40)} liṭparārtham vā . \\
\noindent \textbf{(P\_3,1.40)} liṭparārtham tarhi kṛñaḥ anuprayogavacanam kriyate . \\
\noindent \textbf{(P\_3,1.40)} liṭparasya eva anuprayogaḥ yathā syāt . \\
\noindent \textbf{(P\_3,1.40)} anyaparasya mā bhūt iti . \\
\noindent \textbf{(P\_3,1.40)} kimparasya punaḥ prāpnoti . \\
\noindent \textbf{(P\_3,1.40)} laṭparasya . \\
\noindent \textbf{(P\_3,1.40)} na laṭparasya anuprayogeṇa bhūtakālaḥ viśeṣitaḥ syāt . \\
\noindent \textbf{(P\_3,1.40)} niṣṭhāparasya tarhi . \\
\noindent \textbf{(P\_3,1.40)} naniṣṭhāparasya anuprayogeṇa puruṣopagrahau viśeiṣitau syātām . \\
\noindent \textbf{(P\_3,1.40)} luṅparasya tarhi . \\
\noindent \textbf{(P\_3,1.40)} na luṅparasya anuprayogeṇa anadyatanaḥ bhūtakālaḥ viśeṣitaḥ syāt . \\
\noindent \textbf{(P\_3,1.40)} laṅparasya tarhi . \\
\noindent \textbf{(P\_3,1.40)} na laṅparasya anuprayogeṇa anadyatanaḥ parokṣaḥ kālaḥ viśeṣitaḥ syāt . \\
\noindent \textbf{(P\_3,1.40)} ayam tarhi bhūte parokṣe anadyatane laṅ vidhīyate . \\
\noindent \textbf{(P\_3,1.40)} haśaśvatoḥ laṅ ca iti . \\
\noindent \textbf{(P\_3,1.40)} tatparasya mā bhūt iti . \\
\noindent \textbf{(P\_3,1.40)} atat api na asti prayojanam . \\
\noindent \textbf{(P\_3,1.40)} ekasyāḥ ākṛteḥ caritaḥ prayogaḥ dvitīyasyāḥ tṛtīyasyāḥ ca na bhavati . \\
\noindent \textbf{(P\_3,1.40)} tat yathā goṣu svāmi aśveṣu ca iti . \\
\noindent \textbf{(P\_3,1.40)} na ca bhavati goṣu ca aśvānām ca svāmī iti . \\
\noindent \textbf{(P\_3,1.40)} arthasamāpteḥ vā anuprayogaḥ na syāt . \\
\noindent \textbf{(P\_3,1.40)} arthasamāpteḥ tarhi anuprayogaḥ na syāt . \\
\noindent \textbf{(P\_3,1.40)} āmantena parisamāptaḥ arthaḥ iti kṛtvā anuprayogaḥ na syāt . \\
\noindent \textbf{(P\_3,1.40)} etat api na asti prayojanam . \\
\noindent \textbf{(P\_3,1.40)} idānīm eva uktam āmantam avyaktapadārthakam . \\
\noindent \textbf{(P\_3,1.40)} tena aparisamāptaḥ arthaḥ iti kṛtvā anuprayogaḥ bhaviṣyati iti . \\
\noindent \textbf{(P\_3,1.40)} viparyāsanivṛttyartham vā . \\
\noindent \textbf{(P\_3,1.40)} viparyāsanivṛttyartham tarhi kṛñaḥ anuprayogavacanam kriyate . \\
\noindent \textbf{(P\_3,1.40)} īhām cakre . \\
\noindent \textbf{(P\_3,1.40)} cakre īhām iti mā bhūt . \\
\noindent \textbf{(P\_3,1.40)} vyavahitnivṛttyartham ca . \\
\noindent \textbf{(P\_3,1.40)} vyavahitnivṛttyartham ca kṛñaḥ anuprayogavacanam kriyate . \\
\noindent \textbf{(P\_3,1.40)} anv eva ca anuprayogaḥ yathā syāt . \\
\noindent \textbf{(P\_3,1.40)} īhām cakre . \\
\noindent \textbf{(P\_3,1.40)} vyavahitasya mā bhūt . \\
\noindent \textbf{(P\_3,1.40)} īhām devadattaḥ cakre iti . \\
\noindent \textbf{(P\_3,1.43)} kva ayam cliḥ śrūyate . \\
\noindent \textbf{(P\_3,1.43)} na kva cit śrūyate . \\
\noindent \textbf{(P\_3,1.43)} sijādayaḥ ādeśāḥ ucyante . \\
\noindent \textbf{(P\_3,1.43)} yad na kva cit śrūyate kimarthaḥ tarhi cluḥ utsargaḥ kriyate . \\
\noindent \textbf{(P\_3,1.43)} na sic utsargaḥ eva kartavyaḥ . \\
\noindent \textbf{(P\_3,1.43)} tasya ksādayaḥ apavādāḥ bhaviṣyanti . \\
\noindent \textbf{(P\_3,1.43)} ata uttaram paṭhati . \\
\noindent \textbf{(P\_3,1.43)} clyutsargaḥ sāmānyagrahaṇārthaḥ . \\
\noindent \textbf{(P\_3,1.43)} cliḥ utsargaḥ kriyate sāmānyagrahaṇārthaḥ . \\
\noindent \textbf{(P\_3,1.43)} kva sāmānyagrahaṇārthena arthaḥ . \\
\noindent \textbf{(P\_3,1.43)} mantra ghasahvaraṇaśavṛdahādvṛckṛgamijanibhyaḥ leḥ iti . \\
\noindent \textbf{(P\_3,1.43)} tatra avarataḥ trayāṇām grahaṇam kartavyam syāt . \\
\noindent \textbf{(P\_3,1.43)} caṅaṅoḥ sicaḥ ca . \\
\noindent \textbf{(P\_3,1.43)} ksavidhāne ca aniḍvacane clisampratyayārthaḥ . ksavidhāne ca aniḍvacane clisampratyayārthaḥ cliḥ utsargaḥ kriyate . \\
\noindent \textbf{(P\_3,1.43)} cleḥ aniṭaḥ ksaḥ siddhaḥ bhavati . \\
\noindent \textbf{(P\_3,1.43)} ghasḷbhāve ca . \\
\noindent \textbf{(P\_3,1.43)} ghasḷbhāve ca clav eva kṛte lṛditaḥ iti aṅ siddhaḥ bhavati . \\
\noindent \textbf{(P\_3,1.43)} atha citkaraṇam kimartham . \\
\noindent \textbf{(P\_3,1.43)} cleḥ citkaraṇam viśeṣaṇāṛtham . \\
\noindent \textbf{(P\_3,1.43)} cleḥ citkaraṇam kriyate viśeṣaṇārtham . \\
\noindent \textbf{(P\_3,1.43)} kva viśeṣaṇārthena arthaḥ . \\
\noindent \textbf{(P\_3,1.43)} cleḥ sic iti . \\
\noindent \textbf{(P\_3,1.43)} leḥ sic iti ucyamāne liṅliṭoḥ api prasajyeta . \\
\noindent \textbf{(P\_3,1.43)} na etat asti prayojanam . \\
\noindent \textbf{(P\_3,1.43)} luṅi iti ucyate . \\
\noindent \textbf{(P\_3,1.43)} na ca luṅi liṅliṭau bhavataḥ . \\
\noindent \textbf{(P\_3,1.43)} atha iditkaraṇam kimartham . \\
\noindent \textbf{(P\_3,1.43)} iditkaraṇam sāmānyagrahaṇārtham . iditkaraṇam kriyate ca sāmānyagrahaṇārtham . \\
\noindent \textbf{(P\_3,1.43)} kva sāmānyagrahaṇārthena arthaḥ . \\
\noindent \textbf{(P\_3,1.43)} mantre ghasahvaraṇaśavṛdahādvṛckṛgamijanibhyaḥ leḥ iti āmaḥ iti ca . \\
\noindent \textbf{(P\_3,1.43)} ikāre ca idānīm sāmānyagrahaṇārthe kriyamāṇe avaśyam sāmānyagrahaṇāvighātārthaḥ cakāraḥ kartavyaḥ . \\
\noindent \textbf{(P\_3,1.43)} kva sāmānyagrahaṇāvighātārthena arthaḥ cakāreṇa . \\
\noindent \textbf{(P\_3,1.43)} atra eva . \\
\noindent \textbf{(P\_3,1.43)} yat tāvat ucyate clyutsargaḥ sāmānyagrahaṇārthaḥ iti . \\
\noindent \textbf{(P\_3,1.43)} kriyamāṇe api vai clyutsarge tāni eva trīṇi grahaṇāni bhavanti . \\
\noindent \textbf{(P\_3,1.43)} clu luṅi cleḥ sic leḥ iti . \\
\noindent \textbf{(P\_3,1.43)} yat etat leḥ iti tat parārtham bhaviṣyati . \\
\noindent \textbf{(P\_3,1.43)} katham . \\
\noindent \textbf{(P\_3,1.43)} yat etat gātisthāghupābhūbhyaḥ sicaḥ parasmaipadeṣu iti atra sicaḥ grahaṇam etat leḥ iti vakṣyāmi . \\
\noindent \textbf{(P\_3,1.43)} yadi leḥ iti ucyate dheṭaḥ cātuḥśabdyam prāpnoti . \\
\noindent \textbf{(P\_3,1.43)} adadhat adhāt adhāsīt . \\
\noindent \textbf{(P\_3,1.43)} adadhāt iti api prāpnoti . \\
\noindent \textbf{(P\_3,1.43)} na caṅaḥ luki dvirvacanena bhavitavyam . \\
\noindent \textbf{(P\_3,1.43)} kim kāraṇam . \\
\noindent \textbf{(P\_3,1.43)} caṅi iti ucyate . \\
\noindent \textbf{(P\_3,1.43)} na ca atra caṅam paśyāmaḥ . \\
\noindent \textbf{(P\_3,1.43)} pratyayalakṣaṇena . \\
\noindent \textbf{(P\_3,1.43)} na lumatā tasmin iti pratyayalakṣaṇapratiṣedhaḥ . \\
\noindent \textbf{(P\_3,1.43)} bahuvacane tarhi cātuḥśabdyam prāpnoti . \\
\noindent \textbf{(P\_3,1.43)} adadhan adhuḥ adhāsiṣuḥ . \\
\noindent \textbf{(P\_3,1.43)} adhān iti api prāpnoti . \\
\noindent \textbf{(P\_3,1.43)} na eṣaḥ doṣaḥ . \\
\noindent \textbf{(P\_3,1.43)} ātaḥ iti jusbhāvaḥ bhaviṣyati . \\
\noindent \textbf{(P\_3,1.43)} na sidhyati . \\
\noindent \textbf{(P\_3,1.43)} sijgrahaṇam tatra anuvartate . \\
\noindent \textbf{(P\_3,1.43)} sijgrahaṇam nivartiṣyate . \\
\noindent \textbf{(P\_3,1.43)} yadi nivartate abhūvan iti pratyayalakṣaṇena jusbhāvaḥ prāpnoti . \\
\noindent \textbf{(P\_3,1.43)} evam tarhi luk sijapavādaḥ vijñāsyate . \\
\noindent \textbf{(P\_3,1.43)} yadi luk sijapavādaḥ vijñāyate mā hi dātām mā hi dhātām iti atra ādiḥ sicaḥ anyatarasyām iti eṣaḥ svaraḥ na prāpnoti . \\
\noindent \textbf{(P\_3,1.43)} tasmāt na etat śakyam vaktum luk sijapavādaḥ iti . \\
\noindent \textbf{(P\_3,1.43)} na cet ucyate abhūvan iti pratyayalakṣaṇena jusbhāvaḥ prāpnoti . \\
\noindent \textbf{(P\_3,1.43)} tasmāt ātaḥ iti atra sijgrahaṇam anuvartyam . \\
\noindent \textbf{(P\_3,1.43)} tasmin ca anuvartamāne dheṭaḥ cātuḥśabdyam prāpnoti . \\
\noindent \textbf{(P\_3,1.43)} tasmāt gātisthāghupābhūbhyaḥ sicaḥ parasmaipadeṣu iti atra sicaḥ grahaṇam kartavyam . \\
\noindent \textbf{(P\_3,1.43)} tasmin ca kriyamāṇe tāni eva trīṇi grahaṇāni bhavanti cli luṅi cleḥ sic leḥ iti . \\
\noindent \textbf{(P\_3,1.43)} yat api ucyate ksavidhāne ca aniḍvacane clisampratyayārthaḥ iti . \\
\noindent \textbf{(P\_3,1.43)} dhātum eva atra aniṭvena viśeṣayiṣyāmaḥ . \\
\noindent \textbf{(P\_3,1.43)} dhātoḥ aniṭaḥ iti . \\
\noindent \textbf{(P\_3,1.43)} katham punaḥ dhātuḥ nāma aniṭ syāt . \\
\noindent \textbf{(P\_3,1.43)} dhātuḥ eva aniṭ . \\
\noindent \textbf{(P\_3,1.43)} katham . \\
\noindent \textbf{(P\_3,1.43)} animittam vā iṭaḥ aniṭaḥ na vā tasmāt iṭ asti saḥ ayam aniṭ iti . \\
\noindent \textbf{(P\_3,1.43)} atha dhātau viśeṣyamāṇe kva yaḥ aniṭ iti viśeṣayiṣyasi . \\
\noindent \textbf{(P\_3,1.43)} kim ca ataḥ . \\
\noindent \textbf{(P\_3,1.43)} yadi vijñāyate niṣṭhāyām aniṭaḥ iti bhūyiṣṭhebhyaḥ prāpnoti . \\
\noindent \textbf{(P\_3,1.43)} bhūyiṣṭhāḥ hi śalantāḥ igupadhāḥ niṣṭhāyām aniṭaḥ . \\
\noindent \textbf{(P\_3,1.43)} atha vijñāyate liṭi yaḥ aniṭ iti na kutaḥ cit prāpnoti . \\
\noindent \textbf{(P\_3,1.43)} sarve his śalantāḥ igupadhāḥ liṭi seṭaḥ . \\
\noindent \textbf{(P\_3,1.43)} kim punaḥ kāraṇam dhātau viśeṣyamāṇe etayoḥ viśeṣayoḥ viśeṣayiṣyate . \\
\noindent \textbf{(P\_3,1.43)} na punaḥ atra sāmānyena iṭaḥ vidhipratiṣedhau . \\
\noindent \textbf{(P\_3,1.43)} kva sāmanyena . \\
\noindent \textbf{(P\_3,1.43)} valādau ārdhadhātuke . \\
\noindent \textbf{(P\_3,1.43)} yat api ucyate ghasḷbhāve ca iti . \\
\noindent \textbf{(P\_3,1.43)} ārdhadhātukīyāḥ sāmānyena bhavanti anavasthiteṣu pratyayeṣu . \\
\noindent \textbf{(P\_3,1.43)} tatra ārdhadhātukasāmānye ghasḷbhāve kṛte lṛditaḥ iti aṅ bhaviṣyati . \\
\noindent \textbf{(P\_3,1.44.1)} kimarthaḥ cakāraḥ . \\
\noindent \textbf{(P\_3,1.44.1)} viśeṣaṇārthaḥ . \\
\noindent \textbf{(P\_3,1.44.1)} kva viśeṣaṇārthena arthaḥ . \\
\noindent \textbf{(P\_3,1.44.1)} sici vṛddhiḥ parasmaipadeṣu iti . \\
\noindent \textbf{(P\_3,1.44.1)} sau vṛddhiḥ iti ucyamāne agniḥ vāyuḥ iti atra api prasajyeta . \\
\noindent \textbf{(P\_3,1.44.1)} na etat asti prayojanam . \\
\noindent \textbf{(P\_3,1.44.1)} parasmaipadeṣu iti ucyate . \\
\noindent \textbf{(P\_3,1.44.1)} na ca atra parasmaipadam paśyāmaḥ . \\
\noindent \textbf{(P\_3,1.44.1)} svarārthaḥ tarhi . \\
\noindent \textbf{(P\_3,1.44.1)} citaḥ antaḥ udāttaḥ bhavati iti antodāttatvam yathā syāt . \\
\noindent \textbf{(P\_3,1.44.1)} etat api na asti prayojanam . \\
\noindent \textbf{(P\_3,1.44.1)} anackaḥ ayam . \\
\noindent \textbf{(P\_3,1.44.1)} tatra na arthaḥ svarārthena cakāreṇa anubandhena . \\
\noindent \textbf{(P\_3,1.44.1)} iṭi kṛte sāckaḥ bhaviṣyati . \\
\noindent \textbf{(P\_3,1.44.1)} tatra pratyayādyudāttatvena iṭaḥ udāttatvam bhaviṣyati . \\
\noindent \textbf{(P\_3,1.44.1)} na sidhyati . \\
\noindent \textbf{(P\_3,1.44.1)} āgamāḥ anudāttāḥ bhavanti iti anudāttatvam prāpnoti . \\
\noindent \textbf{(P\_3,1.44.1)} ataḥ uttaram paṭhati . \\
\noindent \textbf{(P\_3,1.44.1)} sicaḥ citkaraṇānarthakyam sthānivatvāt . \\
\noindent \textbf{(P\_3,1.44.1)} sicaḥ citkaraṇam narthayam . \\
\noindent \textbf{(P\_3,1.44.1)} kim kāraṇam . \\
\noindent \textbf{(P\_3,1.44.1)} sthānivatvāt . \\
\noindent \textbf{(P\_3,1.44.1)} sthānivadbhāvāt cit bhaviṣyati . \\
\noindent \textbf{(P\_3,1.44.1)} arthavat tu citkaraṇasāmarthyāt hi iṭaḥ udāttatvam . \\
\noindent \textbf{(P\_3,1.44.1)} arthavat tu citkaraṇam . \\
\noindent \textbf{(P\_3,1.44.1)} kaḥ arthaḥ . \\
\noindent \textbf{(P\_3,1.44.1)} citkaraṇasāmarthyāt hi iṭaḥ udāttatvam bhaviṣyati . \\
\noindent \textbf{(P\_3,1.44.1)} na aprāpte pratyayasvare āgamānudāttatvam ārabhyate . \\
\noindent \textbf{(P\_3,1.44.1)} tat yathā eva pratyayasvaram bādhate evam sthānivadbhāvāt api yā prāptiḥ tām api bādheta . \\
\noindent \textbf{(P\_3,1.44.1)} tasmāt citkaraṇam . \\
\noindent \textbf{(P\_3,1.44.1)} tasmāt cakāraḥ kartavyaḥ . \\
\noindent \textbf{(P\_3,1.44.1)} atha iditkaraṇam kimartham . \\
\noindent \textbf{(P\_3,1.44.1)} iditkaraṇam nakāralopābhāvārtham . \\
\noindent \textbf{(P\_3,1.44.1)} iditkaraṇam kriyate nakāralopaḥ mā bhūt iti . \\
\noindent \textbf{(P\_3,1.44.1)} amaṃsta amaṃsthāḥ . \\
\noindent \textbf{(P\_3,1.44.1)} aniditām halaḥ upadhāyāḥ kṅiti iti . \\
\noindent \textbf{(P\_3,1.44.1)} na vā hanteḥ sicaḥ kitkaraṇam nakāralopābhāvasya . \\
\noindent \textbf{(P\_3,1.44.1)} na vā etat prayojanam asti . \\
\noindent \textbf{(P\_3,1.44.1)} kim kāraṇam . \\
\noindent \textbf{(P\_3,1.44.1)} yat ayam hanaḥ sic iti hanteḥ sicaḥ kittvam śāsti tat jñāpayati ācāryaḥ na sijantasya nakārlopaḥ bhavati iti . \\
\noindent \textbf{(P\_3,1.44.1)} na etat asti jñāpakam . \\
\noindent \textbf{(P\_3,1.44.1)} asti hi anyat etasya vacane prayojanam . \\
\noindent \textbf{(P\_3,1.44.1)} kim . \\
\noindent \textbf{(P\_3,1.44.1)} sici eva nalopaḥ yathā syāt . \\
\noindent \textbf{(P\_3,1.44.1)} parasmin nimitte mā bhūt iti . \\
\noindent \textbf{(P\_3,1.44.1)} kaḥ punaḥ atra viśeṣaḥ sici vā nalope sati parasmin vā nimitte . \\
\noindent \textbf{(P\_3,1.44.1)} ayam asti viśeṣaḥ . \\
\noindent \textbf{(P\_3,1.44.1)} sici nalope sati nalopasya asiddhatvāt akāralopaḥ na bhavati . \\
\noindent \textbf{(P\_3,1.44.1)} parasmin punaḥ nimitte nalope sati akāralopaḥ prāpnoti . \\
\noindent \textbf{(P\_3,1.44.1)} samānāśrayam asiddham vyāśrayam ca idam . \\
\noindent \textbf{(P\_3,1.44.1)} nanu ca parasmin api nimitte nalope sati akāralopaḥ na bhaviṣyati . \\
\noindent \textbf{(P\_3,1.44.1)} katham . \\
\noindent \textbf{(P\_3,1.44.1)} asiddham bahiraṅgalakṣaṇam antaraṅgalakṣaṇe iti ṭat etat hanteḥ sicaḥ kitkaraṇam jñāpakam eva na sijantasya nalopaḥ bhavati iti . \\
\noindent \textbf{(P\_3,1.44.1)} idittvāt vā sthānivattvāt . \\
\noindent \textbf{(P\_3,1.44.1)} atha api anena iditā arthaḥ syāt . \\
\noindent \textbf{(P\_3,1.44.1)} ayam ādeśaḥ sthānivadbhāvāt idit bhaviṣyati. \\
\noindent \textbf{(P\_3,1.44.2)} spṛśamṛśakṛṣatṛpadṛpaḥ sic vā . \\
\noindent \textbf{(P\_3,1.44.2)} spṛśmṛśakṛṣatṛpadṛpaḥ sic vā iti vaktavyam . \\
\noindent \textbf{(P\_3,1.44.2)} spṛśa . \\
\noindent \textbf{(P\_3,1.44.2)} aspṛkṣat asprākṣīt . \\
\noindent \textbf{(P\_3,1.44.2)} spṛśa . \\
\noindent \textbf{(P\_3,1.44.2)} mṛśa . \\
\noindent \textbf{(P\_3,1.44.2)} amṛkṣat amrākṣīt . \\
\noindent \textbf{(P\_3,1.44.2)} mṛśa . \\
\noindent \textbf{(P\_3,1.44.2)} kṛṣa . \\
\noindent \textbf{(P\_3,1.44.2)} akṛkṣat akrākṣīt . \\
\noindent \textbf{(P\_3,1.44.2)} kṛṣa . \\
\noindent \textbf{(P\_3,1.44.2)} tṛpa . \\
\noindent \textbf{(P\_3,1.44.2)} atṛpat atrāpsīt . \\
\noindent \textbf{(P\_3,1.44.2)} tṛpa . \\
\noindent \textbf{(P\_3,1.44.2)} dṛpa . \\
\noindent \textbf{(P\_3,1.44.2)} adṛpat adrapsīt . \\
\noindent \textbf{(P\_3,1.44.2)} kim prayojanam . \\
\noindent \textbf{(P\_3,1.44.2)} sic yathā syāta . \\
\noindent \textbf{(P\_3,1.44.2)} atha ksaḥ siddhaḥ . \\
\noindent \textbf{(P\_3,1.44.2)} siddhaḥ śalaḥ igupadhāt aniṭaḥ iti . \\
\noindent \textbf{(P\_3,1.44.2)} sic api siddhaḥ . \\
\noindent \textbf{(P\_3,1.44.2)} katham . \\
\noindent \textbf{(P\_3,1.44.2)} cleḥ citkaraṇam pratyākhyāyate . \\
\noindent \textbf{(P\_3,1.44.2)} tatra clau eva jhallakṣaṇe amāgame kṛte vihatanimittatvāt ksaḥ na bhaviṣyati . \\
\noindent \textbf{(P\_3,1.44.2)} yadi evam antyasaya sijādayaḥ prāpnuvanti . \\
\noindent \textbf{(P\_3,1.44.2)} siddham tu sicaḥ yāditvāt . \\
\noindent \textbf{(P\_3,1.44.2)} siddham etat . \\
\noindent \textbf{(P\_3,1.44.2)} katham . \\
\noindent \textbf{(P\_3,1.44.2)} yādiḥ sic kariṣyate . \\
\noindent \textbf{(P\_3,1.44.2)} saḥ anekālśit sarvasya iti sarvādeśaḥ bhaviṣyati . \\
\noindent \textbf{(P\_3,1.44.2)} kim na śrūyate yakāraḥ . \\
\noindent \textbf{(P\_3,1.44.2)} luptanirdiṣṭaḥ yakāraḥ . \\
\noindent \textbf{(P\_3,1.44.2)} caṅaṅoḥ katham . \\
\noindent \textbf{(P\_3,1.44.2)} caṅaṅoḥ praśliṣṭanirdeśāt siddham . \\
\noindent \textbf{(P\_3,1.44.2)} caṅaṅoḥ api praśliṣṭanirdeśaḥ ayam : ca aṅ caṅ a aṅ aṅ . \\
\noindent \textbf{(P\_3,1.44.2)} saḥ anekālśit sarvasya iti sarvādeśaḥ bhaviṣyati . \\
\noindent \textbf{(P\_3,1.44.2)} ciṇaḥ katham . \\
\noindent \textbf{(P\_3,1.44.2)} ciṇaḥ anittvāt siddham . \\
\noindent \textbf{(P\_3,1.44.2)} ciṇaḥ anittvāt siddham . \\
\noindent \textbf{(P\_3,1.44.2)} kim idam anittvāt . \\
\noindent \textbf{(P\_3,1.44.2)} antyasya ayam sthāne bhavan na pratyayaḥ syāt . \\
\noindent \textbf{(P\_3,1.44.2)} asatyāyām pratyayasañjñayām itsañjñā na. asatyām itsañjñāyām lopaḥ na . \\
\noindent \textbf{(P\_3,1.44.2)} asati lope anekāl . \\
\noindent \textbf{(P\_3,1.44.2)} yadā anekāl tadā sarvādeśaḥ . \\
\noindent \textbf{(P\_3,1.44.2)} yadā sarvādeśaḥ tadā prayayaḥ . \\
\noindent \textbf{(P\_3,1.44.2)} yadā pratyayaḥ tadā itsañjñā . \\
\noindent \textbf{(P\_3,1.44.2)} yadā itsañjñā tadā lopaḥ . \\
\noindent \textbf{(P\_3,1.44.2)} evam ca tatra vārttikakārasya nirṇayaḥ saprayojanam citkaraṇam iti . \\
\noindent \textbf{(P\_3,1.44.2)} api ca traiśabdyam na prakalpate . \\
\noindent \textbf{(P\_3,1.44.2)} aspṛkṣat asprākṣīt aspārkṣīt iti na sidhyati . \\
\noindent \textbf{(P\_3,1.44.2)} sici punaḥ sati vibhāṣā sic . \\
\noindent \textbf{(P\_3,1.44.2)} sici api jhallakṣaṇaḥ amāgamaḥ vibhāṣā . \\
\noindent \textbf{(P\_3,1.44.2)} yasya khalu api amā nimittam na vihanyate saḥ syāt eva . \\
\noindent \textbf{(P\_3,1.44.2)} tasmāt suṣṭhu ucyate spṛśmṛśakṛṣatṛpadṛpaḥ sic vā iti \\
\noindent \textbf{(P\_3,1.45)} ksavidhāne igupadhābhāvaḥ cleḥ guṇanimittatvāt . \\
\noindent \textbf{(P\_3,1.45)} ksavidhāne igupadhābhāvaḥ . \\
\noindent \textbf{(P\_3,1.45)} kim kāraṇam . \\
\noindent \textbf{(P\_3,1.45)} cleḥ guṇanimittatvāt . \\
\noindent \textbf{(P\_3,1.45)} cliḥ guṇanimittam . \\
\noindent \textbf{(P\_3,1.45)} tatra clau eva guṇe kṛte igupadhāt iti ksaḥ na prāpnoti . \\
\noindent \textbf{(P\_3,1.45)} na vā ksasya anavakāśatvāt apavādaḥ guṇasya . \\
\noindent \textbf{(P\_3,1.45)} na vā eṣaḥ doṣaḥ . \\
\noindent \textbf{(P\_3,1.45)} kim kāraṇam . \\
\noindent \textbf{(P\_3,1.45)} ksasya anavakāśatvāt . \\
\noindent \textbf{(P\_3,1.45)} anavakāśaḥ ksaḥ guṇam bādhiṣyate . \\
\noindent \textbf{(P\_3,1.45)} aniḍvacanam aviśeṣaṇam cleḥ nityādiṣṭatvāt . \\
\noindent \textbf{(P\_3,1.45)} aniḍvacanam aviśeṣaṇam . \\
\noindent \textbf{(P\_3,1.45)} kim kāraṇam . \\
\noindent \textbf{(P\_3,1.45)} cleḥ nityādiṣṭatvāt . \\
\noindent \textbf{(P\_3,1.45)} nityādiṣṭaḥ cliḥ na kva cit śrūyate . \\
\noindent \textbf{(P\_3,1.45)} tatra cleḥ aniṭaḥ iti ksaḥ na prāpnoti . \\
\noindent \textbf{(P\_3,1.45)} na vā ksasya sijapavādatvāt tasya ca aniḍāśrayatvāt aniṭi prasiddhe ksaviddhiḥ . \\
\noindent \textbf{(P\_3,1.45)} na vā eṣaḥ doṣaḥ . \\
\noindent \textbf{(P\_3,1.45)} kim kāraṇam . \\
\noindent \textbf{(P\_3,1.45)} ksasya sijapavādatvāt . \\
\noindent \textbf{(P\_3,1.45)} sijapavādaḥ ksaḥ . \\
\noindent \textbf{(P\_3,1.45)} saḥ ca aniḍāśrayaḥ . \\
\noindent \textbf{(P\_3,1.45)} na ca apavādaviṣaye upasargaḥ abhiniviśate . \\
\noindent \textbf{(P\_3,1.45)} pūrvam hi apavādāḥ abhiniviśante paścāt utsargāḥ . \\
\noindent \textbf{(P\_3,1.45)} prakalpya vā apavādaviṣayam utsargaḥ abhiniviśate . \\
\noindent \textbf{(P\_3,1.45)} tat na tāvat atra kadā cit sic bhavati . \\
\noindent \textbf{(P\_3,1.45)} apavādam ksam pratīkṣate . \\
\noindent \textbf{(P\_3,1.45)} ksasya sijapavādatvāt tasya ca aniḍāśrayatvāt aniṭtvam prasiddham . \\
\noindent \textbf{(P\_3,1.45)} aniṭi prasiddhe ksaviddhiḥ . \\
\noindent \textbf{(P\_3,1.45)} aniṭi prasiddhe ksaḥ bhaviṣyati . \\
\noindent \textbf{(P\_3,1.45)} sic idānīm kva bhaviṣyati . \\
\noindent \textbf{(P\_3,1.45)} śeṣe sijvidhānam . \\
\noindent \textbf{(P\_3,1.45)} śeṣe sijvidhānam bhaviṣyati . \\
\noindent \textbf{(P\_3,1.45)} akoṣīt amoṣīt iti . \\
\noindent \textbf{(P\_3,1.46)} kimartham idam ucyate . \\
\noindent \textbf{(P\_3,1.46)} niyamārtham . \\
\noindent \textbf{(P\_3,1.46)} śliṣaḥ āliṅgane eva ksaḥ yathā syāt . \\
\noindent \textbf{(P\_3,1.46)} iha mā bhūt : upāśliṣat jatu ca kāṣṭham ca . \\
\noindent \textbf{(P\_3,1.46)} samāśliṣat brāhmaṇakulam iti . \\
\noindent \textbf{(P\_3,1.46)} ataḥ uttaram paṭhati . \\
\noindent \textbf{(P\_3,1.46)} śliṣaḥ āliṅgane niyamānupapattiḥ vidheyabhāvāt . \\
\noindent \textbf{(P\_3,1.46)} śliṣaḥ āliṅgane niyamasya anupapattiḥ . \\
\noindent \textbf{(P\_3,1.46)} kim kāraṇam . \\
\noindent \textbf{(P\_3,1.46)} vidheyabhāvāt . \\
\noindent \textbf{(P\_3,1.46)} kaimarthakyāt niyamaḥ bhavati . \\
\noindent \textbf{(P\_3,1.46)} vidheyam na asti iti kṛtvā . \\
\noindent \textbf{(P\_3,1.46)} iha ca asti vidheyam . \\
\noindent \textbf{(P\_3,1.46)} kim . \\
\noindent \textbf{(P\_3,1.46)} puṣādipāṭhāt aṅ prāptaḥ . \\
\noindent \textbf{(P\_3,1.46)} tadbādhanārthaḥ ksaḥ vidheyaḥ . \\
\noindent \textbf{(P\_3,1.46)} tatra apūrvaḥ vidhiḥ astu niyamaḥ vā iti apūrvaḥ eva vidhiḥ syāt na niyamaḥ . \\
\noindent \textbf{(P\_3,1.46)} kim ca syāt yadi ayam niyamaḥ na syāt . \\
\noindent \textbf{(P\_3,1.46)} ātmanepadeṣu āliṅgane ca ksaḥ prasajyeta . \\
\noindent \textbf{(P\_3,1.46)} yathā eva ca ksaḥ aṅam bādhate evam ciṇam api bādheta . \\
\noindent \textbf{(P\_3,1.46)} upāśleṣi kanyā devadattena iti . \\
\noindent \textbf{(P\_3,1.46)} siddham tu śliṣaḥ āliṅgane aciṇviṣaye . \\
\noindent \textbf{(P\_3,1.46)} siddham etat . \\
\noindent \textbf{(P\_3,1.46)} katham . \\
\noindent \textbf{(P\_3,1.46)} śliṣaḥ āliṅgane aciṇviṣaye ksaḥ bhavati iti vaktavyam . \\
\noindent \textbf{(P\_3,1.46)} aṅvidhāne ca śliṣaḥ anāliṅgane . \\
\noindent \textbf{(P\_3,1.46)} aṅvidhāne ca śliṣaḥ anāliṅgane iti vaktavyam . \\
\noindent \textbf{(P\_3,1.46)} sidhyati . \\
\noindent \textbf{(P\_3,1.46)} sūtram tarhi bhidyate . \\
\noindent \textbf{(P\_3,1.46)} yathānyāsam eva astu . \\
\noindent \textbf{(P\_3,1.46)} nanu ca uktam śliṣaḥ āliṅgane niyamānupapattiḥ vidheyabhāvāt iti . \\
\noindent \textbf{(P\_3,1.46)} na eṣaḥ doṣaḥ . \\
\noindent \textbf{(P\_3,1.46)} yogavibhāgāt siddham . \\
\noindent \textbf{(P\_3,1.46)} yogavibhāgaḥ kariṣyate . \\
\noindent \textbf{(P\_3,1.46)} śliṣaḥ . \\
\noindent \textbf{(P\_3,1.46)} śliṣaḥ ksaḥ bhavati . \\
\noindent \textbf{(P\_3,1.46)} kimartham idam . \\
\noindent \textbf{(P\_3,1.46)} puṣādipāṭhāt aṅ prāpnoti . \\
\noindent \textbf{(P\_3,1.46)} tadbādhanārtham . \\
\noindent \textbf{(P\_3,1.46)} tataḥ āliṅgane . \\
\noindent \textbf{(P\_3,1.46)} āliṅgane ca śliṣaḥ ksaḥ bhavati . \\
\noindent \textbf{(P\_3,1.46)} idam idānīm kimartham . \\
\noindent \textbf{(P\_3,1.46)} niyamārtham . \\
\noindent \textbf{(P\_3,1.46)} śliṣaḥ āliṅgane eva . \\
\noindent \textbf{(P\_3,1.46)} kva mā bhūt . \\
\noindent \textbf{(P\_3,1.46)} upāśliṣat jatu ca kāṣṭham ca . \\
\noindent \textbf{(P\_3,1.46)} samāśliṣat brāhmaṇakulam iti . \\
\noindent \textbf{(P\_3,1.46)} yat api ucyate yathā eva ca ksaḥ aṅam bādhate evam ciṇam api bādheta iti . \\
\noindent \textbf{(P\_3,1.46)} purastāt apavādāḥ anantarān vidhīn bādhante na uttarān iti evam ksaḥ aṅam bādhiṣyate . \\
\noindent \textbf{(P\_3,1.46)} ciṇam na bādhiṣyate . \\
\noindent \textbf{(P\_3,1.46)} atha vā tatra vakṣyati : ciṇgrahaṇasya prayojanam ciṇ eva yathā syāt . \\
\noindent \textbf{(P\_3,1.46)} yat anyat prāpnoti tat mā bhūt iti . \\
\noindent \textbf{(P\_3,1.48)} ṇiśridrusruṣu kameḥ upasaṅkhyānam . ṇiśridrusruṣu kameḥ upasaṅkhyānam kartavyam . \\
\noindent \textbf{(P\_3,1.48)} nākam iṣṭamukham yānti suyuktaiḥ vaḍavārathaiḥ . \\
\noindent \textbf{(P\_3,1.48)} atha patkāṣīṇaḥ yānti ye acīkamatabhāṣiṇaḥ . \\
\noindent \textbf{(P\_3,1.48)} karmakartari ca . \\
\noindent \textbf{(P\_3,1.48)} karmakartari ca upasaṅkhyānam kartavyam . \\
\noindent \textbf{(P\_3,1.48)} kārayati kaṭam devadattaḥ . \\
\noindent \textbf{(P\_3,1.48)} acīkarata kaṭaḥ svayam eva . \\
\noindent \textbf{(P\_3,1.48)} ucchrayayati kaṭam devadattaḥ . \\
\noindent \textbf{(P\_3,1.48)} audaśiśriyata kaṭaḥ svayam eva . \\
\noindent \textbf{(P\_3,1.48)} na vā karmaṇi avidhānāt kartṛtvāt ca karmakartuḥ siddham . \\
\noindent \textbf{(P\_3,1.48)} na vā kartavyam . \\
\noindent \textbf{(P\_3,1.48)} kim kāraṇam . \\
\noindent \textbf{(P\_3,1.48)} karmaṇi avidhānāt . \\
\noindent \textbf{(P\_3,1.48)} na hi kaḥ cit karmaṇi vidhīyate yaḥ caṅam bādheta . \\
\noindent \textbf{(P\_3,1.48)} kartṛtvāt ca karmakartuḥ siddham . \\
\noindent \textbf{(P\_3,1.48)} asti ca karmakartari kartṛtvam iti kṛtvā caṅ bhaviṣyati . \\
\noindent \textbf{(P\_3,1.48)} nanu ca ayam karmaṇi vidhīyate . \\
\noindent \textbf{(P\_3,1.48)} ciṇ bhāvakarmaṇoḥ iti . \\
\noindent \textbf{(P\_3,1.48)} pratiṣidhyete tatra yakciṇau . \\
\noindent \textbf{(P\_3,1.48)} yakciṇoḥ pratiṣedhe hetumaṇṇiśribrūñām upasaṅkhyānam iti . \\
\noindent \textbf{(P\_3,1.48)} yaḥ tarhi ahetumaṇṇic . \\
\noindent \textbf{(P\_3,1.48)} udapupucchata gauḥ svayam eva . \\
\noindent \textbf{(P\_3,1.48)} atra api yathā bhāradvājīyāḥ paṭhanti tathā bhavitavyam pratiṣedhena . \\
\noindent \textbf{(P\_3,1.48)} yakciṇoḥ pratiṣedhe ṇiśrigranthibrūñām ātmanepadākarmakāṇām upasaṅkhyānam iti . \\
\noindent \textbf{(P\_3,1.52)} asyatigrahaṇam kimartham . \\
\noindent \textbf{(P\_3,1.52)} asyatigrahaṇam ātmanepadārtham . \\
\noindent \textbf{(P\_3,1.52)} asyatigrahaṇam ātmanepadārtham draṣṭavyam . \\
\noindent \textbf{(P\_3,1.52)} kim ucyate ātmanepadārtham iti . \\
\noindent \textbf{(P\_3,1.52)} na punaḥ parasmaipadārtham api syāt . \\
\noindent \textbf{(P\_3,1.52)} puṣāditvāt . \\
\noindent \textbf{(P\_3,1.52)} puṣādipāṭhāt parasmaipadeṣu aṅ bhaviṣyati . \\
\noindent \textbf{(P\_3,1.52)} karmakartari ca . \\
\noindent \textbf{(P\_3,1.52)} karmakartari ca upasaṅkhyānam kartavyam . \\
\noindent \textbf{(P\_3,1.52)} paryāsthetām kuṇḍale svayam eva . \\
\noindent \textbf{(P\_3,1.52)} atra api na vā karmaṇi avidhānāt kartṛtvāt ca karmakartuḥ siddham iti eva . \\
\noindent \textbf{(P\_3,1.58)} idam lucigrahaṇam gluñcigrahaṇam ca kriyate . \\
\noindent \textbf{(P\_3,1.58)} anyatarat śakyam akartum . \\
\noindent \textbf{(P\_3,1.58)} katham . \\
\noindent \textbf{(P\_3,1.58)} yadi tāvat glucigrahaṇam kriyate gluñcigrahaṇam na kariṣyate . \\
\noindent \textbf{(P\_3,1.58)} tena eva siddham nyaglucat nyaglocīt . \\
\noindent \textbf{(P\_3,1.58)} idam idānīm gluñceḥ rūpam nyagluñcīt . \\
\noindent \textbf{(P\_3,1.58)} atha gluñcigrahaṇam kriyate gluceḥ grahaṇam na kariṣyate . \\
\noindent \textbf{(P\_3,1.58)} tena eva siddham nyaglucat nyagluñcīt . \\
\noindent \textbf{(P\_3,1.58)} idam idānīm gluceḥ rūpam nyaglocīt . \\
\noindent \textbf{(P\_3,1.60)} ayam taśabdaḥ asti eva ātmanepadam asti parasmaipadam asti ekavacanam asti bahuvacanam . \\
\noindent \textbf{(P\_3,1.60)} kasya idam grahaṇam . \\
\noindent \textbf{(P\_3,1.60)} yaḥ padeḥ asti . \\
\noindent \textbf{(P\_3,1.60)} kaḥ ca padeḥ asti . \\
\noindent \textbf{(P\_3,1.60)} padiḥ ayam ātmanepadī . \\
\noindent \textbf{(P\_3,1.66)} ciṇ iti vartamāne punaḥ ciṇgrahaṇam kimartham . \\
\noindent \textbf{(P\_3,1.66)} na iti evam tat abhūt . \\
\noindent \textbf{(P\_3,1.66)} vidhyartham idam . \\
\noindent \textbf{(P\_3,1.66)} atha vā vā iti evam tat abhūt . \\
\noindent \textbf{(P\_3,1.66)} nityārtham idam . \\
\noindent \textbf{(P\_3,1.66)} atha vā ciṇ iti vartamāne punaḥ ciṇgrahaṇasya etat prayojanam . \\
\noindent \textbf{(P\_3,1.66)} ciṇ eva yathā syāt . \\
\noindent \textbf{(P\_3,1.66)} yat anyat prāpnoti tat mā bhūt iti . \\
\noindent \textbf{(P\_3,1.67.1)} iha paśyāmaḥ karmaṇi dvivacanabahuvacanāni udāhriyante . \\
\noindent \textbf{(P\_3,1.67.1)} pacyete* odanau , pacyante odanāḥ iti . \\
\noindent \textbf{(P\_3,1.67.1)} bhāve punaḥ ekavacanam eva : āsyate bhavatā , āsyate bhavadbhyām , āsyate bhavadbhiḥ iti . \\
\noindent \textbf{(P\_3,1.67.1)} kena etat evam bhavati . \\
\noindent \textbf{(P\_3,1.67.1)} karma anekam . \\
\noindent \textbf{(P\_3,1.67.1)} tasya anekatvāt dvivacanabahuvacanāni bhavanti . \\
\noindent \textbf{(P\_3,1.67.1)} bhāvaḥ punaḥ ekaḥ eva . \\
\noindent \textbf{(P\_3,1.67.1)} katham tarhi iha dvivacanabahuvacanāni bhavanti . \\
\noindent \textbf{(P\_3,1.67.1)} pākau pākāḥ iti . \\
\noindent \textbf{(P\_3,1.67.1)} āśrayabhedāt . \\
\noindent \textbf{(P\_3,1.67.1)} yat asau dravyam śritaḥ bhavati bhāvaḥ tasya bhedāt dvivacanabahuvacanāni bhavanti . \\
\noindent \textbf{(P\_3,1.67.1)} iha api tarhi yāvantaḥ tām kriyām kurvanti sarve te tasyāḥ āśrayā bhavanti . \\
\noindent \textbf{(P\_3,1.67.1)} tadbhedāt dvivacanabahuvacanāni prāpnuvanti . \\
\noindent \textbf{(P\_3,1.67.1)} evam tarhi idam tāvat ayam praṣṭavyaḥ . \\
\noindent \textbf{(P\_3,1.67.1)} kim abhisamīkṣya etat prayujyate . \\
\noindent \textbf{(P\_3,1.67.1)} pākau pākāḥ iti . \\
\noindent \textbf{(P\_3,1.67.1)} yadi tāvat pākaviśeṣān abhisamīkṣya yaḥ ca odanasya pākaḥ yaḥ ca guḍasya yaḥ ca tilānām bahavaḥ te śabdāḥ sarūpāḥ ca . \\
\noindent \textbf{(P\_3,1.67.1)} tatra yuktam bahuvacanam ekaśeṣaḥ ca . \\
\noindent \textbf{(P\_3,1.67.1)} tiṅabhihite ca api tadā bhāve bahuvacanam śrūyate . \\
\noindent \textbf{(P\_3,1.67.1)} tat yathā : uṣṭṛāsikā āsyante . \\
\noindent \textbf{(P\_3,1.67.1)} hataśāyikāḥ śayyante iti . \\
\noindent \textbf{(P\_3,1.67.1)} atha kālaviśeṣān abhisamīkṣya yaḥ ca adyatanaḥ pākaḥ yaḥ hyastanaḥ yaḥ śvastanaḥ te api bahavaḥ śabdāḥ sarūpāḥ ca . \\
\noindent \textbf{(P\_3,1.67.1)} tatra yuktam bahuvacanam ekaśeṣaḥ ca . \\
\noindent \textbf{(P\_3,1.67.1)} tiṅabhihite ca api tadā bhāve asārūpyāt ekaśeṣaḥ na bhavati . \\
\noindent \textbf{(P\_3,1.67.1)} āsi āsyate , āsiṣyate . \\
\noindent \textbf{(P\_3,1.67.1)} asti khalu api viśeṣaḥ kṛdabhihitasya bhāvasya tiṅabhihitasya ca . \\
\noindent \textbf{(P\_3,1.67.1)} kṛdabhihitaḥ bhāvaḥ dravyavat bhavati . \\
\noindent \textbf{(P\_3,1.67.1)} kim idam dravyavat iti . \\
\noindent \textbf{(P\_3,1.67.1)} dravyam kriyayā samavāyam gacchati . \\
\noindent \textbf{(P\_3,1.67.1)} kam samavāyam . \\
\noindent \textbf{(P\_3,1.67.1)} dravyam kriyābhinirvṛttau sādhanatvam upaiti . \\
\noindent \textbf{(P\_3,1.67.1)} tadvat ca asya bhāvasya kṛdabhihitasya bhavati . \\
\noindent \textbf{(P\_3,1.67.1)} pākaḥ vartate iti . \\
\noindent \textbf{(P\_3,1.67.1)} kriyāvat na bhavati . \\
\noindent \textbf{(P\_3,1.67.1)} kim idam kriyāvat iti . \\
\noindent \textbf{(P\_3,1.67.1)} kriyā kriyayā samavāyam na gacchati . \\
\noindent \textbf{(P\_3,1.67.1)} pacati paṭhati iti . \\
\noindent \textbf{(P\_3,1.67.1)} tadvac ca asya kṛtabhihitasya na bhavati . \\
\noindent \textbf{(P\_3,1.67.1)} pākaḥ vartate iti . \\
\noindent \textbf{(P\_3,1.67.1)} asti khalu api viśeṣaḥ kṛdabhihitasya bhāvasya tiṅabhihitasya ca . \\
\noindent \textbf{(P\_3,1.67.1)} tiṅabhihitena bhāvena kālapuruṣopagrahāḥ abhivyajyante . \\
\noindent \textbf{(P\_3,1.67.1)} kṛdabhihitena punaḥ na vyajyante . \\
\noindent \textbf{(P\_3,1.67.1)} asti khalu api viśeṣaḥ kṛdabhihitasya bhāvasya tiṅabhihitasya ca . \\
\noindent \textbf{(P\_3,1.67.1)} tiṅabhihitaḥ bhāvaḥ kartrā samprayujyate. kṛdabhihitaḥ punaḥ na samprayujyate . \\
\noindent \textbf{(P\_3,1.67.1)} yāvatā kim cit sāmānyam kaḥ cit viśeṣaḥ yuktam yat ayam api viśeṣaḥ syāt liṅgakṛtaḥ saṅkhyākṛtaḥ ca iti . \\
\noindent \textbf{(P\_3,1.67.2)} idam vicāryate . \\
\noindent \textbf{(P\_3,1.67.2)} bhāvakarmakartāraḥ sārvadhātukārthāḥ vā syuḥ vikaraṇārthāḥ vā iti . \\
\noindent \textbf{(P\_3,1.67.2)} katham ca sārvadhātukārthaḥ syuḥ katham vā vikaraṇārthāḥ . \\
\noindent \textbf{(P\_3,1.67.2)} bhāvakarmavācini sārvadhātuke yak bhavati kartṛvācini śarvadhātuke śap bhavati iti sārvadhātukārthāḥ . \\
\noindent \textbf{(P\_3,1.67.2)} bhāvakarmaṇoḥ yag bhavati sārvadhātuke kartari śap bhavati sārvadhātuke iti vikaraṇārthāḥ . \\
\noindent \textbf{(P\_3,1.67.2)} kaḥ ca atra viśeṣaḥ . \\
\noindent \textbf{(P\_3,1.67.2)} bhāvakarmakartāraḥ sārvadhātukārthāḥ cet ekadvibahuṣu niyamānupapattiḥ atadarthatvāt . \\
\noindent \textbf{(P\_3,1.67.2)} bhāvakarmakartāraḥ sārvadhātukārthāḥ cet ekadvibahuṣu niyamasya anupapattiḥ . \\
\noindent \textbf{(P\_3,1.67.2)} kim kāraṇam . \\
\noindent \textbf{(P\_3,1.67.2)} atadarthatvāt . \\
\noindent \textbf{(P\_3,1.67.2)} na hi tadānīm ekatvādayaḥ eva vibhaktyarthāḥ . \\
\noindent \textbf{(P\_3,1.67.2)} kim tarhi bhāvakarmakartāraḥ api . \\
\noindent \textbf{(P\_3,1.67.2)} santu tarhi vikaraṇārthāḥ . \\
\noindent \textbf{(P\_3,1.67.2)} vikaraṇārthāḥ iti cet kṛtā abhihite vikaraṇābhāvaḥ . \\
\noindent \textbf{(P\_3,1.67.2)} vikaraṇārthāḥ iti cet kṛtā abhihite vikaraṇaḥ na prāpnoti . \\
\noindent \textbf{(P\_3,1.67.2)} dhārayaḥ pārayaḥ iti . \\
\noindent \textbf{(P\_3,1.67.2)} kim ucyate kṛtā abhihite . \\
\noindent \textbf{(P\_3,1.67.2)} na lena api abhidhānam bhavati . \\
\noindent \textbf{(P\_3,1.67.2)} aśakyam lena abhidhānam āśrayitum . \\
\noindent \textbf{(P\_3,1.67.2)} pakṣāntaram idam āsthitam bhāvakarmakartāraḥ sārvadhātukārthāḥ vā syuḥ vikaraṇārthāḥ vā iti . \\
\noindent \textbf{(P\_3,1.67.2)} yadi ca lena api abhidhānam syāt na idam pakṣāntaram syāt . \\
\noindent \textbf{(P\_3,1.67.2)} katham aśakyam yadā bhavān eva āha laḥ karmaṇi ca bhāve ca akarmakebhyaḥ iti . \\
\noindent \textbf{(P\_3,1.67.2)} evam vakṣyāmi . \\
\noindent \textbf{(P\_3,1.67.2)} laḥ karmaṇaḥ bhāvāt ca akarmakebhyaḥ . \\
\noindent \textbf{(P\_3,1.67.2)} yasmin tarhi le vikaraṇāḥ na śrūyante kaḥ tatra bhāvakarmakartṝn abhidhāsyati . \\
\noindent \textbf{(P\_3,1.67.2)} kva ca na śrūyante . \\
\noindent \textbf{(P\_3,1.67.2)} ye ete lugvikaraṇāḥ śluvikaraṇāḥ ca . \\
\noindent \textbf{(P\_3,1.67.2)} atra api ukte kartṛtve luk bhaviṣyati . \\
\noindent \textbf{(P\_3,1.67.2)} yasmin tarhi le vikaraṇāḥ na eva utpadyante kaḥ tatra bhāvakarmakartṝn abhidhāsyati . \\
\noindent \textbf{(P\_3,1.67.2)} kva ca na eva utpadyante . \\
\noindent \textbf{(P\_3,1.67.2)} liṅliṭoḥ . \\
\noindent \textbf{(P\_3,1.67.2)} tasmāt na etat śakyam vaktum . \\
\noindent \textbf{(P\_3,1.67.2)} na lena abhidhānam bhavati iti . \\
\noindent \textbf{(P\_3,1.67.2)} bhavati cet abhihite vikaraṇābhāvaḥ eva . \\
\noindent \textbf{(P\_3,1.67.2)} evam tarhi idam syāt . \\
\noindent \textbf{(P\_3,1.67.2)} yadā bhāvakarmaṇoḥ laḥ tadā kartari vikaraṇāḥ . \\
\noindent \textbf{(P\_3,1.67.2)} yadā kartari laḥ tadā bhāvakarmaṇoḥ vikaraṇāḥ . \\
\noindent \textbf{(P\_3,1.67.2)} idam asya yadi eva svābhāvikam atha api vācanikam : prakṛtipratyayau pratyayārtham saha brūtaḥ iti . \\
\noindent \textbf{(P\_3,1.67.2)} na ca asti sambhavaḥ yat ekasyāḥ prakṛteḥ dvayoḥ nānārthayoḥ yugapat anusahāyībhāvaḥ syāt . \\
\noindent \textbf{(P\_3,1.67.2)} evam ca kṛtvā ekapakṣībhūtam idam bhavati : sārvadhātukārthāḥ eva iti . \\
\noindent \textbf{(P\_3,1.67.2)} nanu ca uktam bhāvakarmakartāraḥ sārvadhātukārthāḥ cet ekadvibahuṣu niyamānupapattiḥ atadarthatvāt iti . \\
\noindent \textbf{(P\_3,1.67.2)} na eṣaḥ doṣaḥ . \\
\noindent \textbf{(P\_3,1.67.2)} supām karmādayaḥ api arthāḥ saṅkhyā ca eva tathā tiṅām . supām saṅkhyā ca eva arthaḥ karmādayaḥ ca . \\
\noindent \textbf{(P\_3,1.67.2)} tathā tiṅām . \\
\noindent \textbf{(P\_3,1.67.2)} prasiddhaḥ niyamaḥ tatra . \\
\noindent \textbf{(P\_3,1.67.2)} prasiddhaḥ tatra niyamaḥ . \\
\noindent \textbf{(P\_3,1.67.2)} niyamaḥ prakṛteṣu vā . \\
\noindent \textbf{(P\_3,1.67.2)} atha vā prakṛtān arthān apekṣya niyamaḥ . \\
\noindent \textbf{(P\_3,1.67.2)} ke ca prakṛtāḥ . \\
\noindent \textbf{(P\_3,1.67.2)} ekatvādayaḥ . \\
\noindent \textbf{(P\_3,1.67.2)} ekasmin eva ekavacanam na dvayoḥ na bahuṣu . \\
\noindent \textbf{(P\_3,1.67.2)} dvayoḥ eva dvivacanam naikasmin na bahuṣu . \\
\noindent \textbf{(P\_3,1.67.2)} bahuṣu eva bahuvacanam na dvayoḥ na ekasmin iti . \\
\noindent \textbf{(P\_3,1.67.3)} bhāvakarmaṇoḥ yagvidhāne karmakartari upasaṅkhyānam . bhāvakarmaṇoḥ yagvidhāne karmakartari upasaṅkhyānam kartavyam . \\
\noindent \textbf{(P\_3,1.67.3)} pacyate svayam eva . \\
\noindent \textbf{(P\_3,1.67.3)} paṭhyate svayam eva . \\
\noindent \textbf{(P\_3,1.67.3)} kim punaḥ kāraṇam na sidhyati . \\
\noindent \textbf{(P\_3,1.67.3)} vipratiṣedhāt hi śapaḥ balīyastvam . vipratiṣedhāt hi śapaḥ balīyastvam prāpnoti . \\
\noindent \textbf{(P\_3,1.67.3)} śapaḥ avakāśaḥ . \\
\noindent \textbf{(P\_3,1.67.3)} pacati paṭhati . \\
\noindent \textbf{(P\_3,1.67.3)} yakaḥ avakāśaḥ . \\
\noindent \textbf{(P\_3,1.67.3)} pacyate odanaḥ devadattena . \\
\noindent \textbf{(P\_3,1.67.3)} paṭhyate vidyā devadattena . \\
\noindent \textbf{(P\_3,1.67.3)} iha ubhayam prāpnoti . \\
\noindent \textbf{(P\_3,1.67.3)} pacyate svayam eva . \\
\noindent \textbf{(P\_3,1.67.3)} paṭhyate svayam eva . \\
\noindent \textbf{(P\_3,1.67.3)} paratvāt śap prāpnoti . \\
\noindent \textbf{(P\_3,1.67.3)} yogavibhāgāt siddham . \\
\noindent \textbf{(P\_3,1.67.3)} yogavibhāgaḥ kariṣyate . \\
\noindent \textbf{(P\_3,1.67.3)} ciṇ bhāvakarmaṇoḥ . \\
\noindent \textbf{(P\_3,1.67.3)} sārvadhātuke yak bhāvakarmaṇoḥ . \\
\noindent \textbf{(P\_3,1.67.3)} tataḥ kartari . \\
\noindent \textbf{(P\_3,1.67.3)} kartari ca yak bhavati bhāvakarmaṇoḥ . \\
\noindent \textbf{(P\_3,1.67.3)} yathā eva tarhi karmaṇi kartari yak bhavati evam bhāve kartari prāpnoti . \\
\noindent \textbf{(P\_3,1.67.3)} eti jīvantam ānandaḥ . \\
\noindent \textbf{(P\_3,1.67.3)} na asya kim cit rujati rogaḥ iti . \\
\noindent \textbf{(P\_3,1.67.3)} dvitīyaḥ yogavibhāgaḥ kariṣyate . \\
\noindent \textbf{(P\_3,1.67.3)} ciṇ bhāve . \\
\noindent \textbf{(P\_3,1.67.3)} tataḥ karmaṇi . \\
\noindent \textbf{(P\_3,1.67.3)} karmaṇi ca ciṇ bhavati . \\
\noindent \textbf{(P\_3,1.67.3)} tataḥ sārvadhātuke yak bhavati bhāve ca karmaṇi ca . \\
\noindent \textbf{(P\_3,1.67.3)} tataḥ kartari . \\
\noindent \textbf{(P\_3,1.67.3)} kartari ca yak bhavati . \\
\noindent \textbf{(P\_3,1.67.3)} karmaṇi iti anuvartate . \\
\noindent \textbf{(P\_3,1.67.3)} bhāve iti nivṛttam . \\
\noindent \textbf{(P\_3,1.67.3)} tataḥ śap . \\
\noindent \textbf{(P\_3,1.67.3)} śap ca bhavati . \\
\noindent \textbf{(P\_3,1.67.3)} kartari iti eva . \\
\noindent \textbf{(P\_3,1.67.3)} karmaṇi iti api nivṛttam . \\
\noindent \textbf{(P\_3,1.67.3)} evam api upasaṅkhyānam kartavyam . \\
\noindent \textbf{(P\_3,1.67.3)} vipratiṣedhāt hi śyanaḥ balīyastvam prāpnoti . \\
\noindent \textbf{(P\_3,1.67.3)} śyanaḥ avakāśaḥ . \\
\noindent \textbf{(P\_3,1.67.3)} dīvyati sīvyati . \\
\noindent \textbf{(P\_3,1.67.3)} yakaḥ avakāśaḥ . \\
\noindent \textbf{(P\_3,1.67.3)} pacyate odanaḥ devadattena . \\
\noindent \textbf{(P\_3,1.67.3)} paṭhyate vidyā devadattena . \\
\noindent \textbf{(P\_3,1.67.3)} iha ubhayam prāpnoti . \\
\noindent \textbf{(P\_3,1.67.3)} dīvyate svayam eva . \\
\noindent \textbf{(P\_3,1.67.3)} sīvyate svayam eva . \\
\noindent \textbf{(P\_3,1.67.3)} paratvāt śyan prāpnoti . \\
\noindent \textbf{(P\_3,1.67.3)} nanu ca etat api yogavibhāgāt eva siddham . \\
\noindent \textbf{(P\_3,1.67.3)} na sidhyati . \\
\noindent \textbf{(P\_3,1.67.3)} anantarā yā praptiḥ sā yogavibhāgena śakyā bādhitum . \\
\noindent \textbf{(P\_3,1.67.3)} kutaḥ etat . \\
\noindent \textbf{(P\_3,1.67.3)} anantarasya vidhiḥ vā bhavati pratiṣedhaḥ vā iti . \\
\noindent \textbf{(P\_3,1.67.3)} parā prāptiḥ apratiṣiddhā . \\
\noindent \textbf{(P\_3,1.67.3)} tayā prāpnoti . \\
\noindent \textbf{(P\_3,1.67.3)} nanu ca iyam prāptiḥ parām prāptim bādheta . \\
\noindent \textbf{(P\_3,1.67.3)} na utsahate pratiṣiddhā satī bādhitum . \\
\noindent \textbf{(P\_3,1.67.3)} evam tarhi śabādeśāḥ śyanādayaḥ kariṣyante . \\
\noindent \textbf{(P\_3,1.67.3)} śap ca syādibhiḥ bādhyate . \\
\noindent \textbf{(P\_3,1.67.3)} tatra divādibhyaḥ yagviṣaye śap eva na asti kutaḥ śyanādayaḥ . \\
\noindent \textbf{(P\_3,1.67.3)} tat tarhi śapaḥ grahaṇam kartavyam . \\
\noindent \textbf{(P\_3,1.67.3)} na kartavyam . \\
\noindent \textbf{(P\_3,1.67.3)} prakṛtam anuvartate . \\
\noindent \textbf{(P\_3,1.67.3)} kva prakṛtam . \\
\noindent \textbf{(P\_3,1.67.3)} kartari śap iti . \\
\noindent \textbf{(P\_3,1.67.3)} tat vai prathamānirdiṣṭam ṣaṣṭhīnirdiṣṭena ca iha arthaḥ . \\
\noindent \textbf{(P\_3,1.67.3)} divādibhyaḥ iti eṣā pañcamī śap iti prathamāyāḥ ṣaṣṭhīm prakalpayiṣyati tasmāt iti uttarasya iti . \\
\noindent \textbf{(P\_3,1.67.3)} pratyayavidhiḥ ayam . \\
\noindent \textbf{(P\_3,1.67.3)} na ca pratyayavidhau pañcamyaḥ prakalpikāḥ bhavanti . \\
\noindent \textbf{(P\_3,1.67.3)} na ayam pratyayavidhiḥ . \\
\noindent \textbf{(P\_3,1.67.3)} vihitaḥ pratyayaḥ . \\
\noindent \textbf{(P\_3,1.67.3)} prakṛtaḥ ca anuvartate . \\
\noindent \textbf{(P\_3,1.67.3)} atha vā bhāvakarmaṇoḥ iti anuvṛttyā eva siddhe sati anivṛttiḥ yakaḥ bhāvāya .iha sārvadhātuke yak iti antareṇa bhāvakarmaṇoḥ iti anuvṛttim siddham . \\
\noindent \textbf{(P\_3,1.67.3)} saḥ ayam evam siddhe sati yat bhāvakarmaṇoḥ iti anuvartayati tasya etat prayojanam . \\
\noindent \textbf{(P\_3,1.67.3)} karmakartari api yathā syāt . \\
\noindent \textbf{(P\_3,1.67.3)} kartari iti ca yogavibhāgaḥ śyanaḥ pūrvavipratiṣedhāvacanāya . \\
\noindent \textbf{(P\_3,1.67.3)} kartari iti yogavibhāgaḥ kartavyaḥ śyanaḥ pūrvavipratiṣedham mā vocam iti . \\
\noindent \textbf{(P\_3,1.67.3)} atha vā karmavadbhāvavacanasāmarthyāt yak bhaviṣyati . \\
\noindent \textbf{(P\_3,1.67.3)} asti anyat karmavadbhāvavacane prayojanam . \\
\noindent \textbf{(P\_3,1.67.3)} kim . \\
\noindent \textbf{(P\_3,1.67.3)} ātmanepadam yathā syāt . \\
\noindent \textbf{(P\_3,1.67.3)} vacanāt ātmanepadam bhaviṣyati . \\
\noindent \textbf{(P\_3,1.67.3)} ciṇ tarhi yathā syāt . \\
\noindent \textbf{(P\_3,1.67.3)} ciṇ api vacanāt bhaviṣyati . \\
\noindent \textbf{(P\_3,1.67.3)} ciṇvadbhāvaḥ tarhi yathā syāt . \\
\noindent \textbf{(P\_3,1.67.3)} na ekam prayojanam yogārambham prayojayati . \\
\noindent \textbf{(P\_3,1.67.3)} tatra karmavadbhāvavacanasāmarthyāt yak bhaviṣyati . \\
\noindent \textbf{(P\_3,1.67.3)} atha vā ācāryapravṛttiḥ jñāpayati bhavati karmakartari yak iti yat ayam na duhasnnunamām yakciṇau iti yakciṇoḥ pratiṣedham śāsti . \\
\noindent \textbf{(P\_3,1.71)} anupasargāt iti kimartham . \\
\noindent \textbf{(P\_3,1.71)} āyasyati prayasyati . \\
\noindent \textbf{(P\_3,1.71)} anupasargāt iti śakyam akartum . \\
\noindent \textbf{(P\_3,1.71)} katham āyasyati prayasyati . \\
\noindent \textbf{(P\_3,1.71)} saṃyasaḥ ca iti etat niyamārtham bhaviṣyati . \\
\noindent \textbf{(P\_3,1.71)} sampūrvāt yasaḥ na anyapūrvāt iti . \\
\noindent \textbf{(P\_3,1.78)} kimarthaḥ śakāraḥ . \\
\noindent \textbf{(P\_3,1.78)} sārvadhātukāṛthaḥ . \\
\noindent \textbf{(P\_3,1.78)} śit sārvadhātukam iti sārvadhātukasañjñā . \\
\noindent \textbf{(P\_3,1.78)} sārvadhātukam apit iti ṅittvam . \\
\noindent \textbf{(P\_3,1.78)} ṅiti iti guṇapratiṣedhaḥ yathā syāt . \\
\noindent \textbf{(P\_3,1.78)} bhinatti chinatti iti . \\
\noindent \textbf{(P\_3,1.78)} na etat asti prayojanam . \\
\noindent \textbf{(P\_3,1.78)} sārvadhātukārdhadhātukayoḥ aṅgasya guṇaḥ ucyate yasmāt ca pratyayavidhiḥ tadādi pratyaye aṅgasañjñam bhavati . \\
\noindent \textbf{(P\_3,1.78)} yasmāt ca atra pratyayavidhiḥ na tat pratyaye parataḥ . \\
\noindent \textbf{(P\_3,1.78)} yat ca pratyaye parataḥ na tasmāt pratyayavidhiḥ . \\
\noindent \textbf{(P\_3,1.78)} idam tarhi prayojanam . \\
\noindent \textbf{(P\_3,1.78)} ārdhadhātukasañjñā mā bhūt iti . \\
\noindent \textbf{(P\_3,1.78)} kim ca syāt . \\
\noindent \textbf{(P\_3,1.78)} valādilakṣaṇaḥ iṭ prasajyeta . \\
\noindent \textbf{(P\_3,1.78)} etat api na asti prayojanam . \\
\noindent \textbf{(P\_3,1.78)} valādeḥ ārdhadhātukasya aṅgasya iṭ ucyate . \\
\noindent \textbf{(P\_3,1.78)} yasmāt ca pratyayavidhiḥ tadādi pratyaye aṅgasañjñam bhavati . \\
\noindent \textbf{(P\_3,1.78)} yasmāt ca atra pratyayavidhiḥ na tat pratyaye parataḥ . \\
\noindent \textbf{(P\_3,1.78)} yat ca pratyaye parataḥ na tasmāt pratyayavidhiḥ . \\
\noindent \textbf{(P\_3,1.78)} ataḥ uttaram paṭhati . \\
\noindent \textbf{(P\_3,1.78)} śnami śitkaraṇam pvādihrasvārtham . \\
\noindent \textbf{(P\_3,1.78)} śnami śitkaraṇam kriyate pvādīnām śiti hrasvatvam yathā syāt . \\
\noindent \textbf{(P\_3,1.78)} pṛṇasi mṛṇasi iti . \\
\noindent \textbf{(P\_3,1.78)} na vā dhātvanyatvāt . \\
\noindent \textbf{(P\_3,1.78)} na vā kartavyam . \\
\noindent \textbf{(P\_3,1.78)} kim kāraṇam . \\
\noindent \textbf{(P\_3,1.78)} dhātvanyatvāt . \\
\noindent \textbf{(P\_3,1.78)} dhātvantaram pṛṇimṛṇī . \\
\noindent \textbf{(P\_3,1.78)} yatra bhūmyām vṛṇase . \\
\noindent \textbf{(P\_3,1.78)} na eṣaḥ śnam . \\
\noindent \textbf{(P\_3,1.78)} śnaḥ etat hrasvatvam . \\
\noindent \textbf{(P\_3,1.78)} yadi śnaḥ hrasvatvam svaraḥ na sidhyati . \\
\noindent \textbf{(P\_3,1.78)} vṛṇase . \\
\noindent \textbf{(P\_3,1.78)} adupadeśāt lasārvadhātukam anudāttam bhavati iti eṣaḥ svaraḥ na prāpnoti . \\
\noindent \textbf{(P\_3,1.78)} tasmāt śnam eṣaḥ . \\
\noindent \textbf{(P\_3,1.78)} yadi śnam snasoḥ allopaḥ iti lopaḥ prāpnoti . \\
\noindent \textbf{(P\_3,1.78)} upadhāyāḥ iti vartate . \\
\noindent \textbf{(P\_3,1.78)} anupadhātvāt na bhaviṣyati . \\
\noindent \textbf{(P\_3,1.78)} na saḥ śakhyaḥ upadhāyāḥ iti vijñātum . \\
\noindent \textbf{(P\_3,1.78)} iha hi doṣaḥ syāt . \\
\noindent \textbf{(P\_3,1.78)} aṅktaḥ añjanti . \\
\noindent \textbf{(P\_3,1.78)} tasmāt śnaḥ eva hrasvatvam . \\
\noindent \textbf{(P\_3,1.78)} svaraḥ katham . \\
\noindent \textbf{(P\_3,1.78)} bahulam pit sārvadhātukam chandasi . sārvadhātukasya bhalulam chandasi pittvam vaktavyam . \\
\noindent \textbf{(P\_3,1.78)} pitaḥ ca apittvam dṛśyate apitaḥ ca pittvam . \\
\noindent \textbf{(P\_3,1.78)} pitaḥ tāvat apittvam . \\
\noindent \textbf{(P\_3,1.78)} mātaram pramiṇīmi janitrīm . \\
\noindent \textbf{(P\_3,1.78)} apitaḥ pittvam . \\
\noindent \textbf{(P\_3,1.78)} śṛṇota grāvāṇaḥ . \\
\noindent \textbf{(P\_3,1.78)} tat tarhi hrasvatvam vaktavyam . \\
\noindent \textbf{(P\_3,1.78)} avaśyam chandasi hrasvatvam vaktavyam upagāyantu mām patnayaḥ garbhiṇayaḥ yuvatayaḥ iti evamartham . \\
\noindent \textbf{(P\_3,1.78)} viśeṣaṇāṛthaḥ tarhi . \\
\noindent \textbf{(P\_3,1.78)} kva viśeṣaṇāṛthena arthaḥ . \\
\noindent \textbf{(P\_3,1.78)} śnāt nalopaḥ iti . \\
\noindent \textbf{(P\_3,1.78)} nāt nalopaḥ iti ucyamāne yajñānām yatnānām iti atra api prasajyeta . \\
\noindent \textbf{(P\_3,1.78)} dīrghatve kṛte na bhaviṣyati . \\
\noindent \textbf{(P\_3,1.78)} idam iha sampradhāryam . \\
\noindent \textbf{(P\_3,1.78)} dīrghatvam kriyatām nalopaḥ iti . \\
\noindent \textbf{(P\_3,1.78)} kim atra kartavyam . \\
\noindent \textbf{(P\_3,1.78)} paratvāt na lopaḥ syāt . \\
\noindent \textbf{(P\_3,1.78)} tasmāt śakāraḥ kartavyaḥ . \\
\noindent \textbf{(P\_3,1.78)} atha kriyamāṇe api śakāre iha kasmāt na bhavati . \\
\noindent \textbf{(P\_3,1.78)} viśnānām praśnānām iti . \\
\noindent \textbf{(P\_3,1.78)} lakṣaṇapratipadoktayoḥ pratipadoktasya eva iti . \\
\noindent \textbf{(P\_3,1.79)} atha kimartham karoteḥ pṛthaggrahaṇam kriyate na tanādibhyaḥ iti eva ucyate . \\
\noindent \textbf{(P\_3,1.79)} anyāni tanotyādikāryāṇi mā bhūvan iti . \\
\noindent \textbf{(P\_3,1.79)} kāni . \\
\noindent \textbf{(P\_3,1.79)} anunāsikalopādīni . \\
\noindent \textbf{(P\_3,1.79)} daivaraktāḥ kiṃsukāḥ . \\
\noindent \textbf{(P\_3,1.79)} anunāsikābhāvāt eva anunāsikalopaḥ na bhaviṣyati . \\
\noindent \textbf{(P\_3,1.79)} idam tarhi tanādikāryam mā bhūt tanādibhyaḥ tathāsoḥ iti . \\
\noindent \textbf{(P\_3,1.79)} nanu ca bhavati eva atra hrasvāt aṅgāt iti . \\
\noindent \textbf{(P\_3,1.79)} tena eva yathā syāt . \\
\noindent \textbf{(P\_3,1.79)} anena mā bhūt iti . \\
\noindent \textbf{(P\_3,1.79)} kaḥ ca atra viśeṣaḥ tena vā sati anena vā . \\
\noindent \textbf{(P\_3,1.79)} tena sati sijlopasya asiddhatvāt ciṇvadbhāvaḥ siddhaḥ bhavati . \\
\noindent \textbf{(P\_3,1.79)} anena punaḥ sati ciṇvadbhāvaḥ na syāt . \\
\noindent \textbf{(P\_3,1.79)} anena api sati ciṇvadbhāvaḥ siddhaḥ . \\
\noindent \textbf{(P\_3,1.79)} katham . \\
\noindent \textbf{(P\_3,1.79)} vibhāṣā luk . \\
\noindent \textbf{(P\_3,1.79)} yadā na luk tadā tena lopaḥ . \\
\noindent \textbf{(P\_3,1.79)} tatra sijlopasya asiddhatvāt ciṇvadbhāvaḥ siddhaḥ bhavati . \\
\noindent \textbf{(P\_3,1.79)} tanāditvāt kṛñaḥ siddham sijlope ca na duṣyati . \\
\noindent \textbf{(P\_3,1.79)} ciṇvadbhāve atra doṣaḥ syāt . \\
\noindent \textbf{(P\_3,1.79)} saḥ api proktaḥ vibhāṣayā . \\
\noindent \textbf{(P\_3,1.80)} kva ayam akāraḥ śrūyate . \\
\noindent \textbf{(P\_3,1.80)} na kva cit śrūyate . \\
\noindent \textbf{(P\_3,1.80)} lopaḥ asya bhavati ataḥ lopaḥ ārdhadhātuke iti . \\
\noindent \textbf{(P\_3,1.80)} yadi na kva cit śrūyate kimartham atvam ucyate na lopaḥ eva ucyate . \\
\noindent \textbf{(P\_3,1.80)} na evam śakyam . \\
\noindent \textbf{(P\_3,1.80)} lope hi sati guṇaḥ prasajyeta . \\
\noindent \textbf{(P\_3,1.80)} nanu ca lope api sati na dhātulope ārdhadhātuke iti pratiṣedhaḥ bhaviṣyati . \\
\noindent \textbf{(P\_3,1.80)} ārdhadhātukanimitte lope saḥ pratiṣedhaḥ . \\
\noindent \textbf{(P\_3,1.80)} na ca eṣaḥ ārdhadhātukanimittaḥ lopaḥ . \\
\noindent \textbf{(P\_3,1.80)} api ca pratyākhyāyate saḥ yogaḥ . \\
\noindent \textbf{(P\_3,1.80)} tasmin pratyākhyāte guṇaḥ syāt eva . \\
\noindent \textbf{(P\_3,1.80)} tasmāt atvam vaktavyam . \\
\noindent \textbf{(P\_3,1.80)} atha kimartham numanuṣaktayoḥ grahaṇam kriyate na dhivikṛvyoḥ iti eva ucyate . \\
\noindent \textbf{(P\_3,1.80)} dhivikṛvyoḥ iti ucyamāne atve kṛte aniṣṭe deśe num prasajyeta . \\
\noindent \textbf{(P\_3,1.80)} idam iha sampradhāryam . \\
\noindent \textbf{(P\_3,1.80)} atvam kriyatām num iti . \\
\noindent \textbf{(P\_3,1.80)} kim atra kartavyam . \\
\noindent \textbf{(P\_3,1.80)} paratvāt numāgamaḥ . \\
\noindent \textbf{(P\_3,1.80)} antaraṅgam atvam . \\
\noindent \textbf{(P\_3,1.80)} kā antaraṅgatā . \\
\noindent \textbf{(P\_3,1.80)} pratyayotpattisanniyogena atvam ucyate . \\
\noindent \textbf{(P\_3,1.80)} utpannepratyaye prakṛtipratyayau āśritya aṅgasya numāgamaḥ . \\
\noindent \textbf{(P\_3,1.80)} num api antaraṅgaḥ . \\
\noindent \textbf{(P\_3,1.80)} katham . \\
\noindent \textbf{(P\_3,1.80)} vakṣyati etat numvidhau upadeśivadvacanam pratyayavidhyartham iti . \\
\noindent \textbf{(P\_3,1.80)} ubhayoḥ antaraṅgayoḥ paratvāt numāgamaḥ . \\
\noindent \textbf{(P\_3,1.80)} tasmāt dhivikṛvyoḥ iti vaktavyam . \\
\noindent \textbf{(P\_3,1.83)} kimarthaḥ śakāraḥ . \\
\noindent \textbf{(P\_3,1.83)} śit sārvadhātukam iti sārvadhātukasañjñā sārvadhātukam apit iti ṅittvam ṅiti iti pratiṣedhaḥ yathā syāt . \\
\noindent \textbf{(P\_3,1.83)} kuṣāṇa puṣāṇa iti . \\
\noindent \textbf{(P\_3,1.83)} ataḥ uttaram paṭhati . \\
\noindent \textbf{(P\_3,1.83)} śnāvikārasya śitkaraṇānarthakyam sthānivatvāt . \\
\noindent \textbf{(P\_3,1.83)} śnāvikārasya śitkaraṇam anarthakam . \\
\noindent \textbf{(P\_3,1.83)} kim kāraṇam . \\
\noindent \textbf{(P\_3,1.83)} sthānivatvāt . \\
\noindent \textbf{(P\_3,1.83)} śitaḥ ayam ādeśaḥ sthānivadbhāvāt śit bhaviṣyati . \\
\noindent \textbf{(P\_3,1.83)} arthavat tu jñāpakam sārvadhātukādeśe anubandhāsthānivattvasya . \\
\noindent \textbf{(P\_3,1.83)} arthavat tu śnāvikārasya śitkaraṇam . \\
\noindent \textbf{(P\_3,1.83)} kaḥ arthaḥ . \\
\noindent \textbf{(P\_3,1.83)} jñāpakārtham . \\
\noindent \textbf{(P\_3,1.83)} kim jñāpyam . \\
\noindent \textbf{(P\_3,1.83)} etat jñāpayati ācāryaḥ sārvadhātukādeśe anubandhāḥ na sthānivat bhavanti iti . \\
\noindent \textbf{(P\_3,1.83)} kim etasya jñapane prayojanam . \\
\noindent \textbf{(P\_3,1.83)} prayojanam hitātaṅoḥ apittvam . \\
\noindent \textbf{(P\_3,1.83)} heḥ pittvam na pratiṣedhyam . \\
\noindent \textbf{(P\_3,1.83)} pitaḥ ayam ādeśaḥ sthānivadbhāvāt pit syāt . \\
\noindent \textbf{(P\_3,1.83)} sārvadhātukādeśe anubandhāḥ na sthānivat bhavanti iti na ayam pit bhaviṣyati . \\
\noindent \textbf{(P\_3,1.83)} tātaṅi ca ṅakāraḥ na uccāryaḥ bhavati . \\
\noindent \textbf{(P\_3,1.83)} pitaḥ ayam ādeśaḥ sthānivadbhāvāt pit syāt . \\
\noindent \textbf{(P\_3,1.83)} sārvadhātukādeśe anubandhāḥ na sthānivat bhavanti iti na ayam pit bhaviṣyati . \\
\noindent \textbf{(P\_3,1.83)} tabādiṣu ca aṅittvam . \\
\noindent \textbf{(P\_3,1.83)} tabādiṣu ca aṅittvam prayojanam . \\
\noindent \textbf{(P\_3,1.83)} śṛṇota grāvāṇaḥ . \\
\noindent \textbf{(P\_3,1.83)} ṅitaḥ ime ādeśāḥ sthānivadbhāvāt ṅitaḥ syuḥ . \\
\noindent \textbf{(P\_3,1.83)} sārvadhātukādeśe anubandhāḥ na sthānivat bhavanti iti na ime ṅitaḥ bhavanti . \\
\noindent \textbf{(P\_3,1.83)} tasya doṣaḥ mipaḥ ādeśe pidabhāvaḥ . \\
\noindent \textbf{(P\_3,1.83)} tasya etasya lakṣaṇasya doṣaḥ mipaḥ ādeśe pitaḥ abhāvaḥ . \\
\noindent \textbf{(P\_3,1.83)} acinavam asunavam akaravam . \\
\noindent \textbf{(P\_3,1.83)} pitaḥ ayam ādeśaḥ sthānivadbhāvāt pit syāt . \\
\noindent \textbf{(P\_3,1.83)} sārvadhātukādeśe anubandhāḥ na sthānivat bhavanti iti na ayam pit syāt . \\
\noindent \textbf{(P\_3,1.83)} atyalpam idam ucyate . \\
\noindent \textbf{(P\_3,1.83)} tipsibmipām ādeśāḥ iti vaktavyam . \\
\noindent \textbf{(P\_3,1.83)} veda vettha . \\
\noindent \textbf{(P\_3,1.83)} videḥ vasoḥ śittvam . \\
\noindent \textbf{(P\_3,1.83)} videḥ uttarasya vasoḥ śittvam vaktavyam . \\
\noindent \textbf{(P\_3,1.83)} śitaḥ ayam ādeśaḥ sthānivadbhāvāt pit syāt . \\
\noindent \textbf{(P\_3,1.83)} sārvadhātukādeśe anubandhāḥ na sthānivat bhavanti iti na ayam śit syāt . \\
\noindent \textbf{(P\_3,1.83)} kitkaraṇāt vā siddham . atha vā avaśyam atra sāmānyagrahaṇāvighātārthaḥ kakāraḥ anubandhaḥ kartavyaḥ . \\
\noindent \textbf{(P\_3,1.83)} kva sāmānyagrahaṇāvighātārthena arthaḥ . \\
\noindent \textbf{(P\_3,1.83)} vasoḥ samprasāraṇam . \\
\noindent \textbf{(P\_3,1.83)} tena eva yatnena guṇaḥ na bhaviṣyati . \\
\noindent \textbf{(P\_3,1.83)} asya jñāpakasya santi doṣāḥ santi prayojanāni . \\
\noindent \textbf{(P\_3,1.83)} samāḥ doṣāḥ bhūyāṃsaḥ vā . \\
\noindent \textbf{(P\_3,1.83)} tasmāt na arthaḥ anena jñāpakena . \\
\noindent \textbf{(P\_3,1.83)} katham yāni prayojanāni . \\
\noindent \textbf{(P\_3,1.83)} tāni kriyante nyāse eva . \\
\noindent \textbf{(P\_3,1.83)} evam api bhavet pitkaraṇasāmarthyāt pitkṛtam syāt ṅitkaraṇasāmarthyāt ṅitkṛtam . \\
\noindent \textbf{(P\_3,1.83)} yat tu khalu piti ṅitkṛtam prāpnoti ṅiti ca pitkṛtam kena tat na syāt . \\
\noindent \textbf{(P\_3,1.83)} tasmāt vaktavyam pit na ṅidvat bhavati ṅit ca na pidvat bhavati iti . \\
\noindent \textbf{(P\_3,1.83)} na vaktavyam . \\
\noindent \textbf{(P\_3,1.83)} evam vakṣyāmi . \\
\noindent \textbf{(P\_3,1.83)} sārvadhātukam ṅit bhavati pit na . \\
\noindent \textbf{(P\_3,1.83)} evam tāvat pitaḥ ṅittvam pratiṣiddham . \\
\noindent \textbf{(P\_3,1.83)} tataḥ asaṃyogāt liṭ kit bhavati iti ṅit ca pit na bhavati . \\
\noindent \textbf{(P\_3,1.83)} evam ṅitaḥ pittvam pratiṣiddham . \\
\noindent \textbf{(P\_3,1.84)} śāyac chandasi sarvatra . \\
\noindent \textbf{(P\_3,1.84)} śāyac chandasi sarvatra iti vaktavyam . \\
\noindent \textbf{(P\_3,1.84)} kva sarvatra . \\
\noindent \textbf{(P\_3,1.84)} hau ca ahau ca . \\
\noindent \textbf{(P\_3,1.84)} kim prayojanam . \\
\noindent \textbf{(P\_3,1.84)} mahīaskabhāyat yaḥ askabhāyat udgṛbhāyata unmathāyata ityartham . \\
\noindent \textbf{(P\_3,1.85)} yogavibhāgaḥ kartavyaḥ . \\
\noindent \textbf{(P\_3,1.85)} vyatyayaḥ bhavati syādīnām iti . \\
\noindent \textbf{(P\_3,1.85)} āṇḍā śuṣṇasya bhṝdati . \\
\noindent \textbf{(P\_3,1.85)} bhinatti iti prāpte . \\
\noindent \textbf{(P\_3,1.85)} saḥ ca na marati . \\
\noindent \textbf{(P\_3,1.85)} miryate iti prāpte . \\
\noindent \textbf{(P\_3,1.85)} tataḥ bahulam. bahulam chandasi viṣaye sarve vidhayaḥ bhavanti iti . \\
\noindent \textbf{(P\_3,1.85)} supām vyatyayaḥ . \\
\noindent \textbf{(P\_3,1.85)} tiṅām vyatyayaḥ . \\
\noindent \textbf{(P\_3,1.85)} varṇavyatyayaḥ . \\
\noindent \textbf{(P\_3,1.85)} liṅgavyatyayaḥ . \\
\noindent \textbf{(P\_3,1.85)} kālavyatyayaḥ . \\
\noindent \textbf{(P\_3,1.85)} puruṣavyatyayaḥ . \\
\noindent \textbf{(P\_3,1.85)} ātmanepadavyatyayaḥ . \\
\noindent \textbf{(P\_3,1.85)} parasmaipadavyatyayaḥ . \\
\noindent \textbf{(P\_3,1.85)} supām vyatyayaḥ . \\
\noindent \textbf{(P\_3,1.85)} yuktā mātā āsīt dhuri dakṣiṇāyāḥ . \\
\noindent \textbf{(P\_3,1.85)} dakṣiṇāyām iti prāpte . \\
\noindent \textbf{(P\_3,1.85)} tiṅām vyatyayaḥ . \\
\noindent \textbf{(P\_3,1.85)} caṣālam ye aśvayūpāya takṣati . \\
\noindent \textbf{(P\_3,1.85)} takṣanti iti prāpte . \\
\noindent \textbf{(P\_3,1.85)} varṇavyatyayaḥ . \\
\noindent \textbf{(P\_3,1.85)} triṣṭubhaujaḥ śubhitam ugravīram . \\
\noindent \textbf{(P\_3,1.85)} suhitam iti prāpte . \\
\noindent \textbf{(P\_3,1.85)} liṅgavyatyayaḥ . \\
\noindent \textbf{(P\_3,1.85)} madhoḥ gṛhṇāti . \\
\noindent \textbf{(P\_3,1.85)} madhoḥ tṛptāḥ iva āsate . \\
\noindent \textbf{(P\_3,1.85)} madhunaḥ iti prāpte . \\
\noindent \textbf{(P\_3,1.85)} kālavyatyayaḥ . \\
\noindent \textbf{(P\_3,1.85)} śvaḥ agnīn ādhāsyamānena . \\
\noindent \textbf{(P\_3,1.85)} śvaḥ somena yakṣyamāṇena . \\
\noindent \textbf{(P\_3,1.85)} ādhātā yaṣṭā iti evam prāpte . \\
\noindent \textbf{(P\_3,1.85)} puruṣavyatyayaḥ . \\
\noindent \textbf{(P\_3,1.85)} adhā saḥ vīraiḥ daśabhiḥ viyūyāḥ . \\
\noindent \textbf{(P\_3,1.85)} viyūyāt iti prāpte . \\
\noindent \textbf{(P\_3,1.85)} ātmanepadavyatyayaḥ . \\
\noindent \textbf{(P\_3,1.85)} brahmnacāriṇam icchate . \\
\noindent \textbf{(P\_3,1.85)} icchati iti prāpte . \\
\noindent \textbf{(P\_3,1.85)} parasmaipadavyatyayaḥ . \\
\noindent \textbf{(P\_3,1.85)} pratīpam anyaḥ ūrmiḥ yudhyati . \\
\noindent \textbf{(P\_3,1.85)} yudhyate iti prāpte . \\
\noindent \textbf{(P\_3,1.85)} suptiṅupagrahaliṅganarāṇām kālhalacsvarakartṛyaṅām ca vyatyayam icchati śāstrakṛt eṣām . \\
\noindent \textbf{(P\_3,1.85)} saḥ api ca sidhyati bāhulakena . \\
\noindent \textbf{(P\_3,1.86)} ayam āśiṣi aṅ vidhīyate . \\
\noindent \textbf{(P\_3,1.86)} tasya kim prayojanam . \\
\noindent \textbf{(P\_3,1.86)} āśiṣi aṅaḥ prayojanam sthāgāgamivacividayaḥ . \\
\noindent \textbf{(P\_3,1.86)} sthā . \\
\noindent \textbf{(P\_3,1.86)} upa stheṣam vṛṣabham . \\
\noindent \textbf{(P\_3,1.86)} sthā . \\
\noindent \textbf{(P\_3,1.86)} gā . \\
\noindent \textbf{(P\_3,1.86)} añjasā satyam upa geṣam . \\
\noindent \textbf{(P\_3,1.86)} gā . \\
\noindent \textbf{(P\_3,1.86)} gami . \\
\noindent \textbf{(P\_3,1.86)} yajñena pratiṣṭhām gameyam . \\
\noindent \textbf{(P\_3,1.86)} gami . \\
\noindent \textbf{(P\_3,1.86)} vaci . \\
\noindent \textbf{(P\_3,1.86)} mantram vocema agnaye . \\
\noindent \textbf{(P\_3,1.86)} vaci . \\
\noindent \textbf{(P\_3,1.86)} vidi . \\
\noindent \textbf{(P\_3,1.86)} videyam enām manasi praviṣṭām . \\
\noindent \textbf{(P\_3,1.86)} śakiruhoḥ ca iti vaktavyam . \\
\noindent \textbf{(P\_3,1.86)} śakema tvā samidham . \\
\noindent \textbf{(P\_3,1.86)} asravantīm ā ruhema svastaye . \\
\noindent \textbf{(P\_3,1.86)} dṛśoḥ ak pitaram ca dṛśeyam mātaram ca . \\
\noindent \textbf{(P\_3,1.86)} dṛśoḥ ak vaktavyaḥ pitaram ca dṛśeyam mātaram ca iti evamartham . \\
\noindent \textbf{(P\_3,1.86)} iha upastheyāma iti āṭ api vaktavyaḥ . \\
\noindent \textbf{(P\_3,1.86)} na hi aṅā eva sidhyati . \\
\noindent \textbf{(P\_3,1.86)} na vaktavyaḥ . \\
\noindent \textbf{(P\_3,1.86)} sārvadhātukatvāt salopaḥ ārdhadhātukatvāt etvam . \\
\noindent \textbf{(P\_3,1.86)} dtatra ubhayaliṅgatvāt siddham . \\
\noindent \textbf{(P\_3,1.87.1)} vatkaraṇam kimartham . \\
\noindent \textbf{(P\_3,1.87.1)} svāśrayam api yathā syāt . \\
\noindent \textbf{(P\_3,1.87.1)} bhidyate kuśūlena iti . \\
\noindent \textbf{(P\_3,1.87.1)} akarmakāṇām bhāve laḥ bhavati iti laḥ yathā syāt . \\
\noindent \textbf{(P\_3,1.87.1)} karmaṇā iti kimartham . \\
\noindent \textbf{(P\_3,1.87.1)} karaṇādhikaraṇābhyām tulyakriyaḥ kartā yaḥ saḥ karmavat mā bhūt . \\
\noindent \textbf{(P\_3,1.87.1)} sādhu asiḥ chinatti . \\
\noindent \textbf{(P\_3,1.87.1)} sādhu sthālī pacati . \\
\noindent \textbf{(P\_3,1.87.1)} tulyakriyaḥ iti kimartham . \\
\noindent \textbf{(P\_3,1.87.1)} pacati odanam devadattaḥ . \\
\noindent \textbf{(P\_3,1.87.1)} tulyakriyaḥ iti ucyamāne api atra prāpnoti . \\
\noindent \textbf{(P\_3,1.87.1)} atra api hi karmaṇā tulyakriyaḥ kartā . \\
\noindent \textbf{(P\_3,1.87.1)} na tulyakriyagrahaṇena samānakriyatvam abhisambadhyate . \\
\noindent \textbf{(P\_3,1.87.1)} kim tarhi . \\
\noindent \textbf{(P\_3,1.87.1)} yasmin karmaṇi kartṛbhūte api tadvat kriya lakṣyate yathā karmaṇi saḥ karmaṇā tulyakriyaḥ kartā karmavat bhavati iti . \\
\noindent \textbf{(P\_3,1.87.2)} karmavat akarmakasya kartā . \\
\noindent \textbf{(P\_3,1.87.2)} akarmakasya kartā karmavat bhavati iti vaktavyam . \\
\noindent \textbf{(P\_3,1.87.2)} kim prayojanam . \\
\noindent \textbf{(P\_3,1.87.2)} sakarmakasya kartā karmavat mā bhūt iti . \\
\noindent \textbf{(P\_3,1.87.2)} bhidyamānaḥ kuśūlaḥ pātrāṇi bhinatti . \\
\noindent \textbf{(P\_3,1.87.2)} tathā karma dṛṣṭaḥ cet samānadhātau . \\
\noindent \textbf{(P\_3,1.87.2)} karma dṛṣṭaḥ cet samānadhātau iti vaktavyam . \\
\noindent \textbf{(P\_3,1.87.2)} iha mā bhūt . \\
\noindent \textbf{(P\_3,1.87.2)} pacati odanam devadattaḥ . \\
\noindent \textbf{(P\_3,1.87.2)} rādhyati odhanaḥ svayam eva . \\
\noindent \textbf{(P\_3,1.87.2)} tathā karmasthabhāvakānam karmasthakriyāṇām ca . \\
\noindent \textbf{(P\_3,1.87.2)} karmasthabhāvakānam karmasthakriyāṇām vā kartā karmavat bhavati iti vaktavyam . \\
\noindent \textbf{(P\_3,1.87.2)} kartṛsthabhāvakānām kartṛsthakriyāṇām vā kartā karmavat mā bhūt iti . \\
\noindent \textbf{(P\_3,1.87.2)} yat tāvat ucyate akarmakasya kartā karmavat bhavati iti vaktavyam iti . \\
\noindent \textbf{(P\_3,1.87.2)} na vaktavyam . \\
\noindent \textbf{(P\_3,1.87.2)} vakṣyati etat . \\
\noindent \textbf{(P\_3,1.87.2)} sakarmakāṇām pratiṣedhaḥ anyonyam āśliṣyataḥ iti . \\
\noindent \textbf{(P\_3,1.87.2)} yat api ucyate karma dṛṣṭaḥ cet samānadhātau iti vaktavyam iti . \\
\noindent \textbf{(P\_3,1.87.2)} na vaktavyam . \\
\noindent \textbf{(P\_3,1.87.2)} dhātoḥ iti vartate . \\
\noindent \textbf{(P\_3,1.87.2)} dhātoḥ karmaṇaḥ katurḥ ayam karmavadbhāvaḥ atidiśyate . \\
\noindent \textbf{(P\_3,1.87.2)} tatra sambandhāt etat gantavyam yasya dhātoḥ yat karma tasya cet kartā syāt iti . \\
\noindent \textbf{(P\_3,1.87.2)} tat yathā dhātoḥ karmaṇi aṇ bhavati iti . \\
\noindent \textbf{(P\_3,1.87.2)} tatra sambandhāt etat gamyate yasya dhātoḥ yat karma iti . \\
\noindent \textbf{(P\_3,1.87.2)} iha mā bhūt . \\
\noindent \textbf{(P\_3,1.87.2)} āhara kumbham karoti kaṭam iti . \\
\noindent \textbf{(P\_3,1.87.2)} yat api ucyate karmasthabhāvakānam karmasthakriyāṇām vā kartā karmavat bhavati iti vaktavyam . \\
\noindent \textbf{(P\_3,1.87.2)} kartṛsthabhāvakānām kartṛsthakriyāṇām vā kartā karmavat mā bhūt iti . \\
\noindent \textbf{(P\_3,1.87.2)} na vaktavyam . \\
\noindent \textbf{(P\_3,1.87.2)} karmasthayā kriyayā ayam kartāram upamimīte . \\
\noindent \textbf{(P\_3,1.87.2)} na ca kartṛsthabhāvakānām kartṛsthakriyāṇām vā karmaṇi kriyāyāḥ pravṛttiḥ asti . \\
\noindent \textbf{(P\_3,1.87.3)} kim punaḥ karmakartari karmāśrayam eva bhavati āhosvit kartrāśrayam api . \\
\noindent \textbf{(P\_3,1.87.3)} kim ca ataḥ . \\
\noindent \textbf{(P\_3,1.87.3)} yadi karmāśrayam eva caṅśapkṛdvidhayaḥ na sidhyanti . \\
\noindent \textbf{(P\_3,1.87.3)} caṅ . \\
\noindent \textbf{(P\_3,1.87.3)} acīkarata kaṭaḥ svayam eva . \\
\noindent \textbf{(P\_3,1.87.3)} śap . \\
\noindent \textbf{(P\_3,1.87.3)} namate daṇḍaḥ svayam eva . \\
\noindent \textbf{(P\_3,1.87.3)} kṛdvidhiḥ . \\
\noindent \textbf{(P\_3,1.87.3)} bhiduram kāṣṭham svayam eva . \\
\noindent \textbf{(P\_3,1.87.3)} atha kartrāśrayam api siddham etat bhavati . \\
\noindent \textbf{(P\_3,1.87.3)} kim tarhi iti . \\
\noindent \textbf{(P\_3,1.87.3)} ātmanepadaśabādividhipratiṣedhaḥ . \\
\noindent \textbf{(P\_3,1.87.3)} ātmanepadam vidheyam śabādīnām ca pratiṣedhaḥ vaktavyaḥ . \\
\noindent \textbf{(P\_3,1.87.3)} ubhayam kriyate nyāse eva . \\
\noindent \textbf{(P\_3,1.87.4)} kimartham punaḥ idam ucyate . \\
\noindent \textbf{(P\_3,1.87.4)} karmakartari kartṛtvam svātantryasya vivakṣitatvāt . karmakartari kartṛtvam asti . \\
\noindent \textbf{(P\_3,1.87.4)} kutaḥ . \\
\noindent \textbf{(P\_3,1.87.4)} svātantryasya vivakṣitatvāt . \\
\noindent \textbf{(P\_3,1.87.4)} svātantryeṇa eva atra kartā vivakṣitaḥ . \\
\noindent \textbf{(P\_3,1.87.4)} kim punaḥ sataḥ svātantryasya vivakṣā āhosvit vivakṣāmātram . \\
\noindent \textbf{(P\_3,1.87.4)} sataḥ iti āha . \\
\noindent \textbf{(P\_3,1.87.4)} katham jñāyate . \\
\noindent \textbf{(P\_3,1.87.4)} bhidyate kuśūlena iti . \\
\noindent \textbf{(P\_3,1.87.4)} na ca anyaḥ kartā dṛśyate kriyā ca upalabhyate . \\
\noindent \textbf{(P\_3,1.87.4)} kim ca bhoḥ vigrahavatā eva kriyāyāḥ kartrā bhavitavyam na punaḥ vātātapakālāḥ api kartāraḥ syuḥ . \\
\noindent \textbf{(P\_3,1.87.4)} bhavet siddham yadi vātātapakālānām anyatamaḥ kartā syāt . \\
\noindent \textbf{(P\_3,1.87.4)} yaḥ tu khalu nivāte nirabhivarṣe acirakālakṛtaḥ kuśūlaḥ bhidyate tasya na anyaḥ kartā bhavati anyat ataḥ kuśūlāt . \\
\noindent \textbf{(P\_3,1.87.4)} yadi api tāvat atra etat śakyate vaktum yatra anyaḥ kartā na asti iha tu katham na syāt lūyate kedāraḥ svayam eva iti yatra asu devadattaḥ dātrahastaḥ samantataḥ viparipatan dṛśyate . \\
\noindent \textbf{(P\_3,1.87.4)} atra api yā asau sukaratā nāma tasyāḥ na anyat kartā bhavati anyat ataḥ kedārāt . \\
\noindent \textbf{(P\_3,1.87.4)} asti prayojanam etat . \\
\noindent \textbf{(P\_3,1.87.4)} kim tarhi iti . \\
\noindent \textbf{(P\_3,1.87.4)} tatra lāntasya karmavadanudeśaḥ . \\
\noindent \textbf{(P\_3,1.87.4)} tatra lāntasya karmavadanudeśaḥ kartavyaḥ . \\
\noindent \textbf{(P\_3,1.87.4)} lāntasya kartā karmavat bhavati iti vaktavyam . \\
\noindent \textbf{(P\_3,1.87.4)} itarathā hi kṛtyaktakhalartheṣu pratiṣedhaḥ . \\
\noindent \textbf{(P\_3,1.87.4)} akriyamāṇe hi lagrahaṇe kṛtyaktakhalartheṣu pratiṣedhaḥ vaktavyaḥ syāt . \\
\noindent \textbf{(P\_3,1.87.4)} kṛtya . \\
\noindent \textbf{(P\_3,1.87.4)} bhettavyaḥ kuśūlaḥ iti karma . \\
\noindent \textbf{(P\_3,1.87.4)} saḥ yadā svātantryeṇa vivakṣitaḥ tadā asya karmavadbhāvaḥ syāt . \\
\noindent \textbf{(P\_3,1.87.4)} tasya pratiṣedhaḥ vaktavyaḥ . \\
\noindent \textbf{(P\_3,1.87.4)} tasmin pratiṣiddhi akarmakāṇām bhāve kṛtyā bhavanti iti bhāve yathā syāt . \\
\noindent \textbf{(P\_3,1.87.4)} bhettavyam kuśūlena iti . \\
\noindent \textbf{(P\_3,1.87.4)} kta . \\
\noindent \textbf{(P\_3,1.87.4)} bhinnaḥ kuśūlaḥ iti karma . \\
\noindent \textbf{(P\_3,1.87.4)} saḥ yadā svātantryeṇa vivakṣitaḥ tadā asya karmavadbhāvaḥ syāt . \\
\noindent \textbf{(P\_3,1.87.4)} tasya pratiṣedhaḥ vaktavyaḥ . \\
\noindent \textbf{(P\_3,1.87.4)} tasmin pratiṣiddhi akarmakāṇām bhāve ktaḥ bhavati iti bhāve ktaḥ yathā syāt . \\
\noindent \textbf{(P\_3,1.87.4)} bhinnam kuśūlena . \\
\noindent \textbf{(P\_3,1.87.4)} khalarthaḥ . \\
\noindent \textbf{(P\_3,1.87.4)} īṣadbhedyaḥ kuśūlaḥ iti karma . \\
\noindent \textbf{(P\_3,1.87.4)} saḥ yadā svātantryeṇa vivakṣitaḥ tadā asya karmavadbhāvaḥ syāt . \\
\noindent \textbf{(P\_3,1.87.4)} tasya pratiṣedhaḥ vaktavyaḥ . \\
\noindent \textbf{(P\_3,1.87.4)} tasmin pratiṣiddhe akarmakāṇām bhāve khal bhavati iti bhāve yathā syāt . \\
\noindent \textbf{(P\_3,1.87.4)} īṣadbhedyam kuśūlena iti . \\
\noindent \textbf{(P\_3,1.87.4)} tat tarhi lagrahaṇam kartavyam . \\
\noindent \textbf{(P\_3,1.87.4)} na kartavyam . \\
\noindent \textbf{(P\_3,1.87.4)} kriyate nyāse eva . \\
\noindent \textbf{(P\_3,1.87.4)} liṅi āśiṣi aṅ iti dvilakārakaḥ nirdeśaḥ . \\
\noindent \textbf{(P\_3,1.87.4)} siddham tu prākṛtakarmatvāt . \\
\noindent \textbf{(P\_3,1.87.4)} siddham etat . \\
\noindent \textbf{(P\_3,1.87.4)} katham . \\
\noindent \textbf{(P\_3,1.87.4)} prākṛtakarmatvāt . \\
\noindent \textbf{(P\_3,1.87.4)} prākṛtam eva etat karma yathā kaṭam karoti śakaṭam karoti . \\
\noindent \textbf{(P\_3,1.87.4)} katham punaḥ jñāyate prākṛtam eva etat karma iti . \\
\noindent \textbf{(P\_3,1.87.4)} ātmasaṃyoge akarmakartuḥ karmadarśanāt . \\
\noindent \textbf{(P\_3,1.87.4)} ātmasaṃyoge akarmakartuḥ karma dṛśyate . \\
\noindent \textbf{(P\_3,1.87.4)} kva . \\
\noindent \textbf{(P\_3,1.87.4)} hanti ātmānam . \\
\noindent \textbf{(P\_3,1.87.4)} hanyata ātmanā iti . \\
\noindent \textbf{(P\_3,1.87.4)} viṣamaḥ upanyāsaḥ . \\
\noindent \textbf{(P\_3,1.87.4)} hanti ātmānam iti karma dṛśyate . \\
\noindent \textbf{(P\_3,1.87.4)} kartā na dṛśyate . \\
\noindent \textbf{(P\_3,1.87.4)} ātmanā hanyate iti kartā dṛśyate . \\
\noindent \textbf{(P\_3,1.87.4)} karma na dṛśyate . \\
\noindent \textbf{(P\_3,1.87.4)} padalopaḥ ca . \\
\noindent \textbf{(P\_3,1.87.4)} padalopaḥ ca draṣṭavyaḥ . \\
\noindent \textbf{(P\_3,1.87.4)} hanti ātmānam ātmanā . \\
\noindent \textbf{(P\_3,1.87.4)} ātmanā hanyate ātmā iti . \\
\noindent \textbf{(P\_3,1.87.4)} kaḥ punaḥ ātmānam hanti kaḥ vā ātmanā hanyate . \\
\noindent \textbf{(P\_3,1.87.4)} dvau ātmānau antarātmā śarīrātmā ca . \\
\noindent \textbf{(P\_3,1.87.4)} antarātmā tat karma karoti yena śarīrātmā sukhaduḥke anubhavati . \\
\noindent \textbf{(P\_3,1.87.4)} śarīrātmā tat karma karoti yena antarātmā sukhaduḥke anubhavati . \\
\noindent \textbf{(P\_3,1.87.5)} sakarmakāṇām pratiṣedhaḥ anyonyam āśliṣyataḥ iti . \\
\noindent \textbf{(P\_3,1.87.5)} sakarmakāṇām pratiṣedhaḥ vaktavyaḥ . \\
\noindent \textbf{(P\_3,1.87.5)} kim prayojanam . \\
\noindent \textbf{(P\_3,1.87.5)} anyonyam āśliṣyataḥ . \\
\noindent \textbf{(P\_3,1.87.5)} anyonyam saṃspṛśataḥ . \\
\noindent \textbf{(P\_3,1.87.5)} anyonyam gṛhṇītaḥ iti . \\
\noindent \textbf{(P\_3,1.87.5)} tapeḥ vā sakarmakasya vacanam niyamārtham . \\
\noindent \textbf{(P\_3,1.87.5)} tapeḥ vā sakarmakasya vacanam niyamārtham bhaviṣyati . \\
\noindent \textbf{(P\_3,1.87.5)} tapeḥ eva sakarmakasya na anyasya sakarmakasya iti . \\
\noindent \textbf{(P\_3,1.87.5)} tasya tarhi anyakarmakasya api prāpnoti . \\
\noindent \textbf{(P\_3,1.87.5)} uttapati suvarṇam suvarṇakāraḥ . \\
\noindent \textbf{(P\_3,1.87.5)} uttapyamānam suvarṇam suvarṇakāram uttapati . \\
\noindent \textbf{(P\_3,1.87.5)} tasya ca tapaḥkarmakasya eva . \\
\noindent \textbf{(P\_3,1.87.5)} tasya ca tapaḥkarmakasya eva kartā karmavat bhavati na anyakarmakasya iti . \\
\noindent \textbf{(P\_3,1.87.5)} kim idam tapaḥ iti . \\
\noindent \textbf{(P\_3,1.87.5)} tapeḥ ayam auṇādikaḥ askāraḥ bhāvasādhanaḥ . \\
\noindent \textbf{(P\_3,1.87.5)} kaḥ prakṛtyarthaḥ kaḥ pratyayārthaḥ . \\
\noindent \textbf{(P\_3,1.87.5)} saḥ eva santapaḥ . \\
\noindent \textbf{(P\_3,1.87.5)} katham punaḥ saḥ eva nāma prakṛtyarthaḥ syāt saḥ eva pratyayārthaḥ . \\
\noindent \textbf{(P\_3,1.87.5)} sāmānyatapeḥ avayavatapiḥ karma bhavati . \\
\noindent \textbf{(P\_3,1.87.5)} tat yathā . \\
\noindent \textbf{(P\_3,1.87.5)} saḥ etān poṣān apuṣyat gopoṣam aśvapoṣam raipoṣam iti . \\
\noindent \textbf{(P\_3,1.87.5)} sāmānyapuṣeḥ avayavipuṣiḥ karma bhavati . \\
\noindent \textbf{(P\_3,1.87.5)} evam iha api sāmānyatapeḥ avayavatapiḥ karma bhavati . \\
\noindent \textbf{(P\_3,1.87.5)} duhipacyoḥ bahulam sakarmakayoḥ . \\
\noindent \textbf{(P\_3,1.87.5)} duhipacyoḥ sakarmakayoḥ kartā bahulam karmavat bhavati iti vaktavyam . \\
\noindent \textbf{(P\_3,1.87.5)} dugdhe gauḥ payaḥ . \\
\noindent \textbf{(P\_3,1.87.5)} tasmāt udumbaraḥ saḥ lohitam phalam pacyate . \\
\noindent \textbf{(P\_3,1.87.5)} bahulavacanam kimartham . \\
\noindent \textbf{(P\_3,1.87.5)} parasmaipadārtham . \\
\noindent \textbf{(P\_3,1.87.5)} yadi evam na arthaḥ bahulavacanena . \\
\noindent \textbf{(P\_3,1.87.5)} na hi parasmaipadam iṣyate . \\
\noindent \textbf{(P\_3,1.87.5)} sṛjiyujyoḥ śyan tu . \\
\noindent \textbf{(P\_3,1.87.5)} sṛjiyujyoḥ sakarmakayoḥ kartā bahulam karmavat bhavati iti vaktavyam . \\
\noindent \textbf{(P\_3,1.87.5)} śyan tu bhavati . \\
\noindent \textbf{(P\_3,1.87.5)} sṛjeḥ śraddhopapanne kartari karmavadbhāvaḥ vācyaḥ ciṇātmanepadārthaḥ . \\
\noindent \textbf{(P\_3,1.87.5)} sṛjyate mālām . \\
\noindent \textbf{(P\_3,1.87.5)} asarji mālām . \\
\noindent \textbf{(P\_3,1.87.5)} yajeḥ tu nyāyye karmakartari yakaḥ abhāvāya . \\
\noindent \textbf{(P\_3,1.87.5)} yujyate brahmacārī yogam . \\
\noindent \textbf{(P\_3,1.87.5)} karaṇena tulyakriyaḥ kartā bahulam . \\
\noindent \textbf{(P\_3,1.87.5)} karaṇena tulyakriyaḥ kartā bahulam karmavat bhavati iti vaktavyam . \\
\noindent \textbf{(P\_3,1.87.5)} parivārayanti kaṇṭakaiḥ vṛkṣam . \\
\noindent \textbf{(P\_3,1.87.5)} parivārayante kaṇṭakāḥ vṛkṣam iti . \\
\noindent \textbf{(P\_3,1.87.5)} sravatyādīnām pratiṣedhaḥ . \\
\noindent \textbf{(P\_3,1.87.5)} sravatyādīnām pratiṣedhaḥ vaktavyaḥ . \\
\noindent \textbf{(P\_3,1.87.5)} sravati kuṇḍikā udakam . \\
\noindent \textbf{(P\_3,1.87.5)} sravati kuṇḍikāyāḥ udakam . \\
\noindent \textbf{(P\_3,1.87.5)} sravanti valīkāni udakam . \\
\noindent \textbf{(P\_3,1.87.5)} sravati valīkebhyaḥ udakam iti . \\
\noindent \textbf{(P\_3,1.87.5)} saḥ tarhi pratiṣedhaḥ vaktavyaḥ . \\
\noindent \textbf{(P\_3,1.87.5)} na vaktavyaḥ . \\
\noindent \textbf{(P\_3,1.87.5)} tulyakriyaḥ iti ucyate . \\
\noindent \textbf{(P\_3,1.87.5)} kriyāntaram ca atra gamyate . \\
\noindent \textbf{(P\_3,1.87.5)} iha tāvat sravati kuṇḍikā udakam iti . \\
\noindent \textbf{(P\_3,1.87.5)} visṛjati iti gamyate . \\
\noindent \textbf{(P\_3,1.87.5)} sravati kuṇḍikāyāḥ udakam iti . \\
\noindent \textbf{(P\_3,1.87.5)} niṣkrāmati iti gamyate . \\
\noindent \textbf{(P\_3,1.87.5)} sravanti valīkāni udakam iti . \\
\noindent \textbf{(P\_3,1.87.5)} visṛjanti iti gamyate . \\
\noindent \textbf{(P\_3,1.87.5)} sravati valīkebhyaḥ udakam iti . \\
\noindent \textbf{(P\_3,1.87.5)} patati iti gamyate . \\
\noindent \textbf{(P\_3,1.87.5)} bhūṣākarmakiratisanām ca anyatra ātmanepadāt . \\
\noindent \textbf{(P\_3,1.87.5)} bhūṣākarmakiratisanām ca pratiṣedhaḥ vaktavyaḥ anyatra ātmanepadāt . \\
\noindent \textbf{(P\_3,1.87.5)} bhūṣayate kanyā svayam eva . \\
\noindent \textbf{(P\_3,1.87.5)} abubhūṣata kanyā svayam eva . \\
\noindent \textbf{(P\_3,1.87.5)} maṇḍayate kanyā svayam eva . \\
\noindent \textbf{(P\_3,1.87.5)} amamaṇḍata kanyā svayam eva . \\
\noindent \textbf{(P\_3,1.87.5)} kirati . \\
\noindent \textbf{(P\_3,1.87.5)} avakirate hastī svayam eva . \\
\noindent \textbf{(P\_3,1.87.5)} avākīrṣṭa hastī svayam eva . \\
\noindent \textbf{(P\_3,1.87.5)} san . \\
\noindent \textbf{(P\_3,1.87.5)} cikīrṣate kaṭaḥ svayam eva . \\
\noindent \textbf{(P\_3,1.87.5)} acikīrṣiṣṭa kaṭaḥ svayam eva . \\
\noindent \textbf{(P\_3,1.89)} yakciṇoḥ pratiṣedhe hetumaṇṇiśribrūñām upasaṅkhyānam . \\
\noindent \textbf{(P\_3,1.89)} yakciṇoḥ pratiṣedhe hetumaṇṇiśribrūñām upasaṅkhyānam kartavyam . \\
\noindent \textbf{(P\_3,1.89)} ṇi . \\
\noindent \textbf{(P\_3,1.89)} kārayate kaṭaḥ svayam eva . \\
\noindent \textbf{(P\_3,1.89)} acīkarata kaṭaḥ svayam eva . \\
\noindent \textbf{(P\_3,1.89)} ṇi . \\
\noindent \textbf{(P\_3,1.89)} śri . \\
\noindent \textbf{(P\_3,1.89)} ucchrayate daṇḍaḥ svayam eva . \\
\noindent \textbf{(P\_3,1.89)} udaśiśriyata daṇḍaḥ svayam eva . \\
\noindent \textbf{(P\_3,1.89)} śri . \\
\noindent \textbf{(P\_3,1.89)} brūñ . \\
\noindent \textbf{(P\_3,1.89)} brūte kathā svayam eva . \\
\noindent \textbf{(P\_3,1.89)} avocata kathā svayam eva . \\
\noindent \textbf{(P\_3,1.89)} bhāradvājīyāḥ paṭhanti . \\
\noindent \textbf{(P\_3,1.89)} yakciṇoḥ pratiṣedhe ṇiśrigranthibrūñātmanepadākarmakāṇām upasaṅkhyānam iti . \\
\noindent \textbf{(P\_3,1.90)} kuṣirajoḥ śyanvidhāne sārvadhātukavacanam .kuṣirajoḥ śyanvidhāne sārvadhātukagrahaṇam kartavyam . \\
\noindent \textbf{(P\_3,1.90)} avacane hi liṅliṭoḥ pratiṣedhaḥ . \\
\noindent \textbf{(P\_3,1.90)} akriyamāṇe hi sārvadhātukagrahaṇe liṅliṭoḥ pratiṣedhaḥ vaktavyaḥ syāt . \\
\noindent \textbf{(P\_3,1.90)} cukuṣe pādaḥ svayam eva . \\
\noindent \textbf{(P\_3,1.90)} rarañje vastram svayam eva . \\
\noindent \textbf{(P\_3,1.90)} koṣiṣīṣṭa pādaḥ svayam eva . \\
\noindent \textbf{(P\_3,1.90)} raṅkṣīṣṭa vastram svayam eva . \\
\noindent \textbf{(P\_3,1.90)} kriyamāṇe api sārvadhātukagrahaṇe iha prāpnoti . \\
\noindent \textbf{(P\_3,1.90)} kati iha kuṣṇāṇāḥ pādāḥ . \\
\noindent \textbf{(P\_3,1.90)} śyanā ca syādīnām bādhanam prāpnoti . \\
\noindent \textbf{(P\_3,1.90)} koṣiṣyate pādaḥ svayam eva . \\
\noindent \textbf{(P\_3,1.90)} raṅkṣyate vastram svayam eva . \\
\noindent \textbf{(P\_3,1.90)} akoṣi pādaḥ svayam eva . \\
\noindent \textbf{(P\_3,1.90)} arañji vastram svayam eva . \\
\noindent \textbf{(P\_3,1.90)} yat tāvat ucyate sārvadhātukagrahaṇam kartavyam iti . \\
\noindent \textbf{(P\_3,1.90)} prakṛtam anuvartate . \\
\noindent \textbf{(P\_3,1.90)} kva prakṛtam . \\
\noindent \textbf{(P\_3,1.90)} sārvadhātuke yak iti . \\
\noindent \textbf{(P\_3,1.90)} yadi tat anuvartate pūrvasmin yoge kim samuccayaḥ . \\
\noindent \textbf{(P\_3,1.90)} le ca sārvadhātuke ca iti . \\
\noindent \textbf{(P\_3,1.90)} āhosvit lagrahaṇam sārvadhātukaviśeṣaṇam . \\
\noindent \textbf{(P\_3,1.90)} kim ca ataḥ . \\
\noindent \textbf{(P\_3,1.90)} yadi samuccayaḥ kati iha bhindānāḥ kuśūlāḥ iti atra api prāpnoti . \\
\noindent \textbf{(P\_3,1.90)} atha lagrahaṇam sārvadhātukaviśeṣaṇam liṅliṭoḥ na sidhyati . \\
\noindent \textbf{(P\_3,1.90)} bibhide kuśūlaḥ svayam eva . \\
\noindent \textbf{(P\_3,1.90)} bhitsīṣṭa kuśūlaḥ svayam eva iti . \\
\noindent \textbf{(P\_3,1.90)} astu lagrahaṇam sārvadhātukaviśeṣaṇam . \\
\noindent \textbf{(P\_3,1.90)} nanu ca uktam liṅliṭoḥ na sidhyati iti . \\
\noindent \textbf{(P\_3,1.90)} liṅliḍgrahaṇam api prakṛtam anuvartate . \\
\noindent \textbf{(P\_3,1.90)} kva prakṛtam . \\
\noindent \textbf{(P\_3,1.90)} kās pratyayāt ām amantre liṭi liṅi āśiṣi āṅ iti . \\
\noindent \textbf{(P\_3,1.90)} evam ca kṛtvā saḥ api adoṣaḥ bhavati yat uktam kati iha kuṣṇāṇāḥ pādāḥ iti prāpnoti . \\
\noindent \textbf{(P\_3,1.90)} atra api laviśiṣṭam sārvadhātukagrahaṇam anuvartate . \\
\noindent \textbf{(P\_3,1.90)} yat api ucyate śyanā ca syādīnām bādhanam prāpnoti . \\
\noindent \textbf{(P\_3,1.90)} yakpratiṣedhasambandhena śyanaṃ vakṣyāmi . \\
\noindent \textbf{(P\_3,1.90)} na duhasnunamām yakciṇau . \\
\noindent \textbf{(P\_3,1.90)} tataḥ kuṣirajoḥ prācām yakciṇau na bhavataḥ . \\
\noindent \textbf{(P\_3,1.90)} tataḥ śyan parasmaipadam ca iti . \\
\noindent \textbf{(P\_3,1.90)} yathā eva tarhi yakaḥ viṣaye śyan bhavati evam ciṇaḥ api viṣaye prāpnoti . \\
\noindent \textbf{(P\_3,1.90)} akoṣi pādaḥ svayam eva . \\
\noindent \textbf{(P\_3,1.90)} arañji vastram svayam eva iti . \\
\noindent \textbf{(P\_3,1.90)} evam tarhi dvitīyaḥ yogavibhāgaḥ kariṣyate . \\
\noindent \textbf{(P\_3,1.90)} na duhasnunamām ciṇ bhavati . \\
\noindent \textbf{(P\_3,1.90)} tataḥ yak . \\
\noindent \textbf{(P\_3,1.90)} yak ca na bhavati duhasnunamām . \\
\noindent \textbf{(P\_3,1.90)} tataḥ kuṣirajoḥ prācām yak na bhavati . \\
\noindent \textbf{(P\_3,1.90)} tataḥ śyan parasmaipadam ca . \\
\noindent \textbf{(P\_3,1.90)} atha vā anuvṛttiḥ kariṣyate . \\
\noindent \textbf{(P\_3,1.90)} syatāsī lṛluṭoḥ . \\
\noindent \textbf{(P\_3,1.90)} cli luṅi . \\
\noindent \textbf{(P\_3,1.90)} cleḥ sic bhavati . \\
\noindent \textbf{(P\_3,1.90)} kartari śap syatāsī lṛluṭoḥ cli luṅi cleḥ sic bhavati . \\
\noindent \textbf{(P\_3,1.90)} kuṣirajoḥ prācām śyan parasmaipadam ca syatāsī lṛluṭoḥ cli luṅi cleḥ sic bhavati iti . \\
\noindent \textbf{(P\_3,1.90)} atha vā antaraṅgāḥ syādayaḥ . \\
\noindent \textbf{(P\_3,1.90)} kā antaraṅgatā . \\
\noindent \textbf{(P\_3,1.90)} lakārāvasthāyām eva syādayaḥ . \\
\noindent \textbf{(P\_3,1.90)} sārvadhātuke śyan . \\
\noindent \textbf{(P\_3,1.91.1)} ā kutaḥ ayam dhātvadhikāraḥ . \\
\noindent \textbf{(P\_3,1.91.1)} kim prāk lādeśāt āhosvit ā tṛtīyādhyāyaparisamāpteḥ . \\
\noindent \textbf{(P\_3,1.91.1)} dhātuvadhikāraḥ prāk lādeśāt . \\
\noindent \textbf{(P\_3,1.91.1)} prāk lādeśāt dhātvadhikāraḥ . \\
\noindent \textbf{(P\_3,1.91.1)} lādeśe hi vyavahitatvāt aprasiddhiḥ . \\
\noindent \textbf{(P\_3,1.91.1)} anuvartamāne hi lādeśe dhātvadhikāre vyavahitavtā aprasiddhiḥ syāt . \\
\noindent \textbf{(P\_3,1.91.1)} kim ca syāt . \\
\noindent \textbf{(P\_3,1.91.1)} ādye yoge na vyavāye tiṅaḥ syuḥ . ādye yoge vikaraṇaiḥ vyavahitatvāt tiṅaḥ na syuḥ . \\
\noindent \textbf{(P\_3,1.91.1)} pacati paṭhati . \\
\noindent \textbf{(P\_3,1.91.1)} idam iha sampradhāryam . \\
\noindent \textbf{(P\_3,1.91.1)} vikaraṇāḥ kriyantām āḍeśāḥ iti . \\
\noindent \textbf{(P\_3,1.91.1)} kim atra kartavyam . \\
\noindent \textbf{(P\_3,1.91.1)} paratvāt ādeśāḥ . \\
\noindent \textbf{(P\_3,1.91.1)} nityāḥ vikaraṇāḥ . \\
\noindent \textbf{(P\_3,1.91.1)} kṛteṣu ādeśeṣu prāpnuvanti akṛteṣu api prāpnuvanti . \\
\noindent \textbf{(P\_3,1.91.1)} nityatvāt vikaraṇeṣu kṛteṣu vikaraṇaiḥ vyavahitatvāt ādeśāḥ na prāpnuvanti . \\
\noindent \textbf{(P\_3,1.91.1)} anavakāśaḥ tarhi ādeśāḥ . \\
\noindent \textbf{(P\_3,1.91.1)} sāvakāśāḥ ādeśāḥ . \\
\noindent \textbf{(P\_3,1.91.1)} kaḥ avakāśaḥ . \\
\noindent \textbf{(P\_3,1.91.1)} ye ete lugvikaraṇāḥ śluvikaraṇāḥ ca liṅliṭau ca . \\
\noindent \textbf{(P\_3,1.91.1)} na syāt etvam ṭeḥ ṭitām yat vidhatte . \\
\noindent \textbf{(P\_3,1.91.1)} yat ca ṭitsañjñānām etvam vidhatte tat ca vikaraṇaiḥ vyavahitatvāt na syāt . \\
\noindent \textbf{(P\_3,1.91.1)} eśaḥ śittvam . ekāraḥ ca śit kartavyaḥ . \\
\noindent \textbf{(P\_3,1.91.1)} kim prayojanam . \\
\noindent \textbf{(P\_3,1.91.1)} śit sarvasya iti sarvādeśaḥ yathā syāt . \\
\noindent \textbf{(P\_3,1.91.1)} akriyamāṇe hi śakāre tasmāt iti uttarasya ādeḥ iti takārasya etve kṛte dvayoḥ ekārayoḥ śravaṇam prasajyeta . \\
\noindent \textbf{(P\_3,1.91.1)} nivṛtte punaḥ lādeśe dhātvadhikāre alaḥ antyasya vidhayaḥ bhavanti iti ekārasya ekārvacanane prayojanam na asti iti kṛtvā antareṇa api śakāram sarvādeśaḥ bhaviṣyati . \\
\noindent \textbf{(P\_3,1.91.1)} yat ca loṭaḥ vidhatte . \\
\noindent \textbf{(P\_3,1.91.1)} tat ca vikaraṇaiḥ vyavahitatvāt na syāt . \\
\noindent \textbf{(P\_3,1.91.1)} kim punaḥ tat . \\
\noindent \textbf{(P\_3,1.91.1)} loṭaḥ laṅvat eḥ uḥ seḥ hi apit ca vā chandasi iti . \\
\noindent \textbf{(P\_3,1.91.1)} yat ca api uktam laṅliṅoḥ tat ca na syāt . \\
\noindent \textbf{(P\_3,1.91.1)} kim punaḥ tat . \\
\noindent \textbf{(P\_3,1.91.1)} nityam ṅitaḥ itaḥ ca tasthasthamipām tāmtamtāmaḥ liṅaḥ sīyuṭ yāsuṭ parasmaipadeṣu udāttaḥ ṅit ca iti . \\
\noindent \textbf{(P\_3,1.91.1)} tasmāt prāk lādeśāt dhātvadhikāraḥ . \\
\noindent \textbf{(P\_3,1.91.1)} yadi prāk lādeśāt dhātvadhikāraḥ akāraḥ śit kartavyaḥ . \\
\noindent \textbf{(P\_3,1.91.1)} kim prayojanam . \\
\noindent \textbf{(P\_3,1.91.1)} śit sarvasya iti sarvādeśaḥ yathā syāt . \\
\noindent \textbf{(P\_3,1.91.1)} anuvartamāne punaḥ lādeśe dhātvadhikāre tasmāt iti uttarasya ādeḥ iti thakārasya atve kṛte dvayoḥ akārayoḥ pararūpeṇa siddham rūpam syāt . \\
\noindent \textbf{(P\_3,1.91.1)} peca yūyam . \\
\noindent \textbf{(P\_3,1.91.1)} cakra yūyam iti . \\
\noindent \textbf{(P\_3,1.91.1)} nanu ca nivṛtte api lādeśe dhātvadhikāre alaḥ antyasya vidhayaḥ bhavanti iti akārasya akārvacanane prayojanam na asti iti kṛtvā antareṇa api śakāram sarvādeśaḥ bhaviṣyati . \\
\noindent \textbf{(P\_3,1.91.1)} asti anyat akārasya akāravacane prayojanam . \\
\noindent \textbf{(P\_3,1.91.1)} kim . \\
\noindent \textbf{(P\_3,1.91.1)} vakṣyati etat tat akārasya akāravacanam samasaṅkhyārtham iti . \\
\noindent \textbf{(P\_3,1.91.1)} ārdhadhātukasañjñāyām dhātugrahaṇam kartavyam dhātoḥ parasya ārdhadhātukasañjñā yathā syāt . \\
\noindent \textbf{(P\_3,1.91.1)} iha mā bhūt : vṛkṣtvam vṛkṣatā iti . \\
\noindent \textbf{(P\_3,1.91.1)} tasmāt lādeśe dhātvadhikāraḥ anuvartyaḥ . \\
\noindent \textbf{(P\_3,1.91.1)} nanu ca uktam ādye yoge na vyavāye tiṅaḥ syuḥ iti . \\
\noindent \textbf{(P\_3,1.91.1)} na eṣaḥ doṣaḥ . \\
\noindent \textbf{(P\_3,1.91.1)} ānupūrvyāt siddham etat . \\
\noindent \textbf{(P\_3,1.91.1)} na atra akṛteṣu ādeśeṣu vikaraṇāḥ prāpnuvanti . \\
\noindent \textbf{(P\_3,1.91.1)} kim kāraṇam . \\
\noindent \textbf{(P\_3,1.91.1)} sārvadhātuke vikaraṇāḥ ucyante . \\
\noindent \textbf{(P\_3,1.91.1)} na ca akṛteṣu ādeśeṣu sārvadhātukatvam bhavati . \\
\noindent \textbf{(P\_3,1.91.1)} ye tarhi na etasmin viśeṣe vidhīyante . \\
\noindent \textbf{(P\_3,1.91.1)} ke punaḥ te . \\
\noindent \textbf{(P\_3,1.91.1)} syādayaḥ . \\
\noindent \textbf{(P\_3,1.91.1)} tatra api vihitaviśeṣaṇam dhātugrahaṇam . \\
\noindent \textbf{(P\_3,1.91.1)} dhātoḥ vihitasya lasya iti . \\
\noindent \textbf{(P\_3,1.91.1)} yadi evam vindati iti ṇalādayaḥ prāpnuvanti . \\
\noindent \textbf{(P\_3,1.91.1)} dhātunā atra vihitam viśeṣayiṣyāmaḥ vidinā ca ānataryam . \\
\noindent \textbf{(P\_3,1.91.1)} dhātoḥ vihitasya lasya videḥ anantarasya iti . \\
\noindent \textbf{(P\_3,1.91.1)} iha tarhi ajakṣiṣyan ajāgairṣyan iti abhyastāt jheḥ jus bhavati iti jusbhāvaḥ prāpnoti . \\
\noindent \textbf{(P\_3,1.91.1)} atra api dhātunā vihitam viśeṣayiṣyāmaḥ abhyastena ānantaryam . \\
\noindent \textbf{(P\_3,1.91.1)} dhātoḥ vihitasya abhastāt anantarasya iti . \\
\noindent \textbf{(P\_3,1.91.1)} ātaḥ iti atra katham viśeṣayiṣyasi . \\
\noindent \textbf{(P\_3,1.91.1)} yadi tāvat dhātugrahaṇam vihitaviśeṣaṇam ākāragrahaṇam ānantaryaviśeṣaṇam alunan iti atra api prāpnoti . \\
\noindent \textbf{(P\_3,1.91.1)} atha ākāragrahaṇam vihitaviśeṣaṇam dhātugrahaṇam ānantaryaviśeṣaṇam apiban iti atra api prāpnoti . \\
\noindent \textbf{(P\_3,1.91.1)} astu tarhi dhātugrahaṇam vihitaviśeṣaṇam ākāragrahaṇam ānantaryaviśeṣaṇam . \\
\noindent \textbf{(P\_3,1.91.1)} nanu ca uktam alunan iti atra api prāpnoti iti . \\
\noindent \textbf{(P\_3,1.91.1)} na eṣaḥ doṣaḥ . \\
\noindent \textbf{(P\_3,1.91.1)} lope kṛte na bhaviṣyati . \\
\noindent \textbf{(P\_3,1.91.1)} na atra lopaḥ prāpnoti . \\
\noindent \textbf{(P\_3,1.91.1)} kim kāraṇam . \\
\noindent \textbf{(P\_3,1.91.1)} ītvena bādhyate . \\
\noindent \textbf{(P\_3,1.91.1)} na atra ītvam prāpnoti . \\
\noindent \textbf{(P\_3,1.91.1)} kim kāraṇam . \\
\noindent \textbf{(P\_3,1.91.1)} antibhāvena bādhyate . \\
\noindent \textbf{(P\_3,1.91.1)} na atra antibhāvaḥ prāpnoti . \\
\noindent \textbf{(P\_3,1.91.1)} kim kāraṇam . \\
\noindent \textbf{(P\_3,1.91.1)} jusbhāvena bādhyate . \\
\noindent \textbf{(P\_3,1.91.1)} na atra jusbhāvaḥ prāpnoti . \\
\noindent \textbf{(P\_3,1.91.1)} kim kāraṇam . \\
\noindent \textbf{(P\_3,1.91.1)} lopena bādhyate . \\
\noindent \textbf{(P\_3,1.91.1)} lopaḥ ītvena ītvam antibhāvena antibhāvaḥ jusbhāvena jusbhāvaḥ lopena iti cakrakam avyavasthā prasajyeta . \\
\noindent \textbf{(P\_3,1.91.1)} na asti cakrakaprasaṅgaḥ . \\
\noindent \textbf{(P\_3,1.91.1)} na hi avyavasthākāriṇā śāstreṇa bhavitavyam . \\
\noindent \textbf{(P\_3,1.91.1)} śāstreṇa nāma vyavasthākāriṇā bhavitavyam . \\
\noindent \textbf{(P\_3,1.91.1)} na ca atra halādinā muhūrtam api śakyam avasthātum . \\
\noindent \textbf{(P\_3,1.91.1)} tāvati eva antibhāvena bhavitavyam . \\
\noindent \textbf{(P\_3,1.91.1)} antibhāve kṛte lopaḥ . \\
\noindent \textbf{(P\_3,1.91.1)} lopena vyavasthā bhaviṣyati . \\
\noindent \textbf{(P\_3,1.91.1)} yat api ucyate eśaḥ śittvam iti . \\
\noindent \textbf{(P\_3,1.91.1)} kriyate nyāse eva . \\
\noindent \textbf{(P\_3,1.91.2)} kāni punaḥ asya yogasya prayojanāni . \\
\noindent \textbf{(P\_3,1.91.2)} prayojanam prātipadikapratiṣedhaḥ . \\
\noindent \textbf{(P\_3,1.91.2)} prātipadikapratiṣedhaḥ prayojanam . \\
\noindent \textbf{(P\_3,1.91.2)} dhātoḥ tavyādayaḥ yathā syuḥ . \\
\noindent \textbf{(P\_3,1.91.2)} prātipadikāt mā bhūvan iti . \\
\noindent \textbf{(P\_3,1.91.2)} na etat asti prayojanam . \\
\noindent \textbf{(P\_3,1.91.2)} sādhane tāvyādayaḥ vidhīyante sādhanam ca kriyāyāḥ . \\
\noindent \textbf{(P\_3,1.91.2)} kriyābhāvāt sādhanābhāvaḥ . \\
\noindent \textbf{(P\_3,1.91.2)} sādhanābhāvāt asati api dhātvadhikāre prātipadikāt tavyādayaḥ na bhaviṣyanti . \\
\noindent \textbf{(P\_3,1.91.2)} svapādiṣu . \\
\noindent \textbf{(P\_3,1.91.2)} svapādiṣu tarhi prayojanam . \\
\noindent \textbf{(P\_3,1.91.2)} svapiti . \\
\noindent \textbf{(P\_3,1.91.2)} supati iti mā bhūt . \\
\noindent \textbf{(P\_3,1.91.2)} aṅgasañjñā ca . \\
\noindent \textbf{(P\_3,1.91.2)} aṅgasañjñā ca prayojanam . \\
\noindent \textbf{(P\_3,1.91.2)} yasmāt pratyayavidhiḥ tadādi pratyaye aṅgam iti dhātoḥ aṅgasañjñā siddhā bhavati . \\
\noindent \textbf{(P\_3,1.91.2)} kṛtsañjñā ca . \\
\noindent \textbf{(P\_3,1.91.2)} kṛtsañjñā ca prayojanam . \\
\noindent \textbf{(P\_3,1.91.2)} dhātuvihitasya pratyayasya kṛtsañjñā siddhā bhavati . \\
\noindent \textbf{(P\_3,1.91.2)} upapadasañjñā ca. upapadasañjñā ca prayojanam . \\
\noindent \textbf{(P\_3,1.91.2)} tatra etasmin dhātvadhikāre saptamīnirdiṣṭam upapadasñjñam bhavati iti upapadasañjñā siddhā bhavati . \\
\noindent \textbf{(P\_3,1.91.2)} kṛdupapadasañjñe tāvan na prayojanam . \\
\noindent \textbf{(P\_3,1.91.2)} adhikārāt api ete siddhe . \\
\noindent \textbf{(P\_3,1.91.2)} svapādiṣu tarhi aṅgasañjñā ca prayojanam . \\
\noindent \textbf{(P\_3,1.91.2)} dhātugrahaṇam anarthakam yaṅvidhau dhātvadhikārāt . \\
\noindent \textbf{(P\_3,1.91.2)} dhātugrahaṇam anarthakam . \\
\noindent \textbf{(P\_3,1.91.2)} kim kāraṇam . \\
\noindent \textbf{(P\_3,1.91.2)} yaṅvidhau dhātvadhikārāt . \\
\noindent \textbf{(P\_3,1.91.2)} yaṅvidhau dhātugrahaṇam prakṛtam anuvartate . \\
\noindent \textbf{(P\_3,1.91.2)} tat ca avaśyam anuvartyam . \\
\noindent \textbf{(P\_3,1.91.2)} anadhikāre hi aṅgasañjñābhāvaḥ . \\
\noindent \textbf{(P\_3,1.91.2)} anadhikāre hi sati aṅgasañjñāyāḥ abhāvaḥ syāt . \\
\noindent \textbf{(P\_3,1.91.2)} kariṣyati hariṣyati iti . \\
\noindent \textbf{(P\_3,1.91.2)} yad tat anuvartate cūrṇacurādibhyaḥ ṇic bhavati dhātoḥ ca iti dhātumātrāt ṇic prāpnoti . \\
\noindent \textbf{(P\_3,1.91.2)} hetumadvacanam tu jñāpakam anyatrābhāvasya . \\
\noindent \textbf{(P\_3,1.91.2)} yat ayam hetumati ca iti āha tat jñāpayati ācāryaḥ na dhātumātrāt ṇic bhavati iti . \\
\noindent \textbf{(P\_3,1.91.2)} iha tarhi kaṇḍvādibhyaḥ yak bhavati dhātoḥ ca iti dhātumātrāt yak prāpnoti . \\
\noindent \textbf{(P\_3,1.91.2)} kaṇḍvādiṣu ca vyapadeśivadvacanāt . \\
\noindent \textbf{(P\_3,1.91.2)} yat ayam kaṇḍvādibhyaḥ yak bhavati iti āha tat jñāpayati ācāryaḥ na dhātumātrāt yak bhavati iti . \\
\noindent \textbf{(P\_3,1.91.2)} atha vā kaṇḍvādibhyaḥ dhātugrahaṇena abhisambhantsyāmaḥ . \\
\noindent \textbf{(P\_3,1.91.2)} kaṇḍvādibhyaḥ dhātubhyaḥ iti . \\
\noindent \textbf{(P\_3,1.92.1)} sthagrahaṇam kimartham . \\
\noindent \textbf{(P\_3,1.92.1)} tatra upapadam saptamī iti iyati ucyamāne yatra eva saptamī śrūyate tatra eva syāt : stamberamaḥ karṇejapaḥ . \\
\noindent \textbf{(P\_3,1.92.1)} yatra vā etena śabdena nirdeśaḥ kriyate . \\
\noindent \textbf{(P\_3,1.92.1)} saptamyām janeḥ ḍaḥ iti . \\
\noindent \textbf{(P\_3,1.92.1)} iha na syāt . \\
\noindent \textbf{(P\_3,1.92.1)} kumbhakāraḥ nagarakāraḥ . \\
\noindent \textbf{(P\_3,1.92.1)} sthagrahaṇe punaḥ kriyamāṇe yatra ca saptamī śrūyate ya ca na śrūyate yatra ca etena śabdena nirdeśaḥ kriyate yatra ca anyena saptamīsthamātre siddham bhavati . \\
\noindent \textbf{(P\_3,1.92.1)} atha tatragrahaṇam kimartham . \\
\noindent \textbf{(P\_3,1.92.1)} tatragrahaṇam viṣayārtham . \\
\noindent \textbf{(P\_3,1.92.1)} viṣayaḥ pratinirdiśyate . \\
\noindent \textbf{(P\_3,1.92.1)} tatra etasmin dhātvadhikāre yat saptamīnirdiṣṭam tat upapadasañjñam bhavati iti upapadasañjñā siddhā bhavati . \\
\noindent \textbf{(P\_3,1.92.2)} upapadasañjñāyām samarthavacanam . upapadasañjñāyām samarthagrahaṇam kartavyam . \\
\noindent \textbf{(P\_3,1.92.2)} samartham upapadam pratyayasya iti vaktavyam . \\
\noindent \textbf{(P\_3,1.92.2)} iha mā bhūt . \\
\noindent \textbf{(P\_3,1.92.2)} āhara kumbham . \\
\noindent \textbf{(P\_3,1.92.2)} karoti kaṭam iti . \\
\noindent \textbf{(P\_3,1.92.2)} kriyamāṇe ca api samarthagrahaṇe mahāntam kumbham karoti iti atra api prāpnoti . \\
\noindent \textbf{(P\_3,1.92.2)} na vā bhavitavyam mahākumbhakāraḥ iti . \\
\noindent \textbf{(P\_3,1.92.2)} bhavaitavyam yadā etat vākyam bhavati . \\
\noindent \textbf{(P\_3,1.92.2)} mahān kumbhaḥ mahākumbhaḥ mahākumbham karoti iti mahākumbhakāraḥ . \\
\noindent \textbf{(P\_3,1.92.2)} yadā tu etat vākyam bhavati mahāntam kumbham karoti iti tadā na bhavitavyam . \\
\noindent \textbf{(P\_3,1.92.2)} tadā ca prāpnoti . \\
\noindent \textbf{(P\_3,1.92.2)} tadā mā bhūt iti . \\
\noindent \textbf{(P\_3,1.92.2)} yat tāvat ucyate samarthagrahaṇam kartavyam iti . \\
\noindent \textbf{(P\_3,1.92.2)} na kartavyam . \\
\noindent \textbf{(P\_3,1.92.2)} dhātoḥ iti vartate . \\
\noindent \textbf{(P\_3,1.92.2)} dhātoḥ karmaṇi aṇ bhavati . \\
\noindent \textbf{(P\_3,1.92.2)} tatra sambandhāt etat gantavyam . \\
\noindent \textbf{(P\_3,1.92.2)} yasya dhātoḥ yat karma iti . \\
\noindent \textbf{(P\_3,1.92.2)} yat api ucyate kriyamāṇe ca api samarthagrahaṇe mahāntam kumbham karoti iti atra api prāpnoti iti . \\
\noindent \textbf{(P\_3,1.92.2)} upapadam iti mahatīiham sañjñā kriyate . \\
\noindent \textbf{(P\_3,1.92.2)} sañjñā ca nāma yataḥ na laghīyaḥ . \\
\noindent \textbf{(P\_3,1.92.2)} kutaḥ etat . \\
\noindent \textbf{(P\_3,1.92.2)} laghvartham hi sañjñākaraṇam . \\
\noindent \textbf{(P\_3,1.92.2)} tatra mahatyāḥ sañjñāyāḥ karaṇe etat prayojanam anvarthasañjñā yathā vijñāyeta : upoccāri padam upapadam . \\
\noindent \textbf{(P\_3,1.92.2)} yat ca atra upoccāri na tat padam yat ca padam na tat upoccāri . \\
\noindent \textbf{(P\_3,1.92.2)} yāvatā ca idānīm padagandhaḥ asti padavidhiḥ ayam bhavati . \\
\noindent \textbf{(P\_3,1.92.2)} padavidhiḥ ca samarthānām bhavati . \\
\noindent \textbf{(P\_3,1.92.2)} tatra asāmārthyān na bhaviṣyati . \\
\noindent \textbf{(P\_3,1.92.2)} atha cvyante upapade kim aṇā bhavitavyam . \\
\noindent \textbf{(P\_3,1.92.2)} akumbham kumbham karoti kumbhīkaroti mṛdam iti . \\
\noindent \textbf{(P\_3,1.92.2)} na bhavitavyam . \\
\noindent \textbf{(P\_3,1.92.2)} kim kāraṇam . \\
\noindent \textbf{(P\_3,1.92.2)} prakṛtivivakṣāyām cviḥ vidhīyate . \\
\noindent \textbf{(P\_3,1.92.2)} tat sāpekṣam . \\
\noindent \textbf{(P\_3,1.92.2)} sāpekṣam ca asamartham bhavati . \\
\noindent \textbf{(P\_3,1.92.2)} na tarhi idānīm idam bhavati : icchāmi aham kāśakaṭīkāram iti . \\
\noindent \textbf{(P\_3,1.92.2)} iṣṭam eva etat gonardīyasya . \\
\noindent \textbf{(P\_3,1.92.2)} nimittopādanam ca . \\
\noindent \textbf{(P\_3,1.92.2)} nimittopādanam ca kartavyam . \\
\noindent \textbf{(P\_3,1.92.2)} nimittam upapadam pratyayasya iti vaktavyam . \\
\noindent \textbf{(P\_3,1.92.2)} anupādāne hi anupapade pratyayaprasaṅgaḥ . \\
\noindent \textbf{(P\_3,1.92.2)} akriyamāṇe hi nimittopādāne anupapade api prasajyeta . \\
\noindent \textbf{(P\_3,1.92.2)} nirdeśaḥ idānīm kimarthaḥ syāt . \\
\noindent \textbf{(P\_3,1.92.2)} nirdeśaḥ sañjñākaraṇārthaḥ . \\
\noindent \textbf{(P\_3,1.92.2)} yadā upapade pratyayaḥ tadā upapadasañjñām vakṣyāmi iti . \\
\noindent \textbf{(P\_3,1.92.2)} tat tarhi nimittopādanam kartavyam . \\
\noindent \textbf{(P\_3,1.92.2)} na kartavyam . \\
\noindent \textbf{(P\_3,1.92.2)} tatravacanam upapadasanniyogārtham . \\
\noindent \textbf{(P\_3,1.92.2)} tatravacanam kriyate . \\
\noindent \textbf{(P\_3,1.92.2)} tat upapadasanniyogārtham bhaviṣyati . \\
\noindent \textbf{(P\_3,1.92.2)} karmaṇi aṇ vidhīyate tatra cet pratyayaḥ bhavati iti . \\
\noindent \textbf{(P\_3,1.92.2)} nanu ca anyat tatragrahaṇasya prayojanam uktam . \\
\noindent \textbf{(P\_3,1.92.2)} kim . \\
\noindent \textbf{(P\_3,1.92.2)} tatragrahaṇam viṣayārtham iti . \\
\noindent \textbf{(P\_3,1.92.2)} adhikārāt api etat siddham . \\
\noindent \textbf{(P\_3,1.93)} atiṅ iti kimartham . \\
\noindent \textbf{(P\_3,1.93)} pacati karoti . \\
\noindent \textbf{(P\_3,1.93)} atiṅ iti śakyam akartum . \\
\noindent \textbf{(P\_3,1.93)} kasmāt na bhavati pacati karoti iti . \\
\noindent \textbf{(P\_3,1.93)} dhātoḥ parasya kṛtsañjñā . \\
\noindent \textbf{(P\_3,1.93)} prāk ca lādeśāt dhātvadhikāraḥ . \\
\noindent \textbf{(P\_3,1.93)} evam api sthānivadbhāvāt kṛtsañjña prāpnoti . \\
\noindent \textbf{(P\_3,1.93)} yathā atiṅ iti ucyamāne yāvatā sthānivadbhāvaḥ katham eva etat sidhyati . \\
\noindent \textbf{(P\_3,1.93)} pratiṣedhavacanasāmarthyāt . \\
\noindent \textbf{(P\_3,1.93)} atha vā tiṅbhāvinaḥ lakārasya kṛtsañjñāpratiṣedhaḥ . \\
\noindent \textbf{(P\_3,1.93)} kim ca syāt yati atra kṛtsañjñā syāt . \\
\noindent \textbf{(P\_3,1.93)} kṛtprātipadikam iti prātipadikasañjñā syāt . \\
\noindent \textbf{(P\_3,1.93)} prātipadikāt iti svādyutpattiḥ prasajyeta . \\
\noindent \textbf{(P\_3,1.93)} na eṣaḥ doṣaḥ . \\
\noindent \textbf{(P\_3,1.93)} ekatvādiṣu artheṣu svādayaḥ vidhīyante . \\
\noindent \textbf{(P\_3,1.93)} te ca atra tiṅoktāḥ ekatvādayaḥ iti kṛtvā uktārthatvāt na bhaviṣyanti . \\
\noindent \textbf{(P\_3,1.93)} ṭābādayaḥ tarhi tiṅantāt mā bhūvan iti . \\
\noindent \textbf{(P\_3,1.93)} striyām ṭābādayaḥ vidhīyante . \\
\noindent \textbf{(P\_3,1.93)} na ca tiṅantasya strītvena yogaḥ asti . \\
\noindent \textbf{(P\_3,1.93)} aṇādayaḥ tarhi tiṅantāt mā bhūvan iti . \\
\noindent \textbf{(P\_3,1.93)} apatyādiṣu artheṣu aṇādayaḥ vidhīyante . \\
\noindent \textbf{(P\_3,1.93)} na ca tiṅantasya apatyādibhiḥ yogaḥ asti . \\
\noindent \textbf{(P\_3,1.93)} atha api katham cit yogaḥ syāt . \\
\noindent \textbf{(P\_3,1.93)} evam api na doṣaḥ . \\
\noindent \textbf{(P\_3,1.93)} ācāryapravṛttiḥ jñāpayati na tiṅantāt aṇādayaḥ bhavanti iti yat ayam kva cit taddhitavidhau tiṅgrahaṇam karoti . \\
\noindent \textbf{(P\_3,1.93)} atiśāyane tamabiṣṭhanau tiṅaḥ ca iti . \\
\noindent \textbf{(P\_3,1.93)} iha tarhi pacati paṭhati iti . \\
\noindent \textbf{(P\_3,1.93)} hrasvasya piti kṛti tuk bhavati iti tuk prāpnoti . \\
\noindent \textbf{(P\_3,1.93)} dhātoḥ iti vartate . \\
\noindent \textbf{(P\_3,1.93)} evam api cikīrṣati iti atra prāpnoti . \\
\noindent \textbf{(P\_3,1.93)} atra api śapā vyavadhānam . \\
\noindent \textbf{(P\_3,1.93)} ekādeśe kṛte na asti vyavadhānam . \\
\noindent \textbf{(P\_3,1.93)} ekādeśaḥ pūrvavidhau sthānivat bhavati iti sthānivadbhāvāt vyavadhānam eva iti . \\
\noindent \textbf{(P\_3,1.94.1)} katham idam vijñāyate . \\
\noindent \textbf{(P\_3,1.94.1)} striyām abhidheyāyām vā asrūpaḥ na bhavati iti āhosvit strīpratyayeṣu iti . \\
\noindent \textbf{(P\_3,1.94.1)} kim ca ataḥ . \\
\noindent \textbf{(P\_3,1.94.1)} yadi striyām abhidheyāyām iti lavyā lavitavyā atra vā asarūpaḥ na prāpnoti . \\
\noindent \textbf{(P\_3,1.94.1)} atha vijñāyate strīpratyayeṣu iti vyāvakrośī vayatikruṣṭiḥ iti na sidhyati . \\
\noindent \textbf{(P\_3,1.94.1)} evam tarhi na evam vijñāyate striyām abhidheyāyām na api strīpratyayeṣu iti . \\
\noindent \textbf{(P\_3,1.94.1)} katham tarhi strīgrahaṇam svarayiṣyate . \\
\noindent \textbf{(P\_3,1.94.1)} tatra svaritena adhikāragatiḥ bhavati iti striyām iti adhikṛtya ye pratayāḥ vihitāḥ teṣām pratiṣedhaḥ vijñāsyate . \\
\noindent \textbf{(P\_3,1.94.2)} kimartham punaḥ idam ucyate . \\
\noindent \textbf{(P\_3,1.94.2)} asarūpasya vāvacanam utsargasya bādhakaviṣaye anivṛttyartham . \\
\noindent \textbf{(P\_3,1.94.2)} asarūpasya vāvacanam kriyate utsargasya bādhakaviṣaye anivṛttiḥ yathā syāt . \\
\noindent \textbf{(P\_3,1.94.2)} tavyattavyānīyaraḥ utsargāḥ . \\
\noindent \textbf{(P\_3,1.94.2)} teṣām ajantāt yat apavādaḥ . \\
\noindent \textbf{(P\_3,1.94.2)} ceyam , cetavyam iti api yathā syāt . \\
\noindent \textbf{(P\_3,1.94.2)} na etat asti prayojanam . \\
\noindent \textbf{(P\_3,1.94.2)} ajantāt yat vidhīyate . \\
\noindent \textbf{(P\_3,1.94.2)} halantāt ṇyat vidhīyate . \\
\noindent \textbf{(P\_3,1.94.2)} etāvantaḥ ca dhātavaḥ yat uta ajantāḥ halantāḥ ca . \\
\noindent \textbf{(P\_3,1.94.2)} ucyante ca tavyādayaḥ . \\
\noindent \textbf{(P\_3,1.94.2)} te vacanāt bhaviṣyanti . \\
\noindent \textbf{(P\_3,1.94.2)} evam tarhi ṇvultṛcau utsargau . \\
\noindent \textbf{(P\_3,1.94.2)} tayoḥ pacādibhyaḥ ac apavādaḥ . \\
\noindent \textbf{(P\_3,1.94.2)} pacati iti pacaḥ . \\
\noindent \textbf{(P\_3,1.94.2)} paktā pācakaḥ iti api yathā syāt . \\
\noindent \textbf{(P\_3,1.94.2)} etat api na asti prayojanam . \\
\noindent \textbf{(P\_3,1.94.2)} vakṣyati etat . \\
\noindent \textbf{(P\_3,1.94.2)} ac api sarvadhātubhyaḥ vaktavyaḥ iti . \\
\noindent \textbf{(P\_3,1.94.2)} evam tarhi ṇvultṛjacaḥ utsargāḥ teṣām igupadhāt kaḥ apavādaḥ . \\
\noindent \textbf{(P\_3,1.94.2)} vikṣipaḥ vilikhaḥ . \\
\noindent \textbf{(P\_3,1.94.2)} vikṣeptā vikṣepakaḥ iti api yathā syāt . \\
\noindent \textbf{(P\_3,1.94.2)} asti prayojanam etat . \\
\noindent \textbf{(P\_3,1.94.2)} kim tarhi iti . \\
\noindent \textbf{(P\_3,1.94.2)} tatra utpattivāprasaṅgaḥ yathā taddhite . \\
\noindent \textbf{(P\_3,1.94.2)} tatra utpattiḥ vibhāṣā prāpnoti yathā taddhite . \\
\noindent \textbf{(P\_3,1.94.2)} astu . \\
\noindent \textbf{(P\_3,1.94.2)} yadā vikṣipaḥ vilikhaḥ iti etat na tadā vikṣeptā vikṣepakaḥ iti etat bhaviṣyati . \\
\noindent \textbf{(P\_3,1.94.2)} yadi etat labhyeta kṛtam syāt . \\
\noindent \textbf{(P\_3,1.94.2)} tat tu na labhyam . \\
\noindent \textbf{(P\_3,1.94.2)} kim kāraṇam . \\
\noindent \textbf{(P\_3,1.94.2)} yathā taddhite iti ucyate . \\
\noindent \textbf{(P\_3,1.94.2)} tadditeṣu ca sarvam eva utsargāpavādam vibhāṣā . \\
\noindent \textbf{(P\_3,1.94.2)} utpadyate vā na vā . \\
\noindent \textbf{(P\_3,1.94.2)} siddham tu asarūpasya bādhakasya vāvacanāt . \\
\noindent \textbf{(P\_3,1.94.2)} siddham etat . \\
\noindent \textbf{(P\_3,1.94.2)} katham . \\
\noindent \textbf{(P\_3,1.94.2)} asarūpasya bādhakasya vāvacanāt . \\
\noindent \textbf{(P\_3,1.94.2)} asarūpaḥ bādhakaḥ vā bādhakaḥ bhavati iti vaktavyam . \\
\noindent \textbf{(P\_3,1.94.2)} sidhyati . \\
\noindent \textbf{(P\_3,1.94.2)} sūtram tarhi bhidyate . \\
\noindent \textbf{(P\_3,1.94.2)} yathānyāsam eva astu . \\
\noindent \textbf{(P\_3,1.94.2)} nanu ca uktam tatra utpattivāprasaṅgaḥ yathā taddhite iti . \\
\noindent \textbf{(P\_3,1.94.2)} na eṣaḥ doṣaḥ . \\
\noindent \textbf{(P\_3,1.94.2)} asti kāraṇam yena taddhite vibhāṣā utpattiḥ bhavati . \\
\noindent \textbf{(P\_3,1.94.2)} kim kāraṇam . \\
\noindent \textbf{(P\_3,1.94.2)} prakṛtiḥ tatra prakṛtyarthe vartate . \\
\noindent \textbf{(P\_3,1.94.2)} anyena śabdena pratyayārthaḥ abhidhīyate . \\
\noindent \textbf{(P\_3,1.94.2)} iha punaḥ na kevalā prakṛtiḥ prakṛtyarthe vartate na ca anyaḥ śabdaḥ asti yaḥ tam artham abhidadhīta iti kṛtvā anutpattiḥ na bhaviṣyati . \\
\noindent \textbf{(P\_3,1.94.2)} atha vā samayaḥ kṛtaḥ . \\
\noindent \textbf{(P\_3,1.94.2)} na kevalā prakṛtiḥ prayoktavyā na ca kevalaḥ pratyayaḥ iti . \\
\noindent \textbf{(P\_3,1.94.2)} etasmāt samayāt anutpattiḥ na bhaviṣyati . \\
\noindent \textbf{(P\_3,1.94.2)} nanu ca yaḥ eva tasya samayasya kartā saḥ eva idam api āha . \\
\noindent \textbf{(P\_3,1.94.2)} yadi asau tatra pramāṇam iha api pramāṇam bhavitum arhati . \\
\noindent \textbf{(P\_3,1.94.2)} pramāṇam asau tatra ca iha ca . \\
\noindent \textbf{(P\_3,1.94.2)} sāmarthyam tu iha draṣṭavyam prayoge . \\
\noindent \textbf{(P\_3,1.94.2)} na ca anutpattau sāmarthyam asti . \\
\noindent \textbf{(P\_3,1.94.2)} tena anutpattiḥ na bhaviṣyati . \\
\noindent \textbf{(P\_3,1.94.2)} katham tarhi taddhiteṣu anutpattau sāmarthyam bhavati . \\
\noindent \textbf{(P\_3,1.94.2)} anyena pratyayena sāmarthyam . \\
\noindent \textbf{(P\_3,1.94.2)} kena . \\
\noindent \textbf{(P\_3,1.94.2)} ṣaṣṭhyā . \\
\noindent \textbf{(P\_3,1.94.2)} atha vā rūpavattām āśritya vāvidhiḥ ucyate . \\
\noindent \textbf{(P\_3,1.94.2)} na ca anutpattiḥ rūpavatī . \\
\noindent \textbf{(P\_3,1.94.2)} tena anutpattiḥ na bhaviṣyati . \\
\noindent \textbf{(P\_3,1.94.2)} evam api kutaḥ etat apavādaḥ vibhāṣā bhaviṣyati na punaḥ utsargaḥ iti . \\
\noindent \textbf{(P\_3,1.94.2)} na ca eva asti viśeṣaḥ yat apavādaḥ vibhāṣā syāt utsargaḥ vā . \\
\noindent \textbf{(P\_3,1.94.2)} api ca sāpekṣaḥ ayam nirdeśaḥ kriyate vā asarūpaḥ iti . \\
\noindent \textbf{(P\_3,1.94.2)} na ca utsargavelāyām kim cit apekṣyam asti . \\
\noindent \textbf{(P\_3,1.94.2)} apavādavelāyām punaḥ utsargaḥ apekṣyate . \\
\noindent \textbf{(P\_3,1.94.2)} tena yaḥ rūpavān anyapūrvakaḥ bādhakaḥ prāpnoti saḥ vā bādhakaḥ bhaviṣyati . \\
\noindent \textbf{(P\_3,1.94.2)} kaḥ punaḥ asau . \\
\noindent \textbf{(P\_3,1.94.2)} apavādaḥ . \\
\noindent \textbf{(P\_3,1.94.2)} yadi yaḥ rūpavān anyapūrvakaḥ bādhakaḥ prāpnoti saḥ vā bādhakaḥ bhavati iti ucyate kvibādiṣu samāveśaḥ na prāpnoti . \\
\noindent \textbf{(P\_3,1.94.2)} grāmaṇīḥ grāmaṇāyaḥ iti . \\
\noindent \textbf{(P\_3,1.94.2)} na hi ete rūpavantaḥ . \\
\noindent \textbf{(P\_3,1.94.2)} ete api rūpavantaḥ . \\
\noindent \textbf{(P\_3,1.94.2)} kasyām avasthāyām . \\
\noindent \textbf{(P\_3,1.94.2)} upadeśāvasthāyām . \\
\noindent \textbf{(P\_3,1.94.2)} yadi evam anubandhabhinneṣu vibhāṣāprasaṅgaḥ . \\
\noindent \textbf{(P\_3,1.94.2)} anubandhabhinneṣu vibhāṣā prāpnoti . \\
\noindent \textbf{(P\_3,1.94.2)} karmaṇi aṇ ātaḥ anupasarge kaḥ iti kaviṣaye aṇ api prāpnoti . \\
\noindent \textbf{(P\_3,1.94.2)} siddham anubandhasya anekāntatvāt . \\
\noindent \textbf{(P\_3,1.94.2)} siddham etat . \\
\noindent \textbf{(P\_3,1.94.2)} katham . \\
\noindent \textbf{(P\_3,1.94.2)} anubandhasya anekāntatvāt . \\
\noindent \textbf{(P\_3,1.94.2)} anekāntāḥ anubandhāḥ . \\
\noindent \textbf{(P\_3,1.94.2)} atha vā prayoge asarūpāṇām vāvidhiḥ nyāyyaḥ . \\
\noindent \textbf{(P\_3,1.94.2)} prayoge cet lādeśeṣu pratiṣedhaḥ . \\
\noindent \textbf{(P\_3,1.94.2)} prayoge cet lādeśeṣu pratiṣedhaḥ vaktavyaḥ . \\
\noindent \textbf{(P\_3,1.94.2)} hyaḥ apacat iti atra luṅ api prāpnoti . \\
\noindent \textbf{(P\_3,1.94.2)} śvaḥ paktā iti atra lṛṭ api prāpnoti . \\
\noindent \textbf{(P\_3,1.94.2)} na eṣaḥ doṣaḥ . \\
\noindent \textbf{(P\_3,1.94.2)} ācāryapravṛttiḥ jñāpayati na lādeśeṣu vā asarūpaḥ bhavati iti yat ayam haśaśvatoḥ laṅ ca iti āha . \\
\noindent \textbf{(P\_3,1.94.2)} atha vā prayoge asarūpāṇām vāvidhau na sarvam iṣṭam saṅgṛhītam iti kṛtvā dvitīyaḥ prayogaḥ upāsyate . \\
\noindent \textbf{(P\_3,1.94.2)} kaḥ asau . \\
\noindent \textbf{(P\_3,1.94.2)} upadeśaḥ nāma . \\
\noindent \textbf{(P\_3,1.94.2)} upadeśe ca ete sarūpāḥ . \\
\noindent \textbf{(P\_3,1.94.2)} nanu ca uktam anubandhabhinneṣu vibhāṣāprasaṅgaḥ iti . \\
\noindent \textbf{(P\_3,1.94.2)} parihṛtam etat . \\
\noindent \textbf{(P\_3,1.94.2)} katham . \\
\noindent \textbf{(P\_3,1.94.2)} siddham anubandhasya anekāntatvāt . \\
\noindent \textbf{(P\_3,1.94.2)} atha ekānte doṣaḥ eva . \\
\noindent \textbf{(P\_3,1.94.2)} ekānte ca na doṣaḥ . \\
\noindent \textbf{(P\_3,1.94.2)} ācāryapravṛttiḥ jñāpayati na anubandhakṛtam asārūpyam bhavati iti yat ayam dadādidadhātyoḥ vibhāṣā śam śāsti . \\
\noindent \textbf{(P\_3,1.94.2)} atha vā asarūpaḥ bādhakaḥ vā bādhakaḥ bhavati iti ucyate . \\
\noindent \textbf{(P\_3,1.94.2)} apavādaḥ nāma anubandhabhinnaḥ vā bhavati rūpānyatvena vā . \\
\noindent \textbf{(P\_3,1.94.2)} tena anena avaśyam kim cit tyājyam kim cit tu saṅgrahītavyam . \\
\noindent \textbf{(P\_3,1.94.2)} tat yat anubandhakṛtam asārūpyam tat na āśrayiṣyāmaḥ yat tu rūpānyatvena asārūpyam tat āśrayiṣyāmaḥ . \\
\noindent \textbf{(P\_3,1.94.2)} atha vā asarūpaḥ bādhakaḥ vā bādhakaḥ bhavati iti ucyate sarvaḥ ca asarūpaḥ . \\
\noindent \textbf{(P\_3,1.94.2)} tatra prakarṣagatiḥ vijñāsyate : sādhīyaḥ yaḥ asarūpaḥ iti . \\
\noindent \textbf{(P\_3,1.94.2)} kaḥ ca sādhīyaḥ . \\
\noindent \textbf{(P\_3,1.94.2)} yaḥ prayoge ca prāk ca prayogāt . \\
\noindent \textbf{(P\_3,1.94.2)} atha vā asarūpaḥ bādhakaḥ vā bādhakaḥ bhavati iti ucyate . \\
\noindent \textbf{(P\_3,1.94.2)} na ca evam kaḥ cit api sarūpaḥ . \\
\noindent \textbf{(P\_3,1.94.2)} te evam vijñāsyāmaḥ : kvat cit ye asarūpāḥ . \\
\noindent \textbf{(P\_3,1.94.2)} anubandhabhinnāḥ ca prayoge sarūpāḥ . \\
\noindent \textbf{(P\_3,1.94.3)} atha katham idam vijñāyate astriyām iti . \\
\noindent \textbf{(P\_3,1.94.3)} kim striyām na bhavati āhosvit prāk striyāḥ bhavati iti . \\
\noindent \textbf{(P\_3,1.94.3)} kaḥ ca atra viśeṣaḥ . \\
\noindent \textbf{(P\_3,1.94.3)} striyām pratiṣedhe ktalyuṭtumunkhalartheṣu vibhāṣāprasaṅgaḥ . \\
\noindent \textbf{(P\_3,1.94.3)} striyām pratiṣedhe ktalyuṭtumunkhalartheṣu vibhāṣā prāpnoti . \\
\noindent \textbf{(P\_3,1.94.3)} kta . \\
\noindent \textbf{(P\_3,1.94.3)} hasitam chātrasys śobhanam . \\
\noindent \textbf{(P\_3,1.94.3)} ghañ api prāpnoti . \\
\noindent \textbf{(P\_3,1.94.3)} lyuṭ . \\
\noindent \textbf{(P\_3,1.94.3)} hasanam chātrasys śobhanam . \\
\noindent \textbf{(P\_3,1.94.3)} ghañ api prāpnoti . \\
\noindent \textbf{(P\_3,1.94.3)} tumun . \\
\noindent \textbf{(P\_3,1.94.3)} icchati bhoktum . \\
\noindent \textbf{(P\_3,1.94.3)} liṅloṭau api prāpnutaḥ . \\
\noindent \textbf{(P\_3,1.94.3)} khalarthaḥ . \\
\noindent \textbf{(P\_3,1.94.3)} īṣatpānaḥ somaḥ bhavatā . \\
\noindent \textbf{(P\_3,1.94.3)} khal api prāpnoti . \\
\noindent \textbf{(P\_3,1.94.3)} evam tarhi striyāḥ prāk iti vakṣyāmi . \\
\noindent \textbf{(P\_3,1.94.3)} striyāḥ prāk iti cet ktvāyām vāvacanam . striyāḥ prāk iti cet ktvāyām vāvacanam kartavyam . \\
\noindent \textbf{(P\_3,1.94.3)} āsitvā bhuṅkte . \\
\noindent \textbf{(P\_3,1.94.3)} āsyate bhoktum iti api yathā syāt . \\
\noindent \textbf{(P\_3,1.94.3)} kālādiṣu tumuni . \\
\noindent \textbf{(P\_3,1.94.3)} kālādiṣu tumuni vāvacanam kartavyam . \\
\noindent \textbf{(P\_3,1.94.3)} kālaḥ bhoktum . \\
\noindent \textbf{(P\_3,1.94.3)} kālaḥ bhojanasya iti api yathā syāt \\
\noindent \textbf{(P\_3,1.94.4)} arhe tṛjvidhānam . \\
\noindent \textbf{(P\_3,1.94.4)} arhe tṛc vidheyaḥ . \\
\noindent \textbf{(P\_3,1.94.4)} ime arhe kṛtyāḥ vidhīyante . \\
\noindent \textbf{(P\_3,1.94.4)} te viśeṣavihitāḥ sāmānyavihitam tṛcam bādheran . \\
\noindent \textbf{(P\_3,1.94.4)} na eṣaḥ doṣaḥ . \\
\noindent \textbf{(P\_3,1.94.4)} bhāvakarmaṇoḥ kṛtyāḥ vidhīyante kartari tṛc . \\
\noindent \textbf{(P\_3,1.94.4)} kaḥ prasaṅgaḥ yat bhāvakarmaṇoḥ kṛtyāḥ kartari tṛcam bādheran . \\
\noindent \textbf{(P\_3,1.94.4)} evam tarhi arhe kṛtyatṛjvidhānam . \\
\noindent \textbf{(P\_3,1.94.4)} arhe kṛtyatṛcaḥ vidheyāḥ . \\
\noindent \textbf{(P\_3,1.94.4)} ayam arhe liṅ vidhīyate . \\
\noindent \textbf{(P\_3,1.94.4)} saḥ viśeṣavihitaḥ sāmānyavihitān kṛtyatṛcaḥ bādheta . \\
\noindent \textbf{(P\_3,1.95)} kṛtyasañjñāyām prāṅṇvulvacanam . \\
\noindent \textbf{(P\_3,1.95)} kṛtyasañjñāyām prāk ṇvulaḥ iti vaktavyam . \\
\noindent \textbf{(P\_3,1.95)} kim prayojanam . \\
\noindent \textbf{(P\_3,1.95)} ṇvulaḥ kṛtyasañjñā mā bhūt . \\
\noindent \textbf{(P\_3,1.95)} arhe kṛtyatrjvacanam tu jñāpakam prāṅṇvulavanānarthyasya . \\
\noindent \textbf{(P\_3,1.95)} yat ayam arhe kṛtyatṛcaḥ ca iti tṛjgrahaṇam karoti tat jñāpayati ācāryaḥ prāk ṇvulaḥ kṛtyasañjñā bhavati iti . \\
\noindent \textbf{(P\_3,1.95)} evam api ṇvulaḥ kṛtyasañjñā prāpnoti . \\
\noindent \textbf{(P\_3,1.95)} yogāpekṣam jñāpakam . \\
\noindent \textbf{(P\_3,1.96)} kelimaraḥ upasaṅkhyānam . \\
\noindent \textbf{(P\_3,1.96)} kelimaraḥ upasaṅkhyānam kartavyam . \\
\noindent \textbf{(P\_3,1.96)} pacelimāḥ māṣāḥ . \\
\noindent \textbf{(P\_3,1.96)} paktavyāḥ . \\
\noindent \textbf{(P\_3,1.96)} bhidelimāḥ saralāḥ . \\
\noindent \textbf{(P\_3,1.96)} bhettavyāḥ . \\
\noindent \textbf{(P\_3,1.96)} vaseḥ tavyat kartari ṇit ca . \\
\noindent \textbf{(P\_3,1.96)} vaseḥ tavyat kartari vaktavyaḥ . \\
\noindent \textbf{(P\_3,1.96)} ṇit ca asau bhavati iti vaktavyam . \\
\noindent \textbf{(P\_3,1.96)} vasati iti vāstavyaḥ . \\
\noindent \textbf{(P\_3,1.96)} taddhitaḥ vā . taddhitaḥ vā punaḥ eṣaḥ bhaviṣyati . \\
\noindent \textbf{(P\_3,1.96)} vāstuni bhavaḥ vāstavyaḥ . \\
\noindent \textbf{(P\_3,1.97.1)} ajgrahaṇam kimartham . \\
\noindent \textbf{(P\_3,1.97.1)} ajantāt yathā syāt . \\
\noindent \textbf{(P\_3,1.97.1)} halantāt mā bhūt iti . \\
\noindent \textbf{(P\_3,1.97.1)} na etat asti prayojanam . \\
\noindent \textbf{(P\_3,1.97.1)} halantāt ṇyat vidhīyate . \\
\noindent \textbf{(P\_3,1.97.1)} saḥ bādhakaḥ bhaviṣyati . \\
\noindent \textbf{(P\_3,1.97.1)} yathā eva tarhi ṇyat yatam bādhate evam tavyādīn api bādheta . \\
\noindent \textbf{(P\_3,1.97.1)} ajgrahaṇe punaḥ kriyamāṇe ajantāt yat vidhīyate halantāt ṇyat . \\
\noindent \textbf{(P\_3,1.97.1)} etāvantaḥ ca dhātavaḥ yat uta ajantāḥ halantāḥ ca . \\
\noindent \textbf{(P\_3,1.97.1)} ucyante ca tavyādayaḥ . \\
\noindent \textbf{(P\_3,1.97.1)} te vacanāt bhaviṣyanti . \\
\noindent \textbf{(P\_3,1.97.1)} na etat asti prayojanam . \\
\noindent \textbf{(P\_3,1.97.1)} vāsarūpeṇa tavyādayaḥ bhaviṣyanti . \\
\noindent \textbf{(P\_3,1.97.1)} idam tarhi prayojanam . \\
\noindent \textbf{(P\_3,1.97.1)} ajantabhūtapūrvamātrāt api yathā syāt . \\
\noindent \textbf{(P\_3,1.97.1)} lavyam pavyam . \\
\noindent \textbf{(P\_3,1.97.1)} ārdhadhātukasāmānye guṇe kṛte yi pratyayasāmānye ca vāntādeśe kṛte halantāt iti ṇyat prāpnoti . \\
\noindent \textbf{(P\_3,1.97.1)} tathā ditsyam dhitsyam . \\
\noindent \textbf{(P\_3,1.97.1)} ārdhadhātukasāmānye akāralope kṛte halantāt iti ṇyat prāpnoti . \\
\noindent \textbf{(P\_3,1.97.1)} ajgrahaṇasāmarthyāt yat eva bhavati . \\
\noindent \textbf{(P\_3,1.97.2)} yati jāteḥ upasaṅkhyānam . \\
\noindent \textbf{(P\_3,1.97.2)} yati jāteḥ upasaṅkhyānam kartavyam . \\
\noindent \textbf{(P\_3,1.97.2)} janyam vatsena . \\
\noindent \textbf{(P\_3,1.97.2)} atyalpam idam ucyate . \\
\noindent \textbf{(P\_3,1.97.2)} takiśasicatiyatijanīnām upasaṅkhyānam iti vaktavyam . \\
\noindent \textbf{(P\_3,1.97.2)} taki takyam : śasi śasyam . \\
\noindent \textbf{(P\_3,1.97.2)} yati yatyam : jani : janyam . \\
\noindent \textbf{(P\_3,1.97.2)} hanaḥ vā vadha ca . \\
\noindent \textbf{(P\_3,1.97.2)} hanaḥ vā yat vaktavyaḥ vadha iti ayam ca ādeśaḥ vaktavyaḥ . \\
\noindent \textbf{(P\_3,1.97.2)} vadhyaḥ ghātyaḥ . \\
\noindent \textbf{(P\_3,1.97.2)} taddhitaḥ vā . \\
\noindent \textbf{(P\_3,1.97.2)} taddhitaḥ vā punaḥ eṣaḥ bhaviṣyati . \\
\noindent \textbf{(P\_3,1.97.2)} vadham arhati vadhyaḥ . \\
\noindent \textbf{(P\_3,1.97.2)} yadi taddhitaḥ samāsaḥ na prāpnoti : asivadhyaḥ , musalavadhyaḥ iti . \\
\noindent \textbf{(P\_3,1.97.2)} yadi punaḥ sati sādhanam kṛtā iti vā pādahārakādyartham iti samāsaḥ siddhaḥ bhavati . \\
\noindent \textbf{(P\_3,1.97.2)} yadi punaḥ asivadhaśabdāt utpattiḥ syāt . \\
\noindent \textbf{(P\_3,1.97.2)} asivadham arhati iti . \\
\noindent \textbf{(P\_3,1.97.2)} na evam śakyam . \\
\noindent \textbf{(P\_3,1.97.2)} svare hi doṣaḥ syāt . \\
\noindent \textbf{(P\_3,1.97.2)} asivadhyaḥ evam svaraḥ prasajyeta . \\
\noindent \textbf{(P\_3,1.97.2)} asivadhyaḥ iti ca iṣyate . \\
\noindent \textbf{(P\_3,1.100)} anupasargāt careḥ āṅi ca agurau . \\
\noindent \textbf{(P\_3,1.100)} anupasargāt careḥ iti atra āṅi ca agurau iti vaktavyam . \\
\noindent \textbf{(P\_3,1.100)} ācaryaḥ deśaḥ . \\
\noindent \textbf{(P\_3,1.100)} agurau iti kimartham . \\
\noindent \textbf{(P\_3,1.100)} ācāryaḥ upanayamānaḥ . \\
\noindent \textbf{(P\_3,1.103)} svāmini antodāttatvam ca . \\
\noindent \textbf{(P\_3,1.103)} svāmini antodāttatvam ca vaktavyam . \\
\noindent \textbf{(P\_3,1.103)} āryaḥ svāmī . \\
\noindent \textbf{(P\_3,1.105)} saṅgatam iti kim pratyudāhriyate . \\
\noindent \textbf{(P\_3,1.105)} ajaraḥ kambalaḥ . \\
\noindent \textbf{(P\_3,1.105)} ajaritā kambalaḥ iti . \\
\noindent \textbf{(P\_3,1.105)} kim punaḥ kāraṇam kartṛsādhanaḥ pratyudāhriyate . \\
\noindent \textbf{(P\_3,1.105)} na bhāvasādhanaḥ pratyudāhāryaḥ . \\
\noindent \textbf{(P\_3,1.105)} evam tarhi ajaryam kartari . \\
\noindent \textbf{(P\_3,1.105)} ajaryam kartari iti vaktavyam . \\
\noindent \textbf{(P\_3,1.105)} tat tarhi vaktavyam . \\
\noindent \textbf{(P\_3,1.105)} na vaktavyam . \\
\noindent \textbf{(P\_3,1.105)} gatyarthānām ktaḥ kartari vidhīyate . \\
\noindent \textbf{(P\_3,1.105)} tena yogāt ajaryam kartari bhaviṣyati . \\
\noindent \textbf{(P\_3,1.105)} gatyarthānām vai ktaḥ karmaṇi api vidhīyate . \\
\noindent \textbf{(P\_3,1.105)} tena yogāt ajaryam karmaṇi api prāpnoti . \\
\noindent \textbf{(P\_3,1.105)} jīryatiḥ akarmakaḥ . \\
\noindent \textbf{(P\_3,1.105)} bhāve tarhi prāpnoti . \\
\noindent \textbf{(P\_3,1.105)} saṅgatagrahaṇam idānīm kimartham syāt . \\
\noindent \textbf{(P\_3,1.105)} kartṛviśeṣaṇam saṅgatagrahaṇam . \\
\noindent \textbf{(P\_3,1.105)} saṅgatam cet kartṛ bhavati iti . \\
\noindent \textbf{(P\_3,1.105)} tat yathā hṛṣeḥ lomasu iti lomāni cet kartṝṇi bhavanti . \\
\noindent \textbf{(P\_3,1.106)} vadaḥ supi anupasargagrahaṇam . \\
\noindent \textbf{(P\_3,1.106)} vadaḥ supi anupasargagrahaṇam kartavyam . \\
\noindent \textbf{(P\_3,1.106)} iha mā bhūt . \\
\noindent \textbf{(P\_3,1.106)} pravādyam apavādyam iti . \\
\noindent \textbf{(P\_3,1.106)} tat tarhi vaktavyam . \\
\noindent \textbf{(P\_3,1.106)} na vaktavyam . \\
\noindent \textbf{(P\_3,1.106)} anupasarge iti vartate . \\
\noindent \textbf{(P\_3,1.106)} evam tarhi anvācaṣṭe anupasarge iti vartate . \\
\noindent \textbf{(P\_3,1.106)} na etat anvākhyeyam adhikārāḥ anuvartante iti . \\
\noindent \textbf{(P\_3,1.106)} eṣaḥ eva nyāyaḥ yat uta adhikārāḥ anuvarteran iti . \\
\noindent \textbf{(P\_3,1.107)} bhāvagrahaṇam kimartham . \\
\noindent \textbf{(P\_3,1.107)} karmaṇi mā bhūt iti . \\
\noindent \textbf{(P\_3,1.107)} na etat asti prayojanam . \\
\noindent \textbf{(P\_3,1.107)} bhavatiḥ ayam akarmaḥ . \\
\noindent \textbf{(P\_3,1.107)} akarmakāḥ api vai dhātavaḥ sopasargāḥ sakarmakāḥ bhavanti . \\
\noindent \textbf{(P\_3,1.107)} tena anubhavyam āmantraṇam iti atra api prāpnoti . \\
\noindent \textbf{(P\_3,1.107)} etat api na asti prayojanam . \\
\noindent \textbf{(P\_3,1.107)} anupasarge iti vartate . \\
\noindent \textbf{(P\_3,1.107)} uttarārtham tarhi bhāvagrahaṇam kartavyam . \\
\noindent \textbf{(P\_3,1.107)} hanaḥ ta ca bhāve yathā syāt . \\
\noindent \textbf{(P\_3,1.107)} śvahatyā vartate . \\
\noindent \textbf{(P\_3,1.107)} kva mā bhūt . \\
\noindent \textbf{(P\_3,1.107)} śvaghātyaḥ vṛṣālaḥ iti \\
\noindent \textbf{(P\_3,1.108)} hanaḥ taḥ cit striyām chandasi . \\
\noindent \textbf{(P\_3,1.108)} hanaḥ taḥ ca iti atra cit striyām chandasi vaktavyaḥ . \\
\noindent \textbf{(P\_3,1.108)} tām bhrūṇahatyām nigṛhya anucaraṇam . \\
\noindent \textbf{(P\_3,1.108)} asyai tvām bhrūṇahatyāyai caturtham pratigṛhāṇa . \\
\noindent \textbf{(P\_3,1.108)} striyām iti kimartham . \\
\noindent \textbf{(P\_3,1.108)} āghnate dasyuhatyāya . \\
\noindent \textbf{(P\_3,1.108)} chandasi iti kimartham . \\
\noindent \textbf{(P\_3,1.108)} dasyuhatyā śvahatyā vartate . \\
\noindent \textbf{(P\_3,1.109)} kyap iti vartamāne punaḥ kyabgrahaṇam kimartham . \\
\noindent \textbf{(P\_3,1.109)} kyap eva yathā syāt . \\
\noindent \textbf{(P\_3,1.109)} anyat yat prāpnoti tat mā bhūt iti . \\
\noindent \textbf{(P\_3,1.109)} kim ca anyat prāpnoti . \\
\noindent \textbf{(P\_3,1.109)} ṇyat . \\
\noindent \textbf{(P\_3,1.109)} oḥ āvaśyake ṇyataḥ stoteḥ kyap pūrvavipratiṣiddham iti vakṣyati . \\
\noindent \textbf{(P\_3,1.109)} saḥ pūrvavipratiṣedhaḥ na paṭhitavyaḥ bhavati . \\
\noindent \textbf{(P\_3,1.109)} atha vā hanaḥ taḥ cit striyām chandasi coditaḥ . \\
\noindent \textbf{(P\_3,1.109)} saḥ na vaktavyaḥ bhavati . \\
\noindent \textbf{(P\_3,1.109)} kyabvidhau vṛñgrahaṇam . kyabvidhau vṛñgrahaṇam kartavyam . \\
\noindent \textbf{(P\_3,1.109)} iha mā bhūt . \\
\noindent \textbf{(P\_3,1.109)} vāryāḥ ṛtvijaḥ iti . \\
\noindent \textbf{(P\_3,1.109)} añjeḥ ca upasaṅkhyānam sañjñāyām . \\
\noindent \textbf{(P\_3,1.109)} sañjñāyām añjeḥ ca upasaṅkhyānam kartavyam . \\
\noindent \textbf{(P\_3,1.109)} ājyam . \\
\noindent \textbf{(P\_3,1.109)} yadi kyap vṛddhiḥ na prāpnoti . \\
\noindent \textbf{(P\_3,1.109)} tasmāt ṇyat eṣaḥ . \\
\noindent \textbf{(P\_3,1.109)} yadi ṇyat upadhālopaḥ na prāpnoti . \\
\noindent \textbf{(P\_3,1.109)} tasmāt kyap eṣaḥ . \\
\noindent \textbf{(P\_3,1.109)} nanu ca uktam vṛddhiḥ na prāpnoti iti . \\
\noindent \textbf{(P\_3,1.109)} āṅpūrvasya eṣaḥ prayogaḥ bhaviṣyati . \\
\noindent \textbf{(P\_3,1.109)} yadi evam avagrahaḥ prāpnoti . \\
\noindent \textbf{(P\_3,1.109)} na lakṣaṇena padakārāḥ anuvartyāḥ . \\
\noindent \textbf{(P\_3,1.109)} padkāraiḥ nāma lakṣaṇam anuvartyam . \\
\noindent \textbf{(P\_3,1.109)} yathālakṣaṇam padam kartavyam . \\
\noindent \textbf{(P\_3,1.111)} dīrghoccāraṇam kimartham na i ca khanaḥ iti eva ucyeta . \\
\noindent \textbf{(P\_3,1.111)} kā rūpasiddhiḥ : kheyam . \\
\noindent \textbf{(P\_3,1.111)} ādguṇena siddham . \\
\noindent \textbf{(P\_3,1.111)} na sidhyati . \\
\noindent \textbf{(P\_3,1.111)} ṣatvatukoḥ asiddhaḥ ekādeśaḥ iti ekādeśasya asiddhatvāt tuk prasajyeta . \\
\noindent \textbf{(P\_3,1.111)} na etat asti . \\
\noindent \textbf{(P\_3,1.111)} padāntapadādyoḥ ekādeśaḥ asiddhaḥ . \\
\noindent \textbf{(P\_3,1.111)} na ca eṣaḥ padāntapadādyoḥ ekādeśaḥ . \\
\noindent \textbf{(P\_3,1.111)} tasmāt i ca khanaḥ iti eva vaktavyam . \\
\noindent \textbf{(P\_3,1.112)} asañjñāyām iti kimartham . \\
\noindent \textbf{(P\_3,1.112)} bhāryā . \\
\noindent \textbf{(P\_3,1.112)} bhṛñaḥ sañjñāpratiṣedhe striyām apratiṣedhaḥ anyena vihitatvāt . \\
\noindent \textbf{(P\_3,1.112)} bhṛñaḥ sañjñāpratiṣedhe striyām apratiṣedhaḥ . \\
\noindent \textbf{(P\_3,1.112)} anarthakaḥ pratiṣedhaḥ apratiṣedhaḥ . \\
\noindent \textbf{(P\_3,1.112)} kim kāraṇam . \\
\noindent \textbf{(P\_3,1.112)} anyena vihitatvāt . \\
\noindent \textbf{(P\_3,1.112)} anyena lakṣaṇena striyām kyap vidhīyate . \\
\noindent \textbf{(P\_3,1.112)} sañjñāyām samajaniṣadanipatamanavidaṣuñśīṅbhṛñiṇaḥ iti . \\
\noindent \textbf{(P\_3,1.112)} pratiṣedhaḥ idānīm kimarthaḥ syāt . \\
\noindent \textbf{(P\_3,1.112)} pratiṣedhaḥ kimarthaḥ iti cet astrīsañjñāpratiṣedhārthaḥ . \\
\noindent \textbf{(P\_3,1.112)} pratiṣedhaḥ kimarthaḥ iti cet astrīsañjñā asti tadarthaḥ pratiṣedhaḥ syāt . \\
\noindent \textbf{(P\_3,1.112)} bhāryāḥ nāma kṣatriyāḥ . \\
\noindent \textbf{(P\_3,1.112)} siddham tu striyām sañjñāpratiṣedhāt . \\
\noindent \textbf{(P\_3,1.112)} siddham etat . \\
\noindent \textbf{(P\_3,1.112)} katham . \\
\noindent \textbf{(P\_3,1.112)} striyām sañjñāpratiṣedhaḥ vaktavyaḥ . \\
\noindent \textbf{(P\_3,1.112)} sañjñāyām samajaniṣadanipatamanavidaṣuñśīṅbhṛñiṇaḥ tataḥ na striyām bhṛñaḥ iti . \\
\noindent \textbf{(P\_3,1.112)} sidhyati . \\
\noindent \textbf{(P\_3,1.112)} sūtram tarhi bhidyate . \\
\noindent \textbf{(P\_3,1.112)} yathānyāsam eva astu . \\
\noindent \textbf{(P\_3,1.112)} nanu ca uktam bhṛñaḥ sañjñāpratiṣedhe striyām apratiṣedhaḥ anyena vihitatvāt iti . \\
\noindent \textbf{(P\_3,1.112)} na eṣaḥ doṣaḥ . \\
\noindent \textbf{(P\_3,1.112)} bhāve iti tatra anuvartate . \\
\noindent \textbf{(P\_3,1.112)} karmasādhanaḥ ca ayam . \\
\noindent \textbf{(P\_3,1.112)} atha vā ye ete sañjñāyām vidhīyante teṣu na evam vijñāyate sañjñāyām abhidheyāyām iti . \\
\noindent \textbf{(P\_3,1.112)} kim tarhi . \\
\noindent \textbf{(P\_3,1.112)} pratyayāntena cet sañjñā gamyate iti . \\
\noindent \textbf{(P\_3,1.112)} aparaḥ āha : sañjñāyām puṃsi dṛṣṭatvāt na te bhāryā prasidhyati . \\
\noindent \textbf{(P\_3,1.112)} sañjñāyām puṃsi dṛṣṭatvāt tava bhāryāśabdaḥ na sidhyati . \\
\noindent \textbf{(P\_3,1.112)} striyām bhāvādhikāraḥ asti tena bhāryā prasidhyati . \\
\noindent \textbf{(P\_3,1.112)} bhāve iti tatra vartate . \\
\noindent \textbf{(P\_3,1.112)} karmasādhanaḥ ca ayam . \\
\noindent \textbf{(P\_3,1.112)} atha vā bahulam kṛtyāḥ sañjñāyām iti tat smṛtam . \\
\noindent \textbf{(P\_3,1.112)} atha vā kṛtyalyuṭaḥ bahulam iti evam atra api ṇyat bhaviṣyati . \\
\noindent \textbf{(P\_3,1.112)} yathā yatyam janyam yathā bhittiḥ tathā eva sā . \\
\noindent \textbf{(P\_3,1.112)} samaḥ ca bahulam . \\
\noindent \textbf{(P\_3,1.112)} samaḥ ca bahulam upasaṅkhyānam kartavyam . \\
\noindent \textbf{(P\_3,1.112)} sambhṛtyāḥ eva sambhārāḥ . \\
\noindent \textbf{(P\_3,1.112)} sambhāryāḥ eva sambhārāḥ . \\
\noindent \textbf{(P\_3,1.114)} sūryarucyāvyathyāḥ kartari . \\
\noindent \textbf{(P\_3,1.114)} sūrya ruci avyathya iti kartari nipātyante . \\
\noindent \textbf{(P\_3,1.114)} kim nipātyate . \\
\noindent \textbf{(P\_3,1.114)} sūryaḥ . \\
\noindent \textbf{(P\_3,1.114)} sūsartibhyām sarteḥ utvam suvateḥ vā ruḍāgamaḥ . saraṇāt vā suvati vā karmaṇi iti sūryaḥ . \\
\noindent \textbf{(P\_3,1.114)} rucya . \\
\noindent \textbf{(P\_3,1.114)} rocate asau rucyaḥ . \\
\noindent \textbf{(P\_3,1.114)} na vyathathe avyathyaḥ . \\
\noindent \textbf{(P\_3,1.114)} kupyam sañjñāyām . \\
\noindent \textbf{(P\_3,1.114)} kupyam sañjñāyām iti vaktavyam . \\
\noindent \textbf{(P\_3,1.114)} gopyam anyat . \\
\noindent \textbf{(P\_3,1.114)} kṛṣṭapacyasya antodāttatvam ca karmakartari ca . \\
\noindent \textbf{(P\_3,1.114)} kṛṣṭapacyasya antodāttatvam ca karmakartari ca iti vaktavyam . \\
\noindent \textbf{(P\_3,1.114)} kṛṣṭe pacyante svayam eva . \\
\noindent \textbf{(P\_3,1.114)} kṛṣṭapacyāḥ ca me akṛṣṭapacyāḥ ca me . \\
\noindent \textbf{(P\_3,1.114)} yaḥ hi kṛṣṭe paktavyaḥ kṣṭapākyaḥ sa bhavati . \\
\noindent \textbf{(P\_3,1.118)} pratyapibhyām graheḥ chandasi . \\
\noindent \textbf{(P\_3,1.118)} pratyapibhyām graheḥ chandasi iti vaktavyam . \\
\noindent \textbf{(P\_3,1.118)} mattasya na pratigṛhyam . \\
\noindent \textbf{(P\_3,1.118)} anṛtam hi mattaḥ bhavati . \\
\noindent \textbf{(P\_3,1.118)} tasmāt na apigṛhyam . \\
\noindent \textbf{(P\_3,1.118)} pratigrāhyam apigrāhyam iti eva anyatra . \\
\noindent \textbf{(P\_3,1.122)} kasya ayam anubandhaḥ . \\
\noindent \textbf{(P\_3,1.122)} pradhānasya . \\
\noindent \textbf{(P\_3,1.122)} yadi pradhānasya amāvasyā evam svaraḥ prasajyeta . \\
\noindent \textbf{(P\_3,1.122)} amāvasyā iti ca iṣyate . \\
\noindent \textbf{(P\_3,1.122)} tathā amāvāsyāgrahaṇena amāvasyāgrahaṇam na prāpnoti . \\
\noindent \textbf{(P\_3,1.122)} evam tarhi nipātanasya . \\
\noindent \textbf{(P\_3,1.122)} yadi tarhi nipātanāni api evañjātīyakāni bhavanti śrotriyan chandaḥ adhīte iti vyapavargābhāvāt ñniti iti ādyudāttatvam na prāpnoti . \\
\noindent \textbf{(P\_3,1.122)} evam tarhi amāvasoḥ aham ṇyatoḥ nipātayāmi avṛddhitām . \\
\noindent \textbf{(P\_3,1.122)} tathā ekavṛttitā tayoḥ svaraḥ ca me prasidhyati . \\
\noindent \textbf{(P\_3,1.123)} niṣṭarkya iti kim nipātyate . \\
\noindent \textbf{(P\_3,1.123)} niṣṭarkye kṛteḥ ādyantaviparyayaḥ chandasi kṛtādyarthaḥ . \\
\noindent \textbf{(P\_3,1.123)} yathā kṛteḥ tarkuḥ kaseḥ sikatāḥ hiṃseḥ siṃhaḥ . \\
\noindent \textbf{(P\_3,1.123)} aparaḥ āha : niṣṭarkye vyatyayam vidyāt nisaḥ ṣatvam nipātanāt . \\
\noindent \textbf{(P\_3,1.123)} ṇyat āyādeśaḥ iti etau upacāyye nipātitau . niṣṭarkyam cinvīta paśukāmaḥ . \\
\noindent \textbf{(P\_3,1.123)} ṇyat ekasmāt caturbhyaḥ kyaP caturbhyaḥ yataḥ vidhiḥ . \\
\noindent \textbf{(P\_3,1.123)} ṇyat ekasmāt yaśabdaḥ ca dvau kyapau ṇyadvidhiḥ catuḥ . ṇyat ekasmāt . \\
\noindent \textbf{(P\_3,1.123)} niṣṭarkyaḥ . \\
\noindent \textbf{(P\_3,1.123)} caturbhyaḥ kyap . \\
\noindent \textbf{(P\_3,1.123)} devahūyaḥ praṇīyaḥ unnīyaḥ ucchiṣyaḥ . \\
\noindent \textbf{(P\_3,1.123)} caturbhyaḥ ca yataḥ vidhiḥ . \\
\noindent \textbf{(P\_3,1.123)} maryaḥ staryā dhvaryaḥ khanyaḥ . \\
\noindent \textbf{(P\_3,1.123)} ṇyat ekasmāt . \\
\noindent \textbf{(P\_3,1.123)} khānyaḥ . \\
\noindent \textbf{(P\_3,1.123)} yaśabdaḥ ca . \\
\noindent \textbf{(P\_3,1.123)} devayajyā . \\
\noindent \textbf{(P\_3,1.123)} dvau kyapau . \\
\noindent \textbf{(P\_3,1.123)} āpṛcchyaḥ pratiṣīvyaḥ . \\
\noindent \textbf{(P\_3,1.123)} ṇyadvidhiḥ catuḥ . \\
\noindent \textbf{(P\_3,1.123)} brahmavādyaḥ bhāvyaḥ stāvyaḥ upacāyyapṛḍam . \\
\noindent \textbf{(P\_3,1.123)} upapūrvāt cinoteḥ āyādeśaḥ nipātyate . \\
\noindent \textbf{(P\_3,1.123)} na hi ṇyatā eva sidhyati . \\
\noindent \textbf{(P\_3,1.123)} hiraṇye iti vaktavyam . \\
\noindent \textbf{(P\_3,1.123)} upaceyapṛḍam eva anyatra . \\
\noindent \textbf{(P\_3,1.124)} pāṇau sṛjeḥ ṇyadvidhiḥ . \\
\noindent \textbf{(P\_3,1.124)} pāṇau sṛjeḥ ṇyat vidheyaḥ . \\
\noindent \textbf{(P\_3,1.124)} pāṇisargyā rajjuḥ . \\
\noindent \textbf{(P\_3,1.124)} samavapūrvāt ca . \\
\noindent \textbf{(P\_3,1.124)} samavapūrvāt ca iti vaktavyam . \\
\noindent \textbf{(P\_3,1.124)} samavasargyaḥ . \\
\noindent \textbf{(P\_3,1.124)} lapidamibhyām ca. lapidamibhyām ca iti vaktavyam . \\
\noindent \textbf{(P\_3,1.124)} apalapyam avadāmyam . \\
\noindent \textbf{(P\_3,1.125.1)} katham idam vijñāyate . \\
\noindent \textbf{(P\_3,1.125.1)} āvaśyake upapade āhosvit dhyotye iti . \\
\noindent \textbf{(P\_3,1.125.1)} kaḥ ca atra viśeṣaḥ . \\
\noindent \textbf{(P\_3,1.125.1)} āvaśyake upapade iti cet dyotye upasaṅkhyānam . \\
\noindent \textbf{(P\_3,1.125.1)} āvaśyake upapade iti cet dyotye upasaṅkhyānam kartavyam . \\
\noindent \textbf{(P\_3,1.125.1)} lāvyam pāvyam . \\
\noindent \textbf{(P\_3,1.125.1)} astu tarhi dyotye . \\
\noindent \textbf{(P\_3,1.125.1)} dyotye iti cet svarasamāsānupapattiḥ . \\
\noindent \textbf{(P\_3,1.125.1)} dyotye iti cet svarasamāsānupapattiḥ . \\
\noindent \textbf{(P\_3,1.125.1)} āvaśyalāvyam āvaśyapāvyam . \\
\noindent \textbf{(P\_3,1.125.1)} na eṣaḥ doṣaḥ . \\
\noindent \textbf{(P\_3,1.125.1)} mayūravyaṃsakāditvāt samāsaḥ viśpaṣṭādivat svaraḥ bhaviṣyati . \\
\noindent \textbf{(P\_3,1.125.2)} oḥ āvaśyake ṇyataḥ stauteḥ kyaP pūrvavipratiṣiddham . oḥ āvaśyake ṇyataḥ stauteḥ kyap bhavati pūrvavipratiṣdhena . \\
\noindent \textbf{(P\_3,1.125.2)} oḥ āvaśyake ṇyat bhavati iti asya avakāśaḥ . \\
\noindent \textbf{(P\_3,1.125.2)} avaśyalāvyam avaśyapāvyam . \\
\noindent \textbf{(P\_3,1.125.2)} kyapaḥ avakāśaḥ . \\
\noindent \textbf{(P\_3,1.125.2)} stutyaḥ . \\
\noindent \textbf{(P\_3,1.125.2)} iha ubhayam prāpnoti . \\
\noindent \textbf{(P\_3,1.125.2)} avaśyastutyaḥ . \\
\noindent \textbf{(P\_3,1.125.2)} kyap bhavati pūrvavipratiṣdhena . \\
\noindent \textbf{(P\_3,1.125.2)} saḥ tarhi pūrvavipratiṣedhaḥ vaktavyaḥ . \\
\noindent \textbf{(P\_3,1.125.2)} na vaktavyaḥ . \\
\noindent \textbf{(P\_3,1.125.2)} uktam tatra kyap iti vartamāne punaḥ kyabgrahaṇasya prayojanam kyap eva yathā syāt . \\
\noindent \textbf{(P\_3,1.125.2)} anyat yat prāpnoti tat mā bhūt iti . \\
\noindent \textbf{(P\_3,1.127)} dakṣiṇāgnau iti vaktavyam . \\
\noindent \textbf{(P\_3,1.127)} āneyaḥ anyaḥ . \\
\noindent \textbf{(P\_3,1.127)} ānāyyaḥ anityaḥ iti cet dakṣiṇāgnau kṛtam bhavet . \\
\noindent \textbf{(P\_3,1.127)} ekayonau tu tam vidyāt . \\
\noindent \textbf{(P\_3,1.127)} āneyaḥ hi anyathā bhavet . \\
\noindent \textbf{(P\_3,1.129)} pāyyanikāyyayoḥ kim nipātyate . \\
\noindent \textbf{(P\_3,1.129)} pāyyanikāyyayoḥ ādipatvakatvanipātanam . pāyyanikāyyayoḥ ādipatvam ādikatvam ca nipātyate . \\
\noindent \textbf{(P\_3,1.129)} meyam niceyam iti eva anyatra . \\
\noindent \textbf{(P\_3,1.130)} kuṇḍapāyye yadvidhiḥ . \\
\noindent \textbf{(P\_3,1.130)} kuṇḍapāyye yat vidheyaḥ . \\
\noindent \textbf{(P\_3,1.130)} kuṇḍapāyyaḥ kratuḥ . \\
\noindent \textbf{(P\_3,1.131)} samūhyaḥ iti anarthakam vacanam sāmānyena kṛtatvāt . \\
\noindent \textbf{(P\_3,1.131)} samūhyaḥ iti vacanam anarthakam . \\
\noindent \textbf{(P\_3,1.131)} kim kāraṇam . \\
\noindent \textbf{(P\_3,1.131)} sāmānyena kṛtatvāt . \\
\noindent \textbf{(P\_3,1.131)} sāmānyena eva ṇyat bhaviṣyati : ṛhaloḥ ṇyat iti . \\
\noindent \textbf{(P\_3,1.131)} vahyartham tarhi nipātanam kartavyam . \\
\noindent \textbf{(P\_3,1.131)} vaheḥ ṇyat yathā syāt . \\
\noindent \textbf{(P\_3,1.131)} vahyartham iti cet ūheḥ tadarthatvāt siddham . \\
\noindent \textbf{(P\_3,1.131)} ūhiḥ api vahyarthe vartate . \\
\noindent \textbf{(P\_3,1.131)} katham punaḥ anyaḥ nāma anyasya arthe vartate . \\
\noindent \textbf{(P\_3,1.131)} katham ūhiḥ vahyarthe vartate . \\
\noindent \textbf{(P\_3,1.131)} bahvarthāḥ api dhātavaḥ bhavanti iti . \\
\noindent \textbf{(P\_3,1.131)} asti punaḥ kva cit anyatra api ūhiḥ vahyarthe vartate . \\
\noindent \textbf{(P\_3,1.131)} asti iti āha . \\
\noindent \textbf{(P\_3,1.131)} ūhivigrahāt ca brāhmaṇe siddham . ūhivigrahāt ca brāhmaṇe siddham etat . \\
\noindent \textbf{(P\_3,1.131)} samūhyam cinvīta paśukāmaḥ . \\
\noindent \textbf{(P\_3,1.131)} paśavaḥ vai purīṣam . \\
\noindent \textbf{(P\_3,1.131)} paśūn eva asmai tat samūhati . \\
\noindent \textbf{(P\_3,1.132)} agnicityā bhāve antodāttaḥ . \\
\noindent \textbf{(P\_3,1.132)} agnicityā iti bhāve antodāttaḥ . \\
\noindent \textbf{(P\_3,1.132)} agnicayanam eva agnicityā . \\
\noindent \textbf{(P\_3,1.133.1)} kimarthaḥ cakāraḥ . \\
\noindent \textbf{(P\_3,1.133.1)} svarārthaḥ . \\
\noindent \textbf{(P\_3,1.133.1)} citaḥ antaḥ udāttaḥ bhavati iti antodāttatvam yathā syāt . \\
\noindent \textbf{(P\_3,1.133.1)} na etat asti prayojanam . \\
\noindent \textbf{(P\_3,1.133.1)} ekāc ayam . \\
\noindent \textbf{(P\_3,1.133.1)} tatra na arthaḥ svarārthena cakāreṇa anubandhena . \\
\noindent \textbf{(P\_3,1.133.1)} pratyayasvareṇa eva siddham . \\
\noindent \textbf{(P\_3,1.133.1)} viśeṣaṇārthaḥ tarhi . \\
\noindent \textbf{(P\_3,1.133.1)} kva viśeṣaṇārthena arthaḥ : aptṛntṛc iti . \\
\noindent \textbf{(P\_3,1.133.1)} tṛ iti ucyamāne mātarau mātaraḥ pitarau pitaraḥ atra api prasajyeta . \\
\noindent \textbf{(P\_3,1.133.1)} svasṛnaptṛgrahaṇam niyamāṛtham bhaviṣyati . \\
\noindent \textbf{(P\_3,1.133.1)} etayoḥ eva yonisambandhayoḥ na anyeṣām yonisambandhānām iti . \\
\noindent \textbf{(P\_3,1.133.1)} sāmānyagrahaṇāvighātārthaḥ tarhi . \\
\noindent \textbf{(P\_3,1.133.1)} kva sāmānyagrahaṇāvighātārthena arthaḥ . \\
\noindent \textbf{(P\_3,1.133.1)} atra eva . \\
\noindent \textbf{(P\_3,1.133.1)} yat etat tṛntṛcoḥ grahaṇam etat tṛ iti vakṣyāmi . \\
\noindent \textbf{(P\_3,1.133.1)} yadi tṛ ici ucyate mātarau mātaraḥ pitarau pitaraḥ atra api prasajyeta . \\
\noindent \textbf{(P\_3,1.133.1)} svasṛnaptṛgrahaṇam niyamāṛtham bhaviṣyati : etayoḥ eva yonisambandhayoḥ na anyeṣām yonisambandhānām iti . \\
\noindent \textbf{(P\_3,1.133.2)} ṇvuli sakarmakagrahaṇam . \\
\noindent \textbf{(P\_3,1.133.2)} ṇvuli sakarmakagrahaṇam kartavyam . \\
\noindent \textbf{(P\_3,1.133.2)} iha mā bhūt . \\
\noindent \textbf{(P\_3,1.133.2)} āsitā śayitā iti . \\
\noindent \textbf{(P\_3,1.133.2)} na vā dhātumātrāt darśanāt ṇvulaḥ . na vā vaktavyam . \\
\noindent \textbf{(P\_3,1.133.2)} kim kāraṇam . \\
\noindent \textbf{(P\_3,1.133.2)} dhātumātrāt ṇvul dṛśyate . \\
\noindent \textbf{(P\_3,1.133.2)} ime asya āsakāḥ ime . \\
\noindent \textbf{(P\_3,1.133.2)} asya śāyakāḥ . \\
\noindent \textbf{(P\_3,1.133.2)} utthitāḥ āsakā vaiśravaṇasya iti . \\
\noindent \textbf{(P\_3,1.133.2)} tṛjādiṣu vartamānakālopādānam adhyāyakavedādhyāyakārtham . \\
\noindent \textbf{(P\_3,1.133.2)} tṛjādiṣu vartamānakālopādānam kartavyam . \\
\noindent \textbf{(P\_3,1.133.2)} kim kāraṇam . \\
\noindent \textbf{(P\_3,1.133.2)} adhyāyakavedādhyāyakārtham . \\
\noindent \textbf{(P\_3,1.133.2)} adhyāyakaḥ vedādhyāyaḥ . \\
\noindent \textbf{(P\_3,1.133.2)} adhītavati adhyeṣyamāṇe vā mā bhūt . \\
\noindent \textbf{(P\_3,1.133.2)} na vā kālamātre darśanāt anyeṣām . \\
\noindent \textbf{(P\_3,1.133.2)} na vā vaktavyam . \\
\noindent \textbf{(P\_3,1.133.2)} kim kāraṇam . \\
\noindent \textbf{(P\_3,1.133.2)} kālamātre darśanāt anyeṣām . \\
\noindent \textbf{(P\_3,1.133.2)} kālamātre hi anye pratyayāḥ dṛśyante . \\
\noindent \textbf{(P\_3,1.133.2)} carcāpāraḥ śamanīpāraḥ . \\
\noindent \textbf{(P\_3,1.134)} ac api sarvadhātubhyaḥ . \\
\noindent \textbf{(P\_3,1.134)} ac api sarvadhātubhyaḥ vaktavyaḥ . \\
\noindent \textbf{(P\_3,1.134)} iha api yathā syāt . \\
\noindent \textbf{(P\_3,1.134)} bhavaḥ śarvaḥ . \\
\noindent \textbf{(P\_3,1.134)} na tarhi idānīm idam pacādyanukramaṇam kartavyam . \\
\noindent \textbf{(P\_3,1.134)} kartavyam ca . \\
\noindent \textbf{(P\_3,1.134)} kim prayojanam . \\
\noindent \textbf{(P\_3,1.134)} pacādyanukramaṇam anubandhāsañjārtham apavādabādhanārtham ca . \\
\noindent \textbf{(P\_3,1.134)} anubandhāsañjanārtham tāvat . \\
\noindent \textbf{(P\_3,1.134)} nadaṭ nadī coraṭ corī . \\
\noindent \textbf{(P\_3,1.134)} apavādabādhanārtham . \\
\noindent \textbf{(P\_3,1.134)} jārabharā śvapacā iti . \\
\noindent \textbf{(P\_3,1.135)} igupadhebhyaḥ upasarge kavidhiḥ meṣādyarthaḥ .igupadhebhyaḥ upasarge kaḥ vidheyaḥ . \\
\noindent \textbf{(P\_3,1.135)} kim prayojanam . \\
\noindent \textbf{(P\_3,1.135)} meṣādyarthaḥ . \\
\noindent \textbf{(P\_3,1.135)} meṣaḥ devaḥ sevaḥ . \\
\noindent \textbf{(P\_3,1.135)} na vā budhādīnām darśanāt anupasarge api . \\
\noindent \textbf{(P\_3,1.135)} na vā vaktavyaḥ . \\
\noindent \textbf{(P\_3,1.135)} kim kāraṇam . \\
\noindent \textbf{(P\_3,1.135)} budhādīnām anupasarge api kaḥ dṛśyate . \\
\noindent \textbf{(P\_3,1.135)} budhaḥ bhidaḥ yudhaḥ sivaḥ iti . \\
\noindent \textbf{(P\_3,1.135)} katham meṣaḥ devaḥ sevaḥ iti . \\
\noindent \textbf{(P\_3,1.135)} pacāciṣu pāṭhaḥ kariṣyate . \\
\noindent \textbf{(P\_3,1.137)} jighraḥ sañjñāyām pratiṣedhaḥ . \\
\noindent \textbf{(P\_3,1.137)} jighraḥ sañjñāyām pratiṣedhaḥ vaktavyaḥ . \\
\noindent \textbf{(P\_3,1.137)} vyājighrati iti vyāghraḥ . \\
\noindent \textbf{(P\_3,1.137)} iha ke cit śasya eva pratiṣedham āhuḥ ke cit jighrabhāvasya . \\
\noindent \textbf{(P\_3,1.137)} kim punaḥ atra nyāyyam . \\
\noindent \textbf{(P\_3,1.137)} śasya eva pratiṣedhaḥ nyāyyaḥ . \\
\noindent \textbf{(P\_3,1.137)} jighrabhāve hi pratiṣiddhe kena śe ākāralopaḥ syāt . \\
\noindent \textbf{(P\_3,1.138)} anupasargāt nau limpeḥ . \\
\noindent \textbf{(P\_3,1.138)} anupasargāt nau limpeḥ iti vaktavyam . \\
\noindent \textbf{(P\_3,1.138)} nilimpāḥ nāma devāḥ . \\
\noindent \textbf{(P\_3,1.138)} gavi ca vindeḥ sañjñāyām . \\
\noindent \textbf{(P\_3,1.138)} gavi ca upapade vindeḥ sañjñāyām upasaṅkhyānam kartavyam . \\
\noindent \textbf{(P\_3,1.138)} govindaḥ iti . \\
\noindent \textbf{(P\_3,1.138)} atyalpam idam ucyate : gavi iti . \\
\noindent \textbf{(P\_3,1.138)} gavādiṣu iti vaktavyam . \\
\noindent \textbf{(P\_3,1.138)} govindaḥ aravindaḥ . \\
\noindent \textbf{(P\_3,1.140)} tanoteḥ ṇaḥ upasaṅkhyānam . tanoteḥ ṇaḥ upasaṅkhyānam kartavyam . \\
\noindent \textbf{(P\_3,1.140)} avatanoti iti avatānaḥ . \\
\noindent \textbf{(P\_3,1.145)} nṛtikhanirañjibhyaḥ iti vaktavyam . \\
\noindent \textbf{(P\_3,1.145)} iha mā bhūt . \\
\noindent \textbf{(P\_3,1.145)} hvāyakaḥ iti . \\
\noindent \textbf{(P\_3,1.149)} prusṛlvaḥ sādhukāriṇi vunvidhānam . \\
\noindent \textbf{(P\_3,1.149)} prusṛlvaḥ sādhukāriṇi vun vidheyaḥ . \\
\noindent \textbf{(P\_3,1.149)} sakṛt api yaḥ suṣṭhu karoti tatra yathā syāt . \\
\noindent \textbf{(P\_3,1.149)} bahuśaḥ yaḥ duṣṭhu karoti tatra mā bhūt . \\
\noindent \textbf{(P\_3,2.1.1)} karmaṇi nirvartyamāṇavikriyamāṇe iti vaktatvyam . \\
\noindent \textbf{(P\_3,2.1.1)} iha mā bhūt . \\
\noindent \textbf{(P\_3,2.1.1)} ādityam paśyati . \\
\noindent \textbf{(P\_3,2.1.1)} hivavantam śrṇoti . \\
\noindent \textbf{(P\_3,2.1.1)} grāmam gacchati iti . \\
\noindent \textbf{(P\_3,2.1.1)} karmaṇi nirvartyamāṇavikriyamāṇe cet vedādhyāyādīnām upasaṅkhyānam . \\
\noindent \textbf{(P\_3,2.1.1)} karmaṇi nirvartyamāṇavikriyamāṇe cet vedādhyāyādīnām upasaṅkhyānam kartavyam . \\
\noindent \textbf{(P\_3,2.1.1)} vedādhyāyaḥ carcāpāraḥ śamanīpāraḥ . \\
\noindent \textbf{(P\_3,2.1.1)} yatra ca niyuktaḥ . yatra ca niyuktaḥ tatra upasaṅkhyānam kartatvyam . \\
\noindent \textbf{(P\_3,2.1.1)} chatradhāraḥ dvārapālaḥ . \\
\noindent \textbf{(P\_3,2.1.1)} hṛgrahinīvahibhyaḥ ca . \\
\noindent \textbf{(P\_3,2.1.1)} hṛgrahinīvahibhyaḥ ca iti vaktavyam . \\
\noindent \textbf{(P\_3,2.1.1)} hṛ . \\
\noindent \textbf{(P\_3,2.1.1)} bhārahāraḥ . \\
\noindent \textbf{(P\_3,2.1.1)} grahi . \\
\noindent \textbf{(P\_3,2.1.1)} kamaṇḍalugrāhaḥ . \\
\noindent \textbf{(P\_3,2.1.1)} nī. uṣṭrapraṇāyaḥ . \\
\noindent \textbf{(P\_3,2.1.1)} vahi . \\
\noindent \textbf{(P\_3,2.1.1)} bhāravāhaḥ . \\
\noindent \textbf{(P\_3,2.1.1)} aparigaṇanam vā . \\
\noindent \textbf{(P\_3,2.1.1)} na vā arthaḥ parigaṇanena . \\
\noindent \textbf{(P\_3,2.1.1)} kasmāt na bhavati : ādityam paśyati , himavantam śrṇoti . \\
\noindent \textbf{(P\_3,2.1.1)} grāmam gacchati iti . \\
\noindent \textbf{(P\_3,2.1.1)} anabhidhānāt . \\
\noindent \textbf{(P\_3,2.1.1)} anabhidhānāt eva na bhaviṣyati . \\
\noindent \textbf{(P\_3,2.1.2)} akārāt anupapadāt karmopapadaḥ vipratiṣedhena . \\
\noindent \textbf{(P\_3,2.1.2)} akārāt anupapadāt karmopapadaḥ bhavati vipratiṣedhena . \\
\noindent \textbf{(P\_3,2.1.2)} anupapadasya avakāśaḥ pacati iti pacaḥ . \\
\noindent \textbf{(P\_3,2.1.2)} karmopapadasya avakāśaḥ kumbhakāraḥ nagarakāraḥ . \\
\noindent \textbf{(P\_3,2.1.2)} odanapāce ubhayam prāpnoti . \\
\noindent \textbf{(P\_3,2.1.2)} karmopadaḥ bhavati vipratiṣedhena . \\
\noindent \textbf{(P\_3,2.1.2)} anupapadasya avakāśaḥ vikṣipaḥ vilikhaḥ . \\
\noindent \textbf{(P\_3,2.1.2)} karmopapadasya saḥ eva . \\
\noindent \textbf{(P\_3,2.1.2)} kāṣṭhabhede ubhayam prāpnoti . \\
\noindent \textbf{(P\_3,2.1.2)} karmopadaḥ bhavati vipratiṣedhena . \\
\noindent \textbf{(P\_3,2.1.2)} anupapadasya avakāśaḥ jānāti iti jñaḥ . \\
\noindent \textbf{(P\_3,2.1.2)} karmopapadasya saḥ eva . \\
\noindent \textbf{(P\_3,2.1.2)} arthajñe ubhayam prāpnoti . \\
\noindent \textbf{(P\_3,2.1.2)} karmopadaḥ bhavati vipratiṣedhena . \\
\noindent \textbf{(P\_3,2.1.2)} na eṣaḥ yuktaḥ vipratiṣedhaḥ . \\
\noindent \textbf{(P\_3,2.1.2)} anupapadaḥ tṛtīyaḥ . \\
\noindent \textbf{(P\_3,2.1.2)} ṇvultṛjacaḥ . \\
\noindent \textbf{(P\_3,2.1.2)} teṣām ṇaḥ . \\
\noindent \textbf{(P\_3,2.1.2)} ṇasya kaḥ . \\
\noindent \textbf{(P\_3,2.1.2)} saḥ yathā eva kaḥ ṇam bādhate evam karmopapadam api bādheta . \\
\noindent \textbf{(P\_3,2.1.2)} karmopapadaḥ api tṛtīyaḥ . \\
\noindent \textbf{(P\_3,2.1.2)} ṇvultṛjacaḥ . \\
\noindent \textbf{(P\_3,2.1.2)} teṣām aṇ . \\
\noindent \textbf{(P\_3,2.1.2)} aṇaḥ kaḥ . \\
\noindent \textbf{(P\_3,2.1.2)} ubhayoḥ tṛtīyayoḥ yuktaḥ vipratiṣedhaḥ . \\
\noindent \textbf{(P\_3,2.1.2)} anupapadasya avakāśaḥ limpati iti limpaḥ . \\
\noindent \textbf{(P\_3,2.1.2)} karmopapadasya saḥ eva . \\
\noindent \textbf{(P\_3,2.1.2)} kuḍyalepe ubhayam prāpnoti . \\
\noindent \textbf{(P\_3,2.1.2)} karmopadaḥ bhavati vipratiṣedhena . \\
\noindent \textbf{(P\_3,2.1.2)} na eṣaḥ yuktaḥ vipratiṣedhaḥ . \\
\noindent \textbf{(P\_3,2.1.2)} anupapadaḥ tṛtīyaḥ . \\
\noindent \textbf{(P\_3,2.1.2)} ṇvultṛjacaḥ . \\
\noindent \textbf{(P\_3,2.1.2)} teṣām kaḥ . \\
\noindent \textbf{(P\_3,2.1.2)} kasya kaḥ . \\
\noindent \textbf{(P\_3,2.1.2)} saḥ yathā eva śaḥ kam bādhate evam karmopapadam api bādheta . \\
\noindent \textbf{(P\_3,2.1.2)} kā tarhi gatiḥ . \\
\noindent \textbf{(P\_3,2.1.2)} madhye apavādāḥ pūrvān vidhīn bādhante iti evam śaḥ kam bādhiṣyate . \\
\noindent \textbf{(P\_3,2.1.2)} karmopapadam na bādhiṣyate . \\
\noindent \textbf{(P\_3,2.1.2)} anupapadasya avakāśaḥ suglaḥ sumlaḥ . \\
\noindent \textbf{(P\_3,2.1.2)} karmopapadasya saḥ eva . \\
\noindent \textbf{(P\_3,2.1.2)} vaḍavāsandāye ubhayam prāpnoti . \\
\noindent \textbf{(P\_3,2.1.2)} karmopadaḥ bhavati vipratiṣedhena . \\
\noindent \textbf{(P\_3,2.1.2)} na eṣaḥ yuktaḥ vipratiṣedhaḥ . \\
\noindent \textbf{(P\_3,2.1.2)} anupapadaḥ tṛtīyaḥ . \\
\noindent \textbf{(P\_3,2.1.2)} ṇvultṛjacaḥ . \\
\noindent \textbf{(P\_3,2.1.2)} teṣām ṇaḥ . \\
\noindent \textbf{(P\_3,2.1.2)} ṇasya kaḥ . \\
\noindent \textbf{(P\_3,2.1.2)} saḥ yathā eva kaḥ ṇam bādhate evam karmopapadam api bādheta . \\
\noindent \textbf{(P\_3,2.1.2)} kā tarhi gatiḥ . \\
\noindent \textbf{(P\_3,2.1.2)} purastād apavādāḥ anantarāñvidhīn bādhante iti evam ayam kaḥ ṇam bādhiṣyate . \\
\noindent \textbf{(P\_3,2.1.2)} karmopapadam na bādhiṣyate . \\
\noindent \textbf{(P\_3,2.1.3)} śīlikāmibhakṣyācaribhyaḥ ṇaḥ pūrvapadaprakṛtisvaratvam ca . \\
\noindent \textbf{(P\_3,2.1.3)} śīlikāmibhakṣyācaribhyaḥ ṇaḥ vaktavyaḥ pūrvapadaprakṛtisvaratvam ca vaktavyam . \\
\noindent \textbf{(P\_3,2.1.3)} śīli . \\
\noindent \textbf{(P\_3,2.1.3)} māṃsaśīlaḥ māṃsaśīlā . \\
\noindent \textbf{(P\_3,2.1.3)} śīli . \\
\noindent \textbf{(P\_3,2.1.3)} kāmi . \\
\noindent \textbf{(P\_3,2.1.3)} māṃsakāmaḥ māṃsakāmā . \\
\noindent \textbf{(P\_3,2.1.3)} kāmi . \\
\noindent \textbf{(P\_3,2.1.3)} bhakṣi . \\
\noindent \textbf{(P\_3,2.1.3)} māṃsabhakṣaḥ māṃsabhakṣā . \\
\noindent \textbf{(P\_3,2.1.3)} bhakṣi ācari . \\
\noindent \textbf{(P\_3,2.1.3)} kalyāṇācāraḥ kalyāṇācārā . \\
\noindent \textbf{(P\_3,2.1.3)} īkṣikṣamibhyām ca . īkṣikṣamibhyām ca iti vaktavyam . \\
\noindent \textbf{(P\_3,2.1.3)} sukhapratīkṣaḥ sukhapratīkṣā . \\
\noindent \textbf{(P\_3,2.1.3)} kalyāṇakṣamaḥ kalyāṇakṣamā . \\
\noindent \textbf{(P\_3,2.1.3)} kimartham idam ucyate . \\
\noindent \textbf{(P\_3,2.1.3)} pūrvapadaprakṛtsvaratvam ca vakṣyāmi īkāraḥ ca mā bhūt iti . \\
\noindent \textbf{(P\_3,2.1.3)} na etat asti prayojanam . \\
\noindent \textbf{(P\_3,2.1.3)} iha yaḥ māṃsam bhakṣayati māṃsam tasya bhakṣaḥ bhavati . \\
\noindent \textbf{(P\_3,2.1.3)} yaḥ asau bhakṣayateḥ ac tadantena bahuvrīhiḥ . \\
\noindent \textbf{(P\_3,2.1.3)} evam tarhi siddhe sati yat karmopapadam ṇam śāsti tat jñāpayati ācāryaḥ samāne arthe kevalam vigrahabhedāt yatra karmopapadaḥ ca prāpnoti bahuvrīhiḥ ca karmopapadaḥ tatra bhavati iti . \\
\noindent \textbf{(P\_3,2.1.3)} kim etasya jñāpane prayojanam . \\
\noindent \textbf{(P\_3,2.1.3)} kāṇḍalāvaḥ . \\
\noindent \textbf{(P\_3,2.1.3)} kāṇḍāni lāvaḥ asya iti bahuvrīhiḥ na bhavati . \\
\noindent \textbf{(P\_3,2.1.3)} bhavati tu bahurvīhiḥ api . \\
\noindent \textbf{(P\_3,2.1.3)} māṃse kāmaḥ asya māṃsakāmaḥ māṃsakāmakaḥ iti vā . \\
\noindent \textbf{(P\_3,2.1.3)} na tu ambhobhigamā . \\
\noindent \textbf{(P\_3,2.1.3)} na tu idam bhavati ambhaḥ abhigamaḥ asyāḥ iti . \\
\noindent \textbf{(P\_3,2.1.3)} kim tarhi . \\
\noindent \textbf{(P\_3,2.1.3)} ambhobhigāmī iti eva bhavati . \\
\noindent \textbf{(P\_3,2.1.3)} kāṇḍalāve api ca vigrahābhāvāt na jñāpakasya prayojanam bhavati iti . \\
\noindent \textbf{(P\_3,2.1.3)} na eṣaḥ asti vigrahaḥ kāṇḍāni lāvaḥ asya iti . \\
\noindent \textbf{(P\_3,2.1.3)} ānnādāya iti ca kṛtām vyatyayaḥ chandasi . \\
\noindent \textbf{(P\_3,2.1.3)} ānnādāya iti ca kṛtām vyatyayaḥ chandasi draṣṭavyaḥ . \\
\noindent \textbf{(P\_3,2.1.3)} annādāya annapataye . \\
\noindent \textbf{(P\_3,2.1.3)} ye āhutim annādīm kṛtvā . \\
\noindent \textbf{(P\_3,2.3)} kavidhau sarvatra prasāraṇibhyaḥ ḍaḥ . \\
\noindent \textbf{(P\_3,2.3)} kavidhau sarvatra prasāraṇibhyaḥ ḍaḥ vaktavyaḥ . \\
\noindent \textbf{(P\_3,2.3)} brahmajyaḥ . \\
\noindent \textbf{(P\_3,2.3)} kim ucyate sarvatra iti . \\
\noindent \textbf{(P\_3,2.3)} anyatra api na avaśyam iha eva . \\
\noindent \textbf{(P\_3,2.3)} hva anyatra . \\
\noindent \textbf{(P\_3,2.3)} āhvaḥ prahvaḥ iti . \\
\noindent \textbf{(P\_3,2.3)} ke hi samprasāraṇaprasaṅgaḥ . \\
\noindent \textbf{(P\_3,2.3)} ke hi sati samprasāraṇam prasajyeta . \\
\noindent \textbf{(P\_3,2.3)} samprasāraṇe kṛte samprasāraṇapūrvatve ca uvaṅādeśe āhuvaḥ iti etat rūpam syāt . \\
\noindent \textbf{(P\_3,2.3)} saḥ tarhi vaktavyaḥ . \\
\noindent \textbf{(P\_3,2.3)} na vaktavyaḥ . \\
\noindent \textbf{(P\_3,2.3)} astu atra samprasāraṇam . \\
\noindent \textbf{(P\_3,2.3)} samprasāraṇe kṛte ākāralopaḥ . \\
\noindent \textbf{(P\_3,2.3)} tasya sthānivadbhāvāt uvaṅādeśaḥ na bhaviṣyati . \\
\noindent \textbf{(P\_3,2.3)} pūrvatve kṛte prāpnoti . \\
\noindent \textbf{(P\_3,2.3)} evam tarhi idam iha sampradhāryam . \\
\noindent \textbf{(P\_3,2.3)} ākāralopaḥ kriyatām pūrvatvam iti . \\
\noindent \textbf{(P\_3,2.3)} kim atra kartavyam . \\
\noindent \textbf{(P\_3,2.3)} paratvāt ākāralopaḥ . \\
\noindent \textbf{(P\_3,2.3)} na sidhyati . \\
\noindent \textbf{(P\_3,2.3)} antaraṅgatvāt pūrvatvam prāpnoti . \\
\noindent \textbf{(P\_3,2.3)} evam tarhi vārṇāt āṅgam balīyaḥ bhavati iti ākāralopaḥ bhaviṣyati . \\
\noindent \textbf{(P\_3,2.3)} evam tarhi idam iha sampradhāryam . \\
\noindent \textbf{(P\_3,2.3)} ākāralopaḥ kriyatām samprasāraṇam iti . \\
\noindent \textbf{(P\_3,2.3)} kim atra kartavyam . \\
\noindent \textbf{(P\_3,2.3)} paratvāt ākāralopaḥ . \\
\noindent \textbf{(P\_3,2.3)} nityam samprasāraṇam . \\
\noindent \textbf{(P\_3,2.3)} kṛte api ākāralope prāpnoti akṛte api prāpnoti . \\
\noindent \textbf{(P\_3,2.3)} ākāralopaḥ api nityaḥ . \\
\noindent \textbf{(P\_3,2.3)} kṛte api samprasāraṇe prāpnoti akṛte api prāpnoti . \\
\noindent \textbf{(P\_3,2.3)} anityaḥ ākāralopaḥ . \\
\noindent \textbf{(P\_3,2.3)} na hi kṛte samprasāraṇe prāpnoti . \\
\noindent \textbf{(P\_3,2.3)} antaraṅgam hi pūrvatvam bhādhate . \\
\noindent \textbf{(P\_3,2.3)} yasya lakṣaṇāntareṇa nimittam vihanyate na tat anityam . \\
\noindent \textbf{(P\_3,2.3)} na ca samprasāraṇam eva ākāralopasya nimittam hanti . \\
\noindent \textbf{(P\_3,2.3)} avaśyam lakṣaṇāntaram pūrvatvam pratīkṣyam . \\
\noindent \textbf{(P\_3,2.3)} ubhayoḥ nityayoḥ paratvāt ākāralopaḥ . \\
\noindent \textbf{(P\_3,2.3)} ākāralope kṛte samprasāraṇam . \\
\noindent \textbf{(P\_3,2.3)} samprasāraṇe kṛte yaṇādeśe siddham rūpam āhvaḥ prahvaḥ iti . \\
\noindent \textbf{(P\_3,2.3)} evam api na sidhyati . \\
\noindent \textbf{(P\_3,2.3)} yaḥ anādiṣṭād acaḥ pūrvaḥ tasya vidhim prati sthānivadbhāvaḥ . \\
\noindent \textbf{(P\_3,2.3)} ādiṣṭāt ca eṣaḥ acaḥ pūrvaḥ bhavati . \\
\noindent \textbf{(P\_3,2.3)} evam tarhi ākāralopasya asiddhatvāt uvaṅādeśaḥ na bhaviṣyati . \\
\noindent \textbf{(P\_3,2.3)} iha api tarhi ākāralopasya asiddhatvāt uvaṅādeśaḥ na syāt . \\
\noindent \textbf{(P\_3,2.3)} juhuvatuḥ jhuhuvuḥ iti . \\
\noindent \textbf{(P\_3,2.3)} asti atra viśeṣaḥ . \\
\noindent \textbf{(P\_3,2.3)} akṛte atra āttve pūrvatvam bhavati . \\
\noindent \textbf{(P\_3,2.3)} idam iha sampradhāryam āttvam kriyatām pūrvatvam iti . \\
\noindent \textbf{(P\_3,2.3)} kim atra kartavyam . \\
\noindent \textbf{(P\_3,2.3)} paratvāt pūrvatvam . \\
\noindent \textbf{(P\_3,2.3)} na sidhyati . \\
\noindent \textbf{(P\_3,2.3)} antaraṅgatvāt āttvam prāpnoti . \\
\noindent \textbf{(P\_3,2.3)} evam tarhi idam iha sampradhāryam . \\
\noindent \textbf{(P\_3,2.3)} āttvam kriyatām samprasāraṇam iti . \\
\noindent \textbf{(P\_3,2.3)} kim atra kartavyam . \\
\noindent \textbf{(P\_3,2.3)} paratvāt āttvam . \\
\noindent \textbf{(P\_3,2.3)} nityam samprasāraṇam . \\
\noindent \textbf{(P\_3,2.3)} kṛte api āttve prāpnoti akṛte api . \\
\noindent \textbf{(P\_3,2.3)} āttvam api nityam . \\
\noindent \textbf{(P\_3,2.3)} kṛte api samprasāraṇe prāpnoti akṛte api prāpnoti . \\
\noindent \textbf{(P\_3,2.3)} anityam āttvam . \\
\noindent \textbf{(P\_3,2.3)} na hi samprasāraṇe kṛte prāpnoti . \\
\noindent \textbf{(P\_3,2.3)} paratvāt pūrvatvena eva bhavitavyam . \\
\noindent \textbf{(P\_3,2.3)} yasya lakṣaṇāntareṇa nimittam vihanyate na tat anityam . \\
\noindent \textbf{(P\_3,2.3)} na ca samprasāraṇam eva āttvasya nimittam vihanti . \\
\noindent \textbf{(P\_3,2.3)} avaśyam lakṣaṇāntaram pūrvatvam pratīkṣyam . \\
\noindent \textbf{(P\_3,2.3)} ubhayoḥ nityayoḥ paratvāt ātttve kṛte samprasāraṇam . \\
\noindent \textbf{(P\_3,2.3)} evam tarhi pūrvatve yogavibhāgaḥ kariṣyate . \\
\noindent \textbf{(P\_3,2.3)} samprasāraṇāt paraḥ pūrvaḥ bhavati . \\
\noindent \textbf{(P\_3,2.3)} tata eṅaḥ . \\
\noindent \textbf{(P\_3,2.3)} eṅaḥ ca samprasāraṇāt pūrvaḥ bhavati . \\
\noindent \textbf{(P\_3,2.3)} kimartham idam . \\
\noindent \textbf{(P\_3,2.3)} akṛte āttve pūrvatvam yathā syāt . \\
\noindent \textbf{(P\_3,2.3)} tataḥ padāntāt ati . \\
\noindent \textbf{(P\_3,2.3)} eṅaḥ iti eva . \\
\noindent \textbf{(P\_3,2.3)} iha api tarhi akṛte āttve pūrvatvam syāt . \\
\noindent \textbf{(P\_3,2.3)} āhvaḥ prahvaḥ iti . \\
\noindent \textbf{(P\_3,2.3)} asti atra viśeṣaḥ . \\
\noindent \textbf{(P\_3,2.3)} ākārāntalakṣaṇaḥ kavidhiḥ . \\
\noindent \textbf{(P\_3,2.3)} tena anena avaśyam āttvam pratīkṣyam . \\
\noindent \textbf{(P\_3,2.3)} liṭ punaḥ aviśeṣeṇa dhātumātrāt vidhīyate . \\
\noindent \textbf{(P\_3,2.3)} nityam prasāraṇam . \\
\noindent \textbf{(P\_3,2.3)} hvaḥ yaṇ . \\
\noindent \textbf{(P\_3,2.3)} vārṇāt āṅgam na pūrvatvam . \\
\noindent \textbf{(P\_3,2.3)} yaḥ anādiṣṭāt acaḥ pūrvaḥ tatkārye sthānivattvam hi provāca bhagavān kātyaḥ . \\
\noindent \textbf{(P\_3,2.3)} tena asiddhiḥ yaṇaḥ te . \\
\noindent \textbf{(P\_3,2.3)} ātaḥ kaḥ . \\
\noindent \textbf{(P\_3,2.3)} liṭ na . \\
\noindent \textbf{(P\_3,2.3)} eṅaḥ pūrvaḥ . \\
\noindent \textbf{(P\_3,2.3)} siddhaḥ āhvaḥ tathā sati . \\
\noindent \textbf{(P\_3,2.4)} supi sthaḥ bhāve ca . \\
\noindent \textbf{(P\_3,2.4)} supi sthaḥ iti atra bhāve ca iti vaktavyam . \\
\noindent \textbf{(P\_3,2.4)} iha api yathā syāt . \\
\noindent \textbf{(P\_3,2.4)} ākhūtthaḥ vartate . \\
\noindent \textbf{(P\_3,2.4)} śyenotthaḥ śalabhotthaḥ . \\
\noindent \textbf{(P\_3,2.4)} tat tarhi vaktavyam . \\
\noindent \textbf{(P\_3,2.4)} na vaktavyam . \\
\noindent \textbf{(P\_3,2.4)} yogavibhāgāt siddham . \\
\noindent \textbf{(P\_3,2.4)} yogavibhāgaḥ kariṣyate . \\
\noindent \textbf{(P\_3,2.4)} ātaḥ anupasarge kaḥ bhavati . \\
\noindent \textbf{(P\_3,2.4)} tataḥ supi . \\
\noindent \textbf{(P\_3,2.4)} supi ca ataḥ kaḥ bhavati . \\
\noindent \textbf{(P\_3,2.4)} kacchena pibati kacchapaḥ . \\
\noindent \textbf{(P\_3,2.4)} kaṭāhena pibati kaṭāhaḥ . \\
\noindent \textbf{(P\_3,2.4)} dvābhyām pibati dvipaḥ . \\
\noindent \textbf{(P\_3,2.4)} tataḥ sthaḥ . \\
\noindent \textbf{(P\_3,2.4)} sthaḥ ca kaḥ bhavati supi iti . \\
\noindent \textbf{(P\_3,2.4)} kimartham idam . \\
\noindent \textbf{(P\_3,2.4)} bhāve yathā syāt . \\
\noindent \textbf{(P\_3,2.4)} kutaḥ nu khalvu etat bhāve bhaviṣyati na punaḥ karmādiṣu kārakeṣu iti . \\
\noindent \textbf{(P\_3,2.4)} yogavibhāgāt ayam kartuḥ apakṛṣyate . \\
\noindent \textbf{(P\_3,2.4)} na ca anyasmin arthe ādiśyate . \\
\noindent \textbf{(P\_3,2.4)} anirdiṣṭārthāḥ pratyayāḥ svārthe bhavanti iti svārthe bhaviṣyanti . \\
\noindent \textbf{(P\_3,2.4)} tat yathā guptijkidbhyaḥ san yāvādibhyaḥ kan . \\
\noindent \textbf{(P\_3,2.4)} saḥ asau svārthe bhavan bhāve bhaviṣyati . \\
\noindent \textbf{(P\_3,2.5)} tundaśokayoḥ parimṛjāpanudoḥ ālasyasukhāharaṇayoḥ . \\
\noindent \textbf{(P\_3,2.5)} tundaśokayoḥ parimṛjāpanudoḥ iti atra ālasyasukhāharaṇayoḥ iti vaktavyam . \\
\noindent \textbf{(P\_3,2.5)} tundaparimṛjaḥ alasaḥ . \\
\noindent \textbf{(P\_3,2.5)} śokāpanudaḥ putraḥ jātaḥ . \\
\noindent \textbf{(P\_3,2.5)} yaḥ hi tundam parimārṣṭi tundaparimārjaḥ saḥ bhavati . \\
\noindent \textbf{(P\_3,2.5)} yaḥ ca śokam apanudati śokāpanodaḥ saḥ bhavati . \\
\noindent \textbf{(P\_3,2.5)} kaprakaraṇe mūliavibhujādibhyaḥ upasaṅkhyānam . \\
\noindent \textbf{(P\_3,2.5)} kaprakaraṇe mūliavibhujādibhyaḥ upasaṅkhyānam kartavyam . \\
\noindent \textbf{(P\_3,2.5)} mūlavibhujaḥ rathaḥ . \\
\noindent \textbf{(P\_3,2.5)} nakhamucāni dhanūṃṣi . \\
\noindent \textbf{(P\_3,2.5)} kākaguhāḥ tilāḥ . \\
\noindent \textbf{(P\_3,2.5)} sarasīruham kumudam . \\
\noindent \textbf{(P\_3,2.8)} surāsīdhvoḥ pibateḥ . \\
\noindent \textbf{(P\_3,2.8)} surāsīdhvoḥ pibateḥ iti vaktavyam . \\
\noindent \textbf{(P\_3,2.8)} iha mā bhūt . \\
\noindent \textbf{(P\_3,2.8)} kṣīrapā brāhmaṇī iti . \\
\noindent \textbf{(P\_3,2.8)} pibateḥ iti kimartham . \\
\noindent \textbf{(P\_3,2.8)} yā hi surām pāti surāpā sā bhavati . \\
\noindent \textbf{(P\_3,2.8)} bahulam taṇi . \\
\noindent \textbf{(P\_3,2.8)} bahulam taṇi iti vaktavyam . \\
\noindent \textbf{(P\_3,2.8)} kim idam taṇi iti . \\
\noindent \textbf{(P\_3,2.8)} sañjñāchandasoḥ grahaṇam . \\
\noindent \textbf{(P\_3,2.8)} yā brāhmaṇī surāpī bhavati na enām devāḥ patilokam nayanti . \\
\noindent \textbf{(P\_3,2.8)} yā brāhmaṇī surāpā bhavati na enām devāḥ patilokam nayanti . \\
\noindent \textbf{(P\_3,2.9)} acprakaraṇe śaktilāṅgalāṅkuśayaṣṭitomaraghaṭaghaṭīdhanuḥṣu ghraheḥ upasaṅkhyānam . \\
\noindent \textbf{(P\_3,2.9)} acprakaraṇe śaktilāṅgalāṅkuśayaṣṭitomaraghaṭaghaṭīdhanuḥṣu ghraheḥ upasaṅkhyānam kartavyam . \\
\noindent \textbf{(P\_3,2.9)} śaktigrahaḥ . \\
\noindent \textbf{(P\_3,2.9)} śakti . \\
\noindent \textbf{(P\_3,2.9)} lāṅgala . \\
\noindent \textbf{(P\_3,2.9)} lāṅgalagrahaḥ . \\
\noindent \textbf{(P\_3,2.9)} lāṅgala . \\
\noindent \textbf{(P\_3,2.9)} āṅkuśa . \\
\noindent \textbf{(P\_3,2.9)} āṅkuśagrahaḥ . \\
\noindent \textbf{(P\_3,2.9)} āṅkuśa . \\
\noindent \textbf{(P\_3,2.9)} yaṣṭi . \\
\noindent \textbf{(P\_3,2.9)} yaṣṭigrahaḥ . \\
\noindent \textbf{(P\_3,2.9)} yaṣṭi . \\
\noindent \textbf{(P\_3,2.9)} tomara . \\
\noindent \textbf{(P\_3,2.9)} tomaragrahaḥ . \\
\noindent \textbf{(P\_3,2.9)} tomara . \\
\noindent \textbf{(P\_3,2.9)} ghaṭa . \\
\noindent \textbf{(P\_3,2.9)} ghaṭagrahaḥ . \\
\noindent \textbf{(P\_3,2.9)} ghaṭa . \\
\noindent \textbf{(P\_3,2.9)} ghaṭī . \\
\noindent \textbf{(P\_3,2.9)} ghaṭīgrahaḥ . \\
\noindent \textbf{(P\_3,2.9)} ghaṭī . \\
\noindent \textbf{(P\_3,2.9)} dhanus . \\
\noindent \textbf{(P\_3,2.9)} dhanurgrahaḥ . \\
\noindent \textbf{(P\_3,2.9)} dhanus . \\
\noindent \textbf{(P\_3,2.9)} sūtre ca dhāryarthe . \\
\noindent \textbf{(P\_3,2.9)} sūtre ca dhāryarthe graheḥ upasaṅkhyānam . \\
\noindent \textbf{(P\_3,2.9)} sūtragrahaḥ . \\
\noindent \textbf{(P\_3,2.9)} dhāryarthe iti kimartham . \\
\noindent \textbf{(P\_3,2.9)} yaḥ hi sūtram gṛhṇāti sūtragrāhaṇ saḥ bhavati . \\
\noindent \textbf{(P\_3,2.13)} stambakarṇayoḥ hastisūcakayoḥ . \\
\noindent \textbf{(P\_3,2.13)} stambakarṇayoḥ iti atra hastisūcakayoḥ iti vaktavyam . \\
\noindent \textbf{(P\_3,2.13)} stamberamaḥ hastī . \\
\noindent \textbf{(P\_3,2.13)} karṇejapaḥ sūcakaḥ . \\
\noindent \textbf{(P\_3,2.13)} sambe rantā karṇe japitā iti eva anyatra . \\
\noindent \textbf{(P\_3,2.14)} dhātugrahaṇam kimartham . \\
\noindent \textbf{(P\_3,2.14)} śami sañjñayām dhātugrahaṇam kṛñaḥ hetvādiṣu ṭapratiṣedhārtham . \\
\noindent \textbf{(P\_3,2.14)} śami sañjñayām dhātugrahaṇam kriyate kṛñaḥ hetvādiṣu ṭaḥ mā bhūt iti . \\
\noindent \textbf{(P\_3,2.14)} śami sañjñayām ac bhavati iti asya avakāśaḥ śamvadaḥ śambhavaḥ . \\
\noindent \textbf{(P\_3,2.14)} ṭasya avakāśaḥ śrāddhakaraḥ piṇḍakaraḥ . \\
\noindent \textbf{(P\_3,2.14)} śaṅkarā nāma parivrājikā . \\
\noindent \textbf{(P\_3,2.14)} śaṅkarā śakunikā tacchilā ca . \\
\noindent \textbf{(P\_3,2.14)} tasyām ubhayam prāpnoti . \\
\noindent \textbf{(P\_3,2.14)} paratvāt ṭaḥ syāt . \\
\noindent \textbf{(P\_3,2.14)} dhātugrahaṇasāmarthyāt ac eva bhavati . \\
\noindent \textbf{(P\_3,2.14)} kuṇaravāḍavaḥ tu āha . \\
\noindent \textbf{(P\_3,2.14)} na eṣā śaṅkarā . \\
\noindent \textbf{(P\_3,2.14)} śaṅgarā eṣā . \\
\noindent \textbf{(P\_3,2.14)} gṛṇātiḥ śabdakarmā . \\
\noindent \textbf{(P\_3,2.14)} tasya eṣaḥ prayogaḥ . \\
\noindent \textbf{(P\_3,2.15)} adhikaraṇe śeteḥ pārśvādiṣu upasaṅkhyānam. adhikaraṇe śeteḥ pārśvādiṣu upasaṅkhyānam kartavyam . \\
\noindent \textbf{(P\_3,2.15)} pārśvaśayaḥ pṛṣṭhaśayaḥ udaraśayaḥ . \\
\noindent \textbf{(P\_3,2.15)} digdhasahapūrvāt ca . \\
\noindent \textbf{(P\_3,2.15)} didghasahapūrvāt ca iti vaktavyam . \\
\noindent \textbf{(P\_3,2.15)} digdhasahaśayaḥ . \\
\noindent \textbf{(P\_3,2.15)} uttānādiṣu kartṛṣu . \\
\noindent \textbf{(P\_3,2.15)} uttānādiṣu kartṛṣu iti vaktavyam . \\
\noindent \textbf{(P\_3,2.15)} uttānaśayaḥ avamūrdhaśayaḥ . \\
\noindent \textbf{(P\_3,2.15)} girau ḍaḥ chandasi . \\
\noindent \textbf{(P\_3,2.15)} girau upapade ḍaḥ chandasi vaktavyaḥ . \\
\noindent \textbf{(P\_3,2.15)} girau śete giriśaḥ . \\
\noindent \textbf{(P\_3,2.15)} taddhitaḥ vā . \\
\noindent \textbf{(P\_3,2.15)} taddhitaḥ vā punaḥ eṣaḥ bhavati . \\
\noindent \textbf{(P\_3,2.15)} girau śete giriśaḥ iti . \\
\noindent \textbf{(P\_3,2.16)} iha kasmāt na bhavati . \\
\noindent \textbf{(P\_3,2.16)} kurūn carati . \\
\noindent \textbf{(P\_3,2.16)} pañcālān carati iti . \\
\noindent \textbf{(P\_3,2.16)} adhikaraṇe iti vartate . \\
\noindent \textbf{(P\_3,2.16)} nanu ca karmaṇi iti api vartate . \\
\noindent \textbf{(P\_3,2.16)} tatra kutaḥ etat . \\
\noindent \textbf{(P\_3,2.16)} adhikaraṇe bhaviṣyati na punaḥ karmaṇi iti . \\
\noindent \textbf{(P\_3,2.16)} careḥ bhikṣāgrahaṇam jñapakam karmaṇi aprasaṅgaḥ . yat ayam bhikṣāsenādāyeṣu ca iti careḥ bhikṣāgrahaṇam karoti tat jñāpayati ācāryaḥ na bhavati karmaṇi iti . \\
\noindent \textbf{(P\_3,2.21)} kiṃyattadbahuṣu kṛñaḥ ajvidhānam . \\
\noindent \textbf{(P\_3,2.21)} kiṃyattadbahuṣu kṛñaḥ ajvidhānam kartavyam . \\
\noindent \textbf{(P\_3,2.21)} kiṅkarā . \\
\noindent \textbf{(P\_3,2.21)} kim . \\
\noindent \textbf{(P\_3,2.21)} yat . \\
\noindent \textbf{(P\_3,2.21)} yatkarā . \\
\noindent \textbf{(P\_3,2.21)} yat . \\
\noindent \textbf{(P\_3,2.21)} tat . \\
\noindent \textbf{(P\_3,2.21)} tatkarā . \\
\noindent \textbf{(P\_3,2.21)} tat . \\
\noindent \textbf{(P\_3,2.21)} bahu . \\
\noindent \textbf{(P\_3,2.21)} bahukarā . \\
\noindent \textbf{(P\_3,2.24)} stambaśakṛtoḥ vrīhivatsayoḥ . \\
\noindent \textbf{(P\_3,2.24)} vrīhivatsayoḥ iti vaktavyam . \\
\noindent \textbf{(P\_3,2.24)} stambakariḥ vrīhiḥ . \\
\noindent \textbf{(P\_3,2.24)} śakṛtkariḥ vatsaḥ . \\
\noindent \textbf{(P\_3,2.26)} ātmambhariḥ iti kim nipātyate . \\
\noindent \textbf{(P\_3,2.26)} ātmanaḥ mum bhṛñaḥ ca inpratyayaḥ . \\
\noindent \textbf{(P\_3,2.26)} atyalpam idam ucyate . \\
\noindent \textbf{(P\_3,2.26)} bhṛñaḥ kukṣyātmanoḥ mum ca . \\
\noindent \textbf{(P\_3,2.26)} bhṛñaḥ kukṣyātmanoḥ mum ca iti vaktavyam . \\
\noindent \textbf{(P\_3,2.26)} kukṣimbharaḥ . \\
\noindent \textbf{(P\_3,2.26)} ātmambhariḥ carati yūtham asevamānaḥ . \\
\noindent \textbf{(P\_3,2.28)} khaśprakaraṇe vātasunītilaśardheṣu ajadheṭtudajahātibhyaḥ . khaśprakaraṇe vātasunītilaśardheṣu ajadheṭtudajahātibhyaḥ iti vaktavyam . \\
\noindent \textbf{(P\_3,2.28)} vātamajāḥ mṛgāḥ . \\
\noindent \textbf{(P\_3,2.28)} vāta . \\
\noindent \textbf{(P\_3,2.28)} śunī . \\
\noindent \textbf{(P\_3,2.28)} śunīndhayaḥ . \\
\noindent \textbf{(P\_3,2.28)} śunī . \\
\noindent \textbf{(P\_3,2.28)} tila . \\
\noindent \textbf{(P\_3,2.28)} tilandtudaḥ . \\
\noindent \textbf{(P\_3,2.28)} tila . \\
\noindent \textbf{(P\_3,2.28)} śardha . \\
\noindent \textbf{(P\_3,2.28)} śardhañjahāḥ māṣāḥ . \\
\noindent \textbf{(P\_3,2.29)} stane dheṭaḥ . \\
\noindent \textbf{(P\_3,2.29)} stane dheṭaḥ iti vaktavyam . \\
\noindent \textbf{(P\_3,2.29)} stanandhayaḥ . \\
\noindent \textbf{(P\_3,2.29)} tataḥ muṣṭau dhmaḥ ca . \\
\noindent \textbf{(P\_3,2.29)} muṣṭau dhmaḥ ca dheṭaḥ ca iti vaktavyam . \\
\noindent \textbf{(P\_3,2.29)} muṣṭindhamaḥ muṣṭidhayaḥ . \\
\noindent \textbf{(P\_3,2.29)} atayalpam idam ucyate . \\
\noindent \textbf{(P\_3,2.29)} nāsikānāḍīmuṣṭighaṭīkhārīṣu iti vaktavyam . \\
\noindent \textbf{(P\_3,2.29)} nāsikandhamaḥ nāsikandhayaḥ . \\
\noindent \textbf{(P\_3,2.29)} nāsika . \\
\noindent \textbf{(P\_3,2.29)} nāḍī . \\
\noindent \textbf{(P\_3,2.29)} nāḍindhamaḥ nāḍindhayaḥ . \\
\noindent \textbf{(P\_3,2.29)} nāḍī . \\
\noindent \textbf{(P\_3,2.29)} muṣṭi . \\
\noindent \textbf{(P\_3,2.29)} muṣṭindhamaḥ muṣṭidhayaḥ . \\
\noindent \textbf{(P\_3,2.29)} muṣṭi . \\
\noindent \textbf{(P\_3,2.29)} ghaṭī . \\
\noindent \textbf{(P\_3,2.29)} ghaṭindhamaḥ ghaṭindhayaḥ . \\
\noindent \textbf{(P\_3,2.29)} ghaṭī . \\
\noindent \textbf{(P\_3,2.29)} khārī . \\
\noindent \textbf{(P\_3,2.29)} khārindhamaḥ khārindhayaḥ . \\
\noindent \textbf{(P\_3,2.29)} khārī . \\
\noindent \textbf{(P\_3,2.38)} khacprakaraṇe gameḥ supi upasaṅkhyānam . \\
\noindent \textbf{(P\_3,2.38)} khacprakaraṇe gameḥ supi upasaṅkhyānam . \\
\noindent \textbf{(P\_3,2.38)} mitaṅgamaḥ . \\
\noindent \textbf{(P\_3,2.38)} mitaṅgamā hastinī . \\
\noindent \textbf{(P\_3,2.38)} vihāyasaḥ viha ca . \\
\noindent \textbf{(P\_3,2.38)} vihāyasaḥ viha iti ayam ādeśaḥ vaktavyaḥ . \\
\noindent \textbf{(P\_3,2.38)} khac ca . \\
\noindent \textbf{(P\_3,2.38)} vihaṅgamaḥ . \\
\noindent \textbf{(P\_3,2.38)} khac ca ḍit vā . \\
\noindent \textbf{(P\_3,2.38)} khac ca ḍit vā vaktavyaḥ . \\
\noindent \textbf{(P\_3,2.38)} vihaṅgaḥ . \\
\noindent \textbf{(P\_3,2.38)} ḍe ca . \\
\noindent \textbf{(P\_3,2.38)} ḍe ca vihāyasaḥ viha iti ayam ādeśaḥ vaktavyaḥ . \\
\noindent \textbf{(P\_3,2.38)} vihagaḥ . \\
\noindent \textbf{(P\_3,2.48)} ḍaprakaraṇe sarvatrapannayoḥ upasaṅkhyānam . \\
\noindent \textbf{(P\_3,2.48)} ḍaprakaraṇe sarvatrapannayoḥ upasaṅkhyānam kartavyam . \\
\noindent \textbf{(P\_3,2.48)} sarvatragaḥ pannagaḥ . \\
\noindent \textbf{(P\_3,2.48)} urasaḥ lopaḥ ca . \\
\noindent \textbf{(P\_3,2.48)} urasaḥ lopaḥ ca vaktavyaḥ . \\
\noindent \textbf{(P\_3,2.48)} uragaḥ . \\
\noindent \textbf{(P\_3,2.48)} suduroḥ adhikaraṇe . \\
\noindent \textbf{(P\_3,2.48)} suduroḥ adhikaraṇe ḍaḥ vaktavyaḥ . \\
\noindent \textbf{(P\_3,2.48)} sugaḥ durgaḥ . \\
\noindent \textbf{(P\_3,2.48)} nisaḥ deśe . \\
\noindent \textbf{(P\_3,2.48)} nisaḥ deśe ḍaḥ vaktavyaḥ . \\
\noindent \textbf{(P\_3,2.48)} nirgaḥ . \\
\noindent \textbf{(P\_3,2.48)} apara āha . \\
\noindent \textbf{(P\_3,2.48)} ḍaprakaraṇe anyeṣu api dṛśyate . \\
\noindent \textbf{(P\_3,2.48)} ḍaprakaraṇe anyeṣu api dṛśyate iti vaktavyam . \\
\noindent \textbf{(P\_3,2.48)} tataḥ stryagāragaḥ . \\
\noindent \textbf{(P\_3,2.48)} aśnute yāvat annāya grāmagaḥ . \\
\noindent \textbf{(P\_3,2.48)} dhvaṃsate gurutalpagaḥ . \\
\noindent \textbf{(P\_3,2.49)} dārau āhanaḥ aṇ antyasya ca ṭaḥ sañjñāyām . \\
\noindent \textbf{(P\_3,2.49)} dārau upapade āṅpūrvāt hanteḥ aṇ vaktavyaḥ antyasya ca ṭaḥ vaktavyaḥ . \\
\noindent \textbf{(P\_3,2.49)} dārvāghāṭaḥ te vanaspatīnām . \\
\noindent \textbf{(P\_3,2.49)} cārau vā . \\
\noindent \textbf{(P\_3,2.49)} cārau upapade āṅpūrvāt hanteḥ aṇ vaktavyaḥ antyasya ca ṭaḥ vā vaktavyaḥ . \\
\noindent \textbf{(P\_3,2.49)} cārvāghāṭaḥ cārvāghātaḥ . \\
\noindent \textbf{(P\_3,2.49)} karmaṇi sami ca . \\
\noindent \textbf{(P\_3,2.49)} karmaṇi upapade sapūrvāt hanteḥ aṇ vaktavyaḥ antyasya ca ṭaḥ vā vaktavyaḥ . \\
\noindent \textbf{(P\_3,2.49)} varṇasaṅghāṭaḥ varṇasaṅghātaḥ padasaṅghāṭaḥ padasaṅghātaḥ . \\
\noindent \textbf{(P\_3,2.52)} katham idam vijñāyate . \\
\noindent \textbf{(P\_3,2.52)} lakṣaṇe kartari iti āhosvit lakṣaṇavati kartari iti . \\
\noindent \textbf{(P\_3,2.52)} kim ca ataḥ . \\
\noindent \textbf{(P\_3,2.52)} yadi vijñāyate lakṣaṇe kartari iti siddham jāyāghnaḥ tilakālakaḥ patighnī pāṇilekhā iti . \\
\noindent \textbf{(P\_3,2.52)} jāyāghnaḥ tilakālakaḥ patighnī pāṇirekhā iti . \\
\noindent \textbf{(P\_3,2.52)} jāyāghnaḥ brāhmaṇaḥ patighnī vrṣalī iti na sidhyati . \\
\noindent \textbf{(P\_3,2.52)} atha vijñāyate lakṣaṇavati kartari iti siddham jāyāghnaḥ brāhmaṇaḥ patighnī vrṣalī iti . \\
\noindent \textbf{(P\_3,2.52)} jāyāghnaḥ tilakālakaḥ patighnī pāṇilekhā iti na sidhyati . \\
\noindent \textbf{(P\_3,2.52)} astu lakṣaṇe kartari iti . \\
\noindent \textbf{(P\_3,2.52)} katham jāyāghnaḥ , brāhmaṇaḥ patighnī vrṣalī iti . \\
\noindent \textbf{(P\_3,2.52)} akāraḥ matvarthīyaḥ . \\
\noindent \textbf{(P\_3,2.52)} jāyāghnaḥ asmin asti iti saḥ ayam jāyāghnaḥ . \\
\noindent \textbf{(P\_3,2.52)} patighnīvṛṣalī iti na sidhyati . \\
\noindent \textbf{(P\_3,2.52)} astu tarhi lakṣaṇavati kartari iti . \\
\noindent \textbf{(P\_3,2.52)} katham jāyāghnaḥ tilakālakaḥ patighnī pāṇilekhā iti . \\
\noindent \textbf{(P\_3,2.52)} amanuṣyakartṛke iti evam bhaviṣyati . \\
\noindent \textbf{(P\_3,2.53)} apraṇikartṛke iti vaktavyam . \\
\noindent \textbf{(P\_3,2.53)} iha mā bhūt . \\
\noindent \textbf{(P\_3,2.53)} nagaraghātaḥ hastī . \\
\noindent \textbf{(P\_3,2.53)} yadi apraṇikartṛke iti ucyate śaśaghnī śakuniḥ iti na sidhyati . \\
\noindent \textbf{(P\_3,2.53)} astu tarhi amanuṣyakartṛke iti eva . \\
\noindent \textbf{(P\_3,2.53)} katham nagaraghātaḥ hastī . \\
\noindent \textbf{(P\_3,2.53)} kṛtyalyuṭaḥ bahulam iti evam atra aṇ bhaviṣyati . \\
\noindent \textbf{(P\_3,2.55)} rājaghe upasaṅkhyānam . \\
\noindent \textbf{(P\_3,2.55)} rājaghe upasaṅkhyānam kartavyam . \\
\noindent \textbf{(P\_3,2.55)} rājaghaḥ . \\
\noindent \textbf{(P\_3,2.56)} khyuni cvipratiṣedhānarthakyam lyuṭkhynoḥ aviśeṣāt . \\
\noindent \textbf{(P\_3,2.56)} khyuni cvipratiṣedhaḥ anarthakaḥ . \\
\noindent \textbf{(P\_3,2.56)} kim kāraṇam . \\
\noindent \textbf{(P\_3,2.56)} lyuṭkhynoḥ aviśeṣāt . \\
\noindent \textbf{(P\_3,2.56)} khyunā mukte lyuṭā bhavitavyam . \\
\noindent \textbf{(P\_3,2.56)} na ca asti viśeṣaḥ cvyante upapade lyuṭaḥ khyunaḥ vā . \\
\noindent \textbf{(P\_3,2.56)} tat eva rūpam saḥ eva svaraḥ . \\
\noindent \textbf{(P\_3,2.56)} ayam asti viśeṣaḥ . \\
\noindent \textbf{(P\_3,2.56)} lyuṭi sati īkāreṇa bhavitavyam . \\
\noindent \textbf{(P\_3,2.56)} khyuni sati na bhavitavyam . \\
\noindent \textbf{(P\_3,2.56)} khyuni api sati bhavatavyam . \\
\noindent \textbf{(P\_3,2.56)} evam hi saunāgāḥ paṭhanti . \\
\noindent \textbf{(P\_3,2.56)} nañsnañīkkhyuṃstaruṇatalunānām upasaṅkhyānam iti . \\
\noindent \textbf{(P\_3,2.56)} ayam tarhi viśeṣaḥ . \\
\noindent \textbf{(P\_3,2.56)} khyuni sati nityasamāsena bhavitavyam . \\
\noindent \textbf{(P\_3,2.56)} upapadasamāsaḥ hi nityasamāsaḥ iti . \\
\noindent \textbf{(P\_3,2.56)} lyuṭi sati na bhavitavyam . \\
\noindent \textbf{(P\_3,2.56)} lyuṭi api bhavitavyam . \\
\noindent \textbf{(P\_3,2.56)} gatisamāsaḥ api hi nityasamāsaḥ . \\
\noindent \textbf{(P\_3,2.56)} cyvantam ca gatisañjñam bhavati . \\
\noindent \textbf{(P\_3,2.56)} mumartham tarhi pratiṣedhaḥ vaktavyaḥ . \\
\noindent \textbf{(P\_3,2.56)} khyuni sati mumā bhavitavyam . \\
\noindent \textbf{(P\_3,2.56)} lyuti sati na bhavitavyam . \\
\noindent \textbf{(P\_3,2.56)} mumartham iti cet na avyayatvāt . \\
\noindent \textbf{(P\_3,2.56)} mumartham iti cet tat na . \\
\noindent \textbf{(P\_3,2.56)} kim kāraṇam . \\
\noindent \textbf{(P\_3,2.56)} avyayatvāt . \\
\noindent \textbf{(P\_3,2.56)} anavyayasya mum ucyate . \\
\noindent \textbf{(P\_3,2.56)} cvyantam ca avyayasañjñam . \\
\noindent \textbf{(P\_3,2.56)} uttarārtham tu . \\
\noindent \textbf{(P\_3,2.56)} uttarārtham tarhi pratiṣedhaḥ vaktavyaḥ . \\
\noindent \textbf{(P\_3,2.56)} kartari bhuvaḥ khiṣṇuckhukañau acvau iti eva . \\
\noindent \textbf{(P\_3,2.56)} āḍhyībhavitā . \\
\noindent \textbf{(P\_3,2.56)} atha idānīm anena mukte tācchīlilaḥ iṣṇuc vidhīyate . \\
\noindent \textbf{(P\_3,2.56)} saḥ atra kasmāt na bhavati . \\
\noindent \textbf{(P\_3,2.56)} rūḍhiśabdaprakārāḥ tacchīlikāḥ . \\
\noindent \textbf{(P\_3,2.56)} na ca rūḍḥiśabdāḥ gatibhiḥ viśeṣyante . \\
\noindent \textbf{(P\_3,2.56)} na hi bhavati pradevadattaḥ iti . \\
\noindent \textbf{(P\_3,2.57)} kimartham khiṣṇuc ikārādiḥ kriyate na ksnuḥ iti eva ucyeta . \\
\noindent \textbf{(P\_3,2.57)} tatra ayam api arthaḥ . \\
\noindent \textbf{(P\_3,2.57)} svarārthaḥ cakāraḥ na kartavyaḥ bhavati . \\
\noindent \textbf{(P\_3,2.57)} kena idānīm ikārāditvam kriyate . \\
\noindent \textbf{(P\_3,2.57)} iṣṇucaḥ ikārāditvam udāttatvāt kṛtam bhuvaḥ . \\
\noindent \textbf{(P\_3,2.57)} bhavateḥ udāttatvāt ikārāditvam bhaviṣyati . \\
\noindent \textbf{(P\_3,2.57)} idam tarhi prayojanam . \\
\noindent \textbf{(P\_3,2.57)} khit ayam kriyate . \\
\noindent \textbf{(P\_3,2.57)} tatra cartve kṛte syāt . \\
\noindent \textbf{(P\_3,2.57)} kit vā khit vā iti . \\
\noindent \textbf{(P\_3,2.57)} sandehamātram etat bhavati . \\
\noindent \textbf{(P\_3,2.57)} sarvasandeheṣu ca idam upatiṣṭhate . \\
\noindent \textbf{(P\_3,2.57)} vyākhyānataḥ viśeṣapratipattiḥ . \\
\noindent \textbf{(P\_3,2.57)} na hi sandehāt alakṣaṇam iti . \\
\noindent \textbf{(P\_3,2.57)} khit iti vyākhyāsyāmaḥ . \\
\noindent \textbf{(P\_3,2.57)} nañaḥ tu svarasiddhyartham ikārāditvam iṣṇucaḥ . \\
\noindent \textbf{(P\_3,2.57)} idam tarhi prayojanam . \\
\noindent \textbf{(P\_3,2.57)} kṛtyokeṣṇuccārvādayaḥ ca iti eṣaḥ svaraḥ yathā syāt . \\
\noindent \textbf{(P\_3,2.57)} etat api na asti prayojanam . \\
\noindent \textbf{(P\_3,2.57)} ayam api iṭi kṛte ṣatve ca iṣṇuc bhaviṣyati . \\
\noindent \textbf{(P\_3,2.57)} na sidhyati . \\
\noindent \textbf{(P\_3,2.57)} lakṣaṇapratipadoktayoḥ pratipadoktasya eva iti . \\
\noindent \textbf{(P\_3,2.57)} atha vā asiddham khalu api ṣatvam . \\
\noindent \textbf{(P\_3,2.57)} ṣatvasya asiddhatvāt isnuc eva bhavati . \\
\noindent \textbf{(P\_3,2.57)} iṣṇucaḥ ikārāditvam udāttatvāt kṛtam bhuvaḥ . \\
\noindent \textbf{(P\_3,2.57)} nañaḥ tu svarasiddhyartham ikārāditvam iṣṇucaḥ . \\
\noindent \textbf{(P\_3,2.58)} kimarthaḥ nakāraḥ . \\
\noindent \textbf{(P\_3,2.58)} ñniti iti ādyudāttatvam yathā syāt . \\
\noindent \textbf{(P\_3,2.58)} na etat asti prayojanam . \\
\noindent \textbf{(P\_3,2.58)} ekācaḥ ayam vidhīyate . \\
\noindent \textbf{(P\_3,2.58)} tatra na arthaḥ svarārthena nakāreṇa anubandhena . \\
\noindent \textbf{(P\_3,2.58)} dhātusvareṇa eva siddham . \\
\noindent \textbf{(P\_3,2.58)} yaḥ tarhi anekāc . \\
\noindent \textbf{(P\_3,2.58)} dadhṛk iti . \\
\noindent \textbf{(P\_3,2.58)} vakṣyati etat dhṛṣeḥ dvirvacanam antodāttatvam ca nipātyate iti . \\
\noindent \textbf{(P\_3,2.58)} viśeṣaṇārthaḥ tarhi . \\
\noindent \textbf{(P\_3,2.58)} kva viśeṣanārthena arthaḥ . \\
\noindent \textbf{(P\_3,2.58)} kvinpratyayasya kuḥ iti . \\
\noindent \textbf{(P\_3,2.58)} kvipratyayasya kuḥ iti ucyamāne sandehaḥ syāt . \\
\noindent \textbf{(P\_3,2.58)} kviḥ vā eṣaḥ pratyayaḥ kvip vā iti . \\
\noindent \textbf{(P\_3,2.58)} sandehamātram etat bhavati . \\
\noindent \textbf{(P\_3,2.58)} sarvasandeheṣu ca idam upatiṣṭhate . \\
\noindent \textbf{(P\_3,2.58)} vyākhyānataḥ viśeṣapratipattiḥ . \\
\noindent \textbf{(P\_3,2.58)} na hi sandehāt alakṣaṇam iti . \\
\noindent \textbf{(P\_3,2.58)} kkvipratyayasya iti vyākhyāsyāmaḥ . \\
\noindent \textbf{(P\_3,2.59)} dadhṛk iti kim nipātyate . \\
\noindent \textbf{(P\_3,2.59)} dhṛṣeḥ dvirvacanam antodāttatvam ca . \\
\noindent \textbf{(P\_3,2.59)} dhṛṣeḥ dvirvacanam antodāttatvam ca nipātyate . \\
\noindent \textbf{(P\_3,2.60.1)} kimarthaḥ ñakāraḥ . \\
\noindent \textbf{(P\_3,2.60.1)} svarārthaḥ . \\
\noindent \textbf{(P\_3,2.60.1)} ñniti iti ādyudāttatvam yathā syāt . \\
\noindent \textbf{(P\_3,2.60.1)} na etat asti prayojanam . \\
\noindent \textbf{(P\_3,2.60.1)} nakāreṇa api eṣaḥ svaraḥ siddhaḥ . \\
\noindent \textbf{(P\_3,2.60.1)} viśeṣaṇārthaḥ tarhi bhaviṣyati . \\
\noindent \textbf{(P\_3,2.60.1)} kva viśeṣaṇārthena arthaḥ . \\
\noindent \textbf{(P\_3,2.60.1)} kañkvarap iti . \\
\noindent \textbf{(P\_3,2.60.1)} kankvarap iti ucyāmāne yācitikā atra api prasajyeta .. \\
\noindent \textbf{(P\_3,2.60.2)} dṛśeḥ samānānyayoḥ ca upasaṅkhyānam . \\
\noindent \textbf{(P\_3,2.60.2)} dṛśeḥ samānānyayoḥ ca upasaṅkhyānam kartavyam . \\
\noindent \textbf{(P\_3,2.60.2)} sadṛk sadṛśaḥ anyādṛk anyādṛśaḥ . \\
\noindent \textbf{(P\_3,2.60.2)} kṛdarthānupapattiḥ tu . \\
\noindent \textbf{(P\_3,2.60.2)} kṛdarthaḥ tu na upapadyate . \\
\noindent \textbf{(P\_3,2.60.2)} dṛśeḥ kartari prāpnoti . \\
\noindent \textbf{(P\_3,2.60.2)} ivārthe tu taddhitaḥ . \\
\noindent \textbf{(P\_3,2.60.2)} ivārthe ayam taddhitaḥ draṣṭavyaḥ . \\
\noindent \textbf{(P\_3,2.60.2)} saḥ iva ayam tādṛk . \\
\noindent \textbf{(P\_3,2.60.2)} anya iva ayam anyādṛk . \\
\noindent \textbf{(P\_3,2.60.2)} atha vā yuktaḥ eva atra kṛdarthaḥ . \\
\noindent \textbf{(P\_3,2.60.2)} karmakartā ayam . \\
\noindent \textbf{(P\_3,2.60.2)} tam iva imam paśyanti janāḥ . \\
\noindent \textbf{(P\_3,2.60.2)} saḥ ayam saḥ iva dṛśyamānaḥ tam iva ātmānam paśyati . \\
\noindent \textbf{(P\_3,2.60.2)} tādṛk . \\
\noindent \textbf{(P\_3,2.60.2)} anyam iva imam paśyanti janāḥ . \\
\noindent \textbf{(P\_3,2.60.2)} saḥ ayam anyaḥ iva dṛśyamānaḥ anyam iva ātmānam paśyati . \\
\noindent \textbf{(P\_3,2.60.2)} anyādṛk iti . \\
\noindent \textbf{(P\_3,2.61)} sadādiṣu subgrahaṇam . \\
\noindent \textbf{(P\_3,2.61)} sadādiṣu subgrahaṇam kartavyam . \\
\noindent \textbf{(P\_3,2.61)} hotā vediṣat . \\
\noindent \textbf{(P\_3,2.61)} atithiḥ curoṇasat . \\
\noindent \textbf{(P\_3,2.61)} na tarhi idānīm upasarge api iti vaktavyam . \\
\noindent \textbf{(P\_3,2.61)} vaktavyam ca . \\
\noindent \textbf{(P\_3,2.61)} kim prayojanam . \\
\noindent \textbf{(P\_3,2.61)} jñāpakārtham . \\
\noindent \textbf{(P\_3,2.61)} kim jñāpyam . \\
\noindent \textbf{(P\_3,2.61)} etat jñāpayati ācāryaḥ . \\
\noindent \textbf{(P\_3,2.61)} anyatra subgrahaṇe upasargagrahaṇam na bhavati iti . \\
\noindent \textbf{(P\_3,2.61)} kim etasya jñāpane prayojanam . \\
\noindent \textbf{(P\_3,2.61)} vadaḥ supi anupasargagrahaṇam coditam . \\
\noindent \textbf{(P\_3,2.61)} tat na vaktavyam bhavati . \\
\noindent \textbf{(P\_3,2.68-69)} KA\_II,108.2-6 Ro\_III,244 {1/9}     kimartham idam ucyate na adaḥ ananne iti eva siddham . \\
\noindent \textbf{(P\_3,2.68-69)} KA\_II,108.2-6 Ro\_III,244 {2/9}     na sidhyati . \\
\noindent \textbf{(P\_3,2.68-69)} KA\_II,108.2-6 Ro\_III,244 {3/9}     chandasi iti etat anuvartate . \\
\noindent \textbf{(P\_3,2.68-69)} KA\_II,108.2-6 Ro\_III,244 {4/9}     bhāṣārthaḥ ayam ārambhaḥ . \\
\noindent \textbf{(P\_3,2.68-69)} KA\_II,108.2-6 Ro\_III,244 {5/9}     pūrvasmin eva yoge chandograhaṇam nivṛttam . \\
\noindent \textbf{(P\_3,2.68-69)} KA\_II,108.2-6 Ro\_III,244 {6/9}     tat ca avaśyam nivartyam amāt iti evamartham . \\
\noindent \textbf{(P\_3,2.68-69)} KA\_II,108.2-6 Ro\_III,244 {7/9}     ataḥ uttaram paṭhati . \\
\noindent \textbf{(P\_3,2.68-69)} KA\_II,108.2-6 Ro\_III,244 {8/9}     adaḥ ananne kravyegrahaṇam vāsarūpanivṛttyartham . \\
\noindent \textbf{(P\_3,2.68-69)} KA\_II,108.2-6 Ro\_III,244 {9/9}     adaḥ ananne kravyegrahaṇam kriyate vāsarūpaḥ mā bhūt . \\
\noindent \textbf{(P\_3,2.71)} śvetavahādīnāṃ ḍas . \\
\noindent \textbf{(P\_3,2.71)} śvetavahādīnāṃ ḍas vaktavyaḥ . \\
\noindent \textbf{(P\_3,2.71)} śvetavāḥ indraḥ . \\
\noindent \textbf{(P\_3,2.71)} padasya ca . \\
\noindent \textbf{(P\_3,2.71)} padasya ca iti vaktavyam . \\
\noindent \textbf{(P\_3,2.71)} iha mā bhūt . \\
\noindent \textbf{(P\_3,2.71)} śvetavāhau śvetavāhaḥ . \\
\noindent \textbf{(P\_3,2.71)} kim prayojanam . \\
\noindent \textbf{(P\_3,2.71)} rvartham . \\
\noindent \textbf{(P\_3,2.71)} ruḥ yathā syāt . \\
\noindent \textbf{(P\_3,2.71)} kriyate rvartham nipātanam . \\
\noindent \textbf{(P\_3,2.71)} avayāḥ śvetavāḥ puroḍāḥ ca iti . \\
\noindent \textbf{(P\_3,2.71)} ātaḥ ca rvartham . \\
\noindent \textbf{(P\_3,2.71)} ukthaśasśabdasya sāmānyena ruḥ siddhaḥ . \\
\noindent \textbf{(P\_3,2.71)} na tasya nipātanam kriyate . \\
\noindent \textbf{(P\_3,2.71)} tat na vaktavyam . \\
\noindent \textbf{(P\_3,2.71)} avaśyam tat vaktavyam dīrghārtham . \\
\noindent \textbf{(P\_3,2.71)} na etat asti prayojanam . \\
\noindent \textbf{(P\_3,2.71)} siddham atra dīrghatvam atvasantasya ca adhātoḥ iti . \\
\noindent \textbf{(P\_3,2.71)} yatra tena na sidhyati tadartham . \\
\noindent \textbf{(P\_3,2.71)} kva ca tena na sidhyati . \\
\noindent \textbf{(P\_3,2.71)} sambuddhau . \\
\noindent \textbf{(P\_3,2.71)} he śvetavāḥ iti . \\
\noindent \textbf{(P\_3,2.71)} na tarhi idānīm ḍas vaktavyaḥ . \\
\noindent \textbf{(P\_3,2.71)} vaktavyaḥ ca . \\
\noindent \textbf{(P\_3,2.71)} kim prayojanam . \\
\noindent \textbf{(P\_3,2.71)} uttarāṛtham . \\
\noindent \textbf{(P\_3,2.71)} śvetavobhyām śvetavobhiḥ . \\
\noindent \textbf{(P\_3,2.77)} kimartham sthaḥ kakvipau ucyete na kvip siddhaḥ anyebhyaḥ api dṛśyate iti kaḥ ca ātaḥ anupasarge kaḥ iti . \\
\noindent \textbf{(P\_3,2.77)} na sidhyati . \\
\noindent \textbf{(P\_3,2.77)} viśeṣvihitaḥ kaḥ sāmānyavihitam kvipam bādhate. vāsarūpeṇa kvip api bhaviṣyati . \\
\noindent \textbf{(P\_3,2.77)} idam tarhi śaṃsthaḥ śaṃsthāḥ . \\
\noindent \textbf{(P\_3,2.77)} uktam etat . \\
\noindent \textbf{(P\_3,2.77)} śami sañjñāyām dhātugrahaṇam kṛñaḥ hetvādiṣu ṭapratiṣedhārtham iti . \\
\noindent \textbf{(P\_3,2.77)} saḥ yathā eva ac ṭam bādhate evam kakvipau api bādheta . \\
\noindent \textbf{(P\_3,2.78)} supi iti vartamāne punaḥ subgrahaṇam kimartham . \\
\noindent \textbf{(P\_3,2.78)} anupasarge iti evam tat abhūt . \\
\noindent \textbf{(P\_3,2.78)} idam submātre yathā syāt . \\
\noindent \textbf{(P\_3,2.78)} pratyāsāriṇyaḥ udāsāriṇyaḥ . \\
\noindent \textbf{(P\_3,2.78)} ṇinvidhau sādhukāriṇi upasaṅkhyānam . \\
\noindent \textbf{(P\_3,2.78)} ṇinvidhau sādhukāriṇi upasaṅkhyānam kartavyam . \\
\noindent \textbf{(P\_3,2.78)} sādhukārī sādhudāyī . \\
\noindent \textbf{(P\_3,2.78)} brahmaṇi vadaḥ . \\
\noindent \textbf{(P\_3,2.78)} brahmaṇi vadaḥ upasaṅkhyānam kartavyam . \\
\noindent \textbf{(P\_3,2.78)} brahmavādinaḥ vadanti . \\
\noindent \textbf{(P\_3,2.80)} kim udāharaṇam . \\
\noindent \textbf{(P\_3,2.80)} aśrāddhabhojī . \\
\noindent \textbf{(P\_3,2.80)} kim yaḥ aśrāddham bhuṅkte saḥ aśrāddhabhojī . \\
\noindent \textbf{(P\_3,2.80)} kim ca ataḥ . \\
\noindent \textbf{(P\_3,2.80)} yadā asau aśrāddham na bhuṅkte tadā asya vratalopaḥ syāt . \\
\noindent \textbf{(P\_3,2.80)} tat yathā : sthāyī yadā na tiṣthati tadā asya vratalopaḥ bhavati . \\
\noindent \textbf{(P\_3,2.80)} evam tarhi ṇinyantena samāsaḥ bhaviṣyati : na śrāddhabhojī aśrāddhabhojī . \\
\noindent \textbf{(P\_3,2.80)} na evam śakyam . \\
\noindent \textbf{(P\_3,2.80)} svare hi doṣaḥ syāt . \\
\noindent \textbf{(P\_3,2.80)} aśrāddhabhojī iti evam svaraḥ prasajyeta . \\
\noindent \textbf{(P\_3,2.80)} aśrāddhabhojī iti ca iṣyate . \\
\noindent \textbf{(P\_3,2.80)} evam tarhi nañaḥ eva ayam bhujipratiṣedhavācinaḥ śrāddhaśabdena asamarthasamāsaḥ : na bhojī śrāddhasya iti . \\
\noindent \textbf{(P\_3,2.80)} saḥ tarhi asamarthasamāsaḥ vaktavyaḥ . \\
\noindent \textbf{(P\_3,2.80)} yadi api vaktavyaḥ atha vā etarhi bahūni prayojanāni . \\
\noindent \textbf{(P\_3,2.80)} kāni . \\
\noindent \textbf{(P\_3,2.80)} asūryampaśyāni mukhāni . \\
\noindent \textbf{(P\_3,2.80)} apūrvageyāḥ ślokāḥ . \\
\noindent \textbf{(P\_3,2.80)} aśrāddhabhojī brāhmaṇaḥ . \\
\noindent \textbf{(P\_3,2.80)} suṭ anapuṃsakasya iti . \\
\noindent \textbf{(P\_3,2.83)} ātmagrahaṇam kimartham . \\
\noindent \textbf{(P\_3,2.83)} paramāne mā bhūt . \\
\noindent \textbf{(P\_3,2.83)} kriyamāṇe api ātmagrahaṇe paramāne prāpnoti . \\
\noindent \textbf{(P\_3,2.83)} kim kāraṇam . \\
\noindent \textbf{(P\_3,2.83)} ātmanaḥ iti iyam kartari ṣaṣṭhī mānaḥ iti akāraḥ bhāve . \\
\noindent \textbf{(P\_3,2.83)} saḥ yadi eva ātmānam manyate atha api param ātmanaḥ eva asau mānaḥ bhavati . \\
\noindent \textbf{(P\_3,2.83)} na eṣaḥ doṣaḥ . \\
\noindent \textbf{(P\_3,2.83)} ātmanaḥ iti karmaṇi ṣaṣṭhī . \\
\noindent \textbf{(P\_3,2.83)} katham . \\
\noindent \textbf{(P\_3,2.83)} kartṛkarmaṇoḥ kṛti iti . \\
\noindent \textbf{(P\_3,2.83)} nanu ca kartari api vai etena eva vidhīyate . \\
\noindent \textbf{(P\_3,2.83)} tatra kutaḥ etat karmaṇi bhaviṣyati na punaḥ kartari iti . \\
\noindent \textbf{(P\_3,2.83)} evam tarhi karmakartari ca . \\
\noindent \textbf{(P\_3,2.83)} karmakartari ca iti vaktavyam . \\
\noindent \textbf{(P\_3,2.83)} tat tarhi vaktavyam . \\
\noindent \textbf{(P\_3,2.83)} na vaktavyam . \\
\noindent \textbf{(P\_3,2.83)} ātmanaḥ iti karmaṇi ṣaṣṭhī . \\
\noindent \textbf{(P\_3,2.83)} katham . \\
\noindent \textbf{(P\_3,2.83)} kartṛkarmaṇoḥ kṛti iti . \\
\noindent \textbf{(P\_3,2.83)} nanu ca uktam kartari api vai etena eva vidhīyate . \\
\noindent \textbf{(P\_3,2.83)} tatra kutaḥ etat karmaṇi bhaviṣyati na punaḥ kartari iti . \\
\noindent \textbf{(P\_3,2.83)} ātmagrahaṇasāmarthyāt karmaṇi vijñāsyate . \\
\noindent \textbf{(P\_3,2.83)} evam api karmakartṛgrahaṇam kartavyam karmāpadiṣṭaḥ yak yathā syāt śyan mā bhūt iti . \\
\noindent \textbf{(P\_3,2.83)} kaḥ ca atra viśeṣaḥ yakaḥ vā śyanaḥ vā . \\
\noindent \textbf{(P\_3,2.83)} yaki sati antodāttatvena bhavitayam śyani sati ādyudāttatvena . \\
\noindent \textbf{(P\_3,2.83)} śyani api sati antodāttatvena eva bhavitavyam . \\
\noindent \textbf{(P\_3,2.83)} katham . \\
\noindent \textbf{(P\_3,2.83)} khaśaḥ svaraḥ śyanaḥ svaram bādhiṣyate . \\
\noindent \textbf{(P\_3,2.83)} sati śiṣṭatvāt śyanaḥ svaraḥ prāpnoti . \\
\noindent \textbf{(P\_3,2.83)} ācāryapravṛttiḥ jñāpayati sati śiṣṭaḥ api vikaraṇasvaraḥ sārvadhātukasvaram na bādhate iti yat ayam tāseḥ parasya lasārvadhātukasya anudāttatvam śāsti . \\
\noindent \textbf{(P\_3,2.83)} lasārvadhātuke etat jñāpakam syāt . \\
\noindent \textbf{(P\_3,2.83)} na iti āha . \\
\noindent \textbf{(P\_3,2.83)} aviśeṣeṇa jñāpakam . \\
\noindent \textbf{(P\_3,2.84)} bhūte iti ucyate . \\
\noindent \textbf{(P\_3,2.84)} kasmin bhūte . \\
\noindent \textbf{(P\_3,2.84)} kāle . \\
\noindent \textbf{(P\_3,2.84)} na vai kālādhikāraḥ asti . \\
\noindent \textbf{(P\_3,2.84)} evam tarhi dhātoḥ iti vartate . \\
\noindent \textbf{(P\_3,2.84)} dhātau bhūte . \\
\noindent \textbf{(P\_3,2.84)} dhātuḥ vai śabdaḥ . \\
\noindent \textbf{(P\_3,2.84)} na ca śabdasya bhūtabhaviṣyadvartamānatāyām sambhavaḥ asti . \\
\noindent \textbf{(P\_3,2.84)} śabde asambhavāt arthe kāryam vijñāsyate . \\
\noindent \textbf{(P\_3,2.84)} kaḥ punaḥ dhātvarthaḥ . \\
\noindent \textbf{(P\_3,2.84)} kriyā . \\
\noindent \textbf{(P\_3,2.84)} kriyāyām bhūtāyām . \\
\noindent \textbf{(P\_3,2.84)} yadi evam niṣṭhāyām itaretarāśrayatvāt aprasiddhiḥ . \\
\noindent \textbf{(P\_3,2.84)} niṣṭhāyām itaretarāśrayatvāt aprasiddhiḥ . \\
\noindent \textbf{(P\_3,2.84)} kā itaretarāśrayatā . \\
\noindent \textbf{(P\_3,2.84)} bhūtakālena śabdena nirdeśaḥ kriyate . \\
\noindent \textbf{(P\_3,2.84)} nirdeśottarakālam ca bhūtakālatā . \\
\noindent \textbf{(P\_3,2.84)} tat etat itaretarāśrayam bhavati . \\
\noindent \textbf{(P\_3,2.84)} itaretarāśrayāṇi ca na prakalpante . \\
\noindent \textbf{(P\_3,2.84)} avyayanirdeśāt siddham . \\
\noindent \textbf{(P\_3,2.84)} avyayavatā śabdena nirdeśaḥ kariṣyate . \\
\noindent \textbf{(P\_3,2.84)} avartamāne abhaviṣyati iti . \\
\noindent \textbf{(P\_3,2.84)} saḥ tarhi avyayavatā śabdena nirdeśaḥ kartavyaḥ . \\
\noindent \textbf{(P\_3,2.84)} na kartavyaḥ . \\
\noindent \textbf{(P\_3,2.84)} avyayam eṣaḥ bhūteśabdaḥ na bhavateḥ niṣṭhā . \\
\noindent \textbf{(P\_3,2.84)} katham avyayatvam . \\
\noindent \textbf{(P\_3,2.84)} vibhaktisvarapratirūpakāḥ ca nipātāḥ bhavanti iti nipātasañjñā . \\
\noindent \textbf{(P\_3,2.84)} nipātam avyayam iti avayayasañjñā . \\
\noindent \textbf{(P\_3,2.84)} atha api bhavateḥ niṣṭhā evam api avayayam eva . \\
\noindent \textbf{(P\_3,2.84)} katham na vyeti iti avyayam . \\
\noindent \textbf{(P\_3,2.84)} kva punaḥ na vyeti . \\
\noindent \textbf{(P\_3,2.84)} etau kālaviśeṣau vartamānabhaviṣyantau . \\
\noindent \textbf{(P\_3,2.84)} svabhāvataḥ bhūte eva vartate . \\
\noindent \textbf{(P\_3,2.84)} yadi tatri na vyeti iti avyayam . \\
\noindent \textbf{(P\_3,2.84)} na vā tadvidhānasya anyatra abhāvāt . \\
\noindent \textbf{(P\_3,2.84)} na vā bhūtādhikāreṇa arthaḥ . \\
\noindent \textbf{(P\_3,2.84)} kim kāraṇam . \\
\noindent \textbf{(P\_3,2.84)} tadvidhānasya anyatra abhāvāt . \\
\noindent \textbf{(P\_3,2.84)} ye api ete itaḥ uttaram pratyayāḥ śiṣyante ete api etau kālaviśeṣau na viyanti vartamānabhaviṣyantau . \\
\noindent \textbf{(P\_3,2.84)} svabhāvataḥ eva te bhūte eva vartante . \\
\noindent \textbf{(P\_3,2.84)} ataḥ uttaram paṭhati . \\
\noindent \textbf{(P\_3,2.84)} bhūtādhikārasya prayojanam kumāraghātī śīrṣaghātī ākhuhā biḍālaḥ sutvānaḥ sunavantaḥ suṣupuṣaḥ anehāḥ agnim ādadhānasya . \\
\noindent \textbf{(P\_3,2.84)} kumāraghātī śīrṣaghātī iti bhaviṣyadvartamānārthaḥ bhūtanivṛttyarthaḥ . \\
\noindent \textbf{(P\_3,2.84)} ākhuhā biḍālaḥ iti bhaviṣyadvartamānārthaḥ . \\
\noindent \textbf{(P\_3,2.84)} itarathā hi brahmādiṣu niyamaḥ triṣu kāleṣu nivartakaḥ syāt . \\
\noindent \textbf{(P\_3,2.84)} sutvānaḥ sunvantaḥ . \\
\noindent \textbf{(P\_3,2.84)} yajñasaṃyoge ṅvanipaḥ triṣu kāleṣu śatā apavādaḥ mā bhūt . \\
\noindent \textbf{(P\_3,2.84)} suṣupuṣaḥ . \\
\noindent \textbf{(P\_3,2.84)} najiṅ sarvakālapavādaḥ mā bhūt . \\
\noindent \textbf{(P\_3,2.84)} anehāḥ iti vartamānakālaḥ eva . \\
\noindent \textbf{(P\_3,2.84)} anyatra anāhantā . \\
\noindent \textbf{(P\_3,2.84)} ādadhānasya . \\
\noindent \textbf{(P\_3,2.84)} kānacaḥ cānaś tācchīlādiṣu sarvakālāpavādaḥ mā bhūt . \\
\noindent \textbf{(P\_3,2.84)} agnim ādadhānasya . \\
\noindent \textbf{(P\_3,2.84)} ādadhānasya iti eva anyatra . \\
\noindent \textbf{(P\_3,2.87)} kimartham brahmādiṣu hanteḥ kvip vidhīyate . \\
\noindent \textbf{(P\_3,2.87)} na kvip ca anyebhyaḥ api dṛśyate iti eva siddham . \\
\noindent \textbf{(P\_3,2.87)} brahmādiṣu hanteḥ kvibvacanam niyamārtham . niyamārthaḥ ayam ārambhaḥ . \\
\noindent \textbf{(P\_3,2.87)} brahmādiṣu eva hanteḥ kvip yathā syāt . \\
\noindent \textbf{(P\_3,2.87)} kim aviśeṣeṇa . \\
\noindent \textbf{(P\_3,2.87)} na iti āha .upapadaviśeṣe etasmin ca viśeṣe . \\
\noindent \textbf{(P\_3,2.87)} atha brahmādiṣu hanteḥ ṇininā bhavitavyam . \\
\noindent \textbf{(P\_3,2.87)} na bhavitavyam . \\
\noindent \textbf{(P\_3,2.87)} kim kāraṇam . \\
\noindent \textbf{(P\_3,2.87)} ubhayataḥ niyamāt . \\
\noindent \textbf{(P\_3,2.87)} ubhayataḥ niyamaḥ ayam . \\
\noindent \textbf{(P\_3,2.87)} brahmādiṣu eva hanteḥ kvip bhavati . \\
\noindent \textbf{(P\_3,2.87)} kvip eva ca brahmādiṣu iti . \\
\noindent \textbf{(P\_3,2.93)} karmaṇi kutsite . \\
\noindent \textbf{(P\_3,2.93)} karmaṇi kutsite iti vaktavyam . \\
\noindent \textbf{(P\_3,2.93)} iha mā bhūt . \\
\noindent \textbf{(P\_3,2.93)} dhānyavikrāyaḥ . \\
\noindent \textbf{(P\_3,2.101)} anyebhyaḥ api dṛśyate iti vaktavyam , iha api yathā syāt . \\
\noindent \textbf{(P\_3,2.101)} ākhā utkhā parikhā . \\
\noindent \textbf{(P\_3,2.102.1)} niṣṭhāyām itaretarāśrayatvāt aprasiddhiḥ . \\
\noindent \textbf{(P\_3,2.102.1)} niṣṭhāyām itaretarāśrayatvāt aprasiddhiḥ . \\
\noindent \textbf{(P\_3,2.102.1)} kā itaretarāśrayatā . \\
\noindent \textbf{(P\_3,2.102.1)} satoḥ ktaktavatvoḥ sañjñayā bhavitavyam sañjñayā ca ktaktavatū bhāvyete . \\
\noindent \textbf{(P\_3,2.102.1)} tat etat itaretarāśrayam bhavati . \\
\noindent \textbf{(P\_3,2.102.1)} itaretarāśrayāṇi ca na prakalpante . \\
\noindent \textbf{(P\_3,2.102.1)} dviḥ vā ktaktavtugrahaṇam . \\
\noindent \textbf{(P\_3,2.102.1)} dviḥ vā ktaktavtugrahaṇam kartavyam . \\
\noindent \textbf{(P\_3,2.102.1)} ktaktavtū bhūte . \\
\noindent \textbf{(P\_3,2.102.1)} ktaktavatū niṣṭhā iti . \\
\noindent \textbf{(P\_3,2.102.1)} yadi punaḥ iha eva niṣṭhāsañjñā api ucyeta : ktaktavtū bhūte . \\
\noindent \textbf{(P\_3,2.102.1)} tataḥ niṣṭhā . \\
\noindent \textbf{(P\_3,2.102.1)} niṣṭhāsañjñau ca ktaktavtū bhavataḥ iti . \\
\noindent \textbf{(P\_3,2.102.1)} kim kṛtam bhavati . \\
\noindent \textbf{(P\_3,2.102.1)} dviḥ vā ktaktavtugrahaṇam na kartavyam bhavati . \\
\noindent \textbf{(P\_3,2.102.1)} evam api tau iti vaktavyam syāt . \\
\noindent \textbf{(P\_3,2.102.1)} vakṣyati hi etat . \\
\noindent \textbf{(P\_3,2.102.1)} tau sat iti vacanam asaṃsargārtham iti . \\
\noindent \textbf{(P\_3,2.102.1)} asaṃsaktayoḥ bhūtena kālena niṣṭhāsañjñā yathā syāt . \\
\noindent \textbf{(P\_3,2.102.1)} ñimidā minnaḥ ñikṣvidā kṣvinnaḥ . \\
\noindent \textbf{(P\_3,2.102.1)} yadi punaḥ adṛṣṭaśrutau eva ktaktavatū gṛhītvā niṣṭhāsañjñā ucyeta . \\
\noindent \textbf{(P\_3,2.102.1)} na evam śakyam . \\
\noindent \textbf{(P\_3,2.102.1)} dṛṣṭaśrutayoḥ na syāt . \\
\noindent \textbf{(P\_3,2.102.1)} ñimidā minnaḥ . \\
\noindent \textbf{(P\_3,2.102.1)} tasmāt na evam śakyam . \\
\noindent \textbf{(P\_3,2.102.1)} na cet evam dviḥ vā ktaktavtugrahaṇam kartavyam itaretarāśrayam vā bhavati . \\
\noindent \textbf{(P\_3,2.102.1)} na eṣaḥ doṣaḥ . \\
\noindent \textbf{(P\_3,2.102.1)} itaretarāśrayamātram etat bhavati . \\
\noindent \textbf{(P\_3,2.102.1)} sarvāṇi ca itaretarāśrayāṇi ekatvena parihṛtāni siddham tu nityaśabdatvāt iti . \\
\noindent \textbf{(P\_3,2.102.1)} na idam tulyam anyaiḥ itaretarāśrayaiḥ . \\
\noindent \textbf{(P\_3,2.102.1)} na hi sañjñā nityā . \\
\noindent \textbf{(P\_3,2.102.1)} evam tarhi bhāvinī sañjñā vijñāsyate . \\
\noindent \textbf{(P\_3,2.102.1)} tat yathā : kaḥ cit kam cit tantuvāyam āha : asya sūtrasya śāṭakam vaya iti . \\
\noindent \textbf{(P\_3,2.102.1)} saḥ paśyati . \\
\noindent \textbf{(P\_3,2.102.1)} yadi śāṭakaḥ na vātavyaḥ atha vātavyaḥ na śāṭakaḥ . \\
\noindent \textbf{(P\_3,2.102.1)} śāṭakaḥ vātavyaḥ iti vipratiṣiddham. bhāvinī khalu asya sañjñā abhipretā . \\
\noindent \textbf{(P\_3,2.102.1)} saḥ manye vātavyaḥ yasmin ute śāṭakaḥ iti etat bhavati iti . \\
\noindent \textbf{(P\_3,2.102.1)} evam iha api tau bhūte kāle bhavataḥ yayoḥ abhinirvṛtayoḥ niṣṭhā iti eṣā sañjñā bhaviṣyati . \\
\noindent \textbf{(P\_3,2.102.2)} ādikarmaṇi niṣṭhā . \\
\noindent \textbf{(P\_3,2.102.2)} ādikarmaṇi niṣṭhā vaktavyā . \\
\noindent \textbf{(P\_3,2.102.2)} prakṛtaḥ kaṭam devadattaḥ . \\
\noindent \textbf{(P\_3,2.102.2)} kim punaḥ kāraṇam na sidhyati . \\
\noindent \textbf{(P\_3,2.102.2)} yat vā bhavantyarthe . \\
\noindent \textbf{(P\_3,2.102.2)} yat vā bhavantyarthe bhāṣyate . \\
\noindent \textbf{(P\_3,2.102.2)} prakṛtaḥ kaṭam devadattaḥ . \\
\noindent \textbf{(P\_3,2.102.2)} prakaroti kaṭam devadattaḥ iti . \\
\noindent \textbf{(P\_3,2.102.2)} nyāyyā tu ādyapavargāt . \\
\noindent \textbf{(P\_3,2.102.2)} nyāyyā tu eṣā bhūtakālatā . \\
\noindent \textbf{(P\_3,2.102.2)} kutaḥ . \\
\noindent \textbf{(P\_3,2.102.2)} ādyapavargāt . \\
\noindent \textbf{(P\_3,2.102.2)} ādiḥ atra apavṛktaḥ . \\
\noindent \textbf{(P\_3,2.102.2)} eṣaḥ ca nāma nyāyyaḥ bhūtakālaḥ yatra kim cit apavṛktam dṛśyate . \\
\noindent \textbf{(P\_3,2.102.2)} vā ca adyatanyām . \\
\noindent \textbf{(P\_3,2.102.2)} vā ca adyatanyām bhāṣyate . \\
\noindent \textbf{(P\_3,2.102.2)} prakṛtaḥ kaṭam devadattaḥ . \\
\noindent \textbf{(P\_3,2.102.2)} prākārṣīt kaṭam devadattaḥ iti . \\
\noindent \textbf{(P\_3,2.102.2)} kim śakyante ete śabdāḥ prayoktum iti ataḥ nyāyyā eṣā bhūtakālatā . \\
\noindent \textbf{(P\_3,2.102.2)} na avaśyam prayogāt eva . \\
\noindent \textbf{(P\_3,2.102.2)} kriyā nāma iyam atyantāparidṛṣṭā anumānagamyā aśakyā piṇḍībhūtā nidarśayitum yathā garbhaḥ nirluṭhitaḥ . \\
\noindent \textbf{(P\_3,2.102.2)} sā asau yena yena śabdena abhisambadhyate tāvati tāvati parisampāpyate . \\
\noindent \textbf{(P\_3,2.102.2)} tat yathā . \\
\noindent \textbf{(P\_3,2.102.2)} kaḥ cit pāṭaliputram jigamiṣuḥ ekam ahaḥ gatvā āha idam adya gatam iti . \\
\noindent \textbf{(P\_3,2.102.2)} na ca tāvatā asya vrajikriyā parisamāptā bhavati . \\
\noindent \textbf{(P\_3,2.102.2)} yat tu gatam tat abhisamīkṣya etat prayujyate idam adya gatam iti . \\
\noindent \textbf{(P\_3,2.102.2)} evam iha api yat kṛtam tat abhisamīkṣya etat prayujyate prakṛtaḥ kaṭam devadattaḥ iti . \\
\noindent \textbf{(P\_3,2.102.2)} yadā hi veṇikāntaḥ kaṭaḥ abhisamīkṣitaḥ bhavati prakaroti kaṭam iti eva tadā bhavati . \\
\noindent \textbf{(P\_3,2.106-107.1)} KA\_II,114.19-115.2 Ro\_III,260 {1/14}     kimartham kānackvasoḥ vāvacanam kriyate . \\
\noindent \textbf{(P\_3,2.106-107.1)} KA\_II,114.19-115.2 Ro\_III,260 {2/14}     kānackvasoḥ vāvacanam chandasi tiṅaḥ darśanāt . \\
\noindent \textbf{(P\_3,2.106-107.1)} KA\_II,114.19-115.2 Ro\_III,260 {3/14}     kānackvasoḥ vāvacanam kriyate chandasi tiṅaḥ darśanāt . \\
\noindent \textbf{(P\_3,2.106-107.1)} KA\_II,114.19-115.2 Ro\_III,260 {4/14}     chandasi tiṅ api dṛśyate . \\
\noindent \textbf{(P\_3,2.106-107.1)} KA\_II,114.19-115.2 Ro\_III,260 {5/14}     aham sūram ubhayataḥ dadarśa . \\
\noindent \textbf{(P\_3,2.106-107.1)} KA\_II,114.19-115.2 Ro\_III,260 {6/14}     aham dyāvāpṛthivī ātatāna . \\
\noindent \textbf{(P\_3,2.106-107.1)} KA\_II,114.19-115.2 Ro\_III,260 {7/14}     na vā anena vihitasya ādeśavacanāt . \\
\noindent \textbf{(P\_3,2.106-107.1)} KA\_II,114.19-115.2 Ro\_III,260 {8/14}     na vā etat prayojanam asti . \\
\noindent \textbf{(P\_3,2.106-107.1)} KA\_II,114.19-115.2 Ro\_III,260 {9/14}     kim kāraṇam . \\
\noindent \textbf{(P\_3,2.106-107.1)} KA\_II,114.19-115.2 Ro\_III,260 {10/14}     anena vihitasya ādeśavacanāt . \\
\noindent \textbf{(P\_3,2.106-107.1)} KA\_II,114.19-115.2 Ro\_III,260 {11/14}     astu anena vihitasya ādeśaḥ . \\
\noindent \textbf{(P\_3,2.106-107.1)} KA\_II,114.19-115.2 Ro\_III,260 {12/14}     kena idānīm chandasi vihitasya liṭaḥ śravaṇam bhaviṣyati . \\
\noindent \textbf{(P\_3,2.106-107.1)} KA\_II,114.19-115.2 Ro\_III,260 {13/14}     chandasi luṅlaṅliṭaḥ iti anena . \\
\noindent \textbf{(P\_3,2.106-107.1)} KA\_II,114.19-115.2 Ro\_III,260 {14/14}     tat etat vāvacanam tiṣṭhatu tāvat sānnyāsikam . \\
\noindent \textbf{(P\_3,2.106-107.2)} KA\_II,115.3-8 Ro\_III,260-261 {1/14}     atha kitkaraṇam kimartham na asaṃyogāt liṭ kit iti eva siddham . \\
\noindent \textbf{(P\_3,2.106-107.2)} KA\_II,115.3-8 Ro\_III,260-261 {2/14}     kitkararaṇam saṃyogārtham . \\
\noindent \textbf{(P\_3,2.106-107.2)} KA\_II,115.3-8 Ro\_III,260-261 {3/14}     kitkararaṇam kriyate saṃyogārtham . \\
\noindent \textbf{(P\_3,2.106-107.2)} KA\_II,115.3-8 Ro\_III,260-261 {4/14}     saṃyogāntāḥ prayojayanti . \\
\noindent \textbf{(P\_3,2.106-107.2)} KA\_II,115.3-8 Ro\_III,260-261 {5/14}     vṛtrasya yat badbadhānasya rodasī . \\
\noindent \textbf{(P\_3,2.106-107.2)} KA\_II,115.3-8 Ro\_III,260-261 {6/14}     tvam arṇavān badbadhānāṃ aramṇāḥ . \\
\noindent \textbf{(P\_3,2.106-107.2)} KA\_II,115.3-8 Ro\_III,260-261 {7/14}     añjeḥ ājivān iti . \\
\noindent \textbf{(P\_3,2.106-107.2)} KA\_II,115.3-8 Ro\_III,260-261 {8/14}     chāndasau kānackvasū . \\
\noindent \textbf{(P\_3,2.106-107.2)} KA\_II,115.3-8 Ro\_III,260-261 {9/14}     liṭ ca chandasi sārvadhātukam api bhavati . \\
\noindent \textbf{(P\_3,2.106-107.2)} KA\_II,115.3-8 Ro\_III,260-261 {10/14}     tatra sārvadhātukam apit ṅit bhavati iti ṅittvāt lupadhālopaḥ bhaviṣyati . ṝkārāntaguṇapratiṣedhārtham vā . \\
\noindent \textbf{(P\_3,2.106-107.2)} KA\_II,115.3-8 Ro\_III,260-261 {11/14}      ṝkārāntaguṇapratiṣedhārtham tarhi kitkaraṇam kartavyam . \\
\noindent \textbf{(P\_3,2.106-107.2)} KA\_II,115.3-8 Ro\_III,260-261 {12/14}     ayam liṭi ṝkārāntānām pratiṣedhaviṣaye guṇaḥ ārabhyate . \\
\noindent \textbf{(P\_3,2.106-107.2)} KA\_II,115.3-8 Ro\_III,260-261 {13/14}     saḥ yathā eva iha pratiṣedham bādhitvā guṇaḥ bhavati teratuḥ teruḥ evam iha api syāt titīrvān tirirāṇaḥ . \\
\noindent \textbf{(P\_3,2.106-107.2)} KA\_II,115.3-8 Ro\_III,260-261 {14/14}     punaḥ kitkaraṇā pratiṣidhyate . \\
\noindent \textbf{(P\_3,2.108)} bhāṣāyām sadādibhyaḥ vā liṭ . \\
\noindent \textbf{(P\_3,2.108)} bhāṣāyām sadādibhyaḥ vā liṭ vaktavyaḥ . \\
\noindent \textbf{(P\_3,2.108)} kim prayojanam . \\
\noindent \textbf{(P\_3,2.108)} tadviṣaye luṅaḥ anivṛttyartham . \\
\noindent \textbf{(P\_3,2.108)} tasya liṭaḥ viṣaye luṅaḥ anivṛttiḥ yathā syāt . \\
\noindent \textbf{(P\_3,2.108)} upasedivān kautsaḥ pāṇinim . \\
\noindent \textbf{(P\_3,2.108)} upāsadat . \\
\noindent \textbf{(P\_3,2.108)} anadyatanaparokṣayoḥ ca . \\
\noindent \textbf{(P\_3,2.108)} anadyatanaparokṣayoḥ ca vā liṭ vaktavyaḥ . \\
\noindent \textbf{(P\_3,2.108)} upasedivān kautsaḥ pāṇinim . \\
\noindent \textbf{(P\_3,2.108)} upāsīdat . \\
\noindent \textbf{(P\_3,2.108)} upasasāda . \\
\noindent \textbf{(P\_3,2.108)} apavādavipratiṣedhāt hi tayoḥ bhāvaḥ . \\
\noindent \textbf{(P\_3,2.108)} apavādavipratiṣedhāt hi tau syātām . \\
\noindent \textbf{(P\_3,2.108)} kau . \\
\noindent \textbf{(P\_3,2.108)} laṅliṭau . \\
\noindent \textbf{(P\_3,2.108)} tasya kvasuḥ aparokṣe nityam . \\
\noindent \textbf{(P\_3,2.108)} tasya liṭaḥ bhāṣāyām kvasuḥ aparokṣe nityam iti vaktavyam . \\
\noindent \textbf{(P\_3,2.108)} aparokṣagrahaṇena na arthaḥ . \\
\noindent \textbf{(P\_3,2.108)} tasya kvasuḥ nityam iti eva . \\
\noindent \textbf{(P\_3,2.108)} kena idānīm liṭaḥ parokṣe śravaṇam bhaviṣyati . \\
\noindent \textbf{(P\_3,2.108)} parokṣe liṭ iti anena . \\
\noindent \textbf{(P\_3,2.108)} tat tarhi vaktavyam . \\
\noindent \textbf{(P\_3,2.108)} na vaktavyam . \\
\noindent \textbf{(P\_3,2.108)} anuvṛttiḥ kariṣyate . \\
\noindent \textbf{(P\_3,2.108)} bhāṣāyām sadādibhyaḥ vā liṭ bhavati liṭaḥ ca kvasuḥ bhavati . \\
\noindent \textbf{(P\_3,2.108)} tataḥ luṅ . \\
\noindent \textbf{(P\_3,2.108)} luṅ bhavati bhūte kāle . \\
\noindent \textbf{(P\_3,2.108)} bhāṣāyām sadādibhyaḥ vā liṭ bhavati liṭaḥ ca kvasuḥ bhavati . \\
\noindent \textbf{(P\_3,2.108)} tataḥ anadyatane laṅ . \\
\noindent \textbf{(P\_3,2.108)} anadyatane bhūte kāle laṅ bhavati . \\
\noindent \textbf{(P\_3,2.108)} bhāṣāyām sadādibhyaḥ vā liṭ bhavati liṭaḥ ca kvasuḥ bhavati . \\
\noindent \textbf{(P\_3,2.108)} parokṣe liṭ bhavati . \\
\noindent \textbf{(P\_3,2.108)} bhāṣāyām sadādibhyaḥ vā liṭ bhavati liṭaḥ ca kvasuḥ bhavati . \\
\noindent \textbf{(P\_3,2.108)} tatra ayam api arthaḥ . \\
\noindent \textbf{(P\_3,2.108)} tasya kvasuḥ aparokṣe nityam iti etat na vaktavyam bhavati . \\
\noindent \textbf{(P\_3,2.109)} kim upeyivān iti nipātanam kriyate . \\
\noindent \textbf{(P\_3,2.109)} upeyuṣi nipātanam iḍartham . \\
\noindent \textbf{(P\_3,2.109)} upeyuṣi nipātanam kriyate iḍartham . \\
\noindent \textbf{(P\_3,2.109)} iṭ yathā syāt . \\
\noindent \textbf{(P\_3,2.109)} na etat asti prayojanam . \\
\noindent \textbf{(P\_3,2.109)} siddhaḥ atra iṭ vasvekākādghasām iti . \\
\noindent \textbf{(P\_3,2.109)} dvirvacane kṛte anekāctvāt na prāpnoti . \\
\noindent \textbf{(P\_3,2.109)} idam iha sampradhāryam . \\
\noindent \textbf{(P\_3,2.109)} dvirvacanam kriyatām iṭ iti . \\
\noindent \textbf{(P\_3,2.109)} kim atra kartavyam . \\
\noindent \textbf{(P\_3,2.109)} paratvāt iḍāgamaḥ . \\
\noindent \textbf{(P\_3,2.109)} nityam dvirvacanam . \\
\noindent \textbf{(P\_3,2.109)} kṛte api iṭi prāpnoti akṛte api prāpnoti . \\
\noindent \textbf{(P\_3,2.109)} iṭ api nityaḥ . \\
\noindent \textbf{(P\_3,2.109)} kṛte api dvirvacane ekādeśe ca prāpnoti akṛte api prāpnoti . \\
\noindent \textbf{(P\_3,2.109)} na atra ekādeśaḥ prāpnoti . \\
\noindent \textbf{(P\_3,2.109)} kim kāraṇam . \\
\noindent \textbf{(P\_3,2.109)} dīrghaḥ iṇaḥ kiti iti dīrghatvena bādhate . \\
\noindent \textbf{(P\_3,2.109)} tat etat upeyuṣi nipātanam iḍartham kriyate . \\
\noindent \textbf{(P\_3,2.109)} upeyuṣi nipātanam iḍartham iti cet ajādau atiprasaṅgaḥ . upeyuṣi nipātanam iḍartham iti cet ajādau atiprasaṅgaḥ bhavati . \\
\noindent \textbf{(P\_3,2.109)} upeyuṣā upeyuṣe upeyuṣaḥ upeyuṣi iti . \\
\noindent \textbf{(P\_3,2.109)} ekādiṣṭasya īybhāvārtham tu . \\
\noindent \textbf{(P\_3,2.109)} ekādiṣṭasya īybhāvārtham tu nipātanam kriyate . \\
\noindent \textbf{(P\_3,2.109)} ekādiṣṭasya īy iti etat rūpam nipātyate . \\
\noindent \textbf{(P\_3,2.109)} nanu ca uktam na atra ekādeśaḥ prāpnoti . \\
\noindent \textbf{(P\_3,2.109)} kim kāraṇam . \\
\noindent \textbf{(P\_3,2.109)} dīrghaḥ iṇaḥ kiti iti dīrghatvena bādhate iti . \\
\noindent \textbf{(P\_3,2.109)} tat hi na suṣṭhu ucyate . \\
\noindent \textbf{(P\_3,2.109)} na hi dīrghatvam ekādeśam bādhate . \\
\noindent \textbf{(P\_3,2.109)} kaḥ tarhi bādhate . \\
\noindent \textbf{(P\_3,2.109)} yaṇādeśaḥ . \\
\noindent \textbf{(P\_3,2.109)} saḥ ca kva bādhate . \\
\noindent \textbf{(P\_3,2.109)} yatra asya nimittam asti . \\
\noindent \textbf{(P\_3,2.109)} yatra hi nimittam na asti niṣpratidvandvaḥ tatra ekādeśaḥ . \\
\noindent \textbf{(P\_3,2.109)} vyañjane yaṇādeśārtham vā . \\
\noindent \textbf{(P\_3,2.109)} atha vā vyañjane eva yaṇādeśaḥ nipātyate . \\
\noindent \textbf{(P\_3,2.109)} yaṇādeśe kṛte ekācaḥ iti iṭ siddhaḥ bhavati . \\
\noindent \textbf{(P\_3,2.109)} aparaḥ āha : na upeyivān nipātyaḥ . \\
\noindent \textbf{(P\_3,2.109)} dvirvacanād iṭ bhaviṣyati paratvāt . \\
\noindent \textbf{(P\_3,2.109)} dvirvacanam kriyatām iṭ iti iṭ bhaviṣyati vipratiṣedhena . \\
\noindent \textbf{(P\_3,2.109)} iha api tarhi dvirvacanāt iṭ syāt bibhidvān cicchidvān . \\
\noindent \textbf{(P\_3,2.109)} anyeṣām ekācām dvirvacanam nityam iti āhuḥ . anyeṣām ekācām nityam dvirvacanam . \\
\noindent \textbf{(P\_3,2.109)} kṛte api iṭi prāpnoti akṛte api prāpnoti . \\
\noindent \textbf{(P\_3,2.109)} asya punaḥ iṭ ca nityaḥ dvivracanam ca . asya punaḥ iṭ ca eva nityaḥ dvivracanam ca . \\
\noindent \textbf{(P\_3,2.109)} dvirvacane ca kṛte ekāc bhavati . \\
\noindent \textbf{(P\_3,2.109)} katham . \\
\noindent \textbf{(P\_3,2.109)} ekādeśe kṛte . \\
\noindent \textbf{(P\_3,2.109)} tasmāt iṭ bādhate dvitvam . \\
\noindent \textbf{(P\_3,2.109)} tasmāt iṭ dvirvacanam bādhate . \\
\noindent \textbf{(P\_3,2.109)} anūcānaḥ kartari . \\
\noindent \textbf{(P\_3,2.109)} anūcānaḥ kartarīti vaktavyam . \\
\noindent \textbf{(P\_3,2.109)} anūktavān anūcānaḥ . \\
\noindent \textbf{(P\_3,2.109)} anūktam iti eva anyatra . \\
\noindent \textbf{(P\_3,2.109)} na upeyivān nipātyaḥ . \\
\noindent \textbf{(P\_3,2.109)} dvirvacanād iṭ bhaviṣyati paratvāt . \\
\noindent \textbf{(P\_3,2.109)} anyeṣām ekācām dvirvacanam nityam iti āhuḥ . \\
\noindent \textbf{(P\_3,2.109)} asya punaḥ iṭ ca nityaḥ dvivracanam ca . \\
\noindent \textbf{(P\_3,2.109)} na vihanyate hi asya . \\
\noindent \textbf{(P\_3,2.109)} dvirvacane ca ekāctvāt . \\
\noindent \textbf{(P\_3,2.109)} tasmāt iṭ bādhate dvitvam . \\
\noindent \textbf{(P\_3,2.110.1)} luṅlṛṭoḥ apavādaprasaṅgaḥ bhūtabhaviṣyatoḥ aviśeṣavacanāt . \\
\noindent \textbf{(P\_3,2.110.1)} luṅlṛṭoḥ apavādaḥ prāpnoti . \\
\noindent \textbf{(P\_3,2.110.1)} agāma ghoṣān . \\
\noindent \textbf{(P\_3,2.110.1)} apāma payaḥ . \\
\noindent \textbf{(P\_3,2.110.1)} aśayiṣmahi pūtīkatṛṇeṣu . \\
\noindent \textbf{(P\_3,2.110.1)} gamiṣyāmaḥ ghoṣān . \\
\noindent \textbf{(P\_3,2.110.1)} pāsyāmaḥ payaḥ . \\
\noindent \textbf{(P\_3,2.110.1)} śayiṣyāmahe pūtīkatṛṇeṣu . \\
\noindent \textbf{(P\_3,2.110.1)} kim kāraṇam . \\
\noindent \textbf{(P\_3,2.110.1)} bhūtabhaviṣyatoḥ aviśeṣavacanāt . \\
\noindent \textbf{(P\_3,2.110.1)} bhūtabhaviṣyatoḥ aviśeṣeṇa vidhīyete luṅlṛṭau . \\
\noindent \textbf{(P\_3,2.110.1)} tayoḥ viśeṣavihitau lanluṭau apavādau prāpnutaḥ . \\
\noindent \textbf{(P\_3,2.110.1)} na vā apavādasya nimittābhāvāt anadyatane hi tayoḥ vidhānam . \\
\noindent \textbf{(P\_3,2.110.1)} na vā eṣaḥ doṣaḥ . \\
\noindent \textbf{(P\_3,2.110.1)} kim kāraṇam . \\
\noindent \textbf{(P\_3,2.110.1)} apavādasya nimittābhāvāt . \\
\noindent \textbf{(P\_3,2.110.1)} na atra apavādasya nimittam asti . \\
\noindent \textbf{(P\_3,2.110.1)} kim kāraṇam . \\
\noindent \textbf{(P\_3,2.110.1)} anadyatane hi tayoḥ vidhānam . \\
\noindent \textbf{(P\_3,2.110.1)} anadyatane hi tau vidhīyete laṅluṭau . \\
\noindent \textbf{(P\_3,2.110.1)} na ca atra anadyatanaḥ kālaḥ vivakṣitaḥ . \\
\noindent \textbf{(P\_3,2.110.1)} kim tarhi . \\
\noindent \textbf{(P\_3,2.110.1)} bhūtakālasāmānyam bhaviṣyatkālasāmānyam ca . \\
\noindent \textbf{(P\_3,2.110.1)} yadi api tāvat atra etat śakyate vaktum gamiṣyāmaḥ ghoṣān pāsyāmaḥ payaḥ śayiṣyāmahe pūtīkatṛṇeṣu iti yatra etat na jñāyate kim kadā iti . \\
\noindent \textbf{(P\_3,2.110.1)} iha tu katham agāma ghoṣān apāma payaḥ aśayiṣmahi pūtīkatṛṇeṣu yatra etat nirjñātam bhavati amuṣmin ahani gatam iti . \\
\noindent \textbf{(P\_3,2.110.1)} atra api na vā apavādasya nimittābhāvāt anadyatane hi tayoḥ vidhānam iti eva . \\
\noindent \textbf{(P\_3,2.110.1)} katham punaḥ sataḥ nāma avivakṣā syāt . \\
\noindent \textbf{(P\_3,2.110.1)} sataḥ api avivakṣā bhavati . \\
\noindent \textbf{(P\_3,2.110.1)} tat yathā . \\
\noindent \textbf{(P\_3,2.110.1)} alomikā eḍakā . \\
\noindent \textbf{(P\_3,2.110.1)} anudarā kanyā . \\
\noindent \textbf{(P\_3,2.110.1)} asataḥ ca vivakṣā bhavati . \\
\noindent \textbf{(P\_3,2.110.1)} tat yathā . \\
\noindent \textbf{(P\_3,2.110.1)} samudraḥ kuṇḍikā . \\
\noindent \textbf{(P\_3,2.110.1)} vindhyaḥ vardhitakam iti . \\
\noindent \textbf{(P\_3,2.110.2)} vaseḥ luṅ rātriśeṣe . \\
\noindent \textbf{(P\_3,2.110.2)} vaseḥ luṅ rātriśeṣe vaktavyaḥ . \\
\noindent \textbf{(P\_3,2.110.2)} nyāyye pratyutthāne pratyutthitam kaḥ cit kam cit pṛcchati . \\
\noindent \textbf{(P\_3,2.110.2)} kva bhavān uṣitaḥ iti . \\
\noindent \textbf{(P\_3,2.110.2)} saḥ āha . \\
\noindent \textbf{(P\_3,2.110.2)} amutra avātsam iti . \\
\noindent \textbf{(P\_3,2.110.2)} amutra avasam iti prāpnoti . \\
\noindent \textbf{(P\_3,2.110.2)} jāgaraṇasantatau . \\
\noindent \textbf{(P\_3,2.110.2)} jāgaraṇasantatau iti vaktavyam . \\
\noindent \textbf{(P\_3,2.110.2)} yaḥ hi muhūrtamātram api svapiti tatra amutra avasam iti eva bhavitavyam . \\
\noindent \textbf{(P\_3,2.111)} anadyatane iti bahuvrīhinirdeśaḥ adya hyaḥ abhukṣmahi iti . \\
\noindent \textbf{(P\_3,2.111)} anadyatane iti bahuvrīhinirdeśaḥ kartavyaḥ . \\
\noindent \textbf{(P\_3,2.111)} avidyamānādyatane anadyatane iti . \\
\noindent \textbf{(P\_3,2.111)} kim prayojanam . \\
\noindent \textbf{(P\_3,2.111)} adya hyaḥ abhukṣmahi iti . \\
\noindent \textbf{(P\_3,2.111)} adya ca hyaḥ ca abhukṣmahi iti vyāmiśre luṅ eva yathā syāt . \\
\noindent \textbf{(P\_3,2.111)} yadi evam adyatane api laṅ prāpnoti . \\
\noindent \textbf{(P\_3,2.111)} na hi adyatane adyatanaḥ vidyate . \\
\noindent \textbf{(P\_3,2.111)} adyatane api adyatanaḥ vidyate . \\
\noindent \textbf{(P\_3,2.111)} katham . \\
\noindent \textbf{(P\_3,2.111)} vyapadeśivadbhāvena . \\
\noindent \textbf{(P\_3,2.111)} parokṣe ca lokavijñāte prayoktuḥ darśanaviṣaye . \\
\noindent \textbf{(P\_3,2.111)} parokṣe ca lokavijñāte prayoktuḥ darśanaviṣaye laṅ vaktavyaḥ . \\
\noindent \textbf{(P\_3,2.111)} aruṇat yavanaḥ sāketam . \\
\noindent \textbf{(P\_3,2.111)} aruṇat yavanaḥ madhamikām . \\
\noindent \textbf{(P\_3,2.111)} parokṣe iti kimartham . \\
\noindent \textbf{(P\_3,2.111)} udagāt ādityaḥ . \\
\noindent \textbf{(P\_3,2.111)} lokavijñate iti kimartham . \\
\noindent \textbf{(P\_3,2.111)} cakāra kaṭam devadattaḥ . \\
\noindent \textbf{(P\_3,2.111)} prayoktuḥ darśanaviṣaye iti kimartham . \\
\noindent \textbf{(P\_3,2.111)} jaghāna kaṃsam kila vāsudevaḥ . \\
\noindent \textbf{(P\_3,2.114)} kim udāharaṇam . \\
\noindent \textbf{(P\_3,2.114)} tatra saktūn pāsyāmaḥ . \\
\noindent \textbf{(P\_3,2.114)} abhijānāsi devadatta tatra saktūn apibāma . \\
\noindent \textbf{(P\_3,2.114)} bhavet pūrvam param ākāṅkṣati iti sākāṅkṣam syāt . \\
\noindent \textbf{(P\_3,2.114)} param tu katham sākāṅkṣam . \\
\noindent \textbf{(P\_3,2.114)} param api sākāṅkṣam . \\
\noindent \textbf{(P\_3,2.114)} katham asti asmin ākāṅkṣā iti ataḥ sākāṅkṣam . \\
\noindent \textbf{(P\_3,2.114)} vibhāṣā sākāṅkṣe sarvatra . \\
\noindent \textbf{(P\_3,2.114)} vibhāṣā sākāṅkṣe sarvatra iti vaktavyam . \\
\noindent \textbf{(P\_3,2.114)} kva sarvatra . \\
\noindent \textbf{(P\_3,2.114)} yadi ca ayadi ca . \\
\noindent \textbf{(P\_3,2.114)} yadi tāvat . \\
\noindent \textbf{(P\_3,2.114)} abhijānāsi devadatta yat kaśmīrān gamiṣyāmaḥ yat kaśmīrān agacchāma yat tatra odanam bhokṣyāmahe yat tatra odanam abhuñjmahi . \\
\noindent \textbf{(P\_3,2.114)} abhijānāsi devadatta kaśmīrān gamiṣyāmaḥ kaśmīrān agacchāma tatra odanam bhokṣyāmahe tatra odanam abhuñjmahi . \\
\noindent \textbf{(P\_3,2.115.1)} parokṣe iti ucyate . \\
\noindent \textbf{(P\_3,2.115.1)} kim parokṣam nāma . \\
\noindent \textbf{(P\_3,2.115.1)} param akṣṇaḥ parokṣam . \\
\noindent \textbf{(P\_3,2.115.1)} akṣi punaḥ kim . \\
\noindent \textbf{(P\_3,2.115.1)} aśnoteḥ ayam auṇādikaḥ karaṇasādhanaḥ si pratyayaḥ . \\
\noindent \textbf{(P\_3,2.115.1)} anute anena iti akṣi . \\
\noindent \textbf{(P\_3,2.115.1)} yadi evam parākṣam iti prāpnoti . \\
\noindent \textbf{(P\_3,2.115.1)} na eṣaḥ doṣaḥ . \\
\noindent \textbf{(P\_3,2.115.1)} parobhāvaḥ parasya akṣe parokṣe liṭi dṛśyatām . paraśabdasya akṣaśabde uttarapade parobhāvaḥ vaktavyaḥ . \\
\noindent \textbf{(P\_3,2.115.1)} utvam vā ādeḥ parāt akṣṇaḥ .atha vā paraśabdāt uttarasya akṣiśabdasya utvam vaktavyam . \\
\noindent \textbf{(P\_3,2.115.1)} siddham vā asmāt nipātanāt . \\
\noindent \textbf{(P\_3,2.115.2)} kasmin punaḥ parokṣe . \\
\noindent \textbf{(P\_3,2.115.2)} kāle . \\
\noindent \textbf{(P\_3,2.115.2)} na vai kālādhikāraḥ asti . \\
\noindent \textbf{(P\_3,2.115.2)} evam tarhi dhātoḥ iti vartate . \\
\noindent \textbf{(P\_3,2.115.2)} dhātau parokṣe . \\
\noindent \textbf{(P\_3,2.115.2)} dhātuḥ vai śabdaḥ . \\
\noindent \textbf{(P\_3,2.115.2)} na ca śabdasya pratyakṣaparokṣatāyām sambhavaḥ asti . \\
\noindent \textbf{(P\_3,2.115.2)} śabde asambhavāt arthe kāryam vijñāsyate . \\
\noindent \textbf{(P\_3,2.115.2)} parokṣe dhātau parokṣe dhātvarthe iti . \\
\noindent \textbf{(P\_3,2.115.2)} kaḥ punaḥ dhātvarthaḥ . \\
\noindent \textbf{(P\_3,2.115.2)} kriyā . \\
\noindent \textbf{(P\_3,2.115.2)} kriyāyām parokṣāyām . \\
\noindent \textbf{(P\_3,2.115.2)} yadi evam hyaḥ apacat iti atra api liṭ prāpnoti . \\
\noindent \textbf{(P\_3,2.115.2)} kim kāraṇam . \\
\noindent \textbf{(P\_3,2.115.2)} kriyā nāma iyam atyantāparidṛṣṭā anumānagamyā aśakyā piṇḍībhūtā nidarśayitum yathā garbhaḥ nirluṭhitaḥ . \\
\noindent \textbf{(P\_3,2.115.2)} evam tarhi sādhaneṣu parokṣeṣu . \\
\noindent \textbf{(P\_3,2.115.2)} sādhaneṣu ca bhavataḥ kaḥ sampratyayaḥ . \\
\noindent \textbf{(P\_3,2.115.2)} yadi sāvad guṇasamudāyaḥ sādhanam sādhanam api anumānagamyam . \\
\noindent \textbf{(P\_3,2.115.2)} atha anyat guṇebhyaḥ sādhanam bhavati prtayakṣaparokṣatāyām sambhavaḥ . \\
\noindent \textbf{(P\_3,2.115.2)} atha yadā anena rathyāyām taṇḍulodakam dṛṣṭam katham tatra bhavitavyam . \\
\noindent \textbf{(P\_3,2.115.2)} yadi tāvat sādhaneṣu parokṣeṣu papāca iti bhavitavyam . \\
\noindent \textbf{(P\_3,2.115.2)} bhavanti hi tasya sādhanāni parokṣāṇi . \\
\noindent \textbf{(P\_3,2.115.2)} atha ye ete kriyākṛtāḥ viśeṣāḥ cītkārāḥ phūtkārāḥ ca teṣu parokṣeṣu evam api papāca iti bhavitavyam . \\
\noindent \textbf{(P\_3,2.115.2)} kathañjātīyakam punaḥ parokṣam nāma . \\
\noindent \textbf{(P\_3,2.115.2)} ke cit tāvat āhuḥ . \\
\noindent \textbf{(P\_3,2.115.2)} varṣaśatavṛttam parokṣam iti . \\
\noindent \textbf{(P\_3,2.115.2)} apare āhuḥ . \\
\noindent \textbf{(P\_3,2.115.2)} kaṭāntaritam parokṣam iti . \\
\noindent \textbf{(P\_3,2.115.2)} apare āhūḥ . \\
\noindent \textbf{(P\_3,2.115.2)} dvyahavṛttam tryahvṛttam ca iti . \\
\noindent \textbf{(P\_3,2.115.2)} sarvathā uttamaḥ na sidhyati . \\
\noindent \textbf{(P\_3,2.115.2)} suptamattoyoḥ iti vaktavyam . \\
\noindent \textbf{(P\_3,2.115.2)} suptaḥ aham kila vilalāpa . \\
\noindent \textbf{(P\_3,2.115.2)} mattaḥ aham kila vilalāpa . \\
\noindent \textbf{(P\_3,2.115.2)} suptaḥ nu aham kila vilalāpa . \\
\noindent \textbf{(P\_3,2.115.2)} mattaḥ nu aham kila vilalāpa . \\
\noindent \textbf{(P\_3,2.115.2)} atha vā bhavati vai kaḥ cit jāgarat api vartamānakālam na upalabhate . \\
\noindent \textbf{(P\_3,2.115.2)} tat yathā vaiyākaraṇānām śākaṭāyanaḥ rathamārge āsīnaḥ śakaṭasāṛtham yāntam na upalabhate . \\
\noindent \textbf{(P\_3,2.115.2)} kim punaḥ kāraṇam jāgarat api vartamānakālam na upalabhate . \\
\noindent \textbf{(P\_3,2.115.2)} manasā saṃyuktāni indriyāṇi upalabdhau kāraṇāni bhavanti . \\
\noindent \textbf{(P\_3,2.115.2)} manasaḥ asānnidhyāt . \\
\noindent \textbf{(P\_3,2.115.3)} parokṣe liṭ atyantāapahnave ca . \\
\noindent \textbf{(P\_3,2.115.3)} parokṣe liṭ iti atra atyantāapahnave ca iti vaktavyam . \\
\noindent \textbf{(P\_3,2.115.3)} no khaṇḍikān jagāma no kaliṅgān jagāma . \\
\noindent \textbf{(P\_3,2.115.3)} na kārisomam prapau . \\
\noindent \textbf{(P\_3,2.115.3)} na dārvajasya pratijagrāha . \\
\noindent \textbf{(P\_3,2.115.3)} kaḥ me manuṣyaḥ praharet vadhāya . \\
\noindent \textbf{(P\_3,2.115.3)} parobhāvaḥ parasya akṣe parokṣe liṭi dṛśyatām . \\
\noindent \textbf{(P\_3,2.115.3)} utvam vā ādeḥ parāt akṣṇaḥ . \\
\noindent \textbf{(P\_3,2.115.3)} siddham vā asmāt nipātanāt . \\
\noindent \textbf{(P\_3,2.118)} sma purā bhūtamātre na sma purā adyatane . \\
\noindent \textbf{(P\_3,2.118)} sma purā bhūtamātre na sma purā adyatane iti vaktavyam . \\
\noindent \textbf{(P\_3,2.118)} kim ayam smādividhiḥ purāntaḥ aviśeṣeṇa bhūtamātre bhavati . \\
\noindent \textbf{(P\_3,2.118)} tatra vaktavyam smalakṣaṇaḥ purālakṣaṇaḥ ca adyatane na bhavataḥ iti . \\
\noindent \textbf{(P\_3,2.118)} āhosvit smalakṣaṇaḥ purālakṣaṇaḥ ca aviśeṣeṇa bhūtamātre bhavataḥ iti . \\
\noindent \textbf{(P\_3,2.118)} tatra smādyartham na sma purā adyatane iti vaktavyam . \\
\noindent \textbf{(P\_3,2.118)} kim ca ataḥ . \\
\noindent \textbf{(P\_3,2.118)} smādividhiḥ purāntaḥ yadi aviśeṣeṇa kim kṛtam bhavati na sma purā adyatane iti bruvatā kātyāyanena iha . \\
\noindent \textbf{(P\_3,2.118)} smādividhiḥ purāntaḥ yadi aviśeṣeṇa bhavati kim vārttikakāraḥ pratiṣedhena karoti na sma purā adyatane iti . \\
\noindent \textbf{(P\_3,2.118)} anuvṛttiḥ anadyatanasya lāt sme iti tatra na asti nañkāryam . laṭ sme iti atra anadyatane iti etat anuvartiṣyate . \\
\noindent \textbf{(P\_3,2.118)} aparokṣānadyatanaḥ nanau ca nanvoḥ ca nivṛttau na purā adyatane iti bhavet etat vācyam . tatra etāvat vaktavyam syāt na purā adyatane iti . \\
\noindent \textbf{(P\_3,2.118)} tatra ca api laṅgrahaṇam . tatra ca api laṅgrahaṇam jñapakam na purālakṣaṇaḥ adyatane bhavati iti . \\
\noindent \textbf{(P\_3,2.118)} atha buddhiḥ aviśeṣād sma purā hetū . \\
\noindent \textbf{(P\_3,2.118)} atha buddhiḥ aviśeṣeṇa sma purā hetū iti . \\
\noindent \textbf{(P\_3,2.118)} tatra ca api śrṇu bhūyaḥ . \\
\noindent \textbf{(P\_3,2.118)} aparokṣe ce iti eṣaḥ prāk purisaṃśabdanāt avinivṛttaḥ sarvatra anadyatanaḥ . \\
\noindent \textbf{(P\_3,2.118)} tathā sati nañā kim iha kāryam . \\
\noindent \textbf{(P\_3,2.118)} smādau aparokṣe ca iti akāryam iti śakyam etat api viddhi . \\
\noindent \textbf{(P\_3,2.118)} śakyam hi nivartayitum parokṣe iti lāt sme iti atra . \\
\noindent \textbf{(P\_3,2.118)} syāt eṣā tava buddhiḥ . \\
\noindent \textbf{(P\_3,2.118)} smalakṣaṇe api evam eva siddham iti . \\
\noindent \textbf{(P\_3,2.118)} laṭ sme iti bhavet na arthaḥ . \\
\noindent \textbf{(P\_3,2.118)} tasmāt kāryam parārtham tu . evam tarhi jñāpayati ācāryaḥ smalakṣaṇaḥ purālakṣaṇaḥ ca anadyatane bhavataḥ iti . \\
\noindent \textbf{(P\_3,2.120)} nanau pṛṣṭaprativacane iti aśiṣyam kriyāsamāpteḥ vivakṣitatvāt . \\
\noindent \textbf{(P\_3,2.120)} nanau pṛṣṭaprativacane iti aśiṣyaḥ laṭ . \\
\noindent \textbf{(P\_3,2.120)} kim kāraṇam . \\
\noindent \textbf{(P\_3,2.120)} kriyāsamāpteḥ vivakṣitatvāt . \\
\noindent \textbf{(P\_3,2.120)} kriyāyāḥ atra asamāptiḥ vivakṣitā . \\
\noindent \textbf{(P\_3,2.120)} eṣaḥ nāma nyāyyaḥ vartamānaḥ kālaḥ yatra kriyāyāḥ asamāptiḥ bhavat . \\
\noindent \textbf{(P\_3,2.120)} tatra vartamāne laṭ iti eva siddham . \\
\noindent \textbf{(P\_3,2.120)} yadi vartamāne laṭ iti eva laṭ bhavati śatṛśānacau api prāpnutaḥ . \\
\noindent \textbf{(P\_3,2.120)} iṣyete śatṛśānacau : nanu mām kurvantam paśya . \\
\noindent \textbf{(P\_3,2.120)} nanu mām kurvāṇam paśya iti . \\
\noindent \textbf{(P\_3,2.122)} haśaśvadbhyām purā . \\
\noindent \textbf{(P\_3,2.122)} haśaśvallakṣaṇāt purālakṣaṇaḥ bhavati vipratiṣedhena . \\
\noindent \textbf{(P\_3,2.122)} haśaśvallakṣaṇasya avakāśaḥ . \\
\noindent \textbf{(P\_3,2.122)} iti ha akarot . \\
\noindent \textbf{(P\_3,2.122)} iha ha cakāra . \\
\noindent \textbf{(P\_3,2.122)} śaśvat akarot . \\
\noindent \textbf{(P\_3,2.122)} śaśvat cakāra . \\
\noindent \textbf{(P\_3,2.122)} purālakṣaṇasya avakāśaḥ . \\
\noindent \textbf{(P\_3,2.122)} rathena ayam purā yāti . \\
\noindent \textbf{(P\_3,2.122)} rathena ayam purā ayāsīt . \\
\noindent \textbf{(P\_3,2.122)} iha ubhayam prāpnoti . \\
\noindent \textbf{(P\_3,2.122)} rathena ha śaśvat purā yāti . \\
\noindent \textbf{(P\_3,2.122)} rathena ha śaśvat purā ayāsīt . \\
\noindent \textbf{(P\_3,2.122)} purālakṣaṇaḥ bhavati vipratiṣedhena . \\
\noindent \textbf{(P\_3,2.122)} smaḥ sarvebhyaḥ vipratiṣedhena . \\
\noindent \textbf{(P\_3,2.122)} smalakṣaṇaḥ sarvebhyaḥ bhavati vipratiṣedhena . \\
\noindent \textbf{(P\_3,2.122)} haśaśvallakṣaṇāt purālakṣaṇāt ca . \\
\noindent \textbf{(P\_3,2.122)} haśaśvallakṣaṇasya avakāśaḥ . \\
\noindent \textbf{(P\_3,2.122)} iti ha akarot . \\
\noindent \textbf{(P\_3,2.122)} iha ha cakāra . \\
\noindent \textbf{(P\_3,2.122)} śaśvat akarot . \\
\noindent \textbf{(P\_3,2.122)} śaśvat cakāra . \\
\noindent \textbf{(P\_3,2.122)} purālakṣaṇasya avakāśaḥ . \\
\noindent \textbf{(P\_3,2.122)} rathena ayam purā yāti . \\
\noindent \textbf{(P\_3,2.122)} rathena ayam purā ayāsīt . \\
\noindent \textbf{(P\_3,2.122)} smalakṣaṇasya avakāśaḥ . \\
\noindent \textbf{(P\_3,2.122)} dharmeṇa sma kuravaḥ yudhyante . \\
\noindent \textbf{(P\_3,2.122)} iha sarvam prāpnoti . \\
\noindent \textbf{(P\_3,2.122)} na ha sma vai purā śaśvat aparaśuvṛkṇam dahati . \\
\noindent \textbf{(P\_3,2.122)} smalakṣaṇaḥ laṭ bhavati vipratiṣedhena . \\
\noindent \textbf{(P\_3,2.123)} pravṛttasya avirāme śiṣyā bhavantī avartamānatvāt . \\
\noindent \textbf{(P\_3,2.123)} pravṛttasya avirāme śiṣyā bhavantī . \\
\noindent \textbf{(P\_3,2.123)} iha adhīmahe . \\
\noindent \textbf{(P\_3,2.123)} iha vasāmaḥ . \\
\noindent \textbf{(P\_3,2.123)} iha puṣyamitram yājayāmaḥ . \\
\noindent \textbf{(P\_3,2.123)} kim punaḥ kāraṇam na sidhyati . \\
\noindent \textbf{(P\_3,2.123)} avartamānatvāt . \\
\noindent \textbf{(P\_3,2.123)} nityapravṛtte ca kālāvibhāgāt . \\
\noindent \textbf{(P\_3,2.123)} nityapravṛtte ca śāsitavyā bhavantī . \\
\noindent \textbf{(P\_3,2.123)} tiṣṭhanti parvatāḥ . \\
\noindent \textbf{(P\_3,2.123)} sravanti nadyaḥ iti . \\
\noindent \textbf{(P\_3,2.123)} kim punaḥ kāraṇam na sidhyati . \\
\noindent \textbf{(P\_3,2.123)} kālāvibhāgāt . \\
\noindent \textbf{(P\_3,2.123)} iha bhūtabhaviṣyatpratidvandvaḥ vartamānaḥ kālaḥ . \\
\noindent \textbf{(P\_3,2.123)} na ca atra bhūtabhaviṣyantau kālau staḥ . \\
\noindent \textbf{(P\_3,2.123)} nyāyyā tu ārambhānapavargāt . \\
\noindent \textbf{(P\_3,2.123)} nyāyyā tu eṣā vartamānakālatā . \\
\noindent \textbf{(P\_3,2.123)} kutaḥ . \\
\noindent \textbf{(P\_3,2.123)} ārambhānapavargāt . \\
\noindent \textbf{(P\_3,2.123)} ārambhaḥ atra anapavṛktaḥ . \\
\noindent \textbf{(P\_3,2.123)} eṣaḥ nām nyāyyaḥ vartamānaḥ kālaḥ yatra ārambhaḥ anapavṛktaḥ . \\
\noindent \textbf{(P\_3,2.123)} asti ca muktasaṃśaye virāmaḥ . \\
\noindent \textbf{(P\_3,2.123)} yam khalu api bhavān muktasaṃśayam vartamānam kālam nyāyyam manyate bhuṅkte devadattaḥ iti tena etat tulyam . \\
\noindent \textbf{(P\_3,2.123)} saḥ api hi avaśyam bhuñjānaḥ hasati vā jalpati vā pānīyam vā pibati . \\
\noindent \textbf{(P\_3,2.123)} yadi atra yuktā vartamānakālatā dṛśyate iha api yuktā dṛśyatām . \\
\noindent \textbf{(P\_3,2.123)} santi ca kālavibhāgāḥ . \\
\noindent \textbf{(P\_3,2.123)} santi khalu api kālavibhāgāḥ . \\
\noindent \textbf{(P\_3,2.123)} tiṣṭhanti parvatāḥ . \\
\noindent \textbf{(P\_3,2.123)} sthāsyanti parvatāḥ . \\
\noindent \textbf{(P\_3,2.123)} tasthuḥ parvatāḥ . \\
\noindent \textbf{(P\_3,2.123)} kim śakyante ete śabdāḥ prayoktum iti ataḥ santi kālavibhāgāḥ . \\
\noindent \textbf{(P\_3,2.123)} na avaśyam prayogāt eva . \\
\noindent \textbf{(P\_3,2.123)} iha bhūtabhaviṣyadvartamānānām rājñām yāḥ kriyāḥ tāḥ tiṣṭhateḥ adhikaraṇam . \\
\noindent \textbf{(P\_3,2.123)} iha tāvat tiṣṭhanti parvatāḥ iti . \\
\noindent \textbf{(P\_3,2.123)} samprati ye rājānaḥ teṣām yāḥ kriyāḥ tāsu vartamānāsu . \\
\noindent \textbf{(P\_3,2.123)} sthāsyanti parvatāḥ iti . \\
\noindent \textbf{(P\_3,2.123)} itaḥ uttaram ye rājānaḥ bhaviṣyanti teṣām yāḥ kriyāḥ tāsu bhaviṣyantīṣu . \\
\noindent \textbf{(P\_3,2.123)} tasthuḥ parvatāḥ iti . \\
\noindent \textbf{(P\_3,2.123)} ye rājānaḥ babhūvuḥ teṣām yāḥ kriyāḥ tāsu vbhūtāsu . \\
\noindent \textbf{(P\_3,2.123)} aparaḥ āha : na asti vartamānaḥ kālaḥ iti . \\
\noindent \textbf{(P\_3,2.123)} api ca atra ślokān udāharanti : na vartate cakram . \\
\noindent \textbf{(P\_3,2.123)} iṣuḥ na pātyate . \\
\noindent \textbf{(P\_3,2.123)} na syandante saritaḥ sāgarāya . \\
\noindent \textbf{(P\_3,2.123)} kūṭasthaḥ ayam lokaḥ na viceṣṭā asti . \\
\noindent \textbf{(P\_3,2.123)} yaḥ hi evam paśyati saḥ api anandhaḥ . \\
\noindent \textbf{(P\_3,2.123)} mīmāṃsakaḥ manyamānaḥ yuvā medhāvisammataḥ kākam sma iha anupṛcchati : kim te patitalakṣaṇam . \\
\noindent \textbf{(P\_3,2.123)} anāgate na patasi atikrānte ca kāka na . \\
\noindent \textbf{(P\_3,2.123)} yadi samprati patasi sarvaḥ lokaḥ patati ayam . \\
\noindent \textbf{(P\_3,2.123)} himavān api gacchati . \\
\noindent \textbf{(P\_3,2.123)} anāgatam atikrāntam vartamānam iti trayam . \\
\noindent \textbf{(P\_3,2.123)} sarvatra gatiḥ na asti . \\
\noindent \textbf{(P\_3,2.123)} gacchati iti kim ucyate . \\
\noindent \textbf{(P\_3,2.123)} kriyāpravṛttau yaḥ hetuḥ tadartham yat viceṣṭitam tat samīkṣya prayuñjīta gacchati iti avicārayan . \\
\noindent \textbf{(P\_3,2.123)} apraḥ āha . \\
\noindent \textbf{(P\_3,2.123)} asti vartamānaḥ kālaḥ iti . \\
\noindent \textbf{(P\_3,2.123)} ādityagativat na upalabhyate . \\
\noindent \textbf{(P\_3,2.123)} api ca atra ślokam udāharanti . \\
\noindent \textbf{(P\_3,2.123)} bisasya vālāḥ iva dahyamānāḥ na lakṣyate vikṛtiḥ sannipāte . \\
\noindent \textbf{(P\_3,2.123)} asti iti tām vedayante tribhabhāvāḥ . \\
\noindent \textbf{(P\_3,2.123)} sūkṣmaḥ hi bhāvaḥ anumitena gamyaḥ . \\
\noindent \textbf{(P\_3,2.124.1)} lasya aprathamāsamānādhikaraṇena ayogāt adeśānupapattiḥ yathā anyatra . \\
\noindent \textbf{(P\_3,2.124.1)} lasya aprathamāsamānādhikaraṇena ayogāt adeśayoḥ nupapattiḥ yathā anyatra . \\
\noindent \textbf{(P\_3,2.124.1)} tat yathā anyatra api lasya aprathamāsamānādhikaraṇena yogaḥ na bhavati . \\
\noindent \textbf{(P\_3,2.124.1)} kva anyatra . \\
\noindent \textbf{(P\_3,2.124.1)} laṅi . \\
\noindent \textbf{(P\_3,2.124.1)} apacat odanam devadattaḥ iti . \\
\noindent \textbf{(P\_3,2.124.1)} yogaḥ iti cet anyatra api yogaḥ syāt . \\
\noindent \textbf{(P\_3,2.124.1)} atha matam etat bhavati yogaḥ iti anyatra api yogaḥ syāt . \\
\noindent \textbf{(P\_3,2.124.1)} kva anyatra . \\
\noindent \textbf{(P\_3,2.124.1)} laṅi . \\
\noindent \textbf{(P\_3,2.124.1)} apacat odanam devadattaḥ iti . \\
\noindent \textbf{(P\_3,2.124.1)} na kva cit yogaḥ iti kṛtvā ataḥ sarvatra yogena bhavitavyam kva cit ayogaḥ iti kṛtvā sarvatra ayogena . \\
\noindent \textbf{(P\_3,2.124.1)} tat yathā . \\
\noindent \textbf{(P\_3,2.124.1)} samānam īhamānānām ca adhīyānānām ca ke cit arthaiḥ yujyante apare na . \\
\noindent \textbf{(P\_3,2.124.1)} na ca idānīm kaḥ cit arthavān iti ataḥ sarvaiḥ arthavadbhiḥ śakyam bhavitum kaḥ cit anarthakaḥ iti sarvaiḥ anarthakaiḥ . \\
\noindent \textbf{(P\_3,2.124.1)} tatra kim asmābhiḥ śakyam kartum . \\
\noindent \textbf{(P\_3,2.124.1)} yat loṭaḥ aprathamāsamānādhikaraṇena yogaḥ bhavati laṅaḥ na svābhāvikam etat . \\
\noindent \textbf{(P\_3,2.124.1)} atha vā ādeśe sāmānādhikaraṇyam dṛṣṭvā anumānāt gantavyam prakṛteḥ api sāmānādhikaraṇyam bhavati iti . \\
\noindent \textbf{(P\_3,2.124.1)} tat yathā dhūmam dṛṣṭvā agniḥ atra iti gamyate triviṣṭabdhakam dṛṣṭvā parivrājakaḥ iti . \\
\noindent \textbf{(P\_3,2.124.1)} viṣamaḥ upanyāsaḥ . \\
\noindent \textbf{(P\_3,2.124.1)} pratyakṣaḥ tena agnidhūmayoḥ abhisambandhaḥ kṛtaḥ bhavati triviṣṭabdhakaparivrājakayoḥ ca . \\
\noindent \textbf{(P\_3,2.124.1)} saḥ tadvideśastham api dṛṣṭvā jānāti agniḥ atra parivrājakaḥ atra iti . \\
\noindent \textbf{(P\_3,2.124.1)} bhavati vai pratyakṣāt api anumānabalīyastvam . \\
\noindent \textbf{(P\_3,2.124.1)} tat yathā alātacakram pratyakṣam dṛśyate anumānāt ca gamyate na etat asti iti . \\
\noindent \textbf{(P\_3,2.124.1)} kasya cit khalu api sakṛt kṛtaḥ abhisambandhaḥ atyantāya kṛtaḥ bhavati . \\
\noindent \textbf{(P\_3,2.124.1)} tat yathā vṛkṣaparṇayoḥ ayam vṛkṣaḥ idam parṇam iti . \\
\noindent \textbf{(P\_3,2.124.1)} saḥ tadvideśastham api dṛṣṭvā jānāti vṛkṣasya idam parṇam iti . \\
\noindent \textbf{(P\_3,2.124.2)} kim punaḥ ayam paryudāsaḥ : yat anyat prathamāsamānādhikaraṇāt iti , āhosvit prasajyapratiṣedhaḥ : prathamāsamānādhikaraṇe na iti . \\
\noindent \textbf{(P\_3,2.124.2)} kaḥ ca atra viśeṣaḥ . \\
\noindent \textbf{(P\_3,2.124.2)} laṭaḥ śatṛśānacau aprathamāsamānādhikaraṇe iti cet pratyayottarapadayoḥ upasaṅkhyānam . \\
\noindent \textbf{(P\_3,2.124.2)} laṭaḥ śatṛśānacau aprathamāsamānādhikaraṇe iti cet pratyayottarapadayoḥ upasaṅkhyānam kartavyam . \\
\noindent \textbf{(P\_3,2.124.2)} kaurvataḥ pācataḥ kurvadbhaktiḥ pacadbhaktiḥ kurvāṇabhaktiḥ pacamānabhaktiḥ . \\
\noindent \textbf{(P\_3,2.124.2)} astu tarhi prasajyapratiṣedhaḥ .prathamāsamānādhikaraṇe na iti . \\
\noindent \textbf{(P\_3,2.124.2)} prasajyapratiṣedheuttarapade ādeśānupapattiḥ . \\
\noindent \textbf{(P\_3,2.124.2)} prasajyapratiṣedheuttarapade ādeśayoḥ anupapattiḥ . \\
\noindent \textbf{(P\_3,2.124.2)} kurvatī ca asau bhaktiḥ ca kurvadbhaktiḥ pacabhaktiḥ kurvāṇabhaktiḥ pacamānabhaktiḥ iti . \\
\noindent \textbf{(P\_3,2.124.2)} ye ca api ete samānādhikaraṇavṛttayaḥ taddhitāḥ tatra ca śatṛśānacau na prāpnutaḥ . \\
\noindent \textbf{(P\_3,2.124.2)} kurvattaraḥ pacattaraḥ kurvāṇataraḥ pacamānataraḥ kurvadrūpaḥ pacadrūpaḥ kurvāṇarūpaḥ pacamānarūpaḥ kurvatkalpaḥ pacatkalpaḥ kurvāṇakalpaḥ pacamānakalpaḥ iti . \\
\noindent \textbf{(P\_3,2.124.2)} siddham tu pratyayottarapadayoḥ ca iti vacanāt . \\
\noindent \textbf{(P\_3,2.124.2)} siddham etat . \\
\noindent \textbf{(P\_3,2.124.2)} katham . \\
\noindent \textbf{(P\_3,2.124.2)} pratyayottarapadayoḥ ca śatṛśānacau bhavataḥ iti vaktavyam . \\
\noindent \textbf{(P\_3,2.124.2)} tatra pratyayasya ādeśanimittatvāt aprasiddhiḥ . \\
\noindent \textbf{(P\_3,2.124.2)} tatra pratyayasya ādeśanimittatvāt aprasiddhiḥ . \\
\noindent \textbf{(P\_3,2.124.2)} ādeśanimittaḥ pratyayaḥ pratyayanimittaḥ ca ādeśaḥ . \\
\noindent \textbf{(P\_3,2.124.2)} tat etat itaretarāśrayam bhavati . \\
\noindent \textbf{(P\_3,2.124.2)} itaretarāśrayāṇi ca na praklpante . \\
\noindent \textbf{(P\_3,2.124.2)} uttarapadasya ca subantanimittatvāt śatṛśānacoḥ aprasiddhiḥ . \\
\noindent \textbf{(P\_3,2.124.2)} uttarapadasya ca subantanimittatvāt śatṛśānacoḥ aprasiddhiḥ . \\
\noindent \textbf{(P\_3,2.124.2)} uttarapadanimittaḥ sup subantanimittam ca uttarapadam . \\
\noindent \textbf{(P\_3,2.124.2)} tat etat itaretarāśrayam bhavati . \\
\noindent \textbf{(P\_3,2.124.2)} itaretarāśrayāṇi ca na praklpante . \\
\noindent \textbf{(P\_3,2.124.2)} na vā lakārasya kṛttvāt prātipadikatvam tadāśrayam pratyayavidhānam . \\
\noindent \textbf{(P\_3,2.124.2)} na vā eṣaḥ doṣaḥ . \\
\noindent \textbf{(P\_3,2.124.2)} kim kāraṇam . \\
\noindent \textbf{(P\_3,2.124.2)} lakārasya kṛttvāt prātipadikatvam . \\
\noindent \textbf{(P\_3,2.124.2)} lakāraḥ kṛt kṛt prātipadikam iti prātipadikasañjñā . \\
\noindent \textbf{(P\_3,2.124.2)} tadāśrayam pratyayavidhānam . \\
\noindent \textbf{(P\_3,2.124.2)} prātipadikāśrayā svādyutpattiḥ bhaviṣyati . \\
\noindent \textbf{(P\_3,2.124.2)} tiṅādeśāt subutpattiḥ . tiṅādeśaḥ kriyatām subutpattiḥ iti paratvāt subutpattiḥ bhaviṣyati . \\
\noindent \textbf{(P\_3,2.124.2)} tasmāt uttarapadaprasiddhiḥ . \\
\noindent \textbf{(P\_3,2.124.2)} tasmāt uttarapadam prasiddham . \\
\noindent \textbf{(P\_3,2.124.2)} uttarapade prasiddhe uttarapade iti śatṛśānacau bhaviṣyataḥ . \\
\noindent \textbf{(P\_3,2.124.2)} iha api tarhi tiṅādeśāt subutpattiḥ syāt pacati paṭhati iti . \\
\noindent \textbf{(P\_3,2.124.2)} asti atra viśeṣaḥ . \\
\noindent \textbf{(P\_3,2.124.2)} nityaḥ atra tiṅādeśaḥ . \\
\noindent \textbf{(P\_3,2.124.2)} utpanne api supi prāpnoti anutpanne api prāpnoti . \\
\noindent \textbf{(P\_3,2.124.2)} nityatvāt tiṅādeśe kṛte subutpattiḥ na bhaviṣyati . \\
\noindent \textbf{(P\_3,2.124.2)} iha api tarhi nityatvāt tiṅādeśaḥ syāt kurvadbhaktiḥ pacadbhaktiḥ pacamānabhaktiḥ iti . \\
\noindent \textbf{(P\_3,2.124.2)} asti atra viśeṣaḥ . \\
\noindent \textbf{(P\_3,2.124.2)} śatṛśānacau tiṅapavādau tau ca nimittavantau . \\
\noindent \textbf{(P\_3,2.124.2)} na ca apavādaviṣaye utsargaḥ abhiniviśate . \\
\noindent \textbf{(P\_3,2.124.2)} pūrvam hi apavādāḥ abhiniviśante paścāt utsargaḥ . \\
\noindent \textbf{(P\_3,2.124.2)} prakalpya vā apavādaviṣayam tataḥ utsargaḥ abhiniviśate . \\
\noindent \textbf{(P\_3,2.124.2)} na tāvat atra kadā cit tiṅ bhavati . \\
\noindent \textbf{(P\_3,2.124.2)} apavādau śatṛśānacau pratīkṣate . \\
\noindent \textbf{(P\_3,2.124.2)} tat etat kva siddham bhavati . \\
\noindent \textbf{(P\_3,2.124.2)} yatra sāmānyāt utpattiḥ . \\
\noindent \textbf{(P\_3,2.124.2)} yatra hi viśeṣāt ataḥ iñ iti itaretarāśrayam eva tatra bhavati . \\
\noindent \textbf{(P\_3,2.124.2)} vīkṣamāṇasya apatyam vaikṣamāṇiḥ iti . \\
\noindent \textbf{(P\_3,2.124.2)} iha ca śatṛśānacau prāpnutaḥ . \\
\noindent \textbf{(P\_3,2.124.2)} pacatitarām jalpatitarām pacatirūpam jalpatirūpam pacatikalpam jalpatikalpam pacati paṭhati iti . \\
\noindent \textbf{(P\_3,2.124.2)} tat etat katham kṛtvā siddham bhavati . \\
\noindent \textbf{(P\_3,2.124.2)} śatṛśānacau yadi laṭaḥ vā . \\
\noindent \textbf{(P\_3,2.124.2)} yadi śatṛśānacau yadi laṭaḥ vā bhavataḥ vyavasthitavibhāṣā ca . \\
\noindent \textbf{(P\_3,2.124.2)} tena iha ca bhaviṣyataḥ kaurvataḥ pācataḥ kurvadbhaktiḥ pacadbhaktiḥ pacamānabhaktiḥ kurvattaraḥ pacattaraḥ pacamānataraḥ kurvadrūpaḥ pacadrūpaḥ pacamānarūpaḥ kurvatkalpaḥ pacatkalpaḥ pacamānakalpaḥ pacan paṭhan iti ca laṭaḥ śatṛśānacau . \\
\noindent \textbf{(P\_3,2.124.2)} iha ca na bhaviṣyataḥ pacatitarām jalpatitarām pacatirūpam jalpatirūpam pacatikalpam jalpatikalpam pacati paṭhati iti ca laṭaḥ śatṛśānacau . \\
\noindent \textbf{(P\_3,2.124.2)} tat tarhi vāvacanam kartavyam . \\
\noindent \textbf{(P\_3,2.124.2)} na kartavyam . \\
\noindent \textbf{(P\_3,2.124.2)} prakṛtam anuvartate . \\
\noindent \textbf{(P\_3,2.124.2)} kva prakṛtam . \\
\noindent \textbf{(P\_3,2.124.2)} nanvoḥ vibhāṣā iti . \\
\noindent \textbf{(P\_3,2.124.2)} yadi tat anuvartate vartamāne laṭa iti laṭ api vibhāṣā prāpnoti . \\
\noindent \textbf{(P\_3,2.124.2)} sambandham anuvartiṣyate . \\
\noindent \textbf{(P\_3,2.124.2)} nanvoḥ vibhāṣā . \\
\noindent \textbf{(P\_3,2.124.2)} puri luṅ ca asme vibhāṣā . \\
\noindent \textbf{(P\_3,2.124.2)} vartamāne laṭ puri luṅ ca asme vibhāṣā . \\
\noindent \textbf{(P\_3,2.124.2)} laṭaḥ śatṛśānacau vibhāṣā . \\
\noindent \textbf{(P\_3,2.124.2)} puri luṅ ca asme iti nivṛttam . \\
\noindent \textbf{(P\_3,2.124.2)} na tarhi idānīm aprathamāsamānādhikaraṇe iti vaktavyam . \\
\noindent \textbf{(P\_3,2.124.2)} vaktavyam ca . \\
\noindent \textbf{(P\_3,2.124.2)} kim prayojanam . \\
\noindent \textbf{(P\_3,2.124.2)} nityāṛtham . \\
\noindent \textbf{(P\_3,2.124.2)} aprathamāsamānādhikaraṇe nityau yathā syātām . \\
\noindent \textbf{(P\_3,2.124.2)} kva tarhi idānīm vibhāṣā . \\
\noindent \textbf{(P\_3,2.124.2)} prathamāsamānādhikaraṇe . \\
\noindent \textbf{(P\_3,2.124.2)} pacan pacati pacamānaḥ pacate iti . \\
\noindent \textbf{(P\_3,2.126)} lakṣaṇahetvoḥ kriyāyāḥ guṇe upasaṅkhyānam . \\
\noindent \textbf{(P\_3,2.126)} lakṣaṇahetvoḥ kriyāyāḥ guṇe upasaṅkhyānam kartavyam . \\
\noindent \textbf{(P\_3,2.126)} tiṣṭhan mūtrayati . \\
\noindent \textbf{(P\_3,2.126)} gacchan bhakṣayati . \\
\noindent \textbf{(P\_3,2.126)} kartuḥ ca lakṣaṇayoḥ paryāyeṇa acayoge . \\
\noindent \textbf{(P\_3,2.126)} kartuḥ ca lakṣaṇayoḥ paryāyeṇa acayoge upasaṅkhyānam kartavyam . \\
\noindent \textbf{(P\_3,2.126)} yaḥ adhīyānaḥ āste saḥ devadatttaḥ . \\
\noindent \textbf{(P\_3,2.126)} yaḥ āsīnaḥ adhīte saḥ devadattaḥ . \\
\noindent \textbf{(P\_3,2.126)} acayoge iti kimartham . \\
\noindent \textbf{(P\_3,2.126)} yaḥ āste ca adhīte ca saḥ caitraḥ . \\
\noindent \textbf{(P\_3,2.126)} tattvānvākhyāne ca . \\
\noindent \textbf{(P\_3,2.126)} tattvānvākhyāne ca upasaṅkhyānam kartavyam . \\
\noindent \textbf{(P\_3,2.126)} śayānā vardhate dūrvā . \\
\noindent \textbf{(P\_3,2.126)} āsīnam vardhate bisam iti . \\
\noindent \textbf{(P\_3,2.126)} sadādayaḥ ca bahulam . \\
\noindent \textbf{(P\_3,2.126)} sadādayaḥ ca bahulam iti vaktavyam . \\
\noindent \textbf{(P\_3,2.126)} san brāhmaṇaḥ asti brāhmaṇaḥ . \\
\noindent \textbf{(P\_3,2.126)} vidyate brāhmaṇaḥ vidyamānaḥ brāhmaṇaḥ iti . \\
\noindent \textbf{(P\_3,2.126)} iṅjuhotyoḥ vāvacanam . \\
\noindent \textbf{(P\_3,2.126)} iṅjuhotyoḥ vā iti vaktavyam . \\
\noindent \textbf{(P\_3,2.126)} adhīte adhīyānaḥ . \\
\noindent \textbf{(P\_3,2.126)} juhoti juhvat . \\
\noindent \textbf{(P\_3,2.126)} māṅi ākrośe . māṅi ākrośe iti vaktavyam . \\
\noindent \textbf{(P\_3,2.126)} mā pacan . \\
\noindent \textbf{(P\_3,2.126)} mā pacamānaḥ . \\
\noindent \textbf{(P\_3,2.126)} tat tarhi vaktavyam . \\
\noindent \textbf{(P\_3,2.126)} na vaktavyam . \\
\noindent \textbf{(P\_3,2.126)} lakṣaṇahetvoḥ kriyāyāḥ iti eva siddham . \\
\noindent \textbf{(P\_3,2.126)} iha tāvat tiṣṭhan mūtrayati iti . \\
\noindent \textbf{(P\_3,2.126)} tiṣṭhatikriyā mūtrayatikriyāyāḥ lakṣaṇam . \\
\noindent \textbf{(P\_3,2.126)} gacchan bhakṣayati iti . \\
\noindent \textbf{(P\_3,2.126)} gacchatikriyā bhakṣayatikriyāyāḥ lakṣaṇam . \\
\noindent \textbf{(P\_3,2.126)} yaḥ adhīyānaḥ āste saḥ devadatttaḥ iti . \\
\noindent \textbf{(P\_3,2.126)} adhyayanakriyā āsanakriyāyāḥ lakṣaṇam . \\
\noindent \textbf{(P\_3,2.126)} yaḥ āsīnaḥ adhīte saḥ devadattaḥ iti . \\
\noindent \textbf{(P\_3,2.126)} āsikriyā adhyayanakriyāyāḥ lakṣaṇam . \\
\noindent \textbf{(P\_3,2.126)} idam tarhi prayojanam . \\
\noindent \textbf{(P\_3,2.126)} acayoge iti vakṣyāmi iti . \\
\noindent \textbf{(P\_3,2.126)} iha mā bhūt . \\
\noindent \textbf{(P\_3,2.126)} yaḥ āste ca adhīte ca saḥ caitraḥ iti . \\
\noindent \textbf{(P\_3,2.126)} etat api na asti prayojanam . \\
\noindent \textbf{(P\_3,2.126)} na etat kriyāyāḥ lakṣaṇam . \\
\noindent \textbf{(P\_3,2.126)} kim tarhi . \\
\noindent \textbf{(P\_3,2.126)} kartṛlakṣaṇam etat . \\
\noindent \textbf{(P\_3,2.126)} śayānā vardhate dūrvā iti . \\
\noindent \textbf{(P\_3,2.126)} śetikriyā vṛddhikriyāyāḥ lakṣaṇam . \\
\noindent \textbf{(P\_3,2.126)} āsīnam vardhate bisam iti . \\
\noindent \textbf{(P\_3,2.126)} āsikriyā vṛddhikriyāyāḥ lakṣaṇam . \\
\noindent \textbf{(P\_3,2.126)} sadādayaḥ ca bahulam iṅjuhotyoḥ vā māṅi ākrośe iti vaktavyam eva. \\
\noindent \textbf{(P\_3,2.127.1)} taugrahaṇam kimartham . \\
\noindent \textbf{(P\_3,2.127.1)} śatṛśānacau pratinirdiśyete . \\
\noindent \textbf{(P\_3,2.127.1)} na etat asti prayojanam . \\
\noindent \textbf{(P\_3,2.127.1)} prakṛtau śatṛśānacau anuvartiṣyete . \\
\noindent \textbf{(P\_3,2.127.1)} kva prakṛtau . \\
\noindent \textbf{(P\_3,2.127.1)} laṭaḥ śatṛsānacau aprathamāsamānādhikaraṇe iti . \\
\noindent \textbf{(P\_3,2.127.1)} ataḥ uttaram paṭhati . \\
\noindent \textbf{(P\_3,2.127.1)} tau sat iti vacanam asaṃsargārtham ṭau grahaṇam kriyate asaṃsargārtham . \\
\noindent \textbf{(P\_3,2.127.1)} asaṃsaktayoḥ etaiḥ viśeṣaiḥ śatṛśānacoḥ sañjñā yathā syāt . \\
\noindent \textbf{(P\_3,2.127.1)} nanu ca ete viśeṣāḥ nivarteran . \\
\noindent \textbf{(P\_3,2.127.1)} yadi api ete viśeṣāḥ nivartante ayam tu khalu vartamānaḥ kālaḥ avaśyam uttarārthaḥ anuvartyaḥ . \\
\noindent \textbf{(P\_3,2.127.1)} tasmin anuvartamāne vartamānakālavihitayoḥ eva śatṛśānacoḥ satsñjñā syāt . \\
\noindent \textbf{(P\_3,2.127.1)} bhūtabhaviṣyakālavihitayoḥ na syāt . \\
\noindent \textbf{(P\_3,2.127.1)} kim punaḥ bhūtabhaviṣyakālavihitayoḥ satsñjñāvacane prayojanam . \\
\noindent \textbf{(P\_3,2.127.1)} pūraṇaguṇasuhitasat iti . \\
\noindent \textbf{(P\_3,2.127.1)} brāhmaṇasya pakṣyan . \\
\noindent \textbf{(P\_3,2.127.1)} brāhmaṇasya pakṣyamāṇaḥ . \\
\noindent \textbf{(P\_3,2.127.1)} atha kriyamāṇe api taugrahaṇe katham eva asaṃsaktayoḥ etaiḥ viśeṣaiḥ sañjñā labhyā . \\
\noindent \textbf{(P\_3,2.127.1)} labhyā iti āha . \\
\noindent \textbf{(P\_3,2.127.1)} katham . \\
\noindent \textbf{(P\_3,2.127.1)} tau iti śabdataḥ sat iti yoge kriyamāṇe tau grahaṇam yogāṅgam jāyate . \\
\noindent \textbf{(P\_3,2.127.1)} sati ca yogāṅge yogavibhāgaḥ kariṣyate . \\
\noindent \textbf{(P\_3,2.127.1)} tau . \\
\noindent \textbf{(P\_3,2.127.1)} tau etau śatṛśānacau dhātumātrāt parasya pratyayasya bhavataḥ . \\
\noindent \textbf{(P\_3,2.127.1)} tataḥ sat . \\
\noindent \textbf{(P\_3,2.127.1)} satsñjñau bhavatau śatṛśānacau . \\
\noindent \textbf{(P\_3,2.127.1)} iha api tarhi prāpnutaḥ . \\
\noindent \textbf{(P\_3,2.127.1)} kārakaḥ hārakaḥ iti . \\
\noindent \textbf{(P\_3,2.127.1)} avadhāraṇam lṛṭi vidhānam . \\
\noindent \textbf{(P\_3,2.127.1)} lṛtaḥ sat vā iti etat niyamārtham bhaviṣyati . \\
\noindent \textbf{(P\_3,2.127.1)} lṛṭaḥ eva dhātumātrāt parasya na anyasya iti . \\
\noindent \textbf{(P\_3,2.127.1)} kaimarthyakyāt niyamaḥ bhavati . \\
\noindent \textbf{(P\_3,2.127.1)} vidheyam na asti iti kṛtvā . \\
\noindent \textbf{(P\_3,2.127.1)} iha ca asti vidheyam . \\
\noindent \textbf{(P\_3,2.127.1)} kim . \\
\noindent \textbf{(P\_3,2.127.1)} nityau śatṛśānacau prāptau . \\
\noindent \textbf{(P\_3,2.127.1)} tau vibhāṣā vidheyau . \\
\noindent \textbf{(P\_3,2.127.1)} tatra apūrvaḥ vidhiḥ astu niyamaḥ astu iti apūrvaḥ eva vidhiḥ bhaviṣyati na niyamaḥ . \\
\noindent \textbf{(P\_3,2.127.1)} yogavibhāgataḥ ca vihitam sat . \\
\noindent \textbf{(P\_3,2.127.1)} evam tarhi yogavibhāgaḥ kariṣyate . \\
\noindent \textbf{(P\_3,2.127.1)} lṛṭaḥ sat . \\
\noindent \textbf{(P\_3,2.127.1)} lṛṭaḥ satsañjñau bhavataḥ . \\
\noindent \textbf{(P\_3,2.127.1)} kimartham idam . \\
\noindent \textbf{(P\_3,2.127.1)} niyamārtham . \\
\noindent \textbf{(P\_3,2.127.1)} lṛṭaḥ eva dhātumātrāt parasya na anyasya iti . \\
\noindent \textbf{(P\_3,2.127.1)} tataḥ vā . \\
\noindent \textbf{(P\_3,2.127.1)} vā ca lṛṭaḥ śatṛśānacau satsañjñau bhavataḥ . \\
\noindent \textbf{(P\_3,2.127.1)} tatra ayam api arthaḥ . \\
\noindent \textbf{(P\_3,2.127.1)} sadvidhiḥ nityam aprathamāsamānādhikaraṇe iti vakṣyati . \\
\noindent \textbf{(P\_3,2.127.1)} tat na vaktavyam bhavati . \\
\noindent \textbf{(P\_3,2.127.2)} atha yau etau uttarau śatānau kim etau lādeśau āhosvit alādeśau . \\
\noindent \textbf{(P\_3,2.127.2)} kaḥ ca atra viśeṣaḥ . \\
\noindent \textbf{(P\_3,2.127.2)} uttarayoḥ lādeśe vāvacanam . \\
\noindent \textbf{(P\_3,2.127.2)} uttarayoḥ lādeśe vā iti vaktavyam . \\
\noindent \textbf{(P\_3,2.127.2)} pavamānaḥ yajamānaḥ . \\
\noindent \textbf{(P\_3,2.127.2)} pavate yajate iti api yathā syāt . \\
\noindent \textbf{(P\_3,2.127.2)} sādhanābhidhānam . \\
\noindent \textbf{(P\_3,2.127.2)} sādhanābhidhānam ca prāpnoti . \\
\noindent \textbf{(P\_3,2.127.2)} laḥ karmaṇi ca bhāve ca akarmakebhyaḥ iti bhāvakarmaṇoḥ api prāpnutaḥ . \\
\noindent \textbf{(P\_3,2.127.2)} svaraḥ . \\
\noindent \textbf{(P\_3,2.127.2)} svaraḥ ca sādhyaḥ . \\
\noindent \textbf{(P\_3,2.127.2)} kati iha pavamānāḥ . \\
\noindent \textbf{(P\_3,2.127.2)} adupadeśāt lasārvadhātukam anudāttam bhavati iti eṣaḥ svaraḥ prāpnoti . \\
\noindent \textbf{(P\_3,2.127.2)} upagrahapratiṣedhaḥ ca . \\
\noindent \textbf{(P\_3,2.127.2)} upagrahasya ca pratiṣedhaḥ vaktavyaḥ . \\
\noindent \textbf{(P\_3,2.127.2)} kati iha nighnānāḥ . \\
\noindent \textbf{(P\_3,2.127.2)} taṅānau ātmanepadam iti ātmanepadasañjñā prāpnoti . \\
\noindent \textbf{(P\_3,2.127.2)} stām tarhi alādeśau . \\
\noindent \textbf{(P\_3,2.127.2)} alādeśe ṣaṣṭhīpratiṣedhaḥ . \\
\noindent \textbf{(P\_3,2.127.2)} alādeśe ṣaṣṭhīpratiṣedhaḥ vaktavyaḥ . \\
\noindent \textbf{(P\_3,2.127.2)} somam pavamānaḥ . \\
\noindent \textbf{(P\_3,2.127.2)} naḍam āghnānaḥ . \\
\noindent \textbf{(P\_3,2.127.2)} adhīyan pārāyaṇam . \\
\noindent \textbf{(P\_3,2.127.2)} laprayoge na iti pratiṣedhaḥ na prāpnoti . \\
\noindent \textbf{(P\_3,2.127.2)} mā bhūt evam . \\
\noindent \textbf{(P\_3,2.127.2)} tṛn iti eva bhaviṣyati . \\
\noindent \textbf{(P\_3,2.127.2)} katham . \\
\noindent \textbf{(P\_3,2.127.2)} tṛn iti na idam pratyayagrahaṇam . \\
\noindent \textbf{(P\_3,2.127.2)} kim tarhi . \\
\noindent \textbf{(P\_3,2.127.2)} pratyāhāragrahaṇam . \\
\noindent \textbf{(P\_3,2.127.2)} kva sanniviṣṭāṇām pratyāhāraḥ . \\
\noindent \textbf{(P\_3,2.127.2)} laṭaḥ śatṛ iti ataḥ ārabhya ā tṛnaḥ nakārāt . \\
\noindent \textbf{(P\_3,2.127.2)} yadi pratyāhāragrahaṇam caurasya dviṣan vṛṣalasya dviṣan atra api prāpnoti . \\
\noindent \textbf{(P\_3,2.127.2)} dviṣaḥ śatuḥ vāvacanam . \\
\noindent \textbf{(P\_3,2.127.2)} dviṣaḥ śatuḥ vā iti vaktavyam . \\
\noindent \textbf{(P\_3,2.127.2)} tat ca avaśyam vaktavyam pratyayagrahaṇe sati pratiṣedhārtham . \\
\noindent \textbf{(P\_3,2.127.2)} tat eva pratyāhāragrahaṇe sati vidhyartham bhaviṣyati. \\
\noindent \textbf{(P\_3,2.135)} tṛnvidhau ṛtvikṣu ca anupasargasya . \\
\noindent \textbf{(P\_3,2.135)} tṛnvidhau ṛtvikṣu ca anupasargasya iti vaktavyam . \\
\noindent \textbf{(P\_3,2.135)} hotā potā . \\
\noindent \textbf{(P\_3,2.135)} anupasargasya iti kimartham . \\
\noindent \textbf{(P\_3,2.135)} praśāstā pratihartā . \\
\noindent \textbf{(P\_3,2.135)} nayateḥ ṣuk ca . \\
\noindent \textbf{(P\_3,2.135)} nayateḥ ṣuk vaktavyaḥ . \\
\noindent \textbf{(P\_3,2.135)} tṛn ca pratyayaḥ vaktavyañ. neṣṭā . \\
\noindent \textbf{(P\_3,2.135)} na vā dhātvanyatvāt . na vā vaktavyam . \\
\noindent \textbf{(P\_3,2.135)} kim kāraṇam . \\
\noindent \textbf{(P\_3,2.135)} dhātvanyatvāt . \\
\noindent \textbf{(P\_3,2.135)} dhātvantaram neṣatiḥ . \\
\noindent \textbf{(P\_3,2.135)} katham jñāyate . \\
\noindent \textbf{(P\_3,2.135)} neṣatu neṣṭāt iti darśanāt . \\
\noindent \textbf{(P\_3,2.135)} neṣatu neṣṭāt iti prayogaḥ dṛśyate . \\
\noindent \textbf{(P\_3,2.135)} indraḥ naḥ tena neṣatu . \\
\noindent \textbf{(P\_3,2.135)} gāḥ vaḥ neṣṭāt . \\
\noindent \textbf{(P\_3,2.135)} tviṣeḥ devatāyām akāraḥ ca upadhāyāḥ aniṭvam ca . \\
\noindent \textbf{(P\_3,2.135)} tviṣeḥ devatāyām tṛn vaktavyaḥ akāraḥ ca upadhāyāḥ aniṭvam ca iti . \\
\noindent \textbf{(P\_3,2.135)} tvaṣṭā . \\
\noindent \textbf{(P\_3,2.135)} kim punaḥ idam tviṣeḥ eva aniṭvam . \\
\noindent \textbf{(P\_3,2.135)} na iti āha . \\
\noindent \textbf{(P\_3,2.135)} yat ca anukrāntam yat ca anukraṃsyate sarvasya eṣaḥ śeṣaḥ aniṭvam iti . \\
\noindent \textbf{(P\_3,2.135)} kṣadeḥ ca yukte . \\
\noindent \textbf{(P\_3,2.135)} kṣadeḥ ca yukte tṛn vaktavyaḥ . \\
\noindent \textbf{(P\_3,2.135)} kṣattā . \\
\noindent \textbf{(P\_3,2.135)} chandasi tṛc ca . \\
\noindent \textbf{(P\_3,2.135)} chandasi tṛc ca tṛn ca vaktavyaḥ . \\
\noindent \textbf{(P\_3,2.135)} kṣattṛbhyaḥ saṅgrahītṛbhyaḥ . \\
\noindent \textbf{(P\_3,2.135)} kṣattṛbhyaḥ saṅgrahītṛbhyaḥ . \\
\noindent \textbf{(P\_3,2.139)} snoḥ kittve sthaḥ īkārpratiṣedhaḥ . \\
\noindent \textbf{(P\_3,2.139)} snoḥ kittve sthaḥ īkārpratiṣedhaḥ vaktavyaḥ . \\
\noindent \textbf{(P\_3,2.139)} sthāsnuḥ iti . \\
\noindent \textbf{(P\_3,2.139)} ghumāsthāgāpājahātisām hali iti īttvam prāpnoti . \\
\noindent \textbf{(P\_3,2.139)} evam tarhi na kit kariṣyate . \\
\noindent \textbf{(P\_3,2.139)} akiti guṇapratiṣedhaḥ . \\
\noindent \textbf{(P\_3,2.139)} yadi akit guṇapratiṣedhaḥ vaktavyaḥ . \\
\noindent \textbf{(P\_3,2.139)} jiṣṇuḥ iti . \\
\noindent \textbf{(P\_3,2.139)} bhuvaḥ iṭpratiṣedhaḥ ca . \\
\noindent \textbf{(P\_3,2.139)} bhuvaḥ iṭpratiṣedhaḥ ca vaktavyaḥ . \\
\noindent \textbf{(P\_3,2.139)} kim ca anyat . \\
\noindent \textbf{(P\_3,2.139)} guṇapratiṣedhaḥ ca . \\
\noindent \textbf{(P\_3,2.139)} bhūṣṇuḥ iti . \\
\noindent \textbf{(P\_3,2.139)} astu tarhi kit. nanu ca uktam snoḥ kittve sthaḥ īkārpratiṣedhaḥ vaktavyaḥ iti . \\
\noindent \textbf{(P\_3,2.139)} na eṣaḥ doṣaḥ . \\
\noindent \textbf{(P\_3,2.139)} sthādaṃśibhyām snuḥ chandasi . \\
\noindent \textbf{(P\_3,2.139)} sthādaṃśibhyām snuḥ chandasi vaktavyaḥ . \\
\noindent \textbf{(P\_3,2.139)} sthāsnu jaṅgamam . \\
\noindent \textbf{(P\_3,2.139)} daṅkṣṇavaḥ paśavaḥ iti . \\
\noindent \textbf{(P\_3,2.139)} saḥ idānīm sthaḥ aviśeṣeṇa vidhāsyate . \\
\noindent \textbf{(P\_3,2.139)} sūtram tarhi bhidyate . \\
\noindent \textbf{(P\_3,2.139)} yathānyāsam astu. nanu ca uktam snoḥ kittve sthaḥ īkārpratiṣedhaḥ vaktavyaḥ iti . \\
\noindent \textbf{(P\_3,2.139)} evam tarhi git kariṣyate . \\
\noindent \textbf{(P\_3,2.139)} sthoḥ gittvāt na sthaḥ īkāraḥ . \\
\noindent \textbf{(P\_3,2.139)} sthoḥ gittvāt sthaḥ īkāraḥ na bhaviṣyati . \\
\noindent \textbf{(P\_3,2.139)} kim kāraṇam . \\
\noindent \textbf{(P\_3,2.139)} kṅitoḥ īttvaśāsanāt . \\
\noindent \textbf{(P\_3,2.139)} kṅitoḥ īttvam śiṣyate . \\
\noindent \textbf{(P\_3,2.139)} iha tarhi jiṣṇuḥ iti guṇaḥ prāpnoti . \\
\noindent \textbf{(P\_3,2.139)} guṇābhāvaḥ triṣu smāryaḥ . \\
\noindent \textbf{(P\_3,2.139)} guṇābhāvaḥ triṣu smartavyaḥ . \\
\noindent \textbf{(P\_3,2.139)} giti kiti ṅiti iti . \\
\noindent \textbf{(P\_3,2.139)} tat gakāragrahaṇam kartavyam . \\
\noindent \textbf{(P\_3,2.139)} na kartavyam . \\
\noindent \textbf{(P\_3,2.139)} kriyate nyāse eva . \\
\noindent \textbf{(P\_3,2.139)} kakāre gakāraḥ cartvabhūtaḥ nirdiśyate . \\
\noindent \textbf{(P\_3,2.139)} kkṅiti ca iti . \\
\noindent \textbf{(P\_3,2.139)} iha tarhi bhūṣṇuḥ iti śryukaḥ kiti iti iṭpratiṣedhaḥ na prāpnoti . \\
\noindent \textbf{(P\_3,2.139)} śryuko aniṭtvam gakoḥ itoḥ . \\
\noindent \textbf{(P\_3,2.139)} śryukaḥ aniṭtvam gakārakakārayoḥ iti vaktavyam . \\
\noindent \textbf{(P\_3,2.139)} tat gakāragrahaṇam kartavyam . \\
\noindent \textbf{(P\_3,2.139)} na kartavyam . \\
\noindent \textbf{(P\_3,2.139)} kriyate nyāsaḥ eva . \\
\noindent \textbf{(P\_3,2.139)} kakāre gakāraḥ cartvabhūtaḥ nirdiśyate . \\
\noindent \textbf{(P\_3,2.139)} śryukaḥ kkiti iti . \\
\noindent \textbf{(P\_3,2.139)} yadi evam cartvasya asiddhatvāt haśi iti uttvam prāpnoti . \\
\noindent \textbf{(P\_3,2.139)} sautraḥ nirdeśaḥ . \\
\noindent \textbf{(P\_3,2.139)} atha vā asaṃhitayā nirdeśaḥ kariṣyate . \\
\noindent \textbf{(P\_3,2.139)} śryukaḥ kkiti iti . \\
\noindent \textbf{(P\_3,2.139)} sthoḥ gittvāt na sthaḥ īkāraḥ kṅitoḥ īttvaśāsanāt . \\
\noindent \textbf{(P\_3,2.139)} guṇābhāvaḥ triṣu smāryaḥ . \\
\noindent \textbf{(P\_3,2.139)} śryuko aniṭtvam gakoḥ itoḥ . \\
\noindent \textbf{(P\_3,2.141)} ghinaṇ ayam vaktavyaḥ . \\
\noindent \textbf{(P\_3,2.141)} ghinuṇi hi sati śaminau śaminaḥ taminau taminaḥ ugidacām sarvanāmasthāne adhātoḥ iti num prasajyeta . \\
\noindent \textbf{(P\_3,2.141)} na eṣaḥ doṣaḥ . \\
\noindent \textbf{(P\_3,2.141)} jhalgrahaṇam tatra codayiṣyati . \\
\noindent \textbf{(P\_3,2.141)} iha tarhi : śaminitarā śaminītarā taminitarā taminītarā : ugitaḥ ghādiṣu nadyāḥ anyatarasyām hrasvaḥ bhavati iti anyaratasyām hrasvatvam prasajyeta . \\
\noindent \textbf{(P\_3,2.141)} iṣyate eva hrasvatvam . \\
\noindent \textbf{(P\_3,2.141)} ghinuṇ akarmakāṇām . \\
\noindent \textbf{(P\_3,2.141)} ghinuṇ akarmakāṇām iti vaktavyam . \\
\noindent \textbf{(P\_3,2.141)} iha mā bhūt . \\
\noindent \textbf{(P\_3,2.141)} sampṛṇakti śākam iti . \\
\noindent \textbf{(P\_3,2.141)} uktam vā . \\
\noindent \textbf{(P\_3,2.141)} kim uktam . \\
\noindent \textbf{(P\_3,2.141)} anabhidhānāt iti . \\
\noindent \textbf{(P\_3,2.146)} kimartham nindādibhyaḥ vuñ vidhīyate na ṇvulā eva siddham . \\
\noindent \textbf{(P\_3,2.146)} na hi asti viśeṣaḥ nindādibhyaḥ ṇvulaḥ vā vuñaḥ vā . \\
\noindent \textbf{(P\_3,2.146)} tat eva rūpam saḥ eva svaraḥ . \\
\noindent \textbf{(P\_3,2.146)} vuñm anekācaḥ . \\
\noindent \textbf{(P\_3,2.146)} vuñm anekācaḥ prayojayanti . \\
\noindent \textbf{(P\_3,2.146)} asūyakaḥ . \\
\noindent \textbf{(P\_3,2.146)} atha ye atra ekācaḥ paṭhyante teṣām grahaṇam kimartham na teṣām ṇvulā eva siddham . \\
\noindent \textbf{(P\_3,2.146)} na sidhyati . \\
\noindent \textbf{(P\_3,2.146)} ayam tacchīlādiṣu tṛn vidhīyate . \\
\noindent \textbf{(P\_3,2.146)} saḥ viśeṣavihitaḥ sāmānyavihitam ṇvulam bādheta . \\
\noindent \textbf{(P\_3,2.146)} vāsarūpanyāyena ṇvul api bhaviṣyati . \\
\noindent \textbf{(P\_3,2.146)} ataḥ uttaram paṭhati . \\
\noindent \textbf{(P\_3,2.146)} nindādibhyaḥ vuñvacanam anyebhyaḥ ṇvulaḥ pratiṣedhārtham . \\
\noindent \textbf{(P\_3,2.146)} nindādibhyaḥ vuñvacanam kriyate jñāpakārtham . \\
\noindent \textbf{(P\_3,2.146)} kim jñāpyam . \\
\noindent \textbf{(P\_3,2.146)} etat jñāpayati ācāryaḥ tacchīlādiṣu vāsarūpanyāyena anyebhyaḥ ṇvul na bhavati iti . \\
\noindent \textbf{(P\_3,2.146)} tṛjādipratiṣedhāṛtham vā . \\
\noindent \textbf{(P\_3,2.146)} atha vā etat jñāpayati ācāryaḥ . \\
\noindent \textbf{(P\_3,2.146)} tacchīlādiṣu sarve eva tṛjāsayaḥ vāsarūpeṇa na bhavanti iti . \\
\noindent \textbf{(P\_3,2.150)} padigrahaṇam anarthakam anudāttetaḥ ca halādeḥ iti siddhatvāt . \\
\noindent \textbf{(P\_3,2.150)} padigrahaṇam anarthakam . \\
\noindent \textbf{(P\_3,2.150)} kim kāraṇam . \\
\noindent \textbf{(P\_3,2.150)} anudāttetaḥ ca halādeḥ iti eva atra yuc siddhaḥ . \\
\noindent \textbf{(P\_3,2.150)} na sidhyati . \\
\noindent \textbf{(P\_3,2.150)} ayam padeḥ ukañ vidhīyate laṣapadapadsthābhūvṛṣahanakamagamaśṝbhyaḥ ukañ iti . \\
\noindent \textbf{(P\_3,2.150)} saḥ viśeṣavihitaḥ sāmānyavihitam yucam bādheta . \\
\noindent \textbf{(P\_3,2.150)} vāsarūpanyāyena yuc api bhaviṣyati . \\
\noindent \textbf{(P\_3,2.150)} asarūpanivṛttyartham tu . \\
\noindent \textbf{(P\_3,2.150)} asarūpanivṛttyartham tarhi padigrahaṇam kriyate . \\
\noindent \textbf{(P\_3,2.150)} etat jñāpayati ācāryaḥ tācchīlikeṣu tācchilikāḥ vāsarūpeṇa na bhavanti iti . \\
\noindent \textbf{(P\_3,2.150)} yadi etat jñāpyate dūdadīpadīkṣaḥ ca iti dīpagrahaṇam anarthakam . \\
\noindent \textbf{(P\_3,2.150)} ayam dīpeḥ raḥ vidhīyate . \\
\noindent \textbf{(P\_3,2.150)} namikampismyajasahiṃsadīpaḥ raḥ iti . \\
\noindent \textbf{(P\_3,2.150)} saḥ viśeṣavihitaḥ sāmānyavihitam yucam bādhiṣyate . \\
\noindent \textbf{(P\_3,2.150)} evam tarhi siddhe sati yat dīpagrahaṇam karoti tat jñāpayati ācāryaḥ bhavati yucaḥ reṇa samāveśaḥ iti . \\
\noindent \textbf{(P\_3,2.150)} kim etasya jñāpane prayojanam . \\
\noindent \textbf{(P\_3,2.150)} kamrā kanyā kamanā kanyā iti etat siddham bhavati . \\
\noindent \textbf{(P\_3,2.158)} kimartham āluc ucyate na luś eva ucyeta . \\
\noindent \textbf{(P\_3,2.158)} kā rūpasiddhiḥ : spṛhayāluḥ gṛhayāluḥ . \\
\noindent \textbf{(P\_3,2.158)} śapi kṛte ataḥ dīrghaḥ yañi iti dīrghatvam bhaviṣyati . \\
\noindent \textbf{(P\_3,2.158)} evam tarhi siddhe sati yat ālucam śāsti tat jñāpayati ācāryaḥ anyebhyaḥ api ayam bhavati iti . \\
\noindent \textbf{(P\_3,2.158)} kim etasya jñāpane prayojanam . \\
\noindent \textbf{(P\_3,2.158)} āluci śīṅgrahaṇam codayiṣyati . \\
\noindent \textbf{(P\_3,2.158)} tat na kartavyam bhavati . \\
\noindent \textbf{(P\_3,2.158)} āluci śīṅgrahaṇam . \\
\noindent \textbf{(P\_3,2.158)} āluci śīṅgrahaṇam kartavyam . \\
\noindent \textbf{(P\_3,2.158)} śayāluḥ . \\
\noindent \textbf{(P\_3,2.171)} kimartham kikinoḥ kittvam ucyate na asaṃyogāt liṭ kit iti eva siddham . \\
\noindent \textbf{(P\_3,2.171)} kikinoḥ kittvam ṝkāraguṇapratiṣedhāṛtham . \\
\noindent \textbf{(P\_3,2.171)} kikinoḥ kittvam kriyate ṝkāraguṇapratiṣedhāṛtham . \\
\noindent \textbf{(P\_3,2.171)} ayam ṝkārārāntānām liṭi guṇaḥ pratiṣedhaviṣaye ārabhyate ṛcchatyṝtrām iti . \\
\noindent \textbf{(P\_3,2.171)} saḥ yathā iha bhavati ātastaratuḥ atastaruḥ iti evam iha api prasajyeta mitrāvaruṇau taturiḥ dūre hi adhvā jaguriḥ iti . \\
\noindent \textbf{(P\_3,2.171)} saḥ punaḥ kittve bādhyate . \\
\noindent \textbf{(P\_3,2.171)} utsargaḥ chandasi sadādibhyaḥ darśanāt . \\
\noindent \textbf{(P\_3,2.171)} utsargaḥ chandasi kikinau vaktavyau . \\
\noindent \textbf{(P\_3,2.171)} kim prayojanam . \\
\noindent \textbf{(P\_3,2.171)} sadādibhyaḥ darśanāt . \\
\noindent \textbf{(P\_3,2.171)} sadādibhyaḥ hi kikinau dṛśyete . \\
\noindent \textbf{(P\_3,2.171)} sadimaniraminamivicīnām sediḥ meniḥ remiḥ . \\
\noindent \textbf{(P\_3,2.171)} nemiḥ cakram iva abhavat . \\
\noindent \textbf{(P\_3,2.171)} vivicim ratnadhātamam . \\
\noindent \textbf{(P\_3,2.171)} bhāṣāyām dhāñkṛsṛjaninimibhyaḥ . bhāṣāyām dhāñkṛsṛjaninimibhyaḥ kikinau vaktavyau . \\
\noindent \textbf{(P\_3,2.171)} dhāñ . \\
\noindent \textbf{(P\_3,2.171)} dadhiḥ . \\
\noindent \textbf{(P\_3,2.171)} dhāñ . \\
\noindent \textbf{(P\_3,2.171)} kṛ . \\
\noindent \textbf{(P\_3,2.171)} cakriḥ . \\
\noindent \textbf{(P\_3,2.171)} kṛ . \\
\noindent \textbf{(P\_3,2.171)} sṛ . \\
\noindent \textbf{(P\_3,2.171)} sasriḥ . \\
\noindent \textbf{(P\_3,2.171)} sṛ . \\
\noindent \textbf{(P\_3,2.171)} jani . \\
\noindent \textbf{(P\_3,2.171)} jajñiḥ . \\
\noindent \textbf{(P\_3,2.171)} jani . \\
\noindent \textbf{(P\_3,2.171)} nami . \\
\noindent \textbf{(P\_3,2.171)} nemiḥ . \\
\noindent \textbf{(P\_3,2.171)} sāsahivāvahicācalipāpatīnām nipātanam . \\
\noindent \textbf{(P\_3,2.171)} sāsahivāvahicācalipāpatīnām nipātanam kartavyam . \\
\noindent \textbf{(P\_3,2.171)} vṛṣā sahamānam sāsahiḥ . \\
\noindent \textbf{(P\_3,2.171)} vāvahiḥ cācaliḥ pāpatiḥ . \\
\noindent \textbf{(P\_3,2.171)} aparaḥ āha : sahivahicalipatibhyaḥ yaṅantebhyaḥ kikinau vaktavyau . \\
\noindent \textbf{(P\_3,2.171)} etāni eva udāharaṇāni . \\
\noindent \textbf{(P\_3,2.174)} bhiyaḥ krukan api vaktavyaḥ . \\
\noindent \textbf{(P\_3,2.174)} bhīrukaḥ . \\
\noindent \textbf{(P\_3,2.178.1)} kimartham idam ucyate na kvip anyebhyaḥ api dṛśyate iti eva siddham . \\
\noindent \textbf{(P\_3,2.178.1)} kvibvidhiḥ anupapadārthaḥ . \\
\noindent \textbf{(P\_3,2.178.1)} anupapadārthaḥ ayam ārambhaḥ . \\
\noindent \textbf{(P\_3,2.178.1)} paceḥ pak . \\
\noindent \textbf{(P\_3,2.178.1)} bhideḥ bhit . \\
\noindent \textbf{(P\_3,2.178.1)} chideḥ chit . \\
\noindent \textbf{(P\_3,2.178.1)} atha yaḥ atra sopasargaḥ tasya grahaṇam kimartham na tena eva siddham . \\
\noindent \textbf{(P\_3,2.178.1)} na sidhyati . \\
\noindent \textbf{(P\_3,2.178.1)} iha ke cit ā kveḥ iti sūtram paṭhanti ke cit prāk kveḥ iti . \\
\noindent \textbf{(P\_3,2.178.1)} tatra ye ā kveḥ iti paṭhanti taiḥ kvip api ākṣiptaḥ bhavati . \\
\noindent \textbf{(P\_3,2.178.1)} tatra tacchīlādiṣu artheṣu kivp yathā syāt . \\
\noindent \textbf{(P\_3,2.178.1)} na etat asti prayojanam . \\
\noindent \textbf{(P\_3,2.178.1)} yaḥ eva asau aviśeṣavihitaḥ saḥ tacchīlādiṣu bhaviṣyati anyatra ca . \\
\noindent \textbf{(P\_3,2.178.1)} na sidhyati . \\
\noindent \textbf{(P\_3,2.178.1)} ayam tacchīlādiṣu tṛn vidhīyate . \\
\noindent \textbf{(P\_3,2.178.1)} saḥ viśeṣavihitaḥ sāmānyavihitam ṇvulam bādheta . \\
\noindent \textbf{(P\_3,2.178.1)} vāsarūpanyāyena ṇvul api bhaviṣyati . \\
\noindent \textbf{(P\_3,2.178.1)} na sidhyati . \\
\noindent \textbf{(P\_3,2.178.1)} idānīm eva hi uktam tacchīlādiṣu vartheṣu vāsarūpeṇa tṛjādayaḥ na bhavanti iti . \\
\noindent \textbf{(P\_3,2.178.1)} kvip ca api tṛjādiḥ . \\
\noindent \textbf{(P\_3,2.178.2)} vacipracchyāyatastukaṭaprujuśrīṇām dīrghaḥ ca . \\
\noindent \textbf{(P\_3,2.178.2)} vacipracchyāyatastukaṭaprujuśrīṇām dīrghatvam ca vaktavyam kvip ca . \\
\noindent \textbf{(P\_3,2.178.2)} vaci . \\
\noindent \textbf{(P\_3,2.178.2)} vāk . \\
\noindent \textbf{(P\_3,2.178.2)} vaci . \\
\noindent \textbf{(P\_3,2.178.2)} pracchi . \\
\noindent \textbf{(P\_3,2.178.2)} śabdaprāṭ pracchi . \\
\noindent \textbf{(P\_3,2.178.2)} āyatastu . \\
\noindent \textbf{(P\_3,2.178.2)} āyatastūḥ . \\
\noindent \textbf{(P\_3,2.178.2)} āyatastu . \\
\noindent \textbf{(P\_3,2.178.2)} kaṭapru . \\
\noindent \textbf{(P\_3,2.178.2)} kaṭaprūḥ . \\
\noindent \textbf{(P\_3,2.178.2)} kaṭapru . \\
\noindent \textbf{(P\_3,2.178.2)} ju . \\
\noindent \textbf{(P\_3,2.178.2)} jūḥ . \\
\noindent \textbf{(P\_3,2.178.2)} ju . \\
\noindent \textbf{(P\_3,2.178.2)} śri . \\
\noindent \textbf{(P\_3,2.178.2)} śrīḥ . \\
\noindent \textbf{(P\_3,2.178.2)} aparaḥ āha : vacipracchyoḥ asamprasāraṇam ca iti vaktavyam . \\
\noindent \textbf{(P\_3,2.178.2)} tat tarhi vaktavyam . \\
\noindent \textbf{(P\_3,2.178.2)} na vaktavyam . \\
\noindent \textbf{(P\_3,2.178.2)} dīrghavacanasāmarthyāt samprasāraṇam na bhaviṣyati . \\
\noindent \textbf{(P\_3,2.178.2)} idam iha sampradhāryam . \\
\noindent \textbf{(P\_3,2.178.2)} dīrghatvam kriyatām samprasāraṇam iti . \\
\noindent \textbf{(P\_3,2.178.2)} kim atra kartavyam . \\
\noindent \textbf{(P\_3,2.178.2)} paratvāt samprasāraṇam . \\
\noindent \textbf{(P\_3,2.178.2)} antaraṅgam dīrghatvam . \\
\noindent \textbf{(P\_3,2.178.2)} kā antaraṅgatā . \\
\noindent \textbf{(P\_3,2.178.2)} pratyayotpattisanniyogena dīrghatvam ucyate utpanne pratyaye samprasāraṇam . \\
\noindent \textbf{(P\_3,2.178.2)} tatra antaraṅgatvāt dīrghatve kṛte samprasāraṇam . \\
\noindent \textbf{(P\_3,2.178.2)} prasāraṇaparapūrvatve kṛte kāryakṛtatvāt punaḥ dīrghatvam na syāt . \\
\noindent \textbf{(P\_3,2.178.2)} tasmāt suṣṭhu ucyate dīrghavacanasāmarthyāt samprasāraṇam na bhaviṣyati iti . \\
\noindent \textbf{(P\_3,2.178.2)} dyutigamijuhotīnām dve ca . \\
\noindent \textbf{(P\_3,2.178.2)} dyutigamijuhotīnām dve ca iti vaktavyam . \\
\noindent \textbf{(P\_3,2.178.2)} didyut. dyuti . \\
\noindent \textbf{(P\_3,2.178.2)} gami . \\
\noindent \textbf{(P\_3,2.178.2)} jagat . \\
\noindent \textbf{(P\_3,2.178.2)} gami . \\
\noindent \textbf{(P\_3,2.178.2)} juhoti . \\
\noindent \textbf{(P\_3,2.178.2)} juhoteḥ dīrghaḥ ca . \\
\noindent \textbf{(P\_3,2.178.2)} juhūḥ . \\
\noindent \textbf{(P\_3,2.178.2)} dṛṇāteḥ hrasvaḥ ca dve ca kvip ca iti vaktavyam . \\
\noindent \textbf{(P\_3,2.178.2)} dadṛt . \\
\noindent \textbf{(P\_3,2.178.2)} juhūḥ juhoteḥ hvayateḥ vā . \\
\noindent \textbf{(P\_3,2.178.2)} dadṛt dṛṇāteḥ dīryateḥ vā . \\
\noindent \textbf{(P\_3,2.178.2)} jūḥ jvarateḥ jīryateḥ vā . \\
\noindent \textbf{(P\_3,2.178.2)} dhāyateḥ samprasāraṇam ca . \\
\noindent \textbf{(P\_3,2.178.2)} dhāyateḥ samprasāraṇam ca kvip ca vaktavyaḥ . \\
\noindent \textbf{(P\_3,2.178.2)} dhīḥ dhyāyateḥ vā dadhāteḥ vā . \\
\noindent \textbf{(P\_3,2.180)} ḍuprakaraṇe mitadrvādibhyaḥ upasaṅkhyānam dhātuvidhitukpratiṣedhārtham . \\
\noindent \textbf{(P\_3,2.180)} ḍuprakaraṇe mitadrvādibhyaḥ upasaṅkhyānam kartavyam . \\
\noindent \textbf{(P\_3,2.180)} kim prayojanam . \\
\noindent \textbf{(P\_3,2.180)} dhātuvidhitukpratiṣedhārtham . \\
\noindent \textbf{(P\_3,2.180)} dhātuvidheḥ tukaḥ ca pratiṣedhaḥ yathā syāt . \\
\noindent \textbf{(P\_3,2.180)} mitradruḥ mitadrū mitadravaḥ . \\
\noindent \textbf{(P\_3,2.180)} aci śnudhātubhruvām iti uvaṅādeśaḥ mā bhūt . \\
\noindent \textbf{(P\_3,2.180)} iha ca mitadrvā mitadrve na ūṅdhātvoḥ iti pratiṣedhaḥ mā bhūt . \\
\noindent \textbf{(P\_3,2.180)} tugvidhiḥ . \\
\noindent \textbf{(P\_3,2.180)} mitadruḥ . \\
\noindent \textbf{(P\_3,2.180)} hrasvasya piti kṛti tuk bhavati iti tuk mā bhūt . \\
\noindent \textbf{(P\_3,2.188)} śīlitaḥ rakṣitaḥ kṣāntaḥ ākruṣṭaḥ juṣṭaḥ iti api ruṣṭaḥ ca ruṣitaḥ ca ubhau abhivyāhṛtaḥ iti api hṛṣṭatuṣṭau tathā kāntaḥ tathā ubhau saṃyotodyatau . \\
\noindent \textbf{(P\_3,2.188)} kaṣṭam bhaviṣyati iti āhuḥ . \\
\noindent \textbf{(P\_3,2.188)} amṛtāḥ pūrvavat smṛtāḥ . na mriyante amṛtāḥ . \\
\noindent \textbf{(P\_3,3.3)} bhaviṣyati iti anadyatane upasaṅkhyānam . \\
\noindent \textbf{(P\_3,3.3)} bhaviṣyati iti anadyatane upasaṅkhyānam kartavyam . \\
\noindent \textbf{(P\_3,3.3)} śvaḥ grāmam gamī . \\
\noindent \textbf{(P\_3,3.3)} kim punaḥ kāraṇam na sidhyati . \\
\noindent \textbf{(P\_3,3.3)} lṛṭā ayam nirdeśaḥ kṛiyate . \\
\noindent \textbf{(P\_3,3.3)} lṛṭ ca anadyatane luṭā bādhyate . \\
\noindent \textbf{(P\_3,3.3)} tena lṛṭaḥ eva viṣaye ete pratyayāḥ syuḥ . \\
\noindent \textbf{(P\_3,3.3)} luṭaḥ viṣaye na syuḥ . \\
\noindent \textbf{(P\_3,3.3)} itaretarāśrayam ca . \\
\noindent \textbf{(P\_3,3.3)} itaretarāśrayam ca bhavati . \\
\noindent \textbf{(P\_3,3.3)} kā itaretarāśrayatā . \\
\noindent \textbf{(P\_3,3.3)} bhaviṣyatkālena śabdena nirdeśaḥ kriyate . \\
\noindent \textbf{(P\_3,3.3)} nirdeśottarakālam ca bhbhaviṣyatkālatā . \\
\noindent \textbf{(P\_3,3.3)} tat etat itaretarāśrayam bhavati . \\
\noindent \textbf{(P\_3,3.3)} itaretarāśrayāṇi ca na prakalpante . \\
\noindent \textbf{(P\_3,3.3)} uktam vā . \\
\noindent \textbf{(P\_3,3.3)} kim uktam . \\
\noindent \textbf{(P\_3,3.3)} ekam tāvat uktam na vā apavādasya nimittābhāvāt anadyatane hi tayoḥ vidhānam iti . \\
\noindent \textbf{(P\_3,3.3)} aparam api uktam avyayanirdeśāt siddham iti. avyayavatā śabdena nirdeśaḥ kariṣyate . \\
\noindent \textbf{(P\_3,3.3)} avartamāne abhūte iti . \\
\noindent \textbf{(P\_3,3.3)} saḥ tarhi avyayavatā śabdena nirdeśaḥ kartavyaḥ . \\
\noindent \textbf{(P\_3,3.3)} na kartavyaḥ . \\
\noindent \textbf{(P\_3,3.3)} avyayam eṣaḥ bhaviṣyatiśabdaḥ na eṣā bhavateḥ lṛṭ . \\
\noindent \textbf{(P\_3,3.3)} katham avyayatvam . \\
\noindent \textbf{(P\_3,3.3)} vibhaktisvarapratirūpakāḥ ca nipātāḥ bhavanti iti nipātasañjñā . \\
\noindent \textbf{(P\_3,3.3)} nipātam avyayam iti avayayasañjñā . \\
\noindent \textbf{(P\_3,3.3)} atha api bhavateḥ lṛṭ evam api avayayam eva . \\
\noindent \textbf{(P\_3,3.3)} katham na vyeti iti avyayam . \\
\noindent \textbf{(P\_3,3.3)} kva punaḥ na vyeti . \\
\noindent \textbf{(P\_3,3.3)} etau kālaviśeṣau bhūtavartamānau . \\
\noindent \textbf{(P\_3,3.3)} svabhāvataḥ bhaviṣyati eva vartate . \\
\noindent \textbf{(P\_3,3.3)} yadi tatri na vyeti iti avyayam . \\
\noindent \textbf{(P\_3,3.3)} na vā tadvidhānasya anyatra abhāvāt . \\
\noindent \textbf{(P\_3,3.3)} na vā bhaviṣyadādhikāreṇa arthaḥ . \\
\noindent \textbf{(P\_3,3.3)} kim kāraṇam . \\
\noindent \textbf{(P\_3,3.3)} tadvidhānasya anyatra abhāvāt . \\
\noindent \textbf{(P\_3,3.3)} ye api ete itaḥ uttaram pratyayāḥ śiṣyante ete api etau kālaviśeṣau na viyanti bhūtavartamānau . \\
\noindent \textbf{(P\_3,3.3)} svabhāvataḥ eva te bhaviṣyati eva vartante . \\
\noindent \textbf{(P\_3,3.3)} ataḥ uttaram paṭhati . \\
\noindent \textbf{(P\_3,3.3)} bhaviṣyadadhikārasya prayojanam yāvat pacati purā pacati iti anapaśabdatvāya . \\
\noindent \textbf{(P\_3,3.4)} yāvatpurādiṣu laḍvidhiḥ luṭaḥ pūrvavipratiṣiddham . yāvatpurādiṣu laḍvidhiḥ bhavati luṭaḥ pūrvavipratiṣedhena . \\
\noindent \textbf{(P\_3,3.4)} yāvatpurānipātayoḥ laṭ bhavati iti asya avakāśaḥ . \\
\noindent \textbf{(P\_3,3.4)} yāvat bhuṅkte . \\
\noindent \textbf{(P\_3,3.4)} purā bhuṅkte . \\
\noindent \textbf{(P\_3,3.4)} luṭaḥ avakāśaḥ . \\
\noindent \textbf{(P\_3,3.4)} śvaḥ kartā . \\
\noindent \textbf{(P\_3,3.4)} śvaḥ adhyetā . \\
\noindent \textbf{(P\_3,3.4)} iha ubhayam prāpnoti . \\
\noindent \textbf{(P\_3,3.4)} yāvat śvaḥ bhuṅkte . \\
\noindent \textbf{(P\_3,3.4)} purā śvaḥ bhuṅkte . \\
\noindent \textbf{(P\_3,3.4)} laṭ bhavati vipratiṣedhena . \\
\noindent \textbf{(P\_3,3.4)} saḥ tarhi pūrvavipratiṣedhaḥ vaktavyaḥ . \\
\noindent \textbf{(P\_3,3.4)} na vaktavyaḥ . \\
\noindent \textbf{(P\_3,3.4)} anadyatane luṭ iti atra yāvatpurānipātayoḥ laṭ iti anuvartiṣyate . \\
\noindent \textbf{(P\_3,3.7)} kimartham idam ucyate na lipsyamānasiddhiḥ api lipsā eva tatra kiṃvṛtte lipsāyām iti eva siddham . \\
\noindent \textbf{(P\_3,3.7)} akiṃvṛttārthaḥ ayam ārambhaḥ . \\
\noindent \textbf{(P\_3,3.7)} yaḥ bhavatām odanam dadāti saḥ svargam lokam gacchati . \\
\noindent \textbf{(P\_3,3.7)} yaḥ bhavatām odanam dāsyati saḥ svargam lokam gamiṣyati . \\
\noindent \textbf{(P\_3,3.10)} kimartham kriyāyām upapade kriyārthāyām ṇvul vidhīyate na aviśeṣeṇa vihitaḥ ṇvul saḥ kriyāyām upapade kriyārthāyām anyatra ca bhaviṣyati . \\
\noindent \textbf{(P\_3,3.10)} ṇvuli sakarmakagrahaṇam coditam . \\
\noindent \textbf{(P\_3,3.10)} akaramakārthaḥ ayam ārambhaḥ . \\
\noindent \textbf{(P\_3,3.10)} āsakaḥ vrajati . \\
\noindent \textbf{(P\_3,3.10)} śāyakaḥ vrajati . \\
\noindent \textbf{(P\_3,3.10)} pratyākhyātam tat na vā dhātumātrāt darśanāt ṇvulaḥ iti . \\
\noindent \textbf{(P\_3,3.10)} evam tarhi tṛjādiṣu vartamānakālopādānam coditam . \\
\noindent \textbf{(P\_3,3.10)} avartamānakālārthaḥ ayam ārambhaḥ . \\
\noindent \textbf{(P\_3,3.10)} tat api pratyākhyātam na vā kālamātre darśanāt anyeṣām iti . \\
\noindent \textbf{(P\_3,3.10)} idam tarhi prayojanam. akenoḥ bhaviṣyadādhamrṇyayoḥ iti atra ṣaṣṭhyāḥ pratiṣedhaḥ uktaḥ . \\
\noindent \textbf{(P\_3,3.10)} saḥ yathā syāt . \\
\noindent \textbf{(P\_3,3.10)} etat api na asti prayojanam . \\
\noindent \textbf{(P\_3,3.10)} yaḥ eva asau aviśeṣavihitaḥ saḥ yadā bhaviṣyati bhaviṣyati tadā asya pratiṣedhaḥ bhaviṣyati . \\
\noindent \textbf{(P\_3,3.10)} evam tarhi bhaviṣyadadhikāravihitasya pratiṣedhaḥ yathā syāt . \\
\noindent \textbf{(P\_3,3.10)} iha mā bhūt : aṅga yajatām . \\
\noindent \textbf{(P\_3,3.10)} lapsyante asya yājakāḥ . \\
\noindent \textbf{(P\_3,3.10)} ye enam yājayayiṣyanti iti . \\
\noindent \textbf{(P\_3,3.10)} na eṣaḥ bhaviṣyatkālaḥ . \\
\noindent \textbf{(P\_3,3.10)} kaḥ tarhi . \\
\noindent \textbf{(P\_3,3.10)} bhūtakālaḥ . \\
\noindent \textbf{(P\_3,3.10)} katham tarhi bhaviṣyatkālatā gamyate . \\
\noindent \textbf{(P\_3,3.10)} dhātusambandhe pratyayāḥ iti . \\
\noindent \textbf{(P\_3,3.10)} yaḥ tarhi na dhātusambandhaḥ . \\
\noindent \textbf{(P\_3,3.10)} ime asya yājakāḥ . \\
\noindent \textbf{(P\_3,3.10)} ime asya lāvakāḥ iti . \\
\noindent \textbf{(P\_3,3.10)} eṣaḥ api bhūtakālaḥ . \\
\noindent \textbf{(P\_3,3.10)} katham tarhi bhaviṣyatkālatā gamyate . \\
\noindent \textbf{(P\_3,3.10)} sambandhāt . \\
\noindent \textbf{(P\_3,3.10)} saḥ ca tāvat taiḥ ayājitaḥ bhavati . \\
\noindent \textbf{(P\_3,3.10)} tasya ca tāvat taiḥ yavāḥ alūnāḥ bhavanti . \\
\noindent \textbf{(P\_3,3.10)} ucyate ca . \\
\noindent \textbf{(P\_3,3.10)} idam tarhi prayojanam . \\
\noindent \textbf{(P\_3,3.10)} ayam kriyāyām upapde kriyārthāyām tumun vidhīyate . \\
\noindent \textbf{(P\_3,3.10)} saḥ viśeṣavihitaḥ sāmānyavihitam ṇvulam bādheta . \\
\noindent \textbf{(P\_3,3.10)} etat api na asti prayojanam . \\
\noindent \textbf{(P\_3,3.10)} bhāve tumun vidhīyate kartari ṇvul . \\
\noindent \textbf{(P\_3,3.10)} tatra kaḥ prasaṅgaḥ yat bhāve vihitaḥ tumun kartari vihitam ṇvulam bādheta . \\
\noindent \textbf{(P\_3,3.10)} lṛṭ tarhi bādheta . \\
\noindent \textbf{(P\_3,3.10)} vāsarūpeṇa bhaviṣyati . \\
\noindent \textbf{(P\_3,3.10)} ataḥ uttaram paṭhati . \\
\noindent \textbf{(P\_3,3.10)} ṇvulaḥ kriyārthopapadasya punarvidhānam tṛjādipratiṣedhārtham . \\
\noindent \textbf{(P\_3,3.10)} ṇvulaḥ kriyārthopapadasya punarvidhānam kriyate jñāpakārtham . \\
\noindent \textbf{(P\_3,3.10)} kim jñāpyam . \\
\noindent \textbf{(P\_3,3.10)} etat jñāpayati ācāryaḥ kriyāyām upapade kriyārthāyām vāsarūpeṇa tṛjādayaḥ na bhavanti iti . \\
\noindent \textbf{(P\_3,3.10)} ṇvul api tṛjādiḥ . \\
\noindent \textbf{(P\_3,3.11)} kimartham idam ucyate na aviśeṣeṇa bhāve pratyayāḥ ye vihitāḥ te kriyāyām upapade kriyārthāyām anyatra ca bhaviṣyanti . \\
\noindent \textbf{(P\_3,3.11)} bhāvavacanānām yathāvihitānām pratipadavidhyartham . \\
\noindent \textbf{(P\_3,3.11)} bhāvavacanānām yathāvihitānām pratipadavidhyarthaḥ ayam ārambhaḥ . \\
\noindent \textbf{(P\_3,3.11)} idānīm eva hi uktam kriyāyām upapde kriyārthāyām vāsarūpeṇa tṛjādayaḥ na bhavanti iti . \\
\noindent \textbf{(P\_3,3.11)} bhāvavacanāḥ ca api tṛjādayaḥ . \\
\noindent \textbf{(P\_3,3.11)} asti prayojanam etat . \\
\noindent \textbf{(P\_3,3.11)} kim tarhi iti . \\
\noindent \textbf{(P\_3,3.11)} yathāvihitāḥ iti tu vaktavyam . \\
\noindent \textbf{(P\_3,3.11)} kim prayojanam . \\
\noindent \textbf{(P\_3,3.11)} iha yābhyaḥ prakṛtibhyaḥ yena viśeṣeṇa bhāve pratyayāḥ vihitāḥ tābhyaḥ prakṛtibhyaḥ tena eva viśeṣeṇa kriyāyām upapde kriyārthāyām yathā syuḥ . \\
\noindent \textbf{(P\_3,3.11)} vyatikaraḥ mā bhūt iti . \\
\noindent \textbf{(P\_3,3.11)} tat tarhi vaktavyam . \\
\noindent \textbf{(P\_3,3.11)} na vaktavyam . \\
\noindent \textbf{(P\_3,3.11)} iha bhāve pratyayāḥ bhavanti iti iyaṭa siddham . \\
\noindent \textbf{(P\_3,3.11)} saḥ ayam evam siddhe sati yat vacanagrahaṇam karoti tasya etat prayojanam vācakāḥ yathā syuḥ iti . \\
\noindent \textbf{(P\_3,3.11)} yadi ca yābhyaḥ prakṛtibhyaḥ yena viśeṣeṇa bhāve pratyayāḥ vihitāḥ tābhyaḥ prakṛtibhyaḥ tena eva viśeṣeṇa kriyāyām upapde kriyārthāyām bhavanti tataḥ amī vācakāḥ kṛtāḥ syuḥ . \\
\noindent \textbf{(P\_3,3.11)} atha hi prakṛtimātrāt vā syuḥ pratyayamātram vā syāt na amīvācakāḥ kṛtāḥ syuḥ . \\
\noindent \textbf{(P\_3,3.12)} kimartham idam ucyate na aviśeṣeṇa karmaṇi aṇ vihitaḥ saḥ kriyāyām upapade kriyārthāyām anyatra ca bhaviṣyati . \\
\noindent \textbf{(P\_3,3.12)} aṇaḥ punarvacanam apavādaviṣaye anivṛttyartham . \\
\noindent \textbf{(P\_3,3.12)} aṇaḥ punarvacanam kriyate apavādaviṣaye anivṛttiḥ yathā syāt . \\
\noindent \textbf{(P\_3,3.12)} godāyaḥ vrajati . \\
\noindent \textbf{(P\_3,3.12)} kambaladāyaḥ vrajati iti . \\
\noindent \textbf{(P\_3,3.12)} kim ucyate apavādaviṣaye anivṛttiḥ yathā syāt iti na punaḥ utsargaviṣaye pratipadavidhyartham syāt . \\
\noindent \textbf{(P\_3,3.12)} idānīm eva hi uktam kriyāyām upapde kriyārthāyām vāsarūpeṇa tṛjādayaḥ na bhavanti iti . \\
\noindent \textbf{(P\_3,3.12)} aṇ ca api tṛjādiḥ . \\
\noindent \textbf{(P\_3,3.12)} evam tarhi ubhayam anena kriyate . \\
\noindent \textbf{(P\_3,3.12)} apavādaviṣaye cānivṛttiḥ utsargaviṣaye pratipadavidhānam . \\
\noindent \textbf{(P\_3,3.12)} katham punaḥ ekena yatnena ubhayam labhyam . \\
\noindent \textbf{(P\_3,3.12)} labhyam iti āha . \\
\noindent \textbf{(P\_3,3.12)} katham . \\
\noindent \textbf{(P\_3,3.12)} karmagrahaṇasāmarthyāt . \\
\noindent \textbf{(P\_3,3.12)} katham punaḥ antareṇa karmagrahaṇam karmaṇi aṇ labhyaḥ . \\
\noindent \textbf{(P\_3,3.12)} vacanagrahaṇam prakṛtam anuvartate . \\
\noindent \textbf{(P\_3,3.12)} asti prayojanam etat . \\
\noindent \textbf{(P\_3,3.12)} kim tarhi iti . \\
\noindent \textbf{(P\_3,3.12)} aparyāyeṇa iti tu vaktavyam . \\
\noindent \textbf{(P\_3,3.12)} kadā cit hi karmaṇi syāt kadā cit kriyāyām upapade kriyārthāyām iti . \\
\noindent \textbf{(P\_3,3.12)} tat tarhi vaktavyam . \\
\noindent \textbf{(P\_3,3.12)} na vaktavyam . \\
\noindent \textbf{(P\_3,3.12)} cena sanniyogaḥ kariṣyate . \\
\noindent \textbf{(P\_3,3.12)} aṇ karmaṇi ca . \\
\noindent \textbf{(P\_3,3.12)} kim ca anyat . \\
\noindent \textbf{(P\_3,3.12)} kriyāyām upapade kriyārthāyām iti . \\
\noindent \textbf{(P\_3,3.12)} evam api pratyekam upapadasañjñā na prāpnoti . \\
\noindent \textbf{(P\_3,3.12)} cena eva sanniyogaḥ kariṣyate . \\
\noindent \textbf{(P\_3,3.12)} pratyekam vākyaparisamāptiḥ dṛṣṭā iti pratyekam upapadasañjña bhaviṣyati . \\
\noindent \textbf{(P\_3,3.13)} śeṣavacanam kimartham . \\
\noindent \textbf{(P\_3,3.13)} lṛṭi śeṣavacanam kriyāyām pratipadavidhyartham . lṛṭi śeṣavacanam kriyate kriyāyām pratipadavidhyartham . \\
\noindent \textbf{(P\_3,3.13)} pratipadavidhiḥ yathā syāt . \\
\noindent \textbf{(P\_3,3.13)} aviśeṣeṇa vidhāne lṛṭaḥ abhāvaḥ pratiṣiddhatvāt . \\
\noindent \textbf{(P\_3,3.13)} aviśeṣeṇa vidhāne lṛṭaḥ abhāvaḥ syāt . \\
\noindent \textbf{(P\_3,3.13)} kariṣyāmi iti vrajati . \\
\noindent \textbf{(P\_3,3.13)} hariṣyāmi iti vrajati iti . \\
\noindent \textbf{(P\_3,3.13)} kim kāraṇam . \\
\noindent \textbf{(P\_3,3.13)} pratiṣiddhatvāt . \\
\noindent \textbf{(P\_3,3.13)} idānīm eva hi uktam kriyāyām upapade kriyārthāyām vāsarūpeṇa tṛjādayaḥ na bhavanti iti . \\
\noindent \textbf{(P\_3,3.13)} lṛṭ ca api tṛjādiḥ . \\
\noindent \textbf{(P\_3,3.13)} asti prayojanam etat . \\
\noindent \textbf{(P\_3,3.13)} kim tarhi iti . \\
\noindent \textbf{(P\_3,3.13)} sādhīyaḥ tu khalu śeṣagrahaṇena kriyārthopapadāt lṛṭ nirbhajyate . \\
\noindent \textbf{(P\_3,3.13)} kim kāraṇam . \\
\noindent \textbf{(P\_3,3.13)} akriyārthopapadatvāt . \\
\noindent \textbf{(P\_3,3.13)} śeṣe iti ucyate . \\
\noindent \textbf{(P\_3,3.13)} śeṣaḥ ca kaḥ . \\
\noindent \textbf{(P\_3,3.13)} yat anyat kriyāyāḥ kriyārthāyāḥ . \\
\noindent \textbf{(P\_3,3.13)} evam tarhi lṛṭi śeṣavacanam kriyāyām pratipadavidhyartham . \\
\noindent \textbf{(P\_3,3.13)} lṛṭi śeṣavacanam kriyate kriyāyām pratipadavidhiḥ yathā syāt . \\
\noindent \textbf{(P\_3,3.13)} lṛṭ śeṣe ca . \\
\noindent \textbf{(P\_3,3.13)} kariṣyati hariṣyati iti . \\
\noindent \textbf{(P\_3,3.13)} kva ca . \\
\noindent \textbf{(P\_3,3.13)} kriyāyām upapade kriyārthāyām iti . \\
\noindent \textbf{(P\_3,3.13)} saḥ tarhi cakāraḥ kartavyaḥ . \\
\noindent \textbf{(P\_3,3.13)} na kartavyaḥ . \\
\noindent \textbf{(P\_3,3.13)} iha lṛṭ bhavati iti iyatā siddham . \\
\noindent \textbf{(P\_3,3.13)} saḥ ayam evam siddhe sati yat śeṣagrahaṇam karoti tasya etat prayojanam yogāṅgam yathā upajāyeta . \\
\noindent \textbf{(P\_3,3.13)} sati ca yogāṅge yogavibhāgaḥ kariṣyate . \\
\noindent \textbf{(P\_3,3.13)} lṛṭ bhavati kriyāyām upapade kriyārthāyām iti . \\
\noindent \textbf{(P\_3,3.13)} tataḥ śeṣe . \\
\noindent \textbf{(P\_3,3.13)} śeṣe ca lṛṭ bhavati iti . \\
\noindent \textbf{(P\_3,3.14)} sadvidhiḥ nityam aprathamāsamānādhikaraṇe . \\
\noindent \textbf{(P\_3,3.14)} sadvidhiḥ aprathamāsamānādhikaraṇe nityam iti vaktavyam . \\
\noindent \textbf{(P\_3,3.14)} pakṣyantam paśya . \\
\noindent \textbf{(P\_3,3.14)} pakṣyamāṇam paśya . \\
\noindent \textbf{(P\_3,3.14)} kva tarhi idānīm vibhāṣā . \\
\noindent \textbf{(P\_3,3.14)} prathamāsamānādhikaraṇe . \\
\noindent \textbf{(P\_3,3.14)} pākṣyan pakṣyati . \\
\noindent \textbf{(P\_3,3.14)} pakṣyamāṇaḥ pakṣyate . \\
\noindent \textbf{(P\_3,3.15.1)} yogavibhāgaḥ kartavyaḥ . \\
\noindent \textbf{(P\_3,3.15.1)} anadyatane lṛṭaḥ satsñjñau bhavataḥ . \\
\noindent \textbf{(P\_3,3.15.1)} śvaḥ agnīn ādhyāsyamānena . \\
\noindent \textbf{(P\_3,3.15.1)} śvaḥ somena yakṣyamāṇena . \\
\noindent \textbf{(P\_3,3.15.1)} tataḥ luṭ . \\
\noindent \textbf{(P\_3,3.15.1)} luṭ bhavati anadyatane . \\
\noindent \textbf{(P\_3,3.15.1)} śvaḥ kartā . \\
\noindent \textbf{(P\_3,3.15.1)} śvaḥ adhyetā . \\
\noindent \textbf{(P\_3,3.15.1)} kena vihitasya anadyatane lṛṭaḥ satsañjñau ucyete . \\
\noindent \textbf{(P\_3,3.15.1)} etat eva jñāpayati bhavati anadyatane lṛṭ iti yat ayam anadyatane lṛṭaḥ satsñjñau śāsti . \\
\noindent \textbf{(P\_3,3.15.1)} evam ca kṛtvā saḥ api adoṣaḥ bhavati yat uktam bhaviṣyati iti anadyatane upasaṅkhyānam . \\
\noindent \textbf{(P\_3,3.15.2)} paridevane śvastanībhaviṣyantyarthe . \\
\noindent \textbf{(P\_3,3.15.2)} paridevane śvastanībhaviṣyantyāḥ arthe iti vaktavyam . \\
\noindent \textbf{(P\_3,3.15.2)} iyam nu kadā gantā yā evam pādau nidadhāti . \\
\noindent \textbf{(P\_3,3.15.2)} ayam nu kadā adhyetā yaḥ evam anabhiyuktaḥ iti . \\
\noindent \textbf{(P\_3,3.15.2)} kālaprakarṣāt tu upamānam . \\
\noindent \textbf{(P\_3,3.15.2)} kālaprakarṣāt tu upamānam . \\
\noindent \textbf{(P\_3,3.15.2)} gantā iva iyam gantā . \\
\noindent \textbf{(P\_3,3.15.2)} na iyam gamiṣyati . \\
\noindent \textbf{(P\_3,3.15.2)} adhyetā iva ayam adhyetā . \\
\noindent \textbf{(P\_3,3.15.2)} na vai tiṅantena upamānam asti . \\
\noindent \textbf{(P\_3,3.15.2)} evam tarhi anadyatane iva anadyatane iti . \\
\noindent \textbf{(P\_3,3.16)} spṛśaḥ upatāpe . \\
\noindent \textbf{(P\_3,3.16)} spṛśaḥ upatāpe iti vaktavyam . \\
\noindent \textbf{(P\_3,3.16)} iha mā bhūt . \\
\noindent \textbf{(P\_3,3.16)} kambalasparśaḥ iti \\
\noindent \textbf{(P\_3,3.17)} vādhimatsyabaleṣu iti vaktavyam . \\
\noindent \textbf{(P\_3,3.17)} atīsāraḥ vyādhiḥ . \\
\noindent \textbf{(P\_3,3.17)} visāraḥ matsyaḥ . \\
\noindent \textbf{(P\_3,3.17)} bale . \\
\noindent \textbf{(P\_3,3.17)} śālasāraḥ khadirasāraḥ . \\
\noindent \textbf{(P\_3,3.18)} bhāve sarvaliṅganirdeśaḥ . \\
\noindent \textbf{(P\_3,3.18)} bhāve sarvaliṅganirdeśaḥ kartavyaḥ . \\
\noindent \textbf{(P\_3,3.18)} bhūtau bhavane bhāve iti . \\
\noindent \textbf{(P\_3,3.18)} kim prayojanam . \\
\noindent \textbf{(P\_3,3.18)} sarvaliṅge bhāṅe ete pratyayāḥ yathā syuḥ iti . \\
\noindent \textbf{(P\_3,3.18)} kim punaḥ kāraṇam na sidhyati . \\
\noindent \textbf{(P\_3,3.18)} puṃliṅgena ayam nirdeśaḥ kriyate ekavacanena ca . \\
\noindent \textbf{(P\_3,3.18)} ten puṃliṅge eva ekavacane ca ete prratyayāḥ syuḥ . \\
\noindent \textbf{(P\_3,3.18)} strīnapuṃsakayoḥ dvivacanabahuvacnayoḥ ca na syuḥ . \\
\noindent \textbf{(P\_3,3.18)} na atra nirdeśaḥ tantram . \\
\noindent \textbf{(P\_3,3.18)} katham punaḥ tena eva ca nāma nirdeśaḥ kriyate tat ca atantram syāt . \\
\noindent \textbf{(P\_3,3.18)} tatkārī ca bhavān taddveṣī ca . \\
\noindent \textbf{(P\_3,3.18)} nāntarīyakatvāt atra puṃliṅgena nirdeśaḥ kriyate ekavacanena ca . \\
\noindent \textbf{(P\_3,3.18)} avaśyam kayā cit vibhaktyā kena cit ca liṅgena nirdeśaḥ kartavyaḥ . \\
\noindent \textbf{(P\_3,3.18)} tat yathā kaḥ cit annārthī śālikalāpam satuṣam sapalālam āharati nāntarīyakatvāt . \\
\noindent \textbf{(P\_3,3.18)} saḥ yāvat ādeyam tāvat ādāya tuṣapalālāni utsṛjati . \\
\noindent \textbf{(P\_3,3.18)} tathā kaḥ cit māṃsārthī matsyān saśakalān sakaṇṭakān āharati nāntarīyakatvāt . \\
\noindent \textbf{(P\_3,3.18)} saḥ yāvat ādeyam tāvat ādāya śakalakaṇṭakān utsṛjati . \\
\noindent \textbf{(P\_3,3.18)} evam iha api nāntarīyakatvāt puṃliṅgena nirdeśaḥ kriyate ekavacanāntena ca . \\
\noindent \textbf{(P\_3,3.18)} na hi atra nirdeśaḥ tantram . \\
\noindent \textbf{(P\_3,3.18)} kayā cit vibhaktyā kena cit ca liṅgena nirdeśaḥ kartavyaḥ . \\
\noindent \textbf{(P\_3,3.18)} atha vā kṛbhvastayaḥ kriyāsāmānyavācinaḥ kriyāviśeṣavācinaḥ pacādayaḥ . \\
\noindent \textbf{(P\_3,3.18)} yat ca atra pacateḥ bhavatiḥ bhavati na tat bhavateḥ pacatiḥ bhavati . \\
\noindent \textbf{(P\_3,3.18)} yat ca bhavateḥ pacatiḥ bhavati na tat pacateḥ bhavatiḥ bhavati . \\
\noindent \textbf{(P\_3,3.18)} kim ca pacateḥ bhavatiḥ bhavati . \\
\noindent \textbf{(P\_3,3.18)} sāmānyam . \\
\noindent \textbf{(P\_3,3.18)} kim ca bhavateḥ pacatiḥ bhavati . \\
\noindent \textbf{(P\_3,3.18)} viśeṣaḥ . \\
\noindent \textbf{(P\_3,3.18)} tat yathā upādhyāyasya śiṣyaḥ mātulasya bhāgineyam gatvā āha . \\
\noindent \textbf{(P\_3,3.18)} upādhyāyam bhavān abhivādayatām iti . \\
\noindent \textbf{(P\_3,3.18)} saḥ gatvā mātulam abhivādayate . \\
\noindent \textbf{(P\_3,3.18)} tathā mātulasya bhāgineyaḥ upādhyāyasya śiṣyam gatvā āha . \\
\noindent \textbf{(P\_3,3.18)} mātulam bhavān abhivādayatām iti . \\
\noindent \textbf{(P\_3,3.18)} saḥ gatvā upādhyāyam abhivādayate . \\
\noindent \textbf{(P\_3,3.18)} evam iha api pacateḥ bhavatau yat tat nirdiśyate . \\
\noindent \textbf{(P\_3,3.19)} kārakagrahaṇam kimartham . \\
\noindent \textbf{(P\_3,3.19)} kārakagrahaṇam anādeśe svārthavijñānāt . \\
\noindent \textbf{(P\_3,3.19)} kārakagrahaṇam anādeśe svārthavijñānāt . \\
\noindent \textbf{(P\_3,3.19)} anirdiṣṭārthāḥ pratyayāḥ svārthe bhavanti iti . \\
\noindent \textbf{(P\_3,3.19)} tat yathā . \\
\noindent \textbf{(P\_3,3.19)} guptijkidbhyaḥ san yāvādibhyaḥ kan iti . \\
\noindent \textbf{(P\_3,3.19)} evam ime api pratyayāḥ svārthe syuḥ . \\
\noindent \textbf{(P\_3,3.19)} svārthe mā bhūvan kārake yathā syuḥ iti evamartham idam ucyate . \\
\noindent \textbf{(P\_3,3.19)} na etat asti prayojanam . \\
\noindent \textbf{(P\_3,3.19)} vihitaḥ pratyayaḥ svārthe bhāve ghañ iti . \\
\noindent \textbf{(P\_3,3.19)} tena atriprasaktam iti kṛtvā niyamārthaḥ ayam vijñāyeta . \\
\noindent \textbf{(P\_3,3.19)} akartari sañjñāyām eva iti . \\
\noindent \textbf{(P\_3,3.19)} asti ca idānīm kaḥ cit sañjñābhūtaḥ bhāvaḥ yadarthaḥ vidhiḥ syāt . \\
\noindent \textbf{(P\_3,3.19)} asti iti āha : āvāhaḥ , vivāhaḥ iti . \\
\noindent \textbf{(P\_3,3.19)} kaimarthakyāt niyamaḥ bhavati . \\
\noindent \textbf{(P\_3,3.19)} vidheyam na asti iti kṛtvā . \\
\noindent \textbf{(P\_3,3.19)} iha ca asti vidheyam . \\
\noindent \textbf{(P\_3,3.19)} akartari ca kārake sañjñāyām ghañ vidheyaḥ . \\
\noindent \textbf{(P\_3,3.19)} tatra apūrvaḥ vidhiḥ astu niyamaḥ astu iti apūrvaḥ eva vidhiḥ bhaviṣyati na niyamaḥ . \\
\noindent \textbf{(P\_3,3.19)} tat eva tarhi prayojanam svārthe mā bhūvan iti . \\
\noindent \textbf{(P\_3,3.19)} nanu ca uktam vihitaḥ pratyayaḥ svārthe bhāve ghañ iti iti . \\
\noindent \textbf{(P\_3,3.19)} anyaḥ saḥ bhāvaḥ bāhyaḥ prakṛtyarthāt . \\
\noindent \textbf{(P\_3,3.19)} anena idānīm ābhyantare bhāve syāt . \\
\noindent \textbf{(P\_3,3.19)} kaḥ punaḥ etayoḥ bhāvayoḥ viśeṣaḥ . \\
\noindent \textbf{(P\_3,3.19)} uktaḥ bhāvabhedaḥ bhāṣye . \\
\noindent \textbf{(P\_3,3.19)} etat api na asti prayojanam . \\
\noindent \textbf{(P\_3,3.19)} nañivayuktam anyasadṛśādhikaraṇe . \\
\noindent \textbf{(P\_3,3.19)} tathā hi arthagatiḥ . \\
\noindent \textbf{(P\_3,3.19)} nañyuktam ivayuktam ca asnyasmin tatsadṛśe kāryam vijñayate . \\
\noindent \textbf{(P\_3,3.19)} tathā hi arthaḥ gamyate . \\
\noindent \textbf{(P\_3,3.19)} tat yathā . \\
\noindent \textbf{(P\_3,3.19)} abrāhmaṇam ānaya iti ukte brāhmaṇasadṛśam puruṣam ānayati . \\
\noindent \textbf{(P\_3,3.19)} na asau loṣṭam ānīya kṛtī bhavati . \\
\noindent \textbf{(P\_3,3.19)} evam iha api akartari iti kartṛpratiṣedhāt anyasmin akartari kartṛsadṛśe kāryam vijñasyate . \\
\noindent \textbf{(P\_3,3.19)} kim ca anyat akartṛ kartṛsadṛśam . \\
\noindent \textbf{(P\_3,3.19)} kārakam . \\
\noindent \textbf{(P\_3,3.19)} uttarārtham tarhi kārakagrahaṇam kartavyam . \\
\noindent \textbf{(P\_3,3.19)} parimāṇākhyāyām sarvebhyaḥ kārake yathā syāt . \\
\noindent \textbf{(P\_3,3.19)} iha mā bhūt . \\
\noindent \textbf{(P\_3,3.19)} ekā tilocchritiḥ . \\
\noindent \textbf{(P\_3,3.19)} deve sṛtī iti . \\
\noindent \textbf{(P\_3,3.19)} ghañanukramaṇam ajabviṣaye avacane hi strīpratyayānām api avādavijñānam iti vakṣyati . \\
\noindent \textbf{(P\_3,3.19)} tat na vaktavyam bhavati . \\
\noindent \textbf{(P\_3,3.19)} etat api na asti prayojanam . \\
\noindent \textbf{(P\_3,3.19)} atra api akartari iti eva anuvartiṣyate . \\
\noindent \textbf{(P\_3,3.19)} sañjñāgrahaṇānarthakyam ca sarvatra ghañaḥ darśanāt . \\
\noindent \textbf{(P\_3,3.19)} sañjñāgrahaṇam ca anarthakam . \\
\noindent \textbf{(P\_3,3.19)} kim kāraṇam . \\
\noindent \textbf{(P\_3,3.19)} sarvatra ghañaḥ darśanāt . \\
\noindent \textbf{(P\_3,3.19)} asañjñāyām api hi ghañ dṛśyate . \\
\noindent \textbf{(P\_3,3.19)} kaḥ bhavatā dāyaḥ dattaḥ . \\
\noindent \textbf{(P\_3,3.19)} kaḥ bhavatā lābhaḥ labdhaḥ iti . \\
\noindent \textbf{(P\_3,3.19)} yadi sañjñāgrahaṇam na kriyate atiprasaṅgaḥ bhavati . \\
\noindent \textbf{(P\_3,3.19)} kṛtaḥ kaṭaḥ iti atra kāraḥ kaṭa iti prāpnoti . \\
\noindent \textbf{(P\_3,3.19)} atiprasaṅgaḥ iti cet abhidhānalakṣaṇatvāt pratyayasya siddham . \\
\noindent \textbf{(P\_3,3.19)} atiprasaṅgaḥ iti cet tat na . \\
\noindent \textbf{(P\_3,3.19)} kim kāraṇam . \\
\noindent \textbf{(P\_3,3.19)} abhidhānalakṣaṇatvāt pratyayasya siddham . \\
\noindent \textbf{(P\_3,3.19)} abhidhānalakṣaṇāḥ kṛttaddhitasamādāḥ . \\
\noindent \textbf{(P\_3,3.19)} anabhidhānāt na bhaviṣyanti . \\
\noindent \textbf{(P\_3,3.20.1)} sarvagrahaṇam kimartham . \\
\noindent \textbf{(P\_3,3.20.1)} sarvebhyaḥ dhātubhyaḥ ghañ yathā syāt ajapoḥ api viṣaye . \\
\noindent \textbf{(P\_3,3.20.1)} ekaḥ taṇḍulaniścāyaḥ . \\
\noindent \textbf{(P\_3,3.20.1)} dvau śūrpaniṣpāvau . \\
\noindent \textbf{(P\_3,3.20.1)} sarvagrahaṇam anarthakam parimāṇākhyāyam iti siddhatvāt . \\
\noindent \textbf{(P\_3,3.20.1)} sarvagrahaṇam anarthakam . \\
\noindent \textbf{(P\_3,3.20.1)} kim kāraṇam . \\
\noindent \textbf{(P\_3,3.20.1)} parimāṇākhyāyam iti siddhatvāt . \\
\noindent \textbf{(P\_3,3.20.1)} parimāṇākhyāyam iti eva ghañ siddhaḥ ajapoḥ api viṣaye . \\
\noindent \textbf{(P\_3,3.20.1)} na arthaḥ sarvagrahaṇena . \\
\noindent \textbf{(P\_3,3.20.1)} asti anyat etasya vacane prayojanam . \\
\noindent \textbf{(P\_3,3.20.1)} kim . \\
\noindent \textbf{(P\_3,3.20.1)} ekaḥ pākaḥ dvau pākau trayaḥ pākāḥ iti . \\
\noindent \textbf{(P\_3,3.20.1)} pūrveṇa api etat siddham . \\
\noindent \textbf{(P\_3,3.20.1)} na sidhyati . \\
\noindent \textbf{(P\_3,3.20.1)} sañjñāyām iti pūrvaḥ yogaḥ . \\
\noindent \textbf{(P\_3,3.20.1)} na ca eṣā sañjñā . \\
\noindent \textbf{(P\_3,3.20.1)} pratyākhyāyate sañjñāgrahaṇam . \\
\noindent \textbf{(P\_3,3.20.1)} atha api kriyate evam api na doṣaḥ . \\
\noindent \textbf{(P\_3,3.20.1)} ajau api sañjñāyām eva . \\
\noindent \textbf{(P\_3,3.20.1)} yathājātīyakaḥ utsargaḥ tathājātīyakena apavādena bhavitavyam . \\
\noindent \textbf{(P\_3,3.20.1)} uttarārtham tarhi . \\
\noindent \textbf{(P\_3,3.20.1)} iṅaḥ ca sarvebhyaḥ api yathā syāt . \\
\noindent \textbf{(P\_3,3.20.1)} nanu ca ayam iṅ ekaḥ eva vaṇṭaraṇḍākalpaḥ . \\
\noindent \textbf{(P\_3,3.20.1)} sarveṣu sādhaneṣu yathā syāt . \\
\noindent \textbf{(P\_3,3.20.1)} upetya adhīyate tasmāt adhyāyaḥ . \\
\noindent \textbf{(P\_3,3.20.1)} adhīyate tasmin adhyāyaḥ . \\
\noindent \textbf{(P\_3,3.20.1)} adhyāyanyāyāyodyāvasaṃhārāvāyāḥ ca iti etat nipātanam na kartavyam bhavati . \\
\noindent \textbf{(P\_3,3.20.1)} etat api na asti prayojanam . \\
\noindent \textbf{(P\_3,3.20.1)} kriyate nyāse eva . \\
\noindent \textbf{(P\_3,3.20.1)} uttarāṛtham eva tarhi vaktavyam . \\
\noindent \textbf{(P\_3,3.20.1)} karmavyatihāre ṇac striyām iti sarvebhyaḥ yathā syāt . \\
\noindent \textbf{(P\_3,3.20.1)} etat api na asti prayojanam . \\
\noindent \textbf{(P\_3,3.20.1)} vakṣyati etat . \\
\noindent \textbf{(P\_3,3.20.1)} karmavyatihāre strīgrahaṇam vyatipākārtham . \\
\noindent \textbf{(P\_3,3.20.1)} pṛthak grahaṇam bādhakabādhanārtham . \\
\noindent \textbf{(P\_3,3.20.1)} vyāvacorīvyāvacarcyartham . \\
\noindent \textbf{(P\_3,3.20.1)} tatra vyatīkṣādiṣu doṣaḥ . \\
\noindent \textbf{(P\_3,3.20.1)} siddham tu prakṛte strīgrahaṇe ṇajgrahaṇam ṇijgrahaṇam ca iti . \\
\noindent \textbf{(P\_3,3.20.1)} uttarārtham tarhi abhividhau bhāve inuṇ sarvebhyaḥ yathā syāt . \\
\noindent \textbf{(P\_3,3.20.1)} sāṃrāviṇam . \\
\noindent \textbf{(P\_3,3.20.1)} etat api na asti prayojanam . \\
\noindent \textbf{(P\_3,3.20.1)} vakṣyati etat . \\
\noindent \textbf{(P\_3,3.20.1)} abhividhau bhāvagrahaṇam napuṃsake ktādinivṛttyartham . \\
\noindent \textbf{(P\_3,3.20.1)} pṛthak grahaṇam bādhakabādhārtham . \\
\noindent \textbf{(P\_3,3.20.1)} na tu lyuṭ iti . \\
\noindent \textbf{(P\_3,3.20.1)} idam tarhi prayojanam . \\
\noindent \textbf{(P\_3,3.20.1)} prakṛtyāśrayaḥ yaḥ apavādaḥ tasya bādhanam yathā syāt . \\
\noindent \textbf{(P\_3,3.20.1)} arthāśrayaḥ yaḥ apavādaḥ tasya bādhanam mā bhūt . \\
\noindent \textbf{(P\_3,3.20.1)} ekā tilocchritiḥ dve sṛtī iti . \\
\noindent \textbf{(P\_3,3.20.1)} ghañanukramaṇam ajabviṣaye . \\
\noindent \textbf{(P\_3,3.20.1)} avacane hi strīpratyayānām api apavādavijñānam iti codayiṣyati . \\
\noindent \textbf{(P\_3,3.20.1)} tat na vaktavyam bhavati . \\
\noindent \textbf{(P\_3,3.20.2)} ghañanukramaṇam ajabviṣaye . \\
\noindent \textbf{(P\_3,3.20.2)} ghañanukramaṇam ajabviṣaye iti vaktavyam . \\
\noindent \textbf{(P\_3,3.20.2)} avacane hi strīpratyayānām api apavādavijñānam . \\
\noindent \textbf{(P\_3,3.20.2)} anucyamāne hi etasmin strīpratyayānām api apavādaḥ ayam vijñayeta . \\
\noindent \textbf{(P\_3,3.20.2)} ekā tilocchritiḥ dve sṛtī iti . dārajārau kartari ṇiluk ca . \\
\noindent \textbf{(P\_3,3.20.2)} dārajārau kartari vaktavyau ṇiluk ca vaktavyaḥ . \\
\noindent \textbf{(P\_3,3.20.2)} dārayanti iti dārāḥ . \\
\noindent \textbf{(P\_3,3.20.2)} jarayanti iti jārāḥ . \\
\noindent \textbf{(P\_3,3.20.2)} karaṇe vā . \\
\noindent \textbf{(P\_3,3.20.2)} karaṇe vā vaktavyau . \\
\noindent \textbf{(P\_3,3.20.2)} dīryate taiḥ dārāḥ . \\
\noindent \textbf{(P\_3,3.20.2)} jīryanti taiḥ jārāḥ . \\
\noindent \textbf{(P\_3,3.21)} iṅaḥ ca iti apādāne striyām upasaṅkhyānam tadantāt ca vā ṅīṣ . \\
\noindent \textbf{(P\_3,3.21)} iṅaḥ ca iti atra apādāne striyām upasaṅkhyānam kartavyam tadantāt ca vā ṅīṣ vaktavyaḥ . \\
\noindent \textbf{(P\_3,3.21)} upetya adhīyate tasyāḥ uapādhyāyī upādhyāyā . \\
\noindent \textbf{(P\_3,3.21)} śṝ vāyuvarṇanivṛteṣu . \\
\noindent \textbf{(P\_3,3.21)} śṝ iti etasmāt vāyuvarṇanivṛteṣu ghañ vaktavyaḥ . \\
\noindent \textbf{(P\_3,3.21)} śāraḥ vāyuḥ . \\
\noindent \textbf{(P\_3,3.21)} śāraḥ varṇaḥ . \\
\noindent \textbf{(P\_3,3.21)} gauḥ iva akṛtanīśāraḥ prāyeṇa śiśire kṛśaḥ . \\
\noindent \textbf{(P\_3,3.36)} sami muṣṭau iti anarthakam vacanam parimāṇākhyāyām iti siddhatvāt . \\
\noindent \textbf{(P\_3,3.36)} sami muṣṭau iti etat vacanam anarthakam . \\
\noindent \textbf{(P\_3,3.36)} kim kāraṇam . \\
\noindent \textbf{(P\_3,3.36)} parimāṇākhyāyām iti siddhatvāt . \\
\noindent \textbf{(P\_3,3.36)} parimāṇākhyāyām iti eva siddham . \\
\noindent \textbf{(P\_3,3.36)} aparimāṇārtham tu . \\
\noindent \textbf{(P\_3,3.36)} aparimāṇārtham tu ayam ārambhaḥ . \\
\noindent \textbf{(P\_3,3.36)} mallasya saṅgrāhaḥ muṣṭikasya saṅgāhaḥ iti . \\
\noindent \textbf{(P\_3,3.36)} udgrābhinigrābhau ca chandasi srugudyamananipātanayoḥ . \\
\noindent \textbf{(P\_3,3.36)} udgrābhaḥ nibrābhaḥ iti imau śabdau chandasi vaktavyau srugudyamananipātanayoḥ . \\
\noindent \textbf{(P\_3,3.36)} udgrābham ca nigrābham ca brahma devāḥ avīvṛdhan . \\
\noindent \textbf{(P\_3,3.43)} strīgrahaṇam kimartham . \\
\noindent \textbf{(P\_3,3.43)} karmavyatihāre strīgrahaṇam vyatipākārtham . karmavyatihāre strīgrahaṇam kriyate vyatipākārtham . \\
\noindent \textbf{(P\_3,3.43)} iha mā bhūt . \\
\noindent \textbf{(P\_3,3.43)} vyatipākaḥ vartate iti . \\
\noindent \textbf{(P\_3,3.43)} atha kimartham pṛthak grahaṇam . \\
\noindent \textbf{(P\_3,3.43)} pṛthak grahaṇam bādhakabādhanārtham . \\
\noindent \textbf{(P\_3,3.43)} pṛthak grahaṇam kriyate bādhakabādhanārtham . \\
\noindent \textbf{(P\_3,3.43)} ye tasya bādhakāḥ tadbādhanārtham . \\
\noindent \textbf{(P\_3,3.43)} kim prayojanam . \\
\noindent \textbf{(P\_3,3.43)} vyāvacorīvyāvacarcyartham . vyāvacorī vartate . \\
\noindent \textbf{(P\_3,3.43)} vyāvacarcī vartate . \\
\noindent \textbf{(P\_3,3.43)} tatra vyatīkṣādiṣu doṣaḥ . \\
\noindent \textbf{(P\_3,3.43)} tatra vyatīkṣādiṣu doṣaḥ bhavati . \\
\noindent \textbf{(P\_3,3.43)} vyatīkṣā vartate . \\
\noindent \textbf{(P\_3,3.43)} vyatīhā vartate . \\
\noindent \textbf{(P\_3,3.43)} siddham tu prakṛte strīgrahaṇe ṇajgrahaṇam ṇijgrahaṇam ca . \\
\noindent \textbf{(P\_3,3.43)} siddham etat . \\
\noindent \textbf{(P\_3,3.43)} katham . \\
\noindent \textbf{(P\_3,3.43)} prakṛte eva strīgrahaṇe ayam yogaḥ kartavyaḥ . \\
\noindent \textbf{(P\_3,3.43)} striyām ktin . \\
\noindent \textbf{(P\_3,3.43)} tataḥ karmavyatihāre ṇac . \\
\noindent \textbf{(P\_3,3.43)} tataḥ ṇicaḥ . \\
\noindent \textbf{(P\_3,3.44)} bhāvagrahaṇam kimartham . \\
\noindent \textbf{(P\_3,3.44)} abhividhau bhāvagrahaṇam napuṃsake ktādinivṛttyartham . \\
\noindent \textbf{(P\_3,3.44)} abhividhau bhāvagrahaṇam kriyate napuṃsake ktādinivṛttyartham . \\
\noindent \textbf{(P\_3,3.44)} napuṃsakaliṅge ktādayaḥ mā bhūvan iti . \\
\noindent \textbf{(P\_3,3.44)} atha kimartham pṛthak grahaṇam . \\
\noindent \textbf{(P\_3,3.44)} pṛthak grahaṇam bādhakabādhārtham . \\
\noindent \textbf{(P\_3,3.44)} pṛthak grahaṇam kriyate bādhakabādhārtham : ye tasya bādhakāḥ tadbādhanārtham . \\
\noindent \textbf{(P\_3,3.44)} na tu lyuṭaḥ . lyuṭaḥ tu bādhanam na iṣyate . \\
\noindent \textbf{(P\_3,3.44)} saṅkūṭanam iti eva bhavati . \\
\noindent \textbf{(P\_3,3.56)} ajvidhau bhayasya upasaṅkhyānam . \\
\noindent \textbf{(P\_3,3.56)} ajvidhau bhayasya upasaṅkhyānam kartavyam . \\
\noindent \textbf{(P\_3,3.56)} bhayam . \\
\noindent \textbf{(P\_3,3.56)} atyalpam idam ucyate : bhayasya iti . \\
\noindent \textbf{(P\_3,3.56)} bhayādīnām iti vaktavyam iha api yathā syāt : bhayam varṣam . \\
\noindent \textbf{(P\_3,3.56)} kim prayojanam . \\
\noindent \textbf{(P\_3,3.56)} napuṃsake ktādinivṛttyartham . \\
\noindent \textbf{(P\_3,3.56)} napuṃsakaliṅge ktādayaḥ mā bhūvan iti . \\
\noindent \textbf{(P\_3,3.56)} kalpādibhyaḥ ca pratiṣedhaḥ . \\
\noindent \textbf{(P\_3,3.56)} kalpādibhyaḥ ca pratiṣedhaḥ vaktavyaḥ . \\
\noindent \textbf{(P\_3,3.56)} kalpaḥ arthaḥ mantraḥ . \\
\noindent \textbf{(P\_3,3.56)} javasavau chandasi . \\
\noindent \textbf{(P\_3,3.56)} javasavau chandasi vaktavyau . \\
\noindent \textbf{(P\_3,3.56)} ūrvoḥ astu me javaḥ . \\
\noindent \textbf{(P\_3,3.56)} ayam me pañcaudanaḥ savaḥ . \\
\noindent \textbf{(P\_3,3.58.1)} kimartham niṣpūrvāt cinoteḥ ap vidhīyate na acā eva siddham . \\
\noindent \textbf{(P\_3,3.58.1)} na hi asti viśeṣaḥ niṣpūrvāt cinoteḥ apaḥ vā acaḥ vā . \\
\noindent \textbf{(P\_3,3.58.1)} tat eva rūpam saḥ eva svaraḥ . \\
\noindent \textbf{(P\_3,3.58.1)} na sidhyati . \\
\noindent \textbf{(P\_3,3.58.1)} hastādāne ceḥ ghañ prāptaḥ . \\
\noindent \textbf{(P\_3,3.58.1)} tadbādhanārtham . \\
\noindent \textbf{(P\_3,3.58.1)} ataḥ uttaram paṭhati . \\
\noindent \textbf{(P\_3,3.58.1)} abvidhau niścigrahaṇam anarthakam steyasya ghañvidhau pratiṣedhāt . \\
\noindent \textbf{(P\_3,3.58.1)} abvidhau niścigrahaṇam anarthakam . \\
\noindent \textbf{(P\_3,3.58.1)} kim kāraṇam . \\
\noindent \textbf{(P\_3,3.58.1)} steyasya ghañvidhau pratiṣedhāt . \\
\noindent \textbf{(P\_3,3.58.1)} steyasya ghañvidhau pratiṣedhaḥ ucyate . \\
\noindent \textbf{(P\_3,3.58.1)} niṣpūrvaḥ cinotiḥ steye vartate . \\
\noindent \textbf{(P\_3,3.58.1)} asteyārtham tarhi idam vaktavyam . \\
\noindent \textbf{(P\_3,3.58.1)} niṣpūrvāt cinoteḥ asteye yathā syāt . \\
\noindent \textbf{(P\_3,3.58.1)} asteyārtham iti cet na aniṣṭatvāt . \\
\noindent \textbf{(P\_3,3.58.1)} asteyārtham iti cet tat na . \\
\noindent \textbf{(P\_3,3.58.1)} kim kāraṇam . \\
\noindent \textbf{(P\_3,3.58.1)} aniṣṭatvāt . \\
\noindent \textbf{(P\_3,3.58.1)} na niṣpūrvāt cinoteḥ asteye ap iṣyate . \\
\noindent \textbf{(P\_3,3.58.1)} kim tarhi ghañ eva iṣyate . \\
\noindent \textbf{(P\_3,3.58.1)} evam tarhi siddhe sati yat niṣpūrvāt cinoteḥ apam śāsti tat jñāpayati ācāryaḥ yat tat antaḥ thāthaghañktājabitṛkāṇām iti tat niṣpūrvāt cinoteḥ na bhavati iti . \\
\noindent \textbf{(P\_3,3.58.1)} kim etasya jñāpane prayojanam . \\
\noindent \textbf{(P\_3,3.58.1)} niścayaḥ . \\
\noindent \textbf{(P\_3,3.58.1)} eṣaḥ svaraḥ siddhaḥ bhavati . \\
\noindent \textbf{(P\_3,3.58.2)} vaśiraṇyoḥ ca upasaṅkhyānam . \\
\noindent \textbf{(P\_3,3.58.2)} vaśiraṇyoḥ ca upasaṅkhyānam . \\
\noindent \textbf{(P\_3,3.58.2)} saḥ vaśam saindhavam . \\
\noindent \textbf{(P\_3,3.58.2)} dhanañjayaḥ raṇe raṇe . \\
\noindent \textbf{(P\_3,3.58.2)} ghañarthe kavidhānam sthāsnāpāvyadhihaniyudhyartham . \\
\noindent \textbf{(P\_3,3.58.2)} ghañarthe kaḥ vidheyaḥ . \\
\noindent \textbf{(P\_3,3.58.2)} kim prayojanam . \\
\noindent \textbf{(P\_3,3.58.2)} sthāsnāpāvyadhihaniyudhyartham . \\
\noindent \textbf{(P\_3,3.58.2)} sthā . \\
\noindent \textbf{(P\_3,3.58.2)} pratiṣṭhante asmin dhānyāni iti prasthaḥ . \\
\noindent \textbf{(P\_3,3.58.2)} prasthe himavataḥ śrṅge . \\
\noindent \textbf{(P\_3,3.58.2)} sthā . \\
\noindent \textbf{(P\_3,3.58.2)} snā . \\
\noindent \textbf{(P\_3,3.58.2)} prasnānti tasmin iti prasnaḥ . \\
\noindent \textbf{(P\_3,3.58.2)} snā . \\
\noindent \textbf{(P\_3,3.58.2)} pā . \\
\noindent \textbf{(P\_3,3.58.2)} prapibanti asyām iti prapā . \\
\noindent \textbf{(P\_3,3.58.2)} pā . \\
\noindent \textbf{(P\_3,3.58.2)} vyadhi . \\
\noindent \textbf{(P\_3,3.58.2)} āvidhyanti tena āvidham . \\
\noindent \textbf{(P\_3,3.58.2)} vyadhi . \\
\noindent \textbf{(P\_3,3.58.2)} hani . \\
\noindent \textbf{(P\_3,3.58.2)} vighnanti tasmin manāṃsi vighnaḥ . \\
\noindent \textbf{(P\_3,3.58.2)} hani . \\
\noindent \textbf{(P\_3,3.58.2)} yudhi . \\
\noindent \textbf{(P\_3,3.58.2)} āyudhyante tena āyudham . \\
\noindent \textbf{(P\_3,3.83)} kasmāt ayam kaḥ vidhīyate . \\
\noindent \textbf{(P\_3,3.83)} hanteḥ iti āha . \\
\noindent \textbf{(P\_3,3.83)} tat hantigrahaṇam kartavyam . \\
\noindent \textbf{(P\_3,3.83)} na kartavyam . \\
\noindent \textbf{(P\_3,3.83)} prakṛtam anuvartate . \\
\noindent \textbf{(P\_3,3.83)} kva prakṛtam . \\
\noindent \textbf{(P\_3,3.83)} hanaḥ ca vadhaḥ . \\
\noindent \textbf{(P\_3,3.83)} tat vai anekena nipātanena vyavacchinnam na śakyam anuvartayitum . \\
\noindent \textbf{(P\_3,3.83)} na etāni nipātanāni . \\
\noindent \textbf{(P\_3,3.83)} hanteḥ ete ādeśāḥ . \\
\noindent \textbf{(P\_3,3.83)} yadi ādeśāḥ ghanasvaraḥ na sidhyati . \\
\noindent \textbf{(P\_3,3.83)} ghanaḥ . \\
\noindent \textbf{(P\_3,3.83)} santu tarhi nipātanāni . \\
\noindent \textbf{(P\_3,3.83)} nanu ca uktam tat vai anekena nipātanena vyavacchinnam na śakyam anuvartayitum iti . \\
\noindent \textbf{(P\_3,3.83)} sambandham anuvartiṣyate . \\
\noindent \textbf{(P\_3,3.83)} atha vā punaḥ santu ādeśāḥ . \\
\noindent \textbf{(P\_3,3.83)} nanu ca uktam svaraḥ na sidhyati iti . \\
\noindent \textbf{(P\_3,3.83)} na eṣaḥ doṣaḥ . \\
\noindent \textbf{(P\_3,3.83)} akārāntaḥ ādeśaḥ . \\
\noindent \textbf{(P\_3,3.83)} atha yadā iṣīkayā stambaḥ hanyate katham tatra bhavitavyam . \\
\noindent \textbf{(P\_3,3.83)} ke cid tāvat āhuḥ . \\
\noindent \textbf{(P\_3,3.83)} stambaghnā iti bhavitavyam . \\
\noindent \textbf{(P\_3,3.83)} apare āhuḥ : stambahetiḥ iti bhavitavyam . \\
\noindent \textbf{(P\_3,3.83)} ūtiyūtijūtisātihetikīrtayaḥ ca iti nipātanam iti . \\
\noindent \textbf{(P\_3,3.83)} apare āhuḥ . \\
\noindent \textbf{(P\_3,3.83)} stambahananīiti bhavitavyam iti . \\
\noindent \textbf{(P\_3,3.83)} vakṣyati etat . \\
\noindent \textbf{(P\_3,3.83)} ajabbhyām strīīkhalanāḥ . \\
\noindent \textbf{(P\_3,3.83)} striyāḥ khalanau vipratiṣedhena iti . \\
\noindent \textbf{(P\_3,3.90)} yajādibhyaḥ nasya ṅittve samprasāraṇapratiṣedhaḥ . \\
\noindent \textbf{(P\_3,3.90)} yajādibhyaḥ nasya ṅittve samprasāraṇapratiṣedhaḥ vaktavyaḥ . \\
\noindent \textbf{(P\_3,3.90)} praśnaḥ iti . \\
\noindent \textbf{(P\_3,3.90)} evan tarhi āṅit kariṣyate . \\
\noindent \textbf{(P\_3,3.90)} aṅiti guṇapratiṣedhaḥ . \\
\noindent \textbf{(P\_3,3.90)} yadi aṅit guṇapratiṣedhaḥ vaktavyaḥ . \\
\noindent \textbf{(P\_3,3.90)} viśnaḥ iti . \\
\noindent \textbf{(P\_3,3.90)} sūtram ca bhidyate . \\
\noindent \textbf{(P\_3,3.90)} yathānyāsam eva astu . \\
\noindent \textbf{(P\_3,3.90)} nanu ca uktam yajādibhyaḥ nasya ṅittve samprasāraṇapratiṣedhaḥ iti . \\
\noindent \textbf{(P\_3,3.90)} na eṣaḥ doṣaḥ . \\
\noindent \textbf{(P\_3,3.90)} nipātanāt etat siddham . \\
\noindent \textbf{(P\_3,3.90)} kim nipātanam . \\
\noindent \textbf{(P\_3,3.90)} praśne ca āsannakāle iti . \\
\noindent \textbf{(P\_3,3.94)} striyām ktin ābādibhyaḥ ca . \\
\noindent \textbf{(P\_3,3.94)} striyām ktin iti atra ābādibhyaḥ ca iti vaktavyam . \\
\noindent \textbf{(P\_3,3.94)} āptiḥ rāddhiḥ dīptiḥ . \\
\noindent \textbf{(P\_3,3.94)} niṣṭhāyām vā seṭaḥ akāravacanāt siddham . \\
\noindent \textbf{(P\_3,3.94)} atha vā niṣṭhāyām seṭaḥ akāraḥ bhavati iti vaktavyam . \\
\noindent \textbf{(P\_3,3.94)} yadi niṣṭhāyām seṭaḥ akāraḥ bhavati iti ucyate sraṃsā dhvaṃsā iti na sidhyati . \\
\noindent \textbf{(P\_3,3.94)} srastiḥ dhvastiḥ iti prāpnoti . \\
\noindent \textbf{(P\_3,3.94)} kim punaḥ idam parigaṇanam trayaḥ eva ābādayaḥ āhosvit udāharaṇamātram . \\
\noindent \textbf{(P\_3,3.94)} kim ca ataḥ . \\
\noindent \textbf{(P\_3,3.94)} yadi parigaṇanam bhedaḥ bhavati . \\
\noindent \textbf{(P\_3,3.94)} atha udāharaṇamātram na asti bhedaḥ . \\
\noindent \textbf{(P\_3,3.94)} srasti dhvastiḥ iti eva bhavitavyam . \\
\noindent \textbf{(P\_3,3.95)} sthādibhyaḥ sarvāpavādaprasaṅgaḥ . \\
\noindent \textbf{(P\_3,3.95)} sthādibhyaḥ sarvāpavādaḥ ktin prāpnoti . \\
\noindent \textbf{(P\_3,3.95)} saḥ yathā eva aṅam bādhate evam ṇvuliñau api bādheta . \\
\noindent \textbf{(P\_3,3.95)} kām tvam sthāyikām asthāḥ . \\
\noindent \textbf{(P\_3,3.95)} kām sthāyim . \\
\noindent \textbf{(P\_3,3.95)} siddham tu aṅvidhāne sthādipratiṣedhāt . \\
\noindent \textbf{(P\_3,3.95)} siddham etat . \\
\noindent \textbf{(P\_3,3.95)} katham . \\
\noindent \textbf{(P\_3,3.95)} aṅvidhāne eva sthādipratiṣedhaḥ vaktavyaḥ . \\
\noindent \textbf{(P\_3,3.95)} pratiṣiddhe tasmin ktin eva bhaviṣyati . \\
\noindent \textbf{(P\_3,3.95)} sidhyati . \\
\noindent \textbf{(P\_3,3.95)} sūtram tarhi bhidyate . \\
\noindent \textbf{(P\_3,3.95)} yathānyāsam eva astu . \\
\noindent \textbf{(P\_3,3.95)} nanu ca uktam sthādibhyaḥ sarvāpavādaprasaṅgaḥ iti . \\
\noindent \textbf{(P\_3,3.95)} na eṣaḥ doṣaḥ . \\
\noindent \textbf{(P\_3,3.95)} purastāt apavādāḥ anantarān vidhīn bādhante iti evam ayam striyām ktin aṅam bādhiṣyate . \\
\noindent \textbf{(P\_3,3.95)} ṇvuliñau na bādhiṣyate . \\
\noindent \textbf{(P\_3,3.95)} śrutijiṣistubhyaḥ karaṇe . \\
\noindent \textbf{(P\_3,3.95)} śrutijiṣistubhyaḥ karaṇe ktin vaktavyaḥ . \\
\noindent \textbf{(P\_3,3.95)} śrūyate anayā śrutiḥ . \\
\noindent \textbf{(P\_3,3.95)} ijyate anayā iṣṭiḥ . \\
\noindent \textbf{(P\_3,3.95)} iṣyate anayā iṣṭiḥ . \\
\noindent \textbf{(P\_3,3.95)} stūyate anayā stutiḥ . \\
\noindent \textbf{(P\_3,3.95)} glājyāhābhyaḥ niḥ . \\
\noindent \textbf{(P\_3,3.95)} glājyāhābhyaḥ niḥ vaktavyaḥ . \\
\noindent \textbf{(P\_3,3.95)} glāniḥ jyāniḥ hāniḥ . \\
\noindent \textbf{(P\_3,3.98)} kyabvidhiḥ adhikaraṇe ca . \\
\noindent \textbf{(P\_3,3.98)} kyabvidhiḥ adhikaraṇe ca iti vaktavyam . \\
\noindent \textbf{(P\_3,3.98)} samajanti tasyām samajyā . \\
\noindent \textbf{(P\_3,3.100)} kṛñaḥ śa ca iti vāvacanam ktinartham . \\
\noindent \textbf{(P\_3,3.100)} kṛñaḥ śa ca iti vāvacanam kartavyam ktin api yathā syāt . \\
\noindent \textbf{(P\_3,3.100)} kṛtiḥ . \\
\noindent \textbf{(P\_3,3.102)} kim nipātyate . \\
\noindent \textbf{(P\_3,3.102)} iṣeḥ śe yagabhāvaḥ . \\
\noindent \textbf{(P\_3,3.102)} atyalpam idam ucyate icchā iti . \\
\noindent \textbf{(P\_3,3.102)} icchāparicaryāparisaryāmṛgayāṭāṭyānām nipātanam kartavyam . \\
\noindent \textbf{(P\_3,3.102)} jāgarteḥ akāraḥ vā . \\
\noindent \textbf{(P\_3,3.102)} jāgarya jāgarā . \\
\noindent \textbf{(P\_3,3.104)} bhidā vidāraṇe . \\
\noindent \textbf{(P\_3,3.104)} bhidā vidāraṇe iti vaktavyam . \\
\noindent \textbf{(P\_3,3.104)} bhittiḥ anyā . \\
\noindent \textbf{(P\_3,3.104)} chidhā dvaidhīkaraṇe . \\
\noindent \textbf{(P\_3,3.104)} chidhā dvaidhīkaraṇe iti vaktavyam . \\
\noindent \textbf{(P\_3,3.104)} chittiḥ anyā . \\
\noindent \textbf{(P\_3,3.104)} ārā śastryām . \\
\noindent \textbf{(P\_3,3.104)} ārā śastryām iti vaktavyam . \\
\noindent \textbf{(P\_3,3.104)} ārtiḥ anyā . \\
\noindent \textbf{(P\_3,3.104)} dhārā prapāte . \\
\noindent \textbf{(P\_3,3.104)} dhārā prapāte iti vaktavyam . \\
\noindent \textbf{(P\_3,3.104)} dhṛtiḥ anyā . \\
\noindent \textbf{(P\_3,3.104)} guhā giryoṣadhyoḥ . guhā giryoṣadhyoḥ iti vaktavyam . \\
\noindent \textbf{(P\_3,3.104)} gūḍhiḥ anyā \\
\noindent \textbf{(P\_3,3.107.1)} kimarthaḥ cakāraḥ . \\
\noindent \textbf{(P\_3,3.107.1)} svarārthaḥ . \\
\noindent \textbf{(P\_3,3.107.1)} citaḥ antaḥ udāttaḥ bhavati iti antodāttatvam yathā syāt . \\
\noindent \textbf{(P\_3,3.107.1)} na etat asti prayojanam . \\
\noindent \textbf{(P\_3,3.107.1)} udāttaḥ iti vartate bhūvīrāḥ udāttaḥ iti . \\
\noindent \textbf{(P\_3,3.107.1)} yadi udāttaḥ iti vartate vajayajoḥ bhāve kyap kimarthaḥ pakāraḥ . \\
\noindent \textbf{(P\_3,3.107.1)} tugarthaḥ . \\
\noindent \textbf{(P\_3,3.107.1)} hrasvasya piti kṛti tuk iti . \\
\noindent \textbf{(P\_3,3.107.1)} udāttaḥ iti vartate . \\
\noindent \textbf{(P\_3,3.107.1)} evam api kutaḥ etat tadantasya udāttatvam bhaviṣyati na punaḥ ādeḥ iti . \\
\noindent \textbf{(P\_3,3.107.1)} udāttaḥ iti anuvartanasāmarthyāt yasya aprāptaḥ svaraḥ tasya bhavati . \\
\noindent \textbf{(P\_3,3.107.1)} kasya ca aprāptaḥ . \\
\noindent \textbf{(P\_3,3.107.1)} antyasya . \\
\noindent \textbf{(P\_3,3.107.1)} sāmānyagrahaṇāvighātārthaḥ tarhi . \\
\noindent \textbf{(P\_3,3.107.1)} kva sāmānyagrahaṇāvighātārthena arthaḥ . \\
\noindent \textbf{(P\_3,3.107.1)} yuvoḥ anākau iti . \\
\noindent \textbf{(P\_3,3.107.1)} etat api na asti prayojanam . \\
\noindent \textbf{(P\_3,3.107.1)} vakṣyati etat . \\
\noindent \textbf{(P\_3,3.107.1)} siddham tu yuvoḥ anunāsikavacanāt iti . \\
\noindent \textbf{(P\_3,3.107.2)} yucprakaraṇe ghaṭṭivandividhibhyaḥ ca upasaṅkhyānam . \\
\noindent \textbf{(P\_3,3.107.2)} yucprakaraṇe ghaṭṭivandividhibhyaḥ ca upasaṅkhyānam kartavyam . \\
\noindent \textbf{(P\_3,3.107.2)} ghaṭṭanā vandanā vedanā . \\
\noindent \textbf{(P\_3,3.107.2)} iṣeḥ anicchārthasya . \\
\noindent \textbf{(P\_3,3.107.2)} iṣeḥ anicchārthasya iti vaktavyam . \\
\noindent \textbf{(P\_3,3.107.2)} anviṣyate anveṣaṇā . \\
\noindent \textbf{(P\_3,3.107.2)} pareḥ vā . \\
\noindent \textbf{(P\_3,3.107.2)} pareḥ vā iti vaktavyam . \\
\noindent \textbf{(P\_3,3.107.2)} anyām parīṣṭim cara . \\
\noindent \textbf{(P\_3,3.107.2)} anyām paryeṣaṇām cara . \\
\noindent \textbf{(P\_3,3.108)} dhātvartanirdeśe ṇvul . \\
\noindent \textbf{(P\_3,3.108)} dhātvartanirdeśe ṇvul vaktavyaḥ . \\
\noindent \textbf{(P\_3,3.108)} kā nāma āsikā anyeṣu īhamāneṣu . \\
\noindent \textbf{(P\_3,3.108)} kā nām śāyikā anyeṣu adhīyāneṣu . \\
\noindent \textbf{(P\_3,3.108)} ikśtipau dhātunirdeśe . \\
\noindent \textbf{(P\_3,3.108)} ikśtipau iti etau pratyayau dhātunirdeśe vaktavyau . \\
\noindent \textbf{(P\_3,3.108)} paceḥ brūhi . \\
\noindent \textbf{(P\_3,3.108)} pacateḥ brūhi . \\
\noindent \textbf{(P\_3,3.108)} varṇāt kāraḥ . \\
\noindent \textbf{(P\_3,3.108)} varṇāt kārapratyayaḥ vaktavyaḥ . \\
\noindent \textbf{(P\_3,3.108)} akāraḥ ikāraḥ . \\
\noindent \textbf{(P\_3,3.108)} rāt iphaḥ . \\
\noindent \textbf{(P\_3,3.108)} rāt iphaḥ vaktavyaḥ . \\
\noindent \textbf{(P\_3,3.108)} rephaḥ . \\
\noindent \textbf{(P\_3,3.108)} matvarthāt chaḥ . \\
\noindent \textbf{(P\_3,3.108)} matvarthāt chaḥ vaktavyaḥ . \\
\noindent \textbf{(P\_3,3.108)} matvarthīyaḥ . \\
\noindent \textbf{(P\_3,3.108)} iṇ ajādibhyaḥ . \\
\noindent \textbf{(P\_3,3.108)} iṇ ajādibhyaḥ vaktavyaḥ . \\
\noindent \textbf{(P\_3,3.108)} ājiḥ ātiḥ ādiḥ . \\
\noindent \textbf{(P\_3,3.108)} iñ vapādibhyaḥ . \\
\noindent \textbf{(P\_3,3.108)} iñ vapādibhyaḥ vaktavyaḥ . \\
\noindent \textbf{(P\_3,3.108)} vāpiḥ vāsiḥ vādiḥ . \\
\noindent \textbf{(P\_3,3.108)} ik kṛṣyādibhyaḥ . \\
\noindent \textbf{(P\_3,3.108)} ik kṛṣyādibhyaḥ vaktavyaḥ . \\
\noindent \textbf{(P\_3,3.108)} kṛṣiḥ kiriḥ giriḥ . \\
\noindent \textbf{(P\_3,3.108)} sampadādibhyaḥ kvip . \\
\noindent \textbf{(P\_3,3.108)} sampadādibhyaḥ kvip vaktavyaḥ . \\
\noindent \textbf{(P\_3,3.108)} sampat vipat pratipat āpat pariṣat . \\
\noindent \textbf{(P\_3,3.108)} KA\_II,154.15-155.10 Ro\_III,34341-342 {30/30} \\
\noindent \textbf{(P\_3,3.113)} kṛtaḥ bahulam iti vaktavyam pādahārakādyartham . \\
\noindent \textbf{(P\_3,3.113)} pādābhyām hriyate pādahārakaḥ . \\
\noindent \textbf{(P\_3,3.113)} gale copyate galecopakaḥ . \\
\noindent \textbf{(P\_3,3.113)} śvaḥ agnīn ādhāsyamānena . \\
\noindent \textbf{(P\_3,3.113)} śvaḥ somena yakṣyamāṇena . \\
\noindent \textbf{(P\_3,3.119)} gocarādīnām agrahaṇam prāyavacanāt yathā kaṣaḥ nikaṣaḥ iti . \\
\noindent \textbf{(P\_3,3.119)} gocarādīnām grahaṇam śakyam akartum . \\
\noindent \textbf{(P\_3,3.119)} ghañ kasmāt na bhavati . \\
\noindent \textbf{(P\_3,3.119)} prāyavacanāt yathā kaṣaḥ nikaṣaḥ iti prāyavacanāt ghañ na bhavati . \\
\noindent \textbf{(P\_3,3.121)} ghañvidhau avahārādhārāvāyānām upasaṅkhyānam . \\
\noindent \textbf{(P\_3,3.121)} ghañvidhau avahārādhārāvāyānām upasaṅkhyānam kartavyam . \\
\noindent \textbf{(P\_3,3.121)} avahriyante asmin avahāraḥ . \\
\noindent \textbf{(P\_3,3.121)} ādhriyante asmin ādhāraḥ . \\
\noindent \textbf{(P\_3,3.121)} etya etasmin vayanti āvāyaḥ . \\
\noindent \textbf{(P\_3,3.123)} kimartham idam ucyate na halaḥ ca iti eva siddham . \\
\noindent \textbf{(P\_3,3.123)} anudake it vakṣyāmi iti . \\
\noindent \textbf{(P\_3,3.123)} iha mā bhūt . \\
\noindent \textbf{(P\_3,3.123)} udakodañcanaḥ . \\
\noindent \textbf{(P\_3,3.123)} udaṅkaḥ anudakagrahaṇānarthakyam ca prāyavacanāt yathā godohanaḥ prasādhanaḥ iti . \\
\noindent \textbf{(P\_3,3.123)} udaṅkaḥ anudakagrahaṇam ca anarthakam . \\
\noindent \textbf{(P\_3,3.123)} ghañ kasmāt na bhavati . \\
\noindent \textbf{(P\_3,3.123)} prāyavacanāt yathā godohanaḥ prasādhanaḥ iti . \\
\noindent \textbf{(P\_3,3.125)} ḍaḥ vaktavyaḥ . \\
\noindent \textbf{(P\_3,3.125)} ākhaḥ . \\
\noindent \textbf{(P\_3,3.125)} ḍaraḥ vaktavyaḥ . \\
\noindent \textbf{(P\_3,3.125)} ākharaḥ . \\
\noindent \textbf{(P\_3,3.125)} ikaḥ vaktavyaḥ . \\
\noindent \textbf{(P\_3,3.125)} ākhanikaḥ . \\
\noindent \textbf{(P\_3,3.125)} ikavakaḥ vaktavyaḥ . \\
\noindent \textbf{(P\_3,3.125)} ākhanikavakaḥ . \\
\noindent \textbf{(P\_3,3.126)} ajabbhyām strīkhalanāḥ . \\
\noindent \textbf{(P\_3,3.126)} ajabbhyām strīkhalanāḥ bhavanti vipratiṣedhena . \\
\noindent \textbf{(P\_3,3.126)} ajapoḥ avakāśaḥ cayaḥ lavaḥ . \\
\noindent \textbf{(P\_3,3.126)} strīpratyayānām avakāśaḥ kṛtiḥ hṛtiḥ . \\
\noindent \textbf{(P\_3,3.126)} iha ubhayam prāpnoti . \\
\noindent \textbf{(P\_3,3.126)} citiḥ stutiḥ . \\
\noindent \textbf{(P\_3,3.126)} khalaḥ avakāśaḥ īṣadbhedaḥ subhedaḥ . \\
\noindent \textbf{(P\_3,3.126)} ajapoḥ saḥ eva . \\
\noindent \textbf{(P\_3,3.126)} iha ubhayam prāpnoti . \\
\noindent \textbf{(P\_3,3.126)} īṣaccayaḥ sucayaḥ īṣallavaḥ sulavaḥ . \\
\noindent \textbf{(P\_3,3.126)} anasya avakāśaḥ idhmapravraścanaḥ . \\
\noindent \textbf{(P\_3,3.126)} ajapoḥ saḥ eva . \\
\noindent \textbf{(P\_3,3.126)} iha ubhayam prāpnoti . \\
\noindent \textbf{(P\_3,3.126)} palāśacayanaḥ avilavanaḥ . \\
\noindent \textbf{(P\_3,3.126)} strīkhalanāḥ bhavanti vipratiṣedhena . \\
\noindent \textbf{(P\_3,3.126)} striyāḥ khalanau vipratiṣedhena . \\
\noindent \textbf{(P\_3,3.126)} striyāḥ khalanau bhavataḥ vipratiṣedhena . \\
\noindent \textbf{(P\_3,3.126)} strīpratyayānām avakāśaḥ kṛtiḥ hṛtiḥ . \\
\noindent \textbf{(P\_3,3.126)} khalaḥ avakāśaḥ īṣadbhedaḥ subhedaḥ . \\
\noindent \textbf{(P\_3,3.126)} iha ubhayam prāpnoti . \\
\noindent \textbf{(P\_3,3.126)} īṣadbhedā subhedā . \\
\noindent \textbf{(P\_3,3.126)} anasya avakāśaḥ idhmapravraścanaḥ . \\
\noindent \textbf{(P\_3,3.126)} strīpratyayānām saḥ eva . \\
\noindent \textbf{(P\_3,3.126)} iha ubhayam prāpnoti . \\
\noindent \textbf{(P\_3,3.126)} saktudhānī tilapīḍanī . \\
\noindent \textbf{(P\_3,3.126)} khalanau bhavataḥ vipratiṣedhena . \\
\noindent \textbf{(P\_3,3.127)} khal kartṛkaraṇayoḥ cvyarthayoḥ . \\
\noindent \textbf{(P\_3,3.127)} khal kartṛkaraṇayoḥ cvyarthayoḥ iti vaktavyam . \\
\noindent \textbf{(P\_3,3.127)} anāḍhyena bhavatā īṣadāḍhyena śakyam bhavitum īṣadāḍhyambhavam bhavatā . \\
\noindent \textbf{(P\_3,3.127)} durāḍhyambhavam svāḍhyambhavam . \\
\noindent \textbf{(P\_3,3.127)} kartṛkarmagrahaṇam ca upapadasañjñārtham . kartṛkarmagrahaṇam ca upapadasañjñārtham draṣṭavyam . \\
\noindent \textbf{(P\_3,3.127)} dveṣyam vijānīyāt : abhidheyayoḥ iti . \\
\noindent \textbf{(P\_3,3.127)} tat ācāryaḥ suhṛt bhūtvā anvācaṣṭe : kartṛkarmagrahaṇam ca upapadasañjñārtham iti \\
\noindent \textbf{(P\_3,3.130)} bhāṣāyām śāsiyudhidṛśidhṛṣibhyaḥ yuc . \\
\noindent \textbf{(P\_3,3.130)} bhāṣāyām śāsiyudhidṛśidhṛṣibhyaḥ yuc vaktavyaḥ . \\
\noindent \textbf{(P\_3,3.130)} duḥśāsanaḥ duryodhanaḥ durdarśanaḥ durdharṣaṇaḥ . \\
\noindent \textbf{(P\_3,3.130)} mṛṣeḥ ca iti vaktavyam . \\
\noindent \textbf{(P\_3,3.130)} durmarṣaṇaḥ . \\
\noindent \textbf{(P\_3,3.131)} vatkaraṇam kimartham . \\
\noindent \textbf{(P\_3,3.131)} vartamānasāmīpye vartmānāḥ vā iti iyati ucyamāne vartamāne ye pratyayāḥ vihitāḥ vartamānasāmīpye dhātumātrāt syuḥ . \\
\noindent \textbf{(P\_3,3.131)} vatkaraṇe punaḥ kriyamāṇe na doṣaḥ bhavati . \\
\noindent \textbf{(P\_3,3.131)} yadi ca yābhyaḥ prakṛtibhyaḥ yena viśeṣeṇa vartamāne pratyayāḥ vihitāḥ tābhyaḥ prakṛtibhyaḥ tena eva viśeṣeṇa vartamānasāmīpye bhavanti tataḥ amīvartamānavat kṛtāḥ syuḥ . \\
\noindent \textbf{(P\_3,3.131)} atha hi prakṛtimātrāt vā syuḥ pratyayamātram vā syāt na amīvartamānavat kṛtāḥ syuḥ . \\
\noindent \textbf{(P\_3,3.131)} iha vartamānasāmīpye vartamānavat vā iti uktvā loṭ eva udāhriyate . \\
\noindent \textbf{(P\_3,3.131)} yadi punaḥ vā laṭ bhavati iti eva ucyeta . \\
\noindent \textbf{(P\_3,3.131)} ataḥ uttaram paṭhati . \\
\noindent \textbf{(P\_3,3.131)} vartamānasāmīpye vartamānavadvacanam śatrādyartham . \\
\noindent \textbf{(P\_3,3.131)} vartamānasāmīpye vartamānavadvacanam kriyate śatrādyartham . \\
\noindent \textbf{(P\_3,3.131)} śatrādyarthaḥ ayam ārambhaḥ . \\
\noindent \textbf{(P\_3,3.131)} eṣaḥ asmi pacan . \\
\noindent \textbf{(P\_3,3.131)} eṣaḥ asmi pacamānaḥ iti . \\
\noindent \textbf{(P\_3,3.131)} na etat asti prayojanam . \\
\noindent \textbf{(P\_3,3.131)} laḍādeśau śatṛśānacau . \\
\noindent \textbf{(P\_3,3.131)} tatra vā laṭ bhavati iti eva siddham . \\
\noindent \textbf{(P\_3,3.131)} yau tarhi alaḍādeśau . \\
\noindent \textbf{(P\_3,3.131)} eṣaḥ asmi pavamānaḥ . \\
\noindent \textbf{(P\_3,3.131)} eṣaḥ asmi yajamānaḥ . \\
\noindent \textbf{(P\_3,3.131)} yau ca api laḍādeśau tau api prayojayataḥ . \\
\noindent \textbf{(P\_3,3.131)} vartamānavihitasya laṭaḥ śatṛśānacau ucyete . \\
\noindent \textbf{(P\_3,3.131)} aviśeṣeṇa vihitaḥ ca ayam yogaḥ . \\
\noindent \textbf{(P\_3,3.131)} śatrādyartham iti khalu api ucyate . \\
\noindent \textbf{(P\_3,3.131)} bahavaḥ ca śatrādayaḥ . \\
\noindent \textbf{(P\_3,3.131)} eṣaḥ asmi alaṅkariṣṇuḥ . \\
\noindent \textbf{(P\_3,3.131)} eṣaḥ asmi prajaniṣṇuḥ . \\
\noindent \textbf{(P\_3,3.132.1)} āśaṃsā nāma bhaviṣyatkālā . \\
\noindent \textbf{(P\_3,3.132.1)} āśaṃsāyām bhūtavadatideśe laṅliṭoḥ pratiṣedhaḥ . \\
\noindent \textbf{(P\_3,3.132.1)} āśaṃsāyām bhūtavadatideśe laṅliṭoḥ pratiṣedhaḥ vaktavyaḥ . \\
\noindent \textbf{(P\_3,3.132.1)} na vā apavādasya nimittābhāvāt anadyatane hi tayoḥ vidhānam . \\
\noindent \textbf{(P\_3,3.132.1)} na vā vaktavyaḥ . \\
\noindent \textbf{(P\_3,3.132.1)} kim kāraṇam . \\
\noindent \textbf{(P\_3,3.132.1)} apavādasya nimittābhāvāt . \\
\noindent \textbf{(P\_3,3.132.1)} na atra apavādasya nimittam asti . \\
\noindent \textbf{(P\_3,3.132.1)} katham . \\
\noindent \textbf{(P\_3,3.132.1)} anadyatane hi tayoḥ vidhānam . \\
\noindent \textbf{(P\_3,3.132.1)} anadyatane hi tau vidhīyete laṅliṭau . \\
\noindent \textbf{(P\_3,3.132.1)} na ca atra anadyatanaḥ kālaḥ vivakṣitaḥ . \\
\noindent \textbf{(P\_3,3.132.1)} kaḥ tarhi . \\
\noindent \textbf{(P\_3,3.132.1)} bhūtakālasāmānyam \\
\noindent \textbf{(P\_3,3.132.2)} āśaṃsāsambhāvanayoḥ aviśeṣāt tadvidhānasya aprāptiḥ . āśaṃsā sambhāvanam iti aviśiṣṭau etau arthau . \\
\noindent \textbf{(P\_3,3.132.2)} āśaṃsāsambhāvanayoḥ aviśeṣāt tadvidhānasya aprāptiḥ . \\
\noindent \textbf{(P\_3,3.132.2)} āśaṃsāyām ye vidhīyante te sambhāvane api prāpnuvanti . \\
\noindent \textbf{(P\_3,3.132.2)} ye ca sambhāvane vidhīyante te āśaṃsāyām api prāpnuvanti . \\
\noindent \textbf{(P\_3,3.132.2)} kim tarhi ucyate aprāptiḥ iti . \\
\noindent \textbf{(P\_3,3.132.2)} na sādhīyaḥ prāptiḥ bhavati . \\
\noindent \textbf{(P\_3,3.132.2)} iṣṭā vyavasthā na prakalpeta . \\
\noindent \textbf{(P\_3,3.132.2)} na sarve sarvatra iṣyante . \\
\noindent \textbf{(P\_3,3.132.2)} na vā sambhāvanāvayavatvāt āśaṃsāyāḥ . \\
\noindent \textbf{(P\_3,3.132.2)} na vā eṣaḥ doṣaḥ . \\
\noindent \textbf{(P\_3,3.132.2)} kim kāraṇam . \\
\noindent \textbf{(P\_3,3.132.2)} sambhāvanāvayavatvāt āśaṃsāyāḥ . \\
\noindent \textbf{(P\_3,3.132.2)} sambhāvanāvayavātmikā āśaṃsā . \\
\noindent \textbf{(P\_3,3.132.2)} āśaṃsā nāma pradhāritaḥ arthaḥ abhinītaḥ ca anabhinītaḥ ca . \\
\noindent \textbf{(P\_3,3.132.2)} sambhāvanam nāma pradhāritaḥ arthaḥ abhinītaḥ eva. arthāsandehaḥ vā alamarthatvāt sambhāvanasya . \\
\noindent \textbf{(P\_3,3.132.2)} atha vā arthāsandehaḥ eva punaḥ asya . \\
\noindent \textbf{(P\_3,3.132.2)} kim kāraṇam . \\
\noindent \textbf{(P\_3,3.132.2)} alamarthatvāt sambhāvanasya . \\
\noindent \textbf{(P\_3,3.132.2)} sambhāvane ālamarthyam gamyate āsaṃśāyām punaḥ anālamarthyam . \\
\noindent \textbf{(P\_3,3.132.2)} ācāryapravṛttiḥ jñāpayati sambhāvane api anālamarthyam gamyate iti yat ayam sambhāvane alam iti āha . \\
\noindent \textbf{(P\_3,3.132.2)} tasmāt suṣṭhu ucyate na vā sambhāvanāvayavatvāt āśaṃsāyāḥ iti . \\
\noindent \textbf{(P\_3,3.133.1)} kṣipravacane lṛaḥ āśaṃsāvacane liṅ vipratiṣedhena . \\
\noindent \textbf{(P\_3,3.133.1)} kṣipravacane lṛaḥ āśaṃsāvacane liṅ bhavati vipratiṣedhena . \\
\noindent \textbf{(P\_3,3.133.1)} kṣipravacane lṛṭ bhavati iti asya avakāśaḥ . \\
\noindent \textbf{(P\_3,3.133.1)} upādhyāyaḥ cet āgataḥ kṣipram adhyeṣyāmahe . \\
\noindent \textbf{(P\_3,3.133.1)} āśaṃsāvacabe liṅ bhavati iti asya avakāśaḥ . \\
\noindent \textbf{(P\_3,3.133.1)} upādhyāyaḥ cet āgataḥ āśaṃse yuktaḥ adhīyīya . \\
\noindent \textbf{(P\_3,3.133.1)} iha ubhayam prāpnoti . \\
\noindent \textbf{(P\_3,3.133.1)} upādhyāyaḥ cet āgataḥ āśaṃse kṣipram adhīyīya . \\
\noindent \textbf{(P\_3,3.133.1)} liṅ bhavati vipratiṣedhena . \\
\noindent \textbf{(P\_3,3.133.2)} aniṣpanne niṣpannaśabdaḥ śiṣyaḥ aniṣpannatvāt . \\
\noindent \textbf{(P\_3,3.133.2)} aniṣpanne niṣpannaśabdaḥ śiṣyaḥ śāsitavyaḥ . \\
\noindent \textbf{(P\_3,3.133.2)} kim kāraṇam . \\
\noindent \textbf{(P\_3,3.133.2)} aniṣpannatvāt . \\
\noindent \textbf{(P\_3,3.133.2)} devaḥ cet vṛṣṭaḥ niṣpannāḥ śālayaḥ . \\
\noindent \textbf{(P\_3,3.133.2)} tatra bhavitavyam sampatsyante śālayaḥ iti . \\
\noindent \textbf{(P\_3,3.133.2)} siddham tu bhaviṣyatpratiṣedhāt . \\
\noindent \textbf{(P\_3,3.133.2)} siddham etat . \\
\noindent \textbf{(P\_3,3.133.2)} katham . \\
\noindent \textbf{(P\_3,3.133.2)} bhaviṣyatpratiṣedhāt . \\
\noindent \textbf{(P\_3,3.133.2)} yat lokaḥ bhaviṣyadvācinaḥ śabdasya prayogam na mṛṣyati . \\
\noindent \textbf{(P\_3,3.133.2)} kaḥ cit āha . \\
\noindent \textbf{(P\_3,3.133.2)} devaḥ cet vṛṣṭaḥ sampatsyante śālayaḥ iti . \\
\noindent \textbf{(P\_3,3.133.2)} saḥ ucyate . \\
\noindent \textbf{(P\_3,3.133.2)} mā evam vocaḥ . \\
\noindent \textbf{(P\_3,3.133.2)} sampannāḥ śālayaḥ iti evam brūhi . \\
\noindent \textbf{(P\_3,3.133.2)} hetubhūtakālasamprekṣitatvāt vā . \\
\noindent \textbf{(P\_3,3.133.2)} hetubhūtakālasamprekṣitatvāt vā punaḥ siddham etat . \\
\noindent \textbf{(P\_3,3.133.2)} hetubhūtakālam varṣam varṣākālā ca kriyā . \\
\noindent \textbf{(P\_3,3.133.2)} yadi tarhi niṣpannaḥ arthaḥ kim niṣpannakāryāṇi na kriyante . \\
\noindent \textbf{(P\_3,3.133.2)} kāni . \\
\noindent \textbf{(P\_3,3.133.2)} bhojanādīni . \\
\noindent \textbf{(P\_3,3.133.2)} anyat idānīm etat ucyate kim niṣpannakāryāṇi na kriyante iti . \\
\noindent \textbf{(P\_3,3.133.2)} yat tu tat niṣpannaḥ arthaḥ na niṣpannaḥ iti . \\
\noindent \textbf{(P\_3,3.133.2)} saḥ niṣpannaḥ arthaḥ . \\
\noindent \textbf{(P\_3,3.133.2)} avaśyam khalu api koṣṭhagateṣu api śāliṣu avahananādīni pratīkṣyāṇi . \\
\noindent \textbf{(P\_3,3.133.2)} evam iha api niṣpannaḥ arthaḥ . \\
\noindent \textbf{(P\_3,3.133.2)} avaśyam tu jananādīni pratīkṣyāṇi . \\
\noindent \textbf{(P\_3,3.133.3)} astyarthānām bhavantyarthe sarvāḥ vibhaktayaḥ . astyarthānām bhavantyarthe sarvāḥ vibhaktayaḥ . \\
\noindent \textbf{(P\_3,3.133.3)} kūpaḥ asti . \\
\noindent \textbf{(P\_3,3.133.3)} kūpaḥ bhaviṣyati . \\
\noindent \textbf{(P\_3,3.133.3)} kūpaḥ bhavitā . \\
\noindent \textbf{(P\_3,3.133.3)} kūpaḥ abhūt . \\
\noindent \textbf{(P\_3,3.133.3)} kūpaḥ āsīt . \\
\noindent \textbf{(P\_3,3.133.3)} kūpaḥ babhūva iti . \\
\noindent \textbf{(P\_3,3.133.3)} katham punaḥ jñāyate bhavantyāḥ eṣaḥ arthaḥ iti . \\
\noindent \textbf{(P\_3,3.133.3)} kartuḥ vidyamānatvāt . \\
\noindent \textbf{(P\_3,3.133.3)} kartā atra vidyate . \\
\noindent \textbf{(P\_3,3.133.3)} katham punaḥ jñāyate kartā atra vidyate iti . \\
\noindent \textbf{(P\_3,3.133.3)} kūpaḥ anena kadā cit dṛṣṭaḥ . \\
\noindent \textbf{(P\_3,3.133.3)} na ca asya kam cid api apāyam paśyati . \\
\noindent \textbf{(P\_3,3.133.3)} saḥ tu tatra buddhyā nityām sattām adhyavasyati . \\
\noindent \textbf{(P\_3,3.133.3)} kūpaḥ asti iti . \\
\noindent \textbf{(P\_3,3.133.3)} siddham tu yathāsvam kālasamuccāraṇāt . \\
\noindent \textbf{(P\_3,3.133.3)} siddham etat . \\
\noindent \textbf{(P\_3,3.133.3)} katham . \\
\noindent \textbf{(P\_3,3.133.3)} yathāsvam etāḥ vibhaktayaḥ sveṣu sveṣu kāleṣu prayujyante iti . \\
\noindent \textbf{(P\_3,3.133.3)} katham punaḥ jñāyate yathāsvam etāḥ vibhaktayaḥ sveṣu sveṣu kāleṣu prayujyante iti . \\
\noindent \textbf{(P\_3,3.133.3)} avātvāt . \\
\noindent \textbf{(P\_3,3.133.3)} yat na vā bhāṣyante . \\
\noindent \textbf{(P\_3,3.133.3)} asiddhaviparyāsaḥ ca . asiddhaḥ ca viparyāsaḥ . \\
\noindent \textbf{(P\_3,3.133.3)} na hi kaḥ cit kūpaḥ asti iti prayoktavye kūpaḥ abhūt iti prayuṅkte . \\
\noindent \textbf{(P\_3,3.133.3)} kim punaḥ kāraṇam . \\
\noindent \textbf{(P\_3,3.133.3)} na vā bhāṣyante asiddhaḥ ca viparyāsaḥ . \\
\noindent \textbf{(P\_3,3.133.3)} iha kim cit indriyakarma kim cit buddhikarma . \\
\noindent \textbf{(P\_3,3.133.3)} indriyakarma samāsādanam buddhikarmavyavasāyaḥ . \\
\noindent \textbf{(P\_3,3.133.3)} evam hi kaḥ cit pāṭaliputram jigamiṣuḥ āha . \\
\noindent \textbf{(P\_3,3.133.3)} yaḥ ayam adhvā gantavyaḥ ā pāṭaliputrāt etasmin kūpaḥ bhaviṣyati . \\
\noindent \textbf{(P\_3,3.133.3)} samāsādya atikramya uṣitvā kūpaḥ āsīt iti . \\
\noindent \textbf{(P\_3,3.133.3)} samāsādya atikramya uṣitvā vismṛtya kūpaḥ babhūva iti . \\
\noindent \textbf{(P\_3,3.133.3)} tat yadā indriyakarma tadā etāḥ vibhaktayaḥ . \\
\noindent \textbf{(P\_3,3.133.3)} yadā hi buddhikarma tadā vartamānā bhaviṣyati . \\
\noindent \textbf{(P\_3,3.135)} kimartham imau dvau pratiṣedhau ucyete na adyatanavat iti eva ucyeta . \\
\noindent \textbf{(P\_3,3.135)} na anadyatanavatpratiṣedhe laṅluṭoḥ pratiṣedhaḥ . \\
\noindent \textbf{(P\_3,3.135)} na anadyatanavatpratiṣedhe laṅluṭoḥ pratiṣedhaḥ draṣṭavyaḥ . \\
\noindent \textbf{(P\_3,3.135)} adyatanavadvacane hi vidhānam . \\
\noindent \textbf{(P\_3,3.135)} adyatanavadvacane hi sati vidhiḥ iyam vijñāyeta . \\
\noindent \textbf{(P\_3,3.135)} tatra kaḥ doṣaḥ . \\
\noindent \textbf{(P\_3,3.135)} tatra laḍvidhiprasaṅgaḥ . \\
\noindent \textbf{(P\_3,3.135)} tatra laḍvidhiḥ prasajyeta . \\
\noindent \textbf{(P\_3,3.135)} luṅlṛṭoḥ ca ayathākālam . luṅlṛṭoḥ ca ayathākālam prayogaḥ prasajyeta . \\
\noindent \textbf{(P\_3,3.135)} luṅaḥ api viṣaye lṛṭ syāt lṛṭaḥ ca viṣaye luṅ syāt . \\
\noindent \textbf{(P\_3,3.135)} adya punaḥ ayam dvau pratiṣedhau uktvā tūṣṇīm āste . \\
\noindent \textbf{(P\_3,3.135)} yathāprāptem eva adyatane bhaviṣyati iti . \\
\noindent \textbf{(P\_3,3.136)} kimartham idam ucyate na na anadyatanavat iti eva siddham . \\
\noindent \textbf{(P\_3,3.136)} bhaviṣyati maryādāvacane avarasmin iti akriyāprabandhārtham . \\
\noindent \textbf{(P\_3,3.136)} akriyāprabandhārthaḥ ayam ārambhaḥ . \\
\noindent \textbf{(P\_3,3.136)} kim ucyate akriyāprabandhaḥ . \\
\noindent \textbf{(P\_3,3.136)} na punaḥ kriyāprabandhārthaḥ api syāt . \\
\noindent \textbf{(P\_3,3.136)} kriyāprabandhārtham iti cet vacanānarthakyam . \\
\noindent \textbf{(P\_3,3.136)} kriyāprabandhārtham iti cet vacanam anarthakam . \\
\noindent \textbf{(P\_3,3.136)} siddham kriyāprabandhe pūrveṇa eva . \\
\noindent \textbf{(P\_3,3.136)} idam tarhi prayojanam . \\
\noindent \textbf{(P\_3,3.136)} anahorātrāṇām iti vakṣyāmi iti . \\
\noindent \textbf{(P\_3,3.136)} iha mā bhūt . \\
\noindent \textbf{(P\_3,3.136)} yaḥ ayam triṃśadrātraḥ āgāmī tasya yaḥ avaraḥ pañcadaśarātraḥ iti . \\
\noindent \textbf{(P\_3,3.136)} ahorātrapratiṣedhārtham iti cet na aniṣṭatvāt . \\
\noindent \textbf{(P\_3,3.136)} ahorātrapratiṣedhārtham iti cet tat na aniṣṭatvāt . \\
\noindent \textbf{(P\_3,3.136)} kim kāraṇam . \\
\noindent \textbf{(P\_3,3.136)} aniṣṭatvāt . \\
\noindent \textbf{(P\_3,3.136)} atra api na anadyatanavat iti eva iṣyate . \\
\noindent \textbf{(P\_3,3.136)} idam tarhi prayojanam : bhaviṣyati iti vakṣyāmi iti . \\
\noindent \textbf{(P\_3,3.136)} iha ma bhūt . \\
\noindent \textbf{(P\_3,3.136)} yaḥ ayam adhvā gataḥ ā pāṭaliputrāt tasya yat avaram sāketāt iti . \\
\noindent \textbf{(P\_3,3.136)} na aniṣṭatvāt . \\
\noindent \textbf{(P\_3,3.136)} atra api na anadyatanavat iti eva iṣyate . \\
\noindent \textbf{(P\_3,3.136)} idam tarhi prayojanam maryādāvacane iti vakṣyāmi iti . \\
\noindent \textbf{(P\_3,3.136)} iha mā bhūt . \\
\noindent \textbf{(P\_3,3.136)} yaḥ ayam adhvā aparimāṇaḥ gantavyaḥ tasya yat avaram sāketāt iti . \\
\noindent \textbf{(P\_3,3.136)} atra api na anadyatanavat iti eva iṣyate . \\
\noindent \textbf{(P\_3,3.136)} tasmāt suṣṭhu ucyate bhaviṣyati maryādāvacane avarasmin iti akriyāprabandhārtham . \\
\noindent \textbf{(P\_3,3.136)} kriyāprabandhārtham iti cet vacanānarthakyam iti . \\
\noindent \textbf{(P\_3,3.137)} anahorātrāṇām iti tadvibhāge pratiṣedhaḥ . \\
\noindent \textbf{(P\_3,3.137)} anahorātrāṇām iti tadvibhāge pratiṣedhaḥ vaktavyaḥ . \\
\noindent \textbf{(P\_3,3.137)} yaḥ ayam triṃśadrātraḥ āgāmītasya yaḥ avaraḥ ardhamāsaḥ . \\
\noindent \textbf{(P\_3,3.137)} taiḥ ca vibhāge . \\
\noindent \textbf{(P\_3,3.137)} taiḥ ca vibhāge iti vaktavyam : yaḥ ayam māsaḥ āgāmītasya yaḥ avaraḥ pañcadaśarātraḥ iti . \\
\noindent \textbf{(P\_3,3.137)} dveṣyam vijānīyāt : ahorātrāṇām eva ahorātraiḥ vibhāge pratiṣedhaḥ iti . \\
\noindent \textbf{(P\_3,3.137)} tat ācāryaḥ suhṛt bhūtvā anvācaṣṭe : anahorātrāṇām iti tadvibhāge pratiṣedhaḥ . \\
\noindent \textbf{(P\_3,3.137)} taiḥ ca vibhāge iti . \\
\noindent \textbf{(P\_3,3.138)} kasmin parasmin . \\
\noindent \textbf{(P\_3,3.138)} kālavibhāge . \\
\noindent \textbf{(P\_3,3.138)} kutaḥ etat . \\
\noindent \textbf{(P\_3,3.138)} yogavibhāgakaraṇasāmarthyāt . \\
\noindent \textbf{(P\_3,3.139)} sādhanātipattau iti api vaktavyam iha api yathā syāt . \\
\noindent \textbf{(P\_3,3.139)} abhokṣyata bhavān māṃsena yadi matsamīpe āsiṣyata iti . \\
\noindent \textbf{(P\_3,3.139)} tat tarhi vaktavyam . \\
\noindent \textbf{(P\_3,3.139)} na vaktavyam . \\
\noindent \textbf{(P\_3,3.139)} na antareṇa sādhanam kriyāyāḥ pravṛttiḥ asti iti sādhanātipattiḥ cet kriyātipattiḥ api bhavati . \\
\noindent \textbf{(P\_3,3.139)} tatra kriyātipattau iti eva siddham . \\
\noindent \textbf{(P\_3,3.140)} bhūte lṛṅ utāpyādiṣu . \\
\noindent \textbf{(P\_3,3.140)} bhūte lṛṅ utāpyādiṣu draṣṭavyaḥ . \\
\noindent \textbf{(P\_3,3.140)} uta adhyāiṣyata . \\
\noindent \textbf{(P\_3,3.140)} api adhyaiṣyata . \\
\noindent \textbf{(P\_3,3.141)} vibhāṣā garhāprabhṛtau prāk utāpibhyām . \\
\noindent \textbf{(P\_3,3.141)} vibhāṣā garhāprabhṛtau prāk utāpibhyām iti vaktavyam . \\
\noindent \textbf{(P\_3,3.141)} vā ā utāpyoḥ iti hi ucyamāne sandehaḥ syāt : prāk vā utāpibhyām saha vā iti . \\
\noindent \textbf{(P\_3,3.141)} tat ācāryaḥ suhṛt bhūtvā anvācaṣṭe vibhāṣā garhāprabhṛtau prāk utāpibhyām iti . \\
\noindent \textbf{(P\_3,3.142)} garhāyām laḍvidhānānarthakyam kriyāsamāptivivakṣitatvāt . \\
\noindent \textbf{(P\_3,3.142)} garhāyām laḍvidhiḥ narthakaḥ . \\
\noindent \textbf{(P\_3,3.142)} kim kāraṇam . \\
\noindent \textbf{(P\_3,3.142)} kriyāsamāptivivakṣitatvāt . \\
\noindent \textbf{(P\_3,3.142)} kriyāyāḥ atra asamāptiḥ gamyate . \\
\noindent \textbf{(P\_3,3.142)} eṣaḥ ca nāma nyāyyaḥ vartamānaḥ kālaḥ yatra kriyā aparisamāptā bhavati . \\
\noindent \textbf{(P\_3,3.142)} tatra vartamāne laṭ iti eva siddham . \\
\noindent \textbf{(P\_3,3.142)} yadi vartamāne laṭ iti evam atra laṭ bhavati śatṛśānacau prāpnutaḥ . \\
\noindent \textbf{(P\_3,3.142)} iṣyete ca śatṛśānacau : api mām yājayantam paśya . \\
\noindent \textbf{(P\_3,3.142)} api mām yājayamānam paśya . \\
\noindent \textbf{(P\_3,3.145)} kiṃvṛttasya anadhikārāt uttaratra akiṃvṛttagrahaṇānarthakyam . \\
\noindent \textbf{(P\_3,3.145)} kiṃvṛttasya anadhikārāt uttaratra akiṃvṛttagrahaṇam anarthakam . \\
\noindent \textbf{(P\_3,3.145)} nivṛttam kiṃvṛtte iti . \\
\noindent \textbf{(P\_3,3.145)} tasmin nivṛtte aviśeṣeṇa kiṃvṛtte akiṃvṛtte ca bhaviṣyati . \\
\noindent \textbf{(P\_3,3.145)} idam tarhi prayojanam upapadasañjñām vakṣyāmi iti . \\
\noindent \textbf{(P\_3,3.145)} upapadasañjñāvacane kim prayojanam . \\
\noindent \textbf{(P\_3,3.145)} upapadam atiṅ iti samāsaḥ yathā syāt . \\
\noindent \textbf{(P\_3,3.145)} atiṅ iti pratiṣedhaḥ prāpnoti . \\
\noindent \textbf{(P\_3,3.145)} yadā tarhi lṛṭaḥ satsañjñau tadā upapadasañjñā bhaviṣyati . \\
\noindent \textbf{(P\_3,3.145)} bhaviṣyadadhikāravihitasya lṛṭaḥ satsñjñau ucyete aviśeṣavihitaḥ ca ayam . \\
\noindent \textbf{(P\_3,3.147)} jātuyadoḥ liṅvidhāne yadāyadyoḥ upasaṅkhyānam . \\
\noindent \textbf{(P\_3,3.147)} jātuyadoḥ liṅvidhāne yadāyadyoḥ upasaṅkhyānam kartavyam . \\
\noindent \textbf{(P\_3,3.147)} yadā bhavadvidhaḥ kṣatriyam yājayet . \\
\noindent \textbf{(P\_3,3.147)} yadi bhavadvidhaḥ kṣatriyam yājayet . \\
\noindent \textbf{(P\_3,3.151)} citrīkaraṇe yadipratiṣedhānarthakyam arthānyatvāt . \\
\noindent \textbf{(P\_3,3.151)} citrīkaraṇe yadipratiṣedhaḥ anarthakaḥ . \\
\noindent \textbf{(P\_3,3.151)} kim kāraṇam . \\
\noindent \textbf{(P\_3,3.151)} arthānyatvāt . \\
\noindent \textbf{(P\_3,3.151)} na hi yadau upapade citrīkaraṇam gamyate . \\
\noindent \textbf{(P\_3,3.151)} kim tarhi . \\
\noindent \textbf{(P\_3,3.151)} sambhāvanam . \\
\noindent \textbf{(P\_3,3.156)} hetuhetumatoḥ liṅ vā . \\
\noindent \textbf{(P\_3,3.156)} hetuhetumatoḥ liṅ vā iti vaktavyam . \\
\noindent \textbf{(P\_3,3.156)} anena cet yāyāt na śakaṭam paryābhavet . \\
\noindent \textbf{(P\_3,3.156)} anena cet yāsyati na śakaṭam paryābhaviṣyati . \\
\noindent \textbf{(P\_3,3.156)} bhaviṣyadadhikāre . \\
\noindent \textbf{(P\_3,3.156)} bhaviṣyadadhikāre iti vaktavyam . \\
\noindent \textbf{(P\_3,3.156)} iha mā bhūt . \\
\noindent \textbf{(P\_3,3.156)} varṣati iti dhāvati . \\
\noindent \textbf{(P\_3,3.156)} hanti iti palayate . \\
\noindent \textbf{(P\_3,3.156)} atha idānīm śatṛśānacau atra kasmāt na bhavataḥ . \\
\noindent \textbf{(P\_3,3.156)} devatrātaḥ galaḥ grāhaḥ itiyoge ca sadvidhiḥ . \\
\noindent \textbf{(P\_3,3.156)} mithaḥ te na vibhāṣyante . \\
\noindent \textbf{(P\_3,3.156)} gavākṣaḥ saṃśitavrataḥ . \\
\noindent \textbf{(P\_3,3.157)} kāmapravedanam cet . \\
\noindent \textbf{(P\_3,3.157)} kāmapravedanam cet gamyate iti vaktavyam . \\
\noindent \textbf{(P\_3,3.157)} iha mā bhūt . \\
\noindent \textbf{(P\_3,3.157)} icchan kaṭam karoti . \\
\noindent \textbf{(P\_3,3.161.1)} vidhyadhīṣṭayoḥ kaḥ viśeṣaḥ . \\
\noindent \textbf{(P\_3,3.161.1)} vidhiḥ nāma preṣaṇam . \\
\noindent \textbf{(P\_3,3.161.1)} adhīṣṭam nām satkārpūrvikā vyāpāraṇā . \\
\noindent \textbf{(P\_3,3.161.1)} atha nimantraṇāmantraṇayoḥ kaḥ viśeṣaḥ . \\
\noindent \textbf{(P\_3,3.161.1)} sannihitena nimantraṇam bhavati asannihitena ca āmantraṇam . \\
\noindent \textbf{(P\_3,3.161.1)} na eṣaḥ asti viśeṣaḥ . \\
\noindent \textbf{(P\_3,3.161.1)} asannihitena api nimantraṇam bhavati sannihitena ca āmantraṇam . \\
\noindent \textbf{(P\_3,3.161.1)} evam tarhi yat niyogataḥ kartavyam tat nimantraṇam . \\
\noindent \textbf{(P\_3,3.161.1)} kim punaḥ tat . \\
\noindent \textbf{(P\_3,3.161.1)} havyam kavyam vā . \\
\noindent \textbf{(P\_3,3.161.1)} brāhmaṇena siddham bhujyatām iti ukte adharmaḥ pratyākhyātuḥ . \\
\noindent \textbf{(P\_3,3.161.1)} āmantraṇe kāmacāraḥ . \\
\noindent \textbf{(P\_3,3.161.2)} katham punaḥ idam vijñāyate . \\
\noindent \textbf{(P\_3,3.161.2)} nimantraṇādīīnām arthe iti āhosvit nimantraṇādiṣu gamyamāneṣu iti . \\
\noindent \textbf{(P\_3,3.161.2)} kaḥ ca atra viśeṣaḥ . \\
\noindent \textbf{(P\_3,3.161.2)} nimantraṇādīīnām arthe iti cet āmantrayai nimantrayai bhavantam iti pratyayānupapattiḥ prakṛtyā abhihitatvāt . \\
\noindent \textbf{(P\_3,3.161.2)} nimantraṇādīīnām arthe iti cet āmantrayai nimantrayai bhavantam iti pratyayānupapattiḥ . \\
\noindent \textbf{(P\_3,3.161.2)} kim kāraṇam . \\
\noindent \textbf{(P\_3,3.161.2)} prakṛtyā abhihitatvāt . \\
\noindent \textbf{(P\_3,3.161.2)} prakṛtyā abhihitaḥ saḥ arthaḥ iti kṛtvā pratyayaḥ na prāpnoti . \\
\noindent \textbf{(P\_3,3.161.2)} dvivacanabahuvanāprasiddhiḥ ca ekārthatvāt . \\
\noindent \textbf{(P\_3,3.161.2)} dvivacanabahuvanayoḥ ca a prasiddhiḥ . \\
\noindent \textbf{(P\_3,3.161.2)} kim kāraṇam . \\
\noindent \textbf{(P\_3,3.161.2)} ekārthatvāt . \\
\noindent \textbf{(P\_3,3.161.2)} ekaḥ ayam arthaḥ nimantraṇam nāma . \\
\noindent \textbf{(P\_3,3.161.2)} tasya ekatvāt ekavacanam eva prāpnoti . \\
\noindent \textbf{(P\_3,3.161.2)} astu tarhi nimantraṇādiṣu gamyamāneṣu . \\
\noindent \textbf{(P\_3,3.161.2)} iha api tarhi prāpnoti . \\
\noindent \textbf{(P\_3,3.161.2)} devadattaḥ bhavantam āmantrayate . \\
\noindent \textbf{(P\_3,3.161.2)} devadattaḥ bhavantam nimantrayate iti . \\
\noindent \textbf{(P\_3,3.161.2)} siddham tu dvitīyākāṅkṣasya prakṛte pratyayārthe pratyayavidhānāt . \\
\noindent \textbf{(P\_3,3.161.2)} siddham etat . \\
\noindent \textbf{(P\_3,3.161.2)} katham . \\
\noindent \textbf{(P\_3,3.161.2)} dvitīyākāṅkṣasya dhātoḥ prakṛte pratyayārthe pratyayaḥ bhavati iti vaktavyam . \\
\noindent \textbf{(P\_3,3.161.2)} ke ca prakṛtāḥ arthāḥ . \\
\noindent \textbf{(P\_3,3.161.2)} bhāvakarmakartāraḥ . \\
\noindent \textbf{(P\_3,3.161.2)} bhavet siddham prāpnotu bhavān āmantraṇam anubhavatu bhavān amantraṇam iti yatra dvitīyaḥ ākāṅkyate . \\
\noindent \textbf{(P\_3,3.161.2)} idam tu na sidhyati āmantrayai nimantrayai iti . \\
\noindent \textbf{(P\_3,3.161.2)} atra api dvitīyaḥ ākāṅkyate . \\
\noindent \textbf{(P\_3,3.161.2)} kaḥ . \\
\noindent \textbf{(P\_3,3.161.2)} nimantriḥ eva . \\
\noindent \textbf{(P\_3,3.161.2)} āmantrayai āmantraṇam . \\
\noindent \textbf{(P\_3,3.161.2)} nimantrayai nimantraṇam . \\
\noindent \textbf{(P\_3,3.161.2)} katham punaḥ nimnatriḥ nimantraṇam ākāṅkṣet . \\
\noindent \textbf{(P\_3,3.161.2)} dṛṣṭaḥ ca bhāvena bhāvayogaḥ . \\
\noindent \textbf{(P\_3,3.161.2)} tat yathā iṣiḥ iṣiṇā yujyate strītvam ca strītvena . \\
\noindent \textbf{(P\_3,3.161.2)} yāvatā atra dvitīyaḥ ākāṅkṣyate asti tarhi nimantraṇādīnām arthe iti . \\
\noindent \textbf{(P\_3,3.161.2)} nanu ca uktam nimantraṇādīīnām arthe iti cet āmantrayai nimantrayai bhavantam iti pratyayānupapattiḥ prakṛtyā abhihitatvāt iti . \\
\noindent \textbf{(P\_3,3.161.2)} na eṣaḥ doṣaḥ . \\
\noindent \textbf{(P\_3,3.161.2)} yaḥ asau dvitīyaḥ ākāṅkṣyate saḥ eva mama pratyayārthaḥ bhaviṣyati . \\
\noindent \textbf{(P\_3,3.161.2)} ayam tarhi doṣaḥ dvivacanabahuvanāprasiddhiḥ ca ekārthatvāt iti . \\
\noindent \textbf{(P\_3,3.161.2)} na eṣaḥ doṣaḥ . \\
\noindent \textbf{(P\_3,3.161.2)} supām karmādayaḥ api arthāḥ saṅkhyā ca eva tathā tiṅām . supām saṅkhyā ca eva arthaḥ karmādayaḥ ca . \\
\noindent \textbf{(P\_3,3.161.2)} tathā tiṅām . \\
\noindent \textbf{(P\_3,3.161.2)} prasiddhaḥ niyamaḥ tatra . \\
\noindent \textbf{(P\_3,3.161.2)} prasiddhaḥ tatra niyamaḥ . \\
\noindent \textbf{(P\_3,3.161.2)} niyamaḥ prakṛteṣu vā . \\
\noindent \textbf{(P\_3,3.161.2)} atha vā prakṛtān arthān apekṣya niyamaḥ . \\
\noindent \textbf{(P\_3,3.161.2)} ke ca prakṛtāḥ . \\
\noindent \textbf{(P\_3,3.161.2)} ekatvādayaḥ . \\
\noindent \textbf{(P\_3,3.161.2)} ekasmin eva ekavacanam na dvayoḥ na bahuṣu . \\
\noindent \textbf{(P\_3,3.161.2)} dvayoḥ eva dvivacanam naikasmin na bahuṣu . \\
\noindent \textbf{(P\_3,3.161.2)} bahuṣu eva bahuvacanam na ekasmin na dvayoḥ iti . \\
\noindent \textbf{(P\_3,3.163)} kimartham praiṣādiṣu artheṣu kṛtyāḥ vidhīyante na aviśeṣeṇa vihitāḥ kṛtyāḥ te praiṣādiṣu bhaviṣyanti anyatra ca . \\
\noindent \textbf{(P\_3,3.163)} praiṣādiṣu kṛtyānām vidhānam niyamārtham . \\
\noindent \textbf{(P\_3,3.163)} niyamārthaḥ ayam ārambhaḥ . \\
\noindent \textbf{(P\_3,3.163)} praiṣādiṣu eva kṛtyāḥ yathā syuḥ iti . \\
\noindent \textbf{(P\_3,3.163)} praiṣādiṣu kṛtyānām vacanam niyamārtham iti cet tat aniṣṭam . praiṣādiṣu kṛtyānām vacanam niyamārtham iti cet tat aniṣṭam prāpnoti . \\
\noindent \textbf{(P\_3,3.163)} na hi praiṣādiṣu eva kṛtyāḥ iṣyante . \\
\noindent \textbf{(P\_3,3.163)} kim tarhi . \\
\noindent \textbf{(P\_3,3.163)} aviśeṣeṇa iṣyante . \\
\noindent \textbf{(P\_3,3.163)} busopendhyam tṛṇopendhyam ghanghātyam . \\
\noindent \textbf{(P\_3,3.163)} vidhyartham tu striyāḥ prāk iti vacanāt . \\
\noindent \textbf{(P\_3,3.163)} vidhyartham tu praiṣādiṣu kṛtyānām vacanam . \\
\noindent \textbf{(P\_3,3.163)} ayam praiṣādiṣu loṭ vidhīyate . \\
\noindent \textbf{(P\_3,3.163)} saḥ viśeṣavihitaḥ sāmānyavihitān kṛtyān bādheta . \\
\noindent \textbf{(P\_3,3.163)} vāsarūpeṇa kṛtyāḥ api bhaviṣyanti . \\
\noindent \textbf{(P\_3,3.163)} na syuḥ . \\
\noindent \textbf{(P\_3,3.163)} kim kāraṇam . \\
\noindent \textbf{(P\_3,3.163)} striyāḥ prāk iti vacanāt . \\
\noindent \textbf{(P\_3,3.163)} prāk striyāḥ vā asarūpaḥ . \\
\noindent \textbf{(P\_3,3.167)} prathamānteṣu iti vaktavyam . \\
\noindent \textbf{(P\_3,3.167)} kim prayojanam . \\
\noindent \textbf{(P\_3,3.167)} iha mā bhūt . \\
\noindent \textbf{(P\_3,3.167)} kāle bhuṅkte . \\
\noindent \textbf{(P\_3,3.167)} tat tarhi vaktavyam . \\
\noindent \textbf{(P\_3,3.167)} na vaktavyam . \\
\noindent \textbf{(P\_3,3.167)} praiṣādiṣu iti vartate . \\
\noindent \textbf{(P\_3,3.167)} tat ca avaśyam praiṣādigrahaṇam anuvartyam . \\
\noindent \textbf{(P\_3,3.167)} prathamānteṣu iti hi ucyamāne iha api prasajyeta . \\
\noindent \textbf{(P\_3,3.167)} kālaḥ pacati bhūtāni kālaḥ saṃharati prajāḥ . \\
\noindent \textbf{(P\_3,4.2)} hisvoḥ parasmaipadātmanepadagrahaṇam lādeśapratiṣedhārtham . \\
\noindent \textbf{(P\_3,4.2)} hisvoḥ parasmaipadātmanepadagrahaṇam kartavyam hiḥ parasmaipadānām yathā syāt svaḥ ātmanepadānām iti . \\
\noindent \textbf{(P\_3,4.2)} kim prayojanam . \\
\noindent \textbf{(P\_3,4.2)} lādeśapratiṣedhārtham . \\
\noindent \textbf{(P\_3,4.2)} lādeśau hisvau mā bhūtām iti . \\
\noindent \textbf{(P\_3,4.2)} kim ca syāt yadi lādeśau hisvau syātām . \\
\noindent \textbf{(P\_3,4.2)} tiṅantam padam iti padasañjñā na syāt . \\
\noindent \textbf{(P\_3,4.2)} māt bhūt evam . \\
\noindent \textbf{(P\_3,4.2)} subantam padam iti padasañjñā bhaviṣyati . \\
\noindent \textbf{(P\_3,4.2)} katham svādyutpattiḥ . \\
\noindent \textbf{(P\_3,4.2)} lakārasya kṛttvāt prātipadikatvam tadāśrayam pratyayavidhānam . \\
\noindent \textbf{(P\_3,4.2)} lakāraḥ kṛt . \\
\noindent \textbf{(P\_3,4.2)} tasya kṛttvāt kṛt prātipadikam iti prātipadikasañjñā . \\
\noindent \textbf{(P\_3,4.2)} prātipadikāśrayā svādyutpattiḥ api bhaviṣyati . \\
\noindent \textbf{(P\_3,4.2)} yadi svādyutpattiḥ supām śravaṇam prāpnoti . \\
\noindent \textbf{(P\_3,4.2)} avyayāt iti subluk bhaviṣyati . \\
\noindent \textbf{(P\_3,4.2)} katham avyayatvam . \\
\noindent \textbf{(P\_3,4.2)} vibhaktisvarapratirūpakāḥ ca nipātāḥ bhavanti iti nipātasañjñā . \\
\noindent \textbf{(P\_3,4.2)} nipātam avyayam iti avyayasañjñā . \\
\noindent \textbf{(P\_3,4.2)} iha tarhi saḥ bhavān lunīhi lunīhi iti eva ayam lunāti tiṅ atiṅaḥ iti nighātaḥ na prāpnoti . \\
\noindent \textbf{(P\_3,4.2)} samasaṅkhyārtham ca . \\
\noindent \textbf{(P\_3,4.2)} samasaṅkhyārtham ca hisvoḥ parasmaipadātmanepadagrahaṇam kartavyam hiḥ parasmaipadānām yathā syāt svaḥ ātmanepadānām . \\
\noindent \textbf{(P\_3,4.2)} vyatikaraḥ mā bhūt iti . \\
\noindent \textbf{(P\_3,4.2)} na vā tadhvamoḥ ādeśavacanam jñāpakam padādeśasya . \\
\noindent \textbf{(P\_3,4.2)} na vā hisvoḥ parasmaipadātmanepadagrahaṇam kartavyam . \\
\noindent \textbf{(P\_3,4.2)} kim kāraṇam . \\
\noindent \textbf{(P\_3,4.2)} tadhvamoḥ ādeśavacanam jñāpakam padādeśasya . \\
\noindent \textbf{(P\_3,4.2)} yat ayam vā ca tadhvamoḥ iti āha tat jñāpayati ācāryaḥ padādeśau hisvau iti . \\
\noindent \textbf{(P\_3,4.2)} tatra padādeśe pittvāṭoḥ pratiṣedhaḥ . \\
\noindent \textbf{(P\_3,4.2)} tatra padādeśe pittvasya āṭaḥ ca pratiṣedhaḥ vaktavyaḥ . \\
\noindent \textbf{(P\_3,4.2)} pittvasya tāvat . \\
\noindent \textbf{(P\_3,4.2)} saḥ bhavān lunīhi lunīhi iti eva ayam lunāti . \\
\noindent \textbf{(P\_3,4.2)} āṭaḥ khalu api . \\
\noindent \textbf{(P\_3,4.2)} saḥ aham lunīhi lunīhi iti evam lunāni . \\
\noindent \textbf{(P\_3,4.2)} pittvasya tāvat na vaktavyaḥ . \\
\noindent \textbf{(P\_3,4.2)} pitpratiṣedhe yogavibhāgaḥ kariṣyate . \\
\noindent \textbf{(P\_3,4.2)} iha seḥ hi bhavati . \\
\noindent \textbf{(P\_3,4.2)} tataḥ apit ca . \\
\noindent \textbf{(P\_3,4.2)} apit ca bhavati yāvān hiḥ nāma . \\
\noindent \textbf{(P\_3,4.2)} āṭaḥ ca api na vaktavyaḥ . \\
\noindent \textbf{(P\_3,4.2)} āṭi kṛte sāṭkasya ādeśaḥ bhaviṣyati . \\
\noindent \textbf{(P\_3,4.2)} idam iha sampradhāryam : āṭ kriyatām ādeśaḥ iti . \\
\noindent \textbf{(P\_3,4.2)} kim atra kartavyam . \\
\noindent \textbf{(P\_3,4.2)} paratvāt āḍāgamaḥ . \\
\noindent \textbf{(P\_3,4.2)} nityaḥ ādeśaḥ . \\
\noindent \textbf{(P\_3,4.2)} kṛte api āṭi prāpnoti akṛte api prāpnoti . \\
\noindent \textbf{(P\_3,4.2)} āṭ api nityaḥ . \\
\noindent \textbf{(P\_3,4.2)} kṛte api ādeśe prāpnoti akṛte api prāpnoti . \\
\noindent \textbf{(P\_3,4.2)} anityaḥ āṭ . \\
\noindent \textbf{(P\_3,4.2)} anyasya kṛte api ādeśe prāpnoti anyasya akṛte api prāpnoti . \\
\noindent \textbf{(P\_3,4.2)} śabdāntarasya ca prāpnuvan vidhiḥ anityaḥ bhavati . \\
\noindent \textbf{(P\_3,4.2)} ādeśaḥ api anityaḥ . \\
\noindent \textbf{(P\_3,4.2)} anyasya kṛte āṭi prāpnoti anyasya akṛte . \\
\noindent \textbf{(P\_3,4.2)} śabdāntarasya ca prāpnuvan vidhiḥ anityaḥ bhavati . \\
\noindent \textbf{(P\_3,4.2)} ubhayoḥ anityayoḥ paratvāt āḍāgamaḥ . \\
\noindent \textbf{(P\_3,4.2)} āṭi kṛte sāṭkasya ādeśaḥ bhaviṣyati . \\
\noindent \textbf{(P\_3,4.2)} idam tarhi saḥ aham bhuṅkṣva bhuṅkṣva iti evam bhunajai iti śnasoḥ allopaḥ iti akāralopaḥ na prāpnoti . \\
\noindent \textbf{(P\_3,4.2)} samasaṅkhyārthatvam ca api aparihṛtam eva. siddham tu loḍmadhyamapuruṣaikavacanasya kriyāsamabhihāre dvirvacanāt . \\
\noindent \textbf{(P\_3,4.2)} siddham etat . \\
\noindent \textbf{(P\_3,4.2)} katham . \\
\noindent \textbf{(P\_3,4.2)} loḍmadhyamapuruṣaikavacanasya kriyāsamabhihāre dve bhavataḥ iti vaktavyam . \\
\noindent \textbf{(P\_3,4.2)} kena vihitasya kriyāsamabhihāre loḍmadhyamapuruṣaikavacanasya dvirvacanam ucyate . \\
\noindent \textbf{(P\_3,4.2)} etat eva jñāpayati ācāryaḥ bhavati kriyāsamabhihāre loṭ iti yat ayam kriyāsamabhihāre loḍmadhyamapuruṣaikavacanasya dvirvacanam śāsti . \\
\noindent \textbf{(P\_3,4.2)} kutaḥ nu khalu etat jñāpakāt atra loṭ bhaviṣyati . \\
\noindent \textbf{(P\_3,4.2)} na punaḥ yaḥ eva asau aviśeṣavihitaḥ saḥ yadā kriyāsamabhihāre bhavati tadā asya dvirvacanam bhavati iti . \\
\noindent \textbf{(P\_3,4.2)} loḍmadhyamapuruṣaikavacane eva khalu api siddham syāt . \\
\noindent \textbf{(P\_3,4.2)} imau ca anyau hisvau sarveṣām puruṣāṇām sarveṣām vacanānām iṣyete . \\
\noindent \textbf{(P\_3,4.2)} sūtram ca bhidyate . \\
\noindent \textbf{(P\_3,4.2)} yathānyāsam eva astu . \\
\noindent \textbf{(P\_3,4.2)} nanu ca uktam hisvoḥ parasmaipadātmanepadagrahaṇam lādeśapratiṣedhārtham . \\
\noindent \textbf{(P\_3,4.2)} samasaṅkhyārtham ca iti . \\
\noindent \textbf{(P\_3,4.2)} na eṣaḥ doṣaḥ . \\
\noindent \textbf{(P\_3,4.2)} yogavibhāgāt siddham . \\
\noindent \textbf{(P\_3,4.2)} yogavibhāgaḥ kariṣyate . \\
\noindent \textbf{(P\_3,4.2)} kriyāsamabhihāre loṭ bhavati . \\
\noindent \textbf{(P\_3,4.2)} tataḥ loṭaḥ hisvau bhavataḥ . \\
\noindent \textbf{(P\_3,4.2)} loṭ iti eva anuvartate . \\
\noindent \textbf{(P\_3,4.2)} loṭaḥ yau hisvau iti . \\
\noindent \textbf{(P\_3,4.2)} katham vā ca tadhvamoḥ iti . \\
\noindent \textbf{(P\_3,4.2)} vā ca tadhvambhāvinaḥ loṭaḥ iti evam etat vijñāyate . \\
\noindent \textbf{(P\_3,4.4)} kimartham idam ucyate . \\
\noindent \textbf{(P\_3,4.4)} anuprayogaḥ yathā syāt . \\
\noindent \textbf{(P\_3,4.4)} na etat asti prayojanam . \\
\noindent \textbf{(P\_3,4.4)} hisvāntam avyaktapadārthakam . \\
\noindent \textbf{(P\_3,4.4)} tena aparisamāptaḥ arthaḥ iti kṛtvā anuprayogaḥ bhaviṣyati . \\
\noindent \textbf{(P\_3,4.4)} idam tarhi prayojanam . \\
\noindent \textbf{(P\_3,4.4)} yathāvidhi iti vakṣyāmi iti . \\
\noindent \textbf{(P\_3,4.4)} etat api na asti prayojanam . \\
\noindent \textbf{(P\_3,4.4)} samuccaye sāmānyavacanasya iti vakṣyati . \\
\noindent \textbf{(P\_3,4.4)} tatra antareṇa vacanam yathāvidhi anuprayogaḥ bhaviṣyati . \\
\noindent \textbf{(P\_3,4.5)} kimartham idam ucyate . \\
\noindent \textbf{(P\_3,4.5)} anuprayogaḥ yathā syāt . \\
\noindent \textbf{(P\_3,4.5)} na etat asti prayojanam . \\
\noindent \textbf{(P\_3,4.5)} hisvāntam avyaktapadārthakam . \\
\noindent \textbf{(P\_3,4.5)} tena aparisamāptaḥ arthaḥ iti kṛtvā anuprayogaḥ bhaviṣyati . \\
\noindent \textbf{(P\_3,4.5)} idam tarhi prayojanam . \\
\noindent \textbf{(P\_3,4.5)} sāmānyavacanasya iti vakṣyāmi iti . \\
\noindent \textbf{(P\_3,4.5)} etat api na asti prayojanam . \\
\noindent \textbf{(P\_3,4.5)} sāmānyavacanasya anuprayogaḥ astu viśeṣavacanasya iti sāmānyavacanasya anuprayogaḥ bhaviṣyati laghutvāt . \\
\noindent \textbf{(P\_3,4.8)} upasaṃvādāśaṅkayoḥ vacanānarthakyam liṅarthatvāt . \\
\noindent \textbf{(P\_3,4.8)} upasaṃvādāśaṅkayoḥ vacanam narthakam . \\
\noindent \textbf{(P\_3,4.8)} kim kāraṇam . \\
\noindent \textbf{(P\_3,4.8)} liṅarthatvāt . \\
\noindent \textbf{(P\_3,4.8)} liṅarthe leṭ iti eva siddham . \\
\noindent \textbf{(P\_3,4.8)} kaḥ punaḥ liṅarthaḥ . \\
\noindent \textbf{(P\_3,4.8)} ke cit tāvat āhuḥ . \\
\noindent \textbf{(P\_3,4.8)} hetuhetumatoḥ liṅ iti . \\
\noindent \textbf{(P\_3,4.8)} apare āhuḥ : vaktavyaḥ eva etasmin viśeṣe liṅ . \\
\noindent \textbf{(P\_3,4.8)} prayujyate hi loke yadi me bhavān idam kuryāt aham api te idam dadyām . \\
\noindent \textbf{(P\_3,4.9)} tumarthe iti ucyate . \\
\noindent \textbf{(P\_3,4.9)} kaḥ tumarthaḥ . \\
\noindent \textbf{(P\_3,4.9)} kartā . \\
\noindent \textbf{(P\_3,4.9)} yadi evam na arthaḥ tumarthagrahaṇena . \\
\noindent \textbf{(P\_3,4.9)} yena eva khalu api hetunā kartari tumun bhavati tena eva hetunā sayādayaḥ api bhaviṣyanti . \\
\noindent \textbf{(P\_3,4.9)} evam tarhi siddhe sati yat tumarthagrahaṇam karoti tat jñāpayati ācāryaḥ asti anyaḥ kartuḥ tumunaḥ arthaḥ iti . \\
\noindent \textbf{(P\_3,4.9)} kaḥ punaḥ asau . \\
\noindent \textbf{(P\_3,4.9)} bhāvaḥ . \\
\noindent \textbf{(P\_3,4.9)} kutaḥ nu khalu etat bhāve tumun bhaviṣyati . \\
\noindent \textbf{(P\_3,4.9)} na punaḥ karmādiṣu kārakeṣu iti . \\
\noindent \textbf{(P\_3,4.9)} jñāpakāt ayam kartuḥ apakṛṣyate . \\
\noindent \textbf{(P\_3,4.9)} na ca anyasmin arthe ādiśyate . \\
\noindent \textbf{(P\_3,4.9)} anirdiṣṭārthāḥ pratyayāḥ svārthe bhavanti iti svārthe bhaviṣyati tat yathā guptijkidbhyaḥ san yāvādibhyaḥ kan iti . \\
\noindent \textbf{(P\_3,4.9)} saḥ asau svārthe bhavan bhāve bhaviṣyati . \\
\noindent \textbf{(P\_3,4.9)} kim etasya jñāpane prayojanam . \\
\noindent \textbf{(P\_3,4.9)} avyayakṛtaḥ bhāve bhavanti iti etat na vaktavyam bhavati . \\
\noindent \textbf{(P\_3,4.19)} kimartham meṅaḥ sānubandhakasya āttvabhūtasya grahaṇam kriyate na udīcām meṅaḥ iti eva ucyeta . \\
\noindent \textbf{(P\_3,4.19)} tatra ayam api arthaḥ . \\
\noindent \textbf{(P\_3,4.19)} udīcām meṅaḥ iti vyatihāragrahaṇam na kartavyam bhavati . \\
\noindent \textbf{(P\_3,4.19)} kim kāraṇam . \\
\noindent \textbf{(P\_3,4.19)} tadviṣayaḥ hi saḥ . \\
\noindent \textbf{(P\_3,4.19)} vaytihāraviṣayaḥ eva mayatiḥ . \\
\noindent \textbf{(P\_3,4.19)} evam tarhi siddhe sati yat meṅaḥ sānubandhakasya āttvabhūtasya grahaṇam karoti tat jñāpayati ācāryaḥ na anubandhakṛtam anejantatvam bhavati iti .kim etasya jñāpane prayojanam . \\
\noindent \textbf{(P\_3,4.19)} tatra asarūpasarvādeśadāppratiṣedhe pṛthaktvanirdeśaḥ anākārāntatvāt iti uktam . \\
\noindent \textbf{(P\_3,4.19)} tat na vaktavyam bhavati . \\
\noindent \textbf{(P\_3,4.19)} kimartham punaḥ idam ucyate na samānakartṛkayoḥ pūrvakāle iti eva siddham . \\
\noindent \textbf{(P\_3,4.19)} apūrvakālārthaḥ ayam ārambhaḥ . \\
\noindent \textbf{(P\_3,4.19)} pūrvam hi asau yācate paścāt apamayate . \\
\noindent \textbf{(P\_3,4.21.1)} iha kasmāt na bhavati : pūrvam bhuṅkte paścāt vrajati . \\
\noindent \textbf{(P\_3,4.21.1)} svaśabdena uktatvāt na bhavati . \\
\noindent \textbf{(P\_3,4.21.1)} na tarhi idānīm idam bhavati : pūrvam bhuktvā tataḥ vrajati iti . \\
\noindent \textbf{(P\_3,4.21.1)} na etat kriyāpaurvakālyam . \\
\noindent \textbf{(P\_3,4.21.1)} kim tarhi . \\
\noindent \textbf{(P\_3,4.21.1)} kartṛpaurvakālyam . \\
\noindent \textbf{(P\_3,4.21.1)} pūrvam hi asau bhuktvā anyebhyaḥ bhoktṛbhyaḥ tataḥ paścāt vrajati anyebhyaḥ vrajitṛbhyaḥ . \\
\noindent \textbf{(P\_3,4.21.1)} iha kasmāt na bhavati : āsyate bhoktum iti . \\
\noindent \textbf{(P\_3,4.21.1)} kutaḥ kasmāt na bhavati . \\
\noindent \textbf{(P\_3,4.21.1)} kim āseḥ āhosvit bhujeḥ . \\
\noindent \textbf{(P\_3,4.21.1)} bhujeḥ kasmāt na bhavati . \\
\noindent \textbf{(P\_3,4.21.1)} apūrvakālatvāt . \\
\noindent \textbf{(P\_3,4.21.1)} āseḥ tarhi kasmāt na bhavati . \\
\noindent \textbf{(P\_3,4.21.1)} yasmāt atra laṭ bhavati . \\
\noindent \textbf{(P\_3,4.21.1)} etat atra praṣṭavyam . \\
\noindent \textbf{(P\_3,4.21.1)} laṭ atra katham bhavati iti . \\
\noindent \textbf{(P\_3,4.21.1)} laṭ ca atra vāsarūpeṇa bhaviṣyati . \\
\noindent \textbf{(P\_3,4.21.2)} samānakartṛkayoḥ iti bahuṣu aprāptiḥ . \\
\noindent \textbf{(P\_3,4.21.2)} samānakartṛkayoḥ iti bahuṣu ktvā na prāpnoti . \\
\noindent \textbf{(P\_3,4.21.2)} snātvā bhuktvā pītvā vrajati iti . \\
\noindent \textbf{(P\_3,4.21.2)} kim puna kāraṇam na sidhyati . \\
\noindent \textbf{(P\_3,4.21.2)} dvivacananirdeśāt . \\
\noindent \textbf{(P\_3,4.21.2)} dvivacanena ayam nirdeśaḥ kriyate . \\
\noindent \textbf{(P\_3,4.21.2)} tena dvayoḥ eva paurvakālye syāt . \\
\noindent \textbf{(P\_3,4.21.2)} bahūnām na syāt . \\
\noindent \textbf{(P\_3,4.21.2)} siddham tu kriyāpradhanatvāt . \\
\noindent \textbf{(P\_3,4.21.2)} siddham etat . \\
\noindent \textbf{(P\_3,4.21.2)} katham . \\
\noindent \textbf{(P\_3,4.21.2)} kriyāpradhanatvāt . \\
\noindent \textbf{(P\_3,4.21.2)} kriyāpradhānaḥ ayam nirdeśaḥ . \\
\noindent \textbf{(P\_3,4.21.2)} na atra nirdeśaḥ tantram . \\
\noindent \textbf{(P\_3,4.21.2)} katham punaḥ tena eva nāma nirdeśaḥ kriyate tat ca atantram syāt . \\
\noindent \textbf{(P\_3,4.21.2)} tatkārī ca bhavān taddveṣī ca . \\
\noindent \textbf{(P\_3,4.21.2)} nāntarīyakatvāt atra dvivacanena nirdeśaḥ kriyate . \\
\noindent \textbf{(P\_3,4.21.2)} avaśyam kayā cit vibhaktyā kena cit vacanena nirdeśaḥ kartavyaḥ . \\
\noindent \textbf{(P\_3,4.21.2)} tat yathā kaḥ cit annārthī śālikalāpam satuṣam sapalālam āharati nāntarīyakatvāt . \\
\noindent \textbf{(P\_3,4.21.2)} saḥ yāvat ādeyam tāvat ādāya tuṣapalālāni utsṛjati . \\
\noindent \textbf{(P\_3,4.21.2)} tathā kaḥ cit māṃsārthī matsyān saśakalān sakaṇṭakān āharati nāntarīyakatvāt . \\
\noindent \textbf{(P\_3,4.21.2)} saḥ yāvat ādeyam tāvat ādāya śakalakaṇṭakān utsṛjati . \\
\noindent \textbf{(P\_3,4.21.2)} evam iha api nāntarīyakatvāt dvivacanena nirdeśaḥ kriyate . \\
\noindent \textbf{(P\_3,4.21.2)} na hi atra nirdeśaḥ tantram . \\
\noindent \textbf{(P\_3,4.21.2)} evam api lokavijñānāt na sidhyati . tat yathā . \\
\noindent \textbf{(P\_3,4.21.2)} loke brāhmaṇānām pūrvam ānīyatām iti ukte sarvapūrvaḥ ānīyate . \\
\noindent \textbf{(P\_3,4.21.2)} evam iha api sarvapūrvāyāḥ kriyāyāḥ prāpnoti . \\
\noindent \textbf{(P\_3,4.21.2)} anantyavacanāt tu siddham . \\
\noindent \textbf{(P\_3,4.21.2)} samānakartṛkayoḥ anantyasya iti vaktavyam . \\
\noindent \textbf{(P\_3,4.21.2)} sidhyati . \\
\noindent \textbf{(P\_3,4.21.2)} sūtram tarhi bhidyate . \\
\noindent \textbf{(P\_3,4.21.2)} yathānyāsam eva astu . \\
\noindent \textbf{(P\_3,4.21.2)} nanu ca uktam samānakartṛkayoḥ iti bahuṣu aprāptiḥ iti . \\
\noindent \textbf{(P\_3,4.21.2)} parihṛtam etat siddham tu kriyāpradhanatvāt iti . \\
\noindent \textbf{(P\_3,4.21.2)} nanu ca uktam evam api lokavijñānāt na sidhyati iti . \\
\noindent \textbf{(P\_3,4.21.2)} na eṣaḥ doṣaḥ sarveṣām atra vrajikriyām prati paurvakālyam . \\
\noindent \textbf{(P\_3,4.21.2)} snātvā vrajati bhuktvā vrajati pītvā vrajati iti . \\
\noindent \textbf{(P\_3,4.21.2)} evam ca kṛtvā prayogaḥ aniyataḥ bhavati . \\
\noindent \textbf{(P\_3,4.21.2)} snātvā bhuktva pītvā vrajati . \\
\noindent \textbf{(P\_3,4.21.2)} pītvā snātvā bhutvā vrajati iti . \\
\noindent \textbf{(P\_3,4.21.3)} vyādāya svapiti iti upasaṅkhyānam apūrvakālatvāt . \\
\noindent \textbf{(P\_3,4.21.3)} vyādāya svapiti iti upasaṅkhyānam kartavyam . \\
\noindent \textbf{(P\_3,4.21.3)} kim punaḥ kāraṇam na sidhyati . \\
\noindent \textbf{(P\_3,4.21.3)} apūrvakālatvāt . \\
\noindent \textbf{(P\_3,4.21.3)} pūrvam hi asau svapiti paścāt vyādadāti . \\
\noindent \textbf{(P\_3,4.21.3)} na vā svapnasya avakālatvāt . \\
\noindent \textbf{(P\_3,4.21.3)} na vā kartavyam . \\
\noindent \textbf{(P\_3,4.21.3)} kim kāraṇam . \\
\noindent \textbf{(P\_3,4.21.3)} svapnasya avakālatvāt . \\
\noindent \textbf{(P\_3,4.21.3)} avarakālaḥ svapnaḥ . \\
\noindent \textbf{(P\_3,4.21.3)} avaśyam asau vyādāya muhurtam api svapiti . \\
\noindent \textbf{(P\_3,4.24)} kim iyam prāpte vibhāṣā āhosvit aprāpte . \\
\noindent \textbf{(P\_3,4.24)} katham ca prāpte katham vā aprāpte . \\
\noindent \textbf{(P\_3,4.24)} ābhīkṣṇye iti vā nitye prāpte anyatra vā aprāpte . \\
\noindent \textbf{(P\_3,4.24)} kim ca ataḥ . \\
\noindent \textbf{(P\_3,4.24)} yadi prāpte ābhīkṣṇye aniṣṭā vibhāṣā prāpnoti anyatra ca iṣṭā na sidhyati . \\
\noindent \textbf{(P\_3,4.24)} atha aprāpte . \\
\noindent \textbf{(P\_3,4.24)} agrādiṣu aprāptavidheḥ samāsapratiṣedhaḥ . \\
\noindent \textbf{(P\_3,4.24)} agrādiṣu aprāptavidheḥ samāsapratiṣedhaḥ vaktavyaḥ . \\
\noindent \textbf{(P\_3,4.24)} saḥ tarhi vaktavyaḥ . \\
\noindent \textbf{(P\_3,4.24)} na vaktavyaḥ . \\
\noindent \textbf{(P\_3,4.24)} uktam etat amā eva avyayena iti atra evakārakaraṇasya prajojanam . \\
\noindent \textbf{(P\_3,4.24)} amā eva avyayena yat tulyavidhānam upapadam tatra samāsaḥ yathā syāt . \\
\noindent \textbf{(P\_3,4.24)} amā ca anyena ca yat tulyavidhānam upapadam tatra mā bhūt iti . \\
\noindent \textbf{(P\_3,4.26.1)} kimartham svādumi makārāntatvam nipātyate na khamuñ prakṛtaḥ saḥ anuvartiṣyate . \\
\noindent \textbf{(P\_3,4.26.1)} svādumi māntanipātanam īkārābhāvārtham . \\
\noindent \textbf{(P\_3,4.26.1)} svādumi māntanipātanam kriyate īkārābhāvārtham . \\
\noindent \textbf{(P\_3,4.26.1)} īkāraḥ mā bhūt iti . \\
\noindent \textbf{(P\_3,4.26.1)} svādvīm kṛtvā yavāgūm bhuṅkte . \\
\noindent \textbf{(P\_3,4.26.1)} svāduṅkāram yavāgūm bhuṅkte . \\
\noindent \textbf{(P\_3,4.26.1)} cvyantasya ca makārāntārtham . \\
\noindent \textbf{(P\_3,4.26.1)} cvyantasya ca makārāntatvam nipātyate . \\
\noindent \textbf{(P\_3,4.26.1)} asvādu svādu kṛtvā bhuṅkte . \\
\noindent \textbf{(P\_3,4.26.1)} svāduṅkāram bhuṅkte . \\
\noindent \textbf{(P\_3,4.26.2)} ā ca tumunaḥ samānādhikaraṇe . \\
\noindent \textbf{(P\_3,4.26.2)} ā ca tumunaḥ pratyayāḥ samānādhikaraṇe vaktavyāḥ . \\
\noindent \textbf{(P\_3,4.26.2)} kena . \\
\noindent \textbf{(P\_3,4.26.2)} anuprayogeṇa . \\
\noindent \textbf{(P\_3,4.26.2)} kim prayojanam . \\
\noindent \textbf{(P\_3,4.26.2)} svāduṅkāram yavāgūḥ bhujyate devadattena iti devadatte tṛtīyā yathā syāt . \\
\noindent \textbf{(P\_3,4.26.2)} kim ca kāraṇam na syāt . \\
\noindent \textbf{(P\_3,4.26.2)} ṇamulā abhihitaḥ kartā iti . \\
\noindent \textbf{(P\_3,4.26.2)} nanu ca bhujipratyayena anabhihitaḥ kartā iti kṛtvā anabhihitāśrayaḥ vidhiḥ bhaviṣyati tṛtīyā . \\
\noindent \textbf{(P\_3,4.26.2)} yadi sati abhidhāne ca anabhidhāne ca kutaḥ cit anabhidhānam iti kṛtvā anabhihitāśrayaḥ vidhiḥ bhaviṣyati tṛtīyā yavāgvām dvitīyā prāpnoti . \\
\noindent \textbf{(P\_3,4.26.2)} kim kāraṇam . \\
\noindent \textbf{(P\_3,4.26.2)} ṇamulā anabhihitam karma iti . \\
\noindent \textbf{(P\_3,4.26.2)} yadi punaḥ ayam karmaṇi vijñāyeta ṇa evam śakyam . \\
\noindent \textbf{(P\_3,4.26.2)} iha hi svāduṅkāram yavāgūm bhuṅkte devadattaḥ iti yavāgvām dvitīyā na syāt . \\
\noindent \textbf{(P\_3,4.26.2)} kim kāraṇam . \\
\noindent \textbf{(P\_3,4.26.2)} ṇamulā abhihitam karma iti . \\
\noindent \textbf{(P\_3,4.26.2)} nanu ca bhujipratyayena anabhihitam karma iti kṛtvā anabhihitāśrayaḥ vidhiḥ bhaviṣyati dvitīyā . \\
\noindent \textbf{(P\_3,4.26.2)} yadi sati abhidhāne ca anabhidhāne ca kutaḥ cit anabhidhānam iti kṛtvā anabhihitāśrayaḥ vidhiḥ bhaviṣyati dvitīyā devadatte tṛtīyā prāpnoti . \\
\noindent \textbf{(P\_3,4.26.2)} kim kāraṇam . \\
\noindent \textbf{(P\_3,4.26.2)} ṇamulā anabhihitaḥ kartā iti . \\
\noindent \textbf{(P\_3,4.26.2)} atha anena ktvāyām arthaḥ : paktvā odanaḥ bhujyate devadattena iti . \\
\noindent \textbf{(P\_3,4.26.2)} bāḍham arthaḥ . \\
\noindent \textbf{(P\_3,4.26.2)} devadatte tṛtīyā yathā syāt . \\
\noindent \textbf{(P\_3,4.26.2)} kim ca kāraṇam na syāt . \\
\noindent \textbf{(P\_3,4.26.2)} ktvayā abhihitaḥ kartā iti . \\
\noindent \textbf{(P\_3,4.26.2)} nanu ca bhujipratyayena anabhihitaḥ kartā iti kṛtvā anabhihitāśrayaḥ vidhiḥ bhaviṣyati tṛtīyā . \\
\noindent \textbf{(P\_3,4.26.2)} yadi sati abhidhāne ca anabhidhāne ca kutaḥ cit anabhidhānam iti kṛtvā anabhihitāśrayaḥ vidhiḥ bhaviṣyati tṛtīyā odane dvitīyā prāpnoti . \\
\noindent \textbf{(P\_3,4.26.2)} kim kāraṇam . \\
\noindent \textbf{(P\_3,4.26.2)} ktvayā anabhihitam karma iti . \\
\noindent \textbf{(P\_3,4.26.2)} yadi punaḥ ayam karmaṇi vijñāyeta ṇa evam śakyam . \\
\noindent \textbf{(P\_3,4.26.2)} iha hi paktvā odanam bhuṅkte devadattaḥ iti odane dvitīyā na syāt . \\
\noindent \textbf{(P\_3,4.26.2)} kim kāraṇam . \\
\noindent \textbf{(P\_3,4.26.2)} ktvayā abhihitam karma iti . \\
\noindent \textbf{(P\_3,4.26.2)} nanu ca bhujipratyayena anabhihitam karma iti kṛtvā anabhihitāśrayaḥ vidhiḥ bhaviṣyati dvitīyā . \\
\noindent \textbf{(P\_3,4.26.2)} yadi sati abhidhāne ca anabhidhāne ca kutaḥ cit anabhidhānam iti kṛtvā anabhihitāśrayaḥ vidhiḥ bhaviṣyati dvitīyā devadatte tṛtīyā prāpnoti . \\
\noindent \textbf{(P\_3,4.26.2)} kim kāraṇam . \\
\noindent \textbf{(P\_3,4.26.2)} ktvayā anabhihitaḥ kartā iti . \\
\noindent \textbf{(P\_3,4.26.2)} atha anena tumuni arthaḥ . \\
\noindent \textbf{(P\_3,4.26.2)} bhoktum odanaḥ pacyate devadattena . \\
\noindent \textbf{(P\_3,4.26.2)} bāḍham arthaḥ . \\
\noindent \textbf{(P\_3,4.26.2)} devadatte tṛtīyā yathā syāt . \\
\noindent \textbf{(P\_3,4.26.2)} kim ca kāraṇam na syāt . \\
\noindent \textbf{(P\_3,4.26.2)} tumunā abhihitaḥ kartā iti . \\
\noindent \textbf{(P\_3,4.26.2)} nanu ca pacipratyayena anabhihitaḥ kartā iti kṛtvā anabhihitāśrayaḥ vidhiḥ bhaviṣyati tṛtīyā . \\
\noindent \textbf{(P\_3,4.26.2)} yadi sati abhidhāne ca anabhidhāne ca kutaḥ cit anabhidhānam iti kṛtvā anabhihitāśrayaḥ vidhiḥ bhaviṣyati tṛtīyā odane dvitīyā prāpnoti . \\
\noindent \textbf{(P\_3,4.26.2)} kim kāraṇam . \\
\noindent \textbf{(P\_3,4.26.2)} tumunā anabhihitam karma iti . \\
\noindent \textbf{(P\_3,4.26.2)} yadi punaḥ ayam karmaṇi vijñāyeta ṇa evam śakyam . \\
\noindent \textbf{(P\_3,4.26.2)} iha hi bhoktum odanam pacati devadattaḥ iti odane dvitīyā na syāt . \\
\noindent \textbf{(P\_3,4.26.2)} kim kāraṇam . \\
\noindent \textbf{(P\_3,4.26.2)} tumunā abhihitam karma iti . \\
\noindent \textbf{(P\_3,4.26.2)} nanu ca pacipratyayena anabhihitam karma iti kṛtvā anabhihitāśrayaḥ vidhiḥ bhaviṣyati dvitīyā . \\
\noindent \textbf{(P\_3,4.26.2)} yadi sati abhidhāne ca anabhidhāne ca kutaḥ cit anabhidhānam iti kṛtvā anabhihitāśrayaḥ vidhiḥ bhaviṣyati dvitīyā devadatte tṛtīyā prāpnoti . \\
\noindent \textbf{(P\_3,4.26.2)} kim kāraṇam . \\
\noindent \textbf{(P\_3,4.26.2)} tumunā anabhihitaḥ kartā iti . \\
\noindent \textbf{(P\_3,4.26.2)} atha anena iha arthaḥ paktvā odanam grāmaḥ gamyate devadattena . \\
\noindent \textbf{(P\_3,4.26.2)} bāḍham arthaḥ . \\
\noindent \textbf{(P\_3,4.26.2)} devadatte tṛtīyā yathā syāt . \\
\noindent \textbf{(P\_3,4.26.2)} kim ca kāraṇam na syāt . \\
\noindent \textbf{(P\_3,4.26.2)} ktvayā abhihitaḥ kartā iti . \\
\noindent \textbf{(P\_3,4.26.2)} nanu ca gamipratyayena anabhihitaḥ kartā iti kṛtvā anabhihitāśrayaḥ vidhiḥ bhaviṣyati tṛtīyā . \\
\noindent \textbf{(P\_3,4.26.2)} yadi sati abhidhāne ca anabhidhāne ca kutaḥ cit anabhidhānam iti kṛtvā anabhihitāśrayaḥ vidhiḥ bhaviṣyati tṛtīyā yat uktam odane dvitīyā prāpnoti iti saḥ doṣaḥ na jāyate . \\
\noindent \textbf{(P\_3,4.26.2)} tat tarhi vaktavyam ā ca tumunaḥ samānādhikaraṇe iti . \\
\noindent \textbf{(P\_3,4.26.2)} na vaktavyam . \\
\noindent \textbf{(P\_3,4.26.2)} avyayakṛtaḥ bhāve bhavanti iti bhāve bhaviṣyanti . \\
\noindent \textbf{(P\_3,4.26.2)} kim vaktavyam etat . \\
\noindent \textbf{(P\_3,4.26.2)} na hi . \\
\noindent \textbf{(P\_3,4.26.2)} katham anucyamānam gaṃsyate . \\
\noindent \textbf{(P\_3,4.26.2)} tumarthe iti vartate . \\
\noindent \textbf{(P\_3,4.26.2)} tumarthaḥ ca kaḥ . \\
\noindent \textbf{(P\_3,4.26.2)} bhāvaḥ . \\
\noindent \textbf{(P\_3,4.32)} ūlopaścāsyānyatarasyāṅgrahaṇam śakyam akartum . \\
\noindent \textbf{(P\_3,4.32)} katham goṣpadam vṛṣṭaḥ devaḥ iti . \\
\noindent \textbf{(P\_3,4.32)} prātiḥ pūraṇakarmā . \\
\noindent \textbf{(P\_3,4.32)} tasmāt eṣaḥ kaḥ . \\
\noindent \textbf{(P\_3,4.32)} yadi kaḥ vibhatīnām śravaṇam prāpnoti . \\
\noindent \textbf{(P\_3,4.32)} śrūyante eva atra vibhaktayaḥ . \\
\noindent \textbf{(P\_3,4.32)} tat yathā ekena goṣpadapreṇa . \\
\noindent \textbf{(P\_3,4.37)} hanaḥ karaṇe anarthakam vacanam hiṃsārthebhyaḥ ṇamulvidhānāt . \\
\noindent \textbf{(P\_3,4.37)} hanaḥ karaṇe anarthakam vacanam . \\
\noindent \textbf{(P\_3,4.37)} kim kāraṇam . \\
\noindent \textbf{(P\_3,4.37)} hiṃsārthebhyaḥ ṇamulvidhānāt . \\
\noindent \textbf{(P\_3,4.37)} hiṃsārthebhyaḥ ṇamulvidhīyate . \\
\noindent \textbf{(P\_3,4.37)} tena eva siddham . \\
\noindent \textbf{(P\_3,4.37)} arthavat tu ahiṃsārthasya vidhānāt . \\
\noindent \textbf{(P\_3,4.37)} arthavat tu hanteḥ ṇamulvacanam . \\
\noindent \textbf{(P\_3,4.37)} kaḥ arthaḥ . \\
\noindent \textbf{(P\_3,4.37)} ahiṃsārthasya vidhānāt . \\
\noindent \textbf{(P\_3,4.37)} ahiṃsārthānām ṇamul yathā syāt . \\
\noindent \textbf{(P\_3,4.37)} asti punaḥ ayam kva cit hantiḥ ahiṃsārthaḥ yadarthaḥ vidhiḥ syāt . \\
\noindent \textbf{(P\_3,4.37)} asti iti āha . \\
\noindent \textbf{(P\_3,4.37)} pāṇyupaghātam vedim hanti . \\
\noindent \textbf{(P\_3,4.37)} nityasamāsārtham ca . \\
\noindent \textbf{(P\_3,4.37)} nityasamāsārtham ca hiṃsārthāt api hanteḥ anena vidhiḥ eṣitavyaḥ . \\
\noindent \textbf{(P\_3,4.37)} katham punaḥ icchatā api hiṃsārthāt hanteḥ anena vidhiḥ labhyaḥ . \\
\noindent \textbf{(P\_3,4.37)} anena astu tena vā iti tena syāt vipratiṣedhena . \\
\noindent \textbf{(P\_3,4.37)} hanteḥ pūrvavipratiṣedhaḥ vārttikena eva jñāpitaḥ . \\
\noindent \textbf{(P\_3,4.37)} yat ayam nityasamāsārtham ca iti āha tat jñāpayati ācāryaḥ hiṃsārthāt api hanteḥ anena vidhiḥ bhavati iti . \\
\noindent \textbf{(P\_3,4.41)} iha kasmāt na bhavati . \\
\noindent \textbf{(P\_3,4.41)} grāme baddhaḥ iti . \\
\noindent \textbf{(P\_3,4.41)} evam vakṣyāmi . \\
\noindent \textbf{(P\_3,4.41)} adhikaraṇe bandhaḥ sañjñāyām . \\
\noindent \textbf{(P\_3,4.41)} tataḥ kartroḥ jīvapuruṣayoḥ naśivahoḥ iti . \\
\noindent \textbf{(P\_3,4.41)} katham aṭṭālikābandham baddhaḥ caṇḍālikābanadham baddhaḥ . \\
\noindent \textbf{(P\_3,4.41)} upamāne karmaṇi ca iti evam bhaviṣyati . \\
\noindent \textbf{(P\_3,4.60)} ayuktaḥ ayam nirdeśaḥ . \\
\noindent \textbf{(P\_3,4.60)} tiraści iti bhavitavyam . \\
\noindent \textbf{(P\_3,4.60)} sautraḥ ayam nirdeśaḥ . \\
\noindent \textbf{(P\_3,4.62)} arthagrahaṇam kimartham . \\
\noindent \textbf{(P\_3,4.62)} nādhāpratyaye iti iyati ucyamāne iha eva syāt dvidhākṛtya . \\
\noindent \textbf{(P\_3,4.62)} iha na syāt dvaidhaṅkṛtya . \\
\noindent \textbf{(P\_3,4.62)} arthagrahaṇe punaḥ kriyamāṇe na doṣaḥ bhavati . \\
\noindent \textbf{(P\_3,4.62)} nādhāpratyaye siddham bhavati yaḥ ca anyaḥ tena samānārthaḥ . \\
\noindent \textbf{(P\_3,4.62)} atha pratyayagrahaṇam kimartham iha mā bhūt hiruk kṛtvā pṛthak kṛtvā \\
\noindent \textbf{(P\_3,4.64)} ayuktaḥ ayam nirdeśaḥ . \\
\noindent \textbf{(P\_3,4.64)} anūci iti bhavitavyam . \\
\noindent \textbf{(P\_3,4.64)} sautraḥ ayam nirdeśaḥ . \\
\noindent \textbf{(P\_3,4.67.1)} kimartham idam ucyate . \\
\noindent \textbf{(P\_3,4.67.1)} kartari kṛdvacanam anādeśe svāṛthavijñānāt . \\
\noindent \textbf{(P\_3,4.67.1)} kartari kṛtaḥ bhavanti iti ucyate anādeśe svāṛthavijñānāt . \\
\noindent \textbf{(P\_3,4.67.1)} anirdiṣṭārthāḥ pratyayāḥ svārthe bhavanti . \\
\noindent \textbf{(P\_3,4.67.1)} tat yathā . \\
\noindent \textbf{(P\_3,4.67.1)} guptijkidbhyaḥ san yāvādibhyaḥ kan iti . \\
\noindent \textbf{(P\_3,4.67.1)} evam ime api pratyayāḥ svārthe syuḥ . \\
\noindent \textbf{(P\_3,4.67.1)} svārthe mā bhūvan kartari yathā syuḥ iti evamartham idam ucyate . \\
\noindent \textbf{(P\_3,4.67.1)} na etat asti prayojanam . \\
\noindent \textbf{(P\_3,4.67.1)} yam icchati svārthe āha tam . \\
\noindent \textbf{(P\_3,4.67.1)} bhāve ghañ bhavati iti . \\
\noindent \textbf{(P\_3,4.67.1)} karmaṇi tarhi mā bhūvan iti . \\
\noindent \textbf{(P\_3,4.67.1)} karmaṇi api yam icchati āha tam . \\
\noindent \textbf{(P\_3,4.67.1)} dhaḥ karmaṇi ṣṭran iti . \\
\noindent \textbf{(P\_3,4.67.1)} karaṇādhikaraṇayoḥ tarhi mā bhūvan iti . \\
\noindent \textbf{(P\_3,4.67.1)} karaṇādhikaraṇayoḥ api yam icchati āha tam . \\
\noindent \textbf{(P\_3,4.67.1)} lyuṭ karaṇādhikaraṇayoḥ bhavati iti . \\
\noindent \textbf{(P\_3,4.67.1)} sampradānāpādānayoḥ tarhi mā bhūvan iti . \\
\noindent \textbf{(P\_3,4.67.1)} sampradānāpādānayoḥ api yam icchati āha tam . \\
\noindent \textbf{(P\_3,4.67.1)} dāśagoghnau sampradāne bhīmādayaḥ apādāne iti . \\
\noindent \textbf{(P\_3,4.67.1)} yaḥ idānīm anyaḥ pratyayaḥ śeṣaḥ saḥ antareṇa vacanam kartari eva bhaviṣyati . \\
\noindent \textbf{(P\_3,4.67.1)} tat eva tarhi prayojanam svārthe mā bhūvan iti . \\
\noindent \textbf{(P\_3,4.67.1)} nanu ca uktam yam icchati svārthe āha tam . \\
\noindent \textbf{(P\_3,4.67.1)} bhāve ghañ bhavati iti . \\
\noindent \textbf{(P\_3,4.67.1)} anyaḥ saḥ bhāvaḥ bāhyaḥ prakṛtyarthāt . \\
\noindent \textbf{(P\_3,4.67.1)} anena idānīm ābhyantare bhāve syuḥ . \\
\noindent \textbf{(P\_3,4.67.1)} tatra mā bhūvan iti kartṛgrahaṇam . \\
\noindent \textbf{(P\_3,4.67.1)} kaḥ punaḥ anayoḥ bhāvayoḥ viśeṣaḥ . \\
\noindent \textbf{(P\_3,4.67.1)} uktaḥ bhāvabhedaḥ bhāṣye . \\
\noindent \textbf{(P\_3,4.67.1)} asti prayojanam etat . \\
\noindent \textbf{(P\_3,4.67.1)} kim tarhi iti . \\
\noindent \textbf{(P\_3,4.67.1)} tatra khyunādipratiṣedhaḥ nānāvākyatvāt . \\
\noindent \textbf{(P\_3,4.67.1)} tatra khyunādīnām pratiṣedhaḥ vaktavyaḥ . \\
\noindent \textbf{(P\_3,4.67.1)} khyunādayaḥ kartari mā bhūvan iti . \\
\noindent \textbf{(P\_3,4.67.1)} nanu ca karaṇe khunādayaḥ vidhīyante . \\
\noindent \textbf{(P\_3,4.67.1)} te kartari na bhaviṣyanti . \\
\noindent \textbf{(P\_3,4.67.1)} tena ca karaṇe syuḥ anena ca kartari . \\
\noindent \textbf{(P\_3,4.67.1)} nanu ca apavādatvāt khyunādayḥ bādhakāḥ syuḥ . \\
\noindent \textbf{(P\_3,4.67.1)} na syuḥ . \\
\noindent \textbf{(P\_3,4.67.1)} kim kāraṇam . \\
\noindent \textbf{(P\_3,4.67.1)} nānāvākyatvāt . \\
\noindent \textbf{(P\_3,4.67.1)} nānāvākyam tat ca idam ca . \\
\noindent \textbf{(P\_3,4.67.1)} samānavākye apavādaiḥ utsargāhḥ bādhyante . \\
\noindent \textbf{(P\_3,4.67.1)} nānāvākyatvāt bādhanam na prāpnoti . \\
\noindent \textbf{(P\_3,4.67.1)} tadvat ca kṛtyeṣu evakārakaraṇam . \\
\noindent \textbf{(P\_3,4.67.1)} evam ca kṛtvā kṛtyeṣu evakāraḥ kriyate . \\
\noindent \textbf{(P\_3,4.67.1)} tayoḥ eva kṛtyaktakhalarthāḥ iti bhāve ca akarmakebhyaḥ iti . \\
\noindent \textbf{(P\_3,4.67.1)} kim prayojanam . \\
\noindent \textbf{(P\_3,4.67.1)} tat ca bhavyādyartham . \\
\noindent \textbf{(P\_3,4.67.1)} bhavyādiṣu samāveśaḥ siddhaḥ bhavati . \\
\noindent \textbf{(P\_3,4.67.1)} geyaḥ māṇavakaḥ sāmnām . \\
\noindent \textbf{(P\_3,4.67.1)} geyāni māṇavakena sāmāni iti . \\
\noindent \textbf{(P\_3,4.67.1)} ṛṣidevatayoḥ tu kṛdbhiḥ samāveśavacanam jñāpakam asamāveśasya . \\
\noindent \textbf{(P\_3,4.67.1)} yat ayam kartari ca ṛṣidevatayoḥ iti siddhe sati samāveśe samāveśārtham cakāram śāsti tat jñāpayati ācāryaḥ na bhavati samāveśaḥ iti . \\
\noindent \textbf{(P\_3,4.67.1)} kimartham tarhi kṛtyeṣu evakāraḥ kriyate . \\
\noindent \textbf{(P\_3,4.67.1)} evakārakaraṇam ca cārthe . \\
\noindent \textbf{(P\_3,4.67.1)} evakārakaraṇam ca cārthe draṣṭavyam . \\
\noindent \textbf{(P\_3,4.67.1)} tayoḥ bhāvakarmaṇoḥ kṛtyā bhavanti bhavyādīnām kartari ca iti . \\
\noindent \textbf{(P\_3,4.67.1)} kim prayojanam . \\
\noindent \textbf{(P\_3,4.67.1)} tat ca bhavyādyartham . \\
\noindent \textbf{(P\_3,4.67.1)} bhavyādiṣu samāveśaḥ siddhaḥ bhavati . \\
\noindent \textbf{(P\_3,4.67.1)} geyaḥ māṇavakaḥ sāmnām . \\
\noindent \textbf{(P\_3,4.67.1)} geyāni māṇavakena sāmāni iti . \\
\noindent \textbf{(P\_3,4.67.1)} yat tāvat ucyate ṛṣidevatayoḥ tu kṛdbhiḥ samāveśavacanam jñāpakam asamāveśasya iti . \\
\noindent \textbf{(P\_3,4.67.1)} na etat jñāpakasādhyam apavādaiḥ utsargāḥ apavādaiḥ bādhyante iti . \\
\noindent \textbf{(P\_3,4.67.1)} eṣaḥ eva nyāyaḥ yat uta apavādaiḥ utsargāḥ bādhyeran . \\
\noindent \textbf{(P\_3,4.67.1)} nanu ca uktam nānāvākyatvāt bādhanam na prāpnoti iti . \\
\noindent \textbf{(P\_3,4.67.1)} na videśastham iti kṛtvā nānāvākyam bhavati . \\
\noindent \textbf{(P\_3,4.67.1)} videśastham api sat ekavākyam bhavati . \\
\noindent \textbf{(P\_3,4.67.1)} tat yathā dvitīye adhyāye luk ucyate . \\
\noindent \textbf{(P\_3,4.67.1)} tasya caturthaṣaṣṭhayoḥ aluk ucyate apavādaḥ . \\
\noindent \textbf{(P\_3,4.67.1)} yat api ucyate evakārakaraṇam ca cārthe iti . \\
\noindent \textbf{(P\_3,4.67.1)} katham punaḥ anyaḥ nāma anyasya arthe vartate . \\
\noindent \textbf{(P\_3,4.67.1)} katham evakāraḥ cārthe vartate . \\
\noindent \textbf{(P\_3,4.67.1)} saḥ eṣaḥ evakāraḥ svārthe vartate . \\
\noindent \textbf{(P\_3,4.67.1)} kim prayojanam . \\
\noindent \textbf{(P\_3,4.67.1)} jñāpakārtham . \\
\noindent \textbf{(P\_3,4.67.1)} etat jñāpayati acāryaḥ itaḥ uttaram samāveśaḥ bhavati iti . \\
\noindent \textbf{(P\_3,4.67.1)} kim etasya jñapane prayojanam . \\
\noindent \textbf{(P\_3,4.67.1)} tat ca bhavyādyartham . \\
\noindent \textbf{(P\_3,4.67.1)} bhavyādiṣu samāveśaḥ siddhaḥ bhavati . \\
\noindent \textbf{(P\_3,4.67.1)} geyaḥ māṇavakaḥ sāmnām . \\
\noindent \textbf{(P\_3,4.67.1)} geyāni māṇavakena sāmāni iti . \\
\noindent \textbf{(P\_3,4.67.1)} yadi etat jñapyate iha api samāveśaḥ prāpnoti dāśagoghnau sampradāne bhīmādayaḥ apādāne iti . \\
\noindent \textbf{(P\_3,4.67.1)} atra api siddham bhavati . \\
\noindent \textbf{(P\_3,4.67.1)} yat ayam ādikarmaṇi ktaḥ kartari ca iti siddhe samāveśe samāveśam śāsti tat jñapayati ācāryaḥ prāk amutaḥ samāveśaḥ bhavati iti . \\
\noindent \textbf{(P\_3,4.67.2)} kim punaḥ ayam pratyayaniyamaḥ : dhātoḥ paraḥ akāraḥ akaśabdaḥ vā niyogataḥ kartāram bruvan kṛtsañjñaḥ ca bhavati pratyayasañjñaḥ ca iti . \\
\noindent \textbf{(P\_3,4.67.2)} āhosvit sañjñāniyamaḥ : dhātoḥ paraḥ akāraḥ akaśabdaḥ vā svabhāvataḥ kartāram bruvan kṛtsañjñaḥ ca bhavati pratyayasañjñaḥ ca iti . \\
\noindent \textbf{(P\_3,4.67.2)} kaḥ ca atra viśeṣaḥ . \\
\noindent \textbf{(P\_3,4.67.2)} tatra pratyayaniyame aniṣṭaprasaṅgaḥ . \\
\noindent \textbf{(P\_3,4.67.2)} tatra pratyayaniyame sati aniṣṭam prāpnoti . \\
\noindent \textbf{(P\_3,4.67.2)} kāṣṭhabhit abrāhmaṇaḥ , balabhit abrāhmaṇaḥ . \\
\noindent \textbf{(P\_3,4.67.2)} eṣaḥ api niyogataḥ kartāram bruvan kṛtsañjñaḥ ca syāt pratyayasañjñaḥ ca . \\
\noindent \textbf{(P\_3,4.67.2)} sañjñāniyame siddham . \\
\noindent \textbf{(P\_3,4.67.2)} sañjñāniyame sati siddham bhavati . \\
\noindent \textbf{(P\_3,4.67.2)} yadi sañjñāniyamaḥ vibhaktādiṣu doṣaḥ . \\
\noindent \textbf{(P\_3,4.67.2)} vibhaktāḥ bhrātaraḥ pītāḥ gāvaḥ iti na sidhyati . \\
\noindent \textbf{(P\_3,4.67.2)} pratyayaniyame punaḥ sati parigaṇitābhyaḥ prakṛtibhyaḥ paraḥ ktaḥ niyogataḥ kartāram āha . \\
\noindent \textbf{(P\_3,4.67.2)} na ca imāḥ tatra parigaṇyante prakṛtayaḥ . \\
\noindent \textbf{(P\_3,4.67.2)} vibhaktādiṣu ca aprāptiḥ prakṛteḥ pratyayaparavacanāt . \\
\noindent \textbf{(P\_3,4.67.2)} vibhaktādiṣu ca pratyayaniyamasya aprāptiḥ . \\
\noindent \textbf{(P\_3,4.67.2)} kim kāraṇam . \\
\noindent \textbf{(P\_3,4.67.2)} prakṛteḥ pratyayaparavacanāt . \\
\noindent \textbf{(P\_3,4.67.2)} parigaṇitābhyaḥ prakṛtibhyaḥ paraḥ ktaḥ svabhāvataḥ kartāram āha . \\
\noindent \textbf{(P\_3,4.67.2)} na ca imāḥ tatra parigaṇyante . \\
\noindent \textbf{(P\_3,4.67.2)} na tarhi idānīm ayam sādhuḥ bhavati . \\
\noindent \textbf{(P\_3,4.67.2)} bhavati sādhuḥ na tu kartari . \\
\noindent \textbf{(P\_3,4.67.2)} katham tarhi idānīm atra kartṛtvam gamyate . \\
\noindent \textbf{(P\_3,4.67.2)} akāraḥ matvarthīyaḥ : vibhaktam eṣām asti vibhaktāḥ . \\
\noindent \textbf{(P\_3,4.67.2)} pītam eṣām asti pitāḥ iti . \\
\noindent \textbf{(P\_3,4.67.2)} atha vā uttarapadalopaḥ atra draṣṭavyaḥ . \\
\noindent \textbf{(P\_3,4.67.2)} vibhaktadhanāḥ vibhaktāḥ . \\
\noindent \textbf{(P\_3,4.67.2)} pītodakāḥ pitāḥ iti . \\
\noindent \textbf{(P\_3,4.69)} kimartham idam ucyate . \\
\noindent \textbf{(P\_3,4.69)} laḥ eṣu sādhaneṣu yathā syāt kartari ca karmaṇi ca bhāve ca akarmakebhyaḥ iti . \\
\noindent \textbf{(P\_3,4.69)} na etat asti prayojanam . \\
\noindent \textbf{(P\_3,4.69)} bhāvakarmaṇoḥ ātmanepadam vidhīyate śeṣāt kartari parasmaipadam . \\
\noindent \textbf{(P\_3,4.69)} etāvān ca laḥ yat uta parasmaipadam ātmanepadam ca . \\
\noindent \textbf{(P\_3,4.69)} saḥ ca ayam evam vihitaḥ . \\
\noindent \textbf{(P\_3,4.69)} ataḥ uttaram paṭhati . \\
\noindent \textbf{(P\_3,4.69)} lagrahaṇam sakarmakanivṛttyartham . \\
\noindent \textbf{(P\_3,4.69)} lagrahaṇam kriyate sakarmakanivṛttyartham . \\
\noindent \textbf{(P\_3,4.69)} sakarmakāṇām bhāve laḥ mā bhūte iti . \\
\noindent \textbf{(P\_3,4.69)} yadi punaḥ tatra eva akarmakagrahaṇam kriyeta . \\
\noindent \textbf{(P\_3,4.69)} tatra akarmakagrahaṇam kartavyam . \\
\noindent \textbf{(P\_3,4.69)} nanu ca iha api kriyate bhāve ca akarmakebhyaḥ iti . \\
\noindent \textbf{(P\_3,4.69)} parārtham etat bhaviṣyati . \\
\noindent \textbf{(P\_3,4.69)} tayoḥ eva kṛtyaktakhalarthāḥ bhāve ca akarmakebhyaḥ . \\
\noindent \textbf{(P\_3,4.69)} yāvat iha lagrahaṇam tāvat tatra akarmakagrahaṇam . \\
\noindent \textbf{(P\_3,4.69)} iha vā lagrahaṇam kriyeta tatra vā akarmakagrahaṇam . \\
\noindent \textbf{(P\_3,4.69)} kaḥ nu atra viśeṣaḥ . \\
\noindent \textbf{(P\_3,4.69)} ayam asti viśeṣaḥ . \\
\noindent \textbf{(P\_3,4.69)} iha lagrahaṇe kriyamāṇe ānaḥ kartari siddhaḥ bhavati . \\
\noindent \textbf{(P\_3,4.69)} tatra punaḥ akarmakagrahaṇe kriyamāṇe ānaḥ kartari na prāpnoti . \\
\noindent \textbf{(P\_3,4.69)} tatra api akarmakagrahaṇe kriyamāṇe ānaḥ kartari siddhaḥ bhavati . \\
\noindent \textbf{(P\_3,4.69)} katham . \\
\noindent \textbf{(P\_3,4.69)} bhāvakarmaṇoḥ iti ataḥ anyat yat ātmanepadānukramaṇam sarvam tat kartrartham . \\
\noindent \textbf{(P\_3,4.69)} vipratiṣedhāt vā ānaḥ kartari . \\
\noindent \textbf{(P\_3,4.69)} vipratiṣedhāt vā ānaḥ kartari bhaviṣyati . \\
\noindent \textbf{(P\_3,4.69)} tatra bhāvakarmaṇoḥ iti etat astu kartari kṛt iti . \\
\noindent \textbf{(P\_3,4.69)} kartari kṛt iti etat bhaviṣyati vipratiṣedhena . \\
\noindent \textbf{(P\_3,4.69)} sarvaprasaṅgaḥ tu . \\
\noindent \textbf{(P\_3,4.69)} sarvebhyaḥ tu dhātubhyaḥ ānaḥ kartari prāpnoti . \\
\noindent \textbf{(P\_3,4.69)} parasmaipadibhyaḥ api . \\
\noindent \textbf{(P\_3,4.69)} na eṣaḥ doṣaḥ . \\
\noindent \textbf{(P\_3,4.69)} anudāttaṅitaḥ iti eṣaḥ yogaḥ niyamārthaḥ bhaviṣyati . \\
\noindent \textbf{(P\_3,4.69)} yadi eṣaḥ yogaḥ niyamārthaḥ vidhiḥ na prakalpate . \\
\noindent \textbf{(P\_3,4.69)} āste śete iti . \\
\noindent \textbf{(P\_3,4.69)} atha vidhyarthaḥ ānasya niyamaḥ na prāpnoti . \\
\noindent \textbf{(P\_3,4.69)} āsīnaḥ śayānaḥ . \\
\noindent \textbf{(P\_3,4.69)} tathā neḥ viśaḥ iti evamādi anukramaṇam yadi niyamāṛthaḥ vidhiḥ na prakalpate . \\
\noindent \textbf{(P\_3,4.69)} atha vidhyarthaḥ ānasya niyamaḥ na prāpnoti . \\
\noindent \textbf{(P\_3,4.69)} astu tarhi niyamārtham . \\
\noindent \textbf{(P\_3,4.69)} nanu ca uktam vidhiḥ na prakalpate iti . \\
\noindent \textbf{(P\_3,4.69)} vidhiḥ ca prakḷptaḥ . \\
\noindent \textbf{(P\_3,4.69)} katham . \\
\noindent \textbf{(P\_3,4.69)} bhāvakarmaṇoḥ iti atra anudāttaṅitaḥ iti etat anuvartiṣyate . \\
\noindent \textbf{(P\_3,4.69)} yadi anuvartate evam api anudāttaṅitaḥ eva bhāvakarmaṇoḥ ātmanepadam prāpnoti . \\
\noindent \textbf{(P\_3,4.69)} evam tarhi yogavibhāgaḥ kariṣyate . \\
\noindent \textbf{(P\_3,4.69)} anudāttaṅitaḥ ātmanepadam bhavati . \\
\noindent \textbf{(P\_3,4.69)} tataḥ bhāvakarmaṇoḥ . \\
\noindent \textbf{(P\_3,4.69)} tataḥ kartari . \\
\noindent \textbf{(P\_3,4.69)} kartari ca ātmanepadam bhavati anudāttaṅitaḥ iti eva . \\
\noindent \textbf{(P\_3,4.69)} bhāvakarmaṇoḥ iti nivṛttam . \\
\noindent \textbf{(P\_3,4.69)} tataḥ karmavyatihāre . \\
\noindent \textbf{(P\_3,4.69)} kartari iti eva anuvartate . \\
\noindent \textbf{(P\_3,4.69)} anudāttaṅitaḥ iti api nivṛttam . \\
\noindent \textbf{(P\_3,4.69)} yat api ucyate neḥ viśaḥ iti evamādi anukramaṇam yadi niyamārtham vidhiḥ na prakalpate . \\
\noindent \textbf{(P\_3,4.69)} atha vidhyarthaḥ ānasya niyamaḥ na prāpnoti iti . \\
\noindent \textbf{(P\_3,4.69)} astu vidhyartham . \\
\noindent \textbf{(P\_3,4.69)} nanu ca uktam ānasya niyamaḥ na prāpnoti iti . \\
\noindent \textbf{(P\_3,4.69)} na eṣaḥ doṣaḥ . \\
\noindent \textbf{(P\_3,4.69)} yathā eva atra aprāptāḥ taṅaḥ bhavanti evam ānaḥ api bhaviṣyati . \\
\noindent \textbf{(P\_3,4.69)} sarvatra aprasaṅgaḥ tu . \\
\noindent \textbf{(P\_3,4.69)} sarveṣu tu sādhaneṣu ānaḥ na prāpnoti . \\
\noindent \textbf{(P\_3,4.69)} vipratiṣedhāt vā ānaḥ kartari iti bhāvakarmaṇoḥ na syāt . \\
\noindent \textbf{(P\_3,4.69)} kartari eva syāt . \\
\noindent \textbf{(P\_3,4.69)} iha punaḥ lagrahaṇe kriyamāṇe kartari kṛt iti etat astu laḥ karmaṇi ca bhāve ca akarmakebhyaḥ iti laḥ karmaṇi ca bhāve ca akarmakebhyaḥ iti etat bhaviṣyat vipratiṣedhena . \\
\noindent \textbf{(P\_3,4.69)} sarvaprasaṅgaḥ tu . \\
\noindent \textbf{(P\_3,4.69)} lādeśaḥ sarveṣu sādhaneṣu prāpnoti . \\
\noindent \textbf{(P\_3,4.69)} śatṛkvasūca bhāvakarmaṇoḥ api prāpnutaḥ . \\
\noindent \textbf{(P\_3,4.69)} na eṣaḥ doṣaḥ . \\
\noindent \textbf{(P\_3,4.69)} śeṣāt parasmaipadam kartari iti evam tau kartāram hriyete . \\
\noindent \textbf{(P\_3,4.77.1)} lādeśe sarvaprasaṅgaḥ aviśeṣāt ḷādeśe sarvaprasaṅgaḥ . \\
\noindent \textbf{(P\_3,4.77.1)} sarvasya lakārasya ādeśaḥ prāpnoti . \\
\noindent \textbf{(P\_3,4.77.1)} asya api prāpnoti : lunāti labhate . \\
\noindent \textbf{(P\_3,4.77.1)} kim kāraṇam . \\
\noindent \textbf{(P\_3,4.77.1)} aviśeṣāt . \\
\noindent \textbf{(P\_3,4.77.1)} na hi kaḥ cit viśeṣaḥ upādīyate : evañjātīyakasya lakārasya ādeśaḥ bhavati iti . \\
\noindent \textbf{(P\_3,4.77.1)} anupādīyamāne viśeṣe sarvaprasaṅgaḥ . \\
\noindent \textbf{(P\_3,4.77.1)} arthavadgrahaṇāt siddham . \\
\noindent \textbf{(P\_3,4.77.1)} arthavataḥ lakārasya grahaṇam na ca eṣaḥ artahvat . \\
\noindent \textbf{(P\_3,4.77.1)} arthavadgrahaṇāt siddham iti cet na varṇagrahaṇeṣu . arthavadgrahaṇāt siddham iti cet tat na . \\
\noindent \textbf{(P\_3,4.77.1)} kim kāraṇam . \\
\noindent \textbf{(P\_3,4.77.1)} varṇagrahaṇam idam . \\
\noindent \textbf{(P\_3,4.77.1)} na ca etat varṇagrahaṇeṣu bhavati arthavadgrahaṇe na anarthakasya iti . \\
\noindent \textbf{(P\_3,4.77.1)} tasmāt viśiṣtagrahaṇam . \\
\noindent \textbf{(P\_3,4.77.1)} tasmāt viśiṣtasya lakārasya grahaṇam kartavyam . \\
\noindent \textbf{(P\_3,4.77.1)} na kartavyam . \\
\noindent \textbf{(P\_3,4.77.1)} dhātoḥ iti vartate . \\
\noindent \textbf{(P\_3,4.77.1)} evam api śālā mālā mallaḥ iti atra prāpnoti . \\
\noindent \textbf{(P\_3,4.77.1)} uṇādayaḥ avyutpannāni prātipadikāni . \\
\noindent \textbf{(P\_3,4.77.1)} evam api nandanaḥ atra prāpnoti . \\
\noindent \textbf{(P\_3,4.77.1)} itsañjñā atra bādhikā bhaviṣyati . \\
\noindent \textbf{(P\_3,4.77.1)} iha api tarhi bādheta . \\
\noindent \textbf{(P\_3,4.77.1)} pacati paṭhati iti . \\
\noindent \textbf{(P\_3,4.77.1)} itkāryābhāvāt atra itsañjñā na bhaviṣyati . \\
\noindent \textbf{(P\_3,4.77.1)} idam asti itkāryam liti pratyayāt pūrvam udāttam bhavati iti eṣaḥ svaraḥ yathā syāt . \\
\noindent \textbf{(P\_3,4.77.1)} liti iti ucyate . \\
\noindent \textbf{(P\_3,4.77.1)} na ca atra litam paśyāmaḥ . \\
\noindent \textbf{(P\_3,4.77.1)} atha api katham cit vacanāt vā anuvartanāt vā itsañjñkānām ādeśaḥ syāt evam api na doṣaḥ . \\
\noindent \textbf{(P\_3,4.77.1)} ācāryapravṛttiḥ jñāpayati na lādeśe litkāryam bhavati iti yat ayam ṇalam litam karoti . \\
\noindent \textbf{(P\_3,4.77.1)} atha api uṇādayaḥ vyutpādyante evam api no doṣaḥ . \\
\noindent \textbf{(P\_3,4.77.1)} kriyate viśiṣṭagrahaṇam lasya iti . \\
\noindent \textbf{(P\_3,4.77.2)} lādeśaḥ varṇavidheḥ pūrvavipratiṣiddham . lādeśaḥ varṇavidheḥ bhavati pūrvavipratiṣedhena . \\
\noindent \textbf{(P\_3,4.77.2)} lādeśasya avakāśaḥ pacatu paṭhatu . \\
\noindent \textbf{(P\_3,4.77.2)} varṇavidheḥ avakāśaḥ dadhyatra madhvatra . \\
\noindent \textbf{(P\_3,4.77.2)} iha ubhayam prāpnoti . \\
\noindent \textbf{(P\_3,4.77.2)} pacatu atra . \\
\noindent \textbf{(P\_3,4.77.2)} paṭhatu atra . \\
\noindent \textbf{(P\_3,4.77.2)} lādeśaḥ bhavati pūrvavipratiṣedhena . \\
\noindent \textbf{(P\_3,4.77.2)} saḥ tarhi pūrvavipratiṣedhaḥ vaktavyaḥ . \\
\noindent \textbf{(P\_3,4.77.2)} na vaktavyaḥ . \\
\noindent \textbf{(P\_3,4.77.2)} uktam vā . \\
\noindent \textbf{(P\_3,4.77.2)} kim uktam . \\
\noindent \textbf{(P\_3,4.77.2)} lādeśaḥ varṇavidheḥ iti . \\
\noindent \textbf{(P\_3,4.79)} ṭitaḥ etve ātmanepadeṣu ānapratiṣedhaḥ . \\
\noindent \textbf{(P\_3,4.79)} ṭitaḥ etve ātmanepadeṣu ānapratiṣedhaḥ vaktavyaḥ . \\
\noindent \textbf{(P\_3,4.79)} pacamānaḥ yajamānaḥ . \\
\noindent \textbf{(P\_3,4.79)} ṭitaḥ iti etvam prāpnoti . \\
\noindent \textbf{(P\_3,4.79)} uktam vā . \\
\noindent \textbf{(P\_3,4.79)} kim uktam . \\
\noindent \textbf{(P\_3,4.79)} jñāpakam vā sānubandhakasya ādeśavacane itkāryābhāvasya iti . \\
\noindent \textbf{(P\_3,4.79)} na etat asti uktam . \\
\noindent \textbf{(P\_3,4.79)} evam kila tat uktam syāt yadi evam vijñāyeta . \\
\noindent \textbf{(P\_3,4.79)} ṭit ātmanepadam ṭidātmanepadam . \\
\noindent \textbf{(P\_3,4.79)} ṭidātmanepadānām iti . \\
\noindent \textbf{(P\_3,4.79)} tat ca na . \\
\noindent \textbf{(P\_3,4.79)} ṭitaḥ lakārasya yāni ātmanepadāni iti evam etat vijñāyate . \\
\noindent \textbf{(P\_3,4.79)} avaśyam ca etat evam vijñeyam . \\
\noindent \textbf{(P\_3,4.79)} ṭit ātmanepadam ṭidātmanepadam . \\
\noindent \textbf{(P\_3,4.79)} ṭidātmanepadānām iti vijñāyamāne akurvi atra api prasajyeta . \\
\noindent \textbf{(P\_3,4.79)} na eṣaḥ ṭit . \\
\noindent \textbf{(P\_3,4.79)} kaḥ tarhi . \\
\noindent \textbf{(P\_3,4.79)} ṭhit . \\
\noindent \textbf{(P\_3,4.79)} saḥ ca avaśyam ṭhit kartavyaḥ ādiḥ mā bhūt iti . \\
\noindent \textbf{(P\_3,4.79)} katham iṭaḥ at iti . \\
\noindent \textbf{(P\_3,4.79)} iṭhaḥ at iti vakṣyāmi iti . \\
\noindent \textbf{(P\_3,4.79)} tat ca avaśyam vaktavyam paryavapādyasya mā bhūt . \\
\noindent \textbf{(P\_3,4.79)} laviṣīṣṭa . \\
\noindent \textbf{(P\_3,4.79)} iha tarhi iṣam ūrjam aham itaḥ ādi ātaḥ lopaḥ iṭi ca iti ākāralopaḥ na prāpnoti . \\
\noindent \textbf{(P\_3,4.79)} tasmāt ṭit eṣaḥ . \\
\noindent \textbf{(P\_3,4.79)} ādiḥ tarhi kasmāt na bhavati . \\
\noindent \textbf{(P\_3,4.79)} saptadaśa ādeśāḥ sthāneyogatvam prayojayanti . \\
\noindent \textbf{(P\_3,4.79)} tān ekaḥ na utsahate vihantum iti kṛtvā ādiḥ na bhaviṣyati . \\
\noindent \textbf{(P\_3,4.79)} paryavapādyasya kasmāt na bhavati . \\
\noindent \textbf{(P\_3,4.79)} laviṣīṣṭa iti . \\
\noindent \textbf{(P\_3,4.79)} asiddham bahiraṅgalakṣaṇam antaraṅgalakṣaṇe iti . \\
\noindent \textbf{(P\_3,4.79)} idam tarhi uktam prākṛtānām ātmanepadānām etvam bhavati iti . \\
\noindent \textbf{(P\_3,4.79)} ke ca prakṛtāḥ . \\
\noindent \textbf{(P\_3,4.79)} tādayaḥ . \\
\noindent \textbf{(P\_3,4.79)} āne muk jñāpakam tu etve ṭittaṅām . \\
\noindent \textbf{(P\_3,4.79)} iśisīricaḥ ḍārauraḥsu . \\
\noindent \textbf{(P\_3,4.79)} ṭit aṭitaḥ . \\
\noindent \textbf{(P\_3,4.79)} prakṛte tat . \\
\noindent \textbf{(P\_3,4.79)} guṇe katham . \\
\noindent \textbf{(P\_3,4.82.1)} ṇalaḥ śitkaraṇam sarvādeśārtham . \\
\noindent \textbf{(P\_3,4.82.1)} ṇal śit kartavyaḥ . \\
\noindent \textbf{(P\_3,4.82.1)} kim prayojanam . \\
\noindent \textbf{(P\_3,4.82.1)} sarvādeśārtham . \\
\noindent \textbf{(P\_3,4.82.1)} śit sarvasya iti sarvādeśaḥ yathā syāt . \\
\noindent \textbf{(P\_3,4.82.1)} akriyamāṇe hi śakāre alaḥ antyasya vidhayaḥ bhavanti iti antyasya prasajyeta . \\
\noindent \textbf{(P\_3,4.82.1)} uktam vā . \\
\noindent \textbf{(P\_3,4.82.1)} kim uktam . \\
\noindent \textbf{(P\_3,4.82.1)} anittvāt siddham iti . \\
\noindent \textbf{(P\_3,4.82.1)} ṇakāraḥ kriyate . \\
\noindent \textbf{(P\_3,4.82.1)} tasya anittvāt siddham . \\
\noindent \textbf{(P\_3,4.82.1)} kaḥ eṣaḥ parihāraḥ nyāyyaḥ . \\
\noindent \textbf{(P\_3,4.82.1)} śakāram asi coditaḥ . \\
\noindent \textbf{(P\_3,4.82.1)} ṇakāram kariṣyāmi śakāram na kariṣyāmi iti . \\
\noindent \textbf{(P\_3,4.82.1)} ṇakāraḥ atra kriyeta śakāraḥ vā kaḥ nu atra viśeṣaḥ . \\
\noindent \textbf{(P\_3,4.82.1)} avaśyam atra ṇakāraḥ vṛddhyarthaḥ kartavyaḥ ṇiti iti vṛddhiḥ yathā syāt . \\
\noindent \textbf{(P\_3,4.82.1)} na arthaḥ vṛddhyarthena ṇakāreṇa . \\
\noindent \textbf{(P\_3,4.82.1)} ṇittve yogavibhāgaḥ kariṣyate . \\
\noindent \textbf{(P\_3,4.82.1)} idam asti gotaḥ ṇit . \\
\noindent \textbf{(P\_3,4.82.1)} tataḥ al . \\
\noindent \textbf{(P\_3,4.82.1)} al ca ṇit bhavati . \\
\noindent \textbf{(P\_3,4.82.1)} tataḥ uttamaḥ vā iti . \\
\noindent \textbf{(P\_3,4.82.1)} evam tarhi lakāraḥ kriyate . \\
\noindent \textbf{(P\_3,4.82.1)} tasya anittvāt siddham . \\
\noindent \textbf{(P\_3,4.82.1)} kaḥ eṣaḥ parihāraḥ nyāyyaḥ . \\
\noindent \textbf{(P\_3,4.82.1)} śakāram asi coditaḥ . \\
\noindent \textbf{(P\_3,4.82.1)} lakāram kariṣyāmi śakāram na kariṣyāmi iti . \\
\noindent \textbf{(P\_3,4.82.1)} lakāraḥ atra kriyeta śakāraḥ vā kaḥ nu atra viśeṣaḥ . \\
\noindent \textbf{(P\_3,4.82.1)} avaśyam eva atra svarārthaḥ lakāraḥ kartavyaḥ liti pratyayāt pūrvam udāttam bhavati iti eṣaḥ svaraḥ yathā syāt . \\
\noindent \textbf{(P\_3,4.82.1)} na etat asti prayojanam . \\
\noindent \textbf{(P\_3,4.82.1)} dhātusvare kṛte dvirvacanam . \\
\noindent \textbf{(P\_3,4.82.1)} tatra āntaryataḥ antodāttasya antodāttaḥ ādeśaḥ bhaviṣyati . \\
\noindent \textbf{(P\_3,4.82.1)} katham punaḥ ayam antodāttaḥ syāt yadā ekāc . \\
\noindent \textbf{(P\_3,4.82.1)} vyapadeśivadbhāvena . \\
\noindent \textbf{(P\_3,4.82.1)} yathā eva tarhi vyapadeśivadbhāvena antodāttaḥ evam ādyudāttaḥ api . \\
\noindent \textbf{(P\_3,4.82.1)} tatra āntaryataḥ ādyudāttasya ādyudāttaḥ ādeśaḥ prasajyeta . \\
\noindent \textbf{(P\_3,4.82.1)} satyam etat . \\
\noindent \textbf{(P\_3,4.82.1)} na tu idam lakṣaṇam asti dhātoḥ ādiḥ udāttaḥ bhavati iti . \\
\noindent \textbf{(P\_3,4.82.1)} idam punaḥ asti dhātoḥ antaḥ udāttaḥ bhavati iti . \\
\noindent \textbf{(P\_3,4.82.1)} saḥ asau lakṣaṇena antodāttaḥ . \\
\noindent \textbf{(P\_3,4.82.1)} tatra āntaryataḥ antodāttasya antodāttaḥ ādeśaḥ bhaviṣyati . \\
\noindent \textbf{(P\_3,4.82.1)} etat api ādeśe na asti ādeśasya antaḥ udāttaḥ bhavati iti . \\
\noindent \textbf{(P\_3,4.82.1)} prakṛtitaḥ anena svaraḥ labhyaḥ . \\
\noindent \textbf{(P\_3,4.82.1)} prakṛtiḥ ca asya yathā eva antodāttā evam ādyudāttā api . \\
\noindent \textbf{(P\_3,4.82.1)} dviḥprayoge ca api dvirvacane ubhayoḥ antodāttatvam prasajyeta . \\
\noindent \textbf{(P\_3,4.82.1)} anudāttam padam ekavarjam iti na asti yaugapadyena sambhavaḥ . \\
\noindent \textbf{(P\_3,4.82.1)} paryāyaḥ prasajyeta . \\
\noindent \textbf{(P\_3,4.82.1)} tasmāt svarārthaḥ lakāraḥ kartavyaḥ . \\
\noindent \textbf{(P\_3,4.82.1)} lakāraḥ kriyate . \\
\noindent \textbf{(P\_3,4.82.1)} tasya anittvāt siddham . \\
\noindent \textbf{(P\_3,4.82.2)} akārasya śitkaraṇam sarvādeśārtham . \\
\noindent \textbf{(P\_3,4.82.2)} akāraḥ śitkartavyaḥ . \\
\noindent \textbf{(P\_3,4.82.2)} kim prayojanam . \\
\noindent \textbf{(P\_3,4.82.2)} sarvādeśārtham . \\
\noindent \textbf{(P\_3,4.82.2)} śit sarvasya iti sarvādeśaḥ yathā syāt . \\
\noindent \textbf{(P\_3,4.82.2)} akriyamāṇe hi śakāre alaḥ antyasya vidhayaḥ bhavanti iti antyasya prasajyeta . \\
\noindent \textbf{(P\_3,4.82.2)} nanu ca akārasya akāravacane prayojanam na asti iti kṛtvā antareṇa śakāram sarvādeśaḥ bhaviṣyati . \\
\noindent \textbf{(P\_3,4.82.2)} asti anyat akārasya akāravacane prayojanam . \\
\noindent \textbf{(P\_3,4.82.2)} kim . \\
\noindent \textbf{(P\_3,4.82.2)} akāravacanam samasaṅkhyārtham . \\
\noindent \textbf{(P\_3,4.82.2)} saṅkhyātānudeśaḥ yathā syāt . \\
\noindent \textbf{(P\_3,4.82.2)} tasmāt śitkaraṇam . \\
\noindent \textbf{(P\_3,4.82.2)} tasmāt śakāraḥ kartavyaḥ . \\
\noindent \textbf{(P\_3,4.82.2)} na kartavyaḥ . \\
\noindent \textbf{(P\_3,4.82.2)} kriyate nyāse eva . \\
\noindent \textbf{(P\_3,4.82.2)} praśliṣṭanirdeśaḥ ayam . \\
\noindent \textbf{(P\_3,4.82.2)} a* a* a . \\
\noindent \textbf{(P\_3,4.82.2)} saḥ anekālśit sarvasya iti sarvādeśaḥ bhaviṣyati . \\
\noindent \textbf{(P\_3,4.85)} laṅvadatideśe jusbhāvapratiṣedhaḥ . \\
\noindent \textbf{(P\_3,4.85)} laṅvadatideśe jusbhāvasya pratiṣedhaḥ vaktavyaḥ . \\
\noindent \textbf{(P\_3,4.85)} yāntu vāntu . \\
\noindent \textbf{(P\_3,4.85)} laṅaḥ śākaṭāyanasya eva iti jusbhāvaḥ prāpnoti . \\
\noindent \textbf{(P\_3,4.85)} utvavacanāt siddham . \\
\noindent \textbf{(P\_3,4.85)} utvam atra bādhakam bhaviṣyati . \\
\noindent \textbf{(P\_3,4.85)} anavakāśāḥ hi vidhayaḥ bādhakāḥ bhavanti . \\
\noindent \textbf{(P\_3,4.85)} sāvakāśam ca utvam . \\
\noindent \textbf{(P\_3,4.85)} kaḥ avakāśaḥ . \\
\noindent \textbf{(P\_3,4.85)} pacatu paṭhatu . \\
\noindent \textbf{(P\_3,4.85)} atra api ikāralopaḥ prāpnoti . \\
\noindent \textbf{(P\_3,4.85)} tat yathā eva utvam ikāralopam bādhate evam jusbhāvam api bādhate . \\
\noindent \textbf{(P\_3,4.85)} na bādhate . \\
\noindent \textbf{(P\_3,4.85)} kim kāraṇam . \\
\noindent \textbf{(P\_3,4.85)} yena na aprāpte tasya bādhanam bhavati . \\
\noindent \textbf{(P\_3,4.85)} na ca aprāpte ikāralope utvam ārabhyate . \\
\noindent \textbf{(P\_3,4.85)} jusbhāve punaḥ prāpte ca aprāpte ca . \\
\noindent \textbf{(P\_3,4.85)} atha vā purastāt apavādāḥ anantarān vidhīn bādhante iti evam utvam ikāralopam bādhate jubhāvam na bādhate . \\
\noindent \textbf{(P\_3,4.85)} evam tarhi vakṣyati tatra laṅgrahaṇasya prayojanam . \\
\noindent \textbf{(P\_3,4.85)} laṅ eva yaḥ laṅ tatra yathā syāt . \\
\noindent \textbf{(P\_3,4.85)} laṅvadbhāvena yaḥ laṅ tatra mā bhūt iti . \\
\noindent \textbf{(P\_3,4.87,} 89) KA\_II,185.3-9 Ro\_III,409 {1/11}     hinyoḥ utvapratiṣedhaḥ . \\
\noindent \textbf{(P\_3,4.87,} 89) KA\_II,185.3-9 Ro\_III,409 {2/11}     hinyoḥ ukārasya pratiṣedhaḥ vaktavyaḥ . \\
\noindent \textbf{(P\_3,4.87,} 89) KA\_II,185.3-9 Ro\_III,409 {3/11}     lunīhi lunāni . \\
\noindent \textbf{(P\_3,4.87,} 89) KA\_II,185.3-9 Ro\_III,409 {4/11}     eḥ uḥ iti utvam prāpnoti . \\
\noindent \textbf{(P\_3,4.87,} 89) KA\_II,185.3-9 Ro\_III,409 {5/11}     na vā uccāraṇasāmarthyāt . \\
\noindent \textbf{(P\_3,4.87,} 89) KA\_II,185.3-9 Ro\_III,409 {6/11}     na vā vaktavyaḥ . \\
\noindent \textbf{(P\_3,4.87,} 89) KA\_II,185.3-9 Ro\_III,409 {7/11}     kim kāraṇam . \\
\noindent \textbf{(P\_3,4.87,} 89) KA\_II,185.3-9 Ro\_III,409 {8/11}     uccāraṇasāmarthyāt atra utvam na bhaviṣyati . \\
\noindent \textbf{(P\_3,4.87,} 89) KA\_II,185.3-9 Ro\_III,409 {9/11}     alaghīyaḥ ca eva hi ikāroccāraṇam ukāroccāraṇāt . \\
\noindent \textbf{(P\_3,4.87,} 89) KA\_II,185.3-9 Ro\_III,409 {10/11}     ikāram ca uccārayati ukāram ca na uccārayati . \\
\noindent \textbf{(P\_3,4.87,} 89) KA\_II,185.3-9 Ro\_III,409 {11/11}     tasya etat prayojanam utvam mā bhūt iti . \\
\noindent \textbf{(P\_3,4.93)} etaḥ aitve ādguṇapratiṣedhaḥ . \\
\noindent \textbf{(P\_3,4.93)} etaḥ aitve ādguṇasya pratiṣedhaḥ vaktavyaḥ . \\
\noindent \textbf{(P\_3,4.93)} pacāva idam (pacāvedam) . \\
\noindent \textbf{(P\_3,4.93)} pacāma idam (pacāmedam) . \\
\noindent \textbf{(P\_3,4.93)} ādguṇe kṛte eta ait iti aitvam prāpnoti . \\
\noindent \textbf{(P\_3,4.93)} na vā bahiraṅgalakṣaṇatvāt . \\
\noindent \textbf{(P\_3,4.93)} na vā vaktavyaḥ . \\
\noindent \textbf{(P\_3,4.93)} kim kāraṇam . \\
\noindent \textbf{(P\_3,4.93)} bahiraṅgalakṣaṇatvāt . \\
\noindent \textbf{(P\_3,4.93)} bahiraṅgalakṣaṇaḥ ādguṇaḥ antaraṅgalakṣaṇam aitvam . \\
\noindent \textbf{(P\_3,4.93)} asiddham bahiraṅgam antaraṅge . \\
\noindent \textbf{(P\_3,4.102)} yāsuḍādeḥ sīyuṭpratiṣedhaḥ . \\
\noindent \textbf{(P\_3,4.102)} yāsuḍādeḥ sīyuṭaḥ pratiṣedhaḥ vaktavyaḥ . \\
\noindent \textbf{(P\_3,4.102)} cinuyuḥ sunuyuḥ . \\
\noindent \textbf{(P\_3,4.102)} liṅaḥ sīyuṭ iti sīyuṭ prāpnoti . \\
\noindent \textbf{(P\_3,4.102)} na vā vākyāpakarṣāt . \\
\noindent \textbf{(P\_3,4.102)} na vā vaktavyaḥ . \\
\noindent \textbf{(P\_3,4.102)} kim kāraṇam . \\
\noindent \textbf{(P\_3,4.102)} vākyāpakarṣāt . \\
\noindent \textbf{(P\_3,4.102)} vākyāpakarṣāt yāsuṭ sīyuṭam bādhiṣyate . \\
\noindent \textbf{(P\_3,4.102)} suṭtithoḥ tu apakarṣavijñānam . \\
\noindent \textbf{(P\_3,4.102)} suṭaḥ tithoḥ tu apakarṣaḥ vijñāyeta . \\
\noindent \textbf{(P\_3,4.102)} kṛṣīṣṭa kṛṣīṣṭhāḥ . \\
\noindent \textbf{(P\_3,4.102)} anādeḥ ca suḍvacanam . \\
\noindent \textbf{(P\_3,4.102)} anādeḥ ca suṭ vaktavyaḥ . \\
\noindent \textbf{(P\_3,4.102)} kṛṣīyāstām kṛṣīyāsthām . \\
\noindent \textbf{(P\_3,4.102)} takārathakārādeḥ liṅaḥ iti suṭ na prāpnoti . \\
\noindent \textbf{(P\_3,4.102)} na vā tithoḥ pradhānabhāvāt tadviśeṣaṇam liṅgrahaṇam . \\
\noindent \textbf{(P\_3,4.102)} na vā vaktavyam . \\
\noindent \textbf{(P\_3,4.102)} kim kāraṇam . \\
\noindent \textbf{(P\_3,4.102)} tithoḥ pradhānabhāvāt . \\
\noindent \textbf{(P\_3,4.102)} tithau eva tatra pradhānam . \\
\noindent \textbf{(P\_3,4.102)} tadviśeṣaṇam liṅgrahaṇam . \\
\noindent \textbf{(P\_3,4.102)} na evam vijñāyate . \\
\noindent \textbf{(P\_3,4.102)} takārathakārayoḥ liṅaḥ iti . \\
\noindent \textbf{(P\_3,4.102)} katham tarhi . \\
\noindent \textbf{(P\_3,4.102)} takārathakārayoḥ suṭ bhavati tau cet liṅaḥ iti . \\
\noindent \textbf{(P\_3,4.103)} kimartham yāsuṭaḥ ṅittvam ucyate . \\
\noindent \textbf{(P\_3,4.103)} yāsuṭaḥ ṅidvacanam pidartham . \\
\noindent \textbf{(P\_3,4.103)} piti vacanāni prayojayanti . \\
\noindent \textbf{(P\_3,4.103)} atha kimartham udāttavacanam kriyate . \\
\noindent \textbf{(P\_3,4.103)} udāttavacanam ca . \\
\noindent \textbf{(P\_3,4.103)} kim . \\
\noindent \textbf{(P\_3,4.103)} pidartham eva . \\
\noindent \textbf{(P\_3,4.103)} āgamānudāttārtham vā . \\
\noindent \textbf{(P\_3,4.103)} atha vā etat jñāpayati ācāryaḥ āgamāḥ anudāttāḥ bhavanti iti . \\
\noindent \textbf{(P\_3,4.103)} asati anyasmin prayojane jñāpakam bhavati . \\
\noindent \textbf{(P\_3,4.103)} uktam ca etat yāsuṭaḥ ṅidvacanam pidartham udāttavacanam ca iti . \\
\noindent \textbf{(P\_3,4.103)} śakyam anena vaktum yāsuṭ parasmaipadeṣu bhavati apit ca liṅ bhavati iti . \\
\noindent \textbf{(P\_3,4.103)} saḥ ayam evam laghīyasā nyāsena siddhe sati yat garīyāṃsam yatnam ārabhate tat jñāpayati ācāryaḥ āgamāḥ anudāttāḥ bhavanti . \\
\noindent \textbf{(P\_3,4.110)} kim idam jusi ākāragrahaṇam niyamārtham āhosvit prāpakam . \\
\noindent \textbf{(P\_3,4.110)} katham ca niyamārtham syāt katham vā prāpakam . \\
\noindent \textbf{(P\_3,4.110)} yadi sijgrahaṇam anuvartate tataḥ niyamārtham . \\
\noindent \textbf{(P\_3,4.110)} atha nivṛttam tataḥ prāpakam . \\
\noindent \textbf{(P\_3,4.110)} kaḥ ca atra viśeṣaḥ . \\
\noindent \textbf{(P\_3,4.110)} jusi ākāragrahaṇam niyamārtham iti cet sijluggrahaṇam . \\
\noindent \textbf{(P\_3,4.110)} jusi ākāragrahaṇam niyamārtham iti cet sijluggrahaṇam kartavyam . \\
\noindent \textbf{(P\_3,4.110)} ātaḥ sijlugantāt iti vaktavyam . \\
\noindent \textbf{(P\_3,4.110)} iha mā bhūt . \\
\noindent \textbf{(P\_3,4.110)} akārṣuḥ ahārṣuḥ . \\
\noindent \textbf{(P\_3,4.110)} astu tarhi prāpakam . \\
\noindent \textbf{(P\_3,4.110)} prāpakam iti cet pratyayalakṣaṇapratiṣedhaḥ . prāpakam iti cet pratyayalakṣaṇapratiṣedhaḥ vaktavyaḥ . \\
\noindent \textbf{(P\_3,4.110)} abhūvan iti pratyayalakṣaṇena jusbhāvaḥ prāpnoti . \\
\noindent \textbf{(P\_3,4.110)} evakārakaraṇam ca . \\
\noindent \textbf{(P\_3,4.110)} evakārakaraṇam ca kartavyam . \\
\noindent \textbf{(P\_3,4.110)} laṅaḥ śākaṭāyanasya eva iti . \\
\noindent \textbf{(P\_3,4.110)} niyamāṛthaḥ punaḥ sati na arthaḥ evakāreṇa . \\
\noindent \textbf{(P\_3,4.110)} nanu ca prāpake api sati siddhi vidhiḥ ārabhyamāṇaḥ antareṇa evakāram niyamārthaḥ bhaviṣyati . \\
\noindent \textbf{(P\_3,4.110)} iṣṭataḥ avadhāraṇārthaḥ tarhi evakāraḥ kartavyaḥ . \\
\noindent \textbf{(P\_3,4.110)} yathā evam vijñāyeta laṅaḥ śākaṭāyanasya eva . \\
\noindent \textbf{(P\_3,4.110)} mā evam vijñāyi laṅaḥ eva śākaṭāyanasya iti . \\
\noindent \textbf{(P\_3,4.110)} kim ca syāt . \\
\noindent \textbf{(P\_3,4.110)} luṅaḥ śākaṭāyanasya na syāt . \\
\noindent \textbf{(P\_3,4.110)} aduḥ apuḥ adhuḥ asthuḥ . \\
\noindent \textbf{(P\_3,4.110)} laṅgrahaṇam ca . \\
\noindent \textbf{(P\_3,4.110)} laṅgrahaṇam ca kartavyam . \\
\noindent \textbf{(P\_3,4.110)} laṅaḥ śākaṭāyanasya eva iti . \\
\noindent \textbf{(P\_3,4.110)} niyamārthe punaḥ sati na arthaḥ laṅgrahaṇena . \\
\noindent \textbf{(P\_3,4.110)} ātaḥ ṅitaḥ iti vartate . \\
\noindent \textbf{(P\_3,4.110)} na ca anyaḥ ākārāt anantaraḥ ṅit asti anyat ataḥ laṅaḥ . \\
\noindent \textbf{(P\_3,4.110)} astu tarhi niyamārthaḥ . \\
\noindent \textbf{(P\_3,4.110)} nanu ca uktam jusi ākāragrahaṇam niyamārtham iti cet sijluggrahaṇam iti . \\
\noindent \textbf{(P\_3,4.110)} na eṣaḥ doṣaḥ . \\
\noindent \textbf{(P\_3,4.110)} tulyajātīyasya niyamaḥ . \\
\noindent \textbf{(P\_3,4.110)} kaḥ ca tulyajātīyaḥ . \\
\noindent \textbf{(P\_3,4.110)} yaḥ dvābhyām anantaraḥ ātaḥ ca sicaḥ ca . \\
\noindent \textbf{(P\_3,4.110)} atha tat evakārakaraṇam na eva kartavyam . \\
\noindent \textbf{(P\_3,4.110)} kartavyam ca . \\
\noindent \textbf{(P\_3,4.110)} kim prayojanam . \\
\noindent \textbf{(P\_3,4.110)} uttarārtham . \\
\noindent \textbf{(P\_3,4.110)} liṭ ca liṅ āśiṣi ārdhadhātukam eva yathā syāt . \\
\noindent \textbf{(P\_3,4.110)} itarathā hi vacanāt ārdhadhātukasañjñā syāt tiṅgrahaṇena ca grahaṇāt sārvadhātukasañjñā . \\
\noindent \textbf{(P\_3,4.110)} atha tat laṅgrahaṇam na eva kartavyam . \\
\noindent \textbf{(P\_3,4.110)} kartavyam ca . \\
\noindent \textbf{(P\_3,4.110)} kim prayojanam . \\
\noindent \textbf{(P\_3,4.110)} laṅ eva yaḥ laṅ tatra yathā syāt . \\
\noindent \textbf{(P\_3,4.110)} laṅvadbhāvena yaḥ laṅ tatra mā bhūt iti . \\
\noindent \textbf{(P\_3,4.114)} ārdhadhātukasañjñāyām dhātugrahaṇam . \\
\noindent \textbf{(P\_3,4.114)} ārdhadhātukasañjñāyām dhātugrahaṇam kartavyam . \\
\noindent \textbf{(P\_3,4.114)} dhātoḥ parasya ārdhadhātukasañjñā yathā syāt . \\
\noindent \textbf{(P\_3,4.114)} iha mā bhūt . \\
\noindent \textbf{(P\_3,4.114)} vṛkṣatvam vṛkṣatā iti . \\
\noindent \textbf{(P\_3,4.114)} kriyamāṇe ca api dhātugrahaṇe svādipratiṣedhaḥ . \\
\noindent \textbf{(P\_3,4.114)} svādīnām pratiṣedhaḥ vaktavyaḥ . \\
\noindent \textbf{(P\_3,4.114)} iha mā bhūt . \\
\noindent \textbf{(P\_3,4.114)} lūbhyām lūbhiḥ iti . \\
\noindent \textbf{(P\_3,4.114)} anukrāntāpekṣam śeṣagrahaṇam .evam api agnikāmpyati vāyukāmyati iti prāpnoti . \\
\noindent \textbf{(P\_3,4.114)} tasmāt dhātugrahaṇam kartavyam . \\
\noindent \textbf{(P\_3,4.114)} na kartavyam . \\
\noindent \textbf{(P\_3,4.114)} ā tṛtīyādhyāyaparisamāpteḥ dhātvadhikāraḥ prakṛtaḥ anuvartate . \\
\noindent \textbf{(P\_3,4.114)} kva prakṛtaḥ . \\
\noindent \textbf{(P\_3,4.114)} dhātoḥ ekācaḥ halādeḥ iti . \\
\noindent \textbf{(P\_3,4.114)} evam api śrīkāmyati bhūkāmyati iti prāpnoti . \\
\noindent \textbf{(P\_3,4.114)} tadvidhānāt siddham . \\
\noindent \textbf{(P\_3,4.114)} vihitaviśeṣaṇam dhātugrahaṇam . \\
\noindent \textbf{(P\_3,4.114)} dhātoḥ yaḥ vihitaḥ iti . \\
\noindent \textbf{(P\_3,4.114)} dhātoḥ eṣaḥ vihitaḥ . \\
\noindent \textbf{(P\_3,4.114)} saṅkīrtya dhātoḥ iti evam yaḥ vihitaḥ iti. \\
\noindent \textbf{(P\_5,1.1)} prāgvacanam kimartham . \\
\noindent \textbf{(P\_5,1.1)} prāgvacane uktam . \\
\noindent \textbf{(P\_5,1.1)} kim uktam . \\
\noindent \textbf{(P\_5,1.1)} tatra tāvat uktam prāgvacanam sakṛdvidhānārtham . \\
\noindent \textbf{(P\_5,1.1)} adhikārāt siddham iti cet apavādaviṣaye aṇprasaṅgaḥ iti . \\
\noindent \textbf{(P\_5,1.1)} iha api prāgvacanam kriyate sakṛdvidhānārtham . \\
\noindent \textbf{(P\_5,1.1)} sakṛt vihitaḥ pratyayaḥ vihitaḥ yathā syāt . \\
\noindent \textbf{(P\_5,1.1)} yoge yoge tasya grahaṇam mā kārṣam iti . \\
\noindent \textbf{(P\_5,1.1)} na etat asti prayojanam . \\
\noindent \textbf{(P\_5,1.1)} adhikārāt api etat siddham . \\
\noindent \textbf{(P\_5,1.1)} adhikāraḥ pratiyogam tasya anirdeśārthaḥ iti yoge yoge upatiṣṭhate . \\
\noindent \textbf{(P\_5,1.1)} adhikārāt siddham iti cet apavādaviṣaye chaprasaṅgaḥ . \\
\noindent \textbf{(P\_5,1.1)} adhikārāt siddham iti cet apavādaviṣaye chaḥ prāpnoti . \\
\noindent \textbf{(P\_5,1.1)} ugavādibhyaḥ yat chaḥ ca iti chaḥ api prāpnoti . \\
\noindent \textbf{(P\_5,1.1)} tasmāt prāgvacanam kartavyam . \\
\noindent \textbf{(P\_5,1.1)} atha kriymāṇe api prāgvacane katham idam vijñāyate . \\
\noindent \textbf{(P\_5,1.1)} prāk krītāt yāḥ prakṛtayaḥ āhosvit prāk krītāt ye arthāḥ iti . \\
\noindent \textbf{(P\_5,1.1)} kim ca ataḥ . \\
\noindent \textbf{(P\_5,1.1)} yadi vijñāyate prāk krītāt yāḥ prakṛtayaḥ iti saḥ eva doṣaḥ apavādaviṣaye api chaprasaṅgaḥ iti . \\
\noindent \textbf{(P\_5,1.1)} atha vijñāyate prāk krītāt ye arthāḥ iti na doṣaḥ bhavati . \\
\noindent \textbf{(P\_5,1.1)} samāne arthe prakṛtiviśeṣāt utpadyamānaḥ yat cham bādhiṣyate . \\
\noindent \textbf{(P\_5,1.1)} yathā na doṣaḥ tathā astu. prāk krītāt ye arthāḥ iti vijñāyate . \\
\noindent \textbf{(P\_5,1.1)} kutaḥ etat . \\
\noindent \textbf{(P\_5,1.1)} tathā hi ayam prādhānyena artham pratinirdiśati . \\
\noindent \textbf{(P\_5,1.1)} itarathā hi bahvyaḥ tatra prakṛtayaḥ paṭhyante . \\
\noindent \textbf{(P\_5,1.1)} tataḥ yām kām cit evam prakṛtim avadhitvena upādadīta . \\
\noindent \textbf{(P\_5,1.1)} atha vā punaḥ astu prāk krītāt yāḥ prakṛtayaḥ iti . \\
\noindent \textbf{(P\_5,1.1)} nanu ca uktam apavādaviṣaye api chaprasaṅgaḥ iti . \\
\noindent \textbf{(P\_5,1.1)} na vā kva cit vāvacanāt . \\
\noindent \textbf{(P\_5,1.1)} na vā eṣaḥ doṣaḥ . \\
\noindent \textbf{(P\_5,1.1)} kim kāraṇam . \\
\noindent \textbf{(P\_5,1.1)} kva cit vāvacanāt . \\
\noindent \textbf{(P\_5,1.1)} yat ayam kvac vāvacanam karoti vibhāṣā havirapūpādibhyaḥ iti tat jñāpayati na apavādaviṣaye chaḥ bhavati iti . \\
\noindent \textbf{(P\_5,1.1)} yadi evam na arthaḥ prāgvacanena . \\
\noindent \textbf{(P\_5,1.1)} adhikārāt siddham . \\
\noindent \textbf{(P\_5,1.1)} nanu ca uktam adhikārāt siddham iti cet apavādaviṣaye chaprasaṅgaḥ iti . \\
\noindent \textbf{(P\_5,1.1)} parihṛtam etat ṇa vā kva cit vāvacanāt iti . \\
\noindent \textbf{(P\_5,1.1)} atha kimartham iyān avadhiḥ gṛhyate na prāk ṭhañaḥ iti eva ucyeta . \\
\noindent \textbf{(P\_5,1.1)} etat jñāpayati ācāryaḥ artheṣu ayam bhavati iti . \\
\noindent \textbf{(P\_5,1.1)} kim etasya jñāpane prayojanam . \\
\noindent \textbf{(P\_5,1.1)} samāne arthe prakṛtiviśeṣāt utpadyamānaḥ yat cham bādhate . \\
\noindent \textbf{(P\_5,1.2.1)} yaññyau añaḥ pūrvavipratiṣiddham sanaṅgūpānahau prayojanam . yaññyau bhavataḥ añaḥ pūrvavipratiṣedhena . \\
\noindent \textbf{(P\_5,1.2.1)} kim prayojanam ṣanaṅgūpānahau prayojanam . \\
\noindent \textbf{(P\_5,1.2.1)} yataḥ avakāśaḥ śaṅkavyam dāru picavyaḥ kārpāsaḥ . \\
\noindent \textbf{(P\_5,1.2.1)} añaḥ avakāśaḥ vārdhram vāratram . \\
\noindent \textbf{(P\_5,1.2.1)} sanaṅguḥ nāma carmavikāraḥ . \\
\noindent \textbf{(P\_5,1.2.1)} tasmāt ubhayam prāpnoti . \\
\noindent \textbf{(P\_5,1.2.1)} sanaṅgavyam carma . \\
\noindent \textbf{(P\_5,1.2.1)} ñyasya avakāśaḥ aupānahyam dāru . \\
\noindent \textbf{(P\_5,1.2.1)} añaḥ saḥ eva . \\
\noindent \textbf{(P\_5,1.2.1)} upānat nāma carmavikāraḥ . \\
\noindent \textbf{(P\_5,1.2.1)} tasmāt ubhayam prāpnoti . \\
\noindent \textbf{(P\_5,1.2.1)} aupānahyam carma . \\
\noindent \textbf{(P\_5,1.2.1)} ḍhañ ca . ḍhañ ca bhavati añaḥ pūrvavipratiṣedhena . \\
\noindent \textbf{(P\_5,1.2.1)} ḍhañaḥ avakāśaḥ chādiṣeyam tṛṇam . \\
\noindent \textbf{(P\_5,1.2.1)} añaḥ saḥ eva . \\
\noindent \textbf{(P\_5,1.2.1)} chadiḥ nām carmavikāraḥ . \\
\noindent \textbf{(P\_5,1.2.1)} tasmāt ubhayam prāpnoti . \\
\noindent \textbf{(P\_5,1.2.1)} chādiṣeyam carma . \\
\noindent \textbf{(P\_5,1.2.1)} ḍhañ bhavati pūrvavipratiṣedhena . \\
\noindent \textbf{(P\_5,1.2.1)} havirapūpādibhyaḥ vibhāṣāyāḥ yat . havirapūpādibhyaḥ vibhāṣāyāḥ yat bhavati pūrvavipratiṣedhena . \\
\noindent \textbf{(P\_5,1.2.1)} havirapūpādibhyaḥ vibhāṣāyāḥ avakāśaḥ āmikṣyam āmikṣīyam puroḍāśyam puroḍāśīyam . \\
\noindent \textbf{(P\_5,1.2.1)} yataḥ saḥ eva . \\
\noindent \textbf{(P\_5,1.2.1)} iha ubhayam prāpnoti . \\
\noindent \textbf{(P\_5,1.2.1)} caravyāḥ taṇḍulāḥ . \\
\noindent \textbf{(P\_5,1.2.1)} yat bhavati pūrvavipratiṣedhena . \\
\noindent \textbf{(P\_5,1.2.1)} annavikārebhyaḥ ca . annavikārebhyaḥ ca vibhāṣāyāḥ yat bhavati pūrvavipratiṣedhena . \\
\noindent \textbf{(P\_5,1.2.1)} annavikārebhyaḥ ca vibhāṣāyāḥ avakāśaḥ suryāḥ surīyāḥ . \\
\noindent \textbf{(P\_5,1.2.1)} yataḥ saḥ eva . \\
\noindent \textbf{(P\_5,1.2.1)} iha ubhayam prāpnoti . \\
\noindent \textbf{(P\_5,1.2.1)} saktavyāḥ dhānāḥ iti . \\
\noindent \textbf{(P\_5,1.2.1)} yat bhavati pūrvavipratiṣedhena . \\
\noindent \textbf{(P\_5,1.2.1)} saḥ tarhi pūrvavipratiṣedhaḥ vaktavyaḥ . \\
\noindent \textbf{(P\_5,1.2.1)} na vaktavyaḥ . \\
\noindent \textbf{(P\_5,1.2.1)} iṣṭavācī paraśabdaḥ . \\
\noindent \textbf{(P\_5,1.2.1)} vipratiṣedhe param yat iṣṭam tat bhavati iti . \\
\noindent \textbf{(P\_5,1.2.2)} ayam nābhiśabdaḥ gavādiṣu paṭhyate. tatra eva ucyate nābhi nabham ca iti . \\
\noindent \textbf{(P\_5,1.2.2)} tatra codyate . \\
\noindent \textbf{(P\_5,1.2.2)} nābheḥ nabhabhāve pratyayānupapattiḥ prakṛtyabhāvāt . nābheḥ nabhabhāve pratyayānupapattiḥ . \\
\noindent \textbf{(P\_5,1.2.2)} kim kāraṇam . \\
\noindent \textbf{(P\_5,1.2.2)} prakṛtyabhāvāt . \\
\noindent \textbf{(P\_5,1.2.2)} vikṛteḥ prakṛtau abhidheyāyām pratyayena bhavitavyam na ca nābhisañjñikāyāḥ vikṛteḥ prakṛtiḥ asti . \\
\noindent \textbf{(P\_5,1.2.2)} yat eva hi tanmaṇḍalacakrāṇām maṇḍalacakram tat nabhyam iti ucyate . \\
\noindent \textbf{(P\_5,1.2.2)} siddham tu śākhādiṣu vacanāt hrasvatvam ca . \\
\noindent \textbf{(P\_5,1.2.2)} siddham etat . \\
\noindent \textbf{(P\_5,1.2.2)} katham . \\
\noindent \textbf{(P\_5,1.2.2)} śākhādiṣu nābhiśabdaḥ paṭhitavyaḥ hrasvatvam ca vaktavyam . \\
\noindent \textbf{(P\_5,1.2.2)} nābhiḥ iva nabhyam iti . \\
\noindent \textbf{(P\_5,1.2.2)} kaḥ punaḥ iha upamārthaḥ . \\
\noindent \textbf{(P\_5,1.2.2)} yat tat akṣadhāraṇam parivartanam vā . \\
\noindent \textbf{(P\_5,1.2.2)} aparaḥ āha : yat tat añjanopāñjanam iti . \\
\noindent \textbf{(P\_5,1.2.2)} na tarhi idānīm idam vaktavyam nābhi nabham ca iti . \\
\noindent \textbf{(P\_5,1.2.2)} vaktavyam ca . \\
\noindent \textbf{(P\_5,1.2.2)} kim prayojanam . \\
\noindent \textbf{(P\_5,1.2.2)} yāni etāni aravanti cakrāṇi tadartham . \\
\noindent \textbf{(P\_5,1.2.2)} tatra nābhisañjñikāyāḥ vikṛteḥ prakṛtiḥ asti . \\
\noindent \textbf{(P\_5,1.2.2)} yāni ca api anaravanti cakrāṇi tadartham api idam vaktavyam . \\
\noindent \textbf{(P\_5,1.2.2)} dṛśyate hi samudāyāt avayavasya pṛthaktvam . \\
\noindent \textbf{(P\_5,1.2.2)} tat yathā vārkṣī śākhā iti . \\
\noindent \textbf{(P\_5,1.2.2)} guṇāntarayogāt ca vikāraśabdaḥ dṛśyate . \\
\noindent \textbf{(P\_5,1.2.2)} tat yathā vaibhītakaḥ yūpaḥ khādiram caṣālam iti . \\
\noindent \textbf{(P\_5,1.2.2)} tatra avayavasamudāye vṛttiḥ bhaviṣyati . \\
\noindent \textbf{(P\_5,1.2.2)} atha yaḥ nabhyārthaḥ vṛkṣaḥ katham tatra bhavitavyam . \\
\noindent \textbf{(P\_5,1.2.2)} nabhyaḥ vṛkṣaḥ nabhyā śiṃśipā iti . \\
\noindent \textbf{(P\_5,1.2.2)} nabhyāt tu lugvacanam . nabhyāt tu luk vaktavyaḥ . \\
\noindent \textbf{(P\_5,1.2.2)} saḥ tarhi vaktavyaḥ . \\
\noindent \textbf{(P\_5,1.2.2)} na vaktavyaḥ . \\
\noindent \textbf{(P\_5,1.2.2)} tādarthyāt tācchabdyam bhaviṣyati . \\
\noindent \textbf{(P\_5,1.2.2)} nabhyārthaḥ nabhyaḥ iti . \\
\noindent \textbf{(P\_5,1.3)} ayam yogaḥ śakyaḥ avaktum . \\
\noindent \textbf{(P\_5,1.3)} katham aśītiśatam kambalyam iti . \\
\noindent \textbf{(P\_5,1.3)} nipātanāt etat siddham . \\
\noindent \textbf{(P\_5,1.3)} kim nipātanam . \\
\noindent \textbf{(P\_5,1.3)} aparimāṇavistācitakambalebhyaḥ na taddhitaluki iti . \\
\noindent \textbf{(P\_5,1.3)} idam tarhi prayojanam sañjñāyām iti vakṣyāmi iti . \\
\noindent \textbf{(P\_5,1.3)} iha mā bhūt . \\
\noindent \textbf{(P\_5,1.3)} kambalīyāḥ ūrṇāḥ . \\
\noindent \textbf{(P\_5,1.3)} etat api na asti prayojanam . \\
\noindent \textbf{(P\_5,1.3)} parimāṇaparyudāsena paryudāse prāpte tatra kambalagrahaṇam kriyate parimāṇārtham parimāṇam ca sañjñā eva . \\
\noindent \textbf{(P\_5,1.4)} kim iyam prāpte vibhāṣā āhosvit aprāpte . \\
\noindent \textbf{(P\_5,1.4)} katham ca prāpte katham vā aprāpte . \\
\noindent \textbf{(P\_5,1.4)} uvarṇāntāt iti vā nitye prāpte anyatra vā aprāpte . \\
\noindent \textbf{(P\_5,1.4)} havirapūpādibhyaḥ aprāpte . \\
\noindent \textbf{(P\_5,1.4)} havirapūpādibhyaḥ aprāpte vibhāṣā . \\
\noindent \textbf{(P\_5,1.4)} prāpte nityaḥ vidhiḥ . \\
\noindent \textbf{(P\_5,1.4)} caravyāḥ taṇḍulāḥ . \\
\noindent \textbf{(P\_5,1.6)} yatprakaraṇe rathāt ca . yatprakaraṇe rathāt ca upasaṅkhyānam . \\
\noindent \textbf{(P\_5,1.6)} rathāya hitā rathyā . \\
\noindent \textbf{(P\_5,1.7)} vṛṣaśabdaḥ ayam akārāntaḥ gṛhyate . \\
\noindent \textbf{(P\_5,1.7)} vṛṣanśabdaḥ api nakārāntaḥ asti . \\
\noindent \textbf{(P\_5,1.7)} tasya upasaṅkhyānam kartavyam vṛṣaśabdaḥ ca ādeśaḥ vaktavyaḥ vṛṣṇe hitam iti vigṛhya vṛṣyam iti eva yathā syāt . \\
\noindent \textbf{(P\_5,1.7)} tathā brahmanśabdaḥ nakārāntaḥ gṛhyate . \\
\noindent \textbf{(P\_5,1.7)} brāhmaṇaśabdaḥ ca akārāntaḥ asti . \\
\noindent \textbf{(P\_5,1.7)} tasya upasaṅkhyānam kartavyam brahmanśabdaḥ ca ādeśaḥ vaktavyaḥ brāhmaṇebhyaḥ hitam iti vigṛhya brahmaṇyam iti eva yathā syāt . \\
\noindent \textbf{(P\_5,1.7)} tat tarhi vaktavyam . \\
\noindent \textbf{(P\_5,1.7)} na vaktavyam . \\
\noindent \textbf{(P\_5,1.7)} samānārthau etau vṛṣaśabdaḥ vṛṣanśabdaḥ ca brahmanśabdaḥ brāhmaṇaśabdaḥ ca . \\
\noindent \textbf{(P\_5,1.7)} ātaḥ ca samānārthau . \\
\noindent \textbf{(P\_5,1.7)} evam hi āha . \\
\noindent \textbf{(P\_5,1.7)} kutaḥ nu carasi brahman . \\
\noindent \textbf{(P\_5,1.7)} kutaḥ nu carasi brāhmaṇa iti . \\
\noindent \textbf{(P\_5,1.7)} tatra dvayoḥ śabdayoḥ samānārthayoḥ ekena vigrahaḥ aparasmāt utpattiḥ bhaviṣyati aviravikanyāyena . \\
\noindent \textbf{(P\_5,1.7)} tat yathā . \\
\noindent \textbf{(P\_5,1.7)} aveḥ māṃsam iti vigṛhya avikaśabdāt utpattiḥ bhavati āvikam iti . \\
\noindent \textbf{(P\_5,1.7)} evam iha api vṛṣāya hitam iti vigṛhya vṛṣyam iti bhaviṣyati . \\
\noindent \textbf{(P\_5,1.7)} vṛṣṇe hitam iti vigṛhya vākyam eva . \\
\noindent \textbf{(P\_5,1.7)} tathā brahmaṇe hitam iti vigṛhya brahmaṇyam iti bhaviṣyati . \\
\noindent \textbf{(P\_5,1.7)} brāhmaṇebhyaḥ hitam iti vigṛhya vākyam eva bhaviṣyati . \\
\noindent \textbf{(P\_5,1.7)} traiśabdyam ca iha sādhyam . \\
\noindent \textbf{(P\_5,1.7)} tat ca evam sati siddham bhavati . \\
\noindent \textbf{(P\_5,1.9.1)} bhogottarapadāt khavidhāne anirdeśaḥ pūrvapadārthahitatvāt . bhogottarapadāt khavidhāne anirdeśaḥ . \\
\noindent \textbf{(P\_5,1.9.1)} agamakaḥ nirdeśaḥ anirdeśaḥ . \\
\noindent \textbf{(P\_5,1.9.1)} kim kāraṇam . \\
\noindent \textbf{(P\_5,1.9.1)} pūrvapadārthahitatvāt . \\
\noindent \textbf{(P\_5,1.9.1)} uttarapadārthapradhānaḥ tatpuruṣaḥ pūrvapadārthapradhāne ca pratyayaḥ iṣyate . \\
\noindent \textbf{(P\_5,1.9.1)} pitṛbhogāya hite prāpnoti pitre ca eva hite iṣyate . \\
\noindent \textbf{(P\_5,1.9.1)} evam tarhi bhogīnarpratyayaḥ vijñāsyate . \\
\noindent \textbf{(P\_5,1.9.1)} bhogīnar iti cet vāvacanam . \\
\noindent \textbf{(P\_5,1.9.1)} bhogīnar iti yadi pratyayaḥ vidhīyate vāvacanam kartavyam mātrīyaḥ pitrīyaḥ iti api yathā syāt . \\
\noindent \textbf{(P\_5,1.9.1)} rājācāryābhyām nityam . rājācāryābhyām nityam iti vaktavyam . \\
\noindent \textbf{(P\_5,1.9.1)} rājabhogīnaḥ . \\
\noindent \textbf{(P\_5,1.9.1)} ācāryāt aṇatvam ca . \\
\noindent \textbf{(P\_5,1.9.1)} ācāryabhogīnaḥ . \\
\noindent \textbf{(P\_5,1.9.1)} kim bhogīnarpratyayaḥ vidhīyate iti ataḥ rājācāryābhyām nityam iti vaktavyam . \\
\noindent \textbf{(P\_5,1.9.1)} na iti āha . \\
\noindent \textbf{(P\_5,1.9.1)} sarvathā rājācāryābhyām nityam iti vaktavyam . \\
\noindent \textbf{(P\_5,1.9.1)} iha ca grāmaṇibhogīnaḥ senānibhogīnaḥ iti uttarapade iti hrasvatvam na prāpnoti . \\
\noindent \textbf{(P\_5,1.9.1)} iha ca abbhoginaḥ iti apaḥ bhi iti tatvam prāpnoti . \\
\noindent \textbf{(P\_5,1.9.1)} sūtram ca bhidyate . \\
\noindent \textbf{(P\_5,1.9.1)} yathānyāsam eva astu . \\
\noindent \textbf{(P\_5,1.9.1)} nanu ca uktam bhogottarapadāt khavidhāne anirdeśaḥ pūrvapadārthahitatvāt iti . \\
\noindent \textbf{(P\_5,1.9.1)} na eṣaḥ doṣaḥ . \\
\noindent \textbf{(P\_5,1.9.1)} ayam bhogaśabdaḥ asti eva dravyapadārthakaḥ . \\
\noindent \textbf{(P\_5,1.9.1)} tat yathā bhogavān ayam deśaḥ iti ucyate yasmin gāvaḥ sasnāni ca vartante . \\
\noindent \textbf{(P\_5,1.9.1)} asti kriyāpadārthakaḥ . \\
\noindent \textbf{(P\_5,1.9.1)} tat yathā bhogavān ayam brāhmaṇaḥ iti ucyate yaḥ samyak snānādīḥ kriyāḥ anubhavati . \\
\noindent \textbf{(P\_5,1.9.1)} tat yaḥ kriyāpadārthakaḥ tasya ayam grahaṇam . \\
\noindent \textbf{(P\_5,1.9.1)} yaḥ ca pitṛsthābhyaḥ kriyābhyaḥ hitaḥ sambandhāt asau pitre api hitaḥ bhavati . \\
\noindent \textbf{(P\_5,1.9.1)} yadi sambandhāt astu dravyapadārthakasya api grahaṇam . \\
\noindent \textbf{(P\_5,1.9.1)} yaḥ api hi pitṛdravyāya hitaḥ sambandhāt asau pitre hitaḥ bhavati . \\
\noindent \textbf{(P\_5,1.9.1)} atha vā bhogaśabdaḥ śarīravācī api drśyate . \\
\noindent \textbf{(P\_5,1.9.1)} tat yathā ahiḥ iva bhogaiḥ paryeti bāhum iti. ahiḥ iva śarīraiḥ iti gamyate . \\
\noindent \textbf{(P\_5,1.9.1)} evam pitṛśarīrāya hitaḥ pitṛbhogīṇaḥ iti . \\
\noindent \textbf{(P\_5,1.9.2)} khavidhāne pañcajanāt upasaṅkhyānam . khavidhāne pañcajanāt upasaṅkhyānam kartavyam . \\
\noindent \textbf{(P\_5,1.9.2)} pañcajanāya hitaḥ pañcajanīnaḥ . \\
\noindent \textbf{(P\_5,1.9.2)} samānādhikaraṇe iti vaktavyam . \\
\noindent \textbf{(P\_5,1.9.2)} yaḥ hi pañcānām janāya hitaḥ pañcajanīyaḥ saḥ bhavati . \\
\noindent \textbf{(P\_5,1.9.2)} sarvajanāt ṭhañ ca . sarvajanāt ṭhañ vaktavyaḥ khaḥ ca . \\
\noindent \textbf{(P\_5,1.9.2)} sarvajanāya hitaḥ sārvajanikaḥ sārvajanīnaḥ . \\
\noindent \textbf{(P\_5,1.9.2)} samānādhikaraṇe iti ca vaktavyam . \\
\noindent \textbf{(P\_5,1.9.2)} yaḥ hi sarveṣām janāya hitaḥ sarvajanīyaḥ saḥ . \\
\noindent \textbf{(P\_5,1.9.2)} mahājanāt nityam . mahājanāt nityam ṭhañ vaktavyaḥ . \\
\noindent \textbf{(P\_5,1.9.2)} mahājanāya hitaḥ māhājanikaḥ . \\
\noindent \textbf{(P\_5,1.9.2)} tatpuruṣe iti vaktavyam bahuvrīhau mā bhūt iti . \\
\noindent \textbf{(P\_5,1.9.2)} mahān janaḥ asya mahājanaḥ mahājanāya hitaḥ mahājanīyaḥ . \\
\noindent \textbf{(P\_5,1.9.2)} yadi tarhi atiprasaṅgāḥ santi iti upādhiḥ kriyate ādyanyāse api upādhiḥ kartavyaḥ . \\
\noindent \textbf{(P\_5,1.9.2)} ātmanviśvajane samānādhikaraṇe iti vaktavyam . \\
\noindent \textbf{(P\_5,1.9.2)} yaḥ his viśveṣām janāya hitaḥ viśvajanīyaḥ saḥ bhavati . \\
\noindent \textbf{(P\_5,1.9.2)} atha matam etat anabhidhānāt ādyanyāse na bhaviṣyati iti iha api na arthaḥ upādhigrahaṇena . \\
\noindent \textbf{(P\_5,1.9.2)} iha api anabhidhānāt na bhaviṣyati . \\
\noindent \textbf{(P\_5,1.10)} sarvāt ṇasya vāvacanam . sarvāt ṇasya vā iti vaktavyam . \\
\noindent \textbf{(P\_5,1.10)} sārvaḥ sarvīyaḥ . \\
\noindent \textbf{(P\_5,1.10)} puruṣāt vadhe . puruṣāt vadhe iti vaktavyam : pauruṣeyaḥ vadhaḥ . \\
\noindent \textbf{(P\_5,1.10)} atyalpam idam ucyate : puruṣāt vadhe iti . \\
\noindent \textbf{(P\_5,1.10)} puruṣāt vadhavikārasamūhatenakṛteṣu iti vaktavyam : pauruṣeyaḥ vadhaḥ , pauruṣeyaḥ vikāraḥ , pauruṣeyaḥ samūhaḥ , tena kṛtam pauruṣeyam . \\
\noindent \textbf{(P\_5,1.12)} tadartham iti kṛtyanāmabhyaḥ ṭhañ . \\
\noindent \textbf{(P\_5,1.12)} tadartham iti kṛtyanāmabhyaḥ ṭhañ vaktavyaḥ . \\
\noindent \textbf{(P\_5,1.12)} indramahārtham aindramahiham gāṅgāmahiham kāśeruyajñikam . \\
\noindent \textbf{(P\_5,1.12)} na vā prayojanena kṛtatvāt . \\
\noindent \textbf{(P\_5,1.12)} na vā vaktavyam . \\
\noindent \textbf{(P\_5,1.12)} kim kāraṇam . \\
\noindent \textbf{(P\_5,1.12)} prayojanena kṛtatvāt . \\
\noindent \textbf{(P\_5,1.12)} yat hi indramahārtham indramahaḥ tasya prayojanam bhavati . \\
\noindent \textbf{(P\_5,1.12)} tatra prayojanam iti eva siddham . \\
\noindent \textbf{(P\_5,1.13)} upadhyartham iti pratyayānupapattiḥ . upadhyartham iti pratyayasya iha anupapattiḥ . \\
\noindent \textbf{(P\_5,1.13)} kim kāraṇam . \\
\noindent \textbf{(P\_5,1.13)} upadhyabhāvāt . \\
\noindent \textbf{(P\_5,1.13)} vikṛteḥ prakṛtau abhidheyāyām pratyayena bhavitavyam . \\
\noindent \textbf{(P\_5,1.13)} na ca upadhisañjñikāyāḥ vikṛteḥ prakṛtiḥ asti . \\
\noindent \textbf{(P\_5,1.13)} yat hi tat rathāṅgam tat aupadheyam iti ucyate . \\
\noindent \textbf{(P\_5,1.13)} siddham tu kṛdantasya svārthe añvacanāt . siddham etat . \\
\noindent \textbf{(P\_5,1.13)} katham . \\
\noindent \textbf{(P\_5,1.13)} kṛdantasya svārthe añ vaktavyaḥ . \\
\noindent \textbf{(P\_5,1.13)} upadhīyate upadheyam . \\
\noindent \textbf{(P\_5,1.13)} upadheyam eva aupadheyam . \\
\noindent \textbf{(P\_5,1.13)} sidhyati . \\
\noindent \textbf{(P\_5,1.13)} sūtram tarhi bhidyate . \\
\noindent \textbf{(P\_5,1.13)} yathānyāsam eva astu . \\
\noindent \textbf{(P\_5,1.13)} nanu ca uktam upadhyartham iti pratyayānupapattiḥ iti . \\
\noindent \textbf{(P\_5,1.13)} na etat asti . \\
\noindent \textbf{(P\_5,1.13)} ayam upadhiśabdaḥ asti eva karmasādhanaḥ . \\
\noindent \textbf{(P\_5,1.13)} upadhīyate upadhiḥ iti . \\
\noindent \textbf{(P\_5,1.13)} asti bhāvasādhanaḥ . \\
\noindent \textbf{(P\_5,1.13)} upadhānam upadhiḥ iti . \\
\noindent \textbf{(P\_5,1.13)} tat yaḥ bhāvasādhanaḥ tasya idam grahaṇam . \\
\noindent \textbf{(P\_5,1.13)} evam api na sidhyati . \\
\noindent \textbf{(P\_5,1.13)} kim kāraṇam . \\
\noindent \textbf{(P\_5,1.13)} vikṛteḥ prakṛtau iti vartate . \\
\noindent \textbf{(P\_5,1.13)} prakṛtivikṛtigrahaṇam nivartiṣyate . \\
\noindent \textbf{(P\_5,1.13)} tat ca avaśyam nivartyam ihārtham uttarātham ca . \\
\noindent \textbf{(P\_5,1.13)} ihārtham tāvat . \\
\noindent \textbf{(P\_5,1.13)} bāleyāḥ taṇḍulāḥ . \\
\noindent \textbf{(P\_5,1.13)} uttarārtham ṛṣabhopānahoḥ ñyaḥ . \\
\noindent \textbf{(P\_5,1.13)} ārṣabhyaḥ vatsaḥ iti . \\
\noindent \textbf{(P\_5,1.13)} atha tadartham iti anuvartate utāho na . \\
\noindent \textbf{(P\_5,1.13)} kim ca arthaḥ anuvṛttyā . \\
\noindent \textbf{(P\_5,1.13)} bāḍham arthaḥ . \\
\noindent \textbf{(P\_5,1.13)} tat asya tat asmin syāt iti tadarthe yathā syāt . \\
\noindent \textbf{(P\_5,1.13)} iha mā bhūt . \\
\noindent \textbf{(P\_5,1.13)} prāsādaḥ devadattasya syāt iti . \\
\noindent \textbf{(P\_5,1.13)} prākāraḥ nagarasya syāt iti . \\
\noindent \textbf{(P\_5,1.13)} yadi tadartham iti anuvartate ṛṣabhopānahoḥ ñyaḥ ṛṣabhārthaḥ ghāsaḥ upānadarthaḥ tilakalkaḥ iti atra api prāpnoti . \\
\noindent \textbf{(P\_5,1.13)} evam tarhi anuvartate prakṛtivikṛtigrahaṇam . \\
\noindent \textbf{(P\_5,1.13)} nanu ca uktam balyṛṣabhayoḥ na sidhyati iti . \\
\noindent \textbf{(P\_5,1.13)} kim punaḥ bhavān vikāram matvā āha balyṛṣabhayoḥ na sidhyati iti . \\
\noindent \textbf{(P\_5,1.13)} yadi tāvat yaḥ prakṛtyupamardena bhavati saḥ vikāraḥ vaibhītakaḥ yūpaḥ khādiram caṣālam iti na sidhyati . \\
\noindent \textbf{(P\_5,1.13)} atha matam etat eva guṇāntarayuktam vikāraḥ iti balyṛṣabhayoḥ api siddham bhavati . \\
\noindent \textbf{(P\_5,1.13)} guṇantarayuktāḥ hi taṇḍulāḥ bāleyāḥ guṇāntarayuktaḥ ca vatsaḥ ārṣabhaḥ . \\
\noindent \textbf{(P\_5,1.13)} aupadheyam tu na sidhyati . \\
\noindent \textbf{(P\_5,1.13)} vacanāt svārthikaḥ bhaviṣyati . \\
\noindent \textbf{(P\_5,1.16)} syādgrahaṇam kimartham . \\
\noindent \textbf{(P\_5,1.16)} iha mā bhūt . \\
\noindent \textbf{(P\_5,1.16)} prāsādaḥ devadattasya prākāraḥ nagarasya iti . \\
\noindent \textbf{(P\_5,1.16)} atha kriyamāṇe api syādgrahaṇe iha kasmāt na bhavati . \\
\noindent \textbf{(P\_5,1.16)} prāsādaḥ devadattasya syāt . \\
\noindent \textbf{(P\_5,1.16)} prākāraḥ nagarasya syāt iti . \\
\noindent \textbf{(P\_5,1.16)} śakyārthe liṅ iti vaktavyam . \\
\noindent \textbf{(P\_5,1.16)} na evam śakyam . \\
\noindent \textbf{(P\_5,1.16)} idānīm eva hi uktam na hi upādheḥ upādhiḥ bhavati viśeṣaṇasya vā viśeṣaṇam iti . \\
\noindent \textbf{(P\_5,1.16)} evam tarhi itikaraṇaḥ kriyate . \\
\noindent \textbf{(P\_5,1.16)} tataḥ cet vivakṣā bhavati . \\
\noindent \textbf{(P\_5,1.16)} vivakṣā ca dvayī . \\
\noindent \textbf{(P\_5,1.16)} asti eva prāyoktrī vivakṣā asti laukikī . \\
\noindent \textbf{(P\_5,1.16)} prayoktā hi mṛdvyā snigdhayā ślakṣṇayā jihvayā mṛdūn snigdhān ślakṣṇān śabdān prayuṅkte . \\
\noindent \textbf{(P\_5,1.16)} laukikī vivakṣā yatra prāyasya sampratyayaḥ . \\
\noindent \textbf{(P\_5,1.16)} prāyaḥ iti lokaḥ vyapadiśyate . \\
\noindent \textbf{(P\_5,1.16)} na ca prāsādaḥ devadattasya syāt prākāraḥ nagarasya syāt iti atra utpadyamānena pratyayena prāyasya sampratyayaḥ syāt . \\
\noindent \textbf{(P\_5,1.16)} yadi evam na arthaḥ syādgrahaṇena . \\
\noindent \textbf{(P\_5,1.16)} na hi prāsādaḥ devadattasya syāt prākāraḥ nagarasya syāt iti atra utpadyamānena pratyayena prāyasya sampratyayaḥ syāt . \\
\noindent \textbf{(P\_5,1.19.1)} kimartham saṅkhyāyāḥ pṛthaggrahaṇam kriyate na saṅkhyā api parimāṇam eva tatra parimāṇaparyudāsena paryudāsaḥ bhaviṣyati . \\
\noindent \textbf{(P\_5,1.19.1)} evam tarhi siddhe sati yat saṅkhyayāḥ pṛthaggrahaṇam karoti tat jñāpayati ācāryaḥ anyā saṅkhyā anyat parimāṇam iti . \\
\noindent \textbf{(P\_5,1.19.1)} kim etasya jñāpane prayojanam . \\
\noindent \textbf{(P\_5,1.19.1)} aparimāṇabistācitakambaelbhyaḥ na taddhitaluki iti dvābhyām śatābhyām krītā dviśatā triśatā parimāṇaparyudāsena na bhavati iti . \\
\noindent \textbf{(P\_5,1.19.1)} yadi etat jñāpyate tat asya parimāṇam saṅkhyāyāḥ sañjñāsaṅghasūtrādhyayaneṣu iti viśeṣaṇam na prakalpate parimāṇam yā saṅkhyā iti . \\
\noindent \textbf{(P\_5,1.19.1)} iha ca krītavat parimāṇat iti saṅkhyāvihitasya pratyayasya atideśaḥ na prakalpate . \\
\noindent \textbf{(P\_5,1.19.1)} śatasya vikāraḥ śatyaḥ śatikaḥ . \\
\noindent \textbf{(P\_5,1.19.1)} sāhasraḥ iti . \\
\noindent \textbf{(P\_5,1.19.1)} yat tāvat ucyate tat jñāpayati ācāryaḥ anyā saṅkhyā anyat parimāṇam iti . \\
\noindent \textbf{(P\_5,1.19.1)} nyāyasiddham eva etat . \\
\noindent \textbf{(P\_5,1.19.1)} bhedamātram saṅkhyā āha . \\
\noindent \textbf{(P\_5,1.19.1)} yat ca iṣīkāntam yat ca aparimāṇam sarvasya saṅkhyā bhedamātram bravīti . \\
\noindent \textbf{(P\_5,1.19.1)} parimāṇam tu sarvataḥ . \\
\noindent \textbf{(P\_5,1.19.1)} sarvataḥ mānam iti ca ataḥ parimāṇam iti . \\
\noindent \textbf{(P\_5,1.19.1)} prasthasya ca samānākṛteḥ na kutaḥ cit viśeṣaḥ gamyate na ca unmānataḥ na parimāṇataḥ na pramāṇataḥ . \\
\noindent \textbf{(P\_5,1.19.1)} kim punaḥ unmānam kim parimāṇam kim pramāṇam . \\
\noindent \textbf{(P\_5,1.19.1)} ūrdhvamānam kila unmānam . \\
\noindent \textbf{(P\_5,1.19.1)} ūrdhvam yat mīyate tat unmānam . \\
\noindent \textbf{(P\_5,1.19.1)} parimāṇam tu sarvataḥ . sarvataḥ mānam iti ca ataḥ parimāṇam . \\
\noindent \textbf{(P\_5,1.19.1)} kuta etat . \\
\noindent \textbf{(P\_5,1.19.1)} pariḥ sarvatobhāve vartate . \\
\noindent \textbf{(P\_5,1.19.1)} āyāmaḥ tu pramāṇam syāt . āyāmavivakṣāyām pramāṇam iti etat bhavati . \\
\noindent \textbf{(P\_5,1.19.1)} saṅkhyā bāhyā tu sarvataḥ . \\
\noindent \textbf{(P\_5,1.19.1)} ātaḥ ca sarvataḥ saṅkhyā bāhyā . \\
\noindent \textbf{(P\_5,1.19.1)} bhedabhāvam bravīti eṣā na eṣā mānam kutaḥ cana . evam ca kṛtvā saṅkhyāyāḥ pṛthaggrahaṇam kriyate . \\
\noindent \textbf{(P\_5,1.19.1)} yat api ucyate tat asya parimāṇam saṅkhyāyāḥ sañjñāsaṅghasūtrādhyayaneṣu iti viśeṣaṇam na prakalpate iti . \\
\noindent \textbf{(P\_5,1.19.1)} āha ayam parimāṇam yā saṅkhyā iti . \\
\noindent \textbf{(P\_5,1.19.1)} na ca asti saṅkhyā parimāṇam . \\
\noindent \textbf{(P\_5,1.19.1)} tatra vacanāt iyatī vivakṣā bhaviṣyati . \\
\noindent \textbf{(P\_5,1.19.1)} yad api ucyate krītavat parimāṇat iti saṅkhyāvihitasya pratyayasya atideśaḥ na prakalpate iti . \\
\noindent \textbf{(P\_5,1.19.1)} saṅkhyāyāḥ iti ca tatra vaktavyam . \\
\noindent \textbf{(P\_5,1.19.2)} kim punaḥ ime ṭhagādayaḥ prāk arhāt bhavanti āhosvit saha arheṇa . \\
\noindent \textbf{(P\_5,1.19.2)} kaḥ ca atra viśeṣaḥ . \\
\noindent \textbf{(P\_5,1.19.2)} ṭhagādayaḥ prāk arhāt cet arhe tadvidhiḥ . \\
\noindent \textbf{(P\_5,1.19.2)} ṭhagādayaḥ prāk arhāt cet arhe tadvidhiḥ . \\
\noindent \textbf{(P\_5,1.19.2)} arhe ṭhagādayaḥ vidheyāḥ . \\
\noindent \textbf{(P\_5,1.19.2)} śatam arhati śatyaḥ śatikaḥ . \\
\noindent \textbf{(P\_5,1.19.2)} sāharaḥ iti . \\
\noindent \textbf{(P\_5,1.19.2)} vasne vasanāt siddham . \\
\noindent \textbf{(P\_5,1.19.2)} iha yaḥ śatam arhati śatam tasya vasnaḥ bhavati . \\
\noindent \textbf{(P\_5,1.19.2)} tatra saḥ asya aṃśavasnabhṛtayaḥ iti eva siddham . \\
\noindent \textbf{(P\_5,1.19.2)} vasne vacanāt siddham iti cet māṃsaudanikādiṣu aprāptiḥ . \\
\noindent \textbf{(P\_5,1.19.2)} vasne vacanāt siddham iti cet māṃsaudanikādiṣu aprāptiḥ . \\
\noindent \textbf{(P\_5,1.19.2)} māṃsaudanikaḥ atithiḥ śvaitacchatrikaḥ kālāyasūpikaḥ . \\
\noindent \textbf{(P\_5,1.19.2)} tathā guṇānām paripraśnaḥ bhavati . \\
\noindent \textbf{(P\_5,1.19.2)} kim ayam brāhmaṇaḥ arhati . \\
\noindent \textbf{(P\_5,1.19.2)} śatam arhati śatyaḥ śatikaḥ sāhasraḥ naiṣkikaḥ iti na sidhyati . \\
\noindent \textbf{(P\_5,1.19.2)} santu tarhi sahārheṇa . \\
\noindent \textbf{(P\_5,1.19.2)} ā arhāt cet bhojanādiṣu atrprasaṅgaḥ . ā arhāt cet bhojanādiṣu atrprasaṅgaḥ bhavati . \\
\noindent \textbf{(P\_5,1.19.2)} bhojanam arhati . \\
\noindent \textbf{(P\_5,1.19.2)} pānam arhati iti . \\
\noindent \textbf{(P\_5,1.19.2)} kim ucyate bhojanādiṣu atrprasaṅgaḥ iti yadā chedādibhyaḥ iti ucyate . \\
\noindent \textbf{(P\_5,1.19.2)} avaśyam māṃsaudanikādyartham yogavibhāgaḥ kartavyaḥ . \\
\noindent \textbf{(P\_5,1.19.2)} tat arhati . \\
\noindent \textbf{(P\_5,1.19.2)} tataḥ chedādibhyaḥ nityam iti . \\
\noindent \textbf{(P\_5,1.19.2)} tasmin kriyamāṇe bhojanādiṣu atrprasaṅgaḥ bhavati . \\
\noindent \textbf{(P\_5,1.19.2)} uktam vā . kim uktam . \\
\noindent \textbf{(P\_5,1.19.2)} anabhidhānāt iti . \\
\noindent \textbf{(P\_5,1.19.2)} anabhidhānāt bhojanādiṣu atiprasaṅgaḥ na bhavati [R: bhaviṣyati] . \\
\noindent \textbf{(P\_5,1.19.2)} atha vā yogavibhāgaḥ na kariṣyate . \\
\noindent \textbf{(P\_5,1.19.2)} katham māṃsaudanikaḥ atithiḥ śvaitacchatrikaḥ kālāyasūpikaḥ . \\
\noindent \textbf{(P\_5,1.19.2)} asmin dīyate asmai iti ca evam etat siddham . \\
\noindent \textbf{(P\_5,1.19.2)} atha vā punaḥ astu prāk arhāt . \\
\noindent \textbf{(P\_5,1.19.2)} nanu ca uktam ṭhagādayaḥ prāk arhāt cet arhe tadvidhiḥ iti . \\
\noindent \textbf{(P\_5,1.19.2)} parihṛtam etat vasne vacanāt siddham iti . \\
\noindent \textbf{(P\_5,1.19.2)} nanu ca uktam vasne vacanāt siddham iti cet māṃsaudanikādiṣu aprāptiḥ . \\
\noindent \textbf{(P\_5,1.19.2)} māṃsaudanikaḥ atithiḥ śvaitacchatrikaḥ kālāyasūpikaḥ . \\
\noindent \textbf{(P\_5,1.19.2)} tathā guṇānām paripraśnaḥ bhavati . \\
\noindent \textbf{(P\_5,1.19.2)} kim ayam brāhmaṇaḥ arhati . \\
\noindent \textbf{(P\_5,1.19.2)} śatam arhati śatyaḥ śatikaḥ sāhasraḥ naiṣkikaḥ iti na sidhyati . \\
\noindent \textbf{(P\_5,1.19.2)} na eṣaḥ doṣaḥ . \\
\noindent \textbf{(P\_5,1.19.2)} asmin dīyate asmai iti ca evam etat siddham . \\
\noindent \textbf{(P\_5,1.20.1)} asamāse iti kimartham . \\
\noindent \textbf{(P\_5,1.20.1)} paramaniṣkeṇa krītam paramanaiṣkikam . \\
\noindent \textbf{(P\_5,1.20.1)} na etat asti . \\
\noindent \textbf{(P\_5,1.20.1)} niṣkaśabdāt pratyayaḥ vidhīyate . \\
\noindent \textbf{(P\_5,1.20.1)} tatra kaḥ prasaṅgaḥ yat paramaniṣkaśabdāt syāt . \\
\noindent \textbf{(P\_5,1.20.1)} na eva prāpnoti na arthaḥ pratiṣedhena . \\
\noindent \textbf{(P\_5,1.20.1)} tadantavidhinā prāpnoti . \\
\noindent \textbf{(P\_5,1.20.1)} grahaṇavatā prātipadikena tadantividhiḥ pratiṣidhyate . \\
\noindent \textbf{(P\_5,1.20.1)} niṣkādiṣu asamāsagrahaṇam jñāpakam pūrvatra tadantāpratiṣedhasya . niṣkādiṣu asamāsagrahaṇam kriyate jñāpakārtham . \\
\noindent \textbf{(P\_5,1.20.1)} kim jñāpyam . \\
\noindent \textbf{(P\_5,1.20.1)} pūrvatra tadantvidheḥ pratiṣedhaḥ na bhavati iti . \\
\noindent \textbf{(P\_5,1.20.1)} kim etasya jñāpane prayojanam . \\
\noindent \textbf{(P\_5,1.20.1)} prāk vateḥ ṭhañ iti atra tadantavidhiḥ siddhaḥ bhavati . \\
\noindent \textbf{(P\_5,1.20.1)} na etat asti prayojanam . \\
\noindent \textbf{(P\_5,1.20.1)} grahaṇavatā prātipadikena tadantividhiḥ pratiṣidhyate . \\
\noindent \textbf{(P\_5,1.20.1)} na ca ṭhañvidhau kā cit prakṛtiḥ gṛhyate . \\
\noindent \textbf{(P\_5,1.20.1)} idam tarhi prayojanam . \\
\noindent \textbf{(P\_5,1.20.1)} ā arhāt agopucchasaṅkhyāparimāṇāt ṭhak . \\
\noindent \textbf{(P\_5,1.20.1)} paramagopucchena krītam pāramagopucchikam . \\
\noindent \textbf{(P\_5,1.20.1)} atra tadantavidhiḥ siddhaḥ bhavati . \\
\noindent \textbf{(P\_5,1.20.1)} etat api na asti prayojanam . \\
\noindent \textbf{(P\_5,1.20.1)} vidhau pratiṣedhaḥ pratiṣedhaḥ ca ayam . \\
\noindent \textbf{(P\_5,1.20.1)} evam tarhi jñāpayati ācāryaḥ itaḥ uttaram tadantavidheḥ pratiṣedhaḥ na bhavati iti . \\
\noindent \textbf{(P\_5,1.20.1)} kim etasya jñāpane prayojanam . \\
\noindent \textbf{(P\_5,1.20.1)} pāryāṇaturāyaṇacāndrāyaṇam vartayati dvaipārāyaṇikaḥ traipārāyaṇikaḥ . \\
\noindent \textbf{(P\_5,1.20.1)} atra tadantavidhiḥ siddhaḥ bhavati . \\
\noindent \textbf{(P\_5,1.20.1)} etat api na asti prayojanam . \\
\noindent \textbf{(P\_5,1.20.1)} vakṣyati etat . \\
\noindent \textbf{(P\_5,1.20.1)} prāk vateḥ saṅkhyāpūrvapadānām tadantagrahaṇam aluki . \\
\noindent \textbf{(P\_5,1.20.1)} pūrvatra eva tarhi prayojanam . \\
\noindent \textbf{(P\_5,1.20.1)} khalayavamāṣatilavṛṣabrahmaṇaḥ ca iti . \\
\noindent \textbf{(P\_5,1.20.1)} kṛṣṇatilebhyaḥ hitaḥ kṛṣṇatilyaḥ . \\
\noindent \textbf{(P\_5,1.20.1)} rājamāṣebhyaḥ hitam rājamāṣyam . \\
\noindent \textbf{(P\_5,1.20.2)} prāk vateḥ saṅkhyāpūrvapadānām tadantagrahaṇam aluki . \\
\noindent \textbf{(P\_5,1.20.2)} prāk vateḥ saṅkhyāpūrvapadānām tadantagrahaṇam aluki kartavyam . \\
\noindent \textbf{(P\_5,1.20.2)} pāryāṇaturāyaṇacāndrāyaṇam vartayati dvaipārāyaṇikaḥ traipārāyaṇikaḥ . \\
\noindent \textbf{(P\_5,1.20.2)} aluki iti kimartham . \\
\noindent \textbf{(P\_5,1.20.2)} dvābhyām śūrpābhyām krītam dviśūrpam . \\
\noindent \textbf{(P\_5,1.20.2)} triśūrpam . \\
\noindent \textbf{(P\_5,1.20.2)} dviśūrpeṇa krītam dvaiśaurpikam . \\
\noindent \textbf{(P\_5,1.20.2)} traiśaurpikam . \\
\noindent \textbf{(P\_5,1.21)} śatapratiṣedhe anyaśatatve apratiṣedhaḥ . \\
\noindent \textbf{(P\_5,1.21)} śatapratiṣedhe anyaśatatve pratiṣedhaḥ na bhavati iti vaktavyam. iha mā bhūt . \\
\noindent \textbf{(P\_5,1.21)} śatena krītam śatyam śāṭakaśatam iti . \\
\noindent \textbf{(P\_5,1.21)} anyaśatatve iti kim . \\
\noindent \textbf{(P\_5,1.21)} śatakam nidānam . \\
\noindent \textbf{(P\_5,1.22)} ḍateḥ ca iti vaktavyam iha api yathā syāt . \\
\noindent \textbf{(P\_5,1.22)} katibhiḥ krītam katikam . \\
\noindent \textbf{(P\_5,1.22)} kim punaḥ kāraṇam na sidhyati . \\
\noindent \textbf{(P\_5,1.22)} tyantāyāḥ na iti pratiṣedhaḥ prāpnoti . \\
\noindent \textbf{(P\_5,1.22)} tipratiṣedhāt ḍatigrahaṇam iti cet arthavadgrahaṇāt siddham . \\
\noindent \textbf{(P\_5,1.22)} arthavataḥ tiśabdasya grahaṇam na ca ḍateḥ tiśabdaḥ arthavān . \\
\noindent \textbf{(P\_5,1.22)} na eṣā paribhāṣā iha śakyā vijñātum . \\
\noindent \textbf{(P\_5,1.22)} na hi kevalena pratyayena arthaḥ gamyate . \\
\noindent \textbf{(P\_5,1.22)} kena tarhi . \\
\noindent \textbf{(P\_5,1.22)} saprakṛtikena . \\
\noindent \textbf{(P\_5,1.22)} kva tarhi eṣā paribhāṣā bhavati . \\
\noindent \textbf{(P\_5,1.22)} yāni etāni śabdasaṅghātagrahaṇāni . \\
\noindent \textbf{(P\_5,1.22)} tat tarhi vaktavyam . \\
\noindent \textbf{(P\_5,1.22)} na vaktavyam . \\
\noindent \textbf{(P\_5,1.22)} arthavadgrahaṇāt siddham . \\
\noindent \textbf{(P\_5,1.22)} nanu ca uktam na eṣā paribhāṣā iha śakyā vijñātum . \\
\noindent \textbf{(P\_5,1.22)} na hi kevalena pratyayena arthaḥ gamyate iti . \\
\noindent \textbf{(P\_5,1.22)} kevalena api pratyayena arthaḥ gamyate . \\
\noindent \textbf{(P\_5,1.22)} katham . \\
\noindent \textbf{(P\_5,1.22)} uktam anvayavyatirekābhyām . \\
\noindent \textbf{(P\_5,1.23)} kasya ayam iṭ vidhīyate . \\
\noindent \textbf{(P\_5,1.23)} kanaḥ iti āha . \\
\noindent \textbf{(P\_5,1.23)} tat kanaḥ grahaṇam kartavyam . \\
\noindent \textbf{(P\_5,1.23)} akriyamāṇe hi kanaḥ grahaṇe pratyayādhikārāt pratyayaḥ ayam vijñāyeta . \\
\noindent \textbf{(P\_5,1.23)} ṭitkaraṇasāmarthyāt ādiḥ bhaviṣyati . \\
\noindent \textbf{(P\_5,1.23)} asti anyat ṭitkaraṇe prayojanam . \\
\noindent \textbf{(P\_5,1.23)} ṭitaḥ iti īkāraḥ yathā syāt . \\
\noindent \textbf{(P\_5,1.23)} akārāntaprakaraṇe īkāraḥ na ca eṣaḥ akārāntaḥ . \\
\noindent \textbf{(P\_5,1.23)} evam api kutaḥ etat ṭitkaraṇasāmarthyāt ādiḥ bhaviṣyati na punaḥ akārāntaprakaraṇe sati anakārāntāt api īkāraḥ syāt . \\
\noindent \textbf{(P\_5,1.23)} tasmāt kaṇaḥ grahaṇam kartavyam . \\
\noindent \textbf{(P\_5,1.23)} na kartavyam . \\
\noindent \textbf{(P\_5,1.23)} prakṛtam anuvartate . \\
\noindent \textbf{(P\_5,1.23)} kva prakṛtam . \\
\noindent \textbf{(P\_5,1.23)} saṅkhyāyāḥ atiśadantāyāḥ kan iti . \\
\noindent \textbf{(P\_5,1.23)} tat vai prathamānirdiṣṭam ṣaṣṭhīnirdiṣṭena ca iha arthaḥ . \\
\noindent \textbf{(P\_5,1.23)} vatoḥ iti eṣā pañcamī kan iti prathamāyāḥ ṣaṣṭhīm prakalpayiṣyati tasmāt iti uttarasya iti . \\
\noindent \textbf{(P\_5,1.23)} pratyayavidhiḥ ayam . \\
\noindent \textbf{(P\_5,1.23)} na ca pratyayavidhau pañcamyāḥ prakalpikāḥ bhavanti . \\
\noindent \textbf{(P\_5,1.23)} na ayam pratyayavidhiḥ . \\
\noindent \textbf{(P\_5,1.23)} vihitaḥ pratyayaḥ prakṛtaḥ ca anuvartate . \\
\noindent \textbf{(P\_5,1.24)} asañjñāyām iti kimartham . \\
\noindent \textbf{(P\_5,1.24)} triṃśatkaḥ viṃśatkaḥ . \\
\noindent \textbf{(P\_5,1.24)} katham ca atra kan bhavati . \\
\noindent \textbf{(P\_5,1.24)} saṅkhyāyāḥ kan bhavati iti . \\
\noindent \textbf{(P\_5,1.24)} atiśadantāyāḥ iti pratiṣedhaḥ prāpnoti . \\
\noindent \textbf{(P\_5,1.24)} evam tarhi ācāryapravṛttiḥ jñāpayati bhavati atra kan iti yat ayam viṃśatikāt khaḥ iti pratyayāntanipātanam karoti . \\
\noindent \textbf{(P\_5,1.24)} viṃśateḥ etat jñāpakam syāt . \\
\noindent \textbf{(P\_5,1.24)} na iti āha . \\
\noindent \textbf{(P\_5,1.24)} yogāpekṣam jñāpakam . \\
\noindent \textbf{(P\_5,1.24)} atha vā yogavibhāgaḥ kariṣyate . \\
\noindent \textbf{(P\_5,1.24)} viṃśatitriṃśabhyām kan bhavati iti . \\
\noindent \textbf{(P\_5,1.24)} tataḥ ḍvun asañjñāyām iti . \\
\noindent \textbf{(P\_5,1.25)} ṭithan ardhāt ca . \\
\noindent \textbf{(P\_5,1.25)} ṭithan ardhāt ca iti vaktavyam . \\
\noindent \textbf{(P\_5,1.25)} ardhikaḥ ardhikī . \\
\noindent \textbf{(P\_5,1.25)} kārṣāpaṇāt vā pratiḥ ca . \\
\noindent \textbf{(P\_5,1.25)} kārṣāpaṇāt ṭiṭhan vaktavyaḥ vā ca pratiḥ ādeśaḥ vaktavyaḥ . \\
\noindent \textbf{(P\_5,1.25)} kārṣāpaṇikaḥ kārṣāpaṇikī pratikaḥ pratikī . \\
\noindent \textbf{(P\_5,1.28)} dvigoḥ luki uktam . \\
\noindent \textbf{(P\_5,1.28)} kim uktam . \\
\noindent \textbf{(P\_5,1.28)} tatra tāvat uktam dvigoḥ luki tannimittagrahaṇam. arthaviśeṣāsampratyaye atannimittāt api iti . \\
\noindent \textbf{(P\_5,1.28)} iha api dvigoḥ luki tannimittagrahaṇam kartavyam . \\
\noindent \textbf{(P\_5,1.28)} dvigoḥ nimittam yaḥ taddhitaḥ tasya luk bhavati iti vaktavyam . \\
\noindent \textbf{(P\_5,1.28)} iha mā bhūt . \\
\noindent \textbf{(P\_5,1.28)} dvābhyām śūrpābhyām krītam dviśūrpam . \\
\noindent \textbf{(P\_5,1.28)} triśūrpam . \\
\noindent \textbf{(P\_5,1.28)} dviśūrpeṇa krītam dvaiśaurpikam . \\
\noindent \textbf{(P\_5,1.28)} traiśaurpikam . \\
\noindent \textbf{(P\_5,1.28)} arthaviśeṣāsampratyaye atannimittāt api . \\
\noindent \textbf{(P\_5,1.28)} arthaviśeṣasya asampratyaye atannimittāt api vaktavyam . \\
\noindent \textbf{(P\_5,1.28)} kim prayojanam . \\
\noindent \textbf{(P\_5,1.28)} dvayoḥ śūrpayoḥ samāhāraḥ dviśūrpī . \\
\noindent \textbf{(P\_5,1.28)} dviśūrpyā krītam iti vigṛhya dviśūrpam iti eva yathā syāt . \\
\noindent \textbf{(P\_5,1.28)} atha kriyamāṇe api tannimittagrahaṇe katham idam vijñāyate . \\
\noindent \textbf{(P\_5,1.28)} tasya nimittam tannimittam tannimittāt iti āhosvit saḥ nimittam asya saḥ ayam tannimittaḥ tannimittāt iti . \\
\noindent \textbf{(P\_5,1.28)} kim ca ataḥ . \\
\noindent \textbf{(P\_5,1.28)} yadi vijñāyate tasya nimittam tannimittam tannimittāt iti kriyamāṇe api tannimittagrahaṇe atra prāpnoti . \\
\noindent \textbf{(P\_5,1.28)} dvābhyām śūrpābhyām krītam dviśūrpam . \\
\noindent \textbf{(P\_5,1.28)} triśūrpam . \\
\noindent \textbf{(P\_5,1.28)} dviśūrpeṇa krītam dvaiśaurpikam . \\
\noindent \textbf{(P\_5,1.28)} traiśaurpikam . \\
\noindent \textbf{(P\_5,1.28)} atha vijñāyate saḥ nimittam asya saḥ ayam tannimittaḥ tannimittāt iti na doṣaḥ bhavati . \\
\noindent \textbf{(P\_5,1.28)} yatha na doṣaḥ tathā astu . \\
\noindent \textbf{(P\_5,1.28)} saḥ nimittam asya saḥ ayam tannimittaḥ tannimittāt iti vijñāyate . \\
\noindent \textbf{(P\_5,1.28)} kutaḥ etat . \\
\noindent \textbf{(P\_5,1.28)} yat ayam āha arthaviśeṣāsampratyaye atannimittāt api iti . \\
\noindent \textbf{(P\_5,1.28)} tat tarhi tannimittagrahaṇam kartavyam . \\
\noindent \textbf{(P\_5,1.28)} na kartavyam . \\
\noindent \textbf{(P\_5,1.28)} dvigoḥ iti na eṣā pañcamī . \\
\noindent \textbf{(P\_5,1.28)} kā tarhi . \\
\noindent \textbf{(P\_5,1.28)} sambandhaṣaṣṭhī . \\
\noindent \textbf{(P\_5,1.28)} dvigoḥ taddhitasya luk bhavati . \\
\noindent \textbf{(P\_5,1.28)} kim ca dvigoḥ taddhitaḥ . \\
\noindent \textbf{(P\_5,1.28)} nimittam . \\
\noindent \textbf{(P\_5,1.28)} yasmin dviguḥ iti etat bhavati . \\
\noindent \textbf{(P\_5,1.28)} kasmin ca etat bhavati . \\
\noindent \textbf{(P\_5,1.28)} pratyaye . \\
\noindent \textbf{(P\_5,1.28)} idam tarhi vaktavyam arthaviśeṣāsampratyaye atannimittāt api iti . \\
\noindent \textbf{(P\_5,1.28)} etat ca na vaktavyam . \\
\noindent \textbf{(P\_5,1.28)} iha asmābhiḥ traiśabdyam sādhyam . \\
\noindent \textbf{(P\_5,1.28)} dvābhyām śūrpābhyām krītam dviśūrpam . \\
\noindent \textbf{(P\_5,1.28)} triśūrpam . \\
\noindent \textbf{(P\_5,1.28)} dviśūrpeṇa krītam dvaiśaurpikam . \\
\noindent \textbf{(P\_5,1.28)} traiśaurpikam iti . \\
\noindent \textbf{(P\_5,1.28)} tatra dvayoḥ śabdayoḥ samānārthayoḥ ekena vigrahaḥ aparasmāt utpattiḥ bhaviṣyati aviravikanyāyena . \\
\noindent \textbf{(P\_5,1.28)} tat yathā . \\
\noindent \textbf{(P\_5,1.28)} aveḥ māṃsam iti vigṛhya avikaśabdāt utpattiḥ bhavati . \\
\noindent \textbf{(P\_5,1.28)} evam iha api dvābhyām śūrpābhyām krītam iti vigṛhya dviśūrpam iti bhaviṣyati . \\
\noindent \textbf{(P\_5,1.28)} dviśūrpyā krītam iti vigṛhya vākyam eva bhaviṣyati . \\
\noindent \textbf{(P\_5,1.28)} atha asañjñāyām iti kimartham . \\
\noindent \textbf{(P\_5,1.28)} pāñcalohitikam pāñcakalāpikam . \\
\noindent \textbf{(P\_5,1.28)} sañjñāpratiṣedhānarthakyam ca tannimittatvāt lopasya . \\
\noindent \textbf{(P\_5,1.28)} sañjñāpratiṣedhaḥ ca anarthakaḥ . \\
\noindent \textbf{(P\_5,1.28)} kim kāraṇam . \\
\noindent \textbf{(P\_5,1.28)} tannimittatvāt lopasya . \\
\noindent \textbf{(P\_5,1.28)} na antareṇa taddhitam taddhitasya ca lukam dviguḥ sañjñā asti . \\
\noindent \textbf{(P\_5,1.28)} yaḥ tasmāt utpadyate na asu tannimittam syāt . \\
\noindent \textbf{(P\_5,1.28)} evam tarhi idam syāt . \\
\noindent \textbf{(P\_5,1.28)} pañcānām lohitānām samāhāraḥ pañcalohitī pañcalohityā krītam . \\
\noindent \textbf{(P\_5,1.28)} atra api pañcalohitam iti eva bhavitavyam . \\
\noindent \textbf{(P\_5,1.28)} katham . \\
\noindent \textbf{(P\_5,1.28)} uktam hi etat arthaviśeṣāsampratyaye atannimittāt api iti . \\
\noindent \textbf{(P\_5,1.28)} uktam saṅkhyātve prayojanam tasmāt iha adhyardhagrahaṇānarthakyam . \\
\noindent \textbf{(P\_5,1.28)} uktam saṅkhyātve adhyardhagrahaṇasya prayojanam . \\
\noindent \textbf{(P\_5,1.28)} kim uktam . \\
\noindent \textbf{(P\_5,1.28)} adhyardhagrahaṇam ca samāsakanvidhyartham luki ca agrahaṇam iti . \\
\noindent \textbf{(P\_5,1.28)} tasmāt iha adhyardhagrahaṇānarthakyam . \\
\noindent \textbf{(P\_5,1.28)} tasmāt iha adhyardhagrahaṇam anarthakam . \\
\noindent \textbf{(P\_5,1.28)} dvigoḥ iti eva luk siddhaḥ . \\
\noindent \textbf{(P\_5,1.29)} kārṣāpaṇasahasrābhyām suvarṇaśatamānayoḥ upasaṅkhyānam kartavyam . \\
\noindent \textbf{(P\_5,1.29)} adhyardhasuvarṇam adhyardhasauvarṇikam adhyardhaśatamānam adhyardhaśātamānam dviśatamānam dviśātamānam \\
\noindent \textbf{(P\_5,1.30-31)} KA\_II,349.12-18 Ro\_IV,37-38 {1/11}     dvitribhyām dvaiyogyam . \\
\noindent \textbf{(P\_5,1.30-31)} KA\_II,349.12-18 Ro\_IV,37-38 {2/11}     dvitribhyām iti yat ucyate dvaiyogyam etat draṣṭavyam . \\
\noindent \textbf{(P\_5,1.30-31)} KA\_II,349.12-18 Ro\_IV,37-38 {3/11}     kim idam dvaiyogyam iti . \\
\noindent \textbf{(P\_5,1.30-31)} KA\_II,349.12-18 Ro\_IV,37-38 {4/11}     dvayoḥ yogayoḥ bhavam dviyogam . \\
\noindent \textbf{(P\_5,1.30-31)} KA\_II,349.12-18 Ro\_IV,37-38 {5/11}     dviyogasya bhāvaḥ dvaiyogyam iti . \\
\noindent \textbf{(P\_5,1.30-31)} KA\_II,349.12-18 Ro\_IV,37-38 {6/11}     dveṣyam vijānīyāt : aviśeṣeṇa itaḥ uttaram dvitribhyām iti . \\
\noindent \textbf{(P\_5,1.30-31)} KA\_II,349.12-18 Ro\_IV,37-38 {7/11}     tat ācāryaḥ suhṛt bhūtvā anvācaṣṭe : dvitribhyām dvaiyogyam iti . \\
\noindent \textbf{(P\_5,1.30-31)} KA\_II,349.12-18 Ro\_IV,37-38 {8/11}     tatra ca bahugrahaṇam . \\
\noindent \textbf{(P\_5,1.30-31)} KA\_II,349.12-18 Ro\_IV,37-38 {9/11}     tatra ca bahugrahaṇam kartavyam . \\
\noindent \textbf{(P\_5,1.30-31)} KA\_II,349.12-18 Ro\_IV,37-38 {10/11}     bahuniṣkam bahunaiṣkikam . \\
\noindent \textbf{(P\_5,1.30-31)} KA\_II,349.12-18 Ro\_IV,37-38 {11/11}     bahubistam bahubaistikam . \\
\noindent \textbf{(P\_5,1.33)} khāryāḥ īkan kevalāyāḥ ca . \\
\noindent \textbf{(P\_5,1.33)} khāryāḥ īkan kevalāyāḥ ca iti vaktavyam . \\
\noindent \textbf{(P\_5,1.33)} khārīkam . \\
\noindent \textbf{(P\_5,1.33)} kākiṇyāḥ ca upasaṅkhyānam .kākiṇyāḥ ca upasaṅkhyānam kartavyam . \\
\noindent \textbf{(P\_5,1.33)} adhyardhakākiṇīkam dvikākiṇīkam . \\
\noindent \textbf{(P\_5,1.33)} kevalāyāḥ ca . \\
\noindent \textbf{(P\_5,1.33)} kevalāyāḥ ca iti vaktavyam . \\
\noindent \textbf{(P\_5,1.33)} kākiṇīkam . \\
\noindent \textbf{(P\_5,1.35)} śataśāṇābhyām vā . \\
\noindent \textbf{(P\_5,1.35)} śataśāṇābhyām vā iti vaktavyam . \\
\noindent \textbf{(P\_5,1.35)} adhyardhaśatam adhyardhaśatyam pañcaśatam pañcaśatyam adhyardhaśāṇam adhyardhaśāṇyam pañcaśāṇam pañcaśāṇyam . \\
\noindent \textbf{(P\_5,1.35)} dvitripūrvāt aṇ ca . \\
\noindent \textbf{(P\_5,1.35)} dvitripūrvāt aṇ ca iti vaktavyam . \\
\noindent \textbf{(P\_5,1.35)} dviśāṇam triśāṇam dvaiśāṇam traiśāṇam dviśāṇyam triśāṇyam . \\
\noindent \textbf{(P\_5,1.37)} tena krītam iti karaṇāt . \\
\noindent \textbf{(P\_5,1.37)} tena krītam iti atra karaṇāt iti vaktavyam . \\
\noindent \textbf{(P\_5,1.37)} iha mā bhūt . \\
\noindent \textbf{(P\_5,1.37)} devadattena krītam . \\
\noindent \textbf{(P\_5,1.37)} yajñadattena krītam iti . \\
\noindent \textbf{(P\_5,1.37)} akartrekāntāt . \\
\noindent \textbf{(P\_5,1.37)} akartrekāntāt iti vaktavyam . \\
\noindent \textbf{(P\_5,1.37)} iha mā bhūt . \\
\noindent \textbf{(P\_5,1.37)} devadattena pāṇinā krītam iti . \\
\noindent \textbf{(P\_5,1.37)} saṅkhyaikavacanāt dvigoḥ ca upasaṅkhyānam . \\
\noindent \textbf{(P\_5,1.37)} saṅkhyāyāḥ iti vaktavyam . \\
\noindent \textbf{(P\_5,1.37)} iha api yathā syāt . \\
\noindent \textbf{(P\_5,1.37)} pañcabhiḥ krītam pañcakam . \\
\noindent \textbf{(P\_5,1.37)} kim punaḥ kāraṇam na sidhyati . \\
\noindent \textbf{(P\_5,1.37)} ekavacanāntāt iti vakṣyati . \\
\noindent \textbf{(P\_5,1.37)} tasya ayam purastāt apakarṣaḥ . \\
\noindent \textbf{(P\_5,1.37)} ekavacanāt . \\
\noindent \textbf{(P\_5,1.37)} ekavacanāntāt iti vaktavyam . \\
\noindent \textbf{(P\_5,1.37)} iha mā bhūt . \\
\noindent \textbf{(P\_5,1.37)} śūrpābhyām krītam . \\
\noindent \textbf{(P\_5,1.37)} śūrpaiḥ krītam iti . \\
\noindent \textbf{(P\_5,1.37)} dvigoḥ ca . \\
\noindent \textbf{(P\_5,1.37)} dvigoḥ ca iti vaktavyam . \\
\noindent \textbf{(P\_5,1.37)} iha api yathā syāt . \\
\noindent \textbf{(P\_5,1.37)} dvābhyām śūrpābhyām krītam dviśūrpam triśūrpam iti . \\
\noindent \textbf{(P\_5,1.37)} yadi ekavacanāntāt iti ucyate mudgaiḥ krītam maudgikam māṣaiḥ krītam māṣikam iti na sidhyati . \\
\noindent \textbf{(P\_5,1.37)} parimāṇasya saṅkhyāyāḥ yat ekavacanam tadantāt iti vaktavyam . \\
\noindent \textbf{(P\_5,1.37)} tat tarhi ekavacanāntāt iti vaktavyam . \\
\noindent \textbf{(P\_5,1.37)} tasmin ca kriyamāṇe bahu vaktavyam bhavati . \\
\noindent \textbf{(P\_5,1.37)} na vaktavyam . \\
\noindent \textbf{(P\_5,1.37)} kasmāt na bhavati śūrpābhyām krītam . \\
\noindent \textbf{(P\_5,1.37)} śūrpaiḥ krītam iti . \\
\noindent \textbf{(P\_5,1.37)} uktam vā . \\
\noindent \textbf{(P\_5,1.37)} kim uktam . \\
\noindent \textbf{(P\_5,1.37)} anabhidhānāt iti . \\
\noindent \textbf{(P\_5,1.37)} yadi evam karaṇāt akartrekāntāt iti api na vaktavyam . \\
\noindent \textbf{(P\_5,1.37)} kartuḥ kartrekāntāt vā kasmāt na bhavati . \\
\noindent \textbf{(P\_5,1.37)} anabhidhānāt . \\
\noindent \textbf{(P\_5,1.38)} saṃyoganipātayoḥ kaḥ viśeṣaḥ . \\
\noindent \textbf{(P\_5,1.38)} saṃyogaḥ nāma saḥ bhavati idam kṛtvā idam avāpyate iti . \\
\noindent \textbf{(P\_5,1.38)} utpātaḥ nāma saḥ bhavati yādṛcchikaḥ bhedaḥ vā chedaḥ vā padmam vā parṇam vā . \\
\noindent \textbf{(P\_5,1.38)} tasyanimittaprakaraṇe vātapittaśleṣmabhyaḥ śamakopanayoḥ upasaṅkhyānam . \\
\noindent \textbf{(P\_5,1.38)} tasyanimittaprakaraṇe vātapittaśleṣmabhyaḥ śamakopanayoḥ upasaṅkhyānam kartavyam . \\
\noindent \textbf{(P\_5,1.38)} vātasya śamanam kopanam vā vātikam . \\
\noindent \textbf{(P\_5,1.38)} paittikam ślaiṣmikam . \\
\noindent \textbf{(P\_5,1.38)} sannipātāt ca . \\
\noindent \textbf{(P\_5,1.38)} sannipātāt ca iti vaktavyam . \\
\noindent \textbf{(P\_5,1.38)} sānnipātikam . \\
\noindent \textbf{(P\_5,1.39)} yatprakaraṇe brahmavarcasāt ca . \\
\noindent \textbf{(P\_5,1.39)} yatprakaraṇe brahmavarcasāt ca upasaṅkhyānam kartavyam . \\
\noindent \textbf{(P\_5,1.39)} brahmavarcasasya nimittam brahmvarcasyaḥ . \\
\noindent \textbf{(P\_5,1.39)} utpātaḥ vā . \\
\noindent \textbf{(P\_5,1.47)} tad asmin dīyate asmai iti ca . \\
\noindent \textbf{(P\_5,1.47)} tad asmin dīyate asmai iti ca iti vaktavyam . \\
\noindent \textbf{(P\_5,1.47)} pañca vṛddhiḥ vā āyaḥ vā lābhaḥ vā śulkaḥ vā upadā vā dīyate asmai pañcakaḥ . \\
\noindent \textbf{(P\_5,1.47)} saptakaḥ . \\
\noindent \textbf{(P\_5,1.47)} aṣṭakaḥ . \\
\noindent \textbf{(P\_5,1.47)} navakaḥ . \\
\noindent \textbf{(P\_5,1.47)} daśakaḥ . \\
\noindent \textbf{(P\_5,1.47)} tat tarhi upasaṅkhyānam kartavyam . \\
\noindent \textbf{(P\_5,1.47)} na kartavyam . \\
\noindent \textbf{(P\_5,1.47)} yat hi yasmai dīyate tasmin api tat dīyate . \\
\noindent \textbf{(P\_5,1.47)} tat asmin dīyate iti eva siddham . \\
\noindent \textbf{(P\_5,1.48)} ṭhanprakaraṇe anantāt upasaṅkhyānam . \\
\noindent \textbf{(P\_5,1.48)} ṭhanprakaraṇe anantāt upasaṅkhyānam kartavyam . \\
\noindent \textbf{(P\_5,1.48)} dvitīyakaḥ . \\
\noindent \textbf{(P\_5,1.48)} tṛtīyakaḥ . \\
\noindent \textbf{(P\_5,1.48)} kim punaḥ kāraṇam na sidhyati . \\
\noindent \textbf{(P\_5,1.48)} pūraṇāt iti ucyate . \\
\noindent \textbf{(P\_5,1.48)} na ca etat pūraṇāntam . \\
\noindent \textbf{(P\_5,1.48)} anā etat paryavapannam . \\
\noindent \textbf{(P\_5,1.48)} pūraṇam nāma arthaḥ . \\
\noindent \textbf{(P\_5,1.48)} tam artham āha tīyaśabdaḥ . \\
\noindent \textbf{(P\_5,1.48)} pūraṇam saḥ asau bhavati . \\
\noindent \textbf{(P\_5,1.48)} pūraṇantāt svārthe bhāge an . \\
\noindent \textbf{(P\_5,1.48)} saḥ api pūraṇam bhavati eva . \\
\noindent \textbf{(P\_5,1.52)} tat pacati iti droṇāt aṇ ca . \\
\noindent \textbf{(P\_5,1.52)} tat pacati iti droṇāt aṇ ca iti vaktavyam . \\
\noindent \textbf{(P\_5,1.52)} droṇam pacati . \\
\noindent \textbf{(P\_5,1.52)} drauṇī . \\
\noindent \textbf{(P\_5,1.52)} drauṇikī . \\
\noindent \textbf{(P\_5,1.55)} kulijāt ca iti siddhe lukkhagrahaṇānarthakyam pūrvsamin trikabhāvāt . \\
\noindent \textbf{(P\_5,1.55)} kulijāt ca iti eva siddham . \\
\noindent \textbf{(P\_5,1.55)} na arthaḥ lukkhagrahaṇena . \\
\noindent \textbf{(P\_5,1.55)} kim kāraṇam . \\
\noindent \textbf{(P\_5,1.55)} pūrvasmin trikabhāvāt . \\
\noindent \textbf{(P\_5,1.55)} pūrvasmin yoge sarvaḥ eṣaḥ trikaḥ nirdiśyate . \\
\noindent \textbf{(P\_5,1.55)} dvyāḍhakī . \\
\noindent \textbf{(P\_5,1.55)} dvyāḍhikī . \\
\noindent \textbf{(P\_5,1.55)} dvyāḍhakīnā . \\
\noindent \textbf{(P\_5,1.57-58.1)} KA\_II,352.18-354.6 Ro\_IV,43-46 {1/73}     sañjñāyām svārthe . \\
\noindent \textbf{(P\_5,1.57-58.1)} KA\_II,352.18-354.6 Ro\_IV,43-46 {2/73}     sañjñāyām svārthe pratyayaḥ utpādyaḥ . \\
\noindent \textbf{(P\_5,1.57-58.1)} KA\_II,352.18-354.6 Ro\_IV,43-46 {3/73}     pañca eva pañcakāḥ śakunayaḥ . \\
\noindent \textbf{(P\_5,1.57-58.1)} KA\_II,352.18-354.6 Ro\_IV,43-46 {4/73}     trikāḥ śālaṅkāyanāḥ . \\
\noindent \textbf{(P\_5,1.57-58.1)} KA\_II,352.18-354.6 Ro\_IV,43-46 {5/73}     saptakāḥ brahmavṛkṣāḥ . \\
\noindent \textbf{(P\_5,1.57-58.1)} KA\_II,352.18-354.6 Ro\_IV,43-46 {6/73}     tataḥ parimāṇini . \\
\noindent \textbf{(P\_5,1.57-58.1)} KA\_II,352.18-354.6 Ro\_IV,43-46 {7/73}     tataḥ paraḥ pratyayaḥ parimāṇini iti vaktavyam . \\
\noindent \textbf{(P\_5,1.57-58.1)} KA\_II,352.18-354.6 Ro\_IV,43-46 {8/73}     pañcakaḥ saṅghaḥ . \\
\noindent \textbf{(P\_5,1.57-58.1)} KA\_II,352.18-354.6 Ro\_IV,43-46 {9/73}     daśakaḥ saṅghaḥ . \\
\noindent \textbf{(P\_5,1.57-58.1)} KA\_II,352.18-354.6 Ro\_IV,43-46 {10/73}     jīvitaparimāṇe ca upasaṅkhyānam . \\
\noindent \textbf{(P\_5,1.57-58.1)} KA\_II,352.18-354.6 Ro\_IV,43-46 {11/73}     jīvitaparimāṇe ca upasaṅkhyānam kartavyam . \\
\noindent \textbf{(P\_5,1.57-58.1)} KA\_II,352.18-354.6 Ro\_IV,43-46 {12/73}     ṣaṣṭiḥ jīvitaparimāṇam asya ṣāṣṭikaḥ . \\
\noindent \textbf{(P\_5,1.57-58.1)} KA\_II,352.18-354.6 Ro\_IV,43-46 {13/73}     sāptatikaḥ . \\
\noindent \textbf{(P\_5,1.57-58.1)} KA\_II,352.18-354.6 Ro\_IV,43-46 {14/73}     jīvitaparimāṇe ca iti anarthakam vacanam kālāt iti siddhatvāt . \\
\noindent \textbf{(P\_5,1.57-58.1)} KA\_II,352.18-354.6 Ro\_IV,43-46 {15/73}     jīvitaparimāṇe ca iti anarthakam vacanam . \\
\noindent \textbf{(P\_5,1.57-58.1)} KA\_II,352.18-354.6 Ro\_IV,43-46 {16/73}     kim kāraṇam . \\
\noindent \textbf{(P\_5,1.57-58.1)} KA\_II,352.18-354.6 Ro\_IV,43-46 {17/73}     kālāt iti siddhatvāt . \\
\noindent \textbf{(P\_5,1.57-58.1)} KA\_II,352.18-354.6 Ro\_IV,43-46 {18/73}     kālāt iti eva siddham . \\
\noindent \textbf{(P\_5,1.57-58.1)} KA\_II,352.18-354.6 Ro\_IV,43-46 {19/73}     iha yasya ṣaṣṭiḥ jīvitaparimāṇam ṣaṣtim asu bhūtaḥ bhavati . \\
\noindent \textbf{(P\_5,1.57-58.1)} KA\_II,352.18-354.6 Ro\_IV,43-46 {20/73}     tatra tam adhīṣṭaḥ bhṛtaḥ bhūtaḥ bhāvī iti eva siddham . \\
\noindent \textbf{(P\_5,1.57-58.1)} KA\_II,352.18-354.6 Ro\_IV,43-46 {21/73}     avaśyam ca etat evam vijñeyam . \\
\noindent \textbf{(P\_5,1.57-58.1)} KA\_II,352.18-354.6 Ro\_IV,43-46 {22/73}     iha vacane hi lukprasaṅgaḥ . \\
\noindent \textbf{(P\_5,1.57-58.1)} KA\_II,352.18-354.6 Ro\_IV,43-46 {23/73}     iha hi kriyamāṇe luk prasajyeta : dviṣāṣṭikaḥ . \\
\noindent \textbf{(P\_5,1.57-58.1)} KA\_II,352.18-354.6 Ro\_IV,43-46 {24/73}     triṣāṣṭikaḥ . \\
\noindent \textbf{(P\_5,1.57-58.1)} KA\_II,352.18-354.6 Ro\_IV,43-46 {25/73}     anena sati luk bhavati . \\
\noindent \textbf{(P\_5,1.57-58.1)} KA\_II,352.18-354.6 Ro\_IV,43-46 {26/73}     tena sati kasmāt na bhavati . \\
\noindent \textbf{(P\_5,1.57-58.1)} KA\_II,352.18-354.6 Ro\_IV,43-46 {27/73}     ā arhāt iti ucyate . \\
\noindent \textbf{(P\_5,1.57-58.1)} KA\_II,352.18-354.6 Ro\_IV,43-46 {28/73}     na sidhyati . \\
\noindent \textbf{(P\_5,1.57-58.1)} KA\_II,352.18-354.6 Ro\_IV,43-46 {29/73}     kim kāraṇam . \\
\noindent \textbf{(P\_5,1.57-58.1)} KA\_II,352.18-354.6 Ro\_IV,43-46 {30/73}     na hi ime kālaśabdāḥ . \\
\noindent \textbf{(P\_5,1.57-58.1)} KA\_II,352.18-354.6 Ro\_IV,43-46 {31/73}     kim tarhi saṅkhyāśabdāḥ ime . \\
\noindent \textbf{(P\_5,1.57-58.1)} KA\_II,352.18-354.6 Ro\_IV,43-46 {32/73}     ime api kālaśabdāḥ . \\
\noindent \textbf{(P\_5,1.57-58.1)} KA\_II,352.18-354.6 Ro\_IV,43-46 {33/73}     katham . \\
\noindent \textbf{(P\_5,1.57-58.1)} KA\_II,352.18-354.6 Ro\_IV,43-46 {34/73}     saṅkhyā saṅkhyeye vartate . \\
\noindent \textbf{(P\_5,1.57-58.1)} KA\_II,352.18-354.6 Ro\_IV,43-46 {35/73}     yadi tarhi yaḥ yaḥ kāle vartate saḥ saḥ kālaśabdaḥ ramaṇīyādiṣu atriprasaṅgaḥ bhavati . \\
\noindent \textbf{(P\_5,1.57-58.1)} KA\_II,352.18-354.6 Ro\_IV,43-46 {36/73}     ramaṇīyam kālam bhūtaḥ . \\
\noindent \textbf{(P\_5,1.57-58.1)} KA\_II,352.18-354.6 Ro\_IV,43-46 {37/73}     śobhanam kālam bhūtaḥ . \\
\noindent \textbf{(P\_5,1.57-58.1)} KA\_II,352.18-354.6 Ro\_IV,43-46 {38/73}     atha matam etat . \\
\noindent \textbf{(P\_5,1.57-58.1)} KA\_II,352.18-354.6 Ro\_IV,43-46 {39/73}     kāle dṛṣṭaḥ śabdaḥ kālaśabdaḥ kālam yaḥ na vyabhicarati iti na ramaṇīyādiṣu atriprasaṅgaḥ bhavati . \\
\noindent \textbf{(P\_5,1.57-58.1)} KA\_II,352.18-354.6 Ro\_IV,43-46 {40/73}     jīvitaparimāṇe tu upasaṅkhyānam kartavyam . \\
\noindent \textbf{(P\_5,1.57-58.1)} KA\_II,352.18-354.6 Ro\_IV,43-46 {41/73}     iha ca upasaṅkhyānam kartavyam . \\
\noindent \textbf{(P\_5,1.57-58.1)} KA\_II,352.18-354.6 Ro\_IV,43-46 {42/73}     vārṣaśatikaḥ . \\
\noindent \textbf{(P\_5,1.57-58.1)} KA\_II,352.18-354.6 Ro\_IV,43-46 {43/73}     vārṣasahasrikaḥ iti . \\
\noindent \textbf{(P\_5,1.57-58.1)} KA\_II,352.18-354.6 Ro\_IV,43-46 {44/73}     kim punaḥ kāraṇam na sidhyati . \\
\noindent \textbf{(P\_5,1.57-58.1)} KA\_II,352.18-354.6 Ro\_IV,43-46 {45/73}     na hi varṣaśataśabdaḥ saṅkhyā . \\
\noindent \textbf{(P\_5,1.57-58.1)} KA\_II,352.18-354.6 Ro\_IV,43-46 {46/73}     kim tarhi saṅkhyeye vartate varṣaśataśabdaḥ . \\
\noindent \textbf{(P\_5,1.57-58.1)} KA\_II,352.18-354.6 Ro\_IV,43-46 {47/73}     evam tarhi anyebhyaḥ api dṛśyate khāraśatādyartham . \\
\noindent \textbf{(P\_5,1.57-58.1)} KA\_II,352.18-354.6 Ro\_IV,43-46 {48/73}     anyebhyaḥ api dṛśyate iti vaktavyam . \\
\noindent \textbf{(P\_5,1.57-58.1)} KA\_II,352.18-354.6 Ro\_IV,43-46 {49/73}     kim prayojanam . \\
\noindent \textbf{(P\_5,1.57-58.1)} KA\_II,352.18-354.6 Ro\_IV,43-46 {50/73}     khāraśatādyartham . \\
\noindent \textbf{(P\_5,1.57-58.1)} KA\_II,352.18-354.6 Ro\_IV,43-46 {51/73}     khāraśatikaḥ rāśiḥ . \\
\noindent \textbf{(P\_5,1.57-58.1)} KA\_II,352.18-354.6 Ro\_IV,43-46 {52/73}     khārasahastrikaḥ rāśiḥ . \\
\noindent \textbf{(P\_5,1.57-58.1)} KA\_II,352.18-354.6 Ro\_IV,43-46 {53/73}     ayam tarhi doṣaḥ . \\
\noindent \textbf{(P\_5,1.57-58.1)} KA\_II,352.18-354.6 Ro\_IV,43-46 {54/73}     iha vacane hi lukprasaṅgaḥ iti . \\
\noindent \textbf{(P\_5,1.57-58.1)} KA\_II,352.18-354.6 Ro\_IV,43-46 {55/73}     na brūmaḥ yatra kriyamāṇe doṣaḥ tatra kartavyam iti . \\
\noindent \textbf{(P\_5,1.57-58.1)} KA\_II,352.18-354.6 Ro\_IV,43-46 {56/73}     kim tarhi . \\
\noindent \textbf{(P\_5,1.57-58.1)} KA\_II,352.18-354.6 Ro\_IV,43-46 {57/73}     yatra kriyamāṇe na doṣaḥ tatra kartavyam . \\
\noindent \textbf{(P\_5,1.57-58.1)} KA\_II,352.18-354.6 Ro\_IV,43-46 {58/73}     kva ca kriyamāṇe na doṣaḥ . \\
\noindent \textbf{(P\_5,1.57-58.1)} KA\_II,352.18-354.6 Ro\_IV,43-46 {59/73}     param arhāt . \\
\noindent \textbf{(P\_5,1.57-58.1)} KA\_II,352.18-354.6 Ro\_IV,43-46 {60/73}     tat tarhi upasaṅkhyānam kartavyam . \\
\noindent \textbf{(P\_5,1.57-58.1)} KA\_II,352.18-354.6 Ro\_IV,43-46 {61/73}     na kartavyam . \\
\noindent \textbf{(P\_5,1.57-58.1)} KA\_II,352.18-354.6 Ro\_IV,43-46 {62/73}     kālāt iti eva siddham . \\
\noindent \textbf{(P\_5,1.57-58.1)} KA\_II,352.18-354.6 Ro\_IV,43-46 {63/73}     nanu ca uktam na ime kālaśabdāḥ . \\
\noindent \textbf{(P\_5,1.57-58.1)} KA\_II,352.18-354.6 Ro\_IV,43-46 {64/73}     kim tarhi . \\
\noindent \textbf{(P\_5,1.57-58.1)} KA\_II,352.18-354.6 Ro\_IV,43-46 {65/73}     saṅkhyāśabdāḥ iti . \\
\noindent \textbf{(P\_5,1.57-58.1)} KA\_II,352.18-354.6 Ro\_IV,43-46 {66/73}     nanu ca uktam ime api kālaśabdāḥ . \\
\noindent \textbf{(P\_5,1.57-58.1)} KA\_II,352.18-354.6 Ro\_IV,43-46 {67/73}     katham . \\
\noindent \textbf{(P\_5,1.57-58.1)} KA\_II,352.18-354.6 Ro\_IV,43-46 {68/73}     saṅkhyā saṅkhyeye vartate . \\
\noindent \textbf{(P\_5,1.57-58.1)} KA\_II,352.18-354.6 Ro\_IV,43-46 {69/73}     nanu ca uktam yadi tarhi yaḥ yaḥ kāle vartate saḥ saḥ kālaśabdaḥ ramaṇīyādiṣu atriprasaṅgaḥ bhavati iti . \\
\noindent \textbf{(P\_5,1.57-58.1)} KA\_II,352.18-354.6 Ro\_IV,43-46 {70/73}     uktam vā . \\
\noindent \textbf{(P\_5,1.57-58.1)} KA\_II,352.18-354.6 Ro\_IV,43-46 {71/73}     kim uktam . \\
\noindent \textbf{(P\_5,1.57-58.1)} KA\_II,352.18-354.6 Ro\_IV,43-46 {72/73}     anabhidhānāt iti . \\
\noindent \textbf{(P\_5,1.57-58.1)} KA\_II,352.18-354.6 Ro\_IV,43-46 {73/73}     anabhidhānāt ramaṇīyādiṣu utpattiḥ na bhaviṣyati . \\
\noindent \textbf{(P\_5,1.57-58.2)} KA\_II,354.7-9 Ro\_IV,46 {1/6}     stome ḍavidhiḥ pañcadaśādyarthaḥ . \\
\noindent \textbf{(P\_5,1.57-58.2)} KA\_II,354.7-9 Ro\_IV,46 {2/6}     stome ḍaḥ vidheyaḥ . \\
\noindent \textbf{(P\_5,1.57-58.2)} KA\_II,354.7-9 Ro\_IV,46 {3/6}     kim prayojanam . \\
\noindent \textbf{(P\_5,1.57-58.2)} KA\_II,354.7-9 Ro\_IV,46 {4/6}     pañcadaśādyarthaḥ . \\
\noindent \textbf{(P\_5,1.57-58.2)} KA\_II,354.7-9 Ro\_IV,46 {5/6}     pañcadaśaḥ stomaḥ . \\
\noindent \textbf{(P\_5,1.57-58.2)} KA\_II,354.7-9 Ro\_IV,46 {6/6}     saptadaśaḥ stomaḥ iti . \\
\noindent \textbf{(P\_5,1.59)} ime viṃśatyādayaḥ saprakṛtikāḥ sapratyayakāḥ nipātyante . \\
\noindent \textbf{(P\_5,1.59)} tatra na jñāyate kā prakṛtiḥ kaḥ pratyayaḥ kaḥ pratyayārthaḥ iti . \\
\noindent \textbf{(P\_5,1.59)} tatra vaktavyam iyam prakṛtiḥ ayam pratyayaḥ ayam pratyayārthaḥ iti . \\
\noindent \textbf{(P\_5,1.59)} ime brūmaḥ dviśabdāt ayam daśadarthābhidāhinaḥ svārthe śaticpratyayaḥ nipātyate vinbhāvaḥ ca . \\
\noindent \textbf{(P\_5,1.59)} dvau daśatau viṃśatiḥ . \\
\noindent \textbf{(P\_5,1.59)} viṃśatyādayaḥ daśāt cet samāsavacanānupapattiḥ . \\
\noindent \textbf{(P\_5,1.59)} viṃśatyādayaḥ daśāt cet samāsaḥ na upapadyate . \\
\noindent \textbf{(P\_5,1.59)} viṃśatigavam iti . \\
\noindent \textbf{(P\_5,1.59)} kim kāraṇam . \\
\noindent \textbf{(P\_5,1.59)} dravyam anabhihitam . \\
\noindent \textbf{(P\_5,1.59)} tasya anabhihitatvāt ṣaṣṭhī prāpnoti . \\
\noindent \textbf{(P\_5,1.59)} ṣaṣṭhyantam ca samāse pūrvam nipatati . \\
\noindent \textbf{(P\_5,1.59)} tatra goviṃśatiḥ iti prāpnoti . \\
\noindent \textbf{(P\_5,1.59)} na ca evam bhavitavyam . \\
\noindent \textbf{(P\_5,1.59)} bhavitavyam ca viṃśatigavam tu na sidhyati . \\
\noindent \textbf{(P\_5,1.59)} iha ca triṃśatpūlī catvāriṃśatpūlī samānādhikaraṇalakṣaṇaḥ samāsaḥ na prāpnoti . \\
\noindent \textbf{(P\_5,1.59)} vacanam ca vidheyam . \\
\noindent \textbf{(P\_5,1.59)} viṃśatiḥ . \\
\noindent \textbf{(P\_5,1.59)} dvitvāt daśatoḥ dvayoḥ dvivacanam iti dvivacanam prāpnoti . \\
\noindent \textbf{(P\_5,1.59)} evam tarhi parimāṇini viṃśatyādayaḥ bhaviṣyanti . \\
\noindent \textbf{(P\_5,1.59)} parimāṇini cet punaḥ svārthe pratyayavidhānam . \\
\noindent \textbf{(P\_5,1.59)} parimāṇini cet punaḥ svārthe pratyayaḥ vidheyaḥ . \\
\noindent \textbf{(P\_5,1.59)} viṃśakaḥ saṅghaḥ . \\
\noindent \textbf{(P\_5,1.59)} ṣaṣṭhīvacanavidhiḥ ca . \\
\noindent \textbf{(P\_5,1.59)} ṣaṣṭhī ca vidheyā . \\
\noindent \textbf{(P\_5,1.59)} gavām viṃśatiḥ . \\
\noindent \textbf{(P\_5,1.59)} dravyam abhihitam . \\
\noindent \textbf{(P\_5,1.59)} tasya abhihitatvāt ṣaṣṭhī na prāpnoti . \\
\noindent \textbf{(P\_5,1.59)} ekavacanam ca vidheyam . \\
\noindent \textbf{(P\_5,1.59)} viṃśatiḥ gāvaḥ . \\
\noindent \textbf{(P\_5,1.59)} gobhiḥ sāmānādhikaraṇyāt bahuṣu bahuvacanam iti bahuvacanam prāpnoti . \\
\noindent \textbf{(P\_5,1.59)} anārambhaḥ vā prātipadikavijñānāt yathā sahasrādiṣu . \\
\noindent \textbf{(P\_5,1.59)} anārambhaḥ vā punaḥ viṃśatyādīnām nyāyyaḥ . \\
\noindent \textbf{(P\_5,1.59)} katham sidhyati . \\
\noindent \textbf{(P\_5,1.59)} prātipadikavijñānāt . \\
\noindent \textbf{(P\_5,1.59)} katham prātipadikavijñānam . \\
\noindent \textbf{(P\_5,1.59)} viṃśatyādayaḥ avyutpannāni prātipadikāni yathā sahasrādiṣu . \\
\noindent \textbf{(P\_5,1.59)} tat yathā sahasram ayutam arbudam iti na ca anugamaḥ kriyate bhavati ca abhidhānam . \\
\noindent \textbf{(P\_5,1.59)} yathā sahasrādiṣu iti ucyate . \\
\noindent \textbf{(P\_5,1.59)} atha sahasrādiṣu api katham bhavitavyam . \\
\noindent \textbf{(P\_5,1.59)} sahasram gavām . \\
\noindent \textbf{(P\_5,1.59)} sahasram gāvaḥ . \\
\noindent \textbf{(P\_5,1.59)} sahasragavam . \\
\noindent \textbf{(P\_5,1.59)} gosahasram iti . \\
\noindent \textbf{(P\_5,1.59)} yāvatā atra api sandehaḥ na asūyā kartavyā yatra anugamaḥ kriyate . \\
\noindent \textbf{(P\_5,1.59)} nanu ca uktam viṃśatyādayaḥ daśāt cet samāsavacanānupapattiḥ . \\
\noindent \textbf{(P\_5,1.59)} parimāṇini cet punaḥ svārthe pratyayavidhānam . \\
\noindent \textbf{(P\_5,1.59)} ṣaṣṭhīvacanavidhiḥ ca iti . \\
\noindent \textbf{(P\_5,1.59)} na eṣaḥ doṣaḥ . \\
\noindent \textbf{(P\_5,1.59)} samudāye viṃśatyādayaḥ bhaviṣyanti . \\
\noindent \textbf{(P\_5,1.59)} kim vaktavyam etat . \\
\noindent \textbf{(P\_5,1.59)} na hi . \\
\noindent \textbf{(P\_5,1.59)} katham anucyamānam gaṃsyate . \\
\noindent \textbf{(P\_5,1.59)} saṅghaḥ iti vartate . \\
\noindent \textbf{(P\_5,1.59)} saṅghaḥ samūhaḥ samudāyaḥ iti anarthāntaram . \\
\noindent \textbf{(P\_5,1.59)} te ete viṃśatyādayaḥ samudāye santaḥ bhāvavacanāḥ bhavanti bhāvavacanāḥ santaḥ guṇavacanāḥ bhavanti guṇavacanāḥ santaḥ aviśiṣṭāḥ bhavanti anyaiḥ guṇavacanaiḥ . \\
\noindent \textbf{(P\_5,1.59)} anyeṣu ca guṇavacaneṣu kadā cit guṇaḥ guṇiviśeṣakaḥ bhavati . \\
\noindent \textbf{(P\_5,1.59)} tat yathā śuklaḥ paṭaḥ iti . \\
\noindent \textbf{(P\_5,1.59)} kadā cit guṇinā guṇaḥ vyapadiśyate : paṭasya śuklaḥ iti . \\
\noindent \textbf{(P\_5,1.59)} tat yadā tāvat ucyate viṃśatyādayaḥ daśāt cet samāsavacanānupapattiḥ iti sāmānādhikaraṇyam tadā guṇaguṇinoḥ . \\
\noindent \textbf{(P\_5,1.59)} vacanaparihāraḥ tiṣṭhatu tāvat . \\
\noindent \textbf{(P\_5,1.59)} parimāṇini cet punaḥ svārthe pratyayavidhānam iti saṃhanane vṛttaḥ saṃhanane vartiṣyate . \\
\noindent \textbf{(P\_5,1.59)} saṅkhyāsaṃhanane vṛttaḥ dravyasaṃhanane vartiṣyate . \\
\noindent \textbf{(P\_5,1.59)} atha ṣaṣṭhī tadā guṇinā guṇaḥ viśeṣyate . \\
\noindent \textbf{(P\_5,1.59)} vacanaparihāraḥ ubhayoḥ api . \\
\noindent \textbf{(P\_5,1.59)} yadi tarhi ime viṃśatyādayaḥ guṇavacanāḥ syuḥ sadharmabhiḥ anyaiḥ guṇavacanaiḥ bhavitavyam . \\
\noindent \textbf{(P\_5,1.59)} anye ca guṇavacanāḥ dravyasya liṅgasaṅkhye anuvartante . \\
\noindent \textbf{(P\_5,1.59)} tat yathā . \\
\noindent \textbf{(P\_5,1.59)} śuklam vastram . \\
\noindent \textbf{(P\_5,1.59)} śuklā śāṭī . \\
\noindent \textbf{(P\_5,1.59)} śuklaḥ kambalaḥ . \\
\noindent \textbf{(P\_5,1.59)} śuklau kambalau . \\
\noindent \textbf{(P\_5,1.59)} śuklāḥ kambalāḥ iti . \\
\noindent \textbf{(P\_5,1.59)} yat asau dravyam śritaḥ guṇaḥ tasya yat liṅgam vacanam ca tat guṇasya api bhavati . \\
\noindent \textbf{(P\_5,1.59)} viṃśatyādayaḥ punaḥ na anuvartante . \\
\noindent \textbf{(P\_5,1.59)} anye api vai guṇavacanāḥ na avaśyam dravyasya liṅgasaṅkhye anuvartante . \\
\noindent \textbf{(P\_5,1.59)} tat yathā . \\
\noindent \textbf{(P\_5,1.59)} gāvaḥ dhanam . \\
\noindent \textbf{(P\_5,1.59)} putrā apatyam . \\
\noindent \textbf{(P\_5,1.59)} indrāgnī devatā . \\
\noindent \textbf{(P\_5,1.59)} viśvedevāḥ devatā . \\
\noindent \textbf{(P\_5,1.59)} yāvantaḥ te vāśitām anuyanti sarve te dakṣiṇā samṛddhyai iti . \\
\noindent \textbf{(P\_5,1.59)} atha atra ananuvṛttau hetuḥ śakyaḥ vaktum . \\
\noindent \textbf{(P\_5,1.59)} bāḍham śakyaḥ vaktum . \\
\noindent \textbf{(P\_5,1.59)} kāmam tarhi ucyatām . \\
\noindent \textbf{(P\_5,1.59)} iha kadā cit guṇaḥ prādhānyena vivakṣitaḥ bhavati . \\
\noindent \textbf{(P\_5,1.59)} tat yathā : pañca uḍupaśatāni tīrṇāni . \\
\noindent \textbf{(P\_5,1.59)} pañca phalakaśatāni tīrṇāni . \\
\noindent \textbf{(P\_5,1.59)} aśvaiḥ yuddham . \\
\noindent \textbf{(P\_5,1.59)} asibhiḥ yuddham iti . \\
\noindent \textbf{(P\_5,1.59)} na ca asayaḥ yudhyante . \\
\noindent \textbf{(P\_5,1.59)} asiguṇāḥ puruṣāḥ yudhyante guṇaḥ tu khalu prādhānyena vivakṣitaḥ . \\
\noindent \textbf{(P\_5,1.59)} iha tāvat gāvaḥ dhanam iti dhinoteḥ dhanam ekaḥ guṇaḥ . \\
\noindent \textbf{(P\_5,1.59)} saḥ prādhānyena vivakṣitaḥ . \\
\noindent \textbf{(P\_5,1.59)} putrāḥ apatyam iti apatanāt apatyam ekaḥ guṇaḥ . \\
\noindent \textbf{(P\_5,1.59)} saḥ prādhānyena vivakṣitaḥ . \\
\noindent \textbf{(P\_5,1.59)} indrāgnī devatā . \\
\noindent \textbf{(P\_5,1.59)} viśvedevāḥ devatā iti diveḥ aiśvaryakarmaṇaḥ devaḥ . \\
\noindent \textbf{(P\_5,1.59)} tasmāt svārthe tal . \\
\noindent \textbf{(P\_5,1.59)} ekaḥ guṇaḥ . \\
\noindent \textbf{(P\_5,1.59)} saḥ prādhānyena vivakṣitaḥ . \\
\noindent \textbf{(P\_5,1.59)} yāvantaḥ te vāśitām anuyanti sarve te dakṣiṇā samṛddhyā iti dakṣeḥ vṛddhikarmaṇaḥ dakṣiṇā ekaḥ guṇaḥ . \\
\noindent \textbf{(P\_5,1.59)} saḥ prādhānyena vivakṣitaḥ . \\
\noindent \textbf{(P\_5,1.59)} tasya ekatvāt ekavacanam bhaviṣyati . \\
\noindent \textbf{(P\_5,1.59)} viṃśatyādiṣu ca api ekaḥ guṇaḥ . \\
\noindent \textbf{(P\_5,1.59)} saḥ prādhānyena vivakṣitaḥ . \\
\noindent \textbf{(P\_5,1.59)} tasya ekatvāt ekavacanam bhaviṣyati . \\
\noindent \textbf{(P\_5,1.59)} ayam tarhi viṃśatyādiṣu bhāvavacaneṣu doṣaḥ . \\
\noindent \textbf{(P\_5,1.59)} goviṃśatiḥ ānīyatām iti bhāvānayane codite dravyānanam na prāpnoti . \\
\noindent \textbf{(P\_5,1.59)} na eṣaḥ doṣaḥ . \\
\noindent \textbf{(P\_5,1.59)} idam tāvat ayam praṣṭavyaḥ : atha iha gauḥ anubandhyaḥ ajaḥ agnīṣomīyaḥ iti katham ākṛtau coditāyām dravye ārambhaṇalambhanaprokṣaṇaviśasanādīni kriyante iti . \\
\noindent \textbf{(P\_5,1.59)} asambhavāt . \\
\noindent \textbf{(P\_5,1.59)} ākṛtau ārambhaṇādīnām sambhavaḥ na asti iti kṛtvā ākṛtisahacarite dravye ārambhaṇādīni kriyante . \\
\noindent \textbf{(P\_5,1.59)} idam api evañjātīyakam eva . \\
\noindent \textbf{(P\_5,1.59)} asambhavāt bhāvānayanasya dravyānayanam bhaviṣyati . \\
\noindent \textbf{(P\_5,1.59)} atha vā avyatirekāt . \\
\noindent \textbf{(P\_5,1.64,} 76) KA\_II,357.11-20 Ro\_IV,56-58 {1/20}     chedādipathibhyaḥ vigrahadarśanāt nityagrahaṇānarthakyam . \\
\noindent \textbf{(P\_5,1.64,} 76) KA\_II,357.11-20 Ro\_IV,56-58 {2/20}     chedādipathibhyaḥ nityagrahaṇam anarthakam . \\
\noindent \textbf{(P\_5,1.64,} 76) KA\_II,357.11-20 Ro\_IV,56-58 {3/20}     kim kāraṇam . \\
\noindent \textbf{(P\_5,1.64,} 76) KA\_II,357.11-20 Ro\_IV,56-58 {4/20}     vigrahadarśanāt . \\
\noindent \textbf{(P\_5,1.64,} 76) KA\_II,357.11-20 Ro\_IV,56-58 {5/20}     vigrahaḥ dṛśyate . \\
\noindent \textbf{(P\_5,1.64,} 76) KA\_II,357.11-20 Ro\_IV,56-58 {6/20}     chedam arhati . \\
\noindent \textbf{(P\_5,1.64,} 76) KA\_II,357.11-20 Ro\_IV,56-58 {7/20}     panthānam gacchati iti . \\
\noindent \textbf{(P\_5,1.64,} 76) KA\_II,357.11-20 Ro\_IV,56-58 {8/20}     vikārārtham tarhi idam nityagrahaṇam kriyate . \\
\noindent \textbf{(P\_5,1.64,} 76) KA\_II,357.11-20 Ro\_IV,56-58 {9/20}     vikāreṇa vigrahaḥ mā bhūt iti . \\
\noindent \textbf{(P\_5,1.64,} 76) KA\_II,357.11-20 Ro\_IV,56-58 {10/20}     virāgaviraṅgam ca . \\
\noindent \textbf{(P\_5,1.64,} 76) KA\_II,357.11-20 Ro\_IV,56-58 {11/20}     panthaḥ ṇa nityam iti . \\
\noindent \textbf{(P\_5,1.64,} 76) KA\_II,357.11-20 Ro\_IV,56-58 {12/20}     vikārārtham iti cet akaṅādibhiḥ tulyam . \\
\noindent \textbf{(P\_5,1.64,} 76) KA\_II,357.11-20 Ro\_IV,56-58 {13/20}     vikārārtham iti cet akaṅādibhiḥ tulyam etat . \\
\noindent \textbf{(P\_5,1.64,} 76) KA\_II,357.11-20 Ro\_IV,56-58 {14/20}     yathā akaṅādibhiḥ vikāraiḥ vigrahaḥ na bhavati evam ābhyām api na bhaviṣyati . \\
\noindent \textbf{(P\_5,1.64,} 76) KA\_II,357.11-20 Ro\_IV,56-58 {15/20}     kim punaḥ iha akartavyam nityagrahaṇam kriyate āhosvit anyatra kartavyam na kriyate . \\
\noindent \textbf{(P\_5,1.64,} 76) KA\_II,357.11-20 Ro\_IV,56-58 {16/20}     iha akartavyam kriyate . \\
\noindent \textbf{(P\_5,1.64,} 76) KA\_II,357.11-20 Ro\_IV,56-58 {17/20}     eṣaḥ eva nyāyaḥ yat uta sanniyogaśiṣṭānām anyatarāpāye ubhayoḥ api abhāvaḥ . \\
\noindent \textbf{(P\_5,1.64,} 76) KA\_II,357.11-20 Ro\_IV,56-58 {18/20}     tat yathā . \\
\noindent \textbf{(P\_5,1.64,} 76) KA\_II,357.11-20 Ro\_IV,56-58 {19/20}     devadattayajñadattābhyām idam kartavyam iti . \\
\noindent \textbf{(P\_5,1.64,} 76) KA\_II,357.11-20 Ro\_IV,56-58 {20/20}     devadattāpāye yajñadattaḥ api na karoti . \\
\noindent \textbf{(P\_5,1.71)} yajñartvigbhyāmtat karma arhati iti upasaṅkhyānam . yajñartvigbhyāmtat karma arhati iti upasaṅkhyānam kartavyam . \\
\noindent \textbf{(P\_5,1.71)} yjañakarma arhati yajñiyaḥ deśaḥ . \\
\noindent \textbf{(P\_5,1.71)} ṛtvikkarma arhati ārtvijīnam brāhmaṇakulam iti . \\
\noindent \textbf{(P\_5,1.72)} tat vartayati iti anirdeśaḥ tatra adarśanāt . \\
\noindent \textbf{(P\_5,1.72)} tat vartayati iti anirdeśaḥ . \\
\noindent \textbf{(P\_5,1.72)} agamakaḥ nirdeśaḥ anirdeśaḥ . \\
\noindent \textbf{(P\_5,1.72)} pārāyaṇam kaḥ vartayati . \\
\noindent \textbf{(P\_5,1.72)} yaḥ parasya karoti . \\
\noindent \textbf{(P\_5,1.72)} turāyaṇam kaḥ vartayati . \\
\noindent \textbf{(P\_5,1.72)} yaḥ carupuroḍāśān nirvapati . \\
\noindent \textbf{(P\_5,1.72)} tatra adarśanāt . \\
\noindent \textbf{(P\_5,1.72)} na ca tatra pratyayaḥ dṛśyate . \\
\noindent \textbf{(P\_5,1.72)} iṅyajyoḥ ca darśanāt . \\
\noindent \textbf{(P\_5,1.72)} iṅyajyoḥ ca pratyayaḥ dṛśyate . \\
\noindent \textbf{(P\_5,1.72)} yaḥ pārāyaṇam adhīte saḥ pārāyaṇikaḥ iti ucyate . \\
\noindent \textbf{(P\_5,1.72)} yaḥ turāyaṇena yajate saḥ taurāyaṇikaḥ iti ucyate . \\
\noindent \textbf{(P\_5,1.72)} yaḥ ca eva adhīte yaḥ parasya karoti ubhau tau vartayataḥ . \\
\noindent \textbf{(P\_5,1.72)} yaḥ ca yajate yaḥ ca yaḥ carupuroḍāśān nirvapati ubhau tau vartayataḥ . \\
\noindent \textbf{(P\_5,1.72)} ubhayatra kasmāt na bhavati . \\
\noindent \textbf{(P\_5,1.72)} anabhidhānāt . \\
\noindent \textbf{(P\_5,1.74)} yojanam gacchati iti krośaśatayojanaśatayoḥ upasaṅkhyānam . \\
\noindent \textbf{(P\_5,1.74)} yojanam gacchati iti krośaśatayojanaśatayoḥ upasaṅkhyānam kartavyam . \\
\noindent \textbf{(P\_5,1.74)} krośaśatam gacchati iti krauśaśatikaḥ . \\
\noindent \textbf{(P\_5,1.74)} yojanaśatam gacchati iti yaujanaśatikaḥ iti . \\
\noindent \textbf{(P\_5,1.74)} tataḥ abhigamanam arhati iti ca . \\
\noindent \textbf{(P\_5,1.74)} tataḥ abhigamanam arhati iti ca krośaśatayojanaśatayoḥ upasaṅkhyānam kartavyam . \\
\noindent \textbf{(P\_5,1.74)} krośaśatāt abhigamanam arhati krauśaśatikaḥ bhikṣuḥ . \\
\noindent \textbf{(P\_5,1.74)} yojanaśatāt abhigamanam arhati yaujanaśatikaḥ guruḥ . \\
\noindent \textbf{(P\_5,1.77)} āhṛtaprakaraṇe vārijaṅgalasthalakāntārapūrvapadāt upasaṅkhyānam . \\
\noindent \textbf{(P\_5,1.77)} āhṛtaprakaraṇe vārijaṅgalasthalakāntārapūrvapadāt upasaṅkhyānam kartavyam . \\
\noindent \textbf{(P\_5,1.77)} vāripathena gacchati vāripathikaḥ . \\
\noindent \textbf{(P\_5,1.77)} vāripathena āhṛtam vāripathikam . \\
\noindent \textbf{(P\_5,1.77)} vāri . \\
\noindent \textbf{(P\_5,1.77)} jaṅgala . \\
\noindent \textbf{(P\_5,1.77)} jaṅgalapathena gacchati jāṅgalapathikaḥ . \\
\noindent \textbf{(P\_5,1.77)} jaṅgalapathena āhṛtam jāṅgalapathikam . \\
\noindent \textbf{(P\_5,1.77)} jaṅgala . \\
\noindent \textbf{(P\_5,1.77)} sthala . \\
\noindent \textbf{(P\_5,1.77)} sthalapathena gacchati sthālapathikaḥ . \\
\noindent \textbf{(P\_5,1.77)} sthalapathena āhṛtam sthālapathikam . \\
\noindent \textbf{(P\_5,1.77)} sthala . \\
\noindent \textbf{(P\_5,1.77)} kāntāra . \\
\noindent \textbf{(P\_5,1.77)} kāntārapathena gacchati kāntārapathikaḥ . \\
\noindent \textbf{(P\_5,1.77)} kāntārapathena āhṛtam kāntārapathikam . \\
\noindent \textbf{(P\_5,1.77)} ajapathaśaṅkupathābhyām ca . \\
\noindent \textbf{(P\_5,1.77)} ajapathaśaṅkupathābhyām ca iti vaktavyam . \\
\noindent \textbf{(P\_5,1.77)} ajapathena gacchati ājapathikaḥ . \\
\noindent \textbf{(P\_5,1.77)} ajapathena āhṛtam ājapathikam. śaṅkupathena gacchati śāṅkupathikaḥ . \\
\noindent \textbf{(P\_5,1.77)} śaṅkupathena āhṛtam śāṅkupathikam . \\
\noindent \textbf{(P\_5,1.77)} madhukamaricayoḥ aṇ sthalāt . \\
\noindent \textbf{(P\_5,1.77)} madhukamaricayoḥ aṇ sthalāt vaktavyaḥ . \\
\noindent \textbf{(P\_5,1.77)} sthālapatham madhukam . \\
\noindent \textbf{(P\_5,1.77)} sthālapatham maricam . \\
\noindent \textbf{(P\_5,1.80)} adhīṣṭabhṛtayoḥ dvitīyānirdeśaḥ anarthakaḥ tatra adarśanāt . \\
\noindent \textbf{(P\_5,1.80)} adhīṣṭabhṛtayoḥ dvitīyānirdeśaḥ anarthakaḥ . \\
\noindent \textbf{(P\_5,1.80)} kim kāraṇam . \\
\noindent \textbf{(P\_5,1.80)} tatra adarśanāt . \\
\noindent \textbf{(P\_5,1.80)} na hi asau māsam adhīṣyate . \\
\noindent \textbf{(P\_5,1.80)} kim tarhi muhūrtam adhīṣṭaḥ māsam tat karma karoti . \\
\noindent \textbf{(P\_5,1.80)} siddham tu caturthīnirdeśāt . \\
\noindent \textbf{(P\_5,1.80)} siddham etat . \\
\noindent \textbf{(P\_5,1.80)} katham . \\
\noindent \textbf{(P\_5,1.80)} caturthīnirdeśāt . \\
\noindent \textbf{(P\_5,1.80)} caturthīnirdeśaḥ kartavyaḥ . \\
\noindent \textbf{(P\_5,1.80)} tasmai adhīṣṭaḥ iti . \\
\noindent \textbf{(P\_5,1.80)} saḥ tarhi caturthīnirdeśaḥ kartavyaḥ . \\
\noindent \textbf{(P\_5,1.80)} na kartavyaḥ . \\
\noindent \textbf{(P\_5,1.80)} tādarthyāt tācchabdyam bhaviṣyati . \\
\noindent \textbf{(P\_5,1.80)} māsārthaḥ muhūrtaḥ māsaḥ . \\
\noindent \textbf{(P\_5,1.84)} avayasi ṭhan ca iti anantarasya anukarṣaḥ . \\
\noindent \textbf{(P\_5,1.84)} avayasi ṭhan ca iti anantarasya anukarṣaḥ draṣṭavyaḥ . \\
\noindent \textbf{(P\_5,1.84)} dveṣyam vijānīyāt : yap api anuvartate iti . \\
\noindent \textbf{(P\_5,1.84)} tat ācāryaḥ suhṛt bhūtvā anvācaṣṭe : avayasi ṭhan ca iti anantarasya anukarṣaḥ iti . \\
\noindent \textbf{(P\_5,1.90)} ṣaṣṭike sañjñāgrahaṇam . \\
\noindent \textbf{(P\_5,1.90)} ṣaṣṭike sañjñāgrahaṇam kartavyam . \\
\noindent \textbf{(P\_5,1.90)} mudgāḥ api hi ṣaṣṭirātreṇe pacyante . \\
\noindent \textbf{(P\_5,1.90)} tatra mā bhūt iti . \\
\noindent \textbf{(P\_5,1.90)} uktam vā . \\
\noindent \textbf{(P\_5,1.90)} kim uktam . \\
\noindent \textbf{(P\_5,1.90)} anabhidhānāt iti . \\
\noindent \textbf{(P\_5,1.94)} tat asya brahmacaryam iti mahānāmnyādibhyaḥ upasaṅkhyānam . \\
\noindent \textbf{(P\_5,1.94)} tat asya brahmacaryam iti mahānāmnyādibhyaḥ upasaṅkhyānam kartavyam . \\
\noindent \textbf{(P\_5,1.94)} mahānāmnīnām brahmacaryam māhānāmnikam . \\
\noindent \textbf{(P\_5,1.94)} ādityavratikam . \\
\noindent \textbf{(P\_5,1.94)} tat carati iti ca . \\
\noindent \textbf{(P\_5,1.94)} tat carati iti ca mahānāmnyādibhyaḥ upasaṅkhyānam kartavyam . \\
\noindent \textbf{(P\_5,1.94)} mahānāmnīḥ carati māhānāmnikaḥ . \\
\noindent \textbf{(P\_5,1.94)} ādityavratikaḥ . \\
\noindent \textbf{(P\_5,1.94)} na eṣaḥ yuktaḥ nirdeśaḥ tat carati iti . \\
\noindent \textbf{(P\_5,1.94)} mahānāmnyaḥ nāma ṛcaḥ . \\
\noindent \textbf{(P\_5,1.94)} na ca tāḥ caryante . \\
\noindent \textbf{(P\_5,1.94)} vratam tāsām caryate . \\
\noindent \textbf{(P\_5,1.94)} na eṣaḥ doṣaḥ . \\
\noindent \textbf{(P\_5,1.94)} sāhacaryāt tācchabyam bhaviṣyati . \\
\noindent \textbf{(P\_5,1.94)} mahānāmnīsahacaritam vratam mahānāmnyaḥ vratam iti . \\
\noindent \textbf{(P\_5,1.94)} avāntaradīkṣādibhyaḥ ḍiniḥ . \\
\noindent \textbf{(P\_5,1.94)} avāntaradīkṣādibhyaḥ ḍiniḥ vaktavyaḥ . \\
\noindent \textbf{(P\_5,1.94)} avāntaradīkṣī . \\
\noindent \textbf{(P\_5,1.94)} tilavratī . \\
\noindent \textbf{(P\_5,1.94)} aṣṭācatvāriṃśataḥ ḍvun ca . \\
\noindent \textbf{(P\_5,1.94)} aṣṭācatvāriṃśataḥ ḍvun ca ḍiniḥ ca vaktavyaḥ . \\
\noindent \textbf{(P\_5,1.94)} aṣṭācatvāriṃśakaḥ . \\
\noindent \textbf{(P\_5,1.94)} aṣṭācatvāriṃśī . \\
\noindent \textbf{(P\_5,1.94)} cāturmāsyānām yalopaḥ ca . \\
\noindent \textbf{(P\_5,1.94)} cāturmāsyānām yalopaḥ ca ḍvun ca ḍiniḥ ca vaktavyaḥ . \\
\noindent \textbf{(P\_5,1.94)} cāturmāsikaḥ . \\
\noindent \textbf{(P\_5,1.94)} cāturmāsī . \\
\noindent \textbf{(P\_5,1.94)} atha kim idam cāturmāsyānām iti . \\
\noindent \textbf{(P\_5,1.94)} caturmāsāt ṇyaḥ yajñe tatra bhave . \\
\noindent \textbf{(P\_5,1.94)} caturmāsāt ṇyaḥ vaktavyaḥ yajñe tatra bhave iti etasmin arthe . \\
\noindent \textbf{(P\_5,1.94)} caturṣu māseṣu bhavāni cāturmāsyāni yajñāḥ . \\
\noindent \textbf{(P\_5,1.94)} sañjñāyām aṇ . sañjñāyām aṇ vaktavyaḥ . \\
\noindent \textbf{(P\_5,1.94)} caturṣu māseṣu bhavā cāturmāsī pauṇamāsī \\
\noindent \textbf{(P\_5,1.95)} ākhyāgrahaṇam kimartham. tasya dakṣiṇā yajñebhyaḥ iti iyati ucyamāne ye ete sañjñībhūtakāḥ yajñāḥ tataḥ utapattiḥ syāt . \\
\noindent \textbf{(P\_5,1.95)} agniṣṭomikyaḥ . \\
\noindent \textbf{(P\_5,1.95)} rājasūyikaḥ . \\
\noindent \textbf{(P\_5,1.95)} vājapeyikyaḥ . \\
\noindent \textbf{(P\_5,1.95)} yatra vā yajñaśabdaḥ asti . \\
\noindent \textbf{(P\_5,1.95)} nāvayajñikyaḥ . \\
\noindent \textbf{(P\_5,1.95)} pākayajñikyaḥ . \\
\noindent \textbf{(P\_5,1.95)} iha na syāt . \\
\noindent \textbf{(P\_5,1.95)} pāñcaudanikyaḥ . \\
\noindent \textbf{(P\_5,1.95)} dāśaudanikyaḥ . \\
\noindent \textbf{(P\_5,1.95)} ākhyāgrahaṇe punaḥ kriyamāṇe na doṣaḥ bhavati . \\
\noindent \textbf{(P\_5,1.95)} ye ca sañjñībhūtakāḥ yatra ca yajñaśabdaḥ asti yatra ca na asti tadākhyāmātrāt siccham bhavati . \\
\noindent \textbf{(P\_5,1.96)} kāryagrahaṇam anarthakam tatrabhavena kṛtatvāt . kāryagrahaṇam anarthakam . \\
\noindent \textbf{(P\_5,1.96)} kim kāraṇam . \\
\noindent \textbf{(P\_5,1.96)} tatrabhavena kṛtatvāt . \\
\noindent \textbf{(P\_5,1.96)} yat hi māse kāryam māse bhavam tat bhavati . \\
\noindent \textbf{(P\_5,1.96)} tatra tatra bhavaḥ iti eva siddham . \\
\noindent \textbf{(P\_5,1.96)} kim idam bhavān kāryagrahaṇam eva pratyācaṣṭe na punaḥ dīyategrahaṇam api . \\
\noindent \textbf{(P\_5,1.96)} yathā eva hi yat māse kāryam tat māse bhavam bhavati evam yat api māse dīyate tat api māse bhavam bhavati . \\
\noindent \textbf{(P\_5,1.96)} tatra tatra bhavaḥ iti eva siddham . \\
\noindent \textbf{(P\_5,1.96)} na sidhyati . \\
\noindent \textbf{(P\_5,1.96)} na tat māse dīyate . \\
\noindent \textbf{(P\_5,1.96)} kim tarhi māse gate . \\
\noindent \textbf{(P\_5,1.96)} evam tarhi aupaśleṣikam adhikaraṇam vijñāsyate . \\
\noindent \textbf{(P\_5,1.96)} evam tarhi yogavibhāgottarakālam idam paṭhitavyam . \\
\noindent \textbf{(P\_5,1.96)} tasya dakṣiṇā yajñākhyebhyaḥ . \\
\noindent \textbf{(P\_5,1.96)} tatra ca dīyate . \\
\noindent \textbf{(P\_5,1.96)} tataḥ kāryam bhavavat kālāt iti . \\
\noindent \textbf{(P\_5,1.97)} aṇprakaraṇe agnipadādibhyaḥ upasaṅkhyānam . \\
\noindent \textbf{(P\_5,1.97)} aṇprakaraṇe agnipadādibhyaḥ upasaṅkhyānam kartavyam . \\
\noindent \textbf{(P\_5,1.97)} trīṇi imāni aṇgrahaṇāni . \\
\noindent \textbf{(P\_5,1.97)} vyuṣṭādibhyaḥ aṇ . \\
\noindent \textbf{(P\_5,1.97)} samayaḥ tat asya prāptam . \\
\noindent \textbf{(P\_5,1.97)} ṛtoḥ aṇ . \\
\noindent \textbf{(P\_5,1.97)} prayojanam . \\
\noindent \textbf{(P\_5,1.97)} viśākhāṣāḍhāt aṇ manthadaṇḍayoḥ iti . \\
\noindent \textbf{(P\_5,1.97)} tatra na jñāyate katarasmin aṇprakaraṇe agnipadādibhyaḥ upasaṅkhyānam . \\
\noindent \textbf{(P\_5,1.97)} aviśeṣāt sarvatra . \\
\noindent \textbf{(P\_5,1.97)} vyuṣṭādibhyaḥ aṇ bhavati iti uktvā agnipadādibhyaḥ ca iti vaktavyam . \\
\noindent \textbf{(P\_5,1.97)} agnipade dīyate kāryam vā āgnipadam . \\
\noindent \textbf{(P\_5,1.97)} pailumūlam . \\
\noindent \textbf{(P\_5,1.97)} samayaḥ tat asya prāptam . \\
\noindent \textbf{(P\_5,1.97)} ṛtoḥ aṇ . \\
\noindent \textbf{(P\_5,1.97)} agnipadādibhyaḥ ca iti vaktavyam . \\
\noindent \textbf{(P\_5,1.97)} upavastā prāptaḥ asya aupavastram . \\
\noindent \textbf{(P\_5,1.97)} prāśitā prāptaḥ asya prāśitram . \\
\noindent \textbf{(P\_5,1.97)} prayojanam . \\
\noindent \textbf{(P\_5,1.97)} viśākhāṣāḍhāt aṇ manthadaṇḍayoḥ .agnipadādibhyaḥ ca iti vaktavyam . \\
\noindent \textbf{(P\_5,1.97)} cūḍā prayojanam asya cauḍam . \\
\noindent \textbf{(P\_5,1.97)} śraddhā prayojanam asya śrāddham . \\
\noindent \textbf{(P\_5,1.111)} chaprakaraṇe viśipūripadiruhiprakṛteḥ anāt sapūrvapadāt upasaṅkhyānam . \\
\noindent \textbf{(P\_5,1.111)} chaprakaraṇe viśipūripadiruhiprakṛteḥ anāt sapūrrvapadāt upasaṅkhyānam kartavyam . \\
\noindent \textbf{(P\_5,1.111)} viśi . \\
\noindent \textbf{(P\_5,1.111)} gehānupraveśanīyam . \\
\noindent \textbf{(P\_5,1.111)} pūri . \\
\noindent \textbf{(P\_5,1.111)} prapāpūraṇīyam . \\
\noindent \textbf{(P\_5,1.111)} padi . \\
\noindent \textbf{(P\_5,1.111)} goprapadanīyam . \\
\noindent \textbf{(P\_5,1.111)} aśvaprapadanīyam . \\
\noindent \textbf{(P\_5,1.111)} ruhi . \\
\noindent \textbf{(P\_5,1.111)} prāśādārohaṇīyam . \\
\noindent \textbf{(P\_5,1.111)} svargādibhyaḥ yat ṣvargādibhyaḥ yat pratyayaḥ bhavati . \\
\noindent \textbf{(P\_5,1.111)} svargyam . \\
\noindent \textbf{(P\_5,1.111)} dhanyam . \\
\noindent \textbf{(P\_5,1.111)} yaśasyam . \\
\noindent \textbf{(P\_5,1.111)} āyuṣyam . \\
\noindent \textbf{(P\_5,1.111)} puṇyāhavācanādibhyaḥ luk . \\
\noindent \textbf{(P\_5,1.111)} puṇyāhavācanādibhyaḥ luk vaktavyaḥ . \\
\noindent \textbf{(P\_5,1.111)} puṇyāhavācanam . \\
\noindent \textbf{(P\_5,1.111)} śāntivācanam . \\
\noindent \textbf{(P\_5,1.111)} svastivācanam . \\
\noindent \textbf{(P\_5,1.113)} ekāgārāt nipātanānarthakyam ṭhañprakaraṇāt . \\
\noindent \textbf{(P\_5,1.113)} ekāgārāt nipātanam anarthakam . \\
\noindent \textbf{(P\_5,1.113)} kim kāraṇam . \\
\noindent \textbf{(P\_5,1.113)} ṭhañprakaraṇāt . \\
\noindent \textbf{(P\_5,1.113)} ṭhañ prakṛtaḥ . \\
\noindent \textbf{(P\_5,1.113)} saḥ anuvartiṣyate . \\
\noindent \textbf{(P\_5,1.113)} idam tarhi prayojanam . \\
\noindent \textbf{(P\_5,1.113)} caure iti vakṣyāmi iti . \\
\noindent \textbf{(P\_5,1.113)} iha mā bhūt . \\
\noindent \textbf{(P\_5,1.113)} ekāgāram prayojanam asya bhikṣoḥ iti . \\
\noindent \textbf{(P\_5,1.113)} yadi etāvat prayojanam syāt ekāgārāt caure iti eva brūyāt . \\
\noindent \textbf{(P\_5,1.114)} ākālāt nipātanānarthakyam ṭhañprakaraṇāt . \\
\noindent \textbf{(P\_5,1.114)} ākālāt nipātanam narthakam . \\
\noindent \textbf{(P\_5,1.114)} kim kāraṇam . \\
\noindent \textbf{(P\_5,1.114)} ṭhañprakaraṇāt . \\
\noindent \textbf{(P\_5,1.114)} ṭhañ prakṛtaḥ . \\
\noindent \textbf{(P\_5,1.114)} saḥ anuvartiṣyate . \\
\noindent \textbf{(P\_5,1.114)} idam tarhi prayojanam . \\
\noindent \textbf{(P\_5,1.114)} etasmin viśeṣe nipātanam kariṣyāmi samānakālasya ādyantavivakṣāyām iti . \\
\noindent \textbf{(P\_5,1.114)} ākālāt ṭhan ca . \\
\noindent \textbf{(P\_5,1.114)} ākālāt ṭhan ca vaktavyaḥ . \\
\noindent \textbf{(P\_5,1.114)} ākālikī . \\
\noindent \textbf{(P\_5,1.114)} ākālikā . \\
\noindent \textbf{(P\_5,1.115)} idam ayuktam vartate . \\
\noindent \textbf{(P\_5,1.115)} kim atra ayuktam . \\
\noindent \textbf{(P\_5,1.115)} yat tat tṛtīyāsamartham kriyā cet sā bhavati iti ucyate . \\
\noindent \textbf{(P\_5,1.115)} katham ca tṛtīyāsamartham nāma kriyā syāt . \\
\noindent \textbf{(P\_5,1.115)} na eṣaḥ doṣaḥ . \\
\noindent \textbf{(P\_5,1.115)} sarve ete śabdāḥ guṇasamudāyeṣu vartante . \\
\noindent \textbf{(P\_5,1.115)} brāhmaṇaḥ kṣatriyaḥ . \\
\noindent \textbf{(P\_5,1.115)} vaiśyaḥ śūdraḥ iti . \\
\noindent \textbf{(P\_5,1.115)} ātaḥ ca guṇasamudāye evam hi āha . \\
\noindent \textbf{(P\_5,1.115)} tapaḥ śrutam ca yoniḥ ca iti etad brāhmaṇakārakam . \\
\noindent \textbf{(P\_5,1.115)} tapaḥśtrutābhyām yaḥ hīnaḥ jātibrāhmaṇaḥ eva saḥ . \\
\noindent \textbf{(P\_5,1.115)} tathā gauraḥ śucyācāraḥ piṅgalaḥ kapilakeśaḥ iti etān api abhyantarān brāhmaṇe guṇan kurvanti . \\
\noindent \textbf{(P\_5,1.115)} samudāyeṣu ca śabdāḥ vṛttāḥ avayaveṣu api vartante . \\
\noindent \textbf{(P\_5,1.115)} tat yathā : pūrve pañcālāḥ , uttare pañcālāḥ , tailam bhuktam , ghṛtam bhuktam , śuklaḥ , nīlaḥ , kṛṣṇaḥ iti . \\
\noindent \textbf{(P\_5,1.115)} evam ayam brāhmaṇaśabdaya samudāye vṛttaḥ avayaveṣu api vartate . \\
\noindent \textbf{(P\_5,1.115)} yadi tarhi tṛtīyāsamartham viśeṣyate pratyayārthaḥ aviśeṣitaḥ bhavati . \\
\noindent \textbf{(P\_5,1.115)} tatra kaḥ doṣaḥ . \\
\noindent \textbf{(P\_5,1.115)} tṛtīyāsamarthāt kriyāvācinaḥ guṇatulye api pratyayaḥ syāt . \\
\noindent \textbf{(P\_5,1.115)} putreṇa tulyaḥ sthūlaḥ . \\
\noindent \textbf{(P\_5,1.115)} putreṇa tulyaḥ piṅgalaḥ . \\
\noindent \textbf{(P\_5,1.115)} astu tarhi pratyayārthaviśeṣaṇam . \\
\noindent \textbf{(P\_5,1.115)} yat tat tulyam kriyā cet sā bhavati iti . \\
\noindent \textbf{(P\_5,1.115)} evam api tṛtīyāsamartham aviśeṣitam bhavati . \\
\noindent \textbf{(P\_5,1.115)} tatra kaḥ doṣaḥ . \\
\noindent \textbf{(P\_5,1.115)} tṛtīyāsamarthāt akriyāvācinaḥ kriyātulye api pratyayaḥ prāpnoti . \\
\noindent \textbf{(P\_5,1.115)} na eṣaḥ doṣaḥ . \\
\noindent \textbf{(P\_5,1.115)} yat tat tulyam kriyā cet sā bhavati iti ucyate . \\
\noindent \textbf{(P\_5,1.115)} tulayā ca sammitam tulyam . \\
\noindent \textbf{(P\_5,1.115)} yadi ca tṛtīyāsamartham api kriyā pratyayārthaḥ api kriyā tataḥ tulayam bhavati . \\
\noindent \textbf{(P\_5,1.115)} atha vā punaḥ astu yat tat tṛtīyāsamartham kriyā cet sā bhavati iti eva . \\
\noindent \textbf{(P\_5,1.115)} nanu ca uktam pratyayārthaḥ aviśeṣitaḥ iti . \\
\noindent \textbf{(P\_5,1.115)} tatra kaḥ doṣaḥ . \\
\noindent \textbf{(P\_5,1.115)} tṛtīyāsamarthāt kriyāvācinaḥ guṇatulye api pratyayaḥ syāt . \\
\noindent \textbf{(P\_5,1.115)} putreṇa tulyaḥ sthūlaḥ . \\
\noindent \textbf{(P\_5,1.115)} putreṇa tulyaḥ piṅgalaḥ iti . \\
\noindent \textbf{(P\_5,1.115)} na eṣaḥ doṣaḥ . \\
\noindent \textbf{(P\_5,1.115)} yat tat tṛtīyāsamartham kriyā cet sā bhavati iti ucyate . \\
\noindent \textbf{(P\_5,1.115)} tulayā ca sammitam tulyam . \\
\noindent \textbf{(P\_5,1.115)} yadi ca tṛtīyāsamartham api kriyā pratyayārthaḥ api kriyā tataḥ tulayam bhavati . \\
\noindent \textbf{(P\_5,1.115)} kim punaḥ atra jyāyaḥ . \\
\noindent \textbf{(P\_5,1.115)} pratyayārthaviśeṣaṇam eva jyāyaḥ . \\
\noindent \textbf{(P\_5,1.115)} kutaḥ etat . \\
\noindent \textbf{(P\_5,1.115)} evam ca eva kṛtvā ācāryeṇa sūtram paṭhitam . \\
\noindent \textbf{(P\_5,1.115)} vatinā sāmānādhikaraṇyam kṛtam . \\
\noindent \textbf{(P\_5,1.115)} api ca vateḥ avyayeṣu pāṭhaḥ na kartavyaḥ bhavati . \\
\noindent \textbf{(P\_5,1.115)} kriyāyām ayam bhavan liṅgasaṅkhyābhyam na yokṣyate . \\
\noindent \textbf{(P\_5,1.116)} kimartham idam ucyate na tena tulyam kriyā cet vatiḥ iti eva siddham . \\
\noindent \textbf{(P\_5,1.116)} na sidhyati . \\
\noindent \textbf{(P\_5,1.116)} tṛtīyāsamarthāt tatra pratyayaḥ yadā anyena kartavyām kriyām anyaḥ karoti tadā pratyayaḥ utpādyate . \\
\noindent \textbf{(P\_5,1.116)} na ca kā cid ivaśabdena yoge tṛtīyā vidhīyate . \\
\noindent \textbf{(P\_5,1.116)} nanu ca sapatamī api na vidhīyate . \\
\noindent \textbf{(P\_5,1.116)} evam tarhi siddhe sati yat ivaśabdena yoge saptamīsamarthāt vatim śāsti tat jñāpayati ācāryaḥ bhavati ivaśabdena yoge saptamī iti . \\
\noindent \textbf{(P\_5,1.116)} kim etasya jñāpane prayojanam . \\
\noindent \textbf{(P\_5,1.116)} deveṣu iva nāma . \\
\noindent \textbf{(P\_5,1.116)} brāhmaṇeṣu iva nāma . \\
\noindent \textbf{(P\_5,1.116)} eṣaḥ prayogaḥ upapannaḥ bhavati . \\
\noindent \textbf{(P\_5,1.117)} kimartham idam ucyate na tena tulyam kriyā cet vatiḥ iti eva siddham . \\
\noindent \textbf{(P\_5,1.117)} na sidhyati . \\
\noindent \textbf{(P\_5,1.117)} tṛtīyāsamarthāt tatra pratyayaḥ yadā anyena kartavyām kriyām anyaḥ karoti tadā pratyayaḥ utpādyate . \\
\noindent \textbf{(P\_5,1.117)} iha punaḥ dvitīyāsamarthāt ātmārhāyām kriyāyām arhatikartari niścitabalādhāne pratyayaḥ utpādyate . \\
\noindent \textbf{(P\_5,1.117)} brāhmaṇavat bhavān vartate . \\
\noindent \textbf{(P\_5,1.117)} etat vṛttam brāhmaṇaḥ arhati iti . \\
\noindent \textbf{(P\_5,1.118.1)} arthagrahaṇam kimartham . \\
\noindent \textbf{(P\_5,1.118.1)} na upasargāt chandasi dhātavu iti eva ucyeta . \\
\noindent \textbf{(P\_5,1.118.1)} dhātuḥ vai śabdaḥ . \\
\noindent \textbf{(P\_5,1.118.1)} śabde kāryasya asambhavāt arthe kāryam vijñāsyate . \\
\noindent \textbf{(P\_5,1.118.1)} kaḥ punaḥ dhātvarthaḥ . \\
\noindent \textbf{(P\_5,1.118.1)} kriyā . \\
\noindent \textbf{(P\_5,1.118.1)} idam tarhi prayojanam . \\
\noindent \textbf{(P\_5,1.118.1)} uttarapadalopaḥ yathā vijñāyeta . \\
\noindent \textbf{(P\_5,1.118.1)} dhātukṛtaḥ arthaḥ dhātvarthaḥ iti . \\
\noindent \textbf{(P\_5,1.118.1)} kaḥ punaḥ dhātukṛtaḥ arthaḥ . \\
\noindent \textbf{(P\_5,1.118.1)} sādhanam . \\
\noindent \textbf{(P\_5,1.118.1)} kim prayojanam . \\
\noindent \textbf{(P\_5,1.118.1)} sādhane ayam bhavan liṅgasaṅkhyābhyam yokṣyate . \\
\noindent \textbf{(P\_5,1.118.1)} udgatāni udvataḥ . \\
\noindent \textbf{(P\_5,1.118.1)} nigatāni nivataḥ iti . \\
\noindent \textbf{(P\_5,1.118.2)} strīpuṃsābhyām vatyupasaṅkhyānam . \\
\noindent \textbf{(P\_5,1.118.2)} strīpuṃsābhyām vatyupasaṅkhyānam kartavyam . \\
\noindent \textbf{(P\_5,1.118.2)} strīvat . \\
\noindent \textbf{(P\_5,1.118.2)} puṃvat iti . \\
\noindent \textbf{(P\_5,1.118.2)} kim punaḥ kāraṇam na sidhyati . \\
\noindent \textbf{(P\_5,1.118.2)} imau nañsnañau prāk bhavanāt iti ucyete . \\
\noindent \textbf{(P\_5,1.118.2)} tau viśeṣavihitau sāmānyavihitam vatim bādheyātām . \\
\noindent \textbf{(P\_5,1.118.2)} na eṣaḥ doṣaḥ . \\
\noindent \textbf{(P\_5,1.118.2)} ācāryapravṛttiḥ jñāpayati na vatyarthe nañsnañau bhavataḥ iti yat ayam striyāḥ puṃvat iti nirdeśam karoti . \\
\noindent \textbf{(P\_5,1.118.2)} evam api strīvat iti na sidhyati . \\
\noindent \textbf{(P\_5,1.118.2)} yopāpekṣam jñāpakam . \\
\noindent \textbf{(P\_5,1.119.1)} strīpuṃsābhyām tvataloḥ upasaṅkhyānam . \\
\noindent \textbf{(P\_5,1.119.1)} strīpuṃsābhyām tvataloḥ upasaṅkhyānam kartavyam . \\
\noindent \textbf{(P\_5,1.119.1)} strībhāvaḥ strītvam strītā . \\
\noindent \textbf{(P\_5,1.119.1)} kim punaḥ kāraṇam na sidhyati . \\
\noindent \textbf{(P\_5,1.119.1)} imau nañsnañau prāk bhavanāt iti ucyete . \\
\noindent \textbf{(P\_5,1.119.1)} tau viśeṣavihitau sāmānyavihitam vatim bādheyātām . \\
\noindent \textbf{(P\_5,1.119.1)} vāvacanam ca . \\
\noindent \textbf{(P\_5,1.119.1)} vāvacanam ca kartavyam . \\
\noindent \textbf{(P\_5,1.119.1)} kim prayojanam . \\
\noindent \textbf{(P\_5,1.119.1)} nañsnañau api yathā syātām . \\
\noindent \textbf{(P\_5,1.119.1)} strībhāvaḥ straiṇam . \\
\noindent \textbf{(P\_5,1.119.1)} pummbhāvaḥ pauṃsnam iti . \\
\noindent \textbf{(P\_5,1.119.1)} apavādasamāveśāt vā siddham . \\
\noindent \textbf{(P\_5,1.119.1)} apavādasamāveśāt vā siddham etat . \\
\noindent \textbf{(P\_5,1.119.1)} tat yathā imanicprabhṛtibhiḥ apavādaiḥ samāveśaḥ bhavati evam ābhyām api bhaviṣyati . \\
\noindent \textbf{(P\_5,1.119.1)} na eva īśvaraḥ ājñāpapayati na api dharmasūtrakārāḥ paṭhanti imanicprabhṛtibhiḥ apavādaiḥ samāveśaḥ bhavati iti . \\
\noindent \textbf{(P\_5,1.119.1)} kim tarhi ā ca tvāt iti etasmāt yatnāt imanicprabhṛtibhiḥ apavādaiḥ samāveśaḥ bhavati . \\
\noindent \textbf{(P\_5,1.119.1)} na ca etau atra abhyantarau . \\
\noindent \textbf{(P\_5,1.119.1)} etau api atra abhyantarau . \\
\noindent \textbf{(P\_5,1.119.1)} katham . \\
\noindent \textbf{(P\_5,1.119.1)} apavādasadeśāḥ apavādāḥ bhavanti iti . \\
\noindent \textbf{(P\_5,1.119.1)} etat ca eva na jānīmaḥ apavādasadeśāḥ apavādāḥ bhavanti iti . \\
\noindent \textbf{(P\_5,1.119.1)} api ca kutaḥ etat etau api atra abhyantarau na punaḥ pūrvau vā syātām parau vā . \\
\noindent \textbf{(P\_5,1.119.1)} evam tarhi vakṣyati ā ca tvāt iti atra cakārakaraṇasya prayojanam . \\
\noindent \textbf{(P\_5,1.119.1)} nañsnañbhyām api samāveśaḥ bhavati iti . \\
\noindent \textbf{(P\_5,1.119.2)} tasya bhāvaḥ iti abhiprāyādiṣu atiprasaṅgaḥ . \\
\noindent \textbf{(P\_5,1.119.2)} tasya bhāvaḥ iti abhiprāyādiṣu atiprasaṅgaḥ bhavati . \\
\noindent \textbf{(P\_5,1.119.2)} iha api prāpnoti . \\
\noindent \textbf{(P\_5,1.119.2)} abhiprāyaḥ devadattasya modakeṣu bhojane . \\
\noindent \textbf{(P\_5,1.119.2)} ye naḥ bhāvāḥ te naḥ bhāvāḥ putrāḥ putraiḥ ceṣṭante iti . \\
\noindent \textbf{(P\_5,1.119.2)} siddham tu yasya guṇasya bhāvāt dravye śabdaniveśaḥ tadabhidhāne tvatalau . \\
\noindent \textbf{(P\_5,1.119.2)} siddham etat . \\
\noindent \textbf{(P\_5,1.119.2)} katham . \\
\noindent \textbf{(P\_5,1.119.2)} yasya guṇasya bhāvāt dravye śabdaniveśaḥ tadabhidhāne tasmin guṇe vaktavye pratyayena bhavitavyam . \\
\noindent \textbf{(P\_5,1.119.2)} na ca abhiprāyādīnām bhāvāt dravye devadattaśabdaḥ vartate . \\
\noindent \textbf{(P\_5,1.119.2)} kim punaḥ dravyam ke punaḥ guṇāḥ . \\
\noindent \textbf{(P\_5,1.119.2)} śabdasparśarūparasagandhāḥ guṇāḥ . \\
\noindent \textbf{(P\_5,1.119.2)} tataḥ anyat dravyam . \\
\noindent \textbf{(P\_5,1.119.2)} kim punaḥ anyat śabdādibhyaḥ dravyam āhosvit ananyat . \\
\noindent \textbf{(P\_5,1.119.2)} guṇasya ayam bhāvāt dravye śabdaniveśam kurvan khyāpayati anyat śabdādibhyaḥ dravyam iti . \\
\noindent \textbf{(P\_5,1.119.2)} ananyat śabdādibhyaḥ dravyam . \\
\noindent \textbf{(P\_5,1.119.2)} na hi anyat upalabhyate . \\
\noindent \textbf{(P\_5,1.119.2)} paśoḥ khalu api viśasitasya parṇaśate nyastasya na anyat śabdādibhyaḥ upalabhyate . \\
\noindent \textbf{(P\_5,1.119.2)} anyat śabdādibhyaḥ dravyam tat tu anumānagamyam . \\
\noindent \textbf{(P\_5,1.119.2)} tat yathā oṣadhivanaspatīnām vṛddhihrāsau . \\
\noindent \textbf{(P\_5,1.119.2)} jyotiṣām gatiḥ iti . \\
\noindent \textbf{(P\_5,1.119.2)} kaḥ asau anumānaḥ . \\
\noindent \textbf{(P\_5,1.119.2)} iha samāne varṣmaṇi pariṇāhe ca anyat tulāgram bhavati lohasya anyat kārpāsānām . \\
\noindent \textbf{(P\_5,1.119.2)} yatkṛtaḥ viśeṣaḥ tat dravyam . \\
\noindent \textbf{(P\_5,1.119.2)} tathā kaḥ cit spṛśan eva chinatti kaḥ cit lambamānaḥ api na chinatti . \\
\noindent \textbf{(P\_5,1.119.2)} yatkṛtaḥ viśeṣaḥ tat dravyam . \\
\noindent \textbf{(P\_5,1.119.2)} kaḥ cit ekena eva prahāreṇa vyapavargam karoti kaḥ cit dvābhyām api an karoti . \\
\noindent \textbf{(P\_5,1.119.2)} yatkṛtaḥ viśeṣaḥ tat dravyam . \\
\noindent \textbf{(P\_5,1.119.2)} atha vā yasya guṇāntareṣu api prādurbhāvatsu tattvam na vihanyate tat dravyam . \\
\noindent \textbf{(P\_5,1.119.2)} kim punaḥ tattvam . \\
\noindent \textbf{(P\_5,1.119.2)} tadbhāvaḥ tattvam . \\
\noindent \textbf{(P\_5,1.119.2)} tat yathā āmalakādīnām phalānām raktādayaḥ pītādayaḥ ca guṇāḥ prāduḥ bhavanti . \\
\noindent \textbf{(P\_5,1.119.2)} āmalakam badaram iti eva bhavati . \\
\noindent \textbf{(P\_5,1.119.2)} anvartham khalu api nirvacanam . \\
\noindent \textbf{(P\_5,1.119.2)} guṇasandrāvaḥ dravyam iti . \\
\noindent \textbf{(P\_5,1.119.2)} yadi tarhi ṣaṣṭhīsamarthāt guṇe pratyayāḥ utapdyante kim iyatā sūtreṇa . \\
\noindent \textbf{(P\_5,1.119.2)} etāvat vaktavyam : ṣaṣṭhīsamarthāt guṇe iti . \\
\noindent \textbf{(P\_5,1.119.2)} ṣaṣṭhīsamarthāt guṇe iti iyati ucyamāne dviguṇā rajjuḥ triguṇā rajjuḥ atra api prāpnoti . \\
\noindent \textbf{(P\_5,1.119.2)} na eṣaḥ doṣaḥ . \\
\noindent \textbf{(P\_5,1.119.2)} guṇaśabdaḥ ayam bahvarthaḥ . \\
\noindent \textbf{(P\_5,1.119.2)} asti eva sameṣu avayaveṣu vartate . \\
\noindent \textbf{(P\_5,1.119.2)} tat yathā dviguṇā rajjuḥ triguṇā rajjuḥ iti . \\
\noindent \textbf{(P\_5,1.119.2)} asti dravyapadārthakaḥ . \\
\noindent \textbf{(P\_5,1.119.2)} tat yathā guṇavān ayam deśaḥ iti ucyate yasmin gāvaḥ sasyāni ca vartante . \\
\noindent \textbf{(P\_5,1.119.2)} asti aprādhānye vartate . \\
\noindent \textbf{(P\_5,1.119.2)} tat yathā yaḥ yatra apradhānam bhavati saḥ āha guṇabhūtāḥ vayam atra iti . \\
\noindent \textbf{(P\_5,1.119.2)} asti ācāre vartate . \\
\noindent \textbf{(P\_5,1.119.2)} tat yathā guṇavān ayam brāhmaṇaḥ iti ucyate yaḥ samyak ācāram karoti . \\
\noindent \textbf{(P\_5,1.119.2)} asti saṃskāre vartate . \\
\noindent \textbf{(P\_5,1.119.2)} tat yathā saṃskṛtam annam guṇavat iti ucyate . \\
\noindent \textbf{(P\_5,1.119.2)} atha vā sarvatra eva ayam guṇaśabdaḥ sameṣu avayaveṣu vartate . \\
\noindent \textbf{(P\_5,1.119.2)} tat yathā dviguṇam adhyayanam triguṇam adhyayanam iti ucyate . \\
\noindent \textbf{(P\_5,1.119.2)} carcāguṇān kramaguṇān ca apekṣya bhavati na saṃhitāguṇān carcāguṇān ca . \\
\noindent \textbf{(P\_5,1.119.2)} yadi evam guṇavat annam iti guṇaśabdaḥ na upapadyate . \\
\noindent \textbf{(P\_5,1.119.2)} na hi annasya sūpādayaḥ guṇāḥ samāḥ bhavanti . \\
\noindent \textbf{(P\_5,1.119.2)} na avaśyam varṣmataḥ parimāṇataḥ eva vā sāmyam bhavati . \\
\noindent \textbf{(P\_5,1.119.2)} kim tarhi yuktitaḥ api . \\
\noindent \textbf{(P\_5,1.119.2)} ātaḥ ca yuktitaḥ . \\
\noindent \textbf{(P\_5,1.119.2)} yaḥ hi mudgaprasthe lavaṇaprastham prakṣipet na adaḥ yuktam syāt . \\
\noindent \textbf{(P\_5,1.119.2)} yadi tāvat adeḥ annam na adaḥ attavyam syāt . \\
\noindent \textbf{(P\_5,1.119.2)} atha aniteḥ annam na adaḥ jagdvhā prāṇyāt . \\
\noindent \textbf{(P\_5,1.119.2)} śuklādiṣu tarhi vartyabhāvāt vṛttiḥ na prāpnoti . \\
\noindent \textbf{(P\_5,1.119.2)} śuklatvam . \\
\noindent \textbf{(P\_5,1.119.2)} śuklatā iti . \\
\noindent \textbf{(P\_5,1.119.2)} kim punaḥ kāraṇam śuklādayaḥ eva udāhriyante na punaḥ vṛkṣādayaḥ api . \\
\noindent \textbf{(P\_5,1.119.2)} vṛkṣatvam vṛkṣatā iti . \\
\noindent \textbf{(P\_5,1.119.2)} asti atra viśeṣaḥ . \\
\noindent \textbf{(P\_5,1.119.2)} ubhayavacanāḥ hi ete dravyam ca āhuḥ guṇam ca . \\
\noindent \textbf{(P\_5,1.119.2)} yataḥ dravyavacanāḥ tataḥ vṛttiḥ bhaviṣyati . \\
\noindent \textbf{(P\_5,1.119.2)} ime api tarhi ubhayavacanāḥ . \\
\noindent \textbf{(P\_5,1.119.2)} katham . \\
\noindent \textbf{(P\_5,1.119.2)} ārabhyate matublopaḥ guṇavacanebhyaḥ matupaḥ luk bhavati iti . \\
\noindent \textbf{(P\_5,1.119.2)} yataḥ dravyavacanāḥ tataḥ vṛttiḥ bhaviṣyati . \\
\noindent \textbf{(P\_5,1.119.2)} ḍitthādiṣu tarhi vartyabhāvāt vṛttiḥ na prāpnoti . \\
\noindent \textbf{(P\_5,1.119.2)} ḍitthatvam . \\
\noindent \textbf{(P\_5,1.119.2)} ḍitthatā . \\
\noindent \textbf{(P\_5,1.119.2)} ḍāmbhiṭṭatvam iti . \\
\noindent \textbf{(P\_5,1.119.2)} atra api kaḥ cit prāthamakalpikaḥ ḍitthaḥ ḍāmbhiṭṭaḥ ca . \\
\noindent \textbf{(P\_5,1.119.2)} tena kṛtām kriyām guṇam vā yaḥ kaḥ cit karoti saḥ ucyate ḍitthatvam te etat ḍāmbhiṭṭatvam te etat . \\
\noindent \textbf{(P\_5,1.119.2)} evam ḍitthāḥ kurvanti . \\
\noindent \textbf{(P\_5,1.119.2)} evam ḍāmbhiṭṭāḥ kurvanti . \\
\noindent \textbf{(P\_5,1.119.2)} yaḥ tarhi prāthamakalpikaḥ ḍitthaḥ ḍāmbhiṭṭaḥ ca tasya vartyabhāvāt vṛttiḥ na prāpnoti . \\
\noindent \textbf{(P\_5,1.119.2)} na eṣaḥ doṣaḥ . \\
\noindent \textbf{(P\_5,1.119.2)} yathā eva tasya kāthañcitkaḥ prayogaḥ evam vṛttiḥ api bhaviṣyati . \\
\noindent \textbf{(P\_5,1.119.2)} yat vā sarve bhāvāḥ svena bhāvena bhavanti saḥ teṣām bhāvaḥ tadabhidhāne . \\
\noindent \textbf{(P\_5,1.119.2)} kim ebhiḥ tribhiḥ bhāvagrahaṇaiḥ kriyate . \\
\noindent \textbf{(P\_5,1.119.2)} ekena śabdaḥ pratinirdiśyate dvābhyām arthaḥ . \\
\noindent \textbf{(P\_5,1.119.2)} yat vā sarve śabdāḥ svena arthena bhavanti . \\
\noindent \textbf{(P\_5,1.119.2)} saḥ teṣām arthaḥ iti tadabhidhāne vā tvatalau bhavataḥ iti vaktavyam . \\
\noindent \textbf{(P\_5,1.119.2)} na evam anyatra bhavati . \\
\noindent \textbf{(P\_5,1.119.2)} na hi tena raktam rāgāt iti atra śabdena rakte pratyayāḥ utpadyante . \\
\noindent \textbf{(P\_5,1.119.2)} śabde asambhavāt arthena rakte pratyayāḥ bhaviṣyanti . \\
\noindent \textbf{(P\_5,1.119.2)} tat tarhi anyatarat kartavyam . \\
\noindent \textbf{(P\_5,1.119.2)} sūtram ca bhidyate . \\
\noindent \textbf{(P\_5,1.119.2)} yathānyāsam eva astu . \\
\noindent \textbf{(P\_5,1.119.2)} nanu ca uktam tasya bhāvaḥ iti abhiprāyādiṣu atiprasaṅgaḥ iti . \\
\noindent \textbf{(P\_5,1.119.2)} uktam vā . \\
\noindent \textbf{(P\_5,1.119.2)} kim uktam . \\
\noindent \textbf{(P\_5,1.119.2)} anabhidhānāt iti . \\
\noindent \textbf{(P\_5,1.119.2)} anabhidhānāt abhiprāyādiṣu utapattiḥ na bhaviṣyati . \\
\noindent \textbf{(P\_5,1.119.3)} tvatalbhyām nañsamāsaḥ pūrvavipratiṣiddham tvataloḥ svarasiddhyartham . \\
\noindent \textbf{(P\_5,1.119.3)} tvatalbhyām nañsamāsaḥ bhavati pūrvavipratiṣdhena . \\
\noindent \textbf{(P\_5,1.119.3)} kim prayojanam . \\
\noindent \textbf{(P\_5,1.119.3)} tvataloḥ svarasiddhyartham . \\
\noindent \textbf{(P\_5,1.119.3)} tvataloḥ svarasiddhiḥ yathā syāt . \\
\noindent \textbf{(P\_5,1.119.3)} tvataloḥ avakāśaḥ bhāvasya vacanam pratiṣedhasya avacanam . \\
\noindent \textbf{(P\_5,1.119.3)} brāhmaṇatvam . \\
\noindent \textbf{(P\_5,1.119.3)} brāhmṇatā . \\
\noindent \textbf{(P\_5,1.119.3)} nañsamāsasya avakāśaḥ pratiṣedhasya vacanam bhāvasya avacanam . \\
\noindent \textbf{(P\_5,1.119.3)} abrāhmaṇaḥ . \\
\noindent \textbf{(P\_5,1.119.3)} avṛṣalaḥ . \\
\noindent \textbf{(P\_5,1.119.3)} ubhayavacane ubhayam prāpnoti . \\
\noindent \textbf{(P\_5,1.119.3)} abrāhmaṇatvam . \\
\noindent \textbf{(P\_5,1.119.3)} abrāhmaṇatā . \\
\noindent \textbf{(P\_5,1.119.3)} nañsamāsaḥ bhavati pūrvavipratiṣdhena . \\
\noindent \textbf{(P\_5,1.119.3)} saḥ tarhi pūrvavipratiṣedhaḥ vaktavyaḥ . \\
\noindent \textbf{(P\_5,1.119.3)} na vaktavyaḥ . \\
\noindent \textbf{(P\_5,1.119.3)} na atra tvatalau prāpnutaḥ . \\
\noindent \textbf{(P\_5,1.119.3)} kim kāraṇam . \\
\noindent \textbf{(P\_5,1.119.3)} asāmarthyāt . \\
\noindent \textbf{(P\_5,1.119.3)} katham asāmarthyam . \\
\noindent \textbf{(P\_5,1.119.3)} sāpekṣam asamartham bhavati iti . \\
\noindent \textbf{(P\_5,1.119.3)} yāvatā brāhmaṇaśabdaḥ pratiṣedham apekṣate . \\
\noindent \textbf{(P\_5,1.119.3)} nañsamāsaḥ api tarhi na prāpnoti . \\
\noindent \textbf{(P\_5,1.119.3)} kim kāraṇam . \\
\noindent \textbf{(P\_5,1.119.3)} asāmarthyāt eva . \\
\noindent \textbf{(P\_5,1.119.3)} katham asāmarthyam . \\
\noindent \textbf{(P\_5,1.119.3)} sāpekṣam asamartham bhavati iti . \\
\noindent \textbf{(P\_5,1.119.3)} yāvatā brāhmaṇaśabdaḥ bhāvam apekṣate . \\
\noindent \textbf{(P\_5,1.119.3)} pradhānam tadā brāhmaṇaśabdaḥ . \\
\noindent \textbf{(P\_5,1.119.3)} bhavati ca pradhānasya sāpekṣasya api samāsaḥ . \\
\noindent \textbf{(P\_5,1.119.3)} idam tarhi prayojanam . \\
\noindent \textbf{(P\_5,1.119.3)} nañsamādāt anyaḥ bhāvavacanaḥ svarottarapadavṛddhyartham iti vakṣyati . \\
\noindent \textbf{(P\_5,1.119.3)} tatra vyavasthārtham idam vaktavyam . \\
\noindent \textbf{(P\_5,1.119.3)} vā chandasi . \\
\noindent \textbf{(P\_5,1.119.3)} vā chandasi nañsamāsaḥ vaktavyaḥ . \\
\noindent \textbf{(P\_5,1.119.3)} nirvīryatām vai yajamānaḥ āśāste apaśutām . \\
\noindent \textbf{(P\_5,1.119.3)} ayonitvāya . \\
\noindent \textbf{(P\_5,1.119.3)} aśithilatvāya . \\
\noindent \textbf{(P\_5,1.119.3)} agotām anapatyatām . \\
\noindent \textbf{(P\_5,1.119.3)} bhavet idam yuktam udāharaṇam . \\
\noindent \textbf{(P\_5,1.119.3)} ayonitvāya . \\
\noindent \textbf{(P\_5,1.119.3)} aśithilatvāya iti . \\
\noindent \textbf{(P\_5,1.119.3)} idam tu ayuktam apaśutām iti . \\
\noindent \textbf{(P\_5,1.119.3)} na hi asau samāsabhāvam āśāste . \\
\noindent \textbf{(P\_5,1.119.3)} kim tarhi . \\
\noindent \textbf{(P\_5,1.119.3)} uttarapadābhāvam āśāste . \\
\noindent \textbf{(P\_5,1.119.3)} na paśoḥ bhāvaḥ iti . \\
\noindent \textbf{(P\_5,1.119.3)} nañsamāsāt anyaḥ bhāvavacanaḥ . \\
\noindent \textbf{(P\_5,1.119.3)} nañsamāsāt anyaḥ bhāvavacanaḥ bhavati vipratiṣedhena . \\
\noindent \textbf{(P\_5,1.119.3)} kim prayojanam . \\
\noindent \textbf{(P\_5,1.119.3)} svarottarapadavṛddhyartham . \\
\noindent \textbf{(P\_5,1.119.3)} svarārtham uttarapadavṛddhyartham ca . \\
\noindent \textbf{(P\_5,1.119.3)} svarārtham tāvat . \\
\noindent \textbf{(P\_5,1.119.3)} aprathimā . \\
\noindent \textbf{(P\_5,1.119.3)} amradimā . \\
\noindent \textbf{(P\_5,1.119.3)} uttarapadavṛddhyartham . \\
\noindent \textbf{(P\_5,1.119.3)} aśauklyam . \\
\noindent \textbf{(P\_5,1.119.3)} akārṣṇyam . \\
\noindent \textbf{(P\_5,1.120)} kimarthaḥ cakāraḥ . \\
\noindent \textbf{(P\_5,1.120)} anukarṣaṇārthaḥ . \\
\noindent \textbf{(P\_5,1.120)} tvatalau anukṛṣyete . \\
\noindent \textbf{(P\_5,1.120)} na etat asti prayojanam . \\
\noindent \textbf{(P\_5,1.120)} prakṛtau tvatalau anuvartiṣyete . \\
\noindent \textbf{(P\_5,1.120)} ataḥ uttaram paṭhati . \\
\noindent \textbf{(P\_5,1.120)} ā ca tvāt iti cakārakaraṇam apavādasamāveśārtham . \\
\noindent \textbf{(P\_5,1.120)} ā ca tvāt iti cakārakaraṇam kriyate apavādasamāveśārtham . \\
\noindent \textbf{(P\_5,1.120)} imanicprabhṛtibhiḥ apavādaiḥ samāveśaḥ yathā syāt . \\
\noindent \textbf{(P\_5,1.120)} na etat asti prayojanam . \\
\noindent \textbf{(P\_5,1.120)} ā tvāt iti evam imanicprabhṛtibhiḥ apavādaiḥ samāveśaḥ bhaviṣyati . \\
\noindent \textbf{(P\_5,1.120)} idam tarhi prayojanam . \\
\noindent \textbf{(P\_5,1.120)} ā tvāt yāḥ prakṛtayaḥ tābhyaḥ ca tvatalau yathā syātām yataḥ ca ucyete . \\
\noindent \textbf{(P\_5,1.120)} etat api na asti prayojanam . \\
\noindent \textbf{(P\_5,1.120)} ā tvāt iti eva yāḥ prakṛtayaḥ tābhyaḥ tvatalau bhaviṣyataḥ yataḥ ca ucyete . \\
\noindent \textbf{(P\_5,1.120)} idam tarhi prayojanam . \\
\noindent \textbf{(P\_5,1.120)} ā tvāt ye arthāḥ tatra tvatalau yathā syātām yatra ca ucyete . \\
\noindent \textbf{(P\_5,1.120)} etat api na asti prayojanam . \\
\noindent \textbf{(P\_5,1.120)} ā tvāt iti eva ā tvāt ye arthāḥ tatra tvatalau bhaviṣyataḥ yatra ca ucyete . \\
\noindent \textbf{(P\_5,1.120)} idam tarhi prayojanam . \\
\noindent \textbf{(P\_5,1.120)} ā tvāt yāḥ prakṛtayaḥ tābhyaḥ ca tvatalau yathā syātām yasyāḥ ca prakṛteḥ atasmin viśeṣe anyaḥ pratyayaḥ utpadyate . \\
\noindent \textbf{(P\_5,1.120)} kim kṛtam bhavati . \\
\noindent \textbf{(P\_5,1.120)} strīpuṃsābhyām tvataloḥ upasaṅkhyānam coditam . \\
\noindent \textbf{(P\_5,1.120)} tat na vaktavyam bhavati . \\
\noindent \textbf{(P\_5,1.121)} kasya ayam pratiṣedhaḥ . \\
\noindent \textbf{(P\_5,1.121)} tvataloḥ iti āha . \\
\noindent \textbf{(P\_5,1.121)} na etat asi prayojanam . \\
\noindent \textbf{(P\_5,1.121)} iṣyete nañpūrvāt tatpuruṣāt tvatalau : abrāhmaṇatvam abrāhmaṇatā iti . \\
\noindent \textbf{(P\_5,1.121)} ataḥ uttaram paṭhati . \\
\noindent \textbf{(P\_5,1.121)} na nañpūrvāt iti uttarasya pratiṣedhaḥ . \\
\noindent \textbf{(P\_5,1.121)} na nañpūrvāt iti uttarasya bhāvapratyayasya pratiṣedhaḥ kriyate . \\
\noindent \textbf{(P\_5,1.121)} na etat asti prayojanam . \\
\noindent \textbf{(P\_5,1.121)} parigaṇitābhyaḥ prakṛtibhyaḥ uttaraḥ bhāvapratyayaḥ vidhīyate . \\
\noindent \textbf{(P\_5,1.121)} na ca tatra kā cit nañpūrvā prakṛtiḥ gṛhyate . \\
\noindent \textbf{(P\_5,1.121)} tadantavidhinā prāpnoti . \\
\noindent \textbf{(P\_5,1.121)} grahaṇavatā prātipadikena tadantavidhiḥ pratiṣidhyate . \\
\noindent \textbf{(P\_5,1.121)} yatra tarhi tadantavidhiḥ asti . \\
\noindent \textbf{(P\_5,1.121)} patyantapurohitādibhyaḥ yak iti . \\
\noindent \textbf{(P\_5,1.121)} yadi etāvat prayojanam syāt tatra eva ayam brūyāt apatyantāt yak bhavati nañpūrvāt tatpuruṣāt iti . \\
\noindent \textbf{(P\_5,1.121)} evam tarhi jñāpayati ācārayaḥ uttaraḥ bhāvapratyayaḥ nañpūrvāt bahuvrīheḥ bhavati iti . \\
\noindent \textbf{(P\_5,1.121)} na iṣyate . \\
\noindent \textbf{(P\_5,1.121)} tvatalau eva iṣyete : avidyamānāḥ pṛthavaḥ asya apṛthuḥ , apṛthoḥ bhāvaḥ apṛthutvam apṛthutā iti . \\
\noindent \textbf{(P\_5,1.121)} evam tarhi jñāpayati ācārayaḥ uttaraḥ bhāvapratyayaḥ anyapūrvāt tatpuruṣāt bhavati iti . \\
\noindent \textbf{(P\_5,1.121)} na iṣyate . \\
\noindent \textbf{(P\_5,1.121)} tvatalau eva iṣyete : paramaḥ pṛthuḥ paramapṛthuḥ , paramapṛthoḥ bhāvaḥ paramapṛthutvam paramapṛthutā . \\
\noindent \textbf{(P\_5,1.121)} evam tarhi jñāpayati ācārayaḥ uttaraḥ bhāvapratyayaḥ sāpekṣāt bhavati iti . \\
\noindent \textbf{(P\_5,1.121)} kim etasya jñāpane prayojanam . \\
\noindent \textbf{(P\_5,1.121)} nañsamāsāt anyaḥ bhāvavacanaḥ svarottarapadavṛddhyartham iti uktam . \\
\noindent \textbf{(P\_5,1.121)} tat upapannam bhavati . \\
\noindent \textbf{(P\_5,1.121)} etat api na asti prayojanam . \\
\noindent \textbf{(P\_5,1.121)} ācāryapravṛttiḥ jñāpayati sarve ete taddhitāḥ sāpekṣāt bhavanti iti yat ayam nañaḥ guṇapratiṣedhe sampādyarhahitālamarthāḥ taddhitāḥ iti āha . \\
\noindent \textbf{(P\_5,1.122)} vāvacanam kimartham . \\
\noindent \textbf{(P\_5,1.122)} vākyam api yathā syāt . \\
\noindent \textbf{(P\_5,1.122)} na etat asti prayojanam . \\
\noindent \textbf{(P\_5,1.122)} prakṛṭā mahāvibhāṣā . \\
\noindent \textbf{(P\_5,1.122)} tayā vākyam api bhaviṣyati . \\
\noindent \textbf{(P\_5,1.122)} idam tarhi prayojanam . \\
\noindent \textbf{(P\_5,1.122)} tvatalau api yathā syātām . \\
\noindent \textbf{(P\_5,1.122)} etat api na asti prayojanam . \\
\noindent \textbf{(P\_5,1.122)} ā ca tvāt iti etasmāt yatnāt tvatalau api bhaviṣyataḥ . \\
\noindent \textbf{(P\_5,1.122)} ataḥ uttaram paṭhati . \\
\noindent \textbf{(P\_5,1.122)} pṛthvādibhyaḥ vāvacanam aṇsamāveśārtham . \\
\noindent \textbf{(P\_5,1.122)} pṛthvādibhyaḥ vāvacanam kriyate aṇsamāveśaḥ yathā syāt . \\
\noindent \textbf{(P\_5,1.122)} pārthavam . \\
\noindent \textbf{(P\_5,1.122)} prathimā . \\
\noindent \textbf{(P\_5,1.124)} brāhmaṇādiṣu cāturvarṇyādīnām upasaṅkhyānam . brāhmaṇādiṣu cāturvarṇyādīnām upasaṅkhyānam kartavyam . cāturvarṇyam . cāturvaidyam . cāturāśramyam . arhataḥ num ca . arhataḥ num ca ṣyañ ca vaktavyaḥ . arhataḥ bhāvaḥ ārhantyam ārhatī . \\
\noindent \textbf{(P\_5,1.125)} kim idam nalope varṇagrahaṇam āhosvit saṅghātagrahaṇam . \\
\noindent \textbf{(P\_5,1.125)} kim ca ataḥ . \\
\noindent \textbf{(P\_5,1.125)} yadi varṇagrahaṇam steyam . \\
\noindent \textbf{(P\_5,1.125)} nalope kṛte ayādeśaḥ prāpnoti . \\
\noindent \textbf{(P\_5,1.125)} atha saṅghātagrahaṇam antyasya lopaḥ kasmāt na bhavati . \\
\noindent \textbf{(P\_5,1.125)} siddhaḥ antyasya lopaḥ yasya iti eva . \\
\noindent \textbf{(P\_5,1.125)} tatra ārambhasāmarthyāt sarvasya bhaviṣyati . \\
\noindent \textbf{(P\_5,1.130)} aṇprakaraṇe śrotriyasya . \\
\noindent \textbf{(P\_5,1.130)} aṇprakaraṇe śrotriyasya upasaṅkhyānam kartavyam ghalopaḥ ca . \\
\noindent \textbf{(P\_5,1.130)} śrotriyasya bhāvaḥ śrautram . \\
\noindent \textbf{(P\_5,2.4)} tilādibhyaḥ khañ ca . \\
\noindent \textbf{(P\_5,2.4)} tilādibhyaḥ khañ ca iti vaktavyam . \\
\noindent \textbf{(P\_5,2.4)} tilyam . \\
\noindent \textbf{(P\_5,2.4)} tailīnam . \\
\noindent \textbf{(P\_5,2.4)} kimartham idam ucyate na yatā mukte dhānyānām bhavane kṣetre khañ iti eva siddham . \\
\noindent \textbf{(P\_5,2.4)} na sidhyati . \\
\noindent \textbf{(P\_5,2.4)} kim kāraṇam . \\
\noindent \textbf{(P\_5,2.4)} umābhaṅgayoḥ adhānyatvāt . \\
\noindent \textbf{(P\_5,2.4)} dhānyānām bhavane kṣetre khañ iti ucyate . \\
\noindent \textbf{(P\_5,2.4)} na ca umābhaṅge dhānye . \\
\noindent \textbf{(P\_5,2.4)} cameṣu yat paṭhyate tat dhānyam . \\
\noindent \textbf{(P\_5,2.4)} na ca ete tatra paṭhyete . \\
\noindent \textbf{(P\_5,2.4)} tat tarhi khañgrahaṇam kartavyam . \\
\noindent \textbf{(P\_5,2.4)} na kartavyam . \\
\noindent \textbf{(P\_5,2.4)} prakṛtam anuvartate . \\
\noindent \textbf{(P\_5,2.4)} kva prakṛtam . \\
\noindent \textbf{(P\_5,2.4)} dhānyānām bhavane kṣetre khañ iti . \\
\noindent \textbf{(P\_5,2.4)} yadi tat anuvartate vrīhiśālayoḥ ḍhak yavayavakaṣaṣṭikāt yat iti khañ ca iti khañ api prāpnoti . \\
\noindent \textbf{(P\_5,2.4)} sambandham anuvartiṣyate . \\
\noindent \textbf{(P\_5,2.4)} dhānyānām bhavane kṣetre khañ . \\
\noindent \textbf{(P\_5,2.4)} vrīhiśālayoḥ ḍhak bhavati . \\
\noindent \textbf{(P\_5,2.4)} dhānyānām bhavane kṣetre khañ . \\
\noindent \textbf{(P\_5,2.4)} yavayavakaṣaṣṭikāt yat bhavati . \\
\noindent \textbf{(P\_5,2.4)} dhānyānām bhavane kṣetre khañ bhavati . \\
\noindent \textbf{(P\_5,2.4)} vibhāṣā tilamāṣomābhaṅgaṇubhyaḥ . \\
\noindent \textbf{(P\_5,2.4)} bhavanekṣetregrahaṇam anuvartate . \\
\noindent \textbf{(P\_5,2.4)} dhānyānām iti nivṛttam . \\
\noindent \textbf{(P\_5,2.4)} atha vā maṇḍūkaplutayaḥ adhikārāḥ . \\
\noindent \textbf{(P\_5,2.4)} yathā maṇḍūkāḥ utplutya utplutya gacchanti tadvat adhikārāḥ . \\
\noindent \textbf{(P\_5,2.4)} atha vā anyavacanāt cakārākaraṇāt prakṛtāpavādaḥ vijñāyate yathā utsargeṇa prasaktasya apavādaḥ . \\
\noindent \textbf{(P\_5,2.4)} anyasya pratyayasya vacanāt cakārasya ca anukarṣaṇārthasya akaraṇāt prakṛtasya khañaḥ ḍhagyatau bādhakau bhaviṣyataḥ yathā utsargeṇa prasaktasya apavādaḥ bādhakaḥ bhavati . \\
\noindent \textbf{(P\_5,2.4)} atha vā etat jñāpayati anuvartante ca nāma vidhayaḥ na ca anuvartanāt eva bhavanti . \\
\noindent \textbf{(P\_5,2.4)} kim tarhi . \\
\noindent \textbf{(P\_5,2.4)} yatnāt bhavanti . \\
\noindent \textbf{(P\_5,2.4)} atha vā yatā mukte dhānyānām bhavane kṣetre khañ iti eva siddham . \\
\noindent \textbf{(P\_5,2.4)} nanu ca uktam . \\
\noindent \textbf{(P\_5,2.4)} na sidhyati . \\
\noindent \textbf{(P\_5,2.4)} kim kāraṇam . \\
\noindent \textbf{(P\_5,2.4)} umābhaṅgayoḥ adhānyatvāt iti . \\
\noindent \textbf{(P\_5,2.4)} na eṣaḥ doṣaḥ . \\
\noindent \textbf{(P\_5,2.4)} dhinoteḥ dhānyam . \\
\noindent \textbf{(P\_5,2.4)} ete ca api dhinutaḥ . \\
\noindent \textbf{(P\_5,2.4)} atha vā śaṇasaptadaśāni dhānyāni . \\
\noindent \textbf{(P\_5,2.6)} sammukha iti kim nipātyate . \\
\noindent \textbf{(P\_5,2.6)} sammukha iti samasya antalopaḥ . \\
\noindent \textbf{(P\_5,2.6)} sammukha iti samasya antalopaḥ nipātyate . \\
\noindent \textbf{(P\_5,2.6)} samamukhasya darśanaḥ sammukhīnaḥ . \\
\noindent \textbf{(P\_5,2.9)} ayānayam neyaḥ iti ucyate . \\
\noindent \textbf{(P\_5,2.9)} tatra na jñāyate kaḥ ayaḥ kaḥ anayaḥ iti . \\
\noindent \textbf{(P\_5,2.9)} ayaḥ pradakṣiṇam . \\
\noindent \textbf{(P\_5,2.9)} anayaḥ prasavyam . \\
\noindent \textbf{(P\_5,2.9)} pradakṣiṇaprasavyagaminām śārāṇām yasmin paraiḥ padānām asamāveśaḥ saḥ ayānayaḥ . \\
\noindent \textbf{(P\_5,2.9)} ayānayam neyaḥ ayānayīnaḥ śāraḥ . \\
\noindent \textbf{(P\_5,2.10)} parovara iti kim nipātyate . \\
\noindent \textbf{(P\_5,2.10)} parovara iti parasotvavacanam . \\
\noindent \textbf{(P\_5,2.10)} parovara iti parasya otvam nipātyate . \\
\noindent \textbf{(P\_5,2.10)} yadi evam parasyautvavacanam iti prāpnoti . \\
\noindent \textbf{(P\_5,2.10)} śakandhunyāyena nirdeśaḥ . \\
\noindent \textbf{(P\_5,2.10)} atha vā na evam vijñāyate parasya otvam nipātyate iti . \\
\noindent \textbf{(P\_5,2.10)} katham tarhi . \\
\noindent \textbf{(P\_5,2.10)} parasya śabdarūpasya ādeḥ utvam nipātyate iti . \\
\noindent \textbf{(P\_5,2.10)} parān ca avarān ca anubhavati parovarīṇaḥ . \\
\noindent \textbf{(P\_5,2.10)} atha parampara iti kim nipātyate . \\
\noindent \textbf{(P\_5,2.10)} paraparatarāṇām paramparabhāvaḥ . paraparatarāṇām paramparabhāvaḥ nipātyate . \\
\noindent \textbf{(P\_5,2.10)} parān ca paratarān ca anubhavati paramparīṇaḥ . \\
\noindent \textbf{(P\_5,2.12)} iha samāṃsamīnā gauḥ supaḥ dhātuprātipadikayoḥ iti subluk prāpnoti . \\
\noindent \textbf{(P\_5,2.12)} samām samām vijāyate iti yalopavacanāt alugvijñānam . \\
\noindent \textbf{(P\_5,2.12)} samām samām vijāyate iti yalopavacanāt alugvijñānam bhaviṣyati . \\
\noindent \textbf{(P\_5,2.12)} yat ayam yalopam śāsti tat jñāpayati ācāryaḥ na atra luk bhavati iti . \\
\noindent \textbf{(P\_5,2.12)} samām samām vijāyate iti yalopavacanāt alugvijñānam iti cet uttarapadasya lugvacanam . \\
\noindent \textbf{(P\_5,2.12)} samām samām vijāyate iti yalopavacanāt alugvijñānam iti cet uttarapadasya luk vaktavyaḥ . \\
\noindent \textbf{(P\_5,2.12)} siddham tu pūrvapadasya yalopavacanāt . \\
\noindent \textbf{(P\_5,2.12)} siddham etat . \\
\noindent \textbf{(P\_5,2.12)} katham . \\
\noindent \textbf{(P\_5,2.12)} pūrvapadasya yalopaḥ vaktavyaḥ . \\
\noindent \textbf{(P\_5,2.12)} anutpattau uttarapadasya ca vāvacanam . \\
\noindent \textbf{(P\_5,2.12)} anutpattau pūrvapadasya uttarapadasya ca yalopaḥ vā vaktavyaḥ . \\
\noindent \textbf{(P\_5,2.12)} samām samām vijāyate . \\
\noindent \textbf{(P\_5,2.12)} samāyām samāyām vijāyate iti . \\
\noindent \textbf{(P\_5,2.14)} āgavīnaḥ iti kim nipātyate . \\
\noindent \textbf{(P\_5,2.14)} goḥ āṅpūrvāt a tasya goḥ pratidānāt kāriṇi khaḥ . \\
\noindent \textbf{(P\_5,2.14)} goḥ āṅpūrvāt a tasya goḥ pratidānāt kāriṇi khaḥ nipātyate . \\
\noindent \textbf{(P\_5,2.14)} a tasya goḥ pratidānāt karmakārī āgavīnaḥ karmakaraḥ . \\
\noindent \textbf{(P\_5,2.20)} kim yaḥ śālāyām adhṛṣṭaḥ saḥ śālīnaḥ kūpe vā yat akāryam tat kaupīnam . \\
\noindent \textbf{(P\_5,2.20)} na iti āha . \\
\noindent \textbf{(P\_5,2.20)} uttarapadalopaḥ atra draṣṭavyaḥ . \\
\noindent \textbf{(P\_5,2.20)} śālāpraveśanam arhati adhṛṣṭaḥ saḥ śālīnaḥ . \\
\noindent \textbf{(P\_5,2.20)} kūpāvataraṇam arhati akāryam tat kaupīnam . \\
\noindent \textbf{(P\_5,2.21)} vrātena jīvati iti ucyate . \\
\noindent \textbf{(P\_5,2.21)} kim vrātam nāma . \\
\noindent \textbf{(P\_5,2.21)} nānājātīyāḥ aniyatavṛttayaḥ utsedhajīvinaḥ saṅghāḥ vrātāḥ . \\
\noindent \textbf{(P\_5,2.21)} teṣām karma vrātam . \\
\noindent \textbf{(P\_5,2.21)} vrātakarmaṇā jīvati iti vrātīnaḥ . \\
\noindent \textbf{(P\_5,2.23)} haiyaṅgavīnam iti kim nipātyate . \\
\noindent \textbf{(P\_5,2.23)} hyogodohasya hiyaṅgvādeśaḥ sañjñāyām tasya vikāre . \\
\noindent \textbf{(P\_5,2.23)} hyogodohasya hiyaṅgvādeśaḥ nipātyate sañjñāyām viṣaye tasya vikāre iti etasmin arthe . \\
\noindent \textbf{(P\_5,2.23)} hyogodohasya vikāraḥ haiyaṅgavīnam ghṛtam . \\
\noindent \textbf{(P\_5,2.23)} sañjñāyām iti kimartham . \\
\noindent \textbf{(P\_5,2.23)} hyogodohasya vikāraḥ udaśvit . \\
\noindent \textbf{(P\_5,2.23)} atra mā bhūt iti . \\
\noindent \textbf{(P\_5,2.27)} iha nānā iti sahārthaḥ gamyeta . \\
\noindent \textbf{(P\_5,2.27)} dvau hi pratiṣedhau prakṛtam artham gamayataḥ . \\
\noindent \textbf{(P\_5,2.27)} na na saḥ saha eva iti . \\
\noindent \textbf{(P\_5,2.27)} na eṣaḥ doṣaḥ . \\
\noindent \textbf{(P\_5,2.27)} na ayam pratyayārthaḥ . \\
\noindent \textbf{(P\_5,2.27)} kim tarhi prakṛtiviśeṣaṇam etat . \\
\noindent \textbf{(P\_5,2.27)} vi nañ iti etābhyām asahavācibhyām nānāñau bhavataḥ . \\
\noindent \textbf{(P\_5,2.27)} kasmin arthe . \\
\noindent \textbf{(P\_5,2.27)} svāṛthe . \\
\noindent \textbf{(P\_5,2.28)} kasmin arthe śālajādayaḥ bhavanti . \\
\noindent \textbf{(P\_5,2.28)} na saha iti vartate . \\
\noindent \textbf{(P\_5,2.28)} bhavet siddham viśāle śṛṅge viśaṅkaṭe śṛṅge iti . \\
\noindent \textbf{(P\_5,2.28)} iha khalu saṅkaṭam iti saṅgatārthaḥ gamyate . \\
\noindent \textbf{(P\_5,2.28)} prakaṭam iti pragarārthaḥ gamyate . \\
\noindent \textbf{(P\_5,2.28)} utkaṭam iti udgatārthaḥ gamyate . \\
\noindent \textbf{(P\_5,2.28)} evam tarhi sādhane śālajādayaḥ bhavanti . \\
\noindent \textbf{(P\_5,2.28)} kim vaktavyam etat . \\
\noindent \textbf{(P\_5,2.28)} na hi . \\
\noindent \textbf{(P\_5,2.28)} katham anucyamānam gaṃsyate . \\
\noindent \textbf{(P\_5,2.28)} upasargebhyaḥ ime vidhīyante . \\
\noindent \textbf{(P\_5,2.28)} upasargāḥ ca punaḥ evamātmakāḥ yatra kaḥ cit kriyāvācai śabdaḥ prayujyate tatra kriyāviśeṣam āhuḥ . \\
\noindent \textbf{(P\_5,2.28)} yatra hi na prayujyate sasādhanam tatra kriyām āhuḥ . \\
\noindent \textbf{(P\_5,2.28)} te ete upasargebhyaḥ vidhīyamānāḥ sasādhanāyām kriyāyām bhaviṣyanti . \\
\noindent \textbf{(P\_5,2.28)} evam api bhavet siddham viśāle śṛṅge iti . \\
\noindent \textbf{(P\_5,2.28)} idam tu na sidhyati . \\
\noindent \textbf{(P\_5,2.28)} viśālaḥ . \\
\noindent \textbf{(P\_5,2.28)} viśaṅkaṭaḥ iti . \\
\noindent \textbf{(P\_5,2.28)} etat api siddham . \\
\noindent \textbf{(P\_5,2.28)} katham . \\
\noindent \textbf{(P\_5,2.28)} akāraḥ matvarthīyaḥ . \\
\noindent \textbf{(P\_5,2.28)} viśāle asya staḥ viśālaḥ . \\
\noindent \textbf{(P\_5,2.28)} viśaṅkaṭe asya staḥ viśaṅkaṭaḥ iti . \\
\noindent \textbf{(P\_5,2.29)} kaṭacprakaraṇe alābūtilomābhyaḥ rajasi upasaṅkhyānam . \\
\noindent \textbf{(P\_5,2.29)} kaṭacprakaraṇe alābūtilomābhyaḥ rajasi abhidheye upasaṅkhyānam kartavyam . \\
\noindent \textbf{(P\_5,2.29)} alābūkaṭaḥ . \\
\noindent \textbf{(P\_5,2.29)} tilakaṭaḥ . \\
\noindent \textbf{(P\_5,2.29)} umākaṭaḥ . \\
\noindent \textbf{(P\_5,2.29)} bhaṅgāyāḥ ca . \\
\noindent \textbf{(P\_5,2.29)} bhaṅgāyāḥ ca iti vaktavyam . \\
\noindent \textbf{(P\_5,2.29)} bhaṅgākaṭaḥ . \\
\noindent \textbf{(P\_5,2.29)} goṣṭhādayaḥ sthānādiṣu paśunāmādibhyaḥ . \\
\noindent \textbf{(P\_5,2.29)} goṣṭhādayaḥ pratyayāḥ sthānādiṣu artheṣu paśunāmādibhyaḥ vaktavyāḥ . \\
\noindent \textbf{(P\_5,2.29)} gogoṣṭham . \\
\noindent \textbf{(P\_5,2.29)} avigoṣṭham . \\
\noindent \textbf{(P\_5,2.29)} kaṭac ca vaktavyaḥ . \\
\noindent \textbf{(P\_5,2.29)} avikaṭaḥ uṣṭrakaṭaḥ . \\
\noindent \textbf{(P\_5,2.29)} paṭac ca vaktavyaḥ . \\
\noindent \textbf{(P\_5,2.29)} avipaṭaḥ uṣṭrapaṭaḥ . \\
\noindent \textbf{(P\_5,2.29)} goyugaśabdaḥ ca pratyayaḥ vaktavyaḥ . \\
\noindent \textbf{(P\_5,2.29)} uṣṭragoyugam . \\
\noindent \textbf{(P\_5,2.29)} kharagoyugam . \\
\noindent \textbf{(P\_5,2.29)} tailaśabdaḥ ca pratyayaḥ vaktavyaḥ . \\
\noindent \textbf{(P\_5,2.29)} iṅgudatailam . \\
\noindent \textbf{(P\_5,2.29)} sarṣapatailam . \\
\noindent \textbf{(P\_5,2.29)} śākaṭaśabdaḥ ca pratyayaḥ vaktavyaḥ . \\
\noindent \textbf{(P\_5,2.29)} ikṣuśākaṭam . \\
\noindent \textbf{(P\_5,2.29)} mūlaśākaṭam . \\
\noindent \textbf{(P\_5,2.29)} śākinaśabdaḥ ca pratyayaḥ vaktavyaḥ . \\
\noindent \textbf{(P\_5,2.29)} ikṣuśākinam mūlaśākinam . \\
\noindent \textbf{(P\_5,2.29)} upamānāt vā siddham . \\
\noindent \textbf{(P\_5,2.29)} upamānāt vā siddham etat . \\
\noindent \textbf{(P\_5,2.29)} gavām sthānam goṣṭham . \\
\noindent \textbf{(P\_5,2.29)} yathā gavam tadvat uṣṭrāṇām . \\
\noindent \textbf{(P\_5,2.29)} kaṭac vaktavyaḥ iti . \\
\noindent \textbf{(P\_5,2.29)} yathā nānādravyāṇām saṅghātaḥ kaṭaḥ evam avayaḥ saṃhatāḥ avikaṭaḥ . \\
\noindent \textbf{(P\_5,2.29)} paṭat ca vaktayaḥ iti . \\
\noindent \textbf{(P\_5,2.29)} yathā paṭaḥ prastīrṇaḥ evam avayaḥ prastīrṇāḥ avipaṭaḥ . \\
\noindent \textbf{(P\_5,2.29)} goyugaśabdaḥ ca pratyayaḥ vaktavyaḥ iti . \\
\noindent \textbf{(P\_5,2.29)} goḥ yugam goyugam . \\
\noindent \textbf{(P\_5,2.29)} yathā goḥ tadvat uṣṭrasya . \\
\noindent \textbf{(P\_5,2.29)} uṣṭragoyugam . \\
\noindent \textbf{(P\_5,2.29)} tailaśabdaḥ ca pratyayaḥ vaktavyaḥ iti . \\
\noindent \textbf{(P\_5,2.29)} prakṛtyantaram tailaśabdaḥ vikāre vartate . \\
\noindent \textbf{(P\_5,2.29)} evam ca kṛtvā tilatailam iti api siddham bhavati . \\
\noindent \textbf{(P\_5,2.29)} śākaṭaśabdaḥ ca pratyayaḥ vaktavyaḥ eva . \\
\noindent \textbf{(P\_5,2.29)} śākinaśabdaḥ ca pratyayaḥ vaktavyaḥ eva . \\
\noindent \textbf{(P\_5,2.33)} inacpiṭackāḥ pratyayāḥ vaktavyāḥ cikacicik iti ete ca prakṛtyādeśāḥ vaktavyāḥ . \\
\noindent \textbf{(P\_5,2.33)} cikinaḥ . \\
\noindent \textbf{(P\_5,2.33)} cipiṭaḥ . \\
\noindent \textbf{(P\_5,2.33)} cikkaḥ . \\
\noindent \textbf{(P\_5,2.33)} klinnasya cilpil laḥ ca asya cakṣuṣī . \\
\noindent \textbf{(P\_5,2.33)} klinnasya cil pil iti etau prakṛtyādeśau vaktavyau laḥ ca pratyayaḥ asya cakṣuṣī iti etasmin arthe . \\
\noindent \textbf{(P\_5,2.33)} klinne asya cakṣuṣī cillaḥ . \\
\noindent \textbf{(P\_5,2.33)} pillaḥ . \\
\noindent \textbf{(P\_5,2.33)} cul ca vaktavyaḥ . \\
\noindent \textbf{(P\_5,2.33)} cullaḥ . \\
\noindent \textbf{(P\_5,2.33)} yadi asya iti ucyate cille cakṣuṣī pille cakṣuṣī iti na sidhyati . \\
\noindent \textbf{(P\_5,2.33)} tasmān na arthaḥ asya grahaṇe . \\
\noindent \textbf{(P\_5,2.33)} katham cillaḥ pillaḥ iti . \\
\noindent \textbf{(P\_5,2.33)} akāraḥ matvarthīyaḥ . \\
\noindent \textbf{(P\_5,2.33)} cille asya staḥ cillaḥ . \\
\noindent \textbf{(P\_5,2.33)} pille asya staḥ pillaḥ iti . \\
\noindent \textbf{(P\_5,2.37)} pramāṇe iti kimayam pratyayārthaḥ . \\
\noindent \textbf{(P\_5,2.37)} pramāṇam pratyayārthaḥ na . pramāṇe iti na ayam pratyayārthaḥ . \\
\noindent \textbf{(P\_5,2.37)} kva tarhi pratyayāḥ bhavanti . \\
\noindent \textbf{(P\_5,2.37)} tadvati . \\
\noindent \textbf{(P\_5,2.37)} kutaḥ etat . \\
\noindent \textbf{(P\_5,2.37)} asya iti vartanāt . \\
\noindent \textbf{(P\_5,2.37)} asya iti vartate . \\
\noindent \textbf{(P\_5,2.37)} kva prakṛtam . \\
\noindent \textbf{(P\_5,2.37)} tad asya sañjātam tārakādibhyaḥ itac iti . \\
\noindent \textbf{(P\_5,2.37)} prathamaḥ ca dvitīyaḥ ca ūrdhvamāne matau mama . \\
\noindent \textbf{(P\_5,2.37)} ūrudvayasam . \\
\noindent \textbf{(P\_5,2.37)} ūrudaghnam . \\
\noindent \textbf{(P\_5,2.37)} pramāṇe laḥ . \\
\noindent \textbf{(P\_5,2.37)} pramāṇe laḥ vaktavyaḥ : śamaḥ , diṣṭiḥ , vitastiḥ . \\
\noindent \textbf{(P\_5,2.37)} dvigoḥ nityam . \\
\noindent \textbf{(P\_5,2.37)} dvigoḥ nityam laḥ vaktavyaḥ . \\
\noindent \textbf{(P\_5,2.37)} dviśatam . \\
\noindent \textbf{(P\_5,2.37)} triśatam . \\
\noindent \textbf{(P\_5,2.37)} dvidiṣṭiḥ . \\
\noindent \textbf{(P\_5,2.37)} tridiṣṭiḥ . \\
\noindent \textbf{(P\_5,2.37)} dvivitastiḥ . \\
\noindent \textbf{(P\_5,2.37)} trivitastiḥ . \\
\noindent \textbf{(P\_5,2.37)} kimartham idam ucyate . \\
\noindent \textbf{(P\_5,2.37)} saṃśaye śrāviṇam vakṣyati yasya asyam purastāt apakarṣaḥ . \\
\noindent \textbf{(P\_5,2.37)} ḍaṭ stome . \\
\noindent \textbf{(P\_5,2.37)} ḍaṭ stome vaktavyaḥ . \\
\noindent \textbf{(P\_5,2.37)} pañcadaśaḥ stomaḥ . \\
\noindent \textbf{(P\_5,2.37)} śacśanoḥ ḍiniḥ . \\
\noindent \textbf{(P\_5,2.37)} śacśanoḥ ḍiniḥ vaktavyaḥ . \\
\noindent \textbf{(P\_5,2.37)} triṃśinaḥ māsāḥ . \\
\noindent \textbf{(P\_5,2.37)} pañcadaśinaḥ ardhamāsāḥ . \\
\noindent \textbf{(P\_5,2.37)} viṃśateḥ ca iti vaktavyam . \\
\noindent \textbf{(P\_5,2.37)} viṃśinaḥ aṅgirasaḥ . \\
\noindent \textbf{(P\_5,2.37)} pramāṇaparimāṇābhyām saṅkhyāyāḥ ca saṃśaye . \\
\noindent \textbf{(P\_5,2.37)} mātrac vaktavyaḥ . \\
\noindent \textbf{(P\_5,2.37)} śamamātram . \\
\noindent \textbf{(P\_5,2.37)} diṣṭimātram . \\
\noindent \textbf{(P\_5,2.37)} vitastimātram . \\
\noindent \textbf{(P\_5,2.37)} kuḍavamātram . \\
\noindent \textbf{(P\_5,2.37)} pañcamātrāḥ . \\
\noindent \textbf{(P\_5,2.37)} daśamātrāḥ . \\
\noindent \textbf{(P\_5,2.37)} vatvantāt svārthe dvayasajmātracau bahulam . \\
\noindent \textbf{(P\_5,2.37)} vatvantāt svārthe dvayasajmātracau bahulam vaktavyau . \\
\noindent \textbf{(P\_5,2.37)} tāvat eva tāvaddvayasam . \\
\noindent \textbf{(P\_5,2.37)} tāvanmātram . \\
\noindent \textbf{(P\_5,2.37)} yāvat eva yāvaddvayasam . \\
\noindent \textbf{(P\_5,2.37)} yāvanmātram . \\
\noindent \textbf{(P\_5,2.37)} [pramāṇam pratyayārthaḥ na tadvati asya iti vartanāt . \\
\noindent \textbf{(P\_5,2.37)} prathamaḥ ca dvitīyaḥ ca ūrdhvamāne matau mama . \\
\noindent \textbf{(P\_5,2.37)} pramāṇe laḥ . \\
\noindent \textbf{(P\_5,2.37)} dvigoḥ nityam . \\
\noindent \textbf{(P\_5,2.37)} ḍaṭ stome . \\
\noindent \textbf{(P\_5,2.37)} śacśanoḥ ḍiniḥ . \\
\noindent \textbf{(P\_5,2.37)} pramāṇaparimāṇābhyām saṅkhyāyāḥ ca saṃśaye (R IV.118)] \\
\noindent \textbf{(P\_5,2.39.1)} kimartham parimāṇe iti ucyate na pramāṇe iti vartate . \\
\noindent \textbf{(P\_5,2.39.1)} evam tarhi siddhe sati yat parimāṇagrahaṇam karoti tat jñāpayati ācāryaḥ anyat pramāṇam anyat parimāṇam iti . \\
\noindent \textbf{(P\_5,2.39.1)} ḍāvatau arthavaiśeṣyāt nirdeśaḥ pṛthak ucyate . \\
\noindent \textbf{(P\_5,2.39.1)} na etat jñāpakasādhyam anyat pramāṇam anyat parimāṇam iti . \\
\noindent \textbf{(P\_5,2.39.1)} uktaḥ atra viśeṣaḥ . \\
\noindent \textbf{(P\_5,2.39.1)} mātrādyapratighātāya . \\
\noindent \textbf{(P\_5,2.39.1)} evam ca kṛtvā mātrādīnām pratighātaḥ na bhavati . \\
\noindent \textbf{(P\_5,2.39.1)} bhāvaḥ siddhaḥ ca ḍāvatoḥ . \\
\noindent \textbf{(P\_5,2.39.1)} ḍāvatvantāt mātrajādīnām bhāvaḥ siddhaḥ bhavati . \\
\noindent \textbf{(P\_5,2.39.2)} vatupprakaraṇe yuṣmadasmadbhyām chandasi sādṛśe upasaṅkhyānam . \\
\noindent \textbf{(P\_5,2.39.2)} vatupprakaraṇe yuṣmadasmadbhyām chandasi sādṛśe upasaṅkhyānam kartavyam . \\
\noindent \textbf{(P\_5,2.39.2)} na tvāvān anyaḥ divyaḥ na parthivaḥ na jātaḥ na janiṣyate . \\
\noindent \textbf{(P\_5,2.39.2)} tvavataḥ purūvaso . \\
\noindent \textbf{(P\_5,2.39.2)} yajñam viprasya mavataḥ . \\
\noindent \textbf{(P\_5,2.39.2)} tvatsadṛśasya . \\
\noindent \textbf{(P\_5,2.39.2)} matsadṛśasya iti . \\
\noindent \textbf{(P\_5,2.39.2)} [ḍāvatau arthavaiśeṣyāt nirdeśaḥ pṛthak ucyate . \\
\noindent \textbf{(P\_5,2.39.2)} mātrādyapratighātāya . \\
\noindent \textbf{(P\_5,2.39.2)} bhāvaḥ siddhaḥ ca ḍāvatoḥ (R IV.120)] \\
\noindent \textbf{(P\_5,2.40)} kena vihitasya kimidambhyām vatupaḥ vaḥ ghatvam ucyate . \\
\noindent \textbf{(P\_5,2.40)} etat eva jñāpayati ācāryaḥ bhavati kimidambhyām vatup iti yat ayam kimidambhyām uttarasya vatupaḥ vaḥ ghatvam śāsti . \\
\noindent \textbf{(P\_5,2.40)} atha vā yogavibhāgaḥ kariṣyate . \\
\noindent \textbf{(P\_5,2.40)} kimidambhyām vatup bhavati . \\
\noindent \textbf{(P\_5,2.40)} tataḥ vaḥ ghaḥ iti . \\
\noindent \textbf{(P\_5,2.40)} vaḥ ca asya ghaḥ bhavati iti . \\
\noindent \textbf{(P\_5,2.41)} bahuṣu iti vaktavyam . \\
\noindent \textbf{(P\_5,2.41)} iha mā bhūt . \\
\noindent \textbf{(P\_5,2.41)} kiyān . \\
\noindent \textbf{(P\_5,2.41)} kiyantau . \\
\noindent \textbf{(P\_5,2.41)} tat tarhi vaktavyam . \\
\noindent \textbf{(P\_5,2.41)} na vaktavyam . \\
\noindent \textbf{(P\_5,2.41)} kim iti etat paripraśne vartate paripraśnaḥ ca anirjñāte anirjñātam ca bahuṣu . \\
\noindent \textbf{(P\_5,2.41)} dvyekayoḥ punaḥ nirjñātam . \\
\noindent \textbf{(P\_5,2.41)} nirjñātatvāt dvyekayoḥ paripraśnaḥ na bhavati . \\
\noindent \textbf{(P\_5,2.41)} paripraśnābhāvāt kim eva tāvat na asti kutaḥ pratyayaḥ . \\
\noindent \textbf{(P\_5,2.42)} iha kasmāt na bhavati . \\
\noindent \textbf{(P\_5,2.42)} bahavaḥ avayavāḥ asyāḥ saṅkhyāyāḥ iti . \\
\noindent \textbf{(P\_5,2.42)} avayave yā saṅkhyā iti ucyate . \\
\noindent \textbf{(P\_5,2.42)} na ca kā cit saṅkhyā asti yasyāḥ bahuśabdaḥ avayavaḥ syāt . \\
\noindent \textbf{(P\_5,2.42)} nanu ca iyam asti saṅkhyā iti eva . \\
\noindent \textbf{(P\_5,2.42)} na eṣā saṅkhyā . \\
\noindent \textbf{(P\_5,2.42)} sañjñā eṣā . \\
\noindent \textbf{(P\_5,2.42)} avayavavidhāne avayavini pratyayaḥ . \\
\noindent \textbf{(P\_5,2.42)} avayavavidhāne avayavini pratyayaḥ bhavati iti vaktavyam . \\
\noindent \textbf{(P\_5,2.42)} iha mā bhūt . \\
\noindent \textbf{(P\_5,2.42)} pañca avayavāḥ . \\
\noindent \textbf{(P\_5,2.42)} daśa avayavāḥ iti . \\
\noindent \textbf{(P\_5,2.42)} atha avayavini iti ucyamāne avayavasvāmini kasmāt na bhavati . \\
\noindent \textbf{(P\_5,2.42)} pañca paśvavayavāḥ devadattasya iti . \\
\noindent \textbf{(P\_5,2.42)} avayavaśabdaḥ ayam guṇaśabdaḥ asya iti ca vartate . \\
\noindent \textbf{(P\_5,2.42)} tena yam prati avayavaḥ guṇaḥ tasmin avayavini pratyayena bhavitavyam . \\
\noindent \textbf{(P\_5,2.42)} kam ca prati avayavaḥ guṇaḥ . \\
\noindent \textbf{(P\_5,2.42)} samudāyam . \\
\noindent \textbf{(P\_5,2.42)} yadi evam avayavini iti api na vaktavyam . \\
\noindent \textbf{(P\_5,2.42)} avayaveṣu kasmāt na bhavati . \\
\noindent \textbf{(P\_5,2.42)} asya iti vartate . \\
\noindent \textbf{(P\_5,2.44)} kimartham udāttaḥ iti ucyate . \\
\noindent \textbf{(P\_5,2.44)} udāttaḥ yathā syāt . \\
\noindent \textbf{(P\_5,2.44)} na etat asti prayojanam . \\
\noindent \textbf{(P\_5,2.44)} pratyayasvareṇa api eṣaḥ svaraḥ siddhaḥ . \\
\noindent \textbf{(P\_5,2.44)} na sidhyati . \\
\noindent \textbf{(P\_5,2.44)} citaḥ antaḥ udāttaḥ bhavati iti antodāttatvam prasajyeta . \\
\noindent \textbf{(P\_5,2.44)} atha udāttaḥ iti ucyamāne kutaḥ etat ādeḥ udāttatvam bhaviṣyati na punaḥ antasya iti . \\
\noindent \textbf{(P\_5,2.44)} udāttavacanasāmarthyāt yasya aprāptaḥ svaraḥ tasya bhaviṣyati . \\
\noindent \textbf{(P\_5,2.44)} kasya ca aprāptaḥ . \\
\noindent \textbf{(P\_5,2.44)} ādeḥ . \\
\noindent \textbf{(P\_5,2.44)} antasya punaḥ citsvareṇa eva siddham . \\
\noindent \textbf{(P\_5,2.45)} iha kasmāt na bhavati . \\
\noindent \textbf{(P\_5,2.45)} ekādaśa māṣāḥ adhikāḥ asmin kārṣāpaṇaśate iti . \\
\noindent \textbf{(P\_5,2.45)} adhike samānajātau . \\
\noindent \textbf{(P\_5,2.45)} samānajātau adhike iṣyate . \\
\noindent \textbf{(P\_5,2.45)} atha iha kasmāt na bhavati . \\
\noindent \textbf{(P\_5,2.45)} ekādaśa kārṣāpaṇāḥ adhikāḥ asyām kārṣāpaṇatriṃsati iti . \\
\noindent \textbf{(P\_5,2.45)} iṣṭam śatasahasrayoḥ . \\
\noindent \textbf{(P\_5,2.45)} śatasahasrayoḥ adhike iṣyate . \\
\noindent \textbf{(P\_5,2.45)} atha ekādaśam śatasahasram iti kasya ādhikye bhavitavyam . \\
\noindent \textbf{(P\_5,2.45)} yasya saṅkhyā tadādhikye ḍaḥ kartavyaḥ mataḥ mama . \\
\noindent \textbf{(P\_5,2.45)} yadi tāvat śatāni saṅkhyāyante śatādhikye bhavitavyam . \\
\noindent \textbf{(P\_5,2.45)} atha sahasrāṇi saṅkhyāyante sahasrādhikye bhavitavyam . \\
\noindent \textbf{(P\_5,2.45)} ḍavidhāne parimāṇaśabdānām ādhikyasya adhikaraṇābhāvāt anirdeśaḥ . \\
\noindent \textbf{(P\_5,2.45)} ḍavidhāne parimāṇaśabdānām ādhikyasya adhikaraṇābhāvāt anirdeśaḥ . \\
\noindent \textbf{(P\_5,2.45)} agamakaḥ nirdeśaḥ anirdeśaḥ . \\
\noindent \textbf{(P\_5,2.45)} na hi ekādaśānām śatam adhikaraṇam . \\
\noindent \textbf{(P\_5,2.45)} siddham tu pañcamīnirdeśāt . \\
\noindent \textbf{(P\_5,2.45)} siddham etat . \\
\noindent \textbf{(P\_5,2.45)} katham . \\
\noindent \textbf{(P\_5,2.45)} pañcamīnirdeśāt . \\
\noindent \textbf{(P\_5,2.45)} pañcamīnirdeśaḥ kartavyaḥ . \\
\noindent \textbf{(P\_5,2.45)} tat asmāt adhikam iti . \\
\noindent \textbf{(P\_5,2.45)} saḥ tarhi pañcamīnirdeśaḥ kartavyaḥ . \\
\noindent \textbf{(P\_5,2.45)} na kartavyaḥ . \\
\noindent \textbf{(P\_5,2.45)} yadi api tāvat vyāpake vaiṣayike vā adhikaraṇe sambhavaḥ na asti aupaśleṣikam adhikaraṇam vijñāsyate . \\
\noindent \textbf{(P\_5,2.45)} ekādaśa kārṣāpaṇāḥ upaśliṣṭāḥ asmin śate ekādaśam śatam . \\
\noindent \textbf{(P\_5,2.46)} kimartham śadgrahaṇe antagrahaṇam . \\
\noindent \textbf{(P\_5,2.46)} śadgrahaṇe antagrahaṇam pratyayagrahaṇe yasmāt saḥ tadādeḥ adhikārtham . \\
\noindent \textbf{(P\_5,2.46)} śadgrahaṇe antagrahaṇam kriyate . \\
\noindent \textbf{(P\_5,2.46)} pratyayagrahaṇe yasmāt saḥ vihitaḥ tadādeḥ tadantasya grahaṇam bhavati iti iha na prāpnoti : ekatriṃśam śatam . \\
\noindent \textbf{(P\_5,2.46)} iṣyate ca atra api syāt iti . \\
\noindent \textbf{(P\_5,2.46)} tat ca antareṇa yatnam na sidhyati iti antagrahaṇam . \\
\noindent \textbf{(P\_5,2.46)} evamartham idam ucyate . \\
\noindent \textbf{(P\_5,2.46)} asti prayojanam etat . \\
\noindent \textbf{(P\_5,2.46)} kim tarhi iti . \\
\noindent \textbf{(P\_5,2.46)} saṅkhyāgrahaṇam ca . \\
\noindent \textbf{(P\_5,2.46)} saṅkhyāgrahaṇam ca kartavyam . \\
\noindent \textbf{(P\_5,2.46)} iha mā bhūt . \\
\noindent \textbf{(P\_5,2.46)} gotriṃśat adhikarm asmin śate iti . \\
\noindent \textbf{(P\_5,2.46)} viṃśateḥ ca . \\
\noindent \textbf{(P\_5,2.46)} viṃśateḥ ca antagrahaṇam kartavyam . \\
\noindent \textbf{(P\_5,2.46)} iha api yathā syāt . \\
\noindent \textbf{(P\_5,2.46)} ekaviṃśam śatam . \\
\noindent \textbf{(P\_5,2.46)} cakārāt saṅkhyāgrahaṇam ca kartavyam . \\
\noindent \textbf{(P\_5,2.46)} iha mā bhūt . \\
\noindent \textbf{(P\_5,2.46)} goviṃśatiḥ adhikam asmin śate iti . \\
\noindent \textbf{(P\_5,2.47)} nimāne guṇini . \\
\noindent \textbf{(P\_5,2.47)} nimāne guṇini iti vaktavyam . \\
\noindent \textbf{(P\_5,2.47)} kim prayojanam . \\
\noindent \textbf{(P\_5,2.47)} guṇeṣu mā bhūt . \\
\noindent \textbf{(P\_5,2.47)} bhūyasaḥ . \\
\noindent \textbf{(P\_5,2.47)} bhūyasaḥ iti ca vaktavyam . \\
\noindent \textbf{(P\_5,2.47)} kim prayojanam . \\
\noindent \textbf{(P\_5,2.47)} bhūyasaḥ vācikāyāḥ saṅkhyāyāḥ utpattiḥ yathā syāt alpīyasaḥ vācikāyāḥ saṅkhyāyāḥ utpattiḥ mā bhūt iti . \\
\noindent \textbf{(P\_5,2.47)} ekaḥ anyataraḥ . \\
\noindent \textbf{(P\_5,2.47)} ekaḥ cet anyataraḥ bhavati iti vaktavyam . \\
\noindent \textbf{(P\_5,2.47)} iha mā bhūt . \\
\noindent \textbf{(P\_5,2.47)} dvau yavānām trayaḥ udaśvitaḥ iti . \\
\noindent \textbf{(P\_5,2.47)} samānānām . \\
\noindent \textbf{(P\_5,2.47)} samānānām ca iti vaktavyam . \\
\noindent \textbf{(P\_5,2.47)} iha mā bhūt . \\
\noindent \textbf{(P\_5,2.47)} ekaḥ yavānām adhyardhaḥ udaśvitaḥ iti . \\
\noindent \textbf{(P\_5,2.47)} tat tarhi bahu vaktavyam . \\
\noindent \textbf{(P\_5,2.47)} na vaktavyam . \\
\noindent \textbf{(P\_5,2.47)} yat tāvat ucyate . \\
\noindent \textbf{(P\_5,2.47)} guṇini iti vaktavyam iti . \\
\noindent \textbf{(P\_5,2.47)} na vaktavyam . \\
\noindent \textbf{(P\_5,2.47)} guṇeṣu kasmāt na bhavati . \\
\noindent \textbf{(P\_5,2.47)} asya iti vartate . \\
\noindent \textbf{(P\_5,2.47)} yat uktam bhūyasaḥ iti vaktavyam iti . \\
\noindent \textbf{(P\_5,2.47)} na vaktavyam . \\
\noindent \textbf{(P\_5,2.47)} alpīyasaḥ vācikāyāḥ saṅkhyāyāḥ utpattiḥ kasmāt na bhavati . \\
\noindent \textbf{(P\_5,2.47)} anabhidhānāt . \\
\noindent \textbf{(P\_5,2.47)} yat uktam ekaḥ cet anyataraḥ bhavati iti vaktavyam iti . \\
\noindent \textbf{(P\_5,2.47)} na vaktavyam . \\
\noindent \textbf{(P\_5,2.47)} kasmāt na bhavati dvau yavānām trayaḥ udaśvitaḥ iti . \\
\noindent \textbf{(P\_5,2.47)} tantram vibhaktinirdeśaḥ . \\
\noindent \textbf{(P\_5,2.47)} yat api ucyate samānānām ca iti vaktavyam iti . \\
\noindent \textbf{(P\_5,2.47)} na vaktavyam . \\
\noindent \textbf{(P\_5,2.47)} kasmāt na bhavati ekaḥ yavānām adhyardhaḥ udaśvitaḥ iti . \\
\noindent \textbf{(P\_5,2.47)} anabhidhānāt . \\
\noindent \textbf{(P\_5,2.47)} nimeye ca api dṛśyate . \\
\noindent \textbf{(P\_5,2.47)} nimeye ca api pratyayaḥ dṛśyate . \\
\noindent \textbf{(P\_5,2.47)} dvimayāḥ yavāḥ . \\
\noindent \textbf{(P\_5,2.47)} trimayāḥ . \\
\noindent \textbf{(P\_5,2.47)} kim punaḥ iha nimānam kim nimeyam yāvatā ubhayam tyajyate . \\
\noindent \textbf{(P\_5,2.47)} satyam evam etat . \\
\noindent \textbf{(P\_5,2.47)} kva cit tu kā cit prasṛtatarā gatiḥ bhavati . \\
\noindent \textbf{(P\_5,2.47)} tat yathā . \\
\noindent \textbf{(P\_5,2.47)} samāne tyāge dhānyam vikrīṇite yavān vikrīṇīte iti ucyate . \\
\noindent \textbf{(P\_5,2.47)} na kaḥ cit āha kārṣāpaṇam vikrīṇite iti . \\
\noindent \textbf{(P\_5,2.47)} atha vā yena adhigamyate tat nimānam . \\
\noindent \textbf{(P\_5,2.47)} yat adhimayate tat nimeyam . \\
\noindent \textbf{(P\_5,2.48)} tasya pūraṇe iti atiprasaṅgaḥ . \\
\noindent \textbf{(P\_5,2.48)} tasya pūraṇe iti atiprasaṅgaḥ bhavati . \\
\noindent \textbf{(P\_5,2.48)} iha api prapnoti . \\
\noindent \textbf{(P\_5,2.48)} pañcānām uṣṭrikāṇām pūraṇaḥ ghaṭaḥ . \\
\noindent \textbf{(P\_5,2.48)} siddham tu saṅkhyāpūraṇe iti vacanāt . \\
\noindent \textbf{(P\_5,2.48)} siddham etat . \\
\noindent \textbf{(P\_5,2.48)} katham . \\
\noindent \textbf{(P\_5,2.48)} saṅkhyāpūraṇe iti vaktavyam . \\
\noindent \textbf{(P\_5,2.48)} evam api ghaṭe prāpnoti . \\
\noindent \textbf{(P\_5,2.48)} saṅkhyeyam hi asau adbhiḥ pūrayati . \\
\noindent \textbf{(P\_5,2.48)} saṅkhyāpūraṇe iti brūmaḥ na saṅkhyeyapūraṇe iti . \\
\noindent \textbf{(P\_5,2.48)} yasya vā bhāvāt anyasaṅkhyātvam tatra . \\
\noindent \textbf{(P\_5,2.48)} atha vā yasya bhāvāt anyā saṅkhyā pravartate tatra iti vaktavyam . \\
\noindent \textbf{(P\_5,2.48)} evam api dvitīye adhyāye aṣṭamaḥ iti prāpnoti . \\
\noindent \textbf{(P\_5,2.48)} sarveṣām hi teṣām bhāvāt saṅkhyā pravartate . \\
\noindent \textbf{(P\_5,2.48)} caramopajāte pūrvasmin ca anapagate iti vaktavyam . \\
\noindent \textbf{(P\_5,2.48)} evam api ekādaśīdvādaśyau sauviṣṭakṛtī . \\
\noindent \textbf{(P\_5,2.48)} idam dvitīyam idam tṛtīyam . \\
\noindent \textbf{(P\_5,2.48)} daśa daśamāni iti na sidhyati . \\
\noindent \textbf{(P\_5,2.48)} sūtram ca bhidyate . \\
\noindent \textbf{(P\_5,2.48)} yathānyāsam eva astu . \\
\noindent \textbf{(P\_5,2.48)} nanu ca uktam tasya pūraṇe iti atiprasaṅgaḥ iti . \\
\noindent \textbf{(P\_5,2.48)} parihṛtam etat siddham saṅkhyāpūraṇe iti vacanāt iti . \\
\noindent \textbf{(P\_5,2.48)} tat tarhi saṅkhyāgrahaṇam kartavyam . \\
\noindent \textbf{(P\_5,2.48)} na kartavyam . \\
\noindent \textbf{(P\_5,2.48)} prakṛtam anuvartate . \\
\noindent \textbf{(P\_5,2.48)} kva prakṛtam . \\
\noindent \textbf{(P\_5,2.48)} saṅkhyāyāḥ guṇasya nimāne mayaṭ iti . \\
\noindent \textbf{(P\_5,2.48)} evam tarhi na iyam vṛttiḥ upālabhyate . \\
\noindent \textbf{(P\_5,2.48)} kim tarhi . \\
\noindent \textbf{(P\_5,2.48)} vṛttisthānam upālabhyate . \\
\noindent \textbf{(P\_5,2.48)} vṛttiḥ eva atra na prāpnoti . \\
\noindent \textbf{(P\_5,2.48)} kim kāraṇam . \\
\noindent \textbf{(P\_5,2.48)} pratyayārthābhāvāt . \\
\noindent \textbf{(P\_5,2.48)} na eṣaḥ doṣaḥ . \\
\noindent \textbf{(P\_5,2.48)} vacanāt svāṛthikaḥ bhaviṣyati . \\
\noindent \textbf{(P\_5,2.48)} atha vā pūrvasyāḥ saṅkhyāyāḥ parāpekṣāyāḥ utpattiḥ vaktavyā uttarā ca sāṅkhyā ādeśaḥ vaktavyaḥ . \\
\noindent \textbf{(P\_5,2.48)} atha vā nyūne ayam kṛtsnaśabdaḥ draṣṭavyaḥ : caturṣu pañcaśabdaḥ . \\
\noindent \textbf{(P\_5,2.48)} atha vā sarve eva dvyādayaḥ anyonyam apekṣante . \\
\noindent \textbf{(P\_5,2.48)} yadi evam dvitīye adhyāye aṣṭamaḥ iti prāpnoti . \\
\noindent \textbf{(P\_5,2.48)} bhavati eva . \\
\noindent \textbf{(P\_5,2.48)} prakṛtyarthāt bahiḥ sarvā vṛttiḥ prāyeṇa lakṣyate . \\
\noindent \textbf{(P\_5,2.48)} pūraṇe syāt katham vṛttiḥ . \\
\noindent \textbf{(P\_5,2.48)} vacanāt iti lakṣyatām . \\
\noindent \textbf{(P\_5,2.48)} tasyāḥ pūrvā tu yā saṅkhyā tasyāḥ [R tasyām ] bhavatu taddhitaḥ . \\
\noindent \textbf{(P\_5,2.48)} ādeśaḥ ca ottarā saṅkhyā tathā nyāyyā bhaviṣyati . \\
\noindent \textbf{(P\_5,2.48)} nyūne vā kṛtsnaśabdaḥ ayam pūrvasyām uttarām yadi sāmarthyam ca tayā tasyāḥ tathā nyāyyā bhaviṣyati . \\
\noindent \textbf{(P\_5,2.48)} anyonyam vā vyapāśritya sarvasmin dvyādayaḥ yadi pravartante tathā nyāyyā vṛttiḥ bhavati pūraṇe . \\
\noindent \textbf{(P\_5,2.48)} bahūnām vācikā saṅkhyā pūraṇaḥ ca ekaḥ iṣyate . \\
\noindent \textbf{(P\_5,2.48)} anyatvāt ubhayoḥ nyāyyā vārkṣī śākhā nidarśanam . \\
\noindent \textbf{(P\_5,2.49)} maḍādiṣu yasya ādiḥ tannirderdeśaḥ . \\
\noindent \textbf{(P\_5,2.49)} maḍādiṣu yasya ādiḥ kriyate tannirderdeśaḥ kartavyaḥ . \\
\noindent \textbf{(P\_5,2.49)} asya ādiḥ bhavati iti vaktavyam . \\
\noindent \textbf{(P\_5,2.49)} akiryamāṇe hi pratyayādhikārāt pratyayaḥ ayam vijñāyeta . \\
\noindent \textbf{(P\_5,2.49)} tatra kaḥ doṣaḥ . \\
\noindent \textbf{(P\_5,2.49)} pratyayāntare hi svare doṣaḥ . \\
\noindent \textbf{(P\_5,2.49)} pratyayāntare hi sati svare doṣaḥ syāt . \\
\noindent \textbf{(P\_5,2.49)} viṃśatitamaḥ . \\
\noindent \textbf{(P\_5,2.49)} eṣaḥ svaraḥ prasajyeta . \\
\noindent \textbf{(P\_5,2.49)} viṃśatitamaḥ iti ca iṣyate . \\
\noindent \textbf{(P\_5,2.49)} saḥ tarhi tathā nirdeśaḥ kartavyaḥ . \\
\noindent \textbf{(P\_5,2.49)} na kartavyaḥ . \\
\noindent \textbf{(P\_5,2.49)} prakṛtam ḍaḍgrahaṇam anuvartate . \\
\noindent \textbf{(P\_5,2.49)} kva prakṛtam . \\
\noindent \textbf{(P\_5,2.49)} tasya pūraṇe ḍaṭ iti . \\
\noindent \textbf{(P\_5,2.49)} tat vai prathamānirdiṣṭam ṣaṣṭhīnirdiṣṭena ca arthaḥ . \\
\noindent \textbf{(P\_5,2.49)} nāntāt iti pañcamī ḍaṭ iti prathamāyāḥ ṣaṣṭhīm prakalpayiṣyati . \\
\noindent \textbf{(P\_5,2.49)} tasmāt iti uttarasya iti . \\
\noindent \textbf{(P\_5,2.49)} pratyayavidhiḥ ayam na ca pratyayavidhau pañcamyaḥ prakalpikāḥ bhavanti . \\
\noindent \textbf{(P\_5,2.49)} na ayam pratyayavidhiḥ . \\
\noindent \textbf{(P\_5,2.49)} vihitaḥ pratyayaḥ prakṛtaḥ ca anuvartate . \\
\noindent \textbf{(P\_5,2.51.1)} caturaḥ chayatau ādyakṣaralopaḥ ca . \\
\noindent \textbf{(P\_5,2.51.1)} caturaḥ chayatau vaktavyau ādyakṣaralopaḥ ca vaktavyaḥ . \\
\noindent \textbf{(P\_5,2.51.1)} turīyam . \\
\noindent \textbf{(P\_5,2.51.1)} turyam . \\
\noindent \textbf{(P\_5,2.51.2)} atha kimartham thaṭthukau pṛthak kriyete na sarvam thaṭ eva syāt thuk eva vā . \\
\noindent \textbf{(P\_5,2.51.2)} thaṭthukoḥ pṛthakkaraṇam padāntavidhipratiṣedhāṛtham . \\
\noindent \textbf{(P\_5,2.51.2)} thaṭthukoḥ pṛthakkaraṇam kriyate padāntavidhipratiṣedhāṛtham . \\
\noindent \textbf{(P\_5,2.51.2)} padāntavidhyartham padāntapratiṣedhāṛtham ca . \\
\noindent \textbf{(P\_5,2.51.2)} padāntavidhyartham tāvat . \\
\noindent \textbf{(P\_5,2.51.2)} parṇamayāni pañcathāni bhavanti . \\
\noindent \textbf{(P\_5,2.51.2)} rathaḥ saptathaḥ . \\
\noindent \textbf{(P\_5,2.51.2)} padantasya iti nalopaḥ yathā syāt . \\
\noindent \textbf{(P\_5,2.51.2)} padāntapratiṣedhāṛtham . \\
\noindent \textbf{(P\_5,2.51.2)} ṣaṣṭhaḥ . \\
\noindent \textbf{(P\_5,2.51.2)} padāntasya iti jaśtvam mā bhūt . \\
\noindent \textbf{(P\_5,2.51.2)} iha caturthaḥ iti padāntasya iti visarjanīyaḥ mā bhūt iti . \\
\noindent \textbf{(P\_5,2.52)} bahukatipayavatūnām liṅgaviśiṣṭāt utpattiḥ . \\
\noindent \textbf{(P\_5,2.52)} bahukatipayavatūnām liṅgaviśiṣṭād utpattiḥ vaktavyā . \\
\noindent \textbf{(P\_5,2.52)} iha api yathā syāt . \\
\noindent \textbf{(P\_5,2.52)} bahvīnām pūraṇī bahutithī . \\
\noindent \textbf{(P\_5,2.52)} katipayānām pūraṇī katipayathī . \\
\noindent \textbf{(P\_5,2.52)} tāvatīnām pūraṇī tāvatithī . \\
\noindent \textbf{(P\_5,2.52)} bahukatipayavatūnām liṅgaviśiṣṭād utpattiḥ siddhā . \\
\noindent \textbf{(P\_5,2.52)} katham . \\
\noindent \textbf{(P\_5,2.52)} prātipadikagrahaṇe liṅgaviśiṣṭasya api grahaṇam bhavati iti . \\
\noindent \textbf{(P\_5,2.52)} puṃvadvacanam ca . puṃvadbhāvaḥ ca vaktavyaḥ . \\
\noindent \textbf{(P\_5,2.52)} bahvīnām pūraṇī bahutithī . \\
\noindent \textbf{(P\_5,2.52)} kimartham na bhasya aḍhe taddhite puṃvat bhavati iti siddham . \\
\noindent \textbf{(P\_5,2.52)} bhasya iti ucyate . \\
\noindent \textbf{(P\_5,2.52)} yajādau ca bham bhavati . \\
\noindent \textbf{(P\_5,2.52)} na ca atra yajādim paśyāmaḥ . \\
\noindent \textbf{(P\_5,2.52)} kim kāraṇam . \\
\noindent \textbf{(P\_5,2.52)} tithukā vyavihitatvāt na prāpnoti . \\
\noindent \textbf{(P\_5,2.52)} idam iha sampradhāryam . \\
\noindent \textbf{(P\_5,2.52)} tithuk kriyatām puṃvadbhāvaḥ iti . \\
\noindent \textbf{(P\_5,2.52)} kim atra kartavyam . \\
\noindent \textbf{(P\_5,2.52)} paratvāt puṃvadbhāvaḥ . \\
\noindent \textbf{(P\_5,2.52)} nityaḥ tithuk . \\
\noindent \textbf{(P\_5,2.52)} kṛte api puṃvadbhāve prāpnoti akṛte api prāpnoti . \\
\noindent \textbf{(P\_5,2.52)} tithuk api anityaḥ . \\
\noindent \textbf{(P\_5,2.52)} anyasya kṛte puṃvadbhāve prāpnoti anyasya akṛte . \\
\noindent \textbf{(P\_5,2.52)} śabdāntarasya ca prāpnuvan vidhiḥ anityaḥ bhavati . \\
\noindent \textbf{(P\_5,2.52)} antaraṅgaḥ tarhi tithuk . \\
\noindent \textbf{(P\_5,2.52)} kā antaraṅgatā . \\
\noindent \textbf{(P\_5,2.52)} utpattisanniyogena tithuk ucyate utpanne pratyayte prakṛtipratyayau āśritya puṃvadbhāvaḥ . \\
\noindent \textbf{(P\_5,2.52)} puṃvadbhāvaḥ api antaraṅgaḥ . \\
\noindent \textbf{(P\_5,2.52)} katham . \\
\noindent \textbf{(P\_5,2.52)} uktam etat siddhaḥ ca pratyayavidhau iti . \\
\noindent \textbf{(P\_5,2.52)} ubhayoḥ antaraṅgayoḥ paratvāt puṃvadbhāvaḥ . \\
\noindent \textbf{(P\_5,2.52)} puṃvadbhāve kṛte punaḥprasaṅgavijñānāt tithuk siddhaḥ : bahutithī . \\
\noindent \textbf{(P\_5,2.58)} asaṅkhyādeḥ iti kimartham . \\
\noindent \textbf{(P\_5,2.58)} iha mā bhūt . \\
\noindent \textbf{(P\_5,2.58)} ekaṣaṣṭaḥ . \\
\noindent \textbf{(P\_5,2.58)} dviṣaṣṭaḥ . \\
\noindent \textbf{(P\_5,2.58)} asaṅkhyādeḥ iti śakyam avaktum . \\
\noindent \textbf{(P\_5,2.58)} kasmāt na bhavati ekaṣaṣṭaḥ . \\
\noindent \textbf{(P\_5,2.58)} dviṣaṣṭaḥ iti . \\
\noindent \textbf{(P\_5,2.58)} ṣaṣṭiśabdāt pratyayaḥ vidhīyate . \\
\noindent \textbf{(P\_5,2.58)} kaḥ prasaṅgaḥ yat ekaṣaṣṭiśabdāt syāt . \\
\noindent \textbf{(P\_5,2.58)} na eva prāpnoti na arthaḥ pratiṣedhena . \\
\noindent \textbf{(P\_5,2.58)} tadantavidhinā prāpnoti . \\
\noindent \textbf{(P\_5,2.58)} grahaṇavatā prātipadikena tadantavidhiḥ pratiṣidhyate . \\
\noindent \textbf{(P\_5,2.58)} evam tarhi jñāpayati ācāryaḥ bhavati iha tadantavidhiḥ iti . \\
\noindent \textbf{(P\_5,2.58)} kim etasya jñāpane prayojanam . \\
\noindent \textbf{(P\_5,2.58)} ekaviṃśatitamaḥ . \\
\noindent \textbf{(P\_5,2.58)} etat siddham bhavati . \\
\noindent \textbf{(P\_5,2.59)} chaprakaraṇe anekapadāt api . \\
\noindent \textbf{(P\_5,2.59)} chaprakaraṇe anekapadāt api iti vaktavyam . \\
\noindent \textbf{(P\_5,2.59)} iha api yathā syāt : asyavāmīyam , kayāśubhīyam . \\
\noindent \textbf{(P\_5,2.59)} kim punaḥ kāraṇam na sidhyati . \\
\noindent \textbf{(P\_5,2.59)} aprātipadikatvāt . \\
\noindent \textbf{(P\_5,2.59)} siddham tu prātipadikavijñānāt . \\
\noindent \textbf{(P\_5,2.59)} siddham etat . \\
\noindent \textbf{(P\_5,2.59)} katham . \\
\noindent \textbf{(P\_5,2.59)} prātipadikavijñānāt . \\
\noindent \textbf{(P\_5,2.59)} katham prātipadikavijñānam . \\
\noindent \textbf{(P\_5,2.59)} svam rūpam śabdasya aśabdasañjñā iti vacanāt . \\
\noindent \textbf{(P\_5,2.59)} svam rūpam śabdasya aśabdasañjñā bhavati iti . \\
\noindent \textbf{(P\_5,2.59)} evam yaḥ asau āmnāye asyavāmaśabdaḥ paṭhyate saḥ asya padārthaḥ . \\
\noindent \textbf{(P\_5,2.59)} kim punaḥ anye āmnāyaśabdāḥ anye ime . \\
\noindent \textbf{(P\_5,2.59)} om iti āha . \\
\noindent \textbf{(P\_5,2.59)} kutaḥ etat . \\
\noindent \textbf{(P\_5,2.59)} āmnāyaśabdānām anyabhāvyam svaravarṇānupūrvīdeśakālaniyatatvāt . \\
\noindent \textbf{(P\_5,2.59)} svaraḥ niyataḥ āmnāye asyavāmaśabdasya . \\
\noindent \textbf{(P\_5,2.59)} varṇānupūrvī khalu api āmnāye niyatā asyavāmaśabdasya . \\
\noindent \textbf{(P\_5,2.59)} deśaḥ khalu api āmnāye niyataḥ . \\
\noindent \textbf{(P\_5,2.59)} śmaśāne na adhyeyam . \\
\noindent \textbf{(P\_5,2.59)} catuṣpathe na adhyeyam iti . \\
\noindent \textbf{(P\_5,2.59)} kālaḥ khalu api āmnāye niyataḥ . \\
\noindent \textbf{(P\_5,2.59)} na amāvāsyāyām na caturdaśāyām iti . \\
\noindent \textbf{(P\_5,2.59)} padaikadeśasubalopadarśanāt ca . \\
\noindent \textbf{(P\_5,2.59)} padaikadeśaḥ khalu api āmnāye dṛśyate . \\
\noindent \textbf{(P\_5,2.59)} asyavāmīyam . \\
\noindent \textbf{(P\_5,2.59)} nanu ca eṣaḥ sublopaḥ syāt . \\
\noindent \textbf{(P\_5,2.59)} subalopadarśanāt ca . \\
\noindent \textbf{(P\_5,2.59)} subalopaḥ khalu api dṛśyate . \\
\noindent \textbf{(P\_5,2.59)} asyavāmīyam iti . \\
\noindent \textbf{(P\_5,2.59)} yadi tarhi anye āmnāyaśabdāḥ anye ime matvarthaḥ na upapadyate . \\
\noindent \textbf{(P\_5,2.59)} asyavāmaśabdaḥ asmin asti iti . \\
\noindent \textbf{(P\_5,2.59)} na sañjñā sañjñinam vyabhicarati . \\
\noindent \textbf{(P\_5,2.60)} adhyāyānuvākābhyām vā luk . \\
\noindent \textbf{(P\_5,2.60)} adhyāyānuvākābhyām vā luk vaktavyaḥ . \\
\noindent \textbf{(P\_5,2.60)} stambhaḥ . \\
\noindent \textbf{(P\_5,2.60)} stambhīyaḥ . \\
\noindent \textbf{(P\_5,2.60)} gardabhāṇḍaḥ . \\
\noindent \textbf{(P\_5,2.60)} gardabhāṇḍīyaḥ . \\
\noindent \textbf{(P\_5,2.60)} anukaḥ . \\
\noindent \textbf{(P\_5,2.60)} anukīyaḥ . \\
\noindent \textbf{(P\_5,2.65)} dhanahiraṇyāt kāmābhidhāne . \\
\noindent \textbf{(P\_5,2.65)} dhanahiraṇyāt kāmābhidhāne iti vaktavyam . \\
\noindent \textbf{(P\_5,2.65)} ṣaṣthyarthe hi aniṣṭaprasaṅgaḥ . \\
\noindent \textbf{(P\_5,2.65)} ṣaṣthyarthe hi sati aniṣṭaḥ prāpnoti . \\
\noindent \textbf{(P\_5,2.65)} dhane kāmaḥ asya iti . \\
\noindent \textbf{(P\_5,2.65)} tat tarhi vaktavyam . \\
\noindent \textbf{(P\_5,2.65)} na vaktavyam . \\
\noindent \textbf{(P\_5,2.65)} kasmāt na bhavati dhane kāmaḥ asya iti . \\
\noindent \textbf{(P\_5,2.65)} anabhidhānāt . \\
\noindent \textbf{(P\_5,2.72)} kim yaḥ śītam karoti saḥ śītakaḥ yaḥ vā uṣṇam karoti sa uṣṇakaḥ . \\
\noindent \textbf{(P\_5,2.72)} kim ca ataḥ . \\
\noindent \textbf{(P\_5,2.72)} tuṣāre āditye ca prāpnoti . \\
\noindent \textbf{(P\_5,2.72)} evam tarhi uttarapadalopaḥ atra draṣṭavyaḥ . \\
\noindent \textbf{(P\_5,2.72)} śītam iva śītam . \\
\noindent \textbf{(P\_5,2.72)} uṣṇam iva uṣṇam . \\
\noindent \textbf{(P\_5,2.72)} yaḥ āśu kartavyān arthān cireṇa karoti saḥ ucyate śītakaḥ iti . \\
\noindent \textbf{(P\_5,2.72)} yaḥ punaḥ āśu kartavyān arthān āśu eva karoti saḥ ucyate uṣṇakaḥ iti . \\
\noindent \textbf{(P\_5,2.73)} adhikam iti kim nipātyate . \\
\noindent \textbf{(P\_5,2.73)} adyārūḍhasya uttarapadalopaḥ ca kan ca pratyayaḥ . \\
\noindent \textbf{(P\_5,2.73)} adhyārūḍham adhikam iti . \\
\noindent \textbf{(P\_5,2.73)} bhavet siddham adhyārūḍhaḥ droṇaḥ khāryām adhikaḥ droṇaḥ khāryām iti . \\
\noindent \textbf{(P\_5,2.73)} idam tu na sidhyati . \\
\noindent \textbf{(P\_5,2.73)} adhyārūḍhā droṇena khārī . \\
\noindent \textbf{(P\_5,2.73)} adhikā droṇena khārī iti . \\
\noindent \textbf{(P\_5,2.73)} gatyarthānām hi ktaḥ kartari vidhīyate . \\
\noindent \textbf{(P\_5,2.73)} gatyarthānām vai ktaḥ karmaṇi api vidhīyate . \\
\noindent \textbf{(P\_5,2.75)} kim yaḥ pārśvena anvicchati saḥ pārvśvakaḥ . \\
\noindent \textbf{(P\_5,2.75)} kim ca ataḥ . \\
\noindent \textbf{(P\_5,2.75)} rājapuruṣe prāpnoti . \\
\noindent \textbf{(P\_5,2.75)} evam tarhi uttarapadalopaḥ atra draṣṭavyaḥ . \\
\noindent \textbf{(P\_5,2.75)} pārśvam iva pārśvam . \\
\noindent \textbf{(P\_5,2.75)} yaḥ ṛjunā upāyena anveṣṭavyān arthān anṛjunā upāyena anvicchati saḥ ucyate pārśvakaḥ iti . \\
\noindent \textbf{(P\_5,2.76)} kim yaḥ ayaḥśulena anvicchati saḥ āyaḥśūlikaḥ . \\
\noindent \textbf{(P\_5,2.76)} kim ca ataḥ . \\
\noindent \textbf{(P\_5,2.76)} śivabhāgavate prāpnoti . \\
\noindent \textbf{(P\_5,2.76)} evam tarhi uttarapadalopaḥ atra draṣṭavyaḥ . \\
\noindent \textbf{(P\_5,2.76)} ayaḥśulam iva ayaḥśulam . \\
\noindent \textbf{(P\_5,2.76)} yaḥ mṛdunā upāyena anveṣṭavyān arthān rabhasena upāyena anvicchati saḥ ucyate āyaḥśūlikaḥ iti . \\
\noindent \textbf{(P\_5,2.77)} tāvatitham grahaṇam iti luk vāvacanānarthakyam vibhāṣāprakaraṇāt . \\
\noindent \textbf{(P\_5,2.77)} tāvatitham grahaṇam iti luk vāvacanam anarthakam . \\
\noindent \textbf{(P\_5,2.77)} kim kāraṇam . \\
\noindent \textbf{(P\_5,2.77)} vibhāṣāprakaraṇāt . \\
\noindent \textbf{(P\_5,2.77)} prakṛtā mahāvibhāṣā . \\
\noindent \textbf{(P\_5,2.77)} tayā etat siddham . \\
\noindent \textbf{(P\_5,2.77)} tāvatithena gṛhṇāti iti luk ca . \\
\noindent \textbf{(P\_5,2.77)} tāvatithena gṛhṇāti iti upasaṅkhyānam kartavya luk ca vaktavyaḥ . \\
\noindent \textbf{(P\_5,2.77)} ṣaṣthena gṛhṇāti ṣaṭkaḥ . \\
\noindent \textbf{(P\_5,2.79)} śṛṅkhalam asya bandhanam karabhe iti anirdeśaḥ . \\
\noindent \textbf{(P\_5,2.79)} śṛṅkhalam asya bandhanam karabhe iti anirdeśaḥ . \\
\noindent \textbf{(P\_5,2.79)} agamakaḥ nirdeśaḥ anirdeśaḥ . \\
\noindent \textbf{(P\_5,2.79)} na hi tasya śṛṅkhalabandhanam . \\
\noindent \textbf{(P\_5,2.79)} śṛṅkhalavatyā asau rajjvā badhyate . \\
\noindent \textbf{(P\_5,2.79)} siddham tu tadvannirdeśāt luk ca . \\
\noindent \textbf{(P\_5,2.79)} siddham etat . \\
\noindent \textbf{(P\_5,2.79)} katham . \\
\noindent \textbf{(P\_5,2.79)} tadvannirdeśaḥ kartavyaḥ luk ca vaktavyaḥ . \\
\noindent \textbf{(P\_5,2.79)} śṛṅkhalavat bandhanam iti . \\
\noindent \textbf{(P\_5,2.79)} saḥ tarhi tadvannirdeśaḥ kartavyaḥ . \\
\noindent \textbf{(P\_5,2.79)} na kartavyaḥ . \\
\noindent \textbf{(P\_5,2.79)} iha yat na antareṇa yasya pravṛttiḥ bhavati tat tasya nimittatvāya kalpate . \\
\noindent \textbf{(P\_5,2.79)} na ca antareṇa śṛṅkhalam bandhanam pravartate . \\
\noindent \textbf{(P\_5,2.79)} atha vā sāhacaryāt tācchabdyam bhaviṣyati . \\
\noindent \textbf{(P\_5,2.79)} śṛṅkhalasahacaritam bandhanam śrṅkhalam bandhanam iti . \\
\noindent \textbf{(P\_5,2.82)} prāye sañjñāyām vaṭakebhyaḥ iniḥ . \\
\noindent \textbf{(P\_5,2.82)} prāye sañjñāyām vaṭakebhyaḥ iniḥ vaktavyaḥ . \\
\noindent \textbf{(P\_5,2.82)} vaṭakinī paurṇamāsī . \\
\noindent \textbf{(P\_5,2.84)} kim nipātyate . \\
\noindent \textbf{(P\_5,2.84)} śrotriyan chandaḥ adhīte iti vākyārthe padavacanam . \\
\noindent \textbf{(P\_5,2.84)} chandaḥ adhīte iti asya vākyasya arthe śrotriyan iti etat padam nipātyate . \\
\noindent \textbf{(P\_5,2.84)} chandasaḥ vā śrtotrabhāvaḥ tat adhīte iti ghan ca . \\
\noindent \textbf{(P\_5,2.84)} chandasaḥ vā śrtotrabhāvaḥ nipātyate tat adhīte iti etasmin arthe ghan ca pratyayaḥ . \\
\noindent \textbf{(P\_5,2.84)} chandaḥ adhīte śṛotriyaḥ . \\
\noindent \textbf{(P\_5,2.85)} iniṭhanoḥ samānakālagrahaṇam . \\
\noindent \textbf{(P\_5,2.85)} iniṭhanoḥ samānakālagrahaṇam kartavyam . \\
\noindent \textbf{(P\_5,2.85)} adya bhukte śraḥ śrāddhikaḥ iti mā bhūt . \\
\noindent \textbf{(P\_5,2.85)} uktam vā . \\
\noindent \textbf{(P\_5,2.85)} kim uktam . \\
\noindent \textbf{(P\_5,2.85)} anabhidhānāt iti . \\
\noindent \textbf{(P\_5,2.91)} sañjñāyām iti kimartham . \\
\noindent \textbf{(P\_5,2.91)} tribhiḥ sākṣāt dṛṣṭam bhavati yaḥ ca dadāti yasmai ca dīyate yaḥ ca upadraṣṭā . \\
\noindent \textbf{(P\_5,2.91)} tatra sarvatra pratyayaḥ prāpnoti . \\
\noindent \textbf{(P\_5,2.91)} sañjñāgrahaṇasāmarthyāt dhanikāntevāsinoḥ na bhavati . \\
\noindent \textbf{(P\_5,2.92)} kim nipātyate . \\
\noindent \textbf{(P\_5,2.92)} kṣetriyaḥ śrotriyavat . kṣetriyaḥ śrotriyavat nipātyate . \\
\noindent \textbf{(P\_5,2.92)} parakṣetre cikitsyaḥ iti etasya vākyasya arthe kṣetriyac iti etat padam nipātyate . \\
\noindent \textbf{(P\_5,2.92)} parakṣetrāt vā tatra cikitsyaḥ iti paralopaḥ ghac ca . \\
\noindent \textbf{(P\_5,2.92)} parakṣetrāt vā tatra cikitsyaḥ iti etasmin arthe paralopaḥ nipātyate ghac ca . \\
\noindent \textbf{(P\_5,2.92)} parakṣetre cikitsyaḥ kṣetriyaḥ . \\
\noindent \textbf{(P\_5,2.94.1)} kimartham imau arthau ubhau nirdiśyete : asya asmin iti na yat yasya bhavati tasmin api tat bhavati yat ca yasmin bhavati tat tasya api bhavati . \\
\noindent \textbf{(P\_5,2.94.1)} na etayoḥ āvaśyakaḥ samāveśaḥ . \\
\noindent \textbf{(P\_5,2.94.1)} bhavanti hi devadattasya gāvaḥ na ca tāḥ tasmin ādhṛtāḥ bhavanti . \\
\noindent \textbf{(P\_5,2.94.1)} bhavanti ca parvate vṛkṣāḥ na ca te tasya bhavanti . \\
\noindent \textbf{(P\_5,2.94.1)} atha astigrahaṇam kimartham . \\
\noindent \textbf{(P\_5,2.94.1)} sattāyām arthe pratyayaḥ yathā syāt . \\
\noindent \textbf{(P\_5,2.94.1)} na etat asti prayojanam . \\
\noindent \textbf{(P\_5,2.94.1)} na sattām padārthaḥ vyabhicarati . \\
\noindent \textbf{(P\_5,2.94.1)} idam tarhi prayojanam : sampratisattāyām yathā syāt , bhūtabhaviṣyatsattāyām mā bhūt : gāvaḥ asya āsan . \\
\noindent \textbf{(P\_5,2.94.1)} gāvaḥ asya bhavitāraḥ iti . \\
\noindent \textbf{(P\_5,2.94.1)} na tarhi idānīm idam bhavati : gomān āsīt . \\
\noindent \textbf{(P\_5,2.94.1)} gomān bhavitā iti . \\
\noindent \textbf{(P\_5,2.94.1)} bhavati na tu etasmin vākye . \\
\noindent \textbf{(P\_5,2.94.1)} yadi etasmin vākye syāt yathā iha asteḥ prayogaḥ na bhavati gomān yavamān iti evam iha api na syāt : gomān āsīt . \\
\noindent \textbf{(P\_5,2.94.1)} gomān bhavitā iti . \\
\noindent \textbf{(P\_5,2.94.1)} sati api asteḥ prayoge yathā iha bahuvacanam śrūyate gāvaḥ asya āsan gāvaḥ asya bhavitāraḥ evam iha api syāt . \\
\noindent \textbf{(P\_5,2.94.1)} gomān āsīt . \\
\noindent \textbf{(P\_5,2.94.1)} gomān bhavitā iti . \\
\noindent \textbf{(P\_5,2.94.1)} kā tarhi iyam vācoyuktiḥ . \\
\noindent \textbf{(P\_5,2.94.1)} gomān āsīt . \\
\noindent \textbf{(P\_5,2.94.1)} gomān bhavitā iti . \\
\noindent \textbf{(P\_5,2.94.1)} eṣā eṣā vācoyuktiḥ . \\
\noindent \textbf{(P\_5,2.94.1)} na eṣā gavām sattā kathyate . \\
\noindent \textbf{(P\_5,2.94.1)} kim tarhi . \\
\noindent \textbf{(P\_5,2.94.1)} gomatsattā eṣā kathyate . \\
\noindent \textbf{(P\_5,2.94.1)} asti atra vartamānakālaḥ astiḥ . \\
\noindent \textbf{(P\_5,2.94.1)} katham tarhi bhūtabhaviṣyatsattā gamyate . \\
\noindent \textbf{(P\_5,2.94.1)} dhātusambandhe pratyayāḥ iti . \\
\noindent \textbf{(P\_5,2.94.1)} idam tarhi prayojanam . \\
\noindent \textbf{(P\_5,2.94.1)} astiyuktāt yathā syāt . \\
\noindent \textbf{(P\_5,2.94.1)} anantarādiyuktāt mā bhūt iti . \\
\noindent \textbf{(P\_5,2.94.1)} gāvaḥ asya anantarāḥ . \\
\noindent \textbf{(P\_5,2.94.1)} gāvaḥ asya samīpe iti . \\
\noindent \textbf{(P\_5,2.94.1)} atha kriyamāṇe api astrigrahaṇe iha kasmāt na bhavati . \\
\noindent \textbf{(P\_5,2.94.1)} gāvaḥ asya santi anantarāḥ . \\
\noindent \textbf{(P\_5,2.94.1)} gāvaḥ asya santi samīpe iti . \\
\noindent \textbf{(P\_5,2.94.1)} asāmarthyāt . \\
\noindent \textbf{(P\_5,2.94.1)} katham asāmarthyam . \\
\noindent \textbf{(P\_5,2.94.1)} sāpekṣam asarmartham bhavati iti . \\
\noindent \textbf{(P\_5,2.94.1)} yathā eva tarhi kriyamāṇe astrigrahaṇe asāmarthyāt anantarādiṣu na bhavanti evam akriyamāṇe api na bhaviṣyati . \\
\noindent \textbf{(P\_5,2.94.1)} asti atra viśeṣaḥ . \\
\noindent \textbf{(P\_5,2.94.1)} kriyamāṇe astrigrahaṇe na antareṇa tṛtīyasya padasya prayogam antarādayaḥ arthāḥ gamyante . \\
\noindent \textbf{(P\_5,2.94.1)} akriyamāṇe punaḥ astrigrahaṇe antareṇa api tṛtīyasya padasya prayogam antarādayaḥ arthāḥ gamyante . \\
\noindent \textbf{(P\_5,2.94.2)} atha iha kasmāt na bhavati . \\
\noindent \textbf{(P\_5,2.94.2)} citraguḥ . \\
\noindent \textbf{(P\_5,2.94.2)} śabalaguḥ iti . \\
\noindent \textbf{(P\_5,2.94.2)} bahuvrīhyuktatvāt matvarthasya . \\
\noindent \textbf{(P\_5,2.94.2)} atha iha kasmāt na bhavati . \\
\noindent \textbf{(P\_5,2.94.2)} citrāḥ gāvaḥ asya santi iti . \\
\noindent \textbf{(P\_5,2.94.2)} kutaḥ kasmāt na bhavati . \\
\noindent \textbf{(P\_5,2.94.2)} kim avayavāt āhosvit samudāyāt . \\
\noindent \textbf{(P\_5,2.94.2)} avayavāt kasmāt na bhavati . \\
\noindent \textbf{(P\_5,2.94.2)} asāmarthyāt . \\
\noindent \textbf{(P\_5,2.94.2)} katham asāmarthyam . \\
\noindent \textbf{(P\_5,2.94.2)} sāpekṣam asarmartham bhavati iti . \\
\noindent \textbf{(P\_5,2.94.2)} samudāyāt tarhi kasmāt na bhavati . \\
\noindent \textbf{(P\_5,2.94.2)} aprātipadikatvāt . \\
\noindent \textbf{(P\_5,2.94.2)} nanu ca bhoḥ ākṛtau śāstrāṇi pravartante . \\
\noindent \textbf{(P\_5,2.94.2)} tat yathā sup supā iti vartamāne anyasya ca anyasya ca samāsaḥ bhavati . \\
\noindent \textbf{(P\_5,2.94.2)} satyam evam etat . \\
\noindent \textbf{(P\_5,2.94.2)} ākṛtiḥ tu pratyekam parisamāpyate . \\
\noindent \textbf{(P\_5,2.94.2)} yāvati etat parisamāpyate ṅyāpprātipadikāt iti tāvataḥ utpattyā bhavitavyam . \\
\noindent \textbf{(P\_5,2.94.2)} avayave ca etat parisamāpyate na samudāye . \\
\noindent \textbf{(P\_5,2.94.2)} atha iha kasmāt na bhavati . \\
\noindent \textbf{(P\_5,2.94.2)} pañca gāvaḥ asya santi pañcaguḥ . \\
\noindent \textbf{(P\_5,2.94.2)} daśaguḥ iti . \\
\noindent \textbf{(P\_5,2.94.2)} pratyekam asāmarthyāt samudāyāt aprātipadikatvāt samāsāt samāsena uktatvāt . \\
\noindent \textbf{(P\_5,2.94.2)} na etat sāram . \\
\noindent \textbf{(P\_5,2.94.2)} ukte api hi pratyayārthe utpadyate dvigoḥ taddhitaḥ . \\
\noindent \textbf{(P\_5,2.94.2)} tat yathā dvaimāturaḥ pāñcanāpitiḥ iti . \\
\noindent \textbf{(P\_5,2.94.2)} na eṣaḥ dviguḥ . \\
\noindent \textbf{(P\_5,2.94.2)} kaḥ tarhi . \\
\noindent \textbf{(P\_5,2.94.2)} bahuvrīhiḥ . \\
\noindent \textbf{(P\_5,2.94.2)} apavādatvāt dviguḥ prāpnoti . \\
\noindent \textbf{(P\_5,2.94.2)} antaraṅgatvāt bahuvrīhiḥ bhaviṣyati . \\
\noindent \textbf{(P\_5,2.94.2)} kā antaraṅgatā . \\
\noindent \textbf{(P\_5,2.94.2)} anyapadārthe bahuvrīhiḥ vartate viśiṣṭe anyapadārthe taddhitārthe dviguḥ . \\
\noindent \textbf{(P\_5,2.94.2)} tasmin ca asya taddhite astigrahaṇam kriyate . \\
\noindent \textbf{(P\_5,2.94.2)} yadi tarhi atiprasaṅgāḥ santi bahuvrīhau api astigrahaṇam kartavyam astiyuktāt yathā syāt anantarādiyuktāt mā bhūt iti . \\
\noindent \textbf{(P\_5,2.94.2)} atha na santi taddhitavidhau api na arthaḥ astigrahaṇena . \\
\noindent \textbf{(P\_5,2.94.2)} satyam evam etat . \\
\noindent \textbf{(P\_5,2.94.2)} kriyate tu idānīm taddhitavidhau astigrahaṇam . \\
\noindent \textbf{(P\_5,2.94.2)} tat vai kriyamāṇam api pratyayavidhyartham na upādhyartham . \\
\noindent \textbf{(P\_5,2.94.2)} astimān iti matup yathā syāt . \\
\noindent \textbf{(P\_5,2.94.2)} kim ca kāraṇam na syāt . \\
\noindent \textbf{(P\_5,2.94.2)} aprātipadikatvāt . \\
\noindent \textbf{(P\_5,2.94.2)} na eṣaḥ doṣaḥ . \\
\noindent \textbf{(P\_5,2.94.2)} avyayam eṣaḥ astiśabdaḥ . \\
\noindent \textbf{(P\_5,2.94.2)} na eṣa asteḥ laṭ . \\
\noindent \textbf{(P\_5,2.94.2)} katham avyayatvam . \\
\noindent \textbf{(P\_5,2.94.2)} vibhaktisvarapratirūpakāḥ ca nipātāḥ bhavanti iti nipātasañjñā . \\
\noindent \textbf{(P\_5,2.94.2)} nipātaḥ avyayam iti avyayasañjñā . \\
\noindent \textbf{(P\_5,2.94.2)} evam api na sidhyati . \\
\noindent \textbf{(P\_5,2.94.2)} kim kāraṇam . \\
\noindent \textbf{(P\_5,2.94.2)} astisāmānādhikaraṇye matup vidhīyate . \\
\noindent \textbf{(P\_5,2.94.2)} na ca asteḥ astinā sāmānādhikaraṇyam . \\
\noindent \textbf{(P\_5,2.94.2)} tat etat kriyamāṇam api pratyayavidhyartham na upādhyartham . \\
\noindent \textbf{(P\_5,2.94.2)} tasmāt dvigoḥ taddhitasya pratiṣedhaḥ vaktavyaḥ yadi tat na asti sarvatra matvarthe pratiṣedhaḥ iti . \\
\noindent \textbf{(P\_5,2.94.2)} sati hi tasmin tena eva siddham . \\
\noindent \textbf{(P\_5,2.94.3)} atha matvarthīyāt matvarthīyena bhavitavyam . \\
\noindent \textbf{(P\_5,2.94.3)} na bhavitavyam . \\
\noindent \textbf{(P\_5,2.94.3)} kim kāraṇam . \\
\noindent \textbf{(P\_5,2.94.3)} arthagatyarthaḥ śabdaprayogaḥ . \\
\noindent \textbf{(P\_5,2.94.3)} artham sampratyāyayiṣyāmi iti śabdaḥ prayujyate . \\
\noindent \textbf{(P\_5,2.94.3)} tatra ekena uktatvāt tasya arthasya dvitīyasya prayogeṇa na bhavitavyam . \\
\noindent \textbf{(P\_5,2.94.3)} kim kāraṇam . \\
\noindent \textbf{(P\_5,2.94.3)} uktārthānām aprayogaḥ iti . \\
\noindent \textbf{(P\_5,2.94.3)} na tarhi idānīm idam bhavati : daṇḍimatī śālā . \\
\noindent \textbf{(P\_5,2.94.3)} hastimatī upapatyakā iti . \\
\noindent \textbf{(P\_5,2.94.3)} bhavati . \\
\noindent \textbf{(P\_5,2.94.3)} arthāntare vṛttāt arthāntare vṛttiḥ . \\
\noindent \textbf{(P\_5,2.94.3)} ṣaṣṭhyarthe vā vṛttam saptamyarthe vartate saptamyarthe vā vṛttam ṣaṣṭhyarthe vartate . \\
\noindent \textbf{(P\_5,2.94.3)} atha matvantāt matupā bhavitavyam : gomantaḥ asya santi . \\
\noindent \textbf{(P\_5,2.94.3)} yavamantaḥ asya santi iti . \\
\noindent \textbf{(P\_5,2.94.3)} na bhavitavyam . \\
\noindent \textbf{(P\_5,2.94.3)} kim kāṛaṇam . \\
\noindent \textbf{(P\_5,2.94.3)} yasya gomantaḥ santi gāvaḥ api tasya santi . \\
\noindent \textbf{(P\_5,2.94.3)} tatra uktaḥ gobhiḥ abhisambandhe pratyayaḥ iti kṛtvā taddhitaḥ na bhaviṣyati . \\
\noindent \textbf{(P\_5,2.94.3)} na tarhi idānīm idam bhavati : daṇḍimatī śālā . \\
\noindent \textbf{(P\_5,2.94.3)} hastimatī upapatyakā iti . \\
\noindent \textbf{(P\_5,2.94.3)} bhavati . \\
\noindent \textbf{(P\_5,2.94.3)} arthāntare vṛttāt arthāntare vṛttiḥ . \\
\noindent \textbf{(P\_5,2.94.3)} ṣaṣṭhyarthe vā vṛttam saptamyarthe vartate saptamyarthe vā vṛttam ṣaṣṭhyarthe vartate . \\
\noindent \textbf{(P\_5,2.94.3)} iha api saptamyarthe vā vṛttam ṣaṣṭhyarthe vartate ṣaṣṭhyarthe vā vṛttam saptamyarthe vartate . \\
\noindent \textbf{(P\_5,2.94.3)} anyathājātīyakaḥ khalu api gobhiḥ abhisambandhe pratyayaḥ anyathājātīyakaḥ tadvatā . \\
\noindent \textbf{(P\_5,2.94.3)} yena eva khalu api hetunā etat vākyam bhavati gomantaḥ asya santi , yavamantaḥ asya santi iti tena eva hetunā vṛttiḥ api prāpnoti . \\
\noindent \textbf{(P\_5,2.94.3)} tasmāt matvarthīyāt matubādeḥ pratiṣedhaḥ vaktavyaḥ . \\
\noindent \textbf{(P\_5,2.94.3)} tam ca api bruvatā samānvṛttau sarūpaḥ iti vaktavyam . \\
\noindent \textbf{(P\_5,2.94.3)} bhavati hi daṇḍimatī śālā hastimatī upapatyakā iti . \\
\noindent \textbf{(P\_5,2.94.3)} śaiṣikāt matubarthīyāt śaiṣikaḥ matubarthīyaḥ sarūpaḥ pratyayaḥ na iṣṭaḥ . \\
\noindent \textbf{(P\_5,2.94.3)} sanantāt na san iṣyate . \\
\noindent \textbf{(P\_5,2.94.4)} kim punaḥ ime matupprabhṛtayaḥ sanmātre bhavanti . \\
\noindent \textbf{(P\_5,2.94.4)} evam bhavitum arhati . \\
\noindent \textbf{(P\_5,2.94.4)} matupprabhṛtayaḥ sanmātre cet atiprasaṅgaḥ . \\
\noindent \textbf{(P\_5,2.94.4)} matupprabhṛtayaḥ sanmātre cet atiprasaṅgaḥ bhavati . \\
\noindent \textbf{(P\_5,2.94.4)} iha api prāpnoti . \\
\noindent \textbf{(P\_5,2.94.4)} vrīhiḥ asya . \\
\noindent \textbf{(P\_5,2.94.4)} yavaḥ asya iti . \\
\noindent \textbf{(P\_5,2.94.4)} tasmāt bhūmādigrahaṇam kartavyam . \\
\noindent \textbf{(P\_5,2.94.4)} ke punaḥ bhūmādayaḥ . \\
\noindent \textbf{(P\_5,2.94.4)} bhūmanindāpraśaṃsāsu nityayoge atiśāyane saṃsarge astivivakṣāyām bhavanti matubādayaḥ . \\
\noindent \textbf{(P\_5,2.94.4)} bhūmni : gomān yavamān . \\
\noindent \textbf{(P\_5,2.94.4)} nindāyām : kakudāvartī saṅkhādakī . \\
\noindent \textbf{(P\_5,2.94.4)} praśaṃsayām : rūpavān varṇavān . \\
\noindent \textbf{(P\_5,2.94.4)} nityayoge : kṣīriṇaḥ vṛkṣāḥ , kaṇṭakinaḥ vṛkṣāḥ iti . \\
\noindent \textbf{(P\_5,2.94.4)} atiśāyane : udariṇī kanyā . \\
\noindent \textbf{(P\_5,2.94.4)} saṃsarge : daṇḍī chatrī . \\
\noindent \textbf{(P\_5,2.94.4)} tat tarhi bhūmādigrahaṇam kartavyam . \\
\noindent \textbf{(P\_5,2.94.4)} na kartavyam . \\
\noindent \textbf{(P\_5,2.94.4)} kasmāt na bhavati . \\
\noindent \textbf{(P\_5,2.94.4)} vrīhiḥ asya . \\
\noindent \textbf{(P\_5,2.94.4)} yavaḥ asya iti . \\
\noindent \textbf{(P\_5,2.94.4)} uktam vā. kim uktam . \\
\noindent \textbf{(P\_5,2.94.4)} anabhidhānāt iti . \\
\noindent \textbf{(P\_5,2.94.4)} itikaraṇaḥ khalu api kriyate . \\
\noindent \textbf{(P\_5,2.94.4)} tataḥ cet vivakṣā . \\
\noindent \textbf{(P\_5,2.94.4)} bhūmādiyuktasya eva ca vivakṣā . \\
\noindent \textbf{(P\_5,2.94.4)} gomān yavamān . \\
\noindent \textbf{(P\_5,2.94.4)} bhūmādiyuktasya eva sattā kathyate . \\
\noindent \textbf{(P\_5,2.94.4)} na hi kasya cit yavaḥ na asti . \\
\noindent \textbf{(P\_5,2.94.4)} saṅkhādakī kakudāvartinī . \\
\noindent \textbf{(P\_5,2.94.4)} nindāuktasya eva sattā kathyate . \\
\noindent \textbf{(P\_5,2.94.4)} na hi kaḥ cit na saṅkhādakī . \\
\noindent \textbf{(P\_5,2.94.4)} rūpavān varṇavān . \\
\noindent \textbf{(P\_5,2.94.4)} praśaṃsāyuktasya eva sattā kathyate . \\
\noindent \textbf{(P\_5,2.94.4)} na hi kasya cit rūpam na asti . \\
\noindent \textbf{(P\_5,2.94.4)} kṣīriṇaḥ vṛkṣāḥ . \\
\noindent \textbf{(P\_5,2.94.4)} kaṇṭakinaḥ vṛkṣāḥ iti . \\
\noindent \textbf{(P\_5,2.94.4)} nityayuktasya eva sattā kathyate . \\
\noindent \textbf{(P\_5,2.94.4)} na hi kasya cit kṣiram na asti . \\
\noindent \textbf{(P\_5,2.94.4)} udariṇī kanyā iti . \\
\noindent \textbf{(P\_5,2.94.4)} atiśāyanayuktasya eva sattā kathyate . \\
\noindent \textbf{(P\_5,2.94.4)} na hi kasya cit udaram na asti . \\
\noindent \textbf{(P\_5,2.94.4)} daṇḍī chatrī . \\
\noindent \textbf{(P\_5,2.94.4)} saṃsargayuktasya eva sattā kathyate . \\
\noindent \textbf{(P\_5,2.94.4)} na hi kasya cit daṇḍaḥ na asti . \\
\noindent \textbf{(P\_5,2.94.4)} yāvatībhiḥ khalu api gobhiḥ vāhadohaprasavāḥ kalpante tāvatīṣu sattā kathyate . \\
\noindent \textbf{(P\_5,2.94.4)} kasya cit tisṛbhiḥ kalpante kasya cit śatena api na prakalpante . \\
\noindent \textbf{(P\_5,2.94.4)} sanmātre ca ṛṣidarśanāt . \\
\noindent \textbf{(P\_5,2.94.4)} sanmātre ca punaḥ ṛṣiḥ darśayati matupam . \\
\noindent \textbf{(P\_5,2.94.4)} yavamatībhiḥ adbhiḥ yūpam prokṣati iti . \\
\noindent \textbf{(P\_5,2.94.5)} guṇavacanebhyaḥ matupaḥ luk . \\
\noindent \textbf{(P\_5,2.94.5)} guṇavacanebhyaḥ matupaḥ luk vaktavyaḥ . \\
\noindent \textbf{(P\_5,2.94.5)} śuklaḥ kṛṣṇaḥ iti . \\
\noindent \textbf{(P\_5,2.94.5)} avyatirekāt siddham . \\
\noindent \textbf{(P\_5,2.94.5)} na guṇaḥ guṇinam vyabhicarati iti . \\
\noindent \textbf{(P\_5,2.94.5)} avyatirekāt siddham iti cet dṛṣṭaḥ vyatirekaḥ . \\
\noindent \textbf{(P\_5,2.94.5)} dṛśyate vyatirekaḥ . \\
\noindent \textbf{(P\_5,2.94.5)} tat yatha paṭasya śuklaḥ iti . \\
\noindent \textbf{(P\_5,2.94.5)} tathā ca liṅgavacanasiddhiḥ . \\
\noindent \textbf{(P\_5,2.94.5)} evam ca kṛtvā liṅgavacanāni siddhāni bhavanti . \\
\noindent \textbf{(P\_5,2.94.5)} śuklam vastram . \\
\noindent \textbf{(P\_5,2.94.5)} śuklā śāṭī . \\
\noindent \textbf{(P\_5,2.94.5)} śuklaḥ kambalaḥ . \\
\noindent \textbf{(P\_5,2.94.5)} śuklau kambalau . \\
\noindent \textbf{(P\_5,2.94.5)} śuklāḥ kambalāḥ iti . \\
\noindent \textbf{(P\_5,2.94.5)} yat asau dravyam śritaḥ bhavati guṇaḥ tasya yat liṅgam vacanam ca tat guṇasya api bhavati . \\
\noindent \textbf{(P\_5,2.95)} kimartham idam ucyate na tat asya asti asmin iti eva matup siddhaḥ . \\
\noindent \textbf{(P\_5,2.95)} rasādibhyaḥ punarvacanam anyanirvṛttyartham . \\
\noindent \textbf{(P\_5,2.95)} rasādibhyaḥ punarvacanam kriyate anyeṣām matvarthīyānām pratiṣedhārtham . \\
\noindent \textbf{(P\_5,2.95)} matup eva yathā syāt . \\
\noindent \textbf{(P\_5,2.95)} ye anye matvarthīyāḥ prāpnuvanti te mā bhūvan iti . \\
\noindent \textbf{(P\_5,2.95)} na etat asti prayojanam . \\
\noindent \textbf{(P\_5,2.95)} dṛśyante hi anye rasādibhyaḥ matvarthīyāḥ . \\
\noindent \textbf{(P\_5,2.95)} rasikaḥ naṭaḥ . \\
\noindent \textbf{(P\_5,2.95)} urvaśī vai rūpiṇī apsarasām . \\
\noindent \textbf{(P\_5,2.95)} sparśikaḥ vāyuḥ iti . \\
\noindent \textbf{(P\_5,2.96)} iha kasmāt na bhavati : cikīrṣā asya asti , jihīrṣā asya asti iti . \\
\noindent \textbf{(P\_5,2.96)} prāṇyaṅgāt iti vaktavyam . \\
\noindent \textbf{(P\_5,2.96)} tat tarhi vaktavyam . \\
\noindent \textbf{(P\_5,2.96)} na vaktavyam . \\
\noindent \textbf{(P\_5,2.96)} kasmāt na bhavati : cikīrṣā asya asti , jihīrṣā asya asti iti . \\
\noindent \textbf{(P\_5,2.96)} anabhidhānāt . \\
\noindent \textbf{(P\_5,2.97)} sidhmādiṣu yāni akārāntāni tebhyaḥ lacā mukte iniṭhanau prapnutaḥ iniṭhanau ca na iṣyete . \\
\noindent \textbf{(P\_5,2.97)} lac anyatarasyām iti samuccayaḥ . \\
\noindent \textbf{(P\_5,2.97)} lac anyatarasyām iti samuccayaḥ ayam na vibhāṣā . \\
\noindent \textbf{(P\_5,2.97)} lac ca matup ca . \\
\noindent \textbf{(P\_5,2.97)} katham punaḥ etat jñāyate lac anyatarasyām iti samuccayaḥ ayam na vibhāṣā iti . \\
\noindent \textbf{(P\_5,2.97)} picchādibhyaḥ tundādīnām nānāyogakaraṇam jñāpakam asamāveśasya . \\
\noindent \textbf{(P\_5,2.97)} yat ayam picchādibhyaḥ tundādīnām nānāyogam karoti tat jñāpayati ācāryaḥ samuccayaḥ ayam na vibhāṣā iti . \\
\noindent \textbf{(P\_5,2.97)} yadi vibhāṣā syāt nānāyogakaraṇam anarthakam syāt . \\
\noindent \textbf{(P\_5,2.97)} tundādīni api picchādiṣu eva paṭhet . \\
\noindent \textbf{(P\_5,2.97)} na etat asti jñāpakam . \\
\noindent \textbf{(P\_5,2.97)} asti hi anyat nānāyogakaraṇe prayojanam . \\
\noindent \textbf{(P\_5,2.97)} kim . \\
\noindent \textbf{(P\_5,2.97)} tundādiṣu yāni anakārāntāni tebhyaḥ iniṭhanau yathā syātām . \\
\noindent \textbf{(P\_5,2.97)} yāni tarhi akārāntāni teṣām pāṭhaḥ kimarthaḥ . \\
\noindent \textbf{(P\_5,2.97)} jñāpakārthaḥ eva . \\
\noindent \textbf{(P\_5,2.97)} aparaḥ āha : picchādibhyaḥ tundādīnām nānāyogakaraṇam jñāpakam asamāveśasya . \\
\noindent \textbf{(P\_5,2.97)} yat ayam tundādibhyaḥ picchādīnām nānāyogam karoti tat jñāpayati ācāryaḥ samuccayaḥ ayam na vibhāṣā iti . \\
\noindent \textbf{(P\_5,2.97)} yadi vibhāṣā syāt nānāyogakaraṇam anarthakam syāt . \\
\noindent \textbf{(P\_5,2.97)} picchādīni api tundādiṣu eva paṭhet . \\
\noindent \textbf{(P\_5,2.97)} na etat asti jñāpakam . \\
\noindent \textbf{(P\_5,2.97)} asti hi anyat nānāyogakaraṇe prayojanam . \\
\noindent \textbf{(P\_5,2.97)} kim . \\
\noindent \textbf{(P\_5,2.97)} picchādiṣu yāni anakārāntāni tebhyaḥ iniṭhanau yathā syātām . \\
\noindent \textbf{(P\_5,2.97)} yāni tarhi akārāntāni teṣām pāṭhaḥ kimarthaḥ . \\
\noindent \textbf{(P\_5,2.97)} jñāpakārthaḥ eva . \\
\noindent \textbf{(P\_5,2.97)} vasya ca punarvacanam sarvavibhāṣārtham . \\
\noindent \textbf{(P\_5,2.97)} vasya khalu api punarvacanam kriyate sarvavibhāṣārtham . \\
\noindent \textbf{(P\_5,2.97)} keśāt vaḥ anyatarasyām iti . \\
\noindent \textbf{(P\_5,2.97)} etat eva jñāpayati ācāryaḥ samuccayaḥ ayam na vibhāṣā iti . \\
\noindent \textbf{(P\_5,2.97)} dyudrubhyām nityārtham eke anyatarasyāṅgrahaṇam icchanti . \\
\noindent \textbf{(P\_5,2.97)} katham . \\
\noindent \textbf{(P\_5,2.97)} vibhāṣāmadhye ayam yogaḥ kriyate . \\
\noindent \textbf{(P\_5,2.97)} vibhāṣāmadhye ye vidhayaḥ nityāḥ te bhavanti iti . \\
\noindent \textbf{(P\_5,2.100)} naprakaraṇe dadrvāḥ hrasvatvam ca . \\
\noindent \textbf{(P\_5,2.100)} naprakaraṇe dadrvāḥ upasaṅkhyānam kartavyam hrasvatvam ca naprakaraṇe dadrvāḥ hrasvatvam ca vaktavyam . \\
\noindent \textbf{(P\_5,2.100)} dadruṇaḥ . \\
\noindent \textbf{(P\_5,2.100)} atyalpam idam ucyate . \\
\noindent \textbf{(P\_5,2.100)} śākīpalālīdadrūṇām hrasvatvam ca iti vaktavyam . \\
\noindent \textbf{(P\_5,2.100)} śākinam , palālinam , dadruṇam . \\
\noindent \textbf{(P\_5,2.100)} viṣvak iti uttarapadalopaḥ ca akṛtasandheḥ . \\
\noindent \textbf{(P\_5,2.100)} viṣvak iti upasaṅkhyānam kartavyam uttarapadalopaḥ ca akṛtasandheḥ vaktavyaḥ . \\
\noindent \textbf{(P\_5,2.100)} viṣvak gatāni asya viṣuṇaḥ . \\
\noindent \textbf{(P\_5,2.101)} vṛtteḥ ca . \\
\noindent \textbf{(P\_5,2.101)} vṛtteḥ ca iti vaktavyam . \\
\noindent \textbf{(P\_5,2.101)} vārttam . \\
\noindent \textbf{(P\_5,2.102-103.1)} KA\_II,396.19-22 Ro\_IV,166 {1/5}     kimartham tapaḥśabdāt vin vidhīyate na asantāt iti eva siddham . \\
\noindent \textbf{(P\_5,2.102-103.1)} KA\_II,396.19-22 Ro\_IV,166 {2/5}     tapasaḥ vinvacanam aṇvidhānāt . \\
\noindent \textbf{(P\_5,2.102-103.1)} KA\_II,396.19-22 Ro\_IV,166 {3/5}     tapasaḥ vinvacanam kriyate . \\
\noindent \textbf{(P\_5,2.102-103.1)} KA\_II,396.19-22 Ro\_IV,166 {4/5}     tapaḥśabdāt an vidhīyate . \\
\noindent \textbf{(P\_5,2.102-103.1)} KA\_II,396.19-22 Ro\_IV,166 {5/5}     saḥ viśeṣavihitaḥ sāmānyvihitam vinam bādheta . \\
\noindent \textbf{(P\_5,2.102-103.2)} KA\_II,397.1-3 Ro\_IV,167 {1/7}     aṇprakaraṇe jyotsnādibhyaḥ upasaṅkhyānam . \\
\noindent \textbf{(P\_5,2.102-103.2)} KA\_II,397.1-3 Ro\_IV,167 {2/7}     aṇprakaraṇe jyotsnādibhyaḥ upasaṅkhyānam kartavyam . \\
\noindent \textbf{(P\_5,2.102-103.2)} KA\_II,397.1-3 Ro\_IV,167 {3/7}     jyautsnaḥ . \\
\noindent \textbf{(P\_5,2.102-103.2)} KA\_II,397.1-3 Ro\_IV,167 {4/7}     tāmisraḥ . \\
\noindent \textbf{(P\_5,2.102-103.2)} KA\_II,397.1-3 Ro\_IV,167 {5/7}     kauṇḍalaḥ . \\
\noindent \textbf{(P\_5,2.102-103.2)} KA\_II,397.1-3 Ro\_IV,167 {6/7}     kautapaḥ . \\
\noindent \textbf{(P\_5,2.102-103.2)} KA\_II,397.1-3 Ro\_IV,167 {7/7}     vaipādikaḥ . \\
\noindent \textbf{(P\_5,2.107.1)} ayam madhuśabdaḥ asti eva dravyapadārthakaḥ asti rasavācī . \\
\noindent \textbf{(P\_5,2.107.1)} ātaḥ ca rasavācī api . \\
\noindent \textbf{(P\_5,2.107.1)} madhuni eva hi madhu idam madhuram iti prasajyate . \\
\noindent \textbf{(P\_5,2.107.1)} tat yaḥ rasavācī tasya idam grahaṇam . \\
\noindent \textbf{(P\_5,2.107.1)} yadi hi dravyapadārthakasya grahaṇam syāt iha api prasajyeta . \\
\noindent \textbf{(P\_5,2.107.1)} madhu asmin ghaṭe asti . \\
\noindent \textbf{(P\_5,2.107.2)} raprakaraṇe khamukhkuñjebhyaḥ upasaṅkhyānam . \\
\noindent \textbf{(P\_5,2.107.2)} kharaḥ . \\
\noindent \textbf{(P\_5,2.107.2)} mukharaḥ . \\
\noindent \textbf{(P\_5,2.107.2)} kuñjaraḥ . \\
\noindent \textbf{(P\_5,2.107.2)} nagāt ca iti vaktavyam . \\
\noindent \textbf{(P\_5,2.107.2)} nagaram . \\
\noindent \textbf{(P\_5,2.109)} vaprakaraṇe maṇihiraṇyābhyām upasaṅkhyānam . \\
\noindent \textbf{(P\_5,2.109)} vaprakaraṇe maṇihiraṇyābhyām upasaṅkhyānam vaktavyam . \\
\noindent \textbf{(P\_5,2.109)} maṇivaḥ . \\
\noindent \textbf{(P\_5,2.109)} hiraṇyavaḥ . \\
\noindent \textbf{(P\_5,2.109)} chandasi īvanipau ca . \\
\noindent \textbf{(P\_5,2.109)} chandasi īvanipau ca vaktavyau vaḥ ca matup ca . \\
\noindent \textbf{(P\_5,2.109)} rathīḥ abhūt mudgalanī gaviṣṭau . \\
\noindent \textbf{(P\_5,2.109)} sumaṅgalīḥ iyam vadhuḥ . \\
\noindent \textbf{(P\_5,2.109)} ṛtavānam . \\
\noindent \textbf{(P\_5,2.109)} maghavānam īmahe . \\
\noindent \textbf{(P\_5,2.109)} ut vā ca udvatī ca . \\
\noindent \textbf{(P\_5,2.109)} medhārathābhyām iraniracau . \\
\noindent \textbf{(P\_5,2.109)} medhārathābhyām iraniracau vaktavyau . \\
\noindent \textbf{(P\_5,2.109)} medhiraḥ . \\
\noindent \textbf{(P\_5,2.109)} rathiraḥ . \\
\noindent \textbf{(P\_5,2.109)} aparaḥ āha : vāprakaraṇe anyebhyaḥ api dṛśyate iti vaktavyam . \\
\noindent \textbf{(P\_5,2.109)} bimbāvam . \\
\noindent \textbf{(P\_5,2.109)} kurarāvam iṣṭakāvam . \\
\noindent \textbf{(P\_5,2.112)} valacprakaraṇe anyebhyaḥ api ḍrśyate . \\
\noindent \textbf{(P\_5,2.112)} valacprakaraṇe anyebhyaḥ api ḍrśyate iti vaktavyam . \\
\noindent \textbf{(P\_5,2.112)} bhrātṛvalaḥ . \\
\noindent \textbf{(P\_5,2.112)} putravalaḥ . \\
\noindent \textbf{(P\_5,2.112)} utsaṅgavalaḥ . \\
\noindent \textbf{(P\_5,2.115)} iniṭhanoḥ ekākṣarāt pratiṣedhaḥ . \\
\noindent \textbf{(P\_5,2.115)} iniṭhanoḥ ekākṣarāt pratiṣedhaḥ vaktavyaḥ : svavān , khavān . \\
\noindent \textbf{(P\_5,2.115)} atyalpam idam ucyate . \\
\noindent \textbf{(P\_5,2.115)} ekākṣarāt kṛtaḥ jāteḥ saptamyām ca na tau smṛtau . \\
\noindent \textbf{(P\_5,2.115)} ekākṣarāt : svavān , khavān . \\
\noindent \textbf{(P\_5,2.115)} kṛtaḥ : kārakavān , hārakavān . \\
\noindent \textbf{(P\_5,2.115)} jāteḥ : vṛkṣavān , plakṣavān , vyāghravān , siṃhavān . \\
\noindent \textbf{(P\_5,2.115)} saptamyām ca na tau . \\
\noindent \textbf{(P\_5,2.115)} daṇḍāḥ asyām śālāyām santi iti . \\
\noindent \textbf{(P\_5,2.115)} yadi kṛtaḥ na iti ucyate kāryī kāryikaḥ iti na sidhyati . \\
\noindent \textbf{(P\_5,2.115)} tathā ca yadi jāteḥ na iti ucyate tuṇḍalī tuṇḍalikaḥ iti na sidhyati . \\
\noindent \textbf{(P\_5,2.115)} evam tarhi na ayam samuccayaḥ kṛtaḥ ca jāteḥ ca iti . \\
\noindent \textbf{(P\_5,2.115)} kim tarhi . \\
\noindent \textbf{(P\_5,2.115)} jātiviśeṣaṇam kṛdgrahaṇam : kṛt yā jātiḥ iti . \\
\noindent \textbf{(P\_5,2.115)} katham kārakavān , hārakavān . \\
\noindent \textbf{(P\_5,2.115)} anabhidhānāt na bhaviṣyati . \\
\noindent \textbf{(P\_5,2.115)} yadi evam na arthaḥ anena . \\
\noindent \textbf{(P\_5,2.115)} katham svavān , vṛkṣavān , siṃhavān , vyāghravān daṇḍāḥ asyām śālāyām santi iti . \\
\noindent \textbf{(P\_5,2.115)} anabhidhānāt na bhaviṣyati . \\
\noindent \textbf{(P\_5,2.116)} śikhādibhyaḥ iniḥ vaktavyaḥ ikan yavakhadādiṣu . \\
\noindent \textbf{(P\_5,2.116)} kim prayojanam . \\
\noindent \textbf{(P\_5,2.116)} niyamārtham . \\
\noindent \textbf{(P\_5,2.116)} iniḥ eva śikhādibhyaḥ ikan eva yavkhadādibhyaḥ . \\
\noindent \textbf{(P\_5,2.116)} śikhāyavakhadādibhyaḥ niyamasya avacanam nivartakatvāt . \\
\noindent \textbf{(P\_5,2.116)} śikhāyavakhadādibhyaḥ niyamasya avacanam . \\
\noindent \textbf{(P\_5,2.116)} kim kāraṇam . \\
\noindent \textbf{(P\_5,2.116)} nivartakatvāt . \\
\noindent \textbf{(P\_5,2.116)} kim nivartakam . \\
\noindent \textbf{(P\_5,2.116)} anabhidhānam . \\
\noindent \textbf{(P\_5,2.118)} nityagrahaṇam kimartham . \\
\noindent \textbf{(P\_5,2.118)} vibhāṣā mā bhūt . \\
\noindent \textbf{(P\_5,2.118)} na etat asti prayojanam . \\
\noindent \textbf{(P\_5,2.118)} pūrvasmin eva yoge vibhāṣāgrahaṇam nivṛttam . \\
\noindent \textbf{(P\_5,2.118)} evam tarhi siddhe sati yat nityagrahaṇam karoti tat jñāpayati ācāryaḥ prāk etasmāt yogāt vibhāṣā iti anuvartate . \\
\noindent \textbf{(P\_5,2.118)} atha ataḥ iti anuvartate utāho na . \\
\noindent \textbf{(P\_5,2.118)} kim ca ataḥ . \\
\noindent \textbf{(P\_5,2.118)} yadi anuvartate ekagavikaḥ na sidhyati . \\
\noindent \textbf{(P\_5,2.118)} samāsānte kṛte bhaviṣyati . \\
\noindent \textbf{(P\_5,2.118)} evam api gauśakaṭikaḥ na sidhyati . \\
\noindent \textbf{(P\_5,2.118)} atha nivṛttam iha api prāpnoti . \\
\noindent \textbf{(P\_5,2.118)} goviṃśatiḥ asya asti iti . \\
\noindent \textbf{(P\_5,2.118)} nivṛttam . \\
\noindent \textbf{(P\_5,2.118)} kasmāt na bhavati : goviṃśatiḥ asya asti iti . \\
\noindent \textbf{(P\_5,2.118)} anabhidhānāt na bhaviṣyati . \\
\noindent \textbf{(P\_5,2.120)} yapprakaraṇe anyebhyaḥ api dṛśyate . \\
\noindent \textbf{(P\_5,2.120)} yapprakaraṇe anyebhyaḥ api dṛśyate iti vaktavyam. himyāḥ parvatāḥ . \\
\noindent \textbf{(P\_5,2.120)} guṇyāḥ brāhmaṇāḥ . \\
\noindent \textbf{(P\_5,2.122)} chandovinprakaraṇe aṣṭṛāmekhalādvayobhayarujāhṛdayānām dīrghaḥ ca . \\
\noindent \textbf{(P\_5,2.122)} chandovinprakaraṇe aṣṭṛāmekhalādvayobhayarujāhṛdayānām dīrghaḥ ca iti vaktavyam . \\
\noindent \textbf{(P\_5,2.122)} aṣṭrāvī . \\
\noindent \textbf{(P\_5,2.122)} mekhalāvī . \\
\noindent \textbf{(P\_5,2.122)} dvayāvī . \\
\noindent \textbf{(P\_5,2.122)} ubhayāvī . \\
\noindent \textbf{(P\_5,2.122)} rujāvī . \\
\noindent \textbf{(P\_5,2.122)} hṛdayāvī . \\
\noindent \textbf{(P\_5,2.122)} marmaṇaḥ ca iti vaktavyam . \\
\noindent \textbf{(P\_5,2.122)} mamāvī . \\
\noindent \textbf{(P\_5,2.122)} sarvatra āmayasya . \\
\noindent \textbf{(P\_5,2.122)} sarvatra āmayasya upasaṅkhyānam kartavyam . \\
\noindent \textbf{(P\_5,2.122)} āmayāvī . \\
\noindent \textbf{(P\_5,2.122)} śṛṅgavṛndābhyām ārakan . \\
\noindent \textbf{(P\_5,2.122)} śṛṅgavṛndābhyām ārakan vaktavyaḥ . \\
\noindent \textbf{(P\_5,2.122)} śṛṅgārakaḥ . \\
\noindent \textbf{(P\_5,2.122)} vṛdārakaḥ . \\
\noindent \textbf{(P\_5,2.122)} phalabarhābhyām inac . \\
\noindent \textbf{(P\_5,2.122)} phalabarhābhyām inac vaktavyaḥ . \\
\noindent \textbf{(P\_5,2.122)} phalinaḥ . \\
\noindent \textbf{(P\_5,2.122)} barhiṇaḥ . \\
\noindent \textbf{(P\_5,2.122)} hṛdayāt cāluḥ anyatarasyām . \\
\noindent \textbf{(P\_5,2.122)} hṛdayāt cāluḥ vaktavyaḥ anyatarasyām . \\
\noindent \textbf{(P\_5,2.122)} hṛdayāluḥ . \\
\noindent \textbf{(P\_5,2.122)} hṛdayī . \\
\noindent \textbf{(P\_5,2.122)} hṛdayikaḥ . \\
\noindent \textbf{(P\_5,2.122)} hṛdayavān . \\
\noindent \textbf{(P\_5,2.122)} śītoṣṇatṛprebhyaḥ tat na sahate . \\
\noindent \textbf{(P\_5,2.122)} śītoṣṇatṛprebhyaḥ tat na sahate iti cāluḥ vaktavyaḥ . \\
\noindent \textbf{(P\_5,2.122)} śītāluḥ . \\
\noindent \textbf{(P\_5,2.122)} uṣṇāluḥ . \\
\noindent \textbf{(P\_5,2.122)} tṛprāluḥ . \\
\noindent \textbf{(P\_5,2.122)} himāt celuḥ . \\
\noindent \textbf{(P\_5,2.122)} himāt celuḥ vaktavyaḥ tat na sahate iti etasmin arthe . \\
\noindent \textbf{(P\_5,2.122)} himeluḥ . \\
\noindent \textbf{(P\_5,2.122)} balāt ca ūlaḥ . \\
\noindent \textbf{(P\_5,2.122)} balāt ca ūlaḥ vaktavyaḥ tat na sahate iti etasmin arthe . \\
\noindent \textbf{(P\_5,2.122)} balūlaḥ . \\
\noindent \textbf{(P\_5,2.122)} vātāt samūhe ca . \\
\noindent \textbf{(P\_5,2.122)} vātāt samūhe ca tat na sahate iti etasmin arthe ūlaḥ vaktavyaḥ . \\
\noindent \textbf{(P\_5,2.122)} vātūlaḥ . \\
\noindent \textbf{(P\_5,2.122)} parvamarudbhyām tap . \\
\noindent \textbf{(P\_5,2.122)} parvamarudbhyām tap vaktavyaḥ . \\
\noindent \textbf{(P\_5,2.122)} parvataḥ . \\
\noindent \textbf{(P\_5,2.122)} maruttaḥ . \\
\noindent \textbf{(P\_5,2.122)} dadātivṛttam vā . \\
\noindent \textbf{(P\_5,2.122)} dadātivṛttam vā punaḥ etat bhaviṣyati . \\
\noindent \textbf{(P\_5,2.122)} marudbhiḥ dattaḥ maruttaḥ . \\
\noindent \textbf{(P\_5,2.125)} kutsite iti vaktavyam . \\
\noindent \textbf{(P\_5,2.125)} yaḥ hi samyak bahu bhāṣate vāgmī iti eva saḥ bhavati . \\
\noindent \textbf{(P\_5,2.125)} tat tarhi vaktavyam . \\
\noindent \textbf{(P\_5,2.125)} na vaktavyam . \\
\noindent \textbf{(P\_5,2.125)} nānāyogakaraṇasāmarthyāt na bhaviṣyati . \\
\noindent \textbf{(P\_5,2.126)} iha kasmāt na bhavati . \\
\noindent \textbf{(P\_5,2.126)} svam asya asti iti . \\
\noindent \textbf{(P\_5,2.126)} na eṣaḥ doṣaḥ . \\
\noindent \textbf{(P\_5,2.126)} na ayam pratyayārthaḥ . \\
\noindent \textbf{(P\_5,2.126)} kim tarhi prakṛtiviśeṣaṇam etat . \\
\noindent \textbf{(P\_5,2.126)} svāmin aiśvarye nipātyate iti . \\
\noindent \textbf{(P\_5,2.129)} piśācāt ca iti vaktavyam . \\
\noindent \textbf{(P\_5,2.129)} piśācakī vaiśravaṇaḥ . \\
\noindent \textbf{(P\_5,2.135)} iniprakaraṇe balāt bāhūrupūrvapadāt upasaṅkhyānam . \\
\noindent \textbf{(P\_5,2.135)} iniprakaraṇe balāt bāhūrupūrvapadāt upasaṅkhyānam kartavyam . \\
\noindent \textbf{(P\_5,2.135)} bāhubalī . \\
\noindent \textbf{(P\_5,2.135)} ūrubalī . \\
\noindent \textbf{(P\_5,2.135)} sarvādeḥ ca . \\
\noindent \textbf{(P\_5,2.135)} sarvādeḥ ca iniḥ vaktavyaḥ . \\
\noindent \textbf{(P\_5,2.135)} sarvadhanī . \\
\noindent \textbf{(P\_5,2.135)} sarvabījī . \\
\noindent \textbf{(P\_5,2.135)} sarvakeśī . \\
\noindent \textbf{(P\_5,2.135)} arthāt ca asannihite . \\
\noindent \textbf{(P\_5,2.135)} arthāt ca asannihite iniḥ vaktavyaḥ . \\
\noindent \textbf{(P\_5,2.135)} arthī . \\
\noindent \textbf{(P\_5,2.135)} asannihite iti kimartham . \\
\noindent \textbf{(P\_5,2.135)} arthavān . \\
\noindent \textbf{(P\_5,2.135)} tadantāt ca . \\
\noindent \textbf{(P\_5,2.135)} tadantāt ca iti vaktavyam . \\
\noindent \textbf{(P\_5,2.135)} dhānyārthī . \\
\noindent \textbf{(P\_5,2.135)} hiraṇyārthī . \\
\noindent \textbf{(P\_5,2.135)} kimartham tadantāt iti ucyate na tadantavidhinā siddham . \\
\noindent \textbf{(P\_5,2.135)} grahaṇavatā prātipadikena tadantavidhiḥ pratiṣidhyate . \\
\noindent \textbf{(P\_5,2.135)} evarm tarhi inantena saha samāsaḥ bhaviṣyati . \\
\noindent \textbf{(P\_5,2.135)} dhānyena arthī dhānyārthī . \\
\noindent \textbf{(P\_5,2.135)} saḥ hi samāsaḥ na prāpnoti . \\
\noindent \textbf{(P\_5,2.135)} yadi punaḥ ayam arthayateḥ ṇiniḥ syāt . \\
\noindent \textbf{(P\_5,2.135)} evam api kriyām eva kurvāṇe syāt . \\
\noindent \textbf{(P\_5,2.135)} tūṣṇīm api āsīnaḥ yaḥ tatsamarthāni ācarati saḥ abhiprāyeṇa gamyate arthyam anena iti . \\
\noindent \textbf{(P\_5,2.135)} evam tarhi ayam arthaśabdaḥ asti eva dravyapadārthakaḥ . \\
\noindent \textbf{(P\_5,2.135)} tat yathā arthavān ayam deśaḥ iti ucyate yasmin gāvaḥ sasyāni ca vartante . \\
\noindent \textbf{(P\_5,2.135)} asti kriyāpadārthakaḥ bhāvasādhanaḥ . \\
\noindent \textbf{(P\_5,2.135)} arthanam arthaḥ iti . \\
\noindent \textbf{(P\_5,2.135)} tat yaḥ kriyāpadārthakaḥ tasya idam grahaṇam . \\
\noindent \textbf{(P\_5,2.135)} evam ca kṛtvā arthikapratyarthikau api siddhau bhavataḥ . \\
\noindent \textbf{(P\_5,3.1)} vibhaktitve kim prayojanam . \\
\noindent \textbf{(P\_5,3.1)} vibhaktitve prayojanam itpratiṣedhaḥ . \\
\noindent \textbf{(P\_5,3.1)} idānīm . \\
\noindent \textbf{(P\_5,3.1)} na vibhaktau tusmāḥ iti itpratiṣedhaḥ siddhaḥ bhavati . \\
\noindent \textbf{(P\_5,3.1)} yadi evam kimaḥ at kva prepsyan dīpyase kva ardhamāsāḥ . \\
\noindent \textbf{(P\_5,3.1)} atra api prāpnoti . \\
\noindent \textbf{(P\_5,3.1)} tau ca uktam . \\
\noindent \textbf{(P\_5,3.1)} kim uktam . \\
\noindent \textbf{(P\_5,3.1)} vibhaktau tavargapratiṣedhaḥ ataddhite iti . \\
\noindent \textbf{(P\_5,3.1)} idamaḥ vibhaktisvaraḥ ca . \\
\noindent \textbf{(P\_5,3.1)} idamaḥ vibhaktisvaraḥ ca prayojanam . \\
\noindent \textbf{(P\_5,3.1)} itaḥ . \\
\noindent \textbf{(P\_5,3.1)} iha . \\
\noindent \textbf{(P\_5,3.1)} idamaḥ tṛtīyādiḥ vibhaktiḥ udāttā bhavati iti eṣaḥ svaraḥ bhavati . \\
\noindent \textbf{(P\_5,3.1)} tyadādividhayaḥ ca . \\
\noindent \textbf{(P\_5,3.1)} tyadādividhayaḥ ca prayojanam . \\
\noindent \textbf{(P\_5,3.1)} yataḥ . \\
\noindent \textbf{(P\_5,3.1)} yatra . \\
\noindent \textbf{(P\_5,3.1)} vibhaktau iti tyadādividhayaḥ siddhāḥ bhavanti . \\
\noindent \textbf{(P\_5,3.2)} bahugrahaṇe saṅkhyāgrahaṇam . \\
\noindent \textbf{(P\_5,3.2)} bahugrahaṇe saṅkhyāgrahaṇam kartavyam . \\
\noindent \textbf{(P\_5,3.2)} iha mā bhūt . \\
\noindent \textbf{(P\_5,3.2)} bahau . \\
\noindent \textbf{(P\_5,3.2)} bahoḥ iti . \\
\noindent \textbf{(P\_5,3.2)} atha kimartham kimaḥ upasaṅkhyānam kriyate na sarvanāmnaḥ iti eva siddham . \\
\noindent \textbf{(P\_5,3.2)} dvyātipratiṣedhāt kimaḥ upasaṅkhyānam . \\
\noindent \textbf{(P\_5,3.2)} dvyātipratiṣedhāt kimaḥ upasaṅkhyānam kriyate . \\
\noindent \textbf{(P\_5,3.2)} advyādibhyaḥ iti pratiṣedhe prāpte kimaḥ upasaṅkhyānam kriyate . \\
\noindent \textbf{(P\_5,3.5.1)} kva ayam nakāraḥ śrūyate . \\
\noindent \textbf{(P\_5,3.5.1)} na kva cit śrūyate . \\
\noindent \textbf{(P\_5,3.5.1)} lopaḥ asya bhavati nalopaḥ prātipadikāntasya iti . \\
\noindent \textbf{(P\_5,3.5.1)} yadi na kva cit śrūyate kimartham uccāryate . \\
\noindent \textbf{(P\_5,3.5.1)} anekālśit sarvasya iti sarvādeśaḥ yathā syāt . \\
\noindent \textbf{(P\_5,3.5.1)} kriyamāṇe api nakāre sarvādeśaḥ na prāpnoti . \\
\noindent \textbf{(P\_5,3.5.1)} kim kāraṇam . \\
\noindent \textbf{(P\_5,3.5.1)} nalope kṛte ekāltvāt . \\
\noindent \textbf{(P\_5,3.5.1)} idam iha sampradhāryam . \\
\noindent \textbf{(P\_5,3.5.1)} nalopaḥ kriyatām sarvādeśaḥ iti . \\
\noindent \textbf{(P\_5,3.5.1)} kim atra kartavyam . \\
\noindent \textbf{(P\_5,3.5.1)} paratvāt nalopaḥ . \\
\noindent \textbf{(P\_5,3.5.1)} asiddhaḥ nalopaḥ . \\
\noindent \textbf{(P\_5,3.5.1)} tasya asiddhatvāt sarvādeśaḥ bhavati . \\
\noindent \textbf{(P\_5,3.5.1)} parigaṇiteṣu kāryeṣu nalopaḥ asiddhaḥ na ca idam tatra parigaṇyate . \\
\noindent \textbf{(P\_5,3.5.1)} evam tarhi ānupūrvyā siddham etat . \\
\noindent \textbf{(P\_5,3.5.1)} na akṛte sarvādeśe prātipadikasañjñā prāpnoti na ca akṛtāyām prātipadikasañjñāyām nalopaḥ prāpnoti . \\
\noindent \textbf{(P\_5,3.5.1)} tat ānupūrvyā siddham . \\
\noindent \textbf{(P\_5,3.5.1)} na etat asti prayojanam . \\
\noindent \textbf{(P\_5,3.5.1)} alaḥ antyasya vidhayaḥ bhavanti iti akārasya akāravacane prayojanam na asti iti kṛtvā antareṇa nakāram sarvādeśaḥ bhaviṣyati . \\
\noindent \textbf{(P\_5,3.5.1)} asti anyat akārasya akāravacane prayojanam . \\
\noindent \textbf{(P\_5,3.5.1)} kim . \\
\noindent \textbf{(P\_5,3.5.1)} ye anye akārādeśāḥ prāpnuvanti tadbādhanārtham . \\
\noindent \textbf{(P\_5,3.5.1)} tat yathā maḥ rāji samaḥ kvau iti makārasya makāravacanasāmarthyāt anusvārādayaḥ na bhavanti . \\
\noindent \textbf{(P\_5,3.5.1)} tasmāt nakāraḥ kartavyaḥ . \\
\noindent \textbf{(P\_5,3.5.1)} na kartavyaḥ . \\
\noindent \textbf{(P\_5,3.5.1)} kriyate nyāse eva .praśliṣṭanirdeśaḥ ayam a a a iti . \\
\noindent \textbf{(P\_5,3.5.1)} saḥ anekālśit sarvasya iti sarvādeśaḥ bhaviṣyati . \\
\noindent \textbf{(P\_5,3.5.2)} etadaḥ iti yogavibhāgaḥ . \\
\noindent \textbf{(P\_5,3.5.2)} etadaḥ iti yogavibhāgaḥ kartavyaḥ . \\
\noindent \textbf{(P\_5,3.5.2)} etadaḥ eta it iti etau ādeśau bhavataḥ tataḥ an . \\
\noindent \textbf{(P\_5,3.5.2)} an ca bhavati etadaḥ iti . \\
\noindent \textbf{(P\_5,3.5.2)} kena vihitena thakāre etadaḥ ādeśaḥ ucyate . \\
\noindent \textbf{(P\_5,3.5.2)} etadaḥ ca thamaḥ upasaṅkhyānam . \\
\noindent \textbf{(P\_5,3.5.2)} etadaḥ ca thamaḥ upasaṅkhyānam kartavyam . \\
\noindent \textbf{(P\_5,3.5.2)} etatprakāram ittham . \\
\noindent \textbf{(P\_5,3.5.2)} tat tarhi upasaṅkhyānam kartavyam . \\
\noindent \textbf{(P\_5,3.5.2)} na kartavyam . \\
\noindent \textbf{(P\_5,3.5.2)} etat jñāpayati bhavati atra thamuḥ iti yat ayam thakārādau ādeśam śāsti . \\
\noindent \textbf{(P\_5,3.5.2)} kutaḥ nu khalu etajjñāpakāt atra thamuḥ bhaviṣyati . \\
\noindent \textbf{(P\_5,3.5.2)} na punaḥ yaḥ eva asau aviśeṣavihitaḥ thakārādiḥ tasmin ādeśaḥ syāt . \\
\noindent \textbf{(P\_5,3.5.2)} idamā thakārādim viśeṣayiṣyāmaḥ . \\
\noindent \textbf{(P\_5,3.5.2)} idamaḥ yaḥ thakārādiḥ iti . \\
\noindent \textbf{(P\_5,3.7,} 10) KA\_II,404.3-23 Ro\_IV,180-182 {1/59}     idam vicāryate : ime tasilādayaḥ vibhaktyādeśaḥ vā syuḥ pare vā iti . \\
\noindent \textbf{(P\_5,3.7,} 10) KA\_II,404.3-23 Ro\_IV,180-182 {2/59}     katham ca ādeśaḥ syuḥ katham vā pare . \\
\noindent \textbf{(P\_5,3.7,} 10) KA\_II,404.3-23 Ro\_IV,180-182 {3/59}     yadi pañcamyāḥ saptamyāḥ iti ṣaṣṭhī tadā ādeśāḥ . \\
\noindent \textbf{(P\_5,3.7,} 10) KA\_II,404.3-23 Ro\_IV,180-182 {4/59}     atha pañcamī tataḥ pare . \\
\noindent \textbf{(P\_5,3.7,} 10) KA\_II,404.3-23 Ro\_IV,180-182 {5/59}     kutaḥ sandehaḥ . \\
\noindent \textbf{(P\_5,3.7,} 10) KA\_II,404.3-23 Ro\_IV,180-182 {6/59}     samānaḥ nirdeśaḥ . \\
\noindent \textbf{(P\_5,3.7,} 10) KA\_II,404.3-23 Ro\_IV,180-182 {7/59}     kaḥ ca atra viśeṣaḥ . \\
\noindent \textbf{(P\_5,3.7,} 10) KA\_II,404.3-23 Ro\_IV,180-182 {8/59}     tasilādayaḥ vibhaktyādeśāḥ cet subluksvaraguṇadīrghaittvauttvasmāyādividhipratiṣedhaḥ . \\
\noindent \textbf{(P\_5,3.7,} 10) KA\_II,404.3-23 Ro\_IV,180-182 {9/59}     tasilādayaḥ vibhaktyādeśāḥ cet subluksvaraguṇadīrghaittvauttvasmāyādividhipratiṣedhaḥ vaktavyaḥ . \\
\noindent \textbf{(P\_5,3.7,} 10) KA\_II,404.3-23 Ro\_IV,180-182 {10/59}     subluk . \\
\noindent \textbf{(P\_5,3.7,} 10) KA\_II,404.3-23 Ro\_IV,180-182 {11/59}     tatastyaḥ . \\
\noindent \textbf{(P\_5,3.7,} 10) KA\_II,404.3-23 Ro\_IV,180-182 {12/59}     yatastyaḥ . \\
\noindent \textbf{(P\_5,3.7,} 10) KA\_II,404.3-23 Ro\_IV,180-182 {13/59}     tatratyaḥ . \\
\noindent \textbf{(P\_5,3.7,} 10) KA\_II,404.3-23 Ro\_IV,180-182 {14/59}     yatratyaḥ . \\
\noindent \textbf{(P\_5,3.7,} 10) KA\_II,404.3-23 Ro\_IV,180-182 {15/59}     supaḥ dhātuprātipadikayoḥ iti subluk prāpnoti . \\
\noindent \textbf{(P\_5,3.7,} 10) KA\_II,404.3-23 Ro\_IV,180-182 {16/59}     subluk . \\
\noindent \textbf{(P\_5,3.7,} 10) KA\_II,404.3-23 Ro\_IV,180-182 {17/59}     svara . \\
\noindent \textbf{(P\_5,3.7,} 10) KA\_II,404.3-23 Ro\_IV,180-182 {18/59}     yadā . \\
\noindent \textbf{(P\_5,3.7,} 10) KA\_II,404.3-23 Ro\_IV,180-182 {19/59}     tadā . \\
\noindent \textbf{(P\_5,3.7,} 10) KA\_II,404.3-23 Ro\_IV,180-182 {20/59}     anudāttau suppitau iti eṣaḥ svaraḥ prāpnoti . \\
\noindent \textbf{(P\_5,3.7,} 10) KA\_II,404.3-23 Ro\_IV,180-182 {21/59}     svara . \\
\noindent \textbf{(P\_5,3.7,} 10) KA\_II,404.3-23 Ro\_IV,180-182 {22/59}     guṇa . \\
\noindent \textbf{(P\_5,3.7,} 10) KA\_II,404.3-23 Ro\_IV,180-182 {23/59}     kasmāt . \\
\noindent \textbf{(P\_5,3.7,} 10) KA\_II,404.3-23 Ro\_IV,180-182 {24/59}     kutaḥ . \\
\noindent \textbf{(P\_5,3.7,} 10) KA\_II,404.3-23 Ro\_IV,180-182 {25/59}     gheḥ ṅiti iti guṇaḥ prāpnoti . \\
\noindent \textbf{(P\_5,3.7,} 10) KA\_II,404.3-23 Ro\_IV,180-182 {26/59}     guṇa . \\
\noindent \textbf{(P\_5,3.7,} 10) KA\_II,404.3-23 Ro\_IV,180-182 {27/59}     dīrgha . \\
\noindent \textbf{(P\_5,3.7,} 10) KA\_II,404.3-23 Ro\_IV,180-182 {28/59}     tasmin . \\
\noindent \textbf{(P\_5,3.7,} 10) KA\_II,404.3-23 Ro\_IV,180-182 {29/59}     tarhi . \\
\noindent \textbf{(P\_5,3.7,} 10) KA\_II,404.3-23 Ro\_IV,180-182 {30/59}     ataḥ dīrghaḥ yañi supi ca iti dīrghatvam prāpnoti . \\
\noindent \textbf{(P\_5,3.7,} 10) KA\_II,404.3-23 Ro\_IV,180-182 {31/59}     dīrgha . \\
\noindent \textbf{(P\_5,3.7,} 10) KA\_II,404.3-23 Ro\_IV,180-182 {32/59}     ettva. teṣu . \\
\noindent \textbf{(P\_5,3.7,} 10) KA\_II,404.3-23 Ro\_IV,180-182 {33/59}     tatra . \\
\noindent \textbf{(P\_5,3.7,} 10) KA\_II,404.3-23 Ro\_IV,180-182 {34/59}     bahuvacane jhali et iti ettvam prāpnoti . \\
\noindent \textbf{(P\_5,3.7,} 10) KA\_II,404.3-23 Ro\_IV,180-182 {35/59}     ettva . \\
\noindent \textbf{(P\_5,3.7,} 10) KA\_II,404.3-23 Ro\_IV,180-182 {36/59}     auttva . \\
\noindent \textbf{(P\_5,3.7,} 10) KA\_II,404.3-23 Ro\_IV,180-182 {37/59}     kasmin . \\
\noindent \textbf{(P\_5,3.7,} 10) KA\_II,404.3-23 Ro\_IV,180-182 {38/59}     kutra . \\
\noindent \textbf{(P\_5,3.7,} 10) KA\_II,404.3-23 Ro\_IV,180-182 {39/59}     idudbhyām aut at ca gheḥ iti auttvam prāpnoti . \\
\noindent \textbf{(P\_5,3.7,} 10) KA\_II,404.3-23 Ro\_IV,180-182 {40/59}     auttva . \\
\noindent \textbf{(P\_5,3.7,} 10) KA\_II,404.3-23 Ro\_IV,180-182 {41/59}     smāyādividhiḥ . \\
\noindent \textbf{(P\_5,3.7,} 10) KA\_II,404.3-23 Ro\_IV,180-182 {42/59}     tasmāt . \\
\noindent \textbf{(P\_5,3.7,} 10) KA\_II,404.3-23 Ro\_IV,180-182 {43/59}     tataḥ . \\
\noindent \textbf{(P\_5,3.7,} 10) KA\_II,404.3-23 Ro\_IV,180-182 {44/59}     tasmin . \\
\noindent \textbf{(P\_5,3.7,} 10) KA\_II,404.3-23 Ro\_IV,180-182 {45/59}     tatra . \\
\noindent \textbf{(P\_5,3.7,} 10) KA\_II,404.3-23 Ro\_IV,180-182 {46/59}     ṅasiṅyoḥ smātsminau it smādayaḥ prāpnuvanti . \\
\noindent \textbf{(P\_5,3.7,} 10) KA\_II,404.3-23 Ro\_IV,180-182 {47/59}     pañcamīnirdeśāt siddham . \\
\noindent \textbf{(P\_5,3.7,} 10) KA\_II,404.3-23 Ro\_IV,180-182 {48/59}     santu pare . \\
\noindent \textbf{(P\_5,3.7,} 10) KA\_II,404.3-23 Ro\_IV,180-182 {49/59}     yadi pare samānaśabdānām pratiṣedhaḥ vaktavyaḥ . \\
\noindent \textbf{(P\_5,3.7,} 10) KA\_II,404.3-23 Ro\_IV,180-182 {50/59}     tasmāt tasyati . \\
\noindent \textbf{(P\_5,3.7,} 10) KA\_II,404.3-23 Ro\_IV,180-182 {51/59}     yasmāt tasyati . \\
\noindent \textbf{(P\_5,3.7,} 10) KA\_II,404.3-23 Ro\_IV,180-182 {52/59}     pañcamyantasya taseḥ tasil bhavati iti tasil prāpnoti . \\
\noindent \textbf{(P\_5,3.7,} 10) KA\_II,404.3-23 Ro\_IV,180-182 {53/59}     anādeśe svārthavijñānāt samānaśabdāpratiṣedhaḥ . \\
\noindent \textbf{(P\_5,3.7,} 10) KA\_II,404.3-23 Ro\_IV,180-182 {54/59}     anādeśe svārthavijñānāt samānaśabdāpratiṣedhaḥ . \\
\noindent \textbf{(P\_5,3.7,} 10) KA\_II,404.3-23 Ro\_IV,180-182 {55/59}     anarthakaḥ pratiṣedhaḥ apratiṣedhaḥ . \\
\noindent \textbf{(P\_5,3.7,} 10) KA\_II,404.3-23 Ro\_IV,180-182 {56/59}     tasil kasmāt na bhavati . \\
\noindent \textbf{(P\_5,3.7,} 10) KA\_II,404.3-23 Ro\_IV,180-182 {57/59}     svārthavijñānāt . \\
\noindent \textbf{(P\_5,3.7,} 10) KA\_II,404.3-23 Ro\_IV,180-182 {58/59}     pañcamyantāt parasya taseḥ svārthe vartamānasya tasilā bhavitavyam . \\
\noindent \textbf{(P\_5,3.7,} 10) KA\_II,404.3-23 Ro\_IV,180-182 {59/59}     na ca atra pañcamyantāt paraḥ tasiḥ svārthe vartate . \\
\noindent \textbf{(P\_5,3.8)} kimartham taseḥ tasil ucyate . \\
\noindent \textbf{(P\_5,3.8)} taseḥ tasilvacanam svarārtham . \\
\noindent \textbf{(P\_5,3.8)} taseḥ tasilvacanam kriyate svarārtham . \\
\noindent \textbf{(P\_5,3.8)} liti pratyayāt pūrvam udāttam bhavati iti eṣaḥ svaraḥ yathā syāt . \\
\noindent \textbf{(P\_5,3.8)} nanu ca ayam tasil tasim bādhiṣyate . \\
\noindent \textbf{(P\_5,3.8)} na sidhyati . \\
\noindent \textbf{(P\_5,3.8)} paratvāt tasiḥ prāpnoti . \\
\noindent \textbf{(P\_5,3.8)} tasilaḥ avakāśaḥ . \\
\noindent \textbf{(P\_5,3.8)} tataḥ hīyate . \\
\noindent \textbf{(P\_5,3.8)} tataḥ avarohati . \\
\noindent \textbf{(P\_5,3.8)} taseḥ avakāśaḥ . \\
\noindent \textbf{(P\_5,3.8)} grāmataḥ āgacchati . \\
\noindent \textbf{(P\_5,3.8)} nagarataḥ āgacchati . \\
\noindent \textbf{(P\_5,3.8)} iha ubhayam prāpnoti . \\
\noindent \textbf{(P\_5,3.8)} tataḥ āgacchati . \\
\noindent \textbf{(P\_5,3.8)} yataḥ āgacchati . \\
\noindent \textbf{(P\_5,3.8)} paratvāt tasiḥ prāpnoti . \\
\noindent \textbf{(P\_5,3.8)} tasmāt suṣthu ucyate taseḥ tasilvacanam svarārtham iti . \\
\noindent \textbf{(P\_5,3.9)} paryabhibhyām ca iti yat ucyate tat sarvobhayārthe draṣṭavyam . \\
\noindent \textbf{(P\_5,3.9)} yāvat sarvataḥ tāvat paritaḥ . \\
\noindent \textbf{(P\_5,3.9)} yāvat ubhayataḥ tāvat abhitaḥ . \\
\noindent \textbf{(P\_5,3.14)} iha kasmāt na bhavati . \\
\noindent \textbf{(P\_5,3.14)} saḥ . \\
\noindent \textbf{(P\_5,3.14)} tau . \\
\noindent \textbf{(P\_5,3.14)} te . \\
\noindent \textbf{(P\_5,3.14)} bhavadādibhiḥ yoge iti vaktavyam . \\
\noindent \textbf{(P\_5,3.14)} ke punaḥ bhavadādayaḥ . \\
\noindent \textbf{(P\_5,3.14)} bhavān . \\
\noindent \textbf{(P\_5,3.14)} dīrghāyuḥ . \\
\noindent \textbf{(P\_5,3.14)} devānāmpriyaḥ . \\
\noindent \textbf{(P\_5,3.14)} āyuṣmān iti . \\
\noindent \textbf{(P\_5,3.14)} saḥ bhavān . \\
\noindent \textbf{(P\_5,3.14)} tatra bhavān . \\
\noindent \textbf{(P\_5,3.14)} tataḥ bhavān . \\
\noindent \textbf{(P\_5,3.14)} tam bhavantam . \\
\noindent \textbf{(P\_5,3.14)} tatra bhavantam . \\
\noindent \textbf{(P\_5,3.14)} tataḥ bhavantam . \\
\noindent \textbf{(P\_5,3.14)} tena bhavatā . \\
\noindent \textbf{(P\_5,3.14)} tatra bhavatā tataḥ bhavatā . \\
\noindent \textbf{(P\_5,3.14)} tasmai bhavate . \\
\noindent \textbf{(P\_5,3.14)} tatra bhavate . \\
\noindent \textbf{(P\_5,3.14)} tataḥ bhavate . \\
\noindent \textbf{(P\_5,3.14)} tasmāt bhavataḥ . \\
\noindent \textbf{(P\_5,3.14)} tatra bhavataḥ . \\
\noindent \textbf{(P\_5,3.14)} tataḥ bhavataḥ . \\
\noindent \textbf{(P\_5,3.14)} tasmin bhavati . \\
\noindent \textbf{(P\_5,3.14)} tatra bhavati . \\
\noindent \textbf{(P\_5,3.14)} tataḥ bhavati . \\
\noindent \textbf{(P\_5,3.14)} saḥ dīrghāyuḥ . \\
\noindent \textbf{(P\_5,3.14)} tatra dīrghāyuḥ . \\
\noindent \textbf{(P\_5,3.14)} tataḥ dīrghāyuḥ . \\
\noindent \textbf{(P\_5,3.14)} tam dīrghāyuṣam . \\
\noindent \textbf{(P\_5,3.14)} tatra dīrghāyuṣam . \\
\noindent \textbf{(P\_5,3.14)} tataḥ dīrghāyuṣam . \\
\noindent \textbf{(P\_5,3.14)} saḥ devānāmpriyaḥ . \\
\noindent \textbf{(P\_5,3.14)} tatra devānāmpriyaḥ . \\
\noindent \textbf{(P\_5,3.14)} tataḥ devānāmpriyaḥ . \\
\noindent \textbf{(P\_5,3.14)} tam devānāmpriyam . \\
\noindent \textbf{(P\_5,3.14)} tatra devānāmpriyam . \\
\noindent \textbf{(P\_5,3.14)} tataḥ devānāmpriyam . \\
\noindent \textbf{(P\_5,3.14)} saḥ āyuṣmān . \\
\noindent \textbf{(P\_5,3.14)} tatra āyuṣmān . \\
\noindent \textbf{(P\_5,3.14)} tataḥ āyuṣmān . \\
\noindent \textbf{(P\_5,3.14)} tam āyuṣmantam . \\
\noindent \textbf{(P\_5,3.14)} tatra āyuṣmantam . \\
\noindent \textbf{(P\_5,3.14)} tataḥ āyuṣmantam \\
\noindent \textbf{(P\_5,3.17)} adhunā iti kim nipātyate . \\
\noindent \textbf{(P\_5,3.17)} idamaḥ aśbhāvaḥ dhunā ca pratyayaḥ idamaḥ vā lopaḥ adhunā ca pratyayaḥ . \\
\noindent \textbf{(P\_5,3.17)} asmin kāle adhuna . \\
\noindent \textbf{(P\_5,3.18)} idānīm . \\
\noindent \textbf{(P\_5,3.18)} idamaḥ tṛtīyādivibhaktiḥ udāttā bhavati iti eṣaḥ svaraḥ prāpnoti . \\
\noindent \textbf{(P\_5,3.18)} dānīm iti nipātanāt svarasiddhiḥ . \\
\noindent \textbf{(P\_5,3.18)} dānīm iti nipātanāt svarasiddhiḥ bhaviṣyati . \\
\noindent \textbf{(P\_5,3.18)} ādyudāttanipātanam kariṣyate . \\
\noindent \textbf{(P\_5,3.18)} saḥ nipātanasvaraḥ vibhaktisvarasya bādhakaḥ bhaviṣyati . \\
\noindent \textbf{(P\_5,3.18)} uktam vā . \\
\noindent \textbf{(P\_5,3.18)} kim uktam . \\
\noindent \textbf{(P\_5,3.18)} ādau siddham iti . \\
\noindent \textbf{(P\_5,3.19)} tadaḥ dāvacanam anarthakam vihitatvāt . \\
\noindent \textbf{(P\_5,3.19)} tadaḥ dāvacanam anarthakam . \\
\noindent \textbf{(P\_5,3.19)} kim kāraṇam . \\
\noindent \textbf{(P\_5,3.19)} vihitatvāt . \\
\noindent \textbf{(P\_5,3.19)} vihitaḥ atra pratyayaḥ sarvaikānyakiṃyattadaḥ kāle dā iti . \\
\noindent \textbf{(P\_5,3.20)} tayoḥ iti prātipadikanirdeśaḥ . \\
\noindent \textbf{(P\_5,3.20)} tayoḥ iti prātipadikanirdeśaḥ draṣṭavyaḥ . \\
\noindent \textbf{(P\_5,3.20)} dveṣyam vijānīyāt :yogayoḥ vā pratyayayoḥ vā iti . \\
\noindent \textbf{(P\_5,3.20)} tat ācāryaḥ suhṛt bhūtvā anvācaṣṭe : tayoḥ iti prātipadikanirdeśaḥ iti . \\
\noindent \textbf{(P\_5,3.22)} sadyaḥ iti kim nipātyate . \\
\noindent \textbf{(P\_5,3.22)} samānasya sabhāvaḥ dyaḥ ca ahani . \\
\noindent \textbf{(P\_5,3.22)} samānasya sabhāvaḥ nipātyate dyaḥ ca pratyayaḥ ahani abhidheye . \\
\noindent \textbf{(P\_5,3.22)} samāne ahani sadyaḥ . \\
\noindent \textbf{(P\_5,3.22)} parut parāri iti kim nipātyate . \\
\noindent \textbf{(P\_5,3.22)} pūrvapūrvatayoḥ parabhāvaḥ udārī ca saṃvatsare . \\
\noindent \textbf{(P\_5,3.22)} pūrvapūrvatayoḥ parabhāvaḥ nipātyate udārī ca pratyayau saṃvatsare abhidheye . \\
\noindent \textbf{(P\_5,3.22)} pūrvasmin saṃvatsare parut . \\
\noindent \textbf{(P\_5,3.22)} pūrvatare saṃvatsare parāri . \\
\noindent \textbf{(P\_5,3.22)} aiṣamaḥ iti kim nipātyate . \\
\noindent \textbf{(P\_5,3.22)} idamaḥ samasaṇ . \\
\noindent \textbf{(P\_5,3.22)} idamaḥ samasaṇ pratyayaḥ nipātyate saṃvatsare abhidheye . \\
\noindent \textbf{(P\_5,3.22)} asmin saṃvatsare aiṣamaḥ . \\
\noindent \textbf{(P\_5,3.22)} paredyavi iti kim nipātyate . \\
\noindent \textbf{(P\_5,3.22)} parasmāt edyavi ahani . \\
\noindent \textbf{(P\_5,3.22)} parasmāt edyavi pratyayaḥ nipātyate ahani abhidheye . \\
\noindent \textbf{(P\_5,3.22)} parasmin ahani paredyavi . \\
\noindent \textbf{(P\_5,3.22)} adya iti kim nipātyate . \\
\noindent \textbf{(P\_5,3.22)} idamaḥ aśbhāvaḥ dyaḥ ca . \\
\noindent \textbf{(P\_5,3.22)} idamaḥ aśbhāvaḥ nipātyate dyaḥ ca pratyayaḥ ahani abhidheye . \\
\noindent \textbf{(P\_5,3.22)} asmin ahani adya . \\
\noindent \textbf{(P\_5,3.22)} pūrvedyuḥ anyedyuḥ anyataredyuḥ itaredyuḥ aparedyuḥ adharedyuḥ ubhayedyuḥ uttaredyuḥ iti kim nipātyate . \\
\noindent \textbf{(P\_5,3.22)} pūrvānyānyatarerāparādharobhayottarebhyaḥ edyusuc . \\
\noindent \textbf{(P\_5,3.22)} pūrvānyānyatarerāparādharobhayottarebhyaḥ edyusuc pratyayaḥ nipātyate ahani abhidheye . \\
\noindent \textbf{(P\_5,3.22)} pūrvasmin ahani pūrvedyuḥ . \\
\noindent \textbf{(P\_5,3.22)} anyasmin ahani anyedyuḥ . \\
\noindent \textbf{(P\_5,3.22)} anyatarasmin ahani anyataredyuḥ . \\
\noindent \textbf{(P\_5,3.22)} itarasmin ahani itaredyuḥ . \\
\noindent \textbf{(P\_5,3.22)} aparasmin ahani aparedyuḥ . \\
\noindent \textbf{(P\_5,3.22)} adharasmin ahani adharedyuḥ . \\
\noindent \textbf{(P\_5,3.22)} ubhayoḥ ahnoḥ ubhayedyuḥ . \\
\noindent \textbf{(P\_5,3.22)} uttarasmin ahani uttaredyuḥ . \\
\noindent \textbf{(P\_5,3.22)} dyuḥ ca ubhayāt . ubhayaśabdāt dyuḥ ca vaktavyaḥ . \\
\noindent \textbf{(P\_5,3.22)} tasmāt manuṣyebhyaḥ ubhayadyuḥ . \\
\noindent \textbf{(P\_5,3.27)} iha kasmāt na bhavati . \\
\noindent \textbf{(P\_5,3.27)} pūrvasmin deśe vasati iti . \\
\noindent \textbf{(P\_5,3.27)} na eṣaḥ deśaḥ . \\
\noindent \textbf{(P\_5,3.27)} deśaviśeṣaṇam etat . \\
\noindent \textbf{(P\_5,3.28)} kimartham atasuc kriyate na tasuc eva kriyate . \\
\noindent \textbf{(P\_5,3.28)} tatra ayam api arthaḥ . \\
\noindent \textbf{(P\_5,3.28)} svarārthaḥ cakāraḥ na kartavyaḥ bhavati . \\
\noindent \textbf{(P\_5,3.28)} pratyayasvareṇa eva siddham . \\
\noindent \textbf{(P\_5,3.28)} kā rūpasiddhiḥ : dakṣiṇataḥ grāmasya . \\
\noindent \textbf{(P\_5,3.28)} uttarataḥ grāmasya . \\
\noindent \textbf{(P\_5,3.28)} dakṣiṇottaraśabdau akārāntau . \\
\noindent \textbf{(P\_5,3.28)} tasuśabdaḥ pratyayaḥ . \\
\noindent \textbf{(P\_5,3.28)} bhavet siddham yadā akārāntau . \\
\noindent \textbf{(P\_5,3.28)} yatu tu khalu ākārāntau tadā na sidhyati . \\
\noindent \textbf{(P\_5,3.28)} tadā api siddham . \\
\noindent \textbf{(P\_5,3.28)} katham . \\
\noindent \textbf{(P\_5,3.28)} puṃvadbhāvena . \\
\noindent \textbf{(P\_5,3.28)} katham puṃvadbhāvaḥ . \\
\noindent \textbf{(P\_5,3.28)} tasilādiṣu ā kṛtvasucaḥ iti . \\
\noindent \textbf{(P\_5,3.28)} na sidhyati . \\
\noindent \textbf{(P\_5,3.28)} bhāṣitapuṃskasya puṃvadbhāvaḥ na ca etau bhāṣitapuṃskau . \\
\noindent \textbf{(P\_5,3.28)} nanu ca bho dakṣiṇaśabdaḥ uttaraśabdaḥ ca puṃsi bhāṣyete . \\
\noindent \textbf{(P\_5,3.28)} samānāyām ākṛtau yat bhāṣitapuṃskam iti ucyate ākṛtyantare ca etau bhāṣitapuṃskau . \\
\noindent \textbf{(P\_5,3.28)} dakṣiṇā uttarā iti dikśabdau . \\
\noindent \textbf{(P\_5,3.28)} dakṣiṇaḥ uttaraḥ iti vyavasthāśabdau . \\
\noindent \textbf{(P\_5,3.28)} yadi punaḥ dikśabdāḥ api vyavasthāśabdāḥ syuḥ . \\
\noindent \textbf{(P\_5,3.28)} katham yāni digapadiṣṭāni kāryāṇi . \\
\noindent \textbf{(P\_5,3.28)} diśaḥ yadā vyavasthām vakṣyanti . \\
\noindent \textbf{(P\_5,3.28)} yadi tari yaḥ yaḥ diśi vartate saḥ saḥ dikśabdaḥ ramaṇīyādiṣu atiprasaṅgaḥ bhavati . \\
\noindent \textbf{(P\_5,3.28)} ramaṇīyā dik śobhanā dik iti . \\
\noindent \textbf{(P\_5,3.28)} atha matam etat diśi dṛṣṭaḥ digdṛṣṭaḥ digdṛṣṭaḥ śabdaḥ dikśabdaḥ diśam yaḥ na vyabhicarati iti ramaṇīyādiṣu atiprasaṅgaḥ na bhavati . \\
\noindent \textbf{(P\_5,3.28)} puṃvadbhāvaḥ tu prāpnoti . \\
\noindent \textbf{(P\_5,3.28)} evam tarhi sarvanāmnaḥ vṛttimātre puṃvadbhāvaḥ vaktavyaḥ dakṣiṇottarapūrvāṇām iti evamartham . \\
\noindent \textbf{(P\_5,3.28)} viśeṣaṇārtham tarhi . \\
\noindent \textbf{(P\_5,3.28)} kva viśeṣaṇārthena arthaḥ . \\
\noindent \textbf{(P\_5,3.28)} ṣaṣṭhī atasarthapratyayena iti . \\
\noindent \textbf{(P\_5,3.28)} ṣaṣṭhī tasarthapratyayene iti ucyamāne iha api syāt . \\
\noindent \textbf{(P\_5,3.28)} tataḥ grāmāt . \\
\noindent \textbf{(P\_5,3.28)} yataḥ grāmāt iti . \\
\noindent \textbf{(P\_5,3.31)} upari upariṣṭāt iti kim nipātyate . \\
\noindent \textbf{(P\_5,3.31)} ūrdhvasya upabhāvaḥ riliṣṭātilau ca . \\
\noindent \textbf{(P\_5,3.31)} ūrdhvasya upabhāvaḥ riliṣṭātilau ca pratyayau nipātyete . \\
\noindent \textbf{(P\_5,3.31)} upari upariṣtāt . \\
\noindent \textbf{(P\_5,3.32)} paścāt iti kim nipātyate . \\
\noindent \textbf{(P\_5,3.32)} aparasya paścabhāvaḥ ātiḥ ca pratayayaḥ . \\
\noindent \textbf{(P\_5,3.32)} aparasya paścabhāvaḥ nipātyate ātiḥ ca pratayayaḥ . \\
\noindent \textbf{(P\_5,3.32)} paścāt . \\
\noindent \textbf{(P\_5,3.32)} dikpūrvapadasya ca . \\
\noindent \textbf{(P\_5,3.32)} dikpūrvapadasya ca aparasya paścabhāvaḥ vaktavyaḥ ātiḥ ca pratayayaḥ . \\
\noindent \textbf{(P\_5,3.32)} dakṣiṇapaścāt . \\
\noindent \textbf{(P\_5,3.32)} uttarapaścāt . \\
\noindent \textbf{(P\_5,3.32)} ardhottarapadasya ca samāse . \\
\noindent \textbf{(P\_5,3.32)} ardhottarapadasya ca samāse aparasya paścabhāvaḥ vaktavyaḥ . \\
\noindent \textbf{(P\_5,3.32)} dakṣiṇapaścārdhaḥ . \\
\noindent \textbf{(P\_5,3.32)} uttarapaścārdhaḥ . \\
\noindent \textbf{(P\_5,3.32)} ardhe ca . \\
\noindent \textbf{(P\_5,3.32)} ardhe ca parataḥ aparasya paścabhāvaḥ vaktavyaḥ . \\
\noindent \textbf{(P\_5,3.32)} paścārdhaḥ . \\
\noindent \textbf{(P\_5,3.35)} apañcamyāḥ iti prāk asaḥ . \\
\noindent \textbf{(P\_5,3.35)} apañcamyāḥ iti yat ucyate prāk asaḥ tat draṣṭavyam . \\
\noindent \textbf{(P\_5,3.35)} dveṣyam vijānīyāt : aviśeṣeṇa uttaram apañcamyāḥ iti . \\
\noindent \textbf{(P\_5,3.35)} tat ācāryaḥ suhṛt bhūtvā anvācaṣṭe : apañcamyāḥ iti prāk asaḥ iti . \\
\noindent \textbf{(P\_5,3.36)} kimarthaḥ cakāraḥ . \\
\noindent \textbf{(P\_5,3.36)} svarārthaḥ . \\
\noindent \textbf{(P\_5,3.36)} citaḥ antaḥ udāttaḥ bhavati iti antodāttatvam yathā syāt . \\
\noindent \textbf{(P\_5,3.36)} na etat asti prayojanam . \\
\noindent \textbf{(P\_5,3.36)} ekāc ayam . \\
\noindent \textbf{(P\_5,3.36)} tatra na arthaḥ svarārthena cakāreṇa anubandhena . \\
\noindent \textbf{(P\_5,3.36)} pratyayasvareṇa eva siddham . \\
\noindent \textbf{(P\_5,3.36)} viśeṣaṇārthaḥ tarhi . \\
\noindent \textbf{(P\_5,3.36)} kva viśeṣaṇārthena arthaḥ . \\
\noindent \textbf{(P\_5,3.36)} anyārāditarartedikśabdāñcūūttarapadājāhiyukte . \\
\noindent \textbf{(P\_5,3.42)} vidhārthe iti ucyate . \\
\noindent \textbf{(P\_5,3.42)} kaḥ vidhārthaḥ nāma . \\
\noindent \textbf{(P\_5,3.42)} vidhāyāḥ arthaḥ vidhārthaḥ . \\
\noindent \textbf{(P\_5,3.42)} yadi evam ekā govidhā . \\
\noindent \textbf{(P\_5,3.42)} ekā hastividhā . \\
\noindent \textbf{(P\_5,3.42)} atra api prāpnoti . \\
\noindent \textbf{(P\_5,3.42)} evam tarhi dhāvidhānam dhātvarthapṛthagbhāve . \\
\noindent \textbf{(P\_5,3.42)} dhāvidhānam dhātvarthapṛthagbhāve iti vaktavyam . \\
\noindent \textbf{(P\_5,3.42)} kaḥ punaḥ dhātvarthapṛthagbhāvaḥ . \\
\noindent \textbf{(P\_5,3.42)} kim yat tat devadattaḥ kaṃsapātryām pāṇinā odanam bhuṅkte iti . \\
\noindent \textbf{(P\_5,3.42)} na iti āha . \\
\noindent \textbf{(P\_5,3.42)} kārakapṛthaktvam etat . \\
\noindent \textbf{(P\_5,3.42)} yat tarhi tat kālye bhuṅke sāyam bhuṅkte iti . \\
\noindent \textbf{(P\_5,3.42)} na iti āha . \\
\noindent \textbf{(P\_5,3.42)} kālapṛthaktvam etat . \\
\noindent \textbf{(P\_5,3.42)} yat tarhi śītam bhuṅkte uṣṇam bhuṅkte iti . \\
\noindent \textbf{(P\_5,3.42)} na iti āha . \\
\noindent \textbf{(P\_5,3.42)} guṇapṛthaktvam etat . \\
\noindent \textbf{(P\_5,3.42)} kaḥ tarhi dhātvarthapṛthagbhāvaḥ . \\
\noindent \textbf{(P\_5,3.42)} kārakāṇām pravṛttiviśeṣaḥ kriyā . \\
\noindent \textbf{(P\_5,3.42)} yadi evam kriyāprakāre ayam bhavati . \\
\noindent \textbf{(P\_5,3.42)} vidhayuktagatāḥ ca prakāre bhavanti . \\
\noindent \textbf{(P\_5,3.42)} evaṃvidham . \\
\noindent \textbf{(P\_5,3.42)} evaṃyuktam . \\
\noindent \textbf{(P\_5,3.42)} evaṅgatam . \\
\noindent \textbf{(P\_5,3.42)} evamprakāram iti . \\
\noindent \textbf{(P\_5,3.44)} sahabhāve dhyamuñ . \\
\noindent \textbf{(P\_5,3.44)} sahabhāve dhyamuñ vaktavyaḥ . \\
\noindent \textbf{(P\_5,3.44)} eikadhyam rāśim kuru . \\
\noindent \textbf{(P\_5,3.44)} saḥ tarhi vaktavyaḥ . \\
\noindent \textbf{(P\_5,3.44)} na vaktavyaḥ . \\
\noindent \textbf{(P\_5,3.44)} adhikaraṇavicāle iti ucyate na ca saḥ eva adhikaraṇavicālaḥ yat ekam anekam kriyate . \\
\noindent \textbf{(P\_5,3.44)} yat api anekam ekam kriyate saḥ api adhikaraṇavicālaḥ . \\
\noindent \textbf{(P\_5,3.45)} dhamuñantāt svārthe ḍadarśanam . \\
\noindent \textbf{(P\_5,3.45)} dhamuñantāt svārthe ḍaḥ dṛśyate saḥ ca vidheyaḥ . \\
\noindent \textbf{(P\_5,3.45)} pathi dvaidhāni . \\
\noindent \textbf{(P\_5,3.45)} saṃśaye dvaidhāni . \\
\noindent \textbf{(P\_5,3.47)} pāśapi kutsitagrahaṇam . \\
\noindent \textbf{(P\_5,3.47)} pāśapi kutsitagrahaṇam kartavyam . \\
\noindent \textbf{(P\_5,3.47)} vaiyākaraṇapāśaḥ . \\
\noindent \textbf{(P\_5,3.47)} yājñikapāśaḥ . \\
\noindent \textbf{(P\_5,3.47)} yaḥ hi yāpayitavyaḥ yāpyaḥ tatra mā bhūt iti . \\
\noindent \textbf{(P\_5,3.47)} atha vaiyākaraṇaḥ śarīreṇa kṛśaḥ vyākaraṇena ca śobhanaḥ kartavyaḥ vaiyākaraṇapāśaḥ iti . \\
\noindent \textbf{(P\_5,3.47)} na kartavyaḥ . \\
\noindent \textbf{(P\_5,3.47)} katham . \\
\noindent \textbf{(P\_5,3.47)} yasya bhāvāt dravye śabdaniveśaḥ tadabhidhāne tadguṇe vaktavye pratyayena bhavitavyam . \\
\noindent \textbf{(P\_5,3.47)} na ca kārśyasya bhāvāt dravye vaiyākaraṇaśabdaḥ . \\
\noindent \textbf{(P\_5,3.48)} pūraṇagrahaṇam śakyam akartum . \\
\noindent \textbf{(P\_5,3.48)} na hi apūraṇaḥ tīyaśabdaḥ asti yatra doṣaḥ syāt . \\
\noindent \textbf{(P\_5,3.48)} nanu ca ayam asti mukhatīyaḥ pārśvatīyaḥ iti . \\
\noindent \textbf{(P\_5,3.48)} arthavadgrahaṇe na anarthakasya iti evam asya na bhaviṣyati . \\
\noindent \textbf{(P\_5,3.48)} uttarārtham tarhi pūraṇagrahaṇam kartavyam . \\
\noindent \textbf{(P\_5,3.48)} prāk ekādaśabhyaḥ acchandasi iti pūraṇāt yathā syāt . \\
\noindent \textbf{(P\_5,3.52)} ekāt ākinici dvibahvarthe pratyayavidhānam . \\
\noindent \textbf{(P\_5,3.52)} ekāt ākinici dvibahvarthe pratyayaḥ vidheyaḥ . \\
\noindent \textbf{(P\_5,3.52)} ekākinau . \\
\noindent \textbf{(P\_5,3.52)} ekākinaḥ iti . \\
\noindent \textbf{(P\_5,3.52)} kim punaḥ kāraṇam na sidhyati . \\
\noindent \textbf{(P\_5,3.52)} ekaśabdaḥ ayam saṅkhyāpadam saṅkhyāyāḥ ca saṅkhyeyam arthaḥ . \\
\noindent \textbf{(P\_5,3.52)} siddham tu saṅkhyādeśavacanāt . \\
\noindent \textbf{(P\_5,3.52)} siddham etat . \\
\noindent \textbf{(P\_5,3.52)} katham . \\
\noindent \textbf{(P\_5,3.52)} dvibahvarthāyāḥ saṅkhyāyāḥ ekaśabdaḥ ādeśaḥ vaktavyaḥ . \\
\noindent \textbf{(P\_5,3.52)} asahāyasya vā . \\
\noindent \textbf{(P\_5,3.52)} asahāyasya vā ekaśabdaḥ ādeśaḥ vaktavyaḥ . \\
\noindent \textbf{(P\_5,3.52)} asahāyaḥ ekākī . \\
\noindent \textbf{(P\_5,3.52)} asahāyau ekākinau . \\
\noindent \textbf{(P\_5,3.52)} asahāyāḥ ekākinaḥ . \\
\noindent \textbf{(P\_5,3.52)} sidhyati . \\
\noindent \textbf{(P\_5,3.52)} sūtram tarhi bhidyate . \\
\noindent \textbf{(P\_5,3.52)} yathānyāsam eva astu . \\
\noindent \textbf{(P\_5,3.52)} nanu ca uktam ekāt ākinici dvibahvarthe pratyayavidhānam iti . \\
\noindent \textbf{(P\_5,3.52)} na eṣaḥ doṣaḥ . \\
\noindent \textbf{(P\_5,3.52)} ayam ekaśabdaḥ asti eva saṅkhyāpadam . \\
\noindent \textbf{(P\_5,3.52)} tat yathā ekaḥ dvau bahavaḥ iti . \\
\noindent \textbf{(P\_5,3.52)} asti anyārthe vartate . \\
\noindent \textbf{(P\_5,3.52)} tat yathā sadhamādaḥ dyumnaḥ ekāḥ tāḥ . \\
\noindent \textbf{(P\_5,3.52)} anyāḥ iti arthaḥ . \\
\noindent \textbf{(P\_5,3.52)} asti asahāyavācī . \\
\noindent \textbf{(P\_5,3.52)} tat yathā . \\
\noindent \textbf{(P\_5,3.52)} ekāgnayaḥ . \\
\noindent \textbf{(P\_5,3.52)} ekahalāni . \\
\noindent \textbf{(P\_5,3.52)} ekākibhiḥ kṣudrakaiḥ jitam iti . \\
\noindent \textbf{(P\_5,3.52)} tat yaḥ asahāyavācī tasya eṣaḥ prayogaḥ . \\
\noindent \textbf{(P\_5,3.55.1)} atiśāyane iti ucyate . \\
\noindent \textbf{(P\_5,3.55.1)} kim idam atiśāyane iti . \\
\noindent \textbf{(P\_5,3.55.1)} deśyāḥ sūtranibandhāḥ kriyante . \\
\noindent \textbf{(P\_5,3.55.1)} yāvat brūyāt prakarṣe atiśaye iti tāvat atiśāyane iti . \\
\noindent \textbf{(P\_5,3.55.1)} kasya punaḥ prakarṣe pratyayaḥ utpadyate . \\
\noindent \textbf{(P\_5,3.55.1)} ṅyāpprātipadikāt iti vartate . \\
\noindent \textbf{(P\_5,3.55.1)} ṅyāpprātipadikasya prakarṣe . \\
\noindent \textbf{(P\_5,3.55.1)} ṅyāpprātipadikam vai śabdaḥ na ca śabdasya prakarṣāpakarṣau staḥ . \\
\noindent \textbf{(P\_5,3.55.1)} śabde asambhavāt arthe kāryam vijñāsyate . \\
\noindent \textbf{(P\_5,3.55.1)} kaḥ punaḥ ṅyāpprātipadikārthaḥ . \\
\noindent \textbf{(P\_5,3.55.1)} dravyam . \\
\noindent \textbf{(P\_5,3.55.1)} na vai dravyasaya prakarṣe iṣyate . \\
\noindent \textbf{(P\_5,3.55.1)} evam tarhi guṇaḥ . \\
\noindent \textbf{(P\_5,3.55.1)} evam api guṇagrahaṇam kartavyam . \\
\noindent \textbf{(P\_5,3.55.1)} dravyam api ṅyāpprātipadikārthaḥ guṇaḥ api . \\
\noindent \textbf{(P\_5,3.55.1)} tatra kutaḥ etat guṇasya prakarṣe bhaviṣyati na punaḥ dravyasya prakarṣe iti . \\
\noindent \textbf{(P\_5,3.55.1)} kriyamāṇe ca api guṇagrahaṇe samānaguṇagrahaṇam kartavyam śuklāt kṛṣṇe mā bhūt iti . \\
\noindent \textbf{(P\_5,3.55.1)} na tarhi idānīm idam bhavati : adhvaryuḥ vai śreyān . \\
\noindent \textbf{(P\_5,3.55.1)} pāpīyān pratiprasthātā . \\
\noindent \textbf{(P\_5,3.55.1)} andhānām kāṇatamaḥ iti . \\
\noindent \textbf{(P\_5,3.55.1)} samānaguṇe eṣā spardhā bhavati . \\
\noindent \textbf{(P\_5,3.55.1)} adhvaryuḥ vai śreyān anyebhyaḥ praśasyebhyaḥ . \\
\noindent \textbf{(P\_5,3.55.1)} pāpīyān pratiprasthātā anyebhyaḥ pāpebhyaḥ . \\
\noindent \textbf{(P\_5,3.55.1)} andhānām kāṇatamaḥ iti kaṇiḥ ayam saukṣmye vartate . \\
\noindent \textbf{(P\_5,3.55.1)} sarve ime kim cit paśyanti . \\
\noindent \textbf{(P\_5,3.55.1)} ayam eṣām kāṇatamaḥ iti . \\
\noindent \textbf{(P\_5,3.55.1)} adūraviprakarṣe iti vaktavyam . \\
\noindent \textbf{(P\_5,3.55.1)} iha mā bhūt . \\
\noindent \textbf{(P\_5,3.55.1)} mahān sarṣapaḥ . \\
\noindent \textbf{(P\_5,3.55.1)} mahān himavān iti . \\
\noindent \textbf{(P\_5,3.55.1)} jāteḥ na iti vaktavyam . \\
\noindent \textbf{(P\_5,3.55.1)} iha mā bhūt : vṛkṣaḥ ayam plakṣaḥ ayam iti . \\
\noindent \textbf{(P\_5,3.55.1)} na tarhi idānīm idam bhavati : gotaraḥ . \\
\noindent \textbf{(P\_5,3.55.1)} gotarā . \\
\noindent \textbf{(P\_5,3.55.1)} aśvataraḥ iti . \\
\noindent \textbf{(P\_5,3.55.1)} na eṣaḥ jāteḥ prakarṣaḥ . \\
\noindent \textbf{(P\_5,3.55.1)} kasya tarhi . \\
\noindent \textbf{(P\_5,3.55.1)} guṇasya . \\
\noindent \textbf{(P\_5,3.55.1)} gauḥ ayam śakaṭam vahati . \\
\noindent \textbf{(P\_5,3.55.1)} gotararḥ ayam yaḥ śakaṭam vahati sīram ca . \\
\noindent \textbf{(P\_5,3.55.1)} gauḥ iyam yā samām samām vijāyate . \\
\noindent \textbf{(P\_5,3.55.1)} gotarā iyam yā samām samām vijāyate strīvatsā ca . \\
\noindent \textbf{(P\_5,3.55.1)} aśvaḥ ayam yaḥ catvāri yojanāni gacchati . \\
\noindent \textbf{(P\_5,3.55.1)} aśvataraḥ ayam yaḥ aṣṭau yojanāni gacchati . \\
\noindent \textbf{(P\_5,3.55.1)} tathā tiṅaḥ ca iti atra kriyāgrahaṇam kartavyam sādhanaprakarṣe mā bhūt . \\
\noindent \textbf{(P\_5,3.55.1)} na eṣaḥ doṣaḥ . \\
\noindent \textbf{(P\_5,3.55.1)} yat tāvat ucyate guṇagrahaṇam kartavyam iti . \\
\noindent \textbf{(P\_5,3.55.1)} na kartavyam . \\
\noindent \textbf{(P\_5,3.55.1)} yasya prakarṣaḥ asti tasya prakarṣe bhaviṣyati . \\
\noindent \textbf{(P\_5,3.55.1)} guṇasya ca eva prakarṣaḥ na dravyasya . \\
\noindent \textbf{(P\_5,3.55.1)} katham jñāyate . \\
\noindent \textbf{(P\_5,3.55.1)} evam hi dṛśyate loke . \\
\noindent \textbf{(P\_5,3.55.1)} iha samāne āyāme vistāre paṭasya anyaḥ arghaḥ bhavati kāśikasya anyaḥ māthurasya . \\
\noindent \textbf{(P\_5,3.55.1)} guṇāntaram khalu api śilpinaḥ utpādayamānāḥ dravyāntareṇa prakṣālayanti . \\
\noindent \textbf{(P\_5,3.55.1)} anyena śuddham dhautakam kurvanti anyena śaiphālikam anyena mādhyamikam . \\
\noindent \textbf{(P\_5,3.55.1)} yat api ucyate kriyamāṇe ca api guṇagrahaṇe samānaguṇagrahaṇam kartavyam śuklāt kṛṣṇe mā bhūt iti . \\
\noindent \textbf{(P\_5,3.55.1)} na kartavyam . \\
\noindent \textbf{(P\_5,3.55.1)} samānaguṇe eva spardhā bhavati . \\
\noindent \textbf{(P\_5,3.55.1)} nahi āḍhyābhirūpau spardhete . \\
\noindent \textbf{(P\_5,3.55.1)} vācakena khalu api utpattavyam na ca śuklāt kṛṣṇe pratyayaḥ utpadyamānaḥ vācakaḥ syāt . \\
\noindent \textbf{(P\_5,3.55.1)} yat api ucyate adūraviprakarṣe iti vaktavyam iti . \\
\noindent \textbf{(P\_5,3.55.1)} na vaktavyam . \\
\noindent \textbf{(P\_5,3.55.1)} adūraviprakarṣe eva spardhā bhavati . \\
\noindent \textbf{(P\_5,3.55.1)} na hi niṣkadhanaḥ śataniṣkadhanena spardhate . \\
\noindent \textbf{(P\_5,3.55.1)} yat api ucyate jāteḥ na iti vaktavyam iti . \\
\noindent \textbf{(P\_5,3.55.1)} na vaktavyam . \\
\noindent \textbf{(P\_5,3.55.1)} jananena yā prāpyate sā jātiḥ na ca etasya arthasya prakarṣāpakarṣau staḥ . \\
\noindent \textbf{(P\_5,3.55.1)} yat api ucyate tiṅaḥ ca iti atra kriyāgrahaṇam kartavyam sādhanaprakarṣe mā bhūt iti . \\
\noindent \textbf{(P\_5,3.55.1)} na kartavyam . \\
\noindent \textbf{(P\_5,3.55.1)} sādhanam vai dravyam na ca dravyasya prakarṣāpakarṣau staḥ . \\
\noindent \textbf{(P\_5,3.55.1)} kim punaḥ ekam śauklyam āhosvit nānā . \\
\noindent \textbf{(P\_5,3.55.1)} kim ca ataḥ . \\
\noindent \textbf{(P\_5,3.55.1)} yadi ekam prakarṣaḥ na upapadyate . \\
\noindent \textbf{(P\_5,3.55.1)} na hi tena eva tasya prakarṣaḥ bhavati . \\
\noindent \textbf{(P\_5,3.55.1)} atha nānā samānaguṇagrahaṇam kartavyam śuklāt kṛṣṇe mā bhūt iti . \\
\noindent \textbf{(P\_5,3.55.1)} asti ekam śauklyam tat tu viśeṣavat . \\
\noindent \textbf{(P\_5,3.55.1)} kiṅkṛtaḥ viśeṣaḥ . \\
\noindent \textbf{(P\_5,3.55.1)} alpatvamahattvakṛtaḥ . \\
\noindent \textbf{(P\_5,3.55.1)} atha vā punaḥ astu ekam nirviśeṣam ca . \\
\noindent \textbf{(P\_5,3.55.1)} nanu ca uktam prakarṣaḥ na upapadyate . \\
\noindent \textbf{(P\_5,3.55.1)} na hi tena eva tasya prakarṣaḥ bhavati iti . \\
\noindent \textbf{(P\_5,3.55.1)} guṇāntareṇa pracchādāt prakarṣaḥ bhaviṣyati . \\
\noindent \textbf{(P\_5,3.55.1)} atha vā punaḥ astu nānā . \\
\noindent \textbf{(P\_5,3.55.1)} nanu ca uktam samānaguṇagrahaṇam kartavyam śuklāt kṛṣṇe mā bhūt iti . \\
\noindent \textbf{(P\_5,3.55.1)} na kartavyam . \\
\noindent \textbf{(P\_5,3.55.1)} samānaguṇe eva spardhā bhavati . \\
\noindent \textbf{(P\_5,3.55.1)} nahi āḍhyābhirūpau spardhete . \\
\noindent \textbf{(P\_5,3.55.1)} vācakena khalu api utpattavyam na ca śuklāt kṛṣṇe pratyayaḥ utpadyamānaḥ vācakaḥ syāt . \\
\noindent \textbf{(P\_5,3.55.2)} kimantāt punaḥ utpattyā bhavitavyam . \\
\noindent \textbf{(P\_5,3.55.2)} dvitīyāntāt atiśayyamānāt . \\
\noindent \textbf{(P\_5,3.55.2)} śuklam atiśete śuklataraḥ . \\
\noindent \textbf{(P\_5,3.55.2)} kṛṣṇam atiśete kṛṣṇataraḥ . \\
\noindent \textbf{(P\_5,3.55.2)} yadi dvitīyāntāt atiśayyamānāt kālaḥ atiśete kālīm kālitaraḥ iti prāpnoti kālataraḥ iti ca iṣyate . \\
\noindent \textbf{(P\_5,3.55.2)} tathā kālī atiśete kālam kālataraḥ iti prāpnoti kālitarā iti ca iṣyate . \\
\noindent \textbf{(P\_5,3.55.2)} tathā gārgyaḥ atiśete gargān gargataraḥ iti prāpnoti gārgyataraḥ iti ca iṣyate . \\
\noindent \textbf{(P\_5,3.55.2)} tathā gargāḥ atiśerate gārgyam gārgyatarāḥ iti prāpnoti gargatarāḥ iti ca iṣyate . \\
\noindent \textbf{(P\_5,3.55.2)} evam tarhi prathamāntāt svārthikaḥ bhaviṣyati . \\
\noindent \textbf{(P\_5,3.55.2)} kālaḥ atiśete kālataraḥ . \\
\noindent \textbf{(P\_5,3.55.2)} kālī atiśete kālitarā . \\
\noindent \textbf{(P\_5,3.55.2)} gārgyaḥ atiśete gārgyataraḥ . \\
\noindent \textbf{(P\_5,3.55.2)} gargāḥ atiśerate gargatarāḥ . \\
\noindent \textbf{(P\_5,3.55.2)} yadi prathamāntāt svārthikaḥ kumāritarā kiśoritarā avyatiriktam vayaḥ iti kṛtvā vayasi prathame iti ṅīp prāpnoti . \\
\noindent \textbf{(P\_5,3.55.2)} tarapā uktatvāt strīpratyayaḥ na bhaviṣyati . \\
\noindent \textbf{(P\_5,3.55.2)} ṭāp api tarhi na prāpnoti . \\
\noindent \textbf{(P\_5,3.55.2)} ukte api hi bhavanti ete ṭābādayaḥ . \\
\noindent \textbf{(P\_5,3.55.2)} uktam etat svārthikāḥ ṭābādayaḥ iti . \\
\noindent \textbf{(P\_5,3.55.2)} ṅīp api tarhi prāpnoti . \\
\noindent \textbf{(P\_5,3.55.2)} evam tarhi guṇaḥ abhidhīyate . \\
\noindent \textbf{(P\_5,3.55.2)} evam api liṅgavacanāni na sidhyanti . \\
\noindent \textbf{(P\_5,3.55.2)} śuklataram . \\
\noindent \textbf{(P\_5,3.55.2)} śuklatarā . \\
\noindent \textbf{(P\_5,3.55.2)} śuklataraḥ . \\
\noindent \textbf{(P\_5,3.55.2)} śuklatarau . \\
\noindent \textbf{(P\_5,3.55.2)} śuklatarāḥ iti . \\
\noindent \textbf{(P\_5,3.55.2)} āśrayataḥ liṅgavacanāni bhaviṣyanti . \\
\noindent \textbf{(P\_5,3.55.2)} guṇavacanānām hi śabdānām āśrayata liṅgavacanāni bhavanti . \\
\noindent \textbf{(P\_5,3.55.2)} śuklam vastram , śuklā śāṭī śuklaḥ kambalaḥ , śuklau kambalau śuklāḥ kambalāḥ iti . \\
\noindent \textbf{(P\_5,3.55.2)} yat asau dravyam śritaḥ bhavati guṇaḥ tasya yat liṅgam vacanam ca tat gu asya api bhaviṣyati . \\
\noindent \textbf{(P\_5,3.55.2)} atha vā kriyā abhidhīyate . \\
\noindent \textbf{(P\_5,3.55.2)} evam api liṅgavacanāni na sidhyanti . \\
\noindent \textbf{(P\_5,3.55.2)} āśrayataḥ liṅgavacanāni bhaviṣyanti . \\
\noindent \textbf{(P\_5,3.55.2)} evam api dvivacanam prāpnoti . \\
\noindent \textbf{(P\_5,3.55.2)} yaḥ ca atiśete yaḥ ca atiśayyate ubhau tau tasya āśrayau bhavataḥ . \\
\noindent \textbf{(P\_5,3.55.2)} na eṣaḥ doṣaḥ . \\
\noindent \textbf{(P\_5,3.55.2)} katham . \\
\noindent \textbf{(P\_5,3.55.2)} śetiḥ akarmakaḥ . \\
\noindent \textbf{(P\_5,3.55.2)} akarmakāḥ api dhātavaḥ sopasargāḥ sakarmakāḥ bhavanti . \\
\noindent \textbf{(P\_5,3.55.2)} karmāpadiṣṭāḥ vidhayaḥ karmasthabhāvakānām karmasthakriyāṇam vā bhavanti kartṛsthabhāvakaḥ ca śetiḥ . \\
\noindent \textbf{(P\_5,3.55.2)} atha yadi eva dvitīyāntāt utpattiḥ prathamāntāt vā svārthikaḥ atha api guṇaḥ abhidhīyate atha api kriyā kim gatam etat iyatā sūtreṇa āhosvit anyatarasmin pakṣe bhūyaḥ sūtram kartavyam . \\
\noindent \textbf{(P\_5,3.55.2)} gatam iti āha . \\
\noindent \textbf{(P\_5,3.55.2)} katham . \\
\noindent \textbf{(P\_5,3.55.2)} yadā tāvat dvitīyāntāt utpattiḥ prathamāntāt vā svārthikaḥ tadā kṛtyalyuṭaḥ bahulam iti evam atra lyuṭ bhaviṣyati . \\
\noindent \textbf{(P\_5,3.55.2)} yadā guṇaḥ abhidhīyate tadā nyāyasiddham eva . \\
\noindent \textbf{(P\_5,3.55.2)} yadā lapi kriyā tadā api nyāyasiddham eva . \\
\noindent \textbf{(P\_5,3.55.2)} atha vā atiśāyayati iti atiśāyanam . \\
\noindent \textbf{(P\_5,3.55.2)} kaḥ prayojyārthaḥ . \\
\noindent \textbf{(P\_5,3.55.2)} guṇāḥ guṇinam prayojayanti guṇī vā guṇān prayojayati . \\
\noindent \textbf{(P\_5,3.55.2)} kaḥ punaḥ iha śetyarthaḥ . \\
\noindent \textbf{(P\_5,3.55.2)} iha yaḥ yatra bhavati śete asau tatra . \\
\noindent \textbf{(P\_5,3.55.2)} guṇāḥ ca guṇini śerate . \\
\noindent \textbf{(P\_5,3.55.2)} śetyarthaḥ kāritārthaḥ vā nirdeśaḥ ayam samīkṣitaḥ . \\
\noindent \textbf{(P\_5,3.55.2)} śetyarthe na asti vaktavyam . \\
\noindent \textbf{(P\_5,3.55.2)} kāritārthe bravīmi te . \\
\noindent \textbf{(P\_5,3.55.2)} guṇī vā guṇasaṃyogāt guṇaḥ vā guṇinā yadi abhivyajyeta saṃyogāt kāritārthaḥ bhaviṣyati . \\
\noindent \textbf{(P\_5,3.55.3)} iha asya api sūkṣmāṇi vastrāṇi asya api sūkṣmāṇi vastrāṇi iti paratvāt ātiśāyikaḥ prāpnoti . \\
\noindent \textbf{(P\_5,3.55.3)} atiśāyane bahuvrīhau uktam . \\
\noindent \textbf{(P\_5,3.55.3)} kim uktam . \\
\noindent \textbf{(P\_5,3.55.3)} pūrvpadātiśaye ātiśāyikāt bahuvrīhiḥ sūkṣmavastratarādyarthaḥ . \\
\noindent \textbf{(P\_5,3.55.3)} uttarapadātiśaye ātiśāyikaḥ bahuvrīheḥ bahvāḍhyatarādyarthaḥ iti . \\
\noindent \textbf{(P\_5,3.55.3)} iha trīṇi śuklāni vastrāṇi prakarṣāpakarṣayuktāni . \\
\noindent \textbf{(P\_5,3.55.3)} tatra pūrvam apekṣya uttare dve tarabante . \\
\noindent \textbf{(P\_5,3.55.3)} tatra dvayoḥ tarabantayoḥ ekasmāt prakarṣayuktāt śuklataraśabdāt utpattiḥ prāpnoti śuklaśabdāt eva ca iṣyate . \\
\noindent \textbf{(P\_5,3.55.3)} śuklatarasya śuklabhāvāt prakṛteḥ pratyayavijñānam . \\
\noindent \textbf{(P\_5,3.55.3)} śuklataraśabde śuklaśabdaḥ asti . \\
\noindent \textbf{(P\_5,3.55.3)} tasmāt utpattiḥ bhaviṣyati . \\
\noindent \textbf{(P\_5,3.55.3)} na etat vivadāmahe śuklataraśabde śuklaśabdaḥ asti na asti iti . \\
\noindent \textbf{(P\_5,3.55.3)} kim tarhi . \\
\noindent \textbf{(P\_5,3.55.3)} śuklataraśabdaḥ api asti . \\
\noindent \textbf{(P\_5,3.55.3)} tataḥ utpattiḥ prāpnoti . \\
\noindent \textbf{(P\_5,3.55.3)} tadantāt ca svārthe chandasi darśanam śreṣṭhamāya iti . \\
\noindent \textbf{(P\_5,3.55.3)} tadantāt ātiśāyikāntāt ca svārthe chandasi ātiśāyikaḥ dṛśyate . \\
\noindent \textbf{(P\_5,3.55.3)} devo vaḥ savita prarpayatu śreṣṭhamāya karmaṇe . \\
\noindent \textbf{(P\_5,3.55.3)} evam tarhi madhyamāt śuklaśabdāt pūrvaparāpekṣāt utpattiḥ vaktavyā . \\
\noindent \textbf{(P\_5,3.55.3)} madhyamaḥ ca śuklaśabdaḥ pūrvam apekṣya prakṛṣṭaḥ param apekṣya nyūnaḥ na ca nyūnaḥ pravartate . \\
\noindent \textbf{(P\_5,3.55.3)} atha vā utpadyatām . \\
\noindent \textbf{(P\_5,3.55.3)} luk bhaviṣyati . \\
\noindent \textbf{(P\_5,3.55.3)} vācakena khalu api utpattavyam na ca śuklataraśabdāt utpadyamānaḥ vācakaḥ syāt . \\
\noindent \textbf{(P\_5,3.55.3)} na khalu api bahūnām prakarṣe tarapā bhavitavyam . \\
\noindent \textbf{(P\_5,3.55.3)} kena tarhi . \\
\noindent \textbf{(P\_5,3.55.3)} tamapā . \\
\noindent \textbf{(P\_5,3.55.3)} pūrveṇa spardhamānaḥ ayam labhate sitaḥ . \\
\noindent \textbf{(P\_5,3.55.3)} parasmin nyūnatām eti na ca nyūnaḥ pravartate . \\
\noindent \textbf{(P\_5,3.55.3)} apekṣya madhyamaḥ pūrvam ādhikyam labhate sitaḥ . \\
\noindent \textbf{(P\_5,3.55.3)} parasmin nyūnatām eti yathā amātyaḥ sthite nṛpe . \\
\noindent \textbf{(P\_5,3.55.3)} astu vā api taraḥ tasmāt . \\
\noindent \textbf{(P\_5,3.55.3)} na apaśabdaḥ bhaviṣyati . \\
\noindent \textbf{(P\_5,3.55.3)} vācakaḥ cet prayoktavyaḥ vācakaḥ cet prayujyatām . \\
\noindent \textbf{(P\_5,3.57)} dvivacane iti ucyate . \\
\noindent \textbf{(P\_5,3.57)} tatra idam na sidhyati . \\
\noindent \textbf{(P\_5,3.57)} dantoṣṭhasya dantāḥ snigdhatarāḥ . \\
\noindent \textbf{(P\_5,3.57)} pāṇipādasya pādau sukumāratarau . \\
\noindent \textbf{(P\_5,3.57)} asmākam ca devadattasya ca devadattaḥ abhirūpataraḥ iti . \\
\noindent \textbf{(P\_5,3.57)} yadi punaḥ dvyarthopapade iti ucyeta . \\
\noindent \textbf{(P\_5,3.57)} tatra ayam api arthaḥ . \\
\noindent \textbf{(P\_5,3.57)} vibhajyopapadagrahaṇam na kartavyam . \\
\noindent \textbf{(P\_5,3.57)} iha api śaṅkāśyakebhyaḥ pāṭaliputrakāḥ abhirūpatarāḥ iti dvyarthopapade iti eva siddham . \\
\noindent \textbf{(P\_5,3.57)} na evañjātīyakā dvyarthatā śakyā vijñātum . \\
\noindent \textbf{(P\_5,3.57)} iha api prasjyeta . \\
\noindent \textbf{(P\_5,3.57)} śaṅkāśyakānām pāṭaliputrakāṇām ca pāṭaliputrakāḥ abhirūpatamāḥ iti . \\
\noindent \textbf{(P\_5,3.57)} avaśyam khalu api vibhajyopapadagrahaṇam kartavyam yaḥ hi bahūnām vibhāgaḥ tadartham . \\
\noindent \textbf{(P\_5,3.57)} śaṅkāśyakebhyaḥ ca pāṭaliputrakebhyaḥ ca māthurāḥ abhirūpatarāḥ iti . \\
\noindent \textbf{(P\_5,3.57)} tat tarhi dvyarthopapade iti vaktavyam . \\
\noindent \textbf{(P\_5,3.57)} na vaktavyam . \\
\noindent \textbf{(P\_5,3.57)} na idam pāribhāṣikasya dvivacanasya grahaṇam . \\
\noindent \textbf{(P\_5,3.57)} kim tarhi anvarthagrahaṇam . \\
\noindent \textbf{(P\_5,3.57)} ucyate vacanam . \\
\noindent \textbf{(P\_5,3.57)} dvayoḥ arthayoḥ vacanam dvivacanam iti . \\
\noindent \textbf{(P\_5,3.57)} evam api tarabīyasunoḥ ekadravyasya utkarṣāpakarṣayoḥ upasaṅkhyānam . \\
\noindent \textbf{(P\_5,3.57)} tarabīyasunoḥ ekadravyasya utkarṣāpakarṣayoḥ upasaṅkhyānam vaktavyam . \\
\noindent \textbf{(P\_5,3.57)} parut bhavān paṭuḥ āsīt . \\
\noindent \textbf{(P\_5,3.57)} paṭutaraḥ ca aiṣamaḥ iti . \\
\noindent \textbf{(P\_5,3.57)} siddham tu guṇapradhānatvāt . \\
\noindent \textbf{(P\_5,3.57)} siddham etat . \\
\noindent \textbf{(P\_5,3.57)} katham . \\
\noindent \textbf{(P\_5,3.57)} guṇapradhānatvāt . \\
\noindent \textbf{(P\_5,3.57)} guṇapradhānaḥ ayam nirdeśaḥ kriyate . \\
\noindent \textbf{(P\_5,3.57)} guṇāntarayogāt ca anyatvam bhavati . \\
\noindent \textbf{(P\_5,3.57)} tat yathā tam eva guṇāntarayuktam vaktāraḥ bhavanti anyaḥ bhavān saṃvṛttaḥ iti . \\
\noindent \textbf{(P\_5,3.58)} evakāraḥ kimarthaḥ . \\
\noindent \textbf{(P\_5,3.58)} niyamārthaḥ . \\
\noindent \textbf{(P\_5,3.58)} na etat asti prayojanam . \\
\noindent \textbf{(P\_5,3.58)} siddhe vidhiḥ ārabhyamāṇaḥ antareṇa evakāram niyamārthaḥ bhaviṣyati . \\
\noindent \textbf{(P\_5,3.58)} iṣṭataḥ avadhāraṇārthaḥ tarhi . \\
\noindent \textbf{(P\_5,3.58)} yathā evam vijñāyeta . \\
\noindent \textbf{(P\_5,3.58)} ajādī guṇavacanāt eva iti . \\
\noindent \textbf{(P\_5,3.58)} mā evam vijñāyi . \\
\noindent \textbf{(P\_5,3.58)} ajādī eva guṇavacanāt iti . \\
\noindent \textbf{(P\_5,3.58)} kim ca syāt na vyañjanādī guṇavacanāt syātām . \\
\noindent \textbf{(P\_5,3.60)} idam ayuktam vartate . \\
\noindent \textbf{(P\_5,3.60)} kim atra ayuktam . \\
\noindent \textbf{(P\_5,3.60)} ajādī guṇavacanāt eva iti uktvā aguṇavacanānām api ajādyoḥ ādeśāḥ ucyante . \\
\noindent \textbf{(P\_5,3.60)} na eṣaḥ doṣaḥ . \\
\noindent \textbf{(P\_5,3.60)} etat eva jñāpayati bhavataḥ etebhyaḥ aguṇavacanebhyaḥ api ajādī iti yat ayam ajādyoḥ parataḥ ādeśān śāsti . \\
\noindent \textbf{(P\_5,3.60)} evam api tayoḥ iti vaktavyam syāt . \\
\noindent \textbf{(P\_5,3.60)} tayoḥ parataḥ iti . \\
\noindent \textbf{(P\_5,3.60)} yadi punaḥ ayam vidhiḥ vijñāyeta . \\
\noindent \textbf{(P\_5,3.60)} na evam śakyam . \\
\noindent \textbf{(P\_5,3.60)} vyañjanādī hi na syātām upādhīnām ca saṅkaraḥ syāt punarvidhānāt ajādyoḥ . \\
\noindent \textbf{(P\_5,3.60)} nanu ca ete viśeṣāḥ anuvarteran . \\
\noindent \textbf{(P\_5,3.60)} yadi api ete anuvarteran vyañjanādī tarhi na syātām . \\
\noindent \textbf{(P\_5,3.60)} evam tarhi ācāryapravṛttiḥ jñāpayati bhavataḥ etebhyaḥ aguṇavacanebhyaḥ api ajādī iti yat ayam ajādyoḥ parataḥ ādeśān śāsti . \\
\noindent \textbf{(P\_5,3.60)} nanu ca uktam tayoḥ iti vaktavyam iti . \\
\noindent \textbf{(P\_5,3.60)} na vaktavyam . \\
\noindent \textbf{(P\_5,3.60)} prakṛtam ajādīgrahaṇam anuvartate . \\
\noindent \textbf{(P\_5,3.60)} kva prakṛtam . \\
\noindent \textbf{(P\_5,3.60)} ajādī guṇavacanāt eva iti . \\
\noindent \textbf{(P\_5,3.60)} tat vai prathamānirdiṣṭam saptamīnirdiṣṭena ca iha arthaḥ . \\
\noindent \textbf{(P\_5,3.60)} arthāt vibhaktivipariṇāmaḥ bhaviṣyati . \\
\noindent \textbf{(P\_5,3.60)} tat yathā . \\
\noindent \textbf{(P\_5,3.60)} uccāni devadattasya gṛhāṇi . \\
\noindent \textbf{(P\_5,3.60)} āmantrayasva enam . \\
\noindent \textbf{(P\_5,3.60)} devadattam iti gamyate . \\
\noindent \textbf{(P\_5,3.60)} devadattasya gāva aśvā hiraṇyam iti . \\
\noindent \textbf{(P\_5,3.60)} āḍhyaḥ vaidhaveyaḥ . \\
\noindent \textbf{(P\_5,3.60)} devadattaḥ iti gamyate . \\
\noindent \textbf{(P\_5,3.60)} purastāt ṣaṣṭhīnirdiṣṭam sat arthāt dvitīyānirdiṣṭam prathamānirdiṣṭam ca bhavati . \\
\noindent \textbf{(P\_5,3.60)} evam iha api purastāt prathamānirdiṣṭam sat arthāt saptamīnirdiṣṭam bhaviṣyati . \\
\noindent \textbf{(P\_5,3.66.1)} strīliṅgena nirdeśaḥ kriyate ekavacanāntena ca . \\
\noindent \textbf{(P\_5,3.66.1)} tena strīliṅgāt eva utpattiḥ syāt ekavacanāntāt ca . \\
\noindent \textbf{(P\_5,3.66.1)} punnapuṃsakaliṅgāt dvivacanabahuvacanāntāt ca na syāt . \\
\noindent \textbf{(P\_5,3.66.1)} na eṣaḥ doṣaḥ . \\
\noindent \textbf{(P\_5,3.66.1)} na ayam pratyayārthaḥ . \\
\noindent \textbf{(P\_5,3.66.1)} kim tarhi . \\
\noindent \textbf{(P\_5,3.66.1)} prakṛtyarthaviśeṣaṇam etat . \\
\noindent \textbf{(P\_5,3.66.1)} praśaṃsāyām yat prātipadikam vartate tasmāt rūpap bhavati . \\
\noindent \textbf{(P\_5,3.66.1)} kasmin arthe . \\
\noindent \textbf{(P\_5,3.66.1)} svārthe iti . \\
\noindent \textbf{(P\_5,3.66.1)} svārthikāḥ ca prakṛtitaḥ liṅgavacanāni anuvartante . \\
\noindent \textbf{(P\_5,3.66.1)} prakṛteḥ liṅgavacanābhāvāt tiṅprakṛteḥ ambhāvavacanam . \\
\noindent \textbf{(P\_5,3.66.1)} prakṛteḥ liṅgavacanābhāvāt tiṅprakṛteḥ rūpapaḥ ambhāvaḥ vaktavyaḥ . \\
\noindent \textbf{(P\_5,3.66.1)} pacatirūpam . \\
\noindent \textbf{(P\_5,3.66.1)} pacatorūpam . \\
\noindent \textbf{(P\_5,3.66.1)} pacantirūpam iti . \\
\noindent \textbf{(P\_5,3.66.1)} siddham tu kriyāpradhānatvāt . \\
\noindent \textbf{(P\_5,3.66.1)} siddham etat . \\
\noindent \textbf{(P\_5,3.66.1)} katham . \\
\noindent \textbf{(P\_5,3.66.1)} kriyāpradhānatvāt . \\
\noindent \textbf{(P\_5,3.66.1)} kriyāpradhānam ākhyātam ekā ca kriyā . \\
\noindent \textbf{(P\_5,3.66.1)} dravyapradhānam nāma . \\
\noindent \textbf{(P\_5,3.66.1)} katham punaḥ jñayate kriyāpradhānam ākhyātam bhavati dravyapradhānam nāma iti . \\
\noindent \textbf{(P\_5,3.66.1)} yat kriyām pṛṣṭaḥ tiṅā ācaṣṭe . \\
\noindent \textbf{(P\_5,3.66.1)} kim devadattaḥ karoti . \\
\noindent \textbf{(P\_5,3.66.1)} pacati iti . \\
\noindent \textbf{(P\_5,3.66.1)} dravyam pṛṣṭaḥ kṛtā ācaṣṭe . \\
\noindent \textbf{(P\_5,3.66.1)} kataraḥ devadattaḥ . \\
\noindent \textbf{(P\_5,3.66.1)} yaḥ kārakaḥ hārakaḥ iti . \\
\noindent \textbf{(P\_5,3.66.1)} yadi tarhi ekā kriyā dvivacanabahuvacanāni na sidyanti . \\
\noindent \textbf{(P\_5,3.66.1)} pacataḥ . \\
\noindent \textbf{(P\_5,3.66.1)} pacanti iti . \\
\noindent \textbf{(P\_5,3.66.1)} na etāni kriyāpekṣāṇi . \\
\noindent \textbf{(P\_5,3.66.1)} kim tarhi sādhanāpekṣāṇi . \\
\noindent \textbf{(P\_5,3.66.1)} iha api tarhi prāpnuvanti . \\
\noindent \textbf{(P\_5,3.66.1)} pacatirūpam . \\
\noindent \textbf{(P\_5,3.66.1)} pacatorūpam . \\
\noindent \textbf{(P\_5,3.66.1)} pacantirūpam iti . \\
\noindent \textbf{(P\_5,3.66.1)} tiṅā uktatvāt tasya abhisambandhasya na bhaviṣyati . \\
\noindent \textbf{(P\_5,3.66.1)} ekavacanam api tarhi na prāpnoti . \\
\noindent \textbf{(P\_5,3.66.1)} samayāt bhaviṣyati . \\
\noindent \textbf{(P\_5,3.66.1)} dvivacanabahuvacanāni api tarhi samayāt prāpnuvanti . \\
\noindent \textbf{(P\_5,3.66.1)} evam tarhi ekavacanam utsargaḥ kariṣyate . \\
\noindent \textbf{(P\_5,3.66.1)} tasya dvibahvoḥ dvivacanabahuvacane apavāvau bhaviṣyataḥ . \\
\noindent \textbf{(P\_5,3.66.1)} evam api napuṃsakatvam vaktavyam . \\
\noindent \textbf{(P\_5,3.66.1)} na vaktavyam . \\
\noindent \textbf{(P\_5,3.66.1)} liṅgam aśiṣyam lokāśrayatvāt liṅgasya . \\
\noindent \textbf{(P\_5,3.66.2)} vṛṣalādibhyaḥ upasaṅkhyānam . \\
\noindent \textbf{(P\_5,3.66.2)} vṛṣalādibhyaḥ upasaṅkhyānam kartavyam . \\
\noindent \textbf{(P\_5,3.66.2)} vṛṣalarūpaḥ . \\
\noindent \textbf{(P\_5,3.66.2)} dasyurūpaḥ . \\
\noindent \textbf{(P\_5,3.66.2)} corarūpaḥ iti . \\
\noindent \textbf{(P\_5,3.66.2)} siddham tu prakṛtyarthavaiśiṣṭyavacanāt . \\
\noindent \textbf{(P\_5,3.66.2)} siddham etat . \\
\noindent \textbf{(P\_5,3.66.2)} katham . \\
\noindent \textbf{(P\_5,3.66.2)} prakṛtyarthasya vaiśiṣṭye iti vaktavyam . \\
\noindent \textbf{(P\_5,3.66.2)} vṛṣalarūpaḥ ayam . \\
\noindent \textbf{(P\_5,3.66.2)} api ayam palāṇḍunā surām pibet . \\
\noindent \textbf{(P\_5,3.66.2)} corarūpaḥ ayam . \\
\noindent \textbf{(P\_5,3.66.2)} api ayam akṣṇoḥ añjanam haret . \\
\noindent \textbf{(P\_5,3.66.2)} dasyurūpaḥ ayam . \\
\noindent \textbf{(P\_5,3.66.2)} api ayam dhāvataḥ lohitam pibet . \\
\noindent \textbf{(P\_5,3.67)} īṣadasamāptaukriyāpradhānatvāt liṅgavacanānupapattiḥ . \\
\noindent \textbf{(P\_5,3.67)} īṣadasamāptau kriyāpradhānatvāt liṅgavacanayoḥ anupapattiḥ . \\
\noindent \textbf{(P\_5,3.67)} paṭukalpaḥ . \\
\noindent \textbf{(P\_5,3.67)} paṭukalpau . \\
\noindent \textbf{(P\_5,3.67)} paṭukalpāḥ iti . \\
\noindent \textbf{(P\_5,3.67)} ekaḥ ayam arthaḥ īṣadasamāptiḥ nāma . \\
\noindent \textbf{(P\_5,3.67)} tasya ekatvāt ekavacanam prāpnoti . \\
\noindent \textbf{(P\_5,3.67)} prakṛtyarthaviśeṣaṇatvād siddham . \\
\noindent \textbf{(P\_5,3.67)} siddham etat . \\
\noindent \textbf{(P\_5,3.67)} katham . \\
\noindent \textbf{(P\_5,3.67)} na ayam pratyayārthaḥ . \\
\noindent \textbf{(P\_5,3.67)} kim tarhi prakṛtyarthaviśeṣaṇam etat . \\
\noindent \textbf{(P\_5,3.67)} īṣadasamāptau yat prātipadikam vartate tasmāt kalpabādayaḥ bhavanti . \\
\noindent \textbf{(P\_5,3.67)} kasmin arthe . \\
\noindent \textbf{(P\_5,3.67)} svārthe iti . \\
\noindent \textbf{(P\_5,3.67)} svārthikāḥ ca prakṛtitaḥ liṅgavacanāni anuvartante . \\
\noindent \textbf{(P\_5,3.67)} prakṛtyarthe cet liṅgavacanānupapattiḥ . \\
\noindent \textbf{(P\_5,3.67)} prakṛtyarthe cet liṅgavacanayoḥ anupapattiḥ . \\
\noindent \textbf{(P\_5,3.67)} guḍakalpā drākṣā . \\
\noindent \textbf{(P\_5,3.67)} tailakalpā prasannā . \\
\noindent \textbf{(P\_5,3.67)} payaskalpā yavāgūḥ iti . \\
\noindent \textbf{(P\_5,3.67)} siddham tu tatsambandhe uttarapadārthe pratyayavacanāt . \\
\noindent \textbf{(P\_5,3.67)} siddham etat . \\
\noindent \textbf{(P\_5,3.67)} katham . \\
\noindent \textbf{(P\_5,3.67)} tatsambandhe īṣadasamāptisambandhe uttarapadārthe pratyayaḥ bhavati iti vaktavyam . \\
\noindent \textbf{(P\_5,3.67)} sidhyati . \\
\noindent \textbf{(P\_5,3.67)} sūtram tarhi bhidyate . \\
\noindent \textbf{(P\_5,3.67)} yathānyāsam eva astu . \\
\noindent \textbf{(P\_5,3.67)} nanu ca uktam īṣadasamāptaukriyāpradhānatvāt liṅgavacanānupapattiḥ iti . \\
\noindent \textbf{(P\_5,3.67)} parihṛtam etat prakṛtyarthaviśeṣaṇatvād siddham iti . \\
\noindent \textbf{(P\_5,3.67)} nanu ca uktam prakṛtyarthe cet liṅgavacanānupapattiḥ iti . \\
\noindent \textbf{(P\_5,3.67)} na eṣaḥ doṣaḥ . \\
\noindent \textbf{(P\_5,3.67)} ācāryapravṛttiḥ jñāpayati svārthikāḥ ativartante api liṅgavacanāni iti yat ayam ṇacaḥ striyām añ iti strīgrahaṇam karoti . \\
\noindent \textbf{(P\_5,3.67)} yadi etat jñāpyate bahuguḍaḥ drākṣā . \\
\noindent \textbf{(P\_5,3.67)} bahutailam prasannā . \\
\noindent \textbf{(P\_5,3.67)} bahupayaḥ yavāgūḥ iti atra api prāpnoti . \\
\noindent \textbf{(P\_5,3.67)} na api ativartante . \\
\noindent \textbf{(P\_5,3.67)} kim punaḥ iha udāharaṇam . \\
\noindent \textbf{(P\_5,3.67)} paṭukalpaḥ . \\
\noindent \textbf{(P\_5,3.67)} mṛdukalpaḥ iti . \\
\noindent \textbf{(P\_5,3.67)} na etat asti . \\
\noindent \textbf{(P\_5,3.67)} nirjñātasya arthasya samāptiḥ vā bhavati visamāptiḥ vā guṇaḥ ca anirjñātaḥ . \\
\noindent \textbf{(P\_5,3.67)} idam tarhi . \\
\noindent \textbf{(P\_5,3.67)} guḍakalpā drākṣā . \\
\noindent \textbf{(P\_5,3.67)} tailakalpā prasannā . \\
\noindent \textbf{(P\_5,3.67)} payaskalpā yavāgūḥ iti . \\
\noindent \textbf{(P\_5,3.67)} dravyam api anirjñātam . \\
\noindent \textbf{(P\_5,3.67)} idam tarhi . \\
\noindent \textbf{(P\_5,3.67)} kṛtakalpam . \\
\noindent \textbf{(P\_5,3.67)} bhuktakalpam . \\
\noindent \textbf{(P\_5,3.67)} pītakalpam iti . \\
\noindent \textbf{(P\_5,3.67)} ktāntāt pratyayavidhānānupapattiḥ ktasya bhūtakālalakṣaṇatvāt kalpādīnām ca asamāptivacanāt . \\
\noindent \textbf{(P\_5,3.67)} ktāntāt pratyayavidhāneḥ anupapattiḥ . \\
\noindent \textbf{(P\_5,3.67)} kim kāraṇam . \\
\noindent \textbf{(P\_5,3.67)} ktasya bhūtakālalakṣaṇatvāt . \\
\noindent \textbf{(P\_5,3.67)} bhūtakālalakṣaṇaḥ ktaḥ . \\
\noindent \textbf{(P\_5,3.67)} kalpādīnām ca asamāptivacanāt . \\
\noindent \textbf{(P\_5,3.67)} visamāptivacanāḥ ca kalpādayaḥ . \\
\noindent \textbf{(P\_5,3.67)} na ca asti sambhavaḥ yat bhūtakālaḥ ca syāt asamāptiḥ ca iti . \\
\noindent \textbf{(P\_5,3.67)} siddham tu āśaṃsāyām bhūtavadvacanāt . \\
\noindent \textbf{(P\_5,3.67)} siddham etat . \\
\noindent \textbf{(P\_5,3.67)} katham . \\
\noindent \textbf{(P\_5,3.67)} āśaṃsāyām bhūtavat ca iti evam atra ktaḥ bhaviṣyati . \\
\noindent \textbf{(P\_5,3.67)} idam ca api udāharaṇam paṭukalpaḥ . \\
\noindent \textbf{(P\_5,3.67)} mṛdukalpaḥ iti . \\
\noindent \textbf{(P\_5,3.67)} nanu ca uktam nirjñātasya arthasya samāptiḥ vā bhavati visamāptiḥ vā guṇaḥ ca anirjñātaḥ iti . \\
\noindent \textbf{(P\_5,3.67)} lokataḥ vyavahāram dṛṣṭvā guṇasya nirjñānam . \\
\noindent \textbf{(P\_5,3.67)} tat yathā paṭuḥ ayam brāhmaṇaḥ iti ucyate yaḥ laghunā upāyena athān sādhayati . \\
\noindent \textbf{(P\_5,3.67)} paṭukalpaḥ ayam iti ucyati yaḥ na tathā sādhayati . \\
\noindent \textbf{(P\_5,3.67)} idam ca api udāharaṇam guḍakalpā drākṣā . \\
\noindent \textbf{(P\_5,3.67)} tailakalpā prasannā . \\
\noindent \textbf{(P\_5,3.67)} payaskalpā yavāgūḥ iti . \\
\noindent \textbf{(P\_5,3.67)} nanu ca uktam dravyam api anirjñātam iti . \\
\noindent \textbf{(P\_5,3.67)} lokataḥ dravyam api nirjñātam . \\
\noindent \textbf{(P\_5,3.68.1)} vibhāṣāgrahaṇam kimartham . \\
\noindent \textbf{(P\_5,3.68.1)} vibhāṣā bahuc yathā syāt . \\
\noindent \textbf{(P\_5,3.68.1)} bahucā mukte vākyam api yathā syāt . \\
\noindent \textbf{(P\_5,3.68.1)} na etat asti prayojanam . \\
\noindent \textbf{(P\_5,3.68.1)} prakṛtā mahāvibhāṣā . \\
\noindent \textbf{(P\_5,3.68.1)} tayā vākyam bhaviṣyati . \\
\noindent \textbf{(P\_5,3.68.1)} idam tarhi prayojanam . \\
\noindent \textbf{(P\_5,3.68.1)} kalpādayaḥ api yathā syuḥ iti . \\
\noindent \textbf{(P\_5,3.68.1)} etat api na asti prayojanam . \\
\noindent \textbf{(P\_5,3.68.1)} bahuc ucyate kalpādayaḥ api . \\
\noindent \textbf{(P\_5,3.68.1)} tat ubhayam vacanāt bhaviṣyati . \\
\noindent \textbf{(P\_5,3.68.1)} na evam śakyam . \\
\noindent \textbf{(P\_5,3.68.1)} akriyamāṇe hi vibhāṣāgrahaṇe anavakāśaḥ bahuc kalpādīn bādheta . \\
\noindent \textbf{(P\_5,3.68.1)} kalpādayaḥ api anavakāśāḥ . \\
\noindent \textbf{(P\_5,3.68.1)} te vacanāt bhaviṣyanti . \\
\noindent \textbf{(P\_5,3.68.1)} sāvakāśāḥ kalpādayaḥ . \\
\noindent \textbf{(P\_5,3.68.1)} kaḥ avakāśaḥ . \\
\noindent \textbf{(P\_5,3.68.1)} tiṅantāni avakāśaḥ . \\
\noindent \textbf{(P\_5,3.68.1)} atha subgrahaṇam kimartham . \\
\noindent \textbf{(P\_5,3.68.1)} subantāt utpattiḥ yathā syāt . \\
\noindent \textbf{(P\_5,3.68.1)} prātipadikāt mā bhūt iti . \\
\noindent \textbf{(P\_5,3.68.1)} na etat asti prayojanam . \\
\noindent \textbf{(P\_5,3.68.1)} na asti atra viśeṣaḥ subantāt utpattau satyām prātipadikāt vā . \\
\noindent \textbf{(P\_5,3.68.1)} yadi evam iha api na arthaḥ subgrahaṇena . \\
\noindent \textbf{(P\_5,3.68.1)} supaḥ ātmanaḥ kyac iti . \\
\noindent \textbf{(P\_5,3.68.1)} iha api na asti atra viśeṣaḥ subantāt utpattau satyām prātipadikāt vā . \\
\noindent \textbf{(P\_5,3.68.1)} ayam asti viśeṣaḥ . \\
\noindent \textbf{(P\_5,3.68.1)} subantāt utpattau satyām padasañjñā siddhā bhavati . \\
\noindent \textbf{(P\_5,3.68.1)} prātipadikāt utpattau satyām padasañjñā na prāpnoti . \\
\noindent \textbf{(P\_5,3.68.1)} nanu ca prātipadikāt api utpattau satyām padasañjñā siddhā . \\
\noindent \textbf{(P\_5,3.68.1)} katham . \\
\noindent \textbf{(P\_5,3.68.1)} ārabhyate naḥ kye iti . \\
\noindent \textbf{(P\_5,3.68.1)} tat ca avaśyam kartavyam subantāt utpattau niyamārtham . \\
\noindent \textbf{(P\_5,3.68.1)} tat eva prātipadikāt utpattau satyām vidhyartham bhaviṣyati . \\
\noindent \textbf{(P\_5,3.68.1)} idam tarhi prayojanam . \\
\noindent \textbf{(P\_5,3.68.1)} subantāt utpattiḥ yathā syāt . \\
\noindent \textbf{(P\_5,3.68.1)} tiṅantāt mā bhūt iti . \\
\noindent \textbf{(P\_5,3.68.1)} etat api na asti prayojanam . \\
\noindent \textbf{(P\_5,3.68.1)} ṅyāpprātipadikāt iti vartate . \\
\noindent \textbf{(P\_5,3.68.1)} ataḥ uttaram paṭhati . \\
\noindent \textbf{(P\_5,3.68.1)} bahuci subgrahaṇāt pūrvatra tiṅaḥ vidhānam . \\
\noindent \textbf{(P\_5,3.68.1)} bahuci subgrahaṇam kriyate pūrvatra tiṅaḥ vidhiḥ yathā vijñāyeta . \\
\noindent \textbf{(P\_5,3.68.1)} na etat asti prayojanam . \\
\noindent \textbf{(P\_5,3.68.1)} prakṛtam tiṅgrahaṇam anuvartate . \\
\noindent \textbf{(P\_5,3.68.1)} kva prakṛtam . \\
\noindent \textbf{(P\_5,3.68.1)} atiśāyane tamabiṣṭhanau . \\
\noindent \textbf{(P\_5,3.68.1)} tiṅaḥ ca iti . \\
\noindent \textbf{(P\_5,3.68.1)} evam tarhi bahuci subgrahaṇam pūrvatra tiṅaḥ vidhānāt . \\
\noindent \textbf{(P\_5,3.68.1)} bahuci subgrahaṇam kriyate . \\
\noindent \textbf{(P\_5,3.68.1)} kim kāraṇam . \\
\noindent \textbf{(P\_5,3.68.1)} pūrvatra tiṅaḥ vidhānāt . \\
\noindent \textbf{(P\_5,3.68.1)} pūrvatra tiṅaḥ ca iti anuvartate . \\
\noindent \textbf{(P\_5,3.68.1)} tat iha api prāpnoti . \\
\noindent \textbf{(P\_5,3.68.1)} nanu ca tiṅgrahaṇam nivarteta . \\
\noindent \textbf{(P\_5,3.68.1)} avaśyam uttarārtham anuvartyam avyayasarvanāmnām akac prāk ṭeḥ iti pacataki jalpataki iti evamartham . \\
\noindent \textbf{(P\_5,3.68.1)} yadi subgrahaṇam kriyate svaraḥ na sidhyati . \\
\noindent \textbf{(P\_5,3.68.1)} bahupaṭavaḥ evam svaraḥ prasajyeta bahupaṭavaḥ iti ca iṣyate . \\
\noindent \textbf{(P\_5,3.68.1)} paṭhiṣyati hi ācāryaḥ citaḥ saprakṛteḥ bahvakajartham iti . \\
\noindent \textbf{(P\_5,3.68.1)} svaraḥ katham . \\
\noindent \textbf{(P\_5,3.68.1)} svaraḥ prātipadikatvāt . \\
\noindent \textbf{(P\_5,3.68.1)} subluki kṛte prātipadikatvāt svaraḥ bhaviṣyati . \\
\noindent \textbf{(P\_5,3.68.1)} atha tugrahaṇam kimartham . \\
\noindent \textbf{(P\_5,3.68.1)} tugrahaṇam nityapūrvārtham . \\
\noindent \textbf{(P\_5,3.68.1)} tugrahaṇam kriyate nityam pūrvaḥ yathā syāt . \\
\noindent \textbf{(P\_5,3.68.1)} vibhāṣā mā bhūt iti . \\
\noindent \textbf{(P\_5,3.68.1)} na etat asti prayojanam . \\
\noindent \textbf{(P\_5,3.68.1)} na vibhāṣāgrahaṇena pūrvam abhisambadhyate . \\
\noindent \textbf{(P\_5,3.68.1)} kim tarhi . \\
\noindent \textbf{(P\_5,3.68.1)} bahuc abhisambadhyate : vibhāṣā bahuc bhavati iti . \\
\noindent \textbf{(P\_5,3.68.1)} yadā ca bhavati tadā pūrvaḥ bhavati . \\
\noindent \textbf{(P\_5,3.68.1)} idam tarhi prayojanam . \\
\noindent \textbf{(P\_5,3.68.1)} prāk utpatteḥ yat liṅgam vacanam ca tat utpanne api pratyaye yathā syāt . \\
\noindent \textbf{(P\_5,3.68.1)} bahuguḍaḥ drākṣā . \\
\noindent \textbf{(P\_5,3.68.1)} bahutailam prasannā . \\
\noindent \textbf{(P\_5,3.68.1)} bahupayaḥ yavāgūḥ iti . \\
\noindent \textbf{(P\_5,3.68.1)} etat api na asti prayojanam . \\
\noindent \textbf{(P\_5,3.68.1)} svāṛthikaḥ ayam svārthikāḥ ca prakṛtitaḥ liṅgavacanāni anuvartante . \\
\noindent \textbf{(P\_5,3.68.1)} evam tarhi siddhe sati yat tugrahaṇam karoti tat jñāpayati ācāryaḥ svārthikāḥ ativartante api liṅgavacanāni iti . \\
\noindent \textbf{(P\_5,3.68.1)} kim etasya jñāpane prayojanam . \\
\noindent \textbf{(P\_5,3.68.1)} guḍakalpā drākṣā , tailakalpā prasannā , payaskalpā yavāgūḥ iti etat siddham bhavati . \\
\noindent \textbf{(P\_5,3.68.2)} tamādibhyaḥ kalpādayaḥ vipratiṣedhena . \\
\noindent \textbf{(P\_5,3.68.2)} tamādibhyaḥ kalpādayaḥ bhavanti vipratiṣedhena . \\
\noindent \textbf{(P\_5,3.68.2)} tamādīnām avakāśaḥ prakarṣasya vacanam īṣadasamāpteḥ avacanam . \\
\noindent \textbf{(P\_5,3.68.2)} paṭutaraḥ . \\
\noindent \textbf{(P\_5,3.68.2)} paṭutamaḥ . \\
\noindent \textbf{(P\_5,3.68.2)} kalpādīnām īṣadasamāpteḥ vacanam prakarṣasya avacanam . \\
\noindent \textbf{(P\_5,3.68.2)} paṭukalpaḥ . \\
\noindent \textbf{(P\_5,3.68.2)} mṛdukalpaḥ . \\
\noindent \textbf{(P\_5,3.68.2)} ubhayavacane ubhayam prāpnoti . \\
\noindent \textbf{(P\_5,3.68.2)} paṭukalpataraḥ . \\
\noindent \textbf{(P\_5,3.68.2)} mṛdukalpataraḥ . \\
\noindent \textbf{(P\_5,3.68.2)} kalpādayaḥ bhavanti vipratiṣedhena . \\
\noindent \textbf{(P\_5,3.68.2)} yadi evam īṣadasamāpteḥ prakarṣe tamādiḥ pratyayaḥ prāpnoti prakṛteḥ eva ca iṣyate . \\
\noindent \textbf{(P\_5,3.68.2)} tamādiḥ īṣatpradhānāt . \\
\noindent \textbf{(P\_5,3.68.2)} tamādiḥ īṣatpradhānāt api bhavati . \\
\noindent \textbf{(P\_5,3.68.2)} asya prakarṣaḥ asti . \\
\noindent \textbf{(P\_5,3.68.2)} tasya prakarṣe bhaviṣyati . \\
\noindent \textbf{(P\_5,3.68.2)} kasya ca prakarṣaḥ asti. prakṛteḥ eva . \\
\noindent \textbf{(P\_5,3.71-72.1)} KA\_II,422.16-21 Ro\_IV,230 {1/10}     kim ayam subantasya prāk ṭeḥ bhavati āhosvit ṅyāpprātipadikasya . \\
\noindent \textbf{(P\_5,3.71-72.1)} KA\_II,422.16-21 Ro\_IV,230 {2/10}     kutaḥ sandehaḥ . \\
\noindent \textbf{(P\_5,3.71-72.1)} KA\_II,422.16-21 Ro\_IV,230 {3/10}     ubhayam prakṛtam . \\
\noindent \textbf{(P\_5,3.71-72.1)} KA\_II,422.16-21 Ro\_IV,230 {4/10}     anyatarat śakyam viśeṣayitum . \\
\noindent \textbf{(P\_5,3.71-72.1)} KA\_II,422.16-21 Ro\_IV,230 {5/10}     kim ca ataḥ . \\
\noindent \textbf{(P\_5,3.71-72.1)} KA\_II,422.16-21 Ro\_IV,230 {6/10}     yadi subantasya yuṣmakābhiḥ asmakābhiḥ yuṣmakāsu asmakāsu yuvakayoḥ āvakayoḥ iti na sidhyati . \\
\noindent \textbf{(P\_5,3.71-72.1)} KA\_II,422.16-21 Ro\_IV,230 {7/10}     atha prātipadikasya tvayakā mayakā tvayaki mayaki iti atra api prāpnoti . \\
\noindent \textbf{(P\_5,3.71-72.1)} KA\_II,422.16-21 Ro\_IV,230 {8/10}     astu subantasya . \\
\noindent \textbf{(P\_5,3.71-72.1)} KA\_II,422.16-21 Ro\_IV,230 {9/10}     katham yuṣmakābhiḥ asmakābhiḥ yuṣmakāsu asmakāsu yuvakayoḥ āvakayoḥ iti . \\
\noindent \textbf{(P\_5,3.71-72.1)} KA\_II,422.16-21 Ro\_IV,230 {10/10}     anokārasakārabhakārādau iti vaktavyam . \\
\noindent \textbf{(P\_5,3.71-72.2)} KA\_II,422.22-423.2 Ro\_IV,230-231 {1/7}     akacprakaraṇe tūṣṇīmaḥ kām . \\
\noindent \textbf{(P\_5,3.71-72.2)} KA\_II,422.22-423.2 Ro\_IV,230-231 {2/7}     akacprakaraṇe tūṣṇīmaḥ kām vaktavyaḥ . \\
\noindent \textbf{(P\_5,3.71-72.2)} KA\_II,422.22-423.2 Ro\_IV,230-231 {3/7}     āsitavyam kila tūṣṇīkām etat paśyataḥ cintitam . \\
\noindent \textbf{(P\_5,3.71-72.2)} KA\_II,422.22-423.2 Ro\_IV,230-231 {4/7}     śīle kaḥ malopaḥ ca . \\
\noindent \textbf{(P\_5,3.71-72.2)} KA\_II,422.22-423.2 Ro\_IV,230-231 {5/7}     śīle kaḥ malopaḥ ca vaktavyaḥ . \\
\noindent \textbf{(P\_5,3.71-72.2)} KA\_II,422.22-423.2 Ro\_IV,230-231 {6/7}     tūṣṇīśīlaḥ . \\
\noindent \textbf{(P\_5,3.71-72.2)} KA\_II,422.22-423.2 Ro\_IV,230-231 {7/7}     tūṣṇīkaḥ . \\
\noindent \textbf{(P\_5,3.71-72.3)} KA\_II,423.3-8 Ro\_IV,231 {1/8}     iha bhinatti chinatti iti śanami kṛte śap prāpnoti . \\
\noindent \textbf{(P\_5,3.71-72.3)} KA\_II,423.3-8 Ro\_IV,231 {2/8}     bahukṛtam bahubhuktam bahupītam iti bahuci kṛte kalpādayaḥ prāpnuvanti . \\
\noindent \textbf{(P\_5,3.71-72.3)} KA\_II,423.3-8 Ro\_IV,231 {3/8}     uccakaiḥ nīcakaiḥ akaci kṛte kādayaḥ prāpnuvanti . \\
\noindent \textbf{(P\_5,3.71-72.3)} KA\_II,423.3-8 Ro\_IV,231 {4/8}     nanu ca śnambahujakacaḥ apavādāḥ te bādhakāḥ bhaviṣyanti . \\
\noindent \textbf{(P\_5,3.71-72.3)} KA\_II,423.3-8 Ro\_IV,231 {5/8}     śnambahujakakṣu nānādeśatvāt utsargapratiṣedhaḥ . \\
\noindent \textbf{(P\_5,3.71-72.3)} KA\_II,423.3-8 Ro\_IV,231 {6/8}     śnambahujakakṣu nānādeśatvāt utsargapratiṣedhaḥ vaktavyaḥ . \\
\noindent \textbf{(P\_5,3.71-72.3)} KA\_II,423.3-8 Ro\_IV,231 {7/8}     samānadeśaiḥ apavādaiḥ utsargāṇām bādhanam bhavati . \\
\noindent \textbf{(P\_5,3.71-72.3)} KA\_II,423.3-8 Ro\_IV,231 {8/8}     nānādeśatvāt na prāpnoti . \\
\noindent \textbf{(P\_5,3.71-72.4)} KA\_II,423.9-24 Ro\_IV,231-232 {1/31}     kavidheḥ tamādayaḥ pūrvavipratiṣiddham . \\
\noindent \textbf{(P\_5,3.71-72.4)} KA\_II,423.9-24 Ro\_IV,231-232 {2/31}     kavidheḥ tamādayaḥ bhavanti pūrvavipratiṣedhena . \\
\noindent \textbf{(P\_5,3.71-72.4)} KA\_II,423.9-24 Ro\_IV,231-232 {3/31}     kavidheḥ avakāśaḥ kutsādīnām vacanam prakarṣasya avacanam . \\
\noindent \textbf{(P\_5,3.71-72.4)} KA\_II,423.9-24 Ro\_IV,231-232 {4/31}     paṭukaḥ . \\
\noindent \textbf{(P\_5,3.71-72.4)} KA\_II,423.9-24 Ro\_IV,231-232 {5/31}     mṛdukaḥ . \\
\noindent \textbf{(P\_5,3.71-72.4)} KA\_II,423.9-24 Ro\_IV,231-232 {6/31}     tamādīnām avakāśaḥ prakarṣasya vacanam kutsādīnām avacanam . \\
\noindent \textbf{(P\_5,3.71-72.4)} KA\_II,423.9-24 Ro\_IV,231-232 {7/31}     paṭutaraḥ . \\
\noindent \textbf{(P\_5,3.71-72.4)} KA\_II,423.9-24 Ro\_IV,231-232 {8/31}     paṭutamaḥ . \\
\noindent \textbf{(P\_5,3.71-72.4)} KA\_II,423.9-24 Ro\_IV,231-232 {9/31}     ubhayavacane ubhayam prāpnoti . \\
\noindent \textbf{(P\_5,3.71-72.4)} KA\_II,423.9-24 Ro\_IV,231-232 {10/31}     paṭutarakaḥ . \\
\noindent \textbf{(P\_5,3.71-72.4)} KA\_II,423.9-24 Ro\_IV,231-232 {11/31}     paṭutamakaḥ . \\
\noindent \textbf{(P\_5,3.71-72.4)} KA\_II,423.9-24 Ro\_IV,231-232 {12/31}     tamādayaḥ bhavanti pūrvavipratiṣedhena . \\
\noindent \textbf{(P\_5,3.71-72.4)} KA\_II,423.9-24 Ro\_IV,231-232 {13/31}     kadā cit chinnakatarādayaḥ . \\
\noindent \textbf{(P\_5,3.71-72.4)} KA\_II,423.9-24 Ro\_IV,231-232 {14/31}     kadā cit chinnakatarādayaḥ bhavanti vipratiṣedhena . \\
\noindent \textbf{(P\_5,3.71-72.4)} KA\_II,423.9-24 Ro\_IV,231-232 {15/31}     chinnakataram . \\
\noindent \textbf{(P\_5,3.71-72.4)} KA\_II,423.9-24 Ro\_IV,231-232 {16/31}     chinnakatamam . ekadeśipradhānaḥ ca samāsaḥ . \\
\noindent \textbf{(P\_5,3.71-72.4)} KA\_II,423.9-24 Ro\_IV,231-232 {17/31}     ekadeśipradhānaḥ ca samāsaḥ kavidheḥ bhavati pūrvavipratiṣedhena . \\
\noindent \textbf{(P\_5,3.71-72.4)} KA\_II,423.9-24 Ro\_IV,231-232 {18/31}     ardhapippalikā . \\
\noindent \textbf{(P\_5,3.71-72.4)} KA\_II,423.9-24 Ro\_IV,231-232 {19/31}     ardhakośātakikā . \\
\noindent \textbf{(P\_5,3.71-72.4)} KA\_II,423.9-24 Ro\_IV,231-232 {20/31}     uttarapadārthapradhānaḥ ca sañjñāyām kanvidhyartham . \\
\noindent \textbf{(P\_5,3.71-72.4)} KA\_II,423.9-24 Ro\_IV,231-232 {21/31}     uttarapadārthapradhānaḥ ca samāsaḥ kavidheḥ bhavati pūrvavipratiṣedhena . \\
\noindent \textbf{(P\_5,3.71-72.4)} KA\_II,423.9-24 Ro\_IV,231-232 {22/31}     kim prayojanam . \\
\noindent \textbf{(P\_5,3.71-72.4)} KA\_II,423.9-24 Ro\_IV,231-232 {23/31}     sañjñāyām kanvidhyartham . \\
\noindent \textbf{(P\_5,3.71-72.4)} KA\_II,423.9-24 Ro\_IV,231-232 {24/31}     sañjñāyām kan yathā syāt . \\
\noindent \textbf{(P\_5,3.71-72.4)} KA\_II,423.9-24 Ro\_IV,231-232 {25/31}     navagrāmakam . \\
\noindent \textbf{(P\_5,3.71-72.4)} KA\_II,423.9-24 Ro\_IV,231-232 {26/31}     navarāṣṭrakam . \\
\noindent \textbf{(P\_5,3.71-72.4)} KA\_II,423.9-24 Ro\_IV,231-232 {27/31}     navanagarakam . \\
\noindent \textbf{(P\_5,3.71-72.4)} KA\_II,423.9-24 Ro\_IV,231-232 {28/31}     kadā cit dvandvaḥ . \\
\noindent \textbf{(P\_5,3.71-72.4)} KA\_II,423.9-24 Ro\_IV,231-232 {29/31}     kadā cit dvandvaḥ kavidheḥ bhavati pūrvavipratiṣedhena . \\
\noindent \textbf{(P\_5,3.71-72.4)} KA\_II,423.9-24 Ro\_IV,231-232 {30/31}     plakṣakanyagrodakau . \\
\noindent \textbf{(P\_5,3.71-72.4)} KA\_II,423.9-24 Ro\_IV,231-232 {31/31}     plakṣanyagrodhakau iti . \\
\noindent \textbf{(P\_5,3.74.)} iha kutsitakaḥ anukampitakaḥ iti svaśabdena uktatvāt tasya arthasya pratyayaḥ na prāpnoti . \\
\noindent \textbf{(P\_5,3.74.)} na eṣaḥ doṣaḥ . \\
\noindent \textbf{(P\_5,3.74.)} kutsitasya anukampāyām bhaviṣyati anukampitasya kutsāyām . \\
\noindent \textbf{(P\_5,3.74.)} atha vā svārtham abhidhāyaḥ śabdaḥ nirapekṣaḥ dravyam āha samavetam . \\
\noindent \textbf{(P\_5,3.74.)} samavetasya ca vacane liṅgam vacanam vibhaktim ca . \\
\noindent \textbf{(P\_5,3.74.)} abhidhāya tān viśeṣān apekṣamāṇaḥ ca kṛtsnamātmānam priyakutsanādiṣu punaḥ pravartate asau vibhaktyantaḥ . \\
\noindent \textbf{(P\_5,3.74.)} katham punaḥ idam vijñāyate : kutsitādīnām arthe iti āhosvit kutsitādisamānādhikaraṇāt iti . \\
\noindent \textbf{(P\_5,3.74.)} kaḥ ca atra viśeṣaḥ . \\
\noindent \textbf{(P\_5,3.74.)} kutsidādīnām arthe cet liṅgavacanānupapattiḥ . \\
\noindent \textbf{(P\_5,3.74.)} kutsidādīnām arthe cet liṅgavacanayoḥ anupapattiḥ . \\
\noindent \textbf{(P\_5,3.74.)} paṭukam . \\
\noindent \textbf{(P\_5,3.74.)} paṭukā . \\
\noindent \textbf{(P\_5,3.74.)} paṭukaḥ . \\
\noindent \textbf{(P\_5,3.74.)} paṭukau . \\
\noindent \textbf{(P\_5,3.74.)} paṭukāḥ iti . \\
\noindent \textbf{(P\_5,3.74.)} ekaḥ ayam arthaḥ kutsitam nāma . \\
\noindent \textbf{(P\_5,3.74.)} tasya ekatvāt ekavacanam eva prāpnoti . \\
\noindent \textbf{(P\_5,3.74.)} asti tarhi kutsitādisamānādhikaraṇāt iti . \\
\noindent \textbf{(P\_5,3.74.)} kutsidādisamānādhikaraṇāt iti cet atiprasaṅgaḥ yathā ṭābādiṣu . \\
\noindent \textbf{(P\_5,3.74.)} kutsidādisamānādhikaraṇāt iti cet atiprasaṅgaḥ bhavati yathā ṭābādiṣu . \\
\noindent \textbf{(P\_5,3.74.)} katham ca ṭābādiṣu . \\
\noindent \textbf{(P\_5,3.74.)} uktam tatra strīsamānādhikaraṇāt iti cet bhūtādiṣu atiprasaṅgaḥ iti . \\
\noindent \textbf{(P\_5,3.74.)} evam iha api kutsidādisamānādhikaraṇāt iti cet atiprasaṅgaḥ bhavati . \\
\noindent \textbf{(P\_5,3.74.)} idam ghṛtakam . \\
\noindent \textbf{(P\_5,3.74.)} idam tailakam . \\
\noindent \textbf{(P\_5,3.74.)} idamśabdāt api prāpnoti . \\
\noindent \textbf{(P\_5,3.74.)} siddham tu yena kutsitādivacanam tadyuktāt svārthe pratyayavidhānāt . \\
\noindent \textbf{(P\_5,3.74.)} siddham etat . \\
\noindent \textbf{(P\_5,3.74.)} katham . \\
\noindent \textbf{(P\_5,3.74.)} yena kutsitādayaḥ arthāḥ gamyante tadyuktāt svārthe pratyayaḥ bhavati iti vaktavyam . \\
\noindent \textbf{(P\_5,3.74.)} sidhyati . \\
\noindent \textbf{(P\_5,3.74.)} sūtram tarhi bhidyate . \\
\noindent \textbf{(P\_5,3.74.)} yathānyāsam eva astu . \\
\noindent \textbf{(P\_5,3.74.)} nanu ca uktam kutsidādīnām arthe cet liṅgavacanānupapattiḥ iti . \\
\noindent \textbf{(P\_5,3.74.)} na eṣaḥ doṣaḥ . \\
\noindent \textbf{(P\_5,3.74.)} na ayam pratyayārthaḥ . \\
\noindent \textbf{(P\_5,3.74.)} kim tarhi prakṛtyarthaviśeṣaṇam etat . \\
\noindent \textbf{(P\_5,3.74.)} kutsitādiṣu yat prātipadikam vartate tasmāt kādayaḥ bhavanti . \\
\noindent \textbf{(P\_5,3.74.)} kasmin arthe . \\
\noindent \textbf{(P\_5,3.74.)} svārthe . \\
\noindent \textbf{(P\_5,3.74.)} svāṛthikāḥ ca prakṛtitaḥ liṅgavacanāni anuvartante . \\
\noindent \textbf{(P\_5,3.83.1)} caturthyāt . \\
\noindent \textbf{(P\_5,3.83.1)} caturthyāt lopaḥ vaktavyaḥ . \\
\noindent \textbf{(P\_5,3.83.1)} bṛhaspatidattakaḥ . \\
\noindent \textbf{(P\_5,3.83.1)} bṛhaspatikaḥ . \\
\noindent \textbf{(P\_5,3.83.1)} prajāpatidattakaḥ . \\
\noindent \textbf{(P\_5,3.83.1)} prajāpatikaḥ . \\
\noindent \textbf{(P\_5,3.83.1)} anajādau ca . \\
\noindent \textbf{(P\_5,3.83.1)} anajādau ca lopaḥ vaktavyaḥ . \\
\noindent \textbf{(P\_5,3.83.1)} devadattakaḥ . \\
\noindent \textbf{(P\_5,3.83.1)} devakaḥ . \\
\noindent \textbf{(P\_5,3.83.1)} yajñadattakaḥ . \\
\noindent \textbf{(P\_5,3.83.1)} yajñakaḥ . \\
\noindent \textbf{(P\_5,3.83.1)} lopaḥ pūrvapadasya ca . \\
\noindent \textbf{(P\_5,3.83.1)} pūrvapadasya ca lopaḥ vaktavyaḥ . \\
\noindent \textbf{(P\_5,3.83.1)} devadattakaḥ . \\
\noindent \textbf{(P\_5,3.83.1)} dattakaḥ . \\
\noindent \textbf{(P\_5,3.83.1)} yajñadattakaḥ . \\
\noindent \textbf{(P\_5,3.83.1)} dattakaḥ . \\
\noindent \textbf{(P\_5,3.83.1)} apratyaye tathā eva iṣṭaḥ . \\
\noindent \textbf{(P\_5,3.83.1)} devadattaḥ . \\
\noindent \textbf{(P\_5,3.83.1)} dattaḥ . \\
\noindent \textbf{(P\_5,3.83.1)} yajñadattaḥ . \\
\noindent \textbf{(P\_5,3.83.1)} dattaḥ . \\
\noindent \textbf{(P\_5,3.83.1)} uvarṇāt laḥ ilasya ca . \\
\noindent \textbf{(P\_5,3.83.1)} uvarṇāt ilasya ca lopaḥ vaktavyaḥ . \\
\noindent \textbf{(P\_5,3.83.1)} bhānudattakaḥ . \\
\noindent \textbf{(P\_5,3.83.1)} bhānulaḥ . \\
\noindent \textbf{(P\_5,3.83.1)} vasudattakaḥ . \\
\noindent \textbf{(P\_5,3.83.1)} vasulaḥ . \\
\noindent \textbf{(P\_5,3.83.2)} atha ṭhaggrahaṇam kimartham na ike kṛte ajādau iti eva siddham . \\
\noindent \textbf{(P\_5,3.83.2)} ṭhaggrahaṇam ukaḥ dvitīyatve kavidhānārtham . \\
\noindent \textbf{(P\_5,3.83.2)} ṭhaggrahaṇam kriyate ukaḥ dvitīyatve kavidhiḥ yathā syāt . \\
\noindent \textbf{(P\_5,3.83.2)} vāyudattakaḥ . \\
\noindent \textbf{(P\_5,3.83.2)} vāyukaḥ . \\
\noindent \textbf{(P\_5,3.83.2)} pitṛdattakaḥ . \\
\noindent \textbf{(P\_5,3.83.2)} pitṛkaḥ . \\
\noindent \textbf{(P\_5,3.83.2)} ajādilakṣaṇe hi māthikādivat prasaṅgaḥ . \\
\noindent \textbf{(P\_5,3.83.2)} ajādilakṣaṇe hi māthikādivat prasajyeta . \\
\noindent \textbf{(P\_5,3.83.2)} tat yatha mathitam paṇyam asya māthitikaḥ iti akāralope kṛte tāntāt iti kādeśaḥ na bhavati evam iha api na syāt . \\
\noindent \textbf{(P\_5,3.83.2)} dvitīyāt acaḥ lope sandhyakṣaradvitīyatve tadādeḥ lopavacanam . \\
\noindent \textbf{(P\_5,3.83.2)} dvitīyāt acaḥ lope kartavye sandhyakṣaradvitīyatve tadādeḥ lopaḥ vaktavyaḥ . \\
\noindent \textbf{(P\_5,3.83.2)} lahoḍaḥ . \\
\noindent \textbf{(P\_5,3.83.2)} lahikaḥ . \\
\noindent \textbf{(P\_5,3.83.2)} kahoḍaḥ . \\
\noindent \textbf{(P\_5,3.83.2)} kahikaḥ . \\
\noindent \textbf{(P\_5,3.84)} varuṇādīnām ca tṛtīyāt saḥ ca akṛtasandhīnām . \\
\noindent \textbf{(P\_5,3.84)} varuṇādīnām ca tṛtīyāt lopaḥ ucyate . \\
\noindent \textbf{(P\_5,3.84)} saḥ ca akṛtasandhīnām vaktavyaḥ . \\
\noindent \textbf{(P\_5,3.84)} suparyāśīrdattaḥ . \\
\noindent \textbf{(P\_5,3.84)} suparikaḥ , supariyaḥ , suparilaḥ . \\
\noindent \textbf{(P\_5,3.84)} iha ṣaḍaṅguliḥ ṣaḍikaḥ iti ajādilope kṛte padasañjñā na prāpnoti . \\
\noindent \textbf{(P\_5,3.84)} tatra kaḥ doṣaḥ . \\
\noindent \textbf{(P\_5,3.84)} jaśtvam na syāt . \\
\noindent \textbf{(P\_5,3.84)} ṣaḍike jaśtve uktam . \\
\noindent \textbf{(P\_5,3.84)} kim uktam . \\
\noindent \textbf{(P\_5,3.84)} siddham acaḥ sthānivatvāt iti . \\
\noindent \textbf{(P\_5,3.84)} yadi evam vācikādiṣu padavṛttapratiṣedhaḥ . \\
\noindent \textbf{(P\_5,3.84)} vācikādiṣu padavṛttasya pratiṣedhaḥ vaktavyaḥ . \\
\noindent \textbf{(P\_5,3.84)} siddham ekākṣarapūrvapadānām uttarapadalopavacanāt . \\
\noindent \textbf{(P\_5,3.84)} siddham etat . \\
\noindent \textbf{(P\_5,3.84)} katham . \\
\noindent \textbf{(P\_5,3.84)} ekākṣarapūrvapadānām uttarapadasya lopaḥ vaktavyaḥ . \\
\noindent \textbf{(P\_5,3.84)} iha api tarhi prāpnoti . \\
\noindent \textbf{(P\_5,3.84)} ṣaḍaṅguliḥ . \\
\noindent \textbf{(P\_5,3.84)} ṣaḍikaḥ iti . \\
\noindent \textbf{(P\_5,3.84)} ṣaṣaḥ ṭhājādivacanāt siddham . \\
\noindent \textbf{(P\_5,3.84)} ṣaṣaḥ ṭhājādivacanāt siddham etat . \\
\noindent \textbf{(P\_5,3.85-86)} KA\_II,426.18-21 Ro\_IV,340 {1/10}     kimartham imau ubhau arthau nirdiśyete na yat alpam hrasvam api tat bhavati yat ca hrasvam alpam api tat bhavati . \\
\noindent \textbf{(P\_5,3.85-86)} KA\_II,426.18-21 Ro\_IV,340 {2/10}     na etayoḥ āvaśyakaḥ samāveśaḥ . \\
\noindent \textbf{(P\_5,3.85-86)} KA\_II,426.18-21 Ro\_IV,340 {3/10}     alpam ghṛtam . \\
\noindent \textbf{(P\_5,3.85-86)} KA\_II,426.18-21 Ro\_IV,340 {4/10}     alpam tailam iti ucyate . \\
\noindent \textbf{(P\_5,3.85-86)} KA\_II,426.18-21 Ro\_IV,340 {5/10}     na kaḥ cit āha hrasvam ghṛtam . \\
\noindent \textbf{(P\_5,3.85-86)} KA\_II,426.18-21 Ro\_IV,340 {6/10}     hrasvam tailam iti . \\
\noindent \textbf{(P\_5,3.85-86)} KA\_II,426.18-21 Ro\_IV,340 {7/10}     tathā hrasvaḥ paṭaḥ . \\
\noindent \textbf{(P\_5,3.85-86)} KA\_II,426.18-21 Ro\_IV,340 {8/10}     hrasvaḥ śāṭakaḥ iti ucyate . \\
\noindent \textbf{(P\_5,3.85-86)} KA\_II,426.18-21 Ro\_IV,340 {9/10}     na kaḥ cit āha alpaḥ paṭaḥ . \\
\noindent \textbf{(P\_5,3.85-86)} KA\_II,426.18-21 Ro\_IV,340 {10/10}     alpaḥ śāṭakaḥ iti . \\
\noindent \textbf{(P\_5,3.88)} kuṭīśamīśuṇḍābhyaḥ pratyayasanniyogena puṃvadbhāvaḥ . \\
\noindent \textbf{(P\_5,3.88)} kuṭīśamīśuṇḍābhyaḥ pratyayasanniyogena puṃvadbhāvaḥ vaktavyaḥ . \\
\noindent \textbf{(P\_5,3.88)} kuṭī . \\
\noindent \textbf{(P\_5,3.88)} kuṭīraḥ . \\
\noindent \textbf{(P\_5,3.88)} śamī . \\
\noindent \textbf{(P\_5,3.88)} śamīraḥ . \\
\noindent \textbf{(P\_5,3.88)} śuṇḍā . \\
\noindent \textbf{(P\_5,3.88)} śuṇḍāraḥ iti . \\
\noindent \textbf{(P\_5,3.88)} kim punaḥ kāraṇam na sidhyati . \\
\noindent \textbf{(P\_5,3.88)} svārthikaḥ ayam svārthikāḥ ca prakṛtitaḥ liṅgavacanāni anuvartante . \\
\noindent \textbf{(P\_5,3.88)} uktam vā . \\
\noindent \textbf{(P\_5,3.88)} kim uktam . \\
\noindent \textbf{(P\_5,3.88)} svāṛthikāḥ ativartante api liṅgavacanāni iti . \\
\noindent \textbf{(P\_5,3.91)} vatsādibhyaḥ tanutve kārśye pratiṣedhaḥ . \\
\noindent \textbf{(P\_5,3.91)} vatsādibhyaḥ tanutve kārśye pratiṣedhaḥ vaktavyaḥ . \\
\noindent \textbf{(P\_5,3.91)} kṛśaḥ vatsaḥ vatsataraḥ iti mā bhūt iti . \\
\noindent \textbf{(P\_5,3.91)} saḥ tarhi pratiṣedhaḥ vaktavyaḥ . \\
\noindent \textbf{(P\_5,3.91)} na vaktavyaḥ . \\
\noindent \textbf{(P\_5,3.91)} yasya guṇasya hi bhāvāt dravye śabdaniveśaḥ tadabhidhāne tasmin guṇe vaktavye pratyayena bhavitavyam . \\
\noindent \textbf{(P\_5,3.91)} na ca kārśyasya sadbhāvāt dravye vatsaśabdaḥ . \\
\noindent \textbf{(P\_5,3.92-93)} KA\_II,427.16-428.4 Ro\_IV, 242 {1/12}     kimādīnām dvibahvarthe pratyayavidhānāt upādhyānarthakyam . \\
\noindent \textbf{(P\_5,3.92-93)} KA\_II,427.16-428.4 Ro\_IV, 242 {2/12}     kimādīnām dvibahvarthe pratyayavidhānāt upādhigrahaṇam anarthakam . \\
\noindent \textbf{(P\_5,3.92-93)} KA\_II,427.16-428.4 Ro\_IV, 242 {3/12}     kim kāraṇam . \\
\noindent \textbf{(P\_5,3.92-93)} KA\_II,427.16-428.4 Ro\_IV, 242 {4/12}     bahirdhāraṇam nirdhāraṇam . \\
\noindent \textbf{(P\_5,3.92-93)} KA\_II,427.16-428.4 Ro\_IV, 242 {5/12}     yāvatā dvayoḥ ekasya eva bahirdhāraṇam bhavati . \\
\noindent \textbf{(P\_5,3.92-93)} KA\_II,427.16-428.4 Ro\_IV, 242 {6/12}     aparaḥ āha : bahūnām jatiparipraśne ḍatamac iti atra bahugrahaṇam anarthakam . \\
\noindent \textbf{(P\_5,3.92-93)} KA\_II,427.16-428.4 Ro\_IV, 242 {7/12}     kim kāraṇam . \\
\noindent \textbf{(P\_5,3.92-93)} KA\_II,427.16-428.4 Ro\_IV, 242 {8/12}     kim iti etat paripraśne vartate paripraśnaḥ ca anirjñāte anirjñātam ca bahuṣu . \\
\noindent \textbf{(P\_5,3.92-93)} KA\_II,427.16-428.4 Ro\_IV, 242 {9/12}     dvyekayoḥ punaḥ nirjñātam . \\
\noindent \textbf{(P\_5,3.92-93)} KA\_II,427.16-428.4 Ro\_IV, 242 {10/12}     nirjñātatvāt dvyekayoḥ paripraśnaḥ na . \\
\noindent \textbf{(P\_5,3.92-93)} KA\_II,427.16-428.4 Ro\_IV, 242 {11/12}     paripraśnābhāvāt kim eva na asti . \\
\noindent \textbf{(P\_5,3.92-93)} KA\_II,427.16-428.4 Ro\_IV, 242 {12/12}     kutaḥ pratyayaḥ . \\
\noindent \textbf{(P\_5,3.94)} prāgvacanam kimartham . \\
\noindent \textbf{(P\_5,3.94)} vibhāṣā yathā syāt . \\
\noindent \textbf{(P\_5,3.94)} prāgvacanānarthakyam ca vibhāṣāprakaraṇāt . \\
\noindent \textbf{(P\_5,3.94)} prāgvacanam anarthakam . \\
\noindent \textbf{(P\_5,3.94)} kim kāraṇam . \\
\noindent \textbf{(P\_5,3.94)} prakṛtā mahāvibhāṣā . \\
\noindent \textbf{(P\_5,3.94)} taya eva siddham . \\
\noindent \textbf{(P\_5,3.95)} avakṣepaṇe kan vidhīyate kutsite kaḥ . \\
\noindent \textbf{(P\_5,3.95)} kaḥ etayoḥ arthayoḥ viśeṣaḥ . \\
\noindent \textbf{(P\_5,3.95)} avakṣepaṇam karaṇam kutsitam karma . \\
\noindent \textbf{(P\_5,3.95)} avakṣepaṇam vai kutsitam karaṇam . \\
\noindent \textbf{(P\_5,3.95)} tena yat kutysyate tat api kutsitam bhavati . \\
\noindent \textbf{(P\_5,3.95)} tatra kutsitam iti eva siddham bhavati . \\
\noindent \textbf{(P\_5,3.95)} evam tarhi yat parasya kutsārtham upādīyate tat iha udāharaṇam . \\
\noindent \textbf{(P\_5,3.95)} vyākaraṇakena nāma ayam garvitaḥ . \\
\noindent \textbf{(P\_5,3.95)} yājñikyena nāma ayam garvitaḥ . \\
\noindent \textbf{(P\_5,3.95)} yat svakutsārtham kutsārtham upādīyate tat tatra udāharaṇam . \\
\noindent \textbf{(P\_5,3.95)} devadattakaḥ . \\
\noindent \textbf{(P\_5,3.95)} yajñadattakaḥ . \\
\noindent \textbf{(P\_5,3.98)} kimartham manuṣye lup ucyate na luk eva ucyeta . \\
\noindent \textbf{(P\_5,3.98)} liṅgasiddhyartham lup manuṣye . \\
\noindent \textbf{(P\_5,3.98)} liṅgasiddhyartham manuṣye lup ucyate . \\
\noindent \textbf{(P\_5,3.98)} cañcā iva cañcā . \\
\noindent \textbf{(P\_5,3.98)} vadhrikā iva vadhrikā . \\
\noindent \textbf{(P\_5,3.98)} kharakuṭī iva kharakuṭī . \\
\noindent \textbf{(P\_5,3.99)} apaṇye iti ucyate . \\
\noindent \textbf{(P\_5,3.99)} tatra idam na sidhyati . \\
\noindent \textbf{(P\_5,3.99)} śivaḥ . \\
\noindent \textbf{(P\_5,3.99)} skandaḥ . \\
\noindent \textbf{(P\_5,3.99)} viśākhaḥ iti . \\
\noindent \textbf{(P\_5,3.99)} kim kāraṇam . \\
\noindent \textbf{(P\_5,3.99)} mauryaiḥ hiraṇyāṛthibhiḥ arcāḥ prakalpitāḥ . \\
\noindent \textbf{(P\_5,3.99)} bhavet tāsu na syāt . \\
\noindent \textbf{(P\_5,3.99)} yāḥ tu etāḥ sampratipūjārthāḥ tāsu bhaviṣyati . \\
\noindent \textbf{(P\_5,3.106)} tat iti anena kim pratinirdiśyate . \\
\noindent \textbf{(P\_5,3.106)} chaḥ . \\
\noindent \textbf{(P\_5,3.106)} katham punaḥ samāsaḥ nāma chaviṣayaḥ syāt . \\
\noindent \textbf{(P\_5,3.106)} evam tarhi ivārthaḥ . \\
\noindent \textbf{(P\_5,3.106)} yadi tarhi samāsaḥ api ivārthe pratyayaḥ api samāsenoktatvāt pratyayaḥ na prāpnoti . \\
\noindent \textbf{(P\_5,3.106)} evam tarhi dvau ivārthau . \\
\noindent \textbf{(P\_5,3.106)} katham . \\
\noindent \textbf{(P\_5,3.106)} kākāgamanam iva tālapatanam iva kākatālam . \\
\noindent \textbf{(P\_5,3.106)} kākatālam iva kākatālīyam . \\
\noindent \textbf{(P\_5,3.118)} aṇaḥ gotrāt gotravacanam . \\
\noindent \textbf{(P\_5,3.118)} aṇaḥ gotrāt gotragrahaṇam kartavyam . \\
\noindent \textbf{(P\_5,3.118)} gotrāt iti vaktavyam . \\
\noindent \textbf{(P\_5,3.118)} iha mā bhūt . \\
\noindent \textbf{(P\_5,3.118)} ābhijitaḥ muhūrtaḥ . \\
\noindent \textbf{(P\_5,3.118)} ābhijitaḥ sthālīpākaḥ iti . \\
\noindent \textbf{(P\_5,3.118)} gotram iti ca vaktavyam . \\
\noindent \textbf{(P\_5,3.118)} kim prayojanam . \\
\noindent \textbf{(P\_5,3.118)} ābhijitakaḥ . \\
\noindent \textbf{(P\_5,3.118)} gotrāśrayaḥ vuñ yathā syāt . \\
\noindent \textbf{(P\_5,3.118)} gotram iti śakyam akartum . \\
\noindent \textbf{(P\_5,3.118)} katham ābhijitakaḥ . \\
\noindent \textbf{(P\_5,3.118)} gotrāt ayam svārthikaḥ gotram eva bhavati . \\
\noindent \textbf{(P\_5,4.1)} pādaśatagrahaṇam anarthakam anyatra api darśanāt . \\
\noindent \textbf{(P\_5,4.1)} pādaśatagrahaṇam anarthakam . \\
\noindent \textbf{(P\_5,4.1)} kim kāraṇam . \\
\noindent \textbf{(P\_5,4.1)} anyatra api darśanāt . \\
\noindent \textbf{(P\_5,4.1)} anyatra api hi vun dṛśyate . \\
\noindent \textbf{(P\_5,4.1)} dvimodikām dadāti . \\
\noindent \textbf{(P\_5,4.3)} kanprakaraṇe cañcadbrhatoḥ upasaṅkhyānam . \\
\noindent \textbf{(P\_5,4.3)} kanprakaraṇe cañcadbrhatoḥ upasaṅkhyānam kartavyam . \\
\noindent \textbf{(P\_5,4.3)} ca~catkaḥ . \\
\noindent \textbf{(P\_5,4.3)} bṛhatkaḥ . \\
\noindent \textbf{(P\_5,4.4)} anatyantagatau ktāntāt tamādayaḥ pūrvavipratiṣiddham . \\
\noindent \textbf{(P\_5,4.4)} anatyantagatau ktāntāt tamādayaḥ bhavanti pūrvavipratiṣedhena . \\
\noindent \textbf{(P\_5,4.4)} anatyantagatau ktāntāt kan bhavati iti asya avakāśaḥ anatyantagateḥ vacanam prakarṣasya avacanam . \\
\noindent \textbf{(P\_5,4.4)} bhinnakam . \\
\noindent \textbf{(P\_5,4.4)} chinnakam . \\
\noindent \textbf{(P\_5,4.4)} tamādīnām avakāśaḥ prakarṣasya vacanam anatyantagateḥ avacanam . \\
\noindent \textbf{(P\_5,4.4)} paṭutaraḥ . \\
\noindent \textbf{(P\_5,4.4)} paṭutamaḥ . \\
\noindent \textbf{(P\_5,4.4)} ubhayavacane ubhayam prāpnoti . \\
\noindent \textbf{(P\_5,4.4)} bhinnatarakam . \\
\noindent \textbf{(P\_5,4.4)} chinnatarakam . \\
\noindent \textbf{(P\_5,4.4)} tamādayaḥ bhavanti pūrvavipratiṣedhena . \\
\noindent \textbf{(P\_5,4.4)} tadantāt ca svārthe kanvacanam . \\
\noindent \textbf{(P\_5,4.4)} tadantāt ca svārthe kan vaktavyaḥ . \\
\noindent \textbf{(P\_5,4.4)} bhinnatarakam . \\
\noindent \textbf{(P\_5,4.5)} sāmivacane pratiṣedhānarthakyam prakṛtyabhihitatvāt . \\
\noindent \textbf{(P\_5,4.5)} sāmivacane pratiṣedhaḥ anarthakaḥ . \\
\noindent \textbf{(P\_5,4.5)} kim kāraṇam . \\
\noindent \textbf{(P\_5,4.5)} prakṛtyabhihitatvāt . \\
\noindent \textbf{(P\_5,4.5)} prakṛtyabhihitaḥ saḥ arthaḥ iti kṛtvā kan na bhaviṣyati . \\
\noindent \textbf{(P\_5,4.7)} adhyuttarapadāt pratyayavidhānānupapattiḥ vigrahābhāvāt . \\
\noindent \textbf{(P\_5,4.7)} adhyuttarapadāt pratyayavidheḥ anupapattiḥ . \\
\noindent \textbf{(P\_5,4.7)} kim kāraṇam . \\
\noindent \textbf{(P\_5,4.7)} vigrahābhāvāt . \\
\noindent \textbf{(P\_5,4.7)} vigrahapūrvikā taddhitotpattiḥ . \\
\noindent \textbf{(P\_5,4.7)} na ca adhyuttarapadena vigrahaḥ dṛśyate . \\
\noindent \textbf{(P\_5,4.7)} tasmāt tatra idam iti sadhīnar . \\
\noindent \textbf{(P\_5,4.7)} tasmāt tatra idam iti sadhīnar pratyayaḥ vaktavyaḥ . \\
\noindent \textbf{(P\_5,4.7)} rājani idam rājādhīnam . \\
\noindent \textbf{(P\_5,4.7)} yadi sadhīnar kriyate sakārasya itsañjñā na prāpnoti . \\
\noindent \textbf{(P\_5,4.7)} iha ca śryadhīnaḥ bhrvadhīnaḥ iti aṅgasya iti iyaṅuvaṅau syātām . \\
\noindent \textbf{(P\_5,4.7)} sūtram ca bhidyate . \\
\noindent \textbf{(P\_5,4.7)} yathānyāsam eva astu . \\
\noindent \textbf{(P\_5,4.7)} nanu ca uktam adhyuttarapadāt pratyayavidhānānupapattiḥ vigrahābhāvāt iti . \\
\noindent \textbf{(P\_5,4.7)} na eṣaḥ doṣaḥ . \\
\noindent \textbf{(P\_5,4.7)} asti kāraṇam yena atra vigrahaḥ na bhavati . \\
\noindent \textbf{(P\_5,4.7)} kim kāraṇam . \\
\noindent \textbf{(P\_5,4.7)} nityapratyayaḥ ayam . \\
\noindent \textbf{(P\_5,4.7)} ke punaḥ nityapratyayāḥ . \\
\noindent \textbf{(P\_5,4.7)} tamādayaḥ prāk kanaḥ ñyādayaḥ prāk vunaḥ āmādayaḥ prāk mayaṭaḥ bṛhatījātyantāḥ samāsāntāḥ ca iti . \\
\noindent \textbf{(P\_5,4.7)} evam tarhi na ayam pratyayavidhiḥ upālabhyate . \\
\noindent \textbf{(P\_5,4.7)} kim tarhi . \\
\noindent \textbf{(P\_5,4.7)} prakṛtiḥ upālabhyate . \\
\noindent \textbf{(P\_5,4.7)} adhyuttarapadā prakṛtiḥ na asti . \\
\noindent \textbf{(P\_5,4.7)} kim kāraṇam . \\
\noindent \textbf{(P\_5,4.7)} vigrahābhāvāt . \\
\noindent \textbf{(P\_5,4.7)} vigrahapūrvikā samāsavṛttiḥ . \\
\noindent \textbf{(P\_5,4.7)} na ca adhinā vigrahaḥ dṛśyate . \\
\noindent \textbf{(P\_5,4.7)} evam tarhi bahuvrīhiḥ bhaviṣyati . \\
\noindent \textbf{(P\_5,4.7)} kim kṛtam bhavati . \\
\noindent \textbf{(P\_5,4.7)} bhavati vai kaḥ cit asvapadavigrahaḥ api bahuvrīhiḥ . \\
\noindent \textbf{(P\_5,4.7)} tat yathā śobhanam mukham asyāḥ sumukhī iti . \\
\noindent \textbf{(P\_5,4.7)} na evam śakyam . \\
\noindent \textbf{(P\_5,4.7)} iha hi mahadadhīnam iti āttvakapau prasajyeyātām . \\
\noindent \textbf{(P\_5,4.7)} evam tarhi avyayībhāvaḥ bhaviṣyati . \\
\noindent \textbf{(P\_5,4.7)} evam api adheḥ pūrvanipātaḥ prāpnoti . \\
\noindent \textbf{(P\_5,4.7)} rājadantādiṣu pāṭhaḥ kariṣyate . \\
\noindent \textbf{(P\_5,4.7)} atha vā saptamīsamāsaḥ ayam . \\
\noindent \textbf{(P\_5,4.7)} adhiḥ śauṇḍādiṣu paṭhyate . \\
\noindent \textbf{(P\_5,4.8)} diggrahaṇam kimartham . \\
\noindent \textbf{(P\_5,4.8)} astriyām iti iyati ucyamāne prācīnā brāhmaṇī avācīnā śikhā iti atra api prasajyeta . \\
\noindent \textbf{(P\_5,4.8)} diggrahaṇe punaḥ kriyamāṇe na doṣaḥ bhavati . \\
\noindent \textbf{(P\_5,4.8)} atha strīgrahaṇam kimartham yāvatā dikśabdaḥ strīviṣayaḥ eva . \\
\noindent \textbf{(P\_5,4.8)} bhavati vai kaḥ cit dikśabdaḥ astrīviṣayaḥ api . \\
\noindent \textbf{(P\_5,4.8)} tat yathā prāk prācīnam . \\
\noindent \textbf{(P\_5,4.8)} pratyak pratīcīnam . \\
\noindent \textbf{(P\_5,4.8)} ucak udīcīnam . \\
\noindent \textbf{(P\_5,4.14)} strīgrahaṇam kimartham na svārthikaḥ ayam svārthikāḥ ca prakṛtitaḥ liṅgavacanāni anuvartante . \\
\noindent \textbf{(P\_5,4.14)} evam tarhi siddhe sati yat strīgrahaṇam karoti tat jñāpayati ācāryaḥ svārthikāḥ ativartante api liṅgavacanāni . \\
\noindent \textbf{(P\_5,4.14)} kim etasya jñāpane prayojanam . \\
\noindent \textbf{(P\_5,4.14)} guḍakalpā drākṣā . \\
\noindent \textbf{(P\_5,4.14)} tailakalpā prasannā payaskalpā yavāgūḥ iti etat siddham bhavati . \\
\noindent \textbf{(P\_5,4.19)} sakṛdādeśe abhyāvṛttigrahaṇam nivartyam . \\
\noindent \textbf{(P\_5,4.19)} kim prayojanam . \\
\noindent \textbf{(P\_5,4.19)} punaḥ punaḥ āvṛttiḥ abhyāvṛttiḥ . \\
\noindent \textbf{(P\_5,4.19)} na ca ekasya punaḥ punaḥ āvṛttiḥ bhavati . \\
\noindent \textbf{(P\_5,4.19)} atha kriyāgrahaṇam anuvartate āhosvit na . \\
\noindent \textbf{(P\_5,4.19)} kim ca arthaḥ anuvṛttyā . \\
\noindent \textbf{(P\_5,4.19)} bāḍham arthaḥ . \\
\noindent \textbf{(P\_5,4.19)} iha mā bhūt : ekaḥ bhuṅkte iti . \\
\noindent \textbf{(P\_5,4.19)} atha anuvartamāne api kriyāgrahaṇe iha kasmāt na bhavati . \\
\noindent \textbf{(P\_5,4.19)} ekaḥ pākaḥ iti . \\
\noindent \textbf{(P\_5,4.19)} pūrvayoḥ ca yogayoḥ kasmāt na bhavati . \\
\noindent \textbf{(P\_5,4.19)} dvau pākau. trayaḥ pākāḥ . \\
\noindent \textbf{(P\_5,4.19)} catvāraḥ pākāḥ . \\
\noindent \textbf{(P\_5,4.19)} pañca pākāḥ . \\
\noindent \textbf{(P\_5,4.19)} daśa pākāḥ iti . \\
\noindent \textbf{(P\_5,4.19)} na etat kriyāgaṇanam . \\
\noindent \textbf{(P\_5,4.19)} kim tarhi . \\
\noindent \textbf{(P\_5,4.19)} dravyagaṇanam etat . \\
\noindent \textbf{(P\_5,4.19)} katham . \\
\noindent \textbf{(P\_5,4.19)} kṛdabhihitaḥ bhāvaḥ dravyavat bhavati iti . \\
\noindent \textbf{(P\_5,4.19)} iha api tarhi dravyagaṇanāt na prāpnoti . \\
\noindent \textbf{(P\_5,4.19)} sakṛt bhuktvā . \\
\noindent \textbf{(P\_5,4.19)} sakṛt bhoktum iti . \\
\noindent \textbf{(P\_5,4.19)} pūrvayoḥ ca yogayoḥ dviḥ bhuktvā dviḥ boktum triḥ bhuktvā triḥ bhoktum pañcakṛtvā bhuktvā pañcakṛtvā bhoktum daśakṛtvā bhuktvā daśakṛtvā bhoktum iti dravyagaṇanān na prāpnoti . \\
\noindent \textbf{(P\_5,4.19)} yadi khalu api punaḥ punaḥ āvṛttiḥ abhyāvṛttiḥ dviḥ āvṛtte sakṛt iti syāt triḥ āvṛtte dviḥ iti . \\
\noindent \textbf{(P\_5,4.19)} evam tarhi anuvartate abhyāvṛttigrahaṇam na tu punaḥ punaḥ āvṛttiḥ abhyāvṛttiḥ . \\
\noindent \textbf{(P\_5,4.19)} kim tarhi abhimukḥī pravṛttiḥ abhyāvṛttiḥ . \\
\noindent \textbf{(P\_5,4.19)} pūrvā ca pare prati abhimukḥī pare ca pūrvām prati abhimukhyau . \\
\noindent \textbf{(P\_5,4.19)} yat api ucyate anuvartamāne api kriyāgrahaṇe iha kasmāt na bhavati . \\
\noindent \textbf{(P\_5,4.19)} ekaḥ pākaḥ iti . \\
\noindent \textbf{(P\_5,4.19)} pūrvayoḥ ca yogayoḥ kasmāt na bhavati . \\
\noindent \textbf{(P\_5,4.19)} dvau pākau. trayaḥ pākāḥ . \\
\noindent \textbf{(P\_5,4.19)} catvāraḥ pākāḥ . \\
\noindent \textbf{(P\_5,4.19)} pañca pākāḥ . \\
\noindent \textbf{(P\_5,4.19)} daśa pākāḥ iti parihṛtam etat . \\
\noindent \textbf{(P\_5,4.19)} na etat kriyāgaṇanam . \\
\noindent \textbf{(P\_5,4.19)} kim tarhi . \\
\noindent \textbf{(P\_5,4.19)} dravyagaṇanam etat . \\
\noindent \textbf{(P\_5,4.19)} katham . \\
\noindent \textbf{(P\_5,4.19)} kṛdabhihitaḥ bhāvaḥ dravyavat bhavati iti . \\
\noindent \textbf{(P\_5,4.19)} nanu ca uktam iha api tarhi dravyagaṇanāt na prāpnoti . \\
\noindent \textbf{(P\_5,4.19)} sakṛt bhuktvā . \\
\noindent \textbf{(P\_5,4.19)} sakṛt bhoktum iti . \\
\noindent \textbf{(P\_5,4.19)} pūrvayoḥ ca yogayoḥ dviḥ bhuktvā dviḥ boktum triḥ bhuktvā triḥ bhoktum pañcakṛtvā bhuktvā pañcakṛtvā bhoktum daśakṛtvā bhuktvā daśakṛtvā bhoktum iti dravyagaṇanān na prāpnoti . \\
\noindent \textbf{(P\_5,4.19)} na eṣaḥ doṣaḥ . \\
\noindent \textbf{(P\_5,4.19)} kriyāgaṇanāt bhaviṣyati . \\
\noindent \textbf{(P\_5,4.19)} katham . \\
\noindent \textbf{(P\_5,4.19)} kṛdabhihitaḥ bhāvaḥ dravyavat api kriyāvat api bhavati . \\
\noindent \textbf{(P\_5,4.24)} devatāntāt iti ucyate . \\
\noindent \textbf{(P\_5,4.24)} tata idam na sidhyati . \\
\noindent \textbf{(P\_5,4.24)} pitṛdevatyam iti . \\
\noindent \textbf{(P\_5,4.24)} kim kāraṇam . \\
\noindent \textbf{(P\_5,4.24)} na hi lpitaraḥ devatā . \\
\noindent \textbf{(P\_5,4.24)} na eṣaḥ doṣaḥ . \\
\noindent \textbf{(P\_5,4.24)} diveḥ aiśvaryakarmaṇaḥ devaḥ . \\
\noindent \textbf{(P\_5,4.24)} tasmāt svārthe tal . \\
\noindent \textbf{(P\_5,4.24)} evam ca kṛtvā devadevatyam api siddham bhavati . \\
\noindent \textbf{(P\_5,4.27)} tali strīliṅgavacanam . \\
\noindent \textbf{(P\_5,4.27)} tali strīliṅgam vaktavyam . \\
\noindent \textbf{(P\_5,4.27)} devatā . \\
\noindent \textbf{(P\_5,4.27)} kim punaḥ kāraṇam na sidhyati . \\
\noindent \textbf{(P\_5,4.27)} devaśabdaḥ ayam puṃliṅgaḥ svārthikaḥ ca ayam . \\
\noindent \textbf{(P\_5,4.27)} svārthikāḥ ca prakṛtitaḥ liṅgavacanāni anuvartante . \\
\noindent \textbf{(P\_5,4.27)} uktam vā . \\
\noindent \textbf{(P\_5,4.27)} kim uktam . \\
\noindent \textbf{(P\_5,4.27)} svārthikāḥ ativartante api liṅgavacanāni iti . \\
\noindent \textbf{(P\_5,4.30)} lohitāt liṅgabādhanam vā . \\
\noindent \textbf{(P\_5,4.30)} lohitāt liṅgabādhanam vā iti vaktavyam . \\
\noindent \textbf{(P\_5,4.30)} lohitika . \\
\noindent \textbf{(P\_5,4.30)} lohinikā . \\
\noindent \textbf{(P\_5,4.30)} akṣarasamūhe chandasaḥ upasaṅkhyānam . \\
\noindent \textbf{(P\_5,4.30)} akṣarasamūhe chandasaḥ upasaṅkhyānam kartavyam . \\
\noindent \textbf{(P\_5,4.30)} o śrāvaya iti caturakṣaram . \\
\noindent \textbf{(P\_5,4.30)} astu śrauṣaṭ iti caturakṣaram . \\
\noindent \textbf{(P\_5,4.30)} ye yajāmahe iti pañcākṣaram . \\
\noindent \textbf{(P\_5,4.30)} yaja iti dvyakṣaram . \\
\noindent \textbf{(P\_5,4.30)} dvyakṣaraḥ vaṣaṭkāraḥ . \\
\noindent \textbf{(P\_5,4.30)} eṣaḥ vai saptadaśākṣaraḥ chandasyaḥ prajñāpatiḥ yajñam anu vihitaḥ . \\
\noindent \textbf{(P\_5,4.30)} chandasi bahubhirvasavyairupasaṅkhyānam . \\
\noindent \textbf{(P\_5,4.30)} chandasi bahubhirvasavyairupasaṅkhyānam . \\
\noindent \textbf{(P\_5,4.30)} hastau pṛṇasva bahuviḥ vasavyaiḥ . \\
\noindent \textbf{(P\_5,4.30)} agnirīśevasavyasya . \\
\noindent \textbf{(P\_5,4.30)} agnirīśevasavyasya upasaṅkhyānam kartavyam . \\
\noindent \textbf{(P\_5,4.30)} uktam vā . \\
\noindent \textbf{(P\_5,4.30)} kim uktam . \\
\noindent \textbf{(P\_5,4.30)} svārthavijñānāt siddham iti . \\
\noindent \textbf{(P\_5,4.30)} apasyaḥ vasānāḥ . \\
\noindent \textbf{(P\_5,4.30)} apaḥ vasānāḥ . \\
\noindent \textbf{(P\_5,4.30)} sve okye . \\
\noindent \textbf{(P\_5,4.30)} sve oke . \\
\noindent \textbf{(P\_5,4.30)} kavyaḥ asi havyasūdana . \\
\noindent \textbf{(P\_5,4.30)} kaviḥ asi . \\
\noindent \textbf{(P\_5,4.30)} raudreṇa anīkena kavyatāyai . \\
\noindent \textbf{(P\_5,4.30)} kavitayai . \\
\noindent \textbf{(P\_5,4.30)} āmuṣyāyaṇasya . \\
\noindent \textbf{(P\_5,4.30)} amuṣyaputrasya . \\
\noindent \textbf{(P\_5,4.30)} kṣemyasya īśe . \\
\noindent \textbf{(P\_5,4.30)} kṣemasya īśe . \\
\noindent \textbf{(P\_5,4.30)} kṣemyam adhyavasyati . \\
\noindent \textbf{(P\_5,4.30)} kṣemam adhyavasyati . \\
\noindent \textbf{(P\_5,4.30)} āyuḥ varcasyam . \\
\noindent \textbf{(P\_5,4.30)} varcaḥ eva varcasyam . \\
\noindent \textbf{(P\_5,4.30)} niṣkevalyam . \\
\noindent \textbf{(P\_5,4.30)} niṣkevalam . \\
\noindent \textbf{(P\_5,4.30)} ukthyam . \\
\noindent \textbf{(P\_5,4.30)} uktham . \\
\noindent \textbf{(P\_5,4.30)} janyam tābhiḥ sajanyam tābhiḥ . \\
\noindent \textbf{(P\_5,4.30)} janam tābhiḥ sajanaṃ tābhiḥ . \\
\noindent \textbf{(P\_5,4.30)} stomaiḥ janayāmi navyam . \\
\noindent \textbf{(P\_5,4.30)} navam . \\
\noindent \textbf{(P\_5,4.30)} pra naḥ navyebhiḥ . \\
\noindent \textbf{(P\_5,4.30)} navaiḥ . \\
\noindent \textbf{(P\_5,4.30)} brahma pūrvyam . \\
\noindent \textbf{(P\_5,4.30)} pāthaḥ pūrvyam . \\
\noindent \textbf{(P\_5,4.30)} tanuṣu pūrvyam . \\
\noindent \textbf{(P\_5,4.30)} pūrvam . \\
\noindent \textbf{(P\_5,4.30)} pūrvyāhaḥ . \\
\noindent \textbf{(P\_5,4.30)} pūrvāhaḥ . \\
\noindent \textbf{(P\_5,4.30)} pūrvyāḥ viśaḥ . \\
\noindent \textbf{(P\_5,4.30)} pūrvāḥ viśaḥ . \\
\noindent \textbf{(P\_5,4.30)} pūrvyāsaḥ . \\
\noindent \textbf{(P\_5,4.30)} pūrvāsaḥ . \\
\noindent \textbf{(P\_5,4.30)} saḥ pra pūrvyaḥ . \\
\noindent \textbf{(P\_5,4.30)} saḥ pra pūrvaḥ . \\
\noindent \textbf{(P\_5,4.30)} agnim vai pūrvyam . \\
\noindent \textbf{(P\_5,4.30)} pūrvam . \\
\noindent \textbf{(P\_5,4.30)} tam juṣasva yaviṣṭhya . \\
\noindent \textbf{(P\_5,4.30)} yaviṣṭha . \\
\noindent \textbf{(P\_5,4.30)} hotravāham yaviṣṭhyam . \\
\noindent \textbf{(P\_5,4.30)} yaviṣṭham . \\
\noindent \textbf{(P\_5,4.30)} tvam ha yat yaviṣṭhya . \\
\noindent \textbf{(P\_5,4.30)} yaviṣṭha . \\
\noindent \textbf{(P\_5,4.30)} samāvat vasati samāvat gṛhṇāti . \\
\noindent \textbf{(P\_5,4.30)} samam vasati samam gṛhṇāti . \\
\noindent \textbf{(P\_5,4.30)} samāvat devayajñe hastau . \\
\noindent \textbf{(P\_5,4.30)} samam . \\
\noindent \textbf{(P\_5,4.30)} samāvat vīryāvahāni . \\
\noindent \textbf{(P\_5,4.30)} samāni . \\
\noindent \textbf{(P\_5,4.30)} samāvat vIryāṇi karoti . \\
\noindent \textbf{(P\_5,4.30)} samāni . \\
\noindent \textbf{(P\_5,4.30)} u īvate u lokam . \\
\noindent \textbf{(P\_5,4.30)} yaḥ īvate brahmaṇe . \\
\noindent \textbf{(P\_5,4.30)} yaḥ iyate . \\
\noindent \textbf{(P\_5,4.30)} navasya nūtnaptanakhāḥ ca . \\
\noindent \textbf{(P\_5,4.30)} navasya nū iti ayam ādeśaḥ vaktavyaḥ tnaptanakhāḥ ca pratyayāḥ vaktavyāḥ . \\
\noindent \textbf{(P\_5,4.30)} nūtnam . \\
\noindent \textbf{(P\_5,4.30)} nūtanam . \\
\noindent \textbf{(P\_5,4.30)} navīnam . \\
\noindent \textbf{(P\_5,4.30)} naḥ ca purāṇe prāt . \\
\noindent \textbf{(P\_5,4.30)} naḥ ca purāṇe prāt vaktavyaḥ tnaptanakhāḥ ca pratyayāḥ vaktavyāḥ . \\
\noindent \textbf{(P\_5,4.30)} praṇam . \\
\noindent \textbf{(P\_5,4.30)} pratnam . \\
\noindent \textbf{(P\_5,4.30)} pratanam . \\
\noindent \textbf{(P\_5,4.30)} prīṇam . \\
\noindent \textbf{(P\_5,4.36)} tat iti anena kim pratinirdiśyate . \\
\noindent \textbf{(P\_5,4.36)} vāk eva . \\
\noindent \textbf{(P\_5,4.36)} yat eva vācā vyavahiryate tat karmaṇā kriyate . \\
\noindent \textbf{(P\_5,4.36)} aṇprakaraṇe kulālavaruḍaniṣādacaṇḍālāmitrebhyaḥ chandasi . \\
\noindent \textbf{(P\_5,4.36)} aṇprakaraṇe kulālavaruḍaniṣādacaṇḍālāmitrebhyaḥ chandasi upasaṅkhyānam kartavyam . \\
\noindent \textbf{(P\_5,4.36)} kaulālaḥ . \\
\noindent \textbf{(P\_5,4.36)} vāṛuḍaḥ . \\
\noindent \textbf{(P\_5,4.36)} naiṣādaḥ . \\
\noindent \textbf{(P\_5,4.36)} cāṇḍālaḥ . \\
\noindent \textbf{(P\_5,4.36)} āmitraḥ . \\
\noindent \textbf{(P\_5,4.36)} bhāgarūpanāmabhyaḥ dheyaḥ . \\
\noindent \textbf{(P\_5,4.36)} bhāgarūpanāmabhyaḥ dheyaḥ vaktavyaḥ . \\
\noindent \textbf{(P\_5,4.36)} bhāgadheyam . \\
\noindent \textbf{(P\_5,4.36)} rūpadheyam . \\
\noindent \textbf{(P\_5,4.36)} nāmadheyam . \\
\noindent \textbf{(P\_5,4.36)} mitrāt chandasi . \\
\noindent \textbf{(P\_5,4.36)} mitrāt chandasi dheyaḥ vaktavyaḥ . \\
\noindent \textbf{(P\_5,4.36)} mitradheye yatasva . \\
\noindent \textbf{(P\_5,4.36)} aṇ amitrāt ca . \\
\noindent \textbf{(P\_5,4.36)} aṇ amitrāt ca iti vaktavyam . \\
\noindent \textbf{(P\_5,4.36)} maitraḥ . \\
\noindent \textbf{(P\_5,4.36)} āmitraḥ . \\
\noindent \textbf{(P\_5,4.36)} sānnāyyānujāvarānuṣūkacātuṣprāśyarākṣoghnavaiyātavaikṛtavārivaskṛtāgrāyaṇāgrahāyaṇasāntapanāni nipātnyante . \\
\noindent \textbf{(P\_5,4.36)} sānnāyyam . \\
\noindent \textbf{(P\_5,4.36)} ānujāvaraḥ . \\
\noindent \textbf{(P\_5,4.36)} ānuṣūkaḥ . \\
\noindent \textbf{(P\_5,4.36)} cātuṣprāśyaḥ . \\
\noindent \textbf{(P\_5,4.36)} rākṣoghnaḥ . \\
\noindent \textbf{(P\_5,4.36)} vaiyātaḥ . \\
\noindent \textbf{(P\_5,4.36)} vaikṛtaḥ . \\
\noindent \textbf{(P\_5,4.36)} vārivaskṛtaḥ . \\
\noindent \textbf{(P\_5,4.36)} āgrāyaṇaḥ . \\
\noindent \textbf{(P\_5,4.36)} āgrahāyaṇaḥ . \\
\noindent \textbf{(P\_5,4.36)} sāntapanaḥ . \\
\noindent \textbf{(P\_5,4.36)} agnīdhrasādhāraṇāt añ . \\
\noindent \textbf{(P\_5,4.36)} agnīdhrasādhāraṇāt añ vaktavyaḥ . \\
\noindent \textbf{(P\_5,4.36)} āgnīdhram . \\
\noindent \textbf{(P\_5,4.36)} sādhāraṇam . \\
\noindent \textbf{(P\_5,4.36)} ayavasamarudbhyām chandasi . \\
\noindent \textbf{(P\_5,4.36)} ayavasamarudbhyām chandasi añ vaktavyaḥ . \\
\noindent \textbf{(P\_5,4.36)} āyavase vardhante . \\
\noindent \textbf{(P\_5,4.36)} mārutam śardhaḥ . \\
\noindent \textbf{(P\_5,4.36)} navasūramartayaviṣṭhebhyaḥ yat . \\
\noindent \textbf{(P\_5,4.36)} navasūramartayaviṣṭhebhyaḥ yat vaktavyaḥ . \\
\noindent \textbf{(P\_5,4.36)} navyaḥ . \\
\noindent \textbf{(P\_5,4.36)} sūryaḥ . \\
\noindent \textbf{(P\_5,4.36)} martyaḥ . \\
\noindent \textbf{(P\_5,4.36)} yaviṣthyaḥ . \\
\noindent \textbf{(P\_5,4.36)} kṣemāt yaḥ . \\
\noindent \textbf{(P\_5,4.36)} kṣemāt yaḥ vaktavyaḥ . \\
\noindent \textbf{(P\_5,4.36)} kṣemyaḥ tiṣṭhan prataraṇaḥ suvīraḥ . \\
\noindent \textbf{(P\_5,4.42)} bahvalpārthāt maṅgalavacanam [R: maṅgalāmaṅgalavacanam] . bahvalpārthāt maṅgalavacanam [maṅgalāmaṅgalavacanam ] kartavyam . \\
\noindent \textbf{(P\_5,4.42)} bahuśaḥ dehi . \\
\noindent \textbf{(P\_5,4.42)} aniṣṭeṣu śrāddhādiṣu mā bhūt . \\
\noindent \textbf{(P\_5,4.42)} iṣṭeṣu prāśitrādiṣu yathā syāt . \\
\noindent \textbf{(P\_5,4.42)} alpaśaḥ dehi . \\
\noindent \textbf{(P\_5,4.42)} iṣṭeṣu prāśitrādiṣu mā bhūt . \\
\noindent \textbf{(P\_5,4.42)} aniṣṭeṣu śrāddhādiṣu yathā syāt . \\
\noindent \textbf{(P\_5,4.44)} tasiprakaraṇe ādyādibhyaḥ upasaṅkhyānam . \\
\noindent \textbf{(P\_5,4.44)} tasiprakaraṇe ādyādibhyaḥ upasaṅkhyānam kartavyam . \\
\noindent \textbf{(P\_5,4.44)} āditaḥ . \\
\noindent \textbf{(P\_5,4.44)} madhyataḥ . \\
\noindent \textbf{(P\_5,4.44)} antataḥ . \\
\noindent \textbf{(P\_5,4.50)} cvividhau abhūtatadbhāvagrahaṇam . \\
\noindent \textbf{(P\_5,4.50)} cvividhau abhūtatadbhāvagrahaṇam kartavyam . \\
\noindent \textbf{(P\_5,4.50)} iha mā bhūt . \\
\noindent \textbf{(P\_5,4.50)} sampadyante yavāḥ . \\
\noindent \textbf{(P\_5,4.50)} sampadyante śālayaḥ iti . \\
\noindent \textbf{(P\_5,4.50)} atha kriyamāṇe api vā abhūtatadbhāvagrahaṇe iha kasmāt na bhavati . \\
\noindent \textbf{(P\_5,4.50)} sampadyante asmin kṣetra śālayaḥ iti . \\
\noindent \textbf{(P\_5,4.50)} prakṛtivivakṣāgrahaṇam ca . \\
\noindent \textbf{(P\_5,4.50)} prakṛtivivakṣāgrahaṇam ca kartavyam . \\
\noindent \textbf{(P\_5,4.50)} samīpādibhyaḥ upasaṅkhyānam . \\
\noindent \textbf{(P\_5,4.50)} samīpādibhyaḥ upasaṅkhyānam kartavyam . \\
\noindent \textbf{(P\_5,4.50)} samīpī bhavati . \\
\noindent \textbf{(P\_5,4.50)} abhyāśī bhavati . \\
\noindent \textbf{(P\_5,4.50)} antikī bhavati . \\
\noindent \textbf{(P\_5,4.50)} kim punaḥ kāraṇam na sidhyati . \\
\noindent \textbf{(P\_5,4.50)} na hi asamīpam samīpam bhavati . \\
\noindent \textbf{(P\_5,4.50)} kim tarhi . \\
\noindent \textbf{(P\_5,4.50)} asamīpastham samīpastham bhavati . \\
\noindent \textbf{(P\_5,4.50)} tat tarhi vaktavyam . \\
\noindent \textbf{(P\_5,4.50)} na vaktavyam . \\
\noindent \textbf{(P\_5,4.50)} tātsthyāt tācchabdyam bhaviṣyati . \\
\noindent \textbf{(P\_5,4.57)} kimarthaḥ cakāraḥ . \\
\noindent \textbf{(P\_5,4.57)} svarārthaḥ . \\
\noindent \textbf{(P\_5,4.57)} citaḥ antaḥ udāttatḥ bhavati iti udāttatvam yathā syāt . \\
\noindent \textbf{(P\_5,4.57)} na etat asti prayojanam . \\
\noindent \textbf{(P\_5,4.57)} ekāc ayam . \\
\noindent \textbf{(P\_5,4.57)} tatra na arthaḥ svarārthena cakāreṇa anubandhena . \\
\noindent \textbf{(P\_5,4.57)} pratyayasvareṇa eva siddham . \\
\noindent \textbf{(P\_5,4.57)} ataḥ uttaram paṭhati . \\
\noindent \textbf{(P\_5,4.57)} ḍāci citkaraṇam viśeṣaṇārtham . \\
\noindent \textbf{(P\_5,4.57)} ḍāci citkaraṇam kriyate viśeṣaṇārtham . \\
\noindent \textbf{(P\_5,4.57)} kva viśeṣaṇāṛthena arthaḥ . \\
\noindent \textbf{(P\_5,4.57)} lohitādiḍājbhyaḥ kyaṣ iti . \\
\noindent \textbf{(P\_5,4.57)} ḍā iti hi ucyamāne iḍā ataḥ api prasajyeta . \\
\noindent \textbf{(P\_5,4.57)} arthavadgrahaṇe na anarthakasya iti evam etasya na bhaviṣyati . \\
\noindent \textbf{(P\_5,4.57)} iha tarhi prāpnoti . \\
\noindent \textbf{(P\_5,4.57)} nābhā pṛthivyāḥ nihitaḥ davidyutat . \\
\noindent \textbf{(P\_5,4.57)} tasmāt cakāraḥ kartavyaḥ . \\
\noindent \textbf{(P\_5,4.67)} bhadrāt ca iti vaktavyam . \\
\noindent \textbf{(P\_5,4.67)} bhadrā karoti . \\
\noindent \textbf{(P\_5,4.68)} antagrahaṇam kimartham . \\
\noindent \textbf{(P\_5,4.68)} antaḥ yathā syāt . \\
\noindent \textbf{(P\_5,4.68)} na eatat asti prayojanam . \\
\noindent \textbf{(P\_5,4.68)} pratyayparatvena api etat siddham . \\
\noindent \textbf{(P\_5,4.68)} idam tarhi prayojanam . \\
\noindent \textbf{(P\_5,4.68)} tadgrahaṇena grahaṇam yathā syāt . \\
\noindent \textbf{(P\_5,4.68)} kāni punaḥ tadgrahaṇasya prayojanāni . \\
\noindent \textbf{(P\_5,4.68)} prayojanam avyayībhāvadvigudvandvatatpuruṣabahuvrīhisañjñāḥ . \\
\noindent \textbf{(P\_5,4.68)} avyayībhāvaḥ prayojanam . \\
\noindent \textbf{(P\_5,4.68)} pratirājam . \\
\noindent \textbf{(P\_5,4.68)} uparājam . \\
\noindent \textbf{(P\_5,4.68)} avyayībhāvaḥ ca samāsaḥ napuṃsakaliṅgaḥ bhavati iti napuṃsakaliṅgatā yathā syāt . \\
\noindent \textbf{(P\_5,4.68)} na etat asti prayojanam . \\
\noindent \textbf{(P\_5,4.68)} liṅgam aśiṣyam lokāśrayatvāt liṅgasya . \\
\noindent \textbf{(P\_5,4.68)} idam tarhi prayojanam . \\
\noindent \textbf{(P\_5,4.68)} na avyayībhāvāt ataḥ am tu apañcamyāḥ iti eṣaḥ vidhiḥ yathā syāt . \\
\noindent \textbf{(P\_5,4.68)} avyayībhāva . \\
\noindent \textbf{(P\_5,4.68)} dvigu . \\
\noindent \textbf{(P\_5,4.68)} dvigusañjñā ca prayojanam . \\
\noindent \textbf{(P\_5,4.68)} pañcagavam . \\
\noindent \textbf{(P\_5,4.68)} daśagavam . \\
\noindent \textbf{(P\_5,4.68)} dviguḥ ca samāsaḥ napuṃsakaliṅgaḥ bhavati iti napuṃsakaliṅgatā yathā syāt . \\
\noindent \textbf{(P\_5,4.68)} na etat asti prayojanam . \\
\noindent \textbf{(P\_5,4.68)} liṅgam aśiṣyam lokāśrayatvāt liṅgasya . \\
\noindent \textbf{(P\_5,4.68)} idam tarhi prayojanam . \\
\noindent \textbf{(P\_5,4.68)} dvipurī . \\
\noindent \textbf{(P\_5,4.68)} tripurī . \\
\noindent \textbf{(P\_5,4.68)} dvigoḥ akārāntāt iti īkāraḥ yathā syāt . \\
\noindent \textbf{(P\_5,4.68)} etat api na asti prayojanam . \\
\noindent \textbf{(P\_5,4.68)} puraśabdaḥ ayam akārāntaḥ . \\
\noindent \textbf{(P\_5,4.68)} tena samāsaḥ bhaviṣyati . \\
\noindent \textbf{(P\_5,4.68)} ātaḥ ca akārāntaḥ iti āha . \\
\noindent \textbf{(P\_5,4.68)} kṣeme subhikṣe kṛtasañcayāni purāṇi vinayanti kopam iti . \\
\noindent \textbf{(P\_5,4.68)} idam tarhi prayojanam . \\
\noindent \textbf{(P\_5,4.68)} dvidhurī tridhurī . \\
\noindent \textbf{(P\_5,4.68)} dvigoḥ akārāntāt iti ṅīp yathā syāt . \\
\noindent \textbf{(P\_5,4.68)} dvigu . \\
\noindent \textbf{(P\_5,4.68)} dvandva . \\
\noindent \textbf{(P\_5,4.68)} dvandvasañjñā ca prayojanam . \\
\noindent \textbf{(P\_5,4.68)} vāktvacam . \\
\noindent \textbf{(P\_5,4.68)} sraktvacam . \\
\noindent \textbf{(P\_5,4.68)} dvandvaḥ ca samāsaḥ napuṃsakaliṅgaḥ bhavati iti napuṃsakaliṅgatā yathā syāt . \\
\noindent \textbf{(P\_5,4.68)} na etat asti prayojanam . \\
\noindent \textbf{(P\_5,4.68)} liṅgam aśiṣyam lokāśrayatvāt liṅgasya . \\
\noindent \textbf{(P\_5,4.68)} idam tarhi prayojanam . \\
\noindent \textbf{(P\_5,4.68)} idam tarhi prayojanam . \\
\noindent \textbf{(P\_5,4.68)} kośaḥ ca niṣat ca kośaniṣadam . \\
\noindent \textbf{(P\_5,4.68)} kośaniṣadinī . \\
\noindent \textbf{(P\_5,4.68)} dvandvopatāpagarhyāt prāṇisthāt iniḥ iti iniḥ yathā syāt . \\
\noindent \textbf{(P\_5,4.68)} dvandva . \\
\noindent \textbf{(P\_5,4.68)} tatpuruṣa . \\
\noindent \textbf{(P\_5,4.68)} tatpuruṣasañjñā ca prayojanam . \\
\noindent \textbf{(P\_5,4.68)} paramadhurā uttamadhurā . \\
\noindent \textbf{(P\_5,4.68)} paravat liṅgam dvandvatatpuruṣayoḥ iti paravalliṅgatā yathā syāt . \\
\noindent \textbf{(P\_5,4.68)} na etat asti prayjonam . \\
\noindent \textbf{(P\_5,4.68)} uttarapadārthapradhānaḥ tatpuruṣaḥ . \\
\noindent \textbf{(P\_5,4.68)} idam tarhi prayojanam . \\
\noindent \textbf{(P\_5,4.68)} ardhadhurā . \\
\noindent \textbf{(P\_5,4.68)} etat api na asti prayojanam . \\
\noindent \textbf{(P\_5,4.68)} liṅgam aśiṣyam lokāśrayatvāt liṅgasya . \\
\noindent \textbf{(P\_5,4.68)} idam tarhi prayojanam . \\
\noindent \textbf{(P\_5,4.68)} idam tarhi . \\
\noindent \textbf{(P\_5,4.68)} nirdhuraḥ . \\
\noindent \textbf{(P\_5,4.68)} avyayam tatpuruṣe prakṛtisvaram bhavati iti eṣaḥ svaraḥ yathā syāt . \\
\noindent \textbf{(P\_5,4.68)} tatpuruṣa . \\
\noindent \textbf{(P\_5,4.68)} bahuvrīhi . \\
\noindent \textbf{(P\_5,4.68)} bahuvrīhisañjñā ca prayojanam . \\
\noindent \textbf{(P\_5,4.68)} uccadhuraḥ . \\
\noindent \textbf{(P\_5,4.68)} nīcadhuraḥ . \\
\noindent \textbf{(P\_5,4.68)} bahuvrīhau prakṛtyā pūrvapadam bhavati iti eṣaḥ svaraḥ yathā syāt . \\
\noindent \textbf{(P\_5,4.69.1)} idam vipratiṣiddham . \\
\noindent \textbf{(P\_5,4.69.1)} kaḥ pratiṣedhaḥ . \\
\noindent \textbf{(P\_5,4.69.1)} parigaṇitābhyaḥ prakṛtibhyaḥ samāsāntaḥ vidhīyate na ca tatra kā cit pūjanāntā prakṛtiḥ nirdiśyate . \\
\noindent \textbf{(P\_5,4.69.1)} na etat vipratiṣiddham . \\
\noindent \textbf{(P\_5,4.69.1)} na evam vijñāyate . \\
\noindent \textbf{(P\_5,4.69.1)} yābhyaḥ prakṛtibhyaḥ samāsāntaḥ vidhīyate na cet tāḥ pūjanāntāḥ bhavanti iti . \\
\noindent \textbf{(P\_5,4.69.1)} katham tarhi . \\
\noindent \textbf{(P\_5,4.69.1)} na cet tāḥ pūjanāt parāḥ bhavanti iti . \\
\noindent \textbf{(P\_5,4.69.2)} pūjāyām svatigrahaṇam . \\
\noindent \textbf{(P\_5,4.69.2)} pūjāyām svatigrahaṇam kartavyam . \\
\noindent \textbf{(P\_5,4.69.2)} surājā . \\
\noindent \textbf{(P\_5,4.69.2)} atirājā . \\
\noindent \textbf{(P\_5,4.69.2)} kva mā bhūt . \\
\noindent \textbf{(P\_5,4.69.2)} paramagavaḥ . \\
\noindent \textbf{(P\_5,4.69.2)} uttamagavaḥ . \\
\noindent \textbf{(P\_5,4.69.2)} prāgbahuvrīhigrahaṇam ca . \\
\noindent \textbf{(P\_5,4.69.2)} prāgbahuvrīhigrahaṇam ca kartavyam . \\
\noindent \textbf{(P\_5,4.69.2)} iha mā bhūt . \\
\noindent \textbf{(P\_5,4.69.2)} svakṣaḥ . \\
\noindent \textbf{(P\_5,4.69.2)} atyakṣaḥ iti . \\
\noindent \textbf{(P\_5,4.70)} kṣepe iti kimartham . \\
\noindent \textbf{(P\_5,4.70)} kasya rājā kiṃrājā . \\
\noindent \textbf{(P\_5,4.70)} kṣepe iti śakyam akartum . \\
\noindent \textbf{(P\_5,4.70)} kasmāt na bhavati kasya rājā kiṃrājā iti . \\
\noindent \textbf{(P\_5,4.70)} lakṣaṇapratipadoktayoḥ pratipadoktasya eva iti . \\
\noindent \textbf{(P\_5,4.73)} ḍacprakaraṇe saṅkhyāyāḥ tatpuruṣasya upasaṅkhyānam nistriṃśādyartham . \\
\noindent \textbf{(P\_5,4.73)} ḍacprakaraṇe saṅkhyāyāḥ tatpuruṣasya upasaṅkhyānam kartavyam . \\
\noindent \textbf{(P\_5,4.73)} kim prayojanam . \\
\noindent \textbf{(P\_5,4.73)} nistriṃśādyartham . \\
\noindent \textbf{(P\_5,4.73)} nistriṃśāni varṣāṇi . \\
\noindent \textbf{(P\_5,4.73)} niścatvāriṃśāni varṣāṇi . \\
\noindent \textbf{(P\_5,4.73)} anyatra adhikalopāt . \\
\noindent \textbf{(P\_5,4.73)} anyatra adhikalopāt iti vaktavyam . \\
\noindent \textbf{(P\_5,4.73)} iha mā bhūt . \\
\noindent \textbf{(P\_5,4.73)} ekādhikā viṃśatiḥ ekaviṃśatiḥ . \\
\noindent \textbf{(P\_5,4.73)} dvyadhikā viṃśatiḥ dvāviṃśatiḥ . \\
\noindent \textbf{(P\_5,4.73)} avyayādeḥ iti vaktavyam . \\
\noindent \textbf{(P\_5,4.73)} iha mā bhūt . \\
\noindent \textbf{(P\_5,4.73)} gotriṃśat . \\
\noindent \textbf{(P\_5,4.73)} gocatvāriṃśat iti . \\
\noindent \textbf{(P\_5,4.73)} tat tarhi vaktavyam . \\
\noindent \textbf{(P\_5,4.73)} yadi api etat ucyate atha vā etarhi anyatra adhikalopāt iti etat na kriyate . \\
\noindent \textbf{(P\_5,4.74)} anakṣe iti katham idam vijñāyate . \\
\noindent \textbf{(P\_5,4.74)} na cet akṣadhūrantaḥ samāsaḥ iti āhosvit na cet akṣaḥ samāsārthaḥ iti . \\
\noindent \textbf{(P\_5,4.74)} kim ca ataḥ . \\
\noindent \textbf{(P\_5,4.74)} yadi vijñāyate na cet akṣadhūrantaḥ samāsaḥ iti siddham akṣasya dhūḥ akṣadhūḥ iti . \\
\noindent \textbf{(P\_5,4.74)} idam tu na sidhyati . \\
\noindent \textbf{(P\_5,4.74)} dṛḍhadhūḥ ayam akṣaḥ . \\
\noindent \textbf{(P\_5,4.74)} astu tarhi na cet akṣaḥ samāsārthaḥ iti . \\
\noindent \textbf{(P\_5,4.74)} siddham dṛḍhadhūḥ akṣaḥ iti . \\
\noindent \textbf{(P\_5,4.74)} idam tu na sidhyati . \\
\noindent \textbf{(P\_5,4.74)} akṣasya dhūḥ akṣadhūḥ iti . \\
\noindent \textbf{(P\_5,4.74)} evam tarhi na evam vijñāyate na cet akṣadhūrantaḥ samāsaḥ iti na api na cet akṣaḥ samāsārthaḥ iti . \\
\noindent \textbf{(P\_5,4.74)} katham tarhi . \\
\noindent \textbf{(P\_5,4.74)} na cet akṣasya dhūḥ iti . \\
\noindent \textbf{(P\_5,4.74)} evam ca kṛtvā na api na cet akṣadhūrantaḥ samāsaḥ iti vijñāyate na api na cet akṣaḥ samāsārthaḥ iti . \\
\noindent \textbf{(P\_5,4.74)} atha ca ubhayoḥ na bhavati . \\
\noindent \textbf{(P\_5,4.76)} adarśanāt iti ucyate . \\
\noindent \textbf{(P\_5,4.76)} tatra idam na sidhyati . \\
\noindent \textbf{(P\_5,4.76)} kavarākṣam . \\
\noindent \textbf{(P\_5,4.76)} adarśanāt iti śakyam akartum . \\
\noindent \textbf{(P\_5,4.76)} katham brāhmaṇakṣi kṣatriyākṣi . \\
\noindent \textbf{(P\_5,4.76)} aprāṇyaṅgāt iti vaktavyam . \\
\noindent \textbf{(P\_5,4.77)} ādyāḥ trayaḥ bahuvrīhayaḥ . \\
\noindent \textbf{(P\_5,4.77)} adraṣṭā caturṇām acaturaḥ . \\
\noindent \textbf{(P\_5,4.77)} vidraṣṭā caturṇām vicaturaḥ . \\
\noindent \textbf{(P\_5,4.77)} sudraṣṭā caturṇām sucaturaḥ . \\
\noindent \textbf{(P\_5,4.77)} tataḥ pare ekādaśa dvandvāḥ . \\
\noindent \textbf{(P\_5,4.77)} strīpuṃsa dhenvanaḍuha ṛksāma vāṅmanasa akṣibhruva dāragava ūrvaṣṭhīva padaṣṭhīva naktandiva rātrindiva ahardiva . \\
\noindent \textbf{(P\_5,4.77)} tataḥ avyayībhāvaḥ . \\
\noindent \textbf{(P\_5,4.77)} saha rajasā sarajasam . \\
\noindent \textbf{(P\_5,4.77)} tataḥ tatpuruṣaḥ . \\
\noindent \textbf{(P\_5,4.77)} niśritam śreyaḥ niḥśreyasam . \\
\noindent \textbf{(P\_5,4.77)} tataḥ ṣaṣṭhīsamāsaḥ . \\
\noindent \textbf{(P\_5,4.77)} puruṣasya āyuḥ puruṣāyuṣam . \\
\noindent \textbf{(P\_5,4.77)} tataḥ dvigū . \\
\noindent \textbf{(P\_5,4.77)} dve āyuṣī dvyāyuṣam . \\
\noindent \textbf{(P\_5,4.77)} trīṇi āyūṃṣi tryāyuṣam . \\
\noindent \textbf{(P\_5,4.77)} tataḥ dvandvaḥ . \\
\noindent \textbf{(P\_5,4.77)} ṛk ca yajuḥ ca rgyajuṣam . \\
\noindent \textbf{(P\_5,4.77)} jātādayaḥ ukṣāntāḥ samānādhikaraṇāḥ . \\
\noindent \textbf{(P\_5,4.77)} jātaḥ ukṣā jātokṣaḥ . \\
\noindent \textbf{(P\_5,4.77)} mahān ukṣā mahokṣaḥ . \\
\noindent \textbf{(P\_5,4.77)} vṛddhaḥ ukṣā vṛddhokṣaḥ . \\
\noindent \textbf{(P\_5,4.77)} tataḥ avyayībhāvaḥ . \\
\noindent \textbf{(P\_5,4.77)} śunaḥ samīpam upaśunam . \\
\noindent \textbf{(P\_5,4.77)} tataḥ saptamīsamāsaḥ goṣṭhe śvā goṣṭhaśvaḥ . \\
\noindent \textbf{(P\_5,4.77)} caturaḥ acprakaraṇe tryupābhyām upasaṅkhyānam . \\
\noindent \textbf{(P\_5,4.77)} caturaḥ acprakaraṇe tryupābhyām upasaṅkhyānam kartavyam . \\
\noindent \textbf{(P\_5,4.77)} tricaturāḥ . \\
\noindent \textbf{(P\_5,4.77)} upacaturāḥ . \\
\noindent \textbf{(P\_5,4.78)} palyarājabhyām ca iti vaktavyam . \\
\noindent \textbf{(P\_5,4.78)} palyavarcasam . \\
\noindent \textbf{(P\_5,4.78)} rājavarcasam . \\
\noindent \textbf{(P\_5,4.87)} ahargrahaṇam dvandvārtham . \\
\noindent \textbf{(P\_5,4.87)} ahargrahaṇam dvandvārtham draṣṭavyam . \\
\noindent \textbf{(P\_5,4.87)} kim ucyate dvandvārtham iti na punaḥ tatpuruṣārtham api syāt . \\
\noindent \textbf{(P\_5,4.87)} tatpuruṣābhāvāt . \\
\noindent \textbf{(P\_5,4.87)} na hi rātryantaḥ aharādiḥ tatpuruṣaḥ asti . \\
\noindent \textbf{(P\_5,4.88)} ahnaḥ ahnavacanānarthakyam ca ahnaḥ ṭakhoḥ niyamavacanāt . \\
\noindent \textbf{(P\_5,4.88)} ahnaḥ ahnavacanam anarthakam . \\
\noindent \textbf{(P\_5,4.88)} kim kāraṇam . \\
\noindent \textbf{(P\_5,4.88)} ahnaḥ ṭakhoḥ niyamavacanāt . \\
\noindent \textbf{(P\_5,4.88)} ahnaḥ ṭakhoḥ eva iti etat niyamārtham bhaviṣyati . \\
\noindent \textbf{(P\_5,4.103)} anasantāt napuṃsakāt chandasi vā iti vaktavyam . \\
\noindent \textbf{(P\_5,4.103)} brahmasāmam . \\
\noindent \textbf{(P\_5,4.103)} brahmasāma . \\
\noindent \textbf{(P\_5,4.103)} devacchandasam . \\
\noindent \textbf{(P\_5,4.103)} devacchandaḥ . \\
\noindent \textbf{(P\_5,4.113)} kimartham ṣac pratyayāntaram vidhīyate na ṭac prakṛtaḥ saḥ anuvartiṣyate . \\
\noindent \textbf{(P\_5,4.113)} ataḥ uttaram paṭhati . \\
\noindent \textbf{(P\_5,4.113)} ṣaci pratyayāntarakaraṇam anantodāttārtham . \\
\noindent \textbf{(P\_5,4.113)} ṣaci pratyayāntaram kriyate . \\
\noindent \textbf{(P\_5,4.113)} kim prayojanam . \\
\noindent \textbf{(P\_5,4.113)} anantodāttārtham . \\
\noindent \textbf{(P\_5,4.113)} anantodāttāḥ prayojayanti . \\
\noindent \textbf{(P\_5,4.113)} cakrasaktham . \\
\noindent \textbf{(P\_5,4.113)} cakrasakthī . \\
\noindent \textbf{(P\_5,4.115)} kimartham mūrdhnaḥ ṣa pratyayāntaram vidhīyate na ṣac prakṛtaḥ saḥ anuvartiṣyate . \\
\noindent \textbf{(P\_5,4.115)} mūrdhnaḥ ca ṣavacanam . \\
\noindent \textbf{(P\_5,4.115)} kim . \\
\noindent \textbf{(P\_5,4.115)} anantodāttārtham iti eva . \\
\noindent \textbf{(P\_5,4.115)} dvimūrdhaḥ . \\
\noindent \textbf{(P\_5,4.115)} trimūrdhaḥ . \\
\noindent \textbf{(P\_5,4.116)} api pradhānapūraṇīgrahaṇam . \\
\noindent \textbf{(P\_5,4.116)} api pradhānapūraṇīgrahaṇam kartavyam . \\
\noindent \textbf{(P\_5,4.116)} pradhānam yā pūraṇī iti vaktavyam . \\
\noindent \textbf{(P\_5,4.116)} iha mā bhūt . \\
\noindent \textbf{(P\_5,4.116)} kalyāṇī pañcamī asya pakṣasya kalyāṇapañcamīkaḥ pakṣaḥ . \\
\noindent \textbf{(P\_5,4.116)} atha iha katham bhavitavyam . \\
\noindent \textbf{(P\_5,4.116)} kalyāṇī pañcamī asām rātrīṇām iti . \\
\noindent \textbf{(P\_5,4.116)} kalyāṇīpañcamāḥ rātrayaḥ iti bhavitavyam . \\
\noindent \textbf{(P\_5,4.116)} ratrayaḥ atra pradhānam . \\
\noindent \textbf{(P\_5,4.116)} netuḥ nakṣatre upasaṅkhyānam . \\
\noindent \textbf{(P\_5,4.116)} netuḥ nakṣatre upasaṅkhyānam kartavyam . \\
\noindent \textbf{(P\_5,4.116)} puṣyanetrāḥ . \\
\noindent \textbf{(P\_5,4.116)} mṛganetrāḥ . \\
\noindent \textbf{(P\_5,4.116)} chandasi ca . \\
\noindent \textbf{(P\_5,4.116)} chandasi ca netuḥ upasaṅkhyānam kartavyam . \\
\noindent \textbf{(P\_5,4.116)} bṛhaspatinetrāḥ . \\
\noindent \textbf{(P\_5,4.116)} somanetrāḥ . \\
\noindent \textbf{(P\_5,4.116)} māsāt bhṛtipratyayapūrvapadāt ṭhajvidhiḥ . \\
\noindent \textbf{(P\_5,4.116)} māsāt bhṛtipratyayapūrvapadāt ṭhac vidheyaḥ . \\
\noindent \textbf{(P\_5,4.116)} pañcakamāsikaḥ . \\
\noindent \textbf{(P\_5,4.116)} ṣaṭkamāsikaḥ . \\
\noindent \textbf{(P\_5,4.116)} daśakamāsikaḥ . \\
\noindent \textbf{(P\_5,4.118)} kharakhurābhyām ca nas vaktavyaḥ . \\
\noindent \textbf{(P\_5,4.118)} kharaṇāḥ . \\
\noindent \textbf{(P\_5,4.118)} khuraṇāḥ . \\
\noindent \textbf{(P\_5,4.118)} śitināḥ arcanāḥ ahināḥ iti naigamāḥ . \\
\noindent \textbf{(P\_5,4.118)} śitināḥ arcanāḥ ahināḥ . \\
\noindent \textbf{(P\_5,4.119)} veḥ graḥ vaktavyaḥ . \\
\noindent \textbf{(P\_5,4.119)} vigraḥ . \\
\noindent \textbf{(P\_5,4.131)} ūdhasaḥ anaṅi strīgrahaṇam . \\
\noindent \textbf{(P\_5,4.131)} ūdhasaḥ anaṅi strīgrahaṇam kartavyam iha mā bhūt . \\
\noindent \textbf{(P\_5,4.131)} mahodhāḥ parjanyaḥ iti . \\
\noindent \textbf{(P\_5,4.135)} gandhasya ittve tadekāntagrahaṇam . \\
\noindent \textbf{(P\_5,4.135)} gandhasya ittve tadekāntagrahaṇam kartavyam iha mā bhūt . \\
\noindent \textbf{(P\_5,4.135)} śobhanāḥ gandhāḥ asya sugandhaḥ āpaṇikaḥ iti . \\
\noindent \textbf{(P\_5,4.135)} atha anulipte katham bhavitavyam . \\
\noindent \textbf{(P\_5,4.135)} yadi tāvat yat anugatam tat abhisamīkṣitam sugandhiḥ iti bhavitavyam . \\
\noindent \textbf{(P\_5,4.135)} atha yat praviśīrṇam sugandhaḥ iti bhavitavyam . \\
\noindent \textbf{(P\_5,4.154)} śeṣāt iti ucyate . \\
\noindent \textbf{(P\_5,4.154)} kaḥ śeṣaḥ nāma . \\
\noindent \textbf{(P\_5,4.154)} yābhyaḥ prakṛtibhyaḥ samāsāntaḥ ni vidhīyate saḥ śeṣaḥ . \\
\noindent \textbf{(P\_5,4.154)} kimartham punaḥ śeṣagrahaṇam kriyate . \\
\noindent \textbf{(P\_5,4.154)} yābhyaḥ prakṛtibhyaḥ samāsāntaḥ vidhīyate tābhyaḥ mā bhūt iti . \\
\noindent \textbf{(P\_5,4.154)} na etat asti prayojanam . \\
\noindent \textbf{(P\_5,4.154)} ye pratipadam vidhīyante te tatra bādhakāḥ bhaviṣyanti . \\
\noindent \textbf{(P\_5,4.154)} anavakāśāḥ hi vidhayaḥ bādhakāḥ bhavanti sāvakāśāḥ ca samāsāntāḥ . \\
\noindent \textbf{(P\_5,4.154)} kaḥ avakāśaḥ . \\
\noindent \textbf{(P\_5,4.154)} vibhāṣā kap . \\
\noindent \textbf{(P\_5,4.154)} yadā na kap saḥ avakāśaḥ . \\
\noindent \textbf{(P\_5,4.154)} kapaḥ prasaṅge ubhayam prāpnoti . \\
\noindent \textbf{(P\_5,4.154)} paratvāt kap prāpnoti . \\
\noindent \textbf{(P\_5,4.154)} tasmāt śeṣagrahaṇam kartavyam . \\
\noindent \textbf{(P\_5,4.154)} kim punaḥ idam śeṣagrahaṇam kabapekṣam yasmāt bahuvrīheḥ kap iti āhosvit samāsāntāpekṣam yasmāt bahuvrīheḥ samāsāntaḥ na vihitaḥ iti . \\
\noindent \textbf{(P\_5,4.154)} kim ca ataḥ . \\
\noindent \textbf{(P\_5,4.154)} yadi vijñāyate kabapekṣam anṛcaḥ bahvṛcaḥ iti atra api prāpnoti . \\
\noindent \textbf{(P\_5,4.154)} atha samāsāntāpekṣam anṛkkam bahvṛkkam sūktam iti na sidhyati . \\
\noindent \textbf{(P\_5,4.154)} astu kabapekṣam . \\
\noindent \textbf{(P\_5,4.154)} katham anṛcaḥ bahvṛcaḥ iti . \\
\noindent \textbf{(P\_5,4.154)} viśeṣe etat vaktavyam . \\
\noindent \textbf{(P\_5,4.154)} anṛcaḥ māṇave bahvṛcaḥ caraṇaśākhāyām iti . \\
\noindent \textbf{(P\_5,4.154)} idam tarhi ūdhasaḥ anaṅi strīgrahaṇam coditam . \\
\noindent \textbf{(P\_5,4.154)} tasmin kriyamāṇe api prāpnoti . \\
\noindent \textbf{(P\_5,4.154)} evam tarhi na eva kabapekṣam śeṣagrahaṇam na api samāsāntāpekṣam . \\
\noindent \textbf{(P\_5,4.154)} kim tarhi anantaraḥ yaḥ bahuvrīhyadhikāraḥ saḥ apekṣyate . \\
\noindent \textbf{(P\_5,4.154)} anantare bahuvrīhyadhikāre yasmāt bahuvrīheḥ samāsāntaḥ na vihitaḥ iti . \\
\noindent \textbf{(P\_5,4.154)} katham anṛcaḥ bahvṛcaḥ iti . \\
\noindent \textbf{(P\_5,4.154)} vaktavyam eva anṛcaḥ māṇave bahvṛcaḥ caraṇaśākhāyām iti . \\
\noindent \textbf{(P\_5,4.156)} īyasaḥ upasarjanadīrghatvam ca . \\
\noindent \textbf{(P\_5,4.156)} īyasaḥ upasarjanadīrghatvam ca vaktavyam . \\
\noindent \textbf{(P\_5,4.156)} bahvyaḥ śreyasyaḥ asya bahuśreyasī . \\
\noindent \textbf{(P\_5,4.156)} vidyamānaśreyasī . \\
\noindent \textbf{(P\_5,4.156)} puṃvadvacanāt siddham . \\
\noindent \textbf{(P\_5,4.156)} puṃvadbhāvaḥ atra bhavati īyasaḥ bahuvrīhau puṃvadvacanam iti . \\
\noindent \textbf{(P\_6,1.1.1)} ekācaḥ iti kim ayam bahuvrīhiḥ . \\
\noindent \textbf{(P\_6,1.1.1)} ekaḥ ac asmin saḥ ekāc . \\
\noindent \textbf{(P\_6,1.1.1)} ekācaḥ iti . \\
\noindent \textbf{(P\_6,1.1.1)} āhosvit tatpuruṣaḥ ayam samānādhikaraṇaḥ . \\
\noindent \textbf{(P\_6,1.1.1)} ekaḥ ac ekāc . \\
\noindent \textbf{(P\_6,1.1.1)} ekācaḥ . \\
\noindent \textbf{(P\_6,1.1.1)} kim ca ataḥ . \\
\noindent \textbf{(P\_6,1.1.1)} yadi bahuvrīhiḥ siddham papāca papāṭha . \\
\noindent \textbf{(P\_6,1.1.1)} iyāya . \\
\noindent \textbf{(P\_6,1.1.1)} āra iti na sidhyati . \\
\noindent \textbf{(P\_6,1.1.1)} atha tatpuruṣaḥ samānādhikaraṇaḥ siddham iyāya . \\
\noindent \textbf{(P\_6,1.1.1)} āra iti . \\
\noindent \textbf{(P\_6,1.1.1)} papāca papāṭha iti na sidhyati . \\
\noindent \textbf{(P\_6,1.1.1)} ataḥ uttaram paṭhati . \\
\noindent \textbf{(P\_6,1.1.1)} ekācaḥ dve prathamasya iti bahuvrīhinirdeśaḥ . \\
\noindent \textbf{(P\_6,1.1.1)} ekācaḥ dve prathamasya iti bahuvrīhinirdeśaḥ ayam . \\
\noindent \textbf{(P\_6,1.1.1)} ekavarṇeṣu katham . \\
\noindent \textbf{(P\_6,1.1.1)} ekavarṇeṣu vyapadeśivadvacanāt . \\
\noindent \textbf{(P\_6,1.1.1)} vyapadeśivat ekasmin kāryam bhavati iti vaktavyam . \\
\noindent \textbf{(P\_6,1.1.1)} evam ekavarṇeṣu dvirvacanam bhaviṣyati . \\
\noindent \textbf{(P\_6,1.1.1)} ekācaḥ dve bhavataḥ iti ucyate . \\
\noindent \textbf{(P\_6,1.1.1)} tatra na jñāyate kasya ekācaḥ dve bhavataḥ iti . \\
\noindent \textbf{(P\_6,1.1.1)} vakṣyati liṭi dhātoḥ anabhyāsasya iti . \\
\noindent \textbf{(P\_6,1.1.1)} tena dhātoḥ ekācaḥ iti vijñāyate . \\
\noindent \textbf{(P\_6,1.1.1)} yadi dhātoḥ ekācaḥ siddham papāca papāṭha . \\
\noindent \textbf{(P\_6,1.1.1)} jajāgāra puputrīyiṣati iti na sidhyati . \\
\noindent \textbf{(P\_6,1.1.1)} dhātoḥ iti na eṣā ekācsamānādhikaraṇā ṣaṣṭhī . \\
\noindent \textbf{(P\_6,1.1.1)} dhātoḥ ekācaḥ iti . \\
\noindent \textbf{(P\_6,1.1.1)} kim tarhi avayavayogā eṣā ṣaṣṭhī . \\
\noindent \textbf{(P\_6,1.1.1)} dhātoḥ yaḥ ekāc avayavaḥ iti . \\
\noindent \textbf{(P\_6,1.1.1)} avayavayogā eṣā ṣaṣṭhī iti cet siddham jajāgāra puputrīyiṣati iti . \\
\noindent \textbf{(P\_6,1.1.1)} papāca papāṭha iti na sidhyati . \\
\noindent \textbf{(P\_6,1.1.1)} eṣaḥ api vyapadeśivadbhāvena dhātoḥ ekāc avayavaḥ bhavati . \\
\noindent \textbf{(P\_6,1.1.1)} ekācaḥ dve prathamasya iti ucyate . \\
\noindent \textbf{(P\_6,1.1.1)} tena yatra eva prathamaḥ ca aprathamaḥ ca tatra dvirvacanam syāt . \\
\noindent \textbf{(P\_6,1.1.1)} jajāgāra puputrīyiṣati iti . \\
\noindent \textbf{(P\_6,1.1.1)} papāca papāṭha iti atra na syāt . \\
\noindent \textbf{(P\_6,1.1.1)} prathamatve ca . \\
\noindent \textbf{(P\_6,1.1.1)} prathamatve ca kim . \\
\noindent \textbf{(P\_6,1.1.1)} vyapadeśivadvacanāt siddham iti eva . \\
\noindent \textbf{(P\_6,1.1.1)} saḥ tarhi vyapadeśivadbhāvaḥ vaktavyaḥ . \\
\noindent \textbf{(P\_6,1.1.1)} na vaktavyaḥ . \\
\noindent \textbf{(P\_6,1.1.1)} uktam vā . \\
\noindent \textbf{(P\_6,1.1.1)} kim uktam . \\
\noindent \textbf{(P\_6,1.1.1)} tatra vyapadeśivadvacanam . \\
\noindent \textbf{(P\_6,1.1.1)} ekācaḥ dve prathamārtham . \\
\noindent \textbf{(P\_6,1.1.1)} ṣatve ca ādeśasampratyayārtham . \\
\noindent \textbf{(P\_6,1.1.1)} avacanāt lokavijñānāt siddham iti eva . \\
\noindent \textbf{(P\_6,1.1.1)} yogavibhāgaḥ vā . \\
\noindent \textbf{(P\_6,1.1.1)} atha vā yogavibhāgaḥ kariṣyate . \\
\noindent \textbf{(P\_6,1.1.1)} ekācaḥ dve bhavataḥ . \\
\noindent \textbf{(P\_6,1.1.1)} kimarthaḥ yogavibhāgaḥ . \\
\noindent \textbf{(P\_6,1.1.1)} ekājmātrasya dvirvacanārthaḥ . \\
\noindent \textbf{(P\_6,1.1.1)} ekājmātrayta dvirvacanam yathā syāt . \\
\noindent \textbf{(P\_6,1.1.1)} iyāya papāca . \\
\noindent \textbf{(P\_6,1.1.1)} tataḥ prathamasya . \\
\noindent \textbf{(P\_6,1.1.1)} prathamasya ekācaḥ dve bhavataḥ . \\
\noindent \textbf{(P\_6,1.1.1)} idam idānīm kimartham . \\
\noindent \textbf{(P\_6,1.1.1)} niyamāṛtham . \\
\noindent \textbf{(P\_6,1.1.1)} yatra prathamaḥ ca aprathamaḥ ca asti tatra prathamasya ekācaḥ dvirvacanam yathā syāt . \\
\noindent \textbf{(P\_6,1.1.1)} aprathamasya mā bhūt . \\
\noindent \textbf{(P\_6,1.1.1)} jajāgāra puputrīyiṣati iti . \\
\noindent \textbf{(P\_6,1.1.1)} ekācaḥ avayavaikāctvāt avayavānām dvirvacanaprasaṅgaḥ . \\
\noindent \textbf{(P\_6,1.1.1)} ekācaḥ avayavaikāctvāt avayavānām dvirvacanam prāpnoti . \\
\noindent \textbf{(P\_6,1.1.1)} nenijati iti atra nijśabdaḥ api ekāc ijśabdaḥ api ekāc ikāraḥ api ekāc niśabdaḥ api . \\
\noindent \textbf{(P\_6,1.1.1)} tatra nijśabdasya dvirvacane rūpam siddham doṣāḥ ca na santi . \\
\noindent \textbf{(P\_6,1.1.1)} ijśabdasya dvirvacane rūpam na sidhyati doṣāḥ ca na santi . \\
\noindent \textbf{(P\_6,1.1.1)} ikārasya dvirvacane rūpam na sidhati doṣāḥ ca na santi . \\
\noindent \textbf{(P\_6,1.1.1)} niśabdasya dvirvacane rūpam siddham doṣāḥ tu santi . \\
\noindent \textbf{(P\_6,1.1.1)} tatra kaḥ doṣaḥ . \\
\noindent \textbf{(P\_6,1.1.1)} tatra jusbhāvavacanam . \\
\noindent \textbf{(P\_6,1.1.1)} tatra jusbhāvaḥ vaktavyaḥ . \\
\noindent \textbf{(P\_6,1.1.1)} anenijuḥ . \\
\noindent \textbf{(P\_6,1.1.1)} paryaveviṣuḥ . \\
\noindent \textbf{(P\_6,1.1.1)} abhyastāt jheḥ jusbhāvaḥ bhavati iti jusbhāvaḥ na prāpnoti jakāreṇavyavadhānāt . \\
\noindent \textbf{(P\_6,1.1.1)} svaraḥ ca . \\
\noindent \textbf{(P\_6,1.1.1)} svaraḥ ca na sidhyati . \\
\noindent \textbf{(P\_6,1.1.1)} nenijati . \\
\noindent \textbf{(P\_6,1.1.1)} yat pariveviṣati iti . \\
\noindent \textbf{(P\_6,1.1.1)} abhyastānām ādiḥ udāttaḥ bhavati ajādau lasārvadhātuke iti eṣaḥ svaraḥ na prāpnoti . \\
\noindent \textbf{(P\_6,1.1.1)} adbhāvaḥ ca . \\
\noindent \textbf{(P\_6,1.1.1)} adbhāvaḥ ca na sidhyati . \\
\noindent \textbf{(P\_6,1.1.1)} nenijati . \\
\noindent \textbf{(P\_6,1.1.1)} pariveviṣati iti . \\
\noindent \textbf{(P\_6,1.1.1)} abhyastāt iti adbhāvaḥ na prāpnoti . \\
\noindent \textbf{(P\_6,1.1.1)} numpratiṣedhaḥ ca . \\
\noindent \textbf{(P\_6,1.1.1)} numpratiṣedhaḥ ca na sidhyati . \\
\noindent \textbf{(P\_6,1.1.1)} nenijat . \\
\noindent \textbf{(P\_6,1.1.1)} pariveviṣat . \\
\noindent \textbf{(P\_6,1.1.1)} na abhyāstāt śatuḥ it numpratiṣedhaḥ na prāpnoti jakāreṇavyavadhānāt . \\
\noindent \textbf{(P\_6,1.1.1)} śāstrahāniḥ ca . \\
\noindent \textbf{(P\_6,1.1.1)} śāstrahāniḥ ca bhavati . \\
\noindent \textbf{(P\_6,1.1.1)} samudāyaikācaḥ śāstram hīyate . \\
\noindent \textbf{(P\_6,1.1.1)} siddham tu tatsamudāyaikāctvāt śāstrāhāneḥ . \\
\noindent \textbf{(P\_6,1.1.1)} siddham etat . \\
\noindent \textbf{(P\_6,1.1.1)} katham . \\
\noindent \textbf{(P\_6,1.1.1)} tatsamudāyaikāctvāt . \\
\noindent \textbf{(P\_6,1.1.1)} kim idam tatsamudāyaikāctvāt iti . \\
\noindent \textbf{(P\_6,1.1.1)} tasya samudāyaḥ tatsamudāyaḥ . \\
\noindent \textbf{(P\_6,1.1.1)} ekājbhāvaḥ ekāctvam . \\
\noindent \textbf{(P\_6,1.1.1)} tatsamudāyasya ekāctvam tatsamudāyaikāctvam . \\
\noindent \textbf{(P\_6,1.1.1)} tatsamudāyaikāctvāt . \\
\noindent \textbf{(P\_6,1.1.1)} tatsamudāyaikācaḥ dvirvacanam bhaviṣyati . \\
\noindent \textbf{(P\_6,1.1.1)} kutaḥ . \\
\noindent \textbf{(P\_6,1.1.1)} śāstrāhāneḥ . \\
\noindent \textbf{(P\_6,1.1.1)} evam hi śāstram ahīnam bhavati . \\
\noindent \textbf{(P\_6,1.1.1)} nanu ca samudāyaikācaḥ dvirvacane kriyamāṇe api avayavaikācaḥ śāstram hīyate . \\
\noindent \textbf{(P\_6,1.1.1)} na hīyate . \\
\noindent \textbf{(P\_6,1.1.1)} kim kāraṇam . \\
\noindent \textbf{(P\_6,1.1.1)} avayavātmakatvāt samudāyasya . \\
\noindent \textbf{(P\_6,1.1.1)} avayavātmakaḥ samudāyaḥ . \\
\noindent \textbf{(P\_6,1.1.1)} abhyantaraḥ hi samudāye avayavaḥ . \\
\noindent \textbf{(P\_6,1.1.1)} tat yathā vṛkṣaḥ pracalan sahāvayavaiḥ pracalati . \\
\noindent \textbf{(P\_6,1.1.1)} tatra bahuvrīhinirdeśe anackasya dvirvacanam anyapadārthatvāt . \\
\noindent \textbf{(P\_6,1.1.1)} tatra bahuvrīhinirdeśe anackasya dvirvacanam prāpnoti . \\
\noindent \textbf{(P\_6,1.1.1)} āṭatuḥ . \\
\noindent \textbf{(P\_6,1.1.1)} āṭuḥ . \\
\noindent \textbf{(P\_6,1.1.1)} kim kāraṇam . \\
\noindent \textbf{(P\_6,1.1.1)} anyapadārthatvāt bahuvrīheḥ . \\
\noindent \textbf{(P\_6,1.1.1)} anyapadārthe bahuvrīhiḥ vartate . \\
\noindent \textbf{(P\_6,1.1.1)} tena yat anyat acaḥ tasya dvirvacanam syāt . \\
\noindent \textbf{(P\_6,1.1.1)} tat yathā citraguḥ ānīyatām iti ukte yasya tāḥ gāvaḥ santi saḥ ānīyate na gāvaḥ . \\
\noindent \textbf{(P\_6,1.1.1)} siddham tu tadguṇasaṃvijñānāt pāṇineḥ yathā loke . \\
\noindent \textbf{(P\_6,1.1.1)} siddham etat . \\
\noindent \textbf{(P\_6,1.1.1)} katham . \\
\noindent \textbf{(P\_6,1.1.1)} tadguṇasaṃvijñānāt bhagavataḥ pāṇineḥ ācāryasya yathā loke . \\
\noindent \textbf{(P\_6,1.1.1)} loke śuklavāsasam ānaya . \\
\noindent \textbf{(P\_6,1.1.1)} lohitoṣṇīṣāḥ pracaranti iti . \\
\noindent \textbf{(P\_6,1.1.1)} tadguṇaḥ ānīyate tadguṇāḥ ca pracaranti . \\
\noindent \textbf{(P\_6,1.1.1)} evam iha api . \\
\noindent \textbf{(P\_6,1.1.2)} atha yasya dvirvacanam ārabhyate kim tasya sthāne bhavati āhosvit dviḥprayogaḥ iti . \\
\noindent \textbf{(P\_6,1.1.2)} kaḥ ca atra viśeṣaḥ . \\
\noindent \textbf{(P\_6,1.1.2)} sthāne dvirvacane ṇilopavacanam samudāyādeśatvāt . \\
\noindent \textbf{(P\_6,1.1.2)} sthāne dvirvacane ṇilopaḥ vaktavyaḥ . \\
\noindent \textbf{(P\_6,1.1.2)} āṭiṭat . \\
\noindent \textbf{(P\_6,1.1.2)} āśiśat . \\
\noindent \textbf{(P\_6,1.1.2)} kim kāraṇam . \\
\noindent \textbf{(P\_6,1.1.2)} samudāyādeśatvāt . \\
\noindent \textbf{(P\_6,1.1.2)} samudāyasya samudāyaḥ ādeśaḥ . \\
\noindent \textbf{(P\_6,1.1.2)} tatra sampramugdhatvāt prakṛtipratyayasamudāyasya naṣṭaḥ ṇiḥ bhavati iti ṇeḥ aniṭi iti ṇilopaḥ na prāpnoti . \\
\noindent \textbf{(P\_6,1.1.2)} idam iha sampradhāryam . \\
\noindent \textbf{(P\_6,1.1.2)} dvirvacanam kriyatām ṇilopaḥ iti kim atra kartavyam . \\
\noindent \textbf{(P\_6,1.1.2)} paratvāt ṇilopaḥ . \\
\noindent \textbf{(P\_6,1.1.2)} nityam dvirvacanam . \\
\noindent \textbf{(P\_6,1.1.2)} kṛte api ṇilope prāpnoti akṛte api . \\
\noindent \textbf{(P\_6,1.1.2)} dvirvacanam api nityam . \\
\noindent \textbf{(P\_6,1.1.2)} anyasya kṛte ṇilope prāpnoti anyasya akṛte . \\
\noindent \textbf{(P\_6,1.1.2)} śabdāntarasya ca prāpnuvan vidhiḥ anityaḥ bhavati . \\
\noindent \textbf{(P\_6,1.1.2)} nityam eva dvirvacanam . \\
\noindent \textbf{(P\_6,1.1.2)} katham . \\
\noindent \textbf{(P\_6,1.1.2)} rūpasya sthānivatvāt . \\
\noindent \textbf{(P\_6,1.1.2)} yat ca sanyaṅantasya dvirvacane . \\
\noindent \textbf{(P\_6,1.1.2)} yat ca sanyaṅantasya dvirvacane codyam tat iha api codyam . \\
\noindent \textbf{(P\_6,1.1.2)} kim punaḥ tat . \\
\noindent \textbf{(P\_6,1.1.2)} sanyaṅantasya cet aśeḥ sani aniṭaḥ . \\
\noindent \textbf{(P\_6,1.1.2)} dīrghakutvaprasāraṇaṣatvam adhikasya dvirvacanāt . \\
\noindent \textbf{(P\_6,1.1.2)} ābṛdhyoḥ ca abhyastavipratiṣedhaḥ . \\
\noindent \textbf{(P\_6,1.1.2)} saṅāśraye ca samudāyasya samudāyādeśatvāt jhalāśraye ca avyapadeśaḥ āmiśratvāt iti . \\
\noindent \textbf{(P\_6,1.1.2)} astu tarhi dviḥprayogaḥ dvirvacanam . \\
\noindent \textbf{(P\_6,1.1.2)} dviḥprayogaḥ iti cet ṇakāraṣakārādeśādeḥ ettvavacanam liṭi . \\
\noindent \textbf{(P\_6,1.1.2)} dviḥprayogaḥ iti cet ṇakāraṣakārādeśādeḥ ettvam liṭi vaktavyam . \\
\noindent \textbf{(P\_6,1.1.2)} nematuḥ . \\
\noindent \textbf{(P\_6,1.1.2)} nemuḥ . \\
\noindent \textbf{(P\_6,1.1.2)} sehe . \\
\noindent \textbf{(P\_6,1.1.2)} sehāte . \\
\noindent \textbf{(P\_6,1.1.2)} sahire . \\
\noindent \textbf{(P\_6,1.1.2)} anādeśādeḥ iti pratiṣedhaḥ prāpnoti . \\
\noindent \textbf{(P\_6,1.1.2)} sthāne punaḥ dvirvacane sati samudāyasya samudāyaḥ ādeśaḥ . \\
\noindent \textbf{(P\_6,1.1.2)} tatra sampramugdhatvāt prakṛtipratyayasya naṣṭaḥ saḥ ādeśādiḥ bhavati . \\
\noindent \textbf{(P\_6,1.1.2)} dviḥprayoge api dvirvacane sati na doṣaḥ . \\
\noindent \textbf{(P\_6,1.1.2)} vakṣyati tatra liḍgrahaṇasya prayojanam . \\
\noindent \textbf{(P\_6,1.1.2)} liṭi yaḥ ādeśādiḥ tadādeḥ na iti . \\
\noindent \textbf{(P\_6,1.1.2)} iḍvacanam ca yaṅlope . \\
\noindent \textbf{(P\_6,1.1.2)} iṭ ca yaṅlope vaktavyaḥ . \\
\noindent \textbf{(P\_6,1.1.2)} bebhiditā . \\
\noindent \textbf{(P\_6,1.1.2)} bebhiditum . \\
\noindent \textbf{(P\_6,1.1.2)} ekācaḥ upadeśe anudāttāt iti iṭpratiṣedhaḥ prāpnoti . \\
\noindent \textbf{(P\_6,1.1.2)} sthāne punaḥ dvirvacane sati samudāyasya samudāyaḥ ādeśaḥ . \\
\noindent \textbf{(P\_6,1.1.2)} tatra sampramugdhatvāt prakṛtipratyayasya naṣṭaḥ saḥ bhavati yaḥ ekāc upadeśe anudāttaḥ . \\
\noindent \textbf{(P\_6,1.1.2)} dviḥprayoge api dvirvacane sati na doṣaḥ . \\
\noindent \textbf{(P\_6,1.1.2)} ekājgrahaṇena aṅgam viśeṣayiṣyāmaḥ . \\
\noindent \textbf{(P\_6,1.1.2)} ekācaḥ aṅgāt iti . \\
\noindent \textbf{(P\_6,1.1.2)} nanu ca ekaikam atra aṅgam . \\
\noindent \textbf{(P\_6,1.1.2)} samudāye yā vākyaparisamāptiḥ tasya aṅgasañjñā bhaviṣyati . \\
\noindent \textbf{(P\_6,1.1.2)} kutaḥ etat . \\
\noindent \textbf{(P\_6,1.1.2)} śāstrāhāneḥ . \\
\noindent \textbf{(P\_6,1.1.2)} evam hi śāstram ahīnam bhavati . iḍdīrghapratiṣedhaḥ ca . \\
\noindent \textbf{(P\_6,1.1.2)} iṭaḥ dīrghatvasya ca pratiṣedhaḥ vaktavyaḥ . \\
\noindent \textbf{(P\_6,1.1.2)} jarīgṛhitā . \\
\noindent \textbf{(P\_6,1.1.2)} jarīgṛhitum . \\
\noindent \textbf{(P\_6,1.1.2)} grahaḥ aliṭi dīrghaḥ iti dīrghatvam prāpnoti . \\
\noindent \textbf{(P\_6,1.1.2)} sthāne punaḥ dvirvacane samudāyasya samudāyaḥ ādeśaḥ . \\
\noindent \textbf{(P\_6,1.1.2)} tatra sampramugdhatvāt prakṛtipratyayasya naṣṭaḥ grahiḥ . \\
\noindent \textbf{(P\_6,1.1.2)} dviḥprayoge api dvirvacane na doṣaḥ . \\
\noindent \textbf{(P\_6,1.1.2)} grahiṇā aṅgam viśeṣayiṣyāmaḥ . \\
\noindent \textbf{(P\_6,1.1.2)} graheḥ aṅgāt iti . \\
\noindent \textbf{(P\_6,1.1.2)} nanu ca ekaikam atra aṅgam . \\
\noindent \textbf{(P\_6,1.1.2)} samudāye yā vākyaparisamāptiḥ tasya aṅgasañjñā bhaviṣyati . \\
\noindent \textbf{(P\_6,1.1.2)} kutaḥ etat . \\
\noindent \textbf{(P\_6,1.1.2)} śāstrāhāneḥ . \\
\noindent \textbf{(P\_6,1.1.2)} evam hi śāstram ahīnam bhavati . padādividhipratiṣedhaḥ ca . \\
\noindent \textbf{(P\_6,1.1.2)} padādilakṣaṇa vidheḥ pratiṣedhaḥ vaktavyaḥ . \\
\noindent \textbf{(P\_6,1.1.2)} siṣeca . \\
\noindent \textbf{(P\_6,1.1.2)} suṣvāpa . \\
\noindent \textbf{(P\_6,1.1.2)} sātpadādyoḥ iti ṣatvapratiṣedhaḥ prāpnoti . \\
\noindent \textbf{(P\_6,1.1.2)} sthāne punaḥ dvirvacane sati samudāyasya samudāyaḥ ādeśaḥ . \\
\noindent \textbf{(P\_6,1.1.2)} tatra sampramugdhatvāt prakṛtipratyayasya naṣṭaḥ saḥ padādiḥ bhavati . \\
\noindent \textbf{(P\_6,1.1.2)} dviḥprayoge ca api dvirvacane na doṣaḥ . \\
\noindent \textbf{(P\_6,1.1.2)} suptiṅbhyām padam viśeṣayiṣyāmaḥ . \\
\noindent \textbf{(P\_6,1.1.2)} yasmāt suptiṅvidhiḥ tadādi subantam tiṅantam ca . \\
\noindent \textbf{(P\_6,1.1.2)} nanu ca ekaikasmāt [api atra (R)] suptiṅvidhiḥ . \\
\noindent \textbf{(P\_6,1.1.2)} samudāye yā vākyaparisamāptiḥ tayā padasañjñā bhaviṣyati . \\
\noindent \textbf{(P\_6,1.1.2)} kutaḥ etat . \\
\noindent \textbf{(P\_6,1.1.2)} śāstrāhāneḥ . \\
\noindent \textbf{(P\_6,1.1.2)} evam hi śāstram ahīnam bhavati . tau eva suptiṅau tataḥ parau sā eva ca prakṛtiḥ ādyā . \\
\noindent \textbf{(P\_6,1.1.2)} ādigrahaṇam prakṛtam . \\
\noindent \textbf{(P\_6,1.1.2)} samudāyapadatvam etena. \\
\noindent \textbf{(P\_6,1.2.1)} dvitīyasya iti avacanam ajādeḥ iti karmadhārayāt pañcamī . \\
\noindent \textbf{(P\_6,1.2.1)} dvitīyasya iti śakyam avaktum . \\
\noindent \textbf{(P\_6,1.2.1)} katham . \\
\noindent \textbf{(P\_6,1.2.1)} ajādeḥ iti na eṣā bahuvrīheḥ ṣaṣṭhī . \\
\noindent \textbf{(P\_6,1.2.1)} ac ādiḥ yasya saḥ ayam ajādiḥ . \\
\noindent \textbf{(P\_6,1.2.1)} ajādeḥ . \\
\noindent \textbf{(P\_6,1.2.1)} kim tarhi karmadhārayāt pañcamī . \\
\noindent \textbf{(P\_6,1.2.1)} ac ādiḥ ajādiḥ . \\
\noindent \textbf{(P\_6,1.2.1)} ajādeḥ parasya iti . \\
\noindent \textbf{(P\_6,1.2.1)} tatra antareṇa dvitīyagrahaṇam dvitīyasya eva bhaviṣyati . \\
\noindent \textbf{(P\_6,1.2.2)} dvitīyadvirvacane prathamanivṛttiḥ prāptatvāt . \\
\noindent \textbf{(P\_6,1.2.2)} dvitīyadvirvacane prathamasya nivṛttiḥ vaktavyā . \\
\noindent \textbf{(P\_6,1.2.2)} aṭiṭiṣati . \\
\noindent \textbf{(P\_6,1.2.2)} aśiśiṣati iti . \\
\noindent \textbf{(P\_6,1.2.2)} kim kāraṇam . \\
\noindent \textbf{(P\_6,1.2.2)} prāptatvāt . \\
\noindent \textbf{(P\_6,1.2.2)} prāpnoti ekācaḥ dve prathamasya iti . \\
\noindent \textbf{(P\_6,1.2.2)} nanu ca dvitīyadvirvacanam prathamadvirvacanam bādhiṣyate . \\
\noindent \textbf{(P\_6,1.2.2)} katham anyasya ucyamānasya bādhakam syāt . \\
\noindent \textbf{(P\_6,1.2.2)} asati khalu api sambhave bādhanam bhavati asti ca sambhavaḥ yat ubhayam syāt . \\
\noindent \textbf{(P\_6,1.2.2)} na vā prathamavijñāne hi dvitīyāprāptiḥ advitīyatvāt . \\
\noindent \textbf{(P\_6,1.2.2)} na vā vaktavyam . \\
\noindent \textbf{(P\_6,1.2.2)} kim kāraṇam . \\
\noindent \textbf{(P\_6,1.2.2)} prathamavijñāne hi sati dvitīyasya aprāptiḥ syāt . \\
\noindent \textbf{(P\_6,1.2.2)} kim kāraṇam . \\
\noindent \textbf{(P\_6,1.2.2)} advitīyatvāt . \\
\noindent \textbf{(P\_6,1.2.2)} na hi idānīm prathamadvirvacane kṛte dvitīyaḥ dvitīyaḥ bhavati . \\
\noindent \textbf{(P\_6,1.2.2)} kaḥ tarhi . \\
\noindent \textbf{(P\_6,1.2.2)} tṛtīyaḥ . \\
\noindent \textbf{(P\_6,1.2.2)} tat yathā dvayoḥ āsīnayoḥ tṛtīye upajāte na dvitīyaḥ dvitīyaḥ bhavati . \\
\noindent \textbf{(P\_6,1.2.2)} kaḥ tarhi . \\
\noindent \textbf{(P\_6,1.2.2)} tṛtīyaḥ . \\
\noindent \textbf{(P\_6,1.2.2)} na hi kim cit ucyate akṛte dvirvacane yaḥ dvitīyaḥ tasya bhavitavyam iti . \\
\noindent \textbf{(P\_6,1.2.2)} kim tarhi kṛte dvirvacane yaḥ dvitīyaḥ tasya bhaviṣyati . \\
\noindent \textbf{(P\_6,1.2.2)} anārambhasamam evam syāt . \\
\noindent \textbf{(P\_6,1.2.2)} aṭeḥ prathamasya dvirvacanam syāt . \\
\noindent \textbf{(P\_6,1.2.2)} halādiśeṣaḥ . \\
\noindent \textbf{(P\_6,1.2.2)} dvitīyasya dvirvacanam . \\
\noindent \textbf{(P\_6,1.2.2)} halādiśeṣaḥ . \\
\noindent \textbf{(P\_6,1.2.2)} trayāṇām akārāṇām pararūpatve aṭiṣati iti evam rūpam syāt . \\
\noindent \textbf{(P\_6,1.2.2)} na anārambhasamam . \\
\noindent \textbf{(P\_6,1.2.2)} aṭeḥ prathamasya dvirvacanam . \\
\noindent \textbf{(P\_6,1.2.2)} halādiśeṣaḥ . \\
\noindent \textbf{(P\_6,1.2.2)} ittvam . \\
\noindent \textbf{(P\_6,1.2.2)} dvitīyasya dvirvacanam . \\
\noindent \textbf{(P\_6,1.2.2)} halādiśeṣaḥ . \\
\noindent \textbf{(P\_6,1.2.2)} ittvam ḍvayoḥ ikārayoḥ savarṇadīrghatvam . \\
\noindent \textbf{(P\_6,1.2.2)} abhyāsasya asavarṇe iti iyaṅādeśaḥ . \\
\noindent \textbf{(P\_6,1.2.2)} iyaṭiṣati iti etat rūpam yathā syāt . \\
\noindent \textbf{(P\_6,1.2.2)} oṇeḥ ca uvaṇiṣati iti . \\
\noindent \textbf{(P\_6,1.2.2)} na aniṣṭārthā śāstrapravṛttiḥ bhavitum arhati . \\
\noindent \textbf{(P\_6,1.2.2)} yathā vā ādivikāre alaḥ antyavikārābhāvaḥ . \\
\noindent \textbf{(P\_6,1.2.2)} yathā vā ādivividhau alaḥ antyavidhiḥ na bhavati evam dvitīyadvirvacane prathamadvirvacanam na bhaviṣyati . \\
\noindent \textbf{(P\_6,1.2.2)} viṣamaḥ upanyāsaḥ . \\
\noindent \textbf{(P\_6,1.2.2)} na aprāpte alaḥ antyavidhau ādividhiḥ ārabhyate . \\
\noindent \textbf{(P\_6,1.2.2)} saḥ tasya bādhakaḥ bhaviṣyati . \\
\noindent \textbf{(P\_6,1.2.2)} idam api evañjātīyakam . \\
\noindent \textbf{(P\_6,1.2.2)} na aprāpte prathamadvirvacane dvitīyadvirvacanamārabhyate . \\
\noindent \textbf{(P\_6,1.2.2)} tat bādhakam bhaviṣyati . \\
\noindent \textbf{(P\_6,1.2.2)} yat api ucyate asati khalu api sambhave bādhanam bhavati asti ca sambhavaḥ yat ubhayam syāt iti . \\
\noindent \textbf{(P\_6,1.2.2)} na etat asti . \\
\noindent \textbf{(P\_6,1.2.2)} sati api sambhave bādhanam bhavati . \\
\noindent \textbf{(P\_6,1.2.2)} tat yathā dadhi brāhmaṇebhyaḥ dīyatām takram kauṇḍinyāya iti . \\
\noindent \textbf{(P\_6,1.2.2)} sati api dadhidānasya sambhave takradānam nivartakam bhavati . \\
\noindent \textbf{(P\_6,1.2.2)} evam iha api sati api sambhave prathamadvirvacanasya dvitīyadvirvacanam bādhiṣyate . \\
\noindent \textbf{(P\_6,1.2.2)} tatra pūrvasya acaḥ nivṛttau vyañjanasya anivṛttiḥ vaktavyā . \\
\noindent \textbf{(P\_6,1.2.2)} aṭiṭiṣati iti . \\
\noindent \textbf{(P\_6,1.2.2)} yathā eva acaḥ nivṛttiḥ bhavati evam vyañjanasya api prāpnoti . \\
\noindent \textbf{(P\_6,1.2.2)} tatra pūrvasya acaḥ nivṛttau vyañjanānivṛttiḥ aśāsanāt pūrvasya . \\
\noindent \textbf{(P\_6,1.2.2)} tatra pūrvasya acaḥ nivṛttau vyañjanasya anivṛttiḥ siddhā . \\
\noindent \textbf{(P\_6,1.2.2)} kutaḥ . \\
\noindent \textbf{(P\_6,1.2.2)} aśāsanāt pūrvasya . \\
\noindent \textbf{(P\_6,1.2.2)} na iha vayam pūrvasya pratiṣedham śiṣmaḥ . \\
\noindent \textbf{(P\_6,1.2.2)} kim tarhi dvitīyasya dvirvacanam ārabhāmahe . \\
\noindent \textbf{(P\_6,1.2.2)} vyañjanāni punaḥ naṭabhāryavat bhavanti . \\
\noindent \textbf{(P\_6,1.2.2)} tat yathā . \\
\noindent \textbf{(P\_6,1.2.2)} naṭānām striyaḥ raṅgam gatāḥ yaḥ yaḥ pṛcchati kasya yūyam kasya yūyam iti tam tam tava tava iti āhuḥ . \\
\noindent \textbf{(P\_6,1.2.2)} evam vyañjanāni yasya yasya acaḥ kāryam ucyate tam tam bhajante . \\
\noindent \textbf{(P\_6,1.2.2)} ndrādipratiṣedhāt ca . \\
\noindent \textbf{(P\_6,1.2.2)} yat ayam na ndrāḥ saṃyogādayaḥ iti pratiṣedham śāsti tat jñāpayati ācāryaḥ pūrvanivṛttau vyañjanasya anivṛttiḥ iti . \\
\noindent \textbf{(P\_6,1.2.2)} tatra dvitīyābhāve prathamādvirvacanam pratiṣiddhatvāt . \\
\noindent \textbf{(P\_6,1.2.2)} tatra dvitīyasya ekācaḥ abhāve prathamasya dvirvacanam na prāpnoti . \\
\noindent \textbf{(P\_6,1.2.2)} āṭatuḥ . \\
\noindent \textbf{(P\_6,1.2.2)} āṭuḥ . \\
\noindent \textbf{(P\_6,1.2.2)} kim kāraṇam . \\
\noindent \textbf{(P\_6,1.2.2)} pratiṣiddhatvāt . \\
\noindent \textbf{(P\_6,1.2.2)} ajādeḥ dvitīyasya iti pratiṣedhāt . \\
\noindent \textbf{(P\_6,1.2.2)} na eṣa doṣaḥ . \\
\noindent \textbf{(P\_6,1.2.2)} sati tasmin pratiṣedhaḥ . \\
\noindent \textbf{(P\_6,1.2.2)} sati dvitīyadvirvacane prathamasya pratiṣedhaḥ . \\
\noindent \textbf{(P\_6,1.2.2)} sati tasmin pratiṣedhaḥ iti cet halādiśeṣe doṣaḥ . \\
\noindent \textbf{(P\_6,1.2.2)} sati tasmin pratiṣedhaḥ iti cet halādiśeṣe doṣaḥ bhavati . \\
\noindent \textbf{(P\_6,1.2.2)} halādiśeṣe sati ādye hali anādyasya lopaḥ syāt . \\
\noindent \textbf{(P\_6,1.2.2)} iha eva syāt . \\
\noindent \textbf{(P\_6,1.2.2)} papāca . \\
\noindent \textbf{(P\_6,1.2.2)} papāṭha iti . \\
\noindent \textbf{(P\_6,1.2.2)} iha na syāt . \\
\noindent \textbf{(P\_6,1.2.2)} āṭatuḥ . \\
\noindent \textbf{(P\_6,1.2.2)} āṭuḥ iti . \\
\noindent \textbf{(P\_6,1.2.2)} lokavat halādiśeṣe . \\
\noindent \textbf{(P\_6,1.2.2)} lokavat halādiśeṣe siddham . \\
\noindent \textbf{(P\_6,1.2.2)} tat yathā loke īśvaraḥ ājñāpayati grāmāt grāmāt manuṣyāḥ ānīyantām prāgāṅgam grāmebhyaḥ brāhmaṇāḥ ānīyantām iti . \\
\noindent \textbf{(P\_6,1.2.2)} yeṣu tatra grāmeṣu brāhmaṇāḥ na santi na tarhi idānīm tataḥ anyasya ānayanam bhavati . \\
\noindent \textbf{(P\_6,1.2.2)} yathā tatra kva cit api brāhmaṇasya sattā (R: sarvatra) abrāhmaṇasya nivarttikā bhavati evam iha api kva cit api hal ādyaḥ san sarvasya anādyasya halaḥ nivartakaḥ bhavati . \\
\noindent \textbf{(P\_6,1.2.2)} kva cit anyatra lopaḥ iti cet dvirvacanam . \\
\noindent \textbf{(P\_6,1.2.2)} kva cit anyatra lopaḥ iti cet dvirvacanam api evam prāpnoti . \\
\noindent \textbf{(P\_6,1.2.2)} kva cit api dvitīyaḥ san sarvasya prathamasya nivartakaḥ syāt . \\
\noindent \textbf{(P\_6,1.2.2)} tasmāt astu sati tasmin pratiṣedhaḥ iti eva. nanu ca uktam sati tasmin pratiṣedhaḥ iti cet halādiśeṣe doṣaḥ iti . \\
\noindent \textbf{(P\_6,1.2.2)} pratividhāsyate halādiśeṣe . \\
\noindent \textbf{(P\_6,1.3)} kimartham idam ucyate . \\
\noindent \textbf{(P\_6,1.3)} ndrādeḥ dvirvacanaprasaṅgaḥ tatra ndrāṇām pratiṣedhaḥ . \\
\noindent \textbf{(P\_6,1.3)} ndrādeḥ ekācaḥ dvirvacanam prāpnoti . \\
\noindent \textbf{(P\_6,1.3)} tatra ndrāṇām saṃyogādīnām pratiṣedhaḥ ucyate . \\
\noindent \textbf{(P\_6,1.3)} īrṣyateḥ tṛtīyasya . \\
\noindent \textbf{(P\_6,1.3)} īrṣyateḥ tṛtīyasya dve bhavataḥ iti vaktavyam . \\
\noindent \textbf{(P\_6,1.3)} ke cit tāvat āhuḥ ekācaḥ iti . \\
\noindent \textbf{(P\_6,1.3)} īrṣyiṣiṣati . \\
\noindent \textbf{(P\_6,1.3)} aparaḥ āha :vyañjanasya iti : īrṣyiyiṣati . \\
\noindent \textbf{(P\_6,1.3)} kaṇḍvādīnām ca . \\
\noindent \textbf{(P\_6,1.3)} kaṇḍvādīnām ca tṛtīyasya dve bhavataḥ iti vaktavyam . \\
\noindent \textbf{(P\_6,1.3)} kaṇḍūyiyiṣati . \\
\noindent \textbf{(P\_6,1.3)} asūyiyiṣati . \\
\noindent \textbf{(P\_6,1.3)} vā nāmadhātūnām . \\
\noindent \textbf{(P\_6,1.3)} vā nāmadhātūnām tṛtīyasya dve bhavataḥ iti vaktavyam . \\
\noindent \textbf{(P\_6,1.3)} aśvīyiyiṣati . \\
\noindent \textbf{(P\_6,1.3)} aśiśvīyiṣati . \\
\noindent \textbf{(P\_6,1.3)} aparaḥ āha yatheṣṭam vā . \\
\noindent \textbf{(P\_6,1.3)} yatheṣṭam vā nāmadhātūnām iti . \\
\noindent \textbf{(P\_6,1.3)} puputrīyiṣati . \\
\noindent \textbf{(P\_6,1.3)} putitrīyiṣati . \\
\noindent \textbf{(P\_6,1.3)} putrīyiyiṣati . \\
\noindent \textbf{(P\_6,1.4)} pūrvaḥ abhyāsaḥ iti ucyate . \\
\noindent \textbf{(P\_6,1.4)} kasya pūrvaḥ abhyāsasañjñaḥ bhavati . \\
\noindent \textbf{(P\_6,1.4)} dve iti vartate . \\
\noindent \textbf{(P\_6,1.4)} dvayoḥ iti vaktavyam . \\
\noindent \textbf{(P\_6,1.4)} saḥ tarhi tathā nirdeśaḥ kartavyaḥ . \\
\noindent \textbf{(P\_6,1.4)} na kartavyaḥ . \\
\noindent \textbf{(P\_6,1.4)} arthāt vibhaktivipariṇāmaḥ bhaviṣyati . \\
\noindent \textbf{(P\_6,1.4)} tat yathā . \\
\noindent \textbf{(P\_6,1.4)} uccāni devadattasya gṛhāṇi . \\
\noindent \textbf{(P\_6,1.4)} āmantrayasva enam . \\
\noindent \textbf{(P\_6,1.4)} devadattam iti gamyate . \\
\noindent \textbf{(P\_6,1.4)} devadattasya gāvaḥ aśvāḥ hiraṇyam iti . \\
\noindent \textbf{(P\_6,1.4)} āḍhyaḥ vaidhaveyaḥ . \\
\noindent \textbf{(P\_6,1.4)} devadattaḥ iti gamyate . \\
\noindent \textbf{(P\_6,1.4)} purastāt ṣaṣṭhīnirdiṣṭam sat arthāt dvitīyānirdiṣṭam prathamānirdiṣṭam ca bhavati . \\
\noindent \textbf{(P\_6,1.4)} evam iha api purastāt prathamānirdiṣṭam sat arthāt ṣaṣṭhīnirdiṣṭam bhaviṣyati . \\
\noindent \textbf{(P\_6,1.5)} abhyastasañjñāyām sahavacanam . \\
\noindent \textbf{(P\_6,1.5)} abhyastasañjñāyām sahagrahaṇam kartavyam . \\
\noindent \textbf{(P\_6,1.5)} kim prayojanam . \\
\noindent \textbf{(P\_6,1.5)} ādyudāttatve pṛthagaprasaṅgārtham . \\
\noindent \textbf{(P\_6,1.5)} ādyudāttatvam saha bhūtayoḥ yathā syāt . \\
\noindent \textbf{(P\_6,1.5)} ekaikasya mā bhūt iti . \\
\noindent \textbf{(P\_6,1.5)} yasmin eva abhyastakārye adoṣaḥ tat eva paṭhitam anudāttam padam ekavarjam iti . \\
\noindent \textbf{(P\_6,1.5)} na asti yaugapadyena sambhavaḥ . \\
\noindent \textbf{(P\_6,1.5)} paryāyaḥ tarhi prasajyeta . \\
\noindent \textbf{(P\_6,1.5)} paryāyaḥ ca . \\
\noindent \textbf{(P\_6,1.5)} pūrvasya tāvat pareṇa rūpeṇa vyavahitatvāt na bhaviṣyati . \\
\noindent \textbf{(P\_6,1.5)} parasya tarhi syāt . \\
\noindent \textbf{(P\_6,1.5)} tatra ācāryapravṛttiḥ jñāpayati na parasya bhavati iti yat ayam bibhetyādīnām piti pratyayāt pūrvam udāttam bhavati iti āha . \\
\noindent \textbf{(P\_6,1.5)} evam vyavadhānāt na pūrvasya jñāpakāt na parasya ucyate ca idam abhyastānām ādiḥ udāttaḥ bhavati iti . \\
\noindent \textbf{(P\_6,1.5)} tatra saḥ eva doṣaḥ paryāyaḥ prasajyeta . \\
\noindent \textbf{(P\_6,1.5)} tasmāt sahagrahaṇam kartavyam . \\
\noindent \textbf{(P\_6,1.5)} na kartavyam . \\
\noindent \textbf{(P\_6,1.5)} ubhegrahaṇam kriyate . \\
\noindent \textbf{(P\_6,1.5)} tat sahārtham vijñāsyate . \\
\noindent \textbf{(P\_6,1.5)} asti anyat ubhegrahaṇasya prayojanam . \\
\noindent \textbf{(P\_6,1.5)} kim . \\
\noindent \textbf{(P\_6,1.5)} ubhegrahaṇam sañjñinirdeśārtham . \\
\noindent \textbf{(P\_6,1.5)} antareṇa api ubhegrahaṇam prakḷptaḥ sañjñinirdeśaḥ . \\
\noindent \textbf{(P\_6,1.5)} katham . \\
\noindent \textbf{(P\_6,1.5)} dve iti vartate . \\
\noindent \textbf{(P\_6,1.5)} idam tarhi prayojanam . \\
\noindent \textbf{(P\_6,1.5)} yatra ubhe śabdarūpe śrūyete tatra abhyastasañjñā yathā syāt . \\
\noindent \textbf{(P\_6,1.5)} iha mā bhūt : īrtsanti , īpsanti , īrtsan , īpsan , airtsan , aipsan . \\
\noindent \textbf{(P\_6,1.5)} kim ca syāt . \\
\noindent \textbf{(P\_6,1.5)} adbhāvaḥ numpratiṣedhaḥ jusbhāvaḥ iti ete vidhayaḥ prasajyeran . \\
\noindent \textbf{(P\_6,1.5)} adbhāve tāvat na doṣaḥ . \\
\noindent \textbf{(P\_6,1.5)} saptame yogavibhāgaḥ kariṣyate . \\
\noindent \textbf{(P\_6,1.5)} idam asti : at abhyastāt . \\
\noindent \textbf{(P\_6,1.5)} tataḥ ātmanepadeṣu . \\
\noindent \textbf{(P\_6,1.5)} ātmanepadeṣu ca at bhavati . \\
\noindent \textbf{(P\_6,1.5)} anataḥ iti ubhayoḥ śeṣaḥ . \\
\noindent \textbf{(P\_6,1.5)} yat api ucyate numpratiṣedhaḥ iti ekādeśe kṛte vyapavargābhāvāt na bhaviṣyati . \\
\noindent \textbf{(P\_6,1.5)} idam iha sampradhāryam . \\
\noindent \textbf{(P\_6,1.5)} numpratiṣedhaḥ kriyatām ekādeśaḥ iti . \\
\noindent \textbf{(P\_6,1.5)} kim atra kartavyam . \\
\noindent \textbf{(P\_6,1.5)} paratvāt numpratiṣedhaḥ . \\
\noindent \textbf{(P\_6,1.5)} nityaḥ ekādeśaḥ . \\
\noindent \textbf{(P\_6,1.5)} kṛte api numpratiṣedhe prāpnoti akṛte api . \\
\noindent \textbf{(P\_6,1.5)} ekādeśaḥ api nityaḥ . \\
\noindent \textbf{(P\_6,1.5)} anyasya kṛte numpratiṣedhe prāpnoti anyasya akṛte . \\
\noindent \textbf{(P\_6,1.5)} śabdāntarasya ca prāpnuvan vidhiḥ anityaḥ bhavati . \\
\noindent \textbf{(P\_6,1.5)} antaraṅgaḥ tarhi ekādeśaḥ . \\
\noindent \textbf{(P\_6,1.5)} kā antaraṅgatā . \\
\noindent \textbf{(P\_6,1.5)} varṇau āśritya ekādeśaḥ . \\
\noindent \textbf{(P\_6,1.5)} vidhiviṣaye numpratiṣedhaḥ vidhiḥ ca numaḥ sarvanāmasthāne prāk tu sarvanāmasthānotpatteḥ ekādeśaḥ . \\
\noindent \textbf{(P\_6,1.5)} tatra nityatvāt ca antaraṅgatvāt ca ekādeśaḥ . \\
\noindent \textbf{(P\_6,1.5)} ekādeśe kṛte vyapavargābhāvāt na bhaviṣyati . \\
\noindent \textbf{(P\_6,1.5)} yat api ucyate jusbhāvaḥ iti ekādeśe kṛte vyapavargābhāvāt na bhaviṣyati . \\
\noindent \textbf{(P\_6,1.5)} ekādeśe iti ucyate kena ca atra ekādeśaḥ . \\
\noindent \textbf{(P\_6,1.5)} antinā . \\
\noindent \textbf{(P\_6,1.5)} na atra antibhāvaḥ prāpnoti . \\
\noindent \textbf{(P\_6,1.5)} kim kāraṇam . \\
\noindent \textbf{(P\_6,1.5)} jusbhāvena bādhyate . \\
\noindent \textbf{(P\_6,1.5)} na atra jusbhāvaḥ prāpnoti . \\
\noindent \textbf{(P\_6,1.5)} kim kāraṇam . \\
\noindent \textbf{(P\_6,1.5)} śapā vyavahitatvāt . \\
\noindent \textbf{(P\_6,1.5)} ekādeśe kṛte na asti vyavadhānam . \\
\noindent \textbf{(P\_6,1.5)} ekādeśaḥ pūrvavidhau sthānivat bhavati iti vyavadhānam eva . \\
\noindent \textbf{(P\_6,1.5)} kim puṇaḥ kāraṇam nimittavān antiḥ ekādeśam tāvat pratīkṣate na punaḥ tāvati eva nimittam asti iti antibhāvena bhāvyam . \\
\noindent \textbf{(P\_6,1.5)} iha api tarhi tāvati eva nimittam asti iti antibhāvaḥ syāt . \\
\noindent \textbf{(P\_6,1.5)} anenijuḥ . \\
\noindent \textbf{(P\_6,1.5)} paryaveviṣuḥ iti . \\
\noindent \textbf{(P\_6,1.5)} astu . \\
\noindent \textbf{(P\_6,1.5)} antibhāve kṛte sthānivadbhāvāt jhigrahaṇena grahaṇāt jusbhāvaḥ bhaviṣyati . \\
\noindent \textbf{(P\_6,1.5)} atha vā yadi api nimittavān antiḥ ayam tasya jusbhāvaḥ apavādaḥ . \\
\noindent \textbf{(P\_6,1.5)} na ca apavādaviṣaye utsargāḥ abhiniviśante . \\
\noindent \textbf{(P\_6,1.5)} pūrvam hi apavādāḥ abhiniviśante paścāt utsargāḥ . \\
\noindent \textbf{(P\_6,1.5)} prakalpya vā apavādaviṣayam tataḥ utasrgaḥ pravartate . \\
\noindent \textbf{(P\_6,1.5)} na tāvat atra kadā cit api antibhāvaḥ bhavati . \\
\noindent \textbf{(P\_6,1.5)} apavādam jusbhāvam pratīkṣate . \\
\noindent \textbf{(P\_6,1.5)} na khalu api kva cit abhyastānām jheḥ ca ānantaryam . \\
\noindent \textbf{(P\_6,1.5)} sarvatra vikaraṇaiḥ vyavadhānam . \\
\noindent \textbf{(P\_6,1.5)} tena anena avaśyam vikaraṇanāśaḥ pratīkṣyaḥ kva cit lukā kva cit ślunā kva cit ekādeśena . \\
\noindent \textbf{(P\_6,1.5)} saḥ yathā ślulukau pratīkṣate evam ekādeśam api pratīkṣate . \\
\noindent \textbf{(P\_6,1.5)} evam tarhi idam iha vyapadeśyam sat ācāryaḥ na vyapadiśati . \\
\noindent \textbf{(P\_6,1.5)} kim . \\
\noindent \textbf{(P\_6,1.5)} sthānivadbhāvam . \\
\noindent \textbf{(P\_6,1.5)} sthānivadbhāvāt vyavadhānam . \\
\noindent \textbf{(P\_6,1.5)} vyavadhānāt na bhaviṣyati . \\
\noindent \textbf{(P\_6,1.5)} pūrvavidhau sthānivadbhāvaḥ na ca ayam pūrvasya vidhiḥ . \\
\noindent \textbf{(P\_6,1.5)} pūrvasmāt api vidhiḥ pūrvavidhiḥ . \\
\noindent \textbf{(P\_6,1.5)} tat etat asati prayojane ubhegrahaṇam sahārtham vijñāsyate . \\
\noindent \textbf{(P\_6,1.5)} katham kṛtvā ekaikasya abhyastasañjñā prāpnoti . \\
\noindent \textbf{(P\_6,1.5)} pratyekam vākyaparisamāptiḥ dṛṣṭā . \\
\noindent \textbf{(P\_6,1.5)} tat yathā . \\
\noindent \textbf{(P\_6,1.5)} vṛddhiguṇasañjñe pratyekam bhavataḥ . \\
\noindent \textbf{(P\_6,1.5)} nanu ca ayam api asti dṛṣṭāntaḥ samudāye vākyaparisamāptiḥ iti . \\
\noindent \textbf{(P\_6,1.5)} tat yathā . \\
\noindent \textbf{(P\_6,1.5)} gargāḥ śatam daṇḍyantām iti . \\
\noindent \textbf{(P\_6,1.5)} arthinaḥ ca rājānaḥ hiraṇyena bhavanti na ca pratyekam daṇḍayanti . \\
\noindent \textbf{(P\_6,1.5)} sati etasmin dṛṣṭānte yadi tatra pratyekam iti ucyate iha api sahagrahaṇam kartavyam . \\
\noindent \textbf{(P\_6,1.5)} atha tatra antareṇa pratyekam iti vacanam pratyekam guṇvṛddhisañjñe bhavataḥ iha api na arthaḥ sahagrahaṇena . \\
\noindent \textbf{(P\_6,1.6)} jakṣityādiṣu saptagrahaṇam vevītyartham . \\
\noindent \textbf{(P\_6,1.6)} jakṣityādiṣu saptagrahaṇam kartavyam . \\
\noindent \textbf{(P\_6,1.6)} sapta jakṣityādayaḥ abhyastasañjñakāḥ bhavanti iti vaktavyam . \\
\noindent \textbf{(P\_6,1.6)} kim prayojanam . \\
\noindent \textbf{(P\_6,1.6)} vevītyartham . \\
\noindent \textbf{(P\_6,1.6)} vevīteḥ abhyastasañjñā yathā syāt . \\
\noindent \textbf{(P\_6,1.6)} vevyate . \\
\noindent \textbf{(P\_6,1.6)} aparigaṇanam vā āgaṇāntatvāt . \\
\noindent \textbf{(P\_6,1.6)} na vā arthaḥ parigaṇanena . \\
\noindent \textbf{(P\_6,1.6)} astu āgaṇāntam abhyastasañjñā . \\
\noindent \textbf{(P\_6,1.6)} iha api tarhi prāpnoti . \\
\noindent \textbf{(P\_6,1.6)} āṅaḥ śāsu . \\
\noindent \textbf{(P\_6,1.6)} astu . \\
\noindent \textbf{(P\_6,1.6)} abhyastakāryāṇi kasmāt na bhavanti . \\
\noindent \textbf{(P\_6,1.6)} bhūyiṣṭhāni parasmaipadeṣu ātmanepadī ca ayam . \\
\noindent \textbf{(P\_6,1.6)} svaraḥ tarhi prāpnoti . \\
\noindent \textbf{(P\_6,1.6)} yatra api asya ātmanepadeṣu abhyastakāryam svaraḥ tatra api anudāttetaḥ param lasārvadhātukam anudāttam bhavati iti anudāttatve kṛte na asti viśeṣaḥ dhātusvareṇa udāttatve sati abhyastasvareṇa vā . \\
\noindent \textbf{(P\_6,1.6)} ṣasivaśī chāndasau . \\
\noindent \textbf{(P\_6,1.6)} dṛṣṭānuvidhiḥ chandasi bhavati . \\
\noindent \textbf{(P\_6,1.6)} carkarītam abhyastam eva . \\
\noindent \textbf{(P\_6,1.6)} hnuṅaḥ tarhi prāpnoti . \\
\noindent \textbf{(P\_6,1.6)} astu . \\
\noindent \textbf{(P\_6,1.6)} abhyastakāryāṇi kasmāt na bhavanti . \\
\noindent \textbf{(P\_6,1.6)} bhūyiṣṭhāni parasmaipadeṣu ātmanepadī ca ayam . \\
\noindent \textbf{(P\_6,1.6)} svaraḥ tarhi prāpnoti . \\
\noindent \textbf{(P\_6,1.6)} ahnviṅoḥ iti pratiṣedhavidhānasāmarthyāt svaraḥ na bhaviṣyati . \\
\noindent \textbf{(P\_6,1.6)} atha vā sapta eva ime dhātavaḥ paṭhyante . \\
\noindent \textbf{(P\_6,1.6)} jakṣ abhyastasañjñaḥ bhavati . \\
\noindent \textbf{(P\_6,1.6)} ityādayaḥ ca ṣaṭ . \\
\noindent \textbf{(P\_6,1.6)} jakṣ ityādayaḥ ṣaṭ iti . \\
\noindent \textbf{(P\_6,1.7)} tujādiṣu chandaḥpratyayagrahaṇam . \\
\noindent \textbf{(P\_6,1.7)} tujādiṣu chandaḥpratyayagrahaṇam kartavyam . \\
\noindent \textbf{(P\_6,1.7)} chandasi tujādīnām dīrghaḥ bhavati iti vaktavyam . \\
\noindent \textbf{(P\_6,1.7)} asmin ca asmin ca pratyaye iti vaktavyam . \\
\noindent \textbf{(P\_6,1.7)} iha mā bhūt . \\
\noindent \textbf{(P\_6,1.7)} tutoja śabalān harān . \\
\noindent \textbf{(P\_6,1.7)} anārambhaḥ vā aparigaṇitatvāt . \\
\noindent \textbf{(P\_6,1.7)} anārambhaḥ vā chandasi dīrghatvasya nyāyyaḥ . \\
\noindent \textbf{(P\_6,1.7)} kutaḥ . \\
\noindent \textbf{(P\_6,1.7)} aparigaṇitatvāt . \\
\noindent \textbf{(P\_6,1.7)} na hi chandasi dīrghatvasya parigaṇanam kartum śakyam . \\
\noindent \textbf{(P\_6,1.7)} kim kāraṇam . \\
\noindent \textbf{(P\_6,1.7)} anyeṣām ca darśanāt . \\
\noindent \textbf{(P\_6,1.7)} yeṣām api dīrghatvam na ārabhyate teṣām api chandasi dīrghatvam dṛśyate . \\
\noindent \textbf{(P\_6,1.7)} tat yathā pūruṣaḥ . \\
\noindent \textbf{(P\_6,1.7)} nārakaḥ iti . \\
\noindent \textbf{(P\_6,1.7)} anekāntatvāt ca . \\
\noindent \textbf{(P\_6,1.7)} yeṣām ca api ārabhyate teṣām api anekāntaḥ . \\
\noindent \textbf{(P\_6,1.7)} yasmin eva ca pratyaye dīrghatvam dṛśyate tasmin eva ca na dṛśyate . \\
\noindent \textbf{(P\_6,1.7)} māmahānaḥ ukthapātram . \\
\noindent \textbf{(P\_6,1.7)} mamahānaḥ iti ca . \\
\noindent \textbf{(P\_6,1.8)} dhātoḥ iti kimartham . \\
\noindent \textbf{(P\_6,1.8)} īhām cakre . \\
\noindent \textbf{(P\_6,1.8)} na etat asti . \\
\noindent \textbf{(P\_6,1.8)} liṭi iti ucyate na ca atra liṭam paśyāmaḥ . \\
\noindent \textbf{(P\_6,1.8)} pratyayalakṣaṇena . \\
\noindent \textbf{(P\_6,1.8)} na lumatā tasmin iti pratyayalakṣaṇapratiṣedhaḥ . \\
\noindent \textbf{(P\_6,1.8)} idam tarhi . \\
\noindent \textbf{(P\_6,1.8)} sasṛvāṃsaḥ viśṛṇvire . \\
\noindent \textbf{(P\_6,1.8)} liṭi dvirvacane jāgarteḥ vāvacanam . \\
\noindent \textbf{(P\_6,1.8)} liṭi dvirvacane jāgarteḥ vā iti vaktavyam . \\
\noindent \textbf{(P\_6,1.8)} yaḥ jāgara tam ṛcaḥ kāmayante . \\
\noindent \textbf{(P\_6,1.8)} yaḥ jajāgāra tam ṛcaḥ kāmayante . \\
\noindent \textbf{(P\_6,1.8)} anabhyāsasya iti kim . \\
\noindent \textbf{(P\_6,1.8)} kṛṣṇaḥ nonāva vṛṣabhaḥ yadi idam . \\
\noindent \textbf{(P\_6,1.8)} nonūyateḥ nonāva . \\
\noindent \textbf{(P\_6,1.8)} abhyāsapratiṣedhānarthakyam ca chandasi vāvacanāt . \\
\noindent \textbf{(P\_6,1.8)} abhyāsapratiṣedhaḥ ca anarthakaḥ . \\
\noindent \textbf{(P\_6,1.8)} kim kāraṇam . \\
\noindent \textbf{(P\_6,1.8)} chandasi vāvacanāt . \\
\noindent \textbf{(P\_6,1.8)} avaśyam chandasi vā dve bhavataḥ iti vaktavyam . \\
\noindent \textbf{(P\_6,1.8)} kim prayojanam . \\
\noindent \textbf{(P\_6,1.8)} prayojanam ādityān yāciṣāmahe . \\
\noindent \textbf{(P\_6,1.8)} yiyāciṣāmahe iti prāpte . \\
\noindent \textbf{(P\_6,1.8)} devatā no dāti priyāṇi . \\
\noindent \textbf{(P\_6,1.8)} dadāti priyāṇi . \\
\noindent \textbf{(P\_6,1.8)} maghavā dātu . \\
\noindent \textbf{(P\_6,1.8)} maghavā dadātu . \\
\noindent \textbf{(P\_6,1.8)} saḥ naḥ stutaḥ vīravat dhātu . \\
\noindent \textbf{(P\_6,1.8)} vīravat dadhātu . \\
\noindent \textbf{(P\_6,1.8)} yāvatā idānīm chandasi vā dve bhavataḥ iti ucyate dhātugrahaṇena api na arthaḥ . \\
\noindent \textbf{(P\_6,1.8)} kasmāt na bhavati sasṛvāṃsaḥ viśṛṇvire iti . \\
\noindent \textbf{(P\_6,1.8)} chandasi vāvacanāt . \\
\noindent \textbf{(P\_6,1.8)} tat etat dhātugrahaṇam sānnyāsikam tiṣṭhatu tāvat . \\
\noindent \textbf{(P\_6,1.9)} kim iyam ṣaṣṭhī āhosvit saptamī . \\
\noindent \textbf{(P\_6,1.9)} kutaḥ sandehaḥ . \\
\noindent \textbf{(P\_6,1.9)} samānaḥ nirdeśaḥ . \\
\noindent \textbf{(P\_6,1.9)} kim ca ataḥ . \\
\noindent \textbf{(P\_6,1.9)} yadi ṣaṣṭhī sanyaṅantasya dvirvacanena bhavitavyam . \\
\noindent \textbf{(P\_6,1.9)} atha saptamī sanyaṅoḥ parataḥ pūrvasya dvirvacanam . \\
\noindent \textbf{(P\_6,1.9)} kaḥ ca atra viśeṣaḥ . \\
\noindent \textbf{(P\_6,1.9)} sanyaṅoḥ parataḥ iti cet iṭaḥ dvirvacanam parāditvāt . \\
\noindent \textbf{(P\_6,1.9)} sanyaṅoḥ parataḥ iti cet iṭaḥ dvirvacanam kartavyam . \\
\noindent \textbf{(P\_6,1.9)} aṭiṭiṣati . \\
\noindent \textbf{(P\_6,1.9)} aśiśiṣati . \\
\noindent \textbf{(P\_6,1.9)} kim punaḥ kāraṇam na sidhyati . \\
\noindent \textbf{(P\_6,1.9)} parāditvāt . \\
\noindent \textbf{(P\_6,1.9)} iṭ parādiḥ . \\
\noindent \textbf{(P\_6,1.9)} hanteḥ ca īṭaḥ . \\
\noindent \textbf{(P\_6,1.9)} hanteḥ ca īṭaḥ dvirvacanam kartavyam . \\
\noindent \textbf{(P\_6,1.9)} jeghnīyate . \\
\noindent \textbf{(P\_6,1.9)} nanu ca yasya api sanyaṅantasya dvirvacanam tasya api sthānivadbhāvaprasaṅgaḥ . \\
\noindent \textbf{(P\_6,1.9)} īṭi sthānivadbhāvāt īṭaḥ dvirvacanam na prāpnoti . \\
\noindent \textbf{(P\_6,1.9)} na eṣaḥ doṣaḥ . \\
\noindent \textbf{(P\_6,1.9)} dvirvacananimitte aci sthānivat iti ucyate na ca asau dvirvacananimittam . \\
\noindent \textbf{(P\_6,1.9)} yasmin api dvirvacanam yasya api dvirvacanam sarvaḥ asau dvirvacananimittam . \\
\noindent \textbf{(P\_6,1.9)} tasmāt īṭaḥ dvirvacanam . \\
\noindent \textbf{(P\_6,1.9)} tasmāt ubhābhyām īṭaḥ dvirvacanam kartavyam . \\
\noindent \textbf{(P\_6,1.9)} yaḥ ca ubhayoḥ doṣaḥ na tam ekaḥ codyaḥ bhavati . \\
\noindent \textbf{(P\_6,1.9)} ekācaḥ upadeśe anudāttāt iti upadeśavacanam udāttaviśeṣaṇam cet sanaḥ iṭpratiṣedhaḥ . \\
\noindent \textbf{(P\_6,1.9)} ekācaḥ upadeśe anudāttāt iti upadeśavacanam udāttaviśeṣaṇam cet sanaḥ iṭpratiṣedhaḥ vaktavyaḥ . \\
\noindent \textbf{(P\_6,1.9)} bibhitsati . \\
\noindent \textbf{(P\_6,1.9)} cicchitsati . \\
\noindent \textbf{(P\_6,1.9)} dvirvacane kṛte upadeśe anudāttāt ekācaḥ śrūyamāṇāt iti iṭpratiṣedhaḥ na prāpnoti . \\
\noindent \textbf{(P\_6,1.9)} astu tarhi sanyaṅantasya . \\
\noindent \textbf{(P\_6,1.9)} sanyaṅantasya iti cet aśeḥ sani aniṭaḥ . \\
\noindent \textbf{(P\_6,1.9)} sanyaṅantasya iti cet aśeḥ sani aniṭaḥ dvirvacanam vaktavyam . \\
\noindent \textbf{(P\_6,1.9)} iyakṣamāṇāḥ bhṛgubhiḥ sajoṣāḥ . \\
\noindent \textbf{(P\_6,1.9)} yasya api sanyaṅoḥ parataḥ dvirvacanam tena api atra avaśyam iḍabhāve yatnaḥ kartavyaḥ . \\
\noindent \textbf{(P\_6,1.9)} kim kāraṇam . \\
\noindent \textbf{(P\_6,1.9)} aśeḥ hi pratipadam iṭ vidhīyate smipūṅrañjvaśām sani iti . \\
\noindent \textbf{(P\_6,1.9)} tena eva dvitīyadvirvacanam api na bhaviṣyati . \\
\noindent \textbf{(P\_6,1.9)} atha vā na etat aśeḥ rūpam . \\
\noindent \textbf{(P\_6,1.9)} yajeḥ eṣaḥ chāndasaḥ varṇalopaḥ . \\
\noindent \textbf{(P\_6,1.9)} tat yathā tubhya idam agne . \\
\noindent \textbf{(P\_6,1.9)} tubhyam idam agne iti prāpte . \\
\noindent \textbf{(P\_6,1.9)} ambānām carum . \\
\noindent \textbf{(P\_6,1.9)} nāmbānām carum iti prāpte . \\
\noindent \textbf{(P\_6,1.9)} āvyādhinīḥ ugaṇāḥ . \\
\noindent \textbf{(P\_6,1.9)} sugaṇāḥ iti prāpte . \\
\noindent \textbf{(P\_6,1.9)} iṣkartaram adhvarasya . \\
\noindent \textbf{(P\_6,1.9)} niṣkartāram iti prāpte . \\
\noindent \textbf{(P\_6,1.9)} śivā udrasya bheṣajī . \\
\noindent \textbf{(P\_6,1.9)} śivā rudrasya bheṣajī iti prāpte . \\
\noindent \textbf{(P\_6,1.9)} aśyarthaḥ vai gamyate . \\
\noindent \textbf{(P\_6,1.9)} kaḥ punaḥ aśeḥ arthaḥ . \\
\noindent \textbf{(P\_6,1.9)} aśnotiḥ vyaptikarmā . \\
\noindent \textbf{(P\_6,1.9)} yajiḥ api aśyarthe vartate . \\
\noindent \textbf{(P\_6,1.9)} katham punaḥ anyaḥ nāma anyasya arthe vartate . \\
\noindent \textbf{(P\_6,1.9)} bahvarthāḥ api dhātavaḥ bhavanti iti . \\
\noindent \textbf{(P\_6,1.9)} tat yathā . \\
\noindent \textbf{(P\_6,1.9)} vapiḥ prakiraṇe dṛṣṭaḥ chedane ca api vartate . \\
\noindent \textbf{(P\_6,1.9)} keśān vapati iti . \\
\noindent \textbf{(P\_6,1.9)} īḍiḥ studicodanāyācñāsu dṛṣṭaḥ īraṇe ca api vartate . \\
\noindent \textbf{(P\_6,1.9)} agniḥ vai itaḥ vṛṣṭim īṭṭe . \\
\noindent \textbf{(P\_6,1.9)} marutaḥ amutaḥ cyāvayanti . \\
\noindent \textbf{(P\_6,1.9)} karotiḥ ayam abhūtaprādurbhāve dṛṣṭaḥ nirmalīkaraṇe ca api vartate . \\
\noindent \textbf{(P\_6,1.9)} pṛṣṭham kuru . \\
\noindent \textbf{(P\_6,1.9)} pādau kuru . \\
\noindent \textbf{(P\_6,1.9)} unmṛdāna iti gamyate . \\
\noindent \textbf{(P\_6,1.9)} nikṣepaṇe ca api dṛśyate . \\
\noindent \textbf{(P\_6,1.9)} kaṭe kuru . \\
\noindent \textbf{(P\_6,1.9)} ghaṭe kuru . \\
\noindent \textbf{(P\_6,1.9)} aśmānam itaḥ kuru . \\
\noindent \textbf{(P\_6,1.9)} sthāpaya iti gamyate . \\
\noindent \textbf{(P\_6,1.9)} evam tarhi dīrghakutvaprasāraṇaṣatvam adhikasya dvirvacanāt . dīrghatvam dvirvacanādhikasya na sidhyati . \\
\noindent \textbf{(P\_6,1.9)} cicīṣati . \\
\noindent \textbf{(P\_6,1.9)} tuṣṭūṣati . \\
\noindent \textbf{(P\_6,1.9)} samudāyasya samudāyaḥ ādeśaḥ . \\
\noindent \textbf{(P\_6,1.9)} tatra sampramugdhatvāt prakṛtipratyayasya naṣṭaḥ san bhavati . \\
\noindent \textbf{(P\_6,1.9)} tatra ajantānām sani iti dīrghatvam na prāpnoti . \\
\noindent \textbf{(P\_6,1.9)} idam iha sampradhāryam dīrghatvam kriyatām dvirvacanam iti kim atra kartavyam . \\
\noindent \textbf{(P\_6,1.9)} paratvāt dīrghatvam . \\
\noindent \textbf{(P\_6,1.9)} nityam dvirvacanam . \\
\noindent \textbf{(P\_6,1.9)} kṛte api dīrghatve prāpnoti akṛte api prāpnoti . \\
\noindent \textbf{(P\_6,1.9)} dīrghatvam api nityam . \\
\noindent \textbf{(P\_6,1.9)} kṛte api dvirvacane prāpnoti akṛte api prāpnoti . \\
\noindent \textbf{(P\_6,1.9)} anityam dīrghatvam . \\
\noindent \textbf{(P\_6,1.9)} na hi kṛte dvirvacane prāpnoti . \\
\noindent \textbf{(P\_6,1.9)} kim kāraṇam . \\
\noindent \textbf{(P\_6,1.9)} samudāyasya samudāyaḥ ādeśaḥ . \\
\noindent \textbf{(P\_6,1.9)} tatra sampramugdhatvāt prakṛtipratyayasya ajantatā na asti iti dīrghatvam na prāpnoti . \\
\noindent \textbf{(P\_6,1.9)} dvirvacanam api anityam . \\
\noindent \textbf{(P\_6,1.9)} anyasya kṛte dīrghatve prāpnoti anyasya akṛte . \\
\noindent \textbf{(P\_6,1.9)} śabdāntarasya ca prāpnuvan vidhiḥ anityaḥ bhavati . \\
\noindent \textbf{(P\_6,1.9)} ubhayoḥ anityayoḥ paratvāt dīrghatvam . \\
\noindent \textbf{(P\_6,1.9)} yat tarhi na akṛte dvirvacane dīrghatvam tat na sidhyati . \\
\noindent \textbf{(P\_6,1.9)} juhūṣati iti . \\
\noindent \textbf{(P\_6,1.9)} kutvam dvirvacanādhikasya na sidhyati . \\
\noindent \textbf{(P\_6,1.9)} jighāṃsati . \\
\noindent \textbf{(P\_6,1.9)} jaṅghanyate . \\
\noindent \textbf{(P\_6,1.9)} kim kāraṇam . \\
\noindent \textbf{(P\_6,1.9)} samudāyasya samudāyaḥ ādeśaḥ . \\
\noindent \textbf{(P\_6,1.9)} tatra sampramugdhatvāt prakṛtipratyayasya naṣṭaḥ hantiḥ bhavati . \\
\noindent \textbf{(P\_6,1.9)} tatra abhyāsāt hantihakārasya iti kutvam na sidhyati . \\
\noindent \textbf{(P\_6,1.9)} samprasāraṇam ca dvirvacanādhikasya na sidhyati . \\
\noindent \textbf{(P\_6,1.9)} juhūṣati. johūyate . \\
\noindent \textbf{(P\_6,1.9)} samudāyasya samudāyaḥ ādeśaḥ . \\
\noindent \textbf{(P\_6,1.9)} tatra sampramugdhatvāt prakṛtipratyayasya naṣṭaḥ havayatiḥ bhavati . \\
\noindent \textbf{(P\_6,1.9)} tatra hvaḥ samprasāraṇam abhyastasya iti samprasāraṇam na prāpnoti . \\
\noindent \textbf{(P\_6,1.9)} na eṣaḥ doṣaḥ . \\
\noindent \textbf{(P\_6,1.9)} vakṣyati hi etat hvaḥ abhyastanimittasya iti . \\
\noindent \textbf{(P\_6,1.9)} yāvatā ca idānīm hvaḥ abhyastanimittasya iti ucyate saḥ api adoṣaḥ bhavati yat uktam yat tarhi na akṛte dvirvacane dīrghatvam tat na sidhyati iti . \\
\noindent \textbf{(P\_6,1.9)} ṣatvam ca dvirvacanādhikasya na sidhyati . \\
\noindent \textbf{(P\_6,1.9)} pipakṣati . \\
\noindent \textbf{(P\_6,1.9)} yiyakṣati . \\
\noindent \textbf{(P\_6,1.9)} samudāyasya samudāyaḥ ādeśaḥ . \\
\noindent \textbf{(P\_6,1.9)} tatra sampramugdhatvāt prakṛtipratyayasya naṣṭaḥ san bhavati . \\
\noindent \textbf{(P\_6,1.9)} tatra iṇkubhyām uttarasya pratyayasakārasya iti ṣatvam na prāpoti . \\
\noindent \textbf{(P\_6,1.9)} idam iha sampradhāryam dvirvacanam kriyatām ṣatvam iti kim atra kartavyam . \\
\noindent \textbf{(P\_6,1.9)} paratvāt ṣatvam . \\
\noindent \textbf{(P\_6,1.9)} pūrvatrāsiddhe ṣatvam siddhāsiddhayoḥ ca na asti sampradhāraṇā . \\
\noindent \textbf{(P\_6,1.9)} ābṛdhyoḥ ca abhyastavidhipratiṣedhaḥ . \\
\noindent \textbf{(P\_6,1.9)} ābṛdhyoḥ ca abhyastāśrayaḥ vidhiḥ prāpnoti . \\
\noindent \textbf{(P\_6,1.9)} saḥ pratiṣedhyaḥ . \\
\noindent \textbf{(P\_6,1.9)} īpsati . \\
\noindent \textbf{(P\_6,1.9)} īrtsati . \\
\noindent \textbf{(P\_6,1.9)} īpsan . \\
\noindent \textbf{(P\_6,1.9)} īrtsan . \\
\noindent \textbf{(P\_6,1.9)} aipsan . \\
\noindent \textbf{(P\_6,1.9)} airtsan . \\
\noindent \textbf{(P\_6,1.9)} kim ca syāt . \\
\noindent \textbf{(P\_6,1.9)} adbhāvaḥ numpratiṣedhaḥ jusbhāvaḥ iti ete vidhayaḥ prasajyeran . \\
\noindent \textbf{(P\_6,1.9)} na eṣaḥ doṣaḥ . \\
\noindent \textbf{(P\_6,1.9)} uktāḥ atra parihārāḥ . \\
\noindent \textbf{(P\_6,1.9)} saṅāśraye ca samudāyasya samudāyādeśatvāt jhalāśraye ca avyapadeśaḥ āmiśratvāt . \\
\noindent \textbf{(P\_6,1.9)} saṅāśraye ca kārye samudāyasya samudāyādeśatvāt jhalāśraye ca avyapadeśaḥ . \\
\noindent \textbf{(P\_6,1.9)} kim kāraṇam . \\
\noindent \textbf{(P\_6,1.9)} āmiśratvāt . \\
\noindent \textbf{(P\_6,1.9)} āmiśrībhūtam idam bhavati . \\
\noindent \textbf{(P\_6,1.9)} tat yathā . \\
\noindent \textbf{(P\_6,1.9)} kṣīrodake sampṛkte . \\
\noindent \textbf{(P\_6,1.9)} āmiśratvāt na jñāyate kiyat kṣīram kiyat udakam iti . \\
\noindent \textbf{(P\_6,1.9)} kasmin avakāśe kṣīram kasmin avakāśe udakam iti . \\
\noindent \textbf{(P\_6,1.9)} evam iha api āmiśratvāt na jñāyate kā prakṛtiḥ kaḥ pratyayaḥ kasmin avakāśe prakṛtiḥ kasmin avakāśe pratyayaḥ iti . \\
\noindent \textbf{(P\_6,1.9)} tatra kaḥ doṣaḥ . \\
\noindent \textbf{(P\_6,1.9)} saṅi jhali iti kutvādīni na sidhyanti . \\
\noindent \textbf{(P\_6,1.9)} idam iha sampradhāryam dvirvacanam kriyatām kutvādīni iti kim atra kartavyam . \\
\noindent \textbf{(P\_6,1.9)} paratvāt kutvādīni . \\
\noindent \textbf{(P\_6,1.9)} pūrvatrāsiddhe kutvādīni siddhāsiddhayoḥ ca na asti sampradhāraṇā . \\
\noindent \textbf{(P\_6,1.9)} evam tarhi pūrvatrāsiddhīyam advirvacane iti vaktavyam. tat ca avaśyam vaktavyam . \\
\noindent \textbf{(P\_6,1.9)} vibhāṣitāḥ prayojayanti . \\
\noindent \textbf{(P\_6,1.9)} drogdhā drogdhā . \\
\noindent \textbf{(P\_6,1.9)} droḍhā droḍhā . \\
\noindent \textbf{(P\_6,1.9)} yāvatā ca idānīm pūrvatrāsiddhīyam advirvacane iti ucyate saḥ api adoṣaḥ bhavati yat uktam ṣatvam na sidhyati . \\
\noindent \textbf{(P\_6,1.9)} iha sthāne dvirvacane ṇilopaḥ aparihṛtaḥ . \\
\noindent \textbf{(P\_6,1.9)} sanyaṅoḥ parataḥ dvirvacane iṭaḥ dvirvacanam vaktavyam . \\
\noindent \textbf{(P\_6,1.9)} sanyaṅantasya dvirvacane hanteḥ kutvam aparihṛtam . \\
\noindent \textbf{(P\_6,1.9)} tatra sanyaṅantasya dvirvacanam dviḥprayogaḥ ca iti eṣaḥ pakṣaḥ nirdoṣaḥ . \\
\noindent \textbf{(P\_6,1.9)} tatra idam aparihṛtam sanaḥ iṭaḥ pratiṣedhaḥ iti . \\
\noindent \textbf{(P\_6,1.9)} etasya api parihāram vakṣyati ubhayaviśeṣaṇatvāt siddham iti . \\
\noindent \textbf{(P\_6,1.9)} katham jeghnīyate . \\
\noindent \textbf{(P\_6,1.9)} vakṣyati etat yaṅprakaraṇe hanteḥ hiṃsāyām ghnī iti . \\
\noindent \textbf{(P\_6,1.12.1)} dāśvān iti kim nipātyate . \\
\noindent \textbf{(P\_6,1.12.1)} dāśeḥ vasau dvitveṭpratiṣedhau . \\
\noindent \textbf{(P\_6,1.12.1)} dāśeḥ vasau dvitveṭpratiṣedhau nipātyete . \\
\noindent \textbf{(P\_6,1.12.1)} dāśvaṃsaḥ dāśuṣaḥ sutam . \\
\noindent \textbf{(P\_6,1.12.1)} dāśvān . \\
\noindent \textbf{(P\_6,1.12.1)} sāhvān iti kim nipātyate . \\
\noindent \textbf{(P\_6,1.12.1)} saheḥ dīrghatvam ca . \\
\noindent \textbf{(P\_6,1.12.1)} kim ca . \\
\noindent \textbf{(P\_6,1.12.1)} dvitveṭpratiṣedhau ca . \\
\noindent \textbf{(P\_6,1.12.1)} sāhvān balāhakaḥ . \\
\noindent \textbf{(P\_6,1.12.1)} sāhvān . \\
\noindent \textbf{(P\_6,1.12.1)} mīḍhvān iti kim nipātyate . \\
\noindent \textbf{(P\_6,1.12.1)} miheḥ ḍhatvam ca . \\
\noindent \textbf{(P\_6,1.12.1)} kim ca . \\
\noindent \textbf{(P\_6,1.12.1)} yat ca pūrvayoḥ . \\
\noindent \textbf{(P\_6,1.12.1)} kim ca pūrvayoḥ . \\
\noindent \textbf{(P\_6,1.12.1)} dvitveṭpratiṣedhau dīrghatvam ca . \\
\noindent \textbf{(P\_6,1.12.1)} mīḍhvaḥ tokaya tanayāya mṛḍaya . \\
\noindent \textbf{(P\_6,1.12.1)} yathā iyam indra mīḍhvaḥ . \\
\noindent \textbf{(P\_6,1.12.1)} mahyarthaḥ vai gamyate . \\
\noindent \textbf{(P\_6,1.12.1)} kaḥ punaḥ mahyarthaḥ . \\
\noindent \textbf{(P\_6,1.12.1)} mahatiḥ dānakarmā . \\
\noindent \textbf{(P\_6,1.12.1)} ataḥ kim . \\
\noindent \textbf{(P\_6,1.12.1)} itvam api nipātyam . \\
\noindent \textbf{(P\_6,1.12.1)} mahyarthaḥ iti cet miheḥ tadarthatvāt siddham . \\
\noindent \textbf{(P\_6,1.12.1)} mahyarthaḥ iti cet mihiḥ api mahyarthe vartate . \\
\noindent \textbf{(P\_6,1.12.1)} katham punaḥ anyaḥ nāma anyasya arthe vartate . \\
\noindent \textbf{(P\_6,1.12.1)} bahvarthāḥ api dhātavaḥ bhavanti iti . \\
\noindent \textbf{(P\_6,1.12.1)} asti punaḥ anyatra api kva cit mihiḥ mahyarthe vartate . \\
\noindent \textbf{(P\_6,1.12.1)} asti iti āha . \\
\noindent \textbf{(P\_6,1.12.1)} miheḥ meghaḥ . \\
\noindent \textbf{(P\_6,1.12.1)} meghaḥ ca kasmāt bhavati . \\
\noindent \textbf{(P\_6,1.12.1)} apaḥ dadāti iti . \\
\noindent \textbf{(P\_6,1.12.2)} dvirvacanaprakaraṇe kṛñādīnām ke . \\
\noindent \textbf{(P\_6,1.12.2)} dvirvacanaprakaraṇe kṛñādīnām ke upasaṅkhyānam kartavyam . \\
\noindent \textbf{(P\_6,1.12.2)} cakram . \\
\noindent \textbf{(P\_6,1.12.2)} ciklidam . \\
\noindent \textbf{(P\_6,1.12.2)} caknam iti . \\
\noindent \textbf{(P\_6,1.12.2)} kādiṣu iti vaktavyam iha api yathā syāt . \\
\noindent \textbf{(P\_6,1.12.2)} babhruḥ . \\
\noindent \textbf{(P\_6,1.12.2)} yayuḥ iti . \\
\noindent \textbf{(P\_6,1.12.2)} caricalipativadīnām aci āk ca abhyāsasya . \\
\noindent \textbf{(P\_6,1.12.2)} caricalipativadīnām aci dve bhavataḥ iti vaktavyam āk ca abhyāsasya . \\
\noindent \textbf{(P\_6,1.12.2)} carācaraḥ . \\
\noindent \textbf{(P\_6,1.12.2)} calācalaḥ . \\
\noindent \textbf{(P\_6,1.12.2)} patāpataḥ . \\
\noindent \textbf{(P\_6,1.12.2)} vadāvadaḥ . \\
\noindent \textbf{(P\_6,1.12.2)} hanteḥ ghaḥ ca . \\
\noindent \textbf{(P\_6,1.12.2)} hanteḥ ghaḥ ca vaktavyaḥ . \\
\noindent \textbf{(P\_6,1.12.2)} aci dve bhavataḥ āk ca abhyāsasya . \\
\noindent \textbf{(P\_6,1.12.2)} ghanāghanaḥ . \\
\noindent \textbf{(P\_6,1.12.2)} pāṭeḥ ṇiluk ca dīrghaḥ ca abhyāsasya ūk ca . \\
\noindent \textbf{(P\_6,1.12.2)} pāṭayateḥ ṇiluk ca vaktavyaḥ . \\
\noindent \textbf{(P\_6,1.12.2)} aci dve bhavataḥ iti vaktavyam . \\
\noindent \textbf{(P\_6,1.12.2)} dīrghaḥ ca abhyāsasya ūk ca āgamaḥ . \\
\noindent \textbf{(P\_6,1.12.2)} pāṭupaṭaḥ . \\
\noindent \textbf{(P\_6,1.12.3)} dvirvacanam yaṇayavāyāvādeśāllopopadhālopaṇilopakikinoruttvebhyaḥ . \\
\noindent \textbf{(P\_6,1.12.3)} yaṇayavāyāvādeśāllopopadhālopaṇilopakikinoruttvebhyaḥ dvirvacanam bhavati vipratiṣedhena . \\
\noindent \textbf{(P\_6,1.12.3)} dvirvacanasya avakāśaḥ . \\
\noindent \textbf{(P\_6,1.12.3)} bibhidatuḥ . \\
\noindent \textbf{(P\_6,1.12.3)} bibhiduḥ . \\
\noindent \textbf{(P\_6,1.12.3)} yaṇādeśasya avakāśaḥ . \\
\noindent \textbf{(P\_6,1.12.3)} dadhi atra . \\
\noindent \textbf{(P\_6,1.12.3)} madhu atra . \\
\noindent \textbf{(P\_6,1.12.3)} iha ubhayam prāpnoti . \\
\noindent \textbf{(P\_6,1.12.3)} cakratuḥ . \\
\noindent \textbf{(P\_6,1.12.3)} cakruḥ . \\
\noindent \textbf{(P\_6,1.12.3)} ayavāyāvādeśānām avakāśaḥ . \\
\noindent \textbf{(P\_6,1.12.3)} cayanam . \\
\noindent \textbf{(P\_6,1.12.3)} cāyakaḥ . \\
\noindent \textbf{(P\_6,1.12.3)} lavanam . \\
\noindent \textbf{(P\_6,1.12.3)} lāvakaḥ . \\
\noindent \textbf{(P\_6,1.12.3)} dvirvacanasya saḥ eva . \\
\noindent \textbf{(P\_6,1.12.3)} iha ubhayam prāpnoti . \\
\noindent \textbf{(P\_6,1.12.3)} cicāya . \\
\noindent \textbf{(P\_6,1.12.3)} cicayitha . \\
\noindent \textbf{(P\_6,1.12.3)} lulāva . \\
\noindent \textbf{(P\_6,1.12.3)} lulavitha . \\
\noindent \textbf{(P\_6,1.12.3)} āllopasya avakāśaḥ . \\
\noindent \textbf{(P\_6,1.12.3)} . godaḥ . \\
\noindent \textbf{(P\_6,1.12.3)} kambaladaḥ . \\
\noindent \textbf{(P\_6,1.12.3)} dvirvacanasya saḥ eva . \\
\noindent \textbf{(P\_6,1.12.3)} iha ubhayam prāpnoti . \\
\noindent \textbf{(P\_6,1.12.3)} yayatuḥ . \\
\noindent \textbf{(P\_6,1.12.3)} yayuḥ . \\
\noindent \textbf{(P\_6,1.12.3)} tasthatuḥ . \\
\noindent \textbf{(P\_6,1.12.3)} tasthuḥ . \\
\noindent \textbf{(P\_6,1.12.3)} upadhālopāsya avakāśaḥ . \\
\noindent \textbf{(P\_6,1.12.3)} śleṣmaghnam madhu . \\
\noindent \textbf{(P\_6,1.12.3)} pittaghnam ghṛtam . \\
\noindent \textbf{(P\_6,1.12.3)} dvirvacanasya saḥ eva . \\
\noindent \textbf{(P\_6,1.12.3)} iha ubhayam prāpnoti . \\
\noindent \textbf{(P\_6,1.12.3)} āṭitat . \\
\noindent \textbf{(P\_6,1.12.3)} āśiśat . \\
\noindent \textbf{(P\_6,1.12.3)} uttvasya avakāśaḥ nipūrtāḥ piṇḍāḥ . \\
\noindent \textbf{(P\_6,1.12.3)} dvirvacanasya saḥ eva . \\
\noindent \textbf{(P\_6,1.12.3)} iha ubhayam prāpnoti . \\
\noindent \textbf{(P\_6,1.12.3)} mitrātvaruṇau taturiḥ . \\
\noindent \textbf{(P\_6,1.12.3)} dūre hyadhvā jaguriḥ . \\
\noindent \textbf{(P\_6,1.12.3)} dvirvacanam bhavati pūrvavipratiṣedhena . \\
\noindent \textbf{(P\_6,1.12.3)} saḥ tarhi pūrvavipratiṣedhaḥ vaktavyaḥ . \\
\noindent \textbf{(P\_6,1.12.3)} na vaktavyaḥ . \\
\noindent \textbf{(P\_6,1.12.3)} iṣṭavācī paraśabdaḥ . \\
\noindent \textbf{(P\_6,1.12.3)} vipratiṣedhe param yat iṣṭam tat bhavati iti . \\
\noindent \textbf{(P\_6,1.12.3)} dvirvacanāt prasāraṇāttvadhātvādivikārarītvettvottvaguṇavṛddhividhayaḥ . \\
\noindent \textbf{(P\_6,1.12.3)} dvirvacanāt prasāraṇāttvadhātvādivikārarītvettvottvaguṇavṛddhividhayaḥ bhavanti vipratiṣedhena . \\
\noindent \textbf{(P\_6,1.12.3)} dvirvacanasya avakāśaḥ . \\
\noindent \textbf{(P\_6,1.12.3)} bibhidatuḥ . \\
\noindent \textbf{(P\_6,1.12.3)} bibhiduḥ . \\
\noindent \textbf{(P\_6,1.12.3)} samprasāraṇasya avakāśaḥ . \\
\noindent \textbf{(P\_6,1.12.3)} iṣṭam . \\
\noindent \textbf{(P\_6,1.12.3)} suptam . \\
\noindent \textbf{(P\_6,1.12.3)} iha ubhayam prāpnoti . \\
\noindent \textbf{(P\_6,1.12.3)} ījatuḥ . \\
\noindent \textbf{(P\_6,1.12.3)} ījuḥ . \\
\noindent \textbf{(P\_6,1.12.3)} na etat asti prayojanam . \\
\noindent \textbf{(P\_6,1.12.3)} astu atra dvirvacanam . \\
\noindent \textbf{(P\_6,1.12.3)} dvirvacane kṛte parasya rūpasya kiti iti bhaviṣyati pūrvasya liṭi abhyāsasya ubhayeṣām iti . \\
\noindent \textbf{(P\_6,1.12.3)} idam tarhi soṣupyate . \\
\noindent \textbf{(P\_6,1.12.3)} idam ca api udāharaṇam . \\
\noindent \textbf{(P\_6,1.12.3)} ījatuḥ , ījuḥ iti . \\
\noindent \textbf{(P\_6,1.12.3)} nanu ca uktam . \\
\noindent \textbf{(P\_6,1.12.3)} astu atra dvirvacanam . \\
\noindent \textbf{(P\_6,1.12.3)} dvirvacane kṛte parasya rūpasya kiti iti bhaviṣyati pūrvasya liṭi abhyāsasya ubhayeṣām iti . \\
\noindent \textbf{(P\_6,1.12.3)} na sidhyati . \\
\noindent \textbf{(P\_6,1.12.3)} na samprasāraṇe samprasāraṇam iti pratiṣedhaḥ prāpnoti . \\
\noindent \textbf{(P\_6,1.12.3)} akāreṇa vyavhitatvāt na bhaviṣyati . \\
\noindent \textbf{(P\_6,1.12.3)} ekādeśe kṛte na asti vyavadhānam . \\
\noindent \textbf{(P\_6,1.12.3)} ekādeśaḥ pūrvavidhau sthānivat bhavati iti sthānivadbhāvāt vyavadhānam eva . \\
\noindent \textbf{(P\_6,1.12.3)} evam tarhi samānāṅgagrahaṇam tatra codayiṣyati . \\
\noindent \textbf{(P\_6,1.12.3)} āttvasya avakāśaḥ . \\
\noindent \textbf{(P\_6,1.12.3)} glātā . \\
\noindent \textbf{(P\_6,1.12.3)} mlātā . \\
\noindent \textbf{(P\_6,1.12.3)} dvirvacanasya saḥ eva . \\
\noindent \textbf{(P\_6,1.12.3)} iha ubhayam prāpnoti . \\
\noindent \textbf{(P\_6,1.12.3)} jagle . \\
\noindent \textbf{(P\_6,1.12.3)} mamle . \\
\noindent \textbf{(P\_6,1.12.3)} dhātvādivikārāṇām avakāśaḥ . \\
\noindent \textbf{(P\_6,1.12.3)} namati . \\
\noindent \textbf{(P\_6,1.12.3)} siñcati . \\
\noindent \textbf{(P\_6,1.12.3)} dvirvacanasya saḥ eva . \\
\noindent \textbf{(P\_6,1.12.3)} iha ubhayam prāpnoti . \\
\noindent \textbf{(P\_6,1.12.3)} nanāma . \\
\noindent \textbf{(P\_6,1.12.3)} siseca . \\
\noindent \textbf{(P\_6,1.12.3)} sasnau . \\
\noindent \textbf{(P\_6,1.12.3)} rītvasya avakāśaḥ . \\
\noindent \textbf{(P\_6,1.12.3)} mātrīyati . \\
\noindent \textbf{(P\_6,1.12.3)} pitrīyati . \\
\noindent \textbf{(P\_6,1.12.3)} dvirvacanasya saḥ eva . \\
\noindent \textbf{(P\_6,1.12.3)} iha ubhayam prāpnoti . \\
\noindent \textbf{(P\_6,1.12.3)} cekrīyate . \\
\noindent \textbf{(P\_6,1.12.3)} jehrīyate . \\
\noindent \textbf{(P\_6,1.12.3)} ītvasya avakāśaḥ . \\
\noindent \textbf{(P\_6,1.12.3)} pīyate . \\
\noindent \textbf{(P\_6,1.12.3)} gīyate . \\
\noindent \textbf{(P\_6,1.12.3)} dvirvacanasya saḥ eva . \\
\noindent \textbf{(P\_6,1.12.3)} iha ubhayam prāpnoti . \\
\noindent \textbf{(P\_6,1.12.3)} pepīyate . \\
\noindent \textbf{(P\_6,1.12.3)} jegīyate . \\
\noindent \textbf{(P\_6,1.12.3)} ittvottvayoḥ avakāśaḥ . \\
\noindent \textbf{(P\_6,1.12.3)} āstīrṇam . \\
\noindent \textbf{(P\_6,1.12.3)} nipūrtāḥ . \\
\noindent \textbf{(P\_6,1.12.3)} dvirvacanasya saḥ eva . \\
\noindent \textbf{(P\_6,1.12.3)} iha ubhayam prāpnoti . \\
\noindent \textbf{(P\_6,1.12.3)} ātestīryate . \\
\noindent \textbf{(P\_6,1.12.3)} nipopūryate . \\
\noindent \textbf{(P\_6,1.12.3)} guṇavṛddhyoḥ avakāśaḥ . \\
\noindent \textbf{(P\_6,1.12.3)} cetā . \\
\noindent \textbf{(P\_6,1.12.3)} gauḥ . \\
\noindent \textbf{(P\_6,1.12.3)} dvirvacanasya saḥ eva . \\
\noindent \textbf{(P\_6,1.12.3)} iha ubhayam prāpnoti . \\
\noindent \textbf{(P\_6,1.12.3)} cicāya . \\
\noindent \textbf{(P\_6,1.12.3)} cicayitha . \\
\noindent \textbf{(P\_6,1.12.3)} lulāva . \\
\noindent \textbf{(P\_6,1.12.3)} lulavitha . \\
\noindent \textbf{(P\_6,1.12.3)} na etat asti prayojanam . \\
\noindent \textbf{(P\_6,1.12.3)} astu atra dvirvacanam . \\
\noindent \textbf{(P\_6,1.12.3)} dvirvacane kṛte parasya rūpasya guṇavṛddhī bhaviṣyataḥ . \\
\noindent \textbf{(P\_6,1.12.3)} idam tarhi prayojanam . \\
\noindent \textbf{(P\_6,1.12.3)} iyāya . \\
\noindent \textbf{(P\_6,1.12.3)} iyayitha . \\
\noindent \textbf{(P\_6,1.12.3)} nanu ca uktam na etat asti prayojanam . \\
\noindent \textbf{(P\_6,1.12.3)} astu atra dvirvacanam . \\
\noindent \textbf{(P\_6,1.12.3)} dvirvacane kṛte parasya rūpasya guṇavṛddhī bhaviṣyataḥ . \\
\noindent \textbf{(P\_6,1.12.3)} na sidhyati . \\
\noindent \textbf{(P\_6,1.12.3)} antaraṅgatvāt savarṇadīrghatvam prāpnoti . \\
\noindent \textbf{(P\_6,1.12.3)} vārṇāt āṅgam balīyaḥ iti guṇavṛddhī bhaviṣyataḥ . \\
\noindent \textbf{(P\_6,1.12.3)} kim vaktavyam etat . \\
\noindent \textbf{(P\_6,1.12.3)} na hi . \\
\noindent \textbf{(P\_6,1.12.3)} katham anucyamānam gaṃsyate . \\
\noindent \textbf{(P\_6,1.12.3)} ācāryapravṛttiḥ jñāpayati vārṇāt āṅgam balīyaḥ bhavati iti yat ayam abhyāsasya asavarṇe iti asavarṇagrahaṇam karoti . \\
\noindent \textbf{(P\_6,1.12.3)} katham kṛtvā jñāpakam . \\
\noindent \textbf{(P\_6,1.12.3)} na hi antareṇa guṇavṛddhī asavarṇaparaḥ abhyāsaḥ bhavati . \\
\noindent \textbf{(P\_6,1.12.3)} na etat asti jñāpakam . \\
\noindent \textbf{(P\_6,1.12.3)} artyartham etat syāt . \\
\noindent \textbf{(P\_6,1.12.3)} iyṛtaḥ . \\
\noindent \textbf{(P\_6,1.12.3)} iyṛthaḥ . \\
\noindent \textbf{(P\_6,1.12.3)} [uvoṇa . \\
\noindent \textbf{(P\_6,1.12.3)} uvoṇithaḥ (R)] . \\
\noindent \textbf{(P\_6,1.12.3)} yat tarhi dīrghaḥ iṇaḥ kiti iti dīrghatvam śāsti . \\
\noindent \textbf{(P\_6,1.12.3)} etasya api asti vacane prayojanam . \\
\noindent \textbf{(P\_6,1.12.3)} kim . \\
\noindent \textbf{(P\_6,1.12.3)} savarṇadīrghabādhanārtham etat syāt . \\
\noindent \textbf{(P\_6,1.12.3)} saḥ yathā eva tarhi savarṇadīrghatvam bādhate evam yaṇādeśam api bādheta . \\
\noindent \textbf{(P\_6,1.12.3)} evam tarhi yaṇādeśe yogavibhāgaḥ kariṣyate . \\
\noindent \textbf{(P\_6,1.12.3)} idam asti iṇaḥ yaṇ bhavati . \\
\noindent \textbf{(P\_6,1.12.3)} tataḥ eḥ anekācaḥ . \\
\noindent \textbf{(P\_6,1.12.3)} eḥ ca anekācaḥ iṇaḥ yaṇ bhavati . \\
\noindent \textbf{(P\_6,1.12.3)} tataḥ asaṃyogapūrvasya . \\
\noindent \textbf{(P\_6,1.12.3)} eḥ anekācaḥ iti eva . \\
\noindent \textbf{(P\_6,1.12.3)} asavarṇagrahaṇam eva tarhi jñāpakam . \\
\noindent \textbf{(P\_6,1.12.3)} nanu ca uktam artyartham etat syāt iti . \\
\noindent \textbf{(P\_6,1.12.3)} na ekam udāharaṇam asavarṇagrahaṇam prayojayati . \\
\noindent \textbf{(P\_6,1.12.3)} evam api sthānivadbhāvāt iyaṅ na prāpnoti . \\
\noindent \textbf{(P\_6,1.12.3)} atha sati api vipratiṣedhe yāvatā sthānivadbhāvaḥ katham eva etat sidhyati . \\
\noindent \textbf{(P\_6,1.12.3)} yaḥ anādiṣṭāt acaḥ pūrvaḥ tasya vidhim prati sthānivadbhāvaḥ . \\
\noindent \textbf{(P\_6,1.12.3)} ādiṣṭāt ca eṣaḥ acaḥ pūrvaḥ bhavati . \\
\noindent \textbf{(P\_6,1.13.1)} ṣyaṅaḥ samprasāraṇe putrapatyoḥ tadādau atiprasaṅgaḥ . \\
\noindent \textbf{(P\_6,1.13.1)} ṣyaṅaḥ samprasāraṇe putrapatyoḥ tadādau atiprasaṅgaḥ bhavati . \\
\noindent \textbf{(P\_6,1.13.1)} putrapatyādau samprasāraṇam prāpnoti . \\
\noindent \textbf{(P\_6,1.13.1)} kārīṣagandhyāputrakulam , kārīṣagandhyāpatikulam . \\
\noindent \textbf{(P\_6,1.13.1)} varṇagrahaṇāt siddham . \\
\noindent \textbf{(P\_6,1.13.1)} varṇagrahaṇe etat bhavati yasmin vidhiḥ tadādau iti na ca idam varṇagrahaṇam . \\
\noindent \textbf{(P\_6,1.13.1)} varṇagrahaṇe iti cet tadantapratiṣedhaḥ . \\
\noindent \textbf{(P\_6,1.13.1)} varṇagrahaṇe iti cet tadantasya pratiṣedhaḥ vaktavyaḥ . \\
\noindent \textbf{(P\_6,1.13.1)} putrapatyante samprasāraṇam prāpnoti . \\
\noindent \textbf{(P\_6,1.13.1)} kārīṣagandhyāparamaputraḥ , kārīṣagandhyāparamapatiḥ . \\
\noindent \textbf{(P\_6,1.13.1)} kaumudagandhyāparamaputraḥ , kaumudagandhyāparamapatiḥ . \\
\noindent \textbf{(P\_6,1.13.1)} kim kāraṇam . \\
\noindent \textbf{(P\_6,1.13.1)} yatra hi tadādividhiḥ na asti tadantavidhinā tatra bhavitavyam . \\
\noindent \textbf{(P\_6,1.13.1)} siddham tu uttarapadavacanāt . \\
\noindent \textbf{(P\_6,1.13.1)} siddham etat . \\
\noindent \textbf{(P\_6,1.13.1)} katham . \\
\noindent \textbf{(P\_6,1.13.1)} uttarapadavacanāt . \\
\noindent \textbf{(P\_6,1.13.1)} putrapatyoḥ uttarapadayoḥ iti vaktavyam . \\
\noindent \textbf{(P\_6,1.13.1)} tat tarhi vaktavyam . \\
\noindent \textbf{(P\_6,1.13.1)} na vaktavyam . \\
\noindent \textbf{(P\_6,1.13.1)} pūrvapadam uttarapadam iti sambandhiśabdau etau . \\
\noindent \textbf{(P\_6,1.13.1)} sati pūrvapade uttarapadam bhavati sati ca uttarapade pūrvapadam iti . \\
\noindent \textbf{(P\_6,1.13.1)} na ca atra putrapatī uttarapade . \\
\noindent \textbf{(P\_6,1.13.1)} iha api tarhi na prāpnoti . \\
\noindent \textbf{(P\_6,1.13.1)} kārīṣagandhīputraḥ , kārīṣagandhīpatiḥ iti . \\
\noindent \textbf{(P\_6,1.13.1)} kim kāraṇam . \\
\noindent \textbf{(P\_6,1.13.1)} pūrvapadam iti ucyate . \\
\noindent \textbf{(P\_6,1.13.1)} na hi atra ṣyaṅ pūrvapadam asti . \\
\noindent \textbf{(P\_6,1.13.1)} ṣyaṅantam etat pūrvapadam . \\
\noindent \textbf{(P\_6,1.13.1)} katham . \\
\noindent \textbf{(P\_6,1.13.1)} pratyayagrahaṇe yasmāt saḥ tadādeḥ grahaṇam bhavati . \\
\noindent \textbf{(P\_6,1.13.1)} yadi pratyayagrahaṇe yasmāt saḥ tadādeḥ grahaṇam bhavati iti ucyate paramakārīṣagandhīputraḥ , paramakārīṣagandhīpatiḥ iti na sidhyati . \\
\noindent \textbf{(P\_6,1.13.1)} pratyayagrahaṇe yasmāt saḥ tadādeḥ grahaṇam bhavati astrīpratyayena iti . \\
\noindent \textbf{(P\_6,1.13.1)} yadi astrīpratyayena iti ucyate atikrāntaḥ kārīṣagandhyām atikārīṣagandhyaḥ , tasya putraḥ atikārīṣagandhyaputraḥ , atikārīṣagandhyapatiḥ iti atra api prāpnoti . \\
\noindent \textbf{(P\_6,1.13.1)} astrīpratyayena anupasarjanena . \\
\noindent \textbf{(P\_6,1.13.1)} yaḥ hi upasarjanam strīpratyayaḥ bhavati eṣā tatra paribhāṣā pratyayagrahaṇe yasmāt saḥ tadādeḥ grahaṇam bhavati iti . \\
\noindent \textbf{(P\_6,1.13.2)} ṣyaṅante yāvantaḥ yaṇaḥ teṣām sarveṣām samprasāraṇam prāpnoti . \\
\noindent \textbf{(P\_6,1.13.2)} vārāhiputraḥ , tārṇakarṇīputraḥ . \\
\noindent \textbf{(P\_6,1.13.2)} tatra apratyayasthasya pratiṣedhaḥ vaktavyaḥ . \\
\noindent \textbf{(P\_6,1.13.2)} yathāgṛhītasya ādeśavacanāt apratyayasthe siddham . \\
\noindent \textbf{(P\_6,1.13.2)} nirdiśyamānasya ādeśāḥ bhavanti iti evam apratyayasthasya na bhaviṣyati . \\
\noindent \textbf{(P\_6,1.13.2)} anantyavikāre antyasadeśasya vā . \\
\noindent \textbf{(P\_6,1.13.2)} atha vā anantyavikāre antyasadeśasya kāryam bhavati iti eṣā paribhāṣā kartavyā . \\
\noindent \textbf{(P\_6,1.13.2)} kaḥ punaḥ atra viśeṣaḥ eṣā vā paribhāṣā kriyeta apratyayasthasya vā pratiṣedhaḥ ucyeta . \\
\noindent \textbf{(P\_6,1.13.2)} avaśyam eṣā paribhāṣā kartavyā . \\
\noindent \textbf{(P\_6,1.13.2)} bahūni etasyāḥ paribhāṣāyāḥ prayojanāni . \\
\noindent \textbf{(P\_6,1.13.2)} kāni . \\
\noindent \textbf{(P\_6,1.13.2)} prayojanam na samprasāraṇe samprasāraṇam . \\
\noindent \textbf{(P\_6,1.13.2)} na samprasāraṇe samprasāraṇam iti etat na vaktavyam bhavati . \\
\noindent \textbf{(P\_6,1.13.2)} katham vyadheḥ viddhaḥ iti . \\
\noindent \textbf{(P\_6,1.13.2)} anantyavikāre antyasadeśasya kāryam bhavati iti na doṣaḥ bhavati . \\
\noindent \textbf{(P\_6,1.13.2)} na etat asti prayojanam . \\
\noindent \textbf{(P\_6,1.13.2)} kriyate nyāse eva . \\
\noindent \textbf{(P\_6,1.13.2)} sāntamahataḥ dīrghatve . \\
\noindent \textbf{(P\_6,1.13.2)} sāntamahataḥ dīrghatve prayojanam . \\
\noindent \textbf{(P\_6,1.13.2)} payāṃsi, yaśāṃsi . \\
\noindent \textbf{(P\_6,1.13.2)} pa iti asya api prāpnoti . \\
\noindent \textbf{(P\_6,1.13.2)} anantyavikāre antyasadeśasya kāryam bhavati iti na doṣaḥ bhavati . \\
\noindent \textbf{(P\_6,1.13.2)} etat api na asti prayojanam . \\
\noindent \textbf{(P\_6,1.13.2)} nopadhāyāḥ iti tatra vartate . \\
\noindent \textbf{(P\_6,1.13.2)} evam api anāṃsi, manāṃsi iti atra api prāpnoti . \\
\noindent \textbf{(P\_6,1.13.2)} na eṣaḥ doṣaḥ . \\
\noindent \textbf{(P\_6,1.13.2)} sāntasaṃyogena nopadhām viśeṣayiṣyāmaḥ . \\
\noindent \textbf{(P\_6,1.13.2)} sāntasaṃyogasya nopādhāyāḥ iti . \\
\noindent \textbf{(P\_6,1.13.2)} evam api haṃsaśirāṃsi , dhvaṃsaśirāṃsai iti atra api prāpnoti . \\
\noindent \textbf{(P\_6,1.13.2)} na eṣaḥ doṣaḥ . \\
\noindent \textbf{(P\_6,1.13.2)} hammateḥ haṃsaḥ . \\
\noindent \textbf{(P\_6,1.13.2)} kaḥ punaḥ āha hammateḥ haṃsaḥ iti . \\
\noindent \textbf{(P\_6,1.13.2)} kim tarhi hanteḥ haṃsaḥ . \\
\noindent \textbf{(P\_6,1.13.2)} hanti adhvānam iti . \\
\noindent \textbf{(P\_6,1.13.2)} evam tarhi sarvanāmasthāne iti vartate . \\
\noindent \textbf{(P\_6,1.13.2)} sarvanāmasthānaparatayā sāntasaṃyogam viśeṣayiṣyāmaḥ . \\
\noindent \textbf{(P\_6,1.13.2)} sarvanāmasthānaparasya sāntasaṃyogasya nopādhāyāḥ iti . \\
\noindent \textbf{(P\_6,1.13.2)} ankārāntasya allope . \\
\noindent \textbf{(P\_6,1.13.2)} ankārāntasya allope prayojanam . \\
\noindent \textbf{(P\_6,1.13.2)} takṣṇā , takṣṇe iti . \\
\noindent \textbf{(P\_6,1.13.2)} ta iti atra api prāpnoti . \\
\noindent \textbf{(P\_6,1.13.2)} anantyavikāre antyasadeśasya kāryam bhavati iti na doṣaḥ bhavati . \\
\noindent \textbf{(P\_6,1.13.2)} etat api na asti prayojanam . \\
\noindent \textbf{(P\_6,1.13.2)} anā akāram viśeṣayiṣyāmaḥ . \\
\noindent \textbf{(P\_6,1.13.2)} anaḥ yaḥ akāraḥ iti . \\
\noindent \textbf{(P\_6,1.13.2)} evam api anasā , anase iti atra api prāpnoti . \\
\noindent \textbf{(P\_6,1.13.2)} ankāreṇa aṅgam viśeṣayiṣyāmaḥ . \\
\noindent \textbf{(P\_6,1.13.2)} ankārāntasya aṅgasya anaḥ yaḥ akāraḥ iti . \\
\noindent \textbf{(P\_6,1.13.2)} evam api anastakṣṇā , anastakṣṇe iti atra api prāpnoti . \\
\noindent \textbf{(P\_6,1.13.2)} evam tarhi kāryakālam sañjñāparibhāṣam . \\
\noindent \textbf{(P\_6,1.13.2)} yatra kāryam tatra upasthitam draṣṭavyam . \\
\noindent \textbf{(P\_6,1.13.2)} bhasya iti upasthitam idam bhavati yaci bham iti . \\
\noindent \textbf{(P\_6,1.13.2)} tatra yajādiparatyā ankāram viśeṣayiṣyāmaḥ anā akāram . \\
\noindent \textbf{(P\_6,1.13.2)} yajādiparasya anaḥ yaḥ akāraḥ iti . \\
\noindent \textbf{(P\_6,1.13.2)} mṛjeḥ vṛddhividhau . \\
\noindent \textbf{(P\_6,1.13.2)} mṛjeḥ vṛddhividhau prayojanam . \\
\noindent \textbf{(P\_6,1.13.2)} nyamārṭ . \\
\noindent \textbf{(P\_6,1.13.2)} aṭaḥ api vṛddhiḥ prāpnoti . \\
\noindent \textbf{(P\_6,1.13.2)} anantyavikāre antyasadeśasya kāryam bhavati iti na doṣaḥ bhavati . \\
\noindent \textbf{(P\_6,1.13.2)} etat api na asti prayojanam . \\
\noindent \textbf{(P\_6,1.13.2)} yathāparibhāṣitam ikaḥ guṇavṛddhī iti ikaḥ eva vṛddhiḥ bhaviṣyati . \\
\noindent \textbf{(P\_6,1.13.2)} evam api mimārjiṣati iti atra prāpnoti . \\
\noindent \textbf{(P\_6,1.13.2)} astu . \\
\noindent \textbf{(P\_6,1.13.2)} abhyāsanirhrāsena hrasvaḥ bhaviṣyati . \\
\noindent \textbf{(P\_6,1.13.2)} vasoḥ samprasāraṇe ca . \\
\noindent \textbf{(P\_6,1.13.2)} vasoḥ samprasāraṇe ca prayojanam . \\
\noindent \textbf{(P\_6,1.13.2)} viduṣaḥ paśya . \\
\noindent \textbf{(P\_6,1.13.2)} vidivakārasya api prāpnoti . \\
\noindent \textbf{(P\_6,1.13.2)} anantyavikāre antyasadeśasya kāryam bhavati iti na doṣaḥ bhavati . \\
\noindent \textbf{(P\_6,1.13.2)} etat api na asti prayojanam . \\
\noindent \textbf{(P\_6,1.13.2)} na samprasāraṇe samprasāraṇam iti pratiṣedhaḥ bhaviṣyati . \\
\noindent \textbf{(P\_6,1.13.2)} dakāreṇa (R: idkāreṇa) vyavahitatvāt na prāpnoti . \\
\noindent \textbf{(P\_6,1.13.2)} evam tarhi nirdiśyamānasya ādeśāḥ bhavanti iti na bhaviṣyati . \\
\noindent \textbf{(P\_6,1.13.2)} yuvādīnām ca . \\
\noindent \textbf{(P\_6,1.13.2)} yuvādīnām ca samprasāraṇe prayojanam . \\
\noindent \textbf{(P\_6,1.13.2)} yūnaḥ , yūnā , yūne . \\
\noindent \textbf{(P\_6,1.13.2)} yakārasya api prāpnoti . \\
\noindent \textbf{(P\_6,1.13.2)} anantyavikāre antyasadeśasya kāryam bhavati iti na doṣaḥ bhavati . \\
\noindent \textbf{(P\_6,1.13.2)} etat api na asti prayojanam . \\
\noindent \textbf{(P\_6,1.13.2)} na samprasāraṇe samprasāraṇam iti na bhaviṣyati . \\
\noindent \textbf{(P\_6,1.13.2)} ukāreṇa vyavahitatvāt na prāpnoti . \\
\noindent \textbf{(P\_6,1.13.2)} ekādeśe kṛte na asti vyavadhānam . \\
\noindent \textbf{(P\_6,1.13.2)} ekādeśaḥ pūrvavidhau sthānivat bhavati iti sthānivadbhāvāt vyavadhānam eva . \\
\noindent \textbf{(P\_6,1.13.2)} evam tarhi samānāṅgagrahaṇam atra codayiṣyati . \\
\noindent \textbf{(P\_6,1.13.2)} rvoḥ upadhāgrahaṇam ca . \\
\noindent \textbf{(P\_6,1.13.2)} rvoḥ upadhāgrahaṇam ca na kartavyam bhavati . \\
\noindent \textbf{(P\_6,1.13.2)} iha kasmāt na bhavati . \\
\noindent \textbf{(P\_6,1.13.2)} abibhaḥ bhavān . \\
\noindent \textbf{(P\_6,1.13.2)} anantyavikāre antyasadeśasya kāryam bhavati iti na doṣaḥ bhavati . \\
\noindent \textbf{(P\_6,1.13.2)} etat api na asti prayojanam . \\
\noindent \textbf{(P\_6,1.13.2)} kriyate nyāse eva . \\
\noindent \textbf{(P\_6,1.13.2)} ādityadādividhisaṃyogādilopakutvaḍhatvabhaṣbhāvaṣatvaṇatveṣu atiprasaṅgaḥ . \\
\noindent \textbf{(P\_6,1.13.2)} ādividhau atiprasaṅgaḥ bhavati . \\
\noindent \textbf{(P\_6,1.13.2)} dhātvādeḥ ṣaḥ saḥ . \\
\noindent \textbf{(P\_6,1.13.2)} ṇaḥ naḥ . \\
\noindent \textbf{(P\_6,1.13.2)} iha eva syāt . \\
\noindent \textbf{(P\_6,1.13.2)} netā , sotā . \\
\noindent \textbf{(P\_6,1.13.2)} iha na syāt . \\
\noindent \textbf{(P\_6,1.13.2)} namati , siñcati . \\
\noindent \textbf{(P\_6,1.13.2)} ādi . \\
\noindent \textbf{(P\_6,1.13.2)} tyadādividhi . \\
\noindent \textbf{(P\_6,1.13.2)} iha eva syāt . \\
\noindent \textbf{(P\_6,1.13.2)} tat , saḥ . \\
\noindent \textbf{(P\_6,1.13.2)} tyat , syaḥ iti atra na syāt . \\
\noindent \textbf{(P\_6,1.13.2)} tyadādividhi . \\
\noindent \textbf{(P\_6,1.13.2)} saṃyogādilopa . \\
\noindent \textbf{(P\_6,1.13.2)} iha eva syāt . \\
\noindent \textbf{(P\_6,1.13.2)} maṅktā . \\
\noindent \textbf{(P\_6,1.13.2)} maṅktavyam iti atra na syāt . \\
\noindent \textbf{(P\_6,1.13.2)} saṃyogādilopa . \\
\noindent \textbf{(P\_6,1.13.2)} kutva . \\
\noindent \textbf{(P\_6,1.13.2)} iha eva syāt . \\
\noindent \textbf{(P\_6,1.13.2)} paktā . \\
\noindent \textbf{(P\_6,1.13.2)} paktavyambhaṣbhāva . \\
\noindent \textbf{(P\_6,1.13.2)} iti atra na syāt . \\
\noindent \textbf{(P\_6,1.13.2)} kutva . \\
\noindent \textbf{(P\_6,1.13.2)} ḍhatva . \\
\noindent \textbf{(P\_6,1.13.2)} iha eva syāt . \\
\noindent \textbf{(P\_6,1.13.2)} leḍhā . \\
\noindent \textbf{(P\_6,1.13.2)} leḍhavyam iti atra na syāt . \\
\noindent \textbf{(P\_6,1.13.2)} ḍhatva . \\
\noindent \textbf{(P\_6,1.13.2)} bhaṣbhāva . \\
\noindent \textbf{(P\_6,1.13.2)} iha eva syāt . \\
\noindent \textbf{(P\_6,1.13.2)} abhutsi . \\
\noindent \textbf{(P\_6,1.13.2)} abhutsātām iti atra na syāt . \\
\noindent \textbf{(P\_6,1.13.2)} bhaṣbhāva . \\
\noindent \textbf{(P\_6,1.13.2)} ṣatva . \\
\noindent \textbf{(P\_6,1.13.2)} iha eva syāt . \\
\noindent \textbf{(P\_6,1.13.2)} draṣṭā . \\
\noindent \textbf{(P\_6,1.13.2)} draṣṭavyam iti atra na syāt . \\
\noindent \textbf{(P\_6,1.13.2)} ṣatva . \\
\noindent \textbf{(P\_6,1.13.2)} ṇatva . \\
\noindent \textbf{(P\_6,1.13.2)} iha eva syāt . \\
\noindent \textbf{(P\_6,1.13.2)} māṣāvāpeṇa . \\
\noindent \textbf{(P\_6,1.13.2)} māṣāvāpāṇām iti atra na syāt . \\
\noindent \textbf{(P\_6,1.13.2)} ṇatva . \\
\noindent \textbf{(P\_6,1.13.2)} ete doṣāḥ samāḥ bhūyāṃsaḥ vā . \\
\noindent \textbf{(P\_6,1.13.2)} tasmāt na arthaḥ anayā paribhāṣayā . \\
\noindent \textbf{(P\_6,1.13.2)} na hi doṣā santi iti paribhāṣā na kartavyā lakṣaṇam vā na praṇeyam . \\
\noindent \textbf{(P\_6,1.13.2)} na hi bhikṣukāḥ santi iti sthālyaḥ na adhiśrīyante na ca mṛgāḥ santi iti yavāḥ na upyante . \\
\noindent \textbf{(P\_6,1.13.2)} doṣāḥ khalu api sākalyena parigaṇitāḥ prayojanānām udāharaṇamātram . \\
\noindent \textbf{(P\_6,1.13.2)} kuta etat . \\
\noindent \textbf{(P\_6,1.13.2)} na hi doṣāṇām lakṣaṇam asti . \\
\noindent \textbf{(P\_6,1.13.2)} tasmāt yāni etasyāḥ paribhāṣāyāḥ prayojanāni tadartham eṣā paribhāṣā kartavyā pratividheyam ca doṣeṣu . \\
\noindent \textbf{(P\_6,1.13.2)} idam pratividhīyate . \\
\noindent \textbf{(P\_6,1.13.2)} udāttanirdeśāt siddham . \\
\noindent \textbf{(P\_6,1.13.2)} yatra eṣā paribhāṣā iṣyate tatra udāttanirdeśaḥ kartavyaḥ . \\
\noindent \textbf{(P\_6,1.13.2)} tataḥ vaktavyam anantyavikāre antyasadeśasya kāryam bhavati udāttanirdeśe iti . \\
\noindent \textbf{(P\_6,1.13.2)} saḥ tarhi udāttanirdeśḥ kartavyaḥ . \\
\noindent \textbf{(P\_6,1.13.2)} na kartavyaḥ . \\
\noindent \textbf{(P\_6,1.13.2)} yatra eva antyasadeśaḥ ca anantyasadeśaḥ ca yugapat samavasthitau tatra eṣā paribhāṣā bhavati . \\
\noindent \textbf{(P\_6,1.13.2)} doṣeṣu ca anyatra antyasadeśaḥ anyatra anantyasadeśaḥ . \\
\noindent \textbf{(P\_6,1.13.2)} prayojaneṣu punaḥ tatra eva antyasadeśaḥ ca anantyasadeśaḥ ca . \\
\noindent \textbf{(P\_6,1.13.2)} tathājātīyakāni khalu api ācāryeṇa prayojanāni paṭhitāni yāni ubhayavanti . \\
\noindent \textbf{(P\_6,1.13.2)} idam ekam yathā doṣaḥ tathā rvoḥ upadhāgrahaṇam iti . \\
\noindent \textbf{(P\_6,1.13.2)} tat ca api kriyate nyāse eva . \\
\noindent \textbf{(P\_6,1.14)} mātac . \\
\noindent \textbf{(P\_6,1.14)} kārīṣagandhyā mātā asya kārīṣagandhīmātaḥ , kārīṣagandhyāmātaḥ . \\
\noindent \textbf{(P\_6,1.14)} mātac . \\
\noindent \textbf{(P\_6,1.14)} mātṛka . \\
\noindent \textbf{(P\_6,1.14)} kārīṣagandhīmātṛkaḥ , kārīṣagandhyāmātṛkaḥ . \\
\noindent \textbf{(P\_6,1.14)} mātṛka . \\
\noindent \textbf{(P\_6,1.14)} mātṛ . \\
\noindent \textbf{(P\_6,1.14)} kārīṣagandhīmātā , kārīṣagandhyāmātā . \\
\noindent \textbf{(P\_6,1.14)} mātṛ . \\
\noindent \textbf{(P\_6,1.16)} vayigrahaṇam kimartham na veñ yajādiṣu paṭhyate veñaḥ ca vayiḥ ādeśaḥ kriyate tatra yajādīnām kiti iti eva siddham . \\
\noindent \textbf{(P\_6,1.16)} tatra etat syāt . \\
\noindent \textbf{(P\_6,1.16)} ṅidarthaḥ ayam ārambhaḥ iti . \\
\noindent \textbf{(P\_6,1.16)} tat ca na . \\
\noindent \textbf{(P\_6,1.16)} liṭi ayam ādeśaḥ liṭ ca kit eva . \\
\noindent \textbf{(P\_6,1.16)} ataḥ uttaram paṭhati . \\
\noindent \textbf{(P\_6,1.16)} vayigrahaṇam veñaḥ pratiṣedhāt . \\
\noindent \textbf{(P\_6,1.16)} vayigrahaṇam kriyate veñaḥ pratiṣedhāt . \\
\noindent \textbf{(P\_6,1.16)} veñaḥ liṭi pratiṣedham vakṣyati . \\
\noindent \textbf{(P\_6,1.16)} saḥ vayeḥ mā bhūt iti . \\
\noindent \textbf{(P\_6,1.16)} yathā eva hi veñgrahaṇāt vidhiḥ prārthyate evam pratiṣedhaḥ api prāpnoti . \\
\noindent \textbf{(P\_6,1.16)} na vā yakārapratiṣedhaḥ jñāpakaḥ apratiṣedhasya . \\
\noindent \textbf{(P\_6,1.16)} na vā eṣaḥ doṣaḥ . \\
\noindent \textbf{(P\_6,1.16)} kim kāraṇam . \\
\noindent \textbf{(P\_6,1.16)} yat ayam liṭi vayaḥ yaḥ iti vayeḥ yakārasya samprasāraṇapratiṣedham śāsti tat jñāpati ācāryaḥ na veñgrahaṇāt samprasāraṇapratiṣedhaḥ bhavati iti . \\
\noindent \textbf{(P\_6,1.16)} na etat asti jñāpakam . \\
\noindent \textbf{(P\_6,1.16)} piti abhyāsārtham etat syāt . \\
\noindent \textbf{(P\_6,1.16)} vayeḥ pitsu vacaneṣu abhyāsasya yakārasya samprasāraṇam mā bhūt iti . \\
\noindent \textbf{(P\_6,1.16)} nanu ca veñgrahaṇāt vayeḥ pitsu api vacaneṣu abhyāsayakārasya samprasāraṇapratiṣedhaḥ siddhaḥ . \\
\noindent \textbf{(P\_6,1.16)} na sidhyati . \\
\noindent \textbf{(P\_6,1.16)} kim kāraṇam . \\
\noindent \textbf{(P\_6,1.16)} kiti iti tatra anuvartate . \\
\noindent \textbf{(P\_6,1.16)} evam api vayeḥ pitsu api vacaneṣu abhyāsayakārasya samprasāraṇam na prāpnoti . \\
\noindent \textbf{(P\_6,1.16)} kim kāraṇam . \\
\noindent \textbf{(P\_6,1.16)} halādiśeṣeṇa bādhyate . \\
\noindent \textbf{(P\_6,1.16)} na atra halādiśeṣaḥ prāpnoti . \\
\noindent \textbf{(P\_6,1.16)} kim kāraṇam . \\
\noindent \textbf{(P\_6,1.16)} vakṣyati hi etat abhyāsasamprasāraṇam halādiśeṣāt vipratiṣedhena iti . \\
\noindent \textbf{(P\_6,1.16)} saḥ eṣaḥ vayeḥ yakārasya samprasāraṇapratiṣedhaḥ piti abhyāsārthaḥ na jñāpakārthaḥ bhavati . \\
\noindent \textbf{(P\_6,1.16)} piti abhyāsārtham iti cet na aviśiṣṭatvāt . \\
\noindent \textbf{(P\_6,1.16)} piti abhyāsārtham iti cet tat na . \\
\noindent \textbf{(P\_6,1.16)} kim kāraṇam . \\
\noindent \textbf{(P\_6,1.16)} aviśiṣṭatvāt . \\
\noindent \textbf{(P\_6,1.16)} aviśeṣeṇa pratiṣedhaḥ . \\
\noindent \textbf{(P\_6,1.16)} nivṛttam tatra kiti iti . \\
\noindent \textbf{(P\_6,1.16)} ātaḥ ca aviśeṣeṇa . \\
\noindent \textbf{(P\_6,1.16)} veñaḥ api hi pitsu vacaneṣu abhyāsasya samprasāraṇam na iṣyate . \\
\noindent \textbf{(P\_6,1.16)} vavau vavitha iti . \\
\noindent \textbf{(P\_6,1.16)} vikṛtigrahaṇam khalu api pratiṣedhe kriyate na ca vikṛtiḥ prakṛtim gṛhṇāti . \\
\noindent \textbf{(P\_6,1.17.1)} grahivṛścatipṛcchatibhṛjjatīnām aviśeṣaḥ . \\
\noindent \textbf{(P\_6,1.17.1)} yat ucyate vṛśceḥ aviśeṣaḥ iti tat na . \\
\noindent \textbf{(P\_6,1.17.1)} yadi atra rephasya samprasāraṇam na syāt vakārasya prasajyeta . \\
\noindent \textbf{(P\_6,1.17.1)} rephasya punaḥ samprasāraṇe sati uḥ adattvasya sthānivadbhāvāt na samprasāraṇe samprasāraṇam iti pratiṣedhaḥ siddhaḥ bhavati . \\
\noindent \textbf{(P\_6,1.17.1)} tasmāt vaktavyam graheḥ aviṣeṣaḥ pṛcchatibhṛjjatyoḥ aviśeṣaḥ iti . \\
\noindent \textbf{(P\_6,1.17.2)} atha ubhayagrahaṇam kimartham . \\
\noindent \textbf{(P\_6,1.17.2)} ubhayeṣām abhyāsasya samprasāraṇam yathā syāt vacisvapiyajādīnām grahādīnām ca . \\
\noindent \textbf{(P\_6,1.17.2)} na etat asti prayojanam . \\
\noindent \textbf{(P\_6,1.17.2)} prakṛtam ubhayeṣām grahaṇam anuvartate . \\
\noindent \textbf{(P\_6,1.17.2)} yadi anuvartate grahijyāvayivyadhivaṣṭivicativṛścatipṛcchatibhṛjjatīnām ṅiti ca iti yajādīnām ṅiti api prāpnoti . \\
\noindent \textbf{(P\_6,1.17.2)} na eṣaḥ doṣaḥ . \\
\noindent \textbf{(P\_6,1.17.2)} sambandham anuvartiṣyate . \\
\noindent \textbf{(P\_6,1.17.2)} vacisvapiyajādīnām kiti . \\
\noindent \textbf{(P\_6,1.17.2)} grahādīnām ṅiti ca vacisvapiyajādīnām kiti . \\
\noindent \textbf{(P\_6,1.17.2)} tataḥ liṭi abhyāsasya ubhayeṣām . \\
\noindent \textbf{(P\_6,1.17.2)} kiti ṅiti iti nivṛttam . \\
\noindent \textbf{(P\_6,1.17.2)} atha vā maṇḍūkagatayaḥ adhikārāḥ . \\
\noindent \textbf{(P\_6,1.17.2)} yathā maṇḍūkāḥ utplutya utplutya gacchanti tadvat adhikārāḥ . \\
\noindent \textbf{(P\_6,1.17.2)} atha vā ekayogaḥ kariṣyate . \\
\noindent \textbf{(P\_6,1.17.2)} vacisvapiyajādīnām kiti grahādīnām ṅiti ca iti . \\
\noindent \textbf{(P\_6,1.17.2)} tataḥ liṭi abhyāsasya iti . \\
\noindent \textbf{(P\_6,1.17.2)} na ca ekayoge anuvṛttiḥ bhavati . \\
\noindent \textbf{(P\_6,1.17.2)} atha vā ubhayam nivṛttam . \\
\noindent \textbf{(P\_6,1.17.2)} tat apekṣiṣyāmahe . \\
\noindent \textbf{(P\_6,1.17.2)} idam tarhi ubhayeṣāṅgrahaṇasya prayojanam . \\
\noindent \textbf{(P\_6,1.17.2)} ubhayeṣām abhyāsasya samprasāraṇam eva yathā syāt . \\
\noindent \textbf{(P\_6,1.17.2)} yat anyat prāpnoti tat mā bhūt iti . \\
\noindent \textbf{(P\_6,1.17.2)} kim ca anyat prāpnoti . \\
\noindent \textbf{(P\_6,1.17.2)} halādiśeṣaḥ . \\
\noindent \textbf{(P\_6,1.17.2)} abhyāsasamprasāraṇam halādiśeṣāt vipratiṣedhena iti vakṣyati . \\
\noindent \textbf{(P\_6,1.17.2)} saḥ pūrvavipratiṣedhaḥ na paṭhitavyaḥ bhavati . \\
\noindent \textbf{(P\_6,1.17.3)} abhyāsasamprasāraṇam halādiśeṣāt vipratiṣedhena . \\
\noindent \textbf{(P\_6,1.17.3)} abhyāsasamprasāraṇam halādiśeṣāt bhavati [bhavati halādiśeṣāt : R] vipratiṣedhena . \\
\noindent \textbf{(P\_6,1.17.3)} abhyāsasamprasāraṇasya avakāśaḥ : iyāja, uvāpa . \\
\noindent \textbf{(P\_6,1.17.3)} halādiśeṣasya avakāśaḥ : bibhidatuḥ , bibhiduḥ . \\
\noindent \textbf{(P\_6,1.17.3)} iha ubhayam prāpnoti vivyādha , vivyadhitha . \\
\noindent \textbf{(P\_6,1.17.3)} abhyāsasamprasāraṇam bhavati pūrvavipratiṣedhena . \\
\noindent \textbf{(P\_6,1.17.3)} saḥ tarhi pūrvavipratiṣedhaḥ vaktavyaḥ . \\
\noindent \textbf{(P\_6,1.17.3)} na vā samprasāraṇāśrayabalīyastvāt anyatra api . \\
\noindent \textbf{(P\_6,1.17.3)} na vā vaktavyaḥ . \\
\noindent \textbf{(P\_6,1.17.3)} kim kāraṇam . \\
\noindent \textbf{(P\_6,1.17.3)} samprasāraṇāśrayabalīyastvāt anyatra api . \\
\noindent \textbf{(P\_6,1.17.3)} samprasāraṇam samprasāraṇāśrayam ca balīyaḥ bhavati iti vaktavyam . \\
\noindent \textbf{(P\_6,1.17.3)} anyatra api na avaśyam iha eva vaktavyam . \\
\noindent \textbf{(P\_6,1.17.3)} kim prayojanam . \\
\noindent \textbf{(P\_6,1.17.3)} prayojanam ramāllopeyiaṅyaṇaḥ . \\
\noindent \textbf{(P\_6,1.17.3)} ram . \\
\noindent \textbf{(P\_6,1.17.3)} bhṛṣṭaḥ, bhṛṣṭavān . \\
\noindent \textbf{(P\_6,1.17.3)} samprasāraṇam ca prāpnoti rambhāvaḥ ca . \\
\noindent \textbf{(P\_6,1.17.3)} paratvāt rambhāvaḥ syāt . \\
\noindent \textbf{(P\_6,1.17.3)} samprasāraṇam balīyaḥ bhavati iti vaktavyam samprasāraṇam yathā syāt . \\
\noindent \textbf{(P\_6,1.17.3)} ram . \\
\noindent \textbf{(P\_6,1.17.3)} āllopaḥ . \\
\noindent \textbf{(P\_6,1.17.3)} juhuvatuḥ , juhuvuḥ . \\
\noindent \textbf{(P\_6,1.17.3)} samprasāraṇam ca prāpnoti āllopaḥ ca . \\
\noindent \textbf{(P\_6,1.17.3)} paratvāt āllopaḥ syāt . \\
\noindent \textbf{(P\_6,1.17.3)} samprasāraṇam balīyaḥ bhavati iti vaktavyam samprasāraṇam yathā syāt . \\
\noindent \textbf{(P\_6,1.17.3)} samprasāraṇe kṛte pūrvatvam ca prāpnoti āllopaḥ ca . \\
\noindent \textbf{(P\_6,1.17.3)} paratvāt āllopaḥ syāt . \\
\noindent \textbf{(P\_6,1.17.3)} samprasāraṇāśrayam balīyaḥ bhavati iti vaktavyam pūrvatvam yathā syāt . \\
\noindent \textbf{(P\_6,1.17.3)} iyaṅ . \\
\noindent \textbf{(P\_6,1.17.3)} śuśuvatuḥ , śuśuvuḥ . \\
\noindent \textbf{(P\_6,1.17.3)} samprasāraṇam ca prāpnoti iyaṅādeśaḥ ca . \\
\noindent \textbf{(P\_6,1.17.3)} paratvāt iyaṅādeśaḥ syāt . \\
\noindent \textbf{(P\_6,1.17.3)} samprasāraṇam balīyaḥ bhavati iti vaktavyam samprasāraṇam yathā syāt . \\
\noindent \textbf{(P\_6,1.17.3)} yaṇ . \\
\noindent \textbf{(P\_6,1.17.3)} samprasāraṇe kṛte pūrvatvam ca prāpnoti yaṇādeśaḥ ca . \\
\noindent \textbf{(P\_6,1.17.3)} paratvāt yaṇādeśaḥ syāt . \\
\noindent \textbf{(P\_6,1.17.3)} samprasāraṇāśrayam balīyaḥ bhavati iti vaktavyam pūrvatvam yathā syāt . \\
\noindent \textbf{(P\_6,1.17.3)} iyaṅ . \\
\noindent \textbf{(P\_6,1.17.3)} na etāni santi prayojanāni . \\
\noindent \textbf{(P\_6,1.17.3)} yat tāvat ucyate ram iti idam iha sampradhāryam : rambhāvaḥ kriyatām samprasāraṇam iti . \\
\noindent \textbf{(P\_6,1.17.3)} kim atra kartavyam . \\
\noindent \textbf{(P\_6,1.17.3)} paratvāt rambhāvaḥ . \\
\noindent \textbf{(P\_6,1.17.3)} nityam samprasāraṇam . \\
\noindent \textbf{(P\_6,1.17.3)} kṛte api rambhābe prāpnoti akṛte api . \\
\noindent \textbf{(P\_6,1.17.3)} rambhāvaḥ api nityaḥ . \\
\noindent \textbf{(P\_6,1.17.3)} kṛte api samprasāraṇe prāpnoti akṛte api . \\
\noindent \textbf{(P\_6,1.17.3)} katham . \\
\noindent \textbf{(P\_6,1.17.3)} yaḥ asau ṛkāre rephaḥ tasya ca upadhāyāḥ ca prāpnoti . \\
\noindent \textbf{(P\_6,1.17.3)} anityaḥ rambhāvaḥ . \\
\noindent \textbf{(P\_6,1.17.3)} na hi kṛte samprasāraṇe prāpnoti . \\
\noindent \textbf{(P\_6,1.17.3)} kim kāraṇam . \\
\noindent \textbf{(P\_6,1.17.3)} upadeśe iti vartate . \\
\noindent \textbf{(P\_6,1.17.3)} tat ca avaśyam upadeśagrahaṇam anuvartyam barībhṛjyate iti evamartham . \\
\noindent \textbf{(P\_6,1.17.3)} āllopeyaṅyaṇaḥ iti . \\
\noindent \textbf{(P\_6,1.17.3)} nityam samprasāraṇam . \\
\noindent \textbf{(P\_6,1.17.3)} antaraṅgam pūrvatvam . \\
\noindent \textbf{(P\_6,1.17.3)} tat etat ananyārtham samprasāraṇāśrayam balīyaḥ bhavati iti vaktavyam pūrvavipratiṣedhaḥ vā vaktavyaḥ . \\
\noindent \textbf{(P\_6,1.17.3)} ubhayam na vaktavyam . \\
\noindent \textbf{(P\_6,1.17.3)} uktam atra ubhayeṣāṅgrahaṇasya prayojanam ubhayeṣām abhyāsasya samprasāraṇam eva yathā syāt . \\
\noindent \textbf{(P\_6,1.17.3)} yat anyat prāpnoti tat mā bhūt iti . \\
\noindent \textbf{(P\_6,1.17.4)} vyaceḥ kuṭāditvam anasi añṇiti samprasāraṇārtham . \\
\noindent \textbf{(P\_6,1.17.4)} vyaceḥ kuṭāditvam anasi iti vaktavyam . \\
\noindent \textbf{(P\_6,1.17.4)} kim prayojanam . \\
\noindent \textbf{(P\_6,1.17.4)} añṇiti samprasāraṇārtham . \\
\noindent \textbf{(P\_6,1.17.4)} añṇiti samprasāraṇam yathā syāt . \\
\noindent \textbf{(P\_6,1.17.4)} udvicitā , udvicitum , udvicitavyam . \\
\noindent \textbf{(P\_6,1.17.4)} anasi iti kimartham . \\
\noindent \textbf{(P\_6,1.17.4)} uruvyacāḥ kaṇṭakaḥ . \\
\noindent \textbf{(P\_6,1.18)} caṅgrahaṇam śakyam akartum . \\
\noindent \textbf{(P\_6,1.18)} katham . \\
\noindent \textbf{(P\_6,1.18)} ṅiti iti vartate na ca anyaḥ svāpeḥ ṅit asti anyat ataḥ caṅaḥ . \\
\noindent \textbf{(P\_6,1.20)} vaśeḥ yaṅi pratiṣedhaḥ . \\
\noindent \textbf{(P\_6,1.20)} vaśeḥ yaṅi pratiṣedhaḥ vaktavyaḥ samprasāraṇasya . \\
\noindent \textbf{(P\_6,1.20)} vāvaśyate . \\
\noindent \textbf{(P\_6,1.20)} kva mā bhūt . \\
\noindent \textbf{(P\_6,1.20)} uṣṭaḥ , uśanti iti . \\
\noindent \textbf{(P\_6,1.20)} saḥ tarhi tathā pratiṣedhaḥ vaktavyaḥ . \\
\noindent \textbf{(P\_6,1.20)} na vaktavyaḥ . \\
\noindent \textbf{(P\_6,1.20)} yaṅi iti vartate . \\
\noindent \textbf{(P\_6,1.20)} evam tarhi anvācaṣṭe yaṅi iti vartate iti . \\
\noindent \textbf{(P\_6,1.20)} na etat anvākhyeyam adhikārāḥ anuvartante iti . \\
\noindent \textbf{(P\_6,1.20)} eṣaḥ eva nyāyaḥ yat uta adhikārāḥ anuvarteran iti . \\
\noindent \textbf{(P\_6,1.27)} kim nipātyate . \\
\noindent \textbf{(P\_6,1.27)} śrāsrapyoḥ śṛbhāvaḥ . \\
\noindent \textbf{(P\_6,1.27)} śrāsrapyoḥ śṛbhāvaḥ nipātyate . \\
\noindent \textbf{(P\_6,1.27)} kṣīrahaviṣoḥ iti vaktavyam . \\
\noindent \textbf{(P\_6,1.27)} śṛtam kṣīram . \\
\noindent \textbf{(P\_6,1.27)} śṛtam haviḥ . \\
\noindent \textbf{(P\_6,1.27)} kva mā bhūt . \\
\noindent \textbf{(P\_6,1.27)} śrāṇā yavāgūḥ , śrapitā yavāgūḥ iti . \\
\noindent \textbf{(P\_6,1.27)} śrapeḥ śṛtam anyatra hetoḥ . \\
\noindent \textbf{(P\_6,1.27)} śrapeḥ śṛtam anyatra hetoḥ iti vaktavyam iha mā bhūt . \\
\noindent \textbf{(P\_6,1.27)} śrapitam kṣīram devadattena yajñadattena iti . \\
\noindent \textbf{(P\_6,1.28)} āṅpūrvāt andhūdhasoḥ . \\
\noindent \textbf{(P\_6,1.28)} āṅpūrvāt andhūdhasoḥ iti vaktavyam . \\
\noindent \textbf{(P\_6,1.28)} āpīnaḥ andhuḥ , āpīnam ūdhaḥ . \\
\noindent \textbf{(P\_6,1.28)} kim prayojanam . \\
\noindent \textbf{(P\_6,1.28)} niyamārtham . \\
\noindent \textbf{(P\_6,1.28)} āṅpūrvāt andhūdhasoḥ eva . \\
\noindent \textbf{(P\_6,1.28)} kva mā bhūt . \\
\noindent \textbf{(P\_6,1.28)} āpyānaḥ candramāḥ iti . \\
\noindent \textbf{(P\_6,1.28)} ubhayataḥ niyamaḥ ca ayam draṣṭavyaḥ . \\
\noindent \textbf{(P\_6,1.28)} āṅpūrvāt eva andhūdhasoḥ , andhūdhasoḥ eva āṅpūrvāt iti . \\
\noindent \textbf{(P\_6,1.28)} kva mā bhūt . \\
\noindent \textbf{(P\_6,1.28)} prapyānaḥ andhuḥ , prapyānam ūdhaḥ . \\
\noindent \textbf{(P\_6,1.28)} āṅpūrvāt ca eṣa niyamaḥ draṣṭavyaḥ . \\
\noindent \textbf{(P\_6,1.28)} bhavati hi pīnam mukham , pīnāḥ śambaṭyaḥ , ślakṣṇapīnamukhī kanyā iti . \\
\noindent \textbf{(P\_6,1.30)} śveḥ liṭi abhyāsalakṣaṇapratiṣedhaḥ . \\
\noindent \textbf{(P\_6,1.30)} śveḥ liṭi abhyāsalakṣaṇam samprasāraṇam nityam prāpnoti . \\
\noindent \textbf{(P\_6,1.30)} tasya pratiṣedhaḥ vaktavyaḥ . \\
\noindent \textbf{(P\_6,1.30)} śiśviyatuḥ , śiśviyuḥ . \\
\noindent \textbf{(P\_6,1.30)} kim ucyate liṭi abhyāsalakṣaṇasya iti na punaḥ killakṣaṇasya api . \\
\noindent \textbf{(P\_6,1.30)} killakṣaṇam api hi nityam atra prāpnoti . \\
\noindent \textbf{(P\_6,1.30)} killakṣaṇam śvayatilakṣaṇam bādhiṣyate . \\
\noindent \textbf{(P\_6,1.30)} yathā eva tarhi killakṣaṇam śvayatilakṣaṇam bādhate evam abhyāsalakṣaṇam api bādheta . \\
\noindent \textbf{(P\_6,1.30)} na brūmaḥ apavādatvāt killakṣaṇam śvayatilakṣaṇam bādhiṣyate iti . \\
\noindent \textbf{(P\_6,1.30)} kim tarhi . \\
\noindent \textbf{(P\_6,1.30)} paratvāt . \\
\noindent \textbf{(P\_6,1.30)} śvayatilakṣaṇasya avakāśaḥ piti vacanāni . \\
\noindent \textbf{(P\_6,1.30)} śuśāva, śuśavitha , śiśvāya, śiśvayitha . \\
\noindent \textbf{(P\_6,1.30)} killakṣaṇasya avakāśaḥ anye kitaḥ . \\
\noindent \textbf{(P\_6,1.30)} śūnaḥ, śūnavān . \\
\noindent \textbf{(P\_6,1.30)} iha ubhayam prāpnoti . \\
\noindent \textbf{(P\_6,1.30)} śiśviyatuḥ , śiśviyuḥ iti . \\
\noindent \textbf{(P\_6,1.30)} śvayatilakṣaṇam bhavati vipratiṣedhena . \\
\noindent \textbf{(P\_6,1.30)} abhyāsalakṣaṇāt api tarhi śvayatilakṣaṇam bhaviṣyati vipratiṣedhena . \\
\noindent \textbf{(P\_6,1.30)} abhyāsalakṣaṇasya avakāśaḥ anye yajādayaḥ . \\
\noindent \textbf{(P\_6,1.30)} iyāja, uvāpa . \\
\noindent \textbf{(P\_6,1.30)} śvayatilakṣaṇasya avakāśaḥ param dhāturūpam . \\
\noindent \textbf{(P\_6,1.30)} śuśuvatuḥ , śuśuvuḥ , śuśuvitha . \\
\noindent \textbf{(P\_6,1.30)} śvayateḥ abhyāsasya ubhayam prāpnoti . \\
\noindent \textbf{(P\_6,1.30)} śiśiviyatuḥ , śiśviyuḥ . \\
\noindent \textbf{(P\_6,1.30)} śvayatilakṣaṇam bhaviṣyati vipratiṣedhena . \\
\noindent \textbf{(P\_6,1.30)} na eṣaḥ yuktaḥ vipratiṣedhaḥ . \\
\noindent \textbf{(P\_6,1.30)} na hi śvayateḥ abhyāsasya anye yajādayaḥ avakāśaḥ . \\
\noindent \textbf{(P\_6,1.30)} śvayateḥ yajādiṣu yaḥ pāṭhaḥ saḥ anavakāśaḥ . \\
\noindent \textbf{(P\_6,1.30)} tasya anavakāśatvāt ayuktaḥ vipratiṣedhaḥ . \\
\noindent \textbf{(P\_6,1.30)} tasmāt suṣṭhu uktam śveḥ liṭi abhyāsalakṣaṇapratiṣedhaḥ iti . \\
\noindent \textbf{(P\_6,1.32-33)} KA\_III,29.8-30.14 Ro\_IV,341-344 {1/65}     hvaḥ samprasāraṇe yogavibhāgaḥ . \\
\noindent \textbf{(P\_6,1.32-33)} KA\_III,29.8-30.14 Ro\_IV,341-344 {2/65}     hvaḥ samprasāraṇe yogavibhāgaḥ kartavyaḥ . \\
\noindent \textbf{(P\_6,1.32-33)} KA\_III,29.8-30.14 Ro\_IV,341-344 {3/65}     hvaḥ samprasāraṇam bhavati ṇau ca saṃścaṅoḥ . \\
\noindent \textbf{(P\_6,1.32-33)} KA\_III,29.8-30.14 Ro\_IV,341-344 {4/65}     tataḥ abhyastasya ca . \\
\noindent \textbf{(P\_6,1.32-33)} KA\_III,29.8-30.14 Ro\_IV,341-344 {5/65}     abhyastasya ca hvaḥ samprasāraṇam bhavati iti . \\
\noindent \textbf{(P\_6,1.32-33)} KA\_III,29.8-30.14 Ro\_IV,341-344 {6/65}     kimarthaḥ yogavibhāgaḥ . \\
\noindent \textbf{(P\_6,1.32-33)} KA\_III,29.8-30.14 Ro\_IV,341-344 {7/65}     ṇau saṃścaṅviṣayāṛthaḥ . \\
\noindent \textbf{(P\_6,1.32-33)} KA\_III,29.8-30.14 Ro\_IV,341-344 {8/65}     ṇau ca saṃścaṅviṣaye hvaḥ samprasāraṇam yathā syāt . \\
\noindent \textbf{(P\_6,1.32-33)} KA\_III,29.8-30.14 Ro\_IV,341-344 {9/65}     juhāvayiṣati , ajūhavat . \\
\noindent \textbf{(P\_6,1.32-33)} KA\_III,29.8-30.14 Ro\_IV,341-344 {10/65}     kim punaḥ kāraṇam na sidhyati . \\
\noindent \textbf{(P\_6,1.32-33)} KA\_III,29.8-30.14 Ro\_IV,341-344 {11/65}     hvaḥ abhyastasya iti ucyate na ca etat hvaḥ abhyastam . \\
\noindent \textbf{(P\_6,1.32-33)} KA\_III,29.8-30.14 Ro\_IV,341-344 {12/65}     kasya tarhi . \\
\noindent \textbf{(P\_6,1.32-33)} KA\_III,29.8-30.14 Ro\_IV,341-344 {13/65}     hvāyayateḥ . \\
\noindent \textbf{(P\_6,1.32-33)} KA\_III,29.8-30.14 Ro\_IV,341-344 {14/65}     hvaḥ etat abhyastam . \\
\noindent \textbf{(P\_6,1.32-33)} KA\_III,29.8-30.14 Ro\_IV,341-344 {15/65}     katham . \\
\noindent \textbf{(P\_6,1.32-33)} KA\_III,29.8-30.14 Ro\_IV,341-344 {16/65}     ekācaḥ dve prathamasya . \\
\noindent \textbf{(P\_6,1.32-33)} KA\_III,29.8-30.14 Ro\_IV,341-344 {17/65}     evam tarhi hvayateḥ abhyastasya iti ucyate na ca atra hvayatiḥ abhyastaḥ . \\
\noindent \textbf{(P\_6,1.32-33)} KA\_III,29.8-30.14 Ro\_IV,341-344 {18/65}     kaḥ tarhi . \\
\noindent \textbf{(P\_6,1.32-33)} KA\_III,29.8-30.14 Ro\_IV,341-344 {19/65}     hvāyayatiḥ . \\
\noindent \textbf{(P\_6,1.32-33)} KA\_III,29.8-30.14 Ro\_IV,341-344 {20/65}     hvayatiḥ eva atra abhyastaḥ . \\
\noindent \textbf{(P\_6,1.32-33)} KA\_III,29.8-30.14 Ro\_IV,341-344 {21/65}     katham . \\
\noindent \textbf{(P\_6,1.32-33)} KA\_III,29.8-30.14 Ro\_IV,341-344 {22/65}      ekācaḥ dve prathamasya iti . \\
\noindent \textbf{(P\_6,1.32-33)} KA\_III,29.8-30.14 Ro\_IV,341-344 {23/65}     evam api abhyastinimitte anabhyastaprasāraṇārtham . \\
\noindent \textbf{(P\_6,1.32-33)} KA\_III,29.8-30.14 Ro\_IV,341-344 {24/65}     abhyastinimitte iti vaktavyam . \\
\noindent \textbf{(P\_6,1.32-33)} KA\_III,29.8-30.14 Ro\_IV,341-344 {25/65}     kim prayojanam . \\
\noindent \textbf{(P\_6,1.32-33)} KA\_III,29.8-30.14 Ro\_IV,341-344 {26/65}     anabhyastaprasāraṇārtham . \\
\noindent \textbf{(P\_6,1.32-33)} KA\_III,29.8-30.14 Ro\_IV,341-344 {27/65}     anabhyastasya prasāraṇam yathā syāt . \\
\noindent \textbf{(P\_6,1.32-33)} KA\_III,29.8-30.14 Ro\_IV,341-344 {28/65}     juhūṣati , johūyate . \\
\noindent \textbf{(P\_6,1.32-33)} KA\_III,29.8-30.14 Ro\_IV,341-344 {29/65}     abhyastaprasāraṇe hi abhyāsaprasāraṇāprāptiḥ . \\
\noindent \textbf{(P\_6,1.32-33)} KA\_III,29.8-30.14 Ro\_IV,341-344 {30/65}     abhyastaprasāraṇe hi abhyāsaprasāraṇasya aprāptiḥ syāt . \\
\noindent \textbf{(P\_6,1.32-33)} KA\_III,29.8-30.14 Ro\_IV,341-344 {31/65}     na samprasāraṇe samprasāraṇam iti pratiṣedhaḥ prasajyeta . \\
\noindent \textbf{(P\_6,1.32-33)} KA\_III,29.8-30.14 Ro\_IV,341-344 {32/65}     na eṣaḥ doṣaḥ . \\
\noindent \textbf{(P\_6,1.32-33)} KA\_III,29.8-30.14 Ro\_IV,341-344 {33/65}     vyavahitatvāt na bhaviṣyati . \\
\noindent \textbf{(P\_6,1.32-33)} KA\_III,29.8-30.14 Ro\_IV,341-344 {34/65}     samānāṅge prasāraṇapratiṣedhāt pratiṣedhaḥ . \\
\noindent \textbf{(P\_6,1.32-33)} KA\_III,29.8-30.14 Ro\_IV,341-344 {35/65}     samānāṅge prasāraṇapratiṣedhāt pratiṣedhaḥ prāpnoti . \\
\noindent \textbf{(P\_6,1.32-33)} KA\_III,29.8-30.14 Ro\_IV,341-344 {36/65}     samānāṅgagrahaṇam tatra codayiṣyati . \\
\noindent \textbf{(P\_6,1.32-33)} KA\_III,29.8-30.14 Ro\_IV,341-344 {37/65}     kṛdantapratiṣedhārtham ca . \\
\noindent \textbf{(P\_6,1.32-33)} KA\_III,29.8-30.14 Ro\_IV,341-344 {38/65}     kṛdantapratiṣedhārtham ca abhyastinimitte iti vaktavyam . \\
\noindent \textbf{(P\_6,1.32-33)} KA\_III,29.8-30.14 Ro\_IV,341-344 {39/65}     kim prayojanam . \\
\noindent \textbf{(P\_6,1.32-33)} KA\_III,29.8-30.14 Ro\_IV,341-344 {40/65}     hvāyakam icchati hvāyakīyati . \\
\noindent \textbf{(P\_6,1.32-33)} KA\_III,29.8-30.14 Ro\_IV,341-344 {41/65}     hvāyakīyateḥ san . \\
\noindent \textbf{(P\_6,1.32-33)} KA\_III,29.8-30.14 Ro\_IV,341-344 {42/65}     jihvāyakīyiṣati . \\
\noindent \textbf{(P\_6,1.32-33)} KA\_III,29.8-30.14 Ro\_IV,341-344 {43/65}     saḥ tarhi nimittaśabdaḥ upādeyaḥ . \\
\noindent \textbf{(P\_6,1.32-33)} KA\_III,29.8-30.14 Ro\_IV,341-344 {44/65}     na hi antareṇa nimittaśabdam nimittārthaḥ gamyate . \\
\noindent \textbf{(P\_6,1.32-33)} KA\_III,29.8-30.14 Ro\_IV,341-344 {45/65}     antareṇa api nimittaśabdam nimittārthaḥ gamyate . \\
\noindent \textbf{(P\_6,1.32-33)} KA\_III,29.8-30.14 Ro\_IV,341-344 {46/65}     tat yathā : dadhitrapusam pratyakṣaḥ jvaraḥ . \\
\noindent \textbf{(P\_6,1.32-33)} KA\_III,29.8-30.14 Ro\_IV,341-344 {47/65}     jvaranimittam iti gamyate . \\
\noindent \textbf{(P\_6,1.32-33)} KA\_III,29.8-30.14 Ro\_IV,341-344 {48/65}     naḍvalodakam pādarogaḥ . \\
\noindent \textbf{(P\_6,1.32-33)} KA\_III,29.8-30.14 Ro\_IV,341-344 {49/65}     pādaroganimittam iti gamyate . \\
\noindent \textbf{(P\_6,1.32-33)} KA\_III,29.8-30.14 Ro\_IV,341-344 {50/65}     āyuḥ ghṛtam . \\
\noindent \textbf{(P\_6,1.32-33)} KA\_III,29.8-30.14 Ro\_IV,341-344 {51/65}     āyuṣaḥ nimittam iti gamyate . \\
\noindent \textbf{(P\_6,1.32-33)} KA\_III,29.8-30.14 Ro\_IV,341-344 {52/65}     atha vā akāraḥ matvarthīyaḥ . \\
\noindent \textbf{(P\_6,1.32-33)} KA\_III,29.8-30.14 Ro\_IV,341-344 {53/65}     abhyastam asmin asti saḥ ayam abhyastaḥ . \\
\noindent \textbf{(P\_6,1.32-33)} KA\_III,29.8-30.14 Ro\_IV,341-344 {54/65}     abhyastasya iti . \\
\noindent \textbf{(P\_6,1.32-33)} KA\_III,29.8-30.14 Ro\_IV,341-344 {55/65}     atha vā abhyastasya iti na eṣā hvayatisamānādhikaraṇā ṣaṣṭhī . \\
\noindent \textbf{(P\_6,1.32-33)} KA\_III,29.8-30.14 Ro\_IV,341-344 {56/65}     kā tarhi . \\
\noindent \textbf{(P\_6,1.32-33)} KA\_III,29.8-30.14 Ro\_IV,341-344 {57/65}     sambandhaṣaṣṭhī . \\
\noindent \textbf{(P\_6,1.32-33)} KA\_III,29.8-30.14 Ro\_IV,341-344 {58/65}     abhyastasya yaḥ hvayatiḥ . \\
\noindent \textbf{(P\_6,1.32-33)} KA\_III,29.8-30.14 Ro\_IV,341-344 {59/65}     kim ca abhyastasya hvayatiḥ . \\
\noindent \textbf{(P\_6,1.32-33)} KA\_III,29.8-30.14 Ro\_IV,341-344 {60/65}     prakṛtiḥ . \\
\noindent \textbf{(P\_6,1.32-33)} KA\_III,29.8-30.14 Ro\_IV,341-344 {61/65}     hvaḥ abhyastasya prakṛteḥ iti . \\
\noindent \textbf{(P\_6,1.32-33)} KA\_III,29.8-30.14 Ro\_IV,341-344 {62/65}     yogavibhāgaḥ tu kartavyaḥ eva . \\
\noindent \textbf{(P\_6,1.32-33)} KA\_III,29.8-30.14 Ro\_IV,341-344 {63/65}     na atra hvayatiḥ abhyastasya prakṛtiḥ . \\
\noindent \textbf{(P\_6,1.32-33)} KA\_III,29.8-30.14 Ro\_IV,341-344 {64/65}     kim tarhi . \\
\noindent \textbf{(P\_6,1.32-33)} KA\_III,29.8-30.14 Ro\_IV,341-344 {65/65}     hvāyayatiḥ . \\
\noindent \textbf{(P\_6,1.36)} apaspṛdhethām iti kim nipātyate . \\
\noindent \textbf{(P\_6,1.36)} spardheḥ laṅi ātmanepadānām madhyamapuruṣasya dvivacane āthāmi dvirvacanam samprasāraṇam akāralopaḥ ca nipātyate . \\
\noindent \textbf{(P\_6,1.36)} indraḥ ca viṣṇo yat apaspṛdhethām . \\
\noindent \textbf{(P\_6,1.36)} aspṛdhethām iti bhāṣāyām . \\
\noindent \textbf{(P\_6,1.36)} aparaḥ āha : apapūrvāt spardheḥ laṅi ātmanepadānām madhyamapuruṣasya dvivacane āthāmi dvirvacanam samprasāraṇam akāralopaḥ ca nipātyate . \\
\noindent \textbf{(P\_6,1.36)} indraḥ ca viṣṇo yat apaspṛdhethām . \\
\noindent \textbf{(P\_6,1.36)} apāspṛdhethām iti bhāṣāyām . \\
\noindent \textbf{(P\_6,1.36)} śrātāḥ śritam iti kim nipātyate . \\
\noindent \textbf{(P\_6,1.36)} śrīṇāteḥ kte śrābhāvaśribhāvau nipātyete . \\
\noindent \textbf{(P\_6,1.36)} kva punaḥ śrābhāvaḥ kva vā śribhāvaḥ . \\
\noindent \textbf{(P\_6,1.36)} some śrābhāvaḥ anyatra śribhāvaḥ . \\
\noindent \textbf{(P\_6,1.36)} na tarhi idānīm idam bhavati : śritaḥ somaḥ iti . \\
\noindent \textbf{(P\_6,1.36)} bahuvacane śrābhāvaḥ . \\
\noindent \textbf{(P\_6,1.36)} na tarhi idānīm idam bhavati : śritāḥ naḥ grahāḥ iti . \\
\noindent \textbf{(P\_6,1.36)} somabahutve śrābhāvaḥ anyatra śribhāvaḥ . \\
\noindent \textbf{(P\_6,1.37.1)} kimartham idam ucyate . \\
\noindent \textbf{(P\_6,1.37.1)} vacispaviyajādīnām grahādīnām ca samprasāraṇam uktam . \\
\noindent \textbf{(P\_6,1.37.1)} tatra yāvantaḥ yaṇaḥ sarveṣām samprasāraṇam prāpnoti . \\
\noindent \textbf{(P\_6,1.37.1)} iṣyate ca parasya yathā syāt na pūrvasya tat ca antareṇa yatnam na sidhyati iti na samprasāraṇe samprasāraṇam . \\
\noindent \textbf{(P\_6,1.37.1)} kim anye api evam vidhayaḥ bhavanti . \\
\noindent \textbf{(P\_6,1.37.1)} ataḥ dīrghaḥ yañi . \\
\noindent \textbf{(P\_6,1.37.1)} supi ca iti . \\
\noindent \textbf{(P\_6,1.37.1)} ghaṭābhyām . \\
\noindent \textbf{(P\_6,1.37.1)} akāramātrasya dīrghatvam kasmāt na bhavati . \\
\noindent \textbf{(P\_6,1.37.1)} asti atra viśeṣaḥ . \\
\noindent \textbf{(P\_6,1.37.1)} iyam atra paribhāṣā upatiṣṭhate . \\
\noindent \textbf{(P\_6,1.37.1)} alaḥ antyasya iti . \\
\noindent \textbf{(P\_6,1.37.1)} nanu ca idānīm etayā paribhāṣayā iha (R: iha api) śakyam upasthātum . \\
\noindent \textbf{(P\_6,1.37.1)} na iti āha . \\
\noindent \textbf{(P\_6,1.37.1)} na hi vacispaviyajādīnām grahādīnām ca antyaḥ yaṇ asti . \\
\noindent \textbf{(P\_6,1.37.1)} evam tarhi anantyavikāre antyasadeśasya kāryam bhavati iti antyasasdeśaḥ yaḥ yaṇ tasya kāryam bhaviṣyati . \\
\noindent \textbf{(P\_6,1.37.1)} na etasyāḥ paribhāṣāyāḥ santi prayojanāni . \\
\noindent \textbf{(P\_6,1.37.1)} evam tarhi ācāryapravṛttiḥ jñāpayati na sarvasya yaṇaḥ samprasāraṇam bhavati iti yat ayam pyāyaḥ pībhāvam śāsti . \\
\noindent \textbf{(P\_6,1.37.1)} katham kṛtvā jñāpakam . \\
\noindent \textbf{(P\_6,1.37.1)} pībhāvavacane etat prayojanam āpīnaḥ andhuḥ , āpīnam ūdhaḥ etat rūpam yathā syāt iti . \\
\noindent \textbf{(P\_6,1.37.1)} yadi ca atra sarvasya yaṇaḥ samprasāraṇam syāt pībhāvavacanam anarthakam syāt . \\
\noindent \textbf{(P\_6,1.37.1)} samprasāraṇe kṛte samprasāraṇaparapūrvatve ca dvayoḥ ikārayoḥ ekādeśe siddham rūpam syāt āpīnaḥ andhuḥ , āpīnam ūdhaḥ iti . \\
\noindent \textbf{(P\_6,1.37.1)} paśyati tu ācāryaḥ na sarvasya yaṇaḥ samprasāraṇam bhavati iti . \\
\noindent \textbf{(P\_6,1.37.1)} tataḥ ayam pyāyaḥ pībhāvam śāsti . \\
\noindent \textbf{(P\_6,1.37.1)} na etat asti jñāpakam . \\
\noindent \textbf{(P\_6,1.37.1)} siddhe hi vidhiḥ ārabhyamāṇaḥ jñāpakārthaḥ bhavati na ca pyāyaḥ samprasāraṇena sidhyati . \\
\noindent \textbf{(P\_6,1.37.1)} samprasāraṇe hi sati antyasya prasajyeta . \\
\noindent \textbf{(P\_6,1.37.1)} evam api jñāpakam eva . \\
\noindent \textbf{(P\_6,1.37.1)} katham . \\
\noindent \textbf{(P\_6,1.37.1)} pyāyaḥ iti na eṣā sthānaṣaṣṭhī . \\
\noindent \textbf{(P\_6,1.37.1)} kā tarhi . \\
\noindent \textbf{(P\_6,1.37.1)} viśeṣaṇaṣaṣṭhī . \\
\noindent \textbf{(P\_6,1.37.1)} pyāyaḥ yaḥ yaṇ iti . \\
\noindent \textbf{(P\_6,1.37.1)} tat etat jñāpayati ācāryaḥ na sarvasya yaṇaḥ samprasāraṇam bhavati iti yat ayam pyāyaḥ pībhāvam śāsti . \\
\noindent \textbf{(P\_6,1.37.1)} evam api anaikāntikam etat . \\
\noindent \textbf{(P\_6,1.37.1)} etāvat jñāpyate na sarvasya yaṇaḥ samprasāraṇam bhavati iti . \\
\noindent \textbf{(P\_6,1.37.1)} tatra kutaḥ etat parasya bhaviṣyati na pūrvasya iti . \\
\noindent \textbf{(P\_6,1.37.1)} ucyamāne api etasmin kutaḥ etat parasya bhaviṣyati na pūrvasya iti . \\
\noindent \textbf{(P\_6,1.37.1)} ekayogakṣaṇam khalu api samprasāraṇam . \\
\noindent \textbf{(P\_6,1.37.1)} tat yadi tāvat param abhinirvṛttam pūrvam api abhinirvṛttam eva . \\
\noindent \textbf{(P\_6,1.37.1)} prasaktasya anabhinirvṛttasya pratiṣedhena nivṛttiḥ śakyā kartum na abhinirvṛttasya . \\
\noindent \textbf{(P\_6,1.37.1)} yaḥ hi bhuktavantam brūyāt mā bhukthāḥ iti kim tena kṛtam syāt . \\
\noindent \textbf{(P\_6,1.37.1)} atha api pūrvam anabhinirvṛttam param api anabhinirvṛttam eva . \\
\noindent \textbf{(P\_6,1.37.1)} tatra nimittasaṃśrayaḥ anupapannaḥ na samprasāraṇe samprasāraṇam iti . \\
\noindent \textbf{(P\_6,1.37.1)} na eṣaḥ doṣaḥ . \\
\noindent \textbf{(P\_6,1.37.1)} yat tāvat ucyate ucyamāne api etasmin kutaḥ etat parasya bhaviṣyati na pūrvasya iti . \\
\noindent \textbf{(P\_6,1.37.1)} iha iṅgitena ceṣṭitena nimiṣitena mahatā vā sūtraprabandhena ācāryāṇām abhiprāyaḥ gamyate . \\
\noindent \textbf{(P\_6,1.37.1)} etat eva jñāpayati parasya bhaviṣyati na pūrvasya iti yat ayam na samprasāraṇe samprasāraṇam iti pratiṣedham śāsti . \\
\noindent \textbf{(P\_6,1.37.1)} yat api ucyate ekayogalakṣaṇam khalu api samprasāraṇam . \\
\noindent \textbf{(P\_6,1.37.1)} tat yadi tāvat param abhinirvṛttam pūrvam api abhinirvṛttam eva . \\
\noindent \textbf{(P\_6,1.37.1)} prasaktasya anabhinirvṛttasya pratiṣedhena nivṛttiḥ śakyā kartum iti . \\
\noindent \textbf{(P\_6,1.37.1)} astu ubhayoḥ abhinirvṛttiḥ . \\
\noindent \textbf{(P\_6,1.37.1)} na vayam pūrvasya pratiṣedham śiṣmaḥ . \\
\noindent \textbf{(P\_6,1.37.1)} kim tarhi . \\
\noindent \textbf{(P\_6,1.37.1)} samprasāraṇāśrayam yat prāpnoti tasya pratiṣedham . \\
\noindent \textbf{(P\_6,1.37.1)} tataḥ pūrvatve pratiṣiddhe yaṇādeśena siddham . \\
\noindent \textbf{(P\_6,1.37.1)} yat api ucyate atha api pūrvam anabhinirvṛttam param api anabhinirvṛttam eva . \\
\noindent \textbf{(P\_6,1.37.1)} tatra nimittasaṃśrayaḥ anupapannaḥ iti . \\
\noindent \textbf{(P\_6,1.37.1)} tādarthyāt tācchabdyam bhaviṣyati . \\
\noindent \textbf{(P\_6,1.37.1)} tat yathā indrārthā sthūṇā indraḥ iti evam iha api samprasāraṇārtham samprasāraṇam . \\
\noindent \textbf{(P\_6,1.37.1)} tat yat prasāraṇārtham prasāraṇam tasmin pratiṣedhaḥ bhaviṣyati . \\
\noindent \textbf{(P\_6,1.37.2)} atha samprasāraṇam iti vartamāne punaḥ samprasāraṇagrahaṇam kimartham . \\
\noindent \textbf{(P\_6,1.37.2)} prasāraṇaprakaraṇe punaḥ prasāraṇagrahaṇam ataḥ anyatra prasāraṇapratiṣedhārtham . \\
\noindent \textbf{(P\_6,1.37.2)} samprasāraṇaprakaraṇe punaḥ prasāraṇagrahaṇe (R: samprasāraṇagrahaṇe) etat prayojanam . \\
\noindent \textbf{(P\_6,1.37.2)} videśastham api yat samprasāraṇam tasya api pratiṣedhaḥ yathā syāt . \\
\noindent \textbf{(P\_6,1.37.2)} vyathaḥ liṭi . \\
\noindent \textbf{(P\_6,1.37.2)} vivyathe . \\
\noindent \textbf{(P\_6,1.37.2)} na etat asti prayojanam . \\
\noindent \textbf{(P\_6,1.37.2)} halādiśeṣāpavādaḥ atra samprasāraṇam . \\
\noindent \textbf{(P\_6,1.37.2)} idam tarhi śvayuvamaghonām ataddhite . \\
\noindent \textbf{(P\_6,1.37.2)} yūnā , yūne . \\
\noindent \textbf{(P\_6,1.37.2)} ucyamāne api etasmin na sidhyati . \\
\noindent \textbf{(P\_6,1.37.2)} kim kāraṇam . \\
\noindent \textbf{(P\_6,1.37.2)} ukāreṇa vyavadhānāt . \\
\noindent \textbf{(P\_6,1.37.2)} ekādeśe kṛte na asti vyavadhānam . \\
\noindent \textbf{(P\_6,1.37.2)} ekādeśaḥ pūrvavidhau sthānivat bhavati iti sthānivadbhāvāt vyavadhānam eva . \\
\noindent \textbf{(P\_6,1.37.2)} evam tarhi samānāṅgagrahaṇam ca . \\
\noindent \textbf{(P\_6,1.37.2)} samānāṅgagrahaṇam ca kartavyam . \\
\noindent \textbf{(P\_6,1.37.2)} na samprasāraṇe samprasāraṇam samānāṅge iti vaktavyam . \\
\noindent \textbf{(P\_6,1.37.2)} tatra upoṣuṣi doṣaḥ . \\
\noindent \textbf{(P\_6,1.37.2)} tatra upoṣuṣi doṣaḥ bhavati . \\
\noindent \textbf{(P\_6,1.37.2)} na vā yasya aṅgasya prasāraṇaprāptiḥ tasmin prāptipratiṣedhāt . \\
\noindent \textbf{(P\_6,1.37.2)} na vā eṣaḥ doṣaḥ . \\
\noindent \textbf{(P\_6,1.37.2)} kim kāraṇam yasya aṅgasya prasāraṇaprāptiḥ tasmin dvitīyā yā prāptiḥ sā pratiṣidhyate . \\
\noindent \textbf{(P\_6,1.37.2)} atra ca vasiḥ kvasau aṅgam kvasantam punaḥ vibhaktau . \\
\noindent \textbf{(P\_6,1.37.2)} atha vā yasya aṅgasya prasāraṇaprāptiḥ iti anena kim kriyate . \\
\noindent \textbf{(P\_6,1.37.2)} yāvat brūyāt prasaktasya anabhinirvṛttasya pratiṣedhena nivṛttiḥ śakyā kartum iti . \\
\noindent \textbf{(P\_6,1.37.2)} atra ca yadā vaseḥ na tadā kvasoḥ yadā ca kvasoḥ abhinirvṛttam tadā vaseḥ bhavati . \\
\noindent \textbf{(P\_6,1.37.2)} atha vā yasya aṅgasya prasāraṇaprāptiḥ iti anena kim kriyate . \\
\noindent \textbf{(P\_6,1.37.2)} yāvat brūyāt asiddham bahiraṅgam antaraṅge iti . \\
\noindent \textbf{(P\_6,1.37.2)} asiddhatvāt bahiraṅgalakṣaṇasya vasausamprasāraṇasya antaraṅgalakṣaṇaḥ pratiṣedhaḥ na bhaviṣyati . \\
\noindent \textbf{(P\_6,1.37.3)} ṛci treḥ uttarapadādilopaḥ chandasi . \\
\noindent \textbf{(P\_6,1.37.3)} ṛci treḥ samprasāraṇam vaktavyam . \\
\noindent \textbf{(P\_6,1.37.3)} uttarapadādilopaḥ chandasi vaktavyaḥ . \\
\noindent \textbf{(P\_6,1.37.3)} tṛcam sūktam . \\
\noindent \textbf{(P\_6,1.37.3)} tṛcam sāma . \\
\noindent \textbf{(P\_6,1.37.3)} chandasi iti kim . \\
\noindent \textbf{(P\_6,1.37.3)} tryṛcāni . \\
\noindent \textbf{(P\_6,1.37.3)} rayeḥ matau bahulam . \\
\noindent \textbf{(P\_6,1.37.3)} rayeḥ matau samprasāraṇam bahulam vaktavyam . \\
\noindent \textbf{(P\_6,1.37.3)} ā revān etu naḥ viśaḥ . \\
\noindent \textbf{(P\_6,1.37.3)} na ca bhavati . \\
\noindent \textbf{(P\_6,1.37.3)} rayiman puṣṭivardhanaḥ . \\
\noindent \textbf{(P\_6,1.37.3)} kakṣyāyāḥ sañjñāyām . \\
\noindent \textbf{(P\_6,1.37.3)} kakṣyāyāḥ sañjñāyām matau samprasāraṇam vaktavyam . \\
\noindent \textbf{(P\_6,1.37.3)} kakṣīvantam yaḥ āśijaḥ . \\
\noindent \textbf{(P\_6,1.37.3)} kaṇvaḥ kakṣīvān . \\
\noindent \textbf{(P\_6,1.37.3)} sañjñāyām iti kim . \\
\noindent \textbf{(P\_6,1.37.3)} kaṣyāvān hastī . \\
\noindent \textbf{(P\_6,1.39)} vaścāsyagrahaṇam śakyam akartum . \\
\noindent \textbf{(P\_6,1.39)} anyatarasyām kiti veñaḥ na samprasāraṇam bhavati iti eva siddham . \\
\noindent \textbf{(P\_6,1.39)} katham . \\
\noindent \textbf{(P\_6,1.39)} samprasāraṇe kṛte uvaṅādeśe ca dvirvacanam savarṇadīrghatvam . \\
\noindent \textbf{(P\_6,1.39)} tena siddham vavatuḥ, vavuḥ , ūvatuḥ, ūvuḥ . \\
\noindent \textbf{(P\_6,1.39)} vayeḥ api nityam yakārasya pratiṣedhaḥ samprasāraṇasya ūyatuḥ , ūyuḥ . \\
\noindent \textbf{(P\_6,1.39)} traiśabyam ca iha sādhyam . \\
\noindent \textbf{(P\_6,1.39)} tat ca evam sati siddham bhavati . \\
\noindent \textbf{(P\_6,1.39)} yadi evam vavau, vavitha iti na sidhyati . \\
\noindent \textbf{(P\_6,1.39)} lyapi ca iti anena cakāreṇa liṭ api anukṛṣyate . \\
\noindent \textbf{(P\_6,1.39)} tasmin nitye prasāraṇapratiṣedhe prāpte iyam kiti vibhāṣā ārabhyate . \\
\noindent \textbf{(P\_6,1.45.1)} katham idam vijñāyate : ec yaḥ upadeśe iti āhosvit ejantantam yat upadeśe iti . \\
\noindent \textbf{(P\_6,1.45.1)} kim ca ataḥ . \\
\noindent \textbf{(P\_6,1.45.1)} yadi vijñāyate : ec yaḥ upadeśe iti ḍhaukitā traukitā iti atra api prāpnoti . \\
\noindent \textbf{(P\_6,1.45.1)} atha vijñāyate : ejantantam yat upadeśe iti na doṣaḥ bhavati . \\
\noindent \textbf{(P\_6,1.45.1)} nanu ca ejantantam yat upadeśe iti api vijñāyamāne atra api prāpnoti . \\
\noindent \textbf{(P\_6,1.45.1)} etat api vyapadeśivadbhāvena ejantam bhavati upadeśe . \\
\noindent \textbf{(P\_6,1.45.1)} arthavatā vyapadeśivadbhāvaḥ . \\
\noindent \textbf{(P\_6,1.45.1)} nanu ca ec yaḥ upadeśe iti api vijñāyamāne na doṣaḥ bhavati . \\
\noindent \textbf{(P\_6,1.45.1)} aśiti iti ucyate na ca atra aśitam paśyāmaḥ . \\
\noindent \textbf{(P\_6,1.45.1)} nanu ca kakāraḥ eva atra aśit . \\
\noindent \textbf{(P\_6,1.45.1)} na kakāre bhavitavyam . \\
\noindent \textbf{(P\_6,1.45.1)} kim kāraṇam . \\
\noindent \textbf{(P\_6,1.45.1)} nañivayuktam anyasadṛśādhikaraṇe . \\
\noindent \textbf{(P\_6,1.45.1)} tathā hi arthagatiḥ . \\
\noindent \textbf{(P\_6,1.45.1)} nañyuktam ivayuktam ca anyasmin tatsadṛśe kāryam vijñāyate . \\
\noindent \textbf{(P\_6,1.45.1)} tathā hi arthaḥ gamyate . \\
\noindent \textbf{(P\_6,1.45.1)} tat yathā loke : abrāhmaṇam ānaya iti ukte brāhmaṇasadṛśam ānayati . \\
\noindent \textbf{(P\_6,1.45.1)} na asau loṣṭam ānīya kṛtī bhavati . \\
\noindent \textbf{(P\_6,1.45.1)} evam iha api aśiti iti śitpratiṣedhāt anyasmin aśiti śitsadṛśe kāryam vijñāsyate . \\
\noindent \textbf{(P\_6,1.45.1)} kim ca anyat śitsadṛśam . \\
\noindent \textbf{(P\_6,1.45.1)} pratyayaḥ . \\
\noindent \textbf{(P\_6,1.45.1)} iha tarhi : glai : glānīyam , mlai : mlānīyam , veñ : vānīyam , śo : niśāmīyam : paratvāt āyādayaḥ prāpnuvanti . \\
\noindent \textbf{(P\_6,1.45.1)} nanu ca ejantantam yat upadeśe iti api vijñāyamāne paratvāt āyādayaḥ prāpnuvanti . \\
\noindent \textbf{(P\_6,1.45.1)} santu . \\
\noindent \textbf{(P\_6,1.45.1)} āyādiṣu kṛteṣu sthānivadbhāvāt ejgrahaṇena grahaṇāt punaḥ āttvam bhaviṣyati . \\
\noindent \textbf{(P\_6,1.45.1)} nanu ca ec yaḥ upadeśe iti api vijñāyamāne paratvāt āyādiṣu kṛteṣu sthānivadbhāvāt ejgrahaṇena grahaṇāt āttvam bhaviṣyati . \\
\noindent \textbf{(P\_6,1.45.1)} na bhaviṣyati . \\
\noindent \textbf{(P\_6,1.45.1)} analvidhau sthānivadbhāvaḥ alvidhiḥ ca ayam . \\
\noindent \textbf{(P\_6,1.45.1)} evam tarhi ejantantam yat upadeśe iti api vijñāyamāne hūtaḥ , hūtavān iti atra api prāpnoti . \\
\noindent \textbf{(P\_6,1.45.1)} bhavatu eva atra āttvam . \\
\noindent \textbf{(P\_6,1.45.1)} śravaṇam kasmāt na bhavati . \\
\noindent \textbf{(P\_6,1.45.1)} pūrvatvam asya bhaviṣyati . \\
\noindent \textbf{(P\_6,1.45.1)} na sidhyati . \\
\noindent \textbf{(P\_6,1.45.1)} idam iha sampradhāryam : āttvam kriyatām pūrvatvam iti . \\
\noindent \textbf{(P\_6,1.45.1)} kim atra kartavyam . \\
\noindent \textbf{(P\_6,1.45.1)} paratvāt pūrvatvam . \\
\noindent \textbf{(P\_6,1.45.1)} evam tarhi idam iha sampradhāryam . \\
\noindent \textbf{(P\_6,1.45.1)} āttvam kriyatām samprasāraṇam iti . \\
\noindent \textbf{(P\_6,1.45.1)} kim atra kartavyam . \\
\noindent \textbf{(P\_6,1.45.1)} paratvāt āttvam . \\
\noindent \textbf{(P\_6,1.45.1)} nityam samprasāraṇam . \\
\noindent \textbf{(P\_6,1.45.1)} kṛte api āttve prāpnoti akṛte api . \\
\noindent \textbf{(P\_6,1.45.1)} āttvam api nityam . \\
\noindent \textbf{(P\_6,1.45.1)} kṛte api samprasāraṇe prāpnoti akṛte api . \\
\noindent \textbf{(P\_6,1.45.1)} anityam āttvam . \\
\noindent \textbf{(P\_6,1.45.1)} na hi kṛte samprasāraṇe prāpnoti . \\
\noindent \textbf{(P\_6,1.45.1)} kim kāraṇam . \\
\noindent \textbf{(P\_6,1.45.1)} antaraṅgam pūrvatvam . \\
\noindent \textbf{(P\_6,1.45.1)} tena bādhyate . \\
\noindent \textbf{(P\_6,1.45.1)} yasya lakṣaṇantareṇa nimittam vihanyate na tat anityam . \\
\noindent \textbf{(P\_6,1.45.1)} na ca samprasāraṇam eva āttvasya nimittam hanti . \\
\noindent \textbf{(P\_6,1.45.1)} avaśyam lakṣaṇāntaram pūrvatvam pratīkṣyam . \\
\noindent \textbf{(P\_6,1.45.1)} ubhayoḥ nityayoḥ paratvāt āttve kṛte samprasāraṇam samprasāraṇapūrvatvam . \\
\noindent \textbf{(P\_6,1.45.1)} kāryakṛtatvāt punaḥ āttvam na bhaviṣyati . \\
\noindent \textbf{(P\_6,1.45.1)} atha api katham cit āttvam anityam syāt evam api na doṣaḥ . \\
\noindent \textbf{(P\_6,1.45.1)} upadeśagrahaṇam na kariṣyate . \\
\noindent \textbf{(P\_6,1.45.1)} yadi na kriyate cetā stotā iti atra api prāpnoti . \\
\noindent \textbf{(P\_6,1.45.1)} na eṣaḥ doṣaḥ . \\
\noindent \textbf{(P\_6,1.45.1)} ācāryapravṛttiḥ jñāpayati na paranimittakasya āttvam bhavati iti yat ayam krīṅjīṇām ṇau āttvam śāsti . \\
\noindent \textbf{(P\_6,1.45.1)} na etat asti jñāpakam . \\
\noindent \textbf{(P\_6,1.45.1)} niyamārtham etat syāt . \\
\noindent \textbf{(P\_6,1.45.1)} krīṅjīṇām ṇau eva iti . \\
\noindent \textbf{(P\_6,1.45.1)} yat tarhi mīnātiminotidīṅām lyapi ca iti atra ejgrahaṇam anuvartayati . \\
\noindent \textbf{(P\_6,1.45.1)} iha tarhi glai glānīyam , mlai mlānīyam , veñ vānīyam , śo niśāmīyam paratvāt āyādayaḥ prāpnuvanti . \\
\noindent \textbf{(P\_6,1.45.1)} atra api ācāryapravṛttiḥ jñāpayati na āyādayaḥ āttvam bādhante iti yat ayam aśiti iti pratiṣedham śāsti . \\
\noindent \textbf{(P\_6,1.45.1)} yadi hi bādheran śiti api bādheran . \\
\noindent \textbf{(P\_6,1.45.1)} atha vā punaḥ astu ec yaḥ upadeśe iti . \\
\noindent \textbf{(P\_6,1.45.1)} nanu ca uktam glai glānīyam , mlai mlānīyam , veñ vānīyam , śo niśāmīyam paratvāt āyādayaḥ prāpnuvanti iti . \\
\noindent \textbf{(P\_6,1.45.1)} atra api śitpratiṣedhaḥ jñāpakaḥ na āyādayaḥ āttvam bādhante iti . \\
\noindent \textbf{(P\_6,1.45.1)} āttve eśi upasaṅkhyānam . \\
\noindent \textbf{(P\_6,1.45.1)} āttve eśi upasaṅkhyānam kartavyam . \\
\noindent \textbf{(P\_6,1.45.1)} jagle mamle . \\
\noindent \textbf{(P\_6,1.45.1)} aśiti iti pratiṣedhaḥ prāpnoti . \\
\noindent \textbf{(P\_6,1.45.1)} na eṣaḥ doṣaḥ . \\
\noindent \textbf{(P\_6,1.45.1)} na evam vijñāyate . \\
\noindent \textbf{(P\_6,1.45.1)} śakāraḥ it yasya saḥ ayam śit . \\
\noindent \textbf{(P\_6,1.45.1)} na śit aśit . \\
\noindent \textbf{(P\_6,1.45.1)} aśiti iti . \\
\noindent \textbf{(P\_6,1.45.1)} katham tarhi . \\
\noindent \textbf{(P\_6,1.45.1)} śakāraḥ it śit . \\
\noindent \textbf{(P\_6,1.45.1)} na śit śit aśit . \\
\noindent \textbf{(P\_6,1.45.1)} aśiti iti . \\
\noindent \textbf{(P\_6,1.45.1)} yadi evam stanandhayaḥ iti atra api prāpnoti . \\
\noindent \textbf{(P\_6,1.45.1)} atra api śap śit bhavati . \\
\noindent \textbf{(P\_6,1.45.2)} kim punaḥ ayam paryudāsaḥ : yat anyat śitaḥ iti āhosvit prasajya ayam pratiṣedhaḥ : śiti na iti . \\
\noindent \textbf{(P\_6,1.45.2)} kaḥ ca atra viśeṣaḥ . \\
\noindent \textbf{(P\_6,1.45.2)} aśiti ekādeśe pratiṣedhaḥ ādivattvāt . \\
\noindent \textbf{(P\_6,1.45.2)} aśiti ekādeśe pratiṣedhaḥ vaktavyaḥ : glāyanti mlayanti . \\
\noindent \textbf{(P\_6,1.45.2)} kim kāraṇam . \\
\noindent \textbf{(P\_6,1.45.2)} ādivattvāt . \\
\noindent \textbf{(P\_6,1.45.2)} śidaśitoḥ ekādeśaḥ ādivat syāt . \\
\noindent \textbf{(P\_6,1.45.2)} asti anyat śitaḥ iti kṛtvā āttvam prāpnoti . \\
\noindent \textbf{(P\_6,1.45.2)} pratyayavidhiḥ . \\
\noindent \textbf{(P\_6,1.45.2)} pratyayavidhiḥ ca na sidhyati . \\
\noindent \textbf{(P\_6,1.45.2)} suglaḥ sumlaḥ . \\
\noindent \textbf{(P\_6,1.45.2)} ākārāntalakṣaṇaḥ pratyayavidhiḥ na prāpnoti . \\
\noindent \textbf{(P\_6,1.45.2)} aniṣṭasya pratyayasya śravaṇam prasajyeta . \\
\noindent \textbf{(P\_6,1.45.2)} abhyāsarūpam ca . \\
\noindent \textbf{(P\_6,1.45.2)} abhyāsarūpam ca na sidhyati : jagle mamle . \\
\noindent \textbf{(P\_6,1.45.2)} ivarṇābhyāsatā prāpnoti . \\
\noindent \textbf{(P\_6,1.45.2)} ayavāyāvām pratiṣedhaḥ ca . \\
\noindent \textbf{(P\_6,1.45.2)} ayavāyāvām ca pratiṣedhaḥ vaktavyaḥ : glai : glānīyam , mlai : mlānīyam , veñ : vānīyam , śo : niśāmīyam : paratvāt āyādayaḥ prāpnuvanti . \\
\noindent \textbf{(P\_6,1.45.2)} astu tarhi prasajya pratiṣedhaḥ śiti na iti . \\
\noindent \textbf{(P\_6,1.45.2)} śiti pratiṣedhe ślulukoḥ upasaṅkhyānam rarīdhvam trādhvam śiśīte . \\
\noindent \textbf{(P\_6,1.45.2)} śiti pratiṣedhe ślulukoḥ upasaṅkhyānam kartavyam . \\
\noindent \textbf{(P\_6,1.45.2)} divaḥ naḥ vṛṣṭim marutaḥ rarīdhvam . \\
\noindent \textbf{(P\_6,1.45.2)} luk . \\
\noindent \textbf{(P\_6,1.45.2)} trādhvam naḥ devā nijuraḥ vṛkasya . \\
\noindent \textbf{(P\_6,1.45.2)} śiśīte śṛṅge rakṣase vinikṣe . \\
\noindent \textbf{(P\_6,1.45.2)} na eṣaḥ doṣaḥ . \\
\noindent \textbf{(P\_6,1.45.2)} iha tāvat divaḥ naḥ vṛṣṭim marutaḥ rarīdhvam iti . \\
\noindent \textbf{(P\_6,1.45.2)} na etat rai iti asya rūpam . \\
\noindent \textbf{(P\_6,1.45.2)} kasya tarhi . \\
\noindent \textbf{(P\_6,1.45.2)} rāteḥ dānakarmaṇaḥ . \\
\noindent \textbf{(P\_6,1.45.2)} śiśīte śṛṅge iti na etat śyateḥ rūpam . \\
\noindent \textbf{(P\_6,1.45.2)} kasya tarhi . \\
\noindent \textbf{(P\_6,1.45.2)} śīṅaḥ . \\
\noindent \textbf{(P\_6,1.45.2)} śyatyarthaḥ vai gamyate . \\
\noindent \textbf{(P\_6,1.45.2)} kaḥ punaḥ śyateḥ arthaḥ . \\
\noindent \textbf{(P\_6,1.45.2)} śyatiḥ niśāne vartate . \\
\noindent \textbf{(P\_6,1.45.2)} śīṅ api śyatyarthe vartate . \\
\noindent \textbf{(P\_6,1.45.2)} katham punaḥ anyaḥ nāma anyasya arthe vartate . \\
\noindent \textbf{(P\_6,1.45.2)} bahvarthāḥ api dhātavaḥ bhavanti iti . \\
\noindent \textbf{(P\_6,1.45.2)} tat yathā . \\
\noindent \textbf{(P\_6,1.45.2)} vapiḥ prakiraṇe dṛṣṭaḥ chedane ca api vartate . \\
\noindent \textbf{(P\_6,1.45.2)} keśān vapati iti . \\
\noindent \textbf{(P\_6,1.45.2)} īḍiḥ studicodanāyācñāsu dṛṣṭaḥ īraṇe ca api vartate . \\
\noindent \textbf{(P\_6,1.45.2)} agniḥ vai itaḥ vṛṣṭim īṭṭe . \\
\noindent \textbf{(P\_6,1.45.2)} marutaḥ amutaḥ cyāvayanti . \\
\noindent \textbf{(P\_6,1.45.2)} karotiḥ ayam abhūtaprādurbhāve dṛṣṭaḥ nirmalīkaraṇe ca api vartate . \\
\noindent \textbf{(P\_6,1.45.2)} pṛṣṭham kuru . \\
\noindent \textbf{(P\_6,1.45.2)} pādau kuru . \\
\noindent \textbf{(P\_6,1.45.2)} unmṛdāna iti gamyate . \\
\noindent \textbf{(P\_6,1.45.2)} nikṣepaṇe ca api dṛśyate . \\
\noindent \textbf{(P\_6,1.45.2)} kaṭe kuru . \\
\noindent \textbf{(P\_6,1.45.2)} ghaṭe kuru . \\
\noindent \textbf{(P\_6,1.45.2)} aśmānam itaḥ kuru . \\
\noindent \textbf{(P\_6,1.45.2)} sthāpaya iti gamyate . \\
\noindent \textbf{(P\_6,1.45.2)} sarveṣām eva parihāraḥ . \\
\noindent \textbf{(P\_6,1.45.2)} śiti iti ucyate na ca atra śitam paśyāmaḥ . \\
\noindent \textbf{(P\_6,1.45.2)} pratyayalakṣaṇena . \\
\noindent \textbf{(P\_6,1.45.2)} na lumatā tasmin iti pratyayalakṣaṇapratiṣedhaḥ . \\
\noindent \textbf{(P\_6,1.45.2)} trādhvam iti luṅi eṣaḥ vyatyayena bhaviṣyati . \\
\noindent \textbf{(P\_6,1.45.2)} atha vā punaḥ astu paryudāsaḥ . \\
\noindent \textbf{(P\_6,1.45.2)} nanu ca uktam aśiti ekādeśe pratiṣedhaḥ ādivattvāt iti . \\
\noindent \textbf{(P\_6,1.45.2)} na eṣaḥ doṣaḥ . \\
\noindent \textbf{(P\_6,1.45.2)} ekādeśaḥ pūrvavidhau sthānivat bhavati iti sthānivadbhāvāt vyavadhānam . \\
\noindent \textbf{(P\_6,1.45.2)} yat api pratyayavidhiḥ iti . \\
\noindent \textbf{(P\_6,1.45.2)} ācāryapravṛttiḥ jñāpayati bhavati ejantebhyaḥ ākārāntalakṣaṇaḥ pratyayavidhiḥ iti yat ayam hvāvāmaḥ ca iti aṇam kabādhanārtham śāsti . \\
\noindent \textbf{(P\_6,1.45.2)} yat api abhyāsarūpam iti : pratyākhyāyate saḥ yogaḥ . \\
\noindent \textbf{(P\_6,1.45.2)} atha api kriyate evam api na doṣaḥ . \\
\noindent \textbf{(P\_6,1.45.2)} katham . \\
\noindent \textbf{(P\_6,1.45.2)} liṭi iti anuvartate . \\
\noindent \textbf{(P\_6,1.45.2)} dvilakārakaḥ ca ayam nirdeśaḥ : liṭi lakārādau iti . \\
\noindent \textbf{(P\_6,1.45.2)} evam ca kṛtvā saḥ api adoṣaḥ bhavati yat uktam āttve eśi upasaṅkhyānam iti . \\
\noindent \textbf{(P\_6,1.45.2)} yat api uktam ayavāyāvām pratiṣedhaḥ ca iti . \\
\noindent \textbf{(P\_6,1.45.2)} śiti pratiṣedhaḥ jñāpakaḥ na ayādayaḥ āttvam bādhante iti . \\
\noindent \textbf{(P\_6,1.45.3)} prātipadikapratiṣedhaḥ . \\
\noindent \textbf{(P\_6,1.45.3)} prātipadikānām pratiṣedhaḥ vaktavyaḥ . \\
\noindent \textbf{(P\_6,1.45.3)} gobhyām , gobhiḥ , naubhyām , naubhiḥ . \\
\noindent \textbf{(P\_6,1.45.3)} saḥ tarhi vaktavyaḥ . \\
\noindent \textbf{(P\_6,1.45.3)} na vaktavyaḥ . \\
\noindent \textbf{(P\_6,1.45.3)} ācāryapravṛttiḥ jñāpayati na prātipadikānām āttvam bhavati iti yat ayam rāyaḥ halaḥ iti āttvam śāsti . \\
\noindent \textbf{(P\_6,1.45.3)} na etat asti jñāpakam . \\
\noindent \textbf{(P\_6,1.45.3)} niyamārtham etat syāt . \\
\noindent \textbf{(P\_6,1.45.3)} rāyaḥ hali eva iti . \\
\noindent \textbf{(P\_6,1.45.3)} yat tarhi ā otaḥ amśasoḥ iti āttvam śāsti . \\
\noindent \textbf{(P\_6,1.45.3)} etasya api asti vacane prayojanam . \\
\noindent \textbf{(P\_6,1.45.3)} ami vṛddhibādhanārtham etat syāt śasi pratiṣedhārtham ca . \\
\noindent \textbf{(P\_6,1.45.3)} tasmāt prātipadikānām pratiṣedhaḥ vaktavyaḥ . \\
\noindent \textbf{(P\_6,1.45.3)} na vaktavyaḥ . \\
\noindent \textbf{(P\_6,1.45.3)} dhātvadhikārāt prātipadikasyāprāptiḥ . \\
\noindent \textbf{(P\_6,1.45.3)} dhātvadhikārāt prātipadikasya āttvam na bhaviṣyati . \\
\noindent \textbf{(P\_6,1.45.3)} dhātoḥ iti vartate . \\
\noindent \textbf{(P\_6,1.45.3)} kva prakṛtam . \\
\noindent \textbf{(P\_6,1.45.3)} liṭi dhātoḥ anabhyāsasya iti . \\
\noindent \textbf{(P\_6,1.45.3)} atha api nivṛttam evam api adoṣaḥ . \\
\noindent \textbf{(P\_6,1.45.3)} upadeśe iti ucyate uddeśaḥ ca prātipadikānām na upadeśaḥ . \\
\noindent \textbf{(P\_6,1.48)} āttve ṇau līyateḥ upasaṅkhyānam pralambhanaśālīnīkaraṇayoḥ . \\
\noindent \textbf{(P\_6,1.48)} āttve ṇau līyateḥ upasaṅkhyānam kartavyam . \\
\noindent \textbf{(P\_6,1.48)} kim prayojanam . \\
\noindent \textbf{(P\_6,1.48)} pralambhane ca arthe śālīnīkaraṇe ca nityam āttvam yathā syāt . \\
\noindent \textbf{(P\_6,1.48)} pralambhane tāvat . \\
\noindent \textbf{(P\_6,1.48)} jaṭābhiḥ ālāpayate . \\
\noindent \textbf{(P\_6,1.48)} śmaśrubhiḥ ālāpayate . \\
\noindent \textbf{(P\_6,1.48)} śālīnīkaraṇe . \\
\noindent \textbf{(P\_6,1.48)} śyenaḥ vārtikam ullāpayate . \\
\noindent \textbf{(P\_6,1.48)} rathī rathinam upalapayate . \\
\noindent \textbf{(P\_6,1.49)} sidhyateḥ ajñānārthasya . \\
\noindent \textbf{(P\_6,1.49)} sidhyateḥ ajñānārthasya iti vaktavyam . \\
\noindent \textbf{(P\_6,1.49)} itarathā hi aniṣṭaprasaṅgaḥ . \\
\noindent \textbf{(P\_6,1.49)} apāralaukike iti ucyamāne aniṣṭam prasajyeta . \\
\noindent \textbf{(P\_6,1.49)} annam sādhayati brāhmaṇebhyaḥ dāsyāmi iti . \\
\noindent \textbf{(P\_6,1.49)} asti punaḥ ayam sidhyatiḥ kva cit anyatra vartate . \\
\noindent \textbf{(P\_6,1.49)} asti iti āha . \\
\noindent \textbf{(P\_6,1.49)} tapaḥ tāpasam sedhayati . \\
\noindent \textbf{(P\_6,1.49)} jñānam asya prakāśayati . \\
\noindent \textbf{(P\_6,1.49)} svāni eva enam karmāṇi sedhayanti . \\
\noindent \textbf{(P\_6,1.49)} jñānam asya prakāśayanti iti arthaḥ . \\
\noindent \textbf{(P\_6,1.50.1)} mīnātyādīnām āttve upadeśavacanam pratyayavidhyartham . \\
\noindent \textbf{(P\_6,1.50.1)} mīnātyādīnām āttve upadeśivadbhāvaḥ vaktavyaḥ . \\
\noindent \textbf{(P\_6,1.50.1)} upadeśāvasthāyām āttvam bhavati iti vaktavyam . \\
\noindent \textbf{(P\_6,1.50.1)} kim prayojanam . \\
\noindent \textbf{(P\_6,1.50.1)} pratyayavidhyartham . \\
\noindent \textbf{(P\_6,1.50.1)} upadeśāvasthāyām āttve kṛte iṣṭaḥ pratyayavidhiḥ yathā syāt . \\
\noindent \textbf{(P\_6,1.50.1)} ke punaḥ pratyayāḥ upadeśivadbhāvam prayojayanti . \\
\noindent \textbf{(P\_6,1.50.1)} kāḥ . \\
\noindent \textbf{(P\_6,1.50.1)} kāḥ tāvat na prayojayanti . \\
\noindent \textbf{(P\_6,1.50.1)} kim kāraṇam . \\
\noindent \textbf{(P\_6,1.50.1)} ecaḥ iti ucyate na ca keṣu ec asti . \\
\noindent \textbf{(P\_6,1.50.1)} ṇaghañyujvidhayaḥ tarhi prayojayanti . \\
\noindent \textbf{(P\_6,1.50.1)} ṇa . \\
\noindent \textbf{(P\_6,1.50.1)} avadāyaḥ . \\
\noindent \textbf{(P\_6,1.50.1)} ātaḥ iti ṇaḥ siddhaḥ bhavati . \\
\noindent \textbf{(P\_6,1.50.1)} ghañ . \\
\noindent \textbf{(P\_6,1.50.1)} avadāyaḥ vartate . \\
\noindent \textbf{(P\_6,1.50.1)} ātaḥ iti ghañ siddhaḥ bhavati . \\
\noindent \textbf{(P\_6,1.50.1)} kim ca bho ātaḥ iti bhañ ucyate . \\
\noindent \textbf{(P\_6,1.50.1)} na khalu api ātaḥ iti ucyate ātaḥ tu vijñāyate . \\
\noindent \textbf{(P\_6,1.50.1)} katham . \\
\noindent \textbf{(P\_6,1.50.1)} aviśeṣeṇa ghañ utsargaḥ . \\
\noindent \textbf{(P\_6,1.50.1)} tasya ivarṇāntāt uvarṇāntāt ca ajapau apavādau . \\
\noindent \textbf{(P\_6,1.50.1)} tatra upadeśāvasthāyām āttve kṛte apavādasya nimittam na asti iti kṛtvā utsargeṇa ghañ siddhaḥ bhavati . \\
\noindent \textbf{(P\_6,1.50.1)} evam ca kṛtvā na ca ātaḥ iti ucyate ātaḥ tu vijñāyate . \\
\noindent \textbf{(P\_6,1.50.1)} yuc . \\
\noindent \textbf{(P\_6,1.50.1)} īṣadavadānam svavadānam . \\
\noindent \textbf{(P\_6,1.50.1)} ātaḥ iti yuc siddhaḥ bhavati . \\
\noindent \textbf{(P\_6,1.50.1)} idam vipratiṣiddham ecaḥ upadeśaḥ iti . \\
\noindent \textbf{(P\_6,1.50.1)} yadi ecaḥ na upadeśe atha upadeśe na ecaḥ . \\
\noindent \textbf{(P\_6,1.50.1)} ecaḥ ca upadeśe ca iti vipratiṣiddham . \\
\noindent \textbf{(P\_6,1.50.1)} na etat vipratiṣiddham . \\
\noindent \textbf{(P\_6,1.50.1)} āha ayam ecaḥ upadeśe iti . \\
\noindent \textbf{(P\_6,1.50.1)} yadi ecaḥ na upadeśe atha upadeśa na ecaḥ . \\
\noindent \textbf{(P\_6,1.50.1)} te vayam viṣayam vijñāsyāmaḥ . \\
\noindent \textbf{(P\_6,1.50.1)} ejviṣaye iti . \\
\noindent \textbf{(P\_6,1.50.1)} tat tarhi upadeśagrahaṇam kartavyam . \\
\noindent \textbf{(P\_6,1.50.1)} na kartavyam . \\
\noindent \textbf{(P\_6,1.50.1)} prakṛtam anuvartate . \\
\noindent \textbf{(P\_6,1.50.1)} kva prakṛtam . \\
\noindent \textbf{(P\_6,1.50.1)} āt ecaḥ upadeśe iti . \\
\noindent \textbf{(P\_6,1.50.1)} tat vai prakṛtiviśeṣaṇam viṣayaviśeṣaṇena ca iha arthaḥ . \\
\noindent \textbf{(P\_6,1.50.1)} na ca anyārtham prakṛtam anyārtham bhavati . \\
\noindent \textbf{(P\_6,1.50.1)} na khalu api anyat prakṛtam anuvartanāt anyat bhavati . \\
\noindent \textbf{(P\_6,1.50.1)} na hi godhā sarpantī sarpaṇāt ahiḥ bhavati . \\
\noindent \textbf{(P\_6,1.50.1)} yat tāvat ucyate na ca anyārtham prakṛtam anyārtham bhavati iti anyārtham api prakṛtam anyārtham bhavati ṭat yathā . \\
\noindent \textbf{(P\_6,1.50.1)} śālyartham kulyāḥ praṇīyante tābhyaḥ ca pāṇīyam pīyate upaśpṛśyate ca śālayaḥ ca bhāvyante . \\
\noindent \textbf{(P\_6,1.50.1)} yat api ucyate na khalu api anyat prakṛtam anuvartanāt anyat bhavati . \\
\noindent \textbf{(P\_6,1.50.1)} na hi godhā sarpantī sarpaṇāt ahiḥ bhavati iti . \\
\noindent \textbf{(P\_6,1.50.1)} bhavet dravyeṣu etat evam syāt . \\
\noindent \textbf{(P\_6,1.50.1)} śabdaḥ tu khalu yena yena viśeṣeṇa abhisambadhyate tasya tasya viśeṣakaḥ bhavati . \\
\noindent \textbf{(P\_6,1.50.1)} tat yatha gauḥ śuklaḥ aśvaḥ ca . \\
\noindent \textbf{(P\_6,1.50.1)} śuklaḥ iti gamyate . \\
\noindent \textbf{(P\_6,1.50.2)} nimimīliyām khalacoḥ pratiṣedhaḥ . \\
\noindent \textbf{(P\_6,1.50.2)} nimimīliyām khalacoḥ pratiṣedhaḥ vaktavyaḥ . \\
\noindent \textbf{(P\_6,1.50.2)} īṣannimayam , sunimayam , nimayaḥ vartate . \\
\noindent \textbf{(P\_6,1.50.2)} mi . \\
\noindent \textbf{(P\_6,1.50.2)} mī . \\
\noindent \textbf{(P\_6,1.50.2)} īṣapramayam , supramayam , pramayaḥ vartate , pramayaḥ . \\
\noindent \textbf{(P\_6,1.50.2)} mī . \\
\noindent \textbf{(P\_6,1.50.2)} lī . \\
\noindent \textbf{(P\_6,1.50.2)} īṣadvilayam , suvilayam , vilayaḥ vartate , vilayaḥ . \\
\noindent \textbf{(P\_6,1.51)} kim idam līyateḥ iti . \\
\noindent \textbf{(P\_6,1.51)} linātilīyatyoḥ yakā nirdeśaḥ . \\
\noindent \textbf{(P\_6,1.56)} hetubhaye iti kimartham . \\
\noindent \textbf{(P\_6,1.56)} kuñcikayā enam bhāyayati . \\
\noindent \textbf{(P\_6,1.56)} ahinā enam bhāyayati . \\
\noindent \textbf{(P\_6,1.56)} hetubhaye iti ucyamāne api atra prāpnoti . \\
\noindent \textbf{(P\_6,1.56)} etat api hi hetubhayam . \\
\noindent \textbf{(P\_6,1.56)} hetubhaye iti na evam vijñāyate . \\
\noindent \textbf{(P\_6,1.56)} hetoḥ bhayam hetubhayam . \\
\noindent \textbf{(P\_6,1.56)} hetubhaye iti . \\
\noindent \textbf{(P\_6,1.56)} katham tarhi . \\
\noindent \textbf{(P\_6,1.56)} hetuḥ eva bhayam hetubhayam . \\
\noindent \textbf{(P\_6,1.56)} hetubhaye iti . \\
\noindent \textbf{(P\_6,1.56)} yadi saḥ eva hetuḥ bhayam bhavati iti . \\
\noindent \textbf{(P\_6,1.58)} ami saṅgrahaṇam . \\
\noindent \textbf{(P\_6,1.58)} ami saṅgrahaṇam . \\
\noindent \textbf{(P\_6,1.58)} kim idam saṅ iti . \\
\noindent \textbf{(P\_6,1.58)} pratyāhāragrahaṇam . \\
\noindent \textbf{(P\_6,1.58)} kva sanniviṣṭānām pratyāhāraḥ . \\
\noindent \textbf{(P\_6,1.58)} sanaḥ prabhṛti ā mahiṅaḥ ṅakārāt . \\
\noindent \textbf{(P\_6,1.58)} kim prayojanam . kvippratiṣedhāṛtham . \\
\noindent \textbf{(P\_6,1.58)} kvibantasya mā bhūt . \\
\noindent \textbf{(P\_6,1.58)} rajjusṛḍbhyām , rajjusṛḍbhiḥ , devadṛgbhyām , devadṛgbhiḥ . \\
\noindent \textbf{(P\_6,1.58)} uktam vā . \\
\noindent \textbf{(P\_6,1.58)} kim uktam . \\
\noindent \textbf{(P\_6,1.58)} dhātoḥ svarūpagrahaṇe tatpratyayavijñānāt siddham iti . \\
\noindent \textbf{(P\_6,1.60)} śīrṣan chandasi prakṛtyantaram . \\
\noindent \textbf{(P\_6,1.60)} śīrṣan chandasi prakṛtyantaram draṣṭavyam . \\
\noindent \textbf{(P\_6,1.60)} kim prayojanam . \\
\noindent \textbf{(P\_6,1.60)} kim prayojanam . \\
\noindent \textbf{(P\_6,1.60)} ādeśapratiṣedhārtham . \\
\noindent \textbf{(P\_6,1.60)} ādeśaḥ mā vijñāyi . \\
\noindent \textbf{(P\_6,1.60)} prakṛtyantaram yathā vijñāyeta . \\
\noindent \textbf{(P\_6,1.60)} kim ca syāt . \\
\noindent \textbf{(P\_6,1.60)} askārāntasya chandasi śravaṇam na syāt . \\
\noindent \textbf{(P\_6,1.60)} śiraḥ me śīryaśaḥ mukham (R: śīryate mukhe ) . \\
\noindent \textbf{(P\_6,1.60)} idam te śiraḥ bhinadmi iti . \\
\noindent \textbf{(P\_6,1.60)} tat vai atharvaṇaḥ śiraḥ . \\
\noindent \textbf{(P\_6,1.61)} ye ca taddhite śirasaḥ ādeśārtham . \\
\noindent \textbf{(P\_6,1.61)} ye ca taddhite iti atra śirasaḥ grahaṇam kartavyam . \\
\noindent \textbf{(P\_6,1.61)} kim prayojanam .ādeśārtham . \\
\noindent \textbf{(P\_6,1.61)} ādeśaḥ yathā vijñāyeta . \\
\noindent \textbf{(P\_6,1.61)} prakṛtyantaram mā vijñāyi . \\
\noindent \textbf{(P\_6,1.61)} kim ca syāt . \\
\noindent \textbf{(P\_6,1.61)} yakārādau taddhite askārāntasya śravaṇam prasajyeta . \\
\noindent \textbf{(P\_6,1.61)} śīrṣaṇyaḥ hi mukhyaḥ bhavati . \\
\noindent \textbf{(P\_6,1.61)} śīrṣaṇyaḥ kharaḥ . \\
\noindent \textbf{(P\_6,1.61)} vā keśeṣu . \\
\noindent \textbf{(P\_6,1.61)} vā keśeṣu śirasaḥ śīrṣanbhāvaḥ vaktavyaḥ . \\
\noindent \textbf{(P\_6,1.61)} śīrṣaṇyāḥ keśāḥ , śirasyāḥ . \\
\noindent \textbf{(P\_6,1.61)} aci śīrṣaḥ . \\
\noindent \textbf{(P\_6,1.61)} aci parataḥ śirasaḥ śīrṣabhāvaḥ vaktavyaḥ . \\
\noindent \textbf{(P\_6,1.61)} hāstiśīrṣiḥ , sthaulyaśīrṣiḥ, pailuśīrṣiḥ . \\
\noindent \textbf{(P\_6,1.61)} chandasi ca . \\
\noindent \textbf{(P\_6,1.61)} chandasi ca śirasaḥ śīrṣabhāvaḥ vaktavyaḥ . \\
\noindent \textbf{(P\_6,1.61)} dve śīrṣe . \\
\noindent \textbf{(P\_6,1.61)} iha hāstiśīrṣyā pailuśīrṣyā iti śirasaḥ grahaṇena grahaṇāt śīrṣanbhāvaḥ prāpnoti . \\
\noindent \textbf{(P\_6,1.61)} astu . \\
\noindent \textbf{(P\_6,1.61)} naḥ taddhite iti ṭilopaḥ bhaviṣyati . \\
\noindent \textbf{(P\_6,1.61)} na sidhyati . \\
\noindent \textbf{(P\_6,1.61)} ye ca abhāvakarmaṇoḥ iti prakṛtibhāvaḥ prasajyeta . \\
\noindent \textbf{(P\_6,1.61)} yadi punaḥ ye aci taddhite iti ucyeta . \\
\noindent \textbf{(P\_6,1.61)} kim kṛtam bhavati . \\
\noindent \textbf{(P\_6,1.61)} iñi śīrṣanbhāve kṛte ṭilopena siddham . \\
\noindent \textbf{(P\_6,1.61)} na evam śakyam . \\
\noindent \textbf{(P\_6,1.61)} iha hi sthūlaśirasaḥ idam sthaulaśīrṣam iti anaṇi iti prakṛtibhāvaḥ prasajyeta . \\
\noindent \textbf{(P\_6,1.61)} tasmāt na evam śakyam . \\
\noindent \textbf{(P\_6,1.61)} na cet evam śirasaḥ grahaṇena grahaṇāt śīrṣanbhāvaḥ prāpnoti . \\
\noindent \textbf{(P\_6,1.61)} pākṣikaḥ eṣaḥ doṣaḥ . \\
\noindent \textbf{(P\_6,1.61)} katarasmin pakṣe . \\
\noindent \textbf{(P\_6,1.61)} ṣyaṅvidhau dvaitam bhavati . \\
\noindent \textbf{(P\_6,1.61)} aṇiñoḥ vā ādeśaḥ ṣyaṅ aṇiñbhyām vā paraḥ iti . \\
\noindent \textbf{(P\_6,1.61)} tat yadā tāvad aṇiñoḥ ādeśaḥ tadā eṣaḥ doṣaḥ . \\
\noindent \textbf{(P\_6,1.61)} yadā hi aṇiñbhyām paraḥ na tadā doṣaḥ bhavati aṇiñbhyām vyavahitatvāt . \\
\noindent \textbf{(P\_6,1.63)} śasprabhṛtiṣu iti ucyate . \\
\noindent \textbf{(P\_6,1.63)} aśasprabhṛtiṣu api dṛśyate . \\
\noindent \textbf{(P\_6,1.63)} śalā doṣaṇī . \\
\noindent \textbf{(P\_6,1.63)} kakut doṣaṇī . \\
\noindent \textbf{(P\_6,1.63)} yācate mahādevaḥ . \\
\noindent \textbf{(P\_6,1.63)} padādiṣu māṃspṛtsnūnām upasaṅkhyānam . \\
\noindent \textbf{(P\_6,1.63)} padādiṣu māṃspṛtsnūnām upasaṅkhyānam kartavyam . \\
\noindent \textbf{(P\_6,1.63)} māṃs . \\
\noindent \textbf{(P\_6,1.63)} yat nīkṣaṇam māṃspacanyāḥ . \\
\noindent \textbf{(P\_6,1.63)} māṃsapacanyāḥ iti prāpte . \\
\noindent \textbf{(P\_6,1.63)} māṃs . \\
\noindent \textbf{(P\_6,1.63)} pṛt. pṛtsu martyam . \\
\noindent \textbf{(P\_6,1.63)} pṛtanāsu martyam iti prāpte . \\
\noindent \textbf{(P\_6,1.63)} pṛt . \\
\noindent \textbf{(P\_6,1.63)} snu . \\
\noindent \textbf{(P\_6,1.63)} na te divaḥ na pṛthivyaḥ adhi snuṣu . \\
\noindent \textbf{(P\_6,1.63)} adhi sānuṣu iti prāpte . \\
\noindent \textbf{(P\_6,1.63)} nas nāsikāyāḥ yattaskṣudreṣu . \\
\noindent \textbf{(P\_6,1.63)} yattaskṣudreṣu parataḥ nāsikāyāḥ nasbhāvaḥ vaktavyaḥ . \\
\noindent \textbf{(P\_6,1.63)} yat . \\
\noindent \textbf{(P\_6,1.63)} nasyam . \\
\noindent \textbf{(P\_6,1.63)} yat . \\
\noindent \textbf{(P\_6,1.63)} tas . \\
\noindent \textbf{(P\_6,1.63)} nastaḥ . \\
\noindent \textbf{(P\_6,1.63)} tas. kṣudra . \\
\noindent \textbf{(P\_6,1.63)} naḥkṣudraḥ . \\
\noindent \textbf{(P\_6,1.63)} avarṇanagarayoḥ iti vaktavyam . \\
\noindent \textbf{(P\_6,1.63)} iha mā bhūt . \\
\noindent \textbf{(P\_6,1.63)} nāsikyaḥ varṇaḥ , nāsikyam nagaram . \\
\noindent \textbf{(P\_6,1.63)} tat tarhi vaktavyam . \\
\noindent \textbf{(P\_6,1.63)} na vaktavyam . \\
\noindent \textbf{(P\_6,1.63)} iha tāvat nāsikyaḥ varṇaḥ iti . \\
\noindent \textbf{(P\_6,1.63)} parimukhādiṣu pāṭhaḥ kariṣyate . \\
\noindent \textbf{(P\_6,1.63)} nāsikyam nagaram iti . \\
\noindent \textbf{(P\_6,1.63)} saṅkāśādiṣu pāṭhaḥ kariṣyate . \\
\noindent \textbf{(P\_6,1.64)} dhātugrahaṇam kimartham . \\
\noindent \textbf{(P\_6,1.64)} iha mā bhūt : ṣoḍan , ṣaṇḍaḥ , ṣoḍikaḥ . \\
\noindent \textbf{(P\_6,1.64)} atha ādigrahaṇam kimartham . \\
\noindent \textbf{(P\_6,1.64)} iha mā bhūt : peṣṭā peṣṭum . \\
\noindent \textbf{(P\_6,1.64)} na etat asti prayojanam . \\
\noindent \textbf{(P\_6,1.64)} astu atra satvam . \\
\noindent \textbf{(P\_6,1.64)} satve kṛte iṇaḥ uttarasya ādeśasakārasya iti ṣatvam bhaviṣyati . \\
\noindent \textbf{(P\_6,1.64)} idam tarhi : laṣitā laṣitum . \\
\noindent \textbf{(P\_6,1.64)} idam ca api udārharaṇam : peṣṭā peṣṭum . \\
\noindent \textbf{(P\_6,1.64)} nanu ca uktam astu atra satvam . \\
\noindent \textbf{(P\_6,1.64)} satve kṛte iṇaḥ uttarasya ādeśasakārasya iti ṣatvam bhaviṣyati iti . \\
\noindent \textbf{(P\_6,1.64)} na evam śakyam . \\
\noindent \textbf{(P\_6,1.64)} iha hi pekṣyati iti ṣatvasya asiddhatvāt ṣaḍhoḥ kaḥ si iti katvam na syāt . \\
\noindent \textbf{(P\_6,1.64)} sādeśe subdhātuṣṭhivuṣvaṣkatīnām pratiṣedhaḥ . \\
\noindent \textbf{(P\_6,1.64)} sādeśe subdhātuṣṭhivuṣvaṣkatīnām pratiṣedhaḥ vaktavyaḥ . \\
\noindent \textbf{(P\_6,1.64)} subdhātu : ṣoḍīyati ṣaṇḍīyati . \\
\noindent \textbf{(P\_6,1.64)} ṣṭhivu : ṣṭhīvati . \\
\noindent \textbf{(P\_6,1.64)} ṣvaṣk : ṣvaṣkate . \\
\noindent \textbf{(P\_6,1.64)} subdhātūnām tāvat na vaktavyaḥ . \\
\noindent \textbf{(P\_6,1.64)} upadeśe iti vartate uddeśaḥ ca prātipadikānām na upadeśaḥ . \\
\noindent \textbf{(P\_6,1.64)} yadi evam na arthaḥ dhātugrahaṇena . \\
\noindent \textbf{(P\_6,1.64)} kasmāt na bhavati ṣoḍan , ṣaṇḍaḥ , ṣoḍikaḥ iti . \\
\noindent \textbf{(P\_6,1.64)} upadeśe iti vartate uddeśaḥ ca prātipadikānām na upadeśaḥ . \\
\noindent \textbf{(P\_6,1.64)} ṣṭhiveḥ api dvitīyaḥ varṇaḥ ṭhakāraḥ . \\
\noindent \textbf{(P\_6,1.64)} yadi ṭhakāraḥ teṣṭhīvyate iti na sidhyati . \\
\noindent \textbf{(P\_6,1.64)} evam tarhi thakāraḥ . \\
\noindent \textbf{(P\_6,1.64)} yadi thakāraḥ ṭuṣṭhyūṣati ṭeṣṭhīvyatie iti na sidhyati . \\
\noindent \textbf{(P\_6,1.64)} evam tarhi dvau imau ṣṭhivū . \\
\noindent \textbf{(P\_6,1.64)} akasya dvitīyaḥ varṇaḥ ṭhakāraḥ aparasya thakāraḥ . \\
\noindent \textbf{(P\_6,1.64)} yasya thakāraḥ tasya satvam prāpnoti . \\
\noindent \textbf{(P\_6,1.64)} evam tarhi dvau imau dviṣakārau ṣṭhivuṣvaṣkatī . \\
\noindent \textbf{(P\_6,1.64)} kim kṛtam bhavati . \\
\noindent \textbf{(P\_6,1.64)} pūrvasya satve kṛte pareṇa sannipāte ṣṭutvam bhaviṣyati . \\
\noindent \textbf{(P\_6,1.64)} na evam śakyam . \\
\noindent \textbf{(P\_6,1.64)} iha hi śvaliṭ ṣṭhīvati madhuliṭ ṣvaṣkate ṣṭutvasya asiddhatvāt ḍaḥ si dhuṭ iti dhuṭ prasajyeta . \\
\noindent \textbf{(P\_6,1.64)} evam tarhi yakārādī dviṣakārau ṣṭhivuṣvaṣkatī . \\
\noindent \textbf{(P\_6,1.64)} kim yakāraḥ na śrūyate . \\
\noindent \textbf{(P\_6,1.64)} luptanirdiṣṭaḥ yakāraḥ . \\
\noindent \textbf{(P\_6,1.64)} atha kimartham ṣakāram upadiśya tasya sakāraḥ ādeśaḥ kriyate na sakāraḥ eva upadiśyeta . \\
\noindent \textbf{(P\_6,1.64)} laghvartham iti āha . \\
\noindent \textbf{(P\_6,1.64)} katham . \\
\noindent \textbf{(P\_6,1.64)} aviśeṣeṇa ayam ṣakāram upadiśya sakāram ādeśam uktvā laghunā upāyena ṣatvam nirvartayati ādeśapratyayayoḥ iti . \\
\noindent \textbf{(P\_6,1.64)} itarathā hi yeṣām ṣatvam iṣyate teṣām tatra grahaṇam kartavyam syāt . \\
\noindent \textbf{(P\_6,1.64)} ke punaḥ ṣopadeśāḥ dhātavaḥ . \\
\noindent \textbf{(P\_6,1.64)} paṭhitavyāḥ . \\
\noindent \textbf{(P\_6,1.64)} kaḥ atra bhavataḥ puruṣakāraḥ . \\
\noindent \textbf{(P\_6,1.64)} yadi antareṇa pāṭham kim cit śakyate vaktum tat ucyatām . \\
\noindent \textbf{(P\_6,1.64)} antareṇa api pāṭham kim cit śakyate vaktum . \\
\noindent \textbf{(P\_6,1.64)} katham . \\
\noindent \textbf{(P\_6,1.64)} ajdantyaparāḥ sādayaḥ ṣopadeśāḥ smiṅsvadisvidisvañjisvapayaḥ ca sṛpisṛjistṛstyāsekṛsṛvarjam . \\
\noindent \textbf{(P\_6,1.65)} atha kimartham ṇakāram upadiśya tasya nakāraḥ ādeśaḥ kriyate na nakāraḥ eva upadiśyeta . \\
\noindent \textbf{(P\_6,1.65)} laghvartham iti āha . \\
\noindent \textbf{(P\_6,1.65)} katham . \\
\noindent \textbf{(P\_6,1.65)} aviśeṣeṇa ayam ṇakāram upadiśya nakāram ādeśam uktvā laghunā upāyena ṇatvam nirvartayati upasargāt asamāse api ṇopadeśasya iti . \\
\noindent \textbf{(P\_6,1.65)} itarathā hi yeṣām ṇatvam iṣyate teṣām tatra grahaṇam kartavyam syāt . \\
\noindent \textbf{(P\_6,1.65)} ke punaḥ ṇopadeśāḥ dhātavaḥ . \\
\noindent \textbf{(P\_6,1.65)} paṭhitavyāḥ . \\
\noindent \textbf{(P\_6,1.65)} kaḥ atra bhavataḥ puruṣakāraḥ . \\
\noindent \textbf{(P\_6,1.65)} yadi antareṇa pāṭham kim cit śakyate vaktum tat ucyatām . \\
\noindent \textbf{(P\_6,1.65)} antareṇa api pāṭham kim cit śakyate vaktum . \\
\noindent \textbf{(P\_6,1.65)} katham . \\
\noindent \textbf{(P\_6,1.65)} sarve nādayaḥ ṇopadeśāḥ nṛtinandinardinakkināṭināthṛnādhṛnṝvarjam . \\
\noindent \textbf{(P\_6,1.66.1)} vyoḥ lope kvau upasaṅkhyānam . \\
\noindent \textbf{(P\_6,1.66.1)} vyoḥ lope kvau upasaṅkhyānam kartavyam iha api yathā syāt . \\
\noindent \textbf{(P\_6,1.66.1)} kaṇḍūyateḥ apratyayaḥ kaṇḍūḥ iti . \\
\noindent \textbf{(P\_6,1.66.1)} kim punaḥ kāraṇam na sidhyati . \\
\noindent \textbf{(P\_6,1.66.1)} vali iti ucyate na ca atra valādim paśyāmaḥ . \\
\noindent \textbf{(P\_6,1.66.1)} nanu ca ayam kvip eva valādiḥ bhavati . \\
\noindent \textbf{(P\_6,1.66.1)} kviblope kṛte valādyabhāvāt na prāpnoti . \\
\noindent \textbf{(P\_6,1.66.1)} idam iha sampradhāryam . \\
\noindent \textbf{(P\_6,1.66.1)} kviblopaḥ kriyatām yalopaḥ iti . \\
\noindent \textbf{(P\_6,1.66.1)} kim atra kartavyam . \\
\noindent \textbf{(P\_6,1.66.1)} paratvāt kviblopaḥ nityatvāt ca . \\
\noindent \textbf{(P\_6,1.66.1)} nityaḥ khalu api kviblopaḥ . \\
\noindent \textbf{(P\_6,1.66.1)} kṛte api yalope prāpnoti akṛte api prāpnoti . \\
\noindent \textbf{(P\_6,1.66.1)} nityatvāt paratvāt ca kviblopaḥ . \\
\noindent \textbf{(P\_6,1.66.1)} kviblope kṛte valādyabhāvāt yalopaḥ na prāpnoti . \\
\noindent \textbf{(P\_6,1.66.1)} evam tarhi pratyayalakṣaṇena bhaviṣyati . \\
\noindent \textbf{(P\_6,1.66.1)} varṇāśraye na asti pratyayalakṣaṇam . \\
\noindent \textbf{(P\_6,1.66.1)} yadi vā kāni cit varṇāśrayāṇi pratyayalakṣaṇena bhavanti tathā idam api bhaviṣyati . \\
\noindent \textbf{(P\_6,1.66.1)} atha vā evam vakṣyāmi . \\
\noindent \textbf{(P\_6,1.66.1)} lopaḥ vyoḥ vali . \\
\noindent \textbf{(P\_6,1.66.1)} tataḥ veḥ . \\
\noindent \textbf{(P\_6,1.66.1)} vyantayoḥ ca vyoḥ lopaḥ bhavati . \\
\noindent \textbf{(P\_6,1.66.1)} tataḥ apṛktasya . \\
\noindent \textbf{(P\_6,1.66.1)} apṛktasya ca lopaḥ bhavati . \\
\noindent \textbf{(P\_6,1.66.1)} veḥ iti eva . \\
\noindent \textbf{(P\_6,1.66.2)} valopāprasiddhiḥ ūḍbhāvavacanāt . \\
\noindent \textbf{(P\_6,1.66.2)} valopasya aprasiddhiḥ . \\
\noindent \textbf{(P\_6,1.66.2)} āsremāṇam , jīradānuḥ iti . \\
\noindent \textbf{(P\_6,1.66.2)} kim kāraṇam . \\
\noindent \textbf{(P\_6,1.66.2)} ūḍbhāvavacanāt . \\
\noindent \textbf{(P\_6,1.66.2)} cchvoḥ śūṭ anunāsike ca iti ūṭh prāpnoti . \\
\noindent \textbf{(P\_6,1.66.2)} atiprasaṅgaḥ vraścādiṣu . \\
\noindent \textbf{(P\_6,1.66.2)} vraścādiṣu ca atiprasaṅgaḥ bhavati . \\
\noindent \textbf{(P\_6,1.66.2)} iha api prāpnoti . \\
\noindent \textbf{(P\_6,1.66.2)} vraścanaḥ , vrīhiḥ , vraṇaḥ iti . \\
\noindent \textbf{(P\_6,1.66.2)} upadeśasāmarthyāt siddham . \\
\noindent \textbf{(P\_6,1.66.2)} upadeśasāmarthyāt vraścādiṣu lopaḥ na bhaviṣyati . \\
\noindent \textbf{(P\_6,1.66.2)} upadeśasāmarthyāt siddham iti cet samprasāraṇahalādiśeṣeṣu sāmarthyam . \\
\noindent \textbf{(P\_6,1.66.2)} upadeśasāmarthyāt siddham iti cet asti anyat upadeśavacane prayojanam . \\
\noindent \textbf{(P\_6,1.66.2)} samprasāraṇahalādiśeṣeṣu kṛteṣu vakārasya śravaṇam yathā syāt . \\
\noindent \textbf{(P\_6,1.66.2)} (R: vṛkṇaḥ ) vṛkṇavān , (R vṛścati ) vivraściṣati iti . \\
\noindent \textbf{(P\_6,1.66.2)} na vā bahiraṅgalakṣaṇatvāt . \\
\noindent \textbf{(P\_6,1.66.2)} na vā etat prayojanam asti . \\
\noindent \textbf{(P\_6,1.66.2)} kim kāraṇam . \\
\noindent \textbf{(P\_6,1.66.2)} bahiraṅgalakṣaṇatvāt . \\
\noindent \textbf{(P\_6,1.66.2)} bahiraṅgāḥ samprasāraṇahalādiśeṣāḥ . \\
\noindent \textbf{(P\_6,1.66.2)} antaraṅgaḥ lopaḥ . \\
\noindent \textbf{(P\_6,1.66.2)} asiddham bahiraṅgam antaraṅge . \\
\noindent \textbf{(P\_6,1.66.2)} anārambhaḥ vā . \\
\noindent \textbf{(P\_6,1.66.2)} anārambhaḥ vā punaḥ valopasya nyāyyaḥ . \\
\noindent \textbf{(P\_6,1.66.2)} katham āsremāṇam , jīradānuḥ iti . \\
\noindent \textbf{(P\_6,1.66.2)} āsremāṇam jīradānuḥ iti varṇalopāt . \\
\noindent \textbf{(P\_6,1.66.2)} āsremāṇam , jīradānuḥ iti chāndasāt varṇalopāt siddham . \\
\noindent \textbf{(P\_6,1.66.2)} yathā saṃsphānaḥ gayasphānaḥ . \\
\noindent \textbf{(P\_6,1.66.2)} tat yathā saṃsphayanaḥ , saṃsphānaḥ , gayasphānaḥ, gayasphānaḥ iti . \\
\noindent \textbf{(P\_6,1.67)} darvijāgṛvyoḥ pratiṣedhaḥ vaktavyaḥ : darviḥ , jāgṛviḥ . \\
\noindent \textbf{(P\_6,1.67)} kim ucyate darvijāgṛvyoḥ pratiṣedhaḥ vaktavyaḥ iti yadā apṛktasya iti ucyati . \\
\noindent \textbf{(P\_6,1.67)} bhavati vai kim cit ācāryāḥ kriyamāṇam api codayanti . \\
\noindent \textbf{(P\_6,1.67)} tat vā kartavyam darvijāgṛvyoḥ vā pratiṣedhaḥ vaktavyaḥ .veḥ lope darvijāgṛvyoḥ apratiṣedhaḥ anunāsikaparatvāt . \\
\noindent \textbf{(P\_6,1.67)} veḥ lope darvijāgṛvyoḥ apratiṣedhaḥ . \\
\noindent \textbf{(P\_6,1.67)} anarthakaḥ pratiṣedhaḥ apratiṣedhaḥ . \\
\noindent \textbf{(P\_6,1.67)} lopaḥ kasmāt na bhavati . \\
\noindent \textbf{(P\_6,1.67)} anunāsikaparatvāt . \\
\noindent \textbf{(P\_6,1.67)} anunāsikaparasya viśabdasya grahaṇam na ca atra anunāsikaparaḥ viśabdaḥ . \\
\noindent \textbf{(P\_6,1.67)} śuddhaparaḥ ca atra viśabdaḥ . \\
\noindent \textbf{(P\_6,1.67)} yadi anunāsikaparasya viśabdasya grahaṇam iti ucyate ghṛtaspṛk, dalaspṛk , atra na prāpnoti . \\
\noindent \textbf{(P\_6,1.67)} na hi etasmāt viśabdāt anunāsikam param paśyāmaḥ . \\
\noindent \textbf{(P\_6,1.67)} anunāsikaparatvāt iti na evam vijñāyate . \\
\noindent \textbf{(P\_6,1.67)} anunāsikaḥ paraḥ asmāt saḥ ayam anunāsikaparaḥ , anunāsikaparatvāt iti . \\
\noindent \textbf{(P\_6,1.67)} katham tarhi . \\
\noindent \textbf{(P\_6,1.67)} anunāsikaḥ paraḥ asmin saḥ ayam anunāsikaparaḥ , anunāsikaparatvāt iti . \\
\noindent \textbf{(P\_6,1.67)} evam api priyadarvi , atra prāpnoti . \\
\noindent \textbf{(P\_6,1.67)} asiddhaḥ atra anunāsikaḥ . \\
\noindent \textbf{(P\_6,1.67)} evam api dhātvantasya pratiṣedhaḥ vaktavyaḥ . \\
\noindent \textbf{(P\_6,1.67)} ivi divi dhivi . \\
\noindent \textbf{(P\_6,1.67)} dhātvantasya ca arthavadgrahaṇāt . \\
\noindent \textbf{(P\_6,1.67)} arthavataḥ viśabdasya grahaṇam . \\
\noindent \textbf{(P\_6,1.67)} na dhātvantaḥ arthavān . \\
\noindent \textbf{(P\_6,1.67)} vasya vā anunāsikatvāt siddham . \\
\noindent \textbf{(P\_6,1.67)} atha vā vakārasya eva idam anunāsikasya grahaṇam . \\
\noindent \textbf{(P\_6,1.67)} santi hi yaṇaḥ sānunāsikāḥ niranunāsikāḥ ca . \\
\noindent \textbf{(P\_6,1.68)} yadi punaḥ ayam apṛktalopaḥ saṃyogāntalopaḥ vijñāyeta . \\
\noindent \textbf{(P\_6,1.68)} kim kṛtam bhavati . \\
\noindent \textbf{(P\_6,1.68)} dvihalapṛktagrahaṇam tisyoḥ ca grahaṇam na kartavyam bhavati . \\
\noindent \textbf{(P\_6,1.68)} halantāt apṛktalopaḥ saṃyogāntalopaḥ cet nalopābhāvaḥ yathā pacan iti . \\
\noindent \textbf{(P\_6,1.68)} halantāt apṛktalopaḥ saṃyogāntalopaḥ cet nalopābhāvaḥ . \\
\noindent \textbf{(P\_6,1.68)} rājā takṣā . \\
\noindent \textbf{(P\_6,1.68)} saṃyogāntalopasya asiddhatvāt nalopaḥ na prāpnoti yathā pacan iti . \\
\noindent \textbf{(P\_6,1.68)} tat yathā pacan, yajan iti atra saṃyogāntalopasya asiddhatvāt nalopaḥ na bhavati . \\
\noindent \textbf{(P\_6,1.68)} na eṣaḥ doṣaḥ . \\
\noindent \textbf{(P\_6,1.68)} ācāryapravṛttiḥ jñāpayati siddhaḥ saṃyogāntalopaḥ nalope iti yat ayam na ṅisambuddhyoḥ iti sambuddhau pratiṣedham śāsti . \\
\noindent \textbf{(P\_6,1.68)} iha api tarhi prāpnoti pacan, yajan . \\
\noindent \textbf{(P\_6,1.68)} tulyajātīyasya jñāpakam bhavati . \\
\noindent \textbf{(P\_6,1.68)} kaḥ ca tulyajātīyaḥ . \\
\noindent \textbf{(P\_6,1.68)} yaḥ sambuddhau anantaraḥ . \\
\noindent \textbf{(P\_6,1.68)} vasvādiṣu datvam saṃyogādilopabalīyastvāt . \\
\noindent \textbf{(P\_6,1.68)} vasvādiṣu datvam na sidhyati . \\
\noindent \textbf{(P\_6,1.68)} ukhāsrat , parṇadvhat . \\
\noindent \textbf{(P\_6,1.68)} kim kāraṇam . \\
\noindent \textbf{(P\_6,1.68)} saṃyogādilopabalīyastvāt . saṃyogāntalopāt saṃyogādilopaḥ balīyān . \\
\noindent \textbf{(P\_6,1.68)} yathā kūṭataṭ iti . \\
\noindent \textbf{(P\_6,1.68)} tat yathā kūṭataṭ , kāṣṭhataṭ iti atra saṃyogāntalopāt saṃyogādilopaḥ balīyān bhavati . \\
\noindent \textbf{(P\_6,1.68)} nanu ca datve kṛte na bhaviṣyati . \\
\noindent \textbf{(P\_6,1.68)} asiddham datvam . \\
\noindent \textbf{(P\_6,1.68)} tasya asiddhatvāt prāpnoti . \\
\noindent \textbf{(P\_6,1.68)} siddhakāṇḍe paṭhitam vasvādiṣu datvam sau dīrghatve iti . \\
\noindent \textbf{(P\_6,1.68)} tatra sau dīrghatvagrahaṇam na kariṣyate . \\
\noindent \textbf{(P\_6,1.68)} vasvādiṣu datvam iti eva . \\
\noindent \textbf{(P\_6,1.68)} evam api apadāntatvāt na prāpnoti . \\
\noindent \textbf{(P\_6,1.68)} atha sau api padam bhavati rājā takṣā nalope kṛte vibhakteḥ śravaṇam prāpnoti . \\
\noindent \textbf{(P\_6,1.68)} sā eṣā ubhayataspāśā rajjuḥ bhavati . \\
\noindent \textbf{(P\_6,1.68)} rāttalopaḥ niyamavacanāt . \\
\noindent \textbf{(P\_6,1.68)} rāt tasya lopaḥ vaktavyaḥ . \\
\noindent \textbf{(P\_6,1.68)} abibhaḥ bhavān . \\
\noindent \textbf{(P\_6,1.68)} ajāgaḥ bhavān . \\
\noindent \textbf{(P\_6,1.68)} kim punaḥ kāraṇam na sidhyati . \\
\noindent \textbf{(P\_6,1.68)} niyamavacanāt . \\
\noindent \textbf{(P\_6,1.68)} rāt sasya iti etasmāt niyamāt na prāpnoti . \\
\noindent \textbf{(P\_6,1.68)} na eṣaḥ doṣaḥ . \\
\noindent \textbf{(P\_6,1.68)} rāt sasya iti atra takāraḥ api nirdiśyate . \\
\noindent \textbf{(P\_6,1.68)} yadi evam kīrtayateḥ apratyayaḥ kīḥ iti prāpnoti . \\
\noindent \textbf{(P\_6,1.68)} kīrt iti ca iṣyate . \\
\noindent \textbf{(P\_6,1.68)} yathālakṣaṇam aprayukte . \\
\noindent \textbf{(P\_6,1.68)} roḥ uttvam ca . \\
\noindent \textbf{(P\_6,1.68)} roḥ uttvam ca vaktavyam . \\
\noindent \textbf{(P\_6,1.68)} abhinaḥ atra, acchinaḥ atra . \\
\noindent \textbf{(P\_6,1.68)} saṃyogāntalopasya asiddhatvāt ataḥ ati iti uttrvam na prāpnoti . \\
\noindent \textbf{(P\_6,1.68)} na vā saṃyogāntalopasya uttve siddhatvāt . \\
\noindent \textbf{(P\_6,1.68)} na vā vaktavyam . \\
\noindent \textbf{(P\_6,1.68)} kim kāraṇam . \\
\noindent \textbf{(P\_6,1.68)} saṃyogāntalopasya uttve siddhatvāt . \\
\noindent \textbf{(P\_6,1.68)} saṃyogāntalopaḥ uttve siddhaḥ bhavati . \\
\noindent \textbf{(P\_6,1.68)} yathā harivaḥ medinam iti . tat yathā harivaḥ medinam tvā iti atra saṃyogāntalopaḥ uttve siddhaḥ bhavati . \\
\noindent \textbf{(P\_6,1.68)} saḥ eva darhi doṣaḥ sā eṣā ubhayataspāśā . \\
\noindent \textbf{(P\_6,1.68)} tasmāt aśakyaḥ apṛktalopaḥ saṃyogāntalopaḥ vijñātum . \\
\noindent \textbf{(P\_6,1.68)} na cet vijñāyate dvihalapṛktagrahaṇam tisyoḥ ca grahaṇam kartavyam eva . \\
\noindent \textbf{(P\_6,1.69.1)} sambuddhilope ḍatarādibhyaḥ pratiṣedhaḥ . \\
\noindent \textbf{(P\_6,1.69.1)} sambuddhilope ḍatarādibhyaḥ pratiṣedhaḥ vaktavyaḥ . \\
\noindent \textbf{(P\_6,1.69.1)} he katarat, he katamat . \\
\noindent \textbf{(P\_6,1.69.1)} kim ucyate ḍatarādibhyaḥ pratiṣedhaḥ vaktavyaḥ iti yadā apṛktasya iti anuvartate . \\
\noindent \textbf{(P\_6,1.69.1)} apṛktādhikārasya nivṛttatvāt . \\
\noindent \textbf{(P\_6,1.69.1)} nivṛttaḥ apṛktādhikāraḥ . \\
\noindent \textbf{(P\_6,1.69.1)} kim ḍatarādibhyaḥ pratiṣedham vakṣyāmi iti apṛktādhikāraḥ nivartyate . \\
\noindent \textbf{(P\_6,1.69.1)} na iti āha . \\
\noindent \textbf{(P\_6,1.69.1)} tat ca amartham . \\
\noindent \textbf{(P\_6,1.69.1)} saḥ ca avaśyam apṛktādhikāraḥ nivartyaḥ . \\
\noindent \textbf{(P\_6,1.69.1)} kimartham . \\
\noindent \textbf{(P\_6,1.69.1)} amartham . \\
\noindent \textbf{(P\_6,1.69.1)} amaḥ lopaḥ yathā syāt . \\
\noindent \textbf{(P\_6,1.69.1)} he kuṇḍa , he pīṭha . \\
\noindent \textbf{(P\_6,1.69.1)} nivṛtte api apṛktādhikāre amaḥ lopaḥ na prāpnoti . \\
\noindent \textbf{(P\_6,1.69.1)} na hi lopaḥ sarvāpahārī . \\
\noindent \textbf{(P\_6,1.69.1)} mā bhūt sarvasya lopaḥ . \\
\noindent \textbf{(P\_6,1.69.1)} alaḥ antyasya vidhayaḥ bhavanti iti antyasya lope kṛte dvayoḥ akārayoḥ pararūpeṇa siddham rūpam syāt he kuṇḍa , he pīṭha iti . \\
\noindent \textbf{(P\_6,1.69.1)} yadi etat labhyeta kṛtam syāt . \\
\noindent \textbf{(P\_6,1.69.1)} tat tu na labhyate . \\
\noindent \textbf{(P\_6,1.69.1)} kim kāraṇam . \\
\noindent \textbf{(P\_6,1.69.1)} atra hi tasmāt iti uttarasya ādeḥ parasya iti akārasya lopaḥ prāpnoti . \\
\noindent \textbf{(P\_6,1.69.1)} akāralope ca sati makāre ataḥ dīrghaḥ yañi supi ca iti dīrghatve he kuṇḍām , he pīṭhām iti etat rūpam prasajyeta . \\
\noindent \textbf{(P\_6,1.69.1)} evam tarhi halaḥ lopaḥ sambuddhilopaḥ . \\
\noindent \textbf{(P\_6,1.69.1)} tat halgrahaṇam kartavyam . \\
\noindent \textbf{(P\_6,1.69.1)} na kartavyam . \\
\noindent \textbf{(P\_6,1.69.1)} prakṛtam anuvartate . \\
\noindent \textbf{(P\_6,1.69.1)} kva prakṛtam . \\
\noindent \textbf{(P\_6,1.69.1)} halṅyābbhyaḥ dīrghāt sutisi apṛktam hal iti . \\
\noindent \textbf{(P\_6,1.69.1)} tat vai prathamānirdiṣṭam ṣaṣṭhīnirdiṣṭena ca iha arthaḥ . \\
\noindent \textbf{(P\_6,1.69.1)} na eṣaḥ doṣaḥ . \\
\noindent \textbf{(P\_6,1.69.1)} eṅ hrasvāt iti eṣā pañcamī hal iti asyāḥ prathamāyāḥ ṣaṣṭhīm prakalpayiṣyati tasmāt iti uttarasya iti . \\
\noindent \textbf{(P\_6,1.69.1)} evam api prathamayoḥ pūrvasavarṇadīrghatve kṛte he pīṭhā iti etat rūpam prasajyeta . \\
\noindent \textbf{(P\_6,1.69.1)} ami pūrvatvam atra bādhakam bhaviṣyati . \\
\noindent \textbf{(P\_6,1.69.1)} ami iti ucyate na ca atra amam paśyāmaḥ . \\
\noindent \textbf{(P\_6,1.69.1)} ekadeśavikṛtam ananyavat bhavati iti . \\
\noindent \textbf{(P\_6,1.69.1)} atha vā idam iha sampradhāryam . \\
\noindent \textbf{(P\_6,1.69.1)} sambuddhilopaḥ kriyatām ekādeśaḥ iti . \\
\noindent \textbf{(P\_6,1.69.1)} kim atra kartavyam . \\
\noindent \textbf{(P\_6,1.69.1)} paratvāt ekādeśaḥ . \\
\noindent \textbf{(P\_6,1.69.1)} evam api ekādeśe kṛte vyapavargābhāvāt sambuddhilopaḥ na prāpnoti . \\
\noindent \textbf{(P\_6,1.69.1)} antādivadbhāvena vyapavargaḥ bhaviṣyati . \\
\noindent \textbf{(P\_6,1.69.1)} ubhayataḥ āśraye na antādivat . \\
\noindent \textbf{(P\_6,1.69.1)} na ubhayataḥ āśrayaḥ kariṣyate . \\
\noindent \textbf{(P\_6,1.69.1)} katham . \\
\noindent \textbf{(P\_6,1.69.1)} na evam vijñāyate . \\
\noindent \textbf{(P\_6,1.69.1)} hrasvāt uttarasyāḥ sambuddheḥ lopaḥ bhavati iti . \\
\noindent \textbf{(P\_6,1.69.1)} katham tarhi . \\
\noindent \textbf{(P\_6,1.69.1)} hrasvāt uttarasya halaḥ lopaḥ bhavati saḥ cet sambuddheḥ iti . \\
\noindent \textbf{(P\_6,1.69.1)} saḥ tarhi pratiṣedhaḥ vaktavyaḥ . \\
\noindent \textbf{(P\_6,1.69.1)} na vaktavyaḥ . \\
\noindent \textbf{(P\_6,1.69.1)} uktam vā . \\
\noindent \textbf{(P\_6,1.69.1)} kim uktam . \\
\noindent \textbf{(P\_6,1.69.1)} siddham anunāsikopadhatvāt iti . \\
\noindent \textbf{(P\_6,1.69.1)} evam api dalopaḥ sādhīyaḥ prāpnoti . \\
\noindent \textbf{(P\_6,1.69.1)} dukkaraṇāt vā . \\
\noindent \textbf{(P\_6,1.69.1)} atha vā duk ḍatarādīnām iti vakṣyāmi . \\
\noindent \textbf{(P\_6,1.69.1)} ḍitkaraṇāt vā . \\
\noindent \textbf{(P\_6,1.69.1)} atha vā ḍit ayam śabdaḥ kariṣyate . \\
\noindent \textbf{(P\_6,1.69.1)} saḥ tarhi ḍakāraḥ kartavyaḥ . \\
\noindent \textbf{(P\_6,1.69.1)} na kartavyaḥ . \\
\noindent \textbf{(P\_6,1.69.1)} kriyate nyāse eva . \\
\noindent \textbf{(P\_6,1.69.1)} dviḍakāraḥ nirdeśaḥ adḍ ḍatarādibhyaḥ iti . \\
\noindent \textbf{(P\_6,1.69.1)} evam api lopaḥ prāpnoti . \\
\noindent \textbf{(P\_6,1.69.1)} vihitaviśeṣaṇam hrasvagrahaṇam . \\
\noindent \textbf{(P\_6,1.69.1)} yasmāt hrasvāt sambuddhiḥ vihitā iti . \\
\noindent \textbf{(P\_6,1.69.2)} apṛktasambuddhilopābhyām luk . \\
\noindent \textbf{(P\_6,1.69.2)} apṛktasambuddhilopābhyām luk bhavati vipratiṣedhena . \\
\noindent \textbf{(P\_6,1.69.2)} apṛktalopasya avakāśaḥ gomān , yavamān . \\
\noindent \textbf{(P\_6,1.69.2)} lukaḥ avakāśaḥ trapu, jatu . \\
\noindent \textbf{(P\_6,1.69.2)} iha ubhayam prāpnoti . \\
\noindent \textbf{(P\_6,1.69.2)} tat brāhmaṇakulam , yat brāhmaṇakulam . \\
\noindent \textbf{(P\_6,1.69.2)} sambuddhilopasya avakāśaḥ he agne, he vāyo . \\
\noindent \textbf{(P\_6,1.69.2)} lukaḥ saḥ eva . \\
\noindent \textbf{(P\_6,1.69.2)} iha ubhayam prāpnoti . \\
\noindent \textbf{(P\_6,1.69.2)} he trapu , he jatu . \\
\noindent \textbf{(P\_6,1.69.2)} luk bhavati vipratiṣedhena . \\
\noindent \textbf{(P\_6,1.69.2)} saḥ tarhi vipratiṣedhaḥ vaktavyaḥ . \\
\noindent \textbf{(P\_6,1.69.2)} na vā lopalukoḥ lugavadhāraṇāt yathā anaḍuhyate iti . \\
\noindent \textbf{(P\_6,1.69.2)} na vā arthaḥ vipratiṣedhena . \\
\noindent \textbf{(P\_6,1.69.2)} kim kāraṇam . \\
\noindent \textbf{(P\_6,1.69.2)} lopalukoḥ lugavadhāraṇāt . \\
\noindent \textbf{(P\_6,1.69.2)} lopalukoḥ hi luk avadhāryate . \\
\noindent \textbf{(P\_6,1.69.2)} luk lopayaṇayavāyāvekādeśebhyaḥ . \\
\noindent \textbf{(P\_6,1.69.2)} yathā anaḍuhyate iti . \\
\noindent \textbf{(P\_6,1.69.2)} tat yathā anaḍvān iva ācarati anaḍuhyate iti atra lopalukoḥ luk avadhāryate evam iha api . \\
\noindent \textbf{(P\_6,1.70)} ayam yogaḥ śakyaḥ avaktum . \\
\noindent \textbf{(P\_6,1.70)} katham agne trī te vajinā trī sadhasthā , ta tā piṇḍānām iti . \\
\noindent \textbf{(P\_6,1.70)} pūrvasavarṇena api etat siddham . \\
\noindent \textbf{(P\_6,1.70)} na sidhyati . \\
\noindent \textbf{(P\_6,1.70)} numā vyavahitatvāt pūrvasavarṇaḥ na prāpnoti . \\
\noindent \textbf{(P\_6,1.70)} chandasi napuṃsakasya puṃvadbhāvaḥ vaktavyaḥ madhoḥ gṛhṇāti , mahoḥ tṛptā iva āsate iti evamartham . \\
\noindent \textbf{(P\_6,1.70)} tatra pūṃvadbhāvena numaḥ nivṛttiḥ . \\
\noindent \textbf{(P\_6,1.70)} numi nivṛtte pūrvasavarṇena siddham . \\
\noindent \textbf{(P\_6,1.70)} bhavet siddham agne trī te vajinā trī ṣadhasthā iti . \\
\noindent \textbf{(P\_6,1.70)} idam tu na sidhyati ta tā piṇḍānām iti . \\
\noindent \textbf{(P\_6,1.70)} idam api siddham . \\
\noindent \textbf{(P\_6,1.70)} katham . \\
\noindent \textbf{(P\_6,1.70)} sāptamike pūrvasavarṇe kṛte punaḥ ṣāṣṭhikaḥ bhaviṣyati . \\
\noindent \textbf{(P\_6,1.70)} evam api jasi guṇaḥ prāpnoti . \\
\noindent \textbf{(P\_6,1.70)} vakṣyati etat . \\
\noindent \textbf{(P\_6,1.70)} jasādiṣu chandovāvacanam prāk ṇau caṅi upadhāyāḥ iti . \\
\noindent \textbf{(P\_6,1.71)} tuki pūrvānte napuṃsakopasarjanahrasvatvam dvigusvaraḥ ca . \\
\noindent \textbf{(P\_6,1.71)} tuki pūrvānte napuṃsakopasarjanahrasvatvam dvigusvaraḥ ca na sidhyati . \\
\noindent \textbf{(P\_6,1.71)} ārāśastri chatram , dhānāśaṣkuli chatram . \\
\noindent \textbf{(P\_6,1.71)} niṣkauśāmbi chatram , nirvārāṇasi chatram . \\
\noindent \textbf{(P\_6,1.71)} pañcāratni chatram , daśāratni chatram . \\
\noindent \textbf{(P\_6,1.71)} tuke kṛte anantyatvāt ete vidhayaḥ na prāpnuvanti . \\
\noindent \textbf{(P\_6,1.71)} na vā bahiraṅgalakṣaṇatvāt . \\
\noindent \textbf{(P\_6,1.71)} na vā eṣaḥ doṣaḥ . \\
\noindent \textbf{(P\_6,1.71)} kim kāraṇam . \\
\noindent \textbf{(P\_6,1.71)} bahiraṅgalakṣaṇatvāt . \\
\noindent \textbf{(P\_6,1.71)} bahiraṅgalakṣaṇaḥ tuk . \\
\noindent \textbf{(P\_6,1.71)} antaraṅgāḥ ete vidhayaḥ . \\
\noindent \textbf{(P\_6,1.71)} asiddham bahiraṅgam antaraṅge . \\
\noindent \textbf{(P\_6,1.71)} idam tarhi grāmaṇiputraḥ , senāniputraḥ iti hrasvatve kṛte tuk prāpnoti . \\
\noindent \textbf{(P\_6,1.71)} grāmaṇiputrādiṣu ca aprāptiḥ . \\
\noindent \textbf{(P\_6,1.71)} grāmaṇiputrādiṣu ca aprāptiḥ . \\
\noindent \textbf{(P\_6,1.71)} kim kāraṇam . \\
\noindent \textbf{(P\_6,1.71)} bahiraṅgalakṣaṇatvāt eva . \\
\noindent \textbf{(P\_6,1.71)} atha vā parādiḥ kariṣyate . \\
\noindent \textbf{(P\_6,1.71)} parādau saṃyogādeḥ iti atiprasaṅgaḥ . \\
\noindent \textbf{(P\_6,1.71)} parādau saṃyogādeḥ iti atiprasaṅgaḥ bhavati . \\
\noindent \textbf{(P\_6,1.71)} apacchāyāt . \\
\noindent \textbf{(P\_6,1.71)} vā anyasya saṃyogādeḥ iti etvam prasajyeta . \\
\noindent \textbf{(P\_6,1.71)} vilopavacanam ca . \\
\noindent \textbf{(P\_6,1.71)} veḥ ca lopaḥ vaktavyaḥ . \\
\noindent \textbf{(P\_6,1.71)} agnicit , somasut . \\
\noindent \textbf{(P\_6,1.71)} apṛktasya iti veḥ lopaḥ na prāpnoti . \\
\noindent \textbf{(P\_6,1.71)} na eṣaḥ doṣaḥ . \\
\noindent \textbf{(P\_6,1.71)} apṛktagrahaṇam na kariṣyate . \\
\noindent \textbf{(P\_6,1.71)} yadi na kriyate darviḥ , jāgṛviḥ , atra api prāpnoti . \\
\noindent \textbf{(P\_6,1.71)} anunāsikaparasya viśabdasya grahaṇam śuddhaparaḥ ca atra viśabdaḥ . \\
\noindent \textbf{(P\_6,1.71)} evam api satukkasya lopaḥ prāpnoti . \\
\noindent \textbf{(P\_6,1.71)} nirdiśyamānasya ādeśāḥ bhavanti iti evam na bhaviṣyati . \\
\noindent \textbf{(P\_6,1.71)} iṭpratiṣedhaḥ ca . \\
\noindent \textbf{(P\_6,1.71)} iṭpratiṣedhaḥ ca vaktavyaḥ . \\
\noindent \textbf{(P\_6,1.71)} parītat . \\
\noindent \textbf{(P\_6,1.71)} satukkasya valādilakṣaṇaḥ iṭ prasajyeta . \\
\noindent \textbf{(P\_6,1.71)} evam tarhi abhaktaḥ . \\
\noindent \textbf{(P\_6,1.71)} abhakte svaraḥ . \\
\noindent \textbf{(P\_6,1.71)} yadi abhaktaḥ tarhi svare doṣaḥ bhavati . \\
\noindent \textbf{(P\_6,1.71)} dadhi chādayati , madhu chādayati . \\
\noindent \textbf{(P\_6,1.71)} tiṅ atiṅaḥ iti nighātaḥ na prāpnoti . \\
\noindent \textbf{(P\_6,1.71)} nanu ca tuk eva atiṅ . \\
\noindent \textbf{(P\_6,1.71)} na tukaḥ parasya nighātaḥ prāpnoti . \\
\noindent \textbf{(P\_6,1.71)} kim kāraṇam . \\
\noindent \textbf{(P\_6,1.71)} nañivayuktam anyasadṛśādhikaraṇe . \\
\noindent \textbf{(P\_6,1.71)} tathā hi arthagatiḥ . \\
\noindent \textbf{(P\_6,1.71)} nañyukte ivayukte vā anyasmin tatsadṛśe kāryam vijñāyate . \\
\noindent \textbf{(P\_6,1.71)} tathā hi arthaḥ gamyate . \\
\noindent \textbf{(P\_6,1.71)} tat yathā . \\
\noindent \textbf{(P\_6,1.71)} abrāhmaṇam ānaya iti ukte brāhmaṇasadṛśam eva ānayati . \\
\noindent \textbf{(P\_6,1.71)} na asau loṣṭam ānīya kṛtī bhavati . \\
\noindent \textbf{(P\_6,1.71)} evam iha api atiṅ iti tiṅpratiṣedhāt anyasmāt atiṅaḥ tiṅsadṛśāt kāryam vijñāsyate . \\
\noindent \textbf{(P\_6,1.71)} kim ca anyat atiṅ tiṅsadṛśam . \\
\noindent \textbf{(P\_6,1.71)} padam . \\
\noindent \textbf{(P\_6,1.72)} ayam yogaḥ śakyaḥ avaktum . \\
\noindent \textbf{(P\_6,1.72)} katham . \\
\noindent \textbf{(P\_6,1.72)} adhikaraṇam nāma triprakāram vyāpakam aupaśleṣikam vaiṣayikam iti . \\
\noindent \textbf{(P\_6,1.72)} śabdasya ca śabdena kaḥ anyaḥ abhisambandhaḥ bhavitum arhati anyat ataḥ upaśleṣāt . \\
\noindent \textbf{(P\_6,1.72)} ikaḥ yaṇ aci . \\
\noindent \textbf{(P\_6,1.72)} aci upaśliṣṭasya iti . \\
\noindent \textbf{(P\_6,1.72)} tatra antareṇa saṃhitāgrahaṇam saṃhitāyām eva bhaviṣyati . \\
\noindent \textbf{(P\_6,1.74)} atha kimartham āṅmāṅoḥ sānubandhakayoḥ nirdeśaḥ . \\
\noindent \textbf{(P\_6,1.74)} āṅmāṅoḥ sānubandhakanirdeśaḥ gatikarmapravacanīyapratiṣedhasampratyayārthaḥ . \\
\noindent \textbf{(P\_6,1.74)} āṅmāṅoḥ sānubandhakayoḥ nirdeśaḥ kriyate āṅaḥ gatikarmapravacanīyasampratyayārthaḥ māṅaḥ pratiṣedhasampratyayārthaḥ . \\
\noindent \textbf{(P\_6,1.74)} iha mā bhūt . \\
\noindent \textbf{(P\_6,1.74)} ā chāyā, āc chāyā . \\
\noindent \textbf{(P\_6,1.74)} pramā chandaḥ , pramāc chandaḥ . \\
\noindent \textbf{(P\_6,1.75-76)} KA\_III,51.20-22 Ro\_IV,389 {1/4}     dīrghāt padāntāt vā viśvajanādīnām chandasi . \\
\noindent \textbf{(P\_6,1.75-76)} KA\_III,51.20-22 Ro\_IV,389 {2/4}     dīrghāt padāntāt vā iti atra viśvajanādīnām chandasi upasaṅkhyanam kartavyam . \\
\noindent \textbf{(P\_6,1.75-76)} KA\_III,51.20-22 Ro\_IV,389 {3/4}     viśvajanasya chatram , viśvajanasya cchatram . \\
\noindent \textbf{(P\_6,1.75-76)} KA\_III,51.20-22 Ro\_IV,389 {4/4}     na chāyām kuravaḥ aparām , nac chāyām kuravaḥ aparām . \\
\noindent \textbf{(P\_6,1.77.1)} iggrahaṇam kimartham . \\
\noindent \textbf{(P\_6,1.77.1)} iha mā bhūt . \\
\noindent \textbf{(P\_6,1.77.1)} agnicit atra , somasut atra . \\
\noindent \textbf{(P\_6,1.77.1)} na etat asti prayojanam . \\
\noindent \textbf{(P\_6,1.77.1)} jaśtvam atra bādhakam bhaviṣyati . \\
\noindent \textbf{(P\_6,1.77.1)} jaśtvam na siddham yaṇam atra paśya . asiddham atra jaśtvam . \\
\noindent \textbf{(P\_6,1.77.1)} tasya asiddhatvāt yaṇādeśaḥ prāpnoti . \\
\noindent \textbf{(P\_6,1.77.1)} yaḥ ca apadāntaḥ hal acaḥ ca pūrvaḥ . \\
\noindent \textbf{(P\_6,1.77.1)} yaḥ ca apadāntaḥ hal acaḥ ca pūrvaḥ tasya prāpnoti . \\
\noindent \textbf{(P\_6,1.77.1)} pacati iti . \\
\noindent \textbf{(P\_6,1.77.1)} evam tarhi dīrghasya yaṇ . \\
\noindent \textbf{(P\_6,1.77.1)} dīrghasya yaṇ ādeśam vakṣyāmi . \\
\noindent \textbf{(P\_6,1.77.1)} tat dīrghagrahaṇam kartavyam . \\
\noindent \textbf{(P\_6,1.77.1)} na kartavyam . \\
\noindent \textbf{(P\_6,1.77.1)} prakṛtam anuvartate . \\
\noindent \textbf{(P\_6,1.77.1)} kva prakṛtam . \\
\noindent \textbf{(P\_6,1.77.1)} dīrghāt padāntāt vā iti . \\
\noindent \textbf{(P\_6,1.77.1)} tat vai pañcamīnirdiṣṭam ṣaṣṭhīnirdiṣṭena ca iha arthaḥ . \\
\noindent \textbf{(P\_6,1.77.1)} aci iti eṣā saptamī dīrghāt iti pañcamyāḥ ṣaṣṭhīm prakalpayiṣyati tasmin iti nirdiṣṭe pūrvasya iti . \\
\noindent \textbf{(P\_6,1.77.1)} bhavet siddham kumārī atra , brahmabandhvartham iti . \\
\noindent \textbf{(P\_6,1.77.1)} idam tu na sidhyati . \\
\noindent \textbf{(P\_6,1.77.1)} dadhi atra , madhu atra iti . \\
\noindent \textbf{(P\_6,1.77.1)} hrasvaḥ iti prtavṛttam . \\
\noindent \textbf{(P\_6,1.77.1)} hrasvagrahaṇam api prakṛtam anuvartate . \\
\noindent \textbf{(P\_6,1.77.1)} kva prakṛtam . \\
\noindent \textbf{(P\_6,1.77.1)} hrasvyasya piti kṛti tuk iti . \\
\noindent \textbf{(P\_6,1.77.1)} yadi tat anuvartate dīrghāt padāntāt vā iti hrasvāt api padāntāt vikalpena prāpnoti . \\
\noindent \textbf{(P\_6,1.77.1)} sambandhavṛttyā . \\
\noindent \textbf{(P\_6,1.77.1)} sambandham anuvartiṣyate . \\
\noindent \textbf{(P\_6,1.77.1)} hrasvyasya piti kṛti tuk . \\
\noindent \textbf{(P\_6,1.77.1)} saṃhitāyām hrasvyasya piti kṛti tuk . \\
\noindent \textbf{(P\_6,1.77.1)} che ca hrasvyasya piti kṛti tuk . \\
\noindent \textbf{(P\_6,1.77.1)} āṅmāṅoḥ ca hrasvyasya piti kṛti tuk . \\
\noindent \textbf{(P\_6,1.77.1)} dīrghāt padāntāt vā hrasvyasya piti kṛti tuk . \\
\noindent \textbf{(P\_6,1.77.1)} tataḥ ikaḥ yaṇ aci . \\
\noindent \textbf{(P\_6,1.77.1)} hrasvyasya iti vartate . \\
\noindent \textbf{(P\_6,1.77.1)} piti kṛti tuk iti nivṛttam . \\
\noindent \textbf{(P\_6,1.77.1)} iha tarhi prāpnoti cayanam , cāyakaḥ , lavaṇam , lāvakaḥ . \\
\noindent \textbf{(P\_6,1.77.1)} ayādayaḥ atra bādhakāḥ bhaviṣyanti . \\
\noindent \textbf{(P\_6,1.77.1)} iha tarhi prāpnoti khaṭvā indraḥ , mālā indraḥ , khaṭvā elakā , mālā elakā . \\
\noindent \textbf{(P\_6,1.77.1)} guṇavṛddhibādhyaḥ . \\
\noindent \textbf{(P\_6,1.77.1)} guṇavṛddhī atra bādhike bhaviṣyataḥ . \\
\noindent \textbf{(P\_6,1.77.1)} idam tarhi prayojanam . \\
\noindent \textbf{(P\_6,1.77.1)} ikaḥ aci yaṇ eva syāt . \\
\noindent \textbf{(P\_6,1.77.1)} yat anyat prāpnoti tat mā bhūt iti . \\
\noindent \textbf{(P\_6,1.77.1)} kim ca anyat prāpnoti . \\
\noindent \textbf{(P\_6,1.77.1)} śākalam . \\
\noindent \textbf{(P\_6,1.77.1)} sinnityasamāsayoḥ śākalapratiṣedham codayiṣyati . \\
\noindent \textbf{(P\_6,1.77.1)} sa na vaktavyaḥ bhavati . \\
\noindent \textbf{(P\_6,1.77.2)} yaṇādeśaḥ plutapūrvasya ca . \\
\noindent \textbf{(P\_6,1.77.2)} yaṇādeśaḥ plutapūrvasya ca iti vaktavyam . \\
\noindent \textbf{(P\_6,1.77.2)} agnā3i* indram , agnā3y indram , paṭā3u* udakam , paṭā3v udakam , agnā3i* āśā , agnā3y āśā , paṭā3u* āśā , paṭā3v āśā . \\
\noindent \textbf{(P\_6,1.77.2)} kim punaḥ kāraṇam na sidhyati . \\
\noindent \textbf{(P\_6,1.77.2)} asiddhaḥ plutaḥ plutavikārau ca imau . \\
\noindent \textbf{(P\_6,1.77.2)} siddhaḥ plutaḥ svarasandhiṣu . \\
\noindent \textbf{(P\_6,1.77.2)} katham jñāyate . \\
\noindent \textbf{(P\_6,1.77.2)} yat ayam plutapragṛhyāḥ aci iti plutasya prakṛtibhāvam śāsti tat jñāpayati ācāryaḥ siddhaḥ plutaḥ svarasandhiṣu iti . \\
\noindent \textbf{(P\_6,1.77.2)} katham kṛtvā jñāpakam . \\
\noindent \textbf{(P\_6,1.77.2)} sataḥ hi kāryiṇaḥ kāryeṇa bhavitavyam . \\
\noindent \textbf{(P\_6,1.77.2)} idam tarhi prayojanam dīrghaśākalapratiṣedhārtham . \\
\noindent \textbf{(P\_6,1.77.2)} dīrghatvam śākalam ca mā bhūt iti . \\
\noindent \textbf{(P\_6,1.77.2)} etat api na asti prayojanam . \\
\noindent \textbf{(P\_6,1.77.2)} ārabhyate plutapūrvasya yaṇādeśaḥ tayoḥ yvau aci saṃhitāyām iti . \\
\noindent \textbf{(P\_6,1.77.2)} tat dīrghaśākalapratiṣedhārtham bhaviṣyati . \\
\noindent \textbf{(P\_6,1.77.2)} tat na vaktavyam bhavati . \\
\noindent \textbf{(P\_6,1.77.2)} nanu ca tasmin api ucyamāne idam na vaktavyam bhavati . \\
\noindent \textbf{(P\_6,1.77.2)} avaśyam idam vaktavyam yau plutapūrvau idutau aplutavikārau tadartham . \\
\noindent \textbf{(P\_6,1.77.2)} bho3i indram , bho3y indram , bho3i iha bho3y iha iti . \\
\noindent \textbf{(P\_6,1.77.2)} yad tarhi asya nibandhanam asti idam eva vaktavyam . \\
\noindent \textbf{(P\_6,1.77.2)} tat na vaktavyam . \\
\noindent \textbf{(P\_6,1.77.2)} tat api avaśyam svarārtham vaktavyam . \\
\noindent \textbf{(P\_6,1.77.2)} anena hi sati udāttasvaritayoḥ yaṇaḥ iti eṣaḥ svaraḥ prasajyeta . \\
\noindent \textbf{(P\_6,1.77.2)} tena punaḥ sati asiddhatvāt na bhaviṣyati . \\
\noindent \textbf{(P\_6,1.77.2)} yadi tarhi tasya nibandhanam asti tat eva vaktavyam . \\
\noindent \textbf{(P\_6,1.77.2)} idam na vaktavyam . \\
\noindent \textbf{(P\_6,1.77.2)} nanu ca uktam idam api avaśyam vaktavyam yau plutapūrvau idutau aplutavikārau tadartham . \\
\noindent \textbf{(P\_6,1.77.2)} bho3i indram , bho3y indram , bho3i iha bho3y iha iti . \\
\noindent \textbf{(P\_6,1.77.2)} chāndasam etat dṛṣṭānuvidhiḥ chandasi bhavati . \\
\noindent \textbf{(P\_6,1.77.2)} yat tarhi na chāndasam bho3y indram , bho3y iha iti sāma gāyati . \\
\noindent \textbf{(P\_6,1.77.2)} eṣaḥ api chandasi dṛṣṭasya anuprayogaḥ kriyate . \\
\noindent \textbf{(P\_6,1.77.2)} jaśtvam na siddham yaṇam atra paśya . \\
\noindent \textbf{(P\_6,1.77.2)} yaḥ ca apadāntaḥ hal acaḥ ca pūrvaḥ . \\
\noindent \textbf{(P\_6,1.77.2)} dīrghasya yaṇ . \\
\noindent \textbf{(P\_6,1.77.2)} hrasvaḥ iti prtavṛttam . \\
\noindent \textbf{(P\_6,1.77.2)} sambandhavṛttyā . \\
\noindent \textbf{(P\_6,1.77.2)} guṇavṛddhibādhyaḥ . \\
\noindent \textbf{(P\_6,1.77.2)} nitye ca yaḥ śākalabhāksamāse tadartham etad bhagavān cakāra . \\
\noindent \textbf{(P\_6,1.77.2)} sāmarthyayogāt na hi kim cit asmin paśyāmi śāstre yat anarthakam syāt . \\
\noindent \textbf{(P\_6,1.79.1)} vāntādeśe sthāninirdeśaḥ . \\
\noindent \textbf{(P\_6,1.79.1)} vāntādeśe sthāninirdeśaḥ kartavyaḥ . \\
\noindent \textbf{(P\_6,1.79.1)} okāraukārayoḥ iti vaktavyam ekāraikārayoḥ mā bhūt iti . \\
\noindent \textbf{(P\_6,1.79.1)} saḥ tarhi kartavyaḥ . \\
\noindent \textbf{(P\_6,1.79.1)} na kartavyaḥ . \\
\noindent \textbf{(P\_6,1.79.1)} vāntagrahaṇam na kariṣyate . \\
\noindent \textbf{(P\_6,1.79.1)} ecaḥ yi pratyaye ayādayaḥ bhavanti iti eva siddham . \\
\noindent \textbf{(P\_6,1.79.1)} yadi vāntagrahaṇam na kriyate ceyam , jeyam iti atra api prāpnoti . \\
\noindent \textbf{(P\_6,1.79.1)} kṣayyajayyau śakyārthe iti etat niyamārtham bhaviṣyati . \\
\noindent \textbf{(P\_6,1.79.1)} kṣijyoḥ eva iti . \\
\noindent \textbf{(P\_6,1.79.1)} tayoḥ tarhi śakyārthāt anyatra api prāpnoti . \\
\noindent \textbf{(P\_6,1.79.1)} kṣeyam pāpam , jeyaḥ vṛṣalaḥ iti . \\
\noindent \textbf{(P\_6,1.79.1)} ubhayataḥ niyamaḥ vijñāsyate . \\
\noindent \textbf{(P\_6,1.79.1)} kṣijyoḥ eva ecaḥ tayoḥ ca śakyārthe eva iti . \\
\noindent \textbf{(P\_6,1.79.1)} iha api tarhi niyamāt na prāpnoti . \\
\noindent \textbf{(P\_6,1.79.1)} lavyam , pavyam , avaśyalāvyam , avaśyapāvyam . \\
\noindent \textbf{(P\_6,1.79.1)} tulyajātīyasya niyamaḥ . \\
\noindent \textbf{(P\_6,1.79.1)} kaḥ ca tulyajātīyaḥ . \\
\noindent \textbf{(P\_6,1.79.1)} yathājātīyakaḥ kṣijyoḥ ec . \\
\noindent \textbf{(P\_6,1.79.1)} kathañjātīyakaḥ kṣijyoḥ ec . \\
\noindent \textbf{(P\_6,1.79.1)} ekāraḥ . \\
\noindent \textbf{(P\_6,1.79.1)} evam api rāyam icchati , raiyati , atra api prāpnoti . \\
\noindent \textbf{(P\_6,1.79.1)} rāyiḥ chāndasaḥ . \\
\noindent \textbf{(P\_6,1.79.1)} dṛṣṭānuvidhiḥ chandasi bhavati . \\
\noindent \textbf{(P\_6,1.79.2)} goḥ yūtau chandasi . \\
\noindent \textbf{(P\_6,1.79.2)} goḥ yūtau chandasi upasaṅkhyānam kartavyam . \\
\noindent \textbf{(P\_6,1.79.2)} a naḥ mitrāvaruṇā ghṛtaiḥ gavyūtim ukṣatam . \\
\noindent \textbf{(P\_6,1.79.2)} goyūtim iti eva anyatra . \\
\noindent \textbf{(P\_6,1.79.2)} adhvaparimāṇe ca . \\
\noindent \textbf{(P\_6,1.79.2)} adhvaparimāṇe ca goḥ yūtau upasaṅkhyānam kartavyam . \\
\noindent \textbf{(P\_6,1.79.2)} gavyūtim adhvānam gataḥ . \\
\noindent \textbf{(P\_6,1.79.2)} goyūtim iti eva anyatra . \\
\noindent \textbf{(P\_6,1.80)} evakāraḥ kimarthaḥ . \\
\noindent \textbf{(P\_6,1.80)} niyamārthaḥ . \\
\noindent \textbf{(P\_6,1.80)} na etat prayojanam . \\
\noindent \textbf{(P\_6,1.80)} na etat asti prayojanam . \\
\noindent \textbf{(P\_6,1.80)} siddhe vidhiḥ ārabhyamāṇaḥ antareṇa evakāram niyamārthaṅ bhaviṣyati . \\
\noindent \textbf{(P\_6,1.80)} iṣṭataḥ avadhāraṇārthaḥ tarhi . \\
\noindent \textbf{(P\_6,1.80)} yathā evam vijñāyeta . \\
\noindent \textbf{(P\_6,1.80)} dhātoḥ tannimittasya eva iti . \\
\noindent \textbf{(P\_6,1.80)} mā evam vijñāyi . \\
\noindent \textbf{(P\_6,1.80)} dhātoḥ eva tannimittasya iti . \\
\noindent \textbf{(P\_6,1.80)} kim ca syāt . \\
\noindent \textbf{(P\_6,1.80)} adhātoḥ tannimittasya na syāt . \\
\noindent \textbf{(P\_6,1.80)} śaṅkavyam dāru , picavyaḥ kārpāsaḥ iti . \\
\noindent \textbf{(P\_6,1.82)} tat iti anena kim pratinirdiśyate . \\
\noindent \textbf{(P\_6,1.82)} saḥ eva krīṇātyarthaḥ . \\
\noindent \textbf{(P\_6,1.82)} iha mā bhūt . \\
\noindent \textbf{(P\_6,1.82)} kreyam naḥ dhānyam . \\
\noindent \textbf{(P\_6,1.82)} na ca asti krayyam iti . \\
\noindent \textbf{(P\_6,1.83)} bhayyādiprakaraṇe hradayyāḥ upasaṅkhyānam . \\
\noindent \textbf{(P\_6,1.83)} bhayyādiprakaraṇe hradayyāḥ upasaṅkhyānam kartavyam . \\
\noindent \textbf{(P\_6,1.83)} hradayyāḥ āpaḥ . \\
\noindent \textbf{(P\_6,1.83)} av śarasya ca . \\
\noindent \textbf{(P\_6,1.83)} śarasya ca hradasya ca ataḥ av vaktayaḥ . \\
\noindent \textbf{(P\_6,1.83)} hradavyāḥ āpaḥ . \\
\noindent \textbf{(P\_6,1.83)} śaravyāḥ vai tejanam . \\
\noindent \textbf{(P\_6,1.83)} śaravyasya paśūn abhighātakaḥ syāt . \\
\noindent \textbf{(P\_6,1.83)} śaruvṛttāt vā siddham . \\
\noindent \textbf{(P\_6,1.83)} śaruvṛttāt vā siddham punaḥ siddham etat . \\
\noindent \textbf{(P\_6,1.83)} ṛñjatī śaruḥ iti api dṛśyate . \\
\noindent \textbf{(P\_6,1.83)} ṛñjatī śaruḥ iti api śaruśabdapravṛttiḥ dṛśyate . \\
\noindent \textbf{(P\_6,1.83)} śaruhastaḥ iti ca loke . \\
\noindent \textbf{(P\_6,1.83)} śaruhastaḥ iti ca loke śarahastam upācaranti . \\
\noindent \textbf{(P\_6,1.84.1)} ekavacanam kimartham . \\
\noindent \textbf{(P\_6,1.84.1)} ekavacanam pṛthak ādeśapratiṣedhārtham . \\
\noindent \textbf{(P\_6,1.84.1)} ekavacanam kriyate ekaḥ ādeśaḥ yathā syāt . \\
\noindent \textbf{(P\_6,1.84.1)} pṛthak ādeśaḥ mā bhūt iti . \\
\noindent \textbf{(P\_6,1.84.1)} na vā dravyavat karmacodanāyām dvayoḥ ekasya abhinirvṛtteḥ . \\
\noindent \textbf{(P\_6,1.84.1)} na vā etat prayojanam asti . \\
\noindent \textbf{(P\_6,1.84.1)} kim kāraṇam . \\
\noindent \textbf{(P\_6,1.84.1)} dravyavat karmacodanāyām dvayoḥ ekasya abhinirvṛtteḥ ekaḥ ādeśaḥ bhaviṣyati . \\
\noindent \textbf{(P\_6,1.84.1)} tat yathā dravyeṣu karmacodanāyām dvayoḥ ekasya abhinirvṛttiḥ bhavati . \\
\noindent \textbf{(P\_6,1.84.1)} anayoḥ pūlayoḥ kaṭam kuru . \\
\noindent \textbf{(P\_6,1.84.1)} anayoḥ mitpiṇḍayoḥ ghaṭam kuru iti . \\
\noindent \textbf{(P\_6,1.84.1)} na ca ucyate ekam iti ekam ca asau karoti . \\
\noindent \textbf{(P\_6,1.84.1)} kim punaḥ kāraṇam dravyeṣu karmacodanāyām dvayoḥ ekasya abhinirvṛttiḥ bhavati . \\
\noindent \textbf{(P\_6,1.84.1)} tat ca ekavākyabhāvāt . \\
\noindent \textbf{(P\_6,1.84.1)} ekavākyabhāvāt dravyeṣu karmacodanāyām dvayoḥ ekasya abhinirvṛttiḥ bhavati . \\
\noindent \textbf{(P\_6,1.84.1)} ātaḥ ca ekavākyabhāvāt . \\
\noindent \textbf{(P\_6,1.84.1)} vyākaraṇe api hi anyatra dvayoḥ sthāninoḥ ekaḥ ādeśaḥ bhavati . \\
\noindent \textbf{(P\_6,1.84.1)} jvaratvarasrivyavimavām upadhāyāḥ ca . \\
\noindent \textbf{(P\_6,1.84.1)} bhrasjaḥ ropadhayoḥ ram anyatarasyām iti . \\
\noindent \textbf{(P\_6,1.84.1)} yat tāvat ucyate ekavākyabhāvāt iti tat na . \\
\noindent \textbf{(P\_6,1.84.1)} arthāt prakaraṇāt vā loke dvayoḥ ekasya abhinirvṛttiḥ bhavati . \\
\noindent \textbf{(P\_6,1.84.1)} ātaḥ ca arthāt prakaraṇāt vā . \\
\noindent \textbf{(P\_6,1.84.1)} vyākaraṇe api hi anyatra dvayoḥ sthāninoḥ dvau ādeśau bhavataḥ . \\
\noindent \textbf{(P\_6,1.84.1)} radābhyām niṣṭhātaḥ naḥ pūrvasya ca daḥ . \\
\noindent \textbf{(P\_6,1.84.1)} ubhau sābhyāsasya iti . \\
\noindent \textbf{(P\_6,1.84.1)} katham yat tat uktam vyākaraṇe api hi anyatra dvayoḥ sthāninoḥ ekaḥ ādeśaḥ bhavati . \\
\noindent \textbf{(P\_6,1.84.1)} jvaratvarasrivyavimavām upadhāyāḥ ca . \\
\noindent \textbf{(P\_6,1.84.1)} bhrasjaḥ ropadhayoḥ ram anyatarasyām iti . \\
\noindent \textbf{(P\_6,1.84.1)} iha tāvat jvaratvarasrivyavimavām upadhāyāḥ ca iti . \\
\noindent \textbf{(P\_6,1.84.1)} stām dvau ūṭhau . \\
\noindent \textbf{(P\_6,1.84.1)} na asti doṣaḥ . \\
\noindent \textbf{(P\_6,1.84.1)} savarṇadīrghatvena siddham . \\
\noindent \textbf{(P\_6,1.84.1)} iha bhrasjaḥ ropadhayoḥ ram anyatarasyām iti . \\
\noindent \textbf{(P\_6,1.84.1)} vakṣyati hi etat . \\
\noindent \textbf{(P\_6,1.84.1)} bhrasjaḥ ropadhayoḥ lopaḥ āgamaḥ ram vidhīyate iti . \\
\noindent \textbf{(P\_6,1.84.1)} yat ucyate arthāt prakaraṇāt vā iti tat na . \\
\noindent \textbf{(P\_6,1.84.1)} kim kāraṇam . \\
\noindent \textbf{(P\_6,1.84.1)} ekavākyabhāvāt eva loke dvayoḥ ekasya abhinirvṛttiḥ bhavati . \\
\noindent \textbf{(P\_6,1.84.1)} ātaḥ ca ekavākyabhāvāt . \\
\noindent \textbf{(P\_6,1.84.1)} aṅga hi bhavān grāmyam pāṃsulapādam aprakaraṇajñam āgatam bravītu anayoḥ pūlayoḥ kaṭam kuru . \\
\noindent \textbf{(P\_6,1.84.1)} anayoḥ mitpiṇḍayoḥ ghaṭam kuru iti . \\
\noindent \textbf{(P\_6,1.84.1)} ekam eva asau kariṣyati . \\
\noindent \textbf{(P\_6,1.84.1)} katham yat uktam vyākaraṇe api hi anyatra dvayoḥ sthāninoḥ dvau ādeśau bhavataḥ . \\
\noindent \textbf{(P\_6,1.84.1)} radābhyām niṣṭhātaḥ naḥ pūrvasya ca daḥ . \\
\noindent \textbf{(P\_6,1.84.1)} ubhau sābhyāsasya iti . \\
\noindent \textbf{(P\_6,1.84.1)} iha tāvat radābhyām niṣṭhātaḥ naḥ pūrvasya ca daḥ iti dve vākye . \\
\noindent \textbf{(P\_6,1.84.1)} katham . \\
\noindent \textbf{(P\_6,1.84.1)} yogavibhāgaḥ kariṣyate . \\
\noindent \textbf{(P\_6,1.84.1)} radābhyām niṣṭhātaḥ naḥ . \\
\noindent \textbf{(P\_6,1.84.1)} tataḥ pūrvasya ca daḥ iti . \\
\noindent \textbf{(P\_6,1.84.1)} iha ubhau sābhyāsasya iti ubhaugrahaṇasāmarthyāt dvau ādeśau bhaviṣyataḥ . \\
\noindent \textbf{(P\_6,1.84.2)} tatra avayave śāstrārthasampratyayaḥ yathā loke ṭatra avayave śāstrārthasampratyayaḥ prāpnoti yathā loke . \\
\noindent \textbf{(P\_6,1.84.2)} tat yathā loke . \\
\noindent \textbf{(P\_6,1.84.2)} vasante brāhmaṇaḥ agnīn ādadhīta iti sakṛt ādhāya kṛtaḥ śāstrārthaḥ iti kṛtvā punaḥ pravṛttiḥ na bhavati . \\
\noindent \textbf{(P\_6,1.84.2)} tathā garbhāṣṭame brāhmaṇaḥ upaneyaḥ iti sakṛt upanīya kṛtaḥ śāstrārthaḥ iti kṛtvā punaḥ pravṛttiḥ na bhavati . \\
\noindent \textbf{(P\_6,1.84.2)} tathā triḥ hṛdayaṅgamābhiḥ adbhiḥ aśabdābhiḥ upaspṛśet iti sakṛt upaspṛśya kṛtaḥ śāstrārthaḥ iti kṛtvā punaḥ pravṛttiḥ na bhavati . \\
\noindent \textbf{(P\_6,1.84.2)} evam iha api khaṭvendre kṛtaḥ śāstrārthaḥ iti kṛtvā mālendrādiṣu na syāt . \\
\noindent \textbf{(P\_6,1.84.2)} siddham tu dharmopadeśane anavayavavijñānāt yathā laukikavaidikeṣu . \\
\noindent \textbf{(P\_6,1.84.2)} siddham etat . \\
\noindent \textbf{(P\_6,1.84.2)} katham . \\
\noindent \textbf{(P\_6,1.84.2)} dharmopadeśanam idam śāstram . \\
\noindent \textbf{(P\_6,1.84.2)} dharmopadeśane ca asmin śāstre anavayavena śāstrārthaḥ sampratīyate yathā laukikeṣu vaidikeṣu ca kṛtānteṣu . \\
\noindent \textbf{(P\_6,1.84.2)} loke tāvat : brāhmaṇaḥ na hantavyaḥ . \\
\noindent \textbf{(P\_6,1.84.2)} surā na peyā iti . \\
\noindent \textbf{(P\_6,1.84.2)} brāhmaṇamātram na hanyate surāmātram ca na pīyate . \\
\noindent \textbf{(P\_6,1.84.2)} yadi ca avayavena śāstrārthasampratyayaḥ syāt ekam ca brāhmaṇam ahatvā ekām ca surām apītvā anyatra kāmacāraḥ syāt . \\
\noindent \textbf{(P\_6,1.84.2)} tathā pūrvavayāḥ brāhmaṇaḥ pratyuttheyaḥ iti pūrvavayomātram pratyutthīyate . \\
\noindent \textbf{(P\_6,1.84.2)} yadi avayavena śāstrārthasampratyayaḥ syāt ekam pūrvavayasam pratyutthāya anyatra kāmacāraḥ syāt . \\
\noindent \textbf{(P\_6,1.84.2)} tathā vede khalu api . \\
\noindent \textbf{(P\_6,1.84.2)} vasante brāhmaṇaḥ agniṣṭomādibhiḥ kratubhiḥ yajeta iti agnyādhānanimittam vasante vasante ijyate . \\
\noindent \textbf{(P\_6,1.84.2)} yadi avayavena śāstrārthasampratyayaḥ syāt sakṛt iṣṭvā punaḥ ijyā na pravarteta . \\
\noindent \textbf{(P\_6,1.84.2)} ubhayathā iha loke dṛśyate . \\
\noindent \textbf{(P\_6,1.84.2)} avayavena api śāstrārthasampratyayaḥ anavayavena api . \\
\noindent \textbf{(P\_6,1.84.2)} katham punaḥ idam ubhayam labhyam . \\
\noindent \textbf{(P\_6,1.84.2)} labhyam iti āha . \\
\noindent \textbf{(P\_6,1.84.2)} katham . \\
\noindent \textbf{(P\_6,1.84.2)} iha tāvat vasante brāhmaṇaḥ agnīn ādadhīta iti . \\
\noindent \textbf{(P\_6,1.84.2)} agnyādhānam yajñamukhaprtipattyartham . \\
\noindent \textbf{(P\_6,1.84.2)} sakṛt ādhāya kṛtaḥ śāstrārthaḥ pratipannam yajñam iti kṛtvā punaḥ pravṛttiḥ na bhavati . \\
\noindent \textbf{(P\_6,1.84.2)} ataḥ atra avayavena śāstrārthaḥ sampratīyate . \\
\noindent \textbf{(P\_6,1.84.2)} tathā garbhāṣṭame brāhmaṇaḥ upaneyaḥ iti . \\
\noindent \textbf{(P\_6,1.84.2)} upanayanam saṃskārārtham . \\
\noindent \textbf{(P\_6,1.84.2)} sakṛt ca asau upanītaḥ saṃskṛtaḥ bhavati . \\
\noindent \textbf{(P\_6,1.84.2)} ataḥ atra api avayavena śāstrārthaḥ sampratīyate . \\
\noindent \textbf{(P\_6,1.84.2)} tathā triḥ hṛdayaṅgamābhiḥ adbhiḥ aśabdābhiḥ upaspṛśet iti . \\
\noindent \textbf{(P\_6,1.84.2)} upasparśanam śaucārtham . \\
\noindent \textbf{(P\_6,1.84.2)} sakṛt ca asau upaspṛśya śuciḥ bhavati . \\
\noindent \textbf{(P\_6,1.84.2)} ataḥ atra api avayavena śāstrārthaḥ sampratīyate . \\
\noindent \textbf{(P\_6,1.84.2)} iha idānīm brāhmaṇaḥ na hantavyaḥ . \\
\noindent \textbf{(P\_6,1.84.2)} surā na peyā iti . \\
\noindent \textbf{(P\_6,1.84.2)} brāhmaṇavadhe surāpāne ca mahān doṣaḥ uktaḥ . \\
\noindent \textbf{(P\_6,1.84.2)} saḥ brāhmaṇavadhamātre surāpānamātre ca prasaktaḥ . \\
\noindent \textbf{(P\_6,1.84.2)} ataḥ atra anavayavena śāstrārthaḥ sampratīyate . \\
\noindent \textbf{(P\_6,1.84.2)} tathā pūrvavayāḥ brāhmaṇaḥ pratyuttheyaḥ iti . \\
\noindent \textbf{(P\_6,1.84.2)} pūrvavayasaḥ apratyutthāne doṣaḥ uktaḥ pratyutthāne ca guṇaḥ . \\
\noindent \textbf{(P\_6,1.84.2)} katham . \\
\noindent \textbf{(P\_6,1.84.2)} ūrdhvam prāṇāḥ hi utkrāmanti yūnaḥ sthavire āyati . \\
\noindent \textbf{(P\_6,1.84.2)} pratyutthānābhābhivādābhyām punaḥ tān pratipadyate iti . \\
\noindent \textbf{(P\_6,1.84.2)} saḥ ca prūrvavayomātre prasaktaḥ . \\
\noindent \textbf{(P\_6,1.84.2)} ataḥ atra api anavayavena śāstrārthaḥ sampratīyate . \\
\noindent \textbf{(P\_6,1.84.2)} tathā vasante brāhmaṇaḥ agniṣṭomādibhiḥ kratubhiḥ yajeta iti ijyāyāḥ kim cit prayojanam uktam . \\
\noindent \textbf{(P\_6,1.84.2)} kim . \\
\noindent \textbf{(P\_6,1.84.2)} svarge loke apsarasaḥ enam jāyāḥ bhūtvā upaśerate iti . \\
\noindent \textbf{(P\_6,1.84.2)} tat ca dvitīyasyāḥ tṛtīyasyāḥ ca ijyāyāḥ bhavitum arhati . \\
\noindent \textbf{(P\_6,1.84.2)} ataḥ atra api anavayavena śāstrārthaḥ sampratīyate . \\
\noindent \textbf{(P\_6,1.84.2)} tathā śabdasya api jñāne prayoge prayojanam uktam . \\
\noindent \textbf{(P\_6,1.84.2)} kim . \\
\noindent \textbf{(P\_6,1.84.2)} ekaḥ śabdaḥ samyak jñātaḥ śāstrānvitaḥ suprayuktaḥ svarge loke kāmadhuk bhavati iti . \\
\noindent \textbf{(P\_6,1.84.2)} yadi ekaḥ śabdaḥ samyak jñātaḥ śāstrānvitaḥ suprayuktaḥ svarge loke kāmadhuk bhavati kimartham dvitīyaḥ tṛtīyaḥ ca prayujyate . \\
\noindent \textbf{(P\_6,1.84.2)} na vai kāmānām tṛptiḥ asti . \\
\noindent \textbf{(P\_6,1.84.3)} atha pūrvagrahaṇam kimartham . \\
\noindent \textbf{(P\_6,1.84.3)} pūrvaparagrahaṇam parasya ādeśapratiṣedhārtham . \\
\noindent \textbf{(P\_6,1.84.3)} pūrvaparagrahaṇam kriyate parasya ādeśapratiṣedhārtham . \\
\noindent \textbf{(P\_6,1.84.3)} parasya ādeśaḥ mā bhūt . \\
\noindent \textbf{(P\_6,1.84.3)} āt guṇaḥ iti . \\
\noindent \textbf{(P\_6,1.84.3)} katham ca prāpnoti . \\
\noindent \textbf{(P\_6,1.84.3)} pañcamīnirdiṣṭāt hi parasya . \\
\noindent \textbf{(P\_6,1.84.3)} pañcamīnirdiṣṭāt hi parasya kāryam ucyate . \\
\noindent \textbf{(P\_6,1.84.3)} tat yathā dvyantarupasargebhyaḥ apaḥ īt iti . \\
\noindent \textbf{(P\_6,1.84.3)} ṣaṣṭhīnirdiṣṭārtham tu . \\
\noindent \textbf{(P\_6,1.84.3)} ṣaṣṭhīnirdiṣṭārtham ca pūrvaparagrahaṇam kriyate . \\
\noindent \textbf{(P\_6,1.84.3)} ṣaṣṭhīnirdeśaḥ yathā pakalpeta . \\
\noindent \textbf{(P\_6,1.84.3)} anirdiṣṭe hi ṣaṣṭhyarthāprasiddhiḥ . \\
\noindent \textbf{(P\_6,1.84.3)} akriyamāṇe hi pūrvaparagrahaṇe ṣaṣṭhyarthasya aprasiddhiḥ syāt . \\
\noindent \textbf{(P\_6,1.84.3)} kasya . \\
\noindent \textbf{(P\_6,1.84.3)} sthāneyogatvasya . \\
\noindent \textbf{(P\_6,1.84.3)} na eṣaḥ doṣaḥ . \\
\noindent \textbf{(P\_6,1.84.3)} āt iti eṣā pañcamī aci iti saptamyāḥ ṣaṣṭhīm prakalpayiṣyati tasmāt iti uttarasya iti . \\
\noindent \textbf{(P\_6,1.84.3)} tathā ca aci iti eṣā saptamī āt iti pañcamyāḥ ṣaṣṭhīm prakalpayiṣyati tasmin iti nirdiṣṭe pūrvasya iti . \\
\noindent \textbf{(P\_6,1.84.3)} evam tarhi siddhe sati yat pūrvagrahaṇam karoti tat jñāpayati ācāryaḥ na ubhe yugapat prakalpike bhavataḥ iti . \\
\noindent \textbf{(P\_6,1.84.3)} kim etasya jñāpane prayojanam . \\
\noindent \textbf{(P\_6,1.84.3)} yat uktam : saptamīpañcamyoḥ ca bhāvāt ubhayatra ṣaṣṭhīprakḷptiḥ tatra ubhayakāryaprasaṅgaḥ iti . \\
\noindent \textbf{(P\_6,1.84.3)} saḥ na doṣaḥ bhavati . \\
\noindent \textbf{(P\_6,1.85.1)} kimartham idam ucyate . \\
\noindent \textbf{(P\_6,1.85.1)} antādivadvacanam āmiśrasya ādeśavacanāt . \\
\noindent \textbf{(P\_6,1.85.1)} antādivat iti ucyate āmiśrasya ādeśavacanāt . \\
\noindent \textbf{(P\_6,1.85.1)} āmiśrasya ayam ādeśaḥ ucyate . \\
\noindent \textbf{(P\_6,1.85.1)} saḥ na eva pūrvagrahaṇena gṛhyate na api paragrahaṇena . \\
\noindent \textbf{(P\_6,1.85.1)} tat yathā kṣīrodake sampṛkte āmiśratvāt na eva kṣīragrahaṇena gṛhyate na api udakagrahaṇena . \\
\noindent \textbf{(P\_6,1.85.1)} iṣyate ca grahaṇam syāt iti . \\
\noindent \textbf{(P\_6,1.85.1)} tat ca antareṇa yatnam na sidhyati iti antādivacvacanam . \\
\noindent \textbf{(P\_6,1.85.1)} evamartham idam ucyate . \\
\noindent \textbf{(P\_6,1.85.1)} asti prayojanam etat . \\
\noindent \textbf{(P\_6,1.85.1)} kim tarhi iti . \\
\noindent \textbf{(P\_6,1.85.1)} tatra yasya antādivat tannirdeśaḥ . \\
\noindent \textbf{(P\_6,1.85.1)} tatra yasya antādivadbhāvaḥ iṣyate tannirdeśaḥ kartavyaḥ . \\
\noindent \textbf{(P\_6,1.85.1)} asya antavat bhavati asya ādivat bhavati iti vaktavyam . \\
\noindent \textbf{(P\_6,1.85.1)} siddham tu pūrvaparādhikārāt . \\
\noindent \textbf{(P\_6,1.85.1)} siddham etat . \\
\noindent \textbf{(P\_6,1.85.1)} katham . \\
\noindent \textbf{(P\_6,1.85.1)} pūrvaparādhikārāt . \\
\noindent \textbf{(P\_6,1.85.1)} pūrvaparayoḥ iti vartate . \\
\noindent \textbf{(P\_6,1.85.1)} pūrvasya kāryam prati antavat bhavati . \\
\noindent \textbf{(P\_6,1.85.1)} parasya kāryam prati ādivat bhavati . \\
\noindent \textbf{(P\_6,1.85.1)} atha yatra ubhayam āśrīyate kim tatra pūrvasya antavat bhavati āhosvit parasya ādivat bhavati . \\
\noindent \textbf{(P\_6,1.85.1)} ubhayataḥ āśraye na antādivat . \\
\noindent \textbf{(P\_6,1.85.1)} kim vaktavyam etat . \\
\noindent \textbf{(P\_6,1.85.1)} na hi . \\
\noindent \textbf{(P\_6,1.85.1)} katham anucyamānam gaṃsyate . \\
\noindent \textbf{(P\_6,1.85.1)} laukikaḥ ayam dṛṣṭāntaḥ . \\
\noindent \textbf{(P\_6,1.85.1)} tat yathā loke yaḥ dvayoḥ tulyabalayoḥ preṣyaḥ bhavati saḥ tayoḥ paryāyeṇa kāryam karoti . \\
\noindent \textbf{(P\_6,1.85.1)} yadā tu tam ubhau yugapat preṣayataḥ nānādikṣu ca kārye bhavataḥ tatra yadi asau avirodhāṛthī bhavati tataḥ ubhayoḥ na karoti . \\
\noindent \textbf{(P\_6,1.85.1)} kim punaḥ kāraṇam ubhayoḥ na karoti . \\
\noindent \textbf{(P\_6,1.85.1)} yaugapadyāsambhavāt . \\
\noindent \textbf{(P\_6,1.85.1)} na asti yaugapadyena sambhavaḥ . \\
\noindent \textbf{(P\_6,1.85.2)} atha antavattve kāni prayojanāni . \\
\noindent \textbf{(P\_6,1.85.2)} antavattve prayojanam bahvacpūrvapadāt ṭhajvidhāne . \\
\noindent \textbf{(P\_6,1.85.2)} antavattve bahvacpūrvapadāt ṭhajvidhāne prayojanam . \\
\noindent \textbf{(P\_6,1.85.2)} dvādaśānyikaḥ . \\
\noindent \textbf{(P\_6,1.85.2)} pūrvapadottarapadayoḥ ekādeśaḥ pūrvapadasaya antavat bhavati yathā śakyeta kartum bahucpūrvapadāt ṭhac bhavati iti . \\
\noindent \textbf{(P\_6,1.85.2)} kva tarhi syāt . \\
\noindent \textbf{(P\_6,1.85.2)} yatra kṛte api ekādeśe bahvacpūrvapadam bhavati . \\
\noindent \textbf{(P\_6,1.85.2)} tarayodaśānyikaḥ . \\
\noindent \textbf{(P\_6,1.85.2)} pratyayaikādeśaḥ pūrvavidhau . \\
\noindent \textbf{(P\_6,1.85.2)} pratyayaikādeśaḥ pūrvavidhau prayojanam . \\
\noindent \textbf{(P\_6,1.85.2)} madhu pibanti . \\
\noindent \textbf{(P\_6,1.85.2)} śidaśitoḥ ekādeśaḥ śitaḥ antavat bhavati yathā śakyeta kartum śiti iti pibādeśaḥ . \\
\noindent \textbf{(P\_6,1.85.2)} kva tarhi syāt . \\
\noindent \textbf{(P\_6,1.85.2)} yatra ekādeśaḥ na bhavati . \\
\noindent \textbf{(P\_6,1.85.2)} pibati . \\
\noindent \textbf{(P\_6,1.85.2)} vaibhaktasya ṇatve . \\
\noindent \textbf{(P\_6,1.85.2)} vaibhaktasya ṇatve prayojanam . \\
\noindent \textbf{(P\_6,1.85.2)} kṣīrapeṇa , surāpeṇa . \\
\noindent \textbf{(P\_6,1.85.2)} uttarapadavibhaktyoḥ ekādeśaḥ uttarapadasya antavat bhavati yathā śakyeta kartum ekājuttarapade ṇaḥ bhavati iti . \\
\noindent \textbf{(P\_6,1.85.2)} kva tarhi syāt . \\
\noindent \textbf{(P\_6,1.85.2)} yatra ekādeśaḥ na bhavati . \\
\noindent \textbf{(P\_6,1.85.2)} kṣīrapāṇām , surāpāṇam . \\
\noindent \textbf{(P\_6,1.85.2)} adasaḥ īttvottve . \\
\noindent \textbf{(P\_6,1.85.2)} adasaḥ īttvottve prayojanam . \\
\noindent \textbf{(P\_6,1.85.2)} amī atra , amī āsate , amū atra , amū āsāte . \\
\noindent \textbf{(P\_6,1.85.2)} adasvibhaktyoḥ ekādeśaḥ adasaḥ antavat bhavati yathā śakyeta kartum adasaḥ aseḥ dāt u daḥ maḥ etaḥ īt bahuvacane iti . \\
\noindent \textbf{(P\_6,1.85.2)} kva tarhi syāt . \\
\noindent \textbf{(P\_6,1.85.2)} yatra ekādeśaḥ na bhavati . \\
\noindent \textbf{(P\_6,1.85.2)} amībhiḥ , amūbhyām . \\
\noindent \textbf{(P\_6,1.85.2)} svaritatve prayojanam . \\
\noindent \textbf{(P\_6,1.85.2)} kāryā , hāryā . \\
\noindent \textbf{(P\_6,1.85.2)} tidatiroḥ ekādeśaḥ titaḥ antavat bhavati yathā śakyeta kartum tit svaritam iti . \\
\noindent \textbf{(P\_6,1.85.2)} kva tarhi syāt . \\
\noindent \textbf{(P\_6,1.85.2)} yatra ekādeśaḥ na bhavati . \\
\noindent \textbf{(P\_6,1.85.2)} kāryaḥ , hāryaḥ . \\
\noindent \textbf{(P\_6,1.85.2)} svaritatvam vipratiṣedhāt . \\
\noindent \textbf{(P\_6,1.85.2)} svaritatvam kriyatām ekādeśaḥ iti kim atra kartavyam . \\
\noindent \textbf{(P\_6,1.85.2)} paratvāt svaritatvam bhaviṣyati vipratiṣedhena . \\
\noindent \textbf{(P\_6,1.85.2)} na eṣaḥ yuktaḥ vipratiṣedhaḥ . \\
\noindent \textbf{(P\_6,1.85.2)} nityaḥ ekādeśaḥ . \\
\noindent \textbf{(P\_6,1.85.2)} kṛte api svaritatve prāpnoti akṛte api . \\
\noindent \textbf{(P\_6,1.85.2)} anityaḥ ekādeśaḥ . \\
\noindent \textbf{(P\_6,1.85.2)} anyathāsvarasya kṛte svaritatve prāpnoti anyathāsvarasya akṛte svaritatve prāpnoti . \\
\noindent \textbf{(P\_6,1.85.2)} svarabhinnasya ca prāpnuvan vidhiḥ anityaḥ bhavati . \\
\noindent \textbf{(P\_6,1.85.2)} antaraṅgaḥ tarhi ekādeśaḥ . \\
\noindent \textbf{(P\_6,1.85.2)} kā antaraṅgatā . \\
\noindent \textbf{(P\_6,1.85.2)} varṇau āśritya ekādeśaḥ padasya svaritatvam . \\
\noindent \textbf{(P\_6,1.85.2)} svaritatvam api antaraṅgam . \\
\noindent \textbf{(P\_6,1.85.2)} katham . \\
\noindent \textbf{(P\_6,1.85.2)} uktam etat padagrahaṇam parimāṇārtham iti . \\
\noindent \textbf{(P\_6,1.85.2)} ubhayoḥ antaraṅgayoḥ paratvāt svaritatvam . \\
\noindent \textbf{(P\_6,1.85.2)} svaritatve kṛte āntaryataḥ svaritānudāttayoḥ ekādeśaḥ svaritaḥ bhaviṣyati . \\
\noindent \textbf{(P\_6,1.85.2)} liṅgaviśiṣtagrahaṇāt vā . \\
\noindent \textbf{(P\_6,1.85.2)} atha vā prātipadikagrahaṇe liṅgaviśiṣtasya api grahaṇm bhavati iti evam atra svaritatvam bhaviṣyati . \\
\noindent \textbf{(P\_6,1.85.2)} pūrvapadāntodāttatvam ca prayojanam . \\
\noindent \textbf{(P\_6,1.85.2)} guḍodakam , mathitodakam . \\
\noindent \textbf{(P\_6,1.85.2)} pūrvapadottarapadayoḥ ekādeśaḥ pūrvapadasya antavat bhavati yathā śakyeta kartum udake akevale pūrvapadasya antaḥ udāttaḥ bhavati iti . \\
\noindent \textbf{(P\_6,1.85.2)} kva tarhi syāt . \\
\noindent \textbf{(P\_6,1.85.2)} yatra akādeśaḥ na bhavati . \\
\noindent \textbf{(P\_6,1.85.2)} udaśvidudakam . \\
\noindent \textbf{(P\_6,1.85.2)} pūrvapadāntodāttatvam ca . \\
\noindent \textbf{(P\_6,1.85.2)} pūrvapadāntodāttatvam ca vipratiṣedhāt . \\
\noindent \textbf{(P\_6,1.85.2)} pūrvapadāntodāttatvam kriyatām ekādeśaḥ iti kim atra kartavyam . \\
\noindent \textbf{(P\_6,1.85.2)} paratvāt pūrvapadāntodāttatvam . \\
\noindent \textbf{(P\_6,1.85.2)} pūrvapadāntodāttatvasya avakāśaḥ udaśvidudakam . \\
\noindent \textbf{(P\_6,1.85.2)} ekādeśasya avakāśaḥ daṇḍāgram , kṣupāgram . \\
\noindent \textbf{(P\_6,1.85.2)} iha ubhayam prāpnoti . \\
\noindent \textbf{(P\_6,1.85.2)} mathitodakam , guḍodakam . \\
\noindent \textbf{(P\_6,1.85.2)} pūrvapadāntodāttatvam bhavati vipratiṣedhena . \\
\noindent \textbf{(P\_6,1.85.2)} saḥ ca avaśyam vipratiṣedhaḥ āśrayayitavyaḥ . \\
\noindent \textbf{(P\_6,1.85.2)} ekādeśe hi svaritāprasiddhiḥ . \\
\noindent \textbf{(P\_6,1.85.2)} ekādeśe hi svaritasya aprasiddhiḥ syāt . \\
\noindent \textbf{(P\_6,1.85.2)} yaḥ hi manyate astu atra ekādeśaḥ ekādeśe kṛte pūrvapadāntodāttatvam bhaviṣyati iti svaritatvam tasya na sidhyati svaritaḥ vā anudātte padādau iti . \\
\noindent \textbf{(P\_6,1.85.2)} mathitodakam , guḍodakam . \\
\noindent \textbf{(P\_6,1.85.2)} kṛdantaprakṛtisvaratvam ca prayojanam . \\
\noindent \textbf{(P\_6,1.85.2)} prāṭitā , prāśitā . \\
\noindent \textbf{(P\_6,1.85.2)} kṛdgatyoḥ ekādeśaḥ gateḥ antavat bhavati yathā śakyeta kartum gatikārakopapadāt kṛdantam uttarapadam prakṛtisvaram bhavati iti . \\
\noindent \textbf{(P\_6,1.85.2)} kva tarhi syāt . \\
\noindent \textbf{(P\_6,1.85.2)} yatra na akādeśaḥ . \\
\noindent \textbf{(P\_6,1.85.2)} prakārakaḥ , prakaraṇam . kṛdantaprakṛtisvaratvam ca . \\
\noindent \textbf{(P\_6,1.85.2)} kṛdantaprakṛtisvaratvam ca vipratiṣedhāt . \\
\noindent \textbf{(P\_6,1.85.2)} kṛdantaprakṛtisvaratvam kriyatām ekādeśaḥ iti kim atra kartavyam . \\
\noindent \textbf{(P\_6,1.85.2)} paratvāt kṛdantaprakṛtisvaratvam bhavati vipratiṣedhena . \\
\noindent \textbf{(P\_6,1.85.2)} kṛdantaprakṛtisvaratvasya avakāśaḥ prakārakaḥ , prakaraṇam . ekādeśasya avakāśaḥ daṇḍāgram , kṣupāgram . \\
\noindent \textbf{(P\_6,1.85.2)} iha ubhayam prāpnoti . \\
\noindent \textbf{(P\_6,1.85.2)} prāṭitā , prāśitā . \\
\noindent \textbf{(P\_6,1.85.2)} kṛdantaprakṛtisvaratvam bhavati vipratiṣedhena . \\
\noindent \textbf{(P\_6,1.85.2)} saḥ ca avaśyam vipratiṣedhaḥ āśrayayitavyaḥ . \\
\noindent \textbf{(P\_6,1.85.2)} ekādeśe hi a prasiddhiḥ uttarapadasya aparatvāt . \\
\noindent \textbf{(P\_6,1.85.2)} yaḥ hi manyate astu atra ekādeśaḥ ekādeśe kṛte kṛdantaprakṛtisvaratvam bhaviṣyati iti kṛdantaprakṛtisvaratvam tasya na sidhyati . \\
\noindent \textbf{(P\_6,1.85.2)} kim kāraṇam . \\
\noindent \textbf{(P\_6,1.85.2)} uttarapadasya aparatvāt . \\
\noindent \textbf{(P\_6,1.85.2)} na hi idānīm ekādeśe kṛte uttarapadam param bhavati . \\
\noindent \textbf{(P\_6,1.85.2)} nanu ca antādivadbhāvena param . \\
\noindent \textbf{(P\_6,1.85.2)} ubhayataḥ āśraye na antādivat . \\
\noindent \textbf{(P\_6,1.85.2)} uttarapadavṛddhiḥ ca ekādeśāt . \\
\noindent \textbf{(P\_6,1.85.2)} uttarapadavṛddhiḥ ca ekādeśāt bhavati vipratiṣedhena . \\
\noindent \textbf{(P\_6,1.85.2)} uttarapadavṛddheḥ avakāśaḥ pūrvatraigartakaḥ, aparatraigartakaḥ . \\
\noindent \textbf{(P\_6,1.85.2)} ekādeśasya avakāśaḥ daṇḍāgram , kṣupāgram . \\
\noindent \textbf{(P\_6,1.85.2)} iha ubhayam prāpnoti . \\
\noindent \textbf{(P\_6,1.85.2)} pūrvaiṣukāmaśamaḥ , aparaiṣukāmaśaḥ . \\
\noindent \textbf{(P\_6,1.85.2)} uttarapadavṛddhiḥ bhavati vipratiṣedhena . \\
\noindent \textbf{(P\_6,1.85.2)} ekādeśaprasaṅgaḥ tu antaraṅgabalīyastvāt . \\
\noindent \textbf{(P\_6,1.85.2)} ekādeśaḥ tu prāpnoti . \\
\noindent \textbf{(P\_6,1.85.2)} kim kāraṇam . \\
\noindent \textbf{(P\_6,1.85.2)} antaraṅgasya balīyastvāt . \\
\noindent \textbf{(P\_6,1.85.2)} antaraṅgam balīyaḥ . \\
\noindent \textbf{(P\_6,1.85.2)} tatra kaḥ doṣaḥ . \\
\noindent \textbf{(P\_6,1.85.2)} tatra vṛddhividhānam . \\
\noindent \textbf{(P\_6,1.85.2)} tatra vṛddhiḥ vidheyā . \\
\noindent \textbf{(P\_6,1.85.2)} na eṣaḥ doṣaḥ . \\
\noindent \textbf{(P\_6,1.85.2)} ācāryapravṛttiḥ jñāpayati pūrvottarapadayoḥ tāvat kāryam bhavati na ekādeśaḥ iti yat ayam na indrasya parasya iti pratiṣedham śāsti . \\
\noindent \textbf{(P\_6,1.85.2)} katham kṛtvā jñāpakam . \\
\noindent \textbf{(P\_6,1.85.2)} indre dvau acau . \\
\noindent \textbf{(P\_6,1.85.2)} tatra ekaḥ yasya īti ca iti lopena hriyate aparaḥ ekādeśena . \\
\noindent \textbf{(P\_6,1.85.2)} tataḥ anackaḥ indraḥ sampannaḥ . \\
\noindent \textbf{(P\_6,1.85.2)} tatra kaḥ prasaṅgaḥ vṛddheḥ . \\
\noindent \textbf{(P\_6,1.85.2)} paśyati tu ācāryaḥ pūrvapadottarapadyoḥ tāvatkāryam bhavati na ekādeśaḥ iti tataḥ na indrasya parasya iti pratiṣedham śāsti . \\
\noindent \textbf{(P\_6,1.85.3)} atha ādivattve kāni prayojanāni . \\
\noindent \textbf{(P\_6,1.85.3)} ādivattve prayojanam pragṛhyasañjñāyām . \\
\noindent \textbf{(P\_6,1.85.3)} ādivattve pragṛhyasañjñāyām prayojanam . \\
\noindent \textbf{(P\_6,1.85.3)} agnī iti , vāyū iti . \\
\noindent \textbf{(P\_6,1.85.3)} dvivacanādvivacanayoḥ ekādeśaḥ dvivacanasya ādivat bhavati yathā śakyeta kartum īdūdet dvivacanam pragṛhyam iti . \\
\noindent \textbf{(P\_6,1.85.3)} kva tarhi syāt . \\
\noindent \textbf{(P\_6,1.85.3)} yatra ekādeśaḥ na bhavati . \\
\noindent \textbf{(P\_6,1.85.3)} trapuṇī iti , jatunī iti . \\
\noindent \textbf{(P\_6,1.85.3)} suptiṅābvidhiṣu . \\
\noindent \textbf{(P\_6,1.85.3)} suptiṅābvidhiṣu prayojanam . \\
\noindent \textbf{(P\_6,1.85.3)} sup . \\
\noindent \textbf{(P\_6,1.85.3)} vṛkṣe tiṣṭhati . \\
\noindent \textbf{(P\_6,1.85.3)} plakṣe tiṣṭhati . \\
\noindent \textbf{(P\_6,1.85.3)} subasupoḥ ekādeśaḥ supaḥ ādivat bhavati yathā śakyeta kartum subantam padam iti . \\
\noindent \textbf{(P\_6,1.85.3)} kva tarhi syāt . \\
\noindent \textbf{(P\_6,1.85.3)} yatra ekādeśaḥ na bhavati . \\
\noindent \textbf{(P\_6,1.85.3)} vṛkṣaḥ tiṣṭhati . \\
\noindent \textbf{(P\_6,1.85.3)} plakṣaḥ tiṣthati . \\
\noindent \textbf{(P\_6,1.85.3)} sup . \\
\noindent \textbf{(P\_6,1.85.3)} tiṅ . \\
\noindent \textbf{(P\_6,1.85.3)} pace, yaje iti . \\
\noindent \textbf{(P\_6,1.85.3)} tiṅatiṅoḥ ekādeśaḥ tiṅaḥ ādivat bhavati yathā śakyeta kartum tiṅantam padam iti . \\
\noindent \textbf{(P\_6,1.85.3)} kva tarhi syāt . \\
\noindent \textbf{(P\_6,1.85.3)} yatra ekādeśaḥ na bhavati . \\
\noindent \textbf{(P\_6,1.85.3)} pacati , yajati . \\
\noindent \textbf{(P\_6,1.85.3)} tiṅ . \\
\noindent \textbf{(P\_6,1.85.3)} āp . \\
\noindent \textbf{(P\_6,1.85.3)} khaṭvā , mālā . \\
\noindent \textbf{(P\_6,1.85.3)} abanāpoḥ ekādeśaḥ āpaḥ ādivat bhavati yathā śakyeta kartum ābantāt soḥ lopaḥ bhavati iti . \\
\noindent \textbf{(P\_6,1.85.3)} kva tarhi syāt . \\
\noindent \textbf{(P\_6,1.85.3)} yatra ekādeśaḥ na bhavati . \\
\noindent \textbf{(P\_6,1.85.3)} kruñcā , uṣṇihā , devadiśā . \\
\noindent \textbf{(P\_6,1.85.3)} āṅgrahaṇe padavidhau . \\
\noindent \textbf{(P\_6,1.85.3)} āṅgrahaṇe padavidhau prayojanam . \\
\noindent \textbf{(P\_6,1.85.3)} adya āhate . \\
\noindent \textbf{(P\_6,1.85.3)} kadā āhate . \\
\noindent \textbf{(P\_6,1.85.3)} āṅanāṅoḥ ekādeśaḥ āṅaḥ ādivat bhavati yathā śakyeta kartum āṅaḥ yamahanaḥ iti ātmanepadam bhavati iti . \\
\noindent \textbf{(P\_6,1.85.3)} kva tarhi syāt . \\
\noindent \textbf{(P\_6,1.85.3)} yatra ekādeśaḥ na bhavati . \\
\noindent \textbf{(P\_6,1.85.3)} āhate . \\
\noindent \textbf{(P\_6,1.85.3)} āṭaḥ ca vṛddhividhau . \\
\noindent \textbf{(P\_6,1.85.3)} āṭaḥ ca vṛddhividhau prayojanam . \\
\noindent \textbf{(P\_6,1.85.3)} adya aihiṣṭa . \\
\noindent \textbf{(P\_6,1.85.3)} kadā aihiṣṭa . \\
\noindent \textbf{(P\_6,1.85.3)} āṭaḥ adyaśabdasya ca ekādeśaḥ āṭaḥ ādivat bhavati yathā śakyeta kartum āṭaḥ ca aci vṛddhiḥ bhavati iti . \\
\noindent \textbf{(P\_6,1.85.3)} kva tarhi syāt . \\
\noindent \textbf{(P\_6,1.85.3)} yatra ekādeśaḥ na . \\
\noindent \textbf{(P\_6,1.85.3)} aihiṣṭa , aikṣiṣṭa . \\
\noindent \textbf{(P\_6,1.85.3)} kṛdantaprātipadikatve ca . \\
\noindent \textbf{(P\_6,1.85.3)} kṛdantaprātipadikatve ca prayojanam . \\
\noindent \textbf{(P\_6,1.85.3)} dhārayaḥ , pārayaḥ . \\
\noindent \textbf{(P\_6,1.85.3)} kṛdakṛtoḥ ekādeśaḥ kṛtaḥ ādivat bhavati yathā śakyeta kartum kṛdantam prātipadikam iti . \\
\noindent \textbf{(P\_6,1.85.3)} kva tarhi syāt . \\
\noindent \textbf{(P\_6,1.85.3)} yatra ekādeśaḥ na . \\
\noindent \textbf{(P\_6,1.85.3)} kārakaḥ , hārakaḥ . \\
\noindent \textbf{(P\_6,1.85.4)} na abhyāsādīnām hrasvatve . \\
\noindent \textbf{(P\_6,1.85.4)} abhyāsādīnām hrasvatve na antādivat bhavati iti vaktavyam . \\
\noindent \textbf{(P\_6,1.85.4)} ke punaḥ abhyāsādayaḥ . \\
\noindent \textbf{(P\_6,1.85.4)} abhyāsohāmbārthanadīnapuṃsakopasarjanahrasvatvāni . \\
\noindent \textbf{(P\_6,1.85.4)} abhyāsahrasvatvam . \\
\noindent \textbf{(P\_6,1.85.4)} upeyāja , upovāpa . \\
\noindent \textbf{(P\_6,1.85.4)} ūheḥ hrasvavam . \\
\noindent \textbf{(P\_6,1.85.4)} upohyate , prohyate , parohyate . \\
\noindent \textbf{(P\_6,1.85.4)} ambārthanadīnapuṃsakopasarjanahrasvatvāni . \\
\noindent \textbf{(P\_6,1.85.4)} amba atra , akka atra . \\
\noindent \textbf{(P\_6,1.85.4)} kumāri idam , kiśori idam . \\
\noindent \textbf{(P\_6,1.85.4)} ārāśastri idam , dhānāśaṣkuli idam . \\
\noindent \textbf{(P\_6,1.85.4)} niṣkauśāmbi idam , nirvārāṇasi idam . \\
\noindent \textbf{(P\_6,1.85.4)} abhyāsohāmbārthanadīnapuṃsakopasarjanagrahaṇena grahaṇāt hrasvatvam prāpnoti . \\
\noindent \textbf{(P\_6,1.85.4)} na vā bahiraṅgalakṣaṇatvāt . \\
\noindent \textbf{(P\_6,1.85.4)} na vā etat vaktavyam . \\
\noindent \textbf{(P\_6,1.85.4)} kim kāraṇam . \\
\noindent \textbf{(P\_6,1.85.4)} bahiraṅgalakṣaṇatvāt . \\
\noindent \textbf{(P\_6,1.85.4)} antaraṅgam hrasvatvam . \\
\noindent \textbf{(P\_6,1.85.4)} bahiraṅgāḥ ete vidhayaḥ . \\
\noindent \textbf{(P\_6,1.85.4)} asiddham bahiraṅgam antaraṅge . \\
\noindent \textbf{(P\_6,1.85.4)} varṇāśrayavidhau ca . \\
\noindent \textbf{(P\_6,1.85.4)} varṇāśrayavidhau ca na antādivat bhavati iti vaktavyam . \\
\noindent \textbf{(P\_6,1.85.4)} kim prayojanam . \\
\noindent \textbf{(P\_6,1.85.4)} prayojanam khaṭvābhiḥ juhāva asyai aśvaḥ iti . \\
\noindent \textbf{(P\_6,1.85.4)} iha khaṭvābhiḥ , mālābhiḥ , ataḥ bhisaḥ ais bhavati iti aisbhāvaḥ prāpnoti . \\
\noindent \textbf{(P\_6,1.85.4)} na eṣaḥ doṣaḥ . \\
\noindent \textbf{(P\_6,1.85.4)} taparakaraṇasāmarthyāt na bhaviṣyati . \\
\noindent \textbf{(P\_6,1.85.4)} asti anyat taparakaraṇe prayojanam . \\
\noindent \textbf{(P\_6,1.85.4)} kim . \\
\noindent \textbf{(P\_6,1.85.4)} kīlālapābhiḥ , śubhaṃyābhiḥ . \\
\noindent \textbf{(P\_6,1.85.4)} juhāva . \\
\noindent \textbf{(P\_6,1.85.4)} ātaḥ au ṇalaḥ iti autvam prāpnoti . \\
\noindent \textbf{(P\_6,1.85.4)} asyai aśvaḥ iti . \\
\noindent \textbf{(P\_6,1.85.4)} eṅaḥ padāntāt ati iti pūrvatvam prāpnoti . \\
\noindent \textbf{(P\_6,1.85.4)} na vā atādrūpyātideśāt . \\
\noindent \textbf{(P\_6,1.85.4)} na vā vaktavyam . \\
\noindent \textbf{(P\_6,1.85.4)} kim kāraṇam . \\
\noindent \textbf{(P\_6,1.85.4)} atādrūpyātideśāt . \\
\noindent \textbf{(P\_6,1.85.4)} na iha tādrūpyam atidiśyate . \\
\noindent \textbf{(P\_6,1.85.4)} rūpāśrayāḥ vai ete vidhayaḥ atādrūpyāt na bhaviṣyanti . \\
\noindent \textbf{(P\_6,1.86.1)} kimartham idam ucyate . \\
\noindent \textbf{(P\_6,1.86.1)} ṣatvatukoḥ asiddhavacanam ādeśalakṣaṇapratiṣedhāṛtham utsargalakṣaṇabhāvārtham ca . \\
\noindent \textbf{(P\_6,1.86.1)} ṣatvatukoḥ asiddhatvam ucyate ādeśalakṣaṇapratiṣedhāṛtham utsargalakṣaṇabhāvārtham ca . \\
\noindent \textbf{(P\_6,1.86.1)} ādeśalakṣaṇapratiṣedhāṛtham tāvat . \\
\noindent \textbf{(P\_6,1.86.1)} kosiñcat . \\
\noindent \textbf{(P\_6,1.86.1)} yosiñcat . \\
\noindent \textbf{(P\_6,1.86.1)} ekādeśe kṛte iṇaḥ iti ṣatvam prāpnoti . \\
\noindent \textbf{(P\_6,1.86.1)} asiddhatvāt na bhavati . \\
\noindent \textbf{(P\_6,1.86.1)} utsargalakṣaṇabhāvārtham ca . \\
\noindent \textbf{(P\_6,1.86.1)} adhītya , pretya . \\
\noindent \textbf{(P\_6,1.86.1)} ekādeśe kṛte hrasvasya iti tuk na prāpnoti . \\
\noindent \textbf{(P\_6,1.86.1)} asiddhatvāt bhavati . \\
\noindent \textbf{(P\_6,1.86.1)} asti prayojanam etat . \\
\noindent \textbf{(P\_6,1.86.1)} kim tarhi iti . \\
\noindent \textbf{(P\_6,1.86.1)} tatra utsargalakṣaṇāprasiddhiḥ utsargābhāvāt . \\
\noindent \textbf{(P\_6,1.86.1)} tatra utsargalakṣaṇasya kāryasya aprasiddhiḥ . \\
\noindent \textbf{(P\_6,1.86.1)} adhītya , pretya iti . \\
\noindent \textbf{(P\_6,1.86.1)} kim kāraṇam . \\
\noindent \textbf{(P\_6,1.86.1)} utsargābhāvāt . \\
\noindent \textbf{(P\_6,1.86.1)} hrasvasya iti ucyate na ca atra hrasvam paśyāmaḥ . \\
\noindent \textbf{(P\_6,1.86.1)} nanu ca atra api asiddhavacanāt siddham . \\
\noindent \textbf{(P\_6,1.86.1)} asiddhavacanāt siddham iti cet na anyasya asiddhavacanāt anyasya bhāvaḥ . \\
\noindent \textbf{(P\_6,1.86.1)} asiddhavacanāt siddham iti cet tat na . \\
\noindent \textbf{(P\_6,1.86.1)} kim kāraṇam . \\
\noindent \textbf{(P\_6,1.86.1)} anyasya asiddhavacanāt anyasya bhāvaḥ . \\
\noindent \textbf{(P\_6,1.86.1)} na hi anyasya asiddhavacanāt anyasya prādurbhāvaḥ bhavati . \\
\noindent \textbf{(P\_6,1.86.1)} na hi devadattasya hantari hate devadattasya prādurbhāvaḥ bhavati . \\
\noindent \textbf{(P\_6,1.86.1)} tasmāt sthānivadvacanam asiddhatvam ca . \\
\noindent \textbf{(P\_6,1.86.1)} tasmāt sthānivadbhāvaḥ vaktavyaḥ asiddhatvam ca . \\
\noindent \textbf{(P\_6,1.86.1)} adhītya , pretya iti sthānivadbhāvaḥ . \\
\noindent \textbf{(P\_6,1.86.1)} kosiñcat , yosiñcat iti atra asiddhatvam . \\
\noindent \textbf{(P\_6,1.86.1)} sthānivadvacanānarthakyam śāstrāsiddhatvāt . \\
\noindent \textbf{(P\_6,1.86.1)} sthānivadvacanam anarthakam . \\
\noindent \textbf{(P\_6,1.86.1)} kim kāraṇam . \\
\noindent \textbf{(P\_6,1.86.1)} śāstrāsiddhatvāt . \\
\noindent \textbf{(P\_6,1.86.1)} na anena kāryāsiddhatvam kriyate . \\
\noindent \textbf{(P\_6,1.86.1)} kim tarhi . \\
\noindent \textbf{(P\_6,1.86.1)} śāstrāsiddhatvam anena kriyate . \\
\noindent \textbf{(P\_6,1.86.1)} ekādeśaśāstram tukśāstre asiddham bhavati iti . \\
\noindent \textbf{(P\_6,1.86.2)} samprasāraṇaṅīṭsu siddhaḥ . \\
\noindent \textbf{(P\_6,1.86.2)} samprasāraṇaṅīṭsu siddhaḥ ekādeśaḥ iti vaktavyam . \\
\noindent \textbf{(P\_6,1.86.2)} śakahūṣu , parivīṣu . \\
\noindent \textbf{(P\_6,1.86.2)} samprasāraṇa . \\
\noindent \textbf{(P\_6,1.86.2)} ṅi . \\
\noindent \textbf{(P\_6,1.86.2)} vṛkṣe cchatram , vṛkṣe chatram . \\
\noindent \textbf{(P\_6,1.86.2)} ṅi . \\
\noindent \textbf{(P\_6,1.86.2)} iṭ . \\
\noindent \textbf{(P\_6,1.86.2)} apace cchatram , apace chatram . \\
\noindent \textbf{(P\_6,1.86.2)} samprasāraṇaṅīṭsu siddhaḥ padāntapadādyoḥ ekādeśasya asiddhavacanāt . samprasāraṇaṅīṭsu siddhaḥ ekādeśaḥ . \\
\noindent \textbf{(P\_6,1.86.2)} kutaḥ . \\
\noindent \textbf{(P\_6,1.86.2)} padāntapadādyoḥ ekādeśasya asiddhavacanāt . \\
\noindent \textbf{(P\_6,1.86.2)} padāntapadādyoḥ ekādeśaḥ asiddhaḥ bhavati iti ucyate na ca eṣaḥ padāntapadādyoḥ ekādeśaḥ . \\
\noindent \textbf{(P\_6,1.86.2)} yadi padāntapadādyoḥ ekādeśaḥ asiddhaḥ susasyāḥ oṣadhīḥ kṛdhi , supippalāḥ oṣadhīḥ kṛdhi , atra ṣatvam prāpnoti . \\
\noindent \textbf{(P\_6,1.86.2)} tugvidhim prati padāntapadādyoḥ ekādeśaḥ asiddhaḥ . \\
\noindent \textbf{(P\_6,1.86.2)} ṣatvam prati ekādeśamātram asiddham bhavati . \\
\noindent \textbf{(P\_6,1.86.2)} yadi ṣatvam prati ekādeśamātram asiddham śakahūṣu , parivīṣu , atra ṣatvam na prāpnoti . \\
\noindent \textbf{(P\_6,1.86.2)} astu tarhi aviśeṣeṇa . \\
\noindent \textbf{(P\_6,1.86.2)} katham susasyāḥ oṣadhīḥ kṛdhi , supippalāḥ oṣadhīḥ kṛdhi iti . \\
\noindent \textbf{(P\_6,1.86.2)} na eṣaḥ doṣaḥ . \\
\noindent \textbf{(P\_6,1.86.2)} bhrātuṣputragrahaṇam jñāpakam ekādeśanimittāt ṣatvapratiṣedhasya . \\
\noindent \textbf{(P\_6,1.86.2)} yat ayam kaskādiṣu bhrātuṣputragrahaṇam karoti tat jñāpayati ācāryaḥ na ekādeśanimittāt ṣatvam bhavati iti . \\
\noindent \textbf{(P\_6,1.86.2)} yadi etat jñāpyate śakahūṣu , parivīṣu iti atra ṣatvam na prāpnoti . \\
\noindent \textbf{(P\_6,1.86.2)} tulyajātīyakasya jñāpakam . \\
\noindent \textbf{(P\_6,1.86.2)} kim ca tuljyajātīyam . \\
\noindent \textbf{(P\_6,1.86.2)} yaḥ kupvoḥ . \\
\noindent \textbf{(P\_6,1.86.2)} yadi evam veñaḥ apratyaye parataḥ uḥ iti prāpnoti ut iti ca iṣyate . \\
\noindent \textbf{(P\_6,1.86.2)} yathālakṣaṇam aprayukte . \\
\noindent \textbf{(P\_6,1.86.2)} atha vā na evam vijñāyate . \\
\noindent \textbf{(P\_6,1.86.2)} pūrvasya ca padādeḥ parasya ca padāntasya iti . \\
\noindent \textbf{(P\_6,1.86.2)} katham tarhi . \\
\noindent \textbf{(P\_6,1.86.2)} parasya ca padādeḥ pūrvasya ca padāntasya iti . \\
\noindent \textbf{(P\_6,1.87.1)} guṇagrahaṇam kimartham na āt ekaḥ bhavati iti eva ucyeta . \\
\noindent \textbf{(P\_6,1.87.1)} āt ekaḥ cet guṇaḥ kena . \\
\noindent \textbf{(P\_6,1.87.1)} āt ekaḥ cet guṇaḥ kena idānīm bhaviṣyati . \\
\noindent \textbf{(P\_6,1.87.1)} khaṭvendraḥ , mālendraḥ , khaṭvodakam , mālodakam . \\
\noindent \textbf{(P\_6,1.87.1)} sthāne antaratamaḥ hi saḥ . \\
\noindent \textbf{(P\_6,1.87.1)} sthāne prāpyamāṇānām antaratamaḥ ādeśaḥ bhavati . \\
\noindent \textbf{(P\_6,1.87.1)} aidautau api tarhi prapnutaḥ . \\
\noindent \textbf{(P\_6,1.87.1)} aidautau na eci tau uktau . aidautau na bhaviṣyataḥ . \\
\noindent \textbf{(P\_6,1.87.1)} kim kāraṇam . \\
\noindent \textbf{(P\_6,1.87.1)} eci hi aidautau ucyete . \\
\noindent \textbf{(P\_6,1.87.1)} iha tarhi khaṭvarśyaḥ , mālarśyaḥ , ṛkāraḥ tarhi prāpnoti . \\
\noindent \textbf{(P\_6,1.87.1)} ṛkāraḥ na ubhayāntaraḥ . \\
\noindent \textbf{(P\_6,1.87.1)} ubhayoḥ yaḥ antaratamaḥ tena bhavitavyam . \\
\noindent \textbf{(P\_6,1.87.1)} na ca ṛkāraḥ ubhayoḥ antaratamaḥ . \\
\noindent \textbf{(P\_6,1.87.1)} ākāraḥ tarhi prāpnoti . \\
\noindent \textbf{(P\_6,1.87.1)} ākāraḥ na ṛti dhātau saḥ . \\
\noindent \textbf{(P\_6,1.87.1)} ākāraḥ na bhaviṣyati . \\
\noindent \textbf{(P\_6,1.87.1)} kim kāraṇam . \\
\noindent \textbf{(P\_6,1.87.1)} ṛti dhātau ākāraḥ ucyate . \\
\noindent \textbf{(P\_6,1.87.1)} tat niyamārtham bhaviṣyati . \\
\noindent \textbf{(P\_6,1.87.1)} ṛkārādau dhātau eva na anyatra iti . \\
\noindent \textbf{(P\_6,1.87.1)} plutaḥ tarhi prāpnoti . \\
\noindent \textbf{(P\_6,1.87.1)} plutaḥ ca viṣaye smṛtaḥ . \\
\noindent \textbf{(P\_6,1.87.1)} viṣaye plutaḥ ucyate . \\
\noindent \textbf{(P\_6,1.87.1)} yadā ca saḥ viṣayaḥ bhavitavyam tadā plutena . \\
\noindent \textbf{(P\_6,1.87.1)} āntaryāt trimātracaturmātrāḥ . idam tarhi prayojanam . \\
\noindent \textbf{(P\_6,1.87.1)} āntaryataḥ trimātracaturmātrāṇām sthāne trimātracaturmātrāḥ ādeśāḥ mā bhūvan iti . \\
\noindent \textbf{(P\_6,1.87.1)} khaṭvā indraḥ khaṭvendraḥ . \\
\noindent \textbf{(P\_6,1.87.1)} khaṭvā udakam khaṭvodakam . \\
\noindent \textbf{(P\_6,1.87.1)} khaṭvā īṣā khaṭveṣā . \\
\noindent \textbf{(P\_6,1.87.1)} khaṭvā ūḍhā khaṭvoḍhā . \\
\noindent \textbf{(P\_6,1.87.1)} khaṭvā elakā khaṭvailakā . \\
\noindent \textbf{(P\_6,1.87.1)} khaṭvā odanaḥ khaṭvaudanaḥ khaṭvā aitikāyanaḥ khaṭvaitikāyanaḥ . \\
\noindent \textbf{(P\_6,1.87.1)} khaṭvā aupagavaḥ khaṭvaupagavaḥ . \\
\noindent \textbf{(P\_6,1.87.1)} atha kriyamāṇe api guṇagrahaṇe kasmāt eva atra trimātracaturmātrāṇām sthāne trimātracaturmātrāḥ ādeśāḥ na bhavanti . \\
\noindent \textbf{(P\_6,1.87.1)} taparatvāt ne te smṛtāḥ . \\
\noindent \textbf{(P\_6,1.87.1)} tapare guṇavṛddhī . \\
\noindent \textbf{(P\_6,1.87.1)} nanu ca bhoḥ taḥ paraḥ yasmāt saḥ ayam taparaḥ . \\
\noindent \textbf{(P\_6,1.87.1)} na iti āha . \\
\noindent \textbf{(P\_6,1.87.1)} tāt api paraḥ taparaḥ iti . \\
\noindent \textbf{(P\_6,1.87.1)} yadi tāt api paraḥ taparaḥ ṝdoḥ ap iti iha eva syāt . \\
\noindent \textbf{(P\_6,1.87.1)} yavaḥ stavaḥ . \\
\noindent \textbf{(P\_6,1.87.1)} lavaḥ pavaḥ iti atra na syāt . \\
\noindent \textbf{(P\_6,1.87.1)} na eṣaḥ takāraḥ . \\
\noindent \textbf{(P\_6,1.87.1)} kaḥ tarhi . \\
\noindent \textbf{(P\_6,1.87.1)} dakāraḥ . \\
\noindent \textbf{(P\_6,1.87.1)} kim dakāre prayojanam . \\
\noindent \textbf{(P\_6,1.87.1)} atha kim takāre prayojanam . \\
\noindent \textbf{(P\_6,1.87.1)} yadi asandehārthaḥ takāraḥ dakāraḥ api . \\
\noindent \textbf{(P\_6,1.87.2)} guṇe ṅiśītām upasaṅkhyānam dīrghatvabādhanārtham . \\
\noindent \textbf{(P\_6,1.87.2)} guṇe ṅiśītām upasaṅkhyānam kartavyam . \\
\noindent \textbf{(P\_6,1.87.2)} ṅi . \\
\noindent \textbf{(P\_6,1.87.2)} vṛkṣe indraḥ , plakṣe indraḥ . \\
\noindent \textbf{(P\_6,1.87.2)} śī . \\
\noindent \textbf{(P\_6,1.87.2)} ye indram , te indram . \\
\noindent \textbf{(P\_6,1.87.2)} iṭ . \\
\noindent \textbf{(P\_6,1.87.2)} apace indram , ayaje indram . \\
\noindent \textbf{(P\_6,1.87.2)} kim prayojanam . \\
\noindent \textbf{(P\_6,1.87.2)} dīrghatvabādhanārtham . \\
\noindent \textbf{(P\_6,1.87.2)} savarṇadīrghatvam mā bhūt iti . \\
\noindent \textbf{(P\_6,1.87.2)} na vā bahiraṅgalakṣaṇatvāt . \\
\noindent \textbf{(P\_6,1.87.2)} na vā kartavyam . \\
\noindent \textbf{(P\_6,1.87.2)} kim kāraṇam . \\
\noindent \textbf{(P\_6,1.87.2)} bahiraṅgalakṣaṇatvāt . \\
\noindent \textbf{(P\_6,1.87.2)} bahiraṅgalakṣaṇam savarṇadīrghatvam . \\
\noindent \textbf{(P\_6,1.87.2)} asiddham bahiraṅgam antaraṅge . \\
\noindent \textbf{(P\_6,1.87.3)} āt ekaḥ cet guṇaḥ kena . \\
\noindent \textbf{(P\_6,1.87.3)} sthāne antaratamaḥ hi saḥ . \\
\noindent \textbf{(P\_6,1.87.3)} aidautau na eci tau uktau . \\
\noindent \textbf{(P\_6,1.87.3)} ṛkāraḥ na ubhayāntaraḥ . \\
\noindent \textbf{(P\_6,1.87.3)} ākāraḥ na ṛti dhātau saḥ . \\
\noindent \textbf{(P\_6,1.87.3)} plutaḥ ca viṣaye smṛtaḥ . \\
\noindent \textbf{(P\_6,1.87.3)} āntaryāt trimātracaturmātrāḥ . \\
\noindent \textbf{(P\_6,1.87.3)} taparatvāt ne te smṛtāḥ . \\
\noindent \textbf{(P\_6,1.89.1)} kim idam etyedhatyoḥ rūpagrahaṇam āhosvit dhātugrahaṇam . \\
\noindent \textbf{(P\_6,1.89.1)} kim ca ataḥ . \\
\noindent \textbf{(P\_6,1.89.1)} yadi rūpagrahaṇam siddham upaiti , praiti . \\
\noindent \textbf{(P\_6,1.89.1)} upaiṣi , praiṣi iti na sidhyati . \\
\noindent \textbf{(P\_6,1.89.1)} atha dhātugrahaṇam siddham etat bhavati . \\
\noindent \textbf{(P\_6,1.89.1)} kim tarhi iti . \\
\noindent \textbf{(P\_6,1.89.1)} iṇi ikārādau vṛddhipratiṣedhaḥ . \\
\noindent \textbf{(P\_6,1.89.1)} iṇi ikārādau vṛddheḥ pratiṣedhaḥ vaktavyaḥ . \\
\noindent \textbf{(P\_6,1.89.1)} upetaḥ pretaḥ iti . \\
\noindent \textbf{(P\_6,1.89.1)} yogavibhāgāt siddham . \\
\noindent \textbf{(P\_6,1.89.1)} yogavibhāgaḥ kariṣyate . \\
\noindent \textbf{(P\_6,1.89.1)} vṛddhiḥ eci . \\
\noindent \textbf{(P\_6,1.89.1)} tataḥ etyedhatyoḥ . \\
\noindent \textbf{(P\_6,1.89.1)} etyedhatyoḥ ca eci vṛddhiḥ bhavati . \\
\noindent \textbf{(P\_6,1.89.1)} tata ūṭhi . \\
\noindent \textbf{(P\_6,1.89.1)} ūṭhi ca vṛddhiḥ bhavati . \\
\noindent \textbf{(P\_6,1.89.1)} evam api ā itaḥ etaḥ . \\
\noindent \textbf{(P\_6,1.89.1)} upetaḥ , pretaḥ iti atra api prāpnoti . \\
\noindent \textbf{(P\_6,1.89.1)} āṅi pararaūpam atra bādhakam bhaviṣyati . \\
\noindent \textbf{(P\_6,1.89.1)} na aprāpte pararūpam iyam vṛddhiḥ ārabhyate . \\
\noindent \textbf{(P\_6,1.89.1)} sā yathā eṅi pararūpam bādhate evam āṅi pararūpam bādheta . \\
\noindent \textbf{(P\_6,1.89.1)} na bādhate . \\
\noindent \textbf{(P\_6,1.89.1)} kim kāraṇam . \\
\noindent \textbf{(P\_6,1.89.1)} yena na aprāpte tasya bādhanam bhavati . \\
\noindent \textbf{(P\_6,1.89.1)} na ca aprāpte eṅi pararūpam iyam vṛddhiḥ ārabhyate . \\
\noindent \textbf{(P\_6,1.89.1)} āṅi pararūpe punaḥ prāpte ca aprāpte ca . \\
\noindent \textbf{(P\_6,1.89.1)} athavā purastāt apavādāḥ anantarān vidhīn bādhante iti iyam vṛddhiḥ eṅi pararūpam bādhiṣyate na āṅi pararūpam . \\
\noindent \textbf{(P\_6,1.89.2)} akṣāt ūhinyām . \\
\noindent \textbf{(P\_6,1.89.2)} akṣāt ūhinyām vṛddhiḥ vaktavyā . \\
\noindent \textbf{(P\_6,1.89.2)} akṣauhiṇī . \\
\noindent \textbf{(P\_6,1.89.2)} prāt ūhoḍhoḍhyeṣaiṣyeṣu . \\
\noindent \textbf{(P\_6,1.89.2)} prāt ūha, ūḍha, ūḍhi, eṣa, eṣya iti eteṣu vṛddhiḥ vaktavyā . \\
\noindent \textbf{(P\_6,1.89.2)} prauhaḥ , prauḍhaḥ , pruḍhiḥ , praiṣaḥ , praiṣyaḥ . \\
\noindent \textbf{(P\_6,1.89.2)} svāt īreriṇoḥ . \\
\noindent \textbf{(P\_6,1.89.2)} svāt īra , īrin iti etayoḥ vṛddhiḥ vaktavyā . \\
\noindent \textbf{(P\_6,1.89.2)} svairaḥ , svairī . \\
\noindent \textbf{(P\_6,1.89.2)} īringrahaṇam śakyam akartum . \\
\noindent \textbf{(P\_6,1.89.2)} katham svarī iti . \\
\noindent \textbf{(P\_6,1.89.2)} ininā etat matvarthīyena siddham . \\
\noindent \textbf{(P\_6,1.89.2)} svairaḥ asya asti iti svairī . \\
\noindent \textbf{(P\_6,1.89.2)} ṛte ca tṛtīyāsamāse . \\
\noindent \textbf{(P\_6,1.89.2)} ṛte ca tṛtīyāsamāse vṛddhiḥ vaktavyā . \\
\noindent \textbf{(P\_6,1.89.2)} sukhārtaḥ , duḥkhārtaḥ . \\
\noindent \textbf{(P\_6,1.89.2)} ṛte iti kim . \\
\noindent \textbf{(P\_6,1.89.2)} sukhetaḥ , duḥkhetaḥ . \\
\noindent \textbf{(P\_6,1.89.2)} tṛtīyāgrahaṇam kim . \\
\noindent \textbf{(P\_6,1.89.2)} paramartaḥ . \\
\noindent \textbf{(P\_6,1.89.2)} samāse iti kim . \\
\noindent \textbf{(P\_6,1.89.2)} sukhenartaḥ . \\
\noindent \textbf{(P\_6,1.89.2)} pravatsatarakambalvasanānām ca ṛṇe . \\
\noindent \textbf{(P\_6,1.89.2)} pravatsatarakambalvasanānām ca ṛṇe vṛddhiḥ vaktavyā . \\
\noindent \textbf{(P\_6,1.89.2)} prārṇam , vatsatarāṇam , vasanārṇam . \\
\noindent \textbf{(P\_6,1.89.2)} ṛṇadaśābhyām ca .ṛṇadaśābhyām ca vṛddhiḥ vaktavyā . \\
\noindent \textbf{(P\_6,1.89.2)} ṛṇārṇam , daśārṇam . \\
\noindent \textbf{(P\_6,1.90)} kimarthaḥ cakāraḥ . \\
\noindent \textbf{(P\_6,1.90)} vṛddheḥ anukarṣaṇārthaḥ . \\
\noindent \textbf{(P\_6,1.90)} na etat asti prayojanam . \\
\noindent \textbf{(P\_6,1.90)} prakṛtā vṛddhiḥ anuvartiṣyate . \\
\noindent \textbf{(P\_6,1.90)} idam tarhi prayojanam . \\
\noindent \textbf{(P\_6,1.90)} ātaḥ aci vṛddhiḥ eva yathā syāt . \\
\noindent \textbf{(P\_6,1.90)} yat anyat prāpnoti tat mā bhūt iti . \\
\noindent \textbf{(P\_6,1.90)} kim ca anyat prāpnoti . \\
\noindent \textbf{(P\_6,1.90)} pararūpam . \\
\noindent \textbf{(P\_6,1.90)} usi omāṅkṣu āṭaḥ pararūpapratiṣedham codayiṣyati . \\
\noindent \textbf{(P\_6,1.90)} saḥ na vaktavyaḥ bhavati . \\
\noindent \textbf{(P\_6,1.91.1)} dhātau iti kimartham . \\
\noindent \textbf{(P\_6,1.91.1)} iha mā bhūt . \\
\noindent \textbf{(P\_6,1.91.1)} prarṣabham vanam . \\
\noindent \textbf{(P\_6,1.91.1)} upasargāt vṛddhividhau dhātugrahaṇe uktam . \\
\noindent \textbf{(P\_6,1.91.1)} kim uktam . \\
\noindent \textbf{(P\_6,1.91.1)} gatyupasargasañjñāḥ kriyāyoge yatkriyāyuktāḥ prādayaḥ tam prati iti vacanam iti . \\
\noindent \textbf{(P\_6,1.91.1)} kriyamāṇe api dhātugrahaṇe prarcchakaḥ iti prāpnoti . \\
\noindent \textbf{(P\_6,1.91.1)} yatkriyāyuktāḥ prādayaḥ tam prati iti vacanāt na bhavati . \\
\noindent \textbf{(P\_6,1.91.1)} idam tarhi prayojanam . \\
\noindent \textbf{(P\_6,1.91.1)} upasargāt ṛti dhātau vṛddhiḥ eva yathā syāt . \\
\noindent \textbf{(P\_6,1.91.1)} yat anyat prāpnoti tat mā bhūt iti . \\
\noindent \textbf{(P\_6,1.91.1)} kim ca anyat prāpnoti . \\
\noindent \textbf{(P\_6,1.91.1)} hrasvatvam . \\
\noindent \textbf{(P\_6,1.91.1)} ṛti akaḥ iti . \\
\noindent \textbf{(P\_6,1.91.1)} ṛti hrasvāt upasargāt vṛddhiḥ pūrvavipratiṣedhena iti codayiṣyati . \\
\noindent \textbf{(P\_6,1.91.1)} saḥ na vaktavyaḥ bhavati . \\
\noindent \textbf{(P\_6,1.91.2)} che tukaḥ sambuddhiguṇaḥ . \\
\noindent \textbf{(P\_6,1.91.2)} che tuk bhavati iti asmāt sambuddhiguṇaḥ bhavati vipratiṣedhena . \\
\noindent \textbf{(P\_6,1.91.2)} che tuk bhavati iti asya avakāśaḥ . \\
\noindent \textbf{(P\_6,1.91.2)} icchati , gacchati . \\
\noindent \textbf{(P\_6,1.91.2)} sambuddhiguṇasya avakāśaḥ . \\
\noindent \textbf{(P\_6,1.91.2)} agne, vāyo . \\
\noindent \textbf{(P\_6,1.91.2)} iha ubhayam prāpnoti . \\
\noindent \textbf{(P\_6,1.91.2)} agnec chatram , agne chatram , vāyoc chatram , vāyo chatram . \\
\noindent \textbf{(P\_6,1.91.2)} sambuddhiguṇaḥ bhavati vipratiṣedhena . \\
\noindent \textbf{(P\_6,1.91.2)} saḥ tarhi vipratiṣedhaḥ vaktavyaḥ . \\
\noindent \textbf{(P\_6,1.91.2)} na vā bahiraṅgalakṣaṇatvāt . na vā vaktavyaḥ . \\
\noindent \textbf{(P\_6,1.91.2)} kim kāṛaṇam . \\
\noindent \textbf{(P\_6,1.91.2)} bahiraṅgalakṣaṇatvāt . \\
\noindent \textbf{(P\_6,1.91.2)} bahiraṅgalakṣaṇaḥ tuk antaraṅgalakṣaṇaḥ sambuddhiguṇaḥ . \\
\noindent \textbf{(P\_6,1.91.2)} asiddham bahiraṅgam antaraṅge . \\
\noindent \textbf{(P\_6,1.91.2)} antareṇa vipratiṣedham antareṇa api ca etām paribhāṣām siddham . \\
\noindent \textbf{(P\_6,1.91.2)} katham . \\
\noindent \textbf{(P\_6,1.91.2)} idam iha sampradhāryam . \\
\noindent \textbf{(P\_6,1.91.2)} sambuddhilopaḥ kriyatām guṇaḥ iti . \\
\noindent \textbf{(P\_6,1.91.2)} kim atra kartavyam . \\
\noindent \textbf{(P\_6,1.91.2)} paratvāt guṇaḥ . \\
\noindent \textbf{(P\_6,1.91.2)} nityaḥ sambuddhilopaḥ . \\
\noindent \textbf{(P\_6,1.91.2)} kṛte api guṇe prāpnoti akṛte api . \\
\noindent \textbf{(P\_6,1.91.2)} guṇaḥ api nityaḥ . \\
\noindent \textbf{(P\_6,1.91.2)} kṛte api samubuddhilope prāpnoti akṛte api . \\
\noindent \textbf{(P\_6,1.91.2)} anityaḥ guṇaḥ . \\
\noindent \textbf{(P\_6,1.91.2)} na hi kṛte sambuddhilope prāpnoti . \\
\noindent \textbf{(P\_6,1.91.2)} tāvati eva chena ānantaryam . \\
\noindent \textbf{(P\_6,1.91.2)} tatra tukā bhavitavyam . \\
\noindent \textbf{(P\_6,1.91.2)} tasmāt suṣṭhu ucyate . \\
\noindent \textbf{(P\_6,1.91.2)} che tukaḥ sambuddhiguṇaḥ . \\
\noindent \textbf{(P\_6,1.91.2)} na vā bahiraṅgalakṣaṇatvāt iti . \\
\noindent \textbf{(P\_6,1.91.2)} samprasāraṇadīrghatvaṇyallopābhyāsaguṇādayaḥ ca . \\
\noindent \textbf{(P\_6,1.91.2)} samprasāraṇadīrghatvaṇyallopābhyāsaguṇādayaḥ ca tukaḥ bhavanti vipratiṣedhena . \\
\noindent \textbf{(P\_6,1.91.2)} samprasāraṇadīrghatvasya avakāśaḥ . \\
\noindent \textbf{(P\_6,1.91.2)} hūtaḥ , jīnaḥ , saṃvītaḥ , śūnaḥ . \\
\noindent \textbf{(P\_6,1.91.2)} tukaḥ avakāśaḥ . \\
\noindent \textbf{(P\_6,1.91.2)} agnicit , somasut . \\
\noindent \textbf{(P\_6,1.91.2)} iha ubhayam prāpnoti . \\
\noindent \textbf{(P\_6,1.91.2)} parivīṣu , śakahūṣu . \\
\noindent \textbf{(P\_6,1.91.2)} ṇilopasya avakāśaḥ .kāraṇā , hāraṇā . \\
\noindent \textbf{(P\_6,1.91.2)} tukaḥ saḥ eva . \\
\noindent \textbf{(P\_6,1.91.2)} iha ubhayam prāpnoti . \\
\noindent \textbf{(P\_6,1.91.2)} prakārya gataḥ . \\
\noindent \textbf{(P\_6,1.91.2)} prahārya gataḥ . \\
\noindent \textbf{(P\_6,1.91.2)} allopasya avakāśaḥ . \\
\noindent \textbf{(P\_6,1.91.2)} cikīrṣitā , jihīrṣitā . \\
\noindent \textbf{(P\_6,1.91.2)} tukaḥ saḥ eva . \\
\noindent \textbf{(P\_6,1.91.2)} iha ubhayam prāpnoti . \\
\noindent \textbf{(P\_6,1.91.2)} pracikīrṣya gataḥ , prajihīrṣya gataḥ . \\
\noindent \textbf{(P\_6,1.91.2)} abhyāsaguṇādayaḥ ca tukaḥ bhavanti vipratiṣedhena . \\
\noindent \textbf{(P\_6,1.91.2)} ke punaḥ abhyāsaguṇādayaḥ . \\
\noindent \textbf{(P\_6,1.91.2)} hrasvatvāttvettvaguṇāḥ . \\
\noindent \textbf{(P\_6,1.91.2)} hrasvatvasya avakāśaḥ . \\
\noindent \textbf{(P\_6,1.91.2)} papatuḥ , papuḥ , tasthatuḥ , tasthuḥ . \\
\noindent \textbf{(P\_6,1.91.2)} tukaḥ saḥ eva . \\
\noindent \textbf{(P\_6,1.91.2)} iha ubhayam prāpnoti . \\
\noindent \textbf{(P\_6,1.91.2)} apacacchatuḥ , apacacchuḥ . \\
\noindent \textbf{(P\_6,1.91.2)} attvasya avakāśaḥ . \\
\noindent \textbf{(P\_6,1.91.2)} cakratuḥ , cakruḥ . \\
\noindent \textbf{(P\_6,1.91.2)} tukaḥ saḥ eva . \\
\noindent \textbf{(P\_6,1.91.2)} iha ubhayam prāpnoti . \\
\noindent \textbf{(P\_6,1.91.2)} apacacchṛdatuḥ , apacacchṛduḥ . \\
\noindent \textbf{(P\_6,1.91.2)} ittvasya avakāśaḥ . \\
\noindent \textbf{(P\_6,1.91.2)} pipakṣati , yiyakṣati . \\
\noindent \textbf{(P\_6,1.91.2)} tukaḥ saḥ eva . \\
\noindent \textbf{(P\_6,1.91.2)} iha ubhayam prāpnoti . \\
\noindent \textbf{(P\_6,1.91.2)} cicchādayiṣati , cicchardayiṣati . \\
\noindent \textbf{(P\_6,1.91.2)} guṇasya avakāśaḥ . \\
\noindent \textbf{(P\_6,1.91.2)} lolūyate , bebhidyate . \\
\noindent \textbf{(P\_6,1.91.2)} tukaḥ saḥ eva . \\
\noindent \textbf{(P\_6,1.91.2)} iha ubhayam prāpnoti . \\
\noindent \textbf{(P\_6,1.91.2)} cecchidyate , cocchupyate . \\
\noindent \textbf{(P\_6,1.91.2)} yaṇadeśāt āt guṇaḥ . \\
\noindent \textbf{(P\_6,1.91.2)} yaṇadeśāt āt guṇaḥ bhavati vipratiṣedhena . \\
\noindent \textbf{(P\_6,1.91.2)} yaṇadeśasya avakāśaḥ . \\
\noindent \textbf{(P\_6,1.91.2)} dadhi atra , madhu atra . \\
\noindent \textbf{(P\_6,1.91.2)} āt guṇasya avakāśaḥ . \\
\noindent \textbf{(P\_6,1.91.2)} khaṭvendraḥ , khaṭvodakam . \\
\noindent \textbf{(P\_6,1.91.2)} iha ubhayam prāpnoti . \\
\noindent \textbf{(P\_6,1.91.2)} vṛkṣaḥ atra, plakśaḥ atra . \\
\noindent \textbf{(P\_6,1.91.2)} irurguṇavṛddhividhayaḥ ca . \\
\noindent \textbf{(P\_6,1.91.2)} irurguṇavṛddhividhayaḥ ca yaṇadeśāt bhavanti vipratiṣedhena . \\
\noindent \textbf{(P\_6,1.91.2)} iruroḥ avakāśaḥ . \\
\noindent \textbf{(P\_6,1.91.2)} āstīrṇam , nipūrtāḥ piṇḍāḥ . \\
\noindent \textbf{(P\_6,1.91.2)} yaṇadeśasya avakāśaḥ . \\
\noindent \textbf{(P\_6,1.91.2)} cakratuḥ , cakruḥ . \\
\noindent \textbf{(P\_6,1.91.2)} iha ubhayam prāpnoti . \\
\noindent \textbf{(P\_6,1.91.2)} dūre hi adhvā jaguriḥ . \\
\noindent \textbf{(P\_6,1.91.2)} mitrāvaruṇau taturiḥ . \\
\noindent \textbf{(P\_6,1.91.2)} kirati , girati . \\
\noindent \textbf{(P\_6,1.91.2)} guṇavṛddhyoḥ avakāśaḥ . \\
\noindent \textbf{(P\_6,1.91.2)} cetā , gauḥ . \\
\noindent \textbf{(P\_6,1.91.2)} yaṇadeśasya saḥ eva . \\
\noindent \textbf{(P\_6,1.91.2)} iha ubhayam prāpnoti . \\
\noindent \textbf{(P\_6,1.91.2)} cayanam , cāyakaḥ , lavanam , lāvakaḥ . \\
\noindent \textbf{(P\_6,1.91.2)} bhalopadhātuprātipadikapratyayasamāsāntodāttanivṛttisvarāḥ ekādeśāt ca . \\
\noindent \textbf{(P\_6,1.91.2)} bhalopadhātuprātipadikapratyayasamāsāntodāttanivṛttisvarāḥ ekādeśāt ca yaṇadeśāt ca bhavanti vipratiṣedhena . \\
\noindent \textbf{(P\_6,1.91.2)} bhalopasya avakāśaḥ . \\
\noindent \textbf{(P\_6,1.91.2)} gārgyaḥ , vātsyaḥ . \\
\noindent \textbf{(P\_6,1.91.2)} ekādeśayaṇādeśayoḥ avakāśaḥ . \\
\noindent \textbf{(P\_6,1.91.2)} dadhīndraḥ , madhūdakam . \\
\noindent \textbf{(P\_6,1.91.2)} dadhi atra , madhu atra. iha ubhayam prāpnoti . \\
\noindent \textbf{(P\_6,1.91.2)} dākṣī , dākṣāyaṇaḥ , plākṣī , plākṣāyaṇaḥ . \\
\noindent \textbf{(P\_6,1.91.2)} aci bhalopaḥ ekādeśāt bhavati viprtiṣedhena . \\
\noindent \textbf{(P\_6,1.91.2)} aci bhalopasya avakāśaḥ . \\
\noindent \textbf{(P\_6,1.91.2)} dākṣī , dākṣāyaṇaḥ , plākṣī , plākṣāyaṇaḥ . \\
\noindent \textbf{(P\_6,1.91.2)} ekādeśasya avakāśaḥ . \\
\noindent \textbf{(P\_6,1.91.2)} daṇḍāgram , kṣupāgram . \\
\noindent \textbf{(P\_6,1.91.2)} iha ubhayam prāpnoti . \\
\noindent \textbf{(P\_6,1.91.2)} gāṅgeyaḥ gāṅgaḥ . \\
\noindent \textbf{(P\_6,1.91.2)} dhātusvarasya avakāśaḥ . \\
\noindent \textbf{(P\_6,1.91.2)} pacati , paṭhati . \\
\noindent \textbf{(P\_6,1.91.2)} ekādeśayaṇādeśayoḥ saḥ eva . \\
\noindent \textbf{(P\_6,1.91.2)} iha ubhayam prāpnoti . \\
\noindent \textbf{(P\_6,1.91.2)} śryartham , śrīṣā . \\
\noindent \textbf{(P\_6,1.91.2)} prātipadikasvarasya avakāśaḥ . \\
\noindent \textbf{(P\_6,1.91.2)} āmraḥ . \\
\noindent \textbf{(P\_6,1.91.2)} ekādeśayaṇādeśayoḥ saḥ eva . \\
\noindent \textbf{(P\_6,1.91.2)} iha ubhayam prāpnoti . \\
\noindent \textbf{(P\_6,1.91.2)} agnyudakam , vṛkṣārtham . \\
\noindent \textbf{(P\_6,1.91.2)} pratyayasvarasya avakāśaḥ . \\
\noindent \textbf{(P\_6,1.91.2)} cikīrṣuḥ , aupagavaḥ . \\
\noindent \textbf{(P\_6,1.91.2)} ekādeśayaṇādeśayoḥ saḥ eva . \\
\noindent \textbf{(P\_6,1.91.2)} iha ubhayam prāpnoti . \\
\noindent \textbf{(P\_6,1.91.2)} cikīrṣuartham , aupagavārtham . \\
\noindent \textbf{(P\_6,1.91.2)} samāsāntodāttasya avakāśaḥ . \\
\noindent \textbf{(P\_6,1.91.2)} rājapuruṣaḥ , brāhmaṇakambalaḥ . \\
\noindent \textbf{(P\_6,1.91.2)} ekādeśayaṇādeśayoḥ saḥ eva . \\
\noindent \textbf{(P\_6,1.91.2)} iha ubhayam prāpnoti . \\
\noindent \textbf{(P\_6,1.91.2)} rājavaidyartham , rājavaidī īhate . \\
\noindent \textbf{(P\_6,1.91.2)} udāttanivṛttisvarasya avakāśaḥ . \\
\noindent \textbf{(P\_6,1.91.2)} nadī , kumārī . \\
\noindent \textbf{(P\_6,1.91.2)} ekādeśayaṇādeśayoḥ saḥ eva . \\
\noindent \textbf{(P\_6,1.91.2)} iha ubhayam prāpnoti . \\
\noindent \textbf{(P\_6,1.91.2)} kumāryartham , kumārī īhate . \\
\noindent \textbf{(P\_6,1.91.2)} allopāllopau ca ārdhadhātuke . \\
\noindent \textbf{(P\_6,1.91.2)} allopāllopau ca ārdhadhātuke ekādeśāt bhavataḥ vipratiṣedhena . \\
\noindent \textbf{(P\_6,1.91.2)} allopasya avakāśaḥ . \\
\noindent \textbf{(P\_6,1.91.2)} cikīrṣitā , jihīrṣitā . \\
\noindent \textbf{(P\_6,1.91.2)} ekādeśasya avakāśaḥ . \\
\noindent \textbf{(P\_6,1.91.2)} pacanti , paṭhanti . \\
\noindent \textbf{(P\_6,1.91.2)} iha ubhayam prāpnoti . \\
\noindent \textbf{(P\_6,1.91.2)} cikīrṣakaḥ , jihīrṣakaḥ . \\
\noindent \textbf{(P\_6,1.91.2)} āllopasya avakāśaḥ . \\
\noindent \textbf{(P\_6,1.91.2)} papiḥ somam , dadiḥ gaḥ . \\
\noindent \textbf{(P\_6,1.91.2)} ekādeśasya avakāśaḥ . \\
\noindent \textbf{(P\_6,1.91.2)} yānti , vānti . \\
\noindent \textbf{(P\_6,1.91.2)} iha ubhayam prāpnoti . \\
\noindent \textbf{(P\_6,1.91.2)} yayatuḥ , yayuḥ . \\
\noindent \textbf{(P\_6,1.91.2)} iyaṅuvaṅguṇavṛddhiṭitkinmitpūrvapadavikārāḥ ca . \\
\noindent \textbf{(P\_6,1.91.2)} iyaṅuvaṅguṇavṛddhiṭitkinmitpūrvapadavikārāḥ ca ekādeśayaṇādeśābhyām bhavanti vipratiṣedhena . \\
\noindent \textbf{(P\_6,1.91.2)} iyaṅuvaṅoḥ avakāśaḥ . \\
\noindent \textbf{(P\_6,1.91.2)} śriyau , śriyaḥ , bhruvau , bhruvaḥ . \\
\noindent \textbf{(P\_6,1.91.2)} ekādeśayaṇādeśayoḥ saḥ eva . \\
\noindent \textbf{(P\_6,1.91.2)} iha ubhayam prāpnoti . \\
\noindent \textbf{(P\_6,1.91.2)} cikṣiyiva , cikṣiyima , luluvatuḥ , luluvuḥ , pupuvatuḥ , pupuvuḥ . \\
\noindent \textbf{(P\_6,1.91.2)} guṇavṛddhyoḥ avakāśaḥ . \\
\noindent \textbf{(P\_6,1.91.2)} cetā , gauḥ . \\
\noindent \textbf{(P\_6,1.91.2)} ekādeśayaṇādeśayoḥ saḥ eva . \\
\noindent \textbf{(P\_6,1.91.2)} iha ubhayam prāpnoti . \\
\noindent \textbf{(P\_6,1.91.2)} sādhucāyī , sucāyī , nagnambhāvukaḥ adhvaryuḥ , śayitā , śayitum . \\
\noindent \textbf{(P\_6,1.91.2)} ṭitaḥ avakāśaḥ . \\
\noindent \textbf{(P\_6,1.91.2)} agnīnām , indūnām . \\
\noindent \textbf{(P\_6,1.91.2)} ekādeśayaṇādeśayoḥ saḥ eva . \\
\noindent \textbf{(P\_6,1.91.2)} iha ubhayam prāpnoti . \\
\noindent \textbf{(P\_6,1.91.2)} vṛkṣāṇām , plakṣāṇām . \\
\noindent \textbf{(P\_6,1.91.2)} kitaḥ avakāśaḥ . \\
\noindent \textbf{(P\_6,1.91.2)} sādhudāyī , suṣṭhudāyī . \\
\noindent \textbf{(P\_6,1.91.2)} ekādeśayaṇādeśayoḥ saḥ eva . \\
\noindent \textbf{(P\_6,1.91.2)} iha ubhayam prāpnoti . \\
\noindent \textbf{(P\_6,1.91.2)} dāyakaḥ , dhāyakaḥ . \\
\noindent \textbf{(P\_6,1.91.2)} mitaḥ avakāśaḥ . \\
\noindent \textbf{(P\_6,1.91.2)} trapuṇī , jatunī . \\
\noindent \textbf{(P\_6,1.91.2)} ekādeśayaṇādeśayoḥ saḥ eva . \\
\noindent \textbf{(P\_6,1.91.2)} iha ubhayam prāpnoti . \\
\noindent \textbf{(P\_6,1.91.2)} asthīni , dadhīni , atisakhīni brāhmaṇakulāni . \\
\noindent \textbf{(P\_6,1.91.2)} pūrvapadavikārāṇām avakāśaḥ . \\
\noindent \textbf{(P\_6,1.91.2)} hotāpotārau . \\
\noindent \textbf{(P\_6,1.91.2)} ekādeśayaṇādeśayoḥ saḥ eva . \\
\noindent \textbf{(P\_6,1.91.2)} iha ubhayam prāpnoti . \\
\noindent \textbf{(P\_6,1.91.2)} neṣṭodgātārau āgnendram . \\
\noindent \textbf{(P\_6,1.91.2)} uttarapadavikārāḥ ca iti vaktavyam . \\
\noindent \textbf{(P\_6,1.91.2)} uttarapadavikārāṇām avakāśaḥ . \\
\noindent \textbf{(P\_6,1.91.2)} samīpam , durīpam . \\
\noindent \textbf{(P\_6,1.91.2)} ekādeśayaṇādeśayoḥ saḥ eva . \\
\noindent \textbf{(P\_6,1.91.2)} iha ubhayam prāpnoti . \\
\noindent \textbf{(P\_6,1.91.2)} prepam , parepam . \\
\noindent \textbf{(P\_6,1.93)} otaḥ tiṅi pratiṣedhaḥ . \\
\noindent \textbf{(P\_6,1.93)} otaḥ tiṅi pratiṣedhaḥ vaktavyaḥ . \\
\noindent \textbf{(P\_6,1.93)} acinavam , asunavam . \\
\noindent \textbf{(P\_6,1.93)} saḥ tarhi pratiṣedhaḥ vaktavyaḥ . \\
\noindent \textbf{(P\_6,1.93)} na vaktavyaḥ . \\
\noindent \textbf{(P\_6,1.93)} gograhaṇam kariṣyate . \\
\noindent \textbf{(P\_6,1.93)} ā gotaḥ iti vaktavyam . \\
\noindent \textbf{(P\_6,1.93)} gograhaṇe dyoḥ upasaṅkhyanam . \\
\noindent \textbf{(P\_6,1.93)} gograhaṇe dyoḥ upasaṅkhyanam kartavyam . \\
\noindent \textbf{(P\_6,1.93)} dyam gaccha . \\
\noindent \textbf{(P\_6,1.93)} samāsāt ca pratiṣedhaḥ . \\
\noindent \textbf{(P\_6,1.93)} samāsāt ca pratiṣedhaḥ vaktavyaḥ . \\
\noindent \textbf{(P\_6,1.93)} citragum paśya . \\
\noindent \textbf{(P\_6,1.93)} śabalagum paśya . \\
\noindent \textbf{(P\_6,1.93)} nanu ca ā otaḥ iti ucyamāne api samāsāt pratiṣedhaḥ vaktavyaḥ . \\
\noindent \textbf{(P\_6,1.93)} na vaktavyaḥ . \\
\noindent \textbf{(P\_6,1.93)} hrasvatve kṛte na bhaviṣyati . \\
\noindent \textbf{(P\_6,1.93)} idam iha sampradhāryam . \\
\noindent \textbf{(P\_6,1.93)} ātvam kriyatām hrasvatvam iti kim atra kartavyam . \\
\noindent \textbf{(P\_6,1.93)} paratvāt ātvam . \\
\noindent \textbf{(P\_6,1.93)} na vā bahiraṅgalakṣaṇatvāt . \\
\noindent \textbf{(P\_6,1.93)} na vā vaktavyaḥ . \\
\noindent \textbf{(P\_6,1.93)} kim kāraṇam . \\
\noindent \textbf{(P\_6,1.93)} bahiraṅgalakṣaṇatvāt . \\
\noindent \textbf{(P\_6,1.93)} bahiraṅgalakṣaṇam ātvam . \\
\noindent \textbf{(P\_6,1.93)} antaraṅgam hrasvatvam . \\
\noindent \textbf{(P\_6,1.93)} asiddham bahiraṅgam antaraṅge . \\
\noindent \textbf{(P\_6,1.93)} nanu ca ā gotaḥ iti ucyamāne api samāsāt pratiṣedhaḥ na vaktavyaḥ . \\
\noindent \textbf{(P\_6,1.93)} katham . \\
\noindent \textbf{(P\_6,1.93)} hrasvatve kṛte na bhaviṣyati . \\
\noindent \textbf{(P\_6,1.93)} sthānivadbhāvāt prāpnoti . \\
\noindent \textbf{(P\_6,1.93)} nanu ca ā otaḥ iti ucyamāne api sthānivadbhāvāt prāpnoti . \\
\noindent \textbf{(P\_6,1.93)} na iti āha . \\
\noindent \textbf{(P\_6,1.93)} analvidhau sthānivadbhāvaḥ . \\
\noindent \textbf{(P\_6,1.93)} ā gotaḥ iti ucyamāne api na doṣaḥ . \\
\noindent \textbf{(P\_6,1.93)} pratiṣidhyate atra sthānivadbhāvaḥ . \\
\noindent \textbf{(P\_6,1.93)} goḥ pūrvaṇitvātvasvareṣu sthānivat na bhavati iti . \\
\noindent \textbf{(P\_6,1.93)} saḥ eva tarhi doṣaḥ . \\
\noindent \textbf{(P\_6,1.93)} gograhaṇe dyoḥ upasaṅkhyanam iti . \\
\noindent \textbf{(P\_6,1.93)} sūtram ca bhidyate . \\
\noindent \textbf{(P\_6,1.93)} yathānyāsam eva astu . \\
\noindent \textbf{(P\_6,1.93)} nanu ca uktam otaḥ tiṅi pratiṣedhaḥ iti . \\
\noindent \textbf{(P\_6,1.93)} subadhikārāt siddham . \\
\noindent \textbf{(P\_6,1.93)} supi iti vartate . \\
\noindent \textbf{(P\_6,1.93)} kva prakṛtam . \\
\noindent \textbf{(P\_6,1.93)} vā supi āpiśaleḥ iti . \\
\noindent \textbf{(P\_6,1.93)} yadi anuvartate iha api vibhāṣā prāpnoti . \\
\noindent \textbf{(P\_6,1.93)} subgrahaṇam anuvartate . \\
\noindent \textbf{(P\_6,1.93)} vāgrahaṇam nivṛttam . \\
\noindent \textbf{(P\_6,1.93)} katham punaḥ ekayoganirdiṣṭayoḥ ekadeśaḥ anuvartate ekadeśaḥ na . \\
\noindent \textbf{(P\_6,1.93)} ekayoge ca ekadeśānuvṛttiḥ anyatra api . \\
\noindent \textbf{(P\_6,1.93)} ekayognirdiṣṭānām api ekadeśānuvṛttiḥ bhavati . \\
\noindent \textbf{(P\_6,1.93)} anyatra api . \\
\noindent \textbf{(P\_6,1.93)} na avaśyam iha eva . \\
\noindent \textbf{(P\_6,1.93)} kva anyatra . \\
\noindent \textbf{(P\_6,1.93)} alugadhikāraḥ prāk ānaṅaḥ . \\
\noindent \textbf{(P\_6,1.93)} uttarapadāhikāraḥ prāk aṅgādhikāraḥ . \\
\noindent \textbf{(P\_6,1.93)} evam api ami upasaṅkhyānam vṛddhibalīyastvāt . \\
\noindent \textbf{(P\_6,1.93)} ami upasaṅkhyānam kartavyam . \\
\noindent \textbf{(P\_6,1.93)} gām paśya . \\
\noindent \textbf{(P\_6,1.93)} kim punaḥ kāraṇam na sidhyati . \\
\noindent \textbf{(P\_6,1.93)} vṛddhibalīyastvāt . \\
\noindent \textbf{(P\_6,1.93)} paratvāt vṛddhiḥ prāpnoti . \\
\noindent \textbf{(P\_6,1.93)} na vā anavakāśatvāt . \\
\noindent \textbf{(P\_6,1.93)} na vā vaktavyam . \\
\noindent \textbf{(P\_6,1.93)} kim kāraṇam . \\
\noindent \textbf{(P\_6,1.93)} anavakāśatvāt . \\
\noindent \textbf{(P\_6,1.93)} anavakāśam ātvam vṛddhim bādhiṣyate . \\
\noindent \textbf{(P\_6,1.93)} sāvakāśam ātvam . \\
\noindent \textbf{(P\_6,1.93)} kaḥ avakāśaḥ . \\
\noindent \textbf{(P\_6,1.93)} dyām gaccha . \\
\noindent \textbf{(P\_6,1.93)} dyoḥ ca sarvanāmasthāne vṛddhividhiḥ . \\
\noindent \textbf{(P\_6,1.93)} dyoḥ ca sarvanāmasthāne vṛddhiḥ vidheyā . \\
\noindent \textbf{(P\_6,1.93)} kim prayojanam . \\
\noindent \textbf{(P\_6,1.93)} yat dyāvaḥ indra iti darśanāt . \\
\noindent \textbf{(P\_6,1.93)} yat dyavaḥ indra te śataṃ śatam bhumīḥ uta syuḥ . \\
\noindent \textbf{(P\_6,1.93)} yāvatā ca idānīm dyoḥ api sarvanāmasthāne vṛddhiḥ ucyate anavakāśam ātvam vṛddhim bādhiṣyate . \\
\noindent \textbf{(P\_6,1.94)} pararūpaprakaraṇe tunvoḥ vi nipāte upasaṅkhyānam . \\
\noindent \textbf{(P\_6,1.94)} pararūpaprakaraṇe tu , nu, iti etayoḥ vakārādau nipāte upasaṅkhyānam kartavyam . \\
\noindent \textbf{(P\_6,1.94)} tu vai tvai , nu vai nvai . \\
\noindent \textbf{(P\_6,1.94)} vakārādau iti kimartham . \\
\noindent \textbf{(P\_6,1.94)} tvāvat , nvāvat . \\
\noindent \textbf{(P\_6,1.94)} nipāte iti kimartham . \\
\noindent \textbf{(P\_6,1.94)} tu vāni , nu vāni . \\
\noindent \textbf{(P\_6,1.94)} na vā nipātaikatvāt . \\
\noindent \textbf{(P\_6,1.94)} na vā kartavyam . \\
\noindent \textbf{(P\_6,1.94)} kim kāraṇam . \\
\noindent \textbf{(P\_6,1.94)} nipātaikatvāt . \\
\noindent \textbf{(P\_6,1.94)} ekaḥ eva ayam nipātaḥ . \\
\noindent \textbf{(P\_6,1.94)} tvai , nvai . \\
\noindent \textbf{(P\_6,1.94)} eve ca aniyoge . \\
\noindent \textbf{(P\_6,1.94)} eve ca aniyoge pararūpam vaktavyam . \\
\noindent \textbf{(P\_6,1.94)} iha eva , iheva . \\
\noindent \textbf{(P\_6,1.94)} adyeva . \\
\noindent \textbf{(P\_6,1.94)} aniyoge iti kimartham . \\
\noindent \textbf{(P\_6,1.94)} ihaiva bhava ma sma gāḥ . \\
\noindent \textbf{(P\_6,1.94)} atraiva tvam iha vayam suśevāḥ . \\
\noindent \textbf{(P\_6,1.94)} śakandhvādiṣu ca . \\
\noindent \textbf{(P\_6,1.94)} śakandhvādiṣu ca pararūpam vaktavyam . \\
\noindent \textbf{(P\_6,1.94)} śaka-andhuḥ śakandhuḥ , kula-aṭā , kulaṭā , sīma-antaḥ sīmantaḥ . \\
\noindent \textbf{(P\_6,1.94)} keśeṣu iti vaktavyam . \\
\noindent \textbf{(P\_6,1.94)} yaḥ hi sīmnaḥ antaḥ sīmāntaḥ saḥ bhavati . \\
\noindent \textbf{(P\_6,1.94)} otvoṣṭhayoḥ samāse vā . otvoṣṭhayoḥ samāse vā pararūpam vaktavyam . \\
\noindent \textbf{(P\_6,1.94)} sthūlautuḥ , sthūlotuḥ . \\
\noindent \textbf{(P\_6,1.94)} bimbauṣthī , bimboṣṭhī . \\
\noindent \textbf{(P\_6,1.94)} emanādiṣu chandasi . \\
\noindent \textbf{(P\_6,1.94)} emanādiṣu chandasi pararūpam vaktavyam . \\
\noindent \textbf{(P\_6,1.94)} apam tveman sādayāmi apam todayan sādayāmi iti . \\
\noindent \textbf{(P\_6,1.95)} kimarthaḥ cakāraḥ . \\
\noindent \textbf{(P\_6,1.95)} eṅi iti anukṛṣyate . \\
\noindent \textbf{(P\_6,1.95)} kim prayojanam . \\
\noindent \textbf{(P\_6,1.95)} iha mā bhūt . \\
\noindent \textbf{(P\_6,1.95)} adya ā ṛśyāt , adyārśyāt , kadārśyāt . \\
\noindent \textbf{(P\_6,1.95)} na etat asti prayojanam . \\
\noindent \textbf{(P\_6,1.95)} adyarśyāt iti eva bhavitavyam . \\
\noindent \textbf{(P\_6,1.95)} evam hi saunāgāḥ paṭhanti . \\
\noindent \textbf{(P\_6,1.95)} caḥ anarthakaḥ anadhikārāt eṅaḥ . \\
\noindent \textbf{(P\_6,1.95)} usyomāṅkṣu āṭaḥ pratiṣedhaḥ . \\
\noindent \textbf{(P\_6,1.95)} usi pararūpe omāṅoḥ ca āṭaḥ pratiṣedhaḥ vaktavyaḥ . \\
\noindent \textbf{(P\_6,1.95)} ausrīyat , auḍhīyat , auṅkārīyat . \\
\noindent \textbf{(P\_6,1.95)} saḥ tarhi pratiṣedhaḥ vaktavyaḥ . \\
\noindent \textbf{(P\_6,1.95)} na vaktavyaḥ . \\
\noindent \textbf{(P\_6,1.95)} uktam ātaḥ ca iti atra cakārasya prayojanam . \\
\noindent \textbf{(P\_6,1.95)} vṛddhiḥ eva yathā syāt . \\
\noindent \textbf{(P\_6,1.95)} yat anyat prāpnoti tat mā bhūt iti . \\
\noindent \textbf{(P\_6,1.96)} apadāntāt iti kimartham . \\
\noindent \textbf{(P\_6,1.96)} kā , usrā , kosrā . \\
\noindent \textbf{(P\_6,1.96)} apadāntāt iti śakyam akartum . \\
\noindent \textbf{(P\_6,1.96)} kasmāt na bhavati kā , usrā , kosrā . \\
\noindent \textbf{(P\_6,1.96)} arthavadgrahaṇe na anarthakasya iti . \\
\noindent \textbf{(P\_6,1.96)} na eṣā paribhāṣā iha śakyā vijñātum . \\
\noindent \textbf{(P\_6,1.96)} iha hi doṣaḥ syāt . \\
\noindent \textbf{(P\_6,1.96)} bhindyā-us , bhindyuḥ , chindyā-us, chindyuḥ . \\
\noindent \textbf{(P\_6,1.96)} evam tarhi lakṣaṇapratipadoktayoḥ pratipadoktasya eva iti evam na bhaviṣyati . \\
\noindent \textbf{(P\_6,1.96)} uttarārtham tarhi apadāntagrahaṇam kartavyam . \\
\noindent \textbf{(P\_6,1.96)} ataḥ guṇe apadāntāt yathā syāt . \\
\noindent \textbf{(P\_6,1.96)} iha mā bhūt . \\
\noindent \textbf{(P\_6,1.96)} daṇḍāgram , kṣupāgram iti . \\
\noindent \textbf{(P\_6,1.98)} itau anekājgrahaṇam śradartham . \\
\noindent \textbf{(P\_6,1.98)} itau anekājgrahaṇam kartavyam . \\
\noindent \textbf{(P\_6,1.98)} kim prayojanam . \\
\noindent \textbf{(P\_6,1.98)} śradartham . \\
\noindent \textbf{(P\_6,1.98)} śrat iti . \\
\noindent \textbf{(P\_6,1.99)} nityam āmreḍite ḍāci . \\
\noindent \textbf{(P\_6,1.99)} nityam āmreḍite ḍāci pararūpam kartavyam . \\
\noindent \textbf{(P\_6,1.99)} paṭapaṭāyati . \\
\noindent \textbf{(P\_6,1.99)} akārantāt anukaraṇāt vā . \\
\noindent \textbf{(P\_6,1.99)} atha vā akārāntam etad udāharaṇam . \\
\noindent \textbf{(P\_6,1.99)} bhavet siddham yadā akārāntam . \\
\noindent \textbf{(P\_6,1.99)} yadā tu khalu acchabdāntam tadā na sidhyati . \\
\noindent \textbf{(P\_6,1.99)} vicitrāḥ taddhitavṛttayaḥ . \\
\noindent \textbf{(P\_6,1.99)} na ataḥ taddhitaḥ utpadyate . \\
\noindent \textbf{(P\_6,1.101)} savarṇadīrghatve ṛti ṛvāvacanam . \\
\noindent \textbf{(P\_6,1.101)} savarṇadīrghatve ṛti ṛ vā bhavati iti vaktavyam . \\
\noindent \textbf{(P\_6,1.101)} hotṛ ṛkāraḥ , hotṝkāraḥ . \\
\noindent \textbf{(P\_6,1.101)} lṛti lṛvāvacanam . \\
\noindent \textbf{(P\_6,1.101)} lṛti ḷ vā bhavati iti vaktavyam . \\
\noindent \textbf{(P\_6,1.101)} hotṛ ḷkāraḥ , hotlḹkāraḥ . \\
\noindent \textbf{(P\_6,1.102.1)} prathamyoḥ iti ucyate . \\
\noindent \textbf{(P\_6,1.102.1)} kayoḥ iha prathamyoḥ grahaṇam . \\
\noindent \textbf{(P\_6,1.102.1)} kim vibhaktyoḥ āhosvit pratyayayoḥ . \\
\noindent \textbf{(P\_6,1.102.1)} vibhaktyoḥ iti āha . \\
\noindent \textbf{(P\_6,1.102.1)} katham jñāyate . \\
\noindent \textbf{(P\_6,1.102.1)} aci iti vartate na ca ajādau prathamau pratyayau staḥ . \\
\noindent \textbf{(P\_6,1.102.1)} nanu ca evam vijñāyate . \\
\noindent \textbf{(P\_6,1.102.1)} ajādī yau prathamau ajādīnām vā yau prathamau iti . \\
\noindent \textbf{(P\_6,1.102.1)} yat tarhi tasmāt śasaḥ naḥ puṃsi iti anukrāntam pūrvasavarṇam pratinirdiśati tat jñāpayati ācāryaḥ vibhaktyoḥ grahaṇam iti . \\
\noindent \textbf{(P\_6,1.102.1)} atha vā supi iti vartate . \\
\noindent \textbf{(P\_6,1.102.1)} atham kimartham pūrvasavarṇadīrghaḥ ami pūrvatvam ca ucyate na prathamyoḥ pūrvasavarṇaḥ iti eva siddham . \\
\noindent \textbf{(P\_6,1.102.1)} na sidhyati . \\
\noindent \textbf{(P\_6,1.102.1)} prathamyoḥ pūrvasavarṇaḥ iti ucyamāne ami api dīrghaḥ prāpnoti . \\
\noindent \textbf{(P\_6,1.102.1)} vṛkṣam , plakṣam . \\
\noindent \textbf{(P\_6,1.102.1)} na eṣaḥ doṣaḥ . \\
\noindent \textbf{(P\_6,1.102.1)} yat pūrvasmin yoge dīrghagrahaṇam tat uttaratra nivṛttam . \\
\noindent \textbf{(P\_6,1.102.1)} evam api idam iha pūrvasavarṇagrahaṇam kriyate . \\
\noindent \textbf{(P\_6,1.102.1)} tena ami api pūrvasavarṇaḥ prasajyeta . \\
\noindent \textbf{(P\_6,1.102.1)} vṛkṣam , plakṣam . \\
\noindent \textbf{(P\_6,1.102.1)} dvimātraḥ prāpnoti . \\
\noindent \textbf{(P\_6,1.102.1)} na eṣaḥ doṣaḥ . \\
\noindent \textbf{(P\_6,1.102.1)} savarṇagrahaṇam na kariṣyate . \\
\noindent \textbf{(P\_6,1.102.1)} yadi savarṇagrahaṇam na kriyate kutaḥ vyavasthā . \\
\noindent \textbf{(P\_6,1.102.1)} āntaryataḥ . \\
\noindent \textbf{(P\_6,1.102.1)} yadi evam agnī vāyū trimātraḥ prāpnoti vṛkṣam , plakṣam dvimātraḥ . \\
\noindent \textbf{(P\_6,1.102.1)} tasmāt savarṇagrahaṇam kartavyam . \\
\noindent \textbf{(P\_6,1.102.1)} tasmin ca kriyamāṇe dīrghagrahaṇam anuvartate . \\
\noindent \textbf{(P\_6,1.102.1)} tasmin anuvartamāne ami pūrvaḥ iti api vaktavyam . \\
\noindent \textbf{(P\_6,1.102.1)} atha kimartham pṛthak ucyate na iha eka eva ucyeta . \\
\noindent \textbf{(P\_6,1.102.1)} prathamayoḥ pūrvasavarṇaḥ ami ca iti . \\
\noindent \textbf{(P\_6,1.102.1)} yadi prathamayoḥ pūrvasavarṇadīrghaḥ ami ca iti ucyate tana ami api dīrghaḥ prasajyeta . \\
\noindent \textbf{(P\_6,1.102.1)} vṛkṣam , plakṣam . \\
\noindent \textbf{(P\_6,1.102.1)} na eṣaḥ doṣaḥ . \\
\noindent \textbf{(P\_6,1.102.1)} dīrghagrahaṇam nivartayiṣyate . \\
\noindent \textbf{(P\_6,1.102.1)} evam api pūrvasavarṇaḥ prasajyeta . \\
\noindent \textbf{(P\_6,1.102.1)} savarṇagrahaṇam na kariṣyate . \\
\noindent \textbf{(P\_6,1.102.1)} yadi savarṇagrahaṇam na kriyate pūrvasmin yoge vipratiṣiddham . \\
\noindent \textbf{(P\_6,1.102.1)} yadi pūrvaḥ na dīrghaḥ atha dīrghaḥ na pūrvaḥ . \\
\noindent \textbf{(P\_6,1.102.1)} pūrvaḥ dīrghaḥ ca iti vipratiṣiddham . \\
\noindent \textbf{(P\_6,1.102.1)} tasmāt ubhayam ārabdhavyam pṛthak ca kartavyam . \\
\noindent \textbf{(P\_6,1.102.2)} prathamayoḥ iti yogavibhāgaḥ savarṇadīrghārthaḥ . \\
\noindent \textbf{(P\_6,1.102.2)} prathamayoḥ iti yogavibhāgaḥ kartavyaḥ . \\
\noindent \textbf{(P\_6,1.102.2)} prathamayoḥ ekaḥ savarṇadīrghaḥ bhavati . \\
\noindent \textbf{(P\_6,1.102.2)} tataḥ pūrvasavarṇaḥ . \\
\noindent \textbf{(P\_6,1.102.2)} pūrvasavarṇadīrghaḥ bhavati ekaḥ prathamayoḥ iti . \\
\noindent \textbf{(P\_6,1.102.2)} kimarthaḥ yogavibhāgaḥ . \\
\noindent \textbf{(P\_6,1.102.2)} savarṇadīrghatvam yathā syāt . \\
\noindent \textbf{(P\_6,1.102.2)} ekayoge hi jaśśahoḥ pararūpaprasaṅgaḥ . \\
\noindent \textbf{(P\_6,1.102.2)} ekayoge hi sati jaśśahoḥ pararūpam prasajyeta . \\
\noindent \textbf{(P\_6,1.102.2)} vṛkṣāḥ , plakṣāḥ , vṛkṣān , plakṣān . \\
\noindent \textbf{(P\_6,1.102.2)} nanu ca pūrvasavarṇadīrghatvam pararūpam bādhiṣyate . \\
\noindent \textbf{(P\_6,1.102.2)} na utsahate bādhitum . \\
\noindent \textbf{(P\_6,1.102.2)} kim kāraṇam . \\
\noindent \textbf{(P\_6,1.102.2)} ādguṇayaṇādeśayoḥ apavādāḥ vṛddhisavarṇadīrghapūrvasavarṇādeśāḥ teṣām pararūpam svarasandhiṣu . \\
\noindent \textbf{(P\_6,1.102.2)} ādguṇayaṇādeśau utsargau . \\
\noindent \textbf{(P\_6,1.102.2)} tayoḥ apavādāḥ vṛddhisavarṇadīrghapūrvasavarṇādeśāḥ teṣām sarveṣām pararūpam apavādaḥ . \\
\noindent \textbf{(P\_6,1.102.2)} tat sarvabādhakam . \\
\noindent \textbf{(P\_6,1.102.2)} sarvabādhakatvāt prāpnoti . \\
\noindent \textbf{(P\_6,1.102.2)} atha kriyamāṇe api yogavibhāge yāvatā pararūpam apavādaḥ kasmāt eva na bādhate . \\
\noindent \textbf{(P\_6,1.102.2)} yogavibhāgaḥ anyaśāstranivṛttiyarthaḥ . \\
\noindent \textbf{(P\_6,1.102.2)} yogavibhāgaḥ anyaśāstranivṛttiyarthaḥ vijñāyate . \\
\noindent \textbf{(P\_6,1.102.2)} yogavibhāgaḥ anyaśāstranivṛttiyarthaḥ cet ami atiprasaṅgaḥ . yogavibhāgaḥ anyaśāstranivṛttiyarthaḥ cet ami atiprasaṅgaḥ bhavati . \\
\noindent \textbf{(P\_6,1.102.2)} vṛkṣam , plakṣam . \\
\noindent \textbf{(P\_6,1.102.2)} yathā eva hi yogavibhāgaḥ pararūpam bādhate evam ami pūrvatvam api bādheta . \\
\noindent \textbf{(P\_6,1.102.2)} nakārābhāvaḥ ca tasmāt iti anantaranirdeśāt . \\
\noindent \textbf{(P\_6,1.102.2)} natvasya ca abhāvaḥ . \\
\noindent \textbf{(P\_6,1.102.2)} vṛkṣān , plakṣān . \\
\noindent \textbf{(P\_6,1.102.2)} kim kāraṇam . \\
\noindent \textbf{(P\_6,1.102.2)} ca tasmāt iti anantaranirdeśāt . \\
\noindent \textbf{(P\_6,1.102.2)} tasmāt iti anena anantaraḥ yogaḥ pratinirdiśyate . \\
\noindent \textbf{(P\_6,1.102.2)} kim punaḥ kāraṇam tasmāt iti anena anantaraḥ yogaḥ pratinirdiśyate . \\
\noindent \textbf{(P\_6,1.102.2)} iha mā bhūt . \\
\noindent \textbf{(P\_6,1.102.2)} etān gāḥ paśya [R: etān gāḥ caturaḥ balivardān paśya] iti . \\
\noindent \textbf{(P\_6,1.102.2)} astu tarhi ekayogaḥ eva . \\
\noindent \textbf{(P\_6,1.102.2)} nanu ca uktam ekayoge hi jaśśahoḥ pararūpaprasaṅgaḥ iti . \\
\noindent \textbf{(P\_6,1.102.2)} na eṣaḥ doṣaḥ . \\
\noindent \textbf{(P\_6,1.102.2)} ijgrahaṇam tu jñāpakam pararūpābhāvasya . \\
\noindent \textbf{(P\_6,1.102.2)} yat ayam na āt ici iti ijgrahaṇam karoti tat jñāpayati ācāryaḥ na jaśśasoḥ pararūpam bhavati iti . \\
\noindent \textbf{(P\_6,1.102.2)} katham kṛtvā jñāpakam . \\
\noindent \textbf{(P\_6,1.102.2)} ijgrahaṇasya idam prayojanam . \\
\noindent \textbf{(P\_6,1.102.2)} iha mā bhūt . \\
\noindent \textbf{(P\_6,1.102.2)} vṛkṣāḥ , plakṣāḥ , vṛkṣān , plakṣān . \\
\noindent \textbf{(P\_6,1.102.2)} yadi ca jaśśasoḥ pararūpam syāt ijgrahaṇam anarthakam syāt . \\
\noindent \textbf{(P\_6,1.102.2)} paśyati tu ācāryaḥ na jaśśasoḥ pararūpam bhavati iti . \\
\noindent \textbf{(P\_6,1.102.2)} tataḥ ijgrahaṇam karoti . \\
\noindent \textbf{(P\_6,1.102.2)} na etat asti jñāpakam . \\
\noindent \textbf{(P\_6,1.102.2)} uttarārtham etat syāt . \\
\noindent \textbf{(P\_6,1.102.2)} dīrghāt jasi ca ici ca iti . \\
\noindent \textbf{(P\_6,1.102.2)} yadi uttarārtham etat syāt atra eva ayam ijdgrahaṇam kurvīta . \\
\noindent \textbf{(P\_6,1.102.2)} iha api tarhi kriyamāṇam yadi uttarārtham na jñāpakam bhavati . \\
\noindent \textbf{(P\_6,1.102.2)} evam tarhi yadi uttarārtham etat syāt na eva ayam ijdgrahaṇam kurvīta na api jasgrahaṇam . \\
\noindent \textbf{(P\_6,1.102.2)} etāvat ayam brūyāt . \\
\noindent \textbf{(P\_6,1.102.2)} dīrghāt śasi pūrvasavarṇaḥ bhavati iti . \\
\noindent \textbf{(P\_6,1.102.2)} tat niyamārtham bhaviṣyati . \\
\noindent \textbf{(P\_6,1.102.2)} dīrghāt śasi eva na anyatra iti . \\
\noindent \textbf{(P\_6,1.102.2)} saḥ ayam evam laghīyasā nyāsena siddhe yat ijgrahaṇam karoti tat jñāpayati ācāryaḥ na jaśśasoḥ pararūpam bhavati iti . \\
\noindent \textbf{(P\_6,1.102.2)} atha vā punaḥ astu yogavibhāgaḥ . \\
\noindent \textbf{(P\_6,1.102.2)} nanu ca uktam yogavibhāgaḥ anyaśāstranivṛttiyarthaḥ cet ami atiprasaṅgaḥ iti . \\
\noindent \textbf{(P\_6,1.102.2)} na eṣaḥ doṣaḥ . \\
\noindent \textbf{(P\_6,1.102.2)} ami api yogavibhāgaḥ kariṣyate . \\
\noindent \textbf{(P\_6,1.102.2)} ami . \\
\noindent \textbf{(P\_6,1.102.2)} ami yat uktam tat na bhavati iti . \\
\noindent \textbf{(P\_6,1.102.2)} tataḥ pūrvaḥ . \\
\noindent \textbf{(P\_6,1.102.2)} pūrvaḥ ca bhavati ami iti . \\
\noindent \textbf{(P\_6,1.102.2)} yat api ucyate nakārābhāvaḥ ca tasmāt iti anantaranirdeśāt iti . \\
\noindent \textbf{(P\_6,1.102.2)} kaḥ punaḥ arhati tasmāt iti anena anantaram yogam pratinirdeṣṭum . \\
\noindent \textbf{(P\_6,1.102.2)} evam kila pratinirdiśyate . \\
\noindent \textbf{(P\_6,1.102.2)} tasmāt pūrvasavarṇadīrghāt iti . \\
\noindent \textbf{(P\_6,1.102.2)} tat ca na . \\
\noindent \textbf{(P\_6,1.102.2)} evam pratinirdiśyate . \\
\noindent \textbf{(P\_6,1.102.2)} tasmāt akaḥ savarṇāt iti . \\
\noindent \textbf{(P\_6,1.102.2)} atha vā tasmāt prathamyoḥ dīrghāt iti . \\
\noindent \textbf{(P\_6,1.102.2)} atha vā punaḥ astu ami ekayogaḥ . \\
\noindent \textbf{(P\_6,1.102.2)} nanu ca uktam yogavibhāgaḥ anyaśāstranivṛttiyarthaḥ cet ami atiprasaṅgaḥ iti . \\
\noindent \textbf{(P\_6,1.102.2)} na eṣaḥ doṣaḥ . \\
\noindent \textbf{(P\_6,1.102.2)} madhye apavādāḥ pūrvān vidhīn bādhante iti evam ayam yogavibhāgaḥ pararūpam bādhiṣyate ami pūrvatvam na bādhiṣyate . \\
\noindent \textbf{(P\_6,1.102.2)} yadi etat asti madhye apavādāḥ purastāt apavādāḥ iti na arthaḥ ekena api yogavibhāgena . \\
\noindent \textbf{(P\_6,1.102.2)} purastāt apavādāḥ anantarān vidhīn bādhante iti evam pararūpam savarṇadīrghatvam bādhiṣyate prathamayoḥ pūrvasavarṇadīrghatvam na bādhiṣyate . \\
\noindent \textbf{(P\_6,1.102.2)} atha vā saptame yogavibhāgaḥ kariṣyate . \\
\noindent \textbf{(P\_6,1.102.2)} idam asti ataḥ dīrghaḥ yañi supi ca iti . \\
\noindent \textbf{(P\_6,1.102.2)} tataḥ vakṣyāmi bahuvacane . \\
\noindent \textbf{(P\_6,1.102.2)} bahuvacane ca ataḥ dīrghaḥ bhavati . \\
\noindent \textbf{(P\_6,1.102.2)} ekāraḥ ca bhavati bahuvacane jhali iti . \\
\noindent \textbf{(P\_6,1.102.2)} iha api tarhi prāpnoti . \\
\noindent \textbf{(P\_6,1.102.2)} vṛṣāṇām , plakṣāṇām . \\
\noindent \textbf{(P\_6,1.102.2)} tatra kaḥ doṣaḥ . \\
\noindent \textbf{(P\_6,1.102.2)} dīrghatve kṛte hrasvāśrayaḥ nuṭ na prāpnoti . \\
\noindent \textbf{(P\_6,1.102.2)} idam iha sampradhāryam . \\
\noindent \textbf{(P\_6,1.102.2)} dīrghatvam kriyatām nuṭ iti kim atra kartavyam . \\
\noindent \textbf{(P\_6,1.102.2)} paratvāt dīrghatvam . \\
\noindent \textbf{(P\_6,1.102.2)} nityam khalu api dīrghatvam . \\
\noindent \textbf{(P\_6,1.102.2)} kṛte api nuṭi prāpnoti akṛte api . \\
\noindent \textbf{(P\_6,1.102.2)} nityatvāt paratvāt ca dīrghatve kṛte hrasvāśrayaḥ nuṭ na prāpnoti . \\
\noindent \textbf{(P\_6,1.102.2)} evam tarhi ādgrahaṇam iha api prakṛtam anuvartate . \\
\noindent \textbf{(P\_6,1.102.2)} kva prakṛtam . \\
\noindent \textbf{(P\_6,1.102.2)} āt jaseḥ asuk iti . \\
\noindent \textbf{(P\_6,1.102.2)} tena kṛte api dīrghatve nuṭ bhaviṣyati . \\
\noindent \textbf{(P\_6,1.102.2)} iha api tarhi prāpnoti . \\
\noindent \textbf{(P\_6,1.102.2)} kīlālapām , śubhaṃyām . \\
\noindent \textbf{(P\_6,1.102.2)} ātaḥ lopaḥ atra bādhakaḥ bhaviṣyati . \\
\noindent \textbf{(P\_6,1.102.2)} idam iha sampradhāryam . \\
\noindent \textbf{(P\_6,1.102.2)} lopaḥ kriyatām nuṭ iti kim atra kartavyam . \\
\noindent \textbf{(P\_6,1.102.2)} paratvāt nuṭ . \\
\noindent \textbf{(P\_6,1.102.2)} evam tarhi hrasvanadyāpaḥ nuṭ iti atra ātaḥ dhātoḥ iti ātaḥ lopaḥ sambandham anuvartiṣyate . \\
\noindent \textbf{(P\_6,1.102.2)} iha api tarhi prāpnoti . \\
\noindent \textbf{(P\_6,1.102.2)} kīlālapānām brāhmaṇakulānām . \\
\noindent \textbf{(P\_6,1.102.2)} napuṃsakasya na iti anuvartiṣyate . \\
\noindent \textbf{(P\_6,1.103)} kim idam natvam puṃsām bahutve bhavati āhosvit puṃśabdāt bahuṣu . \\
\noindent \textbf{(P\_6,1.103)} kaḥ ca atra viśeṣaḥ . \\
\noindent \textbf{(P\_6,1.103)} natvam puṃsām bahutve cet puṃśabdāt iṣyate striyām . \\
\noindent \textbf{(P\_6,1.103)} tat na sidhyati . \\
\noindent \textbf{(P\_6,1.103)} bhrūkuṃsān paśya iti . \\
\noindent \textbf{(P\_6,1.103)} napuṃsake tathā eva iṣṭam . \\
\noindent \textbf{(P\_6,1.103)} tat na sidhyati . \\
\noindent \textbf{(P\_6,1.103)} ṣaṇḍhān paśya . \\
\noindent \textbf{(P\_6,1.103)} paṇḍakān paśya iti . \\
\noindent \textbf{(P\_6,1.103)} strīśabdāt ca prasajyate . strīśabdāt ca prāpnoti : cañcāḥ paśya , vadhrikāḥ paśya , kharakuṭīḥ paśya . \\
\noindent \textbf{(P\_6,1.103)} astu tarhi puṃśabdāt bahuṣu . \\
\noindent \textbf{(P\_6,1.103)} puṃśabāt iti cet iṣṭam sthūrāpatyam na sidhyati . sthūrān paśya iti . \\
\noindent \textbf{(P\_6,1.103)} kuṇḍinyāḥ ararakāyāḥ . \\
\noindent \textbf{(P\_6,1.103)} apatyam ca na sidhyati . \\
\noindent \textbf{(P\_6,1.103)} kuṇḍinān paśya . \\
\noindent \textbf{(P\_6,1.103)} ararakān paśya . \\
\noindent \textbf{(P\_6,1.103)} puṃsprādhānyāt prasidhyati . puṃspradhānā ete śabdāḥ . \\
\noindent \textbf{(P\_6,1.103)} tataḥ natvam bhaviṣyati . \\
\noindent \textbf{(P\_6,1.103)} puṃsprādhānye te eva syuḥ ye doṣāḥ pūrvacoditāḥ . \\
\noindent \textbf{(P\_6,1.103)} bhrūkuṃsān paśya . \\
\noindent \textbf{(P\_6,1.103)} ṣaṇḍhān paśya . \\
\noindent \textbf{(P\_6,1.103)} paṇḍakān paśya . \\
\noindent \textbf{(P\_6,1.103)} cañcāḥ paśya . \\
\noindent \textbf{(P\_6,1.103)} vadhrikāḥ paśya . \\
\noindent \textbf{(P\_6,1.103)} kharakuṭīḥ paśya iti . \\
\noindent \textbf{(P\_6,1.103)} tasmāt yasmin pakṣe alpīyāṃsaḥ doṣāḥ tam āsthāya pratividheyam doṣeṣu . \\
\noindent \textbf{(P\_6,1.107)} vā chandasi iti eva . \\
\noindent \textbf{(P\_6,1.107)} yamīm ca yamyam ca . \\
\noindent \textbf{(P\_6,1.107)} śamīm ca śamyam ca . \\
\noindent \textbf{(P\_6,1.107)} garuīm ca gauryam ca . \\
\noindent \textbf{(P\_6,1.107)} kiśorīm ca kiśoryam ca . \\
\noindent \textbf{(P\_6,1.108.1)} vā chandasi iti eva . \\
\noindent \textbf{(P\_6,1.108.1)} mitrāvaruṇau yajyamānaḥ . \\
\noindent \textbf{(P\_6,1.108.1)} mitrāvaruṇau ijyamānaḥ . \\
\noindent \textbf{(P\_6,1.108.2)} samprasāraṇāt pūrvatve samānāṅgagrahaṇam asamānāṅgapratiṣedhārtham . \\
\noindent \textbf{(P\_6,1.108.2)} samprasāraṇāt pūrvatve samānāṅgagrahaṇam kartavyam . \\
\noindent \textbf{(P\_6,1.108.2)} kim prayojanam . \\
\noindent \textbf{(P\_6,1.108.2)} asamānāṅgapratiṣedhārtham . \\
\noindent \textbf{(P\_6,1.108.2)} asamānāṅgasya mā bhūt iti . \\
\noindent \textbf{(P\_6,1.108.2)} śakahvartham , parivyartham . \\
\noindent \textbf{(P\_6,1.108.2)} siddham asamprasāraṇāt . \\
\noindent \textbf{(P\_6,1.108.2)} siddham etat . \\
\noindent \textbf{(P\_6,1.108.2)} katham . \\
\noindent \textbf{(P\_6,1.108.2)} asamprasāraṇāt . \\
\noindent \textbf{(P\_6,1.108.2)} vākyasya samprasāraṇasañjñā na varṇasya . \\
\noindent \textbf{(P\_6,1.108.2)} atha varṇasya samprasāraṇasañjñāyām doṣaḥ eva . \\
\noindent \textbf{(P\_6,1.108.2)} varṇasya ca samprasāraṇasañjñāyām na doṣaḥ . \\
\noindent \textbf{(P\_6,1.108.2)} katham . \\
\noindent \textbf{(P\_6,1.108.2)} anyaḥ ayam samprasāraṇāsamprasāraṇayoḥ sthāne ekaḥ ādiśyate . \\
\noindent \textbf{(P\_6,1.108.2)} kāryakṛtatvāt vā . \\
\noindent \textbf{(P\_6,1.108.2)} atha vā sakṛt kṛtam pūrvatvam iti kṛtvā punaḥ na bhaviṣyati . \\
\noindent \textbf{(P\_6,1.108.2)} tat yathā vasante brāhmaṇaḥ agnīn ādadhīta iti sakṛt ādhāya kṛtaḥ śāstrārthaḥ iti kṛtvā punaḥ pravṛttiḥ na bhavati . \\
\noindent \textbf{(P\_6,1.108.2)} viṣamaḥ upanyāsaḥ . \\
\noindent \textbf{(P\_6,1.108.2)} yuktam yat tasya punaḥ pravṛttiḥ na bhavati . \\
\noindent \textbf{(P\_6,1.108.2)} yaḥ tu tadāśrayam prāpnoti na tat śakyam bādhitum . \\
\noindent \textbf{(P\_6,1.108.2)} tat yathā vasante brāhmaṇaḥ agniṣṭomādibhiḥ kratubhiḥ yajeta iti agnyādhānanimittam vasante vasante ijyate . \\
\noindent \textbf{(P\_6,1.108.2)} tasmāt pūrvoktaḥ eva parihāraḥ siddham asamprasāraṇāt iti . \\
\noindent \textbf{(P\_6,1.108.2)} yadi tarhi na idam samprasāraṇam hūtaḥ iti dīrghatvam na prāpnoti . \\
\noindent \textbf{(P\_6,1.108.2)} dīrghatvam vacanaprāmāṇyāt . \\
\noindent \textbf{(P\_6,1.108.2)} anavakāśam dīrghatvam . \\
\noindent \textbf{(P\_6,1.108.2)} tat vacanaprāmāṇyāt bhaviṣyati . \\
\noindent \textbf{(P\_6,1.108.2)} antavattvāt vā . \\
\noindent \textbf{(P\_6,1.108.2)} atha vā pūrvasya kāryam prati antavat bhavati iti dīrghatvam bhaviṣyati . \\
\noindent \textbf{(P\_6,1.108.2)} āṭaḥ vṛddheḥ iyaṅ . \\
\noindent \textbf{(P\_6,1.108.2)} āṭaḥ vṛddhiḥ bhavati iti etasmāt iyaṅ bhavati vipratiṣedhena . \\
\noindent \textbf{(P\_6,1.108.2)} āṭaḥ vṛddhiḥ bhavati iti asya avakāśaḥ aikṣiṣṭa , aihiṣṭa . \\
\noindent \textbf{(P\_6,1.108.2)} iyaṅaḥ avakāśaḥ : adhīyāte , adhīyate . \\
\noindent \textbf{(P\_6,1.108.2)} iha ubhayam prāpnoti : adhyaiyātām adhyaiyata . \\
\noindent \textbf{(P\_6,1.108.2)} iyaṅādeśaḥ bhavati vipratiṣedhena . \\
\noindent \textbf{(P\_6,1.108.2)} na eṣaḥ yuktaḥ vipratiṣedhaḥ . \\
\noindent \textbf{(P\_6,1.108.2)} antaraṅgā āṭaḥ vṛddhiḥ . \\
\noindent \textbf{(P\_6,1.108.2)} kā antaraṅgatā . \\
\noindent \textbf{(P\_6,1.108.2)} varṇau āśritya āṭaḥ vṛddhiḥ . \\
\noindent \textbf{(P\_6,1.108.2)} aṅgasya iyaṅ ādeśaḥ . \\
\noindent \textbf{(P\_6,1.108.2)} evam tarhi idam iha sampradhāryam . \\
\noindent \textbf{(P\_6,1.108.2)} āṭ kriyatām iyaṅādeśaḥ iti kim atra kartavyam . \\
\noindent \textbf{(P\_6,1.108.2)} paratvāt iyaṅ . \\
\noindent \textbf{(P\_6,1.108.2)} nityaḥ āṭ āgamaḥ . \\
\noindent \textbf{(P\_6,1.108.2)} kṛte api iyaṅi prāpnoti akṛte api . \\
\noindent \textbf{(P\_6,1.108.2)} iyaṅ api nityaḥ . \\
\noindent \textbf{(P\_6,1.108.2)} kṛte api āṭi prāpnoti akṛte api . \\
\noindent \textbf{(P\_6,1.108.2)} anityaḥ iyaṅ . \\
\noindent \textbf{(P\_6,1.108.2)} na hi kṛte āti prāpnoti . \\
\noindent \textbf{(P\_6,1.108.2)} kim kāraṇam . \\
\noindent \textbf{(P\_6,1.108.2)} antaraṅgā āṭaḥ vṛddhiḥ . \\
\noindent \textbf{(P\_6,1.108.2)} yasya ca lakṣaṇāntareṇa nimittam vihanyate na tat anityam . \\
\noindent \textbf{(P\_6,1.108.2)} na ca atra āt eva iyaṅaḥ nimittam vihanti . \\
\noindent \textbf{(P\_6,1.108.2)} avaśyam lakṣaṇāntaram āṭaḥ vṛddhiḥ pratīkṣyā . \\
\noindent \textbf{(P\_6,1.108.2)} ubhayoḥ nityayoḥ paratvāt iyaṅ ādeśaḥ . \\
\noindent \textbf{(P\_6,1.108.2)} āt guṇāt savarṇadīrghatvam āṅabhyāsayoḥ . \\
\noindent \textbf{(P\_6,1.108.2)} āt guṇāt savarṇadīrghatvam bhavati vipratiṣedhena . \\
\noindent \textbf{(P\_6,1.108.2)} kva . \\
\noindent \textbf{(P\_6,1.108.2)} āṅabhyāsayoḥ . \\
\noindent \textbf{(P\_6,1.108.2)} āt guṇasya avakāśaḥ : khaṭvendraḥ , khaṭvodakam . \\
\noindent \textbf{(P\_6,1.108.2)} savarṇadīrghatvasya avakāśaḥ : daṇḍāgram , kṣupāgram . \\
\noindent \textbf{(P\_6,1.108.2)} iha ubhayam prāpnoti : adya , ā , ūḍhā : adyoḍhā , kadā , ā , ūḍhā : kadoḍhā , upa , i , ijatuḥ : upejatuḥ , upa , u , upatuḥ : upopatuḥ . \\
\noindent \textbf{(P\_6,1.108.2)} savarṇadīrghatvam bhavati vipratiṣedhena . \\
\noindent \textbf{(P\_6,1.108.2)} abhyāsārthena tāvat na arthaḥ . \\
\noindent \textbf{(P\_6,1.108.2)} astu atra āt guṇaḥ ayavau ca halādiśeṣaḥ . \\
\noindent \textbf{(P\_6,1.108.2)} punaḥ āt guṇaḥ bhaviṣyati . \\
\noindent \textbf{(P\_6,1.108.2)} bhavet siddham upejatuḥ , upejatuḥ iti . \\
\noindent \textbf{(P\_6,1.108.2)} idam tu na sidhyati : upopatuḥ , upopuḥ iti . \\
\noindent \textbf{(P\_6,1.108.2)} atra hi āt guṇe kṛte odantaḥ nipātaḥ iti pragṛhyasañjñā , pragṛhyaḥ prakṛtyā iti pragṛhyāśrayaḥ prakṛtibhāvaḥ prāpnoti . \\
\noindent \textbf{(P\_6,1.108.2)} padāntaprakaraṇe prkṛtibhāvaḥ na ca eṣaḥ padāntaḥ . \\
\noindent \textbf{(P\_6,1.108.2)} padāntabhaktaḥ padāntagrahaṇena grāhīṣyate . \\
\noindent \textbf{(P\_6,1.108.2)} evam tarhi etat eva atra na asti odantaḥ nipātaḥ iti . \\
\noindent \textbf{(P\_6,1.108.2)} kim kāraṇam . \\
\noindent \textbf{(P\_6,1.108.2)} lakṣaṇapratipadoktayoḥ pratipadoktasya eva iti . \\
\noindent \textbf{(P\_6,1.108.2)} iha api tarhi adyoḍhā , kadoḍhā iti bhavet rūpam siddham syāt . svare doṣaḥ tu . \\
\noindent \textbf{(P\_6,1.108.2)} svare tu doṣaḥ bhavati . \\
\noindent \textbf{(P\_6,1.108.2)} adyoḍhā* evam svaraḥ prasajyeta. adyoḍhā* iti ca iṣyate . \\
\noindent \textbf{(P\_6,1.108.2)} āṅi pararūpavacanam ca idānīm anarthakam syāt . \\
\noindent \textbf{(P\_6,1.108.2)} na anarthakam . \\
\noindent \textbf{(P\_6,1.108.2)} jñāpakārtham . \\
\noindent \textbf{(P\_6,1.108.2)} kim jñāpyam . \\
\noindent \textbf{(P\_6,1.108.2)} āṅi pararūpavacanam tu jñāpakam antaraṅgabalīyastvāt . \\
\noindent \textbf{(P\_6,1.108.2)} etat jñāpayati ācāryaḥ . \\
\noindent \textbf{(P\_6,1.108.2)} antaraṅgam balīyaḥ bhavati iti . \\
\noindent \textbf{(P\_6,1.108.2)} kim punaḥ iha antaraṅgam kim bahiraṅgam yāvatā dve pade āśritya savarṇadīrghatvam bhavati āt guṇaḥ api . \\
\noindent \textbf{(P\_6,1.108.2)} dhātūpasargayoḥ yat kāryam tat antaraṅgam . \\
\noindent \textbf{(P\_6,1.108.2)} kutaḥ etat . \\
\noindent \textbf{(P\_6,1.108.2)} pūrvam upasargasya dhātuna yogaḥ bhavati na adya śabdena . \\
\noindent \textbf{(P\_6,1.108.2)} kimartham tarhi adyaśabdaḥ prayujyate . \\
\noindent \textbf{(P\_6,1.108.2)} adyaśabdaysa api samudāyena yogaḥ bhavati . \\
\noindent \textbf{(P\_6,1.108.2)} kim etasya jñāpane prayojanam . \\
\noindent \textbf{(P\_6,1.108.2)} prayojanam pūrvasavarṇapūrvatvatahilopaṭenaṅeyyaṅisminṅiṇalautvam antaraṅgam bahiraṅgalakṣaṇāt varṇavikārāt . \\
\noindent \textbf{(P\_6,1.108.2)} pūrvasavarṇaḥ prayojanam . \\
\noindent \textbf{(P\_6,1.108.2)} agnī atra , vāyū atra . \\
\noindent \textbf{(P\_6,1.108.2)} pūrvasavarṇaḥ ca prāpnoti bahiraṅgalakṣaṇaḥ ca varṇavikāraḥ āvādeśaḥ . \\
\noindent \textbf{(P\_6,1.108.2)} pūrvasavarṇadīrghatvam bhavati antaraṅgataḥ . \\
\noindent \textbf{(P\_6,1.108.2)} pūrvatva . \\
\noindent \textbf{(P\_6,1.108.2)} śakahvartham , parivyartham . \\
\noindent \textbf{(P\_6,1.108.2)} pūrvatvam ca prāpnoti bahiraṅgalakṣaṇaḥ ca varṇavikāraḥ savarṇadīrghatvam . \\
\noindent \textbf{(P\_6,1.108.2)} pūrvatvam bhavati antaraṅgataḥ . \\
\noindent \textbf{(P\_6,1.108.2)} tahilopa . \\
\noindent \textbf{(P\_6,1.108.2)} akāri atra . \\
\noindent \textbf{(P\_6,1.108.2)} ahāri atra . \\
\noindent \textbf{(P\_6,1.108.2)} paca idam . \\
\noindent \textbf{(P\_6,1.108.2)} tahilopau ca prāpnutaḥ bahiraṅgalakṣaṇaḥ ca varṇavikāraḥ savarṇadīrghatvam . \\
\noindent \textbf{(P\_6,1.108.2)} tahilopau bhavataḥ antaraṅgataḥ . \\
\noindent \textbf{(P\_6,1.108.2)} ṭena . \\
\noindent \textbf{(P\_6,1.108.2)} vṛkṣeṇa atra , plakṣeṇa atra . \\
\noindent \textbf{(P\_6,1.108.2)} inādeśaḥ ca prāpnoti bahiraṅgalakṣaṇaḥ ca varṇavikāraḥ savarṇadīrghatvam . \\
\noindent \textbf{(P\_6,1.108.2)} inādeśaḥ bhavati antaraṅgataḥ . \\
\noindent \textbf{(P\_6,1.108.2)} ṅerya . \\
\noindent \textbf{(P\_6,1.108.2)} vṛkṣāya atra , plakṣāya atra . \\
\noindent \textbf{(P\_6,1.108.2)} ṅeḥ yādeśaḥ ca prāpnoti bahiraṅgalakṣaṇaḥ ca varṇavikāraḥ eṅaḥ padāntāt ati iti pararūpatvam . \\
\noindent \textbf{(P\_6,1.108.2)} ṅeḥ yādeśaḥ bhavati antaraṅgataḥ . \\
\noindent \textbf{(P\_6,1.108.2)} ṅismin . \\
\noindent \textbf{(P\_6,1.108.2)} yasmin idam , tasmin idam . \\
\noindent \textbf{(P\_6,1.108.2)} sminbhāvaḥ ca prāpnoti bahiraṅgalakṣaṇaḥ ca varṇavikāraḥ savarṇadīrghatvam . \\
\noindent \textbf{(P\_6,1.108.2)} sminbhāvaḥ bhavati antaraṅgataḥ . \\
\noindent \textbf{(P\_6,1.108.2)} ṅiṇalautvam . \\
\noindent \textbf{(P\_6,1.108.2)} agnau idam , yayau atra . \\
\noindent \textbf{(P\_6,1.108.2)} ṅiṇalautvam prāpnoti bahiraṅgalakṣaṇaḥ ca varṇavikāraḥ savarṇadīrghatvam . \\
\noindent \textbf{(P\_6,1.108.2)} autvam bhavati antaraṅgataḥ . \\
\noindent \textbf{(P\_6,1.108.2)} na etāni santi prayojanāni . \\
\noindent \textbf{(P\_6,1.108.2)} vipratiṣedhena api etāni siddhāni . \\
\noindent \textbf{(P\_6,1.108.2)} idam tarhi prayojanam . \\
\noindent \textbf{(P\_6,1.108.2)} vṛkṣāḥ atra . \\
\noindent \textbf{(P\_6,1.108.2)} plakṣāḥ atra . \\
\noindent \textbf{(P\_6,1.108.2)} pūrvasavarṇaḥ ca prāpnoti bahiraṅgalakṣaṇaḥ ca varṇavikāraḥ roḥ aplutāt aplute iti uttvam . \\
\noindent \textbf{(P\_6,1.108.2)} pūrvasavarṇaḥ bhavati antaraṅgataḥ . \\
\noindent \textbf{(P\_6,1.108.2)} na ca avaśyam idam eva prayojanam . \\
\noindent \textbf{(P\_6,1.108.2)} ādye yoge bahūni prayojanāni santi yadartham eṣā paribhāṣā kartavyā . \\
\noindent \textbf{(P\_6,1.108.2)} pratividheyam doṣeṣu . \\
\noindent \textbf{(P\_6,1.112)} kim idam khyatyāt iti . \\
\noindent \textbf{(P\_6,1.112)} sakhipatyoḥ vikṛtagrahaṇam . \\
\noindent \textbf{(P\_6,1.112)} kim punaḥ kāraṇam sakhipatyoḥ vikṛtagrahaṇam kriyate na sakhipatibhyām iti eva ucyeta . \\
\noindent \textbf{(P\_6,1.112)} na evam śakyam . \\
\noindent \textbf{(P\_6,1.112)} garīyān ca eva hi nirdeśaḥ syāt iha ca prasajyeta . \\
\noindent \textbf{(P\_6,1.112)} atisakheḥ āgacchāmi . \\
\noindent \textbf{(P\_6,1.112)} atisakheḥ svam . \\
\noindent \textbf{(P\_6,1.112)} iha ca na syāt : sakhīyateḥ apratyayaḥ sakhyuḥ , patyuḥ . \\
\noindent \textbf{(P\_6,1.112)} lunīyateḥ apratyayaḥ . \\
\noindent \textbf{(P\_6,1.112)} lūnyuḥ , pūnyuḥ . \\
\noindent \textbf{(P\_6,1.113)} kimartham aplutāt aplute iti ucyate . \\
\noindent \textbf{(P\_6,1.113)} plutāt parasya plute vā parataḥ mā bhūt iti . \\
\noindent \textbf{(P\_6,1.113)} plutāt parasya susrotā3 atra nu asi . \\
\noindent \textbf{(P\_6,1.113)} plute parataḥ tiṣṭhatu payaḥ ā3gnidatta . \\
\noindent \textbf{(P\_6,1.113)} ataḥ ati iti ucyate . \\
\noindent \textbf{(P\_6,1.113)} kaḥ prasaṅgaḥ plutāt parasya plute vā parataḥ . \\
\noindent \textbf{(P\_6,1.113)} asiddhaḥ plutaḥ . \\
\noindent \textbf{(P\_6,1.113)} tasya asiddhatvāt prāpnoti . \\
\noindent \textbf{(P\_6,1.113)} atha aplutāt aplute iti ucyamāne yāvatā asiddhaḥ plutaḥ kasmāt eva atra na prāpnoti . \\
\noindent \textbf{(P\_6,1.113)} aplutabhāvinaḥ aplutabhāvini iti evam etat vijñāyate . \\
\noindent \textbf{(P\_6,1.113)} na etat asti prayojanam . \\
\noindent \textbf{(P\_6,1.113)} siddhaḥ plutaḥ svarasandhiṣu . \\
\noindent \textbf{(P\_6,1.113)} katham jñāyate . \\
\noindent \textbf{(P\_6,1.113)} yat ayam plutaḥ prakṛtyā iti plutasya prakṛtibhāvam śāsti . \\
\noindent \textbf{(P\_6,1.113)} sataḥ hi kāryiṇaḥ kāryeṇa bhavitavyam . \\
\noindent \textbf{(P\_6,1.113)} aplutādaplutavacane akārahaśoḥ samānapade pratiṣedhaḥ . \\
\noindent \textbf{(P\_6,1.113)} aplutādaplutavacane akārahaśoḥ samānapade pratiṣedhaḥ vaktavyaḥ . \\
\noindent \textbf{(P\_6,1.113)} payo3ṭ , payo3da . \\
\noindent \textbf{(P\_6,1.113)} na vā bahiraṅgalakṣaṇatvāt . \\
\noindent \textbf{(P\_6,1.113)} na vā vaktavyaḥ . \\
\noindent \textbf{(P\_6,1.113)} kim kāraṇam . \\
\noindent \textbf{(P\_6,1.113)} bahiraṅgalakṣaṇatvāt . \\
\noindent \textbf{(P\_6,1.113)} bahiraṅgaḥ pluraḥ . \\
\noindent \textbf{(P\_6,1.113)} antaraṅgam uttvam . \\
\noindent \textbf{(P\_6,1.113)} asiddham bahiraṅgam antaraṅge . \\
\noindent \textbf{(P\_6,1.113)} iha api tarhi prāpnoti . \\
\noindent \textbf{(P\_6,1.113)} susrotā3 atra nu asi . \\
\noindent \textbf{(P\_6,1.113)} antaraṅgaḥ atra plutaḥ bahiraṅgam uttvam . \\
\noindent \textbf{(P\_6,1.113)} kva punaḥ iha antaraṅgaḥ plutaḥ kva vā bahiraṅgam uttvam uttvam vā antaraṅgam plutaḥ vā bahiraṅgaḥ . \\
\noindent \textbf{(P\_6,1.113)} vākyāntasya vākyādau antaraṅgaḥ plutaḥ bahiraṅgam uttvam . \\
\noindent \textbf{(P\_6,1.113)} samānavākye padāntasya padādau uttvam antaraṅgam bahiraṅgaḥ plutaḥ . \\
\noindent \textbf{(P\_6,1.113)} kim punaḥ kāraṇam bahiraṅgatvam uttve hetuḥ vyapadiśyate na punaḥ asiddhatvam api . \\
\noindent \textbf{(P\_6,1.113)} yathā eva hi ayam bahiraṅgaḥ evam asiddhaḥ api . \\
\noindent \textbf{(P\_6,1.113)} evam manyate . \\
\noindent \textbf{(P\_6,1.113)} asiddhaḥ plutaḥ āśrayāt siddhaḥ bhavati . \\
\noindent \textbf{(P\_6,1.113)} atha vā yasyām na aprāptāyām paribhāṣāyām uttvam ārabhyate sā āśrayāt siddhā syāt . \\
\noindent \textbf{(P\_6,1.113)} kasyām ca na aprāptāyām . \\
\noindent \textbf{(P\_6,1.113)} asiddhaparibhāṣāyām . \\
\noindent \textbf{(P\_6,1.113)} bahiraṅgaparibhāṣāyām punaḥ prāptāyām aprāptāyām ca . \\
\noindent \textbf{(P\_6,1.115)} kasya ayam pratiṣedhaḥ . \\
\noindent \textbf{(P\_6,1.115)} nāntaḥpādam iti sarvapratiṣedhaḥ . \\
\noindent \textbf{(P\_6,1.115)} nāntaḥpādam iti sarvasya ayam pratiṣedhaḥ . \\
\noindent \textbf{(P\_6,1.115)} katham . \\
\noindent \textbf{(P\_6,1.115)} aci iti vartate . \\
\noindent \textbf{(P\_6,1.115)} aci yat prāpnoti tasya pratiṣedhaḥ . \\
\noindent \textbf{(P\_6,1.115)} nāntaḥpādam iti sarvapratiṣedhaḥ cet atiprasaṅgaḥ . \\
\noindent \textbf{(P\_6,1.115)} nāntaḥpādam iti sarvapratiṣedhaḥ cet atiprasaṅgaḥ bhavati . \\
\noindent \textbf{(P\_6,1.115)} iha api prāpnoti . \\
\noindent \textbf{(P\_6,1.115)} anu agniḥ uṣasām agram akhyat , prati agniḥ uṣasām agram akhyat . \\
\noindent \textbf{(P\_6,1.115)} evam tarhi ati iti vartate . \\
\noindent \textbf{(P\_6,1.115)} akārāśrayam yat prāpnoti tasya pratiṣedhaḥ . \\
\noindent \textbf{(P\_6,1.115)} akārāśrayam iti cet uttvavacanam . \\
\noindent \textbf{(P\_6,1.115)} akārāśrayam iti cet uttvam vaktavyam . \\
\noindent \textbf{(P\_6,1.115)} kālaḥ aśvaḥ . \\
\noindent \textbf{(P\_6,1.115)} śatadhāraḥ ayam maṇiḥ . \\
\noindent \textbf{(P\_6,1.115)} ayavoḥ pratiṣedhaḥ ca . \\
\noindent \textbf{(P\_6,1.115)} ayavoḥ ca pratiṣedhaḥ ca vaktavyaḥ . \\
\noindent \textbf{(P\_6,1.115)} sujāte aśvasūnṛte . \\
\noindent \textbf{(P\_6,1.115)} adhvaro adribhiḥ sutam . \\
\noindent \textbf{(P\_6,1.115)} śukram te anyat . \\
\noindent \textbf{(P\_6,1.115)} eṅprakaraṇāt siddham . \\
\noindent \textbf{(P\_6,1.115)} eṅaḥ ati iti vartate . \\
\noindent \textbf{(P\_6,1.115)} eṅaḥ ati yat prāpnoti tasya pratiṣedhaḥ . \\
\noindent \textbf{(P\_6,1.115)} eṅprakaraṇāt siddham cet uttvapratiṣedhaḥ . \\
\noindent \textbf{(P\_6,1.115)} eṅprakaraṇāt siddham cet uttvapratiṣedhaḥ vaktavyaḥ . \\
\noindent \textbf{(P\_6,1.115)} agneḥ atra , vāyoḥ atra. ataḥ roḥ aplutāt aplute eṅaḥ ca iti uttvam prāpnoti . \\
\noindent \textbf{(P\_6,1.115)} punaḥ eṅgrahaṇāt siddham . \\
\noindent \textbf{(P\_6,1.115)} punaḥ eṅgrahaṇam kartavyam . \\
\noindent \textbf{(P\_6,1.115)} tat tarhi kartavyam . \\
\noindent \textbf{(P\_6,1.115)} na kartavyam . \\
\noindent \textbf{(P\_6,1.115)} prakṛtam anuvartate . \\
\noindent \textbf{(P\_6,1.115)} nanu ca uktam eṅprakaraṇāt siddham cet uttvapratiṣedhaḥ iti . \\
\noindent \textbf{(P\_6,1.115)} na eṣaḥ doṣaḥ . \\
\noindent \textbf{(P\_6,1.115)} padāntābhisambaddham eṅgrahaṇam anuvartate na ca eṅaḥ padāntāt paraḥ ruḥ asti . \\
\noindent \textbf{(P\_6,1.123)} goḥ agvacanam gavāgre svarasiddhyartham . \\
\noindent \textbf{(P\_6,1.123)} goḥ ak vaktavyaḥ . \\
\noindent \textbf{(P\_6,1.123)} kim prayojanam . \\
\noindent \textbf{(P\_6,1.123)} gavāgre svarasiddhyartham . \\
\noindent \textbf{(P\_6,1.123)} gavāgre svarasiddhiḥ yathā syāt . \\
\noindent \textbf{(P\_6,1.123)} gavāgram . \\
\noindent \textbf{(P\_6,1.123)} avaṅādeśe hi svare doṣaḥ . \\
\noindent \textbf{(P\_6,1.123)} avaṅādeśe hi svare doṣaḥ syāt . \\
\noindent \textbf{(P\_6,1.123)} antodāttasya āntaryataḥ antodāttaḥ ādeśaḥ prasjyate . \\
\noindent \textbf{(P\_6,1.123)} katham punaḥ ayam antodāttaḥ yadā ekāc . \\
\noindent \textbf{(P\_6,1.123)} vyapadeśivadbhāvena . \\
\noindent \textbf{(P\_6,1.123)} yathā eva tarhi vyapadeśivadbhāvena antodāttaḥ evam ādyudāttaḥ api . \\
\noindent \textbf{(P\_6,1.123)} tatra āntaryataḥ ādyudāttasya ādyudāttaḥ ādeśaḥ bhavati . \\
\noindent \textbf{(P\_6,1.123)} satyam evam etat . \\
\noindent \textbf{(P\_6,1.123)} na tu idam lakṣaṇam asti . \\
\noindent \textbf{(P\_6,1.123)} prātipadikasya ādiḥ udāttaḥ bhavati iti . \\
\noindent \textbf{(P\_6,1.123)} idam punaḥ asti . \\
\noindent \textbf{(P\_6,1.123)} prātipadikasya antaḥ udāttaḥ bhavati iti . \\
\noindent \textbf{(P\_6,1.123)} saḥ asau lakṣaṇena antodāttaḥ . \\
\noindent \textbf{(P\_6,1.123)} tatra āntaryataḥ antodāttasya antodāttaḥ ādeśaḥ prasjyeta . \\
\noindent \textbf{(P\_6,1.123)} yadi punaḥ gameḥ ḍo vidhīyeta . \\
\noindent \textbf{(P\_6,1.123)} kim kṛtam bhavati . \\
\noindent \textbf{(P\_6,1.123)} pratyayādyudāttatve kṛte āntaryataḥ ādyudāttasya ādyudāttaḥ ādeśaḥ bhaviṣyati . \\
\noindent \textbf{(P\_6,1.123)} katham punaḥ ayam ādyudāttaḥ yadā ekāc . \\
\noindent \textbf{(P\_6,1.123)} vyapadeśivadbhāvena . \\
\noindent \textbf{(P\_6,1.123)} yathā eva tarhi vyapadeśivadbhāvena ādyudāttaḥ evam antodāttaḥ api . \\
\noindent \textbf{(P\_6,1.123)} tatra āntaryataḥ antodāttasya antodāttaḥ ādeśaḥ prasjyeta . \\
\noindent \textbf{(P\_6,1.123)} satyam evam etat . \\
\noindent \textbf{(P\_6,1.123)} na tu idam lakṣaṇam asti . \\
\noindent \textbf{(P\_6,1.123)} pratyayasya antaḥ udāttaḥ bhavati iti . \\
\noindent \textbf{(P\_6,1.123)} idam punaḥ asti . \\
\noindent \textbf{(P\_6,1.123)} pratyayasya ādiḥ udāttaḥ bhavati iti . \\
\noindent \textbf{(P\_6,1.123)} saḥ asau lakṣaṇena ādyudāttaḥ . \\
\noindent \textbf{(P\_6,1.123)} tatra āntaryataḥ ādyudāttasya ādyudāttaḥ ādeśaḥ bhaviṣyati .etat api ādeśe na asti . \\
\noindent \textbf{(P\_6,1.123)} ādeśasya ādiḥ udāttaḥ bhavati iti . \\
\noindent \textbf{(P\_6,1.123)} prakṛtitaḥ anena svaraḥ labhyaḥ . \\
\noindent \textbf{(P\_6,1.123)} prakṛtiḥ ca asya yathā eva ādyudāttā evam antodāttā api . \\
\noindent \textbf{(P\_6,1.123)} evam tarhi ādyudāttanipātanam kariṣyate . \\
\noindent \textbf{(P\_6,1.123)} saḥ nipātanasvaraḥ prakṛtisvarasya bādhakaḥ bhaviṣyati . \\
\noindent \textbf{(P\_6,1.123)} evam api upadeśivadbhāvaḥ vaktayaḥ . \\
\noindent \textbf{(P\_6,1.123)} yathā eva nipātanasvaraḥ prakṛtsvarasya bādhakaḥ evam samāsasvarasya api . \\
\noindent \textbf{(P\_6,1.123)} gavāsthi , gavākṣi . \\
\noindent \textbf{(P\_6,1.124)} indrādau iti vaktavyam iha api yathā syāt . \\
\noindent \textbf{(P\_6,1.124)} gavendrayajñe vīhi iti . \\
\noindent \textbf{(P\_6,1.124)} tat tarhi vaktavyam . \\
\noindent \textbf{(P\_6,1.124)} na vaktavyam . \\
\noindent \textbf{(P\_6,1.124)} na evam vijñāyate indre aci iti . \\
\noindent \textbf{(P\_6,1.124)} katham tarhi . \\
\noindent \textbf{(P\_6,1.124)} aci bhavati . \\
\noindent \textbf{(P\_6,1.124)} katarasmin . \\
\noindent \textbf{(P\_6,1.124)} indre aci iti . \\
\noindent \textbf{(P\_6,1.125.1)} nityagrahaṇam kimartham . \\
\noindent \textbf{(P\_6,1.125.1)} vibhāṣā mā bhūt iti . \\
\noindent \textbf{(P\_6,1.125.1)} ma etat asti prayojanam . \\
\noindent \textbf{(P\_6,1.125.1)} pūrvasmin eva yoge vibhāṣāgrahaṇam nivṛttam . \\
\noindent \textbf{(P\_6,1.125.1)} idam tarhi prayojanam . \\
\noindent \textbf{(P\_6,1.125.1)} plutapragṛhyāṇam aci prakṛtibhāvaḥ eva yathā syāt . \\
\noindent \textbf{(P\_6,1.125.1)} yat anyat prāpnoti tat mā bhūt iti . \\
\noindent \textbf{(P\_6,1.125.1)} kim ca anyat prāpnoti . \\
\noindent \textbf{(P\_6,1.125.1)} śākalam . \\
\noindent \textbf{(P\_6,1.125.1)} sinnityasamāsayoḥ śākalapratiṣedham vakṣyati . \\
\noindent \textbf{(P\_6,1.125.1)} saḥ na vaktavyaḥ bhavati . \\
\noindent \textbf{(P\_6,1.125.1)} atha ajgrahaṇam kimartham . \\
\noindent \textbf{(P\_6,1.125.1)} aci prakṛtibhāvaḥ yathā syāt . \\
\noindent \textbf{(P\_6,1.125.1)} plutapragṛhyeṣu ajgrahaṇam anarthakam adhikārāt siddham . \\
\noindent \textbf{(P\_6,1.125.1)} plutapragṛhyeṣu ajgrahaṇam anarthakam . \\
\noindent \textbf{(P\_6,1.125.1)} kim kāraṇam . \\
\noindent \textbf{(P\_6,1.125.1)} adhikārāt eva siddham . \\
\noindent \textbf{(P\_6,1.125.1)} aci iti prakṛtam anuvartate . \\
\noindent \textbf{(P\_6,1.125.1)} kva prakṛtam . \\
\noindent \textbf{(P\_6,1.125.1)} ikaḥ yaṇ aci iti . \\
\noindent \textbf{(P\_6,1.125.1)} tat tu tasmin prakṛtibhāvārtham . \\
\noindent \textbf{(P\_6,1.125.1)} tat tu dvitīyam ajgrahaṇam kartavyam prakṛtibhāvārtham . \\
\noindent \textbf{(P\_6,1.125.1)} tasmin aci pūrvasya prakṛtibhāvaḥ yathā syāt . \\
\noindent \textbf{(P\_6,1.125.1)} iha mā bhūt . \\
\noindent \textbf{(P\_6,1.125.1)} jānu u asya rujati . \\
\noindent \textbf{(P\_6,1.125.1)} jānū asya rujati . \\
\noindent \textbf{(P\_6,1.125.1)} jānv asya rujati iti . \\
\noindent \textbf{(P\_6,1.125.2)} atha kimartham plutasya prakṛtibhāvaḥ ucyate . \\
\noindent \textbf{(P\_6,1.125.2)} svarasandhiḥ mā bhūt iti . \\
\noindent \textbf{(P\_6,1.125.2)} ucyamāne api etasmin svarsandhiḥ prāpnoti . \\
\noindent \textbf{(P\_6,1.125.2)} plute kṛte na bhaviṣyati . \\
\noindent \textbf{(P\_6,1.125.2)} asiddhaḥ plutaḥ . \\
\noindent \textbf{(P\_6,1.125.2)} tasya asiddhatvāt prāpnoti . \\
\noindent \textbf{(P\_6,1.125.2)} plutaprakṛtibhāvavacanam tu jñāpakam ekādeśāt plutaḥ vipratiṣedhena iti . \\
\noindent \textbf{(P\_6,1.125.2)} yat ayam plutaḥ prakṛtyā iti prakṛtibhāvam śāsti tat jñāpayati ācāryaḥ ekādeśāt plutaḥ bhavati vipratiṣedhena iti . \\
\noindent \textbf{(P\_6,1.125.2)} ekādeśāt plutaḥ vipratiṣedhena iti cet śālendre atiprasaṅgaḥ . \\
\noindent \textbf{(P\_6,1.125.2)} ekādeśāt plutaḥ vipratiṣedhena iti cet śālendre atiprasaṅgaḥ bhavati . \\
\noindent \textbf{(P\_6,1.125.2)} śālāyām indraḥ śālendraḥ . \\
\noindent \textbf{(P\_6,1.125.2)} na vā bahiraṅgalakṣaṇatvāt . \\
\noindent \textbf{(P\_6,1.125.2)} na vā atiprasaṅgaḥ . \\
\noindent \textbf{(P\_6,1.125.2)} kim kāraṇam . \\
\noindent \textbf{(P\_6,1.125.2)} bahiraṅgalakṣaṇatvāt . \\
\noindent \textbf{(P\_6,1.125.2)} bahiraṅgaḥ plutaḥ . \\
\noindent \textbf{(P\_6,1.125.2)} antaraṅgaḥ ekādeśaḥ . \\
\noindent \textbf{(P\_6,1.125.2)} asiddham bahiraṅgam antaraṅge . \\
\noindent \textbf{(P\_6,1.126)} āṅaḥ anarthakasya . \\
\noindent \textbf{(P\_6,1.126)} āṅaḥ anarthakasya iti vaktavyam . \\
\noindent \textbf{(P\_6,1.126)} iha mā bhūt . \\
\noindent \textbf{(P\_6,1.126)} indraḥ bāhubhyām ātarat . \\
\noindent \textbf{(P\_6,1.126)} tat tarhi vaktavyam . \\
\noindent \textbf{(P\_6,1.126)} na vaktavyam . \\
\noindent \textbf{(P\_6,1.126)} bahulagrahaṇāt na bhaviṣyati . \\
\noindent \textbf{(P\_6,1.126)} āṅaḥ anunāsikaḥ chandasi bahulam . \\
\noindent \textbf{(P\_6,1.127.1)} kimarthaḥ cakāraḥ . \\
\noindent \textbf{(P\_6,1.127.1)} prakṛtyā iti etat anukṛṣyate . \\
\noindent \textbf{(P\_6,1.127.1)} kim prayojanam . \\
\noindent \textbf{(P\_6,1.127.1)} svarasandhiḥ mā bhūt iti . \\
\noindent \textbf{(P\_6,1.127.1)} na etat asti prayojanam . \\
\noindent \textbf{(P\_6,1.127.1)} hrasvavacanasārmarthyāt na bhaviṣyati . \\
\noindent \textbf{(P\_6,1.127.1)} bhavet dīrghāṇām hrasvavacanasārmarthyāt svarasandhiḥ na syāt . \\
\noindent \textbf{(P\_6,1.127.1)} hrasvānām tu khalu svarasandhiḥ prāpnoti . \\
\noindent \textbf{(P\_6,1.127.1)} hrasvānām api hrasvavacanasāmarthyān svarasandhiḥ na bhaviṣyati . \\
\noindent \textbf{(P\_6,1.127.1)} na hrasvānām hrasvāḥ prāpnuvanti . \\
\noindent \textbf{(P\_6,1.127.1)} na hi bhuktavān punaḥ bhuṅkte . \\
\noindent \textbf{(P\_6,1.127.1)} na ca kṛtaśmaśruḥ punaḥ śmaśrūni kārayati . \\
\noindent \textbf{(P\_6,1.127.1)} nanu ca punaḥpravṛttiḥ api dṛṣṭā . \\
\noindent \textbf{(P\_6,1.127.1)} bhuktavān ca punaḥ bhuṅkte kṛtaśmaśruḥ ca punaḥ śmaśrūni kārayati . \\
\noindent \textbf{(P\_6,1.127.1)} sāmarthyāt tatra punaḥpravṛttiḥ bhavati bhojanaviśeṣāt śilpiviśeṣāt vā . \\
\noindent \textbf{(P\_6,1.127.1)} hrasvāṇām punaḥ hrasvavacane na kim cit prayojanam asti . \\
\noindent \textbf{(P\_6,1.127.1)} akṛtakāri khalu api śāstram agnivat . \\
\noindent \textbf{(P\_6,1.127.1)} tat yathā agniḥ yad adagdham tat dahati . \\
\noindent \textbf{(P\_6,1.127.1)} hrasvāṇām api hrasvavacane etat prayojanam svarasandhiḥ mā bhūt iti . \\
\noindent \textbf{(P\_6,1.127.1)} kṛtakāri khalu api śāstram parjanyavat . \\
\noindent \textbf{(P\_6,1.127.1)} tat yathā parjanyaḥ yāvat ūnam pūrṇam ca sarvam abhivarṣayati . \\
\noindent \textbf{(P\_6,1.127.1)} idam tarhi prayojanam . \\
\noindent \textbf{(P\_6,1.127.1)} plutapragṛhyāḥ anukṛṣyante . \\
\noindent \textbf{(P\_6,1.127.1)} ikaḥ asavarṇe śākalyasya hrasvaḥ ca plutaprgṛhyāḥ ca prakṛtyā . \\
\noindent \textbf{(P\_6,1.127.1)} nityagrahaṇasya api etat prayojanam uktam . \\
\noindent \textbf{(P\_6,1.127.1)} anyatarat śakyam akartum . \\
\noindent \textbf{(P\_6,1.127.2)} sinnityasamāsayoḥ śākalapratiṣedhaḥ . \\
\noindent \textbf{(P\_6,1.127.2)} sinnityasamāsayoḥ śākalapratiṣedhaḥ vaktavyaḥ . \\
\noindent \textbf{(P\_6,1.127.2)} ayam te yoniḥ ṛtviyaḥ . \\
\noindent \textbf{(P\_6,1.127.2)} prajām vindāma ṛtviyām . \\
\noindent \textbf{(P\_6,1.127.2)} vaiyākaraṇaḥ , sauvaśvaḥ . \\
\noindent \textbf{(P\_6,1.127.2)} nityagrahaṇena na arthaḥ . \\
\noindent \textbf{(P\_6,1.127.2)} sitsamāsayoḥ śākalam na bhavati iti eva . \\
\noindent \textbf{(P\_6,1.127.2)} idam api siddham bhavati . \\
\noindent \textbf{(P\_6,1.127.2)} vāpyām āsvaḥ , vāpyaśvaḥ , nadyām ātiḥ , nadyātiḥ . \\
\noindent \textbf{(P\_6,1.127.2)} īṣā akṣādiṣu chandasi prakṛtibhāvamātram . \\
\noindent \textbf{(P\_6,1.127.2)} īṣā akṣādiṣu chandasi prakṛtibhāvamātram draṣṭavyam . \\
\noindent \textbf{(P\_6,1.127.2)} īṣa akṣaḥ . \\
\noindent \textbf{(P\_6,1.127.2)} ka īm are piśaṅgila . \\
\noindent \textbf{(P\_6,1.127.2)} yathā aṅgadaḥ . \\
\noindent \textbf{(P\_6,1.128.1)} kimartham idam ucyate . \\
\noindent \textbf{(P\_6,1.128.1)} ṛti akaḥ savarṇārtham . \\
\noindent \textbf{(P\_6,1.128.1)} savarṇārthaḥ ayam ārambhaḥ . \\
\noindent \textbf{(P\_6,1.128.1)} hotṛ ṛśyaḥ . \\
\noindent \textbf{(P\_6,1.128.1)} anigantārtham ca . \\
\noindent \textbf{(P\_6,1.128.1)} khaṭva ṛśyaḥ , māla ṛśyaḥ . \\
\noindent \textbf{(P\_6,1.128.2)} ṛti hrasvāt upasargāt vṛddhiḥ vipratiṣedhena . ṛti hrasvaḥ bhavati iti etasmāt upasargāt vṛddhiḥ bhavati vipratiṣedhena . \\
\noindent \textbf{(P\_6,1.128.2)} ṛti hrasvaḥ bhavati iti etasya avakāśaḥ khaṭva ṛśyaḥ , māla ṛśyaḥ . \\
\noindent \textbf{(P\_6,1.128.2)} upasargāt vṛddheḥ avakāśaḥ . \\
\noindent \textbf{(P\_6,1.128.2)} vibhāṣā hrasvatvam . \\
\noindent \textbf{(P\_6,1.128.2)} yadā na hrasvatvam tadā avakāśaḥ . \\
\noindent \textbf{(P\_6,1.128.2)} hrasvaprasaṅge ubhayam prāpnoti . \\
\noindent \textbf{(P\_6,1.128.2)} upārdhnoti , prārdhnoti . \\
\noindent \textbf{(P\_6,1.128.2)} upasargāt vṛddhiḥ bhavati vipratiṣedhena . \\
\noindent \textbf{(P\_6,1.128.2)} saḥ tarhi pratiṣedhaḥ vaktavyaḥ . \\
\noindent \textbf{(P\_6,1.128.2)} na vaktavyaḥ . \\
\noindent \textbf{(P\_6,1.128.2)} uktam tatra dhātugrahaṇasya prayojanam . \\
\noindent \textbf{(P\_6,1.128.2)} upasargāt ṛti dhātau vṛddhiḥ eva yathā syāt . \\
\noindent \textbf{(P\_6,1.128.2)} anyat yat prāpnoti tat mā bhūt iti . \\
\noindent \textbf{(P\_6,1.129)} upasthite iti ucyate . \\
\noindent \textbf{(P\_6,1.129)} kim idam upasthitam nāma . \\
\noindent \textbf{(P\_6,1.129)} anārṣaḥ itikaraṇaḥ . \\
\noindent \textbf{(P\_6,1.129)} suślokā3 iti suśloketi . \\
\noindent \textbf{(P\_6,1.129)} atha vadvacanam kimartham . \\
\noindent \textbf{(P\_6,1.129)} vadvacanam plutakāryapratiṣedhārtham . \\
\noindent \textbf{(P\_6,1.129)} vadvacanam kriyate plutakāryapratiṣedhārtham . \\
\noindent \textbf{(P\_6,1.129)} plutakāryam pratiṣidhyate . \\
\noindent \textbf{(P\_6,1.129)} trimātratā na pratiṣidhyate . \\
\noindent \textbf{(P\_6,1.129)} kim ca idānīm trimātratāyāḥ apratiṣedhe prayojanam yāvatā plutakārye pratiṣiddhe svarasandhinā bhavitavyam . \\
\noindent \textbf{(P\_6,1.129)} plutapratiṣedhe hi pragṛhyaplutapratiṣedhaprasaṅgaḥ anyena vihitatvāt . \\
\noindent \textbf{(P\_6,1.129)} plutapratiṣedhe hi sati pragṛhyasya api plutasya trimātratāyāḥ pratiṣedhaḥ prasajyeta . \\
\noindent \textbf{(P\_6,1.129)} agnī3 iti , vāyū3 iti . \\
\noindent \textbf{(P\_6,1.129)} kim ca idānīm tasyāḥ api trimātratāyāḥ apratiṣedhe prayojanam yāvatā plutakārye pratiṣiddhe svarasandhinā bhavitavyam . \\
\noindent \textbf{(P\_6,1.129)} na bhavitavyam . \\
\noindent \textbf{(P\_6,1.129)} kim kāraṇam . \\
\noindent \textbf{(P\_6,1.129)} anyena vihitatvāt . \\
\noindent \textbf{(P\_6,1.129)} anyena hi lakṣaṇena plutapragṛhyasya prakṛtibhāvaḥ ucyate pragṛhyaḥ prakṛtyā iti . \\
\noindent \textbf{(P\_6,1.130)} kimartham idam ucyate . \\
\noindent \textbf{(P\_6,1.130)} ī3 cākravarmaṇasya iti anupasthitārtham . \\
\noindent \textbf{(P\_6,1.130)} anupasthitārthaḥ ayam ārambhaḥ . \\
\noindent \textbf{(P\_6,1.130)} cinu hi3 idam . \\
\noindent \textbf{(P\_6,1.130)} cinu hīdam . \\
\noindent \textbf{(P\_6,1.130)} sunu hi3 idam . \\
\noindent \textbf{(P\_6,1.130)} sunu hīdam . \\
\noindent \textbf{(P\_6,1.130)} īkāragrahaṇena na arthaḥ . \\
\noindent \textbf{(P\_6,1.130)} aviśeṣeṇa cākravarmaṇasya ācāryasya aplutavat bhavati iti eva . \\
\noindent \textbf{(P\_6,1.130)} idam api siddham bhavati . \\
\noindent \textbf{(P\_6,1.130)} vaśa3 iyam , vaśeyam . \\
\noindent \textbf{(P\_6,1.131)} kimarthaḥ takāraḥ . \\
\noindent \textbf{(P\_6,1.131)} taparaḥ tatkālasya iti tatkalaḥ yathā syāt . \\
\noindent \textbf{(P\_6,1.131)} na etat asti prayojanam . \\
\noindent \textbf{(P\_6,1.131)} āntaryataḥ ardhamātrikasya vyañjanasya mātrikaḥ bhaviṣyati . \\
\noindent \textbf{(P\_6,1.131)} na sidhyati . \\
\noindent \textbf{(P\_6,1.131)} ūṭhi kṛte āntaryataḥ dīrghasya dīrghaḥ prāpnoti . \\
\noindent \textbf{(P\_6,1.131)} tadartham taparaḥ kṛtaḥ . \\
\noindent \textbf{(P\_6,1.131)} evamarthaḥ taparaḥ kriyate . \\
\noindent \textbf{(P\_6,1.135.1)} kātpūrvagrahaṇam kimartham . \\
\noindent \textbf{(P\_6,1.135.1)} kāt pūrvaḥ yathā syāt . \\
\noindent \textbf{(P\_6,1.135.1)} saṃskartā , saṃskartum . \\
\noindent \textbf{(P\_6,1.135.1)} na etat asti prayojanam . \\
\noindent \textbf{(P\_6,1.135.1)} suṭ iti ādiliṅgaḥ ayam karotiḥ ca kakārādiḥ . \\
\noindent \textbf{(P\_6,1.135.1)} tatra antareṇa kātpūrvagrahaṇam kāt pūrvaḥ eva bhaviṣyati . \\
\noindent \textbf{(P\_6,1.135.1)} ataḥ uttaram paṭhati suṭi kātpūrvavacanam akakārādau kātpūrvārtham . \\
\noindent \textbf{(P\_6,1.135.1)} suṭi kātpūrvavacanam kriyate akakārādau kātpūrvaḥ yathā syāt . \\
\noindent \textbf{(P\_6,1.135.1)} sañcaskaratuḥ , sañcaskaruḥ . \\
\noindent \textbf{(P\_6,1.135.1)} suṭi kātpūrvavacanam akakārādau kātpūrvārtham iti cet antareṇa api tat siddham . \\
\noindent \textbf{(P\_6,1.135.1)} suṭi kātpūrvavacanam akakārādau kātpūrvārtham iti cet antareṇa api kātpūrvagrahaṇam siddham . \\
\noindent \textbf{(P\_6,1.135.1)} katham . \\
\noindent \textbf{(P\_6,1.135.1)} dvirvacanāt suṭ vipratiṣedhena . \\
\noindent \textbf{(P\_6,1.135.1)} dvirvacanam kriyatām suṭ iti suṭ bhaviṣyati vipratiṣedhena . \\
\noindent \textbf{(P\_6,1.135.1)} tatra dvirvacanam bhavati iti asya avakāśaḥ bibhidatuḥ , bibhiduḥ . \\
\noindent \textbf{(P\_6,1.135.1)} suṭaḥ avakāśaḥ saṃskartā , saṃskartum . \\
\noindent \textbf{(P\_6,1.135.1)} iha ubhayam prāpnoti sañcaskaratuḥ , sañcaskaruḥ . \\
\noindent \textbf{(P\_6,1.135.1)} suṭ bhavati vipratiṣedhena . \\
\noindent \textbf{(P\_6,1.135.1)} dvirvacanāt suṭ vipratiṣedhena iti cet dvirbhūte śabdāntarabhāvāt punaḥ prasaṅgaḥ . \\
\noindent \textbf{(P\_6,1.135.1)} dvirvacanāt suṭ vipratiṣedhena iti cet dvirbhūte śabdāntarasya akṛtaḥ suṭ iti punaḥ suṭ syāt . \\
\noindent \textbf{(P\_6,1.135.1)} dvirbhūte śabdāntarabhāvāt punaḥ prasaṅgaḥ iti cet dvirvacanam . \\
\noindent \textbf{(P\_6,1.135.1)} suṭi kṛte śabdāntarasya akṛtam dvirvacanam iti punaḥ dvirvacanam prāpnoti . \\
\noindent \textbf{(P\_6,1.135.1)} tathā ca anavasthā . \\
\noindent \textbf{(P\_6,1.135.1)} punaḥ suṭ punaḥ dvirvacanam iti cakrakam anavasthā prasajyeta . \\
\noindent \textbf{(P\_6,1.135.1)} na asti cakrakaprasaṅgaḥ . \\
\noindent \textbf{(P\_6,1.135.1)} na hi anavasthākāriṇā śāstreṇa bhavitavyam . \\
\noindent \textbf{(P\_6,1.135.1)} śāstrataḥ hi nāma vyavasthāt . \\
\noindent \textbf{(P\_6,1.135.1)} tatra suṭi kṛte dvirvacanam . \\
\noindent \textbf{(P\_6,1.135.1)} dvirvacanena avasthānam bhaviṣyati . \\
\noindent \textbf{(P\_6,1.135.1)} aḍvyavāye upasaṅkhyānam . \\
\noindent \textbf{(P\_6,1.135.1)} aḍvyavāye upasaṅkhyānam kartavyam . \\
\noindent \textbf{(P\_6,1.135.1)} samaskarot , samaskārṣīt . \\
\noindent \textbf{(P\_6,1.135.1)} abhyāsavyavāye ca . \\
\noindent \textbf{(P\_6,1.135.1)} abhyāsavyavāye ca upasaṅkhyānam kartavyam . \\
\noindent \textbf{(P\_6,1.135.1)} sañcaskaratuḥ , sañcaskaruḥ . \\
\noindent \textbf{(P\_6,1.135.1)} kim ucyate abhyāsavyavāye iti yadā idānīm eva uktam dvirvacanāt suṭ vipratiṣedhena iti . \\
\noindent \textbf{(P\_6,1.135.1)} avipratiṣedhaḥ vā bahiraṅgalakṣaṇatvāt . \\
\noindent \textbf{(P\_6,1.135.1)} avipratiṣedhaḥ vā punaḥ suṭaḥ . \\
\noindent \textbf{(P\_6,1.135.1)} kim kāraṇam . \\
\noindent \textbf{(P\_6,1.135.1)} bahiraṅgalakṣaṇatvāt . \\
\noindent \textbf{(P\_6,1.135.1)} bahiraṅgalakṣaṇaḥ suṭ . \\
\noindent \textbf{(P\_6,1.135.1)} antaraṅgam dvirvacanam . \\
\noindent \textbf{(P\_6,1.135.1)} asiddham bahiraṅgam antaraṅge . \\
\noindent \textbf{(P\_6,1.135.1)} evamartham eva tarhi kātpūrvagrahaṇam kartavyam kāt pūrvaḥ yathā syāt . \\
\noindent \textbf{(P\_6,1.135.1)} kriyamāṇe api vai kātpūrvagrahaṇe atra na sidhyati . \\
\noindent \textbf{(P\_6,1.135.1)} na hi ayam kātpūrvagrahaṇena śakyaḥ madhye praveśayitum . \\
\noindent \textbf{(P\_6,1.135.1)} kim kāraṇam . \\
\noindent \textbf{(P\_6,1.135.1)} ādiliṅgaḥ ayam kriyate karotiḥ ca kakārādiḥ dṛṣṭaḥ ca lpunaḥ ātideśikaḥ karotiḥ akakārādiḥ . \\
\noindent \textbf{(P\_6,1.135.1)} pākṣikaḥ ayam doṣaḥ . \\
\noindent \textbf{(P\_6,1.135.1)} katarasmin pakṣe . \\
\noindent \textbf{(P\_6,1.135.1)} suḍvidhau dvaitam bhavati . \\
\noindent \textbf{(P\_6,1.135.1)} aviśeṣeṇa vā vihitasya suṭaḥ kātpūrvagrahaṇam deśaprakḷptyartham syāt viśeṣeṇa vā vidhiḥ iti . \\
\noindent \textbf{(P\_6,1.135.1)} dvirvacanavidhau ca api dvaitam bhavati . \\
\noindent \textbf{(P\_6,1.135.1)} sthāne dvirvacanam syāt dviḥ prayogaḥ vā dvirvacanam iti . \\
\noindent \textbf{(P\_6,1.135.1)} tat yadā dviḥ prayogaḥ dvirvacanam aviśeṣeṇa vihitasya ca suṭaḥ kātpūrvagrahaṇam deśaprakḷptyartham tadā eṣaḥ doṣaḥ . \\
\noindent \textbf{(P\_6,1.135.1)} yadā hi sthāne dvirvacanam tadā yadi aviśeṣeṇa vihitasya suṭaḥ kātpūrvagrahaṇam deśaprakḷptyartham atha api viśeṣavidhiḥ na tadā doṣaḥ bhavati . \\
\noindent \textbf{(P\_6,1.135.1)} dviḥprayoge ca api dvirvacane na doṣaḥ . \\
\noindent \textbf{(P\_6,1.135.1)} samparibhyām iti na eṣā pañcamī . \\
\noindent \textbf{(P\_6,1.135.1)} kā tarhi . \\
\noindent \textbf{(P\_6,1.135.1)} tṛtīyā . \\
\noindent \textbf{(P\_6,1.135.1)} samparibhyām upasṛṣṭasya iti . \\
\noindent \textbf{(P\_6,1.135.1)} vyavahitaḥ ca api upasṛṣṭaḥ bhavati . \\
\noindent \textbf{(P\_6,1.135.1)} upadeśivadvacanam ca . \\
\noindent \textbf{(P\_6,1.135.1)} upadeśivadbhāvaḥ ca vaktavyaḥ . \\
\noindent \textbf{(P\_6,1.135.1)} kim prayojanam . \\
\noindent \textbf{(P\_6,1.135.1)} liṭiguṇacaṅidīrghapratiṣedhārtham . \\
\noindent \textbf{(P\_6,1.135.1)} liṭi guṇārtham caṅi dīrghapratiṣedhārtham . \\
\noindent \textbf{(P\_6,1.135.1)} liṭi guṇārtham tāvat . \\
\noindent \textbf{(P\_6,1.135.1)} sañcaskaratuḥ , sañcaskaruḥ . \\
\noindent \textbf{(P\_6,1.135.1)} caṅi dīrghapratiṣedhārtham ca . \\
\noindent \textbf{(P\_6,1.135.1)} samaciskarat . \\
\noindent \textbf{(P\_6,1.135.1)} liṭi guṇārthena tāvat na arthaḥ . \\
\noindent \textbf{(P\_6,1.135.1)} vakṣyati etat saṃyogādeḥ guṇavidhāne saṃyogopadhāgrahaṇam kṛñartham iti . \\
\noindent \textbf{(P\_6,1.135.1)} caṅi dīrghapratiṣedhena api na arthaḥ . \\
\noindent \textbf{(P\_6,1.135.1)} padam iti iyam bhagavataḥ kṛtrimā sañjñā . \\
\noindent \textbf{(P\_6,1.135.1)} yuktam iha draṣṭavyam . \\
\noindent \textbf{(P\_6,1.135.1)} kim antaraṅgam kim bahiraṅgam iti . \\
\noindent \textbf{(P\_6,1.135.1)} dhātūpasargayoḥ kāryam yat tat antaraṅgam . \\
\noindent \textbf{(P\_6,1.135.1)} kutaḥ etat . \\
\noindent \textbf{(P\_6,1.135.1)} pūrvam hi dhātuḥ upasargeṇa yujyate paścāt sādhanena . \\
\noindent \textbf{(P\_6,1.135.1)} na etat sāram . \\
\noindent \textbf{(P\_6,1.135.1)} pūrvam dhātuḥ sādhanena yujyate paścāt upasargeṇa . \\
\noindent \textbf{(P\_6,1.135.1)} sādhanam hi kriyām nirvartayati . \\
\noindent \textbf{(P\_6,1.135.1)} tām upasargaḥ viśinaṣti . \\
\noindent \textbf{(P\_6,1.135.1)} abhinirvṛttasya ca arthasya upasargeṇa viśeṣaḥ śakyam kartūm . \\
\noindent \textbf{(P\_6,1.135.1)} satyam evam etat . \\
\noindent \textbf{(P\_6,1.135.1)} yaḥ tu asau dhātūpasargayoḥ abhisambandhaḥ tam abhyantaram kṛtvā dhātuḥ sādhanena yujyate . \\
\noindent \textbf{(P\_6,1.135.1)} avaśyam ca etat evam vijñeyam . \\
\noindent \textbf{(P\_6,1.135.1)} yaḥ hi manyate pūrvam dhātuḥ sādhanena yujyate paścāt upasargeṇa iti tasya āsyate guruṇā iti akarmakaḥ upāsyate guruḥ iti kena sakarmakaḥ syāt . \\
\noindent \textbf{(P\_6,1.135.1)} evam kṛtvā suṭ sarvataḥ antaraṅgatarakaḥ bhavati kātpūrvagrahaṇam ca api śakyam akartum . \\
\noindent \textbf{(P\_6,1.135.2)} yadi punaḥ ayam suṭ kāt pūrvāntaḥ kriyeta . \\
\noindent \textbf{(P\_6,1.135.2)} kāt pūrvāntaḥ iti cet ruvidhipratiṣedhaḥ . \\
\noindent \textbf{(P\_6,1.135.2)} kāt pūrvāntaḥ iti cet kaḥ cit vidheyaḥ kaḥ cit patiṣedhyaḥ . \\
\noindent \textbf{(P\_6,1.135.2)} saṃskartā . \\
\noindent \textbf{(P\_6,1.135.2)} samaḥ vidheyaḥ suṭaḥ pratiṣedhyaḥ . \\
\noindent \textbf{(P\_6,1.135.2)} samaḥ tāvat na vidheyaḥ . \\
\noindent \textbf{(P\_6,1.135.2)} vakṣyati etat sampuṅkānām satvam ruvidhau hi aniṣṭaprasaṅgaḥ iti . \\
\noindent \textbf{(P\_6,1.135.2)} suṭaḥ ca api na pratiṣedhyaḥ . \\
\noindent \textbf{(P\_6,1.135.2)} samaḥ suṭi iti dvisakārakaḥ nirdeśaḥ : suṭi sakārādau iti . \\
\noindent \textbf{(P\_6,1.135.2)} atha vā padādiḥ kariyṣyate . \\
\noindent \textbf{(P\_6,1.135.2)} parādau iḍgrahaṇaprasaṅgaḥ . \\
\noindent \textbf{(P\_6,1.135.2)} yadi parādiḥ iḍguṇau prāpnutaḥ . \\
\noindent \textbf{(P\_6,1.135.2)} saṃskṛṣīṣṭa . \\
\noindent \textbf{(P\_6,1.135.2)} ṛtaḥ ca saṃyogādeḥ iti iṭ prāpnoti . \\
\noindent \textbf{(P\_6,1.135.2)} saṃskriyate . \\
\noindent \textbf{(P\_6,1.135.2)} guṇaḥ artisaṃyogādyoḥ iti guṇaḥ prāpnoti . \\
\noindent \textbf{(P\_6,1.135.2)} evam tarhi abhaktaḥ kariṣyate . \\
\noindent \textbf{(P\_6,1.135.2)} abhakte svaraḥ . \\
\noindent \textbf{(P\_6,1.135.2)} yadi abhaktaḥ svaraḥ na sidhyati . \\
\noindent \textbf{(P\_6,1.135.2)} saṃskaroti . \\
\noindent \textbf{(P\_6,1.135.2)} tiṅ atiṅaḥ iti nighātaḥ na prāpnoti . \\
\noindent \textbf{(P\_6,1.135.2)} nanu ca suṭ eva atiṅ . \\
\noindent \textbf{(P\_6,1.135.2)} na suṭaḥ parasya nighātena bhavitavyam . \\
\noindent \textbf{(P\_6,1.135.2)} kim kāraṇam . \\
\noindent \textbf{(P\_6,1.135.2)} nañivayuktam anyasadṛśādhikaraṇe . \\
\noindent \textbf{(P\_6,1.135.2)} tathā hi arthagatiḥ . \\
\noindent \textbf{(P\_6,1.135.2)} nañyuktam ivayuktam ca anyasmin tatsadṛśe kāryam vijñāyate . \\
\noindent \textbf{(P\_6,1.135.2)} tathā hi arthaḥ gamyate . \\
\noindent \textbf{(P\_6,1.135.2)} tat yathā abrāhmaṇam ānaya iti ukte brāhmaṇasadṛśam kṣatriyam ānayati . \\
\noindent \textbf{(P\_6,1.135.2)} na asau loṣṭam ānīyā kṛtī bhavati . \\
\noindent \textbf{(P\_6,1.135.2)} evam iha api atiṅ iti pratiṣedhāt anyasmāt atiṅsadṛśāt kāryam vijñāyate . \\
\noindent \textbf{(P\_6,1.135.2)} kim ca anyat atiṅ tiṅsadṛśam . \\
\noindent \textbf{(P\_6,1.135.2)} padam . \\
\noindent \textbf{(P\_6,1.142)} kirateḥ harṣajīvikākulāyakaraṇeṣu . \\
\noindent \textbf{(P\_6,1.142)} kirateḥ harṣajīvikākulāyakaraṇeṣu iti vaktavyam . \\
\noindent \textbf{(P\_6,1.142)} apaskirate vṛṣabhaḥ hṛṣṭaḥ . \\
\noindent \textbf{(P\_6,1.142)} apaskirate kukkuṭaḥ bhakṣārthī . \\
\noindent \textbf{(P\_6,1.142)} apaskirate śvā āśrayārthī . \\
\noindent \textbf{(P\_6,1.144)} kim idam sātatye iti . \\
\noindent \textbf{(P\_6,1.144)} santatabhāvaḥ sātatyam . \\
\noindent \textbf{(P\_6,1.144)} yadi evam sāntatye iti bhavitavyam . \\
\noindent \textbf{(P\_6,1.144)} samaḥ hitatatayoḥ vā lopaḥ . \\
\noindent \textbf{(P\_6,1.144)} samaḥ hitatatayoḥ vā lopaḥ vaktavyaḥ . \\
\noindent \textbf{(P\_6,1.144)} saṃhitam , sahitam , santatam , satatam . \\
\noindent \textbf{(P\_6,1.144)} samtumunoḥ kāme . \\
\noindent \textbf{(P\_6,1.144)} samtumunoḥ kāme lopaḥ vaktavyaḥ . \\
\noindent \textbf{(P\_6,1.144)} sakāmaḥ , bhoktukāmaḥ . \\
\noindent \textbf{(P\_6,1.144)} manasi ca iti vaktavyam . \\
\noindent \textbf{(P\_6,1.144)} samanāḥ , bhoktumanāḥ . \\
\noindent \textbf{(P\_6,1.144)} avaśyamaḥ kṛtye . \\
\noindent \textbf{(P\_6,1.144)} avaśyamaḥ kṛtye lopaḥ vaktavyaḥ . \\
\noindent \textbf{(P\_6,1.144)} avaśyabhāvyam . \\
\noindent \textbf{(P\_6,1.145)} idam atibahu kriyate sevite , asevite , pramāṇe iti . \\
\noindent \textbf{(P\_6,1.145)} sevitapramāṇayoḥ iti eva siddham . \\
\noindent \textbf{(P\_6,1.145)} kena idānīm asevite bhaviṣyati . \\
\noindent \textbf{(P\_6,1.145)} nañā sevitapratiṣedham vijñāsyāmaḥ . \\
\noindent \textbf{(P\_6,1.145)} na evam śakyam . \\
\noindent \textbf{(P\_6,1.145)} sevitaprasaṅge eva syāt . \\
\noindent \textbf{(P\_6,1.145)} asevite na syāt . \\
\noindent \textbf{(P\_6,1.145)} asevitagrahaṇe punaḥ kriyamāṇe bahuvrīhiḥ ayam vijñāsyate . \\
\noindent \textbf{(P\_6,1.145)} avidyamānasevite asevite iti . \\
\noindent \textbf{(P\_6,1.145)} tasmāt asevitagrahaṇam kartavyam . \\
\noindent \textbf{(P\_6,1.150)} viṣkiraḥ śakunau vikiraḥ vā . \\
\noindent \textbf{(P\_6,1.150)} viṣkiraḥ śakunau vikiraḥ vā iti vaktavyam . \\
\noindent \textbf{(P\_6,1.150)} śakunau vā iti hi ucyamāne śakunau vā syāt anyatra api nityam . \\
\noindent \textbf{(P\_6,1.150)} tat tarhi vaktavyam . \\
\noindent \textbf{(P\_6,1.150)} na vaktavyam . \\
\noindent \textbf{(P\_6,1.150)} na vāvacanena śakuniḥ abhisambadhyate . \\
\noindent \textbf{(P\_6,1.150)} kim tarhi . \\
\noindent \textbf{(P\_6,1.150)} nipātanam abhisambadhyate : viṣkiraḥ iti etat nipātanam śakunau vā nipātyate iti . \\
\noindent \textbf{(P\_6,1.147)} āścaryam adbhute . \\
\noindent \textbf{(P\_6,1.147)} āścaryam adbhute iti vaktavyam iha api yathā syāt . \\
\noindent \textbf{(P\_6,1.147)} āścaryam uccatā vṛkṣasya . \\
\noindent \textbf{(P\_6,1.147)} āścaryam nīlā dyauḥ . \\
\noindent \textbf{(P\_6,1.147)} āścaryam antarikṣe abandhanāni nakṣatrāṇi na patanti iti . \\
\noindent \textbf{(P\_6,1.147)} tat tarhi vaktavyam . \\
\noindent \textbf{(P\_6,1.147)} na vaktavyam . \\
\noindent \textbf{(P\_6,1.147)} anitye iti eva siddham . \\
\noindent \textbf{(P\_6,1.147)} iha tāvat āścaryam uccatā vṛkṣasya iti . \\
\noindent \textbf{(P\_6,1.147)} āścaryagrahaṇena na vṛkṣaḥ abhisambadhyate . \\
\noindent \textbf{(P\_6,1.147)} kim tarhi . \\
\noindent \textbf{(P\_6,1.147)} uccatā . \\
\noindent \textbf{(P\_6,1.147)} sā ca anityā . \\
\noindent \textbf{(P\_6,1.147)} āścaryam nīlā dyauḥ iti . \\
\noindent \textbf{(P\_6,1.147)} na āścaryagrahaṇena dyauḥ abhisambadhyate . \\
\noindent \textbf{(P\_6,1.147)} kim tarhi . \\
\noindent \textbf{(P\_6,1.147)} nīlatā . \\
\noindent \textbf{(P\_6,1.147)} sā ca anityā . \\
\noindent \textbf{(P\_6,1.147)} āścaryam antarikṣe abandhanāni nakṣatrāṇi na patanti iti . \\
\noindent \textbf{(P\_6,1.147)} na āścaryagrahaṇena nakṣatrāṇi abhisambadhyante . \\
\noindent \textbf{(P\_6,1.147)} kim tarhi . \\
\noindent \textbf{(P\_6,1.147)} patanakriyā . \\
\noindent \textbf{(P\_6,1.147)} sā ca anityā . \\
\noindent \textbf{(P\_6,1.147)} tatra anitye iti eva siddham . \\
\noindent \textbf{(P\_6,1.154)} maskarigrahaṇam śakyam akartum . \\
\noindent \textbf{(P\_6,1.154)} katham maskarī parivrājakaḥ iti . \\
\noindent \textbf{(P\_6,1.154)} ininā etat matvarthīyena siddham . \\
\noindent \textbf{(P\_6,1.154)} maskaraḥ asya asti . \\
\noindent \textbf{(P\_6,1.154)} na vai maskaraḥ asya asti iti maskarī parivrājakaḥ . \\
\noindent \textbf{(P\_6,1.154)} kim tarhi mā kṛta karmāṇi . \\
\noindent \textbf{(P\_6,1.154)} mā kṛta karmāṇi . \\
\noindent \textbf{(P\_6,1.154)} śāntiḥ vaḥ śreyasī iti āha . \\
\noindent \textbf{(P\_6,1.154)} ataḥ maskarī parivrājakaḥ . \\
\noindent \textbf{(P\_6,1.157)} avihitalakṣaṇaḥ suṭ pāraskaraprabhṛtiṣu draṣṭavyaḥ . \\
\noindent \textbf{(P\_6,1.157)} pāraskaraḥ deśaḥ . \\
\noindent \textbf{(P\_6,1.157)} kāraskaraḥ vṛkṣaḥ . \\
\noindent \textbf{(P\_6,1.157)} rathaspā nadī . \\
\noindent \textbf{(P\_6,1.157)} kiṣkindhā guhā . \\
\noindent \textbf{(P\_6,1.157)} kiṣkuḥ . \\
\noindent \textbf{(P\_6,1.157)} tadbṛhatoḥ karapatyoḥ coradevatayoḥ suṭ talopaḥ ca . \\
\noindent \textbf{(P\_6,1.157)} taskaraḥ , bṛhaspatiḥ . \\
\noindent \textbf{(P\_6,1.157)} prāyasya citticittayoḥ suṭ askāraḥ vā . \\
\noindent \textbf{(P\_6,1.157)} prāyaścittiḥ , prāyaścittam . \\
\noindent \textbf{(P\_6,1.158.1)} kim anudāttāni padāni bhavanti ekam padam varjayitvā . \\
\noindent \textbf{(P\_6,1.158.1)} na iti āha . \\
\noindent \textbf{(P\_6,1.158.1)} pade yeṣām udāttaprasaṅgaḥ anudāttāḥ bhavanti ekam acam varjayitvā . \\
\noindent \textbf{(P\_6,1.158.1)} saḥ tarhi tathā nirdeśaḥ kartavyaḥ : anudāttāḥ pade , anudāttāḥ padasya iti vā . \\
\noindent \textbf{(P\_6,1.158.1)} na kartavyaḥ . \\
\noindent \textbf{(P\_6,1.158.1)} anudāttam padam ekavarjam iti eva siddham . \\
\noindent \textbf{(P\_6,1.158.1)} katham . \\
\noindent \textbf{(P\_6,1.158.1)} matublopaḥ atra draṣṭavyaḥ . \\
\noindent \textbf{(P\_6,1.158.1)} tat yathā puṣyakāḥ eṣām puṣyakāḥ , kālakāḥ eṣām kālakāḥ iti . \\
\noindent \textbf{(P\_6,1.158.1)} atha vā akāraḥ matvarthīyaḥ . \\
\noindent \textbf{(P\_6,1.158.1)} tat yathā tundaḥ , ghāṭaḥ iti . \\
\noindent \textbf{(P\_6,1.158.2)} kimartham punaḥ idam ucyate . \\
\noindent \textbf{(P\_6,1.158.2)} āgamasya vikārasya prakṛteḥ pratyayasya ca pṛthak svaranivṛttyartham ekavarjam padasvaraḥ . \\
\noindent \textbf{(P\_6,1.158.2)} āgamasya . \\
\noindent \textbf{(P\_6,1.158.2)} caturanaḍuhoḥ ām udāttaḥ . \\
\noindent \textbf{(P\_6,1.158.2)} catvāraḥ , anaḍvāhaḥ . \\
\noindent \textbf{(P\_6,1.158.2)} vikārasya . \\
\noindent \textbf{(P\_6,1.158.2)} asthidadhisakthyakṣṇām anaṅ udāttaḥ . \\
\noindent \textbf{(P\_6,1.158.2)} asthnā , dadhnā . \\
\noindent \textbf{(P\_6,1.158.2)} prakṛteḥ . \\
\noindent \textbf{(P\_6,1.158.2)} gopāyati , dhūpāyati . \\
\noindent \textbf{(P\_6,1.158.2)} pratyayasya ca . \\
\noindent \textbf{(P\_6,1.158.2)} kartavyam , taittirīyaḥ . \\
\noindent \textbf{(P\_6,1.158.2)} eteṣām pade yugapat svaraḥ prāpnoti . \\
\noindent \textbf{(P\_6,1.158.2)} iṣyate ca ekasya syāt iti . \\
\noindent \textbf{(P\_6,1.158.2)} tat ca antareṇa yatnam na sidhyati iti anudāttam padam ekavarjam . \\
\noindent \textbf{(P\_6,1.158.2)} evamartham idam ucyate . \\
\noindent \textbf{(P\_6,1.158.2)} na etat asti prayojanam . \\
\noindent \textbf{(P\_6,1.158.2)} yaugapadyam tavai siddham . \\
\noindent \textbf{(P\_6,1.158.2)} yat ayam tavai ca antaḥ ca yugapat iti siddhe yaugapadye yaugapadyam śāsti tat jñāpayati ācāryaḥ na yugapat svaraḥ bhavati iti . \\
\noindent \textbf{(P\_6,1.158.2)} paryāyaḥ tarhi prāpnoti . \\
\noindent \textbf{(P\_6,1.158.2)} paryāyaḥ riktaśāsanāt . \\
\noindent \textbf{(P\_6,1.158.2)} yat ayam rikte vibhāṣā iti siddhe paryāye paryāyam śāsti tat jñāpayati ācāryaḥ na paryāyaḥ bhavati iti . \\
\noindent \textbf{(P\_6,1.158.2)} udātte jñāpakam tu etat . \\
\noindent \textbf{(P\_6,1.158.2)} etat udātte jñāpakam syāt . \\
\noindent \textbf{(P\_6,1.158.2)} svaritena samāviśet . \\
\noindent \textbf{(P\_6,1.158.2)} svaritena samāveśaḥ prāpnoti . \\
\noindent \textbf{(P\_6,1.158.2)} svarite api udāttaḥ asti . \\
\noindent \textbf{(P\_6,1.158.2)} tasmāt na arthaḥ anena yogena . \\
\noindent \textbf{(P\_6,1.158.3)} ārabhyamāṇe api etasmin yoge anudātte vipratiṣedhānupapattiḥ ekasmin yugapat sambhavāt . \\
\noindent \textbf{(P\_6,1.158.3)} anudātte vipratiṣedhaḥ na upapadyate . \\
\noindent \textbf{(P\_6,1.158.3)} paṭhiṣyati hi ācāryaḥ vipratiṣedham je dīrghāt bahvacaḥ iti . \\
\noindent \textbf{(P\_6,1.158.3)} saḥ vipratiṣedhaḥ na upapadyate . \\
\noindent \textbf{(P\_6,1.158.3)} kim kāraṇam . \\
\noindent \textbf{(P\_6,1.158.3)} ekasmin yugapat sambhavāt . \\
\noindent \textbf{(P\_6,1.158.3)} asati khalu sambhave vipratiṣedhaḥ bhavati asti ca sambhavaḥ yat ubhayam syāt . \\
\noindent \textbf{(P\_6,1.158.3)} katham sambhavaḥ yadā anudāttam padam ekavarjam iti ucyate . \\
\noindent \textbf{(P\_6,1.158.3)} tat iha na asti . \\
\noindent \textbf{(P\_6,1.158.3)} kim kāraṇam . \\
\noindent \textbf{(P\_6,1.158.3)} na anena udāttatvam pratiṣidhyate . \\
\noindent \textbf{(P\_6,1.158.3)} kim tarhi anudāttatvam anena kriyate asti ca sambhavaḥ yat ubhayoḥ ca udāttatvam syāt anyeṣām ca anudāttatvam . \\
\noindent \textbf{(P\_6,1.158.3)} yadi punaḥ ayam adhikāraḥ vijñāyeta . \\
\noindent \textbf{(P\_6,1.158.3)} kim kṛtam bhavati . \\
\noindent \textbf{(P\_6,1.158.3)} adhikāraḥ pratiyogam tasya anirdeśārthaḥ iti yoge yoge upatiṣṭhate . \\
\noindent \textbf{(P\_6,1.158.3)} je dīrghāntasya ādiḥ udāttaḥ bhavati . \\
\noindent \textbf{(P\_6,1.158.3)} upasthitam idam bhavati anudāttam padam ekavarjam iti . \\
\noindent \textbf{(P\_6,1.158.3)} antyāt pūrvam bahvacaḥ . \\
\noindent \textbf{(P\_6,1.158.3)} upasthitam idam bhavati anudāttam padam ekavarjam iti . \\
\noindent \textbf{(P\_6,1.158.3)} tatra pūrveṇa astu varjyamānatā pareṇa vā iti pareṇa bhaviṣyati paratvāt . \\
\noindent \textbf{(P\_6,1.158.3)} na evam śakyam . \\
\noindent \textbf{(P\_6,1.158.3)} ṣāṣthikaḥ ekaḥ svaraḥ saṅgṛhītaḥ syāt . \\
\noindent \textbf{(P\_6,1.158.3)} ye anye saptādhyāyyām svarāḥ te na saṅgṛhītāḥ syuḥ . \\
\noindent \textbf{(P\_6,1.158.3)} samānodare śayite o ca udāttaḥ . \\
\noindent \textbf{(P\_6,1.158.3)} asthidadhisakthyakṣṇām anaṅ udāttaḥ iti . \\
\noindent \textbf{(P\_6,1.158.3)} siddham tu ekānanudāttatvāt . \\
\noindent \textbf{(P\_6,1.158.3)} siddham etat . \\
\noindent \textbf{(P\_6,1.158.3)} katham . \\
\noindent \textbf{(P\_6,1.158.3)} ekānanudāttatvāt . \\
\noindent \textbf{(P\_6,1.158.3)} ekānanudāttam padam bhavati iti vaktavyam . \\
\noindent \textbf{(P\_6,1.158.3)} kim idam ananudāttatvāt iti . \\
\noindent \textbf{(P\_6,1.158.3)} na udāttaḥ anudāttaḥ . \\
\noindent \textbf{(P\_6,1.158.3)} na anudāttaḥ . \\
\noindent \textbf{(P\_6,1.158.3)} ananudāttaḥ . \\
\noindent \textbf{(P\_6,1.158.3)} ekaḥ ananudāttaḥ asmin tat idam ekānanudāttam . \\
\noindent \textbf{(P\_6,1.158.3)} ekānanudāttatvāt iti . \\
\noindent \textbf{(P\_6,1.158.3)} sidhyati . \\
\noindent \textbf{(P\_6,1.158.3)} sūtram tarhi bhidyate . \\
\noindent \textbf{(P\_6,1.158.3)} yathānyāsam eva astu . \\
\noindent \textbf{(P\_6,1.158.3)} nanu ca uktam anudātte vipratiṣedhānupapattiḥ ekasmin yugapat sambhavāt iti . \\
\noindent \textbf{(P\_6,1.158.3)} na eṣaḥ doṣaḥ . \\
\noindent \textbf{(P\_6,1.158.3)} paribhāṣā iyam . \\
\noindent \textbf{(P\_6,1.158.3)} kim kṛtam bhavati . \\
\noindent \textbf{(P\_6,1.158.3)} kāryakālam hi sañjñāparibhāṣam . \\
\noindent \textbf{(P\_6,1.158.3)} yatra kāryam tatra upasthitam idam draṣṭavyam . \\
\noindent \textbf{(P\_6,1.158.3)} je dīrghāntasya ādiḥ udāttaḥ bhavati . \\
\noindent \textbf{(P\_6,1.158.3)} upasthitam idam bhavati anudāttam padam ekavarjam iti . \\
\noindent \textbf{(P\_6,1.158.3)} antyāt pūrvam bahvacaḥ . \\
\noindent \textbf{(P\_6,1.158.3)} upasthitam idam bhavati anudāttam padam ekavarjam iti . \\
\noindent \textbf{(P\_6,1.158.3)} tatra pūrveṇa astu varjyamānatā pareṇa vā iti pareṇa bhaviṣyati paratvāt . \\
\noindent \textbf{(P\_6,1.158.3)} atha vā na idam pāribhāṣikānudāttasya grahaṇam . \\
\noindent \textbf{(P\_6,1.158.3)} kim tarhi . \\
\noindent \textbf{(P\_6,1.158.3)} anvarthagrahaṇam . \\
\noindent \textbf{(P\_6,1.158.3)} avidyamānodāttam anudāttam iti . \\
\noindent \textbf{(P\_6,1.158.3)} ekavarjam iti ca aprasiddhiḥ sandehāt . \\
\noindent \textbf{(P\_6,1.158.3)} ekavarjam iti ca aprasiddhiḥ . \\
\noindent \textbf{(P\_6,1.158.3)} kutaḥ sandehāt . \\
\noindent \textbf{(P\_6,1.158.3)} na jñāyate kaḥ ekaḥ varjayitavyaḥ iti . \\
\noindent \textbf{(P\_6,1.158.3)} siddham tu yasmin anudātte udāttavacanānarthakyam tadvarjam . \\
\noindent \textbf{(P\_6,1.158.3)} siddham etat . \\
\noindent \textbf{(P\_6,1.158.3)} katham . \\
\noindent \textbf{(P\_6,1.158.3)} yasmin anudātte udāttavacanam anarthakam syāt saḥ ekaḥ varjayitavyaḥ . \\
\noindent \textbf{(P\_6,1.158.3)} prakṛtipratyayayoḥ svarasya sāvakāśatvāt aprasiddhiḥ . \\
\noindent \textbf{(P\_6,1.158.3)} prakṛtipratyayayoḥ svarasya sāvakāśatvāt aprasiddhiḥ syāt . \\
\noindent \textbf{(P\_6,1.158.3)} prakṛtisvarasya avakāśaḥ yatra anudāttaḥ pratyayaḥ . \\
\noindent \textbf{(P\_6,1.158.3)} pacati , paṭhati . \\
\noindent \textbf{(P\_6,1.158.3)} pratyayayasvarasya avakāśaḥ yatra anudāttā prakṛtiḥ . \\
\noindent \textbf{(P\_6,1.158.3)} samatvam , simatvam . \\
\noindent \textbf{(P\_6,1.158.3)} iha ubhayam prāpnoti . \\
\noindent \textbf{(P\_6,1.158.3)} kartavyam, taittirīyaḥ . \\
\noindent \textbf{(P\_6,1.158.3)} vipratiṣedhāt pratyayasvaraḥ . \\
\noindent \textbf{(P\_6,1.158.3)} vipratiṣedhāt pratyayasvaraḥ bhaviṣyati . \\
\noindent \textbf{(P\_6,1.158.3)} na evam . \\
\noindent \textbf{(P\_6,1.158.3)} vipratiṣedhe param kāryam iti ucyate . \\
\noindent \textbf{(P\_6,1.158.3)} na paraḥ pratyayasvaraḥ . \\
\noindent \textbf{(P\_6,1.158.3)} na eṣaḥ doṣaḥ . \\
\noindent \textbf{(P\_6,1.158.3)} iṣṭavācī paraśabdaḥ . \\
\noindent \textbf{(P\_6,1.158.3)} vipratiṣedhe param yat iṣṭam tat bhavati . \\
\noindent \textbf{(P\_6,1.158.3)} vipratiṣedhāt pratyayasvaraḥ iti cet kāmyādiṣu citkaraṇam . \\
\noindent \textbf{(P\_6,1.158.3)} vipratiṣedhāt pratyayasvaraḥ iti cet kāmyādayaḥ citaḥ kartavyāḥ . \\
\noindent \textbf{(P\_6,1.158.3)} putrakāmyati , gopāyati , ṛtīyate . \\
\noindent \textbf{(P\_6,1.158.3)} na eṣaḥ doṣaḥ . \\
\noindent \textbf{(P\_6,1.158.3)} prakṛtisvaraḥ atra bādhakaḥ bhaviṣyati . \\
\noindent \textbf{(P\_6,1.158.3)} prakṛtisvare pratyayasvarābhāvaḥ . \\
\noindent \textbf{(P\_6,1.158.3)} prakṛtisvare pratyayasvarasya abhāvaḥ . \\
\noindent \textbf{(P\_6,1.158.3)} kartavyam, taittirīyaḥ . \\
\noindent \textbf{(P\_6,1.158.3)} siddham tu prakṛtisvarabalīyastvāt pratyayasvarabhāvaḥ . \\
\noindent \textbf{(P\_6,1.158.3)} siddham etat . \\
\noindent \textbf{(P\_6,1.158.3)} katham . \\
\noindent \textbf{(P\_6,1.158.3)} prakṛtisvarāt balīyastvāt pratyayasvarasya bhāvaḥ siddhaḥ . \\
\noindent \textbf{(P\_6,1.158.3)} katham . \\
\noindent \textbf{(P\_6,1.158.3)} prakṛtisvarāt pratyayasvaraḥ balīyān bhavati . \\
\noindent \textbf{(P\_6,1.158.4)} satiśiṣṭasvarabalīyastvam ca . \\
\noindent \textbf{(P\_6,1.158.4)} satiśiṣṭasvaraḥ balīyān bhavati iti vaktavyam . \\
\noindent \textbf{(P\_6,1.158.4)} tat ca anekapratyayasamāsārtham . \\
\noindent \textbf{(P\_6,1.158.4)} tat ca avaśyam satiśiṣṭasvarabalīyastvam vaktavyam . \\
\noindent \textbf{(P\_6,1.158.4)} kim prayojanam . \\
\noindent \textbf{(P\_6,1.158.4)} anekapratyayārtham anekasamāsārtham ca . \\
\noindent \textbf{(P\_6,1.158.4)} anekapratyayārtham tāvat . \\
\noindent \textbf{(P\_6,1.158.4)} aupagavaḥ . \\
\noindent \textbf{(P\_6,1.158.4)} prakṛtisvaram aṇsvaraḥ bādhate . \\
\noindent \textbf{(P\_6,1.158.4)} aupagavatvam . \\
\noindent \textbf{(P\_6,1.158.4)} tvasvaraḥ aṇsvaram bādhate . \\
\noindent \textbf{(P\_6,1.158.4)} aupagavatvakam . \\
\noindent \textbf{(P\_6,1.158.4)} tvasvaram kasvaraḥ bādhate . \\
\noindent \textbf{(P\_6,1.158.4)} anekasamāsārtham . \\
\noindent \textbf{(P\_6,1.158.4)} rājapuruṣaḥ , rājapuruṣaputraḥ , rājapuruṣaputrapuruṣaḥ . \\
\noindent \textbf{(P\_6,1.158.4)} yadi satiśiṣṭasvarabalīyastvam ucyate syādisvaraḥ sārvadhātukasvaram bādheta . \\
\noindent \textbf{(P\_6,1.158.4)} sunutaḥ , cinutaḥ . \\
\noindent \textbf{(P\_6,1.158.4)} syādisvarāprasaṅgaḥ ca tāseḥ parasya anudāttavacanāt . \\
\noindent \textbf{(P\_6,1.158.4)} syādisvarasya ca aprasaṅgaḥ . \\
\noindent \textbf{(P\_6,1.158.4)} kutaḥ . \\
\noindent \textbf{(P\_6,1.158.4)} tāseḥ parasya anudāttavacanāt . \\
\noindent \textbf{(P\_6,1.158.4)} yat ayam tāseḥ parasya lasārvadhātukasya anudāttatvam śāsti tat jñāpayati ācāryaḥ satiśiṣṭaḥ api vikaraṇasvaraḥ lasārvadhātukasvaram na bādhate . \\
\noindent \textbf{(P\_6,1.158.4)} śāstraparavipratiṣedhāniyamāt vā śabdavipratiṣedhāt siddham . \\
\noindent \textbf{(P\_6,1.158.4)} atha vā śāstraparavipratiṣedhe na sarvam iṣṭam saṅgṛhītam bhavati iti kṛtvā śabdavipratiṣedhaḥ vijñāsyate . \\
\noindent \textbf{(P\_6,1.158.4)} yadi śabdavipratiṣedhaḥ bhavati kāmyādayaḥ citaḥ kartavyāḥ . \\
\noindent \textbf{(P\_6,1.158.4)} putrakāmyati , gopāyati , ṛtīyate . \\
\noindent \textbf{(P\_6,1.158.4)} śabdavipratiṣedhaḥ nāma bhavati yatra ubhayoḥ yugapatprasaṅgaḥ na ca kāmyādiṣu yugapatprasaṅgaḥ . \\
\noindent \textbf{(P\_6,1.158.4)} vibhaktisvarāt nañsvaraḥ balīyān . \\
\noindent \textbf{(P\_6,1.158.4)} vibhaktisvarāt nañsvaraḥ balīyān iti vaktavyam . \\
\noindent \textbf{(P\_6,1.158.4)} vibhaktisvarasya avakāśaḥ . \\
\noindent \textbf{(P\_6,1.158.4)} tisraḥ tiṣṭhanti . \\
\noindent \textbf{(P\_6,1.158.4)} nañsvarasya avakāśaḥ . \\
\noindent \textbf{(P\_6,1.158.4)} abrāhmaṇaḥ , avṛṣalaḥ . \\
\noindent \textbf{(P\_6,1.158.4)} iha ubhayam prāpnoti . \\
\noindent \textbf{(P\_6,1.158.4)} atisraḥ . \\
\noindent \textbf{(P\_6,1.158.4)} nañsvaraḥ bhavati . \\
\noindent \textbf{(P\_6,1.158.4)} vibhaktinimittasvarāt ca . \\
\noindent \textbf{(P\_6,1.158.4)} vibhaktinimittasvarāt ca nañsvaraḥ balīyān iti vaktavyam . \\
\noindent \textbf{(P\_6,1.158.4)} vibhaktinimittasvarasya avakāśaḥ . \\
\noindent \textbf{(P\_6,1.158.4)} catvāraḥ , anaḍvāhaḥ . \\
\noindent \textbf{(P\_6,1.158.4)} nañsvarasya saḥ eva . \\
\noindent \textbf{(P\_6,1.158.4)} iha ubhayam prāpnoti . \\
\noindent \textbf{(P\_6,1.158.4)} acatvāraḥ . \\
\noindent \textbf{(P\_6,1.158.4)} ananaḍvāhaḥ . \\
\noindent \textbf{(P\_6,1.158.4)} yat ca upapadam kṛti nañ . \\
\noindent \textbf{(P\_6,1.158.4)} yat ca upapadam kṛti nañ tasya svaraḥ balīyān iti vaktavyam . \\
\noindent \textbf{(P\_6,1.158.4)} akaraṇiḥ hi te vṛṣala . \\
\noindent \textbf{(P\_6,1.158.4)} sahanirdiṣṭasya ca . \\
\noindent \textbf{(P\_6,1.158.4)} sahanirdiṣṭasya ca nañaḥ svaraḥ balīyān iti vaktavyam . \\
\noindent \textbf{(P\_6,1.158.4)} avyathī . \\
\noindent \textbf{(P\_6,1.159)} kimartham kṛṣateḥ vikṛtasya grahaṇam kriyate na kṛṣātvataḥ iti eva ucyeta . \\
\noindent \textbf{(P\_6,1.159)} yasya kṛṣeḥ vikaraṇe etat rūpam tasya yathā syāt . \\
\noindent \textbf{(P\_6,1.159)} iha mā bhūt . \\
\noindent \textbf{(P\_6,1.159)} halasya karṣaḥ iti . \\
\noindent \textbf{(P\_6,1.159)} atha kimartham matupā nirdeśaḥ kriyate na karṣāt iti eva ucyeta . \\
\noindent \textbf{(P\_6,1.159)} karṣāt iti iyati ucyamāne yatra eva ākārāt anantaraḥ ghañ asti tatra eva syāt : dāyaḥ , dhāyaḥ . \\
\noindent \textbf{(P\_6,1.159)} iha na syāt : pākaḥ , pāṭhaḥ . \\
\noindent \textbf{(P\_6,1.159)} na kva cit ākārāt anantaraḥ ghañ asti . \\
\noindent \textbf{(P\_6,1.159)} iha api dāyaḥ , dhāyaḥ iti yukā vyavadhānam . \\
\noindent \textbf{(P\_6,1.159)} evam api vihitaviṣeṣanam ākāragrahaṇam vijñāyeta . \\
\noindent \textbf{(P\_6,1.159)} ākārāt yaḥ vihitaḥ iti . \\
\noindent \textbf{(P\_6,1.159)} matubgrahaṇe punaḥ kriyamāṇe na doṣaḥ bhavati . \\
\noindent \textbf{(P\_6,1.161.1)} anudāttasya iti kimartham . \\
\noindent \textbf{(P\_6,1.161.1)} prāsaṅgam vahati prāsaṅgyaḥ . \\
\noindent \textbf{(P\_6,1.161.1)} udāttalope svaritodāttayoḥ abhāvāt anudāttagrahaṇānarthakyam . \\
\noindent \textbf{(P\_6,1.161.1)} udāttalope svaritodāttayoḥ abhāvāt anudāttagrahaṇam anarthakam . \\
\noindent \textbf{(P\_6,1.161.1)} ha hi kaḥ cit udāttaḥ udātte svarite vā lupyate . \\
\noindent \textbf{(P\_6,1.161.1)} sarvaḥ anudātte eva . \\
\noindent \textbf{(P\_6,1.161.1)} nan ca ayam udāttaḥ svarite lupyate . \\
\noindent \textbf{(P\_6,1.161.1)} prāsaṅgam vahati prāsaṅgyaḥ iti . \\
\noindent \textbf{(P\_6,1.161.1)} eṣaḥ api nighāte kṛte anudātte eva lupyate . \\
\noindent \textbf{(P\_6,1.161.1)} idam iha sampradhāryam . \\
\noindent \textbf{(P\_6,1.161.1)} nighātaḥ kriyatām lopaḥ iti . \\
\noindent \textbf{(P\_6,1.161.1)} kim atra kartavyam . \\
\noindent \textbf{(P\_6,1.161.1)} paratvāt lopaḥ . \\
\noindent \textbf{(P\_6,1.161.1)} evam tarhi ayam adya nighātasvaraḥ sarvasvarāṇām apavādaḥ . \\
\noindent \textbf{(P\_6,1.161.1)} na ca apavādaviṣaye utsargaḥ bhiniviśate . \\
\noindent \textbf{(P\_6,1.161.1)} pūrvam hi apavādāḥ abhiniviśante paścāt utsargāḥ . \\
\noindent \textbf{(P\_6,1.161.1)} prakalpya vā apavādaviṣayam tataḥ utsargaḥ abhiniviśate . \\
\noindent \textbf{(P\_6,1.161.1)} tat na tāvat atra kadā cit thāthādisvaraḥ bhavati . \\
\noindent \textbf{(P\_6,1.161.1)} apavādam nighātam pratīkṣate . \\
\noindent \textbf{(P\_6,1.161.1)} tatra nighātaḥ kriyatām lopaḥ iti yadi api paratvāt lopaḥ saḥ asau avidyamānodāttaḥ anudāttaḥ lupyate . \\
\noindent \textbf{(P\_6,1.161.2)} kim punaḥ anudāttasya antaḥ udāttaḥ bhavati āhosvit ādiḥ . \\
\noindent \textbf{(P\_6,1.161.2)} kaḥ ca atra viśeṣaḥ . \\
\noindent \textbf{(P\_6,1.161.2)} antaḥ iti cet śnamksayuṣmadasmadidaṅkiṃlopeṣu svaraḥ . \\
\noindent \textbf{(P\_6,1.161.2)} antaḥ iti cet śnamksayuṣmadasmadidaṅkiṃlopeṣu svaraḥ na sidhyati . \\
\noindent \textbf{(P\_6,1.161.2)} śnam . \\
\noindent \textbf{(P\_6,1.161.2)} vindate , khindate . \\
\noindent \textbf{(P\_6,1.161.2)} śnam . \\
\noindent \textbf{(P\_6,1.161.2)} ksa . \\
\noindent \textbf{(P\_6,1.161.2)} mā hi dhukṣātām . \\
\noindent \textbf{(P\_6,1.161.2)} mā hi dhukṣāthām . \\
\noindent \textbf{(P\_6,1.161.2)} ksa . \\
\noindent \textbf{(P\_6,1.161.2)} yuṣmadasmad . \\
\noindent \textbf{(P\_6,1.161.2)} yuṣmabhyam , asmabhyam . \\
\noindent \textbf{(P\_6,1.161.2)} idaṅkiṃlopaḥ . \\
\noindent \textbf{(P\_6,1.161.2)} iyān , kiyān . \\
\noindent \textbf{(P\_6,1.161.2)} astu tarhi ādiḥ . \\
\noindent \textbf{(P\_6,1.161.2)} ādiḥ iti cet indhīta dvayam iti antaḥ . \\
\noindent \textbf{(P\_6,1.161.2)} ādiḥ iti cet indhīta dvayam iti antodāttatvam na sidhyati . \\
\noindent \textbf{(P\_6,1.161.2)} indhīta . \\
\noindent \textbf{(P\_6,1.161.2)} dvayam , trayam . \\
\noindent \textbf{(P\_6,1.161.2)} ādau siddham . \\
\noindent \textbf{(P\_6,1.161.2)} astu tarhi ādiḥ udāttaḥ bhavati iti . \\
\noindent \textbf{(P\_6,1.161.2)} nanu ca uktam ādiḥ iti cet indhīta dvayam iti antaḥ iti . \\
\noindent \textbf{(P\_6,1.161.2)} vidīndhikhidibhyaḥ ca lasārvadhātukānudāttapratiṣedhāt liṅi siddham . \\
\noindent \textbf{(P\_6,1.161.2)} vidīndhikhidibhyaḥ ca lasārvadhātukānudāttatvam liṅi na iti vaktavyam . \\
\noindent \textbf{(P\_6,1.161.2)} liṅgrahaṇena na arthaḥ . \\
\noindent \textbf{(P\_6,1.161.2)} aviśeṣeṇa ikhidibhyaḥ ca lasārvadhātukānudāttapratiṣedhāt liṅi siddham . \\
\noindent \textbf{(P\_6,1.161.2)} vidīndhikhidibhyaḥ ca lasārvadhātukānudāttatvam na iti eva . \\
\noindent \textbf{(P\_6,1.161.2)} idam api siddham bhavati . \\
\noindent \textbf{(P\_6,1.161.2)} vindate , khindate . \\
\noindent \textbf{(P\_6,1.161.2)} ayaci katham . \\
\noindent \textbf{(P\_6,1.161.2)} ayaci citkaraṇāt . \\
\noindent \textbf{(P\_6,1.161.2)} ayaci citkaraṇasāmarthyāt antodāttatvam bhaviṣyati . \\
\noindent \textbf{(P\_6,1.162)} kim dhātoḥ antaḥ udāttaḥ bhavati āhosvit ādiḥ iti . \\
\noindent \textbf{(P\_6,1.162)} kaḥ ca atra viśeṣaḥ . \\
\noindent \textbf{(P\_6,1.162)} dhātoḥ antaḥ iti cet anudātte ca bagrahaṇam . \\
\noindent \textbf{(P\_6,1.162)} dhātoḥ antaḥ iti cet anudātte ca bagrahaṇam kartavyam . \\
\noindent \textbf{(P\_6,1.162)} abhyastānām ādiḥ anudātte ca iti vaktavyam . \\
\noindent \textbf{(P\_6,1.162)} bagrahaṇam ca kartavyam . \\
\noindent \textbf{(P\_6,1.162)} bāntaḥ ca pibiḥ ādyudāttaḥ bhavati iti vaktavyam . \\
\noindent \textbf{(P\_6,1.162)} pibati . \\
\noindent \textbf{(P\_6,1.162)} san ca nit . \\
\noindent \textbf{(P\_6,1.162)} san ca nit kartavyaḥ . \\
\noindent \textbf{(P\_6,1.162)} kim prayojanam . \\
\noindent \textbf{(P\_6,1.162)} cikīrṣati jihīrṣati . \\
\noindent \textbf{(P\_6,1.162)} niti iti ādyudāttatvam yathā syāt . \\
\noindent \textbf{(P\_6,1.162)} astu tarhi ādiḥ . \\
\noindent \textbf{(P\_6,1.162)} ādau ūrṇapratyayadhātuṣu antodāttatvam . \\
\noindent \textbf{(P\_6,1.162)} ādau ūrṇapratyayadhātuṣu antodāttatvam na sidhyati . \\
\noindent \textbf{(P\_6,1.162)} ūrṇoti . \\
\noindent \textbf{(P\_6,1.162)} ūrṇu . \\
\noindent \textbf{(P\_6,1.162)} pratyayadhātu . \\
\noindent \textbf{(P\_6,1.162)} gopāyati , dhūpāyati , ṛtīyate . \\
\noindent \textbf{(P\_6,1.162)} antodāttavacanāt siddham . \\
\noindent \textbf{(P\_6,1.162)} astu tarhi antodāttaḥ bhavati iti . \\
\noindent \textbf{(P\_6,1.162)} nanu ca uktam dhātoḥ antaḥ iti cet anudātte ca bagrahaṇam kartavyam iti . \\
\noindent \textbf{(P\_6,1.162)} yat tāvat ucyate . \\
\noindent \textbf{(P\_6,1.162)} anudātte ca grahaṇam kartavyam iti . \\
\noindent \textbf{(P\_6,1.162)} kriyate nyāse eva . \\
\noindent \textbf{(P\_6,1.162)} abhyastānām ādiḥ anudātte ca iti . \\
\noindent \textbf{(P\_6,1.162)} bagrahaṇam kartavyam iti . \\
\noindent \textbf{(P\_6,1.162)} pibau nipātanāt . \\
\noindent \textbf{(P\_6,1.162)} pibau ādyudāttanipātanam kriyate . \\
\noindent \textbf{(P\_6,1.162)} saḥ nipātanasvaraḥ prakṛtisvarasya bādhakaḥ bhaviṣyati . \\
\noindent \textbf{(P\_6,1.162)} san ca nit kartavyaḥ iti . \\
\noindent \textbf{(P\_6,1.162)} avaśyam sanaḥ viśeṣaṇārthaḥ nakāraḥ kartavyaḥ . \\
\noindent \textbf{(P\_6,1.162)} kva viśeṣaṇārthena arthaḥ . \\
\noindent \textbf{(P\_6,1.162)} sanyaṅoḥ iti . \\
\noindent \textbf{(P\_6,1.162)} sayaṅoḥ iti iyati ucyamāne haṃsaḥ , vatsaḥ , atra api prāpnoti . \\
\noindent \textbf{(P\_6,1.162)} arthavadgrahaṇe na anarthakasya iti evam na bhaviṣyati . \\
\noindent \textbf{(P\_6,1.162)} iha api tarhi na prāpnoti . \\
\noindent \textbf{(P\_6,1.162)} jugupsate , mīmāṃsate iti . \\
\noindent \textbf{(P\_6,1.162)} arthavān eṣaḥ . \\
\noindent \textbf{(P\_6,1.162)} na vai kaḥ cit arthaḥ ādiśyate . \\
\noindent \textbf{(P\_6,1.162)} yadi api kaḥ cit arthaḥ na ādiśyate anirdiṣṭārthāḥ svārthe bhavanti iti antataḥ svārthe bhaviṣyati . \\
\noindent \textbf{(P\_6,1.162)} kaḥ ca asya svārthaḥ . \\
\noindent \textbf{(P\_6,1.162)} prakṛtyarthaḥ . \\
\noindent \textbf{(P\_6,1.162)} iha api prāpnoti . \\
\noindent \textbf{(P\_6,1.162)} haṃsaḥ , vatsaḥ iti . \\
\noindent \textbf{(P\_6,1.162)} uṇādayaḥ avyutpannāni prātipadikāni . \\
\noindent \textbf{(P\_6,1.162)} saḥ eṣaḥ ananyārthaḥ nakāraḥ kartavyaḥ . \\
\noindent \textbf{(P\_6,1.162)} na kartavyaḥ . \\
\noindent \textbf{(P\_6,1.162)} kriyate nyāse eva . \\
\noindent \textbf{(P\_6,1.162)} atha vā dhātoḥ iti vartate . \\
\noindent \textbf{(P\_6,1.162)} dhātoḥ saśabdāntasya dve bhavataḥ iti . \\
\noindent \textbf{(P\_6,1.163)} citaḥ saprakṛteḥ bahvakajartham . \\
\noindent \textbf{(P\_6,1.163)} citaḥ saprakṛteḥ iti vaktavyam . \\
\noindent \textbf{(P\_6,1.163)} kim prayojanam . \\
\noindent \textbf{(P\_6,1.163)} bahvakajartham . \\
\noindent \textbf{(P\_6,1.163)} bahujartham akajartham ca . \\
\noindent \textbf{(P\_6,1.163)} bahujartham tāvat . \\
\noindent \textbf{(P\_6,1.163)} bahubhuktam , bahukṛtam . \\
\noindent \textbf{(P\_6,1.163)} akajartham . \\
\noindent \textbf{(P\_6,1.163)} sarvakaiḥ , viśvakaiḥ , uccakaiḥ , nīcakaiḥ , sarvake , viśvake . \\
\noindent \textbf{(P\_6,1.163)} tat tarhi vaktavyam . \\
\noindent \textbf{(P\_6,1.163)} na vaktavyam . \\
\noindent \textbf{(P\_6,1.163)} matublopaḥ atra draṣṭavyaḥ . \\
\noindent \textbf{(P\_6,1.163)} tat yathā puṣyakāḥ eṣām puṣyakāḥ kālakāḥ eṣām kālakāḥ iti . \\
\noindent \textbf{(P\_6,1.163)} atha vā akāraḥ matvarthīyaḥ . \\
\noindent \textbf{(P\_6,1.163)} tat yathā tundaḥ , ghāṭaḥ iti . \\
\noindent \textbf{(P\_6,1.163)} pūrvasūtranirdeśaḥ ca citvān citaḥ iti . \\
\noindent \textbf{(P\_6,1.166)} jasaḥ iti kimartham . \\
\noindent \textbf{(P\_6,1.166)} tisṛkā . \\
\noindent \textbf{(P\_6,1.166)} tisṛbhyaḥ jasgrahaṇānarthakyam anyatra abhāvāt . \\
\noindent \textbf{(P\_6,1.166)} tisṛbhyaḥ jasgrahaṇam anarthakam . \\
\noindent \textbf{(P\_6,1.166)} kim kāraṇam . \\
\noindent \textbf{(P\_6,1.166)} anyatra abhāvāt . \\
\noindent \textbf{(P\_6,1.166)} na hi anyat tisṛśabdāt antodāttatvam prayojayati anyat ataḥ jasaḥ . \\
\noindent \textbf{(P\_6,1.166)} kim kāraṇam . \\
\noindent \textbf{(P\_6,1.166)} bahuvacanaviṣayaḥ eva tisṛśabdaḥ . \\
\noindent \textbf{(P\_6,1.166)} tena ekavacanadvivacane na staḥ . \\
\noindent \textbf{(P\_6,1.166)} śasi bhavitavyam udāttayaṇaḥ halpūrvāt iti . \\
\noindent \textbf{(P\_6,1.166)} anyāḥ sarvāḥ halādayaḥ vibhaktayaḥ . \\
\noindent \textbf{(P\_6,1.166)} tatra ṣaṭtricaturbhyaḥ halādiḥ jhali upottamam iti anena svareṇa bhavitavyam . \\
\noindent \textbf{(P\_6,1.166)} tatra antareṇa jasaḥ grahaṇam jasaḥ eva bhaviṣyati . \\
\noindent \textbf{(P\_6,1.166)} nanu ca idānīm eva udāhṛtam tisṛkā iti . \\
\noindent \textbf{(P\_6,1.166)} nitsvaraḥ atra bādhakaḥ bhaviṣyati . \\
\noindent \textbf{(P\_6,1.166)} na aprāpte anyasvare tisṛsvaraḥ ārabhyate . \\
\noindent \textbf{(P\_6,1.166)} saḥ yathā eva anudāttau suppitau iti etam svaram bādhate evam nitsvaram api bādheta . \\
\noindent \textbf{(P\_6,1.166)} na eṣaḥ doṣaḥ . \\
\noindent \textbf{(P\_6,1.166)} yena na aprāpte tasya bādhanam bhavati . \\
\noindent \textbf{(P\_6,1.166)} na ca aprāpte anudāttau suppitau iti etasmin tisṛsvaraḥ ārabhyate . \\
\noindent \textbf{(P\_6,1.166)} nitsvaraḥ punaḥ prāpte ca aprāpte ca . \\
\noindent \textbf{(P\_6,1.166)} atha vā madhye apavādāḥ pūrvān vidhīn bādhante iti evam tisṛsvaraḥ anudāttau suppitau iti svaram bādhiṣyate nitsvaram na bādhiṣyate . \\
\noindent \textbf{(P\_6,1.166)} upasamastārtham eke jasaḥ grahaṇam icchanti : atitisrau , atitisraḥ . \\
\noindent \textbf{(P\_6,1.167)} śasi striyām pratiṣedhaḥ vaktavyaḥ . \\
\noindent \textbf{(P\_6,1.167)} catasraḥ paśya . \\
\noindent \textbf{(P\_6,1.167)} caturaḥ śasi striyām apratiṣedhaḥ ādyudāttanipātanāt . \\
\noindent \textbf{(P\_6,1.167)} caturaḥ śasi striyām apratiṣedhaḥ . \\
\noindent \textbf{(P\_6,1.167)} anarthakaḥ pratiṣedhaḥ apratiṣedhaḥ . \\
\noindent \textbf{(P\_6,1.167)} śasi svaraḥ kasmāt na bhavati . \\
\noindent \textbf{(P\_6,1.167)} ādyudāttanipātanāt . \\
\noindent \textbf{(P\_6,1.167)} ādyudāttanipātanam kariṣyate . \\
\noindent \textbf{(P\_6,1.167)} saḥ nipātanasvaraḥ śasi svarasya bādhakaḥ bhaviṣyati . \\
\noindent \textbf{(P\_6,1.167)} evam api upadeśivadbhāvaḥ vaktavyaḥ . \\
\noindent \textbf{(P\_6,1.167)} yathā eva nipātanasvaraḥ śasi svaram bādhate evam vibhaktisvaram api bādheta catasṛṇam iti . \\
\noindent \textbf{(P\_6,1.167)} vibhaktisvarabhāvaḥ ca halādigrahaṇāt . \\
\noindent \textbf{(P\_6,1.167)} vibhaktisvarabhāvaḥ ca siddhaḥ . \\
\noindent \textbf{(P\_6,1.167)} kutaḥ . \\
\noindent \textbf{(P\_6,1.167)} halādigrahaṇāt . \\
\noindent \textbf{(P\_6,1.167)} yat ayam ṣaṭtricaturbhyaḥ halādiḥ iti halādigrahaṇam karoti tat jñāpayati ācāryaḥ na nipātanasvaraḥ vibhaktisvaram bādhate iti . \\
\noindent \textbf{(P\_6,1.167)} katham kṛtvā jñāpakam . \\
\noindent \textbf{(P\_6,1.167)} ādyudāttanipātane hi halādigrahaṇānarthakyam . \\
\noindent \textbf{(P\_6,1.167)} ādyudāttanipātane hi sati halādigrahaṇam anarthakam syāt . \\
\noindent \textbf{(P\_6,1.167)} na hi anyat halādigrahaṇam prayojayati anyat ataḥ catasṛśabdāt . \\
\noindent \textbf{(P\_6,1.167)} ṣaṭsañjñāḥ tāvat na prayojayanti . \\
\noindent \textbf{(P\_6,1.167)} kim kāraṇam . \\
\noindent \textbf{(P\_6,1.167)} bahuvacanaviṣayatvāt . \\
\noindent \textbf{(P\_6,1.167)} tena dvivacanaikavacane na staḥ . \\
\noindent \textbf{(P\_6,1.167)} jaśśasī ca atra lupyete . \\
\noindent \textbf{(P\_6,1.167)} anyāḥ sarvāḥ halādayaḥ vibhaktayaḥ . \\
\noindent \textbf{(P\_6,1.167)} triśabdaḥ ca api na prayojayati . \\
\noindent \textbf{(P\_6,1.167)} kim kāraṇam . \\
\noindent \textbf{(P\_6,1.167)} bahuvacanaviṣayatvāt . \\
\noindent \textbf{(P\_6,1.167)} tena dvivacanaikavacane na staḥ . \\
\noindent \textbf{(P\_6,1.167)} asarvanāmasthānam iti vacanāt jasi na bhaviṣyati . \\
\noindent \textbf{(P\_6,1.167)} śasi bhavitavyam ekādeśe udāttena udāttaḥ iti . \\
\noindent \textbf{(P\_6,1.167)} anyāḥ sarvāḥ halādayaḥ vibhaktayaḥ . \\
\noindent \textbf{(P\_6,1.167)} tisṛśabdaḥ ca api na prayojayati . \\
\noindent \textbf{(P\_6,1.167)} kim kāraṇam . \\
\noindent \textbf{(P\_6,1.167)} bahuvacanaviṣayatvāt . \\
\noindent \textbf{(P\_6,1.167)} tena dvivacanaikavacane na staḥ . \\
\noindent \textbf{(P\_6,1.167)} asarvanāmasthānam iti vacanāt jasi na bhavitavyam . \\
\noindent \textbf{(P\_6,1.167)} śasi bhavitavyam udāttayaṇaḥ halpūrvāt iti . \\
\noindent \textbf{(P\_6,1.167)} anyāḥ sarvāḥ halādayaḥ vibhaktayaḥ . \\
\noindent \textbf{(P\_6,1.167)} catuḥśabdaḥ tisṛśabdaḥ ca api na prayojayati . \\
\noindent \textbf{(P\_6,1.167)} kim kāraṇam . \\
\noindent \textbf{(P\_6,1.167)} bahuvacanaviṣayatvāt . \\
\noindent \textbf{(P\_6,1.167)} tena dvivacanaikavacane na staḥ . \\
\noindent \textbf{(P\_6,1.167)} asarvanāmasthānam iti vacanāt jasi na bhavitavyam . \\
\noindent \textbf{(P\_6,1.167)} śasi bhavitavyam caturaḥ śasi iti . \\
\noindent \textbf{(P\_6,1.167)} anyāḥ sarvāḥ halādayaḥ vibhaktayaḥ . \\
\noindent \textbf{(P\_6,1.167)} tatra catasṛśabdāt ekasmāt śas asarvanāmasthānam ajādiḥ vibhaktiḥ asti . \\
\noindent \textbf{(P\_6,1.167)} yadi ca atra nipātanasvaraḥ syāt halādigrahaṇam anarthakam syāt . \\
\noindent \textbf{(P\_6,1.167)} na eva vā punaḥ atra śasisvaraḥ prāpnoti . \\
\noindent \textbf{(P\_6,1.167)} kim kāraṇam . \\
\noindent \textbf{(P\_6,1.167)} yaṇādeśe kṛte śasaḥ pūrvaḥ udāttabhāvī na asti iti kṛtvā . \\
\noindent \textbf{(P\_6,1.167)} avaśiṣṭasya tarhi prāpnoti . \\
\noindent \textbf{(P\_6,1.167)} ṛkāreṇa vyavahitatvāt na bhaviṣyati . \\
\noindent \textbf{(P\_6,1.167)} yaṇādeśe kṛte na asti vyavadhānam . \\
\noindent \textbf{(P\_6,1.167)} sthānivadbhāvāt vyavadhānam eva . \\
\noindent \textbf{(P\_6,1.167)} pratiṣidhyate atra sthānivadbhāvaḥ svaravidhim prati na sthānivat bhavati iti . \\
\noindent \textbf{(P\_6,1.167)} na eṣaḥ asti pratiṣedhaḥ . \\
\noindent \textbf{(P\_6,1.167)} uktam etat pratiṣedhe svaradīrghayalopeṣu lopādādeśaḥ na sthānivat iti . \\
\noindent \textbf{(P\_6,1.168.1)} sau iti kim idam prathamaikavacanasya grahaṇam āhosvit saptamībahuvacanasya . \\
\noindent \textbf{(P\_6,1.168.1)} kutaḥ sandehaḥ . \\
\noindent \textbf{(P\_6,1.168.1)} samānaḥ nirdeśaḥ . \\
\noindent \textbf{(P\_6,1.168.1)} saptamībahuvacanasya grahaṇam . \\
\noindent \textbf{(P\_6,1.168.1)} katham jñāyate . \\
\noindent \textbf{(P\_6,1.168.1)} yat ayam na gośvansāvavarṇa iti gośunoḥ pratiṣedham śāsti . \\
\noindent \textbf{(P\_6,1.168.1)} katham kṛtvā jñāpakam . \\
\noindent \textbf{(P\_6,1.168.1)} yadi prathamaikavacanasya grahaṇam syāt gośunoḥ pratiṣedhavacanam anarthakam syāt . \\
\noindent \textbf{(P\_6,1.168.1)} nanu ca arthasiddhiḥ eva eṣā . \\
\noindent \textbf{(P\_6,1.168.1)} anugṛhītāḥ smaḥ yaiḥ asmābhiḥ prathamaikavacanam āsthāya gośunoḥ pratiṣedhaḥ na vaktavyaḥ bhavati . \\
\noindent \textbf{(P\_6,1.168.1)} bhavet pratiṣedhaḥ na vaktavyaḥ doṣāḥ tu bhavanti . \\
\noindent \textbf{(P\_6,1.168.1)} tatra kaḥ doṣaḥ . \\
\noindent \textbf{(P\_6,1.168.1)} svinā khinā . \\
\noindent \textbf{(P\_6,1.168.1)} antodāttatvam na prāpnoti . \\
\noindent \textbf{(P\_6,1.168.1)} svinkhinau na staḥ . \\
\noindent \textbf{(P\_6,1.168.1)} uktam etat ekākṣarāt kṛtaḥ jāteḥ saptamyām ca na tau smṛtau . \\
\noindent \textbf{(P\_6,1.168.1)} svavān , khavān iti eva bhavitavyam . \\
\noindent \textbf{(P\_6,1.168.1)} iha tarhi yādbhyām , yābhiḥ iti na sidhyati . \\
\noindent \textbf{(P\_6,1.168.1)} tasmāt saptamībahuvacanasya grahaṇam . \\
\noindent \textbf{(P\_6,1.168.2)} sau ekācaḥ udāttatve tvanmadoḥ pratiṣedhaḥ . \\
\noindent \textbf{(P\_6,1.168.2)} sau ekācaḥ udāttatve tvanmadoḥ pratiṣedhaḥ vaktavyaḥ . \\
\noindent \textbf{(P\_6,1.168.2)} tvayā mayā . \\
\noindent \textbf{(P\_6,1.168.2)} siddham tu yasmāt tṛtīyādiḥ tasya abhāvāt sau . \\
\noindent \textbf{(P\_6,1.168.2)} siddham etat . \\
\noindent \textbf{(P\_6,1.168.2)} katham . \\
\noindent \textbf{(P\_6,1.168.2)} yasmāt atra tṛtīyādiḥ vibhaktiḥ na tat sau asti . \\
\noindent \textbf{(P\_6,1.168.2)} yadi api etat sau na asti prakṛtiḥ tu asya sau asti . \\
\noindent \textbf{(P\_6,1.168.2)} prakṛteḥ ca anekāctvāt . \\
\noindent \textbf{(P\_6,1.168.2)} yadi api tasya prakṛtiḥ asti sau anekāc tu sā bhavati . \\
\noindent \textbf{(P\_6,1.169)} uttarapadagrahaṇam kimartham . \\
\noindent \textbf{(P\_6,1.169)} yathā ekājgrahaṇam uttarapadaviśeṣaṇam vijñāyeta . \\
\noindent \textbf{(P\_6,1.169)} ekācaḥ uttarapadāt iti . \\
\noindent \textbf{(P\_6,1.169)} atha akriyamāṇe uttarapadagrahaṇe kasya ekājgrahaṇam viśeṣaṇam syāt . \\
\noindent \textbf{(P\_6,1.169)} samāsaviśeṣaṇam . \\
\noindent \textbf{(P\_6,1.169)} asti ca idānīm kaḥ cit ekāc samāsaḥ yadarthaḥ vidhiḥ syāt . \\
\noindent \textbf{(P\_6,1.169)} asti iti āha : śunaḥ ūrk : śvork , śvorjā , śvorje iti . \\
\noindent \textbf{(P\_6,1.171)} padādiṣu nicantāni prayojayanti . \\
\noindent \textbf{(P\_6,1.171)} anyāni padādīni udāttanivṛttisvareṇa siddhāni . \\
\noindent \textbf{(P\_6,1.171)} ūṭhi upadhāgrahaṇam antyapratiṣedhārtham . ūṭhi upadhāgrahaṇam kartavyam . \\
\noindent \textbf{(P\_6,1.171)} kim prayojanam . \\
\noindent \textbf{(P\_6,1.171)} antyapratiṣedhārtham . \\
\noindent \textbf{(P\_6,1.171)} antyasya mā bhūt . \\
\noindent \textbf{(P\_6,1.171)} akṣadyuvā , akṣadyuve . \\
\noindent \textbf{(P\_6,1.172)} dīrghagrahaṇam kimartham . \\
\noindent \textbf{(P\_6,1.172)} aṣṭasu prakrameṣu brāhmaṇaḥ ādadhīta . \\
\noindent \textbf{(P\_6,1.172)} dīrghāt iti śakyam akartum . \\
\noindent \textbf{(P\_6,1.172)} kasmāt na bhavati aṣṭasu prakrameṣu brāhmaṇaḥ ādadhīta iti . \\
\noindent \textbf{(P\_6,1.172)} ṣaṭsvaraḥ bādhakaḥ bhaviṣyati . \\
\noindent \textbf{(P\_6,1.172)} na aprāpte ṣaṭsvare aṣṭanaḥ svaraḥ ārabhyate . \\
\noindent \textbf{(P\_6,1.172)} saḥ yathā eva dīrghāt bādhate evam hrasvāt api bādheta . \\
\noindent \textbf{(P\_6,1.172)} na dīrghāt ṣaṭsvaraḥ prāpnoti . \\
\noindent \textbf{(P\_6,1.172)} kim kāraṇam . \\
\noindent \textbf{(P\_6,1.172)} āte kṛte ṣaṭsañjñābhāvāt . \\
\noindent \textbf{(P\_6,1.172)} ataḥ uttaram paṭhati aṣṭanaḥ dīrghagrahaṇam ṣaṭsañjñājñāpakam ākārāntasya nuḍartham . \\
\noindent \textbf{(P\_6,1.172)} aṣṭanaḥ dīrghagrahaṇam kriyate jñāpakārtham . \\
\noindent \textbf{(P\_6,1.172)} kim jñāpyam . \\
\noindent \textbf{(P\_6,1.172)} etat jñāpayati ācāryaḥ bhavati ātve kṛte ṣaṭsañjñā iti . \\
\noindent \textbf{(P\_6,1.172)} kim etasya jñapane prayojanam . \\
\noindent \textbf{(P\_6,1.172)} ākārāntasya nuḍartham . \\
\noindent \textbf{(P\_6,1.172)} ākārāntasya nuḍ siddhaḥ bhavati . \\
\noindent \textbf{(P\_6,1.172)} aṣṭānām iti . \\
\noindent \textbf{(P\_6,1.172)} nanu ca nityam ātvam . \\
\noindent \textbf{(P\_6,1.172)} etat eva jñāpayati vibhāṣā ātvam iti yat ayam dīrghagrahaṇam karoti . \\
\noindent \textbf{(P\_6,1.172)} itarathā hi aṣṭanaḥ iti eva brūyāt . \\
\noindent \textbf{(P\_6,1.173)} nadyajādyudāttatve bṛhanmahatoḥ upasaṅkhyānam . \\
\noindent \textbf{(P\_6,1.173)} nadyajādyudāttatve bṛhanmahatoḥ upasaṅkhyānam kartavyam . \\
\noindent \textbf{(P\_6,1.173)} bṛhatī mahatī bṛhatā mahatā . \\
\noindent \textbf{(P\_6,1.174)} halpūrvāt iti kimartham . \\
\noindent \textbf{(P\_6,1.174)} agnaye vāyave . \\
\noindent \textbf{(P\_6,1.174)} udāttayaṇi halgrahaṇam nakārāntārtham . \\
\noindent \textbf{(P\_6,1.174)} udāttayaṇi halgrahaṇam kartavyam . \\
\noindent \textbf{(P\_6,1.174)} kim prayojanam . \\
\noindent \textbf{(P\_6,1.174)} nakārāntārtham . \\
\noindent \textbf{(P\_6,1.174)} nakārāntāt api yathā syāt . \\
\noindent \textbf{(P\_6,1.174)} vākpatnī citpatnī . \\
\noindent \textbf{(P\_6,1.174)} halpūrvagrahaṇānarthakyam ca samudāyādeśatvāt . \\
\noindent \textbf{(P\_6,1.174)} halpūrvagrahaṇam ca anarthakam . \\
\noindent \textbf{(P\_6,1.174)} kim kāṛaṇam . \\
\noindent \textbf{(P\_6,1.174)} samudāyādeśatvāt . \\
\noindent \textbf{(P\_6,1.174)} samudāyaḥ atra ādeśaḥ . \\
\noindent \textbf{(P\_6,1.174)} svaritatve ca avacanāt . \\
\noindent \textbf{(P\_6,1.174)} svaritatve ca halpūrvagrahaṇasya avacanāt manyāmahe halpūrvagrahaṇam anarthakam iti . \\
\noindent \textbf{(P\_6,1.174)} yat tāvat ucyate udāttayaṇi halgrahaṇam nakārāntārtham iti kriyate nyāse eva . \\
\noindent \textbf{(P\_6,1.174)} dvinakārakaḥ nirdeśaḥ . \\
\noindent \textbf{(P\_6,1.174)} udāttayaṇaḥ halpūrvāt na ūṅdhātvoḥ iti . \\
\noindent \textbf{(P\_6,1.174)} yat api ucyate halpūrvagrahaṇānarthakyam ca samudāyādeśatvāt iti . \\
\noindent \textbf{(P\_6,1.174)} ayam asti kevalaḥ ādeśaḥ . \\
\noindent \textbf{(P\_6,1.174)} bahutitavā . \\
\noindent \textbf{(P\_6,1.176)} matubudāttatve regrahaṇam . \\
\noindent \textbf{(P\_6,1.176)} matubudāttatve regrahaṇam kartavyam . \\
\noindent \textbf{(P\_6,1.176)} ā revān etu naḥ viśaḥ . \\
\noindent \textbf{(P\_6,1.176)} tripratiṣedhaḥ ca . \\
\noindent \textbf{(P\_6,1.176)} treḥ ca pratiṣedhaḥ vaktavyaḥ . \\
\noindent \textbf{(P\_6,1.176)} trivatīḥ yājyānuvākyāḥ bhavanti . \\
\noindent \textbf{(P\_6,1.177)} iha kasmāt na bhavati . \\
\noindent \textbf{(P\_6,1.177)} kiśorīṇām , kumārīṇām . \\
\noindent \textbf{(P\_6,1.177)} hrasvāt iti vartate . \\
\noindent \textbf{(P\_6,1.177)} iha api tarhi na prāpnoti . \\
\noindent \textbf{(P\_6,1.177)} agnīnām , vāyūnām . \\
\noindent \textbf{(P\_6,1.177)} kim kāraṇam . \\
\noindent \textbf{(P\_6,1.177)} dīrghatve kṛte hrasvābhāvāt . \\
\noindent \textbf{(P\_6,1.177)} idam iha sampradhāryam . \\
\noindent \textbf{(P\_6,1.177)} dīrghatvam kriyatām svaraḥ iti kim atra kartavyam . \\
\noindent \textbf{(P\_6,1.177)} paratvāt dīrghatvam . \\
\noindent \textbf{(P\_6,1.177)} evam tarhi nāmsvare matau hrasvagrahaṇam . \\
\noindent \textbf{(P\_6,1.177)} nāmsvare matau hrasvagrahaṇam kartavyam . \\
\noindent \textbf{(P\_6,1.177)} matau hrasvāntāt iti . \\
\noindent \textbf{(P\_6,1.177)} tat tarhi vaktavyam . \\
\noindent \textbf{(P\_6,1.177)} na vaktavyam . \\
\noindent \textbf{(P\_6,1.177)} āha ayam hrasvāntāt na ca nāmi hrasvāntaḥ asti . \\
\noindent \textbf{(P\_6,1.177)} tatra bhūtapūrvagatiḥ vijñāsyate : hrasvāntam yat bhūtapūrvam iti . \\
\noindent \textbf{(P\_6,1.177)} sāmpratikābhāve bhūtapūrvagatiḥ vijñāyate ayam ca asti sāmpratikaḥ : tisṛṇām , catasṛṇām iti . \\
\noindent \textbf{(P\_6,1.177)} na etat asti . \\
\noindent \textbf{(P\_6,1.177)} ṣaṭtricaturbhyaḥ halādiḥ iti anena svareṇa bhavitavyam . \\
\noindent \textbf{(P\_6,1.177)} tasmin nitye prāpte iyam vibhāṣā ārabhyate . \\
\noindent \textbf{(P\_6,1.177)} evam tarhi yogavibhāgaḥ kariṣyate . \\
\noindent \textbf{(P\_6,1.177)} ṣaṭtricaturbhyaḥ nām udāttaḥ bhavati . \\
\noindent \textbf{(P\_6,1.177)} tataḥ halādiḥ . \\
\noindent \textbf{(P\_6,1.177)} halādiḥ ca vibhaktiḥ udāttā bhavati ṣaṭtricaturbhyaḥ iti . \\
\noindent \textbf{(P\_6,1.177)} idam tarhi tvam nṛṇam nṛpate jāyase śuciḥ . \\
\noindent \textbf{(P\_6,1.177)} nanu ca atra api nṛ ca anyatarasyām iti eṣaḥ svaraḥ bādhakaḥ bhaviṣyati . \\
\noindent \textbf{(P\_6,1.177)} na sidhyati . \\
\noindent \textbf{(P\_6,1.177)} na sidhyati . \\
\noindent \textbf{(P\_6,1.177)} kim kāraṇam . \\
\noindent \textbf{(P\_6,1.177)} jhalgrahaṇam tatra anuvartate . \\
\noindent \textbf{(P\_6,1.177)} kim punaḥ kāraṇam jhalgrahaṇam tatra anuvartate . \\
\noindent \textbf{(P\_6,1.177)} iha mā bhūt . \\
\noindent \textbf{(P\_6,1.177)} nrā nre . \\
\noindent \textbf{(P\_6,1.177)} udāttayaṇaḥ halpūrvāt iti eṣaḥ svaraḥ atra svaraḥ bādhakaḥ bhaviṣyati . \\
\noindent \textbf{(P\_6,1.177)} idam tarhi nari . \\
\noindent \textbf{(P\_6,1.177)} na ekam udāharam hrasvagrahaṇam prayojayati . \\
\noindent \textbf{(P\_6,1.177)} yadi etāvat prayojanam syāt nām iti eva brūyāt . \\
\noindent \textbf{(P\_6,1.177)} tatra vacanāt bhūtapūrvagatiḥ vijñāsyate . \\
\noindent \textbf{(P\_6,1.177)} hrasvāntam yat bhūtapūrvam iti . \\
\noindent \textbf{(P\_6,1.177)} atha vā na evam vijñāyate . \\
\noindent \textbf{(P\_6,1.177)} nām svarau matau hrasvagrahaṇam kartavyam iti . \\
\noindent \textbf{(P\_6,1.177)} katham tarhi . \\
\noindent \textbf{(P\_6,1.177)} nāmsvare matau hrasvāt iti vartate iti . \\
\noindent \textbf{(P\_6,1.182)} sau iti kim prathamaikavacanasya grahaṇam āhosvit saptamībahuvacanasya . \\
\noindent \textbf{(P\_6,1.182)} kutaḥ sandehaḥ . \\
\noindent \textbf{(P\_6,1.182)} samānaḥ nirdeśaḥ . \\
\noindent \textbf{(P\_6,1.182)} purastāt eṣaḥ nirṇayaḥ saptamībahuvacanasya grahaṇam iti . \\
\noindent \textbf{(P\_6,1.182)} iha api tat eva bhavitum arhati . \\
\noindent \textbf{(P\_6,1.182)} yadi saptamībahuvacanasya grahaṇam tābhyām brāhmaṇābhyām , yābhyām brāhmaṇābhyām atra na prāpnoti . \\
\noindent \textbf{(P\_6,1.182)} vidhiḥ api atra na sidhyati . \\
\noindent \textbf{(P\_6,1.182)} kim kāraṇam . \\
\noindent \textbf{(P\_6,1.182)} na hi etat bhavati yat sau rūpam . \\
\noindent \textbf{(P\_6,1.182)} idam tarhi tebhyaḥ brāhmaṇebhyaḥ , yebhyaḥ brāhmaṇebhyaḥ . \\
\noindent \textbf{(P\_6,1.182)} vidhiḥ ca sidddhaḥ bhavati pratiṣedhaḥ tu na prāpnoti . \\
\noindent \textbf{(P\_6,1.182)} asti punaḥ kim cit sati iṣṭam saṅgṛhītam bhavati āhosvit doṣāntam eva . \\
\noindent \textbf{(P\_6,1.182)} asti iti āha . \\
\noindent \textbf{(P\_6,1.182)} iha yābhyaḥ brāhmaṇībhyaḥ , tābhyaḥ brāhmaṇībhyaḥ iti vidhiḥ ca sidddhaḥ bhavati pratiṣedhaḥ ca . \\
\noindent \textbf{(P\_6,1.182)} asti tarhi prathamaikavacanasya grahaṇam . \\
\noindent \textbf{(P\_6,1.182)} yadi prathamaikavacanasya grahaṇam tena iti svaraḥ puṃsi na sidhyati . \\
\noindent \textbf{(P\_6,1.182)} na ca avaśyam puṃsi eva striyām puṃsi napuṃsake ca . \\
\noindent \textbf{(P\_6,1.182)} tena brāhmaṇena tayā brāhmaṇyā tena kuṇḍena iti . \\
\noindent \textbf{(P\_6,1.182)} saptamībahuvacanasya grahaṇe api eṣaḥ doṣaḥ . \\
\noindent \textbf{(P\_6,1.182)} tasmāt ubhābhyām eva pratiṣedhe yattatadoḥ ca grahaṇam kartavyam . \\
\noindent \textbf{(P\_6,1.182)} na gośvansāvavarṇarāḍaṅkruṅkṛdbhyaḥ yattadoḥ ca iti . \\
\noindent \textbf{(P\_6,1.185)} titi pratyayagrahaṇam . \\
\noindent \textbf{(P\_6,1.185)} titi pratyayagrahaṇam kartavyam . \\
\noindent \textbf{(P\_6,1.185)} iha mā bhūt . \\
\noindent \textbf{(P\_6,1.185)} ṝtaḥ it dhātoḥ . \\
\noindent \textbf{(P\_6,1.185)} kirati , girati . \\
\noindent \textbf{(P\_6,1.185)} tat tarhi vaktavyam . \\
\noindent \textbf{(P\_6,1.185)} na vaktavyam . \\
\noindent \textbf{(P\_6,1.185)} na eṣaḥ takāraḥ . \\
\noindent \textbf{(P\_6,1.185)} kaḥ tarhi . \\
\noindent \textbf{(P\_6,1.185)} dakāraḥ . \\
\noindent \textbf{(P\_6,1.185)} yadi dakāraḥ āntaryataḥ dīrghasya dīrghaḥ prāpnoti . \\
\noindent \textbf{(P\_6,1.185)} bhāvyamānena savarṇānām grahaṇam na iti evam na bhaviṣyati . \\
\noindent \textbf{(P\_6,1.185)} yadi bhāvyamānena savarṇānām grahaṇam na iti ucyate adasaḥ aseḥ dāt u daḥ maḥ , amūbhyām iti atra na prāpnoti . \\
\noindent \textbf{(P\_6,1.185)} evam tarhi ācāryapravṛttiḥ jñāpayati bhavati ukāreṇa bhāvyamānena savarṇānām grahaṇam iti yat ayam divaḥ ut iti ukāram taparam karoti . \\
\noindent \textbf{(P\_6,1.185)} evamartham eva tarhi pratyayagrahaṇam kartavyam atra mā bhūt iti . \\
\noindent \textbf{(P\_6,1.185)} na eṣaḥ takāharaḥ . \\
\noindent \textbf{(P\_6,1.185)} kaḥ tarhi . \\
\noindent \textbf{(P\_6,1.185)} dakāraḥ . \\
\noindent \textbf{(P\_6,1.185)} yadi dakāraḥ na jñāpakam bhavati . \\
\noindent \textbf{(P\_6,1.185)} evam tarhi taparaḥ tatkālasya iti dakāraḥ api cartvabhūtaḥ nirdiśyate . \\
\noindent \textbf{(P\_6,1.185)} yadi evam cartvasya asiddhatvāt haśi ca iti uttvam prāpnoti . \\
\noindent \textbf{(P\_6,1.185)} sautraḥ nirdeśaḥ . \\
\noindent \textbf{(P\_6,1.185)} atha vā asaṃhitayā nirdeśaḥ kariṣyate . \\
\noindent \textbf{(P\_6,1.185)} aṇudit savarṇasya ca apratyayaḥ , ttaparaḥ tatkālasya iti . \\
\noindent \textbf{(P\_6,1.186.1)} adupadeśāt iti kim idam vijñāyate . \\
\noindent \textbf{(P\_6,1.186.1)} akāraḥ yaḥ upadeśaḥ iti āhosvit akārāntam yat upadeśaḥ iti . \\
\noindent \textbf{(P\_6,1.186.1)} kim ca ataḥ . \\
\noindent \textbf{(P\_6,1.186.1)} yadi vijñāyate akāraḥ yaḥ upadeśaḥ iti hataḥ , hathaḥ iti atra api prāpnoti . \\
\noindent \textbf{(P\_6,1.186.1)} atha vijñāyate akārāntam yat upadeśaḥ iti na doṣaḥ bhavati . \\
\noindent \textbf{(P\_6,1.186.1)} nanu ca akārāntam yat upadeśaḥ iti vijñāyamāne api atra api prāpnoti . \\
\noindent \textbf{(P\_6,1.186.1)} etat api hi vyapadeśivadbhāvena akārāntam bhavati upadeśe . \\
\noindent \textbf{(P\_6,1.186.1)} arthavatā vyapadeśivadbhāvaḥ . \\
\noindent \textbf{(P\_6,1.186.1)} yadi tarhi akārāntam yat upadeśaḥ iti vijñāyate mā hi dhukṣātām , mā hi dhuṣāthām atra api prāpnoti . \\
\noindent \textbf{(P\_6,1.186.1)} astu . \\
\noindent \textbf{(P\_6,1.186.1)} anudāttatve kṛte lope udāttanivṛttisvareṇa siddham . \\
\noindent \textbf{(P\_6,1.186.1)} na sidhyati . \\
\noindent \textbf{(P\_6,1.186.1)} idam iha sampradhāryam . \\
\noindent \textbf{(P\_6,1.186.1)} adnudāttatvam kriyatām lopaḥ iti kim atra kartavyam . \\
\noindent \textbf{(P\_6,1.186.1)} paratvāt lopaḥ . \\
\noindent \textbf{(P\_6,1.186.1)} evam tarhi idam adya lasārvadhādukānudāttatvam pratyayasvarasya apavādaḥ . \\
\noindent \textbf{(P\_6,1.186.1)} na ca apavādaviṣaye utsargaḥ abhiniviśate . \\
\noindent \textbf{(P\_6,1.186.1)} pūrvam hi apavādāḥ abhiniviśante paścāt utsargāḥ . \\
\noindent \textbf{(P\_6,1.186.1)} prakalpya vā apavādaviṣayam tataḥ utasrgaḥ abhiniviśate . \\
\noindent \textbf{(P\_6,1.186.1)} tat na tāvat atra kadā cit pratyayasvaraḥ bhavati . \\
\noindent \textbf{(P\_6,1.186.1)} apavādaviṣayam lasārvadhātukānudāttatvam pratīkṣate . \\
\noindent \textbf{(P\_6,1.186.1)} tatra ānudāttatvam kriyatām lopaḥ iti kim atra kartavyam . \\
\noindent \textbf{(P\_6,1.186.1)} paratvāt lopaḥ . \\
\noindent \textbf{(P\_6,1.186.1)} yadi api paratvāt lopaḥ saḥ asau avidyamānodātte anudātte udāttaḥ lupyate . \\
\noindent \textbf{(P\_6,1.186.2)} tāsyādibhyaḥ anudāttatve saptamīnirdeśaḥ abhyastasijarthaḥ . \\
\noindent \textbf{(P\_6,1.186.2)} tāsyādibhyaḥ anudāttatve saptamīnirdeśaḥ kartavyaḥ . \\
\noindent \textbf{(P\_6,1.186.2)} lasārvadhātuke iti vaktavyam . \\
\noindent \textbf{(P\_6,1.186.2)} kim prayojanam . \\
\noindent \textbf{(P\_6,1.186.2)} abhyastasijarthaḥ . \\
\noindent \textbf{(P\_6,1.186.2)} abhyastānām ādiḥ udāttaḥ bhavati lasārvadhātuke . \\
\noindent \textbf{(P\_6,1.186.2)} sijantasya ādiḥ udāttaḥ bhavati lasārvadhātuke . \\
\noindent \textbf{(P\_6,1.186.2)} lasārvadhātukam iti ucyamāne tasya eva ādyudāttatvam syāt . \\
\noindent \textbf{(P\_6,1.186.2)} yadi saptamīnirdeśaḥ kriyate tāsyādīnām eva anudāttatvam prāpnoti . \\
\noindent \textbf{(P\_6,1.186.2)} na eṣaḥ doṣaḥ . \\
\noindent \textbf{(P\_6,1.186.2)} tāsiyādibhyaḥ iti eṣā pañcamī lasārvadhātuke iti saptamyāḥ ṣaṣṭhīm prakalpayiṣyati tasmāt iti uttarasya iti . \\
\noindent \textbf{(P\_6,1.186.3)} citsvarāt tāsyādibhyaḥ anudāttatvam vipratiṣedhena . \\
\noindent \textbf{(P\_6,1.186.3)} citsvarāt tāsyādibhyaḥ anudāttatvam bhavati vipratiṣedhena . \\
\noindent \textbf{(P\_6,1.186.3)} citsvarasya avakāśaḥ calanaḥ , copanaḥ . \\
\noindent \textbf{(P\_6,1.186.3)} tāsyādibhyaḥ anudāttatvasya avakāśaḥ . \\
\noindent \textbf{(P\_6,1.186.3)} āste śete . \\
\noindent \textbf{(P\_6,1.186.3)} iha ubhayam prāpnoti . \\
\noindent \textbf{(P\_6,1.186.3)} āsīnaḥ , śayānaḥ . \\
\noindent \textbf{(P\_6,1.186.3)} tāsyādibhyaḥ anudāttatvam bhavati vipratiṣedhena . \\
\noindent \textbf{(P\_6,1.186.3)} na eṣaḥ yuktaḥ vipratiṣedhaḥ . \\
\noindent \textbf{(P\_6,1.186.3)} kim kāraṇam . \\
\noindent \textbf{(P\_6,1.186.3)} dvikāryayogaḥ hi vipratiṣedhaḥ . \\
\noindent \textbf{(P\_6,1.186.3)} na ca atra ekaḥ dvikāryayuktaḥ . \\
\noindent \textbf{(P\_6,1.186.3)} ādeḥ anudāttatvam antasya udāttatvam . \\
\noindent \textbf{(P\_6,1.186.3)} na avaśyam dvikāryayogaḥ eva vipratiṣedhaḥ . \\
\noindent \textbf{(P\_6,1.186.3)} kim tarhi. asambhavaḥ api . \\
\noindent \textbf{(P\_6,1.186.3)} nanu ca atra api asti sambhavaḥ . \\
\noindent \textbf{(P\_6,1.186.3)} ādeḥ anudāttatvam antasya udāttatvam iti . \\
\noindent \textbf{(P\_6,1.186.3)} asti ca sambhavaḥ yat ubhayam syāt . \\
\noindent \textbf{(P\_6,1.186.3)} na eṣaḥ asti sambhavaḥ . \\
\noindent \textbf{(P\_6,1.186.3)} vakṣyati etat svaravidhau saṅghātaḥ kāryī bhavati iti . \\
\noindent \textbf{(P\_6,1.186.3)} mukaḥ ca upasaṅkhyānam . \\
\noindent \textbf{(P\_6,1.186.3)} mukaḥ ca upasaṅkhyānam kartavyam . \\
\noindent \textbf{(P\_6,1.186.3)} pacamānaḥ , yajamānaḥ . \\
\noindent \textbf{(P\_6,1.186.3)} mukā vyavahitatvāt adupadeśāt lasārvadhātukam anudāttam bhavati iti anudāttatvam na prāpnoti . \\
\noindent \textbf{(P\_6,1.186.3)} nanu ca ayam muk adupadeśabhaktaḥ adupadeśagrahaṇena grāhiṣyate . \\
\noindent \textbf{(P\_6,1.186.3)} na sidhyati . \\
\noindent \textbf{(P\_6,1.186.3)} aṅgasya muk ucyate vikaraṇāntam ca aṅgam . \\
\noindent \textbf{(P\_6,1.186.3)} saḥ asau saṅghātabhaktaḥ aśakyaḥ muk adupadeśagrahaṇena grahītum . \\
\noindent \textbf{(P\_6,1.186.3)} atha ayam adbhaktaḥ syāt gṛhyeta ayam adupadeśagrahaṇena . \\
\noindent \textbf{(P\_6,1.186.3)} bāḍham gṛhyeta . \\
\noindent \textbf{(P\_6,1.186.3)} adbhaktaḥ tarhi bhaviṣyati . \\
\noindent \textbf{(P\_6,1.186.3)} tat katham . \\
\noindent \textbf{(P\_6,1.186.3)} vakṣyati etasya parihāram . \\
\noindent \textbf{(P\_6,1.186.3)} itaḥ ca upasaṅkhyānam . \\
\noindent \textbf{(P\_6,1.186.3)} itaḥ ca upasaṅkhyānam kartavyam . \\
\noindent \textbf{(P\_6,1.186.3)} idbhiḥ ca vyavahitatvāt anudāttatvam na prāpnoti . \\
\noindent \textbf{(P\_6,1.186.3)} pacataḥ , paṭhataḥ . \\
\noindent \textbf{(P\_6,1.186.3)} itaḥ ca anekāntatvāt . \\
\noindent \textbf{(P\_6,1.186.3)} anekāntāḥ anubandhāḥ . \\
\noindent \textbf{(P\_6,1.186.3)} yadi anekāntāḥ anubandhāḥ adiprabhṛtijuhotyādibhyaḥ pratiṣedhaḥ vaktavyaḥ . \\
\noindent \textbf{(P\_6,1.186.3)} attaḥ , juhutaḥ iti . \\
\noindent \textbf{(P\_6,1.186.3)} adupadeśāt iti anudāttatvam prāpnoti . \\
\noindent \textbf{(P\_6,1.186.3)} tatra adiprabhṛtijuhotyādibhyaḥ apratiṣedhaḥ sthānyādeśābhāvāt . \\
\noindent \textbf{(P\_6,1.186.3)} tatra adiprabhṛtibhyaḥ juhotyādibhyaḥ apratiṣedhaḥ . \\
\noindent \textbf{(P\_6,1.186.3)} anarthakaḥ pratiṣedhaḥ apratiṣedhaḥ . \\
\noindent \textbf{(P\_6,1.186.3)} anudāttatvam kasmāt na bhavati . \\
\noindent \textbf{(P\_6,1.186.3)} sthānyādeśābhāvāt . \\
\noindent \textbf{(P\_6,1.186.3)} na eva atra sthāninam na eva ādeśam paśyāmaḥ . \\
\noindent \textbf{(P\_6,1.186.3)} anudāttaṅidgrahaṇāt vā . \\
\noindent \textbf{(P\_6,1.186.3)} atha vā yat ayam anudāttaṅidgrahaṇam karoti tat jñāpayati ācāryaḥ na luptavikaraṇebhyaḥ anudāttatvam bhavati iti . \\
\noindent \textbf{(P\_6,1.186.3)} na etat asti jñāpakam . \\
\noindent \textbf{(P\_6,1.186.3)} śnanartham etat syāt . \\
\noindent \textbf{(P\_6,1.186.3)} vindāte , khindāte . \\
\noindent \textbf{(P\_6,1.186.3)} yat tarhi ṅidgrahaṇam karoti . \\
\noindent \textbf{(P\_6,1.186.3)} na hi śnamvikaraṇaḥ ṅit bhavati . \\
\noindent \textbf{(P\_6,1.186.3)} ṅitaḥ anudāttatve vikaraṇebhyaḥ pratiṣedhaḥ vaktavyaḥ . \\
\noindent \textbf{(P\_6,1.186.3)} cinutaḥ , sunutaḥ , lunītaḥ , punītaḥ . \\
\noindent \textbf{(P\_6,1.186.3)} ṅitaḥ iti anudāttatvam prāpnoti . \\
\noindent \textbf{(P\_6,1.186.3)} ṅitaḥ anudāttatve vikaraṇebhyaḥ apratiṣedhaḥ sarvasya upadeśaviśeṣaṇatvāt . \\
\noindent \textbf{(P\_6,1.186.3)} ṅitaḥ anudāttatve vikaraṇebhyaḥ apratiṣedhaḥ . \\
\noindent \textbf{(P\_6,1.186.3)} anarthakaḥ pratiṣedhaḥ apratiṣedhaḥ . \\
\noindent \textbf{(P\_6,1.186.3)} anudāttatvam kasmāt na bhavati . \\
\noindent \textbf{(P\_6,1.186.3)} sarvasya upadeśaviśeṣaṇatvāt . \\
\noindent \textbf{(P\_6,1.186.3)} sarvam upadeśagrahaṇena viśeṣayiṣyāmaḥ . \\
\noindent \textbf{(P\_6,1.186.3)} upadeśe anudāttetaḥ , upadeśe ṅitaḥ , upadeśe akārāntāt . \\
\noindent \textbf{(P\_6,1.187)} sicaḥ ādyudāttatve aniṭaḥ pitaḥ upasaṅkhyānam . \\
\noindent \textbf{(P\_6,1.187)} sicaḥ ādyudāttatve aniṭaḥ pitaḥ upasaṅkhyānam kartavyam . \\
\noindent \textbf{(P\_6,1.187)} mā hi karṣam , mā hi kārṣam . \\
\noindent \textbf{(P\_6,1.187)} aniṭaḥ iti kimartham . \\
\noindent \textbf{(P\_6,1.187)} mā hi laviṣam . \\
\noindent \textbf{(P\_6,1.188)} svapādīnām vāvacanāt abhyastasvaraḥ vipratiṣedhena . \\
\noindent \textbf{(P\_6,1.188)} svapādīnām vāvacanāt abhyastasvaraḥ bhavati vipratiṣedhena . \\
\noindent \textbf{(P\_6,1.188)} svapādīnām vāvacanasya avakāśaḥ svapanti śvasanti . \\
\noindent \textbf{(P\_6,1.188)} abhyastasvarasya avakāśaḥ dadati , dadhati . \\
\noindent \textbf{(P\_6,1.188)} iha ubhayam prāpnoti . \\
\noindent \textbf{(P\_6,1.188)} jagrati . \\
\noindent \textbf{(P\_6,1.188)} abhyastasvaraḥ bhavati vipratiṣedhena . \\
\noindent \textbf{(P\_6,1.190)} anudātte ca iti bahuvrīhinirdeśaḥ lopayaṇādeśārtham . \\
\noindent \textbf{(P\_6,1.190)} anudātte ca iti bahuvrīhinirdeśaḥ kartavyaḥ . \\
\noindent \textbf{(P\_6,1.190)} avidyamānodātte iti vaktavyam . \\
\noindent \textbf{(P\_6,1.190)} kim prayojanam . \\
\noindent \textbf{(P\_6,1.190)} lopayaṇādeśārtham . \\
\noindent \textbf{(P\_6,1.190)} lopayaṇādeśayoḥ kṛtayoḥ ādyudāttatvam yathā syāt . \\
\noindent \textbf{(P\_6,1.190)} mā hi dadhāt . \\
\noindent \textbf{(P\_6,1.190)} dadhāti atra . \\
\noindent \textbf{(P\_6,1.191.1)} sarvasvaraḥ anackasya . \\
\noindent \textbf{(P\_6,1.191.1)} sarvasvaraḥ anackasya iti vaktavyam . \\
\noindent \textbf{(P\_6,1.191.1)} iha mā bhūt sarvake . \\
\noindent \textbf{(P\_6,1.191.2)} bhyādigrahaṇam kimartham . \\
\noindent \textbf{(P\_6,1.191.2)} iha mā bhūt . \\
\noindent \textbf{(P\_6,1.191.2)} dadāti dadhāti . \\
\noindent \textbf{(P\_6,1.191.2)} na etat asti prayojanam . \\
\noindent \textbf{(P\_6,1.191.2)} abhyastasvaraḥ atra bādhakaḥ bhaviṣyati . \\
\noindent \textbf{(P\_6,1.191.2)} antataḥ ubhayam syāt . \\
\noindent \textbf{(P\_6,1.191.2)} anavakāśāḥ khalu api vidhayaḥ bādhakāḥ bhavanti sāvakāśaḥ ca abhyastasvaraḥ . \\
\noindent \textbf{(P\_6,1.191.2)} kaḥ avakāśaḥ . \\
\noindent \textbf{(P\_6,1.191.2)} mimīte . \\
\noindent \textbf{(P\_6,1.191.2)} atha pratyayagrahaṇam kimartham . \\
\noindent \textbf{(P\_6,1.191.2)} pratyayāt pūrvasya udāttatvam yathā syāt . \\
\noindent \textbf{(P\_6,1.191.2)} āṭaḥ pūrvasya mā bhūt iti . \\
\noindent \textbf{(P\_6,1.191.2)} bibhayāni . \\
\noindent \textbf{(P\_6,1.191.2)} na ca eva asti viśeṣaḥ pratyayāt vā pūrvasya udāttatve sati āṭaḥ vā . \\
\noindent \textbf{(P\_6,1.191.2)} api ca pidbhaktaḥ pidgrahaṇena grāhiṣyate . \\
\noindent \textbf{(P\_6,1.191.2)} idam tarhi prayojanam . \\
\noindent \textbf{(P\_6,1.191.2)} pratyayāt pūrvasya udāttatvam yathā syāt . \\
\noindent \textbf{(P\_6,1.191.2)} āṭaḥ eva mā bhūt iti . \\
\noindent \textbf{(P\_6,1.191.2)} etat api na asti prayojanam . \\
\noindent \textbf{(P\_6,1.191.2)} pidbhaktaḥ pidgrahaṇena grāhiṣyate . \\
\noindent \textbf{(P\_6,1.191.2)} evam tarhi siddhe sati yat pratyayagrahaṇam karoti tat jñāpayati ācāryaḥ svaravidhau saṅghātaḥ kāryī bhavati iti . \\
\noindent \textbf{(P\_6,1.191.2)} kim etasya jñāpane prayojanam . \\
\noindent \textbf{(P\_6,1.191.2)} citsvarāt tāsyādibhyaḥ anudāttatvam vipratiṣedhena iti uktam . \\
\noindent \textbf{(P\_6,1.191.2)} tat upapannam bhavati . \\
\noindent \textbf{(P\_6,1.191.2)} atha pūrvagrahaṇam kimartham na tasmin iti nirdiṣṭe pūrvasya iti pūrvasya eva bhaviṣyati . \\
\noindent \textbf{(P\_6,1.191.2)} evam tarhi siddhe sati yat pūrvagrahaṇam karoti tat jñāpayati ācāryaḥ svaravidhau saptamyaḥ tadantasaptamyaḥ bhavanti iti . \\
\noindent \textbf{(P\_6,1.191.2)} kim etasya jñāpane prayojanam . \\
\noindent \textbf{(P\_6,1.191.2)} upottamam riti ridantasya . \\
\noindent \textbf{(P\_6,1.191.2)} caṅi anyatarasyām caṅantasya . \\
\noindent \textbf{(P\_6,1.191.2)} yadi etat jñāpyate caturaḥ śasi iti śasantasya api prāpnoti . \\
\noindent \textbf{(P\_6,1.191.2)} śasgrahaṇasāmarthyāt na bhaviṣyati . \\
\noindent \textbf{(P\_6,1.191.2)} itarathā hi tatra eva ayam brūyāt ūḍidampadādyappumraidyubhyaḥ caturbhyaḥ ca iti . \\
\noindent \textbf{(P\_6,1.191.2)} atha pidgrahaṇam kimartham . \\
\noindent \textbf{(P\_6,1.191.2)} iha mā bhūt . \\
\noindent \textbf{(P\_6,1.191.2)} jāgrati. na etat asti prayojanam . \\
\noindent \textbf{(P\_6,1.191.2)} bhavati eva atra pūrveṇa . \\
\noindent \textbf{(P\_6,1.191.2)} idam tarhi prayojanam daridrati . \\
\noindent \textbf{(P\_6,1.191.2)} ākāreṇa vyavahitatvāt na bhaviṣyati . \\
\noindent \textbf{(P\_6,1.191.2)} lope kṛte na asti vyavadhānam . \\
\noindent \textbf{(P\_6,1.191.2)} sthānivadbhāvād vyavadhānam eva . \\
\noindent \textbf{(P\_6,1.191.2)} pratiṣidhyate atra sthānivadbhāvaḥ svarasandhim prati na sthānivat iti . \\
\noindent \textbf{(P\_6,1.195)} yaki rapare upasaṅkhyānam . \\
\noindent \textbf{(P\_6,1.195)} yaki rapare upasaṅkhyānamkartavyam . \\
\noindent \textbf{(P\_6,1.195)} stīryate svayam eva . \\
\noindent \textbf{(P\_6,1.195)} upadeśavacanāt siddham . \\
\noindent \textbf{(P\_6,1.195)} upadeśe iti vaktavyam . \\
\noindent \textbf{(P\_6,1.195)} upadeśavacane janādīnām . \\
\noindent \textbf{(P\_6,1.195)} upadeśavacane janādīnām svaraḥ na sidhyati . \\
\noindent \textbf{(P\_6,1.195)} jayate svayam eva . \\
\noindent \textbf{(P\_6,1.195)} jāyate svayam eva . \\
\noindent \textbf{(P\_6,1.195)} yogavibhāgāt siddham . \\
\noindent \textbf{(P\_6,1.195)} yogavibhāgaḥ kariṣyate . \\
\noindent \textbf{(P\_6,1.195)} ajantānām kartṛyaki vā ādiḥ udāttaḥ bhavati . \\
\noindent \textbf{(P\_6,1.195)} cīyate svayam eva . \\
\noindent \textbf{(P\_6,1.195)} ciyate svayam eva . \\
\noindent \textbf{(P\_6,1.195)} jayate svayam eva . \\
\noindent \textbf{(P\_6,1.195)} jāyate svayam eva . \\
\noindent \textbf{(P\_6,1.195)} tataḥ upadeśe . \\
\noindent \textbf{(P\_6,1.195)} upadeśe ca ajantānām kartṛyaki vā ādiḥ udāttaḥ bhavati . \\
\noindent \textbf{(P\_6,1.195)} stīryate svayam eva . \\
\noindent \textbf{(P\_6,1.195)} stīryate svayam eva . \\
\noindent \textbf{(P\_6,1.195)} tat tarhi upadeśagrahaṇam kartavyam . \\
\noindent \textbf{(P\_6,1.195)} na hi antareṇa upadeśagrahaṇam yogāṅgam jāyate . \\
\noindent \textbf{(P\_6,1.195)} na kartavyam . \\
\noindent \textbf{(P\_6,1.195)} prakṛtam anuvartate . \\
\noindent \textbf{(P\_6,1.195)} kva prakṛtam . \\
\noindent \textbf{(P\_6,1.195)} tāsyanudāttenṅidadupadeśāt lasārvadhātukam anudāttam ahnviṅoḥ iti . \\
\noindent \textbf{(P\_6,1.195)} nanu ca uktam upadeśavacane janādīnām svaraḥ na sidhyati iti . \\
\noindent \textbf{(P\_6,1.195)} na eṣaḥ doṣaḥ . \\
\noindent \textbf{(P\_6,1.195)} na evam vijñāyate upadeśavacane janādīnām svaraḥ na sidhyati iti . \\
\noindent \textbf{(P\_6,1.195)} katham tarhi . \\
\noindent \textbf{(P\_6,1.195)} janādīnām api āttve upadeśavacanam kartavyam . \\
\noindent \textbf{(P\_6,1.195)} tat tarhi tatra upadeśagrahaṇam kartavyam . \\
\noindent \textbf{(P\_6,1.195)} na kartavyam . \\
\noindent \textbf{(P\_6,1.195)} prakṛtam anuvartate . \\
\noindent \textbf{(P\_6,1.195)} kva prakṛtam . \\
\noindent \textbf{(P\_6,1.195)} anudāttopadeśavanatitanotyādīnām anunāsikalopaḥ jhali kṅiti iti . \\
\noindent \textbf{(P\_6,1.196)} seḍgrahaṇam kimartham na thali iṭ antaḥ vā iti ucyeta . \\
\noindent \textbf{(P\_6,1.196)} iṭ antaḥ vā iti ucyamāne iha api prasajyeta papaktha . \\
\noindent \textbf{(P\_6,1.196)} na etat asti prayojanam . \\
\noindent \textbf{(P\_6,1.196)} acaḥ iti vartate . \\
\noindent \textbf{(P\_6,1.196)} idam tarhi prayojanam yayātha iti . \\
\noindent \textbf{(P\_6,1.204)} kimartham idam ucyate na ñniti ādiḥ nityam iti eva siddham . \\
\noindent \textbf{(P\_6,1.204)} ñniti iti ucyate na ca atra ñnitam paśyāmaḥ . \\
\noindent \textbf{(P\_6,1.204)} pratyayalakṣaṇena . \\
\noindent \textbf{(P\_6,1.204)} na lumatā tasmin iti pratyayalakṣaṇapratiṣedhaḥ . \\
\noindent \textbf{(P\_6,1.204)} aṅgādhikāroktasya saḥ pratiṣedhaḥ na lumatā aṅgasya iti . \\
\noindent \textbf{(P\_6,1.204)} ataḥ uttaram paṭhati upamānasya ādyudāttavacanam jñāpakam anubandhalakṣaṇe svare pratyayalakṣaṇapratiṣedhasya . \\
\noindent \textbf{(P\_6,1.204)} upamānasya ādyudāttavacanam jñāpakārtham kriyate . \\
\noindent \textbf{(P\_6,1.204)} kim jñāpyate . \\
\noindent \textbf{(P\_6,1.204)} etat jñāpayati ācāryaḥ anubandhalakṣaṇe svare pratyayalakṣaṇam na bhavati iti . \\
\noindent \textbf{(P\_6,1.204)} kim etasya jñāpane prayojanam . \\
\noindent \textbf{(P\_6,1.204)} gargaḥ , vatsaḥ , bidaḥ , urvāḥ , uṣṭragrīvaḥ , vāmarajjuḥ : ñniti iti ādyudāttatvam mā bhūt iti . \\
\noindent \textbf{(P\_6,1.204)} iha ca : atrayaḥ iti : taddhitasya kitaḥ iti antodāttatvam na bhavati . \\
\noindent \textbf{(P\_6,1.204)} yadi anubandhalakṣaṇe iti ucyate pathipriyaḥ , mathipriyaḥ iti : pathimathoḥ sarvanāmasthāne iti ādyudāttatvam prāpnoti . \\
\noindent \textbf{(P\_6,1.204)} evam tarhi ācāryaḥ jñāpayati svare pratyayalakṣaṇam na bhavati iti . \\
\noindent \textbf{(P\_6,1.204)} evam api sarpiḥ āgaccha , sapta āgacchata iti : āmantritasya ca iti ādyudāttatvam na prāpnoti . \\
\noindent \textbf{(P\_6,1.204)} iha ca : ma hi datām , ma hi dhatām : ādiḥ sicaḥ anyatarasyām iti eṣaḥ svaraḥ na prāpnoti . \\
\noindent \textbf{(P\_6,1.204)} evam tarhi jñāpayati ācāryaḥ saptamīnirdiṣṭe svare pratyayalakṣaṇam na bhavati iti . \\
\noindent \textbf{(P\_6,1.204)} evam api sarvastomaḥ , sarvapṛṣṭhaḥ : sarvasya supi iti ādyudāttatvam na prāpnoti . \\
\noindent \textbf{(P\_6,1.204)} astu tarhi anubandhalakṣaṇe iti eva . \\
\noindent \textbf{(P\_6,1.204)} katham pathipriyaḥ , mathipriyaḥ . \\
\noindent \textbf{(P\_6,1.204)} vaktavyam eva etat : pathimathoḥ sarvanāmasthāne luki lumatā lupte pratyayalakṣaṇam na bhavati iti . \\
\noindent \textbf{(P\_6,1.205)} niṣṭhāyām yañi dīrghatve pratiṣedhaḥ . \\
\noindent \textbf{(P\_6,1.205)} niṣṭhāyām yañi dīrghatve pratiṣedhaḥ vaktavyaḥ . \\
\noindent \textbf{(P\_6,1.205)} dattābhyām , guptābhyām . \\
\noindent \textbf{(P\_6,1.205)} na vā bahiraṅgalakṣaṇatvāt . \\
\noindent \textbf{(P\_6,1.205)} na vā vaktavyam . \\
\noindent \textbf{(P\_6,1.205)} kim kāraṇam . \\
\noindent \textbf{(P\_6,1.205)} bahiraṅgalakṣaṇatvāt . \\
\noindent \textbf{(P\_6,1.205)} bahiraṅgaḥ atra dīrghaḥ , antaraṅgaḥ svaraḥ . \\
\noindent \textbf{(P\_6,1.205)} asiddham bahiraṅgam antaraṅge . \\
\noindent \textbf{(P\_6,1.205)} antareṇa pratiṣedham antareṇa ca etām paribhāṣām siddham . \\
\noindent \textbf{(P\_6,1.205)} katham . \\
\noindent \textbf{(P\_6,1.205)} na evam vijñāyate na cet ākārāntā niṣṭhā iti . \\
\noindent \textbf{(P\_6,1.205)} katham tarhi . \\
\noindent \textbf{(P\_6,1.205)} na cet ākārāt parā niṣṭhā iti . \\
\noindent \textbf{(P\_6,1.205)} yadi evam nirdeśaḥ ca eva na upapadyate . \\
\noindent \textbf{(P\_6,1.205)} na hi eṣā ākārāt parā pañcamī yuktā . \\
\noindent \textbf{(P\_6,1.205)} iha ca prāpnoti . \\
\noindent \textbf{(P\_6,1.205)} āptaḥ , rāddhaḥ iti . \\
\noindent \textbf{(P\_6,1.205)} evam tarhi na cet avarṇāt parā niṣṭhā iti . \\
\noindent \textbf{(P\_6,1.205)} bhavet nirdeśaḥ upapannaḥ . \\
\noindent \textbf{(P\_6,1.205)} iha tu prāpnoti . \\
\noindent \textbf{(P\_6,1.205)} āptaḥ , rāddhaḥ iti . \\
\noindent \textbf{(P\_6,1.205)} iha ca na prāpnoti . \\
\noindent \textbf{(P\_6,1.205)} yataḥ , rataḥ . \\
\noindent \textbf{(P\_6,1.205)} evam tarhi vihitaviśeṣaṇam akāragrahaṇam . \\
\noindent \textbf{(P\_6,1.205)} na cet akārārāntāt vihitā niṣṭhā iti . \\
\noindent \textbf{(P\_6,1.205)} evam api dattaḥ , atra na prāpnoti iha ca prāpnoti . \\
\noindent \textbf{(P\_6,1.205)} āptaḥ , rāddhaḥ iti . \\
\noindent \textbf{(P\_6,1.205)} evam tarhi kāryiviśeṣaṇam akāragrahaṇam . \\
\noindent \textbf{(P\_6,1.205)} na cet ākārāraḥ kāryī bhavati . \\
\noindent \textbf{(P\_6,1.205)} evam api adya aṣṭaḥ , kadā aṣṭaḥ , atra na prāpnoti . \\
\noindent \textbf{(P\_6,1.205)} tasmāt suṣṭhu ucyate niṣṭhāyām yañi dīrghatve pratiṣedhaḥ , na vā bahiraṅgalakṣaṇatvāt iti . \\
\noindent \textbf{(P\_6,1.207)} kim nipātyate . \\
\noindent \textbf{(P\_6,1.207)} āśite kartari nipātanam upadhādīrhatvam ādyudāttatvam ca . \\
\noindent \textbf{(P\_6,1.207)} āśitaḥ iti ktaḥ kartari nipātyate upadhādīrhatvam . \\
\noindent \textbf{(P\_6,1.207)} āśitavān āśitaḥ . \\
\noindent \textbf{(P\_6,1.207)} ādyudāttatvam ca nipātyate . \\
\noindent \textbf{(P\_6,1.207)} ādyudāttatvam anipātyam . \\
\noindent \textbf{(P\_6,1.207)} adhikārāt siddham . \\
\noindent \textbf{(P\_6,1.207)} upadhādīrhatvam anipātyam . \\
\noindent \textbf{(P\_6,1.207)} āṅpūrvasya prayogaḥ . \\
\noindent \textbf{(P\_6,1.207)} yadi evam avagrahaḥ prāpnoti . \\
\noindent \textbf{(P\_6,1.207)} na lakṣaṇena padakārāḥ anuvartyāḥ . \\
\noindent \textbf{(P\_6,1.207)} padkāraiḥ nāma lakṣaṇam anuvartyam . \\
\noindent \textbf{(P\_6,1.207)} yathālakṣaṇam padam kartavyam . \\
\noindent \textbf{(P\_6,1.208,} 215) KA\_III,117.22-118.3 Ro\_IV,531 {1/7}     kim iyam prāpte vibhāṣā āhosvit aprāpte . \\
\noindent \textbf{(P\_6,1.208,} 215) KA\_III,117.22-118.3 Ro\_IV,531 {2/7}     katham ca prāpte katham vā aprāpte . \\
\noindent \textbf{(P\_6,1.208,} 215) KA\_III,117.22-118.3 Ro\_IV,531 {3/7}     yadi sañjñāyām upamānam , niṣṭhā ca dvyac anāt iti nitye prāpte ārambhaḥ tata prāpte anyatra vā aprāpte . \\
\noindent \textbf{(P\_6,1.208,} 215) KA\_III,117.22-118.3 Ro\_IV,531 {4/7}     veṇuriktayoḥ aprāpte . \\
\noindent \textbf{(P\_6,1.208,} 215) KA\_III,117.22-118.3 Ro\_IV,531 {5/7}     veṇuriktayoḥ aprāpte vibhāṣā prāpte nityaḥ vidhiḥ . \\
\noindent \textbf{(P\_6,1.208,} 215) KA\_III,117.22-118.3 Ro\_IV,531 {6/7}     veṇuḥ iva veṇuḥ . \\
\noindent \textbf{(P\_6,1.208,} 215) KA\_III,117.22-118.3 Ro\_IV,531 {7/7}     riktaḥ nāma kaḥ cit . \\
\noindent \textbf{(P\_6,1.217)} upottamagrahaṇam kimarthan na riti pūrvam iti eva ucyeta . \\
\noindent \textbf{(P\_6,1.217)} tatra ayam api arthaḥ . \\
\noindent \textbf{(P\_6,1.217)} matoḥ pūrvam āt sañjñāyām striyām iti atra pūrvagrahaṇam na kartavyam bhavati . \\
\noindent \textbf{(P\_6,1.217)} evam tarhi upottamagrahaṇam uttarārtham . \\
\noindent \textbf{(P\_6,1.217)} caṅi anyatarasyām upottamam iti eva . \\
\noindent \textbf{(P\_6,1.217)} iha mā bhūt . \\
\noindent \textbf{(P\_6,1.217)} mā hi sma dadhat . \\
\noindent \textbf{(P\_6,1.220-221)} KA\_III,118.11-15 Ro\_IV,532 {1/9}     kimartham idam ucyate na vatyāḥ iti eva ucyate . \\
\noindent \textbf{(P\_6,1.220-221)} KA\_III,118.11-15 Ro\_IV,532 {2/9}     vatyāḥ iti iyati ucyamāne rājavatī , atra api prasajyeta . \\
\noindent \textbf{(P\_6,1.220-221)} KA\_III,118.11-15 Ro\_IV,532 {3/9}     atha avatyāḥ iti ucyamāne kasmāt eva atra na bhavati . \\
\noindent \textbf{(P\_6,1.220-221)} KA\_III,118.11-15 Ro\_IV,532 {4/9}     asiddhaḥ nalopaḥ . \\
\noindent \textbf{(P\_6,1.220-221)} KA\_III,118.11-15 Ro\_IV,532 {5/9}     tasya asiddhatvāt na eṣaḥ avatīśabdaḥ . \\
\noindent \textbf{(P\_6,1.220-221)} KA\_III,118.11-15 Ro\_IV,532 {6/9}     kaḥ tarhi . \\
\noindent \textbf{(P\_6,1.220-221)} KA\_III,118.11-15 Ro\_IV,532 {7/9}     anvatīśabdaḥ . \\
\noindent \textbf{(P\_6,1.220-221)} KA\_III,118.11-15 Ro\_IV,532 {8/9}     yathā eva tarhi nalopasya asiddhatvāt na avatīśabdaḥ evam vatvasya api asiddhatvāt na avatīśabdaḥ . \\
\noindent \textbf{(P\_6,1.220-221)} KA\_III,118.11-15 Ro\_IV,532 {9/9}     āśrayāt siddhatvam syāt . \\
\noindent \textbf{(P\_6,1.222)} coḥ ataddhite . \\
\noindent \textbf{(P\_6,1.222)} cusvaraḥ ataddhite iti vaktavyam . \\
\noindent \textbf{(P\_6,1.222)} iha mā bhūt . \\
\noindent \textbf{(P\_6,1.222)} dādhīcaḥ , mādhūcaḥ iti . \\
\noindent \textbf{(P\_6,1.222)} tat tarhi vaktavyam . \\
\noindent \textbf{(P\_6,1.222)} na vaktavyam . \\
\noindent \textbf{(P\_6,1.222)} pratyayasvaraḥ atra bādhakaḥ bhaviṣyati . \\
\noindent \textbf{(P\_6,1.222)} sthānāntaraprāptaḥ cusvaraḥ . \\
\noindent \textbf{(P\_6,1.222)} pratyayasvarasya apavādaḥ anudāttau suppitau iti . \\
\noindent \textbf{(P\_6,1.222)} anudāttau suppitau iti asya udāttanivṛttisvaraḥ . \\
\noindent \textbf{(P\_6,1.222)} udāttanivṛttisvarasya cusvaraḥ . \\
\noindent \textbf{(P\_6,1.222)} saḥ yathā eva udāttanivṛttisvaram bādhate evam pratyayasvaram api bādheta . \\
\noindent \textbf{(P\_6,1.222)} na atra udāttanivṛttisvaraḥ prāpnoti . \\
\noindent \textbf{(P\_6,1.222)} kim kāraṇam . \\
\noindent \textbf{(P\_6,1.222)} na gośvansāvavarṇa iti pratiṣedhāt . \\
\noindent \textbf{(P\_6,1.222)} na eṣaḥ udāttanivṛttisvarasya pratiṣedhaḥ . \\
\noindent \textbf{(P\_6,1.222)} kasya tarhi . \\
\noindent \textbf{(P\_6,1.222)} tṛtīyādisvarasya . \\
\noindent \textbf{(P\_6,1.222)} yatra tarhi tṛtīyādisvaraḥ na asti dadhīcaḥ paśya iti . \\
\noindent \textbf{(P\_6,1.222)} evam tarhi na tṛtīyādilakṣaṇasya pratiṣedham ṣiṣmaḥ . \\
\noindent \textbf{(P\_6,1.222)} kim tarhi . \\
\noindent \textbf{(P\_6,1.222)} yena kena cit lakṣaṇena prāptasya vibhaktisvarasya pratiṣedham . \\
\noindent \textbf{(P\_6,1.222)} yadi vibhaktisvarasya pratiṣedhaḥ vṛkṣavān , plakṣavān atra na prāpnoti . \\
\noindent \textbf{(P\_6,1.222)} matubgrahaṇam api prakṛtam anuvartate . \\
\noindent \textbf{(P\_6,1.222)} kva prakṛtam . \\
\noindent \textbf{(P\_6,1.222)} hrasvanuḍbhyām matup iti . \\
\noindent \textbf{(P\_6,1.222)} yadi tat anuvartate vetasvān iti atra prāpnoti . \\
\noindent \textbf{(P\_6,1.222)} matubgrahaṇam anuvartate ḍmatup ca eṣaḥ . \\
\noindent \textbf{(P\_6,1.222)} yadi tari matubgrahaṇe ḍmatupaḥ grahaṇam na bhavati vetasvān iti atra vatvam na prāpnoti . \\
\noindent \textbf{(P\_6,1.222)} sāmānyagrahaṇam vatve iha punaḥ viśiṣṭasya grahaṇam . \\
\noindent \textbf{(P\_6,1.222)} yatra tarhi vibhaktiḥ na asti dadhīcī iti . \\
\noindent \textbf{(P\_6,1.222)} yadi punaḥ ayam udāttanivṛttisvarasya api pratiṣedhaḥ vijñāyeta . \\
\noindent \textbf{(P\_6,1.222)} na evam śakyam . \\
\noindent \textbf{(P\_6,1.222)} iha api prasajyeta kumārī iti . \\
\noindent \textbf{(P\_6,1.222)} satiśiṣṭaḥ khalu api cusvaraḥ . \\
\noindent \textbf{(P\_6,1.222)} katham . \\
\noindent \textbf{(P\_6,1.222)} cau iti ucyate . \\
\noindent \textbf{(P\_6,1.222)} yatra asya etat rūpam . \\
\noindent \textbf{(P\_6,1.222)} ajādau asarvanāmasthāne abhinirvṛtte akāralope nakāralope ca . \\
\noindent \textbf{(P\_6,1.222)} tasmāt suṣthu ucyate coḥ ataddhite iti . \\
\noindent \textbf{(P\_6,1.223)} samāsāntodāttatve vyañjanānteṣu upasaṅkhyānam . \\
\noindent \textbf{(P\_6,1.223)} samāsāntodāttatve vyañjanānteṣu upasaṅkhyānam kartavyam : rājadṛṣat, brāhmaṇasamit . \\
\noindent \textbf{(P\_6,1.223)} halsvaraprāptau vā vyañjanam avidyamānavat . \\
\noindent \textbf{(P\_6,1.223)} atha vā halsvaraprāptau vyañjanam avidyamānavat bhavati iti eṣā paribhāṣā kartavyā . \\
\noindent \textbf{(P\_6,1.223)} kimartham idam ubhayam ucyate na halsvaraprāptau avidyamānavat iti eva ucyate svaraprāptau vyañjanam avidyamānavat bhavati iti vā . \\
\noindent \textbf{(P\_6,1.223)} dvirbaddham subaddham bhavati iti . \\
\noindent \textbf{(P\_6,1.223)} yadi halsvaraprāptau vyañjanam avidyamānavat iti ucyate dadhi , udāttāt anudāttasya svaritaḥ iti svaritatvam na prāpnoti . \\
\noindent \textbf{(P\_6,1.223)} udāttāt ca svaravidhau vyañjanam avidyamānavat bhavati iti eṣā paribhāṣā kartavyā . \\
\noindent \textbf{(P\_6,1.223)} kāni etasyāḥ paribhāṣāyāḥ prayojanāni . \\
\noindent \textbf{(P\_6,1.223)} prayojanam lidādyudāttāntodāttvidhayaḥ . \\
\noindent \textbf{(P\_6,1.223)} liti pratyayāt pūrvam udāttam bhavati iti iha eva syāt : bhaurikividham , bhaulikividham . \\
\noindent \textbf{(P\_6,1.223)} cikīrṣakaḥ , jihīrṣakaḥ iti atra na syāt . \\
\noindent \textbf{(P\_6,1.223)} ñniti ādiḥ nityam iti iha eva syāt : ahicumbukāyaniḥ , āgniveśyaḥ . \\
\noindent \textbf{(P\_6,1.223)} gārgyaḥ , kṛtiḥ iti atra na syāt . \\
\noindent \textbf{(P\_6,1.223)} dhātoḥ antaḥ udāttaḥ bhavati iti iha eva syāt ūrṇoti . \\
\noindent \textbf{(P\_6,1.223)} pacati iti atra na syāt . \\
\noindent \textbf{(P\_6,1.223)} idam tāvat yat ucyate halsvaraprāptau vyañjanam avidyamānavat bhavati iti katham hi halaḥ nāma svaraprāptiḥ syāt . \\
\noindent \textbf{(P\_6,1.223)} tat ca api bruvatā udāttāt ca svaravidhau iti vaktavyam . \\
\noindent \textbf{(P\_6,1.223)} tathā anudāttādeḥ antodāttāt ca yat ucyate tat vyañjanādeḥ vyañjanāntāt ca na prāpnoti . \\
\noindent \textbf{(P\_6,1.223)} yadi punaḥ svaravidhau vyañjanam avidyamānavat bhavati iti ucyeta . \\
\noindent \textbf{(P\_6,1.223)} atha svaravidhau vyañjanam avidyamānavat bhavati iti ucyamāne anudāttādeḥ antodāttāt ca yat ucyate tat kim siddham bhavati vyañjanādeḥ vyañjanāntāt ca . \\
\noindent \textbf{(P\_6,1.223)} bāḍham siddham . \\
\noindent \textbf{(P\_6,1.223)} katham . \\
\noindent \textbf{(P\_6,1.223)} svaravidhiḥ iti sarvavibhaktyantaḥ samāsaḥ : svareṇa vidhiḥ svaravidhiḥ , svarasya vidhiḥ svaravidhiḥ iti . \\
\noindent \textbf{(P\_6,1.223)} na evam śakyam . \\
\noindent \textbf{(P\_6,1.223)} iha hi doṣaḥ syāt . \\
\noindent \textbf{(P\_6,1.223)} udaśvitvān ghoṣaḥ , vidyutvān balāhakaḥ iti . \\
\noindent \textbf{(P\_6,1.223)} hrasvanuḍbhyām matup iti eṣaḥ svaraḥ prasajyeta . \\
\noindent \textbf{(P\_6,1.223)} astu tarhi halsvaraprāptau vyañjanam avidyamānavat bhavati iti . \\
\noindent \textbf{(P\_6,1.223)} nanu ca uktam katham hi halaḥ nāma svaraprāptiḥ syāt . \\
\noindent \textbf{(P\_6,1.223)} uccaiḥ udāttaḥ , nīcaiḥ anudāttaḥ iti atra ṣaṣṭhīnirdiṣṭam ajgrahaṇam nivṛttam . \\
\noindent \textbf{(P\_6,1.223)} tasmin nivṛtte halaḥ api svaraprāptiḥ bhavati . \\
\noindent \textbf{(P\_6,1.223)} yat api ucyate udāttāt ca svaravidhau iti vaktavyam iti . \\
\noindent \textbf{(P\_6,1.223)} na vaktavyam . \\
\noindent \textbf{(P\_6,1.223)} na idam pāribhāṣikasya anudāttasya grahaṇam . \\
\noindent \textbf{(P\_6,1.223)} kim tarhi . \\
\noindent \textbf{(P\_6,1.223)} anvarthagrahaṇam . \\
\noindent \textbf{(P\_6,1.223)} avidyamānodāttam anudāttam . \\
\noindent \textbf{(P\_6,1.223)} tasya svaritaḥ iti . \\
\noindent \textbf{(P\_6,1.223)} yat api ucyate tat vyañjanādeḥ vyañjanāntāt ca na prāpnoti iti . \\
\noindent \textbf{(P\_6,1.223)} ācāryapravṛttiḥ jñāpayati siddham tat bhavati vyañjanādeḥ vyañjanāntāt ca iti yat ayam na uttarapade anudāttādau iti uktvā apṛthivīrudralkpūṣamanthiṣu iti pratiṣedham śāsti . \\
\noindent \textbf{(P\_6,1.223)} sā tarhi eṣā paribhāṣā kartavyā . \\
\noindent \textbf{(P\_6,1.223)} na kartavyā . \\
\noindent \textbf{(P\_6,1.223)} ācāryapravṛttiḥ jñāpayati bhavati eṣā paribhāṣā yat ayam yataḥ anāvaḥ iti nāvaḥ pratiṣedham śāsti . \\
\noindent \textbf{(P\_6,2.1)} kimartham idam ucyate . \\
\noindent \textbf{(P\_6,2.1)} bahuvrīhisvaram śāsti samāsāntavidheḥ sukṛt . \\
\noindent \textbf{(P\_6,2.1)} sukṛt ācāryaḥ samāsāntodāttatve prāpte bahuvrīhisvaram apavādam śāsti . \\
\noindent \textbf{(P\_6,2.1)} na etat asti prayojanam . \\
\noindent \textbf{(P\_6,2.1)} nañsubhyām niyamārtham tu . \\
\noindent \textbf{(P\_6,2.1)} nañsubhyām iti etat niyamārtham bhaviṣyati . \\
\noindent \textbf{(P\_6,2.1)} nañsubhyām eva bahuvrīheḥ antaḥ udāttaḥ bhavati na anyasya iti . \\
\noindent \textbf{(P\_6,2.1)} evam api kutat etat pūrvapadaprakṛtisvaratvam bhaviṣyati na punaḥ parasya iti . \\
\noindent \textbf{(P\_6,2.1)} parasya śitiśāsanāt . \\
\noindent \textbf{(P\_6,2.1)} śiteḥ nityābahvac iti etat niyamārtham bhaviṣyati . \\
\noindent \textbf{(P\_6,2.1)} śiteḥ eva na anyataḥ iti . \\
\noindent \textbf{(P\_6,2.1)} yat tāvat ucyate nañsubhyām niyamārtham iti kṣepe vidhiḥ nañaḥ asiddhaḥ . \\
\noindent \textbf{(P\_6,2.1)} udarāśveṣuṣu kṣepe iti etasmin prāpte tataḥ etat ucyate . \\
\noindent \textbf{(P\_6,2.1)} yat api ucyate parasya śitiśāsanāt iti parasya niyamaḥ bhavet . \\
\noindent \textbf{(P\_6,2.1)} parasya eṣaḥ niyamaḥ syāt . \\
\noindent \textbf{(P\_6,2.1)} śiteḥ nityābahvac iti . \\
\noindent \textbf{(P\_6,2.1)} yadi pūrvapadaprakṛtisvaram samāsāntodāttatvam bādhate capriyaḥ vāpriyaḥ , atra api prāpnoti . \\
\noindent \textbf{(P\_6,2.1)} antaḥ cavāpriye sambhavāt . antodāttatvam cavāpriye siddham . \\
\noindent \textbf{(P\_6,2.1)} kutaḥ . \\
\noindent \textbf{(P\_6,2.1)} sambhavāt . \\
\noindent \textbf{(P\_6,2.1)} asati khalu api sambhave bādhanam bhavati . \\
\noindent \textbf{(P\_6,2.1)} asti ca sambhavaḥ yat ubhayam syāt . \\
\noindent \textbf{(P\_6,2.1)} sati api sambhave bādhanam bhavati . \\
\noindent \textbf{(P\_6,2.1)} tat yathā . \\
\noindent \textbf{(P\_6,2.1)} dadhi brāhmaṇebhyaḥ dīyatām takram kauṇḍinyāya iti sati api sambhave dadhidānasya takradānam nivartakam bhavati . \\
\noindent \textbf{(P\_6,2.1)} evam iha api sati api sambhave pūrvapadaprakṛtisvaram samāsāntodāttatvam bādhiṣyate . \\
\noindent \textbf{(P\_6,2.1)} evam tarhi prakṛtāt vidheḥ . bahuvrīhau prkṛtyā pūrvapadam prakṛtisvaram bhavati iti . \\
\noindent \textbf{(P\_6,2.1)} kim ca prakṛtam . \\
\noindent \textbf{(P\_6,2.1)} udāttaḥ iti ca vartate . \\
\noindent \textbf{(P\_6,2.1)} evam api kāryapriyaḥ , hāryapriyaḥ , atra na prāpnoti . \\
\noindent \textbf{(P\_6,2.1)} svarite api udāttaḥ asti . \\
\noindent \textbf{(P\_6,2.1)} atha vā svaritagrahaṇam api prakṛtam anuvartate . \\
\noindent \textbf{(P\_6,2.1)} kva prakṛtam . \\
\noindent \textbf{(P\_6,2.1)} tit svaritam iti . \\
\noindent \textbf{(P\_6,2.1)} bahuvrīhau ṛte siddham . antareṇa api bahuvrīhigrahaṇam siddham . \\
\noindent \textbf{(P\_6,2.1)} tatpuruṣe kasmāt na bhavati . \\
\noindent \textbf{(P\_6,2.1)} tatpuruṣe tulyārthatṛtīyāsaptamyupamānāvyayadvitīyākṛtyāḥ iti etat niyamārtham bhaviṣyati . \\
\noindent \textbf{(P\_6,2.1)} dvigau tarhi kasmāt na bhavati . \\
\noindent \textbf{(P\_6,2.1)} igante dvigau iti etat niyamārtham bhaviṣyati . \\
\noindent \textbf{(P\_6,2.1)} dvandve tarhi prāpnoti . \\
\noindent \textbf{(P\_6,2.1)} rājanyabahuvacanadvandve andhakavṛṣṇiṣu iti etat niyamārtham bhaviṣyati . \\
\noindent \textbf{(P\_6,2.1)} avyayībhāve tarhi prāpnoti . \\
\noindent \textbf{(P\_6,2.1)} paripratyupāpāḥ varjyamānāhorātrāvayaveṣu iti etat niyamārtham bhaviṣyati . \\
\noindent \textbf{(P\_6,2.1)} evam api kutaḥ etat evam niyamaḥ bhaviṣyati eteṣām eva tatpuruṣādiṣu iti na punaḥ evam niyamaḥ syāt eteṣām tatpuruṣādiṣu eva iti . \\
\noindent \textbf{(P\_6,2.1)} iṣṭataḥ ca avadhāraṇam . \\
\noindent \textbf{(P\_6,2.1)} iṣṭataḥ ca avadhāraṇam bhaviṣyati . \\
\noindent \textbf{(P\_6,2.1)} eteṣām tarhi bahuvrīheḥ ca paryāyaḥ prāpnoti . \\
\noindent \textbf{(P\_6,2.1)} dvipāddiṣṭeḥ vitasteḥ ca paryāyaḥ na prakalpate . yat ayam dvitribhyām pāddanmūrdhasu bahuvrīhau diṣṭivitasyoḥ ca iti siddhe paryāye paryāyam śāsti tat jñāpayati ācāryaḥ na paryāyaḥ bhavati iti . \\
\noindent \textbf{(P\_6,2.1)} udātte jñāpakam tu etat . udātte etat jñāpakam syāt . \\
\noindent \textbf{(P\_6,2.1)} svaritena samāviśet . svaritena samāveśaḥ prāpnoti . \\
\noindent \textbf{(P\_6,2.1)} svarite api udāttaḥ asti . \\
\noindent \textbf{(P\_6,2.1)} bahuvrīhisvaram śāsti samāsāntavidheḥ sukṛt . \\
\noindent \textbf{(P\_6,2.1)} nañsubhyām niyamārtham tu . \\
\noindent \textbf{(P\_6,2.1)} parasya śitiśāsanāt . \\
\noindent \textbf{(P\_6,2.1)} kṣepe vidhiḥ nañaḥ asiddhaḥ . \\
\noindent \textbf{(P\_6,2.1)} parasya niyamaḥ bhavet . \\
\noindent \textbf{(P\_6,2.1)} antaḥ cavāpriye sambhavāt . \\
\noindent \textbf{(P\_6,2.1)} prakṛtāt vidheḥ . \\
\noindent \textbf{(P\_6,2.1)} bahuvrīhau ṛte siddham . \\
\noindent \textbf{(P\_6,2.1)} iṣṭataḥ ca avadhāraṇam . \\
\noindent \textbf{(P\_6,2.1)} dvipāddiṣṭeḥ vitasteḥ ca paryāyaḥ na prakalpate . \\
\noindent \textbf{(P\_6,2.1)} udātte jñāpakam tu etat . \\
\noindent \textbf{(P\_6,2.1)} svaritena samāviśet . \\
\noindent \textbf{(P\_6,2.2)} tatpuruṣe vibhaktiprakṛtisvaratve karmadhāraye pratiṣedhaḥ . \\
\noindent \textbf{(P\_6,2.2)} tatpuruṣe vibhaktiprakṛtisvaratve karmadhāraye pratiṣedhaḥ vaktavyaḥ . \\
\noindent \textbf{(P\_6,2.2)} paramam kārakam paramakārakam paramena kārakeṇa paramakārakeṇa , parame kārake paramakārake . \\
\noindent \textbf{(P\_6,2.2)} siddham tu lakṣaṇapratipadoktayoḥ pratipadoktasya eva grahaṇāt . \\
\noindent \textbf{(P\_6,2.2)} siddham etat . \\
\noindent \textbf{(P\_6,2.2)} katham . \\
\noindent \textbf{(P\_6,2.2)} lakṣaṇapratipadoktayoḥ pratipadoktasya eva iti pratipadam yaḥ dvitīyātṛtīyāsaptamīsamāsaḥ tasya grahaṇam lakṣaṇoktaḥ ca ayam . \\
\noindent \textbf{(P\_6,2.2)} avyaye parigaṇanam kartavyam . \\
\noindent \textbf{(P\_6,2.2)} avyaye nañkunipātānām . \\
\noindent \textbf{(P\_6,2.2)} avyaye nañkunipātānām iti vaktavyam . \\
\noindent \textbf{(P\_6,2.2)} nañ . \\
\noindent \textbf{(P\_6,2.2)} abrāhmaṇaḥ , avṛṣalaḥ . \\
\noindent \textbf{(P\_6,2.2)} nañ . \\
\noindent \textbf{(P\_6,2.2)} ku . \\
\noindent \textbf{(P\_6,2.2)} kubrāhmaṇaḥ , kuvṛṣalaḥ . \\
\noindent \textbf{(P\_6,2.2)} ku . \\
\noindent \textbf{(P\_6,2.2)} nipāta . \\
\noindent \textbf{(P\_6,2.2)} niṣkauśāmbiḥ , nirvārāṇasiḥ . \\
\noindent \textbf{(P\_6,2.2)} kva mā bhūt . \\
\noindent \textbf{(P\_6,2.2)} snātvākālakaḥ , pītvāsthirakaḥ . \\
\noindent \textbf{(P\_6,2.2)} ktvāyām vā pratiṣedhaḥ . \\
\noindent \textbf{(P\_6,2.2)} ktvāyām vā pratiṣedhaḥ vaktavyaḥ . \\
\noindent \textbf{(P\_6,2.2)} snātvākālakaḥ , pītvāsthirakaḥ . \\
\noindent \textbf{(P\_6,2.2)} ubhayam na vaktavyam . \\
\noindent \textbf{(P\_6,2.2)} nipātanāt siddham . \\
\noindent \textbf{(P\_6,2.2)} nipātanāt etat siddham . \\
\noindent \textbf{(P\_6,2.2)} kim nipātanam . \\
\noindent \textbf{(P\_6,2.2)} avaśyam atra samāsārtham lyababhāvārtham ca nipātanam kartavyam . \\
\noindent \textbf{(P\_6,2.2)} tena eva yatnena svaraḥ bhaviṣyati . \\
\noindent \textbf{(P\_6,2.11)} sadṛśagrahaṇam anarthakam tṛtīyāsamāsavacanāt . \\
\noindent \textbf{(P\_6,2.11)} sadṛśagrahaṇam anarthakam . \\
\noindent \textbf{(P\_6,2.11)} kim kāraṇam . \\
\noindent \textbf{(P\_6,2.11)} tṛtīyāsamāsavacanāt . \\
\noindent \textbf{(P\_6,2.11)} sadṛśaśabdena tṛtīyāsamāsaḥ ucyate . \\
\noindent \textbf{(P\_6,2.11)} tatra tṛtīyāpūrvapadam prakṛtisvaram bhavati iti eva siddham . \\
\noindent \textbf{(P\_6,2.11)} ṣaṣṭhyartham tarhi idam vaktavyam . \\
\noindent \textbf{(P\_6,2.11)} pituḥ sadṛśaḥ pitṛsadṛśaḥ iti . \\
\noindent \textbf{(P\_6,2.11)} ṣaṣṭhyartham iti cet tṛtīyāsamāsavacanānarthakyam . \\
\noindent \textbf{(P\_6,2.11)} ṣaṣṭhyartham iti cet tṛtīyāsamāsavacanam anarthakam syāt . \\
\noindent \textbf{(P\_6,2.11)} kim kāraṇam . \\
\noindent \textbf{(P\_6,2.11)} iha asmābhiḥ traiśabdyam sādhyam . \\
\noindent \textbf{(P\_6,2.11)} pitrā sadṛśaḥ pituḥ sadṛśaḥ pitṛsadṛśaḥ iti . \\
\noindent \textbf{(P\_6,2.11)} tatra dvayoḥ śabdayoḥ samānārthayoḥ ekena vigrahaḥ apareṇa samāsaḥ bhaviṣyati aviravikanyāyena . \\
\noindent \textbf{(P\_6,2.11)} tat yathā aveḥ māṃsam iti vigṛhya avikaśabdāt utpattiḥ bhavati , āvikam iti evam pituḥ sadṛśaḥ iti vigṛhya pitṛsadṛśaḥ iti bhaviṣyati pitrā sadṛśaḥ iti vigṛhya vākyam eva . \\
\noindent \textbf{(P\_6,2.11)} avaśyam tṛtīyāsamāsaḥ vaktavyaḥ yatra ṣaṣṭhyarthaḥ na asti tadartham . \\
\noindent \textbf{(P\_6,2.11)} bhojanasadṛśaḥ , adhayayanasadṛśaḥ iti . \\
\noindent \textbf{(P\_6,2.11)} yadi tarhi tasya nibandhanam asti tat eva vaktavyam idam na vaktavyam . \\
\noindent \textbf{(P\_6,2.11)} idam api avaśyam vaktavyam yatra ṣaṣṭhī śrūyate tadartham . \\
\noindent \textbf{(P\_6,2.11)} dāsyāḥsadṛśaḥ , vṛṣalyāḥsadṛśaḥ iti . \\
\noindent \textbf{(P\_6,2.29)} igantaprakṛtisvaratve yaṇguṇayoḥ upasaṅkhyānam . \\
\noindent \textbf{(P\_6,2.29)} igantaprakṛtisvaratve yaṇguṇayoḥ upasaṅkhyānam kartavyam . \\
\noindent \textbf{(P\_6,2.29)} pañcāratnyaḥ , daśāratanyaḥ . \\
\noindent \textbf{(P\_6,2.29)} yaṇguṇayoḥ kṛtayoḥ igante dvigau iti eṣaḥ svaraḥ na prāpnoti . \\
\noindent \textbf{(P\_6,2.29)} na vā bahiraṅgalakṣaṇatvāt . \\
\noindent \textbf{(P\_6,2.29)} na vā vaktavyam . \\
\noindent \textbf{(P\_6,2.29)} kim kāraṇam . \\
\noindent \textbf{(P\_6,2.29)} bahiraṅgalakṣaṇatvāt . \\
\noindent \textbf{(P\_6,2.29)} bahiraṅgau yaṇguṇau . \\
\noindent \textbf{(P\_6,2.29)} antaraṅgaḥ svaraḥ . \\
\noindent \textbf{(P\_6,2.29)} asiddham bahiraṅgam antaraṅge . \\
\noindent \textbf{(P\_6,2.33)} paripratyupāpebhyaḥ vanam samāse vipratiṣedhena . \\
\noindent \textbf{(P\_6,2.33)} paripratyupāpebhyaḥ vanam samāse iti etat bhavati vipratiṣedhena . \\
\noindent \textbf{(P\_6,2.33)} paripratyupāpāḥ varjyamānāhorātāvayaveṣu iti asya avakāśaḥ paritrigartam, parisauvīram . \\
\noindent \textbf{(P\_6,2.33)} vanam samāse iti asya avakāśaḥ pravaṇe yaṣṭavyam . \\
\noindent \textbf{(P\_6,2.33)} iha ubhayam prāpnoti parivanam apavanam . \\
\noindent \textbf{(P\_6,2.33)} vanam samāse iti etat bhavati vipratiṣedhena . \\
\noindent \textbf{(P\_6,2.33)} na vā vanasyāndodāttatvavacanam tadapavādanivṛttyartham . \\
\noindent \textbf{(P\_6,2.33)} na vā arthaḥ vipratiṣedhena . \\
\noindent \textbf{(P\_6,2.33)} kim kāraṇam . \\
\noindent \textbf{(P\_6,2.33)} vanasyāndodāttatvavacanam tadapavādanivṛttyartham . \\
\noindent \textbf{(P\_6,2.33)} siddham atra antodāttatvam utsargeṇa eva . \\
\noindent \textbf{(P\_6,2.33)} tasya punarvacane etat prayojanam . \\
\noindent \textbf{(P\_6,2.33)} ye anye tadapavādāḥ prāpnuvanti tadbādhanārtham . \\
\noindent \textbf{(P\_6,2.33)} saḥ yathā eva tadapavādam avyayasvaram bādhate evam idam api bādhiṣyate . \\
\noindent \textbf{(P\_6,2.36)} ācāryopasarjane anekasya api pūrvapadatvāt sandehaḥ . \\
\noindent \textbf{(P\_6,2.36)} ācāryopasarjane anekasya api pūrvapadatvāt sandehaḥ bhavati . \\
\noindent \textbf{(P\_6,2.36)} āpiśalapāṇinīyavyāḍīyagautamīyāḥ . \\
\noindent \textbf{(P\_6,2.36)} ekam padam varjayitvā sarvāṇi pūrvapadāni . \\
\noindent \textbf{(P\_6,2.36)} tatra na jñāyate kasya pūrvapadasya prakṛtisvareṇa bhavitavyam iti . \\
\noindent \textbf{(P\_6,2.36)} lokavijñānāt siddham . \\
\noindent \textbf{(P\_6,2.36)} tat yathā loke , amīṣām brāhmaṇānām pūrvam ānaya iti yaḥ sarvapūrvaḥ saḥ ānīyate evam iha api yat sarvapūrvapadam tasya prakṛtisvaratvam bhaviṣyati . \\
\noindent \textbf{(P\_6,2.38)} kimartham mahataḥ pravṛddhaśabde uttarapade pūrvapadaprakṛtisvaratvam ucyate na karmadhāraye aniṣṭhā iti eva siddham . \\
\noindent \textbf{(P\_6,2.38)} na sidhyati . \\
\noindent \textbf{(P\_6,2.38)} kim kāraṇam . \\
\noindent \textbf{(P\_6,2.38)} śreṇyādisamāse evat tat iha mā bhūt , mahāniraṣṭaḥ dakṣiṇā dīyate . \\
\noindent \textbf{(P\_6,2.42)} kuruvṛjyoḥ gārhapate . \\
\noindent \textbf{(P\_6,2.42)} kuruvṛjyoḥ gārhapate iti vaktavyam . \\
\noindent \textbf{(P\_6,2.42)} kurugārhapatam , vṛjigāṛhapatam . \\
\noindent \textbf{(P\_6,2.42)} kurugārhapatariktarurvasūtajaratyaślīladṛḍharūpāpārevaḍavātailikadrūḥpaṇyakamabalaḥ dāsībhārādīnām iti vaktavyam . \\
\noindent \textbf{(P\_6,2.42)} iha api yathā syāt . \\
\noindent \textbf{(P\_6,2.42)} devahūtiḥ , devanītiḥ , vasunītiḥ , oṣadhiḥ , candramāḥ . \\
\noindent \textbf{(P\_6,2.42)} tat tarhi vaktavyam . \\
\noindent \textbf{(P\_6,2.42)} na vaktavyam . \\
\noindent \textbf{(P\_6,2.42)} yogavibhāgaḥ kariṣyate . \\
\noindent \textbf{(P\_6,2.42)} kurugārhapatariktarurvasūtajaratyaślīladṛḍharūpāpārevaḍavātailikadrūḥpaṇyakamabalaḥ iti . \\
\noindent \textbf{(P\_6,2.42)} tataḥ dāsībhārāṇām ca iti . \\
\noindent \textbf{(P\_6,2.42)} tatra bahuvacananirdeśāt dāsībhārādīnām iti vijñāsyate . \\
\noindent \textbf{(P\_6,2.42)} paṇyakambalaḥ sañjñāyām . \\
\noindent \textbf{(P\_6,2.42)} paṇyakambalaḥ sañjñāyām iti vaktavyam . \\
\noindent \textbf{(P\_6,2.42)} yaḥ paṇitavyaḥ kambalaḥ paṇyakambalaḥ eva asau bhavati . \\
\noindent \textbf{(P\_6,2.42)} aparaḥ āha : paṇyakambalaḥ eva yathā syāt . \\
\noindent \textbf{(P\_6,2.42)} kva mā bhūt . \\
\noindent \textbf{(P\_6,2.42)} paṇyagavaḥ , paṇyahastī . \\
\noindent \textbf{(P\_6,2.47)} ahīne iti kimartham . \\
\noindent \textbf{(P\_6,2.47)} kāntārātītaḥ , yojanātītaḥ . \\
\noindent \textbf{(P\_6,2.47)} ahīne dvitīyā anupasarge . \\
\noindent \textbf{(P\_6,2.47)} ahīne dvitīyā anupasarge iti vaktavyam . \\
\noindent \textbf{(P\_6,2.47)} iha mā bhūt . \\
\noindent \textbf{(P\_6,2.47)} sukhaprāptaḥ , duḥkhaprāptaḥ . \\
\noindent \textbf{(P\_6,2.47)} tat tarhi vaktavyam . \\
\noindent \textbf{(P\_6,2.47)} yadi api etat ucyate atha vā etarhi ahīnagrahaṇam na kariṣyate . \\
\noindent \textbf{(P\_6,2.47)} iha api kāntārātītaḥ , yojanātītaḥ iti anupasarge iti eva siddham . \\
\noindent \textbf{(P\_6,2.49)} anantaraḥ iti kimartham . \\
\noindent \textbf{(P\_6,2.49)} iha mā bhūt . \\
\noindent \textbf{(P\_6,2.49)} abhyuddhṛtam , upasamāhṛtam . \\
\noindent \textbf{(P\_6,2.49)} gateḥ anantaragrahaṇam anarthakam gatiḥ gatau anudāttavacanāt . \\
\noindent \textbf{(P\_6,2.49)} gateḥ anantaragrahaṇam anarthakam . \\
\noindent \textbf{(P\_6,2.49)} kim kāraṇam . \\
\noindent \textbf{(P\_6,2.49)} gatiḥ gatau anudāttavacanāt . \\
\noindent \textbf{(P\_6,2.49)} gatau parataḥ gateḥ anudāttatvam ucyate . \\
\noindent \textbf{(P\_6,2.49)} tat bādhakam bhaviṣyati . \\
\noindent \textbf{(P\_6,2.49)} tatra yasya aprakṛtisvaratvam tasmāt antodāttaprasaṅgaḥ. tatra yasya gateḥ aprakṛtisvaratvam tasmāt antodāttatvam prāpnoti antaḥ thāthaghañktājabitrakāṇām iti . \\
\noindent \textbf{(P\_6,2.49)} prakṛtisvaravacanāt hi ananodāttatvam . \\
\noindent \textbf{(P\_6,2.49)} prakṛtisvaravacanasāmarthyāt hi antodāttatvam na bhaviṣyati . \\
\noindent \textbf{(P\_6,2.49)} yadi hi syāt prakṛtisvaravacanam idānīm kimartham syāt . \\
\noindent \textbf{(P\_6,2.49)} prakṛtisvaravacanam kimartham iti cet ekagatyartham . \\
\noindent \textbf{(P\_6,2.49)} prakṛtisvaravacanam kimartham iti cet ekagatyartham . \\
\noindent \textbf{(P\_6,2.49)} yatra ekaḥ gatiḥ tadartham etat syāt . \\
\noindent \textbf{(P\_6,2.49)} prakṛtam , prahṛtam . \\
\noindent \textbf{(P\_6,2.49)} evamartham eva tarhi anantagrahaṇam kartavyam atra yathā syāt . \\
\noindent \textbf{(P\_6,2.49)} kriyamāṇe api vai anantagrahaṇe atra na sidhyati . \\
\noindent \textbf{(P\_6,2.49)} kim kāraṇam . \\
\noindent \textbf{(P\_6,2.49)} gatiḥ anantaraḥ pūrvapadam prakṛtisvaram bhavati iti ucyate . \\
\noindent \textbf{(P\_6,2.49)} yaḥ ca atra gatiḥ anantaraḥ na asau pūrvapadam yaḥ ca pūrvapadam na asau anantaraḥ . \\
\noindent \textbf{(P\_6,2.49)} apūrvapadārtham tarhi idam vaktavyam . \\
\noindent \textbf{(P\_6,2.49)} apūrvapadasya api gateḥ prakṛtisvaratvam yathā syāt . \\
\noindent \textbf{(P\_6,2.49)} apūrvapadārtham iti cet kārake atiprasaṅgaḥ .apūrvapadārtham iti cet kārake atiprasaṅgaḥ bhavati . \\
\noindent \textbf{(P\_6,2.49)} āgataḥ , dūrādāgataḥ . \\
\noindent \textbf{(P\_6,2.49)} saḥ yathā eva gatipūrvapadasya bhavati evam kārakapūrvapadasya api prāpnoti . \\
\noindent \textbf{(P\_6,2.49)} siddham tu gateḥ antodāttāprasaṅgāt . \\
\noindent \textbf{(P\_6,2.49)} siddham etat . \\
\noindent \textbf{(P\_6,2.49)} katham .yat tat gateḥ antodāttāprasaṅgāt antaḥ thāthaghañktājabitrakāṇām iti etat gateḥ na prasaṅktavyam . \\
\noindent \textbf{(P\_6,2.49)} kim kṛtam bhavati . \\
\noindent \textbf{(P\_6,2.49)} kṛtsvarāpavādaḥ ayam bhavati . \\
\noindent \textbf{(P\_6,2.49)} tatra gatiḥ anantaraḥ iti asya avakāśaḥ prakṛtam , prahṛtam . \\
\noindent \textbf{(P\_6,2.49)} antaḥ thāthaghañktājabitrakāṇām iti asya avakāśaḥ , dūrādgataḥ , dūrādyātaḥ . \\
\noindent \textbf{(P\_6,2.49)} iha ubhayam prāpnoti . \\
\noindent \textbf{(P\_6,2.49)} āgataḥ , dūrādāgataḥ . \\
\noindent \textbf{(P\_6,2.49)} antaḥ thāthaghañktājabitrakāṇām iti etat bhavati vipratiṣedhena . \\
\noindent \textbf{(P\_6,2.49)} avaśyam gateḥ tat prasaṅktavyam bhedaḥ prabhedaḥ iti evamartham . \\
\noindent \textbf{(P\_6,2.49)} evam tarhi yogavibhāgaḥ kariṣyate . \\
\noindent \textbf{(P\_6,2.49)} antaḥ thāthaghañktājabitrakāṇām iti . \\
\noindent \textbf{(P\_6,2.49)} tataḥ ktaḥ . \\
\noindent \textbf{(P\_6,2.49)} ktāntam uttarapadam antodāttam bhavati . \\
\noindent \textbf{(P\_6,2.49)} atra kārakopapadagrahaṇam . \\
\noindent \textbf{(P\_6,2.49)} anuvartate gatigrahaṇam nivṛttam . \\
\noindent \textbf{(P\_6,2.49)} atha vā upariṣṭād yogavibhāgaḥ kariṣyate . \\
\noindent \textbf{(P\_6,2.49)} idam asti sūpamānāt ktaḥ , sañjñāyām anācitādīnām , pravṛddhādīnām ca iti . \\
\noindent \textbf{(P\_6,2.49)} tataḥ vakṣyāmi kārakāt . \\
\noindent \textbf{(P\_6,2.49)} kārakāt ca ktāntam uttarapadam antodāttam bhavati . \\
\noindent \textbf{(P\_6,2.49)} tataḥ dattaśrutayoḥ eva āśiṣi kārakāt iti . \\
\noindent \textbf{(P\_6,2.49)} evam ca kṛtvā na arthaḥ anantagrahaṇena . \\
\noindent \textbf{(P\_6,2.49)} katham abhyuddhṛtam . \\
\noindent \textbf{(P\_6,2.49)} ut haratikriyam viśinaṣṭi . \\
\noindent \textbf{(P\_6,2.49)} udā viśiṣṭam abhiḥ viśinaṣṭi . \\
\noindent \textbf{(P\_6,2.49)} . tatra gatiḥ anantaraḥ iti ca prāpnoti gatiḥ gatau iti ca . \\
\noindent \textbf{(P\_6,2.49)} gatiḥ anantaraḥ iti asya avakāśaḥ prakṛtam prahṛtam . \\
\noindent \textbf{(P\_6,2.49)} gatiḥ gatau iti asya avakāśaḥ abhi ut harati , upa sam ā dadhāti . \\
\noindent \textbf{(P\_6,2.49)} iha ubhayam prāpnoti , abhyuddhṛtam , upasamāhṛtam . \\
\noindent \textbf{(P\_6,2.49)} gatiḥ gatau iti etat bhavati vipratiṣedhena .evam tarhi siddhe sati yat anantaragrahaṇam karoti tat jñāpayati ācāryaḥ bhavati eṣā paribhāṣā kṛdgrahaṇe gatikārakapūrvasya api iti . \\
\noindent \textbf{(P\_6,2.49)} kim etasya jñāpane prayojanam . \\
\noindent \textbf{(P\_6,2.49)} avataptenakulasthitam te etat , udakeviśīrṇam te etat . \\
\noindent \textbf{(P\_6,2.49)} sagatikena sanakulena samāsaḥ siddhaḥ bhavati . \\
\noindent \textbf{(P\_6,2.50)} kṛdgrahaṇam kimartham . \\
\noindent \textbf{(P\_6,2.50)} yathā takārādigrahṇam kṛdviśeṣaṇam vijñāyeta . \\
\noindent \textbf{(P\_6,2.50)} takārādau niti kṛti iti . \\
\noindent \textbf{(P\_6,2.50)} atha akriyamāṇe kṛdgrahaṇe kasya takārādigrahṇam viśeṣaṇam syāt . \\
\noindent \textbf{(P\_6,2.50)} uttarapadaviśeṣaṇam . \\
\noindent \textbf{(P\_6,2.50)} tatra kaḥ doṣaḥ . \\
\noindent \textbf{(P\_6,2.50)} iha eva syāt prataritā prataritum . \\
\noindent \textbf{(P\_6,2.50)} iha na syāt prakartā prakartum . \\
\noindent \textbf{(P\_6,2.50)} tādau niti kṛdgrahaṇānarthakyam . \\
\noindent \textbf{(P\_6,2.50)} tādau niti kṛdgrahaṇam anarthakam . \\
\noindent \textbf{(P\_6,2.50)} kriyamāṇe api kṛdgrahaṇe aniṣṭam śakyam vijñātum . \\
\noindent \textbf{(P\_6,2.50)} takārādau uttarapade niti kṛti iti . \\
\noindent \textbf{(P\_6,2.50)} akriyamāṇe ca iṣṭam . \\
\noindent \textbf{(P\_6,2.50)} nit yaḥ takārādiḥ tadante uttarapade iti . \\
\noindent \textbf{(P\_6,2.50)} yāvatā kriyamāṇe api aniṣṭam vijñāyate akriyamāṇe ca iṣṭam akriyamāṇe eva iṣṭam vijñāsyāmaḥ . \\
\noindent \textbf{(P\_6,2.50)} kṛdupadeśe vā tādyartham iḍartham .kṛdupadeśe tarhi tādyartham iḍartham kṛdgrahaṇam kartavyam . \\
\noindent \textbf{(P\_6,2.50)} kṛdupadeśe yaḥ takārādiḥ iti evam yathā vijñāyeta . \\
\noindent \textbf{(P\_6,2.50)} kim prayojanam . \\
\noindent \textbf{(P\_6,2.50)} iḍartham . \\
\noindent \textbf{(P\_6,2.50)} iḍādau api siddham bhavati . \\
\noindent \textbf{(P\_6,2.50)} pralavitā pralavitum . \\
\noindent \textbf{(P\_6,2.52.1)} anigantaprakṛtisvaratve yaṇādeśe prakṛtisvarabhāvaprasaṅgaḥ . \\
\noindent \textbf{(P\_6,2.52.1)} anigantaprakṛtisvaratve yaṇādeśe prakṛtisvarabhāvaḥ prāpnoti . \\
\noindent \textbf{(P\_6,2.52.1)} pratyaṅ pratyañcau pratyañcaḥ . \\
\noindent \textbf{(P\_6,2.52.1)} anigantavacanam idānīm kimartham syāt . \\
\noindent \textbf{(P\_6,2.52.1)} anigantavacanam kimartham iti cet ayaṇādiṣṭārtham . \\
\noindent \textbf{(P\_6,2.52.1)} ayaṇādiṣṭārtham etat syāt . \\
\noindent \textbf{(P\_6,2.52.1)} yadā yaṇādeśaḥ na . \\
\noindent \textbf{(P\_6,2.52.1)} kadā ca yaṇādeśaḥ na . \\
\noindent \textbf{(P\_6,2.52.1)} yādā śākalam . \\
\noindent \textbf{(P\_6,2.52.1)} uktam vā . \\
\noindent \textbf{(P\_6,2.52.1)} kim uktam . \\
\noindent \textbf{(P\_6,2.52.1)} samāse śākalam na bhavati iti . \\
\noindent \textbf{(P\_6,2.52.1)} yatra tarhi añcateḥ akāraḥ lupyate : pratīcaḥ pratīicā . \\
\noindent \textbf{(P\_6,2.52.1)} cusvaraḥ tatra bādhakaḥ bhaviṣyati . \\
\noindent \textbf{(P\_6,2.52.1)} ayam eva iṣyate . \\
\noindent \textbf{(P\_6,2.52.1)} vakṣyati hi etat : coḥ anigantaḥ añcatau vapratyaye iti . \\
\noindent \textbf{(P\_6,2.52.1)} yat tarhi nyadhyoḥ prakṛtisvaram śāsti . \\
\noindent \textbf{(P\_6,2.52.1)} eṣaḥ hi yaṇādiṣṭārthaḥ ārambhaḥ . \\
\noindent \textbf{(P\_6,2.52.1)} etat api ayaṇādiṣṭārtham eva syāt . \\
\noindent \textbf{(P\_6,2.52.1)} yadā yaṇādeśaḥ na . \\
\noindent \textbf{(P\_6,2.52.1)} kadā ca yaṇādeśaḥ na . \\
\noindent \textbf{(P\_6,2.52.1)} yādā śākalam . \\
\noindent \textbf{(P\_6,2.52.1)} uktam vā . \\
\noindent \textbf{(P\_6,2.52.1)} kim uktam . \\
\noindent \textbf{(P\_6,2.52.1)} samāse śākalam na bhavati iti . \\
\noindent \textbf{(P\_6,2.52.1)} yatra tarhi añcateḥ akāraḥ lupyate . \\
\noindent \textbf{(P\_6,2.52.1)} adhīcaḥ adhīcā . \\
\noindent \textbf{(P\_6,2.52.1)} cusvaraḥ tatra bādhakaḥ bhaviṣyati . \\
\noindent \textbf{(P\_6,2.52.1)} ayam eva iṣyate . \\
\noindent \textbf{(P\_6,2.52.1)} vakṣyati he etat coḥ anigantaḥ añcatau vapratyaye iti . \\
\noindent \textbf{(P\_6,2.52.1)} yat tarhi neḥ eva prakṛtisvaram śāsti . \\
\noindent \textbf{(P\_6,2.52.1)} eṣaḥ hi yaṇādiṣṭārthaḥ ārambhaḥ . \\
\noindent \textbf{(P\_6,2.52.1)} etat api ayaṇādiṣṭārtham eva syāt . \\
\noindent \textbf{(P\_6,2.52.1)} katham . \\
\noindent \textbf{(P\_6,2.52.1)} akṛte yaṇādeśa pūrvapadaprakṛtisvaratve kṛte udāttasvaritoḥ yaṇaḥ svaritaḥ vā anudāttasya iti eṣaḥ svaraḥ siddhaḥ bhavati . \\
\noindent \textbf{(P\_6,2.52.1)} nyaṅ . \\
\noindent \textbf{(P\_6,2.52.1)} tasmāt suṣṭhu ucyate anigantaprakṛtisvaratve yaṇādeśe prakṛtisvarabhāvaprasaṅgaḥ iti . \\
\noindent \textbf{(P\_6,2.52.2)} coḥ anigantaḥ añcatau vapratyaye . \\
\noindent \textbf{(P\_6,2.52.2)} cusvarāt anigantaḥ añcatau vapratyaye iti eṣaḥ svaraḥ bhavati vipratiṣedhena . \\
\noindent \textbf{(P\_6,2.52.2)} cusvarasya avakāśaḥ dadhīcaḥ paśya . \\
\noindent \textbf{(P\_6,2.52.2)} dadhīcā dadhīce . \\
\noindent \textbf{(P\_6,2.52.2)} anigantaḥ añcatau vapratyaye iti asya avakāśaḥ parāṅ parāñcau parāñcaḥ . \\
\noindent \textbf{(P\_6,2.52.2)} iha ubhayam prāpnoti . \\
\noindent \textbf{(P\_6,2.52.2)} avācā , avāce . \\
\noindent \textbf{(P\_6,2.52.2)} avakāśaḥ iti etat bhavati vipratiṣedhena . \\
\noindent \textbf{(P\_6,2.52.2)} na vā cusvarasya pūrvapadaprakṛtisvarabhāvini pratiṣedhāt itarathā hi sarvāpavādaḥ . \\
\noindent \textbf{(P\_6,2.52.2)} na vā etat vipratiṣedhena api sidhyati . \\
\noindent \textbf{(P\_6,2.52.2)} katham tarhi sidhyati . \\
\noindent \textbf{(P\_6,2.52.2)} cusvarasya pūrvapadaprakṛtisvarabhāvini pratiṣedhāt . \\
\noindent \textbf{(P\_6,2.52.2)} cursvaraḥ pūrvapadaprakṛtisvarabhāvinaḥ pratiṣedhyaḥ . \\
\noindent \textbf{(P\_6,2.52.2)} itarathā hi sarvāpavādaḥ cusvaraḥ . \\
\noindent \textbf{(P\_6,2.52.2)} akriyamāṇe hi pratiṣedhe sarvāpavādaḥ ayam cusvaraḥ . \\
\noindent \textbf{(P\_6,2.52.2)} katham . \\
\noindent \textbf{(P\_6,2.52.2)} pratyayasvarasya apavādaḥ anudāttau suppitau iti . \\
\noindent \textbf{(P\_6,2.52.2)} anudāttau suppitau iti asya udāttanivṛttisvaraḥ . \\
\noindent \textbf{(P\_6,2.52.2)} udāttanivṛttisvarasya cusvaraḥ . \\
\noindent \textbf{(P\_6,2.52.2)} saḥ yathā eva udāttanivṛttisvaram bādhate evam anigantasvaram api bādheta . \\
\noindent \textbf{(P\_6,2.52.2)} yadi tāvat saṅkhyātaḥ sāmyam ayam api caturthaḥ . \\
\noindent \textbf{(P\_6,2.52.2)} samāsāntodāttatvasya apavādaḥ avyayasvaraḥ . \\
\noindent \textbf{(P\_6,2.52.2)} avyayasvarasya kṛtsvaraḥ . \\
\noindent \textbf{(P\_6,2.52.2)} kṛtsvarasya ayam . \\
\noindent \textbf{(P\_6,2.52.2)} ubhayoḥ caturthayoḥ yuktaḥ vipratiṣedhaḥ . \\
\noindent \textbf{(P\_6,2.52.2)} satiśiṣṭaḥ tarhi cusvaraḥ . \\
\noindent \textbf{(P\_6,2.52.2)} katham . \\
\noindent \textbf{(P\_6,2.52.2)} cau iti ucyate . \\
\noindent \textbf{(P\_6,2.52.2)} yatra asya etat rūpam . \\
\noindent \textbf{(P\_6,2.52.2)} ajādau asarvanāmasthāne abhinirvṛtte akāralope nakāralope ca . \\
\noindent \textbf{(P\_6,2.52.2)} tasmāt suṣthu ucyate na vā cusvarasya pūrvapadaprakṛtisvarabhāvini pratiṣedhāt itarathā hi sarvāpavādaḥ iti . \\
\noindent \textbf{(P\_6,2.52.2)} vibhaktīṣatsvarāt kṛtsvaraḥ . \\
\noindent \textbf{(P\_6,2.52.2)} vibhaktisvarāt īṣatsvarāt ca kṛtsvaraḥ bhavati vipratiṣedhena . \\
\noindent \textbf{(P\_6,2.52.2)} vibhaktisvarasya avakāśaḥ akṣaśauṇḍaḥ , strīśauṇḍaḥ . \\
\noindent \textbf{(P\_6,2.52.2)} kṛtsvarasya avakāśaḥ , idhmapravraścanaḥ . \\
\noindent \textbf{(P\_6,2.52.2)} iha ubhayam prāpnoti pūrvāhṇesphoṭakāḥ . \\
\noindent \textbf{(P\_6,2.52.2)} kṛtsvaraḥ bhavati vipratiṣedhena . \\
\noindent \textbf{(P\_6,2.52.2)} īṣatsvarasya avakāśaḥ , īṣatkaḍāraḥ , īṣatpiṅgalaḥ . \\
\noindent \textbf{(P\_6,2.52.2)} kṛtsvarasya saḥ eva . \\
\noindent \textbf{(P\_6,2.52.2)} iha ubhayam prāpnoti , īṣadbhedaḥ . \\
\noindent \textbf{(P\_6,2.52.2)} kṛtsvaraḥ bhavati vipratiṣedhena . \\
\noindent \textbf{(P\_6,2.52.2)} citsvarāt hārisvaraḥ . \\
\noindent \textbf{(P\_6,2.52.2)} citsvarāt hārisvaraḥ bhavati vipratiṣedhena . \\
\noindent \textbf{(P\_6,2.52.2)} citsvarasya avakāśaḥ , calanaḥ , copanaḥ . \\
\noindent \textbf{(P\_6,2.52.2)} hārisvarasya avakāśaḥ , yājñikāśvaḥ , vaiyākaraṇahasī . \\
\noindent \textbf{(P\_6,2.52.2)} iha ubhayam prāpnoti , pitṛgavaḥ, mātṛgavaḥ . \\
\noindent \textbf{(P\_6,2.52.2)} hārisvaraḥ bhavati vipratiṣedhena . \\
\noindent \textbf{(P\_6,2.52.2)} kṛtsvarāt ca . \\
\noindent \textbf{(P\_6,2.52.2)} kṛtsvarāt ca hārisvaraḥ bhavati vipratiṣedhena . \\
\noindent \textbf{(P\_6,2.52.2)} kṛtsvarasya avakāśaḥ , idhmapravraścanaḥ . \\
\noindent \textbf{(P\_6,2.52.2)} hārisvarasya saḥ eva . \\
\noindent \textbf{(P\_6,2.52.2)} iha ubhayam prāpnoti , akṣahṛtaḥ , vāḍavahṛtaḥ . \\
\noindent \textbf{(P\_6,2.52.2)} hārisvaraḥ bhavati vipratiṣedhena . \\
\noindent \textbf{(P\_6,2.52.2)} na vā haraṇapratiṣedhaḥ jñāpakaḥ kṛtsvarābhādhakatavsya . \\
\noindent \textbf{(P\_6,2.52.2)} na vā arthaḥ vipratiṣedhena . \\
\noindent \textbf{(P\_6,2.52.2)} kim kāraṇam . \\
\noindent \textbf{(P\_6,2.52.2)} haraṇapratiṣedhaḥ jñāpakaḥ kṛtsvarābhādhakatvasya . \\
\noindent \textbf{(P\_6,2.52.2)} yat ayam aharaṇe iti pratiṣedham śāsti tat jñāpayati ācāryaḥ na kṛtsvaraḥ hārisvaram bādhate iti . \\
\noindent \textbf{(P\_6,2.52.2)} na etat asti jñāpakam . \\
\noindent \textbf{(P\_6,2.52.2)} anaḥ bhāvakarmavacanaḥ iti etasmin prāpte tata etat ucyate . \\
\noindent \textbf{(P\_6,2.52.2)} yadi evam sādhīyaḥ jñāpakam . \\
\noindent \textbf{(P\_6,2.52.2)} kṛtsvarasya apavādaḥ anaḥ bhāvakarmavacanaḥ iti . \\
\noindent \textbf{(P\_6,2.52.2)} bādhakam kila bādhate kim punaḥ tam . \\
\noindent \textbf{(P\_6,2.52.2)} yuktasvaraḥ ca kṛtsvarāt bhavati vipratiṣedhena . \\
\noindent \textbf{(P\_6,2.52.2)} yuktasvarasya avakāśaḥ , govallavaḥ , aśvavallavaḥ . \\
\noindent \textbf{(P\_6,2.52.2)} kṛtsvarasya saḥ eva. iha ubhayam prāpnoti , gosaṅkhyaḥ , paśūsaṅkhyaḥ , aśvasaṅkhyaḥ . \\
\noindent \textbf{(P\_6,2.52.2)} yuktasvaraḥ bhavati vipratiṣedhena . \\
\noindent \textbf{(P\_6,2.80)} upamānam iti kimartham . \\
\noindent \textbf{(P\_6,2.80)} śabdārthaprakṛtau eva iti iyati ucyamāne pūrveṇa atiprasaktam iti kṛtvā niyamaḥ ayam vijñāyeta . \\
\noindent \textbf{(P\_6,2.80)} tatra kaḥ doṣaḥ . \\
\noindent \textbf{(P\_6,2.80)} iha na syāt . \\
\noindent \textbf{(P\_6,2.80)} puṣphārī phalahārī . \\
\noindent \textbf{(P\_6,2.80)} upamānagrahaṇe punaḥ kriyamāṇe na doṣaḥ bhavati . \\
\noindent \textbf{(P\_6,2.80)} atha śabdārthagrahaṇam kimartham . \\
\noindent \textbf{(P\_6,2.80)} upamānam prakṛtau eva iti iyati ucyamāne iha api prasajyeta . \\
\noindent \textbf{(P\_6,2.80)} vṛkavañcī vṛkaprekṣī . \\
\noindent \textbf{(P\_6,2.80)} śabdāṛthagrahaṇe punaḥ kriyamāṇe na doṣaḥ bhavati . \\
\noindent \textbf{(P\_6,2.80)} atha prakṛtigrahaṇam kimartham . \\
\noindent \textbf{(P\_6,2.80)} śabdārthaprakṛtiḥ eva yaḥ nityam tatra yathā syāt . \\
\noindent \textbf{(P\_6,2.80)} iha mā bhūt . \\
\noindent \textbf{(P\_6,2.80)} kokilabhivyāhārī . \\
\noindent \textbf{(P\_6,2.80)} atha evakāraḥ kimarthaḥ . \\
\noindent \textbf{(P\_6,2.80)} niyamārthaḥ . \\
\noindent \textbf{(P\_6,2.80)} na etat asti prayojanam . \\
\noindent \textbf{(P\_6,2.80)} siddhe vidhiḥ ārabhyamāṇaḥ antareṇa evakāram niyamāṛthaḥ bhaviṣyati . \\
\noindent \textbf{(P\_6,2.80)} iṣṭataḥ avadhāraṇārthaḥ tarhi . \\
\noindent \textbf{(P\_6,2.80)} yathā evam vijñāyete : upamānam śabdārthaprakṛtau eva iti . \\
\noindent \textbf{(P\_6,2.80)} mā evam vijñāyīta : upamānam eva śabdārthaprakṛtau iti . \\
\noindent \textbf{(P\_6,2.80)} śabdārthaprakṛtau hi upamānam ca anupamānam ca ādyudāttam iṣyate : sādhvadhyāī vilambādhyāyī . \\
\noindent \textbf{(P\_6,2.82)} je dīrghāt bahvacaḥ . \\
\noindent \textbf{(P\_6,2.82)} je dīrghāntasya ādiḥ udāttaḥ bhavati iti etasmāt anytāt pūrvam bahvacaḥ iti etat bhavati vipratiṣedhena . \\
\noindent \textbf{(P\_6,2.82)} je dīrghāntasya ādiḥ udāttaḥ bhavati iti asya avakāśaḥ kuṭījaḥ , śamījaḥ . \\
\noindent \textbf{(P\_6,2.82)} anytāt pūrvam bahvacaḥ iti asya avakāśaḥ upasarajaḥ , mandurajaḥ . \\
\noindent \textbf{(P\_6,2.82)} iha ubhayam prāpnoti . \\
\noindent \textbf{(P\_6,2.82)} āmalakījaḥ , balabhījaḥ . \\
\noindent \textbf{(P\_6,2.82)} anytāt pūrvam bahvacaḥ iti etat bhavati vipratiṣedhena . \\
\noindent \textbf{(P\_6,2.91)} ādyudāttaprakaraṇe divodāsādīnām chandasi upasaṅkhyānam . \\
\noindent \textbf{(P\_6,2.91)} ādyudāttaprakaraṇe divodāsādīnām chandasi upasaṅkhyānam kartavyam . \\
\noindent \textbf{(P\_6,2.91)} divodāsāya gāyata vadhryaśvāya dāśuṣe . \\
\noindent \textbf{(P\_6,2.92-93)} KA\_III,132.13-21 Ro\_IV,562-563 {1/14}     sarvagrahaṇam kimartham . \\
\noindent \textbf{(P\_6,2.92-93)} KA\_III,132.13-21 Ro\_IV,562-563 {2/14}     guṇāt kārtsnye iti iyati ucyamāne iha api prasajyeta paramaśuklaḥ , paramakṣṛṇa iti . \\
\noindent \textbf{(P\_6,2.92-93)} KA\_III,132.13-21 Ro\_IV,562-563 {3/14}     sarvagrahaṇe punaḥ kriyamāṇe na doṣaḥ bhavati . \\
\noindent \textbf{(P\_6,2.92-93)} KA\_III,132.13-21 Ro\_IV,562-563 {4/14}     atha guṇagrahaṇam kimartham . \\
\noindent \textbf{(P\_6,2.92-93)} KA\_III,132.13-21 Ro\_IV,562-563 {5/14}     sarvam kārtsnye iti iyati ucyamāne iha api prasajyeta sarvasauvarṇaḥ sarvarājataḥ iti . \\
\noindent \textbf{(P\_6,2.92-93)} KA\_III,132.13-21 Ro\_IV,562-563 {6/14}     guṇagrahaṇe punaḥ kriyamāṇe na doṣaḥ bhavati . \\
\noindent \textbf{(P\_6,2.92-93)} KA\_III,132.13-21 Ro\_IV,562-563 {7/14}     atha kārtsnyagrahaṇam kimartham . \\
\noindent \textbf{(P\_6,2.92-93)} KA\_III,132.13-21 Ro\_IV,562-563 {8/14}     sarvam guṇe iti iyati ucyamāne iha api prasajyeta sarveṣām śvetaḥ sarvaśvetaḥ iti . \\
\noindent \textbf{(P\_6,2.92-93)} KA\_III,132.13-21 Ro\_IV,562-563 {9/14}     katham ca atra samāsaḥ . \\
\noindent \textbf{(P\_6,2.92-93)} KA\_III,132.13-21 Ro\_IV,562-563 {10/14}     ṣaṣṭhīsubantena samasyate iti . \\
\noindent \textbf{(P\_6,2.92-93)} KA\_III,132.13-21 Ro\_IV,562-563 {11/14}     guṇena na iti pratiṣedhaḥ prāpnoti . \\
\noindent \textbf{(P\_6,2.92-93)} KA\_III,132.13-21 Ro\_IV,562-563 {12/14}     guṇāt tareṇa samāsaḥ taralopaḥ ca . \\
\noindent \textbf{(P\_6,2.92-93)} KA\_III,132.13-21 Ro\_IV,562-563 {13/14}     guṇāt tareṇa samāsaḥ taralopaḥ ca vaktavyaḥ . \\
\noindent \textbf{(P\_6,2.92-93)} KA\_III,132.13-21 Ro\_IV,562-563 {14/14}     sarveṣām śvetataraḥ sarvaśvetaḥ . \\
\noindent \textbf{(P\_6,2.105)} ayuktaḥ ayam nirdeśaḥ . \\
\noindent \textbf{(P\_6,2.105)} na hi uttarapadam nāma vṛddhiḥ asti . \\
\noindent \textbf{(P\_6,2.105)} katham tarhi nirdeśaḥ kartavyaḥ . \\
\noindent \textbf{(P\_6,2.105)} vṛddhimati uttarapade iti . \\
\noindent \textbf{(P\_6,2.105)} saḥ tarhi tathā nirdeśaḥ kartavyaḥ . \\
\noindent \textbf{(P\_6,2.105)} na kartavyaḥ . \\
\noindent \textbf{(P\_6,2.105)} na evam vijñāyate . \\
\noindent \textbf{(P\_6,2.105)} uttarapadam vṛddhiḥ uttarapadavṛddhiḥ , uttarapadavṛddhau iti . \\
\noindent \textbf{(P\_6,2.105)} katham tarhi . \\
\noindent \textbf{(P\_6,2.105)} uttarapadasya vṛddhiḥ asmin saḥ ayam uttarapadavṛddhiḥ , uttarapadavṛddhau iti . \\
\noindent \textbf{(P\_6,2.106)} bahuvrīhau viśvasya antodāttāt sañjñāyām mitrājinayoḥ antaḥ .bahuvrīhau viśvasya antodāttāt sañjñāyām mitrājinayoḥ antaḥ iti etat bhavati vipratiṣedhena . \\
\noindent \textbf{(P\_6,2.106)} bahuvrīhau viśvam sañjñāyām iti asya avakāśaḥ , viśvadevaḥ , viśvayaśāḥ . \\
\noindent \textbf{(P\_6,2.106)} sañjñāyām mitrājinayoḥ antaḥ iti asya avakāśaḥ kulamitram , kulājinam . \\
\noindent \textbf{(P\_6,2.106)} iha ubhayam prāpnoti viśvamitraḥ , viśvājinaḥ . \\
\noindent \textbf{(P\_6,2.106)} sañjñāyām mitrājinayoḥ antaḥ iti etat bhavati vipratiṣedhena . \\
\noindent \textbf{(P\_6,2.106)} antodāttaprakaraṇe marudvṛdhādīnām chandasi upasaṅkhyānam . \\
\noindent \textbf{(P\_6,2.106)} antodāttaprakaraṇe marudvṛdhādīnām chandasi upasaṅkhyānam kartavyam . \\
\noindent \textbf{(P\_6,2.106)} marudvṛdhaḥ suvayāḥ upatasthe . \\
\noindent \textbf{(P\_6,2.107-108)} KA\_III,133.16-20 Ro\_IV,564 {1/6}     udarādibhyaḥ nañsubhyām . \\
\noindent \textbf{(P\_6,2.107-108)} KA\_III,133.16-20 Ro\_IV,564 {2/6}     udarāśveṣuṣu kṣepe iti etasmāt nañsubhyām iti etat bhavati vipratiṣedhena . \\
\noindent \textbf{(P\_6,2.107-108)} KA\_III,133.16-20 Ro\_IV,564 {3/6}     udarāśveṣuṣu kṣepe iti asya avakāśaḥ kuṇḍodaraḥ , ghaṭodaraḥ . \\
\noindent \textbf{(P\_6,2.107-108)} KA\_III,133.16-20 Ro\_IV,564 {4/6}     nañsubhyām iti asya avakāśaḥ ayavaḥ , atilaḥ , amāṣaḥ , suyavaḥ , sutilaḥ , sumāṣaḥ . \\
\noindent \textbf{(P\_6,2.107-108)} KA\_III,133.16-20 Ro\_IV,564 {5/6}     iha ubhayam prāpnoti , anudaraḥ , sūdaraḥ . \\
\noindent \textbf{(P\_6,2.107-108)} KA\_III,133.16-20 Ro\_IV,564 {6/6}     nañsubhyām iti etat bhavati vipratiṣedhena . \\
\noindent \textbf{(P\_6,2.117)} soḥ manasoḥ kapi . \\
\noindent \textbf{(P\_6,2.117)} soḥ manasī alomoṣasī iti etasmāt kapi pūrvam iti etat bhavati vipratiṣedhena . \\
\noindent \textbf{(P\_6,2.117)} soḥ manasī alomoṣasī iti etasya avakāśaḥ suśarmāṇam adhi nāvam ruheyam . \\
\noindent \textbf{(P\_6,2.117)} suśarmā asi supratiṣṭhānaḥ . \\
\noindent \textbf{(P\_6,2.117)} susrotāḥ , supayāḥ , suvarcāḥ . \\
\noindent \textbf{(P\_6,2.117)} kapi pūrvam iti asya avakāśaḥ ayavakaḥ . \\
\noindent \textbf{(P\_6,2.117)} iha ubhayam prāpnoti suśarmakaḥ , susrotakaḥ . \\
\noindent \textbf{(P\_6,2.117)} kapi pūrvam iti etat bhavati vipratiṣedhena . \\
\noindent \textbf{(P\_6,2.121)} pūrvādibhyaḥ kūlādīnām ādyudāttatvam . \\
\noindent \textbf{(P\_6,2.121)} pūrvādibhyaḥ kūlādīnām ādyudāttatvam bhavati vipratiṣedhena . \\
\noindent \textbf{(P\_6,2.121)} paripratiupāpāḥ varyajānāhorātrāvayaveṣu iti asya avakāśaḥ paritrigatam , parisauvīram . \\
\noindent \textbf{(P\_6,2.121)} kūlādīnām ādyudāttatvasya avakāśaḥ , atikūlam , anukūlam . \\
\noindent \textbf{(P\_6,2.121)} iha ubhayam prāpnoti parikūlam , kūlādīnām ādyudāttatvam bhavati vipratiṣedhena . \\
\noindent \textbf{(P\_6,2.126,} 130) KA\_III,134.12-17 Ro\_IV,565 {1/9}     celarājyādibhyaḥ avyayam . \\
\noindent \textbf{(P\_6,2.126,} 130) KA\_III,134.12-17 Ro\_IV,565 {2/9}     celarājyādisvarāt avyayayasvaraḥ bhavati vipratiṣedhena . \\
\noindent \textbf{(P\_6,2.126,} 130) KA\_III,134.12-17 Ro\_IV,565 {3/9}     celarājyādisvarasya avakāśaḥ , bhāryācelam , putracelam , brāhmaṇarājyam . \\
\noindent \textbf{(P\_6,2.126,} 130) KA\_III,134.12-17 Ro\_IV,565 {4/9}     avyayayasvarāvakāśaḥ , niṣkauśāmbiḥ , nirvārāṇasiḥ . \\
\noindent \textbf{(P\_6,2.126,} 130) KA\_III,134.12-17 Ro\_IV,565 {5/9}     iha ubhayam prāpnoti kucelam , kurājyam . \\
\noindent \textbf{(P\_6,2.126,} 130) KA\_III,134.12-17 Ro\_IV,565 {6/9}     saḥ tarhi pūrvavipratiṣedhaḥ vaktavyaḥ . \\
\noindent \textbf{(P\_6,2.126,} 130) KA\_III,134.12-17 Ro\_IV,565 {7/9}     na vaktavyaḥ . \\
\noindent \textbf{(P\_6,2.126,} 130) KA\_III,134.12-17 Ro\_IV,565 {8/9}     iṣṭavācī paraśabdaḥ . \\
\noindent \textbf{(P\_6,2.126,} 130) KA\_III,134.12-17 Ro\_IV,565 {9/9}     vipratiṣedhe param yat iṣṭam tat bhavati . \\
\noindent \textbf{(P\_6,2.136)} kuṇḍādyudāttatve tatsamudāyagrahaṇam . \\
\noindent \textbf{(P\_6,2.136)} kuṇḍādyudāttatve tatsamudāyagrahaṇam kartavyam . \\
\noindent \textbf{(P\_6,2.136)} vanasamudāyavācīcet kuṇḍaśabdaḥ bhavati iti vaktavyam . \\
\noindent \textbf{(P\_6,2.136)} iha mā bhūt . \\
\noindent \textbf{(P\_6,2.136)} mṛtkuṇḍam . \\
\noindent \textbf{(P\_6,2.139)} gatikārakopapadāt iti kimartham . \\
\noindent \textbf{(P\_6,2.139)} iha mā bhūt . \\
\noindent \textbf{(P\_6,2.139)} paramam kārakam , paramakārakam . \\
\noindent \textbf{(P\_6,2.139)} gatikārakopapadāt iti ucyamāne api tatra prāpnoti . \\
\noindent \textbf{(P\_6,2.139)} etat hi kārakam . \\
\noindent \textbf{(P\_6,2.139)} idam tarhi . \\
\noindent \textbf{(P\_6,2.139)} devadattasya kārakam , devadattakārakam . \\
\noindent \textbf{(P\_6,2.139)} idam ca api udāharaṇam paramam kārakam , paramakārakam iti . \\
\noindent \textbf{(P\_6,2.139)} na etat kārakam . \\
\noindent \textbf{(P\_6,2.139)} kārakaviśeṣaṇam etat . \\
\noindent \textbf{(P\_6,2.139)} yāvat brūyāt prakṛṣṭam kārakam śobhanam kārakam iti tāvat etat paramakārakam iti . \\
\noindent \textbf{(P\_6,2.139)} atha kṛdgrahaṇam kimartham . \\
\noindent \textbf{(P\_6,2.139)} iha mā bhūt . \\
\noindent \textbf{(P\_6,2.139)} niṣkauśāmbiḥ , nivārāṇasiḥ iti . \\
\noindent \textbf{(P\_6,2.139)} ataḥ uttaram paṭhati gatyādibhyaḥ prakṛtisvaratve kṛdgrahaṇānarthakyam anyasya uttarapadasya abhāvāt . \\
\noindent \textbf{(P\_6,2.139)} gatyādibhyaḥ prakṛtisvaratve kṛdgrahaṇam anarthakam . \\
\noindent \textbf{(P\_6,2.139)} kim kāraṇam . \\
\noindent \textbf{(P\_6,2.139)} anyasya uttarapadasya abhāvāt . \\
\noindent \textbf{(P\_6,2.139)} na hi anyat gatiyādibhyaḥ uttarapadam asti anyat ataḥ kṛtaḥ . \\
\noindent \textbf{(P\_6,2.139)} kim kāraṇam . \\
\noindent \textbf{(P\_6,2.139)} dhātoḥ hi dvaye pratyayāḥ vidhīyante tiṅaḥ kṛtaḥ ca . \\
\noindent \textbf{(P\_6,2.139)} tatra kṛtā saha samāsaḥ bhavati tiṅā ca na bhavati . \\
\noindent \textbf{(P\_6,2.139)} tatra antareṇa kṛdgrahaṇam kṛtaḥ eva bhaviṣyati . \\
\noindent \textbf{(P\_6,2.139)} nanu ca idānīm eva udāhṛtam niṣkauśāmbiḥ , nirvārāṇasiḥ iti . \\
\noindent \textbf{(P\_6,2.139)} yatkriyāyuktāḥ tam prati gatyupasargsañjñe bhavataḥ na ca nisaḥ kauśāmbīśabdam prati kriyāyogaḥ . \\
\noindent \textbf{(P\_6,2.139)} kṛtprakṛtau vā gatitvāt adhikāṛtham kṛdgrahaṇam . \\
\noindent \textbf{(P\_6,2.139)} kṛtprakṛtau tarhi gatitvāt adhikāṛtham kṛdgrahaṇam kartavyam . \\
\noindent \textbf{(P\_6,2.139)} kṛtprakṛtiḥ dhātuḥ . \\
\noindent \textbf{(P\_6,2.139)} dhātum ca prati kriyāyogaḥ . \\
\noindent \textbf{(P\_6,2.139)} tatra yatkriyāyuktāḥ tam prati iti iha eva syāt . \\
\noindent \textbf{(P\_6,2.139)} praṇīḥ , unnīḥ . \\
\noindent \textbf{(P\_6,2.139)} iha na syāt . \\
\noindent \textbf{(P\_6,2.139)} praṇāyakaḥ , unnāyakaḥ . \\
\noindent \textbf{(P\_6,2.139)} etat api na asti prayojanam . \\
\noindent \textbf{(P\_6,2.139)} yatkriyāyuktāḥ iti na evam vijñāyate . \\
\noindent \textbf{(P\_6,2.139)} yasya kriyā yatkriyā yatkriyāyuktāḥ tam prati gatyupasargsañjñe bhavataḥ iti . \\
\noindent \textbf{(P\_6,2.139)} katham tarhi yā kriyā yatkriyā yatkriyāyuktāḥ tam prati gatyupasargsañjñe bhavataḥ iti . \\
\noindent \textbf{(P\_6,2.139)} na ca kaḥ cit kevalaḥ śabdaḥ asti yaḥ tasya arthasya vācakaḥ syāt . \\
\noindent \textbf{(P\_6,2.139)} kevalaḥ tasya arthasya vācakaḥ na asti iti kṛtvā kṛdadhikasya bhaviṣyati . \\
\noindent \textbf{(P\_6,2.139)} nanu ca yayam tasya eva arthasya vācakaḥ praṇīḥ iti . \\
\noindent \textbf{(P\_6,2.139)} eṣaḥ api hi kartṛviśiṣṭasya . \\
\noindent \textbf{(P\_6,2.139)} ayam tarhi tasya eva arthasya vācakaḥ prabhavanam iti . \\
\noindent \textbf{(P\_6,2.139)} tasmāt kṛdgrahaṇam kartavyam . \\
\noindent \textbf{(P\_6,2.139)} yadi kṛdgrahaṇam kriyate āmante svaraḥ na prāpnoti . \\
\noindent \textbf{(P\_6,2.139)} prapacatitarām , prajalpatitarām . \\
\noindent \textbf{(P\_6,2.139)} asati punaḥ kṛdgrahaṇe kriyāpradhānam ākhyātam tasya atiśaye tarap utpadyate tarabantasya svārthe ām . \\
\noindent \textbf{(P\_6,2.139)} tatra yatkriyāyuktāḥ iti bhavati eva saṅghātam prati kriyāyogaḥ . \\
\noindent \textbf{(P\_6,2.139)} na ca kaḥ cit kevalaḥ śabdaḥ asti yaḥ tasya arthasya vācakaḥ syāt . \\
\noindent \textbf{(P\_6,2.139)} kevalaḥ tasya arthasya vācakaḥ na asti iti kṛtvā adhikasya bhaviṣyati . \\
\noindent \textbf{(P\_6,2.139)} nanu ca yayam tasya eva arthasya vācakaḥ prabhavanam iti . \\
\noindent \textbf{(P\_6,2.139)} eṣaḥ api dravyaviśiṣṭasya . \\
\noindent \textbf{(P\_6,2.139)} katham kṛdabhihitaḥ bhāvaḥ dravyavat bhavati kriyāvat api iti . \\
\noindent \textbf{(P\_6,2.143)} kim samāsaya antaḥ udāttaḥ bhavati āhosvit uttarapadasya . \\
\noindent \textbf{(P\_6,2.143)} kutaḥ sandehaḥ . \\
\noindent \textbf{(P\_6,2.143)} ubhayam prakṛtam . \\
\noindent \textbf{(P\_6,2.143)} tatra anyatarat śakyam viśeṣayitum . \\
\noindent \textbf{(P\_6,2.143)} kaḥ ca atra viśeṣaḥ . \\
\noindent \textbf{(P\_6,2.143)} antodāttatvam samāsasya iti cet kapi upasaṅkhyānam . \\
\noindent \textbf{(P\_6,2.143)} antodāttatvam samāsasya iti cet kapi upasaṅkhyānam kartavyam . \\
\noindent \textbf{(P\_6,2.143)} idametattadbhyaḥ prathamapūraṇayoḥ kriyāgaṇane kapi ca iti vaktavyam iha api yathā syāt . \\
\noindent \textbf{(P\_6,2.143)} idamprathamakāḥ . \\
\noindent \textbf{(P\_6,2.143)} astu tarhi uttarapadasya . \\
\noindent \textbf{(P\_6,2.143)} uttarapadāntodāttatve nañsubhyām samāsāntodāttatvam . \\
\noindent \textbf{(P\_6,2.143)} uttarapadāntodāttatve nañsubhyām samāsāntodāttatvam vaktavyam . \\
\noindent \textbf{(P\_6,2.143)} anṛcaḥ , bahvṛcaḥ . \\
\noindent \textbf{(P\_6,2.143)} aparaḥ āha : uttarapadāntodāttatve nañsubhyām samāsāntodāttatvam vaktavyam . \\
\noindent \textbf{(P\_6,2.143)} ajñakaḥ , asvakaḥ . \\
\noindent \textbf{(P\_6,2.143)} kapi pūrvam iti asya apavādaḥ hrasvānte antyāt pūrvam iti . \\
\noindent \textbf{(P\_6,2.143)} tatra hrasvānte antyāt pūrvaḥ udāttabhāvī na asti iti kṛtvā utsargeṇa antodāttatvam prāpnoti . \\
\noindent \textbf{(P\_6,2.143)} na vā kapi pūrvavacanam jñāpakam uttarapadānantodāttatvasya . \\
\noindent \textbf{(P\_6,2.143)} na vā eṣaḥ doṣaḥ . \\
\noindent \textbf{(P\_6,2.143)} kim kāraṇam . \\
\noindent \textbf{(P\_6,2.143)} yat ayam kapi pūrvam iti āha tat jñāpayati ācāryaḥ na uttarapadasya antaḥ udāttatḥ bhavati iti . \\
\noindent \textbf{(P\_6,2.143)} prakaraṇāt ca samāsāntodāttatvam . \\
\noindent \textbf{(P\_6,2.143)} prakṛtam samāsagrahaṇam anuvartate . \\
\noindent \textbf{(P\_6,2.143)} kva prakṛtam . \\
\noindent \textbf{(P\_6,2.143)} cau samāsasya iti . \\
\noindent \textbf{(P\_6,2.143)} nanu ca uktam antodāttatvam samāsasya iti cet kapi upasaṅkhyānam iti . \\
\noindent \textbf{(P\_6,2.143)} na eṣaḥ doṣaḥ . \\
\noindent \textbf{(P\_6,2.143)} uttarapadagrahaṇam api prakṛtam anuvartate . \\
\noindent \textbf{(P\_6,2.143)} kva prakṛtam . \\
\noindent \textbf{(P\_6,2.143)} uttarapadādiḥ iti . \\
\noindent \textbf{(P\_6,2.143)} tatra evam abhisambandhaḥ kariṣyate . \\
\noindent \textbf{(P\_6,2.143)} nañsubhyām samāsasya antaḥ udāttaḥ bhavati . \\
\noindent \textbf{(P\_6,2.143)} idametattadbhyaḥ prathamapūraṇayoḥ kriyāgaṇane uttarapadasya iti . \\
\noindent \textbf{(P\_6,2.148)} kārakāt dattaśrutayoḥ anāśiṣi pratiṣedhaḥ . \\
\noindent \textbf{(P\_6,2.148)} kārakāt dattaśrutayoḥ anāśiṣi pratiṣedhaḥ vaktavyaḥ . \\
\noindent \textbf{(P\_6,2.148)} anāhataḥ nadati devadattaḥ . \\
\noindent \textbf{(P\_6,2.148)} siddham tu ubhayaniyamāt . \\
\noindent \textbf{(P\_6,2.148)} siddham etat . \\
\noindent \textbf{(P\_6,2.148)} katham . \\
\noindent \textbf{(P\_6,2.148)} ubhayaniyamāt . \\
\noindent \textbf{(P\_6,2.148)} ubhayataḥ niyamaḥ āśrayiṣyate . \\
\noindent \textbf{(P\_6,2.148)} kārakāt dattaśrutayoḥ eva āśiṣi. āśiṣi eva kārakāt dattaśrutayoḥ iti . \\
\noindent \textbf{(P\_6,2.165)} ṛṣipratiṣedhaḥ mitre . \\
\noindent \textbf{(P\_6,2.165)} ṛṣipratiṣedhaḥ mitre vaktavyaḥ . \\
\noindent \textbf{(P\_6,2.165)} viśvāmitraḥ ṛṣiḥ . \\
\noindent \textbf{(P\_6,2.175)} kimartham bahoḥ nañvat atideśaḥ kriyate na nañsubahubhyaḥ iti eva ucyeta . \\
\noindent \textbf{(P\_6,2.175)} na evam śakyam . \\
\noindent \textbf{(P\_6,2.175)} uttarapadabhūmni iti vakṣyati . \\
\noindent \textbf{(P\_6,2.175)} tat bahoḥ eva yathā syāt . \\
\noindent \textbf{(P\_6,2.175)} nañsubhyām mā bhūt iti . \\
\noindent \textbf{(P\_6,2.175)} na etat asti prayojanam . \\
\noindent \textbf{(P\_6,2.175)} ekayoge api hi sati yasya uttarapadabhūmā asti tasya bhaviṣyati . \\
\noindent \textbf{(P\_6,2.175)} kasya ca asti . \\
\noindent \textbf{(P\_6,2.175)} bahoḥ eva . \\
\noindent \textbf{(P\_6,2.175)} idam tarhi prayojanam . \\
\noindent \textbf{(P\_6,2.175)} na guṇādayaḥ avyavāḥ iti vakṣyati . \\
\noindent \textbf{(P\_6,2.175)} tat bahoḥ eva yathā syāt . \\
\noindent \textbf{(P\_6,2.175)} nañsubhyām mā bhūt iti . \\
\noindent \textbf{(P\_6,2.175)} etat api na asti prayojanam . \\
\noindent \textbf{(P\_6,2.175)} ekayoge api sati yasya guṇādayaḥ avayavā santi tasya kasya ca santi . \\
\noindent \textbf{(P\_6,2.175)} bahoḥ eva . \\
\noindent \textbf{(P\_6,2.175)} ataḥ uttaram paṭhati bahoḥ nañvat uttarapadādyudāttārtham . \\
\noindent \textbf{(P\_6,2.175)} bahoḥ nañvat atideśaḥ kriayte uttarapadādyudāttārtham . \\
\noindent \textbf{(P\_6,2.175)} uttarapadasya ādyudāttatvam yathā syāt . \\
\noindent \textbf{(P\_6,2.175)} nañaḥ jaramaramitramṛtāḥ . \\
\noindent \textbf{(P\_6,2.175)} ajaraḥ , amaraḥ , bahujaraḥ , bahumitraḥ . \\
\noindent \textbf{(P\_6,2.177)} upasargāt svāṅgam dhruvam mukhasya antodāttatvāt . \\
\noindent \textbf{(P\_6,2.177)} mukhasya antodāttatvāt upasargāt svāṅgam dhruvam iti etat bhavati vipratiṣedhena . \\
\noindent \textbf{(P\_6,2.177)} mukhāntodāttatvasya avakāśaḥ gauramukhaḥ , ślakṣṇamukhaḥ . \\
\noindent \textbf{(P\_6,2.177)} upasargāt svāṅgam iti asya avakāśaḥ prasphik , prodaraḥ . \\
\noindent \textbf{(P\_6,2.177)} iha ubhayam prāpnoti . \\
\noindent \textbf{(P\_6,2.177)} pramukhaḥ . \\
\noindent \textbf{(P\_6,2.177)} upasargāt svāṅgam iti etat bhavati vipratiṣedhena . \\
\noindent \textbf{(P\_6,2.177)} kaḥ punaḥ viśeṣaḥ tena vā sati anena vā . \\
\noindent \textbf{(P\_6,2.177)} sāpavādakaḥ saḥ vidhiḥ ayam punaḥ nirapavādakaḥ . \\
\noindent \textbf{(P\_6,2.177)} avyayāt tasya pratiṣedhaḥ apavādaḥ . \\
\noindent \textbf{(P\_6,2.185-186)} KA\_III,138.14-18 Ro\_IV,576 {1/6}     kimartham idam ucyate na upasargāt svāṅgam dhruvam iti eva siddham . \\
\noindent \textbf{(P\_6,2.185-186)} KA\_III,138.14-18 Ro\_IV,576 {2/6}     abheḥ mukham apāt ca adhruvārtham . \\
\noindent \textbf{(P\_6,2.185-186)} KA\_III,138.14-18 Ro\_IV,576 {3/6}     adhruvārthaḥ ayam ārambhaḥ . \\
\noindent \textbf{(P\_6,2.185-186)} KA\_III,138.14-18 Ro\_IV,576 {4/6}     abhuvrīhyartham vā . \\
\noindent \textbf{(P\_6,2.185-186)} KA\_III,138.14-18 Ro\_IV,576 {5/6}     atha vā bahuvrīheḥ iti vartate . \\
\noindent \textbf{(P\_6,2.185-186)} KA\_III,138.14-18 Ro\_IV,576 {6/6}     abhuvrīhyarthaḥ ayam ārambhaḥ . \\
\noindent \textbf{(P\_6,2.187)} sphigapūtagrahaṇam kimartham na upasargāt svāṅgam dhruvam iti eva siddham . \\
\noindent \textbf{(P\_6,2.187)} sphigapūtagrahaṇam ca . \\
\noindent \textbf{(P\_6,2.187)} kim . \\
\noindent \textbf{(P\_6,2.187)} adhruvārtham abhuvrīhyartham eva vā . \\
\noindent \textbf{(P\_6,2.191)} ateḥ dhātulope . \\
\noindent \textbf{(P\_6,2.191)} ateḥ dhātulope iti vaktavyam . \\
\noindent \textbf{(P\_6,2.191)} akṛtpade iti hi ucyamāne iha ca prasajyeta śobhanaḥ gārgyaḥ atigārgyaḥ , iha ca na syāt , atikāṛakaḥ, atipadā śakvarī . \\
\noindent \textbf{(P\_6,2.197)} kim idam dvitribhyām mūrdhani akārāntagrahaṇam āhosvit nakārāntagrahaṇam . \\
\noindent \textbf{(P\_6,2.197)} kaḥ ca atra viśeṣaḥ . \\
\noindent \textbf{(P\_6,2.197)} dvitribhyām mūrdhani akārāntagrahaṇam cet nakārāntasya upasaṅkhyānam . \\
\noindent \textbf{(P\_6,2.197)} dvitribhyām mūrdhani akārāntagrahaṇam cet nakārāntasya upasaṅkhyānam kartavyam . \\
\noindent \textbf{(P\_6,2.197)} dvimūrdhā trimūrdhā . \\
\noindent \textbf{(P\_6,2.197)} astu tarhi nakārāntagrahaṇam . \\
\noindent \textbf{(P\_6,2.197)} nakārānte akārāntasya . \\
\noindent \textbf{(P\_6,2.197)} nakārānte akārāntasya upasaṅkhyānam kartavyam . \\
\noindent \textbf{(P\_6,2.197)} dvimūrdhaḥ , trimūrdhaḥ . \\
\noindent \textbf{(P\_6,2.197)} udāttalopāt siddham . \\
\noindent \textbf{(P\_6,2.197)} astu tarhi nakārāntagrahaṇam . \\
\noindent \textbf{(P\_6,2.197)} antodāttatve kṛte lopaḥ udāttanivṛttisvareṇa siddham . \\
\noindent \textbf{(P\_6,2.197)} idam iha sampradhāryam . \\
\noindent \textbf{(P\_6,2.197)} antodāttatvam kriyatām lopaḥ iti kim atra kartavyam . \\
\noindent \textbf{(P\_6,2.197)} paratvāt lopaḥ . \\
\noindent \textbf{(P\_6,2.197)} evam tarhi idam iha sampradhāryam . \\
\noindent \textbf{(P\_6,2.197)} antodāttatvam kriyatām samāsāntaḥ iti kim atra kartavyam . \\
\noindent \textbf{(P\_6,2.197)} paratvāt antodāttatvam . \\
\noindent \textbf{(P\_6,2.197)} nityaḥ samāsāntaḥ . \\
\noindent \textbf{(P\_6,2.197)} kṛte api antodāttatve prāpnoti akṛte api . \\
\noindent \textbf{(P\_6,2.197)} antodāttatvam api nityam . \\
\noindent \textbf{(P\_6,2.197)} kṛte api samāsānte prāpnoti akṛte api . \\
\noindent \textbf{(P\_6,2.197)} anityam antodāttatvam . \\
\noindent \textbf{(P\_6,2.197)} na hi kṛte samāsānte prāpnoti . \\
\noindent \textbf{(P\_6,2.197)} paratvāt lopena bhavitavyam . \\
\noindent \textbf{(P\_6,2.197)} yasya ca lakṣaṇāntareṇa nimittam vihanyate na tat anityam . \\
\noindent \textbf{(P\_6,2.197)} na ca samāsāntaḥ eva antodāttatvasya nimittam hanti . \\
\noindent \textbf{(P\_6,2.197)} avaśyam lakṣaṇāntaram lopaḥ pratīkṣyaḥ . \\
\noindent \textbf{(P\_6,2.197)} ubhayoḥ nityayoḥ paratvāt antodāttatvam . \\
\noindent \textbf{(P\_6,2.197)} antodāttatve kṛte samāsāntaḥ , ṭilopaḥ . \\
\noindent \textbf{(P\_6,2.197)} ṭilope kṛte udāttanivṛttisvareṇa siddham . \\
\noindent \textbf{(P\_6,2.197)} yuktam punaḥ idam vicārayitum . \\
\noindent \textbf{(P\_6,2.197)} nan u anena asandigdhena nakārāntasya grahaṇena bhavitavyam yāvatā mūrdhasu iti ucyate . \\
\noindent \textbf{(P\_6,2.197)} yadi hi akārāntasya grahaṇam syāt mūrdheṣu iti brūyāt . \\
\noindent \textbf{(P\_6,2.197)} sā eṣā samāsāntārthā vicāraṇā . \\
\noindent \textbf{(P\_6,2.197)} evam tarhi jñāpayati ācāryaḥ . \\
\noindent \textbf{(P\_6,2.197)} vibhāṣā samāsāntaḥ bhavati iti . \\
\noindent \textbf{(P\_6,2.199)} atyalpam idam ucyate . \\
\noindent \textbf{(P\_6,2.199)} parādiḥ ca parāntaḥ ca pūrvāntaḥ ca dṛśyate .pūrvādayaḥ ca vidyante . \\
\noindent \textbf{(P\_6,2.199)} vyatayaḥ bahulam smṛtaḥ . \\
\noindent \textbf{(P\_6,2.199)} antodāttaprakaraṇe tricakrādīnām chandasi upasaṅkhyānam . \\
\noindent \textbf{(P\_6,2.199)} antodāttaprakaraṇe tricakrādīnām chandasi upasaṅkhyānam kartavyam : tricakreṇa tribandhureṇa trivṛtā rathena . \\
\noindent \textbf{(P\_6,3.1.2)} ekavat ca aluk bhavati iti vaktavyam . \\
\noindent \textbf{(P\_6,3.1.2)} kim prayojanam . \\
\noindent \textbf{(P\_6,3.1.2)} stokābhyām muktaḥ , stokebhyaḥ muktaḥ iti vigṛhya stokānmuktaḥ iti eva yathā syāt . \\
\noindent \textbf{(P\_6,3.1.2)} ekavadvcanam anarthakam . \\
\noindent \textbf{(P\_6,3.1.2)} ekavadbhāvaḥ ca anarthakaḥ . \\
\noindent \textbf{(P\_6,3.1.2)} dvibahvoḥ aluk kasmāt na bhavati . \\
\noindent \textbf{(P\_6,3.1.2)} dvibahuṣu asamāsaḥ . \\
\noindent \textbf{(P\_6,3.1.2)} dvivacanabahuvacanānām asamāsaḥ . \\
\noindent \textbf{(P\_6,3.1.2)} kim vaktavyam etat . \\
\noindent \textbf{(P\_6,3.1.2)} na hi anucyamānaṃ gaṃsyate . \\
\noindent \textbf{(P\_6,3.1.2)} uktam vā . \\
\noindent \textbf{(P\_6,3.1.2)} kim uktam . \\
\noindent \textbf{(P\_6,3.1.2)} anabhidhānāt iti . \\
\noindent \textbf{(P\_6,3.1.2)} tat ca avaśyam anabhidhānam āśrayitavyam . \\
\noindent \textbf{(P\_6,3.1.2)} ekavadvacane hi goṣucare atiprasaṅgaḥ . \\
\noindent \textbf{(P\_6,3.1.2)} ekavadvacane hi goṣucare atiprasaṅgaḥ syāt : goṣucaraḥ . \\
\noindent \textbf{(P\_6,3.1.2)} varṣābhyaḥ ca je . \\
\noindent \textbf{(P\_6,3.1.2)} varṣābhyaḥ ca je atiprasaṅgaḥ bhavati . \\
\noindent \textbf{(P\_6,3.1.2)} varṣāsujaḥ . \\
\noindent \textbf{(P\_6,3.1.2)} apaḥ yoniyanmatiṣu ca . \\
\noindent \textbf{(P\_6,3.1.2)} apaḥ yoniyanmatiṣu ca upasaṅkhyānam kartavyam . \\
\noindent \textbf{(P\_6,3.1.2)} je care ca . \\
\noindent \textbf{(P\_6,3.1.2)} je care ca atiprasaṅgaḥ bhavati . \\
\noindent \textbf{(P\_6,3.1.2)} yoni . \\
\noindent \textbf{(P\_6,3.1.2)} apsuyoniḥ . \\
\noindent \textbf{(P\_6,3.1.2)} yat . \\
\noindent \textbf{(P\_6,3.1.2)} apsavyam . \\
\noindent \textbf{(P\_6,3.1.2)} mati . \\
\noindent \textbf{(P\_6,3.1.2)} apsumatiḥ . \\
\noindent \textbf{(P\_6,3.1.2)} je . \\
\noindent \textbf{(P\_6,3.1.2)} apsujaḥ . \\
\noindent \textbf{(P\_6,3.1.2)} care . \\
\noindent \textbf{(P\_6,3.1.2)} apsucaraḥ gahvareṣṭhāḥ . \\
\noindent \textbf{(P\_6,3.2)} pañcamīprakaraṇe brāhmaṇācchaṃsinaḥ upasaṅkhyānam . \\
\noindent \textbf{(P\_6,3.2)} pañcamīprakaraṇe brāhmaṇācchaṃsinaḥ upasaṅkhyānam kartavyam . \\
\noindent \textbf{(P\_6,3.2)} brāhmaṇācchaṃsī . anyārthe ca . \\
\noindent \textbf{(P\_6,3.2)} anyārthe ca eṣā pañcamī draṣṭavyā . \\
\noindent \textbf{(P\_6,3.2)} brāhmaṇāni śaṃsati iti brāhmaṇācchaṃsī . \\
\noindent \textbf{(P\_6,3.2)} atha vā yuktaḥ eva atra pañcamyarthaḥ . \\
\noindent \textbf{(P\_6,3.2)} brāhmaṇebhyaḥ gṛhītvā , āhṛtya āhṛtya śaṃsati iti brāhmaṇācchaṃsī . \\
\noindent \textbf{(P\_6,3.3)} añjasaḥ upasaṅkhyānam . \\
\noindent \textbf{(P\_6,3.3)} añjasaḥ upasaṅkhyānam kartavyam . \\
\noindent \textbf{(P\_6,3.3)} añjasākṛtam . \\
\noindent \textbf{(P\_6,3.3)} puṃsānujaḥ januṣāndhaḥ vikṛtākṣaḥ iti ca . \\
\noindent \textbf{(P\_6,3.3)} puṃsānujaḥ januṣāndhaḥ vikṛtākṣaḥ iti ca upasaṅkhyānam kartavyam . \\
\noindent \textbf{(P\_6,3.3)} puṃsānujaḥ , januṣāndhaḥ , vikṛtākṣaḥ . \\
\noindent \textbf{(P\_6,3.5)} ātmanaḥ ca pūraṇe . \\
\noindent \textbf{(P\_6,3.5)} ātmanaḥ ca pūraṇe upasaṅkhyānam kartavyam . \\
\noindent \textbf{(P\_6,3.5)} ātmanāpañcamaḥ , ātmanādaśamaḥ . \\
\noindent \textbf{(P\_6,3.5)} anyārthe ca . \\
\noindent \textbf{(P\_6,3.5)} anyārthe ca eṣā tṛtīyā draṣṭavyā . \\
\noindent \textbf{(P\_6,3.5)} ātmā pañcamaḥ asya ātmāpañcamaḥ . \\
\noindent \textbf{(P\_6,3.5)} atha vā yuktaḥ eva atra tṛtīyārthaḥ . \\
\noindent \textbf{(P\_6,3.5)} ātmanā kṛtam tat tasya yena asau pañcamaḥ . \\
\noindent \textbf{(P\_6,3.5)} katham janārdanaḥ tu ātmacaturthaḥ eva iti . \\
\noindent \textbf{(P\_6,3.5)} bahuvrīḥ ayam . \\
\noindent \textbf{(P\_6,3.5)} ātmā caturthaḥ asya iti . \\
\noindent \textbf{(P\_6,3.7-8)} KA\_III,143.11-18 Ro\_IV,586-587 {1/14}     ātmanebhāṣaparasmaibhāṣayoḥ upasaṅkhyānam . \\
\noindent \textbf{(P\_6,3.7-8)} KA\_III,143.11-18 Ro\_IV,586-587 {2/14}     ātmanebhāṣaparasmaibhāṣayoḥ upasaṅkhyānam kartavyam . \\
\noindent \textbf{(P\_6,3.7-8)} KA\_III,143.11-18 Ro\_IV,586-587 {3/14}     ātmanebhāṣaḥ , parasmaibhāṣaḥ . \\
\noindent \textbf{(P\_6,3.7-8)} KA\_III,143.11-18 Ro\_IV,586-587 {4/14}     tat katham kartavyam . \\
\noindent \textbf{(P\_6,3.7-8)} KA\_III,143.11-18 Ro\_IV,586-587 {5/14}     yadi vyākaraṇe bhavā vaiyākaraṇī , vaiyākaraṇī ākhyā vaiyākaraṇākhyā vaiyākaraṇakhyāyām iti . \\
\noindent \textbf{(P\_6,3.7-8)} KA\_III,143.11-18 Ro\_IV,586-587 {6/14}     atha hi vaiyākaraṇānām ākhyā vaiyākaraṇākhyā na arthaḥ upasaṅkhyānena . \\
\noindent \textbf{(P\_6,3.7-8)} KA\_III,143.11-18 Ro\_IV,586-587 {7/14}     yadi api vyākaraṇe bhavā vaiyākaraṇī , vaiyākaraṇī ākhyā vaiyākaraṇākhyā evam api na arthaḥ upasaṅkhyānena . \\
\noindent \textbf{(P\_6,3.7-8)} KA\_III,143.11-18 Ro\_IV,586-587 {8/14}     vacanāt bhaviṣyati . \\
\noindent \textbf{(P\_6,3.7-8)} KA\_III,143.11-18 Ro\_IV,586-587 {9/14}     asti vacane prayojanam . \\
\noindent \textbf{(P\_6,3.7-8)} KA\_III,143.11-18 Ro\_IV,586-587 {10/14}     kim . \\
\noindent \textbf{(P\_6,3.7-8)} KA\_III,143.11-18 Ro\_IV,586-587 {11/14}     ātmanepadam . \\
\noindent \textbf{(P\_6,3.7-8)} KA\_III,143.11-18 Ro\_IV,586-587 {12/14}     nipātanāt etat siddham . \\
\noindent \textbf{(P\_6,3.7-8)} KA\_III,143.11-18 Ro\_IV,586-587 {13/14}     kim nipātanam . \\
\noindent \textbf{(P\_6,3.7-8)} KA\_III,143.11-18 Ro\_IV,586-587 {14/14}     anudāttaṅitaḥ ātmanepadam , śeṣāt kartari parasmaipadam iti . \\
\noindent \textbf{(P\_6,3.9)} hṛddyubhyām ṅeḥ upasaṅkhyānam . \\
\noindent \textbf{(P\_6,3.9)} hṛddyubhyām ṅeḥ upasaṅkhyānam kartavyam . \\
\noindent \textbf{(P\_6,3.9)} hṛdispṛk , divispṛk . \\
\noindent \textbf{(P\_6,3.9)} anyārthe ca . \\
\noindent \textbf{(P\_6,3.9)} anyārthe ca eṣā saptamī draṣṭavyā . \\
\noindent \textbf{(P\_6,3.9)} hṛdayam spṛśati iti hṛdispṛk . \\
\noindent \textbf{(P\_6,3.9)} divam śprśati iti divispṛk . \\
\noindent \textbf{(P\_6,3.9)} haladantādhikāre goḥ upasaṅkhyānam . \\
\noindent \textbf{(P\_6,3.9)} haladantādhikāre goḥ upasaṅkhyānam kartavyam . \\
\noindent \textbf{(P\_6,3.9)} gaviṣthiraḥ . \\
\noindent \textbf{(P\_6,3.9)} na kartavyam . \\
\noindent \textbf{(P\_6,3.9)} lukaḥ avādeśaḥ vipratiṣedhena . \\
\noindent \textbf{(P\_6,3.9)} luk kriyatām avādeśaḥ iti avādeśaḥ bhaviṣyati vipratiṣedhena . \\
\noindent \textbf{(P\_6,3.9)} avādeśe kṛte halantāt iti eva siddham . \\
\noindent \textbf{(P\_6,3.9)} lukaḥ avādeśaḥ vipratiṣedhena iti cet bhūmipāśe atiprasaṅgaḥ . \\
\noindent \textbf{(P\_6,3.9)} lukaḥ avādeśaḥ vipratiṣedhena iti cet bhūmipāśe atiprasaṅgaḥ bhavati . \\
\noindent \textbf{(P\_6,3.9)} bhūmyām pāśaḥ , bhūmipāśaḥ . \\
\noindent \textbf{(P\_6,3.9)} akaḥ ataḥ iti vā sandhyakṣarārtham . \\
\noindent \textbf{(P\_6,3.9)} evam tarhi aviśeṣeṇa saptamyāḥ alukam uktvā akaḥ ataḥ iti vakṣyāmi . \\
\noindent \textbf{(P\_6,3.9)} tat niyamārtham bhaviṣyati . \\
\noindent \textbf{(P\_6,3.9)} akaḥ ataḥ iti eva bhavati na anyataḥ iti . \\
\noindent \textbf{(P\_6,3.9)} tena sandhyakṣarāṇām siddham bhavati . \\
\noindent \textbf{(P\_6,3.9)} sidhyati . \\
\noindent \textbf{(P\_6,3.9)} sūtram tarhi bhidyate . \\
\noindent \textbf{(P\_6,3.9)} yathānyāsam eva astu . \\
\noindent \textbf{(P\_6,3.9)} nanu ca uktam haladantādhikāre goḥ upasaṅkhyānam iti . \\
\noindent \textbf{(P\_6,3.9)} na eṣaḥ doṣaḥ . \\
\noindent \textbf{(P\_6,3.9)} nipātanāt etat siddham . \\
\noindent \textbf{(P\_6,3.9)} kim nipātanam . \\
\noindent \textbf{(P\_6,3.9)} gaviṣṭhiraśabdaḥ vidādiṣu paṭhyate . \\
\noindent \textbf{(P\_6,3.9)} asakṛt khalu api nipātanam kriyate . \\
\noindent \textbf{(P\_6,3.9)} gaviyudhibhyā sthiraḥ iti . \\
\noindent \textbf{(P\_6,3.10)} kim iyam prāpte vibhāṣā āhosvit aprāpte . \\
\noindent \textbf{(P\_6,3.10)} katham ca prāpte katham ca aprāpte . \\
\noindent \textbf{(P\_6,3.10)} yadi sañjñāyām iti vartate tataḥ prāpte . \\
\noindent \textbf{(P\_6,3.10)} atha nivṛttam tataḥ aprāpte . \\
\noindent \textbf{(P\_6,3.10)} kaḥ ca atra viśeṣaḥ . \\
\noindent \textbf{(P\_6,3.10)} kāranāmni vāvacanārtham cet ajādau atiprasaṅgaḥ . \\
\noindent \textbf{(P\_6,3.10)} kāranāmni vāvacanārtham cet ajādau atiprasaṅgaḥ bhavati . \\
\noindent \textbf{(P\_6,3.10)} iha api prāpnoti . \\
\noindent \textbf{(P\_6,3.10)} avikaṭe uraṇaḥ dātavyaḥ avikaṭoraṇaḥ . \\
\noindent \textbf{(P\_6,3.10)} astu tarhi aprāpte . \\
\noindent \textbf{(P\_6,3.10)} aprāpte samāsavidhānam . \\
\noindent \textbf{(P\_6,3.10)} yadi aprāpte samāsaḥ vidheyaḥ . \\
\noindent \textbf{(P\_6,3.10)} prāpte punaḥ sati sañjñāyām iti eva samāsaḥ siddhaḥ . \\
\noindent \textbf{(P\_6,3.10)} na eṣaḥ doṣaḥ . \\
\noindent \textbf{(P\_6,3.10)} etat eva jñāpayati bhavati atra samāsaḥ iti yat ayam kāranāmni saptamyāḥ alukam śāsti . \\
\noindent \textbf{(P\_6,3.10)} yadi api tāvat jñāpakāt samāsaḥ syāt svaraḥ tu na sidhyati . \\
\noindent \textbf{(P\_6,3.10)} yat hi tat saptamīpūrvapadam prakṛtisvaram bhavati iti lakṣaṇapratipadoktayoḥ pratipadoktasya eva iti evam tat . \\
\noindent \textbf{(P\_6,3.10)} na eva atra anena svareṇa bhavitavyam . \\
\noindent \textbf{(P\_6,3.10)} kim tarhi saptamīhāriṇau dharmye aharaṇe iti anena atra svareṇa bhavitavyam .kim ca bhoḥ sañjñāḥ api loke kriyante na lokaḥ sañjñāsu pramāṇam . \\
\noindent \textbf{(P\_6,3.10)} loke ca kāranāma sañjñā . \\
\noindent \textbf{(P\_6,3.10)} nanu ca uktam kāranāmni vāvacanārtham cet ajādau atiprasaṅgaḥ iti . \\
\noindent \textbf{(P\_6,3.10)} na eṣaḥ doṣaḥ . \\
\noindent \textbf{(P\_6,3.10)} yogavibhāgāt siddham . \\
\noindent \textbf{(P\_6,3.10)} yogavibhāgaḥ kariṣyate . \\
\noindent \textbf{(P\_6,3.10)} kāranāmni ca prācām , tataḥ halādau . \\
\noindent \textbf{(P\_6,3.10)} halādau ca kāranāmni saptamyāḥ aluk bhavati . \\
\noindent \textbf{(P\_6,3.10)} idam idānīm kimartham . \\
\noindent \textbf{(P\_6,3.10)} niyamārtham . \\
\noindent \textbf{(P\_6,3.10)} halādau eva kāranāmni na anyatra . \\
\noindent \textbf{(P\_6,3.10)} kva mā bhūt . \\
\noindent \textbf{(P\_6,3.10)} avikaṭe uraṇaḥ dātavyaḥ avikaṭoraṇaḥ . \\
\noindent \textbf{(P\_6,3.11)} gurau antāt ca . \\
\noindent \textbf{(P\_6,3.11)} gurau antāt ca iti vaktavyam . \\
\noindent \textbf{(P\_6,3.11)} anteguruḥ . \\
\noindent \textbf{(P\_6,3.13)} svāṅgagrahaṇam anuvartate utāho na . \\
\noindent \textbf{(P\_6,3.13)} kim ca ataḥ . \\
\noindent \textbf{(P\_6,3.13)} yadi anuvartate siddham hastebandhaḥ , hastabandhaḥ . \\
\noindent \textbf{(P\_6,3.13)} cakrebhandaḥ , cakrabandhaḥ iti na sidhyati . \\
\noindent \textbf{(P\_6,3.13)} atha nivṛttam siddham cakrebhandaḥ , cakrabandhaḥ . \\
\noindent \textbf{(P\_6,3.13)} hastebandhaḥ , hastabandhaḥ iti na sidhyati . \\
\noindent \textbf{(P\_6,3.13)} kim kāraṇam . \\
\noindent \textbf{(P\_6,3.13)} na insiddhabadhnātiṣu iti pratiṣedhaḥ prāpnoti . \\
\noindent \textbf{(P\_6,3.13)} na eṣaḥ doṣaḥ . \\
\noindent \textbf{(P\_6,3.13)} sarvatra eva atra uttarapadādhikare tatpuruṣed kṛti bahulam iti prāpte na insiddhabadhnātiṣu iti pratiṣedhaḥ ucyate . \\
\noindent \textbf{(P\_6,3.13)} tasmin nitye prāpte iyam vibhāṣā ārabhyate . \\
\noindent \textbf{(P\_6,3.13)} evam api na jñāyate kasmin viṣaye vibhāṣā kasmin viṣaye pratiṣedhaḥ iti . \\
\noindent \textbf{(P\_6,3.13)} ghañantasya idam bandhaśabdasya grahaṇam pratiṣedhe punaḥ dhātugrahaṇam . \\
\noindent \textbf{(P\_6,3.13)} ghañante vibhāṣā anyatra pratiṣedhaḥ . \\
\noindent \textbf{(P\_6,3.14)} tatpuruṣe kṛti bahulam akarmadhāraye . \\
\noindent \textbf{(P\_6,3.14)} tatpuruṣe kṛti bahulam iti atra akarmadhāraye iti vaktavyam . \\
\noindent \textbf{(P\_6,3.14)} iha mā bhūt . \\
\noindent \textbf{(P\_6,3.14)} parame kārake paramakārake iti . \\
\noindent \textbf{(P\_6,3.14)} tat tarhi vaktavyam . \\
\noindent \textbf{(P\_6,3.14)} na vaktavyam . \\
\noindent \textbf{(P\_6,3.14)} bahulavacanāt na bhaviṣyati . \\
\noindent \textbf{(P\_6,3.14)} atha kimartham lugaluganukramaṇam kriyate na tatpuruṣe kṛti bahulam iti eva siddham . \\
\noindent \textbf{(P\_6,3.14)} lugaluganukramaṇam bahulavacanasya akṛtsnatvāt . \\
\noindent \textbf{(P\_6,3.14)} lugaluganukramaṇam kriyate akṛtsnam bahulavacanam iti . \\
\noindent \textbf{(P\_6,3.14)} yadi akṛtsnam yat anena kṛtam akṛtam tat . \\
\noindent \textbf{(P\_6,3.14)} evam tarhi na brūmaḥ akṛtsnam iti . \\
\noindent \textbf{(P\_6,3.14)} kṛtsnam ca kārakam ca sādhakam ca nirvartakam ca . \\
\noindent \textbf{(P\_6,3.14)} yat ca anena kṛtam suktṛtam tat . \\
\noindent \textbf{(P\_6,3.14)} kimartham tarhi lugaluganukramaṇam kriyate . \\
\noindent \textbf{(P\_6,3.14)} udāharaṇabhūyastvāt . \\
\noindent \textbf{(P\_6,3.14)} te khalu api vidhayaḥ suparigṛhītāḥ bhavanti yeṣu lakṣaṇam prapañcaḥ ca . \\
\noindent \textbf{(P\_6,3.14)} kevalam lakṣaṇam kevalaḥ prapañcaḥ vā na tathā kārakam bhavati . \\
\noindent \textbf{(P\_6,3.14)} avaśyam khalu asmābhiḥ idam vaktavyam bahulam anyatarasyām ubhayathā vā ekeṣām iti . \\
\noindent \textbf{(P\_6,3.14)} sarvavedapāṛiṣadam hi idam śāstram . \\
\noindent \textbf{(P\_6,3.14)} tatra na ekaḥ panthāḥ śakya āsthātum . \\
\noindent \textbf{(P\_6,3.21)} ṣaṣṭhīprakaraṇe vāgdikpaśyadbhyaḥ yuktidaṇḍahareṣu upasaṅkhyānam . \\
\noindent \textbf{(P\_6,3.21)} ṣaṣṭhīprakaraṇe vāgdikpaśyadbhyaḥ yuktidaṇḍahareṣu upasaṅkhyānam kartavyam . \\
\noindent \textbf{(P\_6,3.21)} vācoyuktiḥ , diśodaṇḍaḥ , paśyatoharaḥ . \\
\noindent \textbf{(P\_6,3.21)} āmuṣyāyaṇāmuṣyputrikā iti upasaṅkhyānam . \\
\noindent \textbf{(P\_6,3.21)} āmuṣyāyaṇāmuṣyputrikā iti upasaṅkhyānam kartavyam . \\
\noindent \textbf{(P\_6,3.21)} āmuṣyāyaṇaḥ , āmuṣyaputrikā . \\
\noindent \textbf{(P\_6,3.21)} āmuṣyakulikā iti ca vaktavyam . \\
\noindent \textbf{(P\_6,3.21)} āmuṣyakulikā . \\
\noindent \textbf{(P\_6,3.21)} devānāmpriyaḥ iti ca . \\
\noindent \textbf{(P\_6,3.21)} devānāmpriyaḥ iti ca upasaṅkhyānam kartavyam . \\
\noindent \textbf{(P\_6,3.21)} devānāmpriyaḥ . \\
\noindent \textbf{(P\_6,3.21)} śepapucchalāṅgūleṣu śunaḥ sañjñāyām . \\
\noindent \textbf{(P\_6,3.21)} śepapucchalāṅgūleṣu śunaḥ sañjñāyām upasaṅkhyānam kartavyam . \\
\noindent \textbf{(P\_6,3.21)} śunaḥśephaḥ , śunaḥpucchaḥ , śunolāṅgūlaḥ . \\
\noindent \textbf{(P\_6,3.21)} divaḥ ca dāse . \\
\noindent \textbf{(P\_6,3.21)} divaḥ ca dāse upasaṅkhyānam kartavyam . \\
\noindent \textbf{(P\_6,3.21)} divodāsāya gāyata . \\
\noindent \textbf{(P\_6,3.23)} vidyāyonisambandhebhyaḥ tatpūrvapadottarapadagrahaṇam . \\
\noindent \textbf{(P\_6,3.23)} vidyāyonisambandhebhyaḥ tatpūrvapadottarapadagrahaṇam kartavyam , vidyāsambandhebhyaḥ vidyāsambandheṣu yathā syāt , yonisambandhebhyaḥ yonisambandheṣu yathā syāt , vyatikaraḥ mā bhūt . \\
\noindent \textbf{(P\_6,3.23)} atha eṣām vyatikareṇa bhavitavyam . \\
\noindent \textbf{(P\_6,3.23)} bāḍham bhavitavyam . \\
\noindent \textbf{(P\_6,3.23)} hotuḥputraḥ , pituḥantevāsī . \\
\noindent \textbf{(P\_6,3.25.1)} kva ayam nakāraḥ śrūyate . \\
\noindent \textbf{(P\_6,3.25.1)} na kva cit śrūyate . \\
\noindent \textbf{(P\_6,3.25.1)} lopaḥ asya bhavati nalopaḥ prātipadikasya iti . \\
\noindent \textbf{(P\_6,3.25.1)} yadi na śrūyate kimartham uccāryate . \\
\noindent \textbf{(P\_6,3.25.1)} raparatvam mā bhūt iti . \\
\noindent \textbf{(P\_6,3.25.1)} kriyamāṇe api vai nakāre raparatvam prāpnoti . \\
\noindent \textbf{(P\_6,3.25.1)} kim kāraṇam . \\
\noindent \textbf{(P\_6,3.25.1)} nalope kṛte eṣaḥ api hi uḥ sthāne aṇ śiṣyate . \\
\noindent \textbf{(P\_6,3.25.1)} na eṣaḥ doṣaḥ . \\
\noindent \textbf{(P\_6,3.25.1)} uḥ sthāne aṇ prasajymānaḥ eva raparaḥ bhavati iti ucyate na ca ayam uḥ sthāne aṇ eva śiṣyate . \\
\noindent \textbf{(P\_6,3.25.1)} kim tarhi aṇ ca anaṇ ca . \\
\noindent \textbf{(P\_6,3.25.2)} katham punaḥ idam vijñāyate : ṛkārāntānām yaḥ dvandvaḥ iti āhosvit dvandve ṛkārasya iti . \\
\noindent \textbf{(P\_6,3.25.2)} kaḥ ca atra viśeṣaḥ . \\
\noindent \textbf{(P\_6,3.25.2)} ṛkārāntānām dvandve putre upasaṅkhyānam . \\
\noindent \textbf{(P\_6,3.25.2)} ṛkārāntānām dvandve putre upasaṅkhyānam kartavyam . \\
\noindent \textbf{(P\_6,3.25.2)} pitāputrau . \\
\noindent \textbf{(P\_6,3.25.2)} kāryī ca anirdiṣṭaḥ . \\
\noindent \textbf{(P\_6,3.25.2)} kāryī ca anirdiṣṭaḥ bhavati . \\
\noindent \textbf{(P\_6,3.25.2)} ṛkārāntānām dvandve na jñāyate kasya ānaṅā bhavitavyam iti . \\
\noindent \textbf{(P\_6,3.25.2)} astu tarhi dvande ṛkārasya iti . \\
\noindent \textbf{(P\_6,3.25.2)} aviśeṣeṇa pitṛpitāmahādiṣu atiprasaṅgaḥ . \\
\noindent \textbf{(P\_6,3.25.2)} aviśeṣeṇa pitṛpitāmahādiṣu atiprasaṅgaḥ bhavati . \\
\noindent \textbf{(P\_6,3.25.2)} pitṛpitāmahau iti . \\
\noindent \textbf{(P\_6,3.25.2)} astu tarhi ṛkārāntānām yaḥ dvandvaḥ iti . \\
\noindent \textbf{(P\_6,3.25.2)} nanu ca uktam ṛkārāntānām dvandve putre upasaṅkhyānam iti . \\
\noindent \textbf{(P\_6,3.25.2)} na eṣaḥ doṣaḥ . \\
\noindent \textbf{(P\_6,3.25.2)} putragrahaṇam api prakṛtam anuvartate . \\
\noindent \textbf{(P\_6,3.25.2)} kva prakṛtam . \\
\noindent \textbf{(P\_6,3.25.2)} putre anyatarasyām iti . \\
\noindent \textbf{(P\_6,3.25.2)} yadi tat anuvartate vibhāṣā svasṛpatyoḥ putre ca iti putre api vibhāṣā prāpnoti . \\
\noindent \textbf{(P\_6,3.25.2)} na eṣaḥ doṣaḥ . \\
\noindent \textbf{(P\_6,3.25.2)} sambandham anuvartiṣyate . \\
\noindent \textbf{(P\_6,3.25.2)} ṣaṣṭhyāḥ ākrośe . \\
\noindent \textbf{(P\_6,3.25.2)} putre anyatarasyām ṣaṣṭhyāḥ ākrośe . \\
\noindent \textbf{(P\_6,3.25.2)} ṛtaḥ vidyāyonisambandhebhyaḥ putre anyatarasyām ṣaṣṭhyāḥ ākrośe . \\
\noindent \textbf{(P\_6,3.25.2)} vibhāṣā svasṛpatyoḥ putre anyatarasyām ṣaṣṭhyāḥ ākrośe . \\
\noindent \textbf{(P\_6,3.25.2)} ānaṅ ṛtaḥ dvandve . \\
\noindent \textbf{(P\_6,3.25.2)} putragrahaṇam anuvartate . \\
\noindent \textbf{(P\_6,3.25.2)} ṣaṣṭhyāḥ ākrośe iti nivṛttam . \\
\noindent \textbf{(P\_6,3.25.2)} yat api ucyate kāryī ca anirdiṣṭaḥ iti . \\
\noindent \textbf{(P\_6,3.25.2)} kāryī ca nirdiṣṭaḥ . \\
\noindent \textbf{(P\_6,3.25.2)} katham . \\
\noindent \textbf{(P\_6,3.25.2)} uttarapade iti vartate . \\
\noindent \textbf{(P\_6,3.25.2)} ṅit ca ayam kriyate . \\
\noindent \textbf{(P\_6,3.25.2)} saḥ antareṇa api kāryinirdeśam ṛkārāntasya eva bhaviṣyati . \\
\noindent \textbf{(P\_6,3.25.2)} putre tarhi kāryī anirdiṣṭaḥ . \\
\noindent \textbf{(P\_6,3.25.2)} putre ca kāryī nirdiṣṭaḥ . \\
\noindent \textbf{(P\_6,3.25.2)} katham . \\
\noindent \textbf{(P\_6,3.25.2)} ṛkāragrahaṇam api prakṛtam anuvartate . \\
\noindent \textbf{(P\_6,3.25.2)} kva prakṛtam . \\
\noindent \textbf{(P\_6,3.25.2)} ṛtaḥ vidyāyonisambandhebhyaḥ . \\
\noindent \textbf{(P\_6,3.25.2)} tat vai pañcamīnirdiṣṭam ṣaṣṭhīnirdiṣṭena ca iha arthaḥ . \\
\noindent \textbf{(P\_6,3.25.2)} putre iti eṣā saptamī ṛtaḥ iti pañcamyāḥ ṣaṣṭhīm prakalpayiṣyati tasmin iti nirdiṣṭe pūrvasya iti . \\
\noindent \textbf{(P\_6,3.26)} devatādvandve ubhayatra vāyoḥ pratiṣedhaḥ . \\
\noindent \textbf{(P\_6,3.26)} devatādvandve ubhayatra vāyoḥ pratiṣedhaḥ vaktavyaḥ . \\
\noindent \textbf{(P\_6,3.26)} vāyvagnī , agnivāyū . \\
\noindent \textbf{(P\_6,3.26)} brahmaprajāpatyādīnām ca . \\
\noindent \textbf{(P\_6,3.26)} brahmaprajāpatyādīnām ca pratiṣedhaḥ vaktavyaḥ . \\
\noindent \textbf{(P\_6,3.26)} brahmaprajāpatī , śivavaiśravaṇau , skandviśākhau . \\
\noindent \textbf{(P\_6,3.26)} saḥ tarhi pratiṣedhaḥ vaktavyaḥ . \\
\noindent \textbf{(P\_6,3.26)} na vaktavyaḥ . \\
\noindent \textbf{(P\_6,3.26)} dvandve iti vartamāne punaḥ dvandragrahaṇasya etat prayojanam lokavedayoḥ yaḥ dvandvaḥ tatra yathā syāt . \\
\noindent \textbf{(P\_6,3.26)} kaḥ ca lokavedayoḥ dvandvaḥ . \\
\noindent \textbf{(P\_6,3.26)} vede ye sahanirvāpanirdiṣṭāḥ na ca ete sahanirvāpanirdiṣṭāḥ . \\
\noindent \textbf{(P\_6,3.28)} id vṛddhau viṣṇoḥ pratiṣedhaḥ . \\
\noindent \textbf{(P\_6,3.28)} id vṛddhau viṣṇoḥ pratiṣedhaḥ vaktavyaḥ . \\
\noindent \textbf{(P\_6,3.28)} āgnāvaiṣṇavam carum nirvapet . \\
\noindent \textbf{(P\_6,3.32-33)} KA\_III,149.8-10 Ro\_IV,598 {1/5}     kim nipātyate . \\
\noindent \textbf{(P\_6,3.32-33)} KA\_III,149.8-10 Ro\_IV,598 {2/5}     pūrvapadottarapadayoḥ ṛkārasya arārau nipātyete . \\
\noindent \textbf{(P\_6,3.32-33)} KA\_III,149.8-10 Ro\_IV,598 {3/5}     mātarapitarau bhojayataḥ . \\
\noindent \textbf{(P\_6,3.32-33)} KA\_III,149.8-10 Ro\_IV,598 {4/5}     mātarapitarau ānaya . \\
\noindent \textbf{(P\_6,3.32-33)} KA\_III,149.8-10 Ro\_IV,598 {5/5}     a mā gantām pitarāmātarā ca a mā somaḥ amṛtatvaya gamyāt . \\
\noindent \textbf{(P\_6,3.34.1)} bhāṣitapuṃskāt iti katham idam vijñāyate . \\
\noindent \textbf{(P\_6,3.34.1)} samānāyām ākṛtau yat bhāṣitapuṃskam āhosvit kva cit bhāṣitapuṃskam iti . \\
\noindent \textbf{(P\_6,3.34.1)} kim ca ataḥ . \\
\noindent \textbf{(P\_6,3.34.1)} yadi vijñāyate samānāyām ākṛtau yat bhāṣitapuṃskam iti garbhibhāryaḥ , prajātabhāryaḥ , prasūtabhāraḥ iti atra na prāpnoti . \\
\noindent \textbf{(P\_6,3.34.1)} atha vijñāyate kva cit bhāṣitapuṃskam iti droṇībhāryaḥ , kuṭībhāryaḥ , pātrībhāryaḥ atra api prāpnoti . \\
\noindent \textbf{(P\_6,3.34.1)} astu samānāyām ākṛtau yat bhāṣitapuṃskam iti . \\
\noindent \textbf{(P\_6,3.34.1)} katham garbhibhāryaḥ , prajātabhāryaḥ , prasūtabhāraḥ iti . \\
\noindent \textbf{(P\_6,3.34.1)} kartavyaḥ atra yatnaḥ . \\
\noindent \textbf{(P\_6,3.34.1)} atha kimartham ūṅaḥ pṛthak pratiṣedhaḥ ucyate na yatra eva anyaḥ pratiṣedhaḥ tatra eva ayam ucyeta . \\
\noindent \textbf{(P\_6,3.34.1)} na kopadhāyāḥ iti uktvā tataḥ ūṅaḥ ca iti ucyeta . \\
\noindent \textbf{(P\_6,3.34.1)} tatra api ayam arthaḥ dviḥ pratiṣedhaḥ na vaktavyaḥ bhavati . \\
\noindent \textbf{(P\_6,3.34.1)} na evam śakyam . \\
\noindent \textbf{(P\_6,3.34.1)} paṭhiṣyati hi ācāryaḥ puṃvat karmadhāraye pratiṣiddhārtham iti . \\
\noindent \textbf{(P\_6,3.34.1)} saḥ puṃvadbhāvaḥ yathā iha bhavati : kārikā vṛndārikā kārakavṛndārikā iti evam iha api syāt : brahmabandhūḥ vṛndārikā brahmabandūvṛndārikā iti . \\
\noindent \textbf{(P\_6,3.34.1)} atha pṛthak pratiṣedhe api ucyamāne yāvatā saḥ pratiṣiddhārthaḥ ārambhaḥ kasmāt eva atra na bhavati . \\
\noindent \textbf{(P\_6,3.34.1)} pṛthakpratiṣedhavacanasāmarthyāt . \\
\noindent \textbf{(P\_6,3.34.1)} atha vā anūṅ iti tatra anuvartiṣyate . \\
\noindent \textbf{(P\_6,3.34.1)} atha vā na ayam prasajyapratiṣedhaḥ . \\
\noindent \textbf{(P\_6,3.34.1)} kim tarhi paryudāsaḥ ayam yat anyat anūṅ iti . \\
\noindent \textbf{(P\_6,3.34.1)} saḥ ca pratiṣedhāṛthaḥ ārambhaḥ . \\
\noindent \textbf{(P\_6,3.34.2)} kim punaḥ idam puṃvadbhāve strīgrhaṇam strīpratyayagrahaṇam āhosvit strīśabdgrahaṇam āhosvit stryarthagrahaṇam . \\
\noindent \textbf{(P\_6,3.34.2)} kaḥ ca atra viśeṣaḥ . \\
\noindent \textbf{(P\_6,3.34.2)} puṃvadbhāve strīgrahaṇam strīpratyayagrahaṇam cet tatra puṃvat iti uttarapade tatpratiṣedhavijñānam . \\
\noindent \textbf{(P\_6,3.34.2)} puṃvadbhāve strīgrahaṇam strīpratyayagrahaṇam cet tatra puṃvat iti uttarapade tatpratiṣedhaḥ ayam vijñāyeta . \\
\noindent \textbf{(P\_6,3.34.2)} kasya . \\
\noindent \textbf{(P\_6,3.34.2)} strīpratyayasya pratiṣedhaḥ . \\
\noindent \textbf{(P\_6,3.34.2)} kim ucyate strīpratyayasya pratiṣedhaḥ iti . \\
\noindent \textbf{(P\_6,3.34.2)} na punaḥ anyat api kim cit puṃsaḥ pratipadam kāryam ucyate yat samānādhikaraṇe uttarapade bhāṣitapuṃskasya atidiśyeta . \\
\noindent \textbf{(P\_6,3.34.2)} anārambhāt puṃsi . \\
\noindent \textbf{(P\_6,3.34.2)} na hi kim cit puṃsaḥ pratipadam kāryam ucyate yat samānādhikaraṇe uttarapade bhāṣitapuṃskasya atidiśyeta . \\
\noindent \textbf{(P\_6,3.34.2)} tatra kim anyat śakyam vijñātum anyat ataḥ strīpratyayapratiṣedhāt . \\
\noindent \textbf{(P\_6,3.34.2)} katham punaḥ puṃvat iti anena strīpratyayasya pratiṣedhaḥ śakyaḥ vijñātum . \\
\noindent \textbf{(P\_6,3.34.2)} vatinirdeśaḥ ayam kāmacāraḥ ca vatinirdeśe vākyaśeṣam samarthayitum . \\
\noindent \textbf{(P\_6,3.34.2)} tat yathā . \\
\noindent \textbf{(P\_6,3.34.2)} uśīnaravat madreṣu yavāḥ . \\
\noindent \textbf{(P\_6,3.34.2)} santi na santi iti . \\
\noindent \textbf{(P\_6,3.34.2)} mātṛvat asyāḥ kalāḥ . \\
\noindent \textbf{(P\_6,3.34.2)} santi na santi . \\
\noindent \textbf{(P\_6,3.34.2)} evam iha api puṃvat bhavati puṃvat na bhavati iti vākyaśeṣam samarthayiṣyāmahe . \\
\noindent \textbf{(P\_6,3.34.2)} yathā puṃsaḥ strīpratyayaḥ na bhavati evam samānādhikaraṇe uttarapade bhāṣitapuṃskasya na bhavati iti . \\
\noindent \textbf{(P\_6,3.34.2)} prātipadikasya ca pratyāpattiḥ . \\
\noindent \textbf{(P\_6,3.34.2)} prātipadikasya ca pratyāpattiḥ vaktavyā . \\
\noindent \textbf{(P\_6,3.34.2)} enī bhāryā asya , etabhāryaḥ , śyetabhāryaḥ . \\
\noindent \textbf{(P\_6,3.34.2)} puṃvadbhāvena kim kriyate . \\
\noindent \textbf{(P\_6,3.34.2)} strīpratyayasya nivṛttiḥ . \\
\noindent \textbf{(P\_6,3.34.2)} arthaḥ anivṛttaḥ strītvam . \\
\noindent \textbf{(P\_6,3.34.2)} tasya anivṛttatvāt kena naśabdaḥ na śrūyeta . \\
\noindent \textbf{(P\_6,3.34.2)} striyām iti ucyamānaḥ prāpnoti . \\
\noindent \textbf{(P\_6,3.34.2)} sthānivatprasaṅgaḥ ca . \\
\noindent \textbf{(P\_6,3.34.2)} sthānivabhāvaḥ ca prāpnoti . \\
\noindent \textbf{(P\_6,3.34.2)} paṭvībhārā asya paṭubhāryaḥ . \\
\noindent \textbf{(P\_6,3.34.2)} puṃvadbhāvena kim kriyate . \\
\noindent \textbf{(P\_6,3.34.2)} strīpratyayasya nivṛttiḥ . \\
\noindent \textbf{(P\_6,3.34.2)} tasya sthānivabhāvāt yaṇādeśaḥ prāpnoti . \\
\noindent \textbf{(P\_6,3.34.2)} kimartham idam ubhayam ucyate na prātipadikasya ca pratyāpattiḥ iti eva sthānivabhāvaḥ api coditaḥ syāt . \\
\noindent \textbf{(P\_6,3.34.2)} purastāt idam ācāryeṇa dṛṣṭam sthānivatprasaṅgaḥ ca iti tat paṭhitam . \\
\noindent \textbf{(P\_6,3.34.2)} tataḥ uttarakālam idam dṛṣṭam prātipadikasya ca pratyāpattiḥ iti . \\
\noindent \textbf{(P\_6,3.34.2)} tat api paṭhitam . \\
\noindent \textbf{(P\_6,3.34.2)} na ca idānīm ācāryāḥ sūtrāṇi kṛtvā nivartayanti . \\
\noindent \textbf{(P\_6,3.34.2)} vataṇḍyādiṣu puṃvadvacanam . \\
\noindent \textbf{(P\_6,3.34.2)} vataṇḍyādiṣu puṃvadbhāvaḥ vaktavyaḥ . \\
\noindent \textbf{(P\_6,3.34.2)} ke punaḥ vataṇḍyādayaḥ . \\
\noindent \textbf{(P\_6,3.34.2)} lugalugastrīviṣayadvistrīpratyayāḥ . \\
\noindent \textbf{(P\_6,3.34.2)} luk . \\
\noindent \textbf{(P\_6,3.34.2)} gārgyaḥ vṛndārikā gargavṛndārikā . \\
\noindent \textbf{(P\_6,3.34.2)} puṃvadbhāvena kim kriyate . \\
\noindent \textbf{(P\_6,3.34.2)} strīpratyayasya nivṛttiḥ . \\
\noindent \textbf{(P\_6,3.34.2)} arthaḥ anivṛttaḥ strītvam . \\
\noindent \textbf{(P\_6,3.34.2)} tasya anivṛttatvāt kena yaśabdaḥ na śrūyeta . \\
\noindent \textbf{(P\_6,3.34.2)} astriyām iti hi luk ucyate . \\
\noindent \textbf{(P\_6,3.34.2)} luk . \\
\noindent \textbf{(P\_6,3.34.2)} aluk . \\
\noindent \textbf{(P\_6,3.34.2)} vataṇḍī vṛndārikā vātaṇḍyavṛndārikā . \\
\noindent \textbf{(P\_6,3.34.2)} puṃvadbhāvena kim kriyate . \\
\noindent \textbf{(P\_6,3.34.2)} strīpratyayasya nivṛttiḥ . \\
\noindent \textbf{(P\_6,3.34.2)} arthaḥ anivṛttaḥ strītvam . \\
\noindent \textbf{(P\_6,3.34.2)} tasya anivṛttatvāt luk striyām vataṇḍāt iti yakārasya luk prāpnoti . \\
\noindent \textbf{(P\_6,3.34.2)} yadi punaḥ ayam īkāre eva luk ucyeta . \\
\noindent \textbf{(P\_6,3.34.2)} tat īkāragrahaṇam kartavyam . \\
\noindent \textbf{(P\_6,3.34.2)} na kartavyam . \\
\noindent \textbf{(P\_6,3.34.2)} kriyate nyāse eva . \\
\noindent \textbf{(P\_6,3.34.2)} praśliṣṭanirdeśaḥ ayam . \\
\noindent \textbf{(P\_6,3.34.2)} strī , ī strī , striyām . \\
\noindent \textbf{(P\_6,3.34.2)} īkāravidhau vai apratyayakasya pāṭhaḥ kriyate vataṇḍa iti . \\
\noindent \textbf{(P\_6,3.34.2)} śārṅgaravādau sapratyayakasya pāṭhaḥ kariṣyate . \\
\noindent \textbf{(P\_6,3.34.2)} saḥ vai sapratyayakasya pāṭhaḥ kartavyaḥ . \\
\noindent \textbf{(P\_6,3.34.2)} antaraṅgatvāt ca luk prāpnoti . \\
\noindent \textbf{(P\_6,3.34.2)} aluk . \\
\noindent \textbf{(P\_6,3.34.2)} astrīviṣaya . \\
\noindent \textbf{(P\_6,3.34.2)} kauṇḍīvṛsī vṛndārikā kauṇḍīvṛsyavṛndārikā . \\
\noindent \textbf{(P\_6,3.34.2)} puṃvadbhāvena kim kriyate . \\
\noindent \textbf{(P\_6,3.34.2)} strīpratyayasya nivṛttiḥ . \\
\noindent \textbf{(P\_6,3.34.2)} arthaḥ anivṛttaḥ strītvam . \\
\noindent \textbf{(P\_6,3.34.2)} tasya anivṛttatvāt kena yaśabdaḥ śrūyeta . \\
\noindent \textbf{(P\_6,3.34.2)} astriyām iti hi ñyaḥ vidhīyate . \\
\noindent \textbf{(P\_6,3.34.2)} astrīviṣaya . \\
\noindent \textbf{(P\_6,3.34.2)} dvistrīpratyaya . \\
\noindent \textbf{(P\_6,3.34.2)} gārgyāyaṇī vṛndārikā gārgyavṛndārikā . \\
\noindent \textbf{(P\_6,3.34.2)} atra puṃvadbhāvaḥ na prāpnoti . \\
\noindent \textbf{(P\_6,3.34.2)} kim kāraṇam . \\
\noindent \textbf{(P\_6,3.34.2)} bhāṣitapuṃskāt anūṅaḥ samānādhikaraṇe uttarapade puṃvadbhāvaḥ bhavati iti ucyate . \\
\noindent \textbf{(P\_6,3.34.2)} yaḥ ca atra bhāṣitapuṃskāt anūṅ na asau uttarapade yaḥ ca uttarapade na asau bhāṣitapuṃskāt anūṅ iti . \\
\noindent \textbf{(P\_6,3.34.2)} astu tarhi strīśabdagrahaṇam . \\
\noindent \textbf{(P\_6,3.34.2)} strīśabdasya puṃśabdātideśaḥ iti cet sarvaprasaṅgaḥ aviśeṣāt . \\
\noindent \textbf{(P\_6,3.34.2)} strīśabdasya puṃśabdātideśaḥ iti cet sarvasya strīśabdasya puṃśabdātideśaḥ prāpnoti . \\
\noindent \textbf{(P\_6,3.34.2)} asya api prāpnoti , aṅgārakāḥ nāma śakunayaḥ . \\
\noindent \textbf{(P\_6,3.34.2)} teṣām kālikāḥ striyaḥ . \\
\noindent \textbf{(P\_6,3.34.2)} kālikāvṛndārikāḥ . \\
\noindent \textbf{(P\_6,3.34.2)} aṅgārakavṛndārikāḥ prāpnuvanti . \\
\noindent \textbf{(P\_6,3.34.2)} kṣemavṛddhayaḥ kṣatriyāḥ . \\
\noindent \textbf{(P\_6,3.34.2)} teṣām tanukeśyaḥ striyaḥ . \\
\noindent \textbf{(P\_6,3.34.2)} tanukeśīvṛndārikāḥ . \\
\noindent \textbf{(P\_6,3.34.2)} kṣemavṛddhivṛndārikāḥ prāpnuvanti . \\
\noindent \textbf{(P\_6,3.34.2)} haṃsasya varaṭā . \\
\noindent \textbf{(P\_6,3.34.2)} kacchapasya ḍulī . \\
\noindent \textbf{(P\_6,3.34.2)} ṛśyasya rohit . \\
\noindent \textbf{(P\_6,3.34.2)} aśvasya vaḍavā . \\
\noindent \textbf{(P\_6,3.34.2)} puruṣasya yoṣit . \\
\noindent \textbf{(P\_6,3.34.2)} kim kāraṇam . \\
\noindent \textbf{(P\_6,3.34.2)} aviśeṣāt . \\
\noindent \textbf{(P\_6,3.34.2)} na hi kaḥ cit viśeṣaḥ upādīyate evañjātīyakasya strīśabdasya puṃśabdātideśaḥ bhavati iti . \\
\noindent \textbf{(P\_6,3.34.2)} anupādīyamāne viśeṣe sarvatra prasaṅgaḥ . \\
\noindent \textbf{(P\_6,3.34.2)} katham ca nāma na upādīyate yāvatā bhāṣitapuṃskāt iti ucyate . \\
\noindent \textbf{(P\_6,3.34.2)} bhāṣitapuṃskānupapattiḥ ca . \\
\noindent \textbf{(P\_6,3.34.2)} hyarthe ca ayam caḥ paṭhitaḥ . \\
\noindent \textbf{(P\_6,3.34.2)} sarvaḥ hi śabdaḥ bhāṣitapuṃskāt paraḥ śakyaḥ kartum . \\
\noindent \textbf{(P\_6,3.34.2)} astu tarhi arthagrahaṇam . \\
\noindent \textbf{(P\_6,3.34.2)} arthātideśe vipratiṣedhānupapattiḥ . \\
\noindent \textbf{(P\_6,3.34.2)} arthātideśe vipratiṣedhaḥ na upapadyate . \\
\noindent \textbf{(P\_6,3.34.2)} paṭhiṣyati hi ācāryaḥ vipratiṣedham puṃvadbhāvāt hrasvatvam khidghādikeṣu iti . \\
\noindent \textbf{(P\_6,3.34.2)} saḥ vipratiṣedhaḥ na upapadyate . \\
\noindent \textbf{(P\_6,3.34.2)} dvikāryayogaḥ hi nāma vipratiṣedhaḥ na ca atra ekaḥ dvikāryayuktaḥ . \\
\noindent \textbf{(P\_6,3.34.2)} śabdasya hrasvatvam arthasya puṃvadbhāvaḥ . \\
\noindent \textbf{(P\_6,3.34.2)} kim ca sarvaprasaṅgaḥ aviśeṣāt iti . \\
\noindent \textbf{(P\_6,3.34.2)} sarvasya strīśabdasya puṃśabdātideśaḥ prāpnoti . \\
\noindent \textbf{(P\_6,3.34.2)} asya api prāpnoti , aṅgārakāḥ nāma śakunayaḥ . \\
\noindent \textbf{(P\_6,3.34.2)} teṣām kālikāḥ striyaḥ . \\
\noindent \textbf{(P\_6,3.34.2)} kālikāvṛndārikāḥ . \\
\noindent \textbf{(P\_6,3.34.2)} aṅgārakavṛndārikāḥ prāpnuvanti . \\
\noindent \textbf{(P\_6,3.34.2)} kṣemavṛddhayaḥ kṣatriyāḥ . \\
\noindent \textbf{(P\_6,3.34.2)} teṣām tanukeśyaḥ striyaḥ . \\
\noindent \textbf{(P\_6,3.34.2)} tanukeśīvṛndārikāḥ . \\
\noindent \textbf{(P\_6,3.34.2)} kṣemavṛddhivṛndārikāḥ prāpnuvanti . \\
\noindent \textbf{(P\_6,3.34.2)} haṃsasya varaṭā . \\
\noindent \textbf{(P\_6,3.34.2)} kacchapasya ḍulī . \\
\noindent \textbf{(P\_6,3.34.2)} ṛśyasya rohit . \\
\noindent \textbf{(P\_6,3.34.2)} aśvasya vaḍavā . \\
\noindent \textbf{(P\_6,3.34.2)} puruṣasya yoṣit . \\
\noindent \textbf{(P\_6,3.34.2)} kim kāraṇam . \\
\noindent \textbf{(P\_6,3.34.2)} aviśeṣāt . \\
\noindent \textbf{(P\_6,3.34.2)} na hi kaḥ cit viśeṣaḥ upādīyate evañjātīyakasya strīśabdasya puṃśabdātideśaḥ bhavati iti . \\
\noindent \textbf{(P\_6,3.34.2)} katham ca nāma na upādīyate yāvatā bhāṣitapuṃskāt iti ucyate . \\
\noindent \textbf{(P\_6,3.34.2)} bhāṣitapuṃskānupapattiḥ hi bhavati . \\
\noindent \textbf{(P\_6,3.34.2)} na hi arthen paurvāparyam asti . \\
\noindent \textbf{(P\_6,3.34.2)} ayam tāvat adoṣaḥ yat ucyate arthātideśe vipratiṣedhānupapattiḥ iti . \\
\noindent \textbf{(P\_6,3.34.2)} na avaśyam dvikārayogaḥ eva vipratiṣedhaḥ . \\
\noindent \textbf{(P\_6,3.34.2)} kim tarhi . \\
\noindent \textbf{(P\_6,3.34.2)} asambhavaḥ api . \\
\noindent \textbf{(P\_6,3.34.2)} saḥ ca atra asti asambhavaḥ . \\
\noindent \textbf{(P\_6,3.34.2)} kaḥ asambhavaḥ . \\
\noindent \textbf{(P\_6,3.34.2)} puṃvadbhāvaḥ abhinirvartamānaḥ hrasvatvasya nimittam vihanti . \\
\noindent \textbf{(P\_6,3.34.2)} hrasvatvam abhinirvartamānam puṃvadbhāvam bādhate . \\
\noindent \textbf{(P\_6,3.34.2)} sati asambhave yuktaḥ vipratiṣedhaḥ . \\
\noindent \textbf{(P\_6,3.34.2)} ayam tarhi doṣaḥ sarvaprasaṅgaḥ aviśeṣāt iti . \\
\noindent \textbf{(P\_6,3.34.2)} tasmāt astu saḥ eva madhyamaḥ pakṣaḥ . \\
\noindent \textbf{(P\_6,3.34.2)} nanu ca uktam strīśabdasya puṃśabdātideśaḥ iti cet sarvaprasaṅgaḥ aviśeṣāt iti . \\
\noindent \textbf{(P\_6,3.34.2)} na eṣaḥ doṣaḥ . \\
\noindent \textbf{(P\_6,3.34.2)} samāsanirdeśaḥ ayam : bhāṣitapuṃskāt anūṅ yasmin saḥ ayam bhāṣitapuṃskādanūṅ iti . \\
\noindent \textbf{(P\_6,3.34.2)} yadi evam luk prāpnoti . \\
\noindent \textbf{(P\_6,3.34.2)} nipātanāt na bhaviṣyati . \\
\noindent \textbf{(P\_6,3.34.2)} atha vā aluk prakṛtaḥ saḥ anuvartiṣyate . \\
\noindent \textbf{(P\_6,3.34.2)} katham punaḥ anūṅ iti anyena strīpratyayagrahaṇam śakyam vijñātum . \\
\noindent \textbf{(P\_6,3.34.2)} nañivayuktam anyasadṛśādhikaraṇe tathā hi arthagatiḥ . \\
\noindent \textbf{(P\_6,3.34.2)} nañyuktam ivayuktam ca anyasmin tatsadṛśe kāryam vijñāyate . \\
\noindent \textbf{(P\_6,3.34.2)} tathā hi arthaḥ gamyate . \\
\noindent \textbf{(P\_6,3.34.2)} tat yathā abrāhmaṇam ānaya iti ukte brāhmaṇasadṛśaḥ ānīyate . \\
\noindent \textbf{(P\_6,3.34.2)} na asau loṣṭam ānīya kṛtī bhavati . \\
\noindent \textbf{(P\_6,3.34.2)} evam iha api anūṅ iti ūṅpratiṣedhāt anyasmin ūṅsadṛśe kāryam vijñāyate . \\
\noindent \textbf{(P\_6,3.34.2)} kim ca anyat anūṅ ūṅsadṛśam . \\
\noindent \textbf{(P\_6,3.34.2)} strīpratyayaḥ . \\
\noindent \textbf{(P\_6,3.34.2)} evam api iḍabiḍ vṛndārikā , aiḍabiḍvṛndārikā , pṛth vṛndārikā , pārthavṛndārikā , darat vṛndārikā , dāradavṛndārikā , uśik vṛndārikā , auśijavṛndārikā , atra puṃvadbhāvaḥ na prāpnoti . \\
\noindent \textbf{(P\_6,3.34.2)} kartavyaḥ atra yatnaḥ . \\
\noindent \textbf{(P\_6,3.34.3)} atha iha katham bhavitavyam . \\
\noindent \textbf{(P\_6,3.34.3)} paṭvīmṛdvyau bhārye asya paṭvīmṛdubhāryaḥ , āhosvit paṭumṛdubhāryaḥ iti . \\
\noindent \textbf{(P\_6,3.34.3)} paṭvīmṛdubhāryaḥ iti bhavitavyam . \\
\noindent \textbf{(P\_6,3.34.3)} puṃvadbhāvaḥ kasmāt na bhavati . \\
\noindent \textbf{(P\_6,3.34.3)} bhāṣitapuṃskāt iti ucyate . \\
\noindent \textbf{(P\_6,3.34.3)} nanu ca bhoḥ paṭuśabdaḥ mṛduśabdaḥ puṃsi bhāṣyete . \\
\noindent \textbf{(P\_6,3.34.3)} samānāyām ākṛtau yat bhāṣitapuṃskam ākṛtyantare ca etau bhāṣitapuṃskau . \\
\noindent \textbf{(P\_6,3.34.3)} samānāyām ākṛtau api etau bhāṣitapuṃskau . \\
\noindent \textbf{(P\_6,3.34.3)} katham . \\
\noindent \textbf{(P\_6,3.34.3)} ārabhyate matublopaḥ . \\
\noindent \textbf{(P\_6,3.34.3)} evam tarhi bhāṣitapuṃskāt anūṅ samānādhikaraṇe uttarapade kṛtaḥ tasya puṃvadbhāvaḥ yasya ca akṛtaḥ na asau bhāṣitapuṃskāt anūṅ samānādhikaraṇe uttarapade . \\
\noindent \textbf{(P\_6,3.34.4)} pūraṇyām pradhānapūraṇīgrahaṇam . \\
\noindent \textbf{(P\_6,3.34.4)} pūraṇyām pradhānapūraṇīgrahaṇam kartavyam . \\
\noindent \textbf{(P\_6,3.34.4)} iha mā bhūt . \\
\noindent \textbf{(P\_6,3.34.4)} kalyāṇī pañcamī asya pakṣasya kalyāṇapañcamīkaḥ pakṣaḥ iti . \\
\noindent \textbf{(P\_6,3.34.4)} atha iha katham bhavitavyam kalyāṇī pañcamī āsām rātrīṇām iti . \\
\noindent \textbf{(P\_6,3.34.4)} kalyāṇīpañcamāḥ rātrayaḥ iti bhavitavyam . \\
\noindent \textbf{(P\_6,3.34.4)} rātrayaḥ atra pradhānam . \\
\noindent \textbf{(P\_6,3.35)} iha ke cit tasilādayaḥ ā kṛtvasucaḥ paṭhyante yeṣu puṃvadbhāvaḥ na iṣyate ke cit ca anyatra paṭhyante yeṣu puṃvadbhāvaḥ iṣyate . \\
\noindent \textbf{(P\_6,3.35)} tatra kim nyāyyam . \\
\noindent \textbf{(P\_6,3.35)} parigaṇanam kartavyam . tasilādī tratasau . \\
\noindent \textbf{(P\_6,3.35)} tratasau tasilādī draṣṭavyau . \\
\noindent \textbf{(P\_6,3.35)} tasyām śālāyām vasati . \\
\noindent \textbf{(P\_6,3.35)} tatra vasati . \\
\noindent \textbf{(P\_6,3.35)} tasyāḥ , tataḥ , yasyām , yatra, yasyāḥ , yataḥ . \\
\noindent \textbf{(P\_6,3.35)} taratamapau . \\
\noindent \textbf{(P\_6,3.35)} taratamapau tasilādī draṣṭavyau . \\
\noindent \textbf{(P\_6,3.35)} darśanīyatarā darśanīyatamā . \\
\noindent \textbf{(P\_6,3.35)} caraḍjātīyarau . \\
\noindent \textbf{(P\_6,3.35)} caraḍjātīyarau tasilādī draṣṭavyau . \\
\noindent \textbf{(P\_6,3.35)} paṭucarī , paṭujātīyā . \\
\noindent \textbf{(P\_6,3.35)} kalpabdeśīyarau . \\
\noindent \textbf{(P\_6,3.35)} kalpabdeśīyarau tasilādī draṣṭavyau . \\
\noindent \textbf{(P\_6,3.35)} darśanīyakalpā , darśanīyadeśīyā . \\
\noindent \textbf{(P\_6,3.35)} rūpappāśapau . \\
\noindent \textbf{(P\_6,3.35)} rūpappāśapau tasilādī draṣṭavyau . \\
\noindent \textbf{(P\_6,3.35)} darśanīyarūpā , darśanīyapāśā . \\
\noindent \textbf{(P\_6,3.35)} thamthālau . \\
\noindent \textbf{(P\_6,3.35)} thamthālau tasilādī draṣṭavyau . \\
\noindent \textbf{(P\_6,3.35)} kayā ākṛtyā katham , yayā yathā . \\
\noindent \textbf{(P\_6,3.35)} dārhilau . \\
\noindent \textbf{(P\_6,3.35)} dārhilau tasilādī draṣṭavyau . \\
\noindent \textbf{(P\_6,3.35)} tasyām velāyām , tadā , tarhi . tilthyanau . \\
\noindent \textbf{(P\_6,3.35)} tilthyanau tasilādī draṣṭavyau . \\
\noindent \textbf{(P\_6,3.35)} vṛkī vṛkatiḥ , ajathyā yūthiḥ . \\
\noindent \textbf{(P\_6,3.35)} śasi bahvalpārthasya . \\
\noindent \textbf{(P\_6,3.35)} śasi bahvalpārthasya puṃvadbhāvaḥ vaktavyaḥ . \\
\noindent \textbf{(P\_6,3.35)} bahvībhyaḥ dehi . \\
\noindent \textbf{(P\_6,3.35)} bahuśaḥ dehi . \\
\noindent \textbf{(P\_6,3.35)} alpaśaḥ . \\
\noindent \textbf{(P\_6,3.35)} tvataloḥ guṇavacanasya . \\
\noindent \textbf{(P\_6,3.35)} tvataloḥ guṇavacanasya puṃvadbhāvaḥ vaktavyaḥ . \\
\noindent \textbf{(P\_6,3.35)} paṭvyāḥ bhāvaḥ paṭutvam , paṭutā . \\
\noindent \textbf{(P\_6,3.35)} guṇavacanasya iti kimartham . \\
\noindent \textbf{(P\_6,3.35)} kaṭhyāḥ bhāvaḥ kaṭhītvam , kaṭhītā . \\
\noindent \textbf{(P\_6,3.35)} bhasya aḍhe taddhite . \\
\noindent \textbf{(P\_6,3.35)} bhasya aḍhe taddhite puṃvadbhāvaḥ vaktavyaḥ . \\
\noindent \textbf{(P\_6,3.35)} hastinīnām samūhaḥ hāstikam . \\
\noindent \textbf{(P\_6,3.35)} aḍhe iti kimartham . \\
\noindent \textbf{(P\_6,3.35)} śyaineyaḥ , rauhiṇeyaḥ . \\
\noindent \textbf{(P\_6,3.35)} yadi aḍhe iti ucyate , agnāyī devatā asya , āgneyaḥ sthālīpākaḥ , atra na prāpnoti . \\
\noindent \textbf{(P\_6,3.35)} iha ca prāpnoti , kauṇḍinyaḥ , sāpatnaḥ iti . \\
\noindent \textbf{(P\_6,3.35)} yadi punaḥ anapatye iti ucyeta . \\
\noindent \textbf{(P\_6,3.35)} na evam śakyam . \\
\noindent \textbf{(P\_6,3.35)} iha hi na syāt . \\
\noindent \textbf{(P\_6,3.35)} gārgyāyaṇyāḥ apatyam māṇavakaḥ gārgaḥ jālmaḥ . \\
\noindent \textbf{(P\_6,3.35)} astu tarhi aḍhe iti eva . \\
\noindent \textbf{(P\_6,3.35)} katham kauṇḍinyaḥ , sāpatnaḥ iti . \\
\noindent \textbf{(P\_6,3.35)} kauṇḍinye nipātanāt siddham . \\
\noindent \textbf{(P\_6,3.35)} kim nipātanam . \\
\noindent \textbf{(P\_6,3.35)} āgastyakauṇḍinyayoḥ iti . \\
\noindent \textbf{(P\_6,3.35)} sāpatnaśabdaḥ prakṛtayataram . \\
\noindent \textbf{(P\_6,3.35)} [R 613: sāpatnaḥ prakṛtyantatatvāt . \\
\noindent \textbf{(P\_6,3.35)} sāpatnaśabdaḥ prakṛtayataram asti .] katham agnāyī devatā asya sthālīpākasya , āgneyaḥ sthālīpākaḥ iti . \\
\noindent \textbf{(P\_6,3.35)} astu tarhi anapatye iti . \\
\noindent \textbf{(P\_6,3.35)} katham gārgaḥ jālmaḥ . \\
\noindent \textbf{(P\_6,3.35)} gārgāgneyau na saṃvadete . \\
\noindent \textbf{(P\_6,3.35)} kartavyaḥ atra yatnaḥ . \\
\noindent \textbf{(P\_6,3.35)} ṭhakchasoḥ ca . \\
\noindent \textbf{(P\_6,3.35)} ṭhakchasoḥ ca puṃvadbhāvaḥ vaktavyaḥ . \\
\noindent \textbf{(P\_6,3.35)} bhavatyāḥ chātrāḥ , bhāvatkāḥ , bhavadīyāḥ . \\
\noindent \textbf{(P\_6,3.35)} ṭhaggrahaṇam kimartham na ike kṛte ajādau iti . \\
\noindent \textbf{(P\_6,3.35)} na evam śakyam . \\
\noindent \textbf{(P\_6,3.35)} ajādilakṣaṇe hi māthitikādivat prasaṅgaḥ . \\
\noindent \textbf{(P\_6,3.35)} ajādilakṣaṇe hi māthikādivat prasajyeta . \\
\noindent \textbf{(P\_6,3.35)} tat yatha mathitam paṇyam asya māthitikaḥ iti akāralope kṛte tāntāt iti kādeśaḥ na bhavati evam iha api na syāt . \\
\noindent \textbf{(P\_6,3.36)} maningrahaṇam kimartham . \\
\noindent \textbf{(P\_6,3.36)} māningrahaṇam astryartham asamānādhikaraṇārtham ca . \\
\noindent \textbf{(P\_6,3.36)} māningrahaṇam kriyate astryartham asamānādhikaraṇārtham ca . \\
\noindent \textbf{(P\_6,3.36)} astryartham tāvat . \\
\noindent \textbf{(P\_6,3.36)} darśanīyām manyate devadattaḥ yajñadattām darśanīyamānī ayam asyāḥ . \\
\noindent \textbf{(P\_6,3.36)} asamānādhikaraṇārtham . \\
\noindent \textbf{(P\_6,3.36)} darśanīyām manyate devadattā yajñadattām darśanīyamāninī iyam asyāḥ . \\
\noindent \textbf{(P\_6,3.37)} kim idam evamādi anukramaṇam ādyasya yogasya viṣaye āhosvit puṃvadbhāvamātrasya . \\
\noindent \textbf{(P\_6,3.37)} kim ca ataḥ . \\
\noindent \textbf{(P\_6,3.37)} yadi ādyasya yogasya viṣaye mādhyamkīyaḥ , śālūkikīyaḥ , atra na prāpnoti . \\
\noindent \textbf{(P\_6,3.37)} vidhīḥ api atra na sidhyati . \\
\noindent \textbf{(P\_6,3.37)} kim kāraṇam . \\
\noindent \textbf{(P\_6,3.37)} bhāṣitapūṃśkāt anūṅ iti ucyate . \\
\noindent \textbf{(P\_6,3.37)} na hi etat bhavati bhāṣitapūṃśkāt anūṅ . \\
\noindent \textbf{(P\_6,3.37)} idam tarhi vilepikāyāḥ dharmyam vailepikam . \\
\noindent \textbf{(P\_6,3.37)} vidhiḥ ca siddhaḥ bhavati pratiṣedhaḥ ca na prāpnoti . \\
\noindent \textbf{(P\_6,3.37)} atha puṃvabhāvamātrasya viṣaye hastinīnām samūhaḥ hāstikam , jātilakṣaṇaḥ puṃvabhāvapratiṣedhaḥ prāpnoti . \\
\noindent \textbf{(P\_6,3.37)} evam tarhi na kopadhāyāḥ iti eṣaḥ yogaḥ puṃvabhāvamātrasya uttaram evamādi anukramaṇam ādyasya yogasya viṣaye . \\
\noindent \textbf{(P\_6,3.37)} na kopadhapratiṣedhe taddhitavugrahaṇam . \\
\noindent \textbf{(P\_6,3.37)} na kopadhapratiṣedhe taddhitavugrahaṇam kartavyam . \\
\noindent \textbf{(P\_6,3.37)} taddhitasya yaḥ kakāraḥ voḥ ca yaḥ kakāraḥ tasya grahaṇam kartavyam . \\
\noindent \textbf{(P\_6,3.37)} iha mā bhūt . \\
\noindent \textbf{(P\_6,3.37)} pākabhāryaḥ , bhekabhāryaḥ . \\
\noindent \textbf{(P\_6,3.40)} svāṅgāt ca ītaḥ amānini . \\
\noindent \textbf{(P\_6,3.40)} svāṅgāt ca ītaḥ amānini iti vaktavyam iha api yathā syāt . \\
\noindent \textbf{(P\_6,3.40)} dīrghamukhamānī , ślakṣṇamukhamāninī . \\
\noindent \textbf{(P\_6,3.40)} yadi amānini iti ucyate dīrghamukhamāninī , ślakṣṇamukhamānininī iti na sidhyati . \\
\noindent \textbf{(P\_6,3.40)} prātipadikagrahaṇe liṅgaviśiṣṭasya api grahaṇam bhavati iti evam bhaviṣyati . \\
\noindent \textbf{(P\_6,3.42.1)} kimartham idam ucyate . \\
\noindent \textbf{(P\_6,3.42.1)} puṃvat karmadhāraye pratiṣiddhārtham . \\
\noindent \textbf{(P\_6,3.42.1)} pratiṣiddhāṛthaḥ ayam ārambhaḥ . \\
\noindent \textbf{(P\_6,3.42.1)} na kopadhāyāḥ iti uktam . \\
\noindent \textbf{(P\_6,3.42.1)} tatra api puṃvat bhavati . \\
\noindent \textbf{(P\_6,3.42.1)} kārikā vṛndārikā kārakavṛndārikā kārakajātīyā kārakadeśīyā . \\
\noindent \textbf{(P\_6,3.42.1)} sañjñāpūraṇayoḥ ca iti uktam . \\
\noindent \textbf{(P\_6,3.42.1)} tatra api puṃvat bhavati . \\
\noindent \textbf{(P\_6,3.42.1)} dattā vṛndārikā dattavṛndārikā dattajātīyā dattadeśīyā . \\
\noindent \textbf{(P\_6,3.42.1)} pañcamī vṛndārikā pañcamavṛndārikā pañcamajātīyā pañcamadeśīyā . \\
\noindent \textbf{(P\_6,3.42.1)} vṛddhinimittasya iti uktam . \\
\noindent \textbf{(P\_6,3.42.1)} tatra api puṃvat bhavati . \\
\noindent \textbf{(P\_6,3.42.1)} sraughnī vṛndārikā sraughnavṛndārikā sraughnajātīyā sraughnadeśīyā . \\
\noindent \textbf{(P\_6,3.42.1)} svāṅgāt ca ītaḥ amānini iti uktam . \\
\noindent \textbf{(P\_6,3.42.1)} tatra api puṃvat bhavati . \\
\noindent \textbf{(P\_6,3.42.1)} ślakṣṇamukhī vṛndārikā ślakṣṇamukhavṛndārikā ślakṣṇamukhajātīyā ślakṣṇamukhadeśīyā . \\
\noindent \textbf{(P\_6,3.42.1)} jāteḥ ca iti uktam . \\
\noindent \textbf{(P\_6,3.42.1)} tatra api puṃvat bhavati . \\
\noindent \textbf{(P\_6,3.42.1)} kaṭhī vṛndārikā kaṭhavṛndārikā kaṭhajātīyā kaṭhadeśīyā . \\
\noindent \textbf{(P\_6,3.42.2)} kukkuṭyādīnām aṇḍādiṣu puṃvadvacanam . \\
\noindent \textbf{(P\_6,3.42.2)} kukkuṭyādīnām aṇḍādiṣu puṃvadbhāvaḥ vaktavyaḥ . \\
\noindent \textbf{(P\_6,3.42.2)} kukkuṭyāḥ aṇḍam kukkuṭāṇḍam , mṛgyāḥ padam mṛgapadam , kākyāḥ śāvaḥ kākaśāvaḥ . \\
\noindent \textbf{(P\_6,3.42.2)} na vā astrīpūrvapadavivakṣitatvāt . \\
\noindent \textbf{(P\_6,3.42.2)} na vā vaktavyam . \\
\noindent \textbf{(P\_6,3.42.2)} kim kāraṇam astrīpūrvapadavivakṣitatvāt . \\
\noindent \textbf{(P\_6,3.42.2)} na atra strīpūrvapadam vivakṣitam . \\
\noindent \textbf{(P\_6,3.42.2)} kim tarhi . \\
\noindent \textbf{(P\_6,3.42.2)} astrīpūrvapadam . \\
\noindent \textbf{(P\_6,3.42.2)} ubhayoḥ aṇḍam ubhayoḥ padam ubhayoḥ śāvaḥ . \\
\noindent \textbf{(P\_6,3.42.2)} yadi api tāvat atra etat śakyate vaktum iha tu katham . \\
\noindent \textbf{(P\_6,3.42.2)} mṛgyāḥ kṣīram mṛghakṣīram iti . \\
\noindent \textbf{(P\_6,3.42.2)} atra api na vā astrīpūrvapadavivakṣitatvāt iti eva . \\
\noindent \textbf{(P\_6,3.42.2)} katham punaḥ sataḥ nāma avāvivakṣā syāt . \\
\noindent \textbf{(P\_6,3.42.2)} sataḥ api avivakṣā bhavati . \\
\noindent \textbf{(P\_6,3.42.2)} tat yathā . \\
\noindent \textbf{(P\_6,3.42.2)} alomikā eḍakā . \\
\noindent \textbf{(P\_6,3.42.2)} anudarā kanyā iti . \\
\noindent \textbf{(P\_6,3.42.2)} asataḥ ca vivakṣā bhavati . \\
\noindent \textbf{(P\_6,3.42.2)} samudraḥ kuṇḍikā . \\
\noindent \textbf{(P\_6,3.42.2)} vindhyaḥ vardhitakam iti . \\
\noindent \textbf{(P\_6,3.42.3)} agneḥ īttvāt varuṇasya vṛddhiḥ vipratiṣedhena . \\
\noindent \textbf{(P\_6,3.42.3)} agneḥ īttvāt varuṇasya vṛddhiḥ bhavati vipratiṣedhena . \\
\noindent \textbf{(P\_6,3.42.3)} agneḥ īttvasya avakāśaḥ , agnīṣomau . \\
\noindent \textbf{(P\_6,3.42.3)} varuṇasya vṛddheḥ avakāśaḥ , vāyuvāruṇam . \\
\noindent \textbf{(P\_6,3.42.3)} iha ubhayam prāpnoti , āgnivāruṇīm anaḍvāhīm ālabheta . \\
\noindent \textbf{(P\_6,3.42.3)} varuṇasya vṛddhiḥ bhavati vipratiṣedhena . \\
\noindent \textbf{(P\_6,3.42.3)} na eṣaḥ yuktaḥ vipratiṣedhaḥ . \\
\noindent \textbf{(P\_6,3.42.3)} dvikāryayogaḥ hi vipratiṣedhaḥ na ca atra ekaḥ dvikāryayuktaḥ . \\
\noindent \textbf{(P\_6,3.42.3)} katham . \\
\noindent \textbf{(P\_6,3.42.3)} agneḥ īttvam varuṇasya vṛddhiḥ . \\
\noindent \textbf{(P\_6,3.42.3)} na avaśyam dvikārayogaḥ eva vipratiṣedhaḥ . \\
\noindent \textbf{(P\_6,3.42.3)} kim tarhi . \\
\noindent \textbf{(P\_6,3.42.3)} asambhavaḥ api . \\
\noindent \textbf{(P\_6,3.42.3)} saḥ ca atra asti asambhavaḥ . \\
\noindent \textbf{(P\_6,3.42.3)} kaḥ asau asambhavaḥ . \\
\noindent \textbf{(P\_6,3.42.3)} agneḥ īttvam abhinirvartamanam varuṇasya vṛddhim bādhate . \\
\noindent \textbf{(P\_6,3.42.3)} varuṇasya vṛddhiḥ abhinirvartamanā agneḥ īttvam bādhate . \\
\noindent \textbf{(P\_6,3.42.3)} eṣaḥ asambhavaḥ . \\
\noindent \textbf{(P\_6,3.42.3)} sati asambhave yuktaḥ vipratiṣedhaḥ . \\
\noindent \textbf{(P\_6,3.42.3)} pūṃvadbhāvāt hrasvatvam khidghādiṣu . \\
\noindent \textbf{(P\_6,3.42.3)} pūṃvadbhāvāt hrasvatvam bhavati vipratiṣedhena khidghādiṣu . \\
\noindent \textbf{(P\_6,3.42.3)} pūṃvadbhāvasya avakāśaḥ , paṭubhāryaḥ , mṛdubhāryaḥ . \\
\noindent \textbf{(P\_6,3.42.3)} khiti hrasvaḥ bhavati iti asya avakāśaḥ , kālimmanyaḥ , hariṇimmanyaḥ . \\
\noindent \textbf{(P\_6,3.42.3)} iha ubhayam prāpnoti , kālimmanyā , hariṇimmanyā . \\
\noindent \textbf{(P\_6,3.42.3)} ghādiṣu nadyāḥ hrasvaḥ bhavati iti asya avakāśaḥ , nartakitarā , nartakitamā . \\
\noindent \textbf{(P\_6,3.42.3)} pūṃvadbhāvasya avakāśaḥ , darśanīyatarā , darśanīyatamā . \\
\noindent \textbf{(P\_6,3.42.3)} iha ubhayam prāpnoti , paṭvitarā , paṭvitamā . \\
\noindent \textbf{(P\_6,3.42.3)} ke hrasvaḥ bhavati iti asya avakāśaḥ , nartakika . \\
\noindent \textbf{(P\_6,3.42.3)} pūṃvadbhāvasya avakāśaḥ , dāradikā . \\
\noindent \textbf{(P\_6,3.42.3)} iha ubhayam prāpnoti , paṭvikā , mṛdvikā . \\
\noindent \textbf{(P\_6,3.42.3)} hrasvatvam bhavati vipratiṣedhena . \\
\noindent \textbf{(P\_6,3.42.3)} atha idānīm hrasvatve kṛte punaḥprasaṅgavijñānāt puṃvadbhāvaḥ kasmāt na bhavati . \\
\noindent \textbf{(P\_6,3.42.3)} sakṛt gatau vipratiṣedhe yat bādhitam tat bādhitam eva iti . \\
\noindent \textbf{(P\_6,3.43)} ṅīgrahaṇam kimartham . \\
\noindent \textbf{(P\_6,3.43)} anekācaḥ hrasvaḥ iti iyati ucyamāne khaṭvātarā mālātarā , atra api prasajyeta . \\
\noindent \textbf{(P\_6,3.43)} na etat asti prayojanam . \\
\noindent \textbf{(P\_6,3.43)} bhāṣitapuṃskāt iti vartate . \\
\noindent \textbf{(P\_6,3.43)} evam api dattātarā guptātarā , atra api prāpnoti . \\
\noindent \textbf{(P\_6,3.43)} ītaḥ iti vartate . \\
\noindent \textbf{(P\_6,3.43)} evam api grāmaṇītaraḥ , senāṇītaraḥ atra api prāpnoti . \\
\noindent \textbf{(P\_6,3.43)} striyām iti vartate . \\
\noindent \textbf{(P\_6,3.43)} evam api grāmaṇītarā , senāṇītarā atra api prāpnoti . \\
\noindent \textbf{(P\_6,3.43)} striyāḥ striyām iti vartate . \\
\noindent \textbf{(P\_6,3.43)} śeṣaprakḷptyartham tarhi ṅīgrahaṇam kartavyam , nadyāḥ śeṣasya anyatarasyām iti . \\
\noindent \textbf{(P\_6,3.43)} kaḥ śeṣaḥ . \\
\noindent \textbf{(P\_6,3.43)} āṅīca yā nadī ṅyantam ca yat ekāc . \\
\noindent \textbf{(P\_6,3.43)} antareṇa api ṅīgrahaṇam kḷptaḥ śeṣaḥ . \\
\noindent \textbf{(P\_6,3.43)} katham . \\
\noindent \textbf{(P\_6,3.43)} ītaḥ iti vartate . \\
\noindent \textbf{(P\_6,3.43)} anīt ca yā nadī , īdantam ca yat ekāc . \\
\noindent \textbf{(P\_6,3.43)} śeṣagrahaṇam ca api śakyam akartum . \\
\noindent \textbf{(P\_6,3.43)} katham . \\
\noindent \textbf{(P\_6,3.43)} aviśeṣeṇa ghādiṣu nadyāḥ anyatarasyām hrasvatvam utsargaḥ . \\
\noindent \textbf{(P\_6,3.43)} tasya anekācaḥ nityam hrasvatvam apavādaḥ . \\
\noindent \textbf{(P\_6,3.43)} tasmin nitye prāpte ugitaḥ vibhāṣā ārabhyate . \\
\noindent \textbf{(P\_6,3.43)} yadi evam lakṣmitarā tantritarā iti na sidhyati . \\
\noindent \textbf{(P\_6,3.43)} lakṣmītarā tantrītarā iti prāpnoti . \\
\noindent \textbf{(P\_6,3.43)} iṣṭam eva etad saṅgṛhītam . \\
\noindent \textbf{(P\_6,3.43)} lakṣmītarā tantrītarā iti eva bhavitavyam . \\
\noindent \textbf{(P\_6,3.43)} evam hi saunāgāḥ paṭhanti . \\
\noindent \textbf{(P\_6,3.43)} ghādiṣu nadyāḥ hrasvatve kṛnnadyāḥ pratiṣedhaḥ . \\
\noindent \textbf{(P\_6,3.46.1)} iha kasmāt na bhavati , amahān mahān sampannaḥ mahadbhūtaḥ candramāḥ iti . \\
\noindent \textbf{(P\_6,3.46.1)} anyaprakṛtiḥ tu amahān mahatprakṛtau mahān mahati eva . anyaḥ mahān . \\
\noindent \textbf{(P\_6,3.46.1)} anyaḥ mahān bhūtaprakṛtau vartate . \\
\noindent \textbf{(P\_6,3.46.1)} mahān mahati eva . \\
\noindent \textbf{(P\_6,3.46.1)} tasmāt āttvam na syāt . tasmāt āttvam na bhaviṣyati . \\
\noindent \textbf{(P\_6,3.46.1)} puṃvattvam tu katham bhavet atra . \\
\noindent \textbf{(P\_6,3.46.1)} puṃvadbhāvaḥ api tarhi na prāpnoti . \\
\noindent \textbf{(P\_6,3.46.1)} amahatī mahatī sampannā mahadbhūtā brāhmaṇī . \\
\noindent \textbf{(P\_6,3.46.1)} evam tarhi amahati mahān hi vṛttaḥ tadvācī ca atra bhūtaśabdaḥ ayam . \\
\noindent \textbf{(P\_6,3.46.1)} amahati hi mahacchabdaḥ vartate tadvācī ca atra bhūtaśabdaḥ prayujyate . \\
\noindent \textbf{(P\_6,3.46.1)} kiṃvācī . \\
\noindent \textbf{(P\_6,3.46.1)} mahadvācī . \\
\noindent \textbf{(P\_6,3.46.1)} tasmāt sidhyati puṃvat . tasmāt sidhyati puṃvadbhāvaḥ . \\
\noindent \textbf{(P\_6,3.46.1)} yadi evam āttvam api prāpnoti . \\
\noindent \textbf{(P\_6,3.46.1)} mahadbhūtaḥ candramāḥ . \\
\noindent \textbf{(P\_6,3.46.1)} nivartyam āttvam tu manyante . \\
\noindent \textbf{(P\_6,3.46.1)} āttvam api prāpnoti . \\
\noindent \textbf{(P\_6,3.46.1)} na eṣaḥ doṣaḥ . \\
\noindent \textbf{(P\_6,3.46.1)} yaḥ tu mahataḥ pratipadam samāsaḥ uktaḥ tadāśrayam hi āttvam kartavyam manyante na lakṣaṇena lakṣaṇoktaḥ ca ayam . evam tarhi lakṣaṇapratipadoktayoḥ pratipadoktasya eva iti pratipadam yaḥ samāsaḥ vihitaḥ tasya grahaṇam . \\
\noindent \textbf{(P\_6,3.46.1)} lakṣaṇoktaḥ ca ayam . \\
\noindent \textbf{(P\_6,3.46.1)} iha api tarhi na prāpnoti . \\
\noindent \textbf{(P\_6,3.46.1)} mahān bāhuḥ asya mahābāhuḥ iti . \\
\noindent \textbf{(P\_6,3.46.1)} śeṣavacanāt tu yaḥ asau pratyārambhāt kṛtaḥ bahuvrīhiḥ tasmāt sidhyati tasmin . \\
\noindent \textbf{(P\_6,3.46.1)} yasmāt śeṣaḥ bahuvrīhiḥ iti siddhe anekam anyapadārthe iti āha tena pratipadam bhavati . \\
\noindent \textbf{(P\_6,3.46.1)} pradhānataḥ vā yataḥ vṛttiḥ . \\
\noindent \textbf{(P\_6,3.46.1)} atha vā gauṇamukhyayoḥ mukhye kāryasampratyayaḥ . \\
\noindent \textbf{(P\_6,3.46.1)} tat yathā gauḥ anubandhyaḥ ajaḥ agnīṣomīyaḥ iti na bāhīkaḥ anubadhyate . \\
\noindent \textbf{(P\_6,3.46.1)} katham tarhi bāhīke vṛddhyāttve bhavataḥ . \\
\noindent \textbf{(P\_6,3.46.1)} gauḥ tiṣṭhati , gām ānaya iti . \\
\noindent \textbf{(P\_6,3.46.1)} arthāśraye etat evam . \\
\noindent \textbf{(P\_6,3.46.1)} yat hi śabdāśrayam śabdamātre tat bhavati . \\
\noindent \textbf{(P\_6,3.46.1)} śabdāśraye ca vṛddhyāttve . \\
\noindent \textbf{(P\_6,3.46.2)} mahadāttve ghāsakaraviśiṣṭeṣu upasaṅkhyānam puṃvadvacanam ca asamānādhikaraṇārtham . \\
\noindent \textbf{(P\_6,3.46.2)} mahadāttve ghāsakaraviśiṣṭeṣu upasaṅkhyānam kartavyam puṃvadbhāvaḥ ca asamānādhikaraṇārthaḥ kartavyaḥ . \\
\noindent \textbf{(P\_6,3.46.2)} mahatyāḥ ghāsaḥ mahāghāsaḥ , mahatyāḥ karaḥ mahākaraḥ , mahatyāḥ viśiṣṭaḥ mahāviśiṣṭaḥ . \\
\noindent \textbf{(P\_6,3.46.2)} aṣṭanaḥ kapāle haviṣi . \\
\noindent \textbf{(P\_6,3.46.2)} aṣṭanaḥ kapāle haviṣi upasaṅkhyānam kartavyam . \\
\noindent \textbf{(P\_6,3.46.2)} aṣṭākapālam carum nirvapet . \\
\noindent \textbf{(P\_6,3.46.2)} haviṣi iti kimartham . \\
\noindent \textbf{(P\_6,3.46.2)} aṣṭakapālam brāhmaṇasya . \\
\noindent \textbf{(P\_6,3.46.2)} gavi ca yukte . \\
\noindent \textbf{(P\_6,3.46.2)} gavi ca yukte upasaṅkhyānam kartavyam . \\
\noindent \textbf{(P\_6,3.46.2)} aṣṭāgavena śakaṭena . \\
\noindent \textbf{(P\_6,3.46.2)} yukte iti kimartham . \\
\noindent \textbf{(P\_6,3.46.2)} aṣṭagavam brāhmaṇasya . \\
\noindent \textbf{(P\_6,3.47)} prāk śatāt iti vaktavyam iha mā bhūt . dviśatam , dvisahasram , aṣtaśatam , aṣṭasahasram . \\
\noindent \textbf{(P\_6,3.48)} sarveṣāṅgrahaṇam kimartham . \\
\noindent \textbf{(P\_6,3.48)} catvāriṃśatprabhṛtau sarveṣām vibhāṣā yathā syāt , dvyaṣṭanoḥ ca treḥ ca . \\
\noindent \textbf{(P\_6,3.48)} na etat asti prayojanam . \\
\noindent \textbf{(P\_6,3.48)} prakṛtam dvyaṣṭangrahaṇamanuvartiṣyate . \\
\noindent \textbf{(P\_6,3.48)} yadi tat anuvartate treḥ trayaḥ dvyaṣṭanoḥ ca iti dvyaṣṭanoḥ api dtrayaḥ ādeśaḥ prāpnoti . \\
\noindent \textbf{(P\_6,3.48)} na eṣaḥ doṣaḥ . \\
\noindent \textbf{(P\_6,3.48)} maṇḍūkagatayaḥ adhikārāḥ . \\
\noindent \textbf{(P\_6,3.48)} yathā maṇḍūkāḥ utplutya utplutya gacchanti tadvat adhikārāḥ . \\
\noindent \textbf{(P\_6,3.48)} atha vā ekayogaḥ kariṣyate . \\
\noindent \textbf{(P\_6,3.48)} dvyaṣṭanaḥ saṅkhyayām abahuvrīhyaśītyoḥ treḥ trayaḥ . \\
\noindent \textbf{(P\_6,3.48)} tataḥ vibhāṣā catvāriṃśatprabhṛtau sarveṣām iti . \\
\noindent \textbf{(P\_6,3.48)} atha vā ubhayam nivṛttam tat apekṣiṣyāmahe . \\
\noindent \textbf{(P\_6,3.50)} yaṇgrahaṇam idam pratyayagrahaṇam . \\
\noindent \textbf{(P\_6,3.50)} tatra pratyayagrahaṇe yasmāt saḥ tadādeḥ grahaṇam bhavati iti yadaṇante prāpnoti . \\
\noindent \textbf{(P\_6,3.50)} yadaṇgrahaṇe rūpagrahaṇam lekhagrahaṇāt . \\
\noindent \textbf{(P\_6,3.50)} yadaṇgrahaṇe rūpagrahaṇam draṣṭavyam . \\
\noindent \textbf{(P\_6,3.50)} kutaḥ . \\
\noindent \textbf{(P\_6,3.50)} lekhagrahaṇāt . \\
\noindent \textbf{(P\_6,3.50)} yat ayam lekhagrahaṇam karoti tat jñāpayati ācāryaḥ na yadaṇante bhavati iti . \\
\noindent \textbf{(P\_6,3.50)} aparaḥ āha : atyalpam idam ucyate . \\
\noindent \textbf{(P\_6,3.50)} sarvatra eva uttarapadādhikāre pratyayagrahaṇe rūpagrahaṇam draṣṭavyam . \\
\noindent \textbf{(P\_6,3.50)} kutaḥ . \\
\noindent \textbf{(P\_6,3.50)} lekhagrahaṇāt eva . \\
\noindent \textbf{(P\_6,3.50)} kim prayojanam . \\
\noindent \textbf{(P\_6,3.50)} kumārī gauritarā . \\
\noindent \textbf{(P\_6,3.50)} ghādiṣu nadyāḥ hrasvaḥ bhavati iti hrasvatvam prasajyeta . \\
\noindent \textbf{(P\_6,3.50)} yadi etat jñāpyate khiti anavyayasya iti khiti eva anantarasya anavyayasya hrasvatvam prāpnoti . \\
\noindent \textbf{(P\_6,3.50)} khiti anantaraḥ hrasvabhāvī na asti iti kṛtvā khidante bhaviṣyati . \\
\noindent \textbf{(P\_6,3.50)} nanu ca ayam asti stanandhayaḥ iti . \\
\noindent \textbf{(P\_6,3.50)} atra api śapā vyavadhānam . \\
\noindent \textbf{(P\_6,3.50)} ekādeśe kṛte na asti vyavadhānam . \\
\noindent \textbf{(P\_6,3.50)} ekādeśaḥ pūrvavidhau sthānivat bhavati iti sthānivadbhāvāt vyavadhānam eva . \\
\noindent \textbf{(P\_6,3.50)} atha vā etat jñāpayati ācāryaḥ khiti anantarasya na bhavati iti yat ayam anavyayasya iti pratiṣedham śāsti . \\
\noindent \textbf{(P\_6,3.50)} na hi khiti anantaram avyayam asti . \\
\noindent \textbf{(P\_6,3.52)} padādeśe antodāttanipātanam [R: padopahatārtham] . padādeśe antodāttanipātanam kartavyam . \\
\noindent \textbf{(P\_6,3.52)} kim prayojanam . \\
\noindent \textbf{(P\_6,3.52)} padopahatārtham . \\
\noindent \textbf{(P\_6,3.52)} pādena upahatam padopahatam . \\
\noindent \textbf{(P\_6,3.52)} tṛtīyā karmaṇi iti prakṛtisvaratve pūrvapadāntodāttatvam yathā syāt . \\
\noindent \textbf{(P\_6,3.52)} upadeśivadvacanam ca svarasiddhyartham . \\
\noindent \textbf{(P\_6,3.52)} upadeśivadbhāvaḥ ca vaktavyaḥ . \\
\noindent \textbf{(P\_6,3.52)} kim prayojanam . \\
\noindent \textbf{(P\_6,3.52)} svarasiddhyartham . \\
\noindent \textbf{(P\_6,3.52)} upadeśāvasthāyām antodāttanipātane kṛte samāsasvareṇa bādhanam yathā syāt . \\
\noindent \textbf{(P\_6,3.52)} padājiḥ , padātiḥ . \\
\noindent \textbf{(P\_6,3.53)} padbhāve ike caratau upasaṅkhyānam . \\
\noindent \textbf{(P\_6,3.53)} padbhāve ike caratau upasaṅkhyānam kartavyam . \\
\noindent \textbf{(P\_6,3.53)} pādābhyām carati padikaḥ . \\
\noindent \textbf{(P\_6,3.56)} niṣke ca upasaṅkhyānam kartavyam . \\
\noindent \textbf{(P\_6,3.56)} panniṣkeṇa pādaniṣkeṇa . \\
\noindent \textbf{(P\_6,3.57)} sañjñāyām uttarapadasya iti vaktavyam iha api yathā syāt . \\
\noindent \textbf{(P\_6,3.57)} lohitodaḥ , kṣīrodaḥ iti . \\
\noindent \textbf{(P\_6,3.59)} ekahalādau iti kimartham . \\
\noindent \textbf{(P\_6,3.59)} udakasthānam . \\
\noindent \textbf{(P\_6,3.59)} ucyamāne api etasmin atra prapnoti . \\
\noindent \textbf{(P\_6,3.59)} etat api ekahalādi . \\
\noindent \textbf{(P\_6,3.59)} kim kāraṇam . \\
\noindent \textbf{(P\_6,3.59)} ekaikavarṇavartitvāt vācaḥ uccaritapradhvaṃsitvāt ca varṇānām . \\
\noindent \textbf{(P\_6,3.59)} ekaikavarṇavartinī vāk . \\
\noindent \textbf{(P\_6,3.59)} na dvau varṇau yugapat uccārayati . \\
\noindent \textbf{(P\_6,3.59)} tat yathā gauḥ iti ukte yāvat gakāre vāk vartate tāvat na aukāre na visarjanīye . \\
\noindent \textbf{(P\_6,3.59)} yāvat auakāre na tāvat gakāre na visarjanīye . \\
\noindent \textbf{(P\_6,3.59)} yāvat visarjanīye na tāvat gakāre na aukāre . \\
\noindent \textbf{(P\_6,3.59)} uccaritapradhvaṃsitvāt ca varṇānām . \\
\noindent \textbf{(P\_6,3.59)} uccaritaḥ varṇaḥ pradhvastaḥ ca . \\
\noindent \textbf{(P\_6,3.59)} atha aparaḥ prayujyate . \\
\noindent \textbf{(P\_6,3.59)} na varṇaḥ varṇasya sahāyaḥ . \\
\noindent \textbf{(P\_6,3.59)} evam tarhi ekahalādau iti ucyate sarvaḥ ca ekahalādiḥ . \\
\noindent \textbf{(P\_6,3.59)} tatra prakarṣagatiḥ vijñāyate : sādhīyaḥ yaḥ ekahalādiḥ iti . \\
\noindent \textbf{(P\_6,3.59)} kaḥ ca sādhīyaḥ . \\
\noindent \textbf{(P\_6,3.59)} yatra ekam halam uccārya ac ucyate . \\
\noindent \textbf{(P\_6,3.61)} ikaḥ hrasvatvam uttarapadamātre . \\
\noindent \textbf{(P\_6,3.61)} ikaḥ hrasvatvam uttarapadamātre vaktavyam iha api yathā syāt , alābukarkandhudṛnbhuphalam iti . \\
\noindent \textbf{(P\_6,3.61)} kim punaḥ kāraṇam na sidhyati . \\
\noindent \textbf{(P\_6,3.61)} sarvānte hi lokavijñānam . \\
\noindent \textbf{(P\_6,3.61)} lokavijñānāt hi yat eva sarvāntam padam tasmin pūrvapadasya hrasvatvam syāt . \\
\noindent \textbf{(P\_6,3.61)} atha vā evam vigrahaḥ kariṣyate . \\
\noindent \textbf{(P\_6,3.61)} alābūḥ ca karkandhūḥ ca , alābukarkandhvau , alābukarkandhvau dṛnbhūḥ ca, alābukarkandhudṛnbhvaḥ , alābukarkandhudṛnbhūnām phalam alābukarkandhudṛnbhuphalam iti . \\
\noindent \textbf{(P\_6,3.61)} yadi evam dṛnbhvāḥ pūrvanipātaḥ prāpnoti . \\
\noindent \textbf{(P\_6,3.61)} rājadantādiṣu pāṭhaḥ kariṣyate . \\
\noindent \textbf{(P\_6,3.61)} atha vā evam vigrahaḥ kariṣyate . \\
\noindent \textbf{(P\_6,3.61)} dṛnbhvāḥ phalam dṛnbhuphalam , karkandhūḥ ca dṛnbhuphalam ca karkandhudṛnbhuphalam , alābūḥ ca karkandhudṛnbhuphalam ca alābukarkandhudṛnbhuphalam iti . \\
\noindent \textbf{(P\_6,3.61)} evam api phalena akṛtaḥ abhisambandhaḥ bhavati . \\
\noindent \textbf{(P\_6,3.61)} pratyekam phalaśabdaḥ parisamāpyate . \\
\noindent \textbf{(P\_6,3.61)} iyaṅuvaṅavyayapratiṣedhaḥ . \\
\noindent \textbf{(P\_6,3.61)} iyaṅuvaṅbhāinām avyayānām ca pratiṣedhaḥ vaktavyaḥ . \\
\noindent \textbf{(P\_6,3.61)} śrīkulam , bhrūkulam , kāṇḍībhūtam vṛṣalakulam , kuḍyībhūtam vṛṣalakulam . \\
\noindent \textbf{(P\_6,3.61)} abhrūkaṃsādīnām iti vaktavyam . \\
\noindent \textbf{(P\_6,3.61)} bhrukuṃsaḥ , bhrukuṭiḥ . \\
\noindent \textbf{(P\_6,3.61)} aparaḥ āha : akāraḥ bhrūkaṃsādīnām iti vaktavyam . \\
\noindent \textbf{(P\_6,3.61)} bhrakuṃsaḥ , bhrakuṭiḥ . \\
\noindent \textbf{(P\_6,3.62)} taddhite kim udāharaṇam . \\
\noindent \textbf{(P\_6,3.62)} ekatvam , ekatā . \\
\noindent \textbf{(P\_6,3.62)} na etat asti prayojanam . \\
\noindent \textbf{(P\_6,3.62)} puṃvadbhāvena api etat siddham . \\
\noindent \textbf{(P\_6,3.62)} katham puṃvadbhāvaḥ . \\
\noindent \textbf{(P\_6,3.62)} tāsilādiṣu ā kṛtvasucaḥ . \\
\noindent \textbf{(P\_6,3.62)} idam tarhi prayojanam . \\
\noindent \textbf{(P\_6,3.62)} ekasyāḥ āgatam ekarūpyam , ekamayam . \\
\noindent \textbf{(P\_6,3.62)} idam ca api udāharaṇam . \\
\noindent \textbf{(P\_6,3.62)} ekatvam , ekatā . \\
\noindent \textbf{(P\_6,3.62)} nanu ca uktam puṃvadbhāvena api etat siddham iti . \\
\noindent \textbf{(P\_6,3.62)} na sidhyati . \\
\noindent \textbf{(P\_6,3.62)} uktam etat tvataloḥ guṇavacanasya iti . \\
\noindent \textbf{(P\_6,3.62)} atha uttarapade kim udāharaṇam . \\
\noindent \textbf{(P\_6,3.62)} ekaśāṭī . \\
\noindent \textbf{(P\_6,3.62)} na etat asti . \\
\noindent \textbf{(P\_6,3.62)} puṃvadbhāvena api etat siddham . \\
\noindent \textbf{(P\_6,3.62)} katham puṃvadbhāvaḥ . \\
\noindent \textbf{(P\_6,3.62)} samānādhikaraṇalakṣaṇaḥ . \\
\noindent \textbf{(P\_6,3.62)} idam tarhi prayojanam . \\
\noindent \textbf{(P\_6,3.62)} ekasyāḥ kṣīram ekaṣīram . \\
\noindent \textbf{(P\_6,3.62)} idam ca api udāharaṇam . \\
\noindent \textbf{(P\_6,3.62)} ekaśāṭī . \\
\noindent \textbf{(P\_6,3.62)} nanu ca uktam puṃvadbhāvena api etat siddham iti . \\
\noindent \textbf{(P\_6,3.62)} na sidhyati . \\
\noindent \textbf{(P\_6,3.62)} na kopadhāyāḥ iti pratiṣedhaḥ prāpnoti . \\
\noindent \textbf{(P\_6,3.62)} na eṣaḥ doṣaḥ . \\
\noindent \textbf{(P\_6,3.62)} uktam etat kopadhapratiṣedhe taddhitavugrahaṇam iti . \\
\noindent \textbf{(P\_6,3.66)} khiti hrasvāprasiddhiḥ anajantatvāt . \\
\noindent \textbf{(P\_6,3.66)} khiti hrasvāprasiddhiḥ . \\
\noindent \textbf{(P\_6,3.66)} kālimmanyā , hariṇimmanyā . \\
\noindent \textbf{(P\_6,3.66)} kim kāraṇam . \\
\noindent \textbf{(P\_6,3.66)} anajantatvāt . \\
\noindent \textbf{(P\_6,3.66)} mumi kṛte anajantatvāt hrasvatvam na prāpnoti . \\
\noindent \textbf{(P\_6,3.66)} siddham tu hrasvāntasya mumvacanāt . \\
\noindent \textbf{(P\_6,3.66)} siddham etat . \\
\noindent \textbf{(P\_6,3.66)} katham . \\
\noindent \textbf{(P\_6,3.66)} hrasvāntasya mum bhavati iti vaktavyam . \\
\noindent \textbf{(P\_6,3.66)} sanniyogāt vā . \\
\noindent \textbf{(P\_6,3.66)} atha vā sanniyogaḥ kariṣyate . \\
\noindent \textbf{(P\_6,3.66)} kaḥ eṣaḥ yatnaḥ codyate sanniyogaḥ nāma . \\
\noindent \textbf{(P\_6,3.66)} cakāraḥ kartavyaḥ . \\
\noindent \textbf{(P\_6,3.66)} mum ca . \\
\noindent \textbf{(P\_6,3.66)} kim ca. yat ca anyat prāpnoti . \\
\noindent \textbf{(P\_6,3.66)} kim ca anyat prāpnoti . \\
\noindent \textbf{(P\_6,3.66)} hrasvatvam . \\
\noindent \textbf{(P\_6,3.66)} sidhyati . \\
\noindent \textbf{(P\_6,3.66)} sūtram tarhi bhidyate . \\
\noindent \textbf{(P\_6,3.66)} yathānyāsam eva astu . \\
\noindent \textbf{(P\_6,3.66)} nanu ca uktam khiti hrasvāprasiddhiḥ anajantatvāt iti . \\
\noindent \textbf{(P\_6,3.66)} parihṛtam etat siddham tu hrasvāntasya mumvacanāt iti . \\
\noindent \textbf{(P\_6,3.66)} tat tarhi hrasvagrahaṇam kartavyam . \\
\noindent \textbf{(P\_6,3.66)} na kartavyam . \\
\noindent \textbf{(P\_6,3.66)} prakṛtam anuvartate . \\
\noindent \textbf{(P\_6,3.66)} kva prakṛtam . \\
\noindent \textbf{(P\_6,3.66)} ikaḥ hrasvaḥ aṅyaḥ gālavasya iti . \\
\noindent \textbf{(P\_6,3.66)} tat vai prathamānirdiṣṭam ṣaṣṭhīnirdiṣṭena ca iha arthaḥ . \\
\noindent \textbf{(P\_6,3.66)} khiti iti eṣā saptamī hrasvaḥ iti prathamāyāḥ ṣaṣṭhīm prakalpayiṣyati tasmin iti nirdiṣṭe pūrvasya iti . \\
\noindent \textbf{(P\_6,3.66)} atha vā khiti hrasvaḥ bhavati iti ucyate . \\
\noindent \textbf{(P\_6,3.66)} khiti hrasvabhāvī na asti iti kṛtvā bhūtapūrvagatiḥ vijñāsyate . \\
\noindent \textbf{(P\_6,3.66)} ajantam yat bhūtapūrvam iti . \\
\noindent \textbf{(P\_6,3.66)} atha vā kāryakālam sañjñāparibhāṣam yatra kāryam tatra draṣṭavyam . \\
\noindent \textbf{(P\_6,3.66)} khiti hrasvaḥ bhavati iti upasthitam idam bhavati acaḥ iti . \\
\noindent \textbf{(P\_6,3.66)} tatra vacanāt anajantasya api bhaviṣyati . \\
\noindent \textbf{(P\_6,3.66)} iha api tarhi vacanāt prāpnoti . \\
\noindent \textbf{(P\_6,3.66)} vāṅmanyaḥ iti . \\
\noindent \textbf{(P\_6,3.66)} na etat asti . \\
\noindent \textbf{(P\_6,3.66)} ikaḥ iti vartate . \\
\noindent \textbf{(P\_6,3.66)} evam api khaṭvammanyaḥ , atra na prāpnoti . \\
\noindent \textbf{(P\_6,3.66)} na eṣaḥ doṣaḥ . \\
\noindent \textbf{(P\_6,3.66)} ābgrahaṇam api prakṛtam anuvartate . \\
\noindent \textbf{(P\_6,3.66)} kva prakṛtam . \\
\noindent \textbf{(P\_6,3.66)} ṅyāpoḥ sañjñācchandasoḥ bahulam iti . \\
\noindent \textbf{(P\_6,3.66)} evam api kīlālapammanyaḥ , śubhaṃyammanyaḥ atra na prāpnoti . \\
\noindent \textbf{(P\_6,3.66)} tasmāt pūrvoktau eva parihārau . \\
\noindent \textbf{(P\_6,3.68.1)} amaḥ pratyayavadanudeśe kim prayojanam . \\
\noindent \textbf{(P\_6,3.68.1)} amaḥ pratyayavadanudeśe prayojanam ātvapūrvasavarṇaguṇeyaṅuvaṅādeśāḥ . \\
\noindent \textbf{(P\_6,3.68.1)} amaḥ pratyayavadanudeśe ātvapūrvasavarṇaguṇeyaṅuvaṅādeśāḥ prayojanam . \\
\noindent \textbf{(P\_6,3.68.1)} ātvam prayojanam . \\
\noindent \textbf{(P\_6,3.68.1)} gāmmanyaḥ . \\
\noindent \textbf{(P\_6,3.68.1)} pūrvasavarṇaḥ prayojanam . \\
\noindent \textbf{(P\_6,3.68.1)} strīmmanyaḥ . \\
\noindent \textbf{(P\_6,3.68.1)} guṇaḥ prayojanam . \\
\noindent \textbf{(P\_6,3.68.1)} narammanyaḥ . \\
\noindent \textbf{(P\_6,3.68.1)} iyaṅuvaṅau prayojanam . \\
\noindent \textbf{(P\_6,3.68.1)} śriyammanyaḥ , bhruvammanyaḥ . \\
\noindent \textbf{(P\_6,3.68.1)} amaḥ pratyayavadanudeśe ātvapūrvasavarṇāprasiddhiḥ aprathamātvāt . \\
\noindent \textbf{(P\_6,3.68.1)} amaḥ pratyayavadanudeśe ātvapūrvasavarṇayoḥ aprasiddhiḥ . \\
\noindent \textbf{(P\_6,3.68.1)} kim kāraṇam . \\
\noindent \textbf{(P\_6,3.68.1)} aprathamātvāt . \\
\noindent \textbf{(P\_6,3.68.1)} prathamayoḥ iti ucyate na ca atra prathamām paśyāmaḥ . \\
\noindent \textbf{(P\_6,3.68.1)} kim ca bhoḥ ātvam prathamayoḥ iti ucyate . \\
\noindent \textbf{(P\_6,3.68.1)} na khalu prathamayoḥ iti ucyate . \\
\noindent \textbf{(P\_6,3.68.1)} prathamayoḥ iti tu vijñāyate . \\
\noindent \textbf{(P\_6,3.68.1)} katham . \\
\noindent \textbf{(P\_6,3.68.1)} amśasoḥ iti ucyate . \\
\noindent \textbf{(P\_6,3.68.1)} te evam vijñāsyāmaḥ . \\
\noindent \textbf{(P\_6,3.68.1)} śassahacaritaḥ yaḥ amśabdaḥ . \\
\noindent \textbf{(P\_6,3.68.1)} kaḥ ca śassahacaritaḥ . \\
\noindent \textbf{(P\_6,3.68.1)} prathamā eva . \\
\noindent \textbf{(P\_6,3.68.1)} nanu ca pratyayavadanudeśāt bhaviṣyati . \\
\noindent \textbf{(P\_6,3.68.1)} na sidhyati . \\
\noindent \textbf{(P\_6,3.68.1)} kim kāraṇam . \\
\noindent \textbf{(P\_6,3.68.1)} sāmānyātideśe [R: hi] viśeṣānatideśaḥ . \\
\noindent \textbf{(P\_6,3.68.1)} sāmanye hi atidiśyamāne viśeṣaḥ na atidiṣṭaḥ bhavati . \\
\noindent \textbf{(P\_6,3.68.1)} tat yathā . \\
\noindent \textbf{(P\_6,3.68.1)} brahmaṇavat asmin kṣatriye vartitavyam iti sāmānyam yat brāhmaṇakāryam tat kṣatriye atidiśyate . \\
\noindent \textbf{(P\_6,3.68.1)} yat viśiṣṭam māṭhare kauṇḍinye vā na tat atidiśyate . \\
\noindent \textbf{(P\_6,3.68.1)} evam iha api sāmānyam yat pratyayakāryam tat atidiśyate yat viśiṣṭam dvitīyaikavacane bhavati prathamayoḥ iti na tat atidiśyate . \\
\noindent \textbf{(P\_6,3.68.1)} siddham tu dvitīyaikavacanavadvacanāt . \\
\noindent \textbf{(P\_6,3.68.1)} siddham etat . \\
\noindent \textbf{(P\_6,3.68.1)} katham . \\
\noindent \textbf{(P\_6,3.68.1)} dvitīyaikavacanavat bhavati iti vaktavyam . \\
\noindent \textbf{(P\_6,3.68.1)} ekaśeṣanirdeśāt vā . \\
\noindent \textbf{(P\_6,3.68.1)} atha vā ekaśeṣanirdeśaḥ ayam . \\
\noindent \textbf{(P\_6,3.68.1)} am ca am ca am . \\
\noindent \textbf{(P\_6,3.68.1)} icaḥ ekācaḥ am bhavati ampratyayavat ca asmin kāryam bhavati iti . \\
\noindent \textbf{(P\_6,3.68.2)} atha iha katham bhavitavyam . \\
\noindent \textbf{(P\_6,3.68.2)} śriyam ātmānam manyate brāhmaṇakulam . \\
\noindent \textbf{(P\_6,3.68.2)} śriyammanyam āhosvit śrimanyam iti . \\
\noindent \textbf{(P\_6,3.68.2)} śriyammanyam iti bhavitavyam . \\
\noindent \textbf{(P\_6,3.68.2)} svamoḥ napuṃsakāt iti luk kasmāt na bhavati . \\
\noindent \textbf{(P\_6,3.68.2)} na aprāpte luki am ārabhyate . \\
\noindent \textbf{(P\_6,3.68.2)} saḥ yathā eva supaḥ dhātuprātipadikayoḥ iti etam bādhate evam svamoḥ napuṃsakāt iti etam ami lukam bādheta . \\
\noindent \textbf{(P\_6,3.68.2)} na bādhate . \\
\noindent \textbf{(P\_6,3.68.2)} kim kāraṇam . \\
\noindent \textbf{(P\_6,3.68.2)} yena na aprāpte tasya bādhanam bhavati . \\
\noindent \textbf{(P\_6,3.68.2)} na ca aprāpte supaḥ dhātuprātipadikayoḥ iti etasmin etat ārabhyate . \\
\noindent \textbf{(P\_6,3.68.2)} svamoḥ napuṃsakāt iti etasmin punaḥ prāpte ca aprāpte ca . \\
\noindent \textbf{(P\_6,3.68.2)} atha vā madhye apavādāḥ pūrvān vidhīn bādhante iti evam supaḥ dhātuprātipadikayoḥ iti etam bādhate . \\
\noindent \textbf{(P\_6,3.68.2)} svamoḥ napuṃsakāt iti etam na bādhiṣyate . \\
\noindent \textbf{(P\_6,3.68.2)} evam tarhi asiddham bahiraṅgam antaraṅge iti asiddhatvāt bahiraṅgalakṣaṇasya amaḥ antaraṅgalakṣaṇaḥ luk na bhaviṣyati . \\
\noindent \textbf{(P\_6,3.68.2)} na eṣā paribhāṣā uttarapadādhikāre śakyā vijñātum . \\
\noindent \textbf{(P\_6,3.68.2)} iha hi doṣaḥ syāt . \\
\noindent \textbf{(P\_6,3.68.2)} dviṣantapaḥ , parantapaḥ . \\
\noindent \textbf{(P\_6,3.68.2)} saṃyogāntalopaḥ na syāt . \\
\noindent \textbf{(P\_6,3.68.2)} tasmāt śrimanyam iti eva bhavitavyam . \\
\noindent \textbf{(P\_6,3.70)} astusatyāgadasya kāre . \\
\noindent \textbf{(P\_6,3.70)} astusatyāgadasya kāre upasaṅkhyānam kartavyam . \\
\noindent \textbf{(P\_6,3.70)} astuṅkāraḥ , satyaṅkāraḥ , agadaṅkāraḥ . \\
\noindent \textbf{(P\_6,3.70)} bhakṣasya chandasi . \\
\noindent \textbf{(P\_6,3.70)} bhakṣasya chandasi upasaṅkhyānam kartavyam . \\
\noindent \textbf{(P\_6,3.70)} tasya te bhakṣaṅkārasya . \\
\noindent \textbf{(P\_6,3.70)} chandasi iti kim . \\
\noindent \textbf{(P\_6,3.70)} bhakṣakārasya tat matam iti . \\
\noindent \textbf{(P\_6,3.70)} dhenoḥ bhavyāyām . \\
\noindent \textbf{(P\_6,3.70)} dhenoḥ bhavyāyām upasaṅkhyānam kartavyam . \\
\noindent \textbf{(P\_6,3.70)} dhenumbhavyā . \\
\noindent \textbf{(P\_6,3.70)} lokasya pṛṇe . \\
\noindent \textbf{(P\_6,3.70)} lokasya pṛṇe upasaṅkhyānam kartavyam . \\
\noindent \textbf{(P\_6,3.70)} lokamprṇasya dhanvinaḥ . \\
\noindent \textbf{(P\_6,3.70)} itye anabhyāśasya . \\
\noindent \textbf{(P\_6,3.70)} itye anabhyāśasya upasaṅkhyānam kartavyam . \\
\noindent \textbf{(P\_6,3.70)} anabhyāśamityaḥ . \\
\noindent \textbf{(P\_6,3.70)} bhrāṣṭrāgnyoḥ indhe . \\
\noindent \textbf{(P\_6,3.70)} bhrāṣṭrāgnyoḥ indhe upasaṅkhyānam kartavyam . \\
\noindent \textbf{(P\_6,3.70)} bhrāṣṭramindhaḥ , agnimindhaḥ . \\
\noindent \textbf{(P\_6,3.70)} gile agilasya . \\
\noindent \textbf{(P\_6,3.70)} gile agilasya upasaṅkhyānam kartavyam . \\
\noindent \textbf{(P\_6,3.70)} timiṅgilaḥ . \\
\noindent \textbf{(P\_6,3.70)} agilasya iti kimartham . \\
\noindent \textbf{(P\_6,3.70)} gilagilaḥ . \\
\noindent \textbf{(P\_6,3.70)} gilagile ca iti vaktavyam . \\
\noindent \textbf{(P\_6,3.70)} timiṅgilagilaḥ . \\
\noindent \textbf{(P\_6,3.70)} uṣṇabhadrayoḥ karaṇe . \\
\noindent \textbf{(P\_6,3.70)} uṣṇabhadrayoḥ karaṇe upasaṅkhyānam kartavyam . \\
\noindent \textbf{(P\_6,3.70)} uṣṇaṅkaraṇam , bhadraṅkaraṇam . \\
\noindent \textbf{(P\_6,3.70)} sūtograrājabhojakulamerubhyaḥ duhituḥ putraṭ vā . \\
\noindent \textbf{(P\_6,3.70)} sūtograrājabhojakulamerubhyaḥ duhituḥ putraṭ vā bhavati iti vaktavyam . \\
\noindent \textbf{(P\_6,3.70)} sūtaputrī , sūtaduhitā , ugraputrī , ugraduhitā , rajaputrī , rājaduhitā , bhojaputrī , bhojaduhitā , kulaputrī , kuladuhitā , meruputrī , meruduhitā . \\
\noindent \textbf{(P\_6,3.72)} kim iyam prāpte vibhāṣā āhosvit aprāpte . \\
\noindent \textbf{(P\_6,3.72)} katham ca prāpte katham vā aprāpte . \\
\noindent \textbf{(P\_6,3.72)} khiti iti vā nitye prāpte anyatra vā aprāpte . \\
\noindent \textbf{(P\_6,3.72)} rātreḥ aprāpte . \\
\noindent \textbf{(P\_6,3.72)} rātreḥ aprāpte vibhāṣā . \\
\noindent \textbf{(P\_6,3.72)} prāpte nityaḥ vidhiḥ . \\
\noindent \textbf{(P\_6,3.72)} rātrimmanyaḥ . \\
\noindent \textbf{(P\_6,3.72)} aprāpte vibhāṣā . \\
\noindent \textbf{(P\_6,3.72)} rātryaṭaḥ , rātrimaṭaḥ . \\
\noindent \textbf{(P\_6,3.73)} kimartham nañaḥ sānubandhakasya grahaṇam kriyate na nasya iti eva ucyeta . \\
\noindent \textbf{(P\_6,3.73)} nasya iti ucyamāne karṇaputraḥ , varṇaputraḥ iti atra api prasajyeta . \\
\noindent \textbf{(P\_6,3.73)} na eṣaḥ doṣaḥ . \\
\noindent \textbf{(P\_6,3.73)} arthavadgrahaṇe na anarthakasya iti evam etasya na bhaviṣyati . \\
\noindent \textbf{(P\_6,3.73)} evam api praśnaputraḥ , viśnaputraḥ iti atra api prāpnoti . \\
\noindent \textbf{(P\_6,3.73)} na eṣaḥ doṣaḥ . \\
\noindent \textbf{(P\_6,3.73)} ananubandhakagrahaṇe na sānubandhakasya iti evam etasya na bhaviṣyati . \\
\noindent \textbf{(P\_6,3.73)} evam api vāmanaputraḥ , pāmanaputraḥ iti atra api prāpnoti . \\
\noindent \textbf{(P\_6,3.73)} tasmāt sānubandhakasya grahaṇam kartavyam . \\
\noindent \textbf{(P\_6,3.73)} nañaḥ nalope avakṣepe tiṅi upasaṅkhyānam . \\
\noindent \textbf{(P\_6,3.73)} nañaḥ nalope avakṣepe tiṅi upasaṅkhyānam kartavyam . \\
\noindent \textbf{(P\_6,3.73)} apacasi vai tvam jālma . \\
\noindent \textbf{(P\_6,3.73)} akaroṣi vai tvam jālma . \\
\noindent \textbf{(P\_6,3.74)} kimartham tasmāt iti ucyate na nuṭ aci iti eva ucyeta . \\
\noindent \textbf{(P\_6,3.74)} nuṭ aci iti ucyamāne nañaḥ eva nuṭ prasajyeta . \\
\noindent \textbf{(P\_6,3.74)} evam tarhi pūrvāntaḥ kariṣyate . \\
\noindent \textbf{(P\_6,3.74)} tatra ayam api arthaḥ . \\
\noindent \textbf{(P\_6,3.74)} tadoḥ saḥ sau anantyayoḥ iti tadoḥ grahaṇam na kartavyam . \\
\noindent \textbf{(P\_6,3.74)} tatra hi tavargānirdeśe etat prayojanam iha mā bhūt . \\
\noindent \textbf{(P\_6,3.74)} aneṣaḥ karoti iti . \\
\noindent \textbf{(P\_6,3.74)} yāvatā pūrvāntaḥ saḥ api adoṣaḥ bhavati . \\
\noindent \textbf{(P\_6,3.74)} na evam śakyam . \\
\noindent \textbf{(P\_6,3.74)} anuṣṇaḥ iti nalopaḥ prātipadikāntasya iti nalopaḥ prasajyeta . \\
\noindent \textbf{(P\_6,3.74)} nugvacanāt na bhaviṣyati . \\
\noindent \textbf{(P\_6,3.74)} ṅamuṭ tarhi prāpnoti . \\
\noindent \textbf{(P\_6,3.74)} tasmāt parādiḥ kartavyaḥ . \\
\noindent \textbf{(P\_6,3.74)} parādau ca kriyamāṇe tasmāt iti vaktavyam . \\
\noindent \textbf{(P\_6,3.76)} kimartham āduk ucyate na aduk eva ucyate . \\
\noindent \textbf{(P\_6,3.76)} kā rūpasiddhiḥ : ekānnaviṃśatiḥ , ekānnaśatam . \\
\noindent \textbf{(P\_6,3.76)} savarṇadīrghatvena siddham . \\
\noindent \textbf{(P\_6,3.76)} na sidhyati . \\
\noindent \textbf{(P\_6,3.76)} ataḥ guṇe iti pararūpatvam prāpnoti . \\
\noindent \textbf{(P\_6,3.76)} evam tarhi aduṭ kariṣyate . \\
\noindent \textbf{(P\_6,3.76)} aduṭ ca aśakyaḥ kartum . \\
\noindent \textbf{(P\_6,3.76)} ānunāsikyam hi na syāt . \\
\noindent \textbf{(P\_6,3.76)} yat hi tat yaraḥ anunāsike anunāsikaḥ va iti padāntasya iti evam tat . \\
\noindent \textbf{(P\_6,3.76)} kim punaḥ kāraṇam padāntasya iti evam tat . \\
\noindent \textbf{(P\_6,3.76)} iha mā bhūt . \\
\noindent \textbf{(P\_6,3.76)} budhnaḥ , bradhnaḥ , badhnāti . \\
\noindent \textbf{(P\_6,3.76)} evam tarhi anuṭ kariṣyate . \\
\noindent \textbf{(P\_6,3.76)} anuṭ ca aśakyaḥ kartum . \\
\noindent \textbf{(P\_6,3.76)} vibhāṣayā ānunāsikyam . \\
\noindent \textbf{(P\_6,3.76)} tena idam eva rūpam syāt ekānnaviṃśatiḥ . \\
\noindent \textbf{(P\_6,3.76)} idam na syāt . \\
\noindent \textbf{(P\_6,3.76)} ekānnaviṃśatiḥ iti . \\
\noindent \textbf{(P\_6,3.76)} astu tarhi aduk eva . \\
\noindent \textbf{(P\_6,3.76)} nanu ca uktam ataḥ guṇe iti pararūpatvam prāpnoti iti . \\
\noindent \textbf{(P\_6,3.76)} na eṣaḥ doṣaḥ . \\
\noindent \textbf{(P\_6,3.76)} akāroccāraṇasāmarthyāt na bhaviṣyati . \\
\noindent \textbf{(P\_6,3.76)} yadi tarhi prāpnuvan vidhiḥ akāroccāraṇasāmarthyāt bādhyate savarṇadīrghatvam api na prāpnoti . \\
\noindent \textbf{(P\_6,3.76)} yam vidhim prati upadeśaḥ anarthakaḥ saḥ vidhiḥ bādhyate . \\
\noindent \textbf{(P\_6,3.76)} yasya tu vidhiḥ nimittam eva na asau bādhyate . \\
\noindent \textbf{(P\_6,3.76)} pararūpam ca prati akāroccāraṇam anarthakam savarṇadīrghatvasya punaḥ nimittam eva . \\
\noindent \textbf{(P\_6,3.78)} sahasya halopavacanam . \\
\noindent \textbf{(P\_6,3.78)} sahasya halopaḥ vaktavyaḥ . \\
\noindent \textbf{(P\_6,3.78)} sādeśe hi svare doṣaḥ . \\
\noindent \textbf{(P\_6,3.78)} sādeśe hi [sati] svare doṣaḥ syāt . \\
\noindent \textbf{(P\_6,3.78)} āntaryataḥ udāttānudāttayoḥ [sthāne] svaritaḥ ādeśaḥ prasajyeta . \\
\noindent \textbf{(P\_6,3.78)} [saputraḥ , sabhāryaḥ .] saḥ tarhi lopaḥ vaktatvyaḥ . \\
\noindent \textbf{(P\_6,3.78)} na vaktatvyaḥ . \\
\noindent \textbf{(P\_6,3.78)} ādyudāttanipātanam kariṣyate . \\
\noindent \textbf{(P\_6,3.78)} saḥ nipātanasvaraḥ prakṛtisvarasya bādhakaḥ bhaviṣyati . \\
\noindent \textbf{(P\_6,3.78)} evam api upadeśivadbhāvaḥ vaktavyaḥ . \\
\noindent \textbf{(P\_6,3.78)} saḥ yathā eva hi nipātanasvaraḥ prakṛtisvaram bādhate evam samāsasvaram api bādheta . \\
\noindent \textbf{(P\_6,3.78)} seṣṭi , sapaśubandham . \\
\noindent \textbf{(P\_6,3.79)} granthānte vacanānarthakyam avyayībhāvena kṛtatvāt . granthānte vacanam anarthakam . \\
\noindent \textbf{(P\_6,3.79)} kim kāraṇam . \\
\noindent \textbf{(P\_6,3.79)} avyayībhāvena kṛtatvāt . \\
\noindent \textbf{(P\_6,3.79)} avyayībhāve ca akāle iti eva siddham . \\
\noindent \textbf{(P\_6,3.79)} yaḥ tarhi kālottarapadaḥ granthāntaḥ tadartham idam vaktavyam . \\
\noindent \textbf{(P\_6,3.79)} sakāṣṭham jyotiṣam adhīte . \\
\noindent \textbf{(P\_6,3.79)} sakalam , samuhūrtam . \\
\noindent \textbf{(P\_6,3.82):} KA\_III,170.19-171.11 Ro\_IV,646-647 {1/27}     upasarjanasya vāvacane sarvaprasaṅgaḥ aviśeṣāt . \\
\noindent \textbf{(P\_6,3.82):} KA\_III,170.19-171.11 Ro\_IV,646-647 {2/27}     upasarjanasya vāvacane sarvaprasaṅgaḥ . \\
\noindent \textbf{(P\_6,3.82):} KA\_III,170.19-171.11 Ro\_IV,646-647 {3/27}     sarvasya upasarjanasya sādeśaḥ prāpnoti . \\
\noindent \textbf{(P\_6,3.82):} KA\_III,170.19-171.11 Ro\_IV,646-647 {4/27}     asya api prāpnoti : sahayudhvā , sahakṛtvā . \\
\noindent \textbf{(P\_6,3.82):} KA\_III,170.19-171.11 Ro\_IV,646-647 {5/27}     kim kāraṇam . \\
\noindent \textbf{(P\_6,3.82):} KA\_III,170.19-171.11 Ro\_IV,646-647 {6/27}     aviśeṣāt . \\
\noindent \textbf{(P\_6,3.82):} KA\_III,170.19-171.11 Ro\_IV,646-647 {7/27}     na hi kaḥ cit viśeṣaḥ upādīyate evañjātīyakasya sādeśaḥ bhavati iti . \\
\noindent \textbf{(P\_6,3.82):} KA\_III,170.19-171.11 Ro\_IV,646-647 {8/27}     anupādīyamāne viśeṣe sarvaprasaṅgaḥ . \\
\noindent \textbf{(P\_6,3.82):} KA\_III,170.19-171.11 Ro\_IV,646-647 {9/27}     siddham tu bahuvrīhinirdeśāt . \\
\noindent \textbf{(P\_6,3.82):} KA\_III,170.19-171.11 Ro\_IV,646-647 {10/27}     siddham etat . \\
\noindent \textbf{(P\_6,3.82):} KA\_III,170.19-171.11 Ro\_IV,646-647 {11/27}     katham . \\
\noindent \textbf{(P\_6,3.82):} KA\_III,170.19-171.11 Ro\_IV,646-647 {12/27}     bahuvrīhinirdeśāt . \\
\noindent \textbf{(P\_6,3.82):} KA\_III,170.19-171.11 Ro\_IV,646-647 {13/27}     bahuvrīhinirdeśaḥ kartavyaḥ . \\
\noindent \textbf{(P\_6,3.82):} KA\_III,170.19-171.11 Ro\_IV,646-647 {14/27}     evam api sahayudhvapriyaḥ , sahakṛtvapriyaḥ iti atra prāpnoti . \\
\noindent \textbf{(P\_6,3.82):} KA\_III,170.19-171.11 Ro\_IV,646-647 {15/27}     bahuvrīhau yat uttarapadam iti evam vijñāsyate . \\
\noindent \textbf{(P\_6,3.82):} KA\_III,170.19-171.11 Ro\_IV,646-647 {16/27}     nanu etat api bahuvrīhau uttarapadam . \\
\noindent \textbf{(P\_6,3.82):} KA\_III,170.19-171.11 Ro\_IV,646-647 {17/27}     evam tarhi bahuvrīhau yat upasarjanam iti evam vijñāsyate . \\
\noindent \textbf{(P\_6,3.82):} KA\_III,170.19-171.11 Ro\_IV,646-647 {18/27}     bahuvrīhau ca yat upasarjanam bahuvrīhim prati ca yat upasarjanam . \\
\noindent \textbf{(P\_6,3.82):} KA\_III,170.19-171.11 Ro\_IV,646-647 {19/27}     saḥ tarhi bahuvrīhinirdeśaḥ kartavyaḥ . \\
\noindent \textbf{(P\_6,3.82):} KA\_III,170.19-171.11 Ro\_IV,646-647 {20/27}     na kartavyaḥ . \\
\noindent \textbf{(P\_6,3.82):} KA\_III,170.19-171.11 Ro\_IV,646-647 {21/27}     iha kaḥ cit pradhānānām eva samāsaḥ kaḥ cit upasarjanānām eva kaḥ cit pradhānopasarjanānām . \\
\noindent \textbf{(P\_6,3.82):} KA\_III,170.19-171.11 Ro\_IV,646-647 {22/27}     tat yaḥ upasarjanānām eva samāsaḥ tat upasarjanam . \\
\noindent \textbf{(P\_6,3.82):} KA\_III,170.19-171.11 Ro\_IV,646-647 {23/27}     atha vā akāraḥ matvarthīyaḥ . \\
\noindent \textbf{(P\_6,3.82):} KA\_III,170.19-171.11 Ro\_IV,646-647 {24/27}      tat yathā tundaḥ ghāṭaḥ iti . \\
\noindent \textbf{(P\_6,3.82):} KA\_III,170.19-171.11 Ro\_IV,646-647 {25/27}     atha vā matublopaḥ atra draṣṭavyaḥ . \\
\noindent \textbf{(P\_6,3.82):} KA\_III,170.19-171.11 Ro\_IV,646-647 {26/27}     tat yathā puṣyakāḥ eṣām te ime puṣyakāḥ . \\
\noindent \textbf{(P\_6,3.82):} KA\_III,170.19-171.11 Ro\_IV,646-647 {27/27}     kālakāḥ eṣām te ime kālakāḥ iti . \\
\noindent \textbf{(P\_6,3.83):} KA\_III,171.13-14 Ro\_IV,648 {1/4}     prakṛtyā āśiṣi agavādiṣu . \\
\noindent \textbf{(P\_6,3.83):} KA\_III,171.13-14 Ro\_IV,648 {2/4}     prakṛtyā āśiṣi agavādiṣu iti vaktavyam . \\
\noindent \textbf{(P\_6,3.83):} KA\_III,171.13-14 Ro\_IV,648 {3/4}     iha mā bhūt . \\
\noindent \textbf{(P\_6,3.83):} KA\_III,171.13-14 Ro\_IV,648 {4/4}     sagave savatsāya sahalāya iti . \\
\noindent \textbf{(P\_6,3.86)} caraṇe kim nipātyate . \\
\noindent \textbf{(P\_6,3.86)} brahmaṇi upapade samānapūrve vrate karmaṇi careḥ ṇiniḥ vratalopaḥ ca . \\
\noindent \textbf{(P\_6,3.86)} brahmaṇi upapade samānapūrve vrate karmaṇi careḥ ṇiniḥ pratyayaḥ vratalopaḥ ca nipātyate . \\
\noindent \textbf{(P\_6,3.86)} samāne brahmaṇi vratam catarti iti sabrahmacārī . \\
\noindent \textbf{(P\_6,3.89)} dṛgdṛśavatuṣu dṛkṣe upasaṅkhyānam . dṛgdṛśavatuṣu dṛkṣe upasaṅkhyānam kartavyam . \\
\noindent \textbf{(P\_6,3.89)} sadṛkṣāsaḥ pratisadṛkṣāsaḥ . \\
\noindent \textbf{(P\_6,3.93,} 116) KA\_III,172.2-11 Ro\_IV,649-650 {1/15}     kimartham añcatinahyādiṣu kvibgrahaṇam kriyate . \\
\noindent \textbf{(P\_6,3.93,} 116) KA\_III,172.2-11 Ro\_IV,649-650 {2/15}     iha mā bhūt . \\
\noindent \textbf{(P\_6,3.93,} 116) KA\_III,172.2-11 Ro\_IV,649-650 {3/15}     samañcanam , upanahanam . \\
\noindent \textbf{(P\_6,3.93,} 116) KA\_III,172.2-11 Ro\_IV,649-650 {4/15}     na etat asti prayojanam . \\
\noindent \textbf{(P\_6,3.93,} 116) KA\_III,172.2-11 Ro\_IV,649-650 {5/15}     uttarapade iti vartate na ca antareṇa kvipam añcatinahyādayaḥ uttarapadāni bhavanti . \\
\noindent \textbf{(P\_6,3.93,} 116) KA\_III,172.2-11 Ro\_IV,649-650 {6/15}     tatra antareṇa kvibgrahaṇam kvibante eva bhaviṣyati . \\
\noindent \textbf{(P\_6,3.93,} 116) KA\_III,172.2-11 Ro\_IV,649-650 {7/15}     tadādividhinā prāpnoti . \\
\noindent \textbf{(P\_6,3.93,} 116) KA\_III,172.2-11 Ro\_IV,649-650 {8/15}     ataḥ uttaram paṭhati añcatinahyādiṣu kvibgrahaṇanārthakyam yasmin vidhiḥ tadādau algrahaṇe . \\
\noindent \textbf{(P\_6,3.93,} 116) KA\_III,172.2-11 Ro\_IV,649-650 {9/15}     añcatinahyādiṣu kvibgrahaṇam anarthakam . \\
\noindent \textbf{(P\_6,3.93,} 116) KA\_III,172.2-11 Ro\_IV,649-650 {10/15}     kim kāraṇam . \\
\noindent \textbf{(P\_6,3.93,} 116) KA\_III,172.2-11 Ro\_IV,649-650 {11/15}     yasmin vidhiḥ tadādau algrahaṇe . \\
\noindent \textbf{(P\_6,3.93,} 116) KA\_III,172.2-11 Ro\_IV,649-650 {12/15}     algrahaṇeṣu etat bhavati na ca idam algrahaṇam . \\
\noindent \textbf{(P\_6,3.93,} 116) KA\_III,172.2-11 Ro\_IV,649-650 {13/15}     evam tarhi siddhe sati yat kvibgrahaṇam karoti tat jñāpayati ācāryaḥ anyatra dhātugrahaṇe tadādividhiḥ bhavati iti . \\
\noindent \textbf{(P\_6,3.93,} 116) KA\_III,172.2-11 Ro\_IV,649-650 {14/15}     kim etasya jñāpane prayojanam . \\
\noindent \textbf{(P\_6,3.93,} 116) KA\_III,172.2-11 Ro\_IV,649-650 {15/15}     ataḥ kṛkami iti atra , ayaskṛt ayaskāra iti api siddham bhavati . \\
\noindent \textbf{(P\_6,3.92)} adrisadhryoḥ antodāttavacanam kṛtsvaranivṛttyartham . \\
\noindent \textbf{(P\_6,3.92)} adrisadhryoḥ antodāttatvam vaktavyam . \\
\noindent \textbf{(P\_6,3.92)} kim prayojanam . \\
\noindent \textbf{(P\_6,3.92)} kṛtsvaranivṛttyartham . \\
\noindent \textbf{(P\_6,3.92)} kṛtsvaraḥ mā bhūt . \\
\noindent \textbf{(P\_6,3.92)} viṣvadryaṅ , viṣvadryañcau , viṣvadryañcaḥ , sadhryaṅ , sadhryañcau , sadhryañcaḥ . \\
\noindent \textbf{(P\_6,3.92)} tatra chandasi striyām pratiṣedhaḥ . \\
\noindent \textbf{(P\_6,3.92)} tatra chandasi striyām pratiṣedhaḥ vaktavyaḥ . \\
\noindent \textbf{(P\_6,3.92)} viśvācī , ghṛtācī . \\
\noindent \textbf{(P\_6,3.92)} yadi chandasi striyām pratiṣedhaḥ ucyate katham sā kadrīcī . \\
\noindent \textbf{(P\_6,3.92)} evam tarhi chandasi striyām bahulam iti vaktavyam . \\
\noindent \textbf{(P\_6,3.97)} samāpaḥ īttvapratiṣedhaḥ . \\
\noindent \textbf{(P\_6,3.97)} samāpaḥ īttvapratiṣedhaḥ vaktavyaḥ . \\
\noindent \textbf{(P\_6,3.97)} samāpam nāma devayajanam . \\
\noindent \textbf{(P\_6,3.97)} aparaḥ āha : īttvam anavarṇāt iti vaktavyam . \\
\noindent \textbf{(P\_6,3.97)} samīpam , antarīpam . \\
\noindent \textbf{(P\_6,3.97)} iha mā bhūt . \\
\noindent \textbf{(P\_6,3.97)} prāpam , parāpam . \\
\noindent \textbf{(P\_6,3.98)} dīrghoccāraṇam kimartham na udanoḥ deśe iti eva ucyeta . \\
\noindent \textbf{(P\_6,3.98)} kā rūpasiddhiḥ : anūpaḥ . \\
\noindent \textbf{(P\_6,3.98)} savarṇadīrghatven siddham . \\
\noindent \textbf{(P\_6,3.98)} na sidhyati . \\
\noindent \textbf{(P\_6,3.98)} avagrahe doṣaḥ syāt . \\
\noindent \textbf{(P\_6,3.99)} aṣaṣṭhyatṛtīyasthasya iti ucyate . \\
\noindent \textbf{(P\_6,3.99)} tatra idam na sidhyati . \\
\noindent \textbf{(P\_6,3.99)} anyasya idam anyadīyam . \\
\noindent \textbf{(P\_6,3.99)} anyasya kārakam anyatkārakam . \\
\noindent \textbf{(P\_6,3.99)} evam tarhi aviśeṣeṇa anyasya duk chakārakayoḥ iti uktvā tataḥ vakṣyāmi aṣaṣṭhyatṛtīyasthasya āśīrāśāsthāsthitotsukotirāgeṣu iti . \\
\noindent \textbf{(P\_6,3.101)} kadbhāve trau upasaṅkhyānam . \\
\noindent \textbf{(P\_6,3.101)} kadbhāve trau upasaṅkhyānam kartavyam . \\
\noindent \textbf{(P\_6,3.101)} kutsitāḥ trayaḥ kattrayaḥ . \\
\noindent \textbf{(P\_6,3.101)} ke vā trayaḥ . \\
\noindent \textbf{(P\_6,3.101)} na bibhṛyuḥ kattrayaḥ \\
\noindent \textbf{(P\_6,3.109.1)} pṛṣodarādīni iti ucyate . \\
\noindent \textbf{(P\_6,3.109.1)} kāni pṛṣodarādīni . \\
\noindent \textbf{(P\_6,3.109.1)} pṛṣodaraprakārāṇi . \\
\noindent \textbf{(P\_6,3.109.1)} kāni punaḥ pṛṣodaraprakārāṇi . \\
\noindent \textbf{(P\_6,3.109.1)} yeṣu lopāgamavikārāḥ śrūyante na ca ucyante . \\
\noindent \textbf{(P\_6,3.109.1)} atha yathā iti kim idam . \\
\noindent \textbf{(P\_6,3.109.1)} prakāravacane thāl . \\
\noindent \textbf{(P\_6,3.109.1)} atha kim idam upadiṣṭāni iti . \\
\noindent \textbf{(P\_6,3.109.1)} uccāritāni . \\
\noindent \textbf{(P\_6,3.109.1)} kutaḥ etat . \\
\noindent \textbf{(P\_6,3.109.1)} diśiḥ uccāraṇakriyaḥ . \\
\noindent \textbf{(P\_6,3.109.1)} uccārya hi varṇān āha updiṣṭāḥ ime varṇāḥ iti . \\
\noindent \textbf{(P\_6,3.109.1)} kaiḥ punaḥ upadiṣṭāḥ . \\
\noindent \textbf{(P\_6,3.109.1)} śiṣṭaiḥ . \\
\noindent \textbf{(P\_6,3.109.1)} ke punaḥ śiṣṭāḥ . \\
\noindent \textbf{(P\_6,3.109.1)} vaiyākaraṇāḥ . \\
\noindent \textbf{(P\_6,3.109.1)} kutaḥ etat . \\
\noindent \textbf{(P\_6,3.109.1)} śāstrapūrvikā hi śiṣṭiḥ vaiyākaraṇāḥ ca śāstrajñāḥ . \\
\noindent \textbf{(P\_6,3.109.1)} yadi tarhi śāstrapūrvikā śiṣṭiḥ śiṣṭipūrvakam ca śāstram tat itaretarāśrayam bhavati . \\
\noindent \textbf{(P\_6,3.109.1)} itaretarāśrayāṇi ca na prakalpante . \\
\noindent \textbf{(P\_6,3.109.1)} evam tarhi nivāsataḥ ācārataḥ ca . \\
\noindent \textbf{(P\_6,3.109.1)} saḥ ca ācāraḥ āryāvartte eva . \\
\noindent \textbf{(P\_6,3.109.1)} kaḥ punaḥ āryāvarttaḥ . \\
\noindent \textbf{(P\_6,3.109.1)} prāk ādarśāt [R adarśanāt] pratyak kālakavanāt dakṣiṇena himavantam uttareṇa pāriyātram . \\
\noindent \textbf{(P\_6,3.109.1)} etasmin āryanivāse ye brāhmaṇāḥ kumbhīdhānyāḥ alolupāḥ agṛhyamāṇakāraṇāḥ kim cit antareṇa kasyāḥ cit vidyāyāḥ pāragāḥ tatrabhavantaḥ śiṣṭāḥ . \\
\noindent \textbf{(P\_6,3.109.1)} yadi tarhi śiṣṭāḥ śabdeṣu pramāṇam kim aṣṭādhyāyyā kriyate . \\
\noindent \textbf{(P\_6,3.109.1)} śiṣṭajñānārthā aṣṭādhyāyī . \\
\noindent \textbf{(P\_6,3.109.1)} katham punaḥ aṣṭādhyāyyā śiṣṭāḥ śakyāḥ vijñātum . \\
\noindent \textbf{(P\_6,3.109.1)} aṣṭādhyāyīm adhīyānaḥ anyam paśyati anadhīyānam ye atra vihitāḥ śabdāḥ tān prayuñjānam . \\
\noindent \textbf{(P\_6,3.109.1)} saḥ paśyati . \\
\noindent \textbf{(P\_6,3.109.1)} nūnam asya daivānugrahaḥ svabhāvaḥ vā yaḥ ayam na ca aṣṭādhyāyīm adhīte ye ca asyam vihitāḥ śabdāḥ tān prayuṅkte . \\
\noindent \textbf{(P\_6,3.109.1)} nūnam ayam anyān api jānāti . \\
\noindent \textbf{(P\_6,3.109.1)} evam eṣā śiṣṭajñānārthā aṣṭādhyāyī . \\
\noindent \textbf{(P\_6,3.109.2)} dikśabdebhyaḥ tīrasya tārabhāvaḥ vā . \\
\noindent \textbf{(P\_6,3.109.2)} dikśabdebhyaḥ tīrasya tārabhāvaḥ vā vaktavyaḥ . \\
\noindent \textbf{(P\_6,3.109.2)} dakṣiṇatīram , dakṣiṇatāram . \\
\noindent \textbf{(P\_6,3.109.2)} vācaḥ vāde ḍatvam valabhāvaḥ ca uttarapadasya iñi . \\
\noindent \textbf{(P\_6,3.109.2)} vācaḥ vāde ḍatvam vaktavyam valabhāvaḥ ca uttarapadasya iñi vaktavyaḥ . \\
\noindent \textbf{(P\_6,3.109.2)} vāgvādasya apatyam vāḍvaliḥ . \\
\noindent \textbf{(P\_6,3.109.2)} ṣaṣaḥ utvam datṛdaśasu uttarapadādeḥ ṣṭutvam ca . \\
\noindent \textbf{(P\_6,3.109.2)} ṣaṣaḥ utvam vaktavyam uttarapadādeḥ ṣṭutvam ca vaktavyam . \\
\noindent \textbf{(P\_6,3.109.2)} ṣoḍaśan , ṣoḍaśa . \\
\noindent \textbf{(P\_6,3.109.2)} dhāsu vā . \\
\noindent \textbf{(P\_6,3.109.2)} dhāsu vā iti vaktavyam uttarapadādeḥ ṣṭutvam ca vaktavyam . \\
\noindent \textbf{(P\_6,3.109.2)} ṣoḍhā ṣaḍḍhā kuru . \\
\noindent \textbf{(P\_6,3.109.2)} atha kimartham bahuvacananirdeśaḥ kriyate na punaḥ dhāyām iti eva ucyate . \\
\noindent \textbf{(P\_6,3.109.2)} nānādhikaraṇavācī yaḥ dhāśabdaḥ tasya grahaṇam yathā vijñāyeta . \\
\noindent \textbf{(P\_6,3.109.2)} iha mā bhūt . \\
\noindent \textbf{(P\_6,3.109.2)} ṣaṭ dadhāti iti ṣaḍdhā iti . \\
\noindent \textbf{(P\_6,3.109.2)} duraḥ dāśanāśadabhadhyeṣu . \\
\noindent \textbf{(P\_6,3.109.2)} duraḥ dāśanāśadabhadhyeṣu utvam vaktavyam uttarapadādeḥ ca ṣṭutvam . \\
\noindent \textbf{(P\_6,3.109.2)} dūḍāśaḥ , dūṇāśaḥ , dūḍabhaḥ , dūḍhyaḥ . \\
\noindent \textbf{(P\_6,3.109.2)} svaro rohatau chandasi . \\
\noindent \textbf{(P\_6,3.109.2)} svaro rohatau chandasi utvam vaktavyam . \\
\noindent \textbf{(P\_6,3.109.2)} ehi tvam jāye svo rohāva . \\
\noindent \textbf{(P\_6,3.109.2)} pīvopavasanādīnām chandasi lopaḥ vaktavyaḥ . \\
\noindent \textbf{(P\_6,3.109.2)} pīvopavasanānām payopavasanānām śriyā idam . \\
\noindent \textbf{(P\_6,3.111)} pūrvagrahaṇam kimartham na tasmin iti nirdiṣṭe pūrvasya iti pūrvasya eva bhaviṣyati . \\
\noindent \textbf{(P\_6,3.111)} na sidhyati . \\
\noindent \textbf{(P\_6,3.111)} na hi ḍhralopena ānantaryam . \\
\noindent \textbf{(P\_6,3.111)} iha kasmāt na bhavati karaṇīyam , haraṇīyam . \\
\noindent \textbf{(P\_6,3.111)} na evam vijñāyate ḍhroḥ lopaḥ ḍhralopaḥ , ḍhralope iti . \\
\noindent \textbf{(P\_6,3.111)} katham tarhi . \\
\noindent \textbf{(P\_6,3.111)} ḍhroḥ lopaḥ asmin saḥ ayam ḍhralopaḥ , ḍhralope iti . \\
\noindent \textbf{(P\_6,3.111)} yadi evam na arthaḥ pūrvagrahaṇena . \\
\noindent \textbf{(P\_6,3.111)} bhavati hi ḍhralopena ānantaryam . \\
\noindent \textbf{(P\_6,3.111)} idam tarhi prayojanam . \\
\noindent \textbf{(P\_6,3.111)} uttarapade iti vartate . \\
\noindent \textbf{(P\_6,3.111)} tena ānantaryamātre yathā syāt . \\
\noindent \textbf{(P\_6,3.111)} audumbariḥ rājā . \\
\noindent \textbf{(P\_6,3.111)} punaḥ rūpāṇi kalpayet . \\
\noindent \textbf{(P\_6,3.112)} varṇagrahaṇam kimartham na sahivahoḥ ot asya iti eva ucyeta . \\
\noindent \textbf{(P\_6,3.112)} vṛddhau api kṛtāyām yathā syāt . \\
\noindent \textbf{(P\_6,3.112)} udavoḍhām , udavoḍham , udavoḍha iti . \\
\noindent \textbf{(P\_6,3.112)} atha avarṇagrahaṇam kimartham . \\
\noindent \textbf{(P\_6,3.112)} iha mā bhūt . \\
\noindent \textbf{(P\_6,3.112)} ūḍhaḥ , ūḍhavān iti . \\
\noindent \textbf{(P\_6,3.112)} na etat asti prayojanam . \\
\noindent \textbf{(P\_6,3.112)} bhavatu atra ottvam . \\
\noindent \textbf{(P\_6,3.112)} śravaṇam kasmāt na bhavati . \\
\noindent \textbf{(P\_6,3.112)} pūrvatvam asya bhaviṣyati . \\
\noindent \textbf{(P\_6,3.112)} idam iha sampradhāryam . \\
\noindent \textbf{(P\_6,3.112)} ottvam kriyatām pūrvatvam iti kim atra kartavyam . \\
\noindent \textbf{(P\_6,3.112)} paratvāt ottvam . \\
\noindent \textbf{(P\_6,3.112)} antaraṅgam pūrvatvam . \\
\noindent \textbf{(P\_6,3.112)} evam tarhi idam iha sampradhāryam . \\
\noindent \textbf{(P\_6,3.112)} ottvam kriyatām samprasāraṇam iti kim atra kartavyam . \\
\noindent \textbf{(P\_6,3.112)} paratvāt ottvam . \\
\noindent \textbf{(P\_6,3.112)} nityam samprasāraṇam . \\
\noindent \textbf{(P\_6,3.112)} kṛte api ottve prāpnoti akṛte api . \\
\noindent \textbf{(P\_6,3.112)} ottvam api nityam . \\
\noindent \textbf{(P\_6,3.112)} kṛte api samprasāraṇe prāpnoti akṛte api . \\
\noindent \textbf{(P\_6,3.112)} anityam ottvam . \\
\noindent \textbf{(P\_6,3.112)} na hi kṛte samprasāraṇe prāpnoti . \\
\noindent \textbf{(P\_6,3.112)} antaraṅgam pūrvatvam . \\
\noindent \textbf{(P\_6,3.112)} yasya ca lakṣaṇāntareṇa nimittam vihanyate na tat anityam . \\
\noindent \textbf{(P\_6,3.112)} na ca samprasāraṇam eva ottvasya nimittam vihanti . \\
\noindent \textbf{(P\_6,3.112)} avaśyam lakṣaṇāntaram pūrvatvam pratīkṣyam . \\
\noindent \textbf{(P\_6,3.112)} ubhayoḥ nityayoḥ paratvāt ottvam . \\
\noindent \textbf{(P\_6,3.112)} ottve kṛte samprasāraṇam samprasāraṇapūrvatvam . \\
\noindent \textbf{(P\_6,3.112)} tatra kāryakṛtatvāt punaḥ ottvam na bhaviṣyati . \\
\noindent \textbf{(P\_6,3.121)} apīlvādīnām iti vaktavyam . \\
\noindent \textbf{(P\_6,3.121)} iha mā bhūt . \\
\noindent \textbf{(P\_6,3.121)} rucivaham , cāruvaham . \\
\noindent \textbf{(P\_6,3.122)} amnuṣyādiṣu iti vaktavyam . \\
\noindent \textbf{(P\_6,3.122)} iha mā bhūt . \\
\noindent \textbf{(P\_6,3.122)} prasevaḥ , prahāraḥ , prasāraḥ . \\
\noindent \textbf{(P\_6,3.122)} sādakārayoḥ kṛtrime . \\
\noindent \textbf{(P\_6,3.122)} sādakārayoḥ kṛtrime iti vaktavyam . \\
\noindent \textbf{(P\_6,3.122)} iha eva yathā syāt . \\
\noindent \textbf{(P\_6,3.122)} prāsādaḥ , prākāraḥ . \\
\noindent \textbf{(P\_6,3.122)} iha mā bhūt . \\
\noindent \textbf{(P\_6,3.122)} eṣaḥ asya prasādaḥ . \\
\noindent \textbf{(P\_6,3.122)} eṣaḥ asya prakāraḥ . \\
\noindent \textbf{(P\_6,3.122)} prativeśādīnām vibhāṣā . \\
\noindent \textbf{(P\_6,3.122)} prativeśādīnām vibhāṣā dīrghatvam vaktavyam . \\
\noindent \textbf{(P\_6,3.122)} prativeśaḥ , pratīveśaḥ , pratikāraḥ , pratīkāraḥ . \\
\noindent \textbf{(P\_6,3.124)} katham idam vijñāyate . \\
\noindent \textbf{(P\_6,3.124)} dā iti etasmin takārādau , āhosvit dā iti etasmin takārānte iti . \\
\noindent \textbf{(P\_6,3.124)} kim ca ataḥ . \\
\noindent \textbf{(P\_6,3.124)} yadi vijñāyate takārādau iti nīttā vittā , atra na prāpnoti . \\
\noindent \textbf{(P\_6,3.124)} atha vijñāyate takārānte iti sudattam pratidattam atra api prāpnoti . \\
\noindent \textbf{(P\_6,3.124)} yathā icchasi tathā astu . \\
\noindent \textbf{(P\_6,3.124)} astu tāvat takārādau iti . \\
\noindent \textbf{(P\_6,3.124)} katham nīttā vittā . \\
\noindent \textbf{(P\_6,3.124)} cartve kṛte bhaviṣyati . \\
\noindent \textbf{(P\_6,3.124)} asiddham cartvam . \\
\noindent \textbf{(P\_6,3.124)} tasya asiddhatvāt na prāpnoti . \\
\noindent \textbf{(P\_6,3.124)} āśrayāt siddhatvam bhaviṣyati . \\
\noindent \textbf{(P\_6,3.124)} atha vā punaḥ astu takārānte iti . \\
\noindent \textbf{(P\_6,3.124)} katham sudattam pratidattam . \\
\noindent \textbf{(P\_6,3.124)} na etat takārāntam . \\
\noindent \textbf{(P\_6,3.124)} thakārāntam etat . \\
\noindent \textbf{(P\_6,3.138)} iha anye ācāryāḥ cau pratyaṅgasya pratiṣedham āhuḥ . \\
\noindent \textbf{(P\_6,3.138)} tat iha api sādhyam . \\
\noindent \textbf{(P\_6,3.138)} na eṣaḥ doṣaḥ . \\
\noindent \textbf{(P\_6,3.138)} etat eva jñāpayati ācāryaḥ na pratyaṅgam bhavati iti yat ayam cau dīrghatvam śāsti . \\
\noindent \textbf{(P\_6,3.139)} ikaḥ hrasvāt samprasāraṇadīrghatvam vipratiṣedhena . \\
\noindent \textbf{(P\_6,3.139)} ikaḥ hrasvāt samprasāraṇadīrghatvam bhavati vipratiṣedhena . \\
\noindent \textbf{(P\_6,3.139)} ikaḥ hrasvasya avakāśaḥ . \\
\noindent \textbf{(P\_6,3.139)} grāmaṇikulam , senānikulam . \\
\noindent \textbf{(P\_6,3.139)} samprasāraṇadīrghatvasya avakāśaḥ . \\
\noindent \textbf{(P\_6,3.139)} vibhāṣā hrasvatvam . \\
\noindent \textbf{(P\_6,3.139)} yadā na hrasvatvam saḥ avakāśaḥ . \\
\noindent \textbf{(P\_6,3.139)} hrasvaprasaṅge ubhayam prāpnoti . \\
\noindent \textbf{(P\_6,3.139)} kārīṣagandhīputraḥ , kaumudagandhīputraḥ . \\
\noindent \textbf{(P\_6,3.139)} samprasāraṇadīrghatvam bhavati vipratiṣedhena . \\
\noindent \textbf{(P\_6,3.139)} atha idānīm dīrghatve kṛte punaḥprasaṅgavijñānāt hrasvatvam kasmāt na bhavati . \\
\noindent \textbf{(P\_6,3.139)} sakṛdgatau vipratiṣedhena yat bādhitam tat bādhitam eva iti . \\
\noindent \textbf{(P\_6,4.1.1)} ā kutaḥ ayam adhikāraḥ . \\
\noindent \textbf{(P\_6,4.1.1)} ā saptamādhyāyaparisamāpteḥ aṅgādhikāraḥ . \\
\noindent \textbf{(P\_6,4.1.1)} yadi ā saptamādhyāyaparisamāpteḥ aṅgādhikāraḥ guṇaḥ yaṅlukoḥ iti yaṅluggrahaṇam kartavyam . \\
\noindent \textbf{(P\_6,4.1.1)} prāk abhyāsavikārebhyaḥ punaḥ aṅgādhikāre sati pratyayalakṣaṇena siddham . \\
\noindent \textbf{(P\_6,4.1.1)} astu tarhi prāk abhyāsavikārebhyaḥ aṅgādhikāraḥ . \\
\noindent \textbf{(P\_6,4.1.1)} yadi prāk abhyāsavikārebhyaḥ aṅgādhikāraḥ vavraśca vakārasya samprasāraṇam prāpnoti . \\
\noindent \textbf{(P\_6,4.1.1)} ā saptamādhyāyaparisamāpteḥ punaḥ aṅgādhikāre sati uḥ adatvasya sthānivadbhāvān na samprasāraṇe samprasāraṇam iti pratiṣedhaḥ siddhaḥ bhavati . \\
\noindent \textbf{(P\_6,4.1.1)} saḥ ca idānīm aparihāraḥ bhavati yat tat uktam aṅgānyatvāt ca siddham iti . \\
\noindent \textbf{(P\_6,4.1.1)} astu tarhi ā saptamādhyāyaparisamāpteḥ aṅgādhikāraḥ . \\
\noindent \textbf{(P\_6,4.1.1)} nanu ca uktam guṇaḥ yaṅlukoḥ iti yaṅluggrahaṇam kartavyam iti . \\
\noindent \textbf{(P\_6,4.1.1)} kriyate nyāse eva . \\
\noindent \textbf{(P\_6,4.1.2)} kim punaḥ iyam sthānṣaṣṭhī , aṅgasya sthāne iti . \\
\noindent \textbf{(P\_6,4.1.2)} evam bhavitum arhati . \\
\noindent \textbf{(P\_6,4.1.2)} aṅgasya iti sthānaṣaṣṭhī cet pañcamyantasya ca adhikāraḥ . \\
\noindent \textbf{(P\_6,4.1.2)} aṅgasya iti sthānaṣaṣṭhī cet pañcamyantasya ca adhikāraḥ kartavyaḥ . \\
\noindent \textbf{(P\_6,4.1.2)} aṅgāt iti api vaktavyam . \\
\noindent \textbf{(P\_6,4.1.2)} anucyamāne hi ataḥ bhisaḥ ais bhavati iti ataḥ iti pañcamī aṅgasya iti sthānaṣaṣṭhī . \\
\noindent \textbf{(P\_6,4.1.2)} tatra aśakyam vivibhaktikatvāt ataḥ iti pañcamyā aṅgam viśeṣayitum . \\
\noindent \textbf{(P\_6,4.1.2)} tatra kaḥ doṣaḥ . \\
\noindent \textbf{(P\_6,4.1.2)} akārāt parasya bhismātrasya ais-bhāvaḥ bhavati iti iha api prasajyeta : brāhmaṇabhissā , odanabhissaṭā iti . \\
\noindent \textbf{(P\_6,4.1.2)} avayavaṣaṣṭhyādīnām ca aprasiddhiḥ . \\
\noindent \textbf{(P\_6,4.1.2)} avayavaṣaṣṭhyādayaḥ ca na sidhyanti . \\
\noindent \textbf{(P\_6,4.1.2)} tatra kaḥ doṣaḥ . \\
\noindent \textbf{(P\_6,4.1.2)} śāsaḥ it aṅhaloḥ iti śāseḥ ca antyasya syāt upadhāmātrasya ca . \\
\noindent \textbf{(P\_6,4.1.2)} ūt upadhāyāḥ gohaḥ iti goheḥ ca antyasya syāt upadhāmātrasya ca . \\
\noindent \textbf{(P\_6,4.1.2)} siddham tu parasparam prati aṅgapratyayasaṅjñābhāvāt . \\
\noindent \textbf{(P\_6,4.1.2)} siddham etat . \\
\noindent \textbf{(P\_6,4.1.2)} katham . \\
\noindent \textbf{(P\_6,4.1.2)} parasparam prati aṅgapratyayasaṅjñe bhavataḥ . \\
\noindent \textbf{(P\_6,4.1.2)} aṅgasañjñām prati pratyayasañjñā pratyayasañjñām prati aṅgasañjñā . \\
\noindent \textbf{(P\_6,4.1.2)} kim ataḥ yat parasparam prati aṅgapratyayasaṅjñe bhavataḥ . \\
\noindent \textbf{(P\_6,4.1.2)} sambandhaṣaṣthīnirdeśaḥ ca . \\
\noindent \textbf{(P\_6,4.1.2)} sambandhaṣaṣthīnirdeśaḥ ca ayam kṛtaḥ bhavati . \\
\noindent \textbf{(P\_6,4.1.2)} aṅgasya yaḥ bhis-śabdaḥ iti . \\
\noindent \textbf{(P\_6,4.1.2)} kim ca aṅgasya bhis-śabdaḥ . \\
\noindent \textbf{(P\_6,4.1.2)} nimittam . \\
\noindent \textbf{(P\_6,4.1.2)} yasmin aṅgam iti etat bhavati . \\
\noindent \textbf{(P\_6,4.1.2)} kasmin ca etat bhavati . \\
\noindent \textbf{(P\_6,4.1.2)} pratyaye . \\
\noindent \textbf{(P\_6,4.1.2)} evam api avayavaṣaṣṭhyādayaḥ aviśeṣitāḥ bhavanti . \\
\noindent \textbf{(P\_6,4.1.2)} avayavaṣaṣṭhyādayaḥ api sambandhe eva . \\
\noindent \textbf{(P\_6,4.1.2)} evam api sthānam aviśeṣitam bhavati . \\
\noindent \textbf{(P\_6,4.1.2)} sthānam api sambandhaḥ eva . \\
\noindent \textbf{(P\_6,4.1.2)} evam api na jñāyate kva sthānaṣaṣṭhī kva viśeṣaṇaṣaṣṭhī iti . \\
\noindent \textbf{(P\_6,4.1.2)} yatra ṣaṣṭhī anyayogam na apekṣate sā sthānaṣaṣṭhī . \\
\noindent \textbf{(P\_6,4.1.2)} yatra hi anyayogam apekṣate sā viśeṣaṇaṣaṣṭhī . \\
\noindent \textbf{(P\_6,4.1.3)} kāni punaḥ aṅgādhikārasya prayojanani . \\
\noindent \textbf{(P\_6,4.1.3)} aṅgādhikārasya prayojanam samprasāraṇadīrghatve . \\
\noindent \textbf{(P\_6,4.1.3)} halaḥ uttarasya samprasāraṇasya dīrghaḥ bhavati . \\
\noindent \textbf{(P\_6,4.1.3)} hūtaḥ , jīnaḥ , saṃvītaḥ , śūnaḥ . \\
\noindent \textbf{(P\_6,4.1.3)} aṅgasya iti kimartham . \\
\noindent \textbf{(P\_6,4.1.3)} nirutam , durutam . \\
\noindent \textbf{(P\_6,4.1.3)} nāmsanoḥ ca . \\
\noindent \textbf{(P\_6,4.1.3)} nāmsanoḥ ca dīrghatve prayojanam . \\
\noindent \textbf{(P\_6,4.1.3)} nāmi dīrghaḥ bhavati . \\
\noindent \textbf{(P\_6,4.1.3)} agnīnām , vāyūnām . \\
\noindent \textbf{(P\_6,4.1.3)} aṅgasya iti kimartham . \\
\noindent \textbf{(P\_6,4.1.3)} krimiṇām paśya . \\
\noindent \textbf{(P\_6,4.1.3)} pāmanām paśya . \\
\noindent \textbf{(P\_6,4.1.3)} sani dīrghaḥ bhavati . \\
\noindent \textbf{(P\_6,4.1.3)} cicīṣati , tuṣṭūṣati . \\
\noindent \textbf{(P\_6,4.1.3)} aṅgasya iti kimartham . \\
\noindent \textbf{(P\_6,4.1.3)} dadhi sanoti . \\
\noindent \textbf{(P\_6,4.1.3)} madhu sanoti . \\
\noindent \textbf{(P\_6,4.1.3)} liṅi etve . \\
\noindent \textbf{(P\_6,4.1.3)} liṅi etve prayojanam . \\
\noindent \textbf{(P\_6,4.1.3)} gleyāt , mleyāt . \\
\noindent \textbf{(P\_6,4.1.3)} aṅgasya iti kimartham . \\
\noindent \textbf{(P\_6,4.1.3)} niryāyāt , nirvāyāt . \\
\noindent \textbf{(P\_6,4.1.3)} ataḥ bhisaḥ aistve . \\
\noindent \textbf{(P\_6,4.1.3)} ataḥ bhisaḥ aistve prayojanam . \\
\noindent \textbf{(P\_6,4.1.3)} vṛkṣaiḥ , plakṣaiḥ . \\
\noindent \textbf{(P\_6,4.1.3)} aṅgasya iti kimartham . \\
\noindent \textbf{(P\_6,4.1.3)} brāhmaṇabhissā , odanabhissaṭā . \\
\noindent \textbf{(P\_6,4.1.3)} luṅādiṣu aḍāṭau . \\
\noindent \textbf{(P\_6,4.1.3)} luṅādiṣu aḍāṭau prayojanam . \\
\noindent \textbf{(P\_6,4.1.3)} akārṣīt , aihiṣṭa . \\
\noindent \textbf{(P\_6,4.1.3)} aṅgasya iti kimartham . \\
\noindent \textbf{(P\_6,4.1.3)} prākarot , upaihiṣṭa . \\
\noindent \textbf{(P\_6,4.1.3)} iyaṅuvaṅyuṣmadasmattātaṅāminuḍānemukkehrasvayidīrghabhitatvāni . \\
\noindent \textbf{(P\_6,4.1.3)} iyaṅuvaṅau prayojanam . \\
\noindent \textbf{(P\_6,4.1.3)} śriyau śriyaḥ , bhruvau bhruvaḥ . \\
\noindent \textbf{(P\_6,4.1.3)} aṅgasya iti kimartham . \\
\noindent \textbf{(P\_6,4.1.3)} śryartham , bhrvartham . \\
\noindent \textbf{(P\_6,4.1.3)} yuṣmadasmadoḥ prayojanam . \\
\noindent \textbf{(P\_6,4.1.3)} sāmaḥ ākam . \\
\noindent \textbf{(P\_6,4.1.3)} yuṣmākam asmākam . \\
\noindent \textbf{(P\_6,4.1.3)} aṅgasya iti kimartham . \\
\noindent \textbf{(P\_6,4.1.3)} yuṣmatsāma , asmatsāma . \\
\noindent \textbf{(P\_6,4.1.3)} tātaṅ prayojanam . \\
\noindent \textbf{(P\_6,4.1.3)} jīvatāt bhavān . \\
\noindent \textbf{(P\_6,4.1.3)} aṅgasya iti kimartham . \\
\noindent \textbf{(P\_6,4.1.3)} paca hi tāvat tvam . \\
\noindent \textbf{(P\_6,4.1.3)} jalpa tu tāvat tvam . \\
\noindent \textbf{(P\_6,4.1.3)} āmi nuṭ prayojanam . \\
\noindent \textbf{(P\_6,4.1.3)} kumārīṇam , kiśorīṇām . \\
\noindent \textbf{(P\_6,4.1.3)} aṅgasya iti kimartham . \\
\noindent \textbf{(P\_6,4.1.3)} kumārī , ām iti āha . \\
\noindent \textbf{(P\_6,4.1.3)} kiśorī , ām iti āha . \\
\noindent \textbf{(P\_6,4.1.3)} āne muk prayojanam . \\
\noindent \textbf{(P\_6,4.1.3)} pacamānaḥ , yajamānaḥ . \\
\noindent \textbf{(P\_6,4.1.3)} aṅgasya iti kimartham . \\
\noindent \textbf{(P\_6,4.1.3)} prāṇaḥ . \\
\noindent \textbf{(P\_6,4.1.3)} ke hrasvaḥ prayojanam . \\
\noindent \textbf{(P\_6,4.1.3)} kiśorikā , kumārikā . \\
\noindent \textbf{(P\_6,4.1.3)} aṅgasya iti kimartham . \\
\noindent \textbf{(P\_6,4.1.3)} kumārī kāyati kumārīkaḥ . \\
\noindent \textbf{(P\_6,4.1.3)} yi dīrghaḥ prayojanam . \\
\noindent \textbf{(P\_6,4.1.3)} cīyate , stūyate . \\
\noindent \textbf{(P\_6,4.1.3)} aṅgasya iti kimartham . \\
\noindent \textbf{(P\_6,4.1.3)} dadhiyānam , madhuyānam . \\
\noindent \textbf{(P\_6,4.1.3)} bhi tatvam prayojanam . \\
\noindent \textbf{(P\_6,4.1.3)} adbhiḥ , adbhyaḥ . \\
\noindent \textbf{(P\_6,4.1.3)} aṅgasya iti kimartham . \\
\noindent \textbf{(P\_6,4.1.3)} abbhāraḥ , abbhakṣaḥ . \\
\noindent \textbf{(P\_6,4.1.3)} na etāni santi prayojanāni . \\
\noindent \textbf{(P\_6,4.1.3)} katham . \\
\noindent \textbf{(P\_6,4.1.3)} arthavadgrahaṇapratyayagrahaṇābhyām siddham . \\
\noindent \textbf{(P\_6,4.1.3)} arthavadgrahaṇapratyayagrahaṇābhyām etāni siddhāni . \\
\noindent \textbf{(P\_6,4.1.3)} kva cit arthavadgrahaṇe na anarthakasya iti evam bhaviṣyati kva cit pratyayāpratyayoḥ grahaṇe pratyayasya eva grahaṇam bhavati iti . \\
\noindent \textbf{(P\_6,4.1.3)} atha vā pratyaye iti prakṛtya aṅgakāryam adhyeṣye . \\
\noindent \textbf{(P\_6,4.1.3)} yadi pratyaye iti prakṛtya aṅgakāryam adhīṣe prākarot , upaihiṣṭa , upasargāt pūrvam aḍāṭau prāpnutaḥ . \\
\noindent \textbf{(P\_6,4.1.3)} siddham tu pratyayagrahaṇe yasmāt saḥ vihitaḥ tadādeḥ tadantasya ca grahaṇam . \\
\noindent \textbf{(P\_6,4.1.3)} siddham etat . \\
\noindent \textbf{(P\_6,4.1.3)} katham . \\
\noindent \textbf{(P\_6,4.1.3)} pratyayagrahaṇe yasmāt saḥ vihitaḥ tadādeḥ tadantasya ca grahaṇam bhavati iti evam upasargāt pūrvam aḍāṭau na bhaviṣyataḥ . \\
\noindent \textbf{(P\_6,4.2)} iha kasmāt na bhavati : tṛtīyaḥ . \\
\noindent \textbf{(P\_6,4.2)} aṇprakaraṇāt ṛkārasya aprāptiḥ . \\
\noindent \textbf{(P\_6,4.2)} aṇprakaraṇāt ṛkārasya dīrghatvam na bhaviṣyati . \\
\noindent \textbf{(P\_6,4.2)} aṇaḥ iti vartate . \\
\noindent \textbf{(P\_6,4.2)} kva prakṛtam . \\
\noindent \textbf{(P\_6,4.2)} ḍhralope pūrvasya dīrghaḥ aṇaḥ iti . \\
\noindent \textbf{(P\_6,4.2)} tat vai ikaḥ kāśe iti anena iggrahaṇena vyavacchinnam na śakyam anuvartayitum . \\
\noindent \textbf{(P\_6,4.2)} iggrahaṇasya ca aṇviśeṣaṇatvāt . \\
\noindent \textbf{(P\_6,4.2)} aṇviśeṣaṇam iggrahaṇam . \\
\noindent \textbf{(P\_6,4.2)} aṇaḥ ikaḥ iti . \\
\noindent \textbf{(P\_6,4.2)} yadi tarhi aṇviśeṣaṇam iggrahaṇam cau dīrghaḥ bhavati iti iha na prāpnoti : avācā , avāce . \\
\noindent \textbf{(P\_6,4.2)} na eṣaḥ doṣaḥ . \\
\noindent \textbf{(P\_6,4.2)} aṇgrahaṇam anuvartate iggrahaṇam nivṛttam . \\
\noindent \textbf{(P\_6,4.2)} evam api kartṝcā kartṝce , atra na prāpnoti . \\
\noindent \textbf{(P\_6,4.2)} yathālakṣaṇam aprayukte . \\
\noindent \textbf{(P\_6,4.2)} atha vā ubhayam nivṛttam . \\
\noindent \textbf{(P\_6,4.2)} kasmāt na bhavati tṛtīyaḥ . \\
\noindent \textbf{(P\_6,4.2)} nipātanāt . \\
\noindent \textbf{(P\_6,4.2)} kim nipātanam . \\
\noindent \textbf{(P\_6,4.2)} dvitīyatṛtīyacaturthaturyāṇi anyatarasyām iti . \\
\noindent \textbf{(P\_6,4.3)} kimartham āmaḥ sanakārasya grahaṇam kriyate na āmi dīrghaḥ iti eva ucyeta . \\
\noindent \textbf{(P\_6,4.3)} kena idānim sanakārake bhaviṣyati . \\
\noindent \textbf{(P\_6,4.3)} nuṭ ayam āmbhaktaḥ āmgrahaṇena grāhiṣyate . \\
\noindent \textbf{(P\_6,4.3)} ataḥ uttaram paṭhati : nāmi dīrghaḥ āmi cet syāt kṛte dīrghe na nuṭ bhavet . \\
\noindent \textbf{(P\_6,4.3)} nāmi dīrghaḥ āmi cet syāt kṛte dīrghatve na nuṭ syāt . \\
\noindent \textbf{(P\_6,4.3)} idam iha sampradhāryam . \\
\noindent \textbf{(P\_6,4.3)} dīrghatvam kriyatām nuṭ iti kim atra kartavyam . \\
\noindent \textbf{(P\_6,4.3)} paratvāt nuṭ . \\
\noindent \textbf{(P\_6,4.3)} nityam dīrghatvam . \\
\noindent \textbf{(P\_6,4.3)} kṛte api nuṭi prāpnoti akṛte api . \\
\noindent \textbf{(P\_6,4.3)} nityatvāt dīrghatve kṛte hrasvāśrayaḥ nuṭ na prāpnoti . \\
\noindent \textbf{(P\_6,4.3)} evam tarhi āha ayam hrasvāntāt nuṭ iti na ca hrasvāntaḥ asti . \\
\noindent \textbf{(P\_6,4.3)} tatra vacanāt bhaviṣyati . \\
\noindent \textbf{(P\_6,4.3)} vacanāt yatra tat na asti . \\
\noindent \textbf{(P\_6,4.3)} na idam vacanāt labhyam . \\
\noindent \textbf{(P\_6,4.3)} asti anyat etasya vacane prayojanam . \\
\noindent \textbf{(P\_6,4.3)} kim . \\
\noindent \textbf{(P\_6,4.3)} yatra dīrghatvam pratiṣidhyate . \\
\noindent \textbf{(P\_6,4.3)} tisṛṇām , catasṛṇām iti . \\
\noindent \textbf{(P\_6,4.3)} na etat asti prayojanam . \\
\noindent \textbf{(P\_6,4.3)} iha tāvat catasṛṇām iti ṣaṭcaturbhyaḥ ca iti evam bhaviṣyati . \\
\noindent \textbf{(P\_6,4.3)} tisṛṇām iti trigrahaṇam api tatra prakṛtam anuvartate . \\
\noindent \textbf{(P\_6,4.3)} kva prakṛtam . \\
\noindent \textbf{(P\_6,4.3)} treḥ trayaḥ iti . \\
\noindent \textbf{(P\_6,4.3)} idam tarhi tvam nṛṇam nṛpate jāyase śuciḥ . \\
\noindent \textbf{(P\_6,4.3)} na ekam udāharaṇam hrasvagrahaṇam prayojayati . \\
\noindent \textbf{(P\_6,4.3)} tatra vacanāt bhūtapūrvagatiḥ vijñāsyate . \\
\noindent \textbf{(P\_6,4.3)} hrasvāntam yat bhūtapūrvam iti . \\
\noindent \textbf{(P\_6,4.3)} uttarāṛtham tarhi sanakāragrahaṇam kartavyam . \\
\noindent \textbf{(P\_6,4.3)} nopadhāyāḥ ca carmaṇām . \\
\noindent \textbf{(P\_6,4.3)} nopadhāyāḥ nāmi yathā syāt . \\
\noindent \textbf{(P\_6,4.3)} iha mā bhūt : carmaṇām , varmaṇām iti . \\
\noindent \textbf{(P\_6,4.3)} nāmi dīrghaḥ āmi cet syāt kṛte dīrghe na nuṭ bhavet . vacanāt yatra tat na asti . nopadhāyāḥ ca carmaṇām . \\
\noindent \textbf{(P\_6,4.12-13)} KA\_III,182.7-183.18 Ro\_IV,673-677 {1/45}     hanaḥ kvau upadhādīrghatvaprasaṅgaḥ . \\
\noindent \textbf{(P\_6,4.12-13)} KA\_III,182.7-183.18 Ro\_IV,673-677 {2/45}     hanaḥ kvau upadhālakṣaṇam dīrghatvam prāpnoti . \\
\noindent \textbf{(P\_6,4.12-13)} KA\_III,182.7-183.18 Ro\_IV,673-677 {3/45}     anunāsikasya kvijhaloḥ kṅiti iti . \\
\noindent \textbf{(P\_6,4.12-13)} KA\_III,182.7-183.18 Ro\_IV,673-677 {4/45}     tasya pratiṣedhaḥ vaktavyaḥ . \\
\noindent \textbf{(P\_6,4.12-13)} KA\_III,182.7-183.18 Ro\_IV,673-677 {5/45}     vṛtrahaṇau vṛtrahaṇaḥ iti . \\
\noindent \textbf{(P\_6,4.12-13)} KA\_III,182.7-183.18 Ro\_IV,673-677 {6/45}     niyamavacanāt siddham . \\
\noindent \textbf{(P\_6,4.12-13)} KA\_III,182.7-183.18 Ro\_IV,673-677 {7/45}     inhanpūṣāryamṇām śau sau ca iti etasmāt niyamavacanāt dīrghatvam na bhaviṣyati . \\
\noindent \textbf{(P\_6,4.12-13)} KA\_III,182.7-183.18 Ro\_IV,673-677 {8/45}     niyamavacanāt siddham iti cet sarvanāmasthānaprakaraṇe niyamavacanāt anyatra aniyamaḥ . \\
\noindent \textbf{(P\_6,4.12-13)} KA\_III,182.7-183.18 Ro\_IV,673-677 {9/45}     niyamavacanāt siddham iti cet sarvanāmasthānaprakaraṇe niyamavacanāt anyatra niyamaḥ na prāpnoti . \\
\noindent \textbf{(P\_6,4.12-13)} KA\_III,182.7-183.18 Ro\_IV,673-677 {10/45}     kva anyatra . \\
\noindent \textbf{(P\_6,4.12-13)} KA\_III,182.7-183.18 Ro\_IV,673-677 {11/45}     vṛtrahaṇi bhrūṇahani . \\
\noindent \textbf{(P\_6,4.12-13)} KA\_III,182.7-183.18 Ro\_IV,673-677 {12/45}     evam tarhi dīrghavidhiḥ yaḥ iha inprabhṛtīnām tam viniyamya suṭi iti suvidvān . \\
\noindent \textbf{(P\_6,4.12-13)} KA\_III,182.7-183.18 Ro\_IV,673-677 {13/45}     dīrghavidhiḥ yaḥ iha inprabhṛtīnām tam sarvanāmasthāne viniyamya , inhanpūṣāryamṇām sarvanāmasthāne dīrghaḥ bhavati . \\
\noindent \textbf{(P\_6,4.12-13)} KA\_III,182.7-183.18 Ro\_IV,673-677 {14/45}     kimartham idam . \\
\noindent \textbf{(P\_6,4.12-13)} KA\_III,182.7-183.18 Ro\_IV,673-677 {15/45}     niyamāṛtham . \\
\noindent \textbf{(P\_6,4.12-13)} KA\_III,182.7-183.18 Ro\_IV,673-677 {16/45}     inhanpūṣāryamṇām sarvanāmasthāne eva na anyatra . \\
\noindent \textbf{(P\_6,4.12-13)} KA\_III,182.7-183.18 Ro\_IV,673-677 {17/45}     śau niyamam punaḥ eva vidadhyāt . \\
\noindent \textbf{(P\_6,4.12-13)} KA\_III,182.7-183.18 Ro\_IV,673-677 {18/45}     tataḥ śau . \\
\noindent \textbf{(P\_6,4.12-13)} KA\_III,182.7-183.18 Ro\_IV,673-677 {19/45}     śau eva sarvanāmasthāne na anyatra . \\
\noindent \textbf{(P\_6,4.12-13)} KA\_III,182.7-183.18 Ro\_IV,673-677 {20/45}     tataḥ sau . \\
\noindent \textbf{(P\_6,4.12-13)} KA\_III,182.7-183.18 Ro\_IV,673-677 {21/45}     sau eva sarvanāmasthāne na anyatra . \\
\noindent \textbf{(P\_6,4.12-13)} KA\_III,182.7-183.18 Ro\_IV,673-677 {22/45}     bhrūṇahani iti tathā asya na duṣyet . \\
\noindent \textbf{(P\_6,4.12-13)} KA\_III,182.7-183.18 Ro\_IV,673-677 {23/45}     tathā asya bhrūṇahani iti na doṣaḥ bhavati . \\
\noindent \textbf{(P\_6,4.12-13)} KA\_III,182.7-183.18 Ro\_IV,673-677 {24/45}     śāsti nivartya suṭi iti aviśeṣe śau niyamam kuru vā api asamīkṣya . \\
\noindent \textbf{(P\_6,4.12-13)} KA\_III,182.7-183.18 Ro\_IV,673-677 {25/45}     atha vā nivṛtte sarvanāmasthānaprakaraṇe aviśeṣeṇa śau niyamam vakṣyāmi . \\
\noindent \textbf{(P\_6,4.12-13)} KA\_III,182.7-183.18 Ro\_IV,673-677 {26/45}     inhanpūṣāryamṇām śau eva . \\
\noindent \textbf{(P\_6,4.12-13)} KA\_III,182.7-183.18 Ro\_IV,673-677 {27/45}     tataḥ sau . \\
\noindent \textbf{(P\_6,4.12-13)} KA\_III,182.7-183.18 Ro\_IV,673-677 {28/45}     sau eva . \\
\noindent \textbf{(P\_6,4.12-13)} KA\_III,182.7-183.18 Ro\_IV,673-677 {29/45}     iha api tarhi niyamāt na prāpnoti : indraḥ vṛtrahāyate . \\
\noindent \textbf{(P\_6,4.12-13)} KA\_III,182.7-183.18 Ro\_IV,673-677 {30/45}     dīrghavidheḥ upadhāniyamāt me hanta yi dīrghavidhau ca na doṣaḥ . \\
\noindent \textbf{(P\_6,4.12-13)} KA\_III,182.7-183.18 Ro\_IV,673-677 {31/45}     upadhālakṣaṇadīrghatvasya niyamaḥ na ca etat upadhālakṣaṇam dīrghatvam . \\
\noindent \textbf{(P\_6,4.12-13)} KA\_III,182.7-183.18 Ro\_IV,673-677 {32/45}     suṭi api vā prakṛte anavakāśaḥ śau niyamaḥ aprakṛtapratiṣedhe . \\
\noindent \textbf{(P\_6,4.12-13)} KA\_III,182.7-183.18 Ro\_IV,673-677 {33/45}     atha vā anuvartamāne sarvanāmasthānagrahaṇe anavakāśaḥ śau niyamaḥ aprakṛtasya api dīrghatvasya niyāmakaḥ bhaviṣyati . \\
\noindent \textbf{(P\_6,4.12-13)} KA\_III,182.7-183.18 Ro\_IV,673-677 {34/45}     katham . \\
\noindent \textbf{(P\_6,4.12-13)} KA\_III,182.7-183.18 Ro\_IV,673-677 {35/45}     yasya hi śau niyamaḥ suṭi na etat tena na tatra bhavet viniyamyam . \\
\noindent \textbf{(P\_6,4.12-13)} KA\_III,182.7-183.18 Ro\_IV,673-677 {36/45}     yasya hi śiḥ sarvanāmasthānam na tasya suṭ . \\
\noindent \textbf{(P\_6,4.12-13)} KA\_III,182.7-183.18 Ro\_IV,673-677 {37/45}     yasya suṭ sarvanāmasthānam na tasya śiḥ . \\
\noindent \textbf{(P\_6,4.12-13)} KA\_III,182.7-183.18 Ro\_IV,673-677 {38/45}     tatra sarvanāmasthānaprakaraṇe niyamyam na asti iti kṛtvā aviśeṣeṇa śau niyamaḥ vijñāsyate . \\
\noindent \textbf{(P\_6,4.12-13)} KA\_III,182.7-183.18 Ro\_IV,673-677 {39/45}     dīrghavidhiḥ yaḥ iha inprabhṛtīnām tam viniyamya suṭi iti suvidvān . \\
\noindent \textbf{(P\_6,4.12-13)} KA\_III,182.7-183.18 Ro\_IV,673-677 {40/45}     śau niyamam punaḥ eva vidadhyāt . \\
\noindent \textbf{(P\_6,4.12-13)} KA\_III,182.7-183.18 Ro\_IV,673-677 {41/45}     bhrūṇahani iti tathā asya na duṣyet . \\
\noindent \textbf{(P\_6,4.12-13)} KA\_III,182.7-183.18 Ro\_IV,673-677 {42/45}     śāsti nivartya suṭi iti aviśeṣe śau niyamam kuru vā api asamīkṣya . \\
\noindent \textbf{(P\_6,4.12-13)} KA\_III,182.7-183.18 Ro\_IV,673-677 {43/45}     dīrghavidheḥ upadhāniyamāt me hanta yi dīrghavidhau ca na doṣaḥ . \\
\noindent \textbf{(P\_6,4.12-13)} KA\_III,182.7-183.18 Ro\_IV,673-677 {44/45}     suṭi api vā prakṛte anavakāśaḥ śau niyamaḥ aprakṛtapratiṣedhe . \\
\noindent \textbf{(P\_6,4.12-13)} KA\_III,182.7-183.18 Ro\_IV,673-677 {45/45}     yasya hi śau niyamaḥ suṭi na etat tena na tatra bhavet viniyamyam . \\
\noindent \textbf{(P\_6,4.14)} atvasantasya dīrghatve pitaḥ upasaṅkhyānam . \\
\noindent \textbf{(P\_6,4.14)} atvasantasya dīrghatve pitaḥ upasaṅkhyānam kartavyam . \\
\noindent \textbf{(P\_6,4.14)} gomān , yavamān . \\
\noindent \textbf{(P\_6,4.14)} kim punaḥ kāraṇam na sidhyati . \\
\noindent \textbf{(P\_6,4.14)} ananubandhakagrahaṇe hi sānubandhakasya grahaṇam na iti evam pitaḥ na prāpnoti . \\
\noindent \textbf{(P\_6,4.14)} ananubandhakagrahaṇe iti ucyate . \\
\noindent \textbf{(P\_6,4.14)} sānubandhakasya idam grahaṇam . \\
\noindent \textbf{(P\_6,4.14)} evam tarhi tadanubandhakagrahaṇe atadanubandhakasya grahaṇam na iti evam pitaḥ na prāpnoti . \\
\noindent \textbf{(P\_6,4.14)} tat tarhi upasaṅkhyānam kartavyam . \\
\noindent \textbf{(P\_6,4.14)} na kartavyam . \\
\noindent \textbf{(P\_6,4.14)} pakāralope kṛte na atubantam bhavati atvantam eva . \\
\noindent \textbf{(P\_6,4.14)} yathā eva tarhi pakāralope kṛte na atubantam evam ukāralope api kṛte na atvantam . \\
\noindent \textbf{(P\_6,4.14)} nanu ca bhūtapūrvagatyā bhaviṣyati atvantam . \\
\noindent \textbf{(P\_6,4.14)} yathā eva tarhi bhūtapūrvagatyā atvantam evam atubantam api . \\
\noindent \textbf{(P\_6,4.14)} evam tarhi āśrīyamāṇe bhūtapūrvagatiḥ atvantam ca āsrīyate na atubantam . \\
\noindent \textbf{(P\_6,4.14)} na sidhyati . \\
\noindent \textbf{(P\_6,4.14)} iha hi vyākaraṇe sarveṣu eva sānubandhakagrahaṇeṣu rūpam āśrīyate : yatra asya etat rūpam iti . \\
\noindent \textbf{(P\_6,4.14)} rūpanirgrahaḥ ca na antareṇa laukikam prayogam . \\
\noindent \textbf{(P\_6,4.14)} tasmin ca laukike prayoge sānubandhakānām prayogaḥ na asti iti kṛtvā dvitīyaḥ prayogaḥ upāsyate . \\
\noindent \textbf{(P\_6,4.14)} kaḥ asau . \\
\noindent \textbf{(P\_6,4.14)} upadeśaḥ nāma . \\
\noindent \textbf{(P\_6,4.14)} upadeśe ca etat atubantam na atvantam . \\
\noindent \textbf{(P\_6,4.14)} yadi punaḥ atśabdam gṛhītvā dīrghatvam ucyeta . \\
\noindent \textbf{(P\_6,4.14)} na evam śakyam . \\
\noindent \textbf{(P\_6,4.14)} iha api prasajyeta : jagat , janagat . \\
\noindent \textbf{(P\_6,4.14)} arthavadgrahaṇe na anarthakasya iti evam etasya na bhaviṣyati . \\
\noindent \textbf{(P\_6,4.14)} iha api tarhi na prāpnoti : kṛtavān , bhuktavān iti . \\
\noindent \textbf{(P\_6,4.14)} kva tarhi syāt . \\
\noindent \textbf{(P\_6,4.14)} pacan , yajan . \\
\noindent \textbf{(P\_6,4.14)} na vai atra iṣyate . \\
\noindent \textbf{(P\_6,4.14)} aniṣṭam ca prāpnoti iṣtam ca na sidhyati . \\
\noindent \textbf{(P\_6,4.14)} tasmāt upasaṅkhyānam kartavyam . \\
\noindent \textbf{(P\_6,4.16.1)} gameḥ dīrghatve iṅgrahaṇam . gameḥ dīrghatve iṅgrahaṇam kartavyam . \\
\noindent \textbf{(P\_6,4.16.1)} iṅgameḥ iti vaktavyam . \\
\noindent \textbf{(P\_6,4.16.1)} iha mā bhūt : sañjigaṃsate vatsaḥ mātrā iti . \\
\noindent \textbf{(P\_6,4.16.1)} agrahaṇe hi anādeśasya api dīrghaprasaṅgaḥ . \\
\noindent \textbf{(P\_6,4.16.1)} akriyamāṇe hi iṅgrahaṇe anādeśasya api dīrghatvam prasajyeta . \\
\noindent \textbf{(P\_6,4.16.1)} sañjigaṃsate vatsaḥ mātrā iti . \\
\noindent \textbf{(P\_6,4.16.1)} na vā chandasi anādeśasya api dīrghatvadarśanāt iṅgrahaṇānarthakyam . \\
\noindent \textbf{(P\_6,4.16.1)} na vā iṅgrahaṇam kartavyam . \\
\noindent \textbf{(P\_6,4.16.1)} kim kāraṇam . \\
\noindent \textbf{(P\_6,4.16.1)} chandasi anādeśasya api dīrghatvadarśanāt . \\
\noindent \textbf{(P\_6,4.16.1)} chandasi anādeśasya api gameḥ dīrghatvam dṛśyate . \\
\noindent \textbf{(P\_6,4.16.1)} svargam lokam sañjigāṃsat . \\
\noindent \textbf{(P\_6,4.16.1)} chandasi anādeśasya api dīrghatvadarśanāt iṅgrahaṇam anarthakam . \\
\noindent \textbf{(P\_6,4.16.1)} yathā eva tarhi chandasi anādeśasya api gameḥ dīrghatvam bhavati evam bhāṣāyām api prāpnoti . \\
\noindent \textbf{(P\_6,4.16.1)} tasmāt iṅgrahaṇam kartavyam . \\
\noindent \textbf{(P\_6,4.16.1)} na kartavyam . \\
\noindent \textbf{(P\_6,4.16.1)} yogavibhāgaḥ kariṣyate . \\
\noindent \textbf{(P\_6,4.16.1)} acaḥ sani . \\
\noindent \textbf{(P\_6,4.16.1)} ajantānām sani dīrghaḥ bhavati . \\
\noindent \textbf{(P\_6,4.16.1)} tataḥ hanigamyoḥ . \\
\noindent \textbf{(P\_6,4.16.1)} hanigamyoḥ ca sani dīrghaḥ bhavati . \\
\noindent \textbf{(P\_6,4.16.1)} acaḥ iti eva . \\
\noindent \textbf{(P\_6,4.16.1)} acaḥ sthāne yau hanigamī . \\
\noindent \textbf{(P\_6,4.16.2)} atha upadhāgrahaṇam anuvartate uta aho na . \\
\noindent \textbf{(P\_6,4.16.2)} kim ca ataḥ . \\
\noindent \textbf{(P\_6,4.16.2)} sani dīrghe upadhādhikāraḥ cet vyañjanapratiṣedhaḥ . \\
\noindent \textbf{(P\_6,4.16.2)} sani dīrghe upadhādhikāraḥ cet vyañjanapratiṣedhaḥ vaktavyaḥ , cicīṣati tuṣṭūṣati iti evam artham . \\
\noindent \textbf{(P\_6,4.16.2)} evam tarhi nivṛttam . \\
\noindent \textbf{(P\_6,4.16.2)} anadhikāre uktam . \\
\noindent \textbf{(P\_6,4.16.2)} kim uktam . \\
\noindent \textbf{(P\_6,4.16.2)} hanigamidīrgheṣu ajgrahaṇam iti . \\
\noindent \textbf{(P\_6,4.16.2)} na eṣaḥ doṣaḥ . \\
\noindent \textbf{(P\_6,4.16.2)} uktam etat hrasvaḥ dīrghaḥ plutaḥ iti yatra brūyāt acaḥ iti etat tatra upasthitam draṣṭavyam iti . \\
\noindent \textbf{(P\_6,4.19.1)} atha , ūṭ ādiḥ kasmān na bhavati . \\
\noindent \textbf{(P\_6,4.19.1)} ādiḥ ṭit bhavati iti prāpnoti . \\
\noindent \textbf{(P\_6,4.19.1)} kasya punaḥ ādiḥ . \\
\noindent \textbf{(P\_6,4.19.1)} vakārasya . \\
\noindent \textbf{(P\_6,4.19.1)} astu . \\
\noindent \textbf{(P\_6,4.19.1)} vakārakasya kā pratipattiḥ . \\
\noindent \textbf{(P\_6,4.19.1)} lopaḥ vyoḥ vali iti lopaḥ bhaviṣyati . \\
\noindent \textbf{(P\_6,4.19.1)} na evam śakyam . \\
\noindent \textbf{(P\_6,4.19.1)} jvaratvarasrivyavimavām upadhāyāḥ ca iti dvau ūṭau syātām . \\
\noindent \textbf{(P\_6,4.19.1)} evam tarhi na eṣaḥ ṭit . \\
\noindent \textbf{(P\_6,4.19.1)} kaḥ tarhi . \\
\noindent \textbf{(P\_6,4.19.1)} ṭhit . \\
\noindent \textbf{(P\_6,4.19.1)} yadi tarhi ṭhit , dhautaḥ paṭaḥ iti etyedhatyūṭsu iti vṛddhiḥ na prāpnoti . \\
\noindent \textbf{(P\_6,4.19.1)} cartve kṛte bhaviṣyati . \\
\noindent \textbf{(P\_6,4.19.1)} asiddham cartvam . \\
\noindent \textbf{(P\_6,4.19.1)} tasya asiddhatvāt na prāpnoti . \\
\noindent \textbf{(P\_6,4.19.1)} āśrayāt siddhatvam bhaviṣyati . \\
\noindent \textbf{(P\_6,4.19.1)} asati anyasmin āśrayāt siddhatvam syāt asti ca anyaḥ siddhaḥ vāhaḥ uṭ iti . \\
\noindent \textbf{(P\_6,4.19.1)} eṣaḥ api ṭhit kariṣyate . \\
\noindent \textbf{(P\_6,4.19.1)} tatra ubhayoḥ cartve kṛte āśrayāt siddhatvam bhaviṣyati . \\
\noindent \textbf{(P\_6,4.19.2)} atha kṅidgrahaṇam anuvartate uta aho na . \\
\noindent \textbf{(P\_6,4.19.2)} kim ca ataḥ . \\
\noindent \textbf{(P\_6,4.19.2)} śūṭtve kṅidadhikāraḥ cet chaḥ ṣatvam . \\
\noindent \textbf{(P\_6,4.19.2)} śūṭtve kṅidadhikāraḥ cet chaḥ ṣatvam vaktavyam . \\
\noindent \textbf{(P\_6,4.19.2)} praṣṭā , praṣṭum , praṣṭavyam . \\
\noindent \textbf{(P\_6,4.19.2)} tukprasaṅgaḥ ca . \\
\noindent \textbf{(P\_6,4.19.2)} tuk ca prāpnoti . \\
\noindent \textbf{(P\_6,4.19.2)} nivṛtte api kṅidgrahaṇe avaśyam atra tugabhāvārthaḥ yatnaḥ kartavyaḥ . \\
\noindent \textbf{(P\_6,4.19.2)} antaraṅgatvāt hi tuk prāpnoti . \\
\noindent \textbf{(P\_6,4.19.2)} cchvoḥ iti sannipātagrahaṇam vijñāyate . \\
\noindent \textbf{(P\_6,4.19.2)} nanu evam api antyasya prāpnoti . \\
\noindent \textbf{(P\_6,4.19.2)} sannipātagrahaṇasāmarthyāt sarvasya bhaviṣyati . \\
\noindent \textbf{(P\_6,4.19.2)} evam api aṅgasya prāpnoti . \\
\noindent \textbf{(P\_6,4.19.2)} nirdiśyamānasya ādeśāḥ bhavanti iti evam aṅgasya na bhaviṣyati . \\
\noindent \textbf{(P\_6,4.19.2)} yadi evam utpucchayateḥ apratyayaḥ utpuṭ iti prāpnoti , utput iti ca iṣyate . \\
\noindent \textbf{(P\_6,4.19.2)} tathā vāñchateḥ apratyayaḥ vān , vāṃśau vāṃśaḥ iti na sidhyati . \\
\noindent \textbf{(P\_6,4.19.2)} yathālakṣaṇam aprayukte . \\
\noindent \textbf{(P\_6,4.19.2)} tatra tu etāvān viśeṣaḥ . \\
\noindent \textbf{(P\_6,4.19.2)} anuvartamāne kṅidgrahaṇe chaḥ ṣatvam vaktavyam tatra ca api sannipātagrahaṇam vijñeyam . \\
\noindent \textbf{(P\_6,4.19.2)} nivṛtte divaḥ ūḍbhāvaḥ . \\
\noindent \textbf{(P\_6,4.19.2)} nivṛtte divaḥ ūḍbhāvaḥ prāpnoti . \\
\noindent \textbf{(P\_6,4.19.2)} dyubhyām , dyubhiḥ . \\
\noindent \textbf{(P\_6,4.19.2)} astu . \\
\noindent \textbf{(P\_6,4.19.2)} katham dyubhyām , dyubhiḥ iti . \\
\noindent \textbf{(P\_6,4.19.2)} ūṭhi kṛte divaḥ ut iti uttvam bhaviṣyati . \\
\noindent \textbf{(P\_6,4.19.2)} na sidhyati . \\
\noindent \textbf{(P\_6,4.19.2)} āntaryataḥ dīrghasya dīrghaḥ prāpnoti . \\
\noindent \textbf{(P\_6,4.19.2)} tadartham taparaḥ kṛtaḥ . \\
\noindent \textbf{(P\_6,4.19.2)} evamartham taparaḥ kriyate . \\
\noindent \textbf{(P\_6,4.19.2)} kva punaḥ kṅidgrahaṇam prakṛtam . \\
\noindent \textbf{(P\_6,4.19.2)} anunāsikasya kvijhaloḥ kṅiti iti . \\
\noindent \textbf{(P\_6,4.19.2)} yadi tat anuvartate ajjhanagamām sani kvijhaloḥ ca iti kvijhaloḥ api dīrghatvam prāpnoti . \\
\noindent \textbf{(P\_6,4.19.2)} jhali tāvat na doṣaḥ . \\
\noindent \textbf{(P\_6,4.19.2)} sanam jhalgrahaṇena viśeṣayiṣyāmaḥ . \\
\noindent \textbf{(P\_6,4.19.2)} sani jhalādau iti . \\
\noindent \textbf{(P\_6,4.19.2)} kvau api ācāryapravṛttiḥ jñāpayati na anena kvau dīrghatvam bhavati iti yat ayam kvibvacipracchyāyatastukaṭaprujuśrīṇām dīrghaḥ asamprasāram ca iti dīrghatvam śāsti . \\
\noindent \textbf{(P\_6,4.22.1)} asiddhavacanam kimartham . \\
\noindent \textbf{(P\_6,4.22.1)} asiddhavacane uktam . \\
\noindent \textbf{(P\_6,4.22.1)} kim uktam . \\
\noindent \textbf{(P\_6,4.22.1)} tatra tāvat uktam ṣatvatukoḥ asiddhavacanam ādeśalakṣaṇapratiṣedhāṛtham utsargalakṣaṇabhāvārtham ca iti . \\
\noindent \textbf{(P\_6,4.22.1)} iha api asiddhavacanam ādeśalakṣaṇapratiṣedhāṛtham utsargalakṣaṇabhāvārtham ca . \\
\noindent \textbf{(P\_6,4.22.1)} ādeśalakṣaṇapratiṣedhāṛtham tāvat . \\
\noindent \textbf{(P\_6,4.22.1)} āgahi jahi gataḥ , gatavān . \\
\noindent \textbf{(P\_6,4.22.1)} anunāsikalope jabhāve ca kṛte ataḥ lopaḥ , ataḥ heḥ iti ca prāpnoti . \\
\noindent \textbf{(P\_6,4.22.1)} asiddhatvāt na bhavati . \\
\noindent \textbf{(P\_6,4.22.1)} utsargalakṣaṇabhāvārtham ca . \\
\noindent \textbf{(P\_6,4.22.1)} edhi śādhi . \\
\noindent \textbf{(P\_6,4.22.1)} astiśāstyoḥ ettvaśābhāvayoḥ kṛtayoḥ jhallakṣaṇam dhitvam na prāpnoti . \\
\noindent \textbf{(P\_6,4.22.1)} asiddhatvāt bhavati . \\
\noindent \textbf{(P\_6,4.22.1)} atha atragrahaṇam kimartham . \\
\noindent \textbf{(P\_6,4.22.1)} atragrahaṇam viṣayārtham . \\
\noindent \textbf{(P\_6,4.22.1)} viṣayaḥ pratinirdiśyate . \\
\noindent \textbf{(P\_6,4.22.1)} atra etasmin ābhācchāstre ābhācchāstram asiddham yathā syāt . \\
\noindent \textbf{(P\_6,4.22.1)} iha mā bhūt : abhāji, rāgaḥ, upabarhaṇam iti . \\
\noindent \textbf{(P\_6,4.22.2)} kāni punaḥ asya yogasya prayojanāni . \\
\noindent \textbf{(P\_6,4.22.2)} prayojanam śaittvam dhitve . \\
\noindent \textbf{(P\_6,4.22.2)} śābhāvaḥ ettvam ca dhitve prayojanam . \\
\noindent \textbf{(P\_6,4.22.2)} edhi śādhi . \\
\noindent \textbf{(P\_6,4.22.2)} astiśāstyoḥ ettvaśābhāvayoḥ kṛtayoḥ jhallakṣaṇam dhitvam na prāpnoti . \\
\noindent \textbf{(P\_6,4.22.2)} asiddhatvāt bhavati . \\
\noindent \textbf{(P\_6,4.22.2)} śābhāvaḥ tāvat na prayojayati . \\
\noindent \textbf{(P\_6,4.22.2)} evam vakṣyāmi . \\
\noindent \textbf{(P\_6,4.22.2)} śās hau śā hau iti . \\
\noindent \textbf{(P\_6,4.22.2)} yatvabhūtaḥ sakāraḥ . \\
\noindent \textbf{(P\_6,4.22.2)} tatra sāt dhitvam dhi ca iti sakārasya lopaḥ . \\
\noindent \textbf{(P\_6,4.22.2)} atha vā , ā hau iti vakṣyāmi . \\
\noindent \textbf{(P\_6,4.22.2)} evam api sakārasya prāpnoti . \\
\noindent \textbf{(P\_6,4.22.2)} upadhāyāḥ iti vartate . \\
\noindent \textbf{(P\_6,4.22.2)} upadhāyāḥ ātve kṛte sāt dhitvam dhi ca iti sakārasya lopaḥ . \\
\noindent \textbf{(P\_6,4.22.2)} atha vā na hau iti vakṣyāmi . \\
\noindent \textbf{(P\_6,4.22.2)} tatra ettve pratiṣiddhe sāt dhitvam dhi ca iti sakārasya lopaḥ . \\
\noindent \textbf{(P\_6,4.22.2)} ettvam api lopāpavādaḥ vijñāsyate na ca sakārasya lopaḥ prāpnoti . \\
\noindent \textbf{(P\_6,4.22.2)} hilopaḥ uttve . \\
\noindent \textbf{(P\_6,4.22.2)} hilopaḥ uttve prayojanam . \\
\noindent \textbf{(P\_6,4.22.2)} kuru iti atra hilope kṛte sārvadhātukapare ukāre iti uttvam na prāpnoti . \\
\noindent \textbf{(P\_6,4.22.2)} asiddhatvāt bhavati . \\
\noindent \textbf{(P\_6,4.22.2)} etat api na asti prayojanam . \\
\noindent \textbf{(P\_6,4.22.2)} vakṣyati tatra sārvadhātukagrahaṇasya prayojanam sārvadhātuke bhūtapūrvamātre yathā styāt uttvam . \\
\noindent \textbf{(P\_6,4.22.2)} tāstilopeṇyaṇādeśāḥ aḍāḍvidhau . \\
\noindent \textbf{(P\_6,4.22.2)} talopaḥ astilopaḥ iṇaḥ ca yaṇādeśaḥ aḍāḍvidhau prayojanam . \\
\noindent \textbf{(P\_6,4.22.2)} akāri , aihī iti . \\
\noindent \textbf{(P\_6,4.22.2)} talope kṛte luṅi iti aḍāṭau na klprāpnutaḥ . \\
\noindent \textbf{(P\_6,4.22.2)} asiddhatvāt bhavataḥ . \\
\noindent \textbf{(P\_6,4.22.2)} astilopaḥ iṇaḥ ca yaṇādeśaḥ prayojanam . \\
\noindent \textbf{(P\_6,4.22.2)} āsan , āyan iti . \\
\noindent \textbf{(P\_6,4.22.2)} iṇastyoḥ yaṇlopayoḥ kṛtayoḥ anajāditvāt āṭ na prāpnoti . \\
\noindent \textbf{(P\_6,4.22.2)} asiddhatvāt bhavati . \\
\noindent \textbf{(P\_6,4.22.2)} astilopaḥ tāvat na prayojayati . \\
\noindent \textbf{(P\_6,4.22.2)} ācāryapravṛttiḥ jñāpayati lopāt āṭ balīyān iti yat ayam śnasoḥ allopaḥ iti taparakaraṇam karoti . \\
\noindent \textbf{(P\_6,4.22.2)} iṇyaṇādeśaḥ ca api na prayojayati . \\
\noindent \textbf{(P\_6,4.22.2)} yaṇādeśe yogavibhāgaḥ kariṣyate . \\
\noindent \textbf{(P\_6,4.22.2)} iṇaḥ yaṇ bhavati . \\
\noindent \textbf{(P\_6,4.22.2)} tataḥ eḥ anekācaḥ . \\
\noindent \textbf{(P\_6,4.22.2)} eḥ ca anekācaḥ iṇaḥ yaṇ bhavati . \\
\noindent \textbf{(P\_6,4.22.2)} tataḥ asaṃyogapūrvasya yaṇ bhavati . \\
\noindent \textbf{(P\_6,4.22.2)} eḥ anekācaḥ iti eva . \\
\noindent \textbf{(P\_6,4.22.2)} sarveṣām eva parihāraḥ . \\
\noindent \textbf{(P\_6,4.22.2)} upadeśaḥ iti vartate . \\
\noindent \textbf{(P\_6,4.22.2)} tatra upadeśāvasthāyām eva aḍāṭau bhavataḥ . \\
\noindent \textbf{(P\_6,4.22.2)} atha vā ārdhadhātuke iti vartate . \\
\noindent \textbf{(P\_6,4.22.2)} atha vā luṅlaṅlṛṅkṣu aṭ iti dvilakārakaḥ nirdeśaḥ : luṅādiṣu lakārādiṣu iti . \\
\noindent \textbf{(P\_6,4.22.2)} sarvathā , aijyata , aupyata iti na sidhyati . \\
\noindent \textbf{(P\_6,4.22.2)} vakṣyati etat ajādīnām aṭā siddham iti . \\
\noindent \textbf{(P\_6,4.22.2)} anunāsikalopaḥ hilopāllopayoḥ jabhāvaḥ ca . \\
\noindent \textbf{(P\_6,4.22.2)} anunāsikalopaḥ hilopāllopayoḥ jabhāvaḥ ca prayojanam . \\
\noindent \textbf{(P\_6,4.22.2)} āgahi jahi gataḥ , gatavān . \\
\noindent \textbf{(P\_6,4.22.2)} anunāsikalope kṛte jabhāve ca ataḥ heḥ ataḥ lopaḥ iti ca lopaḥ prāpnoti . \\
\noindent \textbf{(P\_6,4.22.2)} asiddhatvāt na bhavati . \\
\noindent \textbf{(P\_6,4.22.2)} anunāsikalopaḥ tāvat na prayojayati . \\
\noindent \textbf{(P\_6,4.22.2)} allope upadeśe iti vartate . \\
\noindent \textbf{(P\_6,4.22.2)} yadi upadeśe iti vartate dhinutaḥ , kṛṇutaḥ atra na prāpnoti . \\
\noindent \textbf{(P\_6,4.22.2)} na eṣaḥ doṣaḥ . \\
\noindent \textbf{(P\_6,4.22.2)} na upadeśagrahaṇena prakṛtiḥ abhisambadhyate . \\
\noindent \textbf{(P\_6,4.22.2)} kim tarhi . \\
\noindent \textbf{(P\_6,4.22.2)} ārdhadhātukam abhisambadhyate : ārdhadhātukopadeśe yat akārāntam iti . \\
\noindent \textbf{(P\_6,4.22.2)} jabhāvaḥ ca na prayojayati . \\
\noindent \textbf{(P\_6,4.22.2)} hilope yogavibhāgaḥ kariṣyate . \\
\noindent \textbf{(P\_6,4.22.2)} ataḥ heḥ . \\
\noindent \textbf{(P\_6,4.22.2)} tataḥ utaḥ ca . \\
\noindent \textbf{(P\_6,4.22.2)} utaḥ ca heḥ luk bhavati iti . \\
\noindent \textbf{(P\_6,4.22.2)} tataḥ pratyayāt . \\
\noindent \textbf{(P\_6,4.22.2)} pratyayāt iti ubhayoḥ śeṣaḥ . \\
\noindent \textbf{(P\_6,4.22.2)} atha kimartham anunāsikalopaḥ hilopāllopayoḥ jabhāvaḥ ca iti ucyate na anunāsikalopajabhāvau allopahilopayoḥ iti eva ucyate . \\
\noindent \textbf{(P\_6,4.22.2)} saṅkhyātānudeśaḥ mā bhūt iti . \\
\noindent \textbf{(P\_6,4.22.2)} anunāsikalopaḥ hilope prayojayati . \\
\noindent \textbf{(P\_6,4.22.2)} maṇḍūki tābhiḥ āgahi . \\
\noindent \textbf{(P\_6,4.22.2)} rohitaḥ ca iha a gahi . \\
\noindent \textbf{(P\_6,4.22.2)} marudbhiḥ agne agahi . \\
\noindent \textbf{(P\_6,4.22.2)} samprasāraṇam avarṇalope . \\
\noindent \textbf{(P\_6,4.22.2)} samprasāraṇam avarṇalope prayojanam . \\
\noindent \textbf{(P\_6,4.22.2)} madhonaḥ paśya . \\
\noindent \textbf{(P\_6,4.22.2)} maghonā , maghone . \\
\noindent \textbf{(P\_6,4.22.2)} samprasāraṇe kṛte yasya iti lopaḥ prāpnoti . \\
\noindent \textbf{(P\_6,4.22.2)} asiddhatvāt na bhavati . \\
\noindent \textbf{(P\_6,4.22.2)} na etat asti prayojanam . \\
\noindent \textbf{(P\_6,4.22.2)} vakṣyati etat : maghavanśabdaḥ avyutpannam prātipadikam iti . \\
\noindent \textbf{(P\_6,4.22.2)} rebhāvaḥ āllope . \\
\noindent \textbf{(P\_6,4.22.2)} rebhāvaḥ āllope prayojanam . \\
\noindent \textbf{(P\_6,4.22.2)} kim svit garbham prathamam dadhre āpaḥ . \\
\noindent \textbf{(P\_6,4.22.2)} rebhāve kṛte ātaḥ lopaḥ iṭi ca iti ākāralopaḥ na prāpnoti . \\
\noindent \textbf{(P\_6,4.22.2)} asiddhatvāt bhavati . \\
\noindent \textbf{(P\_6,4.22.2)} etat api na asti prayojanam . \\
\noindent \textbf{(P\_6,4.22.2)} chāndasaḥ rebhāvaḥ liṭ ca chandasi sārvadhātukam api bhavati . \\
\noindent \textbf{(P\_6,4.22.2)} tatra sārvadhātukam apit ṅit bhavati iti ṅitvam . \\
\noindent \textbf{(P\_6,4.22.2)} śnābhyastayoḥ ātaḥ iti ākāralopaḥ bhavati . \\
\noindent \textbf{(P\_6,4.22.3)} yadi tarhi ayam yogaḥ na ārabhyate , ut tu kṛñaḥ katham oḥ vinivṛttau . iha kurvaḥ kurmaḥ kuryāt iti ukāralope kṛte sārvadhātukapare ukāre iti uttvam na prāpnoti . \\
\noindent \textbf{(P\_6,4.22.3)} ṇeḥ api ca iṭi katham vinivṛttiḥ . \\
\noindent \textbf{(P\_6,4.22.3)} iha ca kārayateḥ kāriṣyate ṇeḥ aniṭi iti ṇilopaḥ na prāpnoti . \\
\noindent \textbf{(P\_6,4.22.3)} abruvataḥ tava yogam imam syāt luk ca ciṇaḥ nu katham na tarasya . \\
\noindent \textbf{(P\_6,4.22.3)} iha ca , akāritarām ahāritarām iti ciṇaḥ uttarasya tarasya luk na syāt . \\
\noindent \textbf{(P\_6,4.22.3)} cam bhagavān kṛtavān tu tadartham tena bhavet iṭi ṇeḥ nivṛttiḥ . \\
\noindent \textbf{(P\_6,4.22.3)} iha syasicsīyuṭtāsiṣu bhāvakarmaṇoḥ upadeśe ajjhanagrahadṛśām vā ciṇvat iṭ ca kim ca . \\
\noindent \textbf{(P\_6,4.22.3)} ṇilopaḥ ca . \\
\noindent \textbf{(P\_6,4.22.3)} mvoḥ api ye ca tathā api anuvṛttau . \\
\noindent \textbf{(P\_6,4.22.3)} iha api kurvaḥ kurmaḥ kuryāt iti mvoḥ ye ca iti etat api anuvartiṣyate . \\
\noindent \textbf{(P\_6,4.22.3)} ciṇluki ca kṅitaḥ eva luk syāt . \\
\noindent \textbf{(P\_6,4.22.3)} ciṇluki api prakṛtam kṅidgrahaṇam anuvartate . \\
\noindent \textbf{(P\_6,4.22.3)} kva prakṛtam . \\
\noindent \textbf{(P\_6,4.22.3)} gamahanakhanaghasām lopaḥ kṅiti anaṅi iti . \\
\noindent \textbf{(P\_6,4.22.3)} tat vai saptamīnirdiṣṭam ṣaṣṭhīnirdiṣṭena ca iha arthaḥ . \\
\noindent \textbf{(P\_6,4.22.3)} ciṇaḥ iti eṣā pañcamī kṅiti iti saptamyāḥ ṣaṣṭhīm lprakalpayiṣyati tasmāt iti uttarasya iti . \\
\noindent \textbf{(P\_6,4.22.3)} ut tu kṛñaḥ katham oḥ vinivṛttau . \\
\noindent \textbf{(P\_6,4.22.3)} ṇeḥ api ca iṭi katham vinivṛttiḥ . \\
\noindent \textbf{(P\_6,4.22.3)} abruvataḥ tava yogam imam syāt luk ca ciṇaḥ nu katham na tarasya . \\
\noindent \textbf{(P\_6,4.22.3)} cam bhagavān kṛtavān tu tadartham tena bhavet iṭi ṇeḥ nivṛttiḥ . \\
\noindent \textbf{(P\_6,4.22.3)} mvoḥ api ye ca tathā api anuvṛttau . \\
\noindent \textbf{(P\_6,4.22.3)} ciṇluki ca kṅitaḥ eva luk syāt . \\
\noindent \textbf{(P\_6,4.22.4)} ārabhyamāṇe api etasmin yoge siddham vasusamprasāraṇam ajvidhau . \\
\noindent \textbf{(P\_6,4.22.4)} vasusamprasāraṇam ajvidhau siddham vaktavyam . \\
\noindent \textbf{(P\_6,4.22.4)} kim prayojanam . \\
\noindent \textbf{(P\_6,4.22.4)} papuṣaḥ paśya . \\
\noindent \textbf{(P\_6,4.22.4)} tasthuṣaḥ paśya . \\
\noindent \textbf{(P\_6,4.22.4)} ninyuṣaḥ paśya . \\
\noindent \textbf{(P\_6,4.22.4)} cicyuṣaḥ paśya . \\
\noindent \textbf{(P\_6,4.22.4)} luluvuṣaḥ paśya . \\
\noindent \textbf{(P\_6,4.22.4)} pupuvuṣaḥ paśya iti . \\
\noindent \textbf{(P\_6,4.22.4)} vasoḥ samprasāraṇe kṛte aci iti ākāralopādīni yathā syuḥ iti . \\
\noindent \textbf{(P\_6,4.22.4)} kim punaḥ kāraṇam na sidhyanti . \\
\noindent \textbf{(P\_6,4.22.4)} bahiraṅgalakṣaṇatvāt asiddhatvāt ca . \\
\noindent \textbf{(P\_6,4.22.4)} bahiraṅgalakṣaṇam ca eva hi vasusamprasāraṇam asiddham ca . \\
\noindent \textbf{(P\_6,4.22.4)} āttvam yalopāllopayoḥ paśuṣaḥ na vājān cākhāyitā cākhāyitum . \\
\noindent \textbf{(P\_6,4.22.4)} āttvam yalopāllopayoḥ siddham vaktavyam . \\
\noindent \textbf{(P\_6,4.22.4)} kim prayojanam . \\
\noindent \textbf{(P\_6,4.22.4)} paśuṣaḥ na vājān . \\
\noindent \textbf{(P\_6,4.22.4)} paśuṣaḥ iti ātttvasya asiddhatvāt ātaḥ dhātoḥ iti ākāralopaḥ na prāpnoti . \\
\noindent \textbf{(P\_6,4.22.4)} cākhāyitā cākhāyitum iti āttvasya asiddhatvāt yasya halaḥ iti yalopaḥ prāpnoti . \\
\noindent \textbf{(P\_6,4.22.4)} samānāśrayavacanāt siddham . \\
\noindent \textbf{(P\_6,4.22.4)} samānāśrayam asiddham bhavati vyāśrayam ca etat . \\
\noindent \textbf{(P\_6,4.22.4)} iha tāvat papuṣaḥ paśya , tasthuṣaḥ paśya , ninyuṣaḥ paśya , cicyuṣaḥ paśya , luluvuṣaḥ paśya , pupuvuṣaḥ paśya iti . \\
\noindent \textbf{(P\_6,4.22.4)} vasau ākāralopādīni vasantasya vibhaktau samprasāraṇam . \\
\noindent \textbf{(P\_6,4.22.4)} paśuṣaḥ iti viṭi āttvam viḍantasya vibhaktau ākāralopaḥ . \\
\noindent \textbf{(P\_6,4.22.4)} cākhāyitā cākhāyitum iti yaṅi āttvam yaṅantasya ca ārdhadhātuke lopaḥ iti . \\
\noindent \textbf{(P\_6,4.22.4)} kim vaktavyam etat . \\
\noindent \textbf{(P\_6,4.22.4)} na hi . \\
\noindent \textbf{(P\_6,4.22.4)} katham anucyamānam gaṃsyate . \\
\noindent \textbf{(P\_6,4.22.4)} atragrahaṇasāmarthyāt . \\
\noindent \textbf{(P\_6,4.22.4)} nanu ca anyat atragrahaṇasya prayojanam uktam . \\
\noindent \textbf{(P\_6,4.22.4)} kim uktam . \\
\noindent \textbf{(P\_6,4.22.4)} atragrahaṇam viṣayārtham iti . \\
\noindent \textbf{(P\_6,4.22.4)} adhikārāt api etat siddham . \\
\noindent \textbf{(P\_6,4.22.4)} iha papuṣaḥ , cicyuṣaḥ , luluvuṣaḥ , dvau hetū vypadiṣṭau bahiraṅgalakṣaṇatvam asiddhatvam ca iti . \\
\noindent \textbf{(P\_6,4.22.4)} tatra bhavaet asiddhatvam pratyuktam bahiraṅgalakṣaṇatvam tu na eva pratyuktam . \\
\noindent \textbf{(P\_6,4.22.4)} na eṣaḥ doṣaḥ bahiraṅgam antaraṅgam iti ca pratidvandvibhāvinau etau arthau . \\
\noindent \textbf{(P\_6,4.22.4)} katham . \\
\noindent \textbf{(P\_6,4.22.4)} sati antaraṅge bahiraṅgam sati ca bahiraṅge antaraṅgam . \\
\noindent \textbf{(P\_6,4.22.4)} na ca atra antaraṅgabahiraṅgayoḥ yugapat samavasthānam asti . \\
\noindent \textbf{(P\_6,4.22.4)} na anabhinirvṛtte bahiraṅge antaraṅgam prāpnoti . \\
\noindent \textbf{(P\_6,4.22.4)} tatra nimittam eva bahiraṅgam antaraṅgasya . \\
\noindent \textbf{(P\_6,4.22.4)} hrasvayalopāllopāḥ ca ayādeśe lyapi . \\
\noindent \textbf{(P\_6,4.22.4)} hrasvayalopāllopāḥ ca ayādeśe lyapi siddhāḥ vaktavyāḥ . \\
\noindent \textbf{(P\_6,4.22.4)} praśamayya gataḥ , pratamayya gataḥ . \\
\noindent \textbf{(P\_6,4.22.4)} prabebhidayya gataḥ . \\
\noindent \textbf{(P\_6,4.22.4)} pracecchidayya gataḥ . \\
\noindent \textbf{(P\_6,4.22.4)} prastanayya gataḥ . \\
\noindent \textbf{(P\_6,4.22.4)} pragadayya gataḥ . \\
\noindent \textbf{(P\_6,4.22.4)} hrasvayalopāllopānām asiddhatvāt lyapi laghupūrvāt iti ayādeśaḥ na prāpnoti . \\
\noindent \textbf{(P\_6,4.22.4)} atra api eṣaḥ parihāraḥ samānāśrayavacanāt siddham iti . \\
\noindent \textbf{(P\_6,4.22.4)} katham . \\
\noindent \textbf{(P\_6,4.22.4)} ṇau ete vidhayaḥ ṇeḥ lyapi ayādeśaḥ . \\
\noindent \textbf{(P\_6,4.22.4)} vugyuṭau uvaṅyaṇoḥ . \\
\noindent \textbf{(P\_6,4.22.4)} vugyuṭau uvaṅyaṇoḥ siddhau vaktavyau . \\
\noindent \textbf{(P\_6,4.22.4)} babhūvatuḥ , babhūvuḥ : vukaḥ asiddhatvāt uvaṅādeśaḥ prāpnoti . \\
\noindent \textbf{(P\_6,4.22.4)} upadidīye , upadidīyāte : yuṭaḥ asiddhatvāt yaṇādeśaḥ prāpnoti . \\
\noindent \textbf{(P\_6,4.22.4)} vukaḥ tāvat na vaktavyaḥ . \\
\noindent \textbf{(P\_6,4.22.4)} vukam na vakṣyāmi . \\
\noindent \textbf{(P\_6,4.22.4)} evam vakṣyāmi : bhuvaḥ luṅliṭoḥ ūt upadhāyāḥ iti . \\
\noindent \textbf{(P\_6,4.22.4)} atra uvaṅādeśe kṛte yā upadhā tasyāḥ ūttvam bhaviṣyati . \\
\noindent \textbf{(P\_6,4.22.4)} evam api kutaḥ nu khalu etat uvaṅādeśe kṛte yā upadhā tasyāḥ ūttvam bhaviṣyati na punaḥ sāmpratikī yā upadhā tasyāḥ syāt bhakārasya . \\
\noindent \textbf{(P\_6,4.22.4)} na eṣaḥ doṣaḥ . \\
\noindent \textbf{(P\_6,4.22.4)} oḥ iti vartate . \\
\noindent \textbf{(P\_6,4.22.4)} tena uvarṇasya bhaviṣyati . \\
\noindent \textbf{(P\_6,4.22.4)} bhavet siddham babhūvatuḥ , babhūvuḥ . \\
\noindent \textbf{(P\_6,4.22.4)} idam tu na sidhyati : babhūva babhūvitha iti . \\
\noindent \textbf{(P\_6,4.22.4)} kim kāraṇam . \\
\noindent \textbf{(P\_6,4.22.4)} guṇavṛddhyoḥ kṛtayoḥ uvarṇābhāvāt . \\
\noindent \textbf{(P\_6,4.22.4)} na atra guṇavṛddhī prāpnutaḥ . \\
\noindent \textbf{(P\_6,4.22.4)} kim kāraṇam . \\
\noindent \textbf{(P\_6,4.22.4)} kṅiti ca iti pratiṣedhāt . \\
\noindent \textbf{(P\_6,4.22.4)} katham kittvam . \\
\noindent \textbf{(P\_6,4.22.4)} indhibhavatibhyām ca iti . \\
\noindent \textbf{(P\_6,4.22.4)} tat vai vayam kittvam pratyācakṣmahe vukā . \\
\noindent \textbf{(P\_6,4.22.4)} iha tu kittvena vuk pratyākhyāyate . \\
\noindent \textbf{(P\_6,4.22.4)} kim punaḥ atra nyāyyam . \\
\noindent \textbf{(P\_6,4.22.4)} vugvacanam eva nyāyyam . \\
\noindent \textbf{(P\_6,4.22.4)} sati api hi kittve syātām eva atra guṇavṛddhī . \\
\noindent \textbf{(P\_6,4.22.4)} kim kāraṇam . \\
\noindent \textbf{(P\_6,4.22.4)} iglakṣaṇayoḥ guṇavṛddhyoḥ saḥ pratiṣedhaḥ na ca eṣā iglakṣaṇā vṛddhiḥ . \\
\noindent \textbf{(P\_6,4.22.4)} evam tarhi na arthaḥ vukā na api kittvena . \\
\noindent \textbf{(P\_6,4.22.4)} stām atra guṇavṛddhī . \\
\noindent \textbf{(P\_6,4.22.4)} guṇavṛddhyoḥ kṛtayoḥ avāvoḥ ca kṛtayoḥ yā upadhā tasyāḥ ūttvam bhaviṣyati . \\
\noindent \textbf{(P\_6,4.22.4)} katham . \\
\noindent \textbf{(P\_6,4.22.4)} oḥ iti atra avarṇam api pratinirdiśyate . \\
\noindent \textbf{(P\_6,4.22.4)} iha api tarhi prāpnoti . \\
\noindent \textbf{(P\_6,4.22.4)} kīlālapaḥ paśya . \\
\noindent \textbf{(P\_6,4.22.4)} śubhaṃyaḥ paśya iti . \\
\noindent \textbf{(P\_6,4.22.4)} lopaḥ atra bādhakaḥ bhaviṣyati . \\
\noindent \textbf{(P\_6,4.22.4)} iha tarhi prāpnoti . \\
\noindent \textbf{(P\_6,4.22.4)} kīlālapau kīlālapāḥ iti . \\
\noindent \textbf{(P\_6,4.22.4)} evam tarhi vyoḥ iti vartate . \\
\noindent \textbf{(P\_6,4.22.4)} tena uvarṇam viśeṣayiṣyāmaḥ . \\
\noindent \textbf{(P\_6,4.22.4)} oḥ vyoḥ iti . \\
\noindent \textbf{(P\_6,4.22.4)} iha idānīm oḥ iti anuvartate . \\
\noindent \textbf{(P\_6,4.22.4)} vyoḥ iti nivṛttam . \\
\noindent \textbf{(P\_6,4.22.4)} yuṭaḥ ca api na vaktavyam . \\
\noindent \textbf{(P\_6,4.22.4)} yuḍvacanasāmarthyāt na bhaviṣyati . \\
\noindent \textbf{(P\_6,4.22.4)} asti anyat yuḍvacane prayojanam . \\
\noindent \textbf{(P\_6,4.22.4)} kim . \\
\noindent \textbf{(P\_6,4.22.4)} dvayoḥ yakārayoḥ śravaṇam yathā syāt . \\
\noindent \textbf{(P\_6,4.22.4)} na vyañjanaparasya anekasya ekasya vā yakārasya śravaṇam prati viśeṣaḥ asti . \\
\noindent \textbf{(P\_6,4.22.5)} kim punaḥ prāk bhāt asiddhatvam āhosvit saha tena . \\
\noindent \textbf{(P\_6,4.22.5)} kutaḥ punaḥ ayam sandehaḥ . \\
\noindent \textbf{(P\_6,4.22.5)} āṅā ayam nirdeśaḥ kriyate āṅ ca punaḥ sandeham janayati . \\
\noindent \textbf{(P\_6,4.22.5)} tat yathā : ā pāṭaliputrāt vṛṣṭaḥ devaḥ iti sandehaḥ : kim prāk pāṭaliputrāt saha tena iti . \\
\noindent \textbf{(P\_6,4.22.5)} evam iha api sandehaḥ : prāk bhāt saha tena iti . \\
\noindent \textbf{(P\_6,4.22.5)} kaḥ ca atra viśeṣaḥ . \\
\noindent \textbf{(P\_6,4.22.5)} prāk bhāt iti cet sunāmaghonābhūguṇeṣu upasaṅkhyānam . \\
\noindent \textbf{(P\_6,4.22.5)} prāk bhāt iti cet sunāmaghonābhūguṇeṣu upasaṅkhyānam kartavyam . \\
\noindent \textbf{(P\_6,4.22.5)} śunaḥ paśya . \\
\noindent \textbf{(P\_6,4.22.5)} śunā śune . \\
\noindent \textbf{(P\_6,4.22.5)} samprasāraṇe kṛte allopaḥ anaḥ iti prāpnoti . \\
\noindent \textbf{(P\_6,4.22.5)} yasya punaḥ saha tena asiddhatvam asiddhatvāt tasya na saṃyogāt vamantāt iti pratiṣedhaḥ bhaviṣyati . \\
\noindent \textbf{(P\_6,4.22.5)} yasya api prāk bhāt asiddhatvam tasya api eṣaḥ na doṣaḥ . \\
\noindent \textbf{(P\_6,4.22.5)} katham . \\
\noindent \textbf{(P\_6,4.22.5)} na astri atra viśeṣaḥ allopena vā nivṛttau satyām pūrvatvena vā . \\
\noindent \textbf{(P\_6,4.22.5)} ayam asti viśeṣaḥ . \\
\noindent \textbf{(P\_6,4.22.5)} allopena nivṛttau satyām udāttanivṛttisvaraḥ prasajyeta . \\
\noindent \textbf{(P\_6,4.22.5)} na atra udāttanivṛttisvaraḥ prāpnoti . \\
\noindent \textbf{(P\_6,4.22.5)} kim kāraṇam . \\
\noindent \textbf{(P\_6,4.22.5)} na gośvansāvavarṇa iti pratiṣedhāt . \\
\noindent \textbf{(P\_6,4.22.5)} na eṣaḥ udāttanivṛttisvarasya pratiṣedhaḥ . \\
\noindent \textbf{(P\_6,4.22.5)} kasya tarhi . \\
\noindent \textbf{(P\_6,4.22.5)} tṛtīyādisvarasya . \\
\noindent \textbf{(P\_6,4.22.5)} yatra tarhi tṛtīyādisvaraḥ na asti . \\
\noindent \textbf{(P\_6,4.22.5)} śunaḥ paśya iti . \\
\noindent \textbf{(P\_6,4.22.5)} evam tarhi na vayam lakṣaṇasya pratiṣedham śiṣmaḥ . \\
\noindent \textbf{(P\_6,4.22.5)} kim tarhi yena kena cit lakṣaṇena prāptasya vibhaktisvarasya ayam pratiṣedhaḥ . \\
\noindent \textbf{(P\_6,4.22.5)} yatra tarhi vibhaktisvaraḥ na asti . \\
\noindent \textbf{(P\_6,4.22.5)} bahuśunī iti . \\
\noindent \textbf{(P\_6,4.22.5)} yadi punaḥ ayam udāttanivṛttisvarasya api pratiṣedhaḥ vijñāyeta . \\
\noindent \textbf{(P\_6,4.22.5)} na evam śakyam . \\
\noindent \textbf{(P\_6,4.22.5)} iha api prasjyeta kumārī iti . \\
\noindent \textbf{(P\_6,4.22.5)} evam tarhi ācāryapravṛttiḥ jñāpayati na udāttanivṛttisvaraḥ śuni avatarati iti yat ayam śvanśabdam gaurādiṣu paṭhati . \\
\noindent \textbf{(P\_6,4.22.5)} antodāttārtham yatnam karoti . \\
\noindent \textbf{(P\_6,4.22.5)} siddham hi syāt ṅīpā eva . \\
\noindent \textbf{(P\_6,4.22.5)} maghonaḥ paśya . \\
\noindent \textbf{(P\_6,4.22.5)} maghonā maghone . \\
\noindent \textbf{(P\_6,4.22.5)} samprasāraṇe kṛte yasya iti lopaḥ prāpnoti . \\
\noindent \textbf{(P\_6,4.22.5)} yasya punaḥ saha tena asiddhatvam asiddhatvāt tasya na bhaviṣyati . \\
\noindent \textbf{(P\_6,4.22.5)} yasya api prāk bhāt asiddhatvam tasya api eṣaḥ na doṣaḥ . \\
\noindent \textbf{(P\_6,4.22.5)} katham . \\
\noindent \textbf{(P\_6,4.22.5)} vakṣyati etat maghavan-śabdaḥ avyutpannam prātipadikam iti . \\
\noindent \textbf{(P\_6,4.22.5)} bhūguṇaḥ . \\
\noindent \textbf{(P\_6,4.22.5)} bhūyān . \\
\noindent \textbf{(P\_6,4.22.5)} bhūbhāve kṛte oḥ guṇaḥ prāpnoti . \\
\noindent \textbf{(P\_6,4.22.5)} yasya punaḥ saha tena asiddhatvam asiddhatvāt tasya na bhaviṣyati . \\
\noindent \textbf{(P\_6,4.22.5)} yasya api prāk bhāt asiddhatvam tasya api eṣaḥ na doṣaḥ . \\
\noindent \textbf{(P\_6,4.22.5)} katham . \\
\noindent \textbf{(P\_6,4.22.5)} dīrghoccāraṇasāmarthyāt na bhaviṣyati . \\
\noindent \textbf{(P\_6,4.22.5)} asti dīrghoccāraṇasya prayojanam . \\
\noindent \textbf{(P\_6,4.22.5)} kim . \\
\noindent \textbf{(P\_6,4.22.5)} bhūmā iti . \\
\noindent \textbf{(P\_6,4.22.5)} nipātanāt etat siddham . \\
\noindent \textbf{(P\_6,4.22.5)} kim nipātanam . \\
\noindent \textbf{(P\_6,4.22.5)} bahoḥ nañvat uttarapadabhūmni iti . \\
\noindent \textbf{(P\_6,4.22.5)} atha vā punaḥ astu saha tena iti . \\
\noindent \textbf{(P\_6,4.22.5)} ā bhāt iti cet susamprasāraṇayalopaprasthādīnām pratiṣedhaḥ . \\
\noindent \textbf{(P\_6,4.22.5)} papuṣaḥ paśya . \\
\noindent \textbf{(P\_6,4.22.5)} tasthuṣaḥ , ninyuṣaḥ , cicyuṣaḥ , luluvuṣaḥ , pupuvuṣaḥ iti . \\
\noindent \textbf{(P\_6,4.22.5)} vasoḥ samprasāraṇe kṛte aci iti ākāralopādīni na sidhyanti . \\
\noindent \textbf{(P\_6,4.22.5)} na eṣaḥ doṣaḥ . \\
\noindent \textbf{(P\_6,4.22.5)} uktam etat samānāśrayavacanāt siddham iti . \\
\noindent \textbf{(P\_6,4.22.5)} katham . \\
\noindent \textbf{(P\_6,4.22.5)} vasau ākāralopādīni vasantasya vibhaktau samprasāraṇam . \\
\noindent \textbf{(P\_6,4.22.5)} yalopaḥ . \\
\noindent \textbf{(P\_6,4.22.5)} saurī balākā . \\
\noindent \textbf{(P\_6,4.22.5)} yaḥ asau aṇi akāraḥ lupyate tasya asiddhatvāt īti yalopaḥ na prāpnoti . \\
\noindent \textbf{(P\_6,4.22.5)} atra api eṣaḥ eva parihāraḥ . \\
\noindent \textbf{(P\_6,4.22.5)} samānāśrayavacanāt siddham iti . \\
\noindent \textbf{(P\_6,4.22.5)} katham . \\
\noindent \textbf{(P\_6,4.22.5)} aṇi akāralopaḥ aṇantasya īti lopaḥ . \\
\noindent \textbf{(P\_6,4.22.5)} prasthādiṣu . \\
\noindent \textbf{(P\_6,4.22.5)} preyān , stheyān . \\
\noindent \textbf{(P\_6,4.22.5)} prasthādīnām asiddhatvāt prakṛtyā ekāc iti prakṛtibhāvaḥ na prāpnoti . \\
\noindent \textbf{(P\_6,4.22.5)} na eṣaḥ doṣaḥ . \\
\noindent \textbf{(P\_6,4.22.5)} yathā eva prasthādīnām asiddhatvāt prakṛtibhāvaḥ na prāpnoti evam ṭilopaḥ api na bhaviṣyati . \\
\noindent \textbf{(P\_6,4.23)} atha kimartham śnamaḥ saśakārasya grahaṇam kriyate na nāt nalopaḥ iti eva ucyeta . \\
\noindent \textbf{(P\_6,4.23)} nāt nalopaḥ iti iyati ucyamāne nanditā nandakaḥ iti atra api prasajyeta . \\
\noindent \textbf{(P\_6,4.23)} evam tarhi evam vakṣyāmi nāt nalopaḥ aniditām . \\
\noindent \textbf{(P\_6,4.23)} tataḥ halaḥ upadhāyāḥ kṅiti . \\
\noindent \textbf{(P\_6,4.23)} aniditām iti . \\
\noindent \textbf{(P\_6,4.23)} na evam śakyam . \\
\noindent \textbf{(P\_6,4.23)} iha na syāt : hinasti . \\
\noindent \textbf{(P\_6,4.23)} tasmāt na evam śakyam . \\
\noindent \textbf{(P\_6,4.23)} na cet evam nanditā nandakaḥ iti prāpnoti . \\
\noindent \textbf{(P\_6,4.23)} evam tarhi kṅiti iti vartate . \\
\noindent \textbf{(P\_6,4.23)} evam api hinasti iti atra na prāpnoti . \\
\noindent \textbf{(P\_6,4.23)} na eṣā parasaptamī . \\
\noindent \textbf{(P\_6,4.23)} kā tarhi . \\
\noindent \textbf{(P\_6,4.23)} satsaptamī . \\
\noindent \textbf{(P\_6,4.23)} kṅiti sati . \\
\noindent \textbf{(P\_6,4.23)} evam tarhi naśabdaḥ eva atra kṅittvena viśeṣyate kṅit cet naśabdaḥ bhavati iti . \\
\noindent \textbf{(P\_6,4.23)} evam api yajñānām , yatnānām iti atra na prāpnoti . \\
\noindent \textbf{(P\_6,4.23)} dīrghatvam atra bādhakam bhaviṣyati . \\
\noindent \textbf{(P\_6,4.23)} idam iha sampradhāryam . \\
\noindent \textbf{(P\_6,4.23)} dīrghatvam kriyatām nalopaḥ iti kim atra kartavyam . \\
\noindent \textbf{(P\_6,4.23)} paratvāt nalopaḥ . \\
\noindent \textbf{(P\_6,4.23)} tasmāt saśakārasya grahaṇam kartavyam . \\
\noindent \textbf{(P\_6,4.23)} atha kriyamāṇe api saśakāragrahaṇe iha kasmāt na bhavati viśnānām , praśnānām iti . \\
\noindent \textbf{(P\_6,4.23)} lakṣaṇapratipadoktayoḥ pratipadoktasya eva iti evam na bhaviṣyati . \\
\noindent \textbf{(P\_6,4.24)} aniditām nalope laṅgikampyoḥ upatapaśarīravikārayoḥ upasaṅkhyānam . \\
\noindent \textbf{(P\_6,4.24)} aniditām nalope laṅgikampyoḥ upatapaśarīravikārayoḥ upasaṅkhyānam kartavyam . \\
\noindent \textbf{(P\_6,4.24)} vilagitaḥ , vikapitaḥ . \\
\noindent \textbf{(P\_6,4.24)} upatapaśarīravikārayoḥ iti kimartham . \\
\noindent \textbf{(P\_6,4.24)} vilaṅgitaḥ , vikampitaḥ . \\
\noindent \textbf{(P\_6,4.24)} bṛheḥ aci aniṭi . \\
\noindent \textbf{(P\_6,4.24)} bṛheḥ aci aniṭi upasaṅkhyānam kartavyam . \\
\noindent \textbf{(P\_6,4.24)} nibarhayati nibarhakaḥ . \\
\noindent \textbf{(P\_6,4.24)} aci iti kimartham . \\
\noindent \textbf{(P\_6,4.24)} nibṛṃhyate . \\
\noindent \textbf{(P\_6,4.24)} aniṭi iti kimartham . \\
\noindent \textbf{(P\_6,4.24)} nibṛṃhitā nibṛṃhitum . \\
\noindent \textbf{(P\_6,4.24)} tat tu upasaṅkhyānam kartavyam . \\
\noindent \textbf{(P\_6,4.24)} na kartavyam . \\
\noindent \textbf{(P\_6,4.24)} bṛhiḥ prakṛtyantaram . \\
\noindent \textbf{(P\_6,4.24)} katham jñāyate . \\
\noindent \textbf{(P\_6,4.24)} aci iti lopaḥ ucyate . \\
\noindent \textbf{(P\_6,4.24)} anajādau api dṛśyate : nibṛhyate . \\
\noindent \textbf{(P\_6,4.24)} aniṭi iti ucyate . \\
\noindent \textbf{(P\_6,4.24)} iṭau api dṛśyate : nibarhitum . \\
\noindent \textbf{(P\_6,4.24)} ajādau iti ucyate . \\
\noindent \textbf{(P\_6,4.24)} ajādau api na dṛśyate : nibṛṃhayati nibṛṃhakaḥ . \\
\noindent \textbf{(P\_6,4.24)} rañjeḥ ṇau mṛgamaraṇe upasaṅkhyānam kartavyam . \\
\noindent \textbf{(P\_6,4.24)} rajayati mṛgān . \\
\noindent \textbf{(P\_6,4.24)} mṛgamaraṇe iti kimartham . \\
\noindent \textbf{(P\_6,4.24)} rañjayati vastrāṇi . \\
\noindent \textbf{(P\_6,4.24)} ghinuṇi ca upasaṅkhyānam kartavyam . \\
\noindent \textbf{(P\_6,4.24)} rāgī . \\
\noindent \textbf{(P\_6,4.24)} ghinuṇi nipātanāt siddham . \\
\noindent \textbf{(P\_6,4.24)} kim nipātanam . \\
\noindent \textbf{(P\_6,4.24)} tyajaraja iti . \\
\noindent \textbf{(P\_6,4.24)} aśakyam dhātunirdeśe nipātanam tantram āśrayitum . \\
\noindent \textbf{(P\_6,4.24)} iha hi doṣaḥ syāt : daśahanaḥ karaṇe : daṃṣṭrā . \\
\noindent \textbf{(P\_6,4.24)} na etat dhātunipātanam . \\
\noindent \textbf{(P\_6,4.24)} kim tarhi . \\
\noindent \textbf{(P\_6,4.24)} pratyayāntasya etat rūpam . \\
\noindent \textbf{(P\_6,4.24)} tasmin ca asya pratyaye lopaḥ bhavati . \\
\noindent \textbf{(P\_6,4.24)} daṃśasañjasvañjām śapi iti . \\
\noindent \textbf{(P\_6,4.24)} rajakarajanarajaḥsu upasaṅkhyānam kartavyam . \\
\noindent \textbf{(P\_6,4.24)} rajakaḥ , rajananam , rajaḥ iti . \\
\noindent \textbf{(P\_6,4.24)} rajakarajanarajaḥsu kittvāt siddham . \\
\noindent \textbf{(P\_6,4.24)} kitaḥ eva ete auṇādikāḥ . \\
\noindent \textbf{(P\_6,4.24)} tat yathā rucakaḥ , bhuvanam , śiraḥ iti . \\
\noindent \textbf{(P\_6,4.34)} śāsaḥ ittve āśāsaḥ kvau . \\
\noindent \textbf{(P\_6,4.34)} śāsaḥ ittve āśāsaḥ kvau upasaṅkhyānam kartavyam . \\
\noindent \textbf{(P\_6,4.34)} āśīḥ iti . \\
\noindent \textbf{(P\_6,4.34)} kim punaḥ idam niyamārtham āhosvit vidhyartham . \\
\noindent \textbf{(P\_6,4.34)} katham ca niyamārtham syāt katham vā vidhyartham . \\
\noindent \textbf{(P\_6,4.34)} yadi tāvat śāsimātrasya grahaṇam tataḥ niyamārtham . \\
\noindent \textbf{(P\_6,4.34)} athi hi yasmāt śāsaḥ aṅ vihitaḥ tasya grahaṇam tataḥ vidhyartham . \\
\noindent \textbf{(P\_6,4.34)} yadi api śāsimātrasya grahaṇam evam api vidhyartham eva . \\
\noindent \textbf{(P\_6,4.34)} katham . \\
\noindent \textbf{(P\_6,4.34)} aṅhaloḥ iti ucyate na ca atra halādim paśyāmaḥ . \\
\noindent \textbf{(P\_6,4.34)} nanu ca kvip eva halādiḥ . \\
\noindent \textbf{(P\_6,4.34)} kvipaḥ lope kṛte halādyabhāvāt na prāpnoti . \\
\noindent \textbf{(P\_6,4.34)} idam iha sampradhāryam . \\
\noindent \textbf{(P\_6,4.34)} kviblopaḥ kriyatām aṅhaloḥ itttvam iti kim atra kartavyam . \\
\noindent \textbf{(P\_6,4.34)} paratvāt aṅhaloḥ itttvam . \\
\noindent \textbf{(P\_6,4.34)} nityaḥ kviblopaḥ . \\
\noindent \textbf{(P\_6,4.34)} kṛte api aṅhaloḥ itttve prāpnoti akṛte api . \\
\noindent \textbf{(P\_6,4.34)} nityatvāt kviblope kṛte halādyabhāvāt na prāpnoti . \\
\noindent \textbf{(P\_6,4.34)} evam tarhi pratyayalakṣaṇena bhaviṣyati . \\
\noindent \textbf{(P\_6,4.34)} varṇāśraye na asti pratyayalakṣaṇam . \\
\noindent \textbf{(P\_6,4.34)} yadi vā kāni cit varṇāśrayāṇi api pratyayalakṣaṇena bhavanti tathā ca idam api bhaviṣyati . \\
\noindent \textbf{(P\_6,4.34)} atha vā evam vakṣyāmi . \\
\noindent \textbf{(P\_6,4.34)} śāsaḥ it aṅhaloḥ . \\
\noindent \textbf{(P\_6,4.34)} tataḥ kvau . \\
\noindent \textbf{(P\_6,4.34)} kvau ca śāsaḥ it bhavati . \\
\noindent \textbf{(P\_6,4.34)} āryaśīḥ , mitraśīḥ . \\
\noindent \textbf{(P\_6,4.34)} tataḥ āṅaḥ . \\
\noindent \textbf{(P\_6,4.34)} āṅpūrvāt ca kvau śāsaḥ it bhavati . \\
\noindent \textbf{(P\_6,4.34)} āśīḥ iti . \\
\noindent \textbf{(P\_6,4.34)} idam idānīm kimartham . \\
\noindent \textbf{(P\_6,4.34)} niyamārtham . \\
\noindent \textbf{(P\_6,4.34)} āṅpūrvāt śāsaḥ kvau eva . \\
\noindent \textbf{(P\_6,4.34)} kva mā bhūt . \\
\noindent \textbf{(P\_6,4.34)} āśāsyate , āśāsyamānaḥ iti . \\
\noindent \textbf{(P\_6,4.34)} tat tarhi vaktavyam . \\
\noindent \textbf{(P\_6,4.34)} na vaktavyam . \\
\noindent \textbf{(P\_6,4.34)} aviśeṣeṇa śāsaḥ it bhavati iti uktvā tataḥ aṅi iti vakṣyāmi . \\
\noindent \textbf{(P\_6,4.34)} tat niyamārtham bhaviṣyati . \\
\noindent \textbf{(P\_6,4.34)} aṅi eva ajādau na anyasmin ajādau iti . \\
\noindent \textbf{(P\_6,4.34)} iha api tarhi niyamāt ittvam prāpnoti . \\
\noindent \textbf{(P\_6,4.34)} āśāsyate , āśāsyamānaḥ iti . \\
\noindent \textbf{(P\_6,4.34)} yasmāt śāseḥ aṅ vihitaḥ tasya grahaṇam na ca etasmāt śāseḥ aṅ vihitaḥ . \\
\noindent \textbf{(P\_6,4.34)} katham āśīḥ iti . \\
\noindent \textbf{(P\_6,4.34)} nipātanāt siddham . \\
\noindent \textbf{(P\_6,4.34)} kim nipātanam . \\
\noindent \textbf{(P\_6,4.34)} kṣiyāśīḥpraiṣeṣu tiṅ ākāṅkṣam iti . \\
\noindent \textbf{(P\_6,4.37)} anudāttopadeśe anunāsikalopaḥ lyapi ca . \\
\noindent \textbf{(P\_6,4.37)} anudāttopadeśe anunāsikalopaḥ lyapi ca iti vaktavyam . \\
\noindent \textbf{(P\_6,4.37)} pramatya pratatya . \\
\noindent \textbf{(P\_6,4.37)} tataḥ vā amaḥ . \\
\noindent \textbf{(P\_6,4.37)} vā amaḥ iti vaktavyam . \\
\noindent \textbf{(P\_6,4.37)} prayatya prayamya praratya praramya praṇatya praṇamya . \\
\noindent \textbf{(P\_6,4.40)} gamādīnām iti vaktavyam . \\
\noindent \textbf{(P\_6,4.40)} iha api yathā syāt . \\
\noindent \textbf{(P\_6,4.40)} parītat sahakaṇṭhikā . \\
\noindent \textbf{(P\_6,4.40)} saṃyat , sanut iti . \\
\noindent \textbf{(P\_6,4.40)} ūṅ ca gamādīnām iti vaktavyam . \\
\noindent \textbf{(P\_6,4.40)} agregūḥ , bhrūḥ . \\
\noindent \textbf{(P\_6,4.42.1)} atha kim ayam samuccayaḥ , sani ca jhalādau ca iti , āhosvit sanviśeṣaṇam jhalgrahaṇam , sani jhalādau iti . \\
\noindent \textbf{(P\_6,4.42.1)} kim ca ataḥ . \\
\noindent \textbf{(P\_6,4.42.1)} yadi samuccayaḥ sani ajhalādau api prāpnoti . \\
\noindent \textbf{(P\_6,4.42.1)} sisaniṣati jijaniṣate cikhaniṣati . \\
\noindent \textbf{(P\_6,4.42.1)} atha sanviśeṣaṇam jhalgrahaṇam jātaḥ , jātavān iti atra na prāpnoti . \\
\noindent \textbf{(P\_6,4.42.1)} yathā icchasi tathā astu . \\
\noindent \textbf{(P\_6,4.42.1)} astu tāvat samuccayaḥ . \\
\noindent \textbf{(P\_6,4.42.1)} nanu ca uktam sani ajhalādau api prāpnoti iti . \\
\noindent \textbf{(P\_6,4.42.1)} na eṣaḥ doṣaḥ . \\
\noindent \textbf{(P\_6,4.42.1)} prakṛtam jhalgrahaṇam anuvartate . \\
\noindent \textbf{(P\_6,4.42.1)} tena sanam viśeṣayiṣyāmaḥ . \\
\noindent \textbf{(P\_6,4.42.1)} sani jhalādau iti . \\
\noindent \textbf{(P\_6,4.42.1)} atha vā punaḥ astu sanviśeṣaṇam . \\
\noindent \textbf{(P\_6,4.42.1)} katham jātaḥ , jātavān iti . \\
\noindent \textbf{(P\_6,4.42.1)} prakṛtam jhali kṅiti iti anuvartate . \\
\noindent \textbf{(P\_6,4.42.1)} yadi evam na arthaḥ jhalgrahaṇena . \\
\noindent \textbf{(P\_6,4.42.1)} yogavibhāgaḥ kariṣyate . \\
\noindent \textbf{(P\_6,4.42.1)} janasanakhanām anunāsikasya ākāraḥ bhavati jhali kṅiti . \\
\noindent \textbf{(P\_6,4.42.1)} tataḥ sani . \\
\noindent \textbf{(P\_6,4.42.1)} sani ca janasanakhanām anunāsikasya ākāraḥ bhavati jhali iti eva . \\
\noindent \textbf{(P\_6,4.42.1)} tasmāt na arthaḥ jhalgrahaṇena . \\
\noindent \textbf{(P\_6,4.42.2)} sanoteḥ anunāsikalopāt āttvam vipratiṣedhena . \\
\noindent \textbf{(P\_6,4.42.2)} sanoteḥ anunāsikalopāt āttvam bhavati vipratiṣedhena . \\
\noindent \textbf{(P\_6,4.42.2)} sanoteḥ anunāsikalopasya avakāśaḥ anye tanotyādayaḥ . \\
\noindent \textbf{(P\_6,4.42.2)} āttvasya avakāśaḥ anye janādayaḥ . \\
\noindent \textbf{(P\_6,4.42.2)} sanoteḥ anunāsikasya ubhayam prāpnoti . \\
\noindent \textbf{(P\_6,4.42.2)} sātaḥ sātavān iti . \\
\noindent \textbf{(P\_6,4.42.2)} āttvam bhavati vipratiṣedhena . \\
\noindent \textbf{(P\_6,4.42.2)} na eṣaḥ yuktaḥ vipratiṣedhaḥ . \\
\noindent \textbf{(P\_6,4.42.2)} na hi sanoteḥ anunāsikalopasya anye tanotyādayaḥ avakāśaḥ . \\
\noindent \textbf{(P\_6,4.42.2)} sanoteḥ yaḥ tanotyādiṣu pāṭhaḥ saḥ anavakāśaḥ . \\
\noindent \textbf{(P\_6,4.42.2)} na khalu api āttvasya anye janādayaḥ avakāśaḥ . \\
\noindent \textbf{(P\_6,4.42.2)} sanoteḥ yat āttve grahaṇam tat anavakāśam . \\
\noindent \textbf{(P\_6,4.42.2)} tasya anavakāśatvāt ayuktaḥ vipratiṣedhaḥ . \\
\noindent \textbf{(P\_6,4.42.2)} evam tarhi tanotyādiṣu pāṭhaḥ tāvat sāvakāśaḥ . \\
\noindent \textbf{(P\_6,4.42.2)} kaḥ avakāśaḥ . \\
\noindent \textbf{(P\_6,4.42.2)} anyāni tanotyādikāryāṇi . \\
\noindent \textbf{(P\_6,4.42.2)} tanādibhyaḥ tathāsoḥ iti . \\
\noindent \textbf{(P\_6,4.42.2)} āttve api grahaṇam sāvakāśam . \\
\noindent \textbf{(P\_6,4.42.2)} kaḥ avakāśaḥ . \\
\noindent \textbf{(P\_6,4.42.2)} sani ca ye vibhāṣā ca . \\
\noindent \textbf{(P\_6,4.42.2)} ubhayoḥ sāvakāśayoḥ yuktaḥ vipratiṣedhaḥ . \\
\noindent \textbf{(P\_6,4.42.2)} evam api ayuktaḥ vipratiṣedhaḥ . \\
\noindent \textbf{(P\_6,4.42.2)} paṭhiṣyati hi ācāryaḥ pūrvatra asiddhe na asti vipratiṣedhaḥ abhāvāt uttarasya iti . \\
\noindent \textbf{(P\_6,4.42.2)} ekasya nāma abhāve vipratiṣedhaḥ na syāt kim punaḥ yatra ubhayam na asti . \\
\noindent \textbf{(P\_6,4.42.2)} na eṣaḥ doṣaḥ . \\
\noindent \textbf{(P\_6,4.42.2)} bhavati iha vipratiṣedhaḥ . \\
\noindent \textbf{(P\_6,4.42.2)} kim vaktavyam etat . \\
\noindent \textbf{(P\_6,4.42.2)} na hi . \\
\noindent \textbf{(P\_6,4.42.2)} katham anucyamām gaṃsyate . \\
\noindent \textbf{(P\_6,4.42.2)} ācāryapravṛttiḥ jñāpayati bhavati iha vipratiṣedhaḥ iti yat ayam ghumāśthāgāpājahātisām hali iti halgrahaṇam karoti . \\
\noindent \textbf{(P\_6,4.42.2)} katham kṛtvā jñāpakam . \\
\noindent \textbf{(P\_6,4.42.2)} halgrahaṇasya etat prayojanam halādau īttvam yathā syāt iha mā bhūt , godaḥ , kambaladaḥ iti . \\
\noindent \textbf{(P\_6,4.42.2)} yadi ca atra vipratiṣedhaḥ na syāt halgrahaṇam anarthakam syāt . \\
\noindent \textbf{(P\_6,4.42.2)} astu atra īttvam . \\
\noindent \textbf{(P\_6,4.42.2)} īttvasya asiddhatvāt lopaḥ bhaviṣyati . \\
\noindent \textbf{(P\_6,4.42.2)} paśyati tu ācāryaḥ bhavati iha vipratiṣedhaḥ . \\
\noindent \textbf{(P\_6,4.42.2)} tataḥ halgrahaṇam karoti . \\
\noindent \textbf{(P\_6,4.42.2)} na etat asti jñāpakam . \\
\noindent \textbf{(P\_6,4.42.2)} vyavasthārtham etat syāt . \\
\noindent \textbf{(P\_6,4.42.2)} halādau īttvam yathā syāt ajādau mā bhūt iti . \\
\noindent \textbf{(P\_6,4.42.2)} kim ca syāt . \\
\noindent \textbf{(P\_6,4.42.2)} iyaṅādeśaḥ prasajyeta . \\
\noindent \textbf{(P\_6,4.42.2)} nanu ca asiddhatvāt eva iyaṅādeśaḥ na bhaviṣyati . \\
\noindent \textbf{(P\_6,4.42.2)} na śakyam īttvam iyaṅādeśe asiddham vijñātum . \\
\noindent \textbf{(P\_6,4.42.2)} iha hi doṣaḥ syāt : dhiyau dhiyaḥ piyau piyaḥ iti . \\
\noindent \textbf{(P\_6,4.42.2)} na etat īttvam . \\
\noindent \textbf{(P\_6,4.42.2)} kim tarhi . \\
\noindent \textbf{(P\_6,4.42.2)} dhyāpyoḥ samprasāraṇam etat . \\
\noindent \textbf{(P\_6,4.42.2)} samānāśrayam khalu api asiddham bhavati vyāśram ca etat . \\
\noindent \textbf{(P\_6,4.42.2)} katham . \\
\noindent \textbf{(P\_6,4.42.2)} kvau īttvam kvibantasya vibhaktau iyaṅādeśaḥ . \\
\noindent \textbf{(P\_6,4.42.2)} vyavasthārtham eva tarhi halgrahaṇam kartavyam . \\
\noindent \textbf{(P\_6,4.42.2)} kutaḥ hi etat īttvasya asiddhatvāt lopaḥ na punaḥ lopasya asiddhatvāt īttvam iti . \\
\noindent \textbf{(P\_6,4.42.2)} tatra cakrakam avyavasthā prasajyeta . \\
\noindent \textbf{(P\_6,4.42.2)} na asti cakrakaprasaṅgaḥ . \\
\noindent \textbf{(P\_6,4.42.2)} na hi avyavasthākāriṇa śāstreṇa bhavitavyam . \\
\noindent \textbf{(P\_6,4.42.2)} śāstrataḥ nāma vyavasthā . \\
\noindent \textbf{(P\_6,4.42.2)} tatra īttvasya asiddhatvāt lopaḥ lopena vyavasthānam bhaviṣyati . \\
\noindent \textbf{(P\_6,4.42.2)} na khalu api tasmin tat eva asiddham bhavati . \\
\noindent \textbf{(P\_6,4.42.2)} vyavasthārtham eva tarhi halgrahaṇam kartavyam . \\
\noindent \textbf{(P\_6,4.42.2)} halādau īttvam yathā syāt ajādau mā bhūt iti . \\
\noindent \textbf{(P\_6,4.42.2)} kutaḥ hi etat īttvasya asiddhatvāt lopaḥ lopena avasthānam bhaviṣyati na punaḥ lopasya asiddhatvāt īttvam īttvena vyavasthānam syāt . \\
\noindent \textbf{(P\_6,4.42.2)} tat eva khalu api tasmin asiddham bhavati . \\
\noindent \textbf{(P\_6,4.42.2)} katham . \\
\noindent \textbf{(P\_6,4.42.2)} paṭhiṣyati hi ācāryaḥ ciṇaḥ luki tagrahaṇānarthakyam saṅghātasya apratyayatvāt talopasya ca asiddhatvāt iti . \\
\noindent \textbf{(P\_6,4.42.2)} ciṇaḥ luk ciṇaḥ luki eva asiddhaḥ bhavati . \\
\noindent \textbf{(P\_6,4.42.2)} evam tarhi yadi vyavasthārtham etat syāt na eva ayam halgrahaṇam kurvīta . \\
\noindent \textbf{(P\_6,4.42.2)} aviśeṣeṇa ayam īttvam uktvā tasya ajādau lopam apavādam vidadhīta . \\
\noindent \textbf{(P\_6,4.42.2)} idam asti . \\
\noindent \textbf{(P\_6,4.42.2)} ātaḥ lopaḥ iṭi ca iti . \\
\noindent \textbf{(P\_6,4.42.2)} tataḥ ghumāśthāgāpājahātisām . \\
\noindent \textbf{(P\_6,4.42.2)} lopaḥ bhavati iṭi ca ajādau kṅiti . \\
\noindent \textbf{(P\_6,4.42.2)} kimartham punaḥ idam . \\
\noindent \textbf{(P\_6,4.42.2)} īttvam vakṣyāmi tadbādhanārtham . \\
\noindent \textbf{(P\_6,4.42.2)} tataḥ īt . \\
\noindent \textbf{(P\_6,4.42.2)} īt ca bhavati ghvādīnām . \\
\noindent \textbf{(P\_6,4.42.2)} tataḥ eḥ liṅi . \\
\noindent \textbf{(P\_6,4.42.2)} vā anyasya saṃyogādeḥ . \\
\noindent \textbf{(P\_6,4.42.2)} na lyapi . \\
\noindent \textbf{(P\_6,4.42.2)} mayateḥ it anyatarasyām . \\
\noindent \textbf{(P\_6,4.42.2)} tataḥ yati . \\
\noindent \textbf{(P\_6,4.42.2)} yati ca īt bhavati . \\
\noindent \textbf{(P\_6,4.42.2)} saḥ ayam evam laghīyasā nyāsena siddhe sati yat halgrahaṇam karoti garīyāṃsam yatnam ārabhate tat jñāpayati ācāryaḥ bhavati iha vipratiṣedhaḥ iti . \\
\noindent \textbf{(P\_6,4.45)} iha anyatarasyāṅgrahaṇam śakyam akartum . \\
\noindent \textbf{(P\_6,4.45)} katham . \\
\noindent \textbf{(P\_6,4.45)} sanaḥ ktici lopaḥ ca ātttvam ca vibhāṣā iti . \\
\noindent \textbf{(P\_6,4.45)} aparaḥ āha : sarvaḥ eva ayam yogaḥ śakyaḥ avaktum . \\
\noindent \textbf{(P\_6,4.45)} katham . \\
\noindent \textbf{(P\_6,4.45)} iha lopaḥ api prakṛtaḥ āttvam api prakṛtam vibhāṣāgrahaṇam api prakṛtam . \\
\noindent \textbf{(P\_6,4.45)} tatra kevalam abhisambandhamātram kartavyam : sanaḥ ktici lopaḥ ca āttvam ca vibhāṣā . \\
\noindent \textbf{(P\_6,4.46)} kāni punaḥ ārdhadhātukādhikārasya prayojanāni . \\
\noindent \textbf{(P\_6,4.46)} ataḥ lopaḥ yalopaḥ ca ṇilopaḥ ca prayojanam āllopaḥ īttvam etvam ca ciṇvadbhāvaḥ ca sīyuṭi . \\
\noindent \textbf{(P\_6,4.46)} ataḥ lopaḥ . \\
\noindent \textbf{(P\_6,4.46)} cikīrṣitā cikīrṣitum . \\
\noindent \textbf{(P\_6,4.46)} āradhadhātuke iti kimartham . \\
\noindent \textbf{(P\_6,4.46)} cikīrṣati . \\
\noindent \textbf{(P\_6,4.46)} na etat asti prayojanam . \\
\noindent \textbf{(P\_6,4.46)} astu atra sanaḥ akāralopaḥ . \\
\noindent \textbf{(P\_6,4.46)} śapaḥ akārasya śravaṇam bhaviṣyati . \\
\noindent \textbf{(P\_6,4.46)} śapaḥ eva tarhi mā bhūt . \\
\noindent \textbf{(P\_6,4.46)} etat api na asti prayojanam . \\
\noindent \textbf{(P\_6,4.46)} ācāryapravṛttiḥ jñāpayati na anena śabakārasya lopaḥ bhavati iti yat ayam adiprabhṛtibhyaḥ śapaḥ lukam śāsti . \\
\noindent \textbf{(P\_6,4.46)} na etat asti jñāpakam . \\
\noindent \textbf{(P\_6,4.46)} kāryāṛtham etat syāt . \\
\noindent \textbf{(P\_6,4.46)} vittaḥ , mṛṣṭaḥ iti . \\
\noindent \textbf{(P\_6,4.46)} yat tarhi ākārāntebhyaḥ lukam śāsti . \\
\noindent \textbf{(P\_6,4.46)} idam tarhi prayojanam . \\
\noindent \textbf{(P\_6,4.46)} vṛkṣasya plakṣasya . \\
\noindent \textbf{(P\_6,4.46)} ataḥ lopaḥ . \\
\noindent \textbf{(P\_6,4.46)} prāpnoti . \\
\noindent \textbf{(P\_6,4.46)} yalopaḥ api prayojanam . \\
\noindent \textbf{(P\_6,4.46)} bebhiditā cecchiditā . \\
\noindent \textbf{(P\_6,4.46)} āradhadhātuke iti kimartham . \\
\noindent \textbf{(P\_6,4.46)} bebhidyate cecchidyate . \\
\noindent \textbf{(P\_6,4.46)} ṇilopaḥ . \\
\noindent \textbf{(P\_6,4.46)} pācyate yājyate . \\
\noindent \textbf{(P\_6,4.46)} āradhadhātuke iti kimartham . \\
\noindent \textbf{(P\_6,4.46)} pācayati yājayati . \\
\noindent \textbf{(P\_6,4.46)} āllopaḥ . \\
\noindent \textbf{(P\_6,4.46)} yayatuḥ yayuḥ . \\
\noindent \textbf{(P\_6,4.46)} āradhadhātuke iti kimartham . \\
\noindent \textbf{(P\_6,4.46)} yānti vānti . \\
\noindent \textbf{(P\_6,4.46)} īttvam . \\
\noindent \textbf{(P\_6,4.46)} dīyate , dhīyate . \\
\noindent \textbf{(P\_6,4.46)} āradhadhātuke iti kimartham . \\
\noindent \textbf{(P\_6,4.46)} adātām adhātām . \\
\noindent \textbf{(P\_6,4.46)} etvam . \\
\noindent \textbf{(P\_6,4.46)} sneyāt , mleyāt . \\
\noindent \textbf{(P\_6,4.46)} āradhadhātuke iti kimartham . \\
\noindent \textbf{(P\_6,4.46)} snāyāt . \\
\noindent \textbf{(P\_6,4.46)} ciṇvadbhāvaḥ ca sīyuṭi . \\
\noindent \textbf{(P\_6,4.46)} ciṇvadbhāve sīyuṭi kim udāharaṇam . \\
\noindent \textbf{(P\_6,4.46)} kāriṣīṣṭa hāriṣīṣṭa . \\
\noindent \textbf{(P\_6,4.46)} āradhadhātuke iti kimartham . \\
\noindent \textbf{(P\_6,4.46)} kriyeta hriyeta . \\
\noindent \textbf{(P\_6,4.46)} na etat udāharaṇam . \\
\noindent \textbf{(P\_6,4.46)} yakā vyavahitatvāt na bhaviṣyati . \\
\noindent \textbf{(P\_6,4.46)} idam tarhi udāharaṇam : prasnuvīta . \\
\noindent \textbf{(P\_6,4.46)} idam ca api udāharaṇam : kriyeta hriyeta . \\
\noindent \textbf{(P\_6,4.46)} nanu ca uktam yakā vyavahitatvāt na bhaviṣyati iti . \\
\noindent \textbf{(P\_6,4.46)} yakaḥ eva tarhi mā bhūt iti . \\
\noindent \textbf{(P\_6,4.46)} kim ca syāt . \\
\noindent \textbf{(P\_6,4.46)} vṛddhiḥ . \\
\noindent \textbf{(P\_6,4.46)} vṛddhau ca kṛtāyām yuk prasajyeta . \\
\noindent \textbf{(P\_6,4.47)} ayam ram rephasya sthāne kasmāt na bhavati . \\
\noindent \textbf{(P\_6,4.47)} mit acaḥ antyāt paraḥ iti anena acām antyāt paraḥ kriyate . \\
\noindent \textbf{(P\_6,4.47)} rephasya tarhi śravaṇam kasmāt na bhavati . \\
\noindent \textbf{(P\_6,4.47)} ṣaṣṭhyuccāraṇasāmarthyāt . \\
\noindent \textbf{(P\_6,4.47)} bhāradvājīyāḥ paṭhanti bhrasjaḥ ropadhayoḥ lopaḥ āgamaḥ ram vidhīyate iti . \\
\noindent \textbf{(P\_6,4.47)} bhrasjādeśāt samprasāraṇam vipratiṣedhena . \\
\noindent \textbf{(P\_6,4.47)} bhrasjādeśāt samprasāraṇam bhavati vipratiṣedhena . \\
\noindent \textbf{(P\_6,4.47)} bhrasjādeśasya avakāśaḥ : bharṣṭā bhraṣṭā . \\
\noindent \textbf{(P\_6,4.47)} samprasāraṇasya avakāśaḥ : bhṛjjati . \\
\noindent \textbf{(P\_6,4.47)} iha ubhayam prāpnoti : bhṛṣṭaḥ , bhṛṣṭavān . \\
\noindent \textbf{(P\_6,4.47)} samprasāraṇam bhavati vipratiṣedhena . \\
\noindent \textbf{(P\_6,4.47)} saḥ tarhi pūrvavipratiṣedhaḥ vaktavyaḥ . \\
\noindent \textbf{(P\_6,4.47)} na vaktavyaḥ . \\
\noindent \textbf{(P\_6,4.47)} raseḥ vā ṛvacanāt siddham . \\
\noindent \textbf{(P\_6,4.47)} rasoḥ vā ṛ bhavati iti vakṣyāmi . \\
\noindent \textbf{(P\_6,4.47)} rasoḥ vā ṛvacane sici vṛddheḥ bhrasjādeśaḥ . \\
\noindent \textbf{(P\_6,4.47)} rasoḥ vā ṛvacane sici vṛddheḥ bhrasjādeśaḥ vaktavyaḥ . \\
\noindent \textbf{(P\_6,4.47)} vṛddhau kṛtāyām idam eva rūpam syāt : abhrākṣīt . \\
\noindent \textbf{(P\_6,4.47)} idam na syāt : abhārkṣīt . \\
\noindent \textbf{(P\_6,4.47)} sarvathā vayam pūrvavipratiṣedhāt na mucyāmahe sūtram ca bhidyate . \\
\noindent \textbf{(P\_6,4.47)} yathānyāsam eva astu . \\
\noindent \textbf{(P\_6,4.47)} nanu ca uktam bhrasjādeśāt samprasāraṇam vipratiṣedhena iti . \\
\noindent \textbf{(P\_6,4.47)} idam iha sampradhāryam . \\
\noindent \textbf{(P\_6,4.47)} bhrasjādeśaḥ kriyatām samprasāraṇam iti kim atra kartavyam . \\
\noindent \textbf{(P\_6,4.47)} paratvāt bhrasjādeśaḥ . \\
\noindent \textbf{(P\_6,4.47)} nityatvāt samprasāraṇam . \\
\noindent \textbf{(P\_6,4.47)} kṛte api bhrasjādeśe prāpnoti akṛte api . \\
\noindent \textbf{(P\_6,4.47)} bhrasjādeśaḥ api nityaḥ . \\
\noindent \textbf{(P\_6,4.47)} kṛte api samprasāraṇe prāpnoti akṛte api prāpnoti . \\
\noindent \textbf{(P\_6,4.47)} katham . \\
\noindent \textbf{(P\_6,4.47)} yaḥ asau ṛkāre rephaḥ tasya ca upadhāyāḥ ca kṛte api prāpnoti . \\
\noindent \textbf{(P\_6,4.47)} anityaḥ bhrasjādeśaḥ . \\
\noindent \textbf{(P\_6,4.47)} na hi kṛte samprasāraṇe prāpnoti . \\
\noindent \textbf{(P\_6,4.47)} kim kāraṇam . \\
\noindent \textbf{(P\_6,4.47)} na hi varṇaikadeśāḥ varṇagrahaṇena gṛhyante . \\
\noindent \textbf{(P\_6,4.47)} atha api gṛhyante evam api anityaḥ . \\
\noindent \textbf{(P\_6,4.47)} katham . \\
\noindent \textbf{(P\_6,4.47)} upadeśaḥ iti vartate . \\
\noindent \textbf{(P\_6,4.47)} tat ca avaśyam upadeśagrahaṇam anuvartyam barībhṛjjyataḥ iti evamartham . \\
\noindent \textbf{(P\_6,4.48)} ṇyallopau iyaṅyaṇguṇavṛddhidīrghatvebhyaḥ pūrvavipratiṣiddham . \\
\noindent \textbf{(P\_6,4.48)} ṇyallopau iyaṅyaṇguṇavṛddhidīrghatvebhyaḥ bhavataḥ pūrvavipratiṣedhena . \\
\noindent \textbf{(P\_6,4.48)} ṇilopasya avakāśaḥ : kāryate hāryate . \\
\noindent \textbf{(P\_6,4.48)} iyaṅādeśasya avakāśaḥ : śriyau śriyaḥ . \\
\noindent \textbf{(P\_6,4.48)} iha ubhayam prāpnoti : āṭiṭat , āśiśat . \\
\noindent \textbf{(P\_6,4.48)} nanu ca atra yaṇādeśena bhavitavyam . \\
\noindent \textbf{(P\_6,4.48)} idam tarhi : atatakṣat , ararakṣat . \\
\noindent \textbf{(P\_6,4.48)} yaṇādeśasya avakāśaḥ : ninyatuḥ , ninyuḥ . \\
\noindent \textbf{(P\_6,4.48)} ṇilopasya saḥ eva . \\
\noindent \textbf{(P\_6,4.48)} iha ubhayam prāpnoti : āṭiṭat , āśiśat . \\
\noindent \textbf{(P\_6,4.48)} vṛddeḥ avakāśaḥ : sakhāyau sakhāyaḥ . \\
\noindent \textbf{(P\_6,4.48)} ṇilopasya saḥ eva . \\
\noindent \textbf{(P\_6,4.48)} iha ubhayam prāpnoti : kārayateḥ kārakaḥ , hārayateḥ hārakaḥ . \\
\noindent \textbf{(P\_6,4.48)} guṇasya avakāśaḥ : cetā stotā . \\
\noindent \textbf{(P\_6,4.48)} ṇilopasya avakāśaḥ : āṭiṭat , āśiśat . \\
\noindent \textbf{(P\_6,4.48)} iha ubhayam prāpnoti : kāraṇā hāraṇā . \\
\noindent \textbf{(P\_6,4.48)} dīrghatvasya avakāśaḥ : cīyate , stūyate . \\
\noindent \textbf{(P\_6,4.48)} ṇilopasya avakāśaḥ : kāraṇā hāraṇā . \\
\noindent \textbf{(P\_6,4.48)} iha ubhayam prāpnoti : kāryate hāryati . \\
\noindent \textbf{(P\_6,4.48)} ṇilopaḥ bhavati vipratiṣedhena . \\
\noindent \textbf{(P\_6,4.48)} saḥ tarhi pūrvavipratiṣedhaḥ vaktavyaḥ . \\
\noindent \textbf{(P\_6,4.48)} na vaktavyaḥ . \\
\noindent \textbf{(P\_6,4.48)} santu atra ete vidhayaḥ . \\
\noindent \textbf{(P\_6,4.48)} eteṣu vidhiṣu kṛteṣu sthānivadbhāvāt ṇigrahaṇena grahaṇāt ṇilopaḥ bhaviṣyati . \\
\noindent \textbf{(P\_6,4.48)} na evam śakyam . \\
\noindent \textbf{(P\_6,4.48)} iyaṅādeśe hi doṣaḥ syāt . \\
\noindent \textbf{(P\_6,4.48)} antyasya lopaḥ prasajyeta . \\
\noindent \textbf{(P\_6,4.48)} allopasya iyaṅyaṇoḥ ca na asti sampradhāraṇā . \\
\noindent \textbf{(P\_6,4.48)} vṛddheḥ avakāśaḥ : priyam ācaṣṭe prāpayati . \\
\noindent \textbf{(P\_6,4.48)} allopasya avakāśaḥ : cikīrṣitā cikīrṣitum . \\
\noindent \textbf{(P\_6,4.48)} iha ubhayam prāpnoti : cikīrṣakaḥ , jihīrṣakaḥ . \\
\noindent \textbf{(P\_6,4.48)} guṇasya allopasya ca na asti sampradhāraṇā . \\
\noindent \textbf{(P\_6,4.48)} dīrghatvasya avakāśaḥ : api kākaḥ śyenāyate . \\
\noindent \textbf{(P\_6,4.48)} allopasya saḥ eva . \\
\noindent \textbf{(P\_6,4.48)} iha ubhayam prāpnoti : cikīrṣyate jihīrṣyate . \\
\noindent \textbf{(P\_6,4.48)} allopaḥ bhavati vipratiṣedhena . \\
\noindent \textbf{(P\_6,4.48)} saḥ tarhi pūrvavipratiṣedhaḥ vaktavyaḥ . \\
\noindent \textbf{(P\_6,4.48)} na vaktavyaḥ . \\
\noindent \textbf{(P\_6,4.48)} iṣṭavācī paraśabdaḥ . \\
\noindent \textbf{(P\_6,4.48)} vipratiṣedhe param yat iṣṭam tat bhavati iti . \\
\noindent \textbf{(P\_6,4.49)} kim idam yalope varṇagrahaṇam āhosvit saṅghātagrahaṇam . \\
\noindent \textbf{(P\_6,4.49)} kaḥ ca atra viśeṣaḥ . \\
\noindent \textbf{(P\_6,4.49)} yalope varṇagrahaṇam cet dhātvantasya pratiṣedhaḥ . \\
\noindent \textbf{(P\_6,4.49)} yalope varṇagrahaṇam cet dhātvantasya pratiṣedhaḥ vaktavyaḥ . \\
\noindent \textbf{(P\_6,4.49)} śucyitā śucyitum . \\
\noindent \textbf{(P\_6,4.49)} asti tarjo saṅghātagrahaṇam . \\
\noindent \textbf{(P\_6,4.49)} yadi saṅghātagrahaṇam antyasya lopaḥ prāpnoti . \\
\noindent \textbf{(P\_6,4.49)} siddhaḥ antyasya pūrveṇa eva . \\
\noindent \textbf{(P\_6,4.49)} tatra ārambhasāmarthyāt sarvasya bhaviṣyati . \\
\noindent \textbf{(P\_6,4.49)} evam api tena atiprasaktam iti kṛtvā niyamaḥ vijñāyeta . \\
\noindent \textbf{(P\_6,4.49)} yasya halaḥ eva na anyataḥ . \\
\noindent \textbf{(P\_6,4.49)} kva mā bhūt . \\
\noindent \textbf{(P\_6,4.49)} lolūyitā popūyitā . \\
\noindent \textbf{(P\_6,4.49)} kaimarthakyāt niyamaḥ bhavati . \\
\noindent \textbf{(P\_6,4.49)} vidheyam na asti iti kṛtvā . \\
\noindent \textbf{(P\_6,4.49)} iha ca asti vidheyam . \\
\noindent \textbf{(P\_6,4.49)} kim . \\
\noindent \textbf{(P\_6,4.49)} antyasya lopaḥ prāptaḥ saḥ sarvasya vidheyaḥ . \\
\noindent \textbf{(P\_6,4.49)} tatra apūrvaḥ vidhiḥ astu niyama astu iti apūrvaḥ eva vidhiḥ bhaviṣyati na niyamaḥ . \\
\noindent \textbf{(P\_6,4.49)} evam api antyasya prāpnoti . \\
\noindent \textbf{(P\_6,4.49)} kim kāraṇam . \\
\noindent \textbf{(P\_6,4.49)} na hi lopaḥ sarvāpahārī . \\
\noindent \textbf{(P\_6,4.49)} nanu ca saṅghātagrahaṇasāmarthyāt sarvasya bhaviṣyati . \\
\noindent \textbf{(P\_6,4.49)} saṅghātagrahaṇam cet kyasya vibhāṣāyām doṣaḥ . \\
\noindent \textbf{(P\_6,4.49)} saṅghātagrahaṇam cet kyasya vibhāṣāyām doṣaḥ bhavati . \\
\noindent \textbf{(P\_6,4.49)} samidhitā samidhyitā . \\
\noindent \textbf{(P\_6,4.49)} yadā lopaḥ tadā sarvasya lopaḥ . \\
\noindent \textbf{(P\_6,4.49)} yadā alopaḥ tadā sarvasya alopaḥ prāpnoti . \\
\noindent \textbf{(P\_6,4.49)} ādeḥ paravacanāt siddham . \\
\noindent \textbf{(P\_6,4.49)} halaḥ iti pañcamī . \\
\noindent \textbf{(P\_6,4.49)} tasmāt iti uttarasya ādeḥ parasya iti yakārasya eva bhaviṣyati . \\
\noindent \textbf{(P\_6,4.49)} atha vā punaḥ astu varṇagrahaṇam . \\
\noindent \textbf{(P\_6,4.49)} nanu ca uktam yalope varṇagrahaṇam cet dhātvantasya pratiṣedhaḥ iti . \\
\noindent \textbf{(P\_6,4.49)} na eṣaḥ doṣaḥ . \\
\noindent \textbf{(P\_6,4.49)} aṅgāt iti hi vartate . \\
\noindent \textbf{(P\_6,4.49)} na vā aṅgāt iti pañcamī asti . \\
\noindent \textbf{(P\_6,4.49)} evam tarhi aṅgasya iti sambandhaṣaṣṭhī vijñāsyate . \\
\noindent \textbf{(P\_6,4.49)} aṅgasya yaḥ yakāraḥ . \\
\noindent \textbf{(P\_6,4.49)} kim ca aṅgasya yakāraḥ . \\
\noindent \textbf{(P\_6,4.49)} nimittam . \\
\noindent \textbf{(P\_6,4.49)} yasmin aṅgam iti etat bhavati . \\
\noindent \textbf{(P\_6,4.49)} kasmin ca etat bhavati . \\
\noindent \textbf{(P\_6,4.49)} pratyaye . \\
\noindent \textbf{(P\_6,4.51)} atha aniṭi iti kimartham . \\
\noindent \textbf{(P\_6,4.51)} kārayitā kārayitum . \\
\noindent \textbf{(P\_6,4.51)} aniṭi iti śakyam avaktum . \\
\noindent \textbf{(P\_6,4.51)} kasmāt na bhavati kārayitā kārayitum . \\
\noindent \textbf{(P\_6,4.51)} niṣṭhāyām seṭi iti etat niyamārtham bhaviṣyati . \\
\noindent \textbf{(P\_6,4.51)} niṣṭhāyām eva seṭi ṇeḥ lopaḥ bhavati na ayatra . \\
\noindent \textbf{(P\_6,4.51)} kva mā bhūt . \\
\noindent \textbf{(P\_6,4.51)} kārayitā kārayitum . \\
\noindent \textbf{(P\_6,4.51)} atha vā upariṣṭāt yogavibhāgaḥ kariṣyate . \\
\noindent \textbf{(P\_6,4.51)} idam asti . \\
\noindent \textbf{(P\_6,4.51)} niṣṭhāyām seṭi . \\
\noindent \textbf{(P\_6,4.51)} janita mantra . \\
\noindent \textbf{(P\_6,4.51)} śamitā yajñe . \\
\noindent \textbf{(P\_6,4.51)} tataḥ ay . \\
\noindent \textbf{(P\_6,4.51)} ayādeśaḥ bhavati ṇeḥ seṭi . \\
\noindent \textbf{(P\_6,4.51)} tata āmantālvāyetnviṣṇuṣu ay bhavati iti eva . \\
\noindent \textbf{(P\_6,4.52.1)} atha seḍgrahaṇam kimartham . \\
\noindent \textbf{(P\_6,4.52.1)} niṣṭhāyām seḍgrahaṇam aniṭi pratiṣedhārtham . \\
\noindent \textbf{(P\_6,4.52.1)} niṣṭhāyām seḍgrahaṇam kriyate aniṭi pratiṣedhaḥ yathā syāt iti . \\
\noindent \textbf{(P\_6,4.52.1)} sañjñapitaḥ paśuḥ iti . \\
\noindent \textbf{(P\_6,4.52.1)} niṣṭhāyām seḍgrahaṇam aniṭi pratiṣedhārtham iti cet tat siddham aniḍabhāvāt . \\
\noindent \textbf{(P\_6,4.52.1)} niṣṭhāyām seḍgrahaṇam aniṭi pratiṣedhārtham iti cet antareṇa api seḍgrahaṇam tat siddham . \\
\noindent \textbf{(P\_6,4.52.1)} katham . \\
\noindent \textbf{(P\_6,4.52.1)} aniḍabhāvāt . \\
\noindent \textbf{(P\_6,4.52.1)} nanu ca yasya vibhāṣā iti jñapeḥ iṭpratiṣedhaḥ . \\
\noindent \textbf{(P\_6,4.52.1)} ekācaḥ hi pratiṣedhaḥ . \\
\noindent \textbf{(P\_6,4.52.1)} ekācaḥ hi saḥ pratiṣedhaḥ jñapiḥ ca anekāc . \\
\noindent \textbf{(P\_6,4.52.1)} iḍbhāvārtham tu tannimittatvāt lopasya . \\
\noindent \textbf{(P\_6,4.52.1)} iḍbhāvārtham tarhi seḍgrahaṇam kriyate . \\
\noindent \textbf{(P\_6,4.52.1)} katham punaḥ seṭi iti anena iṭ śakyaḥ bhāvayitum . \\
\noindent \textbf{(P\_6,4.52.1)} tannimittatvāt lopasya . \\
\noindent \textbf{(P\_6,4.52.1)} na atra akṛte iṭi ṇilopena bhavitavyam . \\
\noindent \textbf{(P\_6,4.52.1)} kim kāraṇam . \\
\noindent \textbf{(P\_6,4.52.1)} seṭi iti ucyate . \\
\noindent \textbf{(P\_6,4.52.1)} avacane hi ṇilope iṭpratiṣedhaprasaṅgaḥ . \\
\noindent \textbf{(P\_6,4.52.1)} akriyamāṇe hi seḍgrahaṇe ṇilope kṛte ekācaḥ iti iṭpratiṣedhaḥ prasajyeta . \\
\noindent \textbf{(P\_6,4.52.1)} kāritam , hāritam . \\
\noindent \textbf{(P\_6,4.52.1)} evam tarhi na arthaḥ seḍgrahaṇena na api sūtreṇa . \\
\noindent \textbf{(P\_6,4.52.1)} katham . \\
\noindent \textbf{(P\_6,4.52.1)} saptame yogavibhāgaḥ kariṣyate . \\
\noindent \textbf{(P\_6,4.52.1)} idam asti . \\
\noindent \textbf{(P\_6,4.52.1)} niṣṭhāyām na iṭ bhavati . \\
\noindent \textbf{(P\_6,4.52.1)} tataḥ ṇeḥ . \\
\noindent \textbf{(P\_6,4.52.1)} ṇyantasya niṣṭhāyām na iṭ bhavati . \\
\noindent \textbf{(P\_6,4.52.1)} kāritam , hāritam . \\
\noindent \textbf{(P\_6,4.52.1)} tataḥ vṛttam . \\
\noindent \textbf{(P\_6,4.52.1)} vṛttam iti ca nipātyate . \\
\noindent \textbf{(P\_6,4.52.1)} kim nipātyate . \\
\noindent \textbf{(P\_6,4.52.1)} ṇeḥ niṣṭhāyām lopaḥ nipātyate . \\
\noindent \textbf{(P\_6,4.52.1)} kim prayojanam . \\
\noindent \textbf{(P\_6,4.52.1)} niyamārtham . \\
\noindent \textbf{(P\_6,4.52.1)} atra eva ṇeḥ niṣṭhāyām lopaḥ bhavati na anyatra . \\
\noindent \textbf{(P\_6,4.52.1)} kva mā bhūt . \\
\noindent \textbf{(P\_6,4.52.1)} kāritam , hāritam . \\
\noindent \textbf{(P\_6,4.52.1)} iha api tarhi prāpnoti : vartitam annam , vartitā bhikṣā iti . \\
\noindent \textbf{(P\_6,4.52.1)} tataḥ adhyayane . \\
\noindent \textbf{(P\_6,4.52.1)} adhyayane cet vṛtiḥ vartate iti . \\
\noindent \textbf{(P\_6,4.52.2)} vṛdhiramiśṛdhīnām upasaṅkhyānam sārvadhātukatvāt . \\
\noindent \textbf{(P\_6,4.52.2)} vṛdhiramiśṛdhīnām upasaṅkhyānam kartavyam . \\
\noindent \textbf{(P\_6,4.52.2)} kim kāraṇam . \\
\noindent \textbf{(P\_6,4.52.2)} sārvadhātukatvāt . \\
\noindent \textbf{(P\_6,4.52.2)} vardhantu tvā suṣṭutayaḥ giraḥ me . \\
\noindent \textbf{(P\_6,4.52.2)} vardhayantu iti evam prāpte . \\
\noindent \textbf{(P\_6,4.52.2)} bṛhaspatiḥ tvā sumne ramṇātu . \\
\noindent \textbf{(P\_6,4.52.2)} ramayatu iti evam prāpte . \\
\noindent \textbf{(P\_6,4.52.2)} agne śardha mahate saubhagāya . \\
\noindent \textbf{(P\_6,4.52.2)} śardhaya iti evam prāpte . \\
\noindent \textbf{(P\_6,4.52.2)} tat tarhi vaktavyam . \\
\noindent \textbf{(P\_6,4.52.2)} na vaktavyam . \\
\noindent \textbf{(P\_6,4.52.2)} vṛdhiramiśṛdhīnām ārdhadhātukatvāt siddham . \\
\noindent \textbf{(P\_6,4.52.2)} katham ārdhadhātukatvam . \\
\noindent \textbf{(P\_6,4.52.2)} anye api hi dhātupratyayāḥ ubhayathā chandasi dṛśyante . \\
\noindent \textbf{(P\_6,4.55)} kim punaḥ ayam ktnuḥ āhosvit itnuḥ . \\
\noindent \textbf{(P\_6,4.55)} kaḥ ca atra viśeṣaḥ . \\
\noindent \textbf{(P\_6,4.55)} ktnau iṭi ṇeḥ guṇavacanam . \\
\noindent \textbf{(P\_6,4.55)} ktnau iṭi ṇeḥ guṇaḥ vaktavtyaḥ . \\
\noindent \textbf{(P\_6,4.55)} gadayitnuḥ , stanayitnuḥ . \\
\noindent \textbf{(P\_6,4.55)} astu tarhi itnuḥ . \\
\noindent \textbf{(P\_6,4.55)} itnau pratyayāntarakaraṇam . \\
\noindent \textbf{(P\_6,4.55)} yadi tarhi itnuḥ pratyayāntaram kartavyam . \\
\noindent \textbf{(P\_6,4.55)} ayādeśe ca upasaṅkhyānam . \\
\noindent \textbf{(P\_6,4.55)} ayādeśe ca upasaṅkhyānam kartavyam . \\
\noindent \textbf{(P\_6,4.55)} ubhayam kriyate nyāse eva . \\
\noindent \textbf{(P\_6,4.56)} lyapi laghupūrvasya iti cet vyañjanānteṣu upasaṅkhyānam . \\
\noindent \textbf{(P\_6,4.56)} lyapi laghupūrvasya iti cet vyañjanānteṣu upasaṅkhyānam kartavyam . \\
\noindent \textbf{(P\_6,4.56)} praśamayya gataḥ . \\
\noindent \textbf{(P\_6,4.56)} pratamayya gataḥ . \\
\noindent \textbf{(P\_6,4.56)} allope ca gurupūrvāt pratiṣedhaḥ . \\
\noindent \textbf{(P\_6,4.56)} allope ca gurupūrvāt pratiṣedhaḥ vaktavyaḥ . \\
\noindent \textbf{(P\_6,4.56)} pracikīrṣya gataḥ . \\
\noindent \textbf{(P\_6,4.56)} lyapi laghupūrvāt iti vacanāt siddham . \\
\noindent \textbf{(P\_6,4.56)} lyapi laghupūrvāt iti vaktavyam . \\
\noindent \textbf{(P\_6,4.56)} evam api hrasvayalopāllopānām asiddhatvāt lyapi laghupūrvāt iti ayādeśaḥ na prāpnoti . \\
\noindent \textbf{(P\_6,4.56)} praśamayya gataḥ . \\
\noindent \textbf{(P\_6,4.56)} pratamayya gataḥ . \\
\noindent \textbf{(P\_6,4.56)} prabebhidayya gataḥ . \\
\noindent \textbf{(P\_6,4.56)} pracecchidayya gataḥ . \\
\noindent \textbf{(P\_6,4.56)} pragadayya gataḥ . \\
\noindent \textbf{(P\_6,4.56)} prastanayya gataḥ . \\
\noindent \textbf{(P\_6,4.56)} hrasvādiṣu ca uktam . \\
\noindent \textbf{(P\_6,4.56)} kim uktam . \\
\noindent \textbf{(P\_6,4.56)} samānāśrayatvāt siddham iti . \\
\noindent \textbf{(P\_6,4.56)} katham . \\
\noindent \textbf{(P\_6,4.56)} ṇau ete vidhayaḥ . \\
\noindent \textbf{(P\_6,4.56)} ṇeḥ lyapi ayādeśaḥ . \\
\noindent \textbf{(P\_6,4.57)} iṅādeśasya pratiṣedhaḥ vaktavyaḥ . \\
\noindent \textbf{(P\_6,4.57)} adhyāpya gataḥ . \\
\noindent \textbf{(P\_6,4.57)} āpaḥ sānubandhakanirdeśāt iṅi siddham . \\
\noindent \textbf{(P\_6,4.57)} āpaḥ sānubandhakanirdeśaḥ kariṣyate . \\
\noindent \textbf{(P\_6,4.57)} tena iṅādeśasya na bhaviṣyati . \\
\noindent \textbf{(P\_6,4.57)} saḥ tarhi sānubandhakanirdeśaḥ kartavyaḥ . \\
\noindent \textbf{(P\_6,4.57)} na kartavyaḥ . \\
\noindent \textbf{(P\_6,4.57)} lakṣaṇapratipadoktayoḥ pratipadoktasya eva iti evam na bhaviṣyati . \\
\noindent \textbf{(P\_6,4.62.1)} bhāvakarmaṇoḥ iti katham idam vijñāyate . \\
\noindent \textbf{(P\_6,4.62.1)} bhāvakarmaṇoḥ ye syādayaḥ iti , āhosvit bhāvakarmavācini parataḥ ye syādayaḥ iti . \\
\noindent \textbf{(P\_6,4.62.1)} kim ca ataḥ . \\
\noindent \textbf{(P\_6,4.62.1)} yadi vijñāyate bhāvakarmaṇoḥ ye syādayaḥ iti sīyuṭ viśeṣitaḥ syasictāsayaḥ aviśeṣitāḥ . \\
\noindent \textbf{(P\_6,4.62.1)} atha vijñāyate bhāvakarmavācini parataḥ ye syādayaḥ iti syasictāsayaḥ viśeṣitāḥ sīyuṭ aviśeṣitaḥ . \\
\noindent \textbf{(P\_6,4.62.1)} yathā icchasi tathā astu . \\
\noindent \textbf{(P\_6,4.62.1)} astu tāvat bhāvakarmaṇoḥ ye syādayaḥ iti . \\
\noindent \textbf{(P\_6,4.62.1)} syasictāsayaḥ ca viśeṣitāḥ . \\
\noindent \textbf{(P\_6,4.62.1)} nanu ca uktam sīyuṭ viśeṣitaḥ syasictāsayaḥ aviśeṣitāḥ iti . \\
\noindent \textbf{(P\_6,4.62.1)} syasictāsayaḥ ca viśeṣitāḥ . \\
\noindent \textbf{(P\_6,4.62.1)} katham . \\
\noindent \textbf{(P\_6,4.62.1)} bhāvakarmaṇoḥ yak bhavati iti atra syādayaḥ api anuvartiṣyante . \\
\noindent \textbf{(P\_6,4.62.1)} atha vā punaḥ astu bhāvakarmavācini parataḥ ye syādayaḥ iti . \\
\noindent \textbf{(P\_6,4.62.1)} nanu ca uktam syasictāsayaḥ viśeṣitāḥ sīyuṭ aviśeṣitaḥ iti . \\
\noindent \textbf{(P\_6,4.62.1)} sīyuṭ ca viśeṣitaḥ . \\
\noindent \textbf{(P\_6,4.62.1)} katham . \\
\noindent \textbf{(P\_6,4.62.1)} bhāvakarmavācini parataḥ sīyuṭ na asti iti kṛtva bhāvakarmavācini sīyuṭi kāryam vijñāsyate . \\
\noindent \textbf{(P\_6,4.62.2)} atha iṭ ca iti ucyate . \\
\noindent \textbf{(P\_6,4.62.2)} kasya ayam iṭ bhavati . \\
\noindent \textbf{(P\_6,4.62.2)} aṅgasya iti vartate . \\
\noindent \textbf{(P\_6,4.62.2)} yadi evam āditaḥ iṭ prāpnoti aḍāḍvat . \\
\noindent \textbf{(P\_6,4.62.2)} tat yathā aḍāṭau ṭittvāt āditaḥ bhavataḥ tadvat . \\
\noindent \textbf{(P\_6,4.62.2)} evam tarhi syādīnām eva bhaviṣyanti . \\
\noindent \textbf{(P\_6,4.62.2)} evam api ṣaṣṭhyabhāvāt na prāpnoti . \\
\noindent \textbf{(P\_6,4.62.2)} nanu ca bhāvakarmaṇoḥ iti eṣā ṣaṣṭhī . \\
\noindent \textbf{(P\_6,4.62.2)} na eṣā ṣaṣṭhī . \\
\noindent \textbf{(P\_6,4.62.2)} kim tarhi arthinirdeśe eṣā saptamī : bhāve ca arthe karmaṇi ca iti . \\
\noindent \textbf{(P\_6,4.62.2)} evam tarhi bhāvakarmaṇoḥ iti eṣā saptamī syādiṣu iti saptamyāḥ ṣaṣṭhīm prakalpayiṣyati tasmin iti nirdiṣṭe pūrvasya iti . \\
\noindent \textbf{(P\_6,4.62.2)} evam api na sidhyati . \\
\noindent \textbf{(P\_6,4.62.2)} kim kāraṇam . \\
\noindent \textbf{(P\_6,4.62.2)} na hi arthena paurvāparyam asti . \\
\noindent \textbf{(P\_6,4.62.2)} arthe asambhavāt tadvācini śabde kāryam vijñāsyate . \\
\noindent \textbf{(P\_6,4.62.2)} evam api sīyuṭaḥ na prāpnoti . \\
\noindent \textbf{(P\_6,4.62.2)} evam tarhi saptame yogavibhāgaḥ kariṣyate . \\
\noindent \textbf{(P\_6,4.62.2)} ārdhadhātukasya iṭ . \\
\noindent \textbf{(P\_6,4.62.2)} yāvān iṭ nāma saḥ sarvaḥ ārdhadhātukasya iṭ bhavati . \\
\noindent \textbf{(P\_6,4.62.2)} tataḥ valādeḥ . \\
\noindent \textbf{(P\_6,4.62.2)} valādeḥ ārdhadhātukasya iṭ bhavati iti . \\
\noindent \textbf{(P\_6,4.62.2)} yadi evam syasicsīyuṭtāsiṣu iṭ bhavati ciṇvadbhāvaḥ aviśeṣitaḥ bhavati . \\
\noindent \textbf{(P\_6,4.62.2)} tatra kaḥ doṣaḥ . \\
\noindent \textbf{(P\_6,4.62.2)} syasicsīyuṭtāsiṣu iṭ bhavati ajjhanagrahadṛśām vā ciṇvat iti kva cit eva ciṇvadbhāvaḥ syāt . \\
\noindent \textbf{(P\_6,4.62.2)} evam tarhi syādīn apekṣiṣyāmahe . \\
\noindent \textbf{(P\_6,4.62.2)} syasicsīyuṭtāsiṣu iṭ bhavati ajjhanagrahadṛśām vā ciṇvat syādiṣu iti . \\
\noindent \textbf{(P\_6,4.62.2)} atha ke punaḥ imam iṭam prayojayanti . \\
\noindent \textbf{(P\_6,4.62.2)} ye anudāttāḥ . \\
\noindent \textbf{(P\_6,4.62.2)} atha ye udāttāḥ teṣām katham . \\
\noindent \textbf{(P\_6,4.62.2)} siddham tena eva paratvāt . \\
\noindent \textbf{(P\_6,4.62.2)} udāttebhyaḥ api vā anena eva iṭ eṣitavyaḥ . \\
\noindent \textbf{(P\_6,4.62.2)} kim prayojanam . \\
\noindent \textbf{(P\_6,4.62.2)} kārayateḥ kāriṣyate , hārayateḥ hāriṣyate . \\
\noindent \textbf{(P\_6,4.62.2)} iṭaḥ asiddhatvāt aniṭi iti ṇilopaḥ yathā syāt . \\
\noindent \textbf{(P\_6,4.62.2)} katham punaḥ icchatā api bhavatā udāttebhyaḥ anena eva iṭ labhyaḥ na punaḥ anena astu tena vā iti tena eva syāt vipratiṣedhena . \\
\noindent \textbf{(P\_6,4.62.2)} nanu ca nityaḥ ayam kṛte api tasmin prāpnoti akṛte api prāpnoti . \\
\noindent \textbf{(P\_6,4.62.2)} na tu asmin kṛte api saḥ prāpnoti . \\
\noindent \textbf{(P\_6,4.62.2)} kim kāraṇam . \\
\noindent \textbf{(P\_6,4.62.2)} avalāditvāt . \\
\noindent \textbf{(P\_6,4.62.2)} tasmāt anena eva bhaviṣyati iṭ . \\
\noindent \textbf{(P\_6,4.62.3)} kāni punaḥ asya yogasya prayojanāni . \\
\noindent \textbf{(P\_6,4.62.3)} vṛddhiḥ ciṇvat yuk ca hanteḥ ca ghatvam dīrghaḥ ca uktaḥ yaḥ mitām vā ciṇi iti . \\
\noindent \textbf{(P\_6,4.62.3)} vṛddhiḥ prayojanam . \\
\noindent \textbf{(P\_6,4.62.3)} ceṣyate cāyiṣyate . \\
\noindent \textbf{(P\_6,4.62.3)} yuk ca prayojanam . \\
\noindent \textbf{(P\_6,4.62.3)} glāsyate , glāyiṣyate . \\
\noindent \textbf{(P\_6,4.62.3)} hanteḥ ca ghatvam prayojanam . \\
\noindent \textbf{(P\_6,4.62.3)} haniṣyate ghāniṣyate . \\
\noindent \textbf{(P\_6,4.62.3)} dīrghaḥ ca uktaḥ yaḥ mitām vā ciṇi iti saḥ ca prayojanam . \\
\noindent \textbf{(P\_6,4.62.3)} śamiṣyate śāmiṣyate tamiṣyate tāmiṣyate . \\
\noindent \textbf{(P\_6,4.62.3)} iṭ ca asiddhaḥ tena me lupyate ṇiḥ nityaḥ ca ayam valnimittaḥ vighātī . \\
\noindent \textbf{(P\_6,4.62.3)} iṭaḥ asiddhatvāt ṇeḥ aniṭi iti ṇilopaḥ yathā syāt . \\
\noindent \textbf{(P\_6,4.62.3)} katham punaḥ ayam nityaḥ . \\
\noindent \textbf{(P\_6,4.62.3)} kṛtākṛtaprasaṅgitvāt . \\
\noindent \textbf{(P\_6,4.62.3)} kṛte api tasmin iṭi sāptamike ārdhadhātukasya iṭ valādeḥ iti punaḥ ayam bhavati . \\
\noindent \textbf{(P\_6,4.62.3)} asmin tu vihite valāditvasya nimittasya vihatatvāt sāptamikaḥ na bhavati \\
\noindent \textbf{(P\_6,4.62.4)} atha upadeśagrahaṇam kimartham . \\
\noindent \textbf{(P\_6,4.62.4)} ciṇvadbhāve upadeśavacanam ṛkāraguṇabalīyastvāt . \\
\noindent \textbf{(P\_6,4.62.4)} ciṇvadbhāve upadeśavacanam kriyate ṛkāraguṇasya balīyastvāt . \\
\noindent \textbf{(P\_6,4.62.4)} kāriṣyate . \\
\noindent \textbf{(P\_6,4.62.4)} paratvāt guṇe kṛte raparatve ca anajantatvāt ciṇvadbhāvaḥ na prāpnoti . \\
\noindent \textbf{(P\_6,4.62.4)} upadeśagrahaṇāt bhaviṣyati . \\
\noindent \textbf{(P\_6,4.62.5)} vadhibhāvāt sīyuṭi ciṇvadbhāvaḥ vipratiṣedhena . \\
\noindent \textbf{(P\_6,4.62.5)} vadhibhāvāt sīyuṭi ciṇvadbhāvaḥ bhavati vipratiṣedhena . \\
\noindent \textbf{(P\_6,4.62.5)} vadhibhāvasya avakāśaḥ : vadhyāt , vadhyāstām , vadhyāsuḥ . \\
\noindent \textbf{(P\_6,4.62.5)} ciṇvadbhāvasya avakāśaḥ : ghāniṣyate , aghāniṣyata . \\
\noindent \textbf{(P\_6,4.62.5)} iha ubhayam prāpnoti : ghāniṣīṣṭa ghāniṣīyāstām ghāniṣīran . \\
\noindent \textbf{(P\_6,4.62.5)} ciṇvadbhāvaḥ bhavati vipratiṣedhena . \\
\noindent \textbf{(P\_6,4.62.5)} atha idānīm ciṇvadbhāve kṛte punaḥprasaṅgavijñānāt vadhibhāvaḥ kasmāt na bhavati . \\
\noindent \textbf{(P\_6,4.62.5)} sakṛdgatau vipratiṣedhe yat bādhitam tat bādhitam eva iti . \\
\noindent \textbf{(P\_6,4.62.5)} haniṇiṅādeśapratiṣedhaḥ ca . \\
\noindent \textbf{(P\_6,4.62.5)} haniṇiṅādeśānām ca pratiṣedhaḥ vaktavyaḥ . \\
\noindent \textbf{(P\_6,4.62.5)} haniṣyate , ghāniṣyate , eṣyate , āyiṣyate , adhyeṣyate , adhyāyiṣyate . \\
\noindent \textbf{(P\_6,4.62.5)} luṅi iti haniṇiṅādeśāḥ prāpnuvanti . \\
\noindent \textbf{(P\_6,4.62.5)} aṅgasya iti tu prakaraṇāt aṅgaśāstrātideśāt siddham . \\
\noindent \textbf{(P\_6,4.62.5)} āṅgam yat kāryam tat pratinirdiśyate na ca haniṇiṅādeśāḥ āṅgāḥ .bhavanti iti . \\
\noindent \textbf{(P\_6,4.64)} atha iḍgrahaṇam kimartham . \\
\noindent \textbf{(P\_6,4.64)} iḍgrahaṇam akṅidartham . \\
\noindent \textbf{(P\_6,4.64)} iḍgrahaṇam kriyate akṅiti lopaḥ yathā syāt : papitha tasthitha iti . \\
\noindent \textbf{(P\_6,4.64)} sārvadhātuke ca ādi iti ārdhadhātukādhikārāt upasaṅkhyānam . \\
\noindent \textbf{(P\_6,4.64)} sārvadhātuke ca ādi iti ārdhadhātukādhikārāt upasaṅkhyānam kartavyam . \\
\noindent \textbf{(P\_6,4.64)} iṣam ūrjam aham itaḥ ādi . \\
\noindent \textbf{(P\_6,4.64)} nanu ca kṅiti iti vartamāne yathā eva iḍgrahaṇam akṅidartham evam ārdhadhātuke iti api vartamāne iḍgrahaṇam sārvadhātukārtham bhaviṣyati . \\
\noindent \textbf{(P\_6,4.64)} na sidhyati . \\
\noindent \textbf{(P\_6,4.64)} kim kāraṇam . \\
\noindent \textbf{(P\_6,4.64)} na hi kṅitā ac viśeṣyate : aci bhavati . \\
\noindent \textbf{(P\_6,4.64)} katarasmin . \\
\noindent \textbf{(P\_6,4.64)} kṅiti iti . \\
\noindent \textbf{(P\_6,4.64)} kim tarhi acā kṅit viśeṣyate : kṅiti bhavati . \\
\noindent \textbf{(P\_6,4.64)} katarasmin . \\
\noindent \textbf{(P\_6,4.64)} aci iti . \\
\noindent \textbf{(P\_6,4.64)} kim punaḥ kāraṇam acā kṅit viśeṣyate . \\
\noindent \textbf{(P\_6,4.64)} yathā iṭ api ajgrahaṇena viśeṣyate . \\
\noindent \textbf{(P\_6,4.64)} asti ca idānīm kva cit iṭ anajādiḥ yadarthaḥ vidhiḥ syāt . \\
\noindent \textbf{(P\_6,4.64)} asti iti āha : dāsīya dhāsīya . \\
\noindent \textbf{(P\_6,4.64)} tat tarhi upasaṅkhyānam kartavyam . \\
\noindent \textbf{(P\_6,4.64)} na kartavyam . \\
\noindent \textbf{(P\_6,4.64)} ārdhadhātukagrahaṇāt siddham . \\
\noindent \textbf{(P\_6,4.64)} katham . \\
\noindent \textbf{(P\_6,4.64)} ārdhadhātukatvam . \\
\noindent \textbf{(P\_6,4.64)} ubhayathā chandasi iti vacanāt . \\
\noindent \textbf{(P\_6,4.64)} anye api dhātupratyayāḥ ubhayathā chandasi dṛśyante . \\
\noindent \textbf{(P\_6,4.66)} īttve vakārapratiṣedhaḥ ghṛtam ghṛtapāvānaḥ iti darśanāt . \\
\noindent \textbf{(P\_6,4.66)} īttve vakāre pratiṣedhaḥ vaktavyaḥ . \\
\noindent \textbf{(P\_6,4.66)} kim prayojanam . \\
\noindent \textbf{(P\_6,4.66)} ghṛtam ghṛtapāvānaḥ iti darśanāt . \\
\noindent \textbf{(P\_6,4.66)} iha mā bhūt : ghṛtam ghṛtapāvānaḥ pibata . \\
\noindent \textbf{(P\_6,4.66)} vasām vasapāvānaḥ pibata iti . \\
\noindent \textbf{(P\_6,4.66)} yadi tarhi vakāre pratiṣedhaḥ ucyate katham dīvarī pīvarī iti . \\
\noindent \textbf{(P\_6,4.66)} dhīvarī pīvarī iti ca uktam . \\
\noindent \textbf{(P\_6,4.66)} kim uktam . \\
\noindent \textbf{(P\_6,4.66)} na etat īttvam . \\
\noindent \textbf{(P\_6,4.66)} kim tarhi . \\
\noindent \textbf{(P\_6,4.66)} dhyāpyoḥ etat samprasāraṇam iti . \\
\noindent \textbf{(P\_6,4.66)} saḥ tarhi pratiṣedhaḥ vaktayaḥ . \\
\noindent \textbf{(P\_6,4.66)} na vaktavyaḥ . \\
\noindent \textbf{(P\_6,4.66)} vanip eṣaḥ bhaviṣyati na kvanip iti . \\
\noindent \textbf{(P\_6,4.74)} kasya ayam pratiṣedhaḥ . \\
\noindent \textbf{(P\_6,4.74)} āṭaḥ prāpnoti . \\
\noindent \textbf{(P\_6,4.74)} aṭaḥ api iṣyate . \\
\noindent \textbf{(P\_6,4.74)} tat tarhi aṭaḥ grahaṇam kartavyam . \\
\noindent \textbf{(P\_6,4.74)} na kartavyam . \\
\noindent \textbf{(P\_6,4.74)} prakṛtam anuvartate . \\
\noindent \textbf{(P\_6,4.74)} kva prakṛtam . \\
\noindent \textbf{(P\_6,4.74)} luṅlaṅlṅkṣu aṭ udāttaḥ iti . \\
\noindent \textbf{(P\_6,4.74)} yadi tat anuvartate āṭ ajādīnām aṭ ca iti aṭ api prāpnoti . \\
\noindent \textbf{(P\_6,4.74)} astu . \\
\noindent \textbf{(P\_6,4.74)} aṭi kṛte punaḥ āṭi bhaviṣyati . \\
\noindent \textbf{(P\_6,4.74)} iha api tarhi aṭi kṛte punaḥ āṭ prāpnoti : akārṣīt , ahārṣīt . \\
\noindent \textbf{(P\_6,4.74)} aḍvacanāt na bhaviṣyati . \\
\noindent \textbf{(P\_6,4.74)} iha api tarhi aḍvacanāt na syāt : aihiṣṭa , aikṣiṣta . \\
\noindent \textbf{(P\_6,4.74)} āḍvacanāt bhaviṣyati . \\
\noindent \textbf{(P\_6,4.74)} iha api tarhi āḍvacanāt prāpnoti : akārṣīt , ahārṣīt . \\
\noindent \textbf{(P\_6,4.74)} akṛte aṭi yaḥ ajādiḥ iti evam etat vijñāsyate . \\
\noindent \textbf{(P\_6,4.74)} kim vaktavyam etat . \\
\noindent \textbf{(P\_6,4.74)} na hi . \\
\noindent \textbf{(P\_6,4.74)} katham anucyamānam gaṃsyate . \\
\noindent \textbf{(P\_6,4.74)} ajvacanasāmarthyāt . \\
\noindent \textbf{(P\_6,4.74)} yadi kṛte aṭi yaḥ ajādiḥ tatra syāt ajgrahaṇam anarthakam syāt . \\
\noindent \textbf{(P\_6,4.74)} atha vā upadeśe iti vartate . \\
\noindent \textbf{(P\_6,4.74)} atha vā ārdhadhātuke iti vartate . \\
\noindent \textbf{(P\_6,4.74)} atha vā luṅlaṅlṅkṣu aṭ iti dvilakārakaḥ nirdeśaḥ : luṅādiṣu lakārādiṣu yaḥ ajādiḥ iti . \\
\noindent \textbf{(P\_6,4.74)} sarvathā , aijyata , aupyata iti etat na sidhyati . \\
\noindent \textbf{(P\_6,4.74)} evam tarhi ajādīnām aṭā siddham . \\
\noindent \textbf{(P\_6,4.74)} ajādīnām aṭā eva siddham . \\
\noindent \textbf{(P\_6,4.74)} na arthaḥ āṭā . \\
\noindent \textbf{(P\_6,4.74)} evam tarhi vṛddhyartham āṭ vaktavyaḥ . \\
\noindent \textbf{(P\_6,4.74)} vṛddhyartham iti cet aṭaḥ . \\
\noindent \textbf{(P\_6,4.74)} aṭaḥ vṛddhim vakṣyāmi . \\
\noindent \textbf{(P\_6,4.74)} yadi tarhi aṭaḥ vṛddhiḥ ucyate asvavaḥ hasati iti atra . vṛddhiḥ prapnoti roḥ utve kṛte . \\
\noindent \textbf{(P\_6,4.74)} dhātau vṛddhim aṭaḥ smaret . dhātau aṭaḥ vṛddhim vakṣyāmi . \\
\noindent \textbf{(P\_6,4.74)} tat tarhi dhātugrahaṇam kartavyam . \\
\noindent \textbf{(P\_6,4.74)} na kartavyam . \\
\noindent \textbf{(P\_6,4.74)} yogavibhāgaḥ kariṣyate . \\
\noindent \textbf{(P\_6,4.74)} aṭaḥ aci vṛddhiḥ bhavati . \\
\noindent \textbf{(P\_6,4.74)} tataḥ upasargāt ṛti vṛddhiḥ bhavati . \\
\noindent \textbf{(P\_6,4.74)} tataḥ dhātau . \\
\noindent \textbf{(P\_6,4.74)} dhātau iti ubhayoḥ śeṣaḥ . \\
\noindent \textbf{(P\_6,4.74)} iha tarhi : āṭīt , āśīt iti ataḥ guṇe iti pararūpatvam prāpnoti . \\
\noindent \textbf{(P\_6,4.74)} pararūpam guṇe na aṭaḥ . \\
\noindent \textbf{(P\_6,4.74)} pararūpam guṇe aṭaḥ na iti vakṣyāmi . \\
\noindent \textbf{(P\_6,4.74)} omāṅoḥ usi tat samam . \\
\noindent \textbf{(P\_6,4.74)} yadi api etat ucyate atha vā etarhi usi omāṅkṣu āṭaḥ pararūpapratiṣedhaḥ coditaḥ sa na vaktavyaḥ bhavati . \\
\noindent \textbf{(P\_6,4.74)} chandortham tarhi āṭ vaktavyaḥ . \\
\noindent \textbf{(P\_6,4.74)} araik u kṛṣṇāḥ . \\
\noindent \textbf{(P\_6,4.74)} tritaḥ enam āyunak . \\
\noindent \textbf{(P\_6,4.74)} surucaḥ ven āvaḥ . \\
\noindent \textbf{(P\_6,4.74)} chandortham bahulam dīrgham . \\
\noindent \textbf{(P\_6,4.74)} bahulam chandasi dīrghatvam dṛśyate . \\
\noindent \textbf{(P\_6,4.74)} tat yathā : pūruṣaḥ , nārakaḥ iti . \\
\noindent \textbf{(P\_6,4.74)} evam tarhi āyan , āsan . \\
\noindent \textbf{(P\_6,4.74)} iṇastyoḥ yaṇlopayoḥ kṛtayoḥ anajāditvāt vṛddhiḥ na prāpnoti . \\
\noindent \textbf{(P\_6,4.74)} iṇastyoḥ antaraṅgataḥ . \\
\noindent \textbf{(P\_6,4.74)} antaraṅgatvāt vṛddhiḥ bhaviṣyati . \\
\noindent \textbf{(P\_6,4.74)} tasmāt na arthaḥ āḍgrahaṇena . \\
\noindent \textbf{(P\_6,4.74)} ajādīnām aṭā siddham . \\
\noindent \textbf{(P\_6,4.74)} vṛddhyartham iti cet aṭaḥ . \\
\noindent \textbf{(P\_6,4.74)} asvavaḥ hasati iti atra . \\
\noindent \textbf{(P\_6,4.74)} dhātau vṛddhim aṭaḥ smaret . \\
\noindent \textbf{(P\_6,4.74)} pararūpam guṇe na aṭaḥ . \\
\noindent \textbf{(P\_6,4.74)} omāṅoḥ usi tat samam . \\
\noindent \textbf{(P\_6,4.74)} chandortham bahulam dīrgham . \\
\noindent \textbf{(P\_6,4.74)} iṇastyoḥ antaraṅgataḥ . \\
\noindent \textbf{(P\_6,4.77)} iyaṅādiprakaraṇe tanvādīnām chandasi bahulam . \\
\noindent \textbf{(P\_6,4.77)} iyaṅādiprakaraṇe tanvādīnām chandasi bahulam upasaṅkhyānam kartavyam . \\
\noindent \textbf{(P\_6,4.77)} tanvam puṣema . \\
\noindent \textbf{(P\_6,4.77)} tanuvam puṣema . \\
\noindent \textbf{(P\_6,4.77)} viṣvam paśya . \\
\noindent \textbf{(P\_6,4.77)} viṣuvam paśya . \\
\noindent \textbf{(P\_6,4.77)} svargam lokam . \\
\noindent \textbf{(P\_6,4.77)} suvargam lokam . \\
\noindent \textbf{(P\_6,4.77)} tryambakam yajāmahe . \\
\noindent \textbf{(P\_6,4.77)} triyambakam yajāmahe . \\
\noindent \textbf{(P\_6,4.82)} atha iha kasmāt na bhavati : brāhmaṇasya niyau , brāhmaṇasya niyaḥ . \\
\noindent \textbf{(P\_6,4.82)} aṅgādhikārāt . \\
\noindent \textbf{(P\_6,4.82)} aṅgasya iti anuvartate . \\
\noindent \textbf{(P\_6,4.82)} evam api paramaniyau paramaniyaḥ iti atra prāpnoti . \\
\noindent \textbf{(P\_6,4.82)} gatikārakapūrvasya iṣyate . \\
\noindent \textbf{(P\_6,4.82)} yaṇādeśaḥ svarapadapūrvopadhasya ca . \\
\noindent \textbf{(P\_6,4.82)} yaṇādeśaḥ svarapūrvopadhasya padapūrvopadhasya ca iti vaktavyam . \\
\noindent \textbf{(P\_6,4.82)} svarapūrvopadhasya : ninyatuḥ , ninyuḥ . \\
\noindent \textbf{(P\_6,4.82)} padapūrvopadhasya : unnyau , unnyaḥ , uddhyau , uddhyaḥ . \\
\noindent \textbf{(P\_6,4.82)} ubhayakṛtam: grāmaṇyau , grāmaṇyaḥ , senānyau , senānyaḥ . \\
\noindent \textbf{(P\_6,4.82)} asaṃyogapūrve hi aniṣṭaprasaṅgaḥ . \\
\noindent \textbf{(P\_6,4.82)} asaṃyogapūrvasya iti hi ucyamāne aniṣṭam prasajyeta . \\
\noindent \textbf{(P\_6,4.82)} uddhyau , uddhyaḥ , unnyau , unnyaḥ . \\
\noindent \textbf{(P\_6,4.82)} asaṃyogapūrvasya iti pratiṣedhaḥ prasajyeta . \\
\noindent \textbf{(P\_6,4.82)} tat tarhi vaktavyam . \\
\noindent \textbf{(P\_6,4.82)} na vaktavyam . \\
\noindent \textbf{(P\_6,4.82)} dhātoḥ iti vartate . \\
\noindent \textbf{(P\_6,4.82)} tatra dhātunā saṃyogam viśeṣayiṣyāmaḥ . \\
\noindent \textbf{(P\_6,4.82)} dhātoḥ yaḥ saṃyogaḥ tatpūrvasya na iti . \\
\noindent \textbf{(P\_6,4.82)} upasarjanam vai saṃyogaḥ na ca upasarjanasya viśeṣaṇam asti . \\
\noindent \textbf{(P\_6,4.82)} dhātoḥ iti anuvartanasāmarthyāt upasarjanasya api viśeṣaṇam bhaviṣyati . \\
\noindent \textbf{(P\_6,4.82)} asti anyat dhātoḥ iti anuvartanasya prayojanam . \\
\noindent \textbf{(P\_6,4.82)} kim . \\
\noindent \textbf{(P\_6,4.82)} ivarṇam viśeṣayiṣyāmaḥ . \\
\noindent \textbf{(P\_6,4.82)} na etat asti prayojanam . \\
\noindent \textbf{(P\_6,4.82)} yat hi adhātoḥ ivarṇam bhavitavyam eva tasya yaṇādeśena ikaḥ yaṇ aci iti eva . \\
\noindent \textbf{(P\_6,4.84)} varṣābhūpunarbhvaḥ ca . \\
\noindent \textbf{(P\_6,4.84)} varṣābhū iti atra punarbhvaḥ ca iti vaktavyam : punarbhvau , punarbhvaḥ . \\
\noindent \textbf{(P\_6,4.84)} atyalpam idam ucyate . \\
\noindent \textbf{(P\_6,4.84)} varṣādṛnkārapunaḥpūrvasya bhuvaḥ iti vaktavyam : varṣābhvau , varṣābhvaḥ , dṛnbhvau , dṛnbhvaḥ , kārabhvau , kārabhvaḥ , punarbhvau , punarbhvaḥ . \\
\noindent \textbf{(P\_6,4.87)} huśnugrahaṇam anarthakam . \\
\noindent \textbf{(P\_6,4.87)} kim kāraṇam . \\
\noindent \textbf{(P\_6,4.87)} anyasya abhāvāt . \\
\noindent \textbf{(P\_6,4.87)} na hi anyat sārvadhātuke asti yasya yaṇādeśaḥ syāt . \\
\noindent \textbf{(P\_6,4.87)} nanu ca ayam asti : yāti , vāti iti . \\
\noindent \textbf{(P\_6,4.87)} kṅiti anuvartate . \\
\noindent \textbf{(P\_6,4.87)} iha tarhi : yātaḥ , vātaḥ iti . \\
\noindent \textbf{(P\_6,4.87)} aci iti vartate . \\
\noindent \textbf{(P\_6,4.87)} iha tarhi : yānti , vānti . \\
\noindent \textbf{(P\_6,4.87)} yvoḥ iti vartate . \\
\noindent \textbf{(P\_6,4.87)} evam api dhiyanti , piyanti iti atra prāpnoti . \\
\noindent \textbf{(P\_6,4.87)} oḥ iti vartate . \\
\noindent \textbf{(P\_6,4.87)} evam api suvanti , ruvanti iti atra prāpnoti . \\
\noindent \textbf{(P\_6,4.87)} anekācaḥ iti vartate . \\
\noindent \textbf{(P\_6,4.87)} evam api asuvan , aruvan iti atra prāpnoti . \\
\noindent \textbf{(P\_6,4.87)} etat api aṭaḥ asiddhatvāt ekāc bhavati . \\
\noindent \textbf{(P\_6,4.87)} evam api prorṇuvanti iti atra prāpnoti . \\
\noindent \textbf{(P\_6,4.87)} asaṃyogapūrvasya iti vartate . \\
\noindent \textbf{(P\_6,4.87)} yaṅlugartham tarhi huśnugrahaṇam kartavyam . \\
\noindent \textbf{(P\_6,4.87)} yaṅlugantam anekāc asaṃyogapūrvam uvarṇāntam asti . \\
\noindent \textbf{(P\_6,4.87)} tadartham idam . \\
\noindent \textbf{(P\_6,4.87)} nadam yoyuvatīnām . \\
\noindent \textbf{(P\_6,4.87)} vṛṣabham roruvatīnām . \\
\noindent \textbf{(P\_6,4.87)} yaṅlugartham iti cet ārdhadhātukatvāt siddham . \\
\noindent \textbf{(P\_6,4.87)} yaṅlugartham iti cet tat na . \\
\noindent \textbf{(P\_6,4.87)} kim kāraṇam . \\
\noindent \textbf{(P\_6,4.87)} ārdhadhātukatvāt siddham . \\
\noindent \textbf{(P\_6,4.87)} katham ārdhadhātukatvam . \\
\noindent \textbf{(P\_6,4.87)} ubhayathā chandasi iti vacanāt . \\
\noindent \textbf{(P\_6,4.87)} anye api hi dhātupratyayāḥ ubhayathā chandasi dṛśyante . \\
\noindent \textbf{(P\_6,4.87)} evam tarhi siddhe sati yat huśnugrahaṇam karoti tat jñāpayati ācāryaḥ yaṅluk bhāṣāyām bhavati iti . \\
\noindent \textbf{(P\_6,4.87)} kim etasya jñāpane prayojanam . \\
\noindent \textbf{(P\_6,4.87)} bebhidīti , cecchidīti etat siddham bhavati bhāṣāyām api . \\
\noindent \textbf{(P\_6,4.89)} atha kimartham guheḥ vikṛtasya grahaṇam kriyate na punaḥ guhaḥ iti eva ucyeta . \\
\noindent \textbf{(P\_6,4.89)} gohigrahaṇam viṣayārtham . \\
\noindent \textbf{(P\_6,4.89)} gohigrahaṇam kriyate viṣayārtham . \\
\noindent \textbf{(P\_6,4.89)} viṣayaḥ pratinirdiśyate . \\
\noindent \textbf{(P\_6,4.89)} yatra asya etat rūpam tatra yathā syāt . \\
\noindent \textbf{(P\_6,4.89)} iha mā bhūt : nijuguhatuḥ , nijuguhuḥ iti . \\
\noindent \textbf{(P\_6,4.89)} ayādeśapratiṣedhārtham ca . \\
\noindent \textbf{(P\_6,4.89)} ayādeśapratiṣedhārtham ca vikṛtagrahaṇam kriyate . \\
\noindent \textbf{(P\_6,4.89)} hrasvādeśe hi ayādeśaprasaṅgaḥ ūttvasya asiddhatvāt . \\
\noindent \textbf{(P\_6,4.89)} hrasvādeśe hi sati ayādeśaḥ prasajyeta . \\
\noindent \textbf{(P\_6,4.89)} pragūhya gataḥ . \\
\noindent \textbf{(P\_6,4.89)} kim kāraṇam . \\
\noindent \textbf{(P\_6,4.89)} ūttvasya asiddhatvāt . \\
\noindent \textbf{(P\_6,4.89)} asiddham ūttvam . \\
\noindent \textbf{(P\_6,4.89)} tasya asiddhatvāt lyapi laghupūrvāt iti ayādeśaḥ prasajyeta . \\
\noindent \textbf{(P\_6,4.89)} viṣayārthena tāvat na arthaḥ gohigrahaṇena . \\
\noindent \textbf{(P\_6,4.89)} praśliṣṭanirdeśāt siddham . \\
\noindent \textbf{(P\_6,4.89)} praśliṣṭanirdeśaḥ ayam . \\
\noindent \textbf{(P\_6,4.89)} u-ūt : ūt iti . \\
\noindent \textbf{(P\_6,4.89)} tatra hrasvasya avakāśaḥ : nijuguhatuḥ , nijuguhuḥ . \\
\noindent \textbf{(P\_6,4.89)} guṇasya avakāśaḥ : nigoḍhā , nogoḍhum . \\
\noindent \textbf{(P\_6,4.89)} iha ubhayam prāpnoti : nigūhayati , nigūhakaḥ . \\
\noindent \textbf{(P\_6,4.89)} paratvāt guṇe kṛte āntaryataḥ dīrghasya dīrghaḥ bhaviṣyati . \\
\noindent \textbf{(P\_6,4.89)} ayādeśapratiṣedhārthena api na arthaḥ . \\
\noindent \textbf{(P\_6,4.89)} samānāśrayavacanāt siddham . \\
\noindent \textbf{(P\_6,4.89)} samānāśrayam asiddham bhavati vyāśrayam ca etat . \\
\noindent \textbf{(P\_6,4.89)} katham . \\
\noindent \textbf{(P\_6,4.89)} ṇau ūttvam ṇeḥ lyapi ayādeśaḥ . \\
\noindent \textbf{(P\_6,4.90)} atha kimartham duṣeḥ vikṛtasya grahaṇam kriyate na punaḥ duṣaḥ iti eva ucyeta . \\
\noindent \textbf{(P\_6,4.90)} doṣigrahaṇam ca . \\
\noindent \textbf{(P\_6,4.90)} kim . \\
\noindent \textbf{(P\_6,4.90)} ayādeśapratiṣedhārtham hrasvādeśe hi ayādeśaprasaṅgaḥ ūttvasya asiddhatvāt . \\
\noindent \textbf{(P\_6,4.90)} hrasvādeśe hi sati ayādeśaḥ prasajyeta . \\
\noindent \textbf{(P\_6,4.90)} pradūṣya gataḥ . \\
\noindent \textbf{(P\_6,4.90)} kim kāraṇam . \\
\noindent \textbf{(P\_6,4.90)} ūttvasya asiddhatvāt . \\
\noindent \textbf{(P\_6,4.90)} asiddham ūttvam . \\
\noindent \textbf{(P\_6,4.90)} tasya asiddhatvāt lyapi laghupūrvāt iti ayādeśaḥ prasajyeta . \\
\noindent \textbf{(P\_6,4.90)} atra api samānāśrayavacanāt siddham iti eva . \\
\noindent \textbf{(P\_6,4.93)} ciṇṇamuloḥ ṇijvyavetānām yaṅlope ca upasaṅkhyānam kartavyam . \\
\noindent \textbf{(P\_6,4.93)} śamayantam prayojitavān , aśami , aśāmi , śamam śamam , śāmam śāmam . \\
\noindent \textbf{(P\_6,4.93)} śaṃśamayateḥ : aśaṃśami , aśaṃśāmi , śaṃśamam śaṃśamam , śaṃśāmam śaṃśāmam . \\
\noindent \textbf{(P\_6,4.93)} kim punaḥ kāraṇam na sidhyati . \\
\noindent \textbf{(P\_6,4.93)} ciṇṇamulpare ṇau mitām aṅgānām dīrghaḥ bhavati iti ucyate . \\
\noindent \textbf{(P\_6,4.93)} yaḥ ca atra ciṇṇamulparaḥ na tasmin mit aṅgam yasmin ca mit aṅgam na asau ciṇṇamulparaḥ iti . \\
\noindent \textbf{(P\_6,4.93)} lope kṛte ciṇṇamulparaḥ bhavati . \\
\noindent \textbf{(P\_6,4.93)} sthānivadbhāvāt na ciṇṇamulparaḥ . \\
\noindent \textbf{(P\_6,4.93)} nanu ca pratiṣidhyate atra sthānivadbhāvaḥ dīrghavidhim prati na sthānivat iti . \\
\noindent \textbf{(P\_6,4.93)} evam api asiddhatvāt na prāpnoti . \\
\noindent \textbf{(P\_6,4.93)} evam tarhi ciṇṇamuloḥ ṇijvyavetānām yaṅlope ca antaraṅgalakṣaṇatvāt siddham . \\
\noindent \textbf{(P\_6,4.93)} kim idam antaraṅgalakṣaṇatvāt iti . \\
\noindent \textbf{(P\_6,4.93)} yāvat brūyāt samānāśrayavacanāt siddham iti eva vyāśrayam ca etat . \\
\noindent \textbf{(P\_6,4.93)} katham . \\
\noindent \textbf{(P\_6,4.93)} ṇeḥ ṇau lopaḥ ṇau ciṇṇamulpare mitām aṅgānām dīrghatvam ucyate . \\
\noindent \textbf{(P\_6,4.93)} tasmāt na arthaḥ upasaṅkhyānena iti . \\
\noindent \textbf{(P\_6,4.96)} adviprabhṛtyupasargasya iti vaktavyam iha api yathā syāt : samupābhicchādaḥ iti . \\
\noindent \textbf{(P\_6,4.96)} tat tarhi vaktavyam . \\
\noindent \textbf{(P\_6,4.96)} na vaktavyam . \\
\noindent \textbf{(P\_6,4.96)} yatra triprabhṛtayaḥ santi dvau api tatra staḥ . \\
\noindent \textbf{(P\_6,4.96)} tatra advyupasargasya iti eva siddham . \\
\noindent \textbf{(P\_6,4.96)} na vai eṣaḥ loke sampratyayaḥ . \\
\noindent \textbf{(P\_6,4.96)} na hi dviputraḥ ānīyatām iti ukte triputraḥ ānīyate . \\
\noindent \textbf{(P\_6,4.96)} tasmāt adviprabhṛtyupasargasya iti vaktavyam . \\
\noindent \textbf{(P\_6,4.100)} halgrahaṇam anarthakam anyatra api darśanāt . \\
\noindent \textbf{(P\_6,4.100)} halgrahaṇam anarthakam . \\
\noindent \textbf{(P\_6,4.100)} kim kāraṇam . \\
\noindent \textbf{(P\_6,4.100)} anyatra api darśanāt . \\
\noindent \textbf{(P\_6,4.100)} anyatra api lopaḥ dṛśyate . \\
\noindent \textbf{(P\_6,4.100)} agniḥ tṛṇāni babsati . \\
\noindent \textbf{(P\_6,4.100)} śarāve bapsati caruḥ . \\
\noindent \textbf{(P\_6,4.101)} iṭaḥ pratiṣedhaḥ vaktavyaḥ . \\
\noindent \textbf{(P\_6,4.101)} rudihi svapihi . \\
\noindent \textbf{(P\_6,4.101)} jhalaḥ iti dhitvam prāpnoti . \\
\noindent \textbf{(P\_6,4.101)} heḥ dhitve haladhikārāt iṭaḥ apratiṣedhaḥ . \\
\noindent \textbf{(P\_6,4.101)} heḥ dhitve haladhikārāt iṭaḥ apratiṣedhaḥ . \\
\noindent \textbf{(P\_6,4.101)} anarthakaḥ pratiṣedhaḥ apratiṣedhaḥ . \\
\noindent \textbf{(P\_6,4.101)} dhitvam kasmāt na bhavati . \\
\noindent \textbf{(P\_6,4.101)} haladhikārāt . \\
\noindent \textbf{(P\_6,4.101)} prakṛtam halgrahaṇam anuvartate . \\
\noindent \textbf{(P\_6,4.101)} kva prakṛtam . \\
\noindent \textbf{(P\_6,4.101)} ghasibhasoḥ hali iti . \\
\noindent \textbf{(P\_6,4.101)} tat vai saptamīnirdiṣṭam ṣaṣṭhīnirdiṣṭena ca iha arthaḥ . \\
\noindent \textbf{(P\_6,4.101)} tat vai tatra pratyākhyāyate . \\
\noindent \textbf{(P\_6,4.101)} tatra pratyākhyātam sat yayā vibhaktyā nirdiśyamānam arthavattayā nirdiṣṭam iha anuvartiṣyate . \\
\noindent \textbf{(P\_6,4.101)} atha vā hujhalbhayaḥ iti eṣā pañcamī hali iti saptamyāḥ ṣaṣṭhīm prakalpayiṣyati tasmāt iti uttarasya iti . \\
\noindent \textbf{(P\_6,4.101)} atha vā nirdiśyamānasya ādeśāḥ bhavanti iti evam na bhaviṣyati . \\
\noindent \textbf{(P\_6,4.101)} yaḥ tarhi nirdiśyate tasya kasmāt na bhavati . \\
\noindent \textbf{(P\_6,4.101)} iṭā vyavahitatvāt . \\
\noindent \textbf{(P\_6,4.101)} yadi evam chindhaki bhindhaki iti atra dhitvam na prāpnoti . \\
\noindent \textbf{(P\_6,4.101)} dhitve kṛte akac bhaviṣyati . \\
\noindent \textbf{(P\_6,4.101)} idam iha sampradhāryam . \\
\noindent \textbf{(P\_6,4.101)} dhitvam kriyatām akac iti kim atra kartavyam . \\
\noindent \textbf{(P\_6,4.101)} paratvāt dhitvam . \\
\noindent \textbf{(P\_6,4.101)} nityaḥ akac . \\
\noindent \textbf{(P\_6,4.101)} kṛte api dhitve prāpnoti akṛte api . \\
\noindent \textbf{(P\_6,4.101)} akac api anityaḥ . \\
\noindent \textbf{(P\_6,4.101)} anyasya kṛte dhitve prāpnoti anyasya akṛte śabdāntarasya ca prāpnuvan vidhiḥ anityaḥ bhavati . \\
\noindent \textbf{(P\_6,4.101)} ubhayoḥ anityayoḥ paratvāt dhitve kṛte akac bhaviṣyati . \\
\noindent \textbf{(P\_6,4.101)} atha vā hakārasya eva aśaktijena ikāreṇa grahaṇam . \\
\noindent \textbf{(P\_6,4.104)} ciṇaḥ luki tagrahaṇam kartavyam . \\
\noindent \textbf{(P\_6,4.104)} kim prayojanam . \\
\noindent \textbf{(P\_6,4.104)} iha mā bhūt : akāritarām , ahāritarām iti . \\
\noindent \textbf{(P\_6,4.104)} ciṇaḥ luki tagrahaṇānarthakyam saṅghātasya apratyayatvāt . ciṇaḥ luki tagrahaṇam anarthakam . \\
\noindent \textbf{(P\_6,4.104)} kim kāraṇam . \\
\noindent \textbf{(P\_6,4.104)} saṅghātasya apratyayatvāt . \\
\noindent \textbf{(P\_6,4.104)} saṅghātasya luk kasmāt na bhavati . \\
\noindent \textbf{(P\_6,4.104)} apratyayatvāt . \\
\noindent \textbf{(P\_6,4.104)} pratyayasya lukślulupaḥ bhavanti iti ucyate na ca saṅghātaḥ pratyayaḥ . \\
\noindent \textbf{(P\_6,4.104)} talope tarhi kṛte parasya prāpnoti . \\
\noindent \textbf{(P\_6,4.104)} talopasya ca asiddhatvāt . \\
\noindent \textbf{(P\_6,4.104)} asiddhaḥ talopaḥ . \\
\noindent \textbf{(P\_6,4.104)} tasya asiddhatvāt na bhaviṣyati . \\
\noindent \textbf{(P\_6,4.104)} kāryakṛtatvāt vā . \\
\noindent \textbf{(P\_6,4.104)} atha vā kṛtaḥ ciṇaḥ luk iti kṛtvā punaḥ na bhaviṣyati luk . \\
\noindent \textbf{(P\_6,4.104)} tat yathā vasante brāhmaṇaḥ agnīn ādadhīta iti sakṛt ādhāya kṛtaḥ śāstrārthaḥ iti kṛtvā punaḥ pravṛttiḥ na bhavati . \\
\noindent \textbf{(P\_6,4.104)} viṣamaḥ upanyāsaḥ . \\
\noindent \textbf{(P\_6,4.104)} yuktam yat tasya eva punaḥ pravṛttiḥ na syāt . \\
\noindent \textbf{(P\_6,4.104)} yat tu tadāśrayam prāpnoti na tat śakyam bādhitum . \\
\noindent \textbf{(P\_6,4.104)} tat yathā vasante brāhmaṇaḥ agniṣṭomādibhiḥ kratubhiḥ yajeta iti agnyādhānanimittam vasante vasante ijyate . \\
\noindent \textbf{(P\_6,4.104)} tasmāt pūrvoktau eva parihārau . \\
\noindent \textbf{(P\_6,4.104)} atha vā kṅiti iti vartate . \\
\noindent \textbf{(P\_6,4.104)} kva prakṛtam . \\
\noindent \textbf{(P\_6,4.104)} gamahanajanakhanaghasām lopaḥ kṅiti anaṅi iti . \\
\noindent \textbf{(P\_6,4.104)} tat vai saptamīnirdiṣṭam ṣaṣṭhīnirdiṣṭena ca iha arthaḥ . \\
\noindent \textbf{(P\_6,4.104)} ciṇaḥ luk iti eṣā pañcamī kṅiti iti saptamyāḥ ṣaṣṭhīm prakalpayiṣyati tasmāt iti uttarasya iti . \\
\noindent \textbf{(P\_6,4.106.1)} katham idam vijñāyate . \\
\noindent \textbf{(P\_6,4.106.1)} ukārāt pratyayāt iti āhosvit ukārāntāt pratyayāt iti . \\
\noindent \textbf{(P\_6,4.106.1)} kim ca ataḥ . \\
\noindent \textbf{(P\_6,4.106.1)} yadi vijñāyate ukārāt pratyayāt iti siddham tanu kuru . \\
\noindent \textbf{(P\_6,4.106.1)} cinu sunu iti na sidhyati . \\
\noindent \textbf{(P\_6,4.106.1)} atha vijñāyate ukārāntāt pratyayāt iti siddham cinu sunu iti . \\
\noindent \textbf{(P\_6,4.106.1)} tanu kuru na sidhyati . \\
\noindent \textbf{(P\_6,4.106.1)} tathā asaṃyogapūrvagrahaṇena iha eva paryudāsaḥ syāt : takṣṇuhi , akṣṇuhi . \\
\noindent \textbf{(P\_6,4.106.1)} āpnuhi śaknuhi iti atra na syāt . \\
\noindent \textbf{(P\_6,4.106.1)} yathā icchasi tathā astu . \\
\noindent \textbf{(P\_6,4.106.1)} astu tāvat ukārāt pratyayāt iti . \\
\noindent \textbf{(P\_6,4.106.1)} katham cinu sunu iti . \\
\noindent \textbf{(P\_6,4.106.1)} tadantavidhinā bhaviṣyati . \\
\noindent \textbf{(P\_6,4.106.1)} atha vā punaḥ astu ukārāntāt pratyayāt iti . \\
\noindent \textbf{(P\_6,4.106.1)} katham tanu kuru iti . \\
\noindent \textbf{(P\_6,4.106.1)} vyapdeśivadbhāvena bhaviṣyati . \\
\noindent \textbf{(P\_6,4.106.1)} yat api ucyate tathā asaṃyogapūrvagrahaṇena iha eva paryudāsaḥ syāt : takṣṇuhi , akṣṇuhi . \\
\noindent \textbf{(P\_6,4.106.1)} āpnuhi śaknuhi iti atra na syāt iti . \\
\noindent \textbf{(P\_6,4.106.1)} na asmābhiḥ asaṃyogapūrvagrahaṇena ukārāntam viśeṣyate . \\
\noindent \textbf{(P\_6,4.106.1)} kim tarhi . \\
\noindent \textbf{(P\_6,4.106.1)} ukāraḥ . \\
\noindent \textbf{(P\_6,4.106.1)} ukāraḥ yaḥ asaṃyogapūrvaḥ tadantāt pratyayāt iti . \\
\noindent \textbf{(P\_6,4.106.2)} utaḥ ca pratyayāt chandovāvacanam . \\
\noindent \textbf{(P\_6,4.106.2)} utaḥ ca pratyayāt iti atra chandasi vā iti vaktavyam . \\
\noindent \textbf{(P\_6,4.106.2)} ava sthira tanuhi yātujunām . \\
\noindent \textbf{(P\_6,4.106.2)} dhinuhi yajñam dhinuhi yajñapatim . \\
\noindent \textbf{(P\_6,4.106.2)} tena mā bhāginam kṛṇuhi . \\
\noindent \textbf{(P\_6,4.106.2)} uttarārtham ca . \\
\noindent \textbf{(P\_6,4.106.2)} ke cit tāvat āhuḥ chandograhaṇam kartavyam iti . \\
\noindent \textbf{(P\_6,4.106.2)} apare āhuḥ : vāvacanam kartavyam iti . \\
\noindent \textbf{(P\_6,4.106.2)} lopaḥ ca asya anyaratasyām mvoḥ iti atra anyaratasyāṅgrahaṇam na kartavyam bhavati . \\
\noindent \textbf{(P\_6,4.110)} sārvadhātuke iti kimartham . \\
\noindent \textbf{(P\_6,4.110)} iha mā bhūt : sañcaskaratuḥ , sañcaskaruḥ . \\
\noindent \textbf{(P\_6,4.110)} syāntasya pratiṣedhaḥ vaktavyaḥ : kariṣyati kariṣyataḥ . \\
\noindent \textbf{(P\_6,4.110)} kṛñaḥ uttve ukārāntanirdeśāt syāntasya apratiṣedhaḥ . \\
\noindent \textbf{(P\_6,4.110)} kṛñaḥ uttve ukārāntanirdeśāt syāntasya apratiṣedhaḥ . \\
\noindent \textbf{(P\_6,4.110)} anarthakaḥ pratiṣedhaḥ apratiṣedhaḥ . \\
\noindent \textbf{(P\_6,4.110)} uttvam kasmāt na bhavati . \\
\noindent \textbf{(P\_6,4.110)} ukārāntanirdeśāt . \\
\noindent \textbf{(P\_6,4.110)} aśakyaḥ karotau ukārāntanirdeśaḥ tantram āśrayitum . \\
\noindent \textbf{(P\_6,4.110)} iha samparibhyām bhūṣaṇasamavāyayoḥ krotau iha eva syāt : saṃskaroti . \\
\noindent \textbf{(P\_6,4.110)} saṃskartā saṃsakrtum iti atra na syāt . \\
\noindent \textbf{(P\_6,4.110)} na brūmaḥ asmāt ukārāntanirdeśāt yaḥ ayam karoti iti . \\
\noindent \textbf{(P\_6,4.110)} kim tarhi . \\
\noindent \textbf{(P\_6,4.110)} ukāraprakaraṇāt ukārāntam aṅgam abhisambadhyate . \\
\noindent \textbf{(P\_6,4.110)} utaḥ iti vartate . \\
\noindent \textbf{(P\_6,4.110)} yadi evam na arthaḥ sārvadhātukagrahaṇena . \\
\noindent \textbf{(P\_6,4.110)} kasmāt na bhavati sañcaskaratuḥ , sañcaskaruḥ iti . \\
\noindent \textbf{(P\_6,4.110)} utaḥ iti vartate . \\
\noindent \textbf{(P\_6,4.110)} uttarārtham tarhi sārvadhātukagrahaṇam kartavyam . \\
\noindent \textbf{(P\_6,4.110)} śnasoḥ allopaḥ iti . \\
\noindent \textbf{(P\_6,4.110)} śnam sārvadhātuke eva . \\
\noindent \textbf{(P\_6,4.110)} asteḥ api ārdhadhātuke bhūbhāvena bhavitavyam . \\
\noindent \textbf{(P\_6,4.110)} uttarārtham eva tarhi . \\
\noindent \textbf{(P\_6,4.110)} śnābhyastayoḥ ātaḥ iti . \\
\noindent \textbf{(P\_6,4.110)} śnā sārvadhātuke eva . \\
\noindent \textbf{(P\_6,4.110)} abhyastam api ākārāntam ārdhadhātuke na asti . \\
\noindent \textbf{(P\_6,4.110)} nanu ca idam asti : apsu yāyāvaraḥ pravapeta piṇḍān iti . \\
\noindent \textbf{(P\_6,4.110)} na etat ākārāntam . \\
\noindent \textbf{(P\_6,4.110)} yakārāntam etat . \\
\noindent \textbf{(P\_6,4.110)} uttarārtham eva tarhi . \\
\noindent \textbf{(P\_6,4.110)} ī hali aghoḥ iti . \\
\noindent \textbf{(P\_6,4.110)} tatra api śnābhyastayoḥ iti eva . \\
\noindent \textbf{(P\_6,4.110)} ataḥ api uttarārtham eva tarhi . \\
\noindent \textbf{(P\_6,4.110)} id daridrasya iti . \\
\noindent \textbf{(P\_6,4.110)} vakṣyati etat : daridrāteḥ ārdhadhātuke lopaḥ siddhaḥ ca pratyayavidhau iti . \\
\noindent \textbf{(P\_6,4.110)} ataḥ api uttarārtham . \\
\noindent \textbf{(P\_6,4.110)} bhiyaḥ anyatarasyām . \\
\noindent \textbf{(P\_6,4.110)} abhyastasya iti eva . \\
\noindent \textbf{(P\_6,4.110)} ataḥ api uttarārtham eva . \\
\noindent \textbf{(P\_6,4.110)} jahāteḥ ca . \\
\noindent \textbf{(P\_6,4.110)} abhyastasya iti eva . \\
\noindent \textbf{(P\_6,4.110)} ataḥ api uttarārtham . \\
\noindent \textbf{(P\_6,4.110)} ā ca hau . \\
\noindent \textbf{(P\_6,4.110)} hau iti ucyate . \\
\noindent \textbf{(P\_6,4.110)} abhyastasya iti eva . \\
\noindent \textbf{(P\_6,4.110)} ataḥ api uttarārtham . \\
\noindent \textbf{(P\_6,4.110)} lopaḥ yi . \\
\noindent \textbf{(P\_6,4.110)} abhyastasya iti eva . \\
\noindent \textbf{(P\_6,4.110)} ataḥ api uttarārtham . \\
\noindent \textbf{(P\_6,4.110)} ghavsoḥ et hau abhyāsalopaḥ ca iti . \\
\noindent \textbf{(P\_6,4.110)} hau iti ucyate . \\
\noindent \textbf{(P\_6,4.110)} tat eva tarhi prayojanam . \\
\noindent \textbf{(P\_6,4.110)} śnasoḥ allopaḥ iti . \\
\noindent \textbf{(P\_6,4.110)} nanu ca uktam śnam sārvadhātuke eva . \\
\noindent \textbf{(P\_6,4.110)} asteḥ api ārdhadhātuke bhūbhāvena bhavitavyam iti . \\
\noindent \textbf{(P\_6,4.110)} anuprayoge tu bhuvā astyabādhanam smaranti kartuḥ vacanāt manīṣiṇaḥ . \\
\noindent \textbf{(P\_6,4.110)} anuprayoge tu bhuvā asteḥ abādhanam iṣyate : īhām āsa , īhām āsatuḥ , īhām āsuḥ iti . \\
\noindent \textbf{(P\_6,4.110)} kim ca syāt yadi atra lopaḥ syāt . \\
\noindent \textbf{(P\_6,4.110)} lope dvirvacanāsiddhiḥ . \\
\noindent \textbf{(P\_6,4.110)} lope kṛte anackatvāt dvirvacanam syāt . \\
\noindent \textbf{(P\_6,4.110)} sthānivadbhādāt bhaviṣyati . \\
\noindent \textbf{(P\_6,4.110)} sthānivat iti cet kṛte bhavet dvitve . \\
\noindent \textbf{(P\_6,4.110)} kṛte dvitve lopaḥ prāpnoti . \\
\noindent \textbf{(P\_6,4.110)} asti tarhi parasya lopaḥ . \\
\noindent \textbf{(P\_6,4.110)} abhyāsasya yaḥ akāraḥ tasya dīrghatvam bhaviṣyati . \\
\noindent \textbf{(P\_6,4.110)} na evam sidhyati kasmāt pratyaṅgatvāt bhavet hi pararūpam . \\
\noindent \textbf{(P\_6,4.110)} na evam sidhyati . \\
\noindent \textbf{(P\_6,4.110)} kasmāt . \\
\noindent \textbf{(P\_6,4.110)} pratyaṅgatvāt pararūpam prāpnoti . \\
\noindent \textbf{(P\_6,4.110)} tasmin ca kṛte lopaḥ . \\
\noindent \textbf{(P\_6,4.110)} pararūpe ca kṛte lopaḥ prāpnoti . \\
\noindent \textbf{(P\_6,4.110)} dīrghatvam bādhakam bhavet tatra . \\
\noindent \textbf{(P\_6,4.110)} ataḥ ādeḥ iti dīrghatvam bādhakam bhaviṣyati . \\
\noindent \textbf{(P\_6,4.110)} idam tarhi prayojanam . \\
\noindent \textbf{(P\_6,4.110)} sārvadhātuke bhūtapūrvamātre api yathā syāt : kuru iti . \\
\noindent \textbf{(P\_6,4.111)} atha atra taparakaraṇam kimartham . \\
\noindent \textbf{(P\_6,4.111)} iha mā bhūt : āstām , āsan . \\
\noindent \textbf{(P\_6,4.111)} na etat asti prayojanam . \\
\noindent \textbf{(P\_6,4.111)} āṭaḥ asiddhatvāt na bhaviṣyati . \\
\noindent \textbf{(P\_6,4.114)} daridrāteḥ ārdhadhātuke lopaḥ . \\
\noindent \textbf{(P\_6,4.114)} daridrāteḥ ārdhadhātuke lopaḥ vaktavyaḥ . \\
\noindent \textbf{(P\_6,4.114)} siddhaḥ ca pratyayavidhau . \\
\noindent \textbf{(P\_6,4.114)} saḥ ca siddhaḥ pratyayavidhau . \\
\noindent \textbf{(P\_6,4.114)} kim prayojanam . \\
\noindent \textbf{(P\_6,4.114)} daridrāti iti daridraḥ . \\
\noindent \textbf{(P\_6,4.114)} ākārāntalakṣaṇaḥ pratyayavidhiḥ mā bhūt iti . na daridrāyake lopaḥ daridrāṇe ca na iṣyate . \\
\noindent \textbf{(P\_6,4.114)} didaridrāsasti iti eke didaridriṣati iti vā . \\
\noindent \textbf{(P\_6,4.114)} vā adyatanyām . adyatanyām vā iti vaktavyam . \\
\noindent \textbf{(P\_6,4.114)} adaridrīt , adaridrāsīt . \\
\noindent \textbf{(P\_6,4.120.1)} ṇakāraṣakārādeśādeḥ ettvavacanam liṭi . \\
\noindent \textbf{(P\_6,4.120.1)} ṇakāraṣakārādeśādeḥ ettvam liṭi vaktavyam . \\
\noindent \textbf{(P\_6,4.120.1)} nematuḥ , nemuḥ , sehe, sehāte , sehire . \\
\noindent \textbf{(P\_6,4.120.1)} kim punaḥ kāraṇam na sidhyati . \\
\noindent \textbf{(P\_6,4.120.1)} anādeśādeḥ iti lpratiṣedhaḥ prāpnoti . \\
\noindent \textbf{(P\_6,4.120.1)} tat tarhi vaktavyam . \\
\noindent \textbf{(P\_6,4.120.1)} na vaktavyam . \\
\noindent \textbf{(P\_6,4.120.1)} liṭā atra ādeśādim viśeṣayiṣyāmaḥ . \\
\noindent \textbf{(P\_6,4.120.1)} liṭi yaḥ ādeśādiḥ tadādeḥ na iti . \\
\noindent \textbf{(P\_6,4.120.1)} asti anyat liḍgrahaṇasya prayojanam . \\
\noindent \textbf{(P\_6,4.120.1)} kim . \\
\noindent \textbf{(P\_6,4.120.1)} iha mā bhūt : paktā paktum . \\
\noindent \textbf{(P\_6,4.120.1)} na etat asti prayojanam . \\
\noindent \textbf{(P\_6,4.120.1)} kṅiti iti vartate . \\
\noindent \textbf{(P\_6,4.120.1)} evam api pakvaḥ pakvavān iti atra prāpnoti . \\
\noindent \textbf{(P\_6,4.120.1)} abhyāsalopasanniyogena ettvam ucyate na ca atra abhyāsalopasam paśyāmaḥ . \\
\noindent \textbf{(P\_6,4.120.1)} evam api pāpacyate atra prāpnoti . \\
\noindent \textbf{(P\_6,4.120.1)} dīrghatvam atra bādhakam bhaviṣyati . \\
\noindent \textbf{(P\_6,4.120.1)} na aprāpte abhyāsavikāre ettam arabhyate . \\
\noindent \textbf{(P\_6,4.120.1)} tat yatha anyān abhyāsavikārān bādhate evam dīrghatvam api bādheta . \\
\noindent \textbf{(P\_6,4.120.1)} satyam evam etat . \\
\noindent \textbf{(P\_6,4.120.1)} abhyāsavikāreṣu tu jyeṣṭhamadhyamakanīyāṃsaḥ prakārāḥ bhavanti . \\
\noindent \textbf{(P\_6,4.120.1)} tatra hrasvahalādiśeṣau utsargau . \\
\noindent \textbf{(P\_6,4.120.1)} tayoḥ dīrghatvam apavādaḥ ettvam ca . \\
\noindent \textbf{(P\_6,4.120.1)} apavādavipratiṣedhāt dīrghatvam bhaviṣyati . \\
\noindent \textbf{(P\_6,4.120.1)} iha tarhi babhaṇatuḥ , babhaṇuḥ iti abhyāsādeśasya asiddhatvāt ettvam prāpnoti . \\
\noindent \textbf{(P\_6,4.120.1)} phalibhajigrahaṇam tu jñāpakam abhyāsādeśasiddhatvasya . \\
\noindent \textbf{(P\_6,4.120.1)} yat ayam phalibhajyoḥ grahaṇam karoti tat jñāpayati ācāryaḥ siddhaḥ abhyāsādeśaḥ ettve iti . \\
\noindent \textbf{(P\_6,4.120.1)} yadi evam prathamatṛtīyādīnām ādeśāditvāt ettvābhāvaḥ . \\
\noindent \textbf{(P\_6,4.120.1)} prathamatṛtīyādīnām tarhi ādeśāditvāt ettvam na prāpnoti . \\
\noindent \textbf{(P\_6,4.120.1)} pecatuḥ , pecuḥ , debhatuḥ , debhuḥ . \\
\noindent \textbf{(P\_6,4.120.1)} na vā śasidadyoḥ pratiṣedhaḥ jñāpakaḥ rūpābhede ettvavijñānasya . \\
\noindent \textbf{(P\_6,4.120.1)} na vā eṣaḥ doṣaḥ . \\
\noindent \textbf{(P\_6,4.120.1)} kim kāraṇam . \\
\noindent \textbf{(P\_6,4.120.1)} śasidadyoḥ pratiṣedhaḥ jñāpakaḥ rūpābhede ettvavijñānasya . \\
\noindent \textbf{(P\_6,4.120.1)} yat ayam śasidadyoḥ pratiṣedham śāsti tat jñāpayati ācāryaḥ rūpābhedena yaḥ ādeśādayaḥ na teṣām pratiṣedhaḥ iti . \\
\noindent \textbf{(P\_6,4.120.2)} dambhaḥ ettvam . \\
\noindent \textbf{(P\_6,4.120.2)} dambhaḥ ettvam vaktavyam . \\
\noindent \textbf{(P\_6,4.120.2)} debhatuḥ , debhuḥ . \\
\noindent \textbf{(P\_6,4.120.2)} kim punaḥ kāraṇam na sidhyati . \\
\noindent \textbf{(P\_6,4.120.2)} nalopasya asiddhatvāt . asiddhaḥ nalopaḥ . \\
\noindent \textbf{(P\_6,4.120.2)} tasya asiddhatvāt ettvam na prāpnoti . \\
\noindent \textbf{(P\_6,4.120.2)} naśimanyoḥ aliṭi ettvam . \\
\noindent \textbf{(P\_6,4.120.2)} naśimanyoḥ aliṭi ettvam vaktavyam . \\
\noindent \textbf{(P\_6,4.120.2)} chandasi amipacyoḥ api . \\
\noindent \textbf{(P\_6,4.120.2)} chandasi amipacyoḥ api iti vaktavyam . \\
\noindent \textbf{(P\_6,4.120.2)} kim prayojanam . \\
\noindent \textbf{(P\_6,4.120.2)} aneśam menakā iti etat vyemānam liṅi peciran . \\
\noindent \textbf{(P\_6,4.120.2)} yaj āyeje vap āvepe dambhaḥ ettvam alakṣaṇam . \\
\noindent \textbf{(P\_6,4.120.2)} asiddhatvāt nalopasya dambhaḥ ettvam na sidhyati . \\
\noindent \textbf{(P\_6,4.120.2)} śnasoḥ attve takāreṇa jñāpyate tu ettvaśāsanam . \\
\noindent \textbf{(P\_6,4.120.2)} anityaḥ ayam vidhiḥ iti . \\
\noindent \textbf{(P\_6,4.121)} thalgrahaṇam kimartham . \\
\noindent \textbf{(P\_6,4.121)} thalgrahaṇam akṅidartham . \\
\noindent \textbf{(P\_6,4.121)} thalgrahaṇam kriyate akṅidartham . \\
\noindent \textbf{(P\_6,4.121)} akṅiti ettvam yathā syāt . \\
\noindent \textbf{(P\_6,4.121)} pecitha śekitha . \\
\noindent \textbf{(P\_6,4.121)} na etat asti prayojanam . \\
\noindent \textbf{(P\_6,4.121)} seḍgrahaṇam eve atra akṅidartham bhaviṣyati . \\
\noindent \textbf{(P\_6,4.121)} idam tarhi prayojanam . \\
\noindent \textbf{(P\_6,4.121)} samuccayaḥ yathā vijñāyeta . \\
\noindent \textbf{(P\_6,4.121)} thali ca seṭi kṅiti ca seṭi iti . \\
\noindent \textbf{(P\_6,4.121)} kim prayojanam . \\
\noindent \textbf{(P\_6,4.121)} peciva pecima. tatra pacādibhyaḥ iḍvacanam iti vakṣyati . \\
\noindent \textbf{(P\_6,4.121)} tat na vaktavyam bhavati . \\
\noindent \textbf{(P\_6,4.121)} iha kasmāt na bhavati : lulavitha . \\
\noindent \textbf{(P\_6,4.121)} guṇasya pratiṣedhāt . \\
\noindent \textbf{(P\_6,4.121)} iha api tarhi na prāpnoti : pecitha śekitha . \\
\noindent \textbf{(P\_6,4.121)} guṇasya yaḥ akāraḥ iti evam etat vijñāsyate . \\
\noindent \textbf{(P\_6,4.121)} evam api śaśaritha , atra prāpnoti . \\
\noindent \textbf{(P\_6,4.121)} guṇasya eṣaḥ akāraḥ . \\
\noindent \textbf{(P\_6,4.121)} katham . \\
\noindent \textbf{(P\_6,4.121)} vṛddhiḥ bhavati guṇaḥ bhavati iti rephaśirāḥ guṇavṛddhisañjñakaḥ abhinirvartate . \\
\noindent \textbf{(P\_6,4.121)} atha vā ācāryapravṛttiḥ jñāpayati ne evañjātīyakānām ettvam bhavati iti yat ayam tṝphalabhajatrapaḥ ca iti tṝgrahaṇam karoti . \\
\noindent \textbf{(P\_6,4.123)} rādhādiṣu sthāninirdeśaḥ . \\
\noindent \textbf{(P\_6,4.123)} rādhādiṣu sthāninirdeśaḥ kartavyaḥ . \\
\noindent \textbf{(P\_6,4.123)} na kartavyaḥ . \\
\noindent \textbf{(P\_6,4.123)} ekahalmadhye iti vartate . \\
\noindent \textbf{(P\_6,4.123)} yadi evam tresatuḥ , tresuḥ , ra śabdasya ettvam prāpnoti . \\
\noindent \textbf{(P\_6,4.123)} astu . \\
\noindent \textbf{(P\_6,4.123)} alaḥ antyasya vidhayaḥ bhavanti iti akārasya bhaviṣyati . \\
\noindent \textbf{(P\_6,4.123)} anarthake alaḥ antyavidhiḥ na iti evam na prāpnoti . \\
\noindent \textbf{(P\_6,4.123)} na etasyāḥ paribhāṣāyāḥ santi prayojanāni . \\
\noindent \textbf{(P\_6,4.123)} atha vā ataḥ iti vartate . \\
\noindent \textbf{(P\_6,4.123)} evam api rādheḥ na prāpnoti . \\
\noindent \textbf{(P\_6,4.123)} ākāragrahaṇam api prakṛtam anuvartate . \\
\noindent \textbf{(P\_6,4.123)} kva prakṛtam . \\
\noindent \textbf{(P\_6,4.123)} śnābhyāstayoḥ ātaḥ iti . \\
\noindent \textbf{(P\_6,4.123)} atha vā śnasoḥ allopaḥ iti atra taparakaraṇam pratyākhyāyate . \\
\noindent \textbf{(P\_6,4.123)} tat prakṛtam iha anuvartiṣyate . \\
\noindent \textbf{(P\_6,4.123)} yadi tat anuvartate ataḥ ekahalmadhye anādeśādeḥ liṭi asya ca iti avarṇamātrasya ettvam prāpnoti . \\
\noindent \textbf{(P\_6,4.123)} babādhe . \\
\noindent \textbf{(P\_6,4.123)} akāreṇa tapareṇa avarṇam viśeṣayiṣyāmaḥ . \\
\noindent \textbf{(P\_6,4.123)} asya ātaḥ iti . \\
\noindent \textbf{(P\_6,4.123)} iha idānīm asya iti anuvartate ataḥ iti nivṛttam . \\
\noindent \textbf{(P\_6,4.127-128)} KA\_III,220.11-18 Ro\_IV,772-773 {1/11}     arvaṇas tṛ maghonaḥ ca na śiṣyam chāndasam hi tat . \\
\noindent \textbf{(P\_6,4.127-128)} KA\_III,220.11-18 Ro\_IV,772-773 {2/11}     arvaṇas tṛ maghonaḥ ca na śiṣyam . \\
\noindent \textbf{(P\_6,4.127-128)} KA\_III,220.11-18 Ro\_IV,772-773 {3/11}     kim kāraṇam . \\
\noindent \textbf{(P\_6,4.127-128)} KA\_III,220.11-18 Ro\_IV,772-773 {4/11}     chāndasam hi tat . \\
\noindent \textbf{(P\_6,4.127-128)} KA\_III,220.11-18 Ro\_IV,772-773 {5/11}     dṛṣṭānuvidhiḥ chandasi bhavati . \\
\noindent \textbf{(P\_6,4.127-128)} KA\_III,220.11-18 Ro\_IV,772-773 {6/11}     matubvanyoḥ vidhānāt ca . \\
\noindent \textbf{(P\_6,4.127-128)} KA\_III,220.11-18 Ro\_IV,772-773 {7/11}     matubvanī khalu api chandasi vidhīyete . \\
\noindent \textbf{(P\_6,4.127-128)} KA\_III,220.11-18 Ro\_IV,772-773 {8/11}     chandasi ubhayadarśanāt . \\
\noindent \textbf{(P\_6,4.127-128)} KA\_III,220.11-18 Ro\_IV,772-773 {9/11}     ubhayam khalu api chandasi dṛśyate . \\
\noindent \textbf{(P\_6,4.127-128)} KA\_III,220.11-18 Ro\_IV,772-773 {10/11}     imāni arvaṇaḥ padāni . \\
\noindent \textbf{(P\_6,4.127-128)} KA\_III,220.11-18 Ro\_IV,772-773 {11/11}     anarvaṇam vṛṣabham mandrajihvam . \\
\noindent \textbf{(P\_6,4.130)} pādaḥ upadhāhrasvatvam . \\
\noindent \textbf{(P\_6,4.130)} pādaḥ upadhāhrasvatvam vaktavyam . \\
\noindent \textbf{(P\_6,4.130)} dvipadaḥ paśya . \\
\noindent \textbf{(P\_6,4.130)} ādeśe hi sarvādeśaprasaṅgaḥ . \\
\noindent \textbf{(P\_6,4.130)} ādeśe hi sati sarvādeśaḥ prasajyeta . \\
\noindent \textbf{(P\_6,4.130)} sarvasya dvipācchabdasya tripācchabdasya ca pacchabdādeśaḥ prasajyeta yena vidhiḥ tadantasya iti . \\
\noindent \textbf{(P\_6,4.130)} tat tarhi vaktavyam . \\
\noindent \textbf{(P\_6,4.130)} na vā nirdiśyamānasya ādeśatvāt . \\
\noindent \textbf{(P\_6,4.130)} na vā vaktavyam . \\
\noindent \textbf{(P\_6,4.130)} kim kāraṇam . \\
\noindent \textbf{(P\_6,4.130)} nirdiśyamānasya ādeśāḥ bhavanti iti eṣā paribhāṣā kartavyā . \\
\noindent \textbf{(P\_6,4.130)} kaḥ punaḥ atra viśeṣaḥ eṣā vā paribhāṣā kriyeta upadhāhrasvatvam vā ucyeta . \\
\noindent \textbf{(P\_6,4.130)} avaśyam eṣā paribhāṣā kartavyā . \\
\noindent \textbf{(P\_6,4.130)} bahūni etasyāḥ paribhāṣāyāḥ prayojanāni kāni . \\
\noindent \textbf{(P\_6,4.130)} prayojanam suptiṅādeśe . \\
\noindent \textbf{(P\_6,4.130)} sup . \\
\noindent \textbf{(P\_6,4.130)} kumāryām , kośoryām , khaṭvāyām , mālāyām , tasyām , yasyām . \\
\noindent \textbf{(P\_6,4.130)} āḍyāṭsyāṭsu kṛteṣu sāḍyāṭsyāṭkasya ām prāpnoti . \\
\noindent \textbf{(P\_6,4.130)} nirdiśyamānasya ādeśāḥ bhavanti iti na doṣaḥ bhavati . \\
\noindent \textbf{(P\_6,4.130)} idam iha sampradhāryam . \\
\noindent \textbf{(P\_6,4.130)} āḍyāṭsyāṭaḥ kriyantām ām iti kim atra kartavyam . \\
\noindent \textbf{(P\_6,4.130)} paratvāt ām . \\
\noindent \textbf{(P\_6,4.130)} nityāḥ āḍyāṭsyāṭaḥ . \\
\noindent \textbf{(P\_6,4.130)} kṛte api āmi prapnuvanti akṛte api . \\
\noindent \textbf{(P\_6,4.130)} anityāḥ āḍyāṭsyāṭaḥ . \\
\noindent \textbf{(P\_6,4.130)} anyasya kṛte āmi prapnuvanti anyasya akṛte śabdāntarasya ca prāpnuvantaḥ anityāḥ bhavanti . \\
\noindent \textbf{(P\_6,4.130)} ubhayoḥ anityayoḥ paratvāt ām . \\
\noindent \textbf{(P\_6,4.130)} idam tarhi . \\
\noindent \textbf{(P\_6,4.130)} tasyai yasyai . \\
\noindent \textbf{(P\_6,4.130)} syāṭi kṛte sasyāṭkasya smaibhāvaḥ prāpnoti . \\
\noindent \textbf{(P\_6,4.130)} nirdiśyamānasya ādeśāḥ bhavanti iti na doṣaḥ bhavati . \\
\noindent \textbf{(P\_6,4.130)} yaḥ tarhi nirdiśyate tasya kasmāt na bhavati . \\
\noindent \textbf{(P\_6,4.130)} syāṭā vyavahitatvāt . \\
\noindent \textbf{(P\_6,4.130)} sup . \\
\noindent \textbf{(P\_6,4.130)} tiṅ. aruditām aruditam arudita iti . \\
\noindent \textbf{(P\_6,4.130)} iṭi kṛte seṭkasya tāmtamtāmādeśāḥ prāpnuvanti . \\
\noindent \textbf{(P\_6,4.130)} nirdiśyamānasya ādeśāḥ bhavanti iti na doṣaḥ bhavati . \\
\noindent \textbf{(P\_6,4.130)} idam iha sampradhāryam . \\
\noindent \textbf{(P\_6,4.130)} iṭ kriyatām tāmtamtāmaḥ iti kim atra kartavyam . \\
\noindent \textbf{(P\_6,4.130)} paratvāt iḍāgamaḥ . \\
\noindent \textbf{(P\_6,4.130)} antaraṅgāḥ tāmtamtāmaḥ . \\
\noindent \textbf{(P\_6,4.130)} idam tarhi kriyāstām , kriyāstam , kriyāsta . \\
\noindent \textbf{(P\_6,4.130)} yāsuṭi kṛte sayāsuṭkasya tāmtamtāmādeśāḥ prāpnuvanti . \\
\noindent \textbf{(P\_6,4.130)} nirdiśyamānasya ādeśāḥ bhavanti iti na doṣaḥ bhavati . \\
\noindent \textbf{(P\_6,4.130)} lyabbhāve ca . \\
\noindent \textbf{(P\_6,4.130)} lyabbhāve ca prayojanam . \\
\noindent \textbf{(P\_6,4.130)} prakṛtya prahṛtya . \\
\noindent \textbf{(P\_6,4.130)} ktvāntasya lyap prāpnoti . \\
\noindent \textbf{(P\_6,4.130)} nirdiśyamānasya ādeśāḥ bhavanti iti na doṣaḥ bhavati . \\
\noindent \textbf{(P\_6,4.130)} tricaturyuṣmadasmattyadādivikāreṣu ca . \\
\noindent \textbf{(P\_6,4.130)} tricaturyuṣmadasmattyadādivikāreṣu ca prayojanam . \\
\noindent \textbf{(P\_6,4.130)} atitisraḥ , aticatasraḥ . \\
\noindent \textbf{(P\_6,4.130)} tricaturantasya tisṛcatasṛbhāvaḥ prāpnoti . \\
\noindent \textbf{(P\_6,4.130)} nirdiśyamānasya ādeśāḥ bhavanti iti na doṣaḥ bhavati . \\
\noindent \textbf{(P\_6,4.130)} yuṣmat , asmat . \\
\noindent \textbf{(P\_6,4.130)} atiyūyam ativayam . \\
\noindent \textbf{(P\_6,4.130)} yuṣmadasmadantasya yūyavayau prāpnutaḥ . \\
\noindent \textbf{(P\_6,4.130)} nirdiśyamānasya ādeśāḥ bhavanti iti na doṣaḥ bhavati . \\
\noindent \textbf{(P\_6,4.130)} tyadādivikāra . \\
\noindent \textbf{(P\_6,4.130)} atisyaḥ , uttamasyaḥ , atyasau , uttamāsau . \\
\noindent \textbf{(P\_6,4.130)} tyadādyantasya tyadādivikārāḥ prāpnuvanti . \\
\noindent \textbf{(P\_6,4.130)} kimantasya kādeśaḥ prāpnoti . \\
\noindent \textbf{(P\_6,4.130)} atikaḥ , paramakaḥ . \\
\noindent \textbf{(P\_6,4.130)} nirdiśyamānasya ādeśāḥ bhavanti iti na doṣaḥ bhavati . \\
\noindent \textbf{(P\_6,4.130)} udaḥ pūrvatve . \\
\noindent \textbf{(P\_6,4.130)} udaḥ pūrvatve prayojanam . \\
\noindent \textbf{(P\_6,4.130)} udasthātām . \\
\noindent \textbf{(P\_6,4.130)} aṭi kṛte sāṭkasya pūrvasavarṇaḥ prāpnoti udaḥ sthāstambhoḥ iti . \\
\noindent \textbf{(P\_6,4.130)} nirdiśyamānasya ādeśāḥ bhavanti iti na doṣaḥ bhavati . \\
\noindent \textbf{(P\_6,4.130)} yaḥ tarhi nirdiśyate tasya kasmāt na bhavati . \\
\noindent \textbf{(P\_6,4.130)} aṭā vyavahitatvāt . \\
\noindent \textbf{(P\_6,4.130)} sā tarhi paribhāṣā kartavyā . \\
\noindent \textbf{(P\_6,4.130)} na kartavyā . \\
\noindent \textbf{(P\_6,4.130)} uktam ṣaṣṭhī sthāneyogā iti etasya yogasya vacane prayojanam ṣaṣṭhyantam sthānena yathā yujyeta yataḥ ṣaṣṭhī uccāritā iti . \\
\noindent \textbf{(P\_6,4.132)} ūṭ ādiḥ kāsmāt na bhavati . \\
\noindent \textbf{(P\_6,4.132)} ādiḥ ṭit bhavati iti ādiḥ prāpnoti . \\
\noindent \textbf{(P\_6,4.132)} samprasāraṇam iti anena yaṇaḥ sthānam hriyate . \\
\noindent \textbf{(P\_6,4.132)} yadi evam vāhaḥ ūḍvacanānarthakyam samprasāraṇena kṛtatvāt . \\
\noindent \textbf{(P\_6,4.132)} vāhaḥ ūḍvacanam anarthakam . \\
\noindent \textbf{(P\_6,4.132)} kim kāraṇam . \\
\noindent \textbf{(P\_6,4.132)} samprasāraṇena kṛtatvāt . \\
\noindent \textbf{(P\_6,4.132)} samprasāraṇena eva siddham . \\
\noindent \textbf{(P\_6,4.132)} kā rūpasiddhiḥ . \\
\noindent \textbf{(P\_6,4.132)} praṣṭhauhaḥ paśya . \\
\noindent \textbf{(P\_6,4.132)} guṇaḥ pratyayalakṣaṇatvāt . \\
\noindent \textbf{(P\_6,4.132)} pratyayalakṣaṇena guṇaḥ bhaviṣyati . \\
\noindent \textbf{(P\_6,4.132)} ejgrahaṇāt vṛddhiḥ . \\
\noindent \textbf{(P\_6,4.132)} ejgrahaṇāt vṛddhiḥ bhaviṣyati . \\
\noindent \textbf{(P\_6,4.132)} evam tarhi siddhe sati yat vāhaḥ ūṭham śāsti śāsti tat jñāpayati ācāryaḥ bhavati eṣā paribhāṣā asiddham bahiraṅgalakṣaṇam antaraṅgalakṣaṇe iti . \\
\noindent \textbf{(P\_6,4.132)} kim etasya jñāpane prayojanam . \\
\noindent \textbf{(P\_6,4.132)} pacāva idam , pacāma idam . \\
\noindent \textbf{(P\_6,4.132)} asiddhatvāt bahiraṅgalakṣaṇasya āt guṇasya antaraṅgalakṣaṇam aittvam na bhavati iti . \\
\noindent \textbf{(P\_6,4.133)} śvādīnām prasāraṇe nakārāntagrahaṇam anakārāntapratiṣedhārtham . \\
\noindent \textbf{(P\_6,4.133)} śvādīnām prasāraṇe nakārāntagrahaṇam kartavyam . \\
\noindent \textbf{(P\_6,4.133)} kim prayojanam . \\
\noindent \textbf{(P\_6,4.133)} anakārāntapratiṣedhārtham . \\
\noindent \textbf{(P\_6,4.133)} anakārāntasya mā bhūt . \\
\noindent \textbf{(P\_6,4.133)} mabhavā maghavate . \\
\noindent \textbf{(P\_6,4.133)} tathā prātipadikagrahaṇe liṅgaviśiṣṭasya api grahaṇam bhavati iti yathā iha bhavati , yūnaḥ paśye iti evam yuvatīḥ lpaśya iti atra api syāt iti . \\
\noindent \textbf{(P\_6,4.133)} yat tāvat ucyate nakārāntagrahaṇam kartavyam . \\
\noindent \textbf{(P\_6,4.133)} na kartavyam . \\
\noindent \textbf{(P\_6,4.133)} uktam vā . \\
\noindent \textbf{(P\_6,4.133)} kim uktam . \\
\noindent \textbf{(P\_6,4.133)} arvaṇas tṛ maghonaḥ ca na śiṣyam chāndasam hi tat iti . \\
\noindent \textbf{(P\_6,4.133)} yat api ucyate prātipadikagrahaṇe liṅgaviśiṣṭasya api grahaṇam bhavati iti yathā iha bhavati , yūnaḥ paśye iti evam yuvatīḥ lpaśya iti atra api syāt iti liṅgaviśiṣṭagrahaṇe ca uktam . \\
\noindent \textbf{(P\_6,4.133)} kim uktam . \\
\noindent \textbf{(P\_6,4.133)} na vā vibhaktau liṅgaviśiṣṭāgrahaṇāt iti . \\
\noindent \textbf{(P\_6,4.133)} atha vā upariṣṭāt yogavibhāgaḥ kariṣyate . \\
\noindent \textbf{(P\_6,4.133)} śvayuvamaghonām ataddhite . \\
\noindent \textbf{(P\_6,4.133)} tataḥ allopaḥ . \\
\noindent \textbf{(P\_6,4.133)} akārasya ca lopaḥ bhavati . \\
\noindent \textbf{(P\_6,4.133)} tataḥ anaḥ . \\
\noindent \textbf{(P\_6,4.133)} anaḥ iti ubhayoḥ śeṣaḥ . \\
\noindent \textbf{(P\_6,4.135)} atha kim idam ṣapūrvādīnām punarvacanam allopārtham āhosvit niyamārtham . \\
\noindent \textbf{(P\_6,4.135)} katha ca allopārtham syāt katham vā niyamārtham . \\
\noindent \textbf{(P\_6,4.135)} yadi aviśeṣeṇa allopaṭilopayoḥ saḥ prakṛtibhāvaḥ tataḥ allopārtham . \\
\noindent \textbf{(P\_6,4.135)} atha hi aṇi ṭilopasya eva prakṛtibhāvaḥ tataḥ niyamārtham . \\
\noindent \textbf{(P\_6,4.135)} ṣapūrvādīnām punarvacanam allopārtham . \\
\noindent \textbf{(P\_6,4.135)} ṣapūrvādīnām punarvacanam kriyate allopārtham . \\
\noindent \textbf{(P\_6,4.135)} aviśeṣeṇa allopaṭilopayoḥ saḥ prakṛtibhāvaḥ . \\
\noindent \textbf{(P\_6,4.135)} avadhāraṇe hi anyatra prakṛtibhāve upadhālopaprasaṅgaḥ . avadhāraṇe hi sati anyatra prakṛtibhāve upadhālopaḥ prasajyeta . \\
\noindent \textbf{(P\_6,4.135)} katham . \\
\noindent \textbf{(P\_6,4.135)} yadi tāvat evam niyamaḥ syāt ṣapūrvādīnām eva aṇi iti bhavet iha niyamāt na syāt sāmanaḥ , vaimanaḥ iti . \\
\noindent \textbf{(P\_6,4.135)} tākṣaṇyaḥ iti prāpnoti . \\
\noindent \textbf{(P\_6,4.135)} atha api evam niyamaḥ syāt ṣapūrvādīnām aṇi eva iti evam api bhavet iha niyamāt na syāt tākṣaṇyaḥ iti , sāmanaḥ , vaimanaḥ iti tu prāpnoti . \\
\noindent \textbf{(P\_6,4.135)} atha api ubhayataḥ niyamaḥ syāt ṣapūrvādīnām eva aṇi , aṇi eva ṣapūrvādīnām iti evam api sāmanyaḥ , vemanyaḥ iti prāpnoti . \\
\noindent \textbf{(P\_6,4.135)} tasmāt suṣthu ucyate ṣapūrvādīnām punarvacanam allopārtham . \\
\noindent \textbf{(P\_6,4.135)} avadhāraṇe hi anyatra prakṛtibhāve upadhālopaprasaṅgaḥ iti . \\
\noindent \textbf{(P\_6,4.140)} ātaḥ anāpaḥ . \\
\noindent \textbf{(P\_6,4.140)} ātaḥ anāpaḥ iti vaktavyam . \\
\noindent \textbf{(P\_6,4.140)} iha api yathā syāt : samāse anañpūrve ktvaḥ lyap iti . \\
\noindent \textbf{(P\_6,4.140)} anāpaḥ iti kimartham . \\
\noindent \textbf{(P\_6,4.140)} khaṭvāyām , mālāyām . \\
\noindent \textbf{(P\_6,4.140)} yadi anāpaḥ iti ucyate katham ktvāyām . \\
\noindent \textbf{(P\_6,4.140)} nipātanāt etat siddham . \\
\noindent \textbf{(P\_6,4.140)} kim nipātanam . \\
\noindent \textbf{(P\_6,4.140)} ktvāyām va pratiṣedhaḥ iti . \\
\noindent \textbf{(P\_6,4.140)} yadi evam na arthaḥ anāpaḥ iti anena . \\
\noindent \textbf{(P\_6,4.140)} katham samāse anañpūrve ktvaḥ lyap iti . \\
\noindent \textbf{(P\_6,4.140)} nipātanāt etat siddham . \\
\noindent \textbf{(P\_6,4.140)} katham halaḥ śnaḥ śānac hau iti . \\
\noindent \textbf{(P\_6,4.140)} etat api nipātanāt siddham . \\
\noindent \textbf{(P\_6,4.140)} atha vā yogavibhāgaḥ kariṣyate . \\
\noindent \textbf{(P\_6,4.140)} ātaḥ : ākāralopaḥ bhavati . \\
\noindent \textbf{(P\_6,4.140)} tataḥ dhātoḥ : dhātoḥ ca ākārasya lopaḥ bhavati iti . \\
\noindent \textbf{(P\_6,4.141)} mantreṣu ātmanaḥ pratyayamātraprasaṅgaḥ . \\
\noindent \textbf{(P\_6,4.141)} mantreṣu ātmanaḥ pratyayamātre lopaḥ prasaṅktavyaḥ . \\
\noindent \textbf{(P\_6,4.141)} iha api yathā syāt : tmanyā samañjan . \\
\noindent \textbf{(P\_6,4.141)} tmanoḥ antaḥ asthaḥ iti . \\
\noindent \textbf{(P\_6,4.141)} yadi pratyayamātre lopaḥ ucyate katham ātmanaḥ eva nirmimīṣva iti . \\
\noindent \textbf{(P\_6,4.141)} tasmāt na arthaḥ pratyayamātre lopena . \\
\noindent \textbf{(P\_6,4.141)} katham tmanyā samañjan . \\
\noindent \textbf{(P\_6,4.141)} tmanoḥ antaḥ asthaḥ iti . \\
\noindent \textbf{(P\_6,4.141)} chāndasatvāt siddham . \\
\noindent \textbf{(P\_6,4.141)} chāndasam etat . \\
\noindent \textbf{(P\_6,4.141)} dṛṣṭānuvidhiḥ chandasi bhavati . \\
\noindent \textbf{(P\_6,4.141)} ādigrahaṇānarthakyam ca ākāraprakaraṇāt . \\
\noindent \textbf{(P\_6,4.141)} ādigrahaṇam ca anarthakam . \\
\noindent \textbf{(P\_6,4.141)} kim kāraṇam . \\
\noindent \textbf{(P\_6,4.141)} ākāraprakaraṇāt . \\
\noindent \textbf{(P\_6,4.141)} ātaḥ iti vartate . \\
\noindent \textbf{(P\_6,4.142)} tigrahaṇam kimartham na viṃśateḥ ḍiti lopaḥ iti eva ucyeta . \\
\noindent \textbf{(P\_6,4.142)} na evam śakyam . \\
\noindent \textbf{(P\_6,4.142)} viṃśateḥ ḍiti lopaḥ iti ucyamāne antyasya prasajyeta . \\
\noindent \textbf{(P\_6,4.142)} siddhaḥ antyasya yasyeta lopena . \\
\noindent \textbf{(P\_6,4.142)} tatra ārambhasāmarthyāt tiśabdasya bhaviṣyati . \\
\noindent \textbf{(P\_6,4.142)} kutaḥ nu khalu etat ananyārthe ārambhe tiśabdasya bhaviṣyati na punaḥ aṅgasya iti . \\
\noindent \textbf{(P\_6,4.142)} tasmāt tigrahaṇam kartavyam . \\
\noindent \textbf{(P\_6,4.142)} atha kriyamāṇe api tigrahaṇe antyasya kasmāt na bhavati . \\
\noindent \textbf{(P\_6,4.142)} nirdiśyamānasya ādeśāḥ bhavanti iti evam na bhaviṣyati . \\
\noindent \textbf{(P\_6,4.143)} abhasya upasaṅkhyānam kartavyam . \\
\noindent \textbf{(P\_6,4.143)} iha api yathā syāt : upasarajaḥ , mandurajaḥ iti . \\
\noindent \textbf{(P\_6,4.143)} tat tarhi vaktavyam . \\
\noindent \textbf{(P\_6,4.143)} na vaktavyam . \\
\noindent \textbf{(P\_6,4.143)} katham upasarajaḥ , mandurajaḥ iti . \\
\noindent \textbf{(P\_6,4.143)} ḍiti abhasya api anubandhakaraṇasāmarthyāt . \\
\noindent \textbf{(P\_6,4.143)} abhasya api anubandhakaraṇasāmarthyāt bhaviṣyati . \\
\noindent \textbf{(P\_6,4.144)} nakārantasya ṭilope sabrahmacāripīṭhasarpikalāpikuthumitaitilijājalilāṅgaliśilāliśikhaṇḍisūkarasdmasuparvaṇām upasaṅkhyānam . \\
\noindent \textbf{(P\_6,4.144)} nakārantasya ṭilope sabrahmacārin pīṭhasarpin kalāpin kuthumin taitilin jājalin lāṅgalin śilālin śikhaṇḍin sūkarasdman suparvan iti eteṣām upasaṅkhyānam kartavyam . \\
\noindent \textbf{(P\_6,4.144)} sabrahmacārin . \\
\noindent \textbf{(P\_6,4.144)} sābrahmacārāḥ . \\
\noindent \textbf{(P\_6,4.144)} sabrahmacārin . \\
\noindent \textbf{(P\_6,4.144)} pīṭhasarpin . \\
\noindent \textbf{(P\_6,4.144)} paiṭhasarpāḥ . \\
\noindent \textbf{(P\_6,4.144)} pīṭhasarpin . \\
\noindent \textbf{(P\_6,4.144)} kalāpin . \\
\noindent \textbf{(P\_6,4.144)} kālapāḥ . \\
\noindent \textbf{(P\_6,4.144)} kalāpin . \\
\noindent \textbf{(P\_6,4.144)} kuthumin . \\
\noindent \textbf{(P\_6,4.144)} kauthumāḥ . \\
\noindent \textbf{(P\_6,4.144)} kuthumin . \\
\noindent \textbf{(P\_6,4.144)} taitilin . \\
\noindent \textbf{(P\_6,4.144)} taitilāḥ . \\
\noindent \textbf{(P\_6,4.144)} taitilin . \\
\noindent \textbf{(P\_6,4.144)} jājalin . \\
\noindent \textbf{(P\_6,4.144)} jājalāḥ . \\
\noindent \textbf{(P\_6,4.144)} jājalin . \\
\noindent \textbf{(P\_6,4.144)} lāṅgalin . \\
\noindent \textbf{(P\_6,4.144)} lāṅgalāḥ . \\
\noindent \textbf{(P\_6,4.144)} lāṅgalin . \\
\noindent \textbf{(P\_6,4.144)} śilālin . \\
\noindent \textbf{(P\_6,4.144)} śailālāḥ . \\
\noindent \textbf{(P\_6,4.144)} śilālin . \\
\noindent \textbf{(P\_6,4.144)} śikhaṇḍin . \\
\noindent \textbf{(P\_6,4.144)} śaikhaṇḍāḥ . \\
\noindent \textbf{(P\_6,4.144)} śikhaṇḍin . \\
\noindent \textbf{(P\_6,4.144)} sūkarasdman . \\
\noindent \textbf{(P\_6,4.144)} saukarasadmāḥ . \\
\noindent \textbf{(P\_6,4.144)} sūkarasdman . \\
\noindent \textbf{(P\_6,4.144)} suparvan . \\
\noindent \textbf{(P\_6,4.144)} sauparvāḥ . \\
\noindent \textbf{(P\_6,4.144)} suparvan . \\
\noindent \textbf{(P\_6,4.144)} carmaṇaḥ kośe . \\
\noindent \textbf{(P\_6,4.144)} carmaṇaḥ kośe upasaṅkhyānam kartavyam . \\
\noindent \textbf{(P\_6,4.144)} cārmaḥ kośaḥ . \\
\noindent \textbf{(P\_6,4.144)} āsmanaḥ vikāre . \\
\noindent \textbf{(P\_6,4.144)} āsmanaḥ vikāre upasaṅkhyānam kartavyam . \\
\noindent \textbf{(P\_6,4.144)} aśmanaḥ vikāraḥ āśmaḥ . \\
\noindent \textbf{(P\_6,4.144)} śunaḥ saṅkoce . \\
\noindent \textbf{(P\_6,4.144)} śaunaḥ saṅkocaḥ . \\
\noindent \textbf{(P\_6,4.144)} avyayānām ca . avyayānām ca upasaṅkhyānam kartavyam . \\
\noindent \textbf{(P\_6,4.144)} kim prayojanam . \\
\noindent \textbf{(P\_6,4.144)} sāyampratikādyartham . sāyamprātikaḥ paunaḥpunikaḥ . \\
\noindent \textbf{(P\_6,4.144)} śāśvatike pratiṣedhaḥ vaktavyaḥ . \\
\noindent \textbf{(P\_6,4.144)} na vaktavyaḥ . \\
\noindent \textbf{(P\_6,4.144)} nipātanāt etat siddham . \\
\noindent \textbf{(P\_6,4.144)} kim nipātanam . \\
\noindent \textbf{(P\_6,4.144)} yeṣām ca virodhaḥ śāśvatikaḥ iti . \\
\noindent \textbf{(P\_6,4.144)} evam tarhi śāśvate pratiṣedhaḥ vaktavyaḥ . \\
\noindent \textbf{(P\_6,4.144)} śāśvatam . \\
\noindent \textbf{(P\_6,4.148.1)} ivarṇāntasya iti kim udāharaṇam . \\
\noindent \textbf{(P\_6,4.148.1)} he dākṣi dākṣyā dākṝyaḥ . \\
\noindent \textbf{(P\_6,4.148.1)} he dākṣi iti . \\
\noindent \textbf{(P\_6,4.148.1)} yadi lopaḥ na syāt parasya hrasvatve kṛte savarṇadīrghatvam prasjyeta . \\
\noindent \textbf{(P\_6,4.148.1)} dākṣyā iti . \\
\noindent \textbf{(P\_6,4.148.1)} yadi lopaḥ na syāt parasya yaṇādeśe kṛte pūrvasya śravaṇam prasajyeta . \\
\noindent \textbf{(P\_6,4.148.1)} dākṣeyaḥ iti . \\
\noindent \textbf{(P\_6,4.148.1)} yadi lopaḥ na syāt parasya lope kṛte pūrvasya śravaṇam prasajyeta . \\
\noindent \textbf{(P\_6,4.148.1)} na etāni santi prayojanāni . \\
\noindent \textbf{(P\_6,4.148.1)} savarṇadīrghatvena api etāni siddhāni . \\
\noindent \textbf{(P\_6,4.148.1)} idam tarhi . \\
\noindent \textbf{(P\_6,4.148.1)} atisakheḥ āgacchati . \\
\noindent \textbf{(P\_6,4.148.1)} atisakheḥ svam . \\
\noindent \textbf{(P\_6,4.148.1)} yadi lopaḥ na syāt upasarjanahrasvatve kṛte asakhi iti pratiṣedhaḥ prasajyeta . \\
\noindent \textbf{(P\_6,4.148.2)} yasya ītyādau śyām pratiṣedhaḥ . \\
\noindent \textbf{(P\_6,4.148.2)} yasya ītyādau śyām pratiṣedhaḥ vaktavyaḥ . \\
\noindent \textbf{(P\_6,4.148.2)} kāṇḍe kuḍye . \\
\noindent \textbf{(P\_6,4.148.2)} saurye nāma himavataḥ śṛṅge . \\
\noindent \textbf{(P\_6,4.148.2)} saḥ tarhi pratiṣedhaḥ vaktavyaḥ . \\
\noindent \textbf{(P\_6,4.148.2)} na vaktavyaḥ . \\
\noindent \textbf{(P\_6,4.148.2)} iha śyām iti api prakṛtam na iti api . \\
\noindent \textbf{(P\_6,4.148.2)} tatra abhisambandhamātram kartavyam : yasya ītyādau lopaḥ bhavati śyām na . \\
\noindent \textbf{(P\_6,4.148.2)} iyaṅuvaṅbhyām lopaḥ vipratiṣedhena . \\
\noindent \textbf{(P\_6,4.148.2)} iyaṅuvaṅbhyām lopaḥ bhavati vipratiṣedhena . \\
\noindent \textbf{(P\_6,4.148.2)} iyaṅuvaṅoḥ avakāśaḥ śriyau śriyaḥ , bhruvau bhruvaḥ . \\
\noindent \textbf{(P\_6,4.148.2)} lopasya avakāśaḥ kāmaṇḍaleyaḥ , mādrabāheyaḥ . \\
\noindent \textbf{(P\_6,4.148.2)} iha ubhayam prāpnoti : vatsapreyaḥ , laikhābhreyaḥ . \\
\noindent \textbf{(P\_6,4.148.2)} lopaḥ bhavati vipratiṣedhena . \\
\noindent \textbf{(P\_6,4.148.2)} guṇavṛddhī ca . guṇavṛddhī ca iyaṅuvaṅbhyām bhavataḥ vipratiṣedhena . \\
\noindent \textbf{(P\_6,4.148.2)} guṇavṛddhyoḥ avakāśaḥ : cetā gauḥ . \\
\noindent \textbf{(P\_6,4.148.2)} iyaṅuvaṅoḥ saḥ eva . \\
\noindent \textbf{(P\_6,4.148.2)} iha ubhayam prāpnoti : cayanam , cāyakaḥ , lavanam , lāvakaḥ . \\
\noindent \textbf{(P\_6,4.148.2)} guṇavṛddhī bhavataḥ vipratiṣedhena . \\
\noindent \textbf{(P\_6,4.148.2)} na vā iyaṅuvaṅādeśasya anyaviṣaye vacanāt . \\
\noindent \textbf{(P\_6,4.148.2)} na vā arthaḥ vipratiṣedhena . \\
\noindent \textbf{(P\_6,4.148.2)} kim kāraṇam . \\
\noindent \textbf{(P\_6,4.148.2)} iyaṅuvaṅādeśasya anyaviṣaye vacanāt . \\
\noindent \textbf{(P\_6,4.148.2)} iyaṅuvaṅādeśaḥ anyaviṣaye ārabhyate . \\
\noindent \textbf{(P\_6,4.148.2)} kiṃviṣaye . \\
\noindent \textbf{(P\_6,4.148.2)} yaṇādiviṣaye . \\
\noindent \textbf{(P\_6,4.148.2)} saḥ yathā yaṇādeśam bādhate evam guṇavṛddhī bādheta . \\
\noindent \textbf{(P\_6,4.148.2)} tasmāt tatra guṇavṛddhiviṣaye pratiṣedhaḥ . \\
\noindent \textbf{(P\_6,4.148.2)} tasmāt tatra guṇavṛddhiviṣaye pratiṣedhaḥ vaktavyaḥ . \\
\noindent \textbf{(P\_6,4.148.2)} na vaktavyaḥ . \\
\noindent \textbf{(P\_6,4.148.2)} madhye apavādāḥ pūrvān vidhīn bādhante iti evam iyaṅuvaṅādeśaḥ yaṇādeśam bādhiṣyate guṇavṛddhī na bādhiṣyate . \\
\noindent \textbf{(P\_6,4.149.1)} sūryādīnām aṇante aprasiddhiḥ aṅgānyatvāt . \\
\noindent \textbf{(P\_6,4.149.1)} sūryādīnām aṇante aprasiddhiḥ . \\
\noindent \textbf{(P\_6,4.149.1)} saurī balākā . \\
\noindent \textbf{(P\_6,4.149.1)} kim kāraṇam . \\
\noindent \textbf{(P\_6,4.149.1)} aṅgānyatvāt . \\
\noindent \textbf{(P\_6,4.149.1)} aṇantam etat aṅgam anyat bhavati . \\
\noindent \textbf{(P\_6,4.149.1)} lope kṛte na aṅgānyatvam . \\
\noindent \textbf{(P\_6,4.149.1)} sthānivadbhāvāt aṅgānyatvam bhavati . \\
\noindent \textbf{(P\_6,4.149.1)} siddham tu sthānivatpratiṣedhāt . \\
\noindent \textbf{(P\_6,4.149.1)} siddham etat . \\
\noindent \textbf{(P\_6,4.149.1)} katham . \\
\noindent \textbf{(P\_6,4.149.1)} sthānivatpratiṣedhāt . \\
\noindent \textbf{(P\_6,4.149.1)} pratiṣidhyate atra sthānivadbhāvaḥ yalopavidhim prati na sthānivat bhavati iti . \\
\noindent \textbf{(P\_6,4.149.1)} evam api na sidhyati . \\
\noindent \textbf{(P\_6,4.149.1)} kim kāraṇam . \\
\noindent \textbf{(P\_6,4.149.1)} śabdānyatvāt . \\
\noindent \textbf{(P\_6,4.149.1)} anyaḥ hi śūryaśabdaḥ anyaḥ sauryaśabdaḥ . \\
\noindent \textbf{(P\_6,4.149.1)} na eṣaḥ doṣaḥ . \\
\noindent \textbf{(P\_6,4.149.1)} ekadeśavikṛtam ananyavat bhavati iti bhaviṣyati . \\
\noindent \textbf{(P\_6,4.149.1)} upadhāgrahaṇānarthakyam ca . \\
\noindent \textbf{(P\_6,4.149.1)} sthānivadbhāve ca idānīm pratiṣiddhe upadhāgrahaṇam anarthakam . \\
\noindent \textbf{(P\_6,4.149.1)} kim kāraṇam . \\
\noindent \textbf{(P\_6,4.149.1)} antyaḥ eva hi sūryādīnām yakāraḥ . \\
\noindent \textbf{(P\_6,4.149.1)} kim yātam etat bhavati . \\
\noindent \textbf{(P\_6,4.149.1)} suṣṭhu ca yātam sādhu ca yātam yadi prāk bhāt asiddhatvam . \\
\noindent \textbf{(P\_6,4.149.1)} atha hi saha tena asiddhatvam asiddhatvāt lopasya na antyaḥ yakāraḥ bhavati . \\
\noindent \textbf{(P\_6,4.149.1)} yadi api saha tena asiddhatvam evam api na doṣaḥ . \\
\noindent \textbf{(P\_6,4.149.1)} na evam vijñāyate sūryādīnām aṅgānām yakāralopaḥ iti . \\
\noindent \textbf{(P\_6,4.149.1)} katham tarhi . \\
\noindent \textbf{(P\_6,4.149.1)} aṅgasya yalopaḥ bhavati saḥ cet sūryādīnām yakāraḥ iti . \\
\noindent \textbf{(P\_6,4.149.1)} evam api sūryacarī , atra prāpnoti . \\
\noindent \textbf{(P\_6,4.149.1)} tasmāt upadhāgrahaṇam kartavyam . \\
\noindent \textbf{(P\_6,4.149.2)} viṣayaparigaṇanam ca . \\
\noindent \textbf{(P\_6,4.149.2)} viṣayaparigaṇanam ca kartavyam . \\
\noindent \textbf{(P\_6,4.149.2)} sūryamatsyayoḥ ṅyām . \\
\noindent \textbf{(P\_6,4.149.2)} sūryamatsyayoḥ ṅyām iti vaktavyam . \\
\noindent \textbf{(P\_6,4.149.2)} saurī matsī . \\
\noindent \textbf{(P\_6,4.149.2)} sūryāgastyayoḥ che ca . \\
\noindent \textbf{(P\_6,4.149.2)} sūryāgastyayoḥ che ca ṅyām ca iti vaktavyam . \\
\noindent \textbf{(P\_6,4.149.2)} saurī saurīyaḥ , āgastī , āgastīyaḥ . \\
\noindent \textbf{(P\_6,4.149.2)} tiṣyapuṣyayoḥ nakṣatrāṇi . \\
\noindent \textbf{(P\_6,4.149.2)} tiṣyapuṣyayoḥ nakṣatrāṇi lopaḥ vaktavyaḥ : taiṣam , pauṣam . \\
\noindent \textbf{(P\_6,4.149.2)} antikasya tasi kādilopaḥ ādyudāttatvam ca . \\
\noindent \textbf{(P\_6,4.149.2)} antikasya tasi kādilopaḥ vaktavyaḥ ādyudāttatvam ca vaktavyam . \\
\noindent \textbf{(P\_6,4.149.2)} antitaḥ na dūrāt . \\
\noindent \textbf{(P\_6,4.149.2)} tame tādeḥ ca . \\
\noindent \textbf{(P\_6,4.149.2)} tame tādeḥ ca kādeḥ ca lopaḥ vaktavyaḥ . \\
\noindent \textbf{(P\_6,4.149.2)} agne tvam naḥ antamaḥ . \\
\noindent \textbf{(P\_6,4.149.2)} antitamaḥ avarohati . \\
\noindent \textbf{(P\_6,4.149.2)} tasi iti eṣaḥ na vaktavyaḥ . \\
\noindent \textbf{(P\_6,4.149.2)} dṛṣṭaḥ dāśataye api hi ghau lopaḥ antiṣat iti yatra . antiṣat . \\
\noindent \textbf{(P\_6,4.149.2)} tathā aghau ye antyatharvasu . \\
\noindent \textbf{(P\_6,4.149.2)} anti ye ca dūrake . \\
\noindent \textbf{(P\_6,4.153)} chagrahaṇam śakyam akartum . \\
\noindent \textbf{(P\_6,4.153)} iha kasmāt na bhavati bilvakebhyaḥ . \\
\noindent \textbf{(P\_6,4.153)} bhasya iti vartate . \\
\noindent \textbf{(P\_6,4.153)} evam api bilvakāya , atra prāpnoti . \\
\noindent \textbf{(P\_6,4.153)} taddhitasya iti vartate . \\
\noindent \textbf{(P\_6,4.153)} evam api bilvakasya vikāraḥ avayavaḥ vā bailvakaḥ , atra prāpnoti . \\
\noindent \textbf{(P\_6,4.153)} taddhite taddhitasya iti vartate . \\
\noindent \textbf{(P\_6,4.153)} evam api bilvakīyāyām bhavaḥ bailvakaḥ , bailvakasya kim cit bailvakīyam , atra prāpnoti . \\
\noindent \textbf{(P\_6,4.153)} na saḥ bilvakāt . \\
\noindent \textbf{(P\_6,4.153)} bilvakādibhyaḥ yaḥ vihitaḥ iti ucyate na ca asau bilvakaśabdāt vihitaḥ . \\
\noindent \textbf{(P\_6,4.153)} kim tarhi bilvakīyaśabdāt . \\
\noindent \textbf{(P\_6,4.153)} evam tarhi siddhe sati yat chagrahaṇam karoti tat jñāpayati ācāryaḥ bhavati eṣā paribhāṣā : sanniyogaśiṣṭānām anyatarābhāve ubhayoḥ abhāvaḥ iti . \\
\noindent \textbf{(P\_6,4.153)} tasmāt chagrahaṇam kartavyam chasya eva luk yathā syāt kukaḥ mā bhūt iti . \\
\noindent \textbf{(P\_6,4.154)} tuḥ sarvasya lopaḥ vaktavyaḥ antyasya lopaḥ mā bhūt iti . \\
\noindent \textbf{(P\_6,4.154)} saḥ tarhi vaktavyaḥ . \\
\noindent \textbf{(P\_6,4.154)} na vaktavyaḥ . tuḥ sarvalopavijñānam antyasya vacanānarthakyāt . \\
\noindent \textbf{(P\_6,4.154)} tuḥ sarvalopaḥ vijñāyate . \\
\noindent \textbf{(P\_6,4.154)} kutaḥ . \\
\noindent \textbf{(P\_6,4.154)} antyasya vacanānarthakyāt . \\
\noindent \textbf{(P\_6,4.154)} antyasya lopavacane prayojanam na asti iti kṛtvā sarvasya bhaviṣyati . \\
\noindent \textbf{(P\_6,4.154)} atha vā luk prakṛtaḥ . \\
\noindent \textbf{(P\_6,4.154)} saḥ anuvartiṣyate . \\
\noindent \textbf{(P\_6,4.154)} aśakyaḥ luk anuvartayitum . \\
\noindent \textbf{(P\_6,4.154)} kim kāraṇam . \\
\noindent \textbf{(P\_6,4.154)} vijayiṣṭhakariṣthayoḥ guṇadarśanāt . \\
\noindent \textbf{(P\_6,4.154)} vijayiṣṭhakariṣthayoḥ guṇaḥ dṛśyate . \\
\noindent \textbf{(P\_6,4.154)} vijayiṣṭhaḥ . \\
\noindent \textbf{(P\_6,4.154)} āsutim kariṣṭhaḥ . \\
\noindent \textbf{(P\_6,4.155)} ṇau iṣṭhavat prātipadikasya . \\
\noindent \textbf{(P\_6,4.155)} ṇau prātipadikasya iṣṭhavadbhāvaḥ vaktavyaḥ . \\
\noindent \textbf{(P\_6,4.155)} kim prayojanam . \\
\noindent \textbf{(P\_6,4.155)} puṃvadbhāvarabhāvaṭilopayaṇādiparārtham . \\
\noindent \textbf{(P\_6,4.155)} puṃvadbhāvārtham . \\
\noindent \textbf{(P\_6,4.155)} enīm ācaṣṭe , etayati . \\
\noindent \textbf{(P\_6,4.155)} śyetayati . \\
\noindent \textbf{(P\_6,4.155)} rabhāvārtham . \\
\noindent \textbf{(P\_6,4.155)} pṛthum ācaṣṭe , prathayati . \\
\noindent \textbf{(P\_6,4.155)} mradayati . \\
\noindent \textbf{(P\_6,4.155)} ṭilopārtham . \\
\noindent \textbf{(P\_6,4.155)} paṭum ācaṣṭe paṭayati . \\
\noindent \textbf{(P\_6,4.155)} yaṇādiparārtham . \\
\noindent \textbf{(P\_6,4.155)} sthūlam ācaṣṭe sthavayati . \\
\noindent \textbf{(P\_6,4.155)} davayati . \\
\noindent \textbf{(P\_6,4.155)} kim punaḥ idam parigaṇanam āhosvit udāharaṇamātram . \\
\noindent \textbf{(P\_6,4.155)} udāharaṇamātram iti āha . \\
\noindent \textbf{(P\_6,4.155)} prādayaḥ api hi iṣyante : priyam ācaṣṭe prāpayati . \\
\noindent \textbf{(P\_6,4.155)} bhāradvājīyāḥ paṭhanti : ṇau iṣṭhavat prātipadikasya puṃvadbhāvarabhāvaṭilopayaṇādiparaprādivinmatorlukkanvidhyartham iti . \\
\noindent \textbf{(P\_6,4.159)} kim ayam yiśabdaḥ āhosvit yakāraḥ . \\
\noindent \textbf{(P\_6,4.159)} kim ca ataḥ . \\
\noindent \textbf{(P\_6,4.159)} yadi lopaḥ api anuvartate tatao yiśabdaḥ . \\
\noindent \textbf{(P\_6,4.159)} atha nivṛttam tataḥ yakāraḥ . \\
\noindent \textbf{(P\_6,4.161)} katham idam vijñāyate : halādeḥ aṅgasya iti āhosvit halādeḥ ṛkārasya iti . \\
\noindent \textbf{(P\_6,4.161)} yuktam punaḥ idam vicārayitum . \\
\noindent \textbf{(P\_6,4.161)} nanu anena asandigdhena aṅgaviśeṣaṇena bhavitavyam . \\
\noindent \textbf{(P\_6,4.161)} katham hi ṛkārasya nāma hal ādiḥ syāt anyasya anyaḥ . \\
\noindent \textbf{(P\_6,4.161)} ayam ādiśabdaḥ asti eva avayavavācī . \\
\noindent \textbf{(P\_6,4.161)} tat yathā ṛgādiḥ , ardharcādiḥ , ślokādiḥ iti . \\
\noindent \textbf{(P\_6,4.161)} asti sāmīpye vartate . \\
\noindent \textbf{(P\_6,4.161)} tat yathā . \\
\noindent \textbf{(P\_6,4.161)} dadhibhojanam arthasiddheḥ ādiḥ . \\
\noindent \textbf{(P\_6,4.161)} dadhibhojanasamīpe . \\
\noindent \textbf{(P\_6,4.161)} ghṛtabhojanam ārogyasya ādiḥ . \\
\noindent \textbf{(P\_6,4.161)} ghṛtabhojanasamīpe . \\
\noindent \textbf{(P\_6,4.161)} yāvatā sāmīpye api vartate jāyate vicāraṇā : halsamīpasya ṛkārasya halādeḥ aṅgasya iti . \\
\noindent \textbf{(P\_6,4.161)} kim ca ataḥ . \\
\noindent \textbf{(P\_6,4.161)} yadi vijñāyate halādeḥ aṅgasya iti aprathīyān , atra na prāpnoti . \\
\noindent \textbf{(P\_6,4.161)} atha vijñāyate halādeḥ ṛkārasya iti anṛcīyān , atra api prāpnoti . \\
\noindent \textbf{(P\_6,4.161)} ubhayathā svṛcīyān iti atra prāpnoti . \\
\noindent \textbf{(P\_6,4.161)} astu tāvat halādeḥ aṅgasya iti . \\
\noindent \textbf{(P\_6,4.161)} katham aprathīyān . \\
\noindent \textbf{(P\_6,4.161)} taddhitāntena samāsaḥ bhaviṣyati . \\
\noindent \textbf{(P\_6,4.161)} na prathīyān aprathīyān iti . \\
\noindent \textbf{(P\_6,4.161)} bhavet siddham yadā taddhitāntena samāsaḥ . \\
\noindent \textbf{(P\_6,4.161)} yadā tu khalu samāsāt taddhitotpattiḥ tadā na sidhyati . \\
\noindent \textbf{(P\_6,4.161)} na eva samāsāt taddhitotpattyā bhavitavyam . \\
\noindent \textbf{(P\_6,4.161)} kim kāraṇam . \\
\noindent \textbf{(P\_6,4.161)} bahuvrīhiṇā uktatvāt matvarthasya . \\
\noindent \textbf{(P\_6,4.161)} bhavet yadā bahuvrīhiḥ tadā na syāt . \\
\noindent \textbf{(P\_6,4.161)} yadā tu khalu tatpuruṣaḥ tadā prāpnoti . \\
\noindent \textbf{(P\_6,4.161)} na pṛthuḥ apṛthuḥ . \\
\noindent \textbf{(P\_6,4.161)} ayam api apṛthuḥ . \\
\noindent \textbf{(P\_6,4.161)} ayam api apṛthuḥ . \\
\noindent \textbf{(P\_6,4.161)} ayam anayoḥ aprathīyān iti . \\
\noindent \textbf{(P\_6,4.161)} na samāsāt ajādibhyām bhavitavyam . \\
\noindent \textbf{(P\_6,4.161)} kim kāraṇam . \\
\noindent \textbf{(P\_6,4.161)} guṇavacanāt iti ucyate na ca samāsaḥ guṇavacanaḥ iti . \\
\noindent \textbf{(P\_6,4.161)} yadā tarhi samāsāt vinmatupau vinmatubantāt ajādī tadā prāpnutaḥ . \\
\noindent \textbf{(P\_6,4.161)} avidyamānāḥ pṛthavaḥ apṛthavaḥ . \\
\noindent \textbf{(P\_6,4.161)} apṛthavaḥ asya santi apṛthumān . \\
\noindent \textbf{(P\_6,4.161)} ayam apṛthumān . \\
\noindent \textbf{(P\_6,4.161)} ayam apṛthumān . \\
\noindent \textbf{(P\_6,4.161)} ayam anayoḥ aprathīyān iti . \\
\noindent \textbf{(P\_6,4.161)} na eṣaḥ doṣaḥ . \\
\noindent \textbf{(P\_6,4.161)} apṛthavaḥ eva na santi kutaḥ yasya apṛthavaḥ iti . \\
\noindent \textbf{(P\_6,4.161)} iha kasmāt na bhavati : mātayati , bhrātayati . \\
\noindent \textbf{(P\_6,4.161)} lopaḥ atra bādhakaḥ bhaviṣyati . \\
\noindent \textbf{(P\_6,4.161)} idam iha sampradhāryam . \\
\noindent \textbf{(P\_6,4.161)} ṭilopaḥ kriyatām rabhāvaḥ iti kim atra kartavyam . \\
\noindent \textbf{(P\_6,4.161)} paratvāt rabhāvaḥ . \\
\noindent \textbf{(P\_6,4.161)} yadi punaḥ avaśiṣṭasya rabhāvaḥ ucyeta . \\
\noindent \textbf{(P\_6,4.161)} na evam śakyam . \\
\noindent \textbf{(P\_6,4.161)} iha api prasajyeta : kṛtam ācaṣṭe , kṛtayati iti . \\
\noindent \textbf{(P\_6,4.161)} evam tarhi parigaṇanam kartavyam . \\
\noindent \textbf{(P\_6,4.161)} pṛthumṛdukṛśabhṛśadṛḍhaparivṛḍhānām iti vaktavyam . \\
\noindent \textbf{(P\_6,4.163)} prakṛtyā ekāc iti kim iṣṭheymeyassu āhosvit aviśeṣeṇa . \\
\noindent \textbf{(P\_6,4.163)} kim ca ataḥ . \\
\noindent \textbf{(P\_6,4.163)} yadi aviśeṣeṇa svī khī śauvam adhunā iti atra api prāpnoti . \\
\noindent \textbf{(P\_6,4.163)} svikhinau eva na staḥ . \\
\noindent \textbf{(P\_6,4.163)} katham . \\
\noindent \textbf{(P\_6,4.163)} uktam etat . \\
\noindent \textbf{(P\_6,4.163)} ekākṣarāt kṛtaḥ jāteḥ saptamyām ca na tau smṛtau . \\
\noindent \textbf{(P\_6,4.163)} svavān khavān iti eva bhavitavyam . \\
\noindent \textbf{(P\_6,4.163)} śauvam iti paratvāt aijāgame kṛte ṭilopena bhavitavyam . \\
\noindent \textbf{(P\_6,4.163)} adhunā iti saprakṛtikasya sapratyayakasya sthāne nipātanam kriyate . \\
\noindent \textbf{(P\_6,4.163)} iha tarhi prāpnoti : dravyam . \\
\noindent \textbf{(P\_6,4.163)} yasya īti ādau prakṛtibhāvaḥ . \\
\noindent \textbf{(P\_6,4.163)} yasya īti yasya lopaprāptiḥ tasya prakṛtibhāvaḥ na ca etāni yasya īti ādau . \\
\noindent \textbf{(P\_6,4.163)} evam api śriye hitaḥ śrīyaḥ , jñā devatā asya sthālīpākasya jñaḥ sthālīyāpākaḥ iti atra prāpnoti . \\
\noindent \textbf{(P\_6,4.163)} tasmāt iṣṭheymeyassu prakṛtibhāvaḥ . \\
\noindent \textbf{(P\_6,4.163)} atha iṣṭheymeyassu prakṛtibhāve kim udāharaṇam . \\
\noindent \textbf{(P\_6,4.163)} preyān preṣthaḥ . \\
\noindent \textbf{(P\_6,4.163)} na etat asti . \\
\noindent \textbf{(P\_6,4.163)} prādīnām asiddhatvāt na bhaviṣyati . \\
\noindent \textbf{(P\_6,4.163)} idam tarhi śreyān , śreṣṭhaḥ . \\
\noindent \textbf{(P\_6,4.163)} prakṛtyā ekāc iṣṭheymeyassu cet ekācaḥ uccāraṇasāmarthyāt avacanāt prakṛtibhāvaḥ . \\
\noindent \textbf{(P\_6,4.163)} prakṛtyā ekāc iṣṭheymeyassu cet tat na . \\
\noindent \textbf{(P\_6,4.163)} kim kāraṇam . \\
\noindent \textbf{(P\_6,4.163)} ekācaḥ uccāraṇasāmarthyāt antareṇa api vacanam prakṛtibhāvaḥ bhaviṣyati . \\
\noindent \textbf{(P\_6,4.163)} vinmatoḥ tu lugartham . \\
\noindent \textbf{(P\_6,4.163)} vinmatoḥ tu lugartham prakṛtibhāvaḥ vaktavyaḥ . \\
\noindent \textbf{(P\_6,4.163)} sragvitaraḥ , srajīyān , sragvitamaḥ , srajiṣṭhaḥ . \\
\noindent \textbf{(P\_6,4.163)} srugvattaraḥ , srucīyān , srugvattamaḥ , sruciṣṭhaḥ . \\
\noindent \textbf{(P\_6,4.163)} nanu ca vinmatoḥ luk ṭilopam bādhiṣyate . \\
\noindent \textbf{(P\_6,4.163)} katham anyasya ucyamānasya anyasya bādhakam syāt . \\
\noindent \textbf{(P\_6,4.163)} asati khalu api sambhave bādhanam bhavati . \\
\noindent \textbf{(P\_6,4.163)} asti ca sambhavaḥ yat ubhayam syāt . \\
\noindent \textbf{(P\_6,4.163)} yathā eva khalu api vinmatoḥ luk ṭilopam bādhate eva naḥ taddhite iti etam api bādheta . \\
\noindent \textbf{(P\_6,4.163)} yataraḥ naḥ brahmīyān . \\
\noindent \textbf{(P\_6,4.163)} brahmavattaraḥ iti . \\
\noindent \textbf{(P\_6,4.163)} yat tāvat ucyate katham anyasya ucyamānasya anyasya bādhakam syāt iti . \\
\noindent \textbf{(P\_6,4.163)} idam tāvat ayam praṣṭavyaḥ . \\
\noindent \textbf{(P\_6,4.163)} yadi tarhi vinmatoḥ luk na ucyeta kim iha syāt iti . \\
\noindent \textbf{(P\_6,4.163)} ṭilopaḥ iti āha . \\
\noindent \textbf{(P\_6,4.163)} ṭilopaḥ cet na aprāpte ṭilope vinmatoḥ luk ārabhyate . \\
\noindent \textbf{(P\_6,4.163)} saḥ bādhakaḥ bhaviṣyati . \\
\noindent \textbf{(P\_6,4.163)} yat api ucyate asati khalu api sambhave bādhanam bhavati . \\
\noindent \textbf{(P\_6,4.163)} asti ca sambhavaḥ yat ubhayam syāt iti . \\
\noindent \textbf{(P\_6,4.163)} sati api sambhave bādhanam bhavati . \\
\noindent \textbf{(P\_6,4.163)} tat yathā . \\
\noindent \textbf{(P\_6,4.163)} dadhi brāhmaṇebhyaḥ dīyatām takram kauṇḍinyāya iti sati api sambhave dadhidānasya takradānam nivartakam bhavati . \\
\noindent \textbf{(P\_6,4.163)} evam iha api sati api sambhave vinmatoḥ luk ṭilopam bādhiṣyate . \\
\noindent \textbf{(P\_6,4.163)} yat api ucyate yathā eva khalu api vinmatoḥ luk ṭilopam bādhate eva naḥ taddhite iti etam api bādheta iti . \\
\noindent \textbf{(P\_6,4.163)} na bādhate . \\
\noindent \textbf{(P\_6,4.163)} kim kāraṇam . \\
\noindent \textbf{(P\_6,4.163)} yena na aprāpte tasya bādhanam . \\
\noindent \textbf{(P\_6,4.163)} na aprāpte ṭilope vinmatoḥ luk ārabhyate . \\
\noindent \textbf{(P\_6,4.163)} naḥ taddhite iti etasmin punaḥ prāpte ca aprāpte ca . \\
\noindent \textbf{(P\_6,4.163)} atha vā purastāt apavādāḥ anantarān vidhīn bādhante iti evam vinmatoḥ luk ṭilopam bādhiṣyate . \\
\noindent \textbf{(P\_6,4.163)} naḥ taddhite iti etam na bādhiṣyate . \\
\noindent \textbf{(P\_6,4.163)} yadi tarhi vinmatoḥ luk ṭilopam bādhate payiṣthaḥ iti na sidhyati . \\
\noindent \textbf{(P\_6,4.163)} payasiṣthaḥ iti prāpnoti . \\
\noindent \textbf{(P\_6,4.163)} yathālakṣaṇam aprayukte . \\
\noindent \textbf{(P\_6,4.163)} prakṛtyā ake rājanyamanuṣyayuvānaḥ . \\
\noindent \textbf{(P\_6,4.163)} rājanyamanuṣyayuvānaḥ ake prakṛtyā bhavanti iti vaktavyam . \\
\noindent \textbf{(P\_6,4.163)} rājanyakam , mānuṣyakam , yauvanikā . \\
\noindent \textbf{(P\_6,4.170)} mapūrvāt pratiṣedhe vā hitanāmnaḥ . \\
\noindent \textbf{(P\_6,4.170)} mapūrvāt pratiṣedhe vā hitanāmnaḥ iti vaktavyam . \\
\noindent \textbf{(P\_6,4.170)} samānaḥ haitanāmaḥ , samānaḥ haitanāmanaḥ iti ca . \\
\noindent \textbf{(P\_6,4.171)} atha kim idam brāhmasya ajātau anaḥ lopārtham vacanam āhosvit niyamārtham . \\
\noindent \textbf{(P\_6,4.171)} katha ca lopātham syāt katham va niyamārtham . \\
\noindent \textbf{(P\_6,4.171)} yadi tāvat apatye iti vartate tataḥ niyamārtham . \\
\noindent \textbf{(P\_6,4.171)} atha nivṛttam tataḥ lopārtham . \\
\noindent \textbf{(P\_6,4.171)} ataḥ uttaram paṭhati brāhmasya ajātau lopārtham vacanam . \\
\noindent \textbf{(P\_6,4.171)} brāhmasya ajātau lopārtham vacanam kriyate . \\
\noindent \textbf{(P\_6,4.171)} apatye iti nivṛttam . \\
\noindent \textbf{(P\_6,4.171)} tatra aprāptavidhāne prāptapratiṣedhaḥ . \\
\noindent \textbf{(P\_6,4.171)} tatra aprāptasya ṭilopasya vidhāne prāptasya pratiṣedhaḥ vaktavyaḥ . \\
\noindent \textbf{(P\_6,4.171)} brāhmaṇaḥ . \\
\noindent \textbf{(P\_6,4.171)} na vā paryudāsasāmarthyāt . \\
\noindent \textbf{(P\_6,4.171)} na vā vaktavyaḥ . \\
\noindent \textbf{(P\_6,4.171)} kim kāraṇam . \\
\noindent \textbf{(P\_6,4.171)} paryudāsasāmarthyāt paryudāsaḥ atra bhaviṣyati . \\
\noindent \textbf{(P\_6,4.171)} asti anyat paryudāse prayojanam . \\
\noindent \textbf{(P\_6,4.171)} kim . \\
\noindent \textbf{(P\_6,4.171)} yā jātiḥ eva na apatyam . \\
\noindent \textbf{(P\_6,4.171)} brāhmī oṣadhiḥ iti . \\
\noindent \textbf{(P\_6,4.171)} na vai atra iṣyate . \\
\noindent \textbf{(P\_6,4.171)} aniṣṭam ca prāpnoti iṣṭam ca na sidhyati . \\
\noindent \textbf{(P\_6,4.171)} evam tarhi anuvartate apatye iti na tu apatye iti anena nipātanam abhisambadhyate : brāhmaḥ iti nipātyate apatye ajātau iti . \\
\noindent \textbf{(P\_6,4.171)} kim tarhi . \\
\noindent \textbf{(P\_6,4.171)} pratiṣedhaḥ abhisambadhyate : brāhmaḥ iti nipātyate . \\
\noindent \textbf{(P\_6,4.171)} apatye jātau na iti . \\
\noindent \textbf{(P\_6,4.172)} kimartham idam ucyate na naḥ taddhite iti eva siddham . \\
\noindent \textbf{(P\_6,4.172)} na sidhyati . \\
\noindent \textbf{(P\_6,4.172)} an aṇi iti prakṛtibhāvaḥ prasajyeta . \\
\noindent \textbf{(P\_6,4.172)} aṇi iti ucyate ṇaḥ ca ayam . \\
\noindent \textbf{(P\_6,4.172)} evam tarhi siddhe sati yat nipātanam karoti tat jñāpayati ācāryaḥ tācchīlike ṇe aṇkṛtāni bhavanti . \\
\noindent \textbf{(P\_6,4.172)} kim etasya jñāpane prayojanam . \\
\noindent \textbf{(P\_6,4.172)} caurī tāpasī iti aṇantāt iti īkāraḥ siddhaḥ bhavati . \\
\noindent \textbf{(P\_6,4.174)} atra bhrauṇahatye kim nipātyate . \\
\noindent \textbf{(P\_6,4.174)} yakārādau taddhite tatvam nipātyate . \\
\noindent \textbf{(P\_6,4.174)} bhrauṇahatye tatvanipātanānarthakyam sāmānyena kṛtatvāt . \\
\noindent \textbf{(P\_6,4.174)} bhrauṇahatye tatvanipātanam anarthakam . \\
\noindent \textbf{(P\_6,4.174)} kim kāraṇam . \\
\noindent \textbf{(P\_6,4.174)} sāmānyena kṛtatvāt . \\
\noindent \textbf{(P\_6,4.174)} sāmānyena eva atra tatvam bhaviṣyati . \\
\noindent \textbf{(P\_6,4.174)} hanaḥ taḥ aciṇṇamuloḥ iti . \\
\noindent \textbf{(P\_6,4.174)} jñāpakam tu taddhite tatvapratiṣedhasya . \\
\noindent \textbf{(P\_6,4.174)} evam tarhi jñāpayati ācāryaḥ na taddhite tatvam bhavati iti . \\
\noindent \textbf{(P\_6,4.174)} kim etasya jñāpane prayojanam . \\
\noindent \textbf{(P\_6,4.174)} bhrauṇaghnaḥ , vārtraghnaḥ iti atra tatvam na bhavati . \\
\noindent \textbf{(P\_6,4.174)} aikṣvākasya svarabhedāt nipātanam pṛthaktvena . \\
\noindent \textbf{(P\_6,4.174)} aikṣvākasya svarabhedāt nipātanam pṛthaktvena kartavyam . \\
\noindent \textbf{(P\_6,4.174)} aikṣvākaḥ , aikṣvākaḥ . \\
\noindent \textbf{(P\_6,4.174)} ekaśrutyā nirdeśāt siddham .ekaśrutiḥ svarasarvanāma yathā napuṃsakam liṅgasarvanāma . \\
\noindent \textbf{(P\_6,4.174)} atha maitreye kim nipātyate . \\
\noindent \textbf{(P\_6,4.174)} maitreye ḍhañi yādilopanipātanam . \\
\noindent \textbf{(P\_6,4.174)} maitreye ḍhañi yādilopaḥ nipātyate . \\
\noindent \textbf{(P\_6,4.174)} idam mitrayuśabdasya catuḥ grahaṇam kriyate . \\
\noindent \textbf{(P\_6,4.174)} gṛṣṭyādiṣu pratyayavidhyartham pāṭhaḥ kriyate . \\
\noindent \textbf{(P\_6,4.174)} dvitīye adhyāye yaskādiṣu lugartham grahaṇam kriyate . \\
\noindent \textbf{(P\_6,4.174)} saptame adhyāye iyādeśārtham . \\
\noindent \textbf{(P\_6,4.174)} idam caturtham yādilopārtham . \\
\noindent \textbf{(P\_6,4.174)} dvirgrahaṇam śakyam akartum . \\
\noindent \textbf{(P\_6,4.174)} bidādiṣu pratyayavidhyartham pāṭhaḥ kartavyaḥ . \\
\noindent \textbf{(P\_6,4.174)} tatra na eva arthaḥ lukā na api yādilopena . \\
\noindent \textbf{(P\_6,4.174)} iyādeśena eva siddham . \\
\noindent \textbf{(P\_6,4.174)} na evam śakyam . \\
\noindent \textbf{(P\_6,4.174)} iha hi maitreyakaḥ saṅghaḥ iti saṅghātalakṣaṇeṣu añyañiñām aṇ iti aṇ prasajyeta . \\
\noindent \textbf{(P\_6,4.174)} hiraṇmaye kim nipātyate . \\
\noindent \textbf{(P\_6,4.174)} hiraṇmaye yalopavacanam . \\
\noindent \textbf{(P\_6,4.174)} hiraṇmaye yalopaḥ nipātyate . \\
\noindent \textbf{(P\_6,4.174)} atha hiraṇyaye kim nipātyate . \\
\noindent \textbf{(P\_6,4.174)} hiraṇyayasya chandasi malopavacanāt siddham . \\
\noindent \textbf{(P\_6,4.174)} hiraṇyayasya chandasi malopaḥ nipātyate . \\
\noindent \textbf{(P\_6,4.174)} hiraṇyayī naḥ nayatu . \\
\noindent \textbf{(P\_6,4.174)} hiraṇyayāḥ panthānaḥ āsan . \\
\noindent \textbf{(P\_6,4.174)} hiraṇyayam āsanam . \\
\noindent \textbf{(P\_7,1.1.1)} yuvoḥ anākau iti ucyate kayoḥ yuvoḥ anākau bhavataḥ . \\
\noindent \textbf{(P\_7,1.1.1)} pratyayayoḥ . \\
\noindent \textbf{(P\_7,1.1.1)} katham punaḥ aṅgasya iti anuvartamāne pratyayayoḥ syātām . \\
\noindent \textbf{(P\_7,1.1.1)} yuśabdavuśabdāntam etat vibhaktau aṅgam bhavati . \\
\noindent \textbf{(P\_7,1.1.1)} yadi yuśabdavuśabdāntasya aṅgasya anākau bhavataḥ sarvādeśau prāpnutaḥ . \\
\noindent \textbf{(P\_7,1.1.1)} nirdiśyamānasya ādeśāḥ bhavanti iti evam na bhaviṣyataḥ . \\
\noindent \textbf{(P\_7,1.1.1)} yatra tarhi vibhaktiḥ na asti . \\
\noindent \textbf{(P\_7,1.1.1)} nandanā kārikā iti . \\
\noindent \textbf{(P\_7,1.1.1)} atra api pratyayalakṣaṇena vibhaktiḥ . \\
\noindent \textbf{(P\_7,1.1.1)} yatra tarhi pratyayalakṣaṇam na asti . \\
\noindent \textbf{(P\_7,1.1.1)} nandanapriyaḥ kārakapriyaḥ iti . \\
\noindent \textbf{(P\_7,1.1.1)} mā bhūtām yā asau sāmāsikī vibhaktiḥ tasyām yā asau samāsāt vibhaktiḥ tasyām bhaviṣyataḥ . \\
\noindent \textbf{(P\_7,1.1.1)} na vai tasyām yuśabdavuśabdāntam aṅgam bhavati . \\
\noindent \textbf{(P\_7,1.1.1)} bhavet yaḥ yuśabdavuśabdābhyām aṅgam viśeṣayet tasya ānantyayoḥ na syātām . \\
\noindent \textbf{(P\_7,1.1.1)} vayam khalu aṅgena yuśabdavuśabdau viśeṣayiṣyāmaḥ . \\
\noindent \textbf{(P\_7,1.1.1)} aṅgasya yuvoḥ anākau bhavataḥ yatratatrasthayoḥ iti . \\
\noindent \textbf{(P\_7,1.1.1)} yatra tarhi samāsāt vibhaktiḥ na asti . \\
\noindent \textbf{(P\_7,1.1.1)} nandanadadhi kārakadadhi . \\
\noindent \textbf{(P\_7,1.1.1)} evam tarhi na ca aparam nimittam sañjñā ca pratyayalakṣaṇena . \\
\noindent \textbf{(P\_7,1.1.1)} na ca iha param nimittam āśrīyate : asmin parataḥ yuvoḥ anākau bhavataḥ iti . \\
\noindent \textbf{(P\_7,1.1.1)} kim tarhi aṅgasya yuvoḥ anākau bhavataḥ iti . \\
\noindent \textbf{(P\_7,1.1.1)} aṅgasañjñā ca bhavati pratyayalakṣaṇena . \\
\noindent \textbf{(P\_7,1.1.1)} atha vā tayoḥ eva yat aṅgam tannimittatvena āśrayiṣyāmaḥ . \\
\noindent \textbf{(P\_7,1.1.1)} katham . \\
\noindent \textbf{(P\_7,1.1.1)} aṅgasya iti sambandhasāmānye ṣaṣṭhī vijñāsyate . \\
\noindent \textbf{(P\_7,1.1.1)} aṅgasya yau yuvū . \\
\noindent \textbf{(P\_7,1.1.1)} kim ca aṅgasya yuvū . \\
\noindent \textbf{(P\_7,1.1.1)} nimittam . \\
\noindent \textbf{(P\_7,1.1.1)} yayoḥ yuvoḥ aṅgam iti etat bhavati . \\
\noindent \textbf{(P\_7,1.1.1)} kayoḥ ca etat bhavati . \\
\noindent \textbf{(P\_7,1.1.1)} pratyayayoḥ \\
\noindent \textbf{(P\_7,1.1.2)} yuvoḥ anākau iti cet dhātupratiṣedhaḥ . \\
\noindent \textbf{(P\_7,1.1.2)} yuvoḥ anākau iti cet dhātupratiṣedhaḥ vaktavyaḥ . \\
\noindent \textbf{(P\_7,1.1.2)} yutvā yutaḥ yutavān yutiḥ . \\
\noindent \textbf{(P\_7,1.1.2)} bhujyvādīnām ca . bhujyvādīnām ca pratiṣedhaḥ vaktavyaḥ . \\
\noindent \textbf{(P\_7,1.1.2)} bhujyuḥ kaṃyuḥ śaṃyuḥ iti . anunāsikaparatvāt siddham . anunāsikaparayoḥ yuvoḥ grahaṇam na ca etau anunāsikaparau . \\
\noindent \textbf{(P\_7,1.1.2)} yadi anunāsikaparayoḥ grahaṇam nandanaḥ kārakaḥ atra na prāpnutaḥ na hi etābhyām yuśabdavuśabdābhyām anunāsikam param paśyāmaḥ . \\
\noindent \textbf{(P\_7,1.1.2)} anunāsikaparatvāt iti na evam vijñāyate anunāsikaḥ paraḥ ābhyām tau imau anunāsikaparau anunāsikaparatvāt iti . \\
\noindent \textbf{(P\_7,1.1.2)} katham tarhi . \\
\noindent \textbf{(P\_7,1.1.2)} anunāsikaḥ paraḥ anayoḥ tau imau anunāsikaparau anunāsikaparatvāt iti . \\
\noindent \textbf{(P\_7,1.1.2)} yadi anunāsikaparayoḥ grahaṇam itsañjñā prāpnoti . \\
\noindent \textbf{(P\_7,1.1.2)} tatra kaḥ doṣaḥ . \\
\noindent \textbf{(P\_7,1.1.2)} tatra ṅībnumoḥ pratiṣedhaḥ . ṅībnumaḥ pratiṣedhaḥ vaktavyaḥ . \\
\noindent \textbf{(P\_7,1.1.2)} nandanaḥ kārakaḥ . \\
\noindent \textbf{(P\_7,1.1.2)} nandanā kārikā . \\
\noindent \textbf{(P\_7,1.1.2)} ugillakṣaṇau ṅībnumau prāpnutaḥ . dhātvantasya ca . dhātvantasya ca pratiṣedhaḥ vaktavyaḥ . \\
\noindent \textbf{(P\_7,1.1.2)} divu sivu . \\
\noindent \textbf{(P\_7,1.1.2)} ṣiṭṭitkaraṇam tu jñāpakam ugitkāryābhāvasya . yat ayam yuśabdavuśabdau ṣiṭṭitau karoti śilpini ṣvun ṭyuṭyulau tuṭ ca iti tat jñāpayati ācāryaḥ na yuvoḥ ugitkāryam bhavati iti . \\
\noindent \textbf{(P\_7,1.1.2)} katham kṛtvā jñāpakam ṣiṭṭitkaraṇe etat prayojanam ṣiṭṭitaḥ iti ikāraḥ yathā syāt . \\
\noindent \textbf{(P\_7,1.1.2)} yadi ca atra ugitkāryam syāt ṣiṭṭitkaraṇam anarthakam syāt . \\
\noindent \textbf{(P\_7,1.1.2)} paśyati tu ācāryaḥ na yuvoḥ ugitkāryam bhavati iti tataḥ yuśabdavuśabdau ṣiṭṭitau karoti . \\
\noindent \textbf{(P\_7,1.1.2)} na vā ṣitkaraṇam ṅīṣvidhānārtham . na etat asti jñāpakam . \\
\noindent \textbf{(P\_7,1.1.2)} asti hi anyat etasya vacane prayojanam . \\
\noindent \textbf{(P\_7,1.1.2)} kim . \\
\noindent \textbf{(P\_7,1.1.2)} ṣitkaraṇam kriyate ṅīṣvidhānārtham . \\
\noindent \textbf{(P\_7,1.1.2)} ṣitaḥ iti ṅīṣ yathā syāt . \\
\noindent \textbf{(P\_7,1.1.2)} ṭitkaraṇam anupasarjanārtham . \\
\noindent \textbf{(P\_7,1.1.2)} ṭitkaraṇe api anyat prayojanam asti . \\
\noindent \textbf{(P\_7,1.1.2)} kim . \\
\noindent \textbf{(P\_7,1.1.2)} anupasarjanāt ṭitaḥ iti īkāraḥ yathā syāt . \\
\noindent \textbf{(P\_7,1.1.2)} ṭitaḥ anupasarjanāt bhavati ugitaḥ upasarjanāt ca anupasarjanāt ca . \\
\noindent \textbf{(P\_7,1.1.2)} evam tarhi vipratiṣedhāt tu ṭāpaḥ balīyastvam . \\
\noindent \textbf{(P\_7,1.1.2)} vipratiṣedhāt tu ṭāpaḥ balīyastvam bhaviṣyati . \\
\noindent \textbf{(P\_7,1.1.2)} ṭāpaḥ avakāśaḥ khaṭvā mālā . \\
\noindent \textbf{(P\_7,1.1.2)} ṅīpaḥ avakāśaḥ gomatī yavamatī . \\
\noindent \textbf{(P\_7,1.1.2)} iha ubhayam prāpnoti nandanā kārikā . \\
\noindent \textbf{(P\_7,1.1.2)} ṭāp bhavati vipratiṣedhena . \\
\noindent \textbf{(P\_7,1.1.2)} na eṣaḥ yuktaḥ vipratiṣedhaḥ . \\
\noindent \textbf{(P\_7,1.1.2)} vipratiṣedhe param iti ucyate pūrvaḥ ca ṭāp paraḥ ṅīp . \\
\noindent \textbf{(P\_7,1.1.2)} ṅīpaḥ paraḥ ṭāp kariṣyate . \\
\noindent \textbf{(P\_7,1.1.2)} sūtraviparyāsaḥ kṛtaḥ bhavati . \\
\noindent \textbf{(P\_7,1.1.2)} evam tarhi ugitaḥ ṅīp bhavati iti atra api ataḥ ṭāp iti anuvartiṣyate . \\
\noindent \textbf{(P\_7,1.1.2)} evam api akārāntāt ugitaḥ iha eva syāt nandanā kārikā . \\
\noindent \textbf{(P\_7,1.1.2)} gomatī yavamatī iti atra na syāt . \\
\noindent \textbf{(P\_7,1.1.2)} evam tarhi sambandhānuvṛttiḥ kariṣyate . \\
\noindent \textbf{(P\_7,1.1.2)} ajādyataḥ ṭāp ṛnnebhyaḥ ṅīp ataḥ ṭāp . \\
\noindent \textbf{(P\_7,1.1.2)} ugitaḥ ca ṅīp bhavati ataḥ ṭāp . \\
\noindent \textbf{(P\_7,1.1.2)} vanaḥ ra ca vanaḥ ṅīp bhavati ugitaḥ ataḥ ṭāp . \\
\noindent \textbf{(P\_7,1.1.2)} pādaḥ anyatarasyām ṅīp bhavati ugitaḥ ataḥ ṭāp . \\
\noindent \textbf{(P\_7,1.1.2)} tataḥ ṛci . \\
\noindent \textbf{(P\_7,1.1.2)} ṛci ca ṭāp bhavati . \\
\noindent \textbf{(P\_7,1.1.2)} prakṛtam anuvartate . \\
\noindent \textbf{(P\_7,1.1.2)} sidhyati evam yadi vārttikakāraḥ paṭhati vipratiṣedhāt tu ṭāpaḥ balīyastvam iti etat asaṅgṛhītam bhavati . \\
\noindent \textbf{(P\_7,1.1.2)} etat ca saṅgṛhītam bhavati . \\
\noindent \textbf{(P\_7,1.1.2)} katham . \\
\noindent \textbf{(P\_7,1.1.2)} iṣṭavācī paraśabdaḥ . \\
\noindent \textbf{(P\_7,1.1.2)} vipratiṣedhe param yat iṣṭam tat bhavati iti . dhātvantasya ca arthavadgrahaṇāt . \\
\noindent \textbf{(P\_7,1.1.2)} arthavatoḥ yuvoḥ grahaṇam na ca dhātvantaḥ arthavān \\
\noindent \textbf{(P\_7,1.1.3)} numvidhau jhalgrahaṇam . \\
\noindent \textbf{(P\_7,1.1.3)} numvidhau jhalgrahaṇam kartavyam . \\
\noindent \textbf{(P\_7,1.1.3)} jhalantasya ugitaḥ iṣyate : ugidacām sarvanāmasthāne adhātoḥ jhalaḥ iti . \\
\noindent \textbf{(P\_7,1.1.3)} tat ca avaśyam kartavyam . liṅgaviśiṣṭapratiṣedhārtham . prātipadikagrahaṇe liṅgaviśiṣṭasya api grahaṇam bhavati iti yathā iha bhavati gomān yavamān evam gomatī yavamatī iti atra api syāt . na vā vibhaktau liṅgaviśiṣṭāgrahaṇāt . \\
\noindent \textbf{(P\_7,1.1.3)} na vā vaktavyam . \\
\noindent \textbf{(P\_7,1.1.3)} kim kāraṇam . \\
\noindent \textbf{(P\_7,1.1.3)} vibhaktau liṅgaviśiṣṭagrahaṇam na iti eṣā paribhāṣā kartavyā . \\
\noindent \textbf{(P\_7,1.1.3)} kaḥ punaḥ atra viśeṣaḥ eṣā vā paribhāṣā kriyeta jhalgrahaṇam vā iti . \\
\noindent \textbf{(P\_7,1.1.3)} avaśyam eṣā paribhāṣā kartavyā . \\
\noindent \textbf{(P\_7,1.1.3)} bahūni etasyāḥ paribhāṣāyāḥ prayojanāni . \\
\noindent \textbf{(P\_7,1.1.3)} kāni . prayojanam śunaḥ svare . yathā iha bhavati śunā śunaḥ evam śunyā śunyāḥ iti atra api syāt . yūnaḥ samprasāraṇe . \\
\noindent \textbf{(P\_7,1.1.3)} yūnaḥ samprasāraṇe prayojanam . \\
\noindent \textbf{(P\_7,1.1.3)} yathā iha bhavati yūnaḥ paśya iti evam yuvatīḥ paśya iti atra api syāt . ugidacām numvidhau . ugidacām numvidhau prayojanam . \\
\noindent \textbf{(P\_7,1.1.3)} yathā iha bhavati gomān yavamān evam gomatī yavamatī iti atra api syāt . anaḍuhaḥ ca āmvidhau . \\
\noindent \textbf{(P\_7,1.1.3)} anaḍuhaḥ ca āmvidhau prayojanam . \\
\noindent \textbf{(P\_7,1.1.3)} yathā iha bhavati anaḍvān iti evam anaḍuhī iti atra api syāt . \\
\noindent \textbf{(P\_7,1.1.3)} na vā bhavati anaḍvāhī iti . \\
\noindent \textbf{(P\_7,1.1.3)} bhavati anyena yatnena . \\
\noindent \textbf{(P\_7,1.1.3)} ām anaḍuhaḥ striyām vā iti . \\
\noindent \textbf{(P\_7,1.1.3)} liṅgaviśiṣṭagrahaṇāt īkārāntasya prāpnoti . pathimathoḥ āttve . pathimathoḥ āttve prayojanam . \\
\noindent \textbf{(P\_7,1.1.3)} yathā iha bhavati panthāḥ manthāḥ evam pathī mathī iti atra api prāpnoti . \\
\noindent \textbf{(P\_7,1.1.3)} na kevalaḥ pathiśabdaḥ striyām vartate . \\
\noindent \textbf{(P\_7,1.1.3)} upasamastaḥ tarhi vartate . \\
\noindent \textbf{(P\_7,1.1.3)} supathī iti . puṃsaḥ asuṅvidhau . \\
\noindent \textbf{(P\_7,1.1.3)} puṃsaḥ asuṅvidhau prayojanam . \\
\noindent \textbf{(P\_7,1.1.3)} yathā iha bhavati pumān evam puṃsī iti atra api syāt . \\
\noindent \textbf{(P\_7,1.1.3)} na kevalaḥ puṃśabdaḥ striyām vartate . \\
\noindent \textbf{(P\_7,1.1.3)} upasamastaḥ tarhi vartate . \\
\noindent \textbf{(P\_7,1.1.3)} supuṃsī iti . sakhyuḥ ṇittvānaṅau . sakhyuḥ ṇittvānaṅau prayojanam . \\
\noindent \textbf{(P\_7,1.1.3)} yathā iha bhavati sakhā sakhāyau sakhāyaḥ evam sakhī sakhyau sakhyaḥ iti atra api prāpnoti . bhavadbhagavadaghavatām odbhāve . bhavadbhagavadaghavatām odbhāve prayojanam . \\
\noindent \textbf{(P\_7,1.1.3)} yathā iha bhavati bhoḥ bhagoḥ aghoḥ iti evam bhavati bhagavati aghavati iti atra api syāt . \\
\noindent \textbf{(P\_7,1.1.3)} etāni asyāḥ paribhāṣāyāḥ prayojanāni yadartham eṣā paribhāṣā kartavyā . \\
\noindent \textbf{(P\_7,1.1.3)} etasyām ca satyām na arthaḥ jhalgrahaṇena \\
\noindent \textbf{(P\_7,1.1.4)} tat etat ananyārtham jhalgrahaṇam kartavyam numpratiṣedhaḥ vā vaktavyaḥ . \\
\noindent \textbf{(P\_7,1.1.4)} ubhayam na vaktavyam . \\
\noindent \textbf{(P\_7,1.1.4)} upariṣṭāt jhalgrahaṇam kriyate tat purastāt apakrakṣyate . \\
\noindent \textbf{(P\_7,1.1.4)} evam api sūtraviparyāsaḥ kṛtaḥ bhavati . \\
\noindent \textbf{(P\_7,1.1.4)} evam tarhi yogavibhāgaḥ kariṣyate . \\
\noindent \textbf{(P\_7,1.1.4)} ugidacām sarvanāmasthāne adhātoḥ . \\
\noindent \textbf{(P\_7,1.1.4)} yujeḥ asamāse . \\
\noindent \textbf{(P\_7,1.1.4)} tataḥ napuṃsakasya . \\
\noindent \textbf{(P\_7,1.1.4)} napuṃsakasya num bhavati . \\
\noindent \textbf{(P\_7,1.1.4)} jhalaḥ iti ubhayoḥ śeṣaḥ . \\
\noindent \textbf{(P\_7,1.1.4)} tataḥ acaḥ . \\
\noindent \textbf{(P\_7,1.1.4)} ajantasya ca napuṃsakaliṅgasya num bhavati . \\
\noindent \textbf{(P\_7,1.1.4)} yadi api tāvat etat ugitkāryam parihṛtam idam aparam prāpnoti : śātanitarā pātanitarā . \\
\noindent \textbf{(P\_7,1.1.4)} ugitaḥ nadyāḥ ghādiṣu hrasvaḥ bhavati iti anyatarasyām hrasvatvam prasajyeta nityam ca iṣyate . \\
\noindent \textbf{(P\_7,1.1.4)} ugitaḥ yā nadī evam etat vijñāyate . \\
\noindent \textbf{(P\_7,1.1.4)} ugitaḥ eṣā nadī . \\
\noindent \textbf{(P\_7,1.1.4)} ugitaḥ yā parā . \\
\noindent \textbf{(P\_7,1.1.4)} atra ca eva doṣaḥ bhavati ugitaḥ hi eṣā parā nadī aiṣumatitarāyām ca prāpnoti . \\
\noindent \textbf{(P\_7,1.1.4)} ugitaḥ parā yā vihitā . \\
\noindent \textbf{(P\_7,1.1.4)} ugitaḥ eṣā vihitā . \\
\noindent \textbf{(P\_7,1.1.4)} ugitaḥ iti evam yā vihitā . \\
\noindent \textbf{(P\_7,1.1.4)} evam api bhogavatitarāyām doṣaḥ bhavati . \\
\noindent \textbf{(P\_7,1.1.4)} bhogavatitarā bhogavatītarā . \\
\noindent \textbf{(P\_7,1.1.4)} tasmāt ugitaḥ yā nadī ugitaḥ yā vihitā iti evam etat vijñāsyate . \\
\noindent \textbf{(P\_7,1.1.4)} evam vijñāyamāne śātanitarāyām doṣaḥ eva . siddham tu yuvoḥ anunāsikatvāt siddham etat . \\
\noindent \textbf{(P\_7,1.1.4)} katham . \\
\noindent \textbf{(P\_7,1.1.4)} yakāravakārayoḥ eva idam anunāsikayoḥ grahaṇam . \\
\noindent \textbf{(P\_7,1.1.4)} santi hi yaṇaḥ sānunāsikāḥ niranunāsikāḥ ca . \\
\noindent \textbf{(P\_7,1.2)} āyanādiṣu upadeśivadvacanam svarasiddhyartham . āyanādiṣu upadeśivadbhāvaḥ vaktavyaḥ . \\
\noindent \textbf{(P\_7,1.2)} upadeśāvasthāyām āyanādayaḥ bhavanti iti vaktavyam . \\
\noindent \textbf{(P\_7,1.2)} kim prayojanam . \\
\noindent \textbf{(P\_7,1.2)} svarasiddhyartham . \\
\noindent \textbf{(P\_7,1.2)} upadeśāvasthāyām āyanādiṣu iṣṭaḥ svaraḥ yathā syāt iti . \\
\noindent \textbf{(P\_7,1.2)} śileyam taittirīyaḥ . \\
\noindent \textbf{(P\_7,1.2)} akriyamāṇe hi upadeśivadbhāve pratyayasañjñāsanniyogena ādyudāttatve kṛte āntaryataḥ ādeśāḥ asvarakāṇām asvarakāḥ syuḥ . \\
\noindent \textbf{(P\_7,1.2)} na vā kva cit citkaraṇāt upadeśivadvacanānarthakyam na vā vaktavyam . \\
\noindent \textbf{(P\_7,1.2)} kim kāraṇam . \\
\noindent \textbf{(P\_7,1.2)} kva cit citkaraṇāt . \\
\noindent \textbf{(P\_7,1.2)} yat ayam kva cit ghādīn citaḥ karoti agrāt yat ghacchau ca tat jñāpayati ācāryaḥ upadeśāvasthāyām āyanādayaḥ bhavanti iti . \\
\noindent \textbf{(P\_7,1.2)} katham kṛtvā jñāpakam . \\
\noindent \textbf{(P\_7,1.2)} citkaraṇe etat prayojanam citaḥ iti antodāttatvam yathā syāt iti . \\
\noindent \textbf{(P\_7,1.2)} yadi ca upadeśāvasthāyām āyanādayaḥ bhavanti tataḥ citkaraṇam arthavat bhavati . tatra uṇādipratiṣedhaḥ . tatra uṇādīnām pratiṣedhaḥ vaktavyaḥ . \\
\noindent \textbf{(P\_7,1.2)} śaṅkhaḥ śāṇḍhaḥ iti . dhātoḥ vā īyaṅvacanāt . atha vā yat ayam ṛteḥ īyaṅ iti dhātoḥ īyaṅ śāsti tat jñāpayati ācāryaḥ na dhātupratyayānām āyanādayaḥ bhavanti iti . \\
\noindent \textbf{(P\_7,1.2)} yadi hi syuḥ ṛteḥ chaṅ iti eva brūyāt . \\
\noindent \textbf{(P\_7,1.2)} siddhe vidhiḥ ārabhyamāṇaḥ jñāpakārthaḥ bhavati na ca ṛteḥ chaṅā sidhyati . \\
\noindent \textbf{(P\_7,1.2)} chaṅi sati valādilakṣaṇaḥ iṭ prasajyeta . \\
\noindent \textbf{(P\_7,1.2)} iṭi kṛte anāditvāt ādeśaḥ na syāt . \\
\noindent \textbf{(P\_7,1.2)} idam iha sampradhāryam . \\
\noindent \textbf{(P\_7,1.2)} iṭ kriyatām ādeśaḥ iti kim atra kartavyam . \\
\noindent \textbf{(P\_7,1.2)} paratvāt iḍāgamaḥ . \\
\noindent \textbf{(P\_7,1.2)} nityaḥ ādeśaḥ . \\
\noindent \textbf{(P\_7,1.2)} kṛte api iṭi prāpnoti akṛte api . \\
\noindent \textbf{(P\_7,1.2)} anityaḥ ādeśaḥ na hi kṛte iṭi prāpnoti . \\
\noindent \textbf{(P\_7,1.2)} kim kāraṇam . \\
\noindent \textbf{(P\_7,1.2)} anāditvāt . \\
\noindent \textbf{(P\_7,1.2)} antaraṅgaḥ tarhi ādeśaḥ . \\
\noindent \textbf{(P\_7,1.2)} kā antaraṅgatā . \\
\noindent \textbf{(P\_7,1.2)} idānīm eva hi uktam āyanādiṣu upadeśivadvacanam svarasiddhyartham iti . \\
\noindent \textbf{(P\_7,1.2)} tat etat ṛteḥ īyaṅvacanam jñāpakam eva na dhātupratyayānām āyanādayaḥ bhavanti iti . prātipadikavijñānāt ca pāṇineḥ siddham prātipadikavijñānāt ca bhagavataḥ pāṇineḥ ācāryasya siddham . \\
\noindent \textbf{(P\_7,1.2)} uṇādayaḥ avyutpannāni prātipadikāni . \\
\noindent \textbf{(P\_7,1.3)} jhādeśe dhātvantapratiṣedhaḥ jhādeśe dhātvantasya pratiṣedhaḥ vaktavyaḥ . \\
\noindent \textbf{(P\_7,1.3)} ujjhitā ujjhitum iti . \\
\noindent \textbf{(P\_7,1.3)} pratyayādhikārāt siddham . \\
\noindent \textbf{(P\_7,1.3)} pratyayagrahaṇam prakṛtam anuvartate . \\
\noindent \textbf{(P\_7,1.3)} kva prakṛtam . \\
\noindent \textbf{(P\_7,1.3)} āyaneyīnīyiyaḥ phaḍakhachaghām pratyayādīnām iti . \\
\noindent \textbf{(P\_7,1.3)} pratyayādhikārāt siddham iti cet anādeḥ ādeśavacanam . pratyayādhikārāt siddham iti cet anādeḥ ādeśaḥ vaktavyaḥ . \\
\noindent \textbf{(P\_7,1.3)} api naḥ śvaḥ vijaniṣyamāṇāḥ patibhiḥ saha śayāntai . \\
\noindent \textbf{(P\_7,1.3)} evam tarhi pratyayagrahaṇam anuvartate ādigrahaṇam nivṛttam . \\
\noindent \textbf{(P\_7,1.3)} katham punaḥ samāsanirdiṣṭānām ekadeśaḥ anuvartate ekadeśaḥ vā nivartate . \\
\noindent \textbf{(P\_7,1.3)} asamāsanirdeśāt siddham . asamāsanirdeśaḥ kariṣyate . \\
\noindent \textbf{(P\_7,1.3)} pratyayasya ādīnām iti . \\
\noindent \textbf{(P\_7,1.3)} saḥ tarhi asamāsanirdeśaḥ kartavyaḥ . \\
\noindent \textbf{(P\_7,1.3)} na kartavyaḥ . \\
\noindent \textbf{(P\_7,1.3)} kriyate nyāse eva . \\
\noindent \textbf{(P\_7,1.3)} katham . \\
\noindent \textbf{(P\_7,1.3)} avibhaktikaḥ nirdeśaḥ . \\
\noindent \textbf{(P\_7,1.3)} pratyaya ādīnām iti . \\
\noindent \textbf{(P\_7,1.3)} tatra śayāntai iti anakārāntatvāt aṅgasya ādbhāvapratiṣedhaḥ . tatra etasmin pratyayagrahaṇe anuvartamāne ādigrahaṇe nivṛtte śayāntai iti anakārāntatvāt aṅgasya ādbhāvaḥ prāpnoti tasya pratiṣedhaḥ vaktavyaḥ . \\
\noindent \textbf{(P\_7,1.3)} siddham anānantaryāt anakārāntena adbhāvanivṛttiḥ . siddham etat . \\
\noindent \textbf{(P\_7,1.3)} katham . \\
\noindent \textbf{(P\_7,1.3)} anānantaryāt anakārāntena adbhāvaḥ na bhaviṣyati . \\
\noindent \textbf{(P\_7,1.3)} katham kṛtvā coditam katham kṛtvā parihāraḥ . \\
\noindent \textbf{(P\_7,1.3)} anakārāntagrahaṇam pratyayaviśeṣaṇam iti kṛtvā coditam jhakāraviśeṣaṇam iti kṛtvā parihāraḥ . \\
\noindent \textbf{(P\_7,1.3)} yadi anakārāntagrahaṇam jhakāraviśeṣaṇam śerate atra na prāpnoti . \\
\noindent \textbf{(P\_7,1.3)} tatra ruṭi sanniyogavacanāt siddham . tatra ruṭi sanniyogaḥ kariṣyate . \\
\noindent \textbf{(P\_7,1.3)} kaḥ eṣaḥ yatnaḥ codyate sanniyogaḥ nāma . \\
\noindent \textbf{(P\_7,1.3)} cakāraḥ kartavyaḥ . \\
\noindent \textbf{(P\_7,1.3)} ruṭ ca . \\
\noindent \textbf{(P\_7,1.3)} kim ca . \\
\noindent \textbf{(P\_7,1.3)} yat ca anyat prāpnoti . \\
\noindent \textbf{(P\_7,1.3)} kim ca anyat prāpnoti . \\
\noindent \textbf{(P\_7,1.3)} adbhāvaḥ . \\
\noindent \textbf{(P\_7,1.3)} saḥ tarhi cakāraḥ kartavyaḥ . \\
\noindent \textbf{(P\_7,1.3)} na kartavyaḥ . \\
\noindent \textbf{(P\_7,1.3)} yogavibhāgaḥ kariṣyate . \\
\noindent \textbf{(P\_7,1.3)} śīṅaḥ . \\
\noindent \textbf{(P\_7,1.3)} śīṅaḥ uttarasya jhasya at bhavati . \\
\noindent \textbf{(P\_7,1.3)} tataḥ ruṭ . \\
\noindent \textbf{(P\_7,1.3)} ruṭ ca bhavati śīṅaḥ iti . \\
\noindent \textbf{(P\_7,1.3)} evam api paryāyaḥ prasajyeta . \\
\noindent \textbf{(P\_7,1.3)} evam tarhi acśabdasya ruṭam vakṣyāmi . \\
\noindent \textbf{(P\_7,1.3)} tat acśabdagrahaṇam kartavyam . \\
\noindent \textbf{(P\_7,1.3)} na kartavyam . \\
\noindent \textbf{(P\_7,1.3)} prakṛtam anuvartate . \\
\noindent \textbf{(P\_7,1.3)} kva prakṛtam . \\
\noindent \textbf{(P\_7,1.3)} at abhyastāt iti . \\
\noindent \textbf{(P\_7,1.3)} tat vai prathamānirdiṣṭam ṣaṣṭhīnirdiṣṭena ca iha arthaḥ . \\
\noindent \textbf{(P\_7,1.3)} śīṅaḥ iti eṣā pañcamī at iti prathamāyāḥ ṣaṣṭhī prakalpayiṣyati tasmāt iti uttarasya iti . \\
\noindent \textbf{(P\_7,1.6)} ruṭi dṛśiguṇapratiṣedhaḥ . \\
\noindent \textbf{(P\_7,1.6)} ruṭi dṛśiguṇaḥ prāpnoti . \\
\noindent \textbf{(P\_7,1.6)} adṛśran asya ketavaḥ iti . \\
\noindent \textbf{(P\_7,1.6)} tasya pratiṣedhaḥ vaktavyaḥ . \\
\noindent \textbf{(P\_7,1.6)} parasmin iti kṅiti ca iti pratiṣedhaḥ bhaviṣyati . \\
\noindent \textbf{(P\_7,1.6)} evam api adṛśram asya ketavaḥ iti atra prāpnoti . \\
\noindent \textbf{(P\_7,1.6)} evam tarhi pūrvāntaḥ kariṣyate . \\
\noindent \textbf{(P\_7,1.6)} pūrvānte śīṅaḥ guṇavidhiḥ pūrvānte śīṅaḥ guṇaḥ vidheyaḥ : śerate . \\
\noindent \textbf{(P\_7,1.6)} sūtram ca bhidyate . \\
\noindent \textbf{(P\_7,1.6)} yathānyāsam eva astu . \\
\noindent \textbf{(P\_7,1.6)} nanu ca uktam ruṭi dṛśiguṇapratiṣedhaḥ iti . \\
\noindent \textbf{(P\_7,1.6)} pūrvānte api eṣaḥ doṣaḥ . \\
\noindent \textbf{(P\_7,1.6)} katham . \\
\noindent \textbf{(P\_7,1.6)} ayam dṛśiguṇaḥ pratiṣedhaviṣaye ārabhyate saḥ yathā eva kṅiti ca iti etam pratiṣedham bādhate evam anupadhāyāḥ api prasajyeta . \\
\noindent \textbf{(P\_7,1.6)} tasmāt ubhābhyām dṛśeḥ akpratyayāntaram vaktavyam pitaram ca dṛśeyam mātaram ca dṛśeyam iti evam artham . \\
\noindent \textbf{(P\_7,1.6)} jhādeśāt āṭ leṭi jhādeśāt āṭ leṭi bhavati vipratiṣedhena . \\
\noindent \textbf{(P\_7,1.6)} jhādeśasya avakāśaḥ . \\
\noindent \textbf{(P\_7,1.6)} lunate lunatām alunata . \\
\noindent \textbf{(P\_7,1.6)} āṭaḥ avakāśaḥ . \\
\noindent \textbf{(P\_7,1.6)} patāti didyut . \\
\noindent \textbf{(P\_7,1.6)} udadhim cyāvayāti . \\
\noindent \textbf{(P\_7,1.6)} iha ubhayam prāpnoti . \\
\noindent \textbf{(P\_7,1.6)} api naḥ śvaḥ vijaniṣyamāṇāḥ patibhiḥ saha śayāntai . \\
\noindent \textbf{(P\_7,1.6)} āṭ leṭi bhavati vipratiṣedhena . \\
\noindent \textbf{(P\_7,1.6)} saḥ tarhi pūrvavipratiṣedhaḥ vaktavyaḥ . \\
\noindent \textbf{(P\_7,1.6)} na vā nityatvāt āṭaḥ . \\
\noindent \textbf{(P\_7,1.6)} na vā vaktavyaḥ . \\
\noindent \textbf{(P\_7,1.6)} kim kāraṇam . \\
\noindent \textbf{(P\_7,1.6)} nityatvāt āṭaḥ . \\
\noindent \textbf{(P\_7,1.6)} nitaḥ āḍāgamaḥ . \\
\noindent \textbf{(P\_7,1.6)} saḥ katham nityaḥ . \\
\noindent \textbf{(P\_7,1.6)} yadi anakārāntagrahaṇam jhakāraviśeṣaṇam . \\
\noindent \textbf{(P\_7,1.6)} atha hi pratyayaviśeṣaṇam jhādeśaḥ api nityaḥ . \\
\noindent \textbf{(P\_7,1.6)} antaraṅgalakṣaṇatvāt ca . \\
\noindent \textbf{(P\_7,1.6)} antaraṅgaḥ khalu api āḍāgamaḥ . \\
\noindent \textbf{(P\_7,1.6)} katham antaraṅgaḥ . \\
\noindent \textbf{(P\_7,1.6)} yadi prāk lādeśāt dhātvadhikāraḥ . \\
\noindent \textbf{(P\_7,1.6)} atha hi lādeśe dhātvadhikāraḥ anuvartate ubhayam samānāśrayam . \\
\noindent \textbf{(P\_7,1.6)} yadi eva anakārāntagrahaṇam pratyayaviśeṣaṇam atha api lādeśe dhātvadhikāraḥ anuvartate ubhayathā api pūrvavipratiṣedhena na arthaḥ . \\
\noindent \textbf{(P\_7,1.6)} katham . \\
\noindent \textbf{(P\_7,1.6)} bahulam chandasi iti evam atra śapaḥ luk na bhaviṣyati . \\
\noindent \textbf{(P\_7,1.6)} tatra anataḥ iti pratiṣedhaḥ bhaviṣyati . \\
\noindent \textbf{(P\_7,1.7-10)} KA\_III,244.7-12 Ro\_V,17.11-17 {1/15}     idam bahulam chandasi iti dviḥ kriyate . \\
\noindent \textbf{(P\_7,1.7-10)} KA\_III,244.7-12 Ro\_V,17.11-17 {2/15}     ekam śakyam akartum . \\
\noindent \textbf{(P\_7,1.7-10)} KA\_III,244.7-12 Ro\_V,17.11-17 {3/15}     katham . \\
\noindent \textbf{(P\_7,1.7-10)} KA\_III,244.7-12 Ro\_V,17.11-17 {4/15}     yadi tāvat pūrvam kriyate param na kariṣyate . \\
\noindent \textbf{(P\_7,1.7-10)} KA\_III,244.7-12 Ro\_V,17.11-17 {5/15}     ataḥ bhisaḥ ais iti atra bahulam chandasi iti etat anuvartiṣyate . \\
\noindent \textbf{(P\_7,1.7-10)} KA\_III,244.7-12 Ro\_V,17.11-17 {6/15}     atha param kriyate pūrvam na kariṣyate . \\
\noindent \textbf{(P\_7,1.7-10)} KA\_III,244.7-12 Ro\_V,17.11-17 {7/15}     bahulam chandasi iti atra ruṭ api anuvartiṣyate . \\
\noindent \textbf{(P\_7,1.7-10)} KA\_III,244.7-12 Ro\_V,17.11-17 {8/15}     aparaḥ āha : ubhe bahulagrahaṇe ekam chandograhaṇam śakyam akartum . \\
\noindent \textbf{(P\_7,1.7-10)} KA\_III,244.7-12 Ro\_V,17.11-17 {9/15}     katham . \\
\noindent \textbf{(P\_7,1.7-10)} KA\_III,244.7-12 Ro\_V,17.11-17 {10/15}     idam asti . \\
\noindent \textbf{(P\_7,1.7-10)} KA\_III,244.7-12 Ro\_V,17.11-17 {11/15}     vetteḥ vibhāṣā . \\
\noindent \textbf{(P\_7,1.7-10)} KA\_III,244.7-12 Ro\_V,17.11-17 {12/15}     tataḥ chandasi . \\
\noindent \textbf{(P\_7,1.7-10)} KA\_III,244.7-12 Ro\_V,17.11-17 {13/15}     chandasi ca vibhāṣā . \\
\noindent \textbf{(P\_7,1.7-10)} KA\_III,244.7-12 Ro\_V,17.11-17 {14/15}     tataḥ ataḥ bhisaḥ ais bhavati . \\
\noindent \textbf{(P\_7,1.7-10)} KA\_III,244.7-12 Ro\_V,17.11-17 {15/15}     chandasi vibhāṣā iti . \\
\noindent \textbf{(P\_7,1.9)} iha vṛkṣaiḥ plakṣaiḥ iti paratvāt ettvam prāpnoti . \\
\noindent \textbf{(P\_7,1.9)} aisbhāvaḥ idānīm kva bhaviṣyati . \\
\noindent \textbf{(P\_7,1.9)} kṛte ettve bhautapūrvyāt . \\
\noindent \textbf{(P\_7,1.9)} kṛte ettve bhūtapūrvamakārāntam iti ais bhaviṣyati . \\
\noindent \textbf{(P\_7,1.9)} ais tu nityaḥ tathā sati . \\
\noindent \textbf{(P\_7,1.9)} evam sati nityaḥ aisbhāvaḥ kṛte api ettve prāpnoti akṛte api prāpnoti . \\
\noindent \textbf{(P\_7,1.9)} nityatvāt aistve kṛte vihatanimittatvāt ettvam na bhaviṣyati . \\
\noindent \textbf{(P\_7,1.9)} ettvam bhisi paratvāt cet ataḥ ais kva bhaviṣyati . \\
\noindent \textbf{(P\_7,1.9)} kṛte ettve bhautapūrvyāt ais tu nityaḥ tathā sati . \\
\noindent \textbf{(P\_7,1.11)} imau dvau pratiṣedhau ucyete . \\
\noindent \textbf{(P\_7,1.11)} ubhau śakyau avaktum . \\
\noindent \textbf{(P\_7,1.11)} katham . \\
\noindent \textbf{(P\_7,1.11)} evam vakṣyāmi . \\
\noindent \textbf{(P\_7,1.11)} idamadasoḥ kāt iti . \\
\noindent \textbf{(P\_7,1.11)} tanniyamārtham bhaviṣyati . \\
\noindent \textbf{(P\_7,1.11)} idamadasoḥ kāt eva na anyataḥ iti . \\
\noindent \textbf{(P\_7,1.12)} kimartham inādeśaḥ ucyate na nādeśaḥ eva ucyeta . \\
\noindent \textbf{(P\_7,1.12)} kā rūpasiddhiḥ : vṛkṣeṇa plakṣeṇa . \\
\noindent \textbf{(P\_7,1.12)} ettve yogavibhāgaḥ kariṣyate . \\
\noindent \textbf{(P\_7,1.12)} katham . \\
\noindent \textbf{(P\_7,1.12)} idam asti . \\
\noindent \textbf{(P\_7,1.12)} bahuvacane jhali et osi ca . \\
\noindent \textbf{(P\_7,1.12)} tataḥ āṅi ca . \\
\noindent \textbf{(P\_7,1.12)} āṅi ca parataḥ ataḥ ettvam bhavati . \\
\noindent \textbf{(P\_7,1.12)} vṛkṣeṇa plakṣeṇa . \\
\noindent \textbf{(P\_7,1.12)} tataḥ āpaḥ sambuddhau ca . \\
\noindent \textbf{(P\_7,1.12)} āpaḥ āṅi ca osi ca iti . \\
\noindent \textbf{(P\_7,1.12)} na evam śakyam . \\
\noindent \textbf{(P\_7,1.12)} iha hi anena iti idrūpalopaḥ prasajyeta . \\
\noindent \textbf{(P\_7,1.12)} jhali lopaḥ kariṣyate . \\
\noindent \textbf{(P\_7,1.12)} na śakyaḥ jhali lopaḥ kartum . \\
\noindent \textbf{(P\_7,1.12)} iha hi doṣaḥ syāt . \\
\noindent \textbf{(P\_7,1.12)} ayā viṣṭā iti . \\
\noindent \textbf{(P\_7,1.12)} evam tarhi anlopāpavādaḥ vijñāsyate . \\
\noindent \textbf{(P\_7,1.12)} katham . \\
\noindent \textbf{(P\_7,1.12)} evam vakṣyāmi . \\
\noindent \textbf{(P\_7,1.12)} an ne ca api ca iti . \\
\noindent \textbf{(P\_7,1.12)} tat nakāragrahaṇam kartavyam . \\
\noindent \textbf{(P\_7,1.12)} na kartavyam . \\
\noindent \textbf{(P\_7,1.12)} kriyate nyāse eva . \\
\noindent \textbf{(P\_7,1.12)} lupanirdiṣṭaḥ nakāraḥ . \\
\noindent \textbf{(P\_7,1.12)} yadi evam na upadhāyāḥ iti dīrghatvam prāpnoti . \\
\noindent \textbf{(P\_7,1.12)} sautraḥ nirdeśaḥ . \\
\noindent \textbf{(P\_7,1.12)} atha vā napuṃsakanirdeśaḥ kariṣyate . \\
\noindent \textbf{(P\_7,1.12)} atha kimartham āt ucyate na at eva ucyeta . \\
\noindent \textbf{(P\_7,1.12)} kā rūpasiddhiḥ vṛkṣāt : plakṣāt . \\
\noindent \textbf{(P\_7,1.12)} savarṇadīrghatvena siddham . \\
\noindent \textbf{(P\_7,1.12)} na sidhyati . \\
\noindent \textbf{(P\_7,1.12)} ataḥ guṇe pararūpam iti pararūpatvam prāpnoti . \\
\noindent \textbf{(P\_7,1.12)} akāroccāraṇasāmarthyāt na bhaviṣyati . \\
\noindent \textbf{(P\_7,1.12)} yadi prāpnuvan vidhiḥ uccāraṇasāmarthyāt bādhyate savarṇadīrghatvam api na prāpnoti . \\
\noindent \textbf{(P\_7,1.12)} na eṣaḥ doṣaḥ . \\
\noindent \textbf{(P\_7,1.12)} yam vidhim prati upadeśaḥ anarthakaḥ saḥ vidhiḥ bādhyate yasya tu vidheḥ nimittam eva na asau bādhyate . \\
\noindent \textbf{(P\_7,1.12)} pararūpam prati akāroccāraṇam anarthakam savarṇadīrghatvasya punaḥ nimittam eva . \\
\noindent \textbf{(P\_7,1.13)} kim idam caturthyekavacanasya grahaṇam āhosvit saptamyekavacanasya grahaṇam . \\
\noindent \textbf{(P\_7,1.13)} kutaḥ sandehaḥ . \\
\noindent \textbf{(P\_7,1.13)} samānaḥ nirdeśaḥ . \\
\noindent \textbf{(P\_7,1.13)} caturthyekavacanasya grahaṇam . \\
\noindent \textbf{(P\_7,1.13)} katham jñāyate . \\
\noindent \textbf{(P\_7,1.13)} lakṣaṇapratipadoktayoḥ pratipadoktasya eva iti . \\
\noindent \textbf{(P\_7,1.13)} iha api tarhi caturthyekavacanasya grahaṇam syāt . \\
\noindent \textbf{(P\_7,1.13)} ṅeḥ ām nadyāmnībhyaḥ . \\
\noindent \textbf{(P\_7,1.13)} evam tarhi vyākhyānataḥ viśeṣapratipattiḥ na hi sandehāt alakṣaṇam iti iha caturthyekavacanays agrahaṇam vyākhyāsyāmaḥ tatra saptamyekavacanasya iti . \\
\noindent \textbf{(P\_7,1.14)} aśaḥ ekādiṣṭāt smāyādīnām upasaṅkhyānam aśaḥ ekādiṣṭāt smāyādīnām upasaṅkhyānam kartavyam . \\
\noindent \textbf{(P\_7,1.14)} atha u atra asmai . \\
\noindent \textbf{(P\_7,1.14)} atha u atra asmāt . \\
\noindent \textbf{(P\_7,1.14)} atha u atra asmin iti . \\
\noindent \textbf{(P\_7,1.14)} ekādeśe kṛte ataḥ iti smāyādayaḥ na prāpnuvanti . \\
\noindent \textbf{(P\_7,1.14)} kim punaḥ kāraṇam ekādeśaḥ tāvat bhavati na punaḥ smāyādayaḥ . \\
\noindent \textbf{(P\_7,1.14)} na paratvāt smāyādibhiḥ bhavitavyam . \\
\noindent \textbf{(P\_7,1.14)} na bhavitavyam . \\
\noindent \textbf{(P\_7,1.14)} kim kāraṇam . \\
\noindent \textbf{(P\_7,1.14)} nityatvāt ekādeśaḥ . \\
\noindent \textbf{(P\_7,1.14)} nityaḥ ekādeśaḥ . \\
\noindent \textbf{(P\_7,1.14)} kṛteṣu api smāyādiṣu prāpnoti akṛteṣu api . \\
\noindent \textbf{(P\_7,1.14)} nityatvāt ekādeśe kṛte ataḥ iti smāyādayaḥ na prāpnuvanti . \\
\noindent \textbf{(P\_7,1.14)} kim ucyate aśaḥ iti na iha api kartavyam . \\
\noindent \textbf{(P\_7,1.14)} atra asmai . \\
\noindent \textbf{(P\_7,1.14)} atra asmāt . \\
\noindent \textbf{(P\_7,1.14)} atra asmin iti . \\
\noindent \textbf{(P\_7,1.14)} ekādeśe kṛte ataḥ iti smāyādayaḥ na prāpnuvanti . \\
\noindent \textbf{(P\_7,1.14)} ānupūrvyā siddham etat . \\
\noindent \textbf{(P\_7,1.14)} na atra akṛteṣu smāyādiṣu halādiḥ vibhaktiḥ asti halādau cet rūpalopaḥ na ca akṛtae idrūpalope ekādeśaḥ prāpnoti . \\
\noindent \textbf{(P\_7,1.14)} tat ānupūryā siddham . \\
\noindent \textbf{(P\_7,1.14)} tat tarhi upasaṅkhyānam kartavyam . \\
\noindent \textbf{(P\_7,1.14)} na vā bahiraṅgalakṣaṇatvāt . \\
\noindent \textbf{(P\_7,1.14)} na vā kartavyam . \\
\noindent \textbf{(P\_7,1.14)} kim kāraṇam . \\
\noindent \textbf{(P\_7,1.14)} bahiraṅgalakṣaṇatvāt . \\
\noindent \textbf{(P\_7,1.14)} bahiraṅgalakṣaṇaḥ ekādeśaḥ . \\
\noindent \textbf{(P\_7,1.14)} antaraṅgāḥ smāyādayaḥ . \\
\noindent \textbf{(P\_7,1.14)} asiddham bahiraṅgam antaraṅge . \\
\noindent \textbf{(P\_7,1.17,} 20) KA\_III,246.23-247.3 Ro\_V,21.18-21 {1/9}     kimartham śībhāvaḥ śibhāvaḥ ca ucyate na śibhāvaḥ eva ucyeta . \\
\noindent \textbf{(P\_7,1.17,} 20) KA\_III,246.23-247.3 Ro\_V,21.18-21 {2/9}     kā rūpasiddhiḥ : te ye ke . \\
\noindent \textbf{(P\_7,1.17,} 20) KA\_III,246.23-247.3 Ro\_V,21.18-21 {3/9}     ādguṇena siddham . \\
\noindent \textbf{(P\_7,1.17,} 20) KA\_III,246.23-247.3 Ro\_V,21.18-21 {4/9}     na evam śakyam . \\
\noindent \textbf{(P\_7,1.17,} 20) KA\_III,246.23-247.3 Ro\_V,21.18-21 {5/9}     iha hi trapuṇī jatunī dīrghaśravaṇam na syāt . \\
\noindent \textbf{(P\_7,1.17,} 20) KA\_III,246.23-247.3 Ro\_V,21.18-21 {6/9}     evam tarhi śībhavaḥ eva ucyatām . \\
\noindent \textbf{(P\_7,1.17,} 20) KA\_III,246.23-247.3 Ro\_V,21.18-21 {7/9}     na evam śakyam . \\
\noindent \textbf{(P\_7,1.17,} 20) KA\_III,246.23-247.3 Ro\_V,21.18-21 {8/9}     iha hi kuṇḍāni vanāni iti hrasvasya śravaṇam na syāt . \\
\noindent \textbf{(P\_7,1.17,} 20) KA\_III,246.23-247.3 Ro\_V,21.18-21 {9/9}     tasmāt śībhāvaḥ śibhāvaḥ ca vaktavyaḥ . \\
\noindent \textbf{(P\_7,1.18)} kimarthaḥ ṅakāraḥ . \\
\noindent \textbf{(P\_7,1.18)} sāmānyagrahaṇārthaḥ . \\
\noindent \textbf{(P\_7,1.18)} au , iti ucyamāne prathamādvivacanasya eva syāt . \\
\noindent \textbf{(P\_7,1.18)} atha api auṭ iti ucyate evam api dvitīyādvivacanasya eva syāt . \\
\noindent \textbf{(P\_7,1.18)} asti prayojanam etat . \\
\noindent \textbf{(P\_7,1.18)} kim tarhi iti . \\
\noindent \textbf{(P\_7,1.18)} ṅitkāryam tu prāpnoti . \\
\noindent \textbf{(P\_7,1.18)} khaṭve māle . \\
\noindent \textbf{(P\_7,1.18)} yāṭ āpaḥ iti yāṭ prāpnoti . \\
\noindent \textbf{(P\_7,1.18)} na eṣaḥ doṣaḥ . \\
\noindent \textbf{(P\_7,1.18)} na evam vijñāyate ṅakāraḥ it asya saḥ ayam ṅit ṅiti iti . \\
\noindent \textbf{(P\_7,1.18)} katham tarhi . \\
\noindent \textbf{(P\_7,1.18)} ṅaḥ eva it ṅit ṅiti iti . \\
\noindent \textbf{(P\_7,1.18)} evam sati varṇagrahaṇam idam bhavati varṇagrahaṇeṣu ca etat bhavati yasmin vidhiḥ tadādau algrahaṇe iti . \\
\noindent \textbf{(P\_7,1.18)} na doṣaḥ bhavati . \\
\noindent \textbf{(P\_7,1.18)} atha vā varṇagrahaṇam idam bhavati na ca etat varṇagrahaṇeṣu bhavati : ananubandhakagrahaṇe na sānubandhakasya iti . \\
\noindent \textbf{(P\_7,1.18)} atha vā pūrvasūtranirdeśaḥ ayam pūrvasūtreṣu ca ye anubandhāḥ na taiḥ iha itkāryāṇi kriyante . \\
\noindent \textbf{(P\_7,1.18)} aukāraḥ ayam śīvidhau ṅit gṛhītaḥ ṅit ca asmākam na asti kaḥ ayam prakāraḥ . \\
\noindent \textbf{(P\_7,1.18)} sāmānyārthaḥ tasya ca āsañjane asmin ṅitkāryam te śyām prasaktam saḥ doṣaḥ . \\
\noindent \textbf{(P\_7,1.18)} ṅittve vidyāt varṇanirdeśamātram varṇe yat syāt tat ca vidyāt tadādau . \\
\noindent \textbf{(P\_7,1.18)} varṇaḥ ca ayam tena ṅittve api adoṣaḥ nirdeśaḥ ayam pūrvasūtreṇa vā syāt . \\
\noindent \textbf{(P\_7,1.21)} auśaghau . auśaghau iti vaktavyam . \\
\noindent \textbf{(P\_7,1.21)} kim idam aghau iti . \\
\noindent \textbf{(P\_7,1.21)} anuttarapade iti . \\
\noindent \textbf{(P\_7,1.21)} kim prayojanam . \\
\noindent \textbf{(P\_7,1.21)} iha mā bhūt . \\
\noindent \textbf{(P\_7,1.21)} aṣṭaputraḥ aṣṭabhāryaḥ iti . \\
\noindent \textbf{(P\_7,1.21)} astu luk tatra . \\
\noindent \textbf{(P\_7,1.21)} astu atra auśtvam luk bhaviṣyati . \\
\noindent \textbf{(P\_7,1.21)} ṣaḍbhyaḥ api evam prasajyate . \\
\noindent \textbf{(P\_7,1.21)} iha api tarhi prāpnoti . \\
\noindent \textbf{(P\_7,1.21)} aṣṭau tiṣṭhanti . \\
\noindent \textbf{(P\_7,1.21)} aṣṭau paśya iti . apavādaḥ . \\
\noindent \textbf{(P\_7,1.21)} apavādatvāt atra auśtvam lukam bādhiṣyate . \\
\noindent \textbf{(P\_7,1.21)} iha api tarhi bādheta . \\
\noindent \textbf{(P\_7,1.21)} aṣṭaputraḥ aṣṭabhāryaḥ . \\
\noindent \textbf{(P\_7,1.21)} yasya viṣaye . \\
\noindent \textbf{(P\_7,1.21)} yasya lukaḥ viṣaye auśtvam tasya apavādaḥ . \\
\noindent \textbf{(P\_7,1.21)} yaḥ vā tasmāt anantaraḥ . \\
\noindent \textbf{(P\_7,1.21)} atha vā anantarasya lukaḥ bādhakam bhaviṣyati . \\
\noindent \textbf{(P\_7,1.21)} kutaḥ etat . \\
\noindent \textbf{(P\_7,1.21)} anantarasya vidhiḥ vā bhavati pratiṣedhaḥ vā iti . \\
\noindent \textbf{(P\_7,1.21)} atha iha kasmāt na bhavati auśtvam . \\
\noindent \textbf{(P\_7,1.21)} aṣṭa tiṣṭhanti . \\
\noindent \textbf{(P\_7,1.21)} aṣṭa paśya iti . \\
\noindent \textbf{(P\_7,1.21)} ātvam yatra tu tatra auśtvam . \\
\noindent \textbf{(P\_7,1.21)} yatra eva ātvam tatra eva auśtvena bhavitavyam . \\
\noindent \textbf{(P\_7,1.21)} kutaḥ etat . \\
\noindent \textbf{(P\_7,1.21)} tathā hi asya grahaḥ kṛtaḥ . \\
\noindent \textbf{(P\_7,1.21)} tathā hi asya ātvabhūtasya grahaṇam kriyate . \\
\noindent \textbf{(P\_7,1.21)} aṣṭābhyaḥ iti . \\
\noindent \textbf{(P\_7,1.21)} nanu ca nityam ātvam . \\
\noindent \textbf{(P\_7,1.21)} etat eva jñāpayati ācāryaḥ vibhāṣātvam iti yat ayam ātvabhūtasya grahaṇam karoti . \\
\noindent \textbf{(P\_7,1.21)} aṣṭābhyaḥ iti . \\
\noindent \textbf{(P\_7,1.21)} itarathā hi aṣṭanaḥ iti eva brūyāt \\
\noindent \textbf{(P\_7,1.23)} svamoḥ luk tyadādibhyaḥ ca svamoḥ luk tyadādibhyaḥ ca iti vaktavyam . \\
\noindent \textbf{(P\_7,1.23)} iha api yathā syāt . \\
\noindent \textbf{(P\_7,1.23)} tat brāhmaṇakulam iti . \\
\noindent \textbf{(P\_7,1.23)} kṛte hi atve na luk bhavet . \\
\noindent \textbf{(P\_7,1.23)} atve kṛte luk na prāpnoti . \\
\noindent \textbf{(P\_7,1.23)} idam iha sampradhāryam : atvam kriyatām luk iti kim atra kartavyam . \\
\noindent \textbf{(P\_7,1.23)} paratvāt atvam . \\
\noindent \textbf{(P\_7,1.23)} nityaḥ luk . \\
\noindent \textbf{(P\_7,1.23)} kṛte api atve prāpnoti akṛte api . \\
\noindent \textbf{(P\_7,1.23)} anityaḥ luk na hi kṛte atve prāpnoti . \\
\noindent \textbf{(P\_7,1.23)} ataḥ am iti ambhāvena bhavitavyam . \\
\noindent \textbf{(P\_7,1.23)} tasmāt tyadādibhyaḥ ca iti vaktavyam . \\
\noindent \textbf{(P\_7,1.23)} idam vicāryate : śiśīlugnumvidhiṣu napuṃsakagrahaṇam śabdagrahaṇam vā syāt arthagrahaṇam vā iti . \\
\noindent \textbf{(P\_7,1.23)} kaḥ ca atra viśeṣaḥ . \\
\noindent \textbf{(P\_7,1.23)} śiśīlugnumvidhiṣu napuṃsakagrahaṇam śabdagrahaṇam cet anyapadārthe pratiṣedhaḥ . \\
\noindent \textbf{(P\_7,1.23)} śiśīlugnumvidhiṣu napuṃsakagrahaṇam śabdagrahaṇam cet anyapadārthe pratiṣedhaḥ vaktavyaḥ . \\
\noindent \textbf{(P\_7,1.23)} bahutrapuḥ bahutrapū bahutrapavaḥ iti . \\
\noindent \textbf{(P\_7,1.23)} astu tarhi arthagrahaṇam . \\
\noindent \textbf{(P\_7,1.23)} yadi arthagrahaṇam priyasakthnā brāhmaṇena iti anaṅ na prāpnoti . \\
\noindent \textbf{(P\_7,1.23)} astu tarhi śabdagrahaṇam eva . \\
\noindent \textbf{(P\_7,1.23)} nanu ca uktam śiśīlugnumbidhiṣu napuṃsakagrahaṇam cet anyapadārthe pratiṣedhaḥ iti . \\
\noindent \textbf{(P\_7,1.23)} siddham tu prakṛtārthaviśeṣaṇatvāt . \\
\noindent \textbf{(P\_7,1.23)} siddham etat . \\
\noindent \textbf{(P\_7,1.23)} katham . \\
\noindent \textbf{(P\_7,1.23)} prakṛtasya arthaḥ viśeṣyate . \\
\noindent \textbf{(P\_7,1.23)} kim ca prakṛtam . \\
\noindent \textbf{(P\_7,1.23)} aṅgam . \\
\noindent \textbf{(P\_7,1.23)} aṅgasya śiśīlugnumaḥ bhavanti napuṃsake vartamānasya . \\
\noindent \textbf{(P\_7,1.23)} katham priyasakthnā brāhmaṇena . \\
\noindent \textbf{(P\_7,1.23)} asthyādiṣu śabdagrahaṇam . \\
\noindent \textbf{(P\_7,1.23)} asthyādiṣu napuṃsakagrahaṇam śabdagrahaṇam draṣṭavyam . \\
\noindent \textbf{(P\_7,1.23)} yuktam punaḥ idam vicārayitum . \\
\noindent \textbf{(P\_7,1.23)} nanu anena asandigdhena arthagrahaṇena bhavitavyam na hi napuṃsakam nāma śabdaḥ asti . \\
\noindent \textbf{(P\_7,1.23)} kim tarhi ucyate asthyādiṣu śabdagrahaṇam iti . \\
\noindent \textbf{(P\_7,1.23)} atra api arthagrahaṇam eva . \\
\noindent \textbf{(P\_7,1.23)} atra etāvān sandehaḥ kva prakṛtasya arthaḥ viśeṣyate kva gṛhyamāṇasya iti . \\
\noindent \textbf{(P\_7,1.23)} śiśīlugnumvidhiṣu prakṛtasya arthaḥ viśeṣyate asthyādiṣu gṛhyamāṇasya . \\
\noindent \textbf{(P\_7,1.25)} adbhāve pūrvasavarṇapratiṣedhaḥ . \\
\noindent \textbf{(P\_7,1.25)} adbhāve pūrvasavarṇasya pratiṣedhaḥ vaktavyaḥ . \\
\noindent \textbf{(P\_7,1.25)} katarat tiṣṭhati , katarat paśya . \\
\noindent \textbf{(P\_7,1.25)} siddham anunāsikopadhatvāt . \\
\noindent \textbf{(P\_7,1.25)} siddham etat . \\
\noindent \textbf{(P\_7,1.25)} katham . \\
\noindent \textbf{(P\_7,1.25)} anunāsikopadhaḥ acśabdaḥ kariṣyate . \\
\noindent \textbf{(P\_7,1.25)} dukkaraṇāt vā . \\
\noindent \textbf{(P\_7,1.25)} atha vā dugḍatarādīnām iti vakṣyati . \\
\noindent \textbf{(P\_7,1.25)} ḍitkaraṇāt vā . \\
\noindent \textbf{(P\_7,1.25)} atha vā ḍid acchabdaḥ kariṣyate . \\
\noindent \textbf{(P\_7,1.25)} saḥ tarhi ḍakāraḥ kartavyaḥ . \\
\noindent \textbf{(P\_7,1.25)} na kartavyaḥ . \\
\noindent \textbf{(P\_7,1.25)} kriyate nyāse eva . \\
\noindent \textbf{(P\_7,1.25)} dviḍakārakaḥ nirdeśaḥ . \\
\noindent \textbf{(P\_7,1.25)} adḍḍatarādibhyaḥ iti \\
\noindent \textbf{(P\_7,1.26)} itarāt chandasi pratiṣedhaḥ ekatarāt sarvatra . \\
\noindent \textbf{(P\_7,1.26)} itarāt chandasi pratiṣedhaḥ ekatarāt sarvatra iti vaktavyam . \\
\noindent \textbf{(P\_7,1.26)} ekataram tiṣṭhati , ekataram paśya . \\
\noindent \textbf{(P\_7,1.26)} napuṃsakādeśebhyaḥ yuṣmadasmadoḥ vibhaktyādeśāḥ vipratiṣedhena . napuṃsakādeśebhyaḥ yuṣmadasmadoḥ vibhaktyādeśāḥ bhavanti vipratiṣedhena . \\
\noindent \textbf{(P\_7,1.26)} napuṃsakādeśānām avakāśaḥ . \\
\noindent \textbf{(P\_7,1.26)} trapu , trapuṇī , trapūṇi . \\
\noindent \textbf{(P\_7,1.26)} yuṣmadasmadoḥ vibhaktyādeśānām avakāśaḥ . \\
\noindent \textbf{(P\_7,1.26)} tvam brāhmaṇaḥ , aham brāhmaṇaḥ , yuvām brāhmaṇau , āvām brāhmaṇau , yūyām brāhmaṇāḥ vayam brāhmaṇāḥ . \\
\noindent \textbf{(P\_7,1.26)} iha ubhayam prāpnoti . \\
\noindent \textbf{(P\_7,1.26)} tvam brāhmaṇakulam , aham brāhmaṇakulam , yuvām brāhmaṇakule , āvām brāhmaṇakule , yūyam brāhmaṇakulāni , vayam brāhmaṇakulāni . \\
\noindent \textbf{(P\_7,1.26)} yuṣmadasmadoḥ vibhaktyādeśāḥ bhavanti vipratiṣedhena . \\
\noindent \textbf{(P\_7,1.26)} atha idānīm yuṣmadasmadoḥ vibhaktyādeśeṣu kṛteṣu punaḥprasaṅgāt śiśīlugnumvidhayaḥ kasmāt na bhavanti . \\
\noindent \textbf{(P\_7,1.26)} sakṛdgatau vipratiṣedhena yat bādhitam tat bādhitam eva iti \\
\noindent \textbf{(P\_7,1.27)} kimarthaḥ śakāraḥ . \\
\noindent \textbf{(P\_7,1.27)} sarvādeśārthaḥ . \\
\noindent \textbf{(P\_7,1.27)} śit sarvasya iti sarvādeśaḥ yathā syāt . \\
\noindent \textbf{(P\_7,1.27)} na etat asti prayojanam . \\
\noindent \textbf{(P\_7,1.27)} akriyamāṇe api śakāre alaḥ antyasya vidhayaḥ bhavanti iti antyasya akāre kṛte trayāṇām akārāṇām ataḥ guṇe pararūpatve siddham rūpam syāt : tava svam , mama svam . \\
\noindent \textbf{(P\_7,1.27)} yadi etat labhyeta kṛtam syāt . \\
\noindent \textbf{(P\_7,1.27)} tat tu na labhyam . \\
\noindent \textbf{(P\_7,1.27)} kim kāraṇam . \\
\noindent \textbf{(P\_7,1.27)} atra hi tasmāt iti uttarasya ādeḥ parasya iti akārasya prasajyeta . \\
\noindent \textbf{(P\_7,1.27)} ataḥ uttaram paṭhati . \\
\noindent \textbf{(P\_7,1.27)} ṅasaḥ ādeśe śitkaraṇānarthakyam akārasya akāravacanānarthakyāt . \\
\noindent \textbf{(P\_7,1.27)} ṅasaḥ ādeśe śitkaraṇam anarthakam . \\
\noindent \textbf{(P\_7,1.27)} kim kāraṇam . \\
\noindent \textbf{(P\_7,1.27)} akārasya akāravacanānarthakyāt . \\
\noindent \textbf{(P\_7,1.27)} akārasya akāravacane prayojanam na asti iti kṛtvā antareṇa śakāram sarvādeśaḥ bhaviṣyati . \\
\noindent \textbf{(P\_7,1.27)} arthavattvādeśe lopārtham . \\
\noindent \textbf{(P\_7,1.27)} arthavattvakārasya akāravacanam . \\
\noindent \textbf{(P\_7,1.27)} kaḥ arthaḥ . \\
\noindent \textbf{(P\_7,1.27)} ādeśe lopārtham . \\
\noindent \textbf{(P\_7,1.27)} yaḥ saḥ śeṣe lopaḥ ādeśe saḥ vijñāyate . \\
\noindent \textbf{(P\_7,1.27)} nanu ca ādeśaḥ yā vibhaktiḥ iti evam etat vijñāyate . \\
\noindent \textbf{(P\_7,1.27)} ādeśaḥ eṣā vibhaktiḥ . \\
\noindent \textbf{(P\_7,1.27)} katham . \\
\noindent \textbf{(P\_7,1.27)} sarve sarvapadādeśā dākṣīputrasya pāṇineḥ ekadeśavikāre hi nityatvam na upapadyate . \\
\noindent \textbf{(P\_7,1.27)} tasmāt śitkaraṇam . \\
\noindent \textbf{(P\_7,1.27)} tasmāt śakāraḥ kartavyaḥ . \\
\noindent \textbf{(P\_7,1.27)} na kartavyaḥ . \\
\noindent \textbf{(P\_7,1.27)} kriyate nyāse eva . \\
\noindent \textbf{(P\_7,1.27)} katham . \\
\noindent \textbf{(P\_7,1.27)} praśliṣṭanirdeśaḥ ayam . \\
\noindent \textbf{(P\_7,1.27)} a , a , a , iti . \\
\noindent \textbf{(P\_7,1.27)} saḥ anekāl śit sarvasya iti sarvasya bhaviṣyati \\
\noindent \textbf{(P\_7,1.28)} prathamayoḥ iti ucyate kayoḥ idam prathamayoḥ grahaṇam kim vibhaktyoḥ āhosvit pratyayayoḥ . \\
\noindent \textbf{(P\_7,1.28)} vibhaktyoḥ iti āha . \\
\noindent \textbf{(P\_7,1.28)} katham jñāyate . \\
\noindent \textbf{(P\_7,1.28)} anyatra api hi prathamayoḥ grahaṇe vibhaktyoḥ grahaṇam vijñāyate na pratyayayoḥ . \\
\noindent \textbf{(P\_7,1.28)} kva anyatra . \\
\noindent \textbf{(P\_7,1.28)} prathamayoḥ pūrvasavarṇaḥ iti . \\
\noindent \textbf{(P\_7,1.28)} asti kāraṇam yena tatra vibhaktyoḥ grahaṇam vijñāyate . \\
\noindent \textbf{(P\_7,1.28)} kim kāraṇam . \\
\noindent \textbf{(P\_7,1.28)} aci iti tatra vartate na ca ajādī prathamau staḥ . \\
\noindent \textbf{(P\_7,1.28)} nanu ca evam vijñāyate ajādī yau prathamau ajādīnām vā yau prathamau iti . \\
\noindent \textbf{(P\_7,1.28)} yat tarhi tasmāt śasaḥ naḥ puṃsi iti anukrāntam pūrvasavarṇam pratinirdiśati tajjñāpayati ācāryaḥ vibhaktyoḥ grahaṇam iti . \\
\noindent \textbf{(P\_7,1.28)} iha api ācāryapravṛttiḥ jñāpayati vibhaktyoḥ grahaṇam iti yat ayam śasaḥ na iti pratiṣedham śāsti . \\
\noindent \textbf{(P\_7,1.28)} na eṣaḥ pratiṣedhaḥ . \\
\noindent \textbf{(P\_7,1.28)} natvam etat vidhīyate . \\
\noindent \textbf{(P\_7,1.28)} siddham atra natvam tasmāt śasaḥ naḥ puṃsi iti . \\
\noindent \textbf{(P\_7,1.28)} yatra tena na sidhyati tadartham . \\
\noindent \textbf{(P\_7,1.28)} kva ca tena na sidhyati . \\
\noindent \textbf{(P\_7,1.28)} striyām napuṃsake ca . \\
\noindent \textbf{(P\_7,1.28)} yuṣmān brāhmaṇī paśya , asmān brāhmaṇī paśya , yuṣmān brāhmaṇakulāni paśya , asmān brāhmaṇakulāni paśya iti . \\
\noindent \textbf{(P\_7,1.28)} yat tarhi yuṣmadasmadoḥ anādeśe dvitīyāyām ca iti āha tat jñāpayati ācāryaḥ vibhaktyoḥ grahaṇam iti \\
\noindent \textbf{(P\_7,1.30)} kim ayam bhyamśabdaḥ āhosvit abhyamśabdaḥ . \\
\noindent \textbf{(P\_7,1.30)} kutaḥ sandehaḥ . \\
\noindent \textbf{(P\_7,1.30)} samānaḥ nirdeśaḥ . \\
\noindent \textbf{(P\_7,1.30)} kim ca ataḥ . \\
\noindent \textbf{(P\_7,1.30)} yadi tāvat bhyamśabdaḥ śeṣe lopaḥ ca antyasya etvam prāpnoti . \\
\noindent \textbf{(P\_7,1.30)} atha abhyamśabdaḥ śeṣe lopaḥ ca ṭilopaḥ udāttnivṛttisvaraḥ prāpnoti . \\
\noindent \textbf{(P\_7,1.30)} yathā icchasi tathā astu . \\
\noindent \textbf{(P\_7,1.30)} astu tāvat abhyamśabdaḥ śeṣe lopaḥ ca antyasya . \\
\noindent \textbf{(P\_7,1.30)} nanu ca uktam ettvam prāpnoti iti . \\
\noindent \textbf{(P\_7,1.30)} na eṣaḥ doṣaḥ . \\
\noindent \textbf{(P\_7,1.30)} aṅgavṛtte punaḥ vṛttau avidhiḥ niṣṭhitasya iti na bhaviṣyati . \\
\noindent \textbf{(P\_7,1.30)} atha vā punaḥ astu abhamśabdaḥ śeṣe lopaḥ ca ṭilopaḥ . \\
\noindent \textbf{(P\_7,1.30)} nanu ca uktam udāttanivṛttisvaraḥ prāpnoti iti . \\
\noindent \textbf{(P\_7,1.30)} na eṣaḥ doṣaḥ uktam etat ādau siddham iti \\
\noindent \textbf{(P\_7,1.33)} kimartham āmaḥ sasakārasya grahaṇam kriyate na āmaḥ ākam iti eva ucyeta . \\
\noindent \textbf{(P\_7,1.33)} kena idānīm sasakārasya bhaviṣyati . \\
\noindent \textbf{(P\_7,1.33)} āmaḥ suṭ ayam bhaktaḥ āmgrahaṇena grāhiṣyate . \\
\noindent \textbf{(P\_7,1.33)} ataḥ uttaram paṭhati . \\
\noindent \textbf{(P\_7,1.33)} sāmgrahaṇam yathāgṛhītasya ādeśavacanāt . \\
\noindent \textbf{(P\_7,1.33)} sāmgrahaṇam kriyate . \\
\noindent \textbf{(P\_7,1.33)} nirdiśyamānasya ādeśāḥ bhavanti iti evam sasakārasya na prāpnoti . \\
\noindent \textbf{(P\_7,1.33)} iṣyate ca syāt iti tat ca antareṇa yatnam na sidhyati iti sāmaḥ ākam . \\
\noindent \textbf{(P\_7,1.33)} evamartham idam ucyate . \\
\noindent \textbf{(P\_7,1.33)} na vā dviparyantānām akāravacanāt āmi sakārābhāvaḥ . \\
\noindent \textbf{(P\_7,1.33)} na vā etat prayojanam asti . \\
\noindent \textbf{(P\_7,1.33)} kim kāraṇam . \\
\noindent \textbf{(P\_7,1.33)} dviparyantānām akāravacanāt . \\
\noindent \textbf{(P\_7,1.33)} dviparyantānām hi tyadādīnam atvam ucyate tena āmi sakāraḥ na bhaviṣyati . \\
\noindent \textbf{(P\_7,1.33)} suṭpratiṣedhaḥ tu ādeśe lopavijñānāt . \\
\noindent \textbf{(P\_7,1.33)} suṭpratiṣedhaḥ tu vaktavyaḥ . \\
\noindent \textbf{(P\_7,1.33)} kim kāraṇam . \\
\noindent \textbf{(P\_7,1.33)} ādeśe lopavijñānāt . \\
\noindent \textbf{(P\_7,1.33)} yaḥ saḥ śeṣe lopaḥ ādeśe saḥ vijñāyate . \\
\noindent \textbf{(P\_7,1.33)} na vā ṭilopavacanāt ādeśe ṭāppratiṣedhārtham . \\
\noindent \textbf{(P\_7,1.33)} na vā suṭpratiṣedhaḥ vaktavyaḥ . \\
\noindent \textbf{(P\_7,1.33)} kim kāraṇam . \\
\noindent \textbf{(P\_7,1.33)} ṭilopavacanāt . \\
\noindent \textbf{(P\_7,1.33)} ādeśe yaḥ saḥ śeṣe lopaḥ ṭilopaḥ saḥ vaktavyaḥ . \\
\noindent \textbf{(P\_7,1.33)} kim prayojanam . \\
\noindent \textbf{(P\_7,1.33)} ṭāppratiṣedhārtham . \\
\noindent \textbf{(P\_7,1.33)} ṭāp mā bhūt iti . \\
\noindent \textbf{(P\_7,1.33)} saḥ tarhi ṭilopaḥ vaktavyaḥ . \\
\noindent \textbf{(P\_7,1.33)} na vā liṅgābhāvāt ṭilopavacanānarthakyam . \\
\noindent \textbf{(P\_7,1.33)} na vā vaktavyam . \\
\noindent \textbf{(P\_7,1.33)} kim kāraṇam . \\
\noindent \textbf{(P\_7,1.33)} liṅgābhāvāt . \\
\noindent \textbf{(P\_7,1.33)} aliṅge yuṣmadasmadī . \\
\noindent \textbf{(P\_7,1.33)} kim vaktavyam etat . \\
\noindent \textbf{(P\_7,1.33)} na hi . \\
\noindent \textbf{(P\_7,1.33)} katham anucyamānam gaṃsyate . \\
\noindent \textbf{(P\_7,1.33)} na hi asti viśeṣaḥ yuṣmadasmadoḥ striyām puṃsi napuṃsake vā . \\
\noindent \textbf{(P\_7,1.33)} asti kāraṇam yena etat evam bhavati . \\
\noindent \textbf{(P\_7,1.33)} kim kāraṇam . \\
\noindent \textbf{(P\_7,1.33)} yaḥ asau viśeṣavācī śabdaḥ tadasānnidhyāt . \\
\noindent \textbf{(P\_7,1.33)} aṅga hi bhavān tam uccārayatu gaṃsyate saḥ viśeṣaḥ . \\
\noindent \textbf{(P\_7,1.33)} nanu ca na etena evam bhavitavyam . \\
\noindent \textbf{(P\_7,1.33)} na hi śabdanimittakena nāma arthena bhavitavyam . \\
\noindent \textbf{(P\_7,1.33)} kim tarhi artha nimittakena nāma śabdena bhavitavyam . \\
\noindent \textbf{(P\_7,1.33)} tat etat evam dṛśyatām : artharūpam eva etat evañjātīyakam yena atra viśeṣaḥ na gamyate iti . \\
\noindent \textbf{(P\_7,1.33)} avaśyam ca etat evam vijñeyam . \\
\noindent \textbf{(P\_7,1.33)} yaḥ hi manyate yaḥ asau viśeṣavācī śabdaḥ tadasānnidhyāt atra viśeṣaḥ na gamyate iti iha api tasya viśeṣaḥ na gamyate : dṛṣat samit iti . \\
\noindent \textbf{(P\_7,1.33)} tasmāt suṭpratiṣedhaḥ tasmāt suṭpratiṣedhaḥ vaktavyaḥ sasakāragrahaṇam vā kartavyam . \\
\noindent \textbf{(P\_7,1.33)} atha kriyamāṇe api sasakāragrahaṇe kasmāt eva atra suṭ na bhavati . \\
\noindent \textbf{(P\_7,1.33)} sasakāragrahaṇasāmarthyāt bhāvinaḥ suṭaḥ ādeśaḥ vijñāyate \\
\noindent \textbf{(P\_7,1.34)} iha papau, tasthau iti trīṇi kāryāṇi yugapat prāpnuvanti : dvirvacanam ekādeśaḥ autvam iti . \\
\noindent \textbf{(P\_7,1.34)} tat yadi sarvataḥ autvam labhyeta kṛtam syāt . \\
\noindent \textbf{(P\_7,1.34)} atha api dvirvacanam labhyeta evam api kṛtam syāt . \\
\noindent \textbf{(P\_7,1.34)} tat tu na labhyam . \\
\noindent \textbf{(P\_7,1.34)} kim kāraṇam . \\
\noindent \textbf{(P\_7,1.34)} atra hi paratvāt ekādeśaḥ dvirvacanam bādhate . \\
\noindent \textbf{(P\_7,1.34)} paratvāt autvam . \\
\noindent \textbf{(P\_7,1.34)} nityaḥ ekādeśaḥ autvam bādheta . \\
\noindent \textbf{(P\_7,1.34)} kam punaḥ bhavān autvasya avakāśam matvā āha nityaḥ ekādeśaḥ iti . \\
\noindent \textbf{(P\_7,1.34)} anavakāśam autvam ekādeśam bādhiṣyate . \\
\noindent \textbf{(P\_7,1.34)} autve kṛte dvirvacanam ekādeśaḥ iti yadi api paratvāt ekādeśaḥ sthānivadbhāvāt dvirvacanam bhaviṣyati \\
\noindent \textbf{(P\_7,1.36)} videḥ vasoḥ kittvam . \\
\noindent \textbf{(P\_7,1.36)} videḥ vasoḥ kittvam vaktavyam . \\
\noindent \textbf{(P\_7,1.36)} kim prayojanam . \\
\noindent \textbf{(P\_7,1.36)} vasugrahaṇeṣu liḍādeśasya api grahaṇam yathā syāt . \\
\noindent \textbf{(P\_7,1.36)} kim ca kāraṇam na syāt . \\
\noindent \textbf{(P\_7,1.36)} ananubandhakagrahaṇe hi sānubandhakasya grahaṇam na iti evam liḍādeśasya na prāpnoti . \\
\noindent \textbf{(P\_7,1.36)} sānubandhakaḥ hi saḥ kriyate . \\
\noindent \textbf{(P\_7,1.36)} kim punaḥ kāraṇam saḥ sānubandhakaḥ kriyate . \\
\noindent \textbf{(P\_7,1.36)} ayam ṝkārāntānām liṭi guṅaḥ pratiṣedhaviṣayaḥ ārabhyate saḥ punaḥ kitkaraṇāt bādhyate . \\
\noindent \textbf{(P\_7,1.36)} ātistīrvān , nipupūrvān iti . \\
\noindent \textbf{(P\_7,1.36)} saḥ tarhi asya evamarthaḥ anubandhaḥ kartavyaḥ . \\
\noindent \textbf{(P\_7,1.36)} na kartavyaḥ . \\
\noindent \textbf{(P\_7,1.36)} kriyate nyāse eva . \\
\noindent \textbf{(P\_7,1.36)} dvisakārakaḥ nirdeśaḥ : videḥ śaturvasussamāse anañpūrve ktvaḥ lyap . \\
\noindent \textbf{(P\_7,1.37)} lyabādeśe upadeśivadvacanam . \\
\noindent \textbf{(P\_7,1.37)} lyabādeśe upadeśivadbhāvaḥ vaktavyaḥ . \\
\noindent \textbf{(P\_7,1.37)} upadeśāvasthāyām lyap bhavati iti vaktavyam . \\
\noindent \textbf{(P\_7,1.37)} kim prayojanam . \\
\noindent \textbf{(P\_7,1.37)} anādiṣṭārtham . \\
\noindent \textbf{(P\_7,1.37)} akṛteṣu ādeśeṣu lyap yathā syāt . \\
\noindent \textbf{(P\_7,1.37)} ke punaḥ ādeśāḥ upadeśivadvacanam prayojayanti . \\
\noindent \textbf{(P\_7,1.37)} hitvadattvāttvetvettvadīrghatvaśūḍitaḥ . \\
\noindent \textbf{(P\_7,1.37)} hitvam . \\
\noindent \textbf{(P\_7,1.37)} hitvā , pradhāya . \\
\noindent \textbf{(P\_7,1.37)} hitvam . \\
\noindent \textbf{(P\_7,1.37)} dattvam . \\
\noindent \textbf{(P\_7,1.37)} dattvā , pradāya . \\
\noindent \textbf{(P\_7,1.37)} dattvam . \\
\noindent \textbf{(P\_7,1.37)} āttvam . \\
\noindent \textbf{(P\_7,1.37)} khātvā , prakhanya . \\
\noindent \textbf{(P\_7,1.37)} āttvam . \\
\noindent \textbf{(P\_7,1.37)} itvam . \\
\noindent \textbf{(P\_7,1.37)} sthitvā , prasthāya . \\
\noindent \textbf{(P\_7,1.37)} itvam . \\
\noindent \textbf{(P\_7,1.37)} īttvam . \\
\noindent \textbf{(P\_7,1.37)} pītvā , prapāya . \\
\noindent \textbf{(P\_7,1.37)} īttvam . \\
\noindent \textbf{(P\_7,1.37)} dīrghatvam . \\
\noindent \textbf{(P\_7,1.37)} śāntvā , praśamya . \\
\noindent \textbf{(P\_7,1.37)} dīrghatvam . \\
\noindent \textbf{(P\_7,1.37)} śatvam . \\
\noindent \textbf{(P\_7,1.37)} pṛṣṭvā , āpṛcchya . \\
\noindent \textbf{(P\_7,1.37)} śatvam . \\
\noindent \textbf{(P\_7,1.37)} ūṭh . \\
\noindent \textbf{(P\_7,1.37)} dyūtvā , pradīvya . \\
\noindent \textbf{(P\_7,1.37)} ūṭh . \\
\noindent \textbf{(P\_7,1.37)} iṭ . \\
\noindent \textbf{(P\_7,1.37)} devitvā , pradīvya . \\
\noindent \textbf{(P\_7,1.37)} kim punaḥ kāraṇam ādeśāḥ tāvat bhavanti na punaḥ lyap . \\
\noindent \textbf{(P\_7,1.37)} na paratvāt lyapā bhavitavyam . \\
\noindent \textbf{(P\_7,1.37)} santi ca eva atra ke cit pare ādeśāḥ api ca bahiraṅgalakṣaṇatvāt . \\
\noindent \textbf{(P\_7,1.37)} bahiraṅgaḥ lyap . \\
\noindent \textbf{(P\_7,1.37)} antaraṅgāḥ ādeśāḥ . \\
\noindent \textbf{(P\_7,1.37)} asiddham bahiraṅgam antaraṅge . \\
\noindent \textbf{(P\_7,1.37)} saḥ tarhi upadeśivadbhāvaḥ vaktavyaḥ . \\
\noindent \textbf{(P\_7,1.37)} na vaktavyaḥ . \\
\noindent \textbf{(P\_7,1.37)} ācāryapravṛttiḥ jñāpayati antaraṅgān api vidhīn bahiraṅgaḥ lyap bādhate iti yat ayam adaḥ jagdhiḥ lyapti kiti iti ti kiti iti eva siddhe lyabgrahaṇam karoti . \\
\noindent \textbf{(P\_7,1.37)} snātvākālakādiṣu ca pratiṣedhaḥ snātvākālakādiṣu ca pratiṣedhaḥ vaktavyaḥ . \\
\noindent \textbf{(P\_7,1.37)} snātvākālakaḥ , pītvāsthirakaḥ , bhuktvāsuhitakaḥ iti . \\
\noindent \textbf{(P\_7,1.37)} tadantanirdeśāt siddham . \\
\noindent \textbf{(P\_7,1.37)} tadantanirdeśāt siddham etat . \\
\noindent \textbf{(P\_7,1.37)} katham . \\
\noindent \textbf{(P\_7,1.37)} ktvāntasya lyapā bhavitavyam na ca etat ktvāntam . \\
\noindent \textbf{(P\_7,1.37)} samāsanipātanāt vā . \\
\noindent \textbf{(P\_7,1.37)} atha vā avaśyam atra samāsārtham nipātanam kartavyam tena eva yatnena lyap api na bhaviṣyati . \\
\noindent \textbf{(P\_7,1.37)} anañaḥ vā parasya . \\
\noindent \textbf{(P\_7,1.37)} atha vā anañaḥ parasya lyapā bhavitavyam na ca atra anañam paśyāmaḥ . \\
\noindent \textbf{(P\_7,1.37)} nanu ca dhātuḥ eva anañ . \\
\noindent \textbf{(P\_7,1.37)} na dhātoḥ parasya bhavitavyam . \\
\noindent \textbf{(P\_7,1.37)} kim kāraṇam . \\
\noindent \textbf{(P\_7,1.37)} nañivayuktam anyasadṛśādhikaraṇe tathā hi arthagatiḥ . \\
\noindent \textbf{(P\_7,1.37)} nañyuktam iva yuktam vā anyasmin tatsadṛśe kāryam vijñāsyate . \\
\noindent \textbf{(P\_7,1.37)} kutaḥ etat . \\
\noindent \textbf{(P\_7,1.37)} tathā hi arthaḥ gamyate . \\
\noindent \textbf{(P\_7,1.37)} tat yathā . \\
\noindent \textbf{(P\_7,1.37)} abrāhmaṇam ānaya iti ukte grāhmaṇasadṛśam puruṣam ānayati na asau loṣṭam ānīya kṛtī bhavati . \\
\noindent \textbf{(P\_7,1.37)} evam iha api anañ iti nañpratiṣedhāt anyasmāt anañau nañsadṛśāt kāryam vijñāsyate . \\
\noindent \textbf{(P\_7,1.37)} kim ca anyat anañ nañsadṛśam . \\
\noindent \textbf{(P\_7,1.37)} padam iti āha . \\
\noindent \textbf{(P\_7,1.37)} atha vā pratyayagrahaṇe yasmāt saḥ tadādeḥ grahaṇam bhavati iti evam dhātuḥ api ktvāgrahaṇena grāhiṣyate . \\
\noindent \textbf{(P\_7,1.37)} nanu ca iyam api paribhāṣā asti : kṛdgrahaṇe gatikārakapūrvasya api grahaṇam bhavati iti sā api iha upatiṣṭhate . \\
\noindent \textbf{(P\_7,1.37)} tatra kaḥ doṣaḥ . \\
\noindent \textbf{(P\_7,1.37)} iha na syāt : prakṛtya prahṛtya . \\
\noindent \textbf{(P\_7,1.37)} kva tarhi syāt . \\
\noindent \textbf{(P\_7,1.37)} paramakṛtvā , uttamakṛtvā . \\
\noindent \textbf{(P\_7,1.37)} na vai atra iṣyate . \\
\noindent \textbf{(P\_7,1.37)} aniṣṭam ca prāpnoti iṣṭam ca na sidhyati . \\
\noindent \textbf{(P\_7,1.37)} gatikārakapūrvasya eva iṣyate . \\
\noindent \textbf{(P\_7,1.37)} kutaḥ na khalu etat dvayoḥ paribhāṣayoḥ sāvakāśayoḥ samavasthitayoḥ pratyayagrahaṇe yasmāt saḥ tadādeḥ grahaṇam bhavati kṛdgrahaṇe gatikārakapūrvasya iti ca iyam iha paribhāṣā bhavati pratyayagrahaṇe yasmāt saḥ tadādeḥ grahaṇam bhavati iti iyam na bhavati kṛdgrahaṇe gatikārakapūrvasya api iti . \\
\noindent \textbf{(P\_7,1.37)} ācāryapravṛttiḥ jñāpayati iyam iha paribhāṣā bhavati pratyayagrahaṇe iti iyam na bhavati kṛdgrahaṇe iti yat ayam anañ iti pratiṣedham śāsti . \\
\noindent \textbf{(P\_7,1.37)} katham kṛtvā jñāpakam . \\
\noindent \textbf{(P\_7,1.37)} ayam hi nañ na gatiḥ na ca kārakam tatra kaḥ prasaṅgaḥ yat nañpūrvasya syāt . \\
\noindent \textbf{(P\_7,1.37)} paśyati tu ācāryaḥ iyam iha paribhāṣā bhavati pratyayagrahaṇe iti iyam na bhavati kṛdgrahaṇe iti tataḥ anañ iti pratiṣedham śāsti . \\
\noindent \textbf{(P\_7,1.37)} kim nañaḥ pratiṣedhena na gatiḥ na ca kārakam yāvatā nañi pūrve tu lyabbhāvaḥ na bhaviṣyati . \\
\noindent \textbf{(P\_7,1.37)} pratiṣedhāt tu jānīmaḥ tatpūrvam na iha gṛhyate pratyaygrahaṇe yāvat tāvat bhavitum arhati \\
\noindent \textbf{(P\_7,1.39)} supām ca supaḥ bhavanti iti vaktavyam . \\
\noindent \textbf{(P\_7,1.39)} yuktā mātā āsīt bhuri dakṣiṇāyāḥ dakṣiṇāyām iti prāpte . \\
\noindent \textbf{(P\_7,1.39)} tiṅām ca tiṅaḥ bhavanti iti vaktavyam . \\
\noindent \textbf{(P\_7,1.39)} caṣālam ye , aśvayūpāya takṣati takṣanti iti prāpte . \\
\noindent \textbf{(P\_7,1.39)} luki kim udāharaṇam . \\
\noindent \textbf{(P\_7,1.39)} ārdre carman , lohite carman . \\
\noindent \textbf{(P\_7,1.39)} na etat asti . \\
\noindent \textbf{(P\_7,1.39)} pūrvasavarṇena api etat siddham . \\
\noindent \textbf{(P\_7,1.39)} idam tarhi . \\
\noindent \textbf{(P\_7,1.39)} yat sthavīyasaḥ āvasan uta sapta sākam . \\
\noindent \textbf{(P\_7,1.39)} nanu ca etat api pūrvasavarṇena eva siddham . \\
\noindent \textbf{(P\_7,1.39)} na sidhyati . \\
\noindent \textbf{(P\_7,1.39)} yadi atra pūrvasavarṇaḥ syāt tyadādyatvam prasajyeta . \\
\noindent \textbf{(P\_7,1.39)} idam ca api udāharaṇam . \\
\noindent \textbf{(P\_7,1.39)} ārdre carman , lohite carman . \\
\noindent \textbf{(P\_7,1.39)} nanu ca uktam pūrvasavarṇena api etat siddham iti . \\
\noindent \textbf{(P\_7,1.39)} na sidhyati . \\
\noindent \textbf{(P\_7,1.39)} yadi atra pūrvasavarṇaḥ syāt āntaryataḥ dakāraḥ prasajyeta . \\
\noindent \textbf{(P\_7,1.39)} astu . \\
\noindent \textbf{(P\_7,1.39)} saṃyogāntalopena siddham . \\
\noindent \textbf{(P\_7,1.39)} iyāḍiyājīkārāṇām upasaṅkhyānam . \\
\noindent \textbf{(P\_7,1.39)} iyāḍiyājīkārāṇām upasaṅkhyānam kartavyam . \\
\noindent \textbf{(P\_7,1.39)} dārviyā parijman . \\
\noindent \textbf{(P\_7,1.39)} iyā . \\
\noindent \textbf{(P\_7,1.39)} ḍiyāc . \\
\noindent \textbf{(P\_7,1.39)} sukṣetriyā , sugātuyā . \\
\noindent \textbf{(P\_7,1.39)} ḍiyāc . \\
\noindent \textbf{(P\_7,1.39)} īkāra . \\
\noindent \textbf{(P\_7,1.39)} dṛtim na śuṣkam sarasī śayānam . \\
\noindent \textbf{(P\_7,1.39)} āṅayājayārām ca upasaṅkhyānam kartavyam . \\
\noindent \textbf{(P\_7,1.39)} āṅ , pra bāhavā . \\
\noindent \textbf{(P\_7,1.39)} ayāc . \\
\noindent \textbf{(P\_7,1.39)} svapnayā sacase janam . \\
\noindent \textbf{(P\_7,1.39)} ayāc . \\
\noindent \textbf{(P\_7,1.39)} ayār . \\
\noindent \textbf{(P\_7,1.39)} saḥ naḥ sindhum iva nāvayā \\
\noindent \textbf{(P\_7,1.40)} kimarthaḥ śakāraḥ . \\
\noindent \textbf{(P\_7,1.40)} śit sarvasya iti \\
\noindent \textbf{(P\_7,1.50)} iha : ye pūrvāsaḥ , ye uparāsaḥ : āt jaseḥ asuk iti asuki kṛte jasaḥ grahaṇena grahaṇāt śībhāvaḥ prāpnoti . \\
\noindent \textbf{(P\_7,1.50)} evam tarhi jasi pūrvāntaḥ kariṣyate . \\
\noindent \textbf{(P\_7,1.50)} yadi pūrvāntaḥ kriyate kā rūpasiddhiḥ : brāhmaṇāsaḥ pitaraḥ somyāsaḥ . \\
\noindent \textbf{(P\_7,1.50)} savarṇadīghatvena siddham . \\
\noindent \textbf{(P\_7,1.50)} na sidhyati . \\
\noindent \textbf{(P\_7,1.50)} ataḥ guṇe iti pararūpatvam prāpnoti . \\
\noindent \textbf{(P\_7,1.50)} akāroccāraṇasāmarthyāt na bhaviṣyati . \\
\noindent \textbf{(P\_7,1.50)} yadi tarhi prāpnuvan vidhiḥ uccāraṇasāmarthyāt bādhyate savarṇadīrghatvam api na prāpnoti . \\
\noindent \textbf{(P\_7,1.50)} na eṣaḥ doṣaḥ . \\
\noindent \textbf{(P\_7,1.50)} yam vidhim prati upadeśaḥ anarthakaḥ saḥ vidhiḥ bādhyate yasya tu vidheḥ nimittam na asau bādhyate . \\
\noindent \textbf{(P\_7,1.50)} pararūpam ca prati akāroccāraṇam anarthakam savarṇadīrghatvasya punaḥ nimittam eva . \\
\noindent \textbf{(P\_7,1.50)} atha vā asuṭ kariṣyate . \\
\noindent \textbf{(P\_7,1.50)} evam api ye pūrvāsaḥ ye uparāsaḥ iti asuṭi kṛte jasaḥ grahaṇena grahaṇāt śībhāvaḥ prāpnoti . \\
\noindent \textbf{(P\_7,1.50)} na eṣaḥ doṣaḥ . \\
\noindent \textbf{(P\_7,1.50)} nirdiśyamānasya ādeśāḥ bhavanti iti evam asya na bhaviṣyati . \\
\noindent \textbf{(P\_7,1.50)} yaḥ tarhi nirdiśyate tasya kasmāt na bhavati . \\
\noindent \textbf{(P\_7,1.50)} asuṭā vyavahitatvāt . \\
\noindent \textbf{(P\_7,1.50)} sidhyati . \\
\noindent \textbf{(P\_7,1.50)} sūtram tarhi bhidyate . \\
\noindent \textbf{(P\_7,1.50)} yathānyāsam eva astu . \\
\noindent \textbf{(P\_7,1.50)} nanu ca uktam ye pūrvāsaḥ ye uparāsaḥ asuki kṛte jasaḥ grahaṇena grahaṇāt śībhāvaḥ prāpnoti iti . \\
\noindent \textbf{(P\_7,1.50)} na eṣaḥ doṣaḥ . \\
\noindent \textbf{(P\_7,1.50)} idam iha sampradhāryam . \\
\noindent \textbf{(P\_7,1.50)} śībhāvaḥ kriyatām asuk iti kim atra kartavyam . \\
\noindent \textbf{(P\_7,1.50)} paratvāt asuk . \\
\noindent \textbf{(P\_7,1.50)} atha idānīm asuki kṛte punaḥ prasaṅgavijñānāt śībhāvaḥ kasmāt na bhavati . \\
\noindent \textbf{(P\_7,1.50)} sakṛdgatau vipratiṣedhe yat bādhitam tat bādhitam eva iti . \\
\noindent \textbf{(P\_7,1.51)} aśvavṛṣayoḥ maithunecchāyām . \\
\noindent \textbf{(P\_7,1.51)} aśvavṛṣayoḥ maithunecchāyām iti vaktavyam . \\
\noindent \textbf{(P\_7,1.51)} aśvasyati vaḍavā , vṛṣasyati gauḥ . \\
\noindent \textbf{(P\_7,1.51)} maithunecchāyām iti kimartham . \\
\noindent \textbf{(P\_7,1.51)} aśvīyati , vṛṣīyati . \\
\noindent \textbf{(P\_7,1.51)} kṣīralavaṇayoḥ lālasāyām . \\
\noindent \textbf{(P\_7,1.51)} kṣīrlavaṇayoḥ lālasāyām iti vaktavyam . \\
\noindent \textbf{(P\_7,1.51)} kṣīrasyati māṇavakaḥ , lavaṇasyati uṣṭraḥ iti . \\
\noindent \textbf{(P\_7,1.51)} aparaḥ āha : sarvaprātipadikebhyaḥ lālasāyām iti vaktavyam dadhyasyati madhvasyati iti evamartham . \\
\noindent \textbf{(P\_7,1.51)} aparaḥ āha: suk vaktavyaḥ : dadhisyati madusyati iti evamartham \\
\noindent \textbf{(P\_7,1.52-54)} KA\_III,258.24-260.11 Ro\_V,51.16-54.14 {1/68}     ime bahavaḥ āmśabdāḥ : kāspratyayāt ām amantre liṭi . \\
\noindent \textbf{(P\_7,1.52-54)} KA\_III,258.24-260.11 Ro\_V,51.16-54.14 {2/68}     ṅasosām . \\
\noindent \textbf{(P\_7,1.52-54)} KA\_III,258.24-260.11 Ro\_V,51.16-54.14 {3/68}     kimettiṅavyayaghāt āmu adravyaprakarṣe . \\
\noindent \textbf{(P\_7,1.52-54)} KA\_III,258.24-260.11 Ro\_V,51.16-54.14 {4/68}     ṅeḥ ām nadyāmnībhyaḥ iti . \\
\noindent \textbf{(P\_7,1.52-54)} KA\_III,258.24-260.11 Ro\_V,51.16-54.14 {5/68}     kasya idam grahaṇam . \\
\noindent \textbf{(P\_7,1.52-54)} KA\_III,258.24-260.11 Ro\_V,51.16-54.14 {6/68}     ṣaṣṭhībahuvacanasya grahaṇam . \\
\noindent \textbf{(P\_7,1.52-54)} KA\_III,258.24-260.11 Ro\_V,51.16-54.14 {7/68}     atha asya kasmāt na bhavati kāspratyayāt ām amantre liṭi iti . \\
\noindent \textbf{(P\_7,1.52-54)} KA\_III,258.24-260.11 Ro\_V,51.16-54.14 {8/68}     ananubandhakagrahaṇe hi na sānubandhakasya iti . \\
\noindent \textbf{(P\_7,1.52-54)} KA\_III,258.24-260.11 Ro\_V,51.16-54.14 {9/68}     saḥ tarhi asya evamarthaḥ anubandhaḥ kartavyaḥ iha asya grahaṇam mā bhūt iti . \\
\noindent \textbf{(P\_7,1.52-54)} KA\_III,258.24-260.11 Ro\_V,51.16-54.14 {10/68}     nanu ca avaśyam makārasya itsañjñāparitrāṇārthaḥ anubandhaḥ kartavyaḥ . \\
\noindent \textbf{(P\_7,1.52-54)} KA\_III,258.24-260.11 Ro\_V,51.16-54.14 {11/68}     na arthaḥ itsañjñāparitrāṇārthena . \\
\noindent \textbf{(P\_7,1.52-54)} KA\_III,258.24-260.11 Ro\_V,51.16-54.14 {12/68}     itkāryābhāvāt atra itsañjñā na bhaviṣyati . \\
\noindent \textbf{(P\_7,1.52-54)} KA\_III,258.24-260.11 Ro\_V,51.16-54.14 {13/68}     idam asti itkāryam mit acaḥ antyāt paraḥ iti acām antyāt paraḥ yathā syāt . \\
\noindent \textbf{(P\_7,1.52-54)} KA\_III,258.24-260.11 Ro\_V,51.16-54.14 {14/68}     pratyayāntāt ayam vidhīyate tatra na asti viśeṣaḥ mit acaḥ antyāt paraḥ iti vā paratve pratyayaḥ paraḥ iti vā . \\
\noindent \textbf{(P\_7,1.52-54)} KA\_III,258.24-260.11 Ro\_V,51.16-54.14 {15/68}     yaḥ tarhi na pratyayāntāt ijādeśaḥ ca gurumataḥ anṛcchaḥ iti . \\
\noindent \textbf{(P\_7,1.52-54)} KA\_III,258.24-260.11 Ro\_V,51.16-54.14 {16/68}     atra api āskāsoḥ āmvacanam jñāpakam na ayam acām antyāt paraḥ bhavati iti . \\
\noindent \textbf{(P\_7,1.52-54)} KA\_III,258.24-260.11 Ro\_V,51.16-54.14 {17/68}     katham kṛtvā jñāpakam . \\
\noindent \textbf{(P\_7,1.52-54)} KA\_III,258.24-260.11 Ro\_V,51.16-54.14 {18/68}     na hi asti viśeṣaḥ āmi acām antyāt pare sati asati vā . \\
\noindent \textbf{(P\_7,1.52-54)} KA\_III,258.24-260.11 Ro\_V,51.16-54.14 {19/68}     ayam asti viśeṣaḥ . \\
\noindent \textbf{(P\_7,1.52-54)} KA\_III,258.24-260.11 Ro\_V,51.16-54.14 {20/68}     asati āmi dvirvacanena bhavitavyam sati na bhavitavyam . \\
\noindent \textbf{(P\_7,1.52-54)} KA\_III,258.24-260.11 Ro\_V,51.16-54.14 {21/68}     sati api bhavitavyam . \\
\noindent \textbf{(P\_7,1.52-54)} KA\_III,258.24-260.11 Ro\_V,51.16-54.14 {22/68}     katham . \\
\noindent \textbf{(P\_7,1.52-54)} KA\_III,258.24-260.11 Ro\_V,51.16-54.14 {23/68}     āmaḥ tanmadhyapatitatvāt dhātugrahaṇena grahaṇāt . \\
\noindent \textbf{(P\_7,1.52-54)} KA\_III,258.24-260.11 Ro\_V,51.16-54.14 {24/68}     tat etat kāsāsoḥ āmvacanam jñāpakam eva na ayam acām antyāt paraḥ bhavati iti . \\
\noindent \textbf{(P\_7,1.52-54)} KA\_III,258.24-260.11 Ro\_V,51.16-54.14 {25/68}     atha api katham cit kāryam syāt evam api na doṣaḥ . \\
\noindent \textbf{(P\_7,1.52-54)} KA\_III,258.24-260.11 Ro\_V,51.16-54.14 {26/68}     kriyate nyāse eva . \\
\noindent \textbf{(P\_7,1.52-54)} KA\_III,258.24-260.11 Ro\_V,51.16-54.14 {27/68}     āmaḥ amantre iti . \\
\noindent \textbf{(P\_7,1.52-54)} KA\_III,258.24-260.11 Ro\_V,51.16-54.14 {28/68}     yadi evam āmā antre iti prāpnoti . \\
\noindent \textbf{(P\_7,1.52-54)} KA\_III,258.24-260.11 Ro\_V,51.16-54.14 {29/68}     śakandhunyāyena nirdeśaḥ . \\
\noindent \textbf{(P\_7,1.52-54)} KA\_III,258.24-260.11 Ro\_V,51.16-54.14 {30/68}     atha vā astu asya grahaṇam kaḥ doṣaḥ . \\
\noindent \textbf{(P\_7,1.52-54)} KA\_III,258.24-260.11 Ro\_V,51.16-54.14 {31/68}     iha kārayāñcakāra , harayāñcakāra , cikīrṣāñcakāra , jihīrṣāñcakāra , hrasvandyāpaḥ nuṭ iti nuṭ prasajyeta . \\
\noindent \textbf{(P\_7,1.52-54)} KA\_III,258.24-260.11 Ro\_V,51.16-54.14 {32/68}     lopāyādeśayoḥ kṛtayoḥ na bhaviṣyati . \\
\noindent \textbf{(P\_7,1.52-54)} KA\_III,258.24-260.11 Ro\_V,51.16-54.14 {33/68}     idam iha sampradhāryam . \\
\noindent \textbf{(P\_7,1.52-54)} KA\_III,258.24-260.11 Ro\_V,51.16-54.14 {34/68}     lopāyādeśau kriyetām nuṭ iti kim atra kartavyam . \\
\noindent \textbf{(P\_7,1.52-54)} KA\_III,258.24-260.11 Ro\_V,51.16-54.14 {35/68}     paratvāt nuṭ . \\
\noindent \textbf{(P\_7,1.52-54)} KA\_III,258.24-260.11 Ro\_V,51.16-54.14 {36/68}     nityau lopāyādeśau . \\
\noindent \textbf{(P\_7,1.52-54)} KA\_III,258.24-260.11 Ro\_V,51.16-54.14 {37/68}     kṛte api nuṭi prāpnutaḥ akṛte api . \\
\noindent \textbf{(P\_7,1.52-54)} KA\_III,258.24-260.11 Ro\_V,51.16-54.14 {38/68}     tatra nityatvāt lopāyādeśayoḥ kṛtayoḥ vihatanimittatvāt nuṭ na bhaviṣyati . \\
\noindent \textbf{(P\_7,1.52-54)} KA\_III,258.24-260.11 Ro\_V,51.16-54.14 {39/68}     atha asya kasmāt na bhavati kimettiṅavyayaghāt āmu adravyaprakarṣe iti . \\
\noindent \textbf{(P\_7,1.52-54)} KA\_III,258.24-260.11 Ro\_V,51.16-54.14 {40/68}     ananubandhakagrahaṇe hi na sānubandhakasya iti . \\
\noindent \textbf{(P\_7,1.52-54)} KA\_III,258.24-260.11 Ro\_V,51.16-54.14 {41/68}     saḥ tarhi evamarthaḥ anubandhaḥ kartavyaḥ . \\
\noindent \textbf{(P\_7,1.52-54)} KA\_III,258.24-260.11 Ro\_V,51.16-54.14 {42/68}     nanu ca avaśyam ugitkāryārthaḥ anubandhaḥ kartavyaḥ . \\
\noindent \textbf{(P\_7,1.52-54)} KA\_III,258.24-260.11 Ro\_V,51.16-54.14 {43/68}     na arthaḥ ugitkāryārthena anubandhena . \\
\noindent \textbf{(P\_7,1.52-54)} KA\_III,258.24-260.11 Ro\_V,51.16-54.14 {44/68}     liṅgavibhaktiprakaraṇe sarvam ugitkāryam na ca āmaḥ liṅgavibhaktī staḥ . \\
\noindent \textbf{(P\_7,1.52-54)} KA\_III,258.24-260.11 Ro\_V,51.16-54.14 {45/68}     avyayam eṣaḥ . \\
\noindent \textbf{(P\_7,1.52-54)} KA\_III,258.24-260.11 Ro\_V,51.16-54.14 {46/68}     makārasya tarhi itsañjñāparitrāṇārthaḥ anubandhaḥ kartavyaḥ . \\
\noindent \textbf{(P\_7,1.52-54)} KA\_III,258.24-260.11 Ro\_V,51.16-54.14 {47/68}     itkāryābhāvāt atra itsañjñā na bhaviṣyati . \\
\noindent \textbf{(P\_7,1.52-54)} KA\_III,258.24-260.11 Ro\_V,51.16-54.14 {48/68}     idam asti itkāryam mit acaḥ antyāt paraḥ iti acām antyāt paraḥ yathā syāt . \\
\noindent \textbf{(P\_7,1.52-54)} KA\_III,258.24-260.11 Ro\_V,51.16-54.14 {49/68}     na etat asti . \\
\noindent \textbf{(P\_7,1.52-54)} KA\_III,258.24-260.11 Ro\_V,51.16-54.14 {50/68}     ghāntāt ayam vidhīyate tatra na asti viśeṣaḥ mit acaḥ antyāt paraḥ iti vā paratve pratyayaḥ paraḥ iti vā paratve . \\
\noindent \textbf{(P\_7,1.52-54)} KA\_III,258.24-260.11 Ro\_V,51.16-54.14 {51/68}     atha api katham cit itkāryam syāt evam api na doṣaḥ . \\
\noindent \textbf{(P\_7,1.52-54)} KA\_III,258.24-260.11 Ro\_V,51.16-54.14 {52/68}     kriyate nyāse eva . \\
\noindent \textbf{(P\_7,1.52-54)} KA\_III,258.24-260.11 Ro\_V,51.16-54.14 {53/68}     atha vā astu asya grahaṇam kaḥ doṣaḥ . \\
\noindent \textbf{(P\_7,1.52-54)} KA\_III,258.24-260.11 Ro\_V,51.16-54.14 {54/68}     iha pacatitarām , jalpatitarām , hrasvanadyāpaḥ nuṭ iti nuṭ prasajyeta . \\
\noindent \textbf{(P\_7,1.52-54)} KA\_III,258.24-260.11 Ro\_V,51.16-54.14 {55/68}     lope kṛte na bhaviṣyati . \\
\noindent \textbf{(P\_7,1.52-54)} KA\_III,258.24-260.11 Ro\_V,51.16-54.14 {56/68}     idam iha sampradhāryam . \\
\noindent \textbf{(P\_7,1.52-54)} KA\_III,258.24-260.11 Ro\_V,51.16-54.14 {57/68}     lopaḥ kriyatām nuṭ iti kim atra kartavyam . \\
\noindent \textbf{(P\_7,1.52-54)} KA\_III,258.24-260.11 Ro\_V,51.16-54.14 {58/68}     paratvāt nuṭ . \\
\noindent \textbf{(P\_7,1.52-54)} KA\_III,258.24-260.11 Ro\_V,51.16-54.14 {59/68}     evam tarhi hrasvanadyāpaḥ nuṭ iti atra yasya iti lopaḥ anuvartiṣyate . \\
\noindent \textbf{(P\_7,1.52-54)} KA\_III,258.24-260.11 Ro\_V,51.16-54.14 {60/68}     atha asya kasmāt na bhavati ṅeḥ ām nadyāmnībhyaḥ iti . \\
\noindent \textbf{(P\_7,1.52-54)} KA\_III,258.24-260.11 Ro\_V,51.16-54.14 {61/68}     kim ca syāt . \\
\noindent \textbf{(P\_7,1.52-54)} KA\_III,258.24-260.11 Ro\_V,51.16-54.14 {62/68}     kumāryām , kiśoryām , khaṭvāyām , mālāyām , tasyām , yasyām iti hrasvanadyāpaḥ nuṭ iti nuṭ prasajyeta . \\
\noindent \textbf{(P\_7,1.52-54)} KA\_III,258.24-260.11 Ro\_V,51.16-54.14 {63/68}     āḍyāṭsyāṭaḥ atra bādhakāḥ bhaviṣyanti . \\
\noindent \textbf{(P\_7,1.52-54)} KA\_III,258.24-260.11 Ro\_V,51.16-54.14 {64/68}     idam iha sampradhāryam . \\
\noindent \textbf{(P\_7,1.52-54)} KA\_III,258.24-260.11 Ro\_V,51.16-54.14 {65/68}     āḍyāṭsyāṭaḥ kriyantām nuṭ iti kim atra kartavyam . \\
\noindent \textbf{(P\_7,1.52-54)} KA\_III,258.24-260.11 Ro\_V,51.16-54.14 {66/68}     paratvāt āḍyāṭsyāṭaḥ . \\
\noindent \textbf{(P\_7,1.52-54)} KA\_III,258.24-260.11 Ro\_V,51.16-54.14 {67/68}     atha idānīm āḍyāṭsyāṭsu kṛteṣu punaḥprasaṅgāt nuṭ kasmāt na bhavati . \\
\noindent \textbf{(P\_7,1.52-54)} KA\_III,258.24-260.11 Ro\_V,51.16-54.14 {68/68}     sakṛdgatau vipratiṣedhe yat bādhitam tat bādhitam eva iti . \\
\noindent \textbf{(P\_7,1.56)} ayam yogaḥ śakyaḥ avaktum . \\
\noindent \textbf{(P\_7,1.56)} katham śrīṇām udāraḥ dharuṇaḥ rayīṇām , api tatra sūtagrāmaṇīnām . \\
\noindent \textbf{(P\_7,1.56)} iha tāvat śrīṇām udāraḥ dharuṇaḥ rayīṇām vibhāṣā āmi nadīsañjñā sā chandasi vyavasthitavibhāṣā bhaviṣyati . \\
\noindent \textbf{(P\_7,1.56)} api tatra sūtagrāmaṇīnām iti sūtāḥ ca grāmaṇyaḥ ca sūtagrāmaṇi tatra hrasvanadyāpaḥ nuṭ iti eva siddham \\
\noindent \textbf{(P\_7,1.58)} atha dhātoḥ iti kimartham . \\
\noindent \textbf{(P\_7,1.58)} abhaitsīt , acchaitsīt . \\
\noindent \textbf{(P\_7,1.58)} numvidhau upadeśivadvacanam pratyayavidhyartham . \\
\noindent \textbf{(P\_7,1.58)} numvidhau upadeśivadbhāvaḥ vaktavyaḥ . \\
\noindent \textbf{(P\_7,1.58)} upadeśāvasthāyām num bhavati iti vaktavyam . \\
\noindent \textbf{(P\_7,1.58)} kim prayojanam . \\
\noindent \textbf{(P\_7,1.58)} pratyayavidhyartham . \\
\noindent \textbf{(P\_7,1.58)} upadeśāvasthāyām numi kṛte iṣṭaḥ pratyayavidhiḥ yathā syāt . \\
\noindent \textbf{(P\_7,1.58)} kuṇḍā , huṇḍā iti . \\
\noindent \textbf{(P\_7,1.58)} itarathā hi anakāre pratyayaḥ . \\
\noindent \textbf{(P\_7,1.58)} akriyamāṇe hi upadeśivadbhāve anakāre yaḥ pratyayaḥ prāpnoti saḥ tāvat syāt tasmin avasthite num . \\
\noindent \textbf{(P\_7,1.58)} tatra kaḥ doṣaḥ . \\
\noindent \textbf{(P\_7,1.58)} tatra ayatheṣṭaprasaṅgaḥ . \\
\noindent \textbf{(P\_7,1.58)} tatra ayatheṣṭam prasajyeta . \\
\noindent \textbf{(P\_7,1.58)} aniṣṭe pratyaye avasthite num . \\
\noindent \textbf{(P\_7,1.58)} aniṣṭasya pratyayasya śravaṇam prasajyeta . \\
\noindent \textbf{(P\_7,1.58)} dhātugrahaṇasāmarthyāt vā tadupadeśe numvidhānam dhātugrahaṇasāmarthyāt vā tadupadeśe dhātūpadeśe num bhaviṣyati . \\
\noindent \textbf{(P\_7,1.58)} nanu ca anyat dhātugrahaṇasya prayojanam uktam . \\
\noindent \textbf{(P\_7,1.58)} kim . \\
\noindent \textbf{(P\_7,1.58)} abhaitsīt , acchaitsīt iti . \\
\noindent \textbf{(P\_7,1.58)} na etat asti prayojanam . \\
\noindent \textbf{(P\_7,1.58)} prayojanam nāma tat vaktavyam yat niyogataḥ syāt . \\
\noindent \textbf{(P\_7,1.58)} yat ca atra ikāreṇa kriyate akāreṇa api tat śakyam kartum \\
\noindent \textbf{(P\_7,1.59)} śe tṛmpādīnām . \\
\noindent \textbf{(P\_7,1.59)} śe tṛmpādīnām upasaṅkhyānam kartavyam . \\
\noindent \textbf{(P\_7,1.59)} tṛmpati , tṛmphati . \\
\noindent \textbf{(P\_7,1.59)} kimartham idam . \\
\noindent \textbf{(P\_7,1.59)} na numanuṣaktāḥ eva ete paṭhyante . \\
\noindent \textbf{(P\_7,1.59)} luptanakāratvāt . \\
\noindent \textbf{(P\_7,1.59)} lupyate atra nakāraḥ aniditām halaḥ upadhāyāḥ kṅiti iti . \\
\noindent \textbf{(P\_7,1.59)} yadi punaḥ ime iditaḥ paṭhyeran . \\
\noindent \textbf{(P\_7,1.59)} na evam śakyam . \\
\noindent \textbf{(P\_7,1.59)} iha hi lopaḥ na syāt . \\
\noindent \textbf{(P\_7,1.59)} tṛpitaḥ , dṛpitaḥ iti . \\
\noindent \textbf{(P\_7,1.59)} yadi punaḥ ime mucādiṣu paṭhyeran . \\
\noindent \textbf{(P\_7,1.59)} na doṣaḥ syāt . \\
\noindent \textbf{(P\_7,1.59)} atha vā na evam vijñāyate iditaḥ num dhātoḥ iti . \\
\noindent \textbf{(P\_7,1.59)} katham tarhi . \\
\noindent \textbf{(P\_7,1.59)} iditaḥ num . \\
\noindent \textbf{(P\_7,1.59)} tataḥ dhātoḥ iti \\
\noindent \textbf{(P\_7,1.62)} imau dvau pratiṣedhau ucyete . \\
\noindent \textbf{(P\_7,1.62)} ubhau śakyau avaktum . \\
\noindent \textbf{(P\_7,1.62)} katham . \\
\noindent \textbf{(P\_7,1.62)} evam vakṣyāmi . \\
\noindent \textbf{(P\_7,1.62)} iṭi liṭi radheḥ num bhavati iti . \\
\noindent \textbf{(P\_7,1.62)} tanniyamārtham bhaviṣyati . \\
\noindent \textbf{(P\_7,1.62)} liṭi eva iḍādau na anyasmin iḍādau iti \\
\noindent \textbf{(P\_7,1.65)} iha kasmāt na bhavati : ālabhyate . \\
\noindent \textbf{(P\_7,1.65)} astu . \\
\noindent \textbf{(P\_7,1.65)} aniditām halaḥ upadhāyāḥ kṅiti iti lopaḥ bhaviṣyati . \\
\noindent \textbf{(P\_7,1.65)} iha tarhi ālambhyā gauḥ poḥ adupadhāt iti yati avasthite num . \\
\noindent \textbf{(P\_7,1.65)} tatra kaḥ doṣaḥ . \\
\noindent \textbf{(P\_7,1.65)} ālambhyā* eṣaḥ svaraḥ prasajyeta . \\
\noindent \textbf{(P\_7,1.65)} ālambhyā* iti ca iṣyate . \\
\noindent \textbf{(P\_7,1.65)} na eṣaḥ doṣaḥ . \\
\noindent \textbf{(P\_7,1.65)} uktam etat dhātugrahaṇasāmarthyāt upadeśe numvidhānam iti \\
\noindent \textbf{(P\_7,1.68)} atha kevalagrahaṇam kimartham na na sudurbhyām iti eva ucyeta . \\
\noindent \textbf{(P\_7,1.68)} suduroḥ kevalagrahaṇam anyopasargapratiṣedhārtham . \\
\noindent \textbf{(P\_7,1.68)} suduroḥ kevalagrahaṇam kriyate anyopasṛṣṭāt mā bhūt iti . \\
\noindent \textbf{(P\_7,1.68)} prasulambham . \\
\noindent \textbf{(P\_7,1.68)} na eṣaḥ asti prayogaḥ . \\
\noindent \textbf{(P\_7,1.68)} idam tarhi . \\
\noindent \textbf{(P\_7,1.68)} supralambham . \\
\noindent \textbf{(P\_7,1.68)} preṇa vyavahitatvāt na bhaviṣyati . \\
\noindent \textbf{(P\_7,1.68)} idam tarhi . \\
\noindent \textbf{(P\_7,1.68)} atisulambham . \\
\noindent \textbf{(P\_7,1.68)} karmapravacanīyasañjñā atra bādhikā bhaviṣyati suḥ pūjāyām atiḥ atikramaṇe ca iti . \\
\noindent \textbf{(P\_7,1.68)} yadā tarhi na atikramaṇam na pūjā . \\
\noindent \textbf{(P\_7,1.68)} idam ca api udāharaṇam . \\
\noindent \textbf{(P\_7,1.68)} supralambham . \\
\noindent \textbf{(P\_7,1.68)} nanu ca uktam preṇa vyavahitatvāt na bhaviṣyati iti . \\
\noindent \textbf{(P\_7,1.68)} na eṣaḥ doṣaḥ . \\
\noindent \textbf{(P\_7,1.68)} sudurbhyām iti na eṣā pañcamī . \\
\noindent \textbf{(P\_7,1.68)} kā tarhi . \\
\noindent \textbf{(P\_7,1.68)} tṛtīyā . \\
\noindent \textbf{(P\_7,1.68)} sudurbhyām upasṛṣṭasya iti . \\
\noindent \textbf{(P\_7,1.68)} vyavahitaḥ ca api upasṛṣṭaḥ bhavati \\
\noindent \textbf{(P\_7,1.69)} ciṇṇamuloḥ anupasargasya . \\
\noindent \textbf{(P\_7,1.69)} ciṇṇamuloḥ anupasargasya iti vaktavyam . \\
\noindent \textbf{(P\_7,1.69)} iha mā bhūt . \\
\noindent \textbf{(P\_7,1.69)} prālambhi . \\
\noindent \textbf{(P\_7,1.69)} pralambham pralambham . \\
\noindent \textbf{(P\_7,1.69)} tat tarhi vaktavyam . \\
\noindent \textbf{(P\_7,1.69)} na vaktavyam . \\
\noindent \textbf{(P\_7,1.69)} iha upasargāt iti api prakṛtam na iti api tatra abhisambandhamātram kartavyam : vibhāṣā ciṇṇamuloḥ upasargāt na iti \\
\noindent \textbf{(P\_7,1.70)} adhātoḥ iti kimartham . \\
\noindent \textbf{(P\_7,1.70)} ukhāsrat , parṇadhvat . \\
\noindent \textbf{(P\_7,1.70)} adhātoḥ iti śakyam avaktum . \\
\noindent \textbf{(P\_7,1.70)} kasmāt na bhavati khāsrat , parṇadhvat iti . \\
\noindent \textbf{(P\_7,1.70)} ugiti añcatigrahaṇāt siddham adhātoḥ . \\
\noindent \textbf{(P\_7,1.70)} ugiti añcatigrahaṇāt adhātoḥ siddham . \\
\noindent \textbf{(P\_7,1.70)} añcatigrahaṇam niyamārtham bhaviṣyati . \\
\noindent \textbf{(P\_7,1.70)} añcateḥ eva ugitaḥ dhātoḥ na anyasya ugitaḥ dhātoḥ iti . \\
\noindent \textbf{(P\_7,1.70)} idam tarhi prayojanam adhātubhūtapūrvasya api yathā syāt . \\
\noindent \textbf{(P\_7,1.70)} gomantam icchati gomatyati gomatyateḥ apratyayaḥ gomān iti \\
\noindent \textbf{(P\_7,1.72)} jhalacaḥ numvidhau ugitpratiṣedhaḥ . \\
\noindent \textbf{(P\_7,1.72)} jhalacaḥ numvidhau ugillakṣaṇasya pratiṣedhaḥ vaktavyaḥ . \\
\noindent \textbf{(P\_7,1.72)} gomanti brāhmaṇakulāni , śreyāṃsi , bhūyāṃsi . \\
\noindent \textbf{(P\_7,1.72)} nanu ca jhallakṣaṇaḥ ugillakṣaṇam bādhiṣyate . \\
\noindent \textbf{(P\_7,1.72)} katham anyasya ucyamānam anyasya bādhakam syāt . \\
\noindent \textbf{(P\_7,1.72)} asati khalu api sambhave bādhanam bhavati asti ca sambhavaḥ yat ubhayam syāt . \\
\noindent \textbf{(P\_7,1.72)} kim ca syāt yadi atra ugillakṣaṇaḥ api syāt . \\
\noindent \textbf{(P\_7,1.72)} dvayoḥ nakārayoḥ śravaṇam prasajyeta . \\
\noindent \textbf{(P\_7,1.72)} na vyañjanaparasya ekasya vā aneakasya vā śravaṇam prati viśeṣaḥ asti . \\
\noindent \textbf{(P\_7,1.72)} nanu ca pratijñābhedaḥ bhavati . \\
\noindent \textbf{(P\_7,1.72)} śrutibhede asti kim pratijñābhedaḥ kariṣyati . \\
\noindent \textbf{(P\_7,1.72)} nanu ca śrutikṛtaḥ api bhedaḥ asti . \\
\noindent \textbf{(P\_7,1.72)} iha tāvat śreyāṃsi , bhūyāṃsi iti parasya anusvāre kṛte pūrvasya śravaṇam prāpnoti . \\
\noindent \textbf{(P\_7,1.72)} tathā kurvanti , kṛṣanti iti parasya anusvāraparasavarṇayoḥ kṛtayoḥ pūrvasya ṇatvam prāpnoti . \\
\noindent \textbf{(P\_7,1.72)} atha ekasmin api numi ṇatvam kasmāt na bhavati . \\
\noindent \textbf{(P\_7,1.72)} anusvārībhūtaḥ ṇatvam atikrāmati . \\
\noindent \textbf{(P\_7,1.72)} kṛte tarhi parasavarṇe kasmāt na bhavati . \\
\noindent \textbf{(P\_7,1.72)} asiddhe ca parasavarṇaḥ . \\
\noindent \textbf{(P\_7,1.72)} vipratiṣedhāt siddham . \\
\noindent \textbf{(P\_7,1.72)} vipratiṣedhāt siddham etat . \\
\noindent \textbf{(P\_7,1.72)} jhallakṣaṇaḥ kriyatām ugillakṣaṇaḥ iti jhallakṣaṇaḥ bhaviṣyati vipratiṣedhena . \\
\noindent \textbf{(P\_7,1.72)} jhallakṣaṇasya avakāśaḥ . \\
\noindent \textbf{(P\_7,1.72)} sarpīṃṣi , dhanūṃṣi . \\
\noindent \textbf{(P\_7,1.72)} ugillakṣaṇasya avakāśaḥ . \\
\noindent \textbf{(P\_7,1.72)} gomān , yavamān . \\
\noindent \textbf{(P\_7,1.72)} iha ubhayam prāpnoti . \\
\noindent \textbf{(P\_7,1.72)} gomanti brāhmaṇakulāni , yavamanti brāhmaṇakulāni , śreyāṃsi , bhūyāṃsi iti . \\
\noindent \textbf{(P\_7,1.72)} jhallakṣaṇaḥ bhaviṣyati vipratiṣedhena . \\
\noindent \textbf{(P\_7,1.72)} nanu ca punaḥprasaṅgavijñānāt ugillakṣaṇaḥ prāpnoti . \\
\noindent \textbf{(P\_7,1.72)} punaḥprasaṅgaḥ iti cet amādibhiḥ tulyam . \\
\noindent \textbf{(P\_7,1.72)} punaḥ prasaṅgaḥ iti cet amādibhiḥ tulyam etat bhavati . \\
\noindent \textbf{(P\_7,1.72)} tat yathā . \\
\noindent \textbf{(P\_7,1.72)} yuṣmadasmadoḥ amādiṣu kṛteṣu punaḥprasaṅgāt śaśīlugnumaḥ na bhavanti . \\
\noindent \textbf{(P\_7,1.72)} evam jhallakṣaṇe kṛte punaḥprasaṅgāt ugillakṣaṇaḥ na bhaviṣyati . \\
\noindent \textbf{(P\_7,1.72)} yat api ucyate asati khalu api sambhave bādhanam bhavati asti ca sambhavaḥ yat ubhayam syāt iti sati api sambhave bādhanam bhavati . \\
\noindent \textbf{(P\_7,1.72)} tat yathā . \\
\noindent \textbf{(P\_7,1.72)} dadhi brāhmaṇebhyaḥ dīyatām takram kauṇḍinyāya iti sati api sambhave dadhidānasya takradānam nivartakam bhavati . \\
\noindent \textbf{(P\_7,1.72)} evam iha api sati api sambhave jhallakṣaṇaḥ ugillakṣaṇam bādhiṣyate . \\
\noindent \textbf{(P\_7,1.72)} atha vā astu atra ugillakṣaṇaḥ api . \\
\noindent \textbf{(P\_7,1.72)} nanu ca uktam cvayoḥ nakārayoḥ śravaṇam prasajyeta iti . \\
\noindent \textbf{(P\_7,1.72)} parihṛtam etat na vyañjanaparasya ekasya vā anekasya vā śravaṇam prati viśeṣaḥ asti . \\
\noindent \textbf{(P\_7,1.72)} nanu ca uktam pratijñābhedaḥ bhavati iti . \\
\noindent \textbf{(P\_7,1.72)} śrutibhede asati pratijñābhedaḥ kim kariṣyati . \\
\noindent \textbf{(P\_7,1.72)} nanu ca śrutikṛtaḥ api bhedaḥ uktaḥ iha tāvat śreyāṃsi , bhūyāṃsi iti parasya anusvāre kṛte pūrvasya śravaṇam prasajyeta kurvanti , kṛṣanti iti parasya anusvāraparasavarṇayoḥ kṛtayoḥ pūrvasya ṇatvam prāpnoti iti . \\
\noindent \textbf{(P\_7,1.72)} na eṣaḥ doṣaḥ . \\
\noindent \textbf{(P\_7,1.72)} ayogavāhānām aviśeṣeṇa upadeśaḥ coditaḥ . \\
\noindent \textbf{(P\_7,1.72)} tatra iha tāvat śreyāṃsi , bhūyāṃsi iti parasya anusvāre kṛte tasya jhalgrahaṇena grahaṇāt pūrvasya anusvāraḥ bhaviṣyati kurvanti , kṛṣanti iti parasya anusvāraparasavarṇayoḥ kṛtayoḥ tasya jhalgrahaṇena grahaṇāt pūrvasya anusvāraparasavarṇau bhaviṣyataḥ . \\
\noindent \textbf{(P\_7,1.72)} na eva vā punaḥ atra ugillakṣaṇaḥ prāpnoti . \\
\noindent \textbf{(P\_7,1.72)} kim kāraṇam . \\
\noindent \textbf{(P\_7,1.72)} midacaḥ antyāt paraḥ iti ucyate na ca dvayoḥ mitoḥ acām antyāt paratve sambhavaḥ asti . \\
\noindent \textbf{(P\_7,1.72)} katham tarhi imau dvau mitau acām antyāt parau staḥ . \\
\noindent \textbf{(P\_7,1.72)} bahvanaḍvāṃhi brāhmaṇakulāni iti . \\
\noindent \textbf{(P\_7,1.72)} vinimittau etau . \\
\noindent \textbf{(P\_7,1.72)} tatra bahūrji pratiṣedhaḥ . \\
\noindent \textbf{(P\_7,1.72)} tatra bahūrji pratiṣedhaḥ vaktavyaḥ . \\
\noindent \textbf{(P\_7,1.72)} bahūrji brāhmaṇakulāni iti . \\
\noindent \textbf{(P\_7,1.72)} antyāt pūrvam numam eke . \\
\noindent \textbf{(P\_7,1.72)} antyāt pūrvam numam eke icchanti . \\
\noindent \textbf{(P\_7,1.72)} kim aviśeṣeṇa āhosvit bahūrjau eva . \\
\noindent \textbf{(P\_7,1.72)} kim ca ataḥ . \\
\noindent \textbf{(P\_7,1.72)} yadi aviśeṣeṇa kāṣṭhatakṃṣi iti bhavitavyam . \\
\noindent \textbf{(P\_7,1.72)} atha bahūrjau eva kāṣṭhataṅkṣi iti bhavitavyam . \\
\noindent \textbf{(P\_7,1.72)} evam tarhi bahūrjau eva . \\
\noindent \textbf{(P\_7,1.72)} bahūrñji . \\
\noindent \textbf{(P\_7,1.72)} saḥ tarhi pratiṣedhaḥ vaktavyaḥ . \\
\noindent \textbf{(P\_7,1.72)} na vaktavyaḥ . \\
\noindent \textbf{(P\_7,1.72)} acaḥ iti eṣā pañcamī . \\
\noindent \textbf{(P\_7,1.72)} acaḥ uttaraḥ yaḥ jhal tadantasya napuṃsakasya numā bhavitavyam . \\
\noindent \textbf{(P\_7,1.72)} yaḥ ca atra acaḥ uttaraḥ na asau jhal na api tadantam napuṃsakam yadantam ca napuṃsakam na asau acaḥ uttaraḥ . \\
\noindent \textbf{(P\_7,1.72)} iha api tarhi na prāpnoti . \\
\noindent \textbf{(P\_7,1.72)} kāṣṭhataṅkṣi iti . \\
\noindent \textbf{(P\_7,1.72)} atra yaḥ acaḥ uttaraḥ jhal na tadantam napuṃsakam yadantam ca napuṃsakam na asau acaḥ uttaraḥ . \\
\noindent \textbf{(P\_7,1.72)} na etat asti . \\
\noindent \textbf{(P\_7,1.72)} jhaljātiḥ pratinirdiśyate . \\
\noindent \textbf{(P\_7,1.72)} acaḥ uttarā yā jhaljātiḥ iti . \\
\noindent \textbf{(P\_7,1.72)} yadi pañcamī kuṇḍāni , vanāni iti atra na prāpnoti . \\
\noindent \textbf{(P\_7,1.72)} eva tarhi ikaḥ aci vibhaktau iti atra acaḥ sarvanāmasthāne iti etat anuvartiṣyate . \\
\noindent \textbf{(P\_7,1.72)} evam api ṣaṣṭhyabhāvāt na prāpnoti . \\
\noindent \textbf{(P\_7,1.72)} sarvanāmasthāne iti eṣā saptamī acaḥ iti pañcamyāḥ ṣaṣṭhīm prakalpayiṣyati tasmin iti nirdiṣṭe pūrvasya iti \\
\noindent \textbf{(P\_7,1.73)} ajgrahaṇam kimartham . \\
\noindent \textbf{(P\_7,1.73)} ikaḥ aci vyañjane mā bhūt . \\
\noindent \textbf{(P\_7,1.73)} ikaḥ aci iti ucyate vyañjanādau mā bhūt . \\
\noindent \textbf{(P\_7,1.73)} trapubhyām , trapubhiḥ . \\
\noindent \textbf{(P\_7,1.73)} astu lopaḥ . \\
\noindent \textbf{(P\_7,1.73)} astu atra num . \\
\noindent \textbf{(P\_7,1.73)} nalopaḥ prātipadikāntasya iti nalopaḥ bhaviṣyati . \\
\noindent \textbf{(P\_7,1.73)} svaraḥ katham pañcatrapubhyām , pañcatrapubhyaḥ . \\
\noindent \textbf{(P\_7,1.73)} igante dvigau iti eṣaḥ svaraḥ na prāpnoti . \\
\noindent \textbf{(P\_7,1.73)} svaraḥ vai śrūyamāṇe api . \\
\noindent \textbf{(P\_7,1.73)} śrūyamāṇe api numi svaraḥ bhavati . \\
\noindent \textbf{(P\_7,1.73)} pañcatrapuṇā , pañcatrapuṇaḥ iti . \\
\noindent \textbf{(P\_7,1.73)} lupte kim na bhaviṣyati . \\
\noindent \textbf{(P\_7,1.73)} lupte idānīm kim na bhaviṣyati . \\
\noindent \textbf{(P\_7,1.73)} kim punaḥ kāraṇam śrūyamāne api numi svaraḥ bhavati . \\
\noindent \textbf{(P\_7,1.73)} saṅghātabhaktaḥ asau na utsahate avayavasya igantatām vihantum iti kṛtvā tataḥ śrūyamāṇe api numi svaraḥ bhavati . \\
\noindent \textbf{(P\_7,1.73)} idam tarhi . \\
\noindent \textbf{(P\_7,1.73)} atirābhyām , atirābhiḥ . \\
\noindent \textbf{(P\_7,1.73)} numi kṛte rāyaḥ hali iti ātvam na prāpnoti . \\
\noindent \textbf{(P\_7,1.73)} idam iha sampradhāryam . \\
\noindent \textbf{(P\_7,1.73)} num kriyatām ātvam iti kim atra kartavyam . \\
\noindent \textbf{(P\_7,1.73)} paratvāt ātvam . \\
\noindent \textbf{(P\_7,1.73)} iha tarhi priyatisṛbhyām priyatisṛbhiḥ numi kṛte tisṛbhāvaḥ na prāpnoti . \\
\noindent \textbf{(P\_7,1.73)} idam iha sampradhāryam . \\
\noindent \textbf{(P\_7,1.73)} num kriyatām tisṛbhāvaḥ iti kim atra kartavyam . \\
\noindent \textbf{(P\_7,1.73)} paratvāt tisṛbhāvaḥ . \\
\noindent \textbf{(P\_7,1.73)} atha idānīm tisṛbhāve kṛte punaḥprasaṅgāt num kasmāt na bhavati . \\
\noindent \textbf{(P\_7,1.73)} sakṛdgatau vipratiṣedhe yat bādhitam tat bādhitam eva iti . \\
\noindent \textbf{(P\_7,1.73)} ataḥ uttaram paṭhati . \\
\noindent \textbf{(P\_7,1.73)} ikaḥ aci vibhaktau ajgrahaṇam numnuṭoḥ vipratiṣedhārtham . \\
\noindent \textbf{(P\_7,1.73)} ikaḥ aci vibhaktau ajgrahaṇam kriyate numaḥ nuṭ vipratiṣedhena yathā syāt . \\
\noindent \textbf{(P\_7,1.73)} trapūṇām , jatūnām . \\
\noindent \textbf{(P\_7,1.73)} akriyamāṇe hi ajgrahaṇe nityanimittaḥ num . \\
\noindent \textbf{(P\_7,1.73)} kṛte api nuṭi prāpnoti akṛte api . \\
\noindent \textbf{(P\_7,1.73)} nityanimittatvāt numi kṛte nuṭaḥ abhāvaḥ syāt . \\
\noindent \textbf{(P\_7,1.73)} etat api na asti prayojanam . \\
\noindent \textbf{(P\_7,1.73)} kriyamāṇe api vā ajgrahaṇe avaśyam atra nuḍarthaḥ yatnaḥ kartavyaḥ . \\
\noindent \textbf{(P\_7,1.73)} pūrvavipratiṣedhaḥ vaktavyaḥ . \\
\noindent \textbf{(P\_7,1.73)} idam tarhi prayojanam nuṭi kṛte num mā bhūt iti . \\
\noindent \textbf{(P\_7,1.73)} kim ca syāt . \\
\noindent \textbf{(P\_7,1.73)} trapūṇām , jatūnām . \\
\noindent \textbf{(P\_7,1.73)} itarathā hi numaḥ nityanimittatvāt nuḍabhāvaḥ . \\
\noindent \textbf{(P\_7,1.73)} nāmi iti dīrghatvam na syāt . \\
\noindent \textbf{(P\_7,1.73)} mā bhūt evam . \\
\noindent \textbf{(P\_7,1.73)} nopadhāyāḥ iti evam bhaviṣyati . \\
\noindent \textbf{(P\_7,1.73)} iha tarhi śucīnām inhanpūṣāryamṇām śau sau ca iti asmāt niyamāt na prāpnoti dīrghatvam . \\
\noindent \textbf{(P\_7,1.73)} arthavadgrahaṇe na anarthakasya iti evam na bhaviṣyati . \\
\noindent \textbf{(P\_7,1.73)} na eṣā paribhāṣā iha śakyā vijñātum . \\
\noindent \textbf{(P\_7,1.73)} iha hi doṣaḥ syāt . \\
\noindent \textbf{(P\_7,1.73)} vāgmi iti . \\
\noindent \textbf{(P\_7,1.73)} evam tarhi lakṣaṇapratipadoktayoḥ pratipadoktasya eva iti . \\
\noindent \textbf{(P\_7,1.73)} uttarārtham ca . \\
\noindent \textbf{(P\_7,1.73)} uttarārtham tarhi ajgrahaṇam kartavyam . \\
\noindent \textbf{(P\_7,1.73)} asthidadhisakthyakṣṇām anaṅ udāttaḥ ajādau yathā syāt . \\
\noindent \textbf{(P\_7,1.73)} iha mā bhūt . \\
\noindent \textbf{(P\_7,1.73)} asthibhyām , asthibhiḥ iti . \\
\noindent \textbf{(P\_7,1.73)} yadi uttarārtham syāt tatra eva ayam ajgrahaṇam kurvīta . \\
\noindent \textbf{(P\_7,1.73)} iha kriyamāṇe yadi kim cit prayojanam asti tat ucyatām . \\
\noindent \textbf{(P\_7,1.73)} iha api kriyamāṇe prayojanam asti . \\
\noindent \textbf{(P\_7,1.73)} kim . \\
\noindent \textbf{(P\_7,1.73)} ajādau yathā syāt . \\
\noindent \textbf{(P\_7,1.73)} iha mā bhūt . \\
\noindent \textbf{(P\_7,1.73)} trapu , jatu . \\
\noindent \textbf{(P\_7,1.73)} etat api na asti prayojanam . \\
\noindent \textbf{(P\_7,1.73)} vibhaktau iti ucyate na ca atra vibhaktim paśyāmaḥ . \\
\noindent \textbf{(P\_7,1.73)} pratyayalakṣaṇena . \\
\noindent \textbf{(P\_7,1.73)} na lumatā aṅgasya iti pratyayalakṣaṇasya pratiṣedhaḥ . \\
\noindent \textbf{(P\_7,1.73)} evam tarhi siddhe sati yat ajgrahaṇam karoti tat jñāpayati ācāryaḥ bhavati iha kaḥ cit anyaḥ api prakāraḥ pratyayalakṣaṇam nāma iti . \\
\noindent \textbf{(P\_7,1.73)} kim etasya jñāpane prayojanam . \\
\noindent \textbf{(P\_7,1.73)} he trapu , he trapo . \\
\noindent \textbf{(P\_7,1.73)} atra guṇaḥ siddhaḥ bhavati iti . \\
\noindent \textbf{(P\_7,1.73)} ikaḥ aci vyañjane mā bhūt astu lopaḥ svaraḥ katham . \\
\noindent \textbf{(P\_7,1.73)} svaraḥ vai śrūyamāṇe api lupte kim na bhaviṣyati . \\
\noindent \textbf{(P\_7,1.73)} rāyātvam tisṛbhāvaḥ ca vyavadhānāt numā api nuṭ vācyaḥ uttarārtham tu iha kim cit trapaḥ iti \\
\noindent \textbf{(P\_7,1.74)} kim iha puṃvadbhāvena atidiśyate . \\
\noindent \textbf{(P\_7,1.74)} numpratiṣedhaḥ . \\
\noindent \textbf{(P\_7,1.74)} katham punaḥ puṃvat iti anena numpratiṣedhaḥ śakyaḥ vijñātum . \\
\noindent \textbf{(P\_7,1.74)} vatinirdeśaḥ ayam kāmacāraḥ ca vatinirdeśe vākyaśeṣam samarthayitum . \\
\noindent \textbf{(P\_7,1.74)} tat yathā : uśīnaravat madreṣu yavāḥ . \\
\noindent \textbf{(P\_7,1.74)} santi na santi iti . \\
\noindent \textbf{(P\_7,1.74)} mātṛvat asyāḥ kalāḥ . \\
\noindent \textbf{(P\_7,1.74)} santi na santi iti . \\
\noindent \textbf{(P\_7,1.74)} evam iha api puṃvat bhavati puṃvat na bhavati iti vākyaśeṣam samarthayiṣyāmahe . \\
\noindent \textbf{(P\_7,1.74)} yathā puṃsaḥ na num bhavati evam tṛtīyādiṣu bhāṣitapuṃskasya api na bhavati iti . \\
\noindent \textbf{(P\_7,1.74)} kim ucyate numpratiṣedhaḥ iti na punaḥ anyat api puṃsaḥ pratipadam kāryam ucyate yat tṛtīyādiṣu vibhaktiṣu ajādiṣu bhāṣitapuṃskasya atidiśyeta . \\
\noindent \textbf{(P\_7,1.74)} anārambhāt puṃsi . \\
\noindent \textbf{(P\_7,1.74)} na hi kim cit puṃsaḥ pratipadam kāryam ucyate yat tṛtīyādiṣu ajādiṣu bhāṣitapuṃskasya atidiśyeta . \\
\noindent \textbf{(P\_7,1.74)} num prakṛtaḥ tatra kim anyat śakyam vijñātum anyat ataḥ numpratiṣedhāt . \\
\noindent \textbf{(P\_7,1.74)} puṃvat iti numpratiṣedhaḥ cet guṇanābhāvanuḍauttvapratiṣedhaḥ . \\
\noindent \textbf{(P\_7,1.74)} puṃvat iti njumpratiṣedhaḥ cet guṇanābhāvanuḍauttvānām pratiṣedhaḥ vaktavyaḥ . \\
\noindent \textbf{(P\_7,1.74)} guṇa . \\
\noindent \textbf{(P\_7,1.74)} grāmaṇye brāhmaṇakulāya . \\
\noindent \textbf{(P\_7,1.74)} guṇa . \\
\noindent \textbf{(P\_7,1.74)} nābhāva . \\
\noindent \textbf{(P\_7,1.74)} grāmaṇyā brāhmaṇakulena . \\
\noindent \textbf{(P\_7,1.74)} nābhāva . \\
\noindent \textbf{(P\_7,1.74)} nuṭ . \\
\noindent \textbf{(P\_7,1.74)} grāmaṇyām brāhmaṇakulānām . \\
\noindent \textbf{(P\_7,1.74)} nuṭ . \\
\noindent \textbf{(P\_7,1.74)} auttvam . \\
\noindent \textbf{(P\_7,1.74)} grāmaṇyām brāhmaṇakule . \\
\noindent \textbf{(P\_7,1.74)} hrasvatvam apratiṣiddham hrasvāśrayāḥ ca ete vidhayaḥ prāpnuvanti . \\
\noindent \textbf{(P\_7,1.74)} hrasvābhāvārtham ca . \\
\noindent \textbf{(P\_7,1.74)} kim ca . \\
\noindent \textbf{(P\_7,1.74)} numpratiṣedhārtham ca . \\
\noindent \textbf{(P\_7,1.74)} katham punaḥ atra aprakṛtasya asaṃśabditasya hrasvatvasya pratiṣedhaḥ śakyaḥ vijñātum . \\
\noindent \textbf{(P\_7,1.74)} arthātideśāt siddham . \\
\noindent \textbf{(P\_7,1.74)} na evam vijñāyate bhāṣyate pumān anena śabdena saḥ ayam bhāṣitapuṃskaḥ bhāṣitapuṃskasya śabdasya puṃśabdaḥ bhavati iti . \\
\noindent \textbf{(P\_7,1.74)} katham tarhi . \\
\noindent \textbf{(P\_7,1.74)} bhāṣyate pumān asmin arthe saḥ ayam bhāṣitapuṃskaḥ bhāṣitapuṃskasya arthasya puṃvadarthaḥ bhavati iti . \\
\noindent \textbf{(P\_7,1.74)} taddhitalukpratiṣedhaḥ ca . \\
\noindent \textbf{(P\_7,1.74)} taddhitalukaḥ ca pratiṣedhaḥ vaktavyaḥ . \\
\noindent \textbf{(P\_7,1.74)} pīluḥ vṛkṣaḥ , pīlu phalam . \\
\noindent \textbf{(P\_7,1.74)} pīlunā , pilunaḥ iti . \\
\noindent \textbf{(P\_7,1.74)} na vā samānāyām ākṛtau bhāṣitapuṃskavijñānāt . \\
\noindent \textbf{(P\_7,1.74)} na vā vaktavyam . \\
\noindent \textbf{(P\_7,1.74)} kim kāraṇam . \\
\noindent \textbf{(P\_7,1.74)} samānāyām ākṛtau bhāṣitapuṃskavijñānāt . \\
\noindent \textbf{(P\_7,1.74)} samānāyām ākṛtau yat bhāṣitapuṃskam ākṛtyantare ca etat bhāṣitapuṃskam . \\
\noindent \textbf{(P\_7,1.74)} kim vaktavyam etat . \\
\noindent \textbf{(P\_7,1.74)} na hi . \\
\noindent \textbf{(P\_7,1.74)} katham anucyamānam gaṃsyate . \\
\noindent \textbf{(P\_7,1.74)} etat api arthanirdeśāt siddham \\
\noindent \textbf{(P\_7,1.77)} kim udāharaṇam . \\
\noindent \textbf{(P\_7,1.77)} akṣī te indra piṅgale . \\
\noindent \textbf{(P\_7,1.77)} na etat asti . \\
\noindent \textbf{(P\_7,1.77)} pūrvasavarṇena api etat siddham . \\
\noindent \textbf{(P\_7,1.77)} idam tarhi . \\
\noindent \textbf{(P\_7,1.77)} akṣībhyām te nāsikābhyām . \\
\noindent \textbf{(P\_7,1.77)} idam ca api udāharaṇam . \\
\noindent \textbf{(P\_7,1.77)} akṣī te indra piṅgale . \\
\noindent \textbf{(P\_7,1.77)} nanu ca uktam pūrvasavarṇena api etat siddham iti . \\
\noindent \textbf{(P\_7,1.77)} na sidhyati . \\
\noindent \textbf{(P\_7,1.77)} numā vyavahitatvāt pūrvasavarṇaḥ na prāpnoti . \\
\noindent \textbf{(P\_7,1.77)} chandasi napuṃsakasya puṃvadbhāvaḥ vaktavyaḥ madhoḥ gṛbhṇāmi , madhoḥ tṛptāḥ iva asataḥ iti evamartham . \\
\noindent \textbf{(P\_7,1.77)} puṃvadbhāvena numaḥ nivṛttiḥ numi nivṛtte pūrvasavarṇena eva siddham . \\
\noindent \textbf{(P\_7,1.77)} svarārthaḥ tarhi īkāraḥ vaktavyaḥ . \\
\noindent \textbf{(P\_7,1.77)} udāttasvaraḥ yathā syāt napuṃsakasvaraḥ mā bhūt iti . \\
\noindent \textbf{(P\_7,1.77)} nanu ca puṃvadbhāvātideśāt eva svaraḥ bhaviṣyati . \\
\noindent \textbf{(P\_7,1.77)} aśakyaḥ puṃvadbhāvātideśaḥ svare tantram āśrayitum . \\
\noindent \textbf{(P\_7,1.77)} iha hi doṣaḥ syāt : madhu asmin asti madhuḥ māsaḥ iti . \\
\noindent \textbf{(P\_7,1.77)} saḥ tarhi puṃvadbhāvaḥ vaktavyaḥ . \\
\noindent \textbf{(P\_7,1.77)} na vaktavyaḥ . \\
\noindent \textbf{(P\_7,1.77)} prakṛtam puṃvat iti vartate \\
\noindent \textbf{(P\_7,1.78)} kasya ayam pratiṣedhaḥ . \\
\noindent \textbf{(P\_7,1.78)} numaḥ iti āha . \\
\noindent \textbf{(P\_7,1.78)} tat numaḥ grahaṇam kartavyam . \\
\noindent \textbf{(P\_7,1.78)} na kartavyam . \\
\noindent \textbf{(P\_7,1.78)} prakṛtam anuvartate . \\
\noindent \textbf{(P\_7,1.78)} kva prakṛtam . \\
\noindent \textbf{(P\_7,1.78)} iditaḥ num dhātoḥ iti tat vā anekagrahaṇena vyavacchinnam aśakyam anuvartayitum . \\
\noindent \textbf{(P\_7,1.78)} evam tarhi sarvanāmasthāne iti varate sarvanāmasthāne yat prāpnoti tasya pratiṣedhaḥ . \\
\noindent \textbf{(P\_7,1.78)} tat vai bahutarakeṇa grahaṇena vyavacchinnam aśakyam anuvartayitum . \\
\noindent \textbf{(P\_7,1.78)} atha idānīm vyavahitam api śakyate anuvartayitum num eva anuvartya iha ihārtham uttarārtham ca . \\
\noindent \textbf{(P\_7,1.78)} iha ca eva pratiṣedhaḥ siddhaḥ bhavati iha ca āt śīnadyoḥ num iti numgrahaṇam na kartavyam bhavati \\
\noindent \textbf{(P\_7,1.80)} iha kasmāt na bhavati . \\
\noindent \textbf{(P\_7,1.80)} adatī , ghnatī , lunatī , punatī . \\
\noindent \textbf{(P\_7,1.80)} lope kṛte avarṇābhāvāt . \\
\noindent \textbf{(P\_7,1.80)} kim tarhi asmin yoge udāharaṇam . \\
\noindent \textbf{(P\_7,1.80)} yātī , yāntī . \\
\noindent \textbf{(P\_7,1.80)} atra api ekādeśe kṛte vyapavargābhāvāt na prāpnoti . \\
\noindent \textbf{(P\_7,1.80)} antādivadbhāvena vyapavargaḥ . \\
\noindent \textbf{(P\_7,1.80)} ubhayataḥ āśraye na antādivat . \\
\noindent \textbf{(P\_7,1.80)} na ubhayataḥ āśrayaḥ kariṣyate . \\
\noindent \textbf{(P\_7,1.80)} na evam vijñāyate avarṇāntāt śatuḥ num bhavati iti . \\
\noindent \textbf{(P\_7,1.80)} katham tarhi . \\
\noindent \textbf{(P\_7,1.80)} avarṇāt num bhavati tat cet avarṇam śatuḥ anantaram iti \\
\noindent \textbf{(P\_7,1.82)} anaḍuhaḥ sau āmpratiṣedhaḥ numaḥ anavakāśatvāt . \\
\noindent \textbf{(P\_7,1.82)} anaḍuhaḥ sau āmpratiṣedhaḥ prāpnoti . \\
\noindent \textbf{(P\_7,1.82)} kim kāraṇam . \\
\noindent \textbf{(P\_7,1.82)} numaḥ anavakāśatvāt . \\
\noindent \textbf{(P\_7,1.82)} anavakāśaḥ num āmam bādhate . \\
\noindent \textbf{(P\_7,1.82)} na vā avarṇopadhasya numvacanāt . \\
\noindent \textbf{(P\_7,1.82)} na vā eṣaḥ doṣaḥ . \\
\noindent \textbf{(P\_7,1.82)} kim kāraṇam . \\
\noindent \textbf{(P\_7,1.82)} avarṇopadhasya numvacanāt . \\
\noindent \textbf{(P\_7,1.82)} avarṇopadhasya numam vakṣyāmi . \\
\noindent \textbf{(P\_7,1.82)} tadarthagrahaṇam kartavyam . \\
\noindent \textbf{(P\_7,1.82)} na kartavyam . \\
\noindent \textbf{(P\_7,1.82)} prakṛtam anuvartate . \\
\noindent \textbf{(P\_7,1.82)} kva prakṛtam . \\
\noindent \textbf{(P\_7,1.82)} āt śīnadyoḥ num iti . \\
\noindent \textbf{(P\_7,1.82)} yadi tat anuvartate anaḍuhi yāvanti avarṇāni sarvebhyaḥ paraḥ num prāpnoti . \\
\noindent \textbf{(P\_7,1.82)} na eṣaḥ doṣaḥ . \\
\noindent \textbf{(P\_7,1.82)} mit acaḥ antyāt paraḥ iti anena yat sarvāntyam avarṇam tasmāt paraḥ bhaviṣyati . \\
\noindent \textbf{(P\_7,1.82)} punaḥprasaṅgavijñānāt vā siddham . \\
\noindent \textbf{(P\_7,1.82)} atha vā punaḥprsaṅgāt numi kṛte ām bhaviṣyati . \\
\noindent \textbf{(P\_7,1.82)} yathāttvādiṣu dvirvacanam . \\
\noindent \textbf{(P\_7,1.82)} tat yathā jagle , mamle , ījatuḥ , ījuḥ iti āttvādiṣu kṛteṣu punaḥprasaṅgāt dvirvacanam bhavati evam atra api numi kṛte ām bhaviṣyati . \\
\noindent \textbf{(P\_7,1.82)} na eṣaḥ yuktaḥ parihāraḥ . \\
\noindent \textbf{(P\_7,1.82)} vipratiṣedhe punaḥprasaṅgaḥ vipratiṣedhaḥ ca dvayoḥ sāvakāśayoḥ bhavati . \\
\noindent \textbf{(P\_7,1.82)} iha punaḥ anavakāśaḥ num āmam bādhate . \\
\noindent \textbf{(P\_7,1.82)} evam tarhi vṛttāntāt eṣaḥ parihāraḥ prasthitaḥ . \\
\noindent \textbf{(P\_7,1.82)} kasmāt vṛttāntāt . \\
\noindent \textbf{(P\_7,1.82)} idam ayam codyaḥ bhavati anaḍuhaḥ sau āmpratiṣedhaḥ numaḥ anavakāśatvāt iti . \\
\noindent \textbf{(P\_7,1.82)} tasya parihāraḥ na vā avarṇopadhasya numvacanāt iti . \\
\noindent \textbf{(P\_7,1.82)} tataḥ ayam codyaḥ bhavati yatra tarhi avarṇaprakaraṇam na asti tatra taḥ āmā numaḥ bādhanam prāpnoti bahvanḍvāṃhi brāhmaṇakulāni iti . \\
\noindent \textbf{(P\_7,1.82)} tataḥ uttarakālam idam paṭhitam punaḥprasaṅgavijñānāt vā siddham iti \\
\noindent \textbf{(P\_7,1.84)} divaḥ auttve dhātupratiṣedhaḥ . \\
\noindent \textbf{(P\_7,1.84)} divaḥ auttve dhātoḥ pratiṣedhaḥ vaktavyaḥ . \\
\noindent \textbf{(P\_7,1.84)} akṣadyūḥ iti . \\
\noindent \textbf{(P\_7,1.84)} adhātvadhikārāt siddham . \\
\noindent \textbf{(P\_7,1.84)} adhātoḥ iti vartate . \\
\noindent \textbf{(P\_7,1.84)} kva prakṛtam . \\
\noindent \textbf{(P\_7,1.84)} ugidacām sarvanāmasthāne adhātoḥ iti . \\
\noindent \textbf{(P\_7,1.84)} adhātvadhikārāt siddham iti cet napuṃsake doṣaḥ . \\
\noindent \textbf{(P\_7,1.84)} adhātvadhikārāt siddham iti cet napuṃsake doṣaḥ bhavati . \\
\noindent \textbf{(P\_7,1.84)} kāṣṭhataṅkṣi , kūṭataṅkṣi . \\
\noindent \textbf{(P\_7,1.84)} napuṃsakasya jhal acaḥ adhātoḥ iti pratiṣedhaḥ prāpnoti . \\
\noindent \textbf{(P\_7,1.84)} uktam vā . \\
\noindent \textbf{(P\_7,1.84)} kim uktam . \\
\noindent \textbf{(P\_7,1.84)} ananubandhakagrahaṇe hi na sānubandhakasya iti . \\
\noindent \textbf{(P\_7,1.84)} atha vā sambandham anuvartiṣyate \\
\noindent \textbf{(P\_7,1.86)} itaḥ advacanam anarthakam ākāraprakaraṇāt . \\
\noindent \textbf{(P\_7,1.86)} itaḥ advacanam anarthakam . \\
\noindent \textbf{(P\_7,1.86)} kim kāraṇam . \\
\noindent \textbf{(P\_7,1.86)} ākāraprakaraṇāt . \\
\noindent \textbf{(P\_7,1.86)} āt iti vartate . \\
\noindent \textbf{(P\_7,1.86)} ṣapūrvārtham tu . \\
\noindent \textbf{(P\_7,1.86)} ṣapūrvārtham tarhi at vaktavyaḥ . \\
\noindent \textbf{(P\_7,1.86)} ṛbhukṣāṇam indram , ṛbhukṣaṇam indram \\
\noindent \textbf{(P\_7,1.89)} asuṅi upadeśivadvacanam svarasiddhyartham bahiraṅgalakṣaṇatvāt . \\
\noindent \textbf{(P\_7,1.89)} asuṅi upadeśivadbhāvaḥ vaktavyaḥ . \\
\noindent \textbf{(P\_7,1.89)} upadeśāvasthāyām eva asuṅ bhavati iti vaktavyam . \\
\noindent \textbf{(P\_7,1.89)} kim prayojanam . \\
\noindent \textbf{(P\_7,1.89)} svarasiddhyartham . \\
\noindent \textbf{(P\_7,1.89)} upadeśāvasthāyām asuṅi kṛte iṣṭaḥ svaraḥ yathā syāt . \\
\noindent \textbf{(P\_7,1.89)} paramapumān iti . \\
\noindent \textbf{(P\_7,1.89)} akriyamāṇe hi upadeśivadbhāve samāsāntodāttatve asuṅ āntaryataḥ asvarakasya asvarakaḥ syāt . \\
\noindent \textbf{(P\_7,1.89)} kim punaḥ kāraṇam samāsāntodāttatvam tāvat bhavati na punaḥ asuṅ . \\
\noindent \textbf{(P\_7,1.89)} na paratvāt asuṅā bhavitavyam . \\
\noindent \textbf{(P\_7,1.89)} bahirṅgalakṣaṇatvāt . \\
\noindent \textbf{(P\_7,1.89)} bahiraṅgalakṣaṇaḥ asuṅ . \\
\noindent \textbf{(P\_7,1.89)} antaraṅgaḥ svaraḥ . \\
\noindent \textbf{(P\_7,1.89)} asiddham bahiraṅgam antaraṅge . \\
\noindent \textbf{(P\_7,1.89)} saḥ tarhi upadeśivadbhāvaḥ vaktavyaḥ . \\
\noindent \textbf{(P\_7,1.89)} na vaktavyaḥ . \\
\noindent \textbf{(P\_7,1.89)} ādyudāttanipātanam kariṣyate saḥ nipātanasvaraḥ samāsasvarasya bādhakaḥ bhaviṣyati . \\
\noindent \textbf{(P\_7,1.89)} evam api upadeśivadbhāvaḥ vaktavyaḥ . \\
\noindent \textbf{(P\_7,1.89)} saḥ yathā eva hi nipātanasvaraḥ samāsasvaram bādhate evam prakṛtisvaram api bādheta . \\
\noindent \textbf{(P\_7,1.89)} pumān . \\
\noindent \textbf{(P\_7,1.89)} tasmāt suṣṭhu ucyate asuṅi upadeśivadvacanam svarasiddhyartham bahiraṅgalakṣaṇatvāt iti \\
\noindent \textbf{(P\_7,1.90)} kim idam gotaḥ parasya sarvanāmasthānasya ṅittvam ucyate āhosvit sarvanāmasthāne parataḥ ṅitkāryam atidiśyate . \\
\noindent \textbf{(P\_7,1.90)} kaḥ ca atra viśeṣaḥ . \\
\noindent \textbf{(P\_7,1.90)} gotaḥ sarvanāmasthāne ṅitkāryātideśaḥ . \\
\noindent \textbf{(P\_7,1.90)} gotaḥ sarvanāmasthāne ṅitkāryam atidiśyate . \\
\noindent \textbf{(P\_7,1.90)} sarvanāmasthāne ṇittvavacane hi asampratyayaḥ ṣaṣṭhyanirdeśāt . \\
\noindent \textbf{(P\_7,1.90)} sarvanāmasthānasya ṇidvacane hi asampratyayaḥ syāt . \\
\noindent \textbf{(P\_7,1.90)} kim kāraṇam . \\
\noindent \textbf{(P\_7,1.90)} ṣaṣṭhyabhāvāt . \\
\noindent \textbf{(P\_7,1.90)} ṣaṣṭhīnirdiṣṭasya ādeśāḥ ucyante na ca atra ṣaṣṭhīm paśyāmaḥ . \\
\noindent \textbf{(P\_7,1.90)} evam tarhi vatinirdeśaḥ ayam : gotaḥ ṅidvat bhavati iti . \\
\noindent \textbf{(P\_7,1.90)} saḥ tarhi vatinirdeśaḥ kartavyaḥ na hi antareṇa vatim atideśaḥ gamyate . \\
\noindent \textbf{(P\_7,1.90)} antareṇa api vatim atideśaḥ gamyate . \\
\noindent \textbf{(P\_7,1.90)} tat yathā : eṣaḥ brahmadattaḥ . \\
\noindent \textbf{(P\_7,1.90)} abrahmadattam brahmadattaḥ iti āha . \\
\noindent \textbf{(P\_7,1.90)} te manyāmahe : brahmadattavat ayam bhavati iti . \\
\noindent \textbf{(P\_7,1.90)} evam iha api aṇitam ṇit iti āha ṇidvat iti gamyate . \\
\noindent \textbf{(P\_7,1.90)} atha vā punaḥ astu gotaḥ parasya sarvanāmasthānasya ṇittvam . \\
\noindent \textbf{(P\_7,1.90)} nanu ca uktam sarvanāmasthāne ṇittvavacane hi asampratyayaḥ ṣaṣṭhyanirdeśāt iti . \\
\noindent \textbf{(P\_7,1.90)} na eṣaḥ doṣaḥ . \\
\noindent \textbf{(P\_7,1.90)} gotaḥ iti eṣā pañcamī sarvanāmasthāne iti saptamyāḥ ṣaṣṭhīm prakalpayiṣyati tasmāt iti uttarasya iti . \\
\noindent \textbf{(P\_7,1.90)} atha taparakaraṇam kimartham . \\
\noindent \textbf{(P\_7,1.90)} iha mā bhūt . \\
\noindent \textbf{(P\_7,1.90)} citraguḥ śabalaguḥ iti . \\
\noindent \textbf{(P\_7,1.90)} na etat asti . \\
\noindent \textbf{(P\_7,1.90)} hrasvatve kṛte na bhaviṣyati . \\
\noindent \textbf{(P\_7,1.90)} sthānivadbhāvāt prāpnoti . \\
\noindent \textbf{(P\_7,1.90)} ataḥ uttaram paṭhati . \\
\noindent \textbf{(P\_7,1.90)} taparakaraṇam anarthakam sthānivatpratiṣedhāt . \\
\noindent \textbf{(P\_7,1.90)} taparakarṇam anarthakam . \\
\noindent \textbf{(P\_7,1.90)} kim kāraṇam . \\
\noindent \textbf{(P\_7,1.90)} sthānivatpratiṣedhāt . \\
\noindent \textbf{(P\_7,1.90)} pratiṣidhyate atra sthānivadbhāvaḥ goḥ pūrvaṇittvātvasvareṣu sthānivadbhāvaḥ na bhavati iti . \\
\noindent \textbf{(P\_7,1.90)} saḥ ca avaśyam pratiṣedhaḥ āśrayitavyaḥ . \\
\noindent \textbf{(P\_7,1.90)} itarathā hi sambuddhijasoḥ pratiṣedhaḥ . \\
\noindent \textbf{(P\_7,1.90)} yaḥ hi manyate taparakaraṇasāmarthyāt atra na bhaviṣyati iti sambuddhijasoḥ tena pratiṣedhaḥ vaktavyaḥ syāt : he citrago citragavaḥ iti . \\
\noindent \textbf{(P\_7,1.90)} atha idānīm sati api sthānivadbhāvapratiṣedha guṇe kṛte kasmāt eva atra na bhavati . \\
\noindent \textbf{(P\_7,1.90)} lakṣaṇapratipadoktayoḥ pratipadoktasya eva iti . \\
\noindent \textbf{(P\_7,1.90)} nanu ca idānīm asati api sthānivadbhāvapratiṣedhe etayā paribhāṣayā śakyam upasthātum . \\
\noindent \textbf{(P\_7,1.90)} na iti āha na hi idānīm kva cit api sthānivat syāt \\
\noindent \textbf{(P\_7,1.95-96.1)} KA\_III,272.25-273.19 Ro\_V,81-84 {1/32}     atha atra vibhaktau iti anuvartate utāho na . \\
\noindent \textbf{(P\_7,1.95-96.1)} KA\_III,272.25-273.19 Ro\_V,81-84 {2/32}     kim ca ataḥ . \\
\noindent \textbf{(P\_7,1.95-96.1)} KA\_III,272.25-273.19 Ro\_V,81-84 {3/32}     tṛjvat striyām vibhaktau cet kroṣṭrībhaktiḥ na sidhyati . \\
\noindent \textbf{(P\_7,1.95-96.1)} KA\_III,272.25-273.19 Ro\_V,81-84 {4/32}     tṛjvat striyām vibhaktau cet kroṣṭrībhaktiḥ iti na sidhyati . \\
\noindent \textbf{(P\_7,1.95-96.1)} KA\_III,272.25-273.19 Ro\_V,81-84 {5/32}     evam tarhi īkāre tṛjvadbhāvam vakṣyāmi . \\
\noindent \textbf{(P\_7,1.95-96.1)} KA\_III,272.25-273.19 Ro\_V,81-84 {6/32}     tat īkāragrahaṇam kartavyam . \\
\noindent \textbf{(P\_7,1.95-96.1)} KA\_III,272.25-273.19 Ro\_V,81-84 {7/32}     na kartavyam . \\
\noindent \textbf{(P\_7,1.95-96.1)} KA\_III,272.25-273.19 Ro\_V,81-84 {8/32}     kriyate nyāse eva . \\
\noindent \textbf{(P\_7,1.95-96.1)} KA\_III,272.25-273.19 Ro\_V,81-84 {9/32}     praśliṣṭanirdeśaḥ ayam . \\
\noindent \textbf{(P\_7,1.95-96.1)} KA\_III,272.25-273.19 Ro\_V,81-84 {10/32}     strī , ī strī striyām iti . \\
\noindent \textbf{(P\_7,1.95-96.1)} KA\_III,272.25-273.19 Ro\_V,81-84 {11/32}     īkāre tannimittaḥ saḥ . \\
\noindent \textbf{(P\_7,1.95-96.1)} KA\_III,272.25-273.19 Ro\_V,81-84 {12/32}     īkāre cet tat na . \\
\noindent \textbf{(P\_7,1.95-96.1)} KA\_III,272.25-273.19 Ro\_V,81-84 {13/32}     kim kāraṇam . \\
\noindent \textbf{(P\_7,1.95-96.1)} KA\_III,272.25-273.19 Ro\_V,81-84 {14/32}     tannimittaḥ saḥ . \\
\noindent \textbf{(P\_7,1.95-96.1)} KA\_III,272.25-273.19 Ro\_V,81-84 {15/32}     tṛjvadbhāvanimittaḥ saḥ īkāraḥ . \\
\noindent \textbf{(P\_7,1.95-96.1)} KA\_III,272.25-273.19 Ro\_V,81-84 {16/32}     na akṛte tṛjvadbhāve īkāraḥ prāpnoti . \\
\noindent \textbf{(P\_7,1.95-96.1)} KA\_III,272.25-273.19 Ro\_V,81-84 {17/32}     kim kāraṇam . \\
\noindent \textbf{(P\_7,1.95-96.1)} KA\_III,272.25-273.19 Ro\_V,81-84 {18/32}     ṛnnebhyaḥ ṅīp iti ucyate īkāre ca tṛjvadbhāvaḥ . \\
\noindent \textbf{(P\_7,1.95-96.1)} KA\_III,272.25-273.19 Ro\_V,81-84 {19/32}     tat idam itaretarāśrayam bhavati . \\
\noindent \textbf{(P\_7,1.95-96.1)} KA\_III,272.25-273.19 Ro\_V,81-84 {20/32}     itaretarāśrayāṇi ca na prakalpante . \\
\noindent \textbf{(P\_7,1.95-96.1)} KA\_III,272.25-273.19 Ro\_V,81-84 {21/32}     evam tarhi gaurādiṣu pāṭhāt īkāraḥ bhaviṣyati . \\
\noindent \textbf{(P\_7,1.95-96.1)} KA\_III,272.25-273.19 Ro\_V,81-84 {22/32}     gaurādiṣu na paṭhyate . \\
\noindent \textbf{(P\_7,1.95-96.1)} KA\_III,272.25-273.19 Ro\_V,81-84 {23/32}     na hi kim cit tunantam gaurādiṣu paṭhyate . \\
\noindent \textbf{(P\_7,1.95-96.1)} KA\_III,272.25-273.19 Ro\_V,81-84 {24/32}     evam tarhi etat jñāpayati ācāryaḥ bhavati atra īkāraḥ iti yat ayam īkāre tṛjvadbhāvam śāsti . \\
\noindent \textbf{(P\_7,1.95-96.1)} KA\_III,272.25-273.19 Ro\_V,81-84 {25/32}     tena eva bhāvanam cet syāt aniṣṭaḥ api prasajyate . \\
\noindent \textbf{(P\_7,1.95-96.1)} KA\_III,272.25-273.19 Ro\_V,81-84 {26/32}     yadi api na asti viśeṣaḥ ṅīpaḥ vā ṅīṣaḥ vā ṅīn api tu prāpnoti . \\
\noindent \textbf{(P\_7,1.95-96.1)} KA\_III,272.25-273.19 Ro\_V,81-84 {27/32}     iha ca na prāpnoti . \\
\noindent \textbf{(P\_7,1.95-96.1)} KA\_III,272.25-273.19 Ro\_V,81-84 {28/32}     pañcabhiḥ kroṣṭrībhiḥ krītaiḥ rathaiḥ pañcakroṣṭṛbhiḥ rathaiḥ iti . \\
\noindent \textbf{(P\_7,1.95-96.1)} KA\_III,272.25-273.19 Ro\_V,81-84 {29/32}     evam tarhi na ca aparam nimittam sañjñā ca pratyayalakṣaṇena . \\
\noindent \textbf{(P\_7,1.95-96.1)} KA\_III,272.25-273.19 Ro\_V,81-84 {30/32}     na ca aparam nimittam āśrīyate : asmin parataḥ kroṣṭuḥ tṛjvat bhavati iti . \\
\noindent \textbf{(P\_7,1.95-96.1)} KA\_III,272.25-273.19 Ro\_V,81-84 {31/32}     kim tarhi aṅgasya kroṣṭuḥ tṛjvat bhavati . \\
\noindent \textbf{(P\_7,1.95-96.1)} KA\_III,272.25-273.19 Ro\_V,81-84 {32/32}     aṅgasañjñā ca bhavati pratyayalakṣaṇena \\
\noindent \textbf{(P\_7,1.95-96.2)} KA\_III,273.20-275.22 Ro\_V,84-91 {1/82}     kim punaḥ ayam śāstrātideśaḥ : tṛcaḥ yat śāstram tat atidiśyate . \\
\noindent \textbf{(P\_7,1.95-96.2)} KA\_III,273.20-275.22 Ro\_V,84-91 {2/82}     āhosvit rūpātideśaḥ : tṛcaḥ yat rūpam tat atidiśyate iti . \\
\noindent \textbf{(P\_7,1.95-96.2)} KA\_III,273.20-275.22 Ro\_V,84-91 {3/82}     kaḥ ca atra viśeṣaḥ . \\
\noindent \textbf{(P\_7,1.95-96.2)} KA\_III,273.20-275.22 Ro\_V,84-91 {4/82}     tṛjvat iti śāstrātideśaḥ cet yathā ciṇi tadvat . \\
\noindent \textbf{(P\_7,1.95-96.2)} KA\_III,273.20-275.22 Ro\_V,84-91 {5/82}     tṛjvat iti śāstrātideśaḥ cet yathā ciṇi tadvat prāpnoti . \\
\noindent \textbf{(P\_7,1.95-96.2)} KA\_III,273.20-275.22 Ro\_V,84-91 {6/82}     katham ca ciṇi . \\
\noindent \textbf{(P\_7,1.95-96.2)} KA\_III,273.20-275.22 Ro\_V,84-91 {7/82}     uktam aṅgasya iti tu prakaraṇāt āṇgaśāstrātideśāt siddham iti . \\
\noindent \textbf{(P\_7,1.95-96.2)} KA\_III,273.20-275.22 Ro\_V,84-91 {8/82}     āṅgam yat kāryam tat atidiśyate . \\
\noindent \textbf{(P\_7,1.95-96.2)} KA\_III,273.20-275.22 Ro\_V,84-91 {9/82}     evam iha api anaṅguṇadīrghatvāni atidiṣṭāni raparatvam anatidiṣṭam . \\
\noindent \textbf{(P\_7,1.95-96.2)} KA\_III,273.20-275.22 Ro\_V,84-91 {10/82}     tatra kaḥ doṣaḥ . \\
\noindent \textbf{(P\_7,1.95-96.2)} KA\_III,273.20-275.22 Ro\_V,84-91 {11/82}      tatra raparavacanam . \\
\noindent \textbf{(P\_7,1.95-96.2)} KA\_III,273.20-275.22 Ro\_V,84-91 {12/82}     tatra raparatvam na sidhyati tat vaktavyam . \\
\noindent \textbf{(P\_7,1.95-96.2)} KA\_III,273.20-275.22 Ro\_V,84-91 {13/82}     na eṣaḥ doṣaḥ . \\
\noindent \textbf{(P\_7,1.95-96.2)} KA\_III,273.20-275.22 Ro\_V,84-91 {14/82}     guṇe atidiṣṭe raparatvam api atidiṣṭam bhavati . \\
\noindent \textbf{(P\_7,1.95-96.2)} KA\_III,273.20-275.22 Ro\_V,84-91 {15/82}     katham . \\
\noindent \textbf{(P\_7,1.95-96.2)} KA\_III,273.20-275.22 Ro\_V,84-91 {16/82}     kāryakālam sañjñāparibhāṣam yatra kāryam tatra draṣṭavyam . \\
\noindent \textbf{(P\_7,1.95-96.2)} KA\_III,273.20-275.22 Ro\_V,84-91 {17/82}     ṛtaḥ ṅisarvanāmasthānayoḥ guṅaḥ bhavati . \\
\noindent \textbf{(P\_7,1.95-96.2)} KA\_III,273.20-275.22 Ro\_V,84-91 {18/82}     upasthitam idam bhavati uḥ aṇ raparaḥ iti . \\
\noindent \textbf{(P\_7,1.95-96.2)} KA\_III,273.20-275.22 Ro\_V,84-91 {19/82}     evam tarhi ayam anyaḥ doṣaḥ jāyate . \\
\noindent \textbf{(P\_7,1.95-96.2)} KA\_III,273.20-275.22 Ro\_V,84-91 {20/82}     āhatya tṛcaḥ yat śāstram tat atidiśyeta anāhatya vā iti . \\
\noindent \textbf{(P\_7,1.95-96.2)} KA\_III,273.20-275.22 Ro\_V,84-91 {21/82}     kim ca ataḥ . \\
\noindent \textbf{(P\_7,1.95-96.2)} KA\_III,273.20-275.22 Ro\_V,84-91 {22/82}     yadi āhatya dīrghatvam atidiṣṭam anaṅguṇaraparatvāni anatidiṣṭāni . \\
\noindent \textbf{(P\_7,1.95-96.2)} KA\_III,273.20-275.22 Ro\_V,84-91 {23/82}     atha anāhatya anaṅguṇaraparatvāni atidiṣṭāni dīrghatvam anatidiṣṭam . \\
\noindent \textbf{(P\_7,1.95-96.2)} KA\_III,273.20-275.22 Ro\_V,84-91 {24/82}     astu āhatya . \\
\noindent \textbf{(P\_7,1.95-96.2)} KA\_III,273.20-275.22 Ro\_V,84-91 {25/82}     nanu ca uktam dīrghatvam atidiṣṭam anaṅguṇaraparatvāni anatidiṣṭāni iti . \\
\noindent \textbf{(P\_7,1.95-96.2)} KA\_III,273.20-275.22 Ro\_V,84-91 {26/82}     na eṣaḥ doṣaḥ . \\
\noindent \textbf{(P\_7,1.95-96.2)} KA\_III,273.20-275.22 Ro\_V,84-91 {27/82}     dīrghatve atidiṣṭe anaṅguṇaraparatvāni api atidiṣṭāni bhavanti . \\
\noindent \textbf{(P\_7,1.95-96.2)} KA\_III,273.20-275.22 Ro\_V,84-91 {28/82}     katham . \\
\noindent \textbf{(P\_7,1.95-96.2)} KA\_III,273.20-275.22 Ro\_V,84-91 {29/82}     upadhāyāḥ iti vartate na ca akṛteṣu eteṣu dīrghabhāvini upadhā bhavati . \\
\noindent \textbf{(P\_7,1.95-96.2)} KA\_III,273.20-275.22 Ro\_V,84-91 {30/82}     kutaḥ nu khalu etat eteṣu vidhiṣu kṛteṣu yā upadhā tasyāḥ dīrghatvam bhaviṣyati na punaḥ kroṣṭoḥ yaḥ antaratamaḥ guṇaḥ tasmin kṛte avādeśe ca yā upadhā tasyāḥ dīrghatvam bhaviṣyati . \\
\noindent \textbf{(P\_7,1.95-96.2)} KA\_III,273.20-275.22 Ro\_V,84-91 {31/82}     na ekam udāharaṇam yogārambham prayojayati iti . \\
\noindent \textbf{(P\_7,1.95-96.2)} KA\_III,273.20-275.22 Ro\_V,84-91 {32/82}     tatra tṛjvadvacanasāmarthyāt eteṣu vidhiṣu kṛteṣu yā upadhā tasyāḥ dīrghatvam bhaviṣyati . \\
\noindent \textbf{(P\_7,1.95-96.2)} KA\_III,273.20-275.22 Ro\_V,84-91 {33/82}     atha vā kim naḥ etena āhatya anāhatya vā iti . \\
\noindent \textbf{(P\_7,1.95-96.2)} KA\_III,273.20-275.22 Ro\_V,84-91 {34/82}     āhatya anāhatya ca tṛcaḥ yat śāstram tat atidiśyate . \\
\noindent \textbf{(P\_7,1.95-96.2)} KA\_III,273.20-275.22 Ro\_V,84-91 {35/82}     atha vā punaḥ astu rūpātideśaḥ . \\
\noindent \textbf{(P\_7,1.95-96.2)} KA\_III,273.20-275.22 Ro\_V,84-91 {36/82}     atha etasmin rūpātideśe sati kim prāk ādeśebhyaḥ yat rūpam tat atidiśyate āhosvit kṛteṣu ādeśeṣu . \\
\noindent \textbf{(P\_7,1.95-96.2)} KA\_III,273.20-275.22 Ro\_V,84-91 {37/82}     kim ca ataḥ . \\
\noindent \textbf{(P\_7,1.95-96.2)} KA\_III,273.20-275.22 Ro\_V,84-91 {38/82}     yadi prāk ādeśebhyaḥ yat rūpam tat atidiśyate ṛkāraḥ ekaḥ atidiṣṭaḥ anaṅguṇaraparatvadīrghatvāni anatidiṣṭāni . \\
\noindent \textbf{(P\_7,1.95-96.2)} KA\_III,273.20-275.22 Ro\_V,84-91 {39/82}     atha kṛteṣu ādeśeṣu ṛkāraḥ anatidiṣṭaḥ anaṅguṇaraparatvadīrghatvāni atidiṣṭāni . \\
\noindent \textbf{(P\_7,1.95-96.2)} KA\_III,273.20-275.22 Ro\_V,84-91 {40/82}     ubhayathā ca svaraḥ anatidiṣṭaḥ na hi svaraḥ rūpavān . \\
\noindent \textbf{(P\_7,1.95-96.2)} KA\_III,273.20-275.22 Ro\_V,84-91 {41/82}     astu prāk ādeśebhyaḥ yat rūpam tat atidiśyate . \\
\noindent \textbf{(P\_7,1.95-96.2)} KA\_III,273.20-275.22 Ro\_V,84-91 {42/82}     nanu ca uktam ṛkāraḥ atidiṣṭaḥ anaṅguṇaraparatvadīrghatvāni anatidiṣṭāni iti . \\
\noindent \textbf{(P\_7,1.95-96.2)} KA\_III,273.20-275.22 Ro\_V,84-91 {43/82}     na eṣaḥ doṣaḥ . \\
\noindent \textbf{(P\_7,1.95-96.2)} KA\_III,273.20-275.22 Ro\_V,84-91 {44/82}     ṛkāre atidiṣṭe svāśrayāḥ atra ete vidhayaḥ bhaviṣyanti . \\
\noindent \textbf{(P\_7,1.95-96.2)} KA\_III,273.20-275.22 Ro\_V,84-91 {45/82}     yat api ucyate ubhayathā ca svaraḥ anatidiṣṭaḥ na hi svaraḥ rūpavān iti sacakāragrahaṇasāmarthyāt svaraḥ bhaviṣyati . \\
\noindent \textbf{(P\_7,1.95-96.2)} KA\_III,273.20-275.22 Ro\_V,84-91 {46/82}     rūpātideśaḥ iti cet sarvādeśaprasaṅgaḥ . \\
\noindent \textbf{(P\_7,1.95-96.2)} KA\_III,273.20-275.22 Ro\_V,84-91 {47/82}     rūpātideśaḥ iti cet sarvādeśaḥ prāpnoti . \\
\noindent \textbf{(P\_7,1.95-96.2)} KA\_III,273.20-275.22 Ro\_V,84-91 {48/82}     sarvasya tunantasya tṛśabdaḥ ādeśaḥ prāpnoti . \\
\noindent \textbf{(P\_7,1.95-96.2)} KA\_III,273.20-275.22 Ro\_V,84-91 {49/82}     siddham tu rūpātideśāt . \\
\noindent \textbf{(P\_7,1.95-96.2)} KA\_III,273.20-275.22 Ro\_V,84-91 {50/82}     siddham etat . \\
\noindent \textbf{(P\_7,1.95-96.2)} KA\_III,273.20-275.22 Ro\_V,84-91 {51/82}     katham . \\
\noindent \textbf{(P\_7,1.95-96.2)} KA\_III,273.20-275.22 Ro\_V,84-91 {52/82}     rūpātideśāt . \\
\noindent \textbf{(P\_7,1.95-96.2)} KA\_III,273.20-275.22 Ro\_V,84-91 {53/82}     rūpātideśaḥ ayam . \\
\noindent \textbf{(P\_7,1.95-96.2)} KA\_III,273.20-275.22 Ro\_V,84-91 {54/82}     nanu ca evam eva kṛtvā codyate rūpātideśaḥ iti cet sarvādeśaprasaṅgaḥ iti . \\
\noindent \textbf{(P\_7,1.95-96.2)} KA\_III,273.20-275.22 Ro\_V,84-91 {55/82}     siddham tu pratyayagrahaṇe yasmāt saḥ tadāditadantavijñānāt . \\
\noindent \textbf{(P\_7,1.95-96.2)} KA\_III,273.20-275.22 Ro\_V,84-91 {56/82}     siddham etat . \\
\noindent \textbf{(P\_7,1.95-96.2)} KA\_III,273.20-275.22 Ro\_V,84-91 {57/82}     katham . \\
\noindent \textbf{(P\_7,1.95-96.2)} KA\_III,273.20-275.22 Ro\_V,84-91 {58/82}     pratyayagrahaṇe yasmāt saḥ vihitaḥ tadādeḥ tadantasya ca grahaṇam bhavati iti evam tunantasya tṛjantaḥ ādeśaḥ bhaviṣyati . \\
\noindent \textbf{(P\_7,1.95-96.2)} KA\_III,273.20-275.22 Ro\_V,84-91 {59/82}     evam api kim cit eva tṛjantam prāpnoti . \\
\noindent \textbf{(P\_7,1.95-96.2)} KA\_III,273.20-275.22 Ro\_V,84-91 {60/82}     idam api prāpnoti paktā iti . \\
\noindent \textbf{(P\_7,1.95-96.2)} KA\_III,273.20-275.22 Ro\_V,84-91 {61/82}     āntaratamyāt ca siddham . \\
\noindent \textbf{(P\_7,1.95-96.2)} KA\_III,273.20-275.22 Ro\_V,84-91 {62/82}     kroṣṭoḥ yat antaratamam tat bhaviṣyati . \\
\noindent \textbf{(P\_7,1.95-96.2)} KA\_III,273.20-275.22 Ro\_V,84-91 {63/82}     kim punaḥ tat . \\
\noindent \textbf{(P\_7,1.95-96.2)} KA\_III,273.20-275.22 Ro\_V,84-91 {64/82}     kruśeḥ yaḥ tṛc vihitaḥ tadantam . \\
\noindent \textbf{(P\_7,1.95-96.2)} KA\_III,273.20-275.22 Ro\_V,84-91 {65/82}     tṛjvadvacanam anarthakam tṛjviṣaye tṛcaḥ mṛgavācitvāt . \\
\noindent \textbf{(P\_7,1.95-96.2)} KA\_III,273.20-275.22 Ro\_V,84-91 {66/82}     tṛjviṣaye etat tṛjantam mṛgavāci . \\
\noindent \textbf{(P\_7,1.95-96.2)} KA\_III,273.20-275.22 Ro\_V,84-91 {67/82}     tunaḥ nivṛttyartham tarhi idam vaktavyam . \\
\noindent \textbf{(P\_7,1.95-96.2)} KA\_III,273.20-275.22 Ro\_V,84-91 {68/82}     tunaḥ sarvanāmasthāne nivṛttiḥ yathā syāt . \\
\noindent \textbf{(P\_7,1.95-96.2)} KA\_III,273.20-275.22 Ro\_V,84-91 {69/82}     tunaḥ nivṛttyartham iti cet siddham yathā anyatra api . \\
\noindent \textbf{(P\_7,1.95-96.2)} KA\_III,273.20-275.22 Ro\_V,84-91 {70/82}     tunaḥ nivṛttyartham iti cet tat antareṇa vacanam siddham yathā anyatra api aviśeṣavihitāḥ śabdāḥ niyataviṣayāḥ dṛśyante . \\
\noindent \textbf{(P\_7,1.95-96.2)} KA\_III,273.20-275.22 Ro\_V,84-91 {71/82}     kva anyatra . \\
\noindent \textbf{(P\_7,1.95-96.2)} KA\_III,273.20-275.22 Ro\_V,84-91 {72/82}     tat yathā . \\
\noindent \textbf{(P\_7,1.95-96.2)} KA\_III,273.20-275.22 Ro\_V,84-91 {73/82}     gharatiḥ asmai aviśeṣeṇa upadiṣṭaḥ saḥ ghṛtam , ghṛṇā , gharmaḥ iti evaṃviṣayaḥ . \\
\noindent \textbf{(P\_7,1.95-96.2)} KA\_III,273.20-275.22 Ro\_V,84-91 {74/82}     raśiḥ asmai aviśeṣeṇa upadiṣṭaḥ saḥ rāśiḥ , raśmiḥ , raśanā iti evaṃviṣayaḥ . \\
\noindent \textbf{(P\_7,1.95-96.2)} KA\_III,273.20-275.22 Ro\_V,84-91 {75/82}     luśiḥ asmai aviśeṣeṇa upadiṣṭaḥ saḥ loṣṭaḥ iti evaṃviṣayaḥ . \\
\noindent \textbf{(P\_7,1.95-96.2)} KA\_III,273.20-275.22 Ro\_V,84-91 {76/82}     idam tarhi prayojanam vibhāṣā vakṣyāmi iti . \\
\noindent \textbf{(P\_7,1.95-96.2)} KA\_III,273.20-275.22 Ro\_V,84-91 {77/82}     vibhāṣā tṛtīyādiṣu aci iti . \\
\noindent \textbf{(P\_7,1.95-96.2)} KA\_III,273.20-275.22 Ro\_V,84-91 {78/82}     vāvacanānarthakyam ca svabhāvasiddhatvāt . \\
\noindent \textbf{(P\_7,1.95-96.2)} KA\_III,273.20-275.22 Ro\_V,84-91 {79/82}     vāvacanam ca anarthakam . \\
\noindent \textbf{(P\_7,1.95-96.2)} KA\_III,273.20-275.22 Ro\_V,84-91 {80/82}     kim kāraṇam . \\
\noindent \textbf{(P\_7,1.95-96.2)} KA\_III,273.20-275.22 Ro\_V,84-91 {81/82}     svabhāvasiddhatvāt . \\
\noindent \textbf{(P\_7,1.95-96.2)} KA\_III,273.20-275.22 Ro\_V,84-91 {82/82}     svabhāvataḥ eva tṛtīyādiṣu ajādiṣu vibhaktiṣu tṛjantam ca tunantam ca mṛgavāci iti \\
\noindent \textbf{(P\_7,1.95-96.3)} KA\_III,275.23-276.22 Ro\_V,91-92 {1/62}     guṇavṛddhyauttvatṛjvadbhāvebhyaḥ num pūrvavipratiṣiddham . \\
\noindent \textbf{(P\_7,1.95-96.3)} KA\_III,275.23-276.22 Ro\_V,91-92 {2/62}     guṇavṛddhyauttvatṛjvadbhāvebhyaḥ num bhavati pūrvavipratiṣedhena . \\
\noindent \textbf{(P\_7,1.95-96.3)} KA\_III,275.23-276.22 Ro\_V,91-92 {3/62}     tatra guṇasya avakāśaḥ . \\
\noindent \textbf{(P\_7,1.95-96.3)} KA\_III,275.23-276.22 Ro\_V,91-92 {4/62}     agnaye , vāyave . \\
\noindent \textbf{(P\_7,1.95-96.3)} KA\_III,275.23-276.22 Ro\_V,91-92 {5/62}     numaḥ avakāśaḥ . \\
\noindent \textbf{(P\_7,1.95-96.3)} KA\_III,275.23-276.22 Ro\_V,91-92 {6/62}     trapuṇī , jatunī . \\
\noindent \textbf{(P\_7,1.95-96.3)} KA\_III,275.23-276.22 Ro\_V,91-92 {7/62}     iha ubhayam prāpnoti . \\
\noindent \textbf{(P\_7,1.95-96.3)} KA\_III,275.23-276.22 Ro\_V,91-92 {8/62}     trapuṇe , jatune . \\
\noindent \textbf{(P\_7,1.95-96.3)} KA\_III,275.23-276.22 Ro\_V,91-92 {9/62}     vṛddheḥ avakāśaḥ . \\
\noindent \textbf{(P\_7,1.95-96.3)} KA\_III,275.23-276.22 Ro\_V,91-92 {10/62}     sakhāyai . \\
\noindent \textbf{(P\_7,1.95-96.3)} KA\_III,275.23-276.22 Ro\_V,91-92 {11/62}     sakhāyaḥ . \\
\noindent \textbf{(P\_7,1.95-96.3)} KA\_III,275.23-276.22 Ro\_V,91-92 {12/62}     numaḥ saḥ eva . \\
\noindent \textbf{(P\_7,1.95-96.3)} KA\_III,275.23-276.22 Ro\_V,91-92 {13/62}     iha ubhayam prāpnoti . \\
\noindent \textbf{(P\_7,1.95-96.3)} KA\_III,275.23-276.22 Ro\_V,91-92 {14/62}     atisakhīni brāhmaṇakulāni iti . \\
\noindent \textbf{(P\_7,1.95-96.3)} KA\_III,275.23-276.22 Ro\_V,91-92 {15/62}     auttvasya avakāśaḥ . \\
\noindent \textbf{(P\_7,1.95-96.3)} KA\_III,275.23-276.22 Ro\_V,91-92 {16/62}     agnau , vāyau . \\
\noindent \textbf{(P\_7,1.95-96.3)} KA\_III,275.23-276.22 Ro\_V,91-92 {17/62}     numaḥ saḥ eva . \\
\noindent \textbf{(P\_7,1.95-96.3)} KA\_III,275.23-276.22 Ro\_V,91-92 {18/62}     iha ubhayam prāpnoti . \\
\noindent \textbf{(P\_7,1.95-96.3)} KA\_III,275.23-276.22 Ro\_V,91-92 {19/62}     trapuṇi , jatuni iti . \\
\noindent \textbf{(P\_7,1.95-96.3)} KA\_III,275.23-276.22 Ro\_V,91-92 {20/62}     tṛjvadbhāvasya avakāśaḥ . \\
\noindent \textbf{(P\_7,1.95-96.3)} KA\_III,275.23-276.22 Ro\_V,91-92 {21/62}     kroṣṭunā . \\
\noindent \textbf{(P\_7,1.95-96.3)} KA\_III,275.23-276.22 Ro\_V,91-92 {22/62}     numaḥ saḥ eva . \\
\noindent \textbf{(P\_7,1.95-96.3)} KA\_III,275.23-276.22 Ro\_V,91-92 {23/62}     iha ubhayam prāpnoti . \\
\noindent \textbf{(P\_7,1.95-96.3)} KA\_III,275.23-276.22 Ro\_V,91-92 {24/62}     kṛśakṛoṣṭune arṇyāya . \\
\noindent \textbf{(P\_7,1.95-96.3)} KA\_III,275.23-276.22 Ro\_V,91-92 {25/62}     hitakroṣṭune vṛṣalakulāya . \\
\noindent \textbf{(P\_7,1.95-96.3)} KA\_III,275.23-276.22 Ro\_V,91-92 {26/62}     num bhavati pūrvavipratiṣedhena . \\
\noindent \textbf{(P\_7,1.95-96.3)} KA\_III,275.23-276.22 Ro\_V,91-92 {27/62}     saḥ tarhi pūrvavipratiṣedhaḥ vaktavyaḥ . \\
\noindent \textbf{(P\_7,1.95-96.3)} KA\_III,275.23-276.22 Ro\_V,91-92 {28/62}     na vaktavyaḥ . \\
\noindent \textbf{(P\_7,1.95-96.3)} KA\_III,275.23-276.22 Ro\_V,91-92 {29/62}     iṣṭavācī paraśabdaḥ . \\
\noindent \textbf{(P\_7,1.95-96.3)} KA\_III,275.23-276.22 Ro\_V,91-92 {30/62}     vipratiṣedhe param yat iṣṭam tat bhavati iti . \\
\noindent \textbf{(P\_7,1.95-96.3)} KA\_III,275.23-276.22 Ro\_V,91-92 {31/62}     numaciratṛjvadbhāvebhyaḥ nuṭ . \\
\noindent \textbf{(P\_7,1.95-96.3)} KA\_III,275.23-276.22 Ro\_V,91-92 {32/62}     numaciratṛjvadbhāvebhyaḥ nuṭ pūrvavipratiṣedhena vaktavyaḥ . \\
\noindent \textbf{(P\_7,1.95-96.3)} KA\_III,275.23-276.22 Ro\_V,91-92 {33/62}     numaḥ avakāśaḥ . \\
\noindent \textbf{(P\_7,1.95-96.3)} KA\_III,275.23-276.22 Ro\_V,91-92 {34/62}     trapūṇi , jatūni . \\
\noindent \textbf{(P\_7,1.95-96.3)} KA\_III,275.23-276.22 Ro\_V,91-92 {35/62}     nuṭaḥ avakāśaḥ . \\
\noindent \textbf{(P\_7,1.95-96.3)} KA\_III,275.23-276.22 Ro\_V,91-92 {36/62}     agnīnām , vāyūnām . \\
\noindent \textbf{(P\_7,1.95-96.3)} KA\_III,275.23-276.22 Ro\_V,91-92 {37/62}     iha ubhayam prāpnoti . \\
\noindent \textbf{(P\_7,1.95-96.3)} KA\_III,275.23-276.22 Ro\_V,91-92 {38/62}     trapūṇām , jatūnām . \\
\noindent \textbf{(P\_7,1.95-96.3)} KA\_III,275.23-276.22 Ro\_V,91-92 {39/62}     aci rādeśasya avakāśaḥ . \\
\noindent \textbf{(P\_7,1.95-96.3)} KA\_III,275.23-276.22 Ro\_V,91-92 {40/62}     tisraḥ tiṣṭhanti catasraḥ tiṣṭhanti . \\
\noindent \textbf{(P\_7,1.95-96.3)} KA\_III,275.23-276.22 Ro\_V,91-92 {41/62}     nuṭaḥ saḥ eva . \\
\noindent \textbf{(P\_7,1.95-96.3)} KA\_III,275.23-276.22 Ro\_V,91-92 {42/62}     iha ubhayam prāpnoti . \\
\noindent \textbf{(P\_7,1.95-96.3)} KA\_III,275.23-276.22 Ro\_V,91-92 {43/62}     tisṛṇām , catasṛṇām . \\
\noindent \textbf{(P\_7,1.95-96.3)} KA\_III,275.23-276.22 Ro\_V,91-92 {44/62}     tṛjvadbhāvasya avakāśaḥ . \\
\noindent \textbf{(P\_7,1.95-96.3)} KA\_III,275.23-276.22 Ro\_V,91-92 {45/62}     kroṣṭrā , kroṣṭunā . \\
\noindent \textbf{(P\_7,1.95-96.3)} KA\_III,275.23-276.22 Ro\_V,91-92 {46/62}     nuṭaḥ saḥ eva . \\
\noindent \textbf{(P\_7,1.95-96.3)} KA\_III,275.23-276.22 Ro\_V,91-92 {47/62}     iha ubhayam prāpnoti . \\
\noindent \textbf{(P\_7,1.95-96.3)} KA\_III,275.23-276.22 Ro\_V,91-92 {48/62}     kroṣṭūnām . \\
\noindent \textbf{(P\_7,1.95-96.3)} KA\_III,275.23-276.22 Ro\_V,91-92 {49/62}     nuṭ bhavati pūrvavipratiṣedhena . \\
\noindent \textbf{(P\_7,1.95-96.3)} KA\_III,275.23-276.22 Ro\_V,91-92 {50/62}     saḥ tarhi pūrvavipratiṣedhaḥ vaktavyaḥ . \\
\noindent \textbf{(P\_7,1.95-96.3)} KA\_III,275.23-276.22 Ro\_V,91-92 {51/62}     na vā nuḍviṣaye rapratiṣedhāt . \\
\noindent \textbf{(P\_7,1.95-96.3)} KA\_III,275.23-276.22 Ro\_V,91-92 {52/62}     na vā etat vipratiṣedhena api sidhyati tisṛṇām , catasṛṇām iti . \\
\noindent \textbf{(P\_7,1.95-96.3)} KA\_III,275.23-276.22 Ro\_V,91-92 {53/62}     katham tarhi sidhyati . \\
\noindent \textbf{(P\_7,1.95-96.3)} KA\_III,275.23-276.22 Ro\_V,91-92 {54/62}     nuḍviṣaye rapratiṣedhāt . \\
\noindent \textbf{(P\_7,1.95-96.3)} KA\_III,275.23-276.22 Ro\_V,91-92 {55/62}     nuḍviṣaye rapratiṣedhaḥ vaktavyaḥ . \\
\noindent \textbf{(P\_7,1.95-96.3)} KA\_III,275.23-276.22 Ro\_V,91-92 {56/62}     itarathā hi sarvāpavādaḥ . \\
\noindent \textbf{(P\_7,1.95-96.3)} KA\_III,275.23-276.22 Ro\_V,91-92 {57/62}     itarathā hi sarvāpavādaḥ rādeśaḥ . \\
\noindent \textbf{(P\_7,1.95-96.3)} KA\_III,275.23-276.22 Ro\_V,91-92 {58/62}     saḥ yathā eva guṇapūrvasavarṇau bādhate eva nuṭam api bādheta . \\
\noindent \textbf{(P\_7,1.95-96.3)} KA\_III,275.23-276.22 Ro\_V,91-92 {59/62}     tasmāt nuḍviṣaye rapratiṣedhaḥ . \\
\noindent \textbf{(P\_7,1.95-96.3)} KA\_III,275.23-276.22 Ro\_V,91-92 {60/62}     tasmāt nuḍviṣaye rādeśasya pratiṣedhaḥ vaktavyaḥ . \\
\noindent \textbf{(P\_7,1.95-96.3)} KA\_III,275.23-276.22 Ro\_V,91-92 {61/62}     na vaktavyaḥ . \\
\noindent \textbf{(P\_7,1.95-96.3)} KA\_III,275.23-276.22 Ro\_V,91-92 {62/62}     ācāryapravṛttiḥ jñāpayati na rādeśaḥ nuṭam bādhate iti yat ayam na tisṛcatasṛ , iti pratiṣedham śāsti nāmi dīrghatvasya \\
\noindent \textbf{(P\_7,1.98)} ām anaḍuhaḥ striyām vā . \\
\noindent \textbf{(P\_7,1.98)} ām anaḍuhaḥ striyām vā iti vaktavyam . \\
\noindent \textbf{(P\_7,1.98)} anaḍuhī , anaḍvāhī \\
\noindent \textbf{(P\_7,1.100-102)} KA\_III,277.4-8 Ro\_V,93.4-8 {1/11}     ittvottvābhyām guṇavṛddhī bhavataḥ vipratiṣedhena . \\
\noindent \textbf{(P\_7,1.100-102)} KA\_III,277.4-8 Ro\_V,93.4-8 {2/11}     ittvottvayoḥ avakāśaḥ . \\
\noindent \textbf{(P\_7,1.100-102)} KA\_III,277.4-8 Ro\_V,93.4-8 {3/11}     āstīrṇam , nipūrtāḥ piṇḍāḥ . \\
\noindent \textbf{(P\_7,1.100-102)} KA\_III,277.4-8 Ro\_V,93.4-8 {4/11}     guṇavṛddhyoḥ avakāśaḥ . \\
\noindent \textbf{(P\_7,1.100-102)} KA\_III,277.4-8 Ro\_V,93.4-8 {5/11}     cayanam , cāyakaḥ , lavanam , lāvakaḥ . \\
\noindent \textbf{(P\_7,1.100-102)} KA\_III,277.4-8 Ro\_V,93.4-8 {6/11}     iha ubhayam prāpnoti . \\
\noindent \textbf{(P\_7,1.100-102)} KA\_III,277.4-8 Ro\_V,93.4-8 {7/11}     āstaraṇam , āstārakaḥ , niparaṇam , nipārakaḥ . \\
\noindent \textbf{(P\_7,1.100-102)} KA\_III,277.4-8 Ro\_V,93.4-8 {8/11}     guṇavṛddhī bhavataḥ vipratiṣedhena . \\
\noindent \textbf{(P\_7,1.100-102)} KA\_III,277.4-8 Ro\_V,93.4-8 {9/11}     ayuktaḥ ayam vipratiṣedhaḥ yaḥ ayam guṇasya ittvottayoḥ ca . \\
\noindent \textbf{(P\_7,1.100-102)} KA\_III,277.4-8 Ro\_V,93.4-8 {10/11}     katham . \\
\noindent \textbf{(P\_7,1.100-102)} KA\_III,277.4-8 Ro\_V,93.4-8 {11/11}     nityaḥ guṇaḥ \\
\noindent \textbf{(P\_7,2.1)} sici vṛddhau okārapratiṣedhaḥ . \\
\noindent \textbf{(P\_7,2.1)} sici vṛddhau okārasya pratiṣedhaḥ vaktavyaḥ . \\
\noindent \textbf{(P\_7,2.1)} udavoḍhām , udavoḍham , udavoḍha iti . \\
\noindent \textbf{(P\_7,2.1)} saḥ tarhi pratiṣedhaḥ vaktavyaḥ . \\
\noindent \textbf{(P\_7,2.1)} na vaktavyaḥ . \\
\noindent \textbf{(P\_7,2.1)} okārāt vṛddhiḥ vipratiṣedhena . \\
\noindent \textbf{(P\_7,2.1)} ottvam kriyatām vṛddhiḥ iti . \\
\noindent \textbf{(P\_7,2.1)} vṛddhiḥ bhaviṣyati vipratiṣedhena . \\
\noindent \textbf{(P\_7,2.1)} okārāt vṛddhiḥ vipratiṣedhena iti cet ottvābhāvaḥ . \\
\noindent \textbf{(P\_7,2.1)} okārāt vṛddhiḥ vipratiṣedhena iti cet ottvasya abhāvaḥ . \\
\noindent \textbf{(P\_7,2.1)} udavoḍhām , udavoḍham . \\
\noindent \textbf{(P\_7,2.1)} udavoḍha iti . \\
\noindent \textbf{(P\_7,2.1)} na eṣaḥ doṣaḥ . \\
\noindent \textbf{(P\_7,2.1)} uktam tatra varṇagrahaṇasya prayojanam vṛddhau api kṛtāyām ottvam yathā syāt . \\
\noindent \textbf{(P\_7,2.1)} punaḥprasaṅgavijñānāt vā siddham yathā prasāraṇādiṣu dvirvacanam . \\
\noindent \textbf{(P\_7,2.1)} atha vā punaḥprasaṅgāt atra vṛddhau kṛtāyām ottvam bhaviṣyati . \\
\noindent \textbf{(P\_7,2.1)} sauḍhāmitrau bahiraṅgalakṣaṇatvāt siddham . \\
\noindent \textbf{(P\_7,2.1)} bahiraṅgalakṣaṇā vṛddhiḥ . \\
\noindent \textbf{(P\_7,2.1)} antaraṅgam ottvam . \\
\noindent \textbf{(P\_7,2.1)} asiddham bahiraṅgam antaraṅge \\
\noindent \textbf{(P\_7,2.2)} antagrahaṇam kimartham na ataḥ rlaḥ iti eva ucyate . \\
\noindent \textbf{(P\_7,2.2)} kena idānīm tadantasya bhaviṣyati . \\
\noindent \textbf{(P\_7,2.2)} tadantavidhinā . \\
\noindent \textbf{(P\_7,2.2)} idam tarhi prayojanam . \\
\noindent \textbf{(P\_7,2.2)} ayam antaśabdaḥ asti eva avayavavācī . \\
\noindent \textbf{(P\_7,2.2)} tat yathā . \\
\noindent \textbf{(P\_7,2.2)} vastrāntaḥ , vasanāntaḥ . \\
\noindent \textbf{(P\_7,2.2)} asti sāmīpye vartate . \\
\noindent \textbf{(P\_7,2.2)} tat yathā . \\
\noindent \textbf{(P\_7,2.2)} udakāntam gataḥ . \\
\noindent \textbf{(P\_7,2.2)} udakasamīpam gataḥ iti gamyate . \\
\noindent \textbf{(P\_7,2.2)} tat yaḥ sāmīpye vartate tasya idam grahaṇam yathā vijñāyeta . \\
\noindent \textbf{(P\_7,2.2)} aṅgāntau yau rephalakārau tayoḥ samīpe yaḥ akāraḥ tasya yathā syāt . \\
\noindent \textbf{(P\_7,2.2)} iha mā bhūt . \\
\noindent \textbf{(P\_7,2.2)} aśvallīt , avabhrīt \\
\noindent \textbf{(P\_7,2.3)} halgrahaṇam kimartham . \\
\noindent \textbf{(P\_7,2.3)} samuccayaḥ yathā vijñāyeta . \\
\noindent \textbf{(P\_7,2.3)} vadivrajyoḥ ca halantasya ca acaḥ iti . \\
\noindent \textbf{(P\_7,2.3)} na etat asti prayojanam . \\
\noindent \textbf{(P\_7,2.3)} ajgrahaṇāt eva atra samuccayaḥ bhaviṣyati . \\
\noindent \textbf{(P\_7,2.3)} vadivrajyoḥ ca acaḥ ca iti . \\
\noindent \textbf{(P\_7,2.3)} asti anyat ajgrahaṇe prayojanam vadivrajiviśeṣaṇam yathā vijñāyeta . \\
\noindent \textbf{(P\_7,2.3)} vadivrajyoḥ ecaḥ iti . \\
\noindent \textbf{(P\_7,2.3)} yadi etāvat prayojanam syāt vadivrajyoḥ ataḥ iti evam brūyāt . \\
\noindent \textbf{(P\_7,2.3)} atha vā etat api na brūyāt . \\
\noindent \textbf{(P\_7,2.3)} ataḥ iti vartate . \\
\noindent \textbf{(P\_7,2.3)} idam tarhi prayojanam halantasya yathā syāt ajantasya mā bhūt . \\
\noindent \textbf{(P\_7,2.3)} kasya punaḥ ajantasya prāpnoti . \\
\noindent \textbf{(P\_7,2.3)} akārasya . \\
\noindent \textbf{(P\_7,2.3)} acikīrṣīt , ajihīrṣīt . \\
\noindent \textbf{(P\_7,2.3)} lopaḥ atra bādhakaḥ bhaviṣyati . \\
\noindent \textbf{(P\_7,2.3)} ākārasya tarhi prāpnoti . \\
\noindent \textbf{(P\_7,2.3)} ayāsīt , avāsīt . \\
\noindent \textbf{(P\_7,2.3)} na asti atra viśeṣaḥ satyām vā vṛddhau asatyām vā . \\
\noindent \textbf{(P\_7,2.3)} sandhyakṣarasya tarhi prāpnoti . \\
\noindent \textbf{(P\_7,2.3)} na vai sandhyakṣaram antyam asti . \\
\noindent \textbf{(P\_7,2.3)} nanu ca idam asti ḍhalope kṛte udavoḍhām , udavoḍham , udavoḍha iti . \\
\noindent \textbf{(P\_7,2.3)} asiddhaḥ ḍhalopaḥ tasya asiddhatvāt na etat antyam bhavati . \\
\noindent \textbf{(P\_7,2.3)} ataḥ uttaram paṭhati . \\
\noindent \textbf{(P\_7,2.3)} halgrahaṇam iṭi pratiṣedhārtham . \\
\noindent \textbf{(P\_7,2.3)} halgrahaṇam kriyate iṭi pratiṣedhārtham . \\
\noindent \textbf{(P\_7,2.3)} na iṭi iti pratiṣedham vakṣyati saḥ halantasya yathā syāt ajantasya mā bhūt . \\
\noindent \textbf{(P\_7,2.3)} alāvīt , apāvīt . \\
\noindent \textbf{(P\_7,2.3)} na vā anantarasya pratiṣedhāt . \\
\noindent \textbf{(P\_7,2.3)} na vā etat prayojanam asti . \\
\noindent \textbf{(P\_7,2.3)} kim kāraṇam . \\
\noindent \textbf{(P\_7,2.3)} anantarasya pratiṣedhāt . \\
\noindent \textbf{(P\_7,2.3)} anantaram yat vṛddhividhānam tat pratiṣidhyate . \\
\noindent \textbf{(P\_7,2.3)} kutaḥ etat . \\
\noindent \textbf{(P\_7,2.3)} anantarasya vidhiḥ vā bhavati pratiṣedhaḥ vā iti . \\
\noindent \textbf{(P\_7,2.3)} tat ca anantyārtham . \\
\noindent \textbf{(P\_7,2.3)} tat ca anantaram vṛddhividhānam anantyārtham vijñāyate . \\
\noindent \textbf{(P\_7,2.3)} katham punaḥ anantaram vṛddhividhānam anantyārtham śakyam vijñātum . \\
\noindent \textbf{(P\_7,2.3)} antyasya vacanānarthakyāt . \\
\noindent \textbf{(P\_7,2.3)} antyasya vṛddhividhāne prayojanam na asti iti kṛtvā anantaram vṛddhividhānam anantyārtham vijñāyate . \\
\noindent \textbf{(P\_7,2.3)} ataḥ vibhāṣārtham tarhi idam vaktavyam . \\
\noindent \textbf{(P\_7,2.3)} ataḥ halādeḥ laghoḥ iti vibhāṣā vṛddhim vakṣyati sā halantasya yathā syāt ajantasya mā bhūt . \\
\noindent \textbf{(P\_7,2.3)} acikīrṣīt , ajihīrṣīt . \\
\noindent \textbf{(P\_7,2.3)} ataḥ vibhāṣārtham iti cet siddham vṛddheḥ lopabalīyastvāt . \\
\noindent \textbf{(P\_7,2.3)} ataḥ vibhāṣārtham iti cet tat antareṇa api halgrahaṇam siddham . \\
\noindent \textbf{(P\_7,2.3)} katham . \\
\noindent \textbf{(P\_7,2.3)} vṛddheḥ lopabalīyastvāt . \\
\noindent \textbf{(P\_7,2.3)} vṛddheḥ lopaḥ balīyān bhavati iti . \\
\noindent \textbf{(P\_7,2.3)} idam iha sampradhāryam . \\
\noindent \textbf{(P\_7,2.3)} vṛddhiḥ kriyatām lopaḥ iti kim atra kartavyam . \\
\noindent \textbf{(P\_7,2.3)} paratvāt vṛddhiḥ . \\
\noindent \textbf{(P\_7,2.3)} nityaḥ lopaḥ . \\
\noindent \textbf{(P\_7,2.3)} kṛtāyām api vṛddhau prāpnoti akṛtāyām api . \\
\noindent \textbf{(P\_7,2.3)} anityaḥ lopaḥ na hi kṛtāyām vṛddhau prāpnoti . \\
\noindent \textbf{(P\_7,2.3)} paratvāt sagiḍbhyām bhavitavyam . \\
\noindent \textbf{(P\_7,2.3)} na atra sagiṭau prāpnutaḥ . \\
\noindent \textbf{(P\_7,2.3)} kim kāraṇam . \\
\noindent \textbf{(P\_7,2.3)} ekācaḥ tau vali iti vā . \\
\noindent \textbf{(P\_7,2.3)} ekācaḥ sagiṭau ucyete atha vā vali iti tatra anuvartate . \\
\noindent \textbf{(P\_7,2.3)} kim punaḥ kāraṇam ekācaḥ tau valī iti vā . \\
\noindent \textbf{(P\_7,2.3)} dardrāteḥ mā bhūt iti . \\
\noindent \textbf{(P\_7,2.3)} daridrāteḥ na sagiḍbhyām bhavitavyam . \\
\noindent \textbf{(P\_7,2.3)} uktam etat daridrāteḥ ārdhadhātuke lopaḥ siddhaḥ ca pratyayavidhau iti . \\
\noindent \textbf{(P\_7,2.3)} yaḥ ca idānīm pratyayavidhau siddhaḥ siddhaḥ asau sagiḍvidhau . \\
\noindent \textbf{(P\_7,2.3)} evamartham eva tarhi ekājgrahaṇam anuvartyam atra sagiṭau mā bhūtām iti . \\
\noindent \textbf{(P\_7,2.3)} saḥ eṣaḥ nityaḥ lopaḥ vṛddhim bādhiṣyate . \\
\noindent \textbf{(P\_7,2.3)} kam punaḥ bhavān vṛddheḥ avakāśam matvā āha nityaḥ lopaḥ iti . \\
\noindent \textbf{(P\_7,2.3)} anavakāśā vṛddhiḥ lopam bādhiṣyate . \\
\noindent \textbf{(P\_7,2.3)} sāvakāśā vṛddhiḥ . \\
\noindent \textbf{(P\_7,2.3)} kaḥ avakāśaḥ . \\
\noindent \textbf{(P\_7,2.3)} anantyaḥ : akaṇīt , akāṇīt . \\
\noindent \textbf{(P\_7,2.3)} katham punaḥ sati antye anantyasya vṛddhiḥ syāt . \\
\noindent \textbf{(P\_7,2.3)} bhavet yaḥ atā aṅgam viśeṣayet tasya anantyasya na syāt . \\
\noindent \textbf{(P\_7,2.3)} vayam tu khalu aṅgena akāram viśeṣayiṣyāmaḥ . \\
\noindent \textbf{(P\_7,2.3)} tatra anantyaḥ vṛddheḥ avakāśaḥ antyasya lopaḥ bādhakaḥ bhaviṣyati . \\
\noindent \textbf{(P\_7,2.3)} evam vṛddheḥ lopabalīyastvāt . \\
\noindent \textbf{(P\_7,2.3)} atha vā ārabhyate pūrvavipratiṣedhaḥ ṇyallopau iyaṅyaṇguṇavṛddhidīrghatvebhyaḥ pūrvavipratiṣiddham . \\
\noindent \textbf{(P\_7,2.3)} sā tarhi eṣā anantyārthā vṛddhiḥ halantasya yathā syāt ajantasya mā bhūt . \\
\noindent \textbf{(P\_7,2.3)} apipaṭhiṣīt . \\
\noindent \textbf{(P\_7,2.3)} etat api na asti prayojanam . \\
\noindent \textbf{(P\_7,2.3)} katham . \\
\noindent \textbf{(P\_7,2.3)} halādeḥ iti na eṣā bahuvrīheḥ ṣaṣṭhī : hal ādiḥ yasya saḥ ayam halādiḥ halādeḥ iti . \\
\noindent \textbf{(P\_7,2.3)} kā tarhi . \\
\noindent \textbf{(P\_7,2.3)} karmadhārayāt pañcamī . \\
\noindent \textbf{(P\_7,2.3)} hal ādiḥ halādiḥ halādeḥ parasya iti . \\
\noindent \textbf{(P\_7,2.3)} yadi karmadhārayāt pañcamī acakāsīt atra prāpnoti . \\
\noindent \textbf{(P\_7,2.3)} sicā anantaryam viśeṣayiṣyāmaḥ . \\
\noindent \textbf{(P\_7,2.3)} halādeḥ parasya sici anantarasya iti . \\
\noindent \textbf{(P\_7,2.3)} yadi sicā ānantaryam viśeṣyate akaṇīt , akāṇīt atra na prāpnoti . \\
\noindent \textbf{(P\_7,2.3)} vacanāt bhaviṣyati . \\
\noindent \textbf{(P\_7,2.3)} iha api tarhi vacanāt prāpnoti . \\
\noindent \textbf{(P\_7,2.3)} acakāsīt . \\
\noindent \textbf{(P\_7,2.3)} yena na avyavadhānam tena vyavahite api vacanaprāmāṇyāt . \\
\noindent \textbf{(P\_7,2.3)} kena ca na avyavadhānam . \\
\noindent \textbf{(P\_7,2.3)} varṇena ekena . \\
\noindent \textbf{(P\_7,2.3)} saṅgātena punaḥ vyavadhānam bhavati na ca bhavati . \\
\noindent \textbf{(P\_7,2.3)} yadi sicā ānantaryam viśeṣyate astu bahuvrīheḥ ṣaṣṭhī . \\
\noindent \textbf{(P\_7,2.3)} kasmāt na bhavati . \\
\noindent \textbf{(P\_7,2.3)} apipaṭhiṣīt . \\
\noindent \textbf{(P\_7,2.3)} vyavahitatvāt . \\
\noindent \textbf{(P\_7,2.3)} evam tarhi atidūram eva idam halgrahaṇam anusṛtam . \\
\noindent \textbf{(P\_7,2.3)} halgrahaṇam anantyārtham . \\
\noindent \textbf{(P\_7,2.3)} ajgrahaṇam anigartham \\
\noindent \textbf{(P\_7,2.5)} kimartham jāgarteḥ vṛddhipratiṣedhaḥ ucyate . \\
\noindent \textbf{(P\_7,2.5)} sici vṛddhiḥ mā bhūt iti . \\
\noindent \textbf{(P\_7,2.5)} na etat asti prayojanam . \\
\noindent \textbf{(P\_7,2.5)} jāgarteḥ guṇaḥ ucyate vṛddhiviṣaye pratiṣedhaviṣaye ca saḥ bādhakaḥ bhaviṣyati . \\
\noindent \textbf{(P\_7,2.5)} guṇe tarhi kṛte raparatve ca halantalakṣaṇā vṛddhiḥ prāpnoti . \\
\noindent \textbf{(P\_7,2.5)} na iṭi iti tasyāḥ pratiṣedhaḥ bhaviṣyati . \\
\noindent \textbf{(P\_7,2.5)} iyam tarhi pratiṣedhottarakālā vṛddhiḥ ārabhyate ataḥ rlāntasya iti . \\
\noindent \textbf{(P\_7,2.5)} aparaḥ āha : kakṣyayā kakṣyā nimātavyā . \\
\noindent \textbf{(P\_7,2.5)} sici vṛddhiḥ ca prāpnoti guṇāḥ ca . \\
\noindent \textbf{(P\_7,2.5)} guṇaḥ bhavati . \\
\noindent \textbf{(P\_7,2.5)} guṇe kṛte raparatve ca halantalakṣaṇā vṛddhiḥ prāpnoti na iṭi iti ca tasyāḥ pratiṣedhaḥ bhavati . \\
\noindent \textbf{(P\_7,2.5)} pratiṣedhottarakālam ataḥ halādeḥ laghoḥ iti vibhāṣā vṛddhiḥ prāpnoti na ca kim cit . \\
\noindent \textbf{(P\_7,2.5)} ataḥ rlāntasya iti ca vṛddhiḥ prāpnoti na ca kim cit \\
\noindent \textbf{(P\_7,2.8.1)} kimartham purastāt pratiṣedhaḥ ucyate na vidhyuttarakālaḥ pratiṣedhaḥ kriyeta . \\
\noindent \textbf{(P\_7,2.8.1)} tat yathā anyatra api vidhyuttarakālāḥ pratiṣedhāḥ bhavanti . \\
\noindent \textbf{(P\_7,2.8.1)} kva anyatra . \\
\noindent \textbf{(P\_7,2.8.1)} kartari karmavyatihāre na gatihiṃsārthebhyaḥ iti . \\
\noindent \textbf{(P\_7,2.8.1)} devatādvandve ca na indrasya parasya . \\
\noindent \textbf{(P\_7,2.8.1)} tatra ayam api arthaḥ dviḥ iḍgrahaṇam na kartavyam bhavati prakṛtam anuvartate . \\
\noindent \textbf{(P\_7,2.8.1)} na evam śakyam . \\
\noindent \textbf{(P\_7,2.8.1)} iḍartham sārvadhātukagrahaṇam liṅaḥ salope sannihitam tat vicchidyeta . \\
\noindent \textbf{(P\_7,2.8.1)} yadi punaḥ na vṛdbhyaḥ caturbhyaḥ iti atra eva ucyeta . \\
\noindent \textbf{(P\_7,2.8.1)} kim kṛtam bhavati . \\
\noindent \textbf{(P\_7,2.8.1)} vidhyuttarakālāḥ ca eva pratiṣedhaḥ kṛataḥ bhavati dviḥ ca iḍgrahaṇam na kartavyam iḍartham ca sārvadhātukagrahaṇam liṅaḥ salope sannihitam bhavati . \\
\noindent \textbf{(P\_7,2.8.1)} tatra ayam api arthaḥ dviḥ pratiṣedhaḥ na kartavyaḥ iti etasmāt niyamāt iṭ prasajyeta . \\
\noindent \textbf{(P\_7,2.8.1)} kṛsṛbhṛvṛstudruśrusruvaḥ liṭi iti eṣaḥ yogaḥ pratiṣedhārthaḥ bhaviṣyati . \\
\noindent \textbf{(P\_7,2.8.1)} yadi eṣaḥ yogaḥ pratiṣedhārthaḥ yaḥ etasmāt yogāt iṭ pariprāpyate niyamāt saḥ na sidhyati : peciva , pecima , śekiva, śekima . \\
\noindent \textbf{(P\_7,2.8.1)} evam tarhi kṛsṛbhṛ , iti eteṣām grahaṇam niyamārtham bhaviṣyati studruśrusruvām pratiṣedhārtham vṛṅvṛñoḥ jñāpakārtham . \\
\noindent \textbf{(P\_7,2.8.1)} evam api sāmānyavihitasya eva iṭaḥ pratiṣedhaḥ vijñāyeta viśeṣavihitaḥ ca ayam thali iti . \\
\noindent \textbf{(P\_7,2.8.1)} purastāt punaḥ pratiṣedhe sati anārabhyāpavādaḥ ayam bhavati tena yāvān iṇ nāma tasya sarvasya eva pratiṣedhaḥ siddhaḥ bhavati . \\
\noindent \textbf{(P\_7,2.8.1)} yadi khalu api eṣaḥ abhiprāyaḥ tat na kriyate iti purastāt api pratiṣedhe sati tat na kariṣyate . \\
\noindent \textbf{(P\_7,2.8.1)} katham . \\
\noindent \textbf{(P\_7,2.8.1)} idam asti na iṭ vaśi kṛti iti . \\
\noindent \textbf{(P\_7,2.8.1)} tataḥ vakṣyāmi . \\
\noindent \textbf{(P\_7,2.8.1)} ārdhadhātukasya valādeḥ iti . \\
\noindent \textbf{(P\_7,2.8.1)} iṭ iti anuvartate na iti nivṛttam \\
\noindent \textbf{(P\_7,2.8.2)} atha kṛdgrahaṇam kimartham . \\
\noindent \textbf{(P\_7,2.8.2)} iha mā bhūt . \\
\noindent \textbf{(P\_7,2.8.2)} bibhidiva , bibhidima iti . \\
\noindent \textbf{(P\_7,2.8.2)} na etat asti prayojanam . \\
\noindent \textbf{(P\_7,2.8.2)} kṛsṛbhṛvṛstudruśrusruvaḥ liṭi iti etasmāt niyamāt atra iṭ bhaviṣyati . \\
\noindent \textbf{(P\_7,2.8.2)} na atra tena pariprāpaṇam prāpnoti . \\
\noindent \textbf{(P\_7,2.8.2)} kim kāraṇam . \\
\noindent \textbf{(P\_7,2.8.2)} prakṛtilakṣaṇasya pratiṣedhasya saḥ pratyārambhaḥ pratyayalakṣaṇaḥ ca ayam pratiṣedhaḥ . \\
\noindent \textbf{(P\_7,2.8.2)} ubhayoḥ saḥ pratyārambhaḥ . \\
\noindent \textbf{(P\_7,2.8.2)} katham jñāyate . \\
\noindent \textbf{(P\_7,2.8.2)} vṛṅvṛñoḥ grahaṇāt . \\
\noindent \textbf{(P\_7,2.8.2)} katham kṛtvā jñāpakam . \\
\noindent \textbf{(P\_7,2.8.2)} imau vṛṅvṛñau udāttau tayoḥ prakṛtilakṣaṇaḥ pratyayalakṣaṇaḥ ca . \\
\noindent \textbf{(P\_7,2.8.2)} tataḥ kim . \\
\noindent \textbf{(P\_7,2.8.2)} tulyajātīye asati yathā eva prakṛtilakṣaṇasya niyāmakaḥ bhavati evam pratyayalakṣaṇasya api niyāmakaḥ bhaviṣyati . \\
\noindent \textbf{(P\_7,2.8.2)} idam tarhi prayojanam . \\
\noindent \textbf{(P\_7,2.8.2)} iha mā bhūt . \\
\noindent \textbf{(P\_7,2.8.2)} rudivaḥ , rudimaḥ . \\
\noindent \textbf{(P\_7,2.8.2)} etat api na asti prayojanam . \\
\noindent \textbf{(P\_7,2.8.2)} upariṣṭāt yogavibhāgaḥ kariṣyate . \\
\noindent \textbf{(P\_7,2.8.2)} ārdhadhātukasya . \\
\noindent \textbf{(P\_7,2.8.2)} yat etat anukrāntam etat ārdhadhātukasya draṣṭavyam . \\
\noindent \textbf{(P\_7,2.8.2)} tataḥ iṭ valādeḥ iti . \\
\noindent \textbf{(P\_7,2.8.2)} tatra etāvat draṣṭavyam yadi kim cit tatra anyat api ārdhadhātukagrahaṇasya prayojanam asti . \\
\noindent \textbf{(P\_7,2.8.2)} atha na kim cit iha vā kṛdgrahaṇam kriyeta tatra vā ārdhadhātukagrahaṇam kaḥ nu atra viśeṣaḥ . \\
\noindent \textbf{(P\_7,2.8.2)} na iṭ varam anādau kṛti . \\
\noindent \textbf{(P\_7,2.8.2)} varam anādau kṛti iṭpratiṣedham prayojayati . \\
\noindent \textbf{(P\_7,2.8.2)} īśitā , īśitum , īśvaraḥ . \\
\noindent \textbf{(P\_7,2.8.2)} va . \\
\noindent \textbf{(P\_7,2.8.2)} ra . \\
\noindent \textbf{(P\_7,2.8.2)} dīpitā , dīpitum , dīpraḥ . \\
\noindent \textbf{(P\_7,2.8.2)} ra . \\
\noindent \textbf{(P\_7,2.8.2)} ma . \\
\noindent \textbf{(P\_7,2.8.2)} bhasitā , bhasitum , bhasma . \\
\noindent \textbf{(P\_7,2.8.2)} ma . \\
\noindent \textbf{(P\_7,2.8.2)} na . \\
\noindent \textbf{(P\_7,2.8.2)} yatitā , yatitum , yatnaḥ . \\
\noindent \textbf{(P\_7,2.8.2)} atha anye ye vaśādayaḥ tatra katham . \\
\noindent \textbf{(P\_7,2.8.2)} uṇādayaḥ avyutpannāni prātipadikāni \\
\noindent \textbf{(P\_7,2.9)} titutreṣu agrahādīnām . \\
\noindent \textbf{(P\_7,2.9)} titutreṣu agrahādīnām iti vaktavyam . \\
\noindent \textbf{(P\_7,2.9)} iha mā bhūt . \\
\noindent \textbf{(P\_7,2.9)} nigṛhitiḥ , upasnihitaḥ , nikucitiḥ , nipaṭhitiḥ iti \\
\noindent \textbf{(P\_7,2.10)} ekājgrahaṇam kimartham . \\
\noindent \textbf{(P\_7,2.10)} ekājgrahaṇam jāgartyartham . \\
\noindent \textbf{(P\_7,2.10)} ekājgrahaṇam kriyate jāgarteḥ iṭpratiṣedhaḥ mā bhūt iti . \\
\noindent \textbf{(P\_7,2.10)} jāgaritā , jāgaritum . \\
\noindent \textbf{(P\_7,2.10)} na etat asti prayojanam . \\
\noindent \textbf{(P\_7,2.10)} upadeśe anudāttāt iti ucyate jāgartiḥ ca upadeśe udāttaḥ . \\
\noindent \textbf{(P\_7,2.10)} na brūmaḥ ihārtham jāgartyartham ekājgrahaṇam kartavyam iti . \\
\noindent \textbf{(P\_7,2.10)} kim tarhi . \\
\noindent \textbf{(P\_7,2.10)} uttarārtham . \\
\noindent \textbf{(P\_7,2.10)} śryukaḥ kiti iti iṭpratiṣedham vakṣyati saḥ jāgarteḥ mā bhūt . \\
\noindent \textbf{(P\_7,2.10)} jāgaritaḥ , jāgaritavān iti . \\
\noindent \textbf{(P\_7,2.10)} etat api na asti prayojanam . \\
\noindent \textbf{(P\_7,2.10)} jāgarteḥ guṇaḥ ucyate vṛddhiviṣaye pratiṣedhaviṣaye ca saḥ bādhakaḥ bhaviṣyati . \\
\noindent \textbf{(P\_7,2.10)} tatra guṇe kṛte raparatve ca kṛte anugantatvāt iṭpratiṣedhaḥ na bhaviṣyati . \\
\noindent \textbf{(P\_7,2.10)} nanu ca upadeśādhikārāt prāpnoti . \\
\noindent \textbf{(P\_7,2.10)} upadeśagrahaṇam nivartayiṣyate . \\
\noindent \textbf{(P\_7,2.10)} yadi nivartyate stīrtvā , pūrtvā , ittvottvayoḥ kṛtayoḥ raparatve ca anugantatvāt iṭpratiṣedhaḥ na prāpnoti . \\
\noindent \textbf{(P\_7,2.10)} na eṣaḥ doṣaḥ . \\
\noindent \textbf{(P\_7,2.10)} ānupūrvyā siddham etat . \\
\noindent \textbf{(P\_7,2.10)} na atra akṛte iṭpratiṣedhe ittvottve prāpnutaḥ . \\
\noindent \textbf{(P\_7,2.10)} kim kāraṇam . \\
\noindent \textbf{(P\_7,2.10)} na ktvā seṭ iti kittvapratiṣedhāt . \\
\noindent \textbf{(P\_7,2.10)} idam tarhi . \\
\noindent \textbf{(P\_7,2.10)} ātistīrṣati , nipupūrṣati . \\
\noindent \textbf{(P\_7,2.10)} ittvottvayoḥ kṛtayoḥ raparatve cxa anugantatvāt iṭpratiṣedhaḥ na prāpnoti . \\
\noindent \textbf{(P\_7,2.10)} mā bhūt evam śryukaḥ kiti iti . \\
\noindent \textbf{(P\_7,2.10)} iṭ sani vā iti evam bhaviṣyati . \\
\noindent \textbf{(P\_7,2.10)} idam tarhi . \\
\noindent \textbf{(P\_7,2.10)} āstīrṇam , nipūrtāḥ piṇḍāḥ . \\
\noindent \textbf{(P\_7,2.10)} ittvottvayoḥ kṛtayoḥ raparatve ca anugantatvāt iṭpratiṣedhaḥ na prāpnoti . \\
\noindent \textbf{(P\_7,2.10)} mā bhūt evam . \\
\noindent \textbf{(P\_7,2.10)} iṭ sani vā iti sani vibhāṣā yasya vibhāṣā iti pratiṣedhaḥ bhaviṣyati . \\
\noindent \textbf{(P\_7,2.10)} ihārtham eva tarhi vadhyartham ekājgrahaṇam kartavyam . \\
\noindent \textbf{(P\_7,2.10)} vadhaḥ iṭpratiṣedhaḥ mā bhūt iti . \\
\noindent \textbf{(P\_7,2.10)} vadhiṣīṣṭa iti . \\
\noindent \textbf{(P\_7,2.10)} etat api na asti prayojanam . \\
\noindent \textbf{(P\_7,2.10)} kriyamāṇe api vā ekājgrahaṇe vadhaḥ iṭpratiṣedhaḥ prāpnoti . \\
\noindent \textbf{(P\_7,2.10)} vadhiṣīṣṭa iti . \\
\noindent \textbf{(P\_7,2.10)} kim kāraṇam . \\
\noindent \textbf{(P\_7,2.10)} vadhaḥ iṭpratiṣedhaḥ sannipāte ekāctvāt prakṛteḥ ca anudāttatvāt . \\
\noindent \textbf{(P\_7,2.10)} sannipāte ca eva hi vadhiḥ ekāc śrūyate prakṛtiḥ ca asya anudāttā . \\
\noindent \textbf{(P\_7,2.10)} kim punaḥ kāraṇam evam vijñāyate upadeśe anudāttāt ekācaḥ śrūyamāṇāt iti . \\
\noindent \textbf{(P\_7,2.10)} yaṅlopārtham . \\
\noindent \textbf{(P\_7,2.10)} yaṅlope mā bhūt iti . \\
\noindent \textbf{(P\_7,2.10)} bebhiditā , bebhiditum , cecchiditā , cecchiditum . \\
\noindent \textbf{(P\_7,2.10)} ekācaḥ upadeśe anudāttāt iti upadeśavacanam anudāttaviśeṣaṇam cet kṛñādibhyaḥ liṭi niyamānupapattiḥ aprāptatvāt pratiṣedhasya . \\
\noindent \textbf{(P\_7,2.10)} ekācaḥ upadeśe anudāttāt iti upadeśavacanam anudāttaviśeṣaṇam cet kṛñādibhyaḥ liṭi niyamasya anupapattiḥ . \\
\noindent \textbf{(P\_7,2.10)} kim kāraṇam . \\
\noindent \textbf{(P\_7,2.10)} aprāptatvāt pratiṣedhasya . \\
\noindent \textbf{(P\_7,2.10)} dvirvacane kṛte upadeśe anudāttāt ekācaḥ śrūyamāṇāt iti iṭpratiṣedhaḥ na prāpnoti . \\
\noindent \textbf{(P\_7,2.10)} asati iṭpratiṣedhe niyamaḥ na upapadyate . \\
\noindent \textbf{(P\_7,2.10)} asati niyame kaḥ doṣaḥ . \\
\noindent \textbf{(P\_7,2.10)} tatra pacādibhyaḥ iḍvacanam . \\
\noindent \textbf{(P\_7,2.10)} tatra pacādibhyaḥ iṭ vaktavyaḥ . \\
\noindent \textbf{(P\_7,2.10)} pecima , śekima . \\
\noindent \textbf{(P\_7,2.10)} sanaḥ ca iṭpratiṣedhaḥ . \\
\noindent \textbf{(P\_7,2.10)} sanaḥ ca iṭpratiṣedhaḥ vaktavyaḥ . \\
\noindent \textbf{(P\_7,2.10)} bibhitsati , cicchitsati . \\
\noindent \textbf{(P\_7,2.10)} dvirvacane kṛte upadeśe anudāttāt ekācaḥ śrūyamāṇāt iti iṭpratiṣedhaḥ na prāpnoti . \\
\noindent \textbf{(P\_7,2.10)} iha ca nīttaḥ tatve kṛte anackatvāt iṭpratiṣedhaḥ na prāpnoti . \\
\noindent \textbf{(P\_7,2.10)} na eṣaḥ doṣaḥ . \\
\noindent \textbf{(P\_7,2.10)} ānupūrvyā siddham etat . \\
\noindent \textbf{(P\_7,2.10)} na atra akṛte iṭpratiṣedhe tatvam prāpnoti . \\
\noindent \textbf{(P\_7,2.10)} kim kāraṇam . \\
\noindent \textbf{(P\_7,2.10)} ti kiti iti ucyate . \\
\noindent \textbf{(P\_7,2.10)} yat api ucyate ekācaḥ upadeśe anudāttāt iti upadeśavacanam anudāttaviśeṣaṇam cet kṛñādibhyaḥ liṭi niyamānupapattiḥ aprāptatvāt pratiṣedhasya iti mā bhūt niyamaḥ . \\
\noindent \textbf{(P\_7,2.10)} nanu ca uktam tatra pacādibhyaḥ iḍvacanam iti . \\
\noindent \textbf{(P\_7,2.10)} na eṣaḥ doṣaḥ . \\
\noindent \textbf{(P\_7,2.10)} uktam tatra thalgrahaṇasya prayojanam samuccayaḥ yathā vijñāyeta thali ca seṭi kṅiti ca seṭi iti . \\
\noindent \textbf{(P\_7,2.10)} yat api ucyate sanaḥ ca iṭpratiṣedhaḥ iti . \\
\noindent \textbf{(P\_7,2.10)} ubhayaviśeṣaṇatvāt siddham . \\
\noindent \textbf{(P\_7,2.10)} ubhayam upadeśagrahaṇena viśeṣayiṣyāmaḥ . \\
\noindent \textbf{(P\_7,2.10)} upadeśe anudāttāt upadeśe ekācaḥ iti . \\
\noindent \textbf{(P\_7,2.10)} yaṅlope ca tadantadvirvacanāt . \\
\noindent \textbf{(P\_7,2.10)} sanyaṅantasya sthāne dvirvacanam tatra sampramugdhatvāt prakṛtipratyayasya naṣṭaḥ saḥ bhavati yaḥ saḥ ekājupadeśe anudāttaḥ . \\
\noindent \textbf{(P\_7,2.10)} atha api dviḥprayogaḥ dvirvacanam evam api na doṣaḥ . \\
\noindent \textbf{(P\_7,2.10)} na hi asya bhidyupadeśe upadeśaḥ . \\
\noindent \textbf{(P\_7,2.10)} atha api bhidyupadeśe upadeśaḥ evam api na doṣaḥ . \\
\noindent \textbf{(P\_7,2.10)} akāreṇa vyavahitatvāt na bhaviṣyati . \\
\noindent \textbf{(P\_7,2.10)} nanu ca lope kṛte na asti vyavadhānam . \\
\noindent \textbf{(P\_7,2.10)} sthānivadbhāvāt vyavadhānam eva . \\
\noindent \textbf{(P\_7,2.10)} na sidhyati . \\
\noindent \textbf{(P\_7,2.10)} pūrvavidhau sthānivadbhāvaḥ na ca ayam pūrvavidhiḥ . \\
\noindent \textbf{(P\_7,2.10)} evam tarhi pūrvasmāt api vidhiḥ pūrvavidhiḥ . \\
\noindent \textbf{(P\_7,2.10)} kaḥ punaḥ upadeśaḥ nyāyyaḥ . \\
\noindent \textbf{(P\_7,2.10)} yaḥ kṛtsnaḥ . \\
\noindent \textbf{(P\_7,2.10)} kaḥ ca kṛtsnaḥ . \\
\noindent \textbf{(P\_7,2.10)} yaḥ ubhayoḥ . \\
\noindent \textbf{(P\_7,2.10)} yadi tarhi yaḥ ubhayoḥ saḥ kṛtsnaḥ saḥ ca nyāyyaḥ vadhaḥ iṭpratiṣedhaḥ prāpnoti . \\
\noindent \textbf{(P\_7,2.10)} āvadhiṣīṣṭa iti . \\
\noindent \textbf{(P\_7,2.10)} na eṣaḥ doṣaḥ . \\
\noindent \textbf{(P\_7,2.10)} ādyudāttanipātanam kariṣyate saḥ nipātanasvaraḥ prakṛtisvarasya bādhakaḥ bhaviṣyati . \\
\noindent \textbf{(P\_7,2.10)} evam api upadeśivadbhāvaḥ vaktavyaḥ . \\
\noindent \textbf{(P\_7,2.10)} yathā eva hi saḥ nipātanasvaraḥ prakṛtisvaram bādhate evam pratyayasvaram api bādheta : āvadhiṣīṣṭa iti . \\
\noindent \textbf{(P\_7,2.10)} na eṣaḥ doṣaḥ . \\
\noindent \textbf{(P\_7,2.10)} ārdhadhātukīyāḥ sāmānyena bhavanti anavasthiteṣu pratyayeṣu . \\
\noindent \textbf{(P\_7,2.10)} tatra ārdhadhātukasāmānye vadhibhāve kṛte satiśiṣṭatvāt pratyayasvaraḥ bhaviṣyati . \\
\noindent \textbf{(P\_7,2.10)} atha ke punaḥ anudāttāḥ . \\
\noindent \textbf{(P\_7,2.10)} ādantāḥ , adaridrāḥ . \\
\noindent \textbf{(P\_7,2.10)} ivarṇāntāḥ ca aśviśriḍīśīdīdhīvevīṅaḥ . \\
\noindent \textbf{(P\_7,2.10)} uvarṇāntāḥ yuruṇukṣukṣṇusnūrṇuvarjam . \\
\noindent \textbf{(P\_7,2.10)} ṛdantāḥ ca ajāgṛvṛṅvṛñaḥ . \\
\noindent \textbf{(P\_7,2.10)} śakiḥ kavargāntānām . \\
\noindent \textbf{(P\_7,2.10)} pacivacisicimuciricivicipracchiyajibhajisṛjityajibhujibhrasjibhañjirujiyujiṇijivijisiñjisvañjayaḥ cavargāntānām . \\
\noindent \textbf{(P\_7,2.10)} sadiśadihadicchiditudisvidibhidiskandikṣudikhidyativindividyatirādhiyudhibudhiśudhikrudhirudhisādhivyadhibandhisidhyatihanimanyatayaḥ tavargāntānām . \\
\noindent \textbf{(P\_7,2.10)} tapitipivapiśapicupilupilipisvapyāpikṣipisṛpitṛpidṛpiyabhirabhilabhiyamiraminamigamayaḥ pavargāntānām . \\
\noindent \textbf{(P\_7,2.10)} ruśiriśidiśiviśiliśispṛśidṛśikruśimṛśidaṃśitviṣikṛṣiśliṣiviṣipiṣituṣiduṣidviṣighasivasidahidihivahiduhinahiruhilihimihayaḥ ca uṣmāntānām . \\
\noindent \textbf{(P\_7,2.10)} vasiḥ prasāraṇī \\
\noindent \textbf{(P\_7,2.13)} kṛñaḥ asuṭaḥ . \\
\noindent \textbf{(P\_7,2.13)} kṛñaḥ asuṭaḥ iti vaktavyam . \\
\noindent \textbf{(P\_7,2.13)} iha mā bhūt . \\
\noindent \textbf{(P\_7,2.13)} sañcaskariva , sañcaskarima . \\
\noindent \textbf{(P\_7,2.13)} tat tarhi vaktavyam . \\
\noindent \textbf{(P\_7,2.13)} na vaktavyam . \\
\noindent \textbf{(P\_7,2.13)} guṇe kṛte raparatve ca anugantatvāt iṭpratiṣedhaḥ na bhaviṣyati . \\
\noindent \textbf{(P\_7,2.13)} evam api upadeśādhikārāt prāpnoti . \\
\noindent \textbf{(P\_7,2.13)} tasmāt asuṭaḥ iti vaktavyam \\
\noindent \textbf{(P\_7,2.14)} śvigrahaṇam kimartham na prasāraṇe kṛte prasāraṇapūrvatve ca ugantāt iti eva siddham . \\
\noindent \textbf{(P\_7,2.14)} ataḥ uttaram paṭhati . \\
\noindent \textbf{(P\_7,2.14)} śvigrahaṇam idantatvāt upadeśasya . \\
\noindent \textbf{(P\_7,2.14)} śvigrahaṇam kriyate idantatvāt upadeśasya . \\
\noindent \textbf{(P\_7,2.14)} upadeśaḥ ugantāt iti ucyate śvayatiḥ ca upadeśaḥ idantaḥ \\
\noindent \textbf{(P\_7,2.15)} yasya vibhāṣā avideḥ . \\
\noindent \textbf{(P\_7,2.15)} yasya vibhāṣā avideḥ iti vaktavyam . \\
\noindent \textbf{(P\_7,2.15)} iha mā bhūt . \\
\noindent \textbf{(P\_7,2.15)} viditaḥ , viditavān iti . \\
\noindent \textbf{(P\_7,2.15)} tat tarhi vaktavyam . \\
\noindent \textbf{(P\_7,2.15)} na vaktavyam . \\
\noindent \textbf{(P\_7,2.15)} yadupādheḥ vibhāṣā tadupādheḥ pratiṣedhaḥ . \\
\noindent \textbf{(P\_7,2.15)} śābvikaraṇasya vibhāṣā lugvikaraṇaḥ ca ayam \\
\noindent \textbf{(P\_7,2.16-17)} KA\_III,286.20-287.4 Ro\_V,117.8-118.2 {1/6}     kimarthaḥ yogavibhāgaḥ na āditaḥ vibhāṣā bhāvādikarmaṇoḥ iti eva ucyate . \\
\noindent \textbf{(P\_7,2.16-17)} KA\_III,286.20-287.4 Ro\_V,117.8-118.2 {2/6}     kena idānīm kartari pratiṣedhaḥ bhaviṣyati . \\
\noindent \textbf{(P\_7,2.16-17)} KA\_III,286.20-287.4 Ro\_V,117.8-118.2 {3/6}     yasya vibhāṣā iti anena . \\
\noindent \textbf{(P\_7,2.16-17)} KA\_III,286.20-287.4 Ro\_V,117.8-118.2 {4/6}     evam tarhi siddhe sati yat yogavibhāgam karoti tat jñāpayati ācāryaḥ yadupādheḥ vibhāṣā tadupādheḥ pratiṣedhaḥ iti . \\
\noindent \textbf{(P\_7,2.16-17)} KA\_III,286.20-287.4 Ro\_V,117.8-118.2 {5/6}     kim etasya jñāpane prayojanam . \\
\noindent \textbf{(P\_7,2.16-17)} KA\_III,286.20-287.4 Ro\_V,117.8-118.2 {6/6}     yasya vibhāṣā avideḥ iti uktam tat na vaktavyam bhavati \\
\noindent \textbf{(P\_7,2.18)} kṣubdham manthābhidhāne . \\
\noindent \textbf{(P\_7,2.18)} kṣubdham manthābhidhāne iti vaktavyam . \\
\noindent \textbf{(P\_7,2.18)} kṣubhitam manthena iti eva anyatra . \\
\noindent \textbf{(P\_7,2.18)} kṣubdha . \\
\noindent \textbf{(P\_7,2.18)} svānta . \\
\noindent \textbf{(P\_7,2.18)} svāntam mano'bhidhāne . \\
\noindent \textbf{(P\_7,2.18)} svāntam mano'bhidhāne iti vaktavyam . \\
\noindent \textbf{(P\_7,2.18)} svanitam manasā iti eva anyatra . \\
\noindent \textbf{(P\_7,2.18)} svānta . \\
\noindent \textbf{(P\_7,2.18)} dhvānta .[ dhvāntam tamo'bhidhāne .] dhvāntam tamo'bhidhāne iti vaktavyam . \\
\noindent \textbf{(P\_7,2.18)} dhvanitam tamasā iti eva anyatra . \\
\noindent \textbf{(P\_7,2.18)} [R lagna . \\
\noindent \textbf{(P\_7,2.18)} lagnam saktābhidhāne . \\
\noindent \textbf{(P\_7,2.18)} lagnam saktābhidhāne iti vaktavyam . \\
\noindent \textbf{(P\_7,2.18)} lagitam saktena iti eva anyatra . \\
\noindent \textbf{(P\_7,2.18)} lagna . \\
\noindent \textbf{(P\_7,2.18)} mliṣṭa . \\
\noindent \textbf{(P\_7,2.18)} mliṣṭam avispaṣṭābhidhāne . \\
\noindent \textbf{(P\_7,2.18)} mliṣṭam avispaṣṭābhidhāne iti vaktavyam . \\
\noindent \textbf{(P\_7,2.18)} mlecchitam vispaṣṭena iti eva anyatra . \\
\noindent \textbf{(P\_7,2.18)} mliṣṭā . \\
\noindent \textbf{(P\_7,2.18)} viribdha . \\
\noindent \textbf{(P\_7,2.18)} viribdham svarābhidhāne . \\
\noindent \textbf{(P\_7,2.18)} viribdham svarābhidhāne iti vaktavyam . \\
\noindent \textbf{(P\_7,2.18)} virebhitam svareṇa iti eva anyatra . \\
\noindent \textbf{(P\_7,2.18)} viribdha . \\
\noindent \textbf{(P\_7,2.18)} phāṇṭa . \\
\noindent \textbf{(P\_7,2.18)} phāṇṭam anāyāsābhidhāne . \\
\noindent \textbf{(P\_7,2.18)} phāṇṭam anāyāsābhidhāne iti vaktavyam . \\
\noindent \textbf{(P\_7,2.18)} phaṇitam eva anyatra . \\
\noindent \textbf{(P\_7,2.18)} phāṇṭa . \\
\noindent \textbf{(P\_7,2.18)} bāḍha . \\
\noindent \textbf{(P\_7,2.18)} bāḍham bṛśābhidhāne . \\
\noindent \textbf{(P\_7,2.18)} bāḍham bṛśābhidhāne iti vaktavyam . \\
\noindent \textbf{(P\_7,2.18)} bāhitam eva anyatra . \\
\noindent \textbf{(P\_7,2.18)} ] \\
\noindent \textbf{(P\_7,2.19)} kim idam vaiyātye iti . \\
\noindent \textbf{(P\_7,2.19)} viyātabhāvaḥ vaiyātyam \\
\noindent \textbf{(P\_7,2.20)} dṛḍhanipātanam kimartham na dṛheḥ na iṭ bhavati iti eva ucyeta . \\
\noindent \textbf{(P\_7,2.20)} dṛḍhanipātanam nakārahakāralopārtham parasya ca ḍatvārtham . \\
\noindent \textbf{(P\_7,2.20)} dṛḍhanipātanam kriyate nakārahakāralopārtham . \\
\noindent \textbf{(P\_7,2.20)} nakārahakāralopaḥ yathā syāt . \\
\noindent \textbf{(P\_7,2.20)} parasya ca ḍhatvārtham . \\
\noindent \textbf{(P\_7,2.20)} parasya ca ḍhatvam yathā syāt . \\
\noindent \textbf{(P\_7,2.20)} aniḍvacane hi rabhāvāprasiddhiḥ alaghutvāt . \\
\noindent \textbf{(P\_7,2.20)} aniḍvacane hi rabhāvasya aprasiddhiḥ . \\
\noindent \textbf{(P\_7,2.20)} draḍhīyān . \\
\noindent \textbf{(P\_7,2.20)} kim kāraṇam . \\
\noindent \textbf{(P\_7,2.20)} alaghutvāt . \\
\noindent \textbf{(P\_7,2.20)} nalopavacanam ca . \\
\noindent \textbf{(P\_7,2.20)} nalopaḥ ca vaktavyaḥ . \\
\noindent \textbf{(P\_7,2.20)} iha ca paridraḍhayya gataḥ lyapi laghupūrvasya iti ayādeśaḥ na syāt . \\
\noindent \textbf{(P\_7,2.20)} iha ca pāridṛḍhī kanyā iti gurūpottamalakṣaṇaḥ ṣyaṅ prasajyeta \\
\noindent \textbf{(P\_7,2.21)} parivṛḍhaḥ iti kimartham nipātyate na paripūrvāt vṛheḥ na iṭ bhavati iti eva ucyeta . \\
\noindent \textbf{(P\_7,2.21)} parivṛḍhanipātanam ca . \\
\noindent \textbf{(P\_7,2.21)} kim . \\
\noindent \textbf{(P\_7,2.21)} nakārahakāralopārtham parasya ca ḍhatvārtham aniḍvacane hi rabhāvāprasiddhiḥ alaghutvāt nalopavacanam ca iti eva . \\
\noindent \textbf{(P\_7,2.21)} parivraḍhīyān iti raḥ ṛtaḥ halādeḥ laghoḥ iti rabhāvaḥ na syāt . \\
\noindent \textbf{(P\_7,2.21)} iha ca parivraḍhayya gataḥ iti lyapi laghupūrvasya iti ayādeśaḥ na syāt . \\
\noindent \textbf{(P\_7,2.21)} iha ca pārivṛḍhī kanyā iti gurūpottamalakṣaṇaḥ ṣyaṅ prasajyeta \\
\noindent \textbf{(P\_7,2.23)} kimartham aviśabdane iti ucyate na viśabdane curādiṇicā bhavitavyam . \\
\noindent \textbf{(P\_7,2.23)} evam tarhi siddhe sati yat ayam aviśabdane iti āha tat jñāpayati ācāryaḥ viśabdane ghuṣeḥ vibhāṣā ṇic bhavati iti . \\
\noindent \textbf{(P\_7,2.23)} kim etasya jñāpane prayojanam . \\
\noindent \textbf{(P\_7,2.23)} mahīpālavacaḥ śrutvā jughuṣuḥ puṣyamāṇavāḥ . \\
\noindent \textbf{(P\_7,2.23)} eṣaḥ prayogaḥ upapannaḥ bhavati \\
\noindent \textbf{(P\_7,2.26)} kim idam adhyayanābhidhāyikāyām niṣṭhāyām nipātanam kriyate āhosvit adhyayane cet vṛtiḥ vartate iti . \\
\noindent \textbf{(P\_7,2.26)} kim ca ataḥ . \\
\noindent \textbf{(P\_7,2.26)} yadi adhyayanābhidhāyikāyām niṣṭhāyām nipātanam kriyate siddham vṛttaḥ guṇaḥ vṛttam pārāyaṇam vṛttam guṇasya vṛttam pārāyaṇasya iti na sidhyati . \\
\noindent \textbf{(P\_7,2.26)} atha vijñāyate adhyayane cet vṛtiḥ vartate iti na doṣaḥ bhavati . \\
\noindent \textbf{(P\_7,2.26)} yathā na doṣaḥ tathā astu . \\
\noindent \textbf{(P\_7,2.26)} adhyayane cet vṛtiḥ vartate iti api vai vijñāyamāne na sidhyati . \\
\noindent \textbf{(P\_7,2.26)} kim kāraṇam . \\
\noindent \textbf{(P\_7,2.26)} vṛtiḥ ayam akarmakaḥ . \\
\noindent \textbf{(P\_7,2.26)} akarmakāḥ ca api ṇyantāḥ sakarmakāḥ bhavanti . \\
\noindent \textbf{(P\_7,2.26)} akarmkaḥ ca atra vṛtiḥ . \\
\noindent \textbf{(P\_7,2.26)} katham punaḥ jñāyate akarmakaḥ atra vṛtiḥ iti . \\
\noindent \textbf{(P\_7,2.26)} akarmakāṇām bhāve ktaḥ bhavati iti evam atra bhāve ktaḥ bhavati . \\
\noindent \textbf{(P\_7,2.26)} tatra uditaḥ ktvi vibhāṣā yasya vibhāṣā iti iṭpratiṣedhaḥ bhaviṣyati . \\
\noindent \textbf{(P\_7,2.26)} atha ṇigrahaṇam kimartham . \\
\noindent \textbf{(P\_7,2.26)} vṛttanipātane ṇigrahaṇam aṇyantasya avadhāraṇapratiṣedhārtham . \\
\noindent \textbf{(P\_7,2.26)} vṛttanipātane ṇigrahaṇam kriyate aṇyantasya avadhāraṇam mā bhūt iti . \\
\noindent \textbf{(P\_7,2.26)} kaimarthakyāt niyamaḥ bhavati . \\
\noindent \textbf{(P\_7,2.26)} vidheyam na asti iti kṛtvā . \\
\noindent \textbf{(P\_7,2.26)} iha ca asti vidheyam . \\
\noindent \textbf{(P\_7,2.26)} kim . \\
\noindent \textbf{(P\_7,2.26)} ṇyadhikāt vṛteḥ iṭpratiṣedhaḥ vidheyaḥ . \\
\noindent \textbf{(P\_7,2.26)} tatra apūrvaḥ vidhiḥ astu niyamaḥ astu iti apūrvaḥ eva vidhiḥ bhaviṣyati na niyamaḥ . \\
\noindent \textbf{(P\_7,2.26)} kutaḥ nu khalu etat adhikārthe ārambhe sati ṇyadhikasya bhaviṣyati na punaḥ sanadhikasya vā syāt yaṅadhikasya vā iti . \\
\noindent \textbf{(P\_7,2.26)} tasmāt ṇigrahaṇam kartavyam . \\
\noindent \textbf{(P\_7,2.26)} atha kimartham nipātanam kriyate . \\
\noindent \textbf{(P\_7,2.26)} nipātanam ṇilpeḍguṇapratiṣedhārtham . \\
\noindent \textbf{(P\_7,2.26)} nipātanam kriyate ṇilpārtham iḍguṇapratiṣedhārtham ca \\
\noindent \textbf{(P\_7,2.27)} dāntaśāntayoḥ kim nipātyate . \\
\noindent \textbf{(P\_7,2.27)} dāntaśāntayoḥ upadhādīrghatvam ca . \\
\noindent \textbf{(P\_7,2.27)} kim ca . \\
\noindent \textbf{(P\_7,2.27)} ṇilopeṭpratiṣedhau ca . \\
\noindent \textbf{(P\_7,2.27)} upadhādīrghatvam anipātyam . \\
\noindent \textbf{(P\_7,2.27)} vṛddhyā siddham . \\
\noindent \textbf{(P\_7,2.27)} na sidhyati . \\
\noindent \textbf{(P\_7,2.27)} mitām hrasvaḥ iti hrasvatvena bhavitavyam . \\
\noindent \textbf{(P\_7,2.27)} evam tarhi anunāsikasya kvijhaloḥ kṅiti iti evam atra dīrghatvam bhaviṣyati . \\
\noindent \textbf{(P\_7,2.27)} na sidhyati . \\
\noindent \textbf{(P\_7,2.27)} kim kāraṇam . \\
\noindent \textbf{(P\_7,2.27)} ṇicā vyavahitatvāt . \\
\noindent \textbf{(P\_7,2.27)} ṇilope kṛte na asti vyavadhānam . \\
\noindent \textbf{(P\_7,2.27)} sthānivadbhāvāt vyavadhānam eva . \\
\noindent \textbf{(P\_7,2.27)} pratiṣidhyate atra sthānivadbhāvaḥ dīrghavidhim prati na sthānivat iti . \\
\noindent \textbf{(P\_7,2.27)} atha spaṣṭacchannayoḥ kim nipātyate . \\
\noindent \textbf{(P\_7,2.27)} spaṣṭacchannayoḥ upadhāhrasvatvam ca . \\
\noindent \textbf{(P\_7,2.27)} kim ca . \\
\noindent \textbf{(P\_7,2.27)} ṇilopeṭpratiṣedhau \\
\noindent \textbf{(P\_7,2.28)} ghuṣisvanoḥ vāvacanam iṭpratiṣedhāt vipratiṣedhena . \\
\noindent \textbf{(P\_7,2.28)} ghuṣisvanoḥ vāvacanam iṭpratiṣedhāt bhavati vipratiṣedhena . \\
\noindent \textbf{(P\_7,2.28)} ghuṣeḥ iṭpratiṣedhasya avakāśaḥ asampūrvāt aviśabdanam . \\
\noindent \textbf{(P\_7,2.28)} ghuṣṭā rajjuḥ , ghuṣṭaḥ mārgaḥ . \\
\noindent \textbf{(P\_7,2.28)} vāvacanasya avakāśaḥ sampūrvāt viśabdanam . \\
\noindent \textbf{(P\_7,2.28)} saṅghuṣṭam vākyam , saṅghuṣitam vākyam . \\
\noindent \textbf{(P\_7,2.28)} sampūrvāt aviśabdane ubhayam prāpnoti . \\
\noindent \textbf{(P\_7,2.28)} saṅghuṣṭā rajjuḥ , saṅghuṣitā rajjuḥ . \\
\noindent \textbf{(P\_7,2.28)} vāvacanam bhavati vipratiṣedhena . \\
\noindent \textbf{(P\_7,2.28)} svanaḥ iṭpratiṣedhasya avakāśaḥ anāṅpūrvāt mano'bhidhānam . \\
\noindent \textbf{(P\_7,2.28)} svāntam manaḥ . \\
\noindent \textbf{(P\_7,2.28)} vāvacanasya avakāśaḥ aṅpūrvāt amano'bhidhānam . \\
\noindent \textbf{(P\_7,2.28)} āsvāntaḥ devadattaḥ , āsvanitaḥ devadattaḥ . \\
\noindent \textbf{(P\_7,2.28)} āṅpūrvāt mano'bhidhāne ubhayam prāpnoti . \\
\noindent \textbf{(P\_7,2.28)} āsvāntam manaḥ . \\
\noindent \textbf{(P\_7,2.28)} āsvanitam manaḥ . \\
\noindent \textbf{(P\_7,2.28)} vāvacanam bhavati vipratiṣedhena \\
\noindent \textbf{(P\_7,2.29)} hṛṣeḥ lomakeśakartṛkasya iti vaktavyam . \\
\noindent \textbf{(P\_7,2.29)} hṛṣṭāni lomāni , hṛṣitāni lomāni , hṛṣṭam lomabhiḥ , hṛṣitam lomabhiḥ , hṛṣṭāḥ keśāḥ , hṛṣitāḥ keśāḥ , hṛṣṭam keśaiḥ , hṛṣitam keśaiḥ . \\
\noindent \textbf{(P\_7,2.29)} vismitapratīghātayoḥ iti vaktavyam . \\
\noindent \textbf{(P\_7,2.29)} hṛṣṭaḥ devadattaḥ , hṛṣitaḥ devadattaḥ , hṛṣṭāḥ dantāḥ , hṛṣitāḥ dantāḥ \\
\noindent \textbf{(P\_7,2.30)} apacitaḥ iti kim nipātyate . \\
\noindent \textbf{(P\_7,2.30)} cāyaḥ cibhāvaḥ nipātyate . \\
\noindent \textbf{(P\_7,2.30)} apacitaḥ . \\
\noindent \textbf{(P\_7,2.30)} ktini nityam iti vaktavyam . \\
\noindent \textbf{(P\_7,2.30)} apacitiḥ \\
\noindent \textbf{(P\_7,2.35)} ārdhadhātukagrahaṇam kimartham . \\
\noindent \textbf{(P\_7,2.35)} yathā valādigrahaṇam ārdhadhātukaviśeṣaṇam vijñāyeta . \\
\noindent \textbf{(P\_7,2.35)} valādeḥ ārdhadhātukasya iti . \\
\noindent \textbf{(P\_7,2.35)} atha akriyamāṇe ārdhadhātukagrahaṇe kasya valādigrahaṇam viśeṣaṇam syāt . \\
\noindent \textbf{(P\_7,2.35)} aṅgasya iti vartate . \\
\noindent \textbf{(P\_7,2.35)} aṅgaviśeṣaṇam . \\
\noindent \textbf{(P\_7,2.35)} tatra kaḥ doṣaḥ . \\
\noindent \textbf{(P\_7,2.35)} aṅgasya valādeḥ āditaḥ iṭ prasajyeta āḍāḍvat . \\
\noindent \textbf{(P\_7,2.35)} tat yathā āḍāṭau aṅgasya āditaḥ bhavataḥ tadvat . \\
\noindent \textbf{(P\_7,2.35)} kriyamāṇe api ārdhadhātukagrahaṇe aniṣṭam śakyam vijñātum . \\
\noindent \textbf{(P\_7,2.35)} valādeḥ ārdhadhātukasya yat aṅgam iti . \\
\noindent \textbf{(P\_7,2.35)} akriyamāṇe ca iṣṭam . \\
\noindent \textbf{(P\_7,2.35)} aṅgasya yaḥ valādiḥ iti . \\
\noindent \textbf{(P\_7,2.35)} kim ca aṅgasya valādiḥ . \\
\noindent \textbf{(P\_7,2.35)} nimittam . \\
\noindent \textbf{(P\_7,2.35)} yasmin aṅgam iti etat bhavati . \\
\noindent \textbf{(P\_7,2.35)} kasmin ca etat bhavati . \\
\noindent \textbf{(P\_7,2.35)} pratyaye . \\
\noindent \textbf{(P\_7,2.35)} yāvatā kriyamāṇe ca aniṣṭam vijñāyate akriyamāṇe ca iṣṭam tatra akriyamāṇe eva iṣṭam vijñāsyāmaḥ . \\
\noindent \textbf{(P\_7,2.35)} idam tarhi prayojanam . \\
\noindent \textbf{(P\_7,2.35)} iha mā bhūt . \\
\noindent \textbf{(P\_7,2.35)} āste, śete . \\
\noindent \textbf{(P\_7,2.35)} etat api na asti prayojanam . \\
\noindent \textbf{(P\_7,2.35)} rudādibhyaḥ sārvadhātuke iti etanniyamārtham bhaviṣyati . \\
\noindent \textbf{(P\_7,2.35)} rudādibhyaḥ eva sārvadhātukaḥ iṭ bhavati na anyebhyaḥ iti . \\
\noindent \textbf{(P\_7,2.35)} evam api vṛkṣatvam , vṛkṣatā atra prāpnoti . \\
\noindent \textbf{(P\_7,2.35)} evam tarhi vihitaviśeṣaṇam dhātugrahaṇam . \\
\noindent \textbf{(P\_7,2.35)} dhātoḥ yaḥ vihitaḥ . \\
\noindent \textbf{(P\_7,2.35)} nanu dhātoḥ eva ayam vihitaḥ . \\
\noindent \textbf{(P\_7,2.35)} na ca ayam dhātoḥ iti evam vihitaḥ . \\
\noindent \textbf{(P\_7,2.35)} kva punaḥ dhātugrahaṇam prakṛtam . \\
\noindent \textbf{(P\_7,2.35)} ṛtaḥ it dhātoḥ iti . \\
\noindent \textbf{(P\_7,2.35)} tat vai ṣaṣṭhīnirdiṣṭam pañcamīnirdiṣṭena ca iha vihitaḥ śakyate viśeṣayitum . \\
\noindent \textbf{(P\_7,2.35)} atha idānīm ṣaṣṭhīnirdiṣṭena ca api vihitaḥ śakyate viśeṣayitum śakyam ārdhadhātukagrahaṇam akartum iti \\
\noindent \textbf{(P\_7,2.36)} snukramoḥ anātmanepadanimitte cet kṛti upasaṅkhyānam . \\
\noindent \textbf{(P\_7,2.36)} snukramoḥ anātmanepadanimitte cet kṛti upasaṅkhyānam kartavyam . \\
\noindent \textbf{(P\_7,2.36)} prasnavitā , prasnavitum , prasnavitavyam , prakramitā , prakramitum , prakramitavyam . \\
\noindent \textbf{(P\_7,2.36)} tat tarhi vaktavyam . \\
\noindent \textbf{(P\_7,2.36)} na vaktavyam . \\
\noindent \textbf{(P\_7,2.36)} aviśeṣeṇa snukramoḥ iḍāgamam uktvā ātmanepadapare na iti vakṣyāmi . \\
\noindent \textbf{(P\_7,2.36)} ātmanepadaparapratiṣedhe tatparaparasīyuḍekādeśeṣu pratiṣedhaḥ . \\
\noindent \textbf{(P\_7,2.36)} ātmanepadaparapratiṣedhe ca tatparaparasīyuḍekādeśeṣu pratiṣedhaḥ vaktavyaḥ . \\
\noindent \textbf{(P\_7,2.36)} tatparapare tāvat : prasusnūṣiṣyate , pracikraṃsiṣyate . \\
\noindent \textbf{(P\_7,2.36)} sīyuṭi : prasnoṣīṣta , prakraṃsīṣṭa . \\
\noindent \textbf{(P\_7,2.36)} ekādeśe . \\
\noindent \textbf{(P\_7,2.36)} prasnoṣyante , prakraṃsyante . \\
\noindent \textbf{(P\_7,2.36)} ekādeśe kṛte vyapavargābhāvāt na prāpnoti . \\
\noindent \textbf{(P\_7,2.36)} antādivadbhāvena vyapavargaḥ . \\
\noindent \textbf{(P\_7,2.36)} ubhayataḥ āśraye na antādivat . \\
\noindent \textbf{(P\_7,2.36)} evam tarhi ekādeśaḥ pūrvavidhim prati sthānivat bhavati iti sthānivadbhāvāt vyapavargaḥ . \\
\noindent \textbf{(P\_7,2.36)} tatparaparasīyuṭoḥ tarhi pratiṣedhaḥ vaktavyaḥ . \\
\noindent \textbf{(P\_7,2.36)} siddham tu snoḥ ātmanepadena samānapadasthasya iṭpratiṣedhāt . \\
\noindent \textbf{(P\_7,2.36)} siddham etat . \\
\noindent \textbf{(P\_7,2.36)} katham . \\
\noindent \textbf{(P\_7,2.36)} snoḥ ātmanepadena samānapadasthasya na iṭ bhavati iti vaktavyam . \\
\noindent \textbf{(P\_7,2.36)} yadi snoḥ ātmanepadena samānapadasthasya iṭ na bhavati iti ucyate prasnavitā iva ācarati prasnavitrīyate atra na prāpnoti . \\
\noindent \textbf{(P\_7,2.36)} bahiraṅgalakṣaṇam atra ātmanepadam . \\
\noindent \textbf{(P\_7,2.36)} kramoḥ ca . \\
\noindent \textbf{(P\_7,2.36)} kramoḥ ca ātmanepadena samānapadasthasya iṭ na bhavati iti vaktavyam . \\
\noindent \textbf{(P\_7,2.36)} atha kimartham krameḥ pṛthaggrahaṇam kriyate na snukramibhyām iti eva ucyeta . \\
\noindent \textbf{(P\_7,2.36)} kartari ca ātmanepadaviṣayāt kṛti na iti vakṣyatgi tat krameḥ eva syāt snoḥ mā bhūt . \\
\noindent \textbf{(P\_7,2.36)} vyatiprasnavitārau , vyatiprasnativāraḥ . \\
\noindent \textbf{(P\_7,2.36)} kartari ca ātmanepadaviṣayāt kṛti . \\
\noindent \textbf{(P\_7,2.36)} kartari ca ātmanepadaviṣayāt kṛti pratiṣedhaḥ vaktavyaḥ . \\
\noindent \textbf{(P\_7,2.36)} prakrantā , upakrantā . \\
\noindent \textbf{(P\_7,2.36)} tat tarhi idam bahu vaktavyam . \\
\noindent \textbf{(P\_7,2.36)} snoḥ ātmanepadena samānapadasthasya iṭ na bhavati iti vaktavyam . \\
\noindent \textbf{(P\_7,2.36)} krameḥ ca iti vaktavyam . \\
\noindent \textbf{(P\_7,2.36)} kartari ca ātmanepadaviṣayāt kṛti iti vaktavyam . \\
\noindent \textbf{(P\_7,2.36)} sūtram ca bhidyate . \\
\noindent \textbf{(P\_7,2.36)} yathānyāsam eva astu . \\
\noindent \textbf{(P\_7,2.36)} nanu ca uktam snukramoḥ anātmanepadanimitte cet kṛti upasaṅkhyānam iti . \\
\noindent \textbf{(P\_7,2.36)} na eṣaḥ doṣaḥ . \\
\noindent \textbf{(P\_7,2.36)} snukramī eva ātmanepadanimittatvena viśeṣayiṣyāmaḥ . \\
\noindent \textbf{(P\_7,2.36)} na cet snukramī ātmanepadasya nimitte iti . \\
\noindent \textbf{(P\_7,2.36)} katham punaḥ dhātuḥ nāma ātmanepadasya nimittam syāt . \\
\noindent \textbf{(P\_7,2.36)} dhātuḥ eva nimittam . \\
\noindent \textbf{(P\_7,2.36)} āha hi bhagavān anudāttaṅitaḥ ātmanepadam śeṣāt kartari parasmaipadam iti . \\
\noindent \textbf{(P\_7,2.36)} yatra tarhi dhātuḥ na āśrīyate bhāvakarmaṇoḥ iti . \\
\noindent \textbf{(P\_7,2.36)} atra api dhātuḥ eva āśrīyate . \\
\noindent \textbf{(P\_7,2.36)} bhāvakarmavṛttāt dhātoḥ iti . \\
\noindent \textbf{(P\_7,2.36)} katham prakramitavyam . \\
\noindent \textbf{(P\_7,2.36)} sati ātmanepade nimittaśabdaḥ vartate . \\
\noindent \textbf{(P\_7,2.36)} katham prakrantā , upakrantā . \\
\noindent \textbf{(P\_7,2.36)} tasmāt asati api . \\
\noindent \textbf{(P\_7,2.36)} katham prakramitavyam . \\
\noindent \textbf{(P\_7,2.36)} tasmāt sati eva . \\
\noindent \textbf{(P\_7,2.36)} katham prakrantā , upakrantā . \\
\noindent \textbf{(P\_7,2.36)} vaktavyam eva etat kartari ca ātmanepadaviṣayāt kṛti iti . \\
\noindent \textbf{(P\_7,2.36)} atha vā kṛti iti vartate \\
\noindent \textbf{(P\_7,2.37)} graheḥ dīrghatve iḍgrahaṇam . \\
\noindent \textbf{(P\_7,2.37)} graheḥ dīrghatve iḍgrahaṇam kartavyam . \\
\noindent \textbf{(P\_7,2.37)} iṭaḥ dīrghaḥ iti vaktavyam . \\
\noindent \textbf{(P\_7,2.37)} na vaktavyam . \\
\noindent \textbf{(P\_7,2.37)} prakṛtam anuvartate . \\
\noindent \textbf{(P\_7,2.37)} kva prakṛtam . \\
\noindent \textbf{(P\_7,2.37)} ārdhadhātukasya iṭ valādeḥ iti . \\
\noindent \textbf{(P\_7,2.37)} evam api kartavyam eva . \\
\noindent \textbf{(P\_7,2.37)} agrahaṇe hi asampratyayaḥ ṣaṣṭhyabhāvāt . \\
\noindent \textbf{(P\_7,2.37)} akriyamāṇe hi iḍgrahaṇe asampratyayaḥ syāt . \\
\noindent \textbf{(P\_7,2.37)} kim kāraṇam . \\
\noindent \textbf{(P\_7,2.37)} ṣaṣṭhyabhāvāt . \\
\noindent \textbf{(P\_7,2.37)} ṣaṣṭhīnirdiṣṭasya ādeśāḥ ucyante na ca atra ṣaṣṭhīm paśyāmaḥ . \\
\noindent \textbf{(P\_7,2.37)} kriyamāṇe ca api iḍgrahaṇe . \\
\noindent \textbf{(P\_7,2.37)} ciṇvadiṭaḥ pratiṣedhaḥ . \\
\noindent \textbf{(P\_7,2.37)} ciṅvaditaḥ pratiṣedhaḥ vaktavyaḥ . \\
\noindent \textbf{(P\_7,2.37)} grāhiṣyate . \\
\noindent \textbf{(P\_7,2.37)} yaṅlope ca . \\
\noindent \textbf{(P\_7,2.37)} yaṅlope ca pratiṣedhaḥ vaktavyaḥ . \\
\noindent \textbf{(P\_7,2.37)} jarīgṛhitā , jarīgṛhitum , jarīgṛhitavyam . \\
\noindent \textbf{(P\_7,2.37)} yadi punaḥ iṭ dīrghaḥ āgamāntaram vijñāyeta . \\
\noindent \textbf{(P\_7,2.37)} iṭ dīrghaḥ iti cet vipratiṣiddham . \\
\noindent \textbf{(P\_7,2.37)} iṭ dīrghaḥ iti cet vipratiṣiddham bhavati . \\
\noindent \textbf{(P\_7,2.37)} yadi iṭ na dīrghaḥ . \\
\noindent \textbf{(P\_7,2.37)} atha dīrghaḥ na iṭ . \\
\noindent \textbf{(P\_7,2.37)} iṭ dīrghaḥ ca iti vipratiṣiddham . \\
\noindent \textbf{(P\_7,2.37)} pratiṣiddhasya ca punarvidhāne dīrghatvābhāvaḥ . \\
\noindent \textbf{(P\_7,2.37)} pratiṣiddhasya ca punarvidhāne dīrghatvasya abhāvaḥ . \\
\noindent \textbf{(P\_7,2.37)} vuvūrṣate , vivariṣate , vavarīṣate . \\
\noindent \textbf{(P\_7,2.37)} atra api iṭ dīrghaḥ iti anuvartiṣyate . \\
\noindent \textbf{(P\_7,2.37)} yat tarhi videśastham pratiṣidhya punarvidhānam tat na sidhyati . \\
\noindent \textbf{(P\_7,2.37)} jṝvraścyoḥ ktvi . \\
\noindent \textbf{(P\_7,2.37)} śryukaḥ kiti iti anena pratiṣiddhe dīrghatvam na prāpnoti . \\
\noindent \textbf{(P\_7,2.37)} jaritvā , jarītvā . \\
\noindent \textbf{(P\_7,2.37)} īṭaḥ vidhiḥ iṭaḥ pratiṣedhaḥ . \\
\noindent \textbf{(P\_7,2.37)} yathāprāptaḥ iṭ dīrghaḥ bhaviṣyati . \\
\noindent \textbf{(P\_7,2.37)} yadi tarhi iṭaḥ grahaṇe īṭaḥ grahaṇam na bhavati jarītvā na ktvā seṭ iti kittvapratiṣedhaḥ na prāpnoti . \\
\noindent \textbf{(P\_7,2.37)} iha ca agrhīt iti iṭaḥ īṭi iti sijlopaḥ na prāpnoti . \\
\noindent \textbf{(P\_7,2.37)} iha ca agrahīt na iṭi iti vṛddhipratiṣedhaḥ na prāpnoti . \\
\noindent \textbf{(P\_7,2.37)} mā bhūt evam . \\
\noindent \textbf{(P\_7,2.37)} hmyantānām iti evam bhaviṣyati . \\
\noindent \textbf{(P\_7,2.37)} atra api na iṭ iti eva anuvartate . \\
\noindent \textbf{(P\_7,2.37)} tat ca avaśyam iḍgrahaṇam anuvartyam adhākṣīt iti evamartham . \\
\noindent \textbf{(P\_7,2.37)} tathā agrahīdhvam , agrahīḍhvam vibhāṣā iṭaḥ iti mūrdhanyaḥ na prāpnoti . \\
\noindent \textbf{(P\_7,2.37)} tasmāt na evam śakyam vaktum iṭaḥ grahaṇe īṭaḥ grahaṇam na bhavati iti . \\
\noindent \textbf{(P\_7,2.37)} bhavati cet pratiṣiddhasya ca punarvidhāne dīrghābhāvaḥ iti eva . \\
\noindent \textbf{(P\_7,2.37)} tasmāt aśakyaḥ iṭ dīrghaḥ āgamāntaram vijñātum . \\
\noindent \textbf{(P\_7,2.37)} na cet vijñāyate iṭaḥ grahaṇam kartavyam . \\
\noindent \textbf{(P\_7,2.37)} na kartavyam . \\
\noindent \textbf{(P\_7,2.37)} ārdhadhātukasya iti vartate . \\
\noindent \textbf{(P\_7,2.37)} grahaḥ parasya ārdhadhātukasya dīrghatvam vakṣyāmi . \\
\noindent \textbf{(P\_7,2.37)} iha api tarhi prāpnoti . \\
\noindent \textbf{(P\_7,2.37)} grahaṇam , grahaṇīyam . \\
\noindent \textbf{(P\_7,2.37)} valādeḥ iti vartate . \\
\noindent \textbf{(P\_7,2.37)} evam api grahītā , grahītum atra na prāpnoti . \\
\noindent \textbf{(P\_7,2.37)} bhūtapūrvagatyā bhaviṣyati . \\
\noindent \textbf{(P\_7,2.37)} evam api grāhakaḥ atra prāpnoti . \\
\noindent \textbf{(P\_7,2.37)} kim ca iṭpratīghātena khalu api dīrghatvam ucyamānam iṭam bādhate . \\
\noindent \textbf{(P\_7,2.37)} tasmāt iṭaḥ grahaṇam kartavyam . \\
\noindent \textbf{(P\_7,2.37)} na kartavyam . \\
\noindent \textbf{(P\_7,2.37)} prakṛtam anuvartate . \\
\noindent \textbf{(P\_7,2.37)} kva prakṛtam . \\
\noindent \textbf{(P\_7,2.37)} ārdhadhātukasya iṭ valādeḥ iti . \\
\noindent \textbf{(P\_7,2.37)} nanu ca uktam evam api kartavyam eva agrahaṇe hi asampratyayaḥ ṣaṣṭhyanirdeśāt iti . \\
\noindent \textbf{(P\_7,2.37)} na eṣaḥ doṣaḥ . \\
\noindent \textbf{(P\_7,2.37)} grhaḥ iti eṣā pañcamī iṭ iti prathamāyāḥ ṣaṣṭhīm prakalpayiṣyati tasmāt iti uttarasya iti . \\
\noindent \textbf{(P\_7,2.37)} evam ca kṛtvā saḥ api adoṣaḥ bhavati yat uktam ciṇvadiṭaḥ pratiṣedhaḥ iti . \\
\noindent \textbf{(P\_7,2.37)} katham . \\
\noindent \textbf{(P\_7,2.37)} prakṛtasya iṭaḥ idam dīrghatvam na ca ciṇvadiṭ prakṛtaḥ . \\
\noindent \textbf{(P\_7,2.37)} yaṅlope katham . \\
\noindent \textbf{(P\_7,2.37)} yaṅlope ca uktam iṭi sarvatra . \\
\noindent \textbf{(P\_7,2.37)} kva sarvatra . \\
\noindent \textbf{(P\_7,2.37)} yadi eva prakṛtasya iṭaḥ dīrghatvam atha api iṭ dīrghaḥ āgamāntaram vijñāyeta . \\
\noindent \textbf{(P\_7,2.37)} yaṅlope ca uktam . \\
\noindent \textbf{(P\_7,2.37)} kim uktam . \\
\noindent \textbf{(P\_7,2.37)} tadantadvirvacanāt iti \\
\noindent \textbf{(P\_7,2.44)} atha vā iti vartamāne punaḥ vāvacanam kimartham . \\
\noindent \textbf{(P\_7,2.44)} punaḥ vāvacanam kriyate liṅsicoḥ nivṛttyartham . \\
\noindent \textbf{(P\_7,2.44)} punaḥ vāvacanam kriyate liṅsicoḥ nivṛttyartham . \\
\noindent \textbf{(P\_7,2.44)} atha kimartham sūtisūyatyoḥ pṛthaggrahaṇam kriyate . \\
\noindent \textbf{(P\_7,2.44)} suvateḥ mā bhūt . \\
\noindent \textbf{(P\_7,2.44)} atha kimartham dhūñaḥ sānubandhakasya grahaṇam kriyate . \\
\noindent \textbf{(P\_7,2.44)} dhuvateḥ mā bhūt iti . \\
\noindent \textbf{(P\_7,2.44)} kim punaḥ iyam prāpte vibhāṣā āhosvit aprāpte . \\
\noindent \textbf{(P\_7,2.44)} katham ca prāpte katham vā aprāpte . \\
\noindent \textbf{(P\_7,2.44)} yadi svaratiḥ udāttaḥ tataḥ prāpte . \\
\noindent \textbf{(P\_7,2.44)} atha anudāttaḥ tataḥ aprāpte . \\
\noindent \textbf{(P\_7,2.44)} svaratiḥ udāttaḥ . \\
\noindent \textbf{(P\_7,2.44)} svaratiḥ udāttaḥ paṭhyate . \\
\noindent \textbf{(P\_7,2.44)} kimartham tarhi vāvacanam . \\
\noindent \textbf{(P\_7,2.44)} vāvacanam nivṛttyartham . \\
\noindent \textbf{(P\_7,2.44)} vāvacanam kriyate nivṛttyartham . \\
\noindent \textbf{(P\_7,2.44)} anudātte hi kiti vāprasaṅgaḥ pratiṣidhya punaḥ vidhānāt . \\
\noindent \textbf{(P\_7,2.44)} anudātte hi sati kiti vibhāṣā prasajyeta . \\
\noindent \textbf{(P\_7,2.44)} svṛtvā . \\
\noindent \textbf{(P\_7,2.44)} pratiṣidhya punaḥ vidhānāt . \\
\noindent \textbf{(P\_7,2.44)} pratiṣidhya kila ayam punaḥ vidhīyate . \\
\noindent \textbf{(P\_7,2.44)} saḥ yathā eva ekājlakṣaṇam pratiṣedham bādhate evam śryukaḥ kiti iti etam api bādheta . \\
\noindent \textbf{(P\_7,2.44)} yadi tarhi udāttaḥ svaratiḥ paṭhiṣyati vipratiṣedham svarateḥ veṭtvāt ṛtaḥ sye vipratiṣedhena iti saḥ vipratiṣedhaḥ na upapadyate . \\
\noindent \textbf{(P\_7,2.44)} kim kāraṇam . \\
\noindent \textbf{(P\_7,2.44)} saḥ vidhiḥ ayam pratiṣedhaḥ vidhipratiṣedhayoḥ ca ayuktaḥ vipratiṣedhaḥ . \\
\noindent \textbf{(P\_7,2.44)} saḥ api vidhiḥ na mṛdūnām iva kārpāsānām kṛtaḥ pratiṣedhaviṣaye ārabhyate . \\
\noindent \textbf{(P\_7,2.44)} saḥ yathā eva ekājlakṣaṇam pratiṣedham bādhate evam imam api bādhiṣyate . \\
\noindent \textbf{(P\_7,2.44)} atha vā yena na aprāpte tasya bādhanam bhavati na ca aprāpte valādilakṣaṇe iyam vibhāṣā ārabhyate syalakṣaṇe punaḥ prāpte ca aprāpte ca . \\
\noindent \textbf{(P\_7,2.44)} atha vā madhye apavādāḥ pūrvān vidhīn bādhante iti evam iyam vibhāṣā valādilakṣaṇam iṭam bādhiṣyate syalakṣaṇam na bādhiṣyate . \\
\noindent \textbf{(P\_7,2.44)} atha vā punaḥ astu anudāttaḥ . \\
\noindent \textbf{(P\_7,2.44)} nanu ca uktam anudātte hi kiti vāprasaṅgaḥ pratiṣidhya punaḥ vidhānāt iti . \\
\noindent \textbf{(P\_7,2.44)} na eṣaḥ doṣaḥ . \\
\noindent \textbf{(P\_7,2.44)} yena na aprāpte tasya bādhanam bhavati na ca aprāpte ekājlakṣaṇe pratiṣedhe iyam vibhāṣā ārabhyate śryukaḥ kiti iti etasmin punaḥ prāpte ca aprāpte ca . \\
\noindent \textbf{(P\_7,2.44)} atha vā śryukaḥ kiti iti eṣaḥ yogaḥ udāttārthaḥ ca yebhyaḥ ca anudāttebhyaḥ iṭ prāpyate tadbādhanārthaḥ ca . \\
\noindent \textbf{(P\_7,2.44)} atha vā śryukaḥ kiti iti iha anuvartiṣyate . \\
\noindent \textbf{(P\_7,2.44)} atha vā ācāryapravṛttiḥ jñāpayati na iyam vibhāṣā uglakṣaṇasya pratiṣedhasya viṣaye bhavati iti yat ayam sanīvantardhabhrasjadambhuśrisvṛyūrṇubharajñapisanām iti svṛgrahaṇam karoti \\
\noindent \textbf{(P\_7,2.47)} iṭ iti vartamāne punaḥ iḍgrahaṇam kimartham . \\
\noindent \textbf{(P\_7,2.47)} iḍgrahaṇam nityārtham . \\
\noindent \textbf{(P\_7,2.47)} nityaḥ ayam ārambhaḥ . \\
\noindent \textbf{(P\_7,2.47)} na etat asti prayojanam . \\
\noindent \textbf{(P\_7,2.47)} siddhā atra vibhaṣā pūrveṇa eva tatra ārambhasāmarthyāt nityaḥ vidhiḥ bhaviṣyati . \\
\noindent \textbf{(P\_7,2.47)} na atra pūrveṇa vibhāṣā prāpnoti . \\
\noindent \textbf{(P\_7,2.47)} kim kāraṇam . \\
\noindent \textbf{(P\_7,2.47)} yasya vibhāṣā iti pratiṣedhāt . \\
\noindent \textbf{(P\_7,2.47)} tatra ārambhasāmarthyāt vibhaṣā labhyeta punaḥ iḍgrahaṇāt iṭ eva bhavati \\
\noindent \textbf{(P\_7,2.48)} iṣeḥ takāre śyanpratyayāt pratiṣedhaḥ . \\
\noindent \textbf{(P\_7,2.48)} iṣeḥ takāre śyanpratyayāt pratiṣedhaḥ vaktavyaḥ . \\
\noindent \textbf{(P\_7,2.48)} iha mā bhūt . \\
\noindent \textbf{(P\_7,2.48)} preṣitā , preṣitum , preṣitavyam \\
\noindent \textbf{(P\_7,2.52)} iṭ iti vartamāne punaḥ iḍgrahaṇam kimartham . \\
\noindent \textbf{(P\_7,2.52)} punaḥ iḍgrahaṇam nityārtham . \\
\noindent \textbf{(P\_7,2.52)} iṭ iti vartamāne punaḥ iḍgrahaṇam kriyate nityārtham . \\
\noindent \textbf{(P\_7,2.52)} nityārthaḥ ayam ārambhaḥ \\
\noindent \textbf{(P\_7,2.58)} gameḥ iṭ parasmaipadeṣu cet kṛti upasaṅkhyānam . \\
\noindent \textbf{(P\_7,2.58)} gameḥ iṭ parasmaipadeṣu cet kṛti upasaṅkhyānam kartavyam . \\
\noindent \textbf{(P\_7,2.58)} jigamiṣitā , jigamiṣitum , jigamiṣitavyam . \\
\noindent \textbf{(P\_7,2.58)} tat tarhi upasaṅkhyānam kartavyam . \\
\noindent \textbf{(P\_7,2.58)} na kartavyam . \\
\noindent \textbf{(P\_7,2.58)} aviśeṣeṇa gameḥ iḍāgamam uktvā ātmanepadapare na iti vakṣyāmi . \\
\noindent \textbf{(P\_7,2.58)} ātmanepadaparapratiṣedhe uktam . \\
\noindent \textbf{(P\_7,2.58)} kim uktam . \\
\noindent \textbf{(P\_7,2.58)} ātmanepadaparapratiṣedhe tatparaparasīyuḍekādeśeṣu pratiṣedhaḥ iti . \\
\noindent \textbf{(P\_7,2.58)} iha api ātmanepadaparapratiṣedhe tatparaparasīyuḍekādeśeṣu pratiṣedhaḥ vaktavyaḥ . \\
\noindent \textbf{(P\_7,2.58)} taparapare tāvat . \\
\noindent \textbf{(P\_7,2.58)} saṅjigaṃsiṣyate . \\
\noindent \textbf{(P\_7,2.58)} sīyuṭi . \\
\noindent \textbf{(P\_7,2.58)} saṅgaṃsīṣṭa . \\
\noindent \textbf{(P\_7,2.58)} ekādeśe . \\
\noindent \textbf{(P\_7,2.58)} saṅgaṃsyante . \\
\noindent \textbf{(P\_7,2.58)} ekādeśe kṛte vyapavargābhāvāt na prāpnoti . \\
\noindent \textbf{(P\_7,2.58)} siddham tu gameḥ ātmanepadena samānapadasthasya iṭpratiṣedhāt . \\
\noindent \textbf{(P\_7,2.58)} siddham etat . \\
\noindent \textbf{(P\_7,2.58)} katham . \\
\noindent \textbf{(P\_7,2.58)} gameḥ ātmanepadena samānapadastheṇa na bhavati iti vaktavyam \\
\noindent \textbf{(P\_7,2.59)} vṛtādipratiṣedhe ca . \\
\noindent \textbf{(P\_7,2.59)} kim . \\
\noindent \textbf{(P\_7,2.59)} kṛti upasaṅkhyānam kartavyam : vivṛtsitā vivṛtsitum , vivṛtsitavyam . \\
\noindent \textbf{(P\_7,2.59)} tat tarhi upasaṅkhyānam kartavyam . \\
\noindent \textbf{(P\_7,2.59)} na kartavyam . \\
\noindent \textbf{(P\_7,2.59)} aviśeṣeṇa vṛtādibhyaḥ iṭpratiṣedham uktvā ātmanepadaparaḥ iṭ bhavati iti vakṣyāmi . \\
\noindent \textbf{(P\_7,2.59)} ātmanepadapare iḍvacane tatparaparasīyuḍekādeśeṣu iḍvacanam . \\
\noindent \textbf{(P\_7,2.59)} ātmanepadapare iḍvacane tatparaparasīyuḍekādeśeṣu iṭ vaktavyaḥ . \\
\noindent \textbf{(P\_7,2.59)} taparapare tāvat . \\
\noindent \textbf{(P\_7,2.59)} vivartiṣiṣyate . \\
\noindent \textbf{(P\_7,2.59)} sīyuṭi . \\
\noindent \textbf{(P\_7,2.59)} vartiṣīṣṭa . \\
\noindent \textbf{(P\_7,2.59)} ekādeśe . \\
\noindent \textbf{(P\_7,2.59)} vartiṣyante , vardhiṣyante . \\
\noindent \textbf{(P\_7,2.59)} siddham tu vṛtādīnām ātmanepadena samānapadasthasya iḍvacanāt . \\
\noindent \textbf{(P\_7,2.59)} siddham etat . \\
\noindent \textbf{(P\_7,2.59)} katham . \\
\noindent \textbf{(P\_7,2.59)} vṛtādīnām ātmanepadena samānapdasthasya iṭ bhavati iti vaktavyam . \\
\noindent \textbf{(P\_7,2.59)} catustāsikḷpigrahaṇānarthakyam ca . \\
\noindent \textbf{(P\_7,2.59)} caturgrahaṇam ca anarthakam . \\
\noindent \textbf{(P\_7,2.59)} sarvebhyaḥ hi vṛtādibhyaḥ pratiṣedhaḥ iṣyate . \\
\noindent \textbf{(P\_7,2.59)} tāsigrahaṇam ca anarthakam . \\
\noindent \textbf{(P\_7,2.59)} kim kāraṇam . \\
\noindent \textbf{(P\_7,2.59)} nivṛttatvāt sakārasya . \\
\noindent \textbf{(P\_7,2.59)} nivṛttam sakārādau iti . \\
\noindent \textbf{(P\_7,2.59)} tāsgrahaṇe ca idānīm akriyamāṇe kḷpigrahaṇena api na arthaḥ eṣaḥ api hi vṛtādiḥ pañcamaḥ . \\
\noindent \textbf{(P\_7,2.59)} bhavet kḷpigrahaṇam na kartavyam tāsgrahaṇam tu kartavyam . \\
\noindent \textbf{(P\_7,2.59)} yat hi tat sakārādau iti na tat śakyam nivartayitum tṛci api hi prasajyeta . \\
\noindent \textbf{(P\_7,2.59)} vartitā , vardhitā . \\
\noindent \textbf{(P\_7,2.59)} tāsgrahaṇe ca idānīm kriyamāṇe kḷpigrahaṇam api kartavyam anyebhyaḥ api vṛtādibhyaḥ tāsau mā bhūt iti . \\
\noindent \textbf{(P\_7,2.59)} bhavet tāsgrahaṇam kartavyam kḷpigrahaṇam tu na eva kartavyam . \\
\noindent \textbf{(P\_7,2.59)} anyebhyaḥ api vṛtādibhyaḥ tāsau kasmāt na bhavati . \\
\noindent \textbf{(P\_7,2.59)} parasmaipadeṣu iti vartate kḷpeḥ eva ca tāsparasmaipadaparaḥ na anyebhyaḥ vṛtādibhyaḥ . \\
\noindent \textbf{(P\_7,2.59)} yadi evam tāsgrahaṇena api na arthaḥ . \\
\noindent \textbf{(P\_7,2.59)} tṛci kasmāt na bhavati . \\
\noindent \textbf{(P\_7,2.59)} parasmaipadeṣu iti vartate \\
\noindent \textbf{(P\_7,2.62)} tāsau atvatpratiṣedhe ghaseḥ pratiṣedhaprasaṅgaḥ akāravattvāt . \\
\noindent \textbf{(P\_7,2.62)} tāsau atvatpratiṣedhe ghaseḥ pratiṣedhaḥ prāpnoti . \\
\noindent \textbf{(P\_7,2.62)} jaghasitha . \\
\noindent \textbf{(P\_7,2.62)} kim kāraṇam . \\
\noindent \textbf{(P\_7,2.62)} akāravattvāt . \\
\noindent \textbf{(P\_7,2.62)} saḥ api hi akāravān . \\
\noindent \textbf{(P\_7,2.62)} siddham tu halādigrahaṇāt . \\
\noindent \textbf{(P\_7,2.62)} siddham etat . \\
\noindent \textbf{(P\_7,2.62)} katham . \\
\noindent \textbf{(P\_7,2.62)} halādigrahaṇam kartavyam . \\
\noindent \textbf{(P\_7,2.62)} tat ca avaśyam kartavyam . \\
\noindent \textbf{(P\_7,2.62)} aṭhyaśī prayojayataḥ . \\
\noindent \textbf{(P\_7,2.62)} āthyaśī tāvat na prayojayataḥ . \\
\noindent \textbf{(P\_7,2.62)} kim kāraṇam . \\
\noindent \textbf{(P\_7,2.62)} tāsau aniṭaḥ iti ucyate seṭau ca imau tāsau . \\
\noindent \textbf{(P\_7,2.62)} añjvaśū tarhi prayojayataḥ . \\
\noindent \textbf{(P\_7,2.62)} añjvaśū ca api na prayojayataḥ . \\
\noindent \textbf{(P\_7,2.62)} kim kāraṇam . \\
\noindent \textbf{(P\_7,2.62)} tāsau nityāniṭaḥ iti ucyate vibhāṣiteṭau ca etau . \\
\noindent \textbf{(P\_7,2.62)} adiḥ tarhi prayojayati . \\
\noindent \textbf{(P\_7,2.62)} āditha . \\
\noindent \textbf{(P\_7,2.62)} kriyamāṇe api vai halādigrahaṇe atra prāpnoti . \\
\noindent \textbf{(P\_7,2.62)} jaghasitha . \\
\noindent \textbf{(P\_7,2.62)} eṣaḥ api halādiḥ . \\
\noindent \textbf{(P\_7,2.62)} tasya ca abhāvāt tāsau . \\
\noindent \textbf{(P\_7,2.62)} tāsau aniṭaḥ iti ucyate na ca ghasiḥ tāsau asti . \\
\noindent \textbf{(P\_7,2.62)} nanu ca yaḥ tāsau na asti aniṭ api asau tāsau bhavati . \\
\noindent \textbf{(P\_7,2.62)} na evam vijñāyate yaḥ tāsau aniṭ iti . \\
\noindent \textbf{(P\_7,2.62)} katham tarhi . \\
\noindent \textbf{(P\_7,2.62)} yaḥ tāsau asti aniṭ ca iti . \\
\noindent \textbf{(P\_7,2.62)} kim vaktavyam etat . \\
\noindent \textbf{(P\_7,2.62)} na hi . \\
\noindent \textbf{(P\_7,2.62)} katham anucyamānam gaṃsyate . \\
\noindent \textbf{(P\_7,2.62)} saptamyarthe api vai vatiḥ bhavati . \\
\noindent \textbf{(P\_7,2.62)} tat yathā . \\
\noindent \textbf{(P\_7,2.62)} mathurāyām iva mathurāvat . \\
\noindent \textbf{(P\_7,2.62)} pāṭaliputre iva pāṭaliputravat . \\
\noindent \textbf{(P\_7,2.62)} evam tāsau iva tāsvat \\
\noindent \textbf{(P\_7,2.63)} kimartham idam ucyate na acaḥ tāsvat thali aniṭaḥ nityam iti eva siddham . \\
\noindent \textbf{(P\_7,2.63)} evam tarhi niyamārthaḥ ayam ārambhaḥ . \\
\noindent \textbf{(P\_7,2.63)} ṛtaḥ eva bhāradvājasya na anyataḥ bhāradvājasya iti . \\
\noindent \textbf{(P\_7,2.63)} kva mā bhūt . \\
\noindent \textbf{(P\_7,2.63)} yayitha , vavitha iti . \\
\noindent \textbf{(P\_7,2.63)} ṛtaḥ bhāradvājasya iti niyamānupapattiḥ aprāptatvāt pratiṣedhasya . \\
\noindent \textbf{(P\_7,2.63)} ṛtaḥ bhāradvājasya iti niyamānupapattiḥ . \\
\noindent \textbf{(P\_7,2.63)} kim kāraṇam . \\
\noindent \textbf{(P\_7,2.63)} aprāptatvāt pratiṣedhasya . \\
\noindent \textbf{(P\_7,2.63)} guṇe kṛte raparatve ca anajantatvāt pratiṣedhaḥ na prāpnoti . \\
\noindent \textbf{(P\_7,2.63)} asati niyame kaḥ doṣaḥ . \\
\noindent \textbf{(P\_7,2.63)} tatra pacādibhyaḥ iḍvacanam . \\
\noindent \textbf{(P\_7,2.63)} tatra pacādibhyaḥ iṭ vaktavyaḥ . \\
\noindent \textbf{(P\_7,2.63)} pecitha , śekitha iti . \\
\noindent \textbf{(P\_7,2.63)} yadiḥ punaḥ ayam bhāradvājaḥ purastāt apakṛṣyeta . \\
\noindent \textbf{(P\_7,2.63)} acaḥ tāsvat thali aniṭaḥ nityam bhāradvājasya . \\
\noindent \textbf{(P\_7,2.63)} upadeśe atvataḥ bhāradvājasya . \\
\noindent \textbf{(P\_7,2.63)} tataḥ ṛtaḥ . \\
\noindent \textbf{(P\_7,2.63)} bhāradvājasya iti nivṛttam . \\
\noindent \textbf{(P\_7,2.63)} sidhyate evam ayam tu bhāradvājaḥ svasmāt matāt pracyāvitaḥ bhavati . \\
\noindent \textbf{(P\_7,2.63)} evam tarhi yogavibhāgāt siddham . \\
\noindent \textbf{(P\_7,2.63)} yogavibhāgaḥ kariṣyate . \\
\noindent \textbf{(P\_7,2.63)} acaḥ tāsvat thali aniṭaḥ nityam upadeśe . \\
\noindent \textbf{(P\_7,2.63)} tataḥ atvataḥ . \\
\noindent \textbf{(P\_7,2.63)} atvataḥ ca upadeśe iti \\
\noindent \textbf{(P\_7,2.64)} vṛgrahaṇam kimartham na kṛsṛbhṛvṛstudruśrusruvaḥ liṭi iti eva siddham . \\
\noindent \textbf{(P\_7,2.64)} evam tarhi niyamārthaḥ ayam ārambhaḥ . \\
\noindent \textbf{(P\_7,2.64)} nigamaḥ eva yathā syāt . \\
\noindent \textbf{(P\_7,2.64)} kva mā bhūt . \\
\noindent \textbf{(P\_7,2.64)} vavaritha \\
\noindent \textbf{(P\_7,2.67.1)} kimartham idam ucyate . \\
\noindent \textbf{(P\_7,2.67.1)} vasvekājādghasāṃvacanam niyamārtham . \\
\noindent \textbf{(P\_7,2.67.1)} niyamārthaḥ ayam ārambhaḥ . \\
\noindent \textbf{(P\_7,2.67.1)} vasau ekājāt ghasām eva . \\
\noindent \textbf{(P\_7,2.67.1)} kva mā bhūt . \\
\noindent \textbf{(P\_7,2.67.1)} bibhidvān . \\
\noindent \textbf{(P\_7,2.67.1)} kim ucyate niyamārtham iti na punaḥ vidhyarthaḥ api syāt . \\
\noindent \textbf{(P\_7,2.67.1)} pratiṣedhaḥ api hi atra prāpnoti na iṭ vaśi kṛti iti . \\
\noindent \textbf{(P\_7,2.67.1)} kṛt ca eva hi ayam vaśādiḥ ca . \\
\noindent \textbf{(P\_7,2.67.1)} evam tarhi kṛsṛbhṛvṛstudruśrusruvaḥ liṭi iti etasmāt niyamāt atra iṭ bhaviṣyati . \\
\noindent \textbf{(P\_7,2.67.1)} na atra tena pariprāpaṇam prāpnoti . \\
\noindent \textbf{(P\_7,2.67.1)} kim kāraṇam . \\
\noindent \textbf{(P\_7,2.67.1)} prakṛtilakṣaṇasya pratiṣedhasya saḥ pratyārambhaḥ pratyayalakṣaṇaḥ ca ayam pratiṣedhaḥ . \\
\noindent \textbf{(P\_7,2.67.1)} ubhayoḥ saḥ pratyārambhaḥ . \\
\noindent \textbf{(P\_7,2.67.1)} katham jñāyate . \\
\noindent \textbf{(P\_7,2.67.1)} vṛṅvṛñoḥ grahaṇāt . \\
\noindent \textbf{(P\_7,2.67.1)} katham kṛtvā jñāpakam . \\
\noindent \textbf{(P\_7,2.67.1)} imau vṛṅvṛñau udāttau tayoḥ prakṛtilakṣaṇaḥ pratiṣedhaḥ na prāpnoti . \\
\noindent \textbf{(P\_7,2.67.1)} paśyati tu ācāryaḥ ubhayoḥ saḥ pratyārambhaḥ iti tataḥ vṛṅvṛñoḥ grahaṇam karoti . \\
\noindent \textbf{(P\_7,2.67.1)} na khalu api kaḥ cit ubhayavān pratiṣedhaḥ prakṛtilakṣaṇaḥ pratyayalakṣaṇaḥ ca . \\
\noindent \textbf{(P\_7,2.67.1)} tulyajātīye asati yathā eva prakṛtilakṣaṇasaya niyāmakaḥ bhavati evam pratyayalakṣaṇasya api niyāmakaḥ bhaviṣyati . \\
\noindent \textbf{(P\_7,2.67.1)} atha yāvatā vasau ekājbhyaḥ iṭā bhavitavyam kaḥ nu atra viśeṣaḥ niyamārthe vā sati vidhyarthe vā . \\
\noindent \textbf{(P\_7,2.67.1)} na khalu kaḥ cit viśeṣaḥ . \\
\noindent \textbf{(P\_7,2.67.1)} āhopuruṣikāmātram tu bhavān āha vidhyartham iti . \\
\noindent \textbf{(P\_7,2.67.1)} vayam tu brūmaḥ niyamārtham iti \\
\noindent \textbf{(P\_7,2.67.2)} atha ekājgrahaṇam kimartham . \\
\noindent \textbf{(P\_7,2.67.2)} iha mā bhūt . \\
\noindent \textbf{(P\_7,2.67.2)} bibhidvān , cicchidvān iti . \\
\noindent \textbf{(P\_7,2.67.2)} kriyamāṇe api vā ekājgrahaṇe atra prāpnoti . \\
\noindent \textbf{(P\_7,2.67.2)} eṣaḥ api hi ekāc . \\
\noindent \textbf{(P\_7,2.67.2)} evam tarhi kṛte dvirvacane yaḥ ekāc . \\
\noindent \textbf{(P\_7,2.67.2)} kim vaktavyam etat . \\
\noindent \textbf{(P\_7,2.67.2)} na hi . \\
\noindent \textbf{(P\_7,2.67.2)} katham anucyamānam gaṃsyate . \\
\noindent \textbf{(P\_7,2.67.2)} ekājgrahaṇasāmarthyāt . \\
\noindent \textbf{(P\_7,2.67.2)} na hi kaḥ cit akṛte cirvacane enakāc asti yadartham ekājgrahaṇam kriyate . \\
\noindent \textbf{(P\_7,2.67.2)} nanu ca ayam asti jāgartiḥ . \\
\noindent \textbf{(P\_7,2.67.2)} gāgṛvāṃsaḥ anu gman . \\
\noindent \textbf{(P\_7,2.67.2)} yat tarhi ākāragrahaṇam karoti na hi kaḥ cit akṛte dvirvacane ākārāntaḥ anekāc asti . \\
\noindent \textbf{(P\_7,2.67.2)} nanu ca ayam asti daridrātiḥ . \\
\noindent \textbf{(P\_7,2.67.2)} na daridrāteḥ iṭā bhavitavyam . \\
\noindent \textbf{(P\_7,2.67.2)} kim kāraṇam . \\
\noindent \textbf{(P\_7,2.67.2)} uktam etat daridrāteḥ ārdhadhātuke lopaḥ siddhaḥ ca pratyayavidhau iti . \\
\noindent \textbf{(P\_7,2.67.2)} yaḥ ca idānīm pratyayavidhau siddhaḥ siddhaḥ asau iḍvidhau . \\
\noindent \textbf{(P\_7,2.67.2)} evam api bhūtapūrvagatiḥ vijñāyeta . \\
\noindent \textbf{(P\_7,2.67.2)} ākārāntaḥ yaḥ bhūtapūrvaḥ iti . \\
\noindent \textbf{(P\_7,2.67.2)} ekājgrahaṇam eva tarhi jñāpakam . \\
\noindent \textbf{(P\_7,2.67.2)} nanu ca uktam jāgartyartham etat syāt . \\
\noindent \textbf{(P\_7,2.67.2)} na ekam udāharaṇam ekājgrahaṇam prayojayati . \\
\noindent \textbf{(P\_7,2.67.2)} yadi etāvat prayojanam syāt jāgarteḥ na iti eva bhrūyāt \\
\noindent \textbf{(P\_7,2.67.3)} atha ghasigrahaṇam kimartha na ekāc iti eva siddham . \\
\noindent \textbf{(P\_7,2.67.3)} ghasigrahaṇam anackatvāt . \\
\noindent \textbf{(P\_7,2.67.3)} ghasigrahaṇam kriyate lope kṛte anackatvāt iṭ na prāpnoti . \\
\noindent \textbf{(P\_7,2.67.3)} idam iha sampradhāryam . \\
\noindent \textbf{(P\_7,2.67.3)} iṭ kriyatām lopaḥ iti kim atra kartavyam . \\
\noindent \textbf{(P\_7,2.67.3)} paratvāt iḍāgamaḥ . \\
\noindent \textbf{(P\_7,2.67.3)} nityaḥ lopaḥ . \\
\noindent \textbf{(P\_7,2.67.3)} kṛte api iṭi prāpnoti akṛte api . \\
\noindent \textbf{(P\_7,2.67.3)} iṭ api nityaḥ . \\
\noindent \textbf{(P\_7,2.67.3)} kṛte api lope prāpnoti akṛte api . \\
\noindent \textbf{(P\_7,2.67.3)} anityaḥ iṭ na hi kṛte lope prāpnoti . \\
\noindent \textbf{(P\_7,2.67.3)} kim kāraṇam . \\
\noindent \textbf{(P\_7,2.67.3)} anackatvāt . \\
\noindent \textbf{(P\_7,2.67.3)} evam tarhi dvirvacane kṛte abhyāse yaḥ akāraḥ tadāśrayaḥ iṭ bhaviṣyati . \\
\noindent \textbf{(P\_7,2.67.3)} na sidhyati . \\
\noindent \textbf{(P\_7,2.67.3)} kim kāraṇam . \\
\noindent \textbf{(P\_7,2.67.3)} dvitvāt lopasya paratvāt . \\
\noindent \textbf{(P\_7,2.67.3)} dvirvacanam kriyatām lopaḥ iti kim atra kartavyam . \\
\noindent \textbf{(P\_7,2.67.3)} paratvāt lopaḥ . \\
\noindent \textbf{(P\_7,2.67.3)} lope kṛte anackatvāt dvirvacanam na prāpnoti . \\
\noindent \textbf{(P\_7,2.67.3)} ghasigrahaṇe punaḥ kriyamāṇe na doṣaḥ bhavati . \\
\noindent \textbf{(P\_7,2.67.3)} katham . \\
\noindent \textbf{(P\_7,2.67.3)} vacanāt iṭ bhaviṣyati . \\
\noindent \textbf{(P\_7,2.67.3)} iṭi kṛte dvirvacanam kriyatām lopaḥ iti yadi api paratvāt lopaḥ sthānivadbhāvāt dvirvacanam bhaviṣyati \\
\noindent \textbf{(P\_7,2.68)} dṛśeḥ ca iti vaktavyam . \\
\noindent \textbf{(P\_7,2.68)} dadṛśvān , dadṛśivān . \\
\noindent \textbf{(P\_7,2.68)} tat tarhi vaktavyam . \\
\noindent \textbf{(P\_7,2.68)} na vaktavyam . \\
\noindent \textbf{(P\_7,2.68)} dṛśeḥ iti vartate \\
\noindent \textbf{(P\_7,2.70)} svarateḥ veṭtvāt ṛtaḥ sye vipratiṣedhena . \\
\noindent \textbf{(P\_7,2.70)} svaratilakṣaṇāt vāvacanāt ṛtaḥ sye iti etat bhavati vipratiṣedhena . \\
\noindent \textbf{(P\_7,2.70)} svaratilakṣaṇasya vāvacanasya avakāśaḥ . \\
\noindent \textbf{(P\_7,2.70)} svartā , svaritā . \\
\noindent \textbf{(P\_7,2.70)} ṛtaḥ sye iti asya avakāśaḥ . \\
\noindent \textbf{(P\_7,2.70)} kariṣyate , hariṣyate . \\
\noindent \textbf{(P\_7,2.70)} iha ubhayam prāpnoti . \\
\noindent \textbf{(P\_7,2.70)} svariṣyati , asvariṣyat . \\
\noindent \textbf{(P\_7,2.70)} ṛtaḥ sye iti etat bhavati vipratiṣedhena \\
\noindent \textbf{(P\_7,2.73)} kim udāharaṇam . \\
\noindent \textbf{(P\_7,2.73)} ayaṃsīt , vyaraṃsīt , anaṃsīt , ayāsīt , avāsīt . \\
\noindent \textbf{(P\_7,2.73)} na etat asti . \\
\noindent \textbf{(P\_7,2.73)} na asti atra viśeṣaḥ sati vā iṭi asati vā . \\
\noindent \textbf{(P\_7,2.73)} idam tarhi . \\
\noindent \textbf{(P\_7,2.73)} ayaṃsiṣṭām , ayaṃsiṣuḥ , vyaraṃsiṣṭām , vyaraṃsiṣuḥ , anaṃsiṣṭām , anaṃsiṣuḥ , ayāsiṣṭām , ayāsiṣuḥ , avāsiṣṭām , avāsiṣuḥ . \\
\noindent \textbf{(P\_7,2.73)} idam ca api udāharaṇam . \\
\noindent \textbf{(P\_7,2.73)} ayaṃsīt , vyaraṃsīt , anaṃsīt , ayāsīt , avāsīt . \\
\noindent \textbf{(P\_7,2.73)} nanu ca uktam na asti atra viśeṣaḥ sati vā iṭi asati vā iti . \\
\noindent \textbf{(P\_7,2.73)} ayam asti viśeṣaḥ . \\
\noindent \textbf{(P\_7,2.73)} yadi atra iṭ na syāt vṛddhiḥ prasajyeta . \\
\noindent \textbf{(P\_7,2.73)} iṭi punaḥ sati na iṭi iti pratiṣedhaḥ siddhaḥ bhavati . \\
\noindent \textbf{(P\_7,2.73)} mā bhūt evam hmyantānām iti evam bhaviṣyati . \\
\noindent \textbf{(P\_7,2.73)} atra api na iṭi iti anuvartate . \\
\noindent \textbf{(P\_7,2.73)} tat ca avaśyam iḍgrahaṇam anuvartyam adhākṣīt iti evamartham . \\
\noindent \textbf{(P\_7,2.73)} ākārāntāḥ ca api padapūrvāḥ ekavacane udāharaṇam . \\
\noindent \textbf{(P\_7,2.73)} mā hi yāsīt . \\
\noindent \textbf{(P\_7,2.73)} yadi atra iṭ na syāt anudāttasya īṭaḥ śravaṇam prasajyeta . \\
\noindent \textbf{(P\_7,2.73)} iṭi punaḥ sati uktam etat arthavat tu sicaḥ citkaraṇasāmarthyāt hi iṭaḥ udāttatvam iti tatra ekādeśaḥ udāttena udāttaḥ iti udāttatvam siddham bhavati \\
\noindent \textbf{(P\_7,2.77-78)} KA\_III,302.13-18 Ro\_V,152.13-153.4 {1/13}     kimarthaḥ yogavibhāgaḥ na iśīḍajanām sdhve iti eva ucyeta . \\
\noindent \textbf{(P\_7,2.77-78)} KA\_III,302.13-18 Ro\_V,152.13-153.4 {2/13}     īśaḥ dhve mā bhūt iti . \\
\noindent \textbf{(P\_7,2.77-78)} KA\_III,302.13-18 Ro\_V,152.13-153.4 {3/13}     iṣyate eva : īśidhve iti . \\
\noindent \textbf{(P\_7,2.77-78)} KA\_III,302.13-18 Ro\_V,152.13-153.4 {4/13}     īḍajanoḥ tarhi se mā bhūt iti . \\
\noindent \textbf{(P\_7,2.77-78)} KA\_III,302.13-18 Ro\_V,152.13-153.4 {5/13}     iṣyate eva . \\
\noindent \textbf{(P\_7,2.77-78)} KA\_III,302.13-18 Ro\_V,152.13-153.4 {6/13}     īḍiṣe , janiṣe iti . \\
\noindent \textbf{(P\_7,2.77-78)} KA\_III,302.13-18 Ro\_V,152.13-153.4 {7/13}     īśaḥ tarhi sve mā bhūt iti . \\
\noindent \textbf{(P\_7,2.77-78)} KA\_III,302.13-18 Ro\_V,152.13-153.4 {8/13}     iṣyate eva . \\
\noindent \textbf{(P\_7,2.77-78)} KA\_III,302.13-18 Ro\_V,152.13-153.4 {9/13}     īśiṣva iti . \\
\noindent \textbf{(P\_7,2.77-78)} KA\_III,302.13-18 Ro\_V,152.13-153.4 {10/13}     se tarhi yaḥ svaśabdaḥ tatra yathā syāt kriyāsamabhihāre yaḥ svaśabdaḥ tatra mā bhūt iti . \\
\noindent \textbf{(P\_7,2.77-78)} KA\_III,302.13-18 Ro\_V,152.13-153.4 {11/13}     atra api iṣyate . \\
\noindent \textbf{(P\_7,2.77-78)} KA\_III,302.13-18 Ro\_V,152.13-153.4 {12/13}     saḥ bhavān īśiṣva iti eva ayam īṣṭe iti . \\
\noindent \textbf{(P\_7,2.77-78)} KA\_III,302.13-18 Ro\_V,152.13-153.4 {13/13}     ātaḥ ca iṣyate evam hi āha siddham tu loṇmadhyamapuruṣaikavacanasya kriyāsamabhihāre dvirvacanāt iti . \\
\noindent \textbf{(P\_7,2.80)} kim sārvadhātukagrahaṇam anuvartate utāho na . \\
\noindent \textbf{(P\_7,2.80)} kim ca arthaḥ anuvṛttyā . \\
\noindent \textbf{(P\_7,2.80)} bāḍham arthaḥ yadi akārāt paraḥ yāśabdaḥ ārdhadhātukam asti . \\
\noindent \textbf{(P\_7,2.80)} nanu ca ayam asti . \\
\noindent \textbf{(P\_7,2.80)} cikīrṣyāt , jihīrṣyāt . \\
\noindent \textbf{(P\_7,2.80)} lopaḥ atra bādhakaḥ bhaviṣyati . \\
\noindent \textbf{(P\_7,2.80)} kim tarhi asmin yoge udāharaṇam . \\
\noindent \textbf{(P\_7,2.80)} pacet , yajet . \\
\noindent \textbf{(P\_7,2.80)} atra api ataḥ dīrghaḥ yañi iti dīrghatvam prāpnoti . \\
\noindent \textbf{(P\_7,2.80)} saḥ yathā eva ayādeśaḥ dīrghatvam bādhate evam lopam api bādheta . \\
\noindent \textbf{(P\_7,2.80)} tasmāt sārvadhātukagrahaṇam anuvartyam \\
\noindent \textbf{(P\_7,2.82)} muki svare doṣaḥ . \\
\noindent \textbf{(P\_7,2.82)} muki sati svare doṣaḥ bhavati . \\
\noindent \textbf{(P\_7,2.82)} pacamānaḥ , yajamānaḥ . \\
\noindent \textbf{(P\_7,2.82)} mukā vyavahitatvāt anudāttatvam na prāpnoti . \\
\noindent \textbf{(P\_7,2.82)} nanu ca ayam muk adupadeśabhaktaḥ adupadeśagrahaṇena grāhiṣyate . \\
\noindent \textbf{(P\_7,2.82)} na sidhyati . \\
\noindent \textbf{(P\_7,2.82)} aṅgasya muk ucyate vikaraṇāntam ca aṅgam saḥ ayam saṅghātabhaktaḥ aśakyaḥ adupadeśagrahaṇena grahītum . \\
\noindent \textbf{(P\_7,2.82)} evam tarhi abhaktaḥ kariṣyate . \\
\noindent \textbf{(P\_7,2.82)} abhakte ca . \\
\noindent \textbf{(P\_7,2.82)} kim . \\
\noindent \textbf{(P\_7,2.82)} svare doṣaḥ bhavati . \\
\noindent \textbf{(P\_7,2.82)} pacamānaḥ , yajamānaḥ . \\
\noindent \textbf{(P\_7,2.82)} mukā vyavahitatvāt anudāttatvam na prāpnoti . \\
\noindent \textbf{(P\_7,2.82)} evam tarhi parādiḥ kariṣyate . \\
\noindent \textbf{(P\_7,2.82)} parādau dīrghaprasaṅgaḥ . \\
\noindent \textbf{(P\_7,2.82)} yadi parādiḥ kriyate ataḥ dīrghaḥ yañi iti dīrghatvam prāpnoti . \\
\noindent \textbf{(P\_7,2.82)} na eṣaḥ doṣaḥ . \\
\noindent \textbf{(P\_7,2.82)} tiṅi iti evam tat . \\
\noindent \textbf{(P\_7,2.82)} sidhyati . \\
\noindent \textbf{(P\_7,2.82)} sūtram tarhi bhidyate . \\
\noindent \textbf{(P\_7,2.82)} yathānyāsam eva astu . \\
\noindent \textbf{(P\_7,2.82)} nanu ca uktam muki svare doṣaḥ iti . \\
\noindent \textbf{(P\_7,2.82)} parihṛtam etat adupadeśabhaktaḥ adupadeśagrahaṇena grāhiṣyate . \\
\noindent \textbf{(P\_7,2.82)} nanu ca uktam aṅgasya muk ucyate vikaraṇāntam ca aṅgam saḥ ayam saṅghātabhaktaḥ aśakyaḥ adupadeśagrahaṇena grahītum iti . \\
\noindent \textbf{(P\_7,2.82)} atha ayam adbhaktaḥ syāt gṛhyeta adupadeśagrahaṇena . \\
\noindent \textbf{(P\_7,2.82)} bāḍham gṛhyeta . \\
\noindent \textbf{(P\_7,2.82)} adbhaktaḥ tarhi bhaviṣyati . \\
\noindent \textbf{(P\_7,2.82)} tat katham . \\
\noindent \textbf{(P\_7,2.82)} ataḥ yā iyaḥ iti atra akāragrahaṇam pañcamīnirdiṣṭam aṅgasya iti ca ṣaṣṭhīnirdiṣṭam tatra aśakyam vivibhaktitvāt ataḥ iti pañcamyā aṅgam viśeṣayitum . \\
\noindent \textbf{(P\_7,2.82)} tat prakṛtam iha anuvartiṣyate . \\
\noindent \textbf{(P\_7,2.82)} evam api ṣaṣṭhyabhāvāt na prāpnoti . \\
\noindent \textbf{(P\_7,2.82)} ānaḥ iti eṣā saptamī ataḥ iti pañcamyāḥ ṣaṣṭhīm prakalpayiṣyati tasmin iti nirdiṣṭe pūrvasya iti \\
\noindent \textbf{(P\_7,2.84)} aṣṭanjanādipathimathyātveṣu āntaratamyāt anunāsikaprasaṅgaḥ . \\
\noindent \textbf{(P\_7,2.84)} aṣṭanjanādipathimathyātveṣu āntaratamyāt anunāsikaḥ prāpnoti . \\
\noindent \textbf{(P\_7,2.84)} aṣṭābhiḥ , aṣṭābhyaḥ , jātaḥ , jātavān , panthāḥ , manthāḥ . \\
\noindent \textbf{(P\_7,2.84)} siddham anaṇtvāt . \\
\noindent \textbf{(P\_7,2.84)} siddham etat . \\
\noindent \textbf{(P\_7,2.84)} katham . \\
\noindent \textbf{(P\_7,2.84)} anaṇtvāt . \\
\noindent \textbf{(P\_7,2.84)} katham anaṇtvam . \\
\noindent \textbf{(P\_7,2.84)} aṇsavarṇān gṛhṇāti iti ucyate na ca akāraḥ aṇ . \\
\noindent \textbf{(P\_7,2.84)} uccāraṇasāmarthyāt vā . \\
\noindent \textbf{(P\_7,2.84)} atha vā śuddhoccāraṇasāmarthyāt na bhaviṣyati . \\
\noindent \textbf{(P\_7,2.84)} na etau staḥ parihārau . \\
\noindent \textbf{(P\_7,2.84)} yat tāvat ucyate anaṇtvāt iti na brūmaḥ aṇsavarṇān gṛhṇāti iti . \\
\noindent \textbf{(P\_7,2.84)} katham tarhi taparaḥ tatkālasya iti . \\
\noindent \textbf{(P\_7,2.84)} yat api ucyate uccāraṇasāmarthyāt vā iti asti anyat uccāraṇe prayojanam . \\
\noindent \textbf{(P\_7,2.84)} kim . \\
\noindent \textbf{(P\_7,2.84)} uttarārtham . \\
\noindent \textbf{(P\_7,2.84)} rāyaḥ hali iti . \\
\noindent \textbf{(P\_7,2.84)} evam tarhi na imau pṛthakparihārau . \\
\noindent \textbf{(P\_7,2.84)} ekapariharaḥ ayam . \\
\noindent \textbf{(P\_7,2.84)} siddham anaṇtvāt uccāraṇasāmarthyāt vā iti . \\
\noindent \textbf{(P\_7,2.84)} iha tāvat aṣṭābhiḥ , aṣṭābhyaḥ iti anaṇtvāt siddham . \\
\noindent \textbf{(P\_7,2.84)} jātaḥ , jātavān , panthāḥ , manthāḥ uccāraṇsāmarthyāt siddham . \\
\noindent \textbf{(P\_7,2.84)} yadi evam pṛthakparihārayoḥ api na doṣaḥ . \\
\noindent \textbf{(P\_7,2.84)} yaḥ yatra parihāraḥ saḥ tatra bhaviṣyati \\
\noindent \textbf{(P\_7,2.86)} anādeśagrahaṇam śakyam akartum . \\
\noindent \textbf{(P\_7,2.86)} katham hali iti anuvartate na ca ādeśaḥ halādiḥ asti . \\
\noindent \textbf{(P\_7,2.86)} tat etat anādeśagrahaṇam tiṣṭhatu tāvat sānnyāsikam \\
\noindent \textbf{(P\_7,2.89)} ajgrahaṇam śakyam akartum . \\
\noindent \textbf{(P\_7,2.89)} katham . \\
\noindent \textbf{(P\_7,2.89)} aviśeṣeṇa yatvam utsargaḥ tasya halādau ātvam apavādaḥ \\
\noindent \textbf{(P\_7,2.90)} śeṣagrahaṇam śakyam akartum . \\
\noindent \textbf{(P\_7,2.90)} katham . \\
\noindent \textbf{(P\_7,2.90)} aviśeṣeṇa lopaḥ utsargaḥ tasya ajādau yatvam apavādaḥ halādau ātvam \\
\noindent \textbf{(P\_7,2.91)} parigrahaṇam śakyam akartum . \\
\noindent \textbf{(P\_7,2.91)} māntasya iti eva siddham . \\
\noindent \textbf{(P\_7,2.91)} na sidhyati . \\
\noindent \textbf{(P\_7,2.91)} kim kāraṇam . \\
\noindent \textbf{(P\_7,2.91)} antaśabdasya ubhayārthatvāt . \\
\noindent \textbf{(P\_7,2.91)} katham . \\
\noindent \textbf{(P\_7,2.91)} ayam antaśabdaḥ asti eva saha tena vartate . \\
\noindent \textbf{(P\_7,2.91)} tat yathā . \\
\noindent \textbf{(P\_7,2.91)} maryādāntam devadattasya kṣetram . \\
\noindent \textbf{(P\_7,2.91)} saha maryādayā iti gamyate . \\
\noindent \textbf{(P\_7,2.91)} asti prāk tasmāt vartate . \\
\noindent \textbf{(P\_7,2.91)} tat yathā . \\
\noindent \textbf{(P\_7,2.91)} nadyantam devadattasya kṣetram iti . \\
\noindent \textbf{(P\_7,2.91)} prāk nadyāḥ iti gamyate . \\
\noindent \textbf{(P\_7,2.91)} tat yaḥ saha tena vartate tasya idam grahaṇam yathā vijñāyeta . \\
\noindent \textbf{(P\_7,2.91)} na etat asti prayojanam . \\
\noindent \textbf{(P\_7,2.91)} sarvatra eva antaśabdaḥ saha tena vartate . \\
\noindent \textbf{(P\_7,2.91)} atha katham nadyantam devadattasya kṣetram iti . \\
\noindent \textbf{(P\_7,2.91)} nadyāḥ kṣetratve sambhaḥ na asti iti kṛtvā prāk nadyāḥ iti gamyate . \\
\noindent \textbf{(P\_7,2.91)} avadhidyotanārtham tarhi parigrahaṇam kartavyam . \\
\noindent \textbf{(P\_7,2.91)} māntasya iti iyati ucyamāne yatra eva mānte yuṣmadasmadī tatra eva ādeśāḥ syuḥ . \\
\noindent \textbf{(P\_7,2.91)} kva ca mānte yuṣmadasmadī . \\
\noindent \textbf{(P\_7,2.91)} yuṣmān ācaṣṭe , asmān ācaṣṭe iti yuṣmayateḥ asmayateḥ ca apratyayaḥ \\
\noindent \textbf{(P\_7,2.98.1)} kimartham idam ucyate na tvamau ekavacane iti eva siddham . \\
\noindent \textbf{(P\_7,2.98.1)} na sidhyati . \\
\noindent \textbf{(P\_7,2.98.1)} kim kāraṇam . \\
\noindent \textbf{(P\_7,2.98.1)} ekavacanābhāvāt . \\
\noindent \textbf{(P\_7,2.98.1)} ekavacane iti ucyate na ca atra ekavacanam paśyāmaḥ . \\
\noindent \textbf{(P\_7,2.98.1)} pratyayalakṣaṇena . \\
\noindent \textbf{(P\_7,2.98.1)} na lumatā aṅgasya iti pratyayalakṣaṇasya pratiṣedhaḥ . \\
\noindent \textbf{(P\_7,2.98.1)} evam tarhi idam iha sampradhāryam . \\
\noindent \textbf{(P\_7,2.98.1)} luk kriyatām ādeśau iti kim atra kartavyam . \\
\noindent \textbf{(P\_7,2.98.1)} paratvāt ādeśau . \\
\noindent \textbf{(P\_7,2.98.1)} nityaḥ luk . \\
\noindent \textbf{(P\_7,2.98.1)} kṛtayoḥ api ādeśayoḥ prāpnoti akṛtayoḥ api . \\
\noindent \textbf{(P\_7,2.98.1)} antaraṅgau ādeśau . \\
\noindent \textbf{(P\_7,2.98.1)} evam tarhi siddhe sati yatpratyayottarapadayoḥ tvamau śāsti tat jñāpayati ācāryaḥ antaraṅgān api vidhīn bādhitvā bahiraṅgaḥ luk bhavati iti . \\
\noindent \textbf{(P\_7,2.98.1)} kim etasya jñāpane prayojanam . \\
\noindent \textbf{(P\_7,2.98.1)} gomān priyaḥ asya gomatpriyaḥ , yavamatpriyaḥ gomān iva ācarati gomatyate , yavamatyate antaraṅgān api numādīn bahiraṅgaḥ luk bādhate iti . \\
\noindent \textbf{(P\_7,2.98.1)} na etat asti jñāpakam . \\
\noindent \textbf{(P\_7,2.98.1)} asti anyat etasya vacane prayojanam . \\
\noindent \textbf{(P\_7,2.98.1)} kim . \\
\noindent \textbf{(P\_7,2.98.1)} ye anye ekavacanādeśāḥ prāpnuvanti tadbādhanārtham etat syāt . \\
\noindent \textbf{(P\_7,2.98.1)} tat yathā . \\
\noindent \textbf{(P\_7,2.98.1)} tava putraḥ tvatputraḥ , mama putraḥ matputraḥ , tubhyam hitam tvaddhitam , mahyam hitam maddhitam iti . \\
\noindent \textbf{(P\_7,2.98.1)} yat tarhi maparyantagrahaṇam anuvartayati . \\
\noindent \textbf{(P\_7,2.98.1)} yati atra anye ekavacanādeśāḥ syuḥ mapartyantānuvṛttiḥ anarthikā syāt \\
\noindent \textbf{(P\_7,2.99)} tisṛbhāve sañjñāyām kani upasaṅkhyānam . \\
\noindent \textbf{(P\_7,2.99)} tisṛbhāve sañjñāyām kani upasaṅkhyānam kartavyam . \\
\noindent \textbf{(P\_7,2.99)} tisṛkā nāma grāmaḥ . \\
\noindent \textbf{(P\_7,2.99)} catasari ādyudāttanipātanam ca . \\
\noindent \textbf{(P\_7,2.99)} catasari ādyudāttanipātanam kartavyam . \\
\noindent \textbf{(P\_7,2.99)} tricaturoḥ striyām tisṛcatasṛ . \\
\noindent \textbf{(P\_7,2.99)} kim prayojanam . \\
\noindent \textbf{(P\_7,2.99)} catasraḥ paśya . \\
\noindent \textbf{(P\_7,2.99)} śasi svaraḥ mā bhūt iti . \\
\noindent \textbf{(P\_7,2.99)} kim ca anyat . \\
\noindent \textbf{(P\_7,2.99)} upadeśivadvacanam ca . \\
\noindent \textbf{(P\_7,2.99)} upadeśivadbhāvaḥ ca vaktavyaḥ . \\
\noindent \textbf{(P\_7,2.99)} kim prayojanam . \\
\noindent \textbf{(P\_7,2.99)} svarasiddhyartham . \\
\noindent \textbf{(P\_7,2.99)} upadeśāvasthāyām eva ādyudāttanipātane kṛte vibhaktisvareṇa bādhanam yathā syāt : catasṛṇām iti . \\
\noindent \textbf{(P\_7,2.99)} saḥ tarhi upadeśivadbhāvaḥ vaktavyaḥ . \\
\noindent \textbf{(P\_7,2.99)} na vaktavyaḥ . \\
\noindent \textbf{(P\_7,2.99)} uktam vā . \\
\noindent \textbf{(P\_7,2.99)} kim uktam . \\
\noindent \textbf{(P\_7,2.99)} vibhaktisvarabhāvaḥ ca halādigrahaṇāt ādyudāttanipātane hi halādigrahaṇānarthakyam iti \\
\noindent \textbf{(P\_7,2.100)} aci rādeśe jasi upasaṅkhyānam guṇaparatvāt . \\
\noindent \textbf{(P\_7,2.100)} aci rādeśe jasi upasaṅkhyānam kartavyam . \\
\noindent \textbf{(P\_7,2.100)} tisraḥ tiṣṭhanti , catasraḥ tiṣṭhanti . \\
\noindent \textbf{(P\_7,2.100)} kim punaḥ kāraṇam na sidhyati . \\
\noindent \textbf{(P\_7,2.100)} guṇaparatvāt . \\
\noindent \textbf{(P\_7,2.100)} paratvāt guṇaḥ prāpnoti . \\
\noindent \textbf{(P\_7,2.100)} tat tarhi vaktavyam . \\
\noindent \textbf{(P\_7,2.100)} na vā anavakāśatvāt rasya . \\
\noindent \textbf{(P\_7,2.100)} na vā vaktavyam . \\
\noindent \textbf{(P\_7,2.100)} kim kāraṇam . \\
\noindent \textbf{(P\_7,2.100)} anavakāśatvāt rasya . \\
\noindent \textbf{(P\_7,2.100)} anavakāśaḥ rādeśaḥ guṇam bādhiṣyate . \\
\noindent \textbf{(P\_7,2.100)} sāvakāśaḥ rādeśaḥ . \\
\noindent \textbf{(P\_7,2.100)} kaḥ avakāśaḥ . \\
\noindent \textbf{(P\_7,2.100)} tisraḥ paśya . \\
\noindent \textbf{(P\_7,2.100)} catasraḥ paśya . \\
\noindent \textbf{(P\_7,2.100)} na eṣaḥ asti avakāśaḥ . \\
\noindent \textbf{(P\_7,2.100)} atra api pūrvasavarṇadīrghaḥ prāpnoti . \\
\noindent \textbf{(P\_7,2.100)} saḥ yathā eva pūrvasavarṇam bādhate evam guṇam api bādhiṣyate . \\
\noindent \textbf{(P\_7,2.100)} guṇaḥ api anavakāśaḥ . \\
\noindent \textbf{(P\_7,2.100)} sāvakāśaḥ guṇaḥ . \\
\noindent \textbf{(P\_7,2.100)} kaḥ avakāśaḥ . \\
\noindent \textbf{(P\_7,2.100)} he kartaḥ . \\
\noindent \textbf{(P\_7,2.100)} na eṣaḥ sarvanāmasthāne guṇaḥ . \\
\noindent \textbf{(P\_7,2.100)} kaḥ tarhi . \\
\noindent \textbf{(P\_7,2.100)} sambuddhiguṇaḥ . \\
\noindent \textbf{(P\_7,2.100)} ayam tarhi . \\
\noindent \textbf{(P\_7,2.100)} he mātaḥ . \\
\noindent \textbf{(P\_7,2.100)} eṣaḥ api sambuddhiguṇaḥ eva . \\
\noindent \textbf{(P\_7,2.100)} na atra sambuddhiguṇaḥ prāpnoti . \\
\noindent \textbf{(P\_7,2.100)} kim kāraṇam . \\
\noindent \textbf{(P\_7,2.100)} ambārthanadyoḥ hrasvaḥ iti hrasvatvena bhavitavyam . \\
\noindent \textbf{(P\_7,2.100)} bhavet dīrghāṇām hrasvavacanasāmarthyāt na syāt hrasvānām tu khalu hrasvatvam kriyatām sambuddhiguṇaḥ iti paratvāt sambuddhiguṇena bhavitavyam . \\
\noindent \textbf{(P\_7,2.100)} atha api katham cit sāvakāsaḥ guṇaḥ syāt evam api na doṣaḥ . \\
\noindent \textbf{(P\_7,2.100)} purastāt apavādāḥ anantarān vidhīn bādhante iti evam ayam rādeśaḥ jasi guṇam bādhate sarvanāmasthānaguṇam na bādhiṣyate . \\
\noindent \textbf{(P\_7,2.100)} tasmāt suṣṭhu ucyate aci rādeśe jasi upasaṅkhyānam guṇaparatvāt iti \\
\noindent \textbf{(P\_7,2.101)} numaḥ anaṅjarasau bhavataḥ vipratiṣedhena . \\
\noindent \textbf{(P\_7,2.101)} numaḥ avakāśaḥ . \\
\noindent \textbf{(P\_7,2.101)} trapuṇī , jatunī , tumburuṇī . \\
\noindent \textbf{(P\_7,2.101)} anaṅaḥ avakāśaḥ . \\
\noindent \textbf{(P\_7,2.101)} priyasakthnā brāhmaṇena . \\
\noindent \textbf{(P\_7,2.101)} iha ubhayam prāpnoti . \\
\noindent \textbf{(P\_7,2.101)} dadhnā , sakthnā . \\
\noindent \textbf{(P\_7,2.101)} jarasaḥ avakāśaḥ . \\
\noindent \textbf{(P\_7,2.101)} jarasā , jarase . \\
\noindent \textbf{(P\_7,2.101)} numaḥ avakāśaḥ . \\
\noindent \textbf{(P\_7,2.101)} kuṇḍāni , vanāni . \\
\noindent \textbf{(P\_7,2.101)} iha ubhayam prāpnoti . \\
\noindent \textbf{(P\_7,2.101)} atijarāṃsi brāhmaṇakulāni . \\
\noindent \textbf{(P\_7,2.101)} anaṅjarasau numaḥ bhavataḥ vipratiṣedhena . \\
\noindent \textbf{(P\_7,2.101)} atha iha luk kasmāt na bhavati . \\
\noindent \textbf{(P\_7,2.101)} atijarasam paśya iti . \\
\noindent \textbf{(P\_7,2.101)} kim punaḥ kāraṇam dvitīyaikavacanam eva udāhriyate na punaḥ prathamaikavacanam api . \\
\noindent \textbf{(P\_7,2.101)} atijarasam tiṣṭhati iti . \\
\noindent \textbf{(P\_7,2.101)} asti atra viśeṣaḥ . \\
\noindent \textbf{(P\_7,2.101)} na atra akṛte ambhāve jarasbhāvaḥ prāpnoti . \\
\noindent \textbf{(P\_7,2.101)} kim kāraṇam . \\
\noindent \textbf{(P\_7,2.101)} aci iti ucyate . \\
\noindent \textbf{(P\_7,2.101)} yadā ca jarasbhāvaḥ kṛtaḥ tadā luk na bhaviṣyati sannipātalakṣaṇaḥ vidhiḥ animittam tadvighātasya iti . \\
\noindent \textbf{(P\_7,2.101)} yadi evam atijarasam , atijarasaiḥ iti atra na prāpnoti atijaram , atijaraiḥ iti bhavitavyam . \\
\noindent \textbf{(P\_7,2.101)} gonardīyaḥ āha . \\
\noindent \textbf{(P\_7,2.101)} iṣṭam eva etat saṅgṛhītam bhavati . \\
\noindent \textbf{(P\_7,2.101)} atijaram atijaraiḥ iti bhavitavyam satyām etasyām paribhāṣāyām sannipātalakṣaṇaḥ vidhiḥ animittam tadvighātasya iti \\
\noindent \textbf{(P\_7,2.102)} tyadādīnām dviparyantānām akāravacanam . \\
\noindent \textbf{(P\_7,2.102)} tyadādīnām dviparyantānām atvam vaktavyam . \\
\noindent \textbf{(P\_7,2.102)} kim prayojanam . \\
\noindent \textbf{(P\_7,2.102)} yuṣmadasmadantānām bhavadantānām vā mā bhūt iti . \\
\noindent \textbf{(P\_7,2.102)} tat tarhi vaktavyam . \\
\noindent \textbf{(P\_7,2.102)} na vaktavyam . \\
\noindent \textbf{(P\_7,2.102)} tyadādīnām akāreṇa siddhatvāt yuṣmadasmadoḥ , śeṣe lopasya lopena jñāyate prāk tataḥ at iti . \\
\noindent \textbf{(P\_7,2.102)} yat ayam tyadādīnām atvena siddhe yuṣmadasmadoḥ śeṣe lopam śāsti tat jñāpayati ācāryaḥ prāk tataḥ atvam bhavati iti . \\
\noindent \textbf{(P\_7,2.102)} na sarveṣām iti . \\
\noindent \textbf{(P\_7,2.102)} api vā upasamastārtham atvābhāvāt kṛtam bhavet . \\
\noindent \textbf{(P\_7,2.102)} na etat asti prayojanam . \\
\noindent \textbf{(P\_7,2.102)} upasamastārtham etat syāt : atiyūyam , ativayam . \\
\noindent \textbf{(P\_7,2.102)} upasamastānām hi tyadādīnām atvam na iṣyate : atitat , atitadau , atitadaḥ . \\
\noindent \textbf{(P\_7,2.102)} ṭilopaḥ ṭābabhāvārthaḥ kartavyaḥ iti tat smṛtam . \\
\noindent \textbf{(P\_7,2.102)} yaḥ tu śeṣe lopaḥ ṭilopaḥ saḥ vaktavyaḥ . \\
\noindent \textbf{(P\_7,2.102)} kim prayojanam . \\
\noindent \textbf{(P\_7,2.102)} ṭāppratiṣedhārtham . \\
\noindent \textbf{(P\_7,2.102)} ṭāp mā bhūt iti . \\
\noindent \textbf{(P\_7,2.102)} saḥ tarhi ṭilopaḥ vaktavyaḥ . \\
\noindent \textbf{(P\_7,2.102)} na vaktavyaḥ . \\
\noindent \textbf{(P\_7,2.102)} atha vā śeṣasaptamyā śeṣe lopaḥ vidhīyate . \\
\noindent \textbf{(P\_7,2.102)} iha yuṣmadasmadoḥ lopaḥ iti iyatā antyasya lopaḥ siddhaḥ . \\
\noindent \textbf{(P\_7,2.102)} saḥ ayam evam siddhe sati yat śeṣagrahaṇam karoti tasya etat prayojanam avaśiṣṭasya lopaḥ yathā syāt iti . \\
\noindent \textbf{(P\_7,2.102)} luptaśiṣṭe hi tasya āhuḥ kāryasiddhim manīṣiṇaḥ . \\
\noindent \textbf{(P\_7,2.102)} evam tarhi ācāryapravṛttiḥ jñāpayati na sarveṣām tyadādīnām atvam bhavati iti yat ayam kimaḥ kaḥ iti kādeśam śāsti . \\
\noindent \textbf{(P\_7,2.102)} itarathā hi kimaḥ at bhavati iti eva brūyāt . \\
\noindent \textbf{(P\_7,2.102)} siddhe vidhiḥ ārabhyamāṇaḥ jñāpakārthaḥ bhavati na ca kimaḥ attvena sidhyati . \\
\noindent \textbf{(P\_7,2.102)} attve hi sati antyasya prasajyeta . \\
\noindent \textbf{(P\_7,2.102)} siddham antyasya pūrveṇa eva tatra ārambhasāmarthyāt ikārasya bhaviṣyati . \\
\noindent \textbf{(P\_7,2.102)} kutaḥ nu khalu etat anantyārthe ārambhe sati ikārasya bhaviṣyati na punaḥ kakārasya syāt . \\
\noindent \textbf{(P\_7,2.102)} yat tarhi kimaḥ grahaṇam karoti . \\
\noindent \textbf{(P\_7,2.102)} itarathā hi kaṭ at bhavati iti eva brūyāt . \\
\noindent \textbf{(P\_7,2.102)} evam api kakāramātrāt parasya prāpnoti . \\
\noindent \textbf{(P\_7,2.102)} tyadādīnām iti vartate na ca anyat kimaḥ tyadādiṣu kakāravat asti . \\
\noindent \textbf{(P\_7,2.102)} evam api anaikāntikam jñāpakam . \\
\noindent \textbf{(P\_7,2.102)} etāvat tu jñāpyate na sarveṣām tyadādīnām atvam bhavati iti tatra kutaḥ etat dviparyantānām bhaviṣyati na punaḥ yuṣmadasmadantānām vā syāt bhavadantānām vā . \\
\noindent \textbf{(P\_7,2.102)} kim ca avaśyam khalu api uttarārtham kimaḥ grahaṇam kartavyam . \\
\noindent \textbf{(P\_7,2.102)} ku tihoḥ kva ati iti . \\
\noindent \textbf{(P\_7,2.102)} kādeśaḥ khalu api avaśyam sākackārthaḥ vaktavyaḥ kaḥ kau ke iti evam artham . \\
\noindent \textbf{(P\_7,2.102)} tasmāt dviparyantānām atvam vaktavyam . \\
\noindent \textbf{(P\_7,2.102)} tyadādīnām akāreṇa siddhatvāt yuṣmadasmadoḥ , śeṣe lopasya lopena jñāyate prāk tataḥ at iti . \\
\noindent \textbf{(P\_7,2.102)} api vā upasamastārtham atvābhāvāt kṛtam bhavet , ṭilopaḥ ṭābabhāvārthaḥ kartavyaḥ iti tat smṛtam . \\
\noindent \textbf{(P\_7,2.102)} atha vā śeṣasaptamyā śeṣe lopaḥ vidhīyate , luptaśiṣṭe hi tasya āhuḥ kāryasiddhim manīṣiṇaḥ \\
\noindent \textbf{(P\_7,2.105)} kimartham kvādeśaḥ ucyate na ku tihāt si iti eva ucyate . \\
\noindent \textbf{(P\_7,2.105)} kā rūpasiddhiḥ : kva . \\
\noindent \textbf{(P\_7,2.105)} yaṇādeśena siddham . \\
\noindent \textbf{(P\_7,2.105)} na sidhyati . \\
\noindent \textbf{(P\_7,2.105)} oḥ guṇaḥ prasajyeta \\
\noindent \textbf{(P\_7,2.106)} kimartham anantyayoḥ iti ucyate . \\
\noindent \textbf{(P\_7,2.106)} antyayoḥ mā bhūt iti . \\
\noindent \textbf{(P\_7,2.106)} na etat asti prayojanam . \\
\noindent \textbf{(P\_7,2.106)} atvam antyayoḥ bādhakam bhaviṣyati . \\
\noindent \textbf{(P\_7,2.106)} anavakāśāḥ vidhayaḥ bādhakāḥ bhavanti sāvakāśam ca atvam . \\
\noindent \textbf{(P\_7,2.106)} kaḥ avakāśaḥ . \\
\noindent \textbf{(P\_7,2.106)} dviśabdaḥ . \\
\noindent \textbf{(P\_7,2.106)} satvam api sāvakāśam . \\
\noindent \textbf{(P\_7,2.106)} kaḥ avakāśaḥ . \\
\noindent \textbf{(P\_7,2.106)} anantyaḥ . \\
\noindent \textbf{(P\_7,2.106)} katham punaḥ sati antye anantyasya satvam syāt . \\
\noindent \textbf{(P\_7,2.106)} bhavet yaḥ takāradakārābhyām aṅgam viśeṣayet tasya anantyayoḥ na syāt . \\
\noindent \textbf{(P\_7,2.106)} vayam tu khalu aṅgena takāradakārau viśeṣayiṣyāmaḥ . \\
\noindent \textbf{(P\_7,2.106)} evam api ubhayoḥ sāvakāsaśayoḥ paratvāt satvam prāpnoti . \\
\noindent \textbf{(P\_7,2.106)} kim ca syāt yadi antyayoḥ satvam syāt . \\
\noindent \textbf{(P\_7,2.106)} iha he saḥ iti eṅhrasvāt iti sambuddhilopaḥ na syāt . \\
\noindent \textbf{(P\_7,2.106)} iha ca yā sā ataḥ iti ṭāp na syāt . \\
\noindent \textbf{(P\_7,2.106)} tasmāt anantyayoḥ iti vaktavyam . \\
\noindent \textbf{(P\_7,2.106)} na vaktavyam . \\
\noindent \textbf{(P\_7,2.106)} evam vakṣyāmi . \\
\noindent \textbf{(P\_7,2.106)} tadoḥ saḥ sau . \\
\noindent \textbf{(P\_7,2.106)} tataḥ adasaḥ . \\
\noindent \textbf{(P\_7,2.106)} adasaḥ ca dakārasya saḥ bhavati iti . \\
\noindent \textbf{(P\_7,2.106)} idam idānīm kimartham . \\
\noindent \textbf{(P\_7,2.106)} niyamārtham . \\
\noindent \textbf{(P\_7,2.106)} adasaḥ eva dakārasya na anyasya dakārasya iti . \\
\noindent \textbf{(P\_7,2.106)} yadi niyamaḥ kriyate dvīyateḥ apratyayaḥ dvaḥ iti prāpnoti svaḥ iti ca iṣyate . \\
\noindent \textbf{(P\_7,2.106)} yathālakṣaṇam aprayukte \\
\noindent \textbf{(P\_7,2.107.1)} adasaḥ soḥ bhavet autvam kim sulopaḥ vidhīyate . \\
\noindent \textbf{(P\_7,2.107.1)} adasaḥ eva soḥ bhavet autvam kimartham sulopaḥ vidhīyate . \\
\noindent \textbf{(P\_7,2.107.1)} hrasvāt lupyeta sambuddhiḥ . \\
\noindent \textbf{(P\_7,2.107.1)} iha he asau iti eṅhrasvāt sambuddheḥ iti lopaḥ prasajyeta . \\
\noindent \textbf{(P\_7,2.107.1)} na halaḥ . \\
\noindent \textbf{(P\_7,2.107.1)} halaḥ lopaḥ sambuddhilopaḥ . \\
\noindent \textbf{(P\_7,2.107.1)} tat halgrahaṇam kartavyam . \\
\noindent \textbf{(P\_7,2.107.1)} prakṛtam hi tat . \\
\noindent \textbf{(P\_7,2.107.1)} prakṛtam halgrahaṇam . \\
\noindent \textbf{(P\_7,2.107.1)} kva prakṛtam . \\
\noindent \textbf{(P\_7,2.107.1)} halṅyābbhyaḥ dīrghāt sutisyapṛktam hal iti . \\
\noindent \textbf{(P\_7,2.107.1)} tat vai prathamānirdiṣṭam ṣaṣṭhīnirdiṣṭena ca iha arthaḥ . \\
\noindent \textbf{(P\_7,2.107.1)} hrasvāt iti eṣā pañcamī hal iti prathamāyāḥ ṣaṣṭhīm prakalpayiṣyati tasmāt iti uttarasya iti . \\
\noindent \textbf{(P\_7,2.107.1)} āpaḥ ettvam bhavet tasmin . \\
\noindent \textbf{(P\_7,2.107.1)} iha he asau brāhmaṇi āṅi ca āpaḥ sambuddhau ca iti ettvam prasajyeta . \\
\noindent \textbf{(P\_7,2.107.1)} na jhali iti anuvartanāt . \\
\noindent \textbf{(P\_7,2.107.1)} jhali iti tatra anuvartate . \\
\noindent \textbf{(P\_7,2.107.1)} kva prakṛtam . \\
\noindent \textbf{(P\_7,2.107.1)} supi ca bahuvacane jhali et iti . \\
\noindent \textbf{(P\_7,2.107.1)} pratyayasthāt ca kāt ittvam . \\
\noindent \textbf{(P\_7,2.107.1)} iha ca asakau brāhmaṇī iti pratyaysthāt kāt pūrvasya iti īttvam prasajyeta . \\
\noindent \textbf{(P\_7,2.107.1)} na eṣaḥ doṣaḥ . \\
\noindent \textbf{(P\_7,2.107.1)} praśliṣṭanirdeśaḥ ayam . \\
\noindent \textbf{(P\_7,2.107.1)} ā , āp , āp iti . \\
\noindent \textbf{(P\_7,2.107.1)} iha api tarhi na prāpnoti . \\
\noindent \textbf{(P\_7,2.107.1)} kārike , hārike , iti . \\
\noindent \textbf{(P\_7,2.107.1)} śībhāvaḥ ca prasajyate . \\
\noindent \textbf{(P\_7,2.107.1)} iha ca śībhāvaḥ ca prāpnoti . \\
\noindent \textbf{(P\_7,2.107.1)} asau brāhmaṇī . \\
\noindent \textbf{(P\_7,2.107.1)} āpaḥ uttarasya auṅaḥ śī bhavati iti śībhāvaḥ prāpnoti . \\
\noindent \textbf{(P\_7,2.107.1)} tasmāt soḥ lopaḥ vaktavyaḥ \\
\noindent \textbf{(P\_7,2.107.2)} sau autvapratiṣedhaḥ sākackāt vā sāt utvam ca . \\
\noindent \textbf{(P\_7,2.107.2)} sau autvapratiṣedhaḥ sākackāt vā vaktavyaḥ . \\
\noindent \textbf{(P\_7,2.107.2)} sāt ca parasya utvam vaktavyam . \\
\noindent \textbf{(P\_7,2.107.2)} asakau , asukaḥ . \\
\noindent \textbf{(P\_7,2.107.2)} uttarapadabhūtānām ādeśe upadeśavadvacanam . \\
\noindent \textbf{(P\_7,2.107.2)} uttarapadabhūtānām tyadādīnām ādeśe upadeśivadbhāvaḥ vaktavyaḥ . \\
\noindent \textbf{(P\_7,2.107.2)} paramāham , paramāyam , paramānena . \\
\noindent \textbf{(P\_7,2.107.2)} kim prayojanam . \\
\noindent \textbf{(P\_7,2.107.2)} anādiṣṭārtham . \\
\noindent \textbf{(P\_7,2.107.2)} akṛte ekādeśe ādeśāḥ yathā syuḥ iti . \\
\noindent \textbf{(P\_7,2.107.2)} kim punaḥ kāraṇam ekādeśaḥ tāvat bhavati na punaḥ ādeśāḥ . \\
\noindent \textbf{(P\_7,2.107.2)} na paratvāt ādeśaiḥ bhavitavyam . \\
\noindent \textbf{(P\_7,2.107.2)} bahiraṅgalakṣaṇatvāt . \\
\noindent \textbf{(P\_7,2.107.2)} bahiraṅgāḥ ādeśāḥ . \\
\noindent \textbf{(P\_7,2.107.2)} antaraṅgaḥ ekādeśaḥ . \\
\noindent \textbf{(P\_7,2.107.2)} asiddham bahiraṅgam antaraṅge . \\
\noindent \textbf{(P\_7,2.107.2)} saḥ tarhi upadeśivadbhāvaḥ vaktavyaḥ . \\
\noindent \textbf{(P\_7,2.107.2)} na vaktavyaḥ . \\
\noindent \textbf{(P\_7,2.107.2)} ācāryapravṛttiḥ jñāpayati pūrvapadottarapadayoḥ tāvat kāryam bhavati na ekādeśaḥ iti yat ayam na indrasya parasya iti pratiṣedham śāsti . \\
\noindent \textbf{(P\_7,2.107.2)} katham kṛtvā jñāpakam . \\
\noindent \textbf{(P\_7,2.107.2)} indre dvau acau . \\
\noindent \textbf{(P\_7,2.107.2)} tatra ekaḥ yasya iti lopena apahriyate aparaḥ ekādeśena . \\
\noindent \textbf{(P\_7,2.107.2)} anackaḥ indraḥ saṃvṛttaḥ . \\
\noindent \textbf{(P\_7,2.107.2)} tatra kaḥ vṛddheḥ prasaṅgaḥ . \\
\noindent \textbf{(P\_7,2.107.2)} paśyati tu ācāryaḥ pūrvapadottarapadayoḥ tāvat kāryam bhavati na ekādeśe iti tataḥ na indrasya parasya iti pratiṣedham śāsti \\
\noindent \textbf{(P\_7,2.107.3)} adasaḥ soḥ bhavet autvam kim sulopaḥ vidhīyate , hrasvāt lupyeta sambuddhiḥ na halaḥ prakṛtam hi tat . \\
\noindent \textbf{(P\_7,2.107.3)} āpaḥ ettvam bhavet tasmin na jhali iti anuvartanāt , pratyayasthāt ca kāt ittvam śībhāvaḥ ca prasajyate \\
\noindent \textbf{(P\_7,2.114)} mṛjeḥ vṛddhividhau kvipratiṣedhaḥ . \\
\noindent \textbf{(P\_7,2.114)} mṛjeḥ vṛddhividhau kvyantasya pratiṣedhaḥ vaktavyaḥ . \\
\noindent \textbf{(P\_7,2.114)} kaṃsaparimṛḍbhyām , kaṃsaparimṛḍbhiḥ . \\
\noindent \textbf{(P\_7,2.114)} dhātoḥ svarūpagrahaṇe vā tatpratyayavijñānāt siddham . \\
\noindent \textbf{(P\_7,2.114)} atha vā dhātoḥ svarūpagrahaṇe tatpratyaye kāryavijñānāt siddham . \\
\noindent \textbf{(P\_7,2.114)} dhātupratyaye kāryam bhavati iti eṣā paribhāṣā kartavyā . \\
\noindent \textbf{(P\_7,2.114)} kāni etasyāḥ paribhāṣāyāḥ prayojanāni . \\
\noindent \textbf{(P\_7,2.114)} prayojanam sṛjidṛśimasjinaśihantigiratyartham . \\
\noindent \textbf{(P\_7,2.114)} sṛji . \\
\noindent \textbf{(P\_7,2.114)} rajjusṛḍbhyām , rajjusṛḍbhiḥ . \\
\noindent \textbf{(P\_7,2.114)} sṛji . \\
\noindent \textbf{(P\_7,2.114)} dṛśi . \\
\noindent \textbf{(P\_7,2.114)} devadṛgbhyām , devadṛgbhiḥ . \\
\noindent \textbf{(P\_7,2.114)} dṛśi . \\
\noindent \textbf{(P\_7,2.114)} masji . \\
\noindent \textbf{(P\_7,2.114)} udakamagbhyām , udakamagbhiḥ . \\
\noindent \textbf{(P\_7,2.114)} masji . \\
\noindent \textbf{(P\_7,2.114)} naśi . \\
\noindent \textbf{(P\_7,2.114)} pranaḍbhyām , pranaḍbhiḥ . \\
\noindent \textbf{(P\_7,2.114)} naśi . \\
\noindent \textbf{(P\_7,2.114)} hanti . \\
\noindent \textbf{(P\_7,2.114)} vārtraghnaḥ , bhrauṇaghnaḥ . \\
\noindent \textbf{(P\_7,2.114)} hanti . \\
\noindent \textbf{(P\_7,2.114)} girati . \\
\noindent \textbf{(P\_7,2.114)} devagiraḥ . \\
\noindent \textbf{(P\_7,2.114)} yadi svarūpagrahaṇe iti ucyate prasṛbbhyām , prasṛbbhiḥ , anudāttasya ca ṛdupasya anyatarasyām iti am prāpnoti . \\
\noindent \textbf{(P\_7,2.114)} evam tarhi iyam paribhāṣā kartavyā dhātoḥ kāryam ucyamānam tatpratyaye bhavati iti . \\
\noindent \textbf{(P\_7,2.114)} sā tarhi eṣā paribhāṣā kartavyā . \\
\noindent \textbf{(P\_7,2.114)} na kartavyā . \\
\noindent \textbf{(P\_7,2.114)} ācāryapravṛttiḥ jñāpayati bhavati eṣā paribhāṣā iti yat ayam bhrauṇahatye tatvam śāsti \\
\noindent \textbf{(P\_7,2.115)} vṛddhau ajgrahaṇam go'rtham . \\
\noindent \textbf{(P\_7,2.115)} vṛddhau ajgrahaṇam kriyate gotaḥ vṛdhiḥ yathā syāt . \\
\noindent \textbf{(P\_7,2.115)} gauḥ iti . \\
\noindent \textbf{(P\_7,2.115)} na etat asti prayojanam . \\
\noindent \textbf{(P\_7,2.115)} ṇitkaraṇasāmarthyāt eva atra vṛddhiḥ bhaviṣyati . \\
\noindent \textbf{(P\_7,2.115)} atha yogavibhāgaḥ kimarthaḥ na ñṇiti ataḥ upadhāyāḥ iti eva ucyeta . \\
\noindent \textbf{(P\_7,2.115)} kā rūpasiddhaḥ cāyakaḥ , lāvakaḥ , kārakaḥ . \\
\noindent \textbf{(P\_7,2.115)} guṇe kṛte ayavaḥ raparatve ca ataḥ upadhāyāḥ iti eva siddham . \\
\noindent \textbf{(P\_7,2.115)} yogavibhāgaḥ sakhivyañjanādyarthaḥ . \\
\noindent \textbf{(P\_7,2.115)} yogavibhāgaḥ kriyate sakhyarthaḥ vyañjanādyarthaḥ ca . \\
\noindent \textbf{(P\_7,2.115)} sakhyarthaḥ tāvat . \\
\noindent \textbf{(P\_7,2.115)} sakhāyau , sakhāyaḥ . \\
\noindent \textbf{(P\_7,2.115)} vyañjanādyarthaḥ . \\
\noindent \textbf{(P\_7,2.115)} jaitram , yautram , cyautram . \\
\noindent \textbf{(P\_7,2.115)} yogavibhāge ca idānīm sakhivyañjanādyarthe kriyamāṇe ajgrahaṇam api kartavyam bhavati . \\
\noindent \textbf{(P\_7,2.115)} kim prayojanam . \\
\noindent \textbf{(P\_7,2.115)} gortham . \\
\noindent \textbf{(P\_7,2.115)} nanu ca uktam ṇitkaraṇasāmarthyāt eva atra vṛddhiḥ bhaviṣyati iti . \\
\noindent \textbf{(P\_7,2.115)} asti anyat ṇitkaraṇasaya prayojanam . \\
\noindent \textbf{(P\_7,2.115)} kim . \\
\noindent \textbf{(P\_7,2.115)} gāvau , gāvaḥ . \\
\noindent \textbf{(P\_7,2.115)} avādeśe kṛte ataḥ upadhāyāḥ iti vṛddhiḥ yathā syāt . \\
\noindent \textbf{(P\_7,2.115)} yat tu sau ṇitkaraṇam tat anavakāśam tasya anavakāśatvāt eva vṛddhiḥ bhaviṣyati . \\
\noindent \textbf{(P\_7,2.115)} yathā eva khalu api ṇitkaraṇasāmarthyāt anikaḥ api vṛddhiḥ prārthyate evam tatvam api prāpnoti . \\
\noindent \textbf{(P\_7,2.115)} tatvam api hi ñṇiti iti ucyate . \\
\noindent \textbf{(P\_7,2.115)} tasmāt ajgrahaṇam kartavyam \\
\noindent \textbf{(P\_7,2.117.1)} ajgrahaṇam kartavyam . \\
\noindent \textbf{(P\_7,2.117.1)} nanu ca kriyate eva . \\
\noindent \textbf{(P\_7,2.117.1)} dvitīyam kartavyam yathā acāmādigrahaṇam ajviśeṣaṇam vijñāyeta . \\
\noindent \textbf{(P\_7,2.117.1)} acām ādeḥ acaḥ iti . \\
\noindent \textbf{(P\_7,2.117.1)} atha akriyamāṇe ajgrahaṇe kasya acāmādigrahaṇam viśeṣaṇam syāt . \\
\noindent \textbf{(P\_7,2.117.1)} igviśeṣaṇam iti āha . \\
\noindent \textbf{(P\_7,2.117.1)} acām ādeḥ ikaḥ iti . \\
\noindent \textbf{(P\_7,2.117.1)} tatra kaḥ doṣaḥ . \\
\noindent \textbf{(P\_7,2.117.1)} iha eva syāt . \\
\noindent \textbf{(P\_7,2.117.1)} aitkāyanaḥ , aupagavaḥ . \\
\noindent \textbf{(P\_7,2.117.1)} iha na syāt . \\
\noindent \textbf{(P\_7,2.117.1)} gārgyaḥ , vātsyaḥ iti . \\
\noindent \textbf{(P\_7,2.117.1)} tat tarhi ajgrahaṇam kartavyam . \\
\noindent \textbf{(P\_7,2.117.1)} na kartavyam . \\
\noindent \textbf{(P\_7,2.117.1)} prakṛtam anuvartate . \\
\noindent \textbf{(P\_7,2.117.1)} kva prakṛtam . \\
\noindent \textbf{(P\_7,2.117.1)} acaḥ ñṇiti iti . \\
\noindent \textbf{(P\_7,2.117.1)} yadi tat anuvartate ataḥ upadhāyāḥ acaḥ iti ajmātrasya upadhāyāḥ vṛddhiḥ prasajyeta . \\
\noindent \textbf{(P\_7,2.117.1)} chedakaḥ iti . \\
\noindent \textbf{(P\_7,2.117.1)} akāreṇa tapareṇa acam viśeṣayiṣyāmaḥ . \\
\noindent \textbf{(P\_7,2.117.1)} acaḥ ataḥ iti . \\
\noindent \textbf{(P\_7,2.117.1)} iha idānīm acaḥ iti eva anuvartate ataḥ iti nivṛttam . \\
\noindent \textbf{(P\_7,2.117.1)} atha vā maṇḍūkagatayaḥ adhikārāḥ . \\
\noindent \textbf{(P\_7,2.117.1)} yathā maṇḍūkāḥ utplutya utplutya gacchanti tadvat adhikārāḥ . \\
\noindent \textbf{(P\_7,2.117.1)} atha vā ekayogaḥ kariṣyate . \\
\noindent \textbf{(P\_7,2.117.1)} acaḥ ñṇiti ataḥ upadhāyāḥ . \\
\noindent \textbf{(P\_7,2.117.1)} tataḥ taddhiteṣu acām ādeḥ iti . \\
\noindent \textbf{(P\_7,2.117.1)} na ca ekayoge anuvṛttiḥ bhavati \\
\noindent \textbf{(P\_7,2.117.2)} taddhiteṣu acāmādivṛddhau antyopadhalakṣaṇapratiṣedhaḥ . \\
\noindent \textbf{(P\_7,2.117.2)} taddhiteṣu acāmādivṛddhau antyopadhalakṣaṇāyāḥ vṛddheḥ pratiṣedhaḥ vaktavyaḥ . \\
\noindent \textbf{(P\_7,2.117.2)} krauṣṭuḥ jāgataḥ iti . \\
\noindent \textbf{(P\_7,2.117.2)} nanu ca acāmādivṛddhiḥ antyopadhalakṣaṇām vṛddhim bādhiṣyate . \\
\noindent \textbf{(P\_7,2.117.2)} katham anyasya ucyamānā anyasya bādhikā syāt . \\
\noindent \textbf{(P\_7,2.117.2)} asati khalu api sambhave bādhanam bhavati asti ca sambhavaḥ yat ubhayam syāt . \\
\noindent \textbf{(P\_7,2.117.2)} lokavijñānāt siddham . \\
\noindent \textbf{(P\_7,2.117.2)} sati api sambhave bādhanam bhavati . \\
\noindent \textbf{(P\_7,2.117.2)} tat yathā . \\
\noindent \textbf{(P\_7,2.117.2)} brāhmaṇebhyaḥ dadhi dīyatām takram kauṇḍinyāya iti sati api sambhave dadhidānasya takradānam nivartakam bhavati . \\
\noindent \textbf{(P\_7,2.117.2)} evam iha api sati api sambhave acāmādivṛddhiḥ antyopadhalakṣaṇām vṛddhim bādhiṣyate . \\
\noindent \textbf{(P\_7,2.117.2)} viṣamaḥ upanyāsaḥ . \\
\noindent \textbf{(P\_7,2.117.2)} na aprāpte dadhidāne takradānam ārabhyate tat prāpte ārabhyamāṇam bādhakam bhaviṣyati . \\
\noindent \textbf{(P\_7,2.117.2)} iha punaḥ aprāptāyām antyopadhalakṣaṇāyām vṛddhau acāmādivṛddhiḥ ārabhyate . \\
\noindent \textbf{(P\_7,2.117.2)} suśrut , sauśrutaḥ iti . \\
\noindent \textbf{(P\_7,2.117.2)} puṣkarasadgrahaṇāt vā . \\
\noindent \textbf{(P\_7,2.117.2)} atha vā yat ayam anuśatikādiṣu puṣkarasacśabdam paṭhati tat jñāpayati ācāryaḥ acāmādivṛddhau antyopadhalakṣaṇā vṛddhiḥ na bhavati iti \\
\noindent \textbf{(P\_7,3.1)} devikādiṣu tadādigrahaṇam . \\
\noindent \textbf{(P\_7,3.1)} devikādiṣu tadādigrahaṇam kartavyam . \\
\noindent \textbf{(P\_7,3.1)} devikādyādīnām iti vaktavyam . \\
\noindent \textbf{(P\_7,3.1)} iha api yathā syāt . \\
\noindent \textbf{(P\_7,3.1)} dāvikākulāḥ śālayaḥ , śāṃśapāsthalāḥ devāḥ . \\
\noindent \textbf{(P\_7,3.1)} kim punaḥ kāraṇam na sidhyati . \\
\noindent \textbf{(P\_7,3.1)} anyatra tadgrahaṇāt tadantagrahaṇāt vā . \\
\noindent \textbf{(P\_7,3.1)} anyatra hi tasya vā grahaṇam bhavati tadantasya vā na ca idam tat na api tadantam . \\
\noindent \textbf{(P\_7,3.1)} ādyajviśeṣaṇatvāt siddham . \\
\noindent \textbf{(P\_7,3.1)} ādyajviśeṣaṇam devikādayaḥ . \\
\noindent \textbf{(P\_7,3.1)} na evam vijñāyate devikādīnām aṅgānām acām ādeḥ ākāraḥ bhavati iti . \\
\noindent \textbf{(P\_7,3.1)} katham tarhi . \\
\noindent \textbf{(P\_7,3.1)} ñṇiti aṅgasya acām ādeḥ ākāraḥ bhavati saḥ cet devikādīnām ādjyac bhavati iti . \\
\noindent \textbf{(P\_7,3.1)} āntaratamyanivartakatvāt vā . \\
\noindent \textbf{(P\_7,3.1)} atha vā na anena anantaratamā vṛddhiḥ nirvartyate . \\
\noindent \textbf{(P\_7,3.1)} kim tarhi antaratamā anena nivartyate . \\
\noindent \textbf{(P\_7,3.1)} siddhā atra vṛddhiḥ taddhiteṣu acām ādeḥ iti eva tatra anena antaratamā vṛddhiḥ nivartyate . \\
\noindent \textbf{(P\_7,3.1)} parihārāntaram eva idam matvā paṭhitam katham ca idam parihārāntaram syāt . \\
\noindent \textbf{(P\_7,3.1)} yadi na ādyajviśeṣaṇam devikādayaḥ . \\
\noindent \textbf{(P\_7,3.1)} avaśyam ca etat evam vijñeyam adyajviśeṣaṇam devikādayaḥ iti . \\
\noindent \textbf{(P\_7,3.1)} yadi na ādyajviśeṣaṇam devikādayaḥ syuḥ iha api prāpnoti : sudevikāyām bhavaḥ saudevikaḥ iti . \\
\noindent \textbf{(P\_7,3.1)} atha atra api ādyajviśeṣaṇatvāt iti eva siddham parihārāntaram na bhavati . \\
\noindent \textbf{(P\_7,3.1)} na brūmaḥ yatra kriyamāṇe doṣaḥ tatra kartavyam iti . \\
\noindent \textbf{(P\_7,3.1)} kim tarhi . \\
\noindent \textbf{(P\_7,3.1)} yatra kriyamāṇe na doṣaḥ tatra kartavyam . \\
\noindent \textbf{(P\_7,3.1)} kva ca kriyamāṇe na doṣaḥ . \\
\noindent \textbf{(P\_7,3.1)} sañjñāvidhau . \\
\noindent \textbf{(P\_7,3.1)} vṛddhiḥ āt aic devikādīnām ākāraḥ iti . \\
\noindent \textbf{(P\_7,3.1)} idhyati . \\
\noindent \textbf{(P\_7,3.1)} sūtram tarhi bhidyate . \\
\noindent \textbf{(P\_7,3.1)} yathānyāsam eva astu . \\
\noindent \textbf{(P\_7,3.1)} nanu ca uktam devikādiṣu tadādigrahaṇam anyatra tadgrahaṇāt tadantagrahaṇāt vā iti . \\
\noindent \textbf{(P\_7,3.1)} parihṛtam etat ādyajviśeṣaṇatvāt siddham iti . \\
\noindent \textbf{(P\_7,3.1)} nyagrodhe ca kevalagrahaṇāt . \\
\noindent \textbf{(P\_7,3.1)} nyagrodhe ca kevalagrahaṇāt manyāmahe ādyajviśeṣaṇam devikādayaḥ iti . \\
\noindent \textbf{(P\_7,3.1)} tasya hi kevalagrahaṇasya etat prayojanam iha mā bhūt nyāgrodhamūlāḥ śālayaḥ iti . \\
\noindent \textbf{(P\_7,3.1)} yadi ca ādyajviśeṣaṇam devikādayaḥ tataḥ kevalagrahaṇam arthavat bhavati . \\
\noindent \textbf{(P\_7,3.1)} tat etat katham kṛtvā jñāpakam bhavati . \\
\noindent \textbf{(P\_7,3.1)} yadi nyagrodhaśabdaḥ avyutpannam prātipadikam bhavati . \\
\noindent \textbf{(P\_7,3.1)} atha hi nyagrohati iti nyagrodhaḥ tataḥ niyamārtham padāntaḥ iti kṛtvā na jñāpakam bhavati . \\
\noindent \textbf{(P\_7,3.1)} vahīnarasya idvacanam . \\
\noindent \textbf{(P\_7,3.1)} vahīnarasya ittvam vaktavyam . \\
\noindent \textbf{(P\_7,3.1)} vahīnarasya apatyam vaihīnariḥ . \\
\noindent \textbf{(P\_7,3.1)} kuṇaravāḍavaḥ tu āha . \\
\noindent \textbf{(P\_7,3.1)} na eṣaḥ vahīnaraḥ . \\
\noindent \textbf{(P\_7,3.1)} kaḥ tarhi . \\
\noindent \textbf{(P\_7,3.1)} vihīnaraḥ eṣaḥ . \\
\noindent \textbf{(P\_7,3.1)} vihīnaḥ naraḥ kāmabhogābhyām vihīnaraḥ . \\
\noindent \textbf{(P\_7,3.1)} vihīnarasya apatyam vaihīnariḥ \\
\noindent \textbf{(P\_7,3.3)} yvābhyām parasya avṛddhitvam . \\
\noindent \textbf{(P\_7,3.3)} yvābhyām parasya avṛddhitvam siddham . \\
\noindent \textbf{(P\_7,3.3)} kutaḥ . \\
\noindent \textbf{(P\_7,3.3)} apavādau vṛddheḥ hi tau . \\
\noindent \textbf{(P\_7,3.3)} apavādau hi vṛddheḥ tau aicau ucyete . \\
\noindent \textbf{(P\_7,3.3)} nityau aicau tayoḥ vṛddhiḥ . \\
\noindent \textbf{(P\_7,3.3)} atha vā nityau aicau . \\
\noindent \textbf{(P\_7,3.3)} kṛtāyām api vṛddhau prāpnutaḥ akṛtāyām api . \\
\noindent \textbf{(P\_7,3.3)} nityatvāt aicoḥ kṛtayoḥ yadi api vṛddhiḥ tayoḥ eva . \\
\noindent \textbf{(P\_7,3.3)} kimartham na iti śiṣyate . \\
\noindent \textbf{(P\_7,3.3)} atha kimartham pratiṣedhaḥ ucyate . \\
\noindent \textbf{(P\_7,3.3)} ecoḥ viṣayārtham pratiṣedhasanniyuktavacanam . \\
\noindent \textbf{(P\_7,3.3)} ecoḥ viṣayārtham pratiṣedhasanniyogena aicau ucyete . \\
\noindent \textbf{(P\_7,3.3)} yatra yvābhyām parāvṛddhiḥ tatra adhyaśveḥ yathā na tau . \\
\noindent \textbf{(P\_7,3.3)} yatra yvābhyām parasya avṛddhitvam ucyate tatra aicau yathā syātām . \\
\noindent \textbf{(P\_7,3.3)} iha mā bhūtām . \\
\noindent \textbf{(P\_7,3.3)} ādhyaśviḥ , dādhyaśviḥ , mādhvaśviḥ iti . \\
\noindent \textbf{(P\_7,3.3)} na etat asti prayojanam . \\
\noindent \textbf{(P\_7,3.3)} acām ādeḥ yvābhyām hi tau . \\
\noindent \textbf{(P\_7,3.3)} acām ādinā atra yvau viśeṣayiṣyāmaḥ . \\
\noindent \textbf{(P\_7,3.3)} acām ādeḥ yau yvau iti . \\
\noindent \textbf{(P\_7,3.3)} katham dvyāśītike na tau . \\
\noindent \textbf{(P\_7,3.3)} dvyāśītikaḥ iti atra kasmāt na tau bhavataḥ . \\
\noindent \textbf{(P\_7,3.3)} yatra vṛddhiḥ acām ādeḥ tatra aicau atra ghoḥ hi sā . \\
\noindent \textbf{(P\_7,3.3)} tatra acām ādeḥ iti evam vṛddhiḥ tatra aicau ucyete . \\
\noindent \textbf{(P\_7,3.3)} atra ghoḥ iti evam vṛddhiḥ . \\
\noindent \textbf{(P\_7,3.3)} kim idam ghoḥ iti . \\
\noindent \textbf{(P\_7,3.3)} uttarapadasya iti . \\
\noindent \textbf{(P\_7,3.3)} uttarapadādhikāre api avaśyam aijāgamaḥ anuvartyaḥ pūrvatryalinde bhavaḥ pūrvatrayalindaḥ iti evamartham . \\
\noindent \textbf{(P\_7,3.3)} na eṣaḥ doṣaḥ . \\
\noindent \textbf{(P\_7,3.3)} uttarapadena atra acām ādi viśeṣayiṣyāmaḥ acām ādinā yvau . \\
\noindent \textbf{(P\_7,3.3)} uttarapadasya acām ādeḥ yau yvau iti . \\
\noindent \textbf{(P\_7,3.3)} atha kasmāt padāntābhyām . \\
\noindent \textbf{(P\_7,3.3)} atha kimartham padāntābhyām iti ucyate . \\
\noindent \textbf{(P\_7,3.3)} yathā iṇaḥ na bhavet yaṇaḥ . \\
\noindent \textbf{(P\_7,3.3)} iṇaḥ yaṇādeśe mā bhūt. yataḥ chātrā , yātā iti . \\
\noindent \textbf{(P\_7,3.3)} iha vaiyākaraṇaḥ , sauvaśvaḥ iti śākalam prāpnoti yvoḥ ca sthānivadbhāvāt āyāvau prāpnutaḥ . \\
\noindent \textbf{(P\_7,3.3)} śakalāyāvādeśeṣu ca uktam . \\
\noindent \textbf{(P\_7,3.3)} kim uktam . \\
\noindent \textbf{(P\_7,3.3)} śākale tāvat uktam sinnityasamāsayoḥ śākalapratiṣedhaḥ iti . \\
\noindent \textbf{(P\_7,3.3)} āyāvoḥ kim uktam . \\
\noindent \textbf{(P\_7,3.3)} acaḥ pūrvavijñānāt aicoḥ siddham iti . \\
\noindent \textbf{(P\_7,3.3)} yvābhyām parasya avṛddhitvam apavādau vṛddheḥ hi tau , nityau aicau tayoḥ vṛddhiḥ kimartham na iti śiṣyate . \\
\noindent \textbf{(P\_7,3.3)} yatra yvābhyām parāvṛddhiḥ tatra adhyaśveḥ yathā na tau , acām ādeḥ yvābhyām hi tau katham dvyāśītike na tau . \\
\noindent \textbf{(P\_7,3.3)} yatra vṛddhiḥ acām ādeḥ tatra aicau atra ghoḥ hi sā , atha kasmāt padāntābhyām yathā iṇaḥ na bhavet yaṇaḥ \\
\noindent \textbf{(P\_7,3.4)} atha parasya avṛddhiḥ iti anuvartate utāho na . \\
\noindent \textbf{(P\_7,3.4)} kim ca ataḥ . \\
\noindent \textbf{(P\_7,3.4)} yadi anuvartate śauvam māṃsam ṭilope kṛte aijāgamaḥ na prāpnoti . \\
\noindent \textbf{(P\_7,3.4)} atha nivṛttam svādhyāyaśabdaḥ dvārādiṣu paṭhyate tatra yāvantaḥ yaṇaḥ sarvebhyaḥ pūrvaḥ aijāgamaḥ prāpnoti . \\
\noindent \textbf{(P\_7,3.4)} yatha icchasi tathā astu . \\
\noindent \textbf{(P\_7,3.4)} astu tāvat anuvartate . \\
\noindent \textbf{(P\_7,3.4)} katha śauvam māṃsam . \\
\noindent \textbf{(P\_7,3.4)} ānupūrvyā siddham etat . \\
\noindent \textbf{(P\_7,3.4)} na atra akṛte aijāgame ṭilopaḥ prāpnoti . \\
\noindent \textbf{(P\_7,3.4)} kim kāraṇam . \\
\noindent \textbf{(P\_7,3.4)} prakṛtyā ekāc iti prakṛtibhāvena bhavitavyam . \\
\noindent \textbf{(P\_7,3.4)} tat etat ānupūrvyā siddham bhavati . \\
\noindent \textbf{(P\_7,3.4)} atha vā punaḥ astu nivṛttam . \\
\noindent \textbf{(P\_7,3.4)} nanu ca uktam svādhyāyaśabdaḥ dvārādiṣu paṭhyate tatra yāvantaḥ yaṇaḥ sarvebhyaḥ pūrvaḥ aijāgamaḥ prāpnoti iti . \\
\noindent \textbf{(P\_7,3.4)} kaḥ punaḥ arhati svādhyāyaśabdam dvārādiṣu paṭhitum . \\
\noindent \textbf{(P\_7,3.4)} evam kil paṭhyeta svam adhyayanam svādhyāyaḥ iti . \\
\noindent \textbf{(P\_7,3.4)} tat ca na . \\
\noindent \textbf{(P\_7,3.4)} suṣṭhu vā adhyayanam svādhyāyaḥ śobhanam vā adhyayanam svādhyāyaḥ . \\
\noindent \textbf{(P\_7,3.4)} atha api svam adhyayanam svādhyāyaḥ evam api na doṣaḥ . \\
\noindent \textbf{(P\_7,3.4)} acām ādeḥ iti vartate . \\
\noindent \textbf{(P\_7,3.8)} ayam śvanśabdaḥ dvārādiṣu paṭhyate tatra kaḥ prasaṅgaḥ yat tadādeḥ syāt . \\
\noindent \textbf{(P\_7,3.8)} na eva prāpnoti na arthaḥ pratiṣedhena . \\
\noindent \textbf{(P\_7,3.8)} tadādividhinā prāpnoti . \\
\noindent \textbf{(P\_7,3.8)} na eva tadādividhiḥ asti . \\
\noindent \textbf{(P\_7,3.8)} ataḥ uttaram paṭhati . \\
\noindent \textbf{(P\_7,3.8)} pratiṣedhe śvādigrahaṇam jñāpakam anyatra śvangrahaṇe tadādigrahaṇasya śauvahānādyartham . \\
\noindent \textbf{(P\_7,3.8)} pratiṣedhe śvādigrahaṇam kriyate jñāpakārtham . \\
\noindent \textbf{(P\_7,3.8)} kim jñāpyam . \\
\noindent \textbf{(P\_7,3.8)} etat jñāpayati ācāryaḥ anyatra śvangrahaṇe tadādividhiḥ bhavati iti . \\
\noindent \textbf{(P\_7,3.8)} kim etasya jñāpane prayojanam . \\
\noindent \textbf{(P\_7,3.8)} śauvahānādyartham . \\
\noindent \textbf{(P\_7,3.8)} śauvahānam nāma nagaram . \\
\noindent \textbf{(P\_7,3.8)} śauvādaṃṣṭraḥ maṇiḥ iti . \\
\noindent \textbf{(P\_7,3.8)} ikārādigrahaṇam ca śvāgaṇikādyartham . \\
\noindent \textbf{(P\_7,3.8)} ikārādigrahaṇam ca kartavyam . \\
\noindent \textbf{(P\_7,3.8)} kim prayojanam . \\
\noindent \textbf{(P\_7,3.8)} śvāgaṇikādyartham . \\
\noindent \textbf{(P\_7,3.8)} śvagaṇena carati śvāgaṇikaḥ . \\
\noindent \textbf{(P\_7,3.8)} tadantasya ca anyatra pratiṣedhaḥ . \\
\noindent \textbf{(P\_7,3.8)} tadantasya ca anyatra pratiṣedhaḥ vaktavyaḥ . \\
\noindent \textbf{(P\_7,3.8)} śvābhastreḥ svam śvābhastram \\
\noindent \textbf{(P\_7,3.10)} kimartham idam ucyate . \\
\noindent \textbf{(P\_7,3.10)} avayavāt ṛtoḥ iti vakṣyati taduttarapadasya yathā syāt acām ādeḥ mā bhūt . \\
\noindent \textbf{(P\_7,3.10)} na etat asti prayojanam . \\
\noindent \textbf{(P\_7,3.10)} avayavāt iti pañcamī tatra antareṇa api uttarapadagrahaṇam uttarapadasya eva bhaviṣyati . \\
\noindent \textbf{(P\_7,3.10)} uttarārtham tarhi susarvārdhāt janapadasya iti . \\
\noindent \textbf{(P\_7,3.10)} susarvārdhāt iti pañcamī . \\
\noindent \textbf{(P\_7,3.10)} diśaḥ amadrāṇām . \\
\noindent \textbf{(P\_7,3.10)} diśaḥ iti pañcamī . \\
\noindent \textbf{(P\_7,3.10)} prācām grāmanagarāṇām . \\
\noindent \textbf{(P\_7,3.10)} diśaḥ iti eva . \\
\noindent \textbf{(P\_7,3.10)} saṅkhyāyāḥ saṃvatsarasaṅkhyasya ca . \\
\noindent \textbf{(P\_7,3.10)} saṅkhyāyāḥ iti pañcamī . \\
\noindent \textbf{(P\_7,3.10)} varṣasya abhaviṣyati . \\
\noindent \textbf{(P\_7,3.10)} saṅkhyāyāḥ iti eva . \\
\noindent \textbf{(P\_7,3.10)} parimāṇāntasya asañjñāśāṇayoḥ iti . \\
\noindent \textbf{(P\_7,3.10)} saṅkhyāyāḥ iti eva . \\
\noindent \textbf{(P\_7,3.10)} idam tarhi prayojanam je proṣṭhapadānām uttarapadasya yathā syāt pūrvapadasya mā bhūt . \\
\noindent \textbf{(P\_7,3.10)} proṣṭhapadāsu jātaḥ proṣṭhapādaḥ brāhmaṇaḥ . \\
\noindent \textbf{(P\_7,3.10)} taddhiteṣu acām ādivṛddheḥ uttarapadavṛddhiḥ vipratiṣedhena dvyāśītikādyartham . \\
\noindent \textbf{(P\_7,3.10)} taddhiteṣu acām ādivṛddheḥ uttarapadavṛddhiḥ bhavati vipratiṣedhena . \\
\noindent \textbf{(P\_7,3.10)} kim prayojanam . \\
\noindent \textbf{(P\_7,3.10)} dvyāśītikādyartham . \\
\noindent \textbf{(P\_7,3.10)} acām ādivṛddheḥ avakāśaḥ . \\
\noindent \textbf{(P\_7,3.10)} aitikāyanaḥ , aupagavaḥ . \\
\noindent \textbf{(P\_7,3.10)} uttarapadavṛddheḥ anavakāśaḥ . \\
\noindent \textbf{(P\_7,3.10)} dviṣāṣṭikaḥ , triṣāṣṭikaḥ . \\
\noindent \textbf{(P\_7,3.10)} iha ubhayam prāpnoti . \\
\noindent \textbf{(P\_7,3.10)} dvyāśītikaḥ , tryāśītikaḥ . \\
\noindent \textbf{(P\_7,3.10)} uttarapadavṛddhiḥ bhavati vipratiṣedhena . \\
\noindent \textbf{(P\_7,3.10)} kaḥ punaḥ atra viśeṣaḥ acām ādivṛddhau vā satyām uttarapadavṛddhau vā . \\
\noindent \textbf{(P\_7,3.10)} ayam asti viśeṣaḥ . \\
\noindent \textbf{(P\_7,3.10)} yadi atra acām ādivṛddhiḥ syāt aijāgamaḥ prasajyeta \\
\noindent \textbf{(P\_7,3.14)} nagaragrahaṇam kimartham na prācām grāmāṇām iti eva siddham . \\
\noindent \textbf{(P\_7,3.14)} na sidhyati . \\
\noindent \textbf{(P\_7,3.14)} anyaḥ grāmaḥ anyat nagaram . \\
\noindent \textbf{(P\_7,3.14)} katham jñāyate . \\
\noindent \textbf{(P\_7,3.14)} evam hi kaḥ cit kam cit pṛcchati . \\
\noindent \textbf{(P\_7,3.14)} kutaḥ bhavān āgacchati grāmāt . \\
\noindent \textbf{(P\_7,3.14)} saḥ hi āha . \\
\noindent \textbf{(P\_7,3.14)} na grāmāt nagarāt iti . \\
\noindent \textbf{(P\_7,3.14)} nanu ca bho yaḥ eva grāmaḥ tat nagaram . \\
\noindent \textbf{(P\_7,3.14)} katham jñāyate . \\
\noindent \textbf{(P\_7,3.14)} lokataḥ . \\
\noindent \textbf{(P\_7,3.14)} ye hi grāme vidhayaḥ na iṣyante sādhīyaḥ te nagare na kriyante . \\
\noindent \textbf{(P\_7,3.14)} tat yathā . \\
\noindent \textbf{(P\_7,3.14)} abhakṣyaḥ grāmyakukkuṭaḥ abhakṣyaḥ grāmyaśūkaraḥ iti ukte sutarām nāgaraḥ api na bhakṣyate . \\
\noindent \textbf{(P\_7,3.14)} tathā grāme na adhyeyam iti sādhīyaḥ nagare na adhīyate . \\
\noindent \textbf{(P\_7,3.14)} tasmāt yaḥ eva grāmaḥ tat nagaram . \\
\noindent \textbf{(P\_7,3.14)} katham yat uktam evam hi kaḥ cit kam cit pṛcchati kutaḥ bhavān āgacchati grāmāt saḥ āha na grāmāt nagarāt iti . \\
\noindent \textbf{(P\_7,3.14)} saṃstyāyaviśeṣam asau ācaṣṭe . \\
\noindent \textbf{(P\_7,3.14)} saṃstyāyaviśeṣāḥ hi ete grāmaḥ ghoṣaḥ nagaram saṃvāhaḥ iti . \\
\noindent \textbf{(P\_7,3.14)} evam tarhi siddhe sati yat grāmagrahaṇe nagaragrahaṇam karoti tat jñāpayati ācāryaḥ anyatra grāmagrahaṇe nagaragrahaṇam na bhavati iti . \\
\noindent \textbf{(P\_7,3.14)} kim etasya jñāpane prayojanam . \\
\noindent \textbf{(P\_7,3.14)} viśiṣṭaliṅgaḥ nadīdeśaḥ agrāmāḥ iti atra nagarapratiṣedhaḥ coditaḥ saḥ na vaktavyaḥ bhavati . \\
\noindent \textbf{(P\_7,3.14)} yadi etat jñāpyate udīcyagrāmāt ca bahvacaḥ antodāttāt iti atra nagaragrahaṇam kartavyam . \\
\noindent \textbf{(P\_7,3.14)} bāhīkagrāmebhyaḥ ca nagaragrahaṇam kartavyam . \\
\noindent \textbf{(P\_7,3.14)} dikśabdāḥ grāmajanapadākhyānacānarāṭeṣu nagaragrahaṇam kartavyam . \\
\noindent \textbf{(P\_7,3.14)} idam caturtham jñāpakārtham . \\
\noindent \textbf{(P\_7,3.14)} tatra atinirbandhaḥ na lābhaḥ . \\
\noindent \textbf{(P\_7,3.14)} tasmāt yasmin eva grāmagrahaṇe nagaragrahaṇam na iṣyate tasya pratiṣedhaḥ vaktavyaḥ \\
\noindent \textbf{(P\_7,3.15)} saṃvatsaragrahaṇam anarthakam parimāṇāntasya iti kṛtatvāt . \\
\noindent \textbf{(P\_7,3.15)} saṃvatsaragrahaṇam anarthakam . \\
\noindent \textbf{(P\_7,3.15)} kim kāraṇam . \\
\noindent \textbf{(P\_7,3.15)} parimāṇāntasya iti kṛtatvāt . \\
\noindent \textbf{(P\_7,3.15)} parimāṇāntasya asañjñāśāṇayoḥ iti eva siddham . \\
\noindent \textbf{(P\_7,3.15)} jñāpakam tu kālaparimāṇānām vṛddhipratiṣedhasya . \\
\noindent \textbf{(P\_7,3.15)} evam tarhi jñāpayati ācāryaḥ kālaparimāṇānām vṛddhiḥ na bhavati iti . \\
\noindent \textbf{(P\_7,3.15)} kim etasya jñāpane prayojanam . \\
\noindent \textbf{(P\_7,3.15)} dvairātrikaḥ , trairātrikaḥ , atra vṛddhiḥ na bhavati . \\
\noindent \textbf{(P\_7,3.15)} na etat asti prayojanam . \\
\noindent \textbf{(P\_7,3.15)} na asti atra viśeṣaḥ satyām vā uttarapadavṛddhau asatyām vā . \\
\noindent \textbf{(P\_7,3.15)} idam tarhi . \\
\noindent \textbf{(P\_7,3.15)} dvasamikaḥ , traisamikaḥ . \\
\noindent \textbf{(P\_7,3.15)} idam ca api prayojanam dvairātrikaḥ , trairātrikaḥ . \\
\noindent \textbf{(P\_7,3.15)} nanu ca uktam na asti atra viśeṣaḥ satyām vā uttarapadavṛddhau asatyām vā iti . \\
\noindent \textbf{(P\_7,3.15)} ayam asti viśeṣaḥ . \\
\noindent \textbf{(P\_7,3.15)} yadi atra uttarapadavṛddhiḥ syāt acām ādeḥ vṛddhiḥ na syāt . \\
\noindent \textbf{(P\_7,3.15)} aparaḥ āha : jñāpakam tu kālaparimāṇānām parimāṇāgrahaṇasya . \\
\noindent \textbf{(P\_7,3.15)} evam tarhi jñāpayati ācāryaḥ kālaparimāṇānām parimāṇagrahaṇena grahaṇam na bhavati iti . \\
\noindent \textbf{(P\_7,3.15)} kim etasya jñāpane prayojanam . \\
\noindent \textbf{(P\_7,3.15)} aparimāṇabistācitakambalyebhyaḥ na taddhitaluki dvivarṣā , trivarṣā . \\
\noindent \textbf{(P\_7,3.15)} parimāṇaparyudāsena paryudāsaḥ na bhavati \\
\noindent \textbf{(P\_7,3.28)} parasya vṛddhiḥ na iti anuvartate utāho na . \\
\noindent \textbf{(P\_7,3.28)} kim ca ataḥ . \\
\noindent \textbf{(P\_7,3.28)} yadi anuvartate pravāhaṇeyī bhāryā asya iti pravāhaṇeyībhāryaḥ vṛddhinimittasya iti puṃvadbhāvapratiṣedhaḥ na prāpnoti . \\
\noindent \textbf{(P\_7,3.28)} atha nivṛttam na doṣaḥ bhavati . \\
\noindent \textbf{(P\_7,3.28)} yathā na doṣaḥ tathā astu . \\
\noindent \textbf{(P\_7,3.28)} atha vā punaḥ astu anuvartate . \\
\noindent \textbf{(P\_7,3.28)} nanu ca uktam pravāhaṇeyī bhāryā asya pravāhaṇeyībhāryaḥ vṛddhinimittasya iti puṃvadbhāvapratiṣedhaḥ na prāpnoti iti . \\
\noindent \textbf{(P\_7,3.28)} na eṣaḥ doṣaḥ . \\
\noindent \textbf{(P\_7,3.28)} mā bhūt evam . \\
\noindent \textbf{(P\_7,3.28)} jāteḥ iti evam bhaviṣyati \\
\noindent \textbf{(P\_7,3.31)} ayam yogaḥ śakyaḥ avaktum . \\
\noindent \textbf{(P\_7,3.31)} katham ayāthātathyam , āyathātathyam , ayāthāpuryam , āyathāpuryam . \\
\noindent \textbf{(P\_7,3.31)} yadā tāvat pūrvapadasya vṛddhiḥ tadā evam vigrahaḥ kariṣyate . \\
\noindent \textbf{(P\_7,3.31)} na yathātathā , ayathātathā . \\
\noindent \textbf{(P\_7,3.31)} ayathātathābhāvaḥ āyathātathyam . \\
\noindent \textbf{(P\_7,3.31)} yadā uttarapadasya vṛddhiḥ tadā evam vigrahaḥ kariṣyate . \\
\noindent \textbf{(P\_7,3.31)} yathātathābhāvaḥ yāthātathyam . \\
\noindent \textbf{(P\_7,3.31)} na yāthātathyam ayāthātathyam \\
\noindent \textbf{(P\_7,3.32)} hanteḥ takāre taddhite pratiṣedhaḥ . \\
\noindent \textbf{(P\_7,3.32)} hanteḥ takāre taddhite pratiṣedhaḥ vaktavyaḥ . \\
\noindent \textbf{(P\_7,3.32)} vārtraghnam , bhrauṇaghnam . \\
\noindent \textbf{(P\_7,3.32)} uktam vā . \\
\noindent \textbf{(P\_7,3.32)} kim uktam . \\
\noindent \textbf{(P\_7,3.32)} dhātoḥ svarūpagrahaṇe tatpratyaye kāryavijñānāt siddham iti \\
\noindent \textbf{(P\_7,3.33)} kṛdgrahaṇam kimartham . \\
\noindent \textbf{(P\_7,3.33)} iha mā bhūt . \\
\noindent \textbf{(P\_7,3.33)} dadau , dadhau . \\
\noindent \textbf{(P\_7,3.33)} na etat asti prayojanam . \\
\noindent \textbf{(P\_7,3.33)} aciṇṇaloḥ iti vartate . \\
\noindent \textbf{(P\_7,3.33)} yadi aciṇṇaloḥ iti vartate adāyi , adhāyi iti atra na prāpnoti . \\
\noindent \textbf{(P\_7,3.33)} vacanāt ciṇi bhaviṣyati . \\
\noindent \textbf{(P\_7,3.33)} aciṇṇaloḥ iti vartate . \\
\noindent \textbf{(P\_7,3.33)} evam api cauḍiḥ , bālākiḥ iti atra prāpnoti . \\
\noindent \textbf{(P\_7,3.33)} lopaḥ atra bādhakaḥ bhaviṣyati . \\
\noindent \textbf{(P\_7,3.33)} idam iha sampradhāryam . \\
\noindent \textbf{(P\_7,3.33)} lopaḥ kriyatām yuk iti kim atra kartavyam . \\
\noindent \textbf{(P\_7,3.33)} paratvāt yuk . \\
\noindent \textbf{(P\_7,3.33)} evam tarhi acām ādeḥ iti vartate . \\
\noindent \textbf{(P\_7,3.33)} yatra acām ādiḥ ākāraḥ tatra yuk iti . \\
\noindent \textbf{(P\_7,3.33)} evam api jñā devatā asya sthālīpākasya jñaḥ sthālīpākaḥ , atra prāpnoti . \\
\noindent \textbf{(P\_7,3.33)} tasmāt kṛdgrahaṇam kartavyam \\
\noindent \textbf{(P\_7,3.34)} atyalpam idam ucyate : anācameḥ iti . \\
\noindent \textbf{(P\_7,3.34)} avamikamicamīnām iti vaktavyam : vāmaḥ , kāmaḥ , ācāmaḥ \\
\noindent \textbf{(P\_7,3.37)} ṇicprakaraṇe dhūñprīñoḥ nugvacanam . \\
\noindent \textbf{(P\_7,3.37)} ṇicprakaraṇe dhūñprīñoḥ nuk vaktavyaḥ . \\
\noindent \textbf{(P\_7,3.37)} dhūnayati , prīṇayati . \\
\noindent \textbf{(P\_7,3.37)} pāteḥ lugvacanam . \\
\noindent \textbf{(P\_7,3.37)} pālayati \\
\noindent \textbf{(P\_7,3.44.1)} sthagrahaṇam kimartham . \\
\noindent \textbf{(P\_7,3.44.1)} idam vicārayiṣyate ittve kagrahaṇam saṅghātagrahaṇam vā syāt varṇagrahaṇam vā iti . \\
\noindent \textbf{(P\_7,3.44.1)} tat yadā saṅghatagrahaṇam tadā sthagrahaṇam kartavyam iha api yathā syāt kārikā , hārikā . \\
\noindent \textbf{(P\_7,3.44.1)} yadā hi varṇagrahaṇam tadā kevalaḥ kakāraḥ pratyayaḥ na asti iti kṛtvā vacanāt bhaviṣyati . \\
\noindent \textbf{(P\_7,3.44.1)} atha asupaḥ iti katham idam vijñāyate . \\
\noindent \textbf{(P\_7,3.44.1)} asubvataḥ aṅgasya iti . \\
\noindent \textbf{(P\_7,3.44.1)} āhosvit na cet supaḥ paraḥ āp iti . \\
\noindent \textbf{(P\_7,3.44.1)} kim ca ataḥ . \\
\noindent \textbf{(P\_7,3.44.1)} yadi vijñāyate asubvataḥ aṅgasya iti bahucarmikā atra na prāpnoti . \\
\noindent \textbf{(P\_7,3.44.1)} atha vijñāyate na cet supaḥ paraḥ āp iti na doṣaḥ bhavati . \\
\noindent \textbf{(P\_7,3.44.1)} yathā na doṣaḥ tathā astu \\
\noindent \textbf{(P\_7,3.44.2)} idam vicāryate : ittve kagrahaṇam saṅghātagrahaṇam vā syāt varṇagrahaṇam vā iti . \\
\noindent \textbf{(P\_7,3.44.2)} kaḥ ca atra viśeṣaḥ . \\
\noindent \textbf{(P\_7,3.44.2)} ittve kagrahaṇam saṅghātagrahaṇam cet etikāsu aprāptiḥ . \\
\noindent \textbf{(P\_7,3.44.2)} ittve kagrahaṇam saṅghātagrahaṇam cet etikāsvu aprāptiḥ . \\
\noindent \textbf{(P\_7,3.44.2)} etikāḥ caranti . \\
\noindent \textbf{(P\_7,3.44.2)} vacanāt bhaviṣyati . \\
\noindent \textbf{(P\_7,3.44.2)} asti vacane prayojanam . \\
\noindent \textbf{(P\_7,3.44.2)} kim . \\
\noindent \textbf{(P\_7,3.44.2)} kārikā , hārikā . \\
\noindent \textbf{(P\_7,3.44.2)} astu tarhi varṇagrahaṇam . \\
\noindent \textbf{(P\_7,3.44.2)} varṇagrahaṇam cet vyavahitatvāt aprasiddhiḥ . \\
\noindent \textbf{(P\_7,3.44.2)} varṇagrahaṇam cet vyavahitatvāt na prāpnoti . \\
\noindent \textbf{(P\_7,3.44.2)} kārikā , hārikā . \\
\noindent \textbf{(P\_7,3.44.2)} akāreṇa vyavahitatvāt na prāpnoti . \\
\noindent \textbf{(P\_7,3.44.2)} ekādeśe kṛte na asti vyavadhānam . \\
\noindent \textbf{(P\_7,3.44.2)} ekādeśaḥ pūrvavidhau sthānivat bhavati iti sthānivadbhāvāt vyavadhānam eva . \\
\noindent \textbf{(P\_7,3.44.2)} evam tarhi āha ayam pratyayasthāt kāt pūrvasya iti na kva cit avyavadhānam tatra vacanāt bhaviṣyati . \\
\noindent \textbf{(P\_7,3.44.2)} vacanaprāmāṇyāt iti cet rathakaṭyādiṣu atiprasaṅgaḥ . \\
\noindent \textbf{(P\_7,3.44.2)} vacanaprāmāṇyāt iti cet rathakaṭyādiṣu doṣaḥ bhavati . \\
\noindent \textbf{(P\_7,3.44.2)} rathakaṭyā , gargakāmyā . \\
\noindent \textbf{(P\_7,3.44.2)} na eṣaḥ doṣaḥ . \\
\noindent \textbf{(P\_7,3.44.2)} yena na avyavadhānam tena vyavahite api vacanaprāmāṇyāt . \\
\noindent \textbf{(P\_7,3.44.2)} kena ca na avyavadhānam varṇena ekena . \\
\noindent \textbf{(P\_7,3.44.2)} saṅghātena punaḥ vyavadhānam bhavati na bhavati ca . \\
\noindent \textbf{(P\_7,3.44.2)} atha vā punaḥ astu saṅghātagrahaṇam . \\
\noindent \textbf{(P\_7,3.44.2)} nanu ca uktam ittve kagrahaṇam saṅghātagrahaṇam cet etikāsu aprāptiḥ iti . \\
\noindent \textbf{(P\_7,3.44.2)} parihṛtam etat vacanāt bhaviṣyati iti . \\
\noindent \textbf{(P\_7,3.44.2)} nanu ca uktam asti vacane prayojanam . \\
\noindent \textbf{(P\_7,3.44.2)} kim . \\
\noindent \textbf{(P\_7,3.44.2)} kārikā , hārikā iti . \\
\noindent \textbf{(P\_7,3.44.2)} atra api ekādeśe kṛte vyapavargābhāvāt na prāpnoti . \\
\noindent \textbf{(P\_7,3.44.2)} antādivadbhāvena vyapavargaḥ . \\
\noindent \textbf{(P\_7,3.44.2)} ubhayataḥ āśraye na antādivat . \\
\noindent \textbf{(P\_7,3.44.2)} evam tarhi ekādeśaḥ pūrvavidhau sthānivat bhavati iti sthānivadbhāvāt vyapavargaḥ . \\
\noindent \textbf{(P\_7,3.44.2)} evam tarhi ācāryapravṛttiḥ jñāpayati bhavati evañjātīyakānām api ittvam iti yat ayam na yāsayoḥ iti pratiṣedham śāsti \\
\noindent \textbf{(P\_7,3.44.3)} mamaka(R: māmaka)narakayoḥ upasaṅkhyānam apratyayasthatvāt . \\
\noindent \textbf{(P\_7,3.44.3)} mamaka(R: māmaka)narakayoḥ upasaṅkhyānam kartavyam . \\
\noindent \textbf{(P\_7,3.44.3)} māmikā , narikā . \\
\noindent \textbf{(P\_7,3.44.3)} kim punaḥ kāraṇam na sidhyati . \\
\noindent \textbf{(P\_7,3.44.3)} apratyayasthatvāt . \\
\noindent \textbf{(P\_7,3.44.3)} tyaktyapoḥ ca pratiṣiddhatvāt . \\
\noindent \textbf{(P\_7,3.44.3)} tyaktyapoḥ ca upasaṅkhyānam kartavyam . \\
\noindent \textbf{(P\_7,3.44.3)} dākṣiṇātyikā , amātyikā . \\
\noindent \textbf{(P\_7,3.44.3)} kim punaḥ kāraṇam na sidhyati . \\
\noindent \textbf{(P\_7,3.44.3)} pratiṣiddhatvāt . \\
\noindent \textbf{(P\_7,3.44.3)} udīcām ātaḥ sthāne yakapūrvāyāḥ iti pratiṣiddhatvāt \\
\noindent \textbf{(P\_7,3.45.)} na yattadoḥ iti vaktavyam . \\
\noindent \textbf{(P\_7,3.45.)} iha api yathā syāt . \\
\noindent \textbf{(P\_7,3.45.)} yakām yakām adhīte , takām takām pacāmahe iti . \\
\noindent \textbf{(P\_7,3.45.)} pratiṣedhe tyakanaḥ upasaṅkhyānam . \\
\noindent \textbf{(P\_7,3.45.)} pratiṣedhe tyakanaḥ upasaṅkhyānam kartavyam . \\
\noindent \textbf{(P\_7,3.45.)} upatyakā , adhityakā . \\
\noindent \textbf{(P\_7,3.45.)} tat tarhi upasaṅkhyānam kartavyam . \\
\noindent \textbf{(P\_7,3.45.)} na kartavyam . \\
\noindent \textbf{(P\_7,3.45.)} ācāryapravṛttiḥ jñāpayati na evañjātīyakānām ittvam bhavati iti yat ayam mṛdaḥ tikan iti ittvabhūtam nirdeśam karoti . \\
\noindent \textbf{(P\_7,3.45.)} pāvakādīnām chandasi upasaṅkhyānam . \\
\noindent \textbf{(P\_7,3.45.)} pāvakādīnām chandasi upasaṅkhyānam kartavyam . \\
\noindent \textbf{(P\_7,3.45.)} hiraṇyavarṇāḥ śrucayaḥ pāvakāḥ , ṛkṣakāḥ , alomakāḥ . \\
\noindent \textbf{(P\_7,3.45.)} chandasi iti kimartham . \\
\noindent \textbf{(P\_7,3.45.)} pāvikā , alomikā . \\
\noindent \textbf{(P\_7,3.45.)} āśiṣi ca . \\
\noindent \textbf{(P\_7,3.45.)} āśiṣi ca upasaṅkhyānam kartavyam . \\
\noindent \textbf{(P\_7,3.45.)} jīvatāt jīvakā , nandatāt nandakā , bhavatāt bhavakā . \\
\noindent \textbf{(P\_7,3.45.)} uttarapadalope ca . \\
\noindent \textbf{(P\_7,3.45.)} uttarapadalope ca upasaṅkhyānam kartavyam . \\
\noindent \textbf{(P\_7,3.45.)} devadattikā , devakā , yajñadattikā , yajñakā . \\
\noindent \textbf{(P\_7,3.45.)} kṣipakādīnām ca . \\
\noindent \textbf{(P\_7,3.45.)} kṣipakādīnām ca upasaṅkhyānam kartavyam . \\
\noindent \textbf{(P\_7,3.45.)} kṣipakā , dhruvakā , dhuvakā . \\
\noindent \textbf{(P\_7,3.45.)} tārakā jyotiṣi . \\
\noindent \textbf{(P\_7,3.45.)} tārakā jyotiṣi upasaṅkhyānam kartavyam . \\
\noindent \textbf{(P\_7,3.45.)} tārakā . \\
\noindent \textbf{(P\_7,3.45.)} jyotiṣi iti kimartham . \\
\noindent \textbf{(P\_7,3.45.)} tārikā dāsī . \\
\noindent \textbf{(P\_7,3.45.)} varṇakā tānave . \\
\noindent \textbf{(P\_7,3.45.)} varṇakā tāntave upasaṅkhyānam kartavyam . \\
\noindent \textbf{(P\_7,3.45.)} varṇakā . \\
\noindent \textbf{(P\_7,3.45.)} tāntave iti kimartham . \\
\noindent \textbf{(P\_7,3.45.)} varṇikā bhāgurī lokāyatasya . \\
\noindent \textbf{(P\_7,3.45.)} vartakā śakunau prācām . \\
\noindent \textbf{(P\_7,3.45.)} vartakā śakunau prācām upasaṅkhyānam kartavyam . \\
\noindent \textbf{(P\_7,3.45.)} vartakā śakuniḥ . \\
\noindent \textbf{(P\_7,3.45.)} śakunau iti kimartham . \\
\noindent \textbf{(P\_7,3.45.)} vartikā bhāgurī lokāyatasya . \\
\noindent \textbf{(P\_7,3.45.)} prācām iti kimartham . \\
\noindent \textbf{(P\_7,3.45.)} vartikā . \\
\noindent \textbf{(P\_7,3.45.)} aṣṭakā pitṛdevatye . \\
\noindent \textbf{(P\_7,3.45.)} aṣṭakā pitṛdevatye upasaṅkhyānam kartavyam . \\
\noindent \textbf{(P\_7,3.45.)} aṣṭakā . \\
\noindent \textbf{(P\_7,3.45.)} pitṛdevatye iti kimartham . \\
\noindent \textbf{(P\_7,3.45.)} aṣṭikā khārī . \\
\noindent \textbf{(P\_7,3.45.)} vā sūtakāputrakāvṛndārakāṇām . \\
\noindent \textbf{(P\_7,3.45.)} vā sūtakāputrakāvṛndārakāṇām upasaṅkhyānam kartavyam . \\
\noindent \textbf{(P\_7,3.45.)} sūtakā , sūtikā , putrakā , putrikā , vṛndārakā , vṛndārikā \\
\noindent \textbf{(P\_7,3.46)} kimartham strīliṅganirdeśaḥ kriyate na yakapūrvasya iti eva ucyeta . \\
\noindent \textbf{(P\_7,3.46)} strīviṣayaḥ yaḥ ākāraḥ tasya sthāne yaḥ akāraḥ tasya pratiṣedhaḥ yathā syāt . \\
\noindent \textbf{(P\_7,3.46)} iha mā bhūt . \\
\noindent \textbf{(P\_7,3.46)} śubham yāti iti śubhaṃyāḥ śubhaṃyikā , bhadraṃyikā . \\
\noindent \textbf{(P\_7,3.46)} yakapūrve dhātvantapratiṣedhaḥ . \\
\noindent \textbf{(P\_7,3.46)} yakapūrve dhātvantapratiṣedhaḥ vaktavyaḥ . \\
\noindent \textbf{(P\_7,3.46)} kim prayojanam . \\
\noindent \textbf{(P\_7,3.46)} sunayikā , aśokikā , apākikā \\
\noindent \textbf{(P\_7,3.47)} eṣādve nañpūrve anudāharaṇe asupaḥ iti pratiṣedhāt . \\
\noindent \textbf{(P\_7,3.47)} atha bhastrāgrahaṇam kimartham na abhāṣitapuṃskāt iti eva siddham . \\
\noindent \textbf{(P\_7,3.47)} bhastrāgrahaṇam upasarjanārtham . \\
\noindent \textbf{(P\_7,3.47)} upasarjanārthaḥ ayam ārambhaḥ . \\
\noindent \textbf{(P\_7,3.47)} abhastrikā , abhastrakā . \\
\noindent \textbf{(P\_7,3.47)} nañpūrvagrahaṇānarthakyam ca uttarapadamātrasya idvacanāt . \\
\noindent \textbf{(P\_7,3.47)} nañpūrvagrahaṇam ca anarthakam . \\
\noindent \textbf{(P\_7,3.47)} kim kāraṇam . \\
\noindent \textbf{(P\_7,3.47)} uttarapadamātrasya idvacanāt . \\
\noindent \textbf{(P\_7,3.47)} uttarapadamātrasya ittvam vaktavyam . \\
\noindent \textbf{(P\_7,3.47)} nirbhastrakā , nirbhastrikā , bahubhastrakā , bahubhastrikā \\
\noindent \textbf{(P\_7,3.50)} kim idam ṭhādeśe varṇagrahaṇam āhosvit saṅghātagrahaṇam . \\
\noindent \textbf{(P\_7,3.50)} kaḥ ca atra viśeṣaḥ . \\
\noindent \textbf{(P\_7,3.50)} ṭhādeśe varṇagrahaṇam cet dhātvantasya pratiṣedhaḥ . \\
\noindent \textbf{(P\_7,3.50)} ṭhādeśe varṇagrahaṇam cet dhātvantasya pratiṣedhaḥ vaktavyaḥ . \\
\noindent \textbf{(P\_7,3.50)} paṭhitā , paṭhitum . \\
\noindent \textbf{(P\_7,3.50)} astu tarhi saṅghātagrahaṇam . \\
\noindent \textbf{(P\_7,3.50)} saṅghātagrahaṇam cet aṇādimāthitikādīnām pratiṣedhaḥ . \\
\noindent \textbf{(P\_7,3.50)} saṅghātagrahaṇam cet uṇādimāthitikādīnām pratiṣedhaḥ vaktavyaḥ . \\
\noindent \textbf{(P\_7,3.50)} uṇādīnām tāvat . \\
\noindent \textbf{(P\_7,3.50)} kaṇṭhaḥ , vaṇṭhaḥ , śaṇṭhaḥ . \\
\noindent \textbf{(P\_7,3.50)} iha ca mathitam paṇyam asya māthitikaḥ iti akāralope kṛte tāntāt iti kādeśaḥ prāpnoti . \\
\noindent \textbf{(P\_7,3.50)} varṅagrahaṇe punaḥ sati alvidhiḥ ayam bhavati . \\
\noindent \textbf{(P\_7,3.50)} tasmāt viśiṣṭagrahaṇam . \\
\noindent \textbf{(P\_7,3.50)} tasmāt viśiṣṭasya ṭhakārasya grahaṇam kartavyam . \\
\noindent \textbf{(P\_7,3.50)} na kartavyam . \\
\noindent \textbf{(P\_7,3.50)} astu tāvat varṇagrahaṇam . \\
\noindent \textbf{(P\_7,3.50)} nanu ca uktam ṭhādeśe varṇagrahaṇam cet dhātvantasya pratiṣedhaḥ iti . \\
\noindent \textbf{(P\_7,3.50)} na eṣaḥ doṣaḥ . \\
\noindent \textbf{(P\_7,3.50)} aṅgāt iti vartate . \\
\noindent \textbf{(P\_7,3.50)} na vā aṅgāt iti pañcamī asti . \\
\noindent \textbf{(P\_7,3.50)} evam tarhi pratyayasthasya iti vartate . \\
\noindent \textbf{(P\_7,3.50)} kva prakṛtam . \\
\noindent \textbf{(P\_7,3.50)} pratyayasthāt kāt pūrvasya ataḥ it āpi asupaḥ iti . \\
\noindent \textbf{(P\_7,3.50)} tat vai pañcamīnirdiṣṭam ṣaṣṭhīnirdiṣṭena ca iha arthaḥ . \\
\noindent \textbf{(P\_7,3.50)} arthāt vibhaktivipariṇāmaḥ bhaviṣyati . \\
\noindent \textbf{(P\_7,3.50)} tat yathā . \\
\noindent \textbf{(P\_7,3.50)} uccāni devadattasya gṛhāṇi . \\
\noindent \textbf{(P\_7,3.50)} āmantrayasva enam . \\
\noindent \textbf{(P\_7,3.50)} devadattam iti gamyate . \\
\noindent \textbf{(P\_7,3.50)} devadattasya gāvaḥ aśvāḥ hiraṇyam . \\
\noindent \textbf{(P\_7,3.50)} āḍhyaḥ vaidhaveyaḥ . \\
\noindent \textbf{(P\_7,3.50)} devadattaḥ iti gamyate . \\
\noindent \textbf{(P\_7,3.50)} purastāt ṣaṣṭhīnirdiṣṭam sat arthāt prathamānirdiṣṭam dvitīyānirdiṣṭam ca bhavati . \\
\noindent \textbf{(P\_7,3.50)} evam iha api purastāt pañcamīnirdiṣṭam sat arthāt ṣaṣṭhīnirdiṣṭam bhaviṣyati . \\
\noindent \textbf{(P\_7,3.50)} evam api uṇādīnām pratiṣedhaḥ vaktavyaḥ . \\
\noindent \textbf{(P\_7,3.50)} na vaktavyaḥ . \\
\noindent \textbf{(P\_7,3.50)} uṇādayaḥ avyutpannāni prātipadikāni . \\
\noindent \textbf{(P\_7,3.50)} evam api karmaṭhaḥ iti atra prāpnoti . \\
\noindent \textbf{(P\_7,3.50)} evam tarhi aṅgasya iti sambandhaṣaṣṭhī vijñāsyate . \\
\noindent \textbf{(P\_7,3.50)} aṅgasya yaḥ ṭhakāraḥ . \\
\noindent \textbf{(P\_7,3.50)} kim ca aṅgasya ṭhakāraḥ . \\
\noindent \textbf{(P\_7,3.50)} nimittam . \\
\noindent \textbf{(P\_7,3.50)} yasmin aṅgam iti etat bhavati . \\
\noindent \textbf{(P\_7,3.50)} kasmin ca etat bhavati . \\
\noindent \textbf{(P\_7,3.50)} pratyaye . \\
\noindent \textbf{(P\_7,3.50)} atha vā punaḥ astu saṅghātagrahaṇam . \\
\noindent \textbf{(P\_7,3.50)} nanu ca uktam saṅghātagrahaṇam cet uṇādimāthitikādīnām pratiṣedhaḥ iti uṇādīnām tāvat pratiṣedhaḥ na vaktavyaḥ . \\
\noindent \textbf{(P\_7,3.50)} parihṛtam etat uṇādayaḥ avyutpannāni prātipadikāni iti . \\
\noindent \textbf{(P\_7,3.50)} yat api ucyate iha ca mathitam paṇyam asya māthitikaḥ iti akāralope kṛte tāntāt iti kādeśaḥ prāpnoti iti . \\
\noindent \textbf{(P\_7,3.50)} na eṣaḥ doṣaḥ . \\
\noindent \textbf{(P\_7,3.50)} akāralopasya sthānivadbhāvāt na bhaviṣyati . \\
\noindent \textbf{(P\_7,3.50)} na sidhyati . \\
\noindent \textbf{(P\_7,3.50)} pūrvavidhau sthānivadbhāvaḥ na ca ayam pūrvavidhiḥ . \\
\noindent \textbf{(P\_7,3.50)} ayam api pūrvavidhiḥ . \\
\noindent \textbf{(P\_7,3.50)} pūrvasmāt api vidhiḥ pūrvavidhiḥ iti . \\
\noindent \textbf{(P\_7,3.50)} atha api uṇādayaḥ vyutpādyante evam api na doṣaḥ . \\
\noindent \textbf{(P\_7,3.50)} kriyate nyāse eva viśiṣṭagrahaṇam ṭhasya iti \\
\noindent \textbf{(P\_7,3.51)} iha kasmāt na bhavati . \\
\noindent \textbf{(P\_7,3.51)} āśiṣā tarati āśiṣikaḥ , uṣā tarati auṣikaḥ . \\
\noindent \textbf{(P\_7,3.51)} lakṣaṇapratipadoktayoḥ pratipadoktasya eva iti . \\
\noindent \textbf{(P\_7,3.51)} atha iha katham bhavitavyam . \\
\noindent \textbf{(P\_7,3.51)} dorbhyām tarati . \\
\noindent \textbf{(P\_7,3.51)} dauṣkaḥ iti bhavitavyam . \\
\noindent \textbf{(P\_7,3.51)} katham . \\
\noindent \textbf{(P\_7,3.51)} yadi varṇaikadeśāḥ varṇagrahaṇena gṛhyante \\
\noindent \textbf{(P\_7,3.54)} kim idam ñṇinnakāragrahaṇam hantiviśeṣaṇam : ñṇinnakāraparasya hanteḥ yaḥ hakāraḥ iti . \\
\noindent \textbf{(P\_7,3.54)} āhosvit hakāraviśeṣaṇam : ñṇinnakāraparasya hakārasya saḥ cet hanteḥ iti . \\
\noindent \textbf{(P\_7,3.54)} kaḥ ca atra viśeṣaḥ . \\
\noindent \textbf{(P\_7,3.54)} hanteḥ tatparasya iti cet nakāre aprasiddhiḥ . \\
\noindent \textbf{(P\_7,3.54)} hanteḥ tatparasya iti cet nakāre aprasiddhiḥ . \\
\noindent \textbf{(P\_7,3.54)} ghnanti , ghnantu , aghnan . \\
\noindent \textbf{(P\_7,3.54)} astu tarhi hakāraviśeṣaṇam . \\
\noindent \textbf{(P\_7,3.54)} hakārasya iti cet ñṇiti aprāptiḥ . \\
\noindent \textbf{(P\_7,3.54)} hakārasya iti cet ñṇiti aprāptiḥ . \\
\noindent \textbf{(P\_7,3.54)} ghātayati ghātakaḥ . \\
\noindent \textbf{(P\_7,3.54)} kim kāraṇam . \\
\noindent \textbf{(P\_7,3.54)} nakāreṇa vyavahitatvāt na prāpnoti . \\
\noindent \textbf{(P\_7,3.54)} vacanāt bhaviṣyati . \\
\noindent \textbf{(P\_7,3.54)} iha api vacanāt prāpnoti . \\
\noindent \textbf{(P\_7,3.54)} hananam icchati hananīyate hananīyateḥ ṇvul hananīyakaḥ iti . \\
\noindent \textbf{(P\_7,3.54)} sthānivadbhāvāt ca acaḥ nakāre aprasiddhiḥ . \\
\noindent \textbf{(P\_7,3.54)} sthānivadbhāvāt ca acaḥ nakāre aprasiddhiḥ . \\
\noindent \textbf{(P\_7,3.54)} ghnanti , ghnantu . \\
\noindent \textbf{(P\_7,3.54)} vacanāt bhaviṣyati . \\
\noindent \textbf{(P\_7,3.54)} vacanaprāmāṇyāt iti cet alope pratiṣedhaḥ . \\
\noindent \textbf{(P\_7,3.54)} vacanaprāmāṇyāt iti cet alope pratiṣedhaḥ vaktavyaḥ . \\
\noindent \textbf{(P\_7,3.54)} hantā , hantum . \\
\noindent \textbf{(P\_7,3.54)} nakāragrahaṇasāmarthyāt alope na bhaviṣyati . \\
\noindent \textbf{(P\_7,3.54)} asti anyat nakāragrahaṇasya prayojanam . \\
\noindent \textbf{(P\_7,3.54)} kim . \\
\noindent \textbf{(P\_7,3.54)} śrūyamāṇaviśeṣaṇam . \\
\noindent \textbf{(P\_7,3.54)} yatra nakāraḥ śrūyate tatra yathā syāt . \\
\noindent \textbf{(P\_7,3.54)} iha mā bhūt . \\
\noindent \textbf{(P\_7,3.54)} hataḥ hathaḥ iti . \\
\noindent \textbf{(P\_7,3.54)} siddham tu upadhālope iti vacanāt . \\
\noindent \textbf{(P\_7,3.54)} siddham etat . \\
\noindent \textbf{(P\_7,3.54)} katham . \\
\noindent \textbf{(P\_7,3.54)} upadhālope ca iti vaktavyam . \\
\noindent \textbf{(P\_7,3.54)} sidhyati . \\
\noindent \textbf{(P\_7,3.54)} sūtram tarhi bhidyate . \\
\noindent \textbf{(P\_7,3.54)} yathānyāsam eva astu . \\
\noindent \textbf{(P\_7,3.54)} nanu ca uktam hanteḥ tatparasya iti cet nakāre aprasiddhiḥ iti . \\
\noindent \textbf{(P\_7,3.54)} vacanāt bhaviṣyati . \\
\noindent \textbf{(P\_7,3.54)} atha vā punaḥ astu hakāraviśeṣaṇam . \\
\noindent \textbf{(P\_7,3.54)} nanu ca uktam hakārasya iti cet ñṇiti aprāptiḥ iti . \\
\noindent \textbf{(P\_7,3.54)} vacanāt bhaviṣyati . \\
\noindent \textbf{(P\_7,3.54)} nanu ca uktam iha api vacanāt prāpnoti hananīyakaḥ iti . \\
\noindent \textbf{(P\_7,3.54)} na eṣaḥ doṣaḥ . \\
\noindent \textbf{(P\_7,3.54)} yena na avyavadhānam tena vyavahite api vacanaprāmāṇyāt . \\
\noindent \textbf{(P\_7,3.54)} na ca kva cit dhātvavayavena avyavadhānam etena punaḥ saṅghātena vyavadhānam bhavati na ca bhavati . \\
\noindent \textbf{(P\_7,3.54)} yat api ucyate sthānivadbhāvāt ca acaḥ nakāre aprasiddhiḥ iti vacanāt bhaviṣyati . \\
\noindent \textbf{(P\_7,3.54)} nanu ca uktam vacanaprāmāṇyāt iti cet alope pratiṣedhaḥ iti . \\
\noindent \textbf{(P\_7,3.54)} na eṣaḥ doṣaḥ . \\
\noindent \textbf{(P\_7,3.54)} ānantaryam iha āśrīyate hakārasya nakāraḥ iti . \\
\noindent \textbf{(P\_7,3.54)} kva cit ca sannipātakṛtam ānantaryam śāstrakṛtam anānantaryam kva cit ca na sannipātakṛtam na api śāstrakṛtam . \\
\noindent \textbf{(P\_7,3.54)} lope sannipātakṛtam ānantaryam alope na eva sannipātakṛtam na api śāstrakṛtam . \\
\noindent \textbf{(P\_7,3.54)} yatra kutaḥ cit eva ānantaryam tat āśrayiṣyāmaḥ \\
\noindent \textbf{(P\_7,3.55)} abhyāsāt kutvam asupaḥ . \\
\noindent \textbf{(P\_7,3.55)} abhyāsāt kutvam asupaḥ iti vaktavyam . \\
\noindent \textbf{(P\_7,3.55)} iha mā bhūt . \\
\noindent \textbf{(P\_7,3.55)} hananam icchati hananīyati hananīyateḥ san jihananīyiṣati iti . \\
\noindent \textbf{(P\_7,3.55)} tat tarhi vaktavyam . \\
\noindent \textbf{(P\_7,3.55)} na vaktavyam . \\
\noindent \textbf{(P\_7,3.55)} hanteḥ abhyāsāt iti ucyate na ca eṣaḥ hanteḥ abhyāsaḥ . \\
\noindent \textbf{(P\_7,3.55)} hanteḥ eṣaḥ abhyāsaḥ . \\
\noindent \textbf{(P\_7,3.55)} katham . \\
\noindent \textbf{(P\_7,3.55)} ekācaḥ dve prathamasya iti . \\
\noindent \textbf{(P\_7,3.55)} evam tarhi hanteḥ aṅgasya yaḥ abhyāsaḥ tasmāt iti ucyate na ca eṣaḥ hanteḥ aṅgasya abhyāsaḥ . \\
\noindent \textbf{(P\_7,3.55)} hanteḥ aṅgasya eṣaḥ abhyāsaḥ . \\
\noindent \textbf{(P\_7,3.55)} katham . \\
\noindent \textbf{(P\_7,3.55)} ekācaḥ dve prathamasya iti . \\
\noindent \textbf{(P\_7,3.55)} evam tarhi yasmin hantiḥ aṅgam tasmin yaḥ abhyāsaḥ tasmāt iti ucyate . \\
\noindent \textbf{(P\_7,3.55)} yasmin ca atra hantiḥ aṅgam na tasmin abhyāsaḥ yasmin ca abhyāsaḥ na tasmin hantiḥ aṅgam bhavati \\
\noindent \textbf{(P\_7,3.56)} acaṅi iti kimartham . \\
\noindent \textbf{(P\_7,3.56)} prājīhayat dūtam . \\
\noindent \textbf{(P\_7,3.56)} heḥ caṅi pratiṣedhānarthakyam aṅgānyatvāt . \\
\noindent \textbf{(P\_7,3.56)} heḥ caṅi pratiṣedhaḥ anarthakaḥ . \\
\noindent \textbf{(P\_7,3.56)} kim kāraṇam . \\
\noindent \textbf{(P\_7,3.56)} aṅgānyatvāt . \\
\noindent \textbf{(P\_7,3.56)} ṇyantam etat aṅgam anyat bhavati . \\
\noindent \textbf{(P\_7,3.56)} lope kṛte na aṅgānyatvam . \\
\noindent \textbf{(P\_7,3.56)} sthānivadbhāvāt aṅgānyatvam eva . \\
\noindent \textbf{(P\_7,3.56)} jñāpakam tu anytra ṇyadhikasya kutvavijñānārtham . \\
\noindent \textbf{(P\_7,3.56)} evam tarhi jñāpayati ācāryaḥ anyatra ṇyadhikasya kutvam bhavati iti . \\
\noindent \textbf{(P\_7,3.56)} kim etasya jñāpane prayojanam . \\
\noindent \textbf{(P\_7,3.56)} prajighāyayiṣati iti atra kutvam siddham bhavati \\
\noindent \textbf{(P\_7,3.57)} jigrahaṇe jyaḥ pratiṣedhaḥ . \\
\noindent \textbf{(P\_7,3.57)} jigrahaṇe jyaḥ pratiṣedhaḥ vaktavyaḥ . \\
\noindent \textbf{(P\_7,3.57)} jijyatuḥ , jijyuḥ iti . \\
\noindent \textbf{(P\_7,3.57)} saḥ tarhi pratiṣedhaḥ vaktavyaḥ . \\
\noindent \textbf{(P\_7,3.57)} na vaktavyaḥ . \\
\noindent \textbf{(P\_7,3.57)} lakṣaṇapratipadoktayoḥ pratipadoktasya eva iti evam etasya na bhaviṣyati . \\
\noindent \textbf{(P\_7,3.57)} sā tarhi eṣā paribhaṣā kartavyā . \\
\noindent \textbf{(P\_7,3.57)} avaśyam kartavyā adhyāpya gataḥ iti evamartham \\
\noindent \textbf{(P\_7,3.59)} kvādyajivrajiyācirucīnām apratiṣedhaḥ niṣṭhāyām aniṭaḥ kutvavacanāt . \\
\noindent \textbf{(P\_7,3.59)} kvādyajivrajiyācirucīnām apratiṣedhaḥ . \\
\noindent \textbf{(P\_7,3.59)} anarthakaḥ pratiṣedhaḥ apratiṣedhaḥ . \\
\noindent \textbf{(P\_7,3.59)} kutvam kasmāt na bhavati . \\
\noindent \textbf{(P\_7,3.59)} niṣṭhāyām aniṭaḥ kutvavacanāt . \\
\noindent \textbf{(P\_7,3.59)} niṣṭhāyām aniṭaḥ kutvam vakṣyāmi seṭaḥ ca ete niṣṭhāyām . \\
\noindent \textbf{(P\_7,3.59)} yadi niṣṭhāyām aniṭaḥ kutvam ucyate katham śokaḥ samudraḥ iti . \\
\noindent \textbf{(P\_7,3.59)} śucyubjyoḥ ghañi kutvam . \\
\noindent \textbf{(P\_7,3.59)} śucyubjyoḥ ghañi kutvam vaktavyam . \\
\noindent \textbf{(P\_7,3.59)} katham arkaḥ . \\
\noindent \textbf{(P\_7,3.59)} arceḥ kavidhānāt siddham . \\
\noindent \textbf{(P\_7,3.59)} na etat ghañantam . \\
\noindent \textbf{(P\_7,3.59)} auṇādikaḥ eṣaḥ kaśabdaḥ tasmin āṣṭamikam kutvam \\
\noindent \textbf{(P\_7,3.61)} bhujaḥ pāṇau . \\
\noindent \textbf{(P\_7,3.61)} bhujaḥ pāṇau iti vaktavyam . \\
\noindent \textbf{(P\_7,3.61)} katham nyubjaḥ upatāpe iti . \\
\noindent \textbf{(P\_7,3.61)} nyubjeḥ kartṛtvāt apratiṣedhaḥ . \\
\noindent \textbf{(P\_7,3.61)} anarthakaḥ pratiṣedhaḥ apratiṣedhaḥ . \\
\noindent \textbf{(P\_7,3.61)} kutvam kasmat na bhavati . \\
\noindent \textbf{(P\_7,3.61)} kartṛtvāt . \\
\noindent \textbf{(P\_7,3.61)} na etat ghañantam . \\
\noindent \textbf{(P\_7,3.61)} kartṛpratyayaḥ eṣaḥ . \\
\noindent \textbf{(P\_7,3.61)} nyubjati iti nyubjaḥ . \\
\noindent \textbf{(P\_7,3.61)} adhikaraṇasādhanaḥ vai lakṣyate ghañ . \\
\noindent \textbf{(P\_7,3.61)} nyubjitāḥ śerate asmin nyubjaḥ upatāpe iti . \\
\noindent \textbf{(P\_7,3.61)} eṣaḥ api hi kartṛsādhanaḥ eva . \\
\noindent \textbf{(P\_7,3.61)} nyubjayati iti nyubjaḥ \\
\noindent \textbf{(P\_7,3.66)} pravacigrahaṇam anarthakam vacaḥ aśabdasañjñābhāvāt . \\
\noindent \textbf{(P\_7,3.66)} pravacigrahaṇam anarthakam . \\
\noindent \textbf{(P\_7,3.66)} kim kāraṇam . \\
\noindent \textbf{(P\_7,3.66)} vaco'śabdasañjñābhāvāt . \\
\noindent \textbf{(P\_7,3.66)} vaco'śabdasañjñāyām pratiṣedhaḥ ucyate prapūrvaḥ ca vaciḥ aśabdasañjñāyām vartate . \\
\noindent \textbf{(P\_7,3.66)} upasarganiyamārtham tarhi idam vaktavyam . \\
\noindent \textbf{(P\_7,3.66)} prapūrvasya eva vaceḥ aśabdasañjñāyām pratiṣedhaḥ yathā syāt . \\
\noindent \textbf{(P\_7,3.66)} iha mā bhūt . \\
\noindent \textbf{(P\_7,3.66)} avivākyam iti . \\
\noindent \textbf{(P\_7,3.66)} upasargapūrvaniyamārtham iti cet avivākyasya viśeṣavacanat siddham . \\
\noindent \textbf{(P\_7,3.66)} viśeṣe etat vaktavyam . \\
\noindent \textbf{(P\_7,3.66)} avivākyam ahaḥ iti . \\
\noindent \textbf{(P\_7,3.66)} kva mā bhūt . \\
\noindent \textbf{(P\_7,3.66)} avivācyam eva anyat iti . \\
\noindent \textbf{(P\_7,3.66)} ṇyapratiṣedhe tyajeḥ upasaṅkhyānam . \\
\noindent \textbf{(P\_7,3.66)} ṇyapratiṣedhe tyajeḥ upasaṅkhyānam kartavyam . \\
\noindent \textbf{(P\_7,3.66)} tyājyam \\
\noindent \textbf{(P\_7,3.69)} [bhojyam abhyavahārhye] . \\
\noindent \textbf{(P\_7,3.69)} bhojyam abhyavahārye iti vaktavyam . \\
\noindent \textbf{(P\_7,3.69)} iha api yathā syāt . \\
\noindent \textbf{(P\_7,3.69)} bhojyaḥ sūpaḥ , bhojyā yavāgūḥ iti . \\
\noindent \textbf{(P\_7,3.69)} kim punaḥ kāraṇam na sidhyati . \\
\noindent \textbf{(P\_7,3.69)} bhakṣiḥ ayam kharaviśade vartate tena drave na prāpnoti . \\
\noindent \textbf{(P\_7,3.69)} na avaśyam bhakṣiḥ kharaviśade eva vartate . \\
\noindent \textbf{(P\_7,3.69)} kim tarhi anyatra api vartate . \\
\noindent \textbf{(P\_7,3.69)} tat yathā . \\
\noindent \textbf{(P\_7,3.69)} abbhakṣaḥ , vāyubhakṣaḥ iti \\
\noindent \textbf{(P\_7,3.70)} vā iti śakyam avaktum . \\
\noindent \textbf{(P\_7,3.70)} kasmāt na bhavati . \\
\noindent \textbf{(P\_7,3.70)} tat agniḥ agnaye dadāt . \\
\noindent \textbf{(P\_7,3.70)} astu atra lopaḥ āṭaḥ śravaṇam bhaviṣyati tena ubhayam sidhyati . \\
\noindent \textbf{(P\_7,3.70)} dadhat ratnāni dāśuṣe , dadāt ratnāni dāśuṣe \\
\noindent \textbf{(P\_7,3.71)} [otaḥ śiti] . \\
\noindent \textbf{(P\_7,3.71)} otaḥ śiti iti vaktavyam . \\
\noindent \textbf{(P\_7,3.71)} kim prayojanam . \\
\noindent \textbf{(P\_7,3.71)} uttaratra śidgrahaṇābhāvāya . \\
\noindent \textbf{(P\_7,3.71)} tatra ayam api arthaḥ ṣṭhivuklamvācamām śiti iti śidgrahaṇam na kartavyam bhavati . \\
\noindent \textbf{(P\_7,3.71)} nanu ca bhoḥ śyangrahaṇam api tarhi uttarārtham kartavyam . \\
\noindent \textbf{(P\_7,3.71)} śamām aṣṭānām dīrghaḥ śyani iti śyangrahaṇam na kartavyam bhavati . \\
\noindent \textbf{(P\_7,3.71)} atra api astu śiti iti eva . \\
\noindent \textbf{(P\_7,3.71)} yadi śiti iti ucyate anu tvā indraḥ bhramatu madatu atra api prāpnoti . \\
\noindent \textbf{(P\_7,3.71)} śamādibhiḥ atra śitam viśeṣayiṣyāmaḥ . \\
\noindent \textbf{(P\_7,3.71)} śamādīnām yaḥ śit iti . \\
\noindent \textbf{(P\_7,3.71)} kaḥ ca śamādīnām śit . \\
\noindent \textbf{(P\_7,3.71)} śamādibhyaḥ yaḥ vihitaḥ . \\
\noindent \textbf{(P\_7,3.71)} evam api tasyati , yasyati atra prāpnoti . \\
\noindent \textbf{(P\_7,3.71)} aṣṭānām iti vacanāt na bhaviṣyati \\
\noindent \textbf{(P\_7,3.75)} dīrghatvam āṅi camaḥ . \\
\noindent \textbf{(P\_7,3.75)} dīrghatvam āṅi camaḥ iti vaktavyam . \\
\noindent \textbf{(P\_7,3.75)} ācāmati . \\
\noindent \textbf{(P\_7,3.75)} iha mā bhūt . \\
\noindent \textbf{(P\_7,3.75)} uccamati , vicamati iti \\
\noindent \textbf{(P\_7,3.77)} iṣeḥ chatvam ahali . \\
\noindent \textbf{(P\_7,3.77)} iṣeḥ chatvam ahali iti vaktavyam . \\
\noindent \textbf{(P\_7,3.77)} iha mā bhūt . \\
\noindent \textbf{(P\_7,3.77)} iṣṇāti , iṣyati . \\
\noindent \textbf{(P\_7,3.77)} tat tarhi vaktavyam . \\
\noindent \textbf{(P\_7,3.77)} na vaktavyam . \\
\noindent \textbf{(P\_7,3.77)} aci iti vartate . \\
\noindent \textbf{(P\_7,3.77)} evam api iṣāṇa iti atra prāpnoti . \\
\noindent \textbf{(P\_7,3.77)} atha ahali iti ucyamāne kasmāt eva atra chatvam na bhavati . \\
\noindent \textbf{(P\_7,3.77)} na evam vijñāyate na hal ahal ahali iti . \\
\noindent \textbf{(P\_7,3.77)} katham tarhi . \\
\noindent \textbf{(P\_7,3.77)} avidyamānaḥ hal asmin saḥ ayam ahal ahali iti . \\
\noindent \textbf{(P\_7,3.77)} yadi evam aci iti api vartamāne na doṣaḥ . \\
\noindent \textbf{(P\_7,3.77)} na hi acā śit viśeṣyate . \\
\noindent \textbf{(P\_7,3.77)} śiti bhavati katarsmin aci iti . \\
\noindent \textbf{(P\_7,3.77)} katham tarhi . \\
\noindent \textbf{(P\_7,3.77)} śitā ac viśeṣyate . \\
\noindent \textbf{(P\_7,3.77)} aci bhavati katarsmin śiti iti \\
\noindent \textbf{(P\_7,3.78)} pibeḥ guṇapratiṣedhaḥ . \\
\noindent \textbf{(P\_7,3.78)} pibeḥ guṇapratiṣedhaḥ vaktavyaḥ . \\
\noindent \textbf{(P\_7,3.78)} pibati . \\
\noindent \textbf{(P\_7,3.78)} laghūpadhaguṇaḥ prāpnoti . \\
\noindent \textbf{(P\_7,3.78)} saḥ tarhi pratiṣedhaḥ vaktavyaḥ . \\
\noindent \textbf{(P\_7,3.78)} na vaktavyaḥ . \\
\noindent \textbf{(P\_7,3.78)} guṇaḥ kasmāt na bhavati . \\
\noindent \textbf{(P\_7,3.78)} pibiḥ adantaḥ . \\
\noindent \textbf{(P\_7,3.78)} adante iti cet uktam . \\
\noindent \textbf{(P\_7,3.78)} kim uktam . \\
\noindent \textbf{(P\_7,3.78)} dhātoḥ ante iti cet anudāttecabagrahaṇam iti . \\
\noindent \textbf{(P\_7,3.78)} atha vā aṅgavṛtte punarvṛttau avidhiḥ niṣṭhitasya iti evam na bhaviṣyati \\
\noindent \textbf{(P\_7,3.79)} dīrghoccāraṇam kimartham na jñājanoḥ jaḥ iti eva ucyeta . \\
\noindent \textbf{(P\_7,3.79)} kā rūpasiddhiḥ : jānāti , jāyate . \\
\noindent \textbf{(P\_7,3.79)} ataḥ dīrghaḥ yañi iti dīrghatvam bhaviṣyati . \\
\noindent \textbf{(P\_7,3.79)} evam tarhi siddhe sati yat dīrghoccāraṇam karoti tat jñāpayati ācāryaḥ bhavati eṣā paribhāṣā aṅgavṛtte punarvṛttau avidhiḥ iti . \\
\noindent \textbf{(P\_7,3.79)} kim etasya jñāpane prayojanam . \\
\noindent \textbf{(P\_7,3.79)} pibeḥ guṇapratiṣedhaḥ coditaḥ saḥ na vaktavyaḥ bhavati \\
\noindent \textbf{(P\_7,3.83)} jusi guṇe yāsuṭpratiṣedhaḥ . \\
\noindent \textbf{(P\_7,3.83)} jusi guṇe yāsuḍādau pratiṣedhaḥ vaktavyaḥ . \\
\noindent \textbf{(P\_7,3.83)} cinuyuḥ , sunuyuḥ iti . \\
\noindent \textbf{(P\_7,3.83)} na vaktavyaḥ . \\
\noindent \textbf{(P\_7,3.83)} na evam vijñāyate mideḥ guṇaḥ jusi ca iti . \\
\noindent \textbf{(P\_7,3.83)} katham tarhi . \\
\noindent \textbf{(P\_7,3.83)} mideḥ guṇaḥ ajusi ca iti . \\
\noindent \textbf{(P\_7,3.83)} kim idam ajusi iti . \\
\noindent \textbf{(P\_7,3.83)} ajādau usi ajusi iti . \\
\noindent \textbf{(P\_7,3.83)} iha api tarhi prāpnoti . \\
\noindent \textbf{(P\_7,3.83)} cakruḥ , jahruḥ iti . \\
\noindent \textbf{(P\_7,3.83)} evam tarhi śiti iti vartate . \\
\noindent \textbf{(P\_7,3.83)} evam api ajuhavuḥ , abibhayuḥ iti atra na prāpnoti . \\
\noindent \textbf{(P\_7,3.83)} bhūtapūrvagatyā bhaviṣyati . \\
\noindent \textbf{(P\_7,3.83)} na sidhyati na hi us śidbhūtapūrvaḥ . \\
\noindent \textbf{(P\_7,3.83)} us śidbhūtapūrvaḥ na asti iti kṛtvā usi yaḥ śidbhūtapūrvaḥ tasmin bhaviṣyati . \\
\noindent \textbf{(P\_7,3.83)} atha vā kriyate nyāse eva . \\
\noindent \textbf{(P\_7,3.83)} avibhaktikaḥ nirdeśaḥ . \\
\noindent \textbf{(P\_7,3.83)} na evam vijñāyate mideḥ guṇaḥ jusi ca iti . \\
\noindent \textbf{(P\_7,3.83)} katham tarhi . \\
\noindent \textbf{(P\_7,3.83)} mideḥ guṇaḥ u jusi iti . \\
\noindent \textbf{(P\_7,3.83)} kim idam u jusi iti . \\
\noindent \textbf{(P\_7,3.83)} ukārādau jusi . \\
\noindent \textbf{(P\_7,3.83)} atha vā aci iti vartate tena jusam viśeṣayiṣyāmaḥ . \\
\noindent \textbf{(P\_7,3.83)} ajādau jusi iti \\
\noindent \textbf{(P\_7,3.85)} iha jāgarayati , jāgarakaḥ iti guṇe kṛte raparatve ca ataḥ upadhāyāḥ iti vṛddhiḥ prāpnoti tasyāḥ pratiṣedhaḥ vaktavyaḥ . \\
\noindent \textbf{(P\_7,3.85)} ciṇṇaloḥ pratiṣedhasāmarthyāt anyatra guṇabhūtasya vṛddhipratiṣedhaḥ . \\
\noindent \textbf{(P\_7,3.85)} yat ayam aciṇṇaloḥ iti pratiṣedham śāsti tat jñāpayati ācāryaḥ na guṇābhinirvṛttasya vṛddhiḥ bhavati iti . \\
\noindent \textbf{(P\_7,3.85)} kim punaḥ ayam paryudāsaḥ : yat anyat viciṇṇalṅidbhyaḥ iti . \\
\noindent \textbf{(P\_7,3.85)} āhosvit prasajya ayam pratiṣedhaḥ : viciṇṇalṅitsu na iti . \\
\noindent \textbf{(P\_7,3.85)} kaḥ ca atra viśeṣaḥ . \\
\noindent \textbf{(P\_7,3.85)} prasajyapratiṣedhe jusiguṇapratiṣedhaprasaṅgaḥ . \\
\noindent \textbf{(P\_7,3.85)} prasajyapratiṣedhe jusiguṇapratiṣedhaḥ prāpnoti . \\
\noindent \textbf{(P\_7,3.85)} ajāgaruḥ . \\
\noindent \textbf{(P\_7,3.85)} uttame ca ṇali . \\
\noindent \textbf{(P\_7,3.85)} prasajyapratiṣedhe jusiguṇapratiṣedhaḥ prāpnoti . \\
\noindent \textbf{(P\_7,3.85)} ajāgaruḥ . \\
\noindent \textbf{(P\_7,3.85)} na vā anantarasya pratiṣedhāt . \\
\noindent \textbf{(P\_7,3.85)} na vā eṣaḥ doṣaḥ . \\
\noindent \textbf{(P\_7,3.85)} kim kāraṇam . \\
\noindent \textbf{(P\_7,3.85)} anantarasya pratiṣedhāt . \\
\noindent \textbf{(P\_7,3.85)} anantaram yat guṇavidhānam tasya pratiṣedhaḥ . \\
\noindent \textbf{(P\_7,3.85)} jusi pūrveṇa guṇavidhānam . \\
\noindent \textbf{(P\_7,3.85)} jusi pūrveṇa guṇaḥ vidhīyate jusi ca iti . \\
\noindent \textbf{(P\_7,3.85)} ṇali ca . \\
\noindent \textbf{(P\_7,3.85)} kim . \\
\noindent \textbf{(P\_7,3.85)} na vā anantarasya pratiṣedhāt iti eva . \\
\noindent \textbf{(P\_7,3.85)} ṇali ca pūrveṇa guṇaḥ vidhīyate sārvadhātukārdhadhātukayoḥ iti . \\
\noindent \textbf{(P\_7,3.85)} atha vā punaḥ astu paryudāsaḥ . \\
\noindent \textbf{(P\_7,3.85)} ataḥ anyatra vidhāne vau aguṇatvam . \\
\noindent \textbf{(P\_7,3.85)} ataḥ anyatra vidhāne vau aguṇatvam . \\
\noindent \textbf{(P\_7,3.85)} na vā paryudāsasāmarthyāt . \\
\noindent \textbf{(P\_7,3.85)} na vā vaktavyam . \\
\noindent \textbf{(P\_7,3.85)} kim kāraṇam . \\
\noindent \textbf{(P\_7,3.85)} paryudāsasāmarthyāt atra guṇaḥ na bhaviṣyati . \\
\noindent \textbf{(P\_7,3.85)} asti anyat paryudāse prayojanam . \\
\noindent \textbf{(P\_7,3.85)} kim . \\
\noindent \textbf{(P\_7,3.85)} kvibartham paryudāsaḥ syāt . \\
\noindent \textbf{(P\_7,3.85)} śuddhaparasya viśabdasya pratiṣedhe grahaṇam anunāsikaparaḥ ca kvau viśabdaḥ . \\
\noindent \textbf{(P\_7,3.85)} vasvartham tarhi paryudāsaḥ syāt . \\
\noindent \textbf{(P\_7,3.85)} jāgṛvāṃsaḥ anu gman . \\
\noindent \textbf{(P\_7,3.85)} katham punaḥ veḥ paryudāsaḥ ucyamānaḥ vasvarthaḥ śakyaḥ vijñātum . \\
\noindent \textbf{(P\_7,3.85)} sāmarthyāt vasvartham iti vijñāsyate . \\
\noindent \textbf{(P\_7,3.85)} vasvartham iti cet na sārvadhātukatvāt siddham . \\
\noindent \textbf{(P\_7,3.85)} vasvartham iti cet na . \\
\noindent \textbf{(P\_7,3.85)} kim kāraṇam . \\
\noindent \textbf{(P\_7,3.85)} sārvadhātukatvāt siddham . \\
\noindent \textbf{(P\_7,3.85)} katham sārvadhātukasañjñā . \\
\noindent \textbf{(P\_7,3.85)} chāndasaḥ kvasuḥ . \\
\noindent \textbf{(P\_7,3.85)} liṭ ca chandasi sārvadhātukam api bhavati . \\
\noindent \textbf{(P\_7,3.85)} tatra sārvadhātukam apin ṅit iti ṅittvāt paryudāsaḥ bhaviṣyati . \\
\noindent \textbf{(P\_7,3.85)} atha vā vakārasya eva idam aśaktijena ikāreṇa grahaṇam \\
\noindent \textbf{(P\_7,3.86)} saṃyoge gurusañjñāyām guṇaḥ bhettuḥ na sidhyati . \\
\noindent \textbf{(P\_7,3.86)} saṃyoge gurusañjñāyām bhettā , bhettum iti guṇaḥ na prāpnoti . \\
\noindent \textbf{(P\_7,3.86)} vidhyapekṣam laghoḥ ca asau . \\
\noindent \textbf{(P\_7,3.86)} vidhyapekṣam laghugrahaṇam kṛtam laghoḥ ca asau vihitaḥ . \\
\noindent \textbf{(P\_7,3.86)} katham kuṇḍiḥ na duṣyati . \\
\noindent \textbf{(P\_7,3.86)} kuṇḍitā , huṇḍitā atra kasmāt na bhavati . \\
\noindent \textbf{(P\_7,3.86)} dhātoḥ numaḥ . \\
\noindent \textbf{(P\_7,3.86)} dhātoḥ numvidhau uktam tatra dhātugrahaṇasya prayojanam dhātūpadeśāvasthāyām eva num bhavati iti . \\
\noindent \textbf{(P\_7,3.86)} katham rañjeḥ . \\
\noindent \textbf{(P\_7,3.86)} katham rañjeḥ upadhālakṣaṇā vṛddhiḥ . \\
\noindent \textbf{(P\_7,3.86)} āścaryaḥ rāgaḥ , vicitraḥ rāgaḥ . \\
\noindent \textbf{(P\_7,3.86)} syandiśranthyoḥ nipātanāt . \\
\noindent \textbf{(P\_7,3.86)} yat ayam syandiśranthyoḥ avṛddhyartham nipātanam karoti tat jñāpayati ācāryaḥ bhavati evañjātīyakānām vṛddhiḥ iti . \\
\noindent \textbf{(P\_7,3.86)} anaṅlopaśidīrghatve vidhyapekṣe na sidhyataḥ . \\
\noindent \textbf{(P\_7,3.86)} anaṅlopaḥ . \\
\noindent \textbf{(P\_7,3.86)} dadhnā , sakthnā . \\
\noindent \textbf{(P\_7,3.86)} śidīrghatvam . \\
\noindent \textbf{(P\_7,3.86)} kuṇḍāni , vanāni . \\
\noindent \textbf{(P\_7,3.86)} evam tarhi abhyastasya yat āha aci . \\
\noindent \textbf{(P\_7,3.86)} yat ayam na abhyastasya aci piti sārvadhātuke iti ajgrahaṇam karoti tat jñāpayati ācāryaḥ bhavati evañjātīyakānām guṇaḥ iti . \\
\noindent \textbf{(P\_7,3.86)} laṅartham tat kṛtam bhavet . \\
\noindent \textbf{(P\_7,3.86)} laṅartham etat syāt . \\
\noindent \textbf{(P\_7,3.86)} anena ik . \\
\noindent \textbf{(P\_7,3.86)} knusunoḥ yat kṛtam kittvam jñāpakam syāt laghoḥ guṇe . \\
\noindent \textbf{(P\_7,3.86)} yat ayam trasigṛdhidhṛṣikṣipeḥ knuḥ ikaḥ jhal halantāt ca iti knusanau kitau karoti tat jñāpayati ācāryaḥ bhavati evañjātīyakānām guṇaḥ iti . \\
\noindent \textbf{(P\_7,3.86)} saṃyoge gurusañjñāyām guṇaḥ bhettuḥ na sidhyati , vidhyapekṣam laghoḥ ca asau katham kuṇḍiḥ na duṣyati . \\
\noindent \textbf{(P\_7,3.86)} dhātoḥ numaḥ katham rañjeḥ syandiśranthyoḥ nipātanāt , anaṅlopaśidīrghatve vidhyapekṣe na sidhyataḥ . \\
\noindent \textbf{(P\_7,3.86)} abhyastasya yat āha aci laṅartham tat kṛtam bhavet , knusunoḥ yat kṛtam kittvam jñāpakam syāt laghoḥ guṇe \\
\noindent \textbf{(P\_7,3.87)} abhyastānām upadhāhrasvatvam aci paspaśāte , cākaśīmi , vāvaśatīḥ iti darśanāt . \\
\noindent \textbf{(P\_7,3.87)} abhyastānām upadhāhrasvatvam aci vaktavyam . \\
\noindent \textbf{(P\_7,3.87)} kim prayojanam . \\
\noindent \textbf{(P\_7,3.87)} paspaśāte , cākaśīmi . \\
\noindent \textbf{(P\_7,3.87)} vāvaśatīḥ iti prayogaḥ dṛśyate . \\
\noindent \textbf{(P\_7,3.87)} kapotaḥ śaradam paspaśāte . \\
\noindent \textbf{(P\_7,3.87)} aham bhuvanam cākaśīmi . \\
\noindent \textbf{(P\_7,3.87)} vāvaśatīḥ ut ājat iti . \\
\noindent \textbf{(P\_7,3.87)} bahulam chandasi ānuṣak jujoṣat iti darśanāt . \\
\noindent \textbf{(P\_7,3.87)} bahulam chandasi vaktavyam upadhāhrasvatvam . \\
\noindent \textbf{(P\_7,3.87)} kim prayojanam . \\
\noindent \textbf{(P\_7,3.87)} ānuṣak jujoṣat iti darśanāt . \\
\noindent \textbf{(P\_7,3.87)} yaḥ te ātityam ānuṣak jujoṣat . \\
\noindent \textbf{(P\_7,3.87)} yadi upadhāhrasvatvam ucyate , priyām mayūraḥ pratinarnṛtīti yadvat tvam naravara narnṛtīṣi hṛṣṭaḥ , atra guṇaḥ prāpnoti . \\
\noindent \textbf{(P\_7,3.87)} tasmāt na arthaḥ upadhāhrasvatvena . \\
\noindent \textbf{(P\_7,3.87)} kasmāt na bhavati . \\
\noindent \textbf{(P\_7,3.87)} paspaśāte , cākaśīmi , vāvaśatīḥ iti . \\
\noindent \textbf{(P\_7,3.87)} spaśikaśivaśayaḥ prakṛtyantarāṇi \\
\noindent \textbf{(P\_7,3.88)} bhūsuvoḥ pratiṣedhe ekājgrahaṇam bobhavītyartham . \\
\noindent \textbf{(P\_7,3.88)} bhūsuvoḥ pratiṣedhe ekājgrahaṇam kartavyam . \\
\noindent \textbf{(P\_7,3.88)} kim prayojanam . \\
\noindent \textbf{(P\_7,3.88)} bobhavītyartham . \\
\noindent \textbf{(P\_7,3.88)} iha mā bhūt . \\
\noindent \textbf{(P\_7,3.88)} bobhavīti . \\
\noindent \textbf{(P\_7,3.88)} yadi ekājgrahaṇam kriyate abhūt atra na prāpnoti . \\
\noindent \textbf{(P\_7,3.88)} kva tarhi syāt . \\
\noindent \textbf{(P\_7,3.88)} mā bhūt . \\
\noindent \textbf{(P\_7,3.88)} tasmāt na arthaḥ ekājgrahaṇena . \\
\noindent \textbf{(P\_7,3.88)} kasmāt na bhavati . \\
\noindent \textbf{(P\_7,3.88)} bobhavīti iti . \\
\noindent \textbf{(P\_7,3.88)} bobhūtu iti etat niyamārtham bhaviṣyati . \\
\noindent \textbf{(P\_7,3.88)} atra eva yaṅlugantasya guṇaḥ na bhavati na anyatra iti . \\
\noindent \textbf{(P\_7,3.88)} kva mā bhūt . \\
\noindent \textbf{(P\_7,3.88)} bobhavīti iti \\
\noindent \textbf{(P\_7,3.92)} kimartham tṛhirāgataśnamkaḥ na tṛheḥ im bhavati iti eva ucyeta . \\
\noindent \textbf{(P\_7,3.92)} tṛṇahigrahaṇam śnamimoḥ vyavasthārtham . \\
\noindent \textbf{(P\_7,3.92)} tṛṇahigrahaṇam śnami kṛte im yathā syāt . \\
\noindent \textbf{(P\_7,3.92)} tṛhigrahaṇe hi imviṣaye śnamabhāvaḥ anavakāśatvāt . \\
\noindent \textbf{(P\_7,3.92)} tṛhigrahaṇe hi sati imviṣaye śnamaḥ abhāvaḥ syāt . \\
\noindent \textbf{(P\_7,3.92)} kim kāraṇam . \\
\noindent \textbf{(P\_7,3.92)} anavakāśatvāt . \\
\noindent \textbf{(P\_7,3.92)} anavakāśaḥ im śnamam bādheta . \\
\noindent \textbf{(P\_7,3.92)} idam ayuktam vartate . \\
\noindent \textbf{(P\_7,3.92)} kim atra ayuktam . \\
\noindent \textbf{(P\_7,3.92)} tṛṇahigrahaṇam śnamimoḥ vyavasthārtham iti uktvā tataḥ ucyate tṛhigrahaṇe hi imviṣaye śnamabhāvaḥ anavakāśatvāt iti . \\
\noindent \textbf{(P\_7,3.92)} tatra vaktavyam tṛṇahigrahaṇam śnamimoḥ bhāvāya tṛhigrahaṇe hi imviṣeye śnamabhāvaḥ anavakāśatvāt iti . \\
\noindent \textbf{(P\_7,3.92)} tat tarhi vaktavyam . \\
\noindent \textbf{(P\_7,3.92)} na vaktavyam . \\
\noindent \textbf{(P\_7,3.92)} vyavasthārtham iti eva siddham na hi asataḥ vyavasthā iti \\
\noindent \textbf{(P\_7,3.95)} sārvadhātuke iti vartamāne punaḥ sārvadhātukagrahaṇam kimartham . \\
\noindent \textbf{(P\_7,3.95)} punaḥ sārvadhātukagrahaṇam apidartham . \\
\noindent \textbf{(P\_7,3.95)} apidarthaḥ ayam ārambhaḥ . \\
\noindent \textbf{(P\_7,3.95)} adhrigo śamīdhvam suśami śamīdhva śamīdhvam adhrigo \\
\noindent \textbf{(P\_7,3.103)} ataḥ dīrghāt bahuvacane ettvam vipratiṣedhena . \\
\noindent \textbf{(P\_7,3.103)} ataḥ dīrghāt bahuvacane ettvam bhavati vipratiṣedhena . \\
\noindent \textbf{(P\_7,3.103)} ataḥ dīrghaḥ yañi supi ca iti asya avakāśaḥ . \\
\noindent \textbf{(P\_7,3.103)} vṛkṣābhyām , plakṣābhyām . \\
\noindent \textbf{(P\_7,3.103)} bahuvacane jhali et iti asya avakāśaḥ . \\
\noindent \textbf{(P\_7,3.103)} vṛkṣeṣu , plakṣeṣu . \\
\noindent \textbf{(P\_7,3.103)} iha ubhayam prāpnoti . \\
\noindent \textbf{(P\_7,3.103)} vṛkṣebhyaḥ , plakṣebhyaḥ . \\
\noindent \textbf{(P\_7,3.103)} ettvam bhavati vipratiṣedhena \\
\noindent \textbf{(P\_7,3.107)} ḍalakavatīnām pratiṣedhaḥ vaktavyaḥ . \\
\noindent \textbf{(P\_7,3.107)} ambāḍe , ambāle , ambike . \\
\noindent \textbf{(P\_7,3.107)} talhrasvatvam vā ṅisambuddhyoḥ . \\
\noindent \textbf{(P\_7,3.107)} talhrasvatvam vā ṅisambuddhyoḥ iti vaktavyam . \\
\noindent \textbf{(P\_7,3.107)} devata , devate . \\
\noindent \textbf{(P\_7,3.107)} devatāyām , devate . \\
\noindent \textbf{(P\_7,3.107)} saḥ tarhi pratiṣedhaḥ vaktavyaḥ . \\
\noindent \textbf{(P\_7,3.107)} na vaktavyaḥ . \\
\noindent \textbf{(P\_7,3.107)} saḥ katham na vaktavyaḥ bhavati . \\
\noindent \textbf{(P\_7,3.107)} ambārtham dvyakṣaram yadi . \\
\noindent \textbf{(P\_7,3.107)} yadi ambārtham dvyakṣaram gṛhyate . \\
\noindent \textbf{(P\_7,3.107)} tat tarhi hrasvatvam vaktavyam . \\
\noindent \textbf{(P\_7,3.107)} avaśyam chandasi hrasvatvam vaktavyam upagāyantu mām patnayaḥ garbhiṇayaḥ yuvatayaḥ iti evam artham . \\
\noindent \textbf{(P\_7,3.107)} mātṝṇām mātac putrārtham arhate . \\
\noindent \textbf{(P\_7,3.107)} mātṝṇām mātajādeśaḥ vaktavyaḥ putrārtham arhate . \\
\noindent \textbf{(P\_7,3.107)} gārgīmāta , vātsīmāta \\
\noindent \textbf{(P\_7,3.108)} iha kasmāt na bhavati . \\
\noindent \textbf{(P\_7,3.108)} nadi, kumāri , kiśori , brāhmaṇi , brahmabandhu . \\
\noindent \textbf{(P\_7,3.108)} hrasvavacanasāmarthyāt . \\
\noindent \textbf{(P\_7,3.108)} asti anyat hrasvavacane prayojanam pṛthagvibhaktim mā uccīcaram iti . \\
\noindent \textbf{(P\_7,3.108)} śakyam pṛthagvibhaktiḥ anuccārayitum . \\
\noindent \textbf{(P\_7,3.108)} katham . \\
\noindent \textbf{(P\_7,3.108)} evam ayam brūyāt . \\
\noindent \textbf{(P\_7,3.108)} ambārthānām hrasvaḥ nadīhrasvayoḥ guṇaḥ iti . \\
\noindent \textbf{(P\_7,3.108)} yadi evam ucyate jasi ca iti atra nadyāḥ api guṇaḥ prāpnoti . \\
\noindent \textbf{(P\_7,3.108)} evam tarhi yogavibhāgaḥ kariṣyate . \\
\noindent \textbf{(P\_7,3.108)} ambārthanadyoḥ hrasvaḥ . \\
\noindent \textbf{(P\_7,3.108)} tataḥ hrasvasya . \\
\noindent \textbf{(P\_7,3.108)} hrasvasya ca hrasvaḥ bhavati . \\
\noindent \textbf{(P\_7,3.108)} kimartham idam . \\
\noindent \textbf{(P\_7,3.108)} guṇam vakṣyati tadbādhanārtham . \\
\noindent \textbf{(P\_7,3.108)} tataḥ guṇaḥ . \\
\noindent \textbf{(P\_7,3.108)} guṇaḥ ca bhavati hrasvasya iti . \\
\noindent \textbf{(P\_7,3.108)} atha vā hrasvasya guṇaḥ iti atra ambārthanadyoḥ hrasvaḥ iti etat anuvartiṣyate \\
\noindent \textbf{(P\_7,3.109)} jasādiṣu chandasi vāvacanam prāk ṇau caṅi upadhāyāḥ . \\
\noindent \textbf{(P\_7,3.109)} jasādiṣu chandasi vā iti vaktavyam . \\
\noindent \textbf{(P\_7,3.109)} kim aviśeṣeṇa . \\
\noindent \textbf{(P\_7,3.109)} na iti āha . \\
\noindent \textbf{(P\_7,3.109)} prāk ṇau caṅyupadhāyāḥ . \\
\noindent \textbf{(P\_7,3.109)} kim prayojanam . \\
\noindent \textbf{(P\_7,3.109)} ambe, darvi , śatakratvaḥ , paśve nṛbhyaḥ , kikidīvyā . \\
\noindent \textbf{(P\_7,3.109)} ambe , amba , darvi , darve , śatakravaḥ śatakratavaḥ , paśve , paśave , kikidīvyā , kikidīvinā \\
\noindent \textbf{(P\_7,3.111)} gheḥ ṅiti guṇavidhāne ṅīsārvadhātuke pratiṣedhaḥ . \\
\noindent \textbf{(P\_7,3.111)} gheḥ ṅiti guṇavidhāne ṅīsārvadhātuke pratiṣedhaḥ vaktavyaḥ . \\
\noindent \textbf{(P\_7,3.111)} paṭvī , mṛdvī , kurutaḥ iti . \\
\noindent \textbf{(P\_7,3.111)} subadhikārāt siddham . \\
\noindent \textbf{(P\_7,3.111)} sup iti vartate . \\
\noindent \textbf{(P\_7,3.111)} kva prakṛtam . \\
\noindent \textbf{(P\_7,3.111)} supi ca bahuvacane jhali et iti \\
\noindent \textbf{(P\_7,3.113)} iha atikhaṭvāya , atimālāya iti hrasvatve kṛte sthānivadbhāvāt yāṭ prāpnoti tasya pratiṣedhaḥ vaktavyaḥ . \\
\noindent \textbf{(P\_7,3.113)} na vaktavyaḥ . \\
\noindent \textbf{(P\_7,3.113)} yāḍvidhāne atikhaṭvāya iti apratiṣedhaḥ hrasvādeśatvāt . \\
\noindent \textbf{(P\_7,3.113)} yāḍvidhāne atikhaṭvāya , atimālāya iti apratiṣedhaḥ . \\
\noindent \textbf{(P\_7,3.113)} anarthakaḥ pratiṣedhaḥ apratiṣedhaḥ . \\
\noindent \textbf{(P\_7,3.113)} yāṭ kasmāt na bhavati . \\
\noindent \textbf{(P\_7,3.113)} hrasvādeśatvāt . \\
\noindent \textbf{(P\_7,3.113)} hrasvādeśaḥ ayam . \\
\noindent \textbf{(P\_7,3.113)} uktam etat ṅyābgrahaṇe adīrghaḥ iti . \\
\noindent \textbf{(P\_7,3.113)} atha idānīm asati api sthānivadbhāve dīrghatve kṛte āp ca asau bhūtapūrvaḥ iti kṛtvā yāṭ kasmāt na bhavati . \\
\noindent \textbf{(P\_7,3.113)} lakṣaṇapratipadoktayoḥ pratipadoktasya eva iti . \\
\noindent \textbf{(P\_7,3.113)} nanu ca idānīm sati api sthānivadbhāve etayā paribhāṣayā śakyam upasthātum . \\
\noindent \textbf{(P\_7,3.113)} na iti āha . \\
\noindent \textbf{(P\_7,3.113)} na ca tadānīm kvacit api sthānivadbhāvaḥ syāt \\
\noindent \textbf{(P\_7,3.116)} idudbhyām āmvidhānam auttvasya paratvāt . \\
\noindent \textbf{(P\_7,3.116)} idudbhyām ām vidheyaḥ . \\
\noindent \textbf{(P\_7,3.116)} śakaṭyām , paddhatyām , dhenvām iti . \\
\noindent \textbf{(P\_7,3.116)} kim punaḥ kāraṇam na sidhyati . \\
\noindent \textbf{(P\_7,3.116)} auttvasya paratvāt . \\
\noindent \textbf{(P\_7,3.116)} paratvāt auttvam prāpnoti . \\
\noindent \textbf{(P\_7,3.116)} yogavibhāgāt siddham . \\
\noindent \textbf{(P\_7,3.116)} yogavibhāgaḥ kariṣyate . \\
\noindent \textbf{(P\_7,3.116)} ṅeḥ ām nadyāmnībhyaḥ . \\
\noindent \textbf{(P\_7,3.116)} tataḥ idudbhyām . \\
\noindent \textbf{(P\_7,3.116)} idudbhyām uttarasya ṅeḥ ām bhavati iti . \\
\noindent \textbf{(P\_7,3.116)} śakaṭyām , paddhatyām , dhenvām iti . \\
\noindent \textbf{(P\_7,3.116)} tataḥ aut at ca gheḥ \\
\noindent \textbf{(P\_7,3.118-119)} KA\_III,342.14-343.9 Ro\_V,239.9-241.1 {1/34}     auttve yogavibhāgaḥ . \\
\noindent \textbf{(P\_7,3.118-119)} KA\_III,342.14-343.9 Ro\_V,239.9-241.1 {2/34}     auttve yogavibhāgaḥ kartavyaḥ . \\
\noindent \textbf{(P\_7,3.118-119)} KA\_III,342.14-343.9 Ro\_V,239.9-241.1 {3/34}     aut . \\
\noindent \textbf{(P\_7,3.118-119)} KA\_III,342.14-343.9 Ro\_V,239.9-241.1 {4/34}     aut bhavati idudbhyām . \\
\noindent \textbf{(P\_7,3.118-119)} KA\_III,342.14-343.9 Ro\_V,239.9-241.1 {5/34}     tataḥ at ca gheḥ . \\
\noindent \textbf{(P\_7,3.118-119)} KA\_III,342.14-343.9 Ro\_V,239.9-241.1 {6/34}     akāraḥ ca bhavati gheḥ iti . \\
\noindent \textbf{(P\_7,3.118-119)} KA\_III,342.14-343.9 Ro\_V,239.9-241.1 {7/34}     kimarthaḥ yogavibhāgaḥ . \\
\noindent \textbf{(P\_7,3.118-119)} KA\_III,342.14-343.9 Ro\_V,239.9-241.1 {8/34}     sakhipatibhyām auttvārthaḥ . \\
\noindent \textbf{(P\_7,3.118-119)} KA\_III,342.14-343.9 Ro\_V,239.9-241.1 {9/34}     sakhipatibhyām auttvam yathā syāt . \\
\noindent \textbf{(P\_7,3.118-119)} KA\_III,342.14-343.9 Ro\_V,239.9-241.1 {10/34}     sakhyau , patyau . \\
\noindent \textbf{(P\_7,3.118-119)} KA\_III,342.14-343.9 Ro\_V,239.9-241.1 {11/34}     ekayoge hi aprāptiḥ attvasanniyogāt . \\
\noindent \textbf{(P\_7,3.118-119)} KA\_III,342.14-343.9 Ro\_V,239.9-241.1 {12/34}     ekayoge hi sati auttvasya aprāptiḥ . \\
\noindent \textbf{(P\_7,3.118-119)} KA\_III,342.14-343.9 Ro\_V,239.9-241.1 {13/34}     kim kāraṇam . \\
\noindent \textbf{(P\_7,3.118-119)} KA\_III,342.14-343.9 Ro\_V,239.9-241.1 {14/34}     attvasanniyogāt . \\
\noindent \textbf{(P\_7,3.118-119)} KA\_III,342.14-343.9 Ro\_V,239.9-241.1 {15/34}     attvasanniyogena auttvam ucyate tena yatra eva auttvam syāt . \\
\noindent \textbf{(P\_7,3.118-119)} KA\_III,342.14-343.9 Ro\_V,239.9-241.1 {16/34}     na vā akārasya anvācayavacanāt yathā kyaṅi salopaḥ . \\
\noindent \textbf{(P\_7,3.118-119)} KA\_III,342.14-343.9 Ro\_V,239.9-241.1 {17/34}     na vā artha auttve yogavibhāgena . \\
\noindent \textbf{(P\_7,3.118-119)} KA\_III,342.14-343.9 Ro\_V,239.9-241.1 {18/34}     kim kāraṇam . \\
\noindent \textbf{(P\_7,3.118-119)} KA\_III,342.14-343.9 Ro\_V,239.9-241.1 {19/34}     akārasya anvācayavacanāt . \\
\noindent \textbf{(P\_7,3.118-119)} KA\_III,342.14-343.9 Ro\_V,239.9-241.1 {20/34}     pradhānaśiṣṭam auttvam anvācayaśiṣṭam attvam yathā kyaṅi salopaḥ . \\
\noindent \textbf{(P\_7,3.118-119)} KA\_III,342.14-343.9 Ro\_V,239.9-241.1 {21/34}     tat yathā . \\
\noindent \textbf{(P\_7,3.118-119)} KA\_III,342.14-343.9 Ro\_V,239.9-241.1 {22/34}     pradhānaśiṣṭaḥ kyaṅ prātipadikamātrāt bhavati yatra ca skāraḥ tatra lopaḥ . \\
\noindent \textbf{(P\_7,3.118-119)} KA\_III,342.14-343.9 Ro\_V,239.9-241.1 {23/34}     attve ṭāppratiṣedhaḥ . \\
\noindent \textbf{(P\_7,3.118-119)} KA\_III,342.14-343.9 Ro\_V,239.9-241.1 {24/34}     attve ṭāpaḥ pratiṣedhaḥ vaktavyaḥ . \\
\noindent \textbf{(P\_7,3.118-119)} KA\_III,342.14-343.9 Ro\_V,239.9-241.1 {25/34}     śakaṭau , paddhatau , dhenau . \\
\noindent \textbf{(P\_7,3.118-119)} KA\_III,342.14-343.9 Ro\_V,239.9-241.1 {26/34}     attve kṛte ṭāp prāpnoti . \\
\noindent \textbf{(P\_7,3.118-119)} KA\_III,342.14-343.9 Ro\_V,239.9-241.1 {27/34}     na vā sannipātalakṣaṇasya animittatvāt . \\
\noindent \textbf{(P\_7,3.118-119)} KA\_III,342.14-343.9 Ro\_V,239.9-241.1 {28/34}     na vā vaktavyaḥ . \\
\noindent \textbf{(P\_7,3.118-119)} KA\_III,342.14-343.9 Ro\_V,239.9-241.1 {29/34}     kim kāraṇam . \\
\noindent \textbf{(P\_7,3.118-119)} KA\_III,342.14-343.9 Ro\_V,239.9-241.1 {30/34}     sannipātalakṣaṇasya animittatvāt . \\
\noindent \textbf{(P\_7,3.118-119)} KA\_III,342.14-343.9 Ro\_V,239.9-241.1 {31/34}     sannipātalakṣaṇaḥ vidhiḥ animittam tadvighātasya iti ṭāp na bhaviṣyati . \\
\noindent \textbf{(P\_7,3.118-119)} KA\_III,342.14-343.9 Ro\_V,239.9-241.1 {32/34}     ḍitkaraṇāt vā . \\
\noindent \textbf{(P\_7,3.118-119)} KA\_III,342.14-343.9 Ro\_V,239.9-241.1 {33/34}     atha vā ḍit āukāraḥ kariṣyate . \\
\noindent \textbf{(P\_7,3.118-119)} KA\_III,342.14-343.9 Ro\_V,239.9-241.1 {34/34}     au ḍit ca gheḥ \\
\noindent \textbf{(P\_7,3.120)} kimartham astriyām iti ucyate na āṅaḥ nā puṃsi iti eva ucyeta . \\
\noindent \textbf{(P\_7,3.120)} kā rūpasiddhiḥ : trapuṇā , jatunā . \\
\noindent \textbf{(P\_7,3.120)} numā siddham . \\
\noindent \textbf{(P\_7,3.120)} na evam śakyam . \\
\noindent \textbf{(P\_7,3.120)} iha hi amunā brāhmaṇakulena iti mubhāvasya asiddhatvāt num na syāt . \\
\noindent \textbf{(P\_7,3.120)} astriyām iti punaḥ ucyamāne na doṣaḥ bhavati . \\
\noindent \textbf{(P\_7,3.120)} katham . \\
\noindent \textbf{(P\_7,3.120)} vakṣyati etat na mu ṭādeśe \\
\noindent \textbf{(P\_7,4.1.1)} atha ṇigrahaṇam kimartham na caṅi upadhāyāḥ hrasvaḥ iti eva ucyeta . \\
\noindent \textbf{(P\_7,4.1.1)} caṅi upadhāyāḥ hrasvaḥ iti iyati ucyamāne , alīlavat , apīpavat , ūkārasya eva hrasvatvam prasajyeta . \\
\noindent \textbf{(P\_7,4.1.1)} na etat asti prayojanam . \\
\noindent \textbf{(P\_7,4.1.1)} vṛddhiḥ atra bādhikā bhaviṣyati . \\
\noindent \textbf{(P\_7,4.1.1)} vṛddhau tarhi kṛtāyām aukārasya eva hrasvatvam prasajyeta . \\
\noindent \textbf{(P\_7,4.1.1)} na etat asti . \\
\noindent \textbf{(P\_7,4.1.1)} antaraṅgatvāt atra āvādeśaḥ bhaviṣyati . \\
\noindent \textbf{(P\_7,4.1.1)} na hi idānīm hrasvabhāvinī upadhā bhavati . \\
\noindent \textbf{(P\_7,4.1.1)} tasmāt ṇigrahaṇam kartavyam . \\
\noindent \textbf{(P\_7,4.1.1)} atha caṅgrahaṇam kimartham na ṇau upadhāyāḥ iti eva siddham . \\
\noindent \textbf{(P\_7,4.1.1)} ṇau upadhāyāḥ hrasvaḥ iti iyati ucyamāne , kārayati , hārayati iti atra api prasajyeta . \\
\noindent \textbf{(P\_7,4.1.1)} na etat asti prayojanam . \\
\noindent \textbf{(P\_7,4.1.1)} ācāryapravṛttiḥ jñāpayati na ṇau eva hrasvatvam bhavati iti yat ayam mitām hrasvatvam śāsti . \\
\noindent \textbf{(P\_7,4.1.1)} iha api tarhi na prāpnoti . \\
\noindent \textbf{(P\_7,4.1.1)} acīkarat , ajīharat . \\
\noindent \textbf{(P\_7,4.1.1)} vacanāt bhaviṣyati . \\
\noindent \textbf{(P\_7,4.1.1)} iha api tarhi vacanāt prāpnoti . \\
\noindent \textbf{(P\_7,4.1.1)} kārayati , hārayati . \\
\noindent \textbf{(P\_7,4.1.1)} tasmāt caṅgrahaṇam kartavyam . \\
\noindent \textbf{(P\_7,4.1.1)} atha upadhāgrahaṇam kimartham . \\
\noindent \textbf{(P\_7,4.1.1)} ṇau caṅi upadhāgrahaṇam antyapratiṣedhārtham . \\
\noindent \textbf{(P\_7,4.1.1)} ṇau caṅi upadhāgrahaṇam kriyate antyasya hrastvatvam mā bhūt . \\
\noindent \textbf{(P\_7,4.1.1)} ṇau caṅi hrasvaḥ iti iyati ucyamāne , alīlavat , apīpavat , antyasya eva hrasvatvam prasajyeta . \\
\noindent \textbf{(P\_7,4.1.1)} na etat asti prayojanam . \\
\noindent \textbf{(P\_7,4.1.1)} antaraṅgatvāt atra āvādeśaḥ bhavati . \\
\noindent \textbf{(P\_7,4.1.1)} na hi idānīm hrasvabhāvī antyaḥ asti . \\
\noindent \textbf{(P\_7,4.1.1)} antyaḥ hrasvabhāvī na asti iti kṛtvā vacanāt anantyasya bhaviṣyati . \\
\noindent \textbf{(P\_7,4.1.1)} iha api vacanāt prāpnoti . \\
\noindent \textbf{(P\_7,4.1.1)} acakāṅkṣat , avavāñchat . \\
\noindent \textbf{(P\_7,4.1.1)} yena na avyavadhānam tena vyavahite api vacanaprāmāṇyāt . \\
\noindent \textbf{(P\_7,4.1.1)} kena ca na avyavadhānam . \\
\noindent \textbf{(P\_7,4.1.1)} varṇena . \\
\noindent \textbf{(P\_7,4.1.1)} etena punaḥ saṅghatena vayvadhānam bhavati na bhavati ca . \\
\noindent \textbf{(P\_7,4.1.1)} uttarārtham tarhi upadhāgrahaṇam kartavyam . \\
\noindent \textbf{(P\_7,4.1.1)} lopaḥ pibateḥ ī ca abhyāsasya upadhāyāḥ yathā syāt . \\
\noindent \textbf{(P\_7,4.1.1)} apīpyat , apīpyatām , apīpyan . \\
\noindent \textbf{(P\_7,4.1.1)} atha iha katham bhavitavyam . \\
\noindent \textbf{(P\_7,4.1.1)} mā bhavān aṭiṭat iti . \\
\noindent \textbf{(P\_7,4.1.1)} āhosvit mā bhavān āṭiṭat iti . \\
\noindent \textbf{(P\_7,4.1.1)} mā bhavān āṭiṭat iti bhavitavyam . \\
\noindent \textbf{(P\_7,4.1.1)} hrasvatvam kasmāt na bhavati . \\
\noindent \textbf{(P\_7,4.1.1)} dvirvacane kṛte pareṇa rūpeṇa vyavahitam iti kṛtvā . \\
\noindent \textbf{(P\_7,4.1.1)} idam iha sampradhāryam . \\
\noindent \textbf{(P\_7,4.1.1)} dvirvacanam kriyatām hrasvatvam iti kim atra kartavyam . \\
\noindent \textbf{(P\_7,4.1.1)} paratvāt hrasvatvam . \\
\noindent \textbf{(P\_7,4.1.1)} nityam dvirvacanam . \\
\noindent \textbf{(P\_7,4.1.1)} kṛte hrasvatve prāpnoti akṛte api . \\
\noindent \textbf{(P\_7,4.1.1)} evam tarhi ācāryapravṛttiḥ jñāpayati dvirvacanāt hrasvatvam balīyaḥ iti yat ayam oṇim ṛditam karoti . \\
\noindent \textbf{(P\_7,4.1.1)} katham kṛtvā jñāpakam . \\
\noindent \textbf{(P\_7,4.1.1)} ṛditkaraṇe etat prayojanam ṛditām na iti pratiṣedhaḥ yathā syāt . \\
\noindent \textbf{(P\_7,4.1.1)} yadi ca atra pūrvam dvirvacanam syāt ṛditkaraṇam anarthakam syāt . \\
\noindent \textbf{(P\_7,4.1.1)} dvirvacane kṛte pareṇa vyavahitatvāt hrasvatvam na bhaviṣyati . \\
\noindent \textbf{(P\_7,4.1.1)} paśyati tu ācāryaḥ dvirvacanāt hrasvatvam balīyaḥ iti tataḥ oṇim ṛditam karoti . \\
\noindent \textbf{(P\_7,4.1.1)} tasmāt mā bhavān aṭiṭat iti eva bhavitavyam \\
\noindent \textbf{(P\_7,4.1.2)} upadhāhrasvatve ṇeḥ ṇici upasaṅkhyānāt . \\
\noindent \textbf{(P\_7,4.1.2)} upadhāhrasvatve ṇeḥ ṇici upasaṅkhyānam kartavyam : vāditavantam prayojitavān avīvadadvīṇām parivādakena . \\
\noindent \textbf{(P\_7,4.1.2)} kim punaḥ kāraṇam na sidhyati . \\
\noindent \textbf{(P\_7,4.1.2)} ṇicā vyavahitatvāt . \\
\noindent \textbf{(P\_7,4.1.2)} ṇilope kṛte na asti vyavadhānam . \\
\noindent \textbf{(P\_7,4.1.2)} sthānivadbhāvāt vyavadhānam eva . \\
\noindent \textbf{(P\_7,4.1.2)} pratiṣidhyate atra sthānivadbhāvaḥ caṅparanirhrāse na sthānivat iti . \\
\noindent \textbf{(P\_7,4.1.2)} evam api aglopinām na iti pratiṣedhaḥ prāpnoti . \\
\noindent \textbf{(P\_7,4.1.2)} vṛddhau kṛtāyām lopaḥ tat na aglopai aṅgam bhavati . \\
\noindent \textbf{(P\_7,4.1.2)} idam iha sampradhāryam . \\
\noindent \textbf{(P\_7,4.1.2)} vṛddhiḥ kriyatām lopaḥ iti kim atra kartavyam . \\
\noindent \textbf{(P\_7,4.1.2)} paratvāt vṛddhiḥ . \\
\noindent \textbf{(P\_7,4.1.2)} nityaḥ lopaḥ . \\
\noindent \textbf{(P\_7,4.1.2)} kṛtāyām api vṛddhau prāpnoti akṛtāyām api . \\
\noindent \textbf{(P\_7,4.1.2)} anityaḥ lopaḥ . \\
\noindent \textbf{(P\_7,4.1.2)} anyasya kṛtāyām vṛddhau prāpnoti anyasya akṛtāyām śabdāntarasya ca prāpnuvan vidhiḥ anityaḥ bhavati . \\
\noindent \textbf{(P\_7,4.1.2)} ubhayoḥ anityayoḥ paratvāt vṛddhiḥ . \\
\noindent \textbf{(P\_7,4.1.2)} vṛddhau kṛtāyām lopaḥ tat na aglopi aṅgam bhavati . \\
\noindent \textbf{(P\_7,4.1.2)} evam tarhi ācāryapravṛttiḥ jñāpayati vṛddheḥ lopaḥ balīyān iti yat ayam aglopinām na iti pratiṣedham śāsti . \\
\noindent \textbf{(P\_7,4.1.2)} na etat asti jñāpakam . \\
\noindent \textbf{(P\_7,4.1.2)} asti anyat etasya vacane prayojanam . \\
\noindent \textbf{(P\_7,4.1.2)} kim . \\
\noindent \textbf{(P\_7,4.1.2)} yatra vṛddhau api kṛtāyām eva lupyate . \\
\noindent \textbf{(P\_7,4.1.2)} atyararājat . \\
\noindent \textbf{(P\_7,4.1.2)} yat tarhi pratyāhāragrahaṇam karoti . \\
\noindent \textbf{(P\_7,4.1.2)} itarathā hi alopinām na iti brūyāt . \\
\noindent \textbf{(P\_7,4.1.2)} evam vā vṛddheḥ lopaḥ balīyān iti . \\
\noindent \textbf{(P\_7,4.1.2)} atha vā ārabhyate pūrvavipratiṣedhaḥ ṇyallopāviyaṅyaṇguṇavṛddhidīrghatvebhyaḥ pūrvavipratiṣiddham iti . \\
\noindent \textbf{(P\_7,4.1.2)} tasmāt upasaṅkhyānam kartavyam iti \\
\noindent \textbf{(P\_7,4.2)} aglopipratiṣedhānarthakyam ca sthānivadbhāvāt . \\
\noindent \textbf{(P\_7,4.2)} aglopipratiṣedhaḥ ca anarthakaḥ . \\
\noindent \textbf{(P\_7,4.2)} kim kāraṇam . \\
\noindent \textbf{(P\_7,4.2)} sthānivadbhāvāt . \\
\noindent \textbf{(P\_7,4.2)} sthānivadbhāvāt atra hrasvatvam na bhaviṣyati . \\
\noindent \textbf{(P\_7,4.2)} yatra tarhi sthānivadbhāvaḥ na asti tadartham ayam yogaḥ vaktavyaḥ . \\
\noindent \textbf{(P\_7,4.2)} kva sthānivadbhāvaḥ na asti . \\
\noindent \textbf{(P\_7,4.2)} yaḥ halacoḥ ādeśaḥ . \\
\noindent \textbf{(P\_7,4.2)} atyararājat . \\
\noindent \textbf{(P\_7,4.2)} kim punaḥ kāraṇam halacoḥ ādeśaḥ na sthānivat iti ucyate . \\
\noindent \textbf{(P\_7,4.2)} ajādeśaḥ sthānivat iti ucyate na ca ayam acaḥ eva ādeśaḥ . \\
\noindent \textbf{(P\_7,4.2)} kim tarhi acaḥ anyasya ca . \\
\noindent \textbf{(P\_7,4.2)} aglopinām na iti api tarhi pratiṣedhaḥ na prāpnoti . \\
\noindent \textbf{(P\_7,4.2)} kim kāraṇam . \\
\noindent \textbf{(P\_7,4.2)} aglopinām na iti ucyate na ca atra ac eva lupyate . \\
\noindent \textbf{(P\_7,4.2)} kim tarhi . \\
\noindent \textbf{(P\_7,4.2)} ac ca anyaḥ ca . \\
\noindent \textbf{(P\_7,4.2)} yaḥ atra ac lupyate tadāśrayaḥ pratiṣedhaḥ bhaviṣyati . \\
\noindent \textbf{(P\_7,4.2)} yathā eva tarhi yaḥ atra ac lupyate tadāśrayaḥ pratiṣedhaḥ bhavati evam yaḥ atra ac lupyate tadāśrayaḥ sthānivadbhāvaḥ bhaviṣyati . \\
\noindent \textbf{(P\_7,4.2)} evam tarhi siddhe sati yat aglopinām na iti pratiṣedham śāsti tat jñāpayati ācāryaḥ itaḥ uttaram sthānivadbhāvaḥ na bhavati iti . \\
\noindent \textbf{(P\_7,4.2)} kim etasya jñāpane prayojanam . \\
\noindent \textbf{(P\_7,4.2)} pūrvatra asiddhe na sthānivat iti uktam tat na vaktavyam bhavati . \\
\noindent \textbf{(P\_7,4.2)} yadi etat jñāpyate , ādīdhayateḥ ādīdhakaḥ , āvevayateḥ āvevakaḥ . \\
\noindent \textbf{(P\_7,4.2)} yīvarṇayoḥ dīdhīvevyoḥ iti lopaḥ na prāpnoti . \\
\noindent \textbf{(P\_7,4.2)} iha ca yat pralunīhi atra tiṅi ca udāttavati iti eṣaḥ svaraḥ na prāpnoti . \\
\noindent \textbf{(P\_7,4.2)} na eṣaḥ doṣaḥ . \\
\noindent \textbf{(P\_7,4.2)} yat tāvat ucyate ādīdhayateḥ ādīdhakaḥ , āvevayateḥ āvevakaḥ , yīvarṇayoḥ iti lopaḥ na prāpnoti iti yīvarṇayoḥ iti atra varṇagrahaṇasāmarthyāt bhaviṣyati . \\
\noindent \textbf{(P\_7,4.2)} yat api ucyate yat pralunīhi atra tiṅi ca udāttavati iti eṣaḥ svaraḥ na prāpnoti iti bahiraṅgaḥ yaṇādeśaḥ antaraṅgaḥ svaraḥ asiddham bahiraṅgam antaraṅge \\
\noindent \textbf{(P\_7,4.3)} kāṇyādīnām ca iti vaktavyam . \\
\noindent \textbf{(P\_7,4.3)} ke punaḥ kāṇyādayaḥ . \\
\noindent \textbf{(P\_7,4.3)} kāṇirāṇiśrāṇibhāṇiheṭhilopayaḥ . \\
\noindent \textbf{(P\_7,4.3)} acakāṇat , acīkaṇat , ararāṇat , arīraṇat , aśaśrāṇat , aśīśraṇat , ababhāṇat , abībhaṇat , ajiheṭhat , ajīhiṭhat , alulopat , alūlupat \\
\noindent \textbf{(P\_7,4.9)} iha avadigye , avadigyāte , avadigyare digyādeśe kṛte dvirvacanam prāpnoti tatra sābhyāsasya iti vaktavyam . \\
\noindent \textbf{(P\_7,4.9)} nanu ca dvirvacane kṛte sābhyāsasya digyādeśaḥ bhaviṣyati . \\
\noindent \textbf{(P\_7,4.9)} na sidhyati . \\
\noindent \textbf{(P\_7,4.9)} kim kāraṇam . \\
\noindent \textbf{(P\_7,4.9)} digyādeśasya paratvāt sābhyāsasya ādeśavacanam . \\
\noindent \textbf{(P\_7,4.9)} digyādeśaḥ kriyatām dvirvacanam iti paratvāt digyādeśena bhavitavyam . \\
\noindent \textbf{(P\_7,4.9)} tatra sābhyāsasya iti vaktavyam . \\
\noindent \textbf{(P\_7,4.9)} evam tarhi digyādeśaḥ dvirvacanam bādhiṣyate . \\
\noindent \textbf{(P\_7,4.9)} punaḥprasaṅgavijñānāt dvirvacanam prāpnoti . \\
\noindent \textbf{(P\_7,4.9)} punaḥprasaṅgaḥ iti cet amādibhiḥ tulyam . \\
\noindent \textbf{(P\_7,4.9)} punaḥprasaṅgaḥ iti cet amādibhiḥ tulyam etat bhavati . \\
\noindent \textbf{(P\_7,4.9)} tat yathā . \\
\noindent \textbf{(P\_7,4.9)} amādiṣu kṛteṣu punaḥprasaṅgāt śiśīlugnumaḥ na bhavanti . \\
\noindent \textbf{(P\_7,4.9)} evam digyādeśe kṛte punaḥprasaṅgāt dvirvacanam na bhaviṣyati . \\
\noindent \textbf{(P\_7,4.9)} atha vā vipratiṣedhe punaḥprasaṅgaḥ iti ucyate vipratiṣedhaḥ ca dvayoḥ sāvakaśayoḥ . \\
\noindent \textbf{(P\_7,4.9)} iha punaḥ anavakāśaḥ digyādeśaḥ dvirvacanam bādhiṣyate . \\
\noindent \textbf{(P\_7,4.9)} yadi tarhi anavakāśāḥ vidhayaḥ bādhakāḥ bhavanti , babhūva , bhūbhāvaḥ dvirvacanam bādheta . \\
\noindent \textbf{(P\_7,4.9)} sāvakāśaḥ bhūbhāvaḥ . \\
\noindent \textbf{(P\_7,4.9)} kaḥ avakāśaḥ . \\
\noindent \textbf{(P\_7,4.9)} bhavitā , bhavitum . \\
\noindent \textbf{(P\_7,4.9)} iha tarhi cakṣiṅaḥ khyāñ vā liṭi iti khyāñ dvirvacanam bādheta . \\
\noindent \textbf{(P\_7,4.9)} iha ca api babhūva iti yadi tāvat sthāne dvirvacanam bhūbhāvaḥ sarvādeśaḥ prāpnoti atha dviḥprayogaḥ dvirvacanam parasya bhūbhāve kṛte pūrvasya śravaṇam prāpnoti . \\
\noindent \textbf{(P\_7,4.9)} na eṣaḥ doṣaḥ . \\
\noindent \textbf{(P\_7,4.9)} ārdhadhātukīyāḥ sāmānyena bhavanti anavasthiteṣu pratyayeṣu . \\
\noindent \textbf{(P\_7,4.9)} tatra ārdhadhātukasāmānye bhūbhāve kṛte yaḥ yataḥ pratyayaḥ prāpnoti saḥ tataḥ bhaviṣyati \\
\noindent \textbf{(P\_7,4.10)} saṃyogādeḥ guṇavidhāne saṃyogopadhagrahaṇam kṛñartham . \\
\noindent \textbf{(P\_7,4.10)} saṃyogādeḥ guṇavidhāne saṃyogopadhagrahaṇam kartavyam . \\
\noindent \textbf{(P\_7,4.10)} kimartham . \\
\noindent \textbf{(P\_7,4.10)} kṛñartham . \\
\noindent \textbf{(P\_7,4.10)} iha api yathā syāt . \\
\noindent \textbf{(P\_7,4.10)} sañcaskaratuḥ , sañcaskaruḥ . \\
\noindent \textbf{(P\_7,4.10)} yadi saṃyogopadhagrahaṇam kriyate na arthaḥ saṃyogādigrahaṇena . \\
\noindent \textbf{(P\_7,4.10)} iha api sasvaratuḥ , sasvaruḥ saṃyogopadhasya iti eva siddham . \\
\noindent \textbf{(P\_7,4.10)} bhavet siddham sasvaratuḥ , sasvaruḥ iti idam tu na sidhyati sañcaskaratuḥ , sañcaskaruḥ iti . \\
\noindent \textbf{(P\_7,4.10)} kim kāraṇam . \\
\noindent \textbf{(P\_7,4.10)} suṭaḥ bahiraṅgalakṣaṇatvāt . \\
\noindent \textbf{(P\_7,4.10)} bahiraṅgam suṭ . \\
\noindent \textbf{(P\_7,4.10)} antaraṅgaḥ guṇaḥ . \\
\noindent \textbf{(P\_7,4.10)} asiddham bahiraṅgam antaraṅge . \\
\noindent \textbf{(P\_7,4.10)} saṃyogādigrahaṇe tu kriyamāṇe saṃyogopadhagrahaṇam ananyārtham vijñāyate . \\
\noindent \textbf{(P\_7,4.10)} ṛtaḥ liṭi guṇāt ñṇiti vṛddhiḥ vipratiṣedhena . \\
\noindent \textbf{(P\_7,4.10)} ṛtaḥ liṭi guṇāñ ñṅiti vṛddhiḥ bhavati pūrvavipratiṣedhena . \\
\noindent \textbf{(P\_7,4.10)} ṛtaḥ liṭi guṇasya avakāśaḥ . \\
\noindent \textbf{(P\_7,4.10)} sasvaratuḥ , sasvaruḥ . \\
\noindent \textbf{(P\_7,4.10)} ñṇiti vṛddheḥ avakāśaḥ . \\
\noindent \textbf{(P\_7,4.10)} svārakaḥ , dhvārakaḥ . \\
\noindent \textbf{(P\_7,4.10)} iha ubhayam prāpnoti . \\
\noindent \textbf{(P\_7,4.10)} sasvāra , dadhvāra . \\
\noindent \textbf{(P\_7,4.10)} ñṇiti vṛddhiḥ bhavati pūrvavipratiṣedhena . \\
\noindent \textbf{(P\_7,4.10)} punaḥprasaṅgavijñānāt vā siddham . \\
\noindent \textbf{(P\_7,4.10)} atha vā punaḥprasaṅgāt guṇe kṛte raparatve ca ataḥ upadhāyāḥ iti vṛddhiḥ bhaviṣyati . \\
\noindent \textbf{(P\_7,4.10)} na eṣaḥ yuktaḥ parihāraḥ . \\
\noindent \textbf{(P\_7,4.10)} punaḥprasaṅgaḥ nāma saḥ bhavati yatra tena eva kṛte prāpnoti tena eva ca akṛte . \\
\noindent \textbf{(P\_7,4.10)} atra khalu guṇe kṛte raparatve ca ataḥ upadhāyāḥ iti vṛddhiḥ prāpnoti akṛte ca acaḥ ñṇiti iti . \\
\noindent \textbf{(P\_7,4.10)} tasmāt suṣṭhu ucyate liti guṇāt ñṇiti vṛddhiḥ vipratiṣedhena iti \\
\noindent \textbf{(P\_7,4.12)} kimartham hrasvaḥ vā iti ucyate na guṇaḥ vā iti ucyeta . \\
\noindent \textbf{(P\_7,4.12)} tatra ayam api arthaḥ guṇagrahaṇam na kartavyam bhavati . \\
\noindent \textbf{(P\_7,4.12)} prakṛtam anuvartate . \\
\noindent \textbf{(P\_7,4.12)} kva prakṛtam . \\
\noindent \textbf{(P\_7,4.12)} ṛtaḥ ca saṃyogādeḥ guṇaḥ iti . \\
\noindent \textbf{(P\_7,4.12)} ṛtaḥ hrasvatvam ittvapratiṣedhārtham . \\
\noindent \textbf{(P\_7,4.12)} ṛtaḥ hrasvatvam ucyate ittvapratiṣedhārtham . \\
\noindent \textbf{(P\_7,4.12)} ittvam mā bhūt iti . \\
\noindent \textbf{(P\_7,4.12)} guṇaḥ vā iti iyati ucyamāne guṇena mukte ittvam prasajyeta . \\
\noindent \textbf{(P\_7,4.12)} hrasvaḥ vā iti ucyamāne hrasvena mukte yathāprāptaḥ guṇaḥ bhaviṣyati \\
\noindent \textbf{(P\_7,4.13)} ke aṇaḥ hrasvatve taddhitagrahaṇam kṛnnivṛttyartham . \\
\noindent \textbf{(P\_7,4.13)} ke aṇaḥ hrasvatve taddhitagrahaṇam kartavyam . \\
\noindent \textbf{(P\_7,4.13)} kim prayojanam . \\
\noindent \textbf{(P\_7,4.13)} kṛnnivṛttyartham . \\
\noindent \textbf{(P\_7,4.13)} kṛti mā bhūt . \\
\noindent \textbf{(P\_7,4.13)} rākā , dhākā iti . \\
\noindent \textbf{(P\_7,4.13)} tat tarhi vaktavyam . \\
\noindent \textbf{(P\_7,4.13)} na vaktavyam . \\
\noindent \textbf{(P\_7,4.13)} uṇādayaḥ avyutpannāni prātipadikāni \\
\noindent \textbf{(P\_7,4.23)} iha kasmāt na bhavati . \\
\noindent \textbf{(P\_7,4.23)} prohyate , upohyate . \\
\noindent \textbf{(P\_7,4.23)} ekādeśe kṛte vyapavargābhāvāt . \\
\noindent \textbf{(P\_7,4.23)} evam api ā , ūhyate , ohyate , samohyate . \\
\noindent \textbf{(P\_7,4.23)} aṇaḥ iti vartate \\
\noindent \textbf{(P\_7,4.24)} eteḥ liṅi upasargāt . \\
\noindent \textbf{(P\_7,4.24)} eteḥ liṅi upasargāt iti vaktavyam . \\
\noindent \textbf{(P\_7,4.24)} iha mā bhūt . \\
\noindent \textbf{(P\_7,4.24)} īyāt . \\
\noindent \textbf{(P\_7,4.24)} tat tarhi vaktavyam . \\
\noindent \textbf{(P\_7,4.24)} na vaktavyam . \\
\noindent \textbf{(P\_7,4.24)} upasargāt iti vartate . \\
\noindent \textbf{(P\_7,4.24)} evam tarhi ācāryaḥ anvācaṣṭe upasargāt iti anuvartate iti . \\
\noindent \textbf{(P\_7,4.24)} na etat anvākhyeyam adhikārāḥ anuvartante iti . \\
\noindent \textbf{(P\_7,4.24)} eṣaḥ eva nyāyaḥ yat uta adhikārāḥ anuvarteran \\
\noindent \textbf{(P\_7,4.27)} dīrghoccāraṇam kimartham na riṅ ṛtaḥ iti eva ucyate . \\
\noindent \textbf{(P\_7,4.27)} kā rūpasiddhiḥ : mātrīyati , pitrīyati . \\
\noindent \textbf{(P\_7,4.27)} akṛtsārvadhātukayoḥ iti dīrghatvam bhaviṣyati . \\
\noindent \textbf{(P\_7,4.27)} evam tarhi siddhe sati yat dīrghoccāraṇam karoti tat jñāpayati ācāryaḥ bhavati eṣā paribhāṣā aṅgavṛtte punaḥ vṛttau avidhiḥ niṣṭhitasya iti . \\
\noindent \textbf{(P\_7,4.27)} kim etasya jñāpane prayojanam . \\
\noindent \textbf{(P\_7,4.27)} pibeḥ guṇapratiṣedhaḥ coditaḥ saḥ na vaktavyaḥ bhavati \\
\noindent \textbf{(P\_7,4.30)} yaṅprakaraṇe hanteḥ hiṃsāyām īṭ . \\
\noindent \textbf{(P\_7,4.30)} yaṇprakaraṇe hanteḥ hiṃsāyām īṭ vaktavyaḥ . \\
\noindent \textbf{(P\_7,4.30)} jeghnīyate . \\
\noindent \textbf{(P\_7,4.30)} yadi īṭ abhyāsarūpam na sidhyati . \\
\noindent \textbf{(P\_7,4.30)} evam tarhi yaṅprakaraṇe hanteḥ hiṃsāyām īk . \\
\noindent \textbf{(P\_7,4.30)} evam api upadhālopaḥ na prāpnoti . \\
\noindent \textbf{(P\_7,4.30)} evam tarhi yaṅprakaraṇe hanteḥ hiṃsāyām ghnī \\
\noindent \textbf{(P\_7,4.35)} atyalpam idam ucyate : aputrasya iti . \\
\noindent \textbf{(P\_7,4.35)} aputrādīnām iti vaktavyam iha api yathā syāt : janīyantaḥ nvagravaḥ putrīyantaḥ sudānavaḥ . \\
\noindent \textbf{(P\_7,4.35)} chandasi pratiṣedhe dīrghapratiṣedhaḥ . \\
\noindent \textbf{(P\_7,4.35)} chandasi pratiṣedhe dīrghatvasya pratiṣedhaḥ vaktavyaḥ . \\
\noindent \textbf{(P\_7,4.35)} saṃsvedayuḥ , mitrayuḥ . \\
\noindent \textbf{(P\_7,4.35)} na vā aśvāghasya ādvacanam avadhāraṇārtham . \\
\noindent \textbf{(P\_7,4.35)} na vā vaktavyam . \\
\noindent \textbf{(P\_7,4.35)} kim kāraṇam . \\
\noindent \textbf{(P\_7,4.35)} aśvāghasya ādvacanam avadhāraṇārtham bhaviṣyati aśvāghayoḥ eva chancasi dīrghaḥ bhaviṣyati na anyasya iti \\
\noindent \textbf{(P\_7,4.41)} śyateḥ ittvam vrate nityam . \\
\noindent \textbf{(P\_7,4.41)} śyateḥ ittvam vrate nityam iti vaktavyam . \\
\noindent \textbf{(P\_7,4.41)} saṃśitavrataḥ . \\
\noindent \textbf{(P\_7,4.41)} tat tarhi vaktavyam . \\
\noindent \textbf{(P\_7,4.41)} na vaktavyam . \\
\noindent \textbf{(P\_7,4.41)} devatrātaḥ galaḥ grāhaḥ itiyoge ca sadvidhiḥ , mithaḥ te na vibhāṣyante gavākṣaḥ saṃśitavrataḥ \\
\noindent \textbf{(P\_7,4.46)} avadattam vidattam ca pradattam ca ādikarmaṇi , sudattam anudattam ca nidattam iti ca iṣyate . \\
\noindent \textbf{(P\_7,4.46)} kim punaḥ ayam takārāntaḥ āhosvit dakārāntaḥ uta dhakārāntaḥ atha vā thakārāntaḥ . \\
\noindent \textbf{(P\_7,4.46)} kaḥ ca atra viśeṣaḥ . \\
\noindent \textbf{(P\_7,4.46)} tānte doṣaḥ dīrghatvam syāt . \\
\noindent \textbf{(P\_7,4.46)} yadi takārāntaḥ dasti iti dīrghatvam prāpnoti . \\
\noindent \textbf{(P\_7,4.46)} dānte doṣaḥ niṣṭhānatvam . \\
\noindent \textbf{(P\_7,4.46)} atha dakārāntaḥ radābhyām niṣṭhātaḥ iti natvam prāpnoti . \\
\noindent \textbf{(P\_7,4.46)} dhānte doṣaḥ dhatvaprāptiḥ . \\
\noindent \textbf{(P\_7,4.46)} atha dhāntaḥ jhaṣaḥ tathoḥ dhaḥ adhaḥ iti dhatvam prāpnoti . \\
\noindent \textbf{(P\_7,4.46)} thānte adoṣaḥ tasmāt thāntaḥ . \\
\noindent \textbf{(P\_7,4.46)} atha thakārāntaḥ na doṣaḥ bhavati \\
\noindent \textbf{(P\_7,4.47)} acaḥ upasargāt tatve ākāragrahaṇam . \\
\noindent \textbf{(P\_7,4.47)} acaḥ upasargāt tatve ākāragrahaṇam kartavyam . \\
\noindent \textbf{(P\_7,4.47)} na kartavyam . \\
\noindent \textbf{(P\_7,4.47)} alaḥ antyasya vidhayaḥ bhavanti iti ākārasya bhaviṣyati . \\
\noindent \textbf{(P\_7,4.47)} na sidhyati . \\
\noindent \textbf{(P\_7,4.47)} kim kāraṇam . \\
\noindent \textbf{(P\_7,4.47)} ādeḥ hi parasya . \\
\noindent \textbf{(P\_7,4.47)} atra hi tasmāt iti uttarasya ādeḥ parasya iti dakārasya prāpnoti . \\
\noindent \textbf{(P\_7,4.47)} na eṣaḥ doṣaḥ . \\
\noindent \textbf{(P\_7,4.47)} avarṇaprakaraṇāt siddham . \\
\noindent \textbf{(P\_7,4.47)} asya iti vartate . \\
\noindent \textbf{(P\_7,4.47)} kva prakṛtam . \\
\noindent \textbf{(P\_7,4.47)} asya dvau iti . \\
\noindent \textbf{(P\_7,4.47)} yadi avarṇagrahaṇam anuvartate dadbhāve doṣaḥ bhavati . \\
\noindent \textbf{(P\_7,4.47)} evam tarhi evam vakṣyāmi daḥ adghoḥ iti . \\
\noindent \textbf{(P\_7,4.47)} daḥ yaḥ ākāraḥ tasya at bhavati . \\
\noindent \textbf{(P\_7,4.47)} tataḥ acaḥ upasargāt taḥ . \\
\noindent \textbf{(P\_7,4.47)} asya iti eva . \\
\noindent \textbf{(P\_7,4.47)} evam api sūtrabhedaḥ kṛtaḥ bhavati . \\
\noindent \textbf{(P\_7,4.47)} na asau sūtrabhedaḥ . \\
\noindent \textbf{(P\_7,4.47)} sūtrabhedam kam upācaranti . \\
\noindent \textbf{(P\_7,4.47)} yatra tat eva anyat sūtram kriyate bhūyaḥ vā . \\
\noindent \textbf{(P\_7,4.47)} yat hi tat eva upasaṃhṛtya kriyate na asau sūtrabhedaḥ . \\
\noindent \textbf{(P\_7,4.47)} atha vā dvitakārakaḥ nirdeśaḥ kriyate saḥ anekāl śit sarvasya iti sarvasya bhaviṣyati . \\
\noindent \textbf{(P\_7,4.47)} iha api tarhi prāpnoti . \\
\noindent \textbf{(P\_7,4.47)} adbhiḥ , adbhyaḥ iti . \\
\noindent \textbf{(P\_7,4.47)} acaḥ iti vartate . \\
\noindent \textbf{(P\_7,4.47)} tat ca avaśyam ajgrahaṇam anuvartyam lavābhyām iti evamartham . \\
\noindent \textbf{(P\_7,4.47)} atha vā tritakārakaḥ nirdeśaḥ kariṣyate ihārthau dvau uttarārthaḥ ca ekaḥ . \\
\noindent \textbf{(P\_7,4.47)} dyateḥ ittvāt acaḥ taḥ . \\
\noindent \textbf{(P\_7,4.47)} dyateḥ ittvāt acaḥ taḥ iti etat bhavati vipratiṣedhena . \\
\noindent \textbf{(P\_7,4.47)} dyateḥ ittvasya avakāśaḥ . \\
\noindent \textbf{(P\_7,4.47)} nirditam , nirditavān . \\
\noindent \textbf{(P\_7,4.47)} acaḥ taḥ iti asya avakāśaḥ . \\
\noindent \textbf{(P\_7,4.47)} prattam , avattam . \\
\noindent \textbf{(P\_7,4.47)} iha ubhayam prāpnoti . \\
\noindent \textbf{(P\_7,4.47)} nīttam , vīttam . \\
\noindent \textbf{(P\_7,4.47)} acaḥ taḥ iti etat bhavati vipratiṣedhena \\
\noindent \textbf{(P\_7,4.48)} apaḥ bhi māsaḥ chandasi . \\
\noindent \textbf{(P\_7,4.48)} apaḥ bhi iti atra māsaḥ chandasi upasaṅkhyānam kartavyam : mā adbhiḥ iṣṭvā indraḥ vṛtrahā . \\
\noindent \textbf{(P\_7,4.48)} atyalpam idam ucyate . \\
\noindent \textbf{(P\_7,4.48)} svavassvatavasoḥ māsaḥ uṣasaḥ ca taḥ iṣyate : svavadbhiḥ , svatavadbhiḥ , samuṣadbhiḥ ajāyathāḥ , mā adbhiḥ iṣṭvā indraḥ vṛtrahā \\
\noindent \textbf{(P\_7,4.54)} istvam sani rādhaḥ hiṃsāyām . \\
\noindent \textbf{(P\_7,4.54)} istvam sani rādhaḥ hiṃsāyām iti vaktavyam . \\
\noindent \textbf{(P\_7,4.54)} pratiritsati . \\
\noindent \textbf{(P\_7,4.54)} hiṃsāyām iti kimartham . \\
\noindent \textbf{(P\_7,4.54)} ārirātsati \\
\noindent \textbf{(P\_7,4.55)} jñapeḥ īttvam anantyasya . \\
\noindent \textbf{(P\_7,4.55)} jñapeḥ īttvam anantyasya iti vaktavyam . \\
\noindent \textbf{(P\_7,4.55)} jñīpsati . \\
\noindent \textbf{(P\_7,4.55)} tat tarhi vaktavyam . \\
\noindent \textbf{(P\_7,4.55)} na vaktavyam . \\
\noindent \textbf{(P\_7,4.55)} lopaḥ antyasya bādhakaḥ bhaviṣyati . \\
\noindent \textbf{(P\_7,4.55)} anavakāśāḥ vidhayaḥ bādhakāḥ bhavanti sāvakāśaḥ ca ṇilopaḥ . \\
\noindent \textbf{(P\_7,4.55)} kaḥ avakāśaḥ . \\
\noindent \textbf{(P\_7,4.55)} kāraṇā , hāraṇā . \\
\noindent \textbf{(P\_7,4.55)} evam api īttvam antyasya lopasya bādhakam syāt . \\
\noindent \textbf{(P\_7,4.55)} anavakāśāḥ hi vidhayaḥ bādhakāḥ bhavanti īttvam api sāvakāśam . \\
\noindent \textbf{(P\_7,4.55)} kaḥ avakāśaḥ . \\
\noindent \textbf{(P\_7,4.55)} anantyaḥ . \\
\noindent \textbf{(P\_7,4.55)} katham punaḥ sati antye anantyasya īttvam syāt . \\
\noindent \textbf{(P\_7,4.55)} bhavet yaḥ acā āṅgam viśeṣayet tasya anantyasya na syāt . \\
\noindent \textbf{(P\_7,4.55)} vayam tu khalu aṅena acam viśeṣayiṣyāmaḥ . \\
\noindent \textbf{(P\_7,4.55)} aṅgasya acaḥ yatratatrasthasya iti . \\
\noindent \textbf{(P\_7,4.55)} evam api ubhayoḥ sāvakāśayoḥ paratvāt īttvam prāpnoti . \\
\noindent \textbf{(P\_7,4.55)} tasmāt anantyasya iti vaktavyam \\
\noindent \textbf{(P\_7,4.58)} abhyāsasya anaci . \\
\noindent \textbf{(P\_7,4.58)} abhyāsasya iti yat ucyate tat anaci draṣṭavyam . \\
\noindent \textbf{(P\_7,4.58)} patāpataḥ , carācaraḥ , vadāvadaḥ \\
\noindent \textbf{(P\_7,4.60)} kim ayam ṣaṣṭhīsamāsaḥ : halām ādiḥ halādiḥ halādiḥ śiṣyate iti . \\
\noindent \textbf{(P\_7,4.60)} āhosvit karmadhārayaḥ : hal ādiḥ halādiḥ halādiḥ śiṣyate iti . \\
\noindent \textbf{(P\_7,4.60)} kaḥ ca atra viśeṣaḥ . \\
\noindent \textbf{(P\_7,4.60)} halādiśeṣe ṣaṣṭhīsamāsaḥ iti cet ajādiṣu śeṣaprasaṅgaḥ . \\
\noindent \textbf{(P\_7,4.60)} halādiśeṣe ṣaṣṭhīsamāsaḥ iti cet ajādiṣu śeṣaḥ prāpnoti . \\
\noindent \textbf{(P\_7,4.60)} ānakṣa , ānakṣatuḥ , ānakṣuḥ . \\
\noindent \textbf{(P\_7,4.60)} astu tarhi karmadhārayaḥ . \\
\noindent \textbf{(P\_7,4.60)} karmadhārayaḥ iti cet ādiśeṣanimittatvāt lopasya tadabhāve lopavacanam . \\
\noindent \textbf{(P\_7,4.60)} karmadhārayaḥ iti cet ādiśeṣanimittatvāt lopasya tadabhāve ādyasya halaḥ abhave lopaḥ vaktavyaḥ . \\
\noindent \textbf{(P\_7,4.60)} āṭatuḥ , āṭuḥ . \\
\noindent \textbf{(P\_7,4.60)} tasmāt anādilopaḥ . \\
\noindent \textbf{(P\_7,4.60)} tasmāt anādiḥ hal lupyate iti vaktavyam . \\
\noindent \textbf{(P\_7,4.60)} uktam vā . \\
\noindent \textbf{(P\_7,4.60)} kim uktam . \\
\noindent \textbf{(P\_7,4.60)} pratividhāsyate halādiśeṣaḥ iti . \\
\noindent \textbf{(P\_7,4.60)} ayam idānīm saḥ pratividhānakālaḥ . \\
\noindent \textbf{(P\_7,4.60)} idam pratividhīyate . \\
\noindent \textbf{(P\_7,4.60)} idam prakṛtam atra lopaḥ abhyāsasya iti . \\
\noindent \textbf{(P\_7,4.60)} tataḥ vakṣyāmi . \\
\noindent \textbf{(P\_7,4.60)} hrasvaḥ . \\
\noindent \textbf{(P\_7,4.60)} hrasvaḥ bhavati ādeśaḥ . \\
\noindent \textbf{(P\_7,4.60)} abhyāsasya lopaḥ iti anuvartate . \\
\noindent \textbf{(P\_7,4.60)} tatra hrasvabhāvinām hrasvaḥ lopabhāvinām lopaḥ bhaviṣyati . \\
\noindent \textbf{(P\_7,4.60)} tataḥ halādiḥ śeṣaḥ ca iti . \\
\noindent \textbf{(P\_7,4.60)} atha vā evam vakṣyāmi . \\
\noindent \textbf{(P\_7,4.60)} hrasvaḥ ahal . \\
\noindent \textbf{(P\_7,4.60)} hrasvaḥ bhavati abhyāsasya iti . \\
\noindent \textbf{(P\_7,4.60)} tataḥ ahal . \\
\noindent \textbf{(P\_7,4.60)} ahal ca bhavati abhyāsaḥ . \\
\noindent \textbf{(P\_7,4.60)} tataḥ ādiḥ śeṣaḥ . \\
\noindent \textbf{(P\_7,4.60)} ādiḥ śeṣaḥ bhavati abhyāsasya iti . \\
\noindent \textbf{(P\_7,4.60)} atha vā yogavibhāgaḥ kariṣyate . \\
\noindent \textbf{(P\_7,4.60)} hrasvādeśaḥ bhavati abhyāsasya . \\
\noindent \textbf{(P\_7,4.60)} tataḥ hal . \\
\noindent \textbf{(P\_7,4.60)} hal ca lupyate abhyāsasya . \\
\noindent \textbf{(P\_7,4.60)} tataḥ ādiḥ śeṣaḥ . \\
\noindent \textbf{(P\_7,4.60)} ādiḥ śeṣaḥ ca bhavati abhyāsasya \\
\noindent \textbf{(P\_7,4.61)} śarpūrvaśeṣe kharpūrvagrahaṇam . \\
\noindent \textbf{(P\_7,4.61)} śarpūrvaśeṣe kharpūrvagrahaṇam kartavyam . \\
\noindent \textbf{(P\_7,4.61)} kharpūrvāḥ khayaḥ śiṣyante kharaḥ lupyante iti vaktavyam . \\
\noindent \textbf{(P\_7,4.61)} kim prayojanam . \\
\noindent \textbf{(P\_7,4.61)} ucicchiṣati . \\
\noindent \textbf{(P\_7,4.61)} vyucicchiṣati . \\
\noindent \textbf{(P\_7,4.61)} tukaḥ śravaṇam mā bhūt iti . \\
\noindent \textbf{(P\_7,4.61)} tat tarhi vaktavyam . \\
\noindent \textbf{(P\_7,4.61)} na vaktavyam . \\
\noindent \textbf{(P\_7,4.61)} cartve kṛte tuk na bhaviṣyati . \\
\noindent \textbf{(P\_7,4.61)} asiddham cartvam tasya asiddhatvāt tuk prāpnoti . \\
\noindent \textbf{(P\_7,4.61)} siddhakāṇḍe paṭhitam abhyāsajaśtvacartvam ettvatukoḥ iti . \\
\noindent \textbf{(P\_7,4.61)} evam api antaraṅgatvāt prāpnoti tasmāt kharpūrvagrahaṇam kartavyam . \\
\noindent \textbf{(P\_7,4.61)} na kartavyam . \\
\noindent \textbf{(P\_7,4.61)} ettvatuggrahaṇam na kariṣyate . \\
\noindent \textbf{(P\_7,4.61)} abhyāsajaśtvacartvam siddham iti eva . \\
\noindent \textbf{(P\_7,4.61)} ādiśeṣaprasaṅgaḥ tu . \\
\noindent \textbf{(P\_7,4.61)} ādiśeṣaḥ tu prāpnoti . \\
\noindent \textbf{(P\_7,4.61)} tiṣṭhāsati . \\
\noindent \textbf{(P\_7,4.61)} nanu ca anādiśeṣaḥ ādiśeṣam bādhiṣyate . \\
\noindent \textbf{(P\_7,4.61)} katham anyasya ucyamānam anyasya bādhakam syāt . \\
\noindent \textbf{(P\_7,4.61)} asati khalu api sambhave bādhanam bhavati asti ca sambhavaḥ yat ubhayam syāt . \\
\noindent \textbf{(P\_7,4.61)} yadi ādiśeṣaḥ api bhavati śarpūrvavacanam idānīm kimartham syāt . \\
\noindent \textbf{(P\_7,4.61)} śarpūrvavacanam kimartham iti cet khayām lopapratiṣedhārtham . \\
\noindent \textbf{(P\_7,4.61)} śarpūrvavacanam kimartham iti cet khayām lopaḥ mā bhūt iti . \\
\noindent \textbf{(P\_7,4.61)} vyapakarṣavijñānāt siddham . \\
\noindent \textbf{(P\_7,4.61)} vyapakarṣavijñānāt siddham etat . \\
\noindent \textbf{(P\_7,4.61)} kim idam vyapakarṣavijñānāt iti . \\
\noindent \textbf{(P\_7,4.61)} apavādavijñānāt . \\
\noindent \textbf{(P\_7,4.61)} apavādatvāt atra anādiśeṣaḥ ādiśeṣam bādhiṣyate . \\
\noindent \textbf{(P\_7,4.61)} nanu ca uktam katham anyasya ucyamānam anyasya bādhakam syāt iti . \\
\noindent \textbf{(P\_7,4.61)} idam tāvat ayam praṣṭavyaḥ . \\
\noindent \textbf{(P\_7,4.61)} yadi tat na ucyeta kim iha syāt . \\
\noindent \textbf{(P\_7,4.61)} halādiśeṣaḥ . \\
\noindent \textbf{(P\_7,4.61)} halādiśeṣaḥ cet na aprāpte halādiśeṣe idam ucyate tat bādhakam bhaviṣyati . \\
\noindent \textbf{(P\_7,4.61)} yat api ucyate asati khalu api sambhave bādhanam bhavati asti ca sambhavaḥ yat ubhayam syāt iti sati api sambhave bādhanam bhavati . \\
\noindent \textbf{(P\_7,4.61)} tat yathā . \\
\noindent \textbf{(P\_7,4.61)} dadhi brāhmaṇebhyaḥ dīyatām takram kauṇḍinyāya iti sati api sambhave dadhidānasya takradānam bādhakam bhavati . \\
\noindent \textbf{(P\_7,4.61)} evam iha api sati api sambhave anādiśeṣaḥ ādiśeṣam bādhiṣyate \\
\noindent \textbf{(P\_7,4.65)} dādharti iti kim nipātyate . \\
\noindent \textbf{(P\_7,4.65)} dhārayateḥ ślau abhyāsasya dīrghatvam ṇiluk ca . \\
\noindent \textbf{(P\_7,4.65)} anipātyam . \\
\noindent \textbf{(P\_7,4.65)} tūtujānavadabhyāsasya dīrghatvam parṇaśuṣivat ṇiluk bhaviṣyati . \\
\noindent \textbf{(P\_7,4.65)} dhṛṅaḥ vā abhyāsasya dīrghatvam parasmaipadam ca . \\
\noindent \textbf{(P\_7,4.65)} anipātyam . \\
\noindent \textbf{(P\_7,4.65)} tūtujānavadabhyāsasya dīrghatvam yudhyativat parasmaipadam bhaviṣyati . \\
\noindent \textbf{(P\_7,4.65)} dardharti iti kim nipātyate . \\
\noindent \textbf{(P\_7,4.65)} dhārayateḥ ślau abhyāsasya ruk ṇiluk ca . \\
\noindent \textbf{(P\_7,4.65)} anipātyam . \\
\noindent \textbf{(P\_7,4.65)} devāaduhravadruṭ parṇaśruṣivat ṇiluk bhaviṣyati . \\
\noindent \textbf{(P\_7,4.65)} dhṛṅaḥ vā abhyāsasya ruk parasmaipadam ca . \\
\noindent \textbf{(P\_7,4.65)} anipātyam . \\
\noindent \textbf{(P\_7,4.65)} devāaduhravat ruḍyudhyativat parasmaipadam ca bhaviṣyati . \\
\noindent \textbf{(P\_7,4.65)} bobhūtu iti kim nipātyate . \\
\noindent \textbf{(P\_7,4.65)} bhavateḥ yaṅlugantasya aguṇatvam nipātyate . \\
\noindent \textbf{(P\_7,4.65)} na etat asti prayojanam . \\
\noindent \textbf{(P\_7,4.65)} siddham atra aguṇatvam bhūsuvoḥ tiṅi iti . \\
\noindent \textbf{(P\_7,4.65)} evam tarhi niyamārtham bhaviṣyati . \\
\noindent \textbf{(P\_7,4.65)} atra eva yaṅlugantasya guṇaḥ na bhavati na anyatra iti . \\
\noindent \textbf{(P\_7,4.65)} kva mā bhūt . \\
\noindent \textbf{(P\_7,4.65)} bobhavīti iti . \\
\noindent \textbf{(P\_7,4.65)} tetikte iti kim nipātyate . \\
\noindent \textbf{(P\_7,4.65)} tijeḥ yaṅlugantasya ātmanepadam nipātyate . \\
\noindent \textbf{(P\_7,4.65)} na etat asti prayojanam . \\
\noindent \textbf{(P\_7,4.65)} siddham atra ātmanepadam anudāttaṅitaḥ ātmanepadam iti . \\
\noindent \textbf{(P\_7,4.65)} niyamārtham tarhi bhaviṣyati . \\
\noindent \textbf{(P\_7,4.65)} atra eva yaṅlugantasya ātmanepadam bhavati na anyatra iti . \\
\noindent \textbf{(P\_7,4.65)} kva mā bhūt . \\
\noindent \textbf{(P\_7,4.65)} bebhidi iti cecchidi iti \\
\noindent \textbf{(P\_7,4.67)} kimartham svapeḥ abhyāsasya samprasāraṇam ucyate yadā sarveṣu abhyāsasthāneṣu svapeḥ samprasāraṇam uktam . \\
\noindent \textbf{(P\_7,4.67)} svāpigrahaṇam vyapetārtham . \\
\noindent \textbf{(P\_7,4.67)} svāpigrahaṇam kriyate vyapetārtham . \\
\noindent \textbf{(P\_7,4.67)} vyapetārthaḥ ayam ārambhaḥ . \\
\noindent \textbf{(P\_7,4.67)} suṣvāpayiṣati iti . \\
\noindent \textbf{(P\_7,4.67)} asti prayojanam etat . \\
\noindent \textbf{(P\_7,4.67)} kim tarhi iti . \\
\noindent \textbf{(P\_7,4.67)} tatra kyajante atiprasaṅgaḥ . \\
\noindent \textbf{(P\_7,4.67)} tatra kyajante atiprasaṅgaḥ bhavati . \\
\noindent \textbf{(P\_7,4.67)} iha api prāpnoti . \\
\noindent \textbf{(P\_7,4.67)} svāpakam icchati svāpakīyati svāpakīyateḥ san sisvāpakīyiṣati iti . \\
\noindent \textbf{(P\_7,4.67)} siddham tu ṇigrahaṇāt . \\
\noindent \textbf{(P\_7,4.67)} siddham etat . \\
\noindent \textbf{(P\_7,4.67)} katham . \\
\noindent \textbf{(P\_7,4.67)} ṇigrahaṇam kartavyam . \\
\noindent \textbf{(P\_7,4.67)} na kartavyam . \\
\noindent \textbf{(P\_7,4.67)} nirdeśāt eva hi vyaktam ṇyantasya grahaṇam iti . \\
\noindent \textbf{(P\_7,4.67)} na atra nirdeśaḥ pramāṇam śakyam kartum . \\
\noindent \textbf{(P\_7,4.67)} yathā hi nirdeśaḥ tathā iha api prasajyeta . \\
\noindent \textbf{(P\_7,4.67)} svāpam karoti svāpayati svāpayateḥ san sisvāpayiṣati iti . \\
\noindent \textbf{(P\_7,4.67)} tasmāt ṇigrahaṇam kartavyam \\
\noindent \textbf{(P\_7,4.75)} trigrahaṇānarthakyam gaṇāntatvāt . \\
\noindent \textbf{(P\_7,4.75)} trigrahaṇam anarthakam . \\
\noindent \textbf{(P\_7,4.75)} kim kāraṇam . \\
\noindent \textbf{(P\_7,4.75)} gaṇāntatvāt . \\
\noindent \textbf{(P\_7,4.75)} trayaḥ eva nijādayaḥ . \\
\noindent \textbf{(P\_7,4.75)} uttarārtham tu . \\
\noindent \textbf{(P\_7,4.75)} uttarārtham tarhi trigrahaṇam kartavyam . \\
\noindent \textbf{(P\_7,4.75)} bhṛñām it trayāṇām yathā syāt . \\
\noindent \textbf{(P\_7,4.75)} iha mā bhūt . \\
\noindent \textbf{(P\_7,4.75)} jahāti \\
\noindent \textbf{(P\_7,4.77)} artigrahaṇam kimartham na bahulam chandasi iti eva siddham . \\
\noindent \textbf{(P\_7,4.77)} na hi antareṇa chandaḥ arteḥ śluḥ labhyaḥ . \\
\noindent \textbf{(P\_7,4.77)} evam tarhi siddhe yat artigrahaṇam karoti tat jñāpayati ācāryaḥ bhāṣāyām śluḥ bhavati iti . \\
\noindent \textbf{(P\_7,4.77)} kim etasya jñāpane prayojanam . \\
\noindent \textbf{(P\_7,4.77)} iyarti iti etat siddham bhavati \\
\noindent \textbf{(P\_7,4.82)} aicoḥ yaṅi dīrghaprasaṅgaḥ hrasvāt hi param dīrghatvam . \\
\noindent \textbf{(P\_7,4.82)} aicoḥ yaṅi dīrghatvam prāpnoti . \\
\noindent \textbf{(P\_7,4.82)} ḍoḍhaukyate , totraukyate iti . \\
\noindent \textbf{(P\_7,4.82)} nanu ca hrasvatve kṛte dīrghatvam na bhaviṣyati . \\
\noindent \textbf{(P\_7,4.82)} na sidhyati . \\
\noindent \textbf{(P\_7,4.82)} kim kāraṇam . \\
\noindent \textbf{(P\_7,4.82)} hrasvāt hi param dīrghatvam . \\
\noindent \textbf{(P\_7,4.82)} hrasvatvam kriyatām dīrghatvam iti kim atra kartavyam . \\
\noindent \textbf{(P\_7,4.82)} paratvāt dīrghatvena bhavitavyam . \\
\noindent \textbf{(P\_7,4.82)} na vā abhyāsavikāreṣu apavādasya utsargābādhakatvāt . \\
\noindent \textbf{(P\_7,4.82)} na vā eṣaḥ doṣaḥ . \\
\noindent \textbf{(P\_7,4.82)} kim kāraṇam . \\
\noindent \textbf{(P\_7,4.82)} abhyāsavikāreṣu apavādasya utsargābādhakatvāt . \\
\noindent \textbf{(P\_7,4.82)} abhyāsavikāreṣu apavādāḥ utsargān na bādhante iti eṣā paribhāṣā kartavyā . \\
\noindent \textbf{(P\_7,4.82)} kāni etasyāḥ paribhāṣāyāḥ prayojanāni . \\
\noindent \textbf{(P\_7,4.82)} prayojanam sanvadbhāvasya dīrghatvam . \\
\noindent \textbf{(P\_7,4.82)} acīkarat , ajīharat . \\
\noindent \textbf{(P\_7,4.82)} sanvadbhāvam apavādatvāt dīrghatvam na bādhate . \\
\noindent \textbf{(P\_7,4.82)} mānprabhṛtīnām dīrghatvam ittvasya . \\
\noindent \textbf{(P\_7,4.82)} mānprabhṛtīnām dīrghatvam apavādatvāt ittvam na bādhate . \\
\noindent \textbf{(P\_7,4.82)} gaṇeḥ ītvam halādiśeṣasya . \\
\noindent \textbf{(P\_7,4.82)} gaṇeḥ ītvam apavādatvāt halādiśeṣam na bādhate . \\
\noindent \textbf{(P\_7,4.82)} idam ayuktam vartate . \\
\noindent \textbf{(P\_7,4.82)} kim atra ayuktam . \\
\noindent \textbf{(P\_7,4.82)} aicoḥ yaṅi dīrghaprasaṅgaḥ hrasvāt hi param dīrghatvam iti uktvā tataḥ ucyate na vā abhyāsavikāreṣu apavādasya utsargābādhakatvāt iti . \\
\noindent \textbf{(P\_7,4.82)} tasyāḥ ca paribhāṣāyāḥ prayojanāni nāma ucyante prayojanam sanvadbhāvasya dīrghatvam mānprabhṛtīnām dīrghatvam ittvasya gaṇeḥ ītvam halādiśeṣasya iti ca . \\
\noindent \textbf{(P\_7,4.82)} na ca sanvadbhāvam apavādatvāt dīrghatvam bādhate . \\
\noindent \textbf{(P\_7,4.82)} kim tarhi paratvāt . \\
\noindent \textbf{(P\_7,4.82)} na khalu api mānprabhṛtīnām dīrghatvam apavādatvāt dīrghatvam bādhate . \\
\noindent \textbf{(P\_7,4.82)} kim tarhi antaraṅgatvāt . \\
\noindent \textbf{(P\_7,4.82)} na khalu api gaṇeḥ īttvam apavādatvāt halādiśeṣam bādhate . \\
\noindent \textbf{(P\_7,4.82)} kim tarhi anavakāśatvāt . \\
\noindent \textbf{(P\_7,4.82)} evam tarhi iyam paribhāṣā kartavyā abhyāsavikāreṣu bādhakāḥ na bādhante iti . \\
\noindent \textbf{(P\_7,4.82)} sā tarhi eṣā paribhāṣā kartavyā . \\
\noindent \textbf{(P\_7,4.82)} na kartavyā . \\
\noindent \textbf{(P\_7,4.82)} ācāryapravṛttiḥ jñāpayati bhavati eṣā paribhāṣā iti yat ayam akitaḥ iti pratiṣedham śāsti \\
\noindent \textbf{(P\_7,4.83)} akitaḥ iti kimartham . \\
\noindent \textbf{(P\_7,4.83)} yaṃyamyate , raṃramyate . \\
\noindent \textbf{(P\_7,4.83)} akitaḥ iti śakyam akartum . \\
\noindent \textbf{(P\_7,4.83)} kasmāt na bhavati yaṃyamyate , raṃramyate iti . \\
\noindent \textbf{(P\_7,4.83)} nuki kṛte anajantatvāt . \\
\noindent \textbf{(P\_7,4.83)} ataḥ uttaram paṭhati . \\
\noindent \textbf{(P\_7,4.83)} akidvacanam anyatra kidantasya alaḥ antyanivṛttyartham . \\
\noindent \textbf{(P\_7,4.83)} akidvacanam kriyate jñāpakārtham . \\
\noindent \textbf{(P\_7,4.83)} kim jñāpyam . \\
\noindent \textbf{(P\_7,4.83)} etat jñāpayati ācāryaḥ anyatra kidantasya abhyāsasya alontyavidhiḥ na bhavati iti . \\
\noindent \textbf{(P\_7,4.83)} kim etasya jñāpane prayojanam . \\
\noindent \textbf{(P\_7,4.83)} prayojanam hrasvatvāttvettvaguṇeṣu . \\
\noindent \textbf{(P\_7,4.83)} hrasvatvam . \\
\noindent \textbf{(P\_7,4.83)} avacacchatuḥ , avacacchuḥ . \\
\noindent \textbf{(P\_7,4.83)} attvam . \\
\noindent \textbf{(P\_7,4.83)} cacchṛdatuḥ , cacchṛduḥ . \\
\noindent \textbf{(P\_7,4.83)} ittvam . \\
\noindent \textbf{(P\_7,4.83)} cicchādayiṣati , cicchardayiṣati . \\
\noindent \textbf{(P\_7,4.83)} guṇaḥ . \\
\noindent \textbf{(P\_7,4.83)} cecchidyate , cocchuṣyate . \\
\noindent \textbf{(P\_7,4.83)} tuki kṛte anantyatvāt ete vidhayaḥ na prāpnuvanti . \\
\noindent \textbf{(P\_7,4.83)} vipratiṣedhāt siddham . \\
\noindent \textbf{(P\_7,4.83)} na etāni santi prayojanāni . \\
\noindent \textbf{(P\_7,4.83)} vipratiṣedhena api etāni siddhāni . \\
\noindent \textbf{(P\_7,4.83)} tuk kriyatām ete vidhayaḥ iti kim atra kartavyam . \\
\noindent \textbf{(P\_7,4.83)} paratvāt ete vidhayaḥ iti . \\
\noindent \textbf{(P\_7,4.83)} tadantāgrahaṇāt vā . \\
\noindent \textbf{(P\_7,4.83)} atha vā na evam vijñāyate abhyāsasya ajantasya ṛkārāntasya akārāntasya igantasya iti . \\
\noindent \textbf{(P\_7,4.83)} katham tarhi . \\
\noindent \textbf{(P\_7,4.83)} abhyāse yaḥ ac abhyāse yaḥ ṛkāraḥ abhyāse yaḥ akāraḥ abhyāse yaḥ ic iti . \\
\noindent \textbf{(P\_7,4.83)} evam ca kṛtvā dīrghatvam prāpnoti . \\
\noindent \textbf{(P\_7,4.83)} evam tarhi idam iha vyapadeśyam sat ācāryaḥ na vyapadiśati . \\
\noindent \textbf{(P\_7,4.83)} kim apavādaḥ nuk dīrghatvasya iti . \\
\noindent \textbf{(P\_7,4.83)} evam tarhi siddhe sati yat akitaḥ iti pratiṣedham śāsti tat jñāpayati ācāryaḥ bhavati eṣā paribhāṣā abhyāsavikāreṣu bādhakāḥ na bādhante iti . \\
\noindent \textbf{(P\_7,4.83)} kim etasya jñāpane prayojanam . \\
\noindent \textbf{(P\_7,4.83)} aicoḥ yaṅi dīrghaprasaṅgaḥ hrasvāt hi param dīrghatvam iti uktam saḥ na doṣaḥ bhavati \\
\noindent \textbf{(P\_7,4.85)} nuki yaṃyamyate , raṃramyate iti rūpāsiddhiḥ . \\
\noindent \textbf{(P\_7,4.85)} nuki sati yaṃyamyate , raṃramyate iti rūpam na sidhyati . \\
\noindent \textbf{(P\_7,4.85)} anusvārāgamavacanāt siddham . \\
\noindent \textbf{(P\_7,4.85)} anusvārāgamaḥ vaktavyaḥ . \\
\noindent \textbf{(P\_7,4.85)} evam api idam eva rūpam syāt yaṃyyamyate , idam na syāt yaṃyamyate . \\
\noindent \textbf{(P\_7,4.85)} padāntavat ca . \\
\noindent \textbf{(P\_7,4.85)} padāntāt ca iti vaktavyam . \\
\noindent \textbf{(P\_7,4.85)} vā padāntasya iti \\
\noindent \textbf{(P\_7,4.90)} rīk ṛtvataḥ saṃyogārtham . \\
\noindent \textbf{(P\_7,4.90)} rīk ṛtvataḥ iti vaktavyam . \\
\noindent \textbf{(P\_7,4.90)} kim prayojanam . \\
\noindent \textbf{(P\_7,4.90)} saṃyogārtham . \\
\noindent \textbf{(P\_7,4.90)} saṃyogāntāḥ prayojayanti . \\
\noindent \textbf{(P\_7,4.90)} varīvṛścyate , parīpṛcchyate , barībhṛjjyate \\
\noindent \textbf{(P\_7,4.91)} marmṛjyate , marmṛjyamānāsaḥ iti ca upasaṅkhyānam . \\
\noindent \textbf{(P\_7,4.91)} marmṛjyate , marmṛjyamānāsaḥ iti ca upasaṅkhyānam kartavyam . \\
\noindent \textbf{(P\_7,4.91)} marmṛjyate , marmṛjyamānāsaḥ \\
\noindent \textbf{(P\_7,4.92)} kim idam ṛkāragrahaṇam aṅgaviśeṣaṇam . \\
\noindent \textbf{(P\_7,4.92)} ṛkārāntasya aṅgasya iti . \\
\noindent \textbf{(P\_7,4.92)} āhosvit abhyāsaviśeṣaṇam . \\
\noindent \textbf{(P\_7,4.92)} ṛkārāntasya abhyāsasya iti . \\
\noindent \textbf{(P\_7,4.92)} aṅgaviśeṣaṇam iti āha . \\
\noindent \textbf{(P\_7,4.92)} katham jñāyate . \\
\noindent \textbf{(P\_7,4.92)} yat ayam taparakaraṇam karoti . \\
\noindent \textbf{(P\_7,4.92)} katham kṛtvā jñāpakam . \\
\noindent \textbf{(P\_7,4.92)} na hi kaḥ cit abhyāse dīrghaḥ asti yadartham taparakaraṇam kriyeta . \\
\noindent \textbf{(P\_7,4.92)} atha aṅgaviśeṣaṇe ṛkāragrahaṇe sati taparakaraṇe kim prayojanam . \\
\noindent \textbf{(P\_7,4.92)} iha mā bhūt . \\
\noindent \textbf{(P\_7,4.92)} cākīrti , cākīrtaḥ , cākirati . \\
\noindent \textbf{(P\_7,4.92)} kiratim carkarītāntam pacati iti atra yaḥ nayet , prāptijñam tam aham manye prārabdhaḥ tena saṅgrahaḥ \\
\noindent \textbf{(P\_7,4.93)} iha kasmāt na bhavati . \\
\noindent \textbf{(P\_7,4.93)} ajajāgarat . \\
\noindent \textbf{(P\_7,4.93)} laghuni caṅpare iti ucyate vyavahitam ca atra laghu caṅparam . \\
\noindent \textbf{(P\_7,4.93)} iha api tarhi na prāpnoti . \\
\noindent \textbf{(P\_7,4.93)} acīkarat , ajīharat . \\
\noindent \textbf{(P\_7,4.93)} vacanāt bhaviṣyati . \\
\noindent \textbf{(P\_7,4.93)} iha api vacanāt prāpnoti . \\
\noindent \textbf{(P\_7,4.93)} ajajāgarat . \\
\noindent \textbf{(P\_7,4.93)} yena na avyavadhānam tena vyavahite api vacanaprāmāṇyāt . \\
\noindent \textbf{(P\_7,4.93)} kena ca na avyavadhānam . \\
\noindent \textbf{(P\_7,4.93)} varṇena . \\
\noindent \textbf{(P\_7,4.93)} etena punaḥ saṅghātena vyavadhānam bhavati na bhavati ca . \\
\noindent \textbf{(P\_7,4.93)} evam api acikṣaṇat atra na prāpnoti . \\
\noindent \textbf{(P\_7,4.93)} evam tarhi ācāryapravṛttiḥ jñāpayati bhavati evañjātīyakānām ittvam iti yat ayam atsmṛdṝtvaraprathamradastṝspaśām iti ittvabādhanārtham attvam śāsti . \\
\noindent \textbf{(P\_7,4.93)} sanvadbhāvadīrghatve ṇeḥ ṇici upasaṅkhyānam . \\
\noindent \textbf{(P\_7,4.93)} sanvadbhāvadīrghatve ṇeḥ ṇici upasaṅkhyānam kartavyam : vāditavantam prayojitavān , avīvadat vīṇām parivādakena . \\
\noindent \textbf{(P\_7,4.93)} kim punaḥ kāraṇam na sidhyati . \\
\noindent \textbf{(P\_7,4.93)} ṇicā vyavahitatvāt . \\
\noindent \textbf{(P\_7,4.93)} lope kṛte na asti vyavadhānam . \\
\noindent \textbf{(P\_7,4.93)} sthānivadbhāvāt vyavadhānam eva . \\
\noindent \textbf{(P\_7,4.93)} pratiṣidhyate atra sthānivadbhāvaḥ dīrghavidhim prati na sthānivat iti . \\
\noindent \textbf{(P\_7,4.93)} evam api anaglopaḥ iti pratiṣedham prāpnoti . \\
\noindent \textbf{(P\_7,4.93)} vṛddhau kṛtāyām lopaḥ tat na aglopi aṅgam bhavati . \\
\noindent \textbf{(P\_7,4.93)} evam tarhi idam iha sampradhāryam . \\
\noindent \textbf{(P\_7,4.93)} vṛddhiḥ kriyatām lopaḥ iti kim atra kartavyam . \\
\noindent \textbf{(P\_7,4.93)} paratvāt vṛddhiḥ . \\
\noindent \textbf{(P\_7,4.93)} nityaḥ lopaḥ . \\
\noindent \textbf{(P\_7,4.93)} kṛtāyām api vṛddhau prāpnoti akṛtāyām api . \\
\noindent \textbf{(P\_7,4.93)} lopaḥ api anityaḥ . \\
\noindent \textbf{(P\_7,4.93)} anyasya kṛtāyām vṛddhau prāpnoti akṛtāyām anyasya śabdāntasya ca prāpnuvan vidhiḥ anityaḥ bhavati . \\
\noindent \textbf{(P\_7,4.93)} ubhayoḥ anityayoḥ paratvāt vṛddhiḥ . \\
\noindent \textbf{(P\_7,4.93)} vṛddhau kṛtāyām lopaḥ tat na aglopi aṅgam bhavati . \\
\noindent \textbf{(P\_7,4.93)} mīmādīnām tu lopaprasaṅgaḥ . \\
\noindent \textbf{(P\_7,4.93)} mīmādīnām tu lopaḥ prāpnoti . \\
\noindent \textbf{(P\_7,4.93)} amīmapat . \\
\noindent \textbf{(P\_7,4.93)} siddham tu rūpātideśāt . \\
\noindent \textbf{(P\_7,4.93)} siddham etat . \\
\noindent \textbf{(P\_7,4.93)} katham . \\
\noindent \textbf{(P\_7,4.93)} rūpātideśaḥ ayam . \\
\noindent \textbf{(P\_7,4.93)} sani yādṛśam abhyāsarūpam tat sanvadbhāvena atidiśyate na ca mīmādīnām sani abhyāsarūpam asti . \\
\noindent \textbf{(P\_7,4.93)} aṅgānyatvāt vā siddham . \\
\noindent \textbf{(P\_7,4.93)} atha vā ṇyantam etat aṅgam anyat . \\
\noindent \textbf{(P\_7,4.93)} lope kṛte na aṅgānyatvam . \\
\noindent \textbf{(P\_7,4.93)} sthānivadbhāvāt aṅgam anyat . \\
\noindent \textbf{(P\_7,4.93)} katham ajijñapat . \\
\noindent \textbf{(P\_7,4.93)} atra sani api ṇyantasya eva upādānam āpjñapyṛdhām īt iti . \\
\noindent \textbf{(P\_7,4.93)} atra aṅgānyatvābhāvāt abhyāsalopaḥ syāt . \\
\noindent \textbf{(P\_7,4.93)} tasmāt pūrvaḥ eva parihāraḥ siddham tu rūpātideśāt iti \\
\noindent \textbf{(P\_8,1.1.1)} sarvavacanam kimartham . \\
\noindent \textbf{(P\_8,1.1.1)} sarvavacanam alontyanivṛttyartham . \\
\noindent \textbf{(P\_8,1.1.1)} sarvagrahaṇam kriyate alontyanivṛttyartham . \\
\noindent \textbf{(P\_8,1.1.1)} alaḥ antyasya vidhayaḥ bhavanti iti antyasya dvirvacanam mā bhūt iti . \\
\noindent \textbf{(P\_8,1.1.1)} kva punaḥ alontyanivṛttyarthena arthaḥ sarvagrahaṇena . \\
\noindent \textbf{(P\_8,1.1.1)} nityavīpsayoḥ iti . \\
\noindent \textbf{(P\_8,1.1.1)} nityavīpsayoḥ iti ucyate na ca antyasya dvirvacanena nityatā vīpsā vā gamyate . \\
\noindent \textbf{(P\_8,1.1.1)} iha tarhi pareḥ varjane iti antyasya api dvirvacanena varjyamānatā gamyeta . \\
\noindent \textbf{(P\_8,1.1.1)} ṣaṣṭhīnirdeśārtham ca . ṣaṣṭhīnirdeśārtham ca sarvagrahaṇam kartavyam . \\
\noindent \textbf{(P\_8,1.1.1)} ṣaṣṭhīnirdeśaḥ yathā prakalpeta . \\
\noindent \textbf{(P\_8,1.1.1)} anirdeśe hi ṣaṣṭhyarthāprasiddhiḥ . akriyamāṇe sarvagrahaṇe ṣaṣṭhyarthasya aprasiddhiḥ syāt . \\
\noindent \textbf{(P\_8,1.1.1)} kasya . \\
\noindent \textbf{(P\_8,1.1.1)} sthāneyogatvasya . \\
\noindent \textbf{(P\_8,1.1.1)} kva punaḥ iha ṣaṣṭhīnirdeśārthena arthaḥ sarvagrahaṇena yāvatā sarvatra eva ṣaṣṭhī uccāryate . \\
\noindent \textbf{(P\_8,1.1.1)} parervarjane prasamupodaḥpādapūraṇe uparyadhyadhasaḥsāmīpye vākyāderāmantritasya iti . \\
\noindent \textbf{(P\_8,1.1.1)} iha na kā cit ṣaṣṭhī nityavīpsayoḥ iti . \\
\noindent \textbf{(P\_8,1.1.1)} nanu ca eṣā eva ṣaṣṭhī . \\
\noindent \textbf{(P\_8,1.1.1)} na eṣā ṣaṣṭhī . \\
\noindent \textbf{(P\_8,1.1.1)} kim tarhi . \\
\noindent \textbf{(P\_8,1.1.1)} arthanirdeśaḥ eṣaḥ . \\
\noindent \textbf{(P\_8,1.1.1)} nitye ca arthe vīpsāyām ca iti . \\
\noindent \textbf{(P\_8,1.1.1)} alontyanivṛttyarthena tāvat na arthaḥ sarvagrahaṇena . \\
\noindent \textbf{(P\_8,1.1.1)} idam tāvat ayam praṣṭavyaḥ . \\
\noindent \textbf{(P\_8,1.1.1)} nityavīpsayoḥ dve bhavataḥ iti ucyate dviśabdaḥ ādeśaḥ kasmāt na bhavati . \\
\noindent \textbf{(P\_8,1.1.1)} ācāryapravṛttiḥ jñāpayati na dviśabdaḥ ādeśaḥ bhavati iti yat ayam tasyaparamāmreḍitam anudāttamca iti āha . \\
\noindent \textbf{(P\_8,1.1.1)} katham kṛtvā jñāpakam . \\
\noindent \textbf{(P\_8,1.1.1)} dviśabdaḥ ayam ekāc tasya ekāctvāt tasyaparamāmreḍitam anudāttamca iti etat na asti . \\
\noindent \textbf{(P\_8,1.1.1)} paśyati tu ācāryaḥ na dviśabdaḥ ādeśaḥ bhavati iti tataḥ tasya paramāmreḍitam anudāttamca iti āha . \\
\noindent \textbf{(P\_8,1.1.1)} yadi tarhi na dviśabdaḥ ādeśaḥ bhavati ke tarhi idānīm dve bhavataḥ . \\
\noindent \textbf{(P\_8,1.1.1)} dviśabdena yat ucyate . \\
\noindent \textbf{(P\_8,1.1.1)} kim punaḥ tat . \\
\noindent \textbf{(P\_8,1.1.1)} dviśabdaḥ ayam saṅkhyāpadam saṅkhyāyāḥ ca saṅkhyeyam arthaḥ . \\
\noindent \textbf{(P\_8,1.1.1)} saṅkhyeye dve bhaviṣyataḥ . \\
\noindent \textbf{(P\_8,1.1.1)} ke punaḥ te . \\
\noindent \textbf{(P\_8,1.1.1)} pade vākye mātre vā . \\
\noindent \textbf{(P\_8,1.1.1)} tat yadā tāvat pade vākye vā tadā anekāltvāt sarvādeśaḥ siddhaḥ . \\
\noindent \textbf{(P\_8,1.1.1)} yadā mātre api tadā anekālśitsarvasya iti sarvādeśaḥ bhaviṣyati . \\
\noindent \textbf{(P\_8,1.1.1)} yadā tarhi ardhamātre tadā sarvādeśaḥ na sidhyati . \\
\noindent \textbf{(P\_8,1.1.1)} na eṣaḥ doṣaḥ . \\
\noindent \textbf{(P\_8,1.1.1)} na ca ardhamātre dviḥ ucyete . \\
\noindent \textbf{(P\_8,1.1.1)} kim kāraṇam . \\
\noindent \textbf{(P\_8,1.1.1)} iha vyākareṇe yaḥ sarvālpīyān svaravyavahāraḥ saḥ mātrayā bhavati na ardhamātrayā vyavahāraḥ asti . \\
\noindent \textbf{(P\_8,1.1.1)} tena ardhamātre na bhaviṣyataḥ . \\
\noindent \textbf{(P\_8,1.1.1)} evam api kutaḥ etat pade dve bhaviṣyataḥ iti na punaḥ vākye syātām mātre vā . \\
\noindent \textbf{(P\_8,1.1.1)} nityavīpsayoḥ dve bhavataḥ iti ucyate na ca vākyadvirvacanena mātrādvirvacanena vā nityatā vīpsā vā gamyate . \\
\noindent \textbf{(P\_8,1.1.1)} ṣaṣṭhīnirdeśārtham eva tarhi sarvagrahaṇam kartavyam . \\
\noindent \textbf{(P\_8,1.1.1)} na vā padādhikārāt . na vā vaktavyam . \\
\noindent \textbf{(P\_8,1.1.1)} kim kāraṇam . \\
\noindent \textbf{(P\_8,1.1.1)} padādhikārāt . \\
\noindent \textbf{(P\_8,1.1.1)} padasya iti prakṛtya dvirvacanam vakṣyāmi . \\
\noindent \textbf{(P\_8,1.1.1)} tat ca samāsataddhitavākyanivṛttyartham . tat ca avaśyam padagrahaṇam kartavyam samāsanivṛttyartham taddhitanivṛttyartham vākyanivṛttyartham ca . \\
\noindent \textbf{(P\_8,1.1.1)} samāsanivṛttyartham tāvat . \\
\noindent \textbf{(P\_8,1.1.1)} saptaparṇaḥ aṣṭāpadam . \\
\noindent \textbf{(P\_8,1.1.1)} taddhitanivṛttyartham . \\
\noindent \textbf{(P\_8,1.1.1)} dvipadikā tripadikā . \\
\noindent \textbf{(P\_8,1.1.1)} māṣaśaḥ kārṣāpaṇaśaḥ . \\
\noindent \textbf{(P\_8,1.1.1)} vākyanivṛttyartham . \\
\noindent \textbf{(P\_8,1.1.1)} grāme grāme pānīyam . \\
\noindent \textbf{(P\_8,1.1.1)} māṣam māṣam dehi . \\
\noindent \textbf{(P\_8,1.1.1)} atha kriyamāṇe api vai padagrahaṇe samāsanivṛttyartham iti katham idam vijñāyate . \\
\noindent \textbf{(P\_8,1.1.1)} samasasya nivṛttyartham samāsanivṛttyartham iti . \\
\noindent \textbf{(P\_8,1.1.1)} āhosvit samāse nivṛttyartham samāsanivṛttyartham iti . \\
\noindent \textbf{(P\_8,1.1.1)} kim ca ataḥ . \\
\noindent \textbf{(P\_8,1.1.1)} yadi viñāyate samāsasya nivṛttyartham samāsanivṛttyartham iti siddham saptaparṇaḥ saptaparṇau saptaparṇāḥ iti saptaparṇābhyām saptaparṇebhyaḥ iti atra prāpnoti . \\
\noindent \textbf{(P\_8,1.1.1)} atha vijñāyate samāse nivṛttyartham samāsanivṛttyartham iti saptaparṇaḥ saptaparṇau saptaparṇāḥ iti atra api prāpnoti . \\
\noindent \textbf{(P\_8,1.1.1)} tathā taddhitanivṛttyartham iti . \\
\noindent \textbf{(P\_8,1.1.1)} katham idam vijñāyate . \\
\noindent \textbf{(P\_8,1.1.1)} taddhitasya nivṛttyartham taddhitanivṛttyartham iti . \\
\noindent \textbf{(P\_8,1.1.1)} āhosvit taddhite nivṛttyartham taddhitanivṛttyartham iti . \\
\noindent \textbf{(P\_8,1.1.1)} kim ca ataḥ . \\
\noindent \textbf{(P\_8,1.1.1)} yadi vijñāyate taddhitasya nivṛttyartham taddhitanivṛttyartham iti siddham dvipadikāḥ tripadikāḥ dvipadikābhyām tripadikābhyām māṣaśaḥ kārṣāpaṇaśaḥ iti atra prāpnoti . \\
\noindent \textbf{(P\_8,1.1.1)} atha vijñāyate taddhite nivṛttyartham taddhitanivṛttyartham iti dvipadikāḥ tripadikāḥ iti atra api prāpnoti . \\
\noindent \textbf{(P\_8,1.1.1)} tathā vākyanivṛttyartham iti katham idam vijñāyate . \\
\noindent \textbf{(P\_8,1.1.1)} vākyasya nivṛttyartham vākyanivṛttyartham iti . \\
\noindent \textbf{(P\_8,1.1.1)} āhosvit vākye nivṛttyartham vākyanivṛttyartham iti . \\
\noindent \textbf{(P\_8,1.1.1)} kim ca ataḥ . \\
\noindent \textbf{(P\_8,1.1.1)} yadi vijñāyate vākyasya nivṛttyartham vākyanivṛttyartham iti yadi vākyam vīpsāyuktam bhavitavyam eva dvirvacanena . \\
\noindent \textbf{(P\_8,1.1.1)} atha api avayavaḥ bhavatu eva . \\
\noindent \textbf{(P\_8,1.1.1)} tat etat kriyamāṇe api padagrahaṇe ālūnaviśīrṇam bhavati . \\
\noindent \textbf{(P\_8,1.1.1)} kim cit saṅgṛhītam kim cit asaṅgṛhītam . \\
\noindent \textbf{(P\_8,1.1.1)} sagatigrahaṇam ca . sagatigrahaṇam ca kartavyam . \\
\noindent \textbf{(P\_8,1.1.1)} prapacati prapacati . \\
\noindent \textbf{(P\_8,1.1.1)} prakaroti prakaroti iti . \\
\noindent \textbf{(P\_8,1.1.1)} kim punaḥ kāraṇam na sidhyati . \\
\noindent \textbf{(P\_8,1.1.1)} na hi sagatikam padam bhavati . \\
\noindent \textbf{(P\_8,1.1.1)} samāsanivṛttyarthena tāvat na arthaḥ padagrahaṇena . \\
\noindent \textbf{(P\_8,1.1.1)} samāsena uktatvāt vīpsāyāḥ dvirvacanam na bhaviṣyati . \\
\noindent \textbf{(P\_8,1.1.1)} kim ca bhoḥ samāsaḥ vīpsāyām iti ucyate . \\
\noindent \textbf{(P\_8,1.1.1)} na khalu vīpsāyām iti ucyate gamyate tu saḥ arthaḥ . \\
\noindent \textbf{(P\_8,1.1.1)} tatra uktaḥ samāsena iti kṛtvā dvirvacanam na bhaviṣyati . \\
\noindent \textbf{(P\_8,1.1.1)} yatra ca samāsena anuktā vīpsā bhavati tatra dvirvacanam . \\
\noindent \textbf{(P\_8,1.1.1)} tat yathā . \\
\noindent \textbf{(P\_8,1.1.1)} ekaikavicitāḥ anyonyasahāyāḥ iti . \\
\noindent \textbf{(P\_8,1.1.1)} atha vā yat atra vīpsāyuktam na adaḥ prayujyate . \\
\noindent \textbf{(P\_8,1.1.1)} kim punaḥ tat . \\
\noindent \textbf{(P\_8,1.1.1)} parvaṇi parvaṇi sapta parṇāni asya . \\
\noindent \textbf{(P\_8,1.1.1)} paṅktau paṅktau aṣṭau padāni asya iti . \\
\noindent \textbf{(P\_8,1.1.1)} taddhitanivṛttyarthena ca api na arthaḥ padagrahaṇena . \\
\noindent \textbf{(P\_8,1.1.1)} taddhitena ukatvāt vīpsāyāḥ dvirvacanam na bhaviṣyati . \\
\noindent \textbf{(P\_8,1.1.1)} taddhitaḥ khalu api vīpsāyām iti ucyate . \\
\noindent \textbf{(P\_8,1.1.1)} yatra ca taddhitena anuktā vīpsā bhavati tatra dvirvacanam . \\
\noindent \textbf{(P\_8,1.1.1)} tat yathā . \\
\noindent \textbf{(P\_8,1.1.1)} ekaikaśaḥ dadāti iti . \\
\noindent \textbf{(P\_8,1.1.1)} vākyanivṛttyarthena ca api na arthaḥ padagrahaṇena . \\
\noindent \textbf{(P\_8,1.1.1)} padadvirvacanena uktatvāt vīpsāyāḥ vākyadvirvacanam na bhaviṣyati . \\
\noindent \textbf{(P\_8,1.1.1)} yatra ca padadvirvacanena anuktā vīpsā bhavati tatra dvirvacanam . \\
\noindent \textbf{(P\_8,1.1.1)} tat yathā . \\
\noindent \textbf{(P\_8,1.1.1)} prapacati prapacati . \\
\noindent \textbf{(P\_8,1.1.1)} prakaroti prakaroti . \\
\noindent \textbf{(P\_8,1.1.1)} uttarārtham tarhi padagrahaṇam kartavyam . \\
\noindent \textbf{(P\_8,1.1.1)} tasyaparamāmreḍitam anudāttamca iti vakṣyati tat padadvirvacane yathā syāt vākyadvirvacanemā bhūt . \\
\noindent \textbf{(P\_8,1.1.1)} mahyam grahīṣyati mahyam grahīṣyati . \\
\noindent \textbf{(P\_8,1.1.1)} mām abhivyāhariṣyati mām abhvyāhariṣyati . \\
\noindent \textbf{(P\_8,1.1.1)} katham ca atra dvirvacanam . \\
\noindent \textbf{(P\_8,1.1.1)} chāndasatvāt . \\
\noindent \textbf{(P\_8,1.1.1)} svaraḥ api tarhi chāndasatvāt eva na bhaviṣyati . \\
\noindent \textbf{(P\_8,1.1.1)} uttarārtham eva tarhi padagrahaṇam kartavyam . \\
\noindent \textbf{(P\_8,1.1.1)} padasya padāt iti vakṣyati tat padagrahaṇam na kartavyam bhavati . \\
\noindent \textbf{(P\_8,1.1.1)} sarvagrahaṇam api tarhi uttarārtham . \\
\noindent \textbf{(P\_8,1.1.1)} anudāttaṃsarvamapādādau iti vakṣyati tat sarvagrahaṇam kartavyam bhavati . \\
\noindent \textbf{(P\_8,1.1.1)} ubhayam kriyate tatra eva \\
\noindent \textbf{(P\_8,1.1.2)} ihārtham eva tarhi ṣaṣṭhīnirdeśārtham anyatarat kartavyam . \\
\noindent \textbf{(P\_8,1.1.2)} ṣaṣṭhīnirdiṣṭasya sthāne dvirvacanam yathā syāt dviḥprayogaḥ mā bhūt iti . \\
\noindent \textbf{(P\_8,1.1.2)} kim ca syāt . \\
\noindent \textbf{(P\_8,1.1.2)} ām pacasi devadattā3 āmaekāntaramāmantritamanantike iti ekāntaratā na syāt . \\
\noindent \textbf{(P\_8,1.1.2)} iha ca paunaḥpunyam paunaḥpunikam iti aprātipadikatvāt taddhitotpattiḥ na syāt . \\
\noindent \textbf{(P\_8,1.1.2)} yadi tarhi sthāne dvirvacanam rājā rājā vāk vāk padasya iti nalopādīni na sidhyanti . \\
\noindent \textbf{(P\_8,1.1.2)} idam iha sampradhāryam . \\
\noindent \textbf{(P\_8,1.1.2)} dvirvacanam kriyatām nalopādīni iti kim atra kartavyam . \\
\noindent \textbf{(P\_8,1.1.2)} paratvāt nalopādīni . \\
\noindent \textbf{(P\_8,1.1.2)} pūrvatra asiddhe nalopādīni siddhāsiddhayoḥ ca na asti sampradhāraṇā . \\
\noindent \textbf{(P\_8,1.1.2)} evam tarhi pūrvatra asiddhīyam advirvacane iti vakṣyāmi . \\
\noindent \textbf{(P\_8,1.1.2)} tat ca avaśyam vaktavyam . \\
\noindent \textbf{(P\_8,1.1.2)} kim prayojanam . \\
\noindent \textbf{(P\_8,1.1.2)} vibhāṣitāḥ prayojayanti . \\
\noindent \textbf{(P\_8,1.1.2)} drogdhā drogdhā . \\
\noindent \textbf{(P\_8,1.1.2)} droḍhā droḍhā iti . \\
\noindent \textbf{(P\_8,1.1.2)} iha tarhi bisam bisam musalam musalam ādeśapratyayayoḥ iti ṣatvam prāpnoti . \\
\noindent \textbf{(P\_8,1.1.2)} ādeśaḥ yaḥ sakāraḥ pratayaḥ yaḥ sakāraḥ iti evam etat vijñāyate . \\
\noindent \textbf{(P\_8,1.1.2)} iha tarhi nṛbhiḥ nṛbhiḥ raṣābhyānnoṇaḥsamānapade iti ṇatvam prāpnoti . \\
\noindent \textbf{(P\_8,1.1.2)} samānapade iti ucyate samānam eva yat nityam na ca etat nityam samānapadam eva . \\
\noindent \textbf{(P\_8,1.1.2)} kim vaktavyam etat . \\
\noindent \textbf{(P\_8,1.1.2)} na hi . \\
\noindent \textbf{(P\_8,1.1.2)} katham anucyamānam gaṃsyate . \\
\noindent \textbf{(P\_8,1.1.2)} samānagrahaṇasāmarthyāt . \\
\noindent \textbf{(P\_8,1.1.2)} yadi hi yat samānam ca asamānam ca tatra syāt samānagrahaṇam anarthakaṃ syāt \\
\noindent \textbf{(P\_8,1.1.3)} atha vā punaḥ astu dviḥprayogaḥ dvirvacanam . \\
\noindent \textbf{(P\_8,1.1.3)} nanu ca uktam ām pacasi pacasi devadattā3 āmaḥ ekāntaram āmantritam anantike iti ekāntaratā na prāpnoti . \\
\noindent \textbf{(P\_8,1.1.3)} na eṣaḥ doṣaḥ . \\
\noindent \textbf{(P\_8,1.1.3)} suptiṅbhyām padam viśeṣayiṣyāmaḥ . \\
\noindent \textbf{(P\_8,1.1.3)} suptiṅantampadam . \\
\noindent \textbf{(P\_8,1.1.3)} yasmāt suptiṅvidhiḥ tadādi suptiṅantam ca . \\
\noindent \textbf{(P\_8,1.1.3)} nanu ca ekaikasmāt eva atra suptiṅvidhiḥ . \\
\noindent \textbf{(P\_8,1.1.3)} samudāye yā vākyaparisamāptiḥ tayā padasañjñā . \\
\noindent \textbf{(P\_8,1.1.3)} kutaḥ etat . \\
\noindent \textbf{(P\_8,1.1.3)} śāstrāhāneḥ . \\
\noindent \textbf{(P\_8,1.1.3)} evam hi śāstram ahīnam bhavati . \\
\noindent \textbf{(P\_8,1.1.3)} yat api ucyate iha paunaḥpunyam paunaḥpunikam iti aprātipadikatvāt taddhitotpattiḥ na prāpnoti iti mā bhūt evam . \\
\noindent \textbf{(P\_8,1.1.3)} samarthāt iti evam bhaviṣyati . \\
\noindent \textbf{(P\_8,1.1.3)} atha vā ācāryapravṛttiḥ jñāpayati bhavati evañjātīyakebhyaḥ taddhitotpattiḥ iti yat ayam kaskādiṣu kautaskutaśabdam paṭhati \\
\noindent \textbf{(P\_8,1.4.1)} iha kasmāt na bhavati . \\
\noindent \textbf{(P\_8,1.4.1)} himavān khāṇḍavaḥ pāriyātraḥ samudraḥ iti . \\
\noindent \textbf{(P\_8,1.4.1)} nitye dve bhavataḥ iti prāpnoti . \\
\noindent \textbf{(P\_8,1.4.1)} na eṣaḥ doṣaḥ . \\
\noindent \textbf{(P\_8,1.4.1)} ayam nityaśabdaḥ asti eva kūṭastheṣu avicāliṣu bhāveṣu vartate . \\
\noindent \textbf{(P\_8,1.4.1)} tat yathā : nityā dyauḥ nityā pṛthivī nityam ākāśam iti . \\
\noindent \textbf{(P\_8,1.4.1)} asti ābhīkṣṇye vartate . \\
\noindent \textbf{(P\_8,1.4.1)} tat yathā : nityaprahasitaḥ nityaprajalpitaḥ iti . \\
\noindent \textbf{(P\_8,1.4.1)} tat yaḥ ābhīkṣṇye vartate tasya idam grahaṇam . \\
\noindent \textbf{(P\_8,1.4.2)} atha kim idam vīpsā iti . \\
\noindent \textbf{(P\_8,1.4.2)} āpnoteḥ ayam vipūrvāi icchāyām arthe san vidhīyate . \\
\noindent \textbf{(P\_8,1.4.2)} yadi evam cikīrṣati jihīrṣati iti atra api prāpnoti . \\
\noindent \textbf{(P\_8,1.4.2)} na eṣaḥ doṣaḥ . \\
\noindent \textbf{(P\_8,1.4.2)} na evam vijñāyate vīpsāyām abhidheyāyām iti . \\
\noindent \textbf{(P\_8,1.4.2)} katham tarhi . \\
\noindent \textbf{(P\_8,1.4.2)} kartṛviśeṣaṇam etat . \\
\noindent \textbf{(P\_8,1.4.2)} vīpsati iti vīpsaḥ . \\
\noindent \textbf{(P\_8,1.4.2)} vīpsaḥ cet kartā bhavati iti . \\
\noindent \textbf{(P\_8,1.4.2)} kaḥ punaḥ vīpsārthaḥ . \\
\noindent \textbf{(P\_8,1.4.2)} anavayavābhidhānam vīpsārthaḥ . \\
\noindent \textbf{(P\_8,1.4.2)} anavayavena dravyāṇām abhidhānam eṣaḥ vīpsārthaḥ . \\
\noindent \textbf{(P\_8,1.4.2)} anavayavābhidhānam vīpsārthaḥ iti cet jātyākhyāyām dvirvacanaprasaṅgaḥ . \\
\noindent \textbf{(P\_8,1.4.2)} anavayavābhidhānam vīpsārthaḥ iti cet jātyākhyāyām dvirvacanam prāpnoti . \\
\noindent \textbf{(P\_8,1.4.2)} vrīhibhiḥ yavaiḥ vā iti . \\
\noindent \textbf{(P\_8,1.4.2)} na vā ekārthatvāt jāteḥ . na vā eṣaḥ doṣaḥ . \\
\noindent \textbf{(P\_8,1.4.2)} kim kāraṇam . \\
\noindent \textbf{(P\_8,1.4.2)} ekārthatvāt jāteḥ . \\
\noindent \textbf{(P\_8,1.4.2)} ekārthaḥ hi jātiḥ . \\
\noindent \textbf{(P\_8,1.4.2)} ekam artham pratyāyayiṣyāmi iti jātiśabdaḥ prayujyate . \\
\noindent \textbf{(P\_8,1.4.2)} anekārthāśrayatvāt ca vīpsāyāḥ . anekārthāśrayā ca punaḥ vīpsā . \\
\noindent \textbf{(P\_8,1.4.2)} anekam artham sampratyāyayiṣyāmi iti vīpsā prayujyate . \\
\noindent \textbf{(P\_8,1.4.2)} ekārthatvāt jāteḥ anekārthāśrayatvāt ca vīpsāyāḥ jātyākhyāyām dvirvacanam na bhaviṣyati . \\
\noindent \textbf{(P\_8,1.4.2)} nivartakatvādvā . atha vā na anena dvirvacanam nirvartyate . \\
\noindent \textbf{(P\_8,1.4.2)} kim tarhi advirvacanam anena nivartyate . \\
\noindent \textbf{(P\_8,1.4.2)} yāvantaḥ te arthāḥ tāvatām śabdānām prayogaḥ prāpnoti . \\
\noindent \textbf{(P\_8,1.4.2)} tatra anena nivṛttiḥ kriyate . \\
\noindent \textbf{(P\_8,1.4.2)} nityavīpsayoḥ arthayoḥ dve eva śabdarūpe prayoktavye na atibahu prayoktavyam iti . \\
\noindent \textbf{(P\_8,1.4.2)} sarvapadasagatigrahaṇānarthakyam ca arthābhidhāne dvirvacanavidhānāt . sarvagrahaṇam ca anarthakam . \\
\noindent \textbf{(P\_8,1.4.2)} kim kāraṇam . \\
\noindent \textbf{(P\_8,1.4.2)} sarvasya eva hi dvirvacanena arthaḥ gamyate na avayavasya . \\
\noindent \textbf{(P\_8,1.4.2)} padagrahaṇam ca anarthakam padasya eva hi dvirvacanena arthaḥ gamyate na agatikasya \\
\noindent \textbf{(P\_8,1.4.3)} kim punaḥ idam vīpsāyām sarvam abhidhīyate āhosvit ekam . \\
\noindent \textbf{(P\_8,1.4.3)} kaḥ ca atra viśeṣaḥ . \\
\noindent \textbf{(P\_8,1.4.3)} vīpsāyām sarvābhidhāne vacanāprasiddhiḥ . \\
\noindent \textbf{(P\_8,1.4.3)} vīpsāyām sarvābhidhāne vacanam na sidhyati . \\
\noindent \textbf{(P\_8,1.4.3)} grāmaḥ grāmaḥ . \\
\noindent \textbf{(P\_8,1.4.3)} janapadaḥ janapadaḥ . \\
\noindent \textbf{(P\_8,1.4.3)} bahavaḥ te arthāḥ tatra bahuṣubahuvacanam iti bahuvacanam prāpnoti . \\
\noindent \textbf{(P\_8,1.4.3)} astu tarhi ekam . \\
\noindent \textbf{(P\_8,1.4.3)} ekābhidhāne asarvadravyagatiḥ . ekābhidhāne sarvadravyagatiḥ na sidhyati . \\
\noindent \textbf{(P\_8,1.4.3)} astu tarhi sarvam . \\
\noindent \textbf{(P\_8,1.4.3)} nanu ca uktam vīpsāyām sarvābhidhāne vacanāprasiddhiḥ iti . \\
\noindent \textbf{(P\_8,1.4.3)} na vā padārthatvāt . na vā eṣaḥ doṣaḥ . \\
\noindent \textbf{(P\_8,1.4.3)} kim kāraṇam . \\
\noindent \textbf{(P\_8,1.4.3)} padārthatvāt . \\
\noindent \textbf{(P\_8,1.4.3)} padasya arthaḥ vīpsā . \\
\noindent \textbf{(P\_8,1.4.3)} subantam ca padam ṅyāpprātipadikāt ca ekatvādiṣu artheṣu svādayaḥ vidhīyante na ca etat prātipadikam . \\
\noindent \textbf{(P\_8,1.4.3)} yat tarhi prātipadikam : dṛṣat dṛṣat samit samit iti . \\
\noindent \textbf{(P\_8,1.4.3)} etat api pratyayalakṣaṇena subantam na prātipadikam . \\
\noindent \textbf{(P\_8,1.4.3)} apara āha . \\
\noindent \textbf{(P\_8,1.4.3)} na vā padārthatvāt . \\
\noindent \textbf{(P\_8,1.4.3)} na vā eṣaḥ doṣaḥ . \\
\noindent \textbf{(P\_8,1.4.3)} kim kāraṇam . \\
\noindent \textbf{(P\_8,1.4.3)} padārthatvāt . \\
\noindent \textbf{(P\_8,1.4.3)} padasya arthaḥ vīpsā subantam ca padam ṅyāpprātipadikāt ca ekatvādiṣu artheṣu svādayaḥ vidhīyante na ca etat prātipadikam . \\
\noindent \textbf{(P\_8,1.4.3)} yat tarhi prātipadikam . \\
\noindent \textbf{(P\_8,1.4.3)} dṛṣat dṛṣat samit samit iti . \\
\noindent \textbf{(P\_8,1.4.3)} etat api pratyayalakṣaṇena subantam na prātipadikam \\
\noindent \textbf{(P\_8,1.4.4)} atha iha katham bhavitavyam . \\
\noindent \textbf{(P\_8,1.4.4)} pacati pacatitarām tiṣṭhati . \\
\noindent \textbf{(P\_8,1.4.4)} āhosvit pacatitarām pacatitarām tiṣṭhati iti . \\
\noindent \textbf{(P\_8,1.4.4)} pacati pacatitarām tiṣṭhatīti bhavitavyam . \\
\noindent \textbf{(P\_8,1.4.4)} katham . \\
\noindent \textbf{(P\_8,1.4.4)} dvirvacanam kriyatām ātiśāyikaḥ iti dvirvacanam bhaviṣyati vipratiṣedhena . \\
\noindent \textbf{(P\_8,1.4.4)} iha api tarhi ātiśāyikāt dvirvacanam syāt . \\
\noindent \textbf{(P\_8,1.4.4)} ādyataram ādyataram ānaya iti . \\
\noindent \textbf{(P\_8,1.4.4)} asti atra viśeṣaḥ . \\
\noindent \textbf{(P\_8,1.4.4)} antaraṅgaḥ ātiśāyikaḥ . \\
\noindent \textbf{(P\_8,1.4.4)} kā antaraṅgatā . \\
\noindent \textbf{(P\_8,1.4.4)} ṅyāpprātipadikāt ātiśāyikaḥ padasya dvirvacanam . \\
\noindent \textbf{(P\_8,1.4.4)} ātiśāyikaḥ api na antaraṅgaḥ . \\
\noindent \textbf{(P\_8,1.4.4)} katham . \\
\noindent \textbf{(P\_8,1.4.4)} samarthāt taddhitaḥ asau utpadyate sāmarthyam ca subantena . \\
\noindent \textbf{(P\_8,1.4.4)} atha vā spardhāyām ātiśāyikaḥ vidhīyate na ca antareṇa pratiyoginam spardhā gamyate . \\
\noindent \textbf{(P\_8,1.4.4)} evam tarhi iha dvau arthau vaktavyau nityavīpse ca atiśayaḥ ca na ca ekasya prayoktuḥ anekam artham yugapat vaktum sambhavaḥ asti . \\
\noindent \textbf{(P\_8,1.4.4)} tat etat prayoktari adhīnam bhavati . \\
\noindent \textbf{(P\_8,1.4.4)} etasmin ca prayoktari adhīne kva cit kā cit prasṛtatarā gatiḥ bhavati . \\
\noindent \textbf{(P\_8,1.4.4)} iha tāvat pacati pacatitarām tiṣṭhati iti eṣā prasṛtatarā gatiḥ yat nityam uktvā atiśayaḥ ucyate . \\
\noindent \textbf{(P\_8,1.4.4)} iha idānīm ādyataram ādyataram ānaya iti eṣā prasṛtatarā gatiḥ yat atiśayam uktvā vīpsādvirvacanam ucyate \\
\noindent \textbf{(P\_8,1.5)} pareḥ asamāse . \\
\noindent \textbf{(P\_8,1.5)} pareḥ asamāse iti vaktavyam . \\
\noindent \textbf{(P\_8,1.5)} iha mā bhūt . \\
\noindent \textbf{(P\_8,1.5)} paritrigartam vṛṣṭaḥ devaḥ . \\
\noindent \textbf{(P\_8,1.5)} tat tarhi vaktavyam . \\
\noindent \textbf{(P\_8,1.5)} na vaktavyam . \\
\noindent \textbf{(P\_8,1.5)} pareḥ varjane iti ucyate na ca atra pariḥ varjane vartate . \\
\noindent \textbf{(P\_8,1.5)} kaḥ tarhi . \\
\noindent \textbf{(P\_8,1.5)} samāsaḥ . \\
\noindent \textbf{(P\_8,1.5)} parervarjane vāvacanam . parervarvane vā iti vaktavyam . \\
\noindent \textbf{(P\_8,1.5)} pari trigartebhyaḥ vṛṣṭaḥ devaḥ . \\
\noindent \textbf{(P\_8,1.5)} pari pari trigartebhyaḥ vṛṣṭaḥ devaḥ \\
\noindent \textbf{(P\_8,1.8)} asūyākutsanayoḥ kopabhartsanayoḥ ca ekārthatvāt pṛthaktvanirdeśānarthakyam . \\
\noindent \textbf{(P\_8,1.8)} asūyā kutsanam iti ekaḥ arthaḥ . \\
\noindent \textbf{(P\_8,1.8)} kopaḥ bhartsanam iti ekaḥ arthaḥ . \\
\noindent \textbf{(P\_8,1.8)} asūyākutsanayoḥ kopabhartsanayoḥ ca ekārthatvāt pṛthaktvanirdeśaḥ anarthakaḥ . \\
\noindent \textbf{(P\_8,1.8)} na hi anasūyan kutsayati na ca api akupitaḥ bhartsayate . \\
\noindent \textbf{(P\_8,1.8)} nanu ca bhoḥ akupitāḥ api dṛśyante dārakān bhartsayamānāḥ . \\
\noindent \textbf{(P\_8,1.8)} antataḥ te tām śarīrākṛtim kurvanti yā kupitasya bhavati . \\
\noindent \textbf{(P\_8,1.8)} evam tarhi āha . \\
\noindent \textbf{(P\_8,1.8)} sāmṛtaiḥ pāṇibhiḥ ghnanti guravaḥ na viṣokitaiḥ . \\
\noindent \textbf{(P\_8,1.8)} lāḍanāśrayiṇaḥ doṣāḥ tāḍanāśrayiṇaḥ guṇāḥ . \\
\noindent \textbf{(P\_8,1.8)} ekam bahuvrīhivat \\
\noindent \textbf{(P\_8,1.9)} iha kasmāt bahuvrīhivadbhāvaḥ na bhavati . \\
\noindent \textbf{(P\_8,1.9)} ekaḥ iti . \\
\noindent \textbf{(P\_8,1.9)} ekasya dvirvacanasambandhena bahuvrīhivadbhāvaḥ ucyate na ca atra dvirvacanam paśyāmaḥ . \\
\noindent \textbf{(P\_8,1.9)} ekasya dvirvacanasambandhena iti cet arthanirdeśaḥ . \\
\noindent \textbf{(P\_8,1.9)} ekasya dvirvacanasambandhena iti cet arthanirdeśaḥ kartavyaḥ . \\
\noindent \textbf{(P\_8,1.9)} dvirvacanam api hi atra kasmāt na bhavati . \\
\noindent \textbf{(P\_8,1.9)} tasmāt vācyam asmin arthe dve bhavataḥ bahuvrīhivat ca iti . \\
\noindent \textbf{(P\_8,1.9)} na vā vīpsādhikārāt . na vā vaktavyam . \\
\noindent \textbf{(P\_8,1.9)} kim kāraṇam . \\
\noindent \textbf{(P\_8,1.9)} vīpsādhikārāt . \\
\noindent \textbf{(P\_8,1.9)} nityavīpsayoḥ iti vartate \\
\noindent \textbf{(P\_8,1.9)} atha bahuvrīhivattve kim prayojanam . \\
\noindent \textbf{(P\_8,1.9)} bahuvrīhivattve prayojanam sublopapuṃvadbhāvau . \\
\noindent \textbf{(P\_8,1.9)} sublopaḥ . \\
\noindent \textbf{(P\_8,1.9)} ekaikam . \\
\noindent \textbf{(P\_8,1.9)} puṃvadbhāvaḥ . \\
\noindent \textbf{(P\_8,1.9)} gatagatā . \\
\noindent \textbf{(P\_8,1.9)} yadi evam sarvanāmasvarasamāsānteṣu doṣaḥ . sarvanāmasvarasamāsānteṣu doṣaḥ bhavati . \\
\noindent \textbf{(P\_8,1.9)} sarvanāmavidhau doṣaḥ bhavati . \\
\noindent \textbf{(P\_8,1.9)} ekaikasmai . \\
\noindent \textbf{(P\_8,1.9)} nabahuvrīhau iti pratiṣedhaḥ prāpnoti . \\
\noindent \textbf{(P\_8,1.9)} sarvanāma . \\
\noindent \textbf{(P\_8,1.9)} svara . \\
\noindent \textbf{(P\_8,1.9)} nana susu . \\
\noindent \textbf{(P\_8,1.9)} nañsubhyām iti eṣaḥ svaraḥ prāpnoti . \\
\noindent \textbf{(P\_8,1.9)} svara . \\
\noindent \textbf{(P\_8,1.9)} samāsānta . \\
\noindent \textbf{(P\_8,1.9)} ṛgṛk pūḥpūḥ . \\
\noindent \textbf{(P\_8,1.9)} ṛkpūrabdhūḥpathāmānakṣe iti samāsāntaḥ prāpnoti . \\
\noindent \textbf{(P\_8,1.9)} sarvanāmavidhau tāvat na doṣaḥ . \\
\noindent \textbf{(P\_8,1.9)} uktam tatra bahuvrīhigrahaṇasya prayojanam bahuvrīhiḥ eva yaḥ buhuvrīhiḥ tatra pratiṣedhaḥ yathā syāt bahuvrīhivadbhāvena yaḥ bahuvrīhiḥ tatra mā bhūt iti . \\
\noindent \textbf{(P\_8,1.9)} svarasamāsāntayoḥ api prakṛtam samāsagrahaṇam anuvartate tena eva bahuvrīhim viśeṣayiṣyāmaḥ . \\
\noindent \textbf{(P\_8,1.9)} samāsaḥ yaḥ bahuvrīhiḥ iti \\
\noindent \textbf{(P\_8,1.11)} karmadhārayavattve kāni prayojanāni . \\
\noindent \textbf{(P\_8,1.11)} karmadhārayavattve prayojanam sublopapuṃvadbhāvāntodāttatvāni . \\
\noindent \textbf{(P\_8,1.11)} sublopaḥ . \\
\noindent \textbf{(P\_8,1.11)} paṭupaṭuḥ . \\
\noindent \textbf{(P\_8,1.11)} puṃvadbhāvaḥ . \\
\noindent \textbf{(P\_8,1.11)} paṭupaṭvī . \\
\noindent \textbf{(P\_8,1.11)} antodāttatvam . \\
\noindent \textbf{(P\_8,1.11)} paṭupaṭuḥ \\
\noindent \textbf{(P\_8,1.12.1)} guṇavacanasya iti kimartham . \\
\noindent \textbf{(P\_8,1.12.1)} agniḥ māṇavakaḥ . \\
\noindent \textbf{(P\_8,1.12.1)} gauḥ vāhīkaḥ . \\
\noindent \textbf{(P\_8,1.12.1)} prakāre sarveṣām guṇavacanatvāt sarvaprasaṅgaḥ . \\
\noindent \textbf{(P\_8,1.12.1)} sarve hi śabdāḥ prakāre vartamānāḥ guṇavacanāḥ sampadyante tena iha api prāpnoti . \\
\noindent \textbf{(P\_8,1.12.1)} agniḥ māṇavakaḥ . \\
\noindent \textbf{(P\_8,1.12.1)} gauḥ vāhīkaḥ iti . \\
\noindent \textbf{(P\_8,1.12.1)} siddham tu prakṛtyarthaviśeṣaṇatvāt . siddham etat . \\
\noindent \textbf{(P\_8,1.12.1)} katham . \\
\noindent \textbf{(P\_8,1.12.1)} prakṛtyarthaviśeṣaṇatvāt . \\
\noindent \textbf{(P\_8,1.12.1)} prakṛtyarthaḥ viśeṣyate . \\
\noindent \textbf{(P\_8,1.12.1)} na evam vijñāyate prakāre guṇavacanasya iti . \\
\noindent \textbf{(P\_8,1.12.1)} katham tarhi . \\
\noindent \textbf{(P\_8,1.12.1)} guṇavacanasya śabdasya dve bhavataḥ prakāre vartamānasya iti . \\
\noindent \textbf{(P\_8,1.12.1)} atha vā prakāre guṇavacanasya iti ucyate sarvaḥ ca śabdaḥ prakāre vartamānaḥ guṇavacanaḥ sampadyate tatra prakarṣagatiḥ vijñāsyate : sādhīyaḥ yaḥ guṇavacanaḥ iti . \\
\noindent \textbf{(P\_8,1.12.1)} kaḥ ca sādhīyaḥ . \\
\noindent \textbf{(P\_8,1.12.1)} yaḥ prakāre ca prāk ca prakārāt . \\
\noindent \textbf{(P\_8,1.12.1)} atha vā prakāre guṇavacanasya iti ucyate sarvaḥ ca śabdaḥ prakāre vartamānaḥ guṇavacanaḥ sampadyate te evam vijñāsyāmaḥ prāk prakārāt yaḥ guṇavacanaḥ iti \\
\noindent \textbf{(P\_8,1.12.2)} [ānupūrvye .] ānupūrvye dve bhavataḥ iti vaktavyam . \\
\noindent \textbf{(P\_8,1.12.2)} mūle mūle sthūlāḥ . \\
\noindent \textbf{(P\_8,1.12.2)} agre agre sūkṣmāḥ . \\
\noindent \textbf{(P\_8,1.12.2)} svārthe avadhāryamāṇe anekasmin . svārthe avadhāryamāṇe anekasmin dve bhavataḥ iti vaktavyam . \\
\noindent \textbf{(P\_8,1.12.2)} asmāt kārṣāpaṇāt iha bhavadbhyām māṣam māṣam dehi . \\
\noindent \textbf{(P\_8,1.12.2)} avadhāryamāṇe iti kimartham . \\
\noindent \textbf{(P\_8,1.12.2)} asmāt kārṣāpaṇāt iha bhavadbhyām māṣam dehi dvau dehi trīn dehi . \\
\noindent \textbf{(P\_8,1.12.2)} anekasmin iti kimartham . \\
\noindent \textbf{(P\_8,1.12.2)} asmāt kārṣāpaṇāt iha bhavadbhyām māṣam dehi . \\
\noindent \textbf{(P\_8,1.12.2)} māṣam eva dehi . \\
\noindent \textbf{(P\_8,1.12.2)} kim punaḥ kāraṇam na sidhyati . \\
\noindent \textbf{(P\_8,1.12.2)} anavayavābhidhānam vīpsārthaḥ iti ucyate avayavābhidhānam ca atra gamyate . \\
\noindent \textbf{(P\_8,1.12.2)} ātaḥ ca avayavābhidhānam yaḥ hi ucyate asmāt kārṣāpaṇāt iha bhavadbhyām māṣam māṣam dehi iti māṣam māṣam asau dattvā śeṣam pṛcchati kim anena kriyatām iti. yaḥ punaḥ ucyate imam kārṣāpaṇam iha bhavadbhyām māṣam māṣam dehi iti māṣam māṣam asau dattvā tūṣṇīm āste . \\
\noindent \textbf{(P\_8,1.12.2)} [cāpale .] cāpale dve bhavataḥ iti vaktavyam . \\
\noindent \textbf{(P\_8,1.12.2)} ahiḥ ahiḥ budhyasva budhyasva . \\
\noindent \textbf{(P\_8,1.12.2)} na ca avaśyam dve eva . \\
\noindent \textbf{(P\_8,1.12.2)} yāvadbhiḥ śabdaiḥ saḥ arthaḥ gamyate tāvantaḥ prayokavyāḥ . \\
\noindent \textbf{(P\_8,1.12.2)} ahiḥ ahiḥ ahiḥ budhyasva budhyasva budhyasva iti .kriyāsamabhihāre . kriyāsamabhihāre dve bhavataḥ iti vaktavyam . \\
\noindent \textbf{(P\_8,1.12.2)} saḥ bhavān lunīhi lunīhi iti eva ayam lunāti . \\
\noindent \textbf{(P\_8,1.12.2)} [ābhīkṣṇye .] ābhīkṣṇye dve bhavataḥ iti vaktavyam . \\
\noindent \textbf{(P\_8,1.12.2)} bhuktvā bhuktvā vrajati . \\
\noindent \textbf{(P\_8,1.12.2)} bhojam bhojam vrajati . \\
\noindent \textbf{(P\_8,1.12.2)} ḍāci ca . ḍāci ca dve bhavataḥ iti vaktavyam . \\
\noindent \textbf{(P\_8,1.12.2)} paṭapaṭāyati maṭamaṭāyati . \\
\noindent \textbf{(P\_8,1.12.2)} pūrvaprathamayoḥ arthātiśayavivakṣāyām . pūrvaprathamayoḥ arthātiśayavivakṣāyām dve bhavataḥ iti vaktavyam . \\
\noindent \textbf{(P\_8,1.12.2)} pūrvam pūrvam puṣpyanti . \\
\noindent \textbf{(P\_8,1.12.2)} prathamam prathamam pacyante . \\
\noindent \textbf{(P\_8,1.12.2)} ḍataraḍatamayoḥ samasampradhāraṇāyām strīnigade bhāve . ḍataraḍatamayoḥ samasampradhāraṇāyām strīnigade bhāve dve bhavataḥ iti vaktavyam . \\
\noindent \textbf{(P\_8,1.12.2)} ubhau imau āḍhyau katarā katarā anayoḥ āḍhyatā . \\
\noindent \textbf{(P\_8,1.12.2)} sarve ime āḍhyāḥ katamā katamā eṣām iti . \\
\noindent \textbf{(P\_8,1.12.2)} karmavyatihāre sarvanāmnaḥ samāsavat ca bahulam yadā na samāsavat prathamaikavacanam tadā pūrvapadasya . karmavyatihāre sarvanāmnaḥ dve bhavataḥ iti vaktavyam samāsavat ca bahulam . \\
\noindent \textbf{(P\_8,1.12.2)} yadā na samāsavat prathamaikavacanam bhavati tadā pūrvapadasya . \\
\noindent \textbf{(P\_8,1.12.2)} anyo'nyam ime brāhmaṇāḥ bhojayanti . \\
\noindent \textbf{(P\_8,1.12.2)} anyo'nyasya bhojayanti . \\
\noindent \textbf{(P\_8,1.12.2)} itaretaram bhojayanti . \\
\noindent \textbf{(P\_8,1.12.2)} itaretarasya bhojayanti . \\
\noindent \textbf{(P\_8,1.12.2)} strīnapuṃsakayoḥ uttarapadasya vā ambhāvaḥ . strīnapuṃsakayoḥ uttarapadasya vā ambhāvaḥ vaktavyaḥ . \\
\noindent \textbf{(P\_8,1.12.2)} anyo'nyam ime brāhmaṇyau bhojayataḥ . \\
\noindent \textbf{(P\_8,1.12.2)} anyo'nyām bhojayataḥ . \\
\noindent \textbf{(P\_8,1.12.2)} itaretaram bhojayataḥ . \\
\noindent \textbf{(P\_8,1.12.2)} itaretarām bhojayataḥ . \\
\noindent \textbf{(P\_8,1.12.2)} anyo'nyam ime brāhmaṇakule bhojayataḥ . \\
\noindent \textbf{(P\_8,1.12.2)} anyo'nyām bhojayataḥ . \\
\noindent \textbf{(P\_8,1.12.2)} itaretaram bhojayataḥ . \\
\noindent \textbf{(P\_8,1.12.2)} itaretarām bhojayataḥ \\
\noindent \textbf{(P\_8,1.15)} atyantasahacarite lokavijñāte dvandvam iti upasaṅkhyānam . \\
\noindent \textbf{(P\_8,1.15)} atyantasahacarite lokavijñāte dvandvam iti upasaṅkhyānam kartavyam . \\
\noindent \textbf{(P\_8,1.15)} dvandvam skandaviśākhau . \\
\noindent \textbf{(P\_8,1.15)} dvandvam nāradaparvatau . \\
\noindent \textbf{(P\_8,1.15)} atyantasahacarite iti kimartham . \\
\noindent \textbf{(P\_8,1.15)} dvau yudhiṣṭhirārjunau . \\
\noindent \textbf{(P\_8,1.15)} lokavijñāte iti kimartham . \\
\noindent \textbf{(P\_8,1.15)} dvau devadattayajñadattau \\
\noindent \textbf{(P\_8,1.15)} atha dvandvam iti kim nipātyate . \\
\noindent \textbf{(P\_8,1.15)} dvandvam iti pūrvapadasya ca ambhāvaḥ uttarapadasya ca atvam napuṃsakatvam ca . \\
\noindent \textbf{(P\_8,1.15)} pūrvapadasya ca ambhāvaḥ nipātyate uttarapadasya ca atvam napuṃsakatvam ca . \\
\noindent \textbf{(P\_8,1.15)} uktam vā . kim uktam . \\
\noindent \textbf{(P\_8,1.15)} liṅgam aśiṣyam lokāśrayatvāt liṅgasya iti tatra napuṃsakatvam anipātyam \\
\noindent \textbf{(P\_8,1.16-17)} KA\_III,371.8-372.7 Ro\_V,320-322 {1/21}     ā kutaḥ padādhikāraḥ . \\
\noindent \textbf{(P\_8,1.16-17)} KA\_III,371.8-372.7 Ro\_V,320-322 {2/21}     padādhikāraḥ prāk apadāntādhikārāt . \\
\noindent \textbf{(P\_8,1.16-17)} KA\_III,371.8-372.7 Ro\_V,320-322 {3/21}     apadāntasyamūrdhanyaḥ iti ataḥ prāk padādhikāraḥ . \\
\noindent \textbf{(P\_8,1.16-17)} KA\_III,371.8-372.7 Ro\_V,320-322 {4/21}     atha padāt iti adhikāraḥ ā kutaḥ . \\
\noindent \textbf{(P\_8,1.16-17)} KA\_III,371.8-372.7 Ro\_V,320-322 {5/21}     padāt prāk supi kutsanāt . padāt iti adhikāraḥ prāk supi kutsanāt . \\
\noindent \textbf{(P\_8,1.16-17)} KA\_III,371.8-372.7 Ro\_V,320-322 {6/21}     kutsane ca supy agotrādau iti ataḥ prāk . \\
\noindent \textbf{(P\_8,1.16-17)} KA\_III,371.8-372.7 Ro\_V,320-322 {7/21}     yaṇekādeśasvaraḥ tu ūrdhvam padādhikārāt . yaṇekādeśasvaraḥ tu ūrdhvam padādhikārāt kartavyaḥ . \\
\noindent \textbf{(P\_8,1.16-17)} KA\_III,371.8-372.7 Ro\_V,320-322 {8/21}     iha vacane hi apadāntasya aprāptiḥ . iha hi kriyamāṇe apadāntasya aprāptiḥ syāt . \\
\noindent \textbf{(P\_8,1.16-17)} KA\_III,371.8-372.7 Ro\_V,320-322 {9/21}     udāttasvaritayor yaṇaḥsvaritaḥ anudāttasya iti iha eva syāt kumāryau kiśoryau iha na syāt kumāryaḥ kiśoryaḥ . \\
\noindent \textbf{(P\_8,1.16-17)} KA\_III,371.8-372.7 Ro\_V,320-322 {10/21}     ekādeśe udāttenodāttaḥ iha eva syāt vṛkṣau plakṣau iha na syāt vṛkṣāḥ plakṣāḥ . \\
\noindent \textbf{(P\_8,1.16-17)} KA\_III,371.8-372.7 Ro\_V,320-322 {11/21}     na vā padādhikārasya viśeṣaṇatvāt . na vā ūrdhvam padādhikārāt kartavyaḥ yaṇekādeśasvaraḥ . \\
\noindent \textbf{(P\_8,1.16-17)} KA\_III,371.8-372.7 Ro\_V,320-322 {12/21}     kim kāraṇam . \\
\noindent \textbf{(P\_8,1.16-17)} KA\_III,371.8-372.7 Ro\_V,320-322 {13/21}     padādhikārasya viśeṣaṇatvāt . \\
\noindent \textbf{(P\_8,1.16-17)} KA\_III,371.8-372.7 Ro\_V,320-322 {14/21}     padasya iti na eṣā sthānaṣaṣṭhī . \\
\noindent \textbf{(P\_8,1.16-17)} KA\_III,371.8-372.7 Ro\_V,320-322 {15/21}     kā tarhi . \\
\noindent \textbf{(P\_8,1.16-17)} KA\_III,371.8-372.7 Ro\_V,320-322 {16/21}     viśeṣaṇaṣaṣṭhī . \\
\noindent \textbf{(P\_8,1.16-17)} KA\_III,371.8-372.7 Ro\_V,320-322 {17/21}     kim vaktavyam etat . \\
\noindent \textbf{(P\_8,1.16-17)} KA\_III,371.8-372.7 Ro\_V,320-322 {18/21}     na hi . \\
\noindent \textbf{(P\_8,1.16-17)} KA\_III,371.8-372.7 Ro\_V,320-322 {19/21}     katham anucyamānam gaṃsyate . \\
\noindent \textbf{(P\_8,1.16-17)} KA\_III,371.8-372.7 Ro\_V,320-322 {20/21}     pratyākhyāyate sthanaṣaṣṭhī . \\
\noindent \textbf{(P\_8,1.16-17)} KA\_III,371.8-372.7 Ro\_V,320-322 {21/21}     antagrahaṇāt vā nalope . atha vā yat ayam nalopaḥprātipadikāntasya iti antagrahaṇam karoti tat jñāpayati ācāryaḥ viśeṣaṇaṣaṣṭhī eṣā na sthānaṣaṣṭhī iti \\
\noindent \textbf{(P\_8,1.18.1)} sarvavacanam kimartham . \\
\noindent \textbf{(P\_8,1.18.1)} sarvavacanam anādeḥ anudāttārtham . \\
\noindent \textbf{(P\_8,1.18.1)} sarvagrahaṇam kriyate anādeḥ api anudāttatvam yathā syāt iti . \\
\noindent \textbf{(P\_8,1.18.1)} tiṅatiṅaḥ iha eva syāt devadattaḥ : pacati iti iha na syāt : devadattaḥ karoti iti . \\
\noindent \textbf{(P\_8,1.18.1)} sarvavacanam anādeḥ anudāttārtham iti cet luṭi pratiṣedhāt siddham . \\
\noindent \textbf{(P\_8,1.18.1)} sarvavacanam anādeḥ anudāttārtham iti cet tat na . \\
\noindent \textbf{(P\_8,1.18.1)} kim kāraṇam . \\
\noindent \textbf{(P\_8,1.18.1)} luṭi pratiṣedhāt siddham . \\
\noindent \textbf{(P\_8,1.18.1)} yat ayam luṭi pratiṣedham śāsti naluṭ iti tat jñāpayati ācāryaḥ anādeḥ api anudāttatvam bhavati iti . \\
\noindent \textbf{(P\_8,1.18.1)} katham kṛtvā jñāpakam . \\
\noindent \textbf{(P\_8,1.18.1)} na hi luḍantam ādyudāttam asti . \\
\noindent \textbf{(P\_8,1.18.1)} alo'ntyavidhiprasaṅgaḥ tu . alaḥ antyasya vidhayaḥ bhavanti iti antyasya vidhiḥ prāpnoti . \\
\noindent \textbf{(P\_8,1.18.1)} yatra hi ādividhiḥ na asti alo'ntyavidhinā tatra bhavitavyam . \\
\noindent \textbf{(P\_8,1.18.1)} tatra kaḥ doṣaḥ . \\
\noindent \textbf{(P\_8,1.18.1)} tiṅatiṅaḥ iti iha eva syāt devadattayajñadattau kurutaḥ iha na syāt devadattaḥ karoti iti . \\
\noindent \textbf{(P\_8,1.18.1)} lṛṭi pratiṣedhāt siddham . yat ayam lṛṭi pratiṣedham śāsti tat jñāpayati ācāryaḥ anantyasya api anudāttatvam bhavati iti . \\
\noindent \textbf{(P\_8,1.18.1)} katham kṛtvā jñāpakam . \\
\noindent \textbf{(P\_8,1.18.1)} na hi lṛḍantam antodāttam asti . \\
\noindent \textbf{(P\_8,1.18.1)} nanu ca idam asti bhokṣye iti . \\
\noindent \textbf{(P\_8,1.18.1)} uktam vā . kim uktam . \\
\noindent \textbf{(P\_8,1.18.1)} na vā padādhikārasya viśeṣaṇatvāt iti . \\
\noindent \textbf{(P\_8,1.18.1)} idam tarhi prayojanam yuṣmadasmadoḥṣaṣṭhīcaturthīdvitīyāsthayorvāmnāvau iti vāmnau ādayaḥ savibhaktikasya yathā syuḥ iti . \\
\noindent \textbf{(P\_8,1.18.1)} etat api na asti prayojanam . \\
\noindent \textbf{(P\_8,1.18.1)} padasya iti hi vartate vibhaktyantam ca padam tatra antareṇa sarvagrahaṇam savibhaktikasya bhaviṣyati . \\
\noindent \textbf{(P\_8,1.18.1)} bhavet siddham yatra vibhaktyantam padam yatra tu khalu vibhaktau padam tatra na sidhyati . \\
\noindent \textbf{(P\_8,1.18.1)} grāmaḥ vām dīyate . \\
\noindent \textbf{(P\_8,1.18.1)} grāmaḥ nau dīyate . \\
\noindent \textbf{(P\_8,1.18.1)} janapadaḥ vām dīyate . \\
\noindent \textbf{(P\_8,1.18.1)} janapadaḥ nau dīyate . \\
\noindent \textbf{(P\_8,1.18.1)} nanu ca sthagrahaṇam kriyate tena savibhaktikasya eva bhaviṣyati . \\
\noindent \textbf{(P\_8,1.18.1)} astī anyat sthagrahaṇasya prayaojanam . \\
\noindent \textbf{(P\_8,1.18.1)} kim . \\
\noindent \textbf{(P\_8,1.18.1)} śrūyamāṇavibhakiviśeṣaṇam yathā vijñāyeta . \\
\noindent \textbf{(P\_8,1.18.1)} yatra vibhaktiḥ śrūyate tatra yathā syāt iha mā bhūt iti yuṣmatputraḥ dadāti iti asmatputraḥ dadāti iti \\
\noindent \textbf{(P\_8,1.18.2)} samānavākye nighātayuṣmadasmadādeśāḥ . \\
\noindent \textbf{(P\_8,1.18.2)} samānavākye iti prakṛtya nighātayuṣmadasmadādeśāḥ vaktavyāḥ . \\
\noindent \textbf{(P\_8,1.18.2)} kim prayojanam . \\
\noindent \textbf{(P\_8,1.18.2)} nānāvākye mā bhūvan iti . \\
\noindent \textbf{(P\_8,1.18.2)} ayam daṇḍaḥ hara anena . \\
\noindent \textbf{(P\_8,1.18.2)} odanam paca tava bhaviṣyati mama bhaviṣyati . \\
\noindent \textbf{(P\_8,1.18.2)} paśyārthaiḥ ca pratiṣedhaḥ . paśyārthaiḥ ca pratiṣedhaḥ samānavākye iti prakṛtya vaktavyaḥ . \\
\noindent \textbf{(P\_8,1.18.2)} itarathā hi yatra eva paśyārthānām yuṣmadasmadī sādhanam tatra pratiṣedhaḥ syāt . \\
\noindent \textbf{(P\_8,1.18.2)} grāmaḥ tvām samprekṣya sandṛśya samīkṣya gataḥ . \\
\noindent \textbf{(P\_8,1.18.2)} grāmaḥ mām samprekṣya sandṛśya samīkṣya gataḥ. iha na syāt . \\
\noindent \textbf{(P\_8,1.18.2)} grāmaḥ tava svam samprekṣya sandṛśya samīkṣya gataḥ . \\
\noindent \textbf{(P\_8,1.18.2)} grāmaḥ mama svam samprekṣya sandṛśya samīkṣya gataḥ \\
\noindent \textbf{(P\_8,1.26)} yuṣmadasmadoḥ anyatarasyām ananvādeśe . \\
\noindent \textbf{(P\_8,1.26)} yuṣmadasmadoḥ anyatarasyām ananvādeśe iti vaktavyam . \\
\noindent \textbf{(P\_8,1.26)} grāme kambalaḥ te svam . \\
\noindent \textbf{(P\_8,1.26)} grāme kambalaḥ tava svam . \\
\noindent \textbf{(P\_8,1.26)} grāme kambalaḥ me svam . \\
\noindent \textbf{(P\_8,1.26)} grāme kambalaḥ mama svam . \\
\noindent \textbf{(P\_8,1.26)} ananvādeśe iti kimartham . \\
\noindent \textbf{(P\_8,1.26)} atho grāme kambalaḥ te svam . \\
\noindent \textbf{(P\_8,1.26)} atho grāme kambalaḥ me svam . \\
\noindent \textbf{(P\_8,1.26)} aparaḥ āha : sarve eva vāmnāvādayaḥ ananvādeśe vibhāṣā vaktavyāḥ . \\
\noindent \textbf{(P\_8,1.26)} kambalaḥ te svam . \\
\noindent \textbf{(P\_8,1.26)} kambalaḥ tava svam . \\
\noindent \textbf{(P\_8,1.26)} kambalaḥ me svam . \\
\noindent \textbf{(P\_8,1.26)} kambalaḥ mama svam . \\
\noindent \textbf{(P\_8,1.26)} ananvādeśe iti kimartham . \\
\noindent \textbf{(P\_8,1.26)} atho kambalaḥ te svam . \\
\noindent \textbf{(P\_8,1.26)} atho kambalaḥ me svam . \\
\noindent \textbf{(P\_8,1.26)} na tarhi idānīm idam vaktavyam sapūrvāyāḥ prathamāyāḥ vibhāṣā iti . \\
\noindent \textbf{(P\_8,1.26)} vaktavyam ca . \\
\noindent \textbf{(P\_8,1.26)} kim prayojanam . \\
\noindent \textbf{(P\_8,1.26)} anvādeśārtham . \\
\noindent \textbf{(P\_8,1.26)} anvādeśe vibhāṣā yathā syāt . \\
\noindent \textbf{(P\_8,1.26)} atho grāme kambalaḥ te svam . \\
\noindent \textbf{(P\_8,1.26)} atho grāme kambalaḥ tava svam . \\
\noindent \textbf{(P\_8,1.26)} atho grāme kambalaḥ me svam . \\
\noindent \textbf{(P\_8,1.26)} atho grāme kambalaḥ mama svam \\
\noindent \textbf{(P\_8,1.27)} kim idam tiṅaḥ gotrādiṣu kutsanābhīkṣṇyagrahaṇam pāṭhaviśeṣaṇam . \\
\noindent \textbf{(P\_8,1.27)} kutsanābhīkṣṇyayoḥ arthayoḥ gotrādīni bhavanti tiṅaḥ parāṇi anudāttāni iti . \\
\noindent \textbf{(P\_8,1.27)} āhosvit anudāttaviśeṣaṇam . \\
\noindent \textbf{(P\_8,1.27)} tiṅaḥ parāṇi gotrādīni kutsanābhīkṣṇyayoḥ arthayoḥ anudāttāni bhavanti iti . \\
\noindent \textbf{(P\_8,1.27)} tiṅaḥ gotrādiṣu kutsanābhīkṣṇyagrahaṇam pāṭhaviśeṣaṇam . \\
\noindent \textbf{(P\_8,1.27)} tiṅaḥ gotrādiṣu kutsanābhīkṣṇyagrahaṇam kriyate pāṭhaviśeṣaṇam . \\
\noindent \textbf{(P\_8,1.27)} pāṭhaḥ viśeṣyate . \\
\noindent \textbf{(P\_8,1.27)} anudāttaviśeṣaṇe hi anyatra gotrādigrahaṇe kutsanābhīkṣṇyagrahaṇam . anudāttaviśeṣaṇe hi sati anyatra gotrādigrahaṇe kutsanābhīkṣṇyagrahaṇam kartavyam syāt . \\
\noindent \textbf{(P\_8,1.27)} canacidivagotrāditaddhitāmreḍiteṣvagateḥ iti kutsanābhīkṣṇyayoḥ iti vaktavyam syāt . \\
\noindent \textbf{(P\_8,1.27)} anudāttagrahaṇam vā . atha vā yāni anudāttāni iti vaktavyam syāt . \\
\noindent \textbf{(P\_8,1.27)} tasmāt suṣṭhu ucyate tiṅaḥ gotrādiṣu kutsanābhīkṣṇyagrahaṇam pāṭhaviśeṣaṇam anudāttaviśeṣaṇe hi anyatra gotrādigrahaṇe kutsanābhīkṣṇyagrahaṇam anudāttagrahaṇam vā iti \\
\noindent \textbf{(P\_8,1.28)} atiṅaḥ iti kimartham . \\
\noindent \textbf{(P\_8,1.28)} pacati karoti . \\
\noindent \textbf{(P\_8,1.28)} atiṅvacanam anarthakam samānavākyādhikārāt . \\
\noindent \textbf{(P\_8,1.28)} atiṅvacanam anarthakam . \\
\noindent \textbf{(P\_8,1.28)} kim kāraṇam . \\
\noindent \textbf{(P\_8,1.28)} samānavākyādhikārāt . \\
\noindent \textbf{(P\_8,1.28)} samānavākye iti vartate na ca samānavākye dve tiṅante staḥ \\
\noindent \textbf{(P\_8,1.30.1)} nipātaiḥ iti kimartham . \\
\noindent \textbf{(P\_8,1.30.1)} yat kūjati śakaṭam . \\
\noindent \textbf{(P\_8,1.30.1)} yatī kūjati śakaṭī . \\
\noindent \textbf{(P\_8,1.30.1)} yan rathaḥ kūjati . \\
\noindent \textbf{(P\_8,1.30.1)} nipātaiḥ iti śakyam avaktum . \\
\noindent \textbf{(P\_8,1.30.1)} kasmāt na bhavati . \\
\noindent \textbf{(P\_8,1.30.1)} yat kūjati śakaṭam . \\
\noindent \textbf{(P\_8,1.30.1)} yatī kūjati śakaṭī . \\
\noindent \textbf{(P\_8,1.30.1)} yan rathaḥ kūjati . \\
\noindent \textbf{(P\_8,1.30.1)} lakṣaṇapratipadoktayoḥ pratipadoktasya eva iti . \\
\noindent \textbf{(P\_8,1.30.1)} na eṣā paribhāṣā iha śakyā vijñātum . \\
\noindent \textbf{(P\_8,1.30.1)} iha hi doṣaḥ syāt . \\
\noindent \textbf{(P\_8,1.30.1)} yāvadyathābhyām iha na syāt yāvat asti atra eṣaḥ saraḥ janebhyaḥ kṛṇavat \\
\noindent \textbf{(P\_8,1.30.2)} caṇ ṇidviśiṣṭaḥ cedarthe . \\
\noindent \textbf{(P\_8,1.30.2)} caṇ ṇidviśiṣṭaḥ cedarthe draṣṭavyaḥ . \\
\noindent \textbf{(P\_8,1.30.2)} ayam ca vai mariṣyati . \\
\noindent \textbf{(P\_8,1.30.2)} ayam cet mariṣyati . \\
\noindent \textbf{(P\_8,1.30.2)} na ca pitṛbhyaḥ pūrvebhyaḥ dāsyati . \\
\noindent \textbf{(P\_8,1.30.2)} aprāyaścittikṛtau ca syātām \\
\noindent \textbf{(P\_8,1.35)} anekam iti kim udāharaṇam . \\
\noindent \textbf{(P\_8,1.35)} yadā hi asau mattaḥ bhavati atha yat tapati . \\
\noindent \textbf{(P\_8,1.35)} na etat asti . \\
\noindent \textbf{(P\_8,1.35)} ekam atra hiyuktam aparam yadyuktam tataḥ ubhayoḥ api anighātaḥ . \\
\noindent \textbf{(P\_8,1.35)} idam tarhi . \\
\noindent \textbf{(P\_8,1.35)} anṛtam hi mattaḥ vadati pāpmā enam vipunāti . \\
\noindent \textbf{(P\_8,1.35)} ekam khalu api . \\
\noindent \textbf{(P\_8,1.35)} agniḥ hi pūrvam udajayat tam indraḥ anūdajayat iti \\
\noindent \textbf{(P\_8,1.39)} pūjāyām iti vartamāne punaḥ pūjāgrahaṇam kimartham . \\
\noindent \textbf{(P\_8,1.39)} anighātapratiṣedhābhisambaddham tat . \\
\noindent \textbf{(P\_8,1.39)} yadi tat anuvarteta iha api anighātapratiṣedhaḥ prasajyeta . \\
\noindent \textbf{(P\_8,1.39)} iṣyate ca atra nighātapratiṣedhaḥ . \\
\noindent \textbf{(P\_8,1.39)} yathā punaḥ tatra yāvat yathā iti etābhyām anighāte prāpte anighātapratiṣedhaḥ ucyate iha idānīm kena anighāte prāpte anighātapratiṣedhaḥ ucyeta . \\
\noindent \textbf{(P\_8,1.39)} iha api yadvṛttānnityam iti evamādibhiḥ \\
\noindent \textbf{(P\_8,1.46)} kimartham idam ucyate na gatyarthaloṭā lṛṭ iti eva siddham . \\
\noindent \textbf{(P\_8,1.46)} niyamārthaḥ ayam ārambhaḥ . \\
\noindent \textbf{(P\_8,1.46)} ehi manye prahāse eva yathā syāt . \\
\noindent \textbf{(P\_8,1.46)} kva mā bhūt . \\
\noindent \textbf{(P\_8,1.46)} ehi manye rathena yāsyasi iti \\
\noindent \textbf{(P\_8,1.47)} kim idam apūrvagrahaṇam jātuviśeṣaṇam . \\
\noindent \textbf{(P\_8,1.47)} jātuśabdāt apūrvāt tiṅantam iti . \\
\noindent \textbf{(P\_8,1.47)} āhosvit tiṅantaviśeṣaṇam . \\
\noindent \textbf{(P\_8,1.47)} jātuśabdāt tiṅantam apūrvam iti . \\
\noindent \textbf{(P\_8,1.47)} jātuviśeṣaṇam iti āha . \\
\noindent \textbf{(P\_8,1.47)} katham jñāyate . \\
\noindent \textbf{(P\_8,1.47)} yat ayam kiṃvṛttañcaciduttaram iti āha . \\
\noindent \textbf{(P\_8,1.47)} katham kṛtvā jñāpakam . \\
\noindent \textbf{(P\_8,1.47)} atra api apūrvam iti etat anuvartate na ca asti sambhavaḥ yat kiṃvṛttam ca ciduttaram syāt tiṅantam ca apūrvam . \\
\noindent \textbf{(P\_8,1.47)} atra api tiṅantaviśeṣaṇam eva . \\
\noindent \textbf{(P\_8,1.47)} katham . \\
\noindent \textbf{(P\_8,1.47)} kiṃvṛttāt ciduttarāt tiṅantam apūrvam iti . \\
\noindent \textbf{(P\_8,1.47)} yat tarhi āhoutāhocānantaram iti anantaragrahaṇam karoti . \\
\noindent \textbf{(P\_8,1.47)} etasya api asti vacane prayojanam . \\
\noindent \textbf{(P\_8,1.47)} kim . \\
\noindent \textbf{(P\_8,1.47)} śeṣaprakḷptyartham etat syāt . \\
\noindent \textbf{(P\_8,1.47)} śeṣevibhāṣā kaḥ ca śeṣaḥ . \\
\noindent \textbf{(P\_8,1.47)} sāntaram śeṣaḥ iti . \\
\noindent \textbf{(P\_8,1.47)} antareṇa api anantaragrahaṇam prakḷptaḥ śeṣaḥ . \\
\noindent \textbf{(P\_8,1.47)} katham. apūrvaḥ iti vartate . \\
\noindent \textbf{(P\_8,1.47)} śeṣe vibhāṣā . \\
\noindent \textbf{(P\_8,1.47)} kaḥ ca śeṣaḥ . \\
\noindent \textbf{(P\_8,1.47)} sapūrvaḥ śeṣaḥ iti \\
\noindent \textbf{(P\_8,1.51)} lṛṭaḥ prakṛtibhāve kartuḥ anyatve upasaṅkhyānam kārakānyatvāt . \\
\noindent \textbf{(P\_8,1.51)} lṛṭaḥ prakṛtibhāve kartuḥ yat kārakam anyat tasya anyatve upasaṅkhyānam kartavyam . \\
\noindent \textbf{(P\_8,1.51)} āgaccha devadatta grāmam odanam bhokṣyase . \\
\noindent \textbf{(P\_8,1.51)} kim punaḥ kāraṇam na sidhyati . \\
\noindent \textbf{(P\_8,1.51)} kārakānyatvāt . \\
\noindent \textbf{(P\_8,1.51)} na cet kārakam sarvānyat iti ucyate sarvānyat ca atra kārakam . \\
\noindent \textbf{(P\_8,1.51)} kim punaḥ kāraṇam sarvānyatpratiṣedhena āśrīyate na punaḥ asarvānyadvidhānena āśrīyeta . \\
\noindent \textbf{(P\_8,1.51)} kartā ca atra asarvānyaḥ tataḥ kartṛsāmānyāt siddham . \\
\noindent \textbf{(P\_8,1.51)} kartṛsāmānyāt siddham iti cet tadbhede anyasāmānye prakṛtibhāvaprasaṅgaḥ . kartṛsāmānyāt siddham iti cet tadbhede kartṛbhede anyasmin kārakasāmānye prakṛtibhāvaḥ prāpnoti . \\
\noindent \textbf{(P\_8,1.51)} āhara devadatta śālīn yajñadatta enān bhokṣyate . \\
\noindent \textbf{(P\_8,1.51)} evam tarhi vyaktam eva paṭhitavyam na cet kartā sarvānyaḥ iti . \\
\noindent \textbf{(P\_8,1.51)} na cet kartā sarvānyaḥ iti cet anyābhidhāne pratiṣedham eke . na cet kartā sarvānyaḥ iti cet anyābhidhāne pratiṣedham eke icchanti . \\
\noindent \textbf{(P\_8,1.51)} uhyantām devadattena śālayaḥ yajñadattena bhokṣyante iti prāpnoti bhokṣyante iti ca iṣyate . \\
\noindent \textbf{(P\_8,1.51)} siddham tu tiṅoḥ ekadravyābhidhānāt . siddham etat . \\
\noindent \textbf{(P\_8,1.51)} katham . \\
\noindent \textbf{(P\_8,1.51)} tiṅoḥ ekadravyābhidhānāt . \\
\noindent \textbf{(P\_8,1.51)} yatra tiṅbhyām ekam dravyam abhidhīyate tatra iti vaktavyam \\
\noindent \textbf{(P\_8,1.55)} kasya ayam pratiṣedhaḥ . \\
\noindent \textbf{(P\_8,1.55)} āmaḥ ekāntare aikaśrutyapratiṣedhaḥ . \\
\noindent \textbf{(P\_8,1.55)} āmaḥ ekāntare aikaśrutyasya ayam pratiṣedhaḥ . \\
\noindent \textbf{(P\_8,1.55)} katham punaḥ aprakṛtasya asaṃśabditasya aikaśrutyasya pratiṣedhaḥ śakyaḥ vijñātum . \\
\noindent \textbf{(P\_8,1.55)} anantike iti ucyate . \\
\noindent \textbf{(P\_8,1.55)} anantikam ca kim . \\
\noindent \textbf{(P\_8,1.55)} dūram . \\
\noindent \textbf{(P\_8,1.55)} dūrāt sambuddhau ekaśrutiḥ ucyate . \\
\noindent \textbf{(P\_8,1.55)} asti prayojanam etat . \\
\noindent \textbf{(P\_8,1.55)} kim tarhi iti . \\
\noindent \textbf{(P\_8,1.55)} nighātaprasaṅgaḥ tu . nighātaḥ tu prāpnoti . \\
\noindent \textbf{(P\_8,1.55)} ām bhoḥ devadatta3 . \\
\noindent \textbf{(P\_8,1.55)} āmantritasya anudāttatvam prāpnoti . \\
\noindent \textbf{(P\_8,1.55)} siddham tu pratiṣedhādhikāre pratiṣedhavacanāt . siddham etat . \\
\noindent \textbf{(P\_8,1.55)} katham . \\
\noindent \textbf{(P\_8,1.55)} pratiṣedhādhikāre pratiṣedhavacanasāmarthyāt nighātaḥ na bhaviṣyati . \\
\noindent \textbf{(P\_8,1.55)} na eva vā punaḥ atra aikaśrutyam prāpnoti . \\
\noindent \textbf{(P\_8,1.55)} kim kāraṇam . \\
\noindent \textbf{(P\_8,1.55)} anantike iti ucyate anyat ca dūram anyat anantikam . \\
\noindent \textbf{(P\_8,1.55)} yadi evam plutaḥ api tarhi na prāpnoti plutaḥ api hi dūrāt iti ucyate . \\
\noindent \textbf{(P\_8,1.55)} iṣṭam eva etat saṅgṛhītam . \\
\noindent \textbf{(P\_8,1.55)} ām bhoḥ devadatta iti eva bhavitavyam \\
\noindent \textbf{(P\_8,1.56)} kimartham idam ucyate . \\
\noindent \textbf{(P\_8,1.56)} yadādaiḥ eva sarvaiḥ etaiḥ anighātakāraṇaiḥ yoge anighātaḥ ucyate . \\
\noindent \textbf{(P\_8,1.56)} yathā eva pūrvaiḥ yoge evam paraiḥ api . \\
\noindent \textbf{(P\_8,1.56)} ataḥ uttaram paṭhati . \\
\noindent \textbf{(P\_8,1.56)} yaddhituparasya chandasi anighātaḥ anyaparapratiṣedhārthaḥ . \\
\noindent \textbf{(P\_8,1.56)} yaddhituparasya chandasi anighātaḥ ucyate anyaparapratiṣedhārthaḥ . \\
\noindent \textbf{(P\_8,1.56)} anyaparasya pratiṣedhaḥ mā bhūt iti . \\
\noindent \textbf{(P\_8,1.56)} jāye svaḥ rohāva ehi . \\
\noindent \textbf{(P\_8,1.56)} atha idānīm rohāva iti anena yukte ehi iti asya kasmāt na bhavati . \\
\noindent \textbf{(P\_8,1.56)} loṭ ca gatyarthaloṭā yuktaḥ iti prāpnoti . \\
\noindent \textbf{(P\_8,1.56)} na ruhiḥ gatyarthaḥ . \\
\noindent \textbf{(P\_8,1.56)} katham jñāyate . \\
\noindent \textbf{(P\_8,1.56)} yat ayam gatyarthākarmakaśliṣaśīṅsthāsavasajanaruhajīryatibhyaśca iti pṛthak ruhigrahaṇam karoti . \\
\noindent \textbf{(P\_8,1.56)} yadi na ruhiḥ gatyarthaḥ ārohanti hastinam manuṣyāḥ ārohayati hastī sthalam manuṣyān gatibuddhipratyavasānārthaśabdakarmākarmakāṇāmaṇikartāsaṇau iti karmasañjñā na prāpnoti . \\
\noindent \textbf{(P\_8,1.56)} tasmāt na etat śakyam vaktum na ruhiḥ gatyarthaḥ iti . \\
\noindent \textbf{(P\_8,1.56)} kasmāt tarhi rohāva iti anena yukte ehi iti asya na bhavati . \\
\noindent \textbf{(P\_8,1.56)} chāndasatvāt \\
\noindent \textbf{(P\_8,1.57)} āmreḍiteṣu agateḥ sagatiḥ api tiṅ iti atra gatigrahaṇe upasargagrahaṅam . \\
\noindent \textbf{(P\_8,1.57)} āmreḍiteṣu agateḥ sagatirapitiṅ iti atra gatigrahaṇe upasargagrahaṇam draṣṭavyam . \\
\noindent \textbf{(P\_8,1.57)} iha mā bhūt . \\
\noindent \textbf{(P\_8,1.57)} śuklīkaroti cana . \\
\noindent \textbf{(P\_8,1.57)} krṣṇīkaroti cana . \\
\noindent \textbf{(P\_8,1.57)} yatkāṣṭhā śuklīkaroti . \\
\noindent \textbf{(P\_8,1.57)} yatkāṣṭhā kṛṣṇīkaroti . \\
\noindent \textbf{(P\_8,1.57)} apara āha . \\
\noindent \textbf{(P\_8,1.57)} sarvatra eva āṣṭamike gatigrahaṇe upasargagrahaṇam draṣtavyam gatirgatautiṅicodāttavativarjam iti \\
\noindent \textbf{(P\_8,1.66)} yadvṛttāt iti ucyate tatra idam na sidhyati yaḥ pacati yam pacati iti . \\
\noindent \textbf{(P\_8,1.66)} vṛttagrahaṇena tadvibhaktyantam pratīyāt . \\
\noindent \textbf{(P\_8,1.66)} katham yataraḥ pacati yatamaḥ pacati iti . \\
\noindent \textbf{(P\_8,1.66)} ḍataraḍatamau ca pratīyāt . \\
\noindent \textbf{(P\_8,1.66)} katham yadā dadāti iti . \\
\noindent \textbf{(P\_8,1.66)} eṣaḥ api vibhaktisañjñaḥ . \\
\noindent \textbf{(P\_8,1.66)} katham yāvat asti atra eṣaḥ saraḥ janebhyaḥ kṛṇavat . \\
\noindent \textbf{(P\_8,1.66)} yāvadyathābhyām iti evam bhaviṣyati . \\
\noindent \textbf{(P\_8,1.66)} katham yadryaṅ vāyuḥ pavate yatkāmāḥ te juhumaḥ . \\
\noindent \textbf{(P\_8,1.66)} evam tarhi yat asmin vartate yadvṛttam yadvṛttāt iti evam bhaviṣyati . \\
\noindent \textbf{(P\_8,1.66)} vā yāthākāmye . \\
\noindent \textbf{(P\_8,1.66)} vā yāthākāmye iti vaktavyam . \\
\noindent \textbf{(P\_8,1.66)} yatra kva cana yajate devayajane eva yajate \\
\noindent \textbf{(P\_8,1.67)} pūjitasya anudāttatve kāṣṭhādigrahaṇam . \\
\noindent \textbf{(P\_8,1.67)} pūjitasya anudāttatve kāṣṭhādigrahaṇam kartavyam . \\
\noindent \textbf{(P\_8,1.67)} kāṣṭhādibhyaḥ pūjanāt iti vaktavyam . \\
\noindent \textbf{(P\_8,1.67)} iha mā bhūt . \\
\noindent \textbf{(P\_8,1.67)} śobhanaḥ adhyāpakaḥ . \\
\noindent \textbf{(P\_8,1.67)} malopavacanam ca . malopaḥ ca vaktavyaḥ . \\
\noindent \textbf{(P\_8,1.67)} dāruṇādhyāpakaḥ dāruṇābhirūpaḥ \\
\noindent \textbf{(P\_8,1.68.1)} sagatigrahaṇam kimartham . \\
\noindent \textbf{(P\_8,1.68.1)} sagatigrahaṇam apadatvāt . \\
\noindent \textbf{(P\_8,1.68.1)} sagatigrahaṇam kriyate apadatvāt . \\
\noindent \textbf{(P\_8,1.68.1)} padasya iti vartate na hi sagatikam padam bhavati . \\
\noindent \textbf{(P\_8,1.68.1)} uttarārtham ca . uttarārtham ca sagatigrahaṇam kriyate . \\
\noindent \textbf{(P\_8,1.68.1)} kutsane ca supi agotrādau sagatiḥ api . \\
\noindent \textbf{(P\_8,1.68.1)} prapacati pūti . \\
\noindent \textbf{(P\_8,1.68.1)} atha apigrahaṇam kimartham . \\
\noindent \textbf{(P\_8,1.68.1)} agatikasya api yathā syāt . \\
\noindent \textbf{(P\_8,1.68.1)} yat kāṣṭhā pacati . \\
\noindent \textbf{(P\_8,1.68.1)} na etat asti prayojanam . \\
\noindent \textbf{(P\_8,1.68.1)} siddham pūrveṇa agatikasya . \\
\noindent \textbf{(P\_8,1.68.1)} na sidhyati . \\
\noindent \textbf{(P\_8,1.68.1)} malopābhisambaddham tat . \\
\noindent \textbf{(P\_8,1.68.1)} yadi tat anuvarteta iha api malopaḥ prasajyeta . \\
\noindent \textbf{(P\_8,1.68.1)} dāruṇam pacati iti . \\
\noindent \textbf{(P\_8,1.68.1)} uttarārtham ca apigrahaṇam kriyate . \\
\noindent \textbf{(P\_8,1.68.1)} kutsane ca supi agotrādau agatiḥ api iti . \\
\noindent \textbf{(P\_8,1.68.1)} pacati pūti iti \\
\noindent \textbf{(P\_8,1.68.2)} tiṅnighātāt pūjanāt pūjitam anudāttam vipratiṣedhena . \\
\noindent \textbf{(P\_8,1.68.2)} tiṅnighātāt pūjanāt pūjitam anudāttam iti etat bhavati vipratiṣedhena . \\
\noindent \textbf{(P\_8,1.68.2)} tiṅnighātasya avakāśaḥ . \\
\noindent \textbf{(P\_8,1.68.2)} devadattaḥ pacati . \\
\noindent \textbf{(P\_8,1.68.2)} pūjanāt pūjitam anudāttam iti asya avakāśaḥ . \\
\noindent \textbf{(P\_8,1.68.2)} kāṣṭādhyāpakaḥ . \\
\noindent \textbf{(P\_8,1.68.2)} iha ubhayam prāpnoti . \\
\noindent \textbf{(P\_8,1.68.2)} kāṣṭhā pacati . \\
\noindent \textbf{(P\_8,1.68.2)} pūjanāt pūjitam iti etat bhavati vipratiṣedhena . \\
\noindent \textbf{(P\_8,1.68.2)} kaḥ punaḥ atra viśeṣaḥ tena vā sati anena vā . \\
\noindent \textbf{(P\_8,1.68.2)} ayam asti viśeṣaḥ . \\
\noindent \textbf{(P\_8,1.68.2)} sāpavādakaḥ saḥ vidhiḥ ayam punaḥ nirapavādakaḥ . \\
\noindent \textbf{(P\_8,1.68.2)} yadi hi tena syāt iha na syāt . \\
\noindent \textbf{(P\_8,1.68.2)} yat kāṣṭhā pacati \\
\noindent \textbf{(P\_8,1.69)} supi kutsane kriyāyāḥ makāralopaḥ atiṅi iti ca uktārtham . \\
\noindent \textbf{(P\_8,1.69)} kriyāyāḥ kutsane iti vaktavyam . \\
\noindent \textbf{(P\_8,1.69)} kartuḥ kutsane mā bhūt . \\
\noindent \textbf{(P\_8,1.69)} pacati putiḥ . \\
\noindent \textbf{(P\_8,1.69)} pūtiḥ ca cānubandhaḥ . \\
\noindent \textbf{(P\_8,1.69)} pūtiḥ ca cānubandhaḥ draṣṭavyaḥ . \\
\noindent \textbf{(P\_8,1.69)} pacati pūti . \\
\noindent \textbf{(P\_8,1.69)} vibhāṣitam ca api bahvartham . vibhāṣitam ca api bahvartham draṣṭavyam . \\
\noindent \textbf{(P\_8,1.69)} pacanti puti . \\
\noindent \textbf{(P\_8,1.69)} pacanti puti . \\
\noindent \textbf{(P\_8,1.69)} supi kutsane kriyāyāḥ makāralopaḥ atiṅi iti ca uktārtham . \\
\noindent \textbf{(P\_8,1.69)} pūtiḥ ca cānubandhaḥ . \\
\noindent \textbf{(P\_8,1.69)} vibhāṣitam ca api bahvartham . \\
\noindent \textbf{(P\_8,1.70)} gatau iti kimartham . \\
\noindent \textbf{(P\_8,1.70)} prapacati prakaroti . \\
\noindent \textbf{(P\_8,1.70)} gateḥ anudāttatve gatigrahaṇānarthakyam tiṅi avadhāraṇāt . \\
\noindent \textbf{(P\_8,1.70)} gateḥ anudāttatve gatigrahaṇam anarthakam . \\
\noindent \textbf{(P\_8,1.70)} kim kāraṇam . \\
\noindent \textbf{(P\_8,1.70)} tiṅi avadhāraṇāt . \\
\noindent \textbf{(P\_8,1.70)} tiṅica udāttavati iti etat niyamārtham bhaviṣyati . \\
\noindent \textbf{(P\_8,1.70)} tiṅi udāttavati eva gatiḥ anudāttaḥ bhavati na anyatra iti . \\
\noindent \textbf{(P\_8,1.70)} chandortham tarhi gatigrahaṇam kartavyam . \\
\noindent \textbf{(P\_8,1.70)} chandasi gatau parataḥ anudāttatvam yathā syāt mandraśabde mā bhūt . \\
\noindent \textbf{(P\_8,1.70)} a mandraiḥ indra haribhiḥ yāhi mayuraromabhiḥ . \\
\noindent \textbf{(P\_8,1.70)} chandortham iti cet na agatitvāt . chandortham iti cet tat na . \\
\noindent \textbf{(P\_8,1.70)} kim kāraṇam . \\
\noindent \textbf{(P\_8,1.70)} agatitvāt . \\
\noindent \textbf{(P\_8,1.70)} yatkriyāyuktāḥ tam prati gatyupasargasañjñe bhavataḥ na ca atra āṅaḥ mandraśabdam prati kriyāyogaḥ . \\
\noindent \textbf{(P\_8,1.70)} kim tarhi yāhiśabdam prati . \\
\noindent \textbf{(P\_8,1.70)} iha api tarhi na prāpnoti . \\
\noindent \textbf{(P\_8,1.70)} abhyuddharati upasamādadhāti iti . \\
\noindent \textbf{(P\_8,1.70)} atra api na abheḥ udam prati kriyāyogaḥ . \\
\noindent \textbf{(P\_8,1.70)} kim tarhi haratim prati kriyāyogaḥ . \\
\noindent \textbf{(P\_8,1.70)} na eṣaḥ doṣaḥ . \\
\noindent \textbf{(P\_8,1.70)} udam prati kriyāyogaḥ . \\
\noindent \textbf{(P\_8,1.70)} katham . \\
\noindent \textbf{(P\_8,1.70)} uddharatikriyām viśinaṣṭi . \\
\noindent \textbf{(P\_8,1.70)} udā viśiṣṭām abhiḥ viśinaṣṭi . \\
\noindent \textbf{(P\_8,1.70)} tatra yatkriyāyuktāḥ iti bhavati eva saṅghātam prati kriyāyogaḥ . \\
\noindent \textbf{(P\_8,1.70)} iha api tarhi mandrasādhanā kriyā āṅā vyajyate . \\
\noindent \textbf{(P\_8,1.70)} ā yāhi mandraiḥ iti . \\
\noindent \textbf{(P\_8,1.70)} nanu pūrvam dhātuḥ upasargeṇa yujyate paścāt sādhanena iti . \\
\noindent \textbf{(P\_8,1.70)} na etat sāram . \\
\noindent \textbf{(P\_8,1.70)} pūrvam dhātuḥ sādhanena yujyate paścāt upasargeṇa . \\
\noindent \textbf{(P\_8,1.70)} kim kāraṇam . \\
\noindent \textbf{(P\_8,1.70)} sādhanam hi kriyām nirvartayati tām upasargaḥ viśinaṣṭi abhinirvṛttasya ca arthasya upasargeṇa viśeṣaḥ śakyaḥ vaktum . \\
\noindent \textbf{(P\_8,1.70)} satyam evam etat . \\
\noindent \textbf{(P\_8,1.70)} yaḥ tu asau dhātūpasargayoḥ abhisambandhaḥ tam abhyantare kṛtvā dhātuḥ sādhanena yujyate . \\
\noindent \textbf{(P\_8,1.70)} avaśyam ca etat evam vijñeyam . \\
\noindent \textbf{(P\_8,1.70)} yaḥ hi manyate pūrvam dhātuḥ sādhanena yujyate paścāt upasargeṇa iti āsyate guruṇā iti akarmakaḥ upāsyate guruḥ iti kena sakarmakaḥ syāt . \\
\noindent \textbf{(P\_8,1.70)} gatinā tu viśiṣṭasya gatiḥ eva viśeṣakaḥ . \\
\noindent \textbf{(P\_8,1.70)} sādhane kena te na syāt bāhyam ābhyantaraḥ hi saḥ . \\
\noindent \textbf{(P\_8,1.71)} tiṅgrahaṇam kimartham . \\
\noindent \textbf{(P\_8,1.71)} tiṅgrahaṇam udāttavataḥ parimāṇārtham . \\
\noindent \textbf{(P\_8,1.71)} tiṅgrahaṇam kriyate udāttavataḥ parimāṇārtham . \\
\noindent \textbf{(P\_8,1.71)} tiṅi udāttavati yathā syāt mandraśabde mā bhūt . \\
\noindent \textbf{(P\_8,1.71)} ā mandraiḥ indra haribhiḥ yāhi . \\
\noindent \textbf{(P\_8,1.71)} yadyogāt gatiḥ . \\
\noindent \textbf{(P\_8,1.71)} yatkriyāyuktāḥ tam prati gatyupasargasañjñe bhavataḥ na ca āṅaḥ mandraśabdam prati kriyāyogaḥ . \\
\noindent \textbf{(P\_8,1.71)} kim tarhi yāhiśabdam prati . \\
\noindent \textbf{(P\_8,1.71)} yadyogāt gatiḥ iti cet pratyayodāttatve aprasiddhiḥ . yadyogāt gatiḥ iti cet pratyayodāttatve aprasiddhiḥ syāt . \\
\noindent \textbf{(P\_8,1.71)} yatprakaroti . \\
\noindent \textbf{(P\_8,1.71)} tasmāt tiṅgrahaṇam kartavyam . \\
\noindent \textbf{(P\_8,1.71)} yadi tiṅgrahaṇam kriyate āmante na prāpnoti . \\
\noindent \textbf{(P\_8,1.71)} prapacatitarām . \\
\noindent \textbf{(P\_8,1.71)} prajalpatitarām . \\
\noindent \textbf{(P\_8,1.71)} asati punaḥ tiṅgrahaṇe kriyāpradhānam ākhyātam tasmāt atiśaye tarap utpadyate tarabantāt svārthe ām tatra yatkriyāyuktāḥ tam prati gatyupasargasañjñe bhavataḥ iti bhavati etam saṅghātam prati kriyāyogaḥ . \\
\noindent \textbf{(P\_8,1.71)} tasmāt na arthaḥ tiṅgrahaṇena . \\
\noindent \textbf{(P\_8,1.71)} kasmāt na bhavati . \\
\noindent \textbf{(P\_8,1.71)} a mandraiḥ indra haribhiḥ yāhi mayuraromabhiḥ . \\
\noindent \textbf{(P\_8,1.71)} yadyogāt gatiḥ iti . \\
\noindent \textbf{(P\_8,1.71)} nanu ca uktam yadyogāt gatiḥ iti cet pratyayodāttatve aprasiddhiḥ iti . \\
\noindent \textbf{(P\_8,1.71)} na eṣaḥ doṣaḥ . \\
\noindent \textbf{(P\_8,1.71)} yadkriyāyuktāḥ iti na evam vijñāyate yasya kriyā yatkriyā yatkriyāyuktāḥ tam prati gatyupasargasañjñe bhavataḥ iti . \\
\noindent \textbf{(P\_8,1.71)} katham tarhi . \\
\noindent \textbf{(P\_8,1.71)} yā kriyā yatkriyā yatkriyāyuktāḥ tam prati gatyupasargasañjñe bhavataḥ iti \\
\noindent \textbf{(P\_8,1.72.1)} vatkaraṇam kimartham . \\
\noindent \textbf{(P\_8,1.72.1)} svāśrayam api yathā syāt . \\
\noindent \textbf{(P\_8,1.72.1)} ām bhoḥ devadatta iti atra āmaekāntaramāmantritamanantike iti ekāntaratā yathā syāt \\
\noindent \textbf{(P\_8,1.72.2)} pūrvam prati vidyamānavattvāt uttaratra ānantaryāprasiddhiḥ . \\
\noindent \textbf{(P\_8,1.72.2)} pūrvam prati vidyamānavattvāt uttaratra ānantaryasya aprasiddhiḥ syāt . \\
\noindent \textbf{(P\_8,1.72.2)} imam me gaṅge yamune sarasvati . \\
\noindent \textbf{(P\_8,1.72.2)} gaṅgeśabdaḥ ayam yamuneśabdam prati avidyamānavat bhavati . \\
\noindent \textbf{(P\_8,1.72.2)} tatra āmantritasya padāt parasya iti anudāttatvam na syāt . \\
\noindent \textbf{(P\_8,1.72.2)} siddham tu padapūrvasya iti vacanāt . \\
\noindent \textbf{(P\_8,1.72.2)} siddham etat . \\
\noindent \textbf{(P\_8,1.72.2)} katham . \\
\noindent \textbf{(P\_8,1.72.2)} padapūrvasya iti vacanāt . \\
\noindent \textbf{(P\_8,1.72.2)} padapūrvasya ca āmantritasya avidyamānavadbhāvaḥ bhavati iti vaktavyam \\
\noindent \textbf{(P\_8,1.72.3)} kāni punaḥ asya yogasya prayojanāni . \\
\noindent \textbf{(P\_8,1.72.3)} avidyamānavattve prayojanam āmantritayuṣmadasmattiṅnighātāḥ . \\
\noindent \textbf{(P\_8,1.72.3)} āmantritasya padāt parasya anudāttaḥ bhavati iti iha eva bhavati pacasi devadatta . \\
\noindent \textbf{(P\_8,1.72.3)} devadatta yajñadatta iti atra na bhavati avidyamānavattvāt āmantritasya . \\
\noindent \textbf{(P\_8,1.72.3)} yuṣmadasmadoḥṣaṣṭhīcaturthīdvitīyāsthayorvāmnāvau iti iha eva bhavati grāmaḥ vām svam janapadaḥ nau svam . \\
\noindent \textbf{(P\_8,1.72.3)} devadattayajñadattau yuvayoḥ svam iti atra na bhavati avidyamānavattvāt āmantritasya . \\
\noindent \textbf{(P\_8,1.72.3)} tiṅatiṅaḥ iti iha eva bhavati devadattaḥ pacati . \\
\noindent \textbf{(P\_8,1.72.3)} devadatta pacasi iti atra na bhavati avidyamānavattvāt āmantritasya . \\
\noindent \textbf{(P\_8,1.72.3)} pūjāyām anantarapratiṣedhaḥ . pūjāyām anantarapratiṣedhaḥ prayojanam . \\
\noindent \textbf{(P\_8,1.72.3)} yāvat pacati śobhanam . \\
\noindent \textbf{(P\_8,1.72.3)} yāvat devadatta pacati iti atra api siddham bhavati . \\
\noindent \textbf{(P\_8,1.72.3)} jātu apūrvam . jātu apūrvam prayojanam . \\
\noindent \textbf{(P\_8,1.72.3)} jātu pacati . \\
\noindent \textbf{(P\_8,1.72.3)} devadatta jātu pacasi iti atra api siddham bhavati . \\
\noindent \textbf{(P\_8,1.72.3)} āho utāho ca anantaravidhau . āho utāho ca anantaravidhau prayojanam . \\
\noindent \textbf{(P\_8,1.72.3)} āho pacasi . \\
\noindent \textbf{(P\_8,1.72.3)} āho devadatta pacasi iti atra api siddham bhavati . \\
\noindent \textbf{(P\_8,1.72.3)} utāho pacasi . \\
\noindent \textbf{(P\_8,1.72.3)} utāho devadatta pacasi iti atra api siddham bhavati . \\
\noindent \textbf{(P\_8,1.72.3)} āmaḥ ekāntaravidhau . āmaḥ ekāntaravidhau prayojanam . \\
\noindent \textbf{(P\_8,1.72.3)} ām pacasi devadatta . \\
\noindent \textbf{(P\_8,1.72.3)} ām bhoḥ pacasi devadatta atra api siddham bhavati \\
\noindent \textbf{(P\_8,1.73)} iha kasmāt na bhavati . \\
\noindent \textbf{(P\_8,1.73)} aghnye devi sarasvati iḍe kavye vihavye etani te aghnye namāni . \\
\noindent \textbf{(P\_8,1.73)} yogavibhāgaḥ kariṣyate . \\
\noindent \textbf{(P\_8,1.73)} na āmantrite samānādhikaraṇe sāmānyavacanam . \\
\noindent \textbf{(P\_8,1.73)} tataḥ vibhāṣitam viśeṣavacane iti \\
\noindent \textbf{(P\_8,1.74)} iha kasmāt na bhavati . \\
\noindent \textbf{(P\_8,1.74)} brahmaṇa vaiyākaraṇa . \\
\noindent \textbf{(P\_8,1.74)} bahuvacanam iti vakṣyāmi . \\
\noindent \textbf{(P\_8,1.74)} sāmānyavacanam iti śakyam avaktum . \\
\noindent \textbf{(P\_8,1.74)} katham . \\
\noindent \textbf{(P\_8,1.74)} vibhāṣitam viśeṣavacane iti ucyate tena yat prati viśeṣavacanam iti etat bhavati tasya bhaviṣyati . \\
\noindent \textbf{(P\_8,1.74)} kim ca prati etat bhavati . \\
\noindent \textbf{(P\_8,1.74)} sāmānyavacanam . \\
\noindent \textbf{(P\_8,1.74)} aparaḥ āha : viśeṣavacane iti śakyam avaktum . \\
\noindent \textbf{(P\_8,1.74)} katham . \\
\noindent \textbf{(P\_8,1.74)} sāmānyavacanam vibhāṣitam iti ucyate tena yat prati sāmānyavacanam iti etat bhavati . \\
\noindent \textbf{(P\_8,1.74)} kim ca prati etat bhavati . \\
\noindent \textbf{(P\_8,1.74)} viśeṣavacanam . \\
\noindent \textbf{(P\_8,1.74)} sāmānyavacanam vibhāṣitam viśeṣavacane iti \\
\noindent \textbf{(P\_8,2.1.1)} yā iyam sapādasaptādhyāyī anukrāntā etasyām ayam pādonaḥ adhyāyaḥ asiddhaḥ veditavyaḥ . \\
\noindent \textbf{(P\_8,2.1.1)} yadi sapādāyām saptādhyāyyām ayam pādonaḥ adhyāyaḥ asiddhaḥ iti ucyate yaḥ iha saptamīnirdeśāḥ pañcamīnirdeśāḥ ca ucyante ṣaṣṭhīnirdeśāḥ ca ucyante te api asiddhāḥ syūḥ . \\
\noindent \textbf{(P\_8,2.1.1)} tatra kaḥ doṣaḥ . \\
\noindent \textbf{(P\_8,2.1.1)} jhalojhali hrasvādaṅgāt saṃyogāntasyalopaḥ iti eteṣām nirdeśānām asiddhatvāt tasminnitinirdeṣṭepūrvasya tasmādityuttarasya ṣaṣṭhīsthāneyogā iti etāḥ paribhāṣāḥ na prakalperan . \\
\noindent \textbf{(P\_8,2.1.1)} na eṣaḥ doṣaḥ . \\
\noindent \textbf{(P\_8,2.1.1)} yadi api idam tatra asiddham tat tu iha siddham . \\
\noindent \textbf{(P\_8,2.1.1)} katham . \\
\noindent \textbf{(P\_8,2.1.1)} kāryakālam sañjñāparibhāṣam yatra kāryam tatra draṣṭavyam . \\
\noindent \textbf{(P\_8,2.1.1)} jhalojhali . \\
\noindent \textbf{(P\_8,2.1.1)} hrasvādaṅgāt . \\
\noindent \textbf{(P\_8,2.1.1)} saṃyogāntasyalopaḥ . \\
\noindent \textbf{(P\_8,2.1.1)} upasthitam idam bhavati tasminnitinirdiṣṭepūrvasya tasmādityuttarasya ṣaṣṭhīsthāneyogā iti . \\
\noindent \textbf{(P\_8,2.1.1)} yadi kāryakālam sañjñāparibhāṣam iti ucyate iyam api paribhāṣā asti vipratiṣedhe param iti sā api iha upatiṣṭheta . \\
\noindent \textbf{(P\_8,2.1.1)} tatra kaḥ doṣaḥ . \\
\noindent \textbf{(P\_8,2.1.1)} visphoryam avagoryam iti guṇāt dīrghatvam syāt vipratiṣedhena . \\
\noindent \textbf{(P\_8,2.1.1)} ataḥ uttaram paṭhati . \\
\noindent \textbf{(P\_8,2.1.1)} pūrvatrāsiddhe na asti vipratiṣedhaḥ abhāvāt uttarasya . \\
\noindent \textbf{(P\_8,2.1.1)} pūrvatrāsiddhe na asti vipratiṣedhaḥ . \\
\noindent \textbf{(P\_8,2.1.1)} kim kāraṇam . \\
\noindent \textbf{(P\_8,2.1.1)} abhāvāt uttarasya . \\
\noindent \textbf{(P\_8,2.1.1)} dvayoḥ hi sāvakāśayoḥ samavasthitayoḥ vipratiṣedhaḥ bhavati na ca pūrvatrāsiddhe param pūrvam prati bhavati . \\
\noindent \textbf{(P\_8,2.1.1)} yadi evam dogdhā dogdhum ghatvasya asiddhatvāt ḍhatvam prāpnoti kāṣṭhataṭ kūṭataṭ saṃyogādilopsya asiddhatvāt saṃyogāntalopaḥ prāpnoti . \\
\noindent \textbf{(P\_8,2.1.1)} apavādaḥ vacanaprāmāṇyāt . anavakāśau etau vacanaprāmāṇyāt bhaviṣyataḥ . \\
\noindent \textbf{(P\_8,2.1.1)} tasmāt kāryakālam sañjñāparibhāṣam iti na doṣaḥ \\
\noindent \textbf{(P\_8,2.1.2)} pūrvatrāsiddham adhikāraḥ . \\
\noindent \textbf{(P\_8,2.1.2)} pūrvatrāsiddham iti adhikāraḥ ayam draṣṭavyaḥ . \\
\noindent \textbf{(P\_8,2.1.2)} kim prayojanam . \\
\noindent \textbf{(P\_8,2.1.2)} parasya parasya pūrvatra pūrvatra asiddhavijñānārtham . paraḥ paraḥ yogaḥ pūrvam pūrvam yogam prati asiddhaḥ yathā syāt . \\
\noindent \textbf{(P\_8,2.1.2)} anadhikāre hi samudāye asiddhavijñānam . anadhikāre hi sati samudāyasya samudāye asiddhatvam vijñāyeta . \\
\noindent \textbf{(P\_8,2.1.2)} tatra kaḥ doṣaḥ . \\
\noindent \textbf{(P\_8,2.1.2)} tatra ayatheṣṭaprasaṅgaḥ . tatra ayatheṣṭam prasajyeta . \\
\noindent \textbf{(P\_8,2.1.2)} yodhuṅmān guḍaliṇmān iti . \\
\noindent \textbf{(P\_8,2.1.2)} ghatvaḍhatvayoḥ kṛtayoḥ jhayaḥ iti vatvam prasajyeta . \\
\noindent \textbf{(P\_8,2.1.2)} tasmāt adhikāraḥ . tasmāt adhikāraḥ ayam draṣṭavyaḥ \\
\noindent \textbf{(P\_8,2.1.3)} asiddhavacanam kimartham . \\
\noindent \textbf{(P\_8,2.1.3)} asiddhavacane uktam . \\
\noindent \textbf{(P\_8,2.1.3)} kim uktam . \\
\noindent \textbf{(P\_8,2.1.3)} tatra tāvat uktam ṣatvatukoḥ asiddhavacanam ādeśalakṣaṇapratiṣedhārtham utsargalakṣaṇabhāvārtham ca iti . \\
\noindent \textbf{(P\_8,2.1.3)} evam iha api pūrvatrāsiddhavacanam ādeśalakṣaṇapratiṣedhārtham utsargalakṣaṇabhāvārtham ca . \\
\noindent \textbf{(P\_8,2.1.3)} ādeśalakṣaṇapratiṣedhārtham tāvat . \\
\noindent \textbf{(P\_8,2.1.3)} rājabhiḥ takṣabhiḥ rājabhyām takṣabhyām rājasu takṣasu iti nalope kṛte ataḥ iti aisbhāvādayaḥ prāpnuvanti . \\
\noindent \textbf{(P\_8,2.1.3)} asiddhatvāt na bhavanti . \\
\noindent \textbf{(P\_8,2.1.3)} utsargalakṣaṇabhāvārtham ca . \\
\noindent \textbf{(P\_8,2.1.3)} amuṣmai amuṣmāt amuṣya amuṣmin iti atra mubhāve kṛte ataḥ iti smāyādayaḥ na prāpnuvanti . \\
\noindent \textbf{(P\_8,2.1.3)} asiddhatvāt bhavanti . \\
\noindent \textbf{(P\_8,2.1.3)} suparvāṇau suparvāṇaḥ . \\
\noindent \textbf{(P\_8,2.1.3)} ṇatve kṛte nopadhāyāḥ iti dīrghatvam na prapnoti . \\
\noindent \textbf{(P\_8,2.1.3)} asiddhatvāt bhavati \\
\noindent \textbf{(P\_8,2.2.1)} subvidhim prati nalopaḥ asiddhaḥ bhavati iti ucyate . \\
\noindent \textbf{(P\_8,2.2.1)} bhavet iha rājabhiḥ takṣabhiḥ iti nalope kṛte ataḥ iti aisbhāvaḥ na syāt . \\
\noindent \textbf{(P\_8,2.2.1)} iha tu khalu rājabhyām takṣabhyām rājasu takṣasu iti nalope kṛte dīrghatvaittve prāpnutaḥ. na eṣaḥ doṣaḥ . \\
\noindent \textbf{(P\_8,2.2.1)} subvidhiḥ iti sarvavibhaktyantaḥ samāsaḥ : supaḥ vidhiḥ subvidhiḥ , supi vidhiḥ subvidhiḥ iti \\
\noindent \textbf{(P\_8,2.2.2)} atha sañjñāvidhau kim udāharaṇam . \\
\noindent \textbf{(P\_8,2.2.2)} pañca sapta . \\
\noindent \textbf{(P\_8,2.2.2)} pañca sapta iti atra nalope kṛte ṣṇāntāṣaṭ iti ṣaṭsañjñā na prāpnoti . \\
\noindent \textbf{(P\_8,2.2.2)} asiddhatvāt bhavati . \\
\noindent \textbf{(P\_8,2.2.2)} sañjñāgrahaṇānarthakyam ca tannimittatvāt lopasya . \\
\noindent \textbf{(P\_8,2.2.2)} sañjñāgrahaṇam ca anarthakam . \\
\noindent \textbf{(P\_8,2.2.2)} kim kāraṇam . \\
\noindent \textbf{(P\_8,2.2.2)} tannimittatvāt lopasya . \\
\noindent \textbf{(P\_8,2.2.2)} na akṛtāyām ṣaṭsañjñāyām jaśśasoḥ luk na ca akṛte luki padasañjñā na ca akṛtāyām padasañjñāyām nalopaḥ prāpnoti . \\
\noindent \textbf{(P\_8,2.2.2)} tat etat ānupūrvyā siddham bhavati . \\
\noindent \textbf{(P\_8,2.2.2)} idam tarhi prayojanam pañcabhiḥ saptabhiḥ iti ṣaṭtricaturbhyohalādiḥ jhalyupottamam iti eṣaḥ svaraḥ yathā syāt . \\
\noindent \textbf{(P\_8,2.2.2)} svare avadhāraṇāt ca . svare avadhāraṇāt ca sañjñāgrahaṇam anarthakam . \\
\noindent \textbf{(P\_8,2.2.2)} svare avadhāraṇam kriyate svaravidhim prati iti . \\
\noindent \textbf{(P\_8,2.2.2)} tugvidhau kim udāharaṇam . \\
\noindent \textbf{(P\_8,2.2.2)} vṛtrahabhyām vṛtrahabhiḥ . \\
\noindent \textbf{(P\_8,2.2.2)} nalope kṛte hrasvasyapitikṛtituk iti tuk prāpnoti . \\
\noindent \textbf{(P\_8,2.2.2)} asiddhatvāt na bhavati . \\
\noindent \textbf{(P\_8,2.2.2)} tugvidhau ca uktam . kim uktam . \\
\noindent \textbf{(P\_8,2.2.2)} sannipātalakṣaṇaḥ vidhiḥ animittam tadvighātasya iti . \\
\noindent \textbf{(P\_8,2.2.2)} idam tarhi prayojanam . \\
\noindent \textbf{(P\_8,2.2.2)} kṛti iti vakṣyāmi . \\
\noindent \textbf{(P\_8,2.2.2)} iha mā bhūt . \\
\noindent \textbf{(P\_8,2.2.2)} brahmahacchatram bhrūṇahacchāyā . \\
\noindent \textbf{(P\_8,2.2.2)} na eṣaḥ sannipātalakṣaṇaḥ \\
\noindent \textbf{(P\_8,2.3)} iha ne yat kāryam prāpnoti tat prati mubhāvaḥ na asiddhaḥ iti ucyate nābhāvaḥ ca eva tāvat na prāpnoti . \\
\noindent \textbf{(P\_8,2.3)} evam tarhi na mu ṭādeśe . \\
\noindent \textbf{(P\_8,2.3)} na mu ṭādeśe iti vaktavyam . \\
\noindent \textbf{(P\_8,2.3)} kim idam ṭādeśaḥ iti . \\
\noindent \textbf{(P\_8,2.3)} ṭāyāḥ ādeśaḥ ṭādeśaḥ iti . \\
\noindent \textbf{(P\_8,2.3)} yadi tarhi ṭāyāḥ ādeśe iti ucyate ṭāyām ādeśe aprasiddhiḥ . \\
\noindent \textbf{(P\_8,2.3)} tatra kaḥ doṣaḥ . \\
\noindent \textbf{(P\_8,2.3)} amunā iti atra mubhāvasya asiddhatvāt atodīrghoyañi supica iti dīrghatvam prasajyeta . \\
\noindent \textbf{(P\_8,2.3)} na eṣaḥ doṣaḥ . \\
\noindent \textbf{(P\_8,2.3)} sarvavibhaktyantaḥ samāsaḥ : ṭāyāḥ ādeśaḥ ṭādeśaḥ , ṭāyām ādeśaḥ ṭādeśaḥ iti . \\
\noindent \textbf{(P\_8,2.3)} sidhyati . \\
\noindent \textbf{(P\_8,2.3)} sūtram tarhi bhidyate . \\
\noindent \textbf{(P\_8,2.3)} yathānyāsam eva astu . \\
\noindent \textbf{(P\_8,2.3)} nanu ca uktam ne yat kāryam prāpnoti tasmin mubhāvaḥ na asiddhaḥ iti ucyate nābhāvaḥ ca eva tāvat na prāpnoti iti . \\
\noindent \textbf{(P\_8,2.3)} na eṣaḥ doṣaḥ . \\
\noindent \textbf{(P\_8,2.3)} iha iṅgitena ceṣṭitena nimiṣitena mahatā vā sūtranibandhena ācāryāṇām abhiprāyaḥ lakṣyate . \\
\noindent \textbf{(P\_8,2.3)} etat eva jñāpayati bhavati atra nābhāvaḥ iti yat ayam ne parataḥ asiddhatvapratiṣedham śāsti . \\
\noindent \textbf{(P\_8,2.3)} atha vā dvigatāḥ api hetavaḥ bhavanti . \\
\noindent \textbf{(P\_8,2.3)} tat yathā . \\
\noindent \textbf{(P\_8,2.3)} āmrāḥ ca siktāḥ pitaraḥ ca prīṇitāḥ bhavanti . \\
\noindent \textbf{(P\_8,2.3)} tathā vākyāni api dvigatāni dṛśyante . \\
\noindent \textbf{(P\_8,2.3)} śvetaḥ dhāvati . \\
\noindent \textbf{(P\_8,2.3)} alambusānām yātā iti . \\
\noindent \textbf{(P\_8,2.3)} atha vā vṛddhakumārīvākyavat idam draṣṭavyam . \\
\noindent \textbf{(P\_8,2.3)} tat yathā . \\
\noindent \textbf{(P\_8,2.3)} vṛddhakumārī indreṇa uktā varam vṛṇīṣva iti sā varam avṛṇīta putrāḥ me bahukṣīraghṛtam odanam kāṃsyapātryām bhuñjīran iti . \\
\noindent \textbf{(P\_8,2.3)} na ca tāvat asyāḥ patiḥ bhavati kutaḥ putrāḥ kutaḥ gāvaḥ kutaḥ dhānyam . \\
\noindent \textbf{(P\_8,2.3)} tatrs anayā ekena vākyena patiḥ putrāḥ gāvaḥ dhānyam iti sarvam saṅgṛhītam bhavati . \\
\noindent \textbf{(P\_8,2.3)} evam iha api ne asiddhatvapratiṣedham bruvatā nābhāvaḥ api saṅgṛhītaḥ bhavati \\
\noindent \textbf{(P\_8,2.4)} yaṇsvaraḥ yaṇādeśe svaritayaṇaḥ svaritārtham . \\
\noindent \textbf{(P\_8,2.4)} yaṇsvaraḥ yaṇādeśe siddhaḥ vaktavyaḥ . \\
\noindent \textbf{(P\_8,2.4)} kim prayojanam . \\
\noindent \textbf{(P\_8,2.4)} svaritayaṇaḥ svaritārtham . \\
\noindent \textbf{(P\_8,2.4)} svaritayaṇaḥ svaritatvam yathā syāt . \\
\noindent \textbf{(P\_8,2.4)} khalapvi aṭati . \\
\noindent \textbf{(P\_8,2.4)} khalpvi aśnāti . \\
\noindent \textbf{(P\_8,2.4)} tat tarhi vaktavyam . \\
\noindent \textbf{(P\_8,2.4)} na vaktavyam . \\
\noindent \textbf{(P\_8,2.4)} āha ayam svaritayaṇaḥ iti na ca asti siddhaḥ svaritaḥ tatra āśrayāt siddhatvam bhaviṣyati . \\
\noindent \textbf{(P\_8,2.4)} āśrayāt siddhatvam iti cet udāttāt svarite doṣaḥ . āśrayāt siddhatvam iti cet udāttāt svarite doṣaḥ bhavati . \\
\noindent \textbf{(P\_8,2.4)} dadhyāśa . \\
\noindent \textbf{(P\_8,2.4)} madhvāśa . \\
\noindent \textbf{(P\_8,2.4)} evam tarhi yogavibhāgaḥ kariṣyate . \\
\noindent \textbf{(P\_8,2.4)} udāttayaṇaḥ parasya anudāttasya svaritaḥ bhavati . \\
\noindent \textbf{(P\_8,2.4)} tataḥ svaritayaṇaḥ . \\
\noindent \textbf{(P\_8,2.4)} svaritayaṇaḥ ca parasya anudāttasya svaritaḥ bhavati . \\
\noindent \textbf{(P\_8,2.4)} udāttayaṇaḥ iti eva . \\
\noindent \textbf{(P\_8,2.4)} atha vā svaritagrahaṇam na kariṣyate . \\
\noindent \textbf{(P\_8,2.4)} kena idānīm svaritayaṇaḥ parasya anudāttasya svaritaḥ bhaviṣyati . \\
\noindent \textbf{(P\_8,2.4)} udāttayaṇaḥ iti eva . \\
\noindent \textbf{(P\_8,2.4)} nanu ca svaritayaṇā vyavahitatvāt na prāpnoti . \\
\noindent \textbf{(P\_8,2.4)} svaravidhau vyañjanam avidyamānavat iti na asti vyavadhānam . \\
\noindent \textbf{(P\_8,2.4)} atha vā na evam vijñāyate svaritasya yaṇ svaritayaṇ svaritayaṇaḥ iti . \\
\noindent \textbf{(P\_8,2.4)} katham tarhi . \\
\noindent \textbf{(P\_8,2.4)} svarite yaṇ svaritayaṇ svaritayaṇaḥ iti \\
\noindent \textbf{(P\_8,2.6.1)} svaritagrahaṇam śakyam akartum . \\
\noindent \textbf{(P\_8,2.6.1)} katham . \\
\noindent \textbf{(P\_8,2.6.1)} anudātte parataḥ padādau vā udāttaḥ iti eva siddham . \\
\noindent \textbf{(P\_8,2.6.1)} kena idānīm svaritaḥ bhaviṣyati . \\
\noindent \textbf{(P\_8,2.6.1)} gāṅge anūpe iti . \\
\noindent \textbf{(P\_8,2.6.1)} āntaryataḥ udāttānudāttayoḥ ekādeśaḥ svaritaḥ bhaviṣyati . \\
\noindent \textbf{(P\_8,2.6.1)} idam tarhi prayojanam tena varjyamānatā mā bhūt . \\
\noindent \textbf{(P\_8,2.6.1)} atha kriyamāṇe api svaritagrahaṇe yaḥ siddhaḥ svaritaḥ tena varjyamānatā kasmāt na bhavati . \\
\noindent \textbf{(P\_8,2.6.1)} kanyā anūpe iti . \\
\noindent \textbf{(P\_8,2.6.1)} bahiraṅgalakṣaṇatvāt . \\
\noindent \textbf{(P\_8,2.6.1)} asiddham bahiraṅgam antaraṅge iti evam na bhaviṣyati . \\
\noindent \textbf{(P\_8,2.6.1)} yathā eva tarhi kriyamāṇe svaritagrahaṇe yaḥ siddhaḥ svaritaḥ tena varjyamānatā na bhavati evam akriyamāṇe api na bhaviṣyati . \\
\noindent \textbf{(P\_8,2.6.1)} tasmāt na arthaḥ svaritagrahaṇena . \\
\noindent \textbf{(P\_8,2.6.1)} bahiraṅgalakṣaṇatvāt siddham \\
\noindent \textbf{(P\_8,2.6.2)} ekādeśasvaraḥ antaraṅgaḥ . \\
\noindent \textbf{(P\_8,2.6.2)} ekādeśasvaraḥ antaraṅgaḥ siddhaḥ vaktavyaḥ . \\
\noindent \textbf{(P\_8,2.6.2)} kim prayojanam . \\
\noindent \textbf{(P\_8,2.6.2)} ayavāyāvekādeśaśatṛsvaraikānanudāttasarvānudāttārtham . ay . \\
\noindent \textbf{(P\_8,2.6.2)} vṛkṣe idam plakṣe idam . \\
\noindent \textbf{(P\_8,2.6.2)} udāttānudāttayoḥ ekādeśaḥ . \\
\noindent \textbf{(P\_8,2.6.2)} tasya ekādeśe udāttenodāttaḥ iti etat bhavati . \\
\noindent \textbf{(P\_8,2.6.2)} tasya siddhatvam vaktavyam āntaryataḥ udāttasya udāttaḥ ayādeśaḥ yathā syāt . \\
\noindent \textbf{(P\_8,2.6.2)} avādeśaḥ na asti . \\
\noindent \textbf{(P\_8,2.6.2)} āy . \\
\noindent \textbf{(P\_8,2.6.2)} kumāryai idam . \\
\noindent \textbf{(P\_8,2.6.2)} udāttānudāttayoḥ ekādeśaḥ . \\
\noindent \textbf{(P\_8,2.6.2)} tasya ekādeśe udāttenodāttaḥ iti etat bhavati . \\
\noindent \textbf{(P\_8,2.6.2)} tasya siddhatvam vaktavyam āntaryataḥ udāttasya udāttaḥ āyādeśaḥ yathā syāt . \\
\noindent \textbf{(P\_8,2.6.2)} na etat asti prayojanam . \\
\noindent \textbf{(P\_8,2.6.2)} ekādeśe kṛte udāttayaṇohalpūrvāt iti udāttatvam bhaviṣyati . \\
\noindent \textbf{(P\_8,2.6.2)} idam iha sampradhāryam . \\
\noindent \textbf{(P\_8,2.6.2)} udāttatvam kriyatām ekādeśaḥ iti kim atra kartavyam . \\
\noindent \textbf{(P\_8,2.6.2)} paratvāt udāttatvam . \\
\noindent \textbf{(P\_8,2.6.2)} nityaḥ ekādeśaḥ . \\
\noindent \textbf{(P\_8,2.6.2)} kṛte api udāttatve prāpnoti akṛte api prāpnoti . \\
\noindent \textbf{(P\_8,2.6.2)} ekādeśaḥ api anityaḥ . \\
\noindent \textbf{(P\_8,2.6.2)} anyathāsvarasya kṛte udāttatve prāpnoti anyathāsvarasya akṛte svarabhinnasya ca prāpnuvan vidhiḥ anityaḥ bhavati . \\
\noindent \textbf{(P\_8,2.6.2)} antaraṅgaḥ tarhi ekādeśaḥ . \\
\noindent \textbf{(P\_8,2.6.2)} kā antaraṅgatā . \\
\noindent \textbf{(P\_8,2.6.2)} varṇau āśritya ekādeśaḥ padasya udāttatvam . \\
\noindent \textbf{(P\_8,2.6.2)} evam tarhi idam iha sampradhāryam . \\
\noindent \textbf{(P\_8,2.6.2)} āṭ kriyatām udāttatvam iti kim atra kartavyam . \\
\noindent \textbf{(P\_8,2.6.2)} paratvāt āḍāgamaḥ . \\
\noindent \textbf{(P\_8,2.6.2)} nityam udāttatvam . \\
\noindent \textbf{(P\_8,2.6.2)} kṛte api āṭi prāpnoti akṛte api prāpnoti . \\
\noindent \textbf{(P\_8,2.6.2)} āṭ api nityaḥ . \\
\noindent \textbf{(P\_8,2.6.2)} kṛte api udāttatve prāpnoti akṛte api prāpnoti . \\
\noindent \textbf{(P\_8,2.6.2)} anityaḥ āṭ . \\
\noindent \textbf{(P\_8,2.6.2)} anyathāsvarasya kṛte udāttatve prāpnoti anyathāsvarasya akṛte svarabhinnasya ca prāpnuvan vidhiḥ anityaḥ bhavati . \\
\noindent \textbf{(P\_8,2.6.2)} udāttatvam api anityam . \\
\noindent \textbf{(P\_8,2.6.2)} anyasya kṛte āṭi prāpnoti anyasya akṛte prāpnoti śabdāntarasya ca prāpnuvan vidhiḥ anityaḥ bhavati . \\
\noindent \textbf{(P\_8,2.6.2)} ubhayoḥ anityayoḥ paratvāt āḍāgamaḥ āṭi kṛte antaraṅgaḥ ekādeśaḥ . \\
\noindent \textbf{(P\_8,2.6.2)} āv . \\
\noindent \textbf{(P\_8,2.6.2)} vṛkṣavidam . \\
\noindent \textbf{(P\_8,2.6.2)} udāttānudāttayoḥ ekādeśaḥ . \\
\noindent \textbf{(P\_8,2.6.2)} tasya ekādeśe udāttena udāttaḥ iti etat bhavati . \\
\noindent \textbf{(P\_8,2.6.2)} tasya siddhatvam vaktavyam āntaryataḥ udāttasya udāttaḥ āvādeśaḥ yathā syāt . \\
\noindent \textbf{(P\_8,2.6.2)} ekādeśasvara . \\
\noindent \textbf{(P\_8,2.6.2)} gāṅge anūpe iti . \\
\noindent \textbf{(P\_8,2.6.2)} udāttānudāttayoḥ ekādeśaḥ . \\
\noindent \textbf{(P\_8,2.6.2)} tasya ekādeśe udāttena udāttaḥ iti etat bhavati . \\
\noindent \textbf{(P\_8,2.6.2)} tasya siddhatvam vaktavyam svaritovānudāttepadādau iti etat yathā syāt . \\
\noindent \textbf{(P\_8,2.6.2)} śatṛsvara . \\
\noindent \textbf{(P\_8,2.6.2)} tudati nudati . \\
\noindent \textbf{(P\_8,2.6.2)} udāttānudāttayoḥ ekādeśaḥ . \\
\noindent \textbf{(P\_8,2.6.2)} tasya ekādeśe udāttena udāttaḥ iti etat bhavati . \\
\noindent \textbf{(P\_8,2.6.2)} tasya siddhatvam vaktavyam śatuḥ anumaḥ nadyajādiḥ antodāttāt iti eṣaḥ svaraḥ yathā syāt . \\
\noindent \textbf{(P\_8,2.6.2)} na etat asti prayojanam . \\
\noindent \textbf{(P\_8,2.6.2)} ācāryapravṛttiḥ jñāpayati siddhaḥ ekādeśasvaraḥ śatṛsvaraḥ iti yat ayam anumaḥ iti pratiṣedham śāsti . \\
\noindent \textbf{(P\_8,2.6.2)} katham kṛtvā jñāpakam . \\
\noindent \textbf{(P\_8,2.6.2)} na hi antareṇa udāttānudāttoḥ ekādeśam śatrantam sanumkam antodāttam asti . \\
\noindent \textbf{(P\_8,2.6.2)} nanu ca idam asti yantī vantī . \\
\noindent \textbf{(P\_8,2.6.2)} etat api nighāte kṛte na antareṇa udāttānudāttayoḥ ekādeśam śatrantam sanumkam antodāttam asti . \\
\noindent \textbf{(P\_8,2.6.2)} idam iha sampradhāryam . \\
\noindent \textbf{(P\_8,2.6.2)} nighātaḥ kriyatām ekādeśaḥ iti kim atra kartavyam . \\
\noindent \textbf{(P\_8,2.6.2)} paratvāt nighātaḥ . \\
\noindent \textbf{(P\_8,2.6.2)} nityaḥ ekādeśaḥ . \\
\noindent \textbf{(P\_8,2.6.2)} kṛte api nighāte prāpnoti akṛte api prāpnoti . \\
\noindent \textbf{(P\_8,2.6.2)} ekādeśaḥ api anityaḥ . \\
\noindent \textbf{(P\_8,2.6.2)} anyathāsvarasya kṛte nighāte prāpnoti anyathāsvarasya akṛte nighāte svarabhinnasya ca prāpnuvan vidhiḥ anityaḥ bhavati . \\
\noindent \textbf{(P\_8,2.6.2)} antaraṅgaḥ tarhi ekādeśaḥ . \\
\noindent \textbf{(P\_8,2.6.2)} kā antaraṅgatā . \\
\noindent \textbf{(P\_8,2.6.2)} varṇau āśritya ekādeśaḥ padasya nighātaḥ . \\
\noindent \textbf{(P\_8,2.6.2)} nighātaḥ api antaraṅgaḥ . \\
\noindent \textbf{(P\_8,2.6.2)} katham . \\
\noindent \textbf{(P\_8,2.6.2)} uktam etat padagrahaṇam parimāṇārtham iti . \\
\noindent \textbf{(P\_8,2.6.2)} ubhayoḥ antaraṅgayoḥ paratvāt nighātaḥ nighāte kṛte etat api na antareṇa udāttānudāttayoḥ ekādeśam antodāttam bhavati . \\
\noindent \textbf{(P\_8,2.6.2)} śatṛsvara . \\
\noindent \textbf{(P\_8,2.6.2)} ekānudātta . \\
\noindent \textbf{(P\_8,2.6.2)} tudanti likhanti. udāttānudāttayoḥ ekādeśaḥ . \\
\noindent \textbf{(P\_8,2.6.2)} tasya ekādeśe udāttena udāttaḥ iti etat bhavati . \\
\noindent \textbf{(P\_8,2.6.2)} tasya siddhatvam vaktavyam tena varjyamānatā yathā syāt . \\
\noindent \textbf{(P\_8,2.6.2)} sarvānudātta . \\
\noindent \textbf{(P\_8,2.6.2)} brāhmaṇāḥ tudanti . \\
\noindent \textbf{(P\_8,2.6.2)} brāhmaṇāḥ likhanti . \\
\noindent \textbf{(P\_8,2.6.2)} udāttānudāttayoḥ ekādeśaḥ . \\
\noindent \textbf{(P\_8,2.6.2)} tasya ekādeśe udāttenodāttaḥ iti etat bhavati . \\
\noindent \textbf{(P\_8,2.6.2)} tasya siddhatvam vaktavyam . \\
\noindent \textbf{(P\_8,2.6.2)} kim prayojanam . \\
\noindent \textbf{(P\_8,2.6.2)} tiṅatiṅaḥ iti nighātaḥ yathā syāt . \\
\noindent \textbf{(P\_8,2.6.2)} kim ucyate antaraṅgaḥ iti . \\
\noindent \textbf{(P\_8,2.6.2)} yaḥ hi bahiraṅgaḥ asiddhaḥ eva asau bhavati . \\
\noindent \textbf{(P\_8,2.6.2)} prapacatiti . \\
\noindent \textbf{(P\_8,2.6.2)} somasut pacatiti . \\
\noindent \textbf{(P\_8,2.6.2)} tat tarhi vaktavyam . \\
\noindent \textbf{(P\_8,2.6.2)} na vaktavyam . \\
\noindent \textbf{(P\_8,2.6.2)} sarvatra eva numpratiṣedhaḥ jñāpakaḥ siddhaḥ ekādeśasvaraḥ antaraṅgaḥ iti . \\
\noindent \textbf{(P\_8,2.6.2)} saṃyogāntalopaḥ roḥ uttve harivaḥ medinam tvā . saṃyogāntalopaḥ roḥ uttve siddhaḥ vaktavyaḥ . \\
\noindent \textbf{(P\_8,2.6.2)} kim prayojanam . \\
\noindent \textbf{(P\_8,2.6.2)} harivaḥ medinam tvā . \\
\noindent \textbf{(P\_8,2.6.2)} saṃyogāntalopasya asiddhatvāt haśi iti uttvam na prāpnoti . \\
\noindent \textbf{(P\_8,2.6.2)} plutiḥ ca . plutiḥ ca uttve siddhā vaktavyā . \\
\noindent \textbf{(P\_8,2.6.2)} susrota3 atra nu asi iti atra pluteḥ asiddhatvāt ataḥ ati iti uttvam prāpnoti . \\
\noindent \textbf{(P\_8,2.6.2)} aplutāt aplute iti etat na vaktavyam bhavati . \\
\noindent \textbf{(P\_8,2.6.2)} na etat asti prayojanam . \\
\noindent \textbf{(P\_8,2.6.2)} kriyate nyāse eva . \\
\noindent \textbf{(P\_8,2.6.2)} sijlopaḥ ekādeśe . sijlopaḥ ekādeśe siddhaḥ vaktavyaḥ . \\
\noindent \textbf{(P\_8,2.6.2)} alāvīt apāvīt . \\
\noindent \textbf{(P\_8,2.6.2)} sijlopasya asiddhatvāt savarṇadīrghatvam na prāpnoti . \\
\noindent \textbf{(P\_8,2.6.2)} yadi punaḥ iḍādeḥ sicaḥ lopaḥ ucyeta . \\
\noindent \textbf{(P\_8,2.6.2)} na evam śakyam . \\
\noindent \textbf{(P\_8,2.6.2)} iha hi mā hi lāvit mā hi pāvit yadi atra iṭ na syāt anudāttasya īṭaḥ śravaṇam prasajyeta . \\
\noindent \textbf{(P\_8,2.6.2)} iṭi punaḥ sati uktam etat arthavat tu citkaraṇasāmarthyāt hi iṭaḥ udāttatvam iti tatra ekādeśe udāttenodāttaḥ iti udāttatvam siddham bhavati . \\
\noindent \textbf{(P\_8,2.6.2)} saṃyogādilopaḥ saṃyogāntalope . saṃyogādilopaḥ saṃyogāntasya lope siddhaḥ vaktavyaḥ . \\
\noindent \textbf{(P\_8,2.6.2)} kāṣṭhataṭ kūṭataṭ . \\
\noindent \textbf{(P\_8,2.6.2)} saṃyogādilopasya asiddhatvāt saṃyogāntalopaḥ prāpnoti . \\
\noindent \textbf{(P\_8,2.6.2)} na eṣaḥ doṣaḥ . \\
\noindent \textbf{(P\_8,2.6.2)} ukam etat apavādaḥ vacanaprāmāṇyāt iti . \\
\noindent \textbf{(P\_8,2.6.2)} niṣṭhādeśaḥ ṣatvasvarapratyayeḍvidhiṣu . niṣṭhādeśaḥ ṣatvasvarapratyayeḍvidhiṣu siddhaḥ vaktavyaḥ . \\
\noindent \textbf{(P\_8,2.6.2)} vṛkṇaḥ vṛkṇavān . \\
\noindent \textbf{(P\_8,2.6.2)} niṣṭhādeśasya asiddhatvāt jhali iti ṣatvam prāpnoti . \\
\noindent \textbf{(P\_8,2.6.2)} svara . \\
\noindent \textbf{(P\_8,2.6.2)} kṣivaḥ . \\
\noindent \textbf{(P\_8,2.6.2)} niṣṭhādeśasya asiddhatvāt niṣṭhācadvyajanāt iti eṣaḥ svaraḥ na prāpnoti . \\
\noindent \textbf{(P\_8,2.6.2)} pratyaya . \\
\noindent \textbf{(P\_8,2.6.2)} kṣīveṇa tarati kṣīvikaḥ . \\
\noindent \textbf{(P\_8,2.6.2)} niṣṭhādeśasya asiddhatvāt dvyacaḥ ṭhan iti ṭhan na prāpnoti . \\
\noindent \textbf{(P\_8,2.6.2)} iḍvidhi . \\
\noindent \textbf{(P\_8,2.6.2)} niṣṭhādeśasya asiddhatvāt valādilakṣaṇaḥ iṭ prāpnoti . \\
\noindent \textbf{(P\_8,2.6.2)} nanu ca yaḥ pratyayavidhau siddhaḥ siddhaḥ asau iḍvidhau . \\
\noindent \textbf{(P\_8,2.6.2)} idam tarhi prayojanam . \\
\noindent \textbf{(P\_8,2.6.2)} olasjī lagnaḥ . \\
\noindent \textbf{(P\_8,2.6.2)} niṣṭhādeśaḥ siddhaḥ vaktavyaḥ neḍvaśikṛti iti iṭpratiṣedhaḥ yathā syāt . \\
\noindent \textbf{(P\_8,2.6.2)} īditkaraṇam na kartavyam bhavati . \\
\noindent \textbf{(P\_8,2.6.2)} etat api na asti prayojanam . \\
\noindent \textbf{(P\_8,2.6.2)} kriyate etat nyāse eva . \\
\noindent \textbf{(P\_8,2.6.2)} vasvādiṣu datvam sau dīrghatve . vasvādiṣu datvam sau dīrghatve siddham vaktavyam . \\
\noindent \textbf{(P\_8,2.6.2)} ukhāsrat . \\
\noindent \textbf{(P\_8,2.6.2)} parṇadhvat . \\
\noindent \textbf{(P\_8,2.6.2)} datvasya asiddhatvāt atvasantasya iti dīrghatvam prāpnoti . \\
\noindent \textbf{(P\_8,2.6.2)} adhātoḥ iti na vaktavyam bhavati . \\
\noindent \textbf{(P\_8,2.6.2)} na etat asti prayojanam . \\
\noindent \textbf{(P\_8,2.6.2)} kriyate nyāse eva . \\
\noindent \textbf{(P\_8,2.6.2)} adasaḥ īttvotve svare bahiṣpadalakṣaṇe . adasaḥ īttvotve svare bahiṣpadalakṣaṇe siddhe vaktavye . \\
\noindent \textbf{(P\_8,2.6.2)} amī atra . \\
\noindent \textbf{(P\_8,2.6.2)} amī āsate . \\
\noindent \textbf{(P\_8,2.6.2)} amū atra . \\
\noindent \textbf{(P\_8,2.6.2)} amū āsāte . \\
\noindent \textbf{(P\_8,2.6.2)} īttvotvayoḥ asiddhatvāt ecaḥ iti ayāvekādeśāḥ prāpnuvanti . \\
\noindent \textbf{(P\_8,2.6.2)} kim ucyate bahiṣpadalakṣaṇe iti . \\
\noindent \textbf{(P\_8,2.6.2)} yaḥ hi anyaḥ asiddhaḥ eva asau bhavati . \\
\noindent \textbf{(P\_8,2.6.2)} amuyā amuyoḥ iti . \\
\noindent \textbf{(P\_8,2.6.2)} pragṛhyasañjñāyām ca . pragṛhyasañjñāyām ca siddhe vaktavye . \\
\noindent \textbf{(P\_8,2.6.2)} amī atra . \\
\noindent \textbf{(P\_8,2.6.2)} amī āsate . \\
\noindent \textbf{(P\_8,2.6.2)} amū atra . \\
\noindent \textbf{(P\_8,2.6.2)} amū āsāte . \\
\noindent \textbf{(P\_8,2.6.2)} īttvotvayoḥ asiddhatvāt adasomāt iti pragṛhyasañjñā na prāpnoti . \\
\noindent \textbf{(P\_8,2.6.2)} kim artham idam ubhayam ucyate na pragṛhyasañjñāyām iti eva svare api bahiṣpadalakṣaṇe coditam syāt . \\
\noindent \textbf{(P\_8,2.6.2)} purastāt idam ācāryeṇa dṛṣṭam svare bahiṣpadalakṣaṇe iti tat paṭhitam . \\
\noindent \textbf{(P\_8,2.6.2)} tataḥ uttarakālam idam dṛṣṭam pragṛhyasañjñāyām ca iti tad api paṭhitam . \\
\noindent \textbf{(P\_8,2.6.2)} na ca idānīm ācāryāḥ sūtrāṇi kṛtvā nivartayanti . \\
\noindent \textbf{(P\_8,2.6.2)} plutiḥ tugvidhau che . plutiḥ tugvidhau che siddhā vakavyā . \\
\noindent \textbf{(P\_8,2.6.2)} agna3i cchattram . \\
\noindent \textbf{(P\_8,2.6.2)} paṭa3u cchattram . \\
\noindent \textbf{(P\_8,2.6.2)} pluteḥ asiddhatvāt checa iti tuk na prāpnoti . \\
\noindent \textbf{(P\_8,2.6.2)} kim ucyate che iti . \\
\noindent \textbf{(P\_8,2.6.2)} yaḥ hi anyaḥ asiddhaḥ eva asau bhavati . \\
\noindent \textbf{(P\_8,2.6.2)} agnici3t somasu3t .ścutvam dhuṭtve . ścutvam dhuṭtve siddham vaktavyam . \\
\noindent \textbf{(P\_8,2.6.2)} aṭ ścyotati . \\
\noindent \textbf{(P\_8,2.6.2)} paṭ scyotati . \\
\noindent \textbf{(P\_8,2.6.2)} ścutvasya asiddhatvāt ḍaḥsidhuṭ iti dhuṭ prasajyeta . \\
\noindent \textbf{(P\_8,2.6.2)} abhyāsajaśtvacartvam ettvatukoḥ . abhyāsajaśtvacartvam ettvatukoḥ siddham vaktavyam . \\
\noindent \textbf{(P\_8,2.6.2)} babhaṇatuḥ babhaṇuḥ . \\
\noindent \textbf{(P\_8,2.6.2)} abhyāsādeśasya asiddhatvāt ettvam prāpnoti . \\
\noindent \textbf{(P\_8,2.6.2)} ucicchiṣati . \\
\noindent \textbf{(P\_8,2.6.2)} abhyāsādeśasya asiddhatvāt checa iti tuk prāpnoti . \\
\noindent \textbf{(P\_8,2.6.2)} dvirvacane parasavarṇatvam . dvirvacane parasavarṇatvam siddham vaktavyam . \\
\noindent \textbf{(P\_8,2.6.2)} sayṃyantā savṃvatsaraḥ talṃ lokam yalṃ lokam iti parasavarṇasya asiddhatvāt yaraḥ iti dvirvacanam na prāpnoti . \\
\noindent \textbf{(P\_8,2.6.2)} padādhikāraḥ cet latvaghatvanatvarutvaṣatvaṇatvānunāsikachatvāni . padādhikāraḥ cet latvaghatvanatvarutvaṣatvaṇatvānunāsikachatvāni siddhāni vaktavyāni . \\
\noindent \textbf{(P\_8,2.6.2)} latva garaḥ garaḥ . \\
\noindent \textbf{(P\_8,2.6.2)} galaḥ galaḥ . \\
\noindent \textbf{(P\_8,2.6.2)} latva. ghatva . \\
\noindent \textbf{(P\_8,2.6.2)} drogdhā drogdhā . \\
\noindent \textbf{(P\_8,2.6.2)} droḍhā droḍhā . \\
\noindent \textbf{(P\_8,2.6.2)} ghatva . \\
\noindent \textbf{(P\_8,2.6.2)} natva. nunnaḥ nunnaḥ . \\
\noindent \textbf{(P\_8,2.6.2)} nuttaḥ nuttaḥ . \\
\noindent \textbf{(P\_8,2.6.2)} natva . \\
\noindent \textbf{(P\_8,2.6.2)} rutva . \\
\noindent \textbf{(P\_8,2.6.2)} abhinaḥ abhinaḥ . \\
\noindent \textbf{(P\_8,2.6.2)} abhinat abhinat . \\
\noindent \textbf{(P\_8,2.6.2)} rutva ṣatva . \\
\noindent \textbf{(P\_8,2.6.2)} mātuḥṣvasā mātuḥṣvasā . \\
\noindent \textbf{(P\_8,2.6.2)} mātuḥsvasā mātuḥsvasā pituḥṣvasā pituḥṣvasā . \\
\noindent \textbf{(P\_8,2.6.2)} pituḥsvasā pituḥsvasā . \\
\noindent \textbf{(P\_8,2.6.2)} ṣatva . \\
\noindent \textbf{(P\_8,2.6.2)} ṇatva . \\
\noindent \textbf{(P\_8,2.6.2)} māṣavāpāṇi māṣavāpāṇi . \\
\noindent \textbf{(P\_8,2.6.2)} māṣavāpāni māṣavāpāni . \\
\noindent \textbf{(P\_8,2.6.2)} ṇatva . \\
\noindent \textbf{(P\_8,2.6.2)} anunāsika . \\
\noindent \textbf{(P\_8,2.6.2)} nāṅnayanam nāṅnayanam . \\
\noindent \textbf{(P\_8,2.6.2)} vāgnayanam vāgnayanam . \\
\noindent \textbf{(P\_8,2.6.2)} anunāsika . \\
\noindent \textbf{(P\_8,2.6.2)} chatva . \\
\noindent \textbf{(P\_8,2.6.2)} vākchayanam vākchayanam . \\
\noindent \textbf{(P\_8,2.6.2)} vākśayanam vākśayanam . \\
\noindent \textbf{(P\_8,2.6.2)} ubhayathā ca ayam doṣaḥ yadi api sthāne dvirvacanam atha api dviḥprayogaḥ . \\
\noindent \textbf{(P\_8,2.6.2)} katham . \\
\noindent \textbf{(P\_8,2.6.2)} yadi tāvat sthāne dvirvacanam sampramugdhatvāt prakṛtipratyayasya latvādyabhāvaḥ . \\
\noindent \textbf{(P\_8,2.6.2)} atha dviḥprayogaḥ asiddhatvāt latvādīni nivarteran \\
\noindent \textbf{(P\_8,2.7.1)} antagrahaṇam kimartham . \\
\noindent \textbf{(P\_8,2.7.1)} nalope antagrahaṇam padādhikārasya viśeṣaṇatvāt . \\
\noindent \textbf{(P\_8,2.7.1)} nalope antagrahaṇam kriyate . \\
\noindent \textbf{(P\_8,2.7.1)} kim kāraṇam . \\
\noindent \textbf{(P\_8,2.7.1)} padādhikārasya viśeṣaṇatvāt . \\
\noindent \textbf{(P\_8,2.7.1)} padādhikāraḥ viśeṣaṇam . \\
\noindent \textbf{(P\_8,2.7.1)} katham . \\
\noindent \textbf{(P\_8,2.7.1)} padasya iti na eṣā sthānaṣaṣṭhī . \\
\noindent \textbf{(P\_8,2.7.1)} kā tarhi . \\
\noindent \textbf{(P\_8,2.7.1)} viśeṣaṇaṣaṣṭhī \\
\noindent \textbf{(P\_8,2.7.2)} ahnaḥ nalopapratiṣedhaḥ . \\
\noindent \textbf{(P\_8,2.7.2)} ahnaḥ nalopapratiṣedhaḥ vaktavyaḥ . \\
\noindent \textbf{(P\_8,2.7.2)} ahobhyām ahobhiḥ iti . \\
\noindent \textbf{(P\_8,2.7.2)} saḥ tarhi pratiṣedhaḥ vaktavyaḥ . \\
\noindent \textbf{(P\_8,2.7.2)} na vaktavyaḥ . \\
\noindent \textbf{(P\_8,2.7.2)} ruḥ atra bādhakaḥ bhaviṣyati . \\
\noindent \textbf{(P\_8,2.7.2)} asiddhaḥ ruḥ tasya asiddhatvāt nalopaḥ prāpnoti . \\
\noindent \textbf{(P\_8,2.7.2)} anavakāśaḥ ruḥ nalopam bādhiṣyate . \\
\noindent \textbf{(P\_8,2.7.2)} sāvakāśaḥ ruḥ . \\
\noindent \textbf{(P\_8,2.7.2)} kaḥ avakāśaḥ . \\
\noindent \textbf{(P\_8,2.7.2)} anantyaḥ akāraḥ . \\
\noindent \textbf{(P\_8,2.7.2)} ācāryapravṛttiḥ jñāpayati na anantyasya ruḥ bhavati iti yat ayam ahangrahaṇam karoti . \\
\noindent \textbf{(P\_8,2.7.2)} ahangrahaṇāt iti cet sambuddhyartham vacanam . ahangrahaṇāt iti cet sambuddhyartham etat syāt . \\
\noindent \textbf{(P\_8,2.7.2)} he ahaḥ iti . \\
\noindent \textbf{(P\_8,2.7.2)} yat tarhi rutvam śāsti . \\
\noindent \textbf{(P\_8,2.7.2)} etat api sambuddhyartham eva syāt . \\
\noindent \textbf{(P\_8,2.7.2)} he dīrghāhaḥ atra . \\
\noindent \textbf{(P\_8,2.7.2)} yat tarhi rūparātrirathantareṣu upasaṅkhyānam karoti tat jñāpayati ācāryaḥ na anantyasya ruḥ bhavati iti . \\
\noindent \textbf{(P\_8,2.7.2)} katham kṛtvā jñāpakam . \\
\noindent \textbf{(P\_8,2.7.2)} na hi asti viśeṣaḥ rūparātrirathantareṣu anantyasya rau vā re vā \\
\noindent \textbf{(P\_8,2.8)} na ṅisambuddhyoḥ anuttarapade . \\
\noindent \textbf{(P\_8,2.8)} na ṅisambuddhyoḥ anuttarapade iti vaktavyam . \\
\noindent \textbf{(P\_8,2.8)} iha mā bhūt . \\
\noindent \textbf{(P\_8,2.8)} carmaṇi tilā asya carmatilaḥ iti . \\
\noindent \textbf{(P\_8,2.8)} rājan vṛndāraka rājavṛndāraka iti . \\
\noindent \textbf{(P\_8,2.8)} vā napuṃsakānām . vā napuṃsakānām iti vaktavyam . \\
\noindent \textbf{(P\_8,2.8)} he carma he carman . \\
\noindent \textbf{(P\_8,2.8)} he varma he varman . \\
\noindent \textbf{(P\_8,2.8)} tat tarhi anuttarapade iti vaktavyam . \\
\noindent \textbf{(P\_8,2.8)} na vaktavyam . \\
\noindent \textbf{(P\_8,2.8)} na ṅisambuddhyoḥ iti ucyate na ca atra ṅisambuddhī paśyāmaḥ . \\
\noindent \textbf{(P\_8,2.8)} pratyayalakṣaṇena . \\
\noindent \textbf{(P\_8,2.8)} na lumatā tasmin iti pratyayalakṣaṇasya pratiṣedhaḥ . \\
\noindent \textbf{(P\_8,2.8)} na kvacit ṅiḥ lopena lupyate sarvatra lumatā eva . \\
\noindent \textbf{(P\_8,2.8)} yathā eva iha bhavati ārdre carman lohite carman iti evam iha api syāt carmaṇi tilā asya carmatilaḥ iti . \\
\noindent \textbf{(P\_8,2.8)} tasmāt upasaṅkhyānam kartavyam . \\
\noindent \textbf{(P\_8,2.8)} evam tarhi ṅyarthena tāvat na arthaḥ . \\
\noindent \textbf{(P\_8,2.8)} bhatvāt tu ṅau pratiṣedhānarthakyam . ṅau pratiṣedhaḥ anarthakaḥ . \\
\noindent \textbf{(P\_8,2.8)} kim kāraṇam . \\
\noindent \textbf{(P\_8,2.8)} bhatvāt . \\
\noindent \textbf{(P\_8,2.8)} bhasañjñā atra bhaviṣyati . \\
\noindent \textbf{(P\_8,2.8)} yadi tarhi bhasañjñā atra bhavati rathantare sāman iti atra allopaḥ anaḥ iti allopaḥ prāpnoti . \\
\noindent \textbf{(P\_8,2.8)} na eṣaḥ doṣaḥ . \\
\noindent \textbf{(P\_8,2.8)} uktam ubhayasañjñāni api chandāṃsi dṛśyante tad yathā saḥ suṣṭubhā saḥ ṛkvatā gaṇena padatvāt kutvam bhatvāt jaśtvam na bhavati . \\
\noindent \textbf{(P\_8,2.8)} evam iha api padatvāt allopaḥ na bhatvāt nalopaḥ na bhaviṣyati . \\
\noindent \textbf{(P\_8,2.8)} tasmāt na arthaḥ ṅigrahaṇena . \\
\noindent \textbf{(P\_8,2.8)} sambuddhyarthena ca api na arthaḥ . \\
\noindent \textbf{(P\_8,2.8)} katham . \\
\noindent \textbf{(P\_8,2.8)} sambuddhyantānām asamāsaḥ . \\
\noindent \textbf{(P\_8,2.8)} rājavṛndāraka iti . \\
\noindent \textbf{(P\_8,2.8)} kim vaktavyam etat . \\
\noindent \textbf{(P\_8,2.8)} na hi . \\
\noindent \textbf{(P\_8,2.8)} katham anucyamānam gaṃsyate . \\
\noindent \textbf{(P\_8,2.8)} iha samānārthena vākyena bhavitavyam samāsena ca yaḥ ca iha arthaḥ vākyena gamyate na asau jātu cit samāsena gamyate . \\
\noindent \textbf{(P\_8,2.8)} avayavasambodhanam vākyena gamyate samudāyasambodhanam samāsena . \\
\noindent \textbf{(P\_8,2.8)} vā napuṃsakānām iti etat vaktavyam eva \\
\noindent \textbf{(P\_8,2.9,} 42) KA\_III,395.19-24 Ro\_V,378 {1/10}     anantyayoḥ api niṣṭhāmatupoḥ ādeśaḥ . \\
\noindent \textbf{(P\_8,2.9,} 42) KA\_III,395.19-24 Ro\_V,378 {2/10}     niṣṭhāmatupoḥ ādeśaḥ anantyayoḥ api iti vaktavyam . \\
\noindent \textbf{(P\_8,2.9,} 42) KA\_III,395.19-24 Ro\_V,378 {3/10}     bhinnavantau bhinnavantaḥ . \\
\noindent \textbf{(P\_8,2.9,} 42) KA\_III,395.19-24 Ro\_V,378 {4/10}     vṛkṣavantau vṛkṣavantaḥ . \\
\noindent \textbf{(P\_8,2.9,} 42) KA\_III,395.19-24 Ro\_V,378 {5/10}     na vaktavyam . \\
\noindent \textbf{(P\_8,2.9,} 42) KA\_III,395.19-24 Ro\_V,378 {6/10}     vacanāt bhaviṣyati . \\
\noindent \textbf{(P\_8,2.9,} 42) KA\_III,395.19-24 Ro\_V,378 {7/10}     asti vacane prayojanam . \\
\noindent \textbf{(P\_8,2.9,} 42) KA\_III,395.19-24 Ro\_V,378 {8/10}     kim . \\
\noindent \textbf{(P\_8,2.9,} 42) KA\_III,395.19-24 Ro\_V,378 {9/10}     bhinnavān chinnavān . \\
\noindent \textbf{(P\_8,2.9,} 42) KA\_III,395.19-24 Ro\_V,378 {10/10}     vṛkṣavān plakṣavān \\
\noindent \textbf{(P\_8,2.9)} nārmate pratiṣedhaḥ . \\
\noindent \textbf{(P\_8,2.9)} nārmate pratiṣedhaḥ vaktavyaḥ . \\
\noindent \textbf{(P\_8,2.9)} nṛmataḥ nārmataḥ iti . \\
\noindent \textbf{(P\_8,2.9)} uktam vā . kim uktam . \\
\noindent \textbf{(P\_8,2.9)} niṣṭhāmatupoḥ tāvat uktam na vā padādhikārasya viśeṣaṇatvāt iti . \\
\noindent \textbf{(P\_8,2.9)} nārmate api uktam na vā bahiraṅgalakṣaṇatvāt iti \\
\noindent \textbf{(P\_8,2.11-12)} KA\_III,396.6-13 Ro\_V,379 {1/11}     kim ayam ekayogaḥ āhosvit nānāyogau . \\
\noindent \textbf{(P\_8,2.11-12)} KA\_III,396.6-13 Ro\_V,379 {2/11}     kim ca ataḥ . \\
\noindent \textbf{(P\_8,2.11-12)} KA\_III,396.6-13 Ro\_V,379 {3/11}     yadi ekayogaḥ hīvatī kapīvatī atra na prāpnoti . \\
\noindent \textbf{(P\_8,2.11-12)} KA\_III,396.6-13 Ro\_V,379 {4/11}     atha nānāyogau ikṣumatī drumatī atra api prāpnoti . \\
\noindent \textbf{(P\_8,2.11-12)} KA\_III,396.6-13 Ro\_V,379 {5/11}     yathā icchasi tathā astu . \\
\noindent \textbf{(P\_8,2.11-12)} KA\_III,396.6-13 Ro\_V,379 {6/11}     astu tāvat ekayogaḥ . \\
\noindent \textbf{(P\_8,2.11-12)} KA\_III,396.6-13 Ro\_V,379 {7/11}     katham ahīvatī kapīvatī . \\
\noindent \textbf{(P\_8,2.11-12)} KA\_III,396.6-13 Ro\_V,379 {8/11}     ācāryapravṛttiḥ jñāpayati bhavati evañjātīyakānām vatvam iti yat ayam anto'vatyāḥ īvatyāḥ iti āha . \\
\noindent \textbf{(P\_8,2.11-12)} KA\_III,396.6-13 Ro\_V,379 {9/11}     atha vā punaḥ astu nānāyogau . \\
\noindent \textbf{(P\_8,2.11-12)} KA\_III,396.6-13 Ro\_V,379 {10/11}     nanu ca uktam ikṣumatī drumatī atra api prāpnoti iti . \\
\noindent \textbf{(P\_8,2.11-12)} KA\_III,396.6-13 Ro\_V,379 {11/11}     yavādiṣu pāthaḥ kariṣyate \\
\noindent \textbf{(P\_8,2.15)} chandasi iraḥ iti ucyate tatra te viśvakarmāṇam te saptarṣimantam iti atra api prāpnoti . \\
\noindent \textbf{(P\_8,2.15)} na eṣaḥ doṣaḥ . \\
\noindent \textbf{(P\_8,2.15)} na evam vijñāyate chandasi iraḥ iti . \\
\noindent \textbf{(P\_8,2.15)} katham tarhi . \\
\noindent \textbf{(P\_8,2.15)} chandasi īraḥ iti . \\
\noindent \textbf{(P\_8,2.15)} evam api tviṣīmān patīmān iti atra api prāpnoti . \\
\noindent \textbf{(P\_8,2.15)} na eṣaḥ doṣaḥ . \\
\noindent \textbf{(P\_8,2.15)} vihitaviśeṣaṇam īkāragrahaṅam . \\
\noindent \textbf{(P\_8,2.15)} īkārāntāt yaḥ vihitaḥ iti . \\
\noindent \textbf{(P\_8,2.15)} evam api sūram te dyāvāpṛthivīmantam iti atra api prāpnoti . \\
\noindent \textbf{(P\_8,2.15)} iha ca na prāpnoti trivatīḥ yājyānuvākyāḥ bhavanti iti . \\
\noindent \textbf{(P\_8,2.15)} evam tarhi parigaṇanam kartavyam . \\
\noindent \textbf{(P\_8,2.15)} triharyadhipatyagnire . \\
\noindent \textbf{(P\_8,2.15)} trivatīḥ yājyānuvākyāḥ bhavanti . \\
\noindent \textbf{(P\_8,2.15)} tri . \\
\noindent \textbf{(P\_8,2.15)} hari . \\
\noindent \textbf{(P\_8,2.15)} harivaḥ medinam tvā . \\
\noindent \textbf{(P\_8,2.15)} hari . \\
\noindent \textbf{(P\_8,2.15)} adhipati . \\
\noindent \textbf{(P\_8,2.15)} adhipativatīḥ juhoti . \\
\noindent \textbf{(P\_8,2.15)} adhipati . \\
\noindent \textbf{(P\_8,2.15)} agni . \\
\noindent \textbf{(P\_8,2.15)} caruḥ agnivān iva . \\
\noindent \textbf{(P\_8,2.15)} agni . \\
\noindent \textbf{(P\_8,2.15)} re . \\
\noindent \textbf{(P\_8,2.15)} ā revan etu no viśa iti . \\
\noindent \textbf{(P\_8,2.15)} yadi tarhi parigaṇanam kriyate sarasvatīvān bhāratīvān apūpavān dadhivān caruḥ iti atra na prāpnoti . \\
\noindent \textbf{(P\_8,2.15)} evam tarhi chandasi iraḥ bahulam iti vaktavyam \\
\noindent \textbf{(P\_8,2.16)} yadi punaḥ ayam nuṭ pūrvāntaḥ kriyeta . \\
\noindent \textbf{(P\_8,2.16)} anaḥ nuki vināmaruvidhipratiṣedhaḥ . \\
\noindent \textbf{(P\_8,2.16)} anaḥ nuki sati vināmaḥ vidheyaḥ . \\
\noindent \textbf{(P\_8,2.16)} akṣaṇvān . \\
\noindent \textbf{(P\_8,2.16)} padāntasya na iti pratiṣedhaḥ prāpnoti . \\
\noindent \textbf{(P\_8,2.16)} ruḥ ca pratiṣedhyaḥ . \\
\noindent \textbf{(P\_8,2.16)} supathintaraḥ . \\
\noindent \textbf{(P\_8,2.16)} naśchavyapraśān iti ruḥ prāpnoti . \\
\noindent \textbf{(P\_8,2.16)} astu tarhi parādiḥ . \\
\noindent \textbf{(P\_8,2.16)} parādau vatvapratiṣedhaḥ avagrahaḥ ca . yadi parādiḥ vatvasya pratiṣedhaḥ vaktavyaḥ . \\
\noindent \textbf{(P\_8,2.16)} akṣaṇvān . \\
\noindent \textbf{(P\_8,2.16)} mādupadhāyāścamatorvo'yavādibhyaḥ iti vatvam prāpnoti . \\
\noindent \textbf{(P\_8,2.16)} avagrahaḥ ca aniṣṭe deśe prāpnoti . \\
\noindent \textbf{(P\_8,2.16)} akṣaṇvān . \\
\noindent \textbf{(P\_8,2.16)} astu tarhi pūrvāntaḥ . \\
\noindent \textbf{(P\_8,2.16)} nanu ca uktam anaḥ nuki vināmaruvidhipratiṣedhaḥ iti . \\
\noindent \textbf{(P\_8,2.16)} bhatvāt siddham . bhasañjñā vaktavyā . \\
\noindent \textbf{(P\_8,2.16)} yadi tarhi bhasañjñā allopo'naḥ iti allopaḥ prāpnoti . \\
\noindent \textbf{(P\_8,2.16)} anaḥ tu prakṛtibhāve matubgrahaṇam chandasi . anaḥ tu prakṛtibhāve matubgrahaṇam chandasi vaktavyam . \\
\noindent \textbf{(P\_8,2.16)} iha tarhi supathintaraḥ nāntasya ṭiḥ taddhite lupyate iti lopaḥ prāpnoti . \\
\noindent \textbf{(P\_8,2.16)} ghagrahaṇam ca . ghagrahaṇam ca kartavyam . \\
\noindent \textbf{(P\_8,2.16)} tat tarhi idam bahu vaktavyam . \\
\noindent \textbf{(P\_8,2.16)} nuk vaktavyaḥ . \\
\noindent \textbf{(P\_8,2.16)} bhasañjñā ca vaktavyā . \\
\noindent \textbf{(P\_8,2.16)} anaḥ tu prakṛtibhāve matubgrahaṇam chandasi vaktavyam . \\
\noindent \textbf{(P\_8,2.16)} ghagrahaṇam ca kartavyam iti . \\
\noindent \textbf{(P\_8,2.16)} na kartavyam . \\
\noindent \textbf{(P\_8,2.16)} yat tāvat ucyate nuk vaktavyaḥ iti nukaḥ eṣaḥ parihāraḥ bhatvāt siddham iti . \\
\noindent \textbf{(P\_8,2.16)} bhasañjñā vaktavyā iti kriyate nyāse eva ayasmayādīni chandasi iti . \\
\noindent \textbf{(P\_8,2.16)} yat api ucyate anaḥ tu prakṛtibhāve matubgrahaṇam chandasi ghagrahaṇam ca kartavyam iti na kartavyam . \\
\noindent \textbf{(P\_8,2.16)} ubhayasañjñāni api hi chandāṃsi dṛśyante . \\
\noindent \textbf{(P\_8,2.16)} tat yathā . \\
\noindent \textbf{(P\_8,2.16)} saḥ suṣṭubhā sa ṛkvatā gaṇena . \\
\noindent \textbf{(P\_8,2.16)} padatvāt kutvam bhatvāt jaśtvam na bhavati . \\
\noindent \textbf{(P\_8,2.16)} evam iha api padatvāt allopaṭilopau na bhatvāt vināmaruvidhipratiṣedhau bhaviṣyataḥ . \\
\noindent \textbf{(P\_8,2.16)} sidhyati . \\
\noindent \textbf{(P\_8,2.16)} sūtram tarhi bhidyate . \\
\noindent \textbf{(P\_8,2.16)} yathānyāsam eva astu . \\
\noindent \textbf{(P\_8,2.16)} nanu ca uktam parādau vatvapratiṣedhaḥ avagrahaḥ ca iti . \\
\noindent \textbf{(P\_8,2.16)} yat tāvat ucyate vatvapratiṣedhaḥ iti nirdiśyamānasya ādeśāḥ bhavanti iti evam na bhaviṣyati . \\
\noindent \textbf{(P\_8,2.16)} yaḥ tarhi nirdiśyate tasya na prāpnoti . \\
\noindent \textbf{(P\_8,2.16)} kim kāraṇam . \\
\noindent \textbf{(P\_8,2.16)} nuṭā vyavahitatvāt . \\
\noindent \textbf{(P\_8,2.16)} asiddhaḥ nuṭ tasya asiddhatvāt bhaviṣyati . \\
\noindent \textbf{(P\_8,2.16)} avagrahe api na lakṣaṇena padakārāḥ anuvartyāḥ padakāraiḥ nāma lakṣaṇam anuvartyam . \\
\noindent \textbf{(P\_8,2.16)} yathālakṣaṇam padam kartavyam \\
\noindent \textbf{(P\_8,2.17)} īt rathinaḥ . \\
\noindent \textbf{(P\_8,2.17)} rathinaḥ īt vaktavyaḥ . \\
\noindent \textbf{(P\_8,2.17)} rathītaraḥ . \\
\noindent \textbf{(P\_8,2.17)} bhūridāvnaḥ tuṭ . bhuridāvnaḥ tuṭ vaktavyaḥ . \\
\noindent \textbf{(P\_8,2.17)} bhūridāvattaraḥ janaḥ \\
\noindent \textbf{(P\_8,2.18)} kṛpaṇādīnām pratiṣedhaḥ vaktavyaḥ . \\
\noindent \textbf{(P\_8,2.18)} kṛpaṇaḥ kṛpāṇaḥ kṛpīṭam . \\
\noindent \textbf{(P\_8,2.18)} vālamūlalaghvalamaḍgulīnām vā laḥ ram āpadyate iti vaktavyam . \\
\noindent \textbf{(P\_8,2.18)} aśvavālaḥ aśvavāraḥ . \\
\noindent \textbf{(P\_8,2.18)} mūladevaḥ mūradevaḥ . \\
\noindent \textbf{(P\_8,2.18)} varuṇasya laghusyadaḥ varuṇasya raghusyadaḥ . \\
\noindent \textbf{(P\_8,2.18)} alam bhaktāya aram bhaktāya . \\
\noindent \textbf{(P\_8,2.18)} subāhuḥ svaṅguliḥ subāhuḥ svaṅguriḥ . \\
\noindent \textbf{(P\_8,2.18)} sañjñāchandasoḥ vā kapilakādīnām iti vaktavyam . \\
\noindent \textbf{(P\_8,2.18)} kapirakaḥ kapilakaḥ . \\
\noindent \textbf{(P\_8,2.18)} tilvirīkaḥ tilvilīkaḥ . \\
\noindent \textbf{(P\_8,2.18)} romāṇi lomāni . \\
\noindent \textbf{(P\_8,2.18)} pāṃsuram pāṃsulam . \\
\noindent \textbf{(P\_8,2.18)} karma kalma . \\
\noindent \textbf{(P\_8,2.18)} śukraḥ śuklaḥ \\
\noindent \textbf{(P\_8,2.19)} kim idam ayatigrahaṇam rephaviśeṣaṇam : ayatiparasya rephasya laḥ bhavati saḥ cet upasargasya bhavati iti . \\
\noindent \textbf{(P\_8,2.19)} āhosvit upasargaviśeṣaṇam : ayatiparasya upasargasya yaḥ rephaḥ tasya laḥ bhavati iti . \\
\noindent \textbf{(P\_8,2.19)} kaḥ ca atra viśeṣaḥ . \\
\noindent \textbf{(P\_8,2.19)} rephasya ayatau iti cet pareḥ upasaṅkhyānam . \\
\noindent \textbf{(P\_8,2.19)} rephasya ayatau iti cet pareḥ upasaṅkhyānam kartavyam . \\
\noindent \textbf{(P\_8,2.19)} palyayate . \\
\noindent \textbf{(P\_8,2.19)} vacanāt bhaviṣyati . \\
\noindent \textbf{(P\_8,2.19)} asti vacane prayojanam . \\
\noindent \textbf{(P\_8,2.19)} kim . \\
\noindent \textbf{(P\_8,2.19)} plāyate palāyate . \\
\noindent \textbf{(P\_8,2.19)} astu tarhi upasargaviśeṣaṇam . \\
\noindent \textbf{(P\_8,2.19)} upasargasya iti cet ekādeśe aprasiddhiḥ . upasargasya iti cet ekādeśe aprasiddhiḥ bhavati . \\
\noindent \textbf{(P\_8,2.19)} plāyate palāyate . \\
\noindent \textbf{(P\_8,2.19)} ekādeśe kṛte vyapavargābhāvāt na prāpnoti . \\
\noindent \textbf{(P\_8,2.19)} antādivat bhāvena vyapavargaḥ . \\
\noindent \textbf{(P\_8,2.19)} ubhayataḥ āśraye na antādivat . \\
\noindent \textbf{(P\_8,2.19)} evam tarhi ekādeśaḥ pūrvavidhau sthānivat bhavati iti sthānivadbhāvāt vyapavargaḥ . \\
\noindent \textbf{(P\_8,2.19)} pratiṣidhyate atra sthānivadbhāvaḥ pūrvatrāsiddhe na sthānivat iti . \\
\noindent \textbf{(P\_8,2.19)} doṣāḥ eva ete tasyāḥ paribhāṣāyāḥ tasya doṣaḥ saṃyogādilopalatvaṇatveṣu iti . \\
\noindent \textbf{(P\_8,2.19)} atha vā punaḥ astu rephaviśeṣaṇam . \\
\noindent \textbf{(P\_8,2.19)} nanu ca uktam rephasya ayatau iti cet pareḥ upasaṅkhyānam iti . \\
\noindent \textbf{(P\_8,2.19)} vacanāt bhaviṣyati . \\
\noindent \textbf{(P\_8,2.19)} nanu ca uktam asti vacane prayojanam kim plāyate palāyate iti . \\
\noindent \textbf{(P\_8,2.19)} atra api akāreṇa vyavahitatvāt na prāpnoti . \\
\noindent \textbf{(P\_8,2.19)} ekādeśe kṛte na asti vyavadhānam . \\
\noindent \textbf{(P\_8,2.19)} ekādeśaḥ pūrvavidhau sthānivat bhavati iti sthānivadbhāvāt vyavadhānam eva . \\
\noindent \textbf{(P\_8,2.19)} pratiṣidhyate atra sthānivadbhāvaḥ pūrvatrāsiddhe na sthānivat iti . \\
\noindent \textbf{(P\_8,2.19)} doṣāḥ eva ete tasyāḥ paribhāṣāyāḥ tasya doṣaḥ saṃyogādilopalatvaṇatveṣu iti \\
\noindent \textbf{(P\_8,2.21)} ṇau upasaṅkhyānam kartavyam . \\
\noindent \textbf{(P\_8,2.21)} iha api yathā syāt . \\
\noindent \textbf{(P\_8,2.21)} nigāryate nigālyate . \\
\noindent \textbf{(P\_8,2.21)} kim punaḥ kāraṇam na sidhyati . \\
\noindent \textbf{(P\_8,2.21)} aci iti ucyate na ca atra ajādim paśyāmaḥ . \\
\noindent \textbf{(P\_8,2.21)} pratyayalakṣaṇena . \\
\noindent \textbf{(P\_8,2.21)} varṇāśraye na asti pratyayalakṣaṇam . \\
\noindent \textbf{(P\_8,2.21)} evam tarhi sthānivadbhāvāt bhaviṣyati . \\
\noindent \textbf{(P\_8,2.21)} pratiṣidhyate atra sthānivadbhāvaḥ pūrvatrāsiddhe na sthānivat iti . \\
\noindent \textbf{(P\_8,2.21)} ataḥ uttaram paṭhati . \\
\noindent \textbf{(P\_8,2.21)} girateḥ latve ṇau uktam . \\
\noindent \textbf{(P\_8,2.21)} kim uktam . \\
\noindent \textbf{(P\_8,2.21)} tasya doṣaḥ saṃyogādilopalatvaṇatveṣu iti \\
\noindent \textbf{(P\_8,2.22.1)} yoge ca iti vaktavyam . \\
\noindent \textbf{(P\_8,2.22.1)} iha api yathā syāt . \\
\noindent \textbf{(P\_8,2.22.1)} pariyogaḥ paliyogaḥ \\
\noindent \textbf{(P\_8,2.22.2)} saṅi latvasalopasaṃyogādilopakutvadīrghatvāni . \\
\noindent \textbf{(P\_8,2.22.2)} saṅi iti prakṛtya latvasalopasaṃyogādilopakutvadīrghatvāni vaktavyāni . \\
\noindent \textbf{(P\_8,2.22.2)} kim prayojanam . \\
\noindent \textbf{(P\_8,2.22.2)} prayojanam girau giraḥ payaḥ dhāvati dviṣṭarām dṛṣatsthānam kāṣṭhaśaksthātā kruñcā dhuryaḥ iti . giṛau giraḥ iti atra acivibhāṣā iti latvam prāpnoti . \\
\noindent \textbf{(P\_8,2.22.2)} saṅi iti vacanāt na bhavati . \\
\noindent \textbf{(P\_8,2.22.2)} na etat asti prayojanam . \\
\noindent \textbf{(P\_8,2.22.2)} uktam etat dhātoḥ svarūpagrahaṇe tatpratyaye kāryavijñānāt siddham iti . \\
\noindent \textbf{(P\_8,2.22.2)} payaḥ dhāvati iti atra dhica iti salopaḥ prāpnoti saṅi iti vacanāt na bhavati . \\
\noindent \textbf{(P\_8,2.22.2)} etat api na asti prayojanam . \\
\noindent \textbf{(P\_8,2.22.2)} vakṣyati etat dhisakāre sicaḥ lopaḥ iti . \\
\noindent \textbf{(P\_8,2.22.2)} dviṣṭarām iti atra hrasvāt aṅgāt iti salopaḥ prāpnoti saṅi iti vacanāt na bhavati . \\
\noindent \textbf{(P\_8,2.22.2)} etat api na asti prayojanam . \\
\noindent \textbf{(P\_8,2.22.2)} atra api sicaḥ iti eva anuvartiṣyate . \\
\noindent \textbf{(P\_8,2.22.2)} dṛṣatsthānam iti atra jhalojhali iti salopaḥ prāpnoti saṅi iti vacanāt na bhavati . \\
\noindent \textbf{(P\_8,2.22.2)} etat api na asti prayojanam . \\
\noindent \textbf{(P\_8,2.22.2)} atra api sicaḥ iti eva anuvartiṣyate . \\
\noindent \textbf{(P\_8,2.22.2)} kāṣṭhaśaksthātā iti atra skoḥsaṃyogādyoranteca iti kakāralopaḥ prāpnoti saṅi iti vacanāt na bhavati . \\
\noindent \textbf{(P\_8,2.22.2)} etat api na asti prayojanam . \\
\noindent \textbf{(P\_8,2.22.2)} kāṣṭhaśak eva na asti kutaḥ yaḥ kāṣṭhaśaki tiṣṭhet . \\
\noindent \textbf{(P\_8,2.22.2)} kruñcā iti atra coḥkuḥ jhali iti kutvam prāpnoti saṅi iti vacanāt na bhavati . \\
\noindent \textbf{(P\_8,2.22.2)} etat api na asti prayojanam . \\
\noindent \textbf{(P\_8,2.22.2)} nipātanāt etat siddham . \\
\noindent \textbf{(P\_8,2.22.2)} kim nipātanam . \\
\noindent \textbf{(P\_8,2.22.2)} ṛtvigdadhṛksragdiguṣṇigañcuyujikruñcām iti . \\
\noindent \textbf{(P\_8,2.22.2)} dhuryaḥ iti atra halica iti dīrghatvam prāpnoti saṅi iti vacanāt na bhavati . \\
\noindent \textbf{(P\_8,2.22.2)} etat api na asti prayojanam . \\
\noindent \textbf{(P\_8,2.22.2)} nabhakurchurām iti pratiṣedhaḥ bhaviṣyati . \\
\noindent \textbf{(P\_8,2.23.1)} saṃyogāntasya lope yaṇaḥ pratiṣedhaḥ . saṃyogāntasya lope yaṇaḥ pratiṣedhaḥ vaktavyaḥ . \\
\noindent \textbf{(P\_8,2.23.1)} dadhi atra madhu atra iti . \\
\noindent \textbf{(P\_8,2.23.1)} saṃyogādilope ca yaṇaḥ pratiṣedhaḥ vaktavyaḥ . \\
\noindent \textbf{(P\_8,2.23.1)} kākī artham vāsī artham . \\
\noindent \textbf{(P\_8,2.23.1)} na vā jhalaḥ lopāt . \\
\noindent \textbf{(P\_8,2.23.1)} na vā vaktavyaḥ . \\
\noindent \textbf{(P\_8,2.23.1)} kim kāraṇam . \\
\noindent \textbf{(P\_8,2.23.1)} jhalaḥ lopāt . \\
\noindent \textbf{(P\_8,2.23.1)} jhalaḥ lopaḥ saṃyogāntalopaḥ vaktavyaḥ . \\
\noindent \textbf{(P\_8,2.23.1)} bahiraṅgalakṣaṇatvāt vā . atha vā bahiraṅgaḥ yaṇādeśaḥ antaraṅgaḥ lopaḥ asiddham bahiraṅgam antaraṅge \\
\noindent \textbf{(P\_8,2.23.2)} saṃyogāntalope sagrahaṇam . \\
\noindent \textbf{(P\_8,2.23.2)} saṃyogāntalope sagrahaṇam kartavyam . \\
\noindent \textbf{(P\_8,2.23.2)} saṃyogāntalopaḥ sasya ca iti vaktavyam . \\
\noindent \textbf{(P\_8,2.23.2)} iha api yathā syāt . \\
\noindent \textbf{(P\_8,2.23.2)} śreyān bhūyān jyāyān . \\
\noindent \textbf{(P\_8,2.23.2)} kim punaḥ kāraṇam na sidhyati . \\
\noindent \textbf{(P\_8,2.23.2)} paratvāt ruḥ prāpnoti . \\
\noindent \textbf{(P\_8,2.23.2)} asiddhaḥ ruḥ tasya asiddhatvāt lopaḥ bhaviṣyati . \\
\noindent \textbf{(P\_8,2.23.2)} na sidhyati . \\
\noindent \textbf{(P\_8,2.23.2)} kim kāraṇam . \\
\noindent \textbf{(P\_8,2.23.2)} ruvidhānasya anavakāśatvāt . anavakāsaḥ ruḥ lopam bādheta . \\
\noindent \textbf{(P\_8,2.23.2)} sāvakāśaḥ ruḥ . \\
\noindent \textbf{(P\_8,2.23.2)} kaḥ avakāśaḥ . \\
\noindent \textbf{(P\_8,2.23.2)} payaḥ śiraḥ . \\
\noindent \textbf{(P\_8,2.23.2)} nanu ca atra api jaśtvam prāpnoti . \\
\noindent \textbf{(P\_8,2.23.2)} saḥ yathā eva ruḥ jaśtvam bādhate evam lopam api bādheta . \\
\noindent \textbf{(P\_8,2.23.2)} na bādhate . \\
\noindent \textbf{(P\_8,2.23.2)} kim kāraṇam . \\
\noindent \textbf{(P\_8,2.23.2)} yena nāprāpte tasya bādhanam bhavati na ca aprāpte jaśtve ruḥ ārabhyate lope punaḥ prāpte ca aprāpte ca . \\
\noindent \textbf{(P\_8,2.23.2)} yogavibhāgāt siddham . atha vā yogavibhāgaḥ kariṣyate . \\
\noindent \textbf{(P\_8,2.23.2)} evam vakṣyāmi . \\
\noindent \textbf{(P\_8,2.23.2)} saṃyogāntasya lopaḥ arāt . \\
\noindent \textbf{(P\_8,2.23.2)} saṃyogāntasya lopaḥ bhavati arāt . \\
\noindent \textbf{(P\_8,2.23.2)} tataḥ sasya . \\
\noindent \textbf{(P\_8,2.23.2)} sasya ca lopaḥ bhavati saṃyogāntasya . \\
\noindent \textbf{(P\_8,2.23.2)} kim artham punaḥ idam ucyate . \\
\noindent \textbf{(P\_8,2.23.2)} pratiṣiddhārtham rubādhanārtham ca . \\
\noindent \textbf{(P\_8,2.23.2)} atha vā yat etat rāt sasya iti sagrahaṇam tat purastāt apakrakṣyate . \\
\noindent \textbf{(P\_8,2.23.2)} saṃyogāntasya lopaḥ . \\
\noindent \textbf{(P\_8,2.23.2)} tataḥ sasya . \\
\noindent \textbf{(P\_8,2.23.2)} sasya ca saṃyogāntasya lopaḥ bhavati . \\
\noindent \textbf{(P\_8,2.23.2)} tataḥ rāt . \\
\noindent \textbf{(P\_8,2.23.2)} rāt sasya eva saṃyogāntasya lopaḥ bhavati . \\
\noindent \textbf{(P\_8,2.23.2)} atha vā rāt sasya iti atra saṃyogāntasya lopaḥ iti etat anuvartiṣyate \\
\noindent \textbf{(P\_8,2.25)} dhi sakāre sicaḥ lopaḥ . \\
\noindent \textbf{(P\_8,2.25)} dhi sakāre sicaḥ lopaḥ vaktavyaḥ . \\
\noindent \textbf{(P\_8,2.25)} kim prayojanam . \\
\noindent \textbf{(P\_8,2.25)} cakāddhi iti prayojanam . iha mā bhūt . \\
\noindent \textbf{(P\_8,2.25)} cakāddhi palitam śiraḥ . \\
\noindent \textbf{(P\_8,2.25)} yadi tarhi sicaḥ lopaḥ iti ucyate . \\
\noindent \textbf{(P\_8,2.25)} āśādhvam tu katham te syāt . āśādhvam iti atra na prāpnoti . \\
\noindent \textbf{(P\_8,2.25)} jaśtvam sasya bhaviṣyati . jaśtvam atra sakārasya bhaviṣyati . \\
\noindent \textbf{(P\_8,2.25)} sarvatra evam prasiddham syāt . sarvatra evam jaśtvena siddham syāt . \\
\noindent \textbf{(P\_8,2.25)} iha api āyandhvam arandhvam iti jaśtvena eva siddham . \\
\noindent \textbf{(P\_8,2.25)} śrutiḥ ca api na bhidyate . śrutikṛtaḥ ca api na kaḥ cit bhedaḥ bhavati . \\
\noindent \textbf{(P\_8,2.25)} luṅaḥ ca api na mūrdhanye grahaṇam . tatra ayam api arthaḥ iṇaḥṣīdhvaṃluṅliṭāndho'ṅgāt iti atra luṅgrahaṇam na kartavyam . \\
\noindent \textbf{(P\_8,2.25)} iha api acyoḍḍhvam aploḍḍhvam iti ṣatve sicaḥ dhasya ṣṭutve ca kṛte jaśtvena siddham . \\
\noindent \textbf{(P\_8,2.25)} seṭi duṣyati . seṭi doṣaḥ bhavati . \\
\noindent \textbf{(P\_8,2.25)} idam eva rūpam syāt alaviḍḍhvam idam na syāt alavidhvam iti . \\
\noindent \textbf{(P\_8,2.25)} tasmāt sicaḥ grahaṇam kartavyam . \\
\noindent \textbf{(P\_8,2.25)} yadi tarhi sicaḥ grahaṇam kriyate . \\
\noindent \textbf{(P\_8,2.25)} ghasibhasyoḥ na sidhyet tu . ghasibhasyoḥ na sidhyati . \\
\noindent \textbf{(P\_8,2.25)} sagdhiḥ ca me sapītiḥ ca me, babdhām te harī dhānāḥ iti atra na prāpnoti . \\
\noindent \textbf{(P\_8,2.25)} tasmāt sijgrahaṇam na tat . tasmāt dhica iti atra sicaḥ grahaṇam na kartavyam . \\
\noindent \textbf{(P\_8,2.25)} katham cakāddhi palitam śiraḥ iti . \\
\noindent \textbf{(P\_8,2.25)} evam tarhi sijgrahaṇam kartavyam . \\
\noindent \textbf{(P\_8,2.25)} katham sagdhiḥ ca me sapītiḥ ca me , babdham te harī dhānāḥ iti . \\
\noindent \textbf{(P\_8,2.25)} iha tāvat sagdhiḥ iti na etat ghaseḥ rūpam . \\
\noindent \textbf{(P\_8,2.25)} kim tarhi sagheḥ etat rūpam . \\
\noindent \textbf{(P\_8,2.25)} babdhām te harī dhānāḥ iti na etat bhaseḥ rūpam . \\
\noindent \textbf{(P\_8,2.25)} kim tarhi bandheḥ etat rūpam . \\
\noindent \textbf{(P\_8,2.25)} chāndasaḥ varṇalopaḥ vā yathā iṣkartāramadhvare . atha vā chāndasaḥ varṇalopaḥ bhaviṣyati yathā iṣkartāramadhvare . \\
\noindent \textbf{(P\_8,2.25)} tat yathā . \\
\noindent \textbf{(P\_8,2.25)} tubhyedam agne . \\
\noindent \textbf{(P\_8,2.25)} tubhyam idam agne iti prāpte . \\
\noindent \textbf{(P\_8,2.25)} āmbānām caruḥ . \\
\noindent \textbf{(P\_8,2.25)} nāmbānām caruḥ iti prāpte . \\
\noindent \textbf{(P\_8,2.25)} āvyādhinīḥ ugaṇāḥ . \\
\noindent \textbf{(P\_8,2.25)} āvyādhinīḥ sugaṇāḥ iti prāpte . \\
\noindent \textbf{(P\_8,2.25)} iṣkartāram adhvarasya . \\
\noindent \textbf{(P\_8,2.25)} niṣkartāram iti prāpte . \\
\noindent \textbf{(P\_8,2.25)} śivā udrasya bheṣajī . \\
\noindent \textbf{(P\_8,2.25)} śivā rudrasya bheṣajīti prāpte . \\
\noindent \textbf{(P\_8,2.25)} tasmāt sijgrahaṇam kartavyam . \\
\noindent \textbf{(P\_8,2.25)} na kartavyam . \\
\noindent \textbf{(P\_8,2.25)} yat etat rātsasya iti sakāragrahaṇam tat sicaḥ grahaṇam vijñāsyate . \\
\noindent \textbf{(P\_8,2.25)} katham . \\
\noindent \textbf{(P\_8,2.25)} rātsasya iti ucyate na ca anyaḥ rephāt paraḥ sakāraḥ asti anyat ataḥ sicaḥ . \\
\noindent \textbf{(P\_8,2.25)} nanu ca ayam asti mātuḥ pituḥ iti . \\
\noindent \textbf{(P\_8,2.25)} tasmāt sicaḥ grahaṇam kartavyam . \\
\noindent \textbf{(P\_8,2.25)} na kartavyam . \\
\noindent \textbf{(P\_8,2.25)} kasmāt na bhavati cakāddhi palitam śiraḥ iti . \\
\noindent \textbf{(P\_8,2.25)} iṣṭam eva etat saṅgṛhītam . \\
\noindent \textbf{(P\_8,2.25)} cakādhi iti eva bhavitavyam . \\
\noindent \textbf{(P\_8,2.25)} dhi sakāre sicaḥ lopaḥ . \\
\noindent \textbf{(P\_8,2.25)} cakāddhi iti prayojanam . \\
\noindent \textbf{(P\_8,2.25)} āśādhvam tu katham te syāt . \\
\noindent \textbf{(P\_8,2.25)} jaśtvam sasya bhaviṣyati . \\
\noindent \textbf{(P\_8,2.25)} sarvatra evam prasiddham syāt . \\
\noindent \textbf{(P\_8,2.25)} śrutiḥ ca api na bhidyate . \\
\noindent \textbf{(P\_8,2.25)} luṅaḥ ca api na mūrdhanye grahaṇam . \\
\noindent \textbf{(P\_8,2.25)} seṭi duṣyati . \\
\noindent \textbf{(P\_8,2.25)} ghasibhasyoḥ na sidhyet tu . \\
\noindent \textbf{(P\_8,2.25)} tasmāt sijgrahaṇam na tat . \\
\noindent \textbf{(P\_8,2.25)} chāndasaḥ varṇalopaḥ vā yathā iṣkartāramadhvare . \\
\noindent \textbf{(P\_8,2.25)} dādeḥ dhātoḥ ghaḥ . \\
\noindent \textbf{(P\_8,2.32.1)} iha dogdhā dogdhum iti ghatvasya asiddhatvāt ḍhatvam prāpnoti . \\
\noindent \textbf{(P\_8,2.32.1)} na eṣaḥ doṣaḥ . \\
\noindent \textbf{(P\_8,2.32.1)} uktam etat apavādaḥ vacanaprāmāṇyāt iti . \\
\noindent \textbf{(P\_8,2.32.1)} atha vā evam vakṣyāmi . \\
\noindent \textbf{(P\_8,2.32.1)} haḥ ḍhaḥ adādeḥ . \\
\noindent \textbf{(P\_8,2.32.1)} haḥ ḍhaḥ bhavati adādeḥ . \\
\noindent \textbf{(P\_8,2.32.1)} tataḥ dhātoḥ ghaḥ iti . \\
\noindent \textbf{(P\_8,2.32.1)} dādeḥ iti anuvartate na iti nivṛttam \\
\noindent \textbf{(P\_8,2.32.2)} dādeḥ iti ucyate tatra idam na sidhyati . \\
\noindent \textbf{(P\_8,2.32.2)} adhok . \\
\noindent \textbf{(P\_8,2.32.2)} kva tarhi syāt . \\
\noindent \textbf{(P\_8,2.32.2)} mā sma dhok. na eṣaḥ doṣaḥ . \\
\noindent \textbf{(P\_8,2.32.2)} dhātoḥ iti na eṣā dādisamānādhikaraṇā ṣaṣṭhī . \\
\noindent \textbf{(P\_8,2.32.2)} dādeḥ dhātoḥ iti . \\
\noindent \textbf{(P\_8,2.32.2)} kā tarhi . \\
\noindent \textbf{(P\_8,2.32.2)} avayavayogā eṣā ṣaṣṭhī . \\
\noindent \textbf{(P\_8,2.32.2)} dhātoḥ yaḥ dādiḥ avayavaḥ iti . \\
\noindent \textbf{(P\_8,2.32.2)} sā ca avaśyam avayavayogā ṣaṣṭhī vijñeyā uttarārthā . \\
\noindent \textbf{(P\_8,2.32.2)} kim prayojanam . \\
\noindent \textbf{(P\_8,2.32.2)} ekācaḥ baśaḥ bhaṣ jhaṣantasya sdhvoḥ iti iha api yathā syāt : gardabhayateḥ apratyayaḥ gardhap iti . \\
\noindent \textbf{(P\_8,2.32.2)} yadi avayavayogā ṣaṣṭhī dogdhā dogdhum iti atra na prāpnoti . \\
\noindent \textbf{(P\_8,2.32.2)} eṣaḥ api vyapadeśivadbhāvena dhātoḥ dādiḥ avayavaḥ bhavati . \\
\noindent \textbf{(P\_8,2.32.3)} hṛgrahoḥ bhaḥ chandasi hasya . \\
\noindent \textbf{(P\_8,2.32.3)} hṛgrahoḥ chandasi hasya bhatvam vaktavyam . \\
\noindent \textbf{(P\_8,2.32.3)} gardabhena sambharati . \\
\noindent \textbf{(P\_8,2.32.3)} marut asya grabhītā . \\
\noindent \textbf{(P\_8,2.32.3)} sāmidhenyaḥ jabhrire . \\
\noindent \textbf{(P\_8,2.32.3)} udgrābham ca nigrābham ca brahma devāḥ avīvṛdhan \\
\noindent \textbf{(P\_8,2.38.1)} kimarthaḥ cakāraḥ . \\
\noindent \textbf{(P\_8,2.38.1)} sdhvoḥ iti etat anukṛṣyate . \\
\noindent \textbf{(P\_8,2.38.1)} na etat asti prayojanam . \\
\noindent \textbf{(P\_8,2.38.1)} siddham sdhvoḥ pūrveṇa eva . \\
\noindent \textbf{(P\_8,2.38.1)} na sidhyati . \\
\noindent \textbf{(P\_8,2.38.1)} kim kāraṇam . \\
\noindent \textbf{(P\_8,2.38.1)} abaśāditvāt . \\
\noindent \textbf{(P\_8,2.38.1)} nanu ca jaśtve kṛte baśādiḥ . \\
\noindent \textbf{(P\_8,2.38.1)} asiddham jaśtvam tasya asiddhatvāt na baśādiḥ . \\
\noindent \textbf{(P\_8,2.38.1)} evam tarhi siddhakāṇḍe paṭhitam abhyāsajaśtvacartvam ettvatukoḥ iti . \\
\noindent \textbf{(P\_8,2.38.1)} ettvatukoḥ grahaṇam na kariṣyate . \\
\noindent \textbf{(P\_8,2.38.1)} abhyāsajaśtvacartvam siddham iti eva . \\
\noindent \textbf{(P\_8,2.38.1)} evam api ajhaṣantatvāt na prāpnoti . \\
\noindent \textbf{(P\_8,2.38.1)} lope kṛte jhaṣantaḥ . \\
\noindent \textbf{(P\_8,2.38.1)} sthānivadbhāvāt na jhaṣantaḥ . \\
\noindent \textbf{(P\_8,2.38.1)} ataḥ uttaram paṭhati . \\
\noindent \textbf{(P\_8,2.38.1)} dadhaḥ tathoḥ anukarṣaṇānarthakyam sthānivatpratiṣedhāt . dadhaḥ tathoḥ anukarṣaṇam anarthakam . \\
\noindent \textbf{(P\_8,2.38.1)} kim kāraṇam . \\
\noindent \textbf{(P\_8,2.38.1)} sthānivatpratiṣedhāt . \\
\noindent \textbf{(P\_8,2.38.1)} pratiṣidhyate atra sthanivadbhāvaḥ pūrvatrāsiddhe na sthānivat iti . \\
\noindent \textbf{(P\_8,2.38.1)} sa ca avaśyam pratiṣedhaḥ āśrayitavyaḥ . \\
\noindent \textbf{(P\_8,2.38.1)} itarathā hi alope pratiṣedhaḥ . yaḥ hi manyate anukarṣaṇasāmarthyāt me atra bhavati alope tena pratiṣedhaḥ vaktavyaḥ syāt : dadhāti dadhāsi \\
\noindent \textbf{(P\_8,2.38.2)} tathoḥ ca api grahaṇam śakyam akartum . \\
\noindent \textbf{(P\_8,2.38.2)} katham . \\
\noindent \textbf{(P\_8,2.38.2)} jhali jhaṣantasya iti ucyate tathoḥ ca ayam jhali jhaṣantaḥ bhavati na anyatra \\
\noindent \textbf{(P\_8,2.38.3)} atha api etat na asti pūrvatrāsiddhe na sthānivat iti evam api na eva arthaḥ anukarṣaṇārthena cakāreṇa na api tathoḥ grahaṇena . \\
\noindent \textbf{(P\_8,2.38.3)} ānantaryam iha āśrīyate jhali jhaṣantasya iti . \\
\noindent \textbf{(P\_8,2.38.3)} kva cit ca sannipātakṛtam ānantaryam śāstrakṛtam anānantaryam kva cit na eva sannipātakṛtam na api śāstrakṛtam . \\
\noindent \textbf{(P\_8,2.38.3)} lope sannipātakṛtam ānantaryam śāstrakṛtam anānantaryam alope na eva sannipātakṛtam na api śāstrakṛtam . \\
\noindent \textbf{(P\_8,2.38.3)} yatra kutaḥ cit eva ānantaryam tat āśrayiṣyāmaḥ . \\
\noindent \textbf{(P\_8,2.40)} adhaḥ iti kimartham . \\
\noindent \textbf{(P\_8,2.40)} dhattaḥ . \\
\noindent \textbf{(P\_8,2.40)} dhatthaḥ . \\
\noindent \textbf{(P\_8,2.40)} adhaḥ iti śakyam akartum . \\
\noindent \textbf{(P\_8,2.40)} kasmāt na bhavati dhattaḥ dhatthaḥ iti . \\
\noindent \textbf{(P\_8,2.40)} jaśtve yogavibhāgaḥ kariṣyate . \\
\noindent \textbf{(P\_8,2.40)} idam asti dadhastathośca iti . \\
\noindent \textbf{(P\_8,2.40)} tataḥ vakṣyāmi jhalām jaśaḥ . \\
\noindent \textbf{(P\_8,2.40)} jhalām jaśaḥ bhavanti dadhaḥ tathoḥ . \\
\noindent \textbf{(P\_8,2.40)} tataḥ ante . \\
\noindent \textbf{(P\_8,2.40)} ante ca jhalām jaśaḥ bhavanti . \\
\noindent \textbf{(P\_8,2.40)} tatra jaśtve kṛte ajhaṣantatvāt na bhaviṣyati . \\
\noindent \textbf{(P\_8,2.42.1)} radābhyām iti kimartham . \\
\noindent \textbf{(P\_8,2.42.1)} caritam muditam . \\
\noindent \textbf{(P\_8,2.42.1)} nanu ca radābhyām iti ucyamāne api atra prāpnoti . \\
\noindent \textbf{(P\_8,2.42.1)} atra api rephadakārābhyām parā niṣṭhā . \\
\noindent \textbf{(P\_8,2.42.1)} na rephadakārābhyām niṣṭhā viśeṣyate . \\
\noindent \textbf{(P\_8,2.42.1)} kim tarhi . \\
\noindent \textbf{(P\_8,2.42.1)} takāraḥ viśeṣyate . \\
\noindent \textbf{(P\_8,2.42.1)} raphadakārābhyām uttarasya takārasya naḥ bhavati sa cet niṣṭhāyāḥ iti . \\
\noindent \textbf{(P\_8,2.42.1)} atha pūrvagrahaṇam kimartham . \\
\noindent \textbf{(P\_8,2.42.1)} niṣṭhādeśe pūrvagrahaṇam parasya ādeśapratiṣedhārtham . \\
\noindent \textbf{(P\_8,2.42.1)} niṣṭhādeśe pūrvagrahaṇam kriyate parasya ādeśaḥ mā bhūt iti . \\
\noindent \textbf{(P\_8,2.42.1)} bhinnavadbhyām bhinnavadbhiḥ . \\
\noindent \textbf{(P\_8,2.42.1)} pañcamīnirdiṣṭāt hi parasya . pañcamīnirdiṣṭāt hi parasya iti parasya prāpnoti \\
\noindent \textbf{(P\_8,2.42.2)} vṛddhinimittāt pratiṣedhaḥ . \\
\noindent \textbf{(P\_8,2.42.2)} vṛddhinimittāt pratiṣedhaḥ vaktavyaḥ . \\
\noindent \textbf{(P\_8,2.42.2)} kim prayojanam . \\
\noindent \textbf{(P\_8,2.42.2)} prayojanam kārtikṣaitiphaullayaḥ . kārtiḥ iti vṛddhau kṛtāyām radābhyām iti natvam prāpnoti . \\
\noindent \textbf{(P\_8,2.42.2)} kṣaitiḥ iti vṛddhau kṛtāyām kṣiyodīrghāt iti natvam prāpnoti . \\
\noindent \textbf{(P\_8,2.42.2)} phaulliḥ iti vṛddhau kṛtāyām udupadhatvasanniyogena latvam ucyamānam na prāpnoti . \\
\noindent \textbf{(P\_8,2.42.2)} atha ucyamāne api pratiṣedhe vṛddhinimittāt iti katham idam vijñāyate . \\
\noindent \textbf{(P\_8,2.42.2)} vṛddhiḥ eva nimittam vṛddhinimittam vṛddhinimittāt iti . \\
\noindent \textbf{(P\_8,2.42.2)} āhosvit vṛddhiḥ nimittam asya saḥ ayam vṛddhinimittaḥ vṛddhinimittāt iti . \\
\noindent \textbf{(P\_8,2.42.2)} kim ca ataḥ . \\
\noindent \textbf{(P\_8,2.42.2)} yadi vijñāyate vṛddhiḥ eva nimittam vṛddhinimittam vṛddhinimittāt iti kṣaitiḥ saṅgṛhītaḥ kārtiḥ asaṅgṛhītaḥ . \\
\noindent \textbf{(P\_8,2.42.2)} atha vijñāyate vṛddhiḥ nimittam asya saḥ ayam vṛddhinimittaḥ vṛddhinimittāt iti kārtiḥ saṅgṛhītaḥ kṣaitiḥ asaṅgṛhītaḥ . \\
\noindent \textbf{(P\_8,2.42.2)} ubhayathā ca phaulliḥ asaṅgṛhītaḥ . \\
\noindent \textbf{(P\_8,2.42.2)} yathā icchasi tathā astu . \\
\noindent \textbf{(P\_8,2.42.2)} astu tāvat vṛddhiḥ eva nimittam vṛddhinimittam vṛddhinimittāt iti . \\
\noindent \textbf{(P\_8,2.42.2)} nanu ca uktam kṣaitiḥ saṅgrhītaḥ kārtiḥ asaṅgṛhītaḥ iti . \\
\noindent \textbf{(P\_8,2.42.2)} kārtiḥ ca saṅgṛhītaḥ . \\
\noindent \textbf{(P\_8,2.42.2)} katham . \\
\noindent \textbf{(P\_8,2.42.2)} vṛddhiḥ bhavati guṇaḥ bhavati iti rephaśirāḥ guṇavṛddhisañjñakaḥ abhinirvartate . \\
\noindent \textbf{(P\_8,2.42.2)} atha vā punaḥ astu vṛddhiḥ nimittam asya saḥ ayam vṛddhinimittaḥ vṛddhinimittāt iti . \\
\noindent \textbf{(P\_8,2.42.2)} nanu ca uktam kārtiḥ saṅgṛhītaḥ kṣaitiḥ asaṅgṛhītaḥ iti . \\
\noindent \textbf{(P\_8,2.42.2)} kṣaitiḥ ca saṅgṛhītaḥ . \\
\noindent \textbf{(P\_8,2.42.2)} katham . \\
\noindent \textbf{(P\_8,2.42.2)} yat tat vṛddhiśāstram tasmin vṛddhiśabdaḥ vartate . \\
\noindent \textbf{(P\_8,2.42.2)} saḥ tarhi pratiṣedhaḥ vaktavyaḥ . \\
\noindent \textbf{(P\_8,2.42.2)} na vā bahiraṅgalakṣaṇatvāt . na vā vaktavyam . \\
\noindent \textbf{(P\_8,2.42.2)} kim kāraṇam . \\
\noindent \textbf{(P\_8,2.42.2)} bahiraṅgalakṣaṇatvāt . \\
\noindent \textbf{(P\_8,2.42.2)} bahiraṅgā vṛddhiḥ . \\
\noindent \textbf{(P\_8,2.42.2)} antaraṅgam natvam . \\
\noindent \textbf{(P\_8,2.42.2)} asiddham bahiraṅgam antaraṅge . \\
\noindent \textbf{(P\_8,2.42.2)} evam ca kṛtvā latvam api siddham bhavati phaulliḥ iti \\
\noindent \textbf{(P\_8,2.44)} ṛkāralvādibhyaḥ ktinniṣṭhāvat . \\
\noindent \textbf{(P\_8,2.44)} ṛkāralvādibhyaḥ ktin niṣṭhāvat bhavati iti vaktavyam . \\
\noindent \textbf{(P\_8,2.44)} kīrṇiḥ gīrṇiḥ . \\
\noindent \textbf{(P\_8,2.44)} lūniḥ dhūniḥ . \\
\noindent \textbf{(P\_8,2.44)} dugvoḥ dīrghaḥ ca . dugvoḥ dīrghaḥ ca iti vaktavyam . \\
\noindent \textbf{(P\_8,2.44)} ādūnaḥ vigūnaḥ . \\
\noindent \textbf{(P\_8,2.44)} pūñaḥ vināśe . pūñaḥ vināśe iti vaktavyam . \\
\noindent \textbf{(P\_8,2.44)} pūnāḥ yavāḥ . \\
\noindent \textbf{(P\_8,2.44)} vināśe iti kimartham . \\
\noindent \textbf{(P\_8,2.44)} pūtam dhānyam . \\
\noindent \textbf{(P\_8,2.44)} sinoteḥ grāsakarmakartṛkasya . sinoteḥ grāsakarmakartṛkasya iti vaktavyam . \\
\noindent \textbf{(P\_8,2.44)} sinaḥ grāsaḥ . \\
\noindent \textbf{(P\_8,2.44)} grāsakarmakartṛkasya iti kimartham . \\
\noindent \textbf{(P\_8,2.44)} sitā pāśena sūkarī \\
\noindent \textbf{(P\_8,2.46)} dīrghāt iti kimartham . \\
\noindent \textbf{(P\_8,2.46)} akṣitam asi mā me kṣeṣṭhāḥ . \\
\noindent \textbf{(P\_8,2.46)} dīrghāt iti śakyam akartum . \\
\noindent \textbf{(P\_8,2.46)} kasmāt na bhavati akṣitam asi mā me kṣeṣṭhāḥ iti . \\
\noindent \textbf{(P\_8,2.46)} nirdeśāt eva idam abhivyaktam dīrghasya grahaṇam iti . \\
\noindent \textbf{(P\_8,2.46)} yadi hrasvasya grahaṇam syāt kṣeḥ iti eva brūyāt . \\
\noindent \textbf{(P\_8,2.46)} na atra nirdeśaḥ pramāṇaṃ śakyam kartum . \\
\noindent \textbf{(P\_8,2.46)} yathā eva atra aprāptā vibhaktiḥ evam iyaṅādeśaḥ api . \\
\noindent \textbf{(P\_8,2.46)} na atra aprāptā vibhaktiḥ . \\
\noindent \textbf{(P\_8,2.46)} siddhā atra vibhaktiḥ prātipadikāt iti . \\
\noindent \textbf{(P\_8,2.46)} katham prātipadikasañjñā . \\
\noindent \textbf{(P\_8,2.46)} arthavat prātipadikam iti . \\
\noindent \textbf{(P\_8,2.46)} nanu ca adhātuḥ iti pratiṣedhaḥ prāpnoti . \\
\noindent \textbf{(P\_8,2.46)} na eṣaḥ dhātuḥ dhātoḥ eṣaḥ anukaraṇaḥ . \\
\noindent \textbf{(P\_8,2.46)} yadi anukaraṇaḥ iyaṅādeśaḥ na prāpnoti . \\
\noindent \textbf{(P\_8,2.46)} prakṛtivat anukaraṇam bhavati iti evam iyaṅādeśaḥ bhaviṣyati . \\
\noindent \textbf{(P\_8,2.46)} yadi prakṛtivat anukaraṇam bhavati iti ucyate svādyutpattiḥ na prāpnoti . \\
\noindent \textbf{(P\_8,2.46)} evam tarhi ātideśikānām svāśrayāṇi api na nivartante . \\
\noindent \textbf{(P\_8,2.46)} atha api etat na asti ātideśikānām svāśrayāṇi api na nivartante iti evam api na doṣaḥ . \\
\noindent \textbf{(P\_8,2.46)} avaśyam atra sarvataḥ nairdeśikī vibhaktiḥ vaktavyā . \\
\noindent \textbf{(P\_8,2.46)} tat yathā . \\
\noindent \textbf{(P\_8,2.46)} nerviśaḥ parivyavebhyaḥkriyaḥ viparābhyāñjeḥ iti . \\
\noindent \textbf{(P\_8,2.46)} atha api etat na asti prakṛtivat anukaraṇam bhavati iti evam api na doṣaḥ . \\
\noindent \textbf{(P\_8,2.46)} dhātoḥ ajādau yat rūpam tat anukriyate \\
\noindent \textbf{(P\_8,2.48)} añceḥ natve vyaktapratiṣedhaḥ . \\
\noindent \textbf{(P\_8,2.48)} añceḥ natve vyaktasya pratiṣedhaḥ vaktavyaḥ . \\
\noindent \textbf{(P\_8,2.48)} vyaktam anṛtam kathayati iti . \\
\noindent \textbf{(P\_8,2.48)} añjivijñānāt siddham . na etat añceḥ rūpam . \\
\noindent \textbf{(P\_8,2.48)} añjeḥ etat rūpam . \\
\noindent \textbf{(P\_8,2.48)} añcatyarthaḥ vai gamyate . \\
\noindent \textbf{(P\_8,2.48)} kaḥ punaḥ añcatyarthaḥ . \\
\noindent \textbf{(P\_8,2.48)} añcatiḥ prakāśane vartate . \\
\noindent \textbf{(P\_8,2.48)} añcitam gacchati . \\
\noindent \textbf{(P\_8,2.48)} prakāśayati ātmānam iti gamyate . \\
\noindent \textbf{(P\_8,2.48)} na vai loke añcitam gacchati iti prakāśanam gamyate . \\
\noindent \textbf{(P\_8,2.48)} kim tarhi . \\
\noindent \textbf{(P\_8,2.48)} samādhānam gamyate . \\
\noindent \textbf{(P\_8,2.48)} samāhitaḥ bhūtvā gacchati iti . \\
\noindent \textbf{(P\_8,2.48)} evam tarhi añcateḥ aṅkaḥ aṅkaḥ ca prakāśanam . \\
\noindent \textbf{(P\_8,2.48)} aṅkitāḥ gāvaḥ iti ucyate anyābhyaḥ gobhyaḥ prakāśyante . \\
\noindent \textbf{(P\_8,2.48)} añcatyarthaḥ iti cet añjeḥ tadarthatvāt siddham . añcatyarthaḥ iti cet añjiḥ api añcatyarthe vartate . \\
\noindent \textbf{(P\_8,2.48)} katham punaḥ anyaḥ nāma anyasya arthe vartate . \\
\noindent \textbf{(P\_8,2.48)} katham añjiḥ añcatyarthe vartate . \\
\noindent \textbf{(P\_8,2.48)} anekārthāḥ api dhātavaḥ bhavanti . \\
\noindent \textbf{(P\_8,2.48)} asti punaḥ kva cit anyatra api añjiḥ añcatyarthe vartate . \\
\noindent \textbf{(P\_8,2.48)} asti iti āha . \\
\noindent \textbf{(P\_8,2.48)} añjeḥ añjanam añjanam ca prakāśanam . \\
\noindent \textbf{(P\_8,2.48)} aṅkteṣiṇī iti ucyate yat tat sitam ca asitam ca etat prakāśayati . \\
\noindent \textbf{(P\_8,2.48)} tathā añjeḥ vyañjanam vyañjanam ca prakāśanam . \\
\noindent \textbf{(P\_8,2.48)} yat tat snehena madhureṇa ca jaḍīkṛtānām indriyāṇām svasmin ātmani vyavasthāpanam saḥ rāgaḥ tat vyañjanam . \\
\noindent \textbf{(P\_8,2.48)} anvartham khalu api nirvacanam . \\
\noindent \textbf{(P\_8,2.48)} vyajyate anena iti vyañjanam iti \\
\noindent \textbf{(P\_8,2.50)} avātābhidhāne . \\
\noindent \textbf{(P\_8,2.50)} avātābhidhāne iti vaktavyam . \\
\noindent \textbf{(P\_8,2.50)} iha api yathā syāt . \\
\noindent \textbf{(P\_8,2.50)} nirvāṇaḥ agniḥ vātena . \\
\noindent \textbf{(P\_8,2.50)} nirvāṇaḥ pradīpaḥ vātena iti \\
\noindent \textbf{(P\_8,2.55.1)} anupasargāt iti ucyate tatra idam na sidhyati parikṛśam iti . \\
\noindent \textbf{(P\_8,2.55.1)} kṛśeḥ kaḥ eṣaḥ vihitaḥ igupadhāt . \\
\noindent \textbf{(P\_8,2.55.1)} na etat niṣṭhāntam . \\
\noindent \textbf{(P\_8,2.55.1)} kim tarhi kṛśaḥ eṣaḥ igupadhāt kaḥ vihitaḥ . \\
\noindent \textbf{(P\_8,2.55.1)} svare hi doṣaḥ bhavati parikṛśe . na evam śakyam . \\
\noindent \textbf{(P\_8,2.55.1)} iha hi parikṛśaḥ iti svare doṣaḥ syāt . \\
\noindent \textbf{(P\_8,2.55.1)} antasthāthaghañktājabitrakāṇām . \\
\noindent \textbf{(P\_8,2.55.1)} iti eṣaḥ svaraḥ prasajyeta . \\
\noindent \textbf{(P\_8,2.55.1)} padasya lopaḥ vihitaḥ iti matam . evam tarhi padasya lopaḥ draṣṭavyaḥ . \\
\noindent \textbf{(P\_8,2.55.1)} paryāgataḥ kārśyena parikṛśaḥ . \\
\noindent \textbf{(P\_8,2.55.1)} jagatī anūnā bhavati hi rucirā . \\
\noindent \textbf{(P\_8,2.55.2)} phaleḥ latve utpūrvasya upasaṅkhyānam . \\
\noindent \textbf{(P\_8,2.55.2)} phaleḥ latve utpūrvasya upasaṅkhyānam kartavyam : utphullaḥ anṛtam kathayati . \\
\noindent \textbf{(P\_8,2.55.2)} atyalpam idam ucyate utpūrvāt iti . \\
\noindent \textbf{(P\_8,2.55.2)} utphullasamphullayoḥ iti vaktavyam : utphullaḥ , samphullaḥ . \\
\noindent \textbf{(P\_8,2.55.2)} kṛśeḥ kaḥ eṣaḥ vihitaḥ igupadhāt . \\
\noindent \textbf{(P\_8,2.55.2)} svare hi doṣaḥ bhavati parikṛśe . \\
\noindent \textbf{(P\_8,2.55.2)} padasya lopaḥ vihitaḥ iti matam . \\
\noindent \textbf{(P\_8,2.55.2)} jagati anūnā bhavati hi rucirā . \\
\noindent \textbf{(P\_8,2.56)} kim ayam vidhiḥ āhosvit pratiṣedhaḥ . \\
\noindent \textbf{(P\_8,2.56)} kim ca ataḥ . \\
\noindent \textbf{(P\_8,2.56)} yadi tāvat vidhiḥ nakāragrahaṇam kartavyam . \\
\noindent \textbf{(P\_8,2.56)} na kartavyam . \\
\noindent \textbf{(P\_8,2.56)} prakṛtam anuvartate . \\
\noindent \textbf{(P\_8,2.56)} kva prakṛtam . \\
\noindent \textbf{(P\_8,2.56)} radābhyānniṣṭhātonaḥpūrvasyacadaḥ iti . \\
\noindent \textbf{(P\_8,2.56)} tat vā anekena nipātanena vyavacchinnam na śakyam anuvartayitum . \\
\noindent \textbf{(P\_8,2.56)} atha pratiṣedhaḥ hrīgrahaṇam anarthakam na hi etasmāt vidhiḥ asti . \\
\noindent \textbf{(P\_8,2.56)} yathā icchasi tathā astu . \\
\noindent \textbf{(P\_8,2.56)} astu tāvat vidhiḥ . \\
\noindent \textbf{(P\_8,2.56)} nanu ca uktam nakāragrahaṇam kartavyam . \\
\noindent \textbf{(P\_8,2.56)} na kartavyam . \\
\noindent \textbf{(P\_8,2.56)} prakṛtam anuvartate . \\
\noindent \textbf{(P\_8,2.56)} kva prakṛtam . \\
\noindent \textbf{(P\_8,2.56)} radābhyānniṣṭhātonaḥpūrvasyacadaḥ iti . \\
\noindent \textbf{(P\_8,2.56)} tat vā anekena nipātanena vyavacchinnam na śakyam anuvartayitum iti . \\
\noindent \textbf{(P\_8,2.56)} sambandham anuvartiṣyate . \\
\noindent \textbf{(P\_8,2.56)} atha vā kriyate nyāse eva . \\
\noindent \textbf{(P\_8,2.56)} dvinakārakaḥ nirdeśaḥ . \\
\noindent \textbf{(P\_8,2.56)} nudavidondatrāghrāhrībhyaḥ anyatarasyām n na dhyākhyāpṝmūrchimadām iti . \\
\noindent \textbf{(P\_8,2.56)} atha vā punaḥ astu pratiṣedhaḥ . \\
\noindent \textbf{(P\_8,2.56)} nanu ca uktam hrīgrahaṇam anarthakam na hi etasmāt vidhiḥ asti iti . \\
\noindent \textbf{(P\_8,2.56)} na anarthakam . \\
\noindent \textbf{(P\_8,2.56)} etat eva jñāpayati ācāryaḥ bhavati etasmāt vidhiḥ iti yat ayam hrīgrahaṇam karoti \\
\noindent \textbf{(P\_8,2.58)} bahavaḥ ime vidayaḥ paṭhyante . \\
\noindent \textbf{(P\_8,2.58)} tatra na jñāyate kasya nityam natvam kasya vibhāṣā kasya pratiṣedhaḥ kasya iṭ iti . \\
\noindent \textbf{(P\_8,2.58)} ataḥ uttaram paṭhati . \\
\noindent \textbf{(P\_8,2.58)} yasya videḥ śnaśakau taparatve tanavacane tad u vāpratiṣedhau . \\
\noindent \textbf{(P\_8,2.58)} śnavikaraṇasya vibhāṣā śavikaraṇasya pratiṣedhaḥ . \\
\noindent \textbf{(P\_8,2.58)} śyavikaraṇāt navidhiḥ chiditulyaḥ . śyanvikaraṇāt videḥ navidhiḥ chidinā tulyaḥ . \\
\noindent \textbf{(P\_8,2.58)} lugvikaraṇaḥ vali paryavapannaḥ . lugvikaraṇaḥ vidiḥ valādau paryavapannaḥ . \\
\noindent \textbf{(P\_8,2.58)} eṣa evārthaḥ . \\
\noindent \textbf{(P\_8,2.58)} yayoḥ vidyoḥ śnaśau uktau tayoḥ natvasya vānañau . \\
\noindent \textbf{(P\_8,2.58)} yayoḥ tu śyaṃllukau tābhyām chidivac ca iṭ ca iṣyate . aparaḥ āha : vetteḥ tu viditaḥ niṣṭhā . \\
\noindent \textbf{(P\_8,2.58)} vidyateḥ vinnaḥ iṣyate . \\
\noindent \textbf{(P\_8,2.58)} vintteḥ vinnaḥ ca vittaḥ ca vittaḥ bhogeṣu vindateḥ . \\
\noindent \textbf{(P\_8,2.58)} bhittam śakalam \\
\noindent \textbf{(P\_8,2.59)} bhittam śakalam iti ucyate tatra idam na sidhyati bhittam bhinnam iti . \\
\noindent \textbf{(P\_8,2.59)} na eṣaḥ doṣaḥ . \\
\noindent \textbf{(P\_8,2.59)} sarvatra eva atra bhidiḥ vidāraṇasāmānye vartate tatra avaśyam viśeṣārthinā viśeṣaḥ anuprayoktavyaḥ . \\
\noindent \textbf{(P\_8,2.59)} bhinnam kim bhittam iti . \\
\noindent \textbf{(P\_8,2.59)} tatvam abhidhāyakam cet śakalasya anarthakaḥ prayogaḥ syāt . \\
\noindent \textbf{(P\_8,2.59)} śakalena ca api abhihite na bhavati tatvam nigamayāmaḥ . \\
\noindent \textbf{(P\_8,2.59)} kvin pratyayasya kuḥ \\
\noindent \textbf{(P\_8,2.62)} pratyayagrahaṇam kimartham na kvinaḥ kuḥ iti eva ucyeta . \\
\noindent \textbf{(P\_8,2.62)} kvinaḥ kuḥ iti iyati ucyamāne vakārasya eva kutvam prasajyeta . \\
\noindent \textbf{(P\_8,2.62)} nanu ca lope kṛte na bhaviṣyati . \\
\noindent \textbf{(P\_8,2.62)} anavakāśam kutvam lopam bādheta . \\
\noindent \textbf{(P\_8,2.62)} sāvakāśam kutvam . \\
\noindent \textbf{(P\_8,2.62)} kaḥ avakāśaḥ . \\
\noindent \textbf{(P\_8,2.62)} anantyaḥ . \\
\noindent \textbf{(P\_8,2.62)} katham punaḥ sati antye anantyasya kutvam syāt . \\
\noindent \textbf{(P\_8,2.62)} ācāryapravṛttiḥ jñāpayati nāntyasya kutvam bhavati iti yat ayam kvinaḥ kuḥ iti kavarganirdeśam karoti . \\
\noindent \textbf{(P\_8,2.62)} itarathā hi tadguṇam eva ayam nirdiśet . \\
\noindent \textbf{(P\_8,2.62)} idam tarhi prayojanam yebhyaḥ kvinpratyayaḥ vidhīyate teṣām anyapratyayāntānām api padānte kutvam yathā syāt . \\
\noindent \textbf{(P\_8,2.62)} mā naḥ asrāk . \\
\noindent \textbf{(P\_8,2.62)} mā naḥ adrāk . \\
\noindent \textbf{(P\_8,2.62)} kvinaḥ kuḥ iti vaktavye pratyayagrahaṇam kṛtam . \\
\noindent \textbf{(P\_8,2.62)} kvinpratyayasya sarvatra padānte kutvam iṣyate . \\
\noindent \textbf{(P\_8,2.68)} ruvidhau ahnaḥ rūparātrirathantareṣu upasaṅkhyānam . \\
\noindent \textbf{(P\_8,2.68)} ruvidhau ahnaḥ rūparātrirathantareṣu upasaṅkhyānam kartavyam . \\
\noindent \textbf{(P\_8,2.68)} ahorūpam ahorātraḥ ahorathantaram sāma \\
\noindent \textbf{(P\_8,2.69)} asupi rādeśe upasarjanasamāse pratiṣedhaḥ aluki . \\
\noindent \textbf{(P\_8,2.69)} asupi rādeśe upasarjanasamāse aluki pratiṣedhaḥ vaktavyaḥ . \\
\noindent \textbf{(P\_8,2.69)} dīrghāhā nidāghaḥ iti . \\
\noindent \textbf{(P\_8,2.69)} siddham tu supi pratiṣedhāt . siddham etat . \\
\noindent \textbf{(P\_8,2.69)} katham . \\
\noindent \textbf{(P\_8,2.69)} supi pratiṣedhāt . \\
\noindent \textbf{(P\_8,2.69)} prasajya ayam pratiṣedhaḥ supi na iti . \\
\noindent \textbf{(P\_8,2.69)} iha api tarhi na prāpnoti . \\
\noindent \textbf{(P\_8,2.69)} ahan dadāti . \\
\noindent \textbf{(P\_8,2.69)} ahan bhuṅkte iti . \\
\noindent \textbf{(P\_8,2.69)} luki ca uktam . kim uktam . \\
\noindent \textbf{(P\_8,2.69)} ahnaḥ ravidhau lumatā lupte pratyayalakṣaṇam na bhavati iti \\
\noindent \textbf{(P\_8,2.70)} chandasi bhāṣāyām ca pracetasaḥ rājani upasaṅkhyānam . \\
\noindent \textbf{(P\_8,2.70)} chandasi bhāṣāyām ca pracetasaḥ rājani upasaṅkhyānam kartavyam . \\
\noindent \textbf{(P\_8,2.70)} pracetaḥ rājan . \\
\noindent \textbf{(P\_8,2.70)} pracetar rājan . \\
\noindent \textbf{(P\_8,2.70)} aharādīnām patyādiṣu upasaṅkhyānam kartavyam . \\
\noindent \textbf{(P\_8,2.70)} aharpatiḥ ahaḥpatiḥ . \\
\noindent \textbf{(P\_8,2.70)} aharputraḥ ahaḥputraḥ . \\
\noindent \textbf{(P\_8,2.70)} gīrpatiḥ gīḥpatiḥ \\
\noindent \textbf{(P\_8,2.72)} iha kasmāt na bhavati . \\
\noindent \textbf{(P\_8,2.72)} papivān tasthivān iti . \\
\noindent \textbf{(P\_8,2.72)} sasya iti vartate . \\
\noindent \textbf{(P\_8,2.72)} evam api atra prāpnoti . \\
\noindent \textbf{(P\_8,2.72)} lope kṛte na bhaviṣyati . \\
\noindent \textbf{(P\_8,2.72)} anavakāśam datvam lopam bādheta . \\
\noindent \textbf{(P\_8,2.72)} sāvakāśam datvam . \\
\noindent \textbf{(P\_8,2.72)} kaḥ avakāśaḥ . \\
\noindent \textbf{(P\_8,2.72)} papivadbhyām papivadbhiḥ iti . \\
\noindent \textbf{(P\_8,2.72)} atra api ruḥ prāpnoti . \\
\noindent \textbf{(P\_8,2.72)} tat yathā eva rum bādhate evam lopam api bādheta . \\
\noindent \textbf{(P\_8,2.72)} na bādhate . \\
\noindent \textbf{(P\_8,2.72)} kim kāraṇam . \\
\noindent \textbf{(P\_8,2.72)} yena na aprāpte tasya bādhanam bhavati na ca aprāpte rau datvam ārabhyate lope punaḥ prāpte ca aprāpte ca . \\
\noindent \textbf{(P\_8,2.72)} yadi tarhi sasya iti vartate anaḍudbhyām anaḍudbhiḥ iti atra na prāpnoti . \\
\noindent \textbf{(P\_8,2.72)} vacanāt anaḍuhi bhaviṣyati . \\
\noindent \textbf{(P\_8,2.72)} yadi evam . \\
\noindent \textbf{(P\_8,2.72)} anaḍuhaḥ datve nakārapratiṣedhaḥ . \\
\noindent \textbf{(P\_8,2.72)} anaḍuhaḥ datve nakārasya pratiṣedhaḥ vaktavyaḥ . \\
\noindent \textbf{(P\_8,2.72)} anaḍvān . \\
\noindent \textbf{(P\_8,2.72)} siddham tu pratipadavidhānāt numaḥ . siddham etat . \\
\noindent \textbf{(P\_8,2.72)} katham . \\
\noindent \textbf{(P\_8,2.72)} numaḥ pratipadavidhānasāmarthyāt datvam na bhaviṣyati . \\
\noindent \textbf{(P\_8,2.72)} yadi tarhi yat yat anaḍuhaḥ prāptam tat tat numaḥ pratipadavidhānasāmarthyāt bādhyate rutvam api na prāpnoti . \\
\noindent \textbf{(P\_8,2.72)} anaḍvān tatra iti . \\
\noindent \textbf{(P\_8,2.72)} na eṣaḥ doṣaḥ . \\
\noindent \textbf{(P\_8,2.72)} yam vidhim prati upadeśaḥ anarthakaḥ saḥ vidhiḥ bādhyate yasya tu vidheḥ nimittam eva na asau bādhyate . \\
\noindent \textbf{(P\_8,2.72)} datvam ca prati numaḥ pratipadavidhiḥ anarthakaḥ roḥ punaḥ nimittam eva \\
\noindent \textbf{(P\_8,2.78.1)} kimartham idam ucyate na hali iti eva siddham . \\
\noindent \textbf{(P\_8,2.78.1)} na sidhyati . \\
\noindent \textbf{(P\_8,2.78.1)} dhātoḥ iti tatra vartate tatra rephavakārābhyām dhātuḥ viśeṣyate . \\
\noindent \textbf{(P\_8,2.78.1)} rephavakārāntasya dhātoḥ iti . \\
\noindent \textbf{(P\_8,2.78.1)} kim punaḥ kāraṇam pūrvasmin yoge rephavakārābhyām dhātuḥ viśeṣyate . \\
\noindent \textbf{(P\_8,2.78.1)} iha mā bhūt . \\
\noindent \textbf{(P\_8,2.78.1)} agniḥ vāyuḥ iti . \\
\noindent \textbf{(P\_8,2.78.1)} evam tarhi pūrvasmin yoge yat dhātugrahaṇam tat uttaratra nivṛttam . \\
\noindent \textbf{(P\_8,2.78.1)} evam api kurkuraḥ murmuraḥ iti atra api prāpnoti . \\
\noindent \textbf{(P\_8,2.78.1)} evam tarhi anuvartate tatra dhātugrahaṇam na tu rephavakārābhyām dhātuḥ viśeṣyate . \\
\noindent \textbf{(P\_8,2.78.1)} kim tarhi . \\
\noindent \textbf{(P\_8,2.78.1)} ik viśeṣyate . \\
\noindent \textbf{(P\_8,2.78.1)} rephavakārāntasya ikaḥ dhātoḥ iti . \\
\noindent \textbf{(P\_8,2.78.1)} evam api kurkurīyati murmurīyati iti atra prāpnoti . \\
\noindent \textbf{(P\_8,2.78.1)} tasmāt dhātuḥ eva viśeṣyaḥ dhātau ca viśeṣyamāṇe upadhāyām ca iti vaktavyam \\
\noindent \textbf{(P\_8,2.78.2)} upadhādīrghatve abhyāsajivricaturṇām pratiṣedhaḥ . \\
\noindent \textbf{(P\_8,2.78.2)} upadhādīrghatve abhyāsajivṛicaturṇām pratiṣedhaḥ vaktavyaḥ . \\
\noindent \textbf{(P\_8,2.78.2)} riryatuḥ riryuḥ . \\
\noindent \textbf{(P\_8,2.78.2)} saṃvivyatuḥ saṃvivyuḥ . \\
\noindent \textbf{(P\_8,2.78.2)} jivraḥ . \\
\noindent \textbf{(P\_8,2.78.2)} caturyitā caturyitum . \\
\noindent \textbf{(P\_8,2.78.2)} uṇādipratiṣedhaḥ ca . uṇādīnām ca pratiṣedhaḥ vaktavyaḥ . \\
\noindent \textbf{(P\_8,2.78.2)} kiryoḥ giryoḥ iti . \\
\noindent \textbf{(P\_8,2.78.2)} abhyāsapratiṣedhaḥ tāvat na vaktavyaḥ . \\
\noindent \textbf{(P\_8,2.78.2)} hali iti ucyate na ca atra halādim paśyāmaḥ . \\
\noindent \textbf{(P\_8,2.78.2)} yaṇādeśe kṛte prāpnoti . \\
\noindent \textbf{(P\_8,2.78.2)} sthānivadbhāvāt na bhaviṣyati . \\
\noindent \textbf{(P\_8,2.78.2)} pratiṣidhyate atra sthānivadbhāvaḥ dīrghavidhim prati na sthānivat iti . \\
\noindent \textbf{(P\_8,2.78.2)} na eṣaḥ asti pratiṣedhaḥ . \\
\noindent \textbf{(P\_8,2.78.2)} uktam etat pratiṣedhe svaradirghayalopeṣu lopājādeśaḥ iti . \\
\noindent \textbf{(P\_8,2.78.2)} jivripratiṣedhaḥ ca na vaktavyaḥ . \\
\noindent \textbf{(P\_8,2.78.2)} uṇādayaḥ avyutpannāni prātipadikāni . \\
\noindent \textbf{(P\_8,2.78.2)} caturyitā caturyitum iti supi na iti vartate . \\
\noindent \textbf{(P\_8,2.78.2)} yadi evam gīrbhyām gīrbhiḥ iti aprasiddhiḥ . \\
\noindent \textbf{(P\_8,2.78.2)} na supaḥ vibhaktivipariṇāmāt gīrbhyām gīrbhiḥ iti adoṣaḥ . \\
\noindent \textbf{(P\_8,2.78.2)} uṇādipratiṣedhaḥ vaktavyaḥ iti parihṛtam etat uṇādayaḥ avyutpannāni prātipadikāni iti \\
\noindent \textbf{(P\_8,2.80.1)} adasaḥ anosreḥ . \\
\noindent \textbf{(P\_8,2.80.1)} adasaḥ anosreḥ iti vaktavyam . \\
\noindent \textbf{(P\_8,2.80.1)} kim idam anosreḥ iti . \\
\noindent \textbf{(P\_8,2.80.1)} anokārasya asakārasya arephakasya iti . \\
\noindent \textbf{(P\_8,2.80.1)} anokārasya . \\
\noindent \textbf{(P\_8,2.80.1)} adaḥ atra . \\
\noindent \textbf{(P\_8,2.80.1)} asakārasya . \\
\noindent \textbf{(P\_8,2.80.1)} adasyate . \\
\noindent \textbf{(P\_8,2.80.1)} arephakasya . \\
\noindent \textbf{(P\_8,2.80.1)} adaḥ . \\
\noindent \textbf{(P\_8,2.80.1)} tat tarhi vaktavyam . \\
\noindent \textbf{(P\_8,2.80.1)} na vaktavyam . \\
\noindent \textbf{(P\_8,2.80.1)} kriyate nyāse eva . \\
\noindent \textbf{(P\_8,2.80.1)} avibhaktikaḥ nirdeśaḥ . \\
\noindent \textbf{(P\_8,2.80.1)} adas , o , iti . \\
\noindent \textbf{(P\_8,2.80.1)} okārāt paraḥ patiṣedhaḥ pūrvabhūtaḥ . \\
\noindent \textbf{(P\_8,2.80.1)} tataḥ sakāraḥ . \\
\noindent \textbf{(P\_8,2.80.1)} tataḥ rephaḥ iti . \\
\noindent \textbf{(P\_8,2.80.1)} atha vā na evam vijñāyate adasaḥ asakārasya iti . \\
\noindent \textbf{(P\_8,2.80.1)} katham tarhi . \\
\noindent \textbf{(P\_8,2.80.1)} akāraḥ asya sakārasya saḥ ayam asiḥ aseḥ iti . \\
\noindent \textbf{(P\_8,2.80.1)} yadi evam amumuyaṅ iti na sidhyati adadryaṅ iti prāpnoti . \\
\noindent \textbf{(P\_8,2.80.1)} adamuyaṅ iti bhavitavyam anantyavikāre antyasadeśasya kāryam bhavati iti . \\
\noindent \textbf{(P\_8,2.80.1)} adasaḥ adreḥ pṛthak mutvam kecit icchanti latvavat . \\
\noindent \textbf{(P\_8,2.80.1)} kecit antyasadeśasya . \\
\noindent \textbf{(P\_8,2.80.1)} na iti eke . \\
\noindent \textbf{(P\_8,2.80.1)} aseḥ hi dṛśyate . \\
\noindent \textbf{(P\_8,2.80.1)} . \\
\noindent \textbf{(P\_8,2.80.2)} tatra padādhikārāt apadāntasya aprāptiḥ . \\
\noindent \textbf{(P\_8,2.80.2)} tatra padādhikārāt apadāntasya na prāpnoti . \\
\noindent \textbf{(P\_8,2.80.2)} amuyā amuyoḥ iti . \\
\noindent \textbf{(P\_8,2.80.2)} siddham tu sakārapratiṣedhāt . siddham etat . \\
\noindent \textbf{(P\_8,2.80.2)} katham . \\
\noindent \textbf{(P\_8,2.80.2)} sakārapratiṣedhāt . \\
\noindent \textbf{(P\_8,2.80.2)} yat ayam aseḥ iti pratiṣedham śāsti tat jñāpayati ācāryaḥ apadāntasya api bhavati iti \\
\noindent \textbf{(P\_8,2.80.3)} atha dādgrahaṇam kimartham . \\
\noindent \textbf{(P\_8,2.80.3)} dādgrahaṇam antyapratiṣedhārtham . \\
\noindent \textbf{(P\_8,2.80.3)} dādgrahaṇam kriyate antyapratiṣedhārtham . \\
\noindent \textbf{(P\_8,2.80.3)} alaḥ antyasya mā bhūt iti . \\
\noindent \textbf{(P\_8,2.80.3)} amuyā amuyoḥ iti \\
\noindent \textbf{(P\_8,2.81)} īttvam bahuvacanāntasya . \\
\noindent \textbf{(P\_8,2.81)} īttvam bahuvacanāntasya iti vaktavyam . \\
\noindent \textbf{(P\_8,2.81)} bahuvacane iti iyati ucyamāne iha eva syāt . \\
\noindent \textbf{(P\_8,2.81)} amībhiḥ amīṣu. iha na syāt . \\
\noindent \textbf{(P\_8,2.81)} amī atra . \\
\noindent \textbf{(P\_8,2.81)} amī āsate . \\
\noindent \textbf{(P\_8,2.81)} tat tarhi vaktavyam . \\
\noindent \textbf{(P\_8,2.81)} na vaktavyam . \\
\noindent \textbf{(P\_8,2.81)} na idam pāribhāṣikasya bahuvacanasya grahaṇam . \\
\noindent \textbf{(P\_8,2.81)} kim tarhi . \\
\noindent \textbf{(P\_8,2.81)} anvarthagrahaṇam etat . \\
\noindent \textbf{(P\_8,2.81)} bahūnām arthānām vacanam bahuvacanam bahuvacane iti \\
\noindent \textbf{(P\_8,2.82)} vākyādhikāraḥ kimarthaḥ . \\
\noindent \textbf{(P\_8,2.82)} vākyādhikāraḥ padanivṛttyarthaḥ . \\
\noindent \textbf{(P\_8,2.82)} vākyādhikāraḥ kriyate padanivṛttyarthaḥ . \\
\noindent \textbf{(P\_8,2.82)} padādhikāraḥ nivartyate . \\
\noindent \textbf{(P\_8,2.82)} na hi kākaḥ vāśyate iti adhikārāḥ nivartante . \\
\noindent \textbf{(P\_8,2.82)} doṣaḥ khalu api syāt yadi vākyādhikāraḥ padādhikāram nivartayet . \\
\noindent \textbf{(P\_8,2.82)} iṣyante eva uttaratra padakāryāṇi tāni na sidhyanti . \\
\noindent \textbf{(P\_8,2.82)} naśchavyapraśān iti . \\
\noindent \textbf{(P\_8,2.82)} padanivṛttyartham iti na evam vijñāyate padasya nivṛttyartham padanivṛttyartham iti . \\
\noindent \textbf{(P\_8,2.82)} kim tarhi. pade nivṛttyartham padanivṛttyartham iti . \\
\noindent \textbf{(P\_8,2.82)} vākye yāvanti padāni teṣām sarveṣām ṭeḥ plutaḥ prāpnoti . \\
\noindent \textbf{(P\_8,2.82)} iṣyate ca vākyapadayoḥ antyasya syāt iti tat ca antareṇa yatnam na sidhyati iti evamarthaḥ vākyādhikāraḥ . \\
\noindent \textbf{(P\_8,2.82)} atha ṭigrahaṇam kimartham . \\
\noindent \textbf{(P\_8,2.82)} ṭigrahaṇam alaḥ antyaniyame vyañjanāntārtham . ṭigrahaṇam kriyate alaḥ antyaniyame vyañjanāntasya api yathā syāt . \\
\noindent \textbf{(P\_8,2.82)} agnici3t somasu3t . \\
\noindent \textbf{(P\_8,2.82)} asti prayojanam etat . \\
\noindent \textbf{(P\_8,2.82)} kim tarhi iti . \\
\noindent \textbf{(P\_8,2.82)} sarvādeśaprasaṅgaḥ tu . sarvādeśaḥ tu ṭeḥ plutaḥ prāpnoti . \\
\noindent \textbf{(P\_8,2.82)} kim kāraṇam . \\
\noindent \textbf{(P\_8,2.82)} acaḥ iti vacanāt antyasya na antyasya iti vacanāt acaḥ na ucyate ca plutaḥ saḥ sarvādeśaḥ prāpnoti . \\
\noindent \textbf{(P\_8,2.82)} uktam vā . kim uktam . \\
\noindent \textbf{(P\_8,2.82)} hrasvaḥ dīrghaḥ plutaḥ iti yatra brūyāt acaḥ iti etat tatra upasthitam draṣṭavyam iti \\
\noindent \textbf{(P\_8,2.83.1)} aśūdre iti kimartham . \\
\noindent \textbf{(P\_8,2.83.1)} kuśalī asi tuṣajaka . \\
\noindent \textbf{(P\_8,2.83.1)} atyalpam idam ucyate : asūdre iti . \\
\noindent \textbf{(P\_8,2.83.1)} aśūdrastryasūyakeṣu . aśūdrastryasūyakeṣu iti vaktavyam . \\
\noindent \textbf{(P\_8,2.83.1)} tatra śūdre udāhṛtam . \\
\noindent \textbf{(P\_8,2.83.1)} striyām : gārgī aham , bhoḥ āyuṣmatī bhava gārgi . \\
\noindent \textbf{(P\_8,2.83.1)} asūyake : sthālī aham , bhoḥ āyuṣmān edhi sthāli3n . \\
\noindent \textbf{(P\_8,2.83.1)} na eṣā mama sañjñā sthālī iti . \\
\noindent \textbf{(P\_8,2.83.1)} kim tarhi daṇḍinyāyaḥ mama vivakṣitaḥ . \\
\noindent \textbf{(P\_8,2.83.1)} saḥ vaktavyaḥ . \\
\noindent \textbf{(P\_8,2.83.1)} sthālī aham bhoḥ āyuṣmān edhi sthālin . \\
\noindent \textbf{(P\_8,2.83.1)} na mama daṇḍinyāyaḥ vivakṣitaḥ . \\
\noindent \textbf{(P\_8,2.83.1)} kim tarhi sañjñā mama eṣā . \\
\noindent \textbf{(P\_8,2.83.1)} asūyakaḥ tvam asi jālma na tvam pratyabhivādam arhasi bhidyasva vṛṣala sthālin \\
\noindent \textbf{(P\_8,2.83.2)} bhorājanyaviśām vā . \\
\noindent \textbf{(P\_8,2.83.2)} bhorājanyaviśām vā iti vaktavyam . \\
\noindent \textbf{(P\_8,2.83.2)} devadattaḥ aham bhoḥ āyuṣmān edhi devadatta bho3ḥ . \\
\noindent \textbf{(P\_8,2.83.2)} devadatta bhoḥ . \\
\noindent \textbf{(P\_8,2.83.2)} bhoḥ . \\
\noindent \textbf{(P\_8,2.83.2)} rājanya . \\
\noindent \textbf{(P\_8,2.83.2)} indravarmā aham bhoḥ āyuṣmān edhi indravarma3n . \\
\noindent \textbf{(P\_8,2.83.2)} indravarman . \\
\noindent \textbf{(P\_8,2.83.2)} rājanya . \\
\noindent \textbf{(P\_8,2.83.2)} viṭ . \\
\noindent \textbf{(P\_8,2.83.2)} indrapālitaḥ aham bhoḥ āyuṣmān edhi indrapālita3 . \\
\noindent \textbf{(P\_8,2.83.2)} indrapālita . \\
\noindent \textbf{(P\_8,2.83.2)} aparaḥ āha : . \\
\noindent \textbf{(P\_8,2.83.2)} sarvasya eva nāmnaḥ pratyabhivāde bhoḥśabdaḥ ādeśaḥ vaktavyaḥ . \\
\noindent \textbf{(P\_8,2.83.2)} devadattaḥ aham bhoḥ āyuṣmān edhi bho3ḥ . \\
\noindent \textbf{(P\_8,2.83.2)} āyuṣmān edhi devadatta3 iti vā \\
\noindent \textbf{(P\_8,2.83.3)} iha kasmāt na bhavati . \\
\noindent \textbf{(P\_8,2.83.3)} devadatta kuśalī asi iti . \\
\noindent \textbf{(P\_8,2.83.3)} iha kim cit ucyate kim cit pratyucyate . \\
\noindent \textbf{(P\_8,2.83.3)} apradhānam ucyate pradhānam pratyucyate . \\
\noindent \textbf{(P\_8,2.83.3)} tatra pradhānasthasya ṭisañjñakasya plutyā bhavitavyam na ca atra pradhānastham ṭisañjñam . \\
\noindent \textbf{(P\_8,2.83.3)} iha api tarhi na prāpnoti . \\
\noindent \textbf{(P\_8,2.83.3)} ādheyaḥ agni3ḥ na ādheya3ḥ iti . \\
\noindent \textbf{(P\_8,2.83.3)} na etat vicāryate ādheyaḥ na ādheyaḥ agniḥ cet bhavati iti . \\
\noindent \textbf{(P\_8,2.83.3)} kim tarhi . \\
\noindent \textbf{(P\_8,2.83.3)} iha agnisādhanā kriyā vicāryate ādheyaḥ agniḥ na ādheyaḥ iti . \\
\noindent \textbf{(P\_8,2.83.3)} yadi evam dvitīyaḥ agniśabdasya prayogaḥ prāpnoti . \\
\noindent \textbf{(P\_8,2.83.3)} uktārthānām aprayogaḥ iti na bhaviṣyati . \\
\noindent \textbf{(P\_8,2.83.3)} yadi evam ādheyaśabdasya api tarhi dvitīyasya prayogaḥ na prāpnoti uktārthānām aprayogaḥ nāma bhavati iti . \\
\noindent \textbf{(P\_8,2.83.3)} na eṣaḥ doṣaḥ . \\
\noindent \textbf{(P\_8,2.83.3)} uktārthānām api prayogaḥ dṛśyate . \\
\noindent \textbf{(P\_8,2.83.3)} tat yathā . \\
\noindent \textbf{(P\_8,2.83.3)} apūpau dvau ānaya . \\
\noindent \textbf{(P\_8,2.83.3)} brāhmaṇau dvau ānaya iti \\
\noindent \textbf{(P\_8,2.84)} dūrāt hūte iti ucyate dūraśabdaḥ ca ayam anavasthitapadārthakaḥ . \\
\noindent \textbf{(P\_8,2.84)} tat eva hi kam cit prati dūram kam cit prati antikam bhavati . \\
\noindent \textbf{(P\_8,2.84)} evam hi kaḥ cit kam cit āha . \\
\noindent \textbf{(P\_8,2.84)} eṣaḥ pārśvataḥ karakaḥ tam ānaya iti . \\
\noindent \textbf{(P\_8,2.84)} saḥ āha . \\
\noindent \textbf{(P\_8,2.84)} utthāya gṛhāṇa dūram na śakṣyāmi iti . \\
\noindent \textbf{(P\_8,2.84)} aparaḥ āha : dūram mathurāyāḥ pāṭaliputram iti . \\
\noindent \textbf{(P\_8,2.84)} saḥ āha . \\
\noindent \textbf{(P\_8,2.84)} na dūram idam antikam iti . \\
\noindent \textbf{(P\_8,2.84)} evam eṣaḥ dūraśabdaḥ anavasthitapadārthakaḥ tasya anavasthitapadārthakatvāt na jñāyate kasyām avasthāyām plutyā bhavitavyam iti . \\
\noindent \textbf{(P\_8,2.84)} evam tarhi hvayatinā ayam nirdeśaḥ kriyate . \\
\noindent \textbf{(P\_8,2.84)} hvayatiprasaṅge yat dūram . \\
\noindent \textbf{(P\_8,2.84)} kim punaḥ tat . \\
\noindent \textbf{(P\_8,2.84)} tatra prākṛtāt prayatnāt prayatnaviśeṣe upādīyamāne sandehaḥ bhavati śroṣyati na śroṣyati iti tat dūram iha avagamyate \\
\noindent \textbf{(P\_8,2.85)} haihegrahaṇam kimartham . \\
\noindent \textbf{(P\_8,2.85)} haiheprayoge haihegrahaṇam haihayoḥ plutyartham . \\
\noindent \textbf{(P\_8,2.85)} haiheprayoge ḥaihegrahaṇam kriyate haihayoḥ plutiḥ yathā syāt . \\
\noindent \textbf{(P\_8,2.85)} devadatta hai3 . \\
\noindent \textbf{(P\_8,2.85)} devadatta he3 . \\
\noindent \textbf{(P\_8,2.85)} akriyamāṇe hi haihegrahaṇe tayoḥ prayoge anyasya syāt . \\
\noindent \textbf{(P\_8,2.85)} atha prayogagrahaṇam kimartham . \\
\noindent \textbf{(P\_8,2.85)} prayogagrahaṇam arthavadgrahaṇe anarthakārtham . prayogagrahaṇam kriyate arthavadgrahaṇe anarthakayoḥ api yathā syāt . \\
\noindent \textbf{(P\_8,2.85)} devadatta hai3 . \\
\noindent \textbf{(P\_8,2.85)} devadatta he3 . \\
\noindent \textbf{(P\_8,2.85)} atha punaḥ haihegrahaṇam kimartham . \\
\noindent \textbf{(P\_8,2.85)} punaḥ haihegrahaṇam anantyārtham . punaḥ haihegrahaṇam kriyate anantyayoḥ api yathā syāt . \\
\noindent \textbf{(P\_8,2.85)} hai3 devadatta . \\
\noindent \textbf{(P\_8,2.85)} he3 devadatta iti \\
\noindent \textbf{(P\_8,2.86.1)} guroḥ plutavidhāne laghoḥ antyasya plutaprasaṅgaḥ anyena vihitatvāt . \\
\noindent \textbf{(P\_8,2.86.1)} guroḥ plutavidhāne laghoḥ antyasya plutaḥ prāpnoti . \\
\noindent \textbf{(P\_8,2.86.1)} de3vadatta . \\
\noindent \textbf{(P\_8,2.86.1)} kim kāraṇam . \\
\noindent \textbf{(P\_8,2.86.1)} anyena vihitatvāt . \\
\noindent \textbf{(P\_8,2.86.1)} anyena hi lakṣaṇena laghoḥ antyasya plutaḥ vidhīyate dūrāddhūteca iti . \\
\noindent \textbf{(P\_8,2.86.1)} na vā anantyasya api iti vacanam ubhayanirdeśārtham . na vā eṣaḥ doṣaḥ . \\
\noindent \textbf{(P\_8,2.86.1)} kim kāraṇam . \\
\noindent \textbf{(P\_8,2.86.1)} anantyasya api iti vacanam ubhayanirdeśārtham bhaviṣyati . \\
\noindent \textbf{(P\_8,2.86.1)} anantyasya api guroḥ antyasya api ṭeḥ iti . \\
\noindent \textbf{(P\_8,2.86.1)} nanu ca etat gurvapekṣam syāt . \\
\noindent \textbf{(P\_8,2.86.1)} anantyasya api guroḥ antyasya api guroḥ iti . \\
\noindent \textbf{(P\_8,2.86.1)} na iti āha . \\
\noindent \textbf{(P\_8,2.86.1)} dvyapekṣam etat . \\
\noindent \textbf{(P\_8,2.86.1)} anantyasya api guroḥ antyasya api ṭeḥ iti \\
\noindent \textbf{(P\_8,2.86.2)} atha prāgvacanam kimartham . \\
\noindent \textbf{(P\_8,2.86.2)} prāgvacanam vibhāṣārtham . \\
\noindent \textbf{(P\_8,2.86.2)} prāgvacanam kriyate vibhāṣā yathā syāt . \\
\noindent \textbf{(P\_8,2.86.2)} prāgvacanānarthakyam ca ekaikasya iti vacanāt . prāgvacanam anarthakam . \\
\noindent \textbf{(P\_8,2.86.2)} kim kāraṇam . \\
\noindent \textbf{(P\_8,2.86.2)} ekaikasya iti vacanāt . \\
\noindent \textbf{(P\_8,2.86.2)} ekaikagrahaṇam kriyate tat vibhāṣārtham bhaviṣyati . \\
\noindent \textbf{(P\_8,2.86.2)} asti anyat ekaikagrahaṇasya prayojanam . \\
\noindent \textbf{(P\_8,2.86.2)} kim . \\
\noindent \textbf{(P\_8,2.86.2)} yugapat plutaḥ mā bhūt iti . \\
\noindent \textbf{(P\_8,2.86.2)} anudāttam padam ekavarjam iti vacanāt na asti yaugapadyena sambhavaḥ . \\
\noindent \textbf{(P\_8,2.86.2)} asiddhaḥ plutaḥ tasya asiddhatvāt niyamaḥ na prāpnoti . \\
\noindent \textbf{(P\_8,2.86.2)} na eṣaḥ doṣaḥ . \\
\noindent \textbf{(P\_8,2.86.2)} yadi api idam tatra asiddham tat tu iha siddham . \\
\noindent \textbf{(P\_8,2.86.2)} katham . \\
\noindent \textbf{(P\_8,2.86.2)} kāryakālam sañjñāparibhāṣam iti yatra kāryam tatra upasthitam draṣṭavyam . \\
\noindent \textbf{(P\_8,2.86.2)} guroḥ anṛtaḥ anantyasya api ekaikasya prācām . \\
\noindent \textbf{(P\_8,2.86.2)} upasthitam idam bhavati anudāttampadamekavarjam iti . \\
\noindent \textbf{(P\_8,2.86.2)} iha api tarhi samāveśaḥ na prāpnoti . \\
\noindent \textbf{(P\_8,2.86.2)} devadatta3 . \\
\noindent \textbf{(P\_8,2.86.2)} siddhāsiddhau etau . \\
\noindent \textbf{(P\_8,2.86.2)} yau hi siddhau eva asiddhau eva vā tayoḥ niyamaḥ . \\
\noindent \textbf{(P\_8,2.86.2)} yaḥ tarhi svaritaplutaḥ tena samāveśaḥ prāpnoti . \\
\noindent \textbf{(P\_8,2.86.2)} svaritamāmreḍite'sūyāsammatikopakutsaneṣu iti . \\
\noindent \textbf{(P\_8,2.86.2)} svarite api udāttaḥ asti . \\
\noindent \textbf{(P\_8,2.86.2)} yaḥ tarhi anudāttaplutaḥ tena samāveśaḥ prāpnoti . \\
\noindent \textbf{(P\_8,2.86.2)} anudāttampraśnāntābhipūjitayoḥ iti . \\
\noindent \textbf{(P\_8,2.86.2)} tasmāt prāgvacanam kartavyam \\
\noindent \textbf{(P\_8,2.88)} ye yajñakarmaṇi iti atiprasaṅgaḥ . \\
\noindent \textbf{(P\_8,2.88)} ye yajñakarmaṇi iti atiprasaṅgaḥ bhavati . \\
\noindent \textbf{(P\_8,2.88)} iha api prāpnoti . \\
\noindent \textbf{(P\_8,2.88)} ye devāsaḥ divyekādaśa stha iti . \\
\noindent \textbf{(P\_8,2.88)} siddham tu ye yajāmahe iti brūhyādiṣu upasaṅkhyānāt . siddham etat . \\
\noindent \textbf{(P\_8,2.88)} katham . \\
\noindent \textbf{(P\_8,2.88)} yeyajāmaheśabdaḥ brūhyādiṣu upasaṅkhyeyaḥ \\
\noindent \textbf{(P\_8,2.89)} praṇavaḥ iti ucyate kaḥ praṇavaḥ nāma . \\
\noindent \textbf{(P\_8,2.89)} pādasya vā ardharcasya vā antyam akṣaram upasaṃhṛtya tadādyakṣaraśeṣasya sthāne trimātram oṅkāram trimātram okāram vā vidadhati tam praṇavaḥ iti ācakṣate . \\
\noindent \textbf{(P\_8,2.89)} atha ṭigrahaṇam kimartham . \\
\noindent \textbf{(P\_8,2.89)} ṭigrahaṇam sarvādeśārtham . \\
\noindent \textbf{(P\_8,2.89)} yadā okāraḥ tadā sarvādeśaḥ yathā syāt . \\
\noindent \textbf{(P\_8,2.89)} yadā oṅkāraḥ tadā anekālśitsarvasya iti sarvādeśaḥ bhaviṣyati \\
\noindent \textbf{(P\_8,2.90)} antagrahaṇam kimartham . \\
\noindent \textbf{(P\_8,2.90)} yājyā nāma ṛcaḥ vākyasamudāyaḥ tatra yāvanti vākyāni sarveṣām ṭeḥ plutaḥ prāpnoti . \\
\noindent \textbf{(P\_8,2.90)} iṣyate ca antyasya syāt iti tat ca antareṇa yatnam na sidhyati iti evamartham antagrahaṇam \\
\noindent \textbf{(P\_8,2.92.1)} agnītpreṣaṇe iti atiprasaṅgaḥ . \\
\noindent \textbf{(P\_8,2.92.1)} agnītpreṣaṇe iti atiprasaṅgaḥ bhavati . \\
\noindent \textbf{(P\_8,2.92.1)} iha api prāpnoti . \\
\noindent \textbf{(P\_8,2.92.1)} agnīdagnīnvihara . \\
\noindent \textbf{(P\_8,2.92.1)} siddham tu ośrāvaye parasya ca iti vacanāt . siddham etat . \\
\noindent \textbf{(P\_8,2.92.1)} katham . \\
\noindent \textbf{(P\_8,2.92.1)} ośrāvaye parasya iti vaktavyam . \\
\noindent \textbf{(P\_8,2.92.1)} o3 śrā3vaya . \\
\noindent \textbf{(P\_8,2.92.1)} ā3 śrā3vaya . \\
\noindent \textbf{(P\_8,2.92.1)} aparaḥ āha : ośrāvayāśrāvayayoḥ iti vaktavyam . \\
\noindent \textbf{(P\_8,2.92.1)} o3 śrā3vaya . \\
\noindent \textbf{(P\_8,2.92.1)} a3 śra3vaya \\
\noindent \textbf{(P\_8,2.92.2)} bahulam anyatra iti vaktavyam . \\
\noindent \textbf{(P\_8,2.92.2)} uddhara3 uddhara . \\
\noindent \textbf{(P\_8,2.92.2)} āhara3 āhara . \\
\noindent \textbf{(P\_8,2.92.2)} tat tarhi vaktavyam . \\
\noindent \textbf{(P\_8,2.92.2)} na vaktavyam . \\
\noindent \textbf{(P\_8,2.92.2)} yogavibhāgaḥ kariṣyate . \\
\noindent \textbf{(P\_8,2.92.2)} agnītpreṣaṇe parasya ca vibhāṣā . \\
\noindent \textbf{(P\_8,2.92.2)} tataḥ pṛṣṭhaprativacane heḥ . \\
\noindent \textbf{(P\_8,2.92.2)} vibhāṣā iti eva . \\
\noindent \textbf{(P\_8,2.92.2)} aparaḥ āha : sarvaḥ eva plutaḥ sāhasam anicchatā vibhāṣā vaktavyaḥ \\
\noindent \textbf{(P\_8,2.95)} bhartsane paryāyeṇa . \\
\noindent \textbf{(P\_8,2.95)} bhartsane paryāyeṇa iti vaktavyam . \\
\noindent \textbf{(P\_8,2.95)} caura3 caura . \\
\noindent \textbf{(P\_8,2.95)} caura caura3 . \\
\noindent \textbf{(P\_8,2.95)} kuśīla3 kuśīla . \\
\noindent \textbf{(P\_8,2.95)} kuśīla kuśīla3 \\
\noindent \textbf{(P\_8,2.103)} asūyādiṣu vāvacanam . \\
\noindent \textbf{(P\_8,2.103)} asūyādiṣu vā iti vaktavyam . \\
\noindent \textbf{(P\_8,2.103)} kanye3 kanye . \\
\noindent \textbf{(P\_8,2.103)} kanye kanye . \\
\noindent \textbf{(P\_8,2.103)} śaktike3 śaktike . \\
\noindent \textbf{(P\_8,2.103)} śaktike śaktike \\
\noindent \textbf{(P\_8,2.106)} kimartham idam ucyate . \\
\noindent \textbf{(P\_8,2.106)} aicoḥ ubhayavivṛddhiprasaṅgāt idutoḥ plutavacanam . \\
\noindent \textbf{(P\_8,2.106)} aicoḥ ubhayaviṛddhiprasaṅgāt idutoḥ plutaḥ ucyate . \\
\noindent \textbf{(P\_8,2.106)} kim ucyate ubhayavivṛddhiprasaṅgāt iti yadā nityāḥ śabdāḥ nityeṣu ca śabdeṣu kūṭasthaiḥ avicālibhiḥ varṇaiḥ bhavitavyam anapāyopajanavikāribhiḥ . \\
\noindent \textbf{(P\_8,2.106)} na eṣaḥ doṣaḥ . \\
\noindent \textbf{(P\_8,2.106)} ubhayavivṛddhiprasaṅgāt iti na evam vijñāyate ubhayoḥ vivṛddhiḥ ubhayavivṛddhiḥ ubhayavivṛddhiprasaṅgāt iti . \\
\noindent \textbf{(P\_8,2.106)} katham tarhi . \\
\noindent \textbf{(P\_8,2.106)} ubhayoḥ vivṛddhiḥ asmin saḥ ayam ubhayavivṛddhiḥ ubhayavivṛddhiprasaṅgāt iti . \\
\noindent \textbf{(P\_8,2.106)} imau aicau samāhāravarṇau mātrā avarṇasya mātrā ivarṇovarṇayoḥ iti tayoḥ plutaḥ ucyamāne ubhayavivṛddhiḥ prāpnoti . \\
\noindent \textbf{(P\_8,2.106)} tat yathā . \\
\noindent \textbf{(P\_8,2.106)} abhivardhamānaḥ garbhaḥ sarvāṅgaparipūrṇaḥ vardhate . \\
\noindent \textbf{(P\_8,2.106)} asti prayojanam etat . \\
\noindent \textbf{(P\_8,2.106)} kim tarhi iti . \\
\noindent \textbf{(P\_8,2.106)} tatra ayatheṣṭaprasaṅgaḥ . tatra ayatheṣṭam prasajyeta . \\
\noindent \textbf{(P\_8,2.106)} caturmātraḥ plutaḥ prāpnoti . \\
\noindent \textbf{(P\_8,2.106)} siddham tu idutoḥ dīrghavacanāt . siddham etat . \\
\noindent \textbf{(P\_8,2.106)} katham . \\
\noindent \textbf{(P\_8,2.106)} idutoḥ dīrghaḥ bhavati iti vaktavyam . \\
\noindent \textbf{(P\_8,2.106)} tat etat katham kṛtvā siddham bhavati . \\
\noindent \textbf{(P\_8,2.106)} yadi samaḥ pravibhāgaḥ mātrā avarṇasya mātrā ivarṇovarṇayoḥ . \\
\noindent \textbf{(P\_8,2.106)} atha hi ardhamātrā avarṇasya adhyardhamātrā ivarṇovarṇayoḥ ardhatṛtīyamātraḥ prāpnoti . \\
\noindent \textbf{(P\_8,2.106)} atha hi adhyardhamātrā avarṇasya ardhamātrā ivarṇovarṇayoḥ ardhacaturthamātraḥ prāpnoti . \\
\noindent \textbf{(P\_8,2.106)} sūtram ca bhidyate . \\
\noindent \textbf{(P\_8,2.106)} yathānyāsam eva astu . \\
\noindent \textbf{(P\_8,2.106)} nanu ca uktam tatra ayatheṣṭaprasaṅgaḥ iti . \\
\noindent \textbf{(P\_8,2.106)} tatra sauryabhagavatā uktam aniṣṭijñaḥ vāḍavaḥ paṭhati . \\
\noindent \textbf{(P\_8,2.106)} iṣyate eva caturmātraḥ plutaḥ \\
\noindent \textbf{(P\_8,2.107.1)} ecaḥ plutavikāre padāntagrahaṇam . ecaḥ plutavikāre padāntagrahaṇam kartavyam iha mā bhūt : bhadram karoṣi gau3ḥ iti . \\
\noindent \textbf{(P\_8,2.107.1)} viṣayaparigaṇanam ca . \\
\noindent \textbf{(P\_8,2.107.1)} viṣayaparigaṇanam ca kartavyam . \\
\noindent \textbf{(P\_8,2.107.1)} praśnāntābhipūjitavicāryamāṇapratyabhivādayājyānteṣu iti vaktavyam . \\
\noindent \textbf{(P\_8,2.107.1)} praśnānta . \\
\noindent \textbf{(P\_8,2.107.1)} agama3ḥ pūrva3n grāma3n agnibhūta3i . \\
\noindent \textbf{(P\_8,2.107.1)} paṭa3u . \\
\noindent \textbf{(P\_8,2.107.1)} praśnānta . \\
\noindent \textbf{(P\_8,2.107.1)} abhipūjita . \\
\noindent \textbf{(P\_8,2.107.1)} siddhaḥ asi māṇavaka agnibhūta3i . \\
\noindent \textbf{(P\_8,2.107.1)} paṭa3u . \\
\noindent \textbf{(P\_8,2.107.1)} abhipūjita . \\
\noindent \textbf{(P\_8,2.107.1)} vicāryamāṇa . \\
\noindent \textbf{(P\_8,2.107.1)} hotavyam dīkṣitasya gṛha3i . \\
\noindent \textbf{(P\_8,2.107.1)} vicāryamāṇa . \\
\noindent \textbf{(P\_8,2.107.1)} pratyabhivāda . \\
\noindent \textbf{(P\_8,2.107.1)} āyuṣmān edhi agnibhūta3i . \\
\noindent \textbf{(P\_8,2.107.1)} pratyabhivāda . \\
\noindent \textbf{(P\_8,2.107.1)} yājyānta . \\
\noindent \textbf{(P\_8,2.107.1)} ukṣannāya vaśannāya somapṛṣṭhāya vedhase . \\
\noindent \textbf{(P\_8,2.107.1)} stomaiḥ vidhema agnaya3i \\
\noindent \textbf{(P\_8,2.107.2)} āmantrite chandasi upasaṅkhyānam . \\
\noindent \textbf{(P\_8,2.107.2)} āmantrite chandasi upasaṅkhyānam kartavyam . \\
\noindent \textbf{(P\_8,2.107.2)} agna3i patnīva3ḥ sajuḥ devena tvaṣṭrā somam piba \\
\noindent \textbf{(P\_8,2.108.1)} atha kayoḥ imau yvau ucyete . \\
\noindent \textbf{(P\_8,2.108.1)} idutoḥ iti āha . \\
\noindent \textbf{(P\_8,2.108.1)} tat idutoḥ grahaṇam kartavyam . \\
\noindent \textbf{(P\_8,2.108.1)} na kartavyam . \\
\noindent \textbf{(P\_8,2.108.1)} prkṛtam anuvartate . \\
\noindent \textbf{(P\_8,2.108.1)} kva prakṛtam . \\
\noindent \textbf{(P\_8,2.108.1)} pūrvasya ardhasya aduttarasya idutau iti . \\
\noindent \textbf{(P\_8,2.108.1)} tat vai prathamānirdiṣṭam ṣaṣṭhīnirdiṣṭena ca iha arthaḥ . \\
\noindent \textbf{(P\_8,2.108.1)} aci iti eṣā saptamī idutau iti prathamāyāḥ ṣaṣṭhīm prakalpayiṣyati tasminnitinirdiṣṭepūrvasya iti \\
\noindent \textbf{(P\_8,2.108.2)} kimartham idam ucyate na ikaḥ yaṇ aci iti eva siddham . \\
\noindent \textbf{(P\_8,2.108.2)} na sidhyati . \\
\noindent \textbf{(P\_8,2.108.2)} asiddhaḥ plutaḥ plutavikārau ca imau . \\
\noindent \textbf{(P\_8,2.108.2)} siddhaḥ plutaḥ svarasandhiṣu . \\
\noindent \textbf{(P\_8,2.108.2)} katham jñāyate . \\
\noindent \textbf{(P\_8,2.108.2)} yat ayam plutaḥ prakṛtyā iti plutasya prakṛtibhāvam śāsti . \\
\noindent \textbf{(P\_8,2.108.2)} katham kṛtvā jñāpakam . \\
\noindent \textbf{(P\_8,2.108.2)} sataḥ hi kāryiṇaḥ kāryeṇa bhavitavyam . \\
\noindent \textbf{(P\_8,2.108.2)} idam tarhi prayojanam dīrghaśākalapratiṣedhārtham . \\
\noindent \textbf{(P\_8,2.108.2)} dīrghatvam śākalam ca mā bhūt iti . \\
\noindent \textbf{(P\_8,2.108.2)} agnā3yindram . \\
\noindent \textbf{(P\_8,2.108.2)} paṭā3vudakam . \\
\noindent \textbf{(P\_8,2.108.2)} etat api na asti prayojanam . \\
\noindent \textbf{(P\_8,2.108.2)} ārabhyate plutapūrvasya yaṇādeśaḥ plutapurvasya dīrghaśākalapratiṣedhārtham iti . \\
\noindent \textbf{(P\_8,2.108.2)} tat na vaktavyam bhavati . \\
\noindent \textbf{(P\_8,2.108.2)} avaśyam tat vaktavyam yau plutapūrvau idutau aplutavikārau tadartham . \\
\noindent \textbf{(P\_8,2.108.2)} bho3yindra . \\
\noindent \textbf{(P\_8,2.108.2)} bho3yiha iti . \\
\noindent \textbf{(P\_8,2.108.2)} yadi tarhi tasya nibandhanam asti tat eva vaktavyam idam na vaktavyam . \\
\noindent \textbf{(P\_8,2.108.2)} idam api avaśyam vaktavyam svarārtham . \\
\noindent \textbf{(P\_8,2.108.2)} tena hi sati udāttasvaritayoryaṇaḥsvarito'nudāttasya iti eṣaḥ svaraḥ prasajyeta . \\
\noindent \textbf{(P\_8,2.108.2)} anena punaḥ sati asiddhatvāt na bhaviṣyati . \\
\noindent \textbf{(P\_8,2.108.2)} yadi tarhi asya nibandhanam asti idam eva vaktavyam tat na vaktavyam . \\
\noindent \textbf{(P\_8,2.108.2)} nanu ca uktam tat api avaśyam vaktavyam yau plutapūrvau idutau aplutavikārau tadartham bho3yindra bho3yiha iti . \\
\noindent \textbf{(P\_8,2.108.2)} chāndasam etat dṛṣṭānuvidhiḥ chandasi bhavati . \\
\noindent \textbf{(P\_8,2.108.2)} yat tarhi na chāndasam . \\
\noindent \textbf{(P\_8,2.108.2)} bho3yindram sāma gāyati . \\
\noindent \textbf{(P\_8,2.108.2)} eṣaḥ api chandasi dṛṣṭasya anuprayogaḥ iti . \\
\noindent \textbf{(P\_8,2.108.2)} kim nu yaṇā bhavati iha na siddham yvau idutoḥ yat ayam vidadhāti . \\
\noindent \textbf{(P\_8,2.108.2)} tau ca mama svarasandhiṣu siddhau śākaladīrghavidhī tu nivartyau .1. ik tu yadā bhavati plutapūrvaḥ tasya yaṇam vidadhāti apavādam . \\
\noindent \textbf{(P\_8,2.108.2)} tena tayoḥ ca na śākaladīrghau yaṇsvarabādhanam eva tu hetuḥ \\
\noindent \textbf{(P\_8,3.1.1)} matuvasaḥ rādeśe vanaḥ upasaṅkhyānam . \\
\noindent \textbf{(P\_8,3.1.1)} matuvasaḥ rādeśe vanaḥ upasaṅkhyānam kartavyam . \\
\noindent \textbf{(P\_8,3.1.1)} yaḥ tvā āyantam vasunā prātaritvaḥ . \\
\noindent \textbf{(P\_8,3.1.1)} ibhāṣā bhavadbhagavadaghavatām ot ca avasya . chandasi bhāṣāyām ca bhavat bhagavat aghavat iti eteṣām vibhāṣā ruḥ vaktavyaḥ ot ca avasya vaktavyaḥ . \\
\noindent \textbf{(P\_8,3.1.1)} bhoḥ , bhavan . \\
\noindent \textbf{(P\_8,3.1.1)} bhagoḥ , bhagavan . \\
\noindent \textbf{(P\_8,3.1.1)} aghoḥ , aghavan iti \\
\noindent \textbf{(P\_8,3.1.2)} sambuddhau iti ucyate tatra idam na sidhyati . \\
\noindent \textbf{(P\_8,3.1.2)} bhoḥ brāhmaṇāḥ iti . \\
\noindent \textbf{(P\_8,3.1.2)} tathā vibhaktau liṅgaviśiṣṭagrahaṇam na iti iha na prāpnoti . \\
\noindent \textbf{(P\_8,3.1.2)} bhoḥ brāhmaṇi . \\
\noindent \textbf{(P\_8,3.1.2)} na eṣaḥ doṣaḥ . \\
\noindent \textbf{(P\_8,3.1.2)} avyayam eṣaḥ bhoḥśabdaḥ na eṣā bhavataḥ pravṛttiḥ . \\
\noindent \textbf{(P\_8,3.1.2)} katham avyayatvam . \\
\noindent \textbf{(P\_8,3.1.2)} vibhaktisvarapratirūpakāḥ ca nipātāḥ bhavanti iti nipātasañjñā nipātaḥ avyayam iti avyayasañjñā \\
\noindent \textbf{(P\_8,3.5-6,} 12) KA\_III,424.11-416.8 Ro\_V,432-434 {1/21}     sampuṅkānām satvam . \\
\noindent \textbf{(P\_8,3.5-6,} 12) KA\_III,424.11-416.8 Ro\_V,432-434 {2/21}     sampuṅkānām satvam vaktavyam . \\
\noindent \textbf{(P\_8,3.5-6,} 12) KA\_III,424.11-416.8 Ro\_V,432-434 {3/21}     saṃskartā puṃskāmā kāṃs kān iti . \\
\noindent \textbf{(P\_8,3.5-6,} 12) KA\_III,424.11-416.8 Ro\_V,432-434 {4/21}     ruvidhau hi aniṣṭaprasaṅgaḥ . ruvidhau hi sati aniṣṭam prasajyeta . \\
\noindent \textbf{(P\_8,3.5-6,} 12) KA\_III,424.11-416.8 Ro\_V,432-434 {5/21}     iha tāvat saṃskartā iti vā śari iti prasajyeta . \\
\noindent \textbf{(P\_8,3.5-6,} 12) KA\_III,424.11-416.8 Ro\_V,432-434 {6/21}     puṃskāmā iti idudupadhasya iti ṣatvam prasajyeta . \\
\noindent \textbf{(P\_8,3.5-6,} 12) KA\_III,424.11-416.8 Ro\_V,432-434 {7/21}     kāṃs kān iti kupvoḥ hkaḥ prasajyeta . \\
\noindent \textbf{(P\_8,3.5-6,} 12) KA\_III,424.11-416.8 Ro\_V,432-434 {8/21}     tat tarhi vaktavyam . \\
\noindent \textbf{(P\_8,3.5-6,} 12) KA\_III,424.11-416.8 Ro\_V,432-434 {9/21}     na vaktavyam . \\
\noindent \textbf{(P\_8,3.5-6,} 12) KA\_III,424.11-416.8 Ro\_V,432-434 {10/21}     kriyate nyāse eva . \\
\noindent \textbf{(P\_8,3.5-6,} 12) KA\_III,424.11-416.8 Ro\_V,432-434 {11/21}     samaḥ suṭi iti dvisakārakaḥ nirdeśaḥ : samaḥ suṭi sakāraḥ bhavati . \\
\noindent \textbf{(P\_8,3.5-6,} 12) KA\_III,424.11-416.8 Ro\_V,432-434 {12/21}     tat prakṛtam uttaratra anuvartiṣyate . \\
\noindent \textbf{(P\_8,3.5-6,} 12) KA\_III,424.11-416.8 Ro\_V,432-434 {13/21}     yadi tat anuvartate naśchavyapraśān iti atra api prāpnoti . \\
\noindent \textbf{(P\_8,3.5-6,} 12) KA\_III,424.11-416.8 Ro\_V,432-434 {14/21}     sambandham anuvartiṣyate . \\
\noindent \textbf{(P\_8,3.5-6,} 12) KA\_III,424.11-416.8 Ro\_V,432-434 {15/21}     samaḥsuṭi . \\
\noindent \textbf{(P\_8,3.5-6,} 12) KA\_III,424.11-416.8 Ro\_V,432-434 {16/21}     pumaḥ khayi ampare saḥ bhavati . \\
\noindent \textbf{(P\_8,3.5-6,} 12) KA\_III,424.11-416.8 Ro\_V,432-434 {17/21}     naḥ chavi apraśān ruḥ bhavati pumaḥ khayi ampare sakāraḥ . \\
\noindent \textbf{(P\_8,3.5-6,} 12) KA\_III,424.11-416.8 Ro\_V,432-434 {18/21}     ubhayatharkṣu dīrghādaṭisamānapāde nṝnpe svatavānpāyau ruḥ bhavati pumaḥ khayi ampare sakāraḥ . \\
\noindent \textbf{(P\_8,3.5-6,} 12) KA\_III,424.11-416.8 Ro\_V,432-434 {19/21}     kān āmreḍite sakāraḥ . \\
\noindent \textbf{(P\_8,3.5-6,} 12) KA\_III,424.11-416.8 Ro\_V,432-434 {20/21}     pumaḥ khayi ampare iti nivṛttam . \\
\noindent \textbf{(P\_8,3.5-6,} 12) KA\_III,424.11-416.8 Ro\_V,432-434 {21/21}     samaḥ vā lopam eke icchanti : saṃskartā saṃskartā \\
\noindent \textbf{(P\_8,3.13)} ḍhalope apadāntagrahaṇam . \\
\noindent \textbf{(P\_8,3.13)} ḍhalope apadāntagrahaṇam kartavyam . \\
\noindent \textbf{(P\_8,3.13)} iha mā bhūt . \\
\noindent \textbf{(P\_8,3.13)} śvaliṭ ḍhaukate . \\
\noindent \textbf{(P\_8,3.13)} guḍaliṭ ḍhaukate . \\
\noindent \textbf{(P\_8,3.13)} tat tarhi vaktavyam . \\
\noindent \textbf{(P\_8,3.13)} na vaktavyam . \\
\noindent \textbf{(P\_8,3.13)} jaśtvam atra bādhakam bhaviṣyati . \\
\noindent \textbf{(P\_8,3.13)} jaśbhāvāt iti cet uttaratra ḍhasya abhāvāt apavādaprasaṅgaḥ . jaśbhāvāt iti cet uttaratra ḍhakārasya abhāvāt asiddhatvāt apavādaḥ ayam vijñāyate . \\
\noindent \textbf{(P\_8,3.13)} kasya . \\
\noindent \textbf{(P\_8,3.13)} jaśtvasya . \\
\noindent \textbf{(P\_8,3.13)} tasmāt siddhavacanam . tasmāt siddhatvam vaktavyam . \\
\noindent \textbf{(P\_8,3.13)} kasya . \\
\noindent \textbf{(P\_8,3.13)} ṣṭutvasya . \\
\noindent \textbf{(P\_8,3.13)} saṅgrahaṇam vā . saṅgrahaṇam vā kartavyam . \\
\noindent \textbf{(P\_8,3.13)} saṅi ḍhaḥ iti vaktavyam . \\
\noindent \textbf{(P\_8,3.13)} tattarhi vaktavyam . \\
\noindent \textbf{(P\_8,3.13)} na vaktavyam . \\
\noindent \textbf{(P\_8,3.13)} ānantaryam iha āśrīyate ḍhakārasya ḍhakāre iti . \\
\noindent \textbf{(P\_8,3.13)} kva cit ca sannipātakṛtam ānantaryam śāstrakṛtam anānantaryam kva cit ca na eva sannipātakṛtam na api śāstrakṛtam . \\
\noindent \textbf{(P\_8,3.13)} ṣṭutve sannipātakṛtam ānantaryam jaśtve na eva sannipātakṛtam na api śāstrakṛtam . \\
\noindent \textbf{(P\_8,3.13)} yatra kutaḥ cit eva ānantaryam tat āśrayiṣyāmaḥ \\
\noindent \textbf{(P\_8,3.15)} visarjanīyaḥ anuttarapade . \\
\noindent \textbf{(P\_8,3.15)} visarjanīyaḥ anuttarapade iti vaktavyam . \\
\noindent \textbf{(P\_8,3.15)} iha mā bhūt . \\
\noindent \textbf{(P\_8,3.15)} nārkuṭaḥ nārpatyaḥ iti . \\
\noindent \textbf{(P\_8,3.15)} na vā bahiraṅgalakṣaṇatvāt . na vā vaktavyam . \\
\noindent \textbf{(P\_8,3.15)} kim kāraṇam . \\
\noindent \textbf{(P\_8,3.15)} bahiraṅalakṣaṇatvāt . \\
\noindent \textbf{(P\_8,3.15)} bahiraṅgaḥ rephaḥ . \\
\noindent \textbf{(P\_8,3.15)} antaraṅgaḥ visarjanīyaḥ . \\
\noindent \textbf{(P\_8,3.15)} asiddham bahiraṅgam antaraṅge . \\
\noindent \textbf{(P\_8,3.15)} na eṣaḥ yuktaḥ parihāraḥ . \\
\noindent \textbf{(P\_8,3.15)} antaraṅgam bahiraṅgam iti pratidvandvabhāvinau etau pakṣau . \\
\noindent \textbf{(P\_8,3.15)} sati antaraṅge bahiraṅgam sati bahiraṅge antargam . \\
\noindent \textbf{(P\_8,3.15)} na ca atra antaraṅgabahiraṅgayoḥ yugapat samavasthānam asti . \\
\noindent \textbf{(P\_8,3.15)} kim kāraṇam . \\
\noindent \textbf{(P\_8,3.15)} asiddhatvāt . \\
\noindent \textbf{(P\_8,3.15)} na ca anabhinirvṛtte bahiraṅge antaraṅgam prāpnoti . \\
\noindent \textbf{(P\_8,3.15)} tatra nimittam eva bahiraṅgam antaraṅgasya . \\
\noindent \textbf{(P\_8,3.15)} animittam bahiraṅgam antaraṅgasya . \\
\noindent \textbf{(P\_8,3.15)} kim kāraṇam . \\
\noindent \textbf{(P\_8,3.15)} asiddhatvāt . \\
\noindent \textbf{(P\_8,3.15)} katham asiddhatvam yāvatā pūrvatra asiddham iti asiddhā paribhāṣā . \\
\noindent \textbf{(P\_8,3.15)} asiddham bahiraṅgam antaraṅge . \\
\noindent \textbf{(P\_8,3.15)} katham . \\
\noindent \textbf{(P\_8,3.15)} kāryakālam sañjñāparibhāṣam iti kharavasānayorvisarjanīyaḥ upasthitam idam bhavati asiddham bahiraṅgam antaraṅge iti . \\
\noindent \textbf{(P\_8,3.15)} evam eṣā siddhā paribhāṣā bhavati . \\
\noindent \textbf{(P\_8,3.15)} kutaḥ nu khalu etat dvayoḥ paribhāṣayoḥ sāvakāśayoḥ samavasthitayoḥ pūrvatra asiddham iti ca asiddham bahiraṅgam antaraṅge iti ca pūrvatrāsiddham iti etām upamṛdya asiddham bahiraṅgam antaraṅge iti etayā vyavasthā bhaviṣyati na punaḥ asiddham bahiraṅgam antaraṅge iti etām upamṛdya purvatrāsiddham iti etayā vyavasthā syāt . \\
\noindent \textbf{(P\_8,3.15)} ataḥ kim . \\
\noindent \textbf{(P\_8,3.15)} ataḥ ayuktaḥ parihāraḥ na vā bahiraṅgalakṣaṇatvāt iti \\
\noindent \textbf{(P\_8,3.16)} kimartham idam ucyate na kharavasānayoḥ visarjanīyaḥ iti eva siddham . \\
\noindent \textbf{(P\_8,3.16)} niyamārthaḥ ayam ārambhaḥ . \\
\noindent \textbf{(P\_8,3.16)} roḥ eva supi na anyasya supi . \\
\noindent \textbf{(P\_8,3.16)} kva mā bhūt . \\
\noindent \textbf{(P\_8,3.16)} gīrṣu dhūrṣu \\
\noindent \textbf{(P\_8,3.17)} aśgrahaṇam anarthakam anyatra abhāvāt . \\
\noindent \textbf{(P\_8,3.17)} aśgrahaṇam anarthakam . \\
\noindent \textbf{(P\_8,3.17)} kim kāraṇam . \\
\noindent \textbf{(P\_8,3.17)} anyatra abhāvāt . \\
\noindent \textbf{(P\_8,3.17)} na hi anyatra ruḥ asti anyat ataḥ aśaḥ . \\
\noindent \textbf{(P\_8,3.17)} nanu ca ayam asti. chandaḥsu payaḥsu iti . \\
\noindent \textbf{(P\_8,3.17)} kim punaḥ kāraṇam sukāraparaḥ eva udāhriyate na punaḥ ayam vṛkṣaḥ tatra plakṣaḥ tatra iti . \\
\noindent \textbf{(P\_8,3.17)} asti atra viśeṣaḥ . \\
\noindent \textbf{(P\_8,3.17)} visarjanīye kṛte na bhaviṣyati . \\
\noindent \textbf{(P\_8,3.17)} iha api tarhi visarjanīye kṛte na bhaviṣyati . \\
\noindent \textbf{(P\_8,3.17)} chandaḥsu payaḥsviti . \\
\noindent \textbf{(P\_8,3.17)} sthānivadbhāvāt prāpnoti . \\
\noindent \textbf{(P\_8,3.17)} nanu ca iha api sthānivadbhāvāt prāpnoti . \\
\noindent \textbf{(P\_8,3.17)} vṛkṣaḥ tatra plakṣaḥ tatra iti . \\
\noindent \textbf{(P\_8,3.17)} analvidhau sthānivadbhāvaḥ . \\
\noindent \textbf{(P\_8,3.17)} atha ayam alvidhiḥ syāt śakyam aśgrahaṇam avaktum . \\
\noindent \textbf{(P\_8,3.17)} bāḍham śakyam . \\
\noindent \textbf{(P\_8,3.17)} alvidhiḥ tarhi bhaviṣyati . \\
\noindent \textbf{(P\_8,3.17)} katham . \\
\noindent \textbf{(P\_8,3.17)} idam asti rori iti . \\
\noindent \textbf{(P\_8,3.17)} tataḥ vakṣyāmi kharavasānayoḥ visarjanīyaḥ raḥ . \\
\noindent \textbf{(P\_8,3.17)} tataḥ roḥ supi visarjanīyaḥ raḥ iti eva . \\
\noindent \textbf{(P\_8,3.17)} uttarārtham tarhi aśgrahaṇam kartavyam halisarveṣām hali aśi iti yathā syāt . \\
\noindent \textbf{(P\_8,3.17)} iha mā bhūt : vṛkṣavayateḥ apratyayaḥ vṛkṣav karoti \\
\noindent \textbf{(P\_8,3.20)} kimartham idam ucyate na lopaḥ śākalyasya iti eva siddham . \\
\noindent \textbf{(P\_8,3.20)} okārāt lopavacanam nityārtham . \\
\noindent \textbf{(P\_8,3.20)} okārāt lopavacanam kriyate . \\
\noindent \textbf{(P\_8,3.20)} nityārthaḥ ayam ārambhaḥ \\
\noindent \textbf{(P\_8,3.21)} pade iti kimartham . \\
\noindent \textbf{(P\_8,3.21)} tantre , utam , tantrayutam , tantra*utam . \\
\noindent \textbf{(P\_8,3.21)} pade iti śakyam avaktum . \\
\noindent \textbf{(P\_8,3.21)} kasmāt na bhavati tantre , utam , tantrayutam , tantra*utam iti . \\
\noindent \textbf{(P\_8,3.21)} lakṣaṇapratipadoktayoḥ pratipadoktasya eva iti . \\
\noindent \textbf{(P\_8,3.21)} uttarārtham tarhi padagrahaṇam kartavyam ṅamohrasvādaciṅamuṇnityam iti apade mā bhūt . \\
\noindent \textbf{(P\_8,3.21)} daṇḍinā śakaṭinā \\
\noindent \textbf{(P\_8,3.26)} yavalapare yavalāḥ vā . \\
\noindent \textbf{(P\_8,3.26)} yavalapare hakāre yavalāḥ vā iti vaktavyam . \\
\noindent \textbf{(P\_8,3.26)} kiyhyaḥ kim hyaḥ . \\
\noindent \textbf{(P\_8,3.26)} kivhvalayati kim hvalayati . \\
\noindent \textbf{(P\_8,3.26)} kilhlādayati kim hlādayati \\
\noindent \textbf{(P\_8,3.28-32)} KA\_III,428.5-429.6 Ro\_V,443-444 {1/30}     iha dhuḍādiṣu kecit pūrvāntāḥ kecit parādayaḥ . \\
\noindent \textbf{(P\_8,3.28-32)} KA\_III,428.5-429.6 Ro\_V,443-444 {2/30}     yadi punaḥ sarve eva pūrvāntāḥ syuḥ sarve eva parādayaḥ . \\
\noindent \textbf{(P\_8,3.28-32)} KA\_III,428.5-429.6 Ro\_V,443-444 {3/30}     kaḥ ca atra viśeṣaḥ . \\
\noindent \textbf{(P\_8,3.28-32)} KA\_III,428.5-429.6 Ro\_V,443-444 {4/30}     dhugādiṣu ṣṭutvaṇatvapratiṣedhaḥ . \\
\noindent \textbf{(P\_8,3.28-32)} KA\_III,428.5-429.6 Ro\_V,443-444 {5/30}     dhugādiṣu satsu ṣṭutvaṇatvayoḥ pratiṣedhaḥ vaktavyaḥ . \\
\noindent \textbf{(P\_8,3.28-32)} KA\_III,428.5-429.6 Ro\_V,443-444 {6/30}     ṣṭutvasya tāvat . \\
\noindent \textbf{(P\_8,3.28-32)} KA\_III,428.5-429.6 Ro\_V,443-444 {7/30}     śvaliṭtsāye . \\
\noindent \textbf{(P\_8,3.28-32)} KA\_III,428.5-429.6 Ro\_V,443-444 {8/30}     madhuliṭtsāye . \\
\noindent \textbf{(P\_8,3.28-32)} KA\_III,428.5-429.6 Ro\_V,443-444 {9/30}     ṣṭunāṣṭuḥ iti ṣṭutvam prāpnoti . \\
\noindent \textbf{(P\_8,3.28-32)} KA\_III,428.5-429.6 Ro\_V,443-444 {10/30}     parādau punaḥ sati napadāntāṭṭoranām iti pratiṣedhaḥ siddhaḥ bhavati . \\
\noindent \textbf{(P\_8,3.28-32)} KA\_III,428.5-429.6 Ro\_V,443-444 {11/30}     ṇatvasya . \\
\noindent \textbf{(P\_8,3.28-32)} KA\_III,428.5-429.6 Ro\_V,443-444 {12/30}     kurvannāste . \\
\noindent \textbf{(P\_8,3.28-32)} KA\_III,428.5-429.6 Ro\_V,443-444 {13/30}     kṛṣannāste . \\
\noindent \textbf{(P\_8,3.28-32)} KA\_III,428.5-429.6 Ro\_V,443-444 {14/30}     raṣābhyānnoṇaḥsamānapade it ṇatvam prāpnoti . \\
\noindent \textbf{(P\_8,3.28-32)} KA\_III,428.5-429.6 Ro\_V,443-444 {15/30}     parādau punaḥ sati padāntasya na iti pratiṣedhaḥ siddhaḥ bhavati . \\
\noindent \textbf{(P\_8,3.28-32)} KA\_III,428.5-429.6 Ro\_V,443-444 {16/30}     santu tarhi parādayaḥ . \\
\noindent \textbf{(P\_8,3.28-32)} KA\_III,428.5-429.6 Ro\_V,443-444 {17/30}     parādau chatvaṣatvavidhipratiṣedhaḥ . yadi parādayaḥ chatvam vidheyam ṣatvam ca pratiṣedhyam . \\
\noindent \textbf{(P\_8,3.28-32)} KA\_III,428.5-429.6 Ro\_V,443-444 {18/30}     chatvam vidheyam . \\
\noindent \textbf{(P\_8,3.28-32)} KA\_III,428.5-429.6 Ro\_V,443-444 {19/30}     kurvañcchete . \\
\noindent \textbf{(P\_8,3.28-32)} KA\_III,428.5-429.6 Ro\_V,443-444 {20/30}     kṛṣañcchete . \\
\noindent \textbf{(P\_8,3.28-32)} KA\_III,428.5-429.6 Ro\_V,443-444 {21/30}     yat hi tat śaścho'ṭi iti jhayaḥ padāntāt iti evam tat . \\
\noindent \textbf{(P\_8,3.28-32)} KA\_III,428.5-429.6 Ro\_V,443-444 {22/30}     kim punaḥ kāraṇam jhayaḥ padāntāt iti evam tat . \\
\noindent \textbf{(P\_8,3.28-32)} KA\_III,428.5-429.6 Ro\_V,443-444 {23/30}     iha mā bhūt . \\
\noindent \textbf{(P\_8,3.28-32)} KA\_III,428.5-429.6 Ro\_V,443-444 {24/30}     purā krūrasya visṛpaḥ virapśin iti . \\
\noindent \textbf{(P\_8,3.28-32)} KA\_III,428.5-429.6 Ro\_V,443-444 {25/30}     ṣatvam ca pratiṣedhyam . \\
\noindent \textbf{(P\_8,3.28-32)} KA\_III,428.5-429.6 Ro\_V,443-444 {26/30}     pratyaṅksiñca . \\
\noindent \textbf{(P\_8,3.28-32)} KA\_III,428.5-429.6 Ro\_V,443-444 {27/30}     udaṅksiñca . \\
\noindent \textbf{(P\_8,3.28-32)} KA\_III,428.5-429.6 Ro\_V,443-444 {28/30}     ādeśapratyayayoḥ iti ṣatvam prāpnoti . \\
\noindent \textbf{(P\_8,3.28-32)} KA\_III,428.5-429.6 Ro\_V,443-444 {29/30}     pūrvānte punaḥ sati sātpadādyoḥ iti pratiṣedhaḥ siddhaḥ bhavati . \\
\noindent \textbf{(P\_8,3.28-32)} KA\_III,428.5-429.6 Ro\_V,443-444 {30/30}     tasmāt santu yathānyāsam eva kecit pūrvāntāḥ kecit parādayaḥ . \\
\noindent \textbf{(P\_8,3.32.1)} ayam tu khalu śi tuk chatvārtham niyogataḥ pūrvāntaḥ kartavyaḥ tatra kurvañcchete kṛṣañcchete iti raṣābhyānnoṇaḥsamānapade iti ṇatvam prāpnoti . \\
\noindent \textbf{(P\_8,3.32.1)} na eṣaḥ doṣaḥ . \\
\noindent \textbf{(P\_8,3.32.1)} ścutve yogavibhāgaḥ kariṣyate . \\
\noindent \textbf{(P\_8,3.32.1)} idam asti kṣubhnādiṣu na ṇakāraḥ bhavati . \\
\noindent \textbf{(P\_8,3.32.1)} tataḥ stoḥ ścunā . \\
\noindent \textbf{(P\_8,3.32.1)} stoḥ ścunā sannipāte na ṇakāraḥ bhavati . \\
\noindent \textbf{(P\_8,3.32.1)} tataḥ ścuḥ . \\
\noindent \textbf{(P\_8,3.32.1)} ścuḥ ca bhavati stoḥ ścunā sannipāte \\
\noindent \textbf{(P\_8,3.32.2)} ṅamuṭi padādigrahaṇam . \\
\noindent \textbf{(P\_8,3.32.2)} ṅamuṭi padādigrahaṇam kartavyam . \\
\noindent \textbf{(P\_8,3.32.2)} iha mā bhūt . \\
\noindent \textbf{(P\_8,3.32.2)} daṇḍinā śakaṭinā iti . \\
\noindent \textbf{(P\_8,3.32.2)} tat tarhi vaktavyam . \\
\noindent \textbf{(P\_8,3.32.2)} na vaktavyam . \\
\noindent \textbf{(P\_8,3.32.2)} padāt iti vartate . \\
\noindent \textbf{(P\_8,3.32.2)} evam api paramadaṇḍinā paramacchattriṇā iti prāpnoti . \\
\noindent \textbf{(P\_8,3.32.2)} na eṣaḥ doṣaḥ . \\
\noindent \textbf{(P\_8,3.32.2)} uktam etat uttarapadatve ca apadādividhau lumatā lupte pratyayalakṣaṇam na bhavati iti . \\
\noindent \textbf{(P\_8,3.32.2)} evam api padāt iti vaktavyam . \\
\noindent \textbf{(P\_8,3.32.2)} yat hi tat prakṛtam prāk supi kutsanāt iti evam tat . \\
\noindent \textbf{(P\_8,3.32.2)} evam tarhi ṅamaḥ eva ayam ṅamuṭ kriyate . \\
\noindent \textbf{(P\_8,3.32.2)} katham . \\
\noindent \textbf{(P\_8,3.32.2)} padasya iti vartate ṅamaḥ iti ca na eṣā pañcamī . \\
\noindent \textbf{(P\_8,3.32.2)} kā tarhi . \\
\noindent \textbf{(P\_8,3.32.2)} sambandhaṣaṣṭhī . \\
\noindent \textbf{(P\_8,3.32.2)} padāntasya ṅamaḥ ṅamuṭ bhavati hrasvāt uttarasya aci iti . \\
\noindent \textbf{(P\_8,3.32.2)} yadi ṅamaḥ eva ṅamuṭ kriyate kurvannāste kṛṣannāste raṣābhyānnoṇaḥsamānapade iti ṇatvam prāpnoti . \\
\noindent \textbf{(P\_8,3.32.2)} padāntasya na iti pratiṣedhaḥ bhaviṣyati . \\
\noindent \textbf{(P\_8,3.32.2)} padāntasya iti ucyate na eṣaḥ padāntaḥ . \\
\noindent \textbf{(P\_8,3.32.2)} padāntabhaktaḥ padāntagrahaṇena grāhiṣyate . \\
\noindent \textbf{(P\_8,3.32.2)} evam api na sidhyati . \\
\noindent \textbf{(P\_8,3.32.2)} kim kāraṇam . \\
\noindent \textbf{(P\_8,3.32.2)} uktam etat na vā padādhikārasya viśeṣaṇatvāt iti . \\
\noindent \textbf{(P\_8,3.32.2)} evam tarhi pade iti vartate . \\
\noindent \textbf{(P\_8,3.32.2)} kva prakṛtam . \\
\noindent \textbf{(P\_8,3.32.2)} uñi ca pade iti \\
\noindent \textbf{(P\_8,3.33)} kimartham mayaḥ uttarasya uñaḥ vaḥ vā iti ucyate na ikaḥ yaṇaci iti eva siddham . \\
\noindent \textbf{(P\_8,3.33)} na sidhyati . \\
\noindent \textbf{(P\_8,3.33)} pragṛhyaḥ prakṛtyā iti prakṛtibhāvaḥ prāpnoti . \\
\noindent \textbf{(P\_8,3.33)} yadi punaḥ tatra eva ucyeta ikaḥ yaṇaci mayaḥ uñaḥ vā iti . \\
\noindent \textbf{(P\_8,3.33)} na evam śakyam . \\
\noindent \textbf{(P\_8,3.33)} iha hi doṣaḥ syāt . \\
\noindent \textbf{(P\_8,3.33)} kimvāvapanam mahat . \\
\noindent \textbf{(P\_8,3.33)} maḥ anusvāraḥ hali iti anusvāraḥ prasajyeta . \\
\noindent \textbf{(P\_8,3.33)} vatve punaḥ sati asiddhatvāt na bhaviṣyati \\
\noindent \textbf{(P\_8,3.34)} iha kasmāt na bhavati : vṛkṣaḥ , plakṣaḥ iti . \\
\noindent \textbf{(P\_8,3.34)} saṃhitāyām iti vartate . \\
\noindent \textbf{(P\_8,3.34)} evam api atra prāpnoti . \\
\noindent \textbf{(P\_8,3.34)} kim kāraṇam . \\
\noindent \textbf{(P\_8,3.34)} paraḥ sannikarṣaḥ saṃhitā iti ucyate . \\
\noindent \textbf{(P\_8,3.34)} saḥ yathā eva pareṇa paraḥ sannikarṣaḥ evam pūrveṇa api . \\
\noindent \textbf{(P\_8,3.34)} evam tarhi anavakāśā avasānasañjñā saṃhitāsañjñām bādhiṣyate . \\
\noindent \textbf{(P\_8,3.34)} atha vā saṃhitāsañjñāyām prakarṣagatiḥ vijñāsyate : sādhīyaḥ yaḥ paraḥ sannikarṣaḥ iti . \\
\noindent \textbf{(P\_8,3.34)} kaḥ ca sādhīyaḥ . \\
\noindent \textbf{(P\_8,3.34)} yaḥ pūrvaparayoḥ . \\
\noindent \textbf{(P\_8,3.34)} yadi eva anavakāśā avasānasañjñā saṃhitāsañjñām bādhate atha api saṃhitāsañjñāyām prakarṣagatiḥ vijñāyate ubhayathā doṣaḥ bhavati . \\
\noindent \textbf{(P\_8,3.34)} iṣyante itaḥ uttaram avasāne saṃhitākāryāṇi tāni na sidhyanti . \\
\noindent \textbf{(P\_8,3.34)} aṇaḥ apragṛhyasya anunāsikaḥ iti . \\
\noindent \textbf{(P\_8,3.34)} evam tarhi ācāryapravṛttiḥ jñāpayati na sarvasya visarjanīyasya satvam bhavati iti yat ayam kharavasānayorvisjanīyaḥ iti āha . \\
\noindent \textbf{(P\_8,3.34)} itarathā kharavasānayoḥ saḥ bhavati iti eva brūyāt . \\
\noindent \textbf{(P\_8,3.34)} tat ca laghu bhavati visarjanīyasya saḥ iti etat ca na vaktavyaṃ bhavati . \\
\noindent \textbf{(P\_8,3.34)} avaśyam śarparevisarjanīyaḥ iti atra prakṛtinirdeśārtham visarjanīyagrahaṇam kartavyam . \\
\noindent \textbf{(P\_8,3.34)} atha idānīm etat api rasānnidhyārtham purastāt apakrakṣyate kharavasānayoḥ saḥ iti atra eva evam api kupvoḥ XkkaXppau ca iti evamādinā anukramaṇena vyavacchinnam bhobhagoaghoapūrvasyayo'śi iti atra rugrahaṇam kartavyam syāt . \\
\noindent \textbf{(P\_8,3.34)} evam api ekam visarjanīyagrahaṇam vyājaḥ bhavati . \\
\noindent \textbf{(P\_8,3.34)} saḥ ayam evam laghīyasā nyāsena siddhe sati yat garīyāṃsam yatnam ārabhate tat jñāpayati ācāryaḥ na sarvasya visarjanīyasya satvam bhavati iti . \\
\noindent \textbf{(P\_8,3.34)} evam api anaikāntikam jñāpakam . \\
\noindent \textbf{(P\_8,3.34)} etāvat jñāpyate na sarvasya vijanīyasya satvam bhavati iti tatra kutaḥ etat iha bhaviṣyati vṛkṣaḥ tatra plakṣaḥ tatra iti iha na bhaviṣyati vṛkṣaḥ plakṣaḥ iti . \\
\noindent \textbf{(P\_8,3.34)} evam tarhi ācāryapravṛttiḥ jñāpayati na asya visarjanīyasya satvam bhavati iti yat ayam śarpare visarjanīyaḥ iti āha . \\
\noindent \textbf{(P\_8,3.34)} atha vā hali iti vartate . \\
\noindent \textbf{(P\_8,3.34)} kva prakṛtam . \\
\noindent \textbf{(P\_8,3.34)} hali sarveṣām iti . \\
\noindent \textbf{(P\_8,3.34)} yadi tat anuvartate mayaḥ uñaḥ vaḥ vā hali ca iti hali api vatvam prāpnoti . \\
\noindent \textbf{(P\_8,3.34)} śamu naḥ . \\
\noindent \textbf{(P\_8,3.34)} śamu yoḥ astu . \\
\noindent \textbf{(P\_8,3.34)} evam tarhi visarjanīyasyasaḥ iti atra khari iti anuvartiṣyate . \\
\noindent \textbf{(P\_8,3.34)} atha vā sambandham anuvartiṣyate \\
\noindent \textbf{(P\_8,3.36)} vāśarprakaraṇe kharpare lopaḥ . \\
\noindent \textbf{(P\_8,3.36)} vāśarprakaraṇe kharpare lopaḥ vaktavyaḥ . \\
\noindent \textbf{(P\_8,3.36)} vṛkṣāḥ sthātāraḥ . \\
\noindent \textbf{(P\_8,3.36)} vṛkṣāḥ sthātāraḥ \\
\noindent \textbf{(P\_8,3.37)} sasya kupvoḥ visarjanīyajihvāmūlīyopadhmānīyāḥ . \\
\noindent \textbf{(P\_8,3.37)} sasya kupvoḥ visarjanīyajihvāmūlīyopadhmānīyāḥ vaktavyāḥ . \\
\noindent \textbf{(P\_8,3.37)} visarjanīyādeśe hi śarparayoḥ eva ādeśaprasaṅgaḥ . visarjanīyādeśe hi sati śarparayoḥ eva kupvoḥ hkkahppau syātām . \\
\noindent \textbf{(P\_8,3.37)} adbhiḥ psātam . \\
\noindent \textbf{(P\_8,3.37)} vāsaḥ kṣaumam . \\
\noindent \textbf{(P\_8,3.37)} vacanāt na bhaviṣyataḥ . \\
\noindent \textbf{(P\_8,3.37)} asti vacane prayojanam . \\
\noindent \textbf{(P\_8,3.37)} kim . \\
\noindent \textbf{(P\_8,3.37)} puruṣaḥ tsarukaḥ . \\
\noindent \textbf{(P\_8,3.37)} tat tarhi vaktavyam . \\
\noindent \textbf{(P\_8,3.37)} na vaktavyam . \\
\noindent \textbf{(P\_8,3.37)} yat etat visarjanīyasya saḥ iti atra visarjanīyagrahaṇam etat uttaratra anuvartiṣyate tasmin ca śarpare visarjanīyaḥ asiddhaḥ . \\
\noindent \textbf{(P\_8,3.37)} na asiddhaḥ . \\
\noindent \textbf{(P\_8,3.37)} katham . \\
\noindent \textbf{(P\_8,3.37)} adhikāraḥ nāma triprakāraḥ . \\
\noindent \textbf{(P\_8,3.37)} kaḥ cit ekadeśasthaḥ sarvam śāstram abhijvalayati yathā pradīpaḥ suprajvalitaḥ sarvam veśma abhijvalayati . \\
\noindent \textbf{(P\_8,3.37)} aparaḥ yathā rajjvā ayasā vā baddham kāṣṭham anukṛṣyate tadvat . \\
\noindent \textbf{(P\_8,3.37)} aparaḥ adhikāraḥ pratiyogam tasya anirdeśārthaḥ iti yoge yoge upatiṣṭhate . \\
\noindent \textbf{(P\_8,3.37)} tat yadā eṣaḥ pakṣaḥ adhikāraḥ pratiyogam tasya anirdeśārthaḥ iti tadā hi yat etat visarjanīyasyasaḥ iti atra visarjanīyagrahaṇam etat uttaratra anuvṛttam sat anyat sampadyate tasmin ca śarpare visarjanīyaḥ siddhaḥ . \\
\noindent \textbf{(P\_8,3.37)} evam ca krtvā śarparayoḥ eva kupvoḥ hkkahppau syātām . \\
\noindent \textbf{(P\_8,3.37)} adbhiḥ psātam . \\
\noindent \textbf{(P\_8,3.37)} vāsaḥ kṣaumam iti . \\
\noindent \textbf{(P\_8,3.37)} evam tarhi yogavibhāgaḥ kariṣyate . \\
\noindent \textbf{(P\_8,3.37)} śarpare visarjanīyaḥ . \\
\noindent \textbf{(P\_8,3.37)} vā śari . \\
\noindent \textbf{(P\_8,3.37)} tataḥ kupvoḥ . \\
\noindent \textbf{(P\_8,3.37)} kupvoḥ ca śarparayoḥ visarjanīyasya visarjanīyaḥ bhavati iti . \\
\noindent \textbf{(P\_8,3.37)} kimartham idam . \\
\noindent \textbf{(P\_8,3.37)} kupvoḥ hkkahppau vakṣyati tadbādhanārtham . \\
\noindent \textbf{(P\_8,3.37)} tataḥ hkkahppau bhavataḥ kupvoḥ iti eva . \\
\noindent \textbf{(P\_8,3.37)} śarparayoḥ iti nivṛttam . \\
\noindent \textbf{(P\_8,3.37)} atha vā śarparevisarjanīyaḥ iti etat kupvoḥ hkkahppau ca iti atra anuvartiṣyate \\
\noindent \textbf{(P\_8,3.38)} saḥ apadādau anavyayasya . \\
\noindent \textbf{(P\_8,3.38)} saḥ apadādau anavyayasya iti vaktavyam . \\
\noindent \textbf{(P\_8,3.38)} iha mā bhūt . \\
\noindent \textbf{(P\_8,3.38)} prātaḥkalpam punaḥkalpam . \\
\noindent \textbf{(P\_8,3.38)} roḥ kāmye niyamārtham . roḥ kāmye iti vaktavyam . \\
\noindent \textbf{(P\_8,3.38)} kim prayojanam . \\
\noindent \textbf{(P\_8,3.38)} niyamārtham . \\
\noindent \textbf{(P\_8,3.38)} roḥ eva kāmye na anyasya . \\
\noindent \textbf{(P\_8,3.38)} payaskāmyati . \\
\noindent \textbf{(P\_8,3.38)} kva mā bhūt . \\
\noindent \textbf{(P\_8,3.38)} gīḥkāmyati pūḥkāmyati . \\
\noindent \textbf{(P\_8,3.38)} upadhmānīyasya ca satvam vaktavyam . \\
\noindent \textbf{(P\_8,3.38)} kim prayojanam . \\
\noindent \textbf{(P\_8,3.38)} ayam ubjiḥ upadhmānīyopadhaḥ paṭhyate tasya satve kṛte jaśbhāve ca abhyudgaḥ samudgaḥ iti etat rūpam yathā syāt . \\
\noindent \textbf{(P\_8,3.38)} yadi upadhmānīyopadhaḥ paṭhyate ubjijiṣati iti upadhmānīyasya dvirvacanam prāpnoti . \\
\noindent \textbf{(P\_8,3.38)} dakāropadhe punaḥ sati nandrāḥsaṃyogādayaḥ iti pratiṣedhaḥ siddhaḥ bhavati . \\
\noindent \textbf{(P\_8,3.38)} yadi dakāropadhaḥ paṭhyate kā rūpasiddhiḥ : ubjitā ubjitum iti . \\
\noindent \textbf{(P\_8,3.38)} asiddhe bhaḥ udjeḥ . \\
\noindent \textbf{(P\_8,3.38)} idam asti stoḥścunāścuḥ . \\
\noindent \textbf{(P\_8,3.38)} tataḥ vakṣyāmi . \\
\noindent \textbf{(P\_8,3.38)} bhaḥ udjeḥ . \\
\noindent \textbf{(P\_8,3.38)} udjeḥ ca ścunā sannipāte bhaḥ bhavati iti . \\
\noindent \textbf{(P\_8,3.38)} tat tarhi vaktavyam . \\
\noindent \textbf{(P\_8,3.38)} na vaktavyam . \\
\noindent \textbf{(P\_8,3.38)} nipātanāt etat siddham . \\
\noindent \textbf{(P\_8,3.38)} kim nipātanam . \\
\noindent \textbf{(P\_8,3.38)} bhujanyubjaupāṇyupatāpayoḥ iti . \\
\noindent \textbf{(P\_8,3.38)} iha api prāpnoti . \\
\noindent \textbf{(P\_8,3.38)} abhyudgaḥ samudgaḥ . \\
\noindent \textbf{(P\_8,3.38)} akutvaviṣaye nipātanam . \\
\noindent \textbf{(P\_8,3.38)} atha vā na etat ubjeḥ rūpam . \\
\noindent \textbf{(P\_8,3.38)} kim tarhi gameḥ dvyupasargāt ḍaḥ vidhīyate . \\
\noindent \textbf{(P\_8,3.38)} abhyudgataḥ abhyudgaḥ . \\
\noindent \textbf{(P\_8,3.38)} samudgataḥ samudgaḥ \\
\noindent \textbf{(P\_8,3.39)} kim aviśeṣeṇa satvam uktvā tataḥ iṇaḥ uttarasya sakārasya ṣatvam ucyate āhosvit iṇaḥ uttarasya visarjanīyasya eva ṣatvam vidhīyate . \\
\noindent \textbf{(P\_8,3.39)} kim ca ataḥ . \\
\noindent \textbf{(P\_8,3.39)} yadi aviśeṣeṇa satvam uktvā iṇaḥ uttarasya sakārasya ṣatvam ucyate niṣkṛtam , niṣpītam iti atra satvasya asiddhatvāt ṣatvam na prāpnoti . \\
\noindent \textbf{(P\_8,3.39)} atha iṇaḥ uttarasya visarjanīyasya eva ṣatvam vidhīyate satvam api anuvartate utāho na . \\
\noindent \textbf{(P\_8,3.39)} kim ca ataḥ . \\
\noindent \textbf{(P\_8,3.39)} yadi anuvartate satvam api prāpnoti . \\
\noindent \textbf{(P\_8,3.39)} atha nivṛttam namaspurasorgatyoḥ iti atra sakāragrahaṇam kartavyam . \\
\noindent \textbf{(P\_8,3.39)} tasmin ca kriyamāṇe ṣatvam api anuvartate utāho na . \\
\noindent \textbf{(P\_8,3.39)} kim ca ataḥ . \\
\noindent \textbf{(P\_8,3.39)} yadi anuvartate ṣatvam api prāpnoti . \\
\noindent \textbf{(P\_8,3.39)} atha nivṛttam idudupadhasya ca apratyasya iti atra ṣakāragrahaṇam kartavyam . \\
\noindent \textbf{(P\_8,3.39)} tasmin ca kriyamāṇe satvam api anuvartate utāho na . \\
\noindent \textbf{(P\_8,3.39)} kim ca ataḥ . \\
\noindent \textbf{(P\_8,3.39)} yadi anuvartate satvam api prāpnoti . \\
\noindent \textbf{(P\_8,3.39)} atha nivṛttam tirasaḥ anyatarasyām iti atra sakāragrahaṇam kartavyam . \\
\noindent \textbf{(P\_8,3.39)} tasmin ca kriyamāṇe ṣatvam api anuvartate utāho na . \\
\noindent \textbf{(P\_8,3.39)} kim ca ataḥ . \\
\noindent \textbf{(P\_8,3.39)} yadi anuvartate ṣatvam api prāpnoti . \\
\noindent \textbf{(P\_8,3.39)} atha nivṛttam dvistriścaturitikṛtvo'rthe isusoḥ sāmarthye nityaṃsamāse anuttarapadasthasya iti ṣakāragrahaṇam kartavyam . \\
\noindent \textbf{(P\_8,3.39)} tasmin ca kriyamāṇe satvam api anuvartate utāho na . \\
\noindent \textbf{(P\_8,3.39)} kim ca ataḥ . \\
\noindent \textbf{(P\_8,3.39)} yadi anuvartate satvam api prāpnoti . \\
\noindent \textbf{(P\_8,3.39)} atha nivṛttam ataḥkṛkamikaṃsakumbhapātrakuśākarṇīṣvanavyayasya iti sakāragrahaṇam kartavyam . \\
\noindent \textbf{(P\_8,3.39)} yathā icchasi tathā astu . \\
\noindent \textbf{(P\_8,3.39)} astu tāvat aviśeṣeṇa satvam uktvā iṇaḥ uttarasya sakārasya ṣatvam ucyate . \\
\noindent \textbf{(P\_8,3.39)} nanu ca uktam niṣkṛtam , niṣpītam iti atra satvasya asiddhatvāt ṣatvam na prāpnoti iti . \\
\noindent \textbf{(P\_8,3.39)} na eṣaḥ doṣaḥ . \\
\noindent \textbf{(P\_8,3.39)} ācāryapravṛttiḥ jñāpayati na yoge yogaḥ asiddhaḥ . \\
\noindent \textbf{(P\_8,3.39)} kim tarhi prakaraṇe prakaraṇam asiddham iti yat ayam upasargāt asamāse api ṇopadeśasya iti asamāsepigrahaṇam karoti . \\
\noindent \textbf{(P\_8,3.39)} atha vā punaḥ astu iṇaḥ uttarasya visarjanīyasya ṣatvam vidhīyate . \\
\noindent \textbf{(P\_8,3.39)} nanu ca uktam satvam api anuvartate utāho na kim ca ataḥ yadi anuvartate satvam api prāpnoti iti . \\
\noindent \textbf{(P\_8,3.39)} na eṣaḥ doṣaḥ . \\
\noindent \textbf{(P\_8,3.39)} sambandham anuvartiṣyate . \\
\noindent \textbf{(P\_8,3.39)} saḥ apadādau . \\
\noindent \textbf{(P\_8,3.39)} iṇaḥṣaḥ . \\
\noindent \textbf{(P\_8,3.39)} namaspurasoḥ gatyoḥ sakāraḥ iṇaḥ uttarasya ṣakāraḥ . \\
\noindent \textbf{(P\_8,3.39)} idudupadhasya ca apratyayasya ṣakāraḥ namaspurasoḥ gatyoḥ sakāraḥ . \\
\noindent \textbf{(P\_8,3.39)} tirasaḥ anyatarasyām sakāraḥ idudupadhasya ca apratyayasya ṣakāraḥ . \\
\noindent \textbf{(P\_8,3.39)} dvistriścaturitikṛtvo'rthe isusoḥsāmarthye nityamsamāse'nuttarapadasthasya iti ṣakāraḥ tirasaḥ anyatarasyām sakāraḥ . \\
\noindent \textbf{(P\_8,3.39)} ataḥkṛkamikaṃsakumbhapātrakuśākarṇīṣvanavyayasya . \\
\noindent \textbf{(P\_8,3.39)} sakāraḥ anuvartate ṣakāragrahaṇam nivṛttam \\
\noindent \textbf{(P\_8,3.41.1)} idudupadhasya ca apratyayasya iti cet pummuhusoḥ pratiṣedhaḥ . \\
\noindent \textbf{(P\_8,3.41.1)} idudupadhasya ca apratyayasya iti cet pummuhusoḥ pratiṣedhaḥ vaktavyaḥ . \\
\noindent \textbf{(P\_8,3.41.1)} puṃskāmā muhuḥkāmaḥ iti . \\
\noindent \textbf{(P\_8,3.41.1)} vṛddhibhūtānām ṣatvam vaktavyam . \\
\noindent \textbf{(P\_8,3.41.1)} dauṣkulyam naiṣpuruṣyam . \\
\noindent \textbf{(P\_8,3.41.1)} plutānām tādau ca . plutānām tādau ca kupvoḥ ca iti vaktavyam . \\
\noindent \textbf{(P\_8,3.41.1)} sarpi3ṣṭara . \\
\noindent \textbf{(P\_8,3.41.1)} barhī3ṣṭara ni3ṣkula du3ṣpuruṣa . \\
\noindent \textbf{(P\_8,3.41.1)} na vā bahiraṅgalakṣaṇatvāt . na vā vaktavyam . \\
\noindent \textbf{(P\_8,3.41.1)} kim kāraṇam . \\
\noindent \textbf{(P\_8,3.41.1)} bahirangalakṣaṇatvāt vṛddheḥ . \\
\noindent \textbf{(P\_8,3.41.1)} bahiraṅgalakṣaṇā vṛddhiḥ \\
\noindent \textbf{(P\_8,3.41.2)} iha kasmāt na bhavati . \\
\noindent \textbf{(P\_8,3.41.2)} pituḥ karoti . \\
\noindent \textbf{(P\_8,3.41.2)} mātuḥ karoti . \\
\noindent \textbf{(P\_8,3.41.2)} apratyayavisarjanīyasya iti ṣatvam prasajyeta . \\
\noindent \textbf{(P\_8,3.41.2)} apratyayavisarjanīyasya iti ucyate pratyayavisarjanīyaḥ ca ayam . \\
\noindent \textbf{(P\_8,3.41.2)} lupyate atra pratyayavisarjanīyaḥ rātsasya iti . \\
\noindent \textbf{(P\_8,3.41.2)} evam tarhi . \\
\noindent \textbf{(P\_8,3.41.2)} bhrātuṣputragrahaṇam jñāpakam ekādeśanimittāt ṣatvapratiṣedhasya . \\
\noindent \textbf{(P\_8,3.41.2)} yat ayam kaskādiṣu bhrātuṣputraśabdam paṭhati tat jñāpayati ācāryaḥ na ekādeśanimittāt ṣatvam bhavati iti \\
\noindent \textbf{(P\_8,3.43)} dvistriścaturgrahaṇam kimartham . \\
\noindent \textbf{(P\_8,3.43)} iha mā bhūt . \\
\noindent \textbf{(P\_8,3.43)} pañcakṛtvaḥ karoti . \\
\noindent \textbf{(P\_8,3.43)} atha kṛtvorthagrahaṇam kimartham . \\
\noindent \textbf{(P\_8,3.43)} iha mā bhūt . \\
\noindent \textbf{(P\_8,3.43)} catuṣkapālaḥ catuṣkaṇṭakaḥ iti . \\
\noindent \textbf{(P\_8,3.43)} na etat asti . \\
\noindent \textbf{(P\_8,3.43)} astu etena vibhāṣā pūrveṇa nityaḥ vidhiḥ bhaviṣyati . \\
\noindent \textbf{(P\_8,3.43)} na aprāpte pūrveṇa iyam vibhāṣā ārabhyate sā yathā eva iha bādhikā bhavati catuḥ karoti catuṣkaroti iti evam catuṣkapāle api bādhikā syāt . \\
\noindent \textbf{(P\_8,3.43)} na atra pūrveṇa ṣatvam prāpnoti . \\
\noindent \textbf{(P\_8,3.43)} kim kāraṇam . \\
\noindent \textbf{(P\_8,3.43)} apratyayavisarjanīyasya iti ucyate pratyayavisarjanīyaḥ ca ayam . \\
\noindent \textbf{(P\_8,3.43)} lupyate pratyayavisarjanīyaḥ rātsasya iti . \\
\noindent \textbf{(P\_8,3.43)} tasmāt kṛtvorthagrahaṇam kartavyam . \\
\noindent \textbf{(P\_8,3.43)} dvistriścaturgrahaṇam śakyamavaktum . \\
\noindent \textbf{(P\_8,3.43)} kasmāt na bhavati pañcakṛtvaḥ karoti iti . \\
\noindent \textbf{(P\_8,3.43)} idudupadhasya iti vartate . \\
\noindent \textbf{(P\_8,3.43)} na evam śakyam . \\
\noindent \textbf{(P\_8,3.43)} akriyamāṇe dvistriścaturgrahaṇe kṛtvo'rthagrahaṇena visarjanīyaḥ viśeṣyeta . \\
\noindent \textbf{(P\_8,3.43)} tatra kaḥ doṣaḥ . \\
\noindent \textbf{(P\_8,3.43)} iha eva syāt dviṣkaroti dviḥ karoti . \\
\noindent \textbf{(P\_8,3.43)} iha na syāt catuṣkaroti catuḥ karoti iti . \\
\noindent \textbf{(P\_8,3.43)} dvistriścaturgrahaṇe punaḥ kriyamāṇe kṛtvo'rthagrahaṇe dvistriścaturaḥ viśeṣyante . \\
\noindent \textbf{(P\_8,3.43)} dvistriścaturṇām kṛtro'rthe vartamānānām yaḥ visarjanīyaḥ iti . \\
\noindent \textbf{(P\_8,3.43)} etat api na asti prayojanam . \\
\noindent \textbf{(P\_8,3.43)} padasya iti vartate tat kṛtvo'rthagrahaṇena viśeṣayiṣyaḥ . \\
\noindent \textbf{(P\_8,3.43)} padasya kṛtvo'rthe vartamānasya yaḥ visarjanīyaḥ iti . \\
\noindent \textbf{(P\_8,3.43)} kṛtvasujarthe ṣatvam bravīti . \\
\noindent \textbf{(P\_8,3.43)} kasmāt catuṣkapāle mā . \\
\noindent \textbf{(P\_8,3.43)} ṣatvam vibhāṣayā bhūt . \\
\noindent \textbf{(P\_8,3.43)} nanu siddham tatra pūrveṇa . \\
\noindent \textbf{(P\_8,3.43)} siddhe hi ayam vidhatte caturaḥ ṣatvam tadā api krtvo'rthe . \\
\noindent \textbf{(P\_8,3.43)} lupte kṛtvo'rthīye rephasya visarjanīyaḥ hi . \\
\noindent \textbf{(P\_8,3.43)} evam sati tu idānīm dviḥ triḥ catuḥ iti anena kim kāryam . \\
\noindent \textbf{(P\_8,3.43)} anyaḥ hi na idudupadhaḥ kṛtvo'rthe kaḥ cit api asti . \\
\noindent \textbf{(P\_8,3.43)} akriyamāṇe grahaṇe visarjanīyaḥ tadā viśeṣyeta . \\
\noindent \textbf{(P\_8,3.43)} caturaḥ na sidhyati tadā rephasya visarjanīyaḥ hi . \\
\noindent \textbf{(P\_8,3.43)} tasmin tu gṛhyamāṇe yuktam caturaḥ viśeṣaṇam bhavati . \\
\noindent \textbf{(P\_8,3.43)} prakṛtam padam tadantam tasya api viśeṣaṇam nyāyyam \\
\noindent \textbf{(P\_8,3.45)} anuttarapadasthasya iti kimartham . \\
\noindent \textbf{(P\_8,3.45)} paramasarpiḥkuṇḍikā . \\
\noindent \textbf{(P\_8,3.45)} atha idānīm anena mukte pūrveṇa ṣatvam vibhāṣā kasmāt na bhavati isusoḥ sāmarthye iti . \\
\noindent \textbf{(P\_8,3.45)} nānāpadārthayoḥ vartamānayoḥ khyāyate yadā yogaḥ . \\
\noindent \textbf{(P\_8,3.45)} tasmin ṣatvam kāryam tat yuktam tat ca me na iha . \\
\noindent \textbf{(P\_8,3.45)} vyapekṣāsāmarthye pūrvayogaḥ na ca atra vyapekṣāsāmarthyam . \\
\noindent \textbf{(P\_8,3.45)} kim punaḥ kāraṇam vyapekṣāsāmarthyam āśrīyate na punaḥ ekārthībhāvaḥ yathā anyatra . \\
\noindent \textbf{(P\_8,3.45)} aikārthye sāmarthye vākye ṣatvam na me prasajyeta . aikārthye sāmarthye sati vākye ṣatvam na syāt : sarpiṣ karoti . \\
\noindent \textbf{(P\_8,3.45)} sarpiḥ karoti iti . \\
\noindent \textbf{(P\_8,3.45)} tasmāt iha vyapekṣām sāmarthyam sādhu manyante . \\
\noindent \textbf{(P\_8,3.45)} atha cet kṛdantam etat tataḥ adhike na eva me bhavet prāptiḥ . yadi kṛdantam etat tataḥ adhikasya ṣatvam na prāpnoti . \\
\noindent \textbf{(P\_8,3.45)} kim kāraṇam . \\
\noindent \textbf{(P\_8,3.45)} pratyayagrahaṇe yasmāt saḥ tadādeḥ grahaṇam bhavati iti . \\
\noindent \textbf{(P\_8,3.45)} vākye api tarhi na prāpnoti . \\
\noindent \textbf{(P\_8,3.45)} paramasarpiṣkaroti paramasarpiḥ karoti iti . \\
\noindent \textbf{(P\_8,3.45)} vākye ca me vibhāṣā pratiṣedhaḥ na prakalpeta . yat ayam anuttarapadasthasya iti pratiṣedham śāsti tat jñāpayati ācāryaḥ bhavati vākye vibhāṣā iti . \\
\noindent \textbf{(P\_8,3.45)} atha cet saṃvijñānam nitye ṣatve tataḥ vibhāṣā iyam . atha avyutpannam prātipadikam tataḥ nitye ṣatve prāpte iyam vibhāṣā ārabhyate . \\
\noindent \textbf{(P\_8,3.45)} siddham ca me samāse . \\
\noindent \textbf{(P\_8,3.45)} ṣatvam . \\
\noindent \textbf{(P\_8,3.45)} kimartham tarhi idam ucyate . \\
\noindent \textbf{(P\_8,3.45)} pratiṣedhārthaḥ tu yatnaḥ ayam . anuttarapadasthasya iti pratiṣedham vakṣyāmi iti . \\
\noindent \textbf{(P\_8,3.45)} nānāpadārthayoḥ vartamānayoḥ khyāyate yadā yogaḥ . \\
\noindent \textbf{(P\_8,3.45)} tasmin ṣatvam kāryam tat yuktam tat ca me na iha . \\
\noindent \textbf{(P\_8,3.45)} aikārthye sāmarthye vākye ṣatvam na me prasajyeta . \\
\noindent \textbf{(P\_8,3.45)} tasmāt iha vyapekṣām sāmarthyam sādhu manyante . \\
\noindent \textbf{(P\_8,3.45)} atha cet kṛdantam etat tataḥ adhike na eva me bhavet prāptiḥ . \\
\noindent \textbf{(P\_8,3.45)} vākye ca me vibhāṣā pratiṣedhaḥ na prakalpeta . \\
\noindent \textbf{(P\_8,3.45)} atha cet saṃvijñānam nitye ṣatve tataḥ vibhāṣā iyam . \\
\noindent \textbf{(P\_8,3.45)} siddham ca me samāse pratiṣedhārthaḥ tu yatnaḥ ayam . \\
\noindent \textbf{(P\_8,3.55)} atha mūrdhanyagrahaṇam kimartham na apadāntasya ṣaḥ bhavati iti eva ucyeta . \\
\noindent \textbf{(P\_8,3.55)} tatra ayam api arthaḥ ṣakāragrahaṇam na kartavyam bhavati prakṛtam anuvartate . \\
\noindent \textbf{(P\_8,3.55)} kva prakṛtam . \\
\noindent \textbf{(P\_8,3.55)} iṇaḥṣaḥ iti . \\
\noindent \textbf{(P\_8,3.55)} na evam śakyam . \\
\noindent \textbf{(P\_8,3.55)} avaśyam mūrdhanyagrahaṇam kartavyam ihārtham uttarārtham ca . \\
\noindent \textbf{(P\_8,3.55)} ihārtham tāvat . \\
\noindent \textbf{(P\_8,3.55)} iṇaḥṣīdhvaṃluṅliṭāmdho'ṅgāt iti atra ṃūrdhanyagrahaṇam na kartavyam bhavati . \\
\noindent \textbf{(P\_8,3.55)} uttarārtham ca . \\
\noindent \textbf{(P\_8,3.55)} raṣābhyānnoṇaḥsamānapade iti atra ṇakāragrahaṇam na kartavyam bhavati . \\
\noindent \textbf{(P\_8,3.55)} tatra ayam api arthaḥ padāntasya na iti pratiṣedhaḥ na vaktavyaḥ bhavati . \\
\noindent \textbf{(P\_8,3.55)} apadāntābhisambaddham mūrdhanyagrahaṇam anuvartate . \\
\noindent \textbf{(P\_8,3.56)} sagrahaṇam kimartham na saheḥ sāḍaḥ mūrdhanyaḥ bhavati iti eva ucyeta . \\
\noindent \textbf{(P\_8,3.56)} saheḥ sāḍaḥ mūrdhanyaḥ bhavati iti ucyamāne antyasya prasajyeta . \\
\noindent \textbf{(P\_8,3.56)} nanu ca antyasya mūrdhanyavacane prayojanam na asti iti kṛtvā sakārasya bhaviṣyati . \\
\noindent \textbf{(P\_8,3.56)} kutaḥ nu khalu etat anantyārthe ārambhe sakārasya bhaviṣyati . \\
\noindent \textbf{(P\_8,3.56)} na punaḥ ākārasya syāt . \\
\noindent \textbf{(P\_8,3.56)} sthane antaratamaḥ bhavati iti sakārasya bhaviṣyati . \\
\noindent \textbf{(P\_8,3.56)} bhavet prakṛtitaḥ antaratamanirvṛttau satyām siddham syāt . \\
\noindent \textbf{(P\_8,3.56)} ādeśataḥ tu antaratamanirvṛttau satyām ākārasya prasajyeta . \\
\noindent \textbf{(P\_8,3.56)} tasmāt sakāragrahaṇam kartavyam . \\
\noindent \textbf{(P\_8,3.56)} uttarārtham ca sakāragrahaṇam kriyate . \\
\noindent \textbf{(P\_8,3.56)} ādeśapratyayayoḥ sakārasya yathā syāt . \\
\noindent \textbf{(P\_8,3.56)} iha mā bhūt . \\
\noindent \textbf{(P\_8,3.56)} citam , stutam . \\
\noindent \textbf{(P\_8,3.56)} atha sahigrahaṇam kimartham na sāḍaḥ saḥ bhavati iti eva ucyeta . \\
\noindent \textbf{(P\_8,3.56)} saheḥ eva sāḍrūpam bhavati na anyasya . \\
\noindent \textbf{(P\_8,3.56)} yadi evaṃ . \\
\noindent \textbf{(P\_8,3.56)} sāḍaḥ ṣatve samānaśabdapratiṣedhaḥ . \\
\noindent \textbf{(P\_8,3.56)} sāḍaḥ ṣatve samānaśabdānām pratiṣedhaḥ vaktavyaḥ . \\
\noindent \textbf{(P\_8,3.56)} sāḍaḥ daṇḍaḥ. sāḍaḥ vṛścikaḥ iti . \\
\noindent \textbf{(P\_8,3.56)} arthavadgrahaṇāt siddham . \\
\noindent \textbf{(P\_8,3.56)} arthavataḥ sāḍśabdasya grahaṇam na ca eṣaḥ arthavān . \\
\noindent \textbf{(P\_8,3.56)} arthavadgrahaṇāt siddham iti cet taddhitalope arthavattvāt pratiṣedhaḥ . arthavadgrahaṇāt siddham iti cet taddhitalope arthavattvāt pratiṣedhaḥ vaktavyaḥ . \\
\noindent \textbf{(P\_8,3.56)} saha aḍena sāḍaḥ sāḍasya apatyam sāḍiḥ atra prāpnoti . \\
\noindent \textbf{(P\_8,3.56)} na vaktavyaḥ . \\
\noindent \textbf{(P\_8,3.56)} ṣatvatukoḥ ekādeśasya asiddhatvāt na eṣaḥ sāḍśabdaḥ . \\
\noindent \textbf{(P\_8,3.56)} evam api saha ḍena saḍaḥ saḍasya apatram sāḍiḥ atra prāpnoti . \\
\noindent \textbf{(P\_8,3.56)} tasmāt sahigrahaṇam kartavyam \\
\noindent \textbf{(P\_8,3.57-58)} KA\_III,438.20-439.9 Ro\_V,465-466 {1/24}     numvisarjanīyaśarvyavāye niṃseḥ pratiṣedhaḥ . \\
\noindent \textbf{(P\_8,3.57-58)} KA\_III,438.20-439.9 Ro\_V,465-466 {2/24}     numvisarjanīyaśarvyavāye niṃseḥ pratiṣedaḥ vaktavyaḥ . \\
\noindent \textbf{(P\_8,3.57-58)} KA\_III,438.20-439.9 Ro\_V,465-466 {3/24}     niṃsse niṃssva iti . \\
\noindent \textbf{(P\_8,3.57-58)} KA\_III,438.20-439.9 Ro\_V,465-466 {4/24}     tat tarhi vaktavyam . \\
\noindent \textbf{(P\_8,3.57-58)} KA\_III,438.20-439.9 Ro\_V,465-466 {5/24}     na vaktavyam . \\
\noindent \textbf{(P\_8,3.57-58)} KA\_III,438.20-439.9 Ro\_V,465-466 {6/24}     numā eva vyavāye visarjanīyena eva vyavāye śarā eva vyavāye iti . \\
\noindent \textbf{(P\_8,3.57-58)} KA\_III,438.20-439.9 Ro\_V,465-466 {7/24}     kim vaktavyam etat . \\
\noindent \textbf{(P\_8,3.57-58)} KA\_III,438.20-439.9 Ro\_V,465-466 {8/24}     na hi . \\
\noindent \textbf{(P\_8,3.57-58)} KA\_III,438.20-439.9 Ro\_V,465-466 {9/24}     katham anucyamānam gaṃsyate . \\
\noindent \textbf{(P\_8,3.57-58)} KA\_III,438.20-439.9 Ro\_V,465-466 {10/24}     pratyekam vākyaparisamāptiḥ dṛṣṭā iti . \\
\noindent \textbf{(P\_8,3.57-58)} KA\_III,438.20-439.9 Ro\_V,465-466 {11/24}     tat yathā . \\
\noindent \textbf{(P\_8,3.57-58)} KA\_III,438.20-439.9 Ro\_V,465-466 {12/24}     guṇavṛddhisañjñe pratyekam bhavataḥ . \\
\noindent \textbf{(P\_8,3.57-58)} KA\_III,438.20-439.9 Ro\_V,465-466 {13/24}     nanu ca ayam api asti dṛṣṭāntaḥ samudāye vākyaparisamāptiḥ . \\
\noindent \textbf{(P\_8,3.57-58)} KA\_III,438.20-439.9 Ro\_V,465-466 {14/24}     tat yathā . \\
\noindent \textbf{(P\_8,3.57-58)} KA\_III,438.20-439.9 Ro\_V,465-466 {15/24}     gargāḥ śatam daṇḍyantām . \\
\noindent \textbf{(P\_8,3.57-58)} KA\_III,438.20-439.9 Ro\_V,465-466 {16/24}     arthinaḥ ca rājānaḥ hiraṇyena bhavanti na ca pratyekam daṇḍayanti . \\
\noindent \textbf{(P\_8,3.57-58)} KA\_III,438.20-439.9 Ro\_V,465-466 {17/24}     evam tarhi . \\
\noindent \textbf{(P\_8,3.57-58)} KA\_III,438.20-439.9 Ro\_V,465-466 {18/24}     yogavibhāgāt siddham . yogavibhāgaḥ kariṣyate . \\
\noindent \textbf{(P\_8,3.57-58)} KA\_III,438.20-439.9 Ro\_V,465-466 {19/24}     numvyavāye . \\
\noindent \textbf{(P\_8,3.57-58)} KA\_III,438.20-439.9 Ro\_V,465-466 {20/24}     tataḥ visarjanīyavyavāye . \\
\noindent \textbf{(P\_8,3.57-58)} KA\_III,438.20-439.9 Ro\_V,465-466 {21/24}     tataḥ śarvyavāye . \\
\noindent \textbf{(P\_8,3.57-58)} KA\_III,438.20-439.9 Ro\_V,465-466 {22/24}     saḥ tarhi yogavibhāgaḥ kartavyaḥ . \\
\noindent \textbf{(P\_8,3.57-58)} KA\_III,438.20-439.9 Ro\_V,465-466 {23/24}     na kartavyaḥ . \\
\noindent \textbf{(P\_8,3.57-58)} KA\_III,438.20-439.9 Ro\_V,465-466 {24/24}     pratyekam vyavāyaśabdaḥ parisamāpyate \\
\noindent \textbf{(P\_8,3.59.1)} ādeśapratyayayoḥ ṣatve sarakaḥ pratiṣedhaḥ . ādeśapratyayayoḥ ṣatve sarakaḥ pratiṣedaḥ vaktavyaḥ : kṛsaraḥ , dhūsaraḥ . \\
\noindent \textbf{(P\_8,3.59.1)} atyalpam idam ucyate : sarakaḥ iti . \\
\noindent \textbf{(P\_8,3.59.1)} saragādīnām iti vaktavyam iha api yathā syāt : varsam tarsam iti . \\
\noindent \textbf{(P\_8,3.59.1)} tat tarhi vaktavyam . \\
\noindent \textbf{(P\_8,3.59.1)} na vaktavyam . \\
\noindent \textbf{(P\_8,3.59.1)} uṇādayaḥ avyutpannāni prātipadikāni . \\
\noindent \textbf{(P\_8,3.59.1)} na vai etat ṣatve śakyam vijñātum uṇadayaḥ avyutpannāni prātipadikāni iti . \\
\noindent \textbf{(P\_8,3.59.1)} iha hi na syāt . \\
\noindent \textbf{(P\_8,3.59.1)} sarpiṣaḥ yajuṣaḥ iti . \\
\noindent \textbf{(P\_8,3.59.1)} evam tarhi . \\
\noindent \textbf{(P\_8,3.59.1)} bahulavacanāt siddham . bahulam pratyayasañjñā bhavati \\
\noindent \textbf{(P\_8,3.59.2)} kim punaḥ iyam avayavaṣaṣṭhī : ādeśasya yaḥ sakāraḥ pratyayasya yaḥ sakāraḥ iti . \\
\noindent \textbf{(P\_8,3.59.2)} āhosvit samānādhikaraṇā : ādeśaḥ yaḥ sakāraḥ pratyayaḥ yaḥ sakāraḥ iti . \\
\noindent \textbf{(P\_8,3.59.2)} kaḥ ca atra viśeṣaḥ . \\
\noindent \textbf{(P\_8,3.59.2)} ādeśapratyayayoḥ iti avayavaṣaṣṭhī cet dvirvacane pratiṣedhaḥ . \\
\noindent \textbf{(P\_8,3.59.2)} ādeśapratyayayoḥ iti avayavaṣaṣṭhī cet dvirvacane pratiṣedhaḥ vaktavyaḥ . \\
\noindent \textbf{(P\_8,3.59.2)} bisam bisam . \\
\noindent \textbf{(P\_8,3.59.2)} musalam musalam . \\
\noindent \textbf{(P\_8,3.59.2)} samānādhikaraṇānām ca aprāptiḥ . samānādhikaraṇānām ca ṣatvasya aprāptiḥ . \\
\noindent \textbf{(P\_8,3.59.2)} eṣaḥ , akārṣīt . \\
\noindent \textbf{(P\_8,3.59.2)} astu tarhi samānādhikaraṇā . \\
\noindent \textbf{(P\_8,3.59.2)} yadi samānādhikaraṇā siṣeca suṣvāpa , atra na prāpnoti . \\
\noindent \textbf{(P\_8,3.59.2)} na dhātudvirvacane sthāne dvirvacanam śakyam āsthātum . \\
\noindent \textbf{(P\_8,3.59.2)} iha api hi prasajyeta sarīsṛpyate iti . \\
\noindent \textbf{(P\_8,3.59.2)} tasmāt tatra dviḥprayogaḥ dvirvacanam . \\
\noindent \textbf{(P\_8,3.59.2)} iha tarhi kariṣyati hariṣyati pratyayaḥ yaḥ sakāraḥ iti ṣatvam na prāpnoti . \\
\noindent \textbf{(P\_8,3.59.2)} astu tarhi ādeśaḥ yaḥ sakāraḥ pratyayasya yaḥ sakāraḥ iti . \\
\noindent \textbf{(P\_8,3.59.2)} iha tarhi akārṣīt pratyayasya yaḥ sakāraḥ iti ṣatvam na prāpnoti . \\
\noindent \textbf{(P\_8,3.59.2)} mā bhūt evam ādeśaḥ yaḥ sakāraḥ iti evam bhaviṣyati . \\
\noindent \textbf{(P\_8,3.59.2)} iha tarhi : joṣiṣat , mandiṣat iti pratyayasya yaḥ sakāraḥ iti ṣatvam na prāpnoti . \\
\noindent \textbf{(P\_8,3.59.2)} eṣaḥ api iṭi kṛte pratyayasya sakāraḥ . \\
\noindent \textbf{(P\_8,3.59.2)} iha tarhi indraḥ mā vakṣat saḥ , saḥ devan yakṣat . \\
\noindent \textbf{(P\_8,3.59.2)} nānāvibhaktīnām ca samāsānupapattiḥ . nānāvibhaktīnām ca samāsaḥ na upapadyate ādeśapratyayayoḥ iti . \\
\noindent \textbf{(P\_8,3.59.2)} yogavibhāgāt siddham . yogavibhāgaḥ kariṣyate . \\
\noindent \textbf{(P\_8,3.59.2)} ādeśasya ṣaḥ bhavati iti . \\
\noindent \textbf{(P\_8,3.59.2)} tataḥ pratyayasakārasya ṣaḥ bhavati iti . \\
\noindent \textbf{(P\_8,3.59.2)} sa tarhi yogavibhāgaḥ kartavyaḥ . \\
\noindent \textbf{(P\_8,3.59.2)} na kartavyaḥ . \\
\noindent \textbf{(P\_8,3.59.2)} katham . \\
\noindent \textbf{(P\_8,3.59.2)} astu tāvat avayavaṣaṣṭhī . \\
\noindent \textbf{(P\_8,3.59.2)} nanu ca uktam ādeśapratyayayoḥ iti avayavaṣaṣṭhī cet dvirvacane pratiṣedhaḥ iti . \\
\noindent \textbf{(P\_8,3.59.2)} na eṣaḥ doṣaḥ . \\
\noindent \textbf{(P\_8,3.59.2)} dviḥprayogaḥ dvirvacanam . \\
\noindent \textbf{(P\_8,3.59.2)} yat api ucyate samānādhikaraṇānām ca aprāptiḥ iti vyapadeśivadbhāvena bhaviṣyati . \\
\noindent \textbf{(P\_8,3.59.2)} atha vā punaḥ astu samānādhikaraṇā . \\
\noindent \textbf{(P\_8,3.59.2)} katham kariṣyati hariṣyati . \\
\noindent \textbf{(P\_8,3.59.2)} ācāryapravṛttiḥ jñāpayati bhavati evañjātīyakānām ṣatvam iti yat ayam sātpadādyoḥ iti sātpratiṣedham śāsti . \\
\noindent \textbf{(P\_8,3.59.2)} atha vā punaḥ astu ādeśaḥ yaḥ sakāraḥ pratyayasya yaḥ sakāraḥ iti . \\
\noindent \textbf{(P\_8,3.59.2)} katham indraḥ mā , vakṣat saḥ devān yakṣat . \\
\noindent \textbf{(P\_8,3.59.2)} vyapadeśivadbhāvena bhaviṣyati . \\
\noindent \textbf{(P\_8,3.59.2)} saḥ tarhi vyapadeśivadbhāvaḥ vaktavyaḥ . \\
\noindent \textbf{(P\_8,3.59.2)} na vaktavyaḥ . \\
\noindent \textbf{(P\_8,3.59.2)} uktam vā . kim uktam . \\
\noindent \textbf{(P\_8,3.59.2)} tatra vyapadeśivadvacanam ekācaḥ dve prathamārtham ṣatve ca ādeśasampratyayārtham avacanāt lokavijñanāt siddham iti . \\
\noindent \textbf{(P\_8,3.59.2)} yat api nānāvibhaktīnām ca samāsānupapattiḥ iti ācāryapravṛttiḥ jñāpayati nānāvibhaktyoḥ eṣaḥ samāsaḥ iti yat ayam śāsivasighasīnāñca iti ghasigrahaṇam karoti . \\
\noindent \textbf{(P\_8,3.59.2)} katham kṛtvā jñāpakam . \\
\noindent \textbf{(P\_8,3.59.2)} yadi hi ādeśasya yaḥ sakāraḥ iti evam syāt ghasigrahaṇam anarthakam syāt . \\
\noindent \textbf{(P\_8,3.59.2)} paśyati tu ācāryaḥ ādeśaḥ yaḥ sakāraḥ tasya ṣatvam iti tataḥ ghasigrahaṇam karoti . \\
\noindent \textbf{(P\_8,3.61)} stautiṇigrahaṇam kimartham . \\
\noindent \textbf{(P\_8,3.61)} astautiṇyantānām mā bhūt . \\
\noindent \textbf{(P\_8,3.61)} sisikṣati . \\
\noindent \textbf{(P\_8,3.61)} atha evakāraḥ kimarthaḥ . \\
\noindent \textbf{(P\_8,3.61)} niyamārthaḥ . \\
\noindent \textbf{(P\_8,3.61)} sthautiṇyantānām eva na anyeṣām iti . \\
\noindent \textbf{(P\_8,3.61)} na etat asti prayojanam . \\
\noindent \textbf{(P\_8,3.61)} siddhe vidhiḥ ārabhyamāṇaḥ antareṇa evakārakaraṇam niyamārthaḥ bhaviṣyati . \\
\noindent \textbf{(P\_8,3.61)} iṣṭataḥ avadhāraṇārthaḥ tarhi . \\
\noindent \textbf{(P\_8,3.61)} yathā evam vijñāyeta stautiṇyoḥ eva ṣaṇi iti . \\
\noindent \textbf{(P\_8,3.61)} mā evam vijñāyi stautiṇyoḥ ṣaṇi eva iti . \\
\noindent \textbf{(P\_8,3.61)} iha na syāt tuṣṭāva . \\
\noindent \textbf{(P\_8,3.61)} atha ṣaṇi iti kimartham . \\
\noindent \textbf{(P\_8,3.61)} seṣīvyate . \\
\noindent \textbf{(P\_8,3.61)} kaḥ vinate anurodhaḥ . \\
\noindent \textbf{(P\_8,3.61)} avinate niyamaḥ mā bhūt . \\
\noindent \textbf{(P\_8,3.61)} suṣupsati iti . \\
\noindent \textbf{(P\_8,3.61)} kaḥ sānubandhe anurodhaḥ . \\
\noindent \textbf{(P\_8,3.61)} ṣaśabdamātre niyamaḥ mā bhūt . \\
\noindent \textbf{(P\_8,3.61)} suṣupiṣe indram . \\
\noindent \textbf{(P\_8,3.61)} suṣupiṣe iha iti . \\
\noindent \textbf{(P\_8,3.61)} abhyāsāt iti kimartham . \\
\noindent \textbf{(P\_8,3.61)} abhyāsāt yā prāptiḥ tasyāḥ niyamaḥ yathā syāt upasargāt yā prāptiḥ tasyāḥ niyamaḥ mā bhūt . \\
\noindent \textbf{(P\_8,3.61)} abhiṣiṣikṣati . \\
\noindent \textbf{(P\_8,3.61)} na etat asti . \\
\noindent \textbf{(P\_8,3.61)} asiddham upasargāt ṣatvam tasya asiddhatvāt niyamaḥ na bhaviṣyati . \\
\noindent \textbf{(P\_8,3.61)} idam tarhi prayojanam . \\
\noindent \textbf{(P\_8,3.61)} sani yaḥ abhyāsaḥ tasmāt yā prāptiḥ tasyāḥ niyamaḥ yathā syāt yaṅi yaḥ abhyāsaḥ tasmāt yā prāptiḥ tatra niyamaḥ mā bhūt iti . \\
\noindent \textbf{(P\_8,3.61)} soṣupyateḥ san soṣupiṣate . \\
\noindent \textbf{(P\_8,3.61)} atha vā abhyāsāt yā prāptiḥ tasyāḥ niyamaḥ yathā syāt dhātoḥ yā prāptiḥ tasyāḥ niyamaḥ mā bhūt . \\
\noindent \textbf{(P\_8,3.61)} adhīṣiṣati . \\
\noindent \textbf{(P\_8,3.61)} nanu ca ṣaṇi iti ucyate . \\
\noindent \textbf{(P\_8,3.61)} ṣaṇi iti na eṣā parasaptamī śakyā vijñātum sanyaṅantam hi dvirucyate . \\
\noindent \textbf{(P\_8,3.61)} tasmāt eṣā satsaptamī ṣaṇi sati iti . \\
\noindent \textbf{(P\_8,3.61)} satsaptamī cet prāpnoti \\
\noindent \textbf{(P\_8,3.64)} kimartham idam ucyate . \\
\noindent \textbf{(P\_8,3.64)} sthādiṣu abhyāsavacanam niyamārtham . \\
\noindent \textbf{(P\_8,3.64)} niyamārthaḥ ayam ārambhaḥ . \\
\noindent \textbf{(P\_8,3.64)} sthādiṣu eva abhyāsasya yathā syāt . \\
\noindent \textbf{(P\_8,3.64)} iha mā bhūt . \\
\noindent \textbf{(P\_8,3.64)} abhisusūṣati . \\
\noindent \textbf{(P\_8,3.64)} atha kimartham abhyāsena ca iti ucyate . \\
\noindent \textbf{(P\_8,3.64)} tadvyavāye ca aṣopadeśārtham . tadvyavāye abhyāsavyavāye ca aṣopadeśasya api yathā syāt . \\
\noindent \textbf{(P\_8,3.64)} abhiṣiṣeṇayiṣati . \\
\noindent \textbf{(P\_8,3.64)} avarṇārtham ṣaṇi pratiṣedhārtham ca . avaṇārtham tavat . \\
\noindent \textbf{(P\_8,3.64)} abhitaṣṭhau . \\
\noindent \textbf{(P\_8,3.64)} ṣaṇi pratiṣedhārtham . \\
\noindent \textbf{(P\_8,3.64)} abhiṣiṣikṣati \\
\noindent \textbf{(P\_8,3.65.1)} upasargāt ṣatve nisaḥ upasaṅkhyānam aniṇantatvāt . \\
\noindent \textbf{(P\_8,3.65.1)} upasargāt ṣatve nisaḥ upasaṅkhyānam kartavyam . \\
\noindent \textbf{(P\_8,3.65.1)} niḥṣuṇoti niḥṣiñcati . \\
\noindent \textbf{(P\_8,3.65.1)} kim punaḥ kāraṇam na sidhyati . \\
\noindent \textbf{(P\_8,3.65.1)} aniṇantatvāt . \\
\noindent \textbf{(P\_8,3.65.1)} iṇantāt upasargāt ṣatvam ucyate na ca nis* iṇantaḥ . \\
\noindent \textbf{(P\_8,3.65.1)} na vā varṇāśrayatvāt ṣatvasya tadviśeṣakaḥ upasargaḥ dhātuḥ ca . \\
\noindent \textbf{(P\_8,3.65.1)} na vā vaktavyam . \\
\noindent \textbf{(P\_8,3.65.1)} kim kāraṇam . \\
\noindent \textbf{(P\_8,3.65.1)} varṇāśrayatvāt ṣatvasya . \\
\noindent \textbf{(P\_8,3.65.1)} varṇāśrayam ṣatvam . \\
\noindent \textbf{(P\_8,3.65.1)} tadviśeṣakaḥ upasargaḥ dhātuḥ ca . \\
\noindent \textbf{(P\_8,3.65.1)} na evam vijñāyate iṇantāt upasargāt iti . \\
\noindent \textbf{(P\_8,3.65.1)} katham tarhi . \\
\noindent \textbf{(P\_8,3.65.1)} iṇaḥ uttarasya sakārasya saḥ cet iṇ upasargasya saḥ cet sakāraḥ sunotyādīnām iti . \\
\noindent \textbf{(P\_8,3.65.1)} tatra śarvyavāye iti eva siddham . \\
\noindent \textbf{(P\_8,3.65.1)} yadi evam dhātūpasargayoḥ abhisambandhaḥ akṛtaḥ bhavati . \\
\noindent \textbf{(P\_8,3.65.1)} tatra kaḥ doṣaḥ . \\
\noindent \textbf{(P\_8,3.65.1)} iha api prāpnoti . \\
\noindent \textbf{(P\_8,3.65.1)} vigatāḥ secakāḥ asmāt grāmāt visecakaḥ grāmaḥ . \\
\noindent \textbf{(P\_8,3.65.1)} dhātūpasargayoḥ ca abhisambandhaḥ kṛtaḥ . \\
\noindent \textbf{(P\_8,3.65.1)} katham . \\
\noindent \textbf{(P\_8,3.65.1)} sunotyādibhiḥ atra upasargam viśeṣayiṣyāmaḥ sunotyādīnām yaḥ upasargaḥ tasya yaḥ iṇ iti \\
\noindent \textbf{(P\_8,3.65.2)} sunotyādīnām ṣatve ṇyantasya upasaṅkhyānam adhikatvāt . \\
\noindent \textbf{(P\_8,3.65.2)} sunotyādīnām ṣatve ṇyantasya upasaṅkhyānam kartavyam . \\
\noindent \textbf{(P\_8,3.65.2)} abhiṣāvayati . \\
\noindent \textbf{(P\_8,3.65.2)} kim kāraṇam . \\
\noindent \textbf{(P\_8,3.65.2)} adhikatvāt . \\
\noindent \textbf{(P\_8,3.65.2)} vyatiriktaḥ sunotyādiḥ iti kṛtvā upasargāt sunotyādīnām iti ṣatvam na prāpnoti . \\
\noindent \textbf{(P\_8,3.65.2)} na vā avayavasya ananyatvāt . na vā vaktavyam . \\
\noindent \textbf{(P\_8,3.65.2)} kim kāraṇam . \\
\noindent \textbf{(P\_8,3.65.2)} avayavasya ananyatvāt . \\
\noindent \textbf{(P\_8,3.65.2)} avayavaḥ atra ananyaḥ \\
\noindent \textbf{(P\_8,3.65.3)} nāmadhātoḥ tu pratiṣedhaḥ . \\
\noindent \textbf{(P\_8,3.65.3)} nāmadhātoḥ tu pratiṣedhaḥ vaktavyaḥ . \\
\noindent \textbf{(P\_8,3.65.3)} sāvakam icchati abhisāvakīyati parisāvakīyati . \\
\noindent \textbf{(P\_8,3.65.3)} na vā anupasargatvāt . na vā vaktavyaḥ . \\
\noindent \textbf{(P\_8,3.65.3)} kim kāraṇam . \\
\noindent \textbf{(P\_8,3.65.3)} anupasargatvāt . \\
\noindent \textbf{(P\_8,3.65.3)} yatkriyāyuktāḥ tam prati gatyupasargasañjñe bhavataḥ na ca atra sunotim prati kriyāyogaḥ . \\
\noindent \textbf{(P\_8,3.65.3)} kim tarhi sāvakīyatim prati . \\
\noindent \textbf{(P\_8,3.65.3)} iha api tarhi na prāpnoti : abhiṣāvayati . \\
\noindent \textbf{(P\_8,3.65.3)} atra api na sunotim prati kriyāyogaḥ . \\
\noindent \textbf{(P\_8,3.65.3)} kim tarhi sāvayatim prati . \\
\noindent \textbf{(P\_8,3.65.3)} sunotim prati atra kriyāyogaḥ . \\
\noindent \textbf{(P\_8,3.65.3)} katham . \\
\noindent \textbf{(P\_8,3.65.3)} na asau evam preṣyate sunu abhi iti . \\
\noindent \textbf{(P\_8,3.65.3)} kim tarhi . \\
\noindent \textbf{(P\_8,3.65.3)} upasargaviśiṣṭām asau kriyām preṣyate abhiṣuṇu iti \\
\noindent \textbf{(P\_8,3.67)} aprateḥ iti vartate utāho nivṛttam . \\
\noindent \textbf{(P\_8,3.67)} nivṛttam iti āha . \\
\noindent \textbf{(P\_8,3.67)} katham jñāyate . \\
\noindent \textbf{(P\_8,3.67)} yogavibhāgakaraṇasāmarthyāt . \\
\noindent \textbf{(P\_8,3.67)} itarathā hi sadistambhyoḥ aprateḥ iti eva brūyāt . \\
\noindent \textbf{(P\_8,3.67)} asti anyat yogavibhāgakaraṇe prayojanam . \\
\noindent \textbf{(P\_8,3.67)} kim . \\
\noindent \textbf{(P\_8,3.67)} avāccālambanāvidūrvayoḥ iti vakṣyati tat stambheḥ eva yathā syāt sadeḥ mā bhūt iti . \\
\noindent \textbf{(P\_8,3.67)} na etat asti prayojanam . \\
\noindent \textbf{(P\_8,3.67)} ekayoge api sati yasya ālambanāvidūrye staḥ tasya bhaviṣyati . \\
\noindent \textbf{(P\_8,3.67)} kasya ca ālambanāvidūrye staḥ . \\
\noindent \textbf{(P\_8,3.67)} stambheḥ eva \\
\noindent \textbf{(P\_8,3.72)} atha yaḥ prāṇī aprāṇī ca katham tatra bhavitavyam . \\
\noindent \textbf{(P\_8,3.72)} anuṣyandete matsyodake iti . \\
\noindent \textbf{(P\_8,3.72)} āhosvit anusyandete matsyodake iti . \\
\noindent \textbf{(P\_8,3.72)} yadi tāvat aprāṇī vidhinā āśrīyate asti atra aprāṇī iti kṛtvā bhavitavyam ṣatvena . \\
\noindent \textbf{(P\_8,3.72)} atha prāṇī pratiṣedhena āśrīyate asti atra prāṇī iti kṛtvā bhavitavyam pratiṣeedhena . \\
\noindent \textbf{(P\_8,3.72)} kim punaḥ atra arthasatattvam . \\
\noindent \textbf{(P\_8,3.72)} devāḥ jñātum arhanti \\
\noindent \textbf{(P\_8,3.74)} aniṣṭhāyām iti vartate utāho nivṛttam . \\
\noindent \textbf{(P\_8,3.74)} nivṛttam iti āha . \\
\noindent \textbf{(P\_8,3.74)} katham jñāyate . \\
\noindent \textbf{(P\_8,3.74)} yogavibhāgakaraṇasāmarthyāt . \\
\noindent \textbf{(P\_8,3.74)} itarathā hi viparibhyām ca skandeḥ aniṣṭhāyām iti eva brūyāt \\
\noindent \textbf{(P\_8,3.78-79)} KA\_III,444.5-9 Ro\_V,476 {1/4}     iṇgrahaṇam kimartham . \\
\noindent \textbf{(P\_8,3.78-79)} KA\_III,444.5-9 Ro\_V,476 {2/4}     iṇgrahaṇam ḍhatve kavarganivṛttyartham . \\
\noindent \textbf{(P\_8,3.78-79)} KA\_III,444.5-9 Ro\_V,476 {3/4}     iṇghrahaṇam kriyate kavargāt ḍhatvam mā bhūt iti . \\
\noindent \textbf{(P\_8,3.78-79)} KA\_III,444.5-9 Ro\_V,476 {4/4}     pakṣīdhvam yakṣīdhvam \\
\noindent \textbf{(P\_8,3.79)} kim punaḥ idam iṇgrahaṇam pratyayaviśeṣaṇam : iṇaḥ uttareṣām ṣīdhvaṃluṅliṭām yaḥ dhakāraḥ iti . \\
\noindent \textbf{(P\_8,3.79)} āhosvit dhakāraviśeṣaṇam : iṇaḥ uttarasya dhakārasya saḥ cet ṣīdhvaṃluṅliṭām dhakāraḥ iti . \\
\noindent \textbf{(P\_8,3.79)} kaḥ ca atra viśeṣaḥ . \\
\noindent \textbf{(P\_8,3.79)} tatra pratyayaparatve iṭaḥ liṭi ḍhatvam parāditvāt . tatra pratyayaparatve iṭaḥ liṭi ḍhatvam na prāpnoti . \\
\noindent \textbf{(P\_8,3.79)} luluviḍhve luluvidhve iti . \\
\noindent \textbf{(P\_8,3.79)} kim kāraṇam . \\
\noindent \textbf{(P\_8,3.79)} parāditvāt . \\
\noindent \textbf{(P\_8,3.79)} iṭ parādiḥ . \\
\noindent \textbf{(P\_8,3.79)} vacanāt bhaviṣyati . \\
\noindent \textbf{(P\_8,3.79)} asti vacane prayojanam . \\
\noindent \textbf{(P\_8,3.79)} kim . \\
\noindent \textbf{(P\_8,3.79)} alaviḍhvam alavidhvam . \\
\noindent \textbf{(P\_8,3.79)} astu tarhi dhakāraviśeṣaṇam . \\
\noindent \textbf{(P\_8,3.79)} dhakāraparatve ṣīdhvami ananantaratvāt iṭaḥ vibhāṣābhāvaḥ . \\
\noindent \textbf{(P\_8,3.79)} dhakāraparatve ṣīdhvami ananantaratvāt iṭaḥ vibhāṣā na prāpnoti . \\
\noindent \textbf{(P\_8,3.79)} laviṣīḍhvam laviṣīdhvam iti . \\
\noindent \textbf{(P\_8,3.79)} vacanāt bhaviṣyati . \\
\noindent \textbf{(P\_8,3.79)} asti vacane prayojanam . \\
\noindent \textbf{(P\_8,3.79)} kim . \\
\noindent \textbf{(P\_8,3.79)} luluviḍhve luluvidhve iti . \\
\noindent \textbf{(P\_8,3.79)} iṇgrahaṇasya ca aviśeṣaṇatvāt ṣyādimātre ḍhatvaprasaṅgaḥ . iṇgrahaṇasya ca aviśeṣaṇatvāt ṣyādimātre ḍhatvam prāpnoti : pakṣīdhvam , yakṣīdhvam iti . \\
\noindent \textbf{(P\_8,3.79)} na eṣaḥ doṣaḥ . \\
\noindent \textbf{(P\_8,3.79)} aṅgāt iti vakṣyāmi . \\
\noindent \textbf{(P\_8,3.79)} aṅgagrahaṇāt ca doṣaḥ . \\
\noindent \textbf{(P\_8,3.79)} iha na prāpnoti : upadidīyidhve , upadidīyiḍhve . \\
\noindent \textbf{(P\_8,3.79)} yaḥ hi atra aṅgāntyaḥ iṇ na tasmāt uttaraḥ iṭ yasmāt ca uttaraḥ iṭ na asau aṅgāntyaḥ iṇ iti . \\
\noindent \textbf{(P\_8,3.79)} yathā icchasi tathā astu . \\
\noindent \textbf{(P\_8,3.79)} astu tāvat pratyayaviśeṣaṇam . \\
\noindent \textbf{(P\_8,3.79)} nanu ca uktam tatra pratyayaparatve iṭaḥ liṭi ḍhatvam parāditvāt luluviḍhve luluvidhve iti . \\
\noindent \textbf{(P\_8,3.79)} vacanāt bhaviṣyati . \\
\noindent \textbf{(P\_8,3.79)} nanu ca uktam asti vacane prayojanam . \\
\noindent \textbf{(P\_8,3.79)} kim . \\
\noindent \textbf{(P\_8,3.79)} alaviḍhvam alavidhvam iti . \\
\noindent \textbf{(P\_8,3.79)} yat etasmin yoge liḍgrahaṇam tadanavakāśam tasya anavakāśatvāt vacanāt bhaviṣyati . \\
\noindent \textbf{(P\_8,3.79)} atha vā punaḥ astu dhakāraviśeṣaṇam iti . \\
\noindent \textbf{(P\_8,3.79)} nanu ca uktam dhakāraparatve ṣīdhvami ananantaratvāt iṭaḥ vibhāṣābhāvaḥ laviṣīḍhvam laviṣīdhvam iti . \\
\noindent \textbf{(P\_8,3.79)} vacanāt bhaviṣyati . \\
\noindent \textbf{(P\_8,3.79)} nanu ca uktam asti vacane prayojanam . \\
\noindent \textbf{(P\_8,3.79)} kim . \\
\noindent \textbf{(P\_8,3.79)} luluviḍhve luluvidhve iti . \\
\noindent \textbf{(P\_8,3.79)} yat etasmin yoge ṣīdhvaṅgrahaṇam tat anavakāśam tasya anavakāśatvāt vacanāt bhaviṣyati . \\
\noindent \textbf{(P\_8,3.79)} yat api ucyate iṇgrahaṇasya ca aviśeṣaṇatvāt ṣyādimātre ḍhatvaprasaṅgaḥ iti aṅgāt iti vakṣyāmi . \\
\noindent \textbf{(P\_8,3.79)} nanu ca uktam aṅgagrahaṇāt ca doṣaḥ iti . \\
\noindent \textbf{(P\_8,3.79)} pūrvasmin yoge yat aṅgagrahaṇam tat uttaratra nivṛttam . \\
\noindent \textbf{(P\_8,3.79)} atha vā pūrvasmin yoge iṇgrahaṇam pratyayaviśeṣaṇam uttaratra dhakāraviśeṣaṇam \\
\noindent \textbf{(P\_8,3.82)} agneḥ dīrghāt somasya . \\
\noindent \textbf{(P\_8,3.82)} agneḥ dīrghāt somasya iti vaktavyam . \\
\noindent \textbf{(P\_8,3.82)} agnīṣomau . \\
\noindent \textbf{(P\_8,3.82)} itarathā hi aniṣṭaprasaṅgaḥ . itarathā hi aniṣṭam prasajyeta . \\
\noindent \textbf{(P\_8,3.82)} agnisomau māṇavakau iti . \\
\noindent \textbf{(P\_8,3.82)} tat tarhi vaktavyam . \\
\noindent \textbf{(P\_8,3.82)} na vaktavyam . \\
\noindent \textbf{(P\_8,3.82)} gauṇamukhyayoḥ mukhye sampratiprattiḥ . \\
\noindent \textbf{(P\_8,3.82)} tat yathā . \\
\noindent \textbf{(P\_8,3.82)} gauḥ anubandhyaḥ ajaḥ agnīṣomīyaḥ iti na bāhīkaḥ anubadhyate . \\
\noindent \textbf{(P\_8,3.82)} katham tarhi bāhīke vṛddhyāttve bhavataḥ . \\
\noindent \textbf{(P\_8,3.82)} gauḥ tiṣṭhati . \\
\noindent \textbf{(P\_8,3.82)} gām ānaya iti . \\
\noindent \textbf{(P\_8,3.82)} arthāśraye etat evam bhavati . \\
\noindent \textbf{(P\_8,3.82)} yat hi śabdāśrayam śabdamātre tat bhavati . \\
\noindent \textbf{(P\_8,3.82)} śabdāśraye ca vṛddhyātve \\
\noindent \textbf{(P\_8,3.85)} sāntābhyām ca iti vaktavyam . \\
\noindent \textbf{(P\_8,3.85)} iha api yathā syāt . \\
\noindent \textbf{(P\_8,3.85)} mātuḥṣvasā mātuḥsvasā . \\
\noindent \textbf{(P\_8,3.85)} pituḥṣvasā pituḥsvaseti . \\
\noindent \textbf{(P\_8,3.85)} mātuḥ pituḥ iti sāntagrahaṇānarthakyam ekadeśavikṛtasya ananyatvāt . \\
\noindent \textbf{(P\_8,3.85)} mātuḥ pituḥ iti sāntagrahaṇam anarthakam . \\
\noindent \textbf{(P\_8,3.85)} kim kāraṇam . \\
\noindent \textbf{(P\_8,3.85)} ekadeśavikṛtasya ananyatvāt . \\
\noindent \textbf{(P\_8,3.85)} ekadeśavikṛtam ananyavat bhavati iti sāntasya api bhaviṣyati \\
\noindent \textbf{(P\_8,3.87)} astigrahaṇam kimartham . \\
\noindent \textbf{(P\_8,3.87)} iha mā bhūt . \\
\noindent \textbf{(P\_8,3.87)} anusṛtam , visṛtam iti . \\
\noindent \textbf{(P\_8,3.87)} na etat asti prayojanam . \\
\noindent \textbf{(P\_8,3.87)} yadkriyāyuktāḥ tam prati gatyupasargasañjñe bhavataḥ na ca etam sakāram prati kriyāyogaḥ . \\
\noindent \textbf{(P\_8,3.87)} iha api tarhi na prāpnoti abhiṣanti viṣanti iti na hi astiḥ kriyāvacanaḥ . \\
\noindent \textbf{(P\_8,3.87)} kaḥ punaḥ āha na astiḥ kriyāvacanaḥ iti . \\
\noindent \textbf{(P\_8,3.87)} kriyāvacanaḥ astiḥ . \\
\noindent \textbf{(P\_8,3.87)} ātaḥ ca kriyāvacanaḥ vyatyanuṣate kartarikarmavyatihāre iti anena ātmanepadam bhavati . \\
\noindent \textbf{(P\_8,3.87)} karmavyatihāraḥ ca kaḥ . \\
\noindent \textbf{(P\_8,3.87)} kriyāvyatihāraḥ . \\
\noindent \textbf{(P\_8,3.87)} prāduḥśabdāt tarhi mā bhūt . \\
\noindent \textbf{(P\_8,3.87)} prāduḥśabdaḥ ca niyataviṣayaḥ kṛbhvastiyoge eva vartate . \\
\noindent \textbf{(P\_8,3.87)} upasargāt tarhi syateḥ mā bhūt iti . \\
\noindent \textbf{(P\_8,3.87)} iṣyate upasargāt syateḥ ṣatvam . \\
\noindent \textbf{(P\_8,3.87)} ātaś ca iṣyate . \\
\noindent \textbf{(P\_8,3.87)} evam hi āha : upasargāt sunotisuvatisyatistautistobhatisthāsenayasedhasicasañjasvaṅjām iti . \\
\noindent \textbf{(P\_8,3.87)} prāduḥśabdāt tarhi syateḥ mā bhūt iti . \\
\noindent \textbf{(P\_8,3.87)} prāduḥśabdaḥ ca niyataviṣayaḥ kṛbhvastiyoge eva vartate . \\
\noindent \textbf{(P\_8,3.87)} idam tarhi prayojanam iha mā bhūt : anusūteḥ apratyayaḥ anusūḥ anusvaḥ apatyam ānuseyaḥ \\
\noindent \textbf{(P\_8,3.88)} kimartham svapeḥ supibhūtasya grahaṇam kriyate . \\
\noindent \textbf{(P\_8,3.88)} supeḥ ṣatvam svapeḥ mā bhūt . \\
\noindent \textbf{(P\_8,3.88)} supeḥ ṣatvam ucyate tat svapeḥ mā bhūt iti . \\
\noindent \textbf{(P\_8,3.88)} susvapnaḥ visvapnak iti . \\
\noindent \textbf{(P\_8,3.88)} visuṣvāpa iti kena na . visuṣvāpa iti atra kasmāt na bhavati . \\
\noindent \textbf{(P\_8,3.88)} halādiśeṣāt na supiḥ . halādiśeṣe kṛte na eṣaḥ supiḥ bhavati . \\
\noindent \textbf{(P\_8,3.88)} idam iha sampradhāryam . \\
\noindent \textbf{(P\_8,3.88)} halādiśeṣaḥ kriyatām samprasāraṇam iti kim atra kartavyam . \\
\noindent \textbf{(P\_8,3.88)} paratvāt halādiśeṣaḥ . \\
\noindent \textbf{(P\_8,3.88)} iṣṭam pūrvam prasāraṇam . iṣyate halādiśeṣāt pūrvam samprasāraṇam . \\
\noindent \textbf{(P\_8,3.88)} ātaḥ ca iṣyate . \\
\noindent \textbf{(P\_8,3.88)} evam hi āha : abhyāsasamprasāraṇam halādiśeṣāt vipratiṣedhena iti . \\
\noindent \textbf{(P\_8,3.88)} evam tarhi sthādiṣu abhyāsasya iti etasmāt niyamāt na bhaviṣyati . \\
\noindent \textbf{(P\_8,3.88)} sthādīnām niyamaḥ na atra prāk sitāt uttaraḥ supiḥ . prāk sitasaṃśabdanāt saḥ niyamaḥ uttaraḥ ca supiḥ paṭhyate . \\
\noindent \textbf{(P\_8,3.88)} evam tarhi arthavadgrahaṇe na anarthakasya iti evam etasya na bhaviṣyati . \\
\noindent \textbf{(P\_8,3.88)} anarthake viṣuṣupuḥ . yadi arthavataḥ grahaṇam viṣuṣupuḥ iti na sidhyati . \\
\noindent \textbf{(P\_8,3.88)} na eṣaḥ doṣaḥ . \\
\noindent \textbf{(P\_8,3.88)} katham . \\
\noindent \textbf{(P\_8,3.88)} supibhūtaḥ dviḥ ucyate . supibhutasya dvirvacanam . \\
\noindent \textbf{(P\_8,3.88)} supeḥ ṣatvam svapeḥ mā bhūt . \\
\noindent \textbf{(P\_8,3.88)} visuṣvāpa iti kena na . \\
\noindent \textbf{(P\_8,3.88)} halādiśeṣāt na supiḥ . \\
\noindent \textbf{(P\_8,3.88)} iṣṭam pūrvam prasāraṇam . \\
\noindent \textbf{(P\_8,3.88)} sthādīnām niyamaḥ na atra prāk sitāt uttaraḥ supiḥ . \\
\noindent \textbf{(P\_8,3.88)} anarthake viṣuṣupuḥ supibhūtaḥ dviḥ ucyate . \\
\noindent \textbf{(P\_8,3.91)} kapiṣṭhalaḥ gotraprakṛtau . \\
\noindent \textbf{(P\_8,3.91)} kapiṣṭhalaḥ gotraprakṛtau iti vaktavyam . \\
\noindent \textbf{(P\_8,3.91)} gotre iti ucyamāne iha eva syāt . \\
\noindent \textbf{(P\_8,3.91)} kāpiṣṭhaliḥ . \\
\noindent \textbf{(P\_8,3.91)} iha na syāt . \\
\noindent \textbf{(P\_8,3.91)} kapiṣṭhalaḥ kāpiṣṭhalāyanaḥ . \\
\noindent \textbf{(P\_8,3.91)} tat tarhi vaktavyam . \\
\noindent \textbf{(P\_8,3.91)} na vaktavyam . \\
\noindent \textbf{(P\_8,3.91)} na evam vijñāyate kapiṣṭhalaḥ iti gotre nipātyate iti . \\
\noindent \textbf{(P\_8,3.91)} katham tarhi . \\
\noindent \textbf{(P\_8,3.91)} gotre yaḥ kapiṣṭhalaśabdaḥ tasya ṣatvam nipātyate yatra vā tatra vā iti \\
\noindent \textbf{(P\_8,3.97)} sthaḥ iti kim idam dhātugrahaṇam āhosvit rūpagrahaṇam . \\
\noindent \textbf{(P\_8,3.97)} kim ca ataḥ . \\
\noindent \textbf{(P\_8,3.97)} yadi dhātugrahaṇam gosthānam iti atra prāpnoti . \\
\noindent \textbf{(P\_8,3.97)} atha rūpagrahaṇam savyeṣṭhāḥ , parameṣṭhī , savyeṣṭhā sārathiḥ iti atra na prāpnoti . \\
\noindent \textbf{(P\_8,3.97)} yathā icchasi tathā astu . \\
\noindent \textbf{(P\_8,3.97)} astu tāvat dhātugrahaṇam . \\
\noindent \textbf{(P\_8,3.97)} katham gosthānam iti . \\
\noindent \textbf{(P\_8,3.97)} savanādiṣu pāṭhaḥ kariṣyate . \\
\noindent \textbf{(P\_8,3.97)} atha vā punaḥ astu rūpagrahaṇam . \\
\noindent \textbf{(P\_8,3.97)} katham savyeṣṭhāḥ , parmeṣṭhī , savyeṣṭhā sārathiḥ iti . \\
\noindent \textbf{(P\_8,3.97)} sthaḥ sthāsthinsthṝṇām iti vaktavyam \\
\noindent \textbf{(P\_8,3.98)} avihitalakṣaṇaḥ mūrdhanyaḥ suṣāmādiṣu draṣṭavyaḥ \\
\noindent \textbf{(P\_8,3.101)} hrasvāt tādau tiṅi pratiṣedhaḥ . \\
\noindent \textbf{(P\_8,3.101)} hrasvāt tādau tiṅi pratiṣedhaḥ vaktavyaḥ . \\
\noindent \textbf{(P\_8,3.101)} bhindyustarām chindyustarām iti \\
\noindent \textbf{(P\_8,3.105)} stutastomayoḥ chandasi anarthakam vacanam pūrvapadāt iti siddhatvāt . \\
\noindent \textbf{(P\_8,3.105)} stutastomayoḥ chandasi vacanam anarthakam . \\
\noindent \textbf{(P\_8,3.105)} kim kāraṇam . \\
\noindent \textbf{(P\_8,3.105)} pūrvapadāt iti siddhatvāt . \\
\noindent \textbf{(P\_8,3.105)} pūrvapadāt iti eva siddham \\
\noindent \textbf{(P\_8,3.108)} sanoteḥ anaḥ iti ca . \\
\noindent \textbf{(P\_8,3.108)} kim . \\
\noindent \textbf{(P\_8,3.108)} vacanam anarthakam iti eva . \\
\noindent \textbf{(P\_8,3.108)} kim kāraṇam . \\
\noindent \textbf{(P\_8,3.108)} pūrvapadāt iti siddhatvāt . \\
\noindent \textbf{(P\_8,3.108)} niyamārtham tarhi idam vaktavyam . \\
\noindent \textbf{(P\_8,3.108)} sanoteḥ anakārasya eva yathā syāt . \\
\noindent \textbf{(P\_8,3.108)} iha mā bhūt . \\
\noindent \textbf{(P\_8,3.108)} gosanim iti . \\
\noindent \textbf{(P\_8,3.108)} sanoteḥ anaḥ iti niyamārtham iti cet savanādikṛtatvāt siddham . sanoteḥ anaḥ iti niyamārtham iti cet savanādiṣu pāṭhaḥ kariṣyate . \\
\noindent \textbf{(P\_8,3.108)} sanartham tu . sanartham tu idam vaktavyam . \\
\noindent \textbf{(P\_8,3.108)} sisaniṣati . \\
\noindent \textbf{(P\_8,3.108)} etat api na asti prayojanam . \\
\noindent \textbf{(P\_8,3.108)} stautiṇyorevaṣaṇyabhyāsāt iti etasmāt niyamāt na bhaviṣyati . \\
\noindent \textbf{(P\_8,3.108)} ṇyartham tarhi idam vaktavyam . \\
\noindent \textbf{(P\_8,3.108)} sisānayiṣati . \\
\noindent \textbf{(P\_8,3.108)} katham punaḥ aṇyantasya pratiṣedhe ṇyantaḥ śakyaḥ vijñātum . \\
\noindent \textbf{(P\_8,3.108)} sāmarthyāt . \\
\noindent \textbf{(P\_8,3.108)} aṇyantasya pratiṣedhavacane prayojanam na asti iti kṛtvā ṇyante vijñāsyate . \\
\noindent \textbf{(P\_8,3.108)} atha vā ayam asti aṇyantaḥ : sisaniṣateḥ apratyayaḥ sisanīḥ \\
\noindent \textbf{(P\_8,3.110)} kimartham savādiṣu aśvasaniśabdaḥ paṭhyate . \\
\noindent \textbf{(P\_8,3.110)} pūrvapadāt iti ṣatvam prāpnoti . \\
\noindent \textbf{(P\_8,3.110)} tadbādhanārtham . \\
\noindent \textbf{(P\_8,3.110)} na etat asti prayojanam . \\
\noindent \textbf{(P\_8,3.110)} iṇantāt iti tatra anuvartate aniṇantaḥ ca ayam . \\
\noindent \textbf{(P\_8,3.110)} na eva prāpnoti na arthaḥ pratiṣedhena . \\
\noindent \textbf{(P\_8,3.110)} evam tarhi siddhe sati yat savanādiṣu aśvasaniśabdam paṭhati tat jñāpayati ācāryaḥ aniṇantāt api ṣatvam bhavati iti . \\
\noindent \textbf{(P\_8,3.110)} kim etasya jñāpane prayojanam . \\
\noindent \textbf{(P\_8,3.110)} jalāṣāham māṣaḥ iti etat siddham bhavati . \\
\noindent \textbf{(P\_8,3.110)} atha vā ekadeśavikṛtārthaḥ ayam ārambhaḥ . \\
\noindent \textbf{(P\_8,3.110)} aśvaṣāḥ iti \\
\noindent \textbf{(P\_8,3.112)} upasargāt iti yā prāptiḥ bhavitavyam tasyāḥ pratiṣedhena utāho na . \\
\noindent \textbf{(P\_8,3.112)} na bhavitavyam . \\
\noindent \textbf{(P\_8,3.112)} kim kāraṇam . \\
\noindent \textbf{(P\_8,3.112)} upasargāt ṣatvam pratiṣedhaviṣaye ārabhyate tat yathā eva padādilakṣaṇam pratiṣedham bādhate evam sicaḥ yaṅi iti etam api bādhate . \\
\noindent \textbf{(P\_8,3.112)} na bādhate . \\
\noindent \textbf{(P\_8,3.112)} kim kāraṇam . \\
\noindent \textbf{(P\_8,3.112)} yena na aprāpte tasya bādhanam bhavati na ca aprāpte padādilakṣaṇe pratiṣedhe upasargāt ṣatvam ārabhyate sicaḥ yaṅi iti etasmin punaḥ prāpte ca aprāpte ca . \\
\noindent \textbf{(P\_8,3.112)} atha vā purastāt apavādāḥ anantarān vidhīn bādhante iti evam upasargāt ṣatvam padādilakṣaṇam pratiṣedham bādhiṣyate sicaḥ yaṅi iti etam na bādhiṣyate . \\
\noindent \textbf{(P\_8,3.112)} tasmāt abhisesicyate iti bhavitavyam \\
\noindent \textbf{(P\_8,3.115)} kimartham sahiḥ soḍhabhūtaḥ gṛhyate . \\
\noindent \textbf{(P\_8,3.115)} yatra asya etat rūpam tatra yathā syāt . \\
\noindent \textbf{(P\_8,3.115)} iha mā bhūt . \\
\noindent \textbf{(P\_8,3.115)} pariṣahate iti \\
\noindent \textbf{(P\_8,3.116)} stambhusivusahām caṅi upasargāt . \\
\noindent \textbf{(P\_8,3.116)} stambhusivusahām caṅi upasargāt iti vaktavyam . \\
\noindent \textbf{(P\_8,3.116)} kim prayojanam . \\
\noindent \textbf{(P\_8,3.116)} upasargāt yā prāptiḥ tasyāḥ pratiṣedhaḥ yathā syāt . \\
\noindent \textbf{(P\_8,3.116)} abhyāsāt yā prāptiḥ tasyāḥ pratiṣedhaḥ mā bhūt iti . \\
\noindent \textbf{(P\_8,3.116)} paryasīṣahat \\
\noindent \textbf{(P\_8,3.117)} sani kim udāharaṇam . \\
\noindent \textbf{(P\_8,3.117)} susūṣati . \\
\noindent \textbf{(P\_8,3.117)} na etat asti prayojanam . \\
\noindent \textbf{(P\_8,3.117)} stautiṇyoḥ eva ṣaṇi iti etasmāt niyamāt na bhaviṣyati . \\
\noindent \textbf{(P\_8,3.117)} idam tarhi . \\
\noindent \textbf{(P\_8,3.117)} abhisusūṣati . \\
\noindent \textbf{(P\_8,3.117)} etat api na asti prayojanam . \\
\noindent \textbf{(P\_8,3.117)} sthādiṣvabhyāsenacābhyāsasya iti etasmāt niyamāt na bhaviṣyati . \\
\noindent \textbf{(P\_8,3.117)} idam tarhi prayojanam : abhisusūṣateḥ apratyayaḥ abhisusūḥ \\
\noindent \textbf{(P\_8,3.118)} sadaḥ liṭi pratiṣedhe svañjeḥ upasaṅkhyanam . \\
\noindent \textbf{(P\_8,3.118)} sadaḥ liṭi pratiṣedhe svañjeḥ upasaṅkhyānam kartavyam . \\
\noindent \textbf{(P\_8,3.118)} pariṣasvaje \\
\noindent \textbf{(P\_8,4.1)} raṣābhyām ṇatve ṛkāragrahaṇam . \\
\noindent \textbf{(P\_8,4.1)} raṣābhyām ṇatve ṛkāragrahaṇam kartavyam . \\
\noindent \textbf{(P\_8,4.1)} raṣābhyām naḥ ṇaḥ samānapade ṛkārāt ca iti vaktavyam . \\
\noindent \textbf{(P\_8,4.1)} iha api yathā syāt : mātṝṇām , pitṝṇām iti . \\
\noindent \textbf{(P\_8,4.1)} tat tarhi vaktavyam . \\
\noindent \textbf{(P\_8,4.1)} na vaktavyam . \\
\noindent \textbf{(P\_8,4.1)} yaḥ asau ṛkāre rephaḥ tadāśrayam ṇatvam bhaviṣyati . \\
\noindent \textbf{(P\_8,4.1)} na sidhyati . \\
\noindent \textbf{(P\_8,4.1)} kim kāraṇam . \\
\noindent \textbf{(P\_8,4.1)} na hi varṇaikadeśāḥ varṇagrahaṇena gṛhyante . \\
\noindent \textbf{(P\_8,4.1)} ekadeśe nuḍādiṣu ca uktam . kim uktam . \\
\noindent \textbf{(P\_8,4.1)} agrahaṇam cet nuḍvidhilādeśavināmeṣu ṛkāragrahaṇam iti . \\
\noindent \textbf{(P\_8,4.1)} tasmāt gṛhyante . \\
\noindent \textbf{(P\_8,4.1)} evam api na sidhyati . \\
\noindent \textbf{(P\_8,4.1)} kim kāraṇam . \\
\noindent \textbf{(P\_8,4.1)} ananantaratvāt . \\
\noindent \textbf{(P\_8,4.1)} yat tat rephāt param bhakteḥ tena vyavahitatvāt na prāpnoti . \\
\noindent \textbf{(P\_8,4.1)} aḍvyavāye iti evam bhaviṣyati . \\
\noindent \textbf{(P\_8,4.1)} na sidhyati . \\
\noindent \textbf{(P\_8,4.1)} kim kāraṇam . \\
\noindent \textbf{(P\_8,4.1)} varṇaikadeśāḥ ke varṇagrahaṇena gṛhyante . \\
\noindent \textbf{(P\_8,4.1)} ye vyapavṛktāḥ api varṇāḥ bhavanti . \\
\noindent \textbf{(P\_8,4.1)} yat ca atra rephāt param bhakteḥ na tat kva cit api vyapavṛktam dṛśyate . \\
\noindent \textbf{(P\_8,4.1)} evam tarhi yogavibhāgaḥ kariṣyate . \\
\noindent \textbf{(P\_8,4.1)} raṣābhyām naḥ ṇaḥ samānapade . \\
\noindent \textbf{(P\_8,4.1)} tataḥ vyavāye . \\
\noindent \textbf{(P\_8,4.1)} vyavāye ca raṣābhyām naḥ ṇaḥ bhavati iti . \\
\noindent \textbf{(P\_8,4.1)} tataḥ aṭkupvāṅnumbhiḥ iti . \\
\noindent \textbf{(P\_8,4.1)} idam idānīm kimartham . \\
\noindent \textbf{(P\_8,4.1)} niyamārtham . \\
\noindent \textbf{(P\_8,4.1)} etaiḥ eva akṣarasamāmnāyikaiḥ vyavāye na anyaiḥ iti . \\
\noindent \textbf{(P\_8,4.1)} atha vā ācāryapravṛttiḥ jñāpayati bhavati ṛkārāt ṇatvam iti yat ayam kṣubhnādiṣu nṛnamanaśabdam paṭhati . \\
\noindent \textbf{(P\_8,4.1)} na etat asti jñāpakam . \\
\noindent \textbf{(P\_8,4.1)} vṛddhyartham etat syāt . \\
\noindent \textbf{(P\_8,4.1)} nārnamaniḥ iti . \\
\noindent \textbf{(P\_8,4.1)} yat tarhi tatra eva tṛpnotiśabdam paṭhati . \\
\noindent \textbf{(P\_8,4.1)} yat ca api nṛnamanaśabdam paṭhati . \\
\noindent \textbf{(P\_8,4.1)} nanu ca uktam vṛddhyartham etat syāt iti . \\
\noindent \textbf{(P\_8,4.1)} bahiraṅgā vṛddhiḥ . \\
\noindent \textbf{(P\_8,4.1)} antaraṅgam ṇatvam . \\
\noindent \textbf{(P\_8,4.1)} asiddham bahiraṅgam antaraṅge . \\
\noindent \textbf{(P\_8,4.1)} atha vā upariṣṭāt yogavibhāgaḥ kariṣyate . \\
\noindent \textbf{(P\_8,4.1)} ṛtaḥ naḥ ṇaḥ bhavati . \\
\noindent \textbf{(P\_8,4.1)} tataḥ chandasi avagrahāt ṛtaḥ iti eva \\
\noindent \textbf{(P\_8,4.2.1)} aḍvyavāye ṇatve anyavyavāye pratiṣedhaḥ . \\
\noindent \textbf{(P\_8,4.2.1)} aḍvyavāye ṇatve anyavyavāye pratiṣedhaḥ vaktavyaḥ . \\
\noindent \textbf{(P\_8,4.2.1)} ādarśena akṣadarśena . \\
\noindent \textbf{(P\_8,4.2.1)} na vā anyena vyapetatvāt . na vā vaktavyaḥ . \\
\noindent \textbf{(P\_8,4.2.1)} kim kāraṇam . \\
\noindent \textbf{(P\_8,4.2.1)} anyena vyapetatvāt . \\
\noindent \textbf{(P\_8,4.2.1)} anyena atra vyavāyaḥ . \\
\noindent \textbf{(P\_8,4.2.1)} yadi api atra anyena vyavāyaḥ aṭā api tu vyavāyaḥ asti tatra asti aḍvyavāye iti prāpnoti . \\
\noindent \textbf{(P\_8,4.2.1)} aṭā eva vyavāye bhavati . \\
\noindent \textbf{(P\_8,4.2.1)} kim vaktavyam etat . \\
\noindent \textbf{(P\_8,4.2.1)} na hi . \\
\noindent \textbf{(P\_8,4.2.1)} katham anucyamānam gaṃsyate . \\
\noindent \textbf{(P\_8,4.2.1)} aḍgrahaṇasāmarthyāt . \\
\noindent \textbf{(P\_8,4.2.1)} yadi hi yatra aṭā ca anyena ca vyavāyaḥ tatra syāt aḍgrahaṇam anarthakam syāt . \\
\noindent \textbf{(P\_8,4.2.1)} vyavāye naḥ ṇaḥ bhavati iti eva brūyāt . \\
\noindent \textbf{(P\_8,4.2.1)} asti anyat aḍgrahaṇasya prayojanam . \\
\noindent \textbf{(P\_8,4.2.1)} kim . \\
\noindent \textbf{(P\_8,4.2.1)} yaḥ anirdiṣṭaiḥ eva vyavāyaḥ tatra mā bhūt . \\
\noindent \textbf{(P\_8,4.2.1)} kṛtsnam mṛtsnā iti . \\
\noindent \textbf{(P\_8,4.2.1)} yadi etāvat prayojanam syāt śarvyavāye na iti eva brūyāt \\
\noindent \textbf{(P\_8,4.2.2)} tatsamudāye ṇatvāprasiddhiḥ yathā anyatra . \\
\noindent \textbf{(P\_8,4.2.2)} tatsamudāye vyavāyasamudāye ṇatvasya aprasiddhiḥ . \\
\noindent \textbf{(P\_8,4.2.2)} arkeṇa argheṇa . \\
\noindent \textbf{(P\_8,4.2.2)} yathā anyatra api vyavāyasamudāye kāryam na bhavati . \\
\noindent \textbf{(P\_8,4.2.2)} kva anyatra . \\
\noindent \textbf{(P\_8,4.2.2)} numvisarjanīyaśarvyavāye'pi niṃsse niṃssva iti . \\
\noindent \textbf{(P\_8,4.2.2)} kim punaḥ kāraṇam anyatra api vyavāyasamudāye kāryam na bhavati . \\
\noindent \textbf{(P\_8,4.2.2)} pratyekam vākyaparisamāptiḥ dṛṣṭā iti . \\
\noindent \textbf{(P\_8,4.2.2)} tat yathā . \\
\noindent \textbf{(P\_8,4.2.2)} guṇavṛddhisañjñe pratyekam bhavataḥ . \\
\noindent \textbf{(P\_8,4.2.2)} nanu ca ayam api asti dṛṣṭāntaḥ samudāye vākyaparisamāptiḥ iti . \\
\noindent \textbf{(P\_8,4.2.2)} tat yathā . \\
\noindent \textbf{(P\_8,4.2.2)} gargāḥ śatam daṇḍyantām iti . \\
\noindent \textbf{(P\_8,4.2.2)} arthinaḥ ca rājānaḥ hiraṇyena bhavanti na ca pratyekam daṇḍayanti . \\
\noindent \textbf{(P\_8,4.2.2)} yadi evam ekena vyavāye na prāpnoti . \\
\noindent \textbf{(P\_8,4.2.2)} kiriṇā ririṇā iti . \\
\noindent \textbf{(P\_8,4.2.2)} ubhayathā api vākyaparisamāptiḥ dṛśyate . \\
\noindent \textbf{(P\_8,4.2.2)} tat yathā . \\
\noindent \textbf{(P\_8,4.2.2)} gargaiḥ saha na bhoktavyam iti pratyekam ca na sambhujyate samuditaiḥ ca \\
\noindent \textbf{(P\_8,4.2.3)} kuvyavāye hādeśeṣu pratiṣedhaḥ . \\
\noindent \textbf{(P\_8,4.2.3)} kuvyavāye hādeśeṣu pratiṣedhaḥ vakavyaḥ . \\
\noindent \textbf{(P\_8,4.2.3)} kim prayojanam . \\
\noindent \textbf{(P\_8,4.2.3)} prayojanam vṛtraghnaḥ srughnaḥ prāghāni iti . hanteratpūrvasya iti atpūrvagrahaṇam na kartavyam bhavati . \\
\noindent \textbf{(P\_8,4.2.3)} numvyavāye ṇatve anusvārābhāve pratiṣedhaḥ . numvyavāye ṇatve anusvārābhāve pratiṣedhaḥ vaktavyaḥ . \\
\noindent \textbf{(P\_8,4.2.3)} prenvanam prenvanīyam . \\
\noindent \textbf{(P\_8,4.2.3)} anāgame ca ṇatvam . anāgame ca ṇatvam vaktavyam . \\
\noindent \textbf{(P\_8,4.2.3)} tṛmpaṇīyam . \\
\noindent \textbf{(P\_8,4.2.3)} anusvāravyavāyavacanāt tu siddham . anusvāravyavāye naḥ ṇaḥ bhavati iti vaktavyam . \\
\noindent \textbf{(P\_8,4.2.3)} tadanusvāragrahaṇam kartavyam . \\
\noindent \textbf{(P\_8,4.2.3)} na kartavyam . \\
\noindent \textbf{(P\_8,4.2.3)} kriyate nyāse eva . \\
\noindent \textbf{(P\_8,4.2.3)} nakāre anusvāraḥ parasavarṇībhūtaḥ nirdiśyate . \\
\noindent \textbf{(P\_8,4.2.3)} iha api tarhi prāpnoti . \\
\noindent \textbf{(P\_8,4.2.3)} prenvanam prenvanīyam . \\
\noindent \textbf{(P\_8,4.2.3)} anusvāraviśeṣaṇam numgrahaṇam . \\
\noindent \textbf{(P\_8,4.2.3)} numaḥ yaḥ anusvāraḥ iti . \\
\noindent \textbf{(P\_8,4.2.3)} iha api tarhi na prāpnoti . \\
\noindent \textbf{(P\_8,4.2.3)} tṛmpaṇam tṛmpaṇīyam . \\
\noindent \textbf{(P\_8,4.2.3)} evam tarhi ayogavāhānām aviśeṣeṇa upadeśaḥ coditaḥ tatra anusvāre kṛte aḍvyavāye iti eva siddham . \\
\noindent \textbf{(P\_8,4.2.3)} yadi evam na arthaḥ numgrahaṇena . \\
\noindent \textbf{(P\_8,4.2.3)} anusvāre kṛte aḍvyavāye iti eva siddham \\
\noindent \textbf{(P\_8,4.3.1)} pūrvapadāt sañjñāyām uttarapadagrahaṇam . \\
\noindent \textbf{(P\_8,4.3.1)} pūrvapadāt sañjñāyām uttarapadagrahaṇam kartavyam . \\
\noindent \textbf{(P\_8,4.3.1)} kim prayojanam . \\
\noindent \textbf{(P\_8,4.3.1)} taddhitapūrvapadasthāpratiṣedhārtham . taddhitasthasya pūrvapadasthasya ca pratiṣedhaḥ mā bhūt . \\
\noindent \textbf{(P\_8,4.3.1)} khārapāyaṇaḥ karaṇapriyaḥ .. tat tarhi vaktavyam . \\
\noindent \textbf{(P\_8,4.3.1)} na vaktavyam . \\
\noindent \textbf{(P\_8,4.3.1)} pūrvapadam uttarapadam iti sambandhiśabdau etau . \\
\noindent \textbf{(P\_8,4.3.1)} sati pūrvapade uttarapadam bhavati sati ca uttarapade pūrvapadam bhavati . \\
\noindent \textbf{(P\_8,4.3.1)} tatra sambandhāt etat gantavyam yat prati pūrvapadam iti etat bhavati tatsthasya niyamaḥ iti . \\
\noindent \textbf{(P\_8,4.3.1)} kim ca prati etat bhavati . \\
\noindent \textbf{(P\_8,4.3.1)} uttarapadam prati \\
\noindent \textbf{(P\_8,4.3.2)} sañjñāyām niyamavacane gapratiṣedhāt niyamapratiṣedhaḥ . \\
\noindent \textbf{(P\_8,4.3.2)} sañjñāyām niyamavacane gapratiṣedhāt niyamasya ayam pratiṣedhaḥ vijñāyate agaḥ iti . \\
\noindent \textbf{(P\_8,4.3.2)} tatra kaḥ doṣaḥ . \\
\noindent \textbf{(P\_8,4.3.2)} tatra nityam ṇatvaprasaṅgaḥ . tatra pūrveṇa sañjñāyām ca asañjñāyām ca nityam ṇatvam prāpnoti . \\
\noindent \textbf{(P\_8,4.3.2)} yogavibhāgāt siddham . yogavibhāgaḥ kariṣyate . \\
\noindent \textbf{(P\_8,4.3.2)} pūrvapadāt sañjñāyām . \\
\noindent \textbf{(P\_8,4.3.2)} tataḥ agaḥ . \\
\noindent \textbf{(P\_8,4.3.2)} gāntāt pūrvapadāt yā ca yāvatī ṇatvaprāptiḥ tasyāḥ sarvasyāḥ pratiṣedhaḥ . \\
\noindent \textbf{(P\_8,4.3.2)} apratiṣedhaḥ vā yathā sarvanāmasañjñāyām . na vā arthaḥ pratiṣedhena . \\
\noindent \textbf{(P\_8,4.3.2)} ṇatvam kasmāt na bhavati . \\
\noindent \textbf{(P\_8,4.3.2)} yathā sarvanāmasañjñāyām . \\
\noindent \textbf{(P\_8,4.3.2)} uktam ca sarvanāmasañjñāyām sarvanāmasañjñāyām nipātanāt ṇatvābhāvaḥ iti . \\
\noindent \textbf{(P\_8,4.3.2)} yathā punaḥ tatra nipātanam kriyate sarvādīni sarvanāmāni iti iha idānīm kim nipātanam . \\
\noindent \textbf{(P\_8,4.3.2)} iha api nipātanam asti . \\
\noindent \textbf{(P\_8,4.3.2)} kim . \\
\noindent \textbf{(P\_8,4.3.2)} aṇṛgayanādibhyaḥ iti . \\
\noindent \textbf{(P\_8,4.3.2)} na eva vā punaḥ atra pūrveṇa ṇatvam prāpnoti . \\
\noindent \textbf{(P\_8,4.3.2)} kim kāraṇam . \\
\noindent \textbf{(P\_8,4.3.2)} samānapade iti ucyate na ca etat samānapadam . \\
\noindent \textbf{(P\_8,4.3.2)} samāse kṛte samānapadam . \\
\noindent \textbf{(P\_8,4.3.2)} samānam eva yat nityam na ca etat nityam samānapadam eva . \\
\noindent \textbf{(P\_8,4.3.2)} kim vaktavyam etat . \\
\noindent \textbf{(P\_8,4.3.2)} na hi . \\
\noindent \textbf{(P\_8,4.3.2)} katham anucyamānam gaṃsyate . \\
\noindent \textbf{(P\_8,4.3.2)} samānagrahaṇasāmarthyāt. yadi hi yat samānam ca asamānam ca tatra syāt samanagrahaṇam anarthakam syāt \\
\noindent \textbf{(P\_8,4.6)} dvyakṣaratryakṣarebhyaḥ iti vaktavyam . \\
\noindent \textbf{(P\_8,4.6)} iha mā bhūt . \\
\noindent \textbf{(P\_8,4.6)} devadāruvanam . \\
\noindent \textbf{(P\_8,4.6)} irikādibhyaḥ pratiṣedhaḥ vaktavyaḥ . \\
\noindent \textbf{(P\_8,4.6)} irikāvanam timiravanam \\
\noindent \textbf{(P\_8,4.7)} adantāt adantasya iti vaktavyam . \\
\noindent \textbf{(P\_8,4.7)} iha mā bhūt . \\
\noindent \textbf{(P\_8,4.7)} dīrghāhnī śarat iti . \\
\noindent \textbf{(P\_8,4.7)} tat tarhi vaktavyam . \\
\noindent \textbf{(P\_8,4.7)} na vaktavyam . \\
\noindent \textbf{(P\_8,4.7)} na eṣā ahanśabdāt ṣaṣṭhī . \\
\noindent \textbf{(P\_8,4.7)} kā tarhi . \\
\noindent \textbf{(P\_8,4.7)} ahnaśabdāt prathamā pūrvasūtranirdeśaḥ ca . \\
\noindent \textbf{(P\_8,4.7)} atha vā yuvādiṣu pāṭhaḥ kariṣyate \\
\noindent \textbf{(P\_8,4.8)} āhitopasthitayoḥ iti vaktavyam . \\
\noindent \textbf{(P\_8,4.8)} iha api yathā syāt . \\
\noindent \textbf{(P\_8,4.8)} ikṣuvāhaṇam śaravāhaṇam . \\
\noindent \textbf{(P\_8,4.8)} aparaḥ āha : vāhanam vāhyāt iti vaktavyam . \\
\noindent \textbf{(P\_8,4.8)} yadā hi gargāṇām vāhanam apaviddham tiṣṭhati tadā mā bhūt . \\
\noindent \textbf{(P\_8,4.8)} gargavāhanam iti \\
\noindent \textbf{(P\_8,4.10)} vāprakaraṇe girinadyādīnām upasaṅkhyānam . \\
\noindent \textbf{(P\_8,4.10)} vāprakaraṇe girinadyādīnām upasaṅkhyānam kartavyam . \\
\noindent \textbf{(P\_8,4.10)} giriṇadī girinadī . \\
\noindent \textbf{(P\_8,4.10)} cakraṇitambā cakranitambā \\
\noindent \textbf{(P\_8,4.11)} prātipadikāntasya ṇatve samāsāntagrahaṇam asamāsāntapratiṣedhārtham . \\
\noindent \textbf{(P\_8,4.11)} prātipadikāntasya ṇatve samāsāntagrahaṇam kartavyam . \\
\noindent \textbf{(P\_8,4.11)} kim prayojanam . \\
\noindent \textbf{(P\_8,4.11)} asamāsāntapratiṣedhārtham . \\
\noindent \textbf{(P\_8,4.11)} asamāsāntasya mā bhūt . \\
\noindent \textbf{(P\_8,4.11)} gargabhaginī dakṣabhaginī iti . \\
\noindent \textbf{(P\_8,4.11)} na vā bhavati gargabhagiṇī iti . \\
\noindent \textbf{(P\_8,4.11)} bhavati yadā etat vākyam gargāṇām bhagaḥ gargabhagaḥ gargabhagaḥ asyāḥ asti iti . \\
\noindent \textbf{(P\_8,4.11)} yadā tu etat vākyam bhavati gargāṇām bhaginī gargabhaginī iti tadā na bhavitavyam . \\
\noindent \textbf{(P\_8,4.11)} tadā mā bhūt iti . \\
\noindent \textbf{(P\_8,4.11)} yadi samāsāntagrahaṇam kriyate māṣavāpiṇī vṛīhivāpiṇī atra na prāpnoti . \\
\noindent \textbf{(P\_8,4.11)} liṅgaviśiṣṭagrahaṇe ca uktam . kim uktam . \\
\noindent \textbf{(P\_8,4.11)} gatikārakopapadānām kṛdbhiḥ saha samāsavacanam prāk subutpatteḥ iti \\
\noindent \textbf{(P\_8,4.11)} tatra yuvādipratiṣedhaḥ . \\
\noindent \textbf{(P\_8,4.11)} tatra yuvādīnām pratiṣedaḥ vaktavyaḥ . \\
\noindent \textbf{(P\_8,4.11)} āryayūnā kṣatriyayūnā prapakvāni paripakvāni dīrghāhnī śarat iti \\
\noindent \textbf{(P\_8,4.13)} atha iha katham bhavitavyam . \\
\noindent \textbf{(P\_8,4.13)} māṣakumbhavāpeṇa vrīhikumbhavāpeṇa iti . \\
\noindent \textbf{(P\_8,4.13)} kim nityam ṇatvena bhavitavyam āhosvit vibhāṣayā . \\
\noindent \textbf{(P\_8,4.13)} yadā tāvat etat vākyam bhavati kumbhasya vāpaḥ kumbhavāpaḥ māṣāṇām kumbhavāpaḥ māṣakumbhavāpa iti tadā nityam ṇatvena bhavitavyam . \\
\noindent \textbf{(P\_8,4.13)} yadā tu etat vākyam bhavati māṣāṇām kumbhaḥ māṣakumbhaḥ māṣakumbhasya vāpaḥ māṣakumbhavāpaḥ iti tadā vibhāṣayā bhavitavyam \\
\noindent \textbf{(P\_8,4.14.1)} asamāsagrahaṇam kimartham . \\
\noindent \textbf{(P\_8,4.14.1)} samāse iti vartate asamāse api yathā syāt . \\
\noindent \textbf{(P\_8,4.14.1)} praṇamati pariṇamati . \\
\noindent \textbf{(P\_8,4.14.1)} kva punaḥ samāsagrahaṇam prakṛtam . \\
\noindent \textbf{(P\_8,4.14.1)} pūrvapadātsañjñāyāmagaḥ iti . \\
\noindent \textbf{(P\_8,4.14.1)} katham punaḥ tena samāsagrahaṇam śakyam vijñātum . \\
\noindent \textbf{(P\_8,4.14.1)} pūrvapadagrahaṇasāmarthyāt . \\
\noindent \textbf{(P\_8,4.14.1)} samāse eva etat bhavati pūrvapadam uttarapadam iti . \\
\noindent \textbf{(P\_8,4.14.1)} atha apigrahaṇam kimartham . \\
\noindent \textbf{(P\_8,4.14.1)} samāse api yathā syāt . \\
\noindent \textbf{(P\_8,4.14.1)} praṇāmakaḥ pariṇāmakaḥ . \\
\noindent \textbf{(P\_8,4.14.1)} yadi tarhi samāse ca asamāse ca iṣyate na arthaḥ asamāsepigrahaṇena . \\
\noindent \textbf{(P\_8,4.14.1)} nivṛttam pūrvapadāt iti . \\
\noindent \textbf{(P\_8,4.14.1)} aviśeṣeṇa upasargāt ṇatvam vakṣyāmi . \\
\noindent \textbf{(P\_8,4.14.1)} samāse niyamāt na prāpnoti . \\
\noindent \textbf{(P\_8,4.14.1)} asiddham upasargāt ṇatvam tasya asiddhatvāt niyamaḥ na bhaviṣyati . \\
\noindent \textbf{(P\_8,4.14.1)} evam tarhi siddhe sati yat asamāse apigrahaṇam karoti tat jñāpayati ācāryaḥ na yoge yogaḥ asiddhaḥ . \\
\noindent \textbf{(P\_8,4.14.1)} kim tarhi prakaraṇe prakaraṇam asiddham iti . \\
\noindent \textbf{(P\_8,4.14.1)} kim etasya jñapane prayojanam . \\
\noindent \textbf{(P\_8,4.14.1)} yat tat uktam niṣkṛtam niṣpītam iti atra satvasya asiddhatvāt ṣatvam na prāpnoti iti saḥ na doṣaḥ bhavati \\
\noindent \textbf{(P\_8,4.14.2)} ṇopadeśam prati upasargābhāvāt anirdeśaḥ . \\
\noindent \textbf{(P\_8,4.14.2)} agamakaḥ nirdeśaḥ anirdeśaḥ . \\
\noindent \textbf{(P\_8,4.14.2)} yatkriyāyuktāḥ tam prati gatyupasargasañjñe bhavataḥ na ca ṇopadeśam prati kriyāyogaḥ . \\
\noindent \textbf{(P\_8,4.14.2)} evam tarhi āha ayam upasargāt asamāse api ṇopadeśasya iti na ca ṇopadeśam prati upasargaḥ asti tatra vacanāt bhaviṣyati . \\
\noindent \textbf{(P\_8,4.14.2)} vacanaprāmāṇyāt iti cet padalope pratiṣedhaḥ . vacanaprāmāṇyāt iti cet padalope pratiṣedhaḥ vaktavyaḥ . \\
\noindent \textbf{(P\_8,4.14.2)} pragatāḥ nāyakāḥ asmāt grāmāt pranāyakaḥ grāmaḥ iti . \\
\noindent \textbf{(P\_8,4.14.2)} siddham tu yam prati upasargaḥ tatsthasya iti vacanāt . siddham etat . \\
\noindent \textbf{(P\_8,4.14.2)} katham . \\
\noindent \textbf{(P\_8,4.14.2)} yam prati upasargaḥ tatsthasya ṇaḥ bhavati iti vaktavyam . \\
\noindent \textbf{(P\_8,4.14.2)} sidhyati . \\
\noindent \textbf{(P\_8,4.14.2)} sūtram tarhi bhidyate . \\
\noindent \textbf{(P\_8,4.14.2)} yathānyāsam eva astu . \\
\noindent \textbf{(P\_8,4.14.2)} nanu ca uktam ṇopadeśam prati upasargābhāvāt anirdeśaḥ iti . \\
\noindent \textbf{(P\_8,4.14.2)} na eṣaḥ doṣaḥ . \\
\noindent \textbf{(P\_8,4.14.2)} ṇopadeśaḥ iti na evam vijñāyate ṇaḥ upadeśaḥ ṇopadeśaḥ ṇopadeśasya iti . \\
\noindent \textbf{(P\_8,4.14.2)} katham tarhi . \\
\noindent \textbf{(P\_8,4.14.2)} ṇaḥ upadeśaḥ asya saḥ ayam ṇopadeśaḥ ṇopadeśasya iti \\
\noindent \textbf{(P\_8,4.15)} hinumīnāgrahaṇe vikṛtasya upasaṅkhyānam . \\
\noindent \textbf{(P\_8,4.15)} hinumīnāgrahaṇe vikṛtasya upasaṅkhyānam kartavyam . \\
\noindent \textbf{(P\_8,4.15)} prahiṇoti pramīṇīte . \\
\noindent \textbf{(P\_8,4.15)} vacanāt bhaviṣyati . \\
\noindent \textbf{(P\_8,4.15)} asti vacane prayojanam . \\
\noindent \textbf{(P\_8,4.15)} kim . \\
\noindent \textbf{(P\_8,4.15)} prahiṇutaḥ pramīṇāti . \\
\noindent \textbf{(P\_8,4.15)} siddham acaḥ sthānivattvāt . siddham etat . \\
\noindent \textbf{(P\_8,4.15)} katham . \\
\noindent \textbf{(P\_8,4.15)} acaḥ sthānivattvāt . \\
\noindent \textbf{(P\_8,4.15)} sthānivadbhāvāt atra ṇatvam bhaviṣyati . \\
\noindent \textbf{(P\_8,4.15)} pratiṣidhyate atra sthānivadbhāvaḥ pūrvatrāseddhe na sthānivat iti . \\
\noindent \textbf{(P\_8,4.15)} doṣāḥ eva ete tasyāḥ paribhāṣāyāḥ tasya doṣaḥ saṃyogādilopalatvaṇatveṣu iti \\
\noindent \textbf{(P\_8,4.16)} loṭ iti kimartham . \\
\noindent \textbf{(P\_8,4.16)} prahimāni kulāni . \\
\noindent \textbf{(P\_8,4.16)} pravapāni māṃsāni . \\
\noindent \textbf{(P\_8,4.16)} āni loḍgrahaṇānarthakyam arthavadgrahaṇāt . \\
\noindent \textbf{(P\_8,4.16)} āni loḍgrahaṇam anarthakam . \\
\noindent \textbf{(P\_8,4.16)} kim kāraṇam . \\
\noindent \textbf{(P\_8,4.16)} arthavadgrahaṇāt . \\
\noindent \textbf{(P\_8,4.16)} arthavataḥ āniśabdasya grahaṇam na eṣaḥ arthavān . \\
\noindent \textbf{(P\_8,4.16)} anupasargāt vā . atha vā yatkriyāyuktāḥ tam prati gatyupasargasañjñe bhavataḥ na ca etam āniśabdam prati kriyāyogaḥ . \\
\noindent \textbf{(P\_8,4.16)} iha api tarhi na prāpnoti . \\
\noindent \textbf{(P\_8,4.16)} prayāṇi pariyāṇi iti . \\
\noindent \textbf{(P\_8,4.16)} atra api na āniśabdam prati kriyāyogaḥ . \\
\noindent \textbf{(P\_8,4.16)} āniśabdam prati atra kriyāyogaḥ . \\
\noindent \textbf{(P\_8,4.16)} katham . \\
\noindent \textbf{(P\_8,4.16)} yatkriyayuktāḥ iti na evam vijñāyate yasya kriyā yatkriyā yatkriyāyuktāḥ tam prati gatyupasargasañjñe bhavataḥ iti . \\
\noindent \textbf{(P\_8,4.16)} katham tarhi . \\
\noindent \textbf{(P\_8,4.16)} yā kriyā yatkriyā yatkriyāyuktāḥ tam prati gatyupasargasañjñe bhavataḥ iti \\
\noindent \textbf{(P\_8,4.17)} neḥ gadādiṣu aḍvyavāye upasaṅkhyānam . \\
\noindent \textbf{(P\_8,4.17)} neḥ gadādiṣu aḍvyavāye upasaṅkhyānam kartavyam . \\
\noindent \textbf{(P\_8,4.17)} praṇyagadat pariṇyagadat . \\
\noindent \textbf{(P\_8,4.17)} āṅā ca iti vaktavyam . \\
\noindent \textbf{(P\_8,4.17)} praṇyāgadat . \\
\noindent \textbf{(P\_8,4.17)} nanu ca ayam aṭ gadādibhaktaḥ gadādigrahaṇena grāhiṣyate . \\
\noindent \textbf{(P\_8,4.17)} na sidhyati . \\
\noindent \textbf{(P\_8,4.17)} aṅgasya aṭ ucyate vikaraṇāntam ca aṅgam saḥ asau saṅghātabhaktaḥ aśakyaḥ gadādigrahaṇena grahītum . \\
\noindent \textbf{(P\_8,4.17)} evam tarhi aḍvyavāye iti vartate . \\
\noindent \textbf{(P\_8,4.17)} kva prakṛtam . \\
\noindent \textbf{(P\_8,4.17)} aṭkupvāṅnumvyavāye api iti . \\
\noindent \textbf{(P\_8,4.17)} tat vai kāryiviśeṣaṇam nimittaviśeṣaṇena ca iha arthaḥ . \\
\noindent \textbf{(P\_8,4.17)} tatra api nimittaviśeṣaṇam eva \\
\noindent \textbf{(P\_8,4.19-20)} KA\_III,459.11-22 Ro\_V,500-501 {1/18}     antagrahaṇam kimartham . \\
\noindent \textbf{(P\_8,4.19-20)} KA\_III,459.11-22 Ro\_V,500-501 {2/18}     aniteḥ antagrahaṇam sambuddhyartham . \\
\noindent \textbf{(P\_8,4.19-20)} KA\_III,459.11-22 Ro\_V,500-501 {3/18}     aniteḥ antagrahaṇam kriyate sambuddhyartham . \\
\noindent \textbf{(P\_8,4.19-20)} KA\_III,459.11-22 Ro\_V,500-501 {4/18}     he prāṇ . \\
\noindent \textbf{(P\_8,4.19-20)} KA\_III,459.11-22 Ro\_V,500-501 {5/18}     apraraḥ āha . \\
\noindent \textbf{(P\_8,4.19-20)} KA\_III,459.11-22 Ro\_V,500-501 {6/18}     aniteḥ antaḥ padāntasya . \\
\noindent \textbf{(P\_8,4.19-20)} KA\_III,459.11-22 Ro\_V,500-501 {7/18}     aniteḥ antagrahaṇam kriyate padāntasya na iti pratiṣedhaḥ prāpnoti tadbādhanārtham . \\
\noindent \textbf{(P\_8,4.19-20)} KA\_III,459.11-22 Ro\_V,500-501 {8/18}     yaḥ vā tasmāt anantaraḥ . atha vā ayam antaśabdaḥ asti eva avayavavācī . \\
\noindent \textbf{(P\_8,4.19-20)} KA\_III,459.11-22 Ro\_V,500-501 {9/18}     tat yathā : vastrāntaḥ vasanāntaḥ iti . \\
\noindent \textbf{(P\_8,4.19-20)} KA\_III,459.11-22 Ro\_V,500-501 {10/18}     asti sāmīpye vartate . \\
\noindent \textbf{(P\_8,4.19-20)} KA\_III,459.11-22 Ro\_V,500-501 {11/18}     tat yathā . \\
\noindent \textbf{(P\_8,4.19-20)} KA\_III,459.11-22 Ro\_V,500-501 {12/18}     udakāntam gataḥ . \\
\noindent \textbf{(P\_8,4.19-20)} KA\_III,459.11-22 Ro\_V,500-501 {13/18}     udakasamīpam gataḥ iti gamyate . \\
\noindent \textbf{(P\_8,4.19-20)} KA\_III,459.11-22 Ro\_V,500-501 {14/18}     tat yaḥ sāmīpye vartate tasya grahaṇam vijñāyate . \\
\noindent \textbf{(P\_8,4.19-20)} KA\_III,459.11-22 Ro\_V,500-501 {15/18}     aniteḥ samīpe yaḥ rephaḥ tasmāt nasya yathā syāt . \\
\noindent \textbf{(P\_8,4.19-20)} KA\_III,459.11-22 Ro\_V,500-501 {16/18}     prāṇiti . \\
\noindent \textbf{(P\_8,4.19-20)} KA\_III,459.11-22 Ro\_V,500-501 {17/18}     iha mā bhūt . \\
\noindent \textbf{(P\_8,4.19-20)} KA\_III,459.11-22 Ro\_V,500-501 {18/18}     paryaniti \\
\noindent \textbf{(P\_8,4.21)} sābhyāsasya dvayoḥ iṣṭam . \\
\noindent \textbf{(P\_8,4.21)} sābhyāsasya dvayoḥ ṇatvam iṣyate . \\
\noindent \textbf{(P\_8,4.21)} prāṇiṇiṣati \\
\noindent \textbf{(P\_8,4.22)} atpūrvasya iti kimartham . \\
\noindent \textbf{(P\_8,4.22)} praghnanti parighnanti . \\
\noindent \textbf{(P\_8,4.22)} hanteḥ atpūrvasya vacane uktam . \\
\noindent \textbf{(P\_8,4.22)} kim uktam . \\
\noindent \textbf{(P\_8,4.22)} kuvyavāye hādeśeṣu pratiṣedhaḥ iti \\
\noindent \textbf{(P\_8,4.28)} katham idam vijñayate . \\
\noindent \textbf{(P\_8,4.28)} okārāt paraḥ otparaḥ na otparaḥ anotparaḥ iti . \\
\noindent \textbf{(P\_8,4.28)} āhosvit okāraḥ paraḥ asmāt saḥ ayam otparaḥ na otparaḥ anotparaḥ iti . \\
\noindent \textbf{(P\_8,4.28)} kim ca ataḥ . \\
\noindent \textbf{(P\_8,4.28)} yadi vijñayate okārāt paraḥ otparaḥ na otparaḥ anotparaḥ iti pra naḥ muñcatam atra api prāpnoti . \\
\noindent \textbf{(P\_8,4.28)} atha vijñāyate okāraḥ paraḥ asmāt saḥ ayam otparaḥ na otparaḥ anotparaḥ iti pra ṇaḥ vaniḥ devakṛtā atra na prāpnoti . \\
\noindent \textbf{(P\_8,4.28)} ubhayathā ca prakrame doṣaḥ bhavati . \\
\noindent \textbf{(P\_8,4.28)} pra naḥ* muñcatam , pra naḥ muñcatam . \\
\noindent \textbf{(P\_8,4.28)} pra* u naḥ , pra u naḥ . \\
\noindent \textbf{(P\_8,4.28)} bhāvini api oti na iṣyate . \\
\noindent \textbf{(P\_8,4.28)} bhāvini api okāre ṇatvam na iṣyate . \\
\noindent \textbf{(P\_8,4.28)} evam tarhi upasargāt bahulam iti vaktavyam \\
\noindent \textbf{(P\_8,4.29)} kṛtsthasya ṇatve nirviṇṇasya upasaṅkhyānam . \\
\noindent \textbf{(P\_8,4.29)} kṛtsthasya ṇatve nirviṇṇasya upasaṅkhyānam kartavyam . \\
\noindent \textbf{(P\_8,4.29)} nirviṇṇaḥ aham anena vāsena \\
\noindent \textbf{(P\_8,4.30)} ṇervibhāṣāyām sādhanavyavāye upasaṅkhyānam . \\
\noindent \textbf{(P\_8,4.30)} ṇervibhāṣāyām sādhanavyavāye upasaṅkhyānam kartavyam . \\
\noindent \textbf{(P\_8,4.30)} prāpyamāṇam prāpyamānam . \\
\noindent \textbf{(P\_8,4.30)} tadvidhānāt siddham . vihitaviśeṣaṇam ṇigrahaṇam . \\
\noindent \textbf{(P\_8,4.30)} ṇyantāt yaḥ vihitaḥ iti . \\
\noindent \textbf{(P\_8,4.30)} aḍadhikārāt vā . atha vā aḍvyavāye iti vartate \\
\noindent \textbf{(P\_8,4.32)} kimartham idam ucyate na kṛtyacaḥ iti eva siddham . \\
\noindent \textbf{(P\_8,4.32)} niyamārthaḥ ayam ārambhaḥ . \\
\noindent \textbf{(P\_8,4.32)} ijādeḥ eva ca sanumkāt na anyasmāt sanumkāt iti . \\
\noindent \textbf{(P\_8,4.32)} kva mā bhūt . \\
\noindent \textbf{(P\_8,4.32)} pramaṅkanam parimaṅkanam . \\
\noindent \textbf{(P\_8,4.32)} sanumaḥ ṇatve avadhāraṇāprasiddhiḥ vidheyabhāvāt . \\
\noindent \textbf{(P\_8,4.32)} sanumaḥ ṇatve avadhāraṇasya aprāptiḥ . \\
\noindent \textbf{(P\_8,4.32)} kim kāraṇam . \\
\noindent \textbf{(P\_8,4.32)} vidheyabhāvāt . \\
\noindent \textbf{(P\_8,4.32)} kaimarthakyāt niyamaḥ bhavati . \\
\noindent \textbf{(P\_8,4.32)} vidheyam na asti iti kṛtvā . \\
\noindent \textbf{(P\_8,4.32)} iha ca asti vidheyam . \\
\noindent \textbf{(P\_8,4.32)} kim . \\
\noindent \textbf{(P\_8,4.32)} ṇyantāt vibhāṣā prāptā tatra nityam ṇatvam vidheyam . \\
\noindent \textbf{(P\_8,4.32)} tatra apūrvaḥ vidhiḥ astu niyamaḥ astu iti apūrvaḥ vidhiḥ bhaviṣyati na niyamaḥ . \\
\noindent \textbf{(P\_8,4.32)} siddham tu pratiṣedhādhikāre sanumgrahaṇāt . siddham etat . \\
\noindent \textbf{(P\_8,4.32)} katham . \\
\noindent \textbf{(P\_8,4.32)} pratiṣedhādhikāre sanumgrahaṇāt . \\
\noindent \textbf{(P\_8,4.32)} pratiṣedhādhikāre sanumgrahaṇam kartavyam . \\
\noindent \textbf{(P\_8,4.32)} na bhābhūpūkamigamipyāyivepisanumām iti . \\
\noindent \textbf{(P\_8,4.32)} iha api tarhi na prāpnoti . \\
\noindent \textbf{(P\_8,4.32)} preṅgaṇam preṅgaṇīyam . \\
\noindent \textbf{(P\_8,4.32)} kṛtsthasya ca ṇatve ijādeḥ sanumaḥ grahaṇam . kṛtsthasya ca ṇatve ijādeḥ sanumaḥ grahaṇam kartavyam . \\
\noindent \textbf{(P\_8,4.32)} sidhyati sūtram tarhi bhidyate . \\
\noindent \textbf{(P\_8,4.32)} yathānyāsam eva astu . \\
\noindent \textbf{(P\_8,4.32)} nanu ca uktam sanumaḥ ṇatve avadhāraṇāprasiddhiḥ vidheyabhāvāt iti . \\
\noindent \textbf{(P\_8,4.32)} na eṣaḥ doṣaḥ . \\
\noindent \textbf{(P\_8,4.32)} halaḥ iti vartate . \\
\noindent \textbf{(P\_8,4.32)} kva prakṛtam . \\
\noindent \textbf{(P\_8,4.32)} halaścejupadhāt iti . \\
\noindent \textbf{(P\_8,4.32)} tat vai tatra ādiviśeṣaṇam antaviśeṣaṇena ca iha arthaḥ . \\
\noindent \textbf{(P\_8,4.32)} katham punaḥ jñāyate tatra ādiviśeṣaṇam iti . \\
\noindent \textbf{(P\_8,4.32)} ijupadhāt iti ucyate atra na arthaḥ antaviśeṣaṇena . \\
\noindent \textbf{(P\_8,4.32)} tatra ādiviśeṣaṇam sat iha antaviśeṣaṇam bhaviṣyati . \\
\noindent \textbf{(P\_8,4.32)} katham . \\
\noindent \textbf{(P\_8,4.32)} ijādeḥ iti ucyate tatra na arthaḥ ādiviśeṣaṇena . \\
\noindent \textbf{(P\_8,4.32)} atha vā ijādeḥ sanumaḥ iti atra ṇervibhāṣā iti etat anuvartiṣyate \\
\noindent \textbf{(P\_8,4.34)} bhādiṣu pūñgrahaṇam . \\
\noindent \textbf{(P\_8,4.34)} bhādiṣu pūñgrahaṇam kartavyam . \\
\noindent \textbf{(P\_8,4.34)} iha mā bhūt . \\
\noindent \textbf{(P\_8,4.34)} prapavaṇam somasya iti . \\
\noindent \textbf{(P\_8,4.34)} ṇyantasya ca upasaṅkhyanam . ṇyantasya ca upasaṅkhyānam kartavyam . \\
\noindent \textbf{(P\_8,4.34)} kim pūñaḥ eva . \\
\noindent \textbf{(P\_8,4.34)} na iti āha . \\
\noindent \textbf{(P\_8,4.34)} aviśeṣeṇa . \\
\noindent \textbf{(P\_8,4.34)} prabhāpanam paribhāpanam \\
\noindent \textbf{(P\_8,4.35)} ṣāt padādiparavacanam . \\
\noindent \textbf{(P\_8,4.35)} ṣātpadādiparagrahaṇam kartavyam . \\
\noindent \textbf{(P\_8,4.35)} iha eva yathā syāt . \\
\noindent \textbf{(P\_8,4.35)} niṣpānam duṣpānam . \\
\noindent \textbf{(P\_8,4.35)} iha mā bhūt . \\
\noindent \textbf{(P\_8,4.35)} sasarpiṣkeṇa sayajuṣkeṇa . \\
\noindent \textbf{(P\_8,4.35)} tat tarhi vaktavyam . \\
\noindent \textbf{(P\_8,4.35)} na vaktavyam . \\
\noindent \textbf{(P\_8,4.35)} na evam vijñāyate padasya antaḥ padāntaḥ padāntāt iti . \\
\noindent \textbf{(P\_8,4.35)} katham tarhi . \\
\noindent \textbf{(P\_8,4.35)} pade antaḥ padāntaḥ padāntāt iti \\
\noindent \textbf{(P\_8,4.36)} naśeḥ aśaḥ . \\
\noindent \textbf{(P\_8,4.36)} naśeḥ aśaḥ iti vaktavyam . \\
\noindent \textbf{(P\_8,4.36)} iha api yathā syāt . \\
\noindent \textbf{(P\_8,4.36)} pranaṅkṣyati parinaṅkṣyati . \\
\noindent \textbf{(P\_8,4.36)} tat tarhi vaktavyam . \\
\noindent \textbf{(P\_8,4.36)} na vaktavyam . \\
\noindent \textbf{(P\_8,4.36)} iha naśeḥ ṣaḥ iti iyatā siddham . \\
\noindent \textbf{(P\_8,4.36)} saḥ ayam evam siddhe sati yat antagrahaṇam karoti tasya etat prayojanam ṣāntabhūtapūrvasya api yathā syāt \\
\noindent \textbf{(P\_8,4.38)} padavyavāye ataddhite . \\
\noindent \textbf{(P\_8,4.38)} padavyavāye ataddhite iti vaktavyam . \\
\noindent \textbf{(P\_8,4.38)} iha mā bhūt . \\
\noindent \textbf{(P\_8,4.38)} ārdragomayeṇa śuṣkagomayeṇa iti . \\
\noindent \textbf{(P\_8,4.38)} tat tarhi vaktavyam . \\
\noindent \textbf{(P\_8,4.38)} na vaktavyam . \\
\noindent \textbf{(P\_8,4.38)} na evam vijñāyate padena vyavāye padavyavāye iti . \\
\noindent \textbf{(P\_8,4.38)} katham tarhi . \\
\noindent \textbf{(P\_8,4.38)} pade vyavāyaḥ padavyavāyaḥ padvyavāye iti \\
\noindent \textbf{(P\_8,4.39)} avihitalakṣaṇaḥ ṇatvapratiṣedhaḥ kṣubhnādiṣu draṣṭavyaḥ \\
\noindent \textbf{(P\_8,4.40)} kimartham tṛtīyānirdeśaḥ kriyate na ścau iti eva ucyeta . \\
\noindent \textbf{(P\_8,4.40)} ānantaryamātre ścutvam yathā syāt . \\
\noindent \textbf{(P\_8,4.40)} yajñaḥ rājñaḥ yācñā . \\
\noindent \textbf{(P\_8,4.40)} atha saṅkhyātānudeśaḥ kasmāt na bhavati . \\
\noindent \textbf{(P\_8,4.40)} ācāryapravṛttiḥ jñāpayati saṅkhyātānudeśaḥ na iha iti yat ayam śāt pratiṣedham śāsti \\
\noindent \textbf{(P\_8,4.41)} kimartham trtīyānirdeśaḥ kriyate na ṣṭau iti eva ucyeta . \\
\noindent \textbf{(P\_8,4.41)} ānantaryamātre ṣṭutvam yathā syāt . \\
\noindent \textbf{(P\_8,4.41)} peṣṭā leḍhā . \\
\noindent \textbf{(P\_8,4.41)} atha saṅkhyātānudeśaḥ kasmāt na bhavati . \\
\noindent \textbf{(P\_8,4.41)} ācāryapravṛttiḥ jñāpayati na iha saṅkhyātānudeśaḥ bhavati iti yat ayam toḥṣi iti pratiṣedham śāsti \\
\noindent \textbf{(P\_8,4.42)} anām iti kim . \\
\noindent \textbf{(P\_8,4.42)} ṣaṇṇām bhavati kaśyapaḥ . \\
\noindent \textbf{(P\_8,4.42)} atyalpam idam ucyate anām iti . \\
\noindent \textbf{(P\_8,4.42)} anānnavatinagarīṇām ca iti vaktavyam : ṣaṇṇām , ṣaṇṇavatiḥ , ṣaṇṇagarī \\
\noindent \textbf{(P\_8,4.45)} yaraḥ anunāsike pratyaye bhāṣāyām nityavacanam . \\
\noindent \textbf{(P\_8,4.45)} yaraḥ anunāsike pratyaye bhāṣāyām nityam iti ca vaktavyam . \\
\noindent \textbf{(P\_8,4.45)} vāṅmayam . \\
\noindent \textbf{(P\_8,4.45)} tvaṅmayam iti \\
\noindent \textbf{(P\_8,4.47)} dvirvacane yaṇaḥ mayaḥ . dvirvacane yaṇaḥ mayaḥ iti vaktavyam . \\
\noindent \textbf{(P\_8,4.47)} kim udāharaṇam . \\
\noindent \textbf{(P\_8,4.47)} yadi yaṇaḥ iti pañcamī mayaḥ iti ṣaṣṭhī ulkkā valmmīkam iti udaharaṇam . \\
\noindent \textbf{(P\_8,4.47)} atha mayaḥ iti pañcamī yaṇaḥ iti ṣaṣṭhī dadhyyatra madhvvatra iti udāharaṇam . \\
\noindent \textbf{(P\_8,4.47)} śaraḥ khayaḥ . \\
\noindent \textbf{(P\_8,4.47)} śaraḥ khayaḥ iti vaktavyam . \\
\noindent \textbf{(P\_8,4.47)} kim udāharaṇam . \\
\noindent \textbf{(P\_8,4.47)} yadi śaraḥ iti pañcamī khayaḥ iti ṣaṣṭhī sththālī sththātā iti udāharaṇam . \\
\noindent \textbf{(P\_8,4.47)} atha khayaḥ iti pañcamī śaraḥ iti ṣaṣṭhī vatssaḥ kṣṣīram apssarāḥ iti udāharaṇam . \\
\noindent \textbf{(P\_8,4.47)} avasāne ca . avasāne ca dve bhavataḥ iti vaktavyam . \\
\noindent \textbf{(P\_8,4.47)} vākk vāk . \\
\noindent \textbf{(P\_8,4.47)} tvakk tvak . \\
\noindent \textbf{(P\_8,4.47)} srukk sruk . \\
\noindent \textbf{(P\_8,4.47)} tat tarhi vaktavyam . \\
\noindent \textbf{(P\_8,4.47)} na vaktavyam . \\
\noindent \textbf{(P\_8,4.47)} na ayam prasajyapratiṣedhaḥ . \\
\noindent \textbf{(P\_8,4.47)} aci na iti . \\
\noindent \textbf{(P\_8,4.47)} kim tarhi paryudasaḥ ayam . \\
\noindent \textbf{(P\_8,4.47)} yat anyat acaḥ iti \\
\noindent \textbf{(P\_8,4.48)} na ādini ākrośe putrasya iti tatpare ca . \\
\noindent \textbf{(P\_8,4.48)} na ādini ākrośe putrasya iti atra tatpare ca iti vaktavyam . \\
\noindent \textbf{(P\_8,4.48)} putraputrādini \\
\noindent \textbf{(P\_8,4.48)} vā hatajagdhapare ca . \\
\noindent \textbf{(P\_8,4.48)} vā hatajagdhapare iti vaktavyam . \\
\noindent \textbf{(P\_8,4.48)} putrahatī puttrahatī . \\
\noindent \textbf{(P\_8,4.48)} putrajagdhī puttrajagdhī \\
\noindent \textbf{(P\_8,4.48)} cayaḥ dvitīyāḥ śari pauṣkarasādeḥ . \\
\noindent \textbf{(P\_8,4.48)} cayaḥ dvitīyāḥ bhavanti śari parataḥ pauṣkarasādeḥ ācāryasya matena . \\
\noindent \textbf{(P\_8,4.48)} vathsaḥ , khṣīram , aphsarāḥ \\
\noindent \textbf{(P\_8,4.61)} udaḥ pūrvatve skandeḥ chandasi upasaṅkhyānam . \\
\noindent \textbf{(P\_8,4.61)} udaḥ pūrvatve skandeḥ chandasi upasaṅkhyānam kartavyam . \\
\noindent \textbf{(P\_8,4.61)} aghnye dūram utkanda \\
\noindent \textbf{(P\_8,4.61)} roge ca iti vaktavyam . \\
\noindent \textbf{(P\_8,4.61)} utkandakaḥ rogaḥ \\
\noindent \textbf{(P\_8,4.63)} chatvam ami tacchlokena tacchmaśruṇā iti prayojanam . \\
\noindent \textbf{(P\_8,4.63)} chatvam ami iti vaktavyam . \\
\noindent \textbf{(P\_8,4.63)} kim prayojanam . \\
\noindent \textbf{(P\_8,4.63)} tacchlokena . \\
\noindent \textbf{(P\_8,4.63)} tacchmaśruṇā iti \\
\noindent \textbf{(P\_8,4.65)} savarṇagrahaṇam kimartham . \\
\noindent \textbf{(P\_8,4.65)} jharaḥ jhari savarṇagrahaṇam samasaṅkhyapratiṣedhārtham . \\
\noindent \textbf{(P\_8,4.65)} jharaḥ jhari savarṇagrahaṇam kriyate samasaṅkhyapratiṣedhārtham . \\
\noindent \textbf{(P\_8,4.65)} saṅkhyātānudeśaḥ mā bhūt iti . \\
\noindent \textbf{(P\_8,4.65)} kim ca syāt . \\
\noindent \textbf{(P\_8,4.65)} iha na syāt . \\
\noindent \textbf{(P\_8,4.65)} śiṇḍhi piṇḍhi iti \\
\noindent \textbf{(P\_8,4.68)} kimartham idam ucyate . \\
\noindent \textbf{(P\_8,4.68)} akāraḥ ayam akṣarasamāmnāye vivṛtaḥ upadiṣṭaḥ tasya saṃvṛtatāpratyāpattiḥ kriyate . \\
\noindent \textbf{(P\_8,4.68)} kim punaḥ kāraṇam vivṛtaḥ upadiśyate . \\
\noindent \textbf{(P\_8,4.68)} ādeśārtham savarṇārtham akāraḥ vivṛtaḥ smṛtaḥ . \\
\noindent \textbf{(P\_8,4.68)} ākārasya tathā hrasvaḥ tadartham pāṇineḥ a a . \\
\noindent \textbf{(P\_8,4.68)} ādeśārtham tāvat . \\
\noindent \textbf{(P\_8,4.68)} vṛkṣābhyām , devadattā3 . \\
\noindent \textbf{(P\_8,4.68)} āntaryataḥ vivṛtasya vivṛtau dīrghaplutau yathā syātām . \\
\noindent \textbf{(P\_8,4.68)} savarṇārtham ca . \\
\noindent \textbf{(P\_8,4.68)} akāraḥ savarṇagrahaṇena ākāram api yathā gṛhṇīyāt . \\
\noindent \textbf{(P\_8,4.68)} ākārasya tathā hrasvaḥ . \\
\noindent \textbf{(P\_8,4.68)} tathā ca atikhaṭvaḥ , atimālaḥ iti atra ākārasya hrasvaḥ ucyamānaḥ vivṛtaḥ prāpnoti saḥ saṃvṛtaḥ syāt iti evamarthā prayāpattiḥ . \\
\noindent \textbf{(P\_8,4.68)} asti prayojanam etat . \\
\noindent \textbf{(P\_8,4.68)} kim tarhi iti . \\
\noindent \textbf{(P\_8,4.68)} akārasya pratyāpattau dīrghapratiṣedhaḥ . akārasya pratyāpattau dīrghasya pratiṣedhaḥ vaktavyaḥ . \\
\noindent \textbf{(P\_8,4.68)} khaṭvā mālā . \\
\noindent \textbf{(P\_8,4.68)} na eṣaḥ doṣaḥ . \\
\noindent \textbf{(P\_8,4.68)} yathā eva prakṛtitaḥ savarṇagrahaṇam evam ādeśataḥ api bhavitavyam tatra āntaryataḥ hrasvasya hrasvaḥ dīrghasya dīrghaḥ bhaviṣyati . \\
\noindent \textbf{(P\_8,4.68)} ādeśasya ca anaṇtvāt na savarṇagrahaṇam . ādeśasya ca anaṇtvāt savarṇānām grahaṇam na prāpnoti . \\
\noindent \textbf{(P\_8,4.68)} keṣām . \\
\noindent \textbf{(P\_8,4.68)} udāttānudāttasvaritānunāsikānām . \\
\noindent \textbf{(P\_8,4.68)} siddham tu taparanirdeśāt . siddham etat . \\
\noindent \textbf{(P\_8,4.68)} katham . \\
\noindent \textbf{(P\_8,4.68)} taparanirdeśāt . \\
\noindent \textbf{(P\_8,4.68)} taparanirdeśaḥ kartavyaḥ . \\
\noindent \textbf{(P\_8,4.68)} at a* iti . \\
\noindent \textbf{(P\_8,4.68)} aparaḥ āha : akārasya pratyāpattau dīrghapratiṣedhaḥ . \\
\noindent \textbf{(P\_8,4.68)} akārasya pratyāpattau dīrghasya pratiṣedhaḥ vaktavyaḥ : khaṭvā mālā . \\
\noindent \textbf{(P\_8,4.68)} na eṣaḥ doṣaḥ . \\
\noindent \textbf{(P\_8,4.68)} dīrghoccāraṇasāmarthyāt na bhaviṣyati . \\
\noindent \textbf{(P\_8,4.68)} idam tarhi prayojanam . \\
\noindent \textbf{(P\_8,4.68)} vṛkṣābhyām plakṣābhyām . \\
\noindent \textbf{(P\_8,4.68)} atra api dīrghavacanasāmarthyāt na bhaviṣyati . \\
\noindent \textbf{(P\_8,4.68)} idam tarhi . \\
\noindent \textbf{(P\_8,4.68)} api kākaḥ śyenāyate . \\
\noindent \textbf{(P\_8,4.68)} nanu ca atra api dīrghavacanasamarthyāt eva na bhaviṣyati . \\
\noindent \textbf{(P\_8,4.68)} asti anyat dīrghavacane prayojanam . \\
\noindent \textbf{(P\_8,4.68)} kim . \\
\noindent \textbf{(P\_8,4.68)} dadhīyati madhūyati . \\
\noindent \textbf{(P\_8,4.68)} atra eva ca eṣaḥ doṣaḥ ādeśasya ca anaṇtvāt savarṇānām grahaṇam na prāpnoti . \\
\noindent \textbf{(P\_8,4.68)} keṣām . \\
\noindent \textbf{(P\_8,4.68)} udāttānudāttasvaritānunāsikānām . \\
\noindent \textbf{(P\_8,4.68)} siddham tu taparanirdeśāt . \\
\noindent \textbf{(P\_8,4.68)} siddham etat . \\
\noindent \textbf{(P\_8,4.68)} katham . \\
\noindent \textbf{(P\_8,4.68)} taparanirdeśāt . \\
\noindent \textbf{(P\_8,4.68)} taparanirdeśaḥ kartavyaḥ . \\
\noindent \textbf{(P\_8,4.68)} at at iti . \\
\noindent \textbf{(P\_8,4.68)} ekaśeṣanirdeśāt vā svarabhinnānām bhagavataḥ pāṇineḥ siddham . ekaśeṣanirdeśāt vā svarabhinnānam bhagavataḥ pāṇineḥ ācāryasya siddham . \\
\noindent \textbf{(P\_8,4.68)} ekaśeṣanirdesaḥ ayam . \\
\noindent \textbf{(P\_8,4.68)} a , a , a iti \\
\end{document}